\documentclass[11pt,a4paper]{article}
\usepackage[margin=25mm,headheight=15pt]{geometry}
\usepackage[T1]{fontenc}
\usepackage[utf8]{inputenc}
\usepackage{lmodern}
\usepackage{amsmath,amssymb,amsthm,amscd,mathtools,mathrsfs}
\usepackage{microtype,booktabs,longtable,array,xcolor,graphicx,enumitem,needspace,fancyhdr}
\usepackage[round,authoryear]{natbib}
\usepackage[unicode,colorlinks=true,linkcolor=blue!40!black,citecolor=blue!40!black,urlcolor=blue!40!black]{hyperref}
\usepackage{xurl}
\usepackage{listings}
\usepackage{tabularx}
\lstset{basicstyle=\ttfamily\small,breaklines=true,columns=fullflexible,keepspaces=true}
\usepackage{tikz}
\usepackage{tikz-cd}
\usetikzlibrary{arrows.meta,calc,patterns,positioning}
\usepackage{aliascnt}
\usepackage[nameinlink,noabbrev]{cleveref}
\allowdisplaybreaks[2]
\newcommand{\SsixRepositoryCommit}{83d76022bcd442278b760db9ff438afceed6c6ea}
\newcommand{\SsixPublishRecords}{https://github.com/KokunoYumeto/s6-complex-structure-workbench/tree/1fb9194767aef92abcc06071e3820c2f226b9512/pro_returns/S6_SERIES_PUBLISH_checkpoint_10_20261121T0612Z/}
\setcounter{MaxMatrixCols}{32}
\setlength{\emergencystretch}{3em}
\setlength{\parskip}{2pt}
\newtheorem{theorem}{Theorem}[section]
\newaliascnt{proposition}{theorem}
\newtheorem{proposition}[proposition]{Proposition}
\aliascntresetthe{proposition}
\crefname{proposition}{proposition}{propositions}
\newaliascnt{lemma}{theorem}
\newtheorem{lemma}[lemma]{Lemma}
\aliascntresetthe{lemma}
\crefname{lemma}{lemma}{lemmas}
\newaliascnt{corollary}{theorem}
\newtheorem{corollary}[corollary]{Corollary}
\aliascntresetthe{corollary}
\crefname{corollary}{corollary}{corollaries}
\newaliascnt{claim}{theorem}
\newtheorem{claim}[claim]{Claim}
\aliascntresetthe{claim}
\crefname{claim}{claim}{claims}
\newaliascnt{statement}{theorem}
\newtheorem{statement}[statement]{Statement}
\aliascntresetthe{statement}
\crefname{statement}{statement}{statements}
\theoremstyle{definition}
\newaliascnt{definition}{theorem}
\newtheorem{definition}[definition]{Definition}
\aliascntresetthe{definition}
\crefname{definition}{definition}{definitions}
\newaliascnt{remark}{theorem}
\newtheorem{remark}[remark]{Remark}
\aliascntresetthe{remark}
\crefname{remark}{remark}{remarks}
\newaliascnt{example}{theorem}
\newtheorem{example}[example]{Example}
\aliascntresetthe{example}
\crefname{example}{example}{examples}
\newaliascnt{problem}{theorem}
\newtheorem{problem}[problem]{Open problem}
\aliascntresetthe{problem}
\crefname{problem}{open problem}{open problems}
\newaliascnt{question}{theorem}
\newtheorem{question}[question]{Question}
\aliascntresetthe{question}
\crefname{question}{question}{questions}
\newaliascnt{auditrecord}{theorem}
\newtheorem{auditrecord}[auditrecord]{Verification record}
\aliascntresetthe{auditrecord}
\crefname{auditrecord}{verification record}{verification records}
\newaliascnt{test}{theorem}
\newtheorem{test}[test]{Verification}
\aliascntresetthe{test}
\crefname{test}{verification}{verifications}
\newaliascnt{conjecture}{theorem}
\newtheorem{conjecture}[conjecture]{Conjecture}
\aliascntresetthe{conjecture}
\crefname{conjecture}{conjecture}{conjectures}
\newcommand{\C}{\mathbb C}\newcommand{\R}{\mathbb R}\newcommand{\Z}{\mathbb Z}\newcommand{\Q}{\mathbb Q}\newcommand{\N}{\mathbb N}\newcommand{\F}{\mathbb F}
\newcommand{\CC}{\mathbb C}\newcommand{\RR}{\mathbb R}\newcommand{\ZZ}{\mathbb Z}\newcommand{\QQ}{\mathbb Q}\newcommand{\NN}{\mathbb N}
\newcommand{\PP}{\mathbb P}\newcommand{\Pp}{\mathbb P}\newcommand{\Pone}{\mathbb P^1}\newcommand{\HH}{\mathbb H}
\newcommand{\RP}{\mathbb{RP}}\newcommand{\CP}{\mathbb{CP}}
\newcommand{\OO}{\mathcal O}\newcommand{\cO}{\mathcal O}\newcommand{\cM}{\mathcal M}\newcommand{\cF}{\mathcal F}
\newcommand{\h}{\mathfrak h}\newcommand{\hz}{\mathfrak h_z}\newcommand{\Acon}{\mathfrak A}\newcommand{\wt}[1]{\widehat{#1}}
\newcommand{\red}{\mathrm{red}}\newcommand{\tor}{\mathrm{tor}}\newcommand{\dP}{\mathrm{dP}}
\newcommand{\gh}{\widehat\gamma}\newcommand{\uh}{\widehat u}\newcommand{\wh}{\widehat w}\newcommand{\dhat}{\widehat\delta}\renewcommand{\dh}{\widehat\delta}
\newcommand{\dd}{\mathrm d}\newcommand{\ii}{\mathrm i}\newcommand{\eps}{\varepsilon}\newcommand{\partt}{\mathrm{part}}\newcommand{\Bcirc}{B^\circ}\newcommand{\ddbar}{\partial\bar\partial}
\DeclareMathOperator{\Pic}{Pic}\DeclareMathOperator{\Alb}{Alb}\DeclareMathOperator{\Aut}{Aut}\DeclareMathOperator{\NS}{NS}
\DeclareMathOperator{\Hom}{Hom}\DeclareMathOperator{\End}{End}\DeclareMathOperator{\Ext}{Ext}\DeclareMathOperator{\Tor}{Tor}
\DeclareMathOperator{\Spec}{Spec}\DeclareMathOperator{\Proj}{Proj}\DeclareMathOperator{\Sym}{Sym}\DeclareMathOperator{\Diff}{Diff}
\DeclareMathOperator{\GL}{GL}\DeclareMathOperator{\SL}{SL}\DeclareMathOperator{\PSL}{PSL}\DeclareMathOperator{\Sp}{Sp}
\DeclareMathOperator{\im}{im}\DeclareMathOperator{\coker}{coker}\DeclareMathOperator{\rank}{rank}\DeclareMathOperator{\diag}{diag}
\DeclareMathOperator{\ord}{ord}\DeclareMathOperator{\vol}{vol}\DeclareMathOperator{\pr}{pr}\DeclareMathOperator{\id}{id}
\DeclareMathOperator{\divisor}{div}\DeclareMathOperator{\Tr}{Tr}\DeclareMathOperator{\Nm}{Nm}\DeclareMathOperator{\Sha}{Sha}
\DeclareMathOperator{\Imn}{Im}\DeclareMathOperator{\Ren}{Re}\DeclareMathOperator{\trdeg}{tr.deg}
\newcommand{\Uses}[1]{\par\smallskip\noindent\textit{Uses:} #1.\par\smallskip}
\newcommand{\Verified}{}\newcommand{\AuthorSelfCheck}{}\newcommand{\NeedsCheck}{}\newcommand{\Open}{\textbf{Open.}\ }

\newcommand{\RungZero}[1]{\par\noindent\textbf{Result and mechanism.} #1\par\smallskip}
\newcommand{\RungOne}{\par\Needspace{6\baselineskip}\noindent\textbf{Proof at the first rung.}\par\nobreak\smallskip}
\newcommand{\CurrentChecker}{GPT-6 Astra (OpenAI, ChatGPT Pro mode, effort Pro), 18 November 2026}
\newaliascnt{cited}{theorem}
\newtheorem{cited}[cited]{Cited result}
\aliascntresetthe{cited}
\crefname{cited}{cited result}{cited results}
\newaliascnt{refuted}{theorem}
\newtheorem{refuted}[refuted]{Correction}
\aliascntresetthe{refuted}
\crefname{refuted}{correction}{corrections}
\newaliascnt{computed}{theorem}
\newtheorem{computed}[computed]{Computation}
\aliascntresetthe{computed}
\crefname{computed}{computation}{computations}
\newaliascnt{openq}{theorem}
\newtheorem{openq}[openq]{Open question}
\aliascntresetthe{openq}
\crefname{openq}{open question}{open questions}

\providecommand{\rk}{\operatorname{rk}}
\providecommand{\Ann}{\operatorname{Ann}}
\providecommand{\Log}{\operatorname{Log}}
\providecommand{\vir}{\mathrm{vir}}

\usepackage{etoolbox}
\hypersetup{bookmarksnumbered=true,bookmarksdepth=2}
\setcounter{tocdepth}{2}
\makeatletter\renewcommand{\@pnumwidth}{2.5em}\renewcommand{\@tocrmarg}{3.5em}\renewcommand*{\l@subsection}{\@dottedtocline{2}{1.5em}{3.5em}}\makeatother
\newtheorem{maintheorem}{Theorem}
\renewcommand{\themaintheorem}{\Alph{maintheorem}}
\crefname{maintheorem}{theorem}{theorems}
\newcommand{\StatementRecord}[4]{{\upshape\footnotesize[tag~\texttt{#3}]}\enspace\ignorespaces}
\newcommand{\ExternalPaperRef}[3]{#1#2 (Paper~#3)}
\newcommand{\CompendiumRef}[2]{\ifmmode\text{\upshape S6\_09, Paper~#1, label~\nolinkurl{#2}}\else S6\_09, Paper~#1, label~\nolinkurl{#2}\fi}
\newcommand{\SeriesRefText}[1]{\ifcsname SRef@#1\endcsname\csname SRef@#1\endcsname\else\textbf{UNRESOLVED}\fi}
\newcommand{\SeriesRef}[1]{\ifmmode\text{\SeriesRefText{#1}}\else\SeriesRefText{#1}\fi}


\newcommand{\CurrentPaper}{II}
\expandafter\def\csname SRef@wb:pshift:eq:periods\endcsname{\href{S6_SERIES_GRIND_I.pdf\#nameddest=tag-03CN}{I.{PS1} (p.~349)} }
\expandafter\def\csname SRef@wb:pshift:eq:rank-one-map\endcsname{\href{S6_SERIES_GRIND_I.pdf\#nameddest=tag-03CO}{I.{PS2} (p.~349)} }
\expandafter\def\csname SRef@wb:pshift:eq:finite-data\endcsname{\href{S6_SERIES_GRIND_I.pdf\#nameddest=tag-03C8}{I.{PS3} (p.~349)} }
\expandafter\def\csname SRef@wb:pshift:eq:covariance\endcsname{\href{S6_SERIES_GRIND_I.pdf\#nameddest=tag-03C2}{I.{PS4} (p.~349)} }
\expandafter\def\csname SRef@wb:pshift:eq:integer-shear\endcsname{\href{S6_SERIES_GRIND_I.pdf\#nameddest=tag-03CF}{I.{PS6} (p.~349)} }
\expandafter\def\csname SRef@wb:pshift:eq:centralizer\endcsname{\href{S6_SERIES_GRIND_I.pdf\#nameddest=tag-03BZ}{I.{PS7} (p.~349)} }
\expandafter\def\csname SRef@wb:pshift:eq:period-morphism\endcsname{\href{S6_SERIES_GRIND_I.pdf\#nameddest=tag-03CM}{I.{PS8} (p.~350)} }
\expandafter\def\csname SRef@wb:pshift:eq:norms\endcsname{\href{S6_SERIES_GRIND_I.pdf\#nameddest=tag-03CK}{I.{PS9} (p.~351)} }
\expandafter\def\csname SRef@wb:pshift:eq:finite-equation\endcsname{\href{S6_SERIES_GRIND_I.pdf\#nameddest=tag-03C9}{I.{PS10} (p.~351)} }
\expandafter\def\csname SRef@wb:pshift:eq:finite-norm-obstruction\endcsname{\href{S6_SERIES_GRIND_I.pdf\#nameddest=tag-03CB}{I.{PS11} (p.~351)} }
\expandafter\def\csname SRef@wb:pshift:eq:finite-class\endcsname{\href{S6_SERIES_GRIND_I.pdf\#nameddest=tag-03C7}{I.{PS12} (p.~351)} }
\expandafter\def\csname SRef@wb:pshift:eq:local-corrections\endcsname{\href{S6_SERIES_GRIND_I.pdf\#nameddest=tag-03CH}{I.{PS13} (p.~352)} }
\expandafter\def\csname SRef@wb:pshift:eq:local-matrix-identity\endcsname{\href{S6_SERIES_GRIND_I.pdf\#nameddest=tag-03CI}{I.{PS14} (p.~352)} }
\expandafter\def\csname SRef@wb:pshift:eq:finite-middle-sections\endcsname{\href{S6_SERIES_GRIND_I.pdf\#nameddest=tag-03CA}{I.{PS15} (p.~352)} }
\expandafter\def\csname SRef@wb:pshift:eq:local-affine-cocycle\endcsname{\href{S6_SERIES_GRIND_I.pdf\#nameddest=tag-03CG}{I.{PS16} (p.~352)} }
\expandafter\def\csname SRef@wb:pshift:eq:actual-overlap-difference\endcsname{\href{S6_SERIES_GRIND_I.pdf\#nameddest=tag-03BX}{I.{PS17} (p.~352)} }
\expandafter\def\csname SRef@wb:pshift:eq:cusp-periods\endcsname{\href{S6_SERIES_GRIND_I.pdf\#nameddest=tag-03C5}{I.{PS18} (p.~352)} }
\expandafter\def\csname SRef@wb:pshift:eq:cusp-actions-identical\endcsname{\href{S6_SERIES_GRIND_I.pdf\#nameddest=tag-03C3}{I.{PS19} (p.~352)} }
\expandafter\def\csname SRef@wb:pshift:eq:cocycle-law\endcsname{\href{S6_SERIES_GRIND_I.pdf\#nameddest=tag-03C0}{I.{PS20} (p.~352)} }
\expandafter\def\csname SRef@wb:pshift:eq:cusp-cocycle\endcsname{\href{S6_SERIES_GRIND_I.pdf\#nameddest=tag-03C4}{I.{PS21} (p.~352)} }
\expandafter\def\csname SRef@wb:pshift:eq:cusp-translation\endcsname{\href{S6_SERIES_GRIND_I.pdf\#nameddest=tag-03C6}{I.{PS22} (p.~353)} }
\expandafter\def\csname SRef@wb:pshift:eq:affine-bundle-action\endcsname{\href{S6_SERIES_GRIND_I.pdf\#nameddest=tag-03BY}{I.{PS23} (p.~353)} }
\expandafter\def\csname SRef@wb:pshift:eq:vector-sheaf\endcsname{\href{S6_SERIES_GRIND_I.pdf\#nameddest=tag-03CP}{I.{PS24} (p.~353)} }
\expandafter\def\csname SRef@wb:pshift:eq:no-analytic-obstruction\endcsname{\href{S6_SERIES_GRIND_I.pdf\#nameddest=tag-03CJ}{I.{PS25} (p.~354)} }
\expandafter\def\csname SRef@wb:pshift:eq:global-correction\endcsname{\href{S6_SERIES_GRIND_I.pdf\#nameddest=tag-03CE}{I.{PS26} (p.~354)} }
\expandafter\def\csname SRef@wb:pshift:eq:full-isomorphism\endcsname{\href{S6_SERIES_GRIND_I.pdf\#nameddest=tag-03CD}{I.{PS27} (p.~354)} }
\expandafter\def\csname SRef@wb:pshift:eq:finite-overlap-identity\endcsname{\href{S6_SERIES_GRIND_I.pdf\#nameddest=tag-03CC}{I.{PS28} (p.~354)} }
\expandafter\def\csname SRef@wb:pshift:eq:coherent-shift-action\endcsname{\href{S6_SERIES_GRIND_I.pdf\#nameddest=tag-03C1}{I.{PS28a} (p.~355)} }
\expandafter\def\csname SRef@wb:pshift:eq:parameter-disc\endcsname{\href{S6_SERIES_GRIND_I.pdf\#nameddest=tag-03CL}{I.{PS29} (p.~355)} }
\expandafter\def\csname SRef@mapsI:quotient-eq-2\endcsname{\href{S6_SERIES_GRIND_I.pdf\#nameddest=tag-02OZ}{I.22 (p.~296)} }
\expandafter\def\csname SRef@mapsI:quotient-eq-3\endcsname{\href{S6_SERIES_GRIND_I.pdf\#nameddest=tag-02P0}{I.{{$\langle 1\rangle $1.}} (p.~297)} }
\expandafter\def\csname SRef@mapsI:quotient-eq-4\endcsname{\href{S6_SERIES_GRIND_I.pdf\#nameddest=tag-02P1}{I.{{$\langle 1\rangle $2.}} (p.~297)} }
\expandafter\def\csname SRef@mapsI:quotient-eq-5\endcsname{\href{S6_SERIES_GRIND_I.pdf\#nameddest=tag-02P2}{I.22 (p.~297)} }
\expandafter\def\csname SRef@mapsI:quotient-eq-7\endcsname{\href{S6_SERIES_GRIND_I.pdf\#nameddest=tag-02P3}{I.{{$\langle 1\rangle $1.}} (p.~298)} }
\expandafter\def\csname SRef@app:rederive:fam-src-D.1.2\endcsname{\href{S6_SERIES_GRIND_I.pdf\#nameddest=tag-000F}{I.{{$\langle 1\rangle $1.}} (p.~77)} }
\expandafter\def\csname SRef@app:rederive:fam-src-D.1.6\endcsname{\href{S6_SERIES_GRIND_I.pdf\#nameddest=tag-000G}{I.{{$\langle 1\rangle $1.}} (p.~80)} }
\expandafter\def\csname SRef@app:rederive:fam-src-D.2.3\endcsname{\href{S6_SERIES_GRIND_I.pdf\#nameddest=tag-000H}{I.{{$\langle 1\rangle $1.}} (p.~83)} }
\expandafter\def\csname SRef@mapsI:real-eq-1\endcsname{\href{S6_SERIES_GRIND_I.pdf\#nameddest=tag-02P4}{I.12 (p.~89)} }
\expandafter\def\csname SRef@mapsI:real-eq-2\endcsname{\href{S6_SERIES_GRIND_I.pdf\#nameddest=tag-02P4}{I.12 (p.~89)} }
\expandafter\def\csname SRef@mapsI:real-eq-4\endcsname{\href{S6_SERIES_GRIND_I.pdf\#nameddest=tag-02P5}{I.{{$\langle 1\rangle $1.}} (p.~89)} }
\expandafter\def\csname SRef@mapsI:real-eq-5\endcsname{\href{S6_SERIES_GRIND_I.pdf\#nameddest=tag-02P6}{I.{{$\langle 1\rangle $4.}} (p.~90)} }
\expandafter\def\csname SRef@mapsI:real-eq-6\endcsname{\href{S6_SERIES_GRIND_I.pdf\#nameddest=tag-02P6}{I.{{$\langle 1\rangle $4.}} (p.~90)} }
\expandafter\def\csname SRef@mapsI:real-eq-7\endcsname{\href{S6_SERIES_GRIND_I.pdf\#nameddest=tag-02P7}{I.{{$\langle 1\rangle $8.}} (p.~91)} }
\expandafter\def\csname SRef@mapsI:real-eq-8\endcsname{\href{S6_SERIES_GRIND_I.pdf\#nameddest=tag-02P8}{I.{{$\langle 1\rangle $9.}} (p.~91)} }
\expandafter\def\csname SRef@ari:open:norm\endcsname{\href{S6_SERIES_GRIND_I.pdf\#nameddest=tag-003T}{I.F (p.~406)} }
\expandafter\def\csname SRef@push:src:C.3.5\endcsname{\href{S6_SERIES_GRIND_I.pdf\#nameddest=tag-0322}{I.12.5 (p.~93)} }
\expandafter\def\csname SRef@push:src:B.4.3\endcsname{\href{S6_SERIES_GRIND_I.pdf\#nameddest=tag-030Z}{I.15.16 (p.~144)} }
\expandafter\def\csname SRef@push:src:D.4.6\endcsname{\href{S6_SERIES_GRIND_I.pdf\#nameddest=tag-032O}{I.16.15 (p.~241)} }
\expandafter\def\csname SRef@inv:A1.0\endcsname{\href{S6_SERIES_GRIND_I.pdf\#nameddest=tag-01RQ}{I.18.1 (p.~254)} }
\expandafter\def\csname SRef@wb:lem:cusp-structure-sheaf-audit\endcsname{\href{S6_SERIES_GRIND_I.pdf\#nameddest=tag-01RQ}{I.18.1 (p.~254)} }
\expandafter\def\csname SRef@wb:thm:structure-direct-images-audit\endcsname{\href{S6_SERIES_GRIND_I.pdf\#nameddest=tag-01RQ}{I.18.1 (p.~254)} }
\expandafter\def\csname SRef@inv:C3.6\endcsname{\href{S6_SERIES_GRIND_I.pdf\#nameddest=tag-01SJ}{I.19.10 (p.~273)} }
\expandafter\def\csname SRef@push:src:E.3.1\endcsname{\href{S6_SERIES_GRIND_I.pdf\#nameddest=tag-033I}{I.20.9 (p.~285)} }
\expandafter\def\csname SRef@ari:C.5.1\endcsname{\href{S6_SERIES_GRIND_I.pdf\#nameddest=tag-000U}{I.22.9 (p.~304)} }
\expandafter\def\csname SRef@ari:C.6.R4\endcsname{\href{S6_SERIES_GRIND_I.pdf\#nameddest=tag-0013}{I.23.9 (p.~315)} }
\expandafter\def\csname SRef@ari:C.4.1\endcsname{\href{S6_SERIES_GRIND_I.pdf\#nameddest=tag-000P}{I.5.4 (p.~38)} }
\expandafter\def\csname SRef@ari:C.7.7\endcsname{\href{S6_SERIES_GRIND_I.pdf\#nameddest=tag-001G}{I.E.6 (p.~388)} }
\expandafter\def\csname SRef@ari:C.7.8\endcsname{\href{S6_SERIES_GRIND_I.pdf\#nameddest=tag-001H}{I.E.7 (p.~390)} }
\expandafter\def\csname SRef@push:src:A.4.5\endcsname{\href{S6_SERIES_GRIND_I.pdf\#nameddest=tag-02ZV}{I.14.14 (p.~120)} }
\expandafter\def\csname SRef@push:A45\endcsname{\href{S6_SERIES_GRIND_I.pdf\#nameddest=tag-02ZV}{I.14.14 (p.~120)} }
\expandafter\def\csname SRef@push:src:B.3.5\endcsname{\href{S6_SERIES_GRIND_I.pdf\#nameddest=tag-0309}{I.15.10 (p.~138)} }
\expandafter\def\csname SRef@push:B35\endcsname{\href{S6_SERIES_GRIND_I.pdf\#nameddest=tag-0309}{I.15.10 (p.~138)} }
\expandafter\def\csname SRef@push:src:B.4.1\endcsname{\href{S6_SERIES_GRIND_I.pdf\#nameddest=tag-030D}{I.15.14 (p.~143)} }
\expandafter\def\csname SRef@push:B41\endcsname{\href{S6_SERIES_GRIND_I.pdf\#nameddest=tag-030D}{I.15.14 (p.~143)} }
\expandafter\def\csname SRef@push:src:B.2.2\endcsname{\href{S6_SERIES_GRIND_I.pdf\#nameddest=tag-0303}{I.15.4 (p.~132)} }
\expandafter\def\csname SRef@push:B22\endcsname{\href{S6_SERIES_GRIND_I.pdf\#nameddest=tag-0303}{I.15.4 (p.~132)} }
\expandafter\def\csname SRef@push:src:B.5.1\endcsname{\href{S6_SERIES_GRIND_I.pdf\#nameddest=tag-031A}{I.15.47 (p.~227)} }
\expandafter\def\csname SRef@push:B51\endcsname{\href{S6_SERIES_GRIND_I.pdf\#nameddest=tag-031A}{I.15.47 (p.~227)} }
\expandafter\def\csname SRef@push:src:B.2.3\endcsname{\href{S6_SERIES_GRIND_I.pdf\#nameddest=tag-0304}{I.15.5 (p.~132)} }
\expandafter\def\csname SRef@push:B23\endcsname{\href{S6_SERIES_GRIND_I.pdf\#nameddest=tag-0304}{I.15.5 (p.~132)} }
\expandafter\def\csname SRef@push:src:D.3.2\endcsname{\href{S6_SERIES_GRIND_I.pdf\#nameddest=tag-032I}{I.16.9 (p.~235)} }
\expandafter\def\csname SRef@push:src:E.1.3\endcsname{\href{S6_SERIES_GRIND_I.pdf\#nameddest=tag-032T}{I.9.3 (p.~67)} }
\expandafter\def\csname SRef@fam:eq:cd-beta\endcsname{\href{S6_SERIES_GRIND_I.pdf\#nameddest=tag-01P5}{I.1 (p.~29)} }
\expandafter\def\csname SRef@ari:eq:ca-GHbounds\endcsname{\href{S6_SERIES_GRIND_I.pdf\#nameddest=tag-0024}{I.10 (p.~44)} }
\expandafter\def\csname SRef@eq:P5L\endcsname{\href{S6_SERIES_GRIND_I.pdf\#nameddest=tag-01OG}{I.100 (p.~367)} }
\expandafter\def\csname SRef@eq:P5floor\endcsname{\href{S6_SERIES_GRIND_I.pdf\#nameddest=tag-01OI}{I.101 (p.~367)} }
\expandafter\def\csname SRef@eq:P5normalization\endcsname{\href{S6_SERIES_GRIND_I.pdf\#nameddest=tag-01OJ}{I.102 (p.~368)} }
\expandafter\def\csname SRef@eq:P5exponential\endcsname{\href{S6_SERIES_GRIND_I.pdf\#nameddest=tag-01OH}{I.103 (p.~371)} }
\expandafter\def\csname SRef@eq:P5H2\endcsname{\href{S6_SERIES_GRIND_I.pdf\#nameddest=tag-01OF}{I.104 (p.~372)} }
\expandafter\def\csname SRef@ari:eq:cp-sigma\endcsname{\href{S6_SERIES_GRIND_I.pdf\#nameddest=tag-003K}{I.119 (p.~383)} }
\expandafter\def\csname SRef@ari:eq:cp-klein\endcsname{\href{S6_SERIES_GRIND_I.pdf\#nameddest=tag-003F}{I.120 (p.~383)} }
\expandafter\def\csname SRef@ari:eq:cp-periods\endcsname{\href{S6_SERIES_GRIND_I.pdf\#nameddest=tag-003G}{I.121 (p.~383)} }
\expandafter\def\csname SRef@ari:eq:cp-v\endcsname{\href{S6_SERIES_GRIND_I.pdf\#nameddest=tag-003L}{I.122 (p.~384)} }
\expandafter\def\csname SRef@ari:eq:cp-chi\endcsname{\href{S6_SERIES_GRIND_I.pdf\#nameddest=tag-003E}{I.123 (p.~384)} }
\expandafter\def\csname SRef@ari:eq:ca-derivatives\endcsname{\href{S6_SERIES_GRIND_I.pdf\#nameddest=tag-002B}{I.11 (p.~44)} }
\expandafter\def\csname SRef@ari:eq:cp-MTdim\endcsname{\href{S6_SERIES_GRIND_I.pdf\#nameddest=tag-003D}{I.124 (p.~385)} }
\expandafter\def\csname SRef@ari:eq:cp-positive-centre\endcsname{\href{S6_SERIES_GRIND_I.pdf\#nameddest=tag-003H}{I.125 (p.~386)} }
\expandafter\def\csname SRef@ari:eq:cp-real-linear\endcsname{\href{S6_SERIES_GRIND_I.pdf\#nameddest=tag-003I}{I.126 (p.~387)} }
\expandafter\def\csname SRef@ari:eq:cp-relation\endcsname{\href{S6_SERIES_GRIND_I.pdf\#nameddest=tag-003J}{I.127 (p.~387)} }
\expandafter\def\csname SRef@ari:eq:ct-CMmult\endcsname{\href{S6_SERIES_GRIND_I.pdf\#nameddest=tag-003Q}{I.128 (p.~388)} }
\expandafter\def\csname SRef@ari:eq:ct-Gaussian\endcsname{\href{S6_SERIES_GRIND_I.pdf\#nameddest=tag-003R}{I.129 (p.~389)} }
\expandafter\def\csname SRef@ari:eq:cd-slope\endcsname{\href{S6_SERIES_GRIND_I.pdf\#nameddest=tag-0033}{I.130 (p.~392)} }
\expandafter\def\csname SRef@ari:eq:cd-compatibility\endcsname{\href{S6_SERIES_GRIND_I.pdf\#nameddest=tag-0032}{I.131 (p.~392)} }
\expandafter\def\csname SRef@ari:eq:ce-ribet-r\endcsname{\href{S6_SERIES_GRIND_I.pdf\#nameddest=tag-0036}{I.132 (p.~395)} }
\expandafter\def\csname SRef@ari:eq:ce-zeta\endcsname{\href{S6_SERIES_GRIND_I.pdf\#nameddest=tag-0037}{I.133 (p.~395)} }
\expandafter\def\csname SRef@ari:eq:ca-delta-h\endcsname{\href{S6_SERIES_GRIND_I.pdf\#nameddest=tag-0028}{I.12 (p.~45)} }
\expandafter\def\csname SRef@ari:eq:ce-gap\endcsname{\href{S6_SERIES_GRIND_I.pdf\#nameddest=tag-0035}{I.134 (p.~396)} }
\expandafter\def\csname SRef@ari:eq:ce-explicit-witness\endcsname{\href{S6_SERIES_GRIND_I.pdf\#nameddest=tag-0034}{I.135 (p.~396)} }
\expandafter\def\csname SRef@ari:eq:cc-fixedfield\endcsname{\href{S6_SERIES_GRIND_I.pdf\#nameddest=tag-002Q}{I.136 (p.~399)} }
\expandafter\def\csname SRef@ari:eq:cc-rankcriterion\endcsname{\href{S6_SERIES_GRIND_I.pdf\#nameddest=tag-002X}{I.137 (p.~400)} }
\expandafter\def\csname SRef@ari:eq:division-ratio-model\endcsname{\href{S6_SERIES_GRIND_I.pdf\#nameddest=tag-003S}{I.138 (p.~401)} }
\expandafter\def\csname SRef@ari:eq:ca-ladder\endcsname{\href{S6_SERIES_GRIND_I.pdf\#nameddest=tag-002D}{I.13 (p.~47)} }
\expandafter\def\csname SRef@fam:eq:F-g\endcsname{\href{S6_SERIES_GRIND_I.pdf\#nameddest=tag-01OT}{I.14 (p.~69)} }
\expandafter\def\csname SRef@fam:eq:F-Gamma\endcsname{\href{S6_SERIES_GRIND_I.pdf\#nameddest=tag-01OQ}{I.15 (p.~69)} }
\expandafter\def\csname SRef@fam:eq:F-law\endcsname{\href{S6_SERIES_GRIND_I.pdf\#nameddest=tag-01OV}{I.16 (p.~69)} }
\expandafter\def\csname SRef@fam:eq:F-SLpresentation\endcsname{\href{S6_SERIES_GRIND_I.pdf\#nameddest=tag-01OS}{I.17 (p.~69)} }
\expandafter\def\csname SRef@fam:eq:F-H\endcsname{\href{S6_SERIES_GRIND_I.pdf\#nameddest=tag-01OR}{I.18 (p.~70)} }
\expandafter\def\csname SRef@fam:eq:F-qD\endcsname{\href{S6_SERIES_GRIND_I.pdf\#nameddest=tag-01OW}{I.19 (p.~73)} }
\expandafter\def\csname SRef@fam:eq:cd-wp\endcsname{\href{S6_SERIES_GRIND_I.pdf\#nameddest=tag-01P8}{I.2 (p.~30)} }
\expandafter\def\csname SRef@fam:eq:F-tauaction\endcsname{\href{S6_SERIES_GRIND_I.pdf\#nameddest=tag-01OX}{I.20 (p.~74)} }
\expandafter\def\csname SRef@fam:eq:F-j\endcsname{\href{S6_SERIES_GRIND_I.pdf\#nameddest=tag-01OU}{I.21 (p.~74)} }
\expandafter\def\csname SRef@fam:eq:cd-NS\endcsname{\href{S6_SERIES_GRIND_I.pdf\#nameddest=tag-01P4}{I.22 (p.~79)} }
\expandafter\def\csname SRef@fam:eq:cd-theta-sections\endcsname{\href{S6_SERIES_GRIND_I.pdf\#nameddest=tag-01P7}{I.23 (p.~84)} }
\expandafter\def\csname SRef@mapsI:closing:eq:C-explicit-closing\endcsname{\href{S6_SERIES_GRIND_I.pdf\#nameddest=tag-02OP}{I.24 (p.~93)} }
\expandafter\def\csname SRef@mapsI:closing:eq:C-twelve-scalar\endcsname{\href{S6_SERIES_GRIND_I.pdf\#nameddest=tag-02OQ}{I.25 (p.~93)} }
\expandafter\def\csname SRef@eq:d-period-normal\endcsname{\href{S6_SERIES_GRIND_I.pdf\#nameddest=tag-01ON}{I.26 (p.~230)} }
\expandafter\def\csname SRef@eq:d-limit-group\endcsname{\href{S6_SERIES_GRIND_I.pdf\#nameddest=tag-01OM}{I.27 (p.~233)} }
\expandafter\def\csname SRef@eq:d-dual-exact\endcsname{\href{S6_SERIES_GRIND_I.pdf\#nameddest=tag-01OL}{I.28 (p.~234)} }
\expandafter\def\csname SRef@eq:P2H\endcsname{\href{S6_SERIES_GRIND_I.pdf\#nameddest=tag-01O8}{I.29 (p.~243)} }
\expandafter\def\csname SRef@fam:eq:cd-norm\endcsname{\href{S6_SERIES_GRIND_I.pdf\#nameddest=tag-01P6}{I.3 (p.~31)} }
\expandafter\def\csname SRef@eq:P2zg\endcsname{\href{S6_SERIES_GRIND_I.pdf\#nameddest=tag-01OC}{I.30 (p.~243)} }
\expandafter\def\csname SRef@eq:P2F\endcsname{\href{S6_SERIES_GRIND_I.pdf\#nameddest=tag-01O5}{I.31 (p.~244)} }
\expandafter\def\csname SRef@eq:P2FQ\endcsname{\href{S6_SERIES_GRIND_I.pdf\#nameddest=tag-01O6}{I.32 (p.~244)} }
\expandafter\def\csname SRef@eq:P2FR\endcsname{\href{S6_SERIES_GRIND_I.pdf\#nameddest=tag-01O7}{I.33 (p.~244)} }
\expandafter\def\csname SRef@eq:P2chain\endcsname{\href{S6_SERIES_GRIND_I.pdf\#nameddest=tag-01O9}{I.34 (p.~245)} }
\expandafter\def\csname SRef@eq:P2split\endcsname{\href{S6_SERIES_GRIND_I.pdf\#nameddest=tag-01OB}{I.35 (p.~247)} }
\expandafter\def\csname SRef@eq:P2ratio\endcsname{\href{S6_SERIES_GRIND_I.pdf\#nameddest=tag-01OA}{I.36 (p.~248)} }
\expandafter\def\csname SRef@inv:eq:A-E-line\endcsname{\href{S6_SERIES_GRIND_I.pdf\#nameddest=tag-01SZ}{I.37 (p.~254)} }
\expandafter\def\csname SRef@inv:eq:A-H-images\endcsname{\href{S6_SERIES_GRIND_I.pdf\#nameddest=tag-01T0}{I.38 (p.~255)} }
\expandafter\def\csname SRef@inv:eq:A-N-images\endcsname{\href{S6_SERIES_GRIND_I.pdf\#nameddest=tag-01T1}{I.39 (p.~255)} }
\expandafter\def\csname SRef@ari:eq:cs-isogeny\endcsname{\href{S6_SERIES_GRIND_I.pdf\#nameddest=tag-003O}{I.4 (p.~36)} }
\expandafter\def\csname SRef@inv:eq:A-absolute-table\endcsname{\href{S6_SERIES_GRIND_I.pdf\#nameddest=tag-01T2}{I.40 (p.~256)} }
\expandafter\def\csname SRef@inv:eq:A-s-branch\endcsname{\href{S6_SERIES_GRIND_I.pdf\#nameddest=tag-01T5}{I.41 (p.~257)} }
\expandafter\def\csname SRef@inv:eq:A-modular-K\endcsname{\href{S6_SERIES_GRIND_I.pdf\#nameddest=tag-01T3}{I.42 (p.~257)} }
\expandafter\def\csname SRef@inv:eq:A-mu-value\endcsname{\href{S6_SERIES_GRIND_I.pdf\#nameddest=tag-01T4}{I.43 (p.~258)} }
\expandafter\def\csname SRef@inv:eq:A-DI-branch\endcsname{\href{S6_SERIES_GRIND_I.pdf\#nameddest=tag-01SY}{I.44 (p.~260)} }
\expandafter\def\csname SRef@inv:eq:C-hodge-input\endcsname{\href{S6_SERIES_GRIND_I.pdf\#nameddest=tag-01TF}{I.45 (p.~266)} }
\expandafter\def\csname SRef@inv:eq:C-BC-A\endcsname{\href{S6_SERIES_GRIND_I.pdf\#nameddest=tag-01T6}{I.46 (p.~267)} }
\expandafter\def\csname SRef@inv:eq:C-arrows\endcsname{\href{S6_SERIES_GRIND_I.pdf\#nameddest=tag-01TA}{I.47 (p.~268)} }
\expandafter\def\csname SRef@inv:eq:C-Pic-map\endcsname{\href{S6_SERIES_GRIND_I.pdf\#nameddest=tag-01T8}{I.48 (p.~269)} }
\expandafter\def\csname SRef@inv:eq:C-G-Green\endcsname{\href{S6_SERIES_GRIND_I.pdf\#nameddest=tag-01T7}{I.49 (p.~269)} }
\expandafter\def\csname SRef@ari:eq:cs-division\endcsname{\href{S6_SERIES_GRIND_I.pdf\#nameddest=tag-003N}{I.5 (p.~39)} }
\expandafter\def\csname SRef@inv:eq:C-factor12\endcsname{\href{S6_SERIES_GRIND_I.pdf\#nameddest=tag-01TE}{I.50 (p.~270)} }
\expandafter\def\csname SRef@inv:eq:C-collar-sum\endcsname{\href{S6_SERIES_GRIND_I.pdf\#nameddest=tag-01TB}{I.51 (p.~271)} }
\expandafter\def\csname SRef@inv:eq:C-delta\endcsname{\href{S6_SERIES_GRIND_I.pdf\#nameddest=tag-01TC}{I.52 (p.~273)} }
\expandafter\def\csname SRef@inv:eq:C-Ricci\endcsname{\href{S6_SERIES_GRIND_I.pdf\#nameddest=tag-01T9}{I.53 (p.~274)} }
\expandafter\def\csname SRef@inv:eq:C-eta-negative\endcsname{\href{S6_SERIES_GRIND_I.pdf\#nameddest=tag-01TD}{I.54 (p.~274)} }
\expandafter\def\csname SRef@ari:eq:cs-Li\endcsname{\href{S6_SERIES_GRIND_I.pdf\#nameddest=tag-003M}{I.55 (p.~301)} }
\expandafter\def\csname SRef@ari:eq:ch-global\endcsname{\href{S6_SERIES_GRIND_I.pdf\#nameddest=tag-0039}{I.56 (p.~304)} }
\expandafter\def\csname SRef@ari:eq:ch-local-mult\endcsname{\href{S6_SERIES_GRIND_I.pdf\#nameddest=tag-003B}{I.57 (p.~304)} }
\expandafter\def\csname SRef@ari:eq:ch-local-exact\endcsname{\href{S6_SERIES_GRIND_I.pdf\#nameddest=tag-003A}{I.58 (p.~305)} }
\expandafter\def\csname SRef@ari:eq:ch-bound\endcsname{\href{S6_SERIES_GRIND_I.pdf\#nameddest=tag-0038}{I.59 (p.~308)} }
\expandafter\def\csname SRef@ari:eq:cs-torsionvalue\endcsname{\href{S6_SERIES_GRIND_I.pdf\#nameddest=tag-003P}{I.6 (p.~39)} }
\expandafter\def\csname SRef@ari:eq:ch-wp-bracket\endcsname{\href{S6_SERIES_GRIND_I.pdf\#nameddest=tag-003C}{I.60 (p.~309)} }
\expandafter\def\csname SRef@ari:eq:Cr-first\endcsname{\href{S6_SERIES_GRIND_I.pdf\#nameddest=tag-0020}{I.61 (p.~312)} }
\expandafter\def\csname SRef@ari:eq:Cr-second\endcsname{\href{S6_SERIES_GRIND_I.pdf\#nameddest=tag-0022}{I.62 (p.~312)} }
\expandafter\def\csname SRef@ari:eq:Cr-tails\endcsname{\href{S6_SERIES_GRIND_I.pdf\#nameddest=tag-0023}{I.63 (p.~312)} }
\expandafter\def\csname SRef@ari:eq:Cr-rational-factor\endcsname{\href{S6_SERIES_GRIND_I.pdf\#nameddest=tag-0021}{I.64 (p.~314)} }
\expandafter\def\csname SRef@ari:eq:Cgz-central-integrality\endcsname{\href{S6_SERIES_GRIND_I.pdf\#nameddest=tag-001X}{I.65 (p.~318)} }
\expandafter\def\csname SRef@ari:eq:Cgz-product\endcsname{\href{S6_SERIES_GRIND_I.pdf\#nameddest=tag-001Z}{I.66 (p.~319)} }
\expandafter\def\csname SRef@ari:eq:Cgz-denominator\endcsname{\href{S6_SERIES_GRIND_I.pdf\#nameddest=tag-001Y}{I.67 (p.~319)} }
\expandafter\def\csname SRef@ari:eq:ca-mahler\endcsname{\href{S6_SERIES_GRIND_I.pdf\#nameddest=tag-002E}{I.68 (p.~325)} }
\expandafter\def\csname SRef@ari:eq:ca-wpmean\endcsname{\href{S6_SERIES_GRIND_I.pdf\#nameddest=tag-002H}{I.69 (p.~325)} }
\expandafter\def\csname SRef@ari:eq:ca-sup\endcsname{\href{S6_SERIES_GRIND_I.pdf\#nameddest=tag-002G}{I.7 (p.~43)} }
\expandafter\def\csname SRef@ari:eq:ca-nu\endcsname{\href{S6_SERIES_GRIND_I.pdf\#nameddest=tag-002F}{I.70 (p.~325)} }
\expandafter\def\csname SRef@ari:eq:ca-current\endcsname{\href{S6_SERIES_GRIND_I.pdf\#nameddest=tag-0026}{I.71 (p.~327)} }
\expandafter\def\csname SRef@ari:eq:ca-density0\endcsname{\href{S6_SERIES_GRIND_I.pdf\#nameddest=tag-0029}{I.72 (p.~328)} }
\expandafter\def\csname SRef@ari:eq:ca-density1\endcsname{\href{S6_SERIES_GRIND_I.pdf\#nameddest=tag-002A}{I.73 (p.~328)} }
\expandafter\def\csname SRef@ari:eq:ca-cuspexp\endcsname{\href{S6_SERIES_GRIND_I.pdf\#nameddest=tag-0027}{I.74 (p.~328)} }
\expandafter\def\csname SRef@ari:eq:cc-EZ\endcsname{\href{S6_SERIES_GRIND_I.pdf\#nameddest=tag-002I}{I.75 (p.~330)} }
\expandafter\def\csname SRef@ari:eq:cc-radial\endcsname{\href{S6_SERIES_GRIND_I.pdf\#nameddest=tag-002W}{I.76 (p.~331)} }
\expandafter\def\csname SRef@ari:eq:cc-FR\endcsname{\href{S6_SERIES_GRIND_I.pdf\#nameddest=tag-002J}{I.77 (p.~331)} }
\expandafter\def\csname SRef@ari:eq:cc-cuspbound\endcsname{\href{S6_SERIES_GRIND_I.pdf\#nameddest=tag-002O}{I.78 (p.~332)} }
\expandafter\def\csname SRef@ari:eq:cc-ode\endcsname{\href{S6_SERIES_GRIND_I.pdf\#nameddest=tag-002T}{I.79 (p.~333)} }
\expandafter\def\csname SRef@ari:eq:ca-exp\endcsname{\href{S6_SERIES_GRIND_I.pdf\#nameddest=tag-002C}{I.8 (p.~43)} }
\expandafter\def\csname SRef@ari:eq:cc-h\endcsname{\href{S6_SERIES_GRIND_I.pdf\#nameddest=tag-002R}{I.80 (p.~333)} }
\expandafter\def\csname SRef@ari:eq:cc-Fseries\endcsname{\href{S6_SERIES_GRIND_I.pdf\#nameddest=tag-002L}{I.81 (p.~333)} }
\expandafter\def\csname SRef@ari:eq:cc-Ffactored\endcsname{\href{S6_SERIES_GRIND_I.pdf\#nameddest=tag-002K}{I.82 (p.~334)} }
\expandafter\def\csname SRef@ari:eq:cc-norms\endcsname{\href{S6_SERIES_GRIND_I.pdf\#nameddest=tag-002S}{I.83 (p.~334)} }
\expandafter\def\csname SRef@ari:eq:cc-recurrence\endcsname{\href{S6_SERIES_GRIND_I.pdf\#nameddest=tag-002Y}{I.84 (p.~334)} }
\expandafter\def\csname SRef@ari:eq:cc-power\endcsname{\href{S6_SERIES_GRIND_I.pdf\#nameddest=tag-002V}{I.85 (p.~334)} }
\expandafter\def\csname SRef@ari:eq:cc-constants\endcsname{\href{S6_SERIES_GRIND_I.pdf\#nameddest=tag-002N}{I.86 (p.~334)} }
\expandafter\def\csname SRef@ari:eq:cc-outer\endcsname{\href{S6_SERIES_GRIND_I.pdf\#nameddest=tag-002U}{I.87 (p.~335)} }
\expandafter\def\csname SRef@ari:eq:cc-tail60\endcsname{\href{S6_SERIES_GRIND_I.pdf\#nameddest=tag-0031}{I.88 (p.~335)} }
\expandafter\def\csname SRef@ari:eq:cc-tail100\endcsname{\href{S6_SERIES_GRIND_I.pdf\#nameddest=tag-0030}{I.89 (p.~335)} }
\expandafter\def\csname SRef@ari:eq:ca-H\endcsname{\href{S6_SERIES_GRIND_I.pdf\#nameddest=tag-0025}{I.9 (p.~44)} }
\expandafter\def\csname SRef@ari:eq:cc-finite-table\endcsname{\href{S6_SERIES_GRIND_I.pdf\#nameddest=tag-002P}{I.90 (p.~336)} }
\expandafter\def\csname SRef@ari:eq:cc-scalar\endcsname{\href{S6_SERIES_GRIND_I.pdf\#nameddest=tag-002Z}{I.91 (p.~336)} }
\expandafter\def\csname SRef@ari:eq:cc-P\endcsname{\href{S6_SERIES_GRIND_I.pdf\#nameddest=tag-002M}{I.92 (p.~336)} }
\expandafter\def\csname SRef@fam:eq:ab-root-degree\endcsname{\href{S6_SERIES_GRIND_I.pdf\#nameddest=tag-01P3}{I.93 (p.~356)} }
\expandafter\def\csname SRef@fam:eq:ab-cusp-resolution\endcsname{\href{S6_SERIES_GRIND_I.pdf\#nameddest=tag-01OY}{I.94 (p.~357)} }
\expandafter\def\csname SRef@fam:eq:ab-log-tangent\endcsname{\href{S6_SERIES_GRIND_I.pdf\#nameddest=tag-01P0}{I.95 (p.~358)} }
\expandafter\def\csname SRef@fam:eq:ab-relative-dualizing\endcsname{\href{S6_SERIES_GRIND_I.pdf\#nameddest=tag-01P2}{I.96 (p.~358)} }
\expandafter\def\csname SRef@fam:eq:ab-direct-images\endcsname{\href{S6_SERIES_GRIND_I.pdf\#nameddest=tag-01OZ}{I.97 (p.~358)} }
\expandafter\def\csname SRef@fam:eq:ab-normal-quotient\endcsname{\href{S6_SERIES_GRIND_I.pdf\#nameddest=tag-01P1}{I.98 (p.~359)} }
\expandafter\def\csname SRef@eq:P5tangent\endcsname{\href{S6_SERIES_GRIND_I.pdf\#nameddest=tag-01OK}{I.99 (p.~366)} }
\expandafter\def\csname SRef@fig:I:orbifold\endcsname{\href{S6_SERIES_GRIND_I.pdf\#nameddest=tag-01RL}{I.2 (p.~94)} }
\expandafter\def\csname SRef@fig:I:parameter-disc\endcsname{\href{S6_SERIES_GRIND_I.pdf\#nameddest=tag-01RM}{I.3 (p.~94)} }
\expandafter\def\csname SRef@fig:I:monodromy-domain\endcsname{\href{S6_SERIES_GRIND_I.pdf\#nameddest=tag-01RK}{I.4 (p.~95)} }
\expandafter\def\csname SRef@push:src:A.3.1\endcsname{\href{S6_SERIES_GRIND_I.pdf\#nameddest=tag-02ZO}{I.14.7 (p.~113)} }
\expandafter\def\csname SRef@push:A31\endcsname{\href{S6_SERIES_GRIND_I.pdf\#nameddest=tag-02ZO}{I.14.7 (p.~113)} }
\expandafter\def\csname SRef@push:src:D.1.2\endcsname{\href{S6_SERIES_GRIND_I.pdf\#nameddest=tag-032B}{I.16.2 (p.~228)} }
\expandafter\def\csname SRef@push:src:D.2.1\endcsname{\href{S6_SERIES_GRIND_I.pdf\#nameddest=tag-032D}{I.16.4 (p.~230)} }
\expandafter\def\csname SRef@I:weierstrass-comparison\endcsname{\href{S6_SERIES_GRIND_I.pdf\#nameddest=tag-0002}{I.3.1 (p.~24)} }
\expandafter\def\csname SRef@wb:pshift:lem:centralizer\endcsname{\href{S6_SERIES_GRIND_I.pdf\#nameddest=tag-03CQ}{I.A.4 (p.~349)} }
\expandafter\def\csname SRef@fam:src:F.1.2\endcsname{\href{S6_SERIES_GRIND_I.pdf\#nameddest=tag-01RA}{I.10.2 (p.~70)} }
\expandafter\def\csname SRef@fam:st:F.1.2\endcsname{\href{S6_SERIES_GRIND_I.pdf\#nameddest=tag-01RA}{I.10.2 (p.~70)} }
\expandafter\def\csname SRef@fam:src:F.2.1\endcsname{\href{S6_SERIES_GRIND_I.pdf\#nameddest=tag-01RB}{I.10.3 (p.~71)} }
\expandafter\def\csname SRef@fam:st:F.2.1\endcsname{\href{S6_SERIES_GRIND_I.pdf\#nameddest=tag-01RB}{I.10.3 (p.~71)} }
\expandafter\def\csname SRef@fam:src:F.2.3\endcsname{\href{S6_SERIES_GRIND_I.pdf\#nameddest=tag-01RD}{I.10.5 (p.~73)} }
\expandafter\def\csname SRef@fam:st:F.2.3\endcsname{\href{S6_SERIES_GRIND_I.pdf\#nameddest=tag-01RD}{I.10.5 (p.~73)} }
\expandafter\def\csname SRef@fam:src:F.3.2\endcsname{\href{S6_SERIES_GRIND_I.pdf\#nameddest=tag-01RF}{I.10.7 (p.~75)} }
\expandafter\def\csname SRef@fam:st:F.3.2\endcsname{\href{S6_SERIES_GRIND_I.pdf\#nameddest=tag-01RF}{I.10.7 (p.~75)} }
\expandafter\def\csname SRef@fam:src:F.3.3\endcsname{\href{S6_SERIES_GRIND_I.pdf\#nameddest=tag-01RG}{I.10.8 (p.~75)} }
\expandafter\def\csname SRef@fam:st:F.3.3\endcsname{\href{S6_SERIES_GRIND_I.pdf\#nameddest=tag-01RG}{I.10.8 (p.~75)} }
\expandafter\def\csname SRef@fam:src:D.1.1\endcsname{\href{S6_SERIES_GRIND_I.pdf\#nameddest=tag-01QR}{I.11.1 (p.~76)} }
\expandafter\def\csname SRef@fam:st:D.1.1\endcsname{\href{S6_SERIES_GRIND_I.pdf\#nameddest=tag-01QR}{I.11.1 (p.~76)} }
\expandafter\def\csname SRef@fam:src:C.3.1\endcsname{\href{S6_SERIES_GRIND_I.pdf\#nameddest=tag-01QM}{I.11.11 (p.~84)} }
\expandafter\def\csname SRef@fam:st:C.3.1\endcsname{\href{S6_SERIES_GRIND_I.pdf\#nameddest=tag-01QM}{I.11.11 (p.~84)} }
\expandafter\def\csname SRef@fam:src:C.3.2\endcsname{\href{S6_SERIES_GRIND_I.pdf\#nameddest=tag-01QN}{I.11.12 (p.~85)} }
\expandafter\def\csname SRef@fam:st:C.3.2\endcsname{\href{S6_SERIES_GRIND_I.pdf\#nameddest=tag-01QN}{I.11.12 (p.~85)} }
\expandafter\def\csname SRef@fam:src:C.3.4\endcsname{\href{S6_SERIES_GRIND_I.pdf\#nameddest=tag-01QP}{I.11.13 (p.~86)} }
\expandafter\def\csname SRef@fam:st:C.3.4\endcsname{\href{S6_SERIES_GRIND_I.pdf\#nameddest=tag-01QP}{I.11.13 (p.~86)} }
\expandafter\def\csname SRef@fam:src:D.1.2\endcsname{\href{S6_SERIES_GRIND_I.pdf\#nameddest=tag-01QS}{I.11.2 (p.~77)} }
\expandafter\def\csname SRef@fam:st:D.1.2\endcsname{\href{S6_SERIES_GRIND_I.pdf\#nameddest=tag-01QS}{I.11.2 (p.~77)} }
\expandafter\def\csname SRef@fam:src:D.1.4\endcsname{\href{S6_SERIES_GRIND_I.pdf\#nameddest=tag-01QU}{I.11.4 (p.~79)} }
\expandafter\def\csname SRef@fam:st:D.1.4\endcsname{\href{S6_SERIES_GRIND_I.pdf\#nameddest=tag-01QU}{I.11.4 (p.~79)} }
\expandafter\def\csname SRef@fam:src:D.1.5\endcsname{\href{S6_SERIES_GRIND_I.pdf\#nameddest=tag-01QV}{I.11.5 (p.~79)} }
\expandafter\def\csname SRef@fam:st:D.1.5\endcsname{\href{S6_SERIES_GRIND_I.pdf\#nameddest=tag-01QV}{I.11.5 (p.~79)} }
\expandafter\def\csname SRef@fam:src:D.1.7\endcsname{\href{S6_SERIES_GRIND_I.pdf\#nameddest=tag-01QX}{I.11.7 (p.~81)} }
\expandafter\def\csname SRef@fam:st:D.1.7\endcsname{\href{S6_SERIES_GRIND_I.pdf\#nameddest=tag-01QX}{I.11.7 (p.~81)} }
\expandafter\def\csname SRef@fam:src:D.2.1\endcsname{\href{S6_SERIES_GRIND_I.pdf\#nameddest=tag-01QY}{I.11.8 (p.~82)} }
\expandafter\def\csname SRef@fam:st:D.2.1\endcsname{\href{S6_SERIES_GRIND_I.pdf\#nameddest=tag-01QY}{I.11.8 (p.~82)} }
\expandafter\def\csname SRef@fam:src:D.2.2\endcsname{\href{S6_SERIES_GRIND_I.pdf\#nameddest=tag-01QZ}{I.11.9 (p.~83)} }
\expandafter\def\csname SRef@fam:st:D.2.2\endcsname{\href{S6_SERIES_GRIND_I.pdf\#nameddest=tag-01QZ}{I.11.9 (p.~83)} }
\expandafter\def\csname SRef@mapsI:src:A.3.1\endcsname{\href{S6_SERIES_GRIND_I.pdf\#nameddest=tag-02PE}{I.12.1 (p.~88)} }
\expandafter\def\csname SRef@mapsI:src:A.3.4\endcsname{\href{S6_SERIES_GRIND_I.pdf\#nameddest=tag-02PH}{I.12.4 (p.~93)} }
\expandafter\def\csname SRef@mapsI:real-eq-9\endcsname{\href{S6_SERIES_GRIND_I.pdf\#nameddest=tag-02PH}{I.12.4 (p.~93)} }
\expandafter\def\csname SRef@I:monodromy-fundamental-set\endcsname{\href{S6_SERIES_GRIND_I.pdf\#nameddest=tag-0001}{I.12.6 (p.~94)} }
\expandafter\def\csname SRef@fam:src:E.1.1\endcsname{\href{S6_SERIES_GRIND_I.pdf\#nameddest=tag-01R1}{I.13.1 (p.~96)} }
\expandafter\def\csname SRef@fam:st:E.1.1\endcsname{\href{S6_SERIES_GRIND_I.pdf\#nameddest=tag-01R1}{I.13.1 (p.~96)} }
\expandafter\def\csname SRef@fam:src:E.1.2\endcsname{\href{S6_SERIES_GRIND_I.pdf\#nameddest=tag-01R2}{I.13.2 (p.~97)} }
\expandafter\def\csname SRef@fam:st:E.1.2\endcsname{\href{S6_SERIES_GRIND_I.pdf\#nameddest=tag-01R2}{I.13.2 (p.~97)} }
\expandafter\def\csname SRef@fam:src:E.1.3\endcsname{\href{S6_SERIES_GRIND_I.pdf\#nameddest=tag-01R3}{I.13.3 (p.~98)} }
\expandafter\def\csname SRef@fam:st:E.1.3\endcsname{\href{S6_SERIES_GRIND_I.pdf\#nameddest=tag-01R3}{I.13.3 (p.~98)} }
\expandafter\def\csname SRef@fam:src:E.1.4\endcsname{\href{S6_SERIES_GRIND_I.pdf\#nameddest=tag-01R4}{I.13.4 (p.~99)} }
\expandafter\def\csname SRef@fam:st:E.1.4\endcsname{\href{S6_SERIES_GRIND_I.pdf\#nameddest=tag-01R4}{I.13.4 (p.~99)} }
\expandafter\def\csname SRef@fam:src:E.2.1\endcsname{\href{S6_SERIES_GRIND_I.pdf\#nameddest=tag-01R6}{I.13.5 (p.~101)} }
\expandafter\def\csname SRef@fam:st:E.2.1\endcsname{\href{S6_SERIES_GRIND_I.pdf\#nameddest=tag-01R6}{I.13.5 (p.~101)} }
\expandafter\def\csname SRef@push:src:A.1.1\endcsname{\href{S6_SERIES_GRIND_I.pdf\#nameddest=tag-02ZH}{I.14.1 (p.~103)} }
\expandafter\def\csname SRef@push:A11\endcsname{\href{S6_SERIES_GRIND_I.pdf\#nameddest=tag-02ZH}{I.14.1 (p.~103)} }
\expandafter\def\csname SRef@push:src:A.4.1\endcsname{\href{S6_SERIES_GRIND_I.pdf\#nameddest=tag-02ZR}{I.14.10 (p.~116)} }
\expandafter\def\csname SRef@push:A41\endcsname{\href{S6_SERIES_GRIND_I.pdf\#nameddest=tag-02ZR}{I.14.10 (p.~116)} }
\expandafter\def\csname SRef@push:src:A.4.2\endcsname{\href{S6_SERIES_GRIND_I.pdf\#nameddest=tag-02ZS}{I.14.11 (p.~118)} }
\expandafter\def\csname SRef@push:A42\endcsname{\href{S6_SERIES_GRIND_I.pdf\#nameddest=tag-02ZS}{I.14.11 (p.~118)} }
\expandafter\def\csname SRef@push:src:A.4.3\endcsname{\href{S6_SERIES_GRIND_I.pdf\#nameddest=tag-02ZT}{I.14.12 (p.~119)} }
\expandafter\def\csname SRef@push:A43\endcsname{\href{S6_SERIES_GRIND_I.pdf\#nameddest=tag-02ZT}{I.14.12 (p.~119)} }
\expandafter\def\csname SRef@push:src:A.4.4\endcsname{\href{S6_SERIES_GRIND_I.pdf\#nameddest=tag-02ZU}{I.14.13 (p.~120)} }
\expandafter\def\csname SRef@push:A44\endcsname{\href{S6_SERIES_GRIND_I.pdf\#nameddest=tag-02ZU}{I.14.13 (p.~120)} }
\expandafter\def\csname SRef@push:src:A.4.6\endcsname{\href{S6_SERIES_GRIND_I.pdf\#nameddest=tag-02ZW}{I.14.15 (p.~121)} }
\expandafter\def\csname SRef@push:A46\endcsname{\href{S6_SERIES_GRIND_I.pdf\#nameddest=tag-02ZW}{I.14.15 (p.~121)} }
\expandafter\def\csname SRef@push:src:A.4.9\endcsname{\href{S6_SERIES_GRIND_I.pdf\#nameddest=tag-02ZZ}{I.14.18 (p.~126)} }
\expandafter\def\csname SRef@push:A49\endcsname{\href{S6_SERIES_GRIND_I.pdf\#nameddest=tag-02ZZ}{I.14.18 (p.~126)} }
\expandafter\def\csname SRef@push:src:A.2.1\endcsname{\href{S6_SERIES_GRIND_I.pdf\#nameddest=tag-02ZJ}{I.14.3 (p.~105)} }
\expandafter\def\csname SRef@push:A21\endcsname{\href{S6_SERIES_GRIND_I.pdf\#nameddest=tag-02ZJ}{I.14.3 (p.~105)} }
\expandafter\def\csname SRef@push:src:A.2.5\endcsname{\href{S6_SERIES_GRIND_I.pdf\#nameddest=tag-02ZN}{I.14.6 (p.~111)} }
\expandafter\def\csname SRef@push:A25\endcsname{\href{S6_SERIES_GRIND_I.pdf\#nameddest=tag-02ZN}{I.14.6 (p.~111)} }
\expandafter\def\csname SRef@push:src:A.3.3\endcsname{\href{S6_SERIES_GRIND_I.pdf\#nameddest=tag-02ZQ}{I.14.9 (p.~115)} }
\expandafter\def\csname SRef@push:A33\endcsname{\href{S6_SERIES_GRIND_I.pdf\#nameddest=tag-02ZQ}{I.14.9 (p.~115)} }
\expandafter\def\csname SRef@push:src:B.1.1\endcsname{\href{S6_SERIES_GRIND_I.pdf\#nameddest=tag-0300}{I.15.1 (p.~128)} }
\expandafter\def\csname SRef@push:B11\endcsname{\href{S6_SERIES_GRIND_I.pdf\#nameddest=tag-0300}{I.15.1 (p.~128)} }
\expandafter\def\csname SRef@push:src:B.3.6\endcsname{\href{S6_SERIES_GRIND_I.pdf\#nameddest=tag-030A}{I.15.11 (p.~138)} }
\expandafter\def\csname SRef@push:B36\endcsname{\href{S6_SERIES_GRIND_I.pdf\#nameddest=tag-030A}{I.15.11 (p.~138)} }
\expandafter\def\csname SRef@push:src:B.3.7\endcsname{\href{S6_SERIES_GRIND_I.pdf\#nameddest=tag-030B}{I.15.12 (p.~140)} }
\expandafter\def\csname SRef@push:B37\endcsname{\href{S6_SERIES_GRIND_I.pdf\#nameddest=tag-030B}{I.15.12 (p.~140)} }
\expandafter\def\csname SRef@push:src:B.3.8\endcsname{\href{S6_SERIES_GRIND_I.pdf\#nameddest=tag-030C}{I.15.13 (p.~142)} }
\expandafter\def\csname SRef@push:B38\endcsname{\href{S6_SERIES_GRIND_I.pdf\#nameddest=tag-030C}{I.15.13 (p.~142)} }
\expandafter\def\csname SRef@push:src:B.4.4\endcsname{\href{S6_SERIES_GRIND_I.pdf\#nameddest=tag-0314}{I.15.17 (p.~145)} }
\expandafter\def\csname SRef@push:src:B.4.6\endcsname{\href{S6_SERIES_GRIND_I.pdf\#nameddest=tag-0316}{I.15.19 (p.~148)} }
\expandafter\def\csname SRef@push:B46\endcsname{\href{S6_SERIES_GRIND_I.pdf\#nameddest=tag-0316}{I.15.19 (p.~148)} }
\expandafter\def\csname SRef@push:src:B.4.7\endcsname{\href{S6_SERIES_GRIND_I.pdf\#nameddest=tag-0317}{I.15.20 (p.~151)} }
\expandafter\def\csname SRef@push:src:B.4.8\endcsname{\href{S6_SERIES_GRIND_I.pdf\#nameddest=tag-0318}{I.15.21 (p.~155)} }
\expandafter\def\csname SRef@push:src:B.4.9\endcsname{\href{S6_SERIES_GRIND_I.pdf\#nameddest=tag-0319}{I.15.22 (p.~158)} }
\expandafter\def\csname SRef@push:src:B.4.10\endcsname{\href{S6_SERIES_GRIND_I.pdf\#nameddest=tag-030E}{I.15.23 (p.~160)} }
\expandafter\def\csname SRef@push:src:B.4.11\endcsname{\href{S6_SERIES_GRIND_I.pdf\#nameddest=tag-030F}{I.15.24 (p.~164)} }
\expandafter\def\csname SRef@push:src:B.4.12\endcsname{\href{S6_SERIES_GRIND_I.pdf\#nameddest=tag-030G}{I.15.25 (p.~167)} }
\expandafter\def\csname SRef@push:src:B.4.13\endcsname{\href{S6_SERIES_GRIND_I.pdf\#nameddest=tag-030H}{I.15.26 (p.~169)} }
\expandafter\def\csname SRef@push:src:B.4.14\endcsname{\href{S6_SERIES_GRIND_I.pdf\#nameddest=tag-030I}{I.15.27 (p.~171)} }
\expandafter\def\csname SRef@push:src:B.4.15\endcsname{\href{S6_SERIES_GRIND_I.pdf\#nameddest=tag-030J}{I.15.28 (p.~174)} }
\expandafter\def\csname SRef@push:B415\endcsname{\href{S6_SERIES_GRIND_I.pdf\#nameddest=tag-030J}{I.15.28 (p.~174)} }
\expandafter\def\csname SRef@push:src:B.4.16\endcsname{\href{S6_SERIES_GRIND_I.pdf\#nameddest=tag-030K}{I.15.29 (p.~177)} }
\expandafter\def\csname SRef@push:B416\endcsname{\href{S6_SERIES_GRIND_I.pdf\#nameddest=tag-030K}{I.15.29 (p.~177)} }
\expandafter\def\csname SRef@push:src:B.2.1\endcsname{\href{S6_SERIES_GRIND_I.pdf\#nameddest=tag-0302}{I.15.3 (p.~131)} }
\expandafter\def\csname SRef@push:B21\endcsname{\href{S6_SERIES_GRIND_I.pdf\#nameddest=tag-0302}{I.15.3 (p.~131)} }
\expandafter\def\csname SRef@push:src:B.4.17\endcsname{\href{S6_SERIES_GRIND_I.pdf\#nameddest=tag-030L}{I.15.30 (p.~178)} }
\expandafter\def\csname SRef@push:src:B.4.18\endcsname{\href{S6_SERIES_GRIND_I.pdf\#nameddest=tag-030M}{I.15.31 (p.~180)} }
\expandafter\def\csname SRef@push:B418\endcsname{\href{S6_SERIES_GRIND_I.pdf\#nameddest=tag-030M}{I.15.31 (p.~180)} }
\expandafter\def\csname SRef@push:src:B.4.19\endcsname{\href{S6_SERIES_GRIND_I.pdf\#nameddest=tag-030N}{I.15.32 (p.~182)} }
\expandafter\def\csname SRef@push:B419\endcsname{\href{S6_SERIES_GRIND_I.pdf\#nameddest=tag-030N}{I.15.32 (p.~182)} }
\expandafter\def\csname SRef@push:src:B.4.20\endcsname{\href{S6_SERIES_GRIND_I.pdf\#nameddest=tag-030P}{I.15.33 (p.~185)} }
\expandafter\def\csname SRef@push:B420\endcsname{\href{S6_SERIES_GRIND_I.pdf\#nameddest=tag-030P}{I.15.33 (p.~185)} }
\expandafter\def\csname SRef@push:src:B.4.21\endcsname{\href{S6_SERIES_GRIND_I.pdf\#nameddest=tag-030Q}{I.15.34 (p.~187)} }
\expandafter\def\csname SRef@push:src:B.4.22\endcsname{\href{S6_SERIES_GRIND_I.pdf\#nameddest=tag-030R}{I.15.35 (p.~191)} }
\expandafter\def\csname SRef@push:B422\endcsname{\href{S6_SERIES_GRIND_I.pdf\#nameddest=tag-030R}{I.15.35 (p.~191)} }
\expandafter\def\csname SRef@push:src:B.4.23\endcsname{\href{S6_SERIES_GRIND_I.pdf\#nameddest=tag-030S}{I.15.36 (p.~193)} }
\expandafter\def\csname SRef@push:src:B.4.24\endcsname{\href{S6_SERIES_GRIND_I.pdf\#nameddest=tag-030T}{I.15.37 (p.~195)} }
\expandafter\def\csname SRef@push:B424\endcsname{\href{S6_SERIES_GRIND_I.pdf\#nameddest=tag-030T}{I.15.37 (p.~195)} }
\expandafter\def\csname SRef@push:src:B.4.25\endcsname{\href{S6_SERIES_GRIND_I.pdf\#nameddest=tag-030U}{I.15.38 (p.~198)} }
\expandafter\def\csname SRef@push:B425\endcsname{\href{S6_SERIES_GRIND_I.pdf\#nameddest=tag-030U}{I.15.38 (p.~198)} }
\expandafter\def\csname SRef@push:src:B.4.26\endcsname{\href{S6_SERIES_GRIND_I.pdf\#nameddest=tag-030V}{I.15.39 (p.~200)} }
\expandafter\def\csname SRef@push:B426\endcsname{\href{S6_SERIES_GRIND_I.pdf\#nameddest=tag-030V}{I.15.39 (p.~200)} }
\expandafter\def\csname SRef@push:src:B.4.27\endcsname{\href{S6_SERIES_GRIND_I.pdf\#nameddest=tag-030W}{I.15.40 (p.~202)} }
\expandafter\def\csname SRef@push:B427\endcsname{\href{S6_SERIES_GRIND_I.pdf\#nameddest=tag-030W}{I.15.40 (p.~202)} }
\expandafter\def\csname SRef@push:src:B.4.28\endcsname{\href{S6_SERIES_GRIND_I.pdf\#nameddest=tag-030X}{I.15.41 (p.~204)} }
\expandafter\def\csname SRef@push:src:B.4.29\endcsname{\href{S6_SERIES_GRIND_I.pdf\#nameddest=tag-030Y}{I.15.42 (p.~208)} }
\expandafter\def\csname SRef@push:B429\endcsname{\href{S6_SERIES_GRIND_I.pdf\#nameddest=tag-030Y}{I.15.42 (p.~208)} }
\expandafter\def\csname SRef@push:src:B.4.30\endcsname{\href{S6_SERIES_GRIND_I.pdf\#nameddest=tag-0310}{I.15.43 (p.~210)} }
\expandafter\def\csname SRef@push:B430\endcsname{\href{S6_SERIES_GRIND_I.pdf\#nameddest=tag-0310}{I.15.43 (p.~210)} }
\expandafter\def\csname SRef@push:src:B.4.31\endcsname{\href{S6_SERIES_GRIND_I.pdf\#nameddest=tag-0311}{I.15.44 (p.~216)} }
\expandafter\def\csname SRef@push:src:B.4.32\endcsname{\href{S6_SERIES_GRIND_I.pdf\#nameddest=tag-0312}{I.15.45 (p.~219)} }
\expandafter\def\csname SRef@push:B432\endcsname{\href{S6_SERIES_GRIND_I.pdf\#nameddest=tag-0312}{I.15.45 (p.~219)} }
\expandafter\def\csname SRef@push:src:B.3.2\endcsname{\href{S6_SERIES_GRIND_I.pdf\#nameddest=tag-0306}{I.15.7 (p.~134)} }
\expandafter\def\csname SRef@push:B32\endcsname{\href{S6_SERIES_GRIND_I.pdf\#nameddest=tag-0306}{I.15.7 (p.~134)} }
\expandafter\def\csname SRef@push:src:B.3.3\endcsname{\href{S6_SERIES_GRIND_I.pdf\#nameddest=tag-0307}{I.15.8 (p.~134)} }
\expandafter\def\csname SRef@push:B33\endcsname{\href{S6_SERIES_GRIND_I.pdf\#nameddest=tag-0307}{I.15.8 (p.~134)} }
\expandafter\def\csname SRef@push:src:D.1.1\endcsname{\href{S6_SERIES_GRIND_I.pdf\#nameddest=tag-032A}{I.16.1 (p.~227)} }
\expandafter\def\csname SRef@push:src:D.4.1\endcsname{\href{S6_SERIES_GRIND_I.pdf\#nameddest=tag-032J}{I.16.10 (p.~235)} }
\expandafter\def\csname SRef@push:src:D.4.2\endcsname{\href{S6_SERIES_GRIND_I.pdf\#nameddest=tag-032K}{I.16.11 (p.~236)} }
\expandafter\def\csname SRef@push:src:D.4.3\endcsname{\href{S6_SERIES_GRIND_I.pdf\#nameddest=tag-032L}{I.16.12 (p.~237)} }
\expandafter\def\csname SRef@push:src:D.4.4\endcsname{\href{S6_SERIES_GRIND_I.pdf\#nameddest=tag-032M}{I.16.13 (p.~238)} }
\expandafter\def\csname SRef@push:src:D.4.5\endcsname{\href{S6_SERIES_GRIND_I.pdf\#nameddest=tag-032N}{I.16.14 (p.~238)} }
\expandafter\def\csname SRef@push:src:D.4.7\endcsname{\href{S6_SERIES_GRIND_I.pdf\#nameddest=tag-032P}{I.16.16 (p.~241)} }
\expandafter\def\csname SRef@push:src:D.4.8\endcsname{\href{S6_SERIES_GRIND_I.pdf\#nameddest=tag-032Q}{I.16.17 (p.~242)} }
\expandafter\def\csname SRef@push:src:D.2.2\endcsname{\href{S6_SERIES_GRIND_I.pdf\#nameddest=tag-032E}{I.16.5 (p.~231)} }
\expandafter\def\csname SRef@push:src:D.2.4\endcsname{\href{S6_SERIES_GRIND_I.pdf\#nameddest=tag-032G}{I.16.7 (p.~232)} }
\expandafter\def\csname SRef@inv:A3.2\endcsname{\href{S6_SERIES_GRIND_I.pdf\#nameddest=tag-01RZ}{I.18.10 (p.~264)} }
\expandafter\def\csname SRef@wb:lem:bcpr-eta-extension\endcsname{\href{S6_SERIES_GRIND_I.pdf\#nameddest=tag-01RZ}{I.18.10 (p.~264)} }
\expandafter\def\csname SRef@inv:A4.1\endcsname{\href{S6_SERIES_GRIND_I.pdf\#nameddest=tag-01S0}{I.18.11 (p.~264)} }
\expandafter\def\csname SRef@inv:A4.2\endcsname{\href{S6_SERIES_GRIND_I.pdf\#nameddest=tag-01S1}{I.18.12 (p.~265)} }
\expandafter\def\csname SRef@inv:A1.1\endcsname{\href{S6_SERIES_GRIND_I.pdf\#nameddest=tag-01RR}{I.18.2 (p.~254)} }
\expandafter\def\csname SRef@inv:A1.2\endcsname{\href{S6_SERIES_GRIND_I.pdf\#nameddest=tag-01RS}{I.18.3 (p.~255)} }
\expandafter\def\csname SRef@wb:prop:log-relative-direct-images\endcsname{\href{S6_SERIES_GRIND_I.pdf\#nameddest=tag-01RS}{I.18.3 (p.~255)} }
\expandafter\def\csname SRef@inv:A1.3\endcsname{\href{S6_SERIES_GRIND_I.pdf\#nameddest=tag-01RT}{I.18.4 (p.~256)} }
\expandafter\def\csname SRef@wb:prop:absolute-cotangent-direct-images\endcsname{\href{S6_SERIES_GRIND_I.pdf\#nameddest=tag-01RT}{I.18.4 (p.~256)} }
\expandafter\def\csname SRef@inv:A2.1\endcsname{\href{S6_SERIES_GRIND_I.pdf\#nameddest=tag-01RU}{I.18.5 (p.~257)} }
\expandafter\def\csname SRef@wb:prop:log-ks-divisor\endcsname{\href{S6_SERIES_GRIND_I.pdf\#nameddest=tag-01RU}{I.18.5 (p.~257)} }
\expandafter\def\csname SRef@wb:prop:period-derivative-determinant\endcsname{\href{S6_SERIES_GRIND_I.pdf\#nameddest=tag-01RU}{I.18.5 (p.~257)} }
\expandafter\def\csname SRef@inv:A2.2\endcsname{\href{S6_SERIES_GRIND_I.pdf\#nameddest=tag-01RV}{I.18.6 (p.~258)} }
\expandafter\def\csname SRef@inv:A2.3\endcsname{\href{S6_SERIES_GRIND_I.pdf\#nameddest=tag-01RW}{I.18.7 (p.~260)} }
\expandafter\def\csname SRef@inv:A2.4\endcsname{\href{S6_SERIES_GRIND_I.pdf\#nameddest=tag-01RX}{I.18.8 (p.~262)} }
\expandafter\def\csname SRef@wb:thm:h12-scalar\endcsname{\href{S6_SERIES_GRIND_I.pdf\#nameddest=tag-01RX}{I.18.8 (p.~262)} }
\expandafter\def\csname SRef@inv:A3.1\endcsname{\href{S6_SERIES_GRIND_I.pdf\#nameddest=tag-01RY}{I.18.9 (p.~263)} }
\expandafter\def\csname SRef@wb:cor:exact-hodge-diamond\endcsname{\href{S6_SERIES_GRIND_I.pdf\#nameddest=tag-01RY}{I.18.9 (p.~263)} }
\expandafter\def\csname SRef@wb:thm:parametric-hodge-diamond\endcsname{\href{S6_SERIES_GRIND_I.pdf\#nameddest=tag-01RY}{I.18.9 (p.~263)} }
\expandafter\def\csname SRef@wb:prop:workbench-holomorphic-form-vanishing\endcsname{\href{S6_SERIES_GRIND_I.pdf\#nameddest=tag-01RY}{I.18.9 (p.~263)} }
\expandafter\def\csname SRef@inv:C1.1\endcsname{\href{S6_SERIES_GRIND_I.pdf\#nameddest=tag-01SA}{I.19.1 (p.~266)} }
\expandafter\def\csname SRef@wb:thm:frolicher-e2-candidate\endcsname{\href{S6_SERIES_GRIND_I.pdf\#nameddest=tag-01SA}{I.19.1 (p.~266)} }
\expandafter\def\csname SRef@inv:C3.7\endcsname{\href{S6_SERIES_GRIND_I.pdf\#nameddest=tag-01SK}{I.19.11 (p.~274)} }
\expandafter\def\csname SRef@inv:C1.2\endcsname{\href{S6_SERIES_GRIND_I.pdf\#nameddest=tag-01SB}{I.19.2 (p.~267)} }
\expandafter\def\csname SRef@inv:C2.1\endcsname{\href{S6_SERIES_GRIND_I.pdf\#nameddest=tag-01SC}{I.19.3 (p.~267)} }
\expandafter\def\csname SRef@wb:thm:exact-bott-chern-aeppli\endcsname{\href{S6_SERIES_GRIND_I.pdf\#nameddest=tag-01SC}{I.19.3 (p.~267)} }
\expandafter\def\csname SRef@inv:C2.2\endcsname{\href{S6_SERIES_GRIND_I.pdf\#nameddest=tag-01SD}{I.19.4 (p.~268)} }
\expandafter\def\csname SRef@inv:C3.1\endcsname{\href{S6_SERIES_GRIND_I.pdf\#nameddest=tag-01SE}{I.19.5 (p.~269)} }
\expandafter\def\csname SRef@wb:cor:picard-line-audit\endcsname{\href{S6_SERIES_GRIND_I.pdf\#nameddest=tag-01SE}{I.19.5 (p.~269)} }
\expandafter\def\csname SRef@inv:C3.2\endcsname{\href{S6_SERIES_GRIND_I.pdf\#nameddest=tag-01SF}{I.19.6 (p.~270)} }
\expandafter\def\csname SRef@wb:lem:bcpr-degree-one-input\endcsname{\href{S6_SERIES_GRIND_I.pdf\#nameddest=tag-01SF}{I.19.6 (p.~270)} }
\expandafter\def\csname SRef@inv:C3.3\endcsname{\href{S6_SERIES_GRIND_I.pdf\#nameddest=tag-01SG}{I.19.7 (p.~272)} }
\expandafter\def\csname SRef@wb:thm:bcpr-square-zero-algebra\endcsname{\href{S6_SERIES_GRIND_I.pdf\#nameddest=tag-01SG}{I.19.7 (p.~272)} }
\expandafter\def\csname SRef@wb:prop:bcpr-generators-are-bases\endcsname{\href{S6_SERIES_GRIND_I.pdf\#nameddest=tag-01SG}{I.19.7 (p.~272)} }
\expandafter\def\csname SRef@inv:C3.4\endcsname{\href{S6_SERIES_GRIND_I.pdf\#nameddest=tag-01SH}{I.19.8 (p.~273)} }
\expandafter\def\csname SRef@inv:C3.5\endcsname{\href{S6_SERIES_GRIND_I.pdf\#nameddest=tag-01SI}{I.19.9 (p.~273)} }
\expandafter\def\csname SRef@inv:D4.1\endcsname{\href{S6_SERIES_GRIND_I.pdf\#nameddest=tag-01SO}{I.20.1 (p.~275)} }
\expandafter\def\csname SRef@inv:D4.2\endcsname{\href{S6_SERIES_GRIND_I.pdf\#nameddest=tag-01SP}{I.20.2 (p.~276)} }
\expandafter\def\csname SRef@inv:D4.3\endcsname{\href{S6_SERIES_GRIND_I.pdf\#nameddest=tag-01SQ}{I.20.3 (p.~276)} }
\expandafter\def\csname SRef@inv:D4.4\endcsname{\href{S6_SERIES_GRIND_I.pdf\#nameddest=tag-01SR}{I.20.4 (p.~277)} }
\expandafter\def\csname SRef@inv:D4.5\endcsname{\href{S6_SERIES_GRIND_I.pdf\#nameddest=tag-01SS}{I.20.5 (p.~280)} }
\expandafter\def\csname SRef@push:src:E.2.2\endcsname{\href{S6_SERIES_GRIND_I.pdf\#nameddest=tag-0335}{I.20.7 (p.~284)} }
\expandafter\def\csname SRef@push:src:E.2.3\endcsname{\href{S6_SERIES_GRIND_I.pdf\#nameddest=tag-033B}{I.20.8 (p.~284)} }
\expandafter\def\csname SRef@inv:B1.0\endcsname{\href{S6_SERIES_GRIND_I.pdf\#nameddest=tag-01S2}{I.21.1 (p.~286)} }
\expandafter\def\csname SRef@inv:B1.1\endcsname{\href{S6_SERIES_GRIND_I.pdf\#nameddest=tag-01S3}{I.21.2 (p.~287)} }
\expandafter\def\csname SRef@inv:B1.2\endcsname{\href{S6_SERIES_GRIND_I.pdf\#nameddest=tag-01S4}{I.21.3 (p.~288)} }
\expandafter\def\csname SRef@inv:B2.1\endcsname{\href{S6_SERIES_GRIND_I.pdf\#nameddest=tag-01S5}{I.21.4 (p.~288)} }
\expandafter\def\csname SRef@inv:B2.2\endcsname{\href{S6_SERIES_GRIND_I.pdf\#nameddest=tag-01S6}{I.21.5 (p.~289)} }
\expandafter\def\csname SRef@inv:B3.1\endcsname{\href{S6_SERIES_GRIND_I.pdf\#nameddest=tag-01S7}{I.21.6 (p.~290)} }
\expandafter\def\csname SRef@inv:B3.2\endcsname{\href{S6_SERIES_GRIND_I.pdf\#nameddest=tag-01S8}{I.21.7 (p.~292)} }
\expandafter\def\csname SRef@inv:B3.3\endcsname{\href{S6_SERIES_GRIND_I.pdf\#nameddest=tag-01S9}{I.21.8 (p.~292)} }
\expandafter\def\csname SRef@mapsI:src:A.1.1\endcsname{\href{S6_SERIES_GRIND_I.pdf\#nameddest=tag-02PA}{I.22.1 (p.~295)} }
\expandafter\def\csname SRef@mapsI:quotient-eq-1\endcsname{\href{S6_SERIES_GRIND_I.pdf\#nameddest=tag-02PA}{I.22.1 (p.~295)} }
\expandafter\def\csname SRef@mapsI:src:A.1.2\endcsname{\href{S6_SERIES_GRIND_I.pdf\#nameddest=tag-02PB}{I.22.2 (p.~296)} }
\expandafter\def\csname SRef@mapsI:src:A.1.3\endcsname{\href{S6_SERIES_GRIND_I.pdf\#nameddest=tag-02PC}{I.22.3 (p.~297)} }
\expandafter\def\csname SRef@mapsI:quotient-eq-6\endcsname{\href{S6_SERIES_GRIND_I.pdf\#nameddest=tag-02PC}{I.22.3 (p.~298)} }
\expandafter\def\csname SRef@ari:C.2.1\endcsname{\href{S6_SERIES_GRIND_I.pdf\#nameddest=tag-000L}{I.22.5 (p.~299)} }
\expandafter\def\csname SRef@ari:C.2.2\endcsname{\href{S6_SERIES_GRIND_I.pdf\#nameddest=tag-000M}{I.22.6 (p.~300)} }
\expandafter\def\csname SRef@ari:C.3.1\endcsname{\href{S6_SERIES_GRIND_I.pdf\#nameddest=tag-000N}{I.22.7 (p.~301)} }
\expandafter\def\csname SRef@ari:C.3.2\endcsname{\href{S6_SERIES_GRIND_I.pdf\#nameddest=tag-000O}{I.22.8 (p.~302)} }
\expandafter\def\csname SRef@ari:C.6.1\endcsname{\href{S6_SERIES_GRIND_I.pdf\#nameddest=tag-000V}{I.23.1 (p.~304)} }
\expandafter\def\csname SRef@ari:C.6.R5\endcsname{\href{S6_SERIES_GRIND_I.pdf\#nameddest=tag-0014}{I.23.10 (p.~316)} }
\expandafter\def\csname SRef@ari:C.6.T1\endcsname{\href{S6_SERIES_GRIND_I.pdf\#nameddest=tag-0015}{I.23.11 (p.~322)} }
\expandafter\def\csname SRef@ari:C.6.T2\endcsname{\href{S6_SERIES_GRIND_I.pdf\#nameddest=tag-0016}{I.23.12 (p.~323)} }
\expandafter\def\csname SRef@ari:C.8.6\endcsname{\href{S6_SERIES_GRIND_I.pdf\#nameddest=tag-001T}{I.23.13 (p.~325)} }
\expandafter\def\csname SRef@ari:C.8.7\endcsname{\href{S6_SERIES_GRIND_I.pdf\#nameddest=tag-001U}{I.23.14 (p.~326)} }
\expandafter\def\csname SRef@ari:C.8.8\endcsname{\href{S6_SERIES_GRIND_I.pdf\#nameddest=tag-001V}{I.23.15 (p.~326)} }
\expandafter\def\csname SRef@ari:C.8.9\endcsname{\href{S6_SERIES_GRIND_I.pdf\#nameddest=tag-001W}{I.23.16 (p.~327)} }
\expandafter\def\csname SRef@ari:C.8.10\endcsname{\href{S6_SERIES_GRIND_I.pdf\#nameddest=tag-001K}{I.23.17 (p.~329)} }
\expandafter\def\csname SRef@ari:C.8.11\endcsname{\href{S6_SERIES_GRIND_I.pdf\#nameddest=tag-001L}{I.23.18 (p.~330)} }
\expandafter\def\csname SRef@ari:C.6.2\endcsname{\href{S6_SERIES_GRIND_I.pdf\#nameddest=tag-000W}{I.23.2 (p.~305)} }
\expandafter\def\csname SRef@ari:C.6.3\endcsname{\href{S6_SERIES_GRIND_I.pdf\#nameddest=tag-000X}{I.23.3 (p.~305)} }
\expandafter\def\csname SRef@ari:C.6.4\endcsname{\href{S6_SERIES_GRIND_I.pdf\#nameddest=tag-000Y}{I.23.4 (p.~308)} }
\expandafter\def\csname SRef@ari:C.6.5\endcsname{\href{S6_SERIES_GRIND_I.pdf\#nameddest=tag-000Z}{I.23.5 (p.~309)} }
\expandafter\def\csname SRef@ari:C.6.R1\endcsname{\href{S6_SERIES_GRIND_I.pdf\#nameddest=tag-0010}{I.23.6 (p.~310)} }
\expandafter\def\csname SRef@ari:C.6.R2\endcsname{\href{S6_SERIES_GRIND_I.pdf\#nameddest=tag-0011}{I.23.7 (p.~311)} }
\expandafter\def\csname SRef@ari:C.6.R3\endcsname{\href{S6_SERIES_GRIND_I.pdf\#nameddest=tag-0012}{I.23.8 (p.~314)} }
\expandafter\def\csname SRef@ari:C.8.12\endcsname{\href{S6_SERIES_GRIND_I.pdf\#nameddest=tag-001M}{I.24.1 (p.~331)} }
\expandafter\def\csname SRef@ari:C.8.13\endcsname{\href{S6_SERIES_GRIND_I.pdf\#nameddest=tag-001N}{I.24.2 (p.~332)} }
\expandafter\def\csname SRef@ari:C.8.14\endcsname{\href{S6_SERIES_GRIND_I.pdf\#nameddest=tag-001O}{I.24.3 (p.~332)} }
\expandafter\def\csname SRef@ari:C.8.15\endcsname{\href{S6_SERIES_GRIND_I.pdf\#nameddest=tag-001P}{I.24.4 (p.~337)} }
\expandafter\def\csname SRef@fam:src:B.1.1\endcsname{\href{S6_SERIES_GRIND_I.pdf\#nameddest=tag-01Q7}{I.3.2 (p.~24)} }
\expandafter\def\csname SRef@fam:st:B.1.1\endcsname{\href{S6_SERIES_GRIND_I.pdf\#nameddest=tag-01Q7}{I.3.2 (p.~24)} }
\expandafter\def\csname SRef@fam:src:B.1.2\endcsname{\href{S6_SERIES_GRIND_I.pdf\#nameddest=tag-01Q8}{I.3.3 (p.~26)} }
\expandafter\def\csname SRef@fam:st:B.1.2\endcsname{\href{S6_SERIES_GRIND_I.pdf\#nameddest=tag-01Q8}{I.3.3 (p.~26)} }
\expandafter\def\csname SRef@fam:src:B.1.3\endcsname{\href{S6_SERIES_GRIND_I.pdf\#nameddest=tag-01Q9}{I.3.4 (p.~26)} }
\expandafter\def\csname SRef@fam:st:B.1.3\endcsname{\href{S6_SERIES_GRIND_I.pdf\#nameddest=tag-01Q9}{I.3.4 (p.~26)} }
\expandafter\def\csname SRef@fam:src:B.1.4\endcsname{\href{S6_SERIES_GRIND_I.pdf\#nameddest=tag-01QA}{I.3.5 (p.~28)} }
\expandafter\def\csname SRef@fam:st:B.1.4\endcsname{\href{S6_SERIES_GRIND_I.pdf\#nameddest=tag-01QA}{I.3.5 (p.~28)} }
\expandafter\def\csname SRef@fam:src:C.1.1\endcsname{\href{S6_SERIES_GRIND_I.pdf\#nameddest=tag-01QH}{I.4.1 (p.~29)} }
\expandafter\def\csname SRef@fam:st:C.1.1\endcsname{\href{S6_SERIES_GRIND_I.pdf\#nameddest=tag-01QH}{I.4.1 (p.~29)} }
\expandafter\def\csname SRef@fam:src:C.1.2\endcsname{\href{S6_SERIES_GRIND_I.pdf\#nameddest=tag-01QI}{I.4.2 (p.~30)} }
\expandafter\def\csname SRef@fam:st:C.1.2\endcsname{\href{S6_SERIES_GRIND_I.pdf\#nameddest=tag-01QI}{I.4.2 (p.~30)} }
\expandafter\def\csname SRef@fam:src:C.2.1\endcsname{\href{S6_SERIES_GRIND_I.pdf\#nameddest=tag-01QJ}{I.4.3 (p.~31)} }
\expandafter\def\csname SRef@fam:st:C.2.1\endcsname{\href{S6_SERIES_GRIND_I.pdf\#nameddest=tag-01QJ}{I.4.3 (p.~31)} }
\expandafter\def\csname SRef@fam:src:C.2.2\endcsname{\href{S6_SERIES_GRIND_I.pdf\#nameddest=tag-01QK}{I.4.4 (p.~32)} }
\expandafter\def\csname SRef@fam:st:C.2.2\endcsname{\href{S6_SERIES_GRIND_I.pdf\#nameddest=tag-01QK}{I.4.4 (p.~32)} }
\expandafter\def\csname SRef@ari:C.1.1\endcsname{\href{S6_SERIES_GRIND_I.pdf\#nameddest=tag-000I}{I.5.1 (p.~36)} }
\expandafter\def\csname SRef@ari:C.1.2\endcsname{\href{S6_SERIES_GRIND_I.pdf\#nameddest=tag-000J}{I.5.2 (p.~37)} }
\expandafter\def\csname SRef@ari:C.1.3\endcsname{\href{S6_SERIES_GRIND_I.pdf\#nameddest=tag-000K}{I.5.3 (p.~37)} }
\expandafter\def\csname SRef@ari:C.4.2\endcsname{\href{S6_SERIES_GRIND_I.pdf\#nameddest=tag-000Q}{I.5.5 (p.~39)} }
\expandafter\def\csname SRef@ari:C.4.3\endcsname{\href{S6_SERIES_GRIND_I.pdf\#nameddest=tag-000R}{I.5.6 (p.~39)} }
\expandafter\def\csname SRef@ari:C.4.4\endcsname{\href{S6_SERIES_GRIND_I.pdf\#nameddest=tag-000S}{I.5.7 (p.~40)} }
\expandafter\def\csname SRef@ari:C.4.5\endcsname{\href{S6_SERIES_GRIND_I.pdf\#nameddest=tag-000T}{I.5.8 (p.~42)} }
\expandafter\def\csname SRef@ari:C.8.1\endcsname{\href{S6_SERIES_GRIND_I.pdf\#nameddest=tag-001J}{I.6.1 (p.~42)} }
\expandafter\def\csname SRef@ari:C.8.2\endcsname{\href{S6_SERIES_GRIND_I.pdf\#nameddest=tag-001Q}{I.6.2 (p.~44)} }
\expandafter\def\csname SRef@ari:C.8.3\endcsname{\href{S6_SERIES_GRIND_I.pdf\#nameddest=tag-001R}{I.6.3 (p.~46)} }
\expandafter\def\csname SRef@ari:C.8.4\endcsname{\href{S6_SERIES_GRIND_I.pdf\#nameddest=tag-001S}{I.6.4 (p.~47)} }
\expandafter\def\csname SRef@fam:src:A.1.1\endcsname{\href{S6_SERIES_GRIND_I.pdf\#nameddest=tag-01PZ}{I.7.1 (p.~49)} }
\expandafter\def\csname SRef@fam:st:A.1.1\endcsname{\href{S6_SERIES_GRIND_I.pdf\#nameddest=tag-01PZ}{I.7.1 (p.~49)} }
\expandafter\def\csname SRef@fam:src:A.1.2\endcsname{\href{S6_SERIES_GRIND_I.pdf\#nameddest=tag-01Q0}{I.7.2 (p.~49)} }
\expandafter\def\csname SRef@fam:st:A.1.2\endcsname{\href{S6_SERIES_GRIND_I.pdf\#nameddest=tag-01Q0}{I.7.2 (p.~49)} }
\expandafter\def\csname SRef@fam:src:A.2.1\endcsname{\href{S6_SERIES_GRIND_I.pdf\#nameddest=tag-01Q1}{I.7.3 (p.~50)} }
\expandafter\def\csname SRef@fam:st:A.2.1\endcsname{\href{S6_SERIES_GRIND_I.pdf\#nameddest=tag-01Q1}{I.7.3 (p.~50)} }
\expandafter\def\csname SRef@fam:src:A.3.1\endcsname{\href{S6_SERIES_GRIND_I.pdf\#nameddest=tag-01Q2}{I.7.4 (p.~51)} }
\expandafter\def\csname SRef@fam:st:A.3.1\endcsname{\href{S6_SERIES_GRIND_I.pdf\#nameddest=tag-01Q2}{I.7.4 (p.~51)} }
\expandafter\def\csname SRef@fam:src:A.4.1\endcsname{\href{S6_SERIES_GRIND_I.pdf\#nameddest=tag-01Q3}{I.7.5 (p.~53)} }
\expandafter\def\csname SRef@fam:st:A.4.1\endcsname{\href{S6_SERIES_GRIND_I.pdf\#nameddest=tag-01Q3}{I.7.5 (p.~53)} }
\expandafter\def\csname SRef@fam:src:A.4.2\endcsname{\href{S6_SERIES_GRIND_I.pdf\#nameddest=tag-01Q4}{I.7.6 (p.~54)} }
\expandafter\def\csname SRef@fam:st:A.4.2\endcsname{\href{S6_SERIES_GRIND_I.pdf\#nameddest=tag-01Q4}{I.7.6 (p.~54)} }
\expandafter\def\csname SRef@fam:src:A.4.3\endcsname{\href{S6_SERIES_GRIND_I.pdf\#nameddest=tag-01Q5}{I.7.7 (p.~55)} }
\expandafter\def\csname SRef@fam:st:A.4.3\endcsname{\href{S6_SERIES_GRIND_I.pdf\#nameddest=tag-01Q5}{I.7.7 (p.~55)} }
\expandafter\def\csname SRef@fam:src:B.2.1\endcsname{\href{S6_SERIES_GRIND_I.pdf\#nameddest=tag-01QB}{I.8.1 (p.~57)} }
\expandafter\def\csname SRef@fam:st:B.2.1\endcsname{\href{S6_SERIES_GRIND_I.pdf\#nameddest=tag-01QB}{I.8.1 (p.~57)} }
\expandafter\def\csname SRef@fam:src:B.2.2\endcsname{\href{S6_SERIES_GRIND_I.pdf\#nameddest=tag-01QC}{I.8.2 (p.~59)} }
\expandafter\def\csname SRef@fam:st:B.2.2\endcsname{\href{S6_SERIES_GRIND_I.pdf\#nameddest=tag-01QC}{I.8.2 (p.~59)} }
\expandafter\def\csname SRef@fam:src:B.2.3\endcsname{\href{S6_SERIES_GRIND_I.pdf\#nameddest=tag-01QD}{I.8.3 (p.~60)} }
\expandafter\def\csname SRef@fam:st:B.2.3\endcsname{\href{S6_SERIES_GRIND_I.pdf\#nameddest=tag-01QD}{I.8.3 (p.~60)} }
\expandafter\def\csname SRef@fam:src:B.3.1\endcsname{\href{S6_SERIES_GRIND_I.pdf\#nameddest=tag-01QE}{I.8.4 (p.~60)} }
\expandafter\def\csname SRef@fam:st:B.3.1\endcsname{\href{S6_SERIES_GRIND_I.pdf\#nameddest=tag-01QE}{I.8.4 (p.~60)} }
\expandafter\def\csname SRef@fam:src:B.4.1\endcsname{\href{S6_SERIES_GRIND_I.pdf\#nameddest=tag-01QF}{I.8.5 (p.~62)} }
\expandafter\def\csname SRef@fam:st:B.4.1\endcsname{\href{S6_SERIES_GRIND_I.pdf\#nameddest=tag-01QF}{I.8.5 (p.~62)} }
\expandafter\def\csname SRef@fam:src:B.4.2\endcsname{\href{S6_SERIES_GRIND_I.pdf\#nameddest=tag-01QG}{I.8.6 (p.~62)} }
\expandafter\def\csname SRef@fam:st:B.4.2\endcsname{\href{S6_SERIES_GRIND_I.pdf\#nameddest=tag-01QG}{I.8.6 (p.~62)} }
\expandafter\def\csname SRef@push:src:E.1.1\endcsname{\href{S6_SERIES_GRIND_I.pdf\#nameddest=tag-032R}{I.9.1 (p.~64)} }
\expandafter\def\csname SRef@inv:D1.1\endcsname{\href{S6_SERIES_GRIND_I.pdf\#nameddest=tag-01SM}{I.B.1 (p.~360)} }
\expandafter\def\csname SRef@inv:D1.2\endcsname{\href{S6_SERIES_GRIND_I.pdf\#nameddest=tag-01SN}{I.B.2 (p.~361)} }
\expandafter\def\csname SRef@ari:C.7.1\endcsname{\href{S6_SERIES_GRIND_I.pdf\#nameddest=tag-0017}{I.E.1 (p.~383)} }
\expandafter\def\csname SRef@ari:C.7.11\endcsname{\href{S6_SERIES_GRIND_I.pdf\#nameddest=tag-0019}{I.E.10 (p.~398)} }
\expandafter\def\csname SRef@ari:C.7.12\endcsname{\href{S6_SERIES_GRIND_I.pdf\#nameddest=tag-001A}{I.E.11 (p.~400)} }
\expandafter\def\csname SRef@ari:C.7.2\endcsname{\href{S6_SERIES_GRIND_I.pdf\#nameddest=tag-001B}{I.E.2 (p.~384)} }
\expandafter\def\csname SRef@ari:C.7.3\endcsname{\href{S6_SERIES_GRIND_I.pdf\#nameddest=tag-001C}{I.E.3 (p.~385)} }
\expandafter\def\csname SRef@ari:C.7.4\endcsname{\href{S6_SERIES_GRIND_I.pdf\#nameddest=tag-001D}{I.E.4 (p.~386)} }
\expandafter\def\csname SRef@ari:C.7.6\endcsname{\href{S6_SERIES_GRIND_I.pdf\#nameddest=tag-001F}{I.E.5 (p.~387)} }
\expandafter\def\csname SRef@ari:C.7.9\endcsname{\href{S6_SERIES_GRIND_I.pdf\#nameddest=tag-001I}{I.E.8 (p.~391)} }
\expandafter\def\csname SRef@ari:C.7.10\endcsname{\href{S6_SERIES_GRIND_I.pdf\#nameddest=tag-0018}{I.E.9 (p.~394)} }
\expandafter\def\csname SRef@inv:C3.8\endcsname{\href{S6_SERIES_GRIND_I.pdf\#nameddest=tag-01SL}{I.D.1 (p.~374)} }
\expandafter\def\csname SRef@fam:src:C.3.3\endcsname{\href{S6_SERIES_GRIND_I.pdf\#nameddest=tag-01QO}{I.11.3 (p.~86)} }
\expandafter\def\csname SRef@fam:st:C.3.3\endcsname{\href{S6_SERIES_GRIND_I.pdf\#nameddest=tag-01QO}{I.11.3 (p.~86)} }
\expandafter\def\csname SRef@fam:src:E.1.5\endcsname{\href{S6_SERIES_GRIND_I.pdf\#nameddest=tag-01R5}{I.13.1 (p.~100)} }
\expandafter\def\csname SRef@fam:st:E.1.5\endcsname{\href{S6_SERIES_GRIND_I.pdf\#nameddest=tag-01R5}{I.13.1 (p.~100)} }
\expandafter\def\csname SRef@wb:rem:bcpr-literal-gauge\endcsname{\href{S6_SERIES_GRIND_I.pdf\#nameddest=tag-03CT}{I.D.3 (p.~382)} }
\expandafter\def\csname SRef@wb:rem:bcpr-global-gauge-boundary\endcsname{\href{S6_SERIES_GRIND_I.pdf\#nameddest=tag-03CS}{I.D.3 (p.~382)} }
\expandafter\def\csname SRef@push:src:A.2.3\endcsname{\href{S6_SERIES_GRIND_I.pdf\#nameddest=tag-02ZL}{I.14.2 (p.~107)} }
\expandafter\def\csname SRef@push:A23\endcsname{\href{S6_SERIES_GRIND_I.pdf\#nameddest=tag-02ZL}{I.14.2 (p.~107)} }
\expandafter\def\csname SRef@fam:section:F\endcsname{\href{S6_SERIES_GRIND_I.pdf\#nameddest=tag-01PV}{I.10 (p.~68)} }
\expandafter\def\csname SRef@fam:section:D\endcsname{\href{S6_SERIES_GRIND_I.pdf\#nameddest=tag-01PP}{I.11 (p.~76)} }
\expandafter\def\csname SRef@mapsI:real-forms\endcsname{\href{S6_SERIES_GRIND_I.pdf\#nameddest=tag-02P9}{I.12 (p.~88)} }
\expandafter\def\csname SRef@fam:section:E\endcsname{\href{S6_SERIES_GRIND_I.pdf\#nameddest=tag-01PS}{I.13 (p.~96)} }
\expandafter\def\csname SRef@push:section:A\endcsname{\href{S6_SERIES_GRIND_I.pdf\#nameddest=tag-02YZ}{I.14 (p.~103)} }
\expandafter\def\csname SRef@push:section:B\endcsname{\href{S6_SERIES_GRIND_I.pdf\#nameddest=tag-02Z4}{I.15 (p.~128)} }
\expandafter\def\csname SRef@push:section:D\endcsname{\href{S6_SERIES_GRIND_I.pdf\#nameddest=tag-02Z8}{I.16 (p.~227)} }
\expandafter\def\csname SRef@inv:sec:partA\endcsname{\href{S6_SERIES_GRIND_I.pdf\#nameddest=tag-01TH}{I.18 (p.~254)} }
\expandafter\def\csname SRef@inv:sec:partC\endcsname{\href{S6_SERIES_GRIND_I.pdf\#nameddest=tag-01TI}{I.19 (p.~266)} }
\expandafter\def\csname SRef@def:common-conventions\endcsname{\href{S6_SERIES_GRIND_I.pdf\#nameddest=tag-01O1}{I.2 (p.~21)} }
\expandafter\def\csname SRef@mapsI:circle-quotient\endcsname{\href{S6_SERIES_GRIND_I.pdf\#nameddest=tag-02OO}{I.22 (p.~295)} }
\expandafter\def\csname SRef@inv:notation-dictionary\endcsname{\href{S6_SERIES_GRIND_I.pdf\#nameddest=tag-01TG}{I.3 (p.~21)} }
\expandafter\def\csname SRef@fam:section:C\endcsname{\href{S6_SERIES_GRIND_I.pdf\#nameddest=tag-01PK}{I.4 (p.~29)} }
\expandafter\def\csname SRef@fam:section:A\endcsname{\href{S6_SERIES_GRIND_I.pdf\#nameddest=tag-01P9}{I.7 (p.~48)} }
\expandafter\def\csname SRef@fam:section:B\endcsname{\href{S6_SERIES_GRIND_I.pdf\#nameddest=tag-01PF}{I.8 (p.~57)} }
\expandafter\def\csname SRef@push:section:E\endcsname{\href{S6_SERIES_GRIND_I.pdf\#nameddest=tag-02ZD}{I.9 (p.~64)} }
\expandafter\def\csname SRef@fam:section:F.1\endcsname{\href{S6_SERIES_GRIND_I.pdf\#nameddest=tag-01PW}{I.10.1 (p.~68)} }
\expandafter\def\csname SRef@fam:section:F.2\endcsname{\href{S6_SERIES_GRIND_I.pdf\#nameddest=tag-01PX}{I.10.2 (p.~71)} }
\expandafter\def\csname SRef@fam:section:F.3\endcsname{\href{S6_SERIES_GRIND_I.pdf\#nameddest=tag-01PY}{I.10.3 (p.~74)} }
\expandafter\def\csname SRef@fam:section:D.1\endcsname{\href{S6_SERIES_GRIND_I.pdf\#nameddest=tag-01PQ}{I.11.1 (p.~76)} }
\expandafter\def\csname SRef@fam:section:D.2\endcsname{\href{S6_SERIES_GRIND_I.pdf\#nameddest=tag-01PR}{I.11.2 (p.~82)} }
\expandafter\def\csname SRef@fam:section:C.3\endcsname{\href{S6_SERIES_GRIND_I.pdf\#nameddest=tag-01PN}{I.11.3 (p.~84)} }
\expandafter\def\csname SRef@fam:section:C.4\endcsname{\href{S6_SERIES_GRIND_I.pdf\#nameddest=tag-01PO}{I.11.4 (p.~86)} }
\expandafter\def\csname SRef@fam:section:E.1\endcsname{\href{S6_SERIES_GRIND_I.pdf\#nameddest=tag-01PT}{I.13.1 (p.~96)} }
\expandafter\def\csname SRef@fam:section:E.2\endcsname{\href{S6_SERIES_GRIND_I.pdf\#nameddest=tag-01PU}{I.13.2 (p.~101)} }
\expandafter\def\csname SRef@push:section:A.1\endcsname{\href{S6_SERIES_GRIND_I.pdf\#nameddest=tag-02Z0}{I.14.1 (p.~103)} }
\expandafter\def\csname SRef@push:section:A.2\endcsname{\href{S6_SERIES_GRIND_I.pdf\#nameddest=tag-02Z1}{I.14.2 (p.~105)} }
\expandafter\def\csname SRef@push:section:A.3\endcsname{\href{S6_SERIES_GRIND_I.pdf\#nameddest=tag-02Z2}{I.14.3 (p.~113)} }
\expandafter\def\csname SRef@push:section:A.4\endcsname{\href{S6_SERIES_GRIND_I.pdf\#nameddest=tag-02Z3}{I.14.4 (p.~116)} }
\expandafter\def\csname SRef@push:section:B.1\endcsname{\href{S6_SERIES_GRIND_I.pdf\#nameddest=tag-02Z5}{I.15.1 (p.~128)} }
\expandafter\def\csname SRef@push:section:B.2\endcsname{\href{S6_SERIES_GRIND_I.pdf\#nameddest=tag-02Z6}{I.15.2 (p.~131)} }
\expandafter\def\csname SRef@push:section:B.3\endcsname{\href{S6_SERIES_GRIND_I.pdf\#nameddest=tag-02Z7}{I.15.3 (p.~133)} }
\expandafter\def\csname SRef@push:section:D.1\endcsname{\href{S6_SERIES_GRIND_I.pdf\#nameddest=tag-02Z9}{I.16.1 (p.~227)} }
\expandafter\def\csname SRef@push:section:D.2\endcsname{\href{S6_SERIES_GRIND_I.pdf\#nameddest=tag-02ZA}{I.16.2 (p.~230)} }
\expandafter\def\csname SRef@push:section:D.3\endcsname{\href{S6_SERIES_GRIND_I.pdf\#nameddest=tag-02ZB}{I.16.3 (p.~233)} }
\expandafter\def\csname SRef@push:section:D.4\endcsname{\href{S6_SERIES_GRIND_I.pdf\#nameddest=tag-02ZC}{I.16.4 (p.~235)} }
\expandafter\def\csname SRef@push:section:E.2\endcsname{\href{S6_SERIES_GRIND_I.pdf\#nameddest=tag-02ZF}{I.20.2 (p.~282)} }
\expandafter\def\csname SRef@push:section:E.3\endcsname{\href{S6_SERIES_GRIND_I.pdf\#nameddest=tag-02ZG}{I.20.3 (p.~285)} }
\expandafter\def\csname SRef@sec:e4\endcsname{\href{S6_SERIES_GRIND_I.pdf\#nameddest=tag-034N}{I.24.1 (p.~338)} }
\expandafter\def\csname SRef@push:section:E.4\endcsname{\href{S6_SERIES_GRIND_I.pdf\#nameddest=tag-034N}{I.24.1 (p.~338)} }
\expandafter\def\csname SRef@fam:section:B.1\endcsname{\href{S6_SERIES_GRIND_I.pdf\#nameddest=tag-01PG}{I.3.1 (p.~24)} }
\expandafter\def\csname SRef@fam:section:C.1\endcsname{\href{S6_SERIES_GRIND_I.pdf\#nameddest=tag-01PL}{I.4.1 (p.~29)} }
\expandafter\def\csname SRef@fam:section:C.2\endcsname{\href{S6_SERIES_GRIND_I.pdf\#nameddest=tag-01PM}{I.4.2 (p.~31)} }
\expandafter\def\csname SRef@fam:section:A.1\endcsname{\href{S6_SERIES_GRIND_I.pdf\#nameddest=tag-01PA}{I.7.1 (p.~49)} }
\expandafter\def\csname SRef@fam:section:A.2\endcsname{\href{S6_SERIES_GRIND_I.pdf\#nameddest=tag-01PB}{I.7.2 (p.~50)} }
\expandafter\def\csname SRef@fam:section:A.3\endcsname{\href{S6_SERIES_GRIND_I.pdf\#nameddest=tag-01PC}{I.7.3 (p.~51)} }
\expandafter\def\csname SRef@fam:section:A.4\endcsname{\href{S6_SERIES_GRIND_I.pdf\#nameddest=tag-01PD}{I.7.4 (p.~53)} }
\expandafter\def\csname SRef@fam:section:A.5\endcsname{\href{S6_SERIES_GRIND_I.pdf\#nameddest=tag-01PE}{I.7.5 (p.~56)} }
\expandafter\def\csname SRef@fam:section:B.2\endcsname{\href{S6_SERIES_GRIND_I.pdf\#nameddest=tag-01PH}{I.8.1 (p.~57)} }
\expandafter\def\csname SRef@fam:section:B.3\endcsname{\href{S6_SERIES_GRIND_I.pdf\#nameddest=tag-01PI}{I.8.2 (p.~60)} }
\expandafter\def\csname SRef@fam:section:B.4\endcsname{\href{S6_SERIES_GRIND_I.pdf\#nameddest=tag-01PJ}{I.8.3 (p.~62)} }
\expandafter\def\csname SRef@push:section:E.1\endcsname{\href{S6_SERIES_GRIND_I.pdf\#nameddest=tag-02ZE}{I.9.1 (p.~64)} }
\expandafter\def\csname SRef@wb:sec:canonical-ring-recovery\endcsname{\href{S6_SERIES_GRIND_I.pdf\#nameddest=tag-03CU}{I.A.1 (p.~342)} }
\expandafter\def\csname SRef@wb:sec:retained-period-deformation\endcsname{\href{S6_SERIES_GRIND_I.pdf\#nameddest=tag-03CV}{I.A.2 (p.~344)} }
\expandafter\def\csname SRef@wb:pshift:sec:classification\endcsname{\href{S6_SERIES_GRIND_I.pdf\#nameddest=tag-03CR}{I.A.3 (p.~349)} }
\expandafter\def\csname SRef@app:rederive:fam-src-A.5.1\endcsname{\href{S6_SERIES_GRIND_I.pdf\#nameddest=tag-000C}{I.A.4 (p.~355)} }
\expandafter\def\csname SRef@app:rederive:fam-src-A.5.2\endcsname{\href{S6_SERIES_GRIND_I.pdf\#nameddest=tag-000D}{I.A.5 (p.~356)} }
\expandafter\def\csname SRef@app:rederive:fam-src-B.4.1\endcsname{\href{S6_SERIES_GRIND_I.pdf\#nameddest=tag-000E}{I.A.6 (p.~359)} }
\expandafter\def\csname SRef@fam:src:F.1.1\endcsname{\href{S6_SERIES_GRIND_I.pdf\#nameddest=tag-01R9}{I.10.1 (p.~68)} }
\expandafter\def\csname SRef@fam:st:F.1.1\endcsname{\href{S6_SERIES_GRIND_I.pdf\#nameddest=tag-01R9}{I.10.1 (p.~68)} }
\expandafter\def\csname SRef@cdp:st:D.1.14\endcsname{\href{S6_SERIES_GRIND_I.pdf\#nameddest=tag-01R9}{I.10.1 (p.~68)} }
\expandafter\def\csname SRef@fam:src:F.2.2\endcsname{\href{S6_SERIES_GRIND_I.pdf\#nameddest=tag-01RC}{I.10.4 (p.~72)} }
\expandafter\def\csname SRef@fam:st:F.2.2\endcsname{\href{S6_SERIES_GRIND_I.pdf\#nameddest=tag-01RC}{I.10.4 (p.~72)} }
\expandafter\def\csname SRef@fam:src:F.3.1\endcsname{\href{S6_SERIES_GRIND_I.pdf\#nameddest=tag-01RE}{I.10.6 (p.~74)} }
\expandafter\def\csname SRef@fam:st:F.3.1\endcsname{\href{S6_SERIES_GRIND_I.pdf\#nameddest=tag-01RE}{I.10.6 (p.~74)} }
\expandafter\def\csname SRef@fam:src:D.2.3\endcsname{\href{S6_SERIES_GRIND_I.pdf\#nameddest=tag-01R0}{I.11.10 (p.~83)} }
\expandafter\def\csname SRef@fam:st:D.2.3\endcsname{\href{S6_SERIES_GRIND_I.pdf\#nameddest=tag-01R0}{I.11.10 (p.~83)} }
\expandafter\def\csname SRef@fam:src:C.4.1\endcsname{\href{S6_SERIES_GRIND_I.pdf\#nameddest=tag-01QQ}{I.11.14 (p.~86)} }
\expandafter\def\csname SRef@fam:st:C.4.1\endcsname{\href{S6_SERIES_GRIND_I.pdf\#nameddest=tag-01QQ}{I.11.14 (p.~86)} }
\expandafter\def\csname SRef@fam:src:D.1.3\endcsname{\href{S6_SERIES_GRIND_I.pdf\#nameddest=tag-01QT}{I.11.3 (p.~78)} }
\expandafter\def\csname SRef@fam:st:D.1.3\endcsname{\href{S6_SERIES_GRIND_I.pdf\#nameddest=tag-01QT}{I.11.3 (p.~78)} }
\expandafter\def\csname SRef@fam:src:D.1.6\endcsname{\href{S6_SERIES_GRIND_I.pdf\#nameddest=tag-01QW}{I.11.6 (p.~80)} }
\expandafter\def\csname SRef@fam:st:D.1.6\endcsname{\href{S6_SERIES_GRIND_I.pdf\#nameddest=tag-01QW}{I.11.6 (p.~80)} }
\expandafter\def\csname SRef@mapsI:src:A.3.2\endcsname{\href{S6_SERIES_GRIND_I.pdf\#nameddest=tag-02PF}{I.12.2 (p.~89)} }
\expandafter\def\csname SRef@mapsI:real-eq-3\endcsname{\href{S6_SERIES_GRIND_I.pdf\#nameddest=tag-02PF}{I.12.2 (p.~89)} }
\expandafter\def\csname SRef@mapsI:src:A.3.3\endcsname{\href{S6_SERIES_GRIND_I.pdf\#nameddest=tag-02PG}{I.12.3 (p.~92)} }
\expandafter\def\csname SRef@push:src:A.4.7\endcsname{\href{S6_SERIES_GRIND_I.pdf\#nameddest=tag-02ZX}{I.14.16 (p.~122)} }
\expandafter\def\csname SRef@push:A47\endcsname{\href{S6_SERIES_GRIND_I.pdf\#nameddest=tag-02ZX}{I.14.16 (p.~122)} }
\expandafter\def\csname SRef@push:src:A.4.8\endcsname{\href{S6_SERIES_GRIND_I.pdf\#nameddest=tag-02ZY}{I.14.17 (p.~123)} }
\expandafter\def\csname SRef@push:A48\endcsname{\href{S6_SERIES_GRIND_I.pdf\#nameddest=tag-02ZY}{I.14.17 (p.~123)} }
\expandafter\def\csname SRef@push:src:A.1.2\endcsname{\href{S6_SERIES_GRIND_I.pdf\#nameddest=tag-02ZI}{I.14.2 (p.~103)} }
\expandafter\def\csname SRef@push:A12\endcsname{\href{S6_SERIES_GRIND_I.pdf\#nameddest=tag-02ZI}{I.14.2 (p.~103)} }
\expandafter\def\csname SRef@push:src:A.2.2\endcsname{\href{S6_SERIES_GRIND_I.pdf\#nameddest=tag-02ZK}{I.14.4 (p.~106)} }
\expandafter\def\csname SRef@push:A22\endcsname{\href{S6_SERIES_GRIND_I.pdf\#nameddest=tag-02ZK}{I.14.4 (p.~106)} }
\expandafter\def\csname SRef@push:src:A.2.4\endcsname{\href{S6_SERIES_GRIND_I.pdf\#nameddest=tag-02ZM}{I.14.5 (p.~108)} }
\expandafter\def\csname SRef@push:A24\endcsname{\href{S6_SERIES_GRIND_I.pdf\#nameddest=tag-02ZM}{I.14.5 (p.~108)} }
\expandafter\def\csname SRef@push:src:A.3.2\endcsname{\href{S6_SERIES_GRIND_I.pdf\#nameddest=tag-02ZP}{I.14.8 (p.~114)} }
\expandafter\def\csname SRef@push:A32\endcsname{\href{S6_SERIES_GRIND_I.pdf\#nameddest=tag-02ZP}{I.14.8 (p.~114)} }
\expandafter\def\csname SRef@push:src:B.4.2\endcsname{\href{S6_SERIES_GRIND_I.pdf\#nameddest=tag-030O}{I.15.15 (p.~143)} }
\expandafter\def\csname SRef@push:B42\endcsname{\href{S6_SERIES_GRIND_I.pdf\#nameddest=tag-030O}{I.15.15 (p.~143)} }
\expandafter\def\csname SRef@push:src:B.4.5\endcsname{\href{S6_SERIES_GRIND_I.pdf\#nameddest=tag-0315}{I.15.18 (p.~147)} }
\expandafter\def\csname SRef@push:src:B.1.2\endcsname{\href{S6_SERIES_GRIND_I.pdf\#nameddest=tag-0301}{I.15.2 (p.~129)} }
\expandafter\def\csname SRef@push:B12\endcsname{\href{S6_SERIES_GRIND_I.pdf\#nameddest=tag-0301}{I.15.2 (p.~129)} }
\expandafter\def\csname SRef@push:src:B.4.33\endcsname{\href{S6_SERIES_GRIND_I.pdf\#nameddest=tag-0313}{I.15.46 (p.~222)} }
\expandafter\def\csname SRef@push:B433\endcsname{\href{S6_SERIES_GRIND_I.pdf\#nameddest=tag-0313}{I.15.46 (p.~222)} }
\expandafter\def\csname SRef@push:src:B.3.1\endcsname{\href{S6_SERIES_GRIND_I.pdf\#nameddest=tag-0305}{I.15.6 (p.~133)} }
\expandafter\def\csname SRef@push:B31\endcsname{\href{S6_SERIES_GRIND_I.pdf\#nameddest=tag-0305}{I.15.6 (p.~133)} }
\expandafter\def\csname SRef@push:src:B.3.4\endcsname{\href{S6_SERIES_GRIND_I.pdf\#nameddest=tag-0308}{I.15.9 (p.~135)} }
\expandafter\def\csname SRef@push:B34\endcsname{\href{S6_SERIES_GRIND_I.pdf\#nameddest=tag-0308}{I.15.9 (p.~135)} }
\expandafter\def\csname SRef@push:src:D.1.3\endcsname{\href{S6_SERIES_GRIND_I.pdf\#nameddest=tag-032C}{I.16.3 (p.~229)} }
\expandafter\def\csname SRef@push:src:D.2.3\endcsname{\href{S6_SERIES_GRIND_I.pdf\#nameddest=tag-032F}{I.16.6 (p.~231)} }
\expandafter\def\csname SRef@push:src:D.3.1\endcsname{\href{S6_SERIES_GRIND_I.pdf\#nameddest=tag-032H}{I.16.8 (p.~233)} }
\expandafter\def\csname SRef@night:P2.1\endcsname{\href{S6_SERIES_GRIND_I.pdf\#nameddest=tag-02R9}{I.17.1 (p.~243)} }
\expandafter\def\csname SRef@stmt:P2.1\endcsname{\href{S6_SERIES_GRIND_I.pdf\#nameddest=tag-02R9}{I.17.1 (p.~243)} }
\expandafter\def\csname SRef@night:P2.10\endcsname{\href{S6_SERIES_GRIND_I.pdf\#nameddest=tag-02RA}{I.17.10 (p.~251)} }
\expandafter\def\csname SRef@stmt:P2.10\endcsname{\href{S6_SERIES_GRIND_I.pdf\#nameddest=tag-02RA}{I.17.10 (p.~251)} }
\expandafter\def\csname SRef@night:P2.11\endcsname{\href{S6_SERIES_GRIND_I.pdf\#nameddest=tag-02RB}{I.17.11 (p.~253)} }
\expandafter\def\csname SRef@stmt:P2.11\endcsname{\href{S6_SERIES_GRIND_I.pdf\#nameddest=tag-02RB}{I.17.11 (p.~253)} }
\expandafter\def\csname SRef@night:P2.2\endcsname{\href{S6_SERIES_GRIND_I.pdf\#nameddest=tag-02RC}{I.17.2 (p.~244)} }
\expandafter\def\csname SRef@stmt:P2.2\endcsname{\href{S6_SERIES_GRIND_I.pdf\#nameddest=tag-02RC}{I.17.2 (p.~244)} }
\expandafter\def\csname SRef@night:P2.3\endcsname{\href{S6_SERIES_GRIND_I.pdf\#nameddest=tag-02RD}{I.17.3 (p.~245)} }
\expandafter\def\csname SRef@stmt:P2.3\endcsname{\href{S6_SERIES_GRIND_I.pdf\#nameddest=tag-02RD}{I.17.3 (p.~245)} }
\expandafter\def\csname SRef@night:P2.4\endcsname{\href{S6_SERIES_GRIND_I.pdf\#nameddest=tag-02RE}{I.17.4 (p.~246)} }
\expandafter\def\csname SRef@stmt:P2.4\endcsname{\href{S6_SERIES_GRIND_I.pdf\#nameddest=tag-02RE}{I.17.4 (p.~246)} }
\expandafter\def\csname SRef@night:P2.5\endcsname{\href{S6_SERIES_GRIND_I.pdf\#nameddest=tag-02RF}{I.17.5 (p.~246)} }
\expandafter\def\csname SRef@stmt:P2.5\endcsname{\href{S6_SERIES_GRIND_I.pdf\#nameddest=tag-02RF}{I.17.5 (p.~246)} }
\expandafter\def\csname SRef@night:P2.6\endcsname{\href{S6_SERIES_GRIND_I.pdf\#nameddest=tag-02RG}{I.17.6 (p.~247)} }
\expandafter\def\csname SRef@stmt:P2.6\endcsname{\href{S6_SERIES_GRIND_I.pdf\#nameddest=tag-02RG}{I.17.6 (p.~247)} }
\expandafter\def\csname SRef@night:P2.7\endcsname{\href{S6_SERIES_GRIND_I.pdf\#nameddest=tag-02RH}{I.17.7 (p.~248)} }
\expandafter\def\csname SRef@stmt:P2.7\endcsname{\href{S6_SERIES_GRIND_I.pdf\#nameddest=tag-02RH}{I.17.7 (p.~248)} }
\expandafter\def\csname SRef@night:P2.8\endcsname{\href{S6_SERIES_GRIND_I.pdf\#nameddest=tag-02RI}{I.17.8 (p.~249)} }
\expandafter\def\csname SRef@stmt:P2.8\endcsname{\href{S6_SERIES_GRIND_I.pdf\#nameddest=tag-02RI}{I.17.8 (p.~249)} }
\expandafter\def\csname SRef@night:P2.9\endcsname{\href{S6_SERIES_GRIND_I.pdf\#nameddest=tag-02RJ}{I.17.9 (p.~249)} }
\expandafter\def\csname SRef@stmt:P2.9\endcsname{\href{S6_SERIES_GRIND_I.pdf\#nameddest=tag-02RJ}{I.17.9 (p.~249)} }
\expandafter\def\csname SRef@push:src:E.2.1\endcsname{\href{S6_SERIES_GRIND_I.pdf\#nameddest=tag-032U}{I.20.6 (p.~282)} }
\expandafter\def\csname SRef@mapsI:src:A.2.1\endcsname{\href{S6_SERIES_GRIND_I.pdf\#nameddest=tag-02PD}{I.22.4 (p.~298)} }
\expandafter\def\csname SRef@push:src:E.4.1\endcsname{\href{S6_SERIES_GRIND_I.pdf\#nameddest=tag-033L}{I.24.5 (p.~338)} }
\expandafter\def\csname SRef@push:src:E.4.2\endcsname{\href{S6_SERIES_GRIND_I.pdf\#nameddest=tag-033M}{I.24.6 (p.~341)} }
\expandafter\def\csname SRef@fam:src:C.2.3\endcsname{\href{S6_SERIES_GRIND_I.pdf\#nameddest=tag-01QL}{I.4.5 (p.~34)} }
\expandafter\def\csname SRef@ari:C.8.5\endcsname{\href{S6_SERIES_GRIND_I.pdf\#nameddest=tag-01QL}{I.4.5 (p.~34)} }
\expandafter\def\csname SRef@fam:st:C.2.3\endcsname{\href{S6_SERIES_GRIND_I.pdf\#nameddest=tag-01QL}{I.4.5 (p.~34)} }
\expandafter\def\csname SRef@push:src:E.4.3\endcsname{\href{S6_SERIES_GRIND_I.pdf\#nameddest=tag-01QL}{I.4.5 (p.~34)} }
\expandafter\def\csname SRef@fam:src:A.5.1\endcsname{\href{S6_SERIES_GRIND_I.pdf\#nameddest=tag-01Q6}{I.7.8 (p.~56)} }
\expandafter\def\csname SRef@wb:pshift:thm:classification\endcsname{\href{S6_SERIES_GRIND_I.pdf\#nameddest=tag-01Q6}{I.7.8 (p.~56)} }
\expandafter\def\csname SRef@fam:st:A.5.1\endcsname{\href{S6_SERIES_GRIND_I.pdf\#nameddest=tag-01Q6}{I.7.8 (p.~56)} }
\expandafter\def\csname SRef@fam:st:A.5.2\endcsname{\href{S6_SERIES_GRIND_I.pdf\#nameddest=tag-01RH}{I.7.9 (p.~56)} }
\expandafter\def\csname SRef@fam:src:A.5.2\endcsname{\href{S6_SERIES_GRIND_I.pdf\#nameddest=tag-01RH}{I.7.9 (p.~56)} }
\expandafter\def\csname SRef@inv:D1.3\endcsname{\href{S6_SERIES_GRIND_I.pdf\#nameddest=tag-01RH}{I.7.9 (p.~56)} }
\expandafter\def\csname SRef@inv:D2.1\endcsname{\href{S6_SERIES_GRIND_I.pdf\#nameddest=tag-01RH}{I.7.9 (p.~56)} }
\expandafter\def\csname SRef@inv:D3.1\endcsname{\href{S6_SERIES_GRIND_I.pdf\#nameddest=tag-01RH}{I.7.9 (p.~56)} }
\expandafter\def\csname SRef@cdp:st:E9\endcsname{\href{S6_SERIES_GRIND_I.pdf\#nameddest=tag-01RH}{I.7.9 (p.~56)} }
\expandafter\def\csname SRef@push:src:E.1.2\endcsname{\href{S6_SERIES_GRIND_I.pdf\#nameddest=tag-032S}{I.9.2 (p.~65)} }
\expandafter\def\csname SRef@wb:thm:canonical-ring-recovery\endcsname{\href{S6_SERIES_GRIND_I.pdf\#nameddest=tag-03DM}{I.A.2 (p.~342)} }
\expandafter\def\csname SRef@wb:thm:retained-period-deformation\endcsname{\href{S6_SERIES_GRIND_I.pdf\#nameddest=tag-03DP}{I.A.3 (p.~344)} }
\expandafter\def\csname SRef@night:P5.1\endcsname{\href{S6_SERIES_GRIND_I.pdf\#nameddest=tag-02YM}{I.C.1 (p.~366)} }
\expandafter\def\csname SRef@stmt:P5.1\endcsname{\href{S6_SERIES_GRIND_I.pdf\#nameddest=tag-02YM}{I.C.1 (p.~366)} }
\expandafter\def\csname SRef@night:P5.2\endcsname{\href{S6_SERIES_GRIND_I.pdf\#nameddest=tag-02YN}{I.C.2 (p.~367)} }
\expandafter\def\csname SRef@stmt:P5.2\endcsname{\href{S6_SERIES_GRIND_I.pdf\#nameddest=tag-02YN}{I.C.2 (p.~367)} }
\expandafter\def\csname SRef@night:P5.3\endcsname{\href{S6_SERIES_GRIND_I.pdf\#nameddest=tag-02YO}{I.C.3 (p.~369)} }
\expandafter\def\csname SRef@stmt:P5.3\endcsname{\href{S6_SERIES_GRIND_I.pdf\#nameddest=tag-02YO}{I.C.3 (p.~369)} }
\expandafter\def\csname SRef@night:P5.4\endcsname{\href{S6_SERIES_GRIND_I.pdf\#nameddest=tag-02YP}{I.C.4 (p.~369)} }
\expandafter\def\csname SRef@stmt:P5.4\endcsname{\href{S6_SERIES_GRIND_I.pdf\#nameddest=tag-02YP}{I.C.4 (p.~369)} }
\expandafter\def\csname SRef@night:P5.5\endcsname{\href{S6_SERIES_GRIND_I.pdf\#nameddest=tag-02YQ}{I.C.5 (p.~370)} }
\expandafter\def\csname SRef@stmt:P5.5\endcsname{\href{S6_SERIES_GRIND_I.pdf\#nameddest=tag-02YQ}{I.C.5 (p.~370)} }
\expandafter\def\csname SRef@night:P5.6\endcsname{\href{S6_SERIES_GRIND_I.pdf\#nameddest=tag-02YR}{I.C.6 (p.~372)} }
\expandafter\def\csname SRef@stmt:P5.6\endcsname{\href{S6_SERIES_GRIND_I.pdf\#nameddest=tag-02YR}{I.C.6 (p.~372)} }
\expandafter\def\csname SRef@cdp:eq:A.RC.2\endcsname{\hyperlink{tag-006V}{II.\ref*{cdp:eq:A.RC.2}}}
\expandafter\def\csname SRef@cdp:eq:A.RC.3\endcsname{\hyperlink{tag-006W}{II.\ref*{cdp:eq:A.RC.3}}}
\expandafter\def\csname SRef@cdp:eq:A.RC.4\endcsname{\hyperlink{tag-006X}{II.\ref*{cdp:eq:A.RC.4}}}
\expandafter\def\csname SRef@cdp:eq:A.RC.5\endcsname{\hyperlink{tag-006Y}{II.\ref*{cdp:eq:A.RC.5}}}
\expandafter\def\csname SRef@cdp:eq:A.RC.6\endcsname{\hyperlink{tag-006Z}{II.\ref*{cdp:eq:A.RC.6}}}
\expandafter\def\csname SRef@cdp:eq:A.RC.7\endcsname{\hyperlink{tag-0070}{II.\ref*{cdp:eq:A.RC.7}}}
\expandafter\def\csname SRef@cdp:eq:A.Pi.1\endcsname{\hyperlink{tag-006O}{II.\ref*{cdp:eq:A.Pi.1}}}
\expandafter\def\csname SRef@cdp:eq:A.Pi.2\endcsname{\hyperlink{tag-006P}{II.\ref*{cdp:eq:A.Pi.2}}}
\expandafter\def\csname SRef@cdp:eq:A.Pi.3\endcsname{\hyperlink{tag-006Q}{II.\ref*{cdp:eq:A.Pi.3}}}
\expandafter\def\csname SRef@cdp:eq:A.Pi.4\endcsname{\hyperlink{tag-006R}{II.\ref*{cdp:eq:A.Pi.4}}}
\expandafter\def\csname SRef@cdp:eq:A.CD.1\endcsname{\hyperlink{tag-0058}{II.\ref*{cdp:eq:A.CD.1}}}
\expandafter\def\csname SRef@cdp:eq:A.CD.2\endcsname{\hyperlink{tag-0059}{II.\ref*{cdp:eq:A.CD.2}}}
\expandafter\def\csname SRef@cdp:eq:A.CD.3\endcsname{\hyperlink{tag-005A}{II.\ref*{cdp:eq:A.CD.3}}}
\expandafter\def\csname SRef@cdp:eq:A.CD.4\endcsname{\hyperlink{tag-005B}{II.\ref*{cdp:eq:A.CD.4}}}
\expandafter\def\csname SRef@cdp:eq:A.Q1.1\endcsname{\hyperlink{tag-006S}{II.\ref*{cdp:eq:A.Q1.1}}}
\expandafter\def\csname SRef@cdp:eq:A.Q1.2\endcsname{\hyperlink{tag-006T}{II.\ref*{cdp:eq:A.Q1.2}}}
\expandafter\def\csname SRef@cdp:eq:A.J.1\endcsname{\hyperlink{tag-005C}{II.\ref*{cdp:eq:A.J.1}}}
\expandafter\def\csname SRef@cdp:eq:A.J.2\endcsname{\hyperlink{tag-005N}{II.\ref*{cdp:eq:A.J.2}}}
\expandafter\def\csname SRef@cdp:eq:A.J.3\endcsname{\hyperlink{tag-005Y}{II.\ref*{cdp:eq:A.J.3}}}
\expandafter\def\csname SRef@cdp:eq:A.J.4\endcsname{\hyperlink{tag-006C}{II.\ref*{cdp:eq:A.J.4}}}
\expandafter\def\csname SRef@cdp:eq:A.J.5\endcsname{\hyperlink{tag-006E}{II.\ref*{cdp:eq:A.J.5}}}
\expandafter\def\csname SRef@cdp:eq:A.J.6\endcsname{\hyperlink{tag-006F}{II.\ref*{cdp:eq:A.J.6}}}
\expandafter\def\csname SRef@cdp:eq:A.J.7\endcsname{\hyperlink{tag-006G}{II.\ref*{cdp:eq:A.J.7}}}
\expandafter\def\csname SRef@cdp:eq:A.J.8\endcsname{\hyperlink{tag-006H}{II.\ref*{cdp:eq:A.J.8}}}
\expandafter\def\csname SRef@cdp:eq:A.J.9\endcsname{\hyperlink{tag-006I}{II.\ref*{cdp:eq:A.J.9}}}
\expandafter\def\csname SRef@cdp:eq:A.J.10\endcsname{\hyperlink{tag-005D}{II.\ref*{cdp:eq:A.J.10}}}
\expandafter\def\csname SRef@cdp:eq:A.J.11\endcsname{\hyperlink{tag-005E}{II.\ref*{cdp:eq:A.J.11}}}
\expandafter\def\csname SRef@cdp:eq:A.J.12\endcsname{\hyperlink{tag-005F}{II.\ref*{cdp:eq:A.J.12}}}
\expandafter\def\csname SRef@cdp:eq:A.J.13\endcsname{\hyperlink{tag-005G}{II.\ref*{cdp:eq:A.J.13}}}
\expandafter\def\csname SRef@cdp:eq:A.J.14\endcsname{\hyperlink{tag-005H}{II.\ref*{cdp:eq:A.J.14}}}
\expandafter\def\csname SRef@cdp:eq:A.J.15\endcsname{\hyperlink{tag-005I}{II.\ref*{cdp:eq:A.J.15}}}
\expandafter\def\csname SRef@cdp:eq:A.J.16\endcsname{\hyperlink{tag-005J}{II.\ref*{cdp:eq:A.J.16}}}
\expandafter\def\csname SRef@cdp:eq:A.J.17\endcsname{\hyperlink{tag-005K}{II.\ref*{cdp:eq:A.J.17}}}
\expandafter\def\csname SRef@cdp:eq:A.J.18\endcsname{\hyperlink{tag-005L}{II.\ref*{cdp:eq:A.J.18}}}
\expandafter\def\csname SRef@cdp:eq:A.J.19\endcsname{\hyperlink{tag-005M}{II.\ref*{cdp:eq:A.J.19}}}
\expandafter\def\csname SRef@cdp:eq:A.J.20\endcsname{\hyperlink{tag-005O}{II.\ref*{cdp:eq:A.J.20}}}
\expandafter\def\csname SRef@cdp:eq:A.J.30\endcsname{\hyperlink{tag-005Z}{II.\ref*{cdp:eq:A.J.30}}}
\expandafter\def\csname SRef@cdp:eq:A.J.31\endcsname{\hyperlink{tag-0060}{II.\ref*{cdp:eq:A.J.31}}}
\expandafter\def\csname SRef@cdp:eq:A.J.32\endcsname{\hyperlink{tag-0061}{II.\ref*{cdp:eq:A.J.32}}}
\expandafter\def\csname SRef@cdp:eq:A.J.21\endcsname{\hyperlink{tag-005P}{II.\ref*{cdp:eq:A.J.21}}}
\expandafter\def\csname SRef@cdp:eq:A.J.22\endcsname{\hyperlink{tag-005Q}{II.\ref*{cdp:eq:A.J.22}}}
\expandafter\def\csname SRef@cdp:eq:A.J.23\endcsname{\hyperlink{tag-005R}{II.\ref*{cdp:eq:A.J.23}}}
\expandafter\def\csname SRef@cdp:eq:A.J.24\endcsname{\hyperlink{tag-005S}{II.\ref*{cdp:eq:A.J.24}}}
\expandafter\def\csname SRef@cdp:eq:A.J.25\endcsname{\hyperlink{tag-005T}{II.\ref*{cdp:eq:A.J.25}}}
\expandafter\def\csname SRef@cdp:eq:A.J.26\endcsname{\hyperlink{tag-005U}{II.\ref*{cdp:eq:A.J.26}}}
\expandafter\def\csname SRef@cdp:eq:A.J.27\endcsname{\hyperlink{tag-005V}{II.\ref*{cdp:eq:A.J.27}}}
\expandafter\def\csname SRef@cdp:eq:A.J.28\endcsname{\hyperlink{tag-005W}{II.\ref*{cdp:eq:A.J.28}}}
\expandafter\def\csname SRef@cdp:eq:A.J.29\endcsname{\hyperlink{tag-005X}{II.\ref*{cdp:eq:A.J.29}}}
\expandafter\def\csname SRef@cdp:eq:A.J.33\endcsname{\hyperlink{tag-0062}{II.\ref*{cdp:eq:A.J.33}}}
\expandafter\def\csname SRef@cdp:eq:A.J.34\endcsname{\hyperlink{tag-0065}{II.\ref*{cdp:eq:A.J.34}}}
\expandafter\def\csname SRef@cdp:eq:A.J.33a\endcsname{\hyperlink{tag-0063}{II.\ref*{cdp:eq:A.J.33a}}}
\expandafter\def\csname SRef@cdp:eq:A.J.33b\endcsname{\hyperlink{tag-0064}{II.\ref*{cdp:eq:A.J.33b}}}
\expandafter\def\csname SRef@cdp:eq:A.J.35\endcsname{\hyperlink{tag-0066}{II.\ref*{cdp:eq:A.J.35}}}
\expandafter\def\csname SRef@cdp:eq:A.J.36\endcsname{\hyperlink{tag-0067}{II.\ref*{cdp:eq:A.J.36}}}
\expandafter\def\csname SRef@cdp:eq:A.J.37\endcsname{\hyperlink{tag-0068}{II.\ref*{cdp:eq:A.J.37}}}
\expandafter\def\csname SRef@cdp:eq:A.J.38\endcsname{\hyperlink{tag-0069}{II.\ref*{cdp:eq:A.J.38}}}
\expandafter\def\csname SRef@cdp:eq:A.J.38a\endcsname{\hyperlink{tag-006A}{II.\ref*{cdp:eq:A.J.38a}}}
\expandafter\def\csname SRef@cdp:eq:A.J.39\endcsname{\hyperlink{tag-006B}{II.\ref*{cdp:eq:A.J.39}}}
\expandafter\def\csname SRef@cdp:eq:A.J.40\endcsname{\hyperlink{tag-006D}{II.\ref*{cdp:eq:A.J.40}}}
\expandafter\def\csname SRef@cdp:eq:A.T.1\endcsname{\hyperlink{tag-007M}{II.\ref*{cdp:eq:A.T.1}}}
\expandafter\def\csname SRef@cdp:eq:A.T.2\endcsname{\hyperlink{tag-007X}{II.\ref*{cdp:eq:A.T.2}}}
\expandafter\def\csname SRef@cdp:eq:A.T.3\endcsname{\hyperlink{tag-0082}{II.\ref*{cdp:eq:A.T.3}}}
\expandafter\def\csname SRef@cdp:eq:A.T.4\endcsname{\hyperlink{tag-0083}{II.\ref*{cdp:eq:A.T.4}}}
\expandafter\def\csname SRef@cdp:eq:A.T.5\endcsname{\hyperlink{tag-0084}{II.\ref*{cdp:eq:A.T.5}}}
\expandafter\def\csname SRef@cdp:eq:A.T.6\endcsname{\hyperlink{tag-0085}{II.\ref*{cdp:eq:A.T.6}}}
\expandafter\def\csname SRef@cdp:eq:A.T.7\endcsname{\hyperlink{tag-0086}{II.\ref*{cdp:eq:A.T.7}}}
\expandafter\def\csname SRef@cdp:eq:A.T.8\endcsname{\hyperlink{tag-0087}{II.\ref*{cdp:eq:A.T.8}}}
\expandafter\def\csname SRef@cdp:eq:A.T.9\endcsname{\hyperlink{tag-0088}{II.\ref*{cdp:eq:A.T.9}}}
\expandafter\def\csname SRef@cdp:eq:A.T.10\endcsname{\hyperlink{tag-007N}{II.\ref*{cdp:eq:A.T.10}}}
\expandafter\def\csname SRef@cdp:eq:A.T.11\endcsname{\hyperlink{tag-007O}{II.\ref*{cdp:eq:A.T.11}}}
\expandafter\def\csname SRef@cdp:eq:A.T.12\endcsname{\hyperlink{tag-007P}{II.\ref*{cdp:eq:A.T.12}}}
\expandafter\def\csname SRef@cdp:eq:A.T.13\endcsname{\hyperlink{tag-007Q}{II.\ref*{cdp:eq:A.T.13}}}
\expandafter\def\csname SRef@cdp:eq:A.T.14\endcsname{\hyperlink{tag-007R}{II.\ref*{cdp:eq:A.T.14}}}
\expandafter\def\csname SRef@cdp:eq:A.T.15\endcsname{\hyperlink{tag-007S}{II.\ref*{cdp:eq:A.T.15}}}
\expandafter\def\csname SRef@cdp:eq:A.T.16\endcsname{\hyperlink{tag-007T}{II.\ref*{cdp:eq:A.T.16}}}
\expandafter\def\csname SRef@cdp:eq:A.T.17\endcsname{\hyperlink{tag-007U}{II.\ref*{cdp:eq:A.T.17}}}
\expandafter\def\csname SRef@cdp:eq:A.T.18\endcsname{\hyperlink{tag-007V}{II.\ref*{cdp:eq:A.T.18}}}
\expandafter\def\csname SRef@cdp:eq:A.T.19\endcsname{\hyperlink{tag-007W}{II.\ref*{cdp:eq:A.T.19}}}
\expandafter\def\csname SRef@cdp:eq:A.T.20\endcsname{\hyperlink{tag-007Y}{II.\ref*{cdp:eq:A.T.20}}}
\expandafter\def\csname SRef@cdp:eq:A.T.21\endcsname{\hyperlink{tag-007Z}{II.\ref*{cdp:eq:A.T.21}}}
\expandafter\def\csname SRef@cdp:eq:A.T.22\endcsname{\hyperlink{tag-0080}{II.\ref*{cdp:eq:A.T.22}}}
\expandafter\def\csname SRef@cdp:eq:A.T.23\endcsname{\hyperlink{tag-0081}{II.\ref*{cdp:eq:A.T.23}}}
\expandafter\def\csname SRef@cdp:eq:A.U.1\endcsname{\hyperlink{tag-0089}{II.\ref*{cdp:eq:A.U.1}}}
\expandafter\def\csname SRef@cdp:eq:A.U.2\endcsname{\hyperlink{tag-008A}{II.\ref*{cdp:eq:A.U.2}}}
\expandafter\def\csname SRef@cdp:eq:A.U.3\endcsname{\hyperlink{tag-008B}{II.\ref*{cdp:eq:A.U.3}}}
\expandafter\def\csname SRef@cdp:eq:A.U.4\endcsname{\hyperlink{tag-008C}{II.\ref*{cdp:eq:A.U.4}}}
\expandafter\def\csname SRef@cdp:eq:A.U.5\endcsname{\hyperlink{tag-008D}{II.\ref*{cdp:eq:A.U.5}}}
\expandafter\def\csname SRef@cdp:eq:A.U.6\endcsname{\hyperlink{tag-008E}{II.\ref*{cdp:eq:A.U.6}}}
\expandafter\def\csname SRef@cdp:eq:A.U.7\endcsname{\hyperlink{tag-008F}{II.\ref*{cdp:eq:A.U.7}}}
\expandafter\def\csname SRef@cdp:eq:A.RT.1\endcsname{\hyperlink{tag-0071}{II.\ref*{cdp:eq:A.RT.1}}}
\expandafter\def\csname SRef@cdp:eq:A.RT.2\endcsname{\hyperlink{tag-007C}{II.\ref*{cdp:eq:A.RT.2}}}
\expandafter\def\csname SRef@cdp:eq:A.RT.3\endcsname{\hyperlink{tag-007F}{II.\ref*{cdp:eq:A.RT.3}}}
\expandafter\def\csname SRef@cdp:eq:A.RT.4\endcsname{\hyperlink{tag-007G}{II.\ref*{cdp:eq:A.RT.4}}}
\expandafter\def\csname SRef@cdp:eq:A.RT.5\endcsname{\hyperlink{tag-007H}{II.\ref*{cdp:eq:A.RT.5}}}
\expandafter\def\csname SRef@cdp:eq:A.RT.6\endcsname{\hyperlink{tag-007I}{II.\ref*{cdp:eq:A.RT.6}}}
\expandafter\def\csname SRef@cdp:eq:A.RT.7\endcsname{\hyperlink{tag-007J}{II.\ref*{cdp:eq:A.RT.7}}}
\expandafter\def\csname SRef@cdp:eq:A.RT.8\endcsname{\hyperlink{tag-007K}{II.\ref*{cdp:eq:A.RT.8}}}
\expandafter\def\csname SRef@cdp:eq:A.RT.9\endcsname{\hyperlink{tag-007L}{II.\ref*{cdp:eq:A.RT.9}}}
\expandafter\def\csname SRef@cdp:eq:A.RT.10\endcsname{\hyperlink{tag-0072}{II.\ref*{cdp:eq:A.RT.10}}}
\expandafter\def\csname SRef@cdp:eq:A.RT.11\endcsname{\hyperlink{tag-0073}{II.\ref*{cdp:eq:A.RT.11}}}
\expandafter\def\csname SRef@cdp:eq:A.RT.12\endcsname{\hyperlink{tag-0074}{II.\ref*{cdp:eq:A.RT.12}}}
\expandafter\def\csname SRef@cdp:eq:A.RT.13\endcsname{\hyperlink{tag-0075}{II.\ref*{cdp:eq:A.RT.13}}}
\expandafter\def\csname SRef@cdp:eq:A.RT.14\endcsname{\hyperlink{tag-0076}{II.\ref*{cdp:eq:A.RT.14}}}
\expandafter\def\csname SRef@cdp:eq:A.RT.15\endcsname{\hyperlink{tag-0077}{II.\ref*{cdp:eq:A.RT.15}}}
\expandafter\def\csname SRef@cdp:eq:A.RT.16\endcsname{\hyperlink{tag-0078}{II.\ref*{cdp:eq:A.RT.16}}}
\expandafter\def\csname SRef@cdp:eq:A.RT.17\endcsname{\hyperlink{tag-0079}{II.\ref*{cdp:eq:A.RT.17}}}
\expandafter\def\csname SRef@cdp:eq:A.RT.18\endcsname{\hyperlink{tag-007A}{II.\ref*{cdp:eq:A.RT.18}}}
\expandafter\def\csname SRef@cdp:eq:A.RT.19\endcsname{\hyperlink{tag-007B}{II.\ref*{cdp:eq:A.RT.19}}}
\expandafter\def\csname SRef@cdp:eq:A.RT.20\endcsname{\hyperlink{tag-007D}{II.\ref*{cdp:eq:A.RT.20}}}
\expandafter\def\csname SRef@cdp:eq:A.RT.21\endcsname{\hyperlink{tag-007E}{II.\ref*{cdp:eq:A.RT.21}}}
\expandafter\def\csname SRef@cdp:eq:A.BP.1\endcsname{\hyperlink{tag-004W}{II.\ref*{cdp:eq:A.BP.1}}}
\expandafter\def\csname SRef@cdp:eq:A.BP.4\endcsname{\hyperlink{tag-0052}{II.\ref*{cdp:eq:A.BP.4}}}
\expandafter\def\csname SRef@cdp:eq:A.BP.2\endcsname{\hyperlink{tag-0050}{II.\ref*{cdp:eq:A.BP.2}}}
\expandafter\def\csname SRef@cdp:eq:A.BP.3\endcsname{\hyperlink{tag-0051}{II.\ref*{cdp:eq:A.BP.3}}}
\expandafter\def\csname SRef@cdp:eq:A.BP.5\endcsname{\hyperlink{tag-0053}{II.\ref*{cdp:eq:A.BP.5}}}
\expandafter\def\csname SRef@cdp:eq:A.BP.6\endcsname{\hyperlink{tag-0054}{II.\ref*{cdp:eq:A.BP.6}}}
\expandafter\def\csname SRef@cdp:eq:A.BP.7\endcsname{\hyperlink{tag-0055}{II.\ref*{cdp:eq:A.BP.7}}}
\expandafter\def\csname SRef@cdp:eq:A.BP.8\endcsname{\hyperlink{tag-0056}{II.\ref*{cdp:eq:A.BP.8}}}
\expandafter\def\csname SRef@cdp:eq:A.BP.9\endcsname{\hyperlink{tag-0057}{II.\ref*{cdp:eq:A.BP.9}}}
\expandafter\def\csname SRef@cdp:eq:A.BP.10\endcsname{\hyperlink{tag-004X}{II.\ref*{cdp:eq:A.BP.10}}}
\expandafter\def\csname SRef@cdp:eq:A.BP.11\endcsname{\hyperlink{tag-004Y}{II.\ref*{cdp:eq:A.BP.11}}}
\expandafter\def\csname SRef@cdp:eq:A.BP.12\endcsname{\hyperlink{tag-004Z}{II.\ref*{cdp:eq:A.BP.12}}}
\expandafter\def\csname SRef@cdp:eq:C.NC.4\endcsname{\hyperlink{tag-00S3}{II.\ref*{cdp:eq:C.NC.4}}}
\expandafter\def\csname SRef@cdp:eq:C.NC.5\endcsname{\hyperlink{tag-00S4}{II.\ref*{cdp:eq:C.NC.5}}}
\expandafter\def\csname SRef@cdp:eq:C.NC.6\endcsname{\hyperlink{tag-00S5}{II.\ref*{cdp:eq:C.NC.6}}}
\expandafter\def\csname SRef@cdp:eq:C.NC.7\endcsname{\hyperlink{tag-00S6}{II.\ref*{cdp:eq:C.NC.7}}}
\expandafter\def\csname SRef@cdp:eq:C.NC.8\endcsname{\hyperlink{tag-00S7}{II.\ref*{cdp:eq:C.NC.8}}}
\expandafter\def\csname SRef@cdp:eq:C.NC.9\endcsname{\hyperlink{tag-00S8}{II.\ref*{cdp:eq:C.NC.9}}}
\expandafter\def\csname SRef@cdp:eq:C.NC.10\endcsname{\hyperlink{tag-00S1}{II.\ref*{cdp:eq:C.NC.10}}}
\expandafter\def\csname SRef@cdp:eq:C.NC.11\endcsname{\hyperlink{tag-00S2}{II.\ref*{cdp:eq:C.NC.11}}}
\expandafter\def\csname SRef@cdp:eq:C.J.1\endcsname{\hyperlink{tag-00OV}{II.\ref*{cdp:eq:C.J.1}}}
\expandafter\def\csname SRef@cdp:eq:C.J.2\endcsname{\hyperlink{tag-00OW}{II.\ref*{cdp:eq:C.J.2}}}
\expandafter\def\csname SRef@cdp:eq:C.J.3\endcsname{\hyperlink{tag-00OX}{II.\ref*{cdp:eq:C.J.3}}}
\expandafter\def\csname SRef@cdp:eq:D1\endcsname{\hyperlink{tag-019T}{II.\ref*{cdp:eq:D1}}}
\expandafter\def\csname SRef@cdp:eq:D5\endcsname{\hyperlink{tag-019X}{II.\ref*{cdp:eq:D5}}}
\expandafter\def\csname SRef@cdp:eq:D4\endcsname{\hyperlink{tag-019W}{II.\ref*{cdp:eq:D4}}}
\expandafter\def\csname SRef@cdp:eq:D2\endcsname{\hyperlink{tag-019U}{II.\ref*{cdp:eq:D2}}}
\expandafter\def\csname SRef@cdp:eq:D3\endcsname{\hyperlink{tag-019V}{II.\ref*{cdp:eq:D3}}}
\expandafter\def\csname SRef@cdp:eq:C.SK.1\endcsname{\hyperlink{tag-011H}{II.\ref*{cdp:eq:C.SK.1}}}
\expandafter\def\csname SRef@cdp:eq:C.SK.2\endcsname{\hyperlink{tag-011I}{II.\ref*{cdp:eq:C.SK.2}}}
\expandafter\def\csname SRef@a.-one-dimensional-first-jets\endcsname{\hyperlink{tag-0005}{II.\ref*{a.-one-dimensional-first-jets}}}
\expandafter\def\csname SRef@three-smooth-branches-the-complete-excess-two-classification\endcsname{\hyperlink{tag-0005}{II.\ref*{three-smooth-branches-the-complete-excess-two-classification}}}
\expandafter\def\csname SRef@cdp:eq:C.SK.3\endcsname{\hyperlink{tag-011J}{II.\ref*{cdp:eq:C.SK.3}}}
\expandafter\def\csname SRef@b.-two-dimensional-first-jets\endcsname{\hyperlink{tag-003V}{II.\ref*{b.-two-dimensional-first-jets}}}
\expandafter\def\csname SRef@multiplicity-four-one-cusp-and-two-smooth-branches\endcsname{\hyperlink{tag-02QH}{II.\ref*{multiplicity-four-one-cusp-and-two-smooth-branches}}}
\expandafter\def\csname SRef@a.-a-nonzero-coefficient-of-the-cusp-cubic\endcsname{\hyperlink{tag-0004}{II.\ref*{a.-a-nonzero-coefficient-of-the-cusp-cubic}}}
\expandafter\def\csname SRef@b.-a-zero-coefficient-of-the-cusp-cubic\endcsname{\hyperlink{tag-003U}{II.\ref*{b.-a-zero-coefficient-of-the-cusp-cubic}}}
\expandafter\def\csname SRef@cdp:eq:C.SK.4\endcsname{\hyperlink{tag-011K}{II.\ref*{cdp:eq:C.SK.4}}}
\expandafter\def\csname SRef@every-nonordinary-triple-exceeds-the-budget\endcsname{\hyperlink{tag-01OO}{II.\ref*{every-nonordinary-triple-exceeds-the-budget}}}
\expandafter\def\csname SRef@cdp:eq:C.SK.5\endcsname{\hyperlink{tag-011L}{II.\ref*{cdp:eq:C.SK.5}}}
\expandafter\def\csname SRef@cdp:eq:C.15\endcsname{\hyperlink{tag-008T}{II.\ref*{cdp:eq:C.15}}}
\expandafter\def\csname SRef@cdp:eq:C.16\endcsname{\hyperlink{tag-008U}{II.\ref*{cdp:eq:C.16}}}
\expandafter\def\csname SRef@cdp:eq:C.17\endcsname{\hyperlink{tag-008V}{II.\ref*{cdp:eq:C.17}}}
\expandafter\def\csname SRef@cdp:eq:C.18\endcsname{\hyperlink{tag-008W}{II.\ref*{cdp:eq:C.18}}}
\expandafter\def\csname SRef@cdp:eq:C.SK.6\endcsname{\hyperlink{tag-011M}{II.\ref*{cdp:eq:C.SK.6}}}
\expandafter\def\csname SRef@cdp:eq:C.19\endcsname{\hyperlink{tag-008X}{II.\ref*{cdp:eq:C.19}}}
\expandafter\def\csname SRef@cdp:eq:C.20\endcsname{\hyperlink{tag-008Z}{II.\ref*{cdp:eq:C.20}}}
\expandafter\def\csname SRef@cdp:eq:C.21\endcsname{\hyperlink{tag-0090}{II.\ref*{cdp:eq:C.21}}}
\expandafter\def\csname SRef@cdp:eq:C.22\endcsname{\hyperlink{tag-0091}{II.\ref*{cdp:eq:C.22}}}
\expandafter\def\csname SRef@cdp:eq:C.22b\endcsname{\hyperlink{tag-0093}{II.\ref*{cdp:eq:C.22b}}}
\expandafter\def\csname SRef@cdp:eq:C.22a\endcsname{\hyperlink{tag-0092}{II.\ref*{cdp:eq:C.22a}}}
\expandafter\def\csname SRef@cdp:eq:C.VC.1\endcsname{\hyperlink{tag-015G}{II.\ref*{cdp:eq:C.VC.1}}}
\expandafter\def\csname SRef@cdp:eq:C.VC.2\endcsname{\hyperlink{tag-015H}{II.\ref*{cdp:eq:C.VC.2}}}
\expandafter\def\csname SRef@cdp:eq:C.VC.3\endcsname{\hyperlink{tag-015I}{II.\ref*{cdp:eq:C.VC.3}}}
\expandafter\def\csname SRef@cdp:eq:C.VC.4\endcsname{\hyperlink{tag-015J}{II.\ref*{cdp:eq:C.VC.4}}}
\expandafter\def\csname SRef@cdp:eq:C.PK.1\endcsname{\hyperlink{tag-00V8}{II.\ref*{cdp:eq:C.PK.1}}}
\expandafter\def\csname SRef@cdp:eq:C.TEST.1\endcsname{\hyperlink{tag-012O}{II.\ref*{cdp:eq:C.TEST.1}}}
\expandafter\def\csname SRef@cdp:eq:C.TEST.2\endcsname{\hyperlink{tag-012P}{II.\ref*{cdp:eq:C.TEST.2}}}
\expandafter\def\csname SRef@cdp:eq:C.3\endcsname{\hyperlink{tag-0094}{II.\ref*{cdp:eq:C.3}}}
\expandafter\def\csname SRef@cdp:eq:C.4\endcsname{\hyperlink{tag-0095}{II.\ref*{cdp:eq:C.4}}}
\expandafter\def\csname SRef@cdp:eq:C.5\endcsname{\hyperlink{tag-0096}{II.\ref*{cdp:eq:C.5}}}
\expandafter\def\csname SRef@cdp:eq:C.6\endcsname{\hyperlink{tag-0097}{II.\ref*{cdp:eq:C.6}}}
\expandafter\def\csname SRef@cdp:eq:C.7\endcsname{\hyperlink{tag-0098}{II.\ref*{cdp:eq:C.7}}}
\expandafter\def\csname SRef@cdp:eq:C.8\endcsname{\hyperlink{tag-0099}{II.\ref*{cdp:eq:C.8}}}
\expandafter\def\csname SRef@cdp:eq:C.PK.2\endcsname{\hyperlink{tag-00V9}{II.\ref*{cdp:eq:C.PK.2}}}
\expandafter\def\csname SRef@cdp:eq:C.9\endcsname{\hyperlink{tag-009A}{II.\ref*{cdp:eq:C.9}}}
\expandafter\def\csname SRef@cdp:eq:C.10\endcsname{\hyperlink{tag-008O}{II.\ref*{cdp:eq:C.10}}}
\expandafter\def\csname SRef@cdp:eq:C.11\endcsname{\hyperlink{tag-008P}{II.\ref*{cdp:eq:C.11}}}
\expandafter\def\csname SRef@cdp:eq:C.12\endcsname{\hyperlink{tag-008Q}{II.\ref*{cdp:eq:C.12}}}
\expandafter\def\csname SRef@cdp:eq:C.13\endcsname{\hyperlink{tag-008R}{II.\ref*{cdp:eq:C.13}}}
\expandafter\def\csname SRef@cdp:eq:C.14\endcsname{\hyperlink{tag-008S}{II.\ref*{cdp:eq:C.14}}}
\expandafter\def\csname SRef@cdp:eq:D.18\endcsname{\hyperlink{tag-0195}{II.\ref*{cdp:eq:D.18}}}
\expandafter\def\csname SRef@cdp:eq:D.19\endcsname{\hyperlink{tag-0196}{II.\ref*{cdp:eq:D.19}}}
\expandafter\def\csname SRef@cdp:eq:D.S.1\endcsname{\hyperlink{tag-019O}{II.\ref*{cdp:eq:D.S.1}}}
\expandafter\def\csname SRef@cdp:eq:C.PK.3\endcsname{\hyperlink{tag-00VA}{II.\ref*{cdp:eq:C.PK.3}}}
\expandafter\def\csname SRef@cdp:eq:D.S.2\endcsname{\hyperlink{tag-019P}{II.\ref*{cdp:eq:D.S.2}}}
\expandafter\def\csname SRef@cdp:eq:D.S.3\endcsname{\hyperlink{tag-019Q}{II.\ref*{cdp:eq:D.S.3}}}
\expandafter\def\csname SRef@cdp:eq:D.S.4\endcsname{\hyperlink{tag-019R}{II.\ref*{cdp:eq:D.S.4}}}
\expandafter\def\csname SRef@cdp:eq:D.S.5\endcsname{\hyperlink{tag-019S}{II.\ref*{cdp:eq:D.S.5}}}
\expandafter\def\csname SRef@cdp:eq:D.20\endcsname{\hyperlink{tag-0197}{II.\ref*{cdp:eq:D.20}}}
\expandafter\def\csname SRef@cdp:eq:D.21\endcsname{\hyperlink{tag-0198}{II.\ref*{cdp:eq:D.21}}}
\expandafter\def\csname SRef@cdp:eq:D.22\endcsname{\hyperlink{tag-0199}{II.\ref*{cdp:eq:D.22}}}
\expandafter\def\csname SRef@cdp:eq:D.23\endcsname{\hyperlink{tag-019A}{II.\ref*{cdp:eq:D.23}}}
\expandafter\def\csname SRef@cdp:eq:D.24\endcsname{\hyperlink{tag-019B}{II.\ref*{cdp:eq:D.24}}}
\expandafter\def\csname SRef@cdp:eq:C.PK.4\endcsname{\hyperlink{tag-00VB}{II.\ref*{cdp:eq:C.PK.4}}}
\expandafter\def\csname SRef@cdp:eq:D.25\endcsname{\hyperlink{tag-019C}{II.\ref*{cdp:eq:D.25}}}
\expandafter\def\csname SRef@cdp:eq:D.26\endcsname{\hyperlink{tag-019D}{II.\ref*{cdp:eq:D.26}}}
\expandafter\def\csname SRef@cdp:eq:D.27\endcsname{\hyperlink{tag-019E}{II.\ref*{cdp:eq:D.27}}}
\expandafter\def\csname SRef@cdp:eq:D.28\endcsname{\hyperlink{tag-019F}{II.\ref*{cdp:eq:D.28}}}
\expandafter\def\csname SRef@cdp:eq:D.29\endcsname{\hyperlink{tag-019G}{II.\ref*{cdp:eq:D.29}}}
\expandafter\def\csname SRef@cdp:eq:D.30\endcsname{\hyperlink{tag-019H}{II.\ref*{cdp:eq:D.30}}}
\expandafter\def\csname SRef@cdp:eq:D.31\endcsname{\hyperlink{tag-019I}{II.\ref*{cdp:eq:D.31}}}
\expandafter\def\csname SRef@cdp:eq:D.32\endcsname{\hyperlink{tag-019J}{II.\ref*{cdp:eq:D.32}}}
\expandafter\def\csname SRef@cdp:eq:D.33\endcsname{\hyperlink{tag-019K}{II.\ref*{cdp:eq:D.33}}}
\expandafter\def\csname SRef@cdp:eq:C.PK.5\endcsname{\hyperlink{tag-00VC}{II.\ref*{cdp:eq:C.PK.5}}}
\expandafter\def\csname SRef@cdp:eq:D.34\endcsname{\hyperlink{tag-019L}{II.\ref*{cdp:eq:D.34}}}
\expandafter\def\csname SRef@cdp:eq:D.35\endcsname{\hyperlink{tag-019M}{II.\ref*{cdp:eq:D.35}}}
\expandafter\def\csname SRef@cdp:eq:D.36\endcsname{\hyperlink{tag-019N}{II.\ref*{cdp:eq:D.36}}}
\expandafter\def\csname SRef@cdp:eq:E.16.1\endcsname{\hyperlink{tag-01A2}{II.\ref*{cdp:eq:E.16.1}}}
\expandafter\def\csname SRef@cdp:eq:E.16.2\endcsname{\hyperlink{tag-01A3}{II.\ref*{cdp:eq:E.16.2}}}
\expandafter\def\csname SRef@cdp:eq:E.16.3\endcsname{\hyperlink{tag-01A4}{II.\ref*{cdp:eq:E.16.3}}}
\expandafter\def\csname SRef@cdp:eq:C.PK.6\endcsname{\hyperlink{tag-00VD}{II.\ref*{cdp:eq:C.PK.6}}}
\expandafter\def\csname SRef@cdp:eq:C.PK.7\endcsname{\hyperlink{tag-00VE}{II.\ref*{cdp:eq:C.PK.7}}}
\expandafter\def\csname SRef@cdp:eq:C.PK.8\endcsname{\hyperlink{tag-00VF}{II.\ref*{cdp:eq:C.PK.8}}}
\expandafter\def\csname SRef@cdp:eq:C.IN.1\endcsname{\hyperlink{tag-00NK}{II.\ref*{cdp:eq:C.IN.1}}}
\expandafter\def\csname SRef@invII:eq:E.2\endcsname{\hyperlink{tag-01TR}{II.\ref*{invII:eq:E.2}}}
\expandafter\def\csname SRef@cdp:eq:C.IN.2\endcsname{\hyperlink{tag-00NL}{II.\ref*{cdp:eq:C.IN.2}}}
\expandafter\def\csname SRef@invII:eq:E.3\endcsname{\hyperlink{tag-01TS}{II.\ref*{invII:eq:E.3}}}
\expandafter\def\csname SRef@invII:eq:E.4a\endcsname{\hyperlink{tag-01TT}{II.\ref*{invII:eq:E.4a}}}
\expandafter\def\csname SRef@invII:eq:E.5\endcsname{\hyperlink{tag-01TU}{II.\ref*{invII:eq:E.5}}}
\expandafter\def\csname SRef@invII:eq:E.6\endcsname{\hyperlink{tag-01TV}{II.\ref*{invII:eq:E.6}}}
\expandafter\def\csname SRef@inv:wb:eq:9.19\endcsname{\hyperlink{tag-01TJ}{II.\ref*{inv:wb:eq:9.19}}}
\expandafter\def\csname SRef@inv:wb:eq:9.20\endcsname{\hyperlink{tag-01TK}{II.\ref*{inv:wb:eq:9.20}}}
\expandafter\def\csname SRef@inv:wb:eq:9.21\endcsname{\hyperlink{tag-01TL}{II.\ref*{inv:wb:eq:9.21}}}
\expandafter\def\csname SRef@inv:wb:eq:9.22\endcsname{\hyperlink{tag-01TM}{II.\ref*{inv:wb:eq:9.22}}}
\expandafter\def\csname SRef@inv:wb:eq:9.23\endcsname{\hyperlink{tag-01TN}{II.\ref*{inv:wb:eq:9.23}}}
\expandafter\def\csname SRef@inv:wb:eq:9.24\endcsname{\hyperlink{tag-01TO}{II.\ref*{inv:wb:eq:9.24}}}
\expandafter\def\csname SRef@cdp:eq:C.IN.3\endcsname{\hyperlink{tag-00NM}{II.\ref*{cdp:eq:C.IN.3}}}
\expandafter\def\csname SRef@inv:wb:eq:9.25\endcsname{\hyperlink{tag-01TP}{II.\ref*{inv:wb:eq:9.25}}}
\expandafter\def\csname SRef@wbII:eq:CN:2.2\endcsname{\hyperlink{tag-03DW}{II.\ref*{wbII:eq:CN:2.2}}}
\expandafter\def\csname SRef@wbII:eq:W:5.3\endcsname{\hyperlink{tag-03HS}{II.\ref*{wbII:eq:W:5.3}}}
\expandafter\def\csname SRef@wbII:eq:W:5.4\endcsname{\hyperlink{tag-03HX}{II.\ref*{wbII:eq:W:5.4}}}
\expandafter\def\csname SRef@wbII:eq:W:5.5\endcsname{\hyperlink{tag-03HY}{II.\ref*{wbII:eq:W:5.5}}}
\expandafter\def\csname SRef@wbII:eq:W:5.6\endcsname{\hyperlink{tag-03HZ}{II.\ref*{wbII:eq:W:5.6}}}
\expandafter\def\csname SRef@cdp:eq:C.IN.4\endcsname{\hyperlink{tag-00NN}{II.\ref*{cdp:eq:C.IN.4}}}
\expandafter\def\csname SRef@wbII:eq:W:5.7\endcsname{\hyperlink{tag-03I0}{II.\ref*{wbII:eq:W:5.7}}}
\expandafter\def\csname SRef@wbII:eq:W:5.8\endcsname{\hyperlink{tag-03I1}{II.\ref*{wbII:eq:W:5.8}}}
\expandafter\def\csname SRef@wbII:eq:W:5.9\endcsname{\hyperlink{tag-03I2}{II.\ref*{wbII:eq:W:5.9}}}
\expandafter\def\csname SRef@wbII:eq:W:5.10\endcsname{\hyperlink{tag-03EA}{II.\ref*{wbII:eq:W:5.10}}}
\expandafter\def\csname SRef@wbII:eq:CN:4.3\endcsname{\hyperlink{tag-03DX}{II.\ref*{wbII:eq:CN:4.3}}}
\expandafter\def\csname SRef@wbII:eq:CN:4.4\endcsname{\hyperlink{tag-03DY}{II.\ref*{wbII:eq:CN:4.4}}}
\expandafter\def\csname SRef@wbII:eq:CN:4.5a\endcsname{\hyperlink{tag-03E0}{II.\ref*{wbII:eq:CN:4.5a}}}
\expandafter\def\csname SRef@cdp:eq:C.IN.5\endcsname{\hyperlink{tag-00NO}{II.\ref*{cdp:eq:C.IN.5}}}
\expandafter\def\csname SRef@wbII:eq:CN:4.5\endcsname{\hyperlink{tag-03DZ}{II.\ref*{wbII:eq:CN:4.5}}}
\expandafter\def\csname SRef@wbII:eq:CN:4.6\endcsname{\hyperlink{tag-03E1}{II.\ref*{wbII:eq:CN:4.6}}}
\expandafter\def\csname SRef@wbII:eq:CN:4.7\endcsname{\hyperlink{tag-03E2}{II.\ref*{wbII:eq:CN:4.7}}}
\expandafter\def\csname SRef@wbII:eq:CN:4.8\endcsname{\hyperlink{tag-03E3}{II.\ref*{wbII:eq:CN:4.8}}}
\expandafter\def\csname SRef@wbII:eq:CN:4.9\endcsname{\hyperlink{tag-03E4}{II.\ref*{wbII:eq:CN:4.9}}}
\expandafter\def\csname SRef@wbII:eq:W:5.11\endcsname{\hyperlink{tag-03EB}{II.\ref*{wbII:eq:W:5.11}}}
\expandafter\def\csname SRef@wbII:eq:W:5.15a\endcsname{\hyperlink{tag-03EC}{II.\ref*{wbII:eq:W:5.15a}}}
\expandafter\def\csname SRef@wbII:eq:W:5.16\endcsname{\hyperlink{tag-03ED}{II.\ref*{wbII:eq:W:5.16}}}
\expandafter\def\csname SRef@wbII:eq:W:5.17\endcsname{\hyperlink{tag-03EE}{II.\ref*{wbII:eq:W:5.17}}}
\expandafter\def\csname SRef@cdp:eq:C.IN.6\endcsname{\hyperlink{tag-00NP}{II.\ref*{cdp:eq:C.IN.6}}}
\expandafter\def\csname SRef@wbII:eq:W:5.17-V\endcsname{\hyperlink{tag-03FU}{II.\ref*{wbII:eq:W:5.17-V}}}
\expandafter\def\csname SRef@wbII:eq:W:5.17-T\endcsname{\hyperlink{tag-03FS}{II.\ref*{wbII:eq:W:5.17-T}}}
\expandafter\def\csname SRef@wbII:eq:W:5.17-U\endcsname{\hyperlink{tag-03FT}{II.\ref*{wbII:eq:W:5.17-U}}}
\expandafter\def\csname SRef@wbII:eq:W:5.17a\endcsname{\hyperlink{tag-03FV}{II.\ref*{wbII:eq:W:5.17a}}}
\expandafter\def\csname SRef@wbII:eq:W:5.17b\endcsname{\hyperlink{tag-03GL}{II.\ref*{wbII:eq:W:5.17b}}}
\expandafter\def\csname SRef@wbII:eq:W:5.17c\endcsname{\hyperlink{tag-03GM}{II.\ref*{wbII:eq:W:5.17c}}}
\expandafter\def\csname SRef@wbII:eq:W:5.17d\endcsname{\hyperlink{tag-03GN}{II.\ref*{wbII:eq:W:5.17d}}}
\expandafter\def\csname SRef@wbII:eq:W:5.17e\endcsname{\hyperlink{tag-03GO}{II.\ref*{wbII:eq:W:5.17e}}}
\expandafter\def\csname SRef@wbII:eq:W:5.17g\endcsname{\hyperlink{tag-03GQ}{II.\ref*{wbII:eq:W:5.17g}}}
\expandafter\def\csname SRef@wbII:eq:W:5.17h\endcsname{\hyperlink{tag-03GR}{II.\ref*{wbII:eq:W:5.17h}}}
\expandafter\def\csname SRef@cdp:eq:C.PK.9\endcsname{\hyperlink{tag-00VG}{II.\ref*{cdp:eq:C.PK.9}}}
\expandafter\def\csname SRef@wbII:eq:W:5.17i\endcsname{\hyperlink{tag-03GS}{II.\ref*{wbII:eq:W:5.17i}}}
\expandafter\def\csname SRef@wbII:eq:W:5.17j'\endcsname{\hyperlink{tag-03GU}{II.\ref*{wbII:eq:W:5.17j'}}}
\expandafter\def\csname SRef@wbII:eq:W:5.17j\endcsname{\hyperlink{tag-03GT}{II.\ref*{wbII:eq:W:5.17j}}}
\expandafter\def\csname SRef@wbII:eq:W:5.17k\endcsname{\hyperlink{tag-03GV}{II.\ref*{wbII:eq:W:5.17k}}}
\expandafter\def\csname SRef@wbII:eq:W:5.17l\endcsname{\hyperlink{tag-03GW}{II.\ref*{wbII:eq:W:5.17l}}}
\expandafter\def\csname SRef@wbII:eq:W:5.17l'\endcsname{\hyperlink{tag-03GX}{II.\ref*{wbII:eq:W:5.17l'}}}
\expandafter\def\csname SRef@wbII:eq:W:5.17l''\endcsname{\hyperlink{tag-03GY}{II.\ref*{wbII:eq:W:5.17l''}}}
\expandafter\def\csname SRef@wbII:eq:W:5.17l'''\endcsname{\hyperlink{tag-03GZ}{II.\ref*{wbII:eq:W:5.17l'''}}}
\expandafter\def\csname SRef@wbII:eq:W:5.17m\endcsname{\hyperlink{tag-03H0}{II.\ref*{wbII:eq:W:5.17m}}}
\expandafter\def\csname SRef@wbII:eq:W:5.17m'\endcsname{\hyperlink{tag-03H1}{II.\ref*{wbII:eq:W:5.17m'}}}
\expandafter\def\csname SRef@cdp:eq:C.H.1\endcsname{\hyperlink{tag-00L8}{II.\ref*{cdp:eq:C.H.1}}}
\expandafter\def\csname SRef@wbII:eq:W:5.17m''\endcsname{\hyperlink{tag-03H2}{II.\ref*{wbII:eq:W:5.17m''}}}
\expandafter\def\csname SRef@wbII:eq:W:5.17n\endcsname{\hyperlink{tag-03H3}{II.\ref*{wbII:eq:W:5.17n}}}
\expandafter\def\csname SRef@wbII:eq:W:5.17o\endcsname{\hyperlink{tag-03H4}{II.\ref*{wbII:eq:W:5.17o}}}
\expandafter\def\csname SRef@wbII:eq:W:5.17p\endcsname{\hyperlink{tag-03H5}{II.\ref*{wbII:eq:W:5.17p}}}
\expandafter\def\csname SRef@wbII:eq:W:5.17p0\endcsname{\hyperlink{tag-03H7}{II.\ref*{wbII:eq:W:5.17p0}}}
\expandafter\def\csname SRef@wbII:eq:W:5.17p'\endcsname{\hyperlink{tag-03H6}{II.\ref*{wbII:eq:W:5.17p'}}}
\expandafter\def\csname SRef@wbII:eq:W:5.17q\endcsname{\hyperlink{tag-03H8}{II.\ref*{wbII:eq:W:5.17q}}}
\expandafter\def\csname SRef@wbII:eq:W:5.17r\endcsname{\hyperlink{tag-03H9}{II.\ref*{wbII:eq:W:5.17r}}}
\expandafter\def\csname SRef@wbII:eq:W:5.17s\endcsname{\hyperlink{tag-03HA}{II.\ref*{wbII:eq:W:5.17s}}}
\expandafter\def\csname SRef@wbII:eq:W:5.17t\endcsname{\hyperlink{tag-03HB}{II.\ref*{wbII:eq:W:5.17t}}}
\expandafter\def\csname SRef@cdp:eq:A.1\endcsname{\hyperlink{tag-003Y}{II.\ref*{cdp:eq:A.1}}}
\expandafter\def\csname SRef@wbII:eq:W:5.17u\endcsname{\hyperlink{tag-03HC}{II.\ref*{wbII:eq:W:5.17u}}}
\expandafter\def\csname SRef@wbII:eq:W:5.17v\endcsname{\hyperlink{tag-03HD}{II.\ref*{wbII:eq:W:5.17v}}}
\expandafter\def\csname SRef@wbII:eq:W:5.17w\endcsname{\hyperlink{tag-03HE}{II.\ref*{wbII:eq:W:5.17w}}}
\expandafter\def\csname SRef@wbII:eq:W:5.17x\endcsname{\hyperlink{tag-03HF}{II.\ref*{wbII:eq:W:5.17x}}}
\expandafter\def\csname SRef@wbII:eq:W:5.17y\endcsname{\hyperlink{tag-03HG}{II.\ref*{wbII:eq:W:5.17y}}}
\expandafter\def\csname SRef@wbII:eq:W:5.17z\endcsname{\hyperlink{tag-03HH}{II.\ref*{wbII:eq:W:5.17z}}}
\expandafter\def\csname SRef@wbII:eq:W:5.17aa\endcsname{\hyperlink{tag-03FW}{II.\ref*{wbII:eq:W:5.17aa}}}
\expandafter\def\csname SRef@wbII:eq:W:5.17ab\endcsname{\hyperlink{tag-03FX}{II.\ref*{wbII:eq:W:5.17ab}}}
\expandafter\def\csname SRef@wbII:eq:W:5.17ac\endcsname{\hyperlink{tag-03FY}{II.\ref*{wbII:eq:W:5.17ac}}}
\expandafter\def\csname SRef@wbII:eq:W:5.17ad\endcsname{\hyperlink{tag-03FZ}{II.\ref*{wbII:eq:W:5.17ad}}}
\expandafter\def\csname SRef@cdp:eq:A.2\endcsname{\hyperlink{tag-004B}{II.\ref*{cdp:eq:A.2}}}
\expandafter\def\csname SRef@wbII:eq:W:5.17ae\endcsname{\hyperlink{tag-03G0}{II.\ref*{wbII:eq:W:5.17ae}}}
\expandafter\def\csname SRef@wbII:eq:W:5.17af\endcsname{\hyperlink{tag-03G1}{II.\ref*{wbII:eq:W:5.17af}}}
\expandafter\def\csname SRef@wbII:eq:W:5.17ag\endcsname{\hyperlink{tag-03G2}{II.\ref*{wbII:eq:W:5.17ag}}}
\expandafter\def\csname SRef@wbII:eq:W:5.17ah\endcsname{\hyperlink{tag-03G3}{II.\ref*{wbII:eq:W:5.17ah}}}
\expandafter\def\csname SRef@wbII:eq:W:5.17ai\endcsname{\hyperlink{tag-03G4}{II.\ref*{wbII:eq:W:5.17ai}}}
\expandafter\def\csname SRef@wbII:eq:W:5.17aj\endcsname{\hyperlink{tag-03G5}{II.\ref*{wbII:eq:W:5.17aj}}}
\expandafter\def\csname SRef@wbII:eq:W:5.17ak\endcsname{\hyperlink{tag-03G6}{II.\ref*{wbII:eq:W:5.17ak}}}
\expandafter\def\csname SRef@wbII:eq:W:5.17al\endcsname{\hyperlink{tag-03G7}{II.\ref*{wbII:eq:W:5.17al}}}
\expandafter\def\csname SRef@cdp:eq:A.3\endcsname{\hyperlink{tag-004M}{II.\ref*{cdp:eq:A.3}}}
\expandafter\def\csname SRef@wbII:eq:W:5.17am\endcsname{\hyperlink{tag-03G8}{II.\ref*{wbII:eq:W:5.17am}}}
\expandafter\def\csname SRef@wbII:eq:W:5.17an\endcsname{\hyperlink{tag-03G9}{II.\ref*{wbII:eq:W:5.17an}}}
\expandafter\def\csname SRef@wbII:eq:W:5.17ao\endcsname{\hyperlink{tag-03GA}{II.\ref*{wbII:eq:W:5.17ao}}}
\expandafter\def\csname SRef@wbII:eq:W:5.17ap\endcsname{\hyperlink{tag-03GB}{II.\ref*{wbII:eq:W:5.17ap}}}
\expandafter\def\csname SRef@wbII:eq:W:5.17aq\endcsname{\hyperlink{tag-03GC}{II.\ref*{wbII:eq:W:5.17aq}}}
\expandafter\def\csname SRef@wbII:eq:W:5.17ar\endcsname{\hyperlink{tag-03GD}{II.\ref*{wbII:eq:W:5.17ar}}}
\expandafter\def\csname SRef@wbII:eq:W:5.17as\endcsname{\hyperlink{tag-03GE}{II.\ref*{wbII:eq:W:5.17as}}}
\expandafter\def\csname SRef@wbII:eq:W:5.17at\endcsname{\hyperlink{tag-03GF}{II.\ref*{wbII:eq:W:5.17at}}}
\expandafter\def\csname SRef@wbII:eq:W:5.17au\endcsname{\hyperlink{tag-03GG}{II.\ref*{wbII:eq:W:5.17au}}}
\expandafter\def\csname SRef@wbII:eq:W:5.17av\endcsname{\hyperlink{tag-03GH}{II.\ref*{wbII:eq:W:5.17av}}}
\expandafter\def\csname SRef@cdp:eq:A.4\endcsname{\hyperlink{tag-004O}{II.\ref*{cdp:eq:A.4}}}
\expandafter\def\csname SRef@wbII:eq:W:5.17aw\endcsname{\hyperlink{tag-03GI}{II.\ref*{wbII:eq:W:5.17aw}}}
\expandafter\def\csname SRef@wbII:eq:W:5.17ax\endcsname{\hyperlink{tag-03GJ}{II.\ref*{wbII:eq:W:5.17ax}}}
\expandafter\def\csname SRef@wbII:eq:W:5.17ay\endcsname{\hyperlink{tag-03GK}{II.\ref*{wbII:eq:W:5.17ay}}}
\expandafter\def\csname SRef@wbII:eq:W:5.17-A1\endcsname{\hyperlink{tag-03EF}{II.\ref*{wbII:eq:W:5.17-A1}}}
\expandafter\def\csname SRef@wbII:eq:W:5.17-A2\endcsname{\hyperlink{tag-03EY}{II.\ref*{wbII:eq:W:5.17-A2}}}
\expandafter\def\csname SRef@wbII:eq:W:5.17-A3\endcsname{\hyperlink{tag-03F9}{II.\ref*{wbII:eq:W:5.17-A3}}}
\expandafter\def\csname SRef@wbII:eq:W:5.17-A4\endcsname{\hyperlink{tag-03FM}{II.\ref*{wbII:eq:W:5.17-A4}}}
\expandafter\def\csname SRef@wbII:eq:W:5.17-A5\endcsname{\hyperlink{tag-03FN}{II.\ref*{wbII:eq:W:5.17-A5}}}
\expandafter\def\csname SRef@wbII:eq:W:5.17-A6\endcsname{\hyperlink{tag-03FO}{II.\ref*{wbII:eq:W:5.17-A6}}}
\expandafter\def\csname SRef@wbII:eq:W:5.17-A7\endcsname{\hyperlink{tag-03FP}{II.\ref*{wbII:eq:W:5.17-A7}}}
\expandafter\def\csname SRef@wbII:eq:W:5.17-A8\endcsname{\hyperlink{tag-03FQ}{II.\ref*{wbII:eq:W:5.17-A8}}}
\expandafter\def\csname SRef@wbII:eq:W:5.17-A9\endcsname{\hyperlink{tag-03FR}{II.\ref*{wbII:eq:W:5.17-A9}}}
\expandafter\def\csname SRef@wbII:eq:W:5.17-A10\endcsname{\hyperlink{tag-03EG}{II.\ref*{wbII:eq:W:5.17-A10}}}
\expandafter\def\csname SRef@wbII:eq:W:5.17-A11\endcsname{\hyperlink{tag-03EH}{II.\ref*{wbII:eq:W:5.17-A11}}}
\expandafter\def\csname SRef@wbII:eq:W:5.17-A12\endcsname{\hyperlink{tag-03EI}{II.\ref*{wbII:eq:W:5.17-A12}}}
\expandafter\def\csname SRef@wbII:eq:W:5.17-A13\endcsname{\hyperlink{tag-03EJ}{II.\ref*{wbII:eq:W:5.17-A13}}}
\expandafter\def\csname SRef@wbII:eq:W:5.17-A14\endcsname{\hyperlink{tag-03EK}{II.\ref*{wbII:eq:W:5.17-A14}}}
\expandafter\def\csname SRef@wbII:eq:W:5.17-A14a\endcsname{\hyperlink{tag-03EL}{II.\ref*{wbII:eq:W:5.17-A14a}}}
\expandafter\def\csname SRef@cdp:eq:A.5\endcsname{\hyperlink{tag-004P}{II.\ref*{cdp:eq:A.5}}}
\expandafter\def\csname SRef@wbII:eq:W:5.17-A14c\endcsname{\hyperlink{tag-03EM}{II.\ref*{wbII:eq:W:5.17-A14c}}}
\expandafter\def\csname SRef@wbII:eq:W:5.17-A14d\endcsname{\hyperlink{tag-03EN}{II.\ref*{wbII:eq:W:5.17-A14d}}}
\expandafter\def\csname SRef@wbII:eq:W:5.17-A14e\endcsname{\hyperlink{tag-03EO}{II.\ref*{wbII:eq:W:5.17-A14e}}}
\expandafter\def\csname SRef@wbII:eq:W:5.17-A14f\endcsname{\hyperlink{tag-03EP}{II.\ref*{wbII:eq:W:5.17-A14f}}}
\expandafter\def\csname SRef@wbII:eq:W:5.17-A15\endcsname{\hyperlink{tag-03EQ}{II.\ref*{wbII:eq:W:5.17-A15}}}
\expandafter\def\csname SRef@wbII:eq:W:5.17-A16\endcsname{\hyperlink{tag-03ER}{II.\ref*{wbII:eq:W:5.17-A16}}}
\expandafter\def\csname SRef@wbII:eq:W:5.17-A17\endcsname{\hyperlink{tag-03ES}{II.\ref*{wbII:eq:W:5.17-A17}}}
\expandafter\def\csname SRef@wbII:eq:W:5.17-A18\endcsname{\hyperlink{tag-03ET}{II.\ref*{wbII:eq:W:5.17-A18}}}
\expandafter\def\csname SRef@wbII:eq:W:5.17-A19\endcsname{\hyperlink{tag-03EU}{II.\ref*{wbII:eq:W:5.17-A19}}}
\expandafter\def\csname SRef@wbII:eq:W:5.17-A19a\endcsname{\hyperlink{tag-03EV}{II.\ref*{wbII:eq:W:5.17-A19a}}}
\expandafter\def\csname SRef@cdp:eq:A.6\endcsname{\hyperlink{tag-004Q}{II.\ref*{cdp:eq:A.6}}}
\expandafter\def\csname SRef@wbII:eq:W:5.17-A19b\endcsname{\hyperlink{tag-03EW}{II.\ref*{wbII:eq:W:5.17-A19b}}}
\expandafter\def\csname SRef@wbII:eq:W:5.17-A19c\endcsname{\hyperlink{tag-03EX}{II.\ref*{wbII:eq:W:5.17-A19c}}}
\expandafter\def\csname SRef@wbII:eq:W:5.17-A20\endcsname{\hyperlink{tag-03EZ}{II.\ref*{wbII:eq:W:5.17-A20}}}
\expandafter\def\csname SRef@wbII:eq:W:5.17-A21\endcsname{\hyperlink{tag-03F0}{II.\ref*{wbII:eq:W:5.17-A21}}}
\expandafter\def\csname SRef@wbII:eq:W:5.17-A22\endcsname{\hyperlink{tag-03F1}{II.\ref*{wbII:eq:W:5.17-A22}}}
\expandafter\def\csname SRef@wbII:eq:W:5.17-A23\endcsname{\hyperlink{tag-03F2}{II.\ref*{wbII:eq:W:5.17-A23}}}
\expandafter\def\csname SRef@wbII:eq:W:5.17-A24\endcsname{\hyperlink{tag-03F3}{II.\ref*{wbII:eq:W:5.17-A24}}}
\expandafter\def\csname SRef@wbII:eq:W:5.17-A25\endcsname{\hyperlink{tag-03F4}{II.\ref*{wbII:eq:W:5.17-A25}}}
\expandafter\def\csname SRef@wbII:eq:W:5.17-A26\endcsname{\hyperlink{tag-03F5}{II.\ref*{wbII:eq:W:5.17-A26}}}
\expandafter\def\csname SRef@wbII:eq:W:5.17-A27\endcsname{\hyperlink{tag-03F6}{II.\ref*{wbII:eq:W:5.17-A27}}}
\expandafter\def\csname SRef@cdp:eq:A.7\endcsname{\hyperlink{tag-004R}{II.\ref*{cdp:eq:A.7}}}
\expandafter\def\csname SRef@wbII:eq:W:5.17-A28\endcsname{\hyperlink{tag-03F7}{II.\ref*{wbII:eq:W:5.17-A28}}}
\expandafter\def\csname SRef@wbII:eq:W:5.17-A29\endcsname{\hyperlink{tag-03F8}{II.\ref*{wbII:eq:W:5.17-A29}}}
\expandafter\def\csname SRef@wbII:eq:W:5.17-A30\endcsname{\hyperlink{tag-03FA}{II.\ref*{wbII:eq:W:5.17-A30}}}
\expandafter\def\csname SRef@wbII:eq:W:5.17-A31\endcsname{\hyperlink{tag-03FB}{II.\ref*{wbII:eq:W:5.17-A31}}}
\expandafter\def\csname SRef@wbII:eq:W:5.17-A32\endcsname{\hyperlink{tag-03FC}{II.\ref*{wbII:eq:W:5.17-A32}}}
\expandafter\def\csname SRef@wbII:eq:W:5.17-A33\endcsname{\hyperlink{tag-03FD}{II.\ref*{wbII:eq:W:5.17-A33}}}
\expandafter\def\csname SRef@wbII:eq:W:5.17-A34\endcsname{\hyperlink{tag-03FE}{II.\ref*{wbII:eq:W:5.17-A34}}}
\expandafter\def\csname SRef@wbII:eq:W:5.17-A34a\endcsname{\hyperlink{tag-03FF}{II.\ref*{wbII:eq:W:5.17-A34a}}}
\expandafter\def\csname SRef@wbII:eq:W:5.17-A34b\endcsname{\hyperlink{tag-03FG}{II.\ref*{wbII:eq:W:5.17-A34b}}}
\expandafter\def\csname SRef@wbII:eq:W:5.17-A34c\endcsname{\hyperlink{tag-03FH}{II.\ref*{wbII:eq:W:5.17-A34c}}}
\expandafter\def\csname SRef@cdp:eq:A.8\endcsname{\hyperlink{tag-004S}{II.\ref*{cdp:eq:A.8}}}
\expandafter\def\csname SRef@wbII:eq:W:5.17-A35\endcsname{\hyperlink{tag-03FI}{II.\ref*{wbII:eq:W:5.17-A35}}}
\expandafter\def\csname SRef@wbII:eq:W:5.17-A36\endcsname{\hyperlink{tag-03FJ}{II.\ref*{wbII:eq:W:5.17-A36}}}
\expandafter\def\csname SRef@wbII:eq:W:5.17-A37\endcsname{\hyperlink{tag-03FK}{II.\ref*{wbII:eq:W:5.17-A37}}}
\expandafter\def\csname SRef@wbII:eq:W:5.17-A38\endcsname{\hyperlink{tag-03FL}{II.\ref*{wbII:eq:W:5.17-A38}}}
\expandafter\def\csname SRef@wbII:eq:W:5.17f\endcsname{\hyperlink{tag-03GP}{II.\ref*{wbII:eq:W:5.17f}}}
\expandafter\def\csname SRef@wbII:eq:W:5.18\endcsname{\hyperlink{tag-03HI}{II.\ref*{wbII:eq:W:5.18}}}
\expandafter\def\csname SRef@wbII:eq:W:5.19\endcsname{\hyperlink{tag-03HJ}{II.\ref*{wbII:eq:W:5.19}}}
\expandafter\def\csname SRef@wbII:eq:W:5.20\endcsname{\hyperlink{tag-03HK}{II.\ref*{wbII:eq:W:5.20}}}
\expandafter\def\csname SRef@wbII:eq:CN:5.2\endcsname{\hyperlink{tag-03E5}{II.\ref*{wbII:eq:CN:5.2}}}
\expandafter\def\csname SRef@cdp:eq:A.9\endcsname{\hyperlink{tag-004T}{II.\ref*{cdp:eq:A.9}}}
\expandafter\def\csname SRef@wbII:eq:W:5.23\endcsname{\hyperlink{tag-03HL}{II.\ref*{wbII:eq:W:5.23}}}
\expandafter\def\csname SRef@wbII:eq:W:5.24\endcsname{\hyperlink{tag-03HM}{II.\ref*{wbII:eq:W:5.24}}}
\expandafter\def\csname SRef@wbII:eq:W:5.25\endcsname{\hyperlink{tag-03HN}{II.\ref*{wbII:eq:W:5.25}}}
\expandafter\def\csname SRef@wbII:eq:W:5.26\endcsname{\hyperlink{tag-03HO}{II.\ref*{wbII:eq:W:5.26}}}
\expandafter\def\csname SRef@wbII:eq:W:5.27\endcsname{\hyperlink{tag-03HP}{II.\ref*{wbII:eq:W:5.27}}}
\expandafter\def\csname SRef@wbII:eq:W:5.28\endcsname{\hyperlink{tag-03HQ}{II.\ref*{wbII:eq:W:5.28}}}
\expandafter\def\csname SRef@wbII:eq:W:5.29\endcsname{\hyperlink{tag-03HR}{II.\ref*{wbII:eq:W:5.29}}}
\expandafter\def\csname SRef@wbII:eq:W:5.30\endcsname{\hyperlink{tag-03HT}{II.\ref*{wbII:eq:W:5.30}}}
\expandafter\def\csname SRef@wbII:eq:W:5.31\endcsname{\hyperlink{tag-03HU}{II.\ref*{wbII:eq:W:5.31}}}
\expandafter\def\csname SRef@wbII:eq:W:5.32\endcsname{\hyperlink{tag-03HV}{II.\ref*{wbII:eq:W:5.32}}}
\expandafter\def\csname SRef@cdp:eq:A.9a\endcsname{\hyperlink{tag-004U}{II.\ref*{cdp:eq:A.9a}}}
\expandafter\def\csname SRef@wbII:eq:W:5.33\endcsname{\hyperlink{tag-03HW}{II.\ref*{wbII:eq:W:5.33}}}
\expandafter\def\csname SRef@wbII:eq:CN:6.1\endcsname{\hyperlink{tag-03E6}{II.\ref*{wbII:eq:CN:6.1}}}
\expandafter\def\csname SRef@wbII:eq:CN:6.2\endcsname{\hyperlink{tag-03E7}{II.\ref*{wbII:eq:CN:6.2}}}
\expandafter\def\csname SRef@wbII:eq:CN:6.3\endcsname{\hyperlink{tag-03E8}{II.\ref*{wbII:eq:CN:6.3}}}
\expandafter\def\csname SRef@wbII:eq:CN:6.4\endcsname{\hyperlink{tag-03E9}{II.\ref*{wbII:eq:CN:6.4}}}
\expandafter\def\csname SRef@cdp:eq:A.9b\endcsname{\hyperlink{tag-004V}{II.\ref*{cdp:eq:A.9b}}}
\expandafter\def\csname SRef@cdp:eq:A.NL.1\endcsname{\hyperlink{tag-006J}{II.\ref*{cdp:eq:A.NL.1}}}
\expandafter\def\csname SRef@invII:eq:E.7\endcsname{\hyperlink{tag-01TW}{II.\ref*{invII:eq:E.7}}}
\expandafter\def\csname SRef@cdp:eq:A.NL.2\endcsname{\hyperlink{tag-006K}{II.\ref*{cdp:eq:A.NL.2}}}
\expandafter\def\csname SRef@cdp:eq:A.NL.3\endcsname{\hyperlink{tag-006L}{II.\ref*{cdp:eq:A.NL.3}}}
\expandafter\def\csname SRef@cdp:eq:A.NL.4\endcsname{\hyperlink{tag-006M}{II.\ref*{cdp:eq:A.NL.4}}}
\expandafter\def\csname SRef@cdp:eq:A.NL.5\endcsname{\hyperlink{tag-006N}{II.\ref*{cdp:eq:A.NL.5}}}
\expandafter\def\csname SRef@cdp:eq:B.1\endcsname{\hyperlink{tag-008G}{II.\ref*{cdp:eq:B.1}}}
\expandafter\def\csname SRef@cdp:eq:B.2\endcsname{\hyperlink{tag-008H}{II.\ref*{cdp:eq:B.2}}}
\expandafter\def\csname SRef@cdp:eq:B.3\endcsname{\hyperlink{tag-008I}{II.\ref*{cdp:eq:B.3}}}
\expandafter\def\csname SRef@cdp:eq:B.4\endcsname{\hyperlink{tag-008J}{II.\ref*{cdp:eq:B.4}}}
\expandafter\def\csname SRef@cdp:eq:B.5\endcsname{\hyperlink{tag-008K}{II.\ref*{cdp:eq:B.5}}}
\expandafter\def\csname SRef@cdp:eq:B.6\endcsname{\hyperlink{tag-008L}{II.\ref*{cdp:eq:B.6}}}
\expandafter\def\csname SRef@cdp:eq:B.7\endcsname{\hyperlink{tag-008M}{II.\ref*{cdp:eq:B.7}}}
\expandafter\def\csname SRef@cdp:eq:A.10\endcsname{\hyperlink{tag-003Z}{II.\ref*{cdp:eq:A.10}}}
\expandafter\def\csname SRef@cdp:eq:A.11\endcsname{\hyperlink{tag-0040}{II.\ref*{cdp:eq:A.11}}}
\expandafter\def\csname SRef@cdp:eq:A.12\endcsname{\hyperlink{tag-0041}{II.\ref*{cdp:eq:A.12}}}
\expandafter\def\csname SRef@cdp:eq:A.13\endcsname{\hyperlink{tag-0042}{II.\ref*{cdp:eq:A.13}}}
\expandafter\def\csname SRef@cdp:eq:A.14\endcsname{\hyperlink{tag-0043}{II.\ref*{cdp:eq:A.14}}}
\expandafter\def\csname SRef@cdp:eq:A.15a\endcsname{\hyperlink{tag-0045}{II.\ref*{cdp:eq:A.15a}}}
\expandafter\def\csname SRef@cdp:eq:A.15b\endcsname{\hyperlink{tag-0046}{II.\ref*{cdp:eq:A.15b}}}
\expandafter\def\csname SRef@cdp:eq:A.15\endcsname{\hyperlink{tag-0044}{II.\ref*{cdp:eq:A.15}}}
\expandafter\def\csname SRef@cdp:eq:A.16\endcsname{\hyperlink{tag-0047}{II.\ref*{cdp:eq:A.16}}}
\expandafter\def\csname SRef@cdp:eq:A.17\endcsname{\hyperlink{tag-0048}{II.\ref*{cdp:eq:A.17}}}
\expandafter\def\csname SRef@cdp:eq:A.18\endcsname{\hyperlink{tag-0049}{II.\ref*{cdp:eq:A.18}}}
\expandafter\def\csname SRef@cdp:eq:A.19\endcsname{\hyperlink{tag-004A}{II.\ref*{cdp:eq:A.19}}}
\expandafter\def\csname SRef@cdp:eq:A.20\endcsname{\hyperlink{tag-004C}{II.\ref*{cdp:eq:A.20}}}
\expandafter\def\csname SRef@cdp:eq:A.21\endcsname{\hyperlink{tag-004D}{II.\ref*{cdp:eq:A.21}}}
\expandafter\def\csname SRef@cdp:eq:A.22\endcsname{\hyperlink{tag-004E}{II.\ref*{cdp:eq:A.22}}}
\expandafter\def\csname SRef@cdp:eq:A.23\endcsname{\hyperlink{tag-004F}{II.\ref*{cdp:eq:A.23}}}
\expandafter\def\csname SRef@cdp:eq:A.24\endcsname{\hyperlink{tag-004G}{II.\ref*{cdp:eq:A.24}}}
\expandafter\def\csname SRef@cdp:eq:A.25\endcsname{\hyperlink{tag-004H}{II.\ref*{cdp:eq:A.25}}}
\expandafter\def\csname SRef@cdp:eq:A.26\endcsname{\hyperlink{tag-004I}{II.\ref*{cdp:eq:A.26}}}
\expandafter\def\csname SRef@cdp:eq:A.27\endcsname{\hyperlink{tag-004J}{II.\ref*{cdp:eq:A.27}}}
\expandafter\def\csname SRef@cdp:eq:A.28\endcsname{\hyperlink{tag-004K}{II.\ref*{cdp:eq:A.28}}}
\expandafter\def\csname SRef@cdp:eq:A.29\endcsname{\hyperlink{tag-004L}{II.\ref*{cdp:eq:A.29}}}
\expandafter\def\csname SRef@cdp:eq:A.30\endcsname{\hyperlink{tag-004N}{II.\ref*{cdp:eq:A.30}}}
\expandafter\def\csname SRef@cdp:eq:A.RC.1\endcsname{\hyperlink{tag-006U}{II.\ref*{cdp:eq:A.RC.1}}}
\expandafter\def\csname SRef@push:src:E.2.4\endcsname{\hyperlink{tag-033C}{II.\ref*{push:src:E.2.4}}}
\expandafter\def\csname SRef@push:src:E.2.24\endcsname{\hyperlink{tag-033A}{II.\ref*{push:src:E.2.24}}}
\expandafter\def\csname SRef@wbII:cor:surface-top-base-change\endcsname{\hyperlink{tag-03DV}{II.\ref*{wbII:cor:surface-top-base-change}}}
\expandafter\def\csname SRef@wbII:cor:candidate-pic0-route\endcsname{\hyperlink{tag-03DS}{II.\ref*{wbII:cor:candidate-pic0-route}}}
\expandafter\def\csname SRef@wbII:cor:candidate-conductor-actual\endcsname{\hyperlink{tag-03DR}{II.\ref*{wbII:cor:candidate-conductor-actual}}}
\expandafter\def\csname SRef@wbII:cor:corrected-cdp-relative\endcsname{\hyperlink{tag-03DT}{II.\ref*{wbII:cor:corrected-cdp-relative}}}
\expandafter\def\csname SRef@wbII:cor:relative\endcsname{\hyperlink{tag-03DT}{II.\ref*{wbII:cor:relative}}}
\expandafter\def\csname SRef@wbII:cor:normal-crossing-kernel\endcsname{\hyperlink{tag-03DU}{II.\ref*{wbII:cor:normal-crossing-kernel}}}
\expandafter\def\csname SRef@wbII:test:cdp-comparison\endcsname{\hyperlink{tag-03IV}{II.\ref*{wbII:test:cdp-comparison}}}
\expandafter\def\csname SRef@cdp:A:filtration1\endcsname{\hyperlink{tag-003W}{II.\ref*{cdp:A:filtration1}}}
\expandafter\def\csname SRef@cdp:eq:E13-residue\endcsname{\hyperlink{tag-01A6}{II.\ref*{cdp:eq:E13-residue}}}
\expandafter\def\csname SRef@cdp:eq:E14-elliptic-family\endcsname{\hyperlink{tag-01A8}{II.\ref*{cdp:eq:E14-elliptic-family}}}
\expandafter\def\csname SRef@cdp:eq:E14-canonical\endcsname{\hyperlink{tag-01A7}{II.\ref*{cdp:eq:E14-canonical}}}
\expandafter\def\csname SRef@cdp:eq:E14-log-coho\endcsname{\hyperlink{tag-01AA}{II.\ref*{cdp:eq:E14-log-coho}}}
\expandafter\def\csname SRef@cdp:eq:E14-jet-map\endcsname{\hyperlink{tag-01A9}{II.\ref*{cdp:eq:E14-jet-map}}}
\expandafter\def\csname SRef@cdp:eq:E15-positive\endcsname{\hyperlink{tag-01AB}{II.\ref*{cdp:eq:E15-positive}}}
\expandafter\def\csname SRef@cdp:eq:E15-products\endcsname{\hyperlink{tag-01AC}{II.\ref*{cdp:eq:E15-products}}}
\expandafter\def\csname SRef@wbII:eq:analytic-toric-open-directions\endcsname{\hyperlink{tag-03I5}{II.\ref*{wbII:eq:analytic-toric-open-directions}}}
\expandafter\def\csname SRef@wbII:eq:analytic-residual-descent-factor\endcsname{\hyperlink{tag-03I4}{II.\ref*{wbII:eq:analytic-residual-descent-factor}}}
\expandafter\def\csname SRef@wbII:eq:analytic-factor-horizontal-winding\endcsname{\hyperlink{tag-03I3}{II.\ref*{wbII:eq:analytic-factor-horizontal-winding}}}
\expandafter\def\csname SRef@cdp:A:filtration2\endcsname{\hyperlink{tag-003X}{II.\ref*{cdp:A:filtration2}}}
\expandafter\def\csname SRef@eq:P4chi-resolution\endcsname{\hyperlink{tag-01OD}{II.\ref*{eq:P4chi-resolution}}}
\expandafter\def\csname SRef@eq:P4exceptional-genus\endcsname{\hyperlink{tag-01OE}{II.\ref*{eq:P4exceptional-genus}}}
\expandafter\def\csname SRef@cdp:eq:E-Leray-map\endcsname{\hyperlink{tag-019Y}{II.\ref*{cdp:eq:E-Leray-map}}}
\expandafter\def\csname SRef@cdp:eq:E-surj-embedding\endcsname{\hyperlink{tag-01A0}{II.\ref*{cdp:eq:E-surj-embedding}}}
\expandafter\def\csname SRef@cdp:eq:E-log-period\endcsname{\hyperlink{tag-019Z}{II.\ref*{cdp:eq:E-log-period}}}
\expandafter\def\csname SRef@cdp:eq:E-theta-general\endcsname{\hyperlink{tag-01A1}{II.\ref*{cdp:eq:E-theta-general}}}
\expandafter\def\csname SRef@cdp:eq:E13-connection\endcsname{\hyperlink{tag-01A5}{II.\ref*{cdp:eq:E13-connection}}}
\expandafter\def\csname SRef@wbII:lem:torsion-identification\endcsname{\hyperlink{tag-03IF}{II.\ref*{wbII:lem:torsion-identification}}}
\expandafter\def\csname SRef@wbII:lem:picard-crossed-defect\endcsname{\hyperlink{tag-03IC}{II.\ref*{wbII:lem:picard-crossed-defect}}}
\expandafter\def\csname SRef@wbII:lem:picard-finite-norm\endcsname{\hyperlink{tag-03IE}{II.\ref*{wbII:lem:picard-finite-norm}}}
\expandafter\def\csname SRef@wbII:lem:picard-cusp-defect\endcsname{\hyperlink{tag-03ID}{II.\ref*{wbII:lem:picard-cusp-defect}}}
\expandafter\def\csname SRef@wbII:lem:analytic-support-specialization\endcsname{\hyperlink{tag-03I6}{II.\ref*{wbII:lem:analytic-support-specialization}}}
\expandafter\def\csname SRef@wbII:lem:generic-endpoints\endcsname{\hyperlink{tag-03IA}{II.\ref*{wbII:lem:generic-endpoints}}}
\expandafter\def\csname SRef@wbII:lem:cdp-generic-endpoints\endcsname{\hyperlink{tag-03I7}{II.\ref*{wbII:lem:cdp-generic-endpoints}}}
\expandafter\def\csname SRef@wbII:lem:gluing-bundle\endcsname{\hyperlink{tag-03IB}{II.\ref*{wbII:lem:gluing-bundle}}}
\expandafter\def\csname SRef@wbII:lem:extending-gluing\endcsname{\hyperlink{tag-03I9}{II.\ref*{wbII:lem:extending-gluing}}}
\expandafter\def\csname SRef@wbII:lem:curve-flatness\endcsname{\hyperlink{tag-03I8}{II.\ref*{wbII:lem:curve-flatness}}}
\expandafter\def\csname SRef@inv:E2.3\endcsname{\hyperlink{tag-01SW}{II.\ref*{inv:E2.3}}}
\expandafter\def\csname SRef@cdp:st:D.2.1\endcsname{\hyperlink{tag-01G8}{II.\ref*{cdp:st:D.2.1}}}
\expandafter\def\csname SRef@cdp:st:D.2.3\endcsname{\hyperlink{tag-01GA}{II.\ref*{cdp:st:D.2.3}}}
\expandafter\def\csname SRef@cdp:st:D.2.4\endcsname{\hyperlink{tag-01GB}{II.\ref*{cdp:st:D.2.4}}}
\expandafter\def\csname SRef@cdp:st:D.2.5\endcsname{\hyperlink{tag-01GC}{II.\ref*{cdp:st:D.2.5}}}
\expandafter\def\csname SRef@cdp:st:D.2.6\endcsname{\hyperlink{tag-01GD}{II.\ref*{cdp:st:D.2.6}}}
\expandafter\def\csname SRef@cdp:st:D.2.7\endcsname{\hyperlink{tag-01GE}{II.\ref*{cdp:st:D.2.7}}}
\expandafter\def\csname SRef@cdp:st:D.2.8\endcsname{\hyperlink{tag-01GF}{II.\ref*{cdp:st:D.2.8}}}
\expandafter\def\csname SRef@cdp:st:D.2.2\endcsname{\hyperlink{tag-01G9}{II.\ref*{cdp:st:D.2.2}}}
\expandafter\def\csname SRef@cdp:st:E5\endcsname{\hyperlink{tag-01GU}{II.\ref*{cdp:st:E5}}}
\expandafter\def\csname SRef@cdp:st:E15\endcsname{\hyperlink{tag-01GM}{II.\ref*{cdp:st:E15}}}
\expandafter\def\csname SRef@cdp:st:E16\endcsname{\hyperlink{tag-01GN}{II.\ref*{cdp:st:E16}}}
\expandafter\def\csname SRef@cdp:st:E17\endcsname{\hyperlink{tag-01GO}{II.\ref*{cdp:st:E17}}}
\expandafter\def\csname SRef@cdp:st:E18\endcsname{\hyperlink{tag-01GP}{II.\ref*{cdp:st:E18}}}
\expandafter\def\csname SRef@cdp:st:E6\endcsname{\hyperlink{tag-01GV}{II.\ref*{cdp:st:E6}}}
\expandafter\def\csname SRef@cdp:st:E7\endcsname{\hyperlink{tag-01GW}{II.\ref*{cdp:st:E7}}}
\expandafter\def\csname SRef@cdp:st:E8\endcsname{\hyperlink{tag-01GX}{II.\ref*{cdp:st:E8}}}
\expandafter\def\csname SRef@cdp:st:E10\endcsname{\hyperlink{tag-01GH}{II.\ref*{cdp:st:E10}}}
\expandafter\def\csname SRef@cdp:st:E11\endcsname{\hyperlink{tag-01GI}{II.\ref*{cdp:st:E11}}}
\expandafter\def\csname SRef@cdp:st:E12\endcsname{\hyperlink{tag-01GJ}{II.\ref*{cdp:st:E12}}}
\expandafter\def\csname SRef@cdp:st:E13\endcsname{\hyperlink{tag-01GK}{II.\ref*{cdp:st:E13}}}
\expandafter\def\csname SRef@cdp:st:E14\endcsname{\hyperlink{tag-01GL}{II.\ref*{cdp:st:E14}}}
\expandafter\def\csname SRef@push:src:E.2.13\endcsname{\hyperlink{tag-032Y}{II.\ref*{push:src:E.2.13}}}
\expandafter\def\csname SRef@push:src:E.2.14\endcsname{\hyperlink{tag-032Z}{II.\ref*{push:src:E.2.14}}}
\expandafter\def\csname SRef@push:src:E.2.15\endcsname{\hyperlink{tag-0330}{II.\ref*{push:src:E.2.15}}}
\expandafter\def\csname SRef@push:src:E.2.16\endcsname{\hyperlink{tag-0331}{II.\ref*{push:src:E.2.16}}}
\expandafter\def\csname SRef@push:src:E.2.17\endcsname{\hyperlink{tag-0332}{II.\ref*{push:src:E.2.17}}}
\expandafter\def\csname SRef@push:src:E.2.18\endcsname{\hyperlink{tag-0333}{II.\ref*{push:src:E.2.18}}}
\expandafter\def\csname SRef@push:src:E.2.19\endcsname{\hyperlink{tag-0334}{II.\ref*{push:src:E.2.19}}}
\expandafter\def\csname SRef@push:src:E.2.20\endcsname{\hyperlink{tag-0336}{II.\ref*{push:src:E.2.20}}}
\expandafter\def\csname SRef@push:src:E.2.21\endcsname{\hyperlink{tag-0337}{II.\ref*{push:src:E.2.21}}}
\expandafter\def\csname SRef@push:src:E.2.22\endcsname{\hyperlink{tag-0338}{II.\ref*{push:src:E.2.22}}}
\expandafter\def\csname SRef@push:src:E.2.5\endcsname{\hyperlink{tag-033D}{II.\ref*{push:src:E.2.5}}}
\expandafter\def\csname SRef@push:src:E.5.1\endcsname{\hyperlink{tag-033N}{II.\ref*{push:src:E.5.1}}}
\expandafter\def\csname SRef@push:src:E.5.2\endcsname{\hyperlink{tag-033O}{II.\ref*{push:src:E.5.2}}}
\expandafter\def\csname SRef@push:src:E.2.6\endcsname{\hyperlink{tag-033E}{II.\ref*{push:src:E.2.6}}}
\expandafter\def\csname SRef@push:src:E.2.7\endcsname{\hyperlink{tag-033F}{II.\ref*{push:src:E.2.7}}}
\expandafter\def\csname SRef@push:src:E.2.8\endcsname{\hyperlink{tag-033G}{II.\ref*{push:src:E.2.8}}}
\expandafter\def\csname SRef@push:src:E.2.9\endcsname{\hyperlink{tag-033H}{II.\ref*{push:src:E.2.9}}}
\expandafter\def\csname SRef@push:src:E.2.10\endcsname{\hyperlink{tag-032V}{II.\ref*{push:src:E.2.10}}}
\expandafter\def\csname SRef@push:src:E.2.11\endcsname{\hyperlink{tag-032W}{II.\ref*{push:src:E.2.11}}}
\expandafter\def\csname SRef@push:src:E.2.12\endcsname{\hyperlink{tag-032X}{II.\ref*{push:src:E.2.12}}}
\expandafter\def\csname SRef@inv:E2.2\endcsname{\hyperlink{tag-01SV}{II.\ref*{inv:E2.2}}}
\expandafter\def\csname SRef@inv:E2.4\endcsname{\hyperlink{tag-01SX}{II.\ref*{inv:E2.4}}}
\expandafter\def\csname SRef@wbII:prop:pic0-pullback\endcsname{\hyperlink{tag-03IO}{II.\ref*{wbII:prop:pic0-pullback}}}
\expandafter\def\csname SRef@wbII:prop:all-pic-restriction\endcsname{\hyperlink{tag-03IG}{II.\ref*{wbII:prop:all-pic-restriction}}}
\expandafter\def\csname SRef@wbII:prop:leray-edge-restriction\endcsname{\hyperlink{tag-03IM}{II.\ref*{wbII:prop:leray-edge-restriction}}}
\expandafter\def\csname SRef@wbII:prop:direct-boundary-picard\endcsname{\hyperlink{tag-03IL}{II.\ref*{wbII:prop:direct-boundary-picard}}}
\expandafter\def\csname SRef@wbII:prop:local-cusp-picard-obstruction\endcsname{\hyperlink{tag-03IN}{II.\ref*{wbII:prop:local-cusp-picard-obstruction}}}
\expandafter\def\csname SRef@wbII:prop:analytic-ah-boundary\endcsname{\hyperlink{tag-03IH}{II.\ref*{wbII:prop:analytic-ah-boundary}}}
\expandafter\def\csname SRef@wbII:prop:relative-local\endcsname{\hyperlink{tag-03IP}{II.\ref*{wbII:prop:relative-local}}}
\expandafter\def\csname SRef@wbII:prop:cartier-duality-absolute\endcsname{\hyperlink{tag-03IJ}{II.\ref*{wbII:prop:cartier-duality-absolute}}}
\expandafter\def\csname SRef@wbII:prop:analytic-base-change\endcsname{\hyperlink{tag-03II}{II.\ref*{wbII:prop:analytic-base-change}}}
\expandafter\def\csname SRef@cdp:st:C.1.1\endcsname{\hyperlink{tag-01BN}{II.\ref*{cdp:st:C.1.1}}}
\expandafter\def\csname SRef@cdp:st:C.1.3\endcsname{\hyperlink{tag-01BP}{II.\ref*{cdp:st:C.1.3}}}
\expandafter\def\csname SRef@cdp:sec:E\endcsname{\hyperlink{tag-01BP}{II.\ref*{cdp:sec:E}}}
\expandafter\def\csname SRef@cdp:st:E1\endcsname{\hyperlink{tag-01GG}{II.\ref*{cdp:st:E1}}}
\expandafter\def\csname SRef@cdp:st:D.1.2\endcsname{\hyperlink{tag-01G0}{II.\ref*{cdp:st:D.1.2}}}
\expandafter\def\csname SRef@cdp:st:D.1.4\endcsname{\hyperlink{tag-01G2}{II.\ref*{cdp:st:D.1.4}}}
\expandafter\def\csname SRef@cdp:st:D.1.13\endcsname{\hyperlink{tag-01FY}{II.\ref*{cdp:st:D.1.13}}}
\expandafter\def\csname SRef@cdp:st:D.1.15\endcsname{\hyperlink{tag-01FZ}{II.\ref*{cdp:st:D.1.15}}}
\expandafter\def\csname SRef@cdp:st:D.1.5\endcsname{\hyperlink{tag-01G3}{II.\ref*{cdp:st:D.1.5}}}
\expandafter\def\csname SRef@cdp:st:D.1.6\endcsname{\hyperlink{tag-01G4}{II.\ref*{cdp:st:D.1.6}}}
\expandafter\def\csname SRef@cdp:st:D.1.7\endcsname{\hyperlink{tag-01G5}{II.\ref*{cdp:st:D.1.7}}}
\expandafter\def\csname SRef@cdp:st:D.1.8\endcsname{\hyperlink{tag-01G6}{II.\ref*{cdp:st:D.1.8}}}
\expandafter\def\csname SRef@cdp:st:D.1.9\endcsname{\hyperlink{tag-01G7}{II.\ref*{cdp:st:D.1.9}}}
\expandafter\def\csname SRef@cdp:st:D.1.10\endcsname{\hyperlink{tag-01FV}{II.\ref*{cdp:st:D.1.10}}}
\expandafter\def\csname SRef@cdp:st:D.1.11\endcsname{\hyperlink{tag-01FW}{II.\ref*{cdp:st:D.1.11}}}
\expandafter\def\csname SRef@cdp:st:D.1.12\endcsname{\hyperlink{tag-01FX}{II.\ref*{cdp:st:D.1.12}}}
\expandafter\def\csname SRef@cdp:st:C.3.2\endcsname{\hyperlink{tag-01F8}{II.\ref*{cdp:st:C.3.2}}}
\expandafter\def\csname SRef@cdp:st:C.3.11\endcsname{\hyperlink{tag-01F6}{II.\ref*{cdp:st:C.3.11}}}
\expandafter\def\csname SRef@cdp:st:C.3.5\endcsname{\hyperlink{tag-01FB}{II.\ref*{cdp:st:C.3.5}}}
\expandafter\def\csname SRef@cdp:st:C.3.6\endcsname{\hyperlink{tag-01FC}{II.\ref*{cdp:st:C.3.6}}}
\expandafter\def\csname SRef@cdp:st:C.3.3\endcsname{\hyperlink{tag-01F9}{II.\ref*{cdp:st:C.3.3}}}
\expandafter\def\csname SRef@cdp:st:C.3.4\endcsname{\hyperlink{tag-01FA}{II.\ref*{cdp:st:C.3.4}}}
\expandafter\def\csname SRef@cdp:st:C.3.7\endcsname{\hyperlink{tag-01FD}{II.\ref*{cdp:st:C.3.7}}}
\expandafter\def\csname SRef@cdp:st:C.3.8\endcsname{\hyperlink{tag-01FE}{II.\ref*{cdp:st:C.3.8}}}
\expandafter\def\csname SRef@cdp:st:C.3.9\endcsname{\hyperlink{tag-01FF}{II.\ref*{cdp:st:C.3.9}}}
\expandafter\def\csname SRef@cdp:st:C.3.10\endcsname{\hyperlink{tag-01F5}{II.\ref*{cdp:st:C.3.10}}}
\expandafter\def\csname SRef@cdp:st:C.3.12\endcsname{\hyperlink{tag-01F7}{II.\ref*{cdp:st:C.3.12}}}
\expandafter\def\csname SRef@cdp:st:C.3.1\endcsname{\hyperlink{tag-01F4}{II.\ref*{cdp:st:C.3.1}}}
\expandafter\def\csname SRef@cdp:st:A.1.1\endcsname{\hyperlink{tag-01AD}{II.\ref*{cdp:st:A.1.1}}}
\expandafter\def\csname SRef@wbII:lem:conormal\endcsname{\hyperlink{tag-01AD}{II.\ref*{wbII:lem:conormal}}}
\expandafter\def\csname SRef@wbII:prop:criterion\endcsname{\hyperlink{tag-01AD}{II.\ref*{wbII:prop:criterion}}}
\expandafter\def\csname SRef@wbII:prop:filtration\endcsname{\hyperlink{tag-01AD}{II.\ref*{wbII:prop:filtration}}}
\expandafter\def\csname SRef@wbII:lem:conormal-injective\endcsname{\hyperlink{tag-01AD}{II.\ref*{wbII:lem:conormal-injective}}}
\expandafter\def\csname SRef@wbII:prop:missing-cotangent-filtration\endcsname{\hyperlink{tag-01AD}{II.\ref*{wbII:prop:missing-cotangent-filtration}}}
\expandafter\def\csname SRef@wbII:prop:missing-global-section-criterion\endcsname{\hyperlink{tag-01AD}{II.\ref*{wbII:prop:missing-global-section-criterion}}}
\expandafter\def\csname SRef@cdp:st:B3\endcsname{\hyperlink{tag-01BM}{II.\ref*{cdp:st:B3}}}
\expandafter\def\csname SRef@cdp:st:A.3.1\endcsname{\hyperlink{tag-01AL}{II.\ref*{cdp:st:A.3.1}}}
\expandafter\def\csname SRef@cdp:st:A.3.2\endcsname{\hyperlink{tag-01AM}{II.\ref*{cdp:st:A.3.2}}}
\expandafter\def\csname SRef@cdp:st:A.3.3\endcsname{\hyperlink{tag-01AN}{II.\ref*{cdp:st:A.3.3}}}
\expandafter\def\csname SRef@cdp:st:A.3.4\endcsname{\hyperlink{tag-01AO}{II.\ref*{cdp:st:A.3.4}}}
\expandafter\def\csname SRef@cdp:st:A.3.5\endcsname{\hyperlink{tag-01AP}{II.\ref*{cdp:st:A.3.5}}}
\expandafter\def\csname SRef@cdp:st:A.4.1\endcsname{\hyperlink{tag-01AQ}{II.\ref*{cdp:st:A.4.1}}}
\expandafter\def\csname SRef@cdp:st:A.4.2\endcsname{\hyperlink{tag-01AR}{II.\ref*{cdp:st:A.4.2}}}
\expandafter\def\csname SRef@cdp:st:A.4.3\endcsname{\hyperlink{tag-01AS}{II.\ref*{cdp:st:A.4.3}}}
\expandafter\def\csname SRef@cdp:st:A.4.4\endcsname{\hyperlink{tag-01AT}{II.\ref*{cdp:st:A.4.4}}}
\expandafter\def\csname SRef@cdp:st:A.2.1\endcsname{\hyperlink{tag-01AF}{II.\ref*{cdp:st:A.2.1}}}
\expandafter\def\csname SRef@cdp:st:A.4.5\endcsname{\hyperlink{tag-01AU}{II.\ref*{cdp:st:A.4.5}}}
\expandafter\def\csname SRef@cdp:st:A.4.6\endcsname{\hyperlink{tag-01AV}{II.\ref*{cdp:st:A.4.6}}}
\expandafter\def\csname SRef@cdp:st:A.4.7\endcsname{\hyperlink{tag-01AW}{II.\ref*{cdp:st:A.4.7}}}
\expandafter\def\csname SRef@cdp:st:A.4.8\endcsname{\hyperlink{tag-01AX}{II.\ref*{cdp:st:A.4.8}}}
\expandafter\def\csname SRef@cdp:st:A.4.9\endcsname{\hyperlink{tag-01AY}{II.\ref*{cdp:st:A.4.9}}}
\expandafter\def\csname SRef@cdp:st:A.2.2\endcsname{\hyperlink{tag-01AG}{II.\ref*{cdp:st:A.2.2}}}
\expandafter\def\csname SRef@wbII:lem:conductor\endcsname{\hyperlink{tag-01AG}{II.\ref*{wbII:lem:conductor}}}
\expandafter\def\csname SRef@wbII:lem:conductor-square\endcsname{\hyperlink{tag-01AG}{II.\ref*{wbII:lem:conductor-square}}}
\expandafter\def\csname SRef@cdp:st:A.1.2\endcsname{\hyperlink{tag-01AE}{II.\ref*{cdp:st:A.1.2}}}
\expandafter\def\csname SRef@cdp:st:A.2.4\endcsname{\hyperlink{tag-01AI}{II.\ref*{cdp:st:A.2.4}}}
\expandafter\def\csname SRef@cdp:st:A.2.5\endcsname{\hyperlink{tag-01AJ}{II.\ref*{cdp:st:A.2.5}}}
\expandafter\def\csname SRef@cdp:st:A.2.6\endcsname{\hyperlink{tag-01AK}{II.\ref*{cdp:st:A.2.6}}}
\expandafter\def\csname SRef@cdp:st:B1\endcsname{\hyperlink{tag-01BK}{II.\ref*{cdp:st:B1}}}
\expandafter\def\csname SRef@wbII:thm:nodal-test\endcsname{\hyperlink{tag-01BK}{II.\ref*{wbII:thm:nodal-test}}}
\expandafter\def\csname SRef@wbII:prop:family-geometry\endcsname{\hyperlink{tag-01BK}{II.\ref*{wbII:prop:family-geometry}}}
\expandafter\def\csname SRef@wbII:prop:nodal-counterexample\endcsname{\hyperlink{tag-01BK}{II.\ref*{wbII:prop:nodal-counterexample}}}
\expandafter\def\csname SRef@cdp:st:B2\endcsname{\hyperlink{tag-01BL}{II.\ref*{cdp:st:B2}}}
\expandafter\def\csname SRef@cdp:st:A.5.1\endcsname{\hyperlink{tag-01AZ}{II.\ref*{cdp:st:A.5.1}}}
\expandafter\def\csname SRef@cdp:st:A.5.5\endcsname{\hyperlink{tag-01BF}{II.\ref*{cdp:st:A.5.5}}}
\expandafter\def\csname SRef@cdp:st:A.5.7\endcsname{\hyperlink{tag-01BH}{II.\ref*{cdp:st:A.5.7}}}
\expandafter\def\csname SRef@cdp:st:A.5.8\endcsname{\hyperlink{tag-01BI}{II.\ref*{cdp:st:A.5.8}}}
\expandafter\def\csname SRef@cdp:st:A.5.14\endcsname{\hyperlink{tag-01B4}{II.\ref*{cdp:st:A.5.14}}}
\expandafter\def\csname SRef@cdp:st:A.5.10\endcsname{\hyperlink{tag-01B0}{II.\ref*{cdp:st:A.5.10}}}
\expandafter\def\csname SRef@cdp:st:A.5.20\endcsname{\hyperlink{tag-01BB}{II.\ref*{cdp:st:A.5.20}}}
\expandafter\def\csname SRef@cdp:st:A.5.15\endcsname{\hyperlink{tag-01B5}{II.\ref*{cdp:st:A.5.15}}}
\expandafter\def\csname SRef@cdp:st:A.5.17\endcsname{\hyperlink{tag-01B7}{II.\ref*{cdp:st:A.5.17}}}
\expandafter\def\csname SRef@cdp:st:A.5.18\endcsname{\hyperlink{tag-01B8}{II.\ref*{cdp:st:A.5.18}}}
\expandafter\def\csname SRef@cdp:st:A.5.19\endcsname{\hyperlink{tag-01B9}{II.\ref*{cdp:st:A.5.19}}}
\expandafter\def\csname SRef@cdp:st:A.5.2\endcsname{\hyperlink{tag-01BA}{II.\ref*{cdp:st:A.5.2}}}
\expandafter\def\csname SRef@cdp:st:A.5.16\endcsname{\hyperlink{tag-01B6}{II.\ref*{cdp:st:A.5.16}}}
\expandafter\def\csname SRef@cdp:st:A.5.21\endcsname{\hyperlink{tag-01BC}{II.\ref*{cdp:st:A.5.21}}}
\expandafter\def\csname SRef@cdp:st:A.5.3\endcsname{\hyperlink{tag-01BD}{II.\ref*{cdp:st:A.5.3}}}
\expandafter\def\csname SRef@cdp:st:A.5.4\endcsname{\hyperlink{tag-01BE}{II.\ref*{cdp:st:A.5.4}}}
\expandafter\def\csname SRef@cdp:st:A.5.6\endcsname{\hyperlink{tag-01BG}{II.\ref*{cdp:st:A.5.6}}}
\expandafter\def\csname SRef@cdp:st:A.5.9\endcsname{\hyperlink{tag-01BJ}{II.\ref*{cdp:st:A.5.9}}}
\expandafter\def\csname SRef@cdp:st:A.5.11\endcsname{\hyperlink{tag-01B1}{II.\ref*{cdp:st:A.5.11}}}
\expandafter\def\csname SRef@cdp:st:A.5.12\endcsname{\hyperlink{tag-01B2}{II.\ref*{cdp:st:A.5.12}}}
\expandafter\def\csname SRef@cdp:st:A.5.13\endcsname{\hyperlink{tag-01B3}{II.\ref*{cdp:st:A.5.13}}}
\expandafter\def\csname SRef@cdp:st:C.2.NC.2\endcsname{\hyperlink{tag-01C9}{II.\ref*{cdp:st:C.2.NC.2}}}
\expandafter\def\csname SRef@cdp:st:C.2.NC.3\endcsname{\hyperlink{tag-01CA}{II.\ref*{cdp:st:C.2.NC.3}}}
\expandafter\def\csname SRef@cdp:st:C.2.NC.4\endcsname{\hyperlink{tag-01CB}{II.\ref*{cdp:st:C.2.NC.4}}}
\expandafter\def\csname SRef@cdp:st:C.2.NC.5\endcsname{\hyperlink{tag-01CC}{II.\ref*{cdp:st:C.2.NC.5}}}
\expandafter\def\csname SRef@cdp:st:C2d\endcsname{\hyperlink{tag-01FJ}{II.\ref*{cdp:st:C2d}}}
\expandafter\def\csname SRef@night:P4.1\endcsname{\hyperlink{tag-02WG}{II.\ref*{night:P4.1}}}
\expandafter\def\csname SRef@stmt:P4.1\endcsname{\hyperlink{tag-02WG}{II.\ref*{stmt:P4.1}}}
\expandafter\def\csname SRef@night:P4.10\endcsname{\hyperlink{tag-02WH}{II.\ref*{night:P4.10}}}
\expandafter\def\csname SRef@stmt:P4.10\endcsname{\hyperlink{tag-02WH}{II.\ref*{stmt:P4.10}}}
\expandafter\def\csname SRef@night:P4.11\endcsname{\hyperlink{tag-02WI}{II.\ref*{night:P4.11}}}
\expandafter\def\csname SRef@stmt:P4.11\endcsname{\hyperlink{tag-02WI}{II.\ref*{stmt:P4.11}}}
\expandafter\def\csname SRef@night:P4.12\endcsname{\hyperlink{tag-02WJ}{II.\ref*{night:P4.12}}}
\expandafter\def\csname SRef@stmt:P4.12\endcsname{\hyperlink{tag-02WJ}{II.\ref*{stmt:P4.12}}}
\expandafter\def\csname SRef@night:P4.13\endcsname{\hyperlink{tag-02WK}{II.\ref*{night:P4.13}}}
\expandafter\def\csname SRef@stmt:P4.13\endcsname{\hyperlink{tag-02WK}{II.\ref*{stmt:P4.13}}}
\expandafter\def\csname SRef@night:P4.14\endcsname{\hyperlink{tag-02WL}{II.\ref*{night:P4.14}}}
\expandafter\def\csname SRef@stmt:P4.14\endcsname{\hyperlink{tag-02WL}{II.\ref*{stmt:P4.14}}}
\expandafter\def\csname SRef@night:P4.15\endcsname{\hyperlink{tag-02WM}{II.\ref*{night:P4.15}}}
\expandafter\def\csname SRef@stmt:P4.15\endcsname{\hyperlink{tag-02WM}{II.\ref*{stmt:P4.15}}}
\expandafter\def\csname SRef@night:P4.16\endcsname{\hyperlink{tag-02WN}{II.\ref*{night:P4.16}}}
\expandafter\def\csname SRef@stmt:P4.16\endcsname{\hyperlink{tag-02WN}{II.\ref*{stmt:P4.16}}}
\expandafter\def\csname SRef@night:P4.17\endcsname{\hyperlink{tag-02WO}{II.\ref*{night:P4.17}}}
\expandafter\def\csname SRef@stmt:P4.17\endcsname{\hyperlink{tag-02WO}{II.\ref*{stmt:P4.17}}}
\expandafter\def\csname SRef@night:P4.18\endcsname{\hyperlink{tag-02WP}{II.\ref*{night:P4.18}}}
\expandafter\def\csname SRef@stmt:P4.18\endcsname{\hyperlink{tag-02WP}{II.\ref*{stmt:P4.18}}}
\expandafter\def\csname SRef@night:P4.19\endcsname{\hyperlink{tag-02WQ}{II.\ref*{night:P4.19}}}
\expandafter\def\csname SRef@stmt:P4.19\endcsname{\hyperlink{tag-02WQ}{II.\ref*{stmt:P4.19}}}
\expandafter\def\csname SRef@night:P4.2\endcsname{\hyperlink{tag-02WR}{II.\ref*{night:P4.2}}}
\expandafter\def\csname SRef@stmt:P4.2\endcsname{\hyperlink{tag-02WR}{II.\ref*{stmt:P4.2}}}
\expandafter\def\csname SRef@night:P4.20\endcsname{\hyperlink{tag-02WS}{II.\ref*{night:P4.20}}}
\expandafter\def\csname SRef@stmt:P4.20\endcsname{\hyperlink{tag-02WS}{II.\ref*{stmt:P4.20}}}
\expandafter\def\csname SRef@night:P4.21\endcsname{\hyperlink{tag-02WT}{II.\ref*{night:P4.21}}}
\expandafter\def\csname SRef@stmt:P4.21\endcsname{\hyperlink{tag-02WT}{II.\ref*{stmt:P4.21}}}
\expandafter\def\csname SRef@night:P4.22\endcsname{\hyperlink{tag-02WU}{II.\ref*{night:P4.22}}}
\expandafter\def\csname SRef@stmt:P4.22\endcsname{\hyperlink{tag-02WU}{II.\ref*{stmt:P4.22}}}
\expandafter\def\csname SRef@night:P4.23\endcsname{\hyperlink{tag-02WV}{II.\ref*{night:P4.23}}}
\expandafter\def\csname SRef@stmt:P4.23\endcsname{\hyperlink{tag-02WV}{II.\ref*{stmt:P4.23}}}
\expandafter\def\csname SRef@night:P4.24\endcsname{\hyperlink{tag-02WW}{II.\ref*{night:P4.24}}}
\expandafter\def\csname SRef@stmt:P4.24\endcsname{\hyperlink{tag-02WW}{II.\ref*{stmt:P4.24}}}
\expandafter\def\csname SRef@night:P4.25\endcsname{\hyperlink{tag-02WX}{II.\ref*{night:P4.25}}}
\expandafter\def\csname SRef@stmt:P4.25\endcsname{\hyperlink{tag-02WX}{II.\ref*{stmt:P4.25}}}
\expandafter\def\csname SRef@night:P4.26\endcsname{\hyperlink{tag-02WY}{II.\ref*{night:P4.26}}}
\expandafter\def\csname SRef@stmt:P4.26\endcsname{\hyperlink{tag-02WY}{II.\ref*{stmt:P4.26}}}
\expandafter\def\csname SRef@night:P4.27\endcsname{\hyperlink{tag-02WZ}{II.\ref*{night:P4.27}}}
\expandafter\def\csname SRef@stmt:P4.27\endcsname{\hyperlink{tag-02WZ}{II.\ref*{stmt:P4.27}}}
\expandafter\def\csname SRef@night:P4.28\endcsname{\hyperlink{tag-02X0}{II.\ref*{night:P4.28}}}
\expandafter\def\csname SRef@stmt:P4.28\endcsname{\hyperlink{tag-02X0}{II.\ref*{stmt:P4.28}}}
\expandafter\def\csname SRef@night:P4.29\endcsname{\hyperlink{tag-02X1}{II.\ref*{night:P4.29}}}
\expandafter\def\csname SRef@stmt:P4.29\endcsname{\hyperlink{tag-02X1}{II.\ref*{stmt:P4.29}}}
\expandafter\def\csname SRef@night:P4.3\endcsname{\hyperlink{tag-02X2}{II.\ref*{night:P4.3}}}
\expandafter\def\csname SRef@stmt:P4.3\endcsname{\hyperlink{tag-02X2}{II.\ref*{stmt:P4.3}}}
\expandafter\def\csname SRef@night:P4.30\endcsname{\hyperlink{tag-02X3}{II.\ref*{night:P4.30}}}
\expandafter\def\csname SRef@stmt:P4.30\endcsname{\hyperlink{tag-02X3}{II.\ref*{stmt:P4.30}}}
\expandafter\def\csname SRef@night:P4.31\endcsname{\hyperlink{tag-02X4}{II.\ref*{night:P4.31}}}
\expandafter\def\csname SRef@stmt:P4.31\endcsname{\hyperlink{tag-02X4}{II.\ref*{stmt:P4.31}}}
\expandafter\def\csname SRef@night:P4.32\endcsname{\hyperlink{tag-02X5}{II.\ref*{night:P4.32}}}
\expandafter\def\csname SRef@stmt:P4.32\endcsname{\hyperlink{tag-02X5}{II.\ref*{stmt:P4.32}}}
\expandafter\def\csname SRef@night:P4.33\endcsname{\hyperlink{tag-02X6}{II.\ref*{night:P4.33}}}
\expandafter\def\csname SRef@stmt:P4.33\endcsname{\hyperlink{tag-02X6}{II.\ref*{stmt:P4.33}}}
\expandafter\def\csname SRef@night:P4.34\endcsname{\hyperlink{tag-02X7}{II.\ref*{night:P4.34}}}
\expandafter\def\csname SRef@stmt:P4.34\endcsname{\hyperlink{tag-02X7}{II.\ref*{stmt:P4.34}}}
\expandafter\def\csname SRef@night:P4.35\endcsname{\hyperlink{tag-02X8}{II.\ref*{night:P4.35}}}
\expandafter\def\csname SRef@stmt:P4.35\endcsname{\hyperlink{tag-02X8}{II.\ref*{stmt:P4.35}}}
\expandafter\def\csname SRef@night:P4.36\endcsname{\hyperlink{tag-02X9}{II.\ref*{night:P4.36}}}
\expandafter\def\csname SRef@stmt:P4.36\endcsname{\hyperlink{tag-02X9}{II.\ref*{stmt:P4.36}}}
\expandafter\def\csname SRef@night:P4.37\endcsname{\hyperlink{tag-02XA}{II.\ref*{night:P4.37}}}
\expandafter\def\csname SRef@stmt:P4.37\endcsname{\hyperlink{tag-02XA}{II.\ref*{stmt:P4.37}}}
\expandafter\def\csname SRef@night:P4.38\endcsname{\hyperlink{tag-02XB}{II.\ref*{night:P4.38}}}
\expandafter\def\csname SRef@stmt:P4.38\endcsname{\hyperlink{tag-02XB}{II.\ref*{stmt:P4.38}}}
\expandafter\def\csname SRef@night:P4.39\endcsname{\hyperlink{tag-02XC}{II.\ref*{night:P4.39}}}
\expandafter\def\csname SRef@stmt:P4.39\endcsname{\hyperlink{tag-02XC}{II.\ref*{stmt:P4.39}}}
\expandafter\def\csname SRef@night:P4.4\endcsname{\hyperlink{tag-02XD}{II.\ref*{night:P4.4}}}
\expandafter\def\csname SRef@stmt:P4.4\endcsname{\hyperlink{tag-02XD}{II.\ref*{stmt:P4.4}}}
\expandafter\def\csname SRef@night:P4.40\endcsname{\hyperlink{tag-02XE}{II.\ref*{night:P4.40}}}
\expandafter\def\csname SRef@stmt:P4.40\endcsname{\hyperlink{tag-02XE}{II.\ref*{stmt:P4.40}}}
\expandafter\def\csname SRef@night:P4.41\endcsname{\hyperlink{tag-02XF}{II.\ref*{night:P4.41}}}
\expandafter\def\csname SRef@stmt:P4.41\endcsname{\hyperlink{tag-02XF}{II.\ref*{stmt:P4.41}}}
\expandafter\def\csname SRef@night:P4.42\endcsname{\hyperlink{tag-02XG}{II.\ref*{night:P4.42}}}
\expandafter\def\csname SRef@stmt:P4.42\endcsname{\hyperlink{tag-02XG}{II.\ref*{stmt:P4.42}}}
\expandafter\def\csname SRef@night:P4.43\endcsname{\hyperlink{tag-02XH}{II.\ref*{night:P4.43}}}
\expandafter\def\csname SRef@stmt:P4.43\endcsname{\hyperlink{tag-02XH}{II.\ref*{stmt:P4.43}}}
\expandafter\def\csname SRef@night:P4.44\endcsname{\hyperlink{tag-02XI}{II.\ref*{night:P4.44}}}
\expandafter\def\csname SRef@stmt:P4.44\endcsname{\hyperlink{tag-02XI}{II.\ref*{stmt:P4.44}}}
\expandafter\def\csname SRef@night:P4.45\endcsname{\hyperlink{tag-02XJ}{II.\ref*{night:P4.45}}}
\expandafter\def\csname SRef@stmt:P4.45\endcsname{\hyperlink{tag-02XJ}{II.\ref*{stmt:P4.45}}}
\expandafter\def\csname SRef@night:P4.46\endcsname{\hyperlink{tag-02XK}{II.\ref*{night:P4.46}}}
\expandafter\def\csname SRef@stmt:P4.46\endcsname{\hyperlink{tag-02XK}{II.\ref*{stmt:P4.46}}}
\expandafter\def\csname SRef@night:P4.47\endcsname{\hyperlink{tag-02XL}{II.\ref*{night:P4.47}}}
\expandafter\def\csname SRef@stmt:P4.47\endcsname{\hyperlink{tag-02XL}{II.\ref*{stmt:P4.47}}}
\expandafter\def\csname SRef@night:P4.48\endcsname{\hyperlink{tag-02XM}{II.\ref*{night:P4.48}}}
\expandafter\def\csname SRef@stmt:P4.48\endcsname{\hyperlink{tag-02XM}{II.\ref*{stmt:P4.48}}}
\expandafter\def\csname SRef@night:P4.49\endcsname{\hyperlink{tag-02XN}{II.\ref*{night:P4.49}}}
\expandafter\def\csname SRef@stmt:P4.49\endcsname{\hyperlink{tag-02XN}{II.\ref*{stmt:P4.49}}}
\expandafter\def\csname SRef@night:P4.5\endcsname{\hyperlink{tag-02XO}{II.\ref*{night:P4.5}}}
\expandafter\def\csname SRef@stmt:P4.5\endcsname{\hyperlink{tag-02XO}{II.\ref*{stmt:P4.5}}}
\expandafter\def\csname SRef@night:P4.50\endcsname{\hyperlink{tag-02XP}{II.\ref*{night:P4.50}}}
\expandafter\def\csname SRef@stmt:P4.50\endcsname{\hyperlink{tag-02XP}{II.\ref*{stmt:P4.50}}}
\expandafter\def\csname SRef@night:P4.51\endcsname{\hyperlink{tag-02XQ}{II.\ref*{night:P4.51}}}
\expandafter\def\csname SRef@stmt:P4.51\endcsname{\hyperlink{tag-02XQ}{II.\ref*{stmt:P4.51}}}
\expandafter\def\csname SRef@night:P4.52\endcsname{\hyperlink{tag-02XR}{II.\ref*{night:P4.52}}}
\expandafter\def\csname SRef@stmt:P4.52\endcsname{\hyperlink{tag-02XR}{II.\ref*{stmt:P4.52}}}
\expandafter\def\csname SRef@night:P4.53\endcsname{\hyperlink{tag-02XS}{II.\ref*{night:P4.53}}}
\expandafter\def\csname SRef@stmt:P4.53\endcsname{\hyperlink{tag-02XS}{II.\ref*{stmt:P4.53}}}
\expandafter\def\csname SRef@night:P4.54\endcsname{\hyperlink{tag-02XT}{II.\ref*{night:P4.54}}}
\expandafter\def\csname SRef@stmt:P4.54\endcsname{\hyperlink{tag-02XT}{II.\ref*{stmt:P4.54}}}
\expandafter\def\csname SRef@night:P4.55\endcsname{\hyperlink{tag-02XU}{II.\ref*{night:P4.55}}}
\expandafter\def\csname SRef@stmt:P4.55\endcsname{\hyperlink{tag-02XU}{II.\ref*{stmt:P4.55}}}
\expandafter\def\csname SRef@night:P4.56\endcsname{\hyperlink{tag-02XV}{II.\ref*{night:P4.56}}}
\expandafter\def\csname SRef@stmt:P4.56\endcsname{\hyperlink{tag-02XV}{II.\ref*{stmt:P4.56}}}
\expandafter\def\csname SRef@night:P4.57\endcsname{\hyperlink{tag-02XW}{II.\ref*{night:P4.57}}}
\expandafter\def\csname SRef@stmt:P4.57\endcsname{\hyperlink{tag-02XW}{II.\ref*{stmt:P4.57}}}
\expandafter\def\csname SRef@night:P4.58\endcsname{\hyperlink{tag-02XX}{II.\ref*{night:P4.58}}}
\expandafter\def\csname SRef@stmt:P4.58\endcsname{\hyperlink{tag-02XX}{II.\ref*{stmt:P4.58}}}
\expandafter\def\csname SRef@night:P4R.1\endcsname{\hyperlink{tag-02Y2}{II.\ref*{night:P4R.1}}}
\expandafter\def\csname SRef@stmt:P4R.1\endcsname{\hyperlink{tag-02Y2}{II.\ref*{stmt:P4R.1}}}
\expandafter\def\csname SRef@night:P4.6\endcsname{\hyperlink{tag-02XY}{II.\ref*{night:P4.6}}}
\expandafter\def\csname SRef@stmt:P4.6\endcsname{\hyperlink{tag-02XY}{II.\ref*{stmt:P4.6}}}
\expandafter\def\csname SRef@night:P4R.2\endcsname{\hyperlink{tag-02YD}{II.\ref*{night:P4R.2}}}
\expandafter\def\csname SRef@stmt:P4R.2\endcsname{\hyperlink{tag-02YD}{II.\ref*{stmt:P4R.2}}}
\expandafter\def\csname SRef@night:P4R.3\endcsname{\hyperlink{tag-02YF}{II.\ref*{night:P4R.3}}}
\expandafter\def\csname SRef@stmt:P4R.3\endcsname{\hyperlink{tag-02YF}{II.\ref*{stmt:P4R.3}}}
\expandafter\def\csname SRef@night:P4R.4\endcsname{\hyperlink{tag-02YG}{II.\ref*{night:P4R.4}}}
\expandafter\def\csname SRef@stmt:P4R.4\endcsname{\hyperlink{tag-02YG}{II.\ref*{stmt:P4R.4}}}
\expandafter\def\csname SRef@night:P4R.5\endcsname{\hyperlink{tag-02YH}{II.\ref*{night:P4R.5}}}
\expandafter\def\csname SRef@stmt:P4R.5\endcsname{\hyperlink{tag-02YH}{II.\ref*{stmt:P4R.5}}}
\expandafter\def\csname SRef@night:P4R.6\endcsname{\hyperlink{tag-02YI}{II.\ref*{night:P4R.6}}}
\expandafter\def\csname SRef@stmt:P4R.6\endcsname{\hyperlink{tag-02YI}{II.\ref*{stmt:P4R.6}}}
\expandafter\def\csname SRef@night:P4R.7\endcsname{\hyperlink{tag-02YJ}{II.\ref*{night:P4R.7}}}
\expandafter\def\csname SRef@stmt:P4R.7\endcsname{\hyperlink{tag-02YJ}{II.\ref*{stmt:P4R.7}}}
\expandafter\def\csname SRef@night:P4R.8\endcsname{\hyperlink{tag-02YK}{II.\ref*{night:P4R.8}}}
\expandafter\def\csname SRef@stmt:P4R.8\endcsname{\hyperlink{tag-02YK}{II.\ref*{stmt:P4R.8}}}
\expandafter\def\csname SRef@night:P4R.9\endcsname{\hyperlink{tag-02YL}{II.\ref*{night:P4R.9}}}
\expandafter\def\csname SRef@stmt:P4R.9\endcsname{\hyperlink{tag-02YL}{II.\ref*{stmt:P4R.9}}}
\expandafter\def\csname SRef@night:P4R.10\endcsname{\hyperlink{tag-02Y3}{II.\ref*{night:P4R.10}}}
\expandafter\def\csname SRef@stmt:P4R.10\endcsname{\hyperlink{tag-02Y3}{II.\ref*{stmt:P4R.10}}}
\expandafter\def\csname SRef@night:P4R.11\endcsname{\hyperlink{tag-02Y4}{II.\ref*{night:P4R.11}}}
\expandafter\def\csname SRef@stmt:P4R.11\endcsname{\hyperlink{tag-02Y4}{II.\ref*{stmt:P4R.11}}}
\expandafter\def\csname SRef@night:P4.7\endcsname{\hyperlink{tag-02XZ}{II.\ref*{night:P4.7}}}
\expandafter\def\csname SRef@stmt:P4.7\endcsname{\hyperlink{tag-02XZ}{II.\ref*{stmt:P4.7}}}
\expandafter\def\csname SRef@night:P4R.12\endcsname{\hyperlink{tag-02Y5}{II.\ref*{night:P4R.12}}}
\expandafter\def\csname SRef@stmt:P4R.12\endcsname{\hyperlink{tag-02Y5}{II.\ref*{stmt:P4R.12}}}
\expandafter\def\csname SRef@night:P4R.13\endcsname{\hyperlink{tag-02Y6}{II.\ref*{night:P4R.13}}}
\expandafter\def\csname SRef@stmt:P4R.13\endcsname{\hyperlink{tag-02Y6}{II.\ref*{stmt:P4R.13}}}
\expandafter\def\csname SRef@night:P4R.14\endcsname{\hyperlink{tag-02Y7}{II.\ref*{night:P4R.14}}}
\expandafter\def\csname SRef@stmt:P4R.14\endcsname{\hyperlink{tag-02Y7}{II.\ref*{stmt:P4R.14}}}
\expandafter\def\csname SRef@night:P4R.15\endcsname{\hyperlink{tag-02Y8}{II.\ref*{night:P4R.15}}}
\expandafter\def\csname SRef@stmt:P4R.15\endcsname{\hyperlink{tag-02Y8}{II.\ref*{stmt:P4R.15}}}
\expandafter\def\csname SRef@night:P4R.16\endcsname{\hyperlink{tag-02Y9}{II.\ref*{night:P4R.16}}}
\expandafter\def\csname SRef@stmt:P4R.16\endcsname{\hyperlink{tag-02Y9}{II.\ref*{stmt:P4R.16}}}
\expandafter\def\csname SRef@night:P4R.17\endcsname{\hyperlink{tag-02YA}{II.\ref*{night:P4R.17}}}
\expandafter\def\csname SRef@stmt:P4R.17\endcsname{\hyperlink{tag-02YA}{II.\ref*{stmt:P4R.17}}}
\expandafter\def\csname SRef@bounds-without-an-ordinary-square-assumption\endcsname{\hyperlink{tag-02YA}{II.\ref*{bounds-without-an-ordinary-square-assumption}}}
\expandafter\def\csname SRef@night:P4R.18\endcsname{\hyperlink{tag-02YB}{II.\ref*{night:P4R.18}}}
\expandafter\def\csname SRef@stmt:P4R.18\endcsname{\hyperlink{tag-02YB}{II.\ref*{stmt:P4R.18}}}
\expandafter\def\csname SRef@night:P4R.19\endcsname{\hyperlink{tag-02YC}{II.\ref*{night:P4R.19}}}
\expandafter\def\csname SRef@stmt:P4R.19\endcsname{\hyperlink{tag-02YC}{II.\ref*{stmt:P4R.19}}}
\expandafter\def\csname SRef@night:P4R.20\endcsname{\hyperlink{tag-02YE}{II.\ref*{night:P4R.20}}}
\expandafter\def\csname SRef@stmt:P4R.20\endcsname{\hyperlink{tag-02YE}{II.\ref*{stmt:P4R.20}}}
\expandafter\def\csname SRef@night:P4.8\endcsname{\hyperlink{tag-02Y0}{II.\ref*{night:P4.8}}}
\expandafter\def\csname SRef@stmt:P4.8\endcsname{\hyperlink{tag-02Y0}{II.\ref*{stmt:P4.8}}}
\expandafter\def\csname SRef@night:P4.9\endcsname{\hyperlink{tag-02Y1}{II.\ref*{night:P4.9}}}
\expandafter\def\csname SRef@stmt:P4.9\endcsname{\hyperlink{tag-02Y1}{II.\ref*{stmt:P4.9}}}
\expandafter\def\csname SRef@cdp:st:C4a\endcsname{\hyperlink{tag-01FK}{II.\ref*{cdp:st:C4a}}}
\expandafter\def\csname SRef@cdp:st:C5c\endcsname{\hyperlink{tag-01FT}{II.\ref*{cdp:st:C5c}}}
\expandafter\def\csname SRef@cdp:st:C4b\endcsname{\hyperlink{tag-01FL}{II.\ref*{cdp:st:C4b}}}
\expandafter\def\csname SRef@cdp:st:C4c\endcsname{\hyperlink{tag-01FM}{II.\ref*{cdp:st:C4c}}}
\expandafter\def\csname SRef@cdp:st:C4d\endcsname{\hyperlink{tag-01FN}{II.\ref*{cdp:st:C4d}}}
\expandafter\def\csname SRef@cdp:st:C4e\endcsname{\hyperlink{tag-01FO}{II.\ref*{cdp:st:C4e}}}
\expandafter\def\csname SRef@cdp:st:C4f\endcsname{\hyperlink{tag-01FP}{II.\ref*{cdp:st:C4f}}}
\expandafter\def\csname SRef@cdp:st:C4g\endcsname{\hyperlink{tag-01FQ}{II.\ref*{cdp:st:C4g}}}
\expandafter\def\csname SRef@cdp:st:C5a\endcsname{\hyperlink{tag-01FR}{II.\ref*{cdp:st:C5a}}}
\expandafter\def\csname SRef@cdp:st:C5b\endcsname{\hyperlink{tag-01FS}{II.\ref*{cdp:st:C5b}}}
\expandafter\def\csname SRef@night:P3K.36\endcsname{\hyperlink{tag-02UH}{II.\ref*{night:P3K.36}}}
\expandafter\def\csname SRef@stmt:P3K.36\endcsname{\hyperlink{tag-02UH}{II.\ref*{stmt:P3K.36}}}
\expandafter\def\csname SRef@inv:wb:uniform-nontorsion-provenance\endcsname{\hyperlink{tag-01TQ}{II.\ref*{inv:wb:uniform-nontorsion-provenance}}}
\expandafter\def\csname SRef@wbII:rem:scope\endcsname{\hyperlink{tag-03IS}{II.\ref*{wbII:rem:scope}}}
\expandafter\def\csname SRef@wbII:rem:local-cusp-picard-obstruction\endcsname{\hyperlink{tag-03IQ}{II.\ref*{wbII:rem:local-cusp-picard-obstruction}}}
\expandafter\def\csname SRef@wbII:rem:relative-local-scope\endcsname{\hyperlink{tag-03IR}{II.\ref*{wbII:rem:relative-local-scope}}}
\expandafter\def\csname SRef@sec:II:conventions\endcsname{\hyperlink{tag-034E}{II.\ref*{sec:II:conventions}}}
\expandafter\def\csname SRef@sec:II:13\endcsname{\hyperlink{tag-0349}{II.\ref*{sec:II:13}}}
\expandafter\def\csname SRef@sec:II:14\endcsname{\hyperlink{tag-034A}{II.\ref*{sec:II:14}}}
\expandafter\def\csname SRef@sec:II:push-stability\endcsname{\hyperlink{tag-034G}{II.\ref*{sec:II:push-stability}}}
\expandafter\def\csname SRef@sec:II:combs\endcsname{\hyperlink{tag-034D}{II.\ref*{sec:II:combs}}}
\expandafter\def\csname SRef@sec:II:uniform-nontorsion\endcsname{\hyperlink{tag-034H}{II.\ref*{sec:II:uniform-nontorsion}}}
\expandafter\def\csname SRef@app:II:programs\endcsname{\hyperlink{tag-000B}{II.\ref*{app:II:programs}}}
\expandafter\def\csname SRef@sec:II:01\endcsname{\hyperlink{tag-033X}{II.\ref*{sec:II:01}}}
\expandafter\def\csname SRef@sec:II:02\endcsname{\hyperlink{tag-033Y}{II.\ref*{sec:II:02}}}
\expandafter\def\csname SRef@sec:II:03\endcsname{\hyperlink{tag-033Z}{II.\ref*{sec:II:03}}}
\expandafter\def\csname SRef@sec:II:04\endcsname{\hyperlink{tag-0340}{II.\ref*{sec:II:04}}}
\expandafter\def\csname SRef@sec:II:05\endcsname{\hyperlink{tag-0341}{II.\ref*{sec:II:05}}}
\expandafter\def\csname SRef@sec:II:08\endcsname{\hyperlink{tag-0344}{II.\ref*{sec:II:08}}}
\expandafter\def\csname SRef@sec:II:night-sign\endcsname{\hyperlink{tag-034F}{II.\ref*{sec:II:night-sign}}}
\expandafter\def\csname SRef@sec:II:12\endcsname{\hyperlink{tag-0348}{II.\ref*{sec:II:12}}}
\expandafter\def\csname SRef@wbII:subsec:local-cusp-picard-obstruction\endcsname{\hyperlink{tag-03IU}{II.\ref*{wbII:subsec:local-cusp-picard-obstruction}}}
\expandafter\def\csname SRef@wbII:subsec:analytic-all-picard\endcsname{\hyperlink{tag-03IT}{II.\ref*{wbII:subsec:analytic-all-picard}}}
\expandafter\def\csname SRef@push:src:E.2.23\endcsname{\hyperlink{tag-0339}{II.\ref*{push:src:E.2.23}}}
\expandafter\def\csname SRef@push:src:E.5.3\endcsname{\hyperlink{tag-033P}{II.\ref*{push:src:E.5.3}}}
\expandafter\def\csname SRef@wb:thm:uniform-nontorsion-audit\endcsname{\hyperlink{tag-03DQ}{II.\ref*{wb:thm:uniform-nontorsion-audit}}}
\expandafter\def\csname SRef@wbII:prop:conductor\endcsname{\hyperlink{tag-03IK}{II.\ref*{wbII:prop:conductor}}}
\expandafter\def\csname SRef@wbII:prop:conductor-differential\endcsname{\hyperlink{tag-03IK}{II.\ref*{wbII:prop:conductor-differential}}}
\expandafter\def\csname SRef@wbII:thm:analytic-all-picard\endcsname{\hyperlink{tag-03IW}{II.\ref*{wbII:thm:analytic-all-picard}}}
\expandafter\def\csname SRef@wbII:thm:candidate-conductor\endcsname{\hyperlink{tag-03IX}{II.\ref*{wbII:thm:candidate-conductor}}}
\expandafter\def\csname SRef@wbII:thm:corrected-cdp-global\endcsname{\hyperlink{tag-03J0}{II.\ref*{wbII:thm:corrected-cdp-global}}}
\expandafter\def\csname SRef@wbII:thm:corrected-global\endcsname{\hyperlink{tag-03J0}{II.\ref*{wbII:thm:corrected-global}}}
\expandafter\def\csname SRef@wbII:thm:corrected-cdp-fibrewise\endcsname{\hyperlink{tag-03IZ}{II.\ref*{wbII:thm:corrected-cdp-fibrewise}}}
\expandafter\def\csname SRef@wbII:thm:cartier-duality\endcsname{\hyperlink{tag-03IY}{II.\ref*{wbII:thm:cartier-duality}}}
\expandafter\def\csname SRef@ii:LR:covering\endcsname{\hyperlink{tag-01RO}{II.\ref*{ii:LR:covering}}}
\expandafter\def\csname SRef@maps:eq:D.O9\endcsname{\href{S6_SERIES_GRIND_III.pdf\#nameddest=tag-020U}{III.{D.O9} (p.~85)} }
\expandafter\def\csname SRef@maps:eq:D.FI1\endcsname{\href{S6_SERIES_GRIND_III.pdf\#nameddest=tag-01YT}{III.{D.FI1} (p.~87)} }
\expandafter\def\csname SRef@maps:eq:D.FI2\endcsname{\href{S6_SERIES_GRIND_III.pdf\#nameddest=tag-01YU}{III.{D.FI2} (p.~88)} }
\expandafter\def\csname SRef@maps:eq:D.FI3\endcsname{\href{S6_SERIES_GRIND_III.pdf\#nameddest=tag-01YV}{III.{D.FI3} (p.~88)} }
\expandafter\def\csname SRef@maps:eq:D.FI4\endcsname{\href{S6_SERIES_GRIND_III.pdf\#nameddest=tag-01YW}{III.{D.FI4} (p.~90)} }
\expandafter\def\csname SRef@maps:eq:D.FI5\endcsname{\href{S6_SERIES_GRIND_III.pdf\#nameddest=tag-01YX}{III.{D.FI5} (p.~91)} }
\expandafter\def\csname SRef@maps:eq:D.FG1\endcsname{\href{S6_SERIES_GRIND_III.pdf\#nameddest=tag-01YQ}{III.{D.FG1} (p.~92)} }
\expandafter\def\csname SRef@maps:eq:D.FG2\endcsname{\href{S6_SERIES_GRIND_III.pdf\#nameddest=tag-01YR}{III.{D.FG2} (p.~92)} }
\expandafter\def\csname SRef@maps:eq:D.FG3\endcsname{\href{S6_SERIES_GRIND_III.pdf\#nameddest=tag-01YS}{III.{D.FG3} (p.~92)} }
\expandafter\def\csname SRef@maps:eq:D.IS1\endcsname{\href{S6_SERIES_GRIND_III.pdf\#nameddest=tag-0207}{III.{D.IS1} (p.~93)} }
\expandafter\def\csname SRef@maps:eq:D.IS2\endcsname{\href{S6_SERIES_GRIND_III.pdf\#nameddest=tag-0208}{III.{D.IS2} (p.~93)} }
\expandafter\def\csname SRef@maps:eq:D.SC1\endcsname{\href{S6_SERIES_GRIND_III.pdf\#nameddest=tag-023Z}{III.{D.SC1} (p.~94)} }
\expandafter\def\csname SRef@maps:eq:D.SC2\endcsname{\href{S6_SERIES_GRIND_III.pdf\#nameddest=tag-0240}{III.{D.SC2} (p.~95)} }
\expandafter\def\csname SRef@maps:eq:D.SC3\endcsname{\href{S6_SERIES_GRIND_III.pdf\#nameddest=tag-0241}{III.{D.SC3} (p.~95)} }
\expandafter\def\csname SRef@maps:eq:D.SC4\endcsname{\href{S6_SERIES_GRIND_III.pdf\#nameddest=tag-0242}{III.{D.SC4} (p.~95)} }
\expandafter\def\csname SRef@maps:eq:D.SC5\endcsname{\href{S6_SERIES_GRIND_III.pdf\#nameddest=tag-0243}{III.{D.SC5} (p.~95)} }
\expandafter\def\csname SRef@maps:eq:D.SC6\endcsname{\href{S6_SERIES_GRIND_III.pdf\#nameddest=tag-0244}{III.{D.SC6} (p.~95)} }
\expandafter\def\csname SRef@maps:eq:D.SC7\endcsname{\href{S6_SERIES_GRIND_III.pdf\#nameddest=tag-0245}{III.{D.SC7} (p.~97)} }
\expandafter\def\csname SRef@maps:eq:D.SC8\endcsname{\href{S6_SERIES_GRIND_III.pdf\#nameddest=tag-0246}{III.{D.SC8} (p.~97)} }
\expandafter\def\csname SRef@maps:eq:D.HY1\endcsname{\href{S6_SERIES_GRIND_III.pdf\#nameddest=tag-01ZN}{III.{D.HY1} (p.~98)} }
\expandafter\def\csname SRef@maps:eq:D.HY2\endcsname{\href{S6_SERIES_GRIND_III.pdf\#nameddest=tag-01ZO}{III.{D.HY2} (p.~98)} }
\expandafter\def\csname SRef@maps:eq:D.HY3\endcsname{\href{S6_SERIES_GRIND_III.pdf\#nameddest=tag-01ZP}{III.{D.HY3} (p.~99)} }
\expandafter\def\csname SRef@maps:eq:D.HY4\endcsname{\href{S6_SERIES_GRIND_III.pdf\#nameddest=tag-01ZQ}{III.{D.HY4} (p.~99)} }
\expandafter\def\csname SRef@maps:eq:D.HY5\endcsname{\href{S6_SERIES_GRIND_III.pdf\#nameddest=tag-01ZR}{III.{D.HY5} (p.~99)} }
\expandafter\def\csname SRef@maps:eq:D.HY6\endcsname{\href{S6_SERIES_GRIND_III.pdf\#nameddest=tag-01ZS}{III.{D.HY6} (p.~100)} }
\expandafter\def\csname SRef@maps:eq:D.HY7\endcsname{\href{S6_SERIES_GRIND_III.pdf\#nameddest=tag-01ZT}{III.{D.HY7} (p.~100)} }
\expandafter\def\csname SRef@maps:eq:D.ST1\endcsname{\href{S6_SERIES_GRIND_III.pdf\#nameddest=tag-024V}{III.{D.ST1} (p.~102)} }
\expandafter\def\csname SRef@maps:eq:D.ST2\endcsname{\href{S6_SERIES_GRIND_III.pdf\#nameddest=tag-0250}{III.{D.ST2} (p.~103)} }
\expandafter\def\csname SRef@maps:eq:D.ST3\endcsname{\href{S6_SERIES_GRIND_III.pdf\#nameddest=tag-0251}{III.{D.ST3} (p.~103)} }
\expandafter\def\csname SRef@maps:eq:D.ST4\endcsname{\href{S6_SERIES_GRIND_III.pdf\#nameddest=tag-0252}{III.{D.ST4} (p.~103)} }
\expandafter\def\csname SRef@maps:eq:D.ST5\endcsname{\href{S6_SERIES_GRIND_III.pdf\#nameddest=tag-0253}{III.{D.ST5} (p.~103)} }
\expandafter\def\csname SRef@maps:eq:D.ST6\endcsname{\href{S6_SERIES_GRIND_III.pdf\#nameddest=tag-0254}{III.{D.ST6} (p.~103)} }
\expandafter\def\csname SRef@maps:eq:D.ST7\endcsname{\href{S6_SERIES_GRIND_III.pdf\#nameddest=tag-0255}{III.{D.ST7} (p.~104)} }
\expandafter\def\csname SRef@maps:eq:D.ST8\endcsname{\href{S6_SERIES_GRIND_III.pdf\#nameddest=tag-0256}{III.{D.ST8} (p.~104)} }
\expandafter\def\csname SRef@maps:eq:D.ST9\endcsname{\href{S6_SERIES_GRIND_III.pdf\#nameddest=tag-0257}{III.{D.ST9} (p.~104)} }
\expandafter\def\csname SRef@maps:eq:D.ST10\endcsname{\href{S6_SERIES_GRIND_III.pdf\#nameddest=tag-024W}{III.{D.ST10} (p.~105)} }
\expandafter\def\csname SRef@maps:eq:D.ST11\endcsname{\href{S6_SERIES_GRIND_III.pdf\#nameddest=tag-024X}{III.{D.ST11} (p.~105)} }
\expandafter\def\csname SRef@maps:eq:D.ST12\endcsname{\href{S6_SERIES_GRIND_III.pdf\#nameddest=tag-024Y}{III.{D.ST12} (p.~105)} }
\expandafter\def\csname SRef@maps:eq:D.ST13\endcsname{\href{S6_SERIES_GRIND_III.pdf\#nameddest=tag-024Z}{III.{D.ST13} (p.~106)} }
\expandafter\def\csname SRef@maps:eq:D.UF1\endcsname{\href{S6_SERIES_GRIND_III.pdf\#nameddest=tag-025Z}{III.{D.UF1} (p.~107)} }
\expandafter\def\csname SRef@maps:eq:D.UF2\endcsname{\href{S6_SERIES_GRIND_III.pdf\#nameddest=tag-0261}{III.{D.UF2} (p.~107)} }
\expandafter\def\csname SRef@maps:eq:D.UF3\endcsname{\href{S6_SERIES_GRIND_III.pdf\#nameddest=tag-0262}{III.{D.UF3} (p.~107)} }
\expandafter\def\csname SRef@maps:eq:D.UF4\endcsname{\href{S6_SERIES_GRIND_III.pdf\#nameddest=tag-0263}{III.{D.UF4} (p.~107)} }
\expandafter\def\csname SRef@maps:eq:D.UF5\endcsname{\href{S6_SERIES_GRIND_III.pdf\#nameddest=tag-0264}{III.{D.UF5} (p.~108)} }
\expandafter\def\csname SRef@maps:eq:D.UF6\endcsname{\href{S6_SERIES_GRIND_III.pdf\#nameddest=tag-0265}{III.{D.UF6} (p.~108)} }
\expandafter\def\csname SRef@maps:eq:D.UF7\endcsname{\href{S6_SERIES_GRIND_III.pdf\#nameddest=tag-0266}{III.{D.UF7} (p.~108)} }
\expandafter\def\csname SRef@maps:eq:D.UF8\endcsname{\href{S6_SERIES_GRIND_III.pdf\#nameddest=tag-0267}{III.{D.UF8} (p.~109)} }
\expandafter\def\csname SRef@maps:eq:D.UF9\endcsname{\href{S6_SERIES_GRIND_III.pdf\#nameddest=tag-0268}{III.{D.UF9} (p.~109)} }
\expandafter\def\csname SRef@maps:eq:D.UF10\endcsname{\href{S6_SERIES_GRIND_III.pdf\#nameddest=tag-0260}{III.{D.UF10} (p.~109)} }
\expandafter\def\csname SRef@maps:eq:D.AF1\endcsname{\href{S6_SERIES_GRIND_III.pdf\#nameddest=tag-01XI}{III.{D.AF1} (p.~110)} }
\expandafter\def\csname SRef@maps:eq:D.AF2\endcsname{\href{S6_SERIES_GRIND_III.pdf\#nameddest=tag-01XJ}{III.{D.AF2} (p.~110)} }
\expandafter\def\csname SRef@maps:eq:D.AF3\endcsname{\href{S6_SERIES_GRIND_III.pdf\#nameddest=tag-01XK}{III.{D.AF3} (p.~111)} }
\expandafter\def\csname SRef@maps:eq:D.AF4\endcsname{\href{S6_SERIES_GRIND_III.pdf\#nameddest=tag-01XL}{III.{D.AF4} (p.~111)} }
\expandafter\def\csname SRef@maps:eq:D.AF5\endcsname{\href{S6_SERIES_GRIND_III.pdf\#nameddest=tag-01XM}{III.{D.AF5} (p.~111)} }
\expandafter\def\csname SRef@maps:eq:D.AF6\endcsname{\href{S6_SERIES_GRIND_III.pdf\#nameddest=tag-01XN}{III.{D.AF6} (p.~111)} }
\expandafter\def\csname SRef@maps:eq:D.AF7\endcsname{\href{S6_SERIES_GRIND_III.pdf\#nameddest=tag-01XO}{III.{D.AF7} (p.~111)} }
\expandafter\def\csname SRef@maps:eq:D.AF8\endcsname{\href{S6_SERIES_GRIND_III.pdf\#nameddest=tag-01XP}{III.{D.AF8} (p.~111)} }
\expandafter\def\csname SRef@maps:eq:D.AF9\endcsname{\href{S6_SERIES_GRIND_III.pdf\#nameddest=tag-01XQ}{III.{D.AF9} (p.~112)} }
\expandafter\def\csname SRef@maps:eq:D.OD1\endcsname{\href{S6_SERIES_GRIND_III.pdf\#nameddest=tag-0212}{III.{D.OD1} (p.~113)} }
\expandafter\def\csname SRef@maps:eq:D.OD2\endcsname{\href{S6_SERIES_GRIND_III.pdf\#nameddest=tag-0213}{III.{D.OD2} (p.~113)} }
\expandafter\def\csname SRef@maps:eq:D.OD3\endcsname{\href{S6_SERIES_GRIND_III.pdf\#nameddest=tag-0214}{III.{D.OD3} (p.~113)} }
\expandafter\def\csname SRef@maps:eq:D.OD4\endcsname{\href{S6_SERIES_GRIND_III.pdf\#nameddest=tag-0215}{III.{D.OD4} (p.~114)} }
\expandafter\def\csname SRef@maps:eq:D.OD5\endcsname{\href{S6_SERIES_GRIND_III.pdf\#nameddest=tag-0216}{III.{D.OD5} (p.~114)} }
\expandafter\def\csname SRef@maps:eq:D.OC1\endcsname{\href{S6_SERIES_GRIND_III.pdf\#nameddest=tag-020V}{III.{D.OC1} (p.~116)} }
\expandafter\def\csname SRef@maps:eq:D.OC2\endcsname{\href{S6_SERIES_GRIND_III.pdf\#nameddest=tag-020W}{III.{D.OC2} (p.~116)} }
\expandafter\def\csname SRef@maps:eq:D.OC3\endcsname{\href{S6_SERIES_GRIND_III.pdf\#nameddest=tag-020X}{III.{D.OC3} (p.~116)} }
\expandafter\def\csname SRef@maps:eq:D.OC4\endcsname{\href{S6_SERIES_GRIND_III.pdf\#nameddest=tag-020Y}{III.{D.OC4} (p.~117)} }
\expandafter\def\csname SRef@maps:eq:D.OC5\endcsname{\href{S6_SERIES_GRIND_III.pdf\#nameddest=tag-020Z}{III.{D.OC5} (p.~117)} }
\expandafter\def\csname SRef@maps:eq:D.OC6\endcsname{\href{S6_SERIES_GRIND_III.pdf\#nameddest=tag-0210}{III.{D.OC6} (p.~117)} }
\expandafter\def\csname SRef@maps:eq:D.OC7\endcsname{\href{S6_SERIES_GRIND_III.pdf\#nameddest=tag-0211}{III.{D.OC7} (p.~118)} }
\expandafter\def\csname SRef@maps:eq:D.OM1\endcsname{\href{S6_SERIES_GRIND_III.pdf\#nameddest=tag-0222}{III.{D.OM1} (p.~119)} }
\expandafter\def\csname SRef@maps:eq:D.OM2\endcsname{\href{S6_SERIES_GRIND_III.pdf\#nameddest=tag-0224}{III.{D.OM2} (p.~119)} }
\expandafter\def\csname SRef@maps:eq:D.OM3\endcsname{\href{S6_SERIES_GRIND_III.pdf\#nameddest=tag-0225}{III.{D.OM3} (p.~120)} }
\expandafter\def\csname SRef@maps:eq:D.OM4\endcsname{\href{S6_SERIES_GRIND_III.pdf\#nameddest=tag-0226}{III.{D.OM4} (p.~120)} }
\expandafter\def\csname SRef@maps:eq:D.OM5\endcsname{\href{S6_SERIES_GRIND_III.pdf\#nameddest=tag-0227}{III.{D.OM5} (p.~120)} }
\expandafter\def\csname SRef@maps:eq:D.OM6\endcsname{\href{S6_SERIES_GRIND_III.pdf\#nameddest=tag-0228}{III.{D.OM6} (p.~120)} }
\expandafter\def\csname SRef@maps:eq:D.OM7\endcsname{\href{S6_SERIES_GRIND_III.pdf\#nameddest=tag-0229}{III.{D.OM7} (p.~120)} }
\expandafter\def\csname SRef@maps:eq:D.OM8\endcsname{\href{S6_SERIES_GRIND_III.pdf\#nameddest=tag-022A}{III.{D.OM8} (p.~120)} }
\expandafter\def\csname SRef@maps:eq:D.OM9\endcsname{\href{S6_SERIES_GRIND_III.pdf\#nameddest=tag-022B}{III.{D.OM9} (p.~121)} }
\expandafter\def\csname SRef@maps:eq:B.1.1\endcsname{\href{S6_SERIES_GRIND_III.pdf\#nameddest=tag-01UV}{III.{B.1.1} (p.~14)} }
\expandafter\def\csname SRef@maps:eq:D.OM10\endcsname{\href{S6_SERIES_GRIND_III.pdf\#nameddest=tag-0223}{III.{D.OM10} (p.~121)} }
\expandafter\def\csname SRef@maps:eq:D.WE1\endcsname{\href{S6_SERIES_GRIND_III.pdf\#nameddest=tag-026O}{III.{D.WE1} (p.~122)} }
\expandafter\def\csname SRef@maps:eq:D.WE2\endcsname{\href{S6_SERIES_GRIND_III.pdf\#nameddest=tag-026U}{III.{D.WE2} (p.~122)} }
\expandafter\def\csname SRef@maps:eq:D.WE3\endcsname{\href{S6_SERIES_GRIND_III.pdf\#nameddest=tag-026V}{III.{D.WE3} (p.~122)} }
\expandafter\def\csname SRef@maps:eq:D.WE4\endcsname{\href{S6_SERIES_GRIND_III.pdf\#nameddest=tag-026W}{III.{D.WE4} (p.~122)} }
\expandafter\def\csname SRef@maps:eq:D.WE5\endcsname{\href{S6_SERIES_GRIND_III.pdf\#nameddest=tag-026X}{III.{D.WE5} (p.~122)} }
\expandafter\def\csname SRef@maps:eq:D.WE6\endcsname{\href{S6_SERIES_GRIND_III.pdf\#nameddest=tag-026Y}{III.{D.WE6} (p.~122)} }
\expandafter\def\csname SRef@maps:eq:D.WE7\endcsname{\href{S6_SERIES_GRIND_III.pdf\#nameddest=tag-026Z}{III.{D.WE7} (p.~123)} }
\expandafter\def\csname SRef@maps:eq:D.WE8\endcsname{\href{S6_SERIES_GRIND_III.pdf\#nameddest=tag-0270}{III.{D.WE8} (p.~123)} }
\expandafter\def\csname SRef@maps:eq:D.WE9\endcsname{\href{S6_SERIES_GRIND_III.pdf\#nameddest=tag-0271}{III.{D.WE9} (p.~123)} }
\expandafter\def\csname SRef@maps:eq:B.1.2\endcsname{\href{S6_SERIES_GRIND_III.pdf\#nameddest=tag-01UW}{III.{B.1.2} (p.~14)} }
\expandafter\def\csname SRef@maps:eq:D.WE10\endcsname{\href{S6_SERIES_GRIND_III.pdf\#nameddest=tag-026P}{III.{D.WE10} (p.~123)} }
\expandafter\def\csname SRef@maps:eq:D.WE11\endcsname{\href{S6_SERIES_GRIND_III.pdf\#nameddest=tag-026Q}{III.{D.WE11} (p.~124)} }
\expandafter\def\csname SRef@maps:eq:D.WE12\endcsname{\href{S6_SERIES_GRIND_III.pdf\#nameddest=tag-026R}{III.{D.WE12} (p.~125)} }
\expandafter\def\csname SRef@maps:eq:D.WE13\endcsname{\href{S6_SERIES_GRIND_III.pdf\#nameddest=tag-026S}{III.{D.WE13} (p.~125)} }
\expandafter\def\csname SRef@maps:eq:D.WE14\endcsname{\href{S6_SERIES_GRIND_III.pdf\#nameddest=tag-026T}{III.{D.WE14} (p.~126)} }
\expandafter\def\csname SRef@maps:eq:D.C21\endcsname{\href{S6_SERIES_GRIND_III.pdf\#nameddest=tag-01Y7}{III.{D.C21} (p.~126)} }
\expandafter\def\csname SRef@maps:eq:D.C22\endcsname{\href{S6_SERIES_GRIND_III.pdf\#nameddest=tag-01YF}{III.{D.C22} (p.~126)} }
\expandafter\def\csname SRef@maps:eq:D.C23\endcsname{\href{S6_SERIES_GRIND_III.pdf\#nameddest=tag-01YG}{III.{D.C23} (p.~127)} }
\expandafter\def\csname SRef@maps:eq:D.C24\endcsname{\href{S6_SERIES_GRIND_III.pdf\#nameddest=tag-01YH}{III.{D.C24} (p.~127)} }
\expandafter\def\csname SRef@maps:eq:D.C25\endcsname{\href{S6_SERIES_GRIND_III.pdf\#nameddest=tag-01YI}{III.{D.C25} (p.~127)} }
\expandafter\def\csname SRef@maps:eq:B.1.3\endcsname{\href{S6_SERIES_GRIND_III.pdf\#nameddest=tag-01UX}{III.{B.1.3} (p.~14)} }
\expandafter\def\csname SRef@maps:eq:D.C26\endcsname{\href{S6_SERIES_GRIND_III.pdf\#nameddest=tag-01YJ}{III.{D.C26} (p.~127)} }
\expandafter\def\csname SRef@maps:eq:D.C27\endcsname{\href{S6_SERIES_GRIND_III.pdf\#nameddest=tag-01YK}{III.{D.C27} (p.~127)} }
\expandafter\def\csname SRef@maps:eq:D.C28\endcsname{\href{S6_SERIES_GRIND_III.pdf\#nameddest=tag-01YL}{III.{D.C28} (p.~127)} }
\expandafter\def\csname SRef@maps:eq:D.C29\endcsname{\href{S6_SERIES_GRIND_III.pdf\#nameddest=tag-01YM}{III.{D.C29} (p.~128)} }
\expandafter\def\csname SRef@maps:eq:D.C210\endcsname{\href{S6_SERIES_GRIND_III.pdf\#nameddest=tag-01Y8}{III.{D.C210} (p.~128)} }
\expandafter\def\csname SRef@maps:eq:D.C211\endcsname{\href{S6_SERIES_GRIND_III.pdf\#nameddest=tag-01Y9}{III.{D.C211} (p.~128)} }
\expandafter\def\csname SRef@maps:eq:D.C212\endcsname{\href{S6_SERIES_GRIND_III.pdf\#nameddest=tag-01YA}{III.{D.C212} (p.~129)} }
\expandafter\def\csname SRef@maps:eq:D.C213\endcsname{\href{S6_SERIES_GRIND_III.pdf\#nameddest=tag-01YB}{III.{D.C213} (p.~129)} }
\expandafter\def\csname SRef@maps:eq:D.C214\endcsname{\href{S6_SERIES_GRIND_III.pdf\#nameddest=tag-01YC}{III.{D.C214} (p.~129)} }
\expandafter\def\csname SRef@maps:eq:D.C215\endcsname{\href{S6_SERIES_GRIND_III.pdf\#nameddest=tag-01YD}{III.{D.C215} (p.~129)} }
\expandafter\def\csname SRef@maps:eq:B.2.1\endcsname{\href{S6_SERIES_GRIND_III.pdf\#nameddest=tag-01UY}{III.{B.2.1} (p.~15)} }
\expandafter\def\csname SRef@maps:eq:D.C216\endcsname{\href{S6_SERIES_GRIND_III.pdf\#nameddest=tag-01YE}{III.{D.C216} (p.~130)} }
\expandafter\def\csname SRef@maps:eq:D.OP1\endcsname{\href{S6_SERIES_GRIND_III.pdf\#nameddest=tag-022C}{III.{D.OP1} (p.~131)} }
\expandafter\def\csname SRef@maps:eq:D.OP2\endcsname{\href{S6_SERIES_GRIND_III.pdf\#nameddest=tag-022F}{III.{D.OP2} (p.~131)} }
\expandafter\def\csname SRef@maps:eq:D.OP3\endcsname{\href{S6_SERIES_GRIND_III.pdf\#nameddest=tag-022G}{III.{D.OP3} (p.~131)} }
\expandafter\def\csname SRef@maps:eq:D.OP4\endcsname{\href{S6_SERIES_GRIND_III.pdf\#nameddest=tag-022H}{III.{D.OP4} (p.~131)} }
\expandafter\def\csname SRef@maps:eq:D.OP5\endcsname{\href{S6_SERIES_GRIND_III.pdf\#nameddest=tag-022I}{III.{D.OP5} (p.~132)} }
\expandafter\def\csname SRef@maps:eq:D.OP6\endcsname{\href{S6_SERIES_GRIND_III.pdf\#nameddest=tag-022K}{III.{D.OP6} (p.~132)} }
\expandafter\def\csname SRef@maps:eq:D.OP5a\endcsname{\href{S6_SERIES_GRIND_III.pdf\#nameddest=tag-022J}{III.{D.OP5a} (p.~132)} }
\expandafter\def\csname SRef@maps:eq:D.OP7\endcsname{\href{S6_SERIES_GRIND_III.pdf\#nameddest=tag-022L}{III.{D.OP7} (p.~133)} }
\expandafter\def\csname SRef@maps:eq:D.OP8\endcsname{\href{S6_SERIES_GRIND_III.pdf\#nameddest=tag-022M}{III.{D.OP8} (p.~133)} }
\expandafter\def\csname SRef@maps:eq:B.3.1\endcsname{\href{S6_SERIES_GRIND_III.pdf\#nameddest=tag-01UZ}{III.{B.3.1} (p.~16)} }
\expandafter\def\csname SRef@maps:eq:D.OP9\endcsname{\href{S6_SERIES_GRIND_III.pdf\#nameddest=tag-022N}{III.{D.OP9} (p.~133)} }
\expandafter\def\csname SRef@maps:eq:D.OP10\endcsname{\href{S6_SERIES_GRIND_III.pdf\#nameddest=tag-022D}{III.{D.OP10} (p.~133)} }
\expandafter\def\csname SRef@maps:eq:D.OP11\endcsname{\href{S6_SERIES_GRIND_III.pdf\#nameddest=tag-022E}{III.{D.OP11} (p.~133)} }
\expandafter\def\csname SRef@maps:eq:D.WP1\endcsname{\href{S6_SERIES_GRIND_III.pdf\#nameddest=tag-027O}{III.{D.WP1} (p.~134)} }
\expandafter\def\csname SRef@maps:eq:D.WP2\endcsname{\href{S6_SERIES_GRIND_III.pdf\#nameddest=tag-027S}{III.{D.WP2} (p.~134)} }
\expandafter\def\csname SRef@maps:eq:D.WP3\endcsname{\href{S6_SERIES_GRIND_III.pdf\#nameddest=tag-027T}{III.{D.WP3} (p.~134)} }
\expandafter\def\csname SRef@maps:eq:D.WP4\endcsname{\href{S6_SERIES_GRIND_III.pdf\#nameddest=tag-027U}{III.{D.WP4} (p.~134)} }
\expandafter\def\csname SRef@maps:eq:D.WP5\endcsname{\href{S6_SERIES_GRIND_III.pdf\#nameddest=tag-027V}{III.{D.WP5} (p.~135)} }
\expandafter\def\csname SRef@maps:eq:D.WP6\endcsname{\href{S6_SERIES_GRIND_III.pdf\#nameddest=tag-027W}{III.{D.WP6} (p.~135)} }
\expandafter\def\csname SRef@maps:eq:D.WP7\endcsname{\href{S6_SERIES_GRIND_III.pdf\#nameddest=tag-027X}{III.{D.WP7} (p.~135)} }
\expandafter\def\csname SRef@maps:eq:B.3.2\endcsname{\href{S6_SERIES_GRIND_III.pdf\#nameddest=tag-01V0}{III.{B.3.2} (p.~16)} }
\expandafter\def\csname SRef@maps:eq:D.WP8\endcsname{\href{S6_SERIES_GRIND_III.pdf\#nameddest=tag-027Y}{III.{D.WP8} (p.~135)} }
\expandafter\def\csname SRef@maps:eq:D.WP9\endcsname{\href{S6_SERIES_GRIND_III.pdf\#nameddest=tag-027Z}{III.{D.WP9} (p.~135)} }
\expandafter\def\csname SRef@maps:eq:D.WP10\endcsname{\href{S6_SERIES_GRIND_III.pdf\#nameddest=tag-027P}{III.{D.WP10} (p.~136)} }
\expandafter\def\csname SRef@maps:eq:D.WP11\endcsname{\href{S6_SERIES_GRIND_III.pdf\#nameddest=tag-027Q}{III.{D.WP11} (p.~137)} }
\expandafter\def\csname SRef@maps:eq:D.WP12\endcsname{\href{S6_SERIES_GRIND_III.pdf\#nameddest=tag-027R}{III.{D.WP12} (p.~138)} }
\expandafter\def\csname SRef@maps:eq:D.MP1\endcsname{\href{S6_SERIES_GRIND_III.pdf\#nameddest=tag-0209}{III.{D.MP1} (p.~138)} }
\expandafter\def\csname SRef@maps:eq:D.MP2\endcsname{\href{S6_SERIES_GRIND_III.pdf\#nameddest=tag-020B}{III.{D.MP2} (p.~139)} }
\expandafter\def\csname SRef@maps:eq:D.MP3\endcsname{\href{S6_SERIES_GRIND_III.pdf\#nameddest=tag-020C}{III.{D.MP3} (p.~139)} }
\expandafter\def\csname SRef@maps:eq:D.MP4\endcsname{\href{S6_SERIES_GRIND_III.pdf\#nameddest=tag-020D}{III.{D.MP4} (p.~139)} }
\expandafter\def\csname SRef@maps:eq:D.MP5\endcsname{\href{S6_SERIES_GRIND_III.pdf\#nameddest=tag-020E}{III.{D.MP5} (p.~140)} }
\expandafter\def\csname SRef@maps:eq:B.3.3\endcsname{\href{S6_SERIES_GRIND_III.pdf\#nameddest=tag-01V1}{III.{B.3.3} (p.~16)} }
\expandafter\def\csname SRef@maps:eq:D.MP6\endcsname{\href{S6_SERIES_GRIND_III.pdf\#nameddest=tag-020F}{III.{D.MP6} (p.~140)} }
\expandafter\def\csname SRef@maps:eq:D.MP7\endcsname{\href{S6_SERIES_GRIND_III.pdf\#nameddest=tag-020G}{III.{D.MP7} (p.~140)} }
\expandafter\def\csname SRef@maps:eq:D.MP8\endcsname{\href{S6_SERIES_GRIND_III.pdf\#nameddest=tag-020H}{III.{D.MP8} (p.~140)} }
\expandafter\def\csname SRef@maps:eq:D.MP9\endcsname{\href{S6_SERIES_GRIND_III.pdf\#nameddest=tag-020I}{III.{D.MP9} (p.~141)} }
\expandafter\def\csname SRef@maps:eq:D.MP10\endcsname{\href{S6_SERIES_GRIND_III.pdf\#nameddest=tag-020A}{III.{D.MP10} (p.~141)} }
\expandafter\def\csname SRef@maps:eq:D.T31\endcsname{\href{S6_SERIES_GRIND_III.pdf\#nameddest=tag-025C}{III.{D.T31} (p.~142)} }
\expandafter\def\csname SRef@maps:eq:D.T32\endcsname{\href{S6_SERIES_GRIND_III.pdf\#nameddest=tag-025D}{III.{D.T32} (p.~142)} }
\expandafter\def\csname SRef@maps:eq:D.T33\endcsname{\href{S6_SERIES_GRIND_III.pdf\#nameddest=tag-025E}{III.{D.T33} (p.~142)} }
\expandafter\def\csname SRef@maps:eq:D.T34\endcsname{\href{S6_SERIES_GRIND_III.pdf\#nameddest=tag-025F}{III.{D.T34} (p.~143)} }
\expandafter\def\csname SRef@maps:eq:D.T35\endcsname{\href{S6_SERIES_GRIND_III.pdf\#nameddest=tag-025G}{III.{D.T35} (p.~143)} }
\expandafter\def\csname SRef@maps:eq:B.3.4\endcsname{\href{S6_SERIES_GRIND_III.pdf\#nameddest=tag-01V2}{III.{B.3.4} (p.~16)} }
\expandafter\def\csname SRef@maps:eq:D.T36\endcsname{\href{S6_SERIES_GRIND_III.pdf\#nameddest=tag-025H}{III.{D.T36} (p.~143)} }
\expandafter\def\csname SRef@maps:eq:D.T37\endcsname{\href{S6_SERIES_GRIND_III.pdf\#nameddest=tag-025I}{III.{D.T37} (p.~144)} }
\expandafter\def\csname SRef@maps:eq:D.T38\endcsname{\href{S6_SERIES_GRIND_III.pdf\#nameddest=tag-025J}{III.{D.T38} (p.~144)} }
\expandafter\def\csname SRef@maps:eq:D.T39\endcsname{\href{S6_SERIES_GRIND_III.pdf\#nameddest=tag-025K}{III.{D.T39} (p.~144)} }
\expandafter\def\csname SRef@maps:eq:D.OW1\endcsname{\href{S6_SERIES_GRIND_III.pdf\#nameddest=tag-022W}{III.{D.OW1} (p.~145)} }
\expandafter\def\csname SRef@maps:eq:D.OW2\endcsname{\href{S6_SERIES_GRIND_III.pdf\#nameddest=tag-022X}{III.{D.OW2} (p.~145)} }
\expandafter\def\csname SRef@maps:eq:D.OW3\endcsname{\href{S6_SERIES_GRIND_III.pdf\#nameddest=tag-022Y}{III.{D.OW3} (p.~145)} }
\expandafter\def\csname SRef@maps:eq:D.OW4\endcsname{\href{S6_SERIES_GRIND_III.pdf\#nameddest=tag-022Z}{III.{D.OW4} (p.~145)} }
\expandafter\def\csname SRef@maps:eq:D.OW5\endcsname{\href{S6_SERIES_GRIND_III.pdf\#nameddest=tag-0230}{III.{D.OW5} (p.~145)} }
\expandafter\def\csname SRef@maps:eq:D.OW6\endcsname{\href{S6_SERIES_GRIND_III.pdf\#nameddest=tag-0231}{III.{D.OW6} (p.~146)} }
\expandafter\def\csname SRef@maps:eq:B.3.5\endcsname{\href{S6_SERIES_GRIND_III.pdf\#nameddest=tag-01V3}{III.{B.3.5} (p.~18)} }
\expandafter\def\csname SRef@maps:eq:D.OW7\endcsname{\href{S6_SERIES_GRIND_III.pdf\#nameddest=tag-0232}{III.{D.OW7} (p.~147)} }
\expandafter\def\csname SRef@maps:eq:D.BC1\endcsname{\href{S6_SERIES_GRIND_III.pdf\#nameddest=tag-01XR}{III.{D.BC1} (p.~148)} }
\expandafter\def\csname SRef@maps:eq:D.BC2\endcsname{\href{S6_SERIES_GRIND_III.pdf\#nameddest=tag-01XY}{III.{D.BC2} (p.~148)} }
\expandafter\def\csname SRef@maps:eq:D.BC3\endcsname{\href{S6_SERIES_GRIND_III.pdf\#nameddest=tag-01XZ}{III.{D.BC3} (p.~148)} }
\expandafter\def\csname SRef@maps:eq:D.BC4\endcsname{\href{S6_SERIES_GRIND_III.pdf\#nameddest=tag-01Y0}{III.{D.BC4} (p.~148)} }
\expandafter\def\csname SRef@maps:eq:D.BC5\endcsname{\href{S6_SERIES_GRIND_III.pdf\#nameddest=tag-01Y1}{III.{D.BC5} (p.~148)} }
\expandafter\def\csname SRef@maps:eq:D.BC6\endcsname{\href{S6_SERIES_GRIND_III.pdf\#nameddest=tag-01Y2}{III.{D.BC6} (p.~149)} }
\expandafter\def\csname SRef@maps:eq:D.BC7\endcsname{\href{S6_SERIES_GRIND_III.pdf\#nameddest=tag-01Y3}{III.{D.BC7} (p.~149)} }
\expandafter\def\csname SRef@maps:eq:D.BC8\endcsname{\href{S6_SERIES_GRIND_III.pdf\#nameddest=tag-01Y4}{III.{D.BC8} (p.~149)} }
\expandafter\def\csname SRef@maps:eq:D.BC9\endcsname{\href{S6_SERIES_GRIND_III.pdf\#nameddest=tag-01Y5}{III.{D.BC9} (p.~149)} }
\expandafter\def\csname SRef@maps:eq:B.4.1\endcsname{\href{S6_SERIES_GRIND_III.pdf\#nameddest=tag-01V4}{III.{B.4.1} (p.~19)} }
\expandafter\def\csname SRef@maps:eq:D.BC10\endcsname{\href{S6_SERIES_GRIND_III.pdf\#nameddest=tag-01XS}{III.{D.BC10} (p.~149)} }
\expandafter\def\csname SRef@maps:eq:D.BC11\endcsname{\href{S6_SERIES_GRIND_III.pdf\#nameddest=tag-01XT}{III.{D.BC11} (p.~149)} }
\expandafter\def\csname SRef@maps:eq:D.BC12\endcsname{\href{S6_SERIES_GRIND_III.pdf\#nameddest=tag-01XU}{III.{D.BC12} (p.~149)} }
\expandafter\def\csname SRef@maps:eq:D.BC13\endcsname{\href{S6_SERIES_GRIND_III.pdf\#nameddest=tag-01XV}{III.{D.BC13} (p.~150)} }
\expandafter\def\csname SRef@maps:eq:D.BC14\endcsname{\href{S6_SERIES_GRIND_III.pdf\#nameddest=tag-01XW}{III.{D.BC14} (p.~150)} }
\expandafter\def\csname SRef@maps:eq:D.BC15\endcsname{\href{S6_SERIES_GRIND_III.pdf\#nameddest=tag-01XX}{III.{D.BC15} (p.~150)} }
\expandafter\def\csname SRef@maps:eq:D.MT1\endcsname{\href{S6_SERIES_GRIND_III.pdf\#nameddest=tag-020J}{III.{D.MT1} (p.~151)} }
\expandafter\def\csname SRef@maps:eq:D.MT2\endcsname{\href{S6_SERIES_GRIND_III.pdf\#nameddest=tag-020K}{III.{D.MT2} (p.~152)} }
\expandafter\def\csname SRef@maps:eq:D.MT3\endcsname{\href{S6_SERIES_GRIND_III.pdf\#nameddest=tag-020L}{III.{D.MT3} (p.~153)} }
\expandafter\def\csname SRef@maps:eq:D.GC1\endcsname{\href{S6_SERIES_GRIND_III.pdf\#nameddest=tag-01YY}{III.{D.GC1} (p.~153)} }
\expandafter\def\csname SRef@maps:eq:B.4.2\endcsname{\href{S6_SERIES_GRIND_III.pdf\#nameddest=tag-01V5}{III.{B.4.2} (p.~19)} }
\expandafter\def\csname SRef@maps:eq:D.GC2\endcsname{\href{S6_SERIES_GRIND_III.pdf\#nameddest=tag-01Z4}{III.{D.GC2} (p.~153)} }
\expandafter\def\csname SRef@maps:eq:D.GC3\endcsname{\href{S6_SERIES_GRIND_III.pdf\#nameddest=tag-01Z5}{III.{D.GC3} (p.~154)} }
\expandafter\def\csname SRef@maps:eq:D.GC4\endcsname{\href{S6_SERIES_GRIND_III.pdf\#nameddest=tag-01Z6}{III.{D.GC4} (p.~154)} }
\expandafter\def\csname SRef@maps:eq:D.GC5\endcsname{\href{S6_SERIES_GRIND_III.pdf\#nameddest=tag-01Z7}{III.{D.GC5} (p.~154)} }
\expandafter\def\csname SRef@maps:eq:D.GC6\endcsname{\href{S6_SERIES_GRIND_III.pdf\#nameddest=tag-01Z8}{III.{D.GC6} (p.~154)} }
\expandafter\def\csname SRef@maps:eq:D.GC7\endcsname{\href{S6_SERIES_GRIND_III.pdf\#nameddest=tag-01Z9}{III.{D.GC7} (p.~154)} }
\expandafter\def\csname SRef@maps:eq:D.GC8\endcsname{\href{S6_SERIES_GRIND_III.pdf\#nameddest=tag-01ZA}{III.{D.GC8} (p.~155)} }
\expandafter\def\csname SRef@maps:eq:D.GC9\endcsname{\href{S6_SERIES_GRIND_III.pdf\#nameddest=tag-01ZB}{III.{D.GC9} (p.~155)} }
\expandafter\def\csname SRef@maps:eq:D.GC10\endcsname{\href{S6_SERIES_GRIND_III.pdf\#nameddest=tag-01YZ}{III.{D.GC10} (p.~155)} }
\expandafter\def\csname SRef@maps:eq:D.GC11\endcsname{\href{S6_SERIES_GRIND_III.pdf\#nameddest=tag-01Z0}{III.{D.GC11} (p.~155)} }
\expandafter\def\csname SRef@maps:eq:B.4.3\endcsname{\href{S6_SERIES_GRIND_III.pdf\#nameddest=tag-01V6}{III.{B.4.3} (p.~20)} }
\expandafter\def\csname SRef@maps:eq:D.GC12\endcsname{\href{S6_SERIES_GRIND_III.pdf\#nameddest=tag-01Z1}{III.{D.GC12} (p.~155)} }
\expandafter\def\csname SRef@maps:eq:D.GC13\endcsname{\href{S6_SERIES_GRIND_III.pdf\#nameddest=tag-01Z2}{III.{D.GC13} (p.~155)} }
\expandafter\def\csname SRef@maps:eq:D.GC14\endcsname{\href{S6_SERIES_GRIND_III.pdf\#nameddest=tag-01Z3}{III.{D.GC14} (p.~156)} }
\expandafter\def\csname SRef@maps:eq:D44.1\endcsname{\href{S6_SERIES_GRIND_III.pdf\#nameddest=tag-0282}{III.{D44.1} (p.~157)} }
\expandafter\def\csname SRef@maps:eq:D44.0a\endcsname{\href{S6_SERIES_GRIND_III.pdf\#nameddest=tag-0280}{III.{D44.0a} (p.~157)} }
\expandafter\def\csname SRef@maps:eq:D44.0b\endcsname{\href{S6_SERIES_GRIND_III.pdf\#nameddest=tag-0281}{III.{D44.0b} (p.~157)} }
\expandafter\def\csname SRef@maps:eq:D44.2\endcsname{\href{S6_SERIES_GRIND_III.pdf\#nameddest=tag-0288}{III.{D44.2} (p.~158)} }
\expandafter\def\csname SRef@maps:eq:D44.3\endcsname{\href{S6_SERIES_GRIND_III.pdf\#nameddest=tag-0289}{III.{D44.3} (p.~158)} }
\expandafter\def\csname SRef@maps:eq:D44.4\endcsname{\href{S6_SERIES_GRIND_III.pdf\#nameddest=tag-028A}{III.{D44.4} (p.~159)} }
\expandafter\def\csname SRef@maps:eq:D44.5\endcsname{\href{S6_SERIES_GRIND_III.pdf\#nameddest=tag-028B}{III.{D44.5} (p.~159)} }
\expandafter\def\csname SRef@maps:eq:B.5.1\endcsname{\href{S6_SERIES_GRIND_III.pdf\#nameddest=tag-01V7}{III.{B.5.1} (p.~20)} }
\expandafter\def\csname SRef@maps:eq:D44.6\endcsname{\href{S6_SERIES_GRIND_III.pdf\#nameddest=tag-028C}{III.{D44.6} (p.~159)} }
\expandafter\def\csname SRef@maps:eq:D44.7\endcsname{\href{S6_SERIES_GRIND_III.pdf\#nameddest=tag-028D}{III.{D44.7} (p.~159)} }
\expandafter\def\csname SRef@maps:eq:D44.8\endcsname{\href{S6_SERIES_GRIND_III.pdf\#nameddest=tag-028E}{III.{D44.8} (p.~160)} }
\expandafter\def\csname SRef@maps:eq:D44.9\endcsname{\href{S6_SERIES_GRIND_III.pdf\#nameddest=tag-028F}{III.{D44.9} (p.~160)} }
\expandafter\def\csname SRef@maps:eq:D44.10\endcsname{\href{S6_SERIES_GRIND_III.pdf\#nameddest=tag-0283}{III.{D44.10} (p.~160)} }
\expandafter\def\csname SRef@maps:eq:D44.11\endcsname{\href{S6_SERIES_GRIND_III.pdf\#nameddest=tag-0284}{III.{D44.11} (p.~160)} }
\expandafter\def\csname SRef@maps:eq:D44.12\endcsname{\href{S6_SERIES_GRIND_III.pdf\#nameddest=tag-0285}{III.{D44.12} (p.~160)} }
\expandafter\def\csname SRef@maps:eq:D44.13\endcsname{\href{S6_SERIES_GRIND_III.pdf\#nameddest=tag-0286}{III.{D44.13} (p.~161)} }
\expandafter\def\csname SRef@maps:eq:D44.14\endcsname{\href{S6_SERIES_GRIND_III.pdf\#nameddest=tag-0287}{III.{D44.14} (p.~161)} }
\expandafter\def\csname SRef@maps:eq:D45.1\endcsname{\href{S6_SERIES_GRIND_III.pdf\#nameddest=tag-028G}{III.{D45.1} (p.~162)} }
\expandafter\def\csname SRef@maps:eq:B.5.2\endcsname{\href{S6_SERIES_GRIND_III.pdf\#nameddest=tag-01VH}{III.{B.5.2} (p.~21)} }
\expandafter\def\csname SRef@maps:eq:D45.2\endcsname{\href{S6_SERIES_GRIND_III.pdf\#nameddest=tag-028R}{III.{D45.2} (p.~162)} }
\expandafter\def\csname SRef@maps:eq:D45.3\endcsname{\href{S6_SERIES_GRIND_III.pdf\#nameddest=tag-028V}{III.{D45.3} (p.~162)} }
\expandafter\def\csname SRef@maps:eq:D45.4\endcsname{\href{S6_SERIES_GRIND_III.pdf\#nameddest=tag-028W}{III.{D45.4} (p.~162)} }
\expandafter\def\csname SRef@maps:eq:D45.5\endcsname{\href{S6_SERIES_GRIND_III.pdf\#nameddest=tag-028X}{III.{D45.5} (p.~162)} }
\expandafter\def\csname SRef@maps:eq:D45.6\endcsname{\href{S6_SERIES_GRIND_III.pdf\#nameddest=tag-028Y}{III.{D45.6} (p.~163)} }
\expandafter\def\csname SRef@maps:eq:D45.7\endcsname{\href{S6_SERIES_GRIND_III.pdf\#nameddest=tag-028Z}{III.{D45.7} (p.~163)} }
\expandafter\def\csname SRef@maps:eq:D45.8\endcsname{\href{S6_SERIES_GRIND_III.pdf\#nameddest=tag-0290}{III.{D45.8} (p.~163)} }
\expandafter\def\csname SRef@maps:eq:D45.9\endcsname{\href{S6_SERIES_GRIND_III.pdf\#nameddest=tag-0291}{III.{D45.9} (p.~163)} }
\expandafter\def\csname SRef@maps:eq:D45.10\endcsname{\href{S6_SERIES_GRIND_III.pdf\#nameddest=tag-028H}{III.{D45.10} (p.~164)} }
\expandafter\def\csname SRef@maps:eq:D45.11\endcsname{\href{S6_SERIES_GRIND_III.pdf\#nameddest=tag-028I}{III.{D45.11} (p.~164)} }
\expandafter\def\csname SRef@maps:eq:B.5.3\endcsname{\href{S6_SERIES_GRIND_III.pdf\#nameddest=tag-01VI}{III.{B.5.3} (p.~21)} }
\expandafter\def\csname SRef@maps:eq:D45.12\endcsname{\href{S6_SERIES_GRIND_III.pdf\#nameddest=tag-028J}{III.{D45.12} (p.~164)} }
\expandafter\def\csname SRef@maps:eq:D45.13\endcsname{\href{S6_SERIES_GRIND_III.pdf\#nameddest=tag-028K}{III.{D45.13} (p.~164)} }
\expandafter\def\csname SRef@maps:eq:D45.14\endcsname{\href{S6_SERIES_GRIND_III.pdf\#nameddest=tag-028L}{III.{D45.14} (p.~164)} }
\expandafter\def\csname SRef@maps:eq:D45.15\endcsname{\href{S6_SERIES_GRIND_III.pdf\#nameddest=tag-028M}{III.{D45.15} (p.~164)} }
\expandafter\def\csname SRef@maps:eq:D45.16\endcsname{\href{S6_SERIES_GRIND_III.pdf\#nameddest=tag-028N}{III.{D45.16} (p.~165)} }
\expandafter\def\csname SRef@maps:eq:D45.17\endcsname{\href{S6_SERIES_GRIND_III.pdf\#nameddest=tag-028O}{III.{D45.17} (p.~165)} }
\expandafter\def\csname SRef@maps:eq:D45.18\endcsname{\href{S6_SERIES_GRIND_III.pdf\#nameddest=tag-028P}{III.{D45.18} (p.~165)} }
\expandafter\def\csname SRef@maps:eq:D45.19\endcsname{\href{S6_SERIES_GRIND_III.pdf\#nameddest=tag-028Q}{III.{D45.19} (p.~165)} }
\expandafter\def\csname SRef@maps:eq:D45.20\endcsname{\href{S6_SERIES_GRIND_III.pdf\#nameddest=tag-028S}{III.{D45.20} (p.~165)} }
\expandafter\def\csname SRef@maps:eq:D45.21\endcsname{\href{S6_SERIES_GRIND_III.pdf\#nameddest=tag-028T}{III.{D45.21} (p.~166)} }
\expandafter\def\csname SRef@maps:eq:B.5.4\endcsname{\href{S6_SERIES_GRIND_III.pdf\#nameddest=tag-01VJ}{III.{B.5.4} (p.~21)} }
\expandafter\def\csname SRef@maps:eq:D45.22\endcsname{\href{S6_SERIES_GRIND_III.pdf\#nameddest=tag-028U}{III.{D45.22} (p.~166)} }
\expandafter\def\csname SRef@maps:eq:D.GI1\endcsname{\href{S6_SERIES_GRIND_III.pdf\#nameddest=tag-01ZC}{III.{D.GI1} (p.~167)} }
\expandafter\def\csname SRef@maps:eq:D.GI2\endcsname{\href{S6_SERIES_GRIND_III.pdf\#nameddest=tag-01ZF}{III.{D.GI2} (p.~167)} }
\expandafter\def\csname SRef@maps:eq:D.GI3\endcsname{\href{S6_SERIES_GRIND_III.pdf\#nameddest=tag-01ZG}{III.{D.GI3} (p.~167)} }
\expandafter\def\csname SRef@maps:eq:D.GI4\endcsname{\href{S6_SERIES_GRIND_III.pdf\#nameddest=tag-01ZH}{III.{D.GI4} (p.~167)} }
\expandafter\def\csname SRef@maps:eq:D.GI5\endcsname{\href{S6_SERIES_GRIND_III.pdf\#nameddest=tag-01ZI}{III.{D.GI5} (p.~168)} }
\expandafter\def\csname SRef@maps:eq:D.GI6\endcsname{\href{S6_SERIES_GRIND_III.pdf\#nameddest=tag-01ZJ}{III.{D.GI6} (p.~168)} }
\expandafter\def\csname SRef@maps:eq:D.GI7\endcsname{\href{S6_SERIES_GRIND_III.pdf\#nameddest=tag-01ZK}{III.{D.GI7} (p.~168)} }
\expandafter\def\csname SRef@maps:eq:D.GI8\endcsname{\href{S6_SERIES_GRIND_III.pdf\#nameddest=tag-01ZL}{III.{D.GI8} (p.~168)} }
\expandafter\def\csname SRef@maps:eq:D.GI9\endcsname{\href{S6_SERIES_GRIND_III.pdf\#nameddest=tag-01ZM}{III.{D.GI9} (p.~170)} }
\expandafter\def\csname SRef@maps:eq:B.5.5\endcsname{\href{S6_SERIES_GRIND_III.pdf\#nameddest=tag-01VK}{III.{B.5.5} (p.~22)} }
\expandafter\def\csname SRef@maps:eq:D.GI10\endcsname{\href{S6_SERIES_GRIND_III.pdf\#nameddest=tag-01ZD}{III.{D.GI10} (p.~171)} }
\expandafter\def\csname SRef@maps:eq:D.GI11\endcsname{\href{S6_SERIES_GRIND_III.pdf\#nameddest=tag-01ZE}{III.{D.GI11} (p.~171)} }
\expandafter\def\csname SRef@maps:eq:D47.1\endcsname{\href{S6_SERIES_GRIND_III.pdf\#nameddest=tag-0292}{III.{D47.1} (p.~172)} }
\expandafter\def\csname SRef@maps:eq:D47.2\endcsname{\href{S6_SERIES_GRIND_III.pdf\#nameddest=tag-0297}{III.{D47.2} (p.~172)} }
\expandafter\def\csname SRef@maps:eq:D47.3\endcsname{\href{S6_SERIES_GRIND_III.pdf\#nameddest=tag-0298}{III.{D47.3} (p.~172)} }
\expandafter\def\csname SRef@maps:eq:D47.4\endcsname{\href{S6_SERIES_GRIND_III.pdf\#nameddest=tag-0299}{III.{D47.4} (p.~172)} }
\expandafter\def\csname SRef@maps:eq:D47.5\endcsname{\href{S6_SERIES_GRIND_III.pdf\#nameddest=tag-029A}{III.{D47.5} (p.~173)} }
\expandafter\def\csname SRef@maps:eq:D47.6\endcsname{\href{S6_SERIES_GRIND_III.pdf\#nameddest=tag-029B}{III.{D47.6} (p.~173)} }
\expandafter\def\csname SRef@maps:eq:D47.7\endcsname{\href{S6_SERIES_GRIND_III.pdf\#nameddest=tag-029C}{III.{D47.7} (p.~173)} }
\expandafter\def\csname SRef@maps:eq:D47.8\endcsname{\href{S6_SERIES_GRIND_III.pdf\#nameddest=tag-029D}{III.{D47.8} (p.~173)} }
\expandafter\def\csname SRef@maps:eq:B.5.6\endcsname{\href{S6_SERIES_GRIND_III.pdf\#nameddest=tag-01VL}{III.{B.5.6} (p.~22)} }
\expandafter\def\csname SRef@maps:eq:D47.9\endcsname{\href{S6_SERIES_GRIND_III.pdf\#nameddest=tag-029E}{III.{D47.9} (p.~173)} }
\expandafter\def\csname SRef@maps:eq:D47.10\endcsname{\href{S6_SERIES_GRIND_III.pdf\#nameddest=tag-0293}{III.{D47.10} (p.~174)} }
\expandafter\def\csname SRef@maps:eq:D47.11\endcsname{\href{S6_SERIES_GRIND_III.pdf\#nameddest=tag-0294}{III.{D47.11} (p.~174)} }
\expandafter\def\csname SRef@maps:eq:D47.12\endcsname{\href{S6_SERIES_GRIND_III.pdf\#nameddest=tag-0295}{III.{D47.12} (p.~174)} }
\expandafter\def\csname SRef@maps:eq:D47.13\endcsname{\href{S6_SERIES_GRIND_III.pdf\#nameddest=tag-0296}{III.{D47.13} (p.~174)} }
\expandafter\def\csname SRef@maps:eq:D.WI1\endcsname{\href{S6_SERIES_GRIND_III.pdf\#nameddest=tag-0272}{III.{D.WI1} (p.~175)} }
\expandafter\def\csname SRef@maps:eq:D.WI2\endcsname{\href{S6_SERIES_GRIND_III.pdf\#nameddest=tag-027D}{III.{D.WI2} (p.~175)} }
\expandafter\def\csname SRef@maps:eq:D.WI3\endcsname{\href{S6_SERIES_GRIND_III.pdf\#nameddest=tag-027H}{III.{D.WI3} (p.~175)} }
\expandafter\def\csname SRef@maps:eq:D.WI14\endcsname{\href{S6_SERIES_GRIND_III.pdf\#nameddest=tag-0277}{III.{D.WI14} (p.~175)} }
\expandafter\def\csname SRef@maps:eq:D.WI15\endcsname{\href{S6_SERIES_GRIND_III.pdf\#nameddest=tag-0278}{III.{D.WI15} (p.~175)} }
\expandafter\def\csname SRef@maps:eq:B.5.7\endcsname{\href{S6_SERIES_GRIND_III.pdf\#nameddest=tag-01VM}{III.{B.5.7} (p.~22)} }
\expandafter\def\csname SRef@maps:eq:D.WI16\endcsname{\href{S6_SERIES_GRIND_III.pdf\#nameddest=tag-0279}{III.{D.WI16} (p.~175)} }
\expandafter\def\csname SRef@maps:eq:D.WI4\endcsname{\href{S6_SERIES_GRIND_III.pdf\#nameddest=tag-027I}{III.{D.WI4} (p.~176)} }
\expandafter\def\csname SRef@maps:eq:D.WI5\endcsname{\href{S6_SERIES_GRIND_III.pdf\#nameddest=tag-027J}{III.{D.WI5} (p.~176)} }
\expandafter\def\csname SRef@maps:eq:D.WI6\endcsname{\href{S6_SERIES_GRIND_III.pdf\#nameddest=tag-027K}{III.{D.WI6} (p.~177)} }
\expandafter\def\csname SRef@maps:eq:D.WI7\endcsname{\href{S6_SERIES_GRIND_III.pdf\#nameddest=tag-027L}{III.{D.WI7} (p.~177)} }
\expandafter\def\csname SRef@maps:eq:D.WI8\endcsname{\href{S6_SERIES_GRIND_III.pdf\#nameddest=tag-027M}{III.{D.WI8} (p.~177)} }
\expandafter\def\csname SRef@maps:eq:D.WI9\endcsname{\href{S6_SERIES_GRIND_III.pdf\#nameddest=tag-027N}{III.{D.WI9} (p.~177)} }
\expandafter\def\csname SRef@maps:eq:D.WI10\endcsname{\href{S6_SERIES_GRIND_III.pdf\#nameddest=tag-0273}{III.{D.WI10} (p.~177)} }
\expandafter\def\csname SRef@maps:eq:D.WI11\endcsname{\href{S6_SERIES_GRIND_III.pdf\#nameddest=tag-0274}{III.{D.WI11} (p.~177)} }
\expandafter\def\csname SRef@maps:eq:D.WI12\endcsname{\href{S6_SERIES_GRIND_III.pdf\#nameddest=tag-0275}{III.{D.WI12} (p.~177)} }
\expandafter\def\csname SRef@maps:eq:B.5.8\endcsname{\href{S6_SERIES_GRIND_III.pdf\#nameddest=tag-01VN}{III.{B.5.8} (p.~23)} }
\expandafter\def\csname SRef@maps:eq:D.WI13\endcsname{\href{S6_SERIES_GRIND_III.pdf\#nameddest=tag-0276}{III.{D.WI13} (p.~178)} }
\expandafter\def\csname SRef@maps:eq:D.WI17\endcsname{\href{S6_SERIES_GRIND_III.pdf\#nameddest=tag-027A}{III.{D.WI17} (p.~178)} }
\expandafter\def\csname SRef@maps:eq:D.WI18\endcsname{\href{S6_SERIES_GRIND_III.pdf\#nameddest=tag-027B}{III.{D.WI18} (p.~178)} }
\expandafter\def\csname SRef@maps:eq:D.WI19\endcsname{\href{S6_SERIES_GRIND_III.pdf\#nameddest=tag-027C}{III.{D.WI19} (p.~178)} }
\expandafter\def\csname SRef@maps:eq:D.WI20\endcsname{\href{S6_SERIES_GRIND_III.pdf\#nameddest=tag-027E}{III.{D.WI20} (p.~179)} }
\expandafter\def\csname SRef@maps:eq:D.WI21\endcsname{\href{S6_SERIES_GRIND_III.pdf\#nameddest=tag-027F}{III.{D.WI21} (p.~179)} }
\expandafter\def\csname SRef@maps:eq:D.TP1\endcsname{\href{S6_SERIES_GRIND_III.pdf\#nameddest=tag-025R}{III.{D.TP1} (p.~179)} }
\expandafter\def\csname SRef@maps:eq:D.TP2\endcsname{\href{S6_SERIES_GRIND_III.pdf\#nameddest=tag-025S}{III.{D.TP2} (p.~180)} }
\expandafter\def\csname SRef@maps:eq:D.TP3\endcsname{\href{S6_SERIES_GRIND_III.pdf\#nameddest=tag-025T}{III.{D.TP3} (p.~180)} }
\expandafter\def\csname SRef@maps:eq:B.5.9\endcsname{\href{S6_SERIES_GRIND_III.pdf\#nameddest=tag-01VO}{III.{B.5.9} (p.~24)} }
\expandafter\def\csname SRef@maps:eq:D.TP4\endcsname{\href{S6_SERIES_GRIND_III.pdf\#nameddest=tag-025U}{III.{D.TP4} (p.~180)} }
\expandafter\def\csname SRef@maps:eq:D.TP5\endcsname{\href{S6_SERIES_GRIND_III.pdf\#nameddest=tag-025V}{III.{D.TP5} (p.~180)} }
\expandafter\def\csname SRef@maps:eq:D.TP6\endcsname{\href{S6_SERIES_GRIND_III.pdf\#nameddest=tag-025W}{III.{D.TP6} (p.~180)} }
\expandafter\def\csname SRef@maps:eq:D.TP7\endcsname{\href{S6_SERIES_GRIND_III.pdf\#nameddest=tag-025X}{III.{D.TP7} (p.~180)} }
\expandafter\def\csname SRef@maps:eq:D.TP8\endcsname{\href{S6_SERIES_GRIND_III.pdf\#nameddest=tag-025Y}{III.{D.TP8} (p.~180)} }
\expandafter\def\csname SRef@maps:eq:D.WI22\endcsname{\href{S6_SERIES_GRIND_III.pdf\#nameddest=tag-027G}{III.{D.WI22} (p.~181)} }
\expandafter\def\csname SRef@maps:eq:D.SL1\endcsname{\href{S6_SERIES_GRIND_III.pdf\#nameddest=tag-0247}{III.{D.SL1} (p.~182)} }
\expandafter\def\csname SRef@maps:eq:D.SL2\endcsname{\href{S6_SERIES_GRIND_III.pdf\#nameddest=tag-024I}{III.{D.SL2} (p.~182)} }
\expandafter\def\csname SRef@maps:eq:D.SL3\endcsname{\href{S6_SERIES_GRIND_III.pdf\#nameddest=tag-024O}{III.{D.SL3} (p.~182)} }
\expandafter\def\csname SRef@maps:eq:D.SL4\endcsname{\href{S6_SERIES_GRIND_III.pdf\#nameddest=tag-024P}{III.{D.SL4} (p.~182)} }
\expandafter\def\csname SRef@maps:eq:B.5.10\endcsname{\href{S6_SERIES_GRIND_III.pdf\#nameddest=tag-01V8}{III.{B.5.10} (p.~24)} }
\expandafter\def\csname SRef@maps:eq:D.SL5\endcsname{\href{S6_SERIES_GRIND_III.pdf\#nameddest=tag-024Q}{III.{D.SL5} (p.~182)} }
\expandafter\def\csname SRef@maps:eq:D.SL6\endcsname{\href{S6_SERIES_GRIND_III.pdf\#nameddest=tag-024R}{III.{D.SL6} (p.~183)} }
\expandafter\def\csname SRef@maps:eq:D.SL7\endcsname{\href{S6_SERIES_GRIND_III.pdf\#nameddest=tag-024S}{III.{D.SL7} (p.~183)} }
\expandafter\def\csname SRef@maps:eq:D.SL8\endcsname{\href{S6_SERIES_GRIND_III.pdf\#nameddest=tag-024T}{III.{D.SL8} (p.~183)} }
\expandafter\def\csname SRef@maps:eq:D.SL9\endcsname{\href{S6_SERIES_GRIND_III.pdf\#nameddest=tag-024U}{III.{D.SL9} (p.~183)} }
\expandafter\def\csname SRef@maps:eq:D.SL10\endcsname{\href{S6_SERIES_GRIND_III.pdf\#nameddest=tag-0248}{III.{D.SL10} (p.~183)} }
\expandafter\def\csname SRef@maps:eq:D.SL11\endcsname{\href{S6_SERIES_GRIND_III.pdf\#nameddest=tag-0249}{III.{D.SL11} (p.~183)} }
\expandafter\def\csname SRef@maps:eq:D.SL12\endcsname{\href{S6_SERIES_GRIND_III.pdf\#nameddest=tag-024A}{III.{D.SL12} (p.~184)} }
\expandafter\def\csname SRef@maps:eq:D.SL13\endcsname{\href{S6_SERIES_GRIND_III.pdf\#nameddest=tag-024B}{III.{D.SL13} (p.~184)} }
\expandafter\def\csname SRef@maps:eq:D.SL14\endcsname{\href{S6_SERIES_GRIND_III.pdf\#nameddest=tag-024C}{III.{D.SL14} (p.~184)} }
\expandafter\def\csname SRef@maps:eq:B.5.11\endcsname{\href{S6_SERIES_GRIND_III.pdf\#nameddest=tag-01V9}{III.{B.5.11} (p.~24)} }
\expandafter\def\csname SRef@maps:eq:D.SL15\endcsname{\href{S6_SERIES_GRIND_III.pdf\#nameddest=tag-024D}{III.{D.SL15} (p.~185)} }
\expandafter\def\csname SRef@maps:eq:D.SL16\endcsname{\href{S6_SERIES_GRIND_III.pdf\#nameddest=tag-024E}{III.{D.SL16} (p.~185)} }
\expandafter\def\csname SRef@maps:eq:D.SL17\endcsname{\href{S6_SERIES_GRIND_III.pdf\#nameddest=tag-024F}{III.{D.SL17} (p.~185)} }
\expandafter\def\csname SRef@maps:eq:D.SL18\endcsname{\href{S6_SERIES_GRIND_III.pdf\#nameddest=tag-024G}{III.{D.SL18} (p.~185)} }
\expandafter\def\csname SRef@maps:eq:D.SL19\endcsname{\href{S6_SERIES_GRIND_III.pdf\#nameddest=tag-024H}{III.{D.SL19} (p.~185)} }
\expandafter\def\csname SRef@maps:eq:D.SL20\endcsname{\href{S6_SERIES_GRIND_III.pdf\#nameddest=tag-024J}{III.{D.SL20} (p.~186)} }
\expandafter\def\csname SRef@maps:eq:D.SL21\endcsname{\href{S6_SERIES_GRIND_III.pdf\#nameddest=tag-024K}{III.{D.SL21} (p.~186)} }
\expandafter\def\csname SRef@maps:eq:D.SL22\endcsname{\href{S6_SERIES_GRIND_III.pdf\#nameddest=tag-024L}{III.{D.SL22} (p.~186)} }
\expandafter\def\csname SRef@maps:eq:D.SL24\endcsname{\href{S6_SERIES_GRIND_III.pdf\#nameddest=tag-024N}{III.{D.SL24} (p.~186)} }
\expandafter\def\csname SRef@maps:eq:D.SL23\endcsname{\href{S6_SERIES_GRIND_III.pdf\#nameddest=tag-024M}{III.{D.SL23} (p.~187)} }
\expandafter\def\csname SRef@maps:eq:B.5.12\endcsname{\href{S6_SERIES_GRIND_III.pdf\#nameddest=tag-01VA}{III.{B.5.12} (p.~24)} }
\expandafter\def\csname SRef@maps:eq:D.VC1\endcsname{\href{S6_SERIES_GRIND_III.pdf\#nameddest=tag-026E}{III.{D.VC1} (p.~187)} }
\expandafter\def\csname SRef@maps:eq:D.VC7\endcsname{\href{S6_SERIES_GRIND_III.pdf\#nameddest=tag-026L}{III.{D.VC7} (p.~188)} }
\expandafter\def\csname SRef@maps:eq:D.VC2\endcsname{\href{S6_SERIES_GRIND_III.pdf\#nameddest=tag-026G}{III.{D.VC2} (p.~188)} }
\expandafter\def\csname SRef@maps:eq:D.VC3\endcsname{\href{S6_SERIES_GRIND_III.pdf\#nameddest=tag-026H}{III.{D.VC3} (p.~188)} }
\expandafter\def\csname SRef@maps:eq:D.VC4\endcsname{\href{S6_SERIES_GRIND_III.pdf\#nameddest=tag-026I}{III.{D.VC4} (p.~188)} }
\expandafter\def\csname SRef@maps:eq:D.VC5\endcsname{\href{S6_SERIES_GRIND_III.pdf\#nameddest=tag-026J}{III.{D.VC5} (p.~188)} }
\expandafter\def\csname SRef@maps:eq:D.VC8\endcsname{\href{S6_SERIES_GRIND_III.pdf\#nameddest=tag-026M}{III.{D.VC8} (p.~189)} }
\expandafter\def\csname SRef@maps:eq:D.VC6\endcsname{\href{S6_SERIES_GRIND_III.pdf\#nameddest=tag-026K}{III.{D.VC6} (p.~189)} }
\expandafter\def\csname SRef@maps:eq:D.VC9\endcsname{\href{S6_SERIES_GRIND_III.pdf\#nameddest=tag-026N}{III.{D.VC9} (p.~190)} }
\expandafter\def\csname SRef@maps:eq:D.VC10\endcsname{\href{S6_SERIES_GRIND_III.pdf\#nameddest=tag-026F}{III.{D.VC10} (p.~190)} }
\expandafter\def\csname SRef@maps:eq:B.5.13\endcsname{\href{S6_SERIES_GRIND_III.pdf\#nameddest=tag-01VB}{III.{B.5.13} (p.~25)} }
\expandafter\def\csname SRef@maps:eq:D.IC1\endcsname{\href{S6_SERIES_GRIND_III.pdf\#nameddest=tag-01ZY}{III.{D.IC1} (p.~191)} }
\expandafter\def\csname SRef@maps:eq:D.IC2\endcsname{\href{S6_SERIES_GRIND_III.pdf\#nameddest=tag-01ZZ}{III.{D.IC2} (p.~191)} }
\expandafter\def\csname SRef@maps:eq:D.IC3\endcsname{\href{S6_SERIES_GRIND_III.pdf\#nameddest=tag-0200}{III.{D.IC3} (p.~191)} }
\expandafter\def\csname SRef@maps:eq:D.IC4\endcsname{\href{S6_SERIES_GRIND_III.pdf\#nameddest=tag-0201}{III.{D.IC4} (p.~191)} }
\expandafter\def\csname SRef@maps:eq:D.IC5\endcsname{\href{S6_SERIES_GRIND_III.pdf\#nameddest=tag-0202}{III.{D.IC5} (p.~191)} }
\expandafter\def\csname SRef@maps:eq:D.IC6\endcsname{\href{S6_SERIES_GRIND_III.pdf\#nameddest=tag-0203}{III.{D.IC6} (p.~191)} }
\expandafter\def\csname SRef@maps:eq:D.IC7\endcsname{\href{S6_SERIES_GRIND_III.pdf\#nameddest=tag-0204}{III.{D.IC7} (p.~191)} }
\expandafter\def\csname SRef@maps:eq:D.IC8\endcsname{\href{S6_SERIES_GRIND_III.pdf\#nameddest=tag-0205}{III.{D.IC8} (p.~192)} }
\expandafter\def\csname SRef@maps:eq:D.IC9\endcsname{\href{S6_SERIES_GRIND_III.pdf\#nameddest=tag-0206}{III.{D.IC9} (p.~192)} }
\expandafter\def\csname SRef@maps:eq:D.CI1\endcsname{\href{S6_SERIES_GRIND_III.pdf\#nameddest=tag-01YN}{III.{D.CI1} (p.~195)} }
\expandafter\def\csname SRef@maps:eq:B.5.14\endcsname{\href{S6_SERIES_GRIND_III.pdf\#nameddest=tag-01VC}{III.{B.5.14} (p.~25)} }
\expandafter\def\csname SRef@maps:eq:D.OS1\endcsname{\href{S6_SERIES_GRIND_III.pdf\#nameddest=tag-022O}{III.{D.OS1} (p.~196)} }
\expandafter\def\csname SRef@maps:eq:D.OS2\endcsname{\href{S6_SERIES_GRIND_III.pdf\#nameddest=tag-022P}{III.{D.OS2} (p.~196)} }
\expandafter\def\csname SRef@maps:eq:D.OS3\endcsname{\href{S6_SERIES_GRIND_III.pdf\#nameddest=tag-022Q}{III.{D.OS3} (p.~196)} }
\expandafter\def\csname SRef@maps:eq:D.OS4\endcsname{\href{S6_SERIES_GRIND_III.pdf\#nameddest=tag-022R}{III.{D.OS4} (p.~196)} }
\expandafter\def\csname SRef@maps:eq:D.PC1\endcsname{\href{S6_SERIES_GRIND_III.pdf\#nameddest=tag-0233}{III.{D.PC1} (p.~196)} }
\expandafter\def\csname SRef@maps:eq:D.PC2\endcsname{\href{S6_SERIES_GRIND_III.pdf\#nameddest=tag-0239}{III.{D.PC2} (p.~197)} }
\expandafter\def\csname SRef@maps:eq:D.PC3\endcsname{\href{S6_SERIES_GRIND_III.pdf\#nameddest=tag-023A}{III.{D.PC3} (p.~197)} }
\expandafter\def\csname SRef@maps:eq:D.PC4\endcsname{\href{S6_SERIES_GRIND_III.pdf\#nameddest=tag-023B}{III.{D.PC4} (p.~197)} }
\expandafter\def\csname SRef@maps:eq:D.PC5\endcsname{\href{S6_SERIES_GRIND_III.pdf\#nameddest=tag-023C}{III.{D.PC5} (p.~197)} }
\expandafter\def\csname SRef@maps:eq:D.PC6\endcsname{\href{S6_SERIES_GRIND_III.pdf\#nameddest=tag-023D}{III.{D.PC6} (p.~197)} }
\expandafter\def\csname SRef@maps:eq:B.5.15\endcsname{\href{S6_SERIES_GRIND_III.pdf\#nameddest=tag-01VD}{III.{B.5.15} (p.~25)} }
\expandafter\def\csname SRef@maps:eq:D.PC7\endcsname{\href{S6_SERIES_GRIND_III.pdf\#nameddest=tag-023E}{III.{D.PC7} (p.~197)} }
\expandafter\def\csname SRef@maps:eq:D.PC8\endcsname{\href{S6_SERIES_GRIND_III.pdf\#nameddest=tag-023F}{III.{D.PC8} (p.~197)} }
\expandafter\def\csname SRef@maps:eq:D.PC9\endcsname{\href{S6_SERIES_GRIND_III.pdf\#nameddest=tag-023G}{III.{D.PC9} (p.~198)} }
\expandafter\def\csname SRef@maps:eq:D.PC10\endcsname{\href{S6_SERIES_GRIND_III.pdf\#nameddest=tag-0234}{III.{D.PC10} (p.~198)} }
\expandafter\def\csname SRef@maps:eq:D.PC11\endcsname{\href{S6_SERIES_GRIND_III.pdf\#nameddest=tag-0235}{III.{D.PC11} (p.~198)} }
\expandafter\def\csname SRef@maps:eq:D.PC12\endcsname{\href{S6_SERIES_GRIND_III.pdf\#nameddest=tag-0236}{III.{D.PC12} (p.~198)} }
\expandafter\def\csname SRef@maps:eq:D.PC13\endcsname{\href{S6_SERIES_GRIND_III.pdf\#nameddest=tag-0237}{III.{D.PC13} (p.~198)} }
\expandafter\def\csname SRef@maps:eq:D.PC14\endcsname{\href{S6_SERIES_GRIND_III.pdf\#nameddest=tag-0238}{III.{D.PC14} (p.~198)} }
\expandafter\def\csname SRef@maps:eq:D.OS5\endcsname{\href{S6_SERIES_GRIND_III.pdf\#nameddest=tag-022S}{III.{D.OS5} (p.~199)} }
\expandafter\def\csname SRef@maps:eq:B.5.16\endcsname{\href{S6_SERIES_GRIND_III.pdf\#nameddest=tag-01VE}{III.{B.5.16} (p.~25)} }
\expandafter\def\csname SRef@maps:eq:D.OS6\endcsname{\href{S6_SERIES_GRIND_III.pdf\#nameddest=tag-022T}{III.{D.OS6} (p.~199)} }
\expandafter\def\csname SRef@maps:eq:D.OS7\endcsname{\href{S6_SERIES_GRIND_III.pdf\#nameddest=tag-022U}{III.{D.OS7} (p.~199)} }
\expandafter\def\csname SRef@maps:eq:D.OS8\endcsname{\href{S6_SERIES_GRIND_III.pdf\#nameddest=tag-022V}{III.{D.OS8} (p.~200)} }
\expandafter\def\csname SRef@maps:eq:D.OF1\endcsname{\href{S6_SERIES_GRIND_III.pdf\#nameddest=tag-0217}{III.{D.OF1} (p.~200)} }
\expandafter\def\csname SRef@maps:eq:D.OF2\endcsname{\href{S6_SERIES_GRIND_III.pdf\#nameddest=tag-021I}{III.{D.OF2} (p.~200)} }
\expandafter\def\csname SRef@maps:eq:D.OF3\endcsname{\href{S6_SERIES_GRIND_III.pdf\#nameddest=tag-021J}{III.{D.OF3} (p.~201)} }
\expandafter\def\csname SRef@maps:eq:D.OF4\endcsname{\href{S6_SERIES_GRIND_III.pdf\#nameddest=tag-021K}{III.{D.OF4} (p.~201)} }
\expandafter\def\csname SRef@maps:eq:D.OF5\endcsname{\href{S6_SERIES_GRIND_III.pdf\#nameddest=tag-021L}{III.{D.OF5} (p.~201)} }
\expandafter\def\csname SRef@maps:eq:B.5.17\endcsname{\href{S6_SERIES_GRIND_III.pdf\#nameddest=tag-01VF}{III.{B.5.17} (p.~25)} }
\expandafter\def\csname SRef@maps:eq:D.OF6\endcsname{\href{S6_SERIES_GRIND_III.pdf\#nameddest=tag-021M}{III.{D.OF6} (p.~201)} }
\expandafter\def\csname SRef@maps:eq:D.OF7\endcsname{\href{S6_SERIES_GRIND_III.pdf\#nameddest=tag-021N}{III.{D.OF7} (p.~201)} }
\expandafter\def\csname SRef@maps:eq:D.OF8\endcsname{\href{S6_SERIES_GRIND_III.pdf\#nameddest=tag-021O}{III.{D.OF8} (p.~201)} }
\expandafter\def\csname SRef@maps:eq:D.OF9\endcsname{\href{S6_SERIES_GRIND_III.pdf\#nameddest=tag-021P}{III.{D.OF9} (p.~202)} }
\expandafter\def\csname SRef@maps:eq:D.OF10\endcsname{\href{S6_SERIES_GRIND_III.pdf\#nameddest=tag-0218}{III.{D.OF10} (p.~202)} }
\expandafter\def\csname SRef@maps:eq:D.OF11\endcsname{\href{S6_SERIES_GRIND_III.pdf\#nameddest=tag-0219}{III.{D.OF11} (p.~202)} }
\expandafter\def\csname SRef@maps:eq:D.OF12\endcsname{\href{S6_SERIES_GRIND_III.pdf\#nameddest=tag-021A}{III.{D.OF12} (p.~202)} }
\expandafter\def\csname SRef@maps:eq:D.OF13\endcsname{\href{S6_SERIES_GRIND_III.pdf\#nameddest=tag-021B}{III.{D.OF13} (p.~202)} }
\expandafter\def\csname SRef@maps:eq:D.OF14\endcsname{\href{S6_SERIES_GRIND_III.pdf\#nameddest=tag-021C}{III.{D.OF14} (p.~202)} }
\expandafter\def\csname SRef@maps:eq:B.5.18\endcsname{\href{S6_SERIES_GRIND_III.pdf\#nameddest=tag-01VG}{III.{B.5.18} (p.~26)} }
\expandafter\def\csname SRef@maps:eq:D.OF15\endcsname{\href{S6_SERIES_GRIND_III.pdf\#nameddest=tag-021D}{III.{D.OF15} (p.~203)} }
\expandafter\def\csname SRef@maps:eq:D.OF16\endcsname{\href{S6_SERIES_GRIND_III.pdf\#nameddest=tag-021E}{III.{D.OF16} (p.~203)} }
\expandafter\def\csname SRef@maps:eq:D.OF17\endcsname{\href{S6_SERIES_GRIND_III.pdf\#nameddest=tag-021F}{III.{D.OF17} (p.~203)} }
\expandafter\def\csname SRef@maps:eq:D.OF18\endcsname{\href{S6_SERIES_GRIND_III.pdf\#nameddest=tag-021G}{III.{D.OF18} (p.~203)} }
\expandafter\def\csname SRef@maps:eq:D.OL1\endcsname{\href{S6_SERIES_GRIND_III.pdf\#nameddest=tag-021Q}{III.{D.OL1} (p.~204)} }
\expandafter\def\csname SRef@maps:eq:D.OL2\endcsname{\href{S6_SERIES_GRIND_III.pdf\#nameddest=tag-021U}{III.{D.OL2} (p.~204)} }
\expandafter\def\csname SRef@maps:eq:D.OL3\endcsname{\href{S6_SERIES_GRIND_III.pdf\#nameddest=tag-021V}{III.{D.OL3} (p.~204)} }
\expandafter\def\csname SRef@maps:eq:D.OL4\endcsname{\href{S6_SERIES_GRIND_III.pdf\#nameddest=tag-021W}{III.{D.OL4} (p.~204)} }
\expandafter\def\csname SRef@maps:eq:D.OL5\endcsname{\href{S6_SERIES_GRIND_III.pdf\#nameddest=tag-021X}{III.{D.OL5} (p.~204)} }
\expandafter\def\csname SRef@maps:eq:C.1.1\endcsname{\href{S6_SERIES_GRIND_III.pdf\#nameddest=tag-01XE}{III.{C.1.1} (p.~27)} }
\expandafter\def\csname SRef@maps:eq:D.OL6\endcsname{\href{S6_SERIES_GRIND_III.pdf\#nameddest=tag-021Y}{III.{D.OL6} (p.~204)} }
\expandafter\def\csname SRef@maps:eq:D.OL7\endcsname{\href{S6_SERIES_GRIND_III.pdf\#nameddest=tag-021Z}{III.{D.OL7} (p.~204)} }
\expandafter\def\csname SRef@maps:eq:D.OL8\endcsname{\href{S6_SERIES_GRIND_III.pdf\#nameddest=tag-0220}{III.{D.OL8} (p.~205)} }
\expandafter\def\csname SRef@maps:eq:D.OL9\endcsname{\href{S6_SERIES_GRIND_III.pdf\#nameddest=tag-0221}{III.{D.OL9} (p.~205)} }
\expandafter\def\csname SRef@maps:eq:D.OL10\endcsname{\href{S6_SERIES_GRIND_III.pdf\#nameddest=tag-021R}{III.{D.OL10} (p.~205)} }
\expandafter\def\csname SRef@maps:eq:D.OL11\endcsname{\href{S6_SERIES_GRIND_III.pdf\#nameddest=tag-021S}{III.{D.OL11} (p.~205)} }
\expandafter\def\csname SRef@maps:eq:D.OL12\endcsname{\href{S6_SERIES_GRIND_III.pdf\#nameddest=tag-021T}{III.{D.OL12} (p.~205)} }
\expandafter\def\csname SRef@maps:eq:D.OF19\endcsname{\href{S6_SERIES_GRIND_III.pdf\#nameddest=tag-021H}{III.{D.OF19} (p.~206)} }
\expandafter\def\csname SRef@maps:eq:C.1.2\endcsname{\href{S6_SERIES_GRIND_III.pdf\#nameddest=tag-01XF}{III.{C.1.2} (p.~28)} }
\expandafter\def\csname SRef@maps:eq:C.1.3\endcsname{\href{S6_SERIES_GRIND_III.pdf\#nameddest=tag-01XG}{III.{C.1.3} (p.~30)} }
\expandafter\def\csname SRef@maps:eq:E.1\endcsname{\href{S6_SERIES_GRIND_V.pdf\#nameddest=tag-029F}{V.{E.1} (p.~14)} }
\expandafter\def\csname SRef@maps:eq:E.2\endcsname{\href{S6_SERIES_GRIND_V.pdf\#nameddest=tag-029O}{V.{E.2} (p.~15)} }
\expandafter\def\csname SRef@maps:eq:C.1.4\endcsname{\href{S6_SERIES_GRIND_III.pdf\#nameddest=tag-01XH}{III.{C.1.4} (p.~30)} }
\expandafter\def\csname SRef@maps:eq:E.3\endcsname{\href{S6_SERIES_GRIND_V.pdf\#nameddest=tag-029Q}{V.{E.3} (p.~15)} }
\expandafter\def\csname SRef@maps:eq:E.4\endcsname{\href{S6_SERIES_GRIND_V.pdf\#nameddest=tag-029R}{V.{E.4} (p.~15)} }
\expandafter\def\csname SRef@maps:eq:E.5\endcsname{\href{S6_SERIES_GRIND_V.pdf\#nameddest=tag-029S}{V.{E.5} (p.~15)} }
\expandafter\def\csname SRef@maps:eq:E.6\endcsname{\href{S6_SERIES_GRIND_V.pdf\#nameddest=tag-029T}{V.{E.6} (p.~16)} }
\expandafter\def\csname SRef@maps:eq:E.7\endcsname{\href{S6_SERIES_GRIND_V.pdf\#nameddest=tag-029U}{V.{E.7} (p.~16)} }
\expandafter\def\csname SRef@maps:eq:E.8\endcsname{\href{S6_SERIES_GRIND_V.pdf\#nameddest=tag-029V}{V.{E.8} (p.~18)} }
\expandafter\def\csname SRef@maps:eq:E.9\endcsname{\href{S6_SERIES_GRIND_V.pdf\#nameddest=tag-029W}{V.{E.9} (p.~18)} }
\expandafter\def\csname SRef@maps:eq:E.10\endcsname{\href{S6_SERIES_GRIND_V.pdf\#nameddest=tag-029G}{V.{E.10} (p.~18)} }
\expandafter\def\csname SRef@maps:eq:E.11\endcsname{\href{S6_SERIES_GRIND_V.pdf\#nameddest=tag-029H}{V.{E.11} (p.~18)} }
\expandafter\def\csname SRef@maps:eq:E.I1\endcsname{\href{S6_SERIES_GRIND_V.pdf\#nameddest=tag-02B5}{V.{E.I1} (p.~19)} }
\expandafter\def\csname SRef@maps:eq:D.C1\endcsname{\href{S6_SERIES_GRIND_III.pdf\#nameddest=tag-01Y6}{III.{D.C1} (p.~59)} }
\expandafter\def\csname SRef@maps:eq:E.I2\endcsname{\href{S6_SERIES_GRIND_V.pdf\#nameddest=tag-02B6}{V.{E.I2} (p.~20)} }
\expandafter\def\csname SRef@maps:eq:E.I3\endcsname{\href{S6_SERIES_GRIND_V.pdf\#nameddest=tag-02B7}{V.{E.I3} (p.~20)} }
\expandafter\def\csname SRef@maps:eq:E.I4\endcsname{\href{S6_SERIES_GRIND_V.pdf\#nameddest=tag-02B8}{V.{E.I4} (p.~20)} }
\expandafter\def\csname SRef@maps:eq:E.I5\endcsname{\href{S6_SERIES_GRIND_V.pdf\#nameddest=tag-02B9}{V.{E.I5} (p.~21)} }
\expandafter\def\csname SRef@maps:eq:E.12\endcsname{\href{S6_SERIES_GRIND_V.pdf\#nameddest=tag-029I}{V.{E.12} (p.~21)} }
\expandafter\def\csname SRef@maps:eq:E.13\endcsname{\href{S6_SERIES_GRIND_V.pdf\#nameddest=tag-029J}{V.{E.13} (p.~21)} }
\expandafter\def\csname SRef@maps:eq:E.14\endcsname{\href{S6_SERIES_GRIND_V.pdf\#nameddest=tag-029K}{V.{E.14} (p.~22)} }
\expandafter\def\csname SRef@maps:eq:E.16\endcsname{\href{S6_SERIES_GRIND_V.pdf\#nameddest=tag-029M}{V.{E.16} (p.~23)} }
\expandafter\def\csname SRef@maps:eq:E.17\endcsname{\href{S6_SERIES_GRIND_V.pdf\#nameddest=tag-029N}{V.{E.17} (p.~23)} }
\expandafter\def\csname SRef@maps:eq:E.20\endcsname{\href{S6_SERIES_GRIND_V.pdf\#nameddest=tag-029P}{V.{E.20} (p.~25)} }
\expandafter\def\csname SRef@maps:eq:D.R1\endcsname{\href{S6_SERIES_GRIND_III.pdf\#nameddest=tag-023H}{III.{D.R1} (p.~64)} }
\expandafter\def\csname SRef@maps:eq:E.MV1\endcsname{\href{S6_SERIES_GRIND_V.pdf\#nameddest=tag-02BA}{V.{E.MV1} (p.~26)} }
\expandafter\def\csname SRef@maps:eq:E.MV2\endcsname{\href{S6_SERIES_GRIND_V.pdf\#nameddest=tag-02BB}{V.{E.MV2} (p.~27)} }
\expandafter\def\csname SRef@maps:eq:E.MV3\endcsname{\href{S6_SERIES_GRIND_V.pdf\#nameddest=tag-02BC}{V.{E.MV3} (p.~27)} }
\expandafter\def\csname SRef@maps:eq:E.MV4\endcsname{\href{S6_SERIES_GRIND_V.pdf\#nameddest=tag-02BD}{V.{E.MV4} (p.~27)} }
\expandafter\def\csname SRef@maps:eq:E.MV5\endcsname{\href{S6_SERIES_GRIND_V.pdf\#nameddest=tag-02BE}{V.{E.MV5} (p.~27)} }
\expandafter\def\csname SRef@maps:eq:E.MV6\endcsname{\href{S6_SERIES_GRIND_V.pdf\#nameddest=tag-02BF}{V.{E.MV6} (p.~28)} }
\expandafter\def\csname SRef@maps:eq:E.MV7\endcsname{\href{S6_SERIES_GRIND_V.pdf\#nameddest=tag-02BG}{V.{E.MV7} (p.~28)} }
\expandafter\def\csname SRef@maps:eq:E.MV8\endcsname{\href{S6_SERIES_GRIND_V.pdf\#nameddest=tag-02BH}{V.{E.MV8} (p.~28)} }
\expandafter\def\csname SRef@maps:eq:E.MV9\endcsname{\href{S6_SERIES_GRIND_V.pdf\#nameddest=tag-02BI}{V.{E.MV9} (p.~29)} }
\expandafter\def\csname SRef@maps:eq:E.B1\endcsname{\href{S6_SERIES_GRIND_V.pdf\#nameddest=tag-029X}{V.{E.B1} (p.~30)} }
\expandafter\def\csname SRef@maps:eq:D.R2\endcsname{\href{S6_SERIES_GRIND_III.pdf\#nameddest=tag-023I}{III.{D.R2} (p.~64)} }
\expandafter\def\csname SRef@maps:eq:E.B2\endcsname{\href{S6_SERIES_GRIND_V.pdf\#nameddest=tag-029Y}{V.{E.B2} (p.~30)} }
\expandafter\def\csname SRef@maps:eq:E.B3\endcsname{\href{S6_SERIES_GRIND_V.pdf\#nameddest=tag-029Z}{V.{E.B3} (p.~31)} }
\expandafter\def\csname SRef@maps:eq:E.B4\endcsname{\href{S6_SERIES_GRIND_V.pdf\#nameddest=tag-02A0}{V.{E.B4} (p.~31)} }
\expandafter\def\csname SRef@maps:eq:E.B5\endcsname{\href{S6_SERIES_GRIND_V.pdf\#nameddest=tag-02A1}{V.{E.B5} (p.~31)} }
\expandafter\def\csname SRef@maps:eq:E.B6\endcsname{\href{S6_SERIES_GRIND_V.pdf\#nameddest=tag-02A2}{V.{E.B6} (p.~32)} }
\expandafter\def\csname SRef@maps:eq:E.B7\endcsname{\href{S6_SERIES_GRIND_V.pdf\#nameddest=tag-02A3}{V.{E.B7} (p.~32)} }
\expandafter\def\csname SRef@maps:eq:E.S1\endcsname{\href{S6_SERIES_GRIND_V.pdf\#nameddest=tag-02CC}{V.{E.S1} (p.~33)} }
\expandafter\def\csname SRef@maps:eq:E.S2\endcsname{\href{S6_SERIES_GRIND_V.pdf\#nameddest=tag-02CN}{V.{E.S2} (p.~33)} }
\expandafter\def\csname SRef@maps:eq:E.S3\endcsname{\href{S6_SERIES_GRIND_V.pdf\#nameddest=tag-02CY}{V.{E.S3} (p.~34)} }
\expandafter\def\csname SRef@maps:eq:E.S4\endcsname{\href{S6_SERIES_GRIND_V.pdf\#nameddest=tag-02D9}{V.{E.S4} (p.~34)} }
\expandafter\def\csname SRef@maps:eq:D.R3\endcsname{\href{S6_SERIES_GRIND_III.pdf\#nameddest=tag-023J}{III.{D.R3} (p.~64)} }
\expandafter\def\csname SRef@maps:eq:E.S5\endcsname{\href{S6_SERIES_GRIND_V.pdf\#nameddest=tag-02DK}{V.{E.S5} (p.~34)} }
\expandafter\def\csname SRef@maps:eq:E.S12\endcsname{\href{S6_SERIES_GRIND_V.pdf\#nameddest=tag-02CF}{V.{E.S12} (p.~35)} }
\expandafter\def\csname SRef@maps:eq:E.S6\endcsname{\href{S6_SERIES_GRIND_V.pdf\#nameddest=tag-02DR}{V.{E.S6} (p.~35)} }
\expandafter\def\csname SRef@maps:eq:E.S7\endcsname{\href{S6_SERIES_GRIND_V.pdf\#nameddest=tag-02DS}{V.{E.S7} (p.~35)} }
\expandafter\def\csname SRef@maps:eq:E.S8\endcsname{\href{S6_SERIES_GRIND_V.pdf\#nameddest=tag-02DT}{V.{E.S8} (p.~36)} }
\expandafter\def\csname SRef@maps:eq:E.S9\endcsname{\href{S6_SERIES_GRIND_V.pdf\#nameddest=tag-02DU}{V.{E.S9} (p.~36)} }
\expandafter\def\csname SRef@maps:eq:E.S10\endcsname{\href{S6_SERIES_GRIND_V.pdf\#nameddest=tag-02CD}{V.{E.S10} (p.~36)} }
\expandafter\def\csname SRef@maps:eq:E.S11\endcsname{\href{S6_SERIES_GRIND_V.pdf\#nameddest=tag-02CE}{V.{E.S11} (p.~36)} }
\expandafter\def\csname SRef@maps:eq:E.S13\endcsname{\href{S6_SERIES_GRIND_V.pdf\#nameddest=tag-02CG}{V.{E.S13} (p.~37)} }
\expandafter\def\csname SRef@maps:eq:E.S14\endcsname{\href{S6_SERIES_GRIND_V.pdf\#nameddest=tag-02CH}{V.{E.S14} (p.~37)} }
\expandafter\def\csname SRef@maps:eq:D.R4\endcsname{\href{S6_SERIES_GRIND_III.pdf\#nameddest=tag-023K}{III.{D.R4} (p.~64)} }
\expandafter\def\csname SRef@maps:eq:E.S15\endcsname{\href{S6_SERIES_GRIND_V.pdf\#nameddest=tag-02CI}{V.{E.S15} (p.~37)} }
\expandafter\def\csname SRef@maps:eq:E.S16\endcsname{\href{S6_SERIES_GRIND_V.pdf\#nameddest=tag-02CJ}{V.{E.S16} (p.~37)} }
\expandafter\def\csname SRef@maps:eq:E.S20\endcsname{\href{S6_SERIES_GRIND_V.pdf\#nameddest=tag-02CO}{V.{E.S20} (p.~37)} }
\expandafter\def\csname SRef@maps:eq:E.S17\endcsname{\href{S6_SERIES_GRIND_V.pdf\#nameddest=tag-02CK}{V.{E.S17} (p.~37)} }
\expandafter\def\csname SRef@maps:eq:E.S18\endcsname{\href{S6_SERIES_GRIND_V.pdf\#nameddest=tag-02CL}{V.{E.S18} (p.~37)} }
\expandafter\def\csname SRef@maps:eq:E.S19\endcsname{\href{S6_SERIES_GRIND_V.pdf\#nameddest=tag-02CM}{V.{E.S19} (p.~38)} }
\expandafter\def\csname SRef@maps:eq:E.S21\endcsname{\href{S6_SERIES_GRIND_V.pdf\#nameddest=tag-02CP}{V.{E.S21} (p.~39)} }
\expandafter\def\csname SRef@maps:eq:E.S22\endcsname{\href{S6_SERIES_GRIND_V.pdf\#nameddest=tag-02CQ}{V.{E.S22} (p.~40)} }
\expandafter\def\csname SRef@maps:eq:E.S23\endcsname{\href{S6_SERIES_GRIND_V.pdf\#nameddest=tag-02CR}{V.{E.S23} (p.~40)} }
\expandafter\def\csname SRef@maps:eq:E.S24\endcsname{\href{S6_SERIES_GRIND_V.pdf\#nameddest=tag-02CS}{V.{E.S24} (p.~40)} }
\expandafter\def\csname SRef@maps:eq:D.S1\endcsname{\href{S6_SERIES_GRIND_III.pdf\#nameddest=tag-023L}{III.{D.S1} (p.~65)} }
\expandafter\def\csname SRef@maps:eq:E.S25\endcsname{\href{S6_SERIES_GRIND_V.pdf\#nameddest=tag-02CT}{V.{E.S25} (p.~40)} }
\expandafter\def\csname SRef@maps:eq:E.S26\endcsname{\href{S6_SERIES_GRIND_V.pdf\#nameddest=tag-02CU}{V.{E.S26} (p.~40)} }
\expandafter\def\csname SRef@maps:eq:E.S27\endcsname{\href{S6_SERIES_GRIND_V.pdf\#nameddest=tag-02CV}{V.{E.S27} (p.~41)} }
\expandafter\def\csname SRef@maps:eq:E.S28\endcsname{\href{S6_SERIES_GRIND_V.pdf\#nameddest=tag-02CW}{V.{E.S28} (p.~41)} }
\expandafter\def\csname SRef@maps:eq:E.S29\endcsname{\href{S6_SERIES_GRIND_V.pdf\#nameddest=tag-02CX}{V.{E.S29} (p.~42)} }
\expandafter\def\csname SRef@maps:eq:E.S30\endcsname{\href{S6_SERIES_GRIND_V.pdf\#nameddest=tag-02CZ}{V.{E.S30} (p.~42)} }
\expandafter\def\csname SRef@maps:eq:E.S31\endcsname{\href{S6_SERIES_GRIND_V.pdf\#nameddest=tag-02D0}{V.{E.S31} (p.~42)} }
\expandafter\def\csname SRef@maps:eq:E.S32\endcsname{\href{S6_SERIES_GRIND_V.pdf\#nameddest=tag-02D1}{V.{E.S32} (p.~42)} }
\expandafter\def\csname SRef@maps:eq:E.S33\endcsname{\href{S6_SERIES_GRIND_V.pdf\#nameddest=tag-02D2}{V.{E.S33} (p.~42)} }
\expandafter\def\csname SRef@maps:eq:E.S34\endcsname{\href{S6_SERIES_GRIND_V.pdf\#nameddest=tag-02D3}{V.{E.S34} (p.~43)} }
\expandafter\def\csname SRef@maps:eq:D.S2\endcsname{\href{S6_SERIES_GRIND_III.pdf\#nameddest=tag-023R}{III.{D.S2} (p.~65)} }
\expandafter\def\csname SRef@maps:eq:E.S35\endcsname{\href{S6_SERIES_GRIND_V.pdf\#nameddest=tag-02D4}{V.{E.S35} (p.~43)} }
\expandafter\def\csname SRef@maps:eq:E.S36\endcsname{\href{S6_SERIES_GRIND_V.pdf\#nameddest=tag-02D5}{V.{E.S36} (p.~43)} }
\expandafter\def\csname SRef@maps:eq:E.S37\endcsname{\href{S6_SERIES_GRIND_V.pdf\#nameddest=tag-02D6}{V.{E.S37} (p.~43)} }
\expandafter\def\csname SRef@maps:eq:E.S38\endcsname{\href{S6_SERIES_GRIND_V.pdf\#nameddest=tag-02D7}{V.{E.S38} (p.~43)} }
\expandafter\def\csname SRef@maps:eq:E.S39\endcsname{\href{S6_SERIES_GRIND_V.pdf\#nameddest=tag-02D8}{V.{E.S39} (p.~44)} }
\expandafter\def\csname SRef@maps:eq:E.S40\endcsname{\href{S6_SERIES_GRIND_V.pdf\#nameddest=tag-02DA}{V.{E.S40} (p.~44)} }
\expandafter\def\csname SRef@maps:eq:E.S41\endcsname{\href{S6_SERIES_GRIND_V.pdf\#nameddest=tag-02DB}{V.{E.S41} (p.~44)} }
\expandafter\def\csname SRef@maps:eq:E.S42\endcsname{\href{S6_SERIES_GRIND_V.pdf\#nameddest=tag-02DC}{V.{E.S42} (p.~44)} }
\expandafter\def\csname SRef@maps:eq:E.S43\endcsname{\href{S6_SERIES_GRIND_V.pdf\#nameddest=tag-02DD}{V.{E.S43} (p.~45)} }
\expandafter\def\csname SRef@maps:eq:E.S44\endcsname{\href{S6_SERIES_GRIND_V.pdf\#nameddest=tag-02DE}{V.{E.S44} (p.~45)} }
\expandafter\def\csname SRef@maps:eq:D.S3\endcsname{\href{S6_SERIES_GRIND_III.pdf\#nameddest=tag-023S}{III.{D.S3} (p.~65)} }
\expandafter\def\csname SRef@maps:eq:E.S45\endcsname{\href{S6_SERIES_GRIND_V.pdf\#nameddest=tag-02DF}{V.{E.S45} (p.~45)} }
\expandafter\def\csname SRef@maps:eq:E.S46\endcsname{\href{S6_SERIES_GRIND_V.pdf\#nameddest=tag-02DG}{V.{E.S46} (p.~45)} }
\expandafter\def\csname SRef@maps:eq:E.S47\endcsname{\href{S6_SERIES_GRIND_V.pdf\#nameddest=tag-02DH}{V.{E.S47} (p.~45)} }
\expandafter\def\csname SRef@maps:eq:E.S48\endcsname{\href{S6_SERIES_GRIND_V.pdf\#nameddest=tag-02DI}{V.{E.S48} (p.~46)} }
\expandafter\def\csname SRef@maps:eq:E.S49\endcsname{\href{S6_SERIES_GRIND_V.pdf\#nameddest=tag-02DJ}{V.{E.S49} (p.~46)} }
\expandafter\def\csname SRef@maps:eq:E.S50\endcsname{\href{S6_SERIES_GRIND_V.pdf\#nameddest=tag-02DL}{V.{E.S50} (p.~46)} }
\expandafter\def\csname SRef@maps:eq:E.S51\endcsname{\href{S6_SERIES_GRIND_V.pdf\#nameddest=tag-02DM}{V.{E.S51} (p.~46)} }
\expandafter\def\csname SRef@maps:eq:E.S52\endcsname{\href{S6_SERIES_GRIND_V.pdf\#nameddest=tag-02DN}{V.{E.S52} (p.~46)} }
\expandafter\def\csname SRef@maps:eq:E.S53\endcsname{\href{S6_SERIES_GRIND_V.pdf\#nameddest=tag-02DO}{V.{E.S53} (p.~46)} }
\expandafter\def\csname SRef@maps:eq:E.S54\endcsname{\href{S6_SERIES_GRIND_V.pdf\#nameddest=tag-02DP}{V.{E.S54} (p.~46)} }
\expandafter\def\csname SRef@maps:eq:D.S4\endcsname{\href{S6_SERIES_GRIND_III.pdf\#nameddest=tag-023T}{III.{D.S4} (p.~65)} }
\expandafter\def\csname SRef@maps:eq:E.S55\endcsname{\href{S6_SERIES_GRIND_V.pdf\#nameddest=tag-02DQ}{V.{E.S55} (p.~47)} }
\expandafter\def\csname SRef@maps:eq:E.C1\endcsname{\href{S6_SERIES_GRIND_V.pdf\#nameddest=tag-02A4}{V.{E.C1} (p.~48)} }
\expandafter\def\csname SRef@maps:eq:E.C2\endcsname{\href{S6_SERIES_GRIND_V.pdf\#nameddest=tag-02AG}{V.{E.C2} (p.~48)} }
\expandafter\def\csname SRef@maps:eq:E.C3\endcsname{\href{S6_SERIES_GRIND_V.pdf\#nameddest=tag-02AJ}{V.{E.C3} (p.~48)} }
\expandafter\def\csname SRef@maps:eq:E.C4\endcsname{\href{S6_SERIES_GRIND_V.pdf\#nameddest=tag-02AK}{V.{E.C4} (p.~48)} }
\expandafter\def\csname SRef@maps:eq:E.C19a\endcsname{\href{S6_SERIES_GRIND_V.pdf\#nameddest=tag-02AE}{V.{E.C19a} (p.~48)} }
\expandafter\def\csname SRef@maps:eq:E.C19b\endcsname{\href{S6_SERIES_GRIND_V.pdf\#nameddest=tag-02AF}{V.{E.C19b} (p.~48)} }
\expandafter\def\csname SRef@maps:eq:E.C5\endcsname{\href{S6_SERIES_GRIND_V.pdf\#nameddest=tag-02AL}{V.{E.C5} (p.~49)} }
\expandafter\def\csname SRef@maps:eq:E.C6\endcsname{\href{S6_SERIES_GRIND_V.pdf\#nameddest=tag-02AM}{V.{E.C6} (p.~49)} }
\expandafter\def\csname SRef@maps:eq:E.C7\endcsname{\href{S6_SERIES_GRIND_V.pdf\#nameddest=tag-02AN}{V.{E.C7} (p.~49)} }
\expandafter\def\csname SRef@maps:eq:D.S5\endcsname{\href{S6_SERIES_GRIND_III.pdf\#nameddest=tag-023U}{III.{D.S5} (p.~66)} }
\expandafter\def\csname SRef@maps:eq:E.C8\endcsname{\href{S6_SERIES_GRIND_V.pdf\#nameddest=tag-02AO}{V.{E.C8} (p.~49)} }
\expandafter\def\csname SRef@maps:eq:E.C9\endcsname{\href{S6_SERIES_GRIND_V.pdf\#nameddest=tag-02AP}{V.{E.C9} (p.~49)} }
\expandafter\def\csname SRef@maps:eq:E.C10\endcsname{\href{S6_SERIES_GRIND_V.pdf\#nameddest=tag-02A5}{V.{E.C10} (p.~49)} }
\expandafter\def\csname SRef@maps:eq:E.C11\endcsname{\href{S6_SERIES_GRIND_V.pdf\#nameddest=tag-02A6}{V.{E.C11} (p.~49)} }
\expandafter\def\csname SRef@maps:eq:E.C12\endcsname{\href{S6_SERIES_GRIND_V.pdf\#nameddest=tag-02A7}{V.{E.C12} (p.~50)} }
\expandafter\def\csname SRef@maps:eq:E.C20\endcsname{\href{S6_SERIES_GRIND_V.pdf\#nameddest=tag-02AH}{V.{E.C20} (p.~50)} }
\expandafter\def\csname SRef@maps:eq:E.C13\endcsname{\href{S6_SERIES_GRIND_V.pdf\#nameddest=tag-02A8}{V.{E.C13} (p.~51)} }
\expandafter\def\csname SRef@maps:eq:E.C14\endcsname{\href{S6_SERIES_GRIND_V.pdf\#nameddest=tag-02A9}{V.{E.C14} (p.~51)} }
\expandafter\def\csname SRef@maps:eq:E.C15\endcsname{\href{S6_SERIES_GRIND_V.pdf\#nameddest=tag-02AA}{V.{E.C15} (p.~51)} }
\expandafter\def\csname SRef@maps:eq:E.C16\endcsname{\href{S6_SERIES_GRIND_V.pdf\#nameddest=tag-02AB}{V.{E.C16} (p.~51)} }
\expandafter\def\csname SRef@maps:eq:D.S6\endcsname{\href{S6_SERIES_GRIND_III.pdf\#nameddest=tag-023V}{III.{D.S6} (p.~66)} }
\expandafter\def\csname SRef@maps:eq:E.C17\endcsname{\href{S6_SERIES_GRIND_V.pdf\#nameddest=tag-02AC}{V.{E.C17} (p.~51)} }
\expandafter\def\csname SRef@maps:eq:E.C18\endcsname{\href{S6_SERIES_GRIND_V.pdf\#nameddest=tag-02AD}{V.{E.C18} (p.~52)} }
\expandafter\def\csname SRef@maps:eq:E.C21\endcsname{\href{S6_SERIES_GRIND_V.pdf\#nameddest=tag-02AI}{V.{E.C21} (p.~52)} }
\expandafter\def\csname SRef@maps:eq:E.P1\endcsname{\href{S6_SERIES_GRIND_V.pdf\#nameddest=tag-02BJ}{V.{E.P1} (p.~52)} }
\expandafter\def\csname SRef@maps:eq:E.P2\endcsname{\href{S6_SERIES_GRIND_V.pdf\#nameddest=tag-02BS}{V.{E.P2} (p.~53)} }
\expandafter\def\csname SRef@maps:eq:E.P3\endcsname{\href{S6_SERIES_GRIND_V.pdf\#nameddest=tag-02BT}{V.{E.P3} (p.~53)} }
\expandafter\def\csname SRef@maps:eq:E.P4\endcsname{\href{S6_SERIES_GRIND_V.pdf\#nameddest=tag-02BU}{V.{E.P4} (p.~53)} }
\expandafter\def\csname SRef@maps:eq:E.P5\endcsname{\href{S6_SERIES_GRIND_V.pdf\#nameddest=tag-02BV}{V.{E.P5} (p.~53)} }
\expandafter\def\csname SRef@maps:eq:E.P6\endcsname{\href{S6_SERIES_GRIND_V.pdf\#nameddest=tag-02BW}{V.{E.P6} (p.~53)} }
\expandafter\def\csname SRef@maps:eq:E.P7\endcsname{\href{S6_SERIES_GRIND_V.pdf\#nameddest=tag-02BX}{V.{E.P7} (p.~53)} }
\expandafter\def\csname SRef@maps:eq:D.S7\endcsname{\href{S6_SERIES_GRIND_III.pdf\#nameddest=tag-023W}{III.{D.S7} (p.~67)} }
\expandafter\def\csname SRef@maps:eq:E.P8\endcsname{\href{S6_SERIES_GRIND_V.pdf\#nameddest=tag-02BY}{V.{E.P8} (p.~53)} }
\expandafter\def\csname SRef@maps:eq:E.P9\endcsname{\href{S6_SERIES_GRIND_V.pdf\#nameddest=tag-02BZ}{V.{E.P9} (p.~53)} }
\expandafter\def\csname SRef@maps:eq:E.P10\endcsname{\href{S6_SERIES_GRIND_V.pdf\#nameddest=tag-02BK}{V.{E.P10} (p.~54)} }
\expandafter\def\csname SRef@maps:eq:E.P11\endcsname{\href{S6_SERIES_GRIND_V.pdf\#nameddest=tag-02BL}{V.{E.P11} (p.~54)} }
\expandafter\def\csname SRef@maps:eq:E.P12\endcsname{\href{S6_SERIES_GRIND_V.pdf\#nameddest=tag-02BM}{V.{E.P12} (p.~54)} }
\expandafter\def\csname SRef@maps:eq:E.P13\endcsname{\href{S6_SERIES_GRIND_V.pdf\#nameddest=tag-02BN}{V.{E.P13} (p.~54)} }
\expandafter\def\csname SRef@maps:eq:E.P14\endcsname{\href{S6_SERIES_GRIND_V.pdf\#nameddest=tag-02BO}{V.{E.P14} (p.~55)} }
\expandafter\def\csname SRef@maps:eq:E.P15\endcsname{\href{S6_SERIES_GRIND_V.pdf\#nameddest=tag-02BP}{V.{E.P15} (p.~55)} }
\expandafter\def\csname SRef@maps:eq:E.P16\endcsname{\href{S6_SERIES_GRIND_V.pdf\#nameddest=tag-02BQ}{V.{E.P16} (p.~55)} }
\expandafter\def\csname SRef@maps:eq:E.P17\endcsname{\href{S6_SERIES_GRIND_V.pdf\#nameddest=tag-02BR}{V.{E.P17} (p.~56)} }
\expandafter\def\csname SRef@maps:eq:D.S8\endcsname{\href{S6_SERIES_GRIND_III.pdf\#nameddest=tag-023X}{III.{D.S8} (p.~67)} }
\expandafter\def\csname SRef@maps:eq:E.Q1\endcsname{\href{S6_SERIES_GRIND_V.pdf\#nameddest=tag-02C0}{V.{E.Q1} (p.~57)} }
\expandafter\def\csname SRef@maps:eq:E.Q2\endcsname{\href{S6_SERIES_GRIND_V.pdf\#nameddest=tag-02C4}{V.{E.Q2} (p.~57)} }
\expandafter\def\csname SRef@maps:eq:E.Q3\endcsname{\href{S6_SERIES_GRIND_V.pdf\#nameddest=tag-02C5}{V.{E.Q3} (p.~57)} }
\expandafter\def\csname SRef@maps:eq:E.Q4\endcsname{\href{S6_SERIES_GRIND_V.pdf\#nameddest=tag-02C6}{V.{E.Q4} (p.~57)} }
\expandafter\def\csname SRef@maps:eq:E.Q5\endcsname{\href{S6_SERIES_GRIND_V.pdf\#nameddest=tag-02C7}{V.{E.Q5} (p.~58)} }
\expandafter\def\csname SRef@maps:eq:E.Q6\endcsname{\href{S6_SERIES_GRIND_V.pdf\#nameddest=tag-02C8}{V.{E.Q6} (p.~58)} }
\expandafter\def\csname SRef@maps:eq:E.Q7\endcsname{\href{S6_SERIES_GRIND_V.pdf\#nameddest=tag-02C9}{V.{E.Q7} (p.~58)} }
\expandafter\def\csname SRef@maps:eq:E.Q8\endcsname{\href{S6_SERIES_GRIND_V.pdf\#nameddest=tag-02CA}{V.{E.Q8} (p.~58)} }
\expandafter\def\csname SRef@maps:eq:E.Q9\endcsname{\href{S6_SERIES_GRIND_V.pdf\#nameddest=tag-02CB}{V.{E.Q9} (p.~59)} }
\expandafter\def\csname SRef@maps:eq:E.Q10\endcsname{\href{S6_SERIES_GRIND_V.pdf\#nameddest=tag-02C1}{V.{E.Q10} (p.~59)} }
\expandafter\def\csname SRef@maps:eq:D.S9\endcsname{\href{S6_SERIES_GRIND_III.pdf\#nameddest=tag-023Y}{III.{D.S9} (p.~68)} }
\expandafter\def\csname SRef@maps:eq:E.Q11\endcsname{\href{S6_SERIES_GRIND_V.pdf\#nameddest=tag-02C2}{V.{E.Q11} (p.~59)} }
\expandafter\def\csname SRef@maps:eq:E.Q12\endcsname{\href{S6_SERIES_GRIND_V.pdf\#nameddest=tag-02C3}{V.{E.Q12} (p.~60)} }
\expandafter\def\csname SRef@maps:eq:E.T1\endcsname{\href{S6_SERIES_GRIND_V.pdf\#nameddest=tag-02DV}{V.{E.T1} (p.~60)} }
\expandafter\def\csname SRef@maps:eq:E.T2\endcsname{\href{S6_SERIES_GRIND_V.pdf\#nameddest=tag-02DW}{V.{E.T2} (p.~60)} }
\expandafter\def\csname SRef@maps:eq:E.T3\endcsname{\href{S6_SERIES_GRIND_V.pdf\#nameddest=tag-02DX}{V.{E.T3} (p.~61)} }
\expandafter\def\csname SRef@maps:eq:E.T4\endcsname{\href{S6_SERIES_GRIND_V.pdf\#nameddest=tag-02DY}{V.{E.T4} (p.~61)} }
\expandafter\def\csname SRef@maps:eq:E.T5\endcsname{\href{S6_SERIES_GRIND_V.pdf\#nameddest=tag-02DZ}{V.{E.T5} (p.~61)} }
\expandafter\def\csname SRef@maps:eq:E.T6\endcsname{\href{S6_SERIES_GRIND_V.pdf\#nameddest=tag-02E0}{V.{E.T6} (p.~62)} }
\expandafter\def\csname SRef@maps:eq:E.T7\endcsname{\href{S6_SERIES_GRIND_V.pdf\#nameddest=tag-02E1}{V.{E.T7} (p.~62)} }
\expandafter\def\csname SRef@maps:eq:E.T9\endcsname{\href{S6_SERIES_GRIND_V.pdf\#nameddest=tag-02E3}{V.{E.T9} (p.~62)} }
\expandafter\def\csname SRef@maps:eq:D.S10\endcsname{\href{S6_SERIES_GRIND_III.pdf\#nameddest=tag-023M}{III.{D.S10} (p.~68)} }
\expandafter\def\csname SRef@maps:eq:E.T8\endcsname{\href{S6_SERIES_GRIND_V.pdf\#nameddest=tag-02E2}{V.{E.T8} (p.~62)} }
\expandafter\def\csname SRef@maps:eq:E.U1\endcsname{\href{S6_SERIES_GRIND_V.pdf\#nameddest=tag-02E4}{V.{E.U1} (p.~63)} }
\expandafter\def\csname SRef@maps:eq:E.U2\endcsname{\href{S6_SERIES_GRIND_V.pdf\#nameddest=tag-02E7}{V.{E.U2} (p.~63)} }
\expandafter\def\csname SRef@maps:eq:E.U3\endcsname{\href{S6_SERIES_GRIND_V.pdf\#nameddest=tag-02E8}{V.{E.U3} (p.~63)} }
\expandafter\def\csname SRef@maps:eq:E.U4\endcsname{\href{S6_SERIES_GRIND_V.pdf\#nameddest=tag-02E9}{V.{E.U4} (p.~63)} }
\expandafter\def\csname SRef@maps:eq:E.U5\endcsname{\href{S6_SERIES_GRIND_V.pdf\#nameddest=tag-02EA}{V.{E.U5} (p.~63)} }
\expandafter\def\csname SRef@maps:eq:E.U6\endcsname{\href{S6_SERIES_GRIND_V.pdf\#nameddest=tag-02EB}{V.{E.U6} (p.~64)} }
\expandafter\def\csname SRef@maps:eq:E.U7\endcsname{\href{S6_SERIES_GRIND_V.pdf\#nameddest=tag-02EC}{V.{E.U7} (p.~64)} }
\expandafter\def\csname SRef@maps:eq:E.U8\endcsname{\href{S6_SERIES_GRIND_V.pdf\#nameddest=tag-02ED}{V.{E.U8} (p.~64)} }
\expandafter\def\csname SRef@maps:eq:E.U9\endcsname{\href{S6_SERIES_GRIND_V.pdf\#nameddest=tag-02EE}{V.{E.U9} (p.~64)} }
\expandafter\def\csname SRef@maps:eq:D.S11\endcsname{\href{S6_SERIES_GRIND_III.pdf\#nameddest=tag-023N}{III.{D.S11} (p.~69)} }
\expandafter\def\csname SRef@maps:eq:E.U10\endcsname{\href{S6_SERIES_GRIND_V.pdf\#nameddest=tag-02E5}{V.{E.U10} (p.~64)} }
\expandafter\def\csname SRef@maps:eq:E.U11\endcsname{\href{S6_SERIES_GRIND_V.pdf\#nameddest=tag-02E6}{V.{E.U11} (p.~64)} }
\expandafter\def\csname SRef@maps:eq:E.H0\endcsname{\href{S6_SERIES_GRIND_V.pdf\#nameddest=tag-02AQ}{V.{E.H0} (p.~65)} }
\expandafter\def\csname SRef@maps:eq:E.H1\endcsname{\href{S6_SERIES_GRIND_V.pdf\#nameddest=tag-02AR}{V.{E.H1} (p.~65)} }
\expandafter\def\csname SRef@maps:eq:E.H2\endcsname{\href{S6_SERIES_GRIND_V.pdf\#nameddest=tag-02AW}{V.{E.H2} (p.~66)} }
\expandafter\def\csname SRef@maps:eq:E.H3\endcsname{\href{S6_SERIES_GRIND_V.pdf\#nameddest=tag-02AX}{V.{E.H3} (p.~66)} }
\expandafter\def\csname SRef@maps:eq:E.H4\endcsname{\href{S6_SERIES_GRIND_V.pdf\#nameddest=tag-02AY}{V.{E.H4} (p.~66)} }
\expandafter\def\csname SRef@maps:eq:E.H5\endcsname{\href{S6_SERIES_GRIND_V.pdf\#nameddest=tag-02AZ}{V.{E.H5} (p.~66)} }
\expandafter\def\csname SRef@maps:eq:E.H6\endcsname{\href{S6_SERIES_GRIND_V.pdf\#nameddest=tag-02B0}{V.{E.H6} (p.~66)} }
\expandafter\def\csname SRef@maps:eq:E.H7\endcsname{\href{S6_SERIES_GRIND_V.pdf\#nameddest=tag-02B1}{V.{E.H7} (p.~66)} }
\expandafter\def\csname SRef@maps:eq:D.S12\endcsname{\href{S6_SERIES_GRIND_III.pdf\#nameddest=tag-023O}{III.{D.S12} (p.~69)} }
\expandafter\def\csname SRef@maps:eq:E.H8\endcsname{\href{S6_SERIES_GRIND_V.pdf\#nameddest=tag-02B2}{V.{E.H8} (p.~67)} }
\expandafter\def\csname SRef@maps:eq:E.H9\endcsname{\href{S6_SERIES_GRIND_V.pdf\#nameddest=tag-02B3}{V.{E.H9} (p.~67)} }
\expandafter\def\csname SRef@maps:eq:E.H9a\endcsname{\href{S6_SERIES_GRIND_V.pdf\#nameddest=tag-02B4}{V.{E.H9a} (p.~67)} }
\expandafter\def\csname SRef@maps:eq:E.H10\endcsname{\href{S6_SERIES_GRIND_V.pdf\#nameddest=tag-02AS}{V.{E.H10} (p.~68)} }
\expandafter\def\csname SRef@maps:eq:E.H11\endcsname{\href{S6_SERIES_GRIND_V.pdf\#nameddest=tag-02AT}{V.{E.H11} (p.~68)} }
\expandafter\def\csname SRef@maps:eq:E.H12\endcsname{\href{S6_SERIES_GRIND_V.pdf\#nameddest=tag-02AU}{V.{E.H12} (p.~68)} }
\expandafter\def\csname SRef@maps:eq:E.H13\endcsname{\href{S6_SERIES_GRIND_V.pdf\#nameddest=tag-02AV}{V.{E.H13} (p.~69)} }
\expandafter\def\csname SRef@maps:eq:D.S13\endcsname{\href{S6_SERIES_GRIND_III.pdf\#nameddest=tag-023P}{III.{D.S13} (p.~70)} }
\expandafter\def\csname SRef@maps:eq:D.S14\endcsname{\href{S6_SERIES_GRIND_III.pdf\#nameddest=tag-023Q}{III.{D.S14} (p.~70)} }
\expandafter\def\csname SRef@maps:eq:D.I1\endcsname{\href{S6_SERIES_GRIND_III.pdf\#nameddest=tag-01ZU}{III.{D.I1} (p.~71)} }
\expandafter\def\csname SRef@maps:eq:D.I2\endcsname{\href{S6_SERIES_GRIND_III.pdf\#nameddest=tag-01ZV}{III.{D.I2} (p.~71)} }
\expandafter\def\csname SRef@maps:eq:E.15\endcsname{\href{S6_SERIES_GRIND_III.pdf\#nameddest=tag-029L}{III.{E.15} (p.~317)} }
\expandafter\def\csname SRef@maps:eq:D.I3\endcsname{\href{S6_SERIES_GRIND_III.pdf\#nameddest=tag-01ZW}{III.{D.I3} (p.~71)} }
\expandafter\def\csname SRef@maps:eq:D.I4\endcsname{\href{S6_SERIES_GRIND_III.pdf\#nameddest=tag-01ZX}{III.{D.I4} (p.~72)} }
\expandafter\def\csname SRef@maps:eq:D.T1\endcsname{\href{S6_SERIES_GRIND_III.pdf\#nameddest=tag-0258}{III.{D.T1} (p.~73)} }
\expandafter\def\csname SRef@maps:eq:D.T2\endcsname{\href{S6_SERIES_GRIND_III.pdf\#nameddest=tag-025A}{III.{D.T2} (p.~74)} }
\expandafter\def\csname SRef@maps:eq:D.T3\endcsname{\href{S6_SERIES_GRIND_III.pdf\#nameddest=tag-025B}{III.{D.T3} (p.~74)} }
\expandafter\def\csname SRef@maps:eq:D.T4\endcsname{\href{S6_SERIES_GRIND_III.pdf\#nameddest=tag-025L}{III.{D.T4} (p.~74)} }
\expandafter\def\csname SRef@maps:eq:D.T5\endcsname{\href{S6_SERIES_GRIND_III.pdf\#nameddest=tag-025M}{III.{D.T5} (p.~75)} }
\expandafter\def\csname SRef@maps:eq:D.T6\endcsname{\href{S6_SERIES_GRIND_III.pdf\#nameddest=tag-025N}{III.{D.T6} (p.~75)} }
\expandafter\def\csname SRef@maps:eq:D.T7\endcsname{\href{S6_SERIES_GRIND_III.pdf\#nameddest=tag-025O}{III.{D.T7} (p.~75)} }
\expandafter\def\csname SRef@maps:eq:D.T8\endcsname{\href{S6_SERIES_GRIND_III.pdf\#nameddest=tag-025P}{III.{D.T8} (p.~75)} }
\expandafter\def\csname SRef@maps:eq:D.T9\endcsname{\href{S6_SERIES_GRIND_III.pdf\#nameddest=tag-025Q}{III.{D.T9} (p.~76)} }
\expandafter\def\csname SRef@maps:eq:D.T10\endcsname{\href{S6_SERIES_GRIND_III.pdf\#nameddest=tag-0259}{III.{D.T10} (p.~76)} }
\expandafter\def\csname SRef@maps:eq:D.V1\endcsname{\href{S6_SERIES_GRIND_III.pdf\#nameddest=tag-0269}{III.{D.V1} (p.~78)} }
\expandafter\def\csname SRef@maps:eq:D.V2\endcsname{\href{S6_SERIES_GRIND_III.pdf\#nameddest=tag-026A}{III.{D.V2} (p.~78)} }
\expandafter\def\csname SRef@maps:eq:D.V3\endcsname{\href{S6_SERIES_GRIND_III.pdf\#nameddest=tag-026B}{III.{D.V3} (p.~78)} }
\expandafter\def\csname SRef@maps:eq:D.V4\endcsname{\href{S6_SERIES_GRIND_III.pdf\#nameddest=tag-026C}{III.{D.V4} (p.~79)} }
\expandafter\def\csname SRef@maps:eq:D.V5\endcsname{\href{S6_SERIES_GRIND_III.pdf\#nameddest=tag-026D}{III.{D.V5} (p.~79)} }
\expandafter\def\csname SRef@maps:eq:D.F1\endcsname{\href{S6_SERIES_GRIND_III.pdf\#nameddest=tag-01YO}{III.{D.F1} (p.~81)} }
\expandafter\def\csname SRef@maps:eq:D.F2\endcsname{\href{S6_SERIES_GRIND_III.pdf\#nameddest=tag-01YP}{III.{D.F2} (p.~81)} }
\expandafter\def\csname SRef@maps:eq:D.O1\endcsname{\href{S6_SERIES_GRIND_III.pdf\#nameddest=tag-020M}{III.{D.O1} (p.~82)} }
\expandafter\def\csname SRef@maps:eq:D.O2\endcsname{\href{S6_SERIES_GRIND_III.pdf\#nameddest=tag-020N}{III.{D.O2} (p.~82)} }
\expandafter\def\csname SRef@maps:eq:D.O3\endcsname{\href{S6_SERIES_GRIND_III.pdf\#nameddest=tag-020O}{III.{D.O3} (p.~82)} }
\expandafter\def\csname SRef@maps:eq:D.O4\endcsname{\href{S6_SERIES_GRIND_III.pdf\#nameddest=tag-020P}{III.{D.O4} (p.~83)} }
\expandafter\def\csname SRef@maps:eq:D.O5\endcsname{\href{S6_SERIES_GRIND_III.pdf\#nameddest=tag-020Q}{III.{D.O5} (p.~83)} }
\expandafter\def\csname SRef@maps:eq:D.O6\endcsname{\href{S6_SERIES_GRIND_III.pdf\#nameddest=tag-020R}{III.{D.O6} (p.~83)} }
\expandafter\def\csname SRef@maps:eq:D.O7\endcsname{\href{S6_SERIES_GRIND_III.pdf\#nameddest=tag-020S}{III.{D.O7} (p.~84)} }
\expandafter\def\csname SRef@maps:eq:D.O8\endcsname{\href{S6_SERIES_GRIND_III.pdf\#nameddest=tag-020T}{III.{D.O8} (p.~84)} }
\expandafter\def\csname SRef@maps:src:D.2.2\endcsname{\href{S6_SERIES_GRIND_III.pdf\#nameddest=tag-02KU}{III.5.12 (p.~42)} }
\expandafter\def\csname SRef@maps:src:D.2.17\endcsname{\href{S6_SERIES_GRIND_III.pdf\#nameddest=tag-02KR}{III.5.26 (p.~77)} }
\expandafter\def\csname SRef@maps:src:D.2.28\endcsname{\href{S6_SERIES_GRIND_III.pdf\#nameddest=tag-02L3}{III.5.37 (p.~106)} }
\expandafter\def\csname SRef@wb:cusp-cone-S12-sphere-map-null\endcsname{\href{S6_SERIES_GRIND_III.pdf\#nameddest=tag-035Z}{III.7.14 (p.~270)} }
\expandafter\def\csname SRef@wb:cusp-cone-S24-sphere-map-null\endcsname{\href{S6_SERIES_GRIND_III.pdf\#nameddest=tag-036E}{III.7.15 (p.~270)} }
\expandafter\def\csname SRef@wb:cusp-cone-literal-R4-open-cells\endcsname{\href{S6_SERIES_GRIND_III.pdf\#nameddest=tag-036V}{III.7.16 (p.~271)} }
\expandafter\def\csname SRef@wb:cusp-cone-sharp-global-R4-scope\endcsname{\href{S6_SERIES_GRIND_III.pdf\#nameddest=tag-037J}{III.7.20 (p.~278)} }
\expandafter\def\csname SRef@wb:cusp-s12-all-higher-smashes-null\endcsname{\href{S6_SERIES_GRIND_III.pdf\#nameddest=tag-038T}{III.7.22 (p.~280)} }
\expandafter\def\csname SRef@wb:cusp_s12_carrier_homotopy:L567\endcsname{\href{S6_SERIES_GRIND_III.pdf\#nameddest=tag-03A6}{III.7.24 (p.~286)} }
\expandafter\def\csname SRef@wb:sphere-metric-unit-volume\endcsname{\href{S6_SERIES_GRIND_III.pdf\#nameddest=tag-03DH}{III.7.28 (p.~294)} }
\expandafter\def\csname SRef@wb:h24:ret:q2-detected-quotient\endcsname{\href{S6_SERIES_GRIND_III.pdf\#nameddest=tag-03BE}{III.7.33 (p.~306)} }
\expandafter\def\csname SRef@wb:cusp-cone-target-slice-embeddings\endcsname{\href{S6_SERIES_GRIND_III.pdf\#nameddest=tag-037N}{III.7.9 (p.~261)} }
\expandafter\def\csname SRef@wb:cusp-cover-target-action\endcsname{\href{S6_SERIES_GRIND_III.pdf\#nameddest=tag-0388}{III.1 (p.~238)} }
\expandafter\def\csname SRef@wb:cusp-cover-normalization-diagram\endcsname{\href{S6_SERIES_GRIND_III.pdf\#nameddest=tag-0383}{III.10 (p.~239)} }
\expandafter\def\csname SRef@wb:cusp-cone-products-not-spheres\endcsname{\href{S6_SERIES_GRIND_III.pdf\#nameddest=tag-0378}{III.100 (p.~263)} }
\expandafter\def\csname SRef@wb:cusp-cone-join-dimensions\endcsname{\href{S6_SERIES_GRIND_III.pdf\#nameddest=tag-036T}{III.101 (p.~263)} }
\expandafter\def\csname SRef@wb:cusp-cone-kernel-contained\endcsname{\href{S6_SERIES_GRIND_III.pdf\#nameddest=tag-036U}{III.102 (p.~263)} }
\expandafter\def\csname SRef@wb:cusp-cone-finite-index-orbit\endcsname{\href{S6_SERIES_GRIND_III.pdf\#nameddest=tag-036Q}{III.103 (p.~263)} }
\expandafter\def\csname SRef@wb:cusp-cone-fibre-H4-nonzero\endcsname{\href{S6_SERIES_GRIND_III.pdf\#nameddest=tag-036P}{III.104 (p.~264)} }
\expandafter\def\csname SRef@wb:cusp-cone-sphere-to-sphere-maps\endcsname{\href{S6_SERIES_GRIND_III.pdf\#nameddest=tag-037L}{III.105 (p.~266)} }
\expandafter\def\csname SRef@wb:cusp-cone-S12-S4-radial-map\endcsname{\href{S6_SERIES_GRIND_III.pdf\#nameddest=tag-035R}{III.106 (p.~267)} }
\expandafter\def\csname SRef@wb:cusp-cone-S24-S4-radial-map\endcsname{\href{S6_SERIES_GRIND_III.pdf\#nameddest=tag-0361}{III.107 (p.~267)} }
\expandafter\def\csname SRef@wb:cusp-cone-S12-exact-collapsed-locus\endcsname{\href{S6_SERIES_GRIND_III.pdf\#nameddest=tag-035T}{III.108 (p.~267)} }
\expandafter\def\csname SRef@wb:cusp-cone-S24-exact-collapsed-locus\endcsname{\href{S6_SERIES_GRIND_III.pdf\#nameddest=tag-0362}{III.109 (p.~267)} }
\expandafter\def\csname SRef@wb:cusp-cover-monodromy-conjugacy\endcsname{\href{S6_SERIES_GRIND_III.pdf\#nameddest=tag-0381}{III.11 (p.~239)} }
\expandafter\def\csname SRef@wb:cusp-cone-moving-S4-R4-limits\endcsname{\href{S6_SERIES_GRIND_III.pdf\#nameddest=tag-036Y}{III.110 (p.~267)} }
\expandafter\def\csname SRef@wb:cusp-cone-transported-skeleton-collapse\endcsname{\href{S6_SERIES_GRIND_III.pdf\#nameddest=tag-037O}{III.111 (p.~268)} }
\expandafter\def\csname SRef@wb:cusp-cone-distance-to-skeleton\endcsname{\href{S6_SERIES_GRIND_III.pdf\#nameddest=tag-036M}{III.112 (p.~268)} }
\expandafter\def\csname SRef@wb:cusp-cone-S12-onto-exact-ball\endcsname{\href{S6_SERIES_GRIND_III.pdf\#nameddest=tag-035W}{III.113 (p.~269)} }
\expandafter\def\csname SRef@wb:cusp-cone-S24-onto-exact-ball\endcsname{\href{S6_SERIES_GRIND_III.pdf\#nameddest=tag-0369}{III.114 (p.~270)} }
\expandafter\def\csname SRef@wb:cusp-cone-S24-null-maps\endcsname{\href{S6_SERIES_GRIND_III.pdf\#nameddest=tag-0367}{III.115 (p.~270)} }
\expandafter\def\csname SRef@wb:cusp-cone-S24-null-special-cases\endcsname{\href{S6_SERIES_GRIND_III.pdf\#nameddest=tag-0368}{III.116 (p.~270)} }
\expandafter\def\csname SRef@wb:cusp-cone-S24-smash-square-factorization\endcsname{\href{S6_SERIES_GRIND_III.pdf\#nameddest=tag-036C}{III.117 (p.~271)} }
\expandafter\def\csname SRef@wb:cusp-cone-S24-sixfold-suspension\endcsname{\href{S6_SERIES_GRIND_III.pdf\#nameddest=tag-036B}{III.118 (p.~271)} }
\expandafter\def\csname SRef@wb:cusp-cone-fourth-stable-stem\endcsname{\href{S6_SERIES_GRIND_III.pdf\#nameddest=tag-036R}{III.119 (p.~271)} }
\expandafter\def\csname SRef@wb:cusp-cover-maximal-chart-map\endcsname{\href{S6_SERIES_GRIND_III.pdf\#nameddest=tag-0380}{III.12 (p.~241)} }
\expandafter\def\csname SRef@wb:cusp-cone-U12-homeomorphism\endcsname{\href{S6_SERIES_GRIND_III.pdf\#nameddest=tag-036G}{III.121 (p.~271)} }
\expandafter\def\csname SRef@wb:cusp-cone-U24-homeomorphism\endcsname{\href{S6_SERIES_GRIND_III.pdf\#nameddest=tag-036H}{III.122 (p.~271)} }
\expandafter\def\csname SRef@wb:cusp-cone-D12-embedding\endcsname{\href{S6_SERIES_GRIND_III.pdf\#nameddest=tag-035G}{III.123 (p.~273)} }
\expandafter\def\csname SRef@wb:cusp-cone-D24-embedding\endcsname{\href{S6_SERIES_GRIND_III.pdf\#nameddest=tag-035I}{III.124 (p.~273)} }
\expandafter\def\csname SRef@wb:cusp-cone-D12-slice\endcsname{\href{S6_SERIES_GRIND_III.pdf\#nameddest=tag-035H}{III.125 (p.~273)} }
\expandafter\def\csname SRef@wb:cusp-cone-D24-slice\endcsname{\href{S6_SERIES_GRIND_III.pdf\#nameddest=tag-035J}{III.126 (p.~273)} }
\expandafter\def\csname SRef@wb:cusp-cone-carrier-local-homology\endcsname{\href{S6_SERIES_GRIND_III.pdf\#nameddest=tag-036K}{III.127 (p.~273)} }
\expandafter\def\csname SRef@wb:cusp-cone-A-smash-wedge\endcsname{\href{S6_SERIES_GRIND_III.pdf\#nameddest=tag-0357}{III.128 (p.~274)} }
\expandafter\def\csname SRef@wb:cusp-cone-C-smash-wedge\endcsname{\href{S6_SERIES_GRIND_III.pdf\#nameddest=tag-035F}{III.129 (p.~274)} }
\expandafter\def\csname SRef@wb:cusp-cone-E12-fibres\endcsname{\href{S6_SERIES_GRIND_III.pdf\#nameddest=tag-035K}{III.130 (p.~276)} }
\expandafter\def\csname SRef@wb:cusp-cone-E24-fibres\endcsname{\href{S6_SERIES_GRIND_III.pdf\#nameddest=tag-035L}{III.131 (p.~276)} }
\expandafter\def\csname SRef@wb:cusp-s12-smash-factorization\endcsname{\href{S6_SERIES_GRIND_III.pdf\#nameddest=tag-0390}{III.132 (p.~280)} }
\expandafter\def\csname SRef@wb:cusp-s12-toda-smash-formula\endcsname{\href{S6_SERIES_GRIND_III.pdf\#nameddest=tag-0392}{III.133 (p.~280)} }
\expandafter\def\csname SRef@wb:cusp-s12-three-stem-groups\endcsname{\href{S6_SERIES_GRIND_III.pdf\#nameddest=tag-0391}{III.134 (p.~280)} }
\expandafter\def\csname SRef@wb:cusp-s12-sixfold-suspension-class\endcsname{\href{S6_SERIES_GRIND_III.pdf\#nameddest=tag-038Z}{III.135 (p.~281)} }
\expandafter\def\csname SRef@wb:cusp-s12-regular-inverse-image\endcsname{\href{S6_SERIES_GRIND_III.pdf\#nameddest=tag-038Y}{III.136 (p.~281)} }
\expandafter\def\csname SRef@wb:cusp-null:eq:periods\endcsname{\href{S6_SERIES_GRIND_III.pdf\#nameddest=tag-038N}{III.137 (p.~282)} }
\expandafter\def\csname SRef@wb:cusp-null:eq:action\endcsname{\href{S6_SERIES_GRIND_III.pdf\#nameddest=tag-038E}{III.138 (p.~282)} }
\expandafter\def\csname SRef@wb:cusp-null:eq:g\endcsname{\href{S6_SERIES_GRIND_III.pdf\#nameddest=tag-038J}{III.139 (p.~282)} }
\expandafter\def\csname SRef@wb:cusp-null:eq:chart\endcsname{\href{S6_SERIES_GRIND_III.pdf\#nameddest=tag-038H}{III.140 (p.~282)} }
\expandafter\def\csname SRef@wb:cusp-null:eq:moduli\endcsname{\href{S6_SERIES_GRIND_III.pdf\#nameddest=tag-038L}{III.141 (p.~282)} }
\expandafter\def\csname SRef@wb:cusp-null:eq:invertibility\endcsname{\href{S6_SERIES_GRIND_III.pdf\#nameddest=tag-038K}{III.142 (p.~283)} }
\expandafter\def\csname SRef@wb:cusp-null:eq:straightening\endcsname{\href{S6_SERIES_GRIND_III.pdf\#nameddest=tag-038Q}{III.144 (p.~283)} }
\expandafter\def\csname SRef@wb:cusp-null:eq:torus\endcsname{\href{S6_SERIES_GRIND_III.pdf\#nameddest=tag-038R}{III.145 (p.~283)} }
\expandafter\def\csname SRef@wb:cusp-null:eq:spatial-bound\endcsname{\href{S6_SERIES_GRIND_III.pdf\#nameddest=tag-038P}{III.146 (p.~283)} }
\expandafter\def\csname SRef@wb:cusp-null:eq:nearby-moduli\endcsname{\href{S6_SERIES_GRIND_III.pdf\#nameddest=tag-038M}{III.147 (p.~284)} }
\expandafter\def\csname SRef@wb:cusp-null:eq:J\endcsname{\href{S6_SERIES_GRIND_III.pdf\#nameddest=tag-038B}{III.148 (p.~284)} }
\expandafter\def\csname SRef@wb:cusp-null:eq:actual-frame\endcsname{\href{S6_SERIES_GRIND_III.pdf\#nameddest=tag-038F}{III.149 (p.~284)} }
\expandafter\def\csname SRef@wb:cusp-cover-equivariance-computation\endcsname{\href{S6_SERIES_GRIND_III.pdf\#nameddest=tag-037W}{III.15 (p.~241)} }
\expandafter\def\csname SRef@wb:cusp-null:eq:J0\endcsname{\href{S6_SERIES_GRIND_III.pdf\#nameddest=tag-038C}{III.150 (p.~284)} }
\expandafter\def\csname SRef@wb:cusp-null:eq:P\endcsname{\href{S6_SERIES_GRIND_III.pdf\#nameddest=tag-038D}{III.151 (p.~284)} }
\expandafter\def\csname SRef@wb:cusp-null:eq:q\endcsname{\href{S6_SERIES_GRIND_III.pdf\#nameddest=tag-038O}{III.152 (p.~285)} }
\expandafter\def\csname SRef@wb:cusp-null:eq:bordism\endcsname{\href{S6_SERIES_GRIND_III.pdf\#nameddest=tag-038G}{III.153 (p.~285)} }
\expandafter\def\csname SRef@wb:cusp-null:eq:extended-frame\endcsname{\href{S6_SERIES_GRIND_III.pdf\#nameddest=tag-038I}{III.154 (p.~285)} }
\expandafter\def\csname SRef@wb:sphere-metric-normalized-Q-limit\endcsname{\href{S6_SERIES_GRIND_III.pdf\#nameddest=tag-03DA}{III.155 (p.~287)} }
\expandafter\def\csname SRef@wb:sphere-metric-Hausdorff-collapse\endcsname{\href{S6_SERIES_GRIND_III.pdf\#nameddest=tag-03D2}{III.158 (p.~288)} }
\expandafter\def\csname SRef@wb:sphere-metric-A-fixed-metric\endcsname{\href{S6_SERIES_GRIND_III.pdf\#nameddest=tag-03CW}{III.159 (p.~288)} }
\expandafter\def\csname SRef@wb:cusp-cover-basechange-local-ring\endcsname{\href{S6_SERIES_GRIND_III.pdf\#nameddest=tag-037S}{III.16 (p.~242)} }
\expandafter\def\csname SRef@wb:sphere-metric-BC-fixed-metric\endcsname{\href{S6_SERIES_GRIND_III.pdf\#nameddest=tag-03CY}{III.160 (p.~288)} }
\expandafter\def\csname SRef@wb:sphere-metric-A-slice-exact\endcsname{\href{S6_SERIES_GRIND_III.pdf\#nameddest=tag-03CX}{III.161 (p.~289)} }
\expandafter\def\csname SRef@wb:sphere-metric-BC-slice-exact\endcsname{\href{S6_SERIES_GRIND_III.pdf\#nameddest=tag-03CZ}{III.162 (p.~289)} }
\expandafter\def\csname SRef@wb:sphere-metric-slice-fixed-blowups\endcsname{\href{S6_SERIES_GRIND_III.pdf\#nameddest=tag-03DE}{III.164 (p.~289)} }
\expandafter\def\csname SRef@wb:sphere-metric-ambient-fixed-blowups\endcsname{\href{S6_SERIES_GRIND_III.pdf\#nameddest=tag-03D3}{III.166 (p.~289)} }
\expandafter\def\csname SRef@wb:sphere-metric-scaling-trichotomy-result\endcsname{\href{S6_SERIES_GRIND_III.pdf\#nameddest=tag-03DD}{III.167 (p.~291)} }
\expandafter\def\csname SRef@wb:sphere-metric-vertex-distance\endcsname{\href{S6_SERIES_GRIND_III.pdf\#nameddest=tag-03DI}{III.168 (p.~292)} }
\expandafter\def\csname SRef@wb:sphere-metric-intrinsic-ambient-agreement\endcsname{\href{S6_SERIES_GRIND_III.pdf\#nameddest=tag-03D7}{III.169 (p.~292)} }
\expandafter\def\csname SRef@wb:cusp-cover-invariant-ring\endcsname{\href{S6_SERIES_GRIND_III.pdf\#nameddest=tag-037Z}{III.17 (p.~242)} }
\expandafter\def\csname SRef@wb:sphere-metric-intrinsic-ambient-ball-comparison\endcsname{\href{S6_SERIES_GRIND_III.pdf\#nameddest=tag-03D8}{III.170 (p.~292)} }
\expandafter\def\csname SRef@wb:sphere-metric-arbitrary-unrolling-ambient-bubbles\endcsname{\href{S6_SERIES_GRIND_III.pdf\#nameddest=tag-03D4}{III.172 (p.~292)} }
\expandafter\def\csname SRef@wb:sphere-metric-outward-shortcut-exclusion\endcsname{\href{S6_SERIES_GRIND_III.pdf\#nameddest=tag-03DB}{III.173 (p.~293)} }
\expandafter\def\csname SRef@wb:sphere-metric-C-smash-local-homology\endcsname{\href{S6_SERIES_GRIND_III.pdf\#nameddest=tag-03D1}{III.174 (p.~295)} }
\expandafter\def\csname SRef@wb:sphere-metric-C-boundary-quotient-homeomorphism\endcsname{\href{S6_SERIES_GRIND_III.pdf\#nameddest=tag-03D0}{III.175 (p.~296)} }
\expandafter\def\csname SRef@wb:sphere-metric-boundary-link-local-shift\endcsname{\href{S6_SERIES_GRIND_III.pdf\#nameddest=tag-03D5}{III.176 (p.~296)} }
\expandafter\def\csname SRef@wb:h24:ret:eq:original-periods\endcsname{\href{S6_SERIES_GRIND_III.pdf\#nameddest=tag-03B4}{III.177 (p.~298)} }
\expandafter\def\csname SRef@wb:h24:ret:eq:A\endcsname{\href{S6_SERIES_GRIND_III.pdf\#nameddest=tag-03AF}{III.178 (p.~298)} }
\expandafter\def\csname SRef@wb:h24:ret:eq:deck\endcsname{\href{S6_SERIES_GRIND_III.pdf\#nameddest=tag-03AS}{III.179 (p.~298)} }
\expandafter\def\csname SRef@wb:h24:ret:eq:logx\endcsname{\href{S6_SERIES_GRIND_III.pdf\#nameddest=tag-03B0}{III.180 (p.~298)} }
\expandafter\def\csname SRef@wb:h24:ret:eq:invertibility\endcsname{\href{S6_SERIES_GRIND_III.pdf\#nameddest=tag-03AZ}{III.181 (p.~298)} }
\expandafter\def\csname SRef@wb:h24:ret:eq:native\endcsname{\href{S6_SERIES_GRIND_III.pdf\#nameddest=tag-03B1}{III.182 (p.~298)} }
\expandafter\def\csname SRef@wb:h24:ret:eq:two-skeleton\endcsname{\href{S6_SERIES_GRIND_III.pdf\#nameddest=tag-03B9}{III.183 (p.~298)} }
\expandafter\def\csname SRef@wb:h24:ret:eq:A-target\endcsname{\href{S6_SERIES_GRIND_III.pdf\#nameddest=tag-03AH}{III.184 (p.~299)} }
\expandafter\def\csname SRef@wb:h24:ret:eq:q2\endcsname{\href{S6_SERIES_GRIND_III.pdf\#nameddest=tag-03B5}{III.185 (p.~299)} }
\expandafter\def\csname SRef@wb:h24:ret:eq:H\endcsname{\href{S6_SERIES_GRIND_III.pdf\#nameddest=tag-03AN}{III.186 (p.~299)} }
\expandafter\def\csname SRef@wb:h24:ret:eq:four-cell\endcsname{\href{S6_SERIES_GRIND_III.pdf\#nameddest=tag-03AW}{III.187 (p.~299)} }
\expandafter\def\csname SRef@wb:h24:ret:eq:chosen-value\endcsname{\href{S6_SERIES_GRIND_III.pdf\#nameddest=tag-03AR}{III.188 (p.~299)} }
\expandafter\def\csname SRef@wb:h24:ret:eq:z-chart\endcsname{\href{S6_SERIES_GRIND_III.pdf\#nameddest=tag-03BB}{III.189 (p.~299)} }
\expandafter\def\csname SRef@wb:cusp-cover-triple-fibre-length\endcsname{\href{S6_SERIES_GRIND_III.pdf\#nameddest=tag-038A}{III.19 (p.~242)} }
\expandafter\def\csname SRef@wb:h24:ret:eq:units\endcsname{\href{S6_SERIES_GRIND_III.pdf\#nameddest=tag-03BA}{III.190 (p.~299)} }
\expandafter\def\csname SRef@wb:h24:ret:eq:F\endcsname{\href{S6_SERIES_GRIND_III.pdf\#nameddest=tag-03AL}{III.191 (p.~299)} }
\expandafter\def\csname SRef@wb:h24:ret:eq:Finv\endcsname{\href{S6_SERIES_GRIND_III.pdf\#nameddest=tag-03AM}{III.192 (p.~299)} }
\expandafter\def\csname SRef@wb:h24:ret:eq:M\endcsname{\href{S6_SERIES_GRIND_III.pdf\#nameddest=tag-03AP}{III.193 (p.~300)} }
\expandafter\def\csname SRef@wb:h24:ret:eq:exact-u\endcsname{\href{S6_SERIES_GRIND_III.pdf\#nameddest=tag-03AV}{III.194 (p.~300)} }
\expandafter\def\csname SRef@wb:h24:ret:eq:full-derivative\endcsname{\href{S6_SERIES_GRIND_III.pdf\#nameddest=tag-03AX}{III.195 (p.~300)} }
\expandafter\def\csname SRef@wb:h24:ret:eq:derivative-on-M\endcsname{\href{S6_SERIES_GRIND_III.pdf\#nameddest=tag-03AT}{III.196 (p.~300)} }
\expandafter\def\csname SRef@wb:h24:ret:eq:LQ\endcsname{\href{S6_SERIES_GRIND_III.pdf\#nameddest=tag-03AO}{III.197 (p.~301)} }
\expandafter\def\csname SRef@wb:h24:ret:eq:tangents\endcsname{\href{S6_SERIES_GRIND_III.pdf\#nameddest=tag-03B8}{III.198 (p.~301)} }
\expandafter\def\csname SRef@wb:h24:ret:eq:normal-first\endcsname{\href{S6_SERIES_GRIND_III.pdf\#nameddest=tag-03B2}{III.199 (p.~301)} }
\expandafter\def\csname SRef@wb:cusp-cover-base-change\endcsname{\href{S6_SERIES_GRIND_III.pdf\#nameddest=tag-037Q}{III.2 (p.~238)} }
\expandafter\def\csname SRef@wb:cusp-cover-double-fibre-length\endcsname{\href{S6_SERIES_GRIND_III.pdf\#nameddest=tag-037V}{III.20 (p.~242)} }
\expandafter\def\csname SRef@wb:h24:ret:eq:normal-u\endcsname{\href{S6_SERIES_GRIND_III.pdf\#nameddest=tag-03B3}{III.200 (p.~301)} }
\expandafter\def\csname SRef@wb:h24:ret:eq:differential-table\endcsname{\href{S6_SERIES_GRIND_III.pdf\#nameddest=tag-03AU}{III.201 (p.~301)} }
\expandafter\def\csname SRef@wb:h24:ret:eq:A-homotopy\endcsname{\href{S6_SERIES_GRIND_III.pdf\#nameddest=tag-03AG}{III.202 (p.~301)} }
\expandafter\def\csname SRef@wb:h24:ret:eq:sign\endcsname{\href{S6_SERIES_GRIND_III.pdf\#nameddest=tag-03B7}{III.203 (p.~302)} }
\expandafter\def\csname SRef@wb:h24:ret:eq:D\endcsname{\href{S6_SERIES_GRIND_III.pdf\#nameddest=tag-03AJ}{III.204 (p.~302)} }
\expandafter\def\csname SRef@wb:h24:ret:eq:full-frame\endcsname{\href{S6_SERIES_GRIND_III.pdf\#nameddest=tag-03AY}{III.205 (p.~302)} }
\expandafter\def\csname SRef@wb:h24:ret:eq:D-homotopy\endcsname{\href{S6_SERIES_GRIND_III.pdf\#nameddest=tag-03AK}{III.206 (p.~302)} }
\expandafter\def\csname SRef@wb:h24:ret:eq:quadratic\endcsname{\href{S6_SERIES_GRIND_III.pdf\#nameddest=tag-03B6}{III.207 (p.~303)} }
\expandafter\def\csname SRef@wb:h24:ret:eq:Arf\endcsname{\href{S6_SERIES_GRIND_III.pdf\#nameddest=tag-03AI}{III.208 (p.~303)} }
\expandafter\def\csname SRef@wb:h24:ret:eq:actual-Arf\endcsname{\href{S6_SERIES_GRIND_III.pdf\#nameddest=tag-03AQ}{III.209 (p.~304)} }
\expandafter\def\csname SRef@wb:cusp-cover-additive-fibre-map\endcsname{\href{S6_SERIES_GRIND_III.pdf\#nameddest=tag-037P}{III.21 (p.~243)} }
\expandafter\def\csname SRef@mapsD:repair:cocycle-relation\endcsname{\href{S6_SERIES_GRIND_III.pdf\#nameddest=tag-02OF}{III.210 (p.~310)} }
\expandafter\def\csname SRef@mapsD:repair:canonical-matrices\endcsname{\href{S6_SERIES_GRIND_III.pdf\#nameddest=tag-02OE}{III.211 (p.~312)} }
\expandafter\def\csname SRef@mapsD:repair:free-peripherals\endcsname{\href{S6_SERIES_GRIND_III.pdf\#nameddest=tag-02OH}{III.212 (p.~312)} }
\expandafter\def\csname SRef@mapsD:repair:50-weierstrass\endcsname{\href{S6_SERIES_GRIND_III.pdf\#nameddest=tag-02OD}{III.213 (p.~312)} }
\expandafter\def\csname SRef@mapsD:repair:height-formula\endcsname{\href{S6_SERIES_GRIND_III.pdf\#nameddest=tag-02OI}{III.214 (p.~315)} }
\expandafter\def\csname SRef@wb:cusp-spec-radial-notation\endcsname{\href{S6_SERIES_GRIND_III.pdf\#nameddest=tag-03A3}{III.22 (p.~244)} }
\expandafter\def\csname SRef@wb:cusp-spec-holomorphic-errors\endcsname{\href{S6_SERIES_GRIND_III.pdf\#nameddest=tag-039V}{III.23 (p.~244)} }
\expandafter\def\csname SRef@wb:cusp-spec-full-inverse-transpose\endcsname{\href{S6_SERIES_GRIND_III.pdf\#nameddest=tag-039O}{III.24 (p.~244)} }
\expandafter\def\csname SRef@wb:cusp-spec-global-circle-map\endcsname{\href{S6_SERIES_GRIND_III.pdf\#nameddest=tag-039R}{III.25 (p.~246)} }
\expandafter\def\csname SRef@wb:cusp-spec-global-circle-quotient\endcsname{\href{S6_SERIES_GRIND_III.pdf\#nameddest=tag-039S}{III.26 (p.~246)} }
\expandafter\def\csname SRef@wb:cusp-spec-global-circle-coordinate\endcsname{\href{S6_SERIES_GRIND_III.pdf\#nameddest=tag-039Q}{III.27 (p.~246)} }
\expandafter\def\csname SRef@wb:cusp-spec-global-kernel-family\endcsname{\href{S6_SERIES_GRIND_III.pdf\#nameddest=tag-039T}{III.28 (p.~246)} }
\expandafter\def\csname SRef@wb:cusp-spec-radial-circle-character\endcsname{\href{S6_SERIES_GRIND_III.pdf\#nameddest=tag-03A2}{III.29 (p.~247)} }
\expandafter\def\csname SRef@wb:cusp-cover-source-factor\endcsname{\href{S6_SERIES_GRIND_III.pdf\#nameddest=tag-0387}{III.3 (p.~238)} }
\expandafter\def\csname SRef@wb:cusp-spec-global-circle-circumference\endcsname{\href{S6_SERIES_GRIND_III.pdf\#nameddest=tag-039P}{III.30 (p.~247)} }
\expandafter\def\csname SRef@wb:cusp-spec-fourfold-fibre-product\endcsname{\href{S6_SERIES_GRIND_III.pdf\#nameddest=tag-039N}{III.31 (p.~247)} }
\expandafter\def\csname SRef@wb:cusp-spec-fourfold-circle-map\endcsname{\href{S6_SERIES_GRIND_III.pdf\#nameddest=tag-039M}{III.32 (p.~247)} }
\expandafter\def\csname SRef@wb:cusp-spec-four-circle-fibre\endcsname{\href{S6_SERIES_GRIND_III.pdf\#nameddest=tag-039J}{III.33 (p.~247)} }
\expandafter\def\csname SRef@wb:cusp-spec-four-circle-limit\endcsname{\href{S6_SERIES_GRIND_III.pdf\#nameddest=tag-039K}{III.34 (p.~247)} }
\expandafter\def\csname SRef@wb:cusp-spec-doubled-local-quotient\endcsname{\href{S6_SERIES_GRIND_III.pdf\#nameddest=tag-0398}{III.35 (p.~247)} }
\expandafter\def\csname SRef@wb:cusp-spec-finite-affine-action\endcsname{\href{S6_SERIES_GRIND_III.pdf\#nameddest=tag-039A}{III.36 (p.~248)} }
\expandafter\def\csname SRef@wb:cusp-spec-finite-character-shift\endcsname{\href{S6_SERIES_GRIND_III.pdf\#nameddest=tag-039B}{III.37 (p.~248)} }
\expandafter\def\csname SRef@wb:cusp-spec-finite-Seifert-target\endcsname{\href{S6_SERIES_GRIND_III.pdf\#nameddest=tag-0399}{III.38 (p.~248)} }
\expandafter\def\csname SRef@wb:cusp-spec-finite-fourfold-normalization\endcsname{\href{S6_SERIES_GRIND_III.pdf\#nameddest=tag-039D}{III.39 (p.~248)} }
\expandafter\def\csname SRef@wb:cusp-cover-source-action\endcsname{\href{S6_SERIES_GRIND_III.pdf\#nameddest=tag-0386}{III.4 (p.~238)} }
\expandafter\def\csname SRef@wb:cusp-spec-finite-fourfold-target\endcsname{\href{S6_SERIES_GRIND_III.pdf\#nameddest=tag-039E}{III.40 (p.~248)} }
\expandafter\def\csname SRef@wb:cusp-spec-finite-normalization-map\endcsname{\href{S6_SERIES_GRIND_III.pdf\#nameddest=tag-039F}{III.41 (p.~248)} }
\expandafter\def\csname SRef@wb:cusp-spec-product-not-sphere\endcsname{\href{S6_SERIES_GRIND_III.pdf\#nameddest=tag-03A1}{III.42 (p.~248)} }
\expandafter\def\csname SRef@wb:cusp-spec-four-product-not-sphere\endcsname{\href{S6_SERIES_GRIND_III.pdf\#nameddest=tag-039L}{III.43 (p.~248)} }
\expandafter\def\csname SRef@wb:cusp-spec-smash-spheres\endcsname{\href{S6_SERIES_GRIND_III.pdf\#nameddest=tag-03A4}{III.44 (p.~249)} }
\expandafter\def\csname SRef@wb:cusp-spec-target-smash-map\endcsname{\href{S6_SERIES_GRIND_III.pdf\#nameddest=tag-03A5}{III.45 (p.~249)} }
\expandafter\def\csname SRef@wb:cusp-spec-cusp-kernel\endcsname{\href{S6_SERIES_GRIND_III.pdf\#nameddest=tag-0394}{III.46 (p.~249)} }
\expandafter\def\csname SRef@wb:cusp-spec-invariant-flag-calculation\endcsname{\href{S6_SERIES_GRIND_III.pdf\#nameddest=tag-039X}{III.47 (p.~249)} }
\expandafter\def\csname SRef@wb:cusp-spec-invariant-character-lattice\endcsname{\href{S6_SERIES_GRIND_III.pdf\#nameddest=tag-039W}{III.48 (p.~249)} }
\expandafter\def\csname SRef@wb:cusp-spec-circle-dual-vector\endcsname{\href{S6_SERIES_GRIND_III.pdf\#nameddest=tag-0393}{III.49 (p.~250)} }
\expandafter\def\csname SRef@wb:cusp-cover-toric-power-map\endcsname{\href{S6_SERIES_GRIND_III.pdf\#nameddest=tag-0389}{III.5 (p.~239)} }
\expandafter\def\csname SRef@wb:cusp-spec-no-finite-cover-matrix\endcsname{\href{S6_SERIES_GRIND_III.pdf\#nameddest=tag-039Z}{III.50 (p.~251)} }
\expandafter\def\csname SRef@wb:cusp-spec-no-finite-cover-orbit\endcsname{\href{S6_SERIES_GRIND_III.pdf\#nameddest=tag-03A0}{III.51 (p.~251)} }
\expandafter\def\csname SRef@wb:cusp-spec-finite-target-conjugacy\endcsname{\href{S6_SERIES_GRIND_III.pdf\#nameddest=tag-039I}{III.52 (p.~251)} }
\expandafter\def\csname SRef@wb:cusp-spec-finite-target-cocycle\endcsname{\href{S6_SERIES_GRIND_III.pdf\#nameddest=tag-039H}{III.53 (p.~251)} }
\expandafter\def\csname SRef@wb:cusp-spec-finite-overlap-identity\endcsname{\href{S6_SERIES_GRIND_III.pdf\#nameddest=tag-039G}{III.54 (p.~251)} }
\expandafter\def\csname SRef@wb:cusp-spec-finite-fibre-product-branches\endcsname{\href{S6_SERIES_GRIND_III.pdf\#nameddest=tag-039C}{III.55 (p.~252)} }
\expandafter\def\csname SRef@wb:cusp-spec-cusp-nonextension-position\endcsname{\href{S6_SERIES_GRIND_III.pdf\#nameddest=tag-0397}{III.56 (p.~252)} }
\expandafter\def\csname SRef@wb:cusp-spec-cusp-nonextension-character\endcsname{\href{S6_SERIES_GRIND_III.pdf\#nameddest=tag-0395}{III.57 (p.~252)} }
\expandafter\def\csname SRef@wb:cusp-spec-cusp-nonextension-chart\endcsname{\href{S6_SERIES_GRIND_III.pdf\#nameddest=tag-0396}{III.58 (p.~253)} }
\expandafter\def\csname SRef@wb:cusp-cone-disc\endcsname{\href{S6_SERIES_GRIND_III.pdf\#nameddest=tag-036L}{III.59 (p.~255)} }
\expandafter\def\csname SRef@wb:cusp-cover-completed-map\endcsname{\href{S6_SERIES_GRIND_III.pdf\#nameddest=tag-037T}{III.6 (p.~239)} }
\expandafter\def\csname SRef@wb:cusp-cone-quotient-lattice\endcsname{\href{S6_SERIES_GRIND_III.pdf\#nameddest=tag-037F}{III.60 (p.~255)} }
\expandafter\def\csname SRef@wb:cusp-cone-native-family-map\endcsname{\href{S6_SERIES_GRIND_III.pdf\#nameddest=tag-0371}{III.61 (p.~255)} }
\expandafter\def\csname SRef@wb:cusp-cone-radial-isometry\endcsname{\href{S6_SERIES_GRIND_III.pdf\#nameddest=tag-037G}{III.62 (p.~255)} }
\expandafter\def\csname SRef@wb:cusp-cone-native-restriction\endcsname{\href{S6_SERIES_GRIND_III.pdf\#nameddest=tag-0372}{III.63 (p.~255)} }
\expandafter\def\csname SRef@wb:cusp-cone-native-systole\endcsname{\href{S6_SERIES_GRIND_III.pdf\#nameddest=tag-0373}{III.64 (p.~255)} }
\expandafter\def\csname SRef@wb:cusp-cone-monodromy-image\endcsname{\href{S6_SERIES_GRIND_III.pdf\#nameddest=tag-036X}{III.65 (p.~255)} }
\expandafter\def\csname SRef@wb:cusp-cone-A-definition\endcsname{\href{S6_SERIES_GRIND_III.pdf\#nameddest=tag-0354}{III.66 (p.~256)} }
\expandafter\def\csname SRef@wb:cusp-cone-A-suspension\endcsname{\href{S6_SERIES_GRIND_III.pdf\#nameddest=tag-0358}{III.67 (p.~256)} }
\expandafter\def\csname SRef@wb:cusp-cone-q2\endcsname{\href{S6_SERIES_GRIND_III.pdf\#nameddest=tag-037D}{III.68 (p.~256)} }
\expandafter\def\csname SRef@wb:cusp-cone-q2-formula\endcsname{\href{S6_SERIES_GRIND_III.pdf\#nameddest=tag-037E}{III.69 (p.~256)} }
\expandafter\def\csname SRef@wb:cusp-cover-direct-ramification\endcsname{\href{S6_SERIES_GRIND_III.pdf\#nameddest=tag-037U}{III.7 (p.~239)} }
\expandafter\def\csname SRef@wb:cusp-cone-S12-map\endcsname{\href{S6_SERIES_GRIND_III.pdf\#nameddest=tag-035V}{III.70 (p.~256)} }
\expandafter\def\csname SRef@wb:cusp-cone-S12-source\endcsname{\href{S6_SERIES_GRIND_III.pdf\#nameddest=tag-035Y}{III.71 (p.~256)} }
\expandafter\def\csname SRef@wb:cusp-cone-S12-fibre-map\endcsname{\href{S6_SERIES_GRIND_III.pdf\#nameddest=tag-035U}{III.72 (p.~256)} }
\expandafter\def\csname SRef@wb:cusp-cone-S12-R4-limit\endcsname{\href{S6_SERIES_GRIND_III.pdf\#nameddest=tag-035Q}{III.73 (p.~256)} }
\expandafter\def\csname SRef@wb:cusp-cone-A-homology\endcsname{\href{S6_SERIES_GRIND_III.pdf\#nameddest=tag-0355}{III.74 (p.~256)} }
\expandafter\def\csname SRef@wb:cusp-cone-A-smash-homology\endcsname{\href{S6_SERIES_GRIND_III.pdf\#nameddest=tag-0356}{III.75 (p.~257)} }
\expandafter\def\csname SRef@wb:cusp-cone-B-definition\endcsname{\href{S6_SERIES_GRIND_III.pdf\#nameddest=tag-0359}{III.76 (p.~258)} }
\expandafter\def\csname SRef@wb:cusp-cone-q1-local\endcsname{\href{S6_SERIES_GRIND_III.pdf\#nameddest=tag-037B}{III.77 (p.~258)} }
\expandafter\def\csname SRef@wb:cusp-cone-q1-local-formula\endcsname{\href{S6_SERIES_GRIND_III.pdf\#nameddest=tag-037C}{III.78 (p.~258)} }
\expandafter\def\csname SRef@wb:cusp-cone-S24-local-fibre-map\endcsname{\href{S6_SERIES_GRIND_III.pdf\#nameddest=tag-0364}{III.79 (p.~258)} }
\expandafter\def\csname SRef@wb:cusp-cover-base-changed-target\endcsname{\href{S6_SERIES_GRIND_III.pdf\#nameddest=tag-037R}{III.8 (p.~239)} }
\expandafter\def\csname SRef@wb:cusp-cone-B-homology\endcsname{\href{S6_SERIES_GRIND_III.pdf\#nameddest=tag-035A}{III.80 (p.~258)} }
\expandafter\def\csname SRef@wb:cusp-cone-B-smash-homology\endcsname{\href{S6_SERIES_GRIND_III.pdf\#nameddest=tag-035B}{III.81 (p.~258)} }
\expandafter\def\csname SRef@wb:cusp-cone-pair-of-pants\endcsname{\href{S6_SERIES_GRIND_III.pdf\#nameddest=tag-0377}{III.82 (p.~259)} }
\expandafter\def\csname SRef@wb:cusp-cone-M-definition\endcsname{\href{S6_SERIES_GRIND_III.pdf\#nameddest=tag-035P}{III.83 (p.~259)} }
\expandafter\def\csname SRef@wb:cusp-cone-M-boundary\endcsname{\href{S6_SERIES_GRIND_III.pdf\#nameddest=tag-035N}{III.84 (p.~259)} }
\expandafter\def\csname SRef@wb:cusp-cone-C-definition\endcsname{\href{S6_SERIES_GRIND_III.pdf\#nameddest=tag-035C}{III.85 (p.~259)} }
\expandafter\def\csname SRef@wb:cusp-cone-q1\endcsname{\href{S6_SERIES_GRIND_III.pdf\#nameddest=tag-0379}{III.86 (p.~259)} }
\expandafter\def\csname SRef@wb:cusp-cone-q1-cusp-formula\endcsname{\href{S6_SERIES_GRIND_III.pdf\#nameddest=tag-037A}{III.87 (p.~259)} }
\expandafter\def\csname SRef@wb:cusp-cone-S24-map\endcsname{\href{S6_SERIES_GRIND_III.pdf\#nameddest=tag-0366}{III.88 (p.~260)} }
\expandafter\def\csname SRef@wb:cusp-cone-S24-source\endcsname{\href{S6_SERIES_GRIND_III.pdf\#nameddest=tag-036D}{III.89 (p.~260)} }
\expandafter\def\csname SRef@wb:cusp-cover-relative-factorization\endcsname{\href{S6_SERIES_GRIND_III.pdf\#nameddest=tag-0384}{III.9 (p.~239)} }
\expandafter\def\csname SRef@wb:cusp-cone-S24-fibre-map\endcsname{\href{S6_SERIES_GRIND_III.pdf\#nameddest=tag-0363}{III.90 (p.~260)} }
\expandafter\def\csname SRef@wb:cusp-cone-S24-R4-limit\endcsname{\href{S6_SERIES_GRIND_III.pdf\#nameddest=tag-0360}{III.91 (p.~260)} }
\expandafter\def\csname SRef@wb:cusp-cone-C-homology\endcsname{\href{S6_SERIES_GRIND_III.pdf\#nameddest=tag-035D}{III.92 (p.~260)} }
\expandafter\def\csname SRef@wb:cusp-cone-C-smash-homology\endcsname{\href{S6_SERIES_GRIND_III.pdf\#nameddest=tag-035E}{III.93 (p.~260)} }
\expandafter\def\csname SRef@wb:cusp-cone-Seifert-meridians\endcsname{\href{S6_SERIES_GRIND_III.pdf\#nameddest=tag-036F}{III.94 (p.~261)} }
\expandafter\def\csname SRef@wb:cusp-cone-M-H1\endcsname{\href{S6_SERIES_GRIND_III.pdf\#nameddest=tag-035M}{III.95 (p.~261)} }
\expandafter\def\csname SRef@wb:cusp-cone-M-boundary-map\endcsname{\href{S6_SERIES_GRIND_III.pdf\#nameddest=tag-035O}{III.96 (p.~261)} }
\expandafter\def\csname SRef@wb:cusp-cone-ambient-collapse\endcsname{\href{S6_SERIES_GRIND_III.pdf\#nameddest=tag-036I}{III.97 (p.~262)} }
\expandafter\def\csname SRef@wb:cusp-cone-intrinsic-expansion\endcsname{\href{S6_SERIES_GRIND_III.pdf\#nameddest=tag-036S}{III.98 (p.~262)} }
\expandafter\def\csname SRef@wb:cusp-cone-rank-bound\endcsname{\href{S6_SERIES_GRIND_III.pdf\#nameddest=tag-037H}{III.99 (p.~262)} }
\expandafter\def\csname SRef@maps:src:B.1.1\endcsname{\href{S6_SERIES_GRIND_III.pdf\#nameddest=tag-02J4}{III.3.2 (p.~13)} }
\expandafter\def\csname SRef@maps:src:C.1.2\endcsname{\href{S6_SERIES_GRIND_III.pdf\#nameddest=tag-02K3}{III.4.2 (p.~28)} }
\expandafter\def\csname SRef@maps:src:D.1.10\endcsname{\href{S6_SERIES_GRIND_III.pdf\#nameddest=tag-02K8}{III.5.10 (p.~40)} }
\expandafter\def\csname SRef@wb:h24:ret:chart-lemma\endcsname{\href{S6_SERIES_GRIND_III.pdf\#nameddest=tag-03AE}{III.7.30 (p.~300)} }
\expandafter\def\csname SRef@wb:h24:ret:orientation-lemma\endcsname{\href{S6_SERIES_GRIND_III.pdf\#nameddest=tag-03BD}{III.7.31 (p.~301)} }
\expandafter\def\csname SRef@mapsIII:carrier-three-complex-wedge\endcsname{\href{S6_SERIES_GRIND_III.pdf\#nameddest=tag-02QC}{III.7.34 (p.~306)} }
\expandafter\def\csname SRef@mapsIII:carrier-extension-lemma\endcsname{\href{S6_SERIES_GRIND_III.pdf\#nameddest=tag-02Q8}{III.7.36 (p.~308)} }
\expandafter\def\csname SRef@wb:cusp-cone-native-descent\endcsname{\href{S6_SERIES_GRIND_III.pdf\#nameddest=tag-0370}{III.7.5 (p.~255)} }
\expandafter\def\csname SRef@maps:src:D.2.4\endcsname{\href{S6_SERIES_GRIND_III.pdf\#nameddest=tag-02LG}{III.10.1 (p.~319)} }
\expandafter\def\csname SRef@maps:src:A.2.2\endcsname{\href{S6_SERIES_GRIND_III.pdf\#nameddest=tag-02J3}{III.3.1 (p.~12)} }
\expandafter\def\csname SRef@maps:src:C.1.4\endcsname{\href{S6_SERIES_GRIND_III.pdf\#nameddest=tag-02K5}{III.4.4 (p.~30)} }
\expandafter\def\csname SRef@maps:src:D.2.3\endcsname{\href{S6_SERIES_GRIND_III.pdf\#nameddest=tag-02L5}{III.5.13 (p.~43)} }
\expandafter\def\csname SRef@maps:src:D.2.19\endcsname{\href{S6_SERIES_GRIND_III.pdf\#nameddest=tag-02KT}{III.5.28 (p.~80)} }
\expandafter\def\csname SRef@maps:src:D.1.5\endcsname{\href{S6_SERIES_GRIND_III.pdf\#nameddest=tag-02KE}{III.5.5 (p.~36)} }
\expandafter\def\csname SRef@maps:src:E.1.1\endcsname{\href{S6_SERIES_GRIND_V.pdf\#nameddest=tag-02M4}{V.2.1 (p.~14)} }
\expandafter\def\csname SRef@maps:src:E.2.7\endcsname{\href{S6_SERIES_GRIND_V.pdf\#nameddest=tag-02MR}{V.2.11 (p.~29)} }
\expandafter\def\csname SRef@maps:src:E.2.8\endcsname{\href{S6_SERIES_GRIND_V.pdf\#nameddest=tag-02MS}{V.2.12 (p.~30)} }
\expandafter\def\csname SRef@maps:src:E.2.10\endcsname{\href{S6_SERIES_GRIND_V.pdf\#nameddest=tag-02MA}{V.2.14 (p.~33)} }
\expandafter\def\csname SRef@maps:src:E.2.18\endcsname{\href{S6_SERIES_GRIND_V.pdf\#nameddest=tag-02MI}{V.2.22 (p.~56)} }
\expandafter\def\csname SRef@maps:src:E.2.19\endcsname{\href{S6_SERIES_GRIND_V.pdf\#nameddest=tag-02MJ}{V.2.23 (p.~60)} }
\expandafter\def\csname SRef@maps:src:E.2.20\endcsname{\href{S6_SERIES_GRIND_V.pdf\#nameddest=tag-02ML}{V.2.24 (p.~63)} }
\expandafter\def\csname SRef@maps:src:E.2.4\endcsname{\href{S6_SERIES_GRIND_V.pdf\#nameddest=tag-02MO}{V.2.8 (p.~24)} }
\expandafter\def\csname SRef@maps:src:E.2.5\endcsname{\href{S6_SERIES_GRIND_V.pdf\#nameddest=tag-02MP}{V.2.9 (p.~25)} }
\expandafter\def\csname SRef@wb:cusp-cone-metric-scope\endcsname{\href{S6_SERIES_GRIND_III.pdf\#nameddest=tag-036W}{III.7.10 (p.~262)} }
\expandafter\def\csname SRef@wb:cusp-cone-explicit-carriers\endcsname{\href{S6_SERIES_GRIND_III.pdf\#nameddest=tag-036O}{III.7.17 (p.~273)} }
\expandafter\def\csname SRef@mapsIII:carrier-real-periods\endcsname{\href{S6_SERIES_GRIND_III.pdf\#nameddest=tag-02QB}{III.7.2 (p.~244)} }
\expandafter\def\csname SRef@wb:sphere-metric-tangent-scope\endcsname{\href{S6_SERIES_GRIND_III.pdf\#nameddest=tag-03DF}{III.7.29 (p.~294)} }
\expandafter\def\csname SRef@wb:cusp-spec-native-quotient\endcsname{\href{S6_SERIES_GRIND_III.pdf\#nameddest=tag-039Y}{III.7.3 (p.~245)} }
\expandafter\def\csname SRef@wb:h24:ret:framing-proposition\endcsname{\href{S6_SERIES_GRIND_III.pdf\#nameddest=tag-03BC}{III.7.32 (p.~302)} }
\expandafter\def\csname SRef@mapsaudit:finite-degree-three\endcsname{\href{S6_SERIES_GRIND_III.pdf\#nameddest=tag-02QE}{III.7.35 (p.~307)} }
\expandafter\def\csname SRef@mapsIII:angle-forgotten-one-copy-null\endcsname{\href{S6_SERIES_GRIND_III.pdf\#nameddest=tag-02Q7}{III.7.37 (p.~308)} }
\expandafter\def\csname SRef@wb:cusp-cone-S24-local-route\endcsname{\href{S6_SERIES_GRIND_III.pdf\#nameddest=tag-0365}{III.7.7 (p.~258)} }
\expandafter\def\csname SRef@maps:src:D.3.4\endcsname{\href{S6_SERIES_GRIND_III.pdf\#nameddest=tag-02M3}{III.9.2 (p.~317)} }
\expandafter\def\csname SRef@qa:III:layout:12\endcsname{\href{S6_SERIES_GRIND_III.pdf\#nameddest=tag-033S}{III.3 (p.~13)} }
\expandafter\def\csname SRef@wb:cusp-cover-naive-action-failure\endcsname{\href{S6_SERIES_GRIND_III.pdf\#nameddest=tag-0382}{III.7.1 (p.~243)} }
\expandafter\def\csname SRef@wb:cusp-cover-global-scope\endcsname{\href{S6_SERIES_GRIND_III.pdf\#nameddest=tag-037Y}{III.7.1 (p.~243)} }
\expandafter\def\csname SRef@qa:III:layout:11\endcsname{\href{S6_SERIES_GRIND_III.pdf\#nameddest=tag-033R}{III.{{$\langle 1\rangle $1.}} (p.~271)} }
\expandafter\def\csname SRef@wb:cusp-cone-moving-S4-null-homotopy\endcsname{\href{S6_SERIES_GRIND_III.pdf\#nameddest=tag-036Z}{III.7.3.1 (p.~272)} }
\expandafter\def\csname SRef@wb:cusp-cone-sphere-map-homotopy-status\endcsname{\href{S6_SERIES_GRIND_III.pdf\#nameddest=tag-037K}{III.7.3.1 (p.~273)} }
\expandafter\def\csname SRef@wb:cusp-s12-one-copy-invariant\endcsname{\href{S6_SERIES_GRIND_III.pdf\#nameddest=tag-038X}{III.7.4 (p.~281)} }
\expandafter\def\csname SRef@wb:cusp-s12-nullity-scope\endcsname{\href{S6_SERIES_GRIND_III.pdf\#nameddest=tag-038W}{III.7.4 (p.~281)} }
\expandafter\def\csname SRef@qa:III:layout:1\endcsname{\href{S6_SERIES_GRIND_III.pdf\#nameddest=tag-033Q}{III.{{$\langle 1\rangle $5.}} (p.~283)} }
\expandafter\def\csname SRef@qa:III:layout:2\endcsname{\href{S6_SERIES_GRIND_III.pdf\#nameddest=tag-033T}{III.{{$\langle 1\rangle $6.}} (p.~316)} }
\expandafter\def\csname SRef@qa:III:layout:3\endcsname{\href{S6_SERIES_GRIND_III.pdf\#nameddest=tag-033U}{III.{{$\langle 1\rangle $6.}} (p.~316)} }
\expandafter\def\csname SRef@qa:III:layout:4\endcsname{\href{S6_SERIES_GRIND_III.pdf\#nameddest=tag-033V}{III.8 (p.~316)} }
\expandafter\def\csname SRef@sec:III:conventions\endcsname{\href{S6_SERIES_GRIND_III.pdf\#nameddest=tag-034J}{III.2 (p.~12)} }
\expandafter\def\csname SRef@sec:III:maps\endcsname{\href{S6_SERIES_GRIND_III.pdf\#nameddest=tag-034K}{III.3 (p.~12)} }
\expandafter\def\csname SRef@sec:III:os49-night\endcsname{\href{S6_SERIES_GRIND_III.pdf\#nameddest=tag-034L}{III.6 (p.~215)} }
\expandafter\def\csname SRef@sec:III:carriers\endcsname{\href{S6_SERIES_GRIND_III.pdf\#nameddest=tag-034I}{III.7 (p.~238)} }
\expandafter\def\csname SRef@mapsD:repair:modular-collisions\endcsname{\href{S6_SERIES_GRIND_III.pdf\#nameddest=tag-02OJ}{III.8 (p.~310)} }
\expandafter\def\csname SRef@wb:cusp-cover-section\endcsname{\href{S6_SERIES_GRIND_III.pdf\#nameddest=tag-0385}{III.7.1 (p.~238)} }
\expandafter\def\csname SRef@mapsIII:carrier-radial-prerequisites\endcsname{\href{S6_SERIES_GRIND_III.pdf\#nameddest=tag-02QA}{III.7.2 (p.~244)} }
\expandafter\def\csname SRef@wb:cusp-cone-section\endcsname{\href{S6_SERIES_GRIND_III.pdf\#nameddest=tag-037I}{III.7.3 (p.~255)} }
\expandafter\def\csname SRef@wb:cusp-s12-carrier-homotopy\endcsname{\href{S6_SERIES_GRIND_III.pdf\#nameddest=tag-038U}{III.7.4 (p.~279)} }
\expandafter\def\csname SRef@wb:sphere-metric-limit-note\endcsname{\href{S6_SERIES_GRIND_III.pdf\#nameddest=tag-03D9}{III.7.6 (p.~287)} }
\expandafter\def\csname SRef@wb:h24:ret:section\endcsname{\href{S6_SERIES_GRIND_III.pdf\#nameddest=tag-03BF}{III.7.7 (p.~298)} }
\expandafter\def\csname SRef@mapsIII:carrier-extension-section\endcsname{\href{S6_SERIES_GRIND_III.pdf\#nameddest=tag-02Q9}{III.7.8 (p.~308)} }
\expandafter\def\csname SRef@the-source-section-and-component-divisor\endcsname{\href{S6_SERIES_GRIND_III.pdf\#nameddest=tag-0351}{III.6.0.1 (p.~215)} }
\expandafter\def\csname SRef@the-page-including-the-finite-cap-line\endcsname{\href{S6_SERIES_GRIND_III.pdf\#nameddest=tag-034Z}{III.6.0.10 (p.~231)} }
\expandafter\def\csname SRef@the-actual-graph-line-and-return\endcsname{\href{S6_SERIES_GRIND_III.pdf\#nameddest=tag-034Q}{III.6.0.11 (p.~232)} }
\expandafter\def\csname SRef@all-stabilizers-on-the-actual-root-chart\endcsname{\href{S6_SERIES_GRIND_III.pdf\#nameddest=tag-0008}{III.6.0.12 (p.~233)} }
\expandafter\def\csname SRef@the-full-fundamental-group-presentation\endcsname{\href{S6_SERIES_GRIND_III.pdf\#nameddest=tag-034U}{III.6.0.13 (p.~234)} }
\expandafter\def\csname SRef@the-fixed-exterior-and-the-actual-boundary-circle\endcsname{\href{S6_SERIES_GRIND_III.pdf\#nameddest=tag-034T}{III.6.0.14 (p.~235)} }
\expandafter\def\csname SRef@an-integral-basis-of-the-boundary-two-cohomology\endcsname{\href{S6_SERIES_GRIND_III.pdf\#nameddest=tag-000A}{III.6.0.15 (p.~236)} }
\expandafter\def\csname SRef@determine-the-exterior-restriction-without-guessing-a-splitting\endcsname{\href{S6_SERIES_GRIND_III.pdf\#nameddest=tag-01O2}{III.6.0.16 (p.~236)} }
\expandafter\def\csname SRef@the-two-smith-invariants-and-every-homology-group\endcsname{\href{S6_SERIES_GRIND_III.pdf\#nameddest=tag-0353}{III.6.0.17 (p.~237)} }
\expandafter\def\csname SRef@actual-local-fillings-and-their-gluing\endcsname{\href{S6_SERIES_GRIND_III.pdf\#nameddest=tag-0007}{III.6.0.2 (p.~216)} }
\expandafter\def\csname SRef@marked-monodromy-and-its-completeness\endcsname{\href{S6_SERIES_GRIND_III.pdf\#nameddest=tag-02QG}{III.6.0.3 (p.~217)} }
\expandafter\def\csname SRef@integral-topology\endcsname{\href{S6_SERIES_GRIND_III.pdf\#nameddest=tag-01RP}{III.6.0.4 (p.~218)} }
\expandafter\def\csname SRef@the-exact-line-subfan\endcsname{\href{S6_SERIES_GRIND_III.pdf\#nameddest=tag-034S}{III.6.0.5 (p.~225)} }
\expandafter\def\csname SRef@a-page-carrying-the-actual-source-line\endcsname{\href{S6_SERIES_GRIND_III.pdf\#nameddest=tag-0003}{III.6.0.6 (p.~227)} }
\expandafter\def\csname SRef@the-full-graph-to-page-comparison-by-regular-bands\endcsname{\href{S6_SERIES_GRIND_III.pdf\#nameddest=tag-034V}{III.6.0.7 (p.~228)} }
\expandafter\def\csname SRef@characteristic-calculation-and-the-scope-of-the-result\endcsname{\href{S6_SERIES_GRIND_III.pdf\#nameddest=tag-01GY}{III.6.0.8 (p.~229)} }
\expandafter\def\csname SRef@the-actual-closing-map-and-its-origins\endcsname{\href{S6_SERIES_GRIND_III.pdf\#nameddest=tag-034P}{III.6.0.9 (p.~231)} }
\expandafter\def\csname SRef@wb:cusp-cone-optimal-strengthening\endcsname{\href{S6_SERIES_GRIND_III.pdf\#nameddest=tag-0376}{III.7.3.1 (p.~264)} }
\expandafter\def\csname SRef@maps:src:B.5.2\endcsname{\href{S6_SERIES_GRIND_III.pdf\#nameddest=tag-02JC}{III.3.10 (p.~23)} }
\expandafter\def\csname SRef@maps:src:B.1.2\endcsname{\href{S6_SERIES_GRIND_III.pdf\#nameddest=tag-02J5}{III.3.3 (p.~14)} }
\expandafter\def\csname SRef@maps:src:B.2.1\endcsname{\href{S6_SERIES_GRIND_III.pdf\#nameddest=tag-02J6}{III.3.4 (p.~15)} }
\expandafter\def\csname SRef@maps:src:B.3.1\endcsname{\href{S6_SERIES_GRIND_III.pdf\#nameddest=tag-02J7}{III.3.5 (p.~16)} }
\expandafter\def\csname SRef@maps:src:B.3.2\endcsname{\href{S6_SERIES_GRIND_III.pdf\#nameddest=tag-02J8}{III.3.6 (p.~17)} }
\expandafter\def\csname SRef@maps:src:B.3.3\endcsname{\href{S6_SERIES_GRIND_III.pdf\#nameddest=tag-02J9}{III.3.7 (p.~18)} }
\expandafter\def\csname SRef@maps:src:B.4.1\endcsname{\href{S6_SERIES_GRIND_III.pdf\#nameddest=tag-02JA}{III.3.8 (p.~18)} }
\expandafter\def\csname SRef@maps:src:B.5.1\endcsname{\href{S6_SERIES_GRIND_III.pdf\#nameddest=tag-02JB}{III.3.9 (p.~20)} }
\expandafter\def\csname SRef@maps:src:C.1.1\endcsname{\href{S6_SERIES_GRIND_III.pdf\#nameddest=tag-02K2}{III.4.1 (p.~27)} }
\expandafter\def\csname SRef@maps:src:C.1.3\endcsname{\href{S6_SERIES_GRIND_III.pdf\#nameddest=tag-02K4}{III.4.3 (p.~29)} }
\expandafter\def\csname SRef@maps:src:C.1.5\endcsname{\href{S6_SERIES_GRIND_III.pdf\#nameddest=tag-02K6}{III.4.5 (p.~30)} }
\expandafter\def\csname SRef@maps:src:D.1.1\endcsname{\href{S6_SERIES_GRIND_III.pdf\#nameddest=tag-02K7}{III.5.1 (p.~32)} }
\expandafter\def\csname SRef@maps:src:D.1.12\endcsname{\href{S6_SERIES_GRIND_III.pdf\#nameddest=tag-02KA}{III.5.11 (p.~41)} }
\expandafter\def\csname SRef@maps:src:D.2.5\endcsname{\href{S6_SERIES_GRIND_III.pdf\#nameddest=tag-02LR}{III.5.14 (p.~44)} }
\expandafter\def\csname SRef@maps:src:D.2.6\endcsname{\href{S6_SERIES_GRIND_III.pdf\#nameddest=tag-02LW}{III.5.15 (p.~46)} }
\expandafter\def\csname SRef@maps:src:D.2.7\endcsname{\href{S6_SERIES_GRIND_III.pdf\#nameddest=tag-02LX}{III.5.16 (p.~48)} }
\expandafter\def\csname SRef@maps:src:D.2.8\endcsname{\href{S6_SERIES_GRIND_III.pdf\#nameddest=tag-02LY}{III.5.17 (p.~50)} }
\expandafter\def\csname SRef@maps:src:D.2.9\endcsname{\href{S6_SERIES_GRIND_III.pdf\#nameddest=tag-02LZ}{III.5.18 (p.~53)} }
\expandafter\def\csname SRef@maps:src:D.2.10\endcsname{\href{S6_SERIES_GRIND_III.pdf\#nameddest=tag-02KK}{III.5.19 (p.~55)} }
\expandafter\def\csname SRef@maps:src:D.1.2\endcsname{\href{S6_SERIES_GRIND_III.pdf\#nameddest=tag-02KB}{III.5.2 (p.~32)} }
\expandafter\def\csname SRef@maps:src:D.2.11\endcsname{\href{S6_SERIES_GRIND_III.pdf\#nameddest=tag-02KL}{III.5.20 (p.~59)} }
\expandafter\def\csname SRef@maps:src:D.2.12\endcsname{\href{S6_SERIES_GRIND_III.pdf\#nameddest=tag-02KM}{III.5.21 (p.~62)} }
\expandafter\def\csname SRef@maps:src:D.2.13\endcsname{\href{S6_SERIES_GRIND_III.pdf\#nameddest=tag-02KN}{III.5.22 (p.~65)} }
\expandafter\def\csname SRef@maps:src:D.2.14\endcsname{\href{S6_SERIES_GRIND_III.pdf\#nameddest=tag-02KO}{III.5.23 (p.~68)} }
\expandafter\def\csname SRef@maps:src:D.2.15\endcsname{\href{S6_SERIES_GRIND_III.pdf\#nameddest=tag-02KP}{III.5.24 (p.~70)} }
\expandafter\def\csname SRef@maps:src:D.2.16\endcsname{\href{S6_SERIES_GRIND_III.pdf\#nameddest=tag-02KQ}{III.5.25 (p.~73)} }
\expandafter\def\csname SRef@maps:src:D.2.18\endcsname{\href{S6_SERIES_GRIND_III.pdf\#nameddest=tag-02KS}{III.5.27 (p.~78)} }
\expandafter\def\csname SRef@maps:src:D.2.20\endcsname{\href{S6_SERIES_GRIND_III.pdf\#nameddest=tag-02KV}{III.5.29 (p.~82)} }
\expandafter\def\csname SRef@maps:src:D.1.3\endcsname{\href{S6_SERIES_GRIND_III.pdf\#nameddest=tag-02KC}{III.5.3 (p.~33)} }
\expandafter\def\csname SRef@maps:src:D.2.21\endcsname{\href{S6_SERIES_GRIND_III.pdf\#nameddest=tag-02KW}{III.5.30 (p.~86)} }
\expandafter\def\csname SRef@maps:src:D.2.22\endcsname{\href{S6_SERIES_GRIND_III.pdf\#nameddest=tag-02KX}{III.5.31 (p.~87)} }
\expandafter\def\csname SRef@maps:src:D.2.23\endcsname{\href{S6_SERIES_GRIND_III.pdf\#nameddest=tag-02KY}{III.5.32 (p.~91)} }
\expandafter\def\csname SRef@maps:src:D.2.24\endcsname{\href{S6_SERIES_GRIND_III.pdf\#nameddest=tag-02KZ}{III.5.33 (p.~93)} }
\expandafter\def\csname SRef@maps:src:D.2.25\endcsname{\href{S6_SERIES_GRIND_III.pdf\#nameddest=tag-02L0}{III.5.34 (p.~94)} }
\expandafter\def\csname SRef@maps:src:D.2.26\endcsname{\href{S6_SERIES_GRIND_III.pdf\#nameddest=tag-02L1}{III.5.35 (p.~98)} }
\expandafter\def\csname SRef@maps:src:D.2.27\endcsname{\href{S6_SERIES_GRIND_III.pdf\#nameddest=tag-02L2}{III.5.36 (p.~102)} }
\expandafter\def\csname SRef@maps:src:D.2.29\endcsname{\href{S6_SERIES_GRIND_III.pdf\#nameddest=tag-02L4}{III.5.38 (p.~107)} }
\expandafter\def\csname SRef@maps:src:D.2.30\endcsname{\href{S6_SERIES_GRIND_III.pdf\#nameddest=tag-02L6}{III.5.39 (p.~110)} }
\expandafter\def\csname SRef@maps:src:D.1.4\endcsname{\href{S6_SERIES_GRIND_III.pdf\#nameddest=tag-02KD}{III.5.4 (p.~34)} }
\expandafter\def\csname SRef@maps:src:D.2.31\endcsname{\href{S6_SERIES_GRIND_III.pdf\#nameddest=tag-02L7}{III.5.40 (p.~113)} }
\expandafter\def\csname SRef@maps:src:D.2.32\endcsname{\href{S6_SERIES_GRIND_III.pdf\#nameddest=tag-02L8}{III.5.41 (p.~116)} }
\expandafter\def\csname SRef@maps:src:D.2.33\endcsname{\href{S6_SERIES_GRIND_III.pdf\#nameddest=tag-02L9}{III.5.42 (p.~119)} }
\expandafter\def\csname SRef@maps:src:D.2.34\endcsname{\href{S6_SERIES_GRIND_III.pdf\#nameddest=tag-02LA}{III.5.43 (p.~121)} }
\expandafter\def\csname SRef@maps:src:D.2.35\endcsname{\href{S6_SERIES_GRIND_III.pdf\#nameddest=tag-02LB}{III.5.44 (p.~126)} }
\expandafter\def\csname SRef@maps:src:D.2.36\endcsname{\href{S6_SERIES_GRIND_III.pdf\#nameddest=tag-02LC}{III.5.45 (p.~130)} }
\expandafter\def\csname SRef@maps:src:D.2.37\endcsname{\href{S6_SERIES_GRIND_III.pdf\#nameddest=tag-02LD}{III.5.46 (p.~134)} }
\expandafter\def\csname SRef@maps:src:D.2.38\endcsname{\href{S6_SERIES_GRIND_III.pdf\#nameddest=tag-02LE}{III.5.47 (p.~138)} }
\expandafter\def\csname SRef@maps:src:D.2.39\endcsname{\href{S6_SERIES_GRIND_III.pdf\#nameddest=tag-02LF}{III.5.48 (p.~141)} }
\expandafter\def\csname SRef@maps:src:D.2.40\endcsname{\href{S6_SERIES_GRIND_III.pdf\#nameddest=tag-02LH}{III.5.49 (p.~145)} }
\expandafter\def\csname SRef@maps:src:D.2.41\endcsname{\href{S6_SERIES_GRIND_III.pdf\#nameddest=tag-02LI}{III.5.50 (p.~148)} }
\expandafter\def\csname SRef@maps:src:D.2.42\endcsname{\href{S6_SERIES_GRIND_III.pdf\#nameddest=tag-02LJ}{III.5.51 (p.~151)} }
\expandafter\def\csname SRef@maps:src:D.2.43\endcsname{\href{S6_SERIES_GRIND_III.pdf\#nameddest=tag-02LK}{III.5.52 (p.~153)} }
\expandafter\def\csname SRef@maps:src:D.2.44\endcsname{\href{S6_SERIES_GRIND_III.pdf\#nameddest=tag-02LL}{III.5.53 (p.~156)} }
\expandafter\def\csname SRef@maps:src:D.2.45\endcsname{\href{S6_SERIES_GRIND_III.pdf\#nameddest=tag-02LM}{III.5.54 (p.~161)} }
\expandafter\def\csname SRef@maps:src:D.2.46\endcsname{\href{S6_SERIES_GRIND_III.pdf\#nameddest=tag-02LN}{III.5.55 (p.~167)} }
\expandafter\def\csname SRef@maps:src:D.2.47\endcsname{\href{S6_SERIES_GRIND_III.pdf\#nameddest=tag-02LO}{III.5.56 (p.~172)} }
\expandafter\def\csname SRef@maps:src:D.2.48\endcsname{\href{S6_SERIES_GRIND_III.pdf\#nameddest=tag-02LP}{III.5.57 (p.~175)} }
\expandafter\def\csname SRef@maps:src:D.2.49\endcsname{\href{S6_SERIES_GRIND_III.pdf\#nameddest=tag-02LQ}{III.5.58 (p.~182)} }
\expandafter\def\csname SRef@maps:src:D.2.50\endcsname{\href{S6_SERIES_GRIND_III.pdf\#nameddest=tag-02LS}{III.5.59 (p.~187)} }
\expandafter\def\csname SRef@maps:src:D.1.6\endcsname{\href{S6_SERIES_GRIND_III.pdf\#nameddest=tag-02KF}{III.5.6 (p.~36)} }
\expandafter\def\csname SRef@maps:src:D.2.51\endcsname{\href{S6_SERIES_GRIND_III.pdf\#nameddest=tag-02LT}{III.5.60 (p.~191)} }
\expandafter\def\csname SRef@maps:src:D.2.52\endcsname{\href{S6_SERIES_GRIND_III.pdf\#nameddest=tag-02LU}{III.5.61 (p.~194)} }
\expandafter\def\csname SRef@maps:src:D.2.53\endcsname{\href{S6_SERIES_GRIND_III.pdf\#nameddest=tag-02LV}{III.5.62 (p.~196)} }
\expandafter\def\csname SRef@maps:src:D.3.1\endcsname{\href{S6_SERIES_GRIND_III.pdf\#nameddest=tag-02M0}{III.5.63 (p.~211)} }
\expandafter\def\csname SRef@maps:src:D.3.3\endcsname{\href{S6_SERIES_GRIND_III.pdf\#nameddest=tag-02M2}{III.5.64 (p.~213)} }
\expandafter\def\csname SRef@maps:src:D.1.11\endcsname{\href{S6_SERIES_GRIND_III.pdf\#nameddest=tag-02K9}{III.5.65 (p.~214)} }
\expandafter\def\csname SRef@maps:src:D.1.7\endcsname{\href{S6_SERIES_GRIND_III.pdf\#nameddest=tag-02KG}{III.5.7 (p.~37)} }
\expandafter\def\csname SRef@maps:src:D.1.8\endcsname{\href{S6_SERIES_GRIND_III.pdf\#nameddest=tag-02KH}{III.5.8 (p.~38)} }
\expandafter\def\csname SRef@maps:src:D.1.9\endcsname{\href{S6_SERIES_GRIND_III.pdf\#nameddest=tag-02KI}{III.5.9 (p.~39)} }
\expandafter\def\csname SRef@night:P3F.26\endcsname{\href{S6_SERIES_GRIND_III.pdf\#nameddest=tag-02S2}{III.6.1 (p.~215)} }
\expandafter\def\csname SRef@stmt:P3F.26\endcsname{\href{S6_SERIES_GRIND_III.pdf\#nameddest=tag-02S2}{III.6.1 (p.~215)} }
\expandafter\def\csname SRef@night:P3F.27\endcsname{\href{S6_SERIES_GRIND_III.pdf\#nameddest=tag-02S3}{III.6.2 (p.~220)} }
\expandafter\def\csname SRef@stmt:P3F.27\endcsname{\href{S6_SERIES_GRIND_III.pdf\#nameddest=tag-02S3}{III.6.2 (p.~220)} }
\expandafter\def\csname SRef@night:P3F.28\endcsname{\href{S6_SERIES_GRIND_III.pdf\#nameddest=tag-02S4}{III.6.3 (p.~221)} }
\expandafter\def\csname SRef@stmt:P3F.28\endcsname{\href{S6_SERIES_GRIND_III.pdf\#nameddest=tag-02S4}{III.6.3 (p.~221)} }
\expandafter\def\csname SRef@night:P3F.29\endcsname{\href{S6_SERIES_GRIND_III.pdf\#nameddest=tag-02S5}{III.6.4 (p.~223)} }
\expandafter\def\csname SRef@stmt:P3F.29\endcsname{\href{S6_SERIES_GRIND_III.pdf\#nameddest=tag-02S5}{III.6.4 (p.~223)} }
\expandafter\def\csname SRef@night:P3F.31\endcsname{\href{S6_SERIES_GRIND_III.pdf\#nameddest=tag-02S8}{III.6.5 (p.~225)} }
\expandafter\def\csname SRef@stmt:P3F.31\endcsname{\href{S6_SERIES_GRIND_III.pdf\#nameddest=tag-02S8}{III.6.5 (p.~225)} }
\expandafter\def\csname SRef@night:P3F.32\endcsname{\href{S6_SERIES_GRIND_III.pdf\#nameddest=tag-02S9}{III.6.6 (p.~230)} }
\expandafter\def\csname SRef@stmt:P3F.32\endcsname{\href{S6_SERIES_GRIND_III.pdf\#nameddest=tag-02S9}{III.6.6 (p.~230)} }
\expandafter\def\csname SRef@night:P3F.33\endcsname{\href{S6_SERIES_GRIND_III.pdf\#nameddest=tag-02SA}{III.6.7 (p.~233)} }
\expandafter\def\csname SRef@stmt:P3F.33\endcsname{\href{S6_SERIES_GRIND_III.pdf\#nameddest=tag-02SA}{III.6.7 (p.~233)} }
\expandafter\def\csname SRef@night:P3F.34\endcsname{\href{S6_SERIES_GRIND_III.pdf\#nameddest=tag-02SB}{III.6.8 (p.~234)} }
\expandafter\def\csname SRef@stmt:P3F.34\endcsname{\href{S6_SERIES_GRIND_III.pdf\#nameddest=tag-02SB}{III.6.8 (p.~234)} }
\expandafter\def\csname SRef@maps:src:E.2.6\endcsname{\href{S6_SERIES_GRIND_V.pdf\#nameddest=tag-02MQ}{V.2.10 (p.~26)} }
\expandafter\def\csname SRef@maps:src:E.2.9\endcsname{\href{S6_SERIES_GRIND_V.pdf\#nameddest=tag-02MT}{V.2.13 (p.~31)} }
\expandafter\def\csname SRef@maps:src:E.2.11\endcsname{\href{S6_SERIES_GRIND_V.pdf\#nameddest=tag-02MB}{V.2.15 (p.~35)} }
\expandafter\def\csname SRef@maps:src:E.2.12\endcsname{\href{S6_SERIES_GRIND_V.pdf\#nameddest=tag-02MC}{V.2.16 (p.~37)} }
\expandafter\def\csname SRef@maps:src:E.2.13\endcsname{\href{S6_SERIES_GRIND_V.pdf\#nameddest=tag-02MD}{V.2.17 (p.~40)} }
\expandafter\def\csname SRef@maps:src:E.2.14\endcsname{\href{S6_SERIES_GRIND_V.pdf\#nameddest=tag-02ME}{V.2.18 (p.~42)} }
\expandafter\def\csname SRef@maps:src:E.2.15\endcsname{\href{S6_SERIES_GRIND_V.pdf\#nameddest=tag-02MF}{V.2.19 (p.~45)} }
\expandafter\def\csname SRef@maps:src:E.1.2\endcsname{\href{S6_SERIES_GRIND_V.pdf\#nameddest=tag-02M5}{V.2.2 (p.~14)} }
\expandafter\def\csname SRef@maps:src:E.2.16\endcsname{\href{S6_SERIES_GRIND_V.pdf\#nameddest=tag-02MG}{V.2.20 (p.~48)} }
\expandafter\def\csname SRef@maps:src:E.2.17\endcsname{\href{S6_SERIES_GRIND_V.pdf\#nameddest=tag-02MH}{V.2.21 (p.~52)} }
\expandafter\def\csname SRef@maps:src:E.2.21\endcsname{\href{S6_SERIES_GRIND_V.pdf\#nameddest=tag-02MM}{V.2.25 (p.~65)} }
\expandafter\def\csname SRef@maps:src:E.1.3\endcsname{\href{S6_SERIES_GRIND_V.pdf\#nameddest=tag-02M6}{V.2.3 (p.~16)} }
\expandafter\def\csname SRef@maps:src:E.1.4\endcsname{\href{S6_SERIES_GRIND_V.pdf\#nameddest=tag-02M7}{V.2.4 (p.~18)} }
\expandafter\def\csname SRef@maps:src:E.1.5\endcsname{\href{S6_SERIES_GRIND_V.pdf\#nameddest=tag-02M8}{V.2.5 (p.~19)} }
\expandafter\def\csname SRef@maps:src:E.2.1\endcsname{\href{S6_SERIES_GRIND_V.pdf\#nameddest=tag-02M9}{V.2.6 (p.~21)} }
\expandafter\def\csname SRef@maps:src:E.2.3\endcsname{\href{S6_SERIES_GRIND_V.pdf\#nameddest=tag-02MN}{V.2.7 (p.~23)} }
\expandafter\def\csname SRef@wb:cusp-cover-extension-theorem\endcsname{\href{S6_SERIES_GRIND_III.pdf\#nameddest=tag-037X}{III.7.1 (p.~238)} }
\expandafter\def\csname SRef@wb:cusp-cone-obstructions\endcsname{\href{S6_SERIES_GRIND_III.pdf\#nameddest=tag-0374}{III.7.11 (p.~262)} }
\expandafter\def\csname SRef@wb:cusp-cone-optimal-exponents\endcsname{\href{S6_SERIES_GRIND_III.pdf\#nameddest=tag-0375}{III.7.12 (p.~265)} }
\expandafter\def\csname SRef@wb:cusp-cone-angle-forgotten-sphere-carriers\endcsname{\href{S6_SERIES_GRIND_III.pdf\#nameddest=tag-036J}{III.7.13 (p.~266)} }
\expandafter\def\csname SRef@wb:cusp-cone-target-homotopy-types\endcsname{\href{S6_SERIES_GRIND_III.pdf\#nameddest=tag-037M}{III.7.18 (p.~274)} }
\expandafter\def\csname SRef@wb:cusp-cone-exact-decompactification-space\endcsname{\href{S6_SERIES_GRIND_III.pdf\#nameddest=tag-036N}{III.7.19 (p.~276)} }
\expandafter\def\csname SRef@wb:cusp-s12-full-carrier-null\endcsname{\href{S6_SERIES_GRIND_III.pdf\#nameddest=tag-038V}{III.7.21 (p.~280)} }
\expandafter\def\csname SRef@wb:cusp-null:thm\endcsname{\href{S6_SERIES_GRIND_III.pdf\#nameddest=tag-038S}{III.7.23 (p.~282)} }
\expandafter\def\csname SRef@wb:sphere-metric-uniform-collapse\endcsname{\href{S6_SERIES_GRIND_III.pdf\#nameddest=tag-03DG}{III.7.25 (p.~288)} }
\expandafter\def\csname SRef@wb:sphere-metric-fixed-realization\endcsname{\href{S6_SERIES_GRIND_III.pdf\#nameddest=tag-03D6}{III.7.26 (p.~288)} }
\expandafter\def\csname SRef@wb:sphere-metric-scaling-trichotomy\endcsname{\href{S6_SERIES_GRIND_III.pdf\#nameddest=tag-03DC}{III.7.27 (p.~291)} }
\expandafter\def\csname SRef@wb:cusp-spec-globalization-obstruction\endcsname{\href{S6_SERIES_GRIND_III.pdf\#nameddest=tag-039U}{III.7.4 (p.~246)} }
\expandafter\def\csname SRef@wb:cusp-cone-S12-route\endcsname{\href{S6_SERIES_GRIND_III.pdf\#nameddest=tag-035X}{III.7.6 (p.~256)} }
\expandafter\def\csname SRef@wb:cusp-cone-S24-route\endcsname{\href{S6_SERIES_GRIND_III.pdf\#nameddest=tag-036A}{III.7.8 (p.~259)} }
\expandafter\def\csname SRef@mapsD:repair:collision-realization\endcsname{\href{S6_SERIES_GRIND_III.pdf\#nameddest=tag-02OG}{III.8.1 (p.~310)} }
\expandafter\def\csname SRef@maps:src:D.3.2\endcsname{\href{S6_SERIES_GRIND_III.pdf\#nameddest=tag-02M1}{III.9.1 (p.~316)} }
\expandafter\def\csname SRef@maps:src:E.2.2\endcsname{\href{S6_SERIES_GRIND_III.pdf\#nameddest=tag-02MK}{III.9.3 (p.~317)} }
\expandafter\def\csname SRef@maps:eq:B.6.2\endcsname{\href{S6_SERIES_GRIND_IV.pdf\#nameddest=tag-01VZ}{IV.{B.6.2} (p.~68)} }
\expandafter\def\csname SRef@maps:eq:B.6.3\endcsname{\href{S6_SERIES_GRIND_IV.pdf\#nameddest=tag-01WA}{IV.{B.6.3} (p.~68)} }
\expandafter\def\csname SRef@maps:eq:B.6.4\endcsname{\href{S6_SERIES_GRIND_IV.pdf\#nameddest=tag-01WL}{IV.{B.6.4} (p.~69)} }
\expandafter\def\csname SRef@maps:eq:B.6.5\endcsname{\href{S6_SERIES_GRIND_IV.pdf\#nameddest=tag-01WR}{IV.{B.6.5} (p.~69)} }
\expandafter\def\csname SRef@maps:eq:B.6.6\endcsname{\href{S6_SERIES_GRIND_IV.pdf\#nameddest=tag-01WS}{IV.{B.6.6} (p.~69)} }
\expandafter\def\csname SRef@maps:eq:B.6.7\endcsname{\href{S6_SERIES_GRIND_IV.pdf\#nameddest=tag-01WT}{IV.{B.6.7} (p.~70)} }
\expandafter\def\csname SRef@maps:eq:B.6.8\endcsname{\href{S6_SERIES_GRIND_IV.pdf\#nameddest=tag-01WU}{IV.{B.6.8} (p.~70)} }
\expandafter\def\csname SRef@maps:eq:B.6.9\endcsname{\href{S6_SERIES_GRIND_IV.pdf\#nameddest=tag-01WV}{IV.{B.6.9} (p.~70)} }
\expandafter\def\csname SRef@maps:eq:B.6.10\endcsname{\href{S6_SERIES_GRIND_IV.pdf\#nameddest=tag-01VQ}{IV.{B.6.10} (p.~70)} }
\expandafter\def\csname SRef@maps:eq:B.6.11\endcsname{\href{S6_SERIES_GRIND_IV.pdf\#nameddest=tag-01VR}{IV.{B.6.11} (p.~72)} }
\expandafter\def\csname SRef@maps:eq:B.6.12\endcsname{\href{S6_SERIES_GRIND_IV.pdf\#nameddest=tag-01VS}{IV.{B.6.12} (p.~72)} }
\expandafter\def\csname SRef@maps:eq:B.6.13\endcsname{\href{S6_SERIES_GRIND_IV.pdf\#nameddest=tag-01VT}{IV.{B.6.13} (p.~73)} }
\expandafter\def\csname SRef@maps:eq:B.6.14\endcsname{\href{S6_SERIES_GRIND_IV.pdf\#nameddest=tag-01VU}{IV.{B.6.14} (p.~73)} }
\expandafter\def\csname SRef@maps:eq:B.6.15\endcsname{\href{S6_SERIES_GRIND_IV.pdf\#nameddest=tag-01VV}{IV.{B.6.15} (p.~75)} }
\expandafter\def\csname SRef@maps:eq:B.6.16\endcsname{\href{S6_SERIES_GRIND_IV.pdf\#nameddest=tag-01VW}{IV.{B.6.16} (p.~75)} }
\expandafter\def\csname SRef@maps:eq:B.6.17\endcsname{\href{S6_SERIES_GRIND_IV.pdf\#nameddest=tag-01VX}{IV.{B.6.17} (p.~76)} }
\expandafter\def\csname SRef@maps:eq:B.C1\endcsname{\href{S6_SERIES_GRIND_IV.pdf\#nameddest=tag-01WW}{IV.{B.C1} (p.~77)} }
\expandafter\def\csname SRef@maps:eq:B.C2\endcsname{\href{S6_SERIES_GRIND_IV.pdf\#nameddest=tag-01WX}{IV.{B.C2} (p.~77)} }
\expandafter\def\csname SRef@maps:eq:B.C3\endcsname{\href{S6_SERIES_GRIND_IV.pdf\#nameddest=tag-01WY}{IV.{B.C3} (p.~77)} }
\expandafter\def\csname SRef@maps:eq:B.C4\endcsname{\href{S6_SERIES_GRIND_IV.pdf\#nameddest=tag-01WZ}{IV.{B.C4} (p.~78)} }
\expandafter\def\csname SRef@maps:eq:B.P1\endcsname{\href{S6_SERIES_GRIND_IV.pdf\#nameddest=tag-01X0}{IV.{B.P1} (p.~78)} }
\expandafter\def\csname SRef@maps:eq:B.P2\endcsname{\href{S6_SERIES_GRIND_IV.pdf\#nameddest=tag-01X1}{IV.{B.P2} (p.~79)} }
\expandafter\def\csname SRef@maps:eq:B.P3\endcsname{\href{S6_SERIES_GRIND_IV.pdf\#nameddest=tag-01X2}{IV.{B.P3} (p.~79)} }
\expandafter\def\csname SRef@maps:eq:B.P4\endcsname{\href{S6_SERIES_GRIND_IV.pdf\#nameddest=tag-01X3}{IV.{B.P4} (p.~79)} }
\expandafter\def\csname SRef@maps:eq:B.6.19\endcsname{\href{S6_SERIES_GRIND_IV.pdf\#nameddest=tag-01VY}{IV.{B.6.19} (p.~80)} }
\expandafter\def\csname SRef@maps:eq:B.6.20\endcsname{\href{S6_SERIES_GRIND_IV.pdf\#nameddest=tag-01W0}{IV.{B.6.20} (p.~81)} }
\expandafter\def\csname SRef@maps:eq:B.6.21\endcsname{\href{S6_SERIES_GRIND_IV.pdf\#nameddest=tag-01W1}{IV.{B.6.21} (p.~83)} }
\expandafter\def\csname SRef@maps:eq:B.6.22\endcsname{\href{S6_SERIES_GRIND_IV.pdf\#nameddest=tag-01W2}{IV.{B.6.22} (p.~83)} }
\expandafter\def\csname SRef@maps:eq:B.6.23\endcsname{\href{S6_SERIES_GRIND_IV.pdf\#nameddest=tag-01W3}{IV.{B.6.23} (p.~83)} }
\expandafter\def\csname SRef@maps:eq:B.6.24\endcsname{\href{S6_SERIES_GRIND_IV.pdf\#nameddest=tag-01W4}{IV.{B.6.24} (p.~84)} }
\expandafter\def\csname SRef@maps:eq:B.6.25\endcsname{\href{S6_SERIES_GRIND_IV.pdf\#nameddest=tag-01W5}{IV.{B.6.25} (p.~84)} }
\expandafter\def\csname SRef@maps:eq:B.T1\endcsname{\href{S6_SERIES_GRIND_IV.pdf\#nameddest=tag-01XD}{IV.{B.T1} (p.~87)} }
\expandafter\def\csname SRef@maps:eq:B.6.26\endcsname{\href{S6_SERIES_GRIND_IV.pdf\#nameddest=tag-01W6}{IV.{B.6.26} (p.~87)} }
\expandafter\def\csname SRef@maps:eq:B.6.27\endcsname{\href{S6_SERIES_GRIND_IV.pdf\#nameddest=tag-01W7}{IV.{B.6.27} (p.~88)} }
\expandafter\def\csname SRef@maps:eq:B.6.28\endcsname{\href{S6_SERIES_GRIND_IV.pdf\#nameddest=tag-01W8}{IV.{B.6.28} (p.~89)} }
\expandafter\def\csname SRef@maps:eq:B.6.29\endcsname{\href{S6_SERIES_GRIND_IV.pdf\#nameddest=tag-01W9}{IV.{B.6.29} (p.~89)} }
\expandafter\def\csname SRef@maps:eq:B.6.30\endcsname{\href{S6_SERIES_GRIND_IV.pdf\#nameddest=tag-01WB}{IV.{B.6.30} (p.~89)} }
\expandafter\def\csname SRef@maps:eq:B.6.31\endcsname{\href{S6_SERIES_GRIND_IV.pdf\#nameddest=tag-01WC}{IV.{B.6.31} (p.~90)} }
\expandafter\def\csname SRef@maps:eq:B.Q1\endcsname{\href{S6_SERIES_GRIND_IV.pdf\#nameddest=tag-01X4}{IV.{B.Q1} (p.~91)} }
\expandafter\def\csname SRef@maps:eq:B.Q2\endcsname{\href{S6_SERIES_GRIND_IV.pdf\#nameddest=tag-01X5}{IV.{B.Q2} (p.~91)} }
\expandafter\def\csname SRef@maps:eq:B.Q3\endcsname{\href{S6_SERIES_GRIND_IV.pdf\#nameddest=tag-01X6}{IV.{B.Q3} (p.~91)} }
\expandafter\def\csname SRef@maps:eq:B.Q4\endcsname{\href{S6_SERIES_GRIND_IV.pdf\#nameddest=tag-01X7}{IV.{B.Q4} (p.~92)} }
\expandafter\def\csname SRef@maps:eq:B.Q5\endcsname{\href{S6_SERIES_GRIND_IV.pdf\#nameddest=tag-01X8}{IV.{B.Q5} (p.~92)} }
\expandafter\def\csname SRef@maps:eq:B.Q6\endcsname{\href{S6_SERIES_GRIND_IV.pdf\#nameddest=tag-01X9}{IV.{B.Q6} (p.~92)} }
\expandafter\def\csname SRef@maps:eq:B.Q7\endcsname{\href{S6_SERIES_GRIND_IV.pdf\#nameddest=tag-01XA}{IV.{B.Q7} (p.~92)} }
\expandafter\def\csname SRef@maps:eq:B.Q8\endcsname{\href{S6_SERIES_GRIND_IV.pdf\#nameddest=tag-01XB}{IV.{B.Q8} (p.~92)} }
\expandafter\def\csname SRef@maps:eq:B.Q9\endcsname{\href{S6_SERIES_GRIND_IV.pdf\#nameddest=tag-01XC}{IV.{B.Q9} (p.~93)} }
\expandafter\def\csname SRef@maps:eq:B.6.32\endcsname{\href{S6_SERIES_GRIND_IV.pdf\#nameddest=tag-01WD}{IV.{B.6.32} (p.~93)} }
\expandafter\def\csname SRef@maps:eq:B.6.33\endcsname{\href{S6_SERIES_GRIND_IV.pdf\#nameddest=tag-01WE}{IV.{B.6.33} (p.~94)} }
\expandafter\def\csname SRef@maps:eq:B.6.34\endcsname{\href{S6_SERIES_GRIND_IV.pdf\#nameddest=tag-01WF}{IV.{B.6.34} (p.~94)} }
\expandafter\def\csname SRef@famIII:eq:E9\endcsname{\href{S6_SERIES_GRIND_IV.pdf\#nameddest=tag-01RJ}{IV.{E9} (p.~20)} }
\expandafter\def\csname SRef@maps:eq:B.6.35\endcsname{\href{S6_SERIES_GRIND_IV.pdf\#nameddest=tag-01WG}{IV.{B.6.35} (p.~95)} }
\expandafter\def\csname SRef@maps:eq:B.6.36\endcsname{\href{S6_SERIES_GRIND_IV.pdf\#nameddest=tag-01WH}{IV.{B.6.36} (p.~95)} }
\expandafter\def\csname SRef@maps:eq:B.6.37\endcsname{\href{S6_SERIES_GRIND_IV.pdf\#nameddest=tag-01WI}{IV.{B.6.37} (p.~95)} }
\expandafter\def\csname SRef@maps:eq:B.6.38\endcsname{\href{S6_SERIES_GRIND_IV.pdf\#nameddest=tag-01WJ}{IV.{B.6.38} (p.~99)} }
\expandafter\def\csname SRef@maps:eq:B.6.39\endcsname{\href{S6_SERIES_GRIND_IV.pdf\#nameddest=tag-01WK}{IV.{B.6.39} (p.~100)} }
\expandafter\def\csname SRef@famIII:eq:E10\endcsname{\href{S6_SERIES_GRIND_IV.pdf\#nameddest=tag-01RI}{IV.{E10} (p.~20)} }
\expandafter\def\csname SRef@maps:eq:B.6.40\endcsname{\href{S6_SERIES_GRIND_IV.pdf\#nameddest=tag-01WM}{IV.{B.6.40} (p.~100)} }
\expandafter\def\csname SRef@maps:eq:B.6.41\endcsname{\href{S6_SERIES_GRIND_IV.pdf\#nameddest=tag-01WN}{IV.{B.6.41} (p.~100)} }
\expandafter\def\csname SRef@maps:eq:B.6.42\endcsname{\href{S6_SERIES_GRIND_IV.pdf\#nameddest=tag-01WO}{IV.{B.6.42} (p.~100)} }
\expandafter\def\csname SRef@maps:eq:B.6.43\endcsname{\href{S6_SERIES_GRIND_IV.pdf\#nameddest=tag-01WP}{IV.{B.6.43} (p.~103)} }
\expandafter\def\csname SRef@maps:eq:B.6.44\endcsname{\href{S6_SERIES_GRIND_IV.pdf\#nameddest=tag-01WQ}{IV.{B.6.44} (p.~103)} }
\expandafter\def\csname SRef@maps:eq:B.6.1\endcsname{\href{S6_SERIES_GRIND_IV.pdf\#nameddest=tag-01VP}{IV.{B.6.1} (p.~68)} }
\expandafter\def\csname SRef@push:src:C.3.15\endcsname{\href{S6_SERIES_GRIND_IV.pdf\#nameddest=tag-031M}{IV.2.16 (p.~40)} }
\expandafter\def\csname SRef@push:src:C.3.27\endcsname{\href{S6_SERIES_GRIND_IV.pdf\#nameddest=tag-031Z}{IV.2.28 (p.~63)} }
\expandafter\def\csname SRef@push:src:C.5.1\endcsname{\href{S6_SERIES_GRIND_IV.pdf\#nameddest=tag-0329}{IV.2.31 (p.~65)} }
\expandafter\def\csname SRef@maps:src:B.6.8\endcsname{\href{S6_SERIES_GRIND_IV.pdf\#nameddest=tag-02K0}{IV.2.39 (p.~76)} }
\expandafter\def\csname SRef@maps:src:B.6.25\endcsname{\href{S6_SERIES_GRIND_IV.pdf\#nameddest=tag-02JU}{IV.2.55 (p.~104)} }
\expandafter\def\csname SRef@push:eq:C12-u-flow\endcsname{\href{S6_SERIES_GRIND_IV.pdf\#nameddest=tag-02YV}{IV.1 (p.~33)} }
\expandafter\def\csname SRef@push:eq:C12-limiting-identification\endcsname{\href{S6_SERIES_GRIND_IV.pdf\#nameddest=tag-02YU}{IV.2 (p.~34)} }
\expandafter\def\csname SRef@push:eq:C13-real-overlap\endcsname{\href{S6_SERIES_GRIND_IV.pdf\#nameddest=tag-02YX}{IV.3 (p.~35)} }
\expandafter\def\csname SRef@push:eq:C13-weighted-map\endcsname{\href{S6_SERIES_GRIND_IV.pdf\#nameddest=tag-02YY}{IV.4 (p.~35)} }
\expandafter\def\csname SRef@push:eq:C13-cusp-gauge\endcsname{\href{S6_SERIES_GRIND_IV.pdf\#nameddest=tag-02YW}{IV.5 (p.~37)} }
\expandafter\def\csname SRef@push:src:C.3.9\endcsname{\href{S6_SERIES_GRIND_IV.pdf\#nameddest=tag-0326}{IV.2.10 (p.~29)} }
\expandafter\def\csname SRef@push:src:C.3.22\endcsname{\href{S6_SERIES_GRIND_IV.pdf\#nameddest=tag-031U}{IV.2.23 (p.~54)} }
\expandafter\def\csname SRef@push:src:C.3.8\endcsname{\href{S6_SERIES_GRIND_IV.pdf\#nameddest=tag-0325}{IV.2.9 (p.~28)} }
\expandafter\def\csname SRef@iv:rank_two_specialization\endcsname{\href{S6_SERIES_GRIND_IV.pdf\#nameddest=tag-01UU}{IV.3.20 (p.~131)} }
\expandafter\def\csname SRef@fam:src:E.2.2\endcsname{\href{S6_SERIES_GRIND_IV.pdf\#nameddest=tag-01R7}{IV.2.1 (p.~20)} }
\expandafter\def\csname SRef@push:src:C.3.10\endcsname{\href{S6_SERIES_GRIND_IV.pdf\#nameddest=tag-031H}{IV.2.11 (p.~30)} }
\expandafter\def\csname SRef@push:src:C.3.11\endcsname{\href{S6_SERIES_GRIND_IV.pdf\#nameddest=tag-031I}{IV.2.12 (p.~32)} }
\expandafter\def\csname SRef@push:src:C.3.12\endcsname{\href{S6_SERIES_GRIND_IV.pdf\#nameddest=tag-031J}{IV.2.13 (p.~32)} }
\expandafter\def\csname SRef@push:src:C.3.13\endcsname{\href{S6_SERIES_GRIND_IV.pdf\#nameddest=tag-031K}{IV.2.14 (p.~34)} }
\expandafter\def\csname SRef@push:src:C.3.14\endcsname{\href{S6_SERIES_GRIND_IV.pdf\#nameddest=tag-031L}{IV.2.15 (p.~37)} }
\expandafter\def\csname SRef@push:src:C.3.16\endcsname{\href{S6_SERIES_GRIND_IV.pdf\#nameddest=tag-031N}{IV.2.17 (p.~40)} }
\expandafter\def\csname SRef@push:src:C.3.17\endcsname{\href{S6_SERIES_GRIND_IV.pdf\#nameddest=tag-031O}{IV.2.18 (p.~41)} }
\expandafter\def\csname SRef@push:src:C.3.18\endcsname{\href{S6_SERIES_GRIND_IV.pdf\#nameddest=tag-031P}{IV.2.19 (p.~44)} }
\expandafter\def\csname SRef@push:src:C.3.19\endcsname{\href{S6_SERIES_GRIND_IV.pdf\#nameddest=tag-031Q}{IV.2.20 (p.~47)} }
\expandafter\def\csname SRef@push:src:C.3.20\endcsname{\href{S6_SERIES_GRIND_IV.pdf\#nameddest=tag-031S}{IV.2.21 (p.~49)} }
\expandafter\def\csname SRef@push:src:C.3.21\endcsname{\href{S6_SERIES_GRIND_IV.pdf\#nameddest=tag-031T}{IV.2.22 (p.~51)} }
\expandafter\def\csname SRef@push:src:C.3.23\endcsname{\href{S6_SERIES_GRIND_IV.pdf\#nameddest=tag-031V}{IV.2.24 (p.~55)} }
\expandafter\def\csname SRef@push:src:C.2.2\endcsname{\href{S6_SERIES_GRIND_IV.pdf\#nameddest=tag-031E}{IV.2.3 (p.~22)} }
\expandafter\def\csname SRef@push:src:C.3.25\endcsname{\href{S6_SERIES_GRIND_IV.pdf\#nameddest=tag-031X}{IV.2.26 (p.~58)} }
\expandafter\def\csname SRef@push:src:C.3.26\endcsname{\href{S6_SERIES_GRIND_IV.pdf\#nameddest=tag-031Y}{IV.2.27 (p.~61)} }
\expandafter\def\csname SRef@push:src:C.4.1\endcsname{\href{S6_SERIES_GRIND_IV.pdf\#nameddest=tag-0327}{IV.2.29 (p.~64)} }
\expandafter\def\csname SRef@push:src:C.4.2\endcsname{\href{S6_SERIES_GRIND_IV.pdf\#nameddest=tag-0328}{IV.2.30 (p.~65)} }
\expandafter\def\csname SRef@push:src:C.2.3\endcsname{\href{S6_SERIES_GRIND_IV.pdf\#nameddest=tag-031F}{IV.2.4 (p.~23)} }
\expandafter\def\csname SRef@maps:src:B.6.1\endcsname{\href{S6_SERIES_GRIND_IV.pdf\#nameddest=tag-02JD}{IV.2.32 (p.~66)} }
\expandafter\def\csname SRef@maps:src:B.6.3\endcsname{\href{S6_SERIES_GRIND_IV.pdf\#nameddest=tag-02JV}{IV.2.34 (p.~67)} }
\expandafter\def\csname SRef@maps:src:B.6.4\endcsname{\href{S6_SERIES_GRIND_IV.pdf\#nameddest=tag-02JW}{IV.2.35 (p.~69)} }
\expandafter\def\csname SRef@maps:src:B.6.5\endcsname{\href{S6_SERIES_GRIND_IV.pdf\#nameddest=tag-02JX}{IV.2.36 (p.~71)} }
\expandafter\def\csname SRef@maps:src:B.6.6\endcsname{\href{S6_SERIES_GRIND_IV.pdf\#nameddest=tag-02JY}{IV.2.37 (p.~73)} }
\expandafter\def\csname SRef@maps:src:B.6.7\endcsname{\href{S6_SERIES_GRIND_IV.pdf\#nameddest=tag-02JZ}{IV.2.38 (p.~75)} }
\expandafter\def\csname SRef@maps:src:B.6.10\endcsname{\href{S6_SERIES_GRIND_IV.pdf\#nameddest=tag-02JE}{IV.2.40 (p.~77)} }
\expandafter\def\csname SRef@maps:src:B.6.11\endcsname{\href{S6_SERIES_GRIND_IV.pdf\#nameddest=tag-02JF}{IV.2.41 (p.~78)} }
\expandafter\def\csname SRef@push:src:C.3.2\endcsname{\href{S6_SERIES_GRIND_IV.pdf\#nameddest=tag-031R}{IV.2.5 (p.~25)} }
\expandafter\def\csname SRef@maps:src:B.6.12\endcsname{\href{S6_SERIES_GRIND_IV.pdf\#nameddest=tag-02JG}{IV.2.42 (p.~80)} }
\expandafter\def\csname SRef@maps:src:B.6.13\endcsname{\href{S6_SERIES_GRIND_IV.pdf\#nameddest=tag-02JH}{IV.2.43 (p.~83)} }
\expandafter\def\csname SRef@maps:src:B.6.14\endcsname{\href{S6_SERIES_GRIND_IV.pdf\#nameddest=tag-02JI}{IV.2.44 (p.~86)} }
\expandafter\def\csname SRef@maps:src:B.6.15\endcsname{\href{S6_SERIES_GRIND_IV.pdf\#nameddest=tag-02JJ}{IV.2.45 (p.~87)} }
\expandafter\def\csname SRef@maps:src:B.6.16\endcsname{\href{S6_SERIES_GRIND_IV.pdf\#nameddest=tag-02JK}{IV.2.46 (p.~88)} }
\expandafter\def\csname SRef@maps:src:B.6.17\endcsname{\href{S6_SERIES_GRIND_IV.pdf\#nameddest=tag-02JL}{IV.2.47 (p.~91)} }
\expandafter\def\csname SRef@maps:src:B.6.18\endcsname{\href{S6_SERIES_GRIND_IV.pdf\#nameddest=tag-02JM}{IV.2.48 (p.~93)} }
\expandafter\def\csname SRef@maps:src:B.6.19\endcsname{\href{S6_SERIES_GRIND_IV.pdf\#nameddest=tag-02JN}{IV.2.49 (p.~94)} }
\expandafter\def\csname SRef@push:src:C.3.4\endcsname{\href{S6_SERIES_GRIND_IV.pdf\#nameddest=tag-0321}{IV.2.6 (p.~25)} }
\expandafter\def\csname SRef@maps:src:B.6.22\endcsname{\href{S6_SERIES_GRIND_IV.pdf\#nameddest=tag-02JR}{IV.2.52 (p.~99)} }
\expandafter\def\csname SRef@push:src:C.3.6\endcsname{\href{S6_SERIES_GRIND_IV.pdf\#nameddest=tag-0323}{IV.2.7 (p.~25)} }
\expandafter\def\csname SRef@push:src:C.3.7\endcsname{\href{S6_SERIES_GRIND_IV.pdf\#nameddest=tag-0324}{IV.2.8 (p.~26)} }
\expandafter\def\csname SRef@iv:correction:P3F.25\endcsname{\href{S6_SERIES_GRIND_IV.pdf\#nameddest=tag-01TX}{IV.8.1 (p.~933)} }
\expandafter\def\csname SRef@iv:correction:P3K.29\endcsname{\href{S6_SERIES_GRIND_IV.pdf\#nameddest=tag-01U6}{IV.8.10 (p.~937)} }
\expandafter\def\csname SRef@iv:correction:P3K.30\endcsname{\href{S6_SERIES_GRIND_IV.pdf\#nameddest=tag-01U7}{IV.8.11 (p.~938)} }
\expandafter\def\csname SRef@iv:correction:P3K.31\endcsname{\href{S6_SERIES_GRIND_IV.pdf\#nameddest=tag-01U8}{IV.8.12 (p.~938)} }
\expandafter\def\csname SRef@iv:correction:P3K.32\endcsname{\href{S6_SERIES_GRIND_IV.pdf\#nameddest=tag-01U9}{IV.8.13 (p.~939)} }
\expandafter\def\csname SRef@iv:correction:P3K.33\endcsname{\href{S6_SERIES_GRIND_IV.pdf\#nameddest=tag-01UA}{IV.8.14 (p.~939)} }
\expandafter\def\csname SRef@iv:correction:P3K.34\endcsname{\href{S6_SERIES_GRIND_IV.pdf\#nameddest=tag-01UB}{IV.8.15 (p.~940)} }
\expandafter\def\csname SRef@iv:correction:P3F.53\endcsname{\href{S6_SERIES_GRIND_IV.pdf\#nameddest=tag-01TY}{IV.8.16 (p.~940)} }
\expandafter\def\csname SRef@iv:correction:P3K.35\endcsname{\href{S6_SERIES_GRIND_IV.pdf\#nameddest=tag-01UC}{IV.8.17 (p.~940)} }
\expandafter\def\csname SRef@iv:correction:P3K.36\endcsname{\href{S6_SERIES_GRIND_IV.pdf\#nameddest=tag-01UD}{IV.8.18 (p.~941)} }
\expandafter\def\csname SRef@iv:correction:P3K.37\endcsname{\href{S6_SERIES_GRIND_IV.pdf\#nameddest=tag-01UE}{IV.8.19 (p.~941)} }
\expandafter\def\csname SRef@iv:correction:P3K.6\endcsname{\href{S6_SERIES_GRIND_IV.pdf\#nameddest=tag-01US}{IV.8.2 (p.~934)} }
\expandafter\def\csname SRef@iv:correction:P3K.40\endcsname{\href{S6_SERIES_GRIND_IV.pdf\#nameddest=tag-01UF}{IV.8.20 (p.~942)} }
\expandafter\def\csname SRef@iv:correction:P3K.41\endcsname{\href{S6_SERIES_GRIND_IV.pdf\#nameddest=tag-01UG}{IV.8.21 (p.~942)} }
\expandafter\def\csname SRef@iv:correction:P3K.43\endcsname{\href{S6_SERIES_GRIND_IV.pdf\#nameddest=tag-01UH}{IV.8.22 (p.~943)} }
\expandafter\def\csname SRef@iv:correction:P3K.44\endcsname{\href{S6_SERIES_GRIND_IV.pdf\#nameddest=tag-01UI}{IV.8.23 (p.~943)} }
\expandafter\def\csname SRef@iv:correction:P3K.45\endcsname{\href{S6_SERIES_GRIND_IV.pdf\#nameddest=tag-01UJ}{IV.8.24 (p.~944)} }
\expandafter\def\csname SRef@iv:correction:P3K.46\endcsname{\href{S6_SERIES_GRIND_IV.pdf\#nameddest=tag-01UK}{IV.8.25 (p.~944)} }
\expandafter\def\csname SRef@iv:correction:P3K.50\endcsname{\href{S6_SERIES_GRIND_IV.pdf\#nameddest=tag-01UL}{IV.8.26 (p.~944)} }
\expandafter\def\csname SRef@iv:correction:P3K.52\endcsname{\href{S6_SERIES_GRIND_IV.pdf\#nameddest=tag-01UM}{IV.8.27 (p.~944)} }
\expandafter\def\csname SRef@iv:correction:P3K.54\endcsname{\href{S6_SERIES_GRIND_IV.pdf\#nameddest=tag-01UN}{IV.8.28 (p.~945)} }
\expandafter\def\csname SRef@iv:correction:P3K.56\endcsname{\href{S6_SERIES_GRIND_IV.pdf\#nameddest=tag-01UO}{IV.8.29 (p.~945)} }
\expandafter\def\csname SRef@iv:correction:P3K.10\endcsname{\href{S6_SERIES_GRIND_IV.pdf\#nameddest=tag-01TZ}{IV.8.3 (p.~934)} }
\expandafter\def\csname SRef@iv:correction:P3K.57\endcsname{\href{S6_SERIES_GRIND_IV.pdf\#nameddest=tag-01UP}{IV.8.30 (p.~945)} }
\expandafter\def\csname SRef@iv:correction:P3K.58\endcsname{\href{S6_SERIES_GRIND_IV.pdf\#nameddest=tag-01UQ}{IV.8.31 (p.~946)} }
\expandafter\def\csname SRef@iv:correction:P3K.59\endcsname{\href{S6_SERIES_GRIND_IV.pdf\#nameddest=tag-01UR}{IV.8.32 (p.~946)} }
\expandafter\def\csname SRef@iv:correction:P3K.60\endcsname{\href{S6_SERIES_GRIND_IV.pdf\#nameddest=tag-01UT}{IV.8.33 (p.~947)} }
\expandafter\def\csname SRef@iv:correction:P3K.17\endcsname{\href{S6_SERIES_GRIND_IV.pdf\#nameddest=tag-01U0}{IV.8.4 (p.~935)} }
\expandafter\def\csname SRef@iv:correction:P3K.22\endcsname{\href{S6_SERIES_GRIND_IV.pdf\#nameddest=tag-01U1}{IV.8.5 (p.~935)} }
\expandafter\def\csname SRef@iv:correction:P3K.23\endcsname{\href{S6_SERIES_GRIND_IV.pdf\#nameddest=tag-01U2}{IV.8.6 (p.~936)} }
\expandafter\def\csname SRef@iv:correction:P3K.24\endcsname{\href{S6_SERIES_GRIND_IV.pdf\#nameddest=tag-01U3}{IV.8.7 (p.~936)} }
\expandafter\def\csname SRef@iv:correction:P3K.25\endcsname{\href{S6_SERIES_GRIND_IV.pdf\#nameddest=tag-01U4}{IV.8.8 (p.~937)} }
\expandafter\def\csname SRef@iv:correction:P3K.27\endcsname{\href{S6_SERIES_GRIND_IV.pdf\#nameddest=tag-01U5}{IV.8.9 (p.~937)} }
\expandafter\def\csname SRef@fam:src:E.2.3\endcsname{\href{S6_SERIES_GRIND_IV.pdf\#nameddest=tag-01R8}{IV.2.2 (p.~21)} }
\expandafter\def\csname SRef@construction-and-proof\endcsname{\href{S6_SERIES_GRIND_IV.pdf\#nameddest=tag-01GZ}{IV.3.0.1 (p.~127)} }
\expandafter\def\csname SRef@the-marked-double-and-its-bordism\endcsname{\href{S6_SERIES_GRIND_IV.pdf\#nameddest=tag-034Y}{IV.{{$\langle 1\rangle $1.}} (p.~144)} }
\expandafter\def\csname SRef@removal-of-the-cubic-coordinate-by-an-extending-map\endcsname{\href{S6_SERIES_GRIND_IV.pdf\#nameddest=tag-033W}{IV.{{$\langle 1\rangle $2.}} (p.~144)} }
\expandafter\def\csname SRef@the-surgery-and-its-bounding-eight-manifold\endcsname{\href{S6_SERIES_GRIND_IV.pdf\#nameddest=tag-0352}{IV.{{$\langle 1\rangle $3.}} (p.~144)} }
\expandafter\def\csname SRef@the-invariant-and-the-resulting-homomorphism\endcsname{\href{S6_SERIES_GRIND_IV.pdf\#nameddest=tag-034W}{IV.{{$\langle 1\rangle $4.}} (p.~145)} }
\expandafter\def\csname SRef@proof\endcsname{\href{S6_SERIES_GRIND_IV.pdf\#nameddest=tag-02YT}{IV.3.0.2 (p.~129)} }
\expandafter\def\csname SRef@an-explicit-rational-elliptic-surface-and-primitive-section\endcsname{\href{S6_SERIES_GRIND_IV.pdf\#nameddest=tag-0009}{IV.3.0.3 (p.~135)} }
\expandafter\def\csname SRef@the-linearization-and-every-local-filling\endcsname{\href{S6_SERIES_GRIND_IV.pdf\#nameddest=tag-034X}{IV.3.0.4 (p.~136)} }
\expandafter\def\csname SRef@actual-integral-monodromy-not-an-abstract-substitute-tuple\endcsname{\href{S6_SERIES_GRIND_IV.pdf\#nameddest=tag-0006}{IV.3.0.5 (p.~138)} }
\expandafter\def\csname SRef@fundamental-group-and-every-integral-cohomology-differential\endcsname{\href{S6_SERIES_GRIND_IV.pdf\#nameddest=tag-01RN}{IV.3.0.6 (p.~139)} }
\expandafter\def\csname SRef@the-punctured-holomorphic-family-and-its-closing-map\endcsname{\href{S6_SERIES_GRIND_IV.pdf\#nameddest=tag-0350}{IV.3.0.7 (p.~142)} }
\expandafter\def\csname SRef@the-actual-page-and-its-section-line\endcsname{\href{S6_SERIES_GRIND_IV.pdf\#nameddest=tag-034R}{IV.3.0.8 (p.~142)} }
\expandafter\def\csname SRef@executed-verification\endcsname{\href{S6_SERIES_GRIND_IV.pdf\#nameddest=tag-01OP}{IV.3.0.9 (p.~143)} }
\expandafter\def\csname SRef@night:P3K.75\endcsname{\href{S6_SERIES_GRIND_IV.pdf\#nameddest=tag-02VN}{IV.11.1 (p.~1392)} }
\expandafter\def\csname SRef@night:P3K.117\endcsname{\href{S6_SERIES_GRIND_IV.pdf\#nameddest=tag-02TO}{IV.10.7 (p.~1213)} }
\expandafter\def\csname SRef@night:P3K.118\endcsname{\href{S6_SERIES_GRIND_IV.pdf\#nameddest=tag-02TP}{IV.10.8 (p.~1250)} }
\expandafter\def\csname SRef@night:P3K.119\endcsname{\href{S6_SERIES_GRIND_IV.pdf\#nameddest=tag-02TQ}{IV.10.9 (p.~1278)} }
\expandafter\def\csname SRef@night:P3K.109\endcsname{\href{S6_SERIES_GRIND_IV.pdf\#nameddest=tag-02TF}{IV.10.1 (p.~1116)} }
\expandafter\def\csname SRef@night:P3K.110\endcsname{\href{S6_SERIES_GRIND_IV.pdf\#nameddest=tag-02TH}{IV.10.2 (p.~1127)} }
\expandafter\def\csname SRef@night:P3K.111\endcsname{\href{S6_SERIES_GRIND_IV.pdf\#nameddest=tag-02TI}{IV.10.10 (p.~1305)} }
\expandafter\def\csname SRef@night:P3K.112\endcsname{\href{S6_SERIES_GRIND_IV.pdf\#nameddest=tag-02TJ}{IV.10.11 (p.~1369)} }
\expandafter\def\csname SRef@night:P3K.113\endcsname{\href{S6_SERIES_GRIND_IV.pdf\#nameddest=tag-02TK}{IV.10.3 (p.~1135)} }
\expandafter\def\csname SRef@night:P3K.114\endcsname{\href{S6_SERIES_GRIND_IV.pdf\#nameddest=tag-02TL}{IV.10.4 (p.~1147)} }
\expandafter\def\csname SRef@night:P3K.115\endcsname{\href{S6_SERIES_GRIND_IV.pdf\#nameddest=tag-02TM}{IV.10.5 (p.~1177)} }
\expandafter\def\csname SRef@night:P3K.116\endcsname{\href{S6_SERIES_GRIND_IV.pdf\#nameddest=tag-02TN}{IV.10.6 (p.~1193)} }
\expandafter\def\csname SRef@push:src:C.3.24\endcsname{\href{S6_SERIES_GRIND_IV.pdf\#nameddest=tag-031W}{IV.2.25 (p.~57)} }
\expandafter\def\csname SRef@maps:src:B.6.2\endcsname{\href{S6_SERIES_GRIND_IV.pdf\#nameddest=tag-02JO}{IV.2.33 (p.~67)} }
\expandafter\def\csname SRef@maps:src:B.6.20\endcsname{\href{S6_SERIES_GRIND_IV.pdf\#nameddest=tag-02JP}{IV.2.50 (p.~96)} }
\expandafter\def\csname SRef@maps:src:B.6.21\endcsname{\href{S6_SERIES_GRIND_IV.pdf\#nameddest=tag-02JQ}{IV.2.51 (p.~99)} }
\expandafter\def\csname SRef@maps:src:B.6.23\endcsname{\href{S6_SERIES_GRIND_IV.pdf\#nameddest=tag-02JS}{IV.2.53 (p.~101)} }
\expandafter\def\csname SRef@maps:src:B.6.24\endcsname{\href{S6_SERIES_GRIND_IV.pdf\#nameddest=tag-02JT}{IV.2.54 (p.~102)} }
\expandafter\def\csname SRef@night:P3F.1\endcsname{\href{S6_SERIES_GRIND_IV.pdf\#nameddest=tag-02RK}{IV.3.1 (p.~104)} }
\expandafter\def\csname SRef@night:P3F.10\endcsname{\href{S6_SERIES_GRIND_IV.pdf\#nameddest=tag-02RL}{IV.3.10 (p.~117)} }
\expandafter\def\csname SRef@night:P3F.11\endcsname{\href{S6_SERIES_GRIND_IV.pdf\#nameddest=tag-02RM}{IV.3.11 (p.~119)} }
\expandafter\def\csname SRef@night:P3F.12\endcsname{\href{S6_SERIES_GRIND_IV.pdf\#nameddest=tag-02RN}{IV.3.12 (p.~121)} }
\expandafter\def\csname SRef@night:P3F.13\endcsname{\href{S6_SERIES_GRIND_IV.pdf\#nameddest=tag-02RO}{IV.3.13 (p.~123)} }
\expandafter\def\csname SRef@night:P3F.14\endcsname{\href{S6_SERIES_GRIND_IV.pdf\#nameddest=tag-02RP}{IV.3.14 (p.~124)} }
\expandafter\def\csname SRef@night:P3F.15\endcsname{\href{S6_SERIES_GRIND_IV.pdf\#nameddest=tag-02RQ}{IV.3.15 (p.~125)} }
\expandafter\def\csname SRef@night:P3F.16\endcsname{\href{S6_SERIES_GRIND_IV.pdf\#nameddest=tag-02RR}{IV.3.16 (p.~126)} }
\expandafter\def\csname SRef@night:P3F.17\endcsname{\href{S6_SERIES_GRIND_IV.pdf\#nameddest=tag-02RS}{IV.3.17 (p.~128)} }
\expandafter\def\csname SRef@night:P3F.18\endcsname{\href{S6_SERIES_GRIND_IV.pdf\#nameddest=tag-02RT}{IV.3.18 (p.~129)} }
\expandafter\def\csname SRef@night:P3F.19\endcsname{\href{S6_SERIES_GRIND_IV.pdf\#nameddest=tag-02RU}{IV.3.19 (p.~130)} }
\expandafter\def\csname SRef@night:P3F.2\endcsname{\href{S6_SERIES_GRIND_IV.pdf\#nameddest=tag-02RV}{IV.3.2 (p.~107)} }
\expandafter\def\csname SRef@night:P3F.20\endcsname{\href{S6_SERIES_GRIND_IV.pdf\#nameddest=tag-02RW}{IV.3.21 (p.~134)} }
\expandafter\def\csname SRef@night:P3F.21\endcsname{\href{S6_SERIES_GRIND_IV.pdf\#nameddest=tag-02RX}{IV.3.22 (p.~134)} }
\expandafter\def\csname SRef@night:P3F.22\endcsname{\href{S6_SERIES_GRIND_IV.pdf\#nameddest=tag-02RY}{IV.3.23 (p.~144)} }
\expandafter\def\csname SRef@night:P3F.23\endcsname{\href{S6_SERIES_GRIND_IV.pdf\#nameddest=tag-02RZ}{IV.3.24 (p.~145)} }
\expandafter\def\csname SRef@night:P3F.24\endcsname{\href{S6_SERIES_GRIND_IV.pdf\#nameddest=tag-02S0}{IV.3.25 (p.~147)} }
\expandafter\def\csname SRef@night:P3F.30\endcsname{\href{S6_SERIES_GRIND_IV.pdf\#nameddest=tag-02S7}{IV.3.26 (p.~148)} }
\expandafter\def\csname SRef@night:P3F.35\endcsname{\href{S6_SERIES_GRIND_IV.pdf\#nameddest=tag-02SC}{IV.3.28 (p.~152)} }
\expandafter\def\csname SRef@night:P3F.36\endcsname{\href{S6_SERIES_GRIND_IV.pdf\#nameddest=tag-02SD}{IV.3.29 (p.~154)} }
\expandafter\def\csname SRef@night:P3F.37\endcsname{\href{S6_SERIES_GRIND_IV.pdf\#nameddest=tag-02SE}{IV.3.30 (p.~159)} }
\expandafter\def\csname SRef@night:P3F.3\endcsname{\href{S6_SERIES_GRIND_IV.pdf\#nameddest=tag-02S6}{IV.3.3 (p.~108)} }
\expandafter\def\csname SRef@night:P3F.38\endcsname{\href{S6_SERIES_GRIND_IV.pdf\#nameddest=tag-02SF}{IV.3.31 (p.~164)} }
\expandafter\def\csname SRef@night:P3F.39\endcsname{\href{S6_SERIES_GRIND_IV.pdf\#nameddest=tag-02SG}{IV.3.32 (p.~169)} }
\expandafter\def\csname SRef@night:P3F.40\endcsname{\href{S6_SERIES_GRIND_IV.pdf\#nameddest=tag-02SI}{IV.3.33 (p.~172)} }
\expandafter\def\csname SRef@night:P3F.41\endcsname{\href{S6_SERIES_GRIND_IV.pdf\#nameddest=tag-02SJ}{IV.3.34 (p.~175)} }
\expandafter\def\csname SRef@night:P3F.42\endcsname{\href{S6_SERIES_GRIND_IV.pdf\#nameddest=tag-02SK}{IV.3.35 (p.~180)} }
\expandafter\def\csname SRef@night:P3F.43\endcsname{\href{S6_SERIES_GRIND_IV.pdf\#nameddest=tag-02SL}{IV.3.36 (p.~186)} }
\expandafter\def\csname SRef@night:P3F.44\endcsname{\href{S6_SERIES_GRIND_IV.pdf\#nameddest=tag-02SM}{IV.3.37 (p.~188)} }
\expandafter\def\csname SRef@night:P3F.45\endcsname{\href{S6_SERIES_GRIND_IV.pdf\#nameddest=tag-02SN}{IV.3.38 (p.~191)} }
\expandafter\def\csname SRef@night:P3F.46\endcsname{\href{S6_SERIES_GRIND_IV.pdf\#nameddest=tag-02SO}{IV.3.39 (p.~195)} }
\expandafter\def\csname SRef@night:P3F.47\endcsname{\href{S6_SERIES_GRIND_IV.pdf\#nameddest=tag-02SP}{IV.3.40 (p.~202)} }
\expandafter\def\csname SRef@night:P3F.4\endcsname{\href{S6_SERIES_GRIND_IV.pdf\#nameddest=tag-02SH}{IV.3.4 (p.~109)} }
\expandafter\def\csname SRef@night:P3F.48\endcsname{\href{S6_SERIES_GRIND_IV.pdf\#nameddest=tag-02SQ}{IV.3.41 (p.~207)} }
\expandafter\def\csname SRef@night:P3F.49\endcsname{\href{S6_SERIES_GRIND_IV.pdf\#nameddest=tag-02SR}{IV.3.42 (p.~210)} }
\expandafter\def\csname SRef@night:P3F.50\endcsname{\href{S6_SERIES_GRIND_IV.pdf\#nameddest=tag-02ST}{IV.3.43 (p.~212)} }
\expandafter\def\csname SRef@night:P3F.51\endcsname{\href{S6_SERIES_GRIND_IV.pdf\#nameddest=tag-02SU}{IV.3.44 (p.~216)} }
\expandafter\def\csname SRef@night:P3F.52\endcsname{\href{S6_SERIES_GRIND_IV.pdf\#nameddest=tag-02SV}{IV.3.45 (p.~222)} }
\expandafter\def\csname SRef@night:P3F.54\endcsname{\href{S6_SERIES_GRIND_IV.pdf\#nameddest=tag-02SX}{IV.3.46 (p.~225)} }
\expandafter\def\csname SRef@night:P3F.55\endcsname{\href{S6_SERIES_GRIND_IV.pdf\#nameddest=tag-02SY}{IV.3.47 (p.~229)} }
\expandafter\def\csname SRef@night:P3F.56\endcsname{\href{S6_SERIES_GRIND_IV.pdf\#nameddest=tag-02SZ}{IV.3.48 (p.~234)} }
\expandafter\def\csname SRef@night:P3F.5\endcsname{\href{S6_SERIES_GRIND_IV.pdf\#nameddest=tag-02SS}{IV.3.5 (p.~110)} }
\expandafter\def\csname SRef@night:P3F.6\endcsname{\href{S6_SERIES_GRIND_IV.pdf\#nameddest=tag-02T0}{IV.3.6 (p.~111)} }
\expandafter\def\csname SRef@night:P3F.7\endcsname{\href{S6_SERIES_GRIND_IV.pdf\#nameddest=tag-02T1}{IV.3.7 (p.~112)} }
\expandafter\def\csname SRef@night:P3F.8\endcsname{\href{S6_SERIES_GRIND_IV.pdf\#nameddest=tag-02T2}{IV.3.8 (p.~114)} }
\expandafter\def\csname SRef@night:P3F.9\endcsname{\href{S6_SERIES_GRIND_IV.pdf\#nameddest=tag-02T3}{IV.3.9 (p.~115)} }
\expandafter\def\csname SRef@night:P3R.1\endcsname{\href{S6_SERIES_GRIND_IV.pdf\#nameddest=tag-02WE}{IV.4.1 (p.~243)} }
\expandafter\def\csname SRef@night:P3K.9\endcsname{\href{S6_SERIES_GRIND_IV.pdf\#nameddest=tag-02W3}{IV.4.10 (p.~254)} }
\expandafter\def\csname SRef@night:P3K.11\endcsname{\href{S6_SERIES_GRIND_IV.pdf\#nameddest=tag-02TG}{IV.4.11 (p.~257)} }
\expandafter\def\csname SRef@night:P3K.12\endcsname{\href{S6_SERIES_GRIND_IV.pdf\#nameddest=tag-02TR}{IV.4.12 (p.~260)} }
\expandafter\def\csname SRef@night:P3K.13\endcsname{\href{S6_SERIES_GRIND_IV.pdf\#nameddest=tag-02TS}{IV.4.13 (p.~261)} }
\expandafter\def\csname SRef@night:P3K.14\endcsname{\href{S6_SERIES_GRIND_IV.pdf\#nameddest=tag-02TT}{IV.4.14 (p.~263)} }
\expandafter\def\csname SRef@night:P3K.15\endcsname{\href{S6_SERIES_GRIND_IV.pdf\#nameddest=tag-02TU}{IV.4.15 (p.~264)} }
\expandafter\def\csname SRef@night:P3K.16\endcsname{\href{S6_SERIES_GRIND_IV.pdf\#nameddest=tag-02TV}{IV.4.16 (p.~268)} }
\expandafter\def\csname SRef@night:P3K.18\endcsname{\href{S6_SERIES_GRIND_IV.pdf\#nameddest=tag-02TX}{IV.4.17 (p.~270)} }
\expandafter\def\csname SRef@night:P3K.19\endcsname{\href{S6_SERIES_GRIND_IV.pdf\#nameddest=tag-02TY}{IV.4.18 (p.~271)} }
\expandafter\def\csname SRef@night:P3K.20\endcsname{\href{S6_SERIES_GRIND_IV.pdf\#nameddest=tag-02U0}{IV.4.19 (p.~273)} }
\expandafter\def\csname SRef@night:P3R.2\endcsname{\href{S6_SERIES_GRIND_IV.pdf\#nameddest=tag-02WF}{IV.4.2 (p.~243)} }
\expandafter\def\csname SRef@night:P3K.21\endcsname{\href{S6_SERIES_GRIND_IV.pdf\#nameddest=tag-02U1}{IV.4.20 (p.~276)} }
\expandafter\def\csname SRef@night:P3K.26\endcsname{\href{S6_SERIES_GRIND_IV.pdf\#nameddest=tag-02U6}{IV.4.21 (p.~278)} }
\expandafter\def\csname SRef@night:P3K.28\endcsname{\href{S6_SERIES_GRIND_IV.pdf\#nameddest=tag-02U8}{IV.4.22 (p.~283)} }
\expandafter\def\csname SRef@night:P3K.1\endcsname{\href{S6_SERIES_GRIND_IV.pdf\#nameddest=tag-02T4}{IV.4.3 (p.~244)} }
\expandafter\def\csname SRef@night:P3K.2\endcsname{\href{S6_SERIES_GRIND_IV.pdf\#nameddest=tag-02TZ}{IV.4.4 (p.~245)} }
\expandafter\def\csname SRef@night:P3K.3\endcsname{\href{S6_SERIES_GRIND_IV.pdf\#nameddest=tag-02UA}{IV.4.5 (p.~246)} }
\expandafter\def\csname SRef@night:P3K.4\endcsname{\href{S6_SERIES_GRIND_IV.pdf\#nameddest=tag-02UK}{IV.4.6 (p.~247)} }
\expandafter\def\csname SRef@night:P3K.5\endcsname{\href{S6_SERIES_GRIND_IV.pdf\#nameddest=tag-02UV}{IV.4.7 (p.~248)} }
\expandafter\def\csname SRef@night:P3K.7\endcsname{\href{S6_SERIES_GRIND_IV.pdf\#nameddest=tag-02VH}{IV.4.8 (p.~249)} }
\expandafter\def\csname SRef@night:P3K.8\endcsname{\href{S6_SERIES_GRIND_IV.pdf\#nameddest=tag-02VS}{IV.4.9 (p.~251)} }
\expandafter\def\csname SRef@night:P3K.38\endcsname{\href{S6_SERIES_GRIND_IV.pdf\#nameddest=tag-02UJ}{IV.5.1 (p.~291)} }
\expandafter\def\csname SRef@night:P3K.62\endcsname{\href{S6_SERIES_GRIND_IV.pdf\#nameddest=tag-02V9}{IV.5.10 (p.~333)} }
\expandafter\def\csname SRef@night:P3K.63\endcsname{\href{S6_SERIES_GRIND_IV.pdf\#nameddest=tag-02VA}{IV.5.11 (p.~340)} }
\expandafter\def\csname SRef@night:P3K.64\endcsname{\href{S6_SERIES_GRIND_IV.pdf\#nameddest=tag-02VB}{IV.5.12 (p.~348)} }
\expandafter\def\csname SRef@night:P3K.65\endcsname{\href{S6_SERIES_GRIND_IV.pdf\#nameddest=tag-02VC}{IV.5.13 (p.~352)} }
\expandafter\def\csname SRef@night:P3K.66\endcsname{\href{S6_SERIES_GRIND_IV.pdf\#nameddest=tag-02VD}{IV.5.14 (p.~362)} }
\expandafter\def\csname SRef@night:P3K.67\endcsname{\href{S6_SERIES_GRIND_IV.pdf\#nameddest=tag-02VE}{IV.5.15 (p.~369)} }
\expandafter\def\csname SRef@night:P3K.42\endcsname{\href{S6_SERIES_GRIND_IV.pdf\#nameddest=tag-02UN}{IV.5.2 (p.~298)} }
\expandafter\def\csname SRef@night:P3K.47\endcsname{\href{S6_SERIES_GRIND_IV.pdf\#nameddest=tag-02US}{IV.5.3 (p.~304)} }
\expandafter\def\csname SRef@night:P3K.48\endcsname{\href{S6_SERIES_GRIND_IV.pdf\#nameddest=tag-02UT}{IV.5.4 (p.~307)} }
\expandafter\def\csname SRef@night:P3K.49\endcsname{\href{S6_SERIES_GRIND_IV.pdf\#nameddest=tag-02UU}{IV.5.5 (p.~311)} }
\expandafter\def\csname SRef@night:P3K.51\endcsname{\href{S6_SERIES_GRIND_IV.pdf\#nameddest=tag-02UX}{IV.5.6 (p.~318)} }
\expandafter\def\csname SRef@night:P3K.53\endcsname{\href{S6_SERIES_GRIND_IV.pdf\#nameddest=tag-02UZ}{IV.5.7 (p.~321)} }
\expandafter\def\csname SRef@night:P3K.55\endcsname{\href{S6_SERIES_GRIND_IV.pdf\#nameddest=tag-02V1}{IV.5.8 (p.~325)} }
\expandafter\def\csname SRef@night:P3K.61\endcsname{\href{S6_SERIES_GRIND_IV.pdf\#nameddest=tag-02V8}{IV.5.9 (p.~330)} }
\expandafter\def\csname SRef@night:P3K.68\endcsname{\href{S6_SERIES_GRIND_IV.pdf\#nameddest=tag-02VF}{IV.6.1 (p.~374)} }
\expandafter\def\csname SRef@night:P3K.69\endcsname{\href{S6_SERIES_GRIND_IV.pdf\#nameddest=tag-02VG}{IV.6.2 (p.~380)} }
\expandafter\def\csname SRef@night:P3K.70\endcsname{\href{S6_SERIES_GRIND_IV.pdf\#nameddest=tag-02VI}{IV.6.3 (p.~384)} }
\expandafter\def\csname SRef@night:P3K.71\endcsname{\href{S6_SERIES_GRIND_IV.pdf\#nameddest=tag-02VJ}{IV.6.4 (p.~389)} }
\expandafter\def\csname SRef@night:P3K.72\endcsname{\href{S6_SERIES_GRIND_IV.pdf\#nameddest=tag-02VK}{IV.6.5 (p.~393)} }
\expandafter\def\csname SRef@night:P3K.73\endcsname{\href{S6_SERIES_GRIND_IV.pdf\#nameddest=tag-02VL}{IV.6.6 (p.~400)} }
\expandafter\def\csname SRef@night:P3K.74\endcsname{\href{S6_SERIES_GRIND_IV.pdf\#nameddest=tag-02VM}{IV.6.7 (p.~405)} }
\expandafter\def\csname SRef@night:P3K.76\endcsname{\href{S6_SERIES_GRIND_IV.pdf\#nameddest=tag-02VO}{IV.7.1 (p.~407)} }
\expandafter\def\csname SRef@night:P3K.85\endcsname{\href{S6_SERIES_GRIND_IV.pdf\#nameddest=tag-02VY}{IV.7.10 (p.~494)} }
\expandafter\def\csname SRef@night:P3K.86\endcsname{\href{S6_SERIES_GRIND_IV.pdf\#nameddest=tag-02VZ}{IV.7.11 (p.~499)} }
\expandafter\def\csname SRef@night:P3K.87\endcsname{\href{S6_SERIES_GRIND_IV.pdf\#nameddest=tag-02W0}{IV.7.12 (p.~505)} }
\expandafter\def\csname SRef@night:P3K.88\endcsname{\href{S6_SERIES_GRIND_IV.pdf\#nameddest=tag-02W1}{IV.7.13 (p.~509)} }
\expandafter\def\csname SRef@night:P3K.89\endcsname{\href{S6_SERIES_GRIND_IV.pdf\#nameddest=tag-02W2}{IV.7.14 (p.~513)} }
\expandafter\def\csname SRef@night:P3K.90\endcsname{\href{S6_SERIES_GRIND_IV.pdf\#nameddest=tag-02W4}{IV.7.15 (p.~521)} }
\expandafter\def\csname SRef@night:P3K.91\endcsname{\href{S6_SERIES_GRIND_IV.pdf\#nameddest=tag-02W5}{IV.7.16 (p.~525)} }
\expandafter\def\csname SRef@night:P3K.92\endcsname{\href{S6_SERIES_GRIND_IV.pdf\#nameddest=tag-02W6}{IV.7.17 (p.~531)} }
\expandafter\def\csname SRef@night:P3K.93\endcsname{\href{S6_SERIES_GRIND_IV.pdf\#nameddest=tag-02W7}{IV.7.18 (p.~538)} }
\expandafter\def\csname SRef@night:P3K.94\endcsname{\href{S6_SERIES_GRIND_IV.pdf\#nameddest=tag-02W8}{IV.7.19 (p.~543)} }
\expandafter\def\csname SRef@night:P3K.77\endcsname{\href{S6_SERIES_GRIND_IV.pdf\#nameddest=tag-02VP}{IV.7.2 (p.~414)} }
\expandafter\def\csname SRef@night:P3K.95\endcsname{\href{S6_SERIES_GRIND_IV.pdf\#nameddest=tag-02W9}{IV.7.20 (p.~551)} }
\expandafter\def\csname SRef@night:P3K.96\endcsname{\href{S6_SERIES_GRIND_IV.pdf\#nameddest=tag-02WA}{IV.7.21 (p.~562)} }
\expandafter\def\csname SRef@night:P3K.97\endcsname{\href{S6_SERIES_GRIND_IV.pdf\#nameddest=tag-02WB}{IV.7.22 (p.~569)} }
\expandafter\def\csname SRef@night:P3K.98\endcsname{\href{S6_SERIES_GRIND_IV.pdf\#nameddest=tag-02WC}{IV.7.23 (p.~578)} }
\expandafter\def\csname SRef@night:P3K.99\endcsname{\href{S6_SERIES_GRIND_IV.pdf\#nameddest=tag-02WD}{IV.7.24 (p.~587)} }
\expandafter\def\csname SRef@night:P3K.100\endcsname{\href{S6_SERIES_GRIND_IV.pdf\#nameddest=tag-02T6}{IV.7.25 (p.~662)} }
\expandafter\def\csname SRef@night:P3K.101\endcsname{\href{S6_SERIES_GRIND_IV.pdf\#nameddest=tag-02T7}{IV.7.26 (p.~669)} }
\expandafter\def\csname SRef@night:P3K.102\endcsname{\href{S6_SERIES_GRIND_IV.pdf\#nameddest=tag-02T8}{IV.7.27 (p.~758)} }
\expandafter\def\csname SRef@night:P3K.103\endcsname{\href{S6_SERIES_GRIND_IV.pdf\#nameddest=tag-02T9}{IV.7.28 (p.~768)} }
\expandafter\def\csname SRef@night:P3K.104\endcsname{\href{S6_SERIES_GRIND_IV.pdf\#nameddest=tag-02TA}{IV.7.29 (p.~774)} }
\expandafter\def\csname SRef@night:P3K.78\endcsname{\href{S6_SERIES_GRIND_IV.pdf\#nameddest=tag-02VQ}{IV.7.3 (p.~420)} }
\expandafter\def\csname SRef@night:P3K.105\endcsname{\href{S6_SERIES_GRIND_IV.pdf\#nameddest=tag-02TB}{IV.7.30 (p.~785)} }
\expandafter\def\csname SRef@night:P3K.106\endcsname{\href{S6_SERIES_GRIND_IV.pdf\#nameddest=tag-02TC}{IV.7.31 (p.~798)} }
\expandafter\def\csname SRef@night:P3K.108\endcsname{\href{S6_SERIES_GRIND_IV.pdf\#nameddest=tag-02TE}{IV.7.32 (p.~861)} }
\expandafter\def\csname SRef@iv:i3_all_reduced_caps\endcsname{\href{S6_SERIES_GRIND_IV.pdf\#nameddest=tag-02TE}{IV.7.32 (p.~861)} }
\expandafter\def\csname SRef@night:P3K.107\endcsname{\href{S6_SERIES_GRIND_IV.pdf\#nameddest=tag-02TD}{IV.7.33 (p.~867)} }
\expandafter\def\csname SRef@night:P3K.79\endcsname{\href{S6_SERIES_GRIND_IV.pdf\#nameddest=tag-02VR}{IV.7.4 (p.~426)} }
\expandafter\def\csname SRef@night:P3K.80\endcsname{\href{S6_SERIES_GRIND_IV.pdf\#nameddest=tag-02VT}{IV.7.5 (p.~431)} }
\expandafter\def\csname SRef@night:P3K.81\endcsname{\href{S6_SERIES_GRIND_IV.pdf\#nameddest=tag-02VU}{IV.7.6 (p.~467)} }
\expandafter\def\csname SRef@night:P3K.82\endcsname{\href{S6_SERIES_GRIND_IV.pdf\#nameddest=tag-02VV}{IV.7.7 (p.~475)} }
\expandafter\def\csname SRef@night:P3K.83\endcsname{\href{S6_SERIES_GRIND_IV.pdf\#nameddest=tag-02VW}{IV.7.8 (p.~478)} }
\expandafter\def\csname SRef@night:P3K.84\endcsname{\href{S6_SERIES_GRIND_IV.pdf\#nameddest=tag-02VX}{IV.7.9 (p.~490)} }
\expandafter\def\csname SRef@night:P3F.25\endcsname{\href{S6_SERIES_GRIND_IV.pdf\#nameddest=tag-02S1}{IV.9.1 (p.~947)} }
\expandafter\def\csname SRef@night:P3K.29\endcsname{\href{S6_SERIES_GRIND_IV.pdf\#nameddest=tag-02U9}{IV.9.10 (p.~988)} }
\expandafter\def\csname SRef@night:P3K.30\endcsname{\href{S6_SERIES_GRIND_IV.pdf\#nameddest=tag-02UB}{IV.9.11 (p.~992)} }
\expandafter\def\csname SRef@night:P3K.31\endcsname{\href{S6_SERIES_GRIND_IV.pdf\#nameddest=tag-02UC}{IV.9.12 (p.~997)} }
\expandafter\def\csname SRef@night:P3K.32\endcsname{\href{S6_SERIES_GRIND_IV.pdf\#nameddest=tag-02UD}{IV.9.13 (p.~1000)} }
\expandafter\def\csname SRef@night:P3K.33\endcsname{\href{S6_SERIES_GRIND_IV.pdf\#nameddest=tag-02UE}{IV.9.14 (p.~1010)} }
\expandafter\def\csname SRef@night:P3K.34\endcsname{\href{S6_SERIES_GRIND_IV.pdf\#nameddest=tag-02UF}{IV.9.15 (p.~1015)} }
\expandafter\def\csname SRef@night:P3F.53\endcsname{\href{S6_SERIES_GRIND_IV.pdf\#nameddest=tag-02SW}{IV.9.16 (p.~1020)} }
\expandafter\def\csname SRef@night:P3K.35\endcsname{\href{S6_SERIES_GRIND_IV.pdf\#nameddest=tag-02UG}{IV.9.17 (p.~1025)} }
\expandafter\def\csname SRef@night:P3K.37\endcsname{\href{S6_SERIES_GRIND_IV.pdf\#nameddest=tag-02UI}{IV.9.18 (p.~1032)} }
\expandafter\def\csname SRef@night:P3K.40\endcsname{\href{S6_SERIES_GRIND_IV.pdf\#nameddest=tag-02UL}{IV.9.19 (p.~1036)} }
\expandafter\def\csname SRef@night:P3K.6\endcsname{\href{S6_SERIES_GRIND_IV.pdf\#nameddest=tag-02V6}{IV.9.2 (p.~950)} }
\expandafter\def\csname SRef@night:P3K.41\endcsname{\href{S6_SERIES_GRIND_IV.pdf\#nameddest=tag-02UM}{IV.9.20 (p.~1040)} }
\expandafter\def\csname SRef@night:P3K.43\endcsname{\href{S6_SERIES_GRIND_IV.pdf\#nameddest=tag-02UO}{IV.9.21 (p.~1047)} }
\expandafter\def\csname SRef@night:P3K.44\endcsname{\href{S6_SERIES_GRIND_IV.pdf\#nameddest=tag-02UP}{IV.9.22 (p.~1050)} }
\expandafter\def\csname SRef@night:P3K.45\endcsname{\href{S6_SERIES_GRIND_IV.pdf\#nameddest=tag-02UQ}{IV.9.23 (p.~1055)} }
\expandafter\def\csname SRef@night:P3K.46\endcsname{\href{S6_SERIES_GRIND_IV.pdf\#nameddest=tag-02UR}{IV.9.24 (p.~1058)} }
\expandafter\def\csname SRef@night:P3K.50\endcsname{\href{S6_SERIES_GRIND_IV.pdf\#nameddest=tag-02UW}{IV.9.25 (p.~1061)} }
\expandafter\def\csname SRef@night:P3K.52\endcsname{\href{S6_SERIES_GRIND_IV.pdf\#nameddest=tag-02UY}{IV.9.26 (p.~1064)} }
\expandafter\def\csname SRef@night:P3K.54\endcsname{\href{S6_SERIES_GRIND_IV.pdf\#nameddest=tag-02V0}{IV.9.27 (p.~1067)} }
\expandafter\def\csname SRef@night:P3K.56\endcsname{\href{S6_SERIES_GRIND_IV.pdf\#nameddest=tag-02V2}{IV.9.28 (p.~1069)} }
\expandafter\def\csname SRef@night:P3K.57\endcsname{\href{S6_SERIES_GRIND_IV.pdf\#nameddest=tag-02V3}{IV.9.29 (p.~1072)} }
\expandafter\def\csname SRef@night:P3K.10\endcsname{\href{S6_SERIES_GRIND_IV.pdf\#nameddest=tag-02T5}{IV.9.3 (p.~952)} }
\expandafter\def\csname SRef@night:P3K.58\endcsname{\href{S6_SERIES_GRIND_IV.pdf\#nameddest=tag-02V4}{IV.9.30 (p.~1095)} }
\expandafter\def\csname SRef@night:P3K.59\endcsname{\href{S6_SERIES_GRIND_IV.pdf\#nameddest=tag-02V5}{IV.9.31 (p.~1104)} }
\expandafter\def\csname SRef@night:P3K.60\endcsname{\href{S6_SERIES_GRIND_IV.pdf\#nameddest=tag-02V7}{IV.9.32 (p.~1109)} }
\expandafter\def\csname SRef@night:P3K.17\endcsname{\href{S6_SERIES_GRIND_IV.pdf\#nameddest=tag-02TW}{IV.9.4 (p.~961)} }
\expandafter\def\csname SRef@night:P3K.22\endcsname{\href{S6_SERIES_GRIND_IV.pdf\#nameddest=tag-02U2}{IV.9.5 (p.~965)} }
\expandafter\def\csname SRef@night:P3K.23\endcsname{\href{S6_SERIES_GRIND_IV.pdf\#nameddest=tag-02U3}{IV.9.6 (p.~971)} }
\expandafter\def\csname SRef@night:P3K.24\endcsname{\href{S6_SERIES_GRIND_IV.pdf\#nameddest=tag-02U4}{IV.9.7 (p.~975)} }
\expandafter\def\csname SRef@night:P3K.25\endcsname{\href{S6_SERIES_GRIND_IV.pdf\#nameddest=tag-02U5}{IV.9.8 (p.~980)} }
\expandafter\def\csname SRef@night:P3K.27\endcsname{\href{S6_SERIES_GRIND_IV.pdf\#nameddest=tag-02U7}{IV.9.9 (p.~983)} }
\expandafter\def\csname SRef@maps:eq:F.6\endcsname{\href{S6_SERIES_GRIND_V.pdf\#nameddest=tag-02FA}{V.{F.6} (p.~85)} }
\expandafter\def\csname SRef@maps:eq:F.L2\endcsname{\href{S6_SERIES_GRIND_V.pdf\#nameddest=tag-02GV}{V.{F.L2} (p.~174)} }
\expandafter\def\csname SRef@maps:eq:F.L3\endcsname{\href{S6_SERIES_GRIND_V.pdf\#nameddest=tag-02GW}{V.{F.L3} (p.~174)} }
\expandafter\def\csname SRef@maps:eq:F.L4\endcsname{\href{S6_SERIES_GRIND_V.pdf\#nameddest=tag-02GX}{V.{F.L4} (p.~175)} }
\expandafter\def\csname SRef@maps:eq:F.L5\endcsname{\href{S6_SERIES_GRIND_V.pdf\#nameddest=tag-02GY}{V.{F.L5} (p.~175)} }
\expandafter\def\csname SRef@maps:eq:F.L6\endcsname{\href{S6_SERIES_GRIND_V.pdf\#nameddest=tag-02GZ}{V.{F.L6} (p.~175)} }
\expandafter\def\csname SRef@maps:eq:F.L7\endcsname{\href{S6_SERIES_GRIND_V.pdf\#nameddest=tag-02H0}{V.{F.L7} (p.~175)} }
\expandafter\def\csname SRef@maps:eq:F.L8\endcsname{\href{S6_SERIES_GRIND_V.pdf\#nameddest=tag-02H1}{V.{F.L8} (p.~175)} }
\expandafter\def\csname SRef@maps:eq:F.L9\endcsname{\href{S6_SERIES_GRIND_V.pdf\#nameddest=tag-02H2}{V.{F.L9} (p.~176)} }
\expandafter\def\csname SRef@maps:eq:F.L10\endcsname{\href{S6_SERIES_GRIND_V.pdf\#nameddest=tag-02GU}{V.{F.L10} (p.~176)} }
\expandafter\def\csname SRef@maps:eq:F.T1\endcsname{\href{S6_SERIES_GRIND_V.pdf\#nameddest=tag-02I7}{V.{F.T1} (p.~177)} }
\expandafter\def\csname SRef@maps:eq:F.T2\endcsname{\href{S6_SERIES_GRIND_V.pdf\#nameddest=tag-02IC}{V.{F.T2} (p.~177)} }
\expandafter\def\csname SRef@maps:eq:F.T3\endcsname{\href{S6_SERIES_GRIND_V.pdf\#nameddest=tag-02ID}{V.{F.T3} (p.~177)} }
\expandafter\def\csname SRef@maps:eq:F.T4\endcsname{\href{S6_SERIES_GRIND_V.pdf\#nameddest=tag-02IE}{V.{F.T4} (p.~178)} }
\expandafter\def\csname SRef@maps:eq:F.T5\endcsname{\href{S6_SERIES_GRIND_V.pdf\#nameddest=tag-02IF}{V.{F.T5} (p.~178)} }
\expandafter\def\csname SRef@maps:eq:F.T6\endcsname{\href{S6_SERIES_GRIND_V.pdf\#nameddest=tag-02IG}{V.{F.T6} (p.~178)} }
\expandafter\def\csname SRef@maps:eq:F.T7\endcsname{\href{S6_SERIES_GRIND_V.pdf\#nameddest=tag-02IH}{V.{F.T7} (p.~178)} }
\expandafter\def\csname SRef@maps:eq:F.T8\endcsname{\href{S6_SERIES_GRIND_V.pdf\#nameddest=tag-02II}{V.{F.T8} (p.~178)} }
\expandafter\def\csname SRef@maps:eq:F.T9\endcsname{\href{S6_SERIES_GRIND_V.pdf\#nameddest=tag-02IJ}{V.{F.T9} (p.~179)} }
\expandafter\def\csname SRef@maps:eq:F.T10\endcsname{\href{S6_SERIES_GRIND_V.pdf\#nameddest=tag-02I8}{V.{F.T10} (p.~179)} }
\expandafter\def\csname SRef@maps:eq:F.S1\endcsname{\href{S6_SERIES_GRIND_V.pdf\#nameddest=tag-02HX}{V.{F.S1} (p.~180)} }
\expandafter\def\csname SRef@maps:eq:F.S2\endcsname{\href{S6_SERIES_GRIND_V.pdf\#nameddest=tag-02HY}{V.{F.S2} (p.~180)} }
\expandafter\def\csname SRef@maps:eq:F.S3\endcsname{\href{S6_SERIES_GRIND_V.pdf\#nameddest=tag-02HZ}{V.{F.S3} (p.~180)} }
\expandafter\def\csname SRef@maps:eq:F.S4\endcsname{\href{S6_SERIES_GRIND_V.pdf\#nameddest=tag-02I0}{V.{F.S4} (p.~180)} }
\expandafter\def\csname SRef@maps:eq:F.S5\endcsname{\href{S6_SERIES_GRIND_V.pdf\#nameddest=tag-02I1}{V.{F.S5} (p.~180)} }
\expandafter\def\csname SRef@maps:eq:F.S6\endcsname{\href{S6_SERIES_GRIND_V.pdf\#nameddest=tag-02I2}{V.{F.S6} (p.~181)} }
\expandafter\def\csname SRef@maps:eq:F.S7\endcsname{\href{S6_SERIES_GRIND_V.pdf\#nameddest=tag-02I3}{V.{F.S7} (p.~181)} }
\expandafter\def\csname SRef@maps:eq:F.S8\endcsname{\href{S6_SERIES_GRIND_V.pdf\#nameddest=tag-02I4}{V.{F.S8} (p.~181)} }
\expandafter\def\csname SRef@maps:eq:F.S9\endcsname{\href{S6_SERIES_GRIND_V.pdf\#nameddest=tag-02I5}{V.{F.S9} (p.~181)} }
\expandafter\def\csname SRef@maps:eq:F.R10\endcsname{\href{S6_SERIES_GRIND_V.pdf\#nameddest=tag-02HK}{V.{F.R10} (p.~182)} }
\expandafter\def\csname SRef@maps:eq:F.R11\endcsname{\href{S6_SERIES_GRIND_V.pdf\#nameddest=tag-02HL}{V.{F.R11} (p.~182)} }
\expandafter\def\csname SRef@maps:eq:F.R12\endcsname{\href{S6_SERIES_GRIND_V.pdf\#nameddest=tag-02HM}{V.{F.R12} (p.~182)} }
\expandafter\def\csname SRef@maps:eq:F.R13\endcsname{\href{S6_SERIES_GRIND_V.pdf\#nameddest=tag-02HN}{V.{F.R13} (p.~182)} }
\expandafter\def\csname SRef@maps:eq:F.R14\endcsname{\href{S6_SERIES_GRIND_V.pdf\#nameddest=tag-02HO}{V.{F.R14} (p.~182)} }
\expandafter\def\csname SRef@maps:eq:F.X1\endcsname{\href{S6_SERIES_GRIND_V.pdf\#nameddest=tag-02IV}{V.{F.X1} (p.~183)} }
\expandafter\def\csname SRef@maps:eq:F.X2\endcsname{\href{S6_SERIES_GRIND_V.pdf\#nameddest=tag-02IW}{V.{F.X2} (p.~183)} }
\expandafter\def\csname SRef@maps:eq:F.X3\endcsname{\href{S6_SERIES_GRIND_V.pdf\#nameddest=tag-02IX}{V.{F.X3} (p.~184)} }
\expandafter\def\csname SRef@maps:eq:F.X4\endcsname{\href{S6_SERIES_GRIND_V.pdf\#nameddest=tag-02IY}{V.{F.X4} (p.~184)} }
\expandafter\def\csname SRef@maps:eq:F.X5\endcsname{\href{S6_SERIES_GRIND_V.pdf\#nameddest=tag-02IZ}{V.{F.X5} (p.~184)} }
\expandafter\def\csname SRef@maps:eq:F.T11\endcsname{\href{S6_SERIES_GRIND_V.pdf\#nameddest=tag-02I9}{V.{F.T11} (p.~185)} }
\expandafter\def\csname SRef@maps:eq:F.T12\endcsname{\href{S6_SERIES_GRIND_V.pdf\#nameddest=tag-02IA}{V.{F.T12} (p.~185)} }
\expandafter\def\csname SRef@maps:eq:F.T13\endcsname{\href{S6_SERIES_GRIND_V.pdf\#nameddest=tag-02IB}{V.{F.T13} (p.~185)} }
\expandafter\def\csname SRef@maps:eq:F.HC1\endcsname{\href{S6_SERIES_GRIND_V.pdf\#nameddest=tag-02G9}{V.{F.HC1} (p.~186)} }
\expandafter\def\csname SRef@maps:eq:F.HC2\endcsname{\href{S6_SERIES_GRIND_V.pdf\#nameddest=tag-02GE}{V.{F.HC2} (p.~186)} }
\expandafter\def\csname SRef@maps:eq:F.HC3\endcsname{\href{S6_SERIES_GRIND_V.pdf\#nameddest=tag-02GF}{V.{F.HC3} (p.~186)} }
\expandafter\def\csname SRef@maps:eq:F.HC4\endcsname{\href{S6_SERIES_GRIND_V.pdf\#nameddest=tag-02GG}{V.{F.HC4} (p.~187)} }
\expandafter\def\csname SRef@maps:eq:F.HC5\endcsname{\href{S6_SERIES_GRIND_V.pdf\#nameddest=tag-02GH}{V.{F.HC5} (p.~187)} }
\expandafter\def\csname SRef@maps:eq:F.HC6\endcsname{\href{S6_SERIES_GRIND_V.pdf\#nameddest=tag-02GI}{V.{F.HC6} (p.~187)} }
\expandafter\def\csname SRef@maps:eq:F.HC7\endcsname{\href{S6_SERIES_GRIND_V.pdf\#nameddest=tag-02GJ}{V.{F.HC7} (p.~188)} }
\expandafter\def\csname SRef@maps:eq:F.HC8\endcsname{\href{S6_SERIES_GRIND_V.pdf\#nameddest=tag-02GK}{V.{F.HC8} (p.~188)} }
\expandafter\def\csname SRef@maps:eq:F.HC9\endcsname{\href{S6_SERIES_GRIND_V.pdf\#nameddest=tag-02GL}{V.{F.HC9} (p.~189)} }
\expandafter\def\csname SRef@maps:eq:F.HC10\endcsname{\href{S6_SERIES_GRIND_V.pdf\#nameddest=tag-02GA}{V.{F.HC10} (p.~189)} }
\expandafter\def\csname SRef@maps:eq:F.HC11\endcsname{\href{S6_SERIES_GRIND_V.pdf\#nameddest=tag-02GB}{V.{F.HC11} (p.~189)} }
\expandafter\def\csname SRef@maps:eq:F.HC12\endcsname{\href{S6_SERIES_GRIND_V.pdf\#nameddest=tag-02GC}{V.{F.HC12} (p.~189)} }
\expandafter\def\csname SRef@maps:eq:F.HC13\endcsname{\href{S6_SERIES_GRIND_V.pdf\#nameddest=tag-02GD}{V.{F.HC13} (p.~190)} }
\expandafter\def\csname SRef@maps:eq:F.G1\endcsname{\href{S6_SERIES_GRIND_V.pdf\#nameddest=tag-02FO}{V.{F.G1} (p.~191)} }
\expandafter\def\csname SRef@maps:eq:F.G2\endcsname{\href{S6_SERIES_GRIND_V.pdf\#nameddest=tag-02FR}{V.{F.G2} (p.~191)} }
\expandafter\def\csname SRef@maps:eq:F.G3\endcsname{\href{S6_SERIES_GRIND_V.pdf\#nameddest=tag-02FS}{V.{F.G3} (p.~191)} }
\expandafter\def\csname SRef@maps:eq:F.G4\endcsname{\href{S6_SERIES_GRIND_V.pdf\#nameddest=tag-02FT}{V.{F.G4} (p.~191)} }
\expandafter\def\csname SRef@maps:eq:F.G5\endcsname{\href{S6_SERIES_GRIND_V.pdf\#nameddest=tag-02FU}{V.{F.G5} (p.~191)} }
\expandafter\def\csname SRef@maps:eq:F.G6\endcsname{\href{S6_SERIES_GRIND_V.pdf\#nameddest=tag-02FV}{V.{F.G6} (p.~191)} }
\expandafter\def\csname SRef@maps:eq:F.G7\endcsname{\href{S6_SERIES_GRIND_V.pdf\#nameddest=tag-02FW}{V.{F.G7} (p.~192)} }
\expandafter\def\csname SRef@maps:eq:F.G8\endcsname{\href{S6_SERIES_GRIND_V.pdf\#nameddest=tag-02FX}{V.{F.G8} (p.~192)} }
\expandafter\def\csname SRef@maps:eq:F.G9\endcsname{\href{S6_SERIES_GRIND_V.pdf\#nameddest=tag-02FY}{V.{F.G9} (p.~192)} }
\expandafter\def\csname SRef@maps:eq:F.G10\endcsname{\href{S6_SERIES_GRIND_V.pdf\#nameddest=tag-02FP}{V.{F.G10} (p.~192)} }
\expandafter\def\csname SRef@maps:eq:F.G11\endcsname{\href{S6_SERIES_GRIND_V.pdf\#nameddest=tag-02FQ}{V.{F.G11} (p.~194)} }
\expandafter\def\csname SRef@maps:eq:F.LT1\endcsname{\href{S6_SERIES_GRIND_V.pdf\#nameddest=tag-02H3}{V.{F.LT1} (p.~195)} }
\expandafter\def\csname SRef@maps:eq:F.LT2\endcsname{\href{S6_SERIES_GRIND_V.pdf\#nameddest=tag-02H4}{V.{F.LT2} (p.~195)} }
\expandafter\def\csname SRef@maps:eq:F.LT3\endcsname{\href{S6_SERIES_GRIND_V.pdf\#nameddest=tag-02H5}{V.{F.LT3} (p.~195)} }
\expandafter\def\csname SRef@maps:eq:F.LT4\endcsname{\href{S6_SERIES_GRIND_V.pdf\#nameddest=tag-02H6}{V.{F.LT4} (p.~195)} }
\expandafter\def\csname SRef@maps:eq:F.LT5\endcsname{\href{S6_SERIES_GRIND_V.pdf\#nameddest=tag-02H7}{V.{F.LT5} (p.~195)} }
\expandafter\def\csname SRef@maps:eq:F.LT6\endcsname{\href{S6_SERIES_GRIND_V.pdf\#nameddest=tag-02H8}{V.{F.LT6} (p.~196)} }
\expandafter\def\csname SRef@maps:eq:F.LT7\endcsname{\href{S6_SERIES_GRIND_V.pdf\#nameddest=tag-02H9}{V.{F.LT7} (p.~196)} }
\expandafter\def\csname SRef@maps:eq:F.LT8\endcsname{\href{S6_SERIES_GRIND_V.pdf\#nameddest=tag-02HA}{V.{F.LT8} (p.~196)} }
\expandafter\def\csname SRef@maps:eq:F.LT9\endcsname{\href{S6_SERIES_GRIND_V.pdf\#nameddest=tag-02HB}{V.{F.LT9} (p.~197)} }
\expandafter\def\csname SRef@maps:eq:F.CA1\endcsname{\href{S6_SERIES_GRIND_V.pdf\#nameddest=tag-02FE}{V.{F.CA1} (p.~197)} }
\expandafter\def\csname SRef@maps:eq:F.CA2\endcsname{\href{S6_SERIES_GRIND_V.pdf\#nameddest=tag-02FF}{V.{F.CA2} (p.~198)} }
\expandafter\def\csname SRef@maps:eq:F.CA3\endcsname{\href{S6_SERIES_GRIND_V.pdf\#nameddest=tag-02FG}{V.{F.CA3} (p.~198)} }
\expandafter\def\csname SRef@maps:eq:F.CA4\endcsname{\href{S6_SERIES_GRIND_V.pdf\#nameddest=tag-02FH}{V.{F.CA4} (p.~198)} }
\expandafter\def\csname SRef@maps:eq:F.CA5\endcsname{\href{S6_SERIES_GRIND_V.pdf\#nameddest=tag-02FI}{V.{F.CA5} (p.~198)} }
\expandafter\def\csname SRef@maps:eq:F.CK1\endcsname{\href{S6_SERIES_GRIND_V.pdf\#nameddest=tag-02FJ}{V.{F.CK1} (p.~200)} }
\expandafter\def\csname SRef@maps:eq:F.CK2\endcsname{\href{S6_SERIES_GRIND_V.pdf\#nameddest=tag-02FK}{V.{F.CK2} (p.~201)} }
\expandafter\def\csname SRef@maps:eq:F.GC1\endcsname{\href{S6_SERIES_GRIND_V.pdf\#nameddest=tag-02FZ}{V.{F.GC1} (p.~202)} }
\expandafter\def\csname SRef@maps:eq:F.GC2\endcsname{\href{S6_SERIES_GRIND_V.pdf\#nameddest=tag-02G0}{V.{F.GC2} (p.~203)} }
\expandafter\def\csname SRef@maps:eq:F.GC3\endcsname{\href{S6_SERIES_GRIND_V.pdf\#nameddest=tag-02G1}{V.{F.GC3} (p.~203)} }
\expandafter\def\csname SRef@maps:eq:F.GC4\endcsname{\href{S6_SERIES_GRIND_V.pdf\#nameddest=tag-02G2}{V.{F.GC4} (p.~203)} }
\expandafter\def\csname SRef@maps:eq:F.GC5\endcsname{\href{S6_SERIES_GRIND_V.pdf\#nameddest=tag-02G3}{V.{F.GC5} (p.~204)} }
\expandafter\def\csname SRef@maps:eq:F.GC5a\endcsname{\href{S6_SERIES_GRIND_V.pdf\#nameddest=tag-02G4}{V.{F.GC5a} (p.~204)} }
\expandafter\def\csname SRef@maps:eq:F.GC6\endcsname{\href{S6_SERIES_GRIND_V.pdf\#nameddest=tag-02G5}{V.{F.GC6} (p.~204)} }
\expandafter\def\csname SRef@maps:eq:F.GC7\endcsname{\href{S6_SERIES_GRIND_V.pdf\#nameddest=tag-02G6}{V.{F.GC7} (p.~205)} }
\expandafter\def\csname SRef@maps:eq:F.GC8\endcsname{\href{S6_SERIES_GRIND_V.pdf\#nameddest=tag-02G7}{V.{F.GC8} (p.~205)} }
\expandafter\def\csname SRef@maps:eq:F.GC9\endcsname{\href{S6_SERIES_GRIND_V.pdf\#nameddest=tag-02G8}{V.{F.GC9} (p.~205)} }
\expandafter\def\csname SRef@maps:eq:F.IC1\endcsname{\href{S6_SERIES_GRIND_V.pdf\#nameddest=tag-02GM}{V.{F.IC1} (p.~207)} }
\expandafter\def\csname SRef@maps:eq:F.IC2\endcsname{\href{S6_SERIES_GRIND_V.pdf\#nameddest=tag-02GN}{V.{F.IC2} (p.~207)} }
\expandafter\def\csname SRef@maps:eq:F.IC3\endcsname{\href{S6_SERIES_GRIND_V.pdf\#nameddest=tag-02GO}{V.{F.IC3} (p.~207)} }
\expandafter\def\csname SRef@maps:eq:F.CN1\endcsname{\href{S6_SERIES_GRIND_V.pdf\#nameddest=tag-02FL}{V.{F.CN1} (p.~209)} }
\expandafter\def\csname SRef@maps:eq:F.CN2\endcsname{\href{S6_SERIES_GRIND_V.pdf\#nameddest=tag-02FM}{V.{F.CN2} (p.~209)} }
\expandafter\def\csname SRef@maps:eq:F.CN3\endcsname{\href{S6_SERIES_GRIND_V.pdf\#nameddest=tag-02FN}{V.{F.CN3} (p.~210)} }
\expandafter\def\csname SRef@maps:eq:F.IV1\endcsname{\href{S6_SERIES_GRIND_V.pdf\#nameddest=tag-02GP}{V.{F.IV1} (p.~210)} }
\expandafter\def\csname SRef@maps:eq:F.IV2\endcsname{\href{S6_SERIES_GRIND_V.pdf\#nameddest=tag-02GQ}{V.{F.IV2} (p.~210)} }
\expandafter\def\csname SRef@maps:eq:F.IV3\endcsname{\href{S6_SERIES_GRIND_V.pdf\#nameddest=tag-02GR}{V.{F.IV3} (p.~211)} }
\expandafter\def\csname SRef@maps:eq:F.IV4\endcsname{\href{S6_SERIES_GRIND_V.pdf\#nameddest=tag-02GS}{V.{F.IV4} (p.~211)} }
\expandafter\def\csname SRef@maps:eq:F.SY1\endcsname{\href{S6_SERIES_GRIND_V.pdf\#nameddest=tag-02I6}{V.{F.SY1} (p.~212)} }
\expandafter\def\csname SRef@maps:eq:F.O1\endcsname{\href{S6_SERIES_GRIND_V.pdf\#nameddest=tag-02HC}{V.{F.O1} (p.~216)} }
\expandafter\def\csname SRef@maps:eq:F.O2\endcsname{\href{S6_SERIES_GRIND_V.pdf\#nameddest=tag-02HD}{V.{F.O2} (p.~216)} }
\expandafter\def\csname SRef@maps:eq:F.O3\endcsname{\href{S6_SERIES_GRIND_V.pdf\#nameddest=tag-02HE}{V.{F.O3} (p.~216)} }
\expandafter\def\csname SRef@maps:eq:F.O4\endcsname{\href{S6_SERIES_GRIND_V.pdf\#nameddest=tag-02HF}{V.{F.O4} (p.~217)} }
\expandafter\def\csname SRef@maps:eq:F.O5\endcsname{\href{S6_SERIES_GRIND_V.pdf\#nameddest=tag-02HG}{V.{F.O5} (p.~218)} }
\expandafter\def\csname SRef@maps:eq:F.O6\endcsname{\href{S6_SERIES_GRIND_V.pdf\#nameddest=tag-02HH}{V.{F.O6} (p.~219)} }
\expandafter\def\csname SRef@maps:eq:F.O7\endcsname{\href{S6_SERIES_GRIND_V.pdf\#nameddest=tag-02HI}{V.{F.O7} (p.~220)} }
\expandafter\def\csname SRef@maps:eq:F.7\endcsname{\href{S6_SERIES_GRIND_V.pdf\#nameddest=tag-02FB}{V.{F.7} (p.~104)} }
\expandafter\def\csname SRef@maps:eq:F.8\endcsname{\href{S6_SERIES_GRIND_V.pdf\#nameddest=tag-02FC}{V.{F.8} (p.~104)} }
\expandafter\def\csname SRef@maps:eq:F.9\endcsname{\href{S6_SERIES_GRIND_V.pdf\#nameddest=tag-02FD}{V.{F.9} (p.~105)} }
\expandafter\def\csname SRef@maps:eq:F.10\endcsname{\href{S6_SERIES_GRIND_V.pdf\#nameddest=tag-02EG}{V.{F.10} (p.~105)} }
\expandafter\def\csname SRef@maps:eq:F.11\endcsname{\href{S6_SERIES_GRIND_V.pdf\#nameddest=tag-02EH}{V.{F.11} (p.~106)} }
\expandafter\def\csname SRef@maps:eq:F.12\endcsname{\href{S6_SERIES_GRIND_V.pdf\#nameddest=tag-02EI}{V.{F.12} (p.~106)} }
\expandafter\def\csname SRef@maps:eq:F.13\endcsname{\href{S6_SERIES_GRIND_V.pdf\#nameddest=tag-02EJ}{V.{F.13} (p.~110)} }
\expandafter\def\csname SRef@maps:eq:F.14\endcsname{\href{S6_SERIES_GRIND_V.pdf\#nameddest=tag-02EK}{V.{F.14} (p.~110)} }
\expandafter\def\csname SRef@maps:eq:F.15\endcsname{\href{S6_SERIES_GRIND_V.pdf\#nameddest=tag-02EL}{V.{F.15} (p.~110)} }
\expandafter\def\csname SRef@maps:eq:F.16\endcsname{\href{S6_SERIES_GRIND_V.pdf\#nameddest=tag-02EM}{V.{F.16} (p.~110)} }
\expandafter\def\csname SRef@maps:eq:F.17\endcsname{\href{S6_SERIES_GRIND_V.pdf\#nameddest=tag-02EN}{V.{F.17} (p.~110)} }
\expandafter\def\csname SRef@maps:eq:F.18\endcsname{\href{S6_SERIES_GRIND_V.pdf\#nameddest=tag-02EO}{V.{F.18} (p.~110)} }
\expandafter\def\csname SRef@maps:eq:F.19\endcsname{\href{S6_SERIES_GRIND_V.pdf\#nameddest=tag-02EP}{V.{F.19} (p.~111)} }
\expandafter\def\csname SRef@maps:eq:F.20\endcsname{\href{S6_SERIES_GRIND_V.pdf\#nameddest=tag-02ER}{V.{F.20} (p.~111)} }
\expandafter\def\csname SRef@maps:eq:F.21\endcsname{\href{S6_SERIES_GRIND_V.pdf\#nameddest=tag-02ES}{V.{F.21} (p.~111)} }
\expandafter\def\csname SRef@maps:eq:F.22\endcsname{\href{S6_SERIES_GRIND_V.pdf\#nameddest=tag-02ET}{V.{F.22} (p.~111)} }
\expandafter\def\csname SRef@maps:eq:F.23\endcsname{\href{S6_SERIES_GRIND_V.pdf\#nameddest=tag-02EU}{V.{F.23} (p.~111)} }
\expandafter\def\csname SRef@maps:eq:F.24\endcsname{\href{S6_SERIES_GRIND_V.pdf\#nameddest=tag-02EV}{V.{F.24} (p.~115)} }
\expandafter\def\csname SRef@maps:eq:F.25\endcsname{\href{S6_SERIES_GRIND_V.pdf\#nameddest=tag-02EW}{V.{F.25} (p.~117)} }
\expandafter\def\csname SRef@maps:eq:F.3\endcsname{\href{S6_SERIES_GRIND_V.pdf\#nameddest=tag-02EZ}{V.{F.3} (p.~73)} }
\expandafter\def\csname SRef@maps:eq:F.27\endcsname{\href{S6_SERIES_GRIND_V.pdf\#nameddest=tag-02EX}{V.{F.27} (p.~147)} }
\expandafter\def\csname SRef@maps:eq:F.28\endcsname{\href{S6_SERIES_GRIND_V.pdf\#nameddest=tag-02EY}{V.{F.28} (p.~147)} }
\expandafter\def\csname SRef@maps:eq:F.4\endcsname{\href{S6_SERIES_GRIND_V.pdf\#nameddest=tag-02F7}{V.{F.4} (p.~75)} }
\expandafter\def\csname SRef@maps:eq:F.30\endcsname{\href{S6_SERIES_GRIND_V.pdf\#nameddest=tag-02F0}{V.{F.30} (p.~150)} }
\expandafter\def\csname SRef@maps:eq:F.31\endcsname{\href{S6_SERIES_GRIND_V.pdf\#nameddest=tag-02F1}{V.{F.31} (p.~150)} }
\expandafter\def\csname SRef@maps:eq:F.32\endcsname{\href{S6_SERIES_GRIND_V.pdf\#nameddest=tag-02F2}{V.{F.32} (p.~151)} }
\expandafter\def\csname SRef@maps:eq:F.33\endcsname{\href{S6_SERIES_GRIND_V.pdf\#nameddest=tag-02F3}{V.{F.33} (p.~152)} }
\expandafter\def\csname SRef@maps:eq:F.34\endcsname{\href{S6_SERIES_GRIND_V.pdf\#nameddest=tag-02F4}{V.{F.34} (p.~152)} }
\expandafter\def\csname SRef@maps:eq:F.36\endcsname{\href{S6_SERIES_GRIND_V.pdf\#nameddest=tag-02F5}{V.{F.36} (p.~155)} }
\expandafter\def\csname SRef@maps:eq:F.39\endcsname{\href{S6_SERIES_GRIND_V.pdf\#nameddest=tag-02F6}{V.{F.39} (p.~161)} }
\expandafter\def\csname SRef@maps:eq:F.40\endcsname{\href{S6_SERIES_GRIND_V.pdf\#nameddest=tag-02F8}{V.{F.40} (p.~162)} }
\expandafter\def\csname SRef@maps:eq:F.U1\endcsname{\href{S6_SERIES_GRIND_V.pdf\#nameddest=tag-02IK}{V.{F.U1} (p.~163)} }
\expandafter\def\csname SRef@maps:eq:F.5\endcsname{\href{S6_SERIES_GRIND_V.pdf\#nameddest=tag-02F9}{V.{F.5} (p.~76)} }
\expandafter\def\csname SRef@maps:eq:F.U2\endcsname{\href{S6_SERIES_GRIND_V.pdf\#nameddest=tag-02IO}{V.{F.U2} (p.~163)} }
\expandafter\def\csname SRef@maps:eq:F.U3\endcsname{\href{S6_SERIES_GRIND_V.pdf\#nameddest=tag-02IP}{V.{F.U3} (p.~163)} }
\expandafter\def\csname SRef@maps:eq:F.U4\endcsname{\href{S6_SERIES_GRIND_V.pdf\#nameddest=tag-02IQ}{V.{F.U4} (p.~164)} }
\expandafter\def\csname SRef@maps:eq:F.U5\endcsname{\href{S6_SERIES_GRIND_V.pdf\#nameddest=tag-02IR}{V.{F.U5} (p.~164)} }
\expandafter\def\csname SRef@maps:eq:F.U6\endcsname{\href{S6_SERIES_GRIND_V.pdf\#nameddest=tag-02IS}{V.{F.U6} (p.~164)} }
\expandafter\def\csname SRef@maps:eq:F.U7\endcsname{\href{S6_SERIES_GRIND_V.pdf\#nameddest=tag-02IT}{V.{F.U7} (p.~166)} }
\expandafter\def\csname SRef@maps:eq:F.U9\endcsname{\href{S6_SERIES_GRIND_V.pdf\#nameddest=tag-02IU}{V.{F.U9} (p.~166)} }
\expandafter\def\csname SRef@maps:eq:F.U10\endcsname{\href{S6_SERIES_GRIND_V.pdf\#nameddest=tag-02IL}{V.{F.U10} (p.~166)} }
\expandafter\def\csname SRef@maps:eq:F.U11\endcsname{\href{S6_SERIES_GRIND_V.pdf\#nameddest=tag-02IM}{V.{F.U11} (p.~167)} }
\expandafter\def\csname SRef@maps:eq:F.U12\endcsname{\href{S6_SERIES_GRIND_V.pdf\#nameddest=tag-02IN}{V.{F.U12} (p.~169)} }
\expandafter\def\csname SRef@maps:eq:FT.1\endcsname{\href{S6_SERIES_GRIND_V.pdf\#nameddest=tag-02J0}{V.{FT.1} (p.~81)} }
\expandafter\def\csname SRef@maps:eq:F.R1\endcsname{\href{S6_SERIES_GRIND_V.pdf\#nameddest=tag-02HJ}{V.{F.R1} (p.~170)} }
\expandafter\def\csname SRef@maps:eq:F.R2\endcsname{\href{S6_SERIES_GRIND_V.pdf\#nameddest=tag-02HP}{V.{F.R2} (p.~170)} }
\expandafter\def\csname SRef@maps:eq:F.R3\endcsname{\href{S6_SERIES_GRIND_V.pdf\#nameddest=tag-02HQ}{V.{F.R3} (p.~171)} }
\expandafter\def\csname SRef@maps:eq:F.R4\endcsname{\href{S6_SERIES_GRIND_V.pdf\#nameddest=tag-02HR}{V.{F.R4} (p.~171)} }
\expandafter\def\csname SRef@maps:eq:F.R5\endcsname{\href{S6_SERIES_GRIND_V.pdf\#nameddest=tag-02HS}{V.{F.R5} (p.~171)} }
\expandafter\def\csname SRef@maps:eq:F.R6\endcsname{\href{S6_SERIES_GRIND_V.pdf\#nameddest=tag-02HT}{V.{F.R6} (p.~171)} }
\expandafter\def\csname SRef@maps:eq:F.R7\endcsname{\href{S6_SERIES_GRIND_V.pdf\#nameddest=tag-02HU}{V.{F.R7} (p.~172)} }
\expandafter\def\csname SRef@maps:eq:F.R8\endcsname{\href{S6_SERIES_GRIND_V.pdf\#nameddest=tag-02HV}{V.{F.R8} (p.~172)} }
\expandafter\def\csname SRef@maps:eq:F.R9\endcsname{\href{S6_SERIES_GRIND_V.pdf\#nameddest=tag-02HW}{V.{F.R9} (p.~172)} }
\expandafter\def\csname SRef@maps:eq:F.L1\endcsname{\href{S6_SERIES_GRIND_V.pdf\#nameddest=tag-02GT}{V.{F.L1} (p.~174)} }
\expandafter\def\csname SRef@maps:mapsF:repair:cyclic-inclusion\endcsname{\href{S6_SERIES_GRIND_V.pdf\#nameddest=tag-02J1}{V.{{$\langle 1\rangle $6.}} (p.~199)} }
\expandafter\def\csname SRef@maps:mapsF:repair:log-cylinder-scope\endcsname{\href{S6_SERIES_GRIND_V.pdf\#nameddest=tag-02J2}{V.{{$\langle 1\rangle $5.}} (p.~201)} }
\expandafter\def\csname SRef@sec:fourfold\endcsname{\href{S6_SERIES_GRIND_V.pdf\#nameddest=tag-034O}{V.5 (p.~80)} }
\expandafter\def\csname SRef@maps:src:F.2.21\endcsname{\href{S6_SERIES_GRIND_V.pdf\#nameddest=tag-02NB}{V.11.15 (p.~173)} }
\expandafter\def\csname SRef@maps:src:F.2.30\endcsname{\href{S6_SERIES_GRIND_V.pdf\#nameddest=tag-02NL}{V.11.24 (p.~190)} }
\expandafter\def\csname SRef@maps:eq:F.1\endcsname{\href{S6_SERIES_GRIND_V.pdf\#nameddest=tag-02EF}{V.1 (p.~72)} }
\expandafter\def\csname SRef@maps:eq:F.2\endcsname{\href{S6_SERIES_GRIND_V.pdf\#nameddest=tag-02EQ}{V.2 (p.~72)} }
\expandafter\def\csname SRef@eq:P1Adegree\endcsname{\href{S6_SERIES_GRIND_V.pdf\#nameddest=tag-01O3}{V.3 (p.~134)} }
\expandafter\def\csname SRef@eq:P1Alogcotangent\endcsname{\href{S6_SERIES_GRIND_V.pdf\#nameddest=tag-01O4}{V.4 (p.~135)} }
\expandafter\def\csname SRef@maps:src:F.3.10\endcsname{\href{S6_SERIES_GRIND_V.pdf\#nameddest=tag-02O2}{V.12.5 (p.~218)} }
\expandafter\def\csname SRef@mapsaudit:order-four-torsion-transgression\endcsname{\href{S6_SERIES_GRIND_V.pdf\#nameddest=tag-02QF}{V.5.4 (p.~81)} }
\expandafter\def\csname SRef@maps:src:F.2.17\endcsname{\href{S6_SERIES_GRIND_V.pdf\#nameddest=tag-02N6}{V.11.11 (p.~166)} }
\expandafter\def\csname SRef@maps:src:F.2.8\endcsname{\href{S6_SERIES_GRIND_V.pdf\#nameddest=tag-02NZ}{V.11.2 (p.~147)} }
\expandafter\def\csname SRef@maps:src:F.2.9\endcsname{\href{S6_SERIES_GRIND_V.pdf\#nameddest=tag-02O0}{V.11.3 (p.~148)} }
\expandafter\def\csname SRef@maps:src:F.2.11\endcsname{\href{S6_SERIES_GRIND_V.pdf\#nameddest=tag-02N0}{V.11.5 (p.~151)} }
\expandafter\def\csname SRef@maps:src:F.2.15\endcsname{\href{S6_SERIES_GRIND_V.pdf\#nameddest=tag-02N4}{V.11.9 (p.~162)} }
\expandafter\def\csname SRef@maps:src:F.3.3\endcsname{\href{S6_SERIES_GRIND_V.pdf\#nameddest=tag-02O5}{V.12.1 (p.~215)} }
\expandafter\def\csname SRef@maps:src:F.1.4\endcsname{\href{S6_SERIES_GRIND_V.pdf\#nameddest=tag-02MX}{V.4.2 (p.~72)} }
\expandafter\def\csname SRef@maps:src:F.2.1\endcsname{\href{S6_SERIES_GRIND_V.pdf\#nameddest=tag-02MY}{V.4.3 (p.~73)} }
\expandafter\def\csname SRef@maps:src:F.2.2\endcsname{\href{S6_SERIES_GRIND_V.pdf\#nameddest=tag-02N9}{V.4.4 (p.~73)} }
\expandafter\def\csname SRef@maps:src:F.2.3\endcsname{\href{S6_SERIES_GRIND_V.pdf\#nameddest=tag-02NK}{V.4.5 (p.~74)} }
\expandafter\def\csname SRef@maps:src:F.3.2\endcsname{\href{S6_SERIES_GRIND_V.pdf\#nameddest=tag-02O4}{V.4.7 (p.~75)} }
\expandafter\def\csname SRef@maps:src:F.2.4\endcsname{\href{S6_SERIES_GRIND_V.pdf\#nameddest=tag-02NV}{V.4.8 (p.~75)} }
\expandafter\def\csname SRef@sec:V-known-inputs\endcsname{\href{S6_SERIES_GRIND_V.pdf\#nameddest=tag-034M}{V.3 (p.~72)} }
\expandafter\def\csname SRef@maps:src:F.2.7\endcsname{\href{S6_SERIES_GRIND_V.pdf\#nameddest=tag-02NY}{V.11.1 (p.~146)} }
\expandafter\def\csname SRef@maps:src:F.2.16\endcsname{\href{S6_SERIES_GRIND_V.pdf\#nameddest=tag-02N5}{V.11.10 (p.~163)} }
\expandafter\def\csname SRef@maps:src:F.2.18\endcsname{\href{S6_SERIES_GRIND_V.pdf\#nameddest=tag-02N7}{V.11.12 (p.~167)} }
\expandafter\def\csname SRef@maps:src:F.2.19\endcsname{\href{S6_SERIES_GRIND_V.pdf\#nameddest=tag-02N8}{V.11.13 (p.~169)} }
\expandafter\def\csname SRef@maps:src:F.2.20\endcsname{\href{S6_SERIES_GRIND_V.pdf\#nameddest=tag-02NA}{V.11.14 (p.~170)} }
\expandafter\def\csname SRef@maps:src:F.2.22\endcsname{\href{S6_SERIES_GRIND_V.pdf\#nameddest=tag-02NC}{V.11.16 (p.~174)} }
\expandafter\def\csname SRef@maps:src:F.2.23\endcsname{\href{S6_SERIES_GRIND_V.pdf\#nameddest=tag-02ND}{V.11.17 (p.~176)} }
\expandafter\def\csname SRef@maps:src:F.2.24\endcsname{\href{S6_SERIES_GRIND_V.pdf\#nameddest=tag-02NE}{V.11.18 (p.~177)} }
\expandafter\def\csname SRef@maps:src:F.2.25\endcsname{\href{S6_SERIES_GRIND_V.pdf\#nameddest=tag-02NF}{V.11.19 (p.~179)} }
\expandafter\def\csname SRef@maps:src:F.2.26\endcsname{\href{S6_SERIES_GRIND_V.pdf\#nameddest=tag-02NG}{V.11.20 (p.~181)} }
\expandafter\def\csname SRef@maps:src:F.2.27\endcsname{\href{S6_SERIES_GRIND_V.pdf\#nameddest=tag-02NH}{V.11.21 (p.~183)} }
\expandafter\def\csname SRef@maps:src:F.2.28\endcsname{\href{S6_SERIES_GRIND_V.pdf\#nameddest=tag-02NI}{V.11.22 (p.~185)} }
\expandafter\def\csname SRef@maps:src:F.2.29\endcsname{\href{S6_SERIES_GRIND_V.pdf\#nameddest=tag-02NJ}{V.11.23 (p.~186)} }
\expandafter\def\csname SRef@maps:src:F.2.31\endcsname{\href{S6_SERIES_GRIND_V.pdf\#nameddest=tag-02NM}{V.11.25 (p.~191)} }
\expandafter\def\csname SRef@maps:src:F.2.32\endcsname{\href{S6_SERIES_GRIND_V.pdf\#nameddest=tag-02NN}{V.11.26 (p.~195)} }
\expandafter\def\csname SRef@maps:src:F.2.33\endcsname{\href{S6_SERIES_GRIND_V.pdf\#nameddest=tag-02NO}{V.11.27 (p.~197)} }
\expandafter\def\csname SRef@maps:src:F.2.34\endcsname{\href{S6_SERIES_GRIND_V.pdf\#nameddest=tag-02NP}{V.11.28 (p.~200)} }
\expandafter\def\csname SRef@maps:src:F.2.35\endcsname{\href{S6_SERIES_GRIND_V.pdf\#nameddest=tag-02NQ}{V.11.29 (p.~202)} }
\expandafter\def\csname SRef@maps:src:F.2.36\endcsname{\href{S6_SERIES_GRIND_V.pdf\#nameddest=tag-02NR}{V.11.30 (p.~206)} }
\expandafter\def\csname SRef@maps:src:F.2.37\endcsname{\href{S6_SERIES_GRIND_V.pdf\#nameddest=tag-02NS}{V.11.31 (p.~208)} }
\expandafter\def\csname SRef@maps:src:F.2.38\endcsname{\href{S6_SERIES_GRIND_V.pdf\#nameddest=tag-02NT}{V.11.32 (p.~210)} }
\expandafter\def\csname SRef@maps:src:F.2.39\endcsname{\href{S6_SERIES_GRIND_V.pdf\#nameddest=tag-02NU}{V.11.33 (p.~212)} }
\expandafter\def\csname SRef@night:P1B.1\endcsname{\href{S6_SERIES_GRIND_V.pdf\#nameddest=tag-02R7}{V.11.34 (p.~213)} }
\expandafter\def\csname SRef@night:P1B.2\endcsname{\href{S6_SERIES_GRIND_V.pdf\#nameddest=tag-02R8}{V.11.35 (p.~214)} }
\expandafter\def\csname SRef@maps:src:F.2.10\endcsname{\href{S6_SERIES_GRIND_V.pdf\#nameddest=tag-02MZ}{V.11.4 (p.~150)} }
\expandafter\def\csname SRef@maps:src:F.2.12\endcsname{\href{S6_SERIES_GRIND_V.pdf\#nameddest=tag-02N1}{V.11.6 (p.~152)} }
\expandafter\def\csname SRef@maps:src:F.2.13\endcsname{\href{S6_SERIES_GRIND_V.pdf\#nameddest=tag-02N2}{V.11.7 (p.~155)} }
\expandafter\def\csname SRef@maps:src:F.2.14\endcsname{\href{S6_SERIES_GRIND_V.pdf\#nameddest=tag-02N3}{V.11.8 (p.~161)} }
\expandafter\def\csname SRef@maps:src:F.3.4\endcsname{\href{S6_SERIES_GRIND_V.pdf\#nameddest=tag-02O6}{V.12.2 (p.~215)} }
\expandafter\def\csname SRef@maps:src:F.3.8\endcsname{\href{S6_SERIES_GRIND_V.pdf\#nameddest=tag-02OA}{V.12.3 (p.~216)} }
\expandafter\def\csname SRef@maps:src:F.3.9\endcsname{\href{S6_SERIES_GRIND_V.pdf\#nameddest=tag-02OB}{V.12.4 (p.~217)} }
\expandafter\def\csname SRef@maps:src:F.3.11\endcsname{\href{S6_SERIES_GRIND_V.pdf\#nameddest=tag-02O3}{V.12.6 (p.~219)} }
\expandafter\def\csname SRef@maps:src:F.1.1\endcsname{\href{S6_SERIES_GRIND_V.pdf\#nameddest=tag-02MU}{V.4.1 (p.~72)} }
\expandafter\def\csname SRef@maps:src:F.3.1\endcsname{\href{S6_SERIES_GRIND_V.pdf\#nameddest=tag-02O1}{V.4.6 (p.~74)} }
\expandafter\def\csname SRef@maps:src:F.2.5\endcsname{\href{S6_SERIES_GRIND_V.pdf\#nameddest=tag-02NW}{V.5.1 (p.~76)} }
\expandafter\def\csname SRef@night:P1.1\endcsname{\href{S6_SERIES_GRIND_V.pdf\#nameddest=tag-02QI}{V.5.2 (p.~80)} }
\expandafter\def\csname SRef@night:P1.3\endcsname{\href{S6_SERIES_GRIND_V.pdf\#nameddest=tag-02QU}{V.5.3 (p.~80)} }
\expandafter\def\csname SRef@maps:src:F.2.6\endcsname{\href{S6_SERIES_GRIND_V.pdf\#nameddest=tag-02NX}{V.5.5 (p.~82)} }
\expandafter\def\csname SRef@night:P1.10\endcsname{\href{S6_SERIES_GRIND_V.pdf\#nameddest=tag-02QJ}{V.5.6 (p.~86)} }
\expandafter\def\csname SRef@night:P1.11\endcsname{\href{S6_SERIES_GRIND_V.pdf\#nameddest=tag-02QK}{V.5.7 (p.~87)} }
\expandafter\def\csname SRef@night:P1.12\endcsname{\href{S6_SERIES_GRIND_V.pdf\#nameddest=tag-02QL}{V.5.8 (p.~91)} }
\expandafter\def\csname SRef@night:P1.19\endcsname{\href{S6_SERIES_GRIND_V.pdf\#nameddest=tag-02QS}{V.5.9 (p.~94)} }
\expandafter\def\csname SRef@night:P1.4\endcsname{\href{S6_SERIES_GRIND_V.pdf\#nameddest=tag-02QV}{V.6.1 (p.~103)} }
\expandafter\def\csname SRef@night:P1.5\endcsname{\href{S6_SERIES_GRIND_V.pdf\#nameddest=tag-02QW}{V.6.2 (p.~104)} }
\expandafter\def\csname SRef@maps:src:F.3.5\endcsname{\href{S6_SERIES_GRIND_V.pdf\#nameddest=tag-02O7}{V.6.3 (p.~104)} }
\expandafter\def\csname SRef@night:P1.2\endcsname{\href{S6_SERIES_GRIND_V.pdf\#nameddest=tag-02QT}{V.6.4 (p.~106)} }
\expandafter\def\csname SRef@night:P1.7\endcsname{\href{S6_SERIES_GRIND_V.pdf\#nameddest=tag-02QY}{V.6.5 (p.~108)} }
\expandafter\def\csname SRef@maps:src:F.3.6\endcsname{\href{S6_SERIES_GRIND_V.pdf\#nameddest=tag-02O8}{V.6.6 (p.~109)} }
\expandafter\def\csname SRef@night:P1.8\endcsname{\href{S6_SERIES_GRIND_V.pdf\#nameddest=tag-02QZ}{V.6.7 (p.~112)} }
\expandafter\def\csname SRef@night:P1.9\endcsname{\href{S6_SERIES_GRIND_V.pdf\#nameddest=tag-02R0}{V.6.8 (p.~113)} }
\expandafter\def\csname SRef@maps:src:F.3.7\endcsname{\href{S6_SERIES_GRIND_V.pdf\#nameddest=tag-02O9}{V.6.9 (p.~114)} }
\expandafter\def\csname SRef@night:P1.13\endcsname{\href{S6_SERIES_GRIND_V.pdf\#nameddest=tag-02QM}{V.7.1 (p.~116)} }
\expandafter\def\csname SRef@maps:src:F.4.1\endcsname{\href{S6_SERIES_GRIND_V.pdf\#nameddest=tag-02OC}{V.7.2 (p.~117)} }
\expandafter\def\csname SRef@night:P1.15\endcsname{\href{S6_SERIES_GRIND_V.pdf\#nameddest=tag-02QO}{V.8.1 (p.~119)} }
\expandafter\def\csname SRef@night:P1.16\endcsname{\href{S6_SERIES_GRIND_V.pdf\#nameddest=tag-02QP}{V.8.2 (p.~121)} }
\expandafter\def\csname SRef@night:P1.18\endcsname{\href{S6_SERIES_GRIND_V.pdf\#nameddest=tag-02QR}{V.8.3 (p.~124)} }
\expandafter\def\csname SRef@night:P1A.1\endcsname{\href{S6_SERIES_GRIND_V.pdf\#nameddest=tag-02R1}{V.9.1 (p.~134)} }
\expandafter\def\csname SRef@night:P1A.2\endcsname{\href{S6_SERIES_GRIND_V.pdf\#nameddest=tag-02R2}{V.9.2 (p.~135)} }
\expandafter\def\csname SRef@night:P1A.3\endcsname{\href{S6_SERIES_GRIND_V.pdf\#nameddest=tag-02R3}{V.9.3 (p.~137)} }
\expandafter\def\csname SRef@night:P1A.4\endcsname{\href{S6_SERIES_GRIND_V.pdf\#nameddest=tag-02R4}{V.9.4 (p.~139)} }
\expandafter\def\csname SRef@night:P1.17\endcsname{\href{S6_SERIES_GRIND_V.pdf\#nameddest=tag-02QQ}{V.9.5 (p.~141)} }
\expandafter\def\csname SRef@night:P1A.5\endcsname{\href{S6_SERIES_GRIND_V.pdf\#nameddest=tag-02R5}{V.10.1 (p.~143)} }
\expandafter\def\csname SRef@night:P1A.6\endcsname{\href{S6_SERIES_GRIND_V.pdf\#nameddest=tag-02R6}{V.10.2 (p.~146)} }
\expandafter\def\csname SRef@cdp:eq:C.UN.1\endcsname{\hyperlink{tag-0152}{II.\ref*{cdp:eq:C.UN.1}}}
\expandafter\def\csname SRef@cdp:eq:C.FL.1\endcsname{\hyperlink{tag-00KR}{II.\ref*{cdp:eq:C.FL.1}}}
\expandafter\def\csname SRef@cdp:eq:C.QT.11\endcsname{\hyperlink{tag-00WF}{II.\ref*{cdp:eq:C.QT.11}}}
\expandafter\def\csname SRef@cdp:eq:C.QT.12\endcsname{\hyperlink{tag-00WG}{II.\ref*{cdp:eq:C.QT.12}}}
\expandafter\def\csname SRef@cdp:eq:C.QT.13\endcsname{\hyperlink{tag-00WH}{II.\ref*{cdp:eq:C.QT.13}}}
\expandafter\def\csname SRef@cdp:eq:C.WT.1\endcsname{\hyperlink{tag-017O}{II.\ref*{cdp:eq:C.WT.1}}}
\expandafter\def\csname SRef@cdp:eq:C.WT.2\endcsname{\hyperlink{tag-017Z}{II.\ref*{cdp:eq:C.WT.2}}}
\expandafter\def\csname SRef@cdp:eq:C.WT.3\endcsname{\hyperlink{tag-018B}{II.\ref*{cdp:eq:C.WT.3}}}
\expandafter\def\csname SRef@cdp:eq:C.WT.4\endcsname{\hyperlink{tag-018M}{II.\ref*{cdp:eq:C.WT.4}}}
\expandafter\def\csname SRef@cdp:eq:C.WT.5\endcsname{\hyperlink{tag-018R}{II.\ref*{cdp:eq:C.WT.5}}}
\expandafter\def\csname SRef@cdp:eq:C.WT.6\endcsname{\hyperlink{tag-018S}{II.\ref*{cdp:eq:C.WT.6}}}
\expandafter\def\csname SRef@cdp:eq:C.WT.7\endcsname{\hyperlink{tag-018T}{II.\ref*{cdp:eq:C.WT.7}}}
\expandafter\def\csname SRef@cdp:eq:C.FL.2\endcsname{\hyperlink{tag-00KS}{II.\ref*{cdp:eq:C.FL.2}}}
\expandafter\def\csname SRef@cdp:eq:C.WT.8\endcsname{\hyperlink{tag-018U}{II.\ref*{cdp:eq:C.WT.8}}}
\expandafter\def\csname SRef@cdp:eq:C.WT.9\endcsname{\hyperlink{tag-018V}{II.\ref*{cdp:eq:C.WT.9}}}
\expandafter\def\csname SRef@cdp:eq:C.WT.10\endcsname{\hyperlink{tag-017P}{II.\ref*{cdp:eq:C.WT.10}}}
\expandafter\def\csname SRef@cdp:eq:C.WT.11\endcsname{\hyperlink{tag-017Q}{II.\ref*{cdp:eq:C.WT.11}}}
\expandafter\def\csname SRef@cdp:eq:C.WT.12\endcsname{\hyperlink{tag-017R}{II.\ref*{cdp:eq:C.WT.12}}}
\expandafter\def\csname SRef@cdp:eq:C.WT.13\endcsname{\hyperlink{tag-017S}{II.\ref*{cdp:eq:C.WT.13}}}
\expandafter\def\csname SRef@cdp:eq:C.WT.14\endcsname{\hyperlink{tag-017T}{II.\ref*{cdp:eq:C.WT.14}}}
\expandafter\def\csname SRef@cdp:eq:C.WT.15\endcsname{\hyperlink{tag-017U}{II.\ref*{cdp:eq:C.WT.15}}}
\expandafter\def\csname SRef@cdp:eq:C.WT.16\endcsname{\hyperlink{tag-017V}{II.\ref*{cdp:eq:C.WT.16}}}
\expandafter\def\csname SRef@cdp:eq:C.WT.17\endcsname{\hyperlink{tag-017W}{II.\ref*{cdp:eq:C.WT.17}}}
\expandafter\def\csname SRef@cdp:eq:C.FL.3\endcsname{\hyperlink{tag-00KT}{II.\ref*{cdp:eq:C.FL.3}}}
\expandafter\def\csname SRef@cdp:eq:C.WT.18\endcsname{\hyperlink{tag-017X}{II.\ref*{cdp:eq:C.WT.18}}}
\expandafter\def\csname SRef@cdp:eq:C.WT.19\endcsname{\hyperlink{tag-017Y}{II.\ref*{cdp:eq:C.WT.19}}}
\expandafter\def\csname SRef@cdp:eq:C.WT.20\endcsname{\hyperlink{tag-0180}{II.\ref*{cdp:eq:C.WT.20}}}
\expandafter\def\csname SRef@cdp:eq:C.WT.21\endcsname{\hyperlink{tag-0181}{II.\ref*{cdp:eq:C.WT.21}}}
\expandafter\def\csname SRef@cdp:eq:C.WT.22\endcsname{\hyperlink{tag-0182}{II.\ref*{cdp:eq:C.WT.22}}}
\expandafter\def\csname SRef@cdp:eq:C.WT.23\endcsname{\hyperlink{tag-0183}{II.\ref*{cdp:eq:C.WT.23}}}
\expandafter\def\csname SRef@cdp:eq:C.WT.23a\endcsname{\hyperlink{tag-0184}{II.\ref*{cdp:eq:C.WT.23a}}}
\expandafter\def\csname SRef@cdp:eq:C.WT.24\endcsname{\hyperlink{tag-0185}{II.\ref*{cdp:eq:C.WT.24}}}
\expandafter\def\csname SRef@cdp:eq:C.WT.25\endcsname{\hyperlink{tag-0186}{II.\ref*{cdp:eq:C.WT.25}}}
\expandafter\def\csname SRef@cdp:eq:C.WT.26\endcsname{\hyperlink{tag-0187}{II.\ref*{cdp:eq:C.WT.26}}}
\expandafter\def\csname SRef@cdp:eq:C.FL.4\endcsname{\hyperlink{tag-00KU}{II.\ref*{cdp:eq:C.FL.4}}}
\expandafter\def\csname SRef@cdp:eq:C.WT.27\endcsname{\hyperlink{tag-0188}{II.\ref*{cdp:eq:C.WT.27}}}
\expandafter\def\csname SRef@cdp:eq:C.WT.28\endcsname{\hyperlink{tag-0189}{II.\ref*{cdp:eq:C.WT.28}}}
\expandafter\def\csname SRef@cdp:eq:C.WT.29\endcsname{\hyperlink{tag-018A}{II.\ref*{cdp:eq:C.WT.29}}}
\expandafter\def\csname SRef@cdp:eq:C.WT.30\endcsname{\hyperlink{tag-018C}{II.\ref*{cdp:eq:C.WT.30}}}
\expandafter\def\csname SRef@cdp:eq:C.WT.31\endcsname{\hyperlink{tag-018D}{II.\ref*{cdp:eq:C.WT.31}}}
\expandafter\def\csname SRef@cdp:eq:C.WT.32\endcsname{\hyperlink{tag-018E}{II.\ref*{cdp:eq:C.WT.32}}}
\expandafter\def\csname SRef@cdp:eq:C.WT.33\endcsname{\hyperlink{tag-018F}{II.\ref*{cdp:eq:C.WT.33}}}
\expandafter\def\csname SRef@cdp:eq:C.WT.34\endcsname{\hyperlink{tag-018G}{II.\ref*{cdp:eq:C.WT.34}}}
\expandafter\def\csname SRef@cdp:eq:C.WT.35\endcsname{\hyperlink{tag-018H}{II.\ref*{cdp:eq:C.WT.35}}}
\expandafter\def\csname SRef@cdp:eq:C.WT.36\endcsname{\hyperlink{tag-018I}{II.\ref*{cdp:eq:C.WT.36}}}
\expandafter\def\csname SRef@cdp:eq:C.FL.5\endcsname{\hyperlink{tag-00KV}{II.\ref*{cdp:eq:C.FL.5}}}
\expandafter\def\csname SRef@cdp:eq:C.WT.37\endcsname{\hyperlink{tag-018J}{II.\ref*{cdp:eq:C.WT.37}}}
\expandafter\def\csname SRef@cdp:eq:C.WT.38\endcsname{\hyperlink{tag-018K}{II.\ref*{cdp:eq:C.WT.38}}}
\expandafter\def\csname SRef@cdp:eq:C.WT.39\endcsname{\hyperlink{tag-018L}{II.\ref*{cdp:eq:C.WT.39}}}
\expandafter\def\csname SRef@cdp:eq:C.WT.40\endcsname{\hyperlink{tag-018N}{II.\ref*{cdp:eq:C.WT.40}}}
\expandafter\def\csname SRef@cdp:eq:C.WT.41\endcsname{\hyperlink{tag-018O}{II.\ref*{cdp:eq:C.WT.41}}}
\expandafter\def\csname SRef@cdp:eq:C.WT.42\endcsname{\hyperlink{tag-018P}{II.\ref*{cdp:eq:C.WT.42}}}
\expandafter\def\csname SRef@cdp:eq:C.WT.43\endcsname{\hyperlink{tag-018Q}{II.\ref*{cdp:eq:C.WT.43}}}
\expandafter\def\csname SRef@cdp:eq:C.NM.1\endcsname{\hyperlink{tag-00SR}{II.\ref*{cdp:eq:C.NM.1}}}
\expandafter\def\csname SRef@cdp:eq:C.NM.2\endcsname{\hyperlink{tag-00ST}{II.\ref*{cdp:eq:C.NM.2}}}
\expandafter\def\csname SRef@cdp:eq:C.NM.3\endcsname{\hyperlink{tag-00SU}{II.\ref*{cdp:eq:C.NM.3}}}
\expandafter\def\csname SRef@cdp:eq:C.PG.1\endcsname{\hyperlink{tag-00UZ}{II.\ref*{cdp:eq:C.PG.1}}}
\expandafter\def\csname SRef@cdp:eq:C.NM.4\endcsname{\hyperlink{tag-00SV}{II.\ref*{cdp:eq:C.NM.4}}}
\expandafter\def\csname SRef@cdp:eq:C.NM.5\endcsname{\hyperlink{tag-00SW}{II.\ref*{cdp:eq:C.NM.5}}}
\expandafter\def\csname SRef@cdp:eq:C.NM.6\endcsname{\hyperlink{tag-00SX}{II.\ref*{cdp:eq:C.NM.6}}}
\expandafter\def\csname SRef@cdp:eq:C.NM.7\endcsname{\hyperlink{tag-00SY}{II.\ref*{cdp:eq:C.NM.7}}}
\expandafter\def\csname SRef@cdp:eq:C.NM.8\endcsname{\hyperlink{tag-00SZ}{II.\ref*{cdp:eq:C.NM.8}}}
\expandafter\def\csname SRef@cdp:eq:C.NM.9\endcsname{\hyperlink{tag-00T0}{II.\ref*{cdp:eq:C.NM.9}}}
\expandafter\def\csname SRef@cdp:eq:C.NM.10\endcsname{\hyperlink{tag-00SS}{II.\ref*{cdp:eq:C.NM.10}}}
\expandafter\def\csname SRef@cdp:eq:C.NP.1\endcsname{\hyperlink{tag-00T1}{II.\ref*{cdp:eq:C.NP.1}}}
\expandafter\def\csname SRef@cdp:eq:C.NP.2\endcsname{\hyperlink{tag-00T2}{II.\ref*{cdp:eq:C.NP.2}}}
\expandafter\def\csname SRef@cdp:eq:C.NP.3\endcsname{\hyperlink{tag-00T3}{II.\ref*{cdp:eq:C.NP.3}}}
\expandafter\def\csname SRef@cdp:eq:C.PG.2\endcsname{\hyperlink{tag-00V0}{II.\ref*{cdp:eq:C.PG.2}}}
\expandafter\def\csname SRef@cdp:eq:C.NP.4\endcsname{\hyperlink{tag-00T4}{II.\ref*{cdp:eq:C.NP.4}}}
\expandafter\def\csname SRef@cdp:eq:C.NP.5\endcsname{\hyperlink{tag-00T5}{II.\ref*{cdp:eq:C.NP.5}}}
\expandafter\def\csname SRef@cdp:eq:C.NP.6\endcsname{\hyperlink{tag-00T6}{II.\ref*{cdp:eq:C.NP.6}}}
\expandafter\def\csname SRef@cdp:eq:C.NP.7\endcsname{\hyperlink{tag-00T7}{II.\ref*{cdp:eq:C.NP.7}}}
\expandafter\def\csname SRef@cdp:eq:C.NP.8\endcsname{\hyperlink{tag-00T8}{II.\ref*{cdp:eq:C.NP.8}}}
\expandafter\def\csname SRef@cdp:eq:C.NP.9\endcsname{\hyperlink{tag-00T9}{II.\ref*{cdp:eq:C.NP.9}}}
\expandafter\def\csname SRef@cdp:eq:C.TG.1\endcsname{\hyperlink{tag-012Q}{II.\ref*{cdp:eq:C.TG.1}}}
\expandafter\def\csname SRef@cdp:eq:C.TG.2\endcsname{\hyperlink{tag-012X}{II.\ref*{cdp:eq:C.TG.2}}}
\expandafter\def\csname SRef@cdp:eq:C.TG.3\endcsname{\hyperlink{tag-012Y}{II.\ref*{cdp:eq:C.TG.3}}}
\expandafter\def\csname SRef@cdp:eq:C.TG.4\endcsname{\hyperlink{tag-012Z}{II.\ref*{cdp:eq:C.TG.4}}}
\expandafter\def\csname SRef@cdp:eq:C.PG.3\endcsname{\hyperlink{tag-00V1}{II.\ref*{cdp:eq:C.PG.3}}}
\expandafter\def\csname SRef@cdp:eq:C.TG.5\endcsname{\hyperlink{tag-0130}{II.\ref*{cdp:eq:C.TG.5}}}
\expandafter\def\csname SRef@cdp:eq:C.TG.6\endcsname{\hyperlink{tag-0131}{II.\ref*{cdp:eq:C.TG.6}}}
\expandafter\def\csname SRef@cdp:eq:C.TG.7\endcsname{\hyperlink{tag-0132}{II.\ref*{cdp:eq:C.TG.7}}}
\expandafter\def\csname SRef@cdp:eq:C.TG.8\endcsname{\hyperlink{tag-0133}{II.\ref*{cdp:eq:C.TG.8}}}
\expandafter\def\csname SRef@cdp:eq:C.TG.9\endcsname{\hyperlink{tag-0134}{II.\ref*{cdp:eq:C.TG.9}}}
\expandafter\def\csname SRef@cdp:eq:C.TG.10\endcsname{\hyperlink{tag-012R}{II.\ref*{cdp:eq:C.TG.10}}}
\expandafter\def\csname SRef@cdp:eq:C.TG.11\endcsname{\hyperlink{tag-012S}{II.\ref*{cdp:eq:C.TG.11}}}
\expandafter\def\csname SRef@cdp:eq:C.TG.12\endcsname{\hyperlink{tag-012T}{II.\ref*{cdp:eq:C.TG.12}}}
\expandafter\def\csname SRef@cdp:eq:C.TG.13\endcsname{\hyperlink{tag-012U}{II.\ref*{cdp:eq:C.TG.13}}}
\expandafter\def\csname SRef@cdp:eq:C.TG.14\endcsname{\hyperlink{tag-012V}{II.\ref*{cdp:eq:C.TG.14}}}
\expandafter\def\csname SRef@cdp:eq:C.PG.4\endcsname{\hyperlink{tag-00V2}{II.\ref*{cdp:eq:C.PG.4}}}
\expandafter\def\csname SRef@cdp:eq:C.TG.15\endcsname{\hyperlink{tag-012W}{II.\ref*{cdp:eq:C.TG.15}}}
\expandafter\def\csname SRef@cdp:eq:C.TX.1\endcsname{\hyperlink{tag-013Y}{II.\ref*{cdp:eq:C.TX.1}}}
\expandafter\def\csname SRef@cdp:eq:C.TX.2\endcsname{\hyperlink{tag-0144}{II.\ref*{cdp:eq:C.TX.2}}}
\expandafter\def\csname SRef@cdp:eq:C.TX.3\endcsname{\hyperlink{tag-0145}{II.\ref*{cdp:eq:C.TX.3}}}
\expandafter\def\csname SRef@cdp:eq:C.TX.4\endcsname{\hyperlink{tag-0146}{II.\ref*{cdp:eq:C.TX.4}}}
\expandafter\def\csname SRef@cdp:eq:C.TX.5\endcsname{\hyperlink{tag-0147}{II.\ref*{cdp:eq:C.TX.5}}}
\expandafter\def\csname SRef@cdp:eq:C.TX.6\endcsname{\hyperlink{tag-0148}{II.\ref*{cdp:eq:C.TX.6}}}
\expandafter\def\csname SRef@cdp:eq:C.TX.7\endcsname{\hyperlink{tag-0149}{II.\ref*{cdp:eq:C.TX.7}}}
\expandafter\def\csname SRef@cdp:eq:C.TX.8\endcsname{\hyperlink{tag-014A}{II.\ref*{cdp:eq:C.TX.8}}}
\expandafter\def\csname SRef@cdp:eq:C.TX.9\endcsname{\hyperlink{tag-014B}{II.\ref*{cdp:eq:C.TX.9}}}
\expandafter\def\csname SRef@cdp:eq:C.PG.5\endcsname{\hyperlink{tag-00V3}{II.\ref*{cdp:eq:C.PG.5}}}
\expandafter\def\csname SRef@cdp:eq:C.TX.10\endcsname{\hyperlink{tag-013Z}{II.\ref*{cdp:eq:C.TX.10}}}
\expandafter\def\csname SRef@cdp:eq:C.TX.11\endcsname{\hyperlink{tag-0140}{II.\ref*{cdp:eq:C.TX.11}}}
\expandafter\def\csname SRef@cdp:eq:C.TX.12\endcsname{\hyperlink{tag-0141}{II.\ref*{cdp:eq:C.TX.12}}}
\expandafter\def\csname SRef@cdp:eq:C.TX.13\endcsname{\hyperlink{tag-0142}{II.\ref*{cdp:eq:C.TX.13}}}
\expandafter\def\csname SRef@cdp:eq:C.TX.14\endcsname{\hyperlink{tag-0143}{II.\ref*{cdp:eq:C.TX.14}}}
\expandafter\def\csname SRef@cdp:eq:C.WF.1\endcsname{\hyperlink{tag-0165}{II.\ref*{cdp:eq:C.WF.1}}}
\expandafter\def\csname SRef@cdp:eq:C.WF.2\endcsname{\hyperlink{tag-016G}{II.\ref*{cdp:eq:C.WF.2}}}
\expandafter\def\csname SRef@cdp:eq:C.WF.3\endcsname{\hyperlink{tag-016Q}{II.\ref*{cdp:eq:C.WF.3}}}
\expandafter\def\csname SRef@cdp:eq:C.WF.4\endcsname{\hyperlink{tag-016R}{II.\ref*{cdp:eq:C.WF.4}}}
\expandafter\def\csname SRef@cdp:eq:C.WF.5\endcsname{\hyperlink{tag-016S}{II.\ref*{cdp:eq:C.WF.5}}}
\expandafter\def\csname SRef@cdp:eq:C.UN.2\endcsname{\hyperlink{tag-0153}{II.\ref*{cdp:eq:C.UN.2}}}
\expandafter\def\csname SRef@cdp:eq:C.PG.6\endcsname{\hyperlink{tag-00V4}{II.\ref*{cdp:eq:C.PG.6}}}
\expandafter\def\csname SRef@cdp:eq:C.WF.6\endcsname{\hyperlink{tag-016T}{II.\ref*{cdp:eq:C.WF.6}}}
\expandafter\def\csname SRef@cdp:eq:C.WF.7\endcsname{\hyperlink{tag-016U}{II.\ref*{cdp:eq:C.WF.7}}}
\expandafter\def\csname SRef@cdp:eq:C.WF.8\endcsname{\hyperlink{tag-016V}{II.\ref*{cdp:eq:C.WF.8}}}
\expandafter\def\csname SRef@cdp:eq:C.WF.9\endcsname{\hyperlink{tag-016W}{II.\ref*{cdp:eq:C.WF.9}}}
\expandafter\def\csname SRef@cdp:eq:C.WF.10\endcsname{\hyperlink{tag-0166}{II.\ref*{cdp:eq:C.WF.10}}}
\expandafter\def\csname SRef@cdp:eq:C.WF.11\endcsname{\hyperlink{tag-0167}{II.\ref*{cdp:eq:C.WF.11}}}
\expandafter\def\csname SRef@cdp:eq:C.WF.12\endcsname{\hyperlink{tag-0168}{II.\ref*{cdp:eq:C.WF.12}}}
\expandafter\def\csname SRef@cdp:eq:C.WF.13\endcsname{\hyperlink{tag-0169}{II.\ref*{cdp:eq:C.WF.13}}}
\expandafter\def\csname SRef@cdp:eq:C.WF.14\endcsname{\hyperlink{tag-016A}{II.\ref*{cdp:eq:C.WF.14}}}
\expandafter\def\csname SRef@cdp:eq:C.WF.15\endcsname{\hyperlink{tag-016B}{II.\ref*{cdp:eq:C.WF.15}}}
\expandafter\def\csname SRef@cdp:eq:C.PG.7\endcsname{\hyperlink{tag-00V5}{II.\ref*{cdp:eq:C.PG.7}}}
\expandafter\def\csname SRef@cdp:eq:C.WF.16\endcsname{\hyperlink{tag-016C}{II.\ref*{cdp:eq:C.WF.16}}}
\expandafter\def\csname SRef@cdp:eq:C.WF.17\endcsname{\hyperlink{tag-016D}{II.\ref*{cdp:eq:C.WF.17}}}
\expandafter\def\csname SRef@cdp:eq:C.WF.18\endcsname{\hyperlink{tag-016E}{II.\ref*{cdp:eq:C.WF.18}}}
\expandafter\def\csname SRef@cdp:eq:C.WF.19\endcsname{\hyperlink{tag-016F}{II.\ref*{cdp:eq:C.WF.19}}}
\expandafter\def\csname SRef@cdp:eq:C.WF.20\endcsname{\hyperlink{tag-016H}{II.\ref*{cdp:eq:C.WF.20}}}
\expandafter\def\csname SRef@cdp:eq:C.WF.21\endcsname{\hyperlink{tag-016I}{II.\ref*{cdp:eq:C.WF.21}}}
\expandafter\def\csname SRef@cdp:eq:C.WF.22\endcsname{\hyperlink{tag-016J}{II.\ref*{cdp:eq:C.WF.22}}}
\expandafter\def\csname SRef@cdp:eq:C.WF.23\endcsname{\hyperlink{tag-016K}{II.\ref*{cdp:eq:C.WF.23}}}
\expandafter\def\csname SRef@cdp:eq:C.WF.24\endcsname{\hyperlink{tag-016L}{II.\ref*{cdp:eq:C.WF.24}}}
\expandafter\def\csname SRef@cdp:eq:C.WF.25\endcsname{\hyperlink{tag-016M}{II.\ref*{cdp:eq:C.WF.25}}}
\expandafter\def\csname SRef@cdp:eq:C.PG.8\endcsname{\hyperlink{tag-00V6}{II.\ref*{cdp:eq:C.PG.8}}}
\expandafter\def\csname SRef@cdp:eq:C.WF.26\endcsname{\hyperlink{tag-016N}{II.\ref*{cdp:eq:C.WF.26}}}
\expandafter\def\csname SRef@cdp:eq:C.WF.27\endcsname{\hyperlink{tag-016O}{II.\ref*{cdp:eq:C.WF.27}}}
\expandafter\def\csname SRef@cdp:eq:C.WF.28\endcsname{\hyperlink{tag-016P}{II.\ref*{cdp:eq:C.WF.28}}}
\expandafter\def\csname SRef@cdp:eq:C.RC.1\endcsname{\hyperlink{tag-00XN}{II.\ref*{cdp:eq:C.RC.1}}}
\expandafter\def\csname SRef@cdp:eq:C.RC.2\endcsname{\hyperlink{tag-00XR}{II.\ref*{cdp:eq:C.RC.2}}}
\expandafter\def\csname SRef@cdp:eq:C.RC.3\endcsname{\hyperlink{tag-00XS}{II.\ref*{cdp:eq:C.RC.3}}}
\expandafter\def\csname SRef@cdp:eq:C.RC.4\endcsname{\hyperlink{tag-00XT}{II.\ref*{cdp:eq:C.RC.4}}}
\expandafter\def\csname SRef@cdp:eq:C.RC.5\endcsname{\hyperlink{tag-00XU}{II.\ref*{cdp:eq:C.RC.5}}}
\expandafter\def\csname SRef@cdp:eq:C.RC.6\endcsname{\hyperlink{tag-00XV}{II.\ref*{cdp:eq:C.RC.6}}}
\expandafter\def\csname SRef@cdp:eq:C.RC.7\endcsname{\hyperlink{tag-00XW}{II.\ref*{cdp:eq:C.RC.7}}}
\expandafter\def\csname SRef@cdp:eq:C.PG.9\endcsname{\hyperlink{tag-00V7}{II.\ref*{cdp:eq:C.PG.9}}}
\expandafter\def\csname SRef@cdp:eq:C.RC.8\endcsname{\hyperlink{tag-00XX}{II.\ref*{cdp:eq:C.RC.8}}}
\expandafter\def\csname SRef@cdp:eq:C.RC.9\endcsname{\hyperlink{tag-00XY}{II.\ref*{cdp:eq:C.RC.9}}}
\expandafter\def\csname SRef@cdp:eq:C.RC.10\endcsname{\hyperlink{tag-00XO}{II.\ref*{cdp:eq:C.RC.10}}}
\expandafter\def\csname SRef@cdp:eq:C.RC.11\endcsname{\hyperlink{tag-00XP}{II.\ref*{cdp:eq:C.RC.11}}}
\expandafter\def\csname SRef@cdp:eq:C.RC.12\endcsname{\hyperlink{tag-00XQ}{II.\ref*{cdp:eq:C.RC.12}}}
\expandafter\def\csname SRef@cdp:eq:C.KG.1\endcsname{\hyperlink{tag-00P8}{II.\ref*{cdp:eq:C.KG.1}}}
\expandafter\def\csname SRef@cdp:eq:C.KG.2\endcsname{\hyperlink{tag-00PA}{II.\ref*{cdp:eq:C.KG.2}}}
\expandafter\def\csname SRef@cdp:eq:C.KG.3\endcsname{\hyperlink{tag-00PB}{II.\ref*{cdp:eq:C.KG.3}}}
\expandafter\def\csname SRef@cdp:eq:C.KG.4\endcsname{\hyperlink{tag-00PC}{II.\ref*{cdp:eq:C.KG.4}}}
\expandafter\def\csname SRef@cdp:eq:C.KG.5\endcsname{\hyperlink{tag-00PD}{II.\ref*{cdp:eq:C.KG.5}}}
\expandafter\def\csname SRef@cdp:eq:C.RG.1\endcsname{\hyperlink{tag-00Z3}{II.\ref*{cdp:eq:C.RG.1}}}
\expandafter\def\csname SRef@cdp:eq:C.KG.6\endcsname{\hyperlink{tag-00PE}{II.\ref*{cdp:eq:C.KG.6}}}
\expandafter\def\csname SRef@cdp:eq:C.KG.7\endcsname{\hyperlink{tag-00PF}{II.\ref*{cdp:eq:C.KG.7}}}
\expandafter\def\csname SRef@cdp:eq:C.KG.8\endcsname{\hyperlink{tag-00PG}{II.\ref*{cdp:eq:C.KG.8}}}
\expandafter\def\csname SRef@cdp:eq:C.KG.9\endcsname{\hyperlink{tag-00PH}{II.\ref*{cdp:eq:C.KG.9}}}
\expandafter\def\csname SRef@cdp:eq:C.KG.10\endcsname{\hyperlink{tag-00P9}{II.\ref*{cdp:eq:C.KG.10}}}
\expandafter\def\csname SRef@cdp:eq:C.HM.1\endcsname{\hyperlink{tag-00MG}{II.\ref*{cdp:eq:C.HM.1}}}
\expandafter\def\csname SRef@cdp:eq:C.HM.2\endcsname{\hyperlink{tag-00MI}{II.\ref*{cdp:eq:C.HM.2}}}
\expandafter\def\csname SRef@cdp:eq:C.HM.3\endcsname{\hyperlink{tag-00MJ}{II.\ref*{cdp:eq:C.HM.3}}}
\expandafter\def\csname SRef@cdp:eq:C.HM.4\endcsname{\hyperlink{tag-00MK}{II.\ref*{cdp:eq:C.HM.4}}}
\expandafter\def\csname SRef@cdp:eq:C.HM.5\endcsname{\hyperlink{tag-00ML}{II.\ref*{cdp:eq:C.HM.5}}}
\expandafter\def\csname SRef@cdp:eq:C.RG.2\endcsname{\hyperlink{tag-00Z8}{II.\ref*{cdp:eq:C.RG.2}}}
\expandafter\def\csname SRef@cdp:eq:C.HM.6\endcsname{\hyperlink{tag-00MM}{II.\ref*{cdp:eq:C.HM.6}}}
\expandafter\def\csname SRef@cdp:eq:C.HM.7\endcsname{\hyperlink{tag-00MN}{II.\ref*{cdp:eq:C.HM.7}}}
\expandafter\def\csname SRef@cdp:eq:C.HM.8\endcsname{\hyperlink{tag-00MO}{II.\ref*{cdp:eq:C.HM.8}}}
\expandafter\def\csname SRef@cdp:eq:C.HM.9\endcsname{\hyperlink{tag-00MP}{II.\ref*{cdp:eq:C.HM.9}}}
\expandafter\def\csname SRef@cdp:eq:C.HM.10\endcsname{\hyperlink{tag-00MH}{II.\ref*{cdp:eq:C.HM.10}}}
\expandafter\def\csname SRef@cdp:eq:C.BC.1\endcsname{\hyperlink{tag-00AZ}{II.\ref*{cdp:eq:C.BC.1}}}
\expandafter\def\csname SRef@cdp:eq:C.BC.2\endcsname{\hyperlink{tag-00B4}{II.\ref*{cdp:eq:C.BC.2}}}
\expandafter\def\csname SRef@cdp:eq:C.BC.3\endcsname{\hyperlink{tag-00B5}{II.\ref*{cdp:eq:C.BC.3}}}
\expandafter\def\csname SRef@cdp:eq:C.BC.4\endcsname{\hyperlink{tag-00B6}{II.\ref*{cdp:eq:C.BC.4}}}
\expandafter\def\csname SRef@cdp:eq:C.BC.5\endcsname{\hyperlink{tag-00B7}{II.\ref*{cdp:eq:C.BC.5}}}
\expandafter\def\csname SRef@cdp:eq:C.RG.3\endcsname{\hyperlink{tag-00Z9}{II.\ref*{cdp:eq:C.RG.3}}}
\expandafter\def\csname SRef@cdp:eq:C.BC.6\endcsname{\hyperlink{tag-00B8}{II.\ref*{cdp:eq:C.BC.6}}}
\expandafter\def\csname SRef@cdp:eq:C.BC.7\endcsname{\hyperlink{tag-00B9}{II.\ref*{cdp:eq:C.BC.7}}}
\expandafter\def\csname SRef@cdp:eq:C.BC.8\endcsname{\hyperlink{tag-00BA}{II.\ref*{cdp:eq:C.BC.8}}}
\expandafter\def\csname SRef@cdp:eq:C.BC.9\endcsname{\hyperlink{tag-00BB}{II.\ref*{cdp:eq:C.BC.9}}}
\expandafter\def\csname SRef@cdp:eq:C.BC.10\endcsname{\hyperlink{tag-00B0}{II.\ref*{cdp:eq:C.BC.10}}}
\expandafter\def\csname SRef@cdp:eq:C.BC.11\endcsname{\hyperlink{tag-00B1}{II.\ref*{cdp:eq:C.BC.11}}}
\expandafter\def\csname SRef@cdp:eq:C.BC.12\endcsname{\hyperlink{tag-00B2}{II.\ref*{cdp:eq:C.BC.12}}}
\expandafter\def\csname SRef@cdp:eq:C.BC.13\endcsname{\hyperlink{tag-00B3}{II.\ref*{cdp:eq:C.BC.13}}}
\expandafter\def\csname SRef@cdp:eq:C.EP.1\endcsname{\hyperlink{tag-00I1}{II.\ref*{cdp:eq:C.EP.1}}}
\expandafter\def\csname SRef@cdp:eq:C.EP.2\endcsname{\hyperlink{tag-00I4}{II.\ref*{cdp:eq:C.EP.2}}}
\expandafter\def\csname SRef@cdp:eq:C.RG.4\endcsname{\hyperlink{tag-00ZA}{II.\ref*{cdp:eq:C.RG.4}}}
\expandafter\def\csname SRef@cdp:eq:C.EP.3\endcsname{\hyperlink{tag-00I5}{II.\ref*{cdp:eq:C.EP.3}}}
\expandafter\def\csname SRef@cdp:eq:C.EP.4\endcsname{\hyperlink{tag-00I6}{II.\ref*{cdp:eq:C.EP.4}}}
\expandafter\def\csname SRef@cdp:eq:C.EP.5\endcsname{\hyperlink{tag-00I7}{II.\ref*{cdp:eq:C.EP.5}}}
\expandafter\def\csname SRef@cdp:eq:C.EP.6\endcsname{\hyperlink{tag-00I8}{II.\ref*{cdp:eq:C.EP.6}}}
\expandafter\def\csname SRef@cdp:eq:C.EP.7\endcsname{\hyperlink{tag-00I9}{II.\ref*{cdp:eq:C.EP.7}}}
\expandafter\def\csname SRef@cdp:eq:C.EP.8\endcsname{\hyperlink{tag-00IA}{II.\ref*{cdp:eq:C.EP.8}}}
\expandafter\def\csname SRef@cdp:eq:C.EP.9\endcsname{\hyperlink{tag-00IB}{II.\ref*{cdp:eq:C.EP.9}}}
\expandafter\def\csname SRef@cdp:eq:C.EP.10\endcsname{\hyperlink{tag-00I2}{II.\ref*{cdp:eq:C.EP.10}}}
\expandafter\def\csname SRef@cdp:eq:C.EP.11\endcsname{\hyperlink{tag-00I3}{II.\ref*{cdp:eq:C.EP.11}}}
\expandafter\def\csname SRef@cdp:eq:C.AP.1\endcsname{\hyperlink{tag-009M}{II.\ref*{cdp:eq:C.AP.1}}}
\expandafter\def\csname SRef@cdp:eq:C.RG.5\endcsname{\hyperlink{tag-00ZB}{II.\ref*{cdp:eq:C.RG.5}}}
\expandafter\def\csname SRef@cdp:eq:C.AP.2\endcsname{\hyperlink{tag-009S}{II.\ref*{cdp:eq:C.AP.2}}}
\expandafter\def\csname SRef@cdp:eq:C.AP.3\endcsname{\hyperlink{tag-009T}{II.\ref*{cdp:eq:C.AP.3}}}
\expandafter\def\csname SRef@cdp:eq:C.AP.4\endcsname{\hyperlink{tag-009U}{II.\ref*{cdp:eq:C.AP.4}}}
\expandafter\def\csname SRef@cdp:eq:C.AP.5\endcsname{\hyperlink{tag-009V}{II.\ref*{cdp:eq:C.AP.5}}}
\expandafter\def\csname SRef@cdp:eq:C.AP.6\endcsname{\hyperlink{tag-009W}{II.\ref*{cdp:eq:C.AP.6}}}
\expandafter\def\csname SRef@cdp:eq:C.AP.7\endcsname{\hyperlink{tag-009X}{II.\ref*{cdp:eq:C.AP.7}}}
\expandafter\def\csname SRef@cdp:eq:C.AP.8\endcsname{\hyperlink{tag-009Y}{II.\ref*{cdp:eq:C.AP.8}}}
\expandafter\def\csname SRef@cdp:eq:C.AP.9\endcsname{\hyperlink{tag-009Z}{II.\ref*{cdp:eq:C.AP.9}}}
\expandafter\def\csname SRef@cdp:eq:C.AP.10\endcsname{\hyperlink{tag-009N}{II.\ref*{cdp:eq:C.AP.10}}}
\expandafter\def\csname SRef@cdp:eq:C.AP.11\endcsname{\hyperlink{tag-009O}{II.\ref*{cdp:eq:C.AP.11}}}
\expandafter\def\csname SRef@cdp:eq:C.RG.6\endcsname{\hyperlink{tag-00ZC}{II.\ref*{cdp:eq:C.RG.6}}}
\expandafter\def\csname SRef@cdp:eq:C.AP.12\endcsname{\hyperlink{tag-009P}{II.\ref*{cdp:eq:C.AP.12}}}
\expandafter\def\csname SRef@cdp:eq:C.AP.13\endcsname{\hyperlink{tag-009Q}{II.\ref*{cdp:eq:C.AP.13}}}
\expandafter\def\csname SRef@cdp:eq:C.AP.14\endcsname{\hyperlink{tag-009R}{II.\ref*{cdp:eq:C.AP.14}}}
\expandafter\def\csname SRef@cdp:eq:C.ML.1\endcsname{\hyperlink{tag-00QR}{II.\ref*{cdp:eq:C.ML.1}}}
\expandafter\def\csname SRef@cdp:eq:C.ML.2\endcsname{\hyperlink{tag-00QT}{II.\ref*{cdp:eq:C.ML.2}}}
\expandafter\def\csname SRef@cdp:eq:C.ML.3\endcsname{\hyperlink{tag-00QU}{II.\ref*{cdp:eq:C.ML.3}}}
\expandafter\def\csname SRef@cdp:eq:C.ML.4\endcsname{\hyperlink{tag-00QV}{II.\ref*{cdp:eq:C.ML.4}}}
\expandafter\def\csname SRef@cdp:eq:C.ML.5\endcsname{\hyperlink{tag-00QW}{II.\ref*{cdp:eq:C.ML.5}}}
\expandafter\def\csname SRef@cdp:eq:C.ML.6\endcsname{\hyperlink{tag-00QX}{II.\ref*{cdp:eq:C.ML.6}}}
\expandafter\def\csname SRef@cdp:eq:C.ML.7\endcsname{\hyperlink{tag-00QY}{II.\ref*{cdp:eq:C.ML.7}}}
\expandafter\def\csname SRef@cdp:eq:C.UN.3\endcsname{\hyperlink{tag-0154}{II.\ref*{cdp:eq:C.UN.3}}}
\expandafter\def\csname SRef@cdp:eq:C.RG.7\endcsname{\hyperlink{tag-00ZD}{II.\ref*{cdp:eq:C.RG.7}}}
\expandafter\def\csname SRef@cdp:eq:C.ML.8\endcsname{\hyperlink{tag-00QZ}{II.\ref*{cdp:eq:C.ML.8}}}
\expandafter\def\csname SRef@cdp:eq:C.ML.9\endcsname{\hyperlink{tag-00R0}{II.\ref*{cdp:eq:C.ML.9}}}
\expandafter\def\csname SRef@cdp:eq:C.ML.10\endcsname{\hyperlink{tag-00QS}{II.\ref*{cdp:eq:C.ML.10}}}
\expandafter\def\csname SRef@cdp:eq:C.TI.1\endcsname{\hyperlink{tag-0135}{II.\ref*{cdp:eq:C.TI.1}}}
\expandafter\def\csname SRef@cdp:eq:C.TI.2\endcsname{\hyperlink{tag-013B}{II.\ref*{cdp:eq:C.TI.2}}}
\expandafter\def\csname SRef@cdp:eq:C.TI.3\endcsname{\hyperlink{tag-013C}{II.\ref*{cdp:eq:C.TI.3}}}
\expandafter\def\csname SRef@cdp:eq:C.TI.4\endcsname{\hyperlink{tag-013D}{II.\ref*{cdp:eq:C.TI.4}}}
\expandafter\def\csname SRef@cdp:eq:C.TI.5\endcsname{\hyperlink{tag-013E}{II.\ref*{cdp:eq:C.TI.5}}}
\expandafter\def\csname SRef@cdp:eq:C.TI.6\endcsname{\hyperlink{tag-013F}{II.\ref*{cdp:eq:C.TI.6}}}
\expandafter\def\csname SRef@cdp:eq:C.TI.7\endcsname{\hyperlink{tag-013G}{II.\ref*{cdp:eq:C.TI.7}}}
\expandafter\def\csname SRef@cdp:eq:C.RG.8\endcsname{\hyperlink{tag-00ZE}{II.\ref*{cdp:eq:C.RG.8}}}
\expandafter\def\csname SRef@cdp:eq:C.TI.8\endcsname{\hyperlink{tag-013H}{II.\ref*{cdp:eq:C.TI.8}}}
\expandafter\def\csname SRef@cdp:eq:C.TI.9\endcsname{\hyperlink{tag-013I}{II.\ref*{cdp:eq:C.TI.9}}}
\expandafter\def\csname SRef@cdp:eq:C.TI.10\endcsname{\hyperlink{tag-0136}{II.\ref*{cdp:eq:C.TI.10}}}
\expandafter\def\csname SRef@cdp:eq:C.TI.11\endcsname{\hyperlink{tag-0137}{II.\ref*{cdp:eq:C.TI.11}}}
\expandafter\def\csname SRef@cdp:eq:C.TI.12\endcsname{\hyperlink{tag-0138}{II.\ref*{cdp:eq:C.TI.12}}}
\expandafter\def\csname SRef@cdp:eq:C.TI.13\endcsname{\hyperlink{tag-0139}{II.\ref*{cdp:eq:C.TI.13}}}
\expandafter\def\csname SRef@cdp:eq:C.TI.14\endcsname{\hyperlink{tag-013A}{II.\ref*{cdp:eq:C.TI.14}}}
\expandafter\def\csname SRef@cdp:eq:C.TK.1\endcsname{\hyperlink{tag-013J}{II.\ref*{cdp:eq:C.TK.1}}}
\expandafter\def\csname SRef@cdp:eq:C.TK.2\endcsname{\hyperlink{tag-013K}{II.\ref*{cdp:eq:C.TK.2}}}
\expandafter\def\csname SRef@cdp:eq:C.TK.3\endcsname{\hyperlink{tag-013L}{II.\ref*{cdp:eq:C.TK.3}}}
\expandafter\def\csname SRef@cdp:eq:C.RG.9\endcsname{\hyperlink{tag-00ZF}{II.\ref*{cdp:eq:C.RG.9}}}
\expandafter\def\csname SRef@cdp:eq:C.TK.4\endcsname{\hyperlink{tag-013M}{II.\ref*{cdp:eq:C.TK.4}}}
\expandafter\def\csname SRef@cdp:eq:C.TK.5\endcsname{\hyperlink{tag-013N}{II.\ref*{cdp:eq:C.TK.5}}}
\expandafter\def\csname SRef@cdp:eq:C.TK.6\endcsname{\hyperlink{tag-013O}{II.\ref*{cdp:eq:C.TK.6}}}
\expandafter\def\csname SRef@cdp:eq:C.TK.7\endcsname{\hyperlink{tag-013P}{II.\ref*{cdp:eq:C.TK.7}}}
\expandafter\def\csname SRef@cdp:eq:C.RK.1\endcsname{\hyperlink{tag-00ZQ}{II.\ref*{cdp:eq:C.RK.1}}}
\expandafter\def\csname SRef@cdp:eq:C.RK.2\endcsname{\hyperlink{tag-0101}{II.\ref*{cdp:eq:C.RK.2}}}
\expandafter\def\csname SRef@cdp:eq:C.RK.3\endcsname{\hyperlink{tag-0105}{II.\ref*{cdp:eq:C.RK.3}}}
\expandafter\def\csname SRef@cdp:eq:C.RK.4\endcsname{\hyperlink{tag-0106}{II.\ref*{cdp:eq:C.RK.4}}}
\expandafter\def\csname SRef@cdp:eq:C.RK.5\endcsname{\hyperlink{tag-0107}{II.\ref*{cdp:eq:C.RK.5}}}
\expandafter\def\csname SRef@cdp:eq:C.RK.6\endcsname{\hyperlink{tag-0108}{II.\ref*{cdp:eq:C.RK.6}}}
\expandafter\def\csname SRef@cdp:eq:C.RG.10\endcsname{\hyperlink{tag-00Z4}{II.\ref*{cdp:eq:C.RG.10}}}
\expandafter\def\csname SRef@cdp:eq:C.RK.7\endcsname{\hyperlink{tag-0109}{II.\ref*{cdp:eq:C.RK.7}}}
\expandafter\def\csname SRef@cdp:eq:C.RK.8\endcsname{\hyperlink{tag-010A}{II.\ref*{cdp:eq:C.RK.8}}}
\expandafter\def\csname SRef@cdp:eq:C.RK.18\endcsname{\hyperlink{tag-00ZZ}{II.\ref*{cdp:eq:C.RK.18}}}
\expandafter\def\csname SRef@cdp:eq:C.RK.19\endcsname{\hyperlink{tag-0100}{II.\ref*{cdp:eq:C.RK.19}}}
\expandafter\def\csname SRef@cdp:eq:C.RK.20\endcsname{\hyperlink{tag-0102}{II.\ref*{cdp:eq:C.RK.20}}}
\expandafter\def\csname SRef@cdp:eq:C.RK.9\endcsname{\hyperlink{tag-010B}{II.\ref*{cdp:eq:C.RK.9}}}
\expandafter\def\csname SRef@cdp:eq:C.RK.10\endcsname{\hyperlink{tag-00ZR}{II.\ref*{cdp:eq:C.RK.10}}}
\expandafter\def\csname SRef@cdp:eq:C.RK.11\endcsname{\hyperlink{tag-00ZS}{II.\ref*{cdp:eq:C.RK.11}}}
\expandafter\def\csname SRef@cdp:eq:C.RK.12\endcsname{\hyperlink{tag-00ZT}{II.\ref*{cdp:eq:C.RK.12}}}
\expandafter\def\csname SRef@cdp:eq:C.RK.13\endcsname{\hyperlink{tag-00ZU}{II.\ref*{cdp:eq:C.RK.13}}}
\expandafter\def\csname SRef@cdp:eq:C.RG.11\endcsname{\hyperlink{tag-00Z5}{II.\ref*{cdp:eq:C.RG.11}}}
\expandafter\def\csname SRef@cdp:eq:C.RK.14\endcsname{\hyperlink{tag-00ZV}{II.\ref*{cdp:eq:C.RK.14}}}
\expandafter\def\csname SRef@cdp:eq:C.RK.15\endcsname{\hyperlink{tag-00ZW}{II.\ref*{cdp:eq:C.RK.15}}}
\expandafter\def\csname SRef@cdp:eq:C.RK.16\endcsname{\hyperlink{tag-00ZX}{II.\ref*{cdp:eq:C.RK.16}}}
\expandafter\def\csname SRef@cdp:eq:C.RK.17\endcsname{\hyperlink{tag-00ZY}{II.\ref*{cdp:eq:C.RK.17}}}
\expandafter\def\csname SRef@cdp:eq:C.RK.21\endcsname{\hyperlink{tag-0103}{II.\ref*{cdp:eq:C.RK.21}}}
\expandafter\def\csname SRef@cdp:eq:C.RK.22\endcsname{\hyperlink{tag-0104}{II.\ref*{cdp:eq:C.RK.22}}}
\expandafter\def\csname SRef@cdp:eq:C.NT.1\endcsname{\hyperlink{tag-00TU}{II.\ref*{cdp:eq:C.NT.1}}}
\expandafter\def\csname SRef@cdp:eq:C.NT.2\endcsname{\hyperlink{tag-00TV}{II.\ref*{cdp:eq:C.NT.2}}}
\expandafter\def\csname SRef@cdp:eq:C.NT.3\endcsname{\hyperlink{tag-00TW}{II.\ref*{cdp:eq:C.NT.3}}}
\expandafter\def\csname SRef@cdp:eq:C.NT.4\endcsname{\hyperlink{tag-00TX}{II.\ref*{cdp:eq:C.NT.4}}}
\expandafter\def\csname SRef@cdp:eq:C.RG.12\endcsname{\hyperlink{tag-00Z6}{II.\ref*{cdp:eq:C.RG.12}}}
\expandafter\def\csname SRef@cdp:eq:C.NT.5\endcsname{\hyperlink{tag-00TY}{II.\ref*{cdp:eq:C.NT.5}}}
\expandafter\def\csname SRef@cdp:eq:C.NT.6\endcsname{\hyperlink{tag-00TZ}{II.\ref*{cdp:eq:C.NT.6}}}
\expandafter\def\csname SRef@cdp:eq:C.NT.7\endcsname{\hyperlink{tag-00U0}{II.\ref*{cdp:eq:C.NT.7}}}
\expandafter\def\csname SRef@cdp:eq:C.NT.8\endcsname{\hyperlink{tag-00U1}{II.\ref*{cdp:eq:C.NT.8}}}
\expandafter\def\csname SRef@cdp:eq:C.GP.1\endcsname{\hyperlink{tag-00KY}{II.\ref*{cdp:eq:C.GP.1}}}
\expandafter\def\csname SRef@cdp:eq:C.GP.2\endcsname{\hyperlink{tag-00L0}{II.\ref*{cdp:eq:C.GP.2}}}
\expandafter\def\csname SRef@cdp:eq:C.GP.3\endcsname{\hyperlink{tag-00L1}{II.\ref*{cdp:eq:C.GP.3}}}
\expandafter\def\csname SRef@cdp:eq:C.GP.4\endcsname{\hyperlink{tag-00L2}{II.\ref*{cdp:eq:C.GP.4}}}
\expandafter\def\csname SRef@cdp:eq:C.GP.5\endcsname{\hyperlink{tag-00L3}{II.\ref*{cdp:eq:C.GP.5}}}
\expandafter\def\csname SRef@cdp:eq:C.GP.6\endcsname{\hyperlink{tag-00L4}{II.\ref*{cdp:eq:C.GP.6}}}
\expandafter\def\csname SRef@cdp:eq:C.RG.13\endcsname{\hyperlink{tag-00Z7}{II.\ref*{cdp:eq:C.RG.13}}}
\expandafter\def\csname SRef@cdp:eq:C.GP.7\endcsname{\hyperlink{tag-00L5}{II.\ref*{cdp:eq:C.GP.7}}}
\expandafter\def\csname SRef@cdp:eq:C.GP.8\endcsname{\hyperlink{tag-00L6}{II.\ref*{cdp:eq:C.GP.8}}}
\expandafter\def\csname SRef@cdp:eq:C.GP.9\endcsname{\hyperlink{tag-00L7}{II.\ref*{cdp:eq:C.GP.9}}}
\expandafter\def\csname SRef@cdp:eq:C.GP.10\endcsname{\hyperlink{tag-00KZ}{II.\ref*{cdp:eq:C.GP.10}}}
\expandafter\def\csname SRef@cdp:eq:C.AR.1\endcsname{\hyperlink{tag-00A0}{II.\ref*{cdp:eq:C.AR.1}}}
\expandafter\def\csname SRef@cdp:eq:C.AR.2\endcsname{\hyperlink{tag-00A8}{II.\ref*{cdp:eq:C.AR.2}}}
\expandafter\def\csname SRef@cdp:eq:C.AR.3\endcsname{\hyperlink{tag-00A9}{II.\ref*{cdp:eq:C.AR.3}}}
\expandafter\def\csname SRef@cdp:eq:C.AR.4\endcsname{\hyperlink{tag-00AA}{II.\ref*{cdp:eq:C.AR.4}}}
\expandafter\def\csname SRef@cdp:eq:C.AR.5\endcsname{\hyperlink{tag-00AB}{II.\ref*{cdp:eq:C.AR.5}}}
\expandafter\def\csname SRef@cdp:eq:C.CS.1\endcsname{\hyperlink{tag-00DH}{II.\ref*{cdp:eq:C.CS.1}}}
\expandafter\def\csname SRef@cdp:eq:C.AR.6\endcsname{\hyperlink{tag-00AC}{II.\ref*{cdp:eq:C.AR.6}}}
\expandafter\def\csname SRef@cdp:eq:C.AR.7\endcsname{\hyperlink{tag-00AD}{II.\ref*{cdp:eq:C.AR.7}}}
\expandafter\def\csname SRef@cdp:eq:C.AR.8\endcsname{\hyperlink{tag-00AE}{II.\ref*{cdp:eq:C.AR.8}}}
\expandafter\def\csname SRef@cdp:eq:C.AR.9\endcsname{\hyperlink{tag-00AF}{II.\ref*{cdp:eq:C.AR.9}}}
\expandafter\def\csname SRef@cdp:eq:C.AR.10\endcsname{\hyperlink{tag-00A1}{II.\ref*{cdp:eq:C.AR.10}}}
\expandafter\def\csname SRef@cdp:eq:C.AR.11\endcsname{\hyperlink{tag-00A2}{II.\ref*{cdp:eq:C.AR.11}}}
\expandafter\def\csname SRef@cdp:eq:C.AR.12\endcsname{\hyperlink{tag-00A3}{II.\ref*{cdp:eq:C.AR.12}}}
\expandafter\def\csname SRef@cdp:eq:C.AR.13\endcsname{\hyperlink{tag-00A4}{II.\ref*{cdp:eq:C.AR.13}}}
\expandafter\def\csname SRef@cdp:eq:C.AR.14\endcsname{\hyperlink{tag-00A5}{II.\ref*{cdp:eq:C.AR.14}}}
\expandafter\def\csname SRef@cdp:eq:C.AR.15\endcsname{\hyperlink{tag-00A6}{II.\ref*{cdp:eq:C.AR.15}}}
\expandafter\def\csname SRef@cdp:eq:C.CS.2\endcsname{\hyperlink{tag-00DI}{II.\ref*{cdp:eq:C.CS.2}}}
\expandafter\def\csname SRef@cdp:eq:C.AR.16\endcsname{\hyperlink{tag-00A7}{II.\ref*{cdp:eq:C.AR.16}}}
\expandafter\def\csname SRef@cdp:eq:C.G.1\endcsname{\hyperlink{tag-00KW}{II.\ref*{cdp:eq:C.G.1}}}
\expandafter\def\csname SRef@cdp:eq:C.G.2\endcsname{\hyperlink{tag-00KX}{II.\ref*{cdp:eq:C.G.2}}}
\expandafter\def\csname SRef@cdp:eq:C.CS.3\endcsname{\hyperlink{tag-00DJ}{II.\ref*{cdp:eq:C.CS.3}}}
\expandafter\def\csname SRef@cdp:eq:C.UN.4\endcsname{\hyperlink{tag-0155}{II.\ref*{cdp:eq:C.UN.4}}}
\expandafter\def\csname SRef@cdp:eq:C.CS.4\endcsname{\hyperlink{tag-00DK}{II.\ref*{cdp:eq:C.CS.4}}}
\expandafter\def\csname SRef@cdp:eq:C.CS.5\endcsname{\hyperlink{tag-00DL}{II.\ref*{cdp:eq:C.CS.5}}}
\expandafter\def\csname SRef@cdp:eq:C.CS.6\endcsname{\hyperlink{tag-00DM}{II.\ref*{cdp:eq:C.CS.6}}}
\expandafter\def\csname SRef@cdp:eq:C.EQ.1\endcsname{\hyperlink{tag-00IC}{II.\ref*{cdp:eq:C.EQ.1}}}
\expandafter\def\csname SRef@cdp:eq:C.EQ.2\endcsname{\hyperlink{tag-00IF}{II.\ref*{cdp:eq:C.EQ.2}}}
\expandafter\def\csname SRef@cdp:eq:C.EQ.3\endcsname{\hyperlink{tag-00IG}{II.\ref*{cdp:eq:C.EQ.3}}}
\expandafter\def\csname SRef@cdp:eq:C.EQ.4\endcsname{\hyperlink{tag-00IH}{II.\ref*{cdp:eq:C.EQ.4}}}
\expandafter\def\csname SRef@cdp:eq:C.EQ.5\endcsname{\hyperlink{tag-00II}{II.\ref*{cdp:eq:C.EQ.5}}}
\expandafter\def\csname SRef@cdp:eq:C.EQ.6\endcsname{\hyperlink{tag-00IJ}{II.\ref*{cdp:eq:C.EQ.6}}}
\expandafter\def\csname SRef@cdp:eq:C.EQ.7\endcsname{\hyperlink{tag-00IK}{II.\ref*{cdp:eq:C.EQ.7}}}
\expandafter\def\csname SRef@cdp:eq:C.UN.5\endcsname{\hyperlink{tag-0156}{II.\ref*{cdp:eq:C.UN.5}}}
\expandafter\def\csname SRef@cdp:eq:C.EQ.8\endcsname{\hyperlink{tag-00IL}{II.\ref*{cdp:eq:C.EQ.8}}}
\expandafter\def\csname SRef@cdp:eq:C.EQ.9\endcsname{\hyperlink{tag-00IM}{II.\ref*{cdp:eq:C.EQ.9}}}
\expandafter\def\csname SRef@cdp:eq:C.EQ.10\endcsname{\hyperlink{tag-00ID}{II.\ref*{cdp:eq:C.EQ.10}}}
\expandafter\def\csname SRef@cdp:eq:C.EQ.11\endcsname{\hyperlink{tag-00IE}{II.\ref*{cdp:eq:C.EQ.11}}}
\expandafter\def\csname SRef@cdp:eq:C.CL.1\endcsname{\hyperlink{tag-00BW}{II.\ref*{cdp:eq:C.CL.1}}}
\expandafter\def\csname SRef@cdp:eq:C.CL.2\endcsname{\hyperlink{tag-00C1}{II.\ref*{cdp:eq:C.CL.2}}}
\expandafter\def\csname SRef@cdp:eq:C.CL.3\endcsname{\hyperlink{tag-00C2}{II.\ref*{cdp:eq:C.CL.3}}}
\expandafter\def\csname SRef@cdp:eq:C.CL.4\endcsname{\hyperlink{tag-00C3}{II.\ref*{cdp:eq:C.CL.4}}}
\expandafter\def\csname SRef@cdp:eq:C.CL.5\endcsname{\hyperlink{tag-00C4}{II.\ref*{cdp:eq:C.CL.5}}}
\expandafter\def\csname SRef@cdp:eq:C.CL.6\endcsname{\hyperlink{tag-00C5}{II.\ref*{cdp:eq:C.CL.6}}}
\expandafter\def\csname SRef@cdp:eq:C.UN.6\endcsname{\hyperlink{tag-0157}{II.\ref*{cdp:eq:C.UN.6}}}
\expandafter\def\csname SRef@cdp:eq:C.CL.7\endcsname{\hyperlink{tag-00C6}{II.\ref*{cdp:eq:C.CL.7}}}
\expandafter\def\csname SRef@cdp:eq:C.CL.8\endcsname{\hyperlink{tag-00C7}{II.\ref*{cdp:eq:C.CL.8}}}
\expandafter\def\csname SRef@cdp:eq:C.CL.9\endcsname{\hyperlink{tag-00C8}{II.\ref*{cdp:eq:C.CL.9}}}
\expandafter\def\csname SRef@cdp:eq:C.CL.10\endcsname{\hyperlink{tag-00BX}{II.\ref*{cdp:eq:C.CL.10}}}
\expandafter\def\csname SRef@cdp:eq:C.CL.11\endcsname{\hyperlink{tag-00BY}{II.\ref*{cdp:eq:C.CL.11}}}
\expandafter\def\csname SRef@cdp:eq:C.CL.12\endcsname{\hyperlink{tag-00BZ}{II.\ref*{cdp:eq:C.CL.12}}}
\expandafter\def\csname SRef@cdp:eq:C.CL.13\endcsname{\hyperlink{tag-00C0}{II.\ref*{cdp:eq:C.CL.13}}}
\expandafter\def\csname SRef@cdp:eq:C.HB.1\endcsname{\hyperlink{tag-00LF}{II.\ref*{cdp:eq:C.HB.1}}}
\expandafter\def\csname SRef@cdp:eq:C.HB.2\endcsname{\hyperlink{tag-00LG}{II.\ref*{cdp:eq:C.HB.2}}}
\expandafter\def\csname SRef@cdp:eq:C.HB.3\endcsname{\hyperlink{tag-00LH}{II.\ref*{cdp:eq:C.HB.3}}}
\expandafter\def\csname SRef@cdp:eq:C.UC.1\endcsname{\hyperlink{tag-014C}{II.\ref*{cdp:eq:C.UC.1}}}
\expandafter\def\csname SRef@cdp:eq:C.HB.4\endcsname{\hyperlink{tag-00LI}{II.\ref*{cdp:eq:C.HB.4}}}
\expandafter\def\csname SRef@cdp:eq:C.HB.5\endcsname{\hyperlink{tag-00LJ}{II.\ref*{cdp:eq:C.HB.5}}}
\expandafter\def\csname SRef@cdp:eq:C.HB.6\endcsname{\hyperlink{tag-00LK}{II.\ref*{cdp:eq:C.HB.6}}}
\expandafter\def\csname SRef@cdp:eq:C.HB.7\endcsname{\hyperlink{tag-00LL}{II.\ref*{cdp:eq:C.HB.7}}}
\expandafter\def\csname SRef@cdp:eq:C.HB.8\endcsname{\hyperlink{tag-00LM}{II.\ref*{cdp:eq:C.HB.8}}}
\expandafter\def\csname SRef@cdp:eq:C.HB.9\endcsname{\hyperlink{tag-00LN}{II.\ref*{cdp:eq:C.HB.9}}}
\expandafter\def\csname SRef@cdp:eq:C.ZC.1\endcsname{\hyperlink{tag-018W}{II.\ref*{cdp:eq:C.ZC.1}}}
\expandafter\def\csname SRef@cdp:eq:C.ZC.2\endcsname{\hyperlink{tag-018X}{II.\ref*{cdp:eq:C.ZC.2}}}
\expandafter\def\csname SRef@cdp:eq:C.ZC.3\endcsname{\hyperlink{tag-018Y}{II.\ref*{cdp:eq:C.ZC.3}}}
\expandafter\def\csname SRef@cdp:eq:C.ZC.4\endcsname{\hyperlink{tag-018Z}{II.\ref*{cdp:eq:C.ZC.4}}}
\expandafter\def\csname SRef@cdp:eq:C.UC.2\endcsname{\hyperlink{tag-014E}{II.\ref*{cdp:eq:C.UC.2}}}
\expandafter\def\csname SRef@cdp:eq:C.ZC.5\endcsname{\hyperlink{tag-0190}{II.\ref*{cdp:eq:C.ZC.5}}}
\expandafter\def\csname SRef@cdp:eq:C.ZC.6\endcsname{\hyperlink{tag-0191}{II.\ref*{cdp:eq:C.ZC.6}}}
\expandafter\def\csname SRef@cdp:eq:C.ZC.7\endcsname{\hyperlink{tag-0192}{II.\ref*{cdp:eq:C.ZC.7}}}
\expandafter\def\csname SRef@cdp:eq:C.ZC.8\endcsname{\hyperlink{tag-0193}{II.\ref*{cdp:eq:C.ZC.8}}}
\expandafter\def\csname SRef@cdp:eq:C.ZC.9\endcsname{\hyperlink{tag-0194}{II.\ref*{cdp:eq:C.ZC.9}}}
\expandafter\def\csname SRef@cdp:eq:C.FG.1\endcsname{\hyperlink{tag-00K8}{II.\ref*{cdp:eq:C.FG.1}}}
\expandafter\def\csname SRef@cdp:eq:C.FG.2\endcsname{\hyperlink{tag-00KJ}{II.\ref*{cdp:eq:C.FG.2}}}
\expandafter\def\csname SRef@cdp:eq:C.FG.3\endcsname{\hyperlink{tag-00KK}{II.\ref*{cdp:eq:C.FG.3}}}
\expandafter\def\csname SRef@cdp:eq:C.FG.4\endcsname{\hyperlink{tag-00KL}{II.\ref*{cdp:eq:C.FG.4}}}
\expandafter\def\csname SRef@cdp:eq:C.FG.5\endcsname{\hyperlink{tag-00KM}{II.\ref*{cdp:eq:C.FG.5}}}
\expandafter\def\csname SRef@cdp:eq:C.UC.3\endcsname{\hyperlink{tag-014F}{II.\ref*{cdp:eq:C.UC.3}}}
\expandafter\def\csname SRef@cdp:eq:C.FG.6\endcsname{\hyperlink{tag-00KN}{II.\ref*{cdp:eq:C.FG.6}}}
\expandafter\def\csname SRef@cdp:eq:C.FG.7\endcsname{\hyperlink{tag-00KO}{II.\ref*{cdp:eq:C.FG.7}}}
\expandafter\def\csname SRef@cdp:eq:C.FG.8\endcsname{\hyperlink{tag-00KP}{II.\ref*{cdp:eq:C.FG.8}}}
\expandafter\def\csname SRef@cdp:eq:C.FG.9\endcsname{\hyperlink{tag-00KQ}{II.\ref*{cdp:eq:C.FG.9}}}
\expandafter\def\csname SRef@cdp:eq:C.FG.10\endcsname{\hyperlink{tag-00K9}{II.\ref*{cdp:eq:C.FG.10}}}
\expandafter\def\csname SRef@cdp:eq:C.FG.11\endcsname{\hyperlink{tag-00KA}{II.\ref*{cdp:eq:C.FG.11}}}
\expandafter\def\csname SRef@cdp:eq:C.FG.12\endcsname{\hyperlink{tag-00KB}{II.\ref*{cdp:eq:C.FG.12}}}
\expandafter\def\csname SRef@cdp:eq:C.FG.13\endcsname{\hyperlink{tag-00KC}{II.\ref*{cdp:eq:C.FG.13}}}
\expandafter\def\csname SRef@cdp:eq:C.FG.14\endcsname{\hyperlink{tag-00KD}{II.\ref*{cdp:eq:C.FG.14}}}
\expandafter\def\csname SRef@cdp:eq:C.FG.15\endcsname{\hyperlink{tag-00KE}{II.\ref*{cdp:eq:C.FG.15}}}
\expandafter\def\csname SRef@cdp:eq:C.UC.4\endcsname{\hyperlink{tag-014G}{II.\ref*{cdp:eq:C.UC.4}}}
\expandafter\def\csname SRef@cdp:eq:C.FG.16\endcsname{\hyperlink{tag-00KF}{II.\ref*{cdp:eq:C.FG.16}}}
\expandafter\def\csname SRef@cdp:eq:C.FG.17\endcsname{\hyperlink{tag-00KG}{II.\ref*{cdp:eq:C.FG.17}}}
\expandafter\def\csname SRef@cdp:eq:C.FG.18\endcsname{\hyperlink{tag-00KH}{II.\ref*{cdp:eq:C.FG.18}}}
\expandafter\def\csname SRef@cdp:eq:C.FG.19\endcsname{\hyperlink{tag-00KI}{II.\ref*{cdp:eq:C.FG.19}}}
\expandafter\def\csname SRef@cdp:eq:C.PF.1\endcsname{\hyperlink{tag-00UN}{II.\ref*{cdp:eq:C.PF.1}}}
\expandafter\def\csname SRef@cdp:eq:C.PF.2\endcsname{\hyperlink{tag-00UR}{II.\ref*{cdp:eq:C.PF.2}}}
\expandafter\def\csname SRef@cdp:eq:C.PF.3\endcsname{\hyperlink{tag-00US}{II.\ref*{cdp:eq:C.PF.3}}}
\expandafter\def\csname SRef@cdp:eq:C.PF.4\endcsname{\hyperlink{tag-00UT}{II.\ref*{cdp:eq:C.PF.4}}}
\expandafter\def\csname SRef@cdp:eq:C.PF.5\endcsname{\hyperlink{tag-00UU}{II.\ref*{cdp:eq:C.PF.5}}}
\expandafter\def\csname SRef@cdp:eq:C.PF.6\endcsname{\hyperlink{tag-00UV}{II.\ref*{cdp:eq:C.PF.6}}}
\expandafter\def\csname SRef@cdp:eq:C.UC.5\endcsname{\hyperlink{tag-014H}{II.\ref*{cdp:eq:C.UC.5}}}
\expandafter\def\csname SRef@cdp:eq:C.PF.7\endcsname{\hyperlink{tag-00UW}{II.\ref*{cdp:eq:C.PF.7}}}
\expandafter\def\csname SRef@cdp:eq:C.PF.8\endcsname{\hyperlink{tag-00UX}{II.\ref*{cdp:eq:C.PF.8}}}
\expandafter\def\csname SRef@cdp:eq:C.PF.9\endcsname{\hyperlink{tag-00UY}{II.\ref*{cdp:eq:C.PF.9}}}
\expandafter\def\csname SRef@cdp:eq:C.PF.10\endcsname{\hyperlink{tag-00UO}{II.\ref*{cdp:eq:C.PF.10}}}
\expandafter\def\csname SRef@cdp:eq:C.PF.11\endcsname{\hyperlink{tag-00UP}{II.\ref*{cdp:eq:C.PF.11}}}
\expandafter\def\csname SRef@cdp:eq:C.PF.12\endcsname{\hyperlink{tag-00UQ}{II.\ref*{cdp:eq:C.PF.12}}}
\expandafter\def\csname SRef@cdp:eq:C.UNN.1\endcsname{\hyperlink{tag-015B}{II.\ref*{cdp:eq:C.UNN.1}}}
\expandafter\def\csname SRef@cdp:eq:C.UNN.2\endcsname{\hyperlink{tag-015C}{II.\ref*{cdp:eq:C.UNN.2}}}
\expandafter\def\csname SRef@cdp:eq:C.CR.1\endcsname{\hyperlink{tag-00CV}{II.\ref*{cdp:eq:C.CR.1}}}
\expandafter\def\csname SRef@cdp:eq:C.CR.2\endcsname{\hyperlink{tag-00D6}{II.\ref*{cdp:eq:C.CR.2}}}
\expandafter\def\csname SRef@cdp:eq:C.UC.6\endcsname{\hyperlink{tag-014I}{II.\ref*{cdp:eq:C.UC.6}}}
\expandafter\def\csname SRef@cdp:eq:C.CR.3\endcsname{\hyperlink{tag-00DA}{II.\ref*{cdp:eq:C.CR.3}}}
\expandafter\def\csname SRef@cdp:eq:C.CR.4\endcsname{\hyperlink{tag-00DB}{II.\ref*{cdp:eq:C.CR.4}}}
\expandafter\def\csname SRef@cdp:eq:C.CR.5\endcsname{\hyperlink{tag-00DC}{II.\ref*{cdp:eq:C.CR.5}}}
\expandafter\def\csname SRef@cdp:eq:C.CR.6\endcsname{\hyperlink{tag-00DD}{II.\ref*{cdp:eq:C.CR.6}}}
\expandafter\def\csname SRef@cdp:eq:C.CR.7\endcsname{\hyperlink{tag-00DE}{II.\ref*{cdp:eq:C.CR.7}}}
\expandafter\def\csname SRef@cdp:eq:C.CR.8\endcsname{\hyperlink{tag-00DF}{II.\ref*{cdp:eq:C.CR.8}}}
\expandafter\def\csname SRef@cdp:eq:C.CR.9\endcsname{\hyperlink{tag-00DG}{II.\ref*{cdp:eq:C.CR.9}}}
\expandafter\def\csname SRef@cdp:eq:C.CR.10\endcsname{\hyperlink{tag-00CW}{II.\ref*{cdp:eq:C.CR.10}}}
\expandafter\def\csname SRef@cdp:eq:C.CR.11\endcsname{\hyperlink{tag-00CX}{II.\ref*{cdp:eq:C.CR.11}}}
\expandafter\def\csname SRef@cdp:eq:C.CR.12\endcsname{\hyperlink{tag-00CY}{II.\ref*{cdp:eq:C.CR.12}}}
\expandafter\def\csname SRef@cdp:eq:C.UC.7\endcsname{\hyperlink{tag-014J}{II.\ref*{cdp:eq:C.UC.7}}}
\expandafter\def\csname SRef@cdp:eq:C.CR.13\endcsname{\hyperlink{tag-00CZ}{II.\ref*{cdp:eq:C.CR.13}}}
\expandafter\def\csname SRef@cdp:eq:C.CR.14\endcsname{\hyperlink{tag-00D0}{II.\ref*{cdp:eq:C.CR.14}}}
\expandafter\def\csname SRef@cdp:eq:C.CR.15\endcsname{\hyperlink{tag-00D1}{II.\ref*{cdp:eq:C.CR.15}}}
\expandafter\def\csname SRef@cdp:eq:C.CR.16\endcsname{\hyperlink{tag-00D2}{II.\ref*{cdp:eq:C.CR.16}}}
\expandafter\def\csname SRef@cdp:eq:C.CR.17\endcsname{\hyperlink{tag-00D3}{II.\ref*{cdp:eq:C.CR.17}}}
\expandafter\def\csname SRef@cdp:eq:C.CR.18\endcsname{\hyperlink{tag-00D4}{II.\ref*{cdp:eq:C.CR.18}}}
\expandafter\def\csname SRef@cdp:eq:C.CR.19\endcsname{\hyperlink{tag-00D5}{II.\ref*{cdp:eq:C.CR.19}}}
\expandafter\def\csname SRef@cdp:eq:C.CR.20\endcsname{\hyperlink{tag-00D7}{II.\ref*{cdp:eq:C.CR.20}}}
\expandafter\def\csname SRef@cdp:eq:C.CR.21\endcsname{\hyperlink{tag-00D8}{II.\ref*{cdp:eq:C.CR.21}}}
\expandafter\def\csname SRef@cdp:eq:C.CR.22\endcsname{\hyperlink{tag-00D9}{II.\ref*{cdp:eq:C.CR.22}}}
\expandafter\def\csname SRef@cdp:eq:C.UC.8\endcsname{\hyperlink{tag-014K}{II.\ref*{cdp:eq:C.UC.8}}}
\expandafter\def\csname SRef@cdp:eq:C.UNNODE.1\endcsname{\hyperlink{tag-015D}{II.\ref*{cdp:eq:C.UNNODE.1}}}
\expandafter\def\csname SRef@cdp:eq:C.UNNODE.2\endcsname{\hyperlink{tag-015E}{II.\ref*{cdp:eq:C.UNNODE.2}}}
\expandafter\def\csname SRef@cdp:eq:C.UNNODE.3\endcsname{\hyperlink{tag-015F}{II.\ref*{cdp:eq:C.UNNODE.3}}}
\expandafter\def\csname SRef@cdp:eq:C.UNCR.1\endcsname{\hyperlink{tag-0158}{II.\ref*{cdp:eq:C.UNCR.1}}}
\expandafter\def\csname SRef@cdp:eq:C.UNCR.2\endcsname{\hyperlink{tag-0159}{II.\ref*{cdp:eq:C.UNCR.2}}}
\expandafter\def\csname SRef@cdp:eq:C.UNCR.3\endcsname{\hyperlink{tag-015A}{II.\ref*{cdp:eq:C.UNCR.3}}}
\expandafter\def\csname SRef@cdp:eq:C.IR.1\endcsname{\hyperlink{tag-00OH}{II.\ref*{cdp:eq:C.IR.1}}}
\expandafter\def\csname SRef@cdp:eq:C.IR.2\endcsname{\hyperlink{tag-00OI}{II.\ref*{cdp:eq:C.IR.2}}}
\expandafter\def\csname SRef@cdp:eq:C.IR.4\endcsname{\hyperlink{tag-00OJ}{II.\ref*{cdp:eq:C.IR.4}}}
\expandafter\def\csname SRef@cdp:eq:C.IR.5\endcsname{\hyperlink{tag-00OK}{II.\ref*{cdp:eq:C.IR.5}}}
\expandafter\def\csname SRef@cdp:eq:C.UC.9\endcsname{\hyperlink{tag-014L}{II.\ref*{cdp:eq:C.UC.9}}}
\expandafter\def\csname SRef@cdp:eq:C.IR.6\endcsname{\hyperlink{tag-00OL}{II.\ref*{cdp:eq:C.IR.6}}}
\expandafter\def\csname SRef@cdp:eq:C.IR.9\endcsname{\hyperlink{tag-00OM}{II.\ref*{cdp:eq:C.IR.9}}}
\expandafter\def\csname SRef@cdp:eq:C.VII.1\endcsname{\hyperlink{tag-015K}{II.\ref*{cdp:eq:C.VII.1}}}
\expandafter\def\csname SRef@cdp:eq:C.VII.2\endcsname{\hyperlink{tag-015O}{II.\ref*{cdp:eq:C.VII.2}}}
\expandafter\def\csname SRef@cdp:eq:C.VII.3\endcsname{\hyperlink{tag-015P}{II.\ref*{cdp:eq:C.VII.3}}}
\expandafter\def\csname SRef@cdp:eq:C.VII.4\endcsname{\hyperlink{tag-015Q}{II.\ref*{cdp:eq:C.VII.4}}}
\expandafter\def\csname SRef@cdp:eq:C.VII.5\endcsname{\hyperlink{tag-015R}{II.\ref*{cdp:eq:C.VII.5}}}
\expandafter\def\csname SRef@cdp:eq:C.VII.6\endcsname{\hyperlink{tag-015S}{II.\ref*{cdp:eq:C.VII.6}}}
\expandafter\def\csname SRef@cdp:eq:C.VII.7\endcsname{\hyperlink{tag-015T}{II.\ref*{cdp:eq:C.VII.7}}}
\expandafter\def\csname SRef@cdp:eq:C.VII.8\endcsname{\hyperlink{tag-015U}{II.\ref*{cdp:eq:C.VII.8}}}
\expandafter\def\csname SRef@cdp:eq:C.UC.10\endcsname{\hyperlink{tag-014D}{II.\ref*{cdp:eq:C.UC.10}}}
\expandafter\def\csname SRef@cdp:eq:C.VII.9\endcsname{\hyperlink{tag-015V}{II.\ref*{cdp:eq:C.VII.9}}}
\expandafter\def\csname SRef@cdp:eq:C.VII.10\endcsname{\hyperlink{tag-015L}{II.\ref*{cdp:eq:C.VII.10}}}
\expandafter\def\csname SRef@cdp:eq:C.VII.11\endcsname{\hyperlink{tag-015M}{II.\ref*{cdp:eq:C.VII.11}}}
\expandafter\def\csname SRef@cdp:eq:C.VII.12\endcsname{\hyperlink{tag-015N}{II.\ref*{cdp:eq:C.VII.12}}}
\expandafter\def\csname SRef@cdp:eq:C.RA.1\endcsname{\hyperlink{tag-00WQ}{II.\ref*{cdp:eq:C.RA.1}}}
\expandafter\def\csname SRef@cdp:eq:C.RA.2\endcsname{\hyperlink{tag-00X1}{II.\ref*{cdp:eq:C.RA.2}}}
\expandafter\def\csname SRef@cdp:eq:C.RA.3\endcsname{\hyperlink{tag-00XC}{II.\ref*{cdp:eq:C.RA.3}}}
\expandafter\def\csname SRef@cdp:eq:C.RA.4\endcsname{\hyperlink{tag-00XH}{II.\ref*{cdp:eq:C.RA.4}}}
\expandafter\def\csname SRef@cdp:eq:C.RA.5\endcsname{\hyperlink{tag-00XI}{II.\ref*{cdp:eq:C.RA.5}}}
\expandafter\def\csname SRef@cdp:eq:C.RA.6\endcsname{\hyperlink{tag-00XJ}{II.\ref*{cdp:eq:C.RA.6}}}
\expandafter\def\csname SRef@cdp:eq:C.VR.1\endcsname{\hyperlink{tag-015W}{II.\ref*{cdp:eq:C.VR.1}}}
\expandafter\def\csname SRef@cdp:eq:C.RA.7\endcsname{\hyperlink{tag-00XK}{II.\ref*{cdp:eq:C.RA.7}}}
\expandafter\def\csname SRef@cdp:eq:C.RA.8\endcsname{\hyperlink{tag-00XL}{II.\ref*{cdp:eq:C.RA.8}}}
\expandafter\def\csname SRef@cdp:eq:C.RA.9\endcsname{\hyperlink{tag-00XM}{II.\ref*{cdp:eq:C.RA.9}}}
\expandafter\def\csname SRef@cdp:eq:C.RA.10\endcsname{\hyperlink{tag-00WR}{II.\ref*{cdp:eq:C.RA.10}}}
\expandafter\def\csname SRef@cdp:eq:C.RA.11\endcsname{\hyperlink{tag-00WS}{II.\ref*{cdp:eq:C.RA.11}}}
\expandafter\def\csname SRef@cdp:eq:C.RA.12\endcsname{\hyperlink{tag-00WT}{II.\ref*{cdp:eq:C.RA.12}}}
\expandafter\def\csname SRef@cdp:eq:C.RA.13\endcsname{\hyperlink{tag-00WU}{II.\ref*{cdp:eq:C.RA.13}}}
\expandafter\def\csname SRef@cdp:eq:C.RA.14\endcsname{\hyperlink{tag-00WV}{II.\ref*{cdp:eq:C.RA.14}}}
\expandafter\def\csname SRef@cdp:eq:C.RA.15\endcsname{\hyperlink{tag-00WW}{II.\ref*{cdp:eq:C.RA.15}}}
\expandafter\def\csname SRef@cdp:eq:C.RA.16\endcsname{\hyperlink{tag-00WX}{II.\ref*{cdp:eq:C.RA.16}}}
\expandafter\def\csname SRef@cdp:eq:C.VR.2\endcsname{\hyperlink{tag-015X}{II.\ref*{cdp:eq:C.VR.2}}}
\expandafter\def\csname SRef@cdp:eq:C.RA.17\endcsname{\hyperlink{tag-00WY}{II.\ref*{cdp:eq:C.RA.17}}}
\expandafter\def\csname SRef@cdp:eq:C.RA.18\endcsname{\hyperlink{tag-00WZ}{II.\ref*{cdp:eq:C.RA.18}}}
\expandafter\def\csname SRef@cdp:eq:C.RA.19\endcsname{\hyperlink{tag-00X0}{II.\ref*{cdp:eq:C.RA.19}}}
\expandafter\def\csname SRef@cdp:eq:C.RA.20\endcsname{\hyperlink{tag-00X2}{II.\ref*{cdp:eq:C.RA.20}}}
\expandafter\def\csname SRef@cdp:eq:C.RA.21\endcsname{\hyperlink{tag-00X3}{II.\ref*{cdp:eq:C.RA.21}}}
\expandafter\def\csname SRef@cdp:eq:C.RA.22\endcsname{\hyperlink{tag-00X4}{II.\ref*{cdp:eq:C.RA.22}}}
\expandafter\def\csname SRef@cdp:eq:C.RA.23\endcsname{\hyperlink{tag-00X5}{II.\ref*{cdp:eq:C.RA.23}}}
\expandafter\def\csname SRef@cdp:eq:C.RA.24\endcsname{\hyperlink{tag-00X6}{II.\ref*{cdp:eq:C.RA.24}}}
\expandafter\def\csname SRef@cdp:eq:C.RA.25\endcsname{\hyperlink{tag-00X7}{II.\ref*{cdp:eq:C.RA.25}}}
\expandafter\def\csname SRef@cdp:eq:C.RA.26\endcsname{\hyperlink{tag-00X8}{II.\ref*{cdp:eq:C.RA.26}}}
\expandafter\def\csname SRef@cdp:eq:C.VR.3\endcsname{\hyperlink{tag-015Y}{II.\ref*{cdp:eq:C.VR.3}}}
\expandafter\def\csname SRef@cdp:eq:C.RA.27\endcsname{\hyperlink{tag-00X9}{II.\ref*{cdp:eq:C.RA.27}}}
\expandafter\def\csname SRef@cdp:eq:C.RA.28\endcsname{\hyperlink{tag-00XA}{II.\ref*{cdp:eq:C.RA.28}}}
\expandafter\def\csname SRef@cdp:eq:C.RA.29\endcsname{\hyperlink{tag-00XB}{II.\ref*{cdp:eq:C.RA.29}}}
\expandafter\def\csname SRef@cdp:eq:C.RA.30\endcsname{\hyperlink{tag-00XD}{II.\ref*{cdp:eq:C.RA.30}}}
\expandafter\def\csname SRef@cdp:eq:C.RA.31\endcsname{\hyperlink{tag-00XE}{II.\ref*{cdp:eq:C.RA.31}}}
\expandafter\def\csname SRef@cdp:eq:C.RA.31a\endcsname{\hyperlink{tag-00XF}{II.\ref*{cdp:eq:C.RA.31a}}}
\expandafter\def\csname SRef@cdp:eq:C.RA.32\endcsname{\hyperlink{tag-00XG}{II.\ref*{cdp:eq:C.RA.32}}}
\expandafter\def\csname SRef@cdp:eq:C.RD.1\endcsname{\hyperlink{tag-00XZ}{II.\ref*{cdp:eq:C.RD.1}}}
\expandafter\def\csname SRef@cdp:eq:C.RD.2\endcsname{\hyperlink{tag-00Y9}{II.\ref*{cdp:eq:C.RD.2}}}
\expandafter\def\csname SRef@cdp:eq:C.RD.3\endcsname{\hyperlink{tag-00YK}{II.\ref*{cdp:eq:C.RD.3}}}
\expandafter\def\csname SRef@cdp:eq:C.VR.4\endcsname{\hyperlink{tag-015Z}{II.\ref*{cdp:eq:C.VR.4}}}
\expandafter\def\csname SRef@cdp:eq:C.RD.4\endcsname{\hyperlink{tag-00YT}{II.\ref*{cdp:eq:C.RD.4}}}
\expandafter\def\csname SRef@cdp:eq:C.RD.11\endcsname{\hyperlink{tag-00Y0}{II.\ref*{cdp:eq:C.RD.11}}}
\expandafter\def\csname SRef@cdp:eq:C.RD.12\endcsname{\hyperlink{tag-00Y1}{II.\ref*{cdp:eq:C.RD.12}}}
\expandafter\def\csname SRef@cdp:eq:C.RD.13\endcsname{\hyperlink{tag-00Y2}{II.\ref*{cdp:eq:C.RD.13}}}
\expandafter\def\csname SRef@cdp:eq:C.RD.14\endcsname{\hyperlink{tag-00Y3}{II.\ref*{cdp:eq:C.RD.14}}}
\expandafter\def\csname SRef@cdp:eq:C.RD.15\endcsname{\hyperlink{tag-00Y4}{II.\ref*{cdp:eq:C.RD.15}}}
\expandafter\def\csname SRef@cdp:eq:C.RD.16\endcsname{\hyperlink{tag-00Y5}{II.\ref*{cdp:eq:C.RD.16}}}
\expandafter\def\csname SRef@cdp:eq:C.RD.17\endcsname{\hyperlink{tag-00Y6}{II.\ref*{cdp:eq:C.RD.17}}}
\expandafter\def\csname SRef@cdp:eq:C.RD.18\endcsname{\hyperlink{tag-00Y7}{II.\ref*{cdp:eq:C.RD.18}}}
\expandafter\def\csname SRef@cdp:eq:C.VR.5\endcsname{\hyperlink{tag-0160}{II.\ref*{cdp:eq:C.VR.5}}}
\expandafter\def\csname SRef@cdp:eq:C.RD.19\endcsname{\hyperlink{tag-00Y8}{II.\ref*{cdp:eq:C.RD.19}}}
\expandafter\def\csname SRef@cdp:eq:C.RD.20\endcsname{\hyperlink{tag-00YA}{II.\ref*{cdp:eq:C.RD.20}}}
\expandafter\def\csname SRef@cdp:eq:C.RD.21\endcsname{\hyperlink{tag-00YB}{II.\ref*{cdp:eq:C.RD.21}}}
\expandafter\def\csname SRef@cdp:eq:C.RD.22\endcsname{\hyperlink{tag-00YC}{II.\ref*{cdp:eq:C.RD.22}}}
\expandafter\def\csname SRef@cdp:eq:C.RD.23\endcsname{\hyperlink{tag-00YD}{II.\ref*{cdp:eq:C.RD.23}}}
\expandafter\def\csname SRef@cdp:eq:C.RD.24\endcsname{\hyperlink{tag-00YE}{II.\ref*{cdp:eq:C.RD.24}}}
\expandafter\def\csname SRef@cdp:eq:C.RD.25\endcsname{\hyperlink{tag-00YF}{II.\ref*{cdp:eq:C.RD.25}}}
\expandafter\def\csname SRef@cdp:eq:C.RD.26\endcsname{\hyperlink{tag-00YG}{II.\ref*{cdp:eq:C.RD.26}}}
\expandafter\def\csname SRef@cdp:eq:C.RD.27\endcsname{\hyperlink{tag-00YH}{II.\ref*{cdp:eq:C.RD.27}}}
\expandafter\def\csname SRef@cdp:eq:C.VR.6\endcsname{\hyperlink{tag-0161}{II.\ref*{cdp:eq:C.VR.6}}}
\expandafter\def\csname SRef@cdp:eq:C.RD.28\endcsname{\hyperlink{tag-00YI}{II.\ref*{cdp:eq:C.RD.28}}}
\expandafter\def\csname SRef@cdp:eq:C.RD.29\endcsname{\hyperlink{tag-00YJ}{II.\ref*{cdp:eq:C.RD.29}}}
\expandafter\def\csname SRef@cdp:eq:C.RD.30\endcsname{\hyperlink{tag-00YL}{II.\ref*{cdp:eq:C.RD.30}}}
\expandafter\def\csname SRef@cdp:eq:C.RD.31\endcsname{\hyperlink{tag-00YM}{II.\ref*{cdp:eq:C.RD.31}}}
\expandafter\def\csname SRef@cdp:eq:C.RD.32\endcsname{\hyperlink{tag-00YN}{II.\ref*{cdp:eq:C.RD.32}}}
\expandafter\def\csname SRef@cdp:eq:C.RD.33\endcsname{\hyperlink{tag-00YO}{II.\ref*{cdp:eq:C.RD.33}}}
\expandafter\def\csname SRef@cdp:eq:C.RD.34\endcsname{\hyperlink{tag-00YP}{II.\ref*{cdp:eq:C.RD.34}}}
\expandafter\def\csname SRef@cdp:eq:C.RD.35\endcsname{\hyperlink{tag-00YQ}{II.\ref*{cdp:eq:C.RD.35}}}
\expandafter\def\csname SRef@cdp:eq:C.VR.7\endcsname{\hyperlink{tag-0162}{II.\ref*{cdp:eq:C.VR.7}}}
\expandafter\def\csname SRef@cdp:eq:C.RD.36\endcsname{\hyperlink{tag-00YR}{II.\ref*{cdp:eq:C.RD.36}}}
\expandafter\def\csname SRef@cdp:eq:C.RD.37\endcsname{\hyperlink{tag-00YS}{II.\ref*{cdp:eq:C.RD.37}}}
\expandafter\def\csname SRef@cdp:eq:C.D4.1\endcsname{\hyperlink{tag-00E1}{II.\ref*{cdp:eq:C.D4.1}}}
\expandafter\def\csname SRef@cdp:eq:C.D4.2\endcsname{\hyperlink{tag-00EC}{II.\ref*{cdp:eq:C.D4.2}}}
\expandafter\def\csname SRef@cdp:eq:C.D4.3\endcsname{\hyperlink{tag-00EE}{II.\ref*{cdp:eq:C.D4.3}}}
\expandafter\def\csname SRef@cdp:eq:C.D4.4\endcsname{\hyperlink{tag-00EF}{II.\ref*{cdp:eq:C.D4.4}}}
\expandafter\def\csname SRef@cdp:eq:C.D4.5\endcsname{\hyperlink{tag-00EG}{II.\ref*{cdp:eq:C.D4.5}}}
\expandafter\def\csname SRef@cdp:eq:C.VR.8\endcsname{\hyperlink{tag-0163}{II.\ref*{cdp:eq:C.VR.8}}}
\expandafter\def\csname SRef@cdp:eq:C.D4.6\endcsname{\hyperlink{tag-00EH}{II.\ref*{cdp:eq:C.D4.6}}}
\expandafter\def\csname SRef@cdp:eq:C.D4.7\endcsname{\hyperlink{tag-00EI}{II.\ref*{cdp:eq:C.D4.7}}}
\expandafter\def\csname SRef@cdp:eq:C.D4.8\endcsname{\hyperlink{tag-00EJ}{II.\ref*{cdp:eq:C.D4.8}}}
\expandafter\def\csname SRef@cdp:eq:C.D4.9\endcsname{\hyperlink{tag-00EK}{II.\ref*{cdp:eq:C.D4.9}}}
\expandafter\def\csname SRef@cdp:eq:C.D4.10\endcsname{\hyperlink{tag-00E2}{II.\ref*{cdp:eq:C.D4.10}}}
\expandafter\def\csname SRef@cdp:eq:C.D4.11\endcsname{\hyperlink{tag-00E3}{II.\ref*{cdp:eq:C.D4.11}}}
\expandafter\def\csname SRef@cdp:eq:C.D4.12\endcsname{\hyperlink{tag-00E4}{II.\ref*{cdp:eq:C.D4.12}}}
\expandafter\def\csname SRef@cdp:eq:C.DD.1\endcsname{\hyperlink{tag-00EL}{II.\ref*{cdp:eq:C.DD.1}}}
\expandafter\def\csname SRef@cdp:eq:C.DD.2\endcsname{\hyperlink{tag-00EW}{II.\ref*{cdp:eq:C.DD.2}}}
\expandafter\def\csname SRef@cdp:eq:C.VR.9\endcsname{\hyperlink{tag-0164}{II.\ref*{cdp:eq:C.VR.9}}}
\expandafter\def\csname SRef@cdp:eq:C.DD.3\endcsname{\hyperlink{tag-00EZ}{II.\ref*{cdp:eq:C.DD.3}}}
\expandafter\def\csname SRef@cdp:eq:C.DD.4\endcsname{\hyperlink{tag-00F0}{II.\ref*{cdp:eq:C.DD.4}}}
\expandafter\def\csname SRef@cdp:eq:C.DD.5\endcsname{\hyperlink{tag-00F1}{II.\ref*{cdp:eq:C.DD.5}}}
\expandafter\def\csname SRef@cdp:eq:C.DD.6\endcsname{\hyperlink{tag-00F2}{II.\ref*{cdp:eq:C.DD.6}}}
\expandafter\def\csname SRef@cdp:eq:C.DD.7\endcsname{\hyperlink{tag-00F3}{II.\ref*{cdp:eq:C.DD.7}}}
\expandafter\def\csname SRef@cdp:eq:C.DD.8\endcsname{\hyperlink{tag-00F4}{II.\ref*{cdp:eq:C.DD.8}}}
\expandafter\def\csname SRef@cdp:eq:C.DD.9\endcsname{\hyperlink{tag-00F5}{II.\ref*{cdp:eq:C.DD.9}}}
\expandafter\def\csname SRef@cdp:eq:C.DD.10\endcsname{\hyperlink{tag-00EM}{II.\ref*{cdp:eq:C.DD.10}}}
\expandafter\def\csname SRef@cdp:eq:C.EC.1\endcsname{\hyperlink{tag-00FO}{II.\ref*{cdp:eq:C.EC.1}}}
\expandafter\def\csname SRef@cdp:eq:C.DD.11\endcsname{\hyperlink{tag-00EN}{II.\ref*{cdp:eq:C.DD.11}}}
\expandafter\def\csname SRef@cdp:eq:C.DD.12\endcsname{\hyperlink{tag-00EO}{II.\ref*{cdp:eq:C.DD.12}}}
\expandafter\def\csname SRef@cdp:eq:C.DD.13\endcsname{\hyperlink{tag-00EP}{II.\ref*{cdp:eq:C.DD.13}}}
\expandafter\def\csname SRef@cdp:eq:C.DD.14\endcsname{\hyperlink{tag-00EQ}{II.\ref*{cdp:eq:C.DD.14}}}
\expandafter\def\csname SRef@cdp:eq:C.DD.15\endcsname{\hyperlink{tag-00ER}{II.\ref*{cdp:eq:C.DD.15}}}
\expandafter\def\csname SRef@cdp:eq:C.DD.16\endcsname{\hyperlink{tag-00ES}{II.\ref*{cdp:eq:C.DD.16}}}
\expandafter\def\csname SRef@cdp:eq:C.EC.2\endcsname{\hyperlink{tag-00FP}{II.\ref*{cdp:eq:C.EC.2}}}
\expandafter\def\csname SRef@cdp:eq:C.DD.17\endcsname{\hyperlink{tag-00ET}{II.\ref*{cdp:eq:C.DD.17}}}
\expandafter\def\csname SRef@cdp:eq:C.DD.18\endcsname{\hyperlink{tag-00EU}{II.\ref*{cdp:eq:C.DD.18}}}
\expandafter\def\csname SRef@cdp:eq:C.DD.19\endcsname{\hyperlink{tag-00EV}{II.\ref*{cdp:eq:C.DD.19}}}
\expandafter\def\csname SRef@cdp:eq:C.DD.20\endcsname{\hyperlink{tag-00EX}{II.\ref*{cdp:eq:C.DD.20}}}
\expandafter\def\csname SRef@cdp:eq:C.DD.21\endcsname{\hyperlink{tag-00EY}{II.\ref*{cdp:eq:C.DD.21}}}
\expandafter\def\csname SRef@cdp:eq:C.D4.13\endcsname{\hyperlink{tag-00E5}{II.\ref*{cdp:eq:C.D4.13}}}
\expandafter\def\csname SRef@cdp:eq:C.D4.14\endcsname{\hyperlink{tag-00E6}{II.\ref*{cdp:eq:C.D4.14}}}
\expandafter\def\csname SRef@cdp:eq:C.D4.15\endcsname{\hyperlink{tag-00E7}{II.\ref*{cdp:eq:C.D4.15}}}
\expandafter\def\csname SRef@cdp:eq:C.SC.1\endcsname{\hyperlink{tag-010L}{II.\ref*{cdp:eq:C.SC.1}}}
\expandafter\def\csname SRef@cdp:eq:C.D4.16\endcsname{\hyperlink{tag-00E8}{II.\ref*{cdp:eq:C.D4.16}}}
\expandafter\def\csname SRef@cdp:eq:C.D4.17\endcsname{\hyperlink{tag-00E9}{II.\ref*{cdp:eq:C.D4.17}}}
\expandafter\def\csname SRef@cdp:eq:C.D4.18\endcsname{\hyperlink{tag-00EA}{II.\ref*{cdp:eq:C.D4.18}}}
\expandafter\def\csname SRef@cdp:eq:C.D4.19\endcsname{\hyperlink{tag-00EB}{II.\ref*{cdp:eq:C.D4.19}}}
\expandafter\def\csname SRef@cdp:eq:C.D4.20\endcsname{\hyperlink{tag-00ED}{II.\ref*{cdp:eq:C.D4.20}}}
\expandafter\def\csname SRef@cdp:eq:C.UI.1\endcsname{\hyperlink{tag-014M}{II.\ref*{cdp:eq:C.UI.1}}}
\expandafter\def\csname SRef@cdp:eq:C.UI.2\endcsname{\hyperlink{tag-014U}{II.\ref*{cdp:eq:C.UI.2}}}
\expandafter\def\csname SRef@cdp:eq:C.UI.3\endcsname{\hyperlink{tag-014V}{II.\ref*{cdp:eq:C.UI.3}}}
\expandafter\def\csname SRef@cdp:eq:C.UI.4\endcsname{\hyperlink{tag-014W}{II.\ref*{cdp:eq:C.UI.4}}}
\expandafter\def\csname SRef@cdp:eq:C.SC.2\endcsname{\hyperlink{tag-010M}{II.\ref*{cdp:eq:C.SC.2}}}
\expandafter\def\csname SRef@cdp:eq:C.UI.5\endcsname{\hyperlink{tag-014X}{II.\ref*{cdp:eq:C.UI.5}}}
\expandafter\def\csname SRef@cdp:eq:C.UI.6\endcsname{\hyperlink{tag-014Y}{II.\ref*{cdp:eq:C.UI.6}}}
\expandafter\def\csname SRef@cdp:eq:C.UI.7\endcsname{\hyperlink{tag-014Z}{II.\ref*{cdp:eq:C.UI.7}}}
\expandafter\def\csname SRef@cdp:eq:C.UI.8\endcsname{\hyperlink{tag-0150}{II.\ref*{cdp:eq:C.UI.8}}}
\expandafter\def\csname SRef@cdp:eq:C.UI.9\endcsname{\hyperlink{tag-0151}{II.\ref*{cdp:eq:C.UI.9}}}
\expandafter\def\csname SRef@cdp:eq:C.UI.10\endcsname{\hyperlink{tag-014N}{II.\ref*{cdp:eq:C.UI.10}}}
\expandafter\def\csname SRef@cdp:eq:C.UI.11\endcsname{\hyperlink{tag-014O}{II.\ref*{cdp:eq:C.UI.11}}}
\expandafter\def\csname SRef@cdp:eq:C.UI.12\endcsname{\hyperlink{tag-014P}{II.\ref*{cdp:eq:C.UI.12}}}
\expandafter\def\csname SRef@cdp:eq:C.UI.13\endcsname{\hyperlink{tag-014Q}{II.\ref*{cdp:eq:C.UI.13}}}
\expandafter\def\csname SRef@cdp:eq:C.UI.14\endcsname{\hyperlink{tag-014R}{II.\ref*{cdp:eq:C.UI.14}}}
\expandafter\def\csname SRef@cdp:eq:C.SC.3\endcsname{\hyperlink{tag-010N}{II.\ref*{cdp:eq:C.SC.3}}}
\expandafter\def\csname SRef@cdp:eq:C.UI.15\endcsname{\hyperlink{tag-014S}{II.\ref*{cdp:eq:C.UI.15}}}
\expandafter\def\csname SRef@cdp:eq:C.UI.16\endcsname{\hyperlink{tag-014T}{II.\ref*{cdp:eq:C.UI.16}}}
\expandafter\def\csname SRef@cdp:eq:C.IZ.1\endcsname{\hyperlink{tag-00ON}{II.\ref*{cdp:eq:C.IZ.1}}}
\expandafter\def\csname SRef@cdp:eq:C.IZ.2\endcsname{\hyperlink{tag-00OO}{II.\ref*{cdp:eq:C.IZ.2}}}
\expandafter\def\csname SRef@cdp:eq:C.IZ.3\endcsname{\hyperlink{tag-00OP}{II.\ref*{cdp:eq:C.IZ.3}}}
\expandafter\def\csname SRef@cdp:eq:C.IZ.4\endcsname{\hyperlink{tag-00OQ}{II.\ref*{cdp:eq:C.IZ.4}}}
\expandafter\def\csname SRef@cdp:eq:C.IZ.5\endcsname{\hyperlink{tag-00OR}{II.\ref*{cdp:eq:C.IZ.5}}}
\expandafter\def\csname SRef@cdp:eq:C.IZ.6\endcsname{\hyperlink{tag-00OS}{II.\ref*{cdp:eq:C.IZ.6}}}
\expandafter\def\csname SRef@cdp:eq:C.IZ.7\endcsname{\hyperlink{tag-00OT}{II.\ref*{cdp:eq:C.IZ.7}}}
\expandafter\def\csname SRef@cdp:eq:C.IZ.8\endcsname{\hyperlink{tag-00OU}{II.\ref*{cdp:eq:C.IZ.8}}}
\expandafter\def\csname SRef@cdp:eq:C.SC.4\endcsname{\hyperlink{tag-010O}{II.\ref*{cdp:eq:C.SC.4}}}
\expandafter\def\csname SRef@cdp:eq:C.INT.1\endcsname{\hyperlink{tag-00NQ}{II.\ref*{cdp:eq:C.INT.1}}}
\expandafter\def\csname SRef@cdp:eq:C.INT.2\endcsname{\hyperlink{tag-00O1}{II.\ref*{cdp:eq:C.INT.2}}}
\expandafter\def\csname SRef@cdp:eq:C.INT.3\endcsname{\hyperlink{tag-00OA}{II.\ref*{cdp:eq:C.INT.3}}}
\expandafter\def\csname SRef@cdp:eq:C.INT.4\endcsname{\hyperlink{tag-00OB}{II.\ref*{cdp:eq:C.INT.4}}}
\expandafter\def\csname SRef@cdp:eq:C.INT.5\endcsname{\hyperlink{tag-00OC}{II.\ref*{cdp:eq:C.INT.5}}}
\expandafter\def\csname SRef@cdp:eq:C.INT.6\endcsname{\hyperlink{tag-00OD}{II.\ref*{cdp:eq:C.INT.6}}}
\expandafter\def\csname SRef@cdp:eq:C.INT.7\endcsname{\hyperlink{tag-00OE}{II.\ref*{cdp:eq:C.INT.7}}}
\expandafter\def\csname SRef@cdp:eq:C.INT.8\endcsname{\hyperlink{tag-00OF}{II.\ref*{cdp:eq:C.INT.8}}}
\expandafter\def\csname SRef@cdp:eq:C.INT.9\endcsname{\hyperlink{tag-00OG}{II.\ref*{cdp:eq:C.INT.9}}}
\expandafter\def\csname SRef@cdp:eq:C.INT.10\endcsname{\hyperlink{tag-00NR}{II.\ref*{cdp:eq:C.INT.10}}}
\expandafter\def\csname SRef@cdp:eq:C.SC.5\endcsname{\hyperlink{tag-010P}{II.\ref*{cdp:eq:C.SC.5}}}
\expandafter\def\csname SRef@cdp:eq:C.INT.11\endcsname{\hyperlink{tag-00NS}{II.\ref*{cdp:eq:C.INT.11}}}
\expandafter\def\csname SRef@cdp:eq:C.INT.12\endcsname{\hyperlink{tag-00NT}{II.\ref*{cdp:eq:C.INT.12}}}
\expandafter\def\csname SRef@cdp:eq:C.INT.13\endcsname{\hyperlink{tag-00NU}{II.\ref*{cdp:eq:C.INT.13}}}
\expandafter\def\csname SRef@cdp:eq:C.INT.14\endcsname{\hyperlink{tag-00NV}{II.\ref*{cdp:eq:C.INT.14}}}
\expandafter\def\csname SRef@cdp:eq:C.INT.15\endcsname{\hyperlink{tag-00NW}{II.\ref*{cdp:eq:C.INT.15}}}
\expandafter\def\csname SRef@cdp:eq:C.INT.16\endcsname{\hyperlink{tag-00NX}{II.\ref*{cdp:eq:C.INT.16}}}
\expandafter\def\csname SRef@cdp:eq:C.INT.17\endcsname{\hyperlink{tag-00NY}{II.\ref*{cdp:eq:C.INT.17}}}
\expandafter\def\csname SRef@cdp:eq:C.INT.18\endcsname{\hyperlink{tag-00NZ}{II.\ref*{cdp:eq:C.INT.18}}}
\expandafter\def\csname SRef@cdp:eq:C.INT.19\endcsname{\hyperlink{tag-00O0}{II.\ref*{cdp:eq:C.INT.19}}}
\expandafter\def\csname SRef@cdp:eq:C.INT.20\endcsname{\hyperlink{tag-00O2}{II.\ref*{cdp:eq:C.INT.20}}}
\expandafter\def\csname SRef@cdp:eq:C.SC.6\endcsname{\hyperlink{tag-010Q}{II.\ref*{cdp:eq:C.SC.6}}}
\expandafter\def\csname SRef@cdp:eq:C.INT.21\endcsname{\hyperlink{tag-00O3}{II.\ref*{cdp:eq:C.INT.21}}}
\expandafter\def\csname SRef@cdp:eq:C.INT.22\endcsname{\hyperlink{tag-00O4}{II.\ref*{cdp:eq:C.INT.22}}}
\expandafter\def\csname SRef@cdp:eq:C.INT.23\endcsname{\hyperlink{tag-00O5}{II.\ref*{cdp:eq:C.INT.23}}}
\expandafter\def\csname SRef@cdp:eq:C.INT.24\endcsname{\hyperlink{tag-00O6}{II.\ref*{cdp:eq:C.INT.24}}}
\expandafter\def\csname SRef@cdp:eq:C.INT.25\endcsname{\hyperlink{tag-00O7}{II.\ref*{cdp:eq:C.INT.25}}}
\expandafter\def\csname SRef@cdp:eq:C.INT.26\endcsname{\hyperlink{tag-00O8}{II.\ref*{cdp:eq:C.INT.26}}}
\expandafter\def\csname SRef@cdp:eq:C.INT.27\endcsname{\hyperlink{tag-00O9}{II.\ref*{cdp:eq:C.INT.27}}}
\expandafter\def\csname SRef@cdp:eq:C.MH.1\endcsname{\hyperlink{tag-00QA}{II.\ref*{cdp:eq:C.MH.1}}}
\expandafter\def\csname SRef@cdp:eq:C.MH.2\endcsname{\hyperlink{tag-00QJ}{II.\ref*{cdp:eq:C.MH.2}}}
\expandafter\def\csname SRef@cdp:eq:C.MH.3\endcsname{\hyperlink{tag-00QK}{II.\ref*{cdp:eq:C.MH.3}}}
\expandafter\def\csname SRef@cdp:eq:C.CF.1\endcsname{\hyperlink{tag-00BN}{II.\ref*{cdp:eq:C.CF.1}}}
\expandafter\def\csname SRef@cdp:eq:C.MH.4\endcsname{\hyperlink{tag-00QL}{II.\ref*{cdp:eq:C.MH.4}}}
\expandafter\def\csname SRef@cdp:eq:C.MH.5\endcsname{\hyperlink{tag-00QM}{II.\ref*{cdp:eq:C.MH.5}}}
\expandafter\def\csname SRef@cdp:eq:C.MH.6\endcsname{\hyperlink{tag-00QN}{II.\ref*{cdp:eq:C.MH.6}}}
\expandafter\def\csname SRef@cdp:eq:C.MH.7\endcsname{\hyperlink{tag-00QO}{II.\ref*{cdp:eq:C.MH.7}}}
\expandafter\def\csname SRef@cdp:eq:C.MH.8\endcsname{\hyperlink{tag-00QP}{II.\ref*{cdp:eq:C.MH.8}}}
\expandafter\def\csname SRef@cdp:eq:C.MH.9\endcsname{\hyperlink{tag-00QQ}{II.\ref*{cdp:eq:C.MH.9}}}
\expandafter\def\csname SRef@cdp:eq:C.MH.10\endcsname{\hyperlink{tag-00QB}{II.\ref*{cdp:eq:C.MH.10}}}
\expandafter\def\csname SRef@cdp:eq:C.MH.11\endcsname{\hyperlink{tag-00QC}{II.\ref*{cdp:eq:C.MH.11}}}
\expandafter\def\csname SRef@cdp:eq:C.MH.12\endcsname{\hyperlink{tag-00QD}{II.\ref*{cdp:eq:C.MH.12}}}
\expandafter\def\csname SRef@cdp:eq:C.MH.13\endcsname{\hyperlink{tag-00QE}{II.\ref*{cdp:eq:C.MH.13}}}
\expandafter\def\csname SRef@cdp:eq:C.CF.2\endcsname{\hyperlink{tag-00BO}{II.\ref*{cdp:eq:C.CF.2}}}
\expandafter\def\csname SRef@cdp:eq:C.MH.14\endcsname{\hyperlink{tag-00QF}{II.\ref*{cdp:eq:C.MH.14}}}
\expandafter\def\csname SRef@cdp:eq:C.MH.15\endcsname{\hyperlink{tag-00QG}{II.\ref*{cdp:eq:C.MH.15}}}
\expandafter\def\csname SRef@cdp:eq:C.MH.16\endcsname{\hyperlink{tag-00QH}{II.\ref*{cdp:eq:C.MH.16}}}
\expandafter\def\csname SRef@cdp:eq:C.MH.17\endcsname{\hyperlink{tag-00QI}{II.\ref*{cdp:eq:C.MH.17}}}
\expandafter\def\csname SRef@cdp:eq:C.1\endcsname{\hyperlink{tag-008N}{II.\ref*{cdp:eq:C.1}}}
\expandafter\def\csname SRef@cdp:eq:C.2\endcsname{\hyperlink{tag-008Y}{II.\ref*{cdp:eq:C.2}}}
\expandafter\def\csname SRef@cdp:eq:C.MS.1\endcsname{\hyperlink{tag-00RH}{II.\ref*{cdp:eq:C.MS.1}}}
\expandafter\def\csname SRef@cdp:eq:C.MS.2\endcsname{\hyperlink{tag-00RJ}{II.\ref*{cdp:eq:C.MS.2}}}
\expandafter\def\csname SRef@cdp:eq:C.MS.3\endcsname{\hyperlink{tag-00RK}{II.\ref*{cdp:eq:C.MS.3}}}
\expandafter\def\csname SRef@cdp:eq:C.MS.10\endcsname{\hyperlink{tag-00RI}{II.\ref*{cdp:eq:C.MS.10}}}
\expandafter\def\csname SRef@cdp:eq:C.CF.3\endcsname{\hyperlink{tag-00BP}{II.\ref*{cdp:eq:C.CF.3}}}
\expandafter\def\csname SRef@cdp:eq:C.MS.4\endcsname{\hyperlink{tag-00RL}{II.\ref*{cdp:eq:C.MS.4}}}
\expandafter\def\csname SRef@cdp:eq:C.MS.5\endcsname{\hyperlink{tag-00RM}{II.\ref*{cdp:eq:C.MS.5}}}
\expandafter\def\csname SRef@cdp:eq:C.MS.6\endcsname{\hyperlink{tag-00RN}{II.\ref*{cdp:eq:C.MS.6}}}
\expandafter\def\csname SRef@cdp:eq:C.MS.7\endcsname{\hyperlink{tag-00RO}{II.\ref*{cdp:eq:C.MS.7}}}
\expandafter\def\csname SRef@cdp:eq:C.MS.8\endcsname{\hyperlink{tag-00RP}{II.\ref*{cdp:eq:C.MS.8}}}
\expandafter\def\csname SRef@cdp:eq:C.MS.9\endcsname{\hyperlink{tag-00RQ}{II.\ref*{cdp:eq:C.MS.9}}}
\expandafter\def\csname SRef@cdp:eq:C.K3.1\endcsname{\hyperlink{tag-00OY}{II.\ref*{cdp:eq:C.K3.1}}}
\expandafter\def\csname SRef@cdp:eq:C.K3.2\endcsname{\hyperlink{tag-00P0}{II.\ref*{cdp:eq:C.K3.2}}}
\expandafter\def\csname SRef@cdp:eq:C.K3.3\endcsname{\hyperlink{tag-00P1}{II.\ref*{cdp:eq:C.K3.3}}}
\expandafter\def\csname SRef@cdp:eq:C.K3.4\endcsname{\hyperlink{tag-00P2}{II.\ref*{cdp:eq:C.K3.4}}}
\expandafter\def\csname SRef@cdp:eq:C.CF.4\endcsname{\hyperlink{tag-00BQ}{II.\ref*{cdp:eq:C.CF.4}}}
\expandafter\def\csname SRef@cdp:eq:C.K3.5\endcsname{\hyperlink{tag-00P3}{II.\ref*{cdp:eq:C.K3.5}}}
\expandafter\def\csname SRef@cdp:eq:C.K3.6\endcsname{\hyperlink{tag-00P4}{II.\ref*{cdp:eq:C.K3.6}}}
\expandafter\def\csname SRef@cdp:eq:C.K3.7\endcsname{\hyperlink{tag-00P5}{II.\ref*{cdp:eq:C.K3.7}}}
\expandafter\def\csname SRef@cdp:eq:C.K3.8\endcsname{\hyperlink{tag-00P6}{II.\ref*{cdp:eq:C.K3.8}}}
\expandafter\def\csname SRef@cdp:eq:C.K3.9\endcsname{\hyperlink{tag-00P7}{II.\ref*{cdp:eq:C.K3.9}}}
\expandafter\def\csname SRef@cdp:eq:C.K3.10\endcsname{\hyperlink{tag-00OZ}{II.\ref*{cdp:eq:C.K3.10}}}
\expandafter\def\csname SRef@cdp:eq:C.KX.1\endcsname{\hyperlink{tag-00PW}{II.\ref*{cdp:eq:C.KX.1}}}
\expandafter\def\csname SRef@cdp:eq:C.KX.2\endcsname{\hyperlink{tag-00Q2}{II.\ref*{cdp:eq:C.KX.2}}}
\expandafter\def\csname SRef@cdp:eq:C.KX.3\endcsname{\hyperlink{tag-00Q3}{II.\ref*{cdp:eq:C.KX.3}}}
\expandafter\def\csname SRef@cdp:eq:C.KX.4\endcsname{\hyperlink{tag-00Q4}{II.\ref*{cdp:eq:C.KX.4}}}
\expandafter\def\csname SRef@cdp:eq:C.CF.5\endcsname{\hyperlink{tag-00BR}{II.\ref*{cdp:eq:C.CF.5}}}
\expandafter\def\csname SRef@cdp:eq:C.KX.5\endcsname{\hyperlink{tag-00Q5}{II.\ref*{cdp:eq:C.KX.5}}}
\expandafter\def\csname SRef@cdp:eq:C.KX.6\endcsname{\hyperlink{tag-00Q6}{II.\ref*{cdp:eq:C.KX.6}}}
\expandafter\def\csname SRef@cdp:eq:C.KX.7\endcsname{\hyperlink{tag-00Q7}{II.\ref*{cdp:eq:C.KX.7}}}
\expandafter\def\csname SRef@cdp:eq:C.KX.8\endcsname{\hyperlink{tag-00Q8}{II.\ref*{cdp:eq:C.KX.8}}}
\expandafter\def\csname SRef@cdp:eq:C.KX.9\endcsname{\hyperlink{tag-00Q9}{II.\ref*{cdp:eq:C.KX.9}}}
\expandafter\def\csname SRef@cdp:eq:C.KX.10\endcsname{\hyperlink{tag-00PX}{II.\ref*{cdp:eq:C.KX.10}}}
\expandafter\def\csname SRef@cdp:eq:C.KX.11\endcsname{\hyperlink{tag-00PY}{II.\ref*{cdp:eq:C.KX.11}}}
\expandafter\def\csname SRef@cdp:eq:C.KX.12\endcsname{\hyperlink{tag-00PZ}{II.\ref*{cdp:eq:C.KX.12}}}
\expandafter\def\csname SRef@cdp:eq:C.KX.13\endcsname{\hyperlink{tag-00Q0}{II.\ref*{cdp:eq:C.KX.13}}}
\expandafter\def\csname SRef@cdp:eq:C.KX.14\endcsname{\hyperlink{tag-00Q1}{II.\ref*{cdp:eq:C.KX.14}}}
\expandafter\def\csname SRef@cdp:eq:C.CF.6\endcsname{\hyperlink{tag-00BS}{II.\ref*{cdp:eq:C.CF.6}}}
\expandafter\def\csname SRef@cdp:eq:C.EL.1\endcsname{\hyperlink{tag-00FY}{II.\ref*{cdp:eq:C.EL.1}}}
\expandafter\def\csname SRef@cdp:eq:C.EL.2\endcsname{\hyperlink{tag-00G1}{II.\ref*{cdp:eq:C.EL.2}}}
\expandafter\def\csname SRef@cdp:eq:C.EL.3\endcsname{\hyperlink{tag-00G2}{II.\ref*{cdp:eq:C.EL.3}}}
\expandafter\def\csname SRef@cdp:eq:C.EL.4\endcsname{\hyperlink{tag-00G3}{II.\ref*{cdp:eq:C.EL.4}}}
\expandafter\def\csname SRef@cdp:eq:C.EL.5\endcsname{\hyperlink{tag-00G4}{II.\ref*{cdp:eq:C.EL.5}}}
\expandafter\def\csname SRef@cdp:eq:C.EL.6\endcsname{\hyperlink{tag-00G5}{II.\ref*{cdp:eq:C.EL.6}}}
\expandafter\def\csname SRef@cdp:eq:C.EL.7\endcsname{\hyperlink{tag-00G6}{II.\ref*{cdp:eq:C.EL.7}}}
\expandafter\def\csname SRef@cdp:eq:C.EL.8\endcsname{\hyperlink{tag-00G7}{II.\ref*{cdp:eq:C.EL.8}}}
\expandafter\def\csname SRef@cdp:eq:C.EL.9\endcsname{\hyperlink{tag-00G8}{II.\ref*{cdp:eq:C.EL.9}}}
\expandafter\def\csname SRef@cdp:eq:C.EL.10\endcsname{\hyperlink{tag-00FZ}{II.\ref*{cdp:eq:C.EL.10}}}
\expandafter\def\csname SRef@cdp:eq:C.CF.7\endcsname{\hyperlink{tag-00BT}{II.\ref*{cdp:eq:C.CF.7}}}
\expandafter\def\csname SRef@cdp:eq:C.EL.11\endcsname{\hyperlink{tag-00G0}{II.\ref*{cdp:eq:C.EL.11}}}
\expandafter\def\csname SRef@cdp:eq:C.CZ.1\endcsname{\hyperlink{tag-00DN}{II.\ref*{cdp:eq:C.CZ.1}}}
\expandafter\def\csname SRef@cdp:eq:C.CZ.2\endcsname{\hyperlink{tag-00DT}{II.\ref*{cdp:eq:C.CZ.2}}}
\expandafter\def\csname SRef@cdp:eq:C.CZ.3\endcsname{\hyperlink{tag-00DU}{II.\ref*{cdp:eq:C.CZ.3}}}
\expandafter\def\csname SRef@cdp:eq:C.CZ.4\endcsname{\hyperlink{tag-00DV}{II.\ref*{cdp:eq:C.CZ.4}}}
\expandafter\def\csname SRef@cdp:eq:C.CZ.5\endcsname{\hyperlink{tag-00DW}{II.\ref*{cdp:eq:C.CZ.5}}}
\expandafter\def\csname SRef@cdp:eq:C.CZ.6\endcsname{\hyperlink{tag-00DX}{II.\ref*{cdp:eq:C.CZ.6}}}
\expandafter\def\csname SRef@cdp:eq:C.CZ.7\endcsname{\hyperlink{tag-00DY}{II.\ref*{cdp:eq:C.CZ.7}}}
\expandafter\def\csname SRef@cdp:eq:C.CZ.8\endcsname{\hyperlink{tag-00DZ}{II.\ref*{cdp:eq:C.CZ.8}}}
\expandafter\def\csname SRef@cdp:eq:C.CZ.9\endcsname{\hyperlink{tag-00E0}{II.\ref*{cdp:eq:C.CZ.9}}}
\expandafter\def\csname SRef@cdp:eq:C.CF.8\endcsname{\hyperlink{tag-00BU}{II.\ref*{cdp:eq:C.CF.8}}}
\expandafter\def\csname SRef@cdp:eq:C.CZ.10\endcsname{\hyperlink{tag-00DO}{II.\ref*{cdp:eq:C.CZ.10}}}
\expandafter\def\csname SRef@cdp:eq:C.CZ.11\endcsname{\hyperlink{tag-00DP}{II.\ref*{cdp:eq:C.CZ.11}}}
\expandafter\def\csname SRef@cdp:eq:C.CZ.12\endcsname{\hyperlink{tag-00DQ}{II.\ref*{cdp:eq:C.CZ.12}}}
\expandafter\def\csname SRef@cdp:eq:C.CZ.13\endcsname{\hyperlink{tag-00DR}{II.\ref*{cdp:eq:C.CZ.13}}}
\expandafter\def\csname SRef@cdp:eq:C.CZ.14\endcsname{\hyperlink{tag-00DS}{II.\ref*{cdp:eq:C.CZ.14}}}
\expandafter\def\csname SRef@cdp:eq:C.MR.1\endcsname{\hyperlink{tag-00R1}{II.\ref*{cdp:eq:C.MR.1}}}
\expandafter\def\csname SRef@cdp:eq:C.MR.2\endcsname{\hyperlink{tag-00R9}{II.\ref*{cdp:eq:C.MR.2}}}
\expandafter\def\csname SRef@cdp:eq:C.MR.3\endcsname{\hyperlink{tag-00RA}{II.\ref*{cdp:eq:C.MR.3}}}
\expandafter\def\csname SRef@cdp:eq:C.MR.4\endcsname{\hyperlink{tag-00RB}{II.\ref*{cdp:eq:C.MR.4}}}
\expandafter\def\csname SRef@cdp:eq:C.MR.5\endcsname{\hyperlink{tag-00RC}{II.\ref*{cdp:eq:C.MR.5}}}
\expandafter\def\csname SRef@cdp:eq:C.CF.9\endcsname{\hyperlink{tag-00BV}{II.\ref*{cdp:eq:C.CF.9}}}
\expandafter\def\csname SRef@cdp:eq:C.MR.6\endcsname{\hyperlink{tag-00RD}{II.\ref*{cdp:eq:C.MR.6}}}
\expandafter\def\csname SRef@cdp:eq:C.MR.7\endcsname{\hyperlink{tag-00RE}{II.\ref*{cdp:eq:C.MR.7}}}
\expandafter\def\csname SRef@cdp:eq:C.MR.8\endcsname{\hyperlink{tag-00RF}{II.\ref*{cdp:eq:C.MR.8}}}
\expandafter\def\csname SRef@cdp:eq:C.MR.9\endcsname{\hyperlink{tag-00RG}{II.\ref*{cdp:eq:C.MR.9}}}
\expandafter\def\csname SRef@cdp:eq:C.MR.10\endcsname{\hyperlink{tag-00R2}{II.\ref*{cdp:eq:C.MR.10}}}
\expandafter\def\csname SRef@cdp:eq:C.MR.11\endcsname{\hyperlink{tag-00R3}{II.\ref*{cdp:eq:C.MR.11}}}
\expandafter\def\csname SRef@cdp:eq:C.MR.12\endcsname{\hyperlink{tag-00R4}{II.\ref*{cdp:eq:C.MR.12}}}
\expandafter\def\csname SRef@cdp:eq:C.MR.13\endcsname{\hyperlink{tag-00R5}{II.\ref*{cdp:eq:C.MR.13}}}
\expandafter\def\csname SRef@cdp:eq:C.MR.14\endcsname{\hyperlink{tag-00R6}{II.\ref*{cdp:eq:C.MR.14}}}
\expandafter\def\csname SRef@cdp:eq:C.MR.15\endcsname{\hyperlink{tag-00R7}{II.\ref*{cdp:eq:C.MR.15}}}
\expandafter\def\csname SRef@cdp:eq:C.CA.1\endcsname{\hyperlink{tag-00BC}{II.\ref*{cdp:eq:C.CA.1}}}
\expandafter\def\csname SRef@cdp:eq:C.MR.16\endcsname{\hyperlink{tag-00R8}{II.\ref*{cdp:eq:C.MR.16}}}
\expandafter\def\csname SRef@cdp:eq:C.RH.1\endcsname{\hyperlink{tag-00ZG}{II.\ref*{cdp:eq:C.RH.1}}}
\expandafter\def\csname SRef@cdp:eq:C.RH.2\endcsname{\hyperlink{tag-00ZI}{II.\ref*{cdp:eq:C.RH.2}}}
\expandafter\def\csname SRef@cdp:eq:C.RH.3\endcsname{\hyperlink{tag-00ZJ}{II.\ref*{cdp:eq:C.RH.3}}}
\expandafter\def\csname SRef@cdp:eq:C.RH.4\endcsname{\hyperlink{tag-00ZK}{II.\ref*{cdp:eq:C.RH.4}}}
\expandafter\def\csname SRef@cdp:eq:C.RH.5\endcsname{\hyperlink{tag-00ZL}{II.\ref*{cdp:eq:C.RH.5}}}
\expandafter\def\csname SRef@cdp:eq:C.RH.6\endcsname{\hyperlink{tag-00ZM}{II.\ref*{cdp:eq:C.RH.6}}}
\expandafter\def\csname SRef@cdp:eq:C.RH.7\endcsname{\hyperlink{tag-00ZN}{II.\ref*{cdp:eq:C.RH.7}}}
\expandafter\def\csname SRef@cdp:eq:C.RH.8\endcsname{\hyperlink{tag-00ZO}{II.\ref*{cdp:eq:C.RH.8}}}
\expandafter\def\csname SRef@cdp:eq:C.RH.9\endcsname{\hyperlink{tag-00ZP}{II.\ref*{cdp:eq:C.RH.9}}}
\expandafter\def\csname SRef@cdp:eq:C.CA.2\endcsname{\hyperlink{tag-00BF}{II.\ref*{cdp:eq:C.CA.2}}}
\expandafter\def\csname SRef@cdp:eq:C.RH.10\endcsname{\hyperlink{tag-00ZH}{II.\ref*{cdp:eq:C.RH.10}}}
\expandafter\def\csname SRef@cdp:eq:C.TC.1\endcsname{\hyperlink{tag-012C}{II.\ref*{cdp:eq:C.TC.1}}}
\expandafter\def\csname SRef@cdp:eq:C.TC.2\endcsname{\hyperlink{tag-012G}{II.\ref*{cdp:eq:C.TC.2}}}
\expandafter\def\csname SRef@cdp:eq:C.TC.3\endcsname{\hyperlink{tag-012H}{II.\ref*{cdp:eq:C.TC.3}}}
\expandafter\def\csname SRef@cdp:eq:C.TC.4\endcsname{\hyperlink{tag-012I}{II.\ref*{cdp:eq:C.TC.4}}}
\expandafter\def\csname SRef@cdp:eq:C.TC.5\endcsname{\hyperlink{tag-012J}{II.\ref*{cdp:eq:C.TC.5}}}
\expandafter\def\csname SRef@cdp:eq:C.TC.6\endcsname{\hyperlink{tag-012K}{II.\ref*{cdp:eq:C.TC.6}}}
\expandafter\def\csname SRef@cdp:eq:C.TC.7\endcsname{\hyperlink{tag-012L}{II.\ref*{cdp:eq:C.TC.7}}}
\expandafter\def\csname SRef@cdp:eq:C.TC.8\endcsname{\hyperlink{tag-012M}{II.\ref*{cdp:eq:C.TC.8}}}
\expandafter\def\csname SRef@cdp:eq:C.TC.9\endcsname{\hyperlink{tag-012N}{II.\ref*{cdp:eq:C.TC.9}}}
\expandafter\def\csname SRef@cdp:eq:C.CA.3\endcsname{\hyperlink{tag-00BG}{II.\ref*{cdp:eq:C.CA.3}}}
\expandafter\def\csname SRef@cdp:eq:C.TC.10\endcsname{\hyperlink{tag-012D}{II.\ref*{cdp:eq:C.TC.10}}}
\expandafter\def\csname SRef@cdp:eq:C.TC.11\endcsname{\hyperlink{tag-012E}{II.\ref*{cdp:eq:C.TC.11}}}
\expandafter\def\csname SRef@cdp:eq:C.TC.12\endcsname{\hyperlink{tag-012F}{II.\ref*{cdp:eq:C.TC.12}}}
\expandafter\def\csname SRef@cdp:eq:C.ER.1\endcsname{\hyperlink{tag-00IN}{II.\ref*{cdp:eq:C.ER.1}}}
\expandafter\def\csname SRef@cdp:eq:C.ER.2\endcsname{\hyperlink{tag-00IX}{II.\ref*{cdp:eq:C.ER.2}}}
\expandafter\def\csname SRef@cdp:eq:C.ER.3\endcsname{\hyperlink{tag-00IY}{II.\ref*{cdp:eq:C.ER.3}}}
\expandafter\def\csname SRef@cdp:eq:C.ER.4\endcsname{\hyperlink{tag-00IZ}{II.\ref*{cdp:eq:C.ER.4}}}
\expandafter\def\csname SRef@cdp:eq:C.ER.5\endcsname{\hyperlink{tag-00J0}{II.\ref*{cdp:eq:C.ER.5}}}
\expandafter\def\csname SRef@cdp:eq:C.ER.6\endcsname{\hyperlink{tag-00J1}{II.\ref*{cdp:eq:C.ER.6}}}
\expandafter\def\csname SRef@cdp:eq:C.ER.7\endcsname{\hyperlink{tag-00J2}{II.\ref*{cdp:eq:C.ER.7}}}
\expandafter\def\csname SRef@cdp:eq:C.CA.4\endcsname{\hyperlink{tag-00BH}{II.\ref*{cdp:eq:C.CA.4}}}
\expandafter\def\csname SRef@cdp:eq:C.ER.8\endcsname{\hyperlink{tag-00J3}{II.\ref*{cdp:eq:C.ER.8}}}
\expandafter\def\csname SRef@cdp:eq:C.ER.9\endcsname{\hyperlink{tag-00J4}{II.\ref*{cdp:eq:C.ER.9}}}
\expandafter\def\csname SRef@cdp:eq:C.ER.10\endcsname{\hyperlink{tag-00IO}{II.\ref*{cdp:eq:C.ER.10}}}
\expandafter\def\csname SRef@cdp:eq:C.ER.11\endcsname{\hyperlink{tag-00IP}{II.\ref*{cdp:eq:C.ER.11}}}
\expandafter\def\csname SRef@cdp:eq:C.ER.12\endcsname{\hyperlink{tag-00IQ}{II.\ref*{cdp:eq:C.ER.12}}}
\expandafter\def\csname SRef@cdp:eq:C.ER.13\endcsname{\hyperlink{tag-00IR}{II.\ref*{cdp:eq:C.ER.13}}}
\expandafter\def\csname SRef@cdp:eq:C.ER.14\endcsname{\hyperlink{tag-00IS}{II.\ref*{cdp:eq:C.ER.14}}}
\expandafter\def\csname SRef@cdp:eq:C.ER.15\endcsname{\hyperlink{tag-00IT}{II.\ref*{cdp:eq:C.ER.15}}}
\expandafter\def\csname SRef@cdp:eq:C.ER.16\endcsname{\hyperlink{tag-00IU}{II.\ref*{cdp:eq:C.ER.16}}}
\expandafter\def\csname SRef@cdp:eq:C.CA.5\endcsname{\hyperlink{tag-00BI}{II.\ref*{cdp:eq:C.CA.5}}}
\expandafter\def\csname SRef@cdp:eq:C.ER.17\endcsname{\hyperlink{tag-00IV}{II.\ref*{cdp:eq:C.ER.17}}}
\expandafter\def\csname SRef@cdp:eq:C.ER.18\endcsname{\hyperlink{tag-00IW}{II.\ref*{cdp:eq:C.ER.18}}}
\expandafter\def\csname SRef@cdp:eq:C.CP.1\endcsname{\hyperlink{tag-00C9}{II.\ref*{cdp:eq:C.CP.1}}}
\expandafter\def\csname SRef@cdp:eq:C.CP.2\endcsname{\hyperlink{tag-00CB}{II.\ref*{cdp:eq:C.CP.2}}}
\expandafter\def\csname SRef@cdp:eq:C.CP.3\endcsname{\hyperlink{tag-00CC}{II.\ref*{cdp:eq:C.CP.3}}}
\expandafter\def\csname SRef@cdp:eq:C.CP.4\endcsname{\hyperlink{tag-00CD}{II.\ref*{cdp:eq:C.CP.4}}}
\expandafter\def\csname SRef@cdp:eq:C.CP.5\endcsname{\hyperlink{tag-00CE}{II.\ref*{cdp:eq:C.CP.5}}}
\expandafter\def\csname SRef@cdp:eq:C.CP.6\endcsname{\hyperlink{tag-00CF}{II.\ref*{cdp:eq:C.CP.6}}}
\expandafter\def\csname SRef@cdp:eq:C.CP.7\endcsname{\hyperlink{tag-00CG}{II.\ref*{cdp:eq:C.CP.7}}}
\expandafter\def\csname SRef@cdp:eq:C.CA.6\endcsname{\hyperlink{tag-00BJ}{II.\ref*{cdp:eq:C.CA.6}}}
\expandafter\def\csname SRef@cdp:eq:C.CP.8\endcsname{\hyperlink{tag-00CH}{II.\ref*{cdp:eq:C.CP.8}}}
\expandafter\def\csname SRef@cdp:eq:C.CP.9\endcsname{\hyperlink{tag-00CI}{II.\ref*{cdp:eq:C.CP.9}}}
\expandafter\def\csname SRef@cdp:eq:C.CP.10\endcsname{\hyperlink{tag-00CA}{II.\ref*{cdp:eq:C.CP.10}}}
\expandafter\def\csname SRef@cdp:eq:C.KM.1\endcsname{\hyperlink{tag-00PI}{II.\ref*{cdp:eq:C.KM.1}}}
\expandafter\def\csname SRef@cdp:eq:C.KM.2\endcsname{\hyperlink{tag-00PL}{II.\ref*{cdp:eq:C.KM.2}}}
\expandafter\def\csname SRef@cdp:eq:C.KM.S1\endcsname{\hyperlink{tag-00PT}{II.\ref*{cdp:eq:C.KM.S1}}}
\expandafter\def\csname SRef@cdp:eq:C.KM.3\endcsname{\hyperlink{tag-00PM}{II.\ref*{cdp:eq:C.KM.3}}}
\expandafter\def\csname SRef@cdp:eq:C.KM.4\endcsname{\hyperlink{tag-00PN}{II.\ref*{cdp:eq:C.KM.4}}}
\expandafter\def\csname SRef@cdp:eq:C.KM.5\endcsname{\hyperlink{tag-00PO}{II.\ref*{cdp:eq:C.KM.5}}}
\expandafter\def\csname SRef@cdp:eq:C.KM.6\endcsname{\hyperlink{tag-00PP}{II.\ref*{cdp:eq:C.KM.6}}}
\expandafter\def\csname SRef@cdp:eq:C.CA.7\endcsname{\hyperlink{tag-00BK}{II.\ref*{cdp:eq:C.CA.7}}}
\expandafter\def\csname SRef@cdp:eq:C.KM.7\endcsname{\hyperlink{tag-00PQ}{II.\ref*{cdp:eq:C.KM.7}}}
\expandafter\def\csname SRef@cdp:eq:C.KM.8\endcsname{\hyperlink{tag-00PR}{II.\ref*{cdp:eq:C.KM.8}}}
\expandafter\def\csname SRef@cdp:eq:C.KM.9\endcsname{\hyperlink{tag-00PS}{II.\ref*{cdp:eq:C.KM.9}}}
\expandafter\def\csname SRef@cdp:eq:C.KM.10\endcsname{\hyperlink{tag-00PJ}{II.\ref*{cdp:eq:C.KM.10}}}
\expandafter\def\csname SRef@cdp:eq:C.KM.S2\endcsname{\hyperlink{tag-00PU}{II.\ref*{cdp:eq:C.KM.S2}}}
\expandafter\def\csname SRef@cdp:eq:C.KM.S3\endcsname{\hyperlink{tag-00PV}{II.\ref*{cdp:eq:C.KM.S3}}}
\expandafter\def\csname SRef@cdp:eq:C.KM.11\endcsname{\hyperlink{tag-00PK}{II.\ref*{cdp:eq:C.KM.11}}}
\expandafter\def\csname SRef@cdp:eq:C.PE.1\endcsname{\hyperlink{tag-00U2}{II.\ref*{cdp:eq:C.PE.1}}}
\expandafter\def\csname SRef@cdp:eq:C.PE.2\endcsname{\hyperlink{tag-00UD}{II.\ref*{cdp:eq:C.PE.2}}}
\expandafter\def\csname SRef@cdp:eq:C.PE.3\endcsname{\hyperlink{tag-00UG}{II.\ref*{cdp:eq:C.PE.3}}}
\expandafter\def\csname SRef@cdp:eq:C.CA.8\endcsname{\hyperlink{tag-00BL}{II.\ref*{cdp:eq:C.CA.8}}}
\expandafter\def\csname SRef@cdp:eq:C.PE.4\endcsname{\hyperlink{tag-00UH}{II.\ref*{cdp:eq:C.PE.4}}}
\expandafter\def\csname SRef@cdp:eq:C.PE.5\endcsname{\hyperlink{tag-00UI}{II.\ref*{cdp:eq:C.PE.5}}}
\expandafter\def\csname SRef@cdp:eq:C.PE.6\endcsname{\hyperlink{tag-00UJ}{II.\ref*{cdp:eq:C.PE.6}}}
\expandafter\def\csname SRef@cdp:eq:C.PE.7\endcsname{\hyperlink{tag-00UK}{II.\ref*{cdp:eq:C.PE.7}}}
\expandafter\def\csname SRef@cdp:eq:C.PE.8\endcsname{\hyperlink{tag-00UL}{II.\ref*{cdp:eq:C.PE.8}}}
\expandafter\def\csname SRef@cdp:eq:C.PE.9\endcsname{\hyperlink{tag-00UM}{II.\ref*{cdp:eq:C.PE.9}}}
\expandafter\def\csname SRef@cdp:eq:C.PE.10\endcsname{\hyperlink{tag-00U3}{II.\ref*{cdp:eq:C.PE.10}}}
\expandafter\def\csname SRef@cdp:eq:C.PE.11\endcsname{\hyperlink{tag-00U4}{II.\ref*{cdp:eq:C.PE.11}}}
\expandafter\def\csname SRef@cdp:eq:C.PE.12\endcsname{\hyperlink{tag-00U5}{II.\ref*{cdp:eq:C.PE.12}}}
\expandafter\def\csname SRef@cdp:eq:C.PE.13\endcsname{\hyperlink{tag-00U6}{II.\ref*{cdp:eq:C.PE.13}}}
\expandafter\def\csname SRef@cdp:eq:C.CA.9\endcsname{\hyperlink{tag-00BM}{II.\ref*{cdp:eq:C.CA.9}}}
\expandafter\def\csname SRef@cdp:eq:C.PE.14\endcsname{\hyperlink{tag-00U7}{II.\ref*{cdp:eq:C.PE.14}}}
\expandafter\def\csname SRef@cdp:eq:C.PE.15\endcsname{\hyperlink{tag-00U8}{II.\ref*{cdp:eq:C.PE.15}}}
\expandafter\def\csname SRef@cdp:eq:C.PE.16\endcsname{\hyperlink{tag-00U9}{II.\ref*{cdp:eq:C.PE.16}}}
\expandafter\def\csname SRef@cdp:eq:C.PE.17\endcsname{\hyperlink{tag-00UA}{II.\ref*{cdp:eq:C.PE.17}}}
\expandafter\def\csname SRef@cdp:eq:C.PE.18\endcsname{\hyperlink{tag-00UB}{II.\ref*{cdp:eq:C.PE.18}}}
\expandafter\def\csname SRef@cdp:eq:C.PE.19\endcsname{\hyperlink{tag-00UC}{II.\ref*{cdp:eq:C.PE.19}}}
\expandafter\def\csname SRef@cdp:eq:C.PE.20\endcsname{\hyperlink{tag-00UE}{II.\ref*{cdp:eq:C.PE.20}}}
\expandafter\def\csname SRef@cdp:eq:C.PE.21\endcsname{\hyperlink{tag-00UF}{II.\ref*{cdp:eq:C.PE.21}}}
\expandafter\def\csname SRef@cdp:eq:C.SKD.1\endcsname{\hyperlink{tag-011N}{II.\ref*{cdp:eq:C.SKD.1}}}
\expandafter\def\csname SRef@cdp:eq:C.SKD.2\endcsname{\hyperlink{tag-011U}{II.\ref*{cdp:eq:C.SKD.2}}}
\expandafter\def\csname SRef@cdp:eq:C.CA.10\endcsname{\hyperlink{tag-00BD}{II.\ref*{cdp:eq:C.CA.10}}}
\expandafter\def\csname SRef@cdp:eq:C.SKD.3\endcsname{\hyperlink{tag-011V}{II.\ref*{cdp:eq:C.SKD.3}}}
\expandafter\def\csname SRef@cdp:eq:C.SKD.4\endcsname{\hyperlink{tag-011W}{II.\ref*{cdp:eq:C.SKD.4}}}
\expandafter\def\csname SRef@cdp:eq:C.SKD.5\endcsname{\hyperlink{tag-011X}{II.\ref*{cdp:eq:C.SKD.5}}}
\expandafter\def\csname SRef@cdp:eq:C.SKD.6\endcsname{\hyperlink{tag-011Y}{II.\ref*{cdp:eq:C.SKD.6}}}
\expandafter\def\csname SRef@cdp:eq:C.SKD.7\endcsname{\hyperlink{tag-011Z}{II.\ref*{cdp:eq:C.SKD.7}}}
\expandafter\def\csname SRef@cdp:eq:C.SKD.8\endcsname{\hyperlink{tag-0120}{II.\ref*{cdp:eq:C.SKD.8}}}
\expandafter\def\csname SRef@cdp:eq:C.SKD.9\endcsname{\hyperlink{tag-0121}{II.\ref*{cdp:eq:C.SKD.9}}}
\expandafter\def\csname SRef@cdp:eq:C.SKD.10\endcsname{\hyperlink{tag-011O}{II.\ref*{cdp:eq:C.SKD.10}}}
\expandafter\def\csname SRef@cdp:eq:C.SKD.11\endcsname{\hyperlink{tag-011P}{II.\ref*{cdp:eq:C.SKD.11}}}
\expandafter\def\csname SRef@cdp:eq:C.SKD.14\endcsname{\hyperlink{tag-011S}{II.\ref*{cdp:eq:C.SKD.14}}}
\expandafter\def\csname SRef@cdp:eq:C.CA.11\endcsname{\hyperlink{tag-00BE}{II.\ref*{cdp:eq:C.CA.11}}}
\expandafter\def\csname SRef@cdp:eq:C.SKD.15\endcsname{\hyperlink{tag-011T}{II.\ref*{cdp:eq:C.SKD.15}}}
\expandafter\def\csname SRef@cdp:eq:C.SKD.12\endcsname{\hyperlink{tag-011Q}{II.\ref*{cdp:eq:C.SKD.12}}}
\expandafter\def\csname SRef@cdp:eq:C.SKD.13\endcsname{\hyperlink{tag-011R}{II.\ref*{cdp:eq:C.SKD.13}}}
\expandafter\def\csname SRef@cdp:eq:C.EN.1\endcsname{\hyperlink{tag-00G9}{II.\ref*{cdp:eq:C.EN.1}}}
\expandafter\def\csname SRef@cdp:eq:C.EN.18\endcsname{\hyperlink{tag-00GI}{II.\ref*{cdp:eq:C.EN.18}}}
\expandafter\def\csname SRef@cdp:eq:C.EN.2\endcsname{\hyperlink{tag-00GJ}{II.\ref*{cdp:eq:C.EN.2}}}
\expandafter\def\csname SRef@cdp:eq:C.EN.3\endcsname{\hyperlink{tag-00GK}{II.\ref*{cdp:eq:C.EN.3}}}
\expandafter\def\csname SRef@cdp:eq:C.EN.4\endcsname{\hyperlink{tag-00GL}{II.\ref*{cdp:eq:C.EN.4}}}
\expandafter\def\csname SRef@cdp:eq:C.EN.5\endcsname{\hyperlink{tag-00GM}{II.\ref*{cdp:eq:C.EN.5}}}
\expandafter\def\csname SRef@cdp:eq:C.EN.6\endcsname{\hyperlink{tag-00GN}{II.\ref*{cdp:eq:C.EN.6}}}
\expandafter\def\csname SRef@cdp:eq:C.HQ.1\endcsname{\hyperlink{tag-00MQ}{II.\ref*{cdp:eq:C.HQ.1}}}
\expandafter\def\csname SRef@cdp:eq:C.EN.7\endcsname{\hyperlink{tag-00GO}{II.\ref*{cdp:eq:C.EN.7}}}
\expandafter\def\csname SRef@cdp:eq:C.EN.8\endcsname{\hyperlink{tag-00GP}{II.\ref*{cdp:eq:C.EN.8}}}
\expandafter\def\csname SRef@cdp:eq:C.EN.9\endcsname{\hyperlink{tag-00GQ}{II.\ref*{cdp:eq:C.EN.9}}}
\expandafter\def\csname SRef@cdp:eq:C.EN.10\endcsname{\hyperlink{tag-00GA}{II.\ref*{cdp:eq:C.EN.10}}}
\expandafter\def\csname SRef@cdp:eq:C.EN.11\endcsname{\hyperlink{tag-00GB}{II.\ref*{cdp:eq:C.EN.11}}}
\expandafter\def\csname SRef@cdp:eq:C.EN.12\endcsname{\hyperlink{tag-00GC}{II.\ref*{cdp:eq:C.EN.12}}}
\expandafter\def\csname SRef@cdp:eq:C.EN.13\endcsname{\hyperlink{tag-00GD}{II.\ref*{cdp:eq:C.EN.13}}}
\expandafter\def\csname SRef@cdp:eq:C.EN.14\endcsname{\hyperlink{tag-00GE}{II.\ref*{cdp:eq:C.EN.14}}}
\expandafter\def\csname SRef@cdp:eq:C.EN.15\endcsname{\hyperlink{tag-00GF}{II.\ref*{cdp:eq:C.EN.15}}}
\expandafter\def\csname SRef@cdp:eq:C.EN.16\endcsname{\hyperlink{tag-00GG}{II.\ref*{cdp:eq:C.EN.16}}}
\expandafter\def\csname SRef@cdp:eq:C.HQ.2\endcsname{\hyperlink{tag-00MS}{II.\ref*{cdp:eq:C.HQ.2}}}
\expandafter\def\csname SRef@cdp:eq:C.EN.17\endcsname{\hyperlink{tag-00GH}{II.\ref*{cdp:eq:C.EN.17}}}
\expandafter\def\csname SRef@cdp:eq:C.EN.S1\endcsname{\hyperlink{tag-00H7}{II.\ref*{cdp:eq:C.EN.S1}}}
\expandafter\def\csname SRef@cdp:eq:C.EN.S2\endcsname{\hyperlink{tag-00H8}{II.\ref*{cdp:eq:C.EN.S2}}}
\expandafter\def\csname SRef@cdp:eq:C.EN.S3\endcsname{\hyperlink{tag-00H9}{II.\ref*{cdp:eq:C.EN.S3}}}
\expandafter\def\csname SRef@cdp:eq:C.EN.S4\endcsname{\hyperlink{tag-00HA}{II.\ref*{cdp:eq:C.EN.S4}}}
\expandafter\def\csname SRef@cdp:eq:C.EN.S5\endcsname{\hyperlink{tag-00HB}{II.\ref*{cdp:eq:C.EN.S5}}}
\expandafter\def\csname SRef@cdp:eq:C.EN.S6\endcsname{\hyperlink{tag-00HC}{II.\ref*{cdp:eq:C.EN.S6}}}
\expandafter\def\csname SRef@cdp:eq:C.EN.S7\endcsname{\hyperlink{tag-00HD}{II.\ref*{cdp:eq:C.EN.S7}}}
\expandafter\def\csname SRef@cdp:eq:C.EN.R1\endcsname{\hyperlink{tag-00GV}{II.\ref*{cdp:eq:C.EN.R1}}}
\expandafter\def\csname SRef@cdp:eq:C.EN.R2\endcsname{\hyperlink{tag-00GZ}{II.\ref*{cdp:eq:C.EN.R2}}}
\expandafter\def\csname SRef@cdp:eq:C.HQ.3\endcsname{\hyperlink{tag-00MT}{II.\ref*{cdp:eq:C.HQ.3}}}
\expandafter\def\csname SRef@cdp:eq:C.EN.R3\endcsname{\hyperlink{tag-00H0}{II.\ref*{cdp:eq:C.EN.R3}}}
\expandafter\def\csname SRef@cdp:eq:C.EN.R4\endcsname{\hyperlink{tag-00H1}{II.\ref*{cdp:eq:C.EN.R4}}}
\expandafter\def\csname SRef@cdp:eq:C.EN.R5\endcsname{\hyperlink{tag-00H2}{II.\ref*{cdp:eq:C.EN.R5}}}
\expandafter\def\csname SRef@cdp:eq:C.EN.R6\endcsname{\hyperlink{tag-00H3}{II.\ref*{cdp:eq:C.EN.R6}}}
\expandafter\def\csname SRef@cdp:eq:C.EN.R7\endcsname{\hyperlink{tag-00H4}{II.\ref*{cdp:eq:C.EN.R7}}}
\expandafter\def\csname SRef@cdp:eq:C.EN.R8\endcsname{\hyperlink{tag-00H5}{II.\ref*{cdp:eq:C.EN.R8}}}
\expandafter\def\csname SRef@cdp:eq:C.EN.R9\endcsname{\hyperlink{tag-00H6}{II.\ref*{cdp:eq:C.EN.R9}}}
\expandafter\def\csname SRef@cdp:eq:C.EN.R10\endcsname{\hyperlink{tag-00GW}{II.\ref*{cdp:eq:C.EN.R10}}}
\expandafter\def\csname SRef@cdp:eq:C.EN.R11\endcsname{\hyperlink{tag-00GX}{II.\ref*{cdp:eq:C.EN.R11}}}
\expandafter\def\csname SRef@cdp:eq:C.EN.R12\endcsname{\hyperlink{tag-00GY}{II.\ref*{cdp:eq:C.EN.R12}}}
\expandafter\def\csname SRef@cdp:eq:C.HQ.4\endcsname{\hyperlink{tag-00MU}{II.\ref*{cdp:eq:C.HQ.4}}}
\expandafter\def\csname SRef@cdp:eq:C.EN.G1\endcsname{\hyperlink{tag-00GR}{II.\ref*{cdp:eq:C.EN.G1}}}
\expandafter\def\csname SRef@cdp:eq:C.EN.G2\endcsname{\hyperlink{tag-00GS}{II.\ref*{cdp:eq:C.EN.G2}}}
\expandafter\def\csname SRef@cdp:eq:C.EN.G3\endcsname{\hyperlink{tag-00GT}{II.\ref*{cdp:eq:C.EN.G3}}}
\expandafter\def\csname SRef@cdp:eq:C.EN.G4\endcsname{\hyperlink{tag-00GU}{II.\ref*{cdp:eq:C.EN.G4}}}
\expandafter\def\csname SRef@cdp:eq:C.ET.1\endcsname{\hyperlink{tag-00J5}{II.\ref*{cdp:eq:C.ET.1}}}
\expandafter\def\csname SRef@cdp:eq:C.ET.2\endcsname{\hyperlink{tag-00J9}{II.\ref*{cdp:eq:C.ET.2}}}
\expandafter\def\csname SRef@cdp:eq:C.ET.3\endcsname{\hyperlink{tag-00JA}{II.\ref*{cdp:eq:C.ET.3}}}
\expandafter\def\csname SRef@cdp:eq:C.ET.4\endcsname{\hyperlink{tag-00JB}{II.\ref*{cdp:eq:C.ET.4}}}
\expandafter\def\csname SRef@cdp:eq:C.ET.5\endcsname{\hyperlink{tag-00JC}{II.\ref*{cdp:eq:C.ET.5}}}
\expandafter\def\csname SRef@cdp:eq:C.ET.6\endcsname{\hyperlink{tag-00JD}{II.\ref*{cdp:eq:C.ET.6}}}
\expandafter\def\csname SRef@cdp:eq:C.HQ.5\endcsname{\hyperlink{tag-00MV}{II.\ref*{cdp:eq:C.HQ.5}}}
\expandafter\def\csname SRef@cdp:eq:C.ET.7\endcsname{\hyperlink{tag-00JE}{II.\ref*{cdp:eq:C.ET.7}}}
\expandafter\def\csname SRef@cdp:eq:C.ET.8\endcsname{\hyperlink{tag-00JF}{II.\ref*{cdp:eq:C.ET.8}}}
\expandafter\def\csname SRef@cdp:eq:C.ET.9\endcsname{\hyperlink{tag-00JG}{II.\ref*{cdp:eq:C.ET.9}}}
\expandafter\def\csname SRef@cdp:eq:C.ET.10\endcsname{\hyperlink{tag-00J6}{II.\ref*{cdp:eq:C.ET.10}}}
\expandafter\def\csname SRef@cdp:eq:C.ET.11\endcsname{\hyperlink{tag-00J7}{II.\ref*{cdp:eq:C.ET.11}}}
\expandafter\def\csname SRef@cdp:eq:C.ET.12\endcsname{\hyperlink{tag-00J8}{II.\ref*{cdp:eq:C.ET.12}}}
\expandafter\def\csname SRef@cdp:eq:C.SF.1\endcsname{\hyperlink{tag-010R}{II.\ref*{cdp:eq:C.SF.1}}}
\expandafter\def\csname SRef@cdp:eq:C.SF.2\endcsname{\hyperlink{tag-010W}{II.\ref*{cdp:eq:C.SF.2}}}
\expandafter\def\csname SRef@cdp:eq:C.SF.3\endcsname{\hyperlink{tag-010X}{II.\ref*{cdp:eq:C.SF.3}}}
\expandafter\def\csname SRef@cdp:eq:C.SF.4\endcsname{\hyperlink{tag-010Y}{II.\ref*{cdp:eq:C.SF.4}}}
\expandafter\def\csname SRef@cdp:eq:C.HQ.6\endcsname{\hyperlink{tag-00MW}{II.\ref*{cdp:eq:C.HQ.6}}}
\expandafter\def\csname SRef@cdp:eq:C.SF.5\endcsname{\hyperlink{tag-010Z}{II.\ref*{cdp:eq:C.SF.5}}}
\expandafter\def\csname SRef@cdp:eq:C.SF.6\endcsname{\hyperlink{tag-0110}{II.\ref*{cdp:eq:C.SF.6}}}
\expandafter\def\csname SRef@cdp:eq:C.SF.7\endcsname{\hyperlink{tag-0111}{II.\ref*{cdp:eq:C.SF.7}}}
\expandafter\def\csname SRef@cdp:eq:C.SF.8\endcsname{\hyperlink{tag-0112}{II.\ref*{cdp:eq:C.SF.8}}}
\expandafter\def\csname SRef@cdp:eq:C.SF.9\endcsname{\hyperlink{tag-0113}{II.\ref*{cdp:eq:C.SF.9}}}
\expandafter\def\csname SRef@cdp:eq:C.SF.10\endcsname{\hyperlink{tag-010S}{II.\ref*{cdp:eq:C.SF.10}}}
\expandafter\def\csname SRef@cdp:eq:C.SF.11\endcsname{\hyperlink{tag-010T}{II.\ref*{cdp:eq:C.SF.11}}}
\expandafter\def\csname SRef@cdp:eq:C.SF.12\endcsname{\hyperlink{tag-010U}{II.\ref*{cdp:eq:C.SF.12}}}
\expandafter\def\csname SRef@cdp:eq:C.SF.13\endcsname{\hyperlink{tag-010V}{II.\ref*{cdp:eq:C.SF.13}}}
\expandafter\def\csname SRef@cdp:eq:C.EO.1\endcsname{\hyperlink{tag-00HE}{II.\ref*{cdp:eq:C.EO.1}}}
\expandafter\def\csname SRef@cdp:eq:C.HQ.7\endcsname{\hyperlink{tag-00MX}{II.\ref*{cdp:eq:C.HQ.7}}}
\expandafter\def\csname SRef@cdp:eq:C.EO.2\endcsname{\hyperlink{tag-00HM}{II.\ref*{cdp:eq:C.EO.2}}}
\expandafter\def\csname SRef@cdp:eq:C.EO.3\endcsname{\hyperlink{tag-00HN}{II.\ref*{cdp:eq:C.EO.3}}}
\expandafter\def\csname SRef@cdp:eq:C.EO.4\endcsname{\hyperlink{tag-00HO}{II.\ref*{cdp:eq:C.EO.4}}}
\expandafter\def\csname SRef@cdp:eq:C.EO.5\endcsname{\hyperlink{tag-00HP}{II.\ref*{cdp:eq:C.EO.5}}}
\expandafter\def\csname SRef@cdp:eq:C.EO.6\endcsname{\hyperlink{tag-00HQ}{II.\ref*{cdp:eq:C.EO.6}}}
\expandafter\def\csname SRef@cdp:eq:C.EO.H1\endcsname{\hyperlink{tag-00HU}{II.\ref*{cdp:eq:C.EO.H1}}}
\expandafter\def\csname SRef@cdp:eq:C.EO.H2\endcsname{\hyperlink{tag-00HV}{II.\ref*{cdp:eq:C.EO.H2}}}
\expandafter\def\csname SRef@cdp:eq:C.EO.H3\endcsname{\hyperlink{tag-00HW}{II.\ref*{cdp:eq:C.EO.H3}}}
\expandafter\def\csname SRef@cdp:eq:C.EO.H4\endcsname{\hyperlink{tag-00HX}{II.\ref*{cdp:eq:C.EO.H4}}}
\expandafter\def\csname SRef@cdp:eq:C.EO.H5\endcsname{\hyperlink{tag-00HY}{II.\ref*{cdp:eq:C.EO.H5}}}
\expandafter\def\csname SRef@cdp:eq:C.HQ.8\endcsname{\hyperlink{tag-00MY}{II.\ref*{cdp:eq:C.HQ.8}}}
\expandafter\def\csname SRef@cdp:eq:C.EO.H6\endcsname{\hyperlink{tag-00HZ}{II.\ref*{cdp:eq:C.EO.H6}}}
\expandafter\def\csname SRef@cdp:eq:C.EO.H7\endcsname{\hyperlink{tag-00I0}{II.\ref*{cdp:eq:C.EO.H7}}}
\expandafter\def\csname SRef@cdp:eq:C.EO.7\endcsname{\hyperlink{tag-00HR}{II.\ref*{cdp:eq:C.EO.7}}}
\expandafter\def\csname SRef@cdp:eq:C.EO.8\endcsname{\hyperlink{tag-00HS}{II.\ref*{cdp:eq:C.EO.8}}}
\expandafter\def\csname SRef@cdp:eq:C.EO.9\endcsname{\hyperlink{tag-00HT}{II.\ref*{cdp:eq:C.EO.9}}}
\expandafter\def\csname SRef@cdp:eq:C.EO.10\endcsname{\hyperlink{tag-00HF}{II.\ref*{cdp:eq:C.EO.10}}}
\expandafter\def\csname SRef@cdp:eq:C.EO.11\endcsname{\hyperlink{tag-00HG}{II.\ref*{cdp:eq:C.EO.11}}}
\expandafter\def\csname SRef@cdp:eq:C.EO.12\endcsname{\hyperlink{tag-00HH}{II.\ref*{cdp:eq:C.EO.12}}}
\expandafter\def\csname SRef@cdp:eq:C.EO.13\endcsname{\hyperlink{tag-00HI}{II.\ref*{cdp:eq:C.EO.13}}}
\expandafter\def\csname SRef@cdp:eq:C.EO.14\endcsname{\hyperlink{tag-00HJ}{II.\ref*{cdp:eq:C.EO.14}}}
\expandafter\def\csname SRef@cdp:eq:C.HQ.9\endcsname{\hyperlink{tag-00MZ}{II.\ref*{cdp:eq:C.HQ.9}}}
\expandafter\def\csname SRef@cdp:eq:C.EO.15\endcsname{\hyperlink{tag-00HK}{II.\ref*{cdp:eq:C.EO.15}}}
\expandafter\def\csname SRef@cdp:eq:C.EO.16\endcsname{\hyperlink{tag-00HL}{II.\ref*{cdp:eq:C.EO.16}}}
\expandafter\def\csname SRef@cdp:eq:C.EH.1\endcsname{\hyperlink{tag-00FQ}{II.\ref*{cdp:eq:C.EH.1}}}
\expandafter\def\csname SRef@cdp:eq:C.EH.2\endcsname{\hyperlink{tag-00FR}{II.\ref*{cdp:eq:C.EH.2}}}
\expandafter\def\csname SRef@cdp:eq:C.EH.3\endcsname{\hyperlink{tag-00FS}{II.\ref*{cdp:eq:C.EH.3}}}
\expandafter\def\csname SRef@cdp:eq:C.EH.4\endcsname{\hyperlink{tag-00FT}{II.\ref*{cdp:eq:C.EH.4}}}
\expandafter\def\csname SRef@cdp:eq:C.EH.5\endcsname{\hyperlink{tag-00FU}{II.\ref*{cdp:eq:C.EH.5}}}
\expandafter\def\csname SRef@cdp:eq:C.EH.6\endcsname{\hyperlink{tag-00FV}{II.\ref*{cdp:eq:C.EH.6}}}
\expandafter\def\csname SRef@cdp:eq:C.EH.7\endcsname{\hyperlink{tag-00FW}{II.\ref*{cdp:eq:C.EH.7}}}
\expandafter\def\csname SRef@cdp:eq:C.EH.8\endcsname{\hyperlink{tag-00FX}{II.\ref*{cdp:eq:C.EH.8}}}
\expandafter\def\csname SRef@cdp:eq:C.HQ.10\endcsname{\hyperlink{tag-00MR}{II.\ref*{cdp:eq:C.HQ.10}}}
\expandafter\def\csname SRef@cdp:eq:C.DH.1\endcsname{\hyperlink{tag-00F6}{II.\ref*{cdp:eq:C.DH.1}}}
\expandafter\def\csname SRef@cdp:eq:C.DH.2\endcsname{\hyperlink{tag-00FG}{II.\ref*{cdp:eq:C.DH.2}}}
\expandafter\def\csname SRef@cdp:eq:C.DH.3\endcsname{\hyperlink{tag-00FH}{II.\ref*{cdp:eq:C.DH.3}}}
\expandafter\def\csname SRef@cdp:eq:C.DH.4\endcsname{\hyperlink{tag-00FI}{II.\ref*{cdp:eq:C.DH.4}}}
\expandafter\def\csname SRef@cdp:eq:C.DH.5\endcsname{\hyperlink{tag-00FJ}{II.\ref*{cdp:eq:C.DH.5}}}
\expandafter\def\csname SRef@cdp:eq:C.DH.6\endcsname{\hyperlink{tag-00FK}{II.\ref*{cdp:eq:C.DH.6}}}
\expandafter\def\csname SRef@cdp:eq:C.DH.7\endcsname{\hyperlink{tag-00FL}{II.\ref*{cdp:eq:C.DH.7}}}
\expandafter\def\csname SRef@cdp:eq:C.DH.8\endcsname{\hyperlink{tag-00FM}{II.\ref*{cdp:eq:C.DH.8}}}
\expandafter\def\csname SRef@cdp:eq:C.DH.9\endcsname{\hyperlink{tag-00FN}{II.\ref*{cdp:eq:C.DH.9}}}
\expandafter\def\csname SRef@cdp:eq:C.DH.10\endcsname{\hyperlink{tag-00F7}{II.\ref*{cdp:eq:C.DH.10}}}
\expandafter\def\csname SRef@cdp:eq:C.FD.1\endcsname{\hyperlink{tag-00K2}{II.\ref*{cdp:eq:C.FD.1}}}
\expandafter\def\csname SRef@cdp:eq:C.DH.11\endcsname{\hyperlink{tag-00F8}{II.\ref*{cdp:eq:C.DH.11}}}
\expandafter\def\csname SRef@cdp:eq:C.DH.12\endcsname{\hyperlink{tag-00F9}{II.\ref*{cdp:eq:C.DH.12}}}
\expandafter\def\csname SRef@cdp:eq:C.DH.13\endcsname{\hyperlink{tag-00FA}{II.\ref*{cdp:eq:C.DH.13}}}
\expandafter\def\csname SRef@cdp:eq:C.DH.14\endcsname{\hyperlink{tag-00FB}{II.\ref*{cdp:eq:C.DH.14}}}
\expandafter\def\csname SRef@cdp:eq:C.DH.14a\endcsname{\hyperlink{tag-00FC}{II.\ref*{cdp:eq:C.DH.14a}}}
\expandafter\def\csname SRef@cdp:eq:C.DH.14b\endcsname{\hyperlink{tag-00FD}{II.\ref*{cdp:eq:C.DH.14b}}}
\expandafter\def\csname SRef@cdp:eq:C.DH.15\endcsname{\hyperlink{tag-00FE}{II.\ref*{cdp:eq:C.DH.15}}}
\expandafter\def\csname SRef@cdp:eq:C.DH.16\endcsname{\hyperlink{tag-00FF}{II.\ref*{cdp:eq:C.DH.16}}}
\expandafter\def\csname SRef@cdp:eq:C.NH.1\endcsname{\hyperlink{tag-00S9}{II.\ref*{cdp:eq:C.NH.1}}}
\expandafter\def\csname SRef@cdp:eq:C.NH.2\endcsname{\hyperlink{tag-00SJ}{II.\ref*{cdp:eq:C.NH.2}}}
\expandafter\def\csname SRef@cdp:eq:C.FD.2\endcsname{\hyperlink{tag-00K3}{II.\ref*{cdp:eq:C.FD.2}}}
\expandafter\def\csname SRef@cdp:eq:C.NH.3\endcsname{\hyperlink{tag-00SK}{II.\ref*{cdp:eq:C.NH.3}}}
\expandafter\def\csname SRef@cdp:eq:C.NH.4\endcsname{\hyperlink{tag-00SL}{II.\ref*{cdp:eq:C.NH.4}}}
\expandafter\def\csname SRef@cdp:eq:C.NH.5\endcsname{\hyperlink{tag-00SM}{II.\ref*{cdp:eq:C.NH.5}}}
\expandafter\def\csname SRef@cdp:eq:C.NH.6\endcsname{\hyperlink{tag-00SN}{II.\ref*{cdp:eq:C.NH.6}}}
\expandafter\def\csname SRef@cdp:eq:C.NH.7\endcsname{\hyperlink{tag-00SO}{II.\ref*{cdp:eq:C.NH.7}}}
\expandafter\def\csname SRef@cdp:eq:C.NH.8\endcsname{\hyperlink{tag-00SP}{II.\ref*{cdp:eq:C.NH.8}}}
\expandafter\def\csname SRef@cdp:eq:C.NH.9\endcsname{\hyperlink{tag-00SQ}{II.\ref*{cdp:eq:C.NH.9}}}
\expandafter\def\csname SRef@cdp:eq:C.NH.10\endcsname{\hyperlink{tag-00SA}{II.\ref*{cdp:eq:C.NH.10}}}
\expandafter\def\csname SRef@cdp:eq:C.NH.11\endcsname{\hyperlink{tag-00SB}{II.\ref*{cdp:eq:C.NH.11}}}
\expandafter\def\csname SRef@cdp:eq:C.NH.12\endcsname{\hyperlink{tag-00SC}{II.\ref*{cdp:eq:C.NH.12}}}
\expandafter\def\csname SRef@cdp:eq:C.FD.3\endcsname{\hyperlink{tag-00K4}{II.\ref*{cdp:eq:C.FD.3}}}
\expandafter\def\csname SRef@cdp:eq:C.NH.13\endcsname{\hyperlink{tag-00SD}{II.\ref*{cdp:eq:C.NH.13}}}
\expandafter\def\csname SRef@cdp:eq:C.NH.14\endcsname{\hyperlink{tag-00SE}{II.\ref*{cdp:eq:C.NH.14}}}
\expandafter\def\csname SRef@cdp:eq:C.NH.15\endcsname{\hyperlink{tag-00SF}{II.\ref*{cdp:eq:C.NH.15}}}
\expandafter\def\csname SRef@cdp:eq:C.NH.16\endcsname{\hyperlink{tag-00SG}{II.\ref*{cdp:eq:C.NH.16}}}
\expandafter\def\csname SRef@cdp:eq:C.NH.17\endcsname{\hyperlink{tag-00SH}{II.\ref*{cdp:eq:C.NH.17}}}
\expandafter\def\csname SRef@cdp:eq:C.NH.18\endcsname{\hyperlink{tag-00SI}{II.\ref*{cdp:eq:C.NH.18}}}
\expandafter\def\csname SRef@cdp:eq:C.SH.1\endcsname{\hyperlink{tag-0114}{II.\ref*{cdp:eq:C.SH.1}}}
\expandafter\def\csname SRef@cdp:eq:C.SH.2\endcsname{\hyperlink{tag-0119}{II.\ref*{cdp:eq:C.SH.2}}}
\expandafter\def\csname SRef@cdp:eq:C.SH.3\endcsname{\hyperlink{tag-011A}{II.\ref*{cdp:eq:C.SH.3}}}
\expandafter\def\csname SRef@cdp:eq:C.SH.4\endcsname{\hyperlink{tag-011B}{II.\ref*{cdp:eq:C.SH.4}}}
\expandafter\def\csname SRef@cdp:eq:C.FD.4\endcsname{\hyperlink{tag-00K5}{II.\ref*{cdp:eq:C.FD.4}}}
\expandafter\def\csname SRef@cdp:eq:C.SH.5\endcsname{\hyperlink{tag-011C}{II.\ref*{cdp:eq:C.SH.5}}}
\expandafter\def\csname SRef@cdp:eq:C.SH.6\endcsname{\hyperlink{tag-011D}{II.\ref*{cdp:eq:C.SH.6}}}
\expandafter\def\csname SRef@cdp:eq:C.SH.7\endcsname{\hyperlink{tag-011E}{II.\ref*{cdp:eq:C.SH.7}}}
\expandafter\def\csname SRef@cdp:eq:C.SH.8\endcsname{\hyperlink{tag-011F}{II.\ref*{cdp:eq:C.SH.8}}}
\expandafter\def\csname SRef@cdp:eq:C.SH.9\endcsname{\hyperlink{tag-011G}{II.\ref*{cdp:eq:C.SH.9}}}
\expandafter\def\csname SRef@cdp:eq:C.SH.10\endcsname{\hyperlink{tag-0115}{II.\ref*{cdp:eq:C.SH.10}}}
\expandafter\def\csname SRef@cdp:eq:C.SH.11\endcsname{\hyperlink{tag-0116}{II.\ref*{cdp:eq:C.SH.11}}}
\expandafter\def\csname SRef@cdp:eq:C.SH.12\endcsname{\hyperlink{tag-0117}{II.\ref*{cdp:eq:C.SH.12}}}
\expandafter\def\csname SRef@cdp:eq:C.SH.13\endcsname{\hyperlink{tag-0118}{II.\ref*{cdp:eq:C.SH.13}}}
\expandafter\def\csname SRef@cdp:eq:C.RU.1\endcsname{\hyperlink{tag-010C}{II.\ref*{cdp:eq:C.RU.1}}}
\expandafter\def\csname SRef@cdp:eq:C.FD.5\endcsname{\hyperlink{tag-00K6}{II.\ref*{cdp:eq:C.FD.5}}}
\expandafter\def\csname SRef@cdp:eq:C.RU.2\endcsname{\hyperlink{tag-010D}{II.\ref*{cdp:eq:C.RU.2}}}
\expandafter\def\csname SRef@cdp:eq:C.RU.3\endcsname{\hyperlink{tag-010E}{II.\ref*{cdp:eq:C.RU.3}}}
\expandafter\def\csname SRef@cdp:eq:C.RU.4\endcsname{\hyperlink{tag-010F}{II.\ref*{cdp:eq:C.RU.4}}}
\expandafter\def\csname SRef@cdp:eq:C.RU.5\endcsname{\hyperlink{tag-010G}{II.\ref*{cdp:eq:C.RU.5}}}
\expandafter\def\csname SRef@cdp:eq:C.RU.6\endcsname{\hyperlink{tag-010H}{II.\ref*{cdp:eq:C.RU.6}}}
\expandafter\def\csname SRef@cdp:eq:C.RU.7\endcsname{\hyperlink{tag-010I}{II.\ref*{cdp:eq:C.RU.7}}}
\expandafter\def\csname SRef@cdp:eq:C.RU.8\endcsname{\hyperlink{tag-010J}{II.\ref*{cdp:eq:C.RU.8}}}
\expandafter\def\csname SRef@cdp:eq:C.RU.9\endcsname{\hyperlink{tag-010K}{II.\ref*{cdp:eq:C.RU.9}}}
\expandafter\def\csname SRef@cdp:eq:C.HA.1\endcsname{\hyperlink{tag-00L9}{II.\ref*{cdp:eq:C.HA.1}}}
\expandafter\def\csname SRef@cdp:eq:C.HA.2\endcsname{\hyperlink{tag-00LA}{II.\ref*{cdp:eq:C.HA.2}}}
\expandafter\def\csname SRef@cdp:eq:C.FD.6\endcsname{\hyperlink{tag-00K7}{II.\ref*{cdp:eq:C.FD.6}}}
\expandafter\def\csname SRef@cdp:eq:C.HA.3\endcsname{\hyperlink{tag-00LB}{II.\ref*{cdp:eq:C.HA.3}}}
\expandafter\def\csname SRef@cdp:eq:C.HA.4\endcsname{\hyperlink{tag-00LC}{II.\ref*{cdp:eq:C.HA.4}}}
\expandafter\def\csname SRef@cdp:eq:C.HA.5\endcsname{\hyperlink{tag-00LD}{II.\ref*{cdp:eq:C.HA.5}}}
\expandafter\def\csname SRef@cdp:eq:C.HA.6\endcsname{\hyperlink{tag-00LE}{II.\ref*{cdp:eq:C.HA.6}}}
\expandafter\def\csname SRef@cdp:eq:C.HC.1\endcsname{\hyperlink{tag-00LO}{II.\ref*{cdp:eq:C.HC.1}}}
\expandafter\def\csname SRef@cdp:eq:C.HC.2\endcsname{\hyperlink{tag-00LR}{II.\ref*{cdp:eq:C.HC.2}}}
\expandafter\def\csname SRef@cdp:eq:C.HC.3\endcsname{\hyperlink{tag-00LS}{II.\ref*{cdp:eq:C.HC.3}}}
\expandafter\def\csname SRef@cdp:eq:C.HC.4\endcsname{\hyperlink{tag-00LT}{II.\ref*{cdp:eq:C.HC.4}}}
\expandafter\def\csname SRef@cdp:eq:C.HC.5\endcsname{\hyperlink{tag-00LU}{II.\ref*{cdp:eq:C.HC.5}}}
\expandafter\def\csname SRef@cdp:eq:C.HC.6\endcsname{\hyperlink{tag-00LV}{II.\ref*{cdp:eq:C.HC.6}}}
\expandafter\def\csname SRef@cdp:eq:C.RF.5\endcsname{\hyperlink{tag-00YY}{II.\ref*{cdp:eq:C.RF.5}}}
\expandafter\def\csname SRef@cdp:eq:C.HC.7\endcsname{\hyperlink{tag-00LW}{II.\ref*{cdp:eq:C.HC.7}}}
\expandafter\def\csname SRef@cdp:eq:C.HC.8\endcsname{\hyperlink{tag-00LX}{II.\ref*{cdp:eq:C.HC.8}}}
\expandafter\def\csname SRef@cdp:eq:C.HC.9\endcsname{\hyperlink{tag-00LY}{II.\ref*{cdp:eq:C.HC.9}}}
\expandafter\def\csname SRef@cdp:eq:C.HC.10\endcsname{\hyperlink{tag-00LP}{II.\ref*{cdp:eq:C.HC.10}}}
\expandafter\def\csname SRef@cdp:eq:C.HC.11\endcsname{\hyperlink{tag-00LQ}{II.\ref*{cdp:eq:C.HC.11}}}
\expandafter\def\csname SRef@cdp:eq:C.HG.1\endcsname{\hyperlink{tag-00LZ}{II.\ref*{cdp:eq:C.HG.1}}}
\expandafter\def\csname SRef@cdp:eq:C.HG.2\endcsname{\hyperlink{tag-00M8}{II.\ref*{cdp:eq:C.HG.2}}}
\expandafter\def\csname SRef@cdp:eq:C.HG.3\endcsname{\hyperlink{tag-00M9}{II.\ref*{cdp:eq:C.HG.3}}}
\expandafter\def\csname SRef@cdp:eq:C.HG.4\endcsname{\hyperlink{tag-00MA}{II.\ref*{cdp:eq:C.HG.4}}}
\expandafter\def\csname SRef@cdp:eq:C.HG.5\endcsname{\hyperlink{tag-00MB}{II.\ref*{cdp:eq:C.HG.5}}}
\expandafter\def\csname SRef@cdp:eq:C.RF.1\endcsname{\hyperlink{tag-00YU}{II.\ref*{cdp:eq:C.RF.1}}}
\expandafter\def\csname SRef@cdp:eq:C.HG.6\endcsname{\hyperlink{tag-00MC}{II.\ref*{cdp:eq:C.HG.6}}}
\expandafter\def\csname SRef@cdp:eq:C.HG.7\endcsname{\hyperlink{tag-00MD}{II.\ref*{cdp:eq:C.HG.7}}}
\expandafter\def\csname SRef@cdp:eq:C.HG.8\endcsname{\hyperlink{tag-00ME}{II.\ref*{cdp:eq:C.HG.8}}}
\expandafter\def\csname SRef@cdp:eq:C.HG.9\endcsname{\hyperlink{tag-00MF}{II.\ref*{cdp:eq:C.HG.9}}}
\expandafter\def\csname SRef@cdp:eq:C.HG.10\endcsname{\hyperlink{tag-00M0}{II.\ref*{cdp:eq:C.HG.10}}}
\expandafter\def\csname SRef@cdp:eq:C.HG.11\endcsname{\hyperlink{tag-00M1}{II.\ref*{cdp:eq:C.HG.11}}}
\expandafter\def\csname SRef@cdp:eq:C.HG.12\endcsname{\hyperlink{tag-00M2}{II.\ref*{cdp:eq:C.HG.12}}}
\expandafter\def\csname SRef@cdp:eq:C.HG.13\endcsname{\hyperlink{tag-00M3}{II.\ref*{cdp:eq:C.HG.13}}}
\expandafter\def\csname SRef@cdp:eq:C.HG.14\endcsname{\hyperlink{tag-00M4}{II.\ref*{cdp:eq:C.HG.14}}}
\expandafter\def\csname SRef@cdp:eq:C.HG.15\endcsname{\hyperlink{tag-00M5}{II.\ref*{cdp:eq:C.HG.15}}}
\expandafter\def\csname SRef@cdp:eq:C.RF.2\endcsname{\hyperlink{tag-00YV}{II.\ref*{cdp:eq:C.RF.2}}}
\expandafter\def\csname SRef@cdp:eq:C.HG.16\endcsname{\hyperlink{tag-00M6}{II.\ref*{cdp:eq:C.HG.16}}}
\expandafter\def\csname SRef@cdp:eq:C.HG.17\endcsname{\hyperlink{tag-00M7}{II.\ref*{cdp:eq:C.HG.17}}}
\expandafter\def\csname SRef@cdp:eq:C.NS.1\endcsname{\hyperlink{tag-00TA}{II.\ref*{cdp:eq:C.NS.1}}}
\expandafter\def\csname SRef@cdp:eq:C.NS.2\endcsname{\hyperlink{tag-00TL}{II.\ref*{cdp:eq:C.NS.2}}}
\expandafter\def\csname SRef@cdp:eq:C.NS.3\endcsname{\hyperlink{tag-00TN}{II.\ref*{cdp:eq:C.NS.3}}}
\expandafter\def\csname SRef@cdp:eq:C.NS.4\endcsname{\hyperlink{tag-00TO}{II.\ref*{cdp:eq:C.NS.4}}}
\expandafter\def\csname SRef@cdp:eq:C.NS.5\endcsname{\hyperlink{tag-00TP}{II.\ref*{cdp:eq:C.NS.5}}}
\expandafter\def\csname SRef@cdp:eq:C.NS.6\endcsname{\hyperlink{tag-00TQ}{II.\ref*{cdp:eq:C.NS.6}}}
\expandafter\def\csname SRef@cdp:eq:C.NS.7\endcsname{\hyperlink{tag-00TR}{II.\ref*{cdp:eq:C.NS.7}}}
\expandafter\def\csname SRef@cdp:eq:C.NS.8\endcsname{\hyperlink{tag-00TS}{II.\ref*{cdp:eq:C.NS.8}}}
\expandafter\def\csname SRef@cdp:eq:C.RF.3\endcsname{\hyperlink{tag-00YW}{II.\ref*{cdp:eq:C.RF.3}}}
\expandafter\def\csname SRef@cdp:eq:C.NS.9\endcsname{\hyperlink{tag-00TT}{II.\ref*{cdp:eq:C.NS.9}}}
\expandafter\def\csname SRef@cdp:eq:C.NS.10\endcsname{\hyperlink{tag-00TB}{II.\ref*{cdp:eq:C.NS.10}}}
\expandafter\def\csname SRef@cdp:eq:C.NS.11\endcsname{\hyperlink{tag-00TC}{II.\ref*{cdp:eq:C.NS.11}}}
\expandafter\def\csname SRef@cdp:eq:C.NS.12\endcsname{\hyperlink{tag-00TD}{II.\ref*{cdp:eq:C.NS.12}}}
\expandafter\def\csname SRef@cdp:eq:C.NS.13\endcsname{\hyperlink{tag-00TE}{II.\ref*{cdp:eq:C.NS.13}}}
\expandafter\def\csname SRef@cdp:eq:C.NS.14\endcsname{\hyperlink{tag-00TF}{II.\ref*{cdp:eq:C.NS.14}}}
\expandafter\def\csname SRef@cdp:eq:C.NS.15\endcsname{\hyperlink{tag-00TG}{II.\ref*{cdp:eq:C.NS.15}}}
\expandafter\def\csname SRef@cdp:eq:C.NS.16\endcsname{\hyperlink{tag-00TH}{II.\ref*{cdp:eq:C.NS.16}}}
\expandafter\def\csname SRef@cdp:eq:C.NS.17\endcsname{\hyperlink{tag-00TI}{II.\ref*{cdp:eq:C.NS.17}}}
\expandafter\def\csname SRef@cdp:eq:C.NS.18\endcsname{\hyperlink{tag-00TJ}{II.\ref*{cdp:eq:C.NS.18}}}
\expandafter\def\csname SRef@cdp:eq:C.RF.4\endcsname{\hyperlink{tag-00YX}{II.\ref*{cdp:eq:C.RF.4}}}
\expandafter\def\csname SRef@cdp:eq:C.NS.19\endcsname{\hyperlink{tag-00TK}{II.\ref*{cdp:eq:C.NS.19}}}
\expandafter\def\csname SRef@cdp:eq:C.NS.20\endcsname{\hyperlink{tag-00TM}{II.\ref*{cdp:eq:C.NS.20}}}
\expandafter\def\csname SRef@cdp:eq:C.SN.1\endcsname{\hyperlink{tag-0122}{II.\ref*{cdp:eq:C.SN.1}}}
\expandafter\def\csname SRef@cdp:eq:C.SN.2\endcsname{\hyperlink{tag-0124}{II.\ref*{cdp:eq:C.SN.2}}}
\expandafter\def\csname SRef@cdp:eq:C.SN.3\endcsname{\hyperlink{tag-0125}{II.\ref*{cdp:eq:C.SN.3}}}
\expandafter\def\csname SRef@cdp:eq:C.SN.4\endcsname{\hyperlink{tag-0126}{II.\ref*{cdp:eq:C.SN.4}}}
\expandafter\def\csname SRef@cdp:eq:C.SN.5\endcsname{\hyperlink{tag-0127}{II.\ref*{cdp:eq:C.SN.5}}}
\expandafter\def\csname SRef@cdp:eq:C.SN.6\endcsname{\hyperlink{tag-0128}{II.\ref*{cdp:eq:C.SN.6}}}
\expandafter\def\csname SRef@cdp:eq:C.SN.7\endcsname{\hyperlink{tag-0129}{II.\ref*{cdp:eq:C.SN.7}}}
\expandafter\def\csname SRef@cdp:eq:C.SN.8\endcsname{\hyperlink{tag-012A}{II.\ref*{cdp:eq:C.SN.8}}}
\expandafter\def\csname SRef@cdp:eq:C.RF.9\endcsname{\hyperlink{tag-00Z2}{II.\ref*{cdp:eq:C.RF.9}}}
\expandafter\def\csname SRef@cdp:eq:C.SN.9\endcsname{\hyperlink{tag-012B}{II.\ref*{cdp:eq:C.SN.9}}}
\expandafter\def\csname SRef@cdp:eq:C.SN.10\endcsname{\hyperlink{tag-0123}{II.\ref*{cdp:eq:C.SN.10}}}
\expandafter\def\csname SRef@cdp:eq:C.MX.1\endcsname{\hyperlink{tag-00RR}{II.\ref*{cdp:eq:C.MX.1}}}
\expandafter\def\csname SRef@cdp:eq:C.MX.2\endcsname{\hyperlink{tag-00RT}{II.\ref*{cdp:eq:C.MX.2}}}
\expandafter\def\csname SRef@cdp:eq:C.MX.3\endcsname{\hyperlink{tag-00RU}{II.\ref*{cdp:eq:C.MX.3}}}
\expandafter\def\csname SRef@cdp:eq:C.MX.4\endcsname{\hyperlink{tag-00RV}{II.\ref*{cdp:eq:C.MX.4}}}
\expandafter\def\csname SRef@cdp:eq:C.MX.5\endcsname{\hyperlink{tag-00RW}{II.\ref*{cdp:eq:C.MX.5}}}
\expandafter\def\csname SRef@cdp:eq:C.MX.6\endcsname{\hyperlink{tag-00RX}{II.\ref*{cdp:eq:C.MX.6}}}
\expandafter\def\csname SRef@cdp:eq:C.MX.7\endcsname{\hyperlink{tag-00RY}{II.\ref*{cdp:eq:C.MX.7}}}
\expandafter\def\csname SRef@cdp:eq:C.MX.8\endcsname{\hyperlink{tag-00RZ}{II.\ref*{cdp:eq:C.MX.8}}}
\expandafter\def\csname SRef@cdp:eq:C.RF.6\endcsname{\hyperlink{tag-00YZ}{II.\ref*{cdp:eq:C.RF.6}}}
\expandafter\def\csname SRef@cdp:eq:C.MX.9\endcsname{\hyperlink{tag-00S0}{II.\ref*{cdp:eq:C.MX.9}}}
\expandafter\def\csname SRef@cdp:eq:C.MX.10\endcsname{\hyperlink{tag-00RS}{II.\ref*{cdp:eq:C.MX.10}}}
\expandafter\def\csname SRef@cdp:eq:C.WH.0\endcsname{\hyperlink{tag-016X}{II.\ref*{cdp:eq:C.WH.0}}}
\expandafter\def\csname SRef@cdp:eq:C.WH.1\endcsname{\hyperlink{tag-016Y}{II.\ref*{cdp:eq:C.WH.1}}}
\expandafter\def\csname SRef@cdp:eq:C.WH.2\endcsname{\hyperlink{tag-0177}{II.\ref*{cdp:eq:C.WH.2}}}
\expandafter\def\csname SRef@cdp:eq:C.WH.3\endcsname{\hyperlink{tag-0178}{II.\ref*{cdp:eq:C.WH.3}}}
\expandafter\def\csname SRef@cdp:eq:C.WH.4\endcsname{\hyperlink{tag-0179}{II.\ref*{cdp:eq:C.WH.4}}}
\expandafter\def\csname SRef@cdp:eq:C.WH.5\endcsname{\hyperlink{tag-017A}{II.\ref*{cdp:eq:C.WH.5}}}
\expandafter\def\csname SRef@cdp:eq:C.WH.6\endcsname{\hyperlink{tag-017B}{II.\ref*{cdp:eq:C.WH.6}}}
\expandafter\def\csname SRef@cdp:eq:C.WH.7\endcsname{\hyperlink{tag-017C}{II.\ref*{cdp:eq:C.WH.7}}}
\expandafter\def\csname SRef@cdp:eq:C.RF.7\endcsname{\hyperlink{tag-00Z0}{II.\ref*{cdp:eq:C.RF.7}}}
\expandafter\def\csname SRef@cdp:eq:C.WH.8\endcsname{\hyperlink{tag-017D}{II.\ref*{cdp:eq:C.WH.8}}}
\expandafter\def\csname SRef@cdp:eq:C.WH.9\endcsname{\hyperlink{tag-017E}{II.\ref*{cdp:eq:C.WH.9}}}
\expandafter\def\csname SRef@cdp:eq:C.WH.10\endcsname{\hyperlink{tag-016Z}{II.\ref*{cdp:eq:C.WH.10}}}
\expandafter\def\csname SRef@cdp:eq:C.WH.11\endcsname{\hyperlink{tag-0170}{II.\ref*{cdp:eq:C.WH.11}}}
\expandafter\def\csname SRef@cdp:eq:C.WH.12\endcsname{\hyperlink{tag-0171}{II.\ref*{cdp:eq:C.WH.12}}}
\expandafter\def\csname SRef@cdp:eq:C.WH.13\endcsname{\hyperlink{tag-0172}{II.\ref*{cdp:eq:C.WH.13}}}
\expandafter\def\csname SRef@cdp:eq:C.WH.14\endcsname{\hyperlink{tag-0173}{II.\ref*{cdp:eq:C.WH.14}}}
\expandafter\def\csname SRef@cdp:eq:C.WH.15\endcsname{\hyperlink{tag-0174}{II.\ref*{cdp:eq:C.WH.15}}}
\expandafter\def\csname SRef@cdp:eq:C.WH.16\endcsname{\hyperlink{tag-0175}{II.\ref*{cdp:eq:C.WH.16}}}
\expandafter\def\csname SRef@cdp:eq:C.WH.17\endcsname{\hyperlink{tag-0176}{II.\ref*{cdp:eq:C.WH.17}}}
\expandafter\def\csname SRef@cdp:eq:C.RF.8\endcsname{\hyperlink{tag-00Z1}{II.\ref*{cdp:eq:C.RF.8}}}
\expandafter\def\csname SRef@cdp:eq:C.AS.1\endcsname{\hyperlink{tag-00AG}{II.\ref*{cdp:eq:C.AS.1}}}
\expandafter\def\csname SRef@cdp:eq:C.AS.2\endcsname{\hyperlink{tag-00AR}{II.\ref*{cdp:eq:C.AS.2}}}
\expandafter\def\csname SRef@cdp:eq:C.AS.3\endcsname{\hyperlink{tag-00AS}{II.\ref*{cdp:eq:C.AS.3}}}
\expandafter\def\csname SRef@cdp:eq:C.AS.4\endcsname{\hyperlink{tag-00AT}{II.\ref*{cdp:eq:C.AS.4}}}
\expandafter\def\csname SRef@cdp:eq:C.AS.5\endcsname{\hyperlink{tag-00AU}{II.\ref*{cdp:eq:C.AS.5}}}
\expandafter\def\csname SRef@cdp:eq:C.AS.6\endcsname{\hyperlink{tag-00AV}{II.\ref*{cdp:eq:C.AS.6}}}
\expandafter\def\csname SRef@cdp:eq:C.AS.7\endcsname{\hyperlink{tag-00AW}{II.\ref*{cdp:eq:C.AS.7}}}
\expandafter\def\csname SRef@cdp:eq:C.AS.8\endcsname{\hyperlink{tag-00AX}{II.\ref*{cdp:eq:C.AS.8}}}
\expandafter\def\csname SRef@cdp:eq:C.AS.9\endcsname{\hyperlink{tag-00AY}{II.\ref*{cdp:eq:C.AS.9}}}
\expandafter\def\csname SRef@cdp:eq:C.AS.10\endcsname{\hyperlink{tag-00AH}{II.\ref*{cdp:eq:C.AS.10}}}
\expandafter\def\csname SRef@cdp:eq:C.CQ.1\endcsname{\hyperlink{tag-00CJ}{II.\ref*{cdp:eq:C.CQ.1}}}
\expandafter\def\csname SRef@cdp:eq:C.AS.11\endcsname{\hyperlink{tag-00AI}{II.\ref*{cdp:eq:C.AS.11}}}
\expandafter\def\csname SRef@cdp:eq:C.AS.12\endcsname{\hyperlink{tag-00AJ}{II.\ref*{cdp:eq:C.AS.12}}}
\expandafter\def\csname SRef@cdp:eq:C.AS.13\endcsname{\hyperlink{tag-00AK}{II.\ref*{cdp:eq:C.AS.13}}}
\expandafter\def\csname SRef@cdp:eq:C.AS.14\endcsname{\hyperlink{tag-00AL}{II.\ref*{cdp:eq:C.AS.14}}}
\expandafter\def\csname SRef@cdp:eq:C.AS.15\endcsname{\hyperlink{tag-00AM}{II.\ref*{cdp:eq:C.AS.15}}}
\expandafter\def\csname SRef@cdp:eq:C.AS.16\endcsname{\hyperlink{tag-00AN}{II.\ref*{cdp:eq:C.AS.16}}}
\expandafter\def\csname SRef@cdp:eq:C.AS.17\endcsname{\hyperlink{tag-00AO}{II.\ref*{cdp:eq:C.AS.17}}}
\expandafter\def\csname SRef@cdp:eq:C.AS.18\endcsname{\hyperlink{tag-00AP}{II.\ref*{cdp:eq:C.AS.18}}}
\expandafter\def\csname SRef@cdp:eq:C.AS.19\endcsname{\hyperlink{tag-00AQ}{II.\ref*{cdp:eq:C.AS.19}}}
\expandafter\def\csname SRef@cdp:eq:C.HR.1\endcsname{\hyperlink{tag-00N0}{II.\ref*{cdp:eq:C.HR.1}}}
\expandafter\def\csname SRef@cdp:eq:C.CQ.2\endcsname{\hyperlink{tag-00CN}{II.\ref*{cdp:eq:C.CQ.2}}}
\expandafter\def\csname SRef@cdp:eq:C.HR.2\endcsname{\hyperlink{tag-00N1}{II.\ref*{cdp:eq:C.HR.2}}}
\expandafter\def\csname SRef@cdp:eq:C.HR.3\endcsname{\hyperlink{tag-00N2}{II.\ref*{cdp:eq:C.HR.3}}}
\expandafter\def\csname SRef@cdp:eq:C.HR.4\endcsname{\hyperlink{tag-00N3}{II.\ref*{cdp:eq:C.HR.4}}}
\expandafter\def\csname SRef@cdp:eq:C.HR.5\endcsname{\hyperlink{tag-00N4}{II.\ref*{cdp:eq:C.HR.5}}}
\expandafter\def\csname SRef@cdp:eq:C.HR.6\endcsname{\hyperlink{tag-00N5}{II.\ref*{cdp:eq:C.HR.6}}}
\expandafter\def\csname SRef@cdp:eq:C.HR.7\endcsname{\hyperlink{tag-00N6}{II.\ref*{cdp:eq:C.HR.7}}}
\expandafter\def\csname SRef@cdp:eq:C.HR.8\endcsname{\hyperlink{tag-00N7}{II.\ref*{cdp:eq:C.HR.8}}}
\expandafter\def\csname SRef@cdp:eq:C.TL.1\endcsname{\hyperlink{tag-013Q}{II.\ref*{cdp:eq:C.TL.1}}}
\expandafter\def\csname SRef@cdp:eq:C.TL.2\endcsname{\hyperlink{tag-013R}{II.\ref*{cdp:eq:C.TL.2}}}
\expandafter\def\csname SRef@cdp:eq:C.TL.3\endcsname{\hyperlink{tag-013S}{II.\ref*{cdp:eq:C.TL.3}}}
\expandafter\def\csname SRef@cdp:eq:C.CQ.3\endcsname{\hyperlink{tag-00CO}{II.\ref*{cdp:eq:C.CQ.3}}}
\expandafter\def\csname SRef@cdp:eq:C.TL.4\endcsname{\hyperlink{tag-013T}{II.\ref*{cdp:eq:C.TL.4}}}
\expandafter\def\csname SRef@cdp:eq:C.TL.5\endcsname{\hyperlink{tag-013U}{II.\ref*{cdp:eq:C.TL.5}}}
\expandafter\def\csname SRef@cdp:eq:C.TL.6\endcsname{\hyperlink{tag-013V}{II.\ref*{cdp:eq:C.TL.6}}}
\expandafter\def\csname SRef@cdp:eq:C.TL.7\endcsname{\hyperlink{tag-013W}{II.\ref*{cdp:eq:C.TL.7}}}
\expandafter\def\csname SRef@cdp:eq:C.TL.8\endcsname{\hyperlink{tag-013X}{II.\ref*{cdp:eq:C.TL.8}}}
\expandafter\def\csname SRef@cdp:eq:C.WM.1\endcsname{\hyperlink{tag-017F}{II.\ref*{cdp:eq:C.WM.1}}}
\expandafter\def\csname SRef@cdp:eq:C.WM.2\endcsname{\hyperlink{tag-017G}{II.\ref*{cdp:eq:C.WM.2}}}
\expandafter\def\csname SRef@cdp:eq:C.WM.3\endcsname{\hyperlink{tag-017H}{II.\ref*{cdp:eq:C.WM.3}}}
\expandafter\def\csname SRef@cdp:eq:C.WM.4\endcsname{\hyperlink{tag-017I}{II.\ref*{cdp:eq:C.WM.4}}}
\expandafter\def\csname SRef@cdp:eq:C.WM.5\endcsname{\hyperlink{tag-017J}{II.\ref*{cdp:eq:C.WM.5}}}
\expandafter\def\csname SRef@cdp:eq:C.CQ.4\endcsname{\hyperlink{tag-00CP}{II.\ref*{cdp:eq:C.CQ.4}}}
\expandafter\def\csname SRef@cdp:eq:C.WM.6\endcsname{\hyperlink{tag-017K}{II.\ref*{cdp:eq:C.WM.6}}}
\expandafter\def\csname SRef@cdp:eq:C.WM.7\endcsname{\hyperlink{tag-017L}{II.\ref*{cdp:eq:C.WM.7}}}
\expandafter\def\csname SRef@cdp:eq:C.WM.8\endcsname{\hyperlink{tag-017M}{II.\ref*{cdp:eq:C.WM.8}}}
\expandafter\def\csname SRef@cdp:eq:C.WM.9\endcsname{\hyperlink{tag-017N}{II.\ref*{cdp:eq:C.WM.9}}}
\expandafter\def\csname SRef@cdp:eq:C.AD.1\endcsname{\hyperlink{tag-009B}{II.\ref*{cdp:eq:C.AD.1}}}
\expandafter\def\csname SRef@cdp:eq:C.AD.2\endcsname{\hyperlink{tag-009E}{II.\ref*{cdp:eq:C.AD.2}}}
\expandafter\def\csname SRef@cdp:eq:C.AD.3\endcsname{\hyperlink{tag-009F}{II.\ref*{cdp:eq:C.AD.3}}}
\expandafter\def\csname SRef@cdp:eq:C.AD.4\endcsname{\hyperlink{tag-009G}{II.\ref*{cdp:eq:C.AD.4}}}
\expandafter\def\csname SRef@cdp:eq:C.AD.5\endcsname{\hyperlink{tag-009H}{II.\ref*{cdp:eq:C.AD.5}}}
\expandafter\def\csname SRef@cdp:eq:C.AD.6\endcsname{\hyperlink{tag-009I}{II.\ref*{cdp:eq:C.AD.6}}}
\expandafter\def\csname SRef@cdp:eq:C.CQ.5\endcsname{\hyperlink{tag-00CQ}{II.\ref*{cdp:eq:C.CQ.5}}}
\expandafter\def\csname SRef@cdp:eq:C.AD.7\endcsname{\hyperlink{tag-009J}{II.\ref*{cdp:eq:C.AD.7}}}
\expandafter\def\csname SRef@cdp:eq:C.AD.8\endcsname{\hyperlink{tag-009K}{II.\ref*{cdp:eq:C.AD.8}}}
\expandafter\def\csname SRef@cdp:eq:C.AD.9\endcsname{\hyperlink{tag-009L}{II.\ref*{cdp:eq:C.AD.9}}}
\expandafter\def\csname SRef@cdp:eq:C.AD.10\endcsname{\hyperlink{tag-009C}{II.\ref*{cdp:eq:C.AD.10}}}
\expandafter\def\csname SRef@cdp:eq:C.AD.11\endcsname{\hyperlink{tag-009D}{II.\ref*{cdp:eq:C.AD.11}}}
\expandafter\def\csname SRef@cdp:eq:C.HT.1\endcsname{\hyperlink{tag-00N8}{II.\ref*{cdp:eq:C.HT.1}}}
\expandafter\def\csname SRef@cdp:eq:C.HT.2\endcsname{\hyperlink{tag-00NC}{II.\ref*{cdp:eq:C.HT.2}}}
\expandafter\def\csname SRef@cdp:eq:C.HT.3\endcsname{\hyperlink{tag-00ND}{II.\ref*{cdp:eq:C.HT.3}}}
\expandafter\def\csname SRef@cdp:eq:C.HT.4\endcsname{\hyperlink{tag-00NE}{II.\ref*{cdp:eq:C.HT.4}}}
\expandafter\def\csname SRef@cdp:eq:C.HT.5\endcsname{\hyperlink{tag-00NF}{II.\ref*{cdp:eq:C.HT.5}}}
\expandafter\def\csname SRef@cdp:eq:C.CQ.6\endcsname{\hyperlink{tag-00CR}{II.\ref*{cdp:eq:C.CQ.6}}}
\expandafter\def\csname SRef@cdp:eq:C.HT.6\endcsname{\hyperlink{tag-00NG}{II.\ref*{cdp:eq:C.HT.6}}}
\expandafter\def\csname SRef@cdp:eq:C.HT.7\endcsname{\hyperlink{tag-00NH}{II.\ref*{cdp:eq:C.HT.7}}}
\expandafter\def\csname SRef@cdp:eq:C.HT.8\endcsname{\hyperlink{tag-00NI}{II.\ref*{cdp:eq:C.HT.8}}}
\expandafter\def\csname SRef@cdp:eq:C.HT.9\endcsname{\hyperlink{tag-00NJ}{II.\ref*{cdp:eq:C.HT.9}}}
\expandafter\def\csname SRef@cdp:eq:C.HT.10\endcsname{\hyperlink{tag-00N9}{II.\ref*{cdp:eq:C.HT.10}}}
\expandafter\def\csname SRef@cdp:eq:C.HT.11\endcsname{\hyperlink{tag-00NA}{II.\ref*{cdp:eq:C.HT.11}}}
\expandafter\def\csname SRef@cdp:eq:C.HT.12\endcsname{\hyperlink{tag-00NB}{II.\ref*{cdp:eq:C.HT.12}}}
\expandafter\def\csname SRef@cdp:eq:C.EW.1\endcsname{\hyperlink{tag-00JH}{II.\ref*{cdp:eq:C.EW.1}}}
\expandafter\def\csname SRef@cdp:eq:C.EW.2\endcsname{\hyperlink{tag-00JS}{II.\ref*{cdp:eq:C.EW.2}}}
\expandafter\def\csname SRef@cdp:eq:C.EW.3\endcsname{\hyperlink{tag-00JV}{II.\ref*{cdp:eq:C.EW.3}}}
\expandafter\def\csname SRef@cdp:eq:C.CQ.7\endcsname{\hyperlink{tag-00CS}{II.\ref*{cdp:eq:C.CQ.7}}}
\expandafter\def\csname SRef@cdp:eq:C.EW.4\endcsname{\hyperlink{tag-00JW}{II.\ref*{cdp:eq:C.EW.4}}}
\expandafter\def\csname SRef@cdp:eq:C.EW.5\endcsname{\hyperlink{tag-00JX}{II.\ref*{cdp:eq:C.EW.5}}}
\expandafter\def\csname SRef@cdp:eq:C.EW.6\endcsname{\hyperlink{tag-00JY}{II.\ref*{cdp:eq:C.EW.6}}}
\expandafter\def\csname SRef@cdp:eq:C.EW.7\endcsname{\hyperlink{tag-00JZ}{II.\ref*{cdp:eq:C.EW.7}}}
\expandafter\def\csname SRef@cdp:eq:C.EW.8\endcsname{\hyperlink{tag-00K0}{II.\ref*{cdp:eq:C.EW.8}}}
\expandafter\def\csname SRef@cdp:eq:C.EW.9\endcsname{\hyperlink{tag-00K1}{II.\ref*{cdp:eq:C.EW.9}}}
\expandafter\def\csname SRef@cdp:eq:C.EW.10\endcsname{\hyperlink{tag-00JI}{II.\ref*{cdp:eq:C.EW.10}}}
\expandafter\def\csname SRef@cdp:eq:C.EW.11\endcsname{\hyperlink{tag-00JJ}{II.\ref*{cdp:eq:C.EW.11}}}
\expandafter\def\csname SRef@cdp:eq:C.EW.12\endcsname{\hyperlink{tag-00JK}{II.\ref*{cdp:eq:C.EW.12}}}
\expandafter\def\csname SRef@cdp:eq:C.EW.13\endcsname{\hyperlink{tag-00JL}{II.\ref*{cdp:eq:C.EW.13}}}
\expandafter\def\csname SRef@cdp:eq:C.CQ.8\endcsname{\hyperlink{tag-00CT}{II.\ref*{cdp:eq:C.CQ.8}}}
\expandafter\def\csname SRef@cdp:eq:C.EW.14\endcsname{\hyperlink{tag-00JM}{II.\ref*{cdp:eq:C.EW.14}}}
\expandafter\def\csname SRef@cdp:eq:C.EW.15\endcsname{\hyperlink{tag-00JN}{II.\ref*{cdp:eq:C.EW.15}}}
\expandafter\def\csname SRef@cdp:eq:C.EW.16\endcsname{\hyperlink{tag-00JO}{II.\ref*{cdp:eq:C.EW.16}}}
\expandafter\def\csname SRef@cdp:eq:C.EW.17\endcsname{\hyperlink{tag-00JP}{II.\ref*{cdp:eq:C.EW.17}}}
\expandafter\def\csname SRef@cdp:eq:C.EW.18\endcsname{\hyperlink{tag-00JQ}{II.\ref*{cdp:eq:C.EW.18}}}
\expandafter\def\csname SRef@cdp:eq:C.EW.19\endcsname{\hyperlink{tag-00JR}{II.\ref*{cdp:eq:C.EW.19}}}
\expandafter\def\csname SRef@cdp:eq:C.EW.20\endcsname{\hyperlink{tag-00JT}{II.\ref*{cdp:eq:C.EW.20}}}
\expandafter\def\csname SRef@cdp:eq:C.EW.21\endcsname{\hyperlink{tag-00JU}{II.\ref*{cdp:eq:C.EW.21}}}
\expandafter\def\csname SRef@cdp:eq:C.PR.1\endcsname{\hyperlink{tag-00VH}{II.\ref*{cdp:eq:C.PR.1}}}
\expandafter\def\csname SRef@cdp:eq:C.PR.2\endcsname{\hyperlink{tag-00VR}{II.\ref*{cdp:eq:C.PR.2}}}
\expandafter\def\csname SRef@cdp:eq:C.CQ.9\endcsname{\hyperlink{tag-00CU}{II.\ref*{cdp:eq:C.CQ.9}}}
\expandafter\def\csname SRef@cdp:eq:C.PR.3\endcsname{\hyperlink{tag-00VS}{II.\ref*{cdp:eq:C.PR.3}}}
\expandafter\def\csname SRef@cdp:eq:C.PR.4\endcsname{\hyperlink{tag-00VT}{II.\ref*{cdp:eq:C.PR.4}}}
\expandafter\def\csname SRef@cdp:eq:C.PR.5\endcsname{\hyperlink{tag-00VU}{II.\ref*{cdp:eq:C.PR.5}}}
\expandafter\def\csname SRef@cdp:eq:C.PR.6\endcsname{\hyperlink{tag-00VV}{II.\ref*{cdp:eq:C.PR.6}}}
\expandafter\def\csname SRef@cdp:eq:C.PR.7\endcsname{\hyperlink{tag-00VW}{II.\ref*{cdp:eq:C.PR.7}}}
\expandafter\def\csname SRef@cdp:eq:C.PR.8\endcsname{\hyperlink{tag-00VX}{II.\ref*{cdp:eq:C.PR.8}}}
\expandafter\def\csname SRef@cdp:eq:C.PR.9\endcsname{\hyperlink{tag-00VY}{II.\ref*{cdp:eq:C.PR.9}}}
\expandafter\def\csname SRef@cdp:eq:C.PR.10\endcsname{\hyperlink{tag-00VI}{II.\ref*{cdp:eq:C.PR.10}}}
\expandafter\def\csname SRef@cdp:eq:C.PR.11\endcsname{\hyperlink{tag-00VJ}{II.\ref*{cdp:eq:C.PR.11}}}
\expandafter\def\csname SRef@cdp:eq:C.PR.12\endcsname{\hyperlink{tag-00VK}{II.\ref*{cdp:eq:C.PR.12}}}
\expandafter\def\csname SRef@cdp:eq:C.CQ.10\endcsname{\hyperlink{tag-00CK}{II.\ref*{cdp:eq:C.CQ.10}}}
\expandafter\def\csname SRef@cdp:eq:C.PR.13\endcsname{\hyperlink{tag-00VL}{II.\ref*{cdp:eq:C.PR.13}}}
\expandafter\def\csname SRef@cdp:eq:C.PR.14\endcsname{\hyperlink{tag-00VM}{II.\ref*{cdp:eq:C.PR.14}}}
\expandafter\def\csname SRef@cdp:eq:C.PR.15\endcsname{\hyperlink{tag-00VN}{II.\ref*{cdp:eq:C.PR.15}}}
\expandafter\def\csname SRef@cdp:eq:C.PR.16\endcsname{\hyperlink{tag-00VO}{II.\ref*{cdp:eq:C.PR.16}}}
\expandafter\def\csname SRef@cdp:eq:C.PR.17\endcsname{\hyperlink{tag-00VP}{II.\ref*{cdp:eq:C.PR.17}}}
\expandafter\def\csname SRef@cdp:eq:C.PR.18\endcsname{\hyperlink{tag-00VQ}{II.\ref*{cdp:eq:C.PR.18}}}
\expandafter\def\csname SRef@cdp:eq:C.PT.1\endcsname{\hyperlink{tag-00VZ}{II.\ref*{cdp:eq:C.PT.1}}}
\expandafter\def\csname SRef@cdp:eq:C.PT.2\endcsname{\hyperlink{tag-00W5}{II.\ref*{cdp:eq:C.PT.2}}}
\expandafter\def\csname SRef@cdp:eq:C.PT.3\endcsname{\hyperlink{tag-00W6}{II.\ref*{cdp:eq:C.PT.3}}}
\expandafter\def\csname SRef@cdp:eq:C.PT.4\endcsname{\hyperlink{tag-00W7}{II.\ref*{cdp:eq:C.PT.4}}}
\expandafter\def\csname SRef@cdp:eq:C.CQ.11\endcsname{\hyperlink{tag-00CL}{II.\ref*{cdp:eq:C.CQ.11}}}
\expandafter\def\csname SRef@cdp:eq:C.PT.5\endcsname{\hyperlink{tag-00W8}{II.\ref*{cdp:eq:C.PT.5}}}
\expandafter\def\csname SRef@cdp:eq:C.PT.6\endcsname{\hyperlink{tag-00W9}{II.\ref*{cdp:eq:C.PT.6}}}
\expandafter\def\csname SRef@cdp:eq:C.PT.7\endcsname{\hyperlink{tag-00WA}{II.\ref*{cdp:eq:C.PT.7}}}
\expandafter\def\csname SRef@cdp:eq:C.PT.8\endcsname{\hyperlink{tag-00WB}{II.\ref*{cdp:eq:C.PT.8}}}
\expandafter\def\csname SRef@cdp:eq:C.PT.9\endcsname{\hyperlink{tag-00WC}{II.\ref*{cdp:eq:C.PT.9}}}
\expandafter\def\csname SRef@cdp:eq:C.PT.10\endcsname{\hyperlink{tag-00W0}{II.\ref*{cdp:eq:C.PT.10}}}
\expandafter\def\csname SRef@cdp:eq:C.PT.11\endcsname{\hyperlink{tag-00W1}{II.\ref*{cdp:eq:C.PT.11}}}
\expandafter\def\csname SRef@cdp:eq:C.PT.12\endcsname{\hyperlink{tag-00W2}{II.\ref*{cdp:eq:C.PT.12}}}
\expandafter\def\csname SRef@cdp:eq:C.PT.13\endcsname{\hyperlink{tag-00W3}{II.\ref*{cdp:eq:C.PT.13}}}
\expandafter\def\csname SRef@cdp:eq:C.PT.14\endcsname{\hyperlink{tag-00W4}{II.\ref*{cdp:eq:C.PT.14}}}
\expandafter\def\csname SRef@cdp:eq:C.CQ.12\endcsname{\hyperlink{tag-00CM}{II.\ref*{cdp:eq:C.CQ.12}}}
\expandafter\def\csname SRef@cdp:eq:C.QT.1\endcsname{\hyperlink{tag-00WD}{II.\ref*{cdp:eq:C.QT.1}}}
\expandafter\def\csname SRef@cdp:eq:C.QT.2\endcsname{\hyperlink{tag-00WI}{II.\ref*{cdp:eq:C.QT.2}}}
\expandafter\def\csname SRef@cdp:eq:C.QT.3\endcsname{\hyperlink{tag-00WJ}{II.\ref*{cdp:eq:C.QT.3}}}
\expandafter\def\csname SRef@cdp:eq:C.QT.4\endcsname{\hyperlink{tag-00WK}{II.\ref*{cdp:eq:C.QT.4}}}
\expandafter\def\csname SRef@cdp:eq:C.QT.5\endcsname{\hyperlink{tag-00WL}{II.\ref*{cdp:eq:C.QT.5}}}
\expandafter\def\csname SRef@cdp:eq:C.QT.6\endcsname{\hyperlink{tag-00WM}{II.\ref*{cdp:eq:C.QT.6}}}
\expandafter\def\csname SRef@cdp:eq:C.QT.7\endcsname{\hyperlink{tag-00WN}{II.\ref*{cdp:eq:C.QT.7}}}
\expandafter\def\csname SRef@cdp:eq:C.QT.8\endcsname{\hyperlink{tag-00WO}{II.\ref*{cdp:eq:C.QT.8}}}
\expandafter\def\csname SRef@cdp:eq:C.QT.9\endcsname{\hyperlink{tag-00WP}{II.\ref*{cdp:eq:C.QT.9}}}
\expandafter\def\csname SRef@cdp:eq:C.QT.10\endcsname{\hyperlink{tag-00WE}{II.\ref*{cdp:eq:C.QT.10}}}
\expandafter\def\csname SRef@cdp:st:C.2.UN.1\endcsname{\hyperlink{tag-01CM}{II.\ref*{cdp:st:C.2.UN.1}}}
\expandafter\def\csname SRef@cdp:st:C.2.UN.20\endcsname{\hyperlink{tag-01CY}{II.\ref*{cdp:st:C.2.UN.20}}}
\expandafter\def\csname SRef@cdp:st:C.2.UN.22\endcsname{\hyperlink{tag-01D0}{II.\ref*{cdp:st:C.2.UN.22}}}
\expandafter\def\csname SRef@cdp:st:C.2.UN.23\endcsname{\hyperlink{tag-01D1}{II.\ref*{cdp:st:C.2.UN.23}}}
\expandafter\def\csname SRef@cdp:st:C.2.UN.26\endcsname{\hyperlink{tag-01D4}{II.\ref*{cdp:st:C.2.UN.26}}}
\expandafter\def\csname SRef@cdp:st:C.2.UN.28\endcsname{\hyperlink{tag-01D6}{II.\ref*{cdp:st:C.2.UN.28}}}
\expandafter\def\csname SRef@cdp:st:C.2.UN.29\endcsname{\hyperlink{tag-01D7}{II.\ref*{cdp:st:C.2.UN.29}}}
\expandafter\def\csname SRef@cdp:st:C.2.UN.31\endcsname{\hyperlink{tag-01DA}{II.\ref*{cdp:st:C.2.UN.31}}}
\expandafter\def\csname SRef@cdp:st:C.2.UN.33\endcsname{\hyperlink{tag-01DC}{II.\ref*{cdp:st:C.2.UN.33}}}
\expandafter\def\csname SRef@cdp:st:C.2.UN.36\endcsname{\hyperlink{tag-01DF}{II.\ref*{cdp:st:C.2.UN.36}}}
\expandafter\def\csname SRef@cdp:st:C.2.UN.38\endcsname{\hyperlink{tag-01DH}{II.\ref*{cdp:st:C.2.UN.38}}}
\expandafter\def\csname SRef@cdp:st:C.2.UN.2\endcsname{\hyperlink{tag-01CX}{II.\ref*{cdp:st:C.2.UN.2}}}
\expandafter\def\csname SRef@cdp:st:C.2.UN.39\endcsname{\hyperlink{tag-01DI}{II.\ref*{cdp:st:C.2.UN.39}}}
\expandafter\def\csname SRef@cdp:st:C.2.UN.44\endcsname{\hyperlink{tag-01DO}{II.\ref*{cdp:st:C.2.UN.44}}}
\expandafter\def\csname SRef@cdp:st:C.2.UN.47\endcsname{\hyperlink{tag-01DR}{II.\ref*{cdp:st:C.2.UN.47}}}
\expandafter\def\csname SRef@cdp:st:C.2.UN.3\endcsname{\hyperlink{tag-01D8}{II.\ref*{cdp:st:C.2.UN.3}}}
\expandafter\def\csname SRef@cdp:st:C.2.UN.4\endcsname{\hyperlink{tag-01DJ}{II.\ref*{cdp:st:C.2.UN.4}}}
\expandafter\def\csname SRef@cdp:st:C.2.UN.14\endcsname{\hyperlink{tag-01CR}{II.\ref*{cdp:st:C.2.UN.14}}}
\expandafter\def\csname SRef@cdp:st:C.2.UN.15\endcsname{\hyperlink{tag-01CS}{II.\ref*{cdp:st:C.2.UN.15}}}
\expandafter\def\csname SRef@cdp:st:C.2.UN.16\endcsname{\hyperlink{tag-01CT}{II.\ref*{cdp:st:C.2.UN.16}}}
\expandafter\def\csname SRef@cdp:st:C.2.UN.17\endcsname{\hyperlink{tag-01CU}{II.\ref*{cdp:st:C.2.UN.17}}}
\expandafter\def\csname SRef@cdp:st:C.2.UN.19\endcsname{\hyperlink{tag-01CW}{II.\ref*{cdp:st:C.2.UN.19}}}
\expandafter\def\csname SRef@cdp:st:C.2.UN.5\endcsname{\hyperlink{tag-01DU}{II.\ref*{cdp:st:C.2.UN.5}}}
\expandafter\def\csname SRef@cdp:st:C.2.IRR.1\endcsname{\hyperlink{tag-01C3}{II.\ref*{cdp:st:C.2.IRR.1}}}
\expandafter\def\csname SRef@cdp:st:C.2.VII.1\endcsname{\hyperlink{tag-01F0}{II.\ref*{cdp:st:C.2.VII.1}}}
\expandafter\def\csname SRef@cdp:st:C.2.VII.2\endcsname{\hyperlink{tag-01F1}{II.\ref*{cdp:st:C.2.VII.2}}}
\expandafter\def\csname SRef@cdp:st:C.2.VII.3\endcsname{\hyperlink{tag-01F2}{II.\ref*{cdp:st:C.2.VII.3}}}
\expandafter\def\csname SRef@cdp:st:C.2.VII.4\endcsname{\hyperlink{tag-01F3}{II.\ref*{cdp:st:C.2.VII.4}}}
\expandafter\def\csname SRef@cdp:st:C.2.RAT.1\endcsname{\hyperlink{tag-01CD}{II.\ref*{cdp:st:C.2.RAT.1}}}
\expandafter\def\csname SRef@cdp:st:C.2.RAT.2\endcsname{\hyperlink{tag-01CE}{II.\ref*{cdp:st:C.2.RAT.2}}}
\expandafter\def\csname SRef@cdp:st:C.2.RAT.3\endcsname{\hyperlink{tag-01CF}{II.\ref*{cdp:st:C.2.RAT.3}}}
\expandafter\def\csname SRef@cdp:st:C.2.RAT.4\endcsname{\hyperlink{tag-01CG}{II.\ref*{cdp:st:C.2.RAT.4}}}
\expandafter\def\csname SRef@cdp:st:C.2.RAT.5\endcsname{\hyperlink{tag-01CH}{II.\ref*{cdp:st:C.2.RAT.5}}}
\expandafter\def\csname SRef@cdp:st:C.2.CR.1\endcsname{\hyperlink{tag-01BQ}{II.\ref*{cdp:st:C.2.CR.1}}}
\expandafter\def\csname SRef@cdp:st:C.2.RAT.7\endcsname{\hyperlink{tag-01CJ}{II.\ref*{cdp:st:C.2.RAT.7}}}
\expandafter\def\csname SRef@cdp:st:C.2.RAT.8\endcsname{\hyperlink{tag-01CK}{II.\ref*{cdp:st:C.2.RAT.8}}}
\expandafter\def\csname SRef@cdp:st:C.2.RAT.9\endcsname{\hyperlink{tag-01CL}{II.\ref*{cdp:st:C.2.RAT.9}}}
\expandafter\def\csname SRef@cdp:st:C.2.D4.1\endcsname{\hyperlink{tag-01BW}{II.\ref*{cdp:st:C.2.D4.1}}}
\expandafter\def\csname SRef@cdp:st:C.2.D4.2\endcsname{\hyperlink{tag-01BX}{II.\ref*{cdp:st:C.2.D4.2}}}
\expandafter\def\csname SRef@cdp:st:C.2.D4.3\endcsname{\hyperlink{tag-01BY}{II.\ref*{cdp:st:C.2.D4.3}}}
\expandafter\def\csname SRef@cdp:st:C.2.UN.7\endcsname{\hyperlink{tag-01EG}{II.\ref*{cdp:st:C.2.UN.7}}}
\expandafter\def\csname SRef@cdp:st:C.2.UN.12\endcsname{\hyperlink{tag-01CP}{II.\ref*{cdp:st:C.2.UN.12}}}
\expandafter\def\csname SRef@cdp:st:C.2.INT.1\endcsname{\hyperlink{tag-01C0}{II.\ref*{cdp:st:C.2.INT.1}}}
\expandafter\def\csname SRef@cdp:st:C.2.INT.2\endcsname{\hyperlink{tag-01C1}{II.\ref*{cdp:st:C.2.INT.2}}}
\expandafter\def\csname SRef@cdp:st:C.2.CR.2\endcsname{\hyperlink{tag-01BR}{II.\ref*{cdp:st:C.2.CR.2}}}
\expandafter\def\csname SRef@cdp:st:C.2.INT.3\endcsname{\hyperlink{tag-01C2}{II.\ref*{cdp:st:C.2.INT.3}}}
\expandafter\def\csname SRef@cdp:st:C.2.MH.1\endcsname{\hyperlink{tag-01C4}{II.\ref*{cdp:st:C.2.MH.1}}}
\expandafter\def\csname SRef@cdp:st:C.2.MH.2\endcsname{\hyperlink{tag-01C5}{II.\ref*{cdp:st:C.2.MH.2}}}
\expandafter\def\csname SRef@cdp:st:C.2.MH.3\endcsname{\hyperlink{tag-01C6}{II.\ref*{cdp:st:C.2.MH.3}}}
\expandafter\def\csname SRef@cdp:st:C.2.MH.4\endcsname{\hyperlink{tag-01C7}{II.\ref*{cdp:st:C.2.MH.4}}}
\expandafter\def\csname SRef@cdp:st:C.2.CR.3\endcsname{\hyperlink{tag-01BS}{II.\ref*{cdp:st:C.2.CR.3}}}
\expandafter\def\csname SRef@cdp:st:C.2.CR.4\endcsname{\hyperlink{tag-01BT}{II.\ref*{cdp:st:C.2.CR.4}}}
\expandafter\def\csname SRef@cdp:st:C.2.CR.5\endcsname{\hyperlink{tag-01BU}{II.\ref*{cdp:st:C.2.CR.5}}}
\expandafter\def\csname SRef@cdp:st:C.2.CR.6\endcsname{\hyperlink{tag-01BV}{II.\ref*{cdp:st:C.2.CR.6}}}
\expandafter\def\csname SRef@cdp:st:C.2.UN.9\endcsname{\hyperlink{tag-01EZ}{II.\ref*{cdp:st:C.2.UN.9}}}
\expandafter\def\csname SRef@cdp:st:C.2.UN.10\endcsname{\hyperlink{tag-01CN}{II.\ref*{cdp:st:C.2.UN.10}}}
\expandafter\def\csname SRef@cdp:st:C2a\endcsname{\hyperlink{tag-01FG}{II.\ref*{cdp:st:C2a}}}
\expandafter\def\csname SRef@cdp:st:C.2.UN.67\endcsname{\hyperlink{tag-01ED}{II.\ref*{cdp:st:C.2.UN.67}}}
\expandafter\def\csname SRef@cdp:st:C.2.UN.68\endcsname{\hyperlink{tag-01EE}{II.\ref*{cdp:st:C.2.UN.68}}}
\expandafter\def\csname SRef@cdp:st:C2b\endcsname{\hyperlink{tag-01FH}{II.\ref*{cdp:st:C2b}}}
\expandafter\def\csname SRef@cdp:st:C2c\endcsname{\hyperlink{tag-01FI}{II.\ref*{cdp:st:C2c}}}
\expandafter\def\csname SRef@cdp:st:C.2.UN.11\endcsname{\hyperlink{tag-01CO}{II.\ref*{cdp:st:C.2.UN.11}}}
\expandafter\def\csname SRef@cdp:st:C.2.UN.24\endcsname{\hyperlink{tag-01D2}{II.\ref*{cdp:st:C.2.UN.24}}}
\expandafter\def\csname SRef@cdp:st:C.2.UN.32\endcsname{\hyperlink{tag-01DB}{II.\ref*{cdp:st:C.2.UN.32}}}
\expandafter\def\csname SRef@cdp:st:C.2.UN.64\endcsname{\hyperlink{tag-01EA}{II.\ref*{cdp:st:C.2.UN.64}}}
\expandafter\def\csname SRef@cdp:st:C.2.UN.65\endcsname{\hyperlink{tag-01EB}{II.\ref*{cdp:st:C.2.UN.65}}}
\expandafter\def\csname SRef@cdp:st:C.2.UN.66\endcsname{\hyperlink{tag-01EC}{II.\ref*{cdp:st:C.2.UN.66}}}
\expandafter\def\csname SRef@cdp:st:C.2.UN.6\endcsname{\hyperlink{tag-01E5}{II.\ref*{cdp:st:C.2.UN.6}}}
\expandafter\def\csname SRef@cdp:st:C.2.UN.35\endcsname{\hyperlink{tag-01DE}{II.\ref*{cdp:st:C.2.UN.35}}}
\expandafter\def\csname SRef@cdp:st:C.2.UN.37\endcsname{\hyperlink{tag-01DG}{II.\ref*{cdp:st:C.2.UN.37}}}
\expandafter\def\csname SRef@cdp:st:C.2.UN.41\endcsname{\hyperlink{tag-01DL}{II.\ref*{cdp:st:C.2.UN.41}}}
\expandafter\def\csname SRef@cdp:st:C.2.UN.42\endcsname{\hyperlink{tag-01DM}{II.\ref*{cdp:st:C.2.UN.42}}}
\expandafter\def\csname SRef@cdp:st:C.2.UN.43\endcsname{\hyperlink{tag-01DN}{II.\ref*{cdp:st:C.2.UN.43}}}
\expandafter\def\csname SRef@cdp:st:C.2.UN.45\endcsname{\hyperlink{tag-01DP}{II.\ref*{cdp:st:C.2.UN.45}}}
\expandafter\def\csname SRef@cdp:st:C.2.UN.8\endcsname{\hyperlink{tag-01ER}{II.\ref*{cdp:st:C.2.UN.8}}}
\expandafter\def\csname SRef@cdp:st:C.2.UN.13\endcsname{\hyperlink{tag-01CQ}{II.\ref*{cdp:st:C.2.UN.13}}}
\expandafter\def\csname SRef@cdp:st:C.2.UN.18\endcsname{\hyperlink{tag-01CV}{II.\ref*{cdp:st:C.2.UN.18}}}
\expandafter\def\csname SRef@cdp:st:C.2.UN.21\endcsname{\hyperlink{tag-01CZ}{II.\ref*{cdp:st:C.2.UN.21}}}
\expandafter\def\csname SRef@cdp:st:C.2.UN.25\endcsname{\hyperlink{tag-01D3}{II.\ref*{cdp:st:C.2.UN.25}}}
\expandafter\def\csname SRef@cdp:st:C.2.UN.27\endcsname{\hyperlink{tag-01D5}{II.\ref*{cdp:st:C.2.UN.27}}}
\expandafter\def\csname SRef@cdp:st:C.2.UN.30\endcsname{\hyperlink{tag-01D9}{II.\ref*{cdp:st:C.2.UN.30}}}
\expandafter\def\csname SRef@cdp:st:C.2.UN.34\endcsname{\hyperlink{tag-01DD}{II.\ref*{cdp:st:C.2.UN.34}}}
\expandafter\def\csname SRef@cdp:st:C.2.UN.40\endcsname{\hyperlink{tag-01DK}{II.\ref*{cdp:st:C.2.UN.40}}}
\expandafter\def\csname SRef@cdp:st:C.2.UN.55\endcsname{\hyperlink{tag-01E0}{II.\ref*{cdp:st:C.2.UN.55}}}
\expandafter\def\csname SRef@cdp:st:C.2.UN.56\endcsname{\hyperlink{tag-01E1}{II.\ref*{cdp:st:C.2.UN.56}}}
\expandafter\def\csname SRef@cdp:st:C.2.UN.57\endcsname{\hyperlink{tag-01E2}{II.\ref*{cdp:st:C.2.UN.57}}}
\expandafter\def\csname SRef@cdp:st:C.2.UN.58\endcsname{\hyperlink{tag-01E3}{II.\ref*{cdp:st:C.2.UN.58}}}
\expandafter\def\csname SRef@cdp:st:C.2.UN.59\endcsname{\hyperlink{tag-01E4}{II.\ref*{cdp:st:C.2.UN.59}}}
\expandafter\def\csname SRef@cdp:st:C.2.UN.60\endcsname{\hyperlink{tag-01E6}{II.\ref*{cdp:st:C.2.UN.60}}}
\expandafter\def\csname SRef@cdp:st:C.2.UN.61\endcsname{\hyperlink{tag-01E7}{II.\ref*{cdp:st:C.2.UN.61}}}
\expandafter\def\csname SRef@cdp:st:C.2.UN.62\endcsname{\hyperlink{tag-01E8}{II.\ref*{cdp:st:C.2.UN.62}}}
\expandafter\def\csname SRef@cdp:st:C.2.UN.63\endcsname{\hyperlink{tag-01E9}{II.\ref*{cdp:st:C.2.UN.63}}}
\expandafter\def\csname SRef@cdp:st:C.2.UN.46\endcsname{\hyperlink{tag-01DQ}{II.\ref*{cdp:st:C.2.UN.46}}}
\expandafter\def\csname SRef@cdp:st:C.2.UN.48\endcsname{\hyperlink{tag-01DS}{II.\ref*{cdp:st:C.2.UN.48}}}
\expandafter\def\csname SRef@cdp:st:C.2.UN.49\endcsname{\hyperlink{tag-01DT}{II.\ref*{cdp:st:C.2.UN.49}}}
\expandafter\def\csname SRef@cdp:st:C.2.UN.50\endcsname{\hyperlink{tag-01DV}{II.\ref*{cdp:st:C.2.UN.50}}}
\expandafter\def\csname SRef@cdp:st:C.2.UN.51\endcsname{\hyperlink{tag-01DW}{II.\ref*{cdp:st:C.2.UN.51}}}
\expandafter\def\csname SRef@cdp:st:C.2.UN.52\endcsname{\hyperlink{tag-01DX}{II.\ref*{cdp:st:C.2.UN.52}}}
\expandafter\def\csname SRef@cdp:st:C.2.UN.53\endcsname{\hyperlink{tag-01DY}{II.\ref*{cdp:st:C.2.UN.53}}}
\expandafter\def\csname SRef@cdp:st:C.2.UN.54\endcsname{\hyperlink{tag-01DZ}{II.\ref*{cdp:st:C.2.UN.54}}}
\expandafter\def\csname SRef@cdp:st:C.2.UN.69\endcsname{\hyperlink{tag-01EF}{II.\ref*{cdp:st:C.2.UN.69}}}
\expandafter\def\csname SRef@cdp:st:C.2.UN.80\endcsname{\hyperlink{tag-01ES}{II.\ref*{cdp:st:C.2.UN.80}}}
\expandafter\def\csname SRef@cdp:st:C.2.UN.81\endcsname{\hyperlink{tag-01ET}{II.\ref*{cdp:st:C.2.UN.81}}}
\expandafter\def\csname SRef@cdp:st:C.2.UN.82\endcsname{\hyperlink{tag-01EU}{II.\ref*{cdp:st:C.2.UN.82}}}
\expandafter\def\csname SRef@cdp:st:C.2.UN.83\endcsname{\hyperlink{tag-01EV}{II.\ref*{cdp:st:C.2.UN.83}}}
\expandafter\def\csname SRef@cdp:st:C.2.UN.84\endcsname{\hyperlink{tag-01EW}{II.\ref*{cdp:st:C.2.UN.84}}}
\expandafter\def\csname SRef@cdp:st:C.2.UN.78\endcsname{\hyperlink{tag-01EP}{II.\ref*{cdp:st:C.2.UN.78}}}
\expandafter\def\csname SRef@cdp:st:C.2.UN.72\endcsname{\hyperlink{tag-01EJ}{II.\ref*{cdp:st:C.2.UN.72}}}
\expandafter\def\csname SRef@cdp:st:C.2.UN.85\endcsname{\hyperlink{tag-01EX}{II.\ref*{cdp:st:C.2.UN.85}}}
\expandafter\def\csname SRef@cdp:st:C.2.UN.86\endcsname{\hyperlink{tag-01EY}{II.\ref*{cdp:st:C.2.UN.86}}}
\expandafter\def\csname SRef@cdp:st:C.2.UN.70\endcsname{\hyperlink{tag-01EH}{II.\ref*{cdp:st:C.2.UN.70}}}
\expandafter\def\csname SRef@cdp:st:C.2.UN.71\endcsname{\hyperlink{tag-01EI}{II.\ref*{cdp:st:C.2.UN.71}}}
\expandafter\def\csname SRef@cdp:st:C.2.UN.73\endcsname{\hyperlink{tag-01EK}{II.\ref*{cdp:st:C.2.UN.73}}}
\expandafter\def\csname SRef@cdp:st:C.2.UN.76\endcsname{\hyperlink{tag-01EN}{II.\ref*{cdp:st:C.2.UN.76}}}
\expandafter\def\csname SRef@cdp:st:C.2.UN.74\endcsname{\hyperlink{tag-01EL}{II.\ref*{cdp:st:C.2.UN.74}}}
\expandafter\def\csname SRef@cdp:st:C.2.UN.75\endcsname{\hyperlink{tag-01EM}{II.\ref*{cdp:st:C.2.UN.75}}}
\expandafter\def\csname SRef@cdp:st:C.2.UN.77\endcsname{\hyperlink{tag-01EO}{II.\ref*{cdp:st:C.2.UN.77}}}
\expandafter\def\csname SRef@cdp:st:C.2.UN.79\endcsname{\hyperlink{tag-01EQ}{II.\ref*{cdp:st:C.2.UN.79}}}
\expandafter\def\csname SRef@sec:II:06\endcsname{\hyperlink{tag-0342}{II.\ref*{sec:II:06}}}
\expandafter\def\csname SRef@sec:II:07\endcsname{\hyperlink{tag-0343}{II.\ref*{sec:II:07}}}
\expandafter\def\csname SRef@sec:II:all:08\endcsname{\hyperlink{tag-034B}{II.\ref*{sec:II:all:08}}}
\expandafter\def\csname SRef@sec:II:09\endcsname{\hyperlink{tag-0345}{II.\ref*{sec:II:09}}}
\expandafter\def\csname SRef@sec:II:10\endcsname{\hyperlink{tag-0346}{II.\ref*{sec:II:10}}}
\expandafter\def\csname SRef@sec:II:11\endcsname{\hyperlink{tag-0347}{II.\ref*{sec:II:11}}}
\expandafter\def\csname SRef@sec:II:all:12\endcsname{\hyperlink{tag-034C}{II.\ref*{sec:II:all:12}}}
\expandafter\def\csname SRef@cdp:st:C.2.G.1\endcsname{\hyperlink{tag-01BZ}{II.\ref*{cdp:st:C.2.G.1}}}
\expandafter\def\csname SRef@inv:E1.1\endcsname{\hyperlink{tag-01ST}{II.\ref*{inv:E1.1}}}
\expandafter\def\csname SRef@inv:E2.1\endcsname{\hyperlink{tag-01SU}{II.\ref*{inv:E2.1}}}
\expandafter\def\csname SRef@push:src:C.1.1\endcsname{\href{S6_SERIES_GRIND_III.pdf\#nameddest=tag-031B}{III.? (p.~?)} }
\expandafter\def\csname SRef@push:src:C.1.2\endcsname{\href{S6_SERIES_GRIND_III.pdf\#nameddest=tag-031C}{III.? (p.~?)} }
\expandafter\def\csname SRef@push:src:C.2.1\endcsname{\href{S6_SERIES_GRIND_III.pdf\#nameddest=tag-031D}{III.? (p.~?)} }
\expandafter\def\csname SRef@push:src:C.3.1\endcsname{\href{S6_SERIES_GRIND_III.pdf\#nameddest=tag-031G}{III.? (p.~?)} }
\expandafter\def\csname SRef@push:src:C.3.3\endcsname{\href{S6_SERIES_GRIND_III.pdf\#nameddest=tag-0320}{III.? (p.~?)} }
\expandafter\def\csname SRef@push:src:E.3.2\endcsname{\href{S6_SERIES_GRIND_I.pdf\#nameddest=tag-033J}{I.? (p.~?)} }
\expandafter\def\csname SRef@push:src:E.3.3\endcsname{\href{S6_SERIES_GRIND_I.pdf\#nameddest=tag-033K}{I.? (p.~?)} }
\expandafter\def\csname SRef@cdp:st:A.2.3\endcsname{\hyperlink{tag-01AH}{II.\ref*{cdp:st:A.2.3}}}
\expandafter\def\csname SRef@cdp:st:C.1.2\endcsname{\hyperlink{tag-01BO}{II.\ref*{cdp:st:C.1.2}}}
\expandafter\def\csname SRef@cdp:st:C.2.NC.1\endcsname{\hyperlink{tag-01C8}{II.\ref*{cdp:st:C.2.NC.1}}}
\expandafter\def\csname SRef@cdp:st:C.2.RAT.6\endcsname{\hyperlink{tag-01CI}{II.\ref*{cdp:st:C.2.RAT.6}}}
\expandafter\def\csname SRef@cdp:st:D.1.1\endcsname{\hyperlink{tag-01FU}{II.\ref*{cdp:st:D.1.1}}}
\expandafter\def\csname SRef@cdp:st:D.1.3\endcsname{\hyperlink{tag-01G1}{II.\ref*{cdp:st:D.1.3}}}
\expandafter\def\csname SRef@cdp:st:E2\endcsname{\hyperlink{tag-01GR}{II.\ref*{cdp:st:E2}}}
\expandafter\def\csname SRef@cdp:st:E3\endcsname{\hyperlink{tag-01GS}{II.\ref*{cdp:st:E3}}}
\expandafter\def\csname SRef@cdp:st:E4\endcsname{\hyperlink{tag-01GT}{II.\ref*{cdp:st:E4}}}
\expandafter\def\csname SRef@cdp:st:E19\endcsname{\hyperlink{tag-01GQ}{II.\ref*{cdp:st:E19}}}
\expandafter\def\csname SRef@maps:src:B.6.9\endcsname{\href{S6_SERIES_GRIND_IV.pdf\#nameddest=tag-02K1}{IV.? (p.~?)} }
\expandafter\def\csname SRef@maps:src:D.2.1\endcsname{\href{S6_SERIES_GRIND_III.pdf\#nameddest=tag-02KJ}{III.? (p.~?)} }
\expandafter\def\csname SRef@maps:src:F.1.2\endcsname{\href{S6_SERIES_GRIND_V.pdf\#nameddest=tag-02MV}{V.? (p.~?)} }
\expandafter\def\csname SRef@maps:src:F.1.3\endcsname{\href{S6_SERIES_GRIND_V.pdf\#nameddest=tag-02MW}{V.? (p.~?)} }
\expandafter\def\csname SRef@ari:C.7.5\endcsname{\href{S6_SERIES_GRIND_I.pdf\#nameddest=tag-001E}{I.? (p.~?)} }
\expandafter\def\csname SRef@wb:lem:triangle-generator-orientations\endcsname{\href{S6_SERIES_GRIND_I.pdf\#nameddest=tag-03BM}{I.? (p.~?)} }
\expandafter\def\csname SRef@wb:prop:period-input-audit\endcsname{\href{S6_SERIES_GRIND_I.pdf\#nameddest=tag-03BS}{I.? (p.~?)} }
\expandafter\def\csname SRef@wb:prop:period-exact-admissible-constants\endcsname{\href{S6_SERIES_GRIND_I.pdf\#nameddest=tag-03BR}{I.? (p.~?)} }
\expandafter\def\csname SRef@wb:lem:toric-chart-control\endcsname{\href{S6_SERIES_GRIND_I.pdf\#nameddest=tag-03BK}{I.? (p.~?)} }
\expandafter\def\csname SRef@wb:thm:candidate-cusp-geometry\endcsname{\href{S6_SERIES_GRIND_I.pdf\#nameddest=tag-03DK}{I.? (p.~?)} }
\expandafter\def\csname SRef@wb:prop:candidate-cusp-overlap\endcsname{\href{S6_SERIES_GRIND_I.pdf\#nameddest=tag-03BN}{I.? (p.~?)} }
\expandafter\def\csname SRef@wb:lem:middle-period-proper\endcsname{\href{S6_SERIES_GRIND_I.pdf\#nameddest=tag-03BI}{I.? (p.~?)} }
\expandafter\def\csname SRef@wb:prop:candidate-finite-fillings\endcsname{\href{S6_SERIES_GRIND_I.pdf\#nameddest=tag-03BO}{I.? (p.~?)} }
\expandafter\def\csname SRef@wb:lem:finite-filling-stalk-map\endcsname{\href{S6_SERIES_GRIND_I.pdf\#nameddest=tag-03BG}{I.? (p.~?)} }
\expandafter\def\csname SRef@wb:thm:candidate-global-geometry\endcsname{\href{S6_SERIES_GRIND_I.pdf\#nameddest=tag-03DL}{I.? (p.~?)} }
\expandafter\def\csname SRef@wb:prop:monodromy-certificate\endcsname{\href{S6_SERIES_GRIND_I.pdf\#nameddest=tag-03BQ}{I.? (p.~?)} }
\expandafter\def\csname SRef@wb:thm:cusp-lattices\endcsname{\href{S6_SERIES_GRIND_I.pdf\#nameddest=tag-03DO}{I.? (p.~?)} }
\expandafter\def\csname SRef@wb:thm:coinvariants\endcsname{\href{S6_SERIES_GRIND_I.pdf\#nameddest=tag-03DN}{I.? (p.~?)} }
\expandafter\def\csname SRef@wb:prop:finite-indices\endcsname{\href{S6_SERIES_GRIND_I.pdf\#nameddest=tag-03BP}{I.? (p.~?)} }
\expandafter\def\csname SRef@wb:prop:terminal-group\endcsname{\href{S6_SERIES_GRIND_I.pdf\#nameddest=tag-03BV}{I.? (p.~?)} }
\expandafter\def\csname SRef@wb:prop:terminal-degree-two\endcsname{\href{S6_SERIES_GRIND_I.pdf\#nameddest=tag-03BU}{I.? (p.~?)} }
\expandafter\def\csname SRef@wb:prop:terminal-all-degree\endcsname{\href{S6_SERIES_GRIND_I.pdf\#nameddest=tag-03BT}{I.? (p.~?)} }
\expandafter\def\csname SRef@wb:ffcert:prop:integral-summands\endcsname{\href{S6_SERIES_GRIND_I.pdf\#nameddest=tag-03A8}{I.? (p.~?)} }
\expandafter\def\csname SRef@wb:ffcert:thm:product-covers\endcsname{\href{S6_SERIES_GRIND_I.pdf\#nameddest=tag-03AD}{I.? (p.~?)} }
\expandafter\def\csname SRef@wb:ffcert:prop:quotient-charts\endcsname{\href{S6_SERIES_GRIND_I.pdf\#nameddest=tag-03A9}{I.? (p.~?)} }
\expandafter\def\csname SRef@wb:ffcert:thm:normalization\endcsname{\href{S6_SERIES_GRIND_I.pdf\#nameddest=tag-03AB}{I.? (p.~?)} }
\expandafter\def\csname SRef@wb:ffcert:thm:peripheral-kernel\endcsname{\href{S6_SERIES_GRIND_I.pdf\#nameddest=tag-03AC}{I.? (p.~?)} }
\expandafter\def\csname SRef@wb:ffcert:cor:covering-h1\endcsname{\href{S6_SERIES_GRIND_I.pdf\#nameddest=tag-03A7}{I.? (p.~?)} }
\expandafter\def\csname SRef@wb:ffcert:thm:normal-trivialization\endcsname{\href{S6_SERIES_GRIND_I.pdf\#nameddest=tag-03AA}{I.? (p.~?)} }
\expandafter\def\csname SRef@wb:lem:torus-algebraic-reduction-audit\endcsname{\href{S6_SERIES_GRIND_I.pdf\#nameddest=tag-03BL}{I.? (p.~?)} }
\expandafter\def\csname SRef@wb:prop:very-general-torus-audit\endcsname{\href{S6_SERIES_GRIND_I.pdf\#nameddest=tag-03BW}{I.? (p.~?)} }
\expandafter\def\csname SRef@wb:lem:meromorphic-descent-audit\endcsname{\href{S6_SERIES_GRIND_I.pdf\#nameddest=tag-03BH}{I.? (p.~?)} }
\expandafter\def\csname SRef@wb:thm:algebraic-dimension-audit\endcsname{\href{S6_SERIES_GRIND_I.pdf\#nameddest=tag-03DJ}{I.? (p.~?)} }
\expandafter\def\csname SRef@wb:lem:semidirect-abelianization\endcsname{\href{S6_SERIES_GRIND_I.pdf\#nameddest=tag-03BJ}{I.? (p.~?)} }
\expandafter\def\csname SRef@mapsI:top:thm:p0-integral-specialization\endcsname{\href{S6_SERIES_GRIND_I.pdf\#nameddest=tag-02Q5}{I.? (p.~?)} }
\expandafter\def\csname SRef@mapsI:top:prop:leray-e2\endcsname{\href{S6_SERIES_GRIND_I.pdf\#nameddest=tag-02PU}{I.? (p.~?)} }
\expandafter\def\csname SRef@mapsI:top:lem:first-transgression\endcsname{\href{S6_SERIES_GRIND_I.pdf\#nameddest=tag-02PN}{I.? (p.~?)} }
\expandafter\def\csname SRef@mapsI:top:lem:higher-transgressions\endcsname{\href{S6_SERIES_GRIND_I.pdf\#nameddest=tag-02PO}{I.? (p.~?)} }
\expandafter\def\csname SRef@mapsI:top:thm:h2-vanishing\endcsname{\href{S6_SERIES_GRIND_I.pdf\#nameddest=tag-02Q3}{I.? (p.~?)} }
\expandafter\def\csname SRef@mapsI:top:prop:geometric-pi1\endcsname{\href{S6_SERIES_GRIND_I.pdf\#nameddest=tag-02PT}{I.? (p.~?)} }
\expandafter\def\csname SRef@mapsI:top:cor:candidate-sphere\endcsname{\href{S6_SERIES_GRIND_I.pdf\#nameddest=tag-02PK}{I.? (p.~?)} }
\expandafter\def\csname SRef@mapsI:top:prop:topological-spine\endcsname{\href{S6_SERIES_GRIND_I.pdf\#nameddest=tag-02PV}{I.? (p.~?)} }
\expandafter\def\csname SRef@mapsI:top:lem:collapse-straightening\endcsname{\href{S6_SERIES_GRIND_I.pdf\#nameddest=tag-02PM}{I.? (p.~?)} }
\expandafter\def\csname SRef@mapsI:top:lem:positive-collapse\endcsname{\href{S6_SERIES_GRIND_I.pdf\#nameddest=tag-02PQ}{I.? (p.~?)} }
\expandafter\def\csname SRef@mapsI:top:lem:phase-spread-collapse\endcsname{\href{S6_SERIES_GRIND_I.pdf\#nameddest=tag-02PP}{I.? (p.~?)} }
\expandafter\def\csname SRef@mapsI:top:prop:verified-cellular-collapse\endcsname{\href{S6_SERIES_GRIND_I.pdf\#nameddest=tag-02PW}{I.? (p.~?)} }
\expandafter\def\csname SRef@mapsI:top:prop:verified-collapse-homology\endcsname{\href{S6_SERIES_GRIND_I.pdf\#nameddest=tag-02PX}{I.? (p.~?)} }
\expandafter\def\csname SRef@mapsI:top:thm:collapse-equals-specialization\endcsname{\href{S6_SERIES_GRIND_I.pdf\#nameddest=tag-02Q2}{I.? (p.~?)} }
\expandafter\def\csname SRef@mapsI:top:remark:L600\endcsname{\href{S6_SERIES_GRIND_I.pdf\#nameddest=tag-02Q0}{I.? (p.~?)} }
\expandafter\def\csname SRef@mapsI:top:record:vankampen\endcsname{\href{S6_SERIES_GRIND_I.pdf\#nameddest=tag-02PY}{I.? (p.~?)} }
\expandafter\def\csname SRef@mapsI:top:prop:actual-finite-band-maps\endcsname{\href{S6_SERIES_GRIND_I.pdf\#nameddest=tag-02PR}{I.? (p.~?)} }
\expandafter\def\csname SRef@mapsI:top:prop:cusp-mv-blocks\endcsname{\href{S6_SERIES_GRIND_I.pdf\#nameddest=tag-02PS}{I.? (p.~?)} }
\expandafter\def\csname SRef@mapsI:top:remark:L920\endcsname{\href{S6_SERIES_GRIND_I.pdf\#nameddest=tag-02Q1}{I.? (p.~?)} }
\expandafter\def\csname SRef@mapsI:top:thm:mv-leray-comparison\endcsname{\href{S6_SERIES_GRIND_I.pdf\#nameddest=tag-02Q4}{I.? (p.~?)} }
\expandafter\def\csname SRef@mapsI:top:remark:L1113\endcsname{\href{S6_SERIES_GRIND_I.pdf\#nameddest=tag-02PZ}{I.? (p.~?)} }
\expandafter\def\csname SRef@mapsI:top:lem:cohomological-closure\endcsname{\href{S6_SERIES_GRIND_I.pdf\#nameddest=tag-02PL}{I.? (p.~?)} }
\expandafter\def\csname SRef@mapsI:top:thm:sphere-recognition\endcsname{\href{S6_SERIES_GRIND_I.pdf\#nameddest=tag-02Q6}{I.? (p.~?)} }
\expandafter\def\csname SRef@mapsI:top:cor:abstract-s6\endcsname{\href{S6_SERIES_GRIND_I.pdf\#nameddest=tag-02PI}{I.? (p.~?)} }
\expandafter\def\csname SRef@mapsI:top:cor:candidate-mv-sphere\endcsname{\href{S6_SERIES_GRIND_I.pdf\#nameddest=tag-02PJ}{I.? (p.~?)} }
\expandafter\def\csname SRef@wb:cusp-cone-S12-almost-complex-obstruction\endcsname{\href{S6_SERIES_GRIND_III.pdf\#nameddest=tag-035S}{III.? (p.~?)} }
\expandafter\def\csname SRef@courseIII:almost-complex-structures-and-integrability:section-1\endcsname{\href{S6_SERIES_GRIND_III.pdf\#nameddest=tag-01I4}{III.? (p.~?)} }
\expandafter\def\csname SRef@courseIII:almost-complex-structures-and-integrability:section-2\endcsname{\href{S6_SERIES_GRIND_III.pdf\#nameddest=tag-01I7}{III.? (p.~?)} }
\expandafter\def\csname SRef@courseIII:almost-complex-structures-and-integrability:section-3\endcsname{\href{S6_SERIES_GRIND_III.pdf\#nameddest=tag-01I8}{III.? (p.~?)} }
\expandafter\def\csname SRef@courseIII:almost-complex-structures-and-integrability:section-4\endcsname{\href{S6_SERIES_GRIND_III.pdf\#nameddest=tag-01I9}{III.? (p.~?)} }
\expandafter\def\csname SRef@courseIII:almost-complex-structures-and-integrability:section-5\endcsname{\href{S6_SERIES_GRIND_III.pdf\#nameddest=tag-01IA}{III.? (p.~?)} }
\expandafter\def\csname SRef@courseIII:almost-complex-structures-and-integrability:section-6\endcsname{\href{S6_SERIES_GRIND_III.pdf\#nameddest=tag-01IB}{III.? (p.~?)} }
\expandafter\def\csname SRef@courseIII:almost-complex-structures-and-integrability:section-7\endcsname{\href{S6_SERIES_GRIND_III.pdf\#nameddest=tag-01IC}{III.? (p.~?)} }
\expandafter\def\csname SRef@courseIII:almost-complex-structures-and-integrability:section-8\endcsname{\href{S6_SERIES_GRIND_III.pdf\#nameddest=tag-01ID}{III.? (p.~?)} }
\expandafter\def\csname SRef@courseIII:almost-complex-structures-and-integrability:section-9\endcsname{\href{S6_SERIES_GRIND_III.pdf\#nameddest=tag-01IE}{III.? (p.~?)} }
\expandafter\def\csname SRef@courseIII:almost-complex-structures-and-integrability:section-10\endcsname{\href{S6_SERIES_GRIND_III.pdf\#nameddest=tag-01I5}{III.? (p.~?)} }
\expandafter\def\csname SRef@courseIII:almost-complex-structures-and-integrability:section-11\endcsname{\href{S6_SERIES_GRIND_III.pdf\#nameddest=tag-01I6}{III.? (p.~?)} }
\expandafter\def\csname SRef@mapsI:course:CG-S6-07-HOMOTOPY\endcsname{\href{S6_SERIES_GRIND_I.pdf\#nameddest=tag-02OR}{I.? (p.~?)} }
\expandafter\def\csname SRef@courseIII:integral-homology-and-sphere-recognition:section\endcsname{\href{S6_SERIES_GRIND_III.pdf\#nameddest=tag-01K9}{III.? (p.~?)} }
\expandafter\def\csname SRef@courseIII:sphere-twists-and-the-original-smooth-six-sphere-comparison:twist-actual-tube\endcsname{\href{S6_SERIES_GRIND_III.pdf\#nameddest=tag-01MN}{III.? (p.~?)} }
\expandafter\def\csname SRef@courseIII:sphere-twists-and-the-original-smooth-six-sphere-comparison:twist-signed-action\endcsname{\href{S6_SERIES_GRIND_III.pdf\#nameddest=tag-01MT}{III.? (p.~?)} }
\expandafter\def\csname SRef@courseIII:sphere-twists-and-the-original-smooth-six-sphere-comparison:twist-integral-lagrangians\endcsname{\href{S6_SERIES_GRIND_III.pdf\#nameddest=tag-01MP}{III.? (p.~?)} }
\expandafter\def\csname SRef@courseIII:sphere-twists-and-the-original-smooth-six-sphere-comparison:twist-original-middle-level\endcsname{\href{S6_SERIES_GRIND_III.pdf\#nameddest=tag-01MR}{III.? (p.~?)} }
\expandafter\def\csname SRef@courseIII:sphere-twists-and-the-original-smooth-six-sphere-comparison:twist-original-gluing\endcsname{\href{S6_SERIES_GRIND_III.pdf\#nameddest=tag-01MQ}{III.? (p.~?)} }
\expandafter\def\csname SRef@courseIII:sphere-twists-and-the-original-smooth-six-sphere-comparison:twist-actual-cancellation\endcsname{\href{S6_SERIES_GRIND_III.pdf\#nameddest=tag-01MM}{III.? (p.~?)} }
\expandafter\def\csname SRef@courseIII:sphere-twists-and-the-original-smooth-six-sphere-comparison:twist-original-smooth-map\endcsname{\href{S6_SERIES_GRIND_III.pdf\#nameddest=tag-01MS}{III.? (p.~?)} }
\expandafter\def\csname SRef@courseIII:sphere-twists-and-the-original-smooth-six-sphere-comparison:twist-general-stable-framing\endcsname{\href{S6_SERIES_GRIND_III.pdf\#nameddest=tag-01MO}{III.? (p.~?)} }
\expandafter\def\csname SRef@courseIII:sphere-twists-and-the-original-smooth-six-sphere-comparison:twist-worked-exercises\endcsname{\href{S6_SERIES_GRIND_III.pdf\#nameddest=tag-01MU}{III.? (p.~?)} }
\expandafter\def\csname SRef@mapsaudit:almost-complex-homotopy\endcsname{\href{S6_SERIES_GRIND_III.pdf\#nameddest=tag-02QD}{III.? (p.~?)} }
\expandafter\def\csname SRef@mapsI:hodge:line\endcsname{\href{S6_SERIES_GRIND_I.pdf\#nameddest=tag-02OV}{I.? (p.~?)} }
\expandafter\def\csname SRef@mapsI:chain:lemma:1.1\endcsname{\href{S6_SERIES_GRIND_I.pdf\#nameddest=tag-02OL}{I.? (p.~?)} }
\expandafter\def\csname SRef@mapsI:chain:lemma:2.1\endcsname{\href{S6_SERIES_GRIND_I.pdf\#nameddest=tag-02OM}{I.? (p.~?)} }
\expandafter\def\csname SRef@mapsI:chain:corollary:2.2\endcsname{\href{S6_SERIES_GRIND_I.pdf\#nameddest=tag-02OK}{I.? (p.~?)} }
\expandafter\def\csname SRef@mapsI:chain:lemma:4.1\endcsname{\href{S6_SERIES_GRIND_I.pdf\#nameddest=tag-02ON}{I.? (p.~?)} }
\expandafter\def\csname SRef@mapsI:cw:hurewicz-theorem-result\endcsname{\href{S6_SERIES_GRIND_I.pdf\#nameddest=tag-02OU}{I.? (p.~?)} }
\expandafter\def\csname SRef@mapsI:cw:exercise:5.1\endcsname{\href{S6_SERIES_GRIND_I.pdf\#nameddest=tag-02OS}{I.? (p.~?)} }
\expandafter\def\csname SRef@mapsI:cw:exercise:5.2\endcsname{\href{S6_SERIES_GRIND_I.pdf\#nameddest=tag-02OT}{I.? (p.~?)} }
\expandafter\def\csname SRef@mapsI:pf:theorem:5.1\endcsname{\href{S6_SERIES_GRIND_I.pdf\#nameddest=tag-02OY}{I.? (p.~?)} }
\expandafter\def\csname SRef@mapsI:pf:exercise:6.1\endcsname{\href{S6_SERIES_GRIND_I.pdf\#nameddest=tag-02OW}{I.? (p.~?)} }
\expandafter\def\csname SRef@mapsI:pf:exercise:6.2\endcsname{\href{S6_SERIES_GRIND_I.pdf\#nameddest=tag-02OX}{I.? (p.~?)} }
\expandafter\def\csname SRef@courseIII:thom-euler-for-lesson11:CG-S6-CHAIN-LEMMA-5.1\endcsname{\href{S6_SERIES_GRIND_III.pdf\#nameddest=tag-01N3}{III.? (p.~?)} }
\expandafter\def\csname SRef@courseIII:thom-euler-for-lesson11:CG-S6-CHAIN-THEOREM-5.2\endcsname{\href{S6_SERIES_GRIND_III.pdf\#nameddest=tag-01N7}{III.? (p.~?)} }
\expandafter\def\csname SRef@courseIII:thom-euler-for-lesson11:CG-S6-CHAIN-PROPOSITION-6.1\endcsname{\href{S6_SERIES_GRIND_III.pdf\#nameddest=tag-01N4}{III.? (p.~?)} }
\expandafter\def\csname SRef@courseIII:thom-euler-for-lesson11:CG-S6-CHAIN-THEOREM-7.1\endcsname{\href{S6_SERIES_GRIND_III.pdf\#nameddest=tag-01N8}{III.? (p.~?)} }
\expandafter\def\csname SRef@courseIII:thom-euler-for-lesson11:CG-S6-CHAIN-PROPOSITION-8.1\endcsname{\href{S6_SERIES_GRIND_III.pdf\#nameddest=tag-01N5}{III.? (p.~?)} }
\expandafter\def\csname SRef@courseIII:thom-euler-for-lesson11:CG-S6-CHAIN-PROPOSITION-8.2\endcsname{\href{S6_SERIES_GRIND_III.pdf\#nameddest=tag-01N6}{III.? (p.~?)} }
\expandafter\def\csname SRef@courseIII:projective-splitting-for-lesson11:CG-S6-PROJECTIVE-1\endcsname{\href{S6_SERIES_GRIND_III.pdf\#nameddest=tag-01M7}{III.? (p.~?)} }
\expandafter\def\csname SRef@courseIII:projective-splitting-for-lesson11:CG-S6-PROJECTIVE-2\endcsname{\href{S6_SERIES_GRIND_III.pdf\#nameddest=tag-01M8}{III.? (p.~?)} }
\expandafter\def\csname SRef@courseIII:projective-splitting-for-lesson11:CG-S6-PROJECTIVE-3\endcsname{\href{S6_SERIES_GRIND_III.pdf\#nameddest=tag-01M9}{III.? (p.~?)} }
\expandafter\def\csname SRef@courseIII:projective-splitting-for-lesson11:CG-S6-PROJECTIVE-4\endcsname{\href{S6_SERIES_GRIND_III.pdf\#nameddest=tag-01MA}{III.? (p.~?)} }
\expandafter\def\csname SRef@courseIII:projective-splitting-for-lesson11:CG-S6-PROJECTIVE-5\endcsname{\href{S6_SERIES_GRIND_III.pdf\#nameddest=tag-01MB}{III.? (p.~?)} }
\expandafter\def\csname SRef@courseIII:integral-chern-classes-for-lesson11:CG-S6-CHERN-1\endcsname{\href{S6_SERIES_GRIND_III.pdf\#nameddest=tag-01K6}{III.? (p.~?)} }
\expandafter\def\csname SRef@courseIII:integral-chern-classes-for-lesson11:CG-S6-CHERN-2\endcsname{\href{S6_SERIES_GRIND_III.pdf\#nameddest=tag-01K7}{III.? (p.~?)} }
\expandafter\def\csname SRef@courseIII:integral-chern-classes-for-lesson11:CG-S6-CHERN-3\endcsname{\href{S6_SERIES_GRIND_III.pdf\#nameddest=tag-01K8}{III.? (p.~?)} }
\expandafter\def\csname SRef@courseIII:framed6:1\endcsname{\href{S6_SERIES_GRIND_III.pdf\#nameddest=tag-01J8}{III.? (p.~?)} }
\expandafter\def\csname SRef@courseIII:framed6:2\endcsname{\href{S6_SERIES_GRIND_III.pdf\#nameddest=tag-01JB}{III.? (p.~?)} }
\expandafter\def\csname SRef@courseIII:framed6:3\endcsname{\href{S6_SERIES_GRIND_III.pdf\#nameddest=tag-01JC}{III.? (p.~?)} }
\expandafter\def\csname SRef@courseIII:framed6:4\endcsname{\href{S6_SERIES_GRIND_III.pdf\#nameddest=tag-01JD}{III.? (p.~?)} }
\expandafter\def\csname SRef@courseIII:framed6:5\endcsname{\href{S6_SERIES_GRIND_III.pdf\#nameddest=tag-01JE}{III.? (p.~?)} }
\expandafter\def\csname SRef@courseIII:framed6:6\endcsname{\href{S6_SERIES_GRIND_III.pdf\#nameddest=tag-01JF}{III.? (p.~?)} }
\expandafter\def\csname SRef@courseIII:framed6:7\endcsname{\href{S6_SERIES_GRIND_III.pdf\#nameddest=tag-01JG}{III.? (p.~?)} }
\expandafter\def\csname SRef@courseIII:framed6:8\endcsname{\href{S6_SERIES_GRIND_III.pdf\#nameddest=tag-01JH}{III.? (p.~?)} }
\expandafter\def\csname SRef@courseIII:framed6:9\endcsname{\href{S6_SERIES_GRIND_III.pdf\#nameddest=tag-01JI}{III.? (p.~?)} }
\expandafter\def\csname SRef@courseIII:framed6:10.1\endcsname{\href{S6_SERIES_GRIND_III.pdf\#nameddest=tag-01J9}{III.? (p.~?)} }
\expandafter\def\csname SRef@courseIII:framed6:10.2\endcsname{\href{S6_SERIES_GRIND_III.pdf\#nameddest=tag-01JA}{III.? (p.~?)} }
\expandafter\def\csname SRef@courseIII:whitney-receiving:1\endcsname{\href{S6_SERIES_GRIND_III.pdf\#nameddest=tag-01NG}{III.? (p.~?)} }
\expandafter\def\csname SRef@courseIII:whitney-receiving:2\endcsname{\href{S6_SERIES_GRIND_III.pdf\#nameddest=tag-01NH}{III.? (p.~?)} }
\expandafter\def\csname SRef@courseIII:whitney-receiving:3\endcsname{\href{S6_SERIES_GRIND_III.pdf\#nameddest=tag-01NI}{III.? (p.~?)} }
\expandafter\def\csname SRef@courseIII:whitney-receiving:4\endcsname{\href{S6_SERIES_GRIND_III.pdf\#nameddest=tag-01NJ}{III.? (p.~?)} }
\expandafter\def\csname SRef@courseIII:whitney-receiving:5.1\endcsname{\href{S6_SERIES_GRIND_III.pdf\#nameddest=tag-01NK}{III.? (p.~?)} }
\expandafter\def\csname SRef@courseIII:whitney-receiving:5.2\endcsname{\href{S6_SERIES_GRIND_III.pdf\#nameddest=tag-01NL}{III.? (p.~?)} }
\expandafter\def\csname SRef@courseIII:higher-slides:1\endcsname{\href{S6_SERIES_GRIND_III.pdf\#nameddest=tag-01JY}{III.? (p.~?)} }
\expandafter\def\csname SRef@courseIII:higher-slides:2\endcsname{\href{S6_SERIES_GRIND_III.pdf\#nameddest=tag-01JZ}{III.? (p.~?)} }
\expandafter\def\csname SRef@courseIII:higher-slides:3\endcsname{\href{S6_SERIES_GRIND_III.pdf\#nameddest=tag-01K0}{III.? (p.~?)} }
\expandafter\def\csname SRef@courseIII:higher-slides:4\endcsname{\href{S6_SERIES_GRIND_III.pdf\#nameddest=tag-01K1}{III.? (p.~?)} }
\expandafter\def\csname SRef@courseIII:higher-slides:5\endcsname{\href{S6_SERIES_GRIND_III.pdf\#nameddest=tag-01K2}{III.? (p.~?)} }
\expandafter\def\csname SRef@courseIII:higher-slides:6\endcsname{\href{S6_SERIES_GRIND_III.pdf\#nameddest=tag-01K3}{III.? (p.~?)} }
\expandafter\def\csname SRef@courseIII:higher-slides:8.1\endcsname{\href{S6_SERIES_GRIND_III.pdf\#nameddest=tag-01K4}{III.? (p.~?)} }
\expandafter\def\csname SRef@courseIII:higher-slides:8.2\endcsname{\href{S6_SERIES_GRIND_III.pdf\#nameddest=tag-01K5}{III.? (p.~?)} }
\expandafter\def\csname SRef@courseIII:middle-cancel:1\endcsname{\href{S6_SERIES_GRIND_III.pdf\#nameddest=tag-01KT}{III.? (p.~?)} }
\expandafter\def\csname SRef@courseIII:middle-cancel:2\endcsname{\href{S6_SERIES_GRIND_III.pdf\#nameddest=tag-01KU}{III.? (p.~?)} }
\expandafter\def\csname SRef@courseIII:middle-cancel:3\endcsname{\href{S6_SERIES_GRIND_III.pdf\#nameddest=tag-01KV}{III.? (p.~?)} }
\expandafter\def\csname SRef@courseIII:middle-cancel:4\endcsname{\href{S6_SERIES_GRIND_III.pdf\#nameddest=tag-01KW}{III.? (p.~?)} }
\expandafter\def\csname SRef@courseIII:middle-cancel:5\endcsname{\href{S6_SERIES_GRIND_III.pdf\#nameddest=tag-01KX}{III.? (p.~?)} }
\expandafter\def\csname SRef@courseIII:middle-cancel:6\endcsname{\href{S6_SERIES_GRIND_III.pdf\#nameddest=tag-01KY}{III.? (p.~?)} }
\expandafter\def\csname SRef@courseIII:middle-cancel:7.1\endcsname{\href{S6_SERIES_GRIND_III.pdf\#nameddest=tag-01KZ}{III.? (p.~?)} }
\expandafter\def\csname SRef@courseIII:middle-cancel:7.2\endcsname{\href{S6_SERIES_GRIND_III.pdf\#nameddest=tag-01L0}{III.? (p.~?)} }
\expandafter\def\csname SRef@courseIII:morse-cancel:1\endcsname{\href{S6_SERIES_GRIND_III.pdf\#nameddest=tag-01L5}{III.? (p.~?)} }
\expandafter\def\csname SRef@courseIII:morse-cancel:2.1\endcsname{\href{S6_SERIES_GRIND_III.pdf\#nameddest=tag-01L6}{III.? (p.~?)} }
\expandafter\def\csname SRef@courseIII:morse-cancel:6\endcsname{\href{S6_SERIES_GRIND_III.pdf\#nameddest=tag-01L7}{III.? (p.~?)} }
\expandafter\def\csname SRef@courseIII:morse-cancel:7.1\endcsname{\href{S6_SERIES_GRIND_III.pdf\#nameddest=tag-01L8}{III.? (p.~?)} }
\expandafter\def\csname SRef@courseIII:morse-cancel:7.2\endcsname{\href{S6_SERIES_GRIND_III.pdf\#nameddest=tag-01L9}{III.? (p.~?)} }
\expandafter\def\csname SRef@courseIII:spin6:1\endcsname{\href{S6_SERIES_GRIND_III.pdf\#nameddest=tag-01MV}{III.? (p.~?)} }
\expandafter\def\csname SRef@courseIII:spin6:2\endcsname{\href{S6_SERIES_GRIND_III.pdf\#nameddest=tag-01MW}{III.? (p.~?)} }
\expandafter\def\csname SRef@courseIII:spin6:3\endcsname{\href{S6_SERIES_GRIND_III.pdf\#nameddest=tag-01MX}{III.? (p.~?)} }
\expandafter\def\csname SRef@courseIII:spin6:4\endcsname{\href{S6_SERIES_GRIND_III.pdf\#nameddest=tag-01MY}{III.? (p.~?)} }
\expandafter\def\csname SRef@courseIII:spin6:5\endcsname{\href{S6_SERIES_GRIND_III.pdf\#nameddest=tag-01MZ}{III.? (p.~?)} }
\expandafter\def\csname SRef@courseIII:spin6:6\endcsname{\href{S6_SERIES_GRIND_III.pdf\#nameddest=tag-01N0}{III.? (p.~?)} }
\expandafter\def\csname SRef@courseIII:spin6:7.1\endcsname{\href{S6_SERIES_GRIND_III.pdf\#nameddest=tag-01N1}{III.? (p.~?)} }
\expandafter\def\csname SRef@courseIII:spin6:7.2\endcsname{\href{S6_SERIES_GRIND_III.pdf\#nameddest=tag-01N2}{III.? (p.~?)} }
\expandafter\def\csname SRef@courseIII:integral-stable:1\endcsname{\href{S6_SERIES_GRIND_III.pdf\#nameddest=tag-01KA}{III.? (p.~?)} }
\expandafter\def\csname SRef@courseIII:integral-stable:2\endcsname{\href{S6_SERIES_GRIND_III.pdf\#nameddest=tag-01KB}{III.? (p.~?)} }
\expandafter\def\csname SRef@courseIII:integral-stable:3\endcsname{\href{S6_SERIES_GRIND_III.pdf\#nameddest=tag-01KC}{III.? (p.~?)} }
\expandafter\def\csname SRef@courseIII:integral-stable:4\endcsname{\href{S6_SERIES_GRIND_III.pdf\#nameddest=tag-01KD}{III.? (p.~?)} }
\expandafter\def\csname SRef@courseIII:integral-stable:5\endcsname{\href{S6_SERIES_GRIND_III.pdf\#nameddest=tag-01KE}{III.? (p.~?)} }
\expandafter\def\csname SRef@courseIII:integral-stable:6\endcsname{\href{S6_SERIES_GRIND_III.pdf\#nameddest=tag-01KF}{III.? (p.~?)} }
\expandafter\def\csname SRef@courseIII:integral-stable:7\endcsname{\href{S6_SERIES_GRIND_III.pdf\#nameddest=tag-01KG}{III.? (p.~?)} }
\expandafter\def\csname SRef@courseIII:integral-stable:8.1\endcsname{\href{S6_SERIES_GRIND_III.pdf\#nameddest=tag-01KH}{III.? (p.~?)} }
\expandafter\def\csname SRef@courseIII:integral-stable:8.2\endcsname{\href{S6_SERIES_GRIND_III.pdf\#nameddest=tag-01KI}{III.? (p.~?)} }
\expandafter\def\csname SRef@courseIII:pontryagin-thom:1\endcsname{\href{S6_SERIES_GRIND_III.pdf\#nameddest=tag-01LQ}{III.? (p.~?)} }
\expandafter\def\csname SRef@courseIII:pontryagin-thom:2\endcsname{\href{S6_SERIES_GRIND_III.pdf\#nameddest=tag-01LR}{III.? (p.~?)} }
\expandafter\def\csname SRef@courseIII:pontryagin-thom:3\endcsname{\href{S6_SERIES_GRIND_III.pdf\#nameddest=tag-01LS}{III.? (p.~?)} }
\expandafter\def\csname SRef@courseIII:pontryagin-thom:4\endcsname{\href{S6_SERIES_GRIND_III.pdf\#nameddest=tag-01LT}{III.? (p.~?)} }
\expandafter\def\csname SRef@courseIII:pontryagin-thom:5\endcsname{\href{S6_SERIES_GRIND_III.pdf\#nameddest=tag-01LU}{III.? (p.~?)} }
\expandafter\def\csname SRef@courseIII:pontryagin-thom:6\endcsname{\href{S6_SERIES_GRIND_III.pdf\#nameddest=tag-01LV}{III.? (p.~?)} }
\expandafter\def\csname SRef@courseIII:pontryagin-thom:7\endcsname{\href{S6_SERIES_GRIND_III.pdf\#nameddest=tag-01LW}{III.? (p.~?)} }
\expandafter\def\csname SRef@courseIII:pontryagin-thom:8\endcsname{\href{S6_SERIES_GRIND_III.pdf\#nameddest=tag-01LX}{III.? (p.~?)} }
\expandafter\def\csname SRef@courseIII:pontryagin-thom:9.1\endcsname{\href{S6_SERIES_GRIND_III.pdf\#nameddest=tag-01LY}{III.? (p.~?)} }
\expandafter\def\csname SRef@courseIII:pontryagin-thom:9.2\endcsname{\href{S6_SERIES_GRIND_III.pdf\#nameddest=tag-01LZ}{III.? (p.~?)} }
\expandafter\def\csname SRef@courseIII:relative-morse:1\endcsname{\href{S6_SERIES_GRIND_III.pdf\#nameddest=tag-01MG}{III.? (p.~?)} }
\expandafter\def\csname SRef@courseIII:relative-morse:2\endcsname{\href{S6_SERIES_GRIND_III.pdf\#nameddest=tag-01MH}{III.? (p.~?)} }
\expandafter\def\csname SRef@courseIII:relative-morse:3\endcsname{\href{S6_SERIES_GRIND_III.pdf\#nameddest=tag-01MI}{III.? (p.~?)} }
\expandafter\def\csname SRef@courseIII:relative-morse:4\endcsname{\href{S6_SERIES_GRIND_III.pdf\#nameddest=tag-01MJ}{III.? (p.~?)} }
\expandafter\def\csname SRef@courseIII:relative-morse:5.1\endcsname{\href{S6_SERIES_GRIND_III.pdf\#nameddest=tag-01MK}{III.? (p.~?)} }
\expandafter\def\csname SRef@courseIII:relative-morse:5.2\endcsname{\href{S6_SERIES_GRIND_III.pdf\#nameddest=tag-01ML}{III.? (p.~?)} }
\expandafter\def\csname SRef@courseIII:handle-arrangement:1\endcsname{\href{S6_SERIES_GRIND_III.pdf\#nameddest=tag-01JJ}{III.? (p.~?)} }
\expandafter\def\csname SRef@courseIII:handle-arrangement:2\endcsname{\href{S6_SERIES_GRIND_III.pdf\#nameddest=tag-01JK}{III.? (p.~?)} }
\expandafter\def\csname SRef@courseIII:handle-arrangement:3\endcsname{\href{S6_SERIES_GRIND_III.pdf\#nameddest=tag-01JL}{III.? (p.~?)} }
\expandafter\def\csname SRef@courseIII:handle-arrangement:4\endcsname{\href{S6_SERIES_GRIND_III.pdf\#nameddest=tag-01JM}{III.? (p.~?)} }
\expandafter\def\csname SRef@courseIII:handle-arrangement:5.1\endcsname{\href{S6_SERIES_GRIND_III.pdf\#nameddest=tag-01JN}{III.? (p.~?)} }
\expandafter\def\csname SRef@courseIII:handle-arrangement:5.2\endcsname{\href{S6_SERIES_GRIND_III.pdf\#nameddest=tag-01JO}{III.? (p.~?)} }
\expandafter\def\csname SRef@courseIII:zero-top:1\endcsname{\href{S6_SERIES_GRIND_III.pdf\#nameddest=tag-01NU}{III.? (p.~?)} }
\expandafter\def\csname SRef@courseIII:zero-top:2\endcsname{\href{S6_SERIES_GRIND_III.pdf\#nameddest=tag-01NV}{III.? (p.~?)} }
\expandafter\def\csname SRef@courseIII:zero-top:3\endcsname{\href{S6_SERIES_GRIND_III.pdf\#nameddest=tag-01NW}{III.? (p.~?)} }
\expandafter\def\csname SRef@courseIII:zero-top:4\endcsname{\href{S6_SERIES_GRIND_III.pdf\#nameddest=tag-01NX}{III.? (p.~?)} }
\expandafter\def\csname SRef@courseIII:zero-top:5\endcsname{\href{S6_SERIES_GRIND_III.pdf\#nameddest=tag-01NY}{III.? (p.~?)} }
\expandafter\def\csname SRef@courseIII:zero-top:6.1\endcsname{\href{S6_SERIES_GRIND_III.pdf\#nameddest=tag-01NZ}{III.? (p.~?)} }
\expandafter\def\csname SRef@courseIII:zero-top:6.2\endcsname{\href{S6_SERIES_GRIND_III.pdf\#nameddest=tag-01O0}{III.? (p.~?)} }
\expandafter\def\csname SRef@courseIII:trajectory-lowering:1\endcsname{\href{S6_SERIES_GRIND_III.pdf\#nameddest=tag-01N9}{III.? (p.~?)} }
\expandafter\def\csname SRef@courseIII:trajectory-lowering:2\endcsname{\href{S6_SERIES_GRIND_III.pdf\#nameddest=tag-01NA}{III.? (p.~?)} }
\expandafter\def\csname SRef@courseIII:trajectory-lowering:3\endcsname{\href{S6_SERIES_GRIND_III.pdf\#nameddest=tag-01NB}{III.? (p.~?)} }
\expandafter\def\csname SRef@courseIII:trajectory-lowering:4\endcsname{\href{S6_SERIES_GRIND_III.pdf\#nameddest=tag-01NC}{III.? (p.~?)} }
\expandafter\def\csname SRef@courseIII:trajectory-lowering:5\endcsname{\href{S6_SERIES_GRIND_III.pdf\#nameddest=tag-01ND}{III.? (p.~?)} }
\expandafter\def\csname SRef@courseIII:trajectory-lowering:6.1\endcsname{\href{S6_SERIES_GRIND_III.pdf\#nameddest=tag-01NE}{III.? (p.~?)} }
\expandafter\def\csname SRef@courseIII:trajectory-lowering:6.2\endcsname{\href{S6_SERIES_GRIND_III.pdf\#nameddest=tag-01NF}{III.? (p.~?)} }
\expandafter\def\csname SRef@courseIII:belt-complement:1\endcsname{\href{S6_SERIES_GRIND_III.pdf\#nameddest=tag-01IF}{III.? (p.~?)} }
\expandafter\def\csname SRef@courseIII:belt-complement:2\endcsname{\href{S6_SERIES_GRIND_III.pdf\#nameddest=tag-01IG}{III.? (p.~?)} }
\expandafter\def\csname SRef@courseIII:belt-complement:3\endcsname{\href{S6_SERIES_GRIND_III.pdf\#nameddest=tag-01IH}{III.? (p.~?)} }
\expandafter\def\csname SRef@courseIII:belt-complement:4\endcsname{\href{S6_SERIES_GRIND_III.pdf\#nameddest=tag-01II}{III.? (p.~?)} }
\expandafter\def\csname SRef@courseIII:belt-complement:5\endcsname{\href{S6_SERIES_GRIND_III.pdf\#nameddest=tag-01IJ}{III.? (p.~?)} }
\expandafter\def\csname SRef@courseIII:belt-complement:6\endcsname{\href{S6_SERIES_GRIND_III.pdf\#nameddest=tag-01IK}{III.? (p.~?)} }
\expandafter\def\csname SRef@courseIII:belt-complement:7\endcsname{\href{S6_SERIES_GRIND_III.pdf\#nameddest=tag-01IL}{III.? (p.~?)} }
\expandafter\def\csname SRef@courseIII:whitney-support:1\endcsname{\href{S6_SERIES_GRIND_III.pdf\#nameddest=tag-01NM}{III.? (p.~?)} }
\expandafter\def\csname SRef@courseIII:whitney-support:2\endcsname{\href{S6_SERIES_GRIND_III.pdf\#nameddest=tag-01NN}{III.? (p.~?)} }
\expandafter\def\csname SRef@courseIII:whitney-support:3\endcsname{\href{S6_SERIES_GRIND_III.pdf\#nameddest=tag-01NO}{III.? (p.~?)} }
\expandafter\def\csname SRef@courseIII:whitney-support:4\endcsname{\href{S6_SERIES_GRIND_III.pdf\#nameddest=tag-01NP}{III.? (p.~?)} }
\expandafter\def\csname SRef@courseIII:whitney-support:5\endcsname{\href{S6_SERIES_GRIND_III.pdf\#nameddest=tag-01NQ}{III.? (p.~?)} }
\expandafter\def\csname SRef@courseIII:whitney-support:6\endcsname{\href{S6_SERIES_GRIND_III.pdf\#nameddest=tag-01NR}{III.? (p.~?)} }
\expandafter\def\csname SRef@courseIII:whitney-support:7.1\endcsname{\href{S6_SERIES_GRIND_III.pdf\#nameddest=tag-01NS}{III.? (p.~?)} }
\expandafter\def\csname SRef@courseIII:whitney-support:7.2\endcsname{\href{S6_SERIES_GRIND_III.pdf\#nameddest=tag-01NT}{III.? (p.~?)} }
\expandafter\def\csname SRef@courseIII:handle-chains:1\endcsname{\href{S6_SERIES_GRIND_III.pdf\#nameddest=tag-01JP}{III.? (p.~?)} }
\expandafter\def\csname SRef@courseIII:handle-chains:2\endcsname{\href{S6_SERIES_GRIND_III.pdf\#nameddest=tag-01JQ}{III.? (p.~?)} }
\expandafter\def\csname SRef@courseIII:handle-chains:3\endcsname{\href{S6_SERIES_GRIND_III.pdf\#nameddest=tag-01JR}{III.? (p.~?)} }
\expandafter\def\csname SRef@courseIII:handle-chains:4\endcsname{\href{S6_SERIES_GRIND_III.pdf\#nameddest=tag-01JS}{III.? (p.~?)} }
\expandafter\def\csname SRef@courseIII:handle-chains:5\endcsname{\href{S6_SERIES_GRIND_III.pdf\#nameddest=tag-01JT}{III.? (p.~?)} }
\expandafter\def\csname SRef@courseIII:handle-chains:6.1\endcsname{\href{S6_SERIES_GRIND_III.pdf\#nameddest=tag-01JU}{III.? (p.~?)} }
\expandafter\def\csname SRef@courseIII:handle-chains:6.2\endcsname{\href{S6_SERIES_GRIND_III.pdf\#nameddest=tag-01JV}{III.? (p.~?)} }
\expandafter\def\csname SRef@courseIII:handle-chains:8.1\endcsname{\href{S6_SERIES_GRIND_III.pdf\#nameddest=tag-01JW}{III.? (p.~?)} }
\expandafter\def\csname SRef@courseIII:handle-chains:8.2\endcsname{\href{S6_SERIES_GRIND_III.pdf\#nameddest=tag-01JX}{III.? (p.~?)} }
\expandafter\def\csname SRef@courseIII:pair-creation:1\endcsname{\href{S6_SERIES_GRIND_III.pdf\#nameddest=tag-01LI}{III.? (p.~?)} }
\expandafter\def\csname SRef@courseIII:pair-creation:2\endcsname{\href{S6_SERIES_GRIND_III.pdf\#nameddest=tag-01LJ}{III.? (p.~?)} }
\expandafter\def\csname SRef@courseIII:pair-creation:3\endcsname{\href{S6_SERIES_GRIND_III.pdf\#nameddest=tag-01LK}{III.? (p.~?)} }
\expandafter\def\csname SRef@courseIII:pair-creation:4\endcsname{\href{S6_SERIES_GRIND_III.pdf\#nameddest=tag-01LL}{III.? (p.~?)} }
\expandafter\def\csname SRef@courseIII:pair-creation:5\endcsname{\href{S6_SERIES_GRIND_III.pdf\#nameddest=tag-01LM}{III.? (p.~?)} }
\expandafter\def\csname SRef@courseIII:pair-creation:6\endcsname{\href{S6_SERIES_GRIND_III.pdf\#nameddest=tag-01LN}{III.? (p.~?)} }
\expandafter\def\csname SRef@courseIII:pair-creation:7.1\endcsname{\href{S6_SERIES_GRIND_III.pdf\#nameddest=tag-01LO}{III.? (p.~?)} }
\expandafter\def\csname SRef@courseIII:pair-creation:7.2\endcsname{\href{S6_SERIES_GRIND_III.pdf\#nameddest=tag-01LP}{III.? (p.~?)} }
\expandafter\def\csname SRef@courseIII:low-index-slides:1\endcsname{\href{S6_SERIES_GRIND_III.pdf\#nameddest=tag-01KJ}{III.? (p.~?)} }
\expandafter\def\csname SRef@courseIII:low-index-slides:2\endcsname{\href{S6_SERIES_GRIND_III.pdf\#nameddest=tag-01KK}{III.? (p.~?)} }
\expandafter\def\csname SRef@courseIII:low-index-slides:3\endcsname{\href{S6_SERIES_GRIND_III.pdf\#nameddest=tag-01KL}{III.? (p.~?)} }
\expandafter\def\csname SRef@courseIII:low-index-slides:4\endcsname{\href{S6_SERIES_GRIND_III.pdf\#nameddest=tag-01KM}{III.? (p.~?)} }
\expandafter\def\csname SRef@courseIII:low-index-slides:5\endcsname{\href{S6_SERIES_GRIND_III.pdf\#nameddest=tag-01KN}{III.? (p.~?)} }
\expandafter\def\csname SRef@courseIII:low-index-slides:6\endcsname{\href{S6_SERIES_GRIND_III.pdf\#nameddest=tag-01KO}{III.? (p.~?)} }
\expandafter\def\csname SRef@courseIII:low-index-slides:7\endcsname{\href{S6_SERIES_GRIND_III.pdf\#nameddest=tag-01KP}{III.? (p.~?)} }
\expandafter\def\csname SRef@courseIII:low-index-slides:8\endcsname{\href{S6_SERIES_GRIND_III.pdf\#nameddest=tag-01KQ}{III.? (p.~?)} }
\expandafter\def\csname SRef@courseIII:low-index-slides:9.1\endcsname{\href{S6_SERIES_GRIND_III.pdf\#nameddest=tag-01KR}{III.? (p.~?)} }
\expandafter\def\csname SRef@courseIII:low-index-slides:9.2\endcsname{\href{S6_SERIES_GRIND_III.pdf\#nameddest=tag-01KS}{III.? (p.~?)} }
\expandafter\def\csname SRef@courseIII:adams-arithmetic:1\endcsname{\href{S6_SERIES_GRIND_III.pdf\#nameddest=tag-01HJ}{III.? (p.~?)} }
\expandafter\def\csname SRef@courseIII:adams-arithmetic:2\endcsname{\href{S6_SERIES_GRIND_III.pdf\#nameddest=tag-01HK}{III.? (p.~?)} }
\expandafter\def\csname SRef@courseIII:adams-arithmetic:3\endcsname{\href{S6_SERIES_GRIND_III.pdf\#nameddest=tag-01HL}{III.? (p.~?)} }
\expandafter\def\csname SRef@courseIII:adams-arithmetic:4\endcsname{\href{S6_SERIES_GRIND_III.pdf\#nameddest=tag-01HM}{III.? (p.~?)} }
\expandafter\def\csname SRef@courseIII:adams-arithmetic:5\endcsname{\href{S6_SERIES_GRIND_III.pdf\#nameddest=tag-01HN}{III.? (p.~?)} }
\expandafter\def\csname SRef@courseIII:adams-arithmetic:6\endcsname{\href{S6_SERIES_GRIND_III.pdf\#nameddest=tag-01HO}{III.? (p.~?)} }
\expandafter\def\csname SRef@courseIII:adams-arithmetic:7\endcsname{\href{S6_SERIES_GRIND_III.pdf\#nameddest=tag-01HP}{III.? (p.~?)} }
\expandafter\def\csname SRef@courseIII:adams-arithmetic:8\endcsname{\href{S6_SERIES_GRIND_III.pdf\#nameddest=tag-01HQ}{III.? (p.~?)} }
\expandafter\def\csname SRef@courseIII:adams-arithmetic:9\endcsname{\href{S6_SERIES_GRIND_III.pdf\#nameddest=tag-01HR}{III.? (p.~?)} }
\expandafter\def\csname SRef@courseIII:adams-tower:1\endcsname{\href{S6_SERIES_GRIND_III.pdf\#nameddest=tag-01HS}{III.? (p.~?)} }
\expandafter\def\csname SRef@courseIII:adams-tower:2\endcsname{\href{S6_SERIES_GRIND_III.pdf\#nameddest=tag-01HV}{III.? (p.~?)} }
\expandafter\def\csname SRef@courseIII:adams-tower:3\endcsname{\href{S6_SERIES_GRIND_III.pdf\#nameddest=tag-01HW}{III.? (p.~?)} }
\expandafter\def\csname SRef@courseIII:adams-tower:4\endcsname{\href{S6_SERIES_GRIND_III.pdf\#nameddest=tag-01HX}{III.? (p.~?)} }
\expandafter\def\csname SRef@courseIII:adams-tower:5\endcsname{\href{S6_SERIES_GRIND_III.pdf\#nameddest=tag-01HY}{III.? (p.~?)} }
\expandafter\def\csname SRef@courseIII:adams-tower:6\endcsname{\href{S6_SERIES_GRIND_III.pdf\#nameddest=tag-01HZ}{III.? (p.~?)} }
\expandafter\def\csname SRef@courseIII:adams-tower:7.1\endcsname{\href{S6_SERIES_GRIND_III.pdf\#nameddest=tag-01I0}{III.? (p.~?)} }
\expandafter\def\csname SRef@courseIII:adams-tower:7.2\endcsname{\href{S6_SERIES_GRIND_III.pdf\#nameddest=tag-01I1}{III.? (p.~?)} }
\expandafter\def\csname SRef@courseIII:adams-tower:8\endcsname{\href{S6_SERIES_GRIND_III.pdf\#nameddest=tag-01I2}{III.? (p.~?)} }
\expandafter\def\csname SRef@courseIII:adams-tower:9\endcsname{\href{S6_SERIES_GRIND_III.pdf\#nameddest=tag-01I3}{III.? (p.~?)} }
\expandafter\def\csname SRef@courseIII:adams-tower:10\endcsname{\href{S6_SERIES_GRIND_III.pdf\#nameddest=tag-01HT}{III.? (p.~?)} }
\expandafter\def\csname SRef@courseIII:adams-tower:11\endcsname{\href{S6_SERIES_GRIND_III.pdf\#nameddest=tag-01HU}{III.? (p.~?)} }
\expandafter\def\csname SRef@courseIII:formal-char2:1\endcsname{\href{S6_SERIES_GRIND_III.pdf\#nameddest=tag-01J2}{III.? (p.~?)} }
\expandafter\def\csname SRef@courseIII:formal-char2:2\endcsname{\href{S6_SERIES_GRIND_III.pdf\#nameddest=tag-01J3}{III.? (p.~?)} }
\expandafter\def\csname SRef@courseIII:formal-char2:3\endcsname{\href{S6_SERIES_GRIND_III.pdf\#nameddest=tag-01J4}{III.? (p.~?)} }
\expandafter\def\csname SRef@courseIII:formal-char2:4\endcsname{\href{S6_SERIES_GRIND_III.pdf\#nameddest=tag-01J5}{III.? (p.~?)} }
\expandafter\def\csname SRef@courseIII:formal-char2:4.1\endcsname{\href{S6_SERIES_GRIND_III.pdf\#nameddest=tag-01J6}{III.? (p.~?)} }
\expandafter\def\csname SRef@courseIII:formal-char2:4.2\endcsname{\href{S6_SERIES_GRIND_III.pdf\#nameddest=tag-01J7}{III.? (p.~?)} }
\expandafter\def\csname SRef@courseIII:projective-coop:1\endcsname{\href{S6_SERIES_GRIND_III.pdf\#nameddest=tag-01M0}{III.? (p.~?)} }
\expandafter\def\csname SRef@courseIII:projective-coop:2\endcsname{\href{S6_SERIES_GRIND_III.pdf\#nameddest=tag-01M1}{III.? (p.~?)} }
\expandafter\def\csname SRef@courseIII:projective-coop:3\endcsname{\href{S6_SERIES_GRIND_III.pdf\#nameddest=tag-01M2}{III.? (p.~?)} }
\expandafter\def\csname SRef@courseIII:projective-coop:4\endcsname{\href{S6_SERIES_GRIND_III.pdf\#nameddest=tag-01M3}{III.? (p.~?)} }
\expandafter\def\csname SRef@courseIII:projective-coop:5\endcsname{\href{S6_SERIES_GRIND_III.pdf\#nameddest=tag-01M4}{III.? (p.~?)} }
\expandafter\def\csname SRef@courseIII:projective-coop:6.1\endcsname{\href{S6_SERIES_GRIND_III.pdf\#nameddest=tag-01M5}{III.? (p.~?)} }
\expandafter\def\csname SRef@courseIII:projective-coop:6.2\endcsname{\href{S6_SERIES_GRIND_III.pdf\#nameddest=tag-01M6}{III.? (p.~?)} }
\expandafter\def\csname SRef@courseIII:actual-degree8:1\endcsname{\href{S6_SERIES_GRIND_III.pdf\#nameddest=tag-01H0}{III.? (p.~?)} }
\expandafter\def\csname SRef@courseIII:actual-degree8:2\endcsname{\href{S6_SERIES_GRIND_III.pdf\#nameddest=tag-01H1}{III.? (p.~?)} }
\expandafter\def\csname SRef@courseIII:actual-degree8:3\endcsname{\href{S6_SERIES_GRIND_III.pdf\#nameddest=tag-01H2}{III.? (p.~?)} }
\expandafter\def\csname SRef@courseIII:actual-degree8:4\endcsname{\href{S6_SERIES_GRIND_III.pdf\#nameddest=tag-01H3}{III.? (p.~?)} }
\expandafter\def\csname SRef@courseIII:actual-degree8:5\endcsname{\href{S6_SERIES_GRIND_III.pdf\#nameddest=tag-01H4}{III.? (p.~?)} }
\expandafter\def\csname SRef@courseIII:actual-degree8:6\endcsname{\href{S6_SERIES_GRIND_III.pdf\#nameddest=tag-01H5}{III.? (p.~?)} }
\expandafter\def\csname SRef@courseIII:actual-degree8:7\endcsname{\href{S6_SERIES_GRIND_III.pdf\#nameddest=tag-01H6}{III.? (p.~?)} }
\expandafter\def\csname SRef@courseIII:actual-degree8:8\endcsname{\href{S6_SERIES_GRIND_III.pdf\#nameddest=tag-01H7}{III.? (p.~?)} }
\expandafter\def\csname SRef@courseIII:actual-degree8:9.1\endcsname{\href{S6_SERIES_GRIND_III.pdf\#nameddest=tag-01H8}{III.? (p.~?)} }
\expandafter\def\csname SRef@courseIII:actual-degree8:9.2\endcsname{\href{S6_SERIES_GRIND_III.pdf\#nameddest=tag-01H9}{III.? (p.~?)} }
\expandafter\def\csname SRef@courseIII:pt-relative-embedding\endcsname{\href{S6_SERIES_GRIND_III.pdf\#nameddest=tag-01ME}{III.? (p.~?)} }
\expandafter\def\csname SRef@courseIII:pt-relative-normal-frame\endcsname{\href{S6_SERIES_GRIND_III.pdf\#nameddest=tag-01MF}{III.? (p.~?)} }
\expandafter\def\csname SRef@courseIII:pt-normal-inverse-comparison\endcsname{\href{S6_SERIES_GRIND_III.pdf\#nameddest=tag-01MD}{III.? (p.~?)} }
\expandafter\def\csname SRef@courseIII:pt-abstract-framed-comparison\endcsname{\href{S6_SERIES_GRIND_III.pdf\#nameddest=tag-01MC}{III.? (p.~?)} }
\expandafter\def\csname SRef@courseIII:actual-resolution:1\endcsname{\href{S6_SERIES_GRIND_III.pdf\#nameddest=tag-01HA}{III.? (p.~?)} }
\expandafter\def\csname SRef@courseIII:actual-resolution:2\endcsname{\href{S6_SERIES_GRIND_III.pdf\#nameddest=tag-01HB}{III.? (p.~?)} }
\expandafter\def\csname SRef@courseIII:actual-resolution:3\endcsname{\href{S6_SERIES_GRIND_III.pdf\#nameddest=tag-01HC}{III.? (p.~?)} }
\expandafter\def\csname SRef@courseIII:actual-resolution:4\endcsname{\href{S6_SERIES_GRIND_III.pdf\#nameddest=tag-01HD}{III.? (p.~?)} }
\expandafter\def\csname SRef@courseIII:actual-resolution:5\endcsname{\href{S6_SERIES_GRIND_III.pdf\#nameddest=tag-01HE}{III.? (p.~?)} }
\expandafter\def\csname SRef@courseIII:actual-resolution:6\endcsname{\href{S6_SERIES_GRIND_III.pdf\#nameddest=tag-01HF}{III.? (p.~?)} }
\expandafter\def\csname SRef@courseIII:actual-resolution:7\endcsname{\href{S6_SERIES_GRIND_III.pdf\#nameddest=tag-01HG}{III.? (p.~?)} }
\expandafter\def\csname SRef@courseIII:actual-resolution:8.1\endcsname{\href{S6_SERIES_GRIND_III.pdf\#nameddest=tag-01HH}{III.? (p.~?)} }
\expandafter\def\csname SRef@courseIII:actual-resolution:8.2\endcsname{\href{S6_SERIES_GRIND_III.pdf\#nameddest=tag-01HI}{III.? (p.~?)} }
\expandafter\def\csname SRef@courseIII:milnor-cobar:1.1\endcsname{\href{S6_SERIES_GRIND_III.pdf\#nameddest=tag-01L1}{III.? (p.~?)} }
\expandafter\def\csname SRef@courseIII:milnor-cobar:6.1\endcsname{\href{S6_SERIES_GRIND_III.pdf\#nameddest=tag-01L2}{III.? (p.~?)} }
\expandafter\def\csname SRef@courseIII:milnor-cobar:7.1\endcsname{\href{S6_SERIES_GRIND_III.pdf\#nameddest=tag-01L3}{III.? (p.~?)} }
\expandafter\def\csname SRef@courseIII:milnor-cobar:7.2\endcsname{\href{S6_SERIES_GRIND_III.pdf\#nameddest=tag-01L4}{III.? (p.~?)} }
\expandafter\def\csname SRef@courseIII:odd-coop:1\endcsname{\href{S6_SERIES_GRIND_III.pdf\#nameddest=tag-01LA}{III.? (p.~?)} }
\expandafter\def\csname SRef@courseIII:odd-coop:2\endcsname{\href{S6_SERIES_GRIND_III.pdf\#nameddest=tag-01LB}{III.? (p.~?)} }
\expandafter\def\csname SRef@courseIII:odd-coop:3\endcsname{\href{S6_SERIES_GRIND_III.pdf\#nameddest=tag-01LC}{III.? (p.~?)} }
\expandafter\def\csname SRef@courseIII:odd-coop:4\endcsname{\href{S6_SERIES_GRIND_III.pdf\#nameddest=tag-01LD}{III.? (p.~?)} }
\expandafter\def\csname SRef@courseIII:odd-coop:5\endcsname{\href{S6_SERIES_GRIND_III.pdf\#nameddest=tag-01LE}{III.? (p.~?)} }
\expandafter\def\csname SRef@courseIII:odd-coop:6\endcsname{\href{S6_SERIES_GRIND_III.pdf\#nameddest=tag-01LF}{III.? (p.~?)} }
\expandafter\def\csname SRef@courseIII:odd-coop:7.1\endcsname{\href{S6_SERIES_GRIND_III.pdf\#nameddest=tag-01LG}{III.? (p.~?)} }
\expandafter\def\csname SRef@courseIII:odd-coop:7.2\endcsname{\href{S6_SERIES_GRIND_III.pdf\#nameddest=tag-01LH}{III.? (p.~?)} }
\expandafter\def\csname SRef@courseIII:coefficient-rings:1\endcsname{\href{S6_SERIES_GRIND_III.pdf\#nameddest=tag-01IM}{III.? (p.~?)} }
\expandafter\def\csname SRef@courseIII:coefficient-rings:2\endcsname{\href{S6_SERIES_GRIND_III.pdf\#nameddest=tag-01IN}{III.? (p.~?)} }
\expandafter\def\csname SRef@courseIII:coefficient-rings:3\endcsname{\href{S6_SERIES_GRIND_III.pdf\#nameddest=tag-01IO}{III.? (p.~?)} }
\expandafter\def\csname SRef@courseIII:coefficient-rings:4\endcsname{\href{S6_SERIES_GRIND_III.pdf\#nameddest=tag-01IP}{III.? (p.~?)} }
\expandafter\def\csname SRef@courseIII:coefficient-rings:5\endcsname{\href{S6_SERIES_GRIND_III.pdf\#nameddest=tag-01IQ}{III.? (p.~?)} }
\expandafter\def\csname SRef@courseIII:coefficient-rings:6\endcsname{\href{S6_SERIES_GRIND_III.pdf\#nameddest=tag-01IR}{III.? (p.~?)} }
\expandafter\def\csname SRef@courseIII:coefficient-rings:7.1\endcsname{\href{S6_SERIES_GRIND_III.pdf\#nameddest=tag-01IS}{III.? (p.~?)} }
\expandafter\def\csname SRef@courseIII:coefficient-rings:7.2\endcsname{\href{S6_SERIES_GRIND_III.pdf\#nameddest=tag-01IT}{III.? (p.~?)} }
\expandafter\def\csname SRef@courseIII:duality:2.1\endcsname{\href{S6_SERIES_GRIND_III.pdf\#nameddest=tag-01IU}{III.? (p.~?)} }
\expandafter\def\csname SRef@courseIII:duality:2.2\endcsname{\href{S6_SERIES_GRIND_III.pdf\#nameddest=tag-01IV}{III.? (p.~?)} }
\expandafter\def\csname SRef@courseIII:duality:2.3\endcsname{\href{S6_SERIES_GRIND_III.pdf\#nameddest=tag-01IW}{III.? (p.~?)} }
\expandafter\def\csname SRef@courseIII:duality:3.1\endcsname{\href{S6_SERIES_GRIND_III.pdf\#nameddest=tag-01IX}{III.? (p.~?)} }
\expandafter\def\csname SRef@courseIII:duality:3.2\endcsname{\href{S6_SERIES_GRIND_III.pdf\#nameddest=tag-01IY}{III.? (p.~?)} }
\expandafter\def\csname SRef@courseIII:duality:3.3\endcsname{\href{S6_SERIES_GRIND_III.pdf\#nameddest=tag-01IZ}{III.? (p.~?)} }
\expandafter\def\csname SRef@courseIII:duality:4.1\endcsname{\href{S6_SERIES_GRIND_III.pdf\#nameddest=tag-01J0}{III.? (p.~?)} }
\expandafter\def\csname SRef@courseIII:duality:4.2\endcsname{\href{S6_SERIES_GRIND_III.pdf\#nameddest=tag-01J1}{III.? (p.~?)} }
\expandafter\def\csname SRef@night:P1.6\endcsname{\href{S6_SERIES_GRIND_V.pdf\#nameddest=tag-02QX}{V.? (p.~?)} }
\expandafter\def\csname SRef@night:P1.14\endcsname{\href{S6_SERIES_GRIND_V.pdf\#nameddest=tag-02QN}{V.? (p.~?)} }
\expandafter\def\csname SRef@night:P5.7\endcsname{\href{S6_SERIES_GRIND_I.pdf\#nameddest=tag-02YS}{I.A.1 (p.~342)} }
\expandafter\def\csname SRef@ari:A.1.1\endcsname{\href{S6_SERIES_GRIND_I.pdf\#nameddest=tag-03J1}{I.? (p.~?)} }
\expandafter\def\csname SRef@ari:A.1.2\endcsname{\href{S6_SERIES_GRIND_I.pdf\#nameddest=tag-03J2}{I.? (p.~?)} }
\expandafter\def\csname SRef@ari:A.1.3\endcsname{\href{S6_SERIES_GRIND_I.pdf\#nameddest=tag-03J3}{I.? (p.~?)} }
\expandafter\def\csname SRef@ari:A.1.4\endcsname{\href{S6_SERIES_GRIND_I.pdf\#nameddest=tag-03J4}{I.? (p.~?)} }
\expandafter\def\csname SRef@ari:A.2.1\endcsname{\href{S6_SERIES_GRIND_I.pdf\#nameddest=tag-03J5}{I.? (p.~?)} }
\expandafter\def\csname SRef@ari:A.2.2\endcsname{\href{S6_SERIES_GRIND_I.pdf\#nameddest=tag-03J6}{I.? (p.~?)} }
\expandafter\def\csname SRef@ari:A.2.3\endcsname{\href{S6_SERIES_GRIND_I.pdf\#nameddest=tag-03J7}{I.? (p.~?)} }
\expandafter\def\csname SRef@ari:A.2.4\endcsname{\href{S6_SERIES_GRIND_I.pdf\#nameddest=tag-03J8}{I.? (p.~?)} }
\expandafter\def\csname SRef@ari:A.2.5\endcsname{\href{S6_SERIES_GRIND_I.pdf\#nameddest=tag-03J9}{I.? (p.~?)} }
\expandafter\def\csname SRef@ari:A.3.1\endcsname{\href{S6_SERIES_GRIND_I.pdf\#nameddest=tag-03JA}{I.? (p.~?)} }
\expandafter\def\csname SRef@ari:A.3.2\endcsname{\href{S6_SERIES_GRIND_I.pdf\#nameddest=tag-03JB}{I.? (p.~?)} }
\expandafter\def\csname SRef@ari:A.4.1\endcsname{\href{S6_SERIES_GRIND_I.pdf\#nameddest=tag-03JC}{I.? (p.~?)} }
\expandafter\def\csname SRef@ari:A.4.2\endcsname{\href{S6_SERIES_GRIND_I.pdf\#nameddest=tag-03JD}{I.? (p.~?)} }
\expandafter\def\csname SRef@ari:A.4.3\endcsname{\href{S6_SERIES_GRIND_I.pdf\#nameddest=tag-03JE}{I.? (p.~?)} }
\expandafter\def\csname SRef@ari:A.4.4\endcsname{\href{S6_SERIES_GRIND_I.pdf\#nameddest=tag-03JF}{I.? (p.~?)} }
\expandafter\def\csname SRef@ari:A.4.5\endcsname{\href{S6_SERIES_GRIND_I.pdf\#nameddest=tag-03JG}{I.? (p.~?)} }
\expandafter\def\csname SRef@ari:A.4.6\endcsname{\href{S6_SERIES_GRIND_I.pdf\#nameddest=tag-03JH}{I.? (p.~?)} }
\expandafter\def\csname SRef@ari:A.4.7\endcsname{\href{S6_SERIES_GRIND_I.pdf\#nameddest=tag-03JI}{I.? (p.~?)} }
\expandafter\def\csname SRef@ari:A.4.8\endcsname{\href{S6_SERIES_GRIND_I.pdf\#nameddest=tag-03JJ}{I.? (p.~?)} }
\expandafter\def\csname SRef@ari:A.5.1\endcsname{\href{S6_SERIES_GRIND_I.pdf\#nameddest=tag-03JK}{I.? (p.~?)} }
\expandafter\def\csname SRef@ari:A.5.2\endcsname{\href{S6_SERIES_GRIND_I.pdf\#nameddest=tag-03JL}{I.? (p.~?)} }
\expandafter\def\csname SRef@ari:A.5.3\endcsname{\href{S6_SERIES_GRIND_I.pdf\#nameddest=tag-03JM}{I.? (p.~?)} }
\expandafter\def\csname SRef@ari:A.5.4\endcsname{\href{S6_SERIES_GRIND_I.pdf\#nameddest=tag-03JN}{I.? (p.~?)} }
\expandafter\def\csname SRef@ari:A.5.5\endcsname{\href{S6_SERIES_GRIND_I.pdf\#nameddest=tag-03JO}{I.? (p.~?)} }
\expandafter\def\csname SRef@ari:A.5.6\endcsname{\href{S6_SERIES_GRIND_I.pdf\#nameddest=tag-03JP}{I.? (p.~?)} }
\expandafter\def\csname SRef@ari:A.5.8\endcsname{\href{S6_SERIES_GRIND_I.pdf\#nameddest=tag-03JR}{I.? (p.~?)} }
\expandafter\def\csname SRef@ari:A.5.7\endcsname{\href{S6_SERIES_GRIND_I.pdf\#nameddest=tag-03JQ}{I.? (p.~?)} }
\expandafter\def\csname SRef@series:I:main-theorems\endcsname{\href{S6_SERIES_GRIND_I.pdf\#nameddest=tag-03JS}{I.1.2 (p.~6)} }
\expandafter\def\csname SRef@series:I:A\endcsname{\href{S6_SERIES_GRIND_I.pdf\#nameddest=tag-03JT}{I.A (p.~6)} }
\expandafter\def\csname SRef@series:I:B\endcsname{\href{S6_SERIES_GRIND_I.pdf\#nameddest=tag-03JU}{I.B (p.~8)} }
\expandafter\def\csname SRef@series:I:C\endcsname{\href{S6_SERIES_GRIND_I.pdf\#nameddest=tag-03JV}{I.C (p.~9)} }
\expandafter\def\csname SRef@series:I:D\endcsname{\href{S6_SERIES_GRIND_I.pdf\#nameddest=tag-03JW}{I.D (p.~10)} }
\expandafter\def\csname SRef@series:I:E\endcsname{\href{S6_SERIES_GRIND_I.pdf\#nameddest=tag-03JX}{I.E (p.~11)} }
\expandafter\def\csname SRef@series:I:F\endcsname{\href{S6_SERIES_GRIND_I.pdf\#nameddest=tag-03JY}{I.F (p.~12)} }
\expandafter\def\csname SRef@series:I:G\endcsname{\href{S6_SERIES_GRIND_I.pdf\#nameddest=tag-03JZ}{I.G (p.~12)} }
\expandafter\def\csname SRef@series:I:outlines\endcsname{\href{S6_SERIES_GRIND_I.pdf\#nameddest=tag-03K0}{I.1.3 (p.~15)} }
\expandafter\def\csname SRef@series:II:main-theorems\endcsname{\hyperlink{tag-03K1}{II.\ref*{series:II:main-theorems}}}
\expandafter\def\csname SRef@series:II:A\endcsname{\hyperlink{tag-03K2}{II.\ref*{series:II:A}}}
\expandafter\def\csname SRef@series:II:B\endcsname{\hyperlink{tag-03K3}{II.\ref*{series:II:B}}}
\expandafter\def\csname SRef@series:II:C\endcsname{\hyperlink{tag-03K4}{II.\ref*{series:II:C}}}
\expandafter\def\csname SRef@series:II:D\endcsname{\hyperlink{tag-03K5}{II.\ref*{series:II:D}}}
\expandafter\def\csname SRef@series:II:E\endcsname{\hyperlink{tag-03K6}{II.\ref*{series:II:E}}}
\expandafter\def\csname SRef@series:II:F\endcsname{\hyperlink{tag-03K7}{II.\ref*{series:II:F}}}
\expandafter\def\csname SRef@series:II:G\endcsname{\hyperlink{tag-03K8}{II.\ref*{series:II:G}}}
\expandafter\def\csname SRef@series:II:outline\endcsname{\hyperlink{tag-03K9}{II.\ref*{series:II:outline}}}
\expandafter\def\csname SRef@series:III:main-theorems\endcsname{\href{S6_SERIES_GRIND_III.pdf\#nameddest=tag-03KA}{III.1.2 (p.~3)} }
\expandafter\def\csname SRef@series:III:A\endcsname{\href{S6_SERIES_GRIND_III.pdf\#nameddest=tag-03KB}{III.A (p.~3)} }
\expandafter\def\csname SRef@series:III:B\endcsname{\href{S6_SERIES_GRIND_III.pdf\#nameddest=tag-03KC}{III.B (p.~3)} }
\expandafter\def\csname SRef@series:III:C\endcsname{\href{S6_SERIES_GRIND_III.pdf\#nameddest=tag-03KD}{III.C (p.~4)} }
\expandafter\def\csname SRef@series:III:D\endcsname{\href{S6_SERIES_GRIND_III.pdf\#nameddest=tag-03KE}{III.D (p.~4)} }
\expandafter\def\csname SRef@series:III:E\endcsname{\href{S6_SERIES_GRIND_III.pdf\#nameddest=tag-03KF}{III.E (p.~4)} }
\expandafter\def\csname SRef@series:III:F\endcsname{\href{S6_SERIES_GRIND_III.pdf\#nameddest=tag-03KG}{III.F (p.~5)} }
\expandafter\def\csname SRef@series:III:G\endcsname{\href{S6_SERIES_GRIND_III.pdf\#nameddest=tag-03KH}{III.G (p.~5)} }
\expandafter\def\csname SRef@series:IV:main-theorems\endcsname{\href{S6_SERIES_GRIND_IV.pdf\#nameddest=tag-03KI}{IV.1.2 (p.~3)} }
\expandafter\def\csname SRef@series:IV:A\endcsname{\href{S6_SERIES_GRIND_IV.pdf\#nameddest=tag-03KJ}{IV.A (p.~3)} }
\expandafter\def\csname SRef@series:IV:B\endcsname{\href{S6_SERIES_GRIND_IV.pdf\#nameddest=tag-03KK}{IV.B (p.~4)} }
\expandafter\def\csname SRef@series:IV:C\endcsname{\href{S6_SERIES_GRIND_IV.pdf\#nameddest=tag-03KL}{IV.C (p.~5)} }
\expandafter\def\csname SRef@series:IV:D\endcsname{\href{S6_SERIES_GRIND_IV.pdf\#nameddest=tag-03KM}{IV.D (p.~6)} }
\expandafter\def\csname SRef@series:IV:E\endcsname{\href{S6_SERIES_GRIND_IV.pdf\#nameddest=tag-03KN}{IV.E (p.~7)} }
\expandafter\def\csname SRef@series:IV:F\endcsname{\href{S6_SERIES_GRIND_IV.pdf\#nameddest=tag-03KO}{IV.F (p.~8)} }
\expandafter\def\csname SRef@series:IV:G\endcsname{\href{S6_SERIES_GRIND_IV.pdf\#nameddest=tag-03KP}{IV.G (p.~8)} }
\expandafter\def\csname SRef@series:V:main-theorems\endcsname{\href{S6_SERIES_GRIND_V.pdf\#nameddest=tag-03KQ}{V.1.2 (p.~3)} }
\expandafter\def\csname SRef@series:V:A\endcsname{\href{S6_SERIES_GRIND_V.pdf\#nameddest=tag-03KR}{V.A (p.~3)} }
\expandafter\def\csname SRef@series:V:B\endcsname{\href{S6_SERIES_GRIND_V.pdf\#nameddest=tag-03KS}{V.B (p.~3)} }
\expandafter\def\csname SRef@series:V:C\endcsname{\href{S6_SERIES_GRIND_V.pdf\#nameddest=tag-03KT}{V.C (p.~4)} }
\expandafter\def\csname SRef@series:V:D\endcsname{\href{S6_SERIES_GRIND_V.pdf\#nameddest=tag-03KU}{V.D (p.~4)} }
\expandafter\def\csname SRef@series:V:E\endcsname{\href{S6_SERIES_GRIND_V.pdf\#nameddest=tag-03KV}{V.E (p.~5)} }
\expandafter\def\csname SRef@series:V:F\endcsname{\href{S6_SERIES_GRIND_V.pdf\#nameddest=tag-03KW}{V.F (p.~5)} }
\expandafter\def\csname SRef@layout:V:fresh-gauges\endcsname{\href{S6_SERIES_GRIND_V.pdf\#nameddest=tag-03KX}{V.{{$\langle 2\rangle $1.}} (p.~126)} }
\expandafter\def\csname SRef@layout:V:fresh-minimum\endcsname{\href{S6_SERIES_GRIND_V.pdf\#nameddest=tag-03KY}{V.{{$\langle 1\rangle $4.}} (p.~130)} }
\expandafter\def\csname SRef@layout:V:fresh-bounds\endcsname{\href{S6_SERIES_GRIND_V.pdf\#nameddest=tag-03KZ}{V.{1.18.4} (p.~130)} }
\expandafter\def\csname SRef@layout:V:fresh-marking\endcsname{\href{S6_SERIES_GRIND_V.pdf\#nameddest=tag-03L0}{V.{{$\langle 1\rangle $6.}} (p.~132)} }
\expandafter\def\csname SRef@layout:V:fresh-splitting\endcsname{\href{S6_SERIES_GRIND_V.pdf\#nameddest=tag-03L1}{V.{{$\langle 1\rangle $7.}} (p.~133)} }
\expandafter\def\csname SRef@iv:source_transports\endcsname{\href{S6_SERIES_GRIND_IV.pdf\#nameddest=tag-03L2}{IV.3.27 (p.~150)} }
\expandafter\def\csname SRef@inv:section:coefficient-atlas\endcsname{\href{S6_SERIES_GRIND_I.pdf\#nameddest=tag-03L3}{I.D (p.~374)} }
\expandafter\def\csname SRef@inv:eq:atlas-theta-factor\endcsname{\href{S6_SERIES_GRIND_I.pdf\#nameddest=tag-03L4}{I.105 (p.~375)} }
\expandafter\def\csname SRef@inv:eq:atlas-semiflat\endcsname{\href{S6_SERIES_GRIND_I.pdf\#nameddest=tag-03L5}{I.106 (p.~375)} }
\expandafter\def\csname SRef@inv:eq:atlas-affine-lift\endcsname{\href{S6_SERIES_GRIND_I.pdf\#nameddest=tag-03L6}{I.107 (p.~376)} }
\expandafter\def\csname SRef@inv:eq:atlas-cusp-factor\endcsname{\href{S6_SERIES_GRIND_I.pdf\#nameddest=tag-03L7}{I.108 (p.~377)} }
\expandafter\def\csname SRef@inv:eq:atlas-cusp-unit\endcsname{\href{S6_SERIES_GRIND_I.pdf\#nameddest=tag-03L8}{I.109 (p.~377)} }
\expandafter\def\csname SRef@inv:eq:atlas-connection-inverse\endcsname{\href{S6_SERIES_GRIND_I.pdf\#nameddest=tag-03L9}{I.110 (p.~377)} }
\expandafter\def\csname SRef@inv:eq:atlas-cusp-connection\endcsname{\href{S6_SERIES_GRIND_I.pdf\#nameddest=tag-03LA}{I.111 (p.~377)} }
\expandafter\def\csname SRef@inv:eq:atlas-finite-connection\endcsname{\href{S6_SERIES_GRIND_I.pdf\#nameddest=tag-03LB}{I.112 (p.~378)} }
\expandafter\def\csname SRef@inv:eq:atlas-averaged-connection\endcsname{\href{S6_SERIES_GRIND_I.pdf\#nameddest=tag-03LC}{I.113 (p.~378)} }
\expandafter\def\csname SRef@inv:eq:atlas-cocycle\endcsname{\href{S6_SERIES_GRIND_I.pdf\#nameddest=tag-03LD}{I.114 (p.~378)} }
\expandafter\def\csname SRef@inv:eq:atlas-weights\endcsname{\href{S6_SERIES_GRIND_I.pdf\#nameddest=tag-03LE}{I.115 (p.~379)} }
\expandafter\def\csname SRef@inv:eq:atlas-weight-derivative\endcsname{\href{S6_SERIES_GRIND_I.pdf\#nameddest=tag-03LF}{I.116 (p.~379)} }
\expandafter\def\csname SRef@inv:eq:atlas-final-forms\endcsname{\href{S6_SERIES_GRIND_I.pdf\#nameddest=tag-03LG}{I.117 (p.~379)} }
\expandafter\def\csname SRef@inv:eq:atlas-final-coefficients\endcsname{\href{S6_SERIES_GRIND_I.pdf\#nameddest=tag-03LH}{I.118 (p.~380)} }
\expandafter\def\csname SRef@inv:C3.8:operator-comparison\endcsname{\href{S6_SERIES_GRIND_I.pdf\#nameddest=tag-03LI}{I.D.2 (p.~381)} }
\expandafter\def\csname SRef@iv:open:K111:original_auxiliary_arrays\endcsname{\href{S6_SERIES_GRIND_IV.pdf\#nameddest=tag-03LJ}{IV.12.1 (p.~1396)} }
\expandafter\def\csname SRef@maps:proof:D.OS.refined\endcsname{\href{S6_SERIES_GRIND_III.pdf\#nameddest=tag-03LK}{III.5.2 (p.~207)} }
\expandafter\def\csname SRef@maps:eq:D.OR1\endcsname{\href{S6_SERIES_GRIND_III.pdf\#nameddest=tag-03LL}{III.{D.OR1} (p.~207)} }
\expandafter\def\csname SRef@maps:eq:D.OR2\endcsname{\href{S6_SERIES_GRIND_III.pdf\#nameddest=tag-03LM}{III.{D.OR2} (p.~207)} }
\expandafter\def\csname SRef@maps:eq:D.OR3\endcsname{\href{S6_SERIES_GRIND_III.pdf\#nameddest=tag-03LN}{III.{D.OR3} (p.~207)} }
\expandafter\def\csname SRef@maps:eq:D.OR4\endcsname{\href{S6_SERIES_GRIND_III.pdf\#nameddest=tag-03LO}{III.{D.OR4} (p.~207)} }
\expandafter\def\csname SRef@maps:eq:D.OR5\endcsname{\href{S6_SERIES_GRIND_III.pdf\#nameddest=tag-03LP}{III.{D.OR5} (p.~208)} }
\expandafter\def\csname SRef@maps:eq:D.OR6\endcsname{\href{S6_SERIES_GRIND_III.pdf\#nameddest=tag-03LQ}{III.{D.OR6} (p.~208)} }
\expandafter\def\csname SRef@maps:eq:D.OR7\endcsname{\href{S6_SERIES_GRIND_III.pdf\#nameddest=tag-03LR}{III.{D.OR7} (p.~208)} }
\expandafter\def\csname SRef@maps:eq:D.OR8\endcsname{\href{S6_SERIES_GRIND_III.pdf\#nameddest=tag-03LS}{III.{D.OR8} (p.~208)} }
\expandafter\def\csname SRef@maps:eq:D.OR9\endcsname{\href{S6_SERIES_GRIND_III.pdf\#nameddest=tag-03LT}{III.{D.OR9} (p.~208)} }
\expandafter\def\csname SRef@maps:eq:D.OR10\endcsname{\href{S6_SERIES_GRIND_III.pdf\#nameddest=tag-03LU}{III.{D.OR10} (p.~209)} }
\expandafter\def\csname SRef@maps:eq:D.OR11\endcsname{\href{S6_SERIES_GRIND_III.pdf\#nameddest=tag-03LV}{III.{D.OR11} (p.~209)} }
\expandafter\def\csname SRef@maps:eq:D.OR12\endcsname{\href{S6_SERIES_GRIND_III.pdf\#nameddest=tag-03LW}{III.{D.OR12} (p.~210)} }
\expandafter\def\csname SRef@maps:eq:D.OR13\endcsname{\href{S6_SERIES_GRIND_III.pdf\#nameddest=tag-03LX}{III.{D.OR13} (p.~210)} }
\expandafter\def\csname SRef@maps:eq:D.OR14\endcsname{\href{S6_SERIES_GRIND_III.pdf\#nameddest=tag-03LY}{III.{D.OR14} (p.~211)} }
\expandafter\def\csname SRef@inv:E2.5\endcsname{\hyperlink{tag-03LZ}{II.\ref*{inv:E2.5}}}
\expandafter\def\csname SRef@invII:eq:E12-node\endcsname{\hyperlink{tag-03M0}{II.\ref*{invII:eq:E12-node}}}
\expandafter\def\csname SRef@invII:eq:E12-jets\endcsname{\hyperlink{tag-03M1}{II.\ref*{invII:eq:E12-jets}}}
\expandafter\def\csname SRef@invII:eq:E12-open\endcsname{\hyperlink{tag-03M2}{II.\ref*{invII:eq:E12-open}}}
\expandafter\def\csname SBack@cdp:st:A.1.1\endcsname{\hyperlink{step-II-A-1}{proves step 1 of Theorem A}}
\expandafter\def\csname SBack@cdp:st:B1\endcsname{\hyperlink{step-II-A-2}{proves step 2 of Theorem A}}
\expandafter\def\csname SBack@cdp:st:A.4.7\endcsname{\hyperlink{step-II-A-3}{proves step 3 of Theorem A}}
\expandafter\def\csname SBack@cdp:st:A.4.9\endcsname{\hyperlink{step-II-A-3}{proves step 3 of Theorem A}}
\expandafter\def\csname SBack@cdp:st:C4f\endcsname{\hyperlink{step-II-A-4}{proves step 4 of Theorem A}}
\expandafter\def\csname SBack@cdp:st:E12\endcsname{\hyperlink{step-II-A-5}{proves step 5 of Theorem A}}
\expandafter\def\csname SBack@cdp:st:E17\endcsname{\hyperlink{step-II-A-5}{proves step 5 of Theorem A}}
\expandafter\def\csname SBack@cdp:st:D.1.2\endcsname{\hyperlink{step-II-B-1}{proves step 1 of Theorem B}}
\expandafter\def\csname SBack@cdp:st:C.3.1\endcsname{\hyperlink{step-II-B-2}{proves step 2 of Theorem B}}
\expandafter\def\csname SBack@cdp:st:C.3.7\endcsname{\hyperlink{step-II-B-3}{proves step 3 of Theorem B}}
\expandafter\def\csname SBack@cdp:st:C.3.8\endcsname{\hyperlink{step-II-B-3}{proves step 3 of Theorem B}}
\expandafter\def\csname SBack@cdp:st:C.3.9\endcsname{\hyperlink{step-II-B-3}{proves step 3 of Theorem B}}
\expandafter\def\csname SBack@cdp:st:C.3.10\endcsname{\hyperlink{step-II-B-3}{proves step 3 of Theorem B}}
\expandafter\def\csname SBack@cdp:st:C.3.12\endcsname{\hyperlink{step-II-B-3}{proves step 3 of Theorem B}}
\expandafter\def\csname SBack@cdp:st:C.3.6\endcsname{\hyperlink{step-II-B-4}{proves step 4 of Theorem B}}
\expandafter\def\csname SBack@cdp:st:D.2.7\endcsname{\hyperlink{step-II-B-4}{proves step 4 of Theorem B}; \hyperlink{step-II-C-6}{proves step 6 of Theorem C}}
\expandafter\def\csname SBack@night:P4.58\endcsname{\hyperlink{step-II-C-1}{proves step 1 of Theorem C}}
\expandafter\def\csname SBack@cdp:st:C.2.MH.2\endcsname{\hyperlink{step-II-C-1}{proves step 1 of Theorem C}}
\expandafter\def\csname SBack@cdp:st:C.2.INT.1\endcsname{\hyperlink{step-II-C-2}{proves step 2 of Theorem C}}
\expandafter\def\csname SBack@cdp:st:C.2.INT.2\endcsname{\hyperlink{step-II-C-2}{proves step 2 of Theorem C}}
\expandafter\def\csname SBack@cdp:st:C.2.INT.3\endcsname{\hyperlink{step-II-C-3}{proves step 3 of Theorem C}}
\expandafter\def\csname SBack@cdp:st:C.2.G.1\endcsname{\hyperlink{step-II-C-4}{proves step 4 of Theorem C}}
\expandafter\def\csname SBack@cdp:st:D.2.3\endcsname{\hyperlink{step-II-C-5}{proves step 5 of Theorem C}}
\expandafter\def\csname SBack@cdp:st:D.2.5\endcsname{\hyperlink{step-II-C-5}{proves step 5 of Theorem C}}
\expandafter\def\csname SBack@cdp:st:C.2.UN.1\endcsname{\hyperlink{step-II-D-1}{proves step 1 of Theorem D}; \hyperlink{step-II-D-2}{proves step 2 of Theorem D}}
\expandafter\def\csname SBack@cdp:st:C.2.UN.4\endcsname{\hyperlink{step-II-D-3}{proves step 3 of Theorem D}}
\expandafter\def\csname SBack@cdp:st:C.2.UN.14\endcsname{\hyperlink{step-II-D-4}{proves step 4 of Theorem D}}
\expandafter\def\csname SBack@cdp:st:C.2.UN.2\endcsname{\hyperlink{step-II-D-5}{proves step 5 of Theorem D}}
\expandafter\def\csname SBack@cdp:st:C5b\endcsname{\hyperlink{step-II-E-1}{proves step 1 of Theorem E}; \hyperlink{step-II-E-2}{proves step 2 of Theorem E}}
\expandafter\def\csname SBack@cdp:st:C5c\endcsname{\hyperlink{step-II-E-3}{proves step 3 of Theorem E}; \hyperlink{step-II-E-4}{proves step 4 of Theorem E}}
\expandafter\def\csname SBack@cdp:st:E8\endcsname{\hyperlink{step-II-E-5}{proves step 5 of Theorem E}}
\expandafter\def\csname SBack@cdp:st:E18\endcsname{\hyperlink{step-II-E-5}{proves step 5 of Theorem E}}
\expandafter\def\csname SBack@cdp:st:E5\endcsname{\hyperlink{step-II-F-1}{proves step 1 of Theorem F}}
\expandafter\def\csname SBack@cdp:st:E6\endcsname{\hyperlink{step-II-F-2}{proves step 2 of Theorem F}}
\expandafter\def\csname SBack@cdp:st:E13\endcsname{\hyperlink{step-II-F-3}{proves step 3 of Theorem F}}
\expandafter\def\csname SBack@cdp:st:E14\endcsname{\hyperlink{step-II-F-3}{proves step 3 of Theorem F}}
\expandafter\def\csname SBack@cdp:st:E15\endcsname{\hyperlink{step-II-F-4}{proves step 4 of Theorem F}}
\expandafter\def\csname SBack@push:src:E.2.23\endcsname{\hyperlink{step-II-F-5}{proves step 5 of Theorem F}}
\expandafter\def\csname SBack@push:src:E.2.24\endcsname{\hyperlink{step-II-F-6}{proves step 6 of Theorem F}}
\expandafter\def\csname SBack@inv:E2.5\endcsname{\hyperlink{step-II-F-7}{proves step 7 of Theorem F}}
\expandafter\def\csname SBack@inv:E2.2\endcsname{\hyperlink{step-II-F-7}{proves step 7 of Theorem F}}
\expandafter\def\csname SBack@ii:LR:covering\endcsname{\hyperlink{step-II-G-1}{proves step 1 of Theorem G}; \hyperlink{step-II-G-2}{proves step 2 of Theorem G}; \hyperlink{step-II-G-3}{proves step 3 of Theorem G}; \hyperlink{step-II-G-4}{proves step 4 of Theorem G}}

\begin{document}
\title{The $(3,4,\infty)$ family of complex six-spheres\\[.5em]II. Torus-fibred threefolds and the sign theorem}
\author{AI-produced mathematical text}
\date{21 November 2026}
\maketitle
\begin{abstract}
We study compact complex threefolds fibred over curves with vanishing
first and second Betti numbers. A cotangent-torsion connecting map
omitted from a published vanishing argument gives an explicit nodal
counterexample. We prove the exact normalization-Picard criterion
for generic vanishing. For a torus general fibre, every degeneration
has nonnegative Euler number, with arbitrary multiplicities and
support singularities; globally the Euler number is zero or two.
We describe the equality supports and the additional equality
forced by identity monodromy, using integral component counts,
local parity, conductor comparisons and compact constructions.
We compute tangent twists and deformation spaces, prove persistence
of the fibration for the families described here, and determine
the two-turn covering parameter of the Liu--Rollenske family.
A flat-kernel and conic two-jet calculation gives a nonempty open
family of twelve-tooth smoothings with kernel
$\mathcal O(-8)\oplus\mathcal O(-9)^2$ and Hodge numbers
$h^{1,1}=49$, $h^{1,2}=48$. We state the remaining splitting
problem for general and special comb members.
\end{abstract}
\clearpage
\tableofcontents
\clearpage
\section{Results and their proofs in outline}\label{sec:results}
\subsection{The series and this paper}
This series studies the $(3,4,\infty)$ family of torus-fibred complex
structures on $S^6$ \citep{AC2026,Engel2026}, its periods, maps and smooth
monodromy, and the higher-dimensional complex manifolds built from it.
The table gives the main result of each paper and the papers it uses.
\begin{center}\small
\begin{tabularx}{\textwidth}{@{}lXl@{}}\toprule
Paper & Main result & Uses\\\midrule
I & Exact periods, moduli, fibre strata and the two ends & ---\\
II & Torus-fibred threefolds with $b_1=b_2=0$: the Campana--Demailly--Peternell argument and the sign theorem for all degenerations & I, IV\\
III & Maps from the six-spheres and the manifolds they give & I, II, IV, V\\
IV & Smooth monodromy in $\Theta_7$ & I, II, III, V\\
V & Complex fourfolds on $3\mathbb{HP}^2\#(-\mathbb{HP}^2)$ and a resolved fibre product & I, III, IV\\
\bottomrule\end{tabularx}\end{center}
\subsection{The main theorems}
\hypertarget{tag-03K1}{}\label{series:II:main-theorems}

We write $F_b=f^{-1}(b)$ for a full Cartier fibre, including its
multiplicities, and $S_b=(F_b)_{\rm red}$ for its reduced support.
\begin{definition}[Integral Campana--Demailly--Peternell hypotheses]
\hypertarget{tag-P201}{}\label{def:II:integral-cdp}
\StatementRecord{def:II:integral-cdp}{proved}{P201}{}
The integral CDP hypotheses are that $Y$ is a connected
smooth compact complex threefold, $H^1(Y,\Z)=H^2(Y,\Z)=0$,
$a(Y)=1$, and $f:Y\to\PP^1$ is its holomorphic algebraic reduction
with connected fibres.
\end{definition}

\hypertarget{main-II-A}{}
\begin{maintheorem}[The cotangent term and the exact generic-vanishing criterion]
\hypertarget{tag-03K2}{}\label{series:II:A}
\StatementRecord{series:II:A}{proved}{03K2}{\SeriesRef{cdp:st:A.4.9}, \SeriesRef{cdp:st:E12}, \SeriesRef{cdp:st:E17}, \SeriesRef{cdp:st:A.1.1}, \SeriesRef{cdp:st:B1}, \SeriesRef{cdp:st:C4f}}
Let $S$ be a reduced hypersurface in a smooth complex threefold $Y$,
let $A$ be a holomorphic line bundle on $S$, and put
\[
 N^*=\OO_Y(-S)|_S,\quad V=\Omega_Y^1|_S,\quad
 T=\operatorname{Tors}\Omega_S^1,\quad Q=\Omega_S^1/T,
 \quad K=\ker(V\to Q).
\]
The conormal map and its inverse image of torsion give
\[
 0\to N^*\to K\to T\to0,
 \qquad 0\to K\to V\to Q\to0.
\]
Writing $\delta_A:H^0(T\otimes A)\to H^1(N^*\otimes A)$
for the first connecting map, the two vanishings
$H^0(N^*\otimes A)=H^0(Q\otimes A)=0$ identify
$H^0(V\otimes A)$ canonically with $\ker\delta_A$.
In the family
\[
 \mathcal C=\{y^2z=x^2(x+z)+tz^3\}\subset\PP^2\times\Delta,
 \qquad Y=\mathcal C\times\PP^1,
 \qquad 0<\operatorname{radius}(\Delta)<4/27,
\]
the central support and every nontrivial normalization-gluing line
$A_\lambda$, $\lambda\in\C^*\setminus\{1\}$, have both displayed
endpoint vanishings and $h^0(V\otimes A_\lambda)=1$.

For a map satisfying the integral CDP hypotheses, a general
$L\in\Pic(Y)$ has $H^q(Y,T_Y\otimes L)=0$ for $q\ne1$
exactly when the normalization pullback
\[
 \Pic(Y)\longrightarrow\Pic(T_S),\qquad
 T_S\xrightarrow{\eta_S}S,
\]
is nonzero for every irreducible nonnormal critical support $S$.
A zero pullback gives $H^2(Y,T_Y\otimes L)\ne0$ for every $L$.
This criterion retains arbitrary full-fibre multiplicities and
arbitrary support singularities. In its vanishing case,
$\chi(Y,T_Y\otimes L)\le0$ and $c_3(Y)\le0$.

Two further exact sheaf calculations correct the earlier auxiliary
claims. For an elliptic curve $E$ with origin, its normalized
Poincar\'e line $P$ on $E\times E$, the projection
$g:E\times E\times\PP^1\to E$ to the second factor, and
$V=P\boxtimes\OO_{\PP^1}(-2)$, one has
$g_*V=R^1g_*V=0$ and $R^2g_*V=\C_0$, of length one.
For an effective Cartier divisor $D$ on a compact complex space
$Z$ and locally free $V$, the decreasing spaces
$H^0(Z,V(-kD))\subset H^0(Z,V)$ stabilize at the sections whose
support is disjoint from $|D|$. They eventually vanish when $Z$
has no embedded associated components and $D$ meets every
irreducible component.

\end{maintheorem}

\textit{Mechanism.} The conormal filtration retains the torsion
connecting map, and normalization Picard pullback decides whether
its conductor section survives every twist.
\textit{Relation to literature.} This corrects the sufficiency step
on p.~692 of \cite{CDP2020} and replaces
\cite[Theorem~2.2]{CDP2020}; the torsion-free conclusion of its
Proposition~7.3 is recovered in \ref{cdp:st:A.4.9}.
The final two calculations correct \cite[Lemma~1.5 and
Proposition~1.1]{CDP98} in the printed scopes stated in
\ref{cdp:st:E12} and \ref{cdp:st:E17}.
\hyperlink{outline-II-A}{Proof in outline}; full statements:
\ref{cdp:st:A.1.1}, \ref{cdp:st:B1}, \ref{cdp:st:C4f},
\ref{cdp:st:E12}, \ref{cdp:st:E17}.

\hypertarget{main-II-B}{}
\begin{maintheorem}[Component count and the smooth general fibre]
\hypertarget{tag-03K3}{}\label{series:II:B}
\StatementRecord{series:II:B}{proved}{03K3}{\SeriesRef{cdp:st:D.1.2}, \SeriesRef{cdp:st:C.3.1}}
Let $f:Y\to C$ be a proper holomorphic surjection with connected
fibres from a connected smooth compact complex threefold to a
smooth compact curve. Delete a nonempty finite set $\Delta$
containing all critical values, put $s=|\Delta|$, $g=g(C)$,
let $r$ count the irreducible components over $\Delta$, and set
\[
 t'=\dim_\Q H_1(F,\Q)_{\pi_1(C\setminus\Delta)}
\]
for a smooth fibre $F$. With $b_2(Y)=0$,
\[
 r=2g+s-1+t'-b_1(Y).
\]
Under $H^1(Y,\Z)=H^2(Y,\Z)=0$, the base is $\PP^1$,
$r=s-1+t'$, and the integral coinvariant quotient
$H_1(F,\Z)_{\pi_1(C\setminus\Delta)}$ is free abelian.
For a holomorphic algebraic reduction with $a(Y)=1$ and
$b_1(Y)=b_2(Y)=0$, every smooth fibre satisfies
\[
 c_1(K_F)=0\text{ in }H^2(F,\Q),\qquad
 c_2(F)=0,\qquad \chi(F,\OO_F)=0.
\]
Its very general fibre is a complex two-torus or a primary or
secondary Hopf surface of algebraic dimension zero. Each of these
surviving types occurs in the constructions stated in the body.

\end{maintheorem}

\textit{Mechanism.} The exact homology sequence of the complement
retains monodromy coinvariants; the resulting integral restrictions
and the surface-group comparisons determine the surviving fibres.
\textit{Relation to literature.} The rational component formula is
credited to \cite[\S10.5, p.~91]{AC2026}; the body retains its
integral refinement and the precise general-fibre comparisons.
\hyperlink{outline-II-B}{Proof in outline}; full statements:
\ref{cdp:st:D.1.2}, \ref{cdp:st:C.3.1}.

\hypertarget{main-II-C}{}
\begin{maintheorem}[The sign for every degeneration and its equality criterion]
\hypertarget{tag-03K4}{}\label{series:II:C}
\StatementRecord{series:II:C}{proved}{03K4}{\SeriesRef{cdp:st:D.2.7}, \SeriesRef{cdp:st:C.2.G.1}, \SeriesRef{night:P4.58}, \SeriesRef{cdp:st:D.2.5}}
Let $Y$ be a connected smooth compact complex threefold with
$b_1(Y)=b_2(Y)=0$, and let $f:Y\to B$ be a proper surjective
holomorphic map with connected fibres to a smooth compact curve.
Suppose the smooth general fibre is a complex two-torus. For
arbitrary positive full-fibre multiplicities and arbitrary
reduced-support singularities,
\[
 B\simeq\PP^1,\qquad
 c_3(Y)=e(Y)\in\{0,2\},\qquad b_3(Y)=2-c_3(Y).
\]
Every reduced support has nonnegative Euler number. A singular
irreducible support with rational normalization resolution has
Euler number two; every other irreducible support has Euler
number zero. A normal-crossing support has Euler number equal
to its number of triple points.

The equality $c_3(Y)=0$ holds exactly when every support has its
corresponding description below. A smooth irreducible support
has a finite unramified torus cover. A singular irreducible support
has smooth minimal elliptic ruled or smooth minimal Hopf
normalization. A reducible support has normalization a disjoint
union of smooth minimal Hopf or non-K\"ahler elliptic surfaces;
the reduced conductor image and inverse image are disjoint unions
of smooth elliptic curves, and their finite branch maps are
unramified. These are the normalization and conductor maps of
the actual full fibre, with its multiplicities retained.
For identity monodromy on $H_1(F,\Q)$ one always has
$c_3(Y)=0$ and $b_3(Y)=2$; the compact construction
\ref{cdp:st:D.2.7} attains this value with
$H_1(Y,\Z)=H_2(Y,\Z)=0$.

\end{maintheorem}

\textit{Mechanism.} The normalization and conductor yield a sum of
nonnegative local Euler contributions; identity monodromy supplies
the opposite inequality through bounded logarithmic periods.
\textit{Relation to literature.} The value two for the supplied
sphere family is credited to \cite[Theorem~7.22, p.~60]{AC26}
and \cite[Proposition~7.1]{Engel2026}. The unrestricted sign and
equality statement is proved through the actual fibre comparisons
in \ref{cdp:st:C.2.G.1}; the analytic period input is
\cite[\S2, pp.~234--239]{Steenbrink1976} in its stated scope.
\hyperlink{outline-II-C}{Proof in outline}; full statements:
\ref{night:P4.58}, \ref{cdp:st:C.2.G.1}, \ref{cdp:st:D.2.5}.

\hypertarget{main-II-D}{}
\begin{maintheorem}[Local parity and prescribed Euler numbers]
\hypertarget{tag-03K5}{}\label{series:II:D}
\StatementRecord{series:II:D}{proved}{03K5}{\SeriesRef{cdp:st:C.2.UN.1}, \SeriesRef{cdp:st:C.2.UN.2}, \SeriesRef{cdp:st:C.2.UN.4}, \SeriesRef{cdp:st:C.2.UN.14}}
Let $f:Y\to B$ be a proper holomorphic map from a smooth complex
threefold to a smooth complex curve, with connected fibres and
smooth general fibre a complex two-torus. Every reduced fibre
support has even Euler number, with arbitrary singularities and
multiplicities. This assertion imposes no compactness, K\"ahler
or Betti-number condition on $Y$.
For every fixed $b\in B$ and every $k\in\Z$, a finite composition
of blowups in smooth compact centres over $b$ gives
\[
 \pi:\widetilde Y\to Y,
 \qquad e((\widetilde Y_b)_{\rm red})=2k,
 \qquad \pi|_{B\setminus\{b\}}\text{ a biholomorphism}.
\]
The strict transforms retain their original full-fibre
multiplicities. The operation preserves the finite multiset of
bimeromorphism classes of non-uniruled normalization resolutions,
through the induced normalization and blowdown maps. Every such
filling has zero-Euler refinements with arbitrarily many components.
Every even integer also occurs in a smooth projective total space.

\end{maintheorem}

\textit{Mechanism.} Mapping-torus duality and the symmetric wedge
pairing force evenness; exceptional-divisor replacement gives Euler
changes $+2$ at a point and $-4$ along a plane quartic after a point
blowup.
\textit{Relation to literature.} The proof applies standard pair
duality, the Wang sequence and the blowup comparison in the precise
forms displayed in the body; the prescribed modifications retain
the given punctured family.
\hyperlink{outline-II-D}{Proof in outline}; full statements:
\ref{cdp:st:C.2.UN.1}, \ref{cdp:st:C.2.UN.2},
\ref{cdp:st:C.2.UN.4}, \ref{cdp:st:C.2.UN.14}.

\hypertarget{main-II-E}{}
\begin{maintheorem}[Twisted tangent cohomology of the standard family]
\hypertarget{tag-03K6}{}\label{series:II:E}
\StatementRecord{series:II:E}{proved}{03K6}{\SeriesRef{cdp:st:C5b}, \SeriesRef{cdp:st:C5c}, \SeriesRef{cdp:st:E8}, \SeriesRef{cdp:st:E18}}
Let $X=X_c$ be a member of the standard $(3,4,\infty)$ family,
with the fixed source, physical line $M$, affine attachment data,
period normalization and retained reference cusp fan of Paper~I.
Write $M_z$ for the ambient Picard line with parameter $z\in\C$
in that normalization. For $n\in\Z$, define
\[
 \begin{gathered}
 0\le r_1(n)<3,\ r_1(n)\equiv n\pmod3,\qquad
 0\le r_2(n)<4,\ r_2(n)\equiv3n\pmod4,\\
 d(n)=\frac{n-4r_1(n)-3r_2(n)}{12},\quad
 u(n)=\max\{d(n)+1,0\},\quad v(n)=\max\{-d(n)-1,0\},\\
 \epsilon_n=0\ (3\mid n),\qquad \epsilon_n=1\ (3\nmid n).
 \end{gathered}
\]
Then
\[
\begin{aligned}
 h^0(T_X\otimes M_n)&=u(n)+u(n-1),\\
 h^1(T_X\otimes M_n)&=v(n)+v(n-1)+u(n)+u(n-1)+u(n+2),\\
 h^2(T_X\otimes M_n)&=v(n)+v(n-1)+v(n+2)+u(n+6)+\epsilon_n,\\
 h^3(T_X\otimes M_n)&=v(n+6).
\end{aligned}
\]
The map to $H^2(\mathrm dP_6,K_{\mathrm dP_6})=\C$ has rank one
exactly for $n\in\{-6,-3,-2\}\cup\Z_{\ge0}$, and rank zero
otherwise. At $n=-1$ its source has dimension one.

For $z\notin\Z$, retain the marked period functions
$\tau,\mu,\beta_c$ and let $\mathcal R_z$ be the regular base
points where integers $(n,p,q,r)$ satisfy
\[
 q\tau+r\mu=p,\qquad z=n-6q\mu-r\beta_c.
\]
Put $\rho=\tau(z_1)$,
$\Lambda_1=\Z+(3\beta_c(z_1)+4-2\rho)\Z$ and
$\Lambda_2=\Z+(2\beta_c(z_2)+3-3i)\Z$.
Define $\ell_1(z)=1+\epsilon_n$ at
$z=n+k(3\beta_c(z_1)+4-2\rho)$, $k\ne0$, and zero elsewhere.
Define $\ell_2(z)=1$ at
$z=n+k(2\beta_c(z_2)+3-3i)$, $k\ne0$,
$n\equiv1,2\pmod4$, and zero elsewhere. The set $\mathcal R_z$
is finite. With $J_z=2|\mathcal R_z|+\ell_1(z)+\ell_2(z)$,
\[
 h^\bullet(X,T_X\otimes M_z)=(0,J_z,1+J_z,0),
\]
and its map to the same normalization target is surjective.

For the symmetric period matrix
$\Omega=\left(\begin{smallmatrix}\tau&\mu\\\mu&\beta/6\end{smallmatrix}\right)$,
with Jacobi theta periods $1,\tau$ and
$\beta/6=(i/\pi)\log(\vartheta_1(\mu,\tau)/\eta(\tau))+\mathrm{constant}$,
put $h=\partial_\mu\log\vartheta_1$. The exact identity is
\[
 \det\frac{d\Omega}{d\tau}
 =-\left(\mu'-\frac{i}{2\pi}h\right)^2
   -\frac{\wp(\mu,\tau)}{4\pi^2}.
\]
On the original branch $\wp(\mu,\tau)=0$ it equals
$-g_2^2/(36\pi^2g_3)$. The rank on the smooth root base is one
at the order-three point and two at every other finite point;
the determinant has order two there, and the logarithmic cusp
tensor has rank two.

\end{maintheorem}

\textit{Mechanism.} Root-curve direct images and exact local
resonance modules retain the map to the normalization as well as
its source dimension.
\textit{Relation to literature.} The untwisted tangent tuple is
among the credited DGKX version-one results discussed in Paper~I;
the full twist and map formulas here give the explicit correction
to the generic vanishing invoked in \cite[Theorem~2.2]{CDP2020}.
\hyperlink{outline-II-E}{Proof in outline}; full statements:
\ref{cdp:st:C5b}, \ref{cdp:st:C5c}, \ref{cdp:st:E8},
\ref{cdp:st:E18}.

\hypertarget{main-II-F}{}
\begin{maintheorem}[Fibrations, deformation spaces and comb Hodge numbers]
\hypertarget{tag-03K7}{}\label{series:II:F}
\StatementRecord{series:II:F}{proved}{03K7}{\SeriesRef{inv:E2.4}, \SeriesRef{cdp:st:E6}, \SeriesRef{cdp:st:E13}, \SeriesRef{cdp:st:E14}, \SeriesRef{cdp:st:E15}, \SeriesRef{push:src:E.2.23}, \SeriesRef{push:src:E.2.24}, \SeriesRef{inv:E2.2}, \SeriesRef{inv:E2.5}}
Let $g:Y\to\PP^1$ be proper and surjective with connected fibres
from a smooth connected compact complex threefold, and let its
smooth general fibre be a complex two-torus. Rank two of its
Kodaira--Spencer tensor at one regular value ensures that every
proper deformation $\mathcal Y\to(S,0)$, including one over a
nonreduced complex analytic germ, admits after shrinking $S$ a
map $G:\mathcal Y\to\PP^1\times S$ extending $g$. Every $G_s$
is surjective with connected fibres and smooth general fibre a
two-torus. Under $b_2(Y)=0$ its sufficiently small deformations
have algebraic dimension one.

For the constructed Viaclovsky two-parameter family, with
$\mu$ close to zero and nonzero sufficiently small scale
$\lambda$, and for the Liu--Rollenske family at sufficiently
small nonzero scale, the tangent tuple is $(1,2,2,0)$ and the
Kuranishi germ is a smooth reduced surface. The displayed
Viaclovsky family is locally miniversal; the second
Liu--Rollenske direction splits its cuspidal fibre into two
rank-one nodal fibres. For every positive-comb manifold $X_m$
in the construction of \cite{Viaclovsky2026Combs}, in particular
for its constructed $m\ge12$ range,
\[
 h^\bullet(X_m,T_{X_m})=(2m+1,8m+1,6m+1,0),
\]
and the Kuranishi germ is smooth and reduced of dimension $8m+1$.
Its local period-section family retains the affine attachment
classes, three critical values, multiplicities $3,4$, and the
hexagonal fibre type.
For the original marked section on $\mathbb P(3,4)$, there is a
nonempty open family obtained by embedded smoothings of twelve
admissible deformed twistor conics at distinct ordinary points,
fixing the three marked period planes, whose actual cup kernel is
\[
 K_{12}=\mathcal O_{\mathbb P^1}(-8)
       \oplus\mathcal O_{\mathbb P^1}(-9)^2,
 \qquad \operatorname{coker}\kappa_1=0.
\]
Here $K_{12}$ and $\kappa_1$ are the actual sheaf and cup map of
\eqref{invII:eq:E.6}. These members have
$h^{1,1}=49$ and $h^{1,2}=48$.

In any proper smooth holomorphic family over a connected reduced
complex analytic base containing a constructed sphere member with
coprime ordered multiplicities $(3k,4l)$, every fibre is a
constructed member with that same pair. Two such members are
deformation equivalent through finite chains precisely when their
ordered pairs agree. On the fixed smooth sphere, the structures
with one such pair form an open and closed subset of the
$C^\infty$ space of integrable complex structures, and have the
same separation in the quotients by $\operatorname{Diff}_0(S^6)$
and by orientation-preserving diffeomorphisms.

\end{maintheorem}

\textit{Mechanism.} The fibrewise cup-product map removes map
obstructions, actual period sections realize all tangent
directions, and propagation preserves the ordered multiplicities.
The flat kernel at a comb node gives a cyclic two-jet quotient;
actual conic frames and tooth collapse transport its balanced
cohomology to the smooth section.
\textit{Relation to literature.} The source families retain their
credit to \cite{Viaclovsky2026Two,LiuRollenske2026,Viaclovsky2026Combs}.
The comb vector-field value corrects the bound in
\cite[Corollary~1.3, p.~2]{HKP2000}, as explained in
\ref{inv:E2.4}.
The twelve-tooth calculation gives the stated open-family values
in the remaining-Hodge-number question of
\cite[Section~10(2)]{Viaclovsky2026Combs}; the general and
special-member splitting problem remains in \ref{inv:E2.3}.
\hyperlink{outline-II-F}{Proof in outline}; full statements:
\ref{cdp:st:E6}, \ref{cdp:st:E13}, \ref{cdp:st:E14},
\ref{cdp:st:E15}, \ref{inv:E2.2}, \ref{inv:E2.5},
\ref{push:src:E.2.23}, \ref{push:src:E.2.24}.

\hypertarget{main-II-G}{}
\begin{maintheorem}[The Liu--Rollenske covering parameter]
\hypertarget{tag-03K8}{}\label{series:II:G}
\StatementRecord{series:II:G}{proved}{03K8}{\SeriesRef{ii:LR:covering}, \SeriesRef{night:P3K.36}}
Fix the exact scalar $u$ in the Liu--Rollenske action $B_u=uB_0$,
the good fourth root at the $III^*$ source fibre, its fixed
root-line character, the logarithmic vector $a=e_4/4$, and the
retained periodic cusp fan with its source and line attachment
maps. The construction gives a holomorphic family on the universal
logarithmic cover of a sufficiently small punctured $u$-disc.
One positive $u$-turn reaches a threefold having a different
biholomorphism class. The actual family descends through
$u=w^2$, and this degree two is minimal. Every retained constant
root-line character has the same covering degree.

\end{maintheorem}

\textit{Mechanism.} A positive scalar turn changes the multiple
fibre's Albanese Tate parameter from $Q$ to $-Q$; the explicit
two-turn source-line lift has the required closing norm on every
retained cap.
\textit{Relation to literature.} The source action, good-root
lattice and logarithmic vector are
\cite[Lemma~2.7, Corollary~2.8, pp.~6--7; Proposition~3.6 and
Definition~3.7, pp.~9--10 of the attached-source edition]{LiuRollenske2026}.
The covering comparison is \ref{ii:LR:covering}. The former
fan-independent smooth-class clause is corrected separately by
\ref{night:P3K.36}, with the actual $T_2$ comparison yielding
$26\in\Z/28$.
\hyperlink{outline-II-G}{Proof in outline}; full statement:
\ref{ii:LR:covering}.


\paragraph{Further results.}
\begingroup\small
\begin{longtable}{@{}ll@{\hspace{1cm}}ll@{}}
\toprule Result & Tag & Result & Tag\\\midrule\endhead
\SeriesRef{cdp:st:C.1.1} & \texttt{01BN} & \SeriesRef{push:src:E.2.6} & \texttt{033E}\\
\SeriesRef{cdp:st:C.1.3} & \texttt{01BP} & \SeriesRef{push:src:E.2.7} & \texttt{033F}\\
\SeriesRef{cdp:st:E1} & \texttt{01GG} & \SeriesRef{push:src:E.2.8} & \texttt{033G}\\
\SeriesRef{cdp:st:D.1.5} & \texttt{01G3} & \SeriesRef{push:src:E.2.9} & \texttt{033H}\\
\SeriesRef{cdp:st:D.1.6} & \texttt{01G4} & \SeriesRef{push:src:E.2.10} & \texttt{032V}\\
\SeriesRef{cdp:st:D.1.7} & \texttt{01G5} & \SeriesRef{push:src:E.2.11} & \texttt{032W}\\
\SeriesRef{cdp:st:D.1.8} & \texttt{01G6} & \SeriesRef{push:src:E.2.12} & \texttt{032X}\\
\SeriesRef{cdp:st:D.1.9} & \texttt{01G7} & \SeriesRef{push:src:E.2.13} & \texttt{032Y}\\
\SeriesRef{cdp:st:D.1.10} & \texttt{01FV} & \SeriesRef{push:src:E.2.14} & \texttt{032Z}\\
\SeriesRef{cdp:st:D.1.11} & \texttt{01FW} & \SeriesRef{push:src:E.2.15} & \texttt{0330}\\
\SeriesRef{cdp:st:D.1.12} & \texttt{01FX} & \SeriesRef{push:src:E.2.16} & \texttt{0331}\\
\SeriesRef{cdp:st:D.1.4} & \texttt{01G2} & \SeriesRef{push:src:E.2.17} & \texttt{0332}\\
\SeriesRef{cdp:st:D.1.13} & \texttt{01FY} & \SeriesRef{push:src:E.2.18} & \texttt{0333}\\
\SeriesRef{cdp:st:D.1.15} & \texttt{01FZ} & \SeriesRef{push:src:E.2.19} & \texttt{0334}\\
\SeriesRef{cdp:st:C.3.2} & \texttt{01F8} & \SeriesRef{push:src:E.2.20} & \texttt{0336}\\
\SeriesRef{cdp:st:C.3.3} & \texttt{01F9} & \SeriesRef{push:src:E.2.21} & \texttt{0337}\\
\SeriesRef{cdp:st:C.3.4} & \texttt{01FA} & \SeriesRef{push:src:E.2.22} & \texttt{0338}\\
\SeriesRef{cdp:st:C.3.7} & \texttt{01FD} & \SeriesRef{push:src:E.5.1} & \texttt{033N}\\
\SeriesRef{cdp:st:C.3.8} & \texttt{01FE} & \SeriesRef{push:src:E.5.2} & \texttt{033O}\\
\SeriesRef{cdp:st:C.3.9} & \texttt{01FF} & \SeriesRef{push:src:E.5.3} & \texttt{033P}\\
\SeriesRef{cdp:st:C.3.10} & \texttt{01F5} & \SeriesRef{wb:thm:uniform-nontorsion-audit} & \texttt{03DQ}\\
\SeriesRef{cdp:st:C.3.12} & \texttt{01F7} & \SeriesRef{inv:wb:uniform-nontorsion-provenance} & \texttt{01TQ}\\
\SeriesRef{cdp:st:C.3.11} & \texttt{01F6} & \SeriesRef{wbII:lem:torsion-identification} & \texttt{03IF}\\
\SeriesRef{cdp:st:C.3.5} & \texttt{01FB} & \SeriesRef{wbII:prop:conductor} & \texttt{03IK}\\
\SeriesRef{cdp:st:C.3.6} & \texttt{01FC} & \SeriesRef{wbII:cor:normal-crossing-kernel} & \texttt{03DU}\\
\SeriesRef{cdp:st:A.2.1} & \texttt{01AF} & \SeriesRef{wbII:lem:gluing-bundle} & \texttt{03IB}\\
\SeriesRef{cdp:st:A.2.2} & \texttt{01AG} & \SeriesRef{wbII:lem:extending-gluing} & \texttt{03I9}\\
\SeriesRef{cdp:st:A.1.2} & \texttt{01AE} & \SeriesRef{wbII:thm:cartier-duality} & \texttt{03IY}\\
\SeriesRef{cdp:st:A.2.4} & \texttt{01AI} & \SeriesRef{wbII:prop:cartier-duality-absolute} & \texttt{03IJ}\\
\SeriesRef{cdp:st:A.2.5} & \texttt{01AJ} & \SeriesRef{wbII:lem:curve-flatness} & \texttt{03I8}\\
\SeriesRef{cdp:st:A.2.6} & \texttt{01AK} & \SeriesRef{wbII:prop:analytic-base-change} & \texttt{03II}\\
\SeriesRef{cdp:st:B2} & \texttt{01BL} & \SeriesRef{wbII:cor:surface-top-base-change} & \texttt{03DV}\\
\SeriesRef{cdp:st:B3} & \texttt{01BM} & \SeriesRef{wbII:prop:pic0-pullback} & \texttt{03IO}\\
\SeriesRef{cdp:st:A.3.1} & \texttt{01AL} & \SeriesRef{wbII:prop:all-pic-restriction} & \texttt{03IG}\\
\SeriesRef{cdp:st:A.3.2} & \texttt{01AM} & \SeriesRef{wbII:prop:leray-edge-restriction} & \texttt{03IM}\\
\SeriesRef{cdp:st:A.3.3} & \texttt{01AN} & \SeriesRef{wbII:lem:picard-crossed-defect} & \texttt{03IC}\\
\SeriesRef{cdp:st:A.3.4} & \texttt{01AO} & \SeriesRef{wbII:lem:picard-finite-norm} & \texttt{03IE}\\
\SeriesRef{cdp:st:A.3.5} & \texttt{01AP} & \SeriesRef{wbII:lem:picard-cusp-defect} & \texttt{03ID}\\
\SeriesRef{cdp:st:A.4.1} & \texttt{01AQ} & \SeriesRef{wbII:prop:direct-boundary-picard} & \texttt{03IL}\\
\SeriesRef{cdp:st:A.4.2} & \texttt{01AR} & \SeriesRef{wbII:prop:local-cusp-picard-obstruction} & \texttt{03IN}\\
\SeriesRef{cdp:st:A.4.3} & \texttt{01AS} & \SeriesRef{wbII:prop:analytic-ah-boundary} & \texttt{03IH}\\
\SeriesRef{cdp:st:A.4.4} & \texttt{01AT} & \SeriesRef{wbII:lem:analytic-support-specialization} & \texttt{03I6}\\
\SeriesRef{cdp:st:A.4.5} & \texttt{01AU} & \SeriesRef{wbII:thm:analytic-all-picard} & \texttt{03IW}\\
\SeriesRef{cdp:st:A.4.6} & \texttt{01AV} & \SeriesRef{wbII:cor:candidate-pic0-route} & \texttt{03DS}\\
\SeriesRef{cdp:st:A.4.7} & \texttt{01AW} & \SeriesRef{wbII:thm:candidate-conductor} & \texttt{03IX}\\
\SeriesRef{cdp:st:A.4.8} & \texttt{01AX} & \SeriesRef{wbII:cor:candidate-conductor-actual} & \texttt{03DR}\\
\SeriesRef{cdp:st:A.5.1} & \texttt{01AZ} & \SeriesRef{wbII:lem:generic-endpoints} & \texttt{03IA}\\
\SeriesRef{cdp:st:A.5.2} & \texttt{01BA} & \SeriesRef{wbII:lem:cdp-generic-endpoints} & \texttt{03I7}\\
\SeriesRef{cdp:st:A.5.3} & \texttt{01BD} & \SeriesRef{wbII:thm:corrected-cdp-global} & \texttt{03J0}\\
\SeriesRef{cdp:st:A.5.4} & \texttt{01BE} & \SeriesRef{wbII:cor:corrected-cdp-relative} & \texttt{03DT}\\
\SeriesRef{cdp:st:A.5.6} & \texttt{01BG} & \SeriesRef{wbII:thm:corrected-cdp-fibrewise} & \texttt{03IZ}\\
\SeriesRef{cdp:st:A.5.9} & \texttt{01BJ} & \SeriesRef{wbII:test:cdp-comparison} & \texttt{03IV}\\
\SeriesRef{cdp:st:A.5.11} & \texttt{01B1} & \SeriesRef{wbII:prop:relative-local} & \texttt{03IP}\\
\SeriesRef{cdp:st:A.5.12} & \texttt{01B2} & \SeriesRef{cdp:st:C.2.UN.3} & \texttt{01D8}\\
\SeriesRef{cdp:st:A.5.13} & \texttt{01B3} & \SeriesRef{cdp:st:C.2.UN.15} & \texttt{01CS}\\
\SeriesRef{cdp:st:A.5.5} & \texttt{01BF} & \SeriesRef{cdp:st:C.2.UN.16} & \texttt{01CT}\\
\SeriesRef{cdp:st:A.5.7} & \texttt{01BH} & \SeriesRef{cdp:st:C.2.UN.17} & \texttt{01CU}\\
\SeriesRef{cdp:st:A.5.8} & \texttt{01BI} & \SeriesRef{cdp:st:C.2.UN.19} & \texttt{01CW}\\
\SeriesRef{cdp:st:A.5.14} & \texttt{01B4} & \SeriesRef{cdp:st:C.2.UN.20} & \texttt{01CY}\\
\SeriesRef{cdp:st:A.5.10} & \texttt{01B0} & \SeriesRef{cdp:st:C.2.UN.22} & \texttt{01D0}\\
\SeriesRef{cdp:st:A.5.20} & \texttt{01BB} & \SeriesRef{cdp:st:C.2.UN.23} & \texttt{01D1}\\
\SeriesRef{cdp:st:A.5.15} & \texttt{01B5} & \SeriesRef{cdp:st:C.2.UN.26} & \texttt{01D4}\\
\SeriesRef{cdp:st:A.5.17} & \texttt{01B7} & \SeriesRef{cdp:st:C.2.UN.28} & \texttt{01D6}\\
\SeriesRef{cdp:st:A.5.18} & \texttt{01B8} & \SeriesRef{cdp:st:C.2.UN.29} & \texttt{01D7}\\
\SeriesRef{cdp:st:A.5.19} & \texttt{01B9} & \SeriesRef{cdp:st:C.2.UN.31} & \texttt{01DA}\\
\SeriesRef{cdp:st:A.5.16} & \texttt{01B6} & \SeriesRef{cdp:st:C.2.UN.33} & \texttt{01DC}\\
\SeriesRef{cdp:st:A.5.21} & \texttt{01BC} & \SeriesRef{cdp:st:C.2.UN.36} & \texttt{01DF}\\
\SeriesRef{cdp:st:C.2.NC.2} & \texttt{01C9} & \SeriesRef{cdp:st:C.2.UN.38} & \texttt{01DH}\\
\SeriesRef{cdp:st:C.2.NC.3} & \texttt{01CA} & \SeriesRef{cdp:st:C.2.UN.39} & \texttt{01DI}\\
\SeriesRef{cdp:st:C.2.NC.4} & \texttt{01CB} & \SeriesRef{cdp:st:C.2.UN.44} & \texttt{01DO}\\
\SeriesRef{cdp:st:C.2.NC.5} & \texttt{01CC} & \SeriesRef{cdp:st:C.2.UN.47} & \texttt{01DR}\\
\SeriesRef{cdp:st:C2d} & \texttt{01FJ} & \SeriesRef{cdp:st:C.2.UN.5} & \texttt{01DU}\\
\SeriesRef{night:P4.1} & \texttt{02WG} & \SeriesRef{cdp:st:C.2.CR.1} & \texttt{01BQ}\\
\SeriesRef{night:P4.2} & \texttt{02WR} & \SeriesRef{cdp:st:C.2.CR.2} & \texttt{01BR}\\
\SeriesRef{night:P4.3} & \texttt{02X2} & \SeriesRef{cdp:st:C.2.CR.3} & \texttt{01BS}\\
\SeriesRef{night:P4.4} & \texttt{02XD} & \SeriesRef{cdp:st:C.2.CR.4} & \texttt{01BT}\\
\SeriesRef{night:P4.5} & \texttt{02XO} & \SeriesRef{cdp:st:C.2.CR.5} & \texttt{01BU}\\
\SeriesRef{night:P4.6} & \texttt{02XY} & \SeriesRef{cdp:st:C.2.CR.6} & \texttt{01BV}\\
\SeriesRef{night:P4.7} & \texttt{02XZ} & \SeriesRef{cdp:st:C.2.UN.9} & \texttt{01EZ}\\
\SeriesRef{night:P4.8} & \texttt{02Y0} & \SeriesRef{cdp:st:C.2.UN.10} & \texttt{01CN}\\
\SeriesRef{night:P4.9} & \texttt{02Y1} & \SeriesRef{cdp:st:C.2.IRR.1} & \texttt{01C3}\\
\SeriesRef{night:P4.10} & \texttt{02WH} & \SeriesRef{cdp:st:C.2.VII.1} & \texttt{01F0}\\
\SeriesRef{night:P4.11} & \texttt{02WI} & \SeriesRef{cdp:st:C.2.VII.2} & \texttt{01F1}\\
\SeriesRef{night:P4.12} & \texttt{02WJ} & \SeriesRef{cdp:st:C.2.VII.3} & \texttt{01F2}\\
\SeriesRef{night:P4.13} & \texttt{02WK} & \SeriesRef{cdp:st:C.2.VII.4} & \texttt{01F3}\\
\SeriesRef{night:P4.14} & \texttt{02WL} & \SeriesRef{cdp:st:C.2.RAT.1} & \texttt{01CD}\\
\SeriesRef{night:P4.15} & \texttt{02WM} & \SeriesRef{cdp:st:C.2.RAT.2} & \texttt{01CE}\\
\SeriesRef{night:P4.16} & \texttt{02WN} & \SeriesRef{cdp:st:C.2.RAT.3} & \texttt{01CF}\\
\SeriesRef{night:P4.17} & \texttt{02WO} & \SeriesRef{cdp:st:C.2.RAT.4} & \texttt{01CG}\\
\SeriesRef{night:P4.18} & \texttt{02WP} & \SeriesRef{cdp:st:C.2.RAT.5} & \texttt{01CH}\\
\SeriesRef{night:P4.19} & \texttt{02WQ} & \SeriesRef{cdp:st:C.2.RAT.7} & \texttt{01CJ}\\
\SeriesRef{night:P4.20} & \texttt{02WS} & \SeriesRef{cdp:st:C.2.RAT.8} & \texttt{01CK}\\
\SeriesRef{night:P4.21} & \texttt{02WT} & \SeriesRef{cdp:st:C.2.RAT.9} & \texttt{01CL}\\
\SeriesRef{night:P4.22} & \texttt{02WU} & \SeriesRef{cdp:st:C.2.D4.1} & \texttt{01BW}\\
\SeriesRef{night:P4.23} & \texttt{02WV} & \SeriesRef{cdp:st:C.2.D4.2} & \texttt{01BX}\\
\SeriesRef{night:P4.24} & \texttt{02WW} & \SeriesRef{cdp:st:C.2.D4.3} & \texttt{01BY}\\
\SeriesRef{night:P4.25} & \texttt{02WX} & \SeriesRef{cdp:st:C.2.UN.7} & \texttt{01EG}\\
\SeriesRef{night:P4.26} & \texttt{02WY} & \SeriesRef{cdp:st:C.2.UN.12} & \texttt{01CP}\\
\SeriesRef{night:P4.27} & \texttt{02WZ} & \SeriesRef{cdp:st:C.2.INT.1} & \texttt{01C0}\\
\SeriesRef{night:P4.28} & \texttt{02X0} & \SeriesRef{cdp:st:C.2.INT.2} & \texttt{01C1}\\
\SeriesRef{night:P4.29} & \texttt{02X1} & \SeriesRef{cdp:st:C.2.INT.3} & \texttt{01C2}\\
\SeriesRef{night:P4.30} & \texttt{02X3} & \SeriesRef{cdp:st:C.2.MH.1} & \texttt{01C4}\\
\SeriesRef{night:P4.31} & \texttt{02X4} & \SeriesRef{cdp:st:C.2.MH.2} & \texttt{01C5}\\
\SeriesRef{night:P4.32} & \texttt{02X5} & \SeriesRef{cdp:st:C.2.MH.3} & \texttt{01C6}\\
\SeriesRef{night:P4.33} & \texttt{02X6} & \SeriesRef{cdp:st:C.2.MH.4} & \texttt{01C7}\\
\SeriesRef{night:P4.34} & \texttt{02X7} & \SeriesRef{cdp:st:C2a} & \texttt{01FG}\\
\SeriesRef{night:P4.35} & \texttt{02X8} & \SeriesRef{cdp:st:C2b} & \texttt{01FH}\\
\SeriesRef{night:P4.36} & \texttt{02X9} & \SeriesRef{cdp:st:C2c} & \texttt{01FI}\\
\SeriesRef{night:P4.37} & \texttt{02XA} & \SeriesRef{cdp:st:C.2.UN.11} & \texttt{01CO}\\
\SeriesRef{night:P4.38} & \texttt{02XB} & \SeriesRef{cdp:st:C.2.UN.24} & \texttt{01D2}\\
\SeriesRef{night:P4.39} & \texttt{02XC} & \SeriesRef{cdp:st:C.2.UN.32} & \texttt{01DB}\\
\SeriesRef{night:P4.40} & \texttt{02XE} & \SeriesRef{cdp:st:C.2.UN.64} & \texttt{01EA}\\
\SeriesRef{night:P4.41} & \texttt{02XF} & \SeriesRef{cdp:st:C.2.UN.65} & \texttt{01EB}\\
\SeriesRef{night:P4.42} & \texttt{02XG} & \SeriesRef{cdp:st:C.2.UN.66} & \texttt{01EC}\\
\SeriesRef{night:P4.43} & \texttt{02XH} & \SeriesRef{cdp:st:C.2.UN.67} & \texttt{01ED}\\
\SeriesRef{night:P4.44} & \texttt{02XI} & \SeriesRef{cdp:st:C.2.UN.68} & \texttt{01EE}\\
\SeriesRef{night:P4.45} & \texttt{02XJ} & \SeriesRef{cdp:st:C.2.UN.6} & \texttt{01E5}\\
\SeriesRef{night:P4.46} & \texttt{02XK} & \SeriesRef{cdp:st:C.2.UN.8} & \texttt{01ER}\\
\SeriesRef{night:P4.47} & \texttt{02XL} & \SeriesRef{cdp:st:C.2.UN.13} & \texttt{01CQ}\\
\SeriesRef{night:P4.48} & \texttt{02XM} & \SeriesRef{cdp:st:C.2.UN.18} & \texttt{01CV}\\
\SeriesRef{night:P4.49} & \texttt{02XN} & \SeriesRef{cdp:st:C.2.UN.21} & \texttt{01CZ}\\
\SeriesRef{night:P4.50} & \texttt{02XP} & \SeriesRef{cdp:st:C.2.UN.25} & \texttt{01D3}\\
\SeriesRef{night:P4.51} & \texttt{02XQ} & \SeriesRef{cdp:st:C.2.UN.27} & \texttt{01D5}\\
\SeriesRef{night:P4.52} & \texttt{02XR} & \SeriesRef{cdp:st:C.2.UN.30} & \texttt{01D9}\\
\SeriesRef{night:P4.53} & \texttt{02XS} & \SeriesRef{cdp:st:C.2.UN.34} & \texttt{01DD}\\
\SeriesRef{night:P4.54} & \texttt{02XT} & \SeriesRef{cdp:st:C.2.UN.35} & \texttt{01DE}\\
\SeriesRef{night:P4.55} & \texttt{02XU} & \SeriesRef{cdp:st:C.2.UN.37} & \texttt{01DG}\\
\SeriesRef{night:P4.56} & \texttt{02XV} & \SeriesRef{cdp:st:C.2.UN.41} & \texttt{01DL}\\
\SeriesRef{night:P4.57} & \texttt{02XW} & \SeriesRef{cdp:st:C.2.UN.42} & \texttt{01DM}\\
\SeriesRef{night:P4R.1} & \texttt{02Y2} & \SeriesRef{cdp:st:C.2.UN.43} & \texttt{01DN}\\
\SeriesRef{night:P4R.2} & \texttt{02YD} & \SeriesRef{cdp:st:C.2.UN.45} & \texttt{01DP}\\
\SeriesRef{night:P4R.3} & \texttt{02YF} & \SeriesRef{cdp:st:C.2.UN.40} & \texttt{01DK}\\
\SeriesRef{night:P4R.4} & \texttt{02YG} & \SeriesRef{cdp:st:C.2.UN.46} & \texttt{01DQ}\\
\SeriesRef{night:P4R.5} & \texttt{02YH} & \SeriesRef{cdp:st:C.2.UN.48} & \texttt{01DS}\\
\SeriesRef{night:P4R.6} & \texttt{02YI} & \SeriesRef{cdp:st:C.2.UN.49} & \texttt{01DT}\\
\SeriesRef{night:P4R.7} & \texttt{02YJ} & \SeriesRef{cdp:st:C.2.UN.50} & \texttt{01DV}\\
\SeriesRef{night:P4R.8} & \texttt{02YK} & \SeriesRef{cdp:st:C.2.UN.51} & \texttt{01DW}\\
\SeriesRef{night:P4R.9} & \texttt{02YL} & \SeriesRef{cdp:st:C.2.UN.52} & \texttt{01DX}\\
\SeriesRef{night:P4R.10} & \texttt{02Y3} & \SeriesRef{cdp:st:C.2.UN.53} & \texttt{01DY}\\
\SeriesRef{night:P4R.11} & \texttt{02Y4} & \SeriesRef{cdp:st:C.2.UN.54} & \texttt{01DZ}\\
\SeriesRef{night:P4R.12} & \texttt{02Y5} & \SeriesRef{cdp:st:C.2.UN.55} & \texttt{01E0}\\
\SeriesRef{night:P4R.13} & \texttt{02Y6} & \SeriesRef{cdp:st:C.2.UN.56} & \texttt{01E1}\\
\SeriesRef{night:P4R.14} & \texttt{02Y7} & \SeriesRef{cdp:st:C.2.UN.57} & \texttt{01E2}\\
\SeriesRef{night:P4R.15} & \texttt{02Y8} & \SeriesRef{cdp:st:C.2.UN.58} & \texttt{01E3}\\
\SeriesRef{night:P4R.16} & \texttt{02Y9} & \SeriesRef{cdp:st:C.2.UN.59} & \texttt{01E4}\\
\SeriesRef{night:P4R.17} & \texttt{02YA} & \SeriesRef{cdp:st:C.2.UN.60} & \texttt{01E6}\\
\SeriesRef{night:P4R.18} & \texttt{02YB} & \SeriesRef{cdp:st:C.2.UN.61} & \texttt{01E7}\\
\SeriesRef{night:P4R.19} & \texttt{02YC} & \SeriesRef{cdp:st:C.2.UN.62} & \texttt{01E8}\\
\SeriesRef{night:P4R.20} & \texttt{02YE} & \SeriesRef{cdp:st:C.2.UN.63} & \texttt{01E9}\\
\SeriesRef{cdp:st:C4a} & \texttt{01FK} & \SeriesRef{cdp:st:C.2.UN.69} & \texttt{01EF}\\
\SeriesRef{cdp:st:C4b} & \texttt{01FL} & \SeriesRef{cdp:st:C.2.UN.70} & \texttt{01EH}\\
\SeriesRef{cdp:st:C4c} & \texttt{01FM} & \SeriesRef{cdp:st:C.2.UN.71} & \texttt{01EI}\\
\SeriesRef{cdp:st:C4d} & \texttt{01FN} & \SeriesRef{cdp:st:C.2.UN.73} & \texttt{01EK}\\
\SeriesRef{cdp:st:C4e} & \texttt{01FO} & \SeriesRef{cdp:st:C.2.UN.76} & \texttt{01EN}\\
\SeriesRef{cdp:st:C4g} & \texttt{01FQ} & \SeriesRef{cdp:st:C.2.UN.74} & \texttt{01EL}\\
\SeriesRef{cdp:st:C5a} & \texttt{01FR} & \SeriesRef{cdp:st:C.2.UN.75} & \texttt{01EM}\\
\SeriesRef{cdp:st:D.2.1} & \texttt{01G8} & \SeriesRef{cdp:st:C.2.UN.77} & \texttt{01EO}\\
\SeriesRef{cdp:st:D.2.3} & \texttt{01GA} & \SeriesRef{cdp:st:C.2.UN.79} & \texttt{01EQ}\\
\SeriesRef{cdp:st:D.2.4} & \texttt{01GB} & \SeriesRef{cdp:st:C.2.UN.80} & \texttt{01ES}\\
\SeriesRef{cdp:st:D.2.6} & \texttt{01GD} & \SeriesRef{cdp:st:C.2.UN.81} & \texttt{01ET}\\
\SeriesRef{cdp:st:D.2.8} & \texttt{01GF} & \SeriesRef{cdp:st:C.2.UN.82} & \texttt{01EU}\\
\SeriesRef{cdp:st:D.2.2} & \texttt{01G9} & \SeriesRef{cdp:st:C.2.UN.83} & \texttt{01EV}\\
\SeriesRef{cdp:st:E5} & \texttt{01GU} & \SeriesRef{cdp:st:C.2.UN.84} & \texttt{01EW}\\
\SeriesRef{cdp:st:E7} & \texttt{01GW} & \SeriesRef{cdp:st:C.2.UN.78} & \texttt{01EP}\\
\SeriesRef{cdp:st:E10} & \texttt{01GH} & \SeriesRef{cdp:st:C.2.UN.72} & \texttt{01EJ}\\
\SeriesRef{cdp:st:E11} & \texttt{01GI} & \SeriesRef{cdp:st:C.2.UN.85} & \texttt{01EX}\\
\SeriesRef{cdp:st:E16} & \texttt{01GN} & \SeriesRef{cdp:st:C.2.UN.86} & \texttt{01EY}\\
\SeriesRef{push:src:E.2.4} & \texttt{033C} & \SeriesRef{inv:E2.3} & \texttt{01SW}\\
\SeriesRef{push:src:E.2.5} & \texttt{033D} &  & \\
\bottomrule
\end{longtable}\endgroup

\subsection{The proofs in outline}
\hypertarget{tag-03K9}{}\label{series:II:outline}

\hypertarget{outline-II-A}{}
\subsubsection*{Theorem A}
\begin{enumerate}[label=\arabic*.,leftmargin=*]
\item\hypertarget{step-II-A-1}{} The conormal map is injective at
each generic smooth point of the reduced hypersurface. Its inverse
image of cotangent torsion gives the two exact sequences, and their
cohomology identifies the ambient section space with the stated
connecting-map kernel
(\hyperref[cdp:st:A.1.1]{\ref*{cdp:st:A.1.1}}).
\item\hypertarget{step-II-A-2}{} The node has normalization
$u\mapsto(u^2-1,u(u^2-1))$ and identifies $u=1,-1$. Its gluing-line
complex is $\C\xrightarrow{1-\lambda}\C$, while the inverse
image of torsion is the pushforward of the trivial normalization
line. Taking the product with $\PP^1$ gives the asserted
counterexample dimensions
(\hyperref[cdp:st:B1]{\ref*{cdp:st:B1}}).
\item\hypertarget{step-II-A-3}{} Transgression gives unitary
infinite-order twists on every fibre support. The explicit
fibre-group calculation and neighbourhood comparison make
degree-one local-system cohomology vanish. The canonical
injection then kills the twisted torsion-free forms
(\hyperref[cdp:st:A.4.7]{\ref*{cdp:st:A.4.7}},
\hyperref[cdp:st:A.4.9]{\ref*{cdp:st:A.4.9}}).
\item\hypertarget{step-II-A-4}{} Full Cartier filtrations reduce
generic ambient vanishing to irreducible nonnormal supports.
Normalization Picard pullback determines persistence of their
conductor section, yielding exactly the displayed global
criterion (\hyperref[cdp:st:C4f]{\ref*{cdp:st:C4f}}).
\item\hypertarget{step-II-A-5}{} The Poincar\'e transform and
cohomology of $\OO_{\PP^1}(-2)$ produce the length-one higher
direct image. Stabilization of the finite-dimensional section
spaces and the associated-component calculation give the exact
divisor-twisting statement
(\hyperref[cdp:st:E12]{\ref*{cdp:st:E12}},
\hyperref[cdp:st:E17]{\ref*{cdp:st:E17}}).
\end{enumerate}

\hypertarget{outline-II-B}{}
\subsubsection*{Theorem B}
\begin{enumerate}[label=\arabic*.,leftmargin=*]
\item\hypertarget{step-II-B-1}{} Lifts of the free generators of
$\pi_1(C\setminus\Delta)$ split the punctured bundle group.
Its abelianization is $\Z^{2g+s-1}\oplus H_1(F,\Z)_{\pi_1}$;
the compact-support sequence then gives the component count and
integral freeness (\hyperref[cdp:st:D.1.2]{\ref*{cdp:st:D.1.2}}).
\item\hypertarget{step-II-B-2}{} Fibre adjunction and divisor
Riemann--Roch give the displayed Chern and coherent-Euler
identities. Surface classification supplies the finite necessary
list (\hyperref[cdp:st:C.3.1]{\ref*{cdp:st:C.3.1}}).
\item\hypertarget{step-II-B-3}{} The explicit surface-group and
monodromy comparisons exclude the bielliptic, Kodaira and Inoue
entries; relative algebraic reduction forces the surviving Hopf
fibre to have algebraic dimension zero
(\hyperref[cdp:st:C.3.7]{\ref*{cdp:st:C.3.7}},
\hyperref[cdp:st:C.3.8]{\ref*{cdp:st:C.3.8}},
\hyperref[cdp:st:C.3.9]{\ref*{cdp:st:C.3.9}},
\hyperref[cdp:st:C.3.10]{\ref*{cdp:st:C.3.10}},
\hyperref[cdp:st:C.3.12]{\ref*{cdp:st:C.3.12}}).
\item\hypertarget{step-II-B-4}{} The quotient and filling maps
produce the surviving primary and secondary Hopf types and a fixed
abelian general fibre (\hyperref[cdp:st:C.3.6]{\ref*{cdp:st:C.3.6}},
\hyperref[cdp:st:D.2.7]{\ref*{cdp:st:D.2.7}}).
\end{enumerate}

\hypertarget{outline-II-C}{}
\subsubsection*{Theorem C}
\begin{enumerate}[label=\arabic*.,leftmargin=*]
\item\hypertarget{step-II-C-1}{} Connected-fibre Leray makes the
base rational. Finite-character support comparison and local
conductor maps give the irreducible values zero and two and their
normalization descriptions; the normal-crossing calculation gives
the triple-point count
(\hyperref[night:P4.58]{\ref*{night:P4.58}},
\hyperref[cdp:st:C.2.MH.2]{\ref*{cdp:st:C.2.MH.2}}).
\item\hypertarget{step-II-C-2}{} An adjacent component pulls back
to a nonzero effective divisor with rational Chern class zero.
The surface and Picard comparisons leave the Hopf, non-K\"ahler
elliptic and Enoki possibilities and then exclude the Enoki
interfaces through the explicit local vanishing sphere
(\hyperref[cdp:st:C.2.INT.1]{\ref*{cdp:st:C.2.INT.1}},
\hyperref[cdp:st:C.2.INT.2]{\ref*{cdp:st:C.2.INT.2}}).
\item\hypertarget{step-II-C-3}{} Coherent conductor duality, curve
normalization and exact blowup counts give nonnegative summands
for every surviving reducible support. Equality forces the stated
minimal normalizations and unramified elliptic conductor maps
(\hyperref[cdp:st:C.2.INT.3]{\ref*{cdp:st:C.2.INT.3}}).
\item\hypertarget{step-II-C-4}{} Euler additivity gives
$e(Y)=\sum_b e(S_b)$. The nondegenerate skew pairing on $H^3$
and $e(Y)=2-b_3(Y)$ give the two values and the global equality
criterion (\hyperref[cdp:st:C.2.G.1]{\ref*{cdp:st:C.2.G.1}}).
\item\hypertarget{step-II-C-5}{} For identity monodromy, a
persistent positive-index obstruction has a fixed irreducible
support. The finite \'etale root cover, twisted cotangent
comparison and Cartier-layer filtration exclude this obstruction
on a smooth support. For a singular support, Gauduchon degree
and the nonzero conductor make its normalization have geometric
genus zero; vertical resolution preserves that property on all
components. The logarithmic
period obstruction contradicts identity monodromy, giving
$c_3(Y)\le0$ and hence equality
(\hyperref[cdp:st:D.2.3]{\ref*{cdp:st:D.2.3}},
\hyperref[cdp:st:D.2.5]{\ref*{cdp:st:D.2.5}}).
\item\hypertarget{step-II-C-6}{} The retained three-piece
holomorphic filling and its integral Mayer--Vietoris maps give
$H_1=H_2=0$ and Euler number zero
(\hyperref[cdp:st:D.2.7]{\ref*{cdp:st:D.2.7}}).
\end{enumerate}

\hypertarget{outline-II-D}{}
\subsubsection*{Theorem D}
\begin{enumerate}[label=\arabic*.,leftmargin=*]
\item\hypertarget{step-II-D-1}{} A full fibre tube $U$ retracts
onto the support, and its boundary $M$ is the nearby torus mapping
torus. Pair duality gives
$e(U)\equiv b_0(M)+b_1(M)+b_2(M)$ modulo two
(\hyperref[cdp:st:C.2.UN.1]{\ref*{cdp:st:C.2.UN.1}}).
\item\hypertarget{step-II-D-2}{} Wang identifies that parity with
$\dim\ker(\Lambda^2T-I)$. The wedge pairing and Cayley transform
make this dimension even for every orientation-preserving torus
monodromy (\hyperref[cdp:st:C.2.UN.1]{\ref*{cdp:st:C.2.UN.1}}).
\item\hypertarget{step-II-D-3}{} Replacing a smooth codimension-$r$
centre by its exceptional projective bundle changes support Euler
number by $(r-1)e(C)$. A point adds two; a plane quartic after a
point blowup gives net change minus two. Iteration gives every
prescribed value and cofinal zero-Euler models
(\hyperref[cdp:st:C.2.UN.4]{\ref*{cdp:st:C.2.UN.4}}).
\item\hypertarget{step-II-D-4}{} Strict transforms are
bimeromorphic on normalization resolutions and preserve their
base-function valuations. Every new exceptional surface is
explicitly ruled or rational, proving the multiset assertion
(\hyperref[cdp:st:C.2.UN.14]{\ref*{cdp:st:C.2.UN.14}}).
\item\hypertarget{step-II-D-5}{} The genus-two curve in the
projective product family and disjoint point blowups realize all
values $2(m-k)$ in compact projective threefolds
(\hyperref[cdp:st:C.2.UN.2]{\ref*{cdp:st:C.2.UN.2}}).
\end{enumerate}

\hypertarget{outline-II-E}{}
\subsubsection*{Theorem E}
\begin{enumerate}[label=\arabic*.,leftmargin=*]
\item\hypertarget{step-II-E-1}{} The stabilizer characters on the
root curve give the integral degrees $d(n)$, with direct-image
summands $d(n)$, $d(n-1)$, $d(n+2)$ and $d(n+6)$ and the
specified finite root contribution
(\hyperref[cdp:st:C5b]{\ref*{cdp:st:C5b}}).
\item\hypertarget{step-II-E-2}{} Cohomology of each line on
$\PP^1$ gives the four displayed dimensions. The actual cusp
map has the stated rank locus, including its zero value at
$n=-1$ (\hyperref[cdp:st:C5b]{\ref*{cdp:st:C5b}}).
\item\hypertarget{step-II-E-3}{} The exact translation-character
complex is acyclic off resonance. At a regular resonance its
local modules have two degree-one generators and a length-two
degree-two generator; finite-root characters give
$\ell_1,\ell_2$ (\hyperref[cdp:st:C5c]{\ref*{cdp:st:C5c}}).
\item\hypertarget{step-II-E-4}{} Properness of the direct-image
support makes the regular resonance set finite. The cusp adds
one reduced degree-two skyscraper and its nonzero map to the
normalization target. Leray gives $(0,J_z,1+J_z,0)$
(\hyperref[cdp:st:C5c]{\ref*{cdp:st:C5c}}).
\item\hypertarget{step-II-E-5}{} The identity-induced degree-six
isogeny gives the symmetric period matrix. The theta heat equation
and the eta derivative retain the Weierstrass term in the exact
determinant; the original branch equation then gives its stated
rank locus (\hyperref[cdp:st:E8]{\ref*{cdp:st:E8}},
\hyperref[cdp:st:E18]{\ref*{cdp:st:E18}}).
\end{enumerate}

\hypertarget{outline-II-F}{}
\subsubsection*{Theorem F}
\begin{enumerate}[label=\arabic*.,leftmargin=*]
\item\hypertarget{step-II-F-1}{} At a smooth fibre the boundary
map sends $a\eta_1+b\eta_2$ to
$(bu_1-au_2)\otimes\eta_1\wedge\eta_2$. Rank two makes it
injective and kills the infinitesimal map obstruction over
$\PP^1$ (\hyperref[cdp:st:E5]{\ref*{cdp:st:E5}}).
\item\hypertarget{step-II-F-2}{} The same cup-product comparison
at every Artin order, followed by the relative graph construction,
gives the extending map over the analytic germ, with its
connected torus fibres and stated algebraic dimension
(\hyperref[cdp:st:E6]{\ref*{cdp:st:E6}}).
\item\hypertarget{step-II-F-3}{} Root-curve tangent cohomology
and the actual two-parameter period sections give the two smooth
surface germs and their tuples
(\hyperref[cdp:st:E13]{\ref*{cdp:st:E13}},
\hyperref[cdp:st:E14]{\ref*{cdp:st:E14}}).
\item\hypertarget{step-II-F-4}{} Admissible comb period sections
realize every tangent direction while retaining affine attachments.
Their tangent isomorphism makes the Kuranishi ideal zero and gives
the displayed comb dimensions
(\hyperref[cdp:st:E15]{\ref*{cdp:st:E15}}).
\item\hypertarget{step-II-F-5}{} Disc propagation and finite
curve projections make the fixed ordered-pair locus open and
closed in a connected reduced analytic base. This gives
separation through arbitrary finite deformation chains
(\hyperref[push:src:E.2.23]{\ref*{push:src:E.2.23}}).
\item\hypertarget{step-II-F-6}{} The Kuranishi retraction transports
that separation to the $C^\infty$ space of integrable structures
and to the indicated diffeomorphism quotients
(\hyperref[push:src:E.2.24]{\ref*{push:src:E.2.24}}).
\item\hypertarget{step-II-F-7}{} The actual flat cup kernel in the comb smoothing family
imposes a cyclic two-jet condition at each central node. The explicit
conic frame map realizes a dense open set of these quotients;
twelve teeth attain the balanced cohomology. The exact collapse
and semicontinuity maps give $K_{12}=\mathcal O(-8)\oplus\mathcal O(-9)^2$
on the stated open family, and the actual Hodge formula gives
$(h^{1,1},h^{1,2})=(49,48)$
(\hyperref[inv:E2.5]{\ref*{inv:E2.5}},
\hyperref[inv:E2.2]{\ref*{inv:E2.2}}).
\end{enumerate}

\hypertarget{outline-II-G}{}
\subsubsection*{Theorem G}
\begin{enumerate}[label=\arabic*.,leftmargin=*]
\item\hypertarget{step-II-G-1}{} The exact scalar is $\alpha=cu$
with fixed $c\ne0$. Setting $\tau=\log(cu)/(2\pi i)$ gives
the central lattice generated by $(2,0),(1+i,0),(1,1/2),(0,\tau)$
and affine fourth-root action
$(z,\eta)\mapsto(-iz,\eta+\tau/4+\kappa)$, where
$4\kappa\in\Z$ (\hyperref[ii:LR:covering]{\ref*{ii:LR:covering}}).
\item\hypertarget{step-II-G-2}{} Projection to $\eta$ is the
Albanese map with period lattice
$\tfrac12\Z+(\tau/4+\kappa)\Z$. One $u$-turn sends its
Tate parameter $Q$ to $-Q$; $Q(j(Q)-j(-Q))$ tends to two,
excluding a fibre biholomorphism
(\hyperref[ii:LR:covering]{\ref*{ii:LR:covering}}).
\item\hypertarget{step-II-G-3}{} Ambient $H^2(Y,\R)=0$ makes
every effective divisor vertical. Meromorphic functions therefore
descend through the torus fibration; any endpoint biholomorphism
preserves its unique multiplicity-four fibre. The previous step
excludes every one-turn endpoint
(\hyperref[ii:LR:covering]{\ref*{ii:LR:covering}}).
\item\hypertarget{step-II-G-4}{} The lifted translation by the
marked source section, its rational square root and cap
commutators give the two-turn endpoint map. Its exact two-step
norm equals the prescribed one-gamma gluing, including the cusp
multiplier. The resulting free quotient has parameter
$w=r^2$ and $u=w^2$
(\hyperref[ii:LR:covering]{\ref*{ii:LR:covering}}).
\end{enumerate}
\subsection{Known and new}
\begingroup\small
\begin{longtable}{@{}p{.10\textwidth}p{.85\textwidth}@{}}\toprule Theorem & Exact source relationship\\\midrule\endhead
A & This corrects the sufficiency step
on p.~692 of \cite{CDP2020} and replaces
\cite[Theorem~2.2]{CDP2020}; the torsion-free conclusion of its
Proposition~7.3 is recovered in \ref{cdp:st:A.4.9}.
The final two calculations correct \cite[Lemma~1.5 and
Proposition~1.1]{CDP98} in the printed scopes stated in
\ref{cdp:st:E12} and \ref{cdp:st:E17}.\\
B & The rational component formula is
credited to \cite[\S10.5, p.~91]{AC2026}; the body retains its
integral refinement and the precise general-fibre comparisons.\\
C & The value two for the supplied
sphere family is credited to \cite[Theorem~7.22, p.~60]{AC26}
and \cite[Proposition~7.1]{Engel2026}. The unrestricted sign and
equality statement is proved through the actual fibre comparisons
in \ref{cdp:st:C.2.G.1}; the analytic period input is
\cite[\S2, pp.~234--239]{Steenbrink1976} in its stated scope.\\
D & The proof applies standard pair
duality, the Wang sequence and the blowup comparison in the precise
forms displayed in the body; the prescribed modifications retain
the given punctured family.\\
E & The untwisted tangent tuple is
among the credited DGKX version-one results discussed in Paper~I;
the full twist and map formulas here give the explicit correction
to the generic vanishing invoked in \cite[Theorem~2.2]{CDP2020}.\\
F & The source families retain their
credit to \cite{Viaclovsky2026Two,LiuRollenske2026,Viaclovsky2026Combs}.
The comb vector-field value corrects the bound in
\cite[Corollary~1.3, p.~2]{HKP2000}, as explained in
\ref{inv:E2.4}.
The twelve-tooth calculation gives the stated open-family values
in the remaining-Hodge-number question of
\cite[Section~10(2)]{Viaclovsky2026Combs}; the general and
special-member splitting problem remains in \ref{inv:E2.3}.\\
G & The source action, good-root
lattice and logarithmic vector are
\cite[Lemma~2.7, Corollary~2.8, pp.~6--7; Proposition~3.6 and
Definition~3.7, pp.~9--10 of the attached-source edition]{LiuRollenske2026}.
The covering comparison is \ref{ii:LR:covering}. The former
fan-independent smooth-class clause is corrected separately by
\ref{night:P3K.36}, with the actual $T_2$ comparison yielding
$26\in\Z/28$.\\
\bottomrule\end{longtable}\endgroup
\subsection{Methodology}
This text is AI-produced; no part of it was written by a human.
GPT-5.6 Sol (OpenAI), in Codex at reasoning effort ultra, contributed
to the working edition of 24 August--9 September 2026. The pinned
repository's authorship record also credits GPT-6 Astra (OpenAI),
at reasoning effort ultra, from 5 September.
GPT-6 Astra (OpenAI), in ChatGPT Pro mode, carried out the research
of 12--16 November, the independent re-derivation of 16--18 November,
and the writing of 16--20 November 2026. GPT-6 Astra (OpenAI),
in ChatGPT Pro mode, also undertook the publication revision and
further checks of this paper on 21 November 2026.

Claude Fable 5.2 (Anthropic), at effort high and at effort max,
and Claude Opus 5.5 (Anthropic), at effort max, undertook readings,
intakes, checks and recomputations on 12--20 November 2026.
Their objections and calculations were compared with the statements,
proofs and cited sources, and the resulting corrections were
incorporated into the text. The records specify the passages and
calculations examined by each model.

For the statements printed in this paper, the record contains
932 examinations of proof steps, 34 independent
re-derivations, and 102 computations or formal certificates.
It also records 659 readings of the proofs and applications,
259 uses of background results without a fault found,
and 2 refutations associated with their corrected statements.
Each total counts recorded links from a check to a statement;
one calculation used for several statements can contribute several
links. Every outcome specifies its scope. Re-reading an edited proof
is recorded separately from deriving its steps before comparison.
Exact algebra, numerical enclosures and analytic arguments retain
their respective scopes, and the records identify the implementations
used for independent computational comparisons.
No human has checked the mathematics.

The check records and programs are linked from the
\href{\SsixPublishRecords}{project repository}.
The text, programs and check records are dedicated to the public domain under
\href{https://creativecommons.org/publicdomain/zero/1.0/}{CC0 1.0 Universal}.

\subsection{How the paper is organized}
The body gives the marked constructions, their integral and holomorphic
comparison maps, and the consequences grouped in Theorems A, B, C, D, E, F and G.
The final part states each remaining question with its established
reduction. In the diagram an arrow records reuse of the particular
body argument identified below; a theorem group can contain several
arguments with different inputs. The statement headers give the full
local uses and the links to the corresponding outline steps.
\begin{figure}[htbp]\centering
\begin{tikzpicture}[>=stealth,every node/.style={font=\small}]
\node[draw,rounded corners,minimum width=2.2cm,minimum height=.65cm] (TA) at (0.0,0.0) {\hyperlink{main-II-A}{Theorem A}};
\node[draw,rounded corners,minimum width=2.2cm,minimum height=.65cm] (TB) at (3.5,0.0) {\hyperlink{main-II-B}{Theorem B}};
\node[draw,rounded corners,minimum width=2.2cm,minimum height=.65cm] (TC) at (7.0,0.0) {\hyperlink{main-II-C}{Theorem C}};
\node[draw,rounded corners,minimum width=2.2cm,minimum height=.65cm] (TD) at (0.0,-1.25) {\hyperlink{main-II-D}{Theorem D}};
\node[draw,rounded corners,minimum width=2.2cm,minimum height=.65cm] (TE) at (3.5,-1.25) {\hyperlink{main-II-E}{Theorem E}};
\node[draw,rounded corners,minimum width=2.2cm,minimum height=.65cm] (TF) at (7.0,-1.25) {\hyperlink{main-II-F}{Theorem F}};
\node[draw,rounded corners,minimum width=2.2cm,minimum height=.65cm] (TG) at (0.0,-2.5) {\hyperlink{main-II-G}{Theorem G}};
\draw[->,semithick] (TB) to (TA);
\end{tikzpicture}
\caption{Reuse of the indicated body arguments among the main theorem groups.}
\par\medskip
\begingroup\small
\begin{tabularx}{\textwidth}{@{}p{.10\textwidth}X@{}}
\toprule Arrow & Supplying and receiving arguments\\\midrule
$B\to A$ & \SeriesRef{cdp:st:C.3.1} $\longrightarrow$ \SeriesRef{cdp:st:C4f}. The complete necessary smooth-fibre surface list is used in A.4.9 Step1 for twisted local-system cohomology. C4e Step1 uses C.3.1 plus the stated exclusion propositions to divide into the torus and h1(O)=1 branches; C4f Step2 explicitly applies that fixed-fibre persistence proof.\\
\bottomrule\end{tabularx}\endgroup\end{figure}

\clearpage
\section{Common conventions}
\hypertarget{tag-034E}{}\label{sec:II:conventions}
\begin{definition}[Marked periods, fibres and orientations]
\hypertarget{tag-P202}{}\label{def:II:marked-conventions}
\StatementRecord{def:II:marked-conventions}{proved}{P202}{}
All spaces and coherent sheaves are complex analytic unless a polynomial
ring is specified. We use $B=\PP^1_t$, $(p_1,p_2,p_0)=(0,1,\infty)$,
the ordered lattice
$\Lambda=\Z\langle\widehat\gamma,\widehat u,\widehat w,\widehat\delta\rangle$,
and the marked periods
\[
\Pi_c=\begin{pmatrix}6\mu&\tau&1&0\\\beta^{\rm part}+c&\mu&0&1\end{pmatrix},
\qquad w=e^{-\pi i c},\qquad (\beta^{\rm part}+\tau)(p_0)=0.
\]
Base meridians are counterclockwise. The analytic root decks are clockwise,
$s_j(g_jz)=e^{-2\pi i/m_j}s_j(z)$, with $g_1g_2g_0=1$;
the resulting homology monodromies are $A_1,A_2,M_0$.
The primitive invariant vectors are
$\varepsilon=(1,2,-4,0)^t$ and
$\varepsilon'=(1,3,-3,0)^t$.
For a finite affine datum we write
$v_1=r_1\varepsilon+k_1\widehat\delta$ and
$v_2=r_2\varepsilon'+k_2\widehat\delta$; the coefficients
$r_1,r_2$ are also denoted $l_1,l_2$ in the peripheral presentation.
Define
\[
 D_{\rm part}=\operatorname{Im}\beta^{\rm part}
       -\frac{6(\operatorname{Im}\mu)^2}{\operatorname{Im}\tau},
 \qquad M=\max_{B\setminus\{p_0\}}D_{\rm part}
       =-\frac3\pi\log6.
\]
The admissible parameter inequality is $\operatorname{Im}c<-M$;
the maximum of $D_{\rm part}$ is attained only at $t=1$, as proved in \SeriesRef{fam:src:C.2.3}.
The standard multiplicities are $m_1=3,m_2=4$,
$f^*p_j=m_jS_j$, and $W=f^{-1}(p_0)$ is reduced, with normalization
$\eta:\mathrm dP_6\to W$.
For general finite affine twists we use
$p=12l_0-4l_1-3l_2$, retaining $\pi_1(X)=\Z/|p|$.
The complete marked transport and construction are those of Paper I.
An unqualified fibre is its full Cartier scheme; its reduction is always
specified explicitly. Homology and cohomology retain their displayed
coefficient groups, and complex manifolds carry their complex orientations.
\end{definition}


\section{Index, Picard components and endpoint cohomology}\hypertarget{tag-033X}{}\label{sec:II:01}


\begin{proposition}[Tangent index and Picard components]\hypertarget{tag-01BN}{}\label{cdp:st:C.1.1}
\StatementRecord{cdp:st:C.1.1}{proved}{01BN}{the definitions and constructions stated in the proof}

Let $Y$ be a connected smooth compact complex threefold with
$b_1(Y)=b_2(Y)=0$. Every holomorphic line bundle $L$ on $Y$ satisfies
\[
 \chi(Y,T_Y\otimes L)=\frac12\int_Yc_3(T_Y)
 =1-\frac{b_3(Y)}2.
\]
The orientation in the integral is the complex orientation. The group
$\Pic(Y)$ has finitely many connected components, each a torsor for the
complex vector group $H^1(Y,\OO_Y)$, and
$\pi_0\Pic(Y)\simeq H^2(Y;\Z)$ through the first Chern class.
\end{proposition}


\begin{proof}

\begin{enumerate}[label={$\langle1\rangle$\arabic*.},leftmargin=*]
\item The universal coefficient theorem gives
$H^1(Y;\Z)=\Hom(H_1(Y;\Z),\Z)=0$ from $b_1(Y)=0$.
The group $H^2(Y;\Z)$ is finite because $Y$ has a finite CW model and
$b_2(Y)=0$. In the exponential sequence, normalized by
$\exp(2\pi i\,\cdot)$, the relevant maps are
\[
0=H^1(Y;\Z)\longrightarrow H^1(Y,\OO_Y)
 \xrightarrow{\exp} \Pic(Y)
 \xrightarrow{c_1} H^2(Y;\Z).
\]
They are group homomorphisms; the middle map is holomorphic on the
Picard group. Consequently its image $\Pic^0(Y)$ is isomorphic to the
complex vector space $H^1(Y,\OO_Y)$. The next map in the exponential
sequence is $H^2(Y;\Z)\to H^2(Y,\OO_Y)$; it is zero because its
source is finite and its target is a complex vector space. Thus
$c_1$ is surjective, and the component group is exactly $H^2(Y;\Z)$.

\item Poincar\'e duality gives $H^4(Y;\Q)=0$. Thus
$c_1(T_Y)$, $c_2(T_Y)$ and $c_1(L)$ vanish in rational cohomology.
The degree-six component of
$\operatorname{ch}(T_Y)\exp(c_1(L))\operatorname{td}(T_Y)$ is then
$c_3(T_Y)/2$. The Hirzebruch--Riemann--Roch formula for the Dolbeault
index gives the first equality. Chern--Gauss--Bonnet identifies
$\int_Yc_3(T_Y)$ with $e(Y)$. Since $Y$ is a connected closed oriented
six-manifold, its Betti numbers satisfy
$b_0=b_6=1$, $b_1=b_5=0$, and $b_2=b_4=0$. Hence
$e(Y)=2-b_3(Y)$. The pairing on $H^3(Y;\Q)$ is nondegenerate and
alternating, so $b_3(Y)$ is even, consistently with the integral index.


\end{enumerate}
\end{proof}






\begin{proposition}[A lower bound for twisted tangent cohomology]\hypertarget{tag-01BP}{}\label{cdp:st:C.1.3}
\StatementRecord{cdp:st:C.1.3}{proved}{01BP}{\SeriesRef{cdp:st:C.1.1}}

Let $Y$ be a smooth connected compact complex threefold with
$b_1(Y)=b_2(Y)=0$, carrying a nonzero effective Cartier divisor $D$.
For every $L\in\Pic(Y)$,
\[
 h^2(Y,T_Y\otimes L)\geq
 \max\left\{0,1-\frac{b_3(Y)}2\right\}.
\]
In particular, $b_3(Y)=0$ gives $H^2(Y,T_Y\otimes L)\ne0$ for
every $L$. A threefold in this class possessing a twist with vanishing
second tangent cohomology has $b_3(Y)\geq2$.
\end{proposition}


\begin{proof}
\begin{enumerate}[label={$\langle1\rangle$\arabic*.},leftmargin=*]
\item 
Fix a component $P$ of $\Pic(Y)$ and a representative $L_P$.
Apply \cite[Corollary~1.3, p.~79]{CDP98} to
$T_Y\otimes L_P$ and $T_Y^*\otimes K_Y\otimes L_P^{-1}$.
Tensoring by $L_P$ identifies $\Pic^0(Y)$ with $P$, and
Serre duality identifies the latter degree-zero vanishing with
the degree-three vanishing for the inverse Picard parameter.
The two dense open sets intersect in a dense open subset of $P$.
On this subset, \cref{cdp:st:C.1.1} gives
\[
 h^2(Y,T_Y\otimes L)-h^1(Y,T_Y\otimes L)
 =1-b_3(Y)/2.
\]
The asserted lower bound holds there. By \cref{cdp:st:C.1.1}, the component $P$ is a torsor for
$H^1(Y,\OO_Y)$. Exponentiating additive \v Cech representatives
of this vector space and tensoring by $L_P$ gives a holomorphic
family of line bundles on $Y\times P$. Proper cohomology
semicontinuity makes the locus satisfying the lower bound on
$h^2$ closed analytic in $P$. It contains the dense open subset
above and therefore all of $P$. 

\end{enumerate}
\end{proof}



Statements \ref{cdp:st:C.1.1}--\ref{cdp:st:C.1.3} prove the stated index obstruction with rational Betti-number
vanishing. No integral trivialization of either Chern class is needed
for the index. The existence of the divisor enters in the endpoint
vanishing, and a positive multiple of its actual Cartier divisor
enters the Picard family. These statements supply an index
obstruction; they make no assertion that the failed degree-two
vanishing holds.




\label{cdp:sec:E}

\begin{proposition}[Vertical divisors and the field of meromorphic functions]\hypertarget{tag-01GG}{}\label{cdp:st:E1}
\StatementRecord{cdp:st:E1}{proved}{01GG}{the definitions and constructions stated in the proof}

Let $Y$ be a connected compact complex manifold of dimension $n\geq2$
with $H^2(Y,\mathbb R)=0$. Let $g:Y\to C$ be a surjective holomorphic
map to a connected compact Riemann surface, with connected fibres.
Suppose that its smooth fibres over a nonempty open subset of $C$ are
K\"ahler. Every effective divisor of $Y$ is supported on finitely many
fibres, and pullback induces an isomorphism of fields
\[
 g^*: \mathcal M(C)\xrightarrow{\ \sim\ }\mathcal M(Y).
\]
In particular $a(Y)=1$. 
\end{proposition}


\begin{proof}

\begin{enumerate}[label={$\langle1\rangle$\arabic*.},leftmargin=*]
\item For a prime divisor $D$, the image $g(D)$ is analytic by properness. It
is a point or all of $C$. In the latter case choose a regular value $c$
in the stated open subset for which $F=g^{-1}(c)$ is not contained in
$D$. The canonical section of $\mathcal O_Y(D)$ restricts to a nonzero
section with nonempty effective zero divisor $D_F$ on $F$. For a
K\"ahler form $\omega$ on $F$, the divisor class formula gives
\[
 0=\int_F c_1(\mathcal O_Y(D)|_F)\wedge\omega^{n-2}
   =\int_{D_F}\omega^{n-2}>0.
\]
The first equality follows from the restriction homomorphism
$H^2(Y,\mathbb R)\to H^2(F,\mathbb R)$; the last follows from positivity
on every component with its positive multiplicity. Thus $g(D)$ is a
point. A divisor has finitely many prime components, proving the first
assertion.

\item For $h\in\mathcal M(Y)$ remove from $C$ the finitely many values
supporting its polar divisor. On the inverse image of the complement,
$h$ is holomorphic: a meromorphic function on a manifold having no
codimension-one pole extends across its codimension-two indeterminacy
set by the local normality of the structure sheaf. Its restriction to
each compact connected fibre is constant. The identity
$g_*\mathcal O_Y=\mathcal O_C$, supplied by Stein factorization and
the normality of $C$, therefore gives a holomorphic descended function
there. To check meromorphic extension at the removed values, take the
closure of the graph of $h$ in $Y\times\mathbb P^1$ and apply the proper
map $g\times1$. Its image is a one-dimensional analytic subset of
$C\times\mathbb P^1$, generically the graph of the descended function.
Its component dominating $C$ has generic degree one, and hence is the
graph of a meromorphic function on $C$. Pullback equals $h$ on a dense
open subset and therefore in the meromorphic function field. Pullback
is injective because $g$ is surjective. The field isomorphism preserves transcendence degree.

\end{enumerate}
\end{proof}






\section{Integral specialization, coherent restrictions and component counts}\hypertarget{tag-033Y}{}\label{sec:II:02}


\begin{proposition}[Components of exceptional fibres]\hypertarget{tag-01G0}{}\label{cdp:st:D.1.2}
\StatementRecord{cdp:st:D.1.2}{proved}{01G0}{\SeriesRef{cdp:eq:D1}}

Let $f:Y\to C$ be a proper holomorphic surjection from a connected
compact complex threefold to a compact Riemann surface, with
connected fibres. Let $\Delta$ be a nonempty finite set containing
the critical values, $s=|\Delta|$, $g=g(C)$,
$D=f^{-1}(\Delta)_{\rm red}$, and let $r$ be the number of
irreducible components of $D$. Write $F$ for a smooth fibre and
\[
t'=\dim_{\Q}H_1(F;\Q)_{\pi_1(C\setminus\Delta)}
   =\dim_{\Q}H^1(F;\Q)^{\pi_1(C\setminus\Delta)}.
\]
For $b_2(Y)=0$ the exact component count is
\[
\boxed{\quad r=2g+s-1+t'-b_1(Y).\quad}
\]
Under CDP20's integral assumptions
$H^1(Y;\Z)=H^2(Y;\Z)=0$, the base is $\PP^1$ and
\[
r=s-1+t',\qquad
H_1(F;\Z)_{\pi_1(C\setminus\Delta)}\text{ is free abelian}.
\]
The freeness assertion concerns the displayed coinvariant quotient of
$H_1(F;\Z)$.
\end{proposition}





\begin{proof}
The rational component formula is due to \citet[Section~10.5, p.~91]{AC2026}.

\begin{enumerate}[leftmargin=*]
\item Put $U=Y\setminus D$ and $C^\circ=C\setminus\Delta$.
The map $U\to C^\circ$ is a differentiable fibre bundle. Since
$\pi_2(C^\circ)=0$, its homotopy sequence is
\[
1\longrightarrow\pi_1(F)\longrightarrow\pi_1(U)
 \longrightarrow\pi_1(C^\circ)\longrightarrow1.
\]
The last group is free of rank $2g+s-1$. Lifting its free generators
gives a section of this extension. \citet[Lemma~10.7, p.~90]{AC2026} therefore gives
\begin{equation}
H_1(U;\Z)\simeq
\Z^{2g+s-1}\oplus H_1(F;\Z)_{\pi_1(C^\circ)}.       \hypertarget{tag-019T}{}\label{cdp:eq:D1}
\end{equation}
The second equality defining $t'$ follows by duality: the dual of
the coinvariant quotient is precisely the space of invariant
linear functionals on $H_1(F;\Q)$.

\item The cohomology sequence with compact supports for the open inclusion
$U\subset Y$ contains
\[
H^4(Y;\Q)\longrightarrow H^4(D;\Q)
\longrightarrow H_c^5(U;\Q)
\longrightarrow H^5(Y;\Q)\longrightarrow H^5(D;\Q)=0.
\]
The top-degree cohomology of the compact pure complex surface $D$
is $\Q^r$. One can see this by deleting the singular and intersection
loci, whose real dimensions are at most two, and taking compactly
supported top cohomology of its $r$ oriented smooth components.
Poincar\'e duality identifies $H_c^5(U;\Q)$ with $H_1(U;\Q)$ and
$H^5(Y;\Q)$ with $H_1(Y;\Q)$. Since $H^4(Y;\Q)=0$, taking
dimensions and using \eqref{cdp:eq:D1} proves the formula.

\item For the integral refinement, the universal coefficient theorem shows
first that $H_1(Y;\Z)$ is finite, using $H^1(Y;\Z)=0$, and then that
it is zero, using
$\operatorname{Ext}^1(H_1(Y;\Z),\Z)\subset H^2(Y;\Z)=0$.
The group $H^4(Y;\Z)$ can be finite and nonzero. Its homomorphism
to $H^4(D;\Z)=\Z^r$ is nevertheless zero. The same sequence now
gives $H_1(U;\Z)\simeq\Z^r$, since $H^5(Y;\Z)=H_1(Y;\Z)=0$.
Equation \eqref{cdp:eq:D1} then proves that the coinvariant summand is free.
Finally $H^1(C;\Q)\hookrightarrow H^1(Y;\Q)$ by the Leray edge
map, so $g=0$.

\end{enumerate}
\end{proof}






\begin{proposition}[Two projective torus fibrations]\hypertarget{tag-01G3}{}\label{cdp:st:D.1.5}
\StatementRecord{cdp:st:D.1.5}{proved}{01G3}{\citet[Section~10.5, p.~91]{AC2026}}
Let $E_0=F_0=E'=\C/(\Z+i\Z)$ and $a=1/2\in E_0$.
The involution
\[
(e,z)\longmapsto(e+a,-z)
\]
acts freely on $E_0\times F_0$. Put
$S=(E_0\times F_0)/\langle(a,-1)\rangle$,
$Y=S\times E'$, and $C=E_0/\langle a\rangle$.
The induced map $Y\to C$, denoted by $f$, is a smooth torus fibration with fibre
$F_0\times E'$ and monodromy
$\operatorname{diag}(-I_2,I_2)$ around the nontrivial deck character.
Its integral homology coinvariants are
$\Z^2\oplus(\Z/2\Z)^2$. After removing $s\geq1$ fibres,
\[
H_1(U;\Z)=\Z^{s+3}\oplus(\Z/2\Z)^2.
\]
The sum of the first Betti numbers of the punctured base and the
fibre is $s+5$.

Let $\eta$ be the nontrivial order-two degree-zero line bundle on
$C$ associated to the double cover $E_0\to C$. The coherent direct
images are
\[
R^1f_*\OO_Y=\eta\oplus\OO_C,\qquad R^2f_*\OO_Y=\eta.
\]
The cohomology dimensions are $h^q(Y,\OO_Y)=(1,2,1,0)$.
Restriction from $H^1(Y,\OO_Y)$ to the two-dimensional
$H^1(F_0\times E',\OO)$ has image precisely the $E'$ direction
and has rank one. These homomorphisms are induced by the finite
\'etale quotient and projection, and $b_1(Y)=4$.

For an example over $\PP^1$, let
$h:R\to\PP^1$ be the relatively minimal elliptic surface with
Weierstrass equation
\[
y^2=x^3+(t^4+1)x+t^6.
\]
Its fundamental line bundle is $\OO(1)$ and its discriminant has
twelve simple zeroes. Moreover, $R^1h_*\OO_R=\OO(-1)$.
For $Y'=R\times E'\to\PP^1$, write $f$ for the projection.
The product K\"unneth maps give
\[
R^1f_*\OO_{Y'}=\OO\oplus\OO(-1),\qquad
R^2f_*\OO_{Y'}=\OO(-1),
\]
and restriction in degree one again has rank one. This projective
threefold has $b_1=2$, $b_2=11$ and $a=3$.
\end{proposition}
\begin{proof}
The first construction specializes the bielliptic example of
\citet[Section~10.5, p.~91]{AC2026} to the displayed involution.
Translation by $a$ has no fixed point, so the quotient is smooth.
The map $[e]\mapsto([e,0],0)$ is a section of the fibration. Over
the punctured base its fundamental group is therefore the semidirect
product of the fibre lattice and the free base group of rank $s+1$.
The map $(k,q)\mapsto([q],[k])$ identifies its abelianization with
the direct sum of the base abelianization and the fibre coinvariants:
the relations are $\rho(q)k-k$, which here are twice the two
$F_0$ generators. The same section over the complete elliptic base
gives the stated value of the first Betti number.

On the covering product, the three antiholomorphic one-form
directions have signs $(1,-1,1)$. Descent of the fibre cohomology
gives the two direct images, and taking invariants in the exterior
algebra gives the stated cohomology dimensions. The base direction
restricts to zero on a fibre, while the invariant fibre direction is
exactly that of $E'$. This proves the description of the restriction
map.

For the Weierstrass surface, the coordinates $u=1/t$,
$x'=u^2x$, $y'=u^3y$ give
\[
y'^2=x'^3+(1+u^4)x'+1.
\]
The relative residue differential transforms by
$dx'/(2y')=u^{-1}dx/(2y)$, so the fundamental line is
$\OO_{\PP^1}(1)$. Its discriminant is
\[
-16P(t),\qquad P(t)=31t^{12}+12t^8+12t^4+4=Q(t^4),
\qquad Q(z)=31z^3+12z^2+12z+4.
\]
The identity
\[
\left(\frac5{36}-\frac{31z}{54}\right)Q(z)
+\left(\frac1{27}-\frac{7z}{324}+\frac{31z^2}{162}\right)Q'(z)=1
\]
shows that the roots of $Q$ are simple; its constant term is
nonzero. Thus $P$ has twelve simple zeroes. At infinity the
transformed discriminant is $-16\cdot31$, so the fibre there is
smooth. The twelve singular fibres are nodal cubics. Relative
duality gives the stated first direct image. Leray gives
$h^1(R,\OO_R)=h^2(R,\OO_R)=0$; the twelve nodal fibres give
$e(R)=12$, and hence $b_1(R)=0$ and $b_2(R)=10$.
K\"unneth now gives both direct images and the Betti numbers of
the product. Its degree-one coherent cohomology is the $E'$
summand, giving the restriction map. The product is projective
of dimension three, so its algebraic dimension is three.
\end{proof}


\Uses{the hypotheses, notation and definitions stated here; \CompendiumRef{II}{cdp:st:D.1.1}; \citet[Section~10.5, p.~91]{AC2026}}


The bielliptic-product rank defect is already described in \citet[Section~10.5, p.~91]{AC2026}. The full integral torsion calculation, the coherent direct images and the separate rational-elliptic product below retain their additional precise content.





\begin{proposition}[Components in Wang's torus fibration]\hypertarget{tag-01G4}{}\label{cdp:st:D.1.6}
\StatementRecord{cdp:st:D.1.6}{proved}{01G4}{\SeriesRef{cdp:st:D.1.13}; \cite{Wang2026}}
Wang's threefold $X_2(\lambda)$ has integral homology
$H_0=H_6=\Z$, $H_2=H_3=\Z/2\Z$, and the other reduced
homology groups zero. Its integral cohomology satisfies
$H^1=H^2=0$. Its fibration has four exceptional values, over which
the reduced fibres are irreducible: the multiplicity-three
bielliptic fibre, two rank-one fillings, and the rank-two filling.
In the rank-one filling over the $I_2$ value, the primitive period
exponent is $(2,1)$ and the fibre has one component.

In its cohomology basis the four monodromies are
\[
\begin{aligned}
B_*&=\begin{pmatrix}1&0&-6&-2\\0&-1&-1&1\\0&1&0&-1\\0&0&0&1\end{pmatrix},&\qquad
B_2&=\begin{pmatrix}1&0&6&3\\0&1&2&1\\0&0&1&0\\0&0&0&1\end{pmatrix},\\
B_+&=\begin{pmatrix}1&0&0&0\\0&2&1&0\\0&-1&0&0\\0&0&0&1\end{pmatrix},&\qquad
B_-&=\begin{pmatrix}1&0&0&-1\\0&1&1&0\\0&0&1&0\\0&0&0&1\end{pmatrix}.
\end{aligned}
\]
Their product in the displayed order is the identity, and their
common invariant cohomology group has rank one. On homology the
matrices are $B_i^{-t}$; the Smith diagonal of their difference
matrix is $(1,1,1,0)$. Thus $t'=1$ and the component count is
$r=4-1+1=4$. The expression $s-1+b_1(F)$ equals seven.
The coherent direct images are given in \cref{cdp:st:D.1.13}.
\end{proposition}
\begin{proof}
The construction and its integral homology are
\citet[Theorems~1.1 and~3.1, pp.~2 and~8]{Wang2026}.
The rank-one, rank-two and logarithmic fillings are
\citet[Propositions~2.7, 2.8 and~2.10, pp.~4, 6 and~7]{Wang2026}.
The primitive period vector is given in
\citet[Lemma~2.6, p.~4]{Wang2026}: the monomial change
$(z,v)\mapsto(x,y)=(v,v^2/z)$, with inverse
$(x,y)\mapsto(x^2/y,x)$, sends the vector to $(1,0)$.
The integral universal coefficient sequence gives the asserted
cohomology vanishing, since
$\operatorname{Ext}(H_1,\Z)=\operatorname{Hom}(H_2,\Z)=0$.

The matrices are \citet[Proposition~4.4, equation~(20), p.~11]{Wang2026}.
Direct multiplication gives their ordered product. For a column
$(a,b,c,d)^t$, invariance under $B_-$ gives $c=d=0$, and invariance
under $B_+$ then gives $b=0$. Every matrix fixes the first basis
vector. Preservation of evaluation under transport gives
$B_i^tA_i=I$, so $A_i=B_i^{-t}$ on the homology lattice.
Put $D_i=A_i-I$. Their columns satisfy
\[
D_-e_1=e_4,\qquad D_-e_2=-e_3,\qquad
D_+e_2=-e_2-e_3.
\]
Every image column has first coordinate zero, and these three
columns generate its full kernel. The coinvariant quotient is
therefore identified with $\Z$ by first coordinate, with inverse
$n\mapsto[ne_1]$. This gives the Smith diagonal and the invariant
rank. The component formula of \cref{cdp:st:D.1.2} now gives four,
in agreement with the four irreducible reduced fibres.
\end{proof}


\Uses{the hypotheses, notation and definitions stated here; \cite{Wang2026}; \cref{cdp:st:D.1.13}}





\begin{proposition}[Real cohomology and the structure sheaf]\hypertarget{tag-01G5}{}\label{cdp:st:D.1.7}
\StatementRecord{cdp:st:D.1.7}{proved}{01G5}{the definitions and constructions stated in the proof}

For a connected compact complex space $S$, including a nonreduced
one, the coefficient map
\[
H^1(S;\mathbb R)\longrightarrow H^1(S,\OO_S)
\]
is injective as a map of real vector spaces.
\end{proposition}


\begin{proof}
First take $S$ reduced. Represent a real cohomology class by a
locally constant real \v{C}ech cocycle $a_{ij}$. If its image is zero,
after refining the open cover there are holomorphic functions $u_i$
with $a_{ij}=u_i-u_j$. The imaginary parts of the $u_i$ agree on
overlaps and give a global continuous pluriharmonic function.
Pull it back to a compact resolution of each irreducible component.
The maximum principle makes each such pullback constant; intersection
of components and connectedness of $S$ make the constants equal.
Thus, after subtracting one imaginary constant from all $u_i$,
each $u_i$ is real-valued and holomorphic, hence locally constant.
They form a real coboundary for $a_{ij}$. For nonreduced $S$, a
class zero in $H^1(S,\OO_S)$ also becomes zero after applying
$\OO_S\to\OO_{S_{\rm red}}$. The underlying topological space
and its real cohomology do not change under reduction, so the
reduced assertion proves the result.

\end{proof}






\begin{proposition}[Injectivity of integral generization]\hypertarget{tag-01G6}{}\label{cdp:st:D.1.8}
\StatementRecord{cdp:st:D.1.8}{proved}{01G6}{\SeriesRef{cdp:st:D.1.2}}

Let $f:Y\to C$ be as in \ref{cdp:st:D.1.2}, and let $j:C^\circ\hookrightarrow C$
be the smooth-fibre locus. The natural generization homomorphism
\[
R^1f_*\Z_Y\longrightarrow
j_*\bigl((R^1f_*\Z_Y)|_{C^\circ}\bigr)
\]
is injective. Every nonzero global section of $R^1f_*\Z_Y$ is
nonzero in $H^1(f^{-1}(c);\Z)$ at every $c\in C$.
\end{proposition}


\begin{proof}

\begin{enumerate}[leftmargin=*]
\item It suffices to check a critical value $c$. Represent a stalk class
by $a\in H^1(V;\Z)$ for $V=f^{-1}(\Delta_c)$ and a sufficiently
small disc $\Delta_c$. Put $V^*=f^{-1}(\Delta_c\setminus\{c\})$.
Assume that the restriction of $a$ to a smooth fibre is zero.
The homotopy sequence for $V^*\to\Delta_c^*$ shows that
$a|_{V^*}$ is the pullback of an integer multiple $k$ of the
positive meridian class of the punctured disc.

\item Choose a smooth point of one irreducible component of the central
divisor, away from all its other components. There the map has
the form $t=z^m u$ with $m>0$ and $u$ a holomorphic unit.
A small positive normal circle has base winding $m$ and bounds
its normal disc in $V$. Evaluating $a$ on that circle gives
$km=0$, whence $k=0$.
The inclusion $V^*\hookrightarrow V$ induces a surjection of
fundamental groups: a loop can be perturbed off a complex divisor,
including its higher-codimension singular strata. Therefore
$a|_{V^*}=0$ implies $a=0$.
Taking the direct limit over the discs proves the asserted stalk
injection; proper topological base change identifies its source
with $H^1(f^{-1}(c);\Z)$.

\item If a global section vanished at one point, it would vanish on a
neighbourhood of that point. Its restriction to the connected
smooth locus is a section of a local system, so it would vanish
everywhere on that locus. The just-proved injection then makes
it zero everywhere, as required.

\end{enumerate}
\end{proof}






\begin{proposition}[Line bundles of infinite order on every fibre]\hypertarget{tag-01G7}{}\label{cdp:st:D.1.9}
\StatementRecord{cdp:st:D.1.9}{proved}{01G7}{\SeriesRef{cdp:eq:D5}, \SeriesRef{cdp:st:D.1.8}, \SeriesRef{cdp:st:D.1.7}}

Let $Y$ be a connected compact complex manifold with
$b_1(Y)=b_2(Y)=0$, and let $f:Y\to\PP^1$ be a proper
holomorphic surjection with connected fibres. There is a class
$\beta\in H^1(Y,\OO_Y)$ such that, for every irrational real number
$r$, the line bundle
\[
L_r=\exp(2\pi i\,r\beta)
\]
has infinite order on every scheme-theoretic fibre of $f$.
Every connected component of $\Pic(Y)$ likewise contains a real
one-parameter family with this property at every irrational real
parameter.

\end{proposition}


\begin{proof}

\begin{enumerate}[leftmargin=*]
\item The low-degree integral Leray sequence and connectedness of the
fibres give an injection with finite cokernel
\begin{equation}
d_2:H^0(\PP^1,R^1f_*\Z_Y)
 \hookrightarrow H^2(\PP^1;\Z)=\Z.                   \hypertarget{tag-019X}{}\label{cdp:eq:D5}
\end{equation}
Indeed $H^1(Y;\Z)=0$ and $H^2(Y;\Z)$ is finite.
Thus the image in \eqref{cdp:eq:D5} is $d\Z$ for an integer $d>0$; for the
stronger integral hypothesis $H^2(Y;\Z)=0$ one has $d=1$.
Choose $\gamma$ mapped to $d$. The local meridian argument in
the proof of Proposition~\ref{cdp:st:D.1.8} applies in every complex
dimension: the local equation is $t=z^m u$ with $u$ a unit and $m>0$,
and evaluation on a positive normal circle gives $km=0$, hence $k=0$.
Thus each $\gamma_c\in H^1(f^{-1}(c);\Z)$ is nonzero.
The coefficient morphism $\Z_Y\to\OO_Y$ sends $\gamma$ to a
global section of $R^1f_*\OO_Y$. Coherent Leray and
$H^1(\PP^1,\OO)=H^2(\PP^1,\OO)=0$ identify this space with
$H^1(Y,\OO_Y)$; let $\beta$ be its inverse image.

\item For every full fibre $Y_c$, naturality of the edge maps and of
restriction gives
\[
\beta|_{Y_c}=\iota_c(\gamma_c)\in H^1(Y_c,\OO_{Y_c}), 
\]
where $\iota_c$ is the coefficient map. This is the image under the natural coherent base-change
homomorphism. If $L_r|_{Y_c}$ had finite order, the exponential
sequence on the full fibre would give an integer $n>0$ and
$\lambda\in H^1(Y_c;\Z)$ with
\[
nr\,\iota_c(\gamma_c)=\iota_c(\lambda).
\]
Proposition~\ref{cdp:st:D.1.7} implies $nr\gamma_c=\lambda$ in real cohomology.
The group $H^1(Y_c;\Z)$ is a finitely generated free abelian group,
and $\gamma_c$ is a nonzero element of it. A nonzero integral
coordinate of $\gamma_c$ now implies that $r$ is rational, a
contradiction. This works for every $c$ and proves the simultaneous
assertion in $\Pic^0(Y)$.

\item For a component represented by a line bundle $A$, choose $N>0$
with $A^N\in\Pic^0(Y)$ and write $A^N=\exp(2\pi i b)$.
Then
\[
A_r=A\otimes\exp\bigl(2\pi i(r\beta-b)/N\bigr)
\]
belongs to that component and satisfies $A_r^N=L_r$. A group
element has finite order exactly when any fixed positive power of
it has finite order. Therefore each $A_r|_{Y_c}$ is non-torsion
for every irrational real $r$, proving the componentwise statement.

\end{enumerate}
\end{proof}






\begin{proposition}[Invariant classes in the restriction image]\hypertarget{tag-01FV}{}\label{cdp:st:D.1.10}
\StatementRecord{cdp:st:D.1.10}{proved}{01FV}{\SeriesRef{cdp:st:D.1.2}, \SeriesRef{cdp:st:D.1.7}}

Let $f:Y\to\PP^1$ be a holomorphic algebraic reduction of a compact
complex threefold with $b_1(Y)=b_2(Y)=0$ and $a(Y)=1$.
At a general smooth fibre, the image of the restriction
$H^1(Y,\OO_Y)\to H^1(F,\OO_F)$ contains the complex span of
the coefficient homomorphism
\begin{equation}
\iota_F:H^1(F;\Z)^{\pi_1(C^\circ)}
 \longrightarrow H^1(F,\OO_F),                         \hypertarget{tag-019W}{}\label{cdp:eq:D4}
\end{equation}
which is induced by $\Z_F\hookrightarrow\OO_F$.
The final spanning step of CDP's rank argument reaches the whole
target exactly when this span is $H^1(F,\OO_F)$. This characterization concerns the subspace generated by the
indicated invariant lattice in the restriction image.
\end{proposition}


\begin{proof}

\begin{enumerate}[leftmargin=*]
\item Let $D_1,\ldots,D_r$ be the prime components above $\Delta$, and
choose $N>0$ annihilating the finite group $H^2(Y;\Z)$.
The logarithms, in the exponential identification
$\Pic^0(Y)=H^1(Y,\OO_Y)$, of
$\OO_Y(ND_1),\ldots,\OO_Y(ND_r)$ generate a free abelian
subgroup of rank $r-s+1=t'$.
Here is the rank computation with its maps specified. The map is
$\Z^r\to\Pic^0(Y)$, $(n_i)\mapsto\OO_Y(N\sum n_iD_i)$.
A relation means that $N\sum n_iD_i$ is the divisor of a
meromorphic function. Since $f$ is the algebraic reduction, that
function is $f^*h$ for $h\in\mathcal M(\PP^1)$.
Its divisor is a sum of full fibres with total base degree zero.
Such relations have rank $s-1$, and those given by divisors of
rational functions supported on $\Delta$ already have this rank.
Statement \ref{cdp:st:D.1.2} gives the last equality.

\item Over $U=f^{-1}(C^\circ)$ each of these divisor line bundles has
its canonical trivialization. Exactness of the exponential sequence
therefore places its logarithm in the coefficient image of
$H^1(U;\Z)$. Coherent Leray identifies restriction
$H^1(Y,\OO_Y)\to H^1(U,\OO_U)$ with restriction of global
sections of $E$; this is injective because \cite[Theorem~2.2, pp.~4--5]{Viaclovsky2026Combs} makes $E$ locally free.
Thus the rank-$t'$ group of logarithms remains injective on $U$.

\item Topological Leray on the punctured curve gives
\[
0\to H^1(C^\circ;\Z)\to H^1(U;\Z)
\to H^1(F;\Z)^{\pi_1(C^\circ)}\to0.
\]
The first group's coefficient image vanishes because
$H^1(C^\circ,\OO)=0$. Conversely a class whose coefficient image
in $H^1(U,\OO)$ vanishes restricts to zero in $H^1(F,\OO_F)$;
Proposition~\ref{cdp:st:D.1.7} makes its integral restriction to $F$ zero. Thus the
coefficient image of $H^1(U;\Z)$ identifies exactly with the
invariant lattice in \eqref{cdp:eq:D4}. The logarithm group maps to a subgroup
of the same rank $t'$, hence of finite index, in that lattice.
Its complex span equals the span of the entire lattice, and the
restriction image is a complex vector subspace. This proves the
assertion and the statement about the last spanning step.

\end{enumerate}
\end{proof}






\begin{proposition}[Generic line bundles on every fibre]\hypertarget{tag-01FW}{}\label{cdp:st:D.1.11}
\StatementRecord{cdp:st:D.1.11}{proved}{01FW}{\SeriesRef{cdp:st:D.1.9}, \SeriesRef{cdp:st:D.1.7}, \SeriesRef{cdp:st:D.1.8}}

Let $f:Y\to\PP^1$ be a proper connected-fibre holomorphic map from
a connected compact complex manifold with $b_1(Y)=b_2(Y)=0$.
In every component of $\Pic(Y)$, the line
bundles which are non-torsion on every full fibre form a dense
$G_\delta$ set of full Lebesgue measure in its complex vector-space
coordinates. If evaluation of global sections of $R^1f_*\OO_Y$
has rank one, then $\Pic^0(Y)=\C$ and the exceptional set is
countable. For two-torus fibres this alternative holds whenever
the splitting has a negative summand.
This gives the general-line-bundle conclusion of
\citet[Proposition~5.6]{CDP2020} under the displayed hypotheses.
\end{proposition}


\begin{proof}

\begin{enumerate}[leftmargin=*]
\item Let $E=R^1f_*\OO_Y$, a vector bundle by \cite[Theorem~2.2, pp.~4--5]{Viaclovsky2026Combs}, and use
$\gamma,\beta$ from \ref{cdp:st:D.1.9}. The section of $E$ corresponding to
$\beta$ has no zero. In fact its value maps naturally to the
nonzero class $\iota_c(\gamma_c)$ in the cohomology of every full
fibre; Proposition~\ref{cdp:st:D.1.7} and the dimension-independent local
meridian argument in the proof of Proposition~\ref{cdp:st:D.1.8}
prove nonvanishing. Write
$E=\bigoplus_j\OO(a_j)$ and let $k$ count the nonnegative degrees.
Evaluation has constant rank $k$ onto their direct-summand subbundle
$E_+$. Nonvanishing of $\beta$ gives $k\geq1$. When $k=1$, its
sole nonnegative degree is zero, since a positive-degree line bundle
has no nowhere-zero section. In that case $H^0(E)=\C\beta$.

\item Suppose first $k\geq2$, and put $V=H^1(Y,\OO_Y)=H^0(E)$ and
$N=\dim_{\C}V$. On a simply connected disc in the smooth-fibre
locus where degree-one base change holds, trivialize the integral
local system. For a flat integral
class $\lambda$, the equation
\[
z|_{F_c}=\iota_c(\lambda),\qquad (z,c)\in V\times\Delta
\]
defines an analytic set. The parameter values where
$\iota_c(\lambda)\in(E_+)_c$ form an analytic subset of the disc;
over it the fibres are affine spaces of dimension $N-k$.
The incidence set thus has complex dimension at most
$N-k+1\leq N-1$. Its projection to $V$ has Lebesgue
measure zero. It is also meagre: stratify and exhaust that set
by countably many compact sets; the images are closed sets with
empty interior, by their real dimension at most $2N-2$.
Take countably many discs, integral classes and positive integer
denominators to include every torsion condition.

\item There are finitely many excluded fibres, including the exceptional
values of degree-one base change. At each, the natural map
$E\otimes\C(c)\to H^1(Y_c,\OO_{Y_c})$ is injective. This follows
directly from the long exact sequence for multiplication by a local
base parameter: its cokernel is the parameter-torsion in
$R^2f_*\OO_Y$. The restriction $V\to H^1(Y_c,\OO_{Y_c})$
therefore has rank $k\geq2$. The torsion conditions are contained in
countably many affine subspaces of complex codimension at least
two. Combining these with the preceding projections gives a
meagre null set. The complement is dense and $G_\delta$, as claimed.

\item It remains to prove the case $k=1$.
Now $V=\C\beta$. On a simply connected disc in the same open locus,
the projection of a flat integral class $\lambda$ to
$E/(\OO\beta)$ is holomorphic. Unless it vanishes identically,
its zero set is discrete. Such zeroes give only countably many
parameters $z$ for which $\iota_c(\lambda)=z\beta_c$.
We prove that an identically vanishing projection gives a constant
parameter, including after analytic continuation.

\item Pass to the universal cover of the smooth locus, on which
$H^1(F;\Z)$ is marked. The identity theorem continues the relation
$\iota_c(\lambda)\in\C\iota_c(\gamma)$ everywhere on this cover.
If $\lambda$ and $\gamma$ are real-dependent, their ratio is a
constant rational number. Otherwise let
$W_{\Q}=\Q\gamma+\Q\lambda$. The coefficient map
$H^1(F;\R)\to H^1(F,\OO_F)$ is real-linearly injective by \ref{cdp:st:D.1.7}.
The images of the two real-independent classes span the entire
real two-plane underlying $\C\beta_c$, so $W_{\R}$ is precisely
its inverse image. Analytic continuation around a base loop
sends $\lambda$ to a class whose coefficient image again belongs
to this line. Hence it preserves $W_{\Q}$ and fixes $\gamma$.

\item Choose an integral basis $(\gamma_0,\lambda_0)$ of
$W_{\Q}\cap H^1(F;\Z)$ with $\gamma$ a positive integral multiple
of $\gamma_0$. The induced monodromy on this lattice fixes
$\gamma_0$ and preserves the orientation of $W_{\R}$: the complex
structure induced on that plane by the coefficient map varies
continuously along the loop, so its complex orientation cannot change.
It therefore sends $\lambda_0$ to $\lambda_0+k\gamma_0$ for an
integer $k$. Put
\[
\tau(c)=\frac{\iota_c(\lambda_0)}{\iota_c(\gamma_0)}.
\]
This holomorphic function on the universal cover never takes a
real value, by real injectivity of the coefficient map, and its
imaginary part has constant sign. Monodromy changes $\tau$ by
integers. Accordingly either $\exp(2\pi i\tau)$ or
$\exp(-2\pi i\tau)$ descends to a holomorphic function of absolute
value less than one on $\PP^1$ minus finitely many points. The
removable singularity theorem extends it to $\PP^1$, where it is
constant. Its logarithmic derivative gives $d\tau=0$. Thus the
original ratio is constant as asserted.

\item There are countably many integral classes on a countable disc
cover. The isolated-zero alternatives and the constant alternatives
just proved give a countable exceptional set on smooth fibres.
Each exceptional full fibre contributes another countable set,
because $z\mapsto z\beta_c$ is injective and its intersection with
the rational integral coefficient image is countable. This proves
the assertion in $\Pic^0(Y)$.

\item Finally for any component $A\Pic^0(Y)$, the map taking the $N$th
power, for an integer $N$ killing its component class, is the affine
isomorphism $z\mapsto b+Nz$ of the complex parameter space used
in \ref{cdp:st:D.1.9}. It preserves the torsion condition in every fibre, and
preserves countability, meagreness and measure-zero sets. This
proves all the componentwise assertions.

\end{enumerate}
\end{proof}



The quantifier in \ref{cdp:st:D.1.11} is explicit. In the globally generated
two-torus case, the torsion parameters at even one fixed smooth fibre are
dense: the rational period lattice is dense in
$H^1(F,\OO_F)$, and restriction onto that space is surjective.
Thus the good set in that case cannot contain a nonempty open
subset of the Picard parameter space. A dense $G_\delta$ set of
full measure is the appropriate conclusion of the incidence
argument, rather than a Zariski-open conclusion.




\begin{proposition}[Cup products and cohomological flatness]\hypertarget{tag-01FX}{}\label{cdp:st:D.1.12}
\StatementRecord{cdp:st:D.1.12}{proved}{01FX}{\cite[\S3, equations (9)--(12)]{Borcea1984}}

For a proper connected-fibre holomorphic map $f:Y\to C$ from a
compact complex threefold with $b_2(Y)=0$ and general fibre a
two-torus or a primary Kodaira surface, the cup-product homomorphism
is an isomorphism
\[
\bigwedge^2R^1f_*\OO_Y\xrightarrow{\simeq}R^2f_*\OO_Y.
\]
Every full fibre has structure-sheaf cohomology dimensions
$(h^0,h^1,h^2)=(1,2,1)$, and these two direct images satisfy
base change at every fibre. This proves these two cases of CDP20
Proposition 5.1 by an index and cup-product argument.
\end{proposition}





\begin{proof}
The cup-product index argument for torus-fibred six-spheres is due to
\citet[Theorem~2.3, pp.~5--6]{Viaclovsky2026Combs}. We apply it under
the displayed hypotheses and verify the cup product for primary Kodaira fibres.

\begin{enumerate}[leftmargin=*]
\item Write $E=R^1f_*\OO_Y$, which is locally free of rank two by \cite[Theorem~2.2, pp.~4--5]{Viaclovsky2026Combs},
and $G=R^2f_*\OO_Y$, which has rank one. Put
$d=\deg E$, $g=\deg(G/G_{\rm tor})$, and
$\ell=\operatorname{length}(G_{\rm tor})$.
Since $c_1(Y)=0$ rationally, Hirzebruch--Riemann--Roch gives
$\chi(Y,\OO_Y)=\int_Yc_1c_2/24=0$.
Coherent Leray and Riemann--Roch on $C$ give, with the rank terms
cancelled,
\[
0=\chi(\OO_C)-\chi(E)+\chi(G)=g+\ell-d.
\]
The cup product $\det E\to G$ is an isomorphism on a general
fibre. For a two-torus this is its exterior-algebra description.
For a primary Kodaira surface, use invariant $(1,0)$-forms
$\alpha,\eta$ with $d\alpha=0$ and
$d\eta=c\alpha\wedge\overline\alpha$, $c\ne0$.
Their conjugates are $\overline\partial$-closed. The class
$[\overline\alpha\wedge\overline\eta]$ is nonzero because its
Serre pairing with the holomorphic two-form $\alpha\wedge\eta$
is the nonzero integral of a volume form. This also proves that
$[\overline\alpha],[\overline\eta]$ are independent; since
$h^1(\OO)=2$ and $h^2(\OO)=1$, it proves the required cup-product
isomorphism. The invariant forms and their differential are
the forms in \cite[\S3, equations (9)--(12)]{Borcea1984}.
The composite of $\det E\to G$ with the torsion-free quotient is a
nonzero homomorphism of line bundles, so $g\geq d$. The displayed
identity and $\ell\geq0$ force $g=d$, $\ell=0$. A nonzero map
of line bundles of equal degree has no zero divisor, proving the
claimed global isomorphism.

\item Top-degree base change gives
$G\otimes\C(c)\simeq H^2(Y_c,\OO_{Y_c})$, of dimension one.
The multiplication-parameter exact sequence then identifies
$E\otimes\C(c)$ with $H^1(Y_c,\OO_{Y_c})$, since $G$ has no
parameter-torsion. \cite[Theorem~2.2, pp.~4--5]{Viaclovsky2026Combs} gives $h^0=1$. These arguments retain the
scheme-theoretic fibres throughout.

\end{enumerate}
\end{proof}






\begin{proposition}[Restriction maps and splitting degrees]\hypertarget{tag-01G2}{}\label{cdp:st:D.1.4}
\StatementRecord{cdp:st:D.1.4}{proved}{01G2}{\SeriesRef{cdp:st:D.1.2}}

For any $f:Y\to C$ as in \ref{cdp:st:D.1.2} put $E=R^1f_*\OO_Y$.
The coherent Leray sequence gives
\begin{equation}
0\longrightarrow H^1(C,\OO_C)\xrightarrow{f^*}H^1(Y,\OO_Y)
\xrightarrow{e}H^0(C,E)\longrightarrow0.               \hypertarget{tag-019U}{}\label{cdp:eq:D2}
\end{equation}
At a smooth fibre where cohomological base change holds, the
restriction homomorphism is exactly the composite
\begin{equation}
H^1(Y,\OO_Y)\xrightarrow{e}H^0(C,E)
\xrightarrow{\operatorname{ev}_c}E\otimes\C(c)
\xrightarrow{\simeq}H^1(F_c,\OO_{F_c}).               \hypertarget{tag-019V}{}\label{cdp:eq:D3}
\end{equation}
By \cite[Theorem~2.2, pp.~4--5]{Viaclovsky2026Combs}, $E$ is locally free. For $C=\PP^1$, write
$E=\bigoplus_{j=1}^q\OO(a_j)$.
Then the rank of \eqref{cdp:eq:D3} at a general point is exactly
$\#\{j:a_j\geq0\}$. In particular, the nonnegative splitting
degrees claimed by CDP20 Proposition 5.5 are equivalent to its
claimed general-fibre surjectivity when $E$ has rank two and is
locally free.
\end{proposition}


\begin{proof}
Stein factorization gives $f_*\OO_Y=\OO_C$. A curve has no coherent
cohomology in degree two, so the low-degree Leray sequence is \eqref{cdp:eq:D2}.
The edge map is induced by restriction to inverse images of open
sets in $C$. Restricting these same representatives to $F_c$ gives
\eqref{cdp:eq:D3}, which proves the asserted identification of the maps.
The proof of \cite[Theorem~2.2, pp.~4--5]{Viaclovsky2026Combs} establishes local freeness by resolving the fibre support,
normalizing a cyclic base change with reduced central fibre, and
taking the invariant direct image. For a line bundle
$\OO(a)$ on $\PP^1$, its evaluation has rank one at every point
for $a\geq0$, and its space of global sections is zero for $a<0$.
Applying this separately to the direct summands gives the formula.

\end{proof}






\begin{proposition}[Higher direct images of Wang's threefold]\hypertarget{tag-01FY}{}\label{cdp:st:D.1.13}
\StatementRecord{cdp:st:D.1.13}{proved}{01FY}{\SeriesRef{cdp:st:D.1.6}, \SeriesRef{cdp:st:D.1.12}, \SeriesRef{cdp:st:D.1.9}, \SeriesRef{cdp:st:D.1.4}, \SeriesRef{cdp:st:E1}}


For Wang's threefold $Y=X_2(\lambda)$ in \ref{cdp:st:D.1.6} one has
\[
K_Y\simeq f^*\OO_{\PP^1}(-1),\qquad
R^1f_*\OO_Y\simeq\OO\oplus\OO(-1),\qquad
R^2f_*\OO_Y\simeq\OO(-1).
\]
Consequently $h^q(Y,\OO_Y)=(1,1,0,0)$, and restriction to every
smooth two-torus fibre has rank one. This threefold satisfies
$H^1(Y;\Z)=H^2(Y;\Z)=0$ and $a(Y)=1$. Its rank-one
restriction map and direct-image splitting contradict the conclusions of
\cite[Corollary~5.2 and Proposition~5.5, pp.~688--689]{CDP20}.
For the threefold $Y_0$ of \SeriesRef{maps:src:B.5.2}, one has
$H_2(Y_0;\Z)=0$, whereas $H_2(Y;\Z)=\Z/2\Z$; hence a
biholomorphism $Y_0\to Y$ would induce an impossible isomorphism of
these two homology groups.
\end{proposition}


\begin{proof}

\begin{enumerate}[label={$\langle1\rangle$\arabic*.},leftmargin=*]
\item We compute the canonical line through the local fillings.
The rational elliptic surface $\pi:S\to\PP^1$ in this construction
has Weierstrass equation
\[
y^2=x^3-3(1+t)x+(2+3t+2t^2)
\]
and fundamental line bundle $\OO(1)$; hence
$K_S=\pi^*\OO(-1)$. On the complement of the zero section in
$M=\OO_S(P-O)$, the determinant of the relative cotangent sequence
identifies its canonical line with the pullback of $K_S$.
In a local multiplicative coordinate $v$, this identification sends
a local two-form $\omega_S$ to $\omega_S\wedge d\log v$.
Changing $v$ by a nonzero holomorphic function on $S$ adds a
one-form from $S$ to $d\log v$, whose wedge with $\omega_S$ is zero.
This verifies the coordinate independence of the identification.
The contracting generator covers translation by $-6P$ on the
elliptic fibres. Translation preserves the relative elliptic
differential, and the extra base terms disappear after wedging
with the base differential. Its line-bundle lift changes $d\log v$
by a one-form pulled back from $S$. It therefore preserves this canonical
identification. The latter descends to
$K_Y=f^*\OO(-1)$ on the open quotient.

\item At either rank-one filling, let $q$ be the base coordinate and
$(x,y)$ the two multiplicative coordinates of the construction.
At $t=0$ the explicit change is $x=v$, $y=v^2/z$; its exponent
matrix is unimodular. At $\alpha_+$ no change is required.
Thus the preceding three-form is a nonzero holomorphic unit times
\[
dq\wedge d\log x\wedge d\log y.
\]
The nondegenerating $y$ direction becomes an elliptic curve after
quotienting by its contracting period. Its logarithmic differential
is a nonvanishing relative one-form; changes by a base differential
do not affect the displayed wedge. At a double curve the equation
is $uv=q$, and
$dq\wedge d\log u=-du\wedge dv$. Therefore the three-form extends
without a zero or pole over the full rank-one filling.

\item At $\alpha_-$ the rank-two filling uses a unimodular period lattice
and the standard toric charts $q=z_0z_1z_2$. A pair of independent
fibre characters, together with $q$, is a unimodular system of
characters in each smooth chart. The same form is, up to a nonzero
constant and a holomorphic unit,
\[
q\,d\log z_0\wedge d\log z_1\wedge d\log z_2
=dz_0\wedge dz_1\wedge dz_2.
\]
It extends without zero or pole, including at the triple points.

\item It remains to check the order-three logarithmic filling at infinity.
Put $s=1/t$, $x_s=s^2x$ and $y_s=s^3y$. Its Weierstrass
coefficients have orders three and four in $s$, giving type $IV^*$.
On the cubic base change $s=r^3$, put
$x_s=r^4X$, $y_s=r^6Y$. The good-reduction equation is
\[
Y^2=X^3-3r(r^3+1)X+(2+3r^3+2r^6),
\]
whose central elliptic curve is smooth. The canonical form satisfies
\[
ds\wedge\frac{dx_s}{y_s}
=3\,dr\wedge\frac{dX}{Y}.
\]
It has no zero or pole. For $\zeta=e^{2\pi i/3}$ the root deck
action has weights $r\mapsto\zeta r$, $X\mapsto\zeta^2X$,
$Y\mapsto Y$; the weights on $dr$ and $dX/Y$ multiply to one.
The additional invariant translation used in the logarithmic
transform preserves the holomorphic two-form on the torus cover.
Hence the three-form is invariant and descends through the free
cyclic action as a nonvanishing three-form. This proves that the
canonical identification from the open quotient extends everywhere,
giving $K_Y=f^*\OO(-1)$.

\item The map is flat and Gorenstein, since its source is smooth and its
base is a smooth curve. Relative Serre duality gives
\[
(R^2f_*\OO_Y)^*=f_*\omega_{Y/\PP^1}
=f_*f^*\OO(1)=\OO(1).
\]
Proposition~\ref{cdp:st:D.1.12} gives local freeness and
$\det(R^1f_*\OO_Y)=R^2f_*\OO_Y=\OO(-1)$.
The nowhere-zero section from \ref{cdp:st:D.1.9} embeds $\OO$ as a subbundle
of $R^1f_*\OO_Y$ with quotient $\OO(-1)$. The extension splits
because $H^1(\PP^1,\OO(1))=0$. Coherent Leray now computes all
four displayed structure-sheaf cohomology groups, and \ref{cdp:st:D.1.4}
identifies the degree-one restriction with rank-one evaluation.

\item The homology computation in \ref{cdp:st:D.1.6} and the universal coefficient
theorem give $H^1(Y;\Z)=H^2(Y;\Z)=0$. More explicitly,
$H_1(Y;\Z)=0$ gives
$H^2(Y;\Z)\simeq\operatorname{Hom}(H_2(Y;\Z),\Z)=0$,
and the universal-coefficient injection in degree three is
$\operatorname{Ext}^1(H_2(Y;\Z),\Z)\simeq\Z/2\Z
\hookrightarrow H^3(Y;\Z)$.
The general fibre is K\"ahler and $b_2(Y)=0$, so \ref{cdp:st:E1} proves that
this map is the algebraic reduction and $a(Y)=1$.
The threefold of \SeriesRef{maps:src:B.5.2} has zero second
integral homology, whereas $H_2(Y;\Z)=\Z/2\Z$. A biholomorphism
and its inverse induce mutually inverse homomorphisms on $H_2$;
the displayed groups therefore exclude such a biholomorphism.

\end{enumerate}
\end{proof}






\begin{proposition}[Constant integral monodromy and coherent splitting]\hypertarget{tag-01FZ}{}\label{cdp:st:D.1.15}
\StatementRecord{cdp:st:D.1.15}{proved}{01FZ}{the definitions and constructions stated in the proof}


Two smooth proper families of two-tori over $\PP^1$ have the same
constant marked $H^1(-;\Z)$ local system but have
$R^1f_*\OO=\OO^2$ and $R^1f_*\OO=\OO(1)^2$, respectively.
Thus the integral representation alone does not determine the
coherent splitting for general compact torus fibrations. Both
examples have $b_1=4$.
\end{proposition}


\begin{proof}

\begin{enumerate}[leftmargin=*]
\item Let $V=\OO_{\PP^1}(1)^{\oplus2}$. On the affine coordinate $t$,
four global sections of $V$ are the columns of
\[
P(t)=\begin{pmatrix}1&i&t&it\\-t&it&1&-i\end{pmatrix}.
\]
They are a real basis in every fibre. Indeed, writing
$z_1=a+ib$, $z_2=c+id$, their real linear combination is
$(z_1+t z_2,\overline z_2-t\overline z_1)$.
Its vanishing gives $(1+|t|^2)\overline z_2=0$ and then $z_1=0$.
At infinity, the $\OO(1)$ frame changes $P(t)$ to $t^{-1}P(t)$,
whose leading columns again form a real basis. The lattice
generated by these sections therefore acts freely and properly by
holomorphic fibre translations on $V$. Its quotient $Z$ is a
compact smooth holomorphic torus bundle over $\PP^1$.
The displayed real basis identifies it differentiably with
$\PP^1\times(\R^4/\Z^4)$ and marks its constant integral local
system.

\item The fibrewise period map includes the bundle of holomorphic
one-forms by $P(t)^t$. The Dolbeault coefficient map gives the
exact sequence
\[
0\longrightarrow\OO(-1)^2
 \xrightarrow{P(t)^t}\OO^4
 \longrightarrow R^1f_*\OO_Z\longrightarrow0.
\]
The last map can be represented explicitly by
\[
B(t)=\begin{pmatrix}-it&-t&i&1\\-i&1&-it&t\end{pmatrix}.
\]
The identity $B(t)P(t)^t=0$ holds polynomially. Its two rows have
rank two at every finite $t$, as is seen from their two-by-two
minors, and their leading coefficients have rank two at infinity.
Thus this is the exact complex of vector bundles
\[
0\longrightarrow\OO(-1)^2\longrightarrow\OO^4
 \xrightarrow{B}\OO(1)^2\longrightarrow0,
\]
proving $R^1f_*\OO_Z\simeq\OO(1)^2$. More explicitly, with
\[
J=\begin{pmatrix}0&1&0&0\\-1&0&0&0\\0&0&0&1\\0&0&-1&0\end{pmatrix},
\qquad S=\begin{pmatrix}0&1\\1&0\end{pmatrix},
\]
one has $B=SPJ$ and $PJP^t=0$, so $BP^t=0$.

\item For the constant family $Z_0=Z_{t=0}\times\PP^1$, the same
marked fibre lattice gives the identical constant integral local
system, whereas coherent K\"unneth gives
$R^1f_*\OO_{Z_0}=H^1(Z_{t=0},\OO)\otimes\OO=\OO^2$.
The different splittings are therefore compared by an explicit
identity of the integral representations and explicit coherent
quotient maps.

\end{enumerate}
\end{proof}






\section{Complete smooth-fibre classification and actual Hopf occurrence}\hypertarget{tag-033Z}{}\label{sec:II:03}


\begin{proposition}[A secondary Kodaira surface with torsion-free first homology]\hypertarget{tag-01F8}{}\label{cdp:st:C.3.2}
\StatementRecord{cdp:st:C.3.2}{proved}{01F8}{the definitions and constructions stated in the proof}


There is a secondary Kodaira surface $Q$ with
\[
H_1(Q;\Z)=\Z,\qquad H^2(Q;\Z)=0,
\]
whose canonical bundle has holomorphic order six.
\end{proposition}


\begin{proof}

\begin{enumerate}[label={$\langle1\rangle$\arabic*.},leftmargin=*]
\item Let $H$ be the real Heisenberg group in coordinates $(x,y,z)$ with law
\[
(x,y,z)(x',y',z')
=\left(x+x',y+y',z+z'+\frac{xy'-yx'}2\right).
\]
Let $\Gamma\subset H$ be generated by
$a=(1,0,0)$, $b=(0,1,0)$ and $c=(0,0,1)$. Thus
\[
\Gamma=\{(x,y,z):x,y\in\Z,\ z-xy/2\in\Z\},
\qquad [a,b]=c.
\]
It is a discrete cocompact subgroup. On $H\times\mathbb R_t$ write
\[
\tau=e^{\pi i/3},\quad
\eta=dz-\frac{x\,dy-y\,dx}{2},\quad
\alpha=dx+\tau\,dy,\quad
\beta=\eta-\frac{dx+dy}{2}+i\,dt.
\]
These two complex one-forms and their conjugates are linearly
independent over $\C$. Declaring $\alpha,\beta$ to have type $(1,0)$
defines a left-invariant almost complex structure. It is integrable:
$d\alpha=0$ and $d\beta=-dx\wedge dy$ has type $(1,1)$, so the
ideal generated by the $(1,0)$ forms has no $(0,2)$ obstruction.
The structure descends to
$P=\Gamma\backslash H\times(\mathbb R/\Z)$.
The nonzero form $\alpha\wedge\beta$ is $\bar\partial$-closed and
trivializes $K_P$. The projection induced by $x+\tau y$ is a
holomorphic elliptic bundle over $\C/(\Z+\tau\Z)$, with nonzero
central Chern class $[a,b]=c$. Its first Betti number is three, so
$P$ is a primary Kodaira surface.

\item Define the Lie-group automorphism
\[
\sigma(x,y,z)=(-y,x+y,z-y/2).
\]
It preserves $\Gamma$: the difference between the new central
coordinate and half the product of the new horizontal coordinates
is
\[
(z-xy/2)+xy+y(y-1)/2.
\]
This is integral on $\Gamma$. The inverse preserves the lattice too,
and direct substitution gives $\sigma^6=1$. Moreover
$\sigma^*\alpha=\tau\alpha$ and $\sigma^*\beta=\beta$.
Consequently
\[
g:P\longrightarrow P,\qquad
g([x,y,z],t)=([\sigma(x,y,z)],t+1/6)
\]
is a holomorphic automorphism of order six. Every nonidentity power
is fixed-point-free because of its translation of the $t$-circle.
The quotient $Q=P/\langle g\rangle$ is therefore a compact complex
surface and $P\to Q$ is finite \emph{\'etale} of degree six.

\item The contractible manifold $H\times\mathbb R$ is its universal cover.
The deck group is $\Gamma\rtimes_\sigma\Z$, where the infinite
generator is the lift $(\sigma,t\mapsto t+1/6)$; its sixth power is
the original integral translation in $t$. Since
$\Gamma^{\rm ab}=\Z^2$ and the induced matrix is
\[
A=\begin{pmatrix}0&-1\\1&1\end{pmatrix},\qquad
\det(A-I)=1,
\]
abelianization gives $H_1(Q;\Z)=\Z$.
The finite cover and $e(P)=0$ give $e(Q)=0$. Poincar\'e duality then
gives $b_2(Q)=0$, and the universal coefficient theorem gives
$H^2(Q;\Z)=0$ because $H_1(Q;\Z)$ has no torsion.
Finally $g^*(\alpha\wedge\beta)=\tau\alpha\wedge\beta$.
Every holomorphic section of $K_P^m$ is a constant multiple of
$(\alpha\wedge\beta)^m$, by compactness. Such a section descends
precisely when $6$ divides $m$. Thus $K_Q$ has holomorphic order six.
The surface has $b_1=1$ and a primary Kodaira surface as a finite
\'etale cover, so it is secondary Kodaira.

\end{enumerate}
\end{proof}






\begin{proposition}[A secondary Hopf surface with torsion-free first homology]\hypertarget{tag-01F9}{}\label{cdp:st:C.3.3}
\StatementRecord{cdp:st:C.3.3}{proved}{01F9}{the definitions and constructions stated in the proof}


There is a secondary Hopf surface $Q_H$ with
$H_1(Q_H;\Z)=\Z$ and $H^2(Q_H;\Z)=0$.
\end{proposition}


\begin{proof}

\begin{enumerate}[leftmargin=*]
\item Write a quaternion as $(q_0,q_1,q_2,q_3)$ and put
$\varphi=(1+\sqrt5)/2$. Let $G$ be the following set of 120 unit
quaternions: the eight signed coordinate units; the sixteen points
$(\pm1,\pm1,\pm1,\pm1)/2$; and the 96 points obtained by even
coordinate permutations of
$(0,\pm1,\pm\varphi,\pm\varphi^{-1})/2$, with independent signs.
Quaternion multiplication makes this the binary icosahedral group.
Multiplication of the listed coordinates is closed in $\Q(\sqrt5)^4$,
and the commutator map on this finite group is surjective.
Every element is therefore a commutator, and $G^{\rm ab}=0$.

\item Use the faithful holomorphic representation
\[
(a,b,c,d)\longmapsto
\begin{pmatrix}a+ib&c+id\\-c+id&a-ib\end{pmatrix}\in SU(2)
\]
and let the generator of $\Z$ act on $\C^2\setminus\{0\}$ by
$v\mapsto2v$. The action of $\Z\times G$ is free. An equality
$2^ngv=v$ first gives $n=0$ by norms, and a unit quaternion whose
matrix has eigenvalue one is the identity. Proper discontinuity
follows from the radial scaling and finiteness of $G$.
Polar coordinates identify the quotient smoothly with
\[
Q_H=(\C^2\setminus\{0\})/(2^{\Z}\times G)
\simeq (\mathbb R/(\log2)\Z)\times(S^3/G).
\]
It is compact, and its universal cover is $\C^2\setminus\{0\}$.
The fundamental group is $\Z\times G$, whose abelianization is
$\Z$. The quotient by $2^{\Z}$ alone is a primary Hopf surface
covering $Q_H$ with degree 120. The fundamental group of $Q_H$ is
noncyclic, so $Q_H$ is secondary Hopf. The product description gives
$e(Q_H)=0$, and the same duality and universal-coefficient argument
as before gives $H^2(Q_H;\Z)=0$.

\end{enumerate}
\end{proof}



These two explicit finite-cover maps locate two additional failures
in the surface-level exclusions used in the proof of CDP20,
Lemma 4.3. They do not construct these surfaces as fibres of a
threefold satisfying the standing hypotheses. Such a realization or
a separate exclusion remains part of C.3.




\begin{proposition}[Obstruction to K\"ahler resolutions]\hypertarget{tag-01FA}{}\label{cdp:st:C.3.4}
\StatementRecord{cdp:st:C.3.4}{proved}{01FA}{the definitions and constructions stated in the proof}

Let $Y$ be a smooth compact complex threefold with $b_2(Y)=0$, and let
$f:Y\to C$ be proper, holomorphic, surjective, with connected fibres.
Every irreducible component of a reducible reduced fibre has no
compact K\"ahler resolution. In particular, a fibre all of whose
components admit compact K\"ahler resolutions is irreducible.
\end{proposition}


\begin{proof}
\begin{enumerate}[label={$\langle1\rangle$\arabic*.},leftmargin=*]
\item 
Let $D$ be a component of a reducible fibre. The incidence graph of
the components is connected, so some other component $D'$ meets $D$.
Both are Cartier divisors on the smooth threefold. Let
$\mu:\widehat D\to D\hookrightarrow Y$ be a resolution by a compact
smooth surface. Pullback of the canonical section of $\OO_Y(D')$
is a nonzero holomorphic section of $\mu^*\OO_Y(D')$. Its zero
divisor $E$ is effective and nonzero: the intersection $D\cap D'$
is a nonempty divisor on the irreducible surface $D$, and pullback
to the resolution preserves its nonzero vanishing order along a
component dominating that intersection. On the other hand
$c_1(\mu^*\OO_Y(D'))=0$ in $H^2(\widehat D;\R)$.
A K\"ahler form $\omega$ on $\widehat D$ would give
\[
0=\int_{\widehat D}c_1(\OO_{\widehat D}(E))\wedge\omega
 =\int_E\omega>0,
\]
where positive multiplicities are retained in the last integral.
This contradiction proves the assertion for every component.

\end{enumerate}
\end{proof}






\begin{proposition}[Exclusion of bielliptic fibres]\hypertarget{tag-01FD}{}\label{cdp:st:C.3.7}
\StatementRecord{cdp:st:C.3.7}{proved}{01FD}{\cite[Proposition~2, pp.~605--608]{Kawai1969}; \cite[Theorem~11, p.~297]{Buchdahl1999}}


A connected smooth compact complex threefold $Y$ with $b_1(Y)=b_2(Y)=0$ has no
proper connected-fibre holomorphic map to a compact curve whose
general fibre is bielliptic.
\end{proposition}


\begin{proof}

Suppose that $f:Y\to C$ is a map with the stated properties.
Since $Y$ is reduced, $f$ factors as $f=\iota\circ f_{\mathrm{red}}$
through the inclusion $\iota:C_{\mathrm{red}}\hookrightarrow C$ of the
reduced image curve. Let $\nu:C^\nu\to C_{\mathrm{red}}$ be its
normalization. The reduced closure of the regular-locus lift graph in
$Y\times_{C_{\mathrm{red}}}C^\nu$ is finite and birational over $Y$.
Since $Y$ is normal, its projection to $Y$ is an isomorphism. The other
projection gives a holomorphic map $f^\nu:Y\to C^\nu$ with
$f_{\mathrm{red}}=\nu\circ f^\nu$. This map is surjective because its
proper image contains the regular open set of $C^\nu$. For each
$c\in C_{\mathrm{red}}$, the finitely many closed sets
$(f^\nu)^{-1}(c')$, for $c'\in\nu^{-1}(c)$, partition $f^{-1}(c)$;
each is also open in that fibre. Connectedness of $f^{-1}(c)$ forces
a single nonempty piece, so $f^\nu$ has connected fibres. Over the
regular locus $\nu$ is an isomorphism, and the general bielliptic fibre
is unchanged. Write $f:Y\to C$ for the lifted map $f^\nu:Y\to C^\nu$ in the
remainder of the proof.

\begin{enumerate}[label=\arabic*.,leftmargin=*]
\item A bielliptic surface has
$p_g=0$ and $h^0(\Omega^1)=1$. The relative Albanese construction of
\cite[Proposition~2, pp.~605--608]{Kawai1969} gives a compact complex surface
$B$, a holomorphic map $h:B\to C$, and a meromorphic map
$\Phi:Y\dashrightarrow B$ over $C$. On the complement $C^\circ$
of finitely many values, $B^\circ\to C^\circ$ is a smooth family
of elliptic curves and
\[
 \Phi|_{F_c}:F_c\longrightarrow B_c
\]
is the Albanese morphism of the bielliptic surface, with connected
elliptic fibres. The
construction forms the relative period-lattice quotient, glues its
local Albanese maps by fibre translations, and compactifies at each
missing value by a normalized finite base change and the projective
space associated with the direct images of a divisor which is the
inverse image of a local Albanese section. The sum of the finite orbit of that divisor is invariant under
the covering group. Resolve the normalized base change first if needed, so the
divisor defines a line bundle. The orbit sum restricts on every
general fibre to the pullback of a positive-degree divisor on its
Albanese elliptic curve; sufficiently large powers give the same
Albanese image. Its naturally linearized line bundle and direct
images carry the finite covering-group action, and their quotient
glues the construction to $B^\circ$.

\item Resolve $B$ over $C\setminus C^\circ$ and resolve the induced
meromorphic map by a sequence of blowups with smooth centres
supported over those same values. This produces a diagram
\[
 \widetilde Y\xrightarrow{\ g\ }\widetilde B,
 \qquad \pi:\widetilde Y\longrightarrow Y,
 \qquad (h\circ\rho)\circ g=f\circ\pi,
\]
where $\rho:\widetilde B\to B$ is a resolution and both spaces
are smooth and compact. The map $g$ has connected fibres:
its general fibres are the connected Albanese fibres just described,
and its Stein factorization has finite degree one over the normal
surface $\widetilde B$. Hence
$H^1(\widetilde B;\Q)\hookrightarrow H^1(\widetilde Y;\Q)$.
Blowing up smooth centres does not change first cohomology, so
$b_1(\widetilde B)=0$. The compact-surface K\"ahler criterion
\cite[Theorem~11, p.~297]{Buchdahl1999} therefore gives a K\"ahler form
$\omega_B$ on $\widetilde B$.

\item The blowup cohomology formula gives
\[
 H^2(\widetilde Y;\R)
 =\pi^*H^2(Y;\R)\oplus\bigoplus_j\R[E_j]
 =\bigoplus_j\R[E_j],
\]
where the $E_j$ are the exceptional prime-divisor classes (with
the corresponding total transforms for successive blowups).
Every $E_j$ is supported above one of the removed values in $C$.
Consequently its restriction to a general smooth fibre $F_c$,
identified with its unchanged inverse image in $\widetilde Y$,
is zero. In particular
\[
 [g^*\omega_B]|_{F_c}=0\quad\hbox{in }H^2(F_c;\R).
\]
Choose a K\"ahler form $\omega_F$ on the bielliptic surface $F_c$.
The Albanese morphism $g|_{F_c}$ is a submersion to the elliptic
curve $\widetilde B_c$. Thus its pullback of $\omega_B$ is a
nonzero semipositive $(1,1)$-form and
\[
 \int_{F_c}g^*\omega_B\wedge\omega_F>0.
\]
The vanishing cohomology class makes the same integral zero by
Stokes' theorem, a contradiction.

\end{enumerate}
\end{proof}






\begin{proposition}[Exclusion of elliptic Hopf general fibres]\hypertarget{tag-01FE}{}\label{cdp:st:C.3.8}
\StatementRecord{cdp:st:C.3.8}{proved}{01FE}{\SeriesRef{cdp:st:C.3.1}; \cite[Proposition~3.1 and its proof, pp.~6--7]{SS2002}}


The very general Hopf fibre of the algebraic reduction in \ref{cdp:st:C.3.1}
has algebraic dimension zero. In particular neither primary nor
secondary Hopf surfaces of algebraic dimension one can be its
general fibre.
\end{proposition}


\begin{proof}

\begin{enumerate}[label=\arabic*.,leftmargin=*]
\item A primary Hopf surface $H$ with $a(H)=1$ satisfies
$\kappa(H,K_H^{-1})=1$. The contraction
normal forms are diagonal, or
\[
A(z,w)=(\alpha z+\lambda w^n,\beta w),\qquad
\lambda\ne0,\quad \alpha=\beta^n.
\]
The invariant bivectors $zw\,\partial_z\wedge\partial_w$ in the
diagonal case and $w^{n+1}\,\partial_z\wedge\partial_w$ in the
second case give a nonzero effective anticanonical divisor.
The invariance is an equality under the derivative of the
displayed contraction: its determinant is $\alpha\beta$, and
the displayed coefficient transforms by the same scalar.
The normal forms and elliptic-fibration description are given in \cite[Proposition~3.1 and its proof, pp.~6--7]{SS2002}.

\item The algebraic reduction of this primary elliptic Hopf surface is a
holomorphic map to $\PP^1$, with smooth reductions of all fibres.
Every curve is vertical: its intersection with a general fibre
is zero since $H^2(H;\R)=0$, whereas a horizontal curve would
have positive intersection equal to its degree over the base.
Thus the nonzero anticanonical divisor is a positive sum of
reduced fibres. If $m_iD_i$ is the full fibre supported on
$D_i$, each $\OO(m_iD_i)$ is the pullback of $\OO_{\PP^1}(1)$.
A common multiple of the $m_i$ therefore makes a positive power
of $K_H^{-1}$ the pullback of a positive-degree line bundle on
$\PP^1$. Its Iitaka dimension is one.

\item For a secondary Hopf surface take a normal finite primary cover
$p:H'\to H$. The cover is unramified, so
$K_{H'}^{-1}=p^*K_H^{-1}$, and $a(H')=a(H)=1$ by invariance of
transcendence degree under a finite map. The anticanonical
section ring of $H$ is the invariant subring of that of $H'$.
The latter is integral over the invariant subring: a homogeneous
section $s$ satisfies the monic orbit polynomial
$\prod_{g\in\operatorname{Gal}(H'/H)}(T-gs)$. Their fraction
fields have the same transcendence degree, proving
$\kappa(H,K_H^{-1})=1$ also in this case.

\item For each positive integer $m$, the locus
\[
Z_m=\{c\in C^\circ:
 h^0(Y_c,K_Y^{-m}|_{Y_c})\ge2\}
\]
is a closed analytic subset of the smooth base $C^\circ$.
If the very general Hopf fibre had algebraic dimension one,
the preceding calculation would make the union of these
countably many loci contain the complement of a countable
subset of $C^\circ$. A proper closed analytic subset of a
curve is discrete and countable, so some $Z_m$ is the whole
connected smooth base. Coherent base change then gives
$\operatorname{rank}f_*K_Y^{-m}\ge2$.
The relative linear system defines a meromorphic map of $Y$
to $\PP(f_*K_Y^{-m})$ over $C$ with positive-dimensional image
on a general fibre. Its image has dimension at least two.
The projective bundle over the compact curve is projective,
so pulling back its meromorphic functions gives $a(Y)\ge2$,
contrary to the hypothesis. In fact the same argument shows
that, when the smooth fibres are Hopf surfaces, the values with
$a(Y_c)=1$ form at most a countable set.

\end{enumerate}
\end{proof}






\begin{proposition}[Exclusion of secondary Kodaira general fibres]\hypertarget{tag-01FF}{}\label{cdp:st:C.3.9}
\StatementRecord{cdp:st:C.3.9}{proved}{01FF}{\SeriesRef{cdp:st:D.1.9}, \SeriesRef{cdp:st:D.1.7}, \SeriesRef{cdp:eq:C.SK.1}, \SeriesRef{cdp:eq:C.SK.2}, \SeriesRef{cdp:eq:C.SK.3}, \SeriesRef{cdp:eq:C.SK.4}, \SeriesRef{cdp:eq:C.SK.5}; \cite[Theorem~4 and Remark~1, pp.~377, 380]{Borcea1984}; \cite[\S1, pp.~2--3]{SS2002}}


A connected smooth compact complex threefold $Y$ with $b_1(Y)=b_2(Y)=0$ and
$a(Y)=1$ cannot have a holomorphic algebraic reduction
$f:Y\to\PP^1$ whose general fibre is a secondary Kodaira
surface.
\end{proposition}


\begin{proof}

\begin{enumerate}[label=\arabic*.,leftmargin=*]
\item Suppose that the smooth general fibres $S_c$ are secondary
Kodaira surfaces. They have $h^1(\OO_{S_c})=b_1(S_c)=1$.
The nonzero integral transgression section of \ref{cdp:st:D.1.9} has
coefficient image $\beta\in H^1(Y,\OO_Y)$ and gives a
nowhere-zero section of $\mathscr E=R^1f_*\OO_Y$.
This coherent sheaf is a line bundle by \cite[Theorem~2.2, pp.~4--5]{Viaclovsky2026Combs}. Thus
$\mathscr E=\OO_{\PP^1}$ and
$H^1(Y,\OO_Y)=\C\beta$. For every $z\in\C$ set
\[
 M_z=\exp(2\pi\mathrm i z\beta)\in\Pic^0(Y).
\tag{C.SK.1}\hypertarget{tag-011H}{}\label{cdp:eq:C.SK.1}
\]
Its restriction satisfies $M_{z+1}|_{S_c}\simeq M_z|_{S_c}$,
because $\beta|_{S_c}$ is the coefficient image of an
integral class. It is nonzero as a tangent vector of
$\Pic^0(S_c)$ by \ref{cdp:st:D.1.7}. The exponential sequence therefore
identifies $\Pic^0(S_c)$ with $\C^*$ and makes \eqref{cdp:eq:C.SK.1}, on
each smooth fibre, a nonconstant homomorphism from
$\C/\Z$ to this group.

\item Let $m>1$ be the order of $K_{S_c}$, and let
\[
 p_c:P_c\longrightarrow S_c
\]
be its connected cyclic canonical cover, with deck group
$G=\mu_m$. It is finite \emph{\'etale}. Its canonical bundle
is trivial, and it is non-K\"ahler, because a finite
\emph{\'etale} quotient of a K\"ahler surface is K\"ahler
by averaging a K\"ahler form. The classification of
surfaces with trivial canonical bundle consequently makes
$P_c$ a primary Kodaira surface. Write
$\pi_c:P_c\to E_c$ for its Albanese elliptic bundle, and
$d>0$ for its topological degree. These covers exist in
holomorphic families locally on the smooth base: the
evaluation of a local frame of $f_*K_{Y/\PP^1}^{m}$
trivializes $K_{Y/\PP^1}^{m}$, since it is nowhere zero on
each fibre, and taking its cyclic root gives the required
finite \emph{\'etale} covering family. The integers $m,d$
are locally constant and hence constant on the connected
smooth base.

\item The primary-surface invariant is
\[
 Q_c=\Pic^0(P_c)/\pi_c^*\Pic^0(E_c)\simeq\C^*.
\tag{C.SK.2}\hypertarget{tag-011I}{}\label{cdp:eq:C.SK.2}
\]
The Picard pullback in the denominator is injective:
$\pi_{c*}\OO_{P_c}=\OO_{E_c}$ and the projection formula
show that a line with trivial pullback is trivial. For
any line $L_c$ of degree $d$ on $E_c$, the pullback
$\pi_c^*L_c$ belongs to $\Pic^0(P_c)$, and its class in
$Q_c$ is independent of $L_c$. Choose the group
isomorphism $q_c:Q_c\to\C^*$ that places this
positive-degree class in the punctured unit disc, and
denote its value by $\alpha(c)$.
This fixes the possible inversion ambiguity of the
group coordinate. The descriptions in
\cite[Theorem~4 and Remark~1, pp.~377, 380]{Borcea1984} and
\cite[\S1, pp.~2--3]{SS2002} identify $\alpha$ as the
intrinsic holomorphic punctured-disc coordinate of the
moduli of primary Kodaira surfaces. More explicitly,
$P_c$ is the associated elliptic bundle of a line of
degree $d$ with fibre $\C^*/\langle\alpha(c)\rangle$;
quotienting by $\pi_c^*\Pic^0(E_c)$ removes the choice
of that line within its degree. In a holomorphic
family the local integral marking and the coefficient
map give this quotient a holomorphic $\C^*$ coordinate,
and the chosen positive-degree orientation is locally
constant. Borcea's complete period family gives these local holomorphic
coordinates. Thus $\alpha(c)$ is holomorphic.

\item The locally constructed canonical covers have the same
isomorphism class on overlaps, so this intrinsic
function glues on the smooth base. It has
$0<|\alpha(c)|<1$. It extends across the finitely many
omitted values by the removable singularity theorem
and is constant on $\PP^1$. Write its nonzero constant
value as $\alpha$.

\item Consider the exact homomorphism of connected Picard
groups
\[
 \Pic^0(S_c)\xrightarrow{\ p_c^*\ }\Pic^0(P_c)
          \longrightarrow Q_c\xrightarrow{\ q_c\ }\C^*.
\tag{C.SK.3}\hypertarget{tag-011J}{}\label{cdp:eq:C.SK.3}
\]
Its differential is nonzero. Indeed finite-cover trace
identifies $H^1(S_c,\OO_{S_c})$ with the $G$-invariant
line in $H^1(P_c,\OO_{P_c})$. The other line,
$\pi_c^*H^1(E_c,\OO_{E_c})$, has nontrivial $G$ action:
otherwise the unique holomorphic one-form of $P_c$
would descend to $S_c$, whereas a secondary Kodaira
surface has no holomorphic one-form. Hence the invariant
line has zero intersection with the denominator's
tangent line in \eqref{cdp:eq:C.SK.2}.
Consequently \eqref{cdp:eq:C.SK.3} is a nonzero isogeny of $\C^*$
groups. Composing it with \eqref{cdp:eq:C.SK.1} gives
\[
 q_c\bigl([p_c^*(M_z|_{S_c})]\bigr)
       =\exp(2\pi\mathrm i k z)
       \qquad(k\in\Z\setminus\{0\}).
\tag{C.SK.4}\hypertarget{tag-011K}{}\label{cdp:eq:C.SK.4}
\]
Here integrality follows from the period $z\mapsto z+1$;
every holomorphic group homomorphism $\C^*\to\C^*$ is
an integral power. The integer $k$ is locally constant
in the holomorphic families just described, and is
therefore a single integer on the connected smooth base.

\item Let $\rho_c:S_c\to\PP^1$ be its elliptic algebraic
reduction. The induced finite map $r_c:E_c\to\PP^1$
satisfies $\rho_c p_c=r_c\pi_c$. Write $h=\deg r_c$;
this positive integer is constant. Choose one positive
integer $N$ killing the torsion group $H^2(S_c;\Z)$
and satisfying $d\mid hN$. This choice is uniform
because the smooth fibres are differentiably locally
trivial. Then
$A_c=\rho_c^*\OO_{\PP^1}(N)$ belongs to $\Pic^0(S_c)$,
and \eqref{cdp:eq:C.SK.2} gives
\[
 q_c([p_c^*A_c])=\alpha^{hN/d}.
\tag{C.SK.5}\hypertarget{tag-011L}{}\label{cdp:eq:C.SK.5}
\]
This identity follows by taking degrees on the Albanese elliptic
curve.

\item Choose $z_0\in\C$ with
$\exp(2\pi\mathrm i k z_0)=\alpha^{hN/d}$ and put
$M=M_{z_0}$. By \eqref{cdp:eq:C.SK.4}--\eqref{cdp:eq:C.SK.5}, the difference of
$p_c^*(M|_{S_c})$ and $p_c^*A_c$ is
$\pi_c^*R_c$ for some $R_c\in\Pic^0(E_c)$.
Both of the first two lines come from $S_c$ and are
$G$-invariant. Injectivity of $\pi_c^*$ implies that
$R_c$ is fixed by the induced nontrivial linear action
on $\Pic^0(E_c)$. An elliptic-curve linear automorphism
of finite order other than one has order
$r\in\{2,3,4,6\}$. If it fixes a point $R_c$, then
\[
 rR_c=(1+\zeta+\cdots+\zeta^{r-1})R_c=0
\]
in the elliptic group. Thus $R_c^{\otimes12}$ is
trivial. The norm of the finite cover satisfies
$\operatorname{Nm}_{p_c}(p_c^*L)=L^{\otimes m}$.
Applying it to the difference gives the line-bundle identity
\[
 M^{\otimes12m}|_{S_c}
       \simeq \rho_c^*\OO_{\PP^1}(12mN)
\tag{C.SK.6}\hypertarget{tag-011M}{}\label{cdp:eq:C.SK.6}
\]
for every smooth general fibre.

\item The right side has at least two linearly independent
sections, by pullback from $\PP^1$. Coherent generic
base change gives
$\operatorname{rank}f_*M^{\otimes12m}\ge2$.
Its relative linear system has a positive-dimensional
image on a general fibre and therefore a surface
image in the projective bundle over $\PP^1$.
Pullback of meromorphic functions from this image
gives $a(Y)\ge2$, the desired contradiction.

\end{enumerate}
\end{proof}






\begin{proposition}[Exclusion of primary Kodaira general fibres]\hypertarget{tag-01F5}{}\label{cdp:st:C.3.10}
\StatementRecord{cdp:st:C.3.10}{proved}{01F5}{\SeriesRef{cdp:st:C.3.9}, \SeriesRef{cdp:st:D.1.12}, \SeriesRef{cdp:st:D.1.9}, \SeriesRef{cdp:st:D.1.7}, \SeriesRef{cdp:eq:C.PK.4}, \SeriesRef{cdp:st:D.2.3}, \SeriesRef{cdp:eq:C.PK.5}, \SeriesRef{cdp:eq:C.PK.1}, \SeriesRef{cdp:eq:C.PK.2}, \SeriesRef{cdp:eq:C.PK.7}; \cite[pp.~373--374]{Borcea1984}; \cite[\S1]{SS2002}; \cite[Remark~1]{Borcea1984}; \cite[\S3]{Borcea1984}; \cite[Theorem~2.18 and \S\S2.19--2.21]{Steenbrink1976}}


A smooth compact complex threefold with $b_1(Y)=b_2(Y)=0$ and
$a(Y)=1$ cannot have a holomorphic algebraic reduction
$f:Y\to\PP^1$ whose general fibre is a primary Kodaira
surface. Together with \ref{cdp:st:C.3.9} this excludes both Kodaira
types using only rational Betti-number vanishing.
\end{proposition}


\begin{proof}

\begin{enumerate}[leftmargin=*]
\item Let $B=\PP^1\setminus\Sigma$ be the smooth base and write
$P_c=f^{-1}(c)$ and $\pi_c:P_c\to E_c$ for the canonical
elliptic algebraic reduction of a smooth fibre. The degree
$d>0$ is constant. All the smooth fibres have the fixed
primary Kodaira differentiable type; every complex structure
on that type is again a primary Kodaira structure
\cite[pp.~373--374]{Borcea1984}.
Use the intrinsic quotient and its orientation from \eqref{cdp:eq:C.SK.2}:
\[
 Q_c=\Pic^0(P_c)/\pi_c^*\Pic^0(E_c),\qquad
 q_c:Q_c\xrightarrow{\simeq}\C^*.
\]
The class of the pullback of a degree-$d$ line has coordinate
$\alpha(c)\in\Delta^*$. The holomorphic moduli comparison
in the proof of \ref{cdp:st:C.3.9} applies directly here, so this bounded
function extends to $\PP^1$ and has a single nonzero
constant value $\alpha$.

\item Put $\mathscr E=R^1f_*\OO_Y$ and
$\mathscr G=R^2f_*\OO_Y$. The primary Kodaira case of
\ref{cdp:st:D.1.12} proves, on the original threefold and every full fibre,
\begin{equation}
 \operatorname{rank}\mathscr E=2,\qquad
 \operatorname{rank}\mathscr G=1,\qquad
 \bigwedge^2\mathscr E\xrightarrow{\simeq}\mathscr G.
\hypertarget{tag-00V8}{}\label{cdp:eq:C.PK.1}
\end{equation}
Both sheaves are locally free and satisfy base change.
Let $\gamma$ be the integral transgression section of \ref{cdp:st:D.1.9},
and let $\beta\in H^0(\PP^1,\mathscr E)=H^1(Y,\OO_Y)$
be its coefficient image. This section is nowhere zero.
For $u\in H^1(P_c,\OO_{P_c})$ define the complex-linear
functional $\lambda_c$ by the exact identity
\begin{equation}
 q_c\bigl([\exp(2\pi\mathrm i z u)]\bigr)
     =\exp(2\pi\mathrm i z\lambda_c(u))
     \quad(z\in\C).
\hypertarget{tag-00V9}{}\label{cdp:eq:C.PK.2}
\end{equation}
Its kernel is the line
$A_c=\pi_c^*H^1(E_c,\OO_{E_c})$.
The function $k(c)=\lambda_c(\beta_c)$ takes integral
values, because $\beta_c$ is the coefficient image of
an integral class. It is holomorphic, by the actual relative
Picard quotient in \eqref{cdp:eq:C.SK.2}, hence is a constant integer $k$.

\item When $k\ne0$, choose $n>0$ with $nd\ge2$ and choose
$z\in\C$ with $\exp(2\pi\mathrm i kz)=\alpha^n$.
The ambient line $M=\exp(2\pi\mathrm i z\beta)$ then
has, on every smooth fibre, the same $Q_c$ class as the
pullback of a degree-$nd$ line on $E_c$. Its difference
from that pullback belongs to $\pi_c^*\Pic^0(E_c)$.
Thus
$M|_{P_c}=\pi_c^*L_c$ for a line $L_c$ of degree $nd$,
and $h^0(P_c,M|_{P_c})=nd\ge2$. The relative linear
system gives $a(Y)\ge2$, as at the end of \ref{cdp:st:C.3.9}.
\item For $k=0$, $\beta_c\in A_c$. The coefficient map
$H^1(E_c;\R)\to H^1(E_c,\OO_{E_c})$ is an isomorphism
of real vector spaces. Its pullback, together with the real
injectivity of \ref{cdp:st:D.1.7} on $P_c$, proves the exact implication
\begin{equation}
 \beta_c\in A_c
 \quad\Longrightarrow\quad
 \gamma_c\in\pi_c^*H^1(E_c;\R).
\hypertarget{tag-00VA}{}\label{cdp:eq:C.PK.3}
\end{equation}
The lattice $\pi_c^*H^1(E_c;\Z)$ is saturated in
$H^1(P_c;\Z)$: the elliptic bundle induces a surjection
on fundamental groups, so the cokernel of the induced
injection on integral $H^1$ is torsion free. Thus $\gamma$
lies in this rank-two integral local system. Divide it by
its fixed positive divisibility to obtain a primitive fixed
section $\gamma_0$. On a simply connected open set choose
a complementary integral section $\delta$ with oriented
intersection number one. The ratio of their nonzero
coefficient images is a holomorphic elliptic period
\[
 \tau=\frac{\iota(\delta)}{\iota(\gamma_0)},
 \qquad \operatorname{Im}\tau>0,
\]
after choosing the sign convention of the ordered pair.
Monodromy fixes $\gamma_0$ and replaces $\delta$ by
$\delta+N\gamma_0$, $N\in\Z$. Therefore
$\exp(2\pi\mathrm i\tau)$ is a single-valued bounded
holomorphic function on $B$. It extends to $\PP^1$ and
is constant. Its logarithmic derivative gives $d\tau=0$,
and a nonzero integral translation cannot fix this constant
period. Consequently the Albanese integral $H^1$ local
system has identity monodromy, and $j(E_c)$ is constant.

\item Both moduli $(j,\alpha)$ and the degree $d$ are now
constant. Borcea's classification, in the precise form
\cite[\S1]{SS2002}, makes all smooth fibres biholomorphic
to one primary Kodaira surface $P$. Borcea's complete
holomorphic marked period family \cite[Remark~1]{Borcea1984}
expresses our family, near any smooth base point, as the
pullback by a holomorphic map into his period domain
$\mathcal D\subset\PP(H^2(P;\C))$. By his Theorem~3,
constant isomorphism class means that the image of this
map lies in one orbit of the integral cohomology action.
That orbit is countable, since the action factors through
$\operatorname{GL}(H^2(P;\Z)/\operatorname{torsion})$.
A connected subset of a countable subset of a complex
manifold is one point: a real coordinate separating two
points would have a nontrivial interval as its connected
image. Therefore the period parameter is constant and
its pullback is a product. 
\item Write the fundamental group as
\begin{equation}
 \Gamma=\langle x,y,w,u\mid
       w,u\text{ central},\ [x,y]=w^d\rangle.
\hypertarget{tag-00VB}{}\label{cdp:eq:C.PK.4}
\end{equation}
Its centre $\langle w,u\rangle$ is the fundamental group
of the elliptic fibre, and its quotient is the fundamental
group of the elliptic base. A biholomorphism inducing the
identity on the base $H^1$ induces the identity on this
quotient. The relation in \eqref{cdp:eq:C.PK.4} then forces it to fix
$w$. On an elliptic fibre its linear part is a holomorphic
group automorphism. Such an automorphism acts by a complex
scalar on the universal cover; fixing the nonzero period
$w$ forces that scalar to be one. It therefore fixes both
$w$ and $u$. The remaining possible actions on $x,y$ are
exactly multiplication by central elements.

\item In the real
nilmanifold presentation choose invariant one-forms
$a,b,c,e$ with
\[
 da=db=de=0,\qquad dc=d\,a\wedge b.
\]
The four classes $[a\wedge c],[b\wedge c],
[a\wedge e],[b\wedge e]$ form a basis of
$H^2(P;\R)$, as in \cite[\S3]{Borcea1984}.
An automorphism of \eqref{cdp:eq:C.PK.4} with the stated action extends
to the real nilpotent group: in logarithmic coordinates it
fixes the two central generators and adds central vectors
to the two base generators. Its pullback fixes $a,b$ and
replaces each of $c,e$ by itself plus a real linear
combination of $a,b$. Each of the four displayed classes
therefore changes only by a multiple of the exact form
$a\wedge b$. Since the universal cover of $P$ is
contractible, a map of $P$ and this group-model map with
the same induced fundamental-group map are homotopic.
The original biholomorphism therefore acts as the identity on $H^2$.
In particular every monodromy of our isotrivial family,
whose action on the base $H^1$ was just proved to be the
identity, acts identically on $H^2(P;\C)$.

\item Resolve all exceptional
fibres by successive blowups with smooth centres contained
in those fibres, writing
$\rho:\widetilde Y\to Y$ and $h=f\rho$. For each such
blowup, $\rho_*\OO_{\widetilde Y}=\OO_Y$ and
$R^j\rho_*\OO_{\widetilde Y}=0$ for $j>0$.
One local verification uses the formal neighbourhoods of
the exceptional projective bundle: their successive
conormal layers restrict to nonnegative twists
$\OO_{\PP^{r-1}}(l)$, $l\ge0$, which have no higher
cohomology. The analytic theorem on formal functions
gives the asserted vanishing; normality gives the
degree-zero equality. The Leray edge is consequently the
canonical isomorphism
\begin{equation}
 R^2h_*\OO_{\widetilde Y}\xrightarrow{\simeq}
 R^2f_*\OO_Y=\mathscr G.
\hypertarget{tag-00VC}{}\label{cdp:eq:C.PK.5}
\end{equation}
The same identities hold for a composite by derived
pushforward.

\item Let $D=(h^{-1}(\Sigma))_{\mathrm{red}}$ and use the
relative logarithmic de Rham complex
\[
 \mathscr K^\bullet=
 \Omega_{\widetilde Y}^\bullet(\log D)\big/
 \bigl(h^*\Omega_{\PP^1}^1(\log\Sigma)
       \wedge\Omega_{\widetilde Y}^{\bullet-1}(\log D)\bigr).
\]
Projection onto degree zero is a morphism of complexes
$\mathscr K^\bullet\to\OO_{\widetilde Y}[0]$:
it is the identity in degree zero and zero in every
positive degree. The exact analytic comparison of \ref{cdp:st:D.2.3},
namely \cite[Theorem~2.18 and \S\S2.19--2.21]{Steenbrink1976},
applies to this resolution with all its fibre multiplicities
and without a K\"ahler hypothesis. It gives the logarithmic
de Rham bundle
$\mathscr H^2=\mathbb R^2h_*\mathscr K^\bullet$.
Identity monodromy makes every residue zero, by the
residue interval $[0,1)$ and the exponential comparison.
The global flat frame on $B$ therefore extends across
every point of $\Sigma$, proving
$\mathscr H^2\simeq H^2(P;\C)\otimes\OO_{\PP^1}$.
Taking degree-two direct hypercohomology of the displayed
projection and then \eqref{cdp:eq:C.PK.5} gives a holomorphic map
\begin{equation}
 p:H^2(P;\C)\otimes\OO_{\PP^1}
          \longrightarrow\mathscr G.
\hypertarget{tag-00VD}{}\label{cdp:eq:C.PK.6}
\end{equation}
Over $B$ it is precisely the coefficient projection
$H^2(P_c;\C)\to H^2(P_c,\OO_{P_c})$.

\item In an isotrivial chart this latter map is constant. It is
surjective: the class of
$\overline\alpha\wedge\overline\eta$ used in the proof
of \ref{cdp:st:D.1.12} is represented by a closed invariant two-form
and has nonzero Dolbeault class. The charts also show
that its kernel is a flat subspace. Choose a constant flat
class $v_0\in H^2(P;\C)$ whose projection is nonzero,
and put $s=p(v_0)\in H^0(\PP^1,\mathscr G)$.
This section is nowhere zero on $B$.

\item By \eqref{cdp:eq:C.PK.1} and the nowhere-zero section $\beta$, the
following is an exact sequence of vector bundles on the
entire projective line:
\begin{equation}
 0\longrightarrow\OO_{\PP^1}\beta
  \longrightarrow\mathscr E
  \xrightarrow{\ \beta\wedge-\ }\mathscr G
  \longrightarrow0.
\hypertarget{tag-00VE}{}\label{cdp:eq:C.PK.7}
\end{equation}
Since $H^1(\PP^1,\OO)=0$, choose
$v\in H^0(\PP^1,\mathscr E)=H^1(Y,\OO_Y)$ with
$\beta\wedge v=s$.
In each isotrivial chart,
$\beta$ is the coefficient image of a fixed integral
class, $v_0$ is a fixed de Rham class, and $s=p(v_0)$
is a fixed nonzero class in $H^2(P,\OO_P)$.
The intrinsically oriented quotient $q_c$ in \eqref{cdp:eq:C.PK.2} is also fixed
by the chart: every automorphism preserves its
positive-degree class $\alpha\in\Delta^*$, so it cannot
induce inversion of $\C^*$.
Both linear maps $u\mapsto\beta\wedge u$ and
$u\mapsto\lambda(u)s$ have kernel $A$ and are nonzero.
Thus they differ by a constant $\kappa\in\C^*$ on
that chart. The constants agree on overlaps, because
all the maps and sections used in their definition are
the actual global ones. The connectedness of $B$ gives
one constant $\kappa$, and \eqref{cdp:eq:C.PK.7} yields
\begin{equation}
 \lambda_c(v_c)=\kappa^{-1}\ne0
 \quad\hbox{for every }c\in B.
\hypertarget{tag-00VF}{}\label{cdp:eq:C.PK.8}
\end{equation}

\item Choose $n>0$ with $nd\ge2$ and $z\in\C$ with
$\exp(2\pi\mathrm i z\kappa^{-1})=\alpha^n$.
The single ambient line $M=\exp(2\pi\mathrm i z v)$
has this quotient coordinate on every smooth fibre.
Exactly as in the $k\ne0$ case,
$M|_{P_c}=\pi_c^*L_c$ with $\deg L_c=nd\ge2$.
Its direct image has rank at least two and its relative
linear system gives a surface image over $\PP^1$,
contradicting $a(Y)=1$.

\end{enumerate}
\end{proof}






\begin{proposition}[Exclusion of Inoue general fibres]\hypertarget{tag-01F7}{}\label{cdp:st:C.3.12}
\StatementRecord{cdp:st:C.3.12}{proved}{01F7}{\SeriesRef{cdp:st:C.3.1}, \SeriesRef{cdp:eq:C.IN.2}, \SeriesRef{cdp:eq:C.IN.3}, \SeriesRef{cdp:eq:C.IN.6}; \cite[pp.~274, 278, 280]{Inoue1974}; \cite[Sections 2.1--2.3]{KhaledTeleman2026}; \cite[Theorem 3.2]{BMS2012}; \cite[Section 3.3]{BMS2012}}


Under the hypotheses of \ref{cdp:st:C.3.1}, the very general smooth fibre is
not an Inoue surface. More precisely, a proper holomorphic map from
a compact connected complex manifold to a compact Riemann surface,
whose fibres over the complement of a finite set are Inoue surfaces,
has a class in $H^1(Y;\mathbb R)$ restricting nontrivially to each
such fibre. Thus these families cannot have $b_1(Y)=0$.
This assertion permits arbitrary multiplicities and arbitrary
singularities of the omitted fibres.
\end{proposition}


\begin{proof}

\begin{enumerate}[leftmargin=*]
\item Write $B^\circ$ for the complement of the
finite set and $Y^\circ=f^{-1}(B^\circ)$. Put
$\mathbb H=\{w\in\mathbb C:\operatorname{Im}w>0\}$.
Each Inoue surface is a quotient of $\mathbb H\times\mathbb C$.
In the three standard presentations the first coordinates of all
deck transformations have the form
\begin{equation}
 w\longmapsto a w+b,\qquad a>0,\quad b\in\mathbb R.
\hypertarget{tag-00NK}{}\label{cdp:eq:C.IN.1}
\end{equation}
They include a nonzero real translation and a transformation
$w\mapsto\alpha w$, with $\alpha>1$. Inoue's original
Propositions 2 and 5 establish infinitesimal rigidity for types
$S_M$ and $S^{-}$; his Proposition 4 gives a locally complete
one-parameter family for type $S^{+}$
\cite[pp.~274, 278, 280]{Inoue1974}.
In that family the parameter changes only the second coordinate:
local coordinate transitions have the form
\begin{equation}
\begin{aligned}
 w_i&=\alpha^{m_{ij}}w_j+a_{ij},\\
 z_i&=z_j+b_{ij}w_j+c_{ij}+t\,m_{ij},
\end{aligned}
\hypertarget{tag-00NL}{}\label{cdp:eq:C.IN.2}
\end{equation}
where $m_{ij}\in\mathbb Z$ and the other coefficients are
independent of $t$. In particular the real affine transformations
in the first coordinate are constant throughout this complete
family. The standard deck presentations are also recorded in
\cite[Sections 2.1--2.3]{KhaledTeleman2026}.

\item After shrinking a simply connected base disc $V\subset B^\circ$,
local completeness identifies $f^{-1}(V)$ with the pullback of one
of these families. Its universal cover over $V$ is therefore
$\mathbb H\times\mathbb C\times V$, with holomorphic covering
map to $f^{-1}(V)$. In the $S^{+}$ case the parameter $t$ in
\eqref{cdp:eq:C.IN.2} is a holomorphic function on $V$; the other two cases
are products. The smooth real one-form
\begin{equation}
 \theta_V=d\log\operatorname{Im}w
\hypertarget{tag-00NM}{}\label{cdp:eq:C.IN.3}
\end{equation}
is invariant under every deck transformation, where $d$ is the
exterior derivative on the \emph{whole} covering space. It thus
descends to a closed real one-form on $f^{-1}(V)$.

\item Over a simply connected smaller overlap, a transition
between the two family descriptions lifts holomorphically to the
universal covers. On each fibre its first coordinate is independent
of $z$: composing with a Cayley transform gives a bounded entire
function of $z$, which is constant. Applying the same argument to
the inverse proves that the resulting first coordinate is an
automorphism $h$ of $\mathbb H$. This automorphism conjugates the
first-coordinate deck groups. Their common fixed-point set on
$\partial\mathbb H=\mathbb R\cup\{\infty\}$ is exactly
$\{\infty\}$: the nonzero translation has this as its sole fixed
point, and all their elements fix it. Consequently $h(\infty)=
\infty$ and $h(w)=a w+b$ with $a>0$ and $b\in\mathbb R$.
In a holomorphic family of these transitions, $a$ and $b$ are
holomorphic functions of the base parameter; this follows, for
example, by taking the $w$ derivative and then evaluating at a
fixed $w$. Their real values force them to be locally constant.
Therefore
\begin{equation}
 d\log\operatorname{Im}(a w+b)
 =d\log(a\operatorname{Im}w)
 =d\log\operatorname{Im}w
\hypertarget{tag-00NN}{}\label{cdp:eq:C.IN.4}
\end{equation}
as forms on the total family. Thus the $\theta_V$ glue to an
actual closed form $\theta$ on $Y^\circ$.
The same argument includes a finite \'etale quotient of an Inoue
surface. Locally over a simply connected base disc its finite
cover extends to the family: lift a differentiable marking by
Ehresmann's theorem and equip the resulting covering space with
the pulled-back complex structure. Passing to the Galois closure
when necessary gives a family of finite covers of Inoue surfaces;
these are again Inoue surfaces, as is also seen from their
finite-index lattice presentations. Every deck transformation of
the quotient normalizes a finite-index Inoue deck group on
$\mathbb H\times\mathbb C$. Its first-coordinate group still
contains a nonzero translation and has common boundary fixed
point precisely $\infty$. The preceding conjugation argument
therefore makes each quotient deck transformation real affine in
$w$, and its coefficients are constant over the base disc.
Consequently \eqref{cdp:eq:C.IN.3} descends through the finite quotient as
well. Its pullback to the Inoue cover has the nonzero period
$\log\alpha$, so the descended fibre class is nonzero.
For a fibre loop represented on its universal cover by the deck
transformation $w\mapsto\alpha w$, its integral is
$\log\alpha\ne0$. Hence $[\theta]$ restricts nontrivially to
the fibre.

\item Put
$D=Y\setminus Y^\circ$. Let $\iota:\Delta\to Y$ be a small
holomorphic transverse disc at a smooth point of a reduced
irreducible component of $D$, avoiding all other components.
The pullback of the preceding transverse coordinate atlas on
$\Delta^*$ has constant real affine transitions. On the universal
cover
\[
 \mathbb H_c=\{\operatorname{Im}\tau>c\}
 \longrightarrow\Delta^*,\qquad
 \tau\longmapsto\exp(2\pi\mathrm i\tau),
\]
analytic continuation consequently gives a holomorphic developing
map $W:\mathbb H_c\to\mathbb H$ with
$W(\tau+1)=aW(\tau)+b$ for some $a>0$, $b\in\mathbb R$.
This follows by trivializing the pulled-back constant transition
cocycle on the simply connected covering space.
The period of $\iota^*\theta$ on the positive puncture loop is
$\log a$.

\item Use the Poincar\'e metrics
$|dw|/\operatorname{Im}w$ and
$|d\tau|/(\operatorname{Im}\tau-c)$.
The function $\log\operatorname{Im}w$ has differential of norm
one for the first metric. Schwarz--Pick and the horizontal segment
from $\tau$ to $\tau+1$ therefore give
\begin{equation}
\begin{aligned}
 |\log a|
 &\le d_{\mathbb H}\bigl(W(\tau),W(\tau+1)\bigr)\\
 &\le d_{\mathbb H_c}(\tau,\tau+1)
 \le\frac{1}{\operatorname{Im}\tau-c}.
\end{aligned}
\hypertarget{tag-00NO}{}\label{cdp:eq:C.IN.5}
\end{equation}
Letting $\operatorname{Im}\tau\to+\infty$ proves $a=1$.
Every meridian period of $\theta$ is thus zero. This calculation
applies to the reduced divisor and is unaffected by the
scheme-theoretic multiplicities of its components.

\item The map $H_1(Y^\circ;\mathbb R)\to H_1(Y;\mathbb R)$ is
surjective: loops can be moved away from a complex divisor.
Its kernel is generated by the meridians of the irreducible
components of $D$. Indeed, a bounding real two-chain can be put
in general position, transverse to the smooth divisor locus and
disjoint from its singular locus, which has real codimension at
least four. Removing small discs about its intersection points
leaves a chain in $Y^\circ$ whose additional boundary circles
are meridians. The period homomorphism of $\theta$ therefore factors uniquely
through $H_1(Y;\mathbb R)$ and defines a class
\begin{equation}
 \widetilde\theta\in H^1(Y;\mathbb R),\qquad
 \widetilde\theta|_{Y^\circ}=[\theta].
\hypertarget{tag-00NP}{}\label{cdp:eq:C.IN.6}
\end{equation}
Its restriction to an Inoue fibre is nonzero by \eqref{cdp:eq:C.IN.3}.
Consequently $b_1(Y)>0$.

\item Suppose a very general smooth fibre under the hypotheses of
\ref{cdp:st:C.3.1} is an Inoue surface. Ehresmann's
theorem identifies the diffeomorphism type, and in particular
the fundamental group, of all smooth fibres. The first two
Betti numbers are then $(1,0)$. These surfaces are minimal by
\ref{cdp:st:C.3.1} and non-K\"ahler. The compact-surface classification
gives class VII surfaces or non-K\"ahler elliptic surfaces
\cite[Theorem 3.2]{BMS2012}.
The Inoue group is solvable and is not virtually nilpotent.
For the latter assertion, its normal lattice is $\mathbb Z^3$
in type $S_M$, and a three-dimensional integral Heisenberg group
in types $S^{+}$ and $S^{-}$. Conjugation by a generator of the
cyclic quotient has an eigenvalue $\alpha>1$ on the rational
abelianization of this lattice. Every finite-index subgroup
contains a positive power of that generator and a finite-index
sublattice, so the eigenvalue becomes $\alpha^k>1$. In a
nilpotent group repeated commutators make every such conjugation
unipotent on this rational abelianization. 
\item A Hopf group is virtually cyclic, and a Kodaira group is virtually
nilpotent; neither can equal this group. For a properly elliptic
surface of Euler number zero the canonical bundle formula makes
the base orbifold hyperbolic: its orbifold canonical degree is
positive exactly in Kodaira dimension one. The elliptic
fibration induces a surjection of its fundamental group onto the
orbifold group. That group has a finite-index closed surface
group of genus at least two, and hence is not solvable. This
also cannot be a quotient of an Inoue group. The orbifold exact
sequence is the one used in \cite[Section 3.3]{BMS2012}.
Thus every smooth fibre is an Inoue surface. Applying \eqref{cdp:eq:C.IN.6}
contradicts $b_1(Y)=0$.

\end{enumerate}
\end{proof}






\begin{proposition}[Smooth fibres of an algebraic reduction]\hypertarget{tag-01F4}{}\label{cdp:st:C.3.1}
\StatementRecord{cdp:st:C.3.1}{proved}{01F4}{\SeriesRef{cdp:st:C.3.7}, \SeriesRef{cdp:st:C.3.8}, \SeriesRef{cdp:st:C.3.9}, \SeriesRef{cdp:st:C.3.10}, \SeriesRef{cdp:st:C.3.6}, \SeriesRef{cdp:st:D.2.7}, \SeriesRef{cdp:st:C.3.12}; \cite[Theorem 3.1]{Peters1990}; \cite[p.~2]{Teleman2008}; \cite[Theorem 1]{BBK2019}}

Let $f:Y\to C$ be a holomorphic algebraic reduction of a connected
smooth compact complex threefold with $b_1(Y)=b_2(Y)=0$ and $a(Y)=1$.
The base is $C\simeq\PP^1$. Every smooth fibre $F$ has
\[
 c_1(K_F)=0\text{ in }H^2(F;\Q),\qquad
 c_2(F)=0\text{ in }H^4(F;\Z),\qquad
 \chi(F,\OO_F)=0,
\]
and contains no exceptional curve of the first kind. A very general
smooth fibre has $\kappa(F)\leq0$. The very general fibre's type belongs to the following
list: a complex two-torus, a bielliptic surface, a primary or secondary
Kodaira surface, a primary or secondary Hopf surface, or an Inoue
surface. Every listed type has Euler characteristic zero. This is a
preliminary necessary list. Statement \ref{cdp:st:C.3.7} below excludes the
bielliptic entry, and \ref{cdp:st:C.3.8} makes the very general Hopf fibre have
algebraic dimension zero. \ref{cdp:st:C.3.9} and \ref{cdp:st:C.3.10} exclude both Kodaira
entries. Statement \ref{cdp:st:C.3.6} realizes both primary and secondary Hopf
entries, and \ref{cdp:st:D.2.7} realizes a fixed abelian surface as general
fibre. Statement \ref{cdp:st:C.3.12} excludes the Inoue entry. Thus the
final very general list consists exactly of two-tori and
primary or secondary Hopf surfaces of algebraic dimension zero.
\end{proposition}


\begin{proof}

\begin{enumerate}[leftmargin=*]
\item Connected fibres give $f_*\Q_Y=\Q_C$ and the degree-one edge injection
$H^1(C;\Q)\hookrightarrow H^1(Y;\Q)$. Therefore the genus of $C$ is
zero. The holomorphic tangent sequence on a smooth fibre, with the
last map induced by the differential of $f$, is
\[
0\longrightarrow T_F\longrightarrow T_Y|_F
 \xrightarrow{df} \OO_F\otimes T_{C,f(F)}\longrightarrow0.
\]
The normal factor is a trivial holomorphic line bundle. Thus
$K_F=K_Y|_F$ and $c_2(F)=c_2(Y)|_F$. Rational $H^2(Y)$ and $H^4(Y)$
vanish, and $H^4(F;\Z)=\Z$ has no torsion. These facts give the two
Chern-class assertions. Noether's formula gives $\chi(\OO_F)=0$.
An exceptional curve $E\simeq\PP^1$ with $E^2=-1$ would have
$K_F\cdot E=-1$ by adjunction, contradicting the rational vanishing of
$c_1(K_F)$.

\item Fix $m>0$ and put $G_m=f_*(K_Y^m)$. This is a coherent torsion-free
sheaf on the smooth curve $C$: a section killed by a nonzero base
function vanishes on a dense open subset of $Y$ and hence everywhere.
It is consequently a vector bundle. On the nonempty open set of
cohomological base change its fibre is $H^0(F,K_F^m)$. A rank at least
two would give a meromorphic map
\[
 Y\dashrightarrow \PP_C(G_m^*)
\]
whose restriction to a general fibre is defined by at least two
independent sections of a line bundle and has positive-dimensional
image. Its image has dimension at least two. The target is projective:
a vector bundle on a compact Riemann surface is algebraic, and its
projectivization is a projective variety. The image is therefore
projective and its rational functions pull back to a subfield of
$\mathcal M(Y)$ of transcendence degree at least two. This contradicts
$a(Y)=1$. Thus $\rk G_m\leq1$ for every $m$. Removing the countable
union of the exceptional base-change sets gives the claimed very
general-fibre assertion.

\item The Enriques--Kodaira classification is given in
\cite[Theorem 3.1]{Peters1990}. The surface is minimal by the
adjunction argument. Rational and ruled surfaces have a rational
curve on which the canonical class has nonzero degree, so they are
excluded. K3 and Enriques surfaces have Euler numbers $24$ and $12$,
respectively. Properly elliptic and general-type surfaces have
positive Kodaira dimension. The remaining K\"ahler possibilities are
tori and bielliptic surfaces; the remaining Kodaira-dimension-zero
non-K\"ahler possibilities are the two types of Kodaira surfaces.
For class VII, $b_1(F)=1$, and $e(F)=0$ gives $b_2(F)=0$. The
classification of this case gives Hopf and Inoue surfaces; its
precise scope, including finite covers, is stated in
\cite[p.~2]{Teleman2008} and \cite[Theorem 1]{BBK2019}.

\end{enumerate}
\end{proof}






\begin{proposition}[Full-fibre cohomology and direct images]\hypertarget{tag-01F6}{}\label{cdp:st:C.3.11}
\StatementRecord{cdp:st:C.3.11}{proved}{01F6}{\SeriesRef{cdp:st:C.3.1}, \SeriesRef{cdp:st:C.3.12}, \SeriesRef{cdp:st:C.3.6}, \SeriesRef{cdp:st:D.2.7}, \SeriesRef{cdp:st:D.1.9}, \SeriesRef{cdp:st:C.3.7}, \SeriesRef{cdp:st:C.3.9}, \SeriesRef{cdp:st:C.3.10}, \SeriesRef{cdp:st:C.3.8}, \SeriesRef{cdp:st:D.1.12}}


Under the hypotheses of \ref{cdp:st:C.3.1}, a very general smooth fibre
is a two-torus or a primary or secondary Hopf surface of
algebraic dimension zero. The Inoue exclusion is \ref{cdp:st:C.3.12}.
Both Hopf
entries and the torus entry occur, by \ref{cdp:st:C.3.6} and \ref{cdp:st:D.2.7}.
For the entire smooth locus there are exactly the following
two cohomological alternatives:
\[
\begin{array}{c|c|c}
\text{smooth fibres}&
 (h^0(\OO),h^1(\OO),h^2(\OO))&
 \text{coherent direct images}\\ \hline
\text{two-tori}&(1,2,1)&
 R^2f_*\OO_Y\simeq\det R^1f_*\OO_Y\\
b_1=1,\ b_2=0&(1,1,0)&
 R^1f_*\OO_Y\simeq\OO_{\PP^1},\
 R^2f_*\OO_Y=0 .
\end{array}
\tag{C.PK.9}\hypertarget{tag-00VG}{}\label{cdp:eq:C.PK.9}
\]
The indicated cohomology dimensions hold on every
scheme-theoretic full fibre in the respective branch,
and both direct images satisfy base change there.
In the second branch
$H^q(Y,\OO_Y)=(\C,\C,0,0)$ and
$\Pic^0(Y)=\C\beta$ with the exponential convention
of \ref{cdp:st:D.1.9}.
\end{proposition}


\begin{proof}

\begin{enumerate}[leftmargin=*]
\item \ref{cdp:st:C.3.7}, \ref{cdp:st:C.3.9}, \ref{cdp:st:C.3.10} and \ref{cdp:st:C.3.12} remove the bielliptic,
Kodaira and Inoue types from \ref{cdp:st:C.3.1}; \ref{cdp:st:C.3.8} gives the Hopf algebraic
dimension. All smooth fibres are diffeomorphic over the
connected smooth base. In the torus branch they have
$b_1=4$, $b_2=6$, and are K\"ahler by the compact-surface
criterion. Their rational first Chern class and second
Chern number vanish by \ref{cdp:st:C.3.1}. The surface classification
then makes them two-tori: an effective nonzero canonical
divisor has strictly positive degree against a K\"ahler
form, whereas this rational first Chern class has zero
degree. The minimal K\"ahler surfaces of Kodaira dimension
zero with these Betti numbers are exactly the two-tori.
\ref{cdp:st:D.1.12} supplies this branch's full-fibre assertions.

\item In the Hopf branch every smooth fibre has
$b_1=1$ and $b_2=0$. A nonzero holomorphic two-form
is closed and has strictly positive integral against
its conjugate, so it defines a nonzero de Rham
cohomology class. Thus $b_2=0$ forces $p_g=0$.
Since $\chi(\OO)=0$ by \ref{cdp:st:C.3.1}, one has $q=1$.
By \cite[Theorem~2.2, pp.~4--5]{Viaclovsky2026Combs}, $\mathscr E=R^1f_*\OO_Y$ is a line bundle,
and its nowhere-zero transgression section $\beta$
from \ref{cdp:st:D.1.9} trivializes it. The coherent sheaf
$\mathscr G=R^2f_*\OO_Y$ has rank zero, hence finite
length. Hirzebruch--Riemann--Roch on $Y$ and coherent
Leray give the exact equality
\[
 0=\chi(Y,\OO_Y)
   =\chi(\OO_{\PP^1})-\chi(\mathscr E)+\chi(\mathscr G)
   =1-1+\operatorname{length}(\mathscr G).
\]
Therefore $\mathscr G=0$. Top-degree base change and
the parameter-multiplication exact sequence give
$h^2(\OO_{Y_c})=0$ and $h^1(\OO_{Y_c})=1$ on each
full fibre; \cite[Theorem~2.2, pp.~4--5]{Viaclovsky2026Combs} gives $h^0(\OO_{Y_c})=1$.
Leray now gives the stated global structure-sheaf
cohomology. Finally $H^1(Y;\Z)=0$ and the exponential
sequence identify $\Pic^0(Y)$ with
$H^1(Y,\OO_Y)=\C\beta$.

\end{enumerate}
\end{proof}






\begin{proposition}[Euler numbers of fibre supports]\hypertarget{tag-01FB}{}\label{cdp:st:C.3.5}
\StatementRecord{cdp:st:C.3.5}{proved}{01FB}{\SeriesRef{cdp:st:C.3.1}}

For every threefold and fibration in \ref{cdp:st:C.3.1}, including every type on
its necessary list, the Euler-number identity is
\[
\int_Y c_3(Y)=e(Y)=\sum_{c\in\Delta}e\bigl((Y_c)_{\rm red}\bigr),
\]
where $\Delta$ contains the finitely many critical values.
\end{proposition}


\begin{proof}
All the listed smooth surfaces have Euler number zero by \ref{cdp:st:C.3.1}.
Additivity of compactly supported Euler characteristic for the
stratification by the smooth-fibre locus and the finitely many
exceptional fibres gives
$e(Y)=e_c(C\setminus\Delta)e(F)+\sum_{c\in\Delta}e(Y_c)$.
The first term vanishes. A scheme-theoretic fibre and its reduction
have the same underlying topological space, so their topological
Euler numbers agree. Chern--Gauss--Bonnet gives the first equality.

\end{proof}






\begin{proposition}[Bundles with primary and secondary Hopf fibres]\hypertarget{tag-01FC}{}\label{cdp:st:C.3.6}
\StatementRecord{cdp:st:C.3.6}{proved}{01FC}{\SeriesRef{cdp:eq:C.H.1}}


There are compact complex threefolds with $b_1=b_2=0$ and $a=1$
whose holomorphic algebraic reduction is a smooth bundle over
$\PP^1$ with primary Hopf fibre. There are also such threefolds
with secondary Hopf fibre. In these examples every fibre has Euler
number zero and $c_3=0$.
\end{proposition}


\begin{proof}

\begin{enumerate}[label=\arabic*.,leftmargin=*]
\item Put $a_1=1$, $a_2=\sqrt2$ and
$W=(\C^2\setminus\{0\})_z\times(\C^2\setminus\{0\})_w$.
The additive complex Lie group $\C$ acts holomorphically by
\[
t\cdot(z,w_1,w_2)
 =\bigl(e^t z,e^{ia_1t}w_1,e^{ia_2t}w_2\bigr).
\tag{C.H.1}\hypertarget{tag-00L8}{}\label{cdp:eq:C.H.1}
\]
This action is free. Indeed $e^t z=z$ implies $t=2\pi in$ with
$n\in\Z$; a nonzero $w_j$ then gives $e^{-2\pi a_jn}=1$ and
therefore $n=0$. It is proper as well. If a point and its translate
range in fixed compact subsets of $W$, the $z$-norm bounds
$\operatorname{Re}t$. Writing $s=\operatorname{Im}t$, the identity
\[
\|t\cdot w\|^2
 =e^{-2a_1s}|w_1|^2+e^{-2a_2s}|w_2|^2
\]
and the positive constants $a_1,a_2$ bound $s$ both above and
below, since both $w$-norms are bounded away from zero. The inverse
image of a compact subset under the action graph is consequently
closed and bounded in the group variable and compact in the other
variable.

\item Every orbit meets the real submanifold
$\{\|z\|=\|w\|=1\}=S^3\times S^3$ exactly once.
The first equation fixes $\operatorname{Re}t=-\log\|z\|$.
The displayed squared $w$-norm is strictly decreasing in $s$,
with limits $+\infty$ and $0$, so the second equation has one
solution. Its derivative is strictly negative, and the implicit
function theorem makes the resulting cross-section smooth.
Thus the complex quotient $Y=W/\C$ is a smooth compact complex
threefold diffeomorphic to $S^3\times S^3$. In particular
$b_1(Y)=b_2(Y)=0$ and $e(Y)=0$.

\item The map
\[
f:Y\longrightarrow\PP^1,\qquad [z,w]\longmapsto[z_1:z_2]
\]
is well-defined and holomorphic. Choose a holomorphic nonzero
frame of the tautological line on a small base disc. Scaling the
$z$-coordinate to this frame leaves exactly the transformations
$t=2\pi in$, which act on $w$ by powers of
\[
Q=\operatorname{diag}(q_1,q_2),\qquad
q_j=e^{-2\pi a_j}.
\]
This identifies $f$ over that disc with the product having fibre
$H=(\C^2\setminus\{0\})/\langle Q\rangle$.
The quotient is a primary Hopf surface, and these local products
prove that $f$ is a proper holomorphic submersion with connected
fibres.

\item A meromorphic function on $H$ lifts
to a $Q$-invariant meromorphic function on $\C^2\setminus\{0\}$.
It extends meromorphically across the origin: the closure of its
polar divisor is analytic by the codimension-two extension theorem
for divisors, is locally principal on the smooth ball, and
multiplying by a local defining function reduces extension to
Hartogs' theorem for a holomorphic function. Write the germ as
$g/h$ with relatively prime convergent power series. The identity
$(g\circ Q)h=g(h\circ Q)$ gives
\[
g\circ Q=ug,\qquad h\circ Q=uh
\]
for one holomorphic unit $u$: divisibility follows from relative
primality, and preservation of the vanishing order makes the
quotient a unit. Write $u/u(0)=e^\ell$ with $\ell(0)=0$ and set
\[
v(w)=\exp\left(\sum_{n\ge0}\ell(Q^nw)\right).
\]
The sum converges normally on a sufficiently small ball, because
$Q$ is a strict contraction. It gives
$v(Qw)/v(w)=u(0)/u(w)$. Consequently both $vg$ and $vh$ are
eigenfunctions of pullback by $Q$ with the same eigenvalue $u(0)$.
The monomial $w_1^rw_2^s$ has eigenvalue
$e^{-2\pi(r+s\sqrt2)}$. These numbers are distinct for distinct
$(r,s)\in\Z_{\ge0}^2$. Comparing power-series coefficients shows
that each nonzero eigenfunction is a single monomial and that the
two eigenfunctions are proportional. The meromorphic function is
therefore constant, first near the origin and then everywhere by
analytic continuation. Hence $a(H)=0$.

\item Every meromorphic function on $Y$ is constant on each general
fibre of $f$. On a product chart, restricting to a local section
gives the descended meromorphic function on the base; equality on
the fibres and the identity theorem identify its pullback with
the original function. These local functions agree on overlaps
and extend at the finitely many possible poles on the base.
Thus $\mathcal M(Y)=f^*\mathcal M(\PP^1)$, proving $a(Y)=1$
and proving that $f$ is the holomorphic algebraic reduction.

\item For any integer $m>1$, let $\mu_m$ act on $W$ by
$(z,w)\mapsto(z,\zeta w)$. It commutes with \eqref{cdp:eq:C.H.1} and induces
a free action on $Y$. A fixed orbit would give $t=2\pi in$ from
the $z$-coordinate; equality of norms in a nonzero $w_j$ then
gives $n=0$ and $\zeta=1$. The quotient
$Y_m=Y/\mu_m$ is a smooth compact complex threefold. Transfer
injects $H^j(Y_m;\Q)$ into $H^j(Y;\Q)$, so $b_1(Y_m)=b_2(Y_m)=0$.
The simply connected cover $Y\simeq S^3\times S^3$ gives
$\pi_1(Y_m)=\mu_m$, and the universal coefficient theorem gives
$H^2(Y_m;\Z)\simeq\Z/m\Z$: its free part is zero by the
rational calculation, and its torsion is
$\operatorname{Ext}(H_1(Y_m;\Z),\Z)$. The primary
example has $H^1(Y;\Z)=H^2(Y;\Z)=0$.
The induced map to $\PP^1$ has fibre
\[
H_m=(\C^2\setminus\{0\})/
       \bigl(\langle Q\rangle\times\mu_m\bigr).
\]
Its universal cover is $\C^2\setminus\{0\}$ and its fundamental
group is $\Z\times\mu_m$, which is noncyclic. It is therefore a
secondary Hopf surface, with the displayed primary cover of degree
$m$. Any meromorphic function on $H_m$ pulls back to one on $H$,
so $a(H_m)=0$. The same local descent proves
$\mathcal M(Y_m)=f_m^*\mathcal M(\PP^1)$ and $a(Y_m)=1$.
Transfer of Euler characteristic under the finite covering,
or the smooth-fibration Euler formula, gives $e(Y_m)=0$.

\end{enumerate}
\end{proof}






\section{Cotangent torsion, conductor and the printed endpoint error}\hypertarget{tag-0340}{}\label{sec:II:04}


Throughout this section $i:S\hookrightarrow Y$ is the inclusion of a reduced effective Cartier divisor into a complex manifold of dimension $n+1$, $j:S_{\rm reg}\hookrightarrow S$ is the inclusion of its regular locus, and $A$ is an invertible analytic $\OO_S$-module. We write
\[
 N^*=\OO_Y(-S)|_S,\qquad V=\Omega_Y^1|_S,\qquad
 T=\ker(\Omega_S^1\longrightarrow j_*\Omega^1_{S_{\rm reg}}),\qquad Q=\Omega_S^1/T.
\]
All maps between sheaves below are $\OO_S$-linear maps of coherent analytic sheaves, except for the displayed restriction into $j_*\Omega^1_{S_{\rm reg}}$. The latter is also $\OO_S$-linear; its kernel and image have the coherent descriptions proved below.

\begin{proposition}[The conormal filtration and its exact section criterion]\hypertarget{tag-01AD}{}\label{cdp:st:A.1.1}
\StatementRecord{cdp:st:A.1.1}{proved}{01AD}{the definitions and constructions stated in the proof}
\label{wbII:prop:criterion}\label{wbII:prop:filtration}\label{wbII:lem:conormal}\label{wbII:lem:conormal-injective}\label{wbII:prop:missing-global-section-criterion}\label{wbII:prop:missing-cotangent-filtration}
The conormal map $d:N^*\to V$, defined by $[F]\mapsto dF|_S$ for a local defining equation $F$, is injective. With $K=\ker(V\to Q)$, there are canonical exact sequences
\begin{align}
0&\longrightarrow N^*\xrightarrow{d}K\longrightarrow T\longrightarrow0,
\hypertarget{tag-003W}{}\label{cdp:A:filtration1}\\
0&\longrightarrow K\longrightarrow V\longrightarrow Q\longrightarrow0.
\hypertarget{tag-003X}{}\label{cdp:A:filtration2}
\end{align}
Let $\partial_A:H^0(S,Q\otimes A)\to H^1(S,K\otimes A)$ and $\delta_A:H^0(S,T\otimes A)\to H^1(S,N^*\otimes A)$ be their connecting maps. Then
\[
 H^0(S,V\otimes A)=0
 \quad\Longleftrightarrow\quad
 H^0(S,K\otimes A)=0\ \hbox{and}\ \partial_A\hbox{ is injective}.
\tag{A.1}\hypertarget{tag-003Y}{}\label{cdp:eq:A.1}
\]
Under the explicit hypotheses $H^0(S,N^*\otimes A)=0$ and $H^0(S,Q\otimes A)=0$, $H^0(S,V\otimes A)$ is canonically $\ker\delta_A$.
\end{proposition}


\begin{proof}
\begin{enumerate}[label={$\langle1\rangle$\arabic*.},leftmargin=*]
\item 
\noindent\textbf{Step 1.} Locally $R=\C\{x_0,\ldots,x_n\}/(F)$ is reduced. Its total quotient ring is the product of the fraction fields of its irreducible components. At the generic point of every component, $S$ is smooth and $dF$ is nonzero. Consequently $a\,dF=0$ in $R^{n+1}$ implies that $a$ vanishes in each of those fields. The map from the reduced ring $R$ into their product is injective, so $a=0$. This proves injectivity in the usual right-exact conormal sequence. Taking the inverse image of $T\subset\Omega_S^1$ gives \eqref{cdp:A:filtration1} and \eqref{cdp:A:filtration2}. Tensoring by $A$ is exact because $A$ is locally free. The cohomology sequence of \eqref{cdp:A:filtration2} identifies the image of $H^0(VA)$ in $H^0(QA)$ with $\ker\partial_A$ and its kernel with $H^0(KA)$, proving \eqref{cdp:eq:A.1}. When $H^0(QA)=0$, it identifies $H^0(VA)$ with $H^0(KA)$. The first sequence then identifies this group with $\ker\delta_A$ when $H^0(N^*A)=0$.

\end{enumerate}
\end{proof}





\begin{proposition}[The Jacobian fractional ideal for every reduced hypersurface]\hypertarget{tag-01AF}{}\label{cdp:st:A.2.1}
\StatementRecord{cdp:st:A.2.1}{proved}{01AF}{\SeriesRef{cdp:A:filtration1}}

Let $\mathcal J\subset\OO_S$ be the Jacobian ideal, generated locally by $F_{x_0},\ldots,F_{x_n}$ modulo $F$. In the total quotient sheaf $\mathcal M_S$ define
\[
 \mathcal J^{-1}=\{q\in\mathcal M_S:q\mathcal J\subset\OO_S\}.
\]
There are canonical identifications
\[
 K\simeq N^*\otimes\mathcal J^{-1},\qquad
 T\simeq N^*\otimes(\mathcal J^{-1}/\OO_S),
\tag{A.2}\hypertarget{tag-004B}{}\label{cdp:eq:A.2}
\]
under which the first map in \eqref{cdp:A:filtration1} is induced by $\OO_S\hookrightarrow\mathcal J^{-1}$. Thus (A.1), with this explicitly defined fractional ideal, is an exact necessary and sufficient criterion for every reduced hypersurface.
\end{proposition}


\begin{proof}

\begin{enumerate}[label={$\langle1\rangle$\arabic*.},leftmargin=*]
\item The ideal $J$ contains a nonzerodivisor: it is nonzero at every minimal prime, and prime avoidance supplies an element outside all minimal primes. Every $R$-linear map $\varphi:J\to R$ is multiplication by a unique element of the total quotient ring. Indeed, for a nonzerodivisor $a\in J$, the identity $a\varphi(b)=b\varphi(a)$ for $b\in J$ gives the multiplier $q=\varphi(a)/a$. Hence $J^{-1}=\Hom_R(J,R)$ is coherent.

\item At the generic points, an element $\alpha\in K$ is uniquely $q\,dF$ for a tuple of meromorphic functions $q$. The coefficients of $\alpha$ belong to $R$ exactly when $qF_{x_i}\in R$ for every $i$, which is the condition $q\in J^{-1}$. Conversely such a vector restricts to zero in the cotangent space of every regular component, hence belongs to $K$. The map $[F]\otimes q\mapsto q\,dF|_S$ proves the first identification. Its quotient by $R\,dF$ gives the second. Replacing $F$ by $uF$, with $u$ a holomorphic unit, replaces $dF|_S$ by $u\,dF|_S$ and replaces the conormal generator by $u[F]$. Coordinate changes transform the derivative vector by an invertible Jacobian matrix. These are exactly the transitions in \eqref{cdp:eq:A.2}, so the identifications are canonical.

\end{enumerate}
\end{proof}






\begin{proposition}[Normalization, conductor, and the excess kernel]\hypertarget{tag-01AG}{}\label{cdp:st:A.2.2}
\StatementRecord{cdp:st:A.2.2}{proved}{01AG}{the definitions and constructions stated in the proof}
\label{wbII:lem:conductor}\label{wbII:lem:conductor-square}
Let $\eta:\widetilde S\to S$ be the finite analytic normalization, $\mathcal B=\eta_*\OO_{\widetilde S}$, and $\mathfrak c=\Ann_{\OO_S}(\mathcal B/\OO_S)$ its conductor. There are canonical inclusions
\[
 \mathcal J\subset\mathfrak c,\qquad
 \mathcal B\subset\mathcal J^{-1},\qquad
 \eta_*\eta^*N^*\hookrightarrow K.
\tag{A.3}\hypertarget{tag-004M}{}\label{cdp:eq:A.3}
\]
In a local ring $R$ with normalization $B$, put $D=\Hom_B(JB,B)$, viewed as a fractional ideal. Then
\[
 J^{-1}=\ker\left[D\longrightarrow
 \Hom_R\bigl(J,(B/\mathfrak c)/(R/\mathfrak c)\bigr)\right],
\quad q\longmapsto\bigl(b\mapsto qb\bmod R\bigr).
\tag{A.4}\hypertarget{tag-004O}{}\label{cdp:eq:A.4}
\]
For any locally free $\OO_S$-module $E$, the conductor gives a second exact section criterion:
\begin{multline}
H^0(S,E)=
\ker\left[H^0(\widetilde S,\eta^*E)\oplus H^0(D_c,E|_{D_c})\right.\\
\left.\longrightarrow H^0(\widetilde D_c,\eta^*E|_{\widetilde D_c})\right],
\qquad (a,b)\longmapsto a|_{\widetilde D_c}-\eta^*b,
\tag{A.5}\hypertarget{tag-004P}{}\label{cdp:eq:A.5}
\end{multline}
where $D_c$ and $\widetilde D_c$ are defined by $\mathfrak c$ in $S$ and $\widetilde S$. In particular, the ambient group vanishes exactly when the map in \eqref{cdp:eq:A.5}, for $E=V\otimes A$, is injective.
\end{proposition}


\begin{proof}

\begin{enumerate}[label={$\langle1\rangle$\arabic*.},leftmargin=*]
\item Trace duality gives $JB\subset R$. Choose a generic linear projection for which Weierstrass preparation writes the equation as a unit times a monic polynomial $P(z)$ over $A_0=\C\{x_1,\ldots,x_n\}$. Thus $R=A_0[z]/(P)$ is finite free over $A_0$ and generically separable. The functional taking the coefficient of $z^{m-1}$ in a remainder modulo $P$, where $m=\deg P$, generates $\Hom_{A_0}(R,A_0)$ as an $R$-module: the matrix of the pairing on $1,z,\ldots,z^{m-1}$ has zeros below the reverse diagonal and ones on that diagonal. Over the fraction field, Lagrange interpolation identifies that functional with $r\mapsto\operatorname{Tr}(r/P'(z))$. Therefore the trace dual is $P'(z)^{-1}R$.

\item For $b\in B$ and $r\in R$, the element $br$ is integral over the integrally closed ring $A_0$. Its trace belongs to $A_0$. It follows that $b$ lies in this trace dual, so $P'(z)b\in R$. Modulo $F$, the derivative in the chosen projection direction differs from $P'(z)$ by a unit. The permissible projection directions form a nonempty open set; choosing a linear basis of such directions proves $F_{x_i}B\subset R$ for every $i$. Thus $J\subset\mathfrak c$ and $B\subset J^{-1}$. The last inclusion of \eqref{cdp:eq:A.3} is (A.2) and the projection formula.

\item The module $D$ consists exactly of the meromorphic $q$ satisfying $qJ\subset B$. Requiring that these products descend to $R$ is precisely the vanishing of their classes in $B/R=(B/\mathfrak c)/(R/\mathfrak c)$, proving \eqref{cdp:eq:A.4}. Finally the ring map
\[
 R\longrightarrow B\oplus R/\mathfrak c,
 \qquad r\longmapsto(r,r\bmod\mathfrak c)
\]
has image the kernel of $(b,\bar r)\mapsto b-\bar r\in B/\mathfrak c$. Surjectivity onto $B/\mathfrak c$ comes from the first summand. Tensoring this exact sequence by the free module corresponding to $E$ and taking global sections gives \eqref{cdp:eq:A.5}. All conductor restrictions are scheme-theoretic, so this argument includes nonreduced conductors.

\end{enumerate}
\end{proof}



The fractional ideal appearing in (A.2) has an established interpretation. Granger--Schulze, Proposition~3.4 and Corollary~3.5 of \cite{GS2014}, identify $J^{-1}$ with the module of logarithmic residues and prove the inclusions in \eqref{cdp:eq:A.3}. The preceding proof gives the maps needed here directly. Their Theorems~1.1 and~1.3 also give the following exact extension of the normal-crossings case.




\begin{proposition}[The canonical normal-crossings formula]\hypertarget{tag-01AE}{}\label{cdp:st:A.1.2}
\StatementRecord{cdp:st:A.1.2}{proved}{01AE}{\SeriesRef{cdp:st:A.2.1}}

For a reduced normal-crossings divisor, the map in (A.3) is an isomorphism. In a chart with $F=x_1\cdots x_r$, it is explicitly
\[
 \bigoplus_{i=1}^r R/(x_i)\longrightarrow K,
 \qquad (a_i)_i\longmapsto
 \sum_{i=1}^r a_i\prod_{j\ne i}x_j\,dx_i.
\tag{A.6}\hypertarget{tag-004Q}{}\label{cdp:eq:A.6}
\]
The tensor factor $N^*$ is trivialized here by $[F]$. At a double crossing $T\simeq R/(x,y)$; at a triple crossing,
\[
 T\simeq R^3/\bigl(R(1,1,1)+R(x,0,0)+R(0,y,0)+R(0,0,z)\bigr).
\tag{A.7}\hypertarget{tag-004R}{}\label{cdp:eq:A.7}
\]
\end{proposition}


\begin{proof}

\begin{enumerate}[label={$\langle1\rangle$\arabic*.},leftmargin=*]
\item On the normalization branch $x_i=0$, an ambient form maps to zero in the cotangent bundle of that branch precisely when all coefficients except that of $dx_i$ vanish. Thus a coefficient of $dx_i$ must vanish on every branch $x_j=0$ for $j\ne i$. In the convergent power-series ring it is divisible by $\prod_{j\ne i}x_j$. Its remaining factor is defined modulo $x_i$, which gives exactly \eqref{cdp:eq:A.6}. The coefficients of different $dx_i$ show injectivity. Coefficients of the coordinates not among $x_1,\ldots,x_r$ vanish on every branch and hence vanish in $R$. This proves surjectivity in every dimension.

\item For $r=1$ the map is $a\mapsto a\,dx$. For $r=2$ it is $(a,b)\mapsto ay\,dx+bx\,dy$, and quotienting by $dF=y\,dx+x\,dy$ leaves the generator $y\,dx=-x\,dy$, annihilated by $(x,y)$. For $r=3$ the branch generators are $yz\,dx,xz\,dy,xy\,dz$. Their sum is $dF$ and their respective annihilators are $(x),(y),(z)$, proving \eqref{cdp:eq:A.7}. Multiplication of $F$ by a unit and changes of coordinates were checked in \ref{cdp:st:A.2.1}. Formula \eqref{cdp:eq:A.6} is the restriction of the canonical fractional-ideal map, so it glues without a chosen branch ordering, equation, or coordinate system.

\end{enumerate}
\end{proof}






\begin{proposition}[Pinch points and cuspidal edges]\hypertarget{tag-01AI}{}\label{cdp:st:A.2.4}
\StatementRecord{cdp:st:A.2.4}{proved}{01AI}{\SeriesRef{cdp:st:A.2.1}, \SeriesRef{cdp:st:A.1.2}; \cite[Proposition 3.4, pp.~9--10, and Corollary 3.5, p.~10, author version]{GS2014}; \cite[Theorem 1, p.~1]{Simonic2016Hartogs}}

At the pinch point $F=x^2-y^2z$, with normalization $x=uv,y=u,z=v^2$, write $B=\C\{u,v\}$. Then
\[
 K=B\,dF,\qquad T\simeq B/R\simeq R/(x,y),
\quad \tau=2yz\,dx-2xz\,dy-xy\,dz
\tag{A.8}\hypertarget{tag-004S}{}\label{cdp:eq:A.8}
\]
is a torsion generator. At the cuspidal edge $F=x^2-y^3$, with normalization $x=t^3,y=t^2$ and unchanged coordinate $z$, write $B=\C\{t,z\}$. Then
\[
 K=t^{-1}B\,dF,\qquad
 T\simeq t^{-1}B/R\simeq R/(x,y^2),
 \quad\tau=2y\,dx-3x\,dy.
\tag{A.9}\hypertarget{tag-004T}{}\label{cdp:eq:A.9}
\]
The conductor is $(x,y)$ for the pinch point and $(x,y)$ for the cusp; in the cusp normalization it is $t^2B$. An ordinary surface, meaning a surface whose singularities are double normal crossings, triple normal crossings, and pinch points, satisfies canonically $K=\eta_*\eta^*N^*$ and $T=N^*\otimes(\eta_*\OO_{\widetilde S}/\OO_S)$. For the cuspidal edge the additional quotient $J^{-1}/B$ is generated by $t^{-1}$ over $\C\{z\}$.
\end{proposition}


\begin{proof}

\begin{enumerate}[label={$\langle1\rangle$\arabic*.},leftmargin=*]
\item For the pinch point,
\[
 R=\C\{u,v^2,uv\}=\C\{v^2\}+uB,
 \quad J=(uv,uv^2,u^2)=u(u,v)B.
\]
The equality of the two descriptions of $J$ follows by separating even and odd powers of $v$. An element $q\in J^{-1}$ satisfies $(qu)(u,v)\subset B$. Since $B$ is a regular local ring, $\Hom_B((u,v),B)=B$: a meromorphic function whose products with $u$ and $v$ are holomorphic has no pole along any divisor, and its possible isolated singularity is removable by the holomorphic extension theorem. Hence $qu=h\in B$. The requirements $hv,hv^2\in R$, restricted to $u=0$, say respectively that $h(0,v)v$ and $h(0,v)v^2$ are even power series. The first requires $h(0,v)$ to be odd and the second to be even. Therefore $h(0,v)=0$, so $h\in uB$ and $q\in B$. The reverse inclusion follows either directly or from (A.3). Also $B/R$ is generated by $v$ and its annihilator is $(u,uv)=(y,x)$. Multiplying $dF$ by $v=x/y$ gives the displayed $\tau$.

\item For the cusp, $R=\C\{t^2,t^3,z\}$ and $J=(t^3,t^4)$. A meromorphic $q$ with $qt^3,qt^4\in R$ has no poles away from $t=0$; writing its Laurent expansion in $t$ shows that the lowest allowed exponent is $-1$. Indeed an exponent $-3$ produces a forbidden $t$ term after multiplication by $t^4$, an exponent $-2$ produces one after multiplication by $t^3$, and smaller exponents produce a pole. Conversely every $t^{-1}B$ has both products in $R$. Thus $J^{-1}=t^{-1}B$. Modulo $R$, the quotient has the $\C\{z\}$-basis $t^{-1},t$; multiplication by $y$ sends the first basis vector to the second, while $x$ and $y^2$ kill both. This proves the cyclic presentation and \eqref{cdp:eq:A.9}. The conductor calculations follow directly from the displayed descriptions of $R\subset B$. Formula (A.2) and the three local computations prove the ordinary-surface assertion by gluing.

\end{enumerate}
\end{proof}






\begin{proposition}[Vertical residue poles at a tacnode--cusp intersection]\hypertarget{tag-01AJ}{}\label{cdp:st:A.2.5}
\StatementRecord{cdp:st:A.2.5}{proved}{01AJ}{the definitions and constructions stated in the proof}

For $F=x^2-y^4z^3$, normalization is $x=u^2v^3$, $y=u$, $z=v^2$. Put $A_0=\C\{u,v^2\}$ and $B=\C\{u,v\}$. Then
\[
 \begin{aligned}
 R&=A_0\oplus u^2v^3A_0,&\mathfrak c&=u^2v^2B,\\
 J^{-1}&=A_0\oplus(uv)^{-1}A_0,&
 J^{-1}/B&\simeq A_0/(uv^2).
 \end{aligned}
 \tag{A.9a}\hypertarget{tag-004U}{}\label{cdp:eq:A.9a}
\]
In a trivialization of $N^*$ the torsion module is
\[
 T\simeq R/(x,y^3z^2),\qquad
 \tau=2yz\,dx-4xz\,dy-3xy\,dz.
 \tag{A.9b}\hypertarget{tag-004V}{}\label{cdp:eq:A.9b}
\]
In particular $v^{-1}\in J^{-1}\setminus B$. For the ruling with base coordinate $v$, this is a residue constant on its generic fibres with a pole on the vertical divisor $v=0$, although the horizontal conductor $u=0$ has coefficient two.
\end{proposition}


\begin{proof}

\begin{enumerate}[label={$\langle1\rangle$\arabic*.},leftmargin=*]
\item The element $v=x/(y^2z)$ belongs to the fraction field and satisfies $v^2=z$. Adjoining it gives the regular ring $B$, which is therefore the finite normalization. Separating even and odd powers of $v$ gives the displayed description of $R$. An even monomial $u^av^b$ belongs to the conductor exactly when $a\ge2$ and $b\ge2$; an odd monomial already in $R$ has $a\ge2$, $b\ge3$ and is then in the conductor. Thus $\mathfrak c=u^2v^2B$.

\item The normalized Jacobian generators, with irrelevant nonzero constants removed, are $u^2v^3,u^3v^6,u^4v^4$. If $qJ\subset R$, the first product shows that $q$ has no other polar divisor, with possible orders at most two along $u=0$ and three along $v=0$. A term of order $v^{-3}$ would produce a forbidden $v^1$ term in $qu^4v^4$; a term of order $v^{-2}$ would produce one in $qu^2v^3$. Both coefficients must vanish, so $q$ has order at least $-1$ in $v$. A coefficient $a(v)u^{-2}$ would require both $a(v)v^3$ and $a(v)v^6$ to be even Laurent series, from the $u^0$ and $u^1$ coefficients of the first two products. The first requirement makes $a$ odd and the second makes it even, so $a=0$. Thus $q=t/(uv)$ with $t\in B$.

\item Write $t=a(u,v^2)+v b(u,v^2)$. The first Jacobian product equals $uv^2t$ and belongs to $R$ exactly when $b$ is divisible by $u$. The other two products automatically belong to $R$ under this condition. This gives $J^{-1}=A_0\oplus(uv)^{-1}A_0$. The odd part of $B$ is $(uv)^{-1}(uv^2)A_0$, giving $J^{-1}/B$ in \eqref{cdp:eq:A.9a}. The odd part of $R$ is $(uv)^{-1}u^3v^4A_0$, so $J^{-1}/R$ is generated by $(uv)^{-1}$ with annihilator $(x,y^3z^2)$. Multiplication of $dF$ by $(uv)^{-1}$ gives \eqref{cdp:eq:A.9b}.

\item The particular vertical pole has the completely regular gradient products
\[
 v^{-1}dF=2y^2z\,dx-4xyz\,dy-3xy^2\,dz.
\]
Pulling this expression to $B$ gives zero. It is a nonzero ambient-kernel vector with multiplier $v^{-1}$, which has a pole on the normalization.  

\end{enumerate}
\end{proof}



For ordinary surfaces the exact ambient-vanishing criterion is (A.1) with $K\otimes A=\eta_*\eta^*(N^*\otimes A)$. For surfaces with cuspidal-edge charts it is (A.1) with the locally specified sheaf $N^*\otimes t^{-1}B$. These specify the first group and the target of the connecting homomorphism, including the extra polar term at the cusp. Under $H^0(QA)=0$ they reduce respectively to the vanishing of these explicitly computed kernel sections. The general criterion (A.4)--(A.5) retains the full conductor and applies without a singularity-type restriction. The exact polynomial identities and exponent computations are recorded in \path{code/A_local_algebra.py} and its output.




\begin{proposition}[A cusp along a non-complete-intersection curve]\hypertarget{tag-01AK}{}\label{cdp:st:A.2.6}
\StatementRecord{cdp:st:A.2.6}{proved}{01AK}{\SeriesRef{cdp:st:A.5.5}}

Put
\[
 \begin{aligned}
 I&=(y^2-xz,\ yz^2-x^4,\ z^3-x^3y)\subset\C\{x,y,z\},\\
 F&=4x^7-15x^3yz^2+10x^2y^3z-3xy^5+4z^5.
 \end{aligned}                                                   \tag{A.NL.1}\hypertarget{tag-006J}{}\label{cdp:eq:A.NL.1}
\]
The reduced critical locus of the reduced hypersurface $F=0$ is the
irreducible curve $D$ parametrized by $(x,y,z)=(t^5,t^6,t^7)$.
Its transverse type at every point with $t\ne0$ is the ordinary cusp
$A_2$. The curve is not a local complete intersection at the origin,
and
\[
             F\in I^{(2)}\setminus I^2.                         \tag{A.NL.2}\hypertarget{tag-006K}{}\label{cdp:eq:A.NL.2}
\]
Thus the ordinary-square membership used to define the conormal
quadratic jet in \ref{cdp:st:A.5.5} does not extend to an arbitrary reduced
critical curve, even with generic transverse type $A_2$.

\end{proposition}


\begin{proof}

\begin{enumerate}[label={$\langle1\rangle$\arabic*.},leftmargin=*]
\item Write the three displayed generators of $I$ as $g_1,g_2,g_3$.
The quotient is a finite $\C\{x\}$-module generated by
$1,y,z,yz,z^2$: their products by $y$ and $z$ reduce to this list
using $y^2=xz$, $yz^2=x^4$, and $z^3=x^3y$.
For example these identities also give $y^5=x^6$ and $z^5=x^7$,
which make the analytic quotient finite over $\C\{x\}$.
The parametrization sends the five generators to
$1,t^6,t^7,t^{13},t^{14}$. Their exponents have distinct residues
modulo five, so they are linearly independent over $\C\{t^5\}$.
Consequently this ring homomorphism is injective and has
image $\C\{t^5,t^6,t^7\}$; it proves the stated normalization and
the primeness of $I$. The three generators are minimal. Indeed their
initial terms have degrees two, three, three; after the degree-two
coefficient vanishes, the cubic terms $yz^2$ and $z^3$ are independent
modulo the span of $xg_1,yg_1,zg_1$. A height-two ideal with three
minimal generators is not a complete intersection.

\item The following polynomial identities determine the critical locus:
\[
 \begin{aligned}
 F_y&=-15xg_1^2,\\
 F_x&=(17xyz-3y^3)g_1-28x^2g_2,\\
 F_z&=10x^2yg_1+20zg_3.
 \end{aligned}                                                   \tag{A.NL.3}\hypertarget{tag-006L}{}\label{cdp:eq:A.NL.3}
\]
At $x=0$, the equations $F_x=F_z=0$ give $y=z=0$.
At $x\ne0$, $F_y=0$ gives $g_1=0$, and then $F_x=0$
gives $g_2=0$. The value $z=0$ is impossible there, so $F_z=0$
also gives $g_3=0$. Conversely all three partials vanish on $D$.
Weighted Euler's identity, with weights $(5,6,7)$ and degree $35$,
puts $F$ in its Jacobian ideal. The critical locus is therefore
exactly $D$ as a reduced set. A repeated factor of $F$ would put
a two-dimensional component in this critical locus, which proves
that $F$ is reduced.

\item The Hessian along the normalization has the exact factorization
\[
 \operatorname{Hess}(F)|_D
 =2t^{21}
 \begin{pmatrix}7t^2\\0\\-5\end{pmatrix}
 \begin{pmatrix}7t^2&0&-5\end{pmatrix},
 \qquad \frac{1}{6}F_{yyy}|_D=-20t^{17}.                       \tag{A.NL.4}\hypertarget{tag-006M}{}\label{cdp:eq:A.NL.4}
\]
For $t\ne0$, the plane $x=t^5$ is transverse to $D$.
The restricted Hessian has rank one and its radical is the
$y$ direction; its cubic coefficient in that direction is nonzero.
The holomorphic splitting lemma and a local cube root of a unit
give the transverse cusp equation. At the generic point of $D$,
its conormal is locally free and injects into ambient differentials.
Since $F$ and $dF$ vanish there, $F\in I^2\C\{x,y,z\}_I$,
which is precisely $F\in I^{(2)}$.

\item To decide ordinary-square membership, use the positive weights
$(5,6,7)$. The three generators have weights $12,20,21$.
Among products of two generators, only $g_1^2$ can contribute
to weight $35$: it needs a coefficient of weight $11$, whose
space is spanned by $xy$. Products $g_1g_2$ and $g_1g_3$ would
need weights $3$ and $2$, which do not occur; the other products
already have greater weight. Thus the weight-$35$ part of $I^2$
is exactly $\C\,xyg_1^2$. Its $x^7$ coefficient is zero, whereas
that of $F$ is four. Taking the weight-$35$ component of a
convergent series identity is legitimate because all weights are
positive. This proves $F\notin I^2$ in the analytic local ring.

\item The normalization conductor
is $t^{10}\C\{t\}$: the semigroup gaps are $1,2,3,4,8,9$.
The $B$-module generated by the differential image, where $B=\C\{t\}$, is $t^4B\,dt$.
The smallest valuation of a two-by-two minor of the displayed
generators' Jacobian matrix is $19$. Thus the normalized conormal
determinant correction is $19-4=15$, while the conductor exponent
is $10$; the complete-intersection adjunction comparison cannot
identify these two quantities here.

\item The symbolic-square correction is visible directly on the free
normalized conormal. Set $a=dg_1|_D$ and
$b=dg_2|_D+2t^8a$. Then
\[
 \begin{aligned}
 a&=(-t^7,2t^6,-t^5),& b&=(-6t^{15},5t^{14},0),\\
 dg_3|_D&=-3t^9a+tb,&
 q&=t^{-5}(5t^8a-2b)^2.
 \end{aligned}                                                   \tag{A.NL.5}\hypertarget{tag-006N}{}\label{cdp:eq:A.NL.5}
\]
The first two rows generate the entire normalized conormal image,
and $q$ is the quadratic term, namely half the Hessian.
If $A,B$ are the dual basis, its radical frame $2A+5t^8B$
is represented in the ambient normal quotient by $t^{-6}\partial_y$.
Its intrinsic cubic is therefore $-20t^{-1}$. These are
quadratic and cubic poles of orders five and one. The corresponding
combined invariant $q^3k^2$ has pole order $17$.


\end{enumerate}
\end{proof}




\begin{proposition}[A complete nodal-family computation]\hypertarget{tag-01BK}{}\label{cdp:st:B1}
\StatementRecord{cdp:st:B1}{proved}{01BK}{the definitions and constructions stated in the proof}
\label{wbII:thm:nodal-test}\label{wbII:prop:family-geometry}\label{wbII:prop:nodal-counterexample}
Let $0<\varepsilon<4/27$, let $\Delta=\{|t|<\varepsilon\}$, and define
\[
 \mathcal C=\{y^2z=x^2(x+z)+tz^3\}\subset\PP^2\times\Delta,
 \qquad X=\mathcal C\times\PP^1.
\]
The projection $X\to\Delta$ is proper, its total space is smooth, and its central fibre is $S=C_0\times\PP^1$, where $C_0$ is an irreducible nodal cubic. For each $\lambda\in\C^*$ let $A_\lambda$ be the line bundle whose pullback to the normalization $\nu:\PP^1_u\to C_0$ is trivial and whose sections in this trivialization satisfy $h(1)=\lambda h(-1)$. Put $A=\operatorname{pr}_1^*A_\lambda$ and $d_\lambda=1$ for $\lambda=1$, and $d_\lambda=0$ otherwise. In the order $q=0,1,2$, the coherent cohomology dimensions are
\[
\begin{array}{c|c}
\text{sheaf on }S & (h^0,h^1,h^2)\\\hline
 N^*_{S/X}\otimes A &(d_\lambda,d_\lambda,0)\\
 K_S\otimes A &(1,0,0)\\
 T_S\otimes A &(1,0,0)\\
 \widetilde\Omega^1_S\otimes A &(0,1+d_\lambda,d_\lambda)\\
 \Omega^1_S\otimes A &(1,1+d_\lambda,d_\lambda)\\
 \Omega^1_X|_S\otimes A &(1,1+d_\lambda,d_\lambda).
\end{array}                                                    \tag{B.1}\hypertarget{tag-008G}{}\label{cdp:eq:B.1}
\]
Here $T_S$ denotes the torsion of the cotangent sheaf, and $K_S$ its inverse image in the ambient cotangent sheaf. In particular the two proposed vanishings hold exactly for $\lambda\ne1$, and the ambient space of sections has dimension one for every $\lambda\in\C^*$. After shrinking $\Delta$, every $A_\lambda$ extends to a line bundle on $\mathcal C$ of the form $\OO_{\mathcal C}(P-Q)$ for holomorphic sections $P,Q$.
\end{proposition}



\begin{proof}

\begin{enumerate}[label={$\langle1\rangle$\arabic*.},leftmargin=*]
\item In the affine chart $z=1$, the defining equation is $F=y^2-x^3-x^2-t=0$. Its partial derivative in $t$ is $-1$, so the total space is smooth there. Its only point with $z=0$ is $[0:1:0]$. In the chart $y=1$ the equation is $z-x^3-x^2z-tz^3=0$, whose $z$ derivative is one at this point. Fibre singularities in the affine chart satisfy $y=0$ and $x(3x+2)=0$; they occur at $t=0,-4/27$. Our disc contains only the first. The quadratic term $y^2-x^2$ at the origin has two distinct linear factors, and its smoothing parameter has nonzero differential. The holomorphic Morse lemma with parameter therefore supplies local coordinates with $t=vw$. Properness follows from the closed inclusion into $\PP^2\times\Delta$.

\item The normalization is
\[
 u\longmapsto(x,y)=(u^2-1,u(u^2-1)),
\]
completed by the point at infinity. The two points above the node are $1,-1$. The descent sequence, an exact sequence of coherent $\OO_{C_0}$-modules, is
\[
0\longrightarrow A_\lambda\longrightarrow\nu_*\OO_{\PP^1}
 \xrightarrow{\ h\mapsto h(1)-\lambda h(-1)\ }\C_{\rm node}
 \longrightarrow0.                                      \tag{B.2}\hypertarget{tag-008H}{}\label{cdp:eq:B.2}
\]
Its cohomology is the two-term complex $\C\xrightarrow{1-\lambda}\C$. Thus $h^0(A_\lambda)=h^1(A_\lambda)=d_\lambda$.

\item Write $V=\Omega^1_{\mathcal C}|_{C_0}$. The normalization map induces a surjective $\OO_{C_0}$-linear morphism $\Omega^1_{C_0}\to\nu_*\Omega^1_{\PP^1}$; at a node $vw=0$, it sends $dv$ to $(dv,0)$ and $dw$ to $(0,dw)$. Its kernel is the length-one module generated by $v\,dw=-w\,dv$. Consequently
\[
 T_{C_0}=\C_{\rm node},\qquad
 \widetilde\Omega^1_{C_0}=\nu_*\OO_{\PP^1}(-2).
\]
The kernel of $V\to\nu_*\Omega^1_{\PP^1}$ is $\nu_*\OO_{\PP^1}$. Locally its generators are $w\,dv$ and $v\,dw$ on the two branches; the identification is the normal-crossing conormal identification of Part A with the globally fixed conormal $dt$. Tensoring by $A_\lambda$ preserves exactness, and finite pushforward and the projection formula identify the two resulting normalization terms with $\nu_*\OO$ and $\nu_*\OO(-2)$. Hence
\[
 (h^0,h^1)(C_0,V\otimes A_\lambda)=(1,1),\qquad
 (h^0,h^1)(C_0,\Omega^1_{C_0}\otimes A_\lambda)=(1,1).
\]
The product morphisms give the direct decompositions
\[
 \Omega^1_X|_S\otimes A
 =(V\otimes A_\lambda)\boxtimes\OO
 \ \oplus\ A_\lambda\boxtimes\OO(-2),
\]
and the analogous decomposition for $\Omega^1_S$. The torsion-free quotient is obtained by replacing the first factor by $\nu_*\OO(-2)$. The torsion is $\C_{\rm node}\boxtimes\OO$, its ambient inverse image is $\nu_*\OO\boxtimes\OO$, and $N^*_{S/X}=\OO_S$. The coherent Kunneth isomorphism and $H^0(\PP^1,\OO)=\C$, $H^1(\PP^1,\OO(-2))=\C$ now give every entry of \eqref{cdp:eq:B.1}. All other cohomology of these two line bundles on $\PP^1$ vanishes.

\item We give the extension explicitly. Fix $Q(0)$ at $u=0$, and take $P(0)$ at $u=a=(1-\lambda)/(1+\lambda)$; for $\lambda=-1$ take $a=\infty$. For $\lambda\ne0$ both points are in the smooth locus. With $b=0$, the rational function $r(u)=(u-a)/(u-b)$ has divisor $P(0)-Q(0)$, and a local section in its pullback trivialization descends precisely with scalar
\[
 \frac{r(1)}{r(-1)}
 =\frac{(1-a)(1+b)}{(1-b)(1+a)}=\lambda.
\]
For $a=\infty$ the same expression is its limit, namely $-1$. The projection $\mathcal C\to\Delta$ is a submersion at these points, so the holomorphic implicit-function theorem extends them to sections on a smaller disc. The point at infinity itself is a constant section. For $\lambda=1$ take $P=Q$ identically. Sections through the smooth fibre locus are effective Cartier divisors in the smooth surface $\mathcal C$; their difference therefore defines the claimed extension. Pullback along $X\to\mathcal C$ extends $A$. This proves the claim also for all roots of unity, including the identity.

\end{enumerate}
\end{proof}



The calculation recovers the local example in the supplied correction note, Lemma 4.3 and Theorem 4.4, with the gluing convention fixed by \eqref{cdp:eq:B.2}. The computation applies to roots of unity as well: the only exceptional scalar for the conormal vanishing is $1$. The symbolic computation and the exact gluing matrices are recorded in \texttt{code/B\_examples.py} and its output.




\begin{proposition}[A projective example with twelve length-one contributions]\hypertarget{tag-01BL}{}\label{cdp:st:B2}
\StatementRecord{cdp:st:B2}{proved}{01BL}{\SeriesRef{cdp:eq:B.6}, \SeriesRef{cdp:eq:B.5}, \SeriesRef{cdp:eq:B.7}, \SeriesRef{cdp:st:B1}}

Let $h:E\to\PP^1_t$ be the Weierstrass surface defined on the finite chart by
\[
 y^2z=x^3+(t^4+1)xz^2+t^6z^3,                          \tag{B.3}\hypertarget{tag-008I}{}\label{cdp:eq:B.3}
\]
and at $s=1/t$ by $x_\infty=s^2x$, $y_\infty=s^3y$, $z_\infty=z$. Let $O=[0:1:0]$ and $P=(0,t^3,1)$ be its two sections and put
\[
 M_E=\OO_E(P-O),\quad Y=E\times\PP^1,\quad
 g=h\circ\operatorname{pr}_E,\quad
 L=K_Y\otimes\operatorname{pr}_E^*M_E^{-1},\quad
 \mathcal E=T_Y\otimes L.
\]
The surface $E$ is smooth and projective. Its singular fibres are twelve irreducible nodal cubics, at the distinct roots $t_i$ of
\[
 D(t)=31t^{12}+12t^8+12t^4+4.
\]
For every such fibre $S_i=g^{-1}(t_i)$, with $A=(K_Y\otimes L^{-1})|_{S_i}$, we have
\[
H^0(S_i,N^*_{S_i/Y}\otimes A)=0,
\qquad H^0(S_i,\widetilde\Omega^1_{S_i}\otimes A)=0,
\]
\[
 (h^0,h^1,h^2)(S_i,\mathcal E|_{S_i})=(0,1,1).
\]
As coherent sheaves on the base, with noncanonical identifications of their one-dimensional fibres,
\[
 R^qg_*\mathcal E=0\ (q\ne2),\qquad
 R^2g_*\mathcal E=\bigoplus_{i=1}^{12}\C_{t_i}.        \tag{B.4}\hypertarget{tag-008J}{}\label{cdp:eq:B.4}
\]
In particular $(h^0,h^1,h^2,h^3)(Y,\mathcal E)=(0,0,12,0)$.
\end{proposition}



\begin{proof}

\begin{enumerate}[label=\arabic*.,leftmargin=*]
\item The transition rule realizes \eqref{cdp:eq:B.3} as a cubic in a projective plane bundle over $\PP^1$, with Weierstrass coefficients in $\OO(4)$ and $\OO(6)$; hence its total space is projective. Its discriminant is $-16D(t)$. Exact polynomial division gives
\[
 \gcd(D,D')=1,\qquad \gcd(D,t^4+1)=1.
\]
The leading coefficient $31$ is its nonzero value in the chart at infinity. Each discriminant root therefore has nonzero $c_4=-48(t^4+1)$ and is a nodal cubic. A simple discriminant zero is transverse to the nodal discriminant and makes the surface smooth at that node. Alternatively the Gr\"obner basis of the affine total-space singularity ideal $(F,F_x,F_y,F_t)$ is $(1)$. At the zero section smoothness follows from $F_z=1$ in the chart $y=1$. At infinity the coefficients are $a=1+s^4,b=1$, with discriminant $-16\cdot31$ at $s=0$.

\item The sections $P,O$ are disjoint and meet the smooth locus of every fibre. For $P$, a singular point would require both $a(t)=0$ and $t=0$, which is impossible. Each fibre is an integral arithmetic-genus-one curve with trivial dualizing sheaf. The line bundle $M_E$ has degree zero on every fibre and is nontrivial there. On a smooth fibre, a trivialization would give a meromorphic function with one simple pole and one simple zero, hence a degree-one map to $\PP^1$, contradicting genus one. On a nodal fibre the normalization calculation in \eqref{cdp:eq:B.2} says that two distinct smooth points determine a nontrivial degree-zero line bundle. Thus $H^q(E_t,M_E^{\pm1}|_{E_t})=0$ for both $q=0,1$: a nonzero section of a degree-zero line bundle on an integral curve is a trivialization, and Serre duality gives the assertion in degree one. Grauert's theorem, applied to the proper flat map $h$, yields
\[
 Rh_*M_E=Rh_*M_E^{-1}=0.                               \tag{B.5}\hypertarget{tag-008K}{}\label{cdp:eq:B.5}
\]

\item Let $Z\subset E$ be the reduced set of the twelve critical points. The relative dualizing sheaf is $\omega_{E/\PP^1}=h^*\OO(1)$. Indeed the residue frame $dx/(2y)$ has this transition under $x_\infty=s^2x,y_\infty=s^3y$, and extends as a dualizing frame at a node. It follows that $K_E=h^*\OO(-1)$. There is an exact sequence of coherent sheaves on $E$
\[
0\longrightarrow h^*\Omega^1_{\PP^1}
 \xrightarrow{dh}\Omega^1_E
 \longrightarrow\mathcal I_Z\otimes\omega_{E/\PP^1}
 \longrightarrow0.                                   \tag{B.6}\hypertarget{tag-008L}{}\label{cdp:eq:B.6}
\]
Away from $Z$ this is the ordinary relative cotangent sequence. At a node choose coordinates $vw=t$. The map from $\Omega^1_E$ to the relative dualizing line sends $dv\mapsto v\omega$, $dw\mapsto-w\omega$; its image is $(v,w)\omega$ and its kernel is generated by $w\,dv+v\,dw=dt$. This proves \eqref{cdp:eq:B.6} on every stalk, with the critical scheme reduced.

\item Tensor \eqref{cdp:eq:B.6} by $M_E^{-1}$. By \eqref{cdp:eq:B.5}, projection formula, and the sequence for the ideal of $Z$, its derived direct image has only degree one, namely
\[
 R^1h_*(\Omega^1_E\otimes M_E^{-1})
 =h_*(\omega_{E/\PP^1}\otimes M_E^{-1}|_Z),
 \qquad h_*(\Omega^1_E\otimes M_E^{-1})=0.             \tag{B.7}\hypertarget{tag-008M}{}\label{cdp:eq:B.7}
\]
The right side is the sum of twelve reduced skyscrapers. Each summand in \eqref{cdp:eq:B.7} has length one.

\item Contraction on the smooth surface gives the canonical isomorphism $T_E\otimes K_E\simeq\Omega^1_E$. Consequently the product tangent decomposition identifies
\[
\mathcal E=(\Omega^1_E\otimes M_E^{-1})\boxtimes\OO(-2)
 \ \oplus\ (K_E\otimes M_E^{-1})\boxtimes\OO.
\]
The second summand has zero derived direct image by \eqref{cdp:eq:B.5}. The first has only the tensor product of degree-one terms: \eqref{cdp:eq:B.7} and $H^1(\PP^1,\OO(-2))=\C$. The relative coherent K\"unneth comparison gives \eqref{cdp:eq:B.4}. The Leray spectral sequence over a curve then gives the four stated global dimensions.

\item On a singular fibre $S_i=E_{t_i}\times\PP^1$, the restriction of $M_E$ is a nontrivial $A_\lambda$. Statement \ref{cdp:st:B1} gives the two vanishings and ambient dimensions $(1,1,0)$. Adjunction for the Cartier fibre gives $\omega_{S_i}=K_Y|_{S_i}$, because its normal line is trivial. Serre duality for this compact Gorenstein surface identifies, for $q=0,1,2$,
\[
 H^q(S_i,\mathcal E|_{S_i})^\vee
 =H^{2-q}(S_i,\Omega_Y^1|_{S_i}\otimes A).
\]
The fibre value $h^1=1$ with $R^1g_*\mathcal E=0$ is the degree-one Tor term obtained by tensoring a reduced skyscraper in degree two with the residue field.

\end{enumerate}
\end{proof}






\begin{proposition}[Topology of the Weierstrass product]\hypertarget{tag-01BM}{}\label{cdp:st:B3}
\StatementRecord{cdp:st:B3}{proved}{01BM}{\SeriesRef{cdp:st:B2}; \cite{CDP20}}

The threefold $Y$ of Statement \ref{cdp:st:B2} has $H^1(Y,\Z)=0$, $H^2(Y,\Z)\simeq\Z^{11}$, algebraic dimension $3$, and $e(Y)=24$. Its map $g$ is holomorphic, proper, surjective, and has connected fibres. For the line $A$ and the nodal fibre $S_i$ in Statement \ref{cdp:st:B2}, the two conormal-sequence endpoint groups vanish and $h^0(S_i,\Omega_Y^1|_{S_i}\otimes A)=1$. These are the groups used in the fibrewise inference of \cite[Section~7]{CDP20}. The line $M_E$ restricts to a nontorsion line on each nodal cubic. Theorem~2.2 of that paper assumes $H^2(Y,\Z)=0$ and algebraic dimension one.
\end{proposition}


\begin{proof}
\begin{enumerate}[label=\arabic*.,leftmargin=*]
\item 
The residue calculation above gives $K_E=h^*\OO(-1)$. Relative duality gives $R^1h_*\OO_E=\OO(-1)$, since every fibre is an integral genus-one curve, and $h_*\OO_E=\OO$. Thus $q(E)=p_g(E)=0$ and $H^0(E,K_E^{\otimes m})=0$ for every positive $m$. The invariants $q$ and $P_2$ are unchanged under point blowups. Castelnuovo's criterion \cite[Theorem~2.26, p.~178]{Hassett2009Rational}, applied to a smooth minimal model of $E$, therefore makes $E$ rational. A rational smooth projective surface is obtained from a minimal rational surface by point blowups; its first integral homology vanishes and its second integral cohomology is free. Euler additivity gives $e(E)=12$, since the twelve nodal fibres each have Euler characteristic one and the smooth elliptic fibre has Euler characteristic zero. Poincar\'e duality gives $b_2(E)=10$. The integral K\"unneth isomorphism for $Y=E\times\PP^1$ gives the asserted groups and Euler number. Projectivity and dimension give $a(Y)=3$. Connectedness and surjectivity follow directly from the integral cubic fibres of $h$ and the second product factor.

\item At a nodal fibre write its cubic as $y^2=(x-r)^2(x+2r)$, where $r\ne0$. Its normalization is $x=v^2-2r$, $y=v(v^2-3r)$, and the two preimages of the node are $v=\pm\sqrt{3r}$. At $P$ one has $v_P^2=2r$, while $O$ is the point at infinity. The rational function with divisor $P-O$ on the normalization gives the gluing multiplier $\lambda=(v_P-\sqrt{3r})/(v_P+\sqrt{3r})$, up to inversion. Writing $q=v_P/\sqrt{3r}$ gives $q^2=2/3$ and $\lambda^2+10\lambda+1=0$. Both roots are real and neither is $1$ or $-1$, so $\lambda$ has infinite order. Thus $M_E$ restricts to a nontorsion line on every nodal fibre. The two endpoint vanishings and the one-dimensional ambient cotangent group follow from Statement \ref{cdp:st:B2} and its Serre-duality identification.

\end{enumerate}
\end{proof}



The construction in Statement \ref{cdp:st:B2} supplies the compact local counterexample, including its bundles, cohomology groups and direct-image lengths. The universal assertion under every hypothesis of CDP20 Theorem 2.2 is proved in \ref{cdp:st:C4e} below: positive third Chern number forces one fixed irreducible nonnormal reduced support whose normalization has the required nonzero group for every ambient Picard twist. Its proof uses the full-fibre comparisons of A.3 and the conductor arguments of A.5. The named-family tests are \ref{cdp:st:C4g}.




\begin{proposition}[All thickenings and their lifting maps]\hypertarget{tag-01AL}{}\label{cdp:st:A.3.1}
\StatementRecord{cdp:st:A.3.1}{proved}{01AL}{\cite[Tags 0B3P and 0C4T, Divisors chapter p.~26; Tag 01DZ, Cohomology chapter pp.~2--3]{StacksProject}}

Let $S$ be any effective Cartier divisor in a complex manifold, let $E$ be any locally free sheaf on that manifold, and let $E_k=E|_{kS}$. For each $k\ge2$ there is a canonical exact sequence
\[
0\longrightarrow (E|_S)\otimes(N^*)^{k-1}
 \longrightarrow E_k\longrightarrow E_{k-1}\longrightarrow0.
\tag{A.10}\hypertarget{tag-003Z}{}\label{cdp:eq:A.10}
\]
Write $G_{k-1}=(E|_S)\otimes(N^*)^{k-1}$ and let
$\epsilon_k:H^0((k-1)S,E_{k-1})\to H^1(S,G_{k-1})$
be the connecting map. Then
\[
0\longrightarrow H^0(S,G_{k-1})\longrightarrow H^0(kS,E_k)
 \longrightarrow\ker\epsilon_k\longrightarrow0.
\tag{A.11}\hypertarget{tag-0040}{}\label{cdp:eq:A.11}
\]
In particular $H^0(kS,E_k)=0$ exactly when $H^0(S,G_{k-1})=0$ and $\epsilon_k$ is injective. Starting with $H^0(S,E|_S)$, equations \eqref{cdp:eq:A.10}--\eqref{cdp:eq:A.11} compute the group on $mS$ through all its extension maps. The vanishing of every $H^0(S,G_j)$, $0\le j<m$, is a sufficient criterion. These statements apply to $E=\Omega_Y^1\otimes B$ for every line bundle $B$, and include smooth, normal-crossings, and all other reduced supports.
\end{proposition}


\begin{proof}
\begin{enumerate}[label={$\langle1\rangle$\arabic*.},leftmargin=*]
\item 
\noindent\textbf{Step 1.} With $I=\OO_Y(-S)$, the sequence is the tensor product of $E$ with
$0\to I^{k-1}/I^k\to\OO_Y/I^k\to\OO_Y/I^{k-1}\to0$.
The Cartier property identifies $I^{k-1}/I^k$ with $(N^*)^{k-1}$ on $S$; local freeness makes the tensor product exact. Taking global sections proves \eqref{cdp:eq:A.11} and its equivalence. Induction on $k$ proves the sufficient criterion. No splitting of the thickening, cyclic base change, or assumption about the normal bundle has been used.

\end{enumerate}
\end{proof}






\begin{proposition}[Sections on a multiple abelian fibre]\hypertarget{tag-01AM}{}\label{cdp:st:A.3.2}
\StatementRecord{cdp:st:A.3.2}{proved}{01AM}{the definitions and constructions stated in the proof}

There is a compact smooth complex threefold $g:Y\to\PP^1$ with connected general fibre an abelian surface and a smooth fibre support $S$ of multiplicity $m$, for every $m\ge3$, and a line bundle $B$ on $Y$, such that
\[
H^0(S,(\Omega_Y^1\otimes B)|_S)=0,
\qquad h^0(mS,(\Omega_Y^1\otimes B)|_{mS})=3.
\tag{A.12}\hypertarget{tag-0041}{}\label{cdp:eq:A.12}
\]
Its conormal has exact order $m$.
\end{proposition}


\begin{proof}

\begin{enumerate}[label=\arabic*.,leftmargin=*]
\item Take an abelian surface $F$ and a translation $a:F\to F$ of exact order $m$. Let $\zeta$ be a primitive $m$th root of unity. The action generated by $(p,z)\mapsto(a(p),\zeta z)$ on $F\times\PP^1$ is free, so the quotient $Y$ is a smooth compact threefold and the quotient map is finite \'etale. The invariant map $t=z^m$ gives $g:Y\to\PP^1$. Away from $0$ and $\infty$ its fibre is isomorphic to $F$: the $m$ lifts of a base point are permuted transitively and freely. At $0$ the fibre is $mS$, where $S=F/\langle a\rangle$. Set $B=\OO_Y(-S)$ and $E=\Omega_Y^1\otimes B$.

\item The connected cyclic cover $F\to S$ identifies $N^*=B|_S$ with the character line of weight one. Its order is exactly $m$: otherwise a nontrivial proper character power would have a nonzero invariant trivialization on the connected compact cover, whose holomorphic functions are constant. Pullback to $F$ splits the ambient cotangent bundle into two invariant tangent forms and the normal form $dz$ of weight one. Consequently
\[
 (\Omega_Y^1\otimes B)|_S\simeq (N^*)^{\oplus2}\oplus(N^*)^2.
\]
The nontrivial flat lines on the abelian surface $S$ have no holomorphic sections. Their pullbacks to $F$ are trivial; a section would be a constant function transforming by a nontrivial character, which is zero. Since $m\ge3$, both characters in the preceding splitting are nontrivial, proving the first equality.

\item The inverse image of $mS$ is the ringed space with structure sheaf $\OO_F\otimes\C[z]/(z^m)$. Write $e$ for the local generator of the pulled-back line $B$; it has weight one. All global sections on this inverse image are linear combinations of $z^k\alpha_1e,z^k\alpha_2e,z^kdz\,e$ for $0\le k<m$, where $\alpha_1,\alpha_2$ are a basis of invariant one-forms on $F$. The first two have weight $k+1$ and the third weight $k+2$. The invariant monomials are exactly
$z^{m-1}\alpha_1e,z^{m-1}\alpha_2e,z^{m-2}dz\,e$.
Taking invariants is exact over $\C$, by averaging over the finite group. These three monomials form a basis. Translations act identically on $H^q(F,\OO_F)$, whose dimensions are $\binom{2}{q}$. The same calculation gives all fibre groups: $h^q(S,E|_S)=(0,0,0)$ and $h^q(mS,E|_{mS})=(3,6,3)$ for $q=0,1,2$. On the compact threefold the two tangent summands pull back to $\OO_F\boxtimes\OO_{\PP^1}(-1)$ and have no cohomology. The normal summand is $\OO_F\boxtimes\Omega^1_{\PP^1}(-\{0\})$. Its first cohomology on $\PP^1$ has the \v Cech basis $dz/z,dz$, of weights zero and one, respectively. Indeed sections near zero are multiples of $z\,dz$, and sections near infinity are multiples of $z^{-2}dz$, leaving exactly those two Laurent classes. Its invariant part is one-dimensional. K\"unneth and finite-cover descent therefore give $h^q(Y,E)=(0,1,2,1)$ for $q=0,1,2,3$.

\end{enumerate}
\end{proof}



For this example the complete dimensions on $kS$, $1\le k\le m$, are $0$ for $k\le m-2$, $1$ for $k=m-1$, and $3$ for $k=m$; the same character calculation proves this. More generally the character count with $B=\OO_Y(-bS)$ is
\[
\begin{aligned}
h^0(kS,(\Omega_Y^1\otimes B)|_{kS})
&=2\#\{0\le j<k:j+b\equiv0\pmod m\}\\
&\quad+\#\{0\le j<k:j+b+1\equiv0\pmod m\}.
\end{aligned}
\tag{A.13}\hypertarget{tag-0042}{}\label{cdp:eq:A.13}
\]
The finite counts for $3\le m\le12$ and all characters are checked in \path{code/A_multiple_fibres.py}.




\begin{proposition}[The precise finite-translation step in Remarks 7.2 and 7.4]\hypertarget{tag-01AN}{}\label{cdp:st:A.3.3}
\StatementRecord{cdp:st:A.3.3}{proved}{01AN}{\SeriesRef{cdp:st:A.3.1}}

For a fixed multiple fibre $mS$, suppose that the set of line bundles $B$ for which $H^0(S,(\Omega_Y^1\otimes B)|_S)=0$ contains a dense open subset $U$ of the relevant Picard component. The intersection
\[
\bigcap_{j=0}^{m-1}\{B:B\otimes\OO_Y(-jS)\in U\}
\tag{A.14}\hypertarget{tag-0043}{}\label{cdp:eq:A.14}
\]
is a dense open subset on which the full-fibre ambient group vanishes, whenever the displayed translations preserve that component. The same assertion for a finite union of components follows by using the appropriate dense open subset in each translated component.
\end{proposition}


\begin{proof}
\begin{enumerate}[label={$\langle1\rangle$\arabic*.},leftmargin=*]
\item 
\noindent\textbf{Step 1.} Tensoring by the fixed line bundle $\OO_Y(-jS)$ is a biholomorphic translation between Picard components. Its inverse image of a dense open set is dense and open. The finite intersection retains these properties. Its defining condition is exactly the vanishing of every graded group in \ref{cdp:st:A.3.1}, which proves the result.

\end{enumerate}
\end{proof}



This finite-translation argument supplies the thickening step intended in Remark~7.2 once ambient vanishing on the smooth support is established. On a smooth support, the conormal sequence has $Q=\Omega_S^1$ and $K=N^*$, so the two requisite support vanishings do establish that ambient vanishing. Remark~7.4 requires the same argument on a singular support; a statement solely about $Q=\widetilde\Omega_S^1$ leaves the kernel in \ref{cdp:st:A.1.1} to be checked. Thus the thickening operation itself has a complete exact treatment, while the cited singular-support assertion cannot supply its missing kernel hypothesis. Example \ref{cdp:st:A.3.2} disproves the inference from a single support vanishing even with torsion conormal; it does not refute a generic-vanishing assertion already established for all the finitely translated twists.




\begin{proposition}[The smooth-multiple assertion of Remark 7.2]\hypertarget{tag-01AO}{}\label{cdp:st:A.3.4}
\StatementRecord{cdp:st:A.3.4}{proved}{01AO}{\SeriesRef{cdp:st:D.1.9}, \SeriesRef{cdp:st:A.4.9}, \SeriesRef{cdp:st:A.4.7}, \SeriesRef{cdp:st:A.3.1}}

Under the hypotheses of CDP20 Proposition~7.3, let $mS$, with $m\ge2$, be a full fibre whose reduced support $S$ is smooth. For every real irrational transgression parameter $r$ of \ref{cdp:st:D.1.9},
\[
 H^0\bigl(mS,(\Omega_Y^1\otimes L_r)|_{mS}\bigr)=0.
 \tag{A.15a}\hypertarget{tag-0045}{}\label{cdp:eq:A.15a}
\]
The same simultaneous vanishing on all smooth multiple fibres $(m\ge2)$ holds on the complement of a proper complex analytic subset of $\Pic(Y)$. The assertion includes the entire scheme-theoretic fibre, with every nilpotent layer.
\end{proposition}


\begin{proof}

\begin{enumerate}[label={$\langle1\rangle$\arabic*.},leftmargin=*]
\item Choose a disc $\Delta$ about the fibre value and write $U=g^{-1}(\Delta)$. Form the base change $t=s^m$ and normalize it, obtaining $p:U'\to U$ and $g':U'\to\Delta'$. This map $p$ is finite étale of degree $m$. Indeed, near a point of the smooth support write $g=u z^m$ with $u$ a holomorphic unit. The normalization is given by
\[
 w^m=u,\qquad s=wz.
\]
The derivative $m w^{m-1}$ is invertible, so these equations give an unramified cover of the original smooth chart. They cover the normalized base change, since they coincide away from $z=0$ and are normal. The local constructions glue by uniqueness of normalization. The divisor $S'=(g')^{-1}(0)$ is smooth and reduced, $g'$ is a submersion near it, and $p^{-1}(mS)=mS'$ as Cartier divisors.

\item Pull back the unitary local system (A.22) on $U$ to $U'$; its associated holomorphic line is $A'=p^*L_r$. Every nearby smooth fibre of $g'$ maps isomorphically onto the corresponding nearby fibre of $g$. Its character is therefore the irrational exponential of a nonzero integral class. The smooth-fibre calculation in \ref{cdp:st:A.4.9} gives zero cohomology in degrees zero and one for that local system. Proper smooth local triviality at $S'$ transports the same cohomology statement to $S'$. If a component presentation of the cover is used, the argument applies on each component and the groups add.

\item A holomorphic section of a unitary-flat line on a compact smooth complex manifold is parallel: its squared norm is plurisubharmonic and hence constant by the maximum principle. Thus $H^0(S',A')=0$ follows from the zero degree-zero local-system cohomology. \ref{cdp:st:A.4.7} gives $H^0(S',\Omega_{S'}^1\otimes A')=0$ from degree-one vanishing. The conormal of the reduced full fibre $S'$ is trivial, with global generator $dg'$. Its conormal sequence, tensored by $A'$, now gives
$H^0(S',(\Omega_{U'}^1\otimes A')|_{S'})=0$.

\item All the graded factors for $mS'$ in \ref{cdp:st:A.3.1} equal this same restricted bundle, because the conormal is trivial. Induction in (A.10) proves $H^0(mS',(\Omega_{U'}^1\otimes A')|_{mS'})=0$. Since $p$ is étale, $p^*\Omega_U^1=\Omega_{U'}^1$, and finite faithfully flat pullback injects sections on $mS$ into sections on $mS'$. One can verify that injection locally by tensoring the inclusion of a module into its faithfully flat scalar extension, or by the trace divided by $m$. This proves \eqref{cdp:eq:A.15a}.

\item There are finitely many non-submersion values: the critical locus is a compact analytic set with finitely many irreducible components, and $dg=0$ makes $g$ constant on the smooth locus, hence on the closure, of each component. Fix the root cover above each smooth multiple fibre and write its smooth central surface as $S'_c$. The exponential sequence identifies $\Pic(Y)$ with the complex vector space $H^1(Y,\OO_Y)$. Exponentiating linear combinations of a basis of holomorphic Cech cocycles gives the actual Picard family. On $S'_c$ consider the holomorphic family of bundles
\[
 E_B=(\Omega^1_{U'_c}\otimes p_c^*B)|_{S'_c}.
\]
We give the analytic nonvanishing locus by an explicit finite kernel map. Near a parameter $B_0$, smooth gauges obtained from a partition of unity and holomorphic parameter-dependent logarithms of the transition ratios identify the underlying smooth bundles. Their Dolbeault operators have the form $D_b=D_0+A_b$, where $A_b$ has order zero, depends holomorphically on $b$, vanishes at $B_0$, and satisfies
\[
 D_0A_b+A_bD_0+A_b^2=0.
\]
Choose Hermitian metrics on $S'_c$ and the central bundle. Let $i$ include the finite-dimensional harmonic spaces, let $\pi$ be harmonic projection, let $G$ be the Green operator for the Dolbeault Laplacian, and put $h=D_0^*G$. The classical Green and Hodge identities give
\[
 D_0h+hD_0=1-i\pi,\qquad \pi i=1,
 \qquad \pi h=0,\qquad hi=0;
\]
see \cite[Part~I, Theorem~3.15, author-PDF body p.~15, and Theorems~7.1--7.2, body p.~32]{DemaillyHodgeAuthor2002}. These results apply to compact Hermitian manifolds without a K\"ahler hypothesis. On fixed sufficiently high Sobolev spaces, shrink the parameter neighbourhood so that the Neumann inverses $T_b^q=(1+hA_b)^{-1}$ exist in the degrees and adjacent regularities used below. Define the finite holomorphic matrix
\[
 M_b=\pi^1 A_b T_b^0 i^0:
 \mathcal H^0(E_{B_0})\longrightarrow\mathcal H^1(E_{B_0}).
\]
The two inverse kernel maps are $v\mapsto T_b^0i^0v$ and $u\mapsto\pi^0u$. Indeed, $D_bu=0$ implies
\[
 u-i\pi u=hD_0u=-hA_bu,
 \qquad u=T_b^0i^0\pi^0u.
\]
Conversely, for $u=T_b^0i^0v$, the contraction and integrability identities give
\[
 \begin{aligned}
 D_bu&=i\pi A_bu+hD_0A_bu
       =i\pi A_bu-hA_bD_bu,\\
 D_bu&=T_b^1i^1M_bv,\qquad
 \pi^0T_b^0i^0v=v.
 \end{aligned}
\]
Thus $\ker M_b$ is isomorphic by these actual maps to the section space of $E_B$; elliptic regularity makes every resulting Sobolev solution smooth. The nonzero-section locus is the zero set of the maximal minors expressing failure of full column rank of $M_b$, with the zero-domain and insufficient-row cases interpreted directly. It is therefore a closed complex analytic subset $Z_c$ of the parameter space. The conclusion follows from this kernel calculation, not from numerical upper semicontinuity alone. Step~3 shows that $Z_c$ excludes every real irrational $L_r$. Outside $Z_c$, all the graded factors on $mS'_c$ vanish, and Step~4 gives the full-fibre vanishing after descent. The finite union of the $Z_c$ is proper because the same $L_r$ lies outside every member. This proves the simultaneous generic assertion and the smooth-multiple conclusion of Remark~7.2. For a singular support the identical root cover retains its singularities and the kernel $N^*\otimes\mathcal J^{-1}$; the present proof does not erase that term in Remark~7.4.

\end{enumerate}
\end{proof}






\begin{proposition}[Every reducible full fibre has generic ambient vanishing]\hypertarget{tag-01AP}{}\label{cdp:st:A.3.5}
\StatementRecord{cdp:st:A.3.5}{proved}{01AP}{the definitions and constructions stated in the proof}

Let $Y$ be a smooth compact complex threefold with $H^2(Y,\Q)=0$, and let $F=\sum_i m_iS_i$ be an effective Cartier divisor with connected reducible support. Every connected component of $\Pic(Y)$ contains a dense open subset, with proper closed analytic complement, on which
\[
 H^0(F,(\Omega_Y^1\otimes L)|_F)=0.
 \tag{A.15b}\hypertarget{tag-0046}{}\label{cdp:eq:A.15b}
\]
The multiplicities may be unequal. For a proper connected fibration over a compact curve, the vanishing is simultaneous on all reducible full fibres after intersecting finitely many such open sets. Consequently a full fibre with ambient sections for every Picard twist has irreducible support.
\end{proposition}


\begin{proof}

\begin{enumerate}[label=\arabic*.,leftmargin=*]
\item Every component $S_i$ meets a different component $S_j$, because the support is connected and reducible. The canonical section of $\OO_Y(S_j)$ restricts to a nonzero section on $S_i$. It is a nonzerodivisor in every local ring of $S_i$: otherwise one local branch of $S_i$ would be contained in $S_j$, and analytic continuation on the global irreducible component would give $S_i=S_j$. Its zero scheme is therefore a nonzero effective Cartier divisor $D$ on $S_i$.

\item For any fixed locally free sheaf $E$ on the reduced irreducible compact surface $S_i$, the injections
\[
 H^0(S_i,E(-nD))\hookrightarrow H^0(S_i,E)
\]
form a decreasing sequence of finite-dimensional subspaces. Their intersection is zero. At a point of $D$, a section in the intersection has its germ in every power of the Cartier ideal; the Krull intersection theorem for the Noetherian analytic local ring makes this germ zero. The identity theorem on the reduced irreducible surface, applied in local trivializations of $E$, then makes the global section zero. A descending sequence of finite-dimensional subspaces stabilizes, so $H^0(S_i,E(-nD))=0$ for all sufficiently large $n$.

\item Fix a Picard component and a representative $L_0$ in it. The integral class $c_1(\OO_Y(S_j))$ is torsion because $H^2(Y,\Q)=0$; choose a positive integer $q$ killing it. Then $\OO_Y(-qnS_j)$ belongs to $\Pic^0(Y)$ for every $n$. Apply the preceding argument to $E=(\Omega_Y^1\otimes L_0)|_{S_i}$. It supplies, in the chosen Picard component, a twist $L_0\otimes\OO_Y(-qnS_j)$ with zero ambient group on $S_i$. Local holomorphic Picard families and proper semicontinuity show that the nonvanishing locus on that component is a proper closed analytic subset. Its complement is dense and open. The same argument works in every Picard component.

\item Choose an ordered list $A_1,\ldots,A_N$ of the prime divisors $S_i$, with $S_i$ repeated $m_i$ times, and put $F_0=0$, $F_k=A_1+\cdots+A_k$. The exact Cartier sequences, tensored with the locally free ambient bundle, are
\[
\begin{split}
0\longrightarrow
 (\Omega_Y^1\otimes L\otimes\OO_Y(-F_{k-1}))|_{A_k}
 \longrightarrow (\Omega_Y^1\otimes L)|_{F_k}\\
 \longrightarrow(\Omega_Y^1\otimes L)|_{F_{k-1}}
 \longrightarrow0.
\end{split}
\]
This follows locally from multiplication by an equation of $F_{k-1}$, and remains exact when a component occurs repeatedly. For each $k$ take the dense open set just proved on $A_k$ and translate it by the fixed global line $\OO_Y(-F_{k-1})$, using the appropriate Picard component. Their finite intersection is dense open. Every displayed first term has no sections there, so induction from $F_0=0$ proves \eqref{cdp:eq:A.15b} on the full scheme $F_N=F$.

\item A proper holomorphic map from a smooth compact manifold to a smooth curve has finitely many critical values, since its proper analytic critical-value set does not contain a general point. Its smooth connected fibres are irreducible. Thus there are finitely many reducible fibres, and a further finite intersection gives the asserted simultaneous vanishing.

\end{enumerate}
\end{proof}






\begin{proposition}[Cotangent groups on the rational normalization fibre]\hypertarget{tag-01AQ}{}\label{cdp:st:A.4.1}
\StatementRecord{cdp:st:A.4.1}{proved}{01AQ}{\SeriesRef{cdp:st:A.1.2}, \SeriesRef{cdp:st:A.1.1}}

For the full reduced normal-crossings fibre $W$ in the standard $(3,4,\infty)$ six-sphere member $X$, with normalization $\eta:\operatorname{dP}_6\to W$, and every $L\in\Pic(X)$, put $A=(K_X\otimes L^{-1})|_W$. Then
\[
H^0(W,Q\otimes A)=0,
\qquad H^0(W,\Omega_X^1|_W\otimes A)\simeq\C.
\tag{A.15}\hypertarget{tag-0044}{}\label{cdp:eq:A.15}
\]
The same two conclusions hold with $A=B|_W$ for every $B\in\Pic(X)$.
\end{proposition}



The nonvanishing on the compact example and the torsion-free-cotangent vanishing are contained in \citet[Theorem~10.5 and Section~10.4]{AC2026}.

\begin{proof}

\begin{enumerate}[label=\arabic*.,leftmargin=*]
\item The standard member satisfies $H^2(X,\Z)=0$, and the normalization of $W$ is the blowup of $\PP^2$ at three points. A blowup at a smooth point has $\pi_*\OO=\OO$ and $R^1\pi_*\OO=0$. The latter follows from the theorem on formal functions: successive exceptional-neighborhood quotients are $\OO_{\PP^1}(k)$, $k\ge0$, whose first cohomology is zero. Thus $H^1(\operatorname{dP}_6,\OO)=H^1(\PP^2,\OO)=0$. The exponential sequence makes $c_1:\Pic(\operatorname{dP}_6)\to H^2(\operatorname{dP}_6,\Z)$ injective. The pullback of every line bundle on $X$ has zero first Chern class and is therefore trivial.

\item The differential of normalization gives $\Omega_W^1\to\eta_*\Omega^1_{\operatorname{dP}_6}$. Its target is torsion-free as an $\OO_W$-module, so it factors through $Q$. The resulting map is injective because it is an isomorphism on the dense regular locus and its source is torsion-free. A holomorphic one-form on the blowup restricts away from the exceptional curves to a one-form on $\PP^2$ minus three points. Its coefficients extend over these points by Hartogs' theorem, and $H^0(\PP^2,\Omega^1_{\PP^2})=0$ by the Euler sequence. Hence $H^0(\operatorname{dP}_6,\Omega^1)=0$. The projection formula and the triviality of $\eta^*A$ prove the first equality in \eqref{cdp:eq:A.15}.

\item Because $W$ is the full reduced fibre, its conormal is trivial. \ref{cdp:st:A.1.2} identifies $K\otimes A$ with $\eta_*\eta^*A=\eta_*\OO_{\operatorname{dP}_6}$. The normalization is connected and compact, so its global holomorphic functions are $\C$. The cohomology sequence of \ref{cdp:st:A.1.1} and the first equality in \eqref{cdp:eq:A.15} identify this group with the ambient group. The same argument applies to every $B$.

\end{enumerate}
\end{proof}






\begin{proposition}[Constant twists on an elliptic ruled surface]\hypertarget{tag-01AR}{}\label{cdp:st:A.4.2}
\StatementRecord{cdp:st:A.4.2}{proved}{01AR}{the definitions and constructions stated in the proof}

For an elliptic curve $E$, the minimal ruled surface $S_0=E\times\PP^1$ and the constant family $M_t=\OO_{S_0}$ satisfy $H^0(S_0,\Omega^1_{S_0}\otimes M_t)\ne0$ for every parameter. The ruling has genus one. These data give a counterexample to the genus-at-least-two inference in \cite[Proposition~7.3 and Step~1, p.~693]{CDP20}, where the line-bundle family is permitted to be constant.
\end{proposition}


\begin{proof}
Take $S_0=E\times\PP^1$ for an elliptic curve $E$ and take $M_t=\OO_{S_0}$ for every parameter. Every rational curve on $S_0$ maps constantly to $E$ by Riemann--Hurwitz, so it is a ruling fibre and has self-intersection zero. Thus $S_0$ has no exceptional curve of self-intersection $-1$ and is minimal. The product cotangent decomposition gives
\[
 h^q(S_0,\Omega^1_{S_0})=(1,2,1),\qquad q=0,1,2.
\]
Indeed the $E$ cotangent summand contributes $(1,1,0)$, while the $\PP^1$ cotangent summand contributes $(0,1,1)$ by the K\"unneth formula and $h^1(\PP^1,\OO(-2))=1$. A nonzero invariant form on $E$ provides the asserted sections for every parameter, whereas the ruling has genus one.

\end{proof}






\begin{proposition}[Twisted differentials on a ruled normal-crossings surface]\hypertarget{tag-01AS}{}\label{cdp:st:A.4.3}
\StatementRecord{cdp:st:A.4.3}{proved}{01AS}{the definitions and constructions stated in the proof}

Let $E$ be an elliptic curve, let $P\in\Pic(E)$, and let $T=\PP_E(\OO_E\oplus P)$ with its zero and infinity sections $D_0,D_\infty$. Glue $D_0$ to $D_\infty$ by a translation of $E$ to form the compact normal-crossings surface $S$, and denote its normalization by $\eta:T\to S$. Every line bundle $A$ on $S$ whose pullback $\eta^*A$ is topologically trivial satisfies
\[
H^0(S,\widetilde\Omega_S^1\otimes A)=
\begin{cases}\C,&A\simeq\OO_S,\\0,&A\not\simeq\OO_S.\end{cases}
\tag{A.16}\hypertarget{tag-0047}{}\label{cdp:eq:A.16}
\]
\end{proposition}


\begin{proof}

\begin{enumerate}[label=\arabic*.,leftmargin=*]
\item For the ruling $p:T\to E$, every topologically trivial line bundle on $T$ is $p^*P_A$ for a unique $P_A\in\Pic^0(E)$. Indeed its degree on every ruling fibre is zero, so the adjunction map from the pullback of its rank-one direct image is an isomorphism, checked on each fibre by $H^0(\PP^1,\OO)=\C$ and $H^1(\PP^1,\OO)=0$. The line on $E$ has degree zero because its first Chern class pulls back to zero and $p$ has a section.

\item The relative cotangent sequence and restriction to ruling fibres give
\[
 H^0(T,\Omega_T^1\otimes p^*P_A)
 =p^*H^0(E,\Omega_E^1\otimes P_A).
\]
The relative quotient has no sections because its restriction to every fibre is $\OO_{\PP^1}(-2)$. Since $\Omega_E^1\simeq\OO_E$, the displayed group vanishes for nontrivial $P_A$. The injection into normalization differentials gives the vanishing in \eqref{cdp:eq:A.16}.

\item For $P_A=\OO_E$, choose a trivialization of $\eta^*A$. The descent datum for $A$ is a nonzero constant $\lambda$ identifying the two restrictions to $D_0,D_\infty$, because these are trivial lines on the compact connected curve $E$. Multiplying the chosen trivialization by a constant changes neither $\lambda$ nor its meaning. The local calculation for $xy=0$ identifies the image of $\widetilde\Omega_S^1$ in $\eta_*\Omega_T^1$ with pairs of forms whose tangential restrictions to the double curve agree. The unique normalization form is a scalar multiple of $p^*\alpha$, where $\alpha$ is a nonzero invariant form on $E$. Translation pulls $\alpha$ back to itself. Twisted descent is therefore the equation $\alpha=\lambda\alpha$, so a nonzero section exists exactly when $\lambda=1$. That equality also says precisely that the trivialization descends to $A\simeq\OO_S$.

\end{enumerate}
\end{proof}






\begin{proposition}[Canonical triviality on an irreducible reduced torus fibre]\hypertarget{tag-01AT}{}\label{cdp:st:A.4.4}
\StatementRecord{cdp:st:A.4.4}{proved}{01AT}{\SeriesRef{cdp:st:A.5.2}}

Let $g:Y\to C$ be a proper holomorphic map from a connected smooth compact threefold onto a smooth compact curve, with connected fibres and general fibre a complex two-torus. Every irreducible reduced full fibre $S=g^{-1}(c)$ has $\omega_S\simeq\OO_S$. When $H^2(Y,\Q)=0$, it therefore satisfies $h^q(S,\OO_S)=(1,2,1)$.
\end{proposition}


\begin{proof}

\begin{enumerate}[label=\arabic*.,leftmargin=*]
\item The coherent sheaf $g_*\omega_Y$ is torsion-free: a nonzero function on an open subset of the base cannot annihilate a holomorphic section of a line bundle on the dominating smooth total space. It has generic rank one because the smooth torus fibre has one holomorphic top form; after restricting to a base open set where Grauert base change holds, this identifies its rank. On the smooth curve it is consequently a line bundle $M$. The evaluation morphism $g^*M\to\omega_Y$ is generically nonzero, so its zero scheme is an effective Cartier divisor $D$ and $\omega_Y\simeq g^*M\otimes\OO_Y(D)$. No component of $D$ dominates the base: on a general smooth torus fibre its defining section is a nonzero holomorphic top form and is nowhere zero. Thus $D$ is vertical.

\item All vertical components meeting the specified irreducible full fibre are that fibre itself. Consequently $\OO_Y(D)|_S$ is a power of $\OO_Y(S)|_S$, and the latter is trivial because a coordinate on $C$ defines the full reduced fibre. The pullback $g^*M|_S$ is also trivial. Adjunction gives $\omega_S=(\omega_Y\otimes\OO_Y(S))|_S\simeq\OO_S$. The fibre is reduced and connected, so $h^0(\OO_S)=1$; it is a Cartier divisor, hence Cohen--Macaulay and Gorenstein, so Serre duality gives $h^2(\OO_S)=1$. The divisor Riemann--Roch calculation in \ref{cdp:st:A.5.2} gives $\chi(\OO_S)=0$ under the stated rational cohomology hypothesis, yielding $h^1=2$.

\end{enumerate}
\end{proof}






\begin{proposition}[Twisted cotangent groups of irreducible normal-crossings torus fibres]\hypertarget{tag-01AU}{}\label{cdp:st:A.4.5}
\StatementRecord{cdp:st:A.4.5}{proved}{01AU}{\SeriesRef{cdp:st:A.4.4}, \SeriesRef{cdp:st:A.4.1}, \SeriesRef{cdp:st:A.4.3}, \SeriesRef{cdp:st:A.1.2}; \cite[Definition~2.1]{SS2002}}

Let $g:Y\to C$ have the hypotheses of \ref{cdp:st:A.4.4} and let $S$ be a singular irreducible reduced full fibre with normal crossings. For every nontrivial unitary-flat line bundle $A$ on $S$,
\[
H^0(S,\widetilde\Omega_S^1\otimes A)=0,
\qquad
H^0(S,\Omega_Y^1|_S\otimes A)
 \simeq H^0(\widetilde S,\eta^*A).
\tag{A.17}\hypertarget{tag-0048}{}\label{cdp:eq:A.17}
\]
The latter group has dimension one when $\eta^*A$ is holomorphically trivial, and zero otherwise. The normalization is exactly one of the following: a diagonal Hopf surface with two disjoint elliptic curves identified; a ruled surface over an elliptic curve with two disjoint sections identified by translation; or a rational surface with an anticanonical hexagon as in Schr\"oer--Siebert Proposition~2.6.
\end{proposition}


\begin{proof}

\begin{enumerate}[label=\arabic*.,leftmargin=*]
\item \ref{cdp:st:A.4.4} gives $\omega_S\simeq\OO_S$. Dualizing the locally free conormal resolution of $\Omega_S^1$ yields
\[
 \mathcal T^1_S=\mathcal Ext^1_{\OO_S}(\Omega_S^1,\OO_S)
 \simeq\OO_S/\mathcal J=\OO_{S_{\rm sing,red}}.
\]
Here the conormal is trivial because $S$ is the full reduced fibre; the local identities $\mathcal J=(x,y)$ at a double point and $\mathcal J=(xy,xz,yz)$ at a triple point show that this is the reduced singular locus. Hence $S$ is d-semistable in the exact sense of \cite[Definition~2.1]{SS2002}. Its hypotheses are irreducibility, normal crossings, d-semistability and trivial dualizing sheaf. Propositions~2.4, 2.5 and 2.6 there give the three normalization cases. In the ruled case their Proposition~4.1 says precisely that trivial $\omega_S$ forces the section identification to be a translation. The comparison with their notation is $X_{\mathrm{SS}}=S$ and $S_{\mathrm{SS}}=\widetilde S$.

\item In the rational case $\eta^*A$ is trivial: its first Chern class is torsion because the line is unitary-flat, whereas the integral second cohomology of a smooth rational surface is torsion-free and its Picard group injects into that group by the exponential sequence. Such a surface has no holomorphic one-forms, by birational invariance as proved in \ref{cdp:st:A.4.1}. The injection $\widetilde\Omega_S^1\to\eta_*\Omega^1_{\widetilde S}$ therefore proves the first vanishing.

\item In the elliptic ruled case the proof of \ref{cdp:st:A.4.3} applies to a decomposable ruled surface of any degree: the two disjoint sections express it as $\PP_E(\OO_E\oplus P)$, with $P\in\Pic(E)$. The integral second cohomology of this ruled surface is torsion-free, by the projective-bundle formula. A unitary-flat line has torsion first Chern class, so its pullback is topologically trivial and comes from $\Pic^0(E)$. The cotangent calculation of \ref{cdp:st:A.4.3} applies to every $P\in\Pic(E)$ and gives the vanishing for every nontrivial $A$.

\item In the Hopf case, pull back to a finite primary diagonal Hopf cover. Its universal covering is $\C^2\setminus\{0\}$ with contraction $(z_1,z_2)\mapsto(a_1z_1,a_2z_2)$, where $|a_i|<1$. A holomorphic one-form valued in a unitary-flat line lifts to $f_1(z)dz_1+f_2(z)dz_2$, with each $f_i$ holomorphic on the punctured space and hence on $\C^2$ by Hartogs' theorem. For a contraction character $\rho$ of modulus one, equivariance says that the coefficient of $z_1^{k_1}z_2^{k_2}dz_i$ can be nonzero only when
\[
 a_1^{k_1}a_2^{k_2}a_i=\rho
\]
(or its inverse, according to the convention for the associated line). The modulus on the left is strictly less than one and the character has modulus one. Thus every coefficient is zero. Finite \'etale pullback is injective on sections, so the twisted one-form group on the Hopf normalization also vanishes, and again the normalization injection proves the first assertion on $S$.

\item Proposition~\ref{cdp:st:A.1.2} and the first vanishing identify the ambient group with $H^0(\widetilde S,\eta^*A)$. A section of a unitary-flat line on a compact connected complex manifold is either zero or parallel and nowhere zero: its squared norm is locally a sum of squares of holomorphic functions, hence plurisubharmonic, and the maximum principle makes it constant. It therefore trivializes the line, and in the trivial case the section space is $\C$.

\end{enumerate}
\end{proof}






\begin{proposition}[Euler numbers of irreducible normal-crossings torus fibres]\hypertarget{tag-01AV}{}\label{cdp:st:A.4.6}
\StatementRecord{cdp:st:A.4.6}{proved}{01AV}{\SeriesRef{cdp:st:A.4.5}}

For $S$ in \ref{cdp:st:A.4.5},
\[
e(S)=\begin{cases}
2,&\widetilde S\text{ is rational},\\
0,&\widetilde S\text{ is elliptic ruled or Hopf}.
\end{cases}
\tag{A.18}\hypertarget{tag-0049}{}\label{cdp:eq:A.18}
\]
Thus an arbitrary compact torus fibration whose fibres are irreducible reduced normal-crossings surfaces or multiples of smooth surfaces of Euler number zero has
\[
c_3(Y)=e(Y)=2r,
\tag{A.19}\hypertarget{tag-004A}{}\label{cdp:eq:A.19}
\]
where $r$ is the number of rational-normalization fibres. For $b_1(Y)=b_2(Y)=0$, one has $r\in\{0,1\}$ and $b_3(Y)=2-2r$.
\end{proposition}


\begin{proof}

\begin{enumerate}[label=\arabic*.,leftmargin=*]
\item Write $D\subset\widetilde S$ for the conductor support and $N\subset S$ for the singular curve. Normalization is a homeomorphism off these sets, so additivity of the topological Euler characteristic gives $e(S)=e(\widetilde S)-e(D)+e(N)$. In the Hopf case $e(\widetilde S)=0$, $D$ is a disjoint union of two elliptic curves, and $N$ is one elliptic curve; all three Euler numbers vanish. The ruled case has the same conductor curves and $e(\widetilde S)=e(\PP^1)e(E)=0$.

\item In the rational case the normalization is a Hirzebruch surface blown up twice, so its Euler number is $4+2=6$. The conductor is a cycle of six rational curves, with six nodes, so $e(D)=12-6=6$. Its image consists of three rational curves with exactly two triple points; counting the three preimages of each triple point once gives $e(N)=6-2(3-1)=2$. Hence $e(S)=2$.

\item Away from finitely many critical values the proper holomorphic map is a differentiable fibre bundle, and a two-torus has Euler number zero. Additivity for a stratification by this open base and the exceptional points therefore gives $e(Y)=\sum_c e((Y_c)_{\rm red})$. Nilpotents do not change the underlying topological space. The assumed fibre classes and \eqref{cdp:eq:A.18} give \eqref{cdp:eq:A.19}. The equality with $c_3$ is the complex-oriented Gauss--Bonnet theorem. Poincar\'e duality for the connected oriented six-manifold gives $e(Y)=2-2b_1+2b_2-b_3$. Under the two Betti vanishings this is $2-b_3$, so $b_3=2-2r\ge0$ and $r$ is zero or one.

\end{enumerate}
\end{proof}






\begin{proposition}[Unitary-flat differentials on every compact reduced complex surface]\hypertarget{tag-01AW}{}\label{cdp:st:A.4.7}
\StatementRecord{cdp:st:A.4.7}{proved}{01AW}{the definitions and constructions stated in the proof}

Let $S$ be a compact reduced complex surface, let $\mathbb A$ be a rank-one unitary complex local system on its underlying topological space, and let $A=\mathbb A\otimes_{\C}\OO_S$ be its associated holomorphic line bundle. There is a canonical complex-linear injection
\[
 H^0(S,\widetilde\Omega_S^1\otimes A)
       \longrightarrow H^1(S,\mathbb A).
 \tag{A.20}\hypertarget{tag-004C}{}\label{cdp:eq:A.20}
\]
The map is constructed from continuous local primitives whose pullbacks to the normalization are holomorphic.
\end{proposition}


\begin{proof}

\begin{enumerate}[label=\arabic*.,leftmargin=*]
\item Let $\eta:\widetilde S\to S$ be normalization and choose a resolution $\pi:Z\to\widetilde S$ such that the exceptional curves have normal crossings; the components of $Z$ are compact smooth surfaces. Pullback of Kähler differentials along $h=\eta\circ\pi:Z\to S$ kills the torsion of $\Omega_S^1$, because the resulting form vanishes on the dense inverse image of $S_{\rm reg}$ and the target is locally free. Thus a section $\alpha$ on the left of \eqref{cdp:eq:A.20} gives a holomorphic one-form $\alpha_Z$ with values in the unitary-flat line $h^*A$. This map on sections is injective for the same dense-open reason.

\item Write $\nabla$ for the flat unitary connection. Holomorphicity gives $\nabla\alpha_Z$ type $(2,0)$. The Hermitian pairing identifies the tensor product of the flat line with its conjugate with the trivial local system. Stokes' theorem gives, on each connected component of $Z$,
\[
 \int_Z \nabla\alpha_Z\wedge\overline{\nabla\alpha_Z}
 =\int_Z d\bigl(\alpha_Z\wedge\overline{\nabla\alpha_Z}\bigr)=0.
\]
After multiplication by the fixed dimensional sign, the integrand is the nonnegative squared norm of a $(2,0)$-form times a positive volume form. Hence $\nabla\alpha_Z=0$.

\item Fix a point $x\in S$. Trivialize $\mathbb A$ on a sufficiently small connected neighbourhood $V$ of $x$; complex analytic spaces are locally contractible, so such a flat trivialization exists. The section $\alpha$ has a local lift to $\Omega_S^1\otimes A$, since the latter surjects onto the torsion-free quotient. Its pullback to each reduced fibre of $h$ over $x$ is zero: the morphism of that fibre to $S$ is constant, so the induced morphism on Kähler differentials is zero.

\item For each point $\widetilde x\in\eta^{-1}(x)$, a sufficiently small regular neighbourhood of the exceptional fibre $\pi^{-1}(\widetilde x)$ in the good resolution retracts onto that fibre. This neighbourhood can be constructed by plumbing the normal disc bundles of the exceptional curves at their transverse intersection points; it is a neighbourhood of the entire fibre. A sufficiently small inverse image of a neighbourhood of $\widetilde x$ lies inside it by properness. In the smooth case one instead uses a coordinate ball. Every period of $\alpha_Z$ on the regular neighbourhood is zero, because the form is closed and its restriction to the retracting fibre is zero. Integration from a fixed point therefore produces a single-valued holomorphic primitive: a smooth primitive has differential of type $(1,0)$, so its $\bar\partial$ is zero. Normality and $\pi_*\OO_Z=\OO_{\widetilde S}$ descend that primitive to a holomorphic function $b$ on a neighbourhood of $\widetilde x$. Choose its constant so that $b(\widetilde x)=0$. Perform this independently at every point over $x$ and shrink $V$ so that all its inverse image is covered by these neighbourhoods.

\item These primitives have equal values on every fibre of $\eta$ over $V$. To see this, form the finite analytic fibre product
$W=\widetilde S\times_S\widetilde S$ with projections $p_1,p_2$. On the normalization of every positive-dimensional irreducible component of $W_{\rm red}$, the differential of $p_1^*b-p_2^*b$ is zero on a dense open subset: both differentials are pullbacks of the same local lift of $\alpha$ from $S$. The isolated singularities of $\widetilde S$ can be omitted for this comparison; holomorphicity then extends the equality. Thus the difference is constant on each such component. The same conclusion is automatic on a zero-dimensional component. Shrink $V$ once more so that every component of the finite space $W|_V$ meets the fibre over $x$; components disjoint from that fibre can be removed by properness. Our normalization $b(\widetilde x)=0$ makes all these constants zero. Hence $p_1^*b=p_2^*b$ as continuous functions, proving equality on the fibres of $\eta$.

\item A finite proper surjection is a topological quotient map. The functions $b$ therefore descend to continuous functions on $V$, holomorphic after pullback to normalization, whose differential on the regular locus is $\alpha$. On overlaps, any two such primitives differ by a locally constant function: their difference is constant on each local normalized branch, and continuity identifies the constants whenever branches meet. With the flat transition functions of $\mathbb A$, these differences define a \v Cech class in $H^1(S,\mathbb A)$. Changing the chosen primitives changes it by a locally constant coboundary. This defines the canonical linear map \eqref{cdp:eq:A.20}.

\item If the resulting class is zero, adjust the local primitives by locally constant sections so that they glue. Their pullback is then a global holomorphic section of the unitary-flat line on each compact component of $Z$. Its squared norm is plurisubharmonic, so the maximum principle makes the section parallel; its covariant differential is zero. Consequently $\alpha_Z=0$, and injectivity of pullback gives $\alpha=0$ for every element of the kernel of \eqref{cdp:eq:A.20}.

\end{enumerate}
\end{proof}






\begin{proposition}[Cotangent vanishing for the transgression family]\hypertarget{tag-01AX}{}\label{cdp:st:A.4.8}
\StatementRecord{cdp:st:A.4.8}{proved}{01AX}{\SeriesRef{cdp:st:D.1.9}, \SeriesRef{cdp:st:D.1.8}, \SeriesRef{cdp:eq:A.22}, \SeriesRef{cdp:st:A.4.7}}

Let $g:Y\to\PP^1$ be a proper surjective holomorphic map with connected fibres from a smooth compact complex threefold with $b_1(Y)=b_2(Y)=0$, and suppose its general fibre is a complex two-torus. Let $L_r=\exp(2\pi\mathrm i r\beta)\in\Pic^0(Y)$ be the transgression family constructed in \ref{cdp:st:D.1.9}. For every real irrational $r$ and every reduced fibre support $S=(Y_c)_{\rm red}$,
\[
 H^1(S,\mathbb A_{r,S})=0,
 \qquad
 H^0(S,\widetilde\Omega_S^1\otimes L_r|_S)=0,
 \tag{A.21}\hypertarget{tag-004D}{}\label{cdp:eq:A.21}
\]
where $\mathbb A_{r,S}$ is the unitary local system whose associated holomorphic line is $L_r|_S$. The second group is the torsion-free cotangent group in \cite[Equation~(11), p.~692]{CDP20}.
\end{proposition}


\begin{proof}

\begin{enumerate}[label=\arabic*.,leftmargin=*]
\item Take the transgression section supplied by \ref{cdp:st:D.1.9}. On a sufficiently small disc $\Delta$ about $c$, put $U=g^{-1}(\Delta)$ and $U^*=U\setminus S$. The integral section $\gamma$ of $R^1g_*\Z$ has a lift $\gamma_U\in H^1(U,\Z)$ by Leray over the contractible disc. Its image in $H^1(U,\OO_U)$ equals $\beta|_U$, by the coherent Leray sequence. Thus the restriction of the holomorphic line $L_r$ to $U$ is the line associated to the flat character
\[
 \pi_1(U)\longrightarrow S^1,\qquad
 a\longmapsto\exp\bigl(2\pi\mathrm i r\langle\gamma_U,a\rangle\bigr).
 \tag{A.22}\hypertarget{tag-004E}{}\label{cdp:eq:A.22}
\]
Denote this local system by $\mathbb A_U$. \ref{cdp:st:D.1.8}--\ref{cdp:st:D.1.9} prove that $\gamma_U$ restricts nontrivially to every nearby smooth fibre. Irrationality of $r$ therefore makes \eqref{cdp:eq:A.22} nontrivial on such a fibre $F$.

\item A nontrivial rank-one local system on the real torus underlying $F$ has zero cohomology in every degree. Explicitly, use the standard cellular cochain complex of the product of four circles with coefficients in the character. It is the Koszul complex for the four scalars $\chi(e_j)-1$. At least one scalar is nonzero; contraction in that factor, divided by this scalar, is a contracting homotopy. Therefore $H^q(F,\mathbb A_U|_F)=0$ for every $q$. The map $U^*\to\Delta^*$ is a differentiable fibre bundle. Its local-system Leray spectral sequence has zero fibre-cohomology sheaves, hence
$H^1(U^*,\mathbb A_U|_{U^*})=0$.

\item Inclusion induces a surjection $\pi_1(U^*)\to\pi_1(U)$: a loop in the smooth real six-manifold $U$ can be perturbed off the divisor, whose real codimension is at least two, relative to a base point in $U^*$. The induced map on first cohomology with these local coefficients is injective. For an explicit verification, identify first cohomology with crossed homomorphisms modulo coboundaries. If a crossed homomorphism $a:G\to M$ becomes $a(q(h))=(\chi(q(h))-1)v$ after pullback by a surjection $q:H\to G$, surjectivity gives $a(g)=(\chi(g)-1)v$ for every $g\in G$. Thus its original class was already zero. Applied to the displayed fundamental-group map, this proves $H^1(U,\mathbb A_U)=0$.

\item The argument applies after every sufficiently small shrinking of $\Delta$, with the restrictions of the same local system. Taking the direct limit and using proper base change identifies the resulting stalk with $H^1(S,\mathbb A_U|_S)$; scheme-theoretic nilpotents do not affect this topological local system. The second vanishing in \eqref{cdp:eq:A.21} follows from \ref{cdp:st:A.4.7} applied to the restriction of \eqref{cdp:eq:A.22}. All singular values are treated by the same global parameter $r$, so the assertion is simultaneous.

\end{enumerate}
\end{proof}






\begin{proposition}[CDP20 Proposition 7.3 with a complete replacement proof]\hypertarget{tag-01AY}{}\label{cdp:st:A.4.9}
\StatementRecord{cdp:st:A.4.9}{proved}{01AY}{\SeriesRef{cdp:st:D.1.9}, \SeriesRef{cdp:st:C.3.1}, \SeriesRef{cdp:st:A.4.8}, \SeriesRef{cdp:st:D.1.8}, \SeriesRef{cdp:st:A.4.7}, \SeriesRef{cdp:st:A.4.2}, \SeriesRef{cdp:eq:A.24}; \cite[\S1 and the central-extension diagram in the proof of Lemma~1.1]{SS2002}; \cite[\S2.1.1 and Theorem~2.2]{Albanese2020}; \cite[\S2.1.2 and Proposition~2.3]{Albanese2020}}

Let $Y$ be a connected smooth compact complex threefold with $H^1(Y,\Z)=H^2(Y,\Z)=0$, algebraic dimension one and holomorphic algebraic reduction $g:Y\to C$. There is a dense open subset of $\Pic(Y)$, whose complement is a proper complex analytic subset, such that
\[
 H^0\bigl((Y_c)_{\rm red},
       \widetilde\Omega^1_{(Y_c)_{\rm red}}\otimes L|_{(Y_c)_{\rm red}}\bigr)=0
 \tag{A.23}\hypertarget{tag-004F}{}\label{cdp:eq:A.23}
\]
for every singular reduced fibre support and every $L$ in that subset. In fact the transgression twists $L_r$ of \ref{cdp:st:D.1.9} satisfy \eqref{cdp:eq:A.23} on every reduced fibre support, including smooth and reducible supports, for every real irrational $r$.
\end{proposition}


\begin{proof}

\begin{enumerate}[label={$\langle1\rangle$\arabic*.},leftmargin=*]
\item Choose one very general smooth fibre $F_0$. The classification \ref{cdp:st:C.3.1} places this fibre among complex tori, bielliptic surfaces, primary or secondary Kodaira surfaces, and primary or secondary Hopf surfaces or Inoue surfaces. We first prove that for each surface $F$ in this list and each nonzero $\gamma_F\in H^1(F,\Z)$ the local system
\[
 \mathbb A_{r,F}:\quad a\longmapsto
       \exp(2\pi\mathrm i r\langle\gamma_F,a\rangle)
 \quad(r\in\mathbb R\setminus\Q)
\]
satisfies $H^0(F,\mathbb A_{r,F})=H^1(F,\mathbb A_{r,F})=0$. Its character has infinite image, so the assertion in degree zero is immediate.

\item For a torus the Koszul argument in \ref{cdp:st:A.4.8} applies. A bielliptic surface has a finite étale torus cover. Pullback of the integral class is nonzero: transfer composed with pullback is multiplication by the degree. The character stays nontrivial after pullback, and transfer with local coefficients makes pullback on cohomology injective over $\C$. Thus degree-one vanishing descends from the torus cover.

\item A primary Kodaira surface is a principal elliptic bundle over an elliptic curve, as described in \cite[\S1 and the central-extension diagram in the proof of Lemma~1.1]{SS2002}. Its fundamental group has a central extension
\[
 1\longrightarrow\Z^2\longrightarrow\Gamma\longrightarrow\Z^2
       \longrightarrow1.
 \tag{A.24}\hypertarget{tag-004G}{}\label{cdp:eq:A.24}
\]
For any nontrivial one-dimensional character $\chi$ of $\Gamma$, if its restriction to the central subgroup is nontrivial, the latter's cohomology with these coefficients is zero by the two-variable Koszul contraction. The Hochschild--Serre spectral sequence then gives $H^1(\Gamma,\C_\chi)=0$. If its central restriction is trivial, $\chi$ descends to a nontrivial character of the quotient $\Z^2$. Centrality makes its action on the cohomology of the central subgroup equal to that scalar character, and the same Koszul contraction kills every term on the second page. Again $H^1(\Gamma,\C_\chi)=0$. First cohomology of a connected space with local coefficients is first cohomology of its fundamental group, so this proves the required assertion. A secondary Kodaira surface has a finite primary cover; the same transfer argument as above applies. Hopf surfaces have finite primary covers, and a primary Hopf surface is diffeomorphic to $S^1\times S^3$; the cover statement and topology are recorded in \cite[\S2.1.1 and Theorem~2.2]{Albanese2020}. The nontrivial circle character has an invertible cellular differential, giving the required vanishing and its finite-cover descent.

\item For completeness we also give the Inoue calculation, rather than assuming a generic-vanishing theorem for its cotangent bundle. The explicit constructions in \cite[\S2.1.2 and Proposition~2.3]{Albanese2020} identify every Inoue surface with a mapping torus $N_\varphi$ of a compact three-manifold $N$. In the $S_M$ cases $N$ is a three-torus and the eigenvalues of $\varphi^*$ on $H^1(N,\C)$ are $\alpha,\beta,\bar\beta$, with $\alpha>1$ and $|\beta|=\alpha^{-1/2}<1$, after reversing the circle orientation when necessary. In the $S^+$ and $S^-$ cases $N$ is the nontrivial circle bundle over a two-torus; the nonzero Euler class makes $H^1(N,\C)$ the pullback of $H^1(T^2,\C)$ by its Gysin sequence. The monodromy on it has eigenvalues $\alpha,\alpha^{-1}$ or $\alpha,-\alpha^{-1}$, respectively. These assertions follow as well directly from the conjugation relations displayed in that construction: on the abelian rank-three kernel the matrix is $M$, and on the quotient of the Heisenberg kernel by its centre the matrix is $N\in\mathrm{GL}_2(\Z)$ with precisely those eigenvalues.

\item None of these degree-one eigenvalues has modulus one. In particular there are no invariant degree-one classes on $N$. The ordinary mapping-torus sequence therefore identifies $H^1(F,\Z)$ with the copy of $\Z$ from the circle, since the invariant subgroup of the torsion-free group $H^1(N,\Z)$ is zero. Thus $\mathbb A_{r,F}$ is pulled back from the circle, with monodromy $\lambda\in S^1$ of infinite order. The local-coefficient Leray sequence for $N\to N_\varphi\to S^1$ gives
\[
\begin{split}
0\longrightarrow
 \coker(\lambda-1:H^0(N,\C)\to H^0(N,\C))
 \longrightarrow H^1(F,\mathbb A_{r,F})\\
 \longrightarrow
 \ker(\lambda\varphi^*-1:H^1(N,\C)\to H^1(N,\C))
 \longrightarrow0.
\end{split}
\]
The first term is zero because $\lambda\ne1$, and the last term is zero by the eigenvalue moduli. Using the reciprocal convention for the circle local system replaces $\lambda$ by $\lambda^{-1}$ and gives the same conclusion. This completes the degree-one calculation for every surface type in \ref{cdp:st:C.3.1}.

\item Let $F_1$ be any smooth fibre. The smooth locus in the connected base curve is the complement of finitely many points and is connected. A path from the value of $F_0$ to that of $F_1$, together with proper smooth local triviality, gives an actual diffeomorphism $d:F_0\to F_1$. For any nonzero $\gamma_1\in H^1(F_1,\Z)$, the class $d^*\gamma_1$ is nonzero and the character formula gives the actual local-system isomorphism
\[
 d^*\mathbb A_{r,F_1,\gamma_1}
       \simeq\mathbb A_{r,F_0,d^*\gamma_1}.
\]
The induced cohomology isomorphism identifies degrees zero and one on $F_1$ with the groups just proved zero on $F_0$. This proves both vanishings on every smooth fibre. No constancy of its complex surface type is needed for these maps.

\item Now repeat the neighbourhood argument of \ref{cdp:st:A.4.8} with these two fibre-cohomology vanishings. \ref{cdp:st:D.1.8}--\ref{cdp:st:D.1.9} supply the same nonzero integral class on every nearby smooth fibre and the unitary local system on every fibre neighbourhood. The only terms that can contribute to first cohomology of the punctured neighbourhood are its degree-zero and degree-one fibre cohomology, just proved zero. The fundamental-group surjection and proper base change therefore give $H^1((Y_c)_{\rm red},\mathbb A_{r,c})=0$ for every $c$. Applying \ref{cdp:st:A.4.7} proves \eqref{cdp:eq:A.23} for the single global twist $L_r$, simultaneously on all reduced supports.

\item Finally, the exponential sequence and the integral vanishings identify $\Pic(Y)$ with the complex vector space $H^1(Y,\OO_Y)$. Choose holomorphic Cech cocycles representing a basis; exponentiating their linear combinations gives a holomorphic family of line bundles on $Y\times H^1(Y,\OO_Y)$ whose fibre is the indicated line bundle. There are finitely many singular fibre supports. For each of them, tensor its fixed coherent torsion-free cotangent sheaf by this family and apply proper semicontinuity: the locus where its space of sections is nonzero is a closed complex analytic subset of the parameter space. Each is proper because it excludes $L_r$. Their finite union is proper and has empty interior, giving the asserted dense open set. This proves the statement of Proposition~7.3. The invalid genus-two deduction in its printed proof remains refuted by \ref{cdp:st:A.4.2}; the proof above replaces that argument through the explicit maps (A.20)--\eqref{cdp:eq:A.24}.

\end{enumerate}
\end{proof}






\section{Normalization, conductor genus and local torus exclusions}\hypertarget{tag-0341}{}\label{sec:II:05}


\begin{proposition}[Conormal sections with trivial normalization pullback]\hypertarget{tag-01AZ}{}\label{cdp:st:A.5.1}
\StatementRecord{cdp:st:A.5.1}{proved}{01AZ}{\SeriesRef{cdp:st:A.1.1}, \SeriesRef{cdp:st:A.4.1}}

Let $S$ be a connected reduced full fibre of a proper map from a complex manifold to a smooth curve. For every line bundle $A$ on $S$ with $\eta^*A\simeq\OO_{\widetilde S}$, the group $H^0(S,V\otimes A)$ contains the canonical image of $H^0(\widetilde S,\eta^*(N^*\otimes A))$. A local base coordinate and a choice of trivialization of $\eta^*A$ identify this subspace with $H^0(\widetilde S,\OO_{\widetilde S})$. In particular it is nonzero. This applies to the subcase of \cite[Proposition~7.3, pp.~693--694]{CDP20} in which all $\eta^*(L_t|_S)$ are trivial. It is compatible with the vanishing of $H^0(S,Q\otimes A)$ through the maps in Proposition~\ref{cdp:st:A.1.1}, as \ref{cdp:st:A.4.1} demonstrates.
\end{proposition}


\begin{proof}
\begin{enumerate}[label={$\langle1\rangle$\arabic*.},leftmargin=*]
\item Proposition~\ref{cdp:st:A.2.2} gives the inclusion in
\eqref{cdp:eq:A.3}. Tensoring it by the invertible sheaf $A$ and applying
the finite projection formula gives
\[
 \eta_*\eta^*(N^*\otimes A)\hookrightarrow K\otimes A
                         \hookrightarrow V\otimes A.
\]
Locally, for a defining equation $F$ of $S$, the first map sends
$b\otimes[F]$ to $b\,dF|_S$. The inclusion of the Jacobian ideal in the
conductor makes its coefficients holomorphic on $S$. Under $F'=uF$,
both the conormal frame and $dF|_S$ are multiplied by $u$, so these maps
glue. Taking global sections gives the claimed canonical injection.

\item A coordinate $t$ about the image point on the base defines the full
fibre by $F=t\circ f$. Since this full fibre is reduced, $[F]$ is a frame
of $N^*$. Its pullback, together with the chosen trivialization of
$\eta^*A$, identifies the source with
$H^0(\widetilde S,\OO_{\widetilde S})$. The nonempty compact normalization
has the nonzero constant function $1$ on each connected component.
This proves nonvanishing and applies to each $A=L_t|_S$ in the stated
published subcase.

\item The image lies in $H^0(S,K\otimes A)$, the kernel of
$H^0(S,V\otimes A)\to H^0(S,Q\otimes A)$ by
\eqref{cdp:A:filtration2}. Thus these nonzero sections map to zero in
the quotient group. Proposition~\ref{cdp:st:A.4.1} gives the stated
example with vanishing quotient group.
\end{enumerate}
\end{proof}



The line in the printed proof asserting $H^2(\widetilde S,\omega_{\widetilde S})=0$ also cannot hold for a connected compact normal Cohen--Macaulay surface with its dualizing sheaf: Serre duality identifies its dual with $H^0(\widetilde S,\OO)=\C$. Replacing $\omega_{\widetilde S}$ there by $\OO_{\widetilde S}$ matches the vanishing stated at the beginning of that proof, but does not settle its omitted genus-one case. We have established (11) for $W$ and the exact obstruction to deducing ambient vanishing from it. The full statement of Proposition~7.3 is proved below in \ref{cdp:st:A.4.9} by a different argument through unitary local systems. The following argument establishes Proposition~4.11 independently of the monodromy-dependent restrictions.




\begin{proposition}[Normal Gorenstein surfaces with vanishing first Chern class and Euler characteristic]\hypertarget{tag-01BA}{}\label{cdp:st:A.5.2}
\StatementRecord{cdp:st:A.5.2}{proved}{01BA}{\cite[Theorem~2.2]{Peters2021}}

A compact connected normal Gorenstein complex surface $S$ with $c_1(\omega_S)=0$ in $H^2(S,\Q)$ and $\chi(S,\OO_S)=0$ is smooth. Consequently every singular irreducible reduced divisor in a compact complex threefold $Y$ with $H^2(Y,\Q)=0$ is nonnormal. In particular the assertion of CDP20 Proposition~4.11 holds without its monodromy-dependent restrictions on the smooth fibre.
\end{proposition}


\begin{proof}

\begin{enumerate}[label={$\langle1\rangle$\arabic*.},leftmargin=*]
\item Let $\pi:Z\to S$ be the minimal resolution and set $\ell=h^0(S,R^1\pi_*\OO_Z)$. Normality gives $\pi_*\OO_Z=\OO_S$, and the fibres have dimension at most one, so Leray gives $\chi(Z,\OO_Z)=-\ell$. We first prove $\ell=0$. The classification table in \cite[Theorem~2.2]{Peters2021}, together with Noether's formula and invariance of $\chi(\OO)$ under blowing up, shows that a smooth compact surface of negative holomorphic Euler characteristic is a blown-up ruled surface over a smooth curve $B$ of genus $g\ge2$. Indeed the sums $c_1^2+c_2$ in every other row of that table are nonnegative. Thus, when $\ell>0$, there is a ruling $p:Z\to B$ and $\chi(Z,\OO_Z)=1-g$, giving $\ell=g-1$.

\item An exceptional irreducible curve $E$ of $\pi$ dominating $B$ would have $h^1(E,\OO_E)\ge g$: normalize $E$, apply Riemann--Hurwitz to its finite map to $B$, and use the normalization exact sequence on $E$. The exact sequence $0\to\mathcal I_E\to\OO_Z\to\OO_E\to0$, pushed forward by $\pi$, gives a surjection $R^1\pi_*\OO_Z\to R^1\pi_*\OO_E$ because $R^2\pi_*\mathcal I_E=0$. It follows that $\ell\ge h^1(E,\OO_E)\ge g$, a contradiction. Hence every exceptional curve is vertical for $p$. The ruling is constant on each connected fibre of $\pi$ and descends to a holomorphic map $p_S:S\to B$: local holomorphic coordinate functions descend by $\pi_*\OO_Z=\OO_S$, and properness supplies the required neighborhoods. The map $p$ is the composite of blowups and a $\PP^1$-bundle, so $R^1p_*\OO_Z=0$. The Leray sequence for $p=p_S\circ\pi$ then gives
\[
0=R^1p_*\OO_Z\longrightarrow (p_S)_*R^1\pi_*\OO_Z
 \longrightarrow R^2(p_S)_*\OO_S=0.
\]
A nonzero sheaf supported at finitely many points cannot have zero direct image, so again $\ell=0$, a contradiction. This proves rationality of all singularities.

\item We give the crepancy argument to specify the comparison of canonical bundles. Proper duality applied to $R\pi_*\OO_Z=\OO_S$ identifies $\pi_*\omega_Z$ with the invertible dualizing sheaf $\omega_S$. A local generator of $\omega_S$ therefore pulls back to a holomorphic two-form on $Z$, agreeing with the usual pullback off the exceptional set. Its zero divisor is effective and exceptional. These divisors glue to $D\ge0$ with $K_Z=\pi^*K_S+D$. Every exceptional irreducible curve has arithmetic genus zero by the same surjection to its first cohomology, hence is a smooth rational curve. Minimality of the resolution gives $E^2\le-2$, and adjunction gives $D\cdot E=K_Z\cdot E=-2-E^2\ge0$. The exceptional intersection matrix is negative definite. Thus a nonzero $D$ would satisfy both $D^2<0$ and $D^2=\sum_E d_E(D\cdot E)\ge0$. Therefore $D=0$.

\item Consequently $c_1(K_Z)=0$ in rational cohomology. It follows by adjunction that $Z$ contains no $(-1)$-curve, so $Z$ is a minimal compact surface; Noether's formula gives $c_2(Z)=0$. We exclude rational curves on each surviving row of the classification. A ruled surface is excluded by $K_Z\cdot F=-2$ on a ruling fibre. General-type surfaces have $K_Z^2>0$ and K3 and Enriques surfaces have positive $c_2$, so none occurs. A torus has no rational curve because a map from $\PP^1$ lifts to its vector-space universal cover and is constant; the same holds for a finite etale quotient of a torus. A primary Kodaira surface is an elliptic bundle over an elliptic curve, so a rational curve would map constantly to the base and then constantly to a fibre; finite etale quotients have the same property. A class VII surface here has $b_2=c_2=0$, and the Hopf--Inoue classification gives no rational curves. Equivalently one can use their standard universal covers: $\C^2\setminus\{0\}$ for a primary Hopf cover and $\mathbb H\times\C$ for an Inoue surface. A map from $\PP^1$ lifts and has constant coordinate functions.

\item It remains to check the properly elliptic row. For a minimal elliptic surface, the Kodaira fibre list expresses $c_2$ as a sum of nonnegative fibre Euler numbers, with zero contributions exactly at multiples of smooth elliptic curves. Since $c_2(Z)=0$, all fibre reductions are smooth elliptic curves. The canonical bundle formula on a Kähler elliptic surface with $\chi(\OO_Z)=0$ has rational first Chern class
\[
 c_1(K_Z)=\left(2g(B)-2+\sum_i(1-m_i^{-1})\right)[F].
\]
In Kodaira dimension one the coefficient is positive, and the Kähler class has positive integral over $F$, contradicting $c_1(K_Z)=0$. In the non-Kähler case, a horizontal rational curve would be a multisection $C$. For $D=C+nF$ with $n$ sufficiently large, $D^2>0$ and $(K_Z-mD)\cdot F<0$ for $m>0$. Riemann--Roch therefore gives quadratic growth of $h^0(Z,\OO_Z(mD))$, making $Z$ Moishezon and hence projective as a smooth compact surface. This contradicts the non-Kähler case. Vertical rational curves have already been excluded by the fibre list. Thus $Z$ has no rational curves in any case. The exceptional curves of $\pi$ were rational, so there are none and $S$ is smooth.

\item Finally take an irreducible reduced divisor $S\subset Y$ with $H^2(Y,\Q)=0$. It is a Cartier divisor in the smooth manifold $Y$, hence Gorenstein, and adjunction gives $\omega_S=(K_Y\otimes\OO_Y(S))|_S$, with rational first Chern class zero. The divisor sequence and Hirzebruch--Riemann--Roch give
$\chi(\OO_S)=\chi(\OO_Y)-\chi(\OO_Y(-S))=0$, since $c_1(\OO_Y(S))=0$ rationally. A normal $S$ is therefore smooth by the result just proved.

\end{enumerate}
\end{proof}






\begin{proposition}[The effective conductor divisor on a minimal resolution]\hypertarget{tag-01BD}{}\label{cdp:st:A.5.3}
\StatementRecord{cdp:st:A.5.3}{proved}{01BD}{\SeriesRef{cdp:st:A.2.1}; \cite[Lemma~1.1 and Corollary~1.4]{Furushima1993}}

Let $S$ be a compact reduced hypersurface in a smooth threefold, let $\eta:\coprod_i T_i\to S$ be its normalization, and let $\pi_i:Z_i\to T_i$ be the minimal resolutions. Put $h_i=\eta\pi_i$. There are canonical effective integral divisors $D_i$ on $Z_i$ with
\[
 \omega_{Z_i}(D_i)\simeq h_i^*\omega_S.
 \tag{A.25}\hypertarget{tag-004H}{}\label{cdp:eq:A.25}
\]
The nonexceptional part of $D_i$ is exactly the divisorial conductor, with its scheme-theoretic multiplicities. When $S\subset Y$ and $H^2(Y,\Q)=0$, every $K_{Z_i}+D_i$ has rational first Chern class zero. A ruled $Z_i\to B$ with $g(B)\ge2$ then has
\[
 D_i=2C+V,
 \tag{A.26}\hypertarget{tag-004I}{}\label{cdp:eq:A.26}
\]
where $C$ is a section and $V$ is vertical. In the case where $C$ is $\pi_i$-exceptional, $\eta^*(L|_S)|_{T_i}$ is trivial for every $L\in\Pic(Y)$.
\end{proposition}


\begin{proof}

\begin{enumerate}[label={$\langle1\rangle$\arabic*.},leftmargin=*]
\item Finite duality gives $\omega_{T_i}=\mathfrak c_i\otimes\eta^*\omega_S$ as rank-one reflexive sheaves, where $\mathfrak c_i$ is the conductor viewed as an ideal on $T_i$. Locally this is the equality
$\Hom_R(B,R)=\{q:qB\subset R\}=\mathfrak c$ for a reduced Gorenstein ring $R$ and its finite normalization $B$, tensored with the invertible dualizing module. The equality uses the meromorphic multiplication identification as in \ref{cdp:st:A.2.1}. The conductor is a maximal Cohen--Macaulay module because it is this dual of the maximal Cohen--Macaulay normalization over the Gorenstein ring. On a normal surface a rank-one maximal Cohen--Macaulay module is reflexive, so the displayed equality is determined by its codimension-one orders. Those orders define an effective conductor divisor $D_i^0$ on $T_i$.

\item Off the exceptional set the canonical identification pulls back to $\omega_{Z_i}(D_i^{\rm str})\simeq h_i^*\omega_S$, where $D_i^{\rm str}$ denotes strict transform. The meromorphic extension of this line-bundle identification gives integral exceptional coefficients $a_j$ with
$K_{Z_i}+D_i^{\rm str}+\sum_j a_jE_j=h_i^*K_S$.
We verify their signs. For each exceptional irreducible curve $E_j$, minimality excludes a smooth rational $(-1)$-curve. Adjunction therefore gives
$K_{Z_i}\cdot E_j=2p_a(E_j)-2-E_j^2\ge0$;
also $D_i^{\rm str}\cdot E_j\ge0$. The exceptional intersection matrix $M=(E_j\cdot E_k)$ is negative definite with nonnegative off-diagonal entries, and
\[
 M a=-v,\qquad v_j=(K_{Z_i}+D_i^{\rm str})\cdot E_j\ge0.
\]
To see that $a\ge0$ without an unproved sign convention, write $a=a_+-a_-$ with disjoint nonnegative coordinate supports. If $a_-\ne0$, then
$a_-^tMa=a_-^tMa_+-a_-^tMa_->0$,
whereas $Ma=-v$ gives $a_-^tMa\le0$. This contradiction proves $a_j\ge0$. Thus $D_i=D_i^{\rm str}+\sum_j a_jE_j$ proves \eqref{cdp:eq:A.25}. The conductor duality and effective exceptional comparison also appear, for projective surfaces with trivial dualizing sheaf, in \cite[Lemma~1.1 and Corollary~1.4]{Furushima1993}; the argument here supplies the maps and does not use that paper's additional K3-type hypothesis $h^1(\OO_S)=0$.

\item For a divisor in $Y$, adjunction gives $\omega_S=(K_Y\otimes\OO_Y(S))|_S$, proving the numerical assertion. On a general ruling fibre $f$, $K_{Z_i}\cdot f=-2$, so $D_i\cdot f=2$. Every horizontal irreducible curve has normalization of genus at least $g(B)\ge2$. Such a curve $C_0$ cannot occur in $D_i$ with coefficient one: otherwise
\[
 0=(K_{Z_i}+D_i)\cdot C_0
   =2p_a(C_0)-2+(D_i-C_0)\cdot C_0>0.
\]
Every horizontal coefficient is consequently at least two. The degree-two intersection with $f$ forces exactly one horizontal component, coefficient two and degree one over $B$, giving \eqref{cdp:eq:A.26}. The degree-one finite morphism $C\to B$ is an isomorphism because $B$ is normal; thus $C$ is a smooth section.

\item Finally, integral $H^2$ of a smooth ruled surface and of its blowups is torsion-free, and its topologically trivial line bundles are precisely the pullbacks of $\Pic^0(B)$. These facts follow from the projective-bundle and blowup cohomology formulas and from $p_*\OO_{Z_i}=\OO_B$, $R^1p_*\OO_{Z_i}=0$. Restriction to the section $C$ is inverse to this Picard pullback. The line $h_i^*L$ has torsion, hence zero, integral first Chern class. When $C$ is exceptional, its restriction to $C$ is trivial because $h_i$ maps $C$ to a point. Therefore $h_i^*L$ is trivial on $Z_i$. Normality gives $(\pi_i)_*\OO_{Z_i}=\OO_{T_i}$, and the projection formula descends this triviality to $T_i$. This proves the last assertion.

\end{enumerate}
\end{proof}






\begin{proposition}[The remaining normalization class for a persistent ambient obstruction]\hypertarget{tag-01BE}{}\label{cdp:st:A.5.4}
\StatementRecord{cdp:st:A.5.4}{proved}{01BE}{\SeriesRef{cdp:st:A.5.3}, \SeriesRef{cdp:st:A.4.9}, \SeriesRef{cdp:st:A.1.1}; \cite[Lemma~2.3(1)--(2)]{Teleman2006}; \cite[Lemma~2.3(2)]{Teleman2006}}

Let $Y$ satisfy the CDP hypotheses, and let $S$ be a connected nonnormal reduced fibre support for which
\[
 H^0(S,\Omega_Y^1|_S\otimes L|_S)\ne0
 \quad\hbox{for every }L\in\Pic(Y).
 \tag{A.27}\hypertarget{tag-004J}{}\label{cdp:eq:A.27}
\]
There is a normalization component $T_i$ whose minimal resolution $Z_i$ has Kodaira dimension $-\infty$. Every line from $\Pic(Y)$ pulls back to a unitary-flat line on that $Z_i$. A rational or class VII $Z_i$ has trivial Picard pullback from $Y$, and then
\[
 H^2\!\left(T_i,\omega_{T_i}\otimes
             \eta^*((K_Y^{-1}\otimes L)|_S)|_{T_i}\right)\simeq\C
 \quad\hbox{for every }L\in\Pic(Y).
 \tag{A.28}\hypertarget{tag-004K}{}\label{cdp:eq:A.28}
\]
The same conclusion holds for the exceptional-section case of \ref{cdp:st:A.5.3}. Thus the only normalization type left by this argument, when \eqref{cdp:eq:A.28} has not been obtained, is an irrational ruled surface with nonconstant Picard pullback; in base genus at least two its divisor (A.26) has its horizontal section in the strict conductor.
\end{proposition}


\begin{proof}

\begin{enumerate}[label={$\langle1\rangle$\arabic*.},leftmargin=*]
\item \ref{cdp:st:A.4.9} gives $H^0(S,\widetilde\Omega_S^1\otimes L|_S)=0$ on a dense open subset of $\Pic(Y)$. By \ref{cdp:st:A.1.1} and \eqref{cdp:eq:A.27}, $H^0(S,K\otimes L|_S)$ is nonzero there. Its nonvanishing locus is closed analytic in the holomorphic Picard family, hence is all of $\Pic(Y)$.

\item Pull $K$ to each $Z_i$, quotient the resulting sheaf by its torsion, and take its double dual. The result is a fixed holomorphic line bundle $M_i$, since its generic rank on each component is one and $Z_i$ is a smooth surface. The induced maps give an injection
\[
 H^0(S,K\otimes L|_S)\hookrightarrow
       \bigoplus_i H^0(Z_i,M_i\otimes h_i^*L).
 \tag{A.29}\hypertarget{tag-004L}{}\label{cdp:eq:A.29}
\]
Indeed a section in the kernel vanishes generically on every component, and $K\subset\Omega_Y^1|_S$ is torsion-free over the reduced space. Each summand's nonvanishing locus is closed analytic in the connected complex vector space $\Pic(Y)=H^1(Y,\OO_Y)$. A finite union of proper analytic subsets cannot cover that space. Consequently one component $i$ has
$H^0(Z_i,M_i\otimes h_i^*L)\ne0$ for every $L$.

\item Fix a Gauduchon metric on $Z_i$, with fundamental form $\omega$ satisfying $\partial\bar\partial\omega=0$. The degree of a line bundle is the integral of its real Chern curvature against $\omega$; it is independent of the chosen Hermitian metric and additive under tensor product. A nonzero section has an effective zero divisor, so its line has nonnegative degree. For any $G\in\Pic(Y)$ the preceding nonvanishing applied to $G^{\otimes n}$ therefore gives
\[
 \deg_\omega(M_i)+n\deg_\omega(h_i^*G)\ge0
 \quad\hbox{for every }n\in\Z.
\]
It follows that $\deg_\omega(h_i^*G)=0$ for every $G$ and every Gauduchon metric.

\item We record the precise flatness implication. For a topologically trivial line, the real Chern curvature $F$ is exact. If $\int F\wedge\omega=0$, the Gauduchon scalar elliptic equation supplies a real $u$ with $\Lambda_\omega(F-i\partial\bar\partial u)=0$. The remaining form is closed, exact, real and primitive of type $(1,1)$, hence anti-selfdual. Its square integrates both to zero by exactness and to minus its squared norm. It is therefore zero. Changing the line metric by $e^{-u}$ gives zero Chern curvature, hence a flat unitary connection. This is exactly the curvature argument of \cite[Lemma~2.3(1)--(2)]{Teleman2006}. Since $H^2(Y,\Z)=0$, it applies to every $h_i^*G$.

\item The divisor $D_i$ of \ref{cdp:st:A.5.3} is nonzero. For an irreducible nonnormal support this follows from its nonempty conductor; for a connected reducible support every component meets another component, again giving a nonempty conductor. If $Z_i$ had nonnegative Kodaira dimension, some positive pluricanonical power would have a nonzero section, making $\deg_\omega(K_{Z_i})\ge0$. But $D_i$ is a nonzero effective divisor and therefore has strictly positive degree. Equation (A.25), with $\omega_S=(K_Y\otimes\OO_Y(S))|_S$, would give
\[
 0=\deg_\omega h_i^*(K_Y\otimes\OO_Y(S))
   =\deg_\omega(K_{Z_i})+\deg_\omega(D_i)>0,
\]
a contradiction. Thus $\kappa(Z_i)=-\infty$, whose classification consists of rational, irrational ruled, and class VII surfaces.

\item For a rational $Z_i$, topologically trivial lines are trivial. For a class VII $Z_i$, $h^1(\OO_{Z_i})=1$ and $b_1(Z_i)=1$. With our exponential convention $u\mapsto\exp(2\pi\mathrm i u)$, the unitary-flat subgroup of $\Pic^0(Z_i)$ has real tangent space equal to the injective image of $H^1(Z_i,\mathbb R)$ in $H^1(Z_i,\OO_{Z_i})$. This is \cite[Lemma~2.3(2)]{Teleman2006} after multiplication by the fixed scalar $\mathrm i$, or follows from the pluriharmonic maximum principle. That real tangent space has dimension one. The holomorphic Picard pullback has complex-linear differential and, by the proved flatness, takes values in this real one-dimensional subspace. Its differential is zero, so the homomorphism is constant and equals zero at the trivial line. Every $h_i^*L$ is therefore trivial. Projection formula and $(\pi_i)_*\OO_{Z_i}=\OO_{T_i}$ again descend triviality to $T_i$. The exceptional-section case has the same conclusion by \ref{cdp:st:A.5.3}.

\item In each of these cases, the line multiplying $\omega_{T_i}$ in \eqref{cdp:eq:A.28} is trivial, since it too comes from $Y$. A normal complex surface is Cohen--Macaulay, so its dualizing sheaf and Serre duality give $H^2(T_i,\omega_{T_i})=H^0(T_i,\OO_{T_i})^\vee=\C$. This proves \eqref{cdp:eq:A.28} and the stated exact reduction of the remaining normalization problem.

\end{enumerate}
\end{proof}






\begin{proposition}[Uniform negative-degree vanishing for a fixed vector bundle]\hypertarget{tag-01BG}{}\label{cdp:st:A.5.6}
\StatementRecord{cdp:st:A.5.6}{proved}{01BG}{\SeriesRef{cdp:st:A.5.4}; \cite[Lemma~2.3]{Teleman2006}}

Let $Z$ be a compact smooth complex surface, $\omega$ a Gauduchon metric, and $E$ a fixed holomorphic vector bundle. There is a real constant $b(E,\omega)$ such that
\[
 \deg_\omega P<b(E,\omega)
 \quad\Longrightarrow\quad H^0(Z,E\otimes P)=0
 \tag{A.30}\hypertarget{tag-004N}{}\label{cdp:eq:A.30}
\]
for every holomorphic line bundle $P$. Thus persistence of sections of $E\otimes h^*L$ for every $L\in\Pic(Y)$ forces $\deg_\omega h^*L=0$ for every $L$, for any holomorphic map $h:Z\to Y$. Persistence on one Picard component already forces degree zero on the pullback of $\Pic^0(Y)$. Under $H^2(Y,\Q)=0$, it forces degree zero and unitary flatness for every ambient Picard pullback. On a class VII surface $Z$ with $H^2(Y,\Z)=0$, the entire holomorphic Picard pullback is the zero map.
\end{proposition}


\begin{proof}

\begin{enumerate}[label={$\langle1\rangle$\arabic*.},leftmargin=*]
\item Choose a Hermitian metric on $E$ and let $C_E$ bound the largest eigenvalue of the Hermitian curvature endomorphism $\mathrm i\Lambda_\omega F_E$ on the compact surface. Write $V_\omega=\int_Z\omega^2/2$ and normalize degree by
$\deg_\omega P=(2\pi)^{-1}\int_Z\mathrm iF_P\wedge\omega$.
The scalar Gauduchon elliptic equation changes a Hermitian line metric so that its contracted curvature is the constant
\[
 \mathrm i\Lambda_\omega F_P
       =c_P:=2\pi\deg_\omega(P)/V_\omega.
\]
Indeed the difference from this constant has integral zero, and the image of $u\mapsto\mathrm i\Lambda_\omega\partial\bar\partial u$ is precisely the mean-zero functions for a Gauduchon metric; this is the scalar equation used in \cite[Lemma~2.3]{Teleman2006}. The Hermitian line metric may therefore be rescaled by the corresponding $e^{-u}$.

\item For a holomorphic section $s$ of $E\otimes P$, the Chern-connection identity is
\[
 \mathrm i\Lambda_\omega\partial\bar\partial|s|^2
 =|\nabla^{1,0}s|^2
  -\langle\mathrm i\Lambda_\omega F_{E\otimes P}s,s\rangle.
\]
It follows directly in a holomorphic frame by differentiating the Hermitian norm; thus no Kähler identity is being assumed. At a positive maximum of $|s|^2$, the left side is nonpositive, whereas the right side is at least
$-(C_E+c_P)|s|^2$.
If $c_P<-C_E$, this is strictly positive, a contradiction. We may take $b(E,\omega)=-C_EV_\omega/(2\pi)$, proving \eqref{cdp:eq:A.30}.

\item Apply this bound to $P=h^*(L^{\otimes n})$ for every positive and negative integer $n$. Nonzero sections for all $n$ are incompatible with a nonzero degree, proving the persistence assertion. For persistence on a single component, absorb one representative line into the fixed bundle $E$ and apply the same argument to all powers of every line in $\Pic^0(Y)$. If $H^2(Y,\Q)=0$, a positive power of every ambient line lies in $\Pic^0(Y)$, so additivity of degree gives degree zero for every pullback. Its real first Chern class is zero, even if its integral class is torsion. The Chern curvature is therefore exact, and the exact-curvature argument in \ref{cdp:st:A.5.4} again produces a flat unitary connection. On a class VII surface the differential of the holomorphic Picard pullback then lies in a real one-dimensional unitary tangent space. Complex linearity makes that differential zero. The exponential sequence makes $\Pic(Y)$ a connected quotient of $H^1(Y,\OO_Y)$ when $H^2(Y,\Z)=0$, so under that integral hypothesis the entire map is identically zero.

\end{enumerate}
\end{proof}






\begin{proposition}[The ruled obstruction with a reduced conductor image]\hypertarget{tag-01BJ}{}\label{cdp:st:A.5.9}
\StatementRecord{cdp:st:A.5.9}{proved}{01BJ}{\SeriesRef{cdp:st:A.5.3}, \SeriesRef{cdp:st:A.1.1}, \SeriesRef{cdp:st:A.5.8}, \SeriesRef{cdp:st:A.5.2}, \SeriesRef{cdp:eq:A.RC.3}, \SeriesRef{cdp:eq:A.RC.4}, \SeriesRef{cdp:st:C2d}, \SeriesRef{cdp:eq:A.RC.5}, \SeriesRef{cdp:st:A.2.4}}

Let $S\subset Y$ be an irreducible nonnormal reduced Cartier
divisor in a compact smooth complex threefold with
$H^2(Y,\Z)=0$ and $\chi(\OO_S)=0$.
Let $\eta:T\to S$ be its normalization and
$\pi:Z\to T$ its minimal resolution, and put $h=\eta\pi$.
Suppose that $Z$ is irrational ruled, with ruling $p:Z\to B$,
that every exceptional curve of $\pi$ is vertical, and that
the effective divisor of \ref{cdp:st:A.5.3} has the form
\[
 K_Z+D_Z=h^*\omega_S,\qquad D_Z=2C+V,                 \tag{A.RC.1}\hypertarget{tag-006U}{}\label{cdp:eq:A.RC.1}
\]
where $C$ is a section in the strict conductor and $V$ is
vertical. Suppose also that the conductor image scheme
$\Gamma=\operatorname{Spec}_S(\OO_S/\mathfrak c)$ is reduced,
and that $H^0(S,K\otimes L|_S)\ne0$ for every $L\in\Pic(Y)$,
with $K$ as in \ref{cdp:st:A.1.1}.
Then every ambient line bundle pulls back trivially to $T$.
In particular
\[
 H^2\bigl(T,\omega_T\otimes
       \eta^*((K_Y^{-1}\otimes L)|_S)\bigr)\simeq\C
       \quad\hbox{for every }L\in\Pic(Y).            \tag{A.RC.2}\hypertarget{tag-006V}{}\label{cdp:eq:A.RC.2}
\]
The proof forces $B$ and $\Gamma$ to be smooth elliptic curves,
$Z=T$ to be minimal ruled, and $D_Z=2C$ with $C^2=0$.
Thus no local-complete-intersection assumption on the original
conductor image curve is needed in this reduced-scheme case.
\end{proposition}


\begin{proof}

\begin{enumerate}[label={$\langle1\rangle$\arabic*.},leftmargin=*]
\item Write $\mathcal B=\eta_*\OO_T$ and
$\mathcal Q=\mathcal B/\OO_S$. We first specify conductor
duality, including its twists. Since $S$ is Gorenstein of
dimension two and $\mathcal B$ is maximal Cohen--Macaulay,
its dual is $\mathfrak c=\mathcal Hom_S(\mathcal B,\OO_S)$,
and maximal Cohen--Macaulay biduality gives
$\mathcal Hom_S(\mathfrak c,\OO_S)=\mathcal B$.
The quotient $\OO_\Gamma$ is Cohen--Macaulay of dimension one:
this follows from the depth lemma in
$0\to\mathfrak c\to\OO_S\to\OO_\Gamma\to0$.
Applying $\mathcal Hom_S(-,\OO_S)$ to that sequence, its
initial terms identify the actual normalization inclusion and
its cokernel as
\[
 0\longrightarrow\OO_S\longrightarrow\mathcal B
 \longrightarrow
 \mathcal Ext^1_S(\OO_\Gamma,\OO_S)\longrightarrow0,
 \qquad
 \mathcal Q\simeq\omega_\Gamma\otimes\omega_S^{-1}|_\Gamma.
                                                               \tag{A.RC.3}\hypertarget{tag-006W}{}\label{cdp:eq:A.RC.3}
\]
Adjunction makes $\omega_S$ the restriction of
$K_Y\otimes\OO_Y(S)$. Its degree is zero on every compact
curve, because $H^2(Y,\Z)=0$. Curve Riemann--Roch and Serre
duality therefore give
$\chi(\mathcal Q)=\chi(\omega_\Gamma)=-\chi(\OO_\Gamma)$.
On the other hand, \ref{cdp:st:A.5.8} gives rational singularities of $T$
and $\chi(\OO_T)=1-g(B)$. The normalization sequence and
$\chi(\OO_S)=0$ now give the exact numerical identity
\[
                   \chi(\OO_\Gamma)=g(B)-1.          \tag{A.RC.4}\hypertarget{tag-006X}{}\label{cdp:eq:A.RC.4}
\]

\item We next prove that $\Gamma$ is connected. Contract ruling
$(-1)$-curves from $Z$ to a minimal ruled bundle $Z_0\to B$.
Under a single contraction $\sigma:Z_1\to Z_2$ with exceptional
curve $E$, let $D_2=\sigma_*D_1$, let $a$ be the coefficient
of $E$ in $D_1$, and let $m$ be the multiplicity of $D_2$ at
the contraction point. The canonical divisor formula gives
\[
 K_{Z_1}+D_1=\sigma^*(K_{Z_2}+D_2)+(1+a-m)E.
                                                               \tag{A.RC.5}\hypertarget{tag-006Y}{}\label{cdp:eq:A.RC.5}
\]
Numerical triviality and intersection with $E$ force $a=m-1$.
The same numerical triviality descends to $Z_2$. On $Z_0$
the pushed divisor is $2C_0+V_0$, with $C_0$ a section and
$V_0$ a nonnegative sum of fibres, so its support is connected.
Reverse the contractions. A centre outside the existing support
would produce coefficient $-1$ on its exceptional curve; that
coefficient persists on strict transforms and contradicts the
effectiveness of the final $D_Z$. A centre of multiplicity one
is a smooth point of one coefficient-one component and its
exceptional coefficient is zero, so it does not disconnect the
support. At a centre of multiplicity at least two the exceptional
curve has positive coefficient and connects all strict-transform
branches through that point. Induction proves that
$\operatorname{Supp}D_Z$ is connected.

\item Every positive exceptional component of $D_Z$ maps into
$\Gamma$. Indeed away from $\Gamma$ the surface $T$ equals
the Gorenstein surface $S$ and has rational singularities by
\ref{cdp:st:A.5.8}. Such a singularity is crepant on its minimal resolution:
a generator of its dualizing sheaf pulls back holomorphically
by rational duality; its effective exceptional zero divisor
has nonnegative intersection with each exceptional curve by
adjunction and minimality, so negative definiteness makes it
zero, as in \ref{cdp:st:A.5.2}. Thus \eqref{cdp:eq:A.RC.1} has no positive exceptional
coefficient there. The strict conductor maps onto the whole
conductor support. It follows that
$h(\operatorname{Supp}D_Z)=\operatorname{Supp}\Gamma$,
which proves connectedness.

\item Let $D\subset\Gamma$ be the image of $C$. Localize \eqref{cdp:eq:A.RC.3}
at the generic point of $D$. Because $\Gamma$ is reduced,
$\OO_\Gamma$ there is the function field $k(D)$, and its
dualizing module has dimension one over that field. The
elementary conductor quotient sequence is
\[
 0\longrightarrow\OO_\Gamma\longrightarrow
          \mathcal B/\mathfrak c
   \longrightarrow\mathcal Q\longrightarrow0.
\]
It follows that $\mathcal B/\mathfrak c$ has dimension two
over $k(D)$. The conductor coefficient two on $C$ contributes
$2\deg(C\to D)$ to this dimension. Therefore this degree is
one and there is no other normalization branch over $D$.
The generic transverse ring has conductor $t^2\C\{t\}$ and
is exactly $\C+t^2\C\{t\}=\C\{t^2,t^3\}$, the cusp ring.
Thus $D^\nu\simeq C\simeq B$. This deduction also rules out
any additional vertical branch over the same image curve.

\item A connected compact reduced curve has only constant global
holomorphic functions. From \eqref{cdp:eq:A.RC.4} we therefore obtain
$h^1(\Gamma,\OO_\Gamma)=2-g(B)$.
Restriction to the reduced component $D$ gives a surjection
on $H^1$, because its kernel is a coherent sheaf on a curve
and has no $H^2$. Normalization of $D$ gives
$h^1(D,\OO_D)\ge g(B)$. Hence
\[
             2-g(B)\ge h^1(D,\OO_D)\ge g(B).         \tag{A.RC.6}\hypertarget{tag-006Z}{}\label{cdp:eq:A.RC.6}
\]
As $g(B)\ge1$, all quantities equal one. The normalization
quotient of $D$ has length zero, so $D$ itself is smooth
elliptic. Statement \ref{cdp:st:C2d} applies to this actual smooth curve
in $Y$ and makes the transversal cusp constant along it.
No other singular curve meets $D$. Connectedness of $\Gamma$
therefore gives $\Gamma=D$. Its normalization inverse image
is the single smooth curve $C$, the normalization is smooth
near it, and \ref{cdp:st:C2d} gives $C^2=0$.

\item Adjunction in \eqref{cdp:eq:A.RC.1} yields $V\cdot C=0$. A ruling
$(-1)$-curve outside $V$ cannot meet $C$, since its
intersection with $K_Z+2C+V$ would be at least $-1+2$.
A component of $V$ cannot meet $C$ either. Thus all
contractions to $Z_0$ avoid $C$ and preserve its square zero.
On the minimal bundle, $K_{Z_0}+2C_0$ is numerically trivial,
so its nonnegative vertical divisor $V_0$ is zero.
The first reverse blowup would be away from $C_0$ and have
coefficient $-1$ by \eqref{cdp:eq:A.RC.5}. There are no such blowups.
Thus $Z$ is minimal ruled and $D_Z=2C$.
An exceptional curve for $Z\to T$ would have negative square;
adjunction with $K_Z\equiv-2C$ would make it rational and
therefore vertical. A minimal ruled bundle has only fibres
of square zero, so no such exceptional curve exists and $Z=T$.

\item The constant-cusp computation \ref{cdp:st:A.2.4} now gives the injection
$K\hookrightarrow\eta_*(h^*\OO_Y(-S)\otimes\OO_Z(C))$.
Every ambient line is topologically trivial on $Z$ and is
the pullback of a degree-zero line on the elliptic curve $B$.
The section exact sequence, together with $R^1p_*\OO_Z=0$,
gives
\[
 0\longrightarrow\OO_B\longrightarrow p_*\OO_Z(C)
 \longrightarrow N_{C/Z}\longrightarrow0,
 \qquad\deg N_{C/Z}=0.                              \tag{A.RC.7}\hypertarget{tag-0070}{}\label{cdp:eq:A.RC.7}
\]
After a degree-zero twist $P$, both end terms have zero
sections unless $P=\OO_B$ or $P=N_{C/Z}^{-1}$.
The image of $\Pic(Y)$ in $\Pic^0(B)$ is connected:
$H^2(Y,\Z)=0$ makes $\Pic(Y)$ the connected image of
$H^1(Y,\OO_Y)$ under the exponential map.
All-twist persistence of $K$, and translation by the ambient
line $\OO_Y(-S)$, put this entire image in those at most two
points. It is therefore the identity. This proves triviality
of every ambient pullback. Serre duality on $T$ proves
\eqref{cdp:eq:A.RC.2}.

\end{enumerate}
\end{proof}






\begin{proposition}[The genus of a horizontal conductor image from a local torus group]\hypertarget{tag-01B1}{}\label{cdp:st:A.5.11}
\StatementRecord{cdp:st:A.5.11}{proved}{01B1}{\SeriesRef{cdp:st:A.5.8}, \SeriesRef{cdp:eq:A.Pi.1}, \SeriesRef{cdp:eq:A.Pi.2}, \SeriesRef{cdp:eq:A.Pi.4}; \cite[\S3, Theorem~2, p.~464]{Lojasiewicz1964}}

Let $f:Y\to C_0$ be a proper holomorphic map from a smooth
compact complex threefold to a smooth curve, with connected
fibres and complex two-torus general fibre. Let $S$ be an
irreducible reduced fibre support, $\eta:T\to S$ its
normalization, and $\pi:Z\to T$ a resolution which is
irrational ruled, with ruling $p:Z\to B$.
Suppose every exceptional curve of $\pi$ is vertical, and
the strict conductor has just one horizontal irreducible
component $C$, which is a section of $p$.
Let $D=\eta\pi(C)$ and let $\Sigma=D^\nu$ be its normalization.
Then $g(\Sigma)\le1$.
In particular, when $C\to D$ has degree one, $B$ is elliptic.
This applies to the horizontal-cusp case of $D_Z=2C+V$;
it allows arbitrary multiplicity of the full fibre and
arbitrary singularities or nonreduced conductor image on
the finitely many vertical curves.
\end{proposition}


\begin{proof}

\begin{enumerate}[label={$\langle1\rangle$\arabic*.},leftmargin=*]
\item There is nothing to prove when $\Sigma$ is rational, so
assume $g(\Sigma)\ge1$.
Statement \ref{cdp:st:A.5.8} gives a holomorphic ruling
$\widetilde p:T\to B$ with $p=\widetilde p\pi$.
Identifying $C$ with $B$ by its section map, normalization
of its nonconstant map to $D$ gives a finite surjective
holomorphic map $\phi:B\to\Sigma$.
Put $\psi=\phi\widetilde p:T\to\Sigma$.

\item We describe precisely the finite identifications needed to
make this map descend to $S$. Every vertical conductor
component of $T$ lies over one of finitely many points of
$B$, so $\psi$ is constant on it. Its image under $\eta$
has rational normalization: a vertical component is the
image of a rational curve on $Z$, and a nonconstant map
from that curve to its image normalization is finite.
Such an image cannot equal $D$, because $g(\Sigma)\ge1$.
Consequently these vertical images meet $D$ at finitely
many points. Away from those points and from the finitely
many singular points of $D$, every normalization preimage
of a point of $D$ lies on $C$, and its $\psi$ value is the
unique corresponding point of $\Sigma$.
At a point outside $D$ with more than one normalization
preimage, all those preimages lie on vertical conductor
components. There are no other horizontal conductor
components by hypothesis. Thus there is a finite set
$A\subset\Sigma$ such that, for every $s\in S$, either
$\psi(\eta^{-1}(s))$ is a singleton or it is contained
in $A$. One can take $A$ to contain the constant values
on all vertical conductor components and all points of
$\Sigma$ above the singularities of $D$ and its
intersections with the vertical image curves.

\item Choose an embedded finite tree $J$ in the underlying
oriented surface $\Sigma$ containing $A$, and let
$q:\Sigma\to\Sigma/J$ collapse that tree to one point.
Such a tree is obtained by joining the finite set by
arcs and retaining a spanning tree after subdividing
their finitely many crossings. Triangulate the surface
with $J$ a subcomplex. A contractible subcomplex is a
cofibration; contracting it gives a homotopy equivalence
of this CW pair. Concretely the homotopy extension
property extends a contraction of $J$ and supplies a
homotopy inverse to the quotient map. In particular
\[
              q_*:\pi_1(\Sigma)\xrightarrow{\sim}
                         \pi_1(\Sigma/J).          \tag{A.Pi.1}\hypertarget{tag-006O}{}\label{cdp:eq:A.Pi.1}
\]
The map $q\psi$ is constant on every fibre of the finite
proper surjection $\eta$. A proper continuous surjection
between these compact Hausdorff spaces is a quotient
map, so there is a unique continuous map
\[
 a:S\longrightarrow\Sigma/J,
 \qquad a\eta=q\psi.                              \tag{A.Pi.2}\hypertarget{tag-006P}{}\label{cdp:eq:A.Pi.2}
\]
Let $\delta:\Sigma\to D\subset S$ be the curve normalization
map. On its dense smooth locus, lifting through $C$ gives
$a\delta=q$. Continuity gives this equality everywhere.
Equation \eqref{cdp:eq:A.Pi.1} therefore makes $a_*$ surjective on
fundamental groups, without any assertion that a finite
point-identification quotient has the same fundamental
group as $\Sigma$.

\item We next extend the actual map \eqref{cdp:eq:A.Pi.2} to a neighbourhood
of $S$ in $Y$. The simultaneous analytic triangulation
theorem of \cite[\S3, Theorem~2, p.~464]{Lojasiewicz1964}
triangulates the pair $(Y,S)$, since a complex analytic
subset is real semianalytic. Compactness makes the
subcomplex representing $S$ finite. After barycentric
subdivision it is full: vertices lying in it span a
face in that same subcomplex. Write $\lambda_v$ for
barycentric coordinates and let $L$ be its vertex set.
On the open neighbourhood where $\sum_{v\in L}\lambda_v>0$,
the formula
\[
 r(x)=\frac{\sum_{v\in L}\lambda_v(x)v}
                  {\sum_{v\in L}\lambda_v(x)}       \tag{A.Pi.3}\hypertarget{tag-006Q}{}\label{cdp:eq:A.Pi.3}
\]
defines a continuous retraction to $S$ in triangulation
coordinates. Its value lies in the face spanned by the
$L$ vertices of the same simplex, so these formulas
agree across faces. This supplies an actual open
neighbourhood $N$ and a map $r:N\to S$ fixed on $S$.

\item Let $c=f(S)$. Properness allows a small disc $\Delta$
about $c$ with $U=f^{-1}(\Delta)\subset N$ and no other
critical value. The map $(ar)|_U:U\to\Sigma/J$ is
surjective on fundamental groups: its restriction to
$\delta(\Sigma)\subset S\subset U$ is $q$, whose
fundamental-group map is an isomorphism.
Set $U^*=U\setminus S=f^{-1}(\Delta^*)$.
The inclusion induces a surjection
$\pi_1(U^*)\twoheadrightarrow\pi_1(U)$.
Indeed a loop in the smooth six-manifold $U$ can be
perturbed relative to its base point to be transverse
to an analytic stratification of $S$; every stratum
has real codimension at least two, so the resulting
one-dimensional loop misses every stratum.

\item Over $\Delta^*$, the proper submersion is a smooth
torus bundle. The fundamental-group exact sequence
and $\pi_2(\Delta^*)=0$ give
\[
 1\longrightarrow\Z^4\longrightarrow\pi_1(U^*)
       \longrightarrow\Z\longrightarrow1.
\]
A lift of the generator of $\Z$ splits this group
extension, hence
\[
 \pi_1(U^*)\simeq\Z^4\rtimes_T\Z
 \twoheadrightarrow\pi_1(U)
 \twoheadrightarrow\pi_1(\Sigma/J)
 \simeq\pi_1(\Sigma).                             \tag{A.Pi.4}\hypertarget{tag-006R}{}\label{cdp:eq:A.Pi.4}
\]
The first group is metabelian: its commutator subgroup
is contained in the abelian normal subgroup $\Z^4$.
Every quotient in \eqref{cdp:eq:A.Pi.4} is therefore metabelian.
A closed orientable surface of genus at least two
has a fundamental group surjecting onto the free
group on two generators: in its usual surface
presentation, retain two $a$ generators and send all
$b$ generators and the other $a$ generators to the
identity. That free group surjects onto $S_4$ by
sending its generators to $(1234)$ and $(12)$;
their conjugates supply the adjacent transpositions.
The commutator subgroup of $S_4$ is the nonabelian
group $A_4$, so this quotient is not metabelian.
This contradicts \eqref{cdp:eq:A.Pi.4}. Thus $g(\Sigma)=1$ in the
positive-genus case. When $C\to D$ has degree one,
$\Sigma\simeq C\simeq B$, which proves the final
assertion.

\end{enumerate}
\end{proof}






\begin{proposition}[The arithmetic genus of conductor images from coherent duality]\hypertarget{tag-01B2}{}\label{cdp:st:A.5.12}
\StatementRecord{cdp:st:A.5.12}{proved}{01B2}{\SeriesRef{cdp:eq:A.CD.2}, \SeriesRef{cdp:eq:A.CD.3}, \SeriesRef{cdp:eq:A.CD.1}}

Let $f:Y\to C_0$ be a proper connected-fibre map from a
smooth compact complex threefold with $H^2(Y,\Z)=0$ to
a smooth curve, and let its smooth general fibre have
trivial canonical bundle. This includes complex
two-tori and primary Kodaira surfaces.
Let $mS$ be a full fibre with irreducible
reduced nonnormal support $S$. Write $\eta:T\to S$ for
normalization and $\pi:Z\to T$ for a resolution.
If the pullback homomorphism
$\Pic(Y)\to\Pic(Z)$ is nonconstant, then every reduced
irreducible component $D$ of the conductor image scheme
has arithmetic genus at most one.
In particular every such component whose normalization
has positive genus is a smooth elliptic curve.
The conductor image scheme itself is allowed to be
nonreduced, and no local-complete-intersection assumption
is imposed on its reduced components.
\end{proposition}


\begin{proof}

\begin{enumerate}[label={$\langle1\rangle$\arabic*.},leftmargin=*]
\item First we determine exactly the needed canonical line on
the reduced support, without a multiplicity formula for
the direct image. The sheaf $f_*K_Y$ is torsion-free of
generic rank one on the smooth curve and hence is a
line bundle. Choose a frame near $c=f(S)$ and evaluate
it in $K_Y$. A coordinate $t$ on the base disc supplies
the conormal frame $dt$, and adjunction identifies
$K_Y|_F$ with $K_F$ after dividing by this frame.
Under that line-bundle isomorphism, the evaluated
section restricts on a general smooth fibre to a
nonzero canonical two-form, which is nowhere zero.
Its zero divisor is therefore
vertical. Shrinking the base disc to remove the other
vertical zero divisors, irreducibility of the support
gives an effective divisor $kS$, with $k\ge0$, and an
actual line-bundle comparison on $U=f^{-1}(\Delta)$:
\[
 K_Y|_U\simeq\OO_U(kS),\qquad
 \omega_S\simeq\OO_Y((k+1)S)|_S,\qquad
 \omega_S^{\otimes m}\simeq\OO_S.                 \tag{A.CD.1}\hypertarget{tag-0058}{}\label{cdp:eq:A.CD.1}
\]
The last isomorphism uses the local base equation whose
zero divisor is exactly the full fibre $mS$.
In particular no assertion $k=m-1$ is used.

\item Hirzebruch--Riemann--Roch in
$0\to\OO_Y(-S)\to\OO_Y\to\OO_S\to0$ gives
$\chi(\OO_S)=0$: the rational Chern character of
$\OO_Y(-S)$ equals that of $\OO_Y$, since its first
Chern class is zero. A nonzero section of the torsion
line $\omega_S$ trivializes it. Indeed its $m$th power
is a nonzero holomorphic function on the compact
connected reduced space and hence is a nonzero constant.
Thus, with $e=h^0(S,\omega_S)$, we have
\[
 e\in\{0,1\},\qquad
 e=1\Longleftrightarrow\omega_S\simeq\OO_S,
 \qquad h^1(S,\OO_S)=1+e.                         \tag{A.CD.2}\hypertarget{tag-0059}{}\label{cdp:eq:A.CD.2}
\]
Here $h^0(\OO_S)=1$ and Serre duality identifies
$h^2(\OO_S)$ with $e$.

\item Let $\mathcal Q=\eta_*\OO_T/\OO_S$ and let
$\Gamma\subset S$ be the conductor image scheme.
Both $S$ and $T$ are connected and compact, so their
global holomorphic functions are the same constants.
The normalization sequence consequently gives an
injection and the indicated kernel:
\[
 H^0(S,\mathcal Q)\hookrightarrow H^1(S,\OO_S),
 \qquad
 \operatorname{im}H^0(S,\mathcal Q)
 =\ker\bigl(H^1(S,\OO_S)\to H^1(T,\OO_T)\bigr).
                                                               \tag{A.CD.3}\hypertarget{tag-005A}{}\label{cdp:eq:A.CD.3}
\]
Because $H^2(Y,\Z)=0$, $\Pic(Y)$ is the connected
exponential image of $H^1(Y,\OO_Y)$. A nonconstant
holomorphic group homomorphism out of it has nonzero
differential. The differential of the stated Picard
pullback is the actual composite
\[
 H^1(Y,\OO_Y)\longrightarrow H^1(S,\OO_S)
 \longrightarrow H^1(T,\OO_T)
 \longrightarrow H^1(Z,\OO_Z).
\]
Thus its middle arrow has positive rank. Equations
\eqref{cdp:eq:A.CD.2}--\eqref{cdp:eq:A.CD.3} imply $h^0(\mathcal Q)\le e$.

\item The conductor is Cohen--Macaulay of dimension one,
and the exact duality comparison (A.RC.3) applies
without a reducedness assumption on $\Gamma$:
$\mathcal Q=\omega_\Gamma\otimes\omega_S^{-1}|_\Gamma$.
Serre duality on that curve therefore gives
\[
 h^1(\Gamma,\omega_S|_\Gamma)
       =h^0(S,\mathcal Q)\le e.                   \tag{A.CD.4}\hypertarget{tag-005B}{}\label{cdp:eq:A.CD.4}
\]
For any reduced irreducible component $D$ of $\Gamma$,
restriction of the line gives an exact quotient
$\omega_S|_\Gamma\to\omega_S|_D$. Its kernel is a
coherent sheaf on a one-dimensional compact space and
has no $H^2$, so the corresponding map on $H^1$ is
surjective. Hence
$h^1(D,\omega_S|_D)\le e$.

\item For $e=1$, \eqref{cdp:eq:A.CD.2} makes this precisely
$p_a(D)=h^1(D,\OO_D)\le1$.
For $e=0$, the same inequality makes
$h^1(D,\omega_S|_D)=0$.
The restriction $\omega_S|_D$ is torsion by \eqref{cdp:eq:A.CD.1}
and has degree zero. Curve Riemann--Roch now gives
\[
 0\le h^0(D,\omega_S|_D)
       =\chi(D,\omega_S|_D)=1-p_a(D),
\]
again proving $p_a(D)\le1$.
Finally, for a reduced irreducible curve the
normalization exact sequence gives
$p_a(D)=g(D^\nu)+\operatorname{length}
(\nu_*\OO_{D^\nu}/\OO_D)$.
Positive normalization genus therefore forces genus
one and zero normalization quotient. A normal complex
curve is smooth, which proves the final assertion.

\end{enumerate}
\end{proof}






\begin{proposition}[The conductor bound when the general fibre has $q=1$ and $p_g=0$]\hypertarget{tag-01B3}{}\label{cdp:st:A.5.13}
\StatementRecord{cdp:st:A.5.13}{proved}{01B3}{\SeriesRef{cdp:st:D.1.9}, \SeriesRef{cdp:st:A.5.12}}

Let $f:Y\to\PP^1$ be a proper connected-fibre map from
a smooth compact complex threefold with
$H^1(Y,\Z)=H^2(Y,\Z)=0$. Suppose its smooth general
fibre has $h^1(\OO)=1$ and $h^2(\OO)=0$.
Then
\[
 R^1f_*\OO_Y=\OO_{\PP^1},\qquad R^2f_*\OO_Y=0,
 \qquad\Pic(Y)=\C\beta,                            \tag{A.Q1.1}\hypertarget{tag-006S}{}\label{cdp:eq:A.Q1.1}
\]
where the last identification uses the exponential
coordinates of \ref{cdp:st:D.1.9} and its nonzero transgression
class $\beta$. Every full fibre has structure-sheaf
cohomology dimensions $(1,1,0)$.
For an irreducible reduced support $S$, every
nonconstant ambient Picard pullback to a normalization
resolution forces every reduced conductor component
to have arithmetic genus at most one. In particular
each positive-genus component is smooth elliptic,
without any lci or conductor-reducedness hypothesis.
\end{proposition}


\begin{proof}

\begin{enumerate}[label={$\langle1\rangle$\arabic*.},leftmargin=*]
\item By \cite[Theorem~2.2, pp.~4--5]{Viaclovsky2026Combs}, $E=R^1f_*\OO_Y$ is a line bundle. The section
induced by \ref{cdp:st:D.1.9}'s integral transgression class is
nowhere zero: its image under the natural base-change
map at each point is the nonzero coefficient class on
that full fibre. A zero of the section in $E$ would
have zero image under that map. Thus $E\simeq\OO$.
Coherent Leray now gives $H^1(Y,\OO_Y)=\C\beta$.
The exponential sequence and the two stated integral
cohomology vanishings give the Picard identification
in \eqref{cdp:eq:A.Q1.1}.

\item The coherent sheaf $G=R^2f_*\OO_Y$ has generic rank
zero, hence is a torsion sheaf. Riemann--Roch on $Y$
gives $\chi(\OO_Y)=\int_Yc_1c_2/24=0$, since
$c_1(Y)$ is zero in rational cohomology. The Euler
characteristic of coherent Leray is therefore
\[
 0=\chi(\OO_{\PP^1})-\chi(E)+\chi(G)
      =\operatorname{length}G.
\]
Thus $G=0$. Top-degree base change gives
$H^2(F,\OO_F)=0$ for every full Cartier fibre $F$.
The multiplication-by-parameter sequence then gives
$H^1(F,\OO_F)=E\otimes\C(c)$, of dimension one;
\cite[Theorem~2.2, pp.~4--5]{Viaclovsky2026Combs} gives $h^0(F,\OO_F)=1$.

\item For a full fibre $F=mS$ with irreducible reduced
support, the quotient $\OO_F\to\OO_S$ induces a
surjection on $H^2$, because its kernel has support
of dimension two and has no $H^3$.
Hence $h^2(S,\OO_S)=0$. Riemann--Roch in the Cartier
sequence on $Y$ gives $\chi(\OO_S)=0$ as in \ref{cdp:st:A.5.12},
so $h^1(S,\OO_S)=1$.

\item Put $\mathcal Q=\eta_*\OO_T/\OO_S$ for the
normalization $T$ and write $\Gamma$ for the conductor
image scheme. Nonconstant Picard pullback has nonzero
differential factoring through
$H^1(S,\OO_S)\to H^1(T,\OO_T)$.
The source has dimension one, so this map has zero
kernel. The normalization sequence (A.CD.3) gives
$H^0(\mathcal Q)=0$. Conductor duality and Serre
duality give
\[
 H^1(\Gamma,\omega_S|_\Gamma)=0,
 \qquad H^1(D,\omega_S|_D)=0                      \tag{A.Q1.2}\hypertarget{tag-006T}{}\label{cdp:eq:A.Q1.2}
\]
for each reduced irreducible component $D$, where
the second map is the restriction surjection in the
proof of \ref{cdp:st:A.5.12}. Adjunction makes $\omega_S$ ambient,
so $\deg(\omega_S|_D)=0$ by $H^2(Y,\Z)=0$.
Curve Riemann--Roch yields
$0\le h^0(D,\omega_S|_D)=1-p_a(D)$.
The normalization genus and its nonnegative delta
length then give the asserted smooth elliptic
conclusion whenever that genus is positive.

\end{enumerate}
\end{proof}






\begin{proposition}[Conductor-cleared jets on a complete-intersection critical curve]\hypertarget{tag-01BF}{}\label{cdp:st:A.5.5}
\StatementRecord{cdp:st:A.5.5}{proved}{01BF}{\SeriesRef{cdp:eq:A.J.4}, \SeriesRef{cdp:eq:A.J.5}, \SeriesRef{cdp:eq:A.J.2}, \SeriesRef{cdp:eq:A.J.3}, \SeriesRef{cdp:eq:A.J.7}, \SeriesRef{cdp:eq:A.J.9}, \SeriesRef{cdp:eq:A.J.8}, \SeriesRef{cdp:eq:A.J.10}}

Let $Y$ be a smooth compact complex threefold with
$H^2(Y,\mathbb Q)=0$, let $S\subset Y$ be a reduced divisor,
and let $D\subset\operatorname{Sing}(S)$ be a reduced
irreducible compact curve which is a local complete intersection
in $Y$. Assume that the generic plane-curve transversal to $D$
in $S$ is an ordinary cusp, an ordinary tacnode, or an ordinary
triple point.
If the normalization $C\to D$ has positive genus, then $D$ is
a smooth elliptic curve. The transversal type is the indicated
cusp, tacnode, or triple point at every point of $D$;
in particular no degeneration of that type occurs along $D$.

In the cusp case the proof constructs, on $C$, the actual
meromorphic cubic on the radical of the Hessian and bounds its
pole divisor by twice the conductor. In the tacnode case the
corresponding fourth-order invariant has poles bounded by four
times the conductor. These intermediate jets need not extend
holomorphically before the conductor correction.
\end{proposition}


\begin{proof}

\begin{enumerate}[label={$\langle1\rangle$\arabic*.},leftmargin=*]
\item Let $\nu:C\to D$ be the normalization and $j:C\to Y$
its composite with the inclusion,
$I=\mathcal I_{D/Y}$, and
\[
 N=\bigl(\nu^*(I/I^2)\bigr)^*,\qquad
 M=j^*\mathcal O_Y(S).
\]
These are a rank-two vector bundle and a line bundle on $C$.
Let $\Xi$ be the conductor divisor on $C$, so that
$\nu_*\mathcal O_C(-\Xi)$ is the conductor ideal in
$\mathcal O_D$, and put $c=\deg\Xi\geq0$.
Adjunction for the complete-intersection curve and finite
duality give exact line-bundle comparisons
\[
 \det N=\nu^*\omega_D\otimes j^*K_Y^{-1},\qquad
 \omega_C=\nu^*\omega_D\otimes\mathcal O_C(-\Xi).
                                                        \tag{A.J.1}\hypertarget{tag-005C}{}\label{cdp:eq:A.J.1}
\]
For the second comparison one can work directly in the total
quotient ring: $\operatorname{Hom}_{\mathcal O_D}
(\nu_*\mathcal O_C,\mathcal O_D)$ consists precisely of
elements multiplying the normalization into $\mathcal O_D$,
which is its conductor. Since $D$ is Gorenstein, tensoring
this identity by its invertible dualizing sheaf is finite
duality. Both $M$ and $j^*K_Y$ have degree zero, by the
stated rational cohomology hypothesis. Consequently
\[
 \deg\det N=2g(C)-2+c.                              \tag{A.J.2}\hypertarget{tag-005N}{}\label{cdp:eq:A.J.2}
\]
This is the normal bundle of the complete-intersection ideal,
with its specified normalization pullback. It is the bundle
needed in the following jet calculation.

\item First treat an ordinary cusp. Choose a local defining section
$f$ of $S$, using a local frame of $\mathcal O_Y(S)$.
Since $D\subset\operatorname{Sing}(S)$, its conormal class
maps to zero under
$I/I^2\to\Omega_Y^1|_D$. This homomorphism is injective:
its source is locally free on the reduced curve and the map
is injective at every generic point. Therefore $f\in I^2$.
The associated-graded identity for a regular sequence gives
$I^2/I^3=\operatorname{Sym}^2(I/I^2)$. The class of $f$
is thus a quadratic form
\[
 q\in H^0\bigl(C,\operatorname{Sym}^2N^*\otimes M\bigr).
\]
It has rank one generically. Let $R\subset N$ be its
saturated radical and $Q=N/R$. Both are line bundles.
The quadratic form annihilates $R$ everywhere by holomorphic
continuation, and gives a nonzero section of $M\otimes Q^{-2}$.
In particular
\[
 \deg Q\leq0.                                      \tag{A.J.3}\hypertarget{tag-005Y}{}\label{cdp:eq:A.J.3}
\]

\item On $D_{\rm reg}$, the third-order term on the radical is
intrinsic. Indeed, after writing the quadratic form as
$a\ell^2$, subtract a holomorphic lift of $a\ell^2$ from
$f$ and restrict its cubic class to $R^3$. Changing that lift
adds terms divisible by $\ell$, which vanish on $R^3$.
Changing normal generators has the same effect. Its nonzero
value characterizes the ordinary-cusp transversal among the
rank-one quadratic cases. Denote this cubic by $k$.

\item We now extend precisely this section across $C\setminus
\nu^{-1}(D_{\rm reg})$. Work at a point of $D$ and choose
two generators $X,Y$ of $I$. After a constant invertible linear
change of these generators, the normalized linear factor of
$q$ has the form $X+hY$ on every normalization branch over
that point. Such a choice exists: the saturated quotient
$N\to Q$ is surjective at each of the finitely many points,
so a constant vector can be chosen outside their finitely
many kernels. Thus $h$ is holomorphic on the normalization,
and
\[
 q=a(X+hY)^2,
\]
where $a=q_{XX}$ is the pullback of an actual coefficient
in $\mathcal O_D$. Choose a local conductor generator $d$
on the normalization. Both $d$ and $dh$ belong to
$\mathcal O_D$, since the conductor is an ideal in the
normalization contained in that ring. Choose lifts
$\widehat d,\widehat{dh},\widehat a$ to $\mathcal O_Y$,
and set
\[
 G=\widehat d^{\,2} f-
    \widehat a\bigl(\widehat d X+\widehat{dh}Y\bigr)^2.
                                                        \tag{A.J.4}\hypertarget{tag-006C}{}\label{cdp:eq:A.J.4}
\]
The quadratic class of $G$ is zero after normalization.
It is already zero on $D$: the map from the locally free
sheaf $\operatorname{Sym}^2(I/I^2)$ to its normalization
pullback is injective because $D$ is reduced. Hence
$G\in I^3$.

\item Its class in $I^3/I^4=\operatorname{Sym}^3(I/I^2)$,
pulled back to $C$ and evaluated on $R^3$, is holomorphic.
Dividing that class by $d^2$ gives $k$ on the dense regular
locus. To check this equality, there $d$ is a unit and
\eqref{cdp:eq:A.J.4} is exactly the subtraction of the square used in the
definition of the radical cubic. Different choices of lifts
add cubic terms containing its linear factor, and therefore
have zero restriction to $R^3$. Thus these local expressions
are the same meromorphic section wherever they overlap.
Since $d$ has divisor $\Xi$, the resulting bound is
\[
 0\ne k\in H^0\bigl(C,M\otimes R^{-3}
                               \otimes\mathcal O_C(2\Xi)\bigr).
                                                        \tag{A.J.5}\hypertarget{tag-006E}{}\label{cdp:eq:A.J.5}
\]
This proves the claimed pole control without a splitting of
the formal neighbourhood of the singular curve.

\item Taking degrees in \eqref{cdp:eq:A.J.5} gives $\deg R\leq2c/3$.
Together with \eqref{cdp:eq:A.J.2}--\eqref{cdp:eq:A.J.3}, this yields
\[
 2g(C)-2+c=\deg R+\deg Q\leq\frac{2c}{3},
 \qquad 2g(C)-2+\frac c3\leq0.                     \tag{A.J.6}\hypertarget{tag-006F}{}\label{cdp:eq:A.J.6}
\]
For $g(C)\geq1$, this forces $g(C)=1$ and $c=0$.
Conductor zero means $\mathcal O_D=\nu_*\mathcal O_C$;
a normal complex curve is smooth. Hence $D=C$ is elliptic.
Now \eqref{cdp:eq:A.J.2}--\eqref{cdp:eq:A.J.5} give $\deg R=\deg Q=0$.
The Hessian coefficient and radical cubic are nonzero
sections of degree-zero line bundles on this elliptic curve,
so neither has a zero. Locally the holomorphic Morse lemma
in the nondegenerate normal direction, with the curve
coordinate as parameter, writes $f=y^2+g(x,z)$.
The conditions just proved give $g=x^3u(x,z)$ with $u$ a
holomorphic unit. Replacing $x$ by a local cube root of $u$
times $x$ gives the actual local equation $y^2+x^3=0$.
This proves the asserted constant transversal type.

\item Next treat an ordinary tacnode. Retain $N,M,R,Q,q,a,h,d,P$
from the cusp calculation, with
$P=\widehat d X+\widehat{dh}Y$ and
$G=\widehat d^{\,2}f-\widehat aP^2\in I^3$.
The generic radical cubic is now zero, so its continuation
$[G]_3|_{R^3}$ is zero everywhere on the normalization.
Since $\ell=X+hY$ is monic in $X$, polynomial division gives
an actual holomorphic quadratic $C_2$ on the normalization
such that
\[
 [G]_3=\ell C_2.
\]
Each coefficient of $dC_2$ is in the conductor and hence in
$\mathcal O_D$. Lift those coefficients and let $H$ be the
resulting homogeneous quadratic in $X,Y$. Set
\[
 K=\widehat d^{\,2}G-PH.
\]
Its cubic class pulls back to
$d^2\ell C_2-d\ell\cdot dC_2=0$.
The same locally free graded-piece injection used above gives
$K\in I^4$. There is an exact identity in the local ring of
the ambient smooth threefold:
\[
 4\widehat a\widehat d^{\,6}f
 -(2\widehat a\widehat d^{\,2}P+H)^2
 =4\widehat a\widehat d^{\,2}K-H^2.               \tag{A.J.7}\hypertarget{tag-006G}{}\label{cdp:eq:A.J.7}
\]
On the dense smooth locus, complete the square in the
nondegenerate normal direction. If the relevant terms there
are $a\ell^2+b\ell r^2+er^4$, the fourth-order invariant is
$\Delta=ae-b^2/4$. A substitution
$\ell\mapsto\ell+ur^2$ leaves this expression unchanged;
terms divisible by $\ell^2$ and higher changes in the radical
coordinate do not alter it. Thus it is the intrinsic section
of $M^2\otimes Q^{-2}\otimes R^{-4}$ on that dense locus.
It is nonzero precisely for the ordinary-tacnode transversal
among the rank-one Hessian cases with zero radical cubic.

\item Divide \eqref{cdp:eq:A.J.7} by $4\widehat d^{\,6}$ and take the
fourth-order class on $R^4$. Since
$[H]_2=dC_2$, this expresses exactly that same invariant as
\[
 \Delta=
 \frac{a[K]_4|_{R^4}-(C_2|_{R^2})^2/4}{d^4}.
                                                        \tag{A.J.8}\hypertarget{tag-006H}{}\label{cdp:eq:A.J.8}
\]
Its numerator is holomorphic on the normalization.
Different choices agree on the dense regular locus and hence
give one meromorphic section. Therefore
\[
 0\ne\Delta\in H^0\bigl(C,M^2\otimes Q^{-2}
                    \otimes R^{-4}\otimes\mathcal O_C(4\Xi)\bigr).
                                                        \tag{A.J.9}\hypertarget{tag-006I}{}\label{cdp:eq:A.J.9}
\]
Write $q_0=\deg Q\leq0$. Using
$\deg R=2g(C)-2+c-q_0$, the degree of the line bundle
in \eqref{cdp:eq:A.J.9} is
\[
 -2q_0-4\bigl(2g(C)-2+c-q_0\bigr)+4c
       =-8g(C)+8+2q_0.
\]
Its nonzero section and $g(C)\geq1$ force
$g(C)=1$ and $q_0=0$.

\item The section $q$ of $M\otimes Q^{-2}$ is now nowhere zero.
In the local presentation above, $a=q_{XX}$ is therefore a
unit on every normalization branch. Since it belongs to
$\mathcal O_D$, it is a unit in that local ring as well.
The coefficient $h=q_{XY}/a$ consequently belongs to
$\mathcal O_D$; here $q_{XY}$ denotes the coefficient in
the symmetric matrix, so the cross term is $2q_{XY}XY$.
We can thus repeat the entire construction with $d=1$.
Polynomial division by the monic $X+hY$ then takes place
over $\mathcal O_D$ itself; the zero radical cubic descends
by the same injectivity argument. Formula \eqref{cdp:eq:A.J.8} now has
no poles. The resulting nonzero section of
$M^2\otimes Q^{-2}\otimes R^{-4}$ has degree $-4c$,
which forces $c=0$.

\item Thus again $D=C$ is smooth elliptic. Both the nonzero
quadratic coefficient and $\Delta$ are sections of
degree-zero line bundles and are nowhere zero. The parametric
Morse lemma first gives $f=y^2+g(x,z)$, and the vanished
radical cubic and nonzero fourth invariant give
$g=x^4u(x,z)$ with $u$ a unit. Taking a local fourth root
produces the actual equation $y^2+x^4=0$ at every point of
$D$. This proves the tacnode assertion.

\item For an ordinary triple transversal, the quadratic class
vanishes generically, hence vanishes in the locally free
sheaf $I^2/I^3$. Thus $f\in I^3$, and its initial cubic
is an honest global section
$b\in H^0(C,\operatorname{Sym}^3N^*\otimes M)$.
The discriminant of a binary cubic is homogeneous of degree
four in its coefficients and transforms by the sixth power
of the determinant of a change of variables. Its value is
therefore a section
\[
 0\ne\operatorname{disc}(b)\in
 H^0\bigl(C,M^4\otimes(\det N)^{-6}\bigr).
                                                        \tag{A.J.10}\hypertarget{tag-005D}{}\label{cdp:eq:A.J.10}
\]
It is nonzero because the three generic tangent lines are
distinct. Equations \eqref{cdp:eq:A.J.2} and \eqref{cdp:eq:A.J.10} give
$2g(C)-2+c\leq0$, again forcing $g(C)=1$ and $c=0$
under the genus assumption. The discriminant has degree
zero and is nowhere zero. The three simple tangent roots
extend holomorphically on each small curve disc. The implicit
function theorem gives three smooth local branches of $S$
with those pairwise distinct normal tangents. Choose the
first two branches as the coordinate axes; the third is
$y=xu(x,z)$ with $u$ a unit, and replacing $y$ by $y/u(x,z)$
straightens it to $y=x$. This gives the local equation
$xy(x-y)=0$ up to a unit at every point of $D$.

\end{enumerate}
\end{proof}




The pole allowance in \eqref{cdp:eq:A.J.5} has a concrete local use. Put
$h=z^2-x^3$, $D=(h,w)=0$ in $\mathbb C^3$, and
\[
 f=xw^2+2zwh+x^2h^2+h^3.
\]
On the normalization $x=s^2,z=s^3$ of $D$, its quadratic
class is $s^2(w+sh)^2$. Completing the square on $D_{\rm reg}$
gives
\[
 f=x\left(w+\frac zx h\right)^2+
                  \left(1-\frac1x\right)h^3.
\]
The radical cubic has a pole of order two at $s=0$.
The conductor generator is $d=x=s^2$, and the exact
conductor-cleared identity is
\[
 x^2f-x(xw+zh)^2=x(x-1)h^3\in I_D^3.
\]
Thus the corrected jet comparison applies even in this local
model where holomorphic extension of the uncorrected cubic
would fail. The equalities are checked exactly in
\texttt{code/E\_singular\_curve.py}.

The local calculation below is related to the conductor
derivations and characteristic-zero pole calculation in
Reid~\cite[\S3.9, \S4.4--4.8]{Reid1994Nonnormal}.
We give the canonical-module comparison directly, without
using an anticanonical polarization or a classification of
the surface.




\begin{proposition}[A pure double conductor forces a smooth cusp image]\hypertarget{tag-01BH}{}\label{cdp:st:A.5.7}
\StatementRecord{cdp:st:A.5.7}{proved}{01BH}{\SeriesRef{cdp:eq:A.J.13}, \SeriesRef{cdp:eq:A.J.14}, \SeriesRef{cdp:eq:A.J.12}}

Let $S$ be a reduced Gorenstein complex surface and
$\pi:\overline S\to S$ its finite normalization.
Suppose that, over a neighbourhood of a point $p\in S$,
$\overline S$ is smooth and its conductor ideal is exactly
$\mathcal O_{\overline S}(-2C)$ for a smooth reduced
divisor $C$; its local branches over $p$ are disjoint.
Suppose that the generic transversal of $S$ along the
reduced image $D=\pi(C)$ is an ordinary cusp.
Then $D$ is smooth at $p$, $C\to D$ is an isomorphism
near the preimage of $p$, and the germ of $S$ at $p$ is
analytically a product of an ordinary cusp and a smooth disc.
No complete-intersection hypothesis on $D$ is required.
\end{proposition}


\begin{proof}

\begin{enumerate}[label={$\langle1\rangle$\arabic*.},leftmargin=*]
\item Write $R=\mathcal O_{S,p}$ and let $\overline R$ be its
finite semilocal normalization. It is a product of regular
two-dimensional local analytic rings. Write $\mathfrak c$
for the conductor, and set
\[
 A=R/\mathfrak c,\qquad
 \overline A=\overline R/\mathfrak c,\qquad
 Q=\overline R/R=\overline A/A.
\]
We first give the exact duality comparison needed below.
The $R$-module $\overline R$ is maximal Cohen--Macaulay.
Since $R$ is Gorenstein,
$\operatorname{Ext}_R^j(\overline R,R)=0$ for $j>0$.
Its dual $\operatorname{Hom}_R(\overline R,R)$ is
precisely the conductor, viewed as a fractional ideal:
a homomorphism is multiplication in the total quotient ring,
and multiplication takes $\overline R$ into $R$ exactly
for conductor elements. Applying $\operatorname{Hom}_R(-,R)$
to $0\to R\to\overline R\to Q\to0$ therefore gives
\[
 \operatorname{Ext}_R^1(Q,R)=A.                    \tag{A.J.11}\hypertarget{tag-005E}{}\label{cdp:eq:A.J.11}
\]
The depth lemma shows that $Q$ is a Cohen--Macaulay module
of dimension one. Gorenstein duality on such modules is the
involutive functor $M\mapsto\operatorname{Ext}_R^1(M,R)$.
Consequently $A$ is Cohen--Macaulay and
\[
 Q\simeq\operatorname{Ext}_R^1(A,R)=\omega_A,
 \qquad \operatorname{End}_A(Q)=A.                \tag{A.J.12}\hypertarget{tag-005F}{}\label{cdp:eq:A.J.12}
\]
The final equality is the canonical multiplication map:
duality identifies its endomorphism ring with
$\operatorname{End}_A(A)=A$.

\item The ordinary-cusp hypothesis says that $A$ is generically
reduced. Its nilradical is therefore supported at finitely
many points; a one-dimensional Cohen--Macaulay ring has
no nonzero submodule with that support. Hence $A$ is reduced.
Choose local coordinates $(s_i,z_i)$ on the normalization
branches with $C=(z_i=0)$. Put
\[
 B=\prod_i\mathbb C\{s_i\},\qquad
 K=\prod_i\operatorname{Frac}\bigl(\mathbb C\{s_i\}\bigr).
\]
The specified conductor gives the actual ring comparison
\[
 \overline A=B[\epsilon]/(\epsilon^2),             \tag{A.J.13}\hypertarget{tag-005G}{}\label{cdp:eq:A.J.13}
\]
component by component. The projection
$\overline A\to B$ is injective on $A$, since its kernel
is nilpotent and $A$ is reduced. It identifies $B$ with
the finite normalization of the curve ring $A$.

\item Under \eqref{cdp:eq:A.J.13} the inclusion of $A$ takes the form
\[
 a\longmapsto a+\delta(a)\epsilon,
\]
where $\delta:A\to B$ is a complex-linear analytic
derivation. It extends to the total quotient ring, and on
each normalization branch has the form
$\delta(a)=u_i\,da/ds_i$ for a meromorphic germ $u_i$.
Let $u=(u_i)$ and write derivatives componentwise.
The map
\[
 \Phi:B[\epsilon]/(\epsilon^2)\longrightarrow K,
 \qquad b+c\epsilon\longmapsto c-u b'
                                                        \tag{A.J.14}\hypertarget{tag-005H}{}\label{cdp:eq:A.J.14}
\]
is $A$-linear for this specified action: multiplication by
$a+u a'\epsilon$ changes its value to
$ac+u a'b-u(ab)'=a(c-u b')$. Its kernel at every generic
point is exactly the graph of the extended derivation on
$K$. It therefore induces a generic isomorphism from $Q$
to its image in $K$. The module $Q\simeq\omega_A$ is
torsion-free on the reduced curve, so this induced map is
injective.

\item Differentiation of convergent power-series germs onto
convergent power-series germs is surjective, by termwise
integration. Hence the image in \eqref{cdp:eq:A.J.14} is exactly
\[
 I=B+uB\subset K.                                \tag{A.J.15}\hypertarget{tag-005I}{}\label{cdp:eq:A.J.15}
\]
It is an invertible fractional ideal on the semilocal
normalization $B$. Every $A$-linear endomorphism of this
rank-one module is multiplication by an element of $K$;
preserving $I$ is equivalent, branch by branch, to belonging
to $B$. Thus $\operatorname{End}_A(I)=B$ as subrings of
$K$. Combining this with the actual $A$-linear isomorphism
$Q\xrightarrow{\Phi}I$ and \eqref{cdp:eq:A.J.12} gives $A=B$.
Because $A$ is local, there is a single branch over $p$,
and its ring is regular. This proves that $D$ is smooth
and that $C\to D$ is an isomorphism there.

\item Finally $A=B=\mathbb C\{s\}$ makes $u=\delta(s)$
holomorphic. The analytic coordinate change
$s'=s+u(s)z$ has invertible differential at $z=0$ and
identifies the image of $A$ modulo $z^2$ with the functions
of $s'$ alone. Since $R$ is the inverse image of $A$ under
$\overline R\to\overline A$, it follows that
\[
 R=\mathbb C\{s'\}+z^2\mathbb C\{s',z\}
   =\mathbb C\{s',z^2,z^3\}
   \simeq\mathbb C\{s',x,y\}/(y^2-x^3).
                                                        \tag{A.J.16}\hypertarget{tag-005J}{}\label{cdp:eq:A.J.16}
\]
This is the asserted product germ.

\end{enumerate}
\end{proof}






\begin{proposition}[Rational singularities for a vertically resolved ruled normalization]\hypertarget{tag-01BI}{}\label{cdp:st:A.5.8}
\StatementRecord{cdp:st:A.5.8}{proved}{01BI}{\SeriesRef{cdp:eq:A.J.19}}

Let $\pi:Z\to T$ be a resolution of a compact connected
normal complex surface, where $Z$ is an irrational ruled
surface with ruling $p:Z\to B$. Suppose that every
exceptional curve of $\pi$ is vertical for $p$.
Then there is a unique holomorphic map
$\widetilde p:T\to B$ satisfying
$p=\widetilde p\circ\pi$, and
\[
 R\pi_*\mathcal O_Z=\mathcal O_T,
 \qquad R^1\widetilde p_*\mathcal O_T=0.           \tag{A.J.17}\hypertarget{tag-005K}{}\label{cdp:eq:A.J.17}
\]
In particular all singularities of $T$ are rational. For
$g=g(B)$, its complete structure-sheaf cohomology dimensions
are
\[
 (h^0(T,\mathcal O_T),h^1(T,\mathcal O_T),
                         h^2(T,\mathcal O_T))=(1,g,0).
                                                        \tag{A.J.18}\hypertarget{tag-005L}{}\label{cdp:eq:A.J.18}
\]
\end{proposition}


\begin{proof}

\begin{enumerate}[label={$\langle1\rangle$\arabic*.},leftmargin=*]
\item The exceptional fibre of the proper birational map $\pi$
over a point is connected, because $T$ is normal and
$\pi_*\mathcal O_Z=\mathcal O_T$. Its curve components
are vertical by assumption, so their intersection
connectivity makes $p$ constant on that whole fibre.
It is also constant on the one-point fibres. We verify
holomorphic descent, rather than only descending the
underlying continuous map.

\item Fix $t\in T$ and let $U\subset B$ be a coordinate disc
containing $p(\pi^{-1}(t))$. Properness makes
$\pi(Z\setminus p^{-1}(U))$ closed and disjoint from
$t$. There is therefore an open neighbourhood $W$ of $t$
with $\pi^{-1}(W)\subset p^{-1}(U)$. The coordinate
function of $p$ on $\pi^{-1}(W)$ is a section of
$\pi_*\mathcal O_Z(W)=\mathcal O_T(W)$, so descends
to a holomorphic function on $W$. Its values lie in the
coordinate image of $U$; the inverse coordinate gives a
holomorphic map $W\to U$. These local maps agree after
pullback to the surjective map $\pi$, hence agree on
overlaps. They construct the required unique
$\widetilde p$.

\item Contracting vertical ruling $(-1)$-curves factors $p$ as
\[
 Z\xrightarrow{b} Z_0\xrightarrow{p_0}B,
\]
where $Z_0\to B$ is a $\mathbb P^1$-bundle and $b$
is a finite sequence of point blowups. Every such blowup
has direct image $\mathcal O$ and no higher direct image
of $\mathcal O$: the successive layers on its
infinitesimal exceptional curves are
$\mathcal O_{\mathbb P^1}(k)$ for $k\geq0$, whose
degree-one cohomology is zero, and formal functions gives
the assertion. The bundle $p_0$ has the same direct-image
property by the cohomology of $\mathbb P^1$ and local
triviality. Composition consequently gives
\[
 p_*\mathcal O_Z=\mathcal O_B,
 \qquad R^j p_*\mathcal O_Z=0\quad(j>0).          \tag{A.J.19}\hypertarget{tag-005M}{}\label{cdp:eq:A.J.19}
\]

\item The sheaf $\mathcal T=R^1\pi_*\mathcal O_Z$ has
finite support at the singular points of $T$. The
low-degree Leray sequence for $p=\widetilde p\circ\pi$
contains the actual arrows
\[
 0\longrightarrow R^1\widetilde p_*\mathcal O_T
 \longrightarrow R^1p_*\mathcal O_Z
 \longrightarrow\widetilde p_*\mathcal T
 \longrightarrow R^2\widetilde p_*\mathcal O_T.    \tag{A.J.20}\hypertarget{tag-005O}{}\label{cdp:eq:A.J.20}
\]
The last term is zero: the nonconstant proper map
$\widetilde p$ has fibres of complex dimension one.
The middle term is zero by \eqref{cdp:eq:A.J.19}. Thus
$R^1\widetilde p_*\mathcal O_T=0$ and
$\widetilde p_*\mathcal T=0$. Pushforward is faithful
on a sheaf with finite support, since its stalk sections
at finitely many points cannot disappear under a proper
map. Hence $\mathcal T=0$. Higher $R^j\pi_*$ vanish
for $j>1$ by the dimension of its fibres. This proves
\eqref{cdp:eq:A.J.17} and the defining rational-singularity property.

\item Finally \eqref{cdp:eq:A.J.17} identifies every
$H^q(T,\mathcal O_T)$ with $H^q(Z,\mathcal O_Z)$.
Equation \eqref{cdp:eq:A.J.19} identifies the latter with
$H^q(B,\mathcal O_B)$. A compact connected smooth curve
has dimensions $(1,g)$ in degrees zero and one and no
higher cohomology, proving \eqref{cdp:eq:A.J.18}.

\end{enumerate}
\end{proof}






\begin{proposition}[The genus bound and delta-one curve germs]\hypertarget{tag-01B4}{}\label{cdp:st:A.5.14}
\StatementRecord{cdp:st:A.5.14}{proved}{01B4}{\SeriesRef{cdp:eq:A.J.31}, \SeriesRef{cdp:st:A.5.5}}

Let $f:Y\to C_0$ be a proper connected-fibre map from a
smooth compact complex threefold with $H^2(Y,\Q)=0$ to
a smooth curve. Suppose its smooth general fibre has
trivial canonical bundle. Let $mS$ be a full fibre with
irreducible reduced nonnormal support, and let $D$ be
a reduced irreducible component of its conductor image.
Then
\begin{equation}
                         p_a(D)\le2.                \hypertarget{tag-005Z}{}\label{cdp:eq:A.J.30}
\end{equation}
If the normalization of $D$ has positive genus, then
$D$ is a local complete intersection in $Y$; every
singular germ of $D$ is an ordinary node or an ordinary
cusp, and the sum of its delta invariants is at most
one. In particular, when the generic transverse type
along such a positive-genus $D$ is $A_2$, $A_3$, or an
ordinary triple point, $D$ is smooth elliptic and the
transverse analytic type is constant along all of $D$.
\end{proposition}


\begin{proof}

\begin{enumerate}[label={$\langle1\rangle$\arabic*.},leftmargin=*]
\item Restrict $f$ to a disc about the image of $mS$ and trivialize the
canonical line of the disc. The sheaf $f_*K_Y$ is torsion-free of rank
one, hence locally free there. The evaluation of a generator is
nonzero on the smooth nearby fibres, where their canonical lines are
trivial. Its zero divisor is therefore $aS$ for an integer $a\ge0$.
Thus $K_Y\simeq\OO_Y(aS)$ on this neighbourhood. The full-fibre equation
trivializes $\OO_Y(mS)$, and adjunction gives
$\omega_S\simeq\OO_Y((a+1)S)|_S$, a torsion line.
The Cartier sequence and Riemann--Roch on $Y$ give $\chi(\OO_S)=0$,
since $c_1(\OO_Y(S))=0$ in rational cohomology. A nonzero section of a
torsion line on the compact irreducible surface has a nonzero constant
tensor power; it therefore trivializes the line. Consequently
$e=h^0(S,\omega_S)\in\{0,1\}$, and Serre duality gives
$h^1(\OO_S)=1+e$. For $e=1$, the line $\omega_S$ is trivial.

\item Write $\eta:T\to S$ for normalization,
$\mathcal Q=\eta_*\OO_T/\OO_S$, and $\Gamma$ for
the full conductor image scheme. Equality of the
constant global functions on $S$ and $T$ gives the
injection $H^0(\mathcal Q)\hookrightarrow
H^1(\OO_S)$. To identify the conductor-duality map, put
$\mathcal B=\eta_*\OO_T$ and $\mathfrak c=\operatorname{Ann}(\mathcal Q)$.
The normal surface $T$ is Cohen--Macaulay. Thus $\mathcal B$ is a
maximal Cohen--Macaulay module on the Gorenstein surface $S$.
Multiplication in the total quotient sheaf identifies
$\mathcal Hom(\mathcal B,\omega_S)$ with $\mathfrak c\omega_S$.
Applying $\mathcal Hom(-,\omega_S)$ to
$0\to\OO_S\to\mathcal B\to\mathcal Q\to0$ therefore gives
$\mathcal Ext^1(\mathcal Q,\omega_S)=\omega_S|_\Gamma$ and vanishing
of the other sheaf Ext groups. The local-to-global Ext sequence and
Serre duality identify
$H^0(S,\mathcal Q)^*$ with
$H^1(\Gamma,\omega_S|_\Gamma)$. Taking dimensions gives
\begin{equation}
 h^1(\Gamma,\omega_S|_\Gamma)
       =h^0(S,\mathcal Q)\le1+e.                  \hypertarget{tag-0060}{}\label{cdp:eq:A.J.31}
\end{equation}
The restriction quotient
$\omega_S|_\Gamma\twoheadrightarrow\omega_S|_D$
has a kernel supported in dimension at most one, so
its $H^1$ map is surjective. For $e=1$ it follows that
$p_a(D)=h^1(\OO_D)\le2$. For $e=0$, put
$L=\omega_S|_D$. This torsion line has degree zero,
and \eqref{cdp:eq:A.J.31} gives $h^1(D,L)\le1$. Curve
Riemann--Roch gives
\[
 p_a(D)=1+h^1(D,L)-h^0(D,L)\le2,
\]
proving \eqref{cdp:eq:A.J.30} in both cases.

\item Let $\nu:\widehat D\to D$ be normalization. The
normalization sequence gives the identity
\begin{equation}
 p_a(D)=g(\widehat D)+
        \sum_{x\in D}\dim_{\C}
          (\nu_*\OO_{\widehat D}/\OO_D)_x.       \hypertarget{tag-0061}{}\label{cdp:eq:A.J.32}
\end{equation}
If $g(\widehat D)>0$, every local delta is consequently
at most one and their sum is at most one. We give the
analytic local rings at a point with delta one.
Let $A=\OO_{D,x}$ and let
$B=\bigoplus_{i=1}^r\C\{t_i\}$ be its finite
normalization. Here $\dim_{\C}(B/A)=1$. The residue map
of the local ring $A$ has the same value on every
normalization branch. The resulting quotient of
$\C^r$ by its diagonal has dimension $r-1$, so
$r\le2$.

\item For $r=2$, the subring
\[
 A_0=\{(b_1,b_2)\in\C\{t_1\}\oplus\C\{t_2\}:
                                  b_1(0)=b_2(0)\}
\]
contains $A$ and has colength one in $B$. Hence
$A=A_0$. The explicit homomorphism
$\C\{u,v\}/(uv)\to A_0$ sends
$u$ to $(t_1,0)$, $v$ to $(0,t_2)$, and a constant
to the corresponding equal pair. Expanding a
convergent series shows that it is an isomorphism.
This is the ordinary node germ.

\item For $r=1$, an element of $A$ with valuation one would
be a local coordinate on the normalization disc by
the holomorphic inverse function theorem. Convergent
composition in the analytic algebra $A$ would then
give $\C\{t\}\subset A$, contrary to its colength
one. Therefore every nonconstant element of $A$ has
valuation at least two. We have
$A\subset\C+t^2\C\{t\}$; both rings have
colength one in $\C\{t\}$, so equality holds.
The map
\[
 \C\{u,v\}/(v^2-u^3)\longrightarrow A,
                 \qquad (u,v)\longmapsto(t^2,t^3)
\]
is an isomorphism: every series with no linear term
is uniquely the sum of an even series in $t^2$ and
$t^3$ times a series in $t^2$. This is the ordinary
cusp germ. A delta-zero germ is its smooth
normalization disc.

\item We finally check the embedding into the threefold $Y$. These analytic rings have embedding
dimension at most two. In local coordinates of the
smooth threefold, the tangent space of $D$ therefore
has codimension at least one. There is an element
$z$ of its ideal with $dz(x)\ne0$. The inverse
function theorem makes $z$ an ambient coordinate,
and the ideal of $D$ is the inverse image of a
height-one ideal in the regular analytic ring
$\C\{u,v\}=\OO_{Y,x}/(z)$. This ideal is principal:
a pure reduced curve in a smooth surface is a Cartier
divisor, or equivalently its reduced height-one ideal
is a product of the prime factors in this factorial
ring. Thus $I_{D,x}=(z,h)$, a regular sequence.
This proves the asserted lci embedding, including
at the node and cusp. Proposition~\ref{cdp:st:A.5.5} applies
to the three specified generic transverse types
and proves the final smooth elliptic and constant-type
conclusion.

\end{enumerate}
\end{proof}






\begin{proposition}[Eliminating vertical conductor on a smooth ruled cusp normalization]\hypertarget{tag-01B0}{}\label{cdp:st:A.5.10}
\StatementRecord{cdp:st:A.5.10}{proved}{01B0}{\SeriesRef{cdp:eq:A.J.26}, \SeriesRef{cdp:eq:A.J.23}, \SeriesRef{cdp:eq:A.J.21}, \SeriesRef{cdp:eq:A.J.27}, \SeriesRef{cdp:eq:A.J.22}, \SeriesRef{cdp:st:A.5.7}}

Let $Y$ be a smooth compact complex threefold with
$H^2(Y,\mathbb Q)=0$, let $S\subset Y$ be a reduced
irreducible divisor, and suppose its normalization
$h:Z\to S\subset Y$ is a smooth irrational
ruled surface $p:Z\to B$.
Suppose its conductor divisor is $D_Z=2C+V$, where
$C$ is a ruling section and $V$ is vertical, and
the generic transversal along $h(C)$ is an ordinary cusp.
Then $B$ is elliptic, $V=0$, $C^2=0$, and $Z$ is
already minimal.
The image $h(C)$ is smooth elliptic and $S$ is locally
an ordinary-cusp product along it. No local
complete-intersection hypothesis on the image curve is
required. The conductor-image scheme along the vertical
components may initially be nonreduced.
\end{proposition}


\begin{proof}

\begin{enumerate}[label={$\langle1\rangle$\arabic*.},leftmargin=*]
\item Let $g=g(B)\geq1$, let $F$ be the full ruling fibre
class, and write $V=\sum_i c_i\Gamma_i$ for its
irreducible vertical components. Put
$v=V\cdot C$, $n=C^2$ and $k=2g-2$.
Conductor duality and Cartier adjunction give the exact
line-bundle comparison
\[
 K_Z\otimes\mathcal O_Z(2C+V)
       =h^*(K_Y\otimes\mathcal O_Y(S)).
                                                        \tag{A.J.21}\hypertarget{tag-005P}{}\label{cdp:eq:A.J.21}
\]
Its right side has rational first Chern class zero.
Intersecting with $C$ and using adjunction on the section
gives
\[
 n=2-2g-v=-k-v.                                  \tag{A.J.22}\hypertarget{tag-005Q}{}\label{cdp:eq:A.J.22}
\]
Let $D_p$ be the divisorial common zero of
$dp:T_Z\to p^*T_B$. Its remaining zero scheme is
zero-dimensional, say $\Xi$, so the image is
$I_\Xi\otimes p^*T_B(-D_p)$.
The kernel $T_p$ is an invertible sheaf, by the depth
lemma. Taking determinants gives the actual comparison
\[
 T_p=K_Z^{-1}\otimes p^*K_B\otimes\mathcal O_Z(D_p).
\]
The section makes $dp$ surjective at every point of $C$.
Consequently $D_p$ and $\Xi$ are disjoint from $C$,
and $T_p|_C=N_{C/Z}$.

\item Put $\mathcal F=h^*T_Y$. Poincare duality on the compact
oriented six-manifold $Y$ also gives
$H^4(Y,\mathbb Q)=0$. Therefore
\[
 c_1(\mathcal F)=0,\qquad c_2(\mathcal F)=0
        \quad\hbox{in rational cohomology}.       \tag{A.J.23}\hypertarget{tag-005R}{}\label{cdp:eq:A.J.23}
\]
The differential is generically injective. Let
$L\subset\mathcal F$ be the saturation of the image
of $T_p\xrightarrow{dh}\mathcal F$.
It is invertible: a saturated rank-one subsheaf of a
locally free sheaf on a smooth surface is reflexive.
The map $T_p\to L$ has no vertical divisorial zero.
Indeed $h$ is finite, so no ruling fibre is contracted;
in characteristic zero its derivative along that fibre
is nonzero generically.

\item At the generic point of $C$, an ordinary cusp has
normalization differential equivalent to that of
$(u,s)\mapsto(u^2,u^3,s)$.
The ruling tangent either lies in its kernel along $C$
or does not. The vanishing order in the first case is
exactly one, and in the second case it is zero.
Write $\epsilon=1$ or $\epsilon=0$ for these respective
alternatives. There are no other horizontal conductor
components, and the normalization is an immersion away
from the conductor. Hence the actual line comparison is
\[
 L=T_p\otimes\mathcal O_Z(\epsilon C).
                                                        \tag{A.J.24}\hypertarget{tag-005S}{}\label{cdp:eq:A.J.24}
\]
This checks both positions of the derivative kernel
relative to the ruling tangent.

\item The quotient $E=\mathcal F/L$ is torsion-free of rank
two. The differential induces a generically nonzero map
\[
 T_Z/T_p=I_\Xi\otimes p^*T_B(-D_p)
                    \longrightarrow E\longrightarrow E^{**}.
\]
It extends uniquely to a map from $p^*T_B(-D_p)$ into
$E^{**}$, since both are locally free and the omitted
set has codimension two; locally this is Hartogs
extension of its coefficient functions.
Let $M$ be its saturated image line in the rank-two
bundle $E^{**}$. At the generic point of $C$, the
two-by-two differential minors have common vanishing
order one. Removing the order $\epsilon$ already
assigned to $L$ leaves order $1-\epsilon$ for this map.
At the generic point of $\Gamma_i$, take a plane transversal
to its image and a normalization parameter $u$ for the
corresponding plane-curve branch. If $m_i$ is its
multiplicity, its derivative modulo the nonzero tangent
along $\Gamma_i$ vanishes to order $r_i=m_i-1$.
These are all codimension-one zeros. Put
$R=\sum_i r_i\Gamma_i$ and $r=R\cdot C$. We obtain
\[
 M=p^*T_B(-D_p)\otimes
           \mathcal O_Z\bigl((1-\epsilon)C+R\bigr).
                                                        \tag{A.J.25}\hypertarget{tag-005T}{}\label{cdp:eq:A.J.25}
\]

\item The plane-curve conductor gives the inequalities
\[
 2R\leq V,\qquad 2r\leq v.                      \tag{A.J.26}\hypertarget{tag-005U}{}\label{cdp:eq:A.J.26}
\]
To spell out their local length content, on a single
plane branch of multiplicity $m_i$ the classes
$u,u^2,\ldots,u^{m_i-1}$ are linearly independent in
the normalization quotient. Its delta invariant is thus
at least $m_i-1$. The branch is Gorenstein. Finite
duality identifies the length of its ring modulo its
conductor with the length of its normalization quotient;
the intrinsic conductor exponent is consequently twice
that delta invariant. Other branches add their
nonnegative intersection multiplicities to this exponent.
For the last comparison, writing the plane equation
as $\prod_j f_j$ multiplies the residue denominator on
branch $i$ by $\prod_{j\ne i}f_j$, whose valuation is
exactly the sum of those intersection multiplicities.
This proves \eqref{cdp:eq:A.J.26} for every branch, including cases
with nonreduced conductor-image scheme.

\item Let $\ell_1$ be the length of $E^{**}/E$. Saturation
of $M$ gives a zero-dimensional subscheme $\Xi_M$ and a
line bundle $N$ with
\[
 0\longrightarrow M\longrightarrow E^{**}
       \longrightarrow I_{\Xi_M}\otimes N\longrightarrow0.
\]
Write $\ell_2=\operatorname{length}\Xi_M$.
First Chern classes give $c_1(N)=-c_1(L)-c_1(M)$.
The second Chern classes in these typed quotient
sequences and \eqref{cdp:eq:A.J.23} give the intersection identity
\[
 L^2+L\cdot M+M^2=\ell_1+\ell_2\geq0.            \tag{A.J.27}\hypertarget{tag-005V}{}\label{cdp:eq:A.J.27}
\]
Indeed $c_2(E)=L^2$,
$c_2(E^{**})=L^2-\ell_1$, and
$c_2(E^{**})=-L\cdot M-M^2+\ell_2$.

\item Set $W=V+D_p$ and $B_v=R-D_p$, both vertical
divisors, the first effective. The exact comparisons
above and \eqref{cdp:eq:A.J.21} give the numerical classes
\[
 L\equiv(2+\epsilon)C+W+kF,\qquad
 M\equiv(1-\epsilon)C+B_v-kF.
\]
The intersection form on vertical divisors is negative
semidefinite. One can check this directly by contracting
the ruling blowups: every vertical class is a multiple
of the pullback fibre class plus exceptional classes,
on which the form is negative definite. Hence
\[
 q_v=W^2+W\cdot B_v+B_v^2
      =(W+B_v/2)^2+3B_v^2/4\leq0.
\]
Its vanishing implies $W^2=B_v^2=0$; these classes
are then numerical multiples of the fibre class.

\item For $\epsilon=1$, using $D_p\cdot C=0$ in the
left side of \eqref{cdp:eq:A.J.27} gives
\[
 \ell_1+\ell_2=-6k-3v+3r+q_v
                  \leq-6k-3v/2.
\]
Nonnegativity, with $k\geq0$ and $v\geq0$,
forces $g=1$, $v=0$, $r=0$ and $q_v=0$.

\item For $\epsilon=0$, the same intersection calculation reads
\[
 0\leq\ell_1+\ell_2
       =-6k-2v+4r+q_v\leq-6k.                   \tag{A.J.28}\hypertarget{tag-005W}{}\label{cdp:eq:A.J.28}
\]
It forces $g=1$, $\ell_1=\ell_2=0$, $2r=v$,
and $q_v=0$.
Thus $E=E^{**}$ is locally free, and
$L\to\mathcal F$ is a subbundle inclusion.
Along $C$, the composite
$T_C\to\mathcal F|_C\to E|_C$ vanishes generically,
because $dh$ has rank one there and its nonzero ruling
tangent image is $L|_C$. It therefore vanishes
everywhere. The exact subbundle sequence gives an
actual nonzero holomorphic line-bundle map
\[
 T_C\longrightarrow L|_C=T_p|_C=N_{C/Z}.
                                                        \tag{A.J.29}\hypertarget{tag-005X}{}\label{cdp:eq:A.J.29}
\]
It is nonzero since the finite map $C\to h(C)$ has
nonzero derivative generically. Its Hom line has degree
$n-(2-2g)=-v$ by \eqref{cdp:eq:A.J.22}. A nonzero section on a
compact curve has nonnegative degree, so $v=0$ in this
alternative too.

\item In both alternatives, $v=r=0$ and $q_v=0$.
Thus $W$ is numerically a fibre multiple and
$W\cdot C=V\cdot C+D_p\cdot C=0$, so it is
numerically zero. An effective nonzero divisor on the
projective ruled surface has positive intersection
with an ample class. It follows that $W=0$ and hence
$V=D_p=0$. Now $B_v=R$ is effective, has square zero
and has zero intersection with $C$, so the same
argument gives $R=0$. Equation \eqref{cdp:eq:A.J.22} gives $C^2=0$.
Finally $K_Z+2C$ is numerically trivial. A ruling
$(-1)$-curve $E_0$ would satisfy
$0=(K_Z+2C)\cdot E_0=-1+2C\cdot E_0$, impossible
because $C\cdot E_0$ is an integer. There are therefore
no ruling blowups and $Z$ is minimal.

\item Statement \ref{cdp:st:A.5.7} now applies everywhere along
the conductor: the normalization is smooth and its
conductor is exactly $2C$. It makes the image smooth
and the surface locally the cusp product. The image is
elliptic because $C\to h(C)$ is an isomorphism.

\end{enumerate}
\end{proof}






\begin{proposition}[A noncontracted horizontal cusp has zero Euler number]\hypertarget{tag-01BB}{}\label{cdp:st:A.5.20}
\StatementRecord{cdp:st:A.5.20}{proved}{01BB}{\SeriesRef{cdp:st:A.5.14}, \SeriesRef{cdp:st:A.5.5}, \SeriesRef{cdp:st:C2d}, \SeriesRef{cdp:st:A.5.18}}

Let $f:Y\to C_0$ be a proper holomorphic map from a smooth
compact complex threefold with $b_2(Y)=0$, with connected
fibres and smooth general fibre a complex two-torus.
Let $mS$ be a full fibre with irreducible reduced support,
$\eta:T\to S$ its normalization and $\pi:Z\to T$ an
irrational ruled resolution with ruling $p:Z\to B$.
Suppose every exceptional curve of $\pi$ is vertical
and the resolution conductor is $2C+V$, where $C$
is a noncontracted section and $V$ is vertical.
If the generic transversal at $D=\eta\pi(C)$ is an
ordinary cusp, then $D$ and $C$ are smooth elliptic
curves and $C\to D$ is an isomorphism. The cusp type
is constant along $D$. The normalization $T$ is a
smooth minimal ruled surface over this elliptic curve,
its conductor is twice a section, and $e(S)=0$.
\end{proposition}


\begin{proof}
\begin{enumerate}[label={$\langle1\rangle$\arabic*.},leftmargin=*]
\item The section identifies $C$ with the positive-genus curve $B$.
The proper nonconstant map $C\to D$ is finite. At a general point of
$D$, the cusp normalization
\[
 (u,z)\longmapsto(u^3,u^2,z)
\]
has one branch and degree one on its singular curve. Thus $C\to D$ is
the normalization of $D$.
Proposition~\ref{cdp:st:A.5.14}, applied to the same $Y,S,D$, makes
$D$ smooth elliptic with constant cusp transversal. It follows that
$C\to D$ is an isomorphism and $B$ is elliptic.
Let $C_T\subset T$ be the inverse-image curve.
Proposition~\ref{cdp:st:C2d} gives a smooth neighbourhood of $C_T$,
the isomorphism $C_T\to D$, and
$\deg N_{C_T/T}=0$. Its leading normal terms identify the square and
cube of this normal line with the radical and Hessian quotient lines.

\item Factor the given resolution as
\[
 Z\xrightarrow{\sigma}Z'\xrightarrow{\mu}T,
 \qquad \pi=\mu\sigma,
\]
where $\mu$ is the minimal resolution and $\sigma$ is a sequence of
point blowups. Every $\sigma$-exceptional curve is $\pi$-exceptional
and is vertical for $p$. The map $p$ therefore descends to
$p':Z'\to B$. Explicitly, the proper image of $(\sigma,p)$ in
$Z'\times B$ projects finitely and bimeromorphically to the normal
surface $Z'$ with one point in every fibre. This projection is an
isomorphism and defines $p'$, with $p=p'\sigma$.
The curve $C'=\sigma(C)$ is a section and $\sigma|_C$ is an isomorphism.
Since $\mu$ is an isomorphism near $C_T$, one has
$(C')^2=\deg N_{C_T/T}=0$.

The strict transform of every $\mu$-exceptional curve is
$\pi$-exceptional, so those curves are vertical for $p'$.
Proposition~\ref{cdp:st:A.5.3} supplies the canonical effective divisor
on this minimal resolution. Its nonexceptional components have the
same conductor valuations as on $Z$. Consequently it has the form
$2C'+V'$, with $V'$ effective and vertical, and
\[
 \omega_{Z'}(2C'+V')\simeq(\eta\mu)^*\omega_S,
 \qquad c_1(K_{Z'}+2C'+V')=0\quad\hbox{in }H^2(Z',\Q).
\]
The constant cusp charts and the isomorphism of $\mu$ there show that
$V'$ is disjoint from $C'$.

\item Contract the vertical $(-1)$-curves of the irrational ruled
surface to obtain $\rho:Z'\to Z_0$, a $\mathbf P^1$-bundle over $B$.
Every rational curve is vertical because $B$ has positive genus.
Write $C_0=\rho(C')$ and $V_0=\rho_*V'$.
The image $C_0$ remains a section. The point-blowup formula for its
self-intersection gives $C_0^2=(C')^2+n=n\ge0$, where $n$ counts the
centres lying on the successive strict transforms of the section.
The divisor $V_0$ is an effective sum of fibres.
Pushing the canonical equality forward and intersecting with $C_0$
gives, by adjunction on the elliptic section,
\[
 0=(K_{Z_0}+2C_0+V_0)\cdot C_0
   =C_0^2+V_0\cdot C_0=n+V_0\cdot C_0.
\]
Thus $n=0$ and $V_0=0$.

For completeness, the conductor transform at a point blowup
$\beta:X\to X_0$ is determined by
$K_X=\beta^*K_{X_0}+E$: its exceptional coefficient is
$\operatorname{mult}_p(D_0)-1$. This is the same canonical divisor
identity used above. A ruling $(-1)$-curve meeting $C'$ cannot be a
component of $V'$, since $C'\cap V'=\varnothing$; its intersection
with $K_{Z'}+2C'+V'$ would be at least $-1+2=1$.

Since $n=0$, one has $C'=\rho^*C_0$.
All components of $V'$ are $\rho$-exceptional, and
\[
 K_{Z'}=\rho^*K_{Z_0}+E_\rho
\]
with $E_\rho$ effective and positive on every exceptional prime.
Hence $E_\rho+V'$ is numerically trivial and exceptional.
The total transforms of the successive exceptional curves of a
point-blowup sequence give an orthogonal basis with square $-1$.
The exceptional intersection form is therefore negative definite.
It follows that $E_\rho+V'=0$. Effectiveness gives
$E_\rho=V'=0$, so $\rho$ is an isomorphism and $Z'$ is the minimal
ruled bundle.

\item A vertical irreducible curve of this bundle is a fibre and has
square zero. To compare it with the exceptional curves of $\mu$,
use the local-function form of the intersection argument in
\cite[\S1, p.~6]{Mumford1961}.
For an exceptional fibre over a normal point, choose a nonzero
holomorphic function $g$ vanishing at that point. On the resolution
its divisor is $G+\sum m_iE_i$, with $m_i>0$ and $G$ the strict
transform of the zero curve. Pairing this principal divisor with
$E_j$ gives
$\sum_i m_i(E_i\cdot E_j)=-G\cdot E_j\le0$, with a strict inequality
for some $j$. The strict transform meets the exceptional fibre by
properness, and that fibre is connected because the target is normal.
For $A_{ij}=m_im_j(E_i\cdot E_j)$, its off-diagonal entries are
nonnegative and
\[
 x^{\mathsf T}Ax
 =\sum_i\left(\sum_jA_{ij}\right)x_i^2
     -\sum_{i<j}A_{ij}(x_i-x_j)^2.
\]
A vector in its zero locus has zero coordinate at a strict row, and
connectedness forces all coordinates to vanish. The exceptional
intersection form is negative definite. This argument uses the
displayed local holomorphic divisor and applies to the analytic
surface $T$.
Thus a square-zero fibre cannot be $\mu$-exceptional. The map $\mu$
has no exceptional curves, hence is finite and bimeromorphic; the
normality of $T$ makes it an isomorphism. Therefore $T=Z'$ is smooth
and minimal ruled over the elliptic curve $B$.

\item Finite duality on the smooth normalization identifies the
conductor ideal with the invertible ideal of its divisor, now
$2C_T$. Thus $\eta$ is an isomorphism off $C_T$, and its restriction
$C_T\to D$ is an isomorphism. The normalization is consequently a
proper continuous bijection and a homeomorphism of the underlying
compact spaces. Equivalently, additivity of Euler characteristic
gives
\[
 e(S)=e(T)-e(C_T)+e(D)=0.
\]
Here $e(T)=e(\mathbf P^1)e(B)=2\cdot0$, and both curves have genus one.
The full scheme fibre $mS$ has the same underlying space as $S$,
so its topological Euler characteristic is also zero.
\end{enumerate}
\end{proof}






\begin{proposition}[Hirzebruch germs cannot occur in a torus degeneration]\hypertarget{tag-01B5}{}\label{cdp:st:A.5.15}
\StatementRecord{cdp:st:A.5.15}{proved}{01B5}{\SeriesRef{cdp:eq:A.J.33a}, \SeriesRef{cdp:eq:A.J.33b}, \SeriesRef{cdp:eq:A.J.35}, \SeriesRef{cdp:eq:A.J.34}, \SeriesRef{cdp:eq:A.J.37}; \cite[Theorem L1 and Example 2.4(a)]{Nemethi1991Join}; \cite[Theorem 4.1 and Lemma 4.6]{MPW2005Hirzebruch}; \cite[Theorem 4.2(2)]{MPW2005Hirzebruch}}

Let $f:Y\to C_0$ be a proper holomorphic map from a
smooth complex threefold to a smooth curve, with smooth
nearby fibres complex two-tori. There is no point of
a fibre at which, in holomorphic ambient coordinates
$(x,y,z)$ and a coordinate on the base, the map has
the equation
\[
 f(x,y,z)=u(x,y,z)\bigl(x^a-y^bz^c\bigr)^m,
 \quad a\ge2,\quad b,c,m\ge1,                                    \tag{A.J.33}\hypertarget{tag-0062}{}\label{cdp:eq:A.J.33}
\]
where $u$ is a holomorphic unit. No Betti-number
assumption on $Y$ is needed. This includes the pinch
point, every germ $x^2-y^{2k}z^{2r+1}$ with
$k\ge1$, $r\ge0$, and the reducible germ
$x^2-y^2z^2$.
\end{proposition}


\begin{proof}

\begin{enumerate}[label={$\langle1\rangle$\arabic*.},leftmargin=*]
\item Put $h=x^a-y^bz^c$. We first compare the actual map
with its reduced equation. This polynomial is weighted
homogeneous of degree $a(b+c)$ for weights
$(b+c,a,a)$. On a small simply connected coordinate
neighbourhood choose a holomorphic logarithm of $u$.
Define the holomorphic coordinate change
\[
 \Phi(x,y,z)=\left(
 x u^{1/(ma)},\;
 y u^{1/(m(b+c))},\;
 z u^{1/(m(b+c))}\right).
                                                        \tag{A.J.34}\hypertarget{tag-0065}{}\label{cdp:eq:A.J.34}
\]
Its derivative at the origin is an invertible diagonal
matrix. The holomorphic inverse function theorem and
weighted homogeneity give
$h\circ\Phi=u^{1/m}h$ and therefore
$f=(h\circ\Phi)^m$ exactly. For a sufficiently small
nonzero regular value $t$, a chosen Milnor fibre of
$h$ at one of the $m$ roots of $t$ is consequently a
compact oriented submanifold with boundary of the
smooth global fibre $Y_t$. Denote this Milnor fibre
by $M$.

\item We first give a direct spherical cycle, without an
irreducibility assumption. Put $d=\gcd(b,c)$,
$b'=b/d$ and $c'=c/d$, and consider the smooth affine
curve
\[
 B=\{(x,v)\in\C^2:x^a-v^d=1\}.
\]
The affine surface $H=\{h=1\}$ maps holomorphically
to $B$ by $v=y^{b'}z^{c'}$. The curve $B$ is
connected: projection to the $v$-line has $a$ sheets,
and a small loop about any root of $v^d+1$ cyclically
permutes them. It is smooth because its two partial
derivatives cannot vanish on $x^a-v^d=1$.
The function $v$ has $a$ distinct simple zeros
$P_j=(\zeta_a^j,0)$ on $B$.

\item Choose two of these zeros, $P$ and $Q$, and an
embedded arc $\gamma:[0,1]\to B$ joining them and
avoiding the other zeros. Choose smooth complex
functions $A(t),B_0(t)$ on its interior such that
\[
 A(t)^{b'}B_0(t)^{c'}=v(\gamma(t)),\qquad
 |A(t)|=|v(\gamma(t))|^{1/(2b')},\quad
 |B_0(t)|=|v(\gamma(t))|^{1/(2c')}.
\]
They exist by choosing a continuous argument along
the arc. The continuous map
\[
 (t,\theta)\longmapsto
 \bigl(x(\gamma(t)),A(t)e^{\mathrm i c'\theta},
                 B_0(t)e^{-\mathrm i b'\theta}\bigr)
                                                        \tag{A.J.33a}\hypertarget{tag-0063}{}\label{cdp:eq:A.J.33a}
\]
from $[0,1]\times S^1$ to $H$ collapses each boundary
circle to the corresponding point $P$ or $Q$.
It therefore defines an oriented spherical cycle
$\sigma_\gamma:S^2\to H$. Coprimality of $b'$ and
$c'$ makes each interior circle embedded. The whole
map is injective away from the collapsed ends,
because the base arc is embedded.

\item Let $\gamma'$ be a nearby arc with the same endpoints
and disjoint interior, obtained by an isotopy in a
thin disc neighbourhood of $\gamma$. The isotopy
keeps all other zeros of $v$ outside that disc.
Choose the arguments in \eqref{cdp:eq:A.J.33a} continuously during
the isotopy; both coordinates tend uniformly to zero
at the ends. This gives a continuous homotopy of
maps from $S^2$, so the resulting spherical cycles
are homologous. Their supports meet only at $P$ and
$Q$. Near either endpoint, $(y,z)$ are actual
holomorphic coordinates on $H$, since $a x^{a-1}$
is nonzero there. Arrange the two arcs to be radial
with different arguments in the local $v$-coordinate.
On a sufficiently small sphere in these coordinates,
the links of the two cycles are disjoint parallel
curves of slope $(c',-b')$ on a Clifford torus.
Their oriented linking number is $-b'c'$.
Indeed the torus framing of a $(p,q)$ curve has
linking number $pq$ with its parallel copy: push the
two curves to the opposite sides of the Heegaard
splitting of $S^3$ into its two solid tori and count
meridian discs, obtaining $p q$ signed intersections.
Here $p=c'$ and $q=-b'$.

\item The local intersection of the two coned cycles is
that linking number. This equality follows directly
by replacing the two cones inside the small ball by
transverse relative two-chains with the same boundary
links; their intersection is the defining linking
number in its boundary three-sphere. At $Q$ both
cycle orientations reverse relative to the outgoing
radial coordinate, so the two reversals leave the
local intersection sign unchanged. Consequently
\[
 [\sigma_\gamma]^2
   =[\sigma_\gamma]\cdot[\sigma_{\gamma'}]
   =-2b'c'\ne0.                                  \tag{A.J.33b}\hypertarget{tag-0064}{}\label{cdp:eq:A.J.33b}
\]
Both cycles and their homotopy have compact support.
The weighted homothety
$(x,y,z)\mapsto(s^{b+c}x,s^a y,s^a z)$ maps
$H$ biholomorphically onto $\{h=s^{a(b+c)}\}$.
For sufficiently small positive $s$, it puts those
compact supports into any prescribed Milnor
neighbourhood of the origin. Thus \eqref{cdp:eq:A.J.33b} is a
spherical class of nonzero square in an actual
local Milnor fibre $M$.

\item For an independent check when $\gcd(a,b,c)=1$,
we compute the full homotopy type and the rank of
the intersection form from two primary results.
The join theorem
\cite[Theorem L1 and Example 2.4(a)]{Nemethi1991Join}
applies to arbitrary analytic germs, including
non-isolated ones. Set $d=\gcd(b,c)$. The Milnor
fibre of $y^bz^c$ is a disjoint union of $d$ annuli:
its equations in a Milnor polydisc are
$y^{b/d}z^{c/d}=\zeta\epsilon^{1/d}$, and each
component is parametrized by
$y=A w^{c/d}$, $z=B w^{-b/d}$. Coprimality makes
this parametrization injective, and the two radius
bounds give a nonempty annulus in $w$.
The Milnor fibre of $x^a$ consists of $a$ points.
The join theorem thus identifies $M$, by a homotopy
equivalence, with the union of $a$ cones on a
disjoint union of $d$ circles, all glued along their
common base. Contracting one cone and then choosing
arcs between the $d$ base components gives
\[
 M\simeq
 \bigvee^{(a-1)d}S^2\ \vee\
 \bigvee^{(a-1)(d-1)}S^1.                          \tag{A.J.35}\hypertarget{tag-0066}{}\label{cdp:eq:A.J.35}
\]
In particular, every class in $H_2(M,\Z)$ is a sum
of images of spherical fundamental classes.

\item For the same irreducible germ, the boundary calculation
of \cite[Theorem 4.1 and Lemma 4.6]{MPW2005Hirzebruch}
gives a Seifert three-manifold with base genus
\[
 g_\partial=(a-1)(d-1),\qquad
 e_0=\frac{a}{(b/d)(c/d)}>0.                     \tag{A.J.36}\hypertarget{tag-0067}{}\label{cdp:eq:A.J.36}
\]
Here $e_0$ is its rational Seifert Euler number in
the orientation convention of that paper. Its
nonzero value gives
$b_1(\partial M)=2g_\partial$. To make this last
deduction explicit, the rational abelianized Seifert
presentation has a regular-fibre generator $v$,
exceptional-fibre generators $q_i$, the relations
$\alpha_iq_i+\beta_iv=0$ and
$\sum_iq_i+ev=0$, and $2g_\partial$ free base
generators. Eliminating the $q_i$ gives $e_0v=0$.
Thus all fibre generators vanish over $\Q$, leaving
exactly the $2g_\partial$ base generators.

\item Let $I_M:H_2(M,\Q)\to H_2(M,\partial M;\Q)$ be
the intersection morphism. Poincare--Lefschetz duality
and the homology sequence of the pair give
\[
 \operatorname{rank} I_M
   =b_2(M)+b_1(M)-b_1(\partial M)=a-1>0.          \tag{A.J.37}\hypertarget{tag-0068}{}\label{cdp:eq:A.J.37}
\]
Indeed $H_3(M)=0$ by \eqref{cdp:eq:A.J.35}, so
$H_3(M,\partial M)\simeq H^1(M)$ injects into
$H_2(\partial M)$. Its cokernel is the kernel of
$I_M$, which has dimension
$b_1(\partial M)-b_1(M)$. This proves the displayed
rank formula. For $a=2$, $b=1$, $c=\ell$, the
more precise boundary is $L(2\ell,1)$ by
\cite[Theorem 4.2(2)]{MPW2005Hirzebruch}; the same
pair sequence over $\Z$ makes the rank-one
intersection form have absolute determinant
$2\ell$. The special case $\ell=1$ is the ordinary
quadratic singularity and has determinant two.

\item Finally let $j:M\hookrightarrow Y_t$ be the actual
embedding obtained from \eqref{cdp:eq:A.J.34}. A complex two-torus
has contractible universal cover $\C^2$, so
$\pi_2(Y_t)=0$. In particular the spherical class
constructed in \eqref{cdp:eq:A.J.33a} maps to zero in
$H_2(Y_t,\Z)$. Intersection numbers of compact
two-cycles in the interior of $M$ are preserved by
this oriented codimension-zero embedding: the same
small transverse perturbations compute them in
both manifolds. Its square in $Y_t$ is therefore
both zero and $-2b'c'$, contradicting \eqref{cdp:eq:A.J.33b}.
In the irreducible case, \eqref{cdp:eq:A.J.35}--\eqref{cdp:eq:A.J.37} provide
an additional contradiction on the whole intersection
space.
This proves the exclusion for the actual map,
including every multiplicity in \eqref{cdp:eq:A.J.33}.

\end{enumerate}
\end{proof}



The homotopy and boundary theorems just used are
established results; our application is the displayed
embedding and intersection comparison. The explicit
cycle also treats the reducible monomials. The exact
Seifert arithmetic, homology dimensions and ranks
are checked in \path{code/E_Hirzebruch.py}. No
reduction of an arbitrary nonordinary surface germ
to \eqref{cdp:eq:A.J.33} is asserted.




\begin{proposition}[Isolated singularities are excluded and suspension germs are transverse products]\hypertarget{tag-01B7}{}\label{cdp:st:A.5.17}
\StatementRecord{cdp:st:A.5.17}{proved}{01B7}{\SeriesRef{cdp:st:A.5.15}, \SeriesRef{cdp:eq:A.J.40}, \SeriesRef{cdp:eq:A.J.39}; \cite[Lemma 6.1 and Appendix B]{Milnor1968}}

Let $f:Y\to C_0$ have the hypotheses of \ref{cdp:st:A.5.15}.
No local equation $f=u h^m$, with $u$ a holomorphic
unit and $m\ge1$, can have a reduced equation $h$
with an isolated critical point at the origin.
Suppose its local equation at a point of a fibre is
\[
             f=(x^a-g(y,z))^m,
             \qquad a\ge2,\quad m\ge1,             \tag{A.J.38}\hypertarget{tag-0069}{}\label{cdp:eq:A.J.38}
\]
where $g\ne0$ and $g(0,0)=0$. Then the reduced
plane germ $\{g=0\}_{\rm red}$ is smooth. If
$b$ is its multiplicity in $\operatorname{div}(g)$,
there are local holomorphic coordinates in which
$f=(x^a-y^b)^m$.
Consequently, at every point where the full fibre
is a uniform multiple of its reduced support and
the reduced equation has a nonzero Hessian, its
critical surface germ is a constant
transverse $A_{b-1}$ product. A nonordinary change
of transverse type at a point with nonzero Hessian
is excluded.
\end{proposition}


\begin{proof}

\begin{enumerate}[label={$\langle1\rangle$\arabic*.},leftmargin=*]
\item We first prove the isolated assertion, including its
comparison with the actual global fibre. Replace $h$
by $H=u^{1/m}h$, taking a holomorphic root on a small
coordinate neighbourhood. Then $f=H^m$ exactly.
This replacement preserves an isolated critical point
as a germ. Indeed every local irreducible component
of the critical locus of $H$ has constant $H$-value,
necessarily zero when it contains the origin. On
$H=0$ one has $dH=u^{1/m}dh$. Thus its critical
locus has the same zero-dimensional germ as that
of $h$.

\item Choose a small closed ball $K\subset\C^3$ such that
$H^{-1}(0)$ is smooth along its boundary intersection
and transverse to $\partial K$.
The standard isolated-singularity Milnor-ball
construction supplies this choice; see
\cite[Lemma 6.1 and Appendix B]{Milnor1968}.
The restriction $H|_{\partial K}$ is a submersion
at its zero set. Compactness therefore supplies
$\eta>0$ for which it is a submersion wherever
$|H|<2\eta$ on the boundary. For the linear
deformation
\[
 H_\lambda(w)=H(w)-\lambda_1w_1
                         -\lambda_2w_2-\lambda_3w_3,
                                                        \tag{A.J.38a}\hypertarget{tag-006A}{}\label{cdp:eq:A.J.38a}
\]
this boundary submersion persists for all parameters
in a sufficiently small polydisc $P$ and all values
of modulus less than $\eta$. The gradient map of
$H$ has an isolated zero and consequently is a finite
holomorphic map germ of positive local degree. Choose
a small regular value $\lambda$ of that map. Its
nonempty finite inverse image consists of critical
points of $H_\lambda$ with invertible Hessian. All
these points and their critical values tend to the
origin as $\lambda\to0$.

\item Remove the critical-value discriminant from
$P\times\Delta_\eta$. The relative critical locus
over this domain misses the boundary, so its
projection is proper and its image is a proper
complex analytic subset. Its complement is path
connected. The projection of the corresponding
smooth level sets in $K$ is a proper submersion
with boundary, and hence an Ehresmann bundle.
At any one of the nondegenerate critical points,
the holomorphic Morse lemma gives
$H_\lambda-H_\lambda(p)=X^2+Y^2+Z^2$.
A local Milnor fibre at a sufficiently small regular
value embeds in the corresponding level set in $K$.
The bundle transports it, by an oriented embedding,
into $K\cap H^{-1}(\epsilon_0)$ for a nonzero small
$\epsilon_0$, which is a submanifold of the actual
global torus $Y_{\epsilon_0^m}$.

\item For a positive real local Milnor value, the real
locus $X,Y,Z\in\mathbb R$ is a two-sphere. It is
Lagrangian for the standard K\"ahler form; multiplication
by $\mathrm i$ identifies its tangent bundle with
its normal bundle with the opposite oriented Euler
class. Its self-intersection is therefore
$-\chi(S^2)=-2$. Equivalently this is the case
$a=2$, $b=c=1$ of the explicit cycle in \ref{cdp:st:A.5.15},
after the linear change $Y^2+Z^2=(Y+\mathrm iZ)(Y-\mathrm iZ)$.
The transported sphere has the same square in the
actual torus, whose $\pi_2$ is zero. This is the
required contradiction. No suspension representation
of the original isolated singularity has been used.

\item We prove the comparison of Milnor fibres needed for
the argument. Factor the analytic plane function as
\[
                   g=u\prod_{i=1}^r h_i^{b_i},
\]
where the $h_i$ are distinct reduced irreducible
germs and $u$ is a unit. Deform each $h_i$
holomorphically to $h_{i,s}$ in a finite-dimensional
parameter polydisc $P$, and retain its exponent:
\[
 g_s=u\prod_i h_{i,s}^{b_i},\qquad
 H_s=x^a-g_s.                                    \tag{A.J.39}\hypertarget{tag-006B}{}\label{cdp:eq:A.J.39}
\]
We use only sufficiently small parameters.

\item Choose a small closed ball $K\subset\C^2_{y,z}$
such that each $h_i=0$ meets $\partial K$
transversely in its smooth part, and their boundary
links are disjoint. Choose its representative so
that every critical value of $g$ on $K$ is zero.
This is possible because the critical locus of a
holomorphic function has finitely many local
irreducible components, the function is constant
on each of them, and each component of the germ
contains the origin. The boundary smoothness,
transversality and disjointness persist for
$h_{i,s}$ by compactness and the implicit function
theorem. Near any of those links, the other factors
are units, so there are parameter-dependent local
coordinates $(v,w)$ with $g_s=v^{b_i}$, obtained
by taking a root of that unit. The boundary
$\partial K$ is transverse to the $w$ direction.
Consequently the restriction of $g_s$ to
$\partial K$ is a submersion at every point where
$0<|g_s|<\eta_0$, for a uniform sufficiently small
$\eta_0>0$. Away from these link neighbourhoods
$|g_s|$ has a uniform positive lower bound on the
boundary. This proves the asserted uniform
transversality, including at parameters for which
the germs inside $K$ change their singularities.

\item Take a small disc $D_x$ of radius $\rho$ with
$\rho^a<\eta_0/2$, then shrink $P$ so that every
critical value of $g_s$ on $K$ has modulus less than
$\rho^a/4$. This shrinking follows by compactness:
a sequence of critical points with parameters
tending to zero would otherwise limit to a
nonzero critical value of $g$ in the chosen small
representative. The latter has no such value.
Fix $0<\eta<\rho^a/4$ and put
$W=\overline D_x\times K$, a compact manifold
with corners contained in the original coordinate
chart in $Y$.

\item Every $H_s^{-1}(\epsilon)$ with
$0<|\epsilon|<\eta$ is transverse to this boundary.
On the face $\partial D_x\times K$, the equation
makes $|g_s|>3\rho^a/4$, so $dg_s\ne0$; its plane
directions give the submersion. On the face
$\overline D_x\times\partial K$, the $x$
derivative gives it when $x\ne0$. At $x=0$ we
have $g_s=-\epsilon\ne0$, and the preceding
boundary calculation gives it. At the corner
$|x|=\rho$, that same calculation for $g_s$
supplies both real target directions. The corner
stratum is therefore also transverse. We use
Ehresmann's theorem with corners: local lifts of
base tangent vectors can be chosen tangent to
each boundary face by these submersion statements,
then patched by a partition of unity and
integrated on the compact fibres. Thus no change
of the chosen domain $W$ is required.

\item Inside $P\times\{0<|\epsilon|<\eta\}$ remove the
critical-value discriminant $\mathcal D$ of the
maps $H_s$ on $W$. This is a proper complex
analytic subset. The relative critical locus has
no point on the boundary at such a value, so its
projection is proper over every compact subset
of this parameter domain, and the proper mapping
theorem applies. For fixed $s$ its critical
values are discrete, since a holomorphic function
is constant on each connected critical component.
The complement of a proper complex analytic
subset of a connected complex manifold is path
connected. The transversality just proved and
the implicit function theorem give a proper
smooth submersion, including its boundary,
\[
 \{(s,p):p\in W,\ (s,H_s(p))\in
                    (P\times\Delta_\eta^*)\setminus\mathcal D\}
 \longrightarrow
       (P\times\Delta_\eta^*)\setminus\mathcal D,
 \quad (s,p)\longmapsto(s,H_s(p)).                 \tag{A.J.40}\hypertarget{tag-006D}{}\label{cdp:eq:A.J.40}
\]
Ehresmann's theorem identifies all its compact
fibres by orientation-preserving diffeomorphisms
of manifolds with corners.
In particular a sufficiently small local Milnor
fibre of an interior singularity of some $H_s$,
at a value near zero, embeds into a smooth fibre
in \eqref{cdp:eq:A.J.40}, hence into
$W\cap H_0^{-1}(\epsilon_0)$ for a sufficiently
small nonzero $\epsilon_0$. This last space is
an actual submanifold of the global torus
$Y_{\epsilon_0^m}$ by \eqref{cdp:eq:A.J.38}; no identification
of it with a standard Milnor fibre is needed.
Choose the value to avoid every other
critical value, and choose its local Milnor ball
in the interior of $W$. This is the precise
comparison we use. Arbitrary deformations which
change the exponents along the boundary links
are not used in \eqref{cdp:eq:A.J.40}.

\item Suppose first that $r\ge2$. Choose two distinct
factors $h_i,h_j$. Their map
$(h_i,h_j):(\C^2,0)\to(\C^2,0)$ is finite,
since their common zero germ is zero dimensional.
Its positive local degree is their intersection
multiplicity. A sufficiently small regular value
$(s_i,s_j)$ consequently has nonempty finite
inverse image. At each inverse-image point,
$h_i-s_i$ and $h_j-s_j$ are transverse coordinates.
Choose the remaining small shifts to avoid those
finitely many points. The deformation \eqref{cdp:eq:A.J.39}
then has an interior point at which
\[
            H_s=x^a-U v^{b_i}w^{b_j},
\]
with $U$ a unit independent of $x$. Replacing
$v$ by $vU^{1/b_i}$ makes this exactly
$x^a-v^{b_i}w^{b_j}$ by an invertible holomorphic
coordinate change. \ref{cdp:st:A.5.15} constructs a spherical
class of nonzero square in its local Milnor fibre.
The embedding \eqref{cdp:eq:A.J.40} puts that class in the
actual smooth torus fibre of $f=H_0^m$.
Preservation of intersection
numbers and $\pi_2$-vanishing give the same
contradiction as in \ref{cdp:st:A.5.15}. Thus $r=1$.

\item Suppose its reduced factor $h=h_1$ is singular.
Write $g=u h^b$ and absorb a local $b$th root of
$u$ into $h$. This leaves its reduced zero germ
unchanged. Since $h$ is reduced, its gradient
has an isolated zero at the origin: a critical
curve would have constant value zero and would
be a repeated component of $h$. The finite
gradient map has a positive local degree.
Choose a small regular value
$(\lambda_1,\lambda_2)$ of this map. Then
\[
        h_s=h-\lambda_1y-\lambda_2z-\lambda_0
\]
has nondegenerate critical points near zero;
regularity of the gradient value says exactly
that their Hessians are invertible. Choose
$\lambda_0$ so that their critical values are
nonzero. At one such point $p$ put
$v=h_s(p)\ne0$. The Hessian of $h_s^b$ is
$b v^{b-1}$ times the Hessian of $h_s$, and
is invertible. The holomorphic Morse lemma on
the two plane variables identifies the germ of
$H_s-H_s(p)$ with
\[
                         x^a-Y^2-Z^2.
\]
A linear change of $(Y,Z)$ gives the Hirzebruch
germ $x^a-YZ$, already excluded by \ref{cdp:st:A.5.15}.
Its Milnor fibre embeds by \eqref{cdp:eq:A.J.40}, producing
the same contradiction. Hence $h$ is smooth.
Taking $y=h u^{1/b}$ as a coordinate proves the
asserted transverse product.

\item For the Hessian consequence, take a local $m$th
root of the actual map, absorbing its unit.
Its reduced equation has the same nonzero
Hessian rank as every other reduced equation
at the critical point. The holomorphic splitting
lemma splits off a nondegenerate square and
gives $f=(x^2-g(y,z))^m$ exactly. If $g=0$,
the equation would not be reduced, so $g\ne0$.
At a critical point $dg(0)=0$. The result just
proved gives $g=y^b$ with $b\ge2$ in suitable
holomorphic coordinates, which is the constant
transverse $A_{b-1}$ product.

\end{enumerate}
\end{proof}






\begin{proposition}[Every horizontal tacnode image is elliptic]\hypertarget{tag-01B8}{}\label{cdp:st:A.5.18}
\StatementRecord{cdp:st:A.5.18}{proved}{01B8}{\SeriesRef{cdp:st:A.5.14}, \SeriesRef{cdp:st:A.5.5}, \SeriesRef{cdp:st:A.5.8}, \SeriesRef{cdp:eq:A.T.3}, \SeriesRef{cdp:st:A.5.17}, \SeriesRef{cdp:eq:A.T.4}, \SeriesRef{cdp:eq:A.T.5}, \SeriesRef{cdp:eq:A.T.7}, \SeriesRef{cdp:eq:A.T.8}, \SeriesRef{cdp:st:A.5.15}, \SeriesRef{cdp:eq:A.T.9}, \SeriesRef{cdp:eq:A.T.12}, \SeriesRef{cdp:eq:A.T.13}, \SeriesRef{cdp:eq:A.T.15}, \SeriesRef{cdp:eq:A.T.18}, \SeriesRef{cdp:eq:A.T.20}, \SeriesRef{cdp:eq:A.T.19}, \SeriesRef{cdp:st:C2d}, \SeriesRef{cdp:eq:A.T.22}}

Let $f:Y\to C_0$ be a proper holomorphic map from a smooth
compact complex threefold with $b_2(Y)=0$, with connected
fibres and smooth general fibre a complex two-torus.
Let $mS$ be a full fibre with irreducible reduced support,
$\eta:T\to S$ its normalization, and $\pi:Z\to T$ a
resolution with an irrational ruling $p:Z\to B$.
Suppose every exceptional curve of $\pi$ is vertical and
the conductor divisor on $Z$ is $2C+V$, where $C$ is a
section of $p$ and $V$ is vertical. Let $D=\eta\pi(C)$.
Suppose the generic plane-curve transversal to $D$ in $S$ is an
ordinary tacnode.

Then $D$ is a smooth elliptic curve, $B=C$ is elliptic,
and $C\to D$ is an unramified double cover. The transversal
type along $D$ is everywhere the same ordinary tacnode.
The normalization $T$ is a smooth minimal ruled surface
over this elliptic curve. Its conductor is twice a
section, and $e(S)=0$. In particular its minimal
resolution is the identity.
In particular a rational image of
the horizontal tacnodal conductor cannot occur. The proof
allows vertical conductor and rational singularities of
$T$; it does not assume that the normalization map is
immersive at the special points under examination.
\end{proposition}


\begin{proof}

\begin{enumerate}[label={$\langle1\rangle$\arabic*.},leftmargin=*]
\item If the normalization of $D$ has positive genus, \ref{cdp:st:A.5.14}
and \ref{cdp:st:A.5.5} make $D$ smooth elliptic and its transversal
type constant. There are two generic normalization
branches of a tacnode. A vertical component of $Z$
is rational, so it cannot dominate this elliptic curve.
Thus both branches lie on $C$, and $C\to D$ is an
unramified double cover: the constant local equation
$x^2-y^4=0$ has two separately smooth normalization
branches on every sufficiently small disc in $D$.
Riemann--Hurwitz gives $g(C)=g(B)=1$.
It remains to exclude the case $D^\nu=\Sigma\simeq\PP^1$.

\item The finite map $C\to\Sigma$ has degree at least two,
because $C\simeq B$ has positive genus. Since a generic
tacnode has exactly two normalization branches, this
degree is exactly two. Consequently no vertical
conductor component can also dominate $D$: its generic
points would supply an additional normalization branch.
Write $g=g(C)\ge1$ and let $\mathcal B\subset\Sigma$
be the reduced branch divisor of this double cover.
Riemann--Hurwitz gives
\[
                   \deg\mathcal B=2g+2.             \tag{A.T.1}\hypertarget{tag-007M}{}\label{cdp:eq:A.T.1}
\]

\item We first bound the number of branches of the singular
image curve itself. \ref{cdp:st:A.5.8} supplies the descended ruling
$\widetilde p:T\to B$. Identify $C$ with $B$ using the
section, let $\phi:B\to D$ be its actual map to the
image curve, and put
\[
                    \psi=\phi\widetilde p:T\to D.
\]
Every vertical conductor component has constant $\psi$
value. Its image under $\eta$ differs from $D$, as just
proved, and meets $D$ in finitely many points. Thus
there is a finite set $A\subset D$ such that, for each
$s\in S$, either $\psi(\eta^{-1}(s))$ is a singleton
or all its values belong to $A$. Include in $A$ the
constant values on the vertical conductor components
and the values of $\phi$ over their finitely many
intersections with $D$. Fibres over the remaining
points of $D$ have their preimages on $C$, where all
the $\psi$ values equal that same point of $D$.

\item Choose an embedded finite tree $J\subset D$ containing
$A$, in a triangulation of the compact analytic curve.
It can be constructed by joining these finitely many
points by arcs, subdividing their intersections and
retaining a spanning tree. The quotient
$q:D\to D/J$ is a homotopy equivalence: a contractible
CW subcomplex is a cofibration, and the homotopy
extension property extends its contraction. Since
$q\psi$ is constant on every fibre of the finite
proper quotient $\eta$, there is a continuous map
\[
 a:S\to D/J,\qquad a\eta=q\psi,\qquad a|_D=q.     \tag{A.T.2}\hypertarget{tag-007X}{}\label{cdp:eq:A.T.2}
\]
The last equality holds first on the dense regular
locus of $D$ and then everywhere by continuity.

\item Use the explicit triangulated neighbourhood retraction
$r:N\to S$ of (A.Pi.3). For a sufficiently small disc
$\Delta$ around the fibre value, properness gives
$U=f^{-1}(\Delta)\subset N$. The map $(ar)|_U$ is
surjective on fundamental groups because its restriction
to $D$ is the homotopy equivalence $q$. Perturbing loops
off the real-codimension-two analytic support gives
$\pi_1(U\setminus S)\twoheadrightarrow\pi_1(U)$.
The torus-bundle exact sequence over $\Delta^*$ therefore
gives the actual quotient maps
\[
 \Z^4\rtimes\Z\twoheadrightarrow\pi_1(U)
       \twoheadrightarrow\pi_1(D/J)\simeq\pi_1(D).
                                                        \tag{A.T.3}\hypertarget{tag-0082}{}\label{cdp:eq:A.T.3}
\]
The first group is metabelian, so the last is metabelian.

\item For each singular point $x$ of $D$, let $r_x$ be the
number of its normalization branches and put
$\beta=\sum_x(r_x-1)$. The finite normalization map
identifies exactly those finite sets of points on
$\PP^1$; since it is a compact Hausdorff quotient,
this is also the actual underlying topological quotient.
Adjoining arcs between the points to be identified
gives a CW presentation homotopy equivalent to
$S^2\vee\bigvee^\beta S^1$. Hence
$\pi_1(D)$ is the free group $F_\beta$.
For $\beta\ge2$ this group surjects onto $S_4$ by
sending two free generators to $(1234)$ and $(12)$;
the commutator $A_4$ of that quotient is nonabelian.
Equation \eqref{cdp:eq:A.T.3} therefore forces $\beta\le1$.
By \ref{cdp:st:A.5.14} the total delta invariant $\delta=p_a(D)$
is at most two. If $s$ denotes the number of points
of $\Sigma$ over $\operatorname{Sing}D$, then
\[
 s=\#\operatorname{Sing}D+\beta\le\delta+\beta\le3.
                                                        \tag{A.T.4}\hypertarget{tag-0083}{}\label{cdp:eq:A.T.4}
\]

\item First suppose that $D$ is a local complete intersection.
We retain the exact bundles and jets of \ref{cdp:st:A.5.5}.
Let $\nu:\Sigma\to D$ be normalization, $I=\mathcal I_D$,
\[
 N=(\nu^*(I/I^2))^*,\qquad
 M=\nu^*(\mathcal O_Y(S)|_D).
\]
Thus $M$ is a degree-zero line bundle on $\Sigma$.
Let $\Xi$ be the curve-normalization
conductor divisor, $c=\deg\Xi=2\delta$, and let
$R\subset N$ and $Q=N/R$ be the saturated radical and
quotient of the generic rank-one Hessian. Write
$q_0=\deg Q$. The adjunction and conductor comparisons
(A.J.1) give
\[
 \deg N=-2+c,\quad \deg R=-2+c-q_0,\quad
 0\ne a\in H^0(\Sigma,M\otimes Q^{-2}),\quad q_0\le0.
                                                        \tag{A.T.5}\hypertarget{tag-0084}{}\label{cdp:eq:A.T.5}
\]
Here $a$ denotes the intrinsic quadratic coefficient,
whose zero divisor has degree $-2q_0$.
The actual quartic invariant $\Delta=ae-b^2/4$ of
\ref{cdp:st:A.5.5} is meromorphic in
$L^2$, where $L=M\otimes Q^{-1}\otimes R^{-2}$.
Its pole bound (A.J.9) gives a nonzero holomorphic
section $\sigma$ of the actual square line bundle
$L^2(4\Xi)=(L(2\Xi))^2$. Its zero divisor has degree
\[
               d_\Delta=8+2q_0.                    \tag{A.T.6}\hypertarget{tag-0085}{}\label{cdp:eq:A.T.6}
\]

\item The double cover $C\to\Sigma$ is obtained by adjoining
$\sqrt{-\Delta}$ to the function field, with the line
trivializations specified by $L$. Indeed over the
generic point a tacnode has normal expansion
$a\ell^2+b\ell r^2+er^4$, and the two normalization
branches are distinguished by the two roots of
$a(\ell/r^2)^2+b(\ell/r^2)+e$. Their discriminant
is $-4\Delta$. Changing a local frame of $L$ multiplies
the representing function of $\Delta$ by a square,
so this describes the same degree-two extension on
every overlap. At a point with local parameter $z$,
write that representing function as $z^n u(z)$,
$u(0)\ne0$. A root of the unit and removal of the
even power of $z$ show that its normalization is
ramified exactly when $n$ is odd. Adding $4\Xi$ changes
each order by an even number. Thus $\mathcal B$ is
exactly the odd-multiplicity support of
$\operatorname{div}(\sigma)$, and
\[
                  2g+2\le 8+2q_0.                 \tag{A.T.7}\hypertarget{tag-0086}{}\label{cdp:eq:A.T.7}
\]

\item At a point of $D_{\rm reg}$ where $a\ne0$, \ref{cdp:st:A.5.17}
applies to the actual local $m$th root of the full-fibre
map. Its nonzero Hessian makes the local equation a
constant transverse product, necessarily of type
$A_3$ because this is the generic type along $D$.
The conductor cover is therefore unramified there.
Every branch point outside $\nu^{-1}(\operatorname{Sing}D)$
is consequently a zero of $a$. Equations \eqref{cdp:eq:A.T.4}--\eqref{cdp:eq:A.T.5}
give the second bound
\[
                    2g+2\le -2q_0+s.              \tag{A.T.8}\hypertarget{tag-0087}{}\label{cdp:eq:A.T.8}
\]
Adding \eqref{cdp:eq:A.T.7} and \eqref{cdp:eq:A.T.8} yields
$4g+4\le8+s\le11$. Since $g\ge1$, we have
\[
           g=1,\qquad \deg\mathcal B=4,
                    q_0\in\{-2,-1,0\}.            \tag{A.T.9}\hypertarget{tag-0088}{}\label{cdp:eq:A.T.9}
\]

\item We next prove the exact local obstruction required in
the equality cases. Let a smooth point of $D$ be a
branch point at which $\Delta$ has a simple zero.
The preceding nonzero-Hessian argument gives $a(0)=0$.
In coordinates $D=(x,y)=0$ with curve parameter $z$,
choose the holomorphic saturated Hessian factor as the
first normal coordinate. The quadratic jet is then
$a(z)x^2$. The radical cubic vanishes identically,
so a local reduced equation has the exact expansion
\[
 f_0=a(z)x^2+x^2(xA+yB)+xy^2b(z)+xy^3C+y^4E,
                                                        \tag{A.T.10}\hypertarget{tag-007N}{}\label{cdp:eq:A.T.10}
\]
with holomorphic coefficients. The identity
$\Delta=ae-b^2/4$, $a(0)=0$ and the simple zero of
$\Delta$ imply $b(0)=0$, $a'(0)\ne0$ and $e(0)\ne0$.
Indeed an order at least two for $a$, together with
$b(0)=0$, would give an order at least two for $\Delta$;
and then $\Delta'(0)=a'(0)e(0)$ is nonzero.
Thus $b/a$ is holomorphic and $E(0)\ne0$.

\item Set $X=x+b(z)y^2/(2a(z))$. In this actual analytic
coordinate change the coefficient of $Xy^2$ cancels.
All terms created from $x^2(xA+yB)$ lie in
$(X^2)+(Xy^3)+(y^4)$. Therefore, collecting the
terms of $X$-degree zero and one, gives exactly
\[
 f_0=X^2A_1(X,y,z)+Xy^3B_1(y,z)+y^4C_1(y,z),
 \qquad C_1(0)\ne0.
\]
Choose a fourth root of $C_1$ and define
\[
 Y=yC_1^{1/4}+
                  \frac{XB_1}{4C_1^{3/4}}.
                                                        \tag{A.T.11}\hypertarget{tag-007O}{}\label{cdp:eq:A.T.11}
\]
The map $(X,y,z)\mapsto(X,Y,z)$ is locally
biholomorphic, and direct expansion gives
$f_0-Y^4\in(X^2)$. Hence $f_0=Y^4+X^2H(X,Y,z)$.
Its quadratic jet along $D$ is still $a(z)X^2$,
so $H(0,0,z)=a(z)$ and $\partial H/\partial z(0)\ne0$.
The further coordinate $Z=H$ gives exactly
\[
                         f_0=Y^4+X^2Z.             \tag{A.T.12}\hypertarget{tag-007P}{}\label{cdp:eq:A.T.12}
\]
To apply this to the actual map, first take a local
$m$th root of its unit. This multiplies the reduced
equation by a unit, which changes neither zero order
of $a$ nor of $\Delta$. The same coordinate construction
then makes the full map $(Y^4+X^2Z)^m$.
Statement \ref{cdp:st:A.5.15} excludes it. Consequently every
smooth branch point has order at least three in
$\operatorname{div}(\sigma)$.

\item We also record precisely when the conductor denominator
can be removed at a singular point $x\in D$.
If the intrinsic section $a$ is nonzero at every
normalization point above $x$, choose an actual
constant linear combination of the two lci conormal
generators which is nonzero under the normalized
quotient at all those finitely many points. The
corresponding diagonal entry $a_0$ of the symmetric
Hessian matrix belongs to $\mathcal O_{D,x}$ and
has nonzero value on every normalization branch.
It is therefore a unit in $\mathcal O_{D,x}$.
The other normalized factor coefficient is the
actual quotient $h=q_{XY}/a_0\in\mathcal O_{D,x}$.
Thus square completion and monic polynomial division
in the proof of (A.J.7)--(A.J.8) both take place
over $\mathcal O_{D,x}$ itself. The regular-sequence
graded pieces inject into their normalization
pullbacks, so their cancellations already descend
on $D$. We may repeat that proof with $d=1$.
It gives a holomorphic section $\Delta$ of $L^2$
above $x$, without a conductor pole. At a branch
point above $x$, this proves the sharper inequality
\[
 \operatorname{ord}\sigma\ge
                        4\operatorname{ord}\Xi+1.
                                                        \tag{A.T.13}\hypertarget{tag-007Q}{}\label{cdp:eq:A.T.13}
\]

\item We now exhaust the three integer cases in \eqref{cdp:eq:A.T.9}.
If $q_0=-2$, the divisor of $\sigma$ has degree four
and has four odd-multiplicity points, so each is
simple. At most three lie above singularities of
$D$, by \eqref{cdp:eq:A.T.4}. The remaining smooth branch point
contradicts \eqref{cdp:eq:A.T.12}.
If $q_0=0$, the section $a$ has no zeros and all
four branch points would have to lie above the
singular locus; this contradicts $s\le3$.

\item Finally let $q_0=-1$, so that $\deg\operatorname{div}(a)=2$
and $d_\Delta=6$. If $k$ of the four branch points
lie above $\operatorname{Sing}D$, then the local
obstruction and parity give
\[
               6\ge k+3(4-k)=12-2k.
\]
Hence $k=s=3$, and equality holds: the three singular
preimages have order one in $\operatorname{div}(\sigma)$
and the one smooth branch point has order three.
The relations $s\le\delta+\beta\le3$ now force
$\delta=2$, $\beta=1$ and two singular points of
delta invariant one. The exact delta-one
classification in \ref{cdp:st:A.5.14} identifies them as one
ordinary node and one cusp. Their normalization
conductor orders are respectively $(1,1)$ and $2$;
these follow directly from the displayed node and
cusp rings in that proof.

\item At the cusp preimage, if $a$ were nonzero,
\eqref{cdp:eq:A.T.13} would make the order of $\sigma$ at least
$4\cdot2+1=9$, exceeding its total degree six.
So $a$ vanishes there. It also vanishes at the
smooth branch point. These two distinct points
exhaust the degree-two divisor of $a$. The section
is consequently nonzero at both node preimages.
Apply \eqref{cdp:eq:A.T.13} at the node: each of its two branch
points has order at least $4\cdot1+1=5$ in
$\operatorname{div}(\sigma)$, again exceeding the
total degree six. This excludes every rational lci image.

\item We finish by treating all non-lci image germs permitted
by $\delta\le2$ and $\beta\le1$, with explicit analytic
rings and transverse coordinate maps. For a unibranch
reduced curve ring $A\subset\C\{t\}$ of colength at
most two, an element of order one makes $A=\C\{t\}$.
Otherwise its least positive valuation is two or three:
a least valuation at least four leaves three linearly
independent missing coefficients. If it is three and
the colength is two, the inclusion
$A\subset\C+t^3\C\{t\}$ is equality, since the latter
already has colength two. If it is two, choose a
parameter with $t^2\in A$ by taking a square root of
the leading unit. The even part of every convergent
series belongs to $\C\{t^2\}\subset A$. Subtracting
it shows that the least odd element can be written
$t^{2k+1}v(t^2)$; division by its unit puts
$t^{2k+1}$ in $A$. Consequently
$A=\C\{t^2\}\oplus t^{2k+1}\C\{t^2\}$, of colength
$k\le2$. These are the planar cusp and the planar
$A_4$ curve, besides the smooth case.

\item For two branches, write their image rings as $A_1,A_2$.
The two projections of $A\subset A_1\oplus A_2$ are
surjective. Their kernels define a common finite
quotient $Q_0$ and an exact fibre-product description
$A=A_1\times_{Q_0}A_2$. Thus
$\delta=\delta_1+\delta_2+\operatorname{length}Q_0$.
If both branches are smooth, the quotient is
$\C\{t\}/(t^k)$ with $k\le2$; after changing the
second parameter, the two maps identify the same
$k$-jet. The map
\[
 x\longmapsto(t,s),\qquad
 y\longmapsto(t^k,-s^k)
\]
identifies this fibre product with
$\C\{x,y\}/(y^2-x^{2k})$: write a matching pair as
the sum of its average and $y$ times its difference
divided by $2x^k$. These are the node and tacnode.
If one branch is singular, the delta budget forces
one cusp, one smooth branch, and $Q_0=\C$.
Three branches are excluded by $\beta\le1$.

\item An abstract plane curve ring in the given smooth
threefold is lci by the actual transverse-coordinate
construction in \ref{cdp:st:A.5.14}. Thus the only remaining
non-lci analytic rings, necessarily consuming the
whole delta budget at a single point, are
\[
\begin{array}{ll}
 A=\C\{t^3,t^4,t^5\},
 &I=(y^2-xz,\ yz-x^3,\ z^2-x^2y),\\[2pt]
 A=\C\{t^2,t^3\}\times_{\C}\C\{s\},
 &I=(xz,\ yz,\ y^2-x^3).
\end{array}                                             \tag{A.T.14}\hypertarget{tag-007R}{}\label{cdp:eq:A.T.14}
\]
These are normal forms of the actual embeddings.
Their three displayed coordinate functions form a
basis of the cotangent space of each ring. Lift them
to $\mathcal O_{Y,x}$; their differentials are linearly
independent because the embedding dimension is three.
They are therefore ambient coordinates by the inverse
function theorem. The displayed ideals are their
kernels, obtained respectively by the parametrizations
$(t^3,t^4,t^5)$ and the union of $(t^2,t^3,0)$ with
$(0,0,s)$.

We show that the local reduced equation $f_0$ actually
lies in $I^2$ in each case. In the first case put
\[
 g_1=y^2-xz,\quad g_2=yz-x^3,\quad
 g_3=z^2-x^2y,\quad
 F=x^5+xy^3+z^3-3x^2yz.
\]
The exact identities
\[
 xF=g_2^2-g_1g_3,\qquad
 yF=xg_1^2+g_2g_3,\qquad
 zF=g_3^2+xg_1g_2                              \tag{A.T.15}\hypertarget{tag-007S}{}\label{cdp:eq:A.T.15}
\]
show that the maximal ideal annihilates the class of
$F$ modulo $I^2$. Moreover
$I^{(2)}=I^2+(F)$ in the convergent local ring.
Here is a finite-module verification of that equality.
Let $J=I^2+(F)$. The quotient $\C\{x,y,z\}/J$ is
finite over the discrete valuation ring $\C\{x\}$,
since its reduction is the finite parametrized curve.
Its generic rank is nine: the curve has degree three
over the $x$ coordinate, and the square of its smooth
codimension-two ideal has length three along each
generic normal slice. Modulo $x$, the quotient has
the nine-element basis
\[
 1,y,y^2,y^3,\ z,yz,y^2z,\ z^2,yz^2,
\]
because its ideal is $(y^4,y^3z,y^2z^2,z^3)$.
A finite module over a discrete valuation ring is
the sum of its free and torsion parts; equality of
generic rank and residue-space dimension here makes
the torsion part zero. Thus $J$ is $x$-saturated.
By \eqref{cdp:eq:A.T.15}, $J$ and $I^2$ agree after inverting $x$.
The curve is lci away from the origin, so its symbolic
square differs from its ordinary square only by a
module supported there. Hence
$I^{(2)}=I^2:x^\infty=J$.

On its normalization let
$a_1=dg_1$ and $b_1=dg_2+t\,dg_1$.
They form a basis of the normalized conormal image;
$dg_3=-2t^2a_1+t b_1$.
The quadratic term of $F$ in this basis is
\[
 q_F=\frac3t a_1^2-\frac3{t^2}a_1b_1
                           +\frac1{t^3}b_1^2,
 \qquad \det(q_F)=\frac3{4t^4}.                  \tag{A.T.16}\hypertarget{tag-007T}{}\label{cdp:eq:A.T.16}
\]
Every equation in $I^2$ has a holomorphic quadratic
form in this same basis. If $f_0=cF+I^2$ had
$c(0)\ne0$, the determinant of its quadratic form
would have the noncancellable leading pole
$3c(0)^2/(4t^4)$: the cross terms with a holomorphic
form have pole at most three. It could not have
generic rank one. Thus $c(0)=0$, and \eqref{cdp:eq:A.T.15} gives
$f_0\in I^2$.

For the second ring put $G=y^2-x^3$ and
$P=(z,G)$, $Q_1=(x,y)$. Then $I=P\cap Q_1$ and
\[
 I^{(2)}=P^2\cap Q_1^2=I^2+(zG).                \tag{A.T.17}\hypertarget{tag-007U}{}\label{cdp:eq:A.T.17}
\]
Indeed $P^2=(z^2,zG,G^2)$, while $zG,G^2$ already
belong to $Q_1^2$. Multiplication by $z^2$ is
injective modulo $Q_1^2$, so membership of
$z^2A+zGB+G^2C$ in $Q_1^2$ forces $A\in Q_1^2$.
This gives the asserted generating ideal directly.
Furthermore
$x(zG)=(xz)G$, $y(zG)=(yz)G$, and
$z(zG)=(yz)^2-x(xz)^2$ all lie in $I^2$.
On the cusp branch use $a_1=d(xz)=t^2dz$ and
$b_1=dG=-3t^4dx+2t^3dy$.
The quadratic form of $zG$ is $t^{-2}a_1b_1$,
of determinant $-1/(4t^4)$. The same rank-one
argument forces its coefficient to vanish at the
origin, and the three identities again give
$f_0\in I^2$.

For clarity, the bundle used at these non-lci points
is now defined by its actual morphism. Put
\[
 K=\operatorname{im}\bigl(\nu^*(I/I^2)
                         \longrightarrow\nu^*\Omega_Y^1\bigr),
 \qquad N=K^*.
\]
The coherent image $K$ is locally free of rank two
on the smooth curve $\Sigma$. Let $R_{\rm ram}$ be
the zero divisor of $d\nu$, and let $K_{\rm sat}$
be the kernel of the surjection
$\nu^*\Omega_Y^1\to\Omega_\Sigma^1(-R_{\rm ram})$.
The torsion quotient $K_{\rm sat}/K$ defines a divisor
$J_{\rm jac}$. Taking determinants in these actual
maps gives
\[
 \deg N=-2+\deg(J_{\rm jac}-R_{\rm ram}).         \tag{A.T.18}\hypertarget{tag-007V}{}\label{cdp:eq:A.T.18}
\]
For $(t^3,t^4,t^5)$, the wedge $a_1\wedge b_1$
has coefficients $t^8(3,4t,5t^2)$, while $d\nu$
has common factor $t^2$. Thus the correction is
$8-2=6$. For the cusp branch in the second case,
the wedge is $-t^5(2,3t,0)$ and the tangent common
factor is $t$, giving correction four. On its
smooth branch the wedge is $s^2(dx\wedge dy)$,
giving correction two. In both cases \eqref{cdp:eq:A.T.18}
therefore gives $\deg N=4$.

Because $f_0\in I^2$, its quadratic form is a
holomorphic section of $\operatorname{Sym}^2K\otimes M$,
independent of its expression as a sum of products
of ideal generators. This follows either from its
actual Hessian or from injectivity of
$\operatorname{Sym}^2K\to\nu^*\operatorname{Sym}^2\Omega_Y^1$.
Its saturated radical and quotient again give line
bundles $R,Q$, with $q_0=\deg Q\le0$ and
$\deg R=4-q_0$. The generic quartic invariant
is meromorphic in $L^2$, where
$L=M\otimes Q^{-1}\otimes R^{-2}$; now
\[
                        \deg L^2=-16+2q_0.        \tag{A.T.19}\hypertarget{tag-007W}{}\label{cdp:eq:A.T.19}
\]

We bound its poles using actual transverse affine
coordinate maps, with the curve parameter held
fixed during every normal differentiation. For the
space cusp use, for $t\ne0$,
\[
 (u,v,t)\longmapsto
 \left(t^3,\ t^4+\frac{v}{3t^5},\
                   t^5-\frac{u}{t^3}+\frac{2v}{3t^4}\right).
                                                        \tag{A.T.20}\hypertarget{tag-007Y}{}\label{cdp:eq:A.T.20}
\]
This is locally biholomorphic and its two transverse
derivatives are the dual frame to $a_1,b_1$.
Each coefficient of a derivative has pole at most
five. The first normal derivatives of all three
$g_i$ are holomorphic, since
$dg_1=a_1$, $dg_2=b_1-ta_1$ and
$dg_3=-2t^2a_1+t b_1$.
Second normal derivatives have pole at most ten.
For $f_0=\sum c_{ij}g_i g_j$, the cubic Taylor
coefficients consequently have pole at most ten:
the only surviving terms use a first and a second
derivative of the $g_i$, or a first derivative of
$c_{ij}$ and two first derivatives of the $g_i$.
Every fourth derivative of a holomorphic ambient
function in \eqref{cdp:eq:A.T.20} has pole at most twenty.

Choose a holomorphic change of the $N$ frame adapted
to its saturated radical and quotient. It does not
increase these bounds. The radical cubic is zero
because the generic transversal is $A_3$.
Writing the relevant Taylor terms as
$a\ell^2+b\ell r^2+er^4$ therefore gives a
holomorphic $a$, a mixed cubic $b$ of pole at most
ten, and $e$ of pole at most twenty. The invariant
$\Delta=ae-b^2/4$ has pole at most twenty.
These are the same intrinsic jets on overlaps:
to weighted order four, where $\ell$ has weight
two and $r$ has weight one, an adapted change is
$\ell'=\alpha\ell+u r^2$, $r'=\beta r$.
Direct substitution multiplies $\Delta$ by
$\alpha^{-2}\beta^{-4}$; changes of the defining
section supply the factor $M^2$. Higher weighted
terms and changes of the curve parameter do not
alter these coefficients. This verifies the
transition law for the line $L^2$ in \eqref{cdp:eq:A.T.19}.

There is only one singular preimage $p$ in this case.
Thus $\Delta$ is a holomorphic section of the square
line $L^2(20p)$, whose degree is $4+2q_0$.
As in \eqref{cdp:eq:A.T.7}, the even pole correction preserves
branch parity, so
\[
 \deg\mathcal B\le4+2q_0,
 \qquad \deg\mathcal B\le-2q_0+1.
\]
The second bound again uses \ref{cdp:st:A.5.17} away from that
single singular image preimage. Adding them gives
$2\deg\mathcal B\le5$, contradicting
$\deg\mathcal B=2g+2\ge4$.

For the cusp-plus-line germ, the corresponding
actual affine coordinate maps are
\[
 (u,v,t)\longmapsto
       \left(t^2,t^3+\frac{v}{2t^3},\frac{u}{t^2}\right),
 \qquad
 (u,v,s)\longmapsto\left(\frac us,\frac vs,s\right).
                                                        \tag{A.T.21}\hypertarget{tag-007Z}{}\label{cdp:eq:A.T.21}
\]
Their first generator derivatives are holomorphic
in the respective normalized conormal frames.
The same product differentiation gives cubic pole
bounds six and two, and the fourth derivative bounds
are twelve and four. Thus $\Delta$ has poles bounded
by $12p+4q$ at the two normalization preimages.
The line $L^2(12p+4q)$ is a square and has degree
$2q_0\le0$ by \eqref{cdp:eq:A.T.19}. A nonzero section has no
positive-degree odd divisor, whereas this double
cover has at least four branch points. This is the
last contradiction. All non-lci cases have now
been excluded, and the positive-genus case at the
start gives the asserted elliptic curves and constant
transversal type.

We finally prove the normalization and Euler assertions,
keeping track of any unnecessary blowups in the original
choice of resolution. Factor $\pi$ through the minimal
resolution $\pi_{\min}:Z_{\min}\to T$, with the resulting
sequence of exceptional-curve blowdowns
$q:Z\to Z_{\min}$. Every curve contracted by $q$ is
vertical for $p$, by the stated assumption on $\pi$.
The equality of structure-sheaf direct images for each
blowdown, in local coordinates on $B$, gives a descended
ruling $p_{\min}:Z_{\min}\to B$ with
$p=p_{\min}q$. The image $C_{\min}$ is still a section,
and the pushed-forward effective conductor is
$D_{\min}=2C_{\min}+V_{\min}$ with $V_{\min}$ vertical.
The normalization-canonical comparison gives
$K_{Z_{\min}}+D_{\min}=\pi_{\min}^*\eta^*\omega_S$
as line bundles, and its rational Chern class is zero.

Statement \ref{cdp:st:C2d} gives a smooth normalization near the
inverse image of the elliptic curve $D$, with normal
line of degree zero on its unramified double cover.
The minimal resolution is an isomorphism there, so
$C_{\min}^2=0$. No other conductor or exceptional
curve meets it. Hence $V_{\min}\cdot C_{\min}=0$.
No ruling $(-1)$-curve can meet $C_{\min}$: one in
$V_{\min}$ is disjoint from it, while one outside
$V_{\min}$ would have
\[
 (K_{Z_{\min}}+2C_{\min}+V_{\min})\cdot E
       =-1+2C_{\min}\cdot E+V_{\min}\cdot E\ge1,
\]
contrary to the vanishing Chern class.

Contract ruling $(-1)$-curves to a minimal ruled bundle.
Every contraction therefore avoids the section. For
one blowdown $\sigma:W\to W'$ let $D'=\sigma_*D_W$,
let $a$ be the coefficient of its exceptional curve,
and let $d$ be the multiplicity of $D'$ at the centre.
The actual divisor comparison is
\[
 K_W+D_W=\sigma^*(K_{W'}+D')+(1+a-d)E.
                                                        \tag{A.T.22}\hypertarget{tag-0080}{}\label{cdp:eq:A.T.22}
\]
Its intersection with $E$ gives $a=d-1$, so the
log-canonical line and its zero rational Chern class
descend. The section remains disjoint from the
effective vertical part at every step. On a minimal
ruled bundle every nonzero effective vertical divisor
is a positive sum of fibres and meets the section.
The final vertical part is consequently zero. If
there had been a blowup in the reverse sequence,
its first centre would be away from this section,
and \eqref{cdp:eq:A.T.22} would give coefficient $-1$ on its
exceptional curve. Its strict transform retains that
negative coefficient under subsequent blowups, a
contradiction to effectiveness. Thus no ruling
blowdowns were needed: $Z_{\min}$ is already a
minimal ruled bundle and $D_{\min}=2C_{\min}$.

All exceptional curves of $\pi_{\min}$ would be
vertical. A minimal ruled bundle has irreducible
fibres of square zero, which cannot form the
negative-definite exceptional divisor of a surface
resolution. Hence $\pi_{\min}$ has no exceptional
curve and is an isomorphism. This proves the assertion
about $T$. Finally $\eta$ is an isomorphism away from
its conductor section $C_T$ and the elliptic image
$D$. Additivity of the topological Euler number on
these actual decompositions gives
\[
 e(S)=e(T)-e(C_T)+e(D)
      =e(\PP^1)e(B)-0+0=0.                       \tag{A.T.23}\hypertarget{tag-0081}{}\label{cdp:eq:A.T.23}
\]
The full multiple fibre has the same underlying
topological space and therefore the same Euler number.

\end{enumerate}
\end{proof}



\paragraph{Connected double covers in the two non-lci cases.}
The non-lci calculation also excludes a connected
degree-two conductor cover when its covering curve
is rational, for either of the two image rings in
\eqref{cdp:eq:A.T.14}. In the cusp-plus-line case the cleared
quartic has degree $2q_0\le0$, so it has no branch
divisor at all. In the space-cusp case a rational
double cover has branch degree two. The already
proved inequalities give
\[
                2\le4+2q_0,\qquad 2\le-2q_0+1,
\]
so the integer $q_0$ is $-1$. The cleared quartic
then has degree two and exactly two simple branch
points. Only one point lies above the singular
image germ. The other is a smooth simple branch
point, excluded by the actual coordinate reduction
\eqref{cdp:eq:A.T.10}--\eqref{cdp:eq:A.T.12}. Thus these two non-lci images
cannot support a connected double cover of any
nonnegative genus in a torus degeneration.

The integer branch budgets, the discriminant square
weights and the coordinate divisibility used in
\eqref{cdp:eq:A.T.10}--\eqref{cdp:eq:A.T.12} are checked in
\path{code/E_tacnode_bounds.py}; the non-lci charts,
symbolic generators and cubic/quartic pole orders
are checked in \path{code/E_low_delta_tacnode.py}.
The topological
quotient \eqref{cdp:eq:A.T.2} targets the actual singular curve
$D/J$; replacing it by its normalization would
lose the free-group obstruction in \eqref{cdp:eq:A.T.4}.




\begin{proposition}[Ordinary-triple images have arithmetic genus at most one]\hypertarget{tag-01B9}{}\label{cdp:st:A.5.19}
\StatementRecord{cdp:st:A.5.19}{proved}{01B9}{\SeriesRef{cdp:st:A.5.18}, \SeriesRef{cdp:st:A.5.14}, \SeriesRef{cdp:st:A.5.5}, \SeriesRef{cdp:eq:A.U.2}, \SeriesRef{cdp:eq:A.U.4}, \SeriesRef{cdp:eq:A.U.5}, \SeriesRef{cdp:eq:A.U.1}, \SeriesRef{cdp:st:C2d}}

Let $f:Y\to C_0$ be a proper holomorphic map from a smooth
compact complex threefold with $b_2(Y)=0$, with connected
fibres and smooth general fibre a complex two-torus.
Let $mS$ be a full fibre with irreducible reduced support,
and let $D\subset\operatorname{Sing}S$ be an irreducible
curve with generic ordinary-triple plane transversal.
Then $D$ is a local complete intersection in $Y$ and
$p_a(D)\le1$. The possibilities for its underlying curve
are a smooth rational curve, a rational curve with one
ordinary node or cusp, and a smooth elliptic curve.
In the elliptic case the transversal type is constant.

When $D$ is singular rational, every component of the
normalized conductor curve dominating $D$ is rational,
and its map to $D^\nu\simeq\PP^1$ has degree one.
Consequently, in the irrational ruled setup of \ref{cdp:st:A.5.18}
with the generic transversal replaced by an ordinary
triple point, the image of the noncontracted horizontal
section is either smooth rational or smooth elliptic.
In the elliptic case the horizontal cover is unramified
of degree three, the normalization is smooth minimal
elliptic ruled, and the fibre has Euler number zero.
In the rational case the horizontal cover has degree
two or three; these cases are not excluded by this
statement.
\end{proposition}


\begin{proof}

\begin{enumerate}[label={$\langle1\rangle$\arabic*.},leftmargin=*]
\item By \ref{cdp:st:A.5.14}, $p_a(D)\le2$. If its normalization has
positive genus, that statement and \ref{cdp:st:A.5.5} already give
a smooth elliptic curve of constant transverse type.
Suppose its normalization $\nu:\Sigma\to D$ is
$\PP^1$, with total delta invariant $\delta\le2$.

\item First let $D$ be lci. Put
$N=(\nu^*(I_D/I_D^2))^*$ and
$M=\nu^*(\mathcal O_Y(S)|_D)$.
The quadratic normal jet vanishes generically, so
its vanishing descends on the reduced curve through
the locally free graded piece $I_D^2/I_D^3$.
Thus the initial cubic is an actual section of
$\operatorname{Sym}^3N^*\otimes M$.
Adjunction and the curve conductor give
$\deg N=-2+2\delta$, while $\deg M=0$.
Its nonzero binary-cubic discriminant is a section
of $M^4\otimes(\det N)^{-6}$, of degree
\[
                         12-12\delta.             \tag{A.U.1}\hypertarget{tag-0089}{}\label{cdp:eq:A.U.1}
\]
Hence $\delta\le1$.

\item We exclude every non-lci germ of delta at most two
by its symbolic cube. The one- and two-branch analytic
classification is proved in (A.T.14). There is one
additional possibility when three branches occur:
delta at least two is forced by identifying their
constants, so equality makes the ring exactly
\[
 \{(a,b,c)\in\C\{t_1\}\oplus\C\{t_2\}\oplus\C\{t_3\}:
                      a(0)=b(0)=c(0)\}.
\]
Its actual embedding in the smooth threefold is the
union of the three coordinate axes, with ideal
$I=(xy,xz,yz)$. Indeed the three branch parameters
give a basis of the cotangent space, so their lifts
are ambient coordinates as in \ref{cdp:st:A.5.18}.
Thus there are exactly three non-lci cases to consider;
each consumes the whole delta budget at one point.

\item For the space cusp retain $I=(g_1,g_2,g_3)$ and
$F$ from (A.T.15). Its exact symbolic-cube identity is
\[
                         I^{(3)}=I^3+FI.           \tag{A.U.2}\hypertarget{tag-008A}{}\label{cdp:eq:A.U.2}
\]
Here is a verification in the convergent local ring.
The right-hand side agrees with $I^3$ after inverting
$x$, by (A.T.15). Its quotient is finite over
$\C\{x\}$, of generic rank $3\cdot6=18$: a square
normal plane has six monomials of total degree less
than three. Modulo $x$, its ideal in $\C\{y,z\}$ is
\[
 (y,z)^6+(y^2z^3,yz^4,z^5).
\]
The fifteen monomials of total degree at most four,
together with $y^5,y^4z,y^3z^2$, give residue dimension
eighteen. The same discrete-valuation-module argument
used for (A.T.15) makes the quotient $x$-torsion-free.
It is therefore the saturation of $I^3$, which equals
the symbolic cube since the curve is lci off the origin.
This proves \eqref{cdp:eq:A.U.2}.

\item For the cusp-plus-line ideal $I=(xz,yz,G)$ with
$G=y^2-x^3$, one has
\[
                   I^{(3)}=I^3+(zG)I.             \tag{A.U.3}\hypertarget{tag-008B}{}\label{cdp:eq:A.U.3}
\]
To verify this directly, write $P=(z,G)$ and
$Q=(x,y)$. Their powers are primary because these
prime ideals are complete intersections. The
symbolic cube is $P^3\cap Q^3$.
Write an element of $P^3$ as
$z^3A+z^2GB+zG^2C+G^3E$.
The last two terms are already in $Q^3$.
Reducing the other terms modulo $Q^3$ gives
$zA+GB=0$ there. The constant-in-$(x,y)$ coefficient
of $B$ must therefore be divisible by $z$; write
$B=zB'+B_Q$, $B_Q\in Q$.
It follows that $A+GB'\in Q^3$.
Consequently
\[
 P^3\cap Q^3=(z^3Q^3,z^2GQ,zG^2,G^3)
                        =I^3+(zG)I,
\]
as claimed.

\item For the three axes, membership of
$x^a y^b z^c$ in the symbolic cube means
$a+b,a+c,b+c\ge3$.
If one exponent is zero, the other two are at
least three and the monomial belongs to $I^3$.
If all are positive, either at least two are at
least two, in which case it belongs to $(xyz)I$,
or two are one, which violates one of the three
inequalities. Hence
\[
                         I^{(3)}=I^3+(xyz)I.       \tag{A.U.4}\hypertarget{tag-008C}{}\label{cdp:eq:A.U.4}
\]

\item In all three cases the generic ordinary-triple
equation $f_0$ belongs to $I^{(3)}$: on the smooth
locus of each curve branch its order in the two
normal variables is at least three. Define the
bundle by its actual conormal image,
\[
 K=\operatorname{im}\bigl(\nu^*(I/I^2)
                           \to\nu^*\Omega_Y^1\bigr),
 \qquad N=K^*.
\]
For the first two cases (A.T.18) gives $\deg N=4$.
For the three axes, on the branch $(t,0,0)$ the
frame is $t\,dy,t\,dz$, with conormal determinant
order two and no tangent ramification. Each of the
three branches has the same correction, for total
six, so again $\deg N=-2+6=4$.

\item Use the actual affine maps (A.T.20)--(A.T.21).
For the three axes add
\[
              (u,v,t)\longmapsto(t,u/t,v/t)
                                                        \tag{A.U.5}\hypertarget{tag-008D}{}\label{cdp:eq:A.U.5}
\]
on the first branch, and the two analogous permuted
maps on the other branches. In these charts the
first normal derivatives of the ideal generators
are holomorphic in the specified conormal frame.
An $I^3$ term consequently has holomorphic initial
cubic. The extra term in each of \eqref{cdp:eq:A.U.2}--\eqref{cdp:eq:A.U.4}
has initial cubic equal to the quadratic term of
its displayed symbolic-square generator, multiplied
by the first derivative of an ideal generator.

\item For the space cusp that quadratic term is (A.T.16),
so the cubic pole is at most three. For the
cusp-plus-line germ it is $t^{-2}uv$ on the cusp
and $s^{-1}v^2$ on the smooth branch, so the pole
bounds are two and one. For the three axes it is
$t^{-1}uv$ in \eqref{cdp:eq:A.U.5}, giving pole bound one on
each branch. Coefficients multiplying these terms
are holomorphic; their derivatives do not enter
the initial cubic because the lower normal jets
of the products vanish.

\item The initial cubic is intrinsic on the dense smooth
part of $D$, as all lower jets are zero. Its
transition under a normal-coordinate change is
therefore the actual cubic tensor in
$\operatorname{Sym}^3N^*\otimes M$; higher-order
coordinate terms make no contribution. The binary
discriminant is homogeneous of degree four in
these coefficients. It is a meromorphic section
of $M^4\otimes(\det N)^{-6}$, of degree $-24$,
with pole divisor of total degree at most twelve
in each of the three cases. After this exact pole
correction its line has negative degree, at most
$-12$. It cannot have the required nonzero
holomorphic section. This excludes all non-lci
cases and proves that $D$ is lci with $p_a(D)\le1$.
The list of underlying curves follows from the
delta-one classification in \ref{cdp:st:A.5.14}.

\item We give the finite-cover comparison at a singular
rational image. Here $\delta=1$, so \eqref{cdp:eq:A.U.1} has
degree zero and its nonzero section is nowhere zero
on $\Sigma=\PP^1$. The cubic defines a subscheme
$R_3\subset\PP(N)$ of tangent-root lines. Every
fibre consists of three distinct points. In a
projective coordinate chart, its derivative with
respect to the root coordinate is nonzero, so the
implicit function theorem makes
$R_3\to\Sigma$ a finite etale map of degree three.
There is no component in a whole projective fibre,
since its nonzero discriminant forbids an identically
zero cubic there.

\item Over the generic point of $D$, normalization of an
ordinary triple transversal separates exactly its
three tangent lines. This supplies an isomorphism
between the function-field algebra of $R_3$ and
that of the normalized conductor curves dominating
$D$, after their maps are factored through $\Sigma$.
Both are normal finite curves over $\Sigma$.
Taking the integral closure of $\mathcal O_\Sigma$
in this same algebra therefore identifies the two
finite covers everywhere. Every connected finite
etale cover of the simply connected Riemann sphere
has degree one. Thus each of these conductor
components is rational and maps isomorphically to
$\Sigma$, proving the second assertion.

\item In the irrational ruled setting the horizontal
section has positive genus. Its finite map to a
rational image normalization must have degree at
least two. The preceding comparison excludes a
singular rational image, and the generic triple
transversal bounds its degree above by three.
The remaining rational image is therefore smooth,
with degree two or three. If the image is elliptic,
no vertical rational curve can dominate it, so all
three branches belong to the horizontal section.
Statement \ref{cdp:st:C2d} gives an unramified degree-three
cover, normal degree zero and constant type.
The exact minimal-resolution factorization and
blowdown argument (A.T.22) applies without change
to this coefficient-two section: its proof uses
only those properties. It makes the normalization
smooth minimal elliptic ruled with conductor twice
the section. The actual normalization Euler formula
(A.T.23) then gives zero.

\end{enumerate}
\end{proof}



The singular rational possibilities cannot be
removed by a blanket local spherical-cycle argument.
For example
\[
                     f=x(yz)(x-yz)
                                                        \tag{A.U.6}\hypertarget{tag-008E}{}\label{cdp:eq:A.U.6}
\]
has ordinary-triple transversals along the nodal curve
$D=(x,yz)=0$ and an additional ordinary-double
critical line $y=z=0$. For $t\ne0$, put
$B_t=\{(x,v):xv(x-v)=t\}$. The exact maps
\[
 B_t\times\C^*\longrightarrow f^{-1}(t),
 \quad (x,v,y)\longmapsto(x,y,v/y),\qquad
 (x,y,z)\longmapsto(x,yz,y)
                                                        \tag{A.U.7}\hypertarget{tag-008F}{}\label{cdp:eq:A.U.7}
\]
are inverse holomorphic maps. Nonzero $t$ forces
$v,y$ and $z$ to be nonzero. The smooth affine
curve $B_t$ and $\C^*$ are aspherical, so this
regular affine fibre has no spherical homotopy
class. This is a local model, not a compact
$b_2=0$ example. A deformation smoothing $yz$
changes the additional critical line's transverse
type on the boundary and does not satisfy the
comparison used in \ref{cdp:st:A.5.17}.

Exact ideal, normal-coordinate and discriminant
calculations are retained in
\path{code/E_triple_images.py} and
\path{code/E_triple_images.out}; the independently
computed symbolic saturations are also in
\path{code/A_low_delta_cusp.py}.




\begin{proposition}[Rational normalization bounds with the exact constant-sheaf map]\hypertarget{tag-01B6}{}\label{cdp:st:A.5.16}
\StatementRecord{cdp:st:A.5.16}{proved}{01B6}{\SeriesRef{cdp:st:A.5.17}, \SeriesRef{cdp:st:A.5.12}, \SeriesRef{cdp:eq:A.RT.5}, \SeriesRef{cdp:st:A.5.9}, \SeriesRef{cdp:eq:A.RT.9}, \SeriesRef{cdp:st:A.5.3}, \SeriesRef{cdp:eq:A.RT.8}}

Let $f:Y\to B$ be a proper holomorphic map from a smooth compact
complex threefold with $b_2(Y)=0$, with connected fibres and smooth
general fibre a two-torus. Let $mS$ be a full fibre with irreducible
nonnormal reduced support. Write $\eta:T\to S$ for normalization,
$\pi:Z\to T$ for its minimal resolution, and $\Gamma\subset S$
for the conductor-image scheme. Suppose that $Z$ is rational.
Put $Q=\eta_*\OO_T/\OO_S$ and $e=h^0(S,\omega_S)$.
Then $e\in\{0,1\}$, and
\[
\begin{gathered}
 R^1\pi_*\OO_Z=0,\qquad
 H^1(T;\mathbb C)=H^1(T,\OO_T)=H^2(T,\OO_T)=0,\\
 H^0(S,Q)\xrightarrow{\ \simeq\ }H^1(S,\OO_S),\qquad
 h^0(S,Q)=1+e,\quad h^1(S,Q)=e,\quad
 \chi(\Gamma,\OO_\Gamma)=-1.
\end{gathered}
\tag{A.RT.1}\hypertarget{tag-0071}{}\label{cdp:eq:A.RT.1}
\]
In particular $T$ has rational singularities. The coefficient map
\[
 H^1(S;\mathbb C)\longrightarrow H^1(S,\OO_S)
\tag{A.RT.2}\hypertarget{tag-007C}{}\label{cdp:eq:A.RT.2}
\]
is injective, so $b_1(S)\le1+e\le2$.

There is an effective integral divisor $D$ on $Z$ with
$\OO_Z(D)\simeq K_Z^{-1}$, whose strict-transform part is the
normalization conductor. It satisfies
\[
 h^0(D,\OO_D)=h^1(D,\OO_D)=1,
 \quad D_{\rm red}\text{ connected},
 \quad p_a(D_{\rm red})\le1.
\tag{A.RT.3}\hypertarget{tag-007F}{}\label{cdp:eq:A.RT.3}
\]
The conductor image is connected,
$p_a(\Gamma_{\rm red})\le2$, and
$e(\Gamma_{\rm red})\ge0$. Set-theoretically
$D_{\rm red}=\pi^{-1}(\eta^{-1}(\Gamma_{\rm red}))$.
Exceptional curves over points outside the conductor have
coefficient zero in $D$; they are not to be counted as conductor
branches. With $E_p=(\pi^{-1}(p))_{\rm red}$ the exact Euler
formula is
\[
 e(S)=e(Z)-e(D_{\rm red})+e(\Gamma_{\rm red})
       -\sum_{p\in\operatorname{Sing}T\setminus\eta^{-1}(\Gamma)}
                  \bigl(e(E_p)-1\bigr).
\tag{A.RT.4}\hypertarget{tag-007G}{}\label{cdp:eq:A.RT.4}
\]
In the torus smoothing of the statement, \ref{cdp:st:A.5.17} excludes
every point in this sum. Hence $T$ is smooth outside the
normalization conductor, and the sum is zero.
\end{proposition}


\begin{proof}

\begin{enumerate}[label={$\langle1\rangle$\arabic*.},leftmargin=*]
\item The canonical evaluation argument in \ref{cdp:st:A.5.12} gives
$\omega_S^m\simeq\OO_S$ and $\chi(S,\OO_S)=0$.
Thus $e=0$ or $1$, with $e=1$ precisely when
$\omega_S\simeq\OO_S$, and
$h^1(S,\OO_S)=1+e$.

\item We first prove a useful fact without the rationality assumption.
Finite duality identifies the normal dualizing sheaf with
\[
 \omega_T=\mathfrak c_T\otimes\eta^*\omega_S
       \ \subset\ \eta^*\omega_S,
\tag{A.RT.5}\hypertarget{tag-007H}{}\label{cdp:eq:A.RT.5}
\]
where $\mathfrak c_T$ is the conductor ideal on $T$.
A nonzero global section of $\omega_T$ is therefore a nonzero
section of the torsion line $\eta^*\omega_S$ vanishing on its
nonempty conductor. Its $m$th tensor power would be a nonzero
holomorphic function on compact connected $T$ vanishing there,
which is impossible. Consequently
$H^2(T,\OO_T)=H^0(T,\omega_T)^\vee=0$.

\item For the proper resolution $\pi$, normality gives
$\pi_*\OO_Z=\OO_T$, and its fibres have dimension at most one.
The low-degree Leray sequence is
\[
\begin{split}
0\longrightarrow H^1(T,\OO_T)\longrightarrow H^1(Z,\OO_Z)
 \longrightarrow H^0(T,R^1\pi_*\OO_Z)\\
 \longrightarrow H^2(T,\OO_T)\longrightarrow H^2(Z,\OO_Z)
 \longrightarrow0.
\end{split}
\tag{A.RT.6}\hypertarget{tag-007I}{}\label{cdp:eq:A.RT.6}
\]
Rationality of $Z$ makes both its positive-degree groups zero.
Together with \eqref{cdp:eq:A.RT.5}, this gives $H^1(T,\OO_T)=0$ and
$R^1\pi_*\OO_Z=0$, since this last sheaf has finite support.
The fibres of $\pi$ are connected. Proper base change for the
constant sheaf gives $R^0\pi_*\mathbb C_Z=\mathbb C_T$.
The degree-one topological Leray edge map is consequently an
injection $H^1(T;\mathbb C)\hookrightarrow H^1(Z;\mathbb C)=0$.

\item In the normalization exact sequence
\[
 0\longrightarrow\OO_S\longrightarrow\eta_*\OO_T
       \longrightarrow Q\longrightarrow0,
\tag{A.RT.7}\hypertarget{tag-007J}{}\label{cdp:eq:A.RT.7}
\]
the map on $H^0$ is the identity on constants. The preceding
vanishings therefore make its connecting map
$H^0(Q)\to H^1(\OO_S)$ an isomorphism and its next connecting
map $H^1(Q)\to H^2(\OO_S)$ an isomorphism. These are the
maps and dimensions in \eqref{cdp:eq:A.RT.1}. In particular $\chi(Q)=1$.
The exact conductor-duality isomorphism proved in \ref{cdp:st:A.5.9} is
\[
 Q\simeq\omega_\Gamma\otimes\omega_S^{-1}|_\Gamma.
\tag{A.RT.8}\hypertarget{tag-007K}{}\label{cdp:eq:A.RT.8}
\]
The conductor image is a proper Cohen--Macaulay curve, including
any nilpotent structure. Serre duality on it gives
$\chi(Q)=-\chi(\omega_S|_\Gamma)$. A torsion line bundle has
degree zero on every reduced component, and therefore also has
total degree zero on this possibly nonreduced curve. Curve
Riemann--Roch gives
$\chi(\omega_S|_\Gamma)=\chi(\OO_\Gamma)$.
This proves $\chi(\OO_\Gamma)=-1$.

\item Here is the asserted injectivity of the complex coefficient map.
In the category of sheaves of complex vector spaces set
\[
 \mathcal V=\operatorname{coker}
     (\mathbb C_S\longrightarrow\eta_*\mathbb C_T).
\]
The diagram from constants to holomorphic functions induces the
specific sheaf map
\[
 \mathcal V\longrightarrow Q,
 \quad[(a_1,\ldots,a_r)]\longmapsto[(a_1,\ldots,a_r)] .
\tag{A.RT.9}\hypertarget{tag-007L}{}\label{cdp:eq:A.RT.9}
\]
At a point $s\in S$, the $r$ entries refer to the connected
normalization germs above $s$. A tuple of branchwise constants
that comes from an element of $\OO_{S,s}$ has equal values on
all these germs, because evaluation at their common image point
is one scalar. It is therefore a diagonal tuple. This proves
that \eqref{cdp:eq:A.RT.9} is injective on every stalk. No $\OO_S$-module
structure on $\mathcal V$ is asserted or needed.

\item A finite map has zero positive-degree direct images of the
constant sheaf, by proper base change on its finite fibres.
Since $H^1(T;\mathbb C)=0$ and the map on constants is an
isomorphism, the connecting map from
$0\to\mathbb C_S\to\eta_*\mathbb C_T\to\mathcal V\to0$
identifies $H^0(S,\mathcal V)$ with $H^1(S;\mathbb C)$.
Naturality of the connecting maps gives the commutative square
\[
\begin{array}{ccc}
 H^0(S,\mathcal V)&\xrightarrow{\ \simeq\ }&H^1(S;\mathbb C)\\
 \big\downarrow&&\big\downarrow\\
 H^0(S,Q)&\xrightarrow{\ \simeq\ }&H^1(S,\OO_S).
\end{array}
\tag{A.RT.10}\hypertarget{tag-0072}{}\label{cdp:eq:A.RT.10}
\]
Its left arrow is injective by \eqref{cdp:eq:A.RT.9}, proving \eqref{cdp:eq:A.RT.2}.

\ref{cdp:st:A.5.3} gives the effective integral divisor $D$ and the exact
line-bundle relation
$K_Z\otimes\OO_Z(D)\simeq(\eta\pi)^*\omega_S$.
The torsion line on the right is trivial on the smooth rational
surface: its Picard group is torsion-free and has zero identity
component. Hence $\OO_Z(-D)\simeq K_Z$. The divisor sequence
$0\to K_Z\to\OO_Z\to\OO_D\to0$, rationality of $Z$ and
Serre duality prove \eqref{cdp:eq:A.RT.3}. In detail the first positive
cohomology of $K_Z$ is zero, and $H^2(K_Z)=\mathbb C$;
the positive cohomology of $\OO_Z$ is zero. Thus $D$ is
connected. Restriction $\OO_D\to\OO_{D_{\rm red}}$ is
surjective on $H^1$, because the kernel is supported in
dimension one. Consequently $p_a(D_{\rm red})\le1$.

\item We identify its support, rather than discarding exceptional
components. On a minimal resolution every exceptional
irreducible curve $E$ has $K_Z\cdot E\ge0$. Indeed
$E^2<0$, and adjunction gives
$K_Z\cdot E=2p_a(E)-2-E^2$; a smooth rational
$(-1)$-curve is excluded by minimality. If $E$ has coefficient
zero in $D$, then
$0=(K_Z+D)\cdot E$ forces $K_Z\cdot E=D\cdot E=0$.
Thus $E$ is disjoint from $D_{\rm red}$.
Every exceptional fibre over the normalization conductor meets
its strict transform. Connectedness of that fibre then forces
all its exceptional curves to have positive coefficient in $D$.
Conversely, a positive-coefficient exceptional curve over a point
outside the conductor would lie in a component of $D_{\rm red}$
disjoint from every strict conductor curve. The connectedness
just proved, and the nonempty strict conductor, exclude this.
This proves the stated support equality and the zero-coefficient
assertion. It also proves connectedness of $\Gamma_{\rm red}$.

\item For completeness the bound on its arithmetic genus follows
directly from \eqref{cdp:eq:A.RT.8}. One has
$h^1(\Gamma,\omega_S|_\Gamma)=h^0(Q)=1+e$,
and restriction to $\Gamma_{\rm red}$ is surjective on $H^1$.
For $e=1$ the line is trivial, giving
$p_a(\Gamma_{\rm red})\le2$. For $e=0$ its $H^1$ has
dimension at most one, and degree-zero Riemann--Roch gives
$p_a(\Gamma_{\rm red})=1-h^0(\omega_S|_{\Gamma_{\rm red}})
+h^1(\omega_S|_{\Gamma_{\rm red}})\le2$.
Finally normalization and resolution are isomorphisms off the
conductor and the finitely many exceptional fibres.
Stratifying by these disjoint closed sets gives \eqref{cdp:eq:A.RT.4} by
additivity of the compactly supported Euler characteristic.
At a point outside the conductor, $T\to S$ is an
isomorphism and $S$ is normal. A normal complex surface
has isolated singularities. The reduced local equation
therefore has an isolated critical point whenever that
point is singular. Since the full fibre has equation
$u h^m$, the isolated clause of \ref{cdp:st:A.5.17} excludes it.
This proves the final assertion and eliminates the displayed
sum by an actual local comparison with the nearby torus.

\item We prove the asserted Euler bound for the conductor image,
including its possible positive-genus component. Put
$C=\Gamma_{\rm red}$. The normalization exact sequence for
the connected reduced curve $D_{\rm red}$ shows that the sum
of its normalized component genera is at most
$p_a(D_{\rm red})\le1$. Any component with genus one is
smooth, is unique, and meets the remaining rational curves
only in trees; a second attachment path joining distinct
points of that elliptic component would increase the
arithmetic genus to at least two. Each component of $C$ is
the finite image of a strict conductor component of
$D_{\rm red}$. Riemann--Hurwitz therefore shows that its
normalization is rational or elliptic, with at most one
elliptic-normalized component.

\item Write $n_0$ for the number of rational-normalized components
of $C$, $n_1\in\{0,1\}$ for the number of its
elliptic-normalized components, and $r_p,\delta_p$ for the
number of branches and the delta invariant at a singular
point. The exact coherent and topological normalization maps
give
\[
 \chi(\OO_C)=n_0-\sum_p\delta_p,
 \qquad e(C)=2n_0-\sum_p(r_p-1).
\]
For each reduced curve germ, independent constants on its
normalized branches give $\delta_p\ge r_p-1$.
The already proved arithmetic-genus bound and connectedness
give $\chi(\OO_C)\ge-1$, hence $e(C)\ge n_0-1$.
If $n_0\ge1$ this is nonnegative. If $n_0=0$, there is
one elliptic-normalized component and
$\sum\delta_p\le1$. A negative value could therefore only
be $-1$, from one ordinary node. Indeed its curve ring has
colength one in the direct sum of its two normalization
rings, and is contained in the subring of pairs with equal
constant value; equality of colengths makes these rings
equal. The two points over that node are distinct on the
elliptic normalization.

\item Here the hypersurface embedding excludes this last case.
At least one point $t\in T$ over the node receives both
local conductor branches. Otherwise partition the local
irreducible hypersurface branches according to which of
the two conductor branches they receive. Both parts are
nonempty. Two surface branches from opposite parts meet
in a curve, since their distinct equations in the regular
three-dimensional ambient local ring generate an ideal
of height two. That intersection curve belongs to the
conductor image, so is one of its two branches. Both
surface branches receive it, contradicting the partition.
For a locally irreducible hypersurface there is only one
normalization point and the assertion is immediate.

\item The connected resolution fibre $\pi^{-1}(t)$ lies in
$D_{\rm red}$ by the support equality proved above.
The strict conductor curves over the two branches are
elliptic, because a rational curve cannot map
nonconstantly to an elliptic curve. There is at most one
such curve on $D_{\rm red}$. If the two strict curves
were different, they would already give total genus two.
If they coincide, they meet $\pi^{-1}(t)$ at distinct
points: the map of that normal elliptic curve factors
through the normalization of its image, so one point
cannot map to both points over the node. The connected
resolution fibre then supplies an additional path joining
those points, outside the elliptic curve. Its components
are rational, since $T$ has rational singularities.
This produces a graph cycle in addition to the elliptic
genus, again forcing $p_a(D_{\rm red})\ge2$.
Both alternatives are impossible. Thus $e(C)\ge0$.

\end{enumerate}
\end{proof}



We record the exact minimal-model calculation that follows from
\ref{cdp:st:A.5.16}. This is a reduction of the outstanding Euler calculation;
the configurations with negative complement Euler number are
retained in the calculation.
Choose any sequence of contractions of smooth rational
$(-1)$-curves from $Z$ to a minimal rational surface $Z_0$,
and let $D_0$ be the pushforward of $D$. At one contraction
$\sigma:Z_i\to Z_{i-1}$, let $E$ be the exceptional curve,
$a$ its coefficient in $D_i$, and $r$ the multiplicity of
$D_{i-1}$ at the centre. The actual divisor comparison is
\[
 K_{Z_i}+D_i
   =\sigma^*(K_{Z_{i-1}}+D_{i-1})+(1+a-r)E.
 \tag{A.RT.11}\hypertarget{tag-0073}{}\label{cdp:eq:A.RT.11}
\]
Both canonical-plus-boundary line bundles are trivial, so
intersection with $E$ gives $a=r-1$. In particular the centre
lies on $D_{i-1}$, and effectivity is retained at every stage.
If $r=1$, the centre is a smooth point of a reduced boundary
component, $a=0$, and the strict boundary is isomorphic to the
old boundary. The blowup increases its complement Euler number
by one. If $r\ge2$, the new reduced boundary is exactly the
full inverse image of the old reduced boundary. Both the
surface and the boundary replace the centre by $\mathbf P^1$,
so the complement Euler number is unchanged. Thus
\[
 e(Z\setminus D_{\rm red})
 =e(Z_0\setminus (D_0)_{\rm red})+N_0,
 \qquad N_0\ge0,
 \tag{A.RT.12}\hypertarget{tag-0074}{}\label{cdp:eq:A.RT.12}
\]
where $N_0$ is the number of multiplicity-one centres in this
specified sequence of reverse blowups. Formula \eqref{cdp:eq:A.RT.12} concerns
the resolution complement. For the actual torus smoothing,
the isolated clause of \ref{cdp:st:A.5.17} has already removed the
off-conductor exceptional correction in \eqref{cdp:eq:A.RT.4}. This local
argument is needed: on an abstract minimal $\mathbb F_2$,
its exceptional $(-2)$-section can be disjoint from an
anticanonical divisor without having arisen from one of the
$N_0$ blowups.

On $\mathbf P^2$, the divisor $D_0$ is an effective cubic.
Its possible reduced supports and their complement Euler
numbers are
\[
\begin{array}{c|c}
 \text{reduced support}&e(\mathbf P^2\setminus(D_0)_{\rm red})\\ \hline
 \text{smooth cubic, nodal cubic, cuspidal cubic}&3,\ 2,\ 1\\
 \text{line and conic, transverse or tangent}&1,\ 0\\
 \text{three lines forming a triangle, or concurrent}&0,\ -1\\
 \text{two distinct lines, or one line}&0,\ 1.
\end{array}
\tag{A.RT.13}\hypertarget{tag-0075}{}\label{cdp:eq:A.RT.13}
\]
For example the union of three concurrent lines has Euler
number $6-3+1=4$, whereas $e(\mathbf P^2)=3$.
The last row retains respectively a double line plus a line,
and a triple line. These are all cubic possibilities: an
irreducible singular cubic has arithmetic genus one and
therefore exactly a node or a cusp; the other degrees partition
three as $2+1$ or $1+1+1$.

On $\mathbb F_n$, with $n=0$ or $n\ge2$, write $C_0^2=-n$
for a minimal section and $F$ for a ruling fibre. Then
\[
 D_0\sim 2C_0+(n+2)F.
 \tag{A.RT.14}\hypertarget{tag-0076}{}\label{cdp:eq:A.RT.14}
\]
Its horizontal part has degree two over the ruling base.
If that part consists of two distinct sections $C_1+C_2$,
write $V$ for the vertical part, $d$ for its total degree and
$k=C_1\cdot C_2$. Adjunction on either section gives
$k+d=2$. Let $q$ be the number of distinct intersection
points of $C_1$ and $C_2$, and let $r$ count those points
whose ruling fibre occurs in $V_{\rm red}$. A vertical fibre
meets the two sections in two points unless they coincide at
that fibre. Additivity therefore gives the exact formula
\[
 e((D_0)_{\rm red})=4-q+r,
 \qquad e(\mathbb F_n\setminus(D_0)_{\rm red})=q-r\ge0.
 \tag{A.RT.15}\hypertarget{tag-0077}{}\label{cdp:eq:A.RT.15}
\]
This includes two sections and a fibre through their common
point: that union has Euler number four and complement Euler
number zero.

If the horizontal part is an irreducible bisection $C$, the
same adjunction calculation gives $p_a(C)=1-d$, so $d=0$
or $1$. For $d=0$, its Euler number is at most two, and the
complement Euler number is at least two. For $d=1$, the
bisection is a smooth rational curve, the vertical part is one
reduced fibre, and their set-theoretic intersection has one or
two points. The complement Euler number is respectively one
or two.

The remaining horizontal part is $2C$, where $C$ is a section.
Write $C\sim C_0+kF$. The vertical divisor has total degree
$d=n+2-2k$, and, with $v$ its number of distinct fibres,
\[
 e((D_0)_{\rm red})=2+v,
 \qquad e(\mathbb F_n\setminus(D_0)_{\rm red})=2-v.
 \tag{A.RT.16}\hypertarget{tag-0078}{}\label{cdp:eq:A.RT.16}
\]
If $C\ne C_0$, positivity of intersection with $C_0$ gives
$k\ge n$; hence $d\le2-n$ and $v\le2$.
For $n=0$ one also has $d\le2$. Consequently the only
negative values on a minimal ruled surface occur in the
explicit configurations
\[
 Z_0=\mathbb F_n,\quad n\ge2,\qquad
 D_0=2C_0+\sum_{j=1}^{v}a_jF_j,\qquad
 a_j\ge1,\quad \sum_j a_j=n+2,\quad v\ge3.
 \tag{A.RT.17}\hypertarget{tag-0079}{}\label{cdp:eq:A.RT.17}
\]
The fibres $F_j$ here are distinct. Equations
\eqref{cdp:eq:A.RT.11}--\eqref{cdp:eq:A.RT.17}, together with \eqref{cdp:eq:A.RT.4}, isolate the
remaining contributions exactly; an arbitrary identification
of their conductor branches has not been substituted for an
actual hypersurface normalization map.

There is a further fundamental-group restriction on precisely
these negative minimal-model configurations. Their reduced
supports are unions of rational curves with simply connected
underlying spaces: three concurrent lines on $\mathbf P^2$,
or the section with its attached ruling fibres in \eqref{cdp:eq:A.RT.17}.
Every reverse blowup in \eqref{cdp:eq:A.RT.11} preserves this property.
A multiplicity-one blowup leaves the support isomorphic to
the preceding one. A positive-coefficient blowup at a smooth
point adds a rational leaf, at an ordinary double point
subdivides the incidence edge, and at the initial concurrent
point replaces that vertex by an exceptional $\mathbf P^1$
with three attached branches. Subsequent supports have only
the resulting ordinary double points. None of these explicit
operations introduces a topological loop. Consequently
$D_{\rm red}$ is simply connected throughout these cases.

The actual map $h=\eta\pi$ expresses $S$ as the topological
pushout of
\[
 D_{\rm red}\hookrightarrow Z,
 \qquad D_{\rm red}\xrightarrow{\ h|_{D_{\rm red}}\ }
                           \Gamma_{\rm red}.
 \tag{A.RT.18}\hypertarget{tag-007A}{}\label{cdp:eq:A.RT.18}
\]
Indeed $h$ is an isomorphism off these closed sets and the
whole preimage of $\Gamma_{\rm red}$ is $D_{\rm red}$.
The continuous map from this pushout to $S$ is a bijection
from a compact space to a Hausdorff space, and hence is a
homeomorphism. A triangulation of the pair
$(Z,D_{\rm red})$ makes its inclusion a cofibration.
The double mapping-cylinder model for \eqref{cdp:eq:A.RT.18} therefore
has the same homotopy type as the pushout, and van Kampen
gives
\[
 \pi_1(S)\simeq
 \pi_1(Z)*_{\pi_1(D_{\rm red})}\pi_1(\Gamma_{\rm red})
 \simeq\pi_1(\Gamma_{\rm red}).
 \tag{A.RT.19}\hypertarget{tag-007B}{}\label{cdp:eq:A.RT.19}
\]
All normalizations of the image components are $\mathbf P^1$,
by Riemann--Hurwitz. Gluing their finitely many points over
singularities identifies the topological curve with a graph
of simply connected spheres. Its fundamental group is a
free group, of rank $b_1(\Gamma_{\rm red})$.

That free group is a quotient of a metabelian group. To see
the exact comparison, triangulate $(Y,S)$ and retract a small
regular neighbourhood $N$ onto $S$, as in \ref{cdp:st:A.5.11}. Properness
puts a sufficiently small full fibre tube $U=f^{-1}(\Delta)$
inside $N$. Its retraction is surjective on fundamental
groups because it restricts to the identity on $S$.
General position makes
$\pi_1(U\setminus S)\to\pi_1(U)$ surjective, and the
torus bundle over $\Delta^*$ has
$\pi_1(U\setminus S)=\mathbb Z^4\rtimes\mathbb Z$.
This last group has abelian commutator subgroup, so its
quotient \eqref{cdp:eq:A.RT.19} is metabelian. A free group of rank at
least two is not metabelian: it maps onto $S_4$, whose
commutator group $A_4$ is not abelian. We obtain
\[
 \pi_1(S)\simeq\pi_1(\Gamma_{\rm red})\in\{1,\mathbb Z\},
 \qquad b_1(\Gamma_{\rm red})\le1.
 \tag{A.RT.20}\hypertarget{tag-007D}{}\label{cdp:eq:A.RT.20}
\]
Under the full CDP Betti hypotheses, the nonzero class on
every support from \ref{cdp:st:D.1.9} excludes the trivial group in
\eqref{cdp:eq:A.RT.20}. This last use of the global transgression is
separate from the preceding $b_2(Y)=0$ calculations.

For these two boundary classes one can further exclude a
one-component image with first Betti number one. Its normalization
is $\mathbf P^1$, and exactly one point $q$ identifies two distinct
normalization points $\alpha,\beta$. Other singular points are
unibranch and do not change this topological quotient. The universal
cover of the image is an infinite chain of normalized spheres $C_j$,
with $\beta_j$ identified with $\alpha_{j+1}$.

The actual map $h|_{D_{\rm red}}$ lifts to this cover because
$D_{\rm red}$ is simply connected. Each noncontracted rational
component factors through the normalization of the image and maps
surjectively onto it; its lift therefore covers one whole $C_j$.
Each contracted component lifts to a point. There are finitely many
components, so their connected lifted image is a finite chain
$C_m,\ldots,C_M$.

Choose a preimage in $D_{\rm red}$ of the left endpoint node and
let $t_-$ be its image under $\pi$ in the normalization $T$ of $S$.
The resolution fibre $\pi^{-1}(t_-)$ is connected and lies in
$D_{\rm red}=h^{-1}(\Gamma_{\rm red})$. Its lift is a single point
of the discrete fibre over $q$. Every local conductor branch of $T$
through $t_-$ has a strict transform meeting this resolution fibre.
It can receive only the branch directed into $C_m$: receiving the
other branch would force the lifted image to contain $C_{m-1}$.
At the right endpoint the same argument produces a point $t_+$
receiving only the opposite branch. The two points are distinct,
because their connected resolution fibres lift to different endpoint
nodes, including when $m=M$.

Distinct normalization points over $q$ correspond to distinct local
irreducible hypersurface branches of $S$ in the smooth threefold.
Their two distinct equations have an intersection curve: the ideal
they generate has height two. Every such curve is an actual
conductor-image branch received by both normalization points.
There are only the two image branches at $q$, whereas $t_-$ and
$t_+$ receive disjoint ones. This contradiction excludes the
one-component, first-Betti-number-one image.

Since every image component is rational-normalized, write
$n_\Gamma$ for their number and $\beta_\Gamma\le1$ for the
first Betti number. Their exact topological normalization
formula now yields, on these negative boundary classes,
\[
 e(\Gamma_{\rm red})=n_\Gamma+1-\beta_\Gamma\ge2.
 \tag{A.RT.21}\hypertarget{tag-007E}{}\label{cdp:eq:A.RT.21}
\]
For one component the excluded case is the only value below
two; for at least two components the inequality follows
directly from $\beta_\Gamma\le1$.




\begin{proposition}[A homogeneous binary-product grid is excluded]\hypertarget{tag-01BC}{}\label{cdp:st:A.5.21}
\StatementRecord{cdp:st:A.5.21}{proved}{01BC}{\SeriesRef{cdp:st:A.5.15}, \SeriesRef{cdp:eq:A.BP.2}, \SeriesRef{cdp:eq:A.BP.8}, \SeriesRef{cdp:eq:A.BP.6}}

Let $f:Y\to C_0$ be a proper holomorphic map from a smooth
complex threefold with smooth nearby fibres complex
two-tori. Let $P(X,V)$ be a homogeneous polynomial of
degree $d\ge2$, a product of distinct complex linear
forms, and suppose $P(1,0)\ne0$.
There is no point at which the actual fibre map has
the local equation
\[
             f=u(x,y,z)\,P(x,yz)^m,
             \qquad m\ge1,                         \tag{A.BP.1}\hypertarget{tag-004W}{}\label{cdp:eq:A.BP.1}
\]
with $u$ a holomorphic unit. In particular this excludes
$\prod_{i=1}^d(x-a_i yz)$ for distinct constants $a_i$,
and every positive full-fibre multiplicity of that germ.
No Betti-number assumption on $Y$ is required.

The same exclusion holds for the variable-coefficient
equation $f=u h^m$ with
\[
 h=\sum_{i=0}^d a_i(x,y,z)x^{d-i}(yz)^i,
 \qquad P_0(X,V)=\sum_{i=0}^d a_i(0)X^{d-i}V^i,
                                                        \tag{A.BP.4}\hypertarget{tag-0052}{}\label{cdp:eq:A.BP.4}
\]
provided $P_0$ has $d$ distinct projective roots and
$P_0(1,0)\ne0$. The comparison with its constant
initial form is the explicit analytic family below.
\end{proposition}


\begin{proof}

\begin{enumerate}[label={$\langle1\rangle$\arabic*.},leftmargin=*]
\item Put $h=P(x,yz)$. It is weighted homogeneous of degree
$2d$ for the positive weights $(2,1,1)$.
On a small simply connected coordinate neighbourhood
take a holomorphic logarithm of the unit and set
\[
 \Phi(x,y,z)=\bigl(xu^{1/(md)},
                  yu^{1/(2md)},zu^{1/(2md)}\bigr).
                                                        \tag{A.BP.2}\hypertarget{tag-0050}{}\label{cdp:eq:A.BP.2}
\]
Its derivative at the origin is invertible, and direct
homogeneity gives $h\Phi=u^{1/m}h$. Thus
$f=(h\Phi)^m$ exactly. It suffices to construct a
compact spherical cycle of nonzero square in $h=1$,
because weighted homotheties then put it inside an
arbitrarily small actual fibre of \eqref{cdp:eq:A.BP.1}.

\item The curve $B=\{P(x,v)=1\}\subset\C^2$ is smooth:
Euler's identity $xP_x+vP_v=dP$ precludes a common
zero of its two partial derivatives on that level.
It is also connected. The map
\[
 B\longrightarrow\PP^1\setminus\{P=0\},
                       \quad(x,v)\longmapsto[x:v]
\]
is an etale cover of degree $d$. Locally a projective
representative has nonzero $P$-value, and its $d$
rescalings with $P=1$ give the sheets. Around any
one of the simple projective zeros of $P$, that
value has a simple zero, so its inverse $d$th root
has monodromy a full $d$-cycle. This monodromy is
transitive on the sheets, proving connectedness.

\item The function $v$ on $B$ has exactly $d$ distinct
simple zeros, since
$P(x,0)=P(1,0)x^d$ and $P_x(x,0)=d/x$ at those
points. Choose an embedded arc $\gamma$ in $B$
joining two of its zeros and avoiding all other
zeros. On the interior choose a continuous argument
of $v(\gamma(t))$ and two functions $A(t),B_0(t)$
with $A B_0=v(\gamma(t))$ and
$|A|=|B_0|=|v(\gamma(t))|^{1/2}$.
The map
\[
 (t,\theta)\longmapsto
   \bigl(x(\gamma(t)),A(t)e^{\mathrm i\theta},
                        B_0(t)e^{-\mathrm i\theta}\bigr)
                                                        \tag{A.BP.3}\hypertarget{tag-0051}{}\label{cdp:eq:A.BP.3}
\]
lies in $h=1$ and collapses both boundary circles
of $[0,1]\times S^1$ to points. It is therefore a
spherical cycle, with embedded interior circles.

\item A parallel-arc isotopy with the same endpoints gives
a homologous spherical cycle whose support meets
the first only at those endpoints. At each endpoint,
$P_x\ne0$, so $(y,z)$ are actual holomorphic
coordinates on $h=1$. The links of the two cycles
are parallel curves of slope $(1,-1)$ on a Clifford
torus in the local three-sphere. Their linking
number is $-1$. As in the explicit relative-chain
comparison in \ref{cdp:st:A.5.15}, that number is the local
intersection of the two coned cycles. Both
orientations reverse at the second endpoint,
so its intersection sign is unchanged. The
spherical class consequently has square $-2$.

\item The weighted homothety
$(x,y,z)\mapsto(\lambda^2x,\lambda y,\lambda z)$
maps $h=1$ biholomorphically to $h=\lambda^{2d}$
and shrinks the compact cycles and their isotopy
to the origin. Together with \eqref{cdp:eq:A.BP.2}, this puts
them inside an actual smooth global fibre of $f$.
That fibre is a complex two-torus and has zero
$\pi_2$, so the spherical class has zero homology
image. The oriented codimension-zero embedding
preserves the intersection of the compact cycles,
giving both zero and $-2$ for its square. This
contradiction proves the statement.

\item We now prove the variable-coefficient assertion,
including the boundary comparison. Absorb a local
$m$th root of the unit into $h$; this preserves the
form \eqref{cdp:eq:A.BP.4} and its two initial-form conditions.
For a complex parameter $s$ set
\[
 h_s=\sum_{i=0}^d a_i(s^2x,sy,sz)x^{d-i}(yz)^i.
                                                        \tag{A.BP.5}\hypertarget{tag-0053}{}\label{cdp:eq:A.BP.5}
\]
This is holomorphic also at $s=0$, with
$h_0=P_0(x,yz)$. For $s\ne0$ there is the exact
identity
\[
       h(s^2x,sy,sz)=s^{2d}h_s(x,y,z).           \tag{A.BP.6}\hypertarget{tag-0054}{}\label{cdp:eq:A.BP.6}
\]
The leading coefficient $a_0$ is a unit. The implicit
function theorem applied to its degree-$d$ root
polynomial gives distinct holomorphic functions
$b_i(s,x,y,z)$ and the exact factorization
\[
       h_s=a_0(s^2x,sy,sz)
                     \prod_{i=1}^d(x-b_i(s,x,y,z)yz)
                                                        \tag{A.BP.7}\hypertarget{tag-0055}{}\label{cdp:eq:A.BP.7}
\]
on one small parameter and coordinate neighbourhood.
Each factor defines a smooth surface, since its
$x$ derivative is a unit. For $i\ne j$ the ideal
generated by the two factors equals $(x,yz)$:
their difference is the unit $b_j-b_i$ times $yz$.
Thus the zero level has precisely the same two
critical curves $D=(x,yz)=0$ for every $s$.

\item Choose a small closed ball $K$ in the coordinate
chart such that the two axes of $D$ meet its boundary
transversely and the zero-level surfaces away from
those axes are transverse to the boundary. At
$s=0$ this follows directly for the smooth graph
surfaces $x=b_i(0)yz$ and sufficiently small radii.
It persists for small $s$ by compactness away from
neighbourhoods of the two boundary links.
Near either link, use normal coordinates $(U,V)$
for the smooth critical curve and a parameter $w$
along it. The normal initial term of $h_s$ is a
degree-$d$ binary form with distinct roots.
On the compact parameterized link its normal
gradient therefore satisfies a uniform bound
\[
 \|\partial_{U,V}h_s\|\ge c\|(U,V)\|^{d-1},
 \qquad |\partial_w h_s|\le C\|(U,V)\|^d
                                                        \tag{A.BP.8}\hypertarget{tag-0056}{}\label{cdp:eq:A.BP.8}
\]
after the neighbourhood is shrunk. The first bound
follows because the gradient of a homogeneous
reduced binary form has no zero on the unit sphere;
higher normal terms have one additional order.
The second follows from $h_s\in(U,V)^d$.

\item Transversality of the critical curve to $\partial K$
allows a boundary coordinate with one real component
of $w$ solved as a smooth function of $(U,V)$ and
the remaining component. The correction to normal
differentiation tangent to this boundary is a
bounded multiple of $\partial_w h_s$.
By \eqref{cdp:eq:A.BP.8} it is smaller by one normal order than
the surjective complex normal derivative. Thus
$h_s|_{\partial K}$ is a real submersion at every
nearby point with $h_s\ne0$, uniformly in $s$.
Away from the two link neighbourhoods, compactness
and the preceding zero-level transversality give
the same conclusion for all sufficiently small
nonzero values.

\item Remove the proper critical-value discriminant from
the product of the parameter disc and a sufficiently
small punctured value disc. The critical locus over
that domain cannot meet the boundary, so its
projection is proper and its image is complex
analytic. Its complement is connected. The level
sets in $K$ form a proper submersion with boundary
there, and Ehresmann transport identifies their
compact fibres by oriented diffeomorphisms.
Choose a small sphere cycle of the already proved
model $h_0=P_0(x,yz)$, and transport it to a fibre
of $h_{s_0}$ for $s_0\ne0$. Equation \eqref{cdp:eq:A.BP.6}
then embeds that fibre into the actual fibre
of $h^m$ with value $s_0^{2dm}\epsilon^m$.
It carries the spherical class of square $-2$
into a global torus. The same intersection and
asphericity contradiction proves \eqref{cdp:eq:A.BP.4}.

\end{enumerate}
\end{proof}



The condition $P(1,0)\ne0$ is used to produce the
distinct collapsing circles. The exact product
model (A.U.6), for which $P(X,V)=XV(X-V)$, has
$P(1,0)=0$ and is outside this statement.
The variable-coefficient extension uses only the
specific comparison \eqref{cdp:eq:A.BP.5}--\eqref{cdp:eq:A.BP.8}, with its
unchanged boundary critical curves. The weighted scaling,
simple-zero and endpoint-square identities are
checked in \path{code/E_binary_products.py}.

\paragraph{Two exact scalar-zero comparisons.}
The homogeneous germ
\[
                  h=xy\{z(x-y)+xy\}
\tag{A.BP.9}\hypertarget{tag-0057}{}\label{cdp:eq:A.BP.9}
\]
has a cubic normal tensor along $(x,y)=0$ which
vanishes identically at $z=0$. On $xy\ne0$ put
\[
 v=x-y,\qquad Z=y(y+v)z+y^2(2y+v).
\]
These are actual coordinates with $y$, with inverse
$x=y+v$, $z=(Z-y^2(2y+v))/(y(y+v))$, and
\[
                              h=vZ+y^4.
\tag{A.BP.10}\hypertarget{tag-004X}{}\label{cdp:eq:A.BP.10}
\]
Choose an arc between two roots of $y^4=1$ avoiding
$y=0$ and the other roots. The matching circles
$v=\rho A e^{\mathrm i\theta}$,
$Z=\rho^{-1}B e^{-\mathrm i\theta}$, where
$AB=1-y^4$ and $|A|=|B|=|1-y^4|^{1/2}$,
collapse at its endpoints. Choose the positive
constant $\rho$ small enough that $|v|<|y|/2$
on the arc. The entire spherical cycle then lies
in the indicated coordinate chart. Its two endpoint
linkings are $-1$, as in \eqref{cdp:eq:A.BP.3}, so its square
is $-2$. Ordinary homothety in $(x,y,z)$ shrinks
it because \eqref{cdp:eq:A.BP.9} has degree four. The same
weighted unit map as \eqref{cdp:eq:A.BP.2}, with all three
weights equal to one, includes every local unit
and every positive full-fibre multiplicity.
Thus this particular scalar-zero germ is excluded.

The following different scalar-zero model records
the limit of that argument:
\[
                     h=xz(y^2-x^2z).
\tag{A.BP.11}\hypertarget{tag-004Y}{}\label{cdp:eq:A.BP.11}
\]
For $h=t\ne0$, the substitution $u=y/x$ is
invertible and gives $h=x^3 z(u^2-z)$.
The map forgetting $x$ is a finite etale triple
cover of $E=\{z(u^2-z)\ne0\}$.
There is an explicit finite etale double cover
\[
 \C^*\times(\C\setminus\{1,-1\})\longrightarrow E,
                  \qquad(s,v)\longmapsto(s^2,sv)
\tag{A.BP.12}\hypertarget{tag-004Z}{}\label{cdp:eq:A.BP.12}
\]
in coordinates $(z,u)$, with deck involution
$(s,v)\mapsto(-s,-v)$. Its source is aspherical,
so both $E$ and the regular affine fibre of
\eqref{cdp:eq:A.BP.11} are aspherical by these actual covering
maps. The quadratic surface branch has normalization
$(v,u)\mapsto(v,vu,u^2)$ and conductor
$2(v=0)+2(u=0)$ in the full reduced surface.
The other two normalization branches have conductors
$2(y=0)+(z=0)$ on $x=0$, and
$(x=0)+2(y=0)$ on $z=0$.
Thus their reduced conductor supports have only
normal-crossing pairs. This is an exact local
model, not a construction of a compact $b_2=0$
fibre. It shows why a vanishing cubic coefficient
alone does not justify a local spherical-cycle
exclusion. The identities and the two finite-cover
formulas are checked in \path{code/E_cubic_double.py}.



\section{Normal-crossing topology and canonical-flat component comparisons}\hypertarget{tag-0344}{}\label{sec:II:08}


\begin{proposition}[The normalized double curves and the coherent Euler formula]\hypertarget{tag-01C9}{}\label{cdp:st:C.2.NC.2}
\StatementRecord{cdp:st:C.2.NC.2}{proved}{01C9}{\SeriesRef{cdp:eq:C.NC.6}, \SeriesRef{cdp:eq:C.NC.7}}

Let $S$ be a connected compact normal-crossings divisor in a smooth compact complex threefold $Y$ with $H^2(Y,\Q)=0$. Write $\nu:\coprod_iT_i\to S$ for normalization, $D_i\subset T_i$ for its conductor, $C_\alpha$ for the connected components of the normalization of the double locus, and $\tau$ for the number of triple points. Let $R_\alpha$ be the reduced divisor on $C_\alpha$ of preimages of triple points, counted separately for the three double-curve branches at each triple point. Then
\[
 e(C_\alpha)=\deg R_\alpha,\qquad
 \sum_\alpha e(C_\alpha)=3\tau,
 \qquad
 \sum_i\chi(T_i,\OO_{T_i})=\frac{\tau}{2}.
 \tag{C.NC.4}\hypertarget{tag-00S3}{}\label{cdp:eq:C.NC.4}
\]
In particular every normalized double curve has genus at most one, and $\tau$ is even. The remaining topological identities are
\[
\begin{aligned}
 \sum_i e(D_i)&=2\sum_\alpha e(C_\alpha)-3\tau,\\
 e(S)&=\sum_i e(T_i)-\sum_\alpha e(C_\alpha)+\tau.
\end{aligned}
 \tag{C.NC.5}\hypertarget{tag-00S4}{}\label{cdp:eq:C.NC.5}
\]
\end{proposition}


\begin{proof}

\begin{enumerate}[label=\arabic*.,leftmargin=*]
\item Adjunction and the local residue comparison give a canonical holomorphic line-bundle isomorphism
\[
 \nu^*\omega_S|_{T_i}\simeq\omega_{T_i}(D_i),
 \qquad \omega_S=(\omega_Y\otimes\OO_Y(S))|_S.
 \tag{C.NC.6}\hypertarget{tag-00S5}{}\label{cdp:eq:C.NC.6}
\]
For example, the residue of $dx\wedge dy\wedge dz/(xy)$ on the branch $x=0$ is $\pm dy\wedge dz/y$, and at $xyz=0$ it is $\pm dy\wedge dz/(yz)$. These expressions agree under changes of local equations by the usual determinant and unit factors of the displayed line bundles. Since both line bundles on the right in the expression for $\omega_S$ come from $Y$, its first Chern class and that of $\omega_{T_i}(D_i)$ vanish rationally.

\item Normalize every irreducible curve in every $D_i$. The resulting disjoint union $\widehat D$ maps to $\coprod_\alpha C_\alpha$ by an étale map of total degree two over each $C_\alpha$. Away from a triple point this follows from the two branches of $xy=0$. At a triple point, each of the three normalized double branches has exactly two preimages, and their maps in the remaining coordinate are local isomorphisms. The covering over a given $C_\alpha$ may be connected.

\item Adjunction for a nodal divisor on a smooth surface, pulled back to $\widehat D$, identifies the restriction of $\omega_{T_i}(D_i)$ with the canonical line of that normalized conductor curve twisted by all its preimages of nodes of $D_i$. Such nodes are exactly the preimages of triple points of $S$. This is the elementary local dualizing-form identity $dw/w$ on each branch of a node. The resulting degree is zero by \eqref{cdp:eq:C.NC.6}. Since the covering of $C_\alpha$ is étale and its marked divisor is the pullback of $R_\alpha$, it follows that
$\deg\omega_{C_\alpha}+\deg R_\alpha=0$.
Thus $e(C_\alpha)=\deg R_\alpha$. There are exactly three marks over every triple point, even when two or three belong to one $C_\alpha$.

\item For the coherent calculation, the alternating branch-resolution sequence is
\[
 0\longrightarrow\OO_S\longrightarrow\nu_*\OO_{\coprod T_i}
 \longrightarrow\bigoplus_\alpha(j_\alpha)_*\varepsilon_\alpha
 \longrightarrow\bigoplus_{p\in S^{[3]}}\C_p
 \longrightarrow0.
 \tag{C.NC.7}\hypertarget{tag-00S6}{}\label{cdp:eq:C.NC.7}
\]
Here $j_\alpha:C_\alpha\to S$ is the finite normalization map of that double curve, and $\varepsilon_\alpha$ is the sign line of its two-sheeted conductor covering. Locally order the branches and send $(f_i)$ to $(f_j-f_i)_{i<j}$; the next map is the alternating sum of the three values at a triple point. The sequence is exact at $xy=0$ because a pair of functions glues precisely when its restrictions agree, and at $xyz=0$ because pairwise differences satisfy exactly the single alternating relation. Equivalently, compare coefficients of convergent power series on the three coordinate planes; equality on their coordinate axes is exactly the condition for these series to come from one series modulo $xyz$. Changing a local branch ordering multiplies the appropriate difference by $-1$, producing $\varepsilon_\alpha$. At a triple point this sign convention only chooses a basis of the one-dimensional last term. Hence the sequence is globally well-defined with the indicated local systems, including self-intersections.

\item The sign line satisfies $\varepsilon_\alpha^{\otimes2}\simeq\OO_{C_\alpha}$, so it has degree zero and $\chi(C_\alpha,\varepsilon_\alpha)=1-g(C_\alpha)=e(C_\alpha)/2$. The divisor sequence on $Y$ and Hirzebruch--Riemann--Roch give $\chi(S,\OO_S)=0$, since $c_1(\OO_Y(S))=0$ rationally. Taking Euler characteristics in \eqref{cdp:eq:C.NC.7} now gives
\[
 0=\sum_i\chi(T_i,\OO_{T_i})-\frac12\sum_\alpha e(C_\alpha)+\tau.
\]

\item $\widehat D$ has Euler number $2\sum_\alpha e(C_\alpha)$. Each triple point produces three nodes in $\coprod D_i$, and identifying the two preimages of each node subtracts one. The double locus in $S$ has normalization $\coprod C_\alpha$ and three preimages over each triple point, so its Euler number is $\sum_\alpha e(C_\alpha)-2\tau$. Normalization is an isomorphism off the conductor. Additivity of Euler characteristic therefore gives the second identity in \eqref{cdp:eq:C.NC.5}.

\end{enumerate}
\end{proof}






\begin{proposition}[Uniform multiple normal-crossings fibres]\hypertarget{tag-01CA}{}\label{cdp:st:C.2.NC.3}
\StatementRecord{cdp:st:C.2.NC.3}{proved}{01CA}{the definitions and constructions stated in the proof}

Let $g:Y\to C$ be a proper holomorphic map from a connected smooth compact complex threefold with $H^2(Y,\Q)=0$ to a smooth compact curve, with connected fibres and smooth general fibre of Euler characteristic zero. A fibre $mS$ with reduced normal-crossings support has
\[
 e(S)=\tau(S)\ge0.
 \tag{C.NC.8}\hypertarget{tag-00S7}{}\label{cdp:eq:C.NC.8}
\]
The normalized base change $t=s^m$ is finite étale of degree $m$ over a neighbourhood of $S$ and has reduced normal-crossings central fibre $S'$. The induced cover of reduced fibres satisfies
$e(S')=m e(S)$ and $\tau(S')=m\tau(S)$.
\end{proposition}


\begin{proof}

\begin{enumerate}[label=\arabic*.,leftmargin=*]
\item The local formula \CompendiumRef{II}{cdp:st:C.2.NC.1} gives
\[
 0=e(Y_t)=m\sum_i e(T_i\setminus D_i).
\]
By \eqref{cdp:eq:C.NC.5},
\[
 e(S)-\sum_i e(T_i\setminus D_i)
 =\sum_\alpha e(C_\alpha)-2\tau=\tau.
\]

\item Locally write $t=u(z_1\cdots z_r)^m$ with $u$ nowhere zero. The normalization of $s^m=u(z_1\cdots z_r)^m$ is given by
$w^m=u$ and $s=wz_1\cdots z_r$.
The derivative of $w^m-u$ with respect to $w$ is nonzero because $w$ is a unit. Thus its map to the original chart is étale of degree $m$, and the lifted central fibre has equation $z_1\cdots z_r=0$. These local normalizations glue by the universal property of normalization; their degree and étaleness are local properties. Restriction is therefore an étale degree-$m$ covering $S'\to S$, and each triple point has exactly $m$ preimages. A lifted triangulation proves the Euler-number comparison. If the cover has several connected components, the degrees and Euler numbers add over them.

\end{enumerate}
\end{proof}






\begin{proposition}[Reducible normal-crossings fibres with unequal multiplicities]\hypertarget{tag-01CB}{}\label{cdp:st:C.2.NC.4}
\StatementRecord{cdp:st:C.2.NC.4}{proved}{01CB}{\SeriesRef{cdp:st:C.2.NC.3}, \SeriesRef{cdp:st:C.2.NC.2}; \cite[Theorem~2.2, p.~843]{Peters2021}; \cite[Theorem~3.2, p.~16]{BMS2012}}

Under the hypotheses on $g:Y\to C$ in \ref{cdp:st:C.2.NC.3}, let a fibre be $\sum_i m_iS_i$ with connected reducible normal-crossings support $S$. Its multiplicities may be unequal. Then
\[
 \tau(S)=0,\qquad e(S)=0.
 \tag{C.NC.9}\hypertarget{tag-00S8}{}\label{cdp:eq:C.NC.9}
\]
Every normalized conductor is a disjoint union of smooth elliptic curves, and every normalization component has Euler number zero. More generally, any singular normal-crossings fibre support under these hypotheses has either no triple point or exactly two triple points; in the latter case it is irreducible and its normalization is rational.
\end{proposition}


\begin{proof}

\begin{enumerate}[label=\arabic*.,leftmargin=*]
\item First consider a normalization component $T_i$ whose conductor $D_i$ is nonempty. If $T_i$ is Kähler and has nonnegative Kodaira dimension, choose a nonzero section of $\omega_{T_i}^{\otimes n}$ and denote its effective zero divisor by $E$. The effective divisor $E+nD_i$ has rational Chern class zero by \eqref{cdp:eq:C.NC.6}. Its integral against a Kähler form is positive because $D_i$ is nonempty, whereas its cohomology class gives integral zero. This is impossible. By \cite[Theorem~2.2, p.~843]{Peters2021}, a Kähler $T_i$ is consequently rational or a blown-up ruled surface. For a ruling $T_i\to B$ and a general ruling fibre $F$, one has $K_{T_i}\cdot F=-2$, hence $D_i\cdot F=2$. Some conductor curve therefore dominates $B$. Its normalization has genus at most one by \ref{cdp:st:C.2.NC.2}, so Riemann--Hurwitz forces $g(B)\le1$. The Kähler possibilities have
$\chi(\OO_{T_i})=1$ in the rational case and $0$ in the elliptic ruled case; their Euler numbers are nonnegative.

\item For a non-Kähler $T_i$, the Enriques--Kodaira table \cite[Theorem~2.2, p.~843]{Peters2021} leaves class VII, Kodaira and properly elliptic minimal models. Class VII has $c_1^2=-b_2$ and $c_2=b_2$, so Noether's formula gives $\chi(\OO)=0$. Kodaira surfaces have both Chern numbers zero. A minimal non-Kähler elliptic surface has topological Euler number zero by Kodaira's theorem in the precise form \cite[Theorem~3.2, p.~16]{BMS2012}; its canonical square is zero by the same classification table, so Noether's formula again gives $\chi(\OO)=0$. Blowing up preserves the holomorphic Euler characteristic and adds one to the topological Euler characteristic. Thus every non-Kähler $T_i$ has $\chi(\OO_{T_i})=0$ and $e(T_i)\ge0$.

\item A rational normalization component cannot meet a different global component of $S$. Indeed, for $j\ne i$, pull back the canonical section of $\OO_Y(S_j)$ to $T_i$. It is a nonzero section whose zeros are precisely the inverse image of $S_i\cap S_j$. Its line bundle has zero rational first Chern class. For a smooth rational surface $T_i$, $H^1(T_i,\OO)=0$ and $H^2(T_i,\Z)$ is torsion-free, by the corresponding facts for $\PP^2$ or a Hirzebruch surface and the blowup formula. The exponential sequence therefore makes that pulled-back line bundle trivial. A nonzero section of a trivial line bundle on a connected compact surface is a nonzero constant and has no zeros. Hence $S_i\cap S_j$ is empty.

\item The intersection graph of the global components of a connected reducible fibre is connected. Each component has nonempty conductor, and the preceding paragraph excludes rational normalizations. Thus every $\chi(\OO_{T_i})$ is zero, and \eqref{cdp:eq:C.NC.4} gives $\tau=0$. Each $C_\alpha$ is now elliptic and has no marked point. The conductor covering is étale, so the normalized conductor components are elliptic; there are no nodes of $D_i$, since these would be preimages of triple points. Consequently $e(D_i)=0$. The nearby Euler function is $m_i$ along the smooth part of $S_i$. At a double crossing the local map is $x^{m_i}y^{m_j}$, with equal exponents for a self-intersection of one global component; its transverse nearby fibre is a disjoint union of annuli and has Euler number zero. Euler integration gives $e(Y_t)=\sum_i m_i e(T_i\setminus D_i)$. Since $e(D_i)=0$, this is
\[
 0=e(Y_t)=\sum_i m_i e(T_i).
\]
All summands are nonnegative and every $m_i$ is positive, so each $e(T_i)=0$. Formula \eqref{cdp:eq:C.NC.5} gives $e(S)=0$.

\item For a singular irreducible support the conductor is nonempty and the same classification argument gives $\chi(\OO_{T_1})\in\{0,1\}$, with value one exactly in the rational case. Formula \eqref{cdp:eq:C.NC.4} then gives $\tau=0$ or $2$, respectively. Thus any triple point lies on an irreducible support with rational normalization.

\end{enumerate}
\end{proof}






\begin{proposition}[Global positivity and the all-twist normalization obstruction in the normal-crossings class]\hypertarget{tag-01CC}{}\label{cdp:st:C.2.NC.5}
\StatementRecord{cdp:st:C.2.NC.5}{proved}{01CC}{\SeriesRef{cdp:st:C.2.NC.3}, \SeriesRef{cdp:st:C.3.5}, \SeriesRef{cdp:st:C.3.1}, \SeriesRef{cdp:st:C.2.NC.4}, \SeriesRef{cdp:st:A.1.2}}

Let $g:Y\to C$ have the hypotheses of \ref{cdp:st:C.2.NC.3}, and assume every fibre reduction has normal crossings; smooth reductions are included. Then
\[
 c_3(Y)=e(Y)=\sum_{c\in C}\tau((Y_c)_{\rm red})=2r\ge0,
 \tag{C.NC.10}\hypertarget{tag-00S1}{}\label{cdp:eq:C.NC.10}
\]
where $r$ is the number of fibre supports with rational normalization. Every such support is irreducible, nonnormal, and has exactly two triple points. Equality $c_3(Y)=0$ holds exactly when every fibre support is free of triple points. If also $b_1(Y)=0$, then $r\in\{0,1\}$ and $b_3(Y)=2-2r$.

In the positive case there is an irreducible nonnormal reduced support $S$ with rational normalization $\eta:T\to S$ such that, for every $L\in\Pic(Y)$,
\[
 H^2\!\left(T,K_T\otimes\eta^*((K_Y^{-1}\otimes L)|_S)\right)\simeq\C,
 \qquad
 H^0\!\left(S,\Omega_Y^1|_S\otimes(K_Y\otimes L^{-1})|_S\right)\simeq\C.
 \tag{C.NC.11}\hypertarget{tag-00S2}{}\label{cdp:eq:C.NC.11}
\]
The scheme-theoretic fibre may be a multiple of $S$. By \ref{cdp:st:C.3.5}, all these conclusions apply to every fibration under the CDP hypotheses whose fibre reductions have normal crossings, for every smooth-fibre type in \ref{cdp:st:C.3.1}.
\end{proposition}


\begin{proof}

\begin{enumerate}[label=\arabic*.,leftmargin=*]
\item There are finitely many critical values: the critical locus is analytic, its image is analytic by properness, and it does not contain a general point of the curve. Away from these values $g$ is a differentiable fibre bundle whose fibre has Euler number zero. Euler characteristic is additive over this open base and the finite exceptional set, so $e(Y)=\sum_c e((Y_c)_{\rm red})$. Nilpotents leave the underlying space unchanged. An irreducible support has a single multiplicity, so \ref{cdp:st:C.2.NC.3} applies; a reducible support is covered by \ref{cdp:st:C.2.NC.4}. A smooth support is connected and irreducible, so the same uniform-multiplicity formula gives Euler number zero. Thus every summand equals its triple-point count. \ref{cdp:st:C.2.NC.4} identifies the positive summands and their value two. Gauss--Bonnet gives $c_3(Y)=e(Y)$. Poincaré duality gives $e(Y)=2-2b_1(Y)+2b_2(Y)-b_3(Y)$. Under the additional hypothesis $b_1(Y)=0$, this gives $b_3(Y)=2-2r\ge0$, hence $r\in\{0,1\}$. Under the CDP hypotheses, \ref{cdp:st:C.3.5} gives Euler number zero for every smooth fibre.

\item For a positive support, every line bundle on $Y$ pulls back trivially to its smooth rational normalization: its rational first Chern class is zero, and the exponential sequence identifies $\Pic(T)$ with the torsion-free group $H^2(T,\Z)$ by the first Chern class, as in \ref{cdp:st:C.2.NC.4}. In particular $\eta^*((K_Y^{-1}\otimes L)|_S)$ is trivial. Serre duality on the connected smooth compact surface $T$ identifies the first group in \eqref{cdp:eq:C.NC.11} with $H^0(T,\OO_T)^\vee\simeq\C$.

\item For the ambient group set $A=(K_Y\otimes L^{-1})|_S$ and $N^*=\OO_Y(-S)|_S$. Both $\eta^*A$ and $\eta^*N^*$ are trivial, even when the full fibre is $mS$. The differential pullback has kernel $\operatorname{Tor}(\Omega_S^1)$ and induces $\widetilde\Omega_S^1\hookrightarrow\eta_*\Omega_T^1$. Tensoring by $A$ and applying the projection formula gives
$H^0(S,\widetilde\Omega_S^1\otimes A)\hookrightarrow H^0(T,\Omega_T^1)=0$.
Put $K=\ker(\Omega_Y^1|_S\to\widetilde\Omega_S^1)$. The exact sequence in \ref{cdp:st:A.1.1} and the isomorphism $K\simeq\eta_*\eta^*N^*$ of \ref{cdp:st:A.1.2}, followed by the projection formula, give
\[
 H^0(S,\Omega_Y^1|_S\otimes A)
 =H^0(T,\eta^*(N^*\otimes A))=\C.
\]


\end{enumerate}
\end{proof}






\begin{proposition}[The normal jets along cusp, tacnode and triple curves]\hypertarget{tag-01FJ}{}\label{cdp:st:C2d}
\StatementRecord{cdp:st:C2d}{proved}{01FJ}{\SeriesRef{cdp:eq:C.J.2}, \SeriesRef{cdp:eq:C.J.3}, \SeriesRef{cdp:st:A.2.4}}

Let $Y$ be a smooth complex threefold with $b_2(Y)=0$, let $S$ be a reduced Cartier divisor, and let $D\subset\operatorname{Sing}S$ be a smooth compact connected curve. Suppose the general plane section transverse to $D$ has one of the equations
\begin{equation}
 x^2-y^3=0,\qquad x^2-y^4=0,\qquad xy(x-y)=0.
 \hypertarget{tag-00OV}{}\label{cdp:eq:C.J.1}
\end{equation}
Then $g(D)\le1$. For $g(D)=1$, the indicated transverse analytic type is constant along all of $D$. The inverse image of $D$ in the normalization is respectively an isomorphic copy, a finite \'etale cover of degree two, or a finite \'etale cover of degree three. Each connected component $C$ of that cover has genus one and normal line of degree zero in the normalization, which is smooth near $C$. At these three transverse types, the inverse Jacobian ideal pulls back into the line allowing at most a simple pole on each normalization branch.
\end{proposition}


\begin{proof}

\begin{enumerate}[label={$\langle1\rangle$\arabic*.},leftmargin=*]
\item Put $N=N_{D/Y}$ and $B=\OO_Y(S)|_D$. Since every ambient line bundle has zero rational first Chern class, $\deg B=\deg K_Y|_D=0$. Adjunction for the smooth curve gives
\begin{equation}
 \det N=K_D\otimes K_Y^{-1}|_D,\qquad
 \deg N=2g(D)-2.                                      \hypertarget{tag-00OW}{}\label{cdp:eq:C.J.2}
\end{equation}
A local defining section of $S$ belongs to $I_D^2$. Its quadratic normal jet is a well-defined section of $\operatorname{Sym}^2N^*\otimes B$. In the first two cases its generic rank is one. Its saturated radical $R\subset N$ is a line subbundle on the smooth curve, and $Q=N/R$ is a line bundle. The quadratic jet gives a nonzero section of $B\otimes Q^{-2}$, so
$\deg Q\le0$.

\item In the cusp case, the cubic jet on the radical is an intrinsic nonzero section of $B\otimes R^{-3}$. To check its independence, choose coordinates $(x,y,z)$ with $D=(x=y=0)$ and $R$ the $y$ direction. A change of lift of a radical vector changes the cubic by a Hessian term with one radical argument, which is zero. Changing the local retraction to $D$ differentiates the identity that the quadratic jet vanishes on $R$, again giving zero. Thus the local cubics glue across all points, including possible zeros of the quadratic coefficient. Consequently $\deg R\le0$. Equation \eqref{cdp:eq:C.J.2} proves $g(D)\le1$.

\item In the tacnode case the radical cubic vanishes identically. Write the relevant normal terms as
\[
 a x^2+bxy^2+cy^4.
\]
Completing the square has invariant numerator
\begin{equation}
 \Delta=ac-b^2/4\in
 H^0(D,B^2\otimes Q^{-2}\otimes R^{-4}).               \hypertarget{tag-00OX}{}\label{cdp:eq:C.J.3}
\end{equation}
This numerator is holomorphic even at a zero of $a$. Under $x\mapsto x+uy^2$, the coefficients become $b+2au$ and $c+bu+au^2$, and \eqref{cdp:eq:C.J.3} is unchanged. Linear changes preserving the radical give precisely the displayed line-bundle weights. Changes of the retraction to $D$ add terms containing the radical cubic or a radical Hessian argument, so also leave this quartic numerator unchanged. These checks give its global definition. It is nonzero by the generic tacnode assumption. Hence
\[
 \deg Q\le0,\qquad \deg Q+2\deg R\le0,
 \qquad \deg N\le\tfrac12\deg Q\le0.
\]
This proves the same genus bound.

\item For the ordinary triple case the quadratic jet is identically zero. The binary cubic normal jet is a section of $\operatorname{Sym}^3N^*\otimes B$. Its discriminant is a nonzero section of
$B^4\otimes(\det N^*)^6$. The discriminant has degree four in the cubic coefficients and weight six under a linear change of the two normal coordinates; these weights also follow by writing it as the square of the product of the three pairwise determinants of the linear factors. Its nonnegative divisor degree gives $\deg N\le0$ and the genus bound again.

\item When $g(D)=1$, all the inequalities are equalities. In the cusp case both the quadratic and radical cubic sections have no zeros. In the tacnode case $a$ and $\Delta$ have no zeros. In the triple case the cubic discriminant has no zeros. The holomorphic splitting lemma with parameter $z$ completes the square in the first two cases. The remaining one-variable term is respectively a unit times $y^3$ or $y^4$; taking a local root of the unit gives the corresponding equation in \eqref{cdp:eq:C.J.1}. In the triple case its three distinct tangent lines split after an \'etale ordering cover; the holomorphic implicit function theorem gives three smooth branches meeting transversely in each normal plane. Straightening two branches to the axes and rescaling the third branch's unit gives the third equation. Thus these are constant transverse types throughout $D$. In particular no other singular curve meets $D$.

\item In the cusp normalization, $x=u^3,y=u^2$, up to the choice of which coordinate is the radical. The leading terms $y=u^2$ and $x=u^3$ identify
$P^2$ with $R$ and $P^3$ with $Q$, where
$P=N_{C/\widetilde S}$. Adapted coordinate changes preserve these
leading terms through the corresponding line transitions. Thus
$\deg P=0$. For the tacnode, the two branches have $x=\pm y^2$ and normal line the pullback of $R$, again of degree zero. Their permutation cover is \'etale because $\Delta$ is nowhere zero. For the ordinary triple, order the three tangent lines by a finite \'etale cover. Three pairwise transverse line subbundles in a rank-two bundle make it a sum $P\oplus P$: the third is a graph of an isomorphism between the first two. Since the determinant has degree zero, $\deg P=0$. Each normalization branch has this tangent-root normal line. \'{E}tale descent gives the claimed degrees and covers before ordering.

\item Compute the inverse Jacobian ideal in the parameter-preserving
normal forms, with coefficients holomorphic in the curve parameter $z$.
A biholomorphic coordinate change transforms the gradient by its
invertible Jacobian matrix; multiplication of the defining equation
by a unit multiplies the gradient by that unit modulo the equation.
These maps transport the following fractional ideals to $S$.

For the cusp set $B=\C\{u,z\}$ and
$A=\C\{u^2,u^3,z\}=\C\{z\}+u^2B$.
The Jacobian ideal has images $u^3,u^4$. For a multiplier $h$, put
$a=u^3h$. The conditions $a,ua\in A$ kill respectively the linear
and constant coefficients of $a$. Hence
$\mathcal J^{-1}=u^{-1}B$, agreeing with
Proposition~\ref{cdp:st:A.2.4} under the displayed normalization map.

For the tacnode the two maps are $x=\pm u^2,y=u$, and the image ring is
\[
 A=\{(f,g)\in B^2:f-g\in u^2B\}.
\]
Indeed, the inverse expression is
$(f+g)/2+x(f-g)/(2u^2)$. The Jacobian generators have images
$(u^2,-u^2)$ and $(u^3,u^3)$, up to nonzero scalar factors.
For $a=u^2h_1$ and $b=u^2h_2$ the membership conditions are
$a+b\in u^2B$ and $a-b\in uB$. Their constant coefficients vanish,
so each $h_i$ has at most a simple pole. Their residues are opposite:
\[
 \mathcal J^{-1}=B^2+\C\{z\}(u^{-1},-u^{-1}).
\]

For the ordinary triple point use the branches
$(0,u),(u,0),(u,u)$. The image ring consists of triples with a common
constant and with third linear coefficient equal to the sum of the
first two. All coefficients of degree at least two are independent:
for $k\ge2$ the polynomials
$y^k-xy^{k-1}$, $x^k-x^{k-1}y$, and $x^{k-1}y$ have respectively
the three individual order-$k$ branch values.
The Jacobian generators have images
$(-u^2,0,u^2)$ and $(0,u^2,-u^2)$.
For $a_i=u^2h_i$, the common-constant condition on both products forces
all three constant coefficients to vanish. The linear relations
give residues $(a,a,-a)$. Thus
\[
 \mathcal J^{-1}=B^3+\C\{z\}(u^{-1},u^{-1},-u^{-1}).
\]
For the tacnode and triple point, regularity of the products on each
normalization branch excludes poles of order greater than two before
these coefficient conditions are imposed. The three displayed modules are contained in the
invertible sheaf allowing a simple pole on the reduced inverse-image
curve. This proves the claimed bound at every point.

\end{enumerate}
\end{proof}







\section{Irreducible supports: topology, conductor and equality}
\hypertarget{tag-034F}{}\label{sec:II:night-sign}
The following proofs retain every Part~4 result of the night run. The final classification is Theorem~\ref{night:P4.58}; its exact normalization maps and its local exclusions are proved below.
\paragraph{Standing torus-fibration setting.} The ambient space is a smooth compact complex threefold $Y$ with $b_1(Y)=b_2(Y)=0$, and $f:Y\to B$ is proper and surjective with connected fibres and smooth general fibre a complex two-torus. Every statement below concerning an irreducible support uses its actual full fibre $mS$, normalization and resolution. The global assertions expressly retaining the hypothesis that every support is irreducible or has normal crossings are used with that hypothesis. Constructions and comparison examples with different Betti numbers state their own assumptions.
\subsection{The explicit threefold with identity monodromy}

Throughout this part we orient complex spaces by their complex structures.
All homology groups in the example have coefficients in $\Z$. A meridian is
counterclockwise in its indicated holomorphic coordinate; at infinity this
coordinate is $1/t$. We write
\[
 E=\C/(\Z+i\Z)=\C^*/q^\Z,\qquad q=e^{-2\pi},\qquad
 A=E^2,\qquad \Lambda=\Z e_1\oplus\Z e_2\oplus\Z e_3\oplus\Z e_4,
\]
where $(e_1,e_2,e_3,e_4)=((1,0),(i,0),(0,1),(0,i))$. The comparison
$E\longrightarrow\C^*/q^\Z$ is $z\mapsto e^{2\pi iz}$.

\begin{proposition}[An elliptic bundle degeneration]\hypertarget{tag-02WG}{}\label{night:P4.1}\label{stmt:P4.1}
\StatementRecord{night:P4.1}{proved}{02WG}{\SeriesRef{night:P4.3}}

Let $V=\operatorname{Bl}_{(0,0)}(B\times\Delta)$, where $B=E$, let $R$ be
the exceptional curve, and put $L=\OO_V(R)$. The holomorphic quotient
$U=L^\times/q^\Z$ is a smooth Hausdorff threefold, proper over $\Delta$.
Its noncentral fibres are $E\times B$. Its central fibre is the reduced
union of a primary Kodaira surface and the scalar Hopf surface
$(\C^2\setminus\{0\})/q^\Z$, meeting transversally in one elliptic curve.
\end{proposition}
\begin{proof}

\begin{enumerate}[label={$\langle1\rangle$\arabic*.},leftmargin=*]
\item We use coordinates $(z,t)$ at the blown-up point. The blowup has charts
\[
 (z,t)=(u,uv),\qquad (z,t)=(rs,s),\qquad r=v^{-1},\quad s=uv.
\]
The frames $u^{-1}$ and $s^{-1}$ for $L$ give fibre coordinates related by
$\lambda_B=v\lambda_A$. Each nonzero fibre coordinate, modulo multiplication
by $q$, gives a product chart with fibre $E$. The transitions are holomorphic
translations on $E$. The scalar action is free and properly discontinuous:
the logarithm of the norm of a nonzero fibre coordinate changes by the
nonzero constant $\log q$, and compact subsets meet only finitely many
translates. Thus the quotient is Hausdorff and smooth. These product charts
have compact fibre $E$, so $U\to V$ is proper. The blowup and $B\times\Delta
\to\Delta$ are proper, proving properness of their composite.

\item The canonical section of $\OO(R)$ has coefficients $u$ and $s$. Over
$t\ne0$ it never vanishes. The ratios $\lambda_A/u=\lambda_B/s$ therefore
give the holomorphic product trivialization over $\Delta^*$.

\item The central divisor in $V$ is $B'\cup R$, with $R^2=-1$ and $R\cdot B'=1$.
Thus $L|_R=\OO_{\Pp^1}(-1)$ and $L|_{B'}$ has degree one. Over $B'$ the
quotient is a principal elliptic bundle with primitive nonzero circle Chern
class. Its canonical bundle is trivial: the base differential wedged with
$d\lambda/\lambda$ is invariant under the line-bundle transition, because
the additional term is pulled back from the one-dimensional base. The
circle bundle over $B'$ has primitive Euler class one: in its
fundamental-group presentation the commutator of the two base generators
is the fibre generator. Abelianization therefore gives $H_1=\Z^2$.
The logarithm of the fibre norm modulo $\log q$ gives the further
circle factor, so the surface has $b_1=3$. The classification
\cite[Chapter~VI, Theorem~1.1 and Table~10, pp.~243--244]{BHPV2004E}
identifies this surface, whose canonical bundle is trivial, as a primary
Kodaira surface.

\item The tautological evaluation morphism
$\OO_{\Pp^1}(-1)^\times\to\C^2\setminus\{0\}$ is a biholomorphism: a
nonzero vector uniquely determines its line. It intertwines both scalar
actions, and hence identifies the other component with the stated Hopf
surface. Over $B'\cap R$ the components meet in $E$. In local product
coordinates the map is $(u,v,c)\mapsto uv$; this proves transversality,
multiplicity one, and the absence of other critical points.

\end{enumerate}
\end{proof}

\begin{proposition}[A compact torus fibration with three singular fibres]\hypertarget{tag-02WR}{}\label{night:P4.2}\label{stmt:P4.2}
\StatementRecord{night:P4.2}{proved}{02WR}{\SeriesRef{night:P4.1}}

Let $E=\C/(\Z+i\Z)$ and $A=E^2$. Glue a principal $A$-bundle
over $\Pp^1$ by
$a_{\rm out}=a_{\rm in}+(\ell(t),0)$, where
$\ell(t)=\log(t)/(2\pi i)\in E$. Over disjoint discs about
$0,1,\infty$, replace its product models by the degeneration of
Proposition~\ref{night:P4.1}, using, respectively, the markings
\[
 C_i\times B_i\longrightarrow A,\qquad
 (z,w)\longmapsto(w,z),\quad(iz,iz+w),\quad(z,iz+w),
\]
and the punctured-disc product trivialization from the canonical section
of $\OO(R)$. These transitions define a smooth compact complex
threefold $f:Y\to\Pp^1$ with connected fibres, exactly three singular
fibres of the form in Proposition~\ref{night:P4.1}, and identity
monodromy on $H_1$ of the regular fibre.
\end{proposition}
\begin{proof}

\begin{enumerate}[label={$\langle1\rangle$\arabic*.},leftmargin=*]
\item The map $\ell(t)=(\log t)/(2\pi i)$ with values in $E$ is single-valued and
holomorphic on an annulus: changing the logarithm adds an integer.
Consequently $a_{\rm out}=a_{\rm in}+(\ell(t),0)$ is a holomorphic transition
for a principal $A$-bundle on $\Pp^1$. Product charts with compact fibre
prove that this bundle is Hausdorff and compact.

\item We choose disjoint discs about $0,1,\infty$ and replace their product models
by Proposition~\ref{night:P4.1} using the holomorphic group isomorphisms
\[
 C_i\times B_i\longrightarrow A:\qquad
 (z,w)\mapsto(w,z),\quad(iz,iz+w),\quad(z,iz+w),
\]
respectively. Multiplication by $i$ preserves $\Z+i\Z$. The three inverse maps
are $(x,y)\mapsto(y,x)$, $(x,y)\mapsto(-ix,y-x)$ and
$(x,y)\mapsto(x,y-ix)$, respectively. Their Gaussian-integer matrices
have determinants $-1,i,1$ and hence their real lattice determinants
are all one. These maps, $\lambda_B=v\lambda_A$, the canonical
punctured-disc trivialization in Proposition~\ref{night:P4.1}, and
the clutch $\ell$ specify all transitions.
Different base points separate in the base. Points over the same base
point lie in one Hausdorff local model, which separates them. The gluing is
therefore Hausdorff. Properness holds over each base chart and hence over
the base; compactness follows.

\item The two central components have connected elliptic fibres over connected
curves and meet each other. All remaining fibres equal $A$. The equations
$uv=t$ give exactly the indicated critical values. The punctured-disc
trivializations identify all fibre lattices with $\Lambda$, and the remaining
transition is a translation on $A$. Translation induces the identity on
$H_1(A;\Z)$, giving the asserted monodromy representation.

\end{enumerate}
\end{proof}

\begin{proposition}[Integral homology of the local degeneration]\hypertarget{tag-02X2}{}\label{night:P4.3}\label{stmt:P4.3}
\StatementRecord{night:P4.3}{proved}{02X2}{the definitions and constructions stated in the proof}

For the degeneration $U$ of Proposition~\ref{night:P4.1}, order the
regular fibre lattice as $(a,c,b,d)$, where $(a,c)$ are the cycles
$(1,i)$ in $C_i$. The regular-fibre inclusions induce the quotient maps
\[
 H_1(A)=\Lambda\longrightarrow\Lambda/\Z a=H_1(U),\qquad
 H_2(A)=\bigwedge^2\Lambda\longrightarrow
 \bigwedge^2\Lambda/\Z(a\wedge c)=H_2(U).
\]
Both target groups are free. The meridian and every product of the
meridian with a fibre one-cycle map to zero in these degrees.
\end{proposition}
\begin{proof}

\begin{enumerate}[label={$\langle1\rangle$\arabic*.},leftmargin=*]
\item A Hermitian metric gives a diffeomorphism
$L^\times/q^\Z\cong S(L)\times S^1_c$, with factors the circle bundle and
the logarithm of the norm modulo $\log q$. We first specify a homotopy
model for $V$. Blowing up a four-ball replaces it by the disc bundle of
$\OO(-1)$ along its Hopf boundary. The complement of the ball in
$B\times\Delta$ has homotopy type $B\vee S^3$; the boundary sphere is
the $S^3$ summand, since its map to $B$ is nullhomotopic. The resulting
homotopy pushout is $B\vee S^2$. Equivalently, on simply connected
covers the replacement adds one sphere at each lift of the centre and
no other homology; the comparison is an isomorphism on $\pi_1$ and on
the homology of those covers, and hence a homotopy equivalence. We may
choose the $B$ summand to be a regular $B_t$ and the sphere to be $R$,
joining their base points by a path.

\item The circle Chern class is zero on $B_t$ and $-1$ on $R$. Therefore
$S(L)$ is represented by the homotopy pushout
\[
 T^3_{(a,b,d)}\longleftarrow S^1_a\longrightarrow S^3.
\]
The second map is a Hopf fibre and is nullhomotopic. Thus this pushout
has the homotopy type $(T^3\cup_aD^2)\vee S^3$. We multiply by $S^1_c$.
The resulting integral cellular complex starts with the exterior-cell
complex of $T^4$, ordered by $(a,c,b,d)$, and adds cells $r$ in degree
two, $rc,h$ in degree three, and $hc$ in degree four. Its only nonzero
differentials are
\[
 \partial r=a,\qquad \partial(rc)=a\wedge c.
\]
We choose the orientation of $r$ to give the first equality and use the
product orientation for $rc$. The pairs $(r,a)$ and $(rc,a\wedge c)$
are integral contractible summands: the degree-one homotopy sends
$a$ to $r$ and $a\wedge c$ to $rc$, and is zero on the remaining
basis cells. Deleting these pairs leaves the zero differential and
free homology of ranks $(1,3,5,5,2)$. The inclusion of the torus
subcomplex kills exactly $a$ in degree one and $a\wedge c$ in
degree two, giving the asserted quotient maps.

\item We represent a base point, and both base generators $b,d$, away from
$0\in B$. Their products with the disc miss $R$, so the canonical section
trivializes $L$ over the resulting chains. The meridian and the products
with $b,d$ consequently bound. The products with $a,c$ bound over one
such base point. This proves all the remaining assertions on the maps.

\end{enumerate}
\end{proof}

The lattice matrices, with columns $(a_i,c_i,b_i,d_i)$ in the common basis,
are
\[
S_1=\begin{pmatrix}0&0&1&0\\0&0&0&1\\1&0&0&0\\0&1&0&0\end{pmatrix},\quad
S_2=\begin{pmatrix}0&-1&0&0\\1&0&0&0\\0&-1&1&0\\1&0&0&1\end{pmatrix},\quad
S_3=\begin{pmatrix}1&0&0&0\\0&1&0&0\\0&-1&1&0\\1&0&0&1\end{pmatrix}.
\]
Put $Q_i=S_i^{-1}$. Deleting the first row of $Q_i$ gives $R_{1,i}$.
The row indices of $\bigwedge^2Q_i$ are $(ac,ab,ad,cb,cd,bd)$ and its
column indices are $(12,13,14,23,24,34)$; deleting its first row gives
$R_{2,i}$. Thus these are the integer matrices of the local quotient maps,
with target bases $(c_i,b_i,d_i)$ and $(a_ib_i,a_id_i,c_ib_i,c_id_i,b_id_i)$.

\begin{proposition}[P4.4]\hypertarget{tag-02XD}{}\label{night:P4.4}\label{stmt:P4.4}
\StatementRecord{night:P4.4}{proved}{02XD}{the definitions and constructions stated in the proof}

The threefold $Y$ has $H_1(Y;\Z)=H_2(Y;\Z)=0$ and $\pi_1(Y)=1$.
\end{proposition}
\begin{proof}

\begin{enumerate}[label={$\langle1\rangle$\arabic*.},leftmargin=*]
\item Let $P$ be the complement of three small open base discs. The remaining
bundle $W=f^{-1}(P)$ is trivial. We use $W=A\times P$, let $s_1,s_2$
generate $H_1(P)$, and set $s_3=-s_1-s_2$. The overlaps are
$N_i=A\times S^1$. Positive continuation of $\ell$ around zero adds $e_1$.
With
\[
 g_1=e_1,\quad g_2=g_3=0,\qquad
 v_1=\binom10,\quad v_2=\binom01,\quad v_3=\binom{-1}{-1},
\]
the overlap maps into $W$ are
\[
 x\mapsto x,\quad s_i\mapsto g_i+v_i,\qquad
 x\wedge y\mapsto x\wedge y,\quad
 x\wedge s_i\mapsto x\wedge g_i+x\otimes v_i.
\]
These are the maps represented by the products of the cycles, or by the
induced maps into the aspherical space $A\times P$; in particular they
retain the degree-two contribution of the clutch.

\item We use the ordered basis $(e_1,e_2,e_3,e_4,s_1,s_2)$ for $H_1(W)$ and
append the three bases $(c_i,b_i,d_i)$ for the local groups. On the source
we concatenate the three bases $(e_1,e_2,e_3,e_4,s_i)$. The resulting
$15\times15$ Mayer--Vietoris matrix $F_1$ has, in its $i$th five-column
block, upper block
\[
 \begin{pmatrix}I_4&g_i\\0_{2\times4}&v_i\end{pmatrix}
\]
and local block $(-R_{1,i},0)$ in the $i$th summand; its other local
blocks vanish. The sign is the usual difference of the two inclusions.

\item For $H_2(W)$ we use
\[
 (12,13,14,23,24,34,e_1s_1,\ldots,e_4s_1,e_1s_2,\ldots,e_4s_2)
\]
and append the three local five-element bases. The source is the
concatenation of $(12,13,14,23,24,34,e_1s_i,\ldots,e_4s_i)$.
The $29\times30$ matrix $F_2$ has top six rows
$[I_6,(e_j\wedge g_i)_{j=1}^4]$ in its $i$th block; its next eight rows
have the coefficients of $e_j\otimes v_i$ in the displayed order, and
its $i$th local block is $[-R_{2,i},0]$. All other blocks vanish.
Finally
\[
 F_0=\begin{pmatrix}1&1&1\\-1&0&0\\0&-1&0\\0&0&-1\end{pmatrix}
\]
has zero kernel. Thus the required exact sequence is
\[
 \Z^{30}\xrightarrow{F_2}\Z^{29}\longrightarrow H_2(Y)
 \longrightarrow\Z^{15}\xrightarrow{F_1}\Z^{15}\longrightarrow H_1(Y)
 \longrightarrow\Z^3\xrightarrow{F_0}\Z^4.
\]

\item We verify integral surjectivity by hand. Eliminating the local $H_1$
groups and two independent base meridians leaves $\Lambda$ modulo
$e_1,a_1,a_2,a_3$. The relation matrix is
\[
 R_{H_1}=\begin{pmatrix}1&0&0&1\\0&0&1&0\\0&1&0&0\\0&0&1&1\end{pmatrix},
 \qquad\det R_{H_1}=-1.
\]
Hence $F_1$ is an integral isomorphism. In degree two the fibre
$\bigwedge^2\Lambda$ columns eliminate the fifteen local generators,
leaving the three relations $\eta_i=a_i\wedge c_i$. The second meridian
columns kill $\Lambda\otimes s_2$, the third then kill
$\Lambda\otimes s_1$, and the first leave $e_1\wedge\Lambda$. The six
relations, in the order $e_{12},e_{13},e_{14},\eta_1,\eta_2,\eta_3$, have
matrix
\[
R_{H_2}=\begin{pmatrix}
1&0&0&0&1&1\\0&1&0&0&0&-1\\0&0&1&0&1&0\\
0&0&0&0&-1&0\\0&0&0&0&0&-1\\0&0&0&1&1&1
\end{pmatrix},\qquad \det R_{H_2}=1.
\]
This proves surjectivity of $F_2$ over $\Z$ and excludes torsion. The
executed independent Smith calculation gives fifteen unit factors for
$F_1$ and twenty-nine unit factors for $F_2$, and computes $\det F_1=-1$.

\item For the fundamental group the local pushout gives
$\pi_1(U_i)=\Lambda/\Z a_i$, with trivial meridian. Van Kampen starts
from $\pi_1(W)=\Lambda\times F(s_1,s_2)$ and imposes
\[
 a_1=a_2=a_3=0,\qquad s_1+e_1=0,\quad s_2=0,\quad s_1+s_2=0.
\]
Both base generators and $e_1$ vanish. The same unimodular matrix
$R_{H_1}$ then kills the fibre lattice, proving $\pi_1(Y)=1$.

\end{enumerate}
\end{proof}

The program \texttt{code/part4/example\_mayer\_vietoris.py} recomputes all
entries, differentials, determinants and Smith forms exactly. Its output
and machine-readable matrices are supplied alongside it. For completeness
we print every integer matrix, including the zero differentials:
% Generated by code/part4/matrices_to_tex.py from example_matrices.json.
\begingroup
\setlength{\arraycolsep}{1.7pt}
\renewcommand{\arraystretch}{0.88}
\par\smallskip\noindent
$S_1$ has size $4\times 4$; its full integer matrix is
\[
\resizebox{!}{11.2mm}{$\displaystyle S_1=\begin{pmatrix}
0 & 0 & 1 & 0\\
0 & 0 & 0 & 1\\
1 & 0 & 0 & 0\\
0 & 1 & 0 & 0\\
\end{pmatrix}$}
\]
\par\smallskip\noindent
$S_2$ has size $4\times 4$; its full integer matrix is
\[
\resizebox{!}{11.2mm}{$\displaystyle S_2=\begin{pmatrix}
0 & -1 & 0 & 0\\
1 & 0 & 0 & 0\\
0 & -1 & 1 & 0\\
1 & 0 & 0 & 1\\
\end{pmatrix}$}
\]
\par\smallskip\noindent
$S_3$ has size $4\times 4$; its full integer matrix is
\[
\resizebox{!}{11.2mm}{$\displaystyle S_3=\begin{pmatrix}
1 & 0 & 0 & 0\\
0 & 1 & 0 & 0\\
0 & -1 & 1 & 0\\
1 & 0 & 0 & 1\\
\end{pmatrix}$}
\]
\par\smallskip\noindent
$R_{1,1}$ has size $3\times 4$; its full integer matrix is
\[
\resizebox{!}{8.399999999999999mm}{$\displaystyle R_{1,1}=\begin{pmatrix}
0 & 0 & 0 & 1\\
1 & 0 & 0 & 0\\
0 & 1 & 0 & 0\\
\end{pmatrix}$}
\]
\par\smallskip\noindent
$R_{1,2}$ has size $3\times 4$; its full integer matrix is
\[
\resizebox{!}{8.399999999999999mm}{$\displaystyle R_{1,2}=\begin{pmatrix}
-1 & 0 & 0 & 0\\
-1 & 0 & 1 & 0\\
0 & -1 & 0 & 1\\
\end{pmatrix}$}
\]
\par\smallskip\noindent
$R_{1,3}$ has size $3\times 4$; its full integer matrix is
\[
\resizebox{!}{8.399999999999999mm}{$\displaystyle R_{1,3}=\begin{pmatrix}
0 & 1 & 0 & 0\\
0 & 1 & 1 & 0\\
-1 & 0 & 0 & 1\\
\end{pmatrix}$}
\]
\par\smallskip\noindent
$R_{2,1}$ has size $5\times 6$; its full integer matrix is
\[
\resizebox{!}{14.0mm}{$\displaystyle R_{2,1}=\begin{pmatrix}
0 & -1 & 0 & 0 & 0 & 0\\
0 & 0 & 0 & -1 & 0 & 0\\
0 & 0 & -1 & 0 & 0 & 0\\
0 & 0 & 0 & 0 & -1 & 0\\
1 & 0 & 0 & 0 & 0 & 0\\
\end{pmatrix}$}
\]
\par\smallskip\noindent
$R_{2,2}$ has size $5\times 6$; its full integer matrix is
\[
\resizebox{!}{14.0mm}{$\displaystyle R_{2,2}=\begin{pmatrix}
1 & 0 & 0 & 1 & 0 & 0\\
0 & 0 & 0 & 0 & 1 & 0\\
0 & -1 & 0 & 0 & 0 & 0\\
1 & 0 & -1 & 0 & 0 & 0\\
1 & 0 & -1 & 1 & 0 & 1\\
\end{pmatrix}$}
\]
\par\smallskip\noindent
$R_{2,3}$ has size $5\times 6$; its full integer matrix is
\[
\resizebox{!}{14.0mm}{$\displaystyle R_{2,3}=\begin{pmatrix}
1 & 1 & 0 & 0 & 0 & 0\\
0 & 0 & 1 & 0 & 0 & 0\\
0 & 0 & 0 & 1 & 0 & 0\\
1 & 0 & 0 & 0 & 1 & 0\\
1 & 1 & 0 & 0 & 1 & 1\\
\end{pmatrix}$}
\]
\par\smallskip\noindent
$\partial_1$ has size $1\times 4$; its full integer matrix is
\[
\resizebox{!}{2.8mm}{$\displaystyle \partial_1=\begin{pmatrix}
0 & 0 & 0 & 0\\
\end{pmatrix}$}
\]
\par\smallskip\noindent
$\partial_2$ has size $4\times 7$; its full integer matrix is
\[
\resizebox{!}{11.2mm}{$\displaystyle \partial_2=\begin{pmatrix}
0 & 0 & 0 & 0 & 0 & 0 & 1\\
0 & 0 & 0 & 0 & 0 & 0 & 0\\
0 & 0 & 0 & 0 & 0 & 0 & 0\\
0 & 0 & 0 & 0 & 0 & 0 & 0\\
\end{pmatrix}$}
\]
\par\smallskip\noindent
$\partial_3$ has size $7\times 6$; its full integer matrix is
\[
\resizebox{!}{19.599999999999998mm}{$\displaystyle \partial_3=\begin{pmatrix}
0 & 0 & 0 & 0 & 1 & 0\\
0 & 0 & 0 & 0 & 0 & 0\\
0 & 0 & 0 & 0 & 0 & 0\\
0 & 0 & 0 & 0 & 0 & 0\\
0 & 0 & 0 & 0 & 0 & 0\\
0 & 0 & 0 & 0 & 0 & 0\\
0 & 0 & 0 & 0 & 0 & 0\\
\end{pmatrix}$}
\]
\par\smallskip\noindent
$\partial_4$ has size $6\times 2$; its full integer matrix is
\[
\resizebox{!}{16.799999999999997mm}{$\displaystyle \partial_4=\begin{pmatrix}
0 & 0\\
0 & 0\\
0 & 0\\
0 & 0\\
0 & 0\\
0 & 0\\
\end{pmatrix}$}
\]
\par\smallskip\noindent
$F_0$ has size $4\times 3$; its full integer matrix is
\[
\resizebox{!}{11.2mm}{$\displaystyle F_0=\begin{pmatrix}
1 & 1 & 1\\
-1 & 0 & 0\\
0 & -1 & 0\\
0 & 0 & -1\\
\end{pmatrix}$}
\]
\par\smallskip\noindent
$F_1$ has size $15\times 15$; its full integer matrix is
\[
\resizebox{0.94\linewidth}{!}{$\displaystyle F_1=\begin{pmatrix}
1 & 0 & 0 & 0 & 1 & 1 & 0 & 0 & 0 & 0 & 1 & 0 & 0 & 0 & 0\\
0 & 1 & 0 & 0 & 0 & 0 & 1 & 0 & 0 & 0 & 0 & 1 & 0 & 0 & 0\\
0 & 0 & 1 & 0 & 0 & 0 & 0 & 1 & 0 & 0 & 0 & 0 & 1 & 0 & 0\\
0 & 0 & 0 & 1 & 0 & 0 & 0 & 0 & 1 & 0 & 0 & 0 & 0 & 1 & 0\\
0 & 0 & 0 & 0 & 1 & 0 & 0 & 0 & 0 & 0 & 0 & 0 & 0 & 0 & -1\\
0 & 0 & 0 & 0 & 0 & 0 & 0 & 0 & 0 & 1 & 0 & 0 & 0 & 0 & -1\\
0 & 0 & 0 & -1 & 0 & 0 & 0 & 0 & 0 & 0 & 0 & 0 & 0 & 0 & 0\\
-1 & 0 & 0 & 0 & 0 & 0 & 0 & 0 & 0 & 0 & 0 & 0 & 0 & 0 & 0\\
0 & -1 & 0 & 0 & 0 & 0 & 0 & 0 & 0 & 0 & 0 & 0 & 0 & 0 & 0\\
0 & 0 & 0 & 0 & 0 & 1 & 0 & 0 & 0 & 0 & 0 & 0 & 0 & 0 & 0\\
0 & 0 & 0 & 0 & 0 & 1 & 0 & -1 & 0 & 0 & 0 & 0 & 0 & 0 & 0\\
0 & 0 & 0 & 0 & 0 & 0 & 1 & 0 & -1 & 0 & 0 & 0 & 0 & 0 & 0\\
0 & 0 & 0 & 0 & 0 & 0 & 0 & 0 & 0 & 0 & 0 & -1 & 0 & 0 & 0\\
0 & 0 & 0 & 0 & 0 & 0 & 0 & 0 & 0 & 0 & 0 & -1 & -1 & 0 & 0\\
0 & 0 & 0 & 0 & 0 & 0 & 0 & 0 & 0 & 0 & 1 & 0 & 0 & -1 & 0\\
\end{pmatrix}$}
\]
\par\smallskip\noindent
$F_2$ has size $29\times 30$; its full integer matrix is
\[
\resizebox{0.94\linewidth}{!}{$\displaystyle F_2=\begin{pmatrix}
1 & 0 & 0 & 0 & 0 & 0 & 0 & -1 & 0 & 0 & 1 & 0 & 0 & 0 & 0 & 0 & 0 & 0 & 0 & 0 & 1 & 0 & 0 & 0 & 0 & 0 & 0 & 0 & 0 & 0\\
0 & 1 & 0 & 0 & 0 & 0 & 0 & 0 & -1 & 0 & 0 & 1 & 0 & 0 & 0 & 0 & 0 & 0 & 0 & 0 & 0 & 1 & 0 & 0 & 0 & 0 & 0 & 0 & 0 & 0\\
0 & 0 & 1 & 0 & 0 & 0 & 0 & 0 & 0 & -1 & 0 & 0 & 1 & 0 & 0 & 0 & 0 & 0 & 0 & 0 & 0 & 0 & 1 & 0 & 0 & 0 & 0 & 0 & 0 & 0\\
0 & 0 & 0 & 1 & 0 & 0 & 0 & 0 & 0 & 0 & 0 & 0 & 0 & 1 & 0 & 0 & 0 & 0 & 0 & 0 & 0 & 0 & 0 & 1 & 0 & 0 & 0 & 0 & 0 & 0\\
0 & 0 & 0 & 0 & 1 & 0 & 0 & 0 & 0 & 0 & 0 & 0 & 0 & 0 & 1 & 0 & 0 & 0 & 0 & 0 & 0 & 0 & 0 & 0 & 1 & 0 & 0 & 0 & 0 & 0\\
0 & 0 & 0 & 0 & 0 & 1 & 0 & 0 & 0 & 0 & 0 & 0 & 0 & 0 & 0 & 1 & 0 & 0 & 0 & 0 & 0 & 0 & 0 & 0 & 0 & 1 & 0 & 0 & 0 & 0\\
0 & 0 & 0 & 0 & 0 & 0 & 1 & 0 & 0 & 0 & 0 & 0 & 0 & 0 & 0 & 0 & 0 & 0 & 0 & 0 & 0 & 0 & 0 & 0 & 0 & 0 & -1 & 0 & 0 & 0\\
0 & 0 & 0 & 0 & 0 & 0 & 0 & 1 & 0 & 0 & 0 & 0 & 0 & 0 & 0 & 0 & 0 & 0 & 0 & 0 & 0 & 0 & 0 & 0 & 0 & 0 & 0 & -1 & 0 & 0\\
0 & 0 & 0 & 0 & 0 & 0 & 0 & 0 & 1 & 0 & 0 & 0 & 0 & 0 & 0 & 0 & 0 & 0 & 0 & 0 & 0 & 0 & 0 & 0 & 0 & 0 & 0 & 0 & -1 & 0\\
0 & 0 & 0 & 0 & 0 & 0 & 0 & 0 & 0 & 1 & 0 & 0 & 0 & 0 & 0 & 0 & 0 & 0 & 0 & 0 & 0 & 0 & 0 & 0 & 0 & 0 & 0 & 0 & 0 & -1\\
0 & 0 & 0 & 0 & 0 & 0 & 0 & 0 & 0 & 0 & 0 & 0 & 0 & 0 & 0 & 0 & 1 & 0 & 0 & 0 & 0 & 0 & 0 & 0 & 0 & 0 & -1 & 0 & 0 & 0\\
0 & 0 & 0 & 0 & 0 & 0 & 0 & 0 & 0 & 0 & 0 & 0 & 0 & 0 & 0 & 0 & 0 & 1 & 0 & 0 & 0 & 0 & 0 & 0 & 0 & 0 & 0 & -1 & 0 & 0\\
0 & 0 & 0 & 0 & 0 & 0 & 0 & 0 & 0 & 0 & 0 & 0 & 0 & 0 & 0 & 0 & 0 & 0 & 1 & 0 & 0 & 0 & 0 & 0 & 0 & 0 & 0 & 0 & -1 & 0\\
0 & 0 & 0 & 0 & 0 & 0 & 0 & 0 & 0 & 0 & 0 & 0 & 0 & 0 & 0 & 0 & 0 & 0 & 0 & 1 & 0 & 0 & 0 & 0 & 0 & 0 & 0 & 0 & 0 & -1\\
0 & 1 & 0 & 0 & 0 & 0 & 0 & 0 & 0 & 0 & 0 & 0 & 0 & 0 & 0 & 0 & 0 & 0 & 0 & 0 & 0 & 0 & 0 & 0 & 0 & 0 & 0 & 0 & 0 & 0\\
0 & 0 & 0 & 1 & 0 & 0 & 0 & 0 & 0 & 0 & 0 & 0 & 0 & 0 & 0 & 0 & 0 & 0 & 0 & 0 & 0 & 0 & 0 & 0 & 0 & 0 & 0 & 0 & 0 & 0\\
0 & 0 & 1 & 0 & 0 & 0 & 0 & 0 & 0 & 0 & 0 & 0 & 0 & 0 & 0 & 0 & 0 & 0 & 0 & 0 & 0 & 0 & 0 & 0 & 0 & 0 & 0 & 0 & 0 & 0\\
0 & 0 & 0 & 0 & 1 & 0 & 0 & 0 & 0 & 0 & 0 & 0 & 0 & 0 & 0 & 0 & 0 & 0 & 0 & 0 & 0 & 0 & 0 & 0 & 0 & 0 & 0 & 0 & 0 & 0\\
-1 & 0 & 0 & 0 & 0 & 0 & 0 & 0 & 0 & 0 & 0 & 0 & 0 & 0 & 0 & 0 & 0 & 0 & 0 & 0 & 0 & 0 & 0 & 0 & 0 & 0 & 0 & 0 & 0 & 0\\
0 & 0 & 0 & 0 & 0 & 0 & 0 & 0 & 0 & 0 & -1 & 0 & 0 & -1 & 0 & 0 & 0 & 0 & 0 & 0 & 0 & 0 & 0 & 0 & 0 & 0 & 0 & 0 & 0 & 0\\
0 & 0 & 0 & 0 & 0 & 0 & 0 & 0 & 0 & 0 & 0 & 0 & 0 & 0 & -1 & 0 & 0 & 0 & 0 & 0 & 0 & 0 & 0 & 0 & 0 & 0 & 0 & 0 & 0 & 0\\
0 & 0 & 0 & 0 & 0 & 0 & 0 & 0 & 0 & 0 & 0 & 1 & 0 & 0 & 0 & 0 & 0 & 0 & 0 & 0 & 0 & 0 & 0 & 0 & 0 & 0 & 0 & 0 & 0 & 0\\
0 & 0 & 0 & 0 & 0 & 0 & 0 & 0 & 0 & 0 & -1 & 0 & 1 & 0 & 0 & 0 & 0 & 0 & 0 & 0 & 0 & 0 & 0 & 0 & 0 & 0 & 0 & 0 & 0 & 0\\
0 & 0 & 0 & 0 & 0 & 0 & 0 & 0 & 0 & 0 & -1 & 0 & 1 & -1 & 0 & -1 & 0 & 0 & 0 & 0 & 0 & 0 & 0 & 0 & 0 & 0 & 0 & 0 & 0 & 0\\
0 & 0 & 0 & 0 & 0 & 0 & 0 & 0 & 0 & 0 & 0 & 0 & 0 & 0 & 0 & 0 & 0 & 0 & 0 & 0 & -1 & -1 & 0 & 0 & 0 & 0 & 0 & 0 & 0 & 0\\
0 & 0 & 0 & 0 & 0 & 0 & 0 & 0 & 0 & 0 & 0 & 0 & 0 & 0 & 0 & 0 & 0 & 0 & 0 & 0 & 0 & 0 & -1 & 0 & 0 & 0 & 0 & 0 & 0 & 0\\
0 & 0 & 0 & 0 & 0 & 0 & 0 & 0 & 0 & 0 & 0 & 0 & 0 & 0 & 0 & 0 & 0 & 0 & 0 & 0 & 0 & 0 & 0 & -1 & 0 & 0 & 0 & 0 & 0 & 0\\
0 & 0 & 0 & 0 & 0 & 0 & 0 & 0 & 0 & 0 & 0 & 0 & 0 & 0 & 0 & 0 & 0 & 0 & 0 & 0 & -1 & 0 & 0 & 0 & -1 & 0 & 0 & 0 & 0 & 0\\
0 & 0 & 0 & 0 & 0 & 0 & 0 & 0 & 0 & 0 & 0 & 0 & 0 & 0 & 0 & 0 & 0 & 0 & 0 & 0 & -1 & -1 & 0 & 0 & -1 & -1 & 0 & 0 & 0 & 0\\
\end{pmatrix}$}
\]
\par\smallskip\noindent
$R_{H_1}$ has size $4\times 4$; its full integer matrix is
\[
\resizebox{!}{11.2mm}{$\displaystyle R_{H_1}=\begin{pmatrix}
1 & 0 & 0 & 1\\
0 & 0 & 1 & 0\\
0 & 1 & 0 & 0\\
0 & 0 & 1 & 1\\
\end{pmatrix}$}
\]
\par\smallskip\noindent
$R_{H_2}$ has size $6\times 6$; its full integer matrix is
\[
\resizebox{!}{16.799999999999997mm}{$\displaystyle R_{H_2}=\begin{pmatrix}
1 & 0 & 0 & 0 & 1 & 1\\
0 & 1 & 0 & 0 & 0 & -1\\
0 & 0 & 1 & 0 & 1 & 0\\
0 & 0 & 0 & 0 & -1 & 0\\
0 & 0 & 0 & 0 & 0 & -1\\
0 & 0 & 0 & 1 & 1 & 1\\
\end{pmatrix}$}
\]
\endgroup


\begin{proposition}[Invariants of the compact torus fibration]\hypertarget{tag-02XO}{}\label{night:P4.5}\label{stmt:P4.5}
\StatementRecord{night:P4.5}{proved}{02XO}{\SeriesRef{night:P4.4}}

For the threefold $Y$ of Proposition~\ref{night:P4.2},
\[
 a(Y)=1,\qquad c_3[Y]=0,\qquad b_3(Y)=2,
\]
and the integral homology is free of ranks $(1,0,0,2,0,0,1)$.
\end{proposition}
\begin{proof}

\begin{enumerate}[label={$\langle1\rangle$\arabic*.},leftmargin=*]
\item The regular fibre has Euler number zero. Both components of each special
fibre are elliptic bundles over compact curves, and their intersection is
elliptic, so additivity gives zero for each special support. Stratifying
the base and applying multiplicativity to the smooth bundle gives
$e(Y)=0$. The Euler class of the real bundle underlying a complex
rank-three tangent bundle is its top Chern class, with the complex
orientation. Hence $c_3[Y]=0$. Duality and Proposition~\ref{night:P4.4} give $b_4=b_5=0$ and
$e(Y)=2-b_3$, so $b_3=2$. Integral duality and the universal-coefficient
theorem identify the torsion of $H_3$ with that of $H_2$, which is zero.
All other groups have already been determined.

\item We also have $H^2(Y;\Z)=0$. A horizontal effective divisor would restrict
to a nonzero effective divisor on the Kähler torus $A$ with zero first
Chern class. Its pairing with a Kähler class would be positive by
integration over the divisor and zero by the Chern-class computation.
Thus every divisor on $Y$ is vertical. The polar divisor of a meromorphic
function is consequently vertical; on a general fibre the function is
holomorphic and therefore constant. Local holomorphic sections give a
descended meromorphic function on the regular part of the base. Near
each exceptional base value a sufficiently high power of the base
coordinate removes the vertical poles. Connected fibres then give
holomorphic descent there, and dividing by that power gives meromorphic
descent. Therefore $\C(Y)=f^*\C(\Pp^1)$ and $a(Y)=1$.

\end{enumerate}
\end{proof}

\subsection{The normal-crossings sign theorem}

\begin{proposition}[Euler characteristic of a normal crossing divisor]\hypertarget{tag-02XY}{}\label{night:P4.6}\label{stmt:P4.6}
\StatementRecord{night:P4.6}{proved}{02XY}{the definitions and constructions stated in the proof}

Let $X$ be a compact complex threefold with $b_2(X)=0$, and let $F$ be a
reduced compact normal-crossings divisor. We allow different local
branches of one global irreducible component to meet. If $T(F)$ is the
number of ordinary triple points, then
\[
 e(F)=T(F),\qquad T(F)\text{ is even}.
\]
For a fibre $\sum_j m_jF_j$ with this reduced support, the divisor
in the formula is $F=\sum_jF_j$.
\end{proposition}
\begin{proof}

\begin{enumerate}[label={$\langle1\rangle$\arabic*.},leftmargin=*]
\item Write $\nu:S=\coprod S_i\to X$ for the normalization immersion, and $D_i$
for the conductor divisor on $S_i$. The surfaces $S_i$ are smooth and
the curves $D_i$ are reduced nodal divisors. The quotient
$N_i=\nu_i^*T_X/T_{S_i}$ is a line bundle. A local equation of $F$ gives
\[
 N_i\simeq \nu_i^*\OO_X(F)\otimes\OO_{S_i}(-D_i).
\]
Indeed, for the coordinate equation $x_1\cdots x_k=0$ the normal
derivative on $x_1=0$ is $x_2\cdots x_k$, whose zero divisor is $D_i$.
These are morphisms of holomorphic vector bundles, including when
$S_i$ is not embedded globally.

\item Since $b_2(X)=0$, both $c_1(T_X)$ and $c_1(\OO_X(F))$ vanish in real
cohomology. Duality gives $b_4(X)=0$, so $c_2(T_X)$ also vanishes. The
tangent exact sequence yields
\[
 K_{S_i}\equiv-D_i,\qquad c_2(S_i)=K_{S_i}^2.
\]
If $a=\sum_i\chi(\OO_{S_i})$, the surface Riemann--Roch formula
$12\chi(\OO_{S_i})=K_{S_i}^2+e(S_i)$ consequently gives
$\sum_i e(S_i)=6a$.

\item Let $\Gamma$ be the normalization of the double-curve locus, with one
sheet for each unordered pair of local branches. The normalization of
$\coprod D_i$ is an unramified degree-two cover of $\Gamma$, recording
the choice of branch. Each triple point contributes three nodes to
$\coprod D_i$. Adjunction gives
$\chi(\OO_{D_i})=-D_i\cdot(D_i+K_{S_i})/2=0$. The normalization
sequence for these nodal curves therefore gives
\[
 0=2\chi(\OO_\Gamma)-3T(F)=e(\Gamma)-3T(F).
\]
The normalization complex of $F$ is the exact alternating restriction
complex of the local coordinate hyperplanes. Permuting the branches
globally replaces its double-curve term by $\OO_\Gamma$ tensored with
the sign line of the two-sheeted branch cover. That line has order at
most two, and hence degree zero on every component; it has the same
coherent Euler characteristic as $\OO_\Gamma$. The triple term has
length one at each triple point. It follows that
\[
 \chi(\OO_F)=a-\chi(\OO_\Gamma)+T(F)=a-T(F)/2.
\]
The exact sequence for the Cartier divisor $F$ and analytic
Hirzebruch--Riemann--Roch give $\chi(\OO_F)=0$, since
$\operatorname{ch}(\OO_X(-F))=\operatorname{ch}(\OO_X)$ in rational
cohomology. Thus $T(F)=2a$. Finally ordinary Euler additivity gives
\[
 e(F)=\sum_i e(S_i)-e(\Gamma)+T(F)=6a-2T(F)=T(F).
\]
At a point with $r=1,2,3$ local branches, the alternating
contribution is $r-\binom r2+\binom r3=1$. This gives the
displayed Euler additivity formula and proves the identity and parity.

\end{enumerate}
\end{proof}

\begin{proposition}[The Chern number of a normal crossing fibration]\hypertarget{tag-02XZ}{}\label{night:P4.7}\label{stmt:P4.7}
\StatementRecord{night:P4.7}{proved}{02XZ}{\SeriesRef{night:P4.6}}

Let $f:X\to B$ be a proper surjective holomorphic map with connected
fibres from a compact complex threefold satisfying $b_1(X)=b_2(X)=0$
to a smooth compact curve. Suppose all reduced fibre supports have
normal crossings. Then
\[
 B=\Pp^1,\qquad
 c_3[X]=\sum_{b\in B}T((f^{-1}(b))_{\mathrm{red}})\in\{0,2\},\qquad
 b_3(X)=2-c_3[X].
\]
\end{proposition}
\begin{proof}
\begin{enumerate}[label={$\langle1\rangle$\arabic*.},leftmargin=*]
\item Connected fibres give $f_*\R=\R$. The beginning of the constant-sheaf
Leray sequence injects $H^1(B;\R)$ into $H^1(X;\R)$, so $B$ has genus
zero. A smooth fibre $G$ has trivial normal line bundle. Its tangent
sequence gives $c_2(G)=c_2(X)|_G=0$, hence $e(G)=0$ for any such fibre.
The critical values are finite. Ehresmann's theorem gives a smooth
bundle over their complement. Euler additivity and Proposition~\ref{night:P4.6} give the
displayed sum. Duality gives $e(X)=2-b_3(X)$, and the nondegenerate
alternating cup product on $H^3(X;\R)$ makes $b_3$ even. The
nonnegative sum is therefore zero or two.

\end{enumerate}
\end{proof}

\begin{proposition}[Normal crossing torus fibres and their normalizations]\hypertarget{tag-02Y0}{}\label{night:P4.8}\label{stmt:P4.8}
\StatementRecord{night:P4.8}{proved}{02Y0}{\SeriesRef{night:P4.12}, \SeriesRef{night:P4.6}, \SeriesRef{night:P4.7}}

For a fibration in Proposition~\ref{night:P4.7} with general fibre
a complex two-torus, an irreducible singular support with
rational normalization has Euler number two; every other irreducible
support has Euler number zero. The supports of Euler number zero have the following forms: a smooth support has a finite unramified
torus cover; a singular irreducible support has smooth minimal
elliptic-ruled or Hopf normalization; a reducible support has smooth
minimal Hopf or non-Kähler elliptic normalization components, with
smooth elliptic double curves and finite unramified branch maps.
Thus $c_3[X]=0$ precisely when every fibre belongs to these zero cases.
\end{proposition}
\begin{proof}

\begin{enumerate}[label={$\langle1\rangle$\arabic*.},leftmargin=*]
\item We apply \cite[Chapter~VI, Theorem~1.1 and Table~10, pp.~243--244]{BHPV2004E} to the smooth compact
normalization components, using the identity $e(S_i)=K_{S_i}^2$ proved
in Proposition~\ref{night:P4.6}. The relevant non-Kähler classification is stated in
\cite[arXiv version~3, Theorem~4.2, pp.~17--18; Theorem~5.1 and proof, pp.~27--29]{OVV2020}: a minimal non-Kähler
surface is elliptic or of class VII, and the reduced fibres of a
minimal non-Kähler elliptic surface are smooth elliptic curves.

\item For a component with nonempty conductor, a Kähler form $\omega$ would
give $K_{S_i}\cdot\omega=-D_i\cdot\omega<0$. No positive power of
the canonical bundle can then have a nonzero section. The classification
makes a Kähler component rational or ruled. If it is obtained by $k$
blowups of a ruled surface over a genus-$g$ curve, then
\[
 e=4(1-g)+k,\qquad K^2=8(1-g)-k.
\]
Equality gives $k=2(1-g)$, so $g\le1$ and the genus-one case is minimal.
A rational component has $\chi(\OO)=1$, hence $e=6$ here, including
the case with minimal model $\Pp^2$.

\item A minimal non-Kähler elliptic surface has $e=K^2=0$, by its smooth
elliptic reduced fibres. A minimal class-VII surface has $b_1=1$ and
$p_g=0$, so $\chi(\OO)=0$ and $K^2=-e$, with $e=b_2\ge0$.
Blowups preserve $\chi(\OO)$ and change $(e,K^2)$ by $(1,-1)$.
Consequently $e=K^2$ forces minimality and $e=0$ in all these
non-Kähler cases. The class-VII case has $b_2=0$. The finite-cover classification
\cite[arXiv version~3, Theorem~1, p.~1]{BBK2019} gives a finite
unramified map $p:\widehat S_i\to S_i$ with $\widehat S_i$ Hopf
or Inoue. The inverse image $p^{-1}(D_i)$ is a nonempty compact
curve: surjectivity makes it nonempty, and an unramified map preserves
its local codimension. Inoue surfaces have no compact curves
\cite[Theorem~3.1, p.~303]{KhaledTeleman2026Journal}; the type-III
family is covered by the explicit unramified type-II double cover
of \cite[proof of Remark~3.5, p.~305]{KhaledTeleman2026Journal}.
Thus $\widehat S_i$ is Hopf. Its universal covering
$\C^2\setminus\{0\}\to\widehat S_i$, composed with $p$, is a
universal covering of $S_i$, which is therefore a Hopf surface.

\item For an irreducible singular support there is one normalization
component. Proposition~\ref{night:P4.6} gives $T=2\chi(\OO_{S_1})$, which is two in the
rational case and zero in the other cases. Proposition~\ref{night:P4.12} proves that the
dualizing sheaf of an irreducible support is torsion. Adjunction on
the normalization therefore makes $K_{S_1}+D_1$ torsion. Since
$D_1\ne0$, every positive plurigenus of $S_1$ vanishes: multiplying
a pluricanonical section by an appropriate positive power of the
conductor section would give a section of a torsion line bundle with
a zero, whose sufficiently high power is a nonzero constant. Thus a
non-Kähler normalization in this irreducible case is of class VII,
and hence is the Hopf case already obtained.

\item For a reducible connected fibre, each normalization component meets
another global irreducible component. Pulling back the divisor of
that component gives a nonzero effective divisor on $S_i$ with zero
real first Chern class. A Kähler form would pair positively with it,
contradicting this zero class. Hence every component is non-Kähler,
has $\chi(\OO)=0$, and $T=0$. The conductor components are disjoint
and smooth, because all their nodes would come from triple points.
For any conductor component $C$, adjunction and $K+D\equiv0$ give
$2g(C)-2=(K+C)\cdot C=0$, since the other conductor components
are disjoint. Thus the double curves are smooth elliptic, and their
branch maps are the unramified maps already described in Proposition~\ref{night:P4.6}.

\item Finally let the support $G$ of a multiple fibre $mG$ be smooth. The
normalized cyclic base change $t=u^m$ is locally a disjoint union of
smooth charts: the original equation is $t=z^m$ times a holomorphic
unit, whose root can be extracted locally. These charts show that
the central fibre of the base-changed family maps finitely and
unramifiedly to $G$, and that its map to the $u$-disc is a proper submersion. Its nearby
fibres are the original tori. Catanese's theorem
\cite[definition of deformation equivalence, p.~261; Theorem~2.1 and proof, pp.~264--268]{Catanese2004} states that every
deformation of a complex torus is a complex torus; its definition
uses precisely a proper holomorphic submersion over a connected
disc. The normalized cyclic base change constructed above is such
a proper holomorphic submersion, so its central fibre is a torus;
Ehresmann's theorem also shows that it is connected. This gives
the finite unramified torus cover of $G$.

\end{enumerate}
\end{proof}

\begin{remark}[Euler values of irreducible fibre supports]\hypertarget{tag-02Y1}{}\label{night:P4.9}\label{stmt:P4.9}
\StatementRecord{night:P4.9}{proved}{02Y1}{\SeriesRef{night:P4.26}, \SeriesRef{night:P4.32}, \SeriesRef{night:P4.35}, \SeriesRef{night:P4.37}, \SeriesRef{night:P4.54}, \SeriesRef{night:P4.57}, \SeriesRef{night:P4.58}, \SeriesRef{night:P4R.18}, \SeriesRef{night:P4R.20}}

The sign and equality classifications are given by Theorem~\ref{night:P4.58}
for the same morphism $f:Y\to B$ and the same reduced fibre supports.
In particular, each irreducible singular support whose normalization has
a rational resolution has Euler number two, and every other irreducible
support has Euler number zero. This is the identity specialization of
that theorem to the irreducible supports of its fibration.
\end{remark}

We checked CDP 2020, Lemma~4.6, Propositions~4.7,
4.11 and~7.3, deriving the sign and classification independently
from its local conductor identities. P4.14 records the precise
incorrect curve-classification inference in Proposition~7.3.
Earlier proof obligations and unsuccessful reductions are retained
in the working notes; the final classification is P4.58.
The original stated-value row also claimed a sign for arbitrary
reducible supports without normal crossings. Such supports are
outside the two classes explicitly requested in Part~4.2; the
proved global assertion retains the stated irreducible-or-NC
support hypothesis.

\subsection{Explicit checks on the hypotheses}

\begin{proposition}[A projective torus fibration with Chern number twenty-four]\hypertarget{tag-02WH}{}\label{night:P4.10}\label{stmt:P4.10}
\StatementRecord{night:P4.10}{proved}{02WH}{the definitions and constructions stated in the proof}

There is a smooth projective threefold $Z$ with $b_1(Z)=0$ and a
connected torus fibration $Z\to\Pp^1$ whose reduced special fibres
are irreducible normal-crossings surfaces, but $c_3[Z]=24$.
This example has $b_2(Z)>0$ and the fibre types of
Proposition~\ref{night:P4.7}; its value $c_3[Z]=24$ shows that the
bound in that proposition depends on the hypothesis $b_2=0$.
\end{proposition}
\begin{proof}

\begin{enumerate}[label={$\langle1\rangle$\arabic*.},leftmargin=*]
\item On the two affine charts of $\Pp^1$, with $s=t^{-1}$, we take
\[
 y^2z=x^3-3xz^2+t^6z^3,\qquad
 y_\infty^2z=x_\infty^3-3s^4x_\infty z^2+z^3,
\]
glued by $x_\infty=s^2x$ and $y_\infty=s^3y$. These are the
Weierstrass cubics in $\Pp(\OO\oplus\OO(2)\oplus\OO(3))$ for
the line bundle $\OO_{\Pp^1}(1)$. They define an elliptic surface
$\pi:S\to\Pp^1$ with zero section $[0:1:0]$. The discriminant
is a nonzero constant times $t^{12}-4$, or $1-4s^{12}$ on the
other chart. There are twelve simple zeros and a smooth fibre at
infinity. At a discriminant zero the cubic has a double root
$x=1$ or $x=-1$ and a distinct third root, so the fibre has an
ordinary node. Since $6t^5\ne0$ there, the total surface is smooth.
At the zero section the derivative in $z$ is nonzero. The map near
each node is therefore $xy=t$ after holomorphic coordinate changes
and translation of the base.

\item The relative dualizing differential $dx/y$ has transitions
$\pi^*\OO(1)$ and generates the dualizing sheaf also at the nodes.
Thus $\omega_S=\pi^*\OO(-1)$. Relative duality gives
$R^1\pi_*\OO_S=\OO(-1)$, and connected fibres give
$\pi_*\OO_S=\OO$. Leray yields $h^1(\OO_S)=0$, while all
positive plurigenera vanish. The classification \cite[Chapter~VI, Theorem~1.1, p.~243]{BHPV2004E} makes $S$
rational; in the ruled case the base genus equals $h^1(\OO_S)=0$.
In particular $\pi_1(S)=1$, since the minimal rational models
$\Pp^2$ and the Hirzebruch surfaces are simply connected and a
point blowup preserves the fundamental group.

\item The projective fibre product $W=S\times_{\Pp^1}S$ has exactly
twelve singularities, at the pair of nodes in each nodal fibre.
Their equations are $xy=uv$, with base function $xy=uv$. We blow
up the coherent ideal of the diagonal $\Delta_S\subset W$. This
divisor is Cartier away from the twelve points, so the blowup is
the identity there. Put $r=x-u$, $h=y-v$ at a singular point.
The equation becomes $ry+uh=0$ and the diagonal ideal is $(r,h)$.
The blowup has charts
\label{layout:II:1}
\[
\begin{aligned}
h&=ar, & y&=-au, & x&=u+r, & f&=-aux,\\
r&=bh, & u&=-by, & x&=-bv, & f&=-byv.
\end{aligned}
\]
The independent coordinates are $(a,u,x)$ and $(b,y,v)$.
Both charts are smooth, with $b=a^{-1}$ on the overlap, and the
exceptional fibre is $\Pp^1$. Thus the resulting $Z$ is smooth
and projective. The displayed equations make every special fibre
reduced and normal crossings. Its support remains irreducible:
the original product of two irreducible nodal cubics is irreducible,
and the exceptional locus has dimension one, so introduces no
surface component. Its normalization is rational, by the
birational map from $\Pp^1\times\Pp^1$.

\item A nodal cubic has Euler number $2-1=1$, so its self-product has
Euler number one. The resolution replaces one point by $\Pp^1$;
each resolved special support therefore has Euler number two.
All other fibres are $E_t\times E_t$ with Euler number zero.
Hence $e(Z)=12\cdot2=24=c_3[Z]$.

\item We compute $\pi_1(Z)$. Every loop can be moved
off the fibre divisors, so the fundamental group of their complement
surjects onto $\pi_1(Z)$. The complement is a bundle with fibre
$E\times E$ and a section over a twelve-punctured sphere. The
section extends over $\Pp^1$, killing the base generators. The
two maps
\[
 S\longrightarrow Z:\quad p\mapsto(p,O(\pi(p))),\qquad
 p\mapsto(O(\pi(p)),p)
\]
are defined because their images miss every conifold point, and
extend the two factor inclusions of $E$. Since $S$ is simply
connected they kill both fibre fundamental groups. Thus
$\pi_1(Z)=1$ and $b_1(Z)=0$. A projective Kähler class is nonzero
in $H^2(Z;\R)$, because its cube has positive integral. Thus $b_2(Z)>0$, as asserted.

\end{enumerate}
\end{proof}

\begin{proposition}[A primary Hopf bundle on a product of spheres]\hypertarget{tag-02WI}{}\label{night:P4.11}\label{stmt:P4.11}
\StatementRecord{night:P4.11}{proved}{02WI}{\SeriesRef{night:P4.7}}

There is a smooth compact complex threefold $C$ with
$b_1(C)=b_2(C)=0$, $e(C)=0$, and a holomorphic submersion
$C\to\Pp^1$ whose fibres are primary Hopf surfaces. No fibre has
a finite unramified torus cover. The covering obstruction gives a counterexample to the smooth-support equality description of \SeriesRef{cdp:st:C.2.G.1} when its torus-fibre hypothesis is omitted. The numerical normal-crossings sign theorem \ref{night:P4.7} applies to this fibration.
\end{proposition}
\begin{proof}

\begin{enumerate}[label=\arabic*.,leftmargin=*]
\item The additive group $\C$ acts on
$V=(\C^2\setminus\{0\})\times(\C^2\setminus\{0\})$ by
\[
 \tau\cdot(z,w)=(e^\tau z,e^{i\tau}w).
\]
For $\tau=a+ib$, the logarithms of the two Euclidean norms change
by $a$ and $-b$. The action is free. It is proper: when both a
point and its translate lie in prescribed compact subsets of $V$,
the two norm ratios bound $a$ and $b$. Each orbit has a unique
representative with both norms one, obtained with
$a=-\log\|z\|$ and $b=\log\|w\|$. This gives an explicit
diffeomorphism $C=V/\C\cong S^3\times S^3$ and proves compactness.
Holomorphic transverse slices, for example $z_j=1$ after choosing
a local logarithm, give quotient charts with holomorphic
transitions. Hence $C$ is a smooth compact complex threefold.
The product cell complex gives free integral homology of ranks
$(1,0,0,2,0,0,1)$ and therefore the stated Betti and Euler numbers.

\item The invariant map $[z,w]\mapsto[z_1:z_2]$ is a holomorphic
submersion to $\Pp^1$. In an affine base chart write
$z=\lambda(1,\xi)$ and set $\lambda=1$ using the action.
The choices of logarithm differ by $2\pi ik$, changing $w$ by
$e^{-2\pi k}$. Thus we obtain the explicit local fibre bundle
with fibre
$H=(\C^2\setminus\{0\})/\langle w\mapsto e^{-2\pi}w\rangle$.
Its universal cover retracts onto the simply connected $S^3$, so
$\pi_1(H)=\Z$. A connected finite unramified torus cover would
inject $\Z^4$ as a finite-index subgroup of $\Z$. Every subgroup
of $\Z$ is cyclic, whereas $\Z^4$ has rank four, which is
impossible.

\end{enumerate}
\end{proof}

The executed program \texttt{code/part4/necessity\_examples.py}
checks the degree and squarefreeness of the discriminant, both
resolution substitutions, the Euler count $24$, and the product
cell ranks used above. Its complete output is supplied.

\subsection{Irreducible supports: canonical torsion and the normal case}

\begin{proposition}[Irreducible fibres of torus fibrations]\hypertarget{tag-02WJ}{}\label{night:P4.12}\label{stmt:P4.12}
\StatementRecord{night:P4.12}{proved}{02WJ}{\SeriesRef{night:P4.6}}

Let $Y$ be a connected smooth compact complex threefold with
$b_1(Y)=b_2(Y)=0$, and let $f:Y\to B$ be a proper surjective
holomorphic map with connected fibres to a smooth compact curve.
Assume that the smooth general fibre is a complex two-torus, and let
$f^{-1}(b)=mS$ have irreducible reduced support. Then
\[
 \omega_S^{\otimes m}\simeq\OO_S,\qquad
 \chi(\OO_S)=0,\qquad h^1(\OO_S)\in\{1,2\}.
\]
Every normal such support is smooth. For a singular, necessarily
nonnormal support, a minimal resolved model of its normalization
has Kodaira dimension $-\infty$.
\end{proposition}
\begin{proof}
\noindent\textbf{Step 1. Canonical torsion.}
Choose a disc $\Delta$ with no other critical value, with coordinate
$t$ vanishing at $b$, and write $U=f^{-1}(\Delta)$. The coherent
direct image $f_*\omega_U$ is torsion-free of rank one: nonzero
base functions act injectively on sections, and a nearby torus has
one-dimensional canonical space. It is therefore a line bundle
on the disc. A generator gives a section $v$ of $\omega_U$,
nonzero and nowhere vanishing on every noncentral fibre. Its
divisor is $aS$ for an integer $a\ge0$. If $a\ge m$, then $v/t$
would be holomorphic, contradicting primitivity of the generator
at zero. We have $0\le a<m$ and
\[
 \omega_U\simeq\OO_U(aS),\qquad\OO_U(mS)\simeq\OO_U.
\]
Adjunction gives $\omega_S\simeq\OO_U((a+1)S)|_S$, proving
$\omega_S^{\otimes m}\simeq\OO_S$. The ambient
Riemann--Roch computation in \ref{night:P4.6}, applied to $S$, gives
$\chi(\OO_S)=0$. Reduced irreducibility gives $h^0(\OO_S)=1$.
A nonzero section of a torsion line bundle has a power equal to
a nonzero constant, and hence is nowhere zero; its space of
sections consequently has dimension at most one. Serre duality
on the Gorenstein surface $S$ gives
$h^2(\OO_S)=h^0(\omega_S)\in\{0,1\}$, proving the assertion
about $h^1$.

\noindent\textbf{Step 2. Exceptional curves and discrepancies.}
For a minimal resolution $\pi:T\to S$ of a normal Gorenstein
surface write $K_T=\pi^*K_S+\sum a_jE_j$. Its exceptional
intersection matrix $M$ has nonnegative off-diagonal entries
and is negative definite. Indeed,
pulling back a generic holomorphic function vanishing at a
contracted point gives $H+\sum m_jE_j$, with $m_j>0$ and $H$
the strict transform. Intersecting gives $Mm\le0$, strictly in
at least one entry of each connected block. After scaling by
$\operatorname{diag}(m_j)$, the negative quadratic form is a
sum of nonnegative terms $c_{ij}(z_i-z_j)^2$ and $d_i z_i^2$,
with connected support and some $d_i>0$ in every block. It is
positive definite. This is the quadratic identity of
\cite[\S I, unnumbered assertion and proof, p.~6]{Mumford1961}.

Adjunction gives
$K_T\cdot E_j=2p_a(E_j)-2-E_j^2\ge0$: the sole negative
possibility is a smooth rational $(-1)$-curve, excluded by
minimality of the resolution. The inverse of a negative definite
matrix with nonnegative off-diagonal entries has nonpositive
entries. Indeed, put $Q=-M$ and choose $c$ larger than all its
eigenvalues and diagonal entries; then
$Q^{-1}=c^{-1}\sum_{k\ge0}(I-Q/c)^k$ is entrywise nonnegative.
The discrepancy equations $Ma=(K_T\cdot E_j)_j$ imply
$a_j\le0$. Hence $K_T=\pi^*K_S-E$ with $E$ an effective
exceptional integral divisor.

Put $p=\dim H^0(S,R^1\pi_*\OO_T)$. Leray, with this sheaf
supported at finitely many points, gives
\begin{equation}\hypertarget{tag-01OD}{}\label{eq:P4chi-resolution}
 \chi(\OO_T)=\chi(\OO_S)-p=-p.
\end{equation}
For an effective exceptional cycle $Z$, put $\mathcal F=I_Z$.
Fix $s\in S$ and its maximal ideal $\mathfrak m_s$. Each coherent
quotient $\mathcal F/\mathfrak m_s^k\mathcal F$ is supported on a
compact infinitesimal fibre of dimension at most one. Its cohomology
vanishes in degrees at least two. To see this, filter by the nilradical
and reduce to coherent sheaves on the reduced compact curve. Removing
point-supported torsion and pulling back to its finite normalization
leaves vector bundles on compact Riemann surfaces, up to kernels and
cokernels supported at finitely many points. Their Dolbeault resolutions
have length one, while point-supported sheaves have no higher
cohomology. The finite filtration therefore gives the asserted vanishing
on every infinitesimal fibre.
Analytic formal functions and coherence
\cite[\S6, Haupts\"atze~I and~IIa, p.~59]{Grauert1960D} identify
the completed stalk of $R^2\pi_*\mathcal F$ with the inverse limit of
these zero groups. Nakayama's lemma gives $R^2\pi_*I_Z=0$.
The exact sequence of $Z$ gives a surjection
$R^1\pi_*\OO_T\to R^1\pi_*\OO_Z$, and therefore
\begin{equation}\hypertarget{tag-01OE}{}\label{eq:P4exceptional-genus}
 p\ge h^1(Z,\OO_Z).
\end{equation}

\noindent\textbf{Step 3. Exclusion of normal singular supports.}
Suppose first that $E=0$. Then $K_T$ is torsion and the
classification \cite[Chapter~VI, Theorem~1.1, p.~243]{BHPV2004E} gives Kodaira dimension zero and nonnegative
holomorphic Euler characteristic. Equation
\eqref{eq:P4chi-resolution} forces $p=0$ and $\chi(\OO_T)=0$.
The remaining surfaces are tori, bielliptic surfaces and Kodaira
surfaces. They have no rational curves: a map from $\Pp^1$ to
a torus is constant, bielliptic surfaces have finite unramified
torus covers, and a primary Kodaira surface has an elliptic
fibration with smooth elliptic reduced fibres over an elliptic
curve. A map from $\Pp^1$ to a secondary Kodaira surface lifts
to its finite unramified primary Kodaira cover. On a torsion
canonical surface a negative curve is a smooth rational
$(-2)$-curve by adjunction. Thus there are no exceptional
curves, and $S$ is smooth.

Suppose next that $E\ne0$. Since $K_T+E$ is torsion, all
positive plurigenera vanish, so $T$ is rational, ruled, or a
blowup of a class-VII surface. Rationality contradicts
\eqref{eq:P4chi-resolution}, since $\chi(\OO_T)=1$. In class
VII that equation gives $p=0$, whereas
\[
 \chi(\OO_E)=-E\cdot(E+K_T)/2=0.
\]
The nonempty effective compact divisor $E$ has $h^0(\OO_E)\ge1$,
so $h^1(\OO_E)\ge1$, contradicting
\eqref{eq:P4exceptional-genus}.

Finally suppose $T$ is ruled, after point blowups, over a curve
$C$ of genus $g\ge1$. Equation \eqref{eq:P4chi-resolution}
gives $p=g-1$. A general ruling fibre $F$ has $K_T\cdot F=-2$,
so $E\cdot F=2$. An exceptional component $E_j$ therefore
maps onto $C$. Its normalization has genus at least $g$ by
Riemann--Hurwitz, so $h^1(\OO_{E_j})\ge g$. Taking $Z=E_j$
in \eqref{eq:P4exceptional-genus} gives $p\ge g$, a contradiction.

\noindent\textbf{Step 4. The resolved normalization of a nonnormal support.}
Write $\nu:N\to S$ for the finite normalization and $\pi:T\to N$
for its minimal resolution. The residue comparison of dualizing
sheaves gives a pole divisor along the conductor; its strict
transform $D$ on $T$ is nonzero and effective. Pull back a local generator of $\omega_S$ as a meromorphic
two-form. Its canonical divisor has the form $-D+\sum a_jE_j$.
The equations for the exceptional coefficients have right side
$(K_T+D)\cdot E_j\ge0$, because both $K_T\cdot E_j$ and
$D\cdot E_j$ are nonnegative. The same inverse-matrix argument
gives $a_j\le0$. Taking the $m$th power to trivialize the pulled
back $\omega_S$ consequently gives
\[
 mK_T+D'\sim0
\]
for a nonzero effective integral divisor $D'$. This is also the
conductor formula of \cite[Lemma~4.6, p.~7 of the attached-source edition]{CDP20}. Any
positive pluricanonical section on $T$, multiplied by the
appropriate effective section and raised to a power, would give
a nonzero holomorphic function on compact $T$ with a zero.
It follows that $\kappa(T)=-\infty$. Contraction of exceptional
curves of the first kind preserves the plurigenera, so the minimal
model has the same Kodaira dimension.

\end{proof}

The ruled-genus and nonnormal-conductor arguments needed after
this reduction are proved below and completed in P4.58.

\begin{proposition}[Monodromy and the smooth locus of an irreducible fibre]\hypertarget{tag-02WK}{}\label{night:P4.13}\label{stmt:P4.13}
\StatementRecord{night:P4.13}{proved}{02WK}{the definitions and constructions stated in the proof}

Let a proper holomorphic degeneration of two-dimensional complex
tori have smooth total threefold and central fibre $mS$ with
irreducible reduced support. If $A$ is the monodromy on $H_1$ of
a nearby torus, then
\[
 m e(S_{\mathrm{reg}})=\det(I-A^m),\qquad
 0\le e(S_{\mathrm{reg}})\le16/m.
\]
\end{proposition}
\begin{proof}

\begin{enumerate}[label=\arabic*.,leftmargin=*]
\item Choose a proper holomorphic resolution of the pair
$\pi:U'\to U$ which is an isomorphism away from $S$ and over
$S_{\mathrm{reg}}$. The total central divisor has normal crossings.
The strict transform has multiplicity $m$, and each exceptional
component has multiplicity at least $2m$. Indeed, its first centre
lies over a singular point of the reduced divisor, whose order is
at least two; subsequent centres either retain such an exceptional
multiplicity or add further positive multiplicities.

\item At a point on $k$ branches,
with local equation $z_1^{n_1}\cdots z_k^{n_k}=t$, the nearby
fibre retracts onto a disjoint union of real $(k-1)$-tori.
Monodromy permutes these components and translates any component
that it preserves. For $k\ge2$ the alternating cohomological
trace is zero. On a single branch of multiplicity $n$ there are
$n$ contractible nearby components in a cyclic permutation;
the $d$th iterate has trace $n$ when $n\mid d$, and zero otherwise.
Let $F$ be the nearby torus and $T$ its induced cohomology monodromy. Alternating traces in the nearby-cycle spectral sequence, with a
finite stratification of the compact central divisor, give
\[
 L(T^d;H^*(F;\C))=\sum_{n_j\mid d}n_j e(E_j^\circ).
\]
This is the proper analytic trace formula of \citet[Theorem~1, p.~235; \S\S2--3, pp.~237--245]{ACampo1975}, obtained from the displayed local traces and the nearby-cycle spectral sequence.
For $d=m$ the strict transform is the sole contribution, and its
open stratum is $S_{\mathrm{reg}}$.

\item For a common multiple $N$
of the $n_j$, the $N$th power of monodromy acts trivially on all local
nearby-cohomology sheaves, hence on the $E_2$ page. This $N$th power
acts unipotently on the filtered abutment. Every
eigenvalue of $A$ is a root of unity. Under the evaluation pairing $H^1(F;\C)=H_1(F;\C)^\vee$, the induced cohomology action is $A^T$, since $(T\alpha)(c)=\alpha(Ac)$. The cohomology of the nearby torus $F$ is the exterior algebra on $H^1(F;\C)$, so its Lefschetz number is
$\det(I-(A^m)^T)=\det(I-A^m)$. Complex conjugate eigenvalues contribute
$|1-\lambda^m|^2$; real eigenvalues are $1$ or $-1$ and
contribute zero or two. Thus this determinant is nonnegative
and at most $2^4$.

\end{enumerate}
\end{proof}

\begin{proposition}[Genus-one curves with nonreduced conductor]\hypertarget{tag-02WL}{}\label{night:P4.14}\label{stmt:P4.14}
\StatementRecord{night:P4.14}{proved}{02WL}{the definitions and constructions stated in the proof}

The reducible-curve inference in \cite[proof of Proposition~7.3,
p.~16 of the attached-source edition]{CDP20} omits the two plane cubics
\[
 Z(Y^2-XZ)=0,\qquad XY(X-Y)=0.
\]
They are reduced Gorenstein curves of arithmetic genus one with
trivial dualizing sheaf and rational normalization, but have
respectively two tangent components and three concurrent
components. Their conductor order is two on every normalization
branch. These curves refute the asserted nodal-cycle
classification at that step.
\end{proposition}
\begin{proof}

\begin{enumerate}[leftmargin=*]
\item Plane adjunction gives $\omega_C=\OO_C(3-3)=\OO_C$ and
$p_a(C)=1$. The first cubic has a line tangent to a smooth
conic at $[1:0:0]$. Its local equation is $y(y-x^2)=0$.
The normalization restrictions are $(p(t,0),p(t,t^2))$.
Their image is exactly the pairs $(a,b)$ with $b-a$ divisible
by $t^2$: the inverse expression is
$p(x,y)=a(x)+y(b(x)-a(x))/x^2$. Thus the normalization
quotient has length two and the conductor exponent is two on
each branch.

\item For the second cubic the branches are $(t,0),(0,t),(t,t)$.
All constant coefficients must agree, and the third linear
coefficient is the sum of the first two. These three conditions
are sufficient. In every degree $k\ge2$, the monomials
$x^k-x^{k-1}y$, $y^k-xy^{k-1}$ and $x^{k-1}y$ independently
realize the coefficients on the three branches. Hence the
normalization quotient has length three and each conductor
exponent is again two. The normalization Euler characteristics
are two and three; subtracting these quotient lengths gives
$\chi(\OO_C)=0$ in both cases.

\item Let $C_2$ and $C_3$ be the two cubics above. The normalization
of $C_2$ is the disjoint union of two copies of $\Pp^1$, parametrized
by $[s:t]\mapsto[s:t:0]$ on the line and
$[s:t]\mapsto[s^2:st:t^2]$ on the conic. The parameters $t/s$ agree
at the two preimages of the tacnode. Exchanging these copies with
these parameters preserves their local ring
$\{(a,b):b-a\in(t^2)\}$ and descends to an involution $\sigma_2$
of $C_2$. In the local coordinates above it is
$(x,y)\mapsto(x,x^2-y)$. The projective automorphism
\[
 \sigma_3[X:Y:Z]=[-Y:X-Y:Z]
\]
has order three and cyclically permutes the components of $C_3$.
For $r=2,3$, take $E'=\C/(\Z+i\Z)$ and let $\tau_r$ be
translation by $1/r$. Its order is exactly $r$, and every nonidentity
power acts freely. Thus the diagonal action defines the compact
surface and its map to an elliptic curve
\[
 X_r=(C_r\times E')/\langle(\sigma_r,\tau_r)\rangle
       \longrightarrow E=E'/\langle\tau_r\rangle.
\]
The cyclic group acts transitively on the components of
$C_r\times E'$, so $X_r$ is irreducible. Let $\widetilde\sigma_r$
be the lift to the normalization $\nu_r:C_r^\nu\to C_r$ induced
by the displayed parameter maps, so
$\nu_r\widetilde\sigma_r=\sigma_r\nu_r$. The normalization of $X_r$ is
$(\coprod_{j=1}^r\Pp^1\times E')/\langle(\widetilde\sigma_r,\tau_r)\rangle$.
Choose a component $C_{r,0}$ of $C_r^\nu$. Restriction of the quotient
map to $C_{r,0}\times E'$ is an isomorphism. Its inverse sends a class
represented by $(z,e)$ with $z\in\widetilde\sigma_r^j C_{r,0}$ to
$(\widetilde\sigma_r^{-j}z,\tau_r^{-j}e)$; this expression is unchanged
on each orbit and is holomorphic on every component.
The reduced inverse-image conductor is the single section
$\Gamma'=\{p\}\times E'$, its conductor divisor is $2\Gamma'$, and
its map to the singular curve $E$ is the finite unramified cover
$E'\to E$ of degree $r$. Analytically locally on $E$, the image
conductor has algebra $\OO_E[u]/(u^2)$ for $r=2$ and
$\OO_E[u,v]/(u,v)^2$ for $r=3$. These local algebras follow from
the normalization constraints in the first two steps and remain
nonreduced after the free quotient. The constructions give a single
horizontal conductor component of multiplicity two with generic
degree two or three over the image conductor.

\end{enumerate}
\end{proof}

The branch permutations in P4.14 have explicit local realizations.
Take an elliptic curve $C$ with a free translation of order two or
three. The map $(x,y)\mapsto(x,x^2-y)$ exchanges the two tacnode
branches and preserves $y(y-x^2)$. The order-three map
$(x,y)\mapsto(-y,x-y)$ cycles the three branches and preserves
$xy(x-y)$. The diagonal action on $C$ times a small invariant
neighbourhood of the origin is free. Its quotient is a smooth
threefold; the divisor is irreducible because the action permutes
its branches transitively. Its normalization conductor is one
elliptic curve mapping to $C/\langle\tau\rangle$ with degree two
or three. These specify the finite morphisms in P4.14 explicitly;
they make no assertion of a compact torus smoothing.

\begin{proposition}[Finite rank drop of a normalization]\hypertarget{tag-02WM}{}\label{night:P4.15}\label{stmt:P4.15}
\StatementRecord{night:P4.15}{proved}{02WM}{\SeriesRef{cdp:st:A.5.3}, \SeriesRef{night:P4.12}, \SeriesRef{night:P4.13}}

Let $f:Y\to B$ be a proper surjective holomorphic map with
connected fibres from a connected smooth compact complex threefold
with $b_1(Y)=b_2(Y)=0$ to a smooth compact curve, and suppose that
its smooth general fibre is a complex two-torus. Let $mS$ be a full
fibre whose reduced support $S$ is irreducible and singular.
Suppose its normalization $\nu:T\to S$ has smooth domain $T$. Set
$j:T\xrightarrow{\nu}S\hookrightarrow Y$ and assume that
$dj:T_T\to j^*T_Y$ has rank two away from finitely many points.
Let $\ell$ be the length of its maximal-minor scheme. Then
\[
 e(T)=K_T^2-\ell.
\]
The alternatives for $T$ are a rational surface with
$e(T)=6-\ell/2$ and $0\leq\ell\leq6$, a minimal ruled surface
over an elliptic curve with $\ell=0$, or a Hopf surface with
$\ell=0$. In either of the last two alternatives $e(S)=0$.
In those last two alternatives, the reduced conductor $R\subset T$
and $C=\operatorname{Sing}(S)_{\mathrm{red}}$ are disjoint unions of
smooth elliptic curves, and the restriction $\nu|_R:R\to C$ is
finite and unramified.
\end{proposition}
\begin{proof}

\begin{enumerate}[label={$\langle1\rangle$\arabic*.},leftmargin=*]
\item The differential is injective as a sheaf map: its kernel vanishes
on a dense open set, and its locally free source is torsion-free.
Its rank-one cokernel $Q$ is locally free away from finitely many
points. The depth lemma applied to the length-one locally free
resolution gives depth at least one at those points. Hence $Q$
has no finite-length subsheaf, and all its possible torsion would
be supported there. Thus $Q$ is torsion-free, and $L=Q^{**}$ is a
line bundle. In local frames the signed maximal-minor row of $dj$
annihilates its two columns. Its determinant transition functions
give a map to $L=\det(j^*T_Y)\otimes\det(T_T)^{-1}$.
The induced map $Q\to L$ is generically injective, hence injective
because $Q$ is torsion-free. Its image is the ideal generated by
those minors, tensored with $L$. We obtain the exact sequence
\[
 0\longrightarrow T_T\xrightarrow{dj}j^*T_Y
   \longrightarrow\mathcal I_Z\otimes L\longrightarrow0,
 \qquad L=K_T\otimes j^*(\det T_Y),\qquad
 \operatorname{length}(Z)=\ell.
\]
Poincar\'e duality gives $b_4(Y)=b_2(Y)=0$, so
$c_1(T_Y)=c_2(T_Y)=0$ in rational cohomology. Writing $k=c_1(K_T)$,
the Chern-character identity for this sequence is
\[
 3=\left(2-k+\frac{k^2-2c_2(T_T)}2\right)
     +\left(1+k+\frac{k^2}2-[Z]\right).
\]
Its degree-four part gives $e(T)=K_T^2-\ell$, and Noether's formula
gives $2e(T)+\ell=12\chi(\mathcal O_T)$.

Let $D$ be the full conductor divisor on $T$, and put $R=D_{\mathrm{red}}$.
Propositions~\ref{cdp:st:A.5.3} and~\ref{night:P4.12}, with the
identity as the resolution of $T$, give
$\omega_T(D)\simeq\nu^*\omega_S$ and $\omega_S^{\otimes m}\simeq\mathcal O_S$.
In particular $K_T+D\equiv0$. Proposition~\ref{night:P4.12}
also gives $\kappa(T)=-\infty$ and nonnormality of the singular $S$.
Since the hypersurface $S$ satisfies $S_2$, its nonnormal locus
contains a curve; thus $D$ is nonempty.
Let $q:T\to T_0$ contract $r$ point blowups to a minimal model.
Pluricanonical sections descend across each blowup by Hartogs'
extension on the complement of its centre, so $\kappa(T_0)=-\infty$
and $p_g(T_0)=0$.

For a K\"ahler $T_0$, the equality $p_g=0$ makes every real
degree-two cohomology class of type $(1,1)$. A sufficiently close
rational perturbation of a K\"ahler class, followed by integral
scaling, gives an integral K\"ahler class. Kodaira's embedding
theorem \cite[Main Theorem, p.~314]{Kodaira1954} therefore gives an
embedding $\Phi:T_0\hookrightarrow\mathbb P^N$. Apply
\cite[VI.17--VI.18, pp.~111--112; V.10, pp.~79--80; III.10, p.~40]{Beauville1978}
to $\Phi(T_0)$ and pull its ruling back by $\Phi$: the model is
rational or a projective-line bundle over a positive-genus curve.

For a non-K\"ahler $T_0$, \cite[Theorem~5.1, pp.~27--29]{OVV2020}
gives an isotrivial elliptic surface or a class-VII surface, with
overlap permitted. In the elliptic alternative, the local maps in
\cite[Theorem~4.2 and Remark~4.4, pp.~17--21]{OVV2020} are
\[
 (E\times\Delta_s)/\langle(z,s)\mapsto(z+e,\zeta_m s)\rangle
 \longrightarrow\Delta_t,\qquad t=s^m,
\]
where $e$ has exact order $m$. The elliptic action of
\cite[Proposition~4.5, pp.~21--22]{OVV2020} gives a nowhere-zero
vertical vector field. For its ruling $p:T_0\to B_0$, the quotient
of $T_{T_0}$ by this vector field maps to $p^*T_{B_0}$ with divisor
$R_p=\sum(m_i-1)E_i$, as follows from $dt=m_i s^{m_i-1}ds$.
Consequently $K_{T_0}=p^*K_{B_0}+R_p$, so $K_{T_0}^2=0$;
the vector field gives $c_2(T_0)=0$. Noether's formula and $p_g=0$
give $h^1(\mathcal O_{T_0})=1$.
The identity $b_1=2h^1(\mathcal O)-1$ in
\cite[Theorem~2.5, pp.~7--9]{OVV2020} gives $b_1(T_0)=1$.
Thus this elliptic alternative also has the class-VII invariants
$b_1=1$ and $\kappa=-\infty$.

For a projective-line bundle over a genus-$g$ curve, followed by
$r$ point blowups, the bundle and blowup formulas
\cite[II.3, pp.~17--18; III.18, pp.~46--47]{Beauville1978} give
\[
 e(T)=4(1-g)+r,\qquad K_T^2=8(1-g)-r.
\]
The Chern identity becomes $2r+\ell=4(1-g)$. For $g\ge1$ it forces
$g=1$ and $r=\ell=0$.
For a class-VII minimal model, write $b_2(T_0)=k_0\ge0$.
The equalities $b_1=1$, $p_g=0$ and
$b_1=2h^1(\mathcal O)-1$ give $\chi(\mathcal O_{T_0})=0$.
Thus $e(T)=k_0+r$ and $K_T^2=-k_0-r$, and the Chern identity gives
$2(k_0+r)+\ell=0$. Hence $T$ is minimal with $b_2=0$ and $\ell=0$.
By \cite[Theorem~1, p.~1]{BBK2019} it has a finite \emph{\'etale}
cover $a:T'\to T$ with $T'$ Hopf or Inoue. The inverse image
$a^{-1}(R)$ is a nonempty curve. Inoue surfaces contain no curves
by \cite[Proposition~2(i), pp.~274--275; Proposition~3(i), p.~277;
paragraph following (21), p.~279]{Inoue1974}, so $T'$ is Hopf.
The composite of its universal cover
$\mathbb C^2\setminus\{0\}\to T'$ with $a$ is the universal cover
of $T$, since $\mathbb C^2\setminus\{0\}$ is simply connected.
Therefore $T$ is Hopf.

For rational $T$, $\chi(\mathcal O_T)=1$ and $e(T)\ge3$ by the
same rational minimal-model and blowup formulas. Hence
$e(T)=6-\ell/2$ and $0\le\ell\le6$.
When $\ell=0$, the maximal-minor ideal is the unit ideal, so $dj$
has rank two everywhere.

\item Consider the two latter alternatives. First let $T$ be a
minimal ruled surface over an elliptic curve $E$. Its ruling
$p:T\to E$ is a $\mathbb P^1$-bundle by
\cite[III.10, p.~40]{Beauville1978}. For a ruling fibre $F$,
the equality $K_T+D\equiv0$ gives $D\cdot F=2$.
If a horizontal component $H$ occurs with multiplicity one,
adjunction gives
\[
 2p_a(H)-2=(K_T+H)H=-(D-H)H\leq0.
\]
The map to $E$ gives geometric genus at least one. Equality
therefore holds, $H$ is smooth elliptic, and $H$ is disjoint
from every other component of $D$. There is no vertical
component because every ruling fibre meets $H$. The possibilities
are one elliptic bisection, of degree two over $E$, or two
disjoint elliptic sections.

\item The other possibility is $D=2H+\sum_j a_jF_j$, where $H$ is a
section, $F_j$ are distinct ruling fibres and $a_j>0$.
The normalization of $H$ maps isomorphically to $E$ by its
degree-one ruling map. Its map into $T$ is a section of $p$,
so its image $H$ is smooth elliptic.
The reduced conductor $R$ has Euler number equal to the number $v$
of these fibres: $e(H)=0$, $e(F_j)=2$, and each meets $H$ once.
The normalization restricts to an isomorphism
$T\setminus R\simeq S_{\mathrm{reg}}$. Hence
\[
 e(S_{\mathrm{reg}})=e(T)-e(R)=-v.
\]
Proposition~\ref{night:P4.13} forces $v=0$. Thus in every case the reduced conductor is
a disjoint union of smooth elliptic curves.

\item On a Hopf surface every irreducible curve $H$ satisfies
$H^2=K_TH=0$, so $p_a(H)=1$. A rational curve would lift from
$\mathbb P^1$ to the universal cover $\mathbb C^2\setminus\{0\}$,
where both coordinate functions must be constant. Thus $H$
is smooth elliptic. Distinct curves have intersection zero and
are disjoint, proving the same conductor description.

\item At a point of $C=\operatorname{Sing}(S)_{\mathrm{red}}$,
the finite immersive normalization presents the germ as distinct
smooth surface branches $S_i$. The reduced conductor on each branch is a single
smooth curve germ $C_i$. Any two branches intersect in a curve
in the smooth ambient threefold, and its support lies in both
$C_i$ and $C_j$. Their irreducibility forces those curve germs
to agree. Thus the image singular locus is one smooth curve
germ, and the conductor maps locally isomorphically to it.
Globally the restriction $\nu|_R:R\to C$ is finite and unramified.
For a component of degree $d$, Riemann--Hurwitz gives
$0=d(2g(C_i)-2)$, so every image component is elliptic.
Euler additivity for the normalization and its conductor gives
\[
 e(S)=e(T)-e(R)+e(C)=0.
\]

\end{enumerate}
\end{proof}

\begin{proposition}[A spherical intersection obstruction]\hypertarget{tag-02WN}{}\label{night:P4.16}\label{stmt:P4.16}
\StatementRecord{night:P4.16}{proved}{02WN}{the definitions and constructions stated in the proof}

For every integer $n\geq1$, the reduced hypersurface germ
$x^2-zy^n=0$ cannot occur in a fibre of a proper holomorphic
degeneration of complex two-tori with smooth total threefold,
even with any positive common multiplicity on that germ.
In particular the Whitney umbrella ($n=2$) is excluded.
The total threefold may have arbitrary second Betti number.
\end{proposition}
\begin{proof}

\begin{enumerate}[label={$\langle1\rangle$\arabic*.},leftmargin=*]
\item Put $F_n=\{x^2-zy^n=1\}\subset\mathbb C^3$. We exhibit its
compact spherical homology class and compute its square.
The open subset $U=\{y\ne0\}$ is
$\mathbb C^*\times\mathbb C$ in coordinates $(y,x)$.
For $0<\delta<1$, take the two disjoint graph neighbourhoods
\[
 N_\pm:\quad x=\pm\sqrt{1+zy^n},\quad |zy^n|<\delta.
\]
Their $(y,z)$ domains contract by simultaneous real scaling.
Each $N_\pm\cap U$ has coordinates $(y,zy^n)$ in
$\mathbb C^*\times\{|w|<\delta\}$. Its map to $U$ takes
the generating $y$-loop to the generating $y$-loop. These three
sets cover $F_n$. Van Kampen gives $\pi_1(F_n)=0$; see \cite[Theorem~1.20, p.~43]{Hatcher2002}.
Mayer--Vietoris \cite[\S2.2, pp.~149--150]{Hatcher2002} identifies $H_2(F_n;\mathbb Z)$ with the
kernel of $(a,b)\mapsto a+b$ on the two intersection
$H_1$ groups, hence with $\mathbb Z$. The Hurewicz theorem
\cite[Theorem~4.32, pp.~366--367]{Hatcher2002} represents its
generator by a map $S^2\to F_n$.

\item We compute the intersection form by compactification. Blow up
$\mathbb P^1_x\times\mathbb P^1_y$ successively $n$ times above
each point $(1,0)$ and $(-1,0)$. Near either point let
$a=x^2-1$, $b=y$. The successive charts at the centres are
\[
 a=b^j a_j,\qquad z=\frac{a_j}{b^{n-j}}.
\]
For $j<n$, blow up the next centre $a_j=b=0$.
At $j=n$, $z=a_n$ is regular. The final exceptional curve
minus its point on the preceding boundary is precisely the
line $y=0,x=\pm1$ in $F_n$. Every earlier exceptional curve
is a pole divisor of $z$. Consequently $F_n=T_n\setminus B$,
where $T_n$ is this smooth projective blown-up surface and
$B$ consists of the two infinity divisors, the strict transform
of $y=0$, and the first $n-1$ exceptional curves on each chain.

\item In the total-transform basis $H_x,H_y,E_{\pm,j}$ the nonzero
products are $H_xH_y=1$ and $E_{\pm,j}^2=-1$. The boundary
classes are
\[
 H_x,\ H_y,\ H_y-E_{+,1}-E_{-,1},\quad
 E_{+,j}-E_{+,j+1},\ E_{-,j}-E_{-,j+1}\quad(1\leq j<n).
\]
They form a connected tree of rational curves; thus $H_1(B)=0$.
Compact-support duality \cite[Theorem~3.35, p.~245]{Hatcher2002}
and the sequence of $(T_n,B)$ identify
the image of $H_2(F_n)$ in $H_2(T_n)$ with the kernel of
intersection against these boundary classes. Orthogonality to
$H_x,H_y$ removes both ruling coefficients. Orthogonality to
the exceptional differences makes the coefficients constant
along each chain. The strict transform of $y=0$ makes the two
constants opposite. Thus the kernel is generated by the
primitive integral vector
\[
 v=\sum_{j=1}^n E_{+,j}-\sum_{j=1}^n E_{-,j},\qquad v^2=-2n.
\]
Since $H^1(B;\mathbb Z)=0$, the pair sequence injects
$H^2(T_n,B;\mathbb Z)$ into $H^2(T_n;\mathbb Z)$ with
precisely this integral kernel.
Therefore the compact spherical generator has square $-2n$.
The corresponding computation forms the intersection matrix from
these boundary classes and calculates its primitive integral kernel;
its source is listed in Table~\ref{tab:II:programs}.

\item The defining polynomial is weighted homogeneous of degree
$d=2n+4$ with positive weights $(n+2,2,4)$. Scaling by
\[
 (x,y,z)\longmapsto
 (\varepsilon^{n+2}x,\varepsilon^2y,\varepsilon^4z)
\]
identifies $F_n$ with the fibre of value $\varepsilon^{2n+4}$.
It carries any compact representative of $v$, together with a
compact transverse representative for its square, into an
arbitrarily small neighbourhood of the origin.

\item Write the local defining function as $u(x^2-zy^n)^m$, with $u$
a unit, and choose a holomorphic root $h^m=u$. Then
the local analytic isomorphism
\[
 \Phi(x,y,z)=
 \bigl(h^{(n+2)/d}x,h^{2/d}y,h^{4/d}z\bigr)
\]
satisfies $g\circ\Phi=hg$, where $g=x^2-zy^n$; its derivative
at the origin is invertible. Thus the local map to the
base is $(g\circ\Phi)^m$. The small local fibre of $g$ embeds
in the corresponding smooth compact torus fibre.

\item The spherical representative is null-homotopic in that torus,
whose universal cover is the contractible space $\mathbb C^2$.
It is therefore zero in $H_2$, and its self-intersection there
is zero. The two compact representatives already chosen compute
the same intersection in the open local fibre and in the torus.
Their intersection is $-2n$, which is a contradiction.

\end{enumerate}
\end{proof}

P4.16 excludes the displayed pinching and ramification germs.
The subsequent arguments retain singular normalizations and
nonimmersive normalization maps and complete their classification
in P4.58.

\begin{proposition}[The lattice of $x^a-zy^n=1$]\hypertarget{tag-02WO}{}\label{night:P4.17}\label{stmt:P4.17}
\StatementRecord{night:P4.17}{proved}{02WO}{\SeriesRef{night:P4.16}}

For all integers $a\geq2$ and $n\geq1$, the germ $x^a-zy^n=0$ is
excluded in the setting of Proposition~\ref{night:P4.16}. Its affine nearby fibre has
$\pi_1=0$ and compact second-homology intersection lattice
$A_{a-1}(-n)$. This includes the order-three cyclic-ramification
germ $x^3-zy^3$ and examples with singular normalization.
\end{proposition}
\begin{proof}

\begin{enumerate}[label={$\langle1\rangle$\arabic*.},leftmargin=*]
\item Put $F_{a,n}=\{x^a-zy^n=1\}$. Replace the two roots of
$x^2=1$ in Proposition~\ref{night:P4.16} by the $a$ distinct
roots of $x^a=1$. The open set $y\ne0$ is still
$\mathbb C^*\times\mathbb C$. There are $a$ contractible graph
neighbourhoods, each meeting it in $\mathbb C^*\times\Delta$
with the same generating $y$-loop. Thus van Kampen gives
$\pi_1=0$, and Mayer--Vietoris gives
\[
 H_2(F_{a,n};\mathbb Z)=\ker\bigl(\mathbb Z^a
   \xrightarrow{(c_i)\mapsto\sum_i c_i}\mathbb Z\bigr).
\]
At each root, $x^a-1$ is a local coordinate. The same $n$
successive blowup charts give a compactification with $a$
exceptional chains and boundary classes
\[
 H_x,\ H_y,\ H_y-\sum_i E_{i,1},\qquad
 E_{i,j}-E_{i,j+1}\quad(1\leq i\leq a,\ 1\leq j<n).
\]
This boundary is a tree of rational curves. The exact sequence
of the pair identifies compact $H_2$ integrally with the kernel
of intersection against these classes. Its primitive basis is
\[
 v_i=\sum_{j=1}^n E_{i,j}-\sum_{j=1}^n E_{i+1,j}
       \quad(1\leq i<a).
\]
The products are $v_i^2=-2n$, $v_iv_{i+1}=n$, with all other
off-diagonal products zero. This is $A_{a-1}(-n)$.
Hurewicz makes every class spherical. Positive weights
$(n+1,a,a)$, of degree $a(n+1)$, scale compact representatives
into arbitrarily small nearby fibres. The unit-removal map and torus argument in
Proposition~\ref{night:P4.16}, with these weights, exclude $v_1$,
whose square is $-2n$. The computation of the integral kernel
from the displayed boundary classes is listed in
Table~\ref{tab:II:programs}.

\item We specify the normalization morphism. Let $\mu_a$ act on
$\mathbb C^2_{u,v}$ by $(u,v)\mapsto(\zeta u,\zeta^{-n}v)$.
Then
\[
 x=u^nv,\qquad y=u^a,\qquad z=v^a
\]
gives a finite map to $x^a=zy^n$: both $u$ and $v$ satisfy
monic equations over its coordinate ring. The function-field
degree is $a$, with the displayed faithful $\mu_a$ action.
Consequently the induced quotient map is finite and birational.
The invariant ring is integrally closed, so this is the
normalization. For $(a,n)=(2,3)$ it has coordinates
$U=u^2,V=uv,W=v^2$ with $UW=V^2$, and its map is
$(x,y,z)=(UV,U,W)$. Thus it has an $A_1$ singularity.
For $n=a$ the quotient is smooth with coordinates $(y,v)$,
and the reduced conductor $\{y=0\}$ maps to $\{x=y=0\}$
by $z=v^a$. A positive loop about $z=0$ sends $v$ to
$e^{2\pi i/a}v$, cyclically permuting the $a$ branches.

\end{enumerate}
\end{proof}

\begin{proposition}[Spherical classes for binomial surface germs]\hypertarget{tag-02WP}{}\label{night:P4.18}\label{stmt:P4.18}
\StatementRecord{night:P4.18}{proved}{02WP}{\SeriesRef{night:P4.15}, \SeriesRef{night:P4.16}}

For $a\geq2$ and $b,c\geq1$, the reduced germ
$x^a-y^bz^c=0$ cannot occur with a common positive multiplicity
in a proper holomorphic degeneration of complex two-tori with
smooth total threefold. Its affine
nearby fibre contains a spherical class of nonzero square.
In particular this excludes $x^2-z^3y^4$, whose smooth
normalization has a divisorial differential rank drop.
\end{proposition}
\begin{proof}

\begin{enumerate}[label={$\langle1\rangle$\arabic*.},leftmargin=*]
\item Write $d=\gcd(b,c)$, $b=dB$, $c=dC$, so $\gcd(B,C)=1$.
The nearby fibre $F=\{x^a-y^bz^c=1\}$ maps to the smooth
affine curve $A=\{x^a-w^d=1\}$ by $w=y^Bz^C$.
The function $w$ on $A$ has $a$ simple zeros. Choose an
embedded arc between two zeros avoiding the others, and a
small open disc $D$ around it. In a holomorphic coordinate
$\xi$ on $D$, write $w=u(\xi)(\xi-\alpha)(\xi-\beta)$
with $u$ a unit. Set
\[
 X=\frac{2\xi-\alpha-\beta}{\beta-\alpha},\qquad
 h_0(\xi)=\frac{(\beta-\alpha)^2}{4}u(\xi),\qquad
 y_{\mathrm{new}}=\frac{y}{h_0(\xi)^{1/B}}.
\]
A holomorphic $B$-th root exists on the disc. In these coordinates,
writing $y$ for $y_{\mathrm{new}}$, the preimage of $D$ is
\[
 W=\{y^Bz^C=X^2-1,\ X\in D'\},
\]
where $D'$ is a disc containing both $-1$ and $1$.

\item Choose $Bp+Cq=1$. With $h=X^2-1$, the explicit trivialization
away from the two special fibres is
\[
 y=h^pt^C,\quad z=h^qt^{-B},\quad t=y^qz^{-p}.
\]
Its fundamental group is $F(\beta_-,\beta_+)\times\mathbb Z\tau$.
The neighbourhood of each special fibre is the graph
$X=\pm\sqrt{1+y^Bz^C}$ on $|y^Bz^C|<\delta$, a contractible
domain by real scaling. Its punctured part has two generating
loops in $y,z$, whose images in $\langle\beta_\pm,\tau\rangle$
are the columns of
\[
 M=\begin{pmatrix}B&C\\q&-p\end{pmatrix},\qquad
 M^{-1}=\begin{pmatrix}p&C\\q&-B\end{pmatrix}.
\]
The equality $Bp+Cq=1$ proves this integral inverse.
Thus adjoining each neighbourhood kills $\beta_\pm$ and $\tau$.
Van Kampen gives $\pi_1(W)=0$.

\item The map
\[
 \pi:W\longrightarrow W'=\{y^BZ=X^2-1,\ X\in D'\},
 \qquad (X,y,z)\longmapsto(X,y,z^C)
\]
is proper finite of degree $C$. The compact class $v$ in
Proposition~\ref{night:P4.16}, with $n=B$ and square $-2B$,
has a representative supported over $D'$.
Indeed, the cover just given with $C=1$ has the identical
Mayer--Vietoris maps before and after enlarging $D'$ to the
whole plane: both punctured bases have the same two meridians,
and the same two graph neighbourhoods are attached. The
inclusion is therefore an isomorphism on $H_2$ and preserves
the primitive class and its square.
Let $\alpha\in H_c^2(W';\mathbb Z)$ be its compact-support dual
under \cite[Theorem~3.35, p.~245]{Hatcher2002}. Proper pullback
preserves compact support, and the degree formula gives
\[
 \int_W(\pi^*\alpha)^2=C\int_{W'}\alpha^2=-2BC.
\]
One may verify this by smooth compactly supported forms:
outside the measure-zero branch divisor the map is an oriented
$C$-sheeted covering. The dual of $\pi^*\alpha$ is spherical
by $\pi_1(W)=0$ and \cite[Theorem~4.32, pp.~366--367]{Hatcher2002}. Its compact representatives
embed in $F$ with the same nonzero square.

\item The polynomial has positive weights $(b+c,a,a)$. Weighted
scaling places these representatives in arbitrarily small
nearby fibres. The unit-removal map in
Proposition~\ref{night:P4.16}, with these weights, handles
the unit and common multiplicity. The spherical class becomes zero
in a torus, contradicting its square $-2BC$.

\item For the mixed example the normalization is the finite
birational map
\[
 (u,v)\longmapsto(u^3v^2,v,u^2).
\]
Finiteness follows from $u^2=z$, and $u=x/(zy^2)$ gives
birationality. The image ring contains all even powers of $u$
and the odd monomials $u^iv^j$ with $i\geq3,j\geq2$.
Since the normalization is generated over the image ring by $1,u$,
its conductor consists of the series $h$ for which both $h$ and $uh$
lie in that ring. The even--odd monomial description therefore gives
$u^2v^2\mathbb C\{u,v\}$, with divisor
$2\{u=0\}+2\{v=0\}$. Its differential minors are
$3u^2v^2,-4u^4v,-2u$, and generate the ideal $(u)$. Thus the rank-drop locus is exactly
the divisor $\{u=0\}$, which maps by $v\mapsto(0,v,0)$
to a line. The finite rank-drop hypothesis of
Proposition~\ref{night:P4.15} fails for this morphism.
The exponent-matrix, transfer-pairing and normalization computations
are listed in Table~\ref{tab:II:programs}.

\end{enumerate}
\end{proof}

\begin{proposition}[A ramified conductor with vanishing nearby homology]\hypertarget{tag-02WQ}{}\label{night:P4.19}\label{stmt:P4.19}
\StatementRecord{night:P4.19}{proved}{02WQ}{\SeriesRef{night:P4.16}, \SeriesRef{night:P4.18}}

For the reduced hypersurface germ $y(x^2-zy^2)=0$, the normalization
conductor has multiplicity two on each component and its reduced
map to the singular locus is ramified. The affine nearby fibre is
$\mathbb C^*\times\mathbb C$. Every compact spherical class
in an open subset of this affine fibre has self-intersection zero.
\end{proposition}
\begin{proof}

\begin{enumerate}[label={$\langle1\rangle$\arabic*.},leftmargin=*]
\item One normalization component is the plane $y=0$; the other is
the Whitney normalization $(u,v)\mapsto(uv,v,u^2)$.
The singular locus is the smooth $z$-axis. The plane conductor
is $x^2=0$, obtained by restricting the other component's
equation. Write $R_B=\mathbb C\{v,uv,u^2\}$ and
$N_B=\mathbb C\{u,v\}$. The module $N_B$ is generated over
$R_B$ by $1,u$. A convergent series lies in $R_B$ precisely when
its odd powers of $u$ have coefficients divisible by $v$.
Thus $h,uh\in R_B$ is equivalent to $h\in vN_B$, proving
that the intrinsic conductor is $vN_B$.
The union ring is the fibre product of the component rings over
$\mathbb C\{x,z\}/(x^2)$. On the Whitney component, a
conductor element $h$ must satisfy $hN_B\subset vR_B$; the
intrinsic conductor calculation gives $h\in v^2N_B$.
On the plane, the corresponding kernel is $(x^2)$.
These are the two claimed conductor ideals.
The reduced conductor cover is the identity sheet from the
plane together with $u\mapsto u^2$ from the Whitney component.
The latter is ramified.

\item Put $F=\{y(x^2-zy^2)=1\}$. Here $y\ne0$, and the inverse maps
\[
 (x,y,z)\longmapsto(y,x),\qquad
 (y,x)\longmapsto\left(x,y,\frac{yx^2-1}{y^3}\right)
\]
identify this fibre with $\mathbb C^*\times\mathbb C$.
The retraction $(y,x)\mapsto(y/|y|,0)$ is a homotopy equivalence
with $S^1$, so $\pi_2$ and $H_2$ vanish. For an open inclusion
$j:U\hookrightarrow F$, compact-support duality identifies the
homology map $j_*$ with extension by zero $j_!$ on dual classes.
For the compact-support dual $\alpha$ of a compact spherical class,
\[
 \int_U\alpha^2=\int_F(j_!\alpha)^2=0.
\]
This proves the asserted zero self-intersection in $U$.

\end{enumerate}
\end{proof}

\begin{proposition}[Ordinary double strata and normal crossings]\hypertarget{tag-02WS}{}\label{night:P4.20}\label{stmt:P4.20}
\StatementRecord{night:P4.20}{proved}{02WS}{\SeriesRef{night:P4.16}, \SeriesRef{night:P4.6}, \SeriesRef{night:P4.8}}

Let $mS$ be a fibre with irreducible reduced support $S$ in a proper
holomorphic degeneration of complex two-tori with smooth total threefold.
Suppose the normalization of $S$ is smooth and, outside
finitely many points, $S$ is smooth or has the ordinary-double
germ $xy=0$. Then $S$ has normal crossings everywhere.
Propositions~\ref{night:P4.6}--\ref{night:P4.8} apply to
fibrations in the standing compact setting whose reduced supports
all satisfy these local hypotheses or already have normal crossings.
\end{proposition}
\begin{proof}

\begin{enumerate}[label={$\langle1\rangle$\arabic*.},leftmargin=*]
\item Take a small ball $B$ around an exceptional point, with
boundary transverse to the smooth and double strata away
from that point. Write the local defining function as $a g_0^m$,
where $g_0$ is reduced and $a$ is a unit. Choose a holomorphic
root $h^m=a$ and set $g=h g_0$. The local defining function
is then $g^m$.

\item Consider a small holomorphic deformation $G_s$ of $g$ whose
zero set near $\partial B$ retains its smoothly moving smooth
and ordinary-double strata, transverse to $\partial B$.
There is a uniform $\delta>0$ such that
$G_s|_{\partial B}$ is submersive over $0<|t|<\delta$.
At a boundary point on a double curve, choose coordinates
$a,b,c$ with $G_s=ab$, absorbing a unit in $a$.
Transversality of the double curve to $\partial B$ allows one
real $c$-coordinate to be solved from the boundary equation.
The complex coordinates $a,b$ remain free there, so $ab$ is
submersive when nonzero. At a smooth point use $G_s=a$.
Compactness gives uniform bounds; points outside these
neighbourhoods stay away from zero.

\item For each fixed small $s$ there are finitely many nonzero
critical values in this disc. After shrinking, the analytic
critical set in a neighbourhood of the compact closed ball
has finitely many relevant local components. The function is constant on
each analytic critical component. A component at a nonzero
small value cannot meet $\partial B$, so it is compact;
the maximum principle excludes a positive-dimensional
compact analytic subset of a ball. Thus the relevant critical
points are isolated and finite. Paths avoiding these values
join a small regular value to one fixed regular value $t_0$.
Proper Ehresmann transport on the closed ball, with the
boundary submersion from the preceding step, identifies the fibres along
such paths. For sufficiently small deformation the fibres
$G_s=t_0$ also form a smooth proper family to $s=0$, by
compactness. Hence every compact spherical class of nonzero
square in a local nearby fibre of $G_s$ transports into
$g=t_0$. Its inclusion in the compact torus fibre is impossible.

\item At a normalization preimage write the map germ as
$\phi:(\mathbb C^2,0)\to(\mathbb C^3,0)$.
Away from the exceptional point it is immersive with transverse
multiple images. The isolated-instability criterion in
\cite[p.~5]{GCP2024} makes this monogerm finitely determined.
Write $I_2(d\phi)$ for the ideal of its maximal minors.
A rank drop gives a nonempty finite scheme, and
\cite[Lemma~2.3, p.~5]{GCP2024} identifies its positive length
\[
 C(\phi)=\dim_{\mathbb C}\mathbb C\{u,v\}/I_2(d\phi)
\]
with the number of cross-caps in a stable deformation.
Apply this formula separately at each normalization preimage.
Generic polynomial perturbations vary the branch jets independently.
The stable multigerms listed in \cite[p.~4]{GCP2024} place each
cross-cap on a single sheet. Concretely, the incidence of its
zero-dimensional critical set with another two-dimensional image
in the three-dimensional target has negative expected dimension;
generic relative perturbations avoid it.

\item For the cross-cap coordinates, choose the first target
coordinate at a transverse rank-one point as a source
coordinate $u$, and write $\phi=(u,F(u,v),H(u,v))$.
Transversality of the equations $F_v=H_v=0$ gives a linear
combination with nonzero second $v$-derivative. The
holomorphic Morse lemma with parameter $u$, followed by
subtracting the critical value, makes the second coordinate
$v^2$. Subtract the even part of the third coordinate to
write it as $vA(u,v^2)$. Transversality gives $A_u(0,0)\ne0$.
The target change $u\mapsto A(u,Y)$ and the matching source
change give exactly $(u,v)\mapsto(u,v^2,uv)$.
The target coordinate permutation $(x,y,z)=(uv,u,v^2)$ gives
$x^2-zy^2=0$. Proposition~\ref{night:P4.16}, with $n=2$,
gives a compact spherical class of square $-4$ in its nearby fibre.

\item The deformed images have holomorphic equations $G_s$ with
$G_0=g$. To see this, a generic projection of the original
finite normalization to two target coordinates is finite;
Weierstrass preparation preserves this after shrinking the
deformation. The total image is an analytic hypersurface
in $\mathbb C^3\times\Delta$. Its central equation has
multiplicity one, as seen at a smooth point of $S$; a unit
adjustment gives $G_0=g$. Outside a smaller ball the original
map is a transverse immersion, and this compact boundary
data persist under small deformation. We therefore have
the boundary property above; compare the stable-image discussion
in \cite[\S5.3, p.~33]{GCP2024}. Transporting the
crosscap class contradicts its torus inclusion. The
normalization is consequently immersive everywhere.

\item Write its image germ as a unit times $g_1\cdots g_r$, where
each branch is smooth and $v_i=dg_i(0)\ne0$. Absorb the
unit in one factor. Independent gradients give $r\leq3$
and normal-crossing coordinates by the inverse function
theorem. Suppose instead $\sum_i c_iv_i=0$ is a nonzero
relation, and let $I$ be its nonzero support. Set
$a_i=\varepsilon/c_i$ for $i\in I$ and $a_j=\eta$ for
$j\notin I$, where $0<|\varepsilon|\ll|\eta|\ll1$.
Perturb constants so that $h_i(0)=a_i$. Keep all gradients
except one $i_0\in I$, and change that gradient by
\[
 \Delta v_{i_0}=-a_{i_0}\sum_{j\notin I}\frac{v_j}{a_j}.
\]
This is arbitrarily small and makes $\sum_i dh_i(0)/a_i=0$.
Thus $H=\prod_i h_i$ is critical at the origin and has
nonzero critical value $\prod_i a_i$. Its Hessian divided
by that value is
\[
 \sum_i\frac{\operatorname{Hess}(h_i)}{a_i}
 -\sum_i\frac{dh_i\otimes dh_i}{a_i^2}.
\]
An arbitrarily small generic quadratic change to $h_{i_0}$
makes this matrix invertible: its determinant is a nonzero
polynomial in the freely variable symmetric Hessian.
The branch functions remain smooth and retain their
ordinary-double/smooth boundary strata. We have constructed
a holomorphic Morse point off the divisor.

\item Its nearby fibre contains the real vanishing sphere in
$w_1^2+w_2^2+w_3^2=\epsilon$. Its normal bundle is
$iTS^2$, with the opposite orientation to $TS^2$ under the
complex surface orientation. Its self-intersection is
therefore $-e(S^2)=-2$. Boundary transport again puts this
class in the original torus fibre, a contradiction. Hence
the original branch gradients are independent everywhere,
which proves the normal-crossing assertion.

\end{enumerate}
\end{proof}


\paragraph{Exact source locations used here.}
The classification inputs are \cite[Chapter~VI, Theorem~1.1 and
Table~10, pp.~243--244]{BHPV2004E},
\cite[Theorems~4.2 and~5.1, with proofs, pp.~17--21 and~27--29]{OVV2020},
\cite[Theorem~1, p.~1]{BBK2019}, and
\cite[Theorem~2.1]{Catanese2004}.
The no-curve theorem is credited to Inoue by
\cite[Theorem~3.1, p.~303]{KhaledTeleman2026Journal}; the actual
unramified double cover is in its proof of Remark~3.5, p.~305.
The exceptional intersection identity is
\cite[\S I, p.~6]{Mumford1961}; the proof above supplies its
local analytic function and its exact exceptional-divisor map.
The current source records distinguish theorem-statement openings
from complete source-proof readings. The sign and equality assertions
are proved through the displayed normalization, conductor, filling
and intersection maps, with the stated irreducible-or-NC scope for
the global conclusion.

\subsection{The remaining smooth ruled normalizations}

\begin{proposition}[Positive Euler number for rational normalization]\hypertarget{tag-02WT}{}\label{night:P4.21}\label{stmt:P4.21}
\StatementRecord{night:P4.21}{proved}{02WT}{\SeriesRef{cdp:st:A.5.3}, \SeriesRef{night:P4.12}, \SeriesRef{night:P4.13}, \SeriesRef{night:P4.20}, \SeriesRef{night:P4.6}, \SeriesRef{night:P4R.1}}

Let $f:Y\to B$ be a proper surjective holomorphic map with connected
fibres from a smooth compact complex threefold with $b_1(Y)=b_2(Y)=0$
to a smooth compact curve, with smooth general fibre a complex two-torus.
Let $mS$ be a fibre with irreducible singular reduced support, and let
$\nu:T\to S$ be its smooth rational normalization. Then $e(S)>0$.
\end{proposition}
\begin{proof}

\begin{enumerate}[label={$\langle1\rangle$\arabic*.},leftmargin=*]
\item Let $D$ be the conductor divisor on $T$, with its scheme
multiplicities, and put $C=\operatorname{Sing}(S)_{\mathrm{red}}$.
Propositions~\ref{cdp:st:A.5.3} and~\ref{night:P4.12} give
$\omega_T(D)\simeq\nu^*\omega_S$, a torsion line bundle.
The minimal resolution in the first proposition is the identity of the
smooth surface $T$. Pullback to a common point-blowup resolution of
$T$ and $\mathbb P^2$ injects $\operatorname{Pic}(T)$ into a free
abelian Picard group; see \cite[II.3, pp.~17--18; II.12, p.~24]{Beauville1978}.
Thus $D\sim-K_T$. Contract $T$ to a minimal rational surface $T_0$,
which is $\mathbb P^2$ or a Hirzebruch surface $\mathbb F_n$ with
$n\ne1$ \cite[V.10, pp.~79--80]{Beauville1978}, and push $D$ forward
to $D_0$. For each reversed point blowup $q$ with exceptional curve $E$,
the canonical formula $K_{\mathrm{new}}=q^*K_{\mathrm{old}}+E$
gives $D_{\mathrm{new}}=q^*D_{\mathrm{old}}-E$.
Its exceptional coefficient is $\mu-1$, where $\mu$ is the
conductor multiplicity at the centre. Effectivity therefore puts
every centre on the conductor.

\item On $\mathbb P^2$, the divisor $D_0$ is a plane cubic. An
irreducible cubic is smooth elliptic, rational nodal, or
rational cuspidal. A reducible cubic is a smooth conic
and a line, three lines, a double line and another line,
or a triple line. The conic and line meet at two
transverse points or at one tangency; three distinct
lines form a triangle or are concurrent. Thus either
$D_0$ is reduced nodal and has fundamental group
$\mathbb Z$ or $\mathbb Z^2$, or its reduced support is
simply connected and all normalized components are
rational.

\item On $\mathbb F_n$, write $h$ for the tautological class
and $F$ for a ruling fibre. The formulas
$h^2=n$, $h\cdot F=1$, $F^2=0$, and
$K_{\mathbb F_n}=-2h+(n-2)F$ are
\cite[III.18, pp.~46--47; IV.1, pp.~53--54]{Beauville1978}.
Thus $D_0\cdot F=2$, and its horizontal part is a double
section, two distinct reduced sections, or a reduced
irreducible bisection. A double section together with
vertical fibres is a rational tree. For the other cases
write
\[
 D_0=H_1+H_2+\sum_j a_jF_j
 \quad\hbox{or}\quad
 D_0=H+\sum_j a_jF_j.
\]
In the first case adjunction on either section gives
\[
 H_1\cdot H_2+\sum_j a_j=2.
\]
The support is therefore a rational tree or a single
rational cycle. A cycle occurs only with all coefficients
one and ordinary double intersection points. This
includes the cases of two section-intersection points,
one section-intersection point and one vertical fibre,
or two vertical fibres. Coincident intersections or a
vertical fibre through the section-intersection point
instead give trees.

\item For a bisection, adjunction gives
\[
 p_a(H)=1-\sum_j a_j.
\]
There is either no vertical component and $H$ has
arithmetic genus one, or one coefficient-one vertical
fibre and a smooth rational bisection. In the former
case $H$ is smooth elliptic, rational nodal, or rational
cuspidal. In the latter it meets the fibre at two distinct
points or at one tangency. Again the only nontrivial
fundamental groups are $\mathbb Z$ and $\mathbb Z^2$,
and those cases have reduced nodal conductor.

\item We follow these alternatives through the blowups. At a
point of conductor multiplicity at least two, the new
conductor contains the exceptional $\mathbb P^1$.
Its contraction is the quotient which collapses that
connected simply connected subcomplex to the old
point. Van Kampen therefore makes the induced map on
fundamental groups an isomorphism. At a point of
multiplicity one the conductor is smooth and reduced,
and its strict transform is isomorphic to it. Thus
$\pi_1(D_{\mathrm{red}})$ is unchanged throughout.
If the initial group was nontrivial, the conductor was
reduced nodal. A smooth-point blowup adds no conductor
component, and a node blowup inserts one component of
coefficient one. It remains reduced nodal.

\item In this latter case the generic transverse conductor
exponent of every normalization branch is one. The
exponent is its intrinsic conductor exponent plus the
intersection multiplicities with the other branches.
A singular branch has intrinsic exponent at least two.
Exponent one therefore means two smooth branches
meeting transversely. Consequently $S$ has ordinary double curves outside finitely
many points. Proposition~\ref{night:P4.20} makes $S$ normal crossing
everywhere. Write $t$ for its number of ordinary triple points.
On each normal-crossing chart the composite $T\to S\subset Y$
is the inclusion of a smooth coordinate plane, hence is immersive.
The exact sequence
\[
0\longrightarrow T_T\longrightarrow\nu^*(T_Y|_S)
 \longrightarrow L\longrightarrow0
\]
has line quotient. Since $H^2(Y,\mathbb Q)=H^4(Y,\mathbb Q)=0$,
the pulled-back first and second Chern classes vanish. Thus
$c_1(L)=K_T$ and $e(T)=K_T^2$. Rationality gives
$\chi(\OO_T)=1$ \cite[III.20--21, pp.~48--50]{Beauville1978};
Noether's formula therefore gives $e(T)=6$.
The reduced nodal anticanonical conductor is either smooth elliptic
or a cycle of rational curves. Adjunction on a rational component
gives exactly two node preimages. Its number of components therefore
equals its number of nodes, and its Euler number equals that number;
the smooth elliptic case has Euler number zero and no nodes. Every triple point
of $S$ has three node preimages on $D$, so $e(D)=3t$.
Over the smooth part of $C$ the conductor has two points per fibre,
whereas a triple point has three. Hence
\[
e(D)=2e(C\setminus\{\text{triple points}\})+3t
     =2e(C)+t,
\]
and $e(C)=t$. Additivity gives $e(S)=6-2t$.
Proposition~\ref{night:P4.6}, applied to this reduced divisor $S$,
gives $e(S)=t$. Therefore $t=2$ and $e(S)=2$.

\item In the remaining case $D_{\mathrm{red}}$ is connected
and simply connected, and all its normalized components
are rational. Triangulate the analytic pair $(T,D_{\mathrm{red}})$ and the finite
map $D_{\mathrm{red}}\to C$. The inclusion is a cofibration, so
its pushout has the fundamental group of the corresponding mapping
cylinder. Van Kampen gives
\[
 \pi_1(S)=\pi_1(T)*_{\pi_1(D_{\mathrm{red}})}\pi_1(C)
          =\pi_1(C).
\]
Every component of $C$ has rational normalization,
because it is the finite nonconstant image of a
rational component of $D_{\mathrm{red}}$. Proposition~\ref{night:P4R.1} makes
$\pi_1(C)$ virtually abelian. For a connected curve
with rational normalized components, this is the free
group $F_b$ of its point-identification graph, so
$b\le1$. If $c\ge1$ is its number of components, then
\[
 e(C)=c+1-b\ge c\ge1.
\]
Proposition~\ref{night:P4.13} gives $e(S_{\mathrm{reg}})\ge0$. Additivity proves
$e(S)=e(S_{\mathrm{reg}})+e(C)>0$.

\end{enumerate}
\end{proof}

\begin{proposition}[Euler number of an elliptic-ruled normalization]\hypertarget{tag-02WU}{}\label{night:P4.22}\label{stmt:P4.22}
\StatementRecord{night:P4.22}{proved}{02WU}{\SeriesRef{cdp:st:A.5.3}, \SeriesRef{night:P4.12}, \SeriesRef{night:P4.13}, \SeriesRef{night:P4.15}, \SeriesRef{night:P4R.1}, \SeriesRef{night:P4R.3}}

In the standing torus-fibration setting of this section, an irreducible singular support
$S$ whose smooth normalization $T$ is birationally ruled
over an elliptic curve has $e(S)\ge0$. Equality holds
exactly when $T$ is minimal.
\end{proposition}
\begin{proof}

\begin{enumerate}[label={$\langle1\rangle$\arabic*.},leftmargin=*]
\item Write $\nu:T\to S$ for normalization, $D$ for its full
conductor divisor, and $C=\operatorname{Sing}(S)_{\mathrm{red}}$.
Propositions~\ref{cdp:st:A.5.3} and~\ref{night:P4.12}, with the
identity as the resolution of $T$, give the torsion line
$\omega_T(D)\simeq\nu^*\omega_S$.
Push $D$ to the minimal ruled surface $T_0\to E$, where $E$ is
elliptic. The projective-bundle and canonical formulas are
\cite[III.10, p.~40; III.18, pp.~46--47]{Beauville1978}.
They give $K_{T_0}+D_0\equiv0$ and $D_0\cdot F=2$ for a fibre $F$. If a horizontal component $H$
has coefficient one, adjunction gives $p_a(H)\le1$,
whereas its nonconstant map to $E$ gives geometric
genus at least one. Thus $H$ is smooth elliptic and
$(D_0-H)\cdot H=0$. There is no vertical component.
The horizontal part is one elliptic bisection or two
disjoint elliptic sections, all with coefficient one.
All subsequent blowups have smooth coefficient-one
conductor centres; their exceptional curves are absent
from the conductor. The coefficient-one branch conductors give two smooth
transverse branches at a general point of each image curve:
the sum of the branch exponents is $2\delta$, while
$\delta\ge q-1$ for $q$ branches, so $q=2$ and $\delta=1$.
Thus the normalization differential has no divisorial rank drop;
its analytic rank-drop locus is finite.
Proposition~\ref{night:P4.15} makes $T$ minimal,
makes the differential immersive everywhere, and
gives $e(S)=0$.

\item The other possibility is
\[
 D_0=2H+\sum_j a_jF_j,
\]
where $H$ is a smooth elliptic section. The canonical blowup
formula in \cite[II.3, pp.~17--18]{Beauville1978}, together with
$K+D\equiv0$, gives conductor coefficient $\mu-1$ on the exceptional
curve. These blowups make $D_{\mathrm{red}}$ a copy of $H$ with rational trees attached:
a smooth-point blowup adds a leaf or no conductor component,
and a node blowup subdivides an edge. Collapsing those rational
trees induces an isomorphism to $\pi_1(H)$. The section
$H\to T_0$ is inverse on fundamental groups to $T_0\to E$,
and every point blowup preserves this inclusion map. Hence
$\pi_1(D_{\mathrm{red}})\to\pi_1(T)$ is an isomorphism.
Apply the finite-normalization pushout and curve-graph calculation
in the proof of Proposition~\ref{night:P4R.3} to the inclusion
$D_{\mathrm{red}}\subset T$ and the finite map
$D_{\mathrm{red}}\to C$. Their pushout is $S$, so
$\pi_1(S)\simeq\pi_1(C)$.
Proposition~\ref{night:P4R.1} makes this group virtually abelian.
Every rational conductor component has rationally normalized image;
only $H$ can dominate an elliptically normalized component.
The curve-graph formula gives $e(C)\ge0$.
Proposition~\ref{night:P4.13} and additivity therefore give $e(S)\ge0$.

\item If $e(S)=0$, the curve calculation forces $C$ to have
one component, genus-one smooth normalization, and
no point-identification graph cycles. Each conductor
component maps finitely and nonconstantly onto $C$.
Riemann--Hurwitz excludes rational conductor
components. Thus $D_{\mathrm{red}}=H$ and no blowup
occurred: the first blowup on $2H$ would create a
vertical conductor component of coefficient one.
Therefore $T$ is minimal.

\item Conversely suppose $T$ is minimal. In the double-section
case let $v$ be the number of distinct vertical fibres
in the conductor. Since $e(T)=e(H)=0$ and each vertical
$\mathbb P^1$ meets $H$ once,
\[
 e(S_{\mathrm{reg}})=e(T\setminus D_{\mathrm{red}})=-v.
\]
Proposition~\ref{night:P4.13} forces $v=0$, hence $D=2H$. At a general point
of its image curve, every normalization branch has
conductor exponent two. The branch exponent equals
its intrinsic conductor exponent plus the sum of its
intersections with the other branches. A singular
branch together with another branch would have
exponent greater than two. Thus either there is one
ordinary cusp branch, or all branches are smooth:
two with intersection multiplicity two, or three
with pairwise intersection multiplicity one. Four
smooth branches would have exponent at least three,
so the list is exhaustive.

\item In the cusp case $H\to C$ has generic degree one.
The normalization of $C$ is therefore $H$, an elliptic
curve. Virtual abelianness of $\pi_1(C)$ excludes any
additional point-identification graph cycle, so
$e(C)=0$. Additivity gives $e(S)=0$. In the other two
cases the normalization differential is immersive at
the generic point of $H$. There is no other conductor
component, so its rank-drop set is finite. Proposition~\ref{night:P4.15}
applies and gives $e(S)=0$.

\end{enumerate}
\end{proof}


\subsection{Class-VII minimal models with positive second Betti number}

\begin{proposition}[P4.23]\hypertarget{tag-02WV}{}\label{night:P4.23}\label{stmt:P4.23}
\StatementRecord{night:P4.23}{proved}{02WV}{\SeriesRef{night:P4.12}, \SeriesRef{night:P4.13}, \SeriesRef{night:P4.20}, \SeriesRef{night:P4.21}, \SeriesRef{night:P4.6}, \SeriesRef{night:P4R.1}}

In the standing torus-fibration setting of this section, let an irreducible singular support
$S$ have smooth normalization $T$. If its minimal model $T_0$
belongs to class VII and $b_2(T_0)>0$, then $e(S)>0$.
\end{proposition}
\begin{proof}

\begin{enumerate}[label={$\langle1\rangle$\arabic*.},leftmargin=*]
\item Let $D$ be the conductor. P4.12 makes $K_T+D$ torsion;
pushing down the point blowups gives $K_{T_0}+D_0$ torsion,
with $D_0$ effective and nonzero. This is a numerically
anticanonical divisor in the sense of Dloussky,
Definition~2.20 of \cite{Dloussky2006}. His Theorem~0.2
therefore gives a global spherical shell. Proposition~2.27
and the square-zero-divisor paragraph after Lemma~2.22
leave the parabolic Inoue or Inoue--Hirzebruch cases,
or a connected rational cycle with attached rational
trees. The latter configuration and $\pi_1(T_0)=\mathbb Z$
are stated explicitly in \cite[\S1, pp.~284--285]{DOT2003}.
We verify the consequences needed
for our conductor.

\item For an Inoue--Hirzebruch surface, the reduced cycle or
union of two cycles is numerically anticanonical. The
curve intersection matrix is negative definite. Subtracting
that divisor from $D_0$ and intersecting with every curve
therefore makes every coefficient difference zero.
Thus $D_0$ is reduced nodal. In the parabolic Inoue case
there is an elliptic curve $E$ and a disjoint rational
cycle $\Gamma$, with
\[
 K_{T_0}+E+\Gamma\sim0,\qquad \Gamma\equiv0.
\]
These relations are stated after Dloussky's Lemma~2.22.
Intersections with $E$ and the rational curves give
$D_0=E+k\Gamma$: the cycle matrix has kernel generated
by the sum of its components. Torsion of $K_{T_0}+D_0$
then makes $\mathcal O((k-1)\Gamma)$ torsion. A nonzero
effective divisor cannot have torsion associated line
bundle on a compact connected reduced complex space:
a power of its defining section would be a nonzero
constant and could have no zero. Hence $k=1$ and
$D_0$ is again reduced nodal.

\item The blowup coefficient $\mu-1$ preserves the reduced
nodal property, as in P4.21. Thus $S$ has ordinary
double curves outside finitely many points. P4.20
makes it normal crossing everywhere. Its normalization
is then immersive and P4.6 gives $e(T)=K_T^2$.
For a class-VII surface, however,
\[
 e(T)=b_2(T)>0,\qquad K_T^2=-b_2(T),
\]
a contradiction. These two exceptional minimal-model
cases are excluded.

\item In the remaining case the support of the conductor is
a connected rational cycle with attached rational trees.
It is the whole connected curve configuration: Dloussky's
Lemma~2.22(2) puts every curve meeting the anticanonical
support in that support, and Proposition~2.27 gives its
connectedness. Its fundamental group is $\mathbb Z$.
The blowups preserve this group by the quotient argument
of P4.21, and preserve the surface fundamental group.
Consequently
\[
 \pi_1(D_{\mathrm{red}})=\mathbb Z,\qquad
 \pi_1(T)=\mathbb Z.
\]
Every normalized component of
$C=\operatorname{Sing}(S)_{\mathrm{red}}$ is rational,
being the finite nonconstant image of a rational conductor
component. The curve is connected. Write $\pi_1(C)=F_b$
for the free group of its point-identification graph.
The conductor pushout gives
\[
 \pi_1(S)=
 \langle t,c_1,\ldots,c_b\mid t^k=w(c_1,\ldots,c_b)\rangle.
\]
Here $k$ records the conductor-cycle map to $\pi_1(T)$,
and $w$ records its map to $\pi_1(C)$. Neither map is
assumed to be an isomorphism.

\item Virtual abelianness from P4R.1 forces $b\le1$.
If $k\ne0$ and $w\ne1$, both cyclic amalgam maps are
injective; alternating nonidentity coset representatives
give the reduced normal form and show that $F_b$ injects
in the amalgam. If $w=1$, the group is
$(\mathbb Z/k\mathbb Z)*F_b$, with
$\mathbb Z/0\mathbb Z=\mathbb Z$, and $F_b$ again
injects. A nonabelian free subgroup is impossible in a
virtually abelian group. Finally, if $k=0$ and $w\ne1$,
the group is
\[
 \mathbb Z*\bigl(F_b/\langle\!\langle w\rangle\!\rangle\bigr).
\]
A free product of an infinite cyclic group and a nontrivial
group contains a nonabelian free subgroup: for a nontrivial
$g$ in the second factor, $t$ and $gtg^{-1}$ have no
nonempty reduced relation. Thus the second factor must
be trivial. Its abelianization has rank at least $b-1$,
which again gives $b\le1$.

\item For the number $c\ge1$ of components of $C$, we obtain
\[
 e(C)=c+1-b\ge c\ge1.
\]
By P4.13, $e(S_{\mathrm{reg}})\ge0$, so
$e(S)=e(S_{\mathrm{reg}})+e(C)>0$.

\end{enumerate}
\end{proof}

\subsection{Minimality of a smooth Hopf normalization}

\begin{proposition}[P4.24]\hypertarget{tag-02WW}{}\label{night:P4.24}\label{stmt:P4.24}
\StatementRecord{night:P4.24}{proved}{02WW}{the definitions and constructions stated in the proof}

In the standing torus-fibration setting of this section, let $S$ be an irreducible
singular fibre support with smooth normalization $T$.
If the minimal model of $T$ is a Hopf surface, then
$T$ is already minimal. The reduced conductor on $T$
and its reduced image in $S$ are disjoint unions of
smooth elliptic curves. Each map between their
components is a finite unramified covering, and
$e(S)=0$.
\end{proposition}
\begin{proof}

\begin{enumerate}[label={$\langle1\rangle$\arabic*.},leftmargin=*]
\item Write $\nu:T\to S\hookrightarrow Y$ for the finite
holomorphic normalization map followed by the closed
embedding. We use the effective integral conductor
divisor $D$ with $K_T+D$ torsion. Let $T\to T_0$
be the sequence of $k$ point blowups giving the
minimal Hopf model. At each step the conductor
coefficient on the exceptional curve is $\mu-1$,
where $\mu$ is the multiplicity of the conductor
at the centre. Effectiveness therefore puts every
centre on the current conductor.

\item There is a finite unramified primary Hopf cover
$T'_0\to T_0$. Pulling back every blowup gives a
finite unramified cover $T'\to T$ and a blowup
presentation of $T'$ over $T'_0$. A centre has
exactly the degree of the covering many lifts,
so the new number of blowups is that degree times
$k$. The composite $T'\to Y$ is finite and is
immersive away from the pulled-back conductor.
The rational Chern classes of its pulled-back
ambient tangent bundle vanish. Indeed $b_2(Y)=0$
gives $c_1(T_Y)=c_2(T_Y)=0$ in rational cohomology,
using Poincar\'e duality in degree four. The
argument below applies to this finite map, without
requiring it to be a normalization. Proving
minimality of $T'$ proves $k=0$. We therefore use
$T,T_0,\nu,k$ for the covering objects during the
Chern calculation.

\item We first construct on $T_0$ a nonsingular
holomorphic foliation tangent to every compact
curve. Khaled, \emph{Real structures on primary
Hopf surfaces}, Theorem~2.1(1) and formula~(1),
gives
\[
 T_0=(\mathbb C^2\setminus\{0\})/\langle F\rangle,
 \qquad F(z,w)=(\alpha z+\lambda w^n,\beta w),
\]
where $0<|\alpha|\le|\beta|<1$ and
$\lambda(\alpha-\beta^n)=0$. We verify the needed
curve statement from these coordinates. Every
compact irreducible curve is smooth elliptic:
adjunction and $b_2(T_0)=0$ give arithmetic genus
one, and a rational normalization would lift to
$\mathbb C^2\setminus\{0\}$ and have constant
coordinate functions.

\item The map from the fundamental group of such an
elliptic curve to $\pi_1(T_0)=\mathbb Z$ has
nonzero image, since zero image would give the
same impossible compact lift. Its full inverse
image has finitely many connected components.
The Remmert--Stein extension across the origin
is an analytic curve germ; a power $F^h$ fixes
one of its branches. On the normalization of
that branch the induced contraction can be
written $t\mapsto\eta t$, $0<|\eta|<1$.
For completeness, if the initial coordinate
has return map $g(t)=\eta t+O(t^2)$, choose
$|\eta|<\rho<\sqrt{|\eta|}$. On a sufficiently
small disc, $|g(t)|\le\rho|t|$, and the
successive differences in $g^j(t)/\eta^j$
are bounded by a constant times
$(\rho^2/|\eta|)^j|t|^2$. Their convergent sum
gives the asserted linearizing coordinate.

\item When $\lambda=0$, comparison of coefficients in
$z(\eta t)=\alpha^h z(t)$ and
$w(\eta t)=\beta^h w(t)$ shows that each
nonzero coordinate is a single monomial.
Thus every branch away from the coordinate axes
has parametrization $(a t^p,b t^q)$. Its ratio
$p/q=\log|\alpha|/\log|\beta|$ is independent
of the branch. If such a branch exists, take
coprime positive integers $p,q$ in this ratio
and the vector field
$p z\partial_z+q w\partial_w$.
If no such branch exists, use
$z\partial_z+w\partial_w$.
Both are invariant under $F$, have no zero on
$\mathbb C^2\setminus\{0\}$, and are tangent
to every compact curve.

\item When $\lambda\ne0$, we have $\alpha=\beta^n$ and
\[
 F^h(z,w)=
 (\alpha^h z+h\lambda\alpha^{h-1}w^n,\beta^h w).
\]
If a normalized invariant branch had
$w(t)=b t^q\ne0$, the coefficient of $t^{nq}$
in the first coordinate identity would give
$0=h\lambda\alpha^{h-1}b^n$, a contradiction.
Thus the only curve is the image of $w=0$.
The vector field $n z\partial_z+w\partial_w$
is invariant, nowhere zero, and tangent to it.
This proves the foliation assertion.

\item Pull this foliation through the point blowups
and saturate its tangent sheaf. It is a line
bundle $L\subset T_T$. Every strict transform
of a conductor component and every exceptional
curve is invariant. The only singularities of
the foliation are simple saddles. To verify
this assertion and its multiplicities, at a
regular point use coordinates $(x,y)$ with
leaves $y=\text{constant}$. In the chart
$x=u,y=uv$ the saturated generator is
$u\partial_u-v\partial_v$; there is one
saddle, and the other blowup chart is regular
away from it. At an existing saddle use
$b x\partial_x-a y\partial_y$, with positive
integers $a,b$. The two charts have generators
\[
 b u\partial_u-(a+b)v\partial_v,\qquad
 (a+b)u\partial_u-a v\partial_v.
\]
One saddle is replaced by two. Each has
zero ideal $(u,v)$ and length one.
Consequently there are exactly $k$ saddles.

\item Let $V$ be the divisor measuring the
divisorial zero of the pulled-back transverse
one-form before saturation. Numerically
\[
 c_1(L)=-K_T+V.
\]
The omitted pullback line bundle is topologically
trivial because it comes from $T_0$.
The preceding charts give, at every blowup $\mu$,
\[
 K'=\mu^*K+E,\qquad
 V'=\mu^*V+\varepsilon E,\qquad
 \varepsilon=
 \begin{cases}0&\text{at a regular point},\\
 1&\text{at a saddle}.
 \end{cases}
\]
Starting with $V=0$ and $K_{T_0}^2=0$, we obtain
\[
 K_T^2=-k,\qquad K_TV-V^2=0.
\]
The second identity is preserved because its
change is $\varepsilon^2-\varepsilon=0$.

\item Put $B=\nu^*T_Y$ and compose the holomorphic
maps $L\to T_T\to B$. This map has no
divisorial zero. Away from the conductor this
follows from immersivity. At the generic point
of a conductor component the foliation is
tangent to that curve, and finiteness of
$\nu$ makes its differential along the curve
nonzero. A nonconstant holomorphic map of
curves has nonzero derivative at the generic
point. Thus the cokernel $Q=B/L$ is a
torsionfree sheaf of rank two. Its double
dual $Q^{**}$ is locally free, and
\[
 z_0=\operatorname{length}(Q^{**}/Q)
\]
is finite. At a point, write the map $L\to B$
as a column $(f_1,f_2,f_3)$ in trivializations.
If this column vanishes, the ideal
$(f_1,f_2,f_3)$ has finite colength. Dualizing
$0\to\mathcal O\to\mathcal O^3\to Q\to0$
gives $\mathcal E xt^1(Q,\mathcal O)
=\mathcal O/(f_1,f_2,f_3)$.
Dualizing $0\to Q\to Q^{**}\to Q^{**}/Q\to0$
and using the two-dimensional regular local
ring identifies this with
$\mathcal E xt^2(Q^{**}/Q,\mathcal O)$,
which has the same length. At each foliation
saddle all three coefficients vanish.
It follows that
\[
 z_0\ge k.
\]

\item Let $E_B\subset B$ be the saturation of the
generic rank-two image of $d\nu:T_T\to B$.
Its determinant differs from that of $T_T$
by an effective integral divisor $R$, the
divisorial rank-drop divisor:
\[
 c_1(E_B)=-K_T+R.
\]
The support of $R$ is contained in the
conductor. Saturate $E_B/L$ in $Q^{**}$.
The resulting rank-one reflexive subsheaf
is a line bundle $J$, with
$c_1(J)=R-V$. Its torsionfree rank-one
quotient has the form
$\mathcal I_{Z_1}\otimes J_1$ with
$c_1(J_1)=K_T-R$ and
$z_1=\operatorname{length}(Z_1)\ge0$.
All these assertions concern sheaves on the
smooth surface $T$; determinants are computed
off finite sets and therefore determine their
line-bundle classes.

\item The identities $c_1(B)=c_2(B)=0$ give
\[
 c_2(Q)=(-K_T+V)^2,\qquad
 c_2(Q^{**})=(-K_T+V)^2-z_0.
\]
The exact sequence with $J$ gives another
expression for the same number:
\[
 (-K_T+V)^2-z_0=(R-V)(K_T-R)+z_1.
\]
Using $K_T^2=-k$ and $K_TV-V^2=0$, this becomes
\[
 G(R):=R^2-K_TR-VR=k+z_0+z_1\ge2k.
\]

\item We compare this expression with the actual
blowup sequence. Push $R$ and $V$ down at
each step; $R$ has integral coefficients.
Writing $R'=\mu^*R+cE$ gives $c\in\mathbb Z$,
without any positivity requirement on $c$.
The intersection calculation is
\[
 G(R')-G(R)=-c^2+(1+\varepsilon)c.
\]
For $\varepsilon=0$ this is at most zero
for every integer $c$; for $\varepsilon=1$
it is $1-(c-1)^2\le1$. On $T_0$ every
divisor is numerically zero, so $G(R_0)=0$.
Consequently
\[
 G(R)\le\#\{\text{blowups at saddles}\}\le k.
\]
The lower bound $G(R)\ge2k$ forces
$k=z_0=z_1=0$. Returning through the finite
unramified cover proves that the original
normalization $T$ is minimal.

\item It remains to identify the conductor maps,
rather than only their Euler number.
The reduced conductor on a minimal Hopf
surface consists of disjoint smooth elliptic
curves. Since $z_0=0$ on the primary cover,
the differential of $\nu$ along each such
curve is nowhere zero, also before that
cover. Thus every branch of its image is
smooth. At an image point, two different
curve germs would have to come from two
different normalization points. Their
surface branches in the smooth threefold
intersect in a curve germ. This germ is
contained in the conductor image of both
branches; the unique conductor germ at
each normalization point forces the two
image germs to be equal. Therefore the
image conductor is locally a single smooth
curve. Its components are disjoint, and
each conductor map is finite unramified.
Riemann--Hurwitz makes every image component
elliptic. Finally, normalization is an
isomorphism off these curves, so
\[
 e(S)=e(T)-e(D_{\mathrm{red}})
                 +e(\operatorname{Sing}S)=0.
\]

\end{enumerate}
\end{proof}


\subsection{The reduced resolved-conductor case}

\begin{proposition}[Semi-log-canonical fibre supports]\hypertarget{tag-02WX}{}\label{night:P4.25}\label{stmt:P4.25}
\StatementRecord{night:P4.25}{proved}{02WX}{\SeriesRef{night:P4.16}, \SeriesRef{night:P4.20}, \SeriesRef{night:P4.6}, \SeriesRef{night:P4.8}}

Let $f:Y\to B$ be a proper surjective holomorphic map with connected
fibres from a connected smooth compact complex threefold to a smooth
compact curve. Assume $b_1(Y)=b_2(Y)=0$ and that the smooth general
fibre is a complex two-torus. Every semi-log-canonical reduced fibre
support of $f$ has normal crossings.

For an irreducible reduced support $S$, let $\nu:N\to S$ be its
normalization, let $\pi:T\to N$ be a resolution, and put
$\rho=\nu\circ\pi$. For local residue generators $\Omega$ of $\omega_S$,
define the integral divisor $D=-\operatorname{div}(\rho^*\Omega)$, so that
$\omega_T(D)\simeq\rho^*\omega_S$ with this identification.
A reduced normal-crossing divisor $D$ gives normal crossings on $S$
and a smooth normalization $N$. Proposition~\ref{night:P4.6} applies
to each such support. Propositions~\ref{night:P4.7}--\ref{night:P4.8}
apply to fibrations whose reduced supports are all semi-log-canonical.
\end{proposition}
\begin{proof}

\begin{enumerate}[label={$\langle1\rangle$\arabic*.},leftmargin=*]
\item At a generic point of the conductor, an exponent one on every normalization
branch forces the transverse plane curve to be an
ordinary node. Indeed the conductor exponent on
branch $i$ is
$2\delta_i+\sum_{j\ne i}I_{ij}$; positivity of
intersection multiplicities leaves exactly two
smooth transverse branches. The support $S$ is a
Cartier hypersurface, hence is Cohen--Macaulay and
has invertible dualizing sheaf. On its normalization $N$, the pair consisting of $N$
and its reduced conductor is log canonical.
To check this last assertion, pull its canonical
divisor to $T$. The resulting boundary is exactly
$D$, by the displayed isomorphism and the
divisorial orders of a local generator. A smooth
surface with reduced normal-crossing boundary is
log canonical: blowing up a point gives new
boundary coefficient $\mu-1\le1$, and this
bound persists under successive point blowups.
Discrepancy comparison therefore gives the
assertion on $N$. Definition~2.3 and Lemma~2.2 of
\cite[p.~189]{LiuRollenske2012SLC} give semi-log-canonicity in
the complex analytic category.

\item The analytic classification is
\cite[Theorem~1.1, p.~186; Table~1, p.~187;
Propositions~3.1 and~4.1, pp.~192--200]{LiuRollenske2012SLC}.
An isomorphism of the singular hypersurface germs lifts to an ambient
biholomorphism: lift its three coordinate classes to convergent functions;
the induced cotangent isomorphism makes their linear part invertible.
Thus, in the displayed coordinates, the full local fibre function is
$U\prod_i g_i^{m_i}$, with $U$ a unit and the original component
multiplicities $m_i$. The normal singular germs in the table are isolated.
Each is irreducible, since distinct surface branches would meet along
a singular curve. Write its full function as $Ug^m$, choose a holomorphic
$m$th root of $U$, and set $h=U^{1/m}g$, so that $f=h^m$. Any such germ is excluded
in our family: a small generic linear
perturbation $h-c\cdot z$ has Morse critical points.
Their number is the positive multiplicity of the gradient map
\cite[Appendix~B, Lemmas~B.1--B.3 and the positivity proof,
pp.~111--114]{Milnor1968}. It is also the colength of the gradient ideal:
the map $\mathbb C\{u\}\to\mathbb C\{z\}$,
$u_i\mapsto h_{z_i}$, is finite; its three parameters form a regular
sequence, so its target is a free module of that colength. A regular
nearby fibre consists of that many reduced points, each of local degree
one. The boundary degree is unchanged under the small perturbation. A constant perturbation
puts their critical values away from zero.
The original zero divisor is smooth near
a small transverse boundary sphere, so it
retains that boundary property under these
perturbations. The proper boundary
transport in the proof of Proposition~\ref{night:P4.20} carries a
vanishing sphere of square $-2$ into a
nearby torus fibre, which is impossible.

\item The ordinary double point and $xyz=0$ have normal crossings.
The change $z\mapsto-z$ identifies the pinch point $x^2+y^2z=0$
with the case $n=2$ of Proposition~\ref{night:P4.16}.
For a weighted homogeneous full function $G$ of positive weights
$(w_1,w_2,w_3)$ and degree $W$, the unit is absorbed by the
biholomorphism
\[
 (z_1,z_2,z_3)\longmapsto
 (U(z)^{w_1/W}z_1,U(z)^{w_2/W}z_2,U(z)^{w_3/W}z_3),
\]
using one local logarithm of $U$; it satisfies $G\circ\Psi=UG$.
The weights used below are positive. Scaling the coordinates by
$\lambda^{w_i}$ places every compact cycle in $G^{-1}(1)$ in
$G^{-1}(\lambda^W)$ inside an arbitrarily small neighborhood.
For a deformation $G_s$ of an irreducible reduced equation with full
multiplicity $m$, use the full map $G_s^m$. At nonzero values its
levels are the disjoint union of the levels $G_s=\zeta t^{1/m}$,
where $\zeta^m=1$, and
$d(G_s^m)=mG_s^{m-1}dG_s$. Thus the boundary transport on these
levels carries the original full multiplicity.
The remaining double-point equations are
\[
 x^2+y^2z^2,\qquad
 x^2+y^2(z^2+y^{q-2}),\quad q\ge3.
\]
For the first equation the two normalization
branches are the smooth graphs $x=\pm i yz$.
Their tangent planes coincide at the origin,
whereas they meet transversely off that point
along the two axes. Put $U_1=x-iyz$ and $U_2=x+iyz$, and let their full fibre
multiplicities be $a,b>0$. For $d=\gcd(a,b)$ write
$a'=a/d$, $b'=b/d$. On the component $U_1^{a'}U_2^{b'}=1$
of $U_1^aU_2^b=1$, set
\[
 U_1=t^{b'},\qquad U_2=t^{-a'},\qquad
 yz=\frac{t^{-a'}-t^{b'}}{2i}.
\]
The inverse parameter is $t=U_1^rU_2^s$ for integers
$rb'-sa'=1$. The last equation has $a'+b'\ge2$ simple nonzero
zeros in its right-hand side. Over an arc joining two of them,
the circles $|y|=|z|$ collapse at the endpoints and form a smooth
embedded sphere. The circle action $(y,z)\mapsto(e^{i\theta}y,
e^{-i\theta}z)$ has moment map proportional to $|y|^2-|z|^2$,
so the sphere is Lagrangian. Its oriented normal Euler number is
$-\chi(S^2)=-2$: the order $(e_1,e_2,Je_1,Je_2)$ differs by one
transposition from the complex orientation $(e_1,Je_1,e_2,Je_2)$.
Weights $(2,1,1)$ place this sphere in the full nearby fibre with
the original multiplicities. Every sphere in a complex two-torus
lifts to $\mathbb C^2$ and contracts there, so its homology class,
and hence its self-intersection, is zero. This excludes the germ.

\item For finite $q$, the equation is analytically irreducible:
its weighted homogeneous polynomial of weights $(q,2,q-2)$ is
irreducible, and lowest weighted terms of an analytic factorization
would factor that polynomial. Thus its full multiplicity is a common
integer $m$. After the unit map above, put
\[
 G_s=x^2+y^2(z^2+y^{q-2}+s).
\]
For $s\ne0$ small, at each point
$(0,0,z_0)$ with $z_0^2=-s$ the function
$Z=z^2+y^{q-2}+s$ is a local coordinate:
its $z$ derivative is $2z_0\ne0$.
The reduced equation in these coordinates is
$x^2+y^2Z$, a pinch point. At the original
boundary double axis $x=y=0$, $z\ne0$,
the transverse Hessian determinant is
$4(z^2+s)$, which stays nonzero on the
compact boundary when $s$ is small.
All boundary smooth points remain smooth.
Thus the zero divisor retains the same
boundary smooth and double strata.
The spherical class of square $-4$
from the pinch-point nearby fibre transports,
by the boundary transport in Proposition~\ref{night:P4.20}, to
the original nearby torus fibre.
This excludes every finite $q$.

\item The remaining triple-point equations are
\[
 xyz+x^p,\quad p\ge3,\qquad
 xyz+x^p+y^q,\quad q\ge p\ge3.
\]
For the first put
$G_s=xyz+x^p+s x^2
=x(yz+x^{p-1}+s x)$.
When $s\ne0$, its two local normalization
branches are smooth: the derivative in $x$
of the second factor at the origin is $s$.
Their tangent planes coincide. The two branches of $uv=0$ have
distinct tangent planes, and an ambient biholomorphism has an invertible
derivative preserving that distinction. Hence these two tangent branches
cannot be a normal crossing. Along
the boundary $y$ axis the transverse
Hessian determinant is $-y^2$, and along
the boundary $z$ axis it is $-z^2$.
Both axes persist as ordinary-double
strata. For full component multiplicities $a,b$, weights $(2,p-1,p-1)$
give degree $2a+2(p-1)b$ and absorb the original unit. Use the deformation
$x^a(yz+x^{p-1}+sx)^b$. At its boundary double strata, free normal
coordinates $u,v$ give $u^av^b$, which is submersive at nonzero
values; thus the same boundary transport applies with these multiplicities.
For fixed nonzero $s$, the coordinate $X=x(s+x^{p-2})$ is
biholomorphic near the origin. The full deformed function becomes
a unit times $X^a(X+yz)^b$. The coordinates
\[
 x'=X+yz/2,\qquad y'=y,\qquad z'=z/(2i)
\]
satisfy $x'-iy'z'=X$ and $x'+iy'z'=X+yz$.
After absorbing the unit with weights $(2,1,1)$, the preceding
two-branch case supplies a sphere of square $-2$. Boundary transport
therefore gives the contradiction in the original full fibre.

\item For the second equation, positive weights
$(q,p,pq-p-q)$ and the primitive degree-one expression in $z$
show analytic irreducibility by the same initial-form argument.
It therefore has one common full multiplicity. Put
\[
 G_s=xyz+x^p+y^q+s x^2+s y^3,
 \qquad 0<|s|\ll1.
\]
Its transverse Hessian along the boundary
$z$ axis has determinant $-z^2$,
so its sole boundary double stratum
persists. We check the smooth normalization
of its germ at the origin explicitly.
Since $\partial^2G_s/\partial x^2(0)=2s$,
the critical equation in $x$ has a unique
analytic solution $x=\xi(y,z)$ with
$\xi$ divisible by $yz$. The parametric
holomorphic Morse change in $x$ writes
\[
 G_s=X^2+y^2H(y,z).
\]
Divisibility follows directly on substituting
$\xi$: each of $\xi yz,\xi^p,s\xi^2$
is divisible by $y^2$, as are the two
pure $y$ terms. Moreover
\[
 H_y(0,0)=
 \begin{cases}1+s&q=3,\\s&q>3,\end{cases}
 \qquad H(0,z)=-z^2/(4s).
\]
The finite
birational map
\[
 \{v^2+H(y,z)=0\}\longrightarrow\{G_s=0\},
 \qquad (v,y,z)\longmapsto(X=yv,y,z),
\]
followed by the inverse Morse coordinate
change, is the normalization. Its source
is smooth because $H_y(0,0)\ne0$,
and the map is finite because $v$
satisfies the displayed monic equation.
It is birational on $y\ne0$ and is
therefore the normalization by normality
of its source. At the origin its
differential has rank one. The normalization of a normal-crossing
germ consists of coordinate-plane inclusions, all of differential rank
two. The normalization isomorphism induced by an ambient biholomorphism
preserves these ranks, so this rank-one map excludes normal crossings. The local construction in Proposition~\ref{night:P4.20} produces
a spherical class of nonzero square, and its boundary transport gives
the contradiction. The transverse Hessians and the restrictions
$H(0,z)$ and $H_y(0,0)$ follow by collecting the degree-two terms
and substituting the critical equation
$yz+p\xi^{p-1}+2s\xi=0$.

\item The remaining germs in Table~1 are smooth, ordinary double,
and triple normal crossing. The linear coordinates
$(u,v)=(x+iy,x-iy)$ give $uv=0$ for the double germ;
the triple germ is $xyz=0$. Their normalization morphisms are the
inclusions of their disjoint smooth coordinate-plane branches.
Proposition~\ref{night:P4.6} applies to each compact support so obtained.
When every reduced support is semi-log-canonical, the same local
conclusion at every fibre supplies the hypotheses of
Propositions~\ref{night:P4.7}--\ref{night:P4.8}.

\end{enumerate}
\end{proof}


\subsection{Numerical sign for arbitrary irreducible supports}

\begin{proposition}[P4.26]\hypertarget{tag-02WY}{}\label{night:P4.26}\label{stmt:P4.26}
\StatementRecord{night:P4.26}{proved}{02WY}{\SeriesRef{night:P4.12}, \SeriesRef{night:P4.13}, \SeriesRef{night:P4.21}, \SeriesRef{night:P4.23}, \SeriesRef{night:P4.25}, \SeriesRef{night:P4.6}, \SeriesRef{night:P4.7}, \SeriesRef{night:P4.8}, \SeriesRef{night:P4R.1}, \SeriesRef{night:P4R.6}, \SeriesRef{night:P4R.7}, \SeriesRef{night:P4R.8}, \SeriesRef{night:P4R.9}}

In Part~4.2 every irreducible reduced fibre support
has nonnegative topological Euler number. For a
singular irreducible support, rational resolution
of the normalization gives strictly positive
Euler number; so does a resolution whose minimal
model is of class VII with positive second
Betti number. These assertions permit singular
normalization.

Every normal-crossing support has the
nonnegative even Euler number of P4.6.
Consequently, for a fibration all of whose
reduced supports are irreducible or have
normal crossings,
\[
 B=\mathbb P^1,\qquad c_3(Y)\in\{0,2\},
 \qquad b_3(Y)=2-c_3(Y).
\]
This last conclusion does not assume a
classification of the equality cases for
irreducible nonnormal supports.
\end{proposition}
\begin{proof}

\begin{enumerate}[label={$\langle1\rangle$\arabic*.},leftmargin=*]
\item For an irreducible singular support use the
minimal resolution of its normalization
$\rho:T\to N\to S$, and put
$A=\rho^{-1}(\operatorname{Sing}S)_{\rm red}$,
$C=\operatorname{Sing}(S)_{\rm red}$.
P4R.6 gives an effective integral divisor
$D$ with support $A$ and $K_T+D$ torsion.
The resolution $T$ has Kodaira dimension
$-\infty$, by P4.12. We explain the two
remaining transfers from smooth
normalization to this resolution.

\item Suppose first that $T$ is rational.
Push $D$ to a minimal rational surface.
The complete anticanonical-support calculation
in P4.21 has the following exact alternatives.
The reduced support is connected; either
its fundamental group is trivial and all
its normalized components are rational,
or the divisor itself is reduced with
normal crossings. On $\mathbb P^2$ these
are the plane cubic configurations
enumerated there. On a ruled minimal
surface they are the degree-two horizontal
configurations there. Those arguments use
only the effective integral divisor and
$K+D\equiv0$, so apply to the present
resolved divisor.

\item We check their preservation through the
actual blowup morphisms. A blowup at
conductor multiplicity $\mu$ gives
exceptional coefficient $\mu-1$.
For $\mu=1$, the centre is a smooth
point of the reduced support and its
strict transform is isomorphic to that
support. For $\mu\ge2$, the whole
exceptional $\mathbb P^1$ belongs to
the new support. Its contraction is
the blowup map on supports. Collapsing
this connected simply connected
subcomplex does not change the
fundamental group, by van Kampen.
The reduced normal-crossing case stays
reduced normal crossing: a smooth
point gives coefficient zero and a
double point gives coefficient one.
Thus the same alternatives hold for
the actual divisor $D$ on $T$.

\item In the reduced normal-crossing alternative,
P4.25 makes $S$ normal crossing. Its
normalization is smooth, so this minimal
normalization resolution is the
normalization itself. The rational case
of P4.6--P4.8 gives $e(S)=2$.
In the other alternative $\pi_1(A)=1$
and $\pi_1(T)=1$. The proper-map pushout
$T\cup_A C\to S$ is a homeomorphism:
it is a continuous compact-to-Hausdorff
bijection, since $\rho$ is an isomorphism
off $A$ and $A$ is the full inverse
image of $C$. Compatible triangulation
makes $A\to T$ a cofibration.
Van Kampen therefore gives
$\pi_1(S)\simeq\pi_1(C)$.
Every normalized component of $C$ is
rational: a nonconstant map from a
normalized component of $A$ onto it
is a finite map from $\mathbb P^1$.
Contracted components of $A$ contribute
no exception to this statement.
The curve $C$ is connected, and its
fundamental group is a free group
$F_b$ on the cycles of its
point-identification graph. P4R.1
implies $b\le1$. If $c\ge1$ is the
number of its components, then
\[
 e(C)=c+1-b\ge1.
\]
P4.13 gives $e(S_{\rm reg})\ge0$,
so $e(S)>0$.

\item Now suppose that the minimal model
$T_0$ is of class VII with
$b_2(T_0)>0$. The divisor pushed to
$T_0$ satisfies the hypotheses of
Dloussky's numerically anticanonical
theorem exactly as in P4.23.
The verified sources and their
normalization-independent hypotheses
are specified there. The
Inoue--Hirzebruch and parabolic Inoue
alternatives give a reduced
normal-crossing divisor, preserved
under the blowups just checked.
P4.25 then makes $S$ normal crossing
and $T=N$ smooth. The immersive
normalization Chern formula would
give $e(T)=K_T^2$, whereas
$e(T)=b_2(T)>0$ and
$K_T^2=-b_2(T)$. These alternatives
cannot occur.

\item In the intermediate alternative
$A$ is a connected rational cycle
with attached rational trees, and
\[
 \pi_1(A)=\mathbb Z,\qquad
 \pi_1(T)=\mathbb Z.
\]
Both statements survive the actual
blowups; they concern the resolution
and do not require $N$ to be smooth.
Every normalized component of $C$
is again rational and $C$ is connected.
With $\pi_1(C)=F_b$, the pushout
presentation is exactly
\[
 \pi_1(S)=
 \langle t,c_1,\ldots,c_b
 \mid t^k=w(c_1,\ldots,c_b)\rangle
\]
for the actual two conductor-group
homomorphisms. A component contracted
to a point merely induces the zero
homomorphism on its group. The
amalgam and free-product normal-form
argument in P4.23, which allowed
arbitrary $k$ and $w$, gives $b\le1$
from virtual abelianness.
Thus $e(C)=c+1-b\ge1$, and
$e(S)>0$ follows again from P4.13.

\item This exhausts these two resolution
classes. Ruled resolutions of positive
base genus are covered by P4R.7
and P4R.9, and Hopf minimal models
by P4R.8. The class-VII minimal
model with $b_2=0$ is Hopf or Inoue,
by the verified classification in
P4.8. An Inoue model has no curves,
so it cannot carry the nonempty
effective divisor pushed down from
$D$. Therefore it does not occur
here. All irreducible singular
supports have now been covered.

\item A normal irreducible support is
smooth by P4.12. The cyclic
base-change comparison and torus
deformation theorem in P4.8 show
that it has a finite unramified
torus cover, hence Euler number
zero. This proves the full
irreducible assertion.

\item For the final global statement,
$B=\mathbb P^1$ follows from
$b_1(Y)=0$ and connected fibres,
as proved in P4.7. Euler
additivity over the finitely many
exceptional base values gives
\[
 e(Y)=\sum_{b\in B}
 e\bigl((f^{-1}(b))_{\rm red}\bigr),
\]
because a smooth torus has Euler
number zero and the complement
is a proper smooth fibre bundle.
Every summand is nonnegative
under the stated support
hypothesis. Poincar\'e duality
and $b_1=b_2=0$ give
$e(Y)=2-b_3(Y)$; the nondegenerate
skew-symmetric cup-product
pairing on $H^3(Y;\mathbb Q)$
makes $b_3(Y)$ even. Thus
$0\le e(Y)\le2$ and $e(Y)$ is
even, which proves the asserted
values and their Chern-number
interpretation.

\end{enumerate}
\end{proof}


\subsection{Equality in the Hopf-resolution case}

\begin{proposition}[P4.27]\hypertarget{tag-02WZ}{}\label{night:P4.27}\label{stmt:P4.27}
\StatementRecord{night:P4.27}{proved}{02WZ}{\SeriesRef{night:P4.24}, \SeriesRef{night:P4R.1}, \SeriesRef{night:P4R.6}, \SeriesRef{night:P4R.8}}

Let \(S\) be an irreducible singular reduced fibre support in Part 4.2.
Suppose that the minimal resolution \(T\to N\) of its normalization has a
Hopf minimal model \(T_0\). Then
\[
 e(S)=0\quad\Longleftrightarrow\quad
 N\text{ is a smooth minimal Hopf surface}.
\]
In this case the reduced conductor and its image are disjoint unions of
smooth elliptic curves, and all component maps are finite unramified covers.
\end{proposition}
\begin{proof}

\begin{enumerate}[label={$\langle1\rangle$\arabic*.},leftmargin=*]
\item Write \(\rho:T\to N\to S\), \(A=\rho^{-1}(\operatorname{Sing}S)_{\rm red}\),
and \(C=\operatorname{Sing}(S)_{\rm red}\). Let \(k\) be the number of
blowups in \(\pi:T\to T_0\). P4R.6 and P4R.8 give an effective integral pole
divisor whose reduced support \(A\) has \(r\) connected components. Each is
a smooth elliptic root with a rational tree attached. Let \(v\) be the
number of its rational components. Every such component is one of the
\(k\) blowup exceptional curves. We have
\[
 e(S)=k-v+e(C),\qquad k-v\ge0,\qquad e(C)\ge0.
\]
Thus equality forces \(k=v\) and \(e(C)=0\).

\item Let \(c,s,g,b\) denote the numbers of irreducible components of \(C\),
connected components of \(C\), elliptic normalized components of \(C\), and
the sum of its point-identification graph cycle ranks. P4R.8 gives
\[
 e(C)=c+s-2g-b,\qquad
 b+r-s\le1,\qquad c\ge g,\quad r\ge g.
\]
Consequently \(e(C)=0\) implies \(c+r-2g\le1\).

\item First suppose \(c=g\). Every normalized component of \(C\) is elliptic.
Every rational component of \(A\) is contracted by \(T\to N\), since a
finite nonconstant map from \(\mathbb P^1\) to an elliptic normalization
does not exist. A root cannot be contracted. Indeed all rational curves
attached to it are already contracted, so its whole rooted connected
component would be a connected exceptional block disjoint from the strict
conductor. P4R.6 excludes precisely such a block. Hence all roots are
noncontracted.

\item The equality \(k=v\) now says that every blowup exceptional curve of
\(\pi:T\to T_0\) belongs to \(A\) and is contracted by \(T\to N\).
Conversely every exceptional curve of \(T\to N\) belongs to \(A\), by
P4R.6, and cannot be a root. The two proper birational maps therefore have
exactly the same exceptional curves. Their fibres are identical: each
nontrivial fibre is a connected component of the exceptional locus, and
their fibres are connected by Stein factorization and normality of the
targets.

\item Here is the holomorphic comparison, rather than only a topological
identification. Take the image of
\[
 (\pi,\rho_N):T\longrightarrow T_0\times N
\]
and then its normalization \(G\). The two projections \(G\to T_0,N\)
are proper birational maps with zero-dimensional fibres, because the
original two maps have identical fibres. They are finite. A finite
birational morphism to a normal complex space is an isomorphism, as follows
on local rings from integral closedness. Thus \(N\simeq T_0\)
biholomorphically. It is a smooth minimal Hopf surface.

\item The only other possibility is \(c=g+1,r=g\). There is one rational
normalized component of \(C\), and all \(g=r\) roots are noncontracted and
map to distinct elliptic normalized components of \(C\). Each such map is
an isogeny. Also
\[
 b+r-s=1.
\]
We exclude this possibility by a quotient of the exact fundamental-group
pushout; the quotient is constructed explicitly below.

\item Represent \(C\) as its disjoint normalized smooth curves with a finite
graph of point identifications. This is a graph of spaces with elliptic
vertex groups \(\mathbb Z^2\), trivial rational vertex groups, and trivial
edge groups. Its graph has \(s\) connected components and total cycle rank
\(b\). Add a vertex for \(T\), and add one root edge for each connected
component of \(A\). The root edge group is \(\mathbb Z^2\), since its
rational tree is simply connected. Its map to the \(C\) side factors
through the elliptic normalization receiving that root; its map to the
\(T\) side is the actual inclusion. This graph-of-spaces presentation is
the ordinary van Kampen presentation of the homotopy pushout
\(T\cup_A C=S\). It has total graph cycle rank \(b+r-s=1\).

\item Let \(\Gamma=\pi_1(T)=\pi_1(T_0)\). It is virtually infinite cyclic, and
\(b_1(T)=1\). The primitive map from its abelianization to its free part
therefore gives a surjection
\[
 h:\Gamma\longrightarrow\mathbb Z
\]
with finite kernel. Indeed a cyclic subgroup of finite index maps
nontrivially to \(\mathbb Z\), so the kernel meets that subgroup trivially
and is finite.

\item For root \(H_i\), write
\[
 a_i:\pi_1(H_i)=\mathbb Z^2\longrightarrow\Gamma,\qquad
 q_i:\pi_1(H_i)=\mathbb Z^2\longrightarrow
       \pi_1(\widetilde C_i)=\mathbb Z^2.
\]
The second map is injective of finite index because it is induced by an
isogeny. The composite \(h a_i\) is nonzero. To see this, a finite image
of \(a_i\) would become trivial after a finite cover of \(H_i\); that
compact elliptic cover would lift holomorphically to the universal Hopf
cover \(\mathbb C^2\setminus\{0\}\), where its coordinate functions are
constant. Thus \(a_i\) has infinite image; since \(\ker h\) is finite,
\(h a_i\ne0\).

\item Put \(K_i=\ker(h a_i)\), a primitive rank-one subgroup of
\(\mathbb Z^2\). Quotient the \(i\)-th elliptic vertex group by the
saturation of \(q_i(K_i)\). The resulting vertex group is \(\mathbb Z\).
Quotient the root edge group by \(K_i\); it too becomes \(\mathbb Z\).
Its two induced maps into the \(T\) vertex \(\mathbb Z\) and the elliptic
vertex \(\mathbb Z\) are nonzero multiplications and hence injective.
The latter assertion follows because an element outside \(K_i\) cannot
map under the injective rational map \(q_i\otimes\mathbb Q\) into the
rational span of \(q_i(K_i)\). This verifies compatibility of the chosen
quotients with both edge homomorphisms.

\item We have constructed a surjective homomorphism from \(\pi_1(S)\) onto the
fundamental group \(G\) of a graph of groups. Its nontrivial vertex groups
are infinite cyclic, its root edges are infinite cyclic with injective
maps, and all point-identification edges are trivial. The root edges
form a star: the roots map to distinct elliptic normalizations. They
therefore contain no graph cycle. The unique graph cycle of the complete
graph contains a point-identification edge, whose edge group is trivial.

\item Delete one such edge on the cycle. The remaining graph is a connected
tree, and the normal form for iterated amalgamated products along
injective edge maps embeds its \(T\) vertex group \(\mathbb Z\) into its
fundamental group \(G_0\). Restoring the removed edge is an HNN extension
over the trivial group, hence
\[
 G\simeq G_0*\mathbb Z.
\]
Choosing a nontrivial element of the embedded \(\mathbb Z\subset G_0\)
and the new stable letter gives an embedded free group \(F_2\), by reduced
free-product normal form. But every quotient of a virtually abelian group
is virtually abelian, and a virtually abelian group has no \(F_2\)
subgroup. This contradicts P4R.1. The case \(c=g+1,r=g\) is impossible.

\item We have proved that equality makes \(N\) smooth minimal Hopf. Conversely
P4.24 applies to such an actual normalization and proves the unramified
elliptic-conductor description and \(e(S)=0\). This completes both
directions.

\end{enumerate}
\end{proof}


\subsection{Cuspidal transverse type and the ruled equality case}

\begin{proposition}[P4.28]\hypertarget{tag-02X0}{}\label{night:P4.28}\label{stmt:P4.28}
\StatementRecord{night:P4.28}{proved}{02X0}{\SeriesRef{night:P4.16}, \SeriesRef{night:P4.18}, \SeriesRef{night:P4.20}}

Let $p$ lie on a reduced fibre support of Part~4.2.
Suppose its reduced singular locus is a smooth curve
near $p$, and its generic transverse plane-curve germ
is the cusp $a^2+b^3=0$. Then the actual germ at $p$
is analytically isomorphic to
$\{a^2+b^3=0\}\times(\mathbb C,0)$.
In particular its normalization is smooth.
\end{proposition}
\begin{proof}

\begin{enumerate}[label={$\langle1\rangle$\arabic*.},leftmargin=*]
\item We extend the boundary argument of P4.20 in the
precise form used here. Suppose a deformation retains
the smooth critical curve near a fixed transverse
boundary sphere and retains transverse type $A_2$.
Near a boundary point, the holomorphic Morse lemma
in the nondegenerate transverse variable followed
by extraction of a cube root of a unit gives
coordinates with defining function $a^2+b^3$.
The third coordinate parametrizes the critical
curve. Its transversality to the boundary lets
one real third-coordinate variable be solved
from the boundary equation; the complex
coordinates $a,b$ remain free. Thus the
restriction to the boundary is submersive
at every nonzero level. Compactness and
the smooth part of the boundary give the
same uniform nonzero-level submersion as
in P4.20. Its finite-critical-value and
proper Ehresmann transport proof applies
unchanged. A compact spherical class of
nonzero square in a deformed nearby
fibre would therefore lie in an actual
nearby torus fibre, which is impossible.
All deformations below retain this
boundary transverse type.

\item Choose holomorphic coordinates $(x,y,z)$
so the smooth critical curve is $x=y=0$.
The transverse Hessian along it is a
symmetric rank-one matrix over the
discrete valuation ring $\mathbb C\{z\}$.
It is not the zero matrix generically.
It can be diagonalized holomorphically
to $h(z)x^2$. Here is the elementary
ring calculation. The determinant is
zero. Choose a diagonal entry of
minimal order; one exists since the
square of the off-diagonal entry is
the product of the two diagonal entries.
The ratio of the off-diagonal entry
to the chosen entry is holomorphic.
Completing its square gives the desired
invertible coordinate change. The
coefficient $h$ is not identically zero.

\item First suppose that the Hessian has
rank one at $p$. The Morse lemma with
parameters $(y,z)$ writes the reduced
equation as
\[
 g=X^2+y^3h_1(y,z).
\]
Indeed after removing the $X$ direction,
the remaining function and its first
two $y$ derivatives vanish on the
critical curve; the third derivative
is not identically zero there. If
$h_1(0,0)\ne0$, extracting its cube
root in the $y$ coordinate proves the
asserted product form.

\item Otherwise $h_1(0,z)$ has a zero of
finite positive order. Use
$g_\eta=X^2+y^3(h_1+\eta)$.
For arbitrarily small generic $\eta$,
the function $h_1(0,z)+\eta$ has a
simple zero $z_\eta$ near zero.
This follows directly from one-variable
Weierstrass preparation, avoiding the
finite critical values of $h_1(0,z)$.
At $(0,0,z_\eta)$ the function
$Z=h_1(y,z)+\eta$ is a coordinate,
since its $z$ derivative is nonzero.
The actual germ is
$X^2+y^3Z$, which is excluded by
P4.16: its nearby fibre has a
spherical class of square $-6$.
On a fixed sufficiently small boundary,
$h_1(0,z)$ is nonzero at the critical
curve, and this persists for small
$\eta$. The deformed transverse type
is therefore still $A_2$ there.
The preceding boundary transport gives
the contradiction.

\item It remains to exclude Hessian rank
zero at $p$, so $h(0)=0$. In the
diagonalizing coordinates add
$\varepsilon y^3$ to $g$. The
critical curve and its transverse
Hessian do not change. The cubic
coefficient in the transverse kernel
remains nonzero near the original
compact boundary for small
$\varepsilon$, so this deformation
has the boundary property just proved.
Choose $\varepsilon$ so that the
coefficient of $y^3$ at the origin
is nonzero.

\item On a smaller neighbourhood, Weierstrass
preparation in $y$ gives a unit times
\[
 P=y^3+a(x,z)y^2+b(x,z)y+c(x,z).
\]
Restriction to $x=0$ is a unit times
$y^3$, so all three coefficients
vanish at $x=0$. The critical-curve
equation $g_x(0,0,z)=0$ first gives
$c_x(0,z)=0$. In the identity
$g=U P$ the mixed derivative on
the critical curve is then
$g_{xy}=U b_x$, since the terms
containing $c,c_x,b$ vanish there.
The zero $xy$ Hessian entry therefore
gives $b_x(0,z)=0$. Thus
$x^2\mid b$, $x^2\mid c$, and
$x\mid a$.
The Tschirnhaus change $y=Y-a/3$
therefore gives exactly
\[
 P=Y^3+x^2H(x,Y,z).
\]
The coefficient $H(0,0,z)$ is a unit
multiple of the original Hessian
coefficient $h(z)$. It is not
identically zero but vanishes at
zero. Perturb it to $H+\eta$.
At a simple root of
$H(0,0,z)+\eta$, the coordinates
$(x,Y,Z=H+\eta)$ give
$Y^3+x^2Z$. P4.18 excludes this
binomial germ by its compact
spherical class of nonzero square.

\item There are two nested transport
neighbourhoods in this argument.
First fix the original ball for
$g\to g+\varepsilon y^3$.
Then choose a smaller ball in the
Weierstrass coordinate neighbourhood,
where $g+\varepsilon y^3$ has only
the original critical curve and its
boundary has transverse type $A_2$.
The second perturbation transports
its spherical class into a nearby
fibre on this smaller ball. Inclusion
puts it in a fibre on the original
ball; the first boundary transport
puts it in a fibre of $g$.
This does not require extending the
second perturbation outside its
coordinate neighbourhood.
Units in the reduced equations are
retained during the perturbations.
At a binomial special point a unit
can be absorbed by the positive
weighted coordinate scaling of
P4.16--P4.18.

\item We have ruled out both a zero
Hessian and a vanishing cubic
coefficient. The product form
follows. Its normalization is
the finite map
\[
 (u,z)\longmapsto (a=u^3,b=-u^2,z),
\]
whose source is smooth, whose
inverse rational parameter off
$b=0$ is $u=-a/b$, and whose
restriction to the conductor is
the identity on the $z$ coordinate.

\end{enumerate}
\end{proof}

\begin{proposition}[P4.29]\hypertarget{tag-02X1}{}\label{night:P4.29}\label{stmt:P4.29}
\StatementRecord{night:P4.29}{proved}{02X1}{\SeriesRef{night:P4.12}, \SeriesRef{night:P4.13}, \SeriesRef{night:P4.22}, \SeriesRef{night:P4.28}, \SeriesRef{night:P4R.7}, \SeriesRef{night:P4R.9}}

Let $S$ be an irreducible singular support
in Part~4.2 whose normalization has ruled
resolution, and suppose $e(S)=0$.
In the double-horizontal-divisor case,
when the generic transverse type is
cuspidal $A_2$, the normalization is
a smooth minimal elliptic ruled surface.
The image conductor is a smooth elliptic
curve and every point along it has the
product cusp germ of P4.28.
\end{proposition}
\begin{proof}

\begin{enumerate}[label={$\langle1\rangle$\arabic*.},leftmargin=*]
\item Use the notation of P4R.7 and P4R.9.
The actual resolved pole support is
a section $H$ with rational trees,
and its inclusion in the ruled
resolution $T$ induces an isomorphism
on fundamental groups. Thus
$\pi_1(S)=\pi_1(C)$, where
$C=\operatorname{Sing}(S)_{\rm red}$.
The virtually abelian curve-group
calculation gives $e(C)\ge0$, and
P4.13 gives $e(S_{\rm reg})\ge0$.
Equality makes both numbers zero.
The same connected-curve calculation
shows that $C$ has exactly one
normalized component, of genus one,
and has no point-identification
cycles.

\item All rational components of the
resolved pole support are contracted
to points: a nonconstant proper map
from a rational curve to the
elliptic normalization of $C$ is
impossible. Hence the strict
conductor consists only of $H$.
The generic cusp has one
normalization branch, so
$H\to C$ has generic degree one
and is the normalization of $C$.
In particular $H$ and the ruling
base both have genus one.

\item Every exceptional curve of
$T\to N$ is now rational and
vertical for the ruling, since
$H$ is not contracted. The ruling
$r:T\to H$ descends to a holomorphic
map $\bar r:N\to H$. For an explicit
descent, the image of
$(\rho_N,r):T\to N\times H$
projects properly and bijectively
to $N$, because $r$ is constant
on each exceptional fibre. Its
normalized graph projects finitely
and birationally to the normal
space $N$, hence is isomorphic
to $N$.

\item The normalization singularities
are rational. Indeed a sequence
of point blowups of a
$\mathbb P^1$-bundle has
$R^1r_*\mathcal O_T=0$.
For a blowup this follows from
the two affine charts, or from
the exact sequence of the
exceptional $\mathbb P^1$ with
$H^1(\mathcal O_{\mathbb P^1})=0$;
the assertion for the whole
sequence follows by Leray.
Since $\rho_{N*}\mathcal O_T
=\mathcal O_N$, the low-degree
Leray sequence for
$r=\bar r\circ\rho_N$ gives
\[
 \bar r_*R^1\rho_{N*}\mathcal O_T=0.
\]
Here $R^2\bar r_*\mathcal O_N=0$
because $\bar r$ has one-dimensional
fibres. The sheaf
$R^1\rho_{N*}\mathcal O_T$ is
supported at finitely many
normalization singular points,
and direct image of a nonzero
such sheaf to $H$ is nonzero.
Thus it vanishes. In particular
$\chi(\mathcal O_N)
=\chi(\mathcal O_T)=0$.

\item We next check the scheme structure
of the conductor, to prove that
its image curve is smooth.
Write $I$ for the conductor ideal
on $S$ and on $N$, and let
$C_{\rm sch}\subset S$ and
$D_{\rm sch}\subset N$ be its
two zero schemes. Locally
$\mathcal O_S$ is Gorenstein
of dimension two and
$\nu_*\mathcal O_N$ is a
maximal Cohen--Macaulay module.
Finite duality identifies
its dual module with $I$,
which is again maximal
Cohen--Macaulay. The depth
lemma in
$0\to I\to\mathcal O_S
\to\mathcal O_{C_{\rm sch}}\to0$
therefore gives depth one
for the conductor curve.
At a generic cusp its local
length in the transverse
direction is one. A
one-dimensional Cohen--Macaulay
scheme which is generically
reduced is reduced: any
nilpotent ideal would have
zero-dimensional support,
contradicting the absence of
embedded associated points.
Consequently $C_{\rm sch}=C$.

\item Put $L=\nu^*\omega_S$, a
torsion line bundle by P4.12.
Finite duality on $N$ gives
\[
 I=\omega_N\otimes L^{-1}.
\]
Serre duality on the compact
Cohen--Macaulay surface $N$
gives $\chi(I)=\chi(L)$.
The rational-singularity
comparison just proved and
the projection formula give
$\chi(L)=\chi(\rho_N^*L)$.
Riemann--Roch on the smooth
surface $T$, with the rational
Chern class of this torsion
line bundle equal to zero,
gives $\chi(\rho_N^*L)
=\chi(\mathcal O_T)=0$.
Therefore
$\chi(\mathcal O_{D_{\rm sch}})
=\chi(\mathcal O_N)-\chi(I)=0$.

\item The exact conductor square is
\[
 0\longrightarrow\mathcal O_S
 \longrightarrow
 \nu_*\mathcal O_N\oplus\mathcal O_C
 \longrightarrow\nu_*\mathcal O_{D_{\rm sch}}
 \longrightarrow0.
\]
Taking Euler characteristics and
using $\chi(\mathcal O_S)=0$
from P4.12 yields
$\chi(\mathcal O_C)=0$.
Its normalization is the
smooth elliptic curve $H$,
so the normalization sequence
gives $\chi(\mathcal O_C)
 =-\sum_p\delta(C,p)$.
All delta invariants vanish;
thus $C$ is smooth.

\item Every point of $C$ now satisfies
P4.28, so $S$ has the product
cusp germ there and $N$ is
smooth along the conductor.
It was already smooth away
from the conductor. P4.22
therefore applies and makes
this elliptic ruled
normalization minimal.

\end{enumerate}
\end{proof}


\subsection{The conormal calculation and ordinary-triple equality}

\begin{proposition}[P4.30]\hypertarget{tag-02X3}{}\label{night:P4.30}\label{stmt:P4.30}
\StatementRecord{night:P4.30}{proved}{02X3}{the definitions and constructions stated in the proof}

Let $C$ be a compact irreducible reduced curve in the smooth
threefold $Y$ of Part~4.2. Suppose its normalization is an elliptic
curve $E$, and every point of $C$ is unibranch. Write $P$ for its
ideal in $Y$ and $P^{(2)}$ for the ideal of functions vanishing to
order at least two at the generic point of $C$. Then
\[
 \chi(\mathcal O_Y/P^{(2)})\le0.
\]
Equality forces $C$ to be smooth. The assertion uses
$c_1(TY)\cdot C=0$, which follows from $b_2(Y)=0$.
\end{proposition}
\begin{proof}

\begin{enumerate}[label={$\langle1\rangle$\arabic*.},leftmargin=*]
\item Let $j:E\to Y$ be the normalization map followed by inclusion.
The derivative map
\[
 j^*\Omega_Y\longrightarrow\Omega_E
\]
has image $\Omega_E(-R)$ for an effective divisor $R$.
Its kernel $K$ is locally free of rank two: dividing the derivative
row by its common zero gives a surjection onto a line bundle.
Since $\deg j^*\Omega_Y=0$ and $\deg\Omega_E=0$, we have
$\deg K=\deg R$.

\item Differentiation induces an injection
\[
 Q:=P/P^{(2)}\longrightarrow j_*K.
\]
Indeed its kernel consists precisely of functions with zero first
normal derivative at a generic smooth point of $C$, which is the
defining order-two condition for $P^{(2)}$. Both sides agree at the
generic point. Their quotient has finite length, denoted by $\tau$.
Riemann--Roch on the smooth elliptic curve gives
\[
 \chi(Q)=\deg R-\tau.
\]
We prove $\tau\ge\deg R$ by local equations, including at a
singular point of $C$.

\item Parametrize its sole branch by $t$. Choose ambient coordinates
vanishing at the point, and let
$m=\min_i\operatorname{ord}_t x_i(t)$ be the multiplicity.
The common order of the three derivative entries is $m-1$,
so the coefficient of $R$ is $m-1$. The saturated kernel $K$
is a direct summand of the free rank-three module over
$\mathbb C\{t\}$. Consequently a generator of $\bigwedge^2K$
has at least one unit Pl\"ucker coordinate. If $L$ is the
$\mathbb C\{t\}$-span of the differentiated defining equations,
then the colength of $L$ in $K$ is the minimum order of their
two-by-two minors. This is the diagonal elementary-divisor
calculation for a rank-two submodule over a discrete valuation
ring. The local length of $j_*K/Q$ is at least this colength,
since $Q\subset L\subset K$.

\item At a smooth point $m=1$ and there is nothing to prove.
At a singular point of embedding dimension three, every defining
equation belongs to the square of the ambient maximal ideal.
Every entry of its gradient therefore has order at least $m$
on the parametrization, and every two-by-two minor has order
at least $2m$. Hence the local quotient length is at least
$2m\ge m-1$.

\item At a singular point of embedding dimension two, an invertible
ambient coordinate change writes the ideal as $(z,h(x,y))$,
where $h$ is a reduced irreducible plane equation without a
linear term. The conormal image has a free $dz$ summand.
Its other generator is $dh$; both its plane coefficients have
order at least $m$. The primitive rank-one kernel of the plane
derivative row has a unit coefficient, so $dh$ is its generator
multiplied by a function of order at least $m$. The quotient
length is therefore at least $m\ge m-1$. Embedding dimension
one is the smooth case. These cases prove the local bound and
therefore $\chi(Q)\le0$.

\item The normalization exact sequence gives
$\chi(\mathcal O_C)=-\sum_p\delta(C,p)$.
Finally the exact sequence
\[
 0\longrightarrow P/P^{(2)}\longrightarrow
 \mathcal O_Y/P^{(2)}\longrightarrow\mathcal O_C
 \longrightarrow0
\]
gives
\[
 \chi(\mathcal O_Y/P^{(2)})
 =-\sum_p\delta(C,p)+\deg R-\tau\le0.
\]
Equality forces all delta invariants to vanish, so the
normalization is an isomorphism and $C$ is smooth.

\end{enumerate}
\end{proof}

\begin{proposition}[P4.31]\hypertarget{tag-02X4}{}\label{night:P4.31}\label{stmt:P4.31}
\StatementRecord{night:P4.31}{proved}{02X4}{\SeriesRef{night:P4.22}, \SeriesRef{night:P4.29}, \SeriesRef{night:P4.30}, \SeriesRef{night:P4R.6}, \SeriesRef{night:P4R.7}, \SeriesRef{night:P4R.9}}

Let $S$ be an irreducible singular support in Part~4.2 whose
normalization has ruled resolution, and suppose $e(S)=0$.
In the double-horizontal-divisor case, when the generic
transverse type is the ordinary triple point $xy(x-y)=0$,
the normalization is a smooth minimal elliptic ruled surface.
Its conductor image is a smooth elliptic curve, and its
reduced conductor is an unramified cover of that curve of
degree three. Every germ along the image is analytically
the product of an ordinary plane triple point with a disc.
\end{proposition}
\begin{proof}

\begin{enumerate}[label={$\langle1\rangle$\arabic*.},leftmargin=*]
\item Use the minimal resolution $T\to N\to S$ and the section $H$
of P4R.7 and P4R.9. Put $g=g(H)$ and
$C=\operatorname{Sing}(S)_{\rm red}$.
The same fundamental-group and Euler argument as in P4.29
shows that $C$ is connected, has just one normalized
component, that this normalization $E$ has genus one, and
that its point-identification graph has no cycles. Thus
every point of $C$ is unibranch. All rational components
of the resolved pole support are contracted, and the
only strict conductor component is $H$. Its finite map
to $C$ has degree three, since an ordinary transverse
triple point has three normalization branches.
In particular $g\ge1$ by Riemann--Hurwitz for $H\to E$.

\item All exceptional curves of $T\to N$ are rational and
vertical for the ruling. Exactly the normalized-graph
descent and Leray calculation of P4.29 give a morphism
$N\to H$ and
\[
 R^1\rho_{N*}\mathcal O_T=0,\qquad
 \chi(\mathcal O_N)=\chi(\mathcal O_T)=1-g.
\]
Thus the normalization singularities are rational; no
smoothness assertion is used here.

\item Let $I$ be the conductor ideal, and let $C_{\rm sch}$
and $D_{\rm sch}$ be its zero schemes in $S$ and $N$.
Finite duality and the depth argument of P4.29 show that
$\mathcal O_{C_{\rm sch}}$ is Cohen--Macaulay of pure
dimension one. If $P\subset\mathcal O_Y$ is the reduced
ideal of $C$, the inverse image $J$ of $I$ in
$\mathcal O_Y$ therefore has no embedded associated
prime. Its only associated prime locally is $P$.
The generic ordinary triple point has conductor equal
to the square of its maximal ideal in the transverse
plane. To check this directly, normalize its three
lines. Constant functions on them have a common value;
their linear terms are the evaluations of one linear
form on three distinct directions, a two-dimensional
space. Quadratic forms evaluate surjectively on those
three directions. All terms of degree at least two
can therefore be chosen independently on the three
branches, and a nonzero constant or linear term cannot
belong to the conductor. This proves the stated square
ideal, including its scheme structure.

\item Consequently $J_P=P_P^2$, and in fact
\[
 J=P^{(2)}.
\]
For completeness, if $a\in P^{(2)}$, some $s\notin P$
satisfies $sa\in J$. Multiplication by $s$ is injective
on $\mathcal O_Y/J$, because its only associated prime
is $P$, so $a\in J$. The reverse inclusion follows by
localizing at $P$. Thus the conductor image is the
actual symbolic double curve to which P4.30 applies.

\item Put $L=\nu^*\omega_S$, which is torsion. Finite duality
gives $I=\omega_N\otimes L^{-1}$. Serre duality, the
rational-singularity comparison and Riemann--Roch on
$T$ give
\[
 \chi(I)=\chi(L)=\chi(\rho_N^*L)
   =\chi(\mathcal O_T)=1-g.
\]
Hence $\chi(\mathcal O_{D_{\rm sch}})=0$.
Taking Euler characteristics in the exact conductor
square, and using $\chi(\mathcal O_S)=0$, yields
\[
 \chi(\mathcal O_{C_{\rm sch}})=g-1.
\]
P4.30 now gives $g-1\le0$. Since $g\ge1$, equality
holds, $g=1$, and $C$ is smooth. Riemann--Hurwitz then
makes the degree-three map $H\to C$ unramified.

\item We finish by proving normalization smoothness from
the actual local morphisms. Over a point of $C$ the
three distinct points of $H$ remain distinct in $N$:
every exceptional block is vertical for the ruling
$T\to H$ and therefore cannot join different points
of $H$. The finite normalization consequently gives
three local irreducible branches of $S$. Each branch
is a hypersurface in $Y$ and is smooth at the generic
point of $C$. Its singularities are isolated, so the
hypersurface Cohen--Macaulay property and regularity
in codimension one make it normal. The corresponding
component of $N$ is that branch itself.

\item Write its three local equations as $f_1,f_2,f_3$.
The intersection of any two branches is supported on
$C$ and has generic multiplicity one, because the
generic transverse triple point is ordinary. On a
normal branch the other equation is a nonzero divisor.
Its quotient is a one-dimensional Cohen--Macaulay
ring, generically reduced with support $C$, so it is
reduced and is exactly $\mathcal O_C$. Thus $C$ is a
smooth Cartier curve in each branch. A local ring of
dimension two with a nonzero divisor whose quotient
is a regular one-dimensional local ring has embedding
dimension at most two: lift its one maximal-ideal
generator and add the nonzero divisor. It is regular.
All three branches are therefore smooth.

\item Any two branches meet transversely along $C$.
Choose their equations as coordinates $x,y$, with
$z$ a coordinate on $C$. The third equation is
$a(x,y,z)x+b(x,y,z)y$, where both $a$ and $b$ are
units by the pairwise transversality just proved.
The invertible coordinate change $X=ax$, $Y=-by$
writes the reduced union exactly as
$XY(X-Y)=0$. This is the asserted product triple
point, and its normalization is smooth. Away from
the conductor, normalization singularities were
already excluded by P4R.6. Finally P4.22 applies to
this smooth elliptic ruled normalization and proves
its minimality.

\end{enumerate}
\end{proof}


\subsection{Complete zero-Euler classification in the requested classes}

\begin{proposition}[P4.32]\hypertarget{tag-02X5}{}\label{night:P4.32}\label{stmt:P4.32}
\StatementRecord{night:P4.32}{proved}{02X5}{\SeriesRef{night:P4.12}, \SeriesRef{night:P4.22}, \SeriesRef{night:P4.24}, \SeriesRef{night:P4.26}, \SeriesRef{night:P4.27}, \SeriesRef{night:P4.29}, \SeriesRef{night:P4.31}, \SeriesRef{night:P4.8}, \SeriesRef{night:P4R.11}, \SeriesRef{night:P4R.7}, \SeriesRef{night:P4R.9}}

An irreducible support in Part~4.2 has Euler number zero exactly
in the following cases: it is smooth with a finite unramified
torus cover, or it is singular with smooth minimal elliptic ruled
normalization, or it is singular with smooth minimal Hopf
normalization. In the ruled double-horizontal case the conductor
image is a smooth elliptic curve; the local germs are product
cusps, product tacnodes, or product ordinary plane triple points,
and the reduced conductor maps have degrees one, two, or three
respectively and are unramified.

For a fibration whose supports are all irreducible or normal
crossing, $c_3(Y)=0$ exactly when every irreducible support has
one of these forms and every reducible support has the
zero-Euler normal-crossing form of P4.8.
\end{proposition}
\begin{proof}

\begin{enumerate}[label={$\langle1\rangle$\arabic*.},leftmargin=*]
\item Normal irreducible supports are smooth with the stated finite
unramified cover by P4.8 and P4.12. For a singular irreducible
support, use its minimal normalization resolution $T$.
P4.26 excludes rational resolutions and class-VII minimal
models with positive second Betti number from equality.
The Inoue minimal model has no conductor curves and is
excluded there as well. A Hopf minimal model is covered
completely by P4.27: equality makes the actual normalization
a smooth minimal Hopf surface, and the converse is P4.24.

\item It remains to use the positive-genus ruled resolution.
The reduced-horizontal alternatives of P4R.9 already give
smooth minimal elliptic ruled normalization. In the remaining
case the pole support is a section $H$ with rational trees,
and the coefficient of $H$ is two. Its inclusion in $T$
induces an isomorphism on fundamental groups. The conductor
pushout and the curve-group calculation of P4R.7--P4R.9 give
$\pi_1(S)=\pi_1(C)$ and $e(C)\ge0$ for the reduced image
conductor $C$. Equality in
$e(S)=e(S_{\rm reg})+e(C)$ makes $C$ a connected curve
with elliptic normalization and no point-identification
cycles. Every rational pole component is contracted, since
there is no finite nonconstant map from a rational curve
to an elliptic curve. Thus $H$ is the sole strict conductor
component, and every branch of a generic transverse
plane-curve singularity has conductor exponent two.

\item We enumerate those transverse germs explicitly. For a branch
$B_i$ of a reduced plane curve, its conductor exponent is
\[
 c_i=2\delta(B_i)+\sum_{j\ne i}I(B_i,B_j).
\]
If a branch is singular, $c_i=2$ forces
$\delta(B_i)=1$ and no other branch. The one-gap
parametrization of a delta-one irreducible plane curve
has orders two and three and gives the ordinary cusp.
If all branches are smooth, two branches have intersection
multiplicity two and give a tacnode, or three branches
have all pairwise intersection multiplicities one and
give an ordinary triple point. Four branches would give
each conductor exponent at least three. One smooth branch
is nonsingular and has exponent zero. These exhaust the
possibilities.

\item P4.29, P4R.11, and P4.31 respectively treat the cusp,
tacnode, and ordinary-triple cases, without assuming that
the initial normalization or image conductor was smooth.
They give the asserted smooth minimal elliptic ruled
normalization and the precise local conductor maps.
The converse for a smooth minimal elliptic ruled
normalization is P4.22. This proves the entire irreducible
zero-Euler classification.

\item Finally P4.26 expresses $c_3(Y)$ as a finite sum of
nonnegative fibre Euler numbers. Its vanishing is therefore
equivalent to vanishing of every summand. The irreducible
summands have just been classified, and P4.8 gives the
complete reducible normal-crossing zero case. This proves
the final assertion.

\end{enumerate}
\end{proof}


\subsection{Normalization cohomology for the remaining positive cases}

\begin{proposition}[P4.33]\hypertarget{tag-02X6}{}\label{night:P4.33}\label{stmt:P4.33}
\StatementRecord{night:P4.33}{proved}{02X6}{\SeriesRef{night:P4.12}, \SeriesRef{night:P4.26}, \SeriesRef{night:P4.29}}

Let $S$ be an irreducible singular support in Part~4.2, and let
$\rho:T\to N\to S$ be its minimal normalization resolution.
Then $H^2(N,\mathcal O_N)=0$. If
$p=\dim H^0(N,R^1\rho_{N*}\mathcal O_T)$ is the sum of the
geometric genera of the normalization singularities, then
\[
 p=h^1(T,\mathcal O_T)-h^1(N,\mathcal O_N).
\]
In particular a rational resolution forces every normalization
singularity to be rational and gives $\chi(\mathcal O_N)=1$.
For the scheme-theoretic image conductor $C_{\rm sch}$ we have
\[
 \chi(\mathcal O_{C_{\rm sch}})
   =-\chi(\mathcal O_N)=h^1(N,\mathcal O_N)-1.
\]

In the ruled double-horizontal case, with section $H$ of genus
$g\ge1$, a contracted $H$ gives $h^1(N,\mathcal O_N)=0$ and
$p=g$. A noncontracted $H$ gives $p=0$ and
$h^1(N,\mathcal O_N)=g$.
\end{proposition}
\begin{proof}

\begin{enumerate}[label={$\langle1\rangle$\arabic*.},leftmargin=*]
\item Write $I_S\subset\mathcal O_S$ for the conductor ideal.
Finite duality gives
\[
 \nu_*\omega_N=I_S\otimes\omega_S.
\]
The dualizing line bundle $\omega_S$ is torsion by P4.12.
A nonzero section of a torsion line bundle on a compact
irreducible reduced complex space has no zero: raising it
to a trivializing positive power gives a nonzero holomorphic
function, hence a nonzero constant. Such a section cannot
belong to the proper conductor ideal along the nonempty
conductor. Therefore $H^0(N,\omega_N)=0$. Serre duality
on the compact Cohen--Macaulay normal surface gives
$H^2(N,\mathcal O_N)=0$.

\item The resolution is birationally ruled or of class VII and
has $H^2(T,\mathcal O_T)=0$, by P4.12 and the classification
already used in P4.26. The low-degree Leray sequence,
using $\rho_{N*}\mathcal O_T=\mathcal O_N$ and the
finite support of $R^1\rho_{N*}\mathcal O_T$, is therefore
\[
 0\longrightarrow H^1(N,\mathcal O_N)
 \longrightarrow H^1(T,\mathcal O_T)
 \longrightarrow H^0(N,R^1\rho_{N*}\mathcal O_T)
 \longrightarrow0.
\]
This proves the formula for $p$. If $T$ is rational,
$h^1(\mathcal O_T)=0$, so $p=0$. Each nonnegative local
geometric genus is zero, giving the asserted rational
singularities and Euler characteristic.

\item For completeness the conductor Euler identity does not
require rational singularities. Put $L=\nu^*\omega_S$,
and write $I_N=\omega_N\otimes L^{-1}$ for the conductor
on $N$. Serre duality gives $\chi(I_N)=\chi(L)$.
Riemann--Roch on $T$, applied to the torsion line bundle
$\rho_N^*L$, gives
$\chi(\rho_N^*L)=\chi(\mathcal O_T)$.
The projection formula makes
$R^1\rho_{N*}\rho_N^*L=L\otimes R^1\rho_{N*}\mathcal O_T$;
its finite length is still $p$. Leray therefore gives
$\chi(L)=\chi(\mathcal O_N)$.
Consequently the conductor scheme $D_{\rm sch}\subset N$
has $\chi(\mathcal O_{D_{\rm sch}})=0$.
The exact conductor square, together with
$\chi(\mathcal O_S)=0$, proves
$\chi(\mathcal O_{C_{\rm sch}})=-\chi(\mathcal O_N)$.
Since $N$ is compact connected and $H^2(\mathcal O_N)=0$,
the remaining equality follows.

\item Finally suppose the ruled pole support is $H$ with
rational vertical trees. If $H$ is contracted, the
composition $H\to T\to N$ is constant. Pullback from
$H^1(N,\mathcal O_N)$ to $H^1(H,\mathcal O_H)$ is zero.
But restriction
$H^1(T,\mathcal O_T)\to H^1(H,\mathcal O_H)$ is an
isomorphism: the ruling has section $H$, a
$\mathbb P^1$-bundle has zero first higher direct image
of its structure sheaf, and the actual point blowups
preserve these groups. The Leray injection above
therefore forces $H^1(N,\mathcal O_N)=0$, so $p=g$.
If $H$ is not contracted, every exceptional curve is
vertical. The ruling descends by the normalized-graph
comparison of P4.29. Its Leray calculation gives $p=0$,
and the same exact sequence gives $h^1(\mathcal O_N)=g$.

\end{enumerate}
\end{proof}


\subsection{The value one in the rational ordinary-triple sector}

\begin{proposition}[P4.34]\hypertarget{tag-02X7}{}\label{night:P4.34}\label{stmt:P4.34}
\StatementRecord{night:P4.34}{proved}{02X7}{\SeriesRef{night:P4.26}, \SeriesRef{night:P4.30}, \SeriesRef{night:P4.31}, \SeriesRef{night:P4.33}, \SeriesRef{night:P4.8}}

An irreducible singular support with rationally resolved
normalization and generic transverse ordinary triple point
cannot have Euler number one. For a fibration all of whose
supports are irreducible or normal crossing, its Euler number
is therefore two.
\end{proposition}
\begin{proof}

\begin{enumerate}[label={$\langle1\rangle$\arabic*.},leftmargin=*]
\item Suppose its Euler number were one. The rational
normal-crossing case has Euler number two by P4.8, so we
are in the simply connected resolved-conductor alternative
of P4.26. That proof gives $\pi_1(S)=\pi_1(C)=F_b$,
$b\le1$, and
$e(C)=c+1-b$ for the number $c$ of normalized rational
components. Since $e(S_{\rm reg})\ge0$, the value one
forces $c=b=1$ and $e(S_{\rm reg})=0$.
Thus $C$ has normalization $E=\mathbb P^1$, exactly one
point with two analytic branches, and only unibranch
singular points elsewhere.

\item P4.33 proves that $N$ has rational singularities and
that its image conductor scheme has
$\chi(\mathcal O_{C_{\rm sch}})=-1$.
The generic ordinary-triple conductor and pure
Cohen--Macaulay argument of P4.31 identify this scheme
with the symbolic double curve $\mathcal O_Y/P^{(2)}$.
At a point with two analytic branches, the symbolic ideal
means order at least two along each branch. The same
argument applies: a nonzero finite-support defect in the
conductor quotient would violate depth one.

\item We compute its Euler characteristic and the possible
local defects. On $E$ use the bundles $K$ and the
ramification divisor $R$ of P4.30. This time
$\deg K=2+\deg R$, since $E=\mathbb P^1$.
Write $\tau_p$ for the local length of the conormal
quotient $j_*K/(P/P^{(2)})$ and $r_p$ for the sum of
the coefficients of $R$ over $p$. Normalization and
Riemann--Roch give
\[
 \chi(\mathcal O_Y/P^{(2)})
 =5-\sum_p d_p,
 \qquad d_p=\delta(C,p)+\tau_p-r_p.
\]
We prove the precise gap needed here: at an ordinary
node $d_p=5$; at every planar singularity
$d_p=5\delta(C,p)$; and at every nonplanar singularity
having at most two branches $d_p\ge7$.

\item First consider a planar curve, with ideal $(z,h(x,y))$.
Its conormal is free over $\mathcal O_C$, generated by
$dz,dh$. On each normalized branch, write
$dh=a_i u_i$ for a primitive plane kernel generator,
and put $v_i=\operatorname{ord}a_i$.
The inclusion of each of these two free summands into
its normalized module has length $\delta(C,p)$.
The second summand has the additional vanishing orders
$v_i$. Hence
\[
 \tau_p=2\delta(C,p)+\sum_i v_i.
\]
We verify $\sum_i(v_i-r_i)=2\delta(C,p)$ directly by
curve duality. Let $A$ be the reduced plane-curve ring,
$B$ its normalization, and $I=\operatorname{Hom}_A(B,A)$
its conductor. Applying $\operatorname{Hom}_A(-,A)$ to
$0\to A\to B\to B/A\to0$ gives
$A/I\simeq\operatorname{Ext}^1_A(B/A,A)$.
Here $\operatorname{Ext}^1_A(B,A)=0$ because $B$ is
maximal Cohen--Macaulay over the one-dimensional
Gorenstein ring $A$. The dual of a finite-length module
has the same length, so $\operatorname{length}(A/I)=\delta$,
and $\operatorname{length}(B/I)=2\delta$.
If the conductor exponent on branch $i$ is $c_i$,
this says $\sum_i c_i=2\delta$.
Choose a generic plane coordinate $x$ with minimal
valuation on every normalized branch. The dualizing
generator $dx/h_y$ then pulls back with
order $r_i-v_i$. The identity
$\omega_B=I\omega_A$ therefore gives
$c_i=v_i-r_i$. Combining these equalities proves
$\tau_p-r_p=4\delta$ and the stated planar formula.

\item Now consider a nonplanar unibranch curve of multiplicity
$m$. Its embedding dimension is three. As in P4.30,
all gradient entries have valuation at least $m$, so
$\tau_p\ge2m$, while $r_p=m-1$ and
$\delta\ge m-1$. A nonplanar unibranch curve has
$m\ge3$. To see the last assertion, if $x=t^2$, write
each other coordinate as $a(x)+t b(x)$. After subtracting
the even parts, choose the nonzero $b$ of least order.
The other odd coordinate is an analytic multiple of it
over $\mathbb C\{x\}$. This is a smooth ambient equation
containing the curve, so its embedding dimension is at
most two. The case where both odd parts vanish contradicts
the birational normalization parametrization.

\item The estimates give $d_p\ge2m\ge6$. Equality at six
would require $m=3$, $\delta=2$, and $\tau_p=6$.
The valuation semigroup of a unibranch curve has exactly
$\delta$ gaps, as follows by filtering its local ring
and its normalization by order. Thus the two gaps here
are one and two, and we can choose generators with
initial orders three, four, and five. They form an
ambient coordinate system: their classes generate the
three-dimensional cotangent space, since the square
of the maximal ideal has order at least six.
In a defining relation, the quadratic monomials
$x^2$ and $xy$ have respectively the unique least
orders six and seven, so their coefficients vanish.
Every gradient with valuation three consequently has
its initial vector in the single $z$ direction
(from a possible $xz$ term). Two such initial
gradient vectors are dependent. Every two-by-two
minor therefore has valuation at least seven,
contradicting $\tau_p=6$. Hence $d_p\ge7$ in this case.

\item Finally consider a nonplanar point with two branches
of multiplicities $m_1,m_2$. Two smooth curve branches
in a smooth threefold are contained in a smooth
surface. Indeed choose a common immersive coordinate
$x$ and write them as graphs
$(y,z)=(a_i(x),b_i(x))$. Divide their difference pair
by its common power of $x$; the equation
\[
 \frac{b_1-b_2}{x^k}(y-a_1)
 -\frac{a_1-a_2}{x^k}(z-b_1)=0
\]
has a nonzero linear derivative and contains both.
Thus in the nonplanar case $m_1+m_2\ge3$.
The absence of a linear defining equation gives,
on the two branches separately,
$\tau_p\ge2(m_1+m_2)$.
The reduced-curve normalization sequence gives
$\delta(C,p)=\delta_1+\delta_2+
\operatorname{length}\mathcal O_Y/(P_1+P_2)
\ge m_1+m_2-1$.
Also $r_p=m_1+m_2-2$.
It follows that
$d_p\ge2(m_1+m_2)+1\ge7$.
This proves the gap in all cases.

\item At the unique two-branch point, a planar singularity
has defect five if its delta invariant is one, and
at least ten otherwise. A nonplanar point has defect
at least seven. Every other singular point is
unibranch and contributes at least five if planar,
and at least seven if nonplanar. Smooth points
contribute zero. Thus the sum of the defects can
be five or at least seven, but never six.
The required value $\chi(\mathcal O_{C_{\rm sch}})=-1$
would give exactly six. This contradiction excludes
$e(S)=1$.

\item The final statement uses P4.26: the sum of all the
nonnegative fibre Euler numbers is at most two.
The rational-resolution case is strictly positive,
and we have just excluded one, leaving exactly two.

\end{enumerate}
\end{proof}


\subsection{The conductor graph excludes Euler number one in the rational case}

\begin{proposition}[P4.35]\hypertarget{tag-02X8}{}\label{night:P4.35}\label{stmt:P4.35}
\StatementRecord{night:P4.35}{proved}{02X8}{\SeriesRef{night:P4.26}, \SeriesRef{night:P4.33}, \SeriesRef{night:P4.8}, \SeriesRef{night:P4R.6}}

No irreducible support with rationally resolved normalization
has Euler number one. Consequently, in a fibration whose
supports are all irreducible or normal crossing, every
irreducible singular support with rationally resolved
normalization has Euler number exactly two.

More generally, an irreducible support cannot have reduced
image conductor equal to a rational curve with exactly one
point-identification cycle when its reduced normalization
conductor has no graph cycles. This excludes that conductor
case for every ruled-resolution double-horizontal pole tree.
\end{proposition}
\begin{proof}

\begin{enumerate}[label={$\langle1\rangle$\arabic*.},leftmargin=*]
\item We first prove the graph assertion independently of
the transverse singularity type and of ramification.
Let $D\subset N$ be the reduced normalization conductor,
and suppose its point-identification graph is a tree.
Let $c$ be its number of normalized compact components.
Suppose its image $C$ is irreducible with rational
normalization and exactly one point-identification
cycle. There is then a unique point $p\in C$ with
two analytic branches; all other points are unibranch.
Denote the two preimages of $p$ on $C^\nu$ by
$p_+,p_-$. Every normalized component of $D$ maps
finitely and surjectively to $C^\nu$.

\item The finite set $\nu^{-1}(p)$ has $n$ points and
parametrizes the local irreducible hypersurface
branches of $S$ at $p$. Every pair of such branches
intersects along a curve: two nonunit local
hypersurface equations in a regular local ring of
dimension three generate an ideal of height at most
two. Their common curve lies in $C$. Consequently
the nonempty subsets of $\{+,-\}$ recording which
branches of $C$ are contained in each hypersurface
branch pairwise intersect. Such a family of subsets
of a two-element set has a common element. If a
singleton occurs its element belongs to every set;
if no singleton occurs every set is the whole set.
Interchanging the labels, every point of
$\nu^{-1}(p)$ therefore has at least one conductor
branch mapping to the plus branch of $C$.

\item Let $r_+$ and $r_-$ count the points on $D^\nu$
over $p_+$ and $p_-$ respectively. The preceding
paragraph gives $r_+\ge n$. Surjectivity of the
map from each of the $c$ normalized components
gives $r_-\ge c$. The contribution of the points
over $p$ to the point-identification graph of $D$
is exactly
\[
 \sum_{q\in\nu^{-1}(p)}(r_q-1)
       =r_++r_--n\ge c.
\]
For a connected reduced compact curve with $c$
normalized components, the cycle rank of this graph
is $1-c+\sum_q(r_q-1)$. All other contributions
are nonnegative. Its rank is therefore at least
one, contradicting the tree hypothesis.
This proves the graph assertion.

\item Now let the normalization have rational resolution
and suppose $e(S)=1$. The normal-crossing alternative
has Euler number two by P4.8. In the other alternative
P4.26 gives a connected simply connected resolved
pole support $A\subset T$, and
\[
 \pi_1(S)=\pi_1(C)=F_b,\qquad b\le1,
 \qquad e(C)=c_C+1-b.
\]
Its normalized components are rational.
Since $e(S_{\rm reg})\ge0$, Euler number one forces
$c_C=b=1$ and $e(S_{\rm reg})=0$.
It remains to verify the tree property for the
actual conductor $D\subset N$, allowing singular $N$.

\item P4.33 makes every singularity of $N$ rational.
Every connected exceptional fibre of $T\to N$
is simply connected. Here is the precise curve
calculation. Over a small Stein neighbourhood,
$R^1\rho_{N*}\mathcal O_T=0$ gives
$H^1(\mathcal O_T)=0$ on its inverse image.
This proper map has every fibre of dimension at
most one over that Stein neighbourhood. Its coherent
cohomological dimension is consequently at most
one, so the second higher direct image of every
coherent sheaf vanishes. Applied to the ideal of the reduced
exceptional curve $E$, Leray and the Stein property
give $H^2(\mathcal O_T(-E))=0$ there. Its ideal
sequence consequently gives $H^1(E,\mathcal O_E)=0$.
The curve normalization exact sequence now implies
that all normalized components are rational and
that their point-identification graph is a tree.
Thus $E$ is simply connected.

\item By P4R.6, $A$ is the full inverse image of $D$.
The proper map $A\to D$ is the quotient which
collapses these connected exceptional curves to
points and is an isomorphism elsewhere.
Compatible triangulation makes these curve
inclusions cofibrations. Collapsing a connected
simply connected subcomplex preserves fundamental
groups by van Kampen. Collapsing the finitely many
exceptional fibres therefore gives $\pi_1(D)=0$.
Each normalized component of $D$ is rational,
being birational to a strict pole component in
$T$. Hence its point-identification graph is a
tree. The already proved graph assertion contradicts
the description of $C$. This excludes $e(S)=1$.

\item For the final more general application, the
resolved pole support in the ruled double-section
case is $H$ with rational trees attached.
Every exceptional fibre is connected and is a
union of its components. If it does not contain
$H$, it is a connected rational subtree. If it
contains $H$, all the remaining components form
rational trees attached to that component once.
Contracting such a connected subcurve cannot
create a cycle in the graph: contraction of a
connected subgraph of a tree again has a tree
as quotient. The only positive-genus vertex,
if contracted, is replaced by a point, and this
also introduces no graph edge. Thus the image
conductor $D$ has no graph cycles, and the same
contradiction applies.

\item Finally, for a fibration in the two requested
support classes, P4.26 makes the total sum of
nonnegative fibre Euler numbers at most two.
The rational-resolution case is strictly positive;
excluding one leaves exactly two.

\end{enumerate}
\end{proof}


\subsection{Eliminating all positive-Euler Hopf resolutions}

\begin{proposition}[P4.36]\hypertarget{tag-02X9}{}\label{night:P4.36}\label{stmt:P4.36}
\StatementRecord{night:P4.36}{proved}{02X9}{\SeriesRef{night:P4.24}, \SeriesRef{night:P4R.6}}

Every irreducible singular support in Part~4.2 whose
normalization has a Hopf minimal model has smooth minimal
Hopf normalization and Euler number zero. The assertion
does not assume initially that the normalization is smooth
or that its Euler number vanishes.
\end{proposition}
\begin{proof}

\begin{enumerate}[label={$\langle1\rangle$\arabic*.},leftmargin=*]
\item Let $T\to N\to S\subset Y$ be the minimal normalization
resolution, and let $\pi:T\to T_0$ contract its $k$ blowups
to a minimal Hopf surface. We first take a finite unramified
primary Hopf cover $T'_0\to T_0$ and pull back the actual
sequence of point blowups. This gives a finite unramified
map $T'\to T$ and a blowup morphism $T'\to T'_0$.
Compose with $T\to Y$. All Chern calculations below take
place on this explicit smooth cover; no cover of the
singular normalization $N$ is asserted. To avoid a second
set of symbols we write $T,T_0,k,\rho$ for these covered
objects until the last paragraph.

\item The primary Hopf normal form and proof of P4.24 give a
nowhere-zero holomorphic foliation on $T_0$ tangent to
every compact curve. Its saturated pullback to $T$ is a
line subbundle away from finitely many simple saddle
zeros. It is tangent to every strict conductor component
and every created exceptional curve. The actual conductor
centres lie on these invariant curves. Every blowup centre
is either a regular point of this foliation or one of its
saddles. Let $\varepsilon=0$ in the first case and
$\varepsilon=1$ in the second. The same explicit two-chart
calculation as in P4.24 gives
\[
 c_1(L)=-K_T+V,\qquad
 V'=\mu^*V+\varepsilon E,
 \qquad K_T^2=-k,\qquad K_TV-V^2=0.
\]
Topologically trivial line-bundle factors from $T_0$ do
not change these integral intersection calculations.

\item Put $B=\rho^*TY$. Its first and second rational Chern
classes vanish: $b_2(Y)=0$ gives $c_1(TY)=0$ rationally,
and Poincar\'e duality gives $H^4(Y;\mathbb Q)=0$ and
hence $c_2(TY)=0$. The derivative $L\to B$ can now have
a divisorial zero, since $\rho$ contracts exceptional
curves. Denote that effective integral divisor by $Z$.
Dividing by it gives the saturated line injection in
codimension one
\[
 L'=L(Z)\longrightarrow B,
 \qquad c_1(L')=-K_T+V+Z.
\]
The quotient $Q=B/L'$ is torsionfree of rank two.
Write $z_0\ge0$ for the length of $Q^{**}/Q$.
The double dual is locally free on the smooth surface,
and the Chern sequence gives
\[
 c_2(Q^{**})=(-K_T+V+Z)^2-z_0.
\]

\item Let $E_B\subset B$ be the saturated rank-two image
of $d\rho:TT\to B$. Its determinant is
$\det E_B=\omega_T^{-1}(R)$, where the effective
integral divisor $R$ is the divisorial zero of the
determinant map. Since $E_B$ is saturated in $B$,
the divided foliation line $L'$ is contained in it.
The image of $E_B/L'$ in $Q^{**}$ is saturated in
codimension one. Its rank-one double dual
$J=(E_B/L')^{**}$ is a line bundle. The map into
$Q^{**}$ extends across the finitely many omitted
points by Hartogs, since both sheaves are locally
free there. Its saturation is therefore a line bundle
of first Chern class $R-V-Z$; the rank-one quotient
is an ideal sheaf of a zero-dimensional scheme
times a line bundle of first Chern class $K_T-R$.
Let its length be $z_1\ge0$. Thus
\[
 (-K_T+V+Z)^2-z_0
       =(R-V-Z)(K_T-R)+z_1.
\]
Using $K_T^2=-k$ and $K_TV-V^2=0$, this is precisely
\[
 \mathcal B(R,Z)
 :=R^2-K_TR-VR-ZR-K_TZ+2VZ+Z^2
 =k+z_0+z_1.                                      \tag{P4.36.1}
\]

\item We now evaluate the left side through the actual
point blowups, retaining all integral coefficients.
Push $R$ and $Z$ down along the chosen sequence.
At one reversed blowup write
\[
 R'=\mu^*R+cE,\quad Z'=\mu^*Z+dE,
 \quad K'=\mu^*K+E,\quad V'=\mu^*V+\varepsilon E.
\]
The integers $c,d$ are not assumed positive.
Since $E^2=-1$ and a pulled-back divisor has zero
intersection with $E$, direct expansion gives
\[
 \mathcal B(R',Z')-\mathcal B(R,Z)
 =-c^2+(1+\varepsilon)c+cd+(1-2\varepsilon)d-d^2.
\]
Set $a=c-1$ and $b=d-(1-\varepsilon)$. For each
$\varepsilon\in\{0,1\}$ this is exactly
\[
 1-(a^2+b^2-ab)\le1.                            \tag{P4.36.2}
\]
The quadratic form is
$\frac12(a^2+b^2+(a-b)^2)$; it is nonnegative
and is zero exactly when $a=b=0$.
All divisor intersections on the initial Hopf
surface are zero because $b_2(T_0)=0$.
Summing (P4.36.2) therefore gives
$\mathcal B(R,Z)\le k$.
Equation (P4.36.1) forces equality, $z_0=z_1=0$,
and at every single blowup
\[
 c=1,\qquad d=1-\varepsilon.                    \tag{P4.36.3}
\]

\item Every created exceptional curve has positive
coefficient in $Z$. At a regular centre, (P4.36.3)
adds one to its nonnegative pulled-back coefficient.
At a saddle centre it adds zero, but at least one
previously created exceptional curve passes through
that saddle. The original Hopf foliation has no
zero; the two-chart blowup calculation places each
new saddle on the new exceptional curve. Induction
therefore gives a positive pulled-back coefficient
at every later saddle blowup as well.

\item The foliation is tangent to each exceptional curve.
If its image under $\rho$ were a curve, the derivative
along its tangent line would be nonzero at a generic
point: a nonconstant holomorphic map of smooth
normalized curves in characteristic zero is
generically immersive. Positive divisorial order
of $d\rho|_L$ contradicts that fact. Thus every
created exceptional curve is contracted by $\rho$,
and hence by the normalization resolution, since
the normalization map is finite.

\item Return to the original, possibly secondary Hopf
surface. Every exceptional curve downstairs has
an exceptional lift in the finite unramified cover;
the preceding argument makes its image a point.
Thus all exceptional curves of the original
$\pi:T\to T_0$ would be contracted by the minimal
resolution $T\to N$. If $k>0$, the exceptional curve
of its last blowup is a smooth rational curve of
self-intersection $-1$, which a minimal resolution
cannot contract. Hence $k=0$.
There is no further exceptional contraction from
a minimal Hopf surface: every compact curve has
self-intersection zero by $b_2=0$, whereas a
nonempty exceptional configuration has negative
definite intersection matrix, as proved in P4R.6.
Therefore $T=N=T_0$ is a smooth minimal Hopf
surface. P4.24 gives $e(S)=0$ and its complete
elliptic-conductor description.

\end{enumerate}
\end{proof}


\subsection{Excluding the remaining class-VII minimal models}

\begin{proposition}[P4.37]\hypertarget{tag-02XA}{}\label{night:P4.37}\label{stmt:P4.37}
\StatementRecord{night:P4.37}{proved}{02XA}{\SeriesRef{night:P4.23}, \SeriesRef{night:P4.26}, \SeriesRef{night:P4.36}}

No irreducible singular support in Part~4.2 has normalization
whose minimal resolution has a class-VII minimal model with
positive second Betti number. Consequently every class-VII
normalization occurring here is a smooth minimal Hopf surface
and has Euler number zero.
\end{proposition}
\begin{proof}

\begin{enumerate}[label={$\langle1\rangle$\arabic*.},leftmargin=*]
\item The numerically anticanonical classification used in P4.23
and P4.26 reduces this question to an intermediate Kato
minimal model $T_0$: the Inoue--Hirzebruch and parabolic
alternatives were excluded there by their reduced
normal-crossing conductor and the Chern identity.
Write $n=b_2(T_0)>0$, let $A_0$ be its maximal reduced
curve divisor, and let $D_0$ be the effective integral
pole divisor pushed down from the normalization resolution.
Thus $K_{T_0}+D_0$ is numerically zero. The support of
$D_0$ is $A_0$.

\item We state the precise additional primary input.
Favre, \emph{Dynamique des applications rationnelles},
Theorem~1.2.24(1)--(2), printed p.~35, says that the
$n$ rational curve classes span $H^2(T_0;\mathbb Q)$
and that their intersection form is negative definite.
Theorem~1.2.30, printed p.~42, and its proof on
pp.~44 and~46 give a holomorphic foliation with conormal
line
\[
 N_{\mathcal F}^*=\mathcal O_{T_0}(-A_0)
                       \otimes\text{a flat line}.
\]
Theorem~1.2.32 and its proof on p.~46 show that every
component of $A_0$ is invariant and that its singularities
are precisely the intersection points, with local
first integrals $z^\alpha w^\beta$, $\alpha,\beta>0$.
The invariant-curve and order-one assertions are also
Dloussky--Oeljeklaus, Theorem~3.1, pp.~1528--1529.
The actual source URLs and exact theorem locators are
recorded in the bibliography and literature log.

\item It follows from the determinant of tangent and normal
lines that the foliation tangent line has numerical
first Chern class $-K_{T_0}-A_0$. The local generator
at an intersection is
$\beta z\partial_z-\alpha w\partial_w$.
Blowing up a regular point or one of these saddles
again gives only regular points and saddles with
opposite-sign nonzero eigenvalues. Thus every actual
blowup centre has the same two possibilities
$\varepsilon=0,1$ as in P4.36, with the same formulas.

\item Let $T\to T_0$ contain $k$ blowups. Write
$c_1(L)=-K_T+V$ for the saturated foliation line;
initially $V_0=-A_0$, and
$V'=\mu^*V+\varepsilon E$ at a blowup.
The reduced divisor $A_0$ is a rational cycle with
rational trees attached. Its normalization has
$n$ rational components and $n$ nodes, including
self-nodes, so $\chi(\mathcal O_{A_0})=0$.
Adjunction, or the divisor exact sequence on $T_0$,
therefore gives
\[
 (K_{T_0}+A_0)A_0=0.
\]
Hence $K_{T_0}V_0-V_0^2=0$.
Under each blowup this quantity changes by
$\varepsilon^2-\varepsilon=0$.
Since $K_{T_0}^2=-n$, we obtain
\[
 K_TV-V^2=0,\qquad K_T^2=-n-k.
\]

\item Define $B=\rho^*TY$, the divisorial determinant zero
$R$, the divisorial foliation zero $Z$, and the
nonnegative lengths $z_0,z_1$ exactly as in P4.36.
Its saturated sheaf calculation, including the
Hartogs extension of the rank-one double dual,
now gives
\[
 \mathcal B(R,Z)
 =R^2-K_TR-VR-ZR-K_TZ+2VZ+Z^2
 =n+k+z_0+z_1.                                  \tag{P4.37.1}
\]
Every reversed blowup changes $\mathcal B$ by at
most one, by the unrestricted integer identity
(P4.36.2). It remains to evaluate the initial term.

\item Push $R,Z$ down to effective divisors $R_0,Z_0$
on $T_0$. Their supports are contained in $A_0$:
the derivative is generically injective away from
the resolved conductor, and in any event all
compact curves of the intermediate surface are
components of its maximal divisor. Put $D=D_0$,
$A=A_0$. Using $K=-D$ numerically and $V=-A$ gives
\[
 \mathcal B_0
 =R_0^2+DR_0+AR_0-Z_0R_0+DZ_0-2AZ_0+Z_0^2.
\]
For divisor classes define
$q(X,Y)=X^2-XY+Y^2$.
Adjunction above gives $DA=A^2$, while $D^2=-n$.
Expansion consequently gives the exact completion
\[
 \mathcal B_0
       =n+q(R_0+D,Z_0+D-A).                     \tag{P4.37.2}
\]
The form on divisor classes is negative definite.
Moreover
\[
 q(X,Y)=\tfrac12\bigl(X^2+Y^2+(X-Y)^2\bigr),
\]
so it is nonpositive, and it vanishes only for
$X=Y=0$. The first argument $R_0+D$ is nonzero:
it is a nonzero effective combination of the
$n$ independent curve classes. Thus
$\mathcal B_0<n$.
After the $k$ blowups we still have
$\mathcal B(R,Z)<n+k$, contradicting (P4.37.1).
This excludes the intermediate minimal model.

\item All class-VII cases with positive second Betti
number have now been excluded. The zero-Betti
Inoue case has no conductor curves and was
excluded in P4.26. P4.36 treats the remaining
Hopf case and proves the final assertion.

\end{enumerate}
\end{proof}


\subsection{Two exact local exclusions for ordinary-double curves}

\begin{proposition}[P4.38]\hypertarget{tag-02XB}{}\label{night:P4.38}\label{stmt:P4.38}
\StatementRecord{night:P4.38}{proved}{02XB}{\SeriesRef{night:P4.16}, \SeriesRef{night:P4.20}}

An actual irreducible support in Part~4.2 cannot have a singular
planar reduced critical curve at a point at which every generic
transverse germ is an ordinary double point. Nor can its reduced
critical curve there be the transverse union of a plane cusp
and a smooth branch, again with ordinary-double transverse germs
on both components. The latter assertion also holds with the
cusp replaced by a plane tacnode.
These exclusions allow singular normalization.
\end{proposition}
\begin{proof}

\begin{enumerate}[label={$\langle1\rangle$\arabic*.},leftmargin=*]
\item We retain the boundary transport of P4.20 throughout. At a fixed
compact boundary, ordinary-double transverse germs form an open
condition. A sufficiently small deformation retaining that
condition gives the same nearby nonzero-level boundary
submersion. The finite-critical-value argument in that proof
then transports any compact vanishing sphere from the deformed
germ into the preceding nearby fibre. When a coordinate change
or a unit requires a smaller ball, we perform this construction
on that ball first, and then transport through the earlier
larger ball. Thus every sphere constructed below reaches the
original nearby four-torus.

\item First suppose that the reduced critical curve has ideal
$P=(z,h(x,y))$, where $h$ is reduced and singular at the
origin. The equation $f$ lies in $P^{(2)}$: its generic
transverse ordinary double point vanishes to order two on
each component. The complete-intersection square $P^2$ is
unmixed, so $P^{(2)}=P^2$. We therefore have the exact expression
\[
 f=A z^2+Bzh+Ch^2.
\]
Adding a small constant multiple of $z^2$ preserves the critical
curve and the ordinary-double boundary condition. We choose it
so that the $z^2$ coefficient at the origin is nonzero. The
Morse lemma in $z$, with parameters $x,y$, then gives
\[
 f=Z^2+h(x,y)^2u(x,y).
\]
Indeed the critical $z$ coordinate vanishes on $h=0$ and
is divisible by $h$, so its substitution has the displayed
square factor. Add a small constant to $u$ and shrink the
ball so that $u$ is a unit. Absorbing its square root in $h$
gives $Z^2+H(x,y)^2$, with the same singular reduced zero
curve. A reduced singular plane curve has an isolated critical
point. A generic small morsification $H_s$ has a Morse critical
point, and a small additive constant makes all these critical
values nonzero. At any such point $Z^2+H_s^2$ has a
nondegenerate isolated critical point: its plane Hessian is
$2H_s\operatorname{Hess}H_s$, and its $Z$ Hessian is $2$.
Its vanishing sphere has self-intersection $-2$.
The moving curve $H_s=0$ remains transverse to the chosen
boundary, with nondegenerate normal quadratic form there.
The preceding boundary transport therefore contradicts
$\pi_2(T^4)=0$.

\item For the second assertion write
\[
 h=y^2-x^k,\qquad P=(xz,yz,h),\qquad k=3\text{ or }4.
\]
This is the union of the plane curve $z=h=0$ and the
transverse line $x=y=0$. Its symbolic square is
\[
 P^{(2)}=(z,h)^2\cap(x,y)^2
  =(x^2z^2,xyz^2,y^2z^2,zh,h^2).
\]
For $k=4$, the first equality uses the two smooth branches of
$h=0$ separately; their complete-intersection square is unmixed
and hence is still $(z,h)^2$. To verify the intersection
formula directly, write an element of $(z,h)^2$ as
$az^2+bzh+ch^2$. The last two terms belong to $(x,y)^2$,
and $z$ is a nonzero divisor modulo this ideal. The first term
belongs to it exactly when $a\in(x,y)^2$.
Consequently
\[
 f=Azh+z^2(Bx^2+Cxy+Dy^2)+Gh^2.
\]
Small constant changes in $A,B$ retain the boundary condition
and make $A(0)B(0)\ne0$. Replace $h$ in this expression by
$h_s=y^2-x^k-sx^2$. The new critical curve contains
$z=h_s=0$ and $x=y=0$, and its boundary ordinary-double
germs persist. On the latter axis the transverse quadratic
matrix has coefficients
\[
 a=-sAz+Bz^2,\qquad b=Cz^2,\qquad c=Az+Dz^2,
\]
with off-diagonal entry $b/2$. Here the four coefficient
functions are restricted to that axis. Its determinant
divided by $z^2$ is
\[
 F(s,z)=(-sA+Bz)(A+Dz)-C^2z^2/4.
\]
At $(s,z)=(0,0)$ we have $F_z=AB\ne0$ and
$F_s=-A^2$. The implicit-function theorem gives a simple root
$z(s)=sA(0)/B(0)+O(s^2)$, nonzero for small nonzero $s$.
There the $y^2$ coefficient is nonzero, the Hessian has rank
one, and its determinant has a simple zero in the axis
direction. The parametrized Morse lemma gives
$Y^2+x^2H(x,z)$, with $H_z\ne0$. The exact coordinate
change $Z=H(x,z)$ yields $Y^2+x^2Z$.
P4.16 gives its compact spherical class of square $-4$.
Nested boundary transport once more places this class in
the original nearby four-torus, a contradiction.

\end{enumerate}
\end{proof}

The exact quadratic matrix, determinant, and the two nonzero
derivatives in this proof are verified by
\texttt{code/part4/critical\_curve\_node.py}; its saved output
reports PASS.


\subsection{A coefficient-two exceptional curve over an ordinary-triple germ}

\begin{proposition}[P4.39]\hypertarget{tag-02XC}{}\label{night:P4.39}\label{stmt:P4.39}
\StatementRecord{night:P4.39}{proved}{02XC}{\SeriesRef{night:P4R.6}}

Let $S=(f=0)$ be a reduced hypersurface germ in a smooth
threefold, with one-dimensional reduced critical locus $C$.
Suppose that every generic transverse germ along $C$ is an
ordinary triple point. On the minimal resolution $T$ of its
normalization, let $D$ be the integral total pole divisor
defined by the residue of an ambient volume form divided by
$f$. No exceptional curve of positive genus has coefficient
two in $D$.

In particular, the contracted horizontal section of a
positive-genus ruled resolution in Part~4.2 cannot lie over
a point where all local critical-curve branches have generic
type $D_4$. This assertion does not replace that local
hypothesis by a mixture of transverse types.
\end{proposition}
\begin{proof}

\begin{enumerate}[label={$\langle1\rangle$\arabic*.},leftmargin=*]
\item We blow up the ambient origin, write $S_1$ for the strict
transform and $N_1$ for its normalization, and take the
minimal resolution $T_1$ of $N_1$. The proper map $N_1\to S$
lifts to the normalization $N$ by its universal property.
Thus $T_1\to N$ is a resolution and dominates $T$ by a
sequence of point blowups. The strict transform of a proposed
positive-genus exceptional curve $H$ remains present and has
the same pole coefficient two.

\item Put $d=\operatorname{mult}_0f$, and denote the ambient
exceptional divisor by $E$. Actual adjunction gives
\[
 \omega_{S_1}=\pi^*\omega_S\otimes
        \mathcal O((2-d)E)|_{S_1}.
\]
Let $D_1$ be the total conductor divisor on $T_1$ relative
to $\omega_{S_1}$, and let $e$ be the Cartier pullback of
$E$. With the specified residue comparison, the total pole
of the original germ on $T_1$ is
\[
                    D_1+(d-2)e.                       \tag{4.39a}
\]
The divisor $D_1$ is effective. Its strict conductor
coefficients are positive integers. For exceptional curves
of the minimal resolution of $N_1$, the negative-matrix
argument of P4R.6 applies locally: $K_{T_1}F\ge0$, since
there is no smooth rational exceptional $(-1)$-curve;
adding the strict conductor retains nonnegative intersection.
The inverse exceptional intersection matrix has nonpositive
entries, so all exceptional coefficients are nonnegative.

\item If $H$ is exceptional over $N_1$, its coefficient in $D_1$
is strictly positive. A zero coefficient would give
\[
 0=(K_{T_1}+D_1)H
   =2p_a(H)-2-H^2+D_1H>0.
\]
If that coefficient is one, adjunction gives
\[
 0=2p_a(H)-2+(D_1-H)H.
\]
The normalization of $H$ has positive genus, so $p_a(H)\ge g(H^\nu)\ge1$. Both terms in the last equality are nonnegative, hence $p_a(H)=1$ and $(D_1-H)H=0$. The normalization formula $p_a(H)=g(H^\nu)+\sum\delta(H,p)$ now gives $g(H^\nu)=1$ and zero delta invariant at every point. Thus $H$ is smooth elliptic and is disjoint from every other positive component of $D_1$. A zero-coefficient exceptional
curve cannot meet a positive one, by its own intersection
equation and $K_{T_1}F\ge0$. The connected exceptional
fibre is consequently $H$ alone and meets no strict conductor.
The entire germ of $S_1$ at its image is normal with an
isolated simple elliptic singularity. Indeed, any additional
hypersurface branch would intersect this branch in a curve
and would give strict conductor meeting $H$.

\item Set $n=\operatorname{ord}_H(e)>0$ and
$a=\operatorname{coeff}_H(D_1)\ge0$. Formula (4.39a)
gives the exact identity
\[
                       2=(d-2)n+a.                    \tag{4.39b}
\]
Thus $d\ge5$ is impossible. For $d=4$ we have
$n=1,a=0$, and $H$ is not exceptional over $N_1$.
At its generic point the normalization is an isomorphism,
so $H$ is birational to a reduced component of the projective
quartic tangent cone. Every second derivative of $f$ vanishes
on $C$. Its leading term along any analytic arc of $C$
therefore makes that quartic have a point of multiplicity
at least three. An irreducible quartic has arithmetic genus
three and loses at least three there. A smooth elliptic cubic
component together with the remaining line could have
multiplicity at most two at the point. All components of
degree at most two are rational. Thus every reduced tangent
component is rational, excluding $H$ for $d=4$.

\item We are left with $d=3$ and
\[
                   (n,a)=(2,0)\quad\hbox{or}\quad(1,1).
\]
Choose a coordinate $u$ with nonzero pure cubic initial
coefficient, apply Weierstrass preparation, and depress the
monic cubic. Multiplication by the preparation unit does not
affect Hessian vanishing on $C$, because $f=df=0$ there.
We obtain
\[
 f=\text{unit}\cdot\bigl(u^3+b(x,z)u+c(x,z)\bigr).
\]
The second $u$ derivative is $6u$, so $C$ lies in the smooth
surface $u=0$. Write its reduced ideal as $(u,h(x,z))$.
The complete-intersection cube is unmixed and equals its
symbolic cube. Generic order three along $C$ and coefficient
comparison therefore give the exact expression
\[
 f=\text{unit}\cdot
       \bigl(u^3+h^2B(x,z)u+h^3C(x,z)\bigr).           \tag{4.39c}
\]
If $C$ is singular, $\operatorname{ord}h\ge2$ and the
tangent cone is the triple plane $u^3$. The degree-four
term restricts to zero on this plane. After the ambient
blowup the entire reduced exceptional line is critical,
so none of its points is an isolated normal singularity.
If $C$ is smooth, take $x=h$. Every homogeneous term $f_j$
then has degree at least three in $u,x$. The cubic tangent
cone is a binary cubic in these two variables, hence three
lines, counted with multiplicities, through the tangent
direction of $C$.

\item First suppose that $H$ dominates a tangent-cone line.
For $(n,a)=(2,0)$ the normalization is generically an
isomorphism, and $H$ is birational to that line. For
$(n,a)=(1,1)$, a transverse normalization branch has conductor
exponent one. It is smooth and meets exactly one other smooth
branch transversely: the intrinsic conductor exponent of a
singular plane branch is at least two, and each other branch
contributes a positive intersection multiplicity. Thus the
generic transverse germ of $S_1$ is a node. A multiplicity-one
tangent line is generically smooth and cannot have $a=1$.
For a multiplicity-three line, the two geometric branches of
that node have exceptional orders whose weighted sum is three.
A connected double cover would give equal orders and an even
sum. Thus the order-one branch has degree one over the line
and is rational.

\item The remaining possibility is a double line together with a
distinct simple line. Here $C$ is smooth by (4.39c). Make a
linear change so that $f_3=xu^2$. On the $x$ chart of the
blowup put $u=xv,z=xw$. The strict equation is
\[
 v^2+x f_4(1,v,w)+x^2f_5(1,v,w)+\cdots.
\]
An ordinary double curve on $x=v=0$ requires
$f_4(1,0,w)=0$ identically. Its transverse quadratic form is
\[
 v^2+xvB(w)+x^2C(w),
 \qquad \deg B\le1,\quad\deg C\le2.                  \tag{4.39d}
\]
The degree bounds follow from the degree-at-least-three
condition in $u,x$ on every $f_j$: here
$B=\partial_v f_4(1,0,w)$ and $C=f_5(1,0,w)$.
The normalization over the exceptional line adjoins a root
of $q^2+B(w)q+C(w)$. Its discriminant has degree at most
two. The resulting quadratic extension of $\mathbb C(w)$
is branched at at most two points, including infinity.
Riemann--Hurwitz makes its compact normalization rational;
the components of a split cover are rational as well.

\item Finally suppose that $H$ is exceptional over a point of $N_1$.
We must have $n=a=1$, and the first part of the proof would
make $S_1$ normal with an isolated simple elliptic singularity
there. This cannot happen. For singular $C$ the whole
exceptional line is critical, as already proved. For smooth
$C$ its strict transform passes through the common point of
the tangent lines, so that point is not an isolated normal
singularity. Away from it, at a line of multiplicity
$e=2$ or $3$, divide by the nonzero factor from any other
line. The strict equation has the form
\[
                         v^e+xG(v,w,x),
\]
where $G(0,w,0)$ is affine linear in $w$, by the same
normal-degree bound. If this function is nonzero at a point,
the surface is smooth. If it is identically zero, the entire
line is critical. At a simple zero, $G_w$ is a unit and the
exact coordinate change $W=G(v,w,x)$ gives $v^e+xW$, a
rational double point of type $A_{e-1}$. A multiplicity-one
line is smooth away from intersections. There is therefore
no isolated simple elliptic singularity. All possibilities
for $H$ have been excluded.

\end{enumerate}
\end{proof}

The positive-integer solutions of (4.39b), the monomial degree
bounds in (4.39d), and the forced fourth-order factor of the
homogeneous quadratic discriminant are checked in
\texttt{code/part4/d4\_pole.py}. Its saved output reports PASS.


\subsection{The excess-one curve calculation and the ordinary-double sector}

\begin{proposition}[P4.40]\hypertarget{tag-02XE}{}\label{night:P4.40}\label{stmt:P4.40}
\StatementRecord{night:P4.40}{proved}{02XE}{\SeriesRef{night:P4.33}, \SeriesRef{night:P4.35}, \SeriesRef{night:P4.38}, \SeriesRef{night:P4R.4}}

In the scope of Part~4.2 where every fibre support is
irreducible or normal crossing, no irreducible support with
positive-genus ruled resolution has generic transverse type
$A_1$ on every component of its critical curve.
\end{proposition}
\begin{proof}

\begin{enumerate}[label={$\langle1\rangle$\arabic*.},leftmargin=*]
\item We first give the exact local curve calculation needed in
the proof. Let $A$ be a reduced complex curve local ring with
$r$ branches, let $B=\bigoplus_{i=1}^r\mathbb C\{t_i\}$
be its normalization, and let $R\subset B$ be the ring whose
constant terms agree. Then
\[
 \operatorname{length}(R/A)=\delta(A)-(r-1).
\]
Suppose this length is one. The simple $A$-module $R/A$
is annihilated by $\mathfrak m_A$, so
$\mathfrak m_A R\subset A$. Choose
$v\in\mathfrak m_R$ spanning $R/A$ and write
$v^2=a+cv$, where $a\in\mathfrak m_A$ and $c\in\mathbb C$.
If $c\ne0$, the convergent square-root series solves
$b^2-cb-a=0$ with $b\in\mathfrak m_A$. Then
$(v-b)(v+b-c)=0$ and the second factor is a unit, forcing
$v=b$, a contradiction. Thus $v^2\in A$ and
\[
                         \mathfrak m_R^2\subset A.
\]
Consequently $A$ is exactly the inverse image of a hyperplane
\[
                         \sum_i a_i f_i'(0)=0
\]
in $\mathfrak m_R/\mathfrak m_R^2=\mathbb C^r$.
Rescaling the parameters with nonzero $a_i$ gives actual
analytic ring isomorphisms with the following models.

\item With one nonzero coefficient, that branch has ring
$\mathbb C\{t^2,t^3\}$ and the others are independent smooth
branches joined at their constant terms. The embedding
dimension is $r+1$: for $r=1$ this is a plane cusp, for
$r=2$ a cusp and transverse line, and for $r=3$ it does
not embed in a smooth threefold. With two nonzero
coefficients, the pair has equal first derivatives. Its
ring is generated by $(t,t)$ and $(t^2,-t^2)$, giving the
plane tacnode $y^2=x^4$. The other $r-2$ branches are
independent transverse lines. For $r=3$ and three nonzero
coefficients, the first jets are the restrictions of linear
forms to three distinct lines in a plane. Their products
span all three second-jet coordinates, so the ring is
exactly the ordinary plane triple point. Equivalently,
the latter ring has precisely the same constant and
first-derivative conditions and colength three in $B$.
These cases exhaust $r\le3$. The analytic ring
isomorphisms lift to ambient coordinate changes: choose
minimal algebra generators, lift their independent linear
parts, and apply the inverse-function theorem. In embedding
dimension two we first straighten the smooth containing
surface. Every realizable model is therefore one of the
local classes excluded in P4.38 for ordinary-double
transverse type.

\item Now suppose the stated ruled support exists. Its horizontal
pole section has coefficient two. It cannot map
nonconstantly to the conductor, since an ordinary double
point has conductor exponent one on each normalization
branch. The section is contracted. P4.33 gives
$\chi(\mathcal O_{C_{\rm sch}})=-1$.
The conductor image is Cohen--Macaulay and generically
reduced in type $A_1$. It is therefore reduced, so
$C_{\rm sch}=C$ and $\chi(\mathcal O_C)=-1$.
All normalized components of $C$ are rational.

\item The already proved group and Euler calculations leave
exactly two normalized rational components and one
point-identification cycle. Indeed the rational-tree case
has Euler number at least three by P4R.4; the
one-component cycle is excluded by P4.35; and the global
nonnegative Euler sum is at most two. The normalization
sequence and the graph formula consequently give
\[
 \sum_p\delta_p=3,\qquad
 \sum_p(r_p-1)=2,
 \qquad
 \sum_p\bigl(\delta_p-(r_p-1)\bigr)=1.
\]
Every summand in the last sum is nonnegative. There is
exactly one excess-one point and every point has at most
three local branches. The local calculation classifies that
point by an actual analytic isomorphism, and P4.38 excludes
each possible model. This is the required contradiction.

\end{enumerate}
\end{proof}

The first-jet cases, the second-jet product determinant, and
the global excess computation are verified in
\texttt{code/part4/curve\_excess.py}; its saved output reports
PASS. The proof retains the ordinary-double boundary
condition and makes no substitution of transverse type
$A_2$ for $A_1$.


\subsection{The ordinary-node defect in transverse type \texorpdfstring{$A_3$}{A3}}

\begin{proposition}[P4.41]\hypertarget{tag-02XF}{}\label{night:P4.41}\label{stmt:P4.41}
\StatementRecord{night:P4.41}{proved}{02XF}{\SeriesRef{night:P4.28}, \SeriesRef{night:P4R.10}, \SeriesRef{night:P4R.14}}

At a point of an actual support in Part~4.2, suppose the
entire reduced critical curve is the ordinary node
$C=\{z=xy=0\}$ and both branches have generic transverse
type $A_3$. Let $d_i,b_i,q$ be the Hessian zero orders,
projected conormal colengths, and normalization descent
defect of P4R.10. Then
\[
 \delta(C)+\sum_{i=1}^2\left(\frac{d_i}{2}+b_i\right)+q
                         \ge\frac92.                 \tag{4.41a}
\]
Here $R=0$ and $\delta(C)=1$. Equality makes the normalization
of the hypersurface smooth over the point.
\end{proposition}
\begin{proof}

\begin{enumerate}[label={$\langle1\rangle$\arabic*.},leftmargin=*]
\item Put $P=(z,xy)$, and let $I$ be the ambient ideal of the
scheme-theoretic conductor image. We have $P^2=P^{(2)}\subset I$,
and $J=P/I$ is torsionfree of rank one on both branches.
After dividing its normalized image by the two branchwise
vanishing orders, it is an $A=\mathcal O_C$ submodule of
$B=\mathbb C\{x\}\oplus\mathbb C\{y\}$ generating both
summands. The conductor of $A$ in $B$ is $\mathfrak m_A$,
so $\mathfrak m_AB\subset J$. Thus $J$ is determined by
a subspace of $\mathbb C^2$ projecting onto both coordinates.
If that subspace is a line, rescaling gives $J\simeq A$ and
$q=1$. If it is the whole space, $J\simeq B$ and $q=0$.
These are actual module isomorphisms. The divided orders are
the $b_i$ already in the formula.

\item We use the ring condition of the actual conductor:
\[
 \operatorname{Hom}_{\mathcal O_S}(I,I)=
 \operatorname{Hom}_{\mathcal O_S}(I,\mathcal O_S).
                                                               \tag{4.41b}
\]
The right side is the normalization algebra by finite duality,
and $I$ is an ideal of that algebra. For the determinant
presentation of a Cohen--Macaulay codimension-two conductor
inside a Gorenstein hypersurface, this condition is equivalent
to membership of the adjugate entries in $I$. This is precisely
de Jong--van Straten, Theorem~1.12, printed p.~532
\cite{DeJongVanStraten1990Normalization}. Its hypotheses hold
here: the ambient space is regular, the reduced support is a
hypersurface, and its conductor quotient is Cohen--Macaulay.
We give all the matrices used below.

\item Suppose first that $q=1$. The surjection
$P/P^2\to J$ is a surjection $A^2\to A$, with primitive
kernel generator $g$. Taking a complementary generator $k$
of $P$ gives $I=(g,k^2)$. A determinant presentation
$f=ag+bk^2$ has two-by-two adjugate containing $a,b,g,k^2$.
Condition (4.41b) says $a,b\in I$, so the exact equation is
\[
                   f=Ag^2+Bgk^2+Ck^4.                \tag{4.41c}
\]
We claim that $A(0)$ cannot be a unit. If $g_z(0)\ne0$,
straighten $g$, take $k$ to be a nodal equation in $g=0$,
and apply the parameter Morse lemma in $g$. The residual
is $k^4$ times a unit after a small change in its fourth-order
coefficient. A morsification of the node and a shift of its
critical value away from zero give an isolated Morse point
of nonzero critical value.

\item If $g_z(0)=0$, its $xy$ coefficient is a unit and we take
$k=z$. Completing the square in (4.41c), including the
square root of $A$ in its first factor, gives exactly
\[
                         f=G^2+z^4U,
\]
with $G(x,y,0)=xy$ times a unit and $G_z(0)=0$.
Add a small constant to $U$, then take the actual coordinate
$Z=zU^{1/4}$. The parameter Morse lemma for $G$ in $x,y$
gives $G=XY+a(Z)$, where $a\in(Z^2)$. A small change
$a\mapsto a+\eta Z^2$ makes $a''(0)\ne0$. Replacing
$G$ by $G-\epsilon$ produces at the origin a critical point
of value $\epsilon^2$. The $XY$ Hessian entry is
$-2\epsilon$, and the $ZZ$ entry is
$-2\epsilon a''(0)$, so the point is Morse.
These deformations retain type $A_3$ along the moving
critical curve at each fixed boundary. The nested boundary
transport of P4.28 and P4R.14 places its square-$-2$ sphere
in the actual nearby torus, a contradiction. Hence $A(0)=0$.

\item If $g_z(0)$ is a unit, both saturated Hessian lines approach
$dz$, so $b_1=b_2=1$ and $d_i\ge1$. If $g_z(0)=0$,
the differential of $g$ on each branch has exactly the
common factor $t_i$ and a remaining nonzero coefficient
in the other normal direction. Thus $b_i=0,d_i\ge3$.
In either case the left side of (4.41a) is at least five,
including $q=1$.

\item Now let $q=0$. Choose actual generators $a,b$ of $P$ mapping
to $(1,0),(0,1)$ in $J\simeq B$. Their coefficient matrix
is invertible modulo the maximal ideal, hence invertible
locally. We obtain
\[
 I/P^2=(ya,xb),\qquad
 f=xU(x)b^2+yV(y)a^2\pmod {P^3}.                      \tag{4.41d}
\]
The second identity records the Hessian kernel $b$ on the
$x$ branch and $a$ on the $y$ branch. Write the exact
relation $xy=\alpha a+\beta b$. Near a general point of
the $x$ axis, solve this for $y$ using coordinates $x,a,b$.
The null cubic coefficient contains
$V(0)\alpha(x,0)/x$; the remaining contribution from
$P^3$ is holomorphic at the node. Its identical vanishing
in type $A_3$ forces $V(0)\alpha(0)=0$. The other branch
gives $U(0)\beta(0)=0$.

\item If both $a_z(0),b_z(0)$ are nonzero, the relation for $xy$
has $\alpha(0)\beta(0)\ne0$. Hence $U(0)=V(0)=0$,
both $b_i=1$, and both $d_i\ge2$. The defect is at least
five. Otherwise interchange the branches so that $a_z(0)$
is a unit and $b_z(0)=0$. Then $\alpha(0)=0$ and
$\beta(0)$ is a unit, forcing $U(0)=0$. On the $x$ axis,
$db$ has exactly the common factor $x$ and remaining
$dy$ coefficient a unit. On the $y$ axis, $da$ is primitive
with a unit $dz$ coefficient. Thus
\[
 b_x=0,\quad d_x=3+\operatorname{ord}_xU\ge4,
 \qquad
 b_y=1,\quad d_y=1+\operatorname{ord}_yV\ge1.
\]
This proves (4.41a). Equality requires
$\operatorname{ord}_xU=1$ and $V(0)\ne0$.

\item We identify the normalization in this exact equality case.
Use $a$ as the ambient coordinate $z$. The generator relation
with $\beta$ a unit gives
\[
                    I=(a^2,ya,x^2y-x\alpha a).
\]
Modulo its first two generators, replace $\alpha$ by
$\alpha(x,0,0)=x\ell(x)$. The actual coordinate
$Y=y-\ell(x)a$ puts the ideal into
$I=(z^2,Yz,x^2Y)$. Renaming $Y$ as $y$, its Hilbert--Burch
matrix and determinant extension are
\[
 \begin{pmatrix}y&0\\-z&x^2\\0&-z\end{pmatrix},
 \qquad
 M=\begin{pmatrix}y&0&a\\-z&x^2&b\\0&-z&c\end{pmatrix},
 \qquad
 \det M=az^2+byz+cx^2y.
\]
The new two-by-two minors are
$yb+za,yc,-zc,-x^2a,-za,x^2c+zb$, up to signs.
Their membership in $I$ gives exactly
\[
             a=yA,\qquad c=zC,\qquad b=-x^2C+zB
                          \pmod I.
\]
Therefore $f\in(yz^2)+I^2$, so
\[
           f=Uyz^2+Vz^4+Wx^2y^2z+Tx^4y^2.            \tag{4.41e}
\]
The equality orders $d_y=1,d_x=4$ say precisely
$U(0)T(0)\ne0$.

\item Substitute $y=z^2v$ and divide (4.41e) by $z^4$.
At the finite root its $v$ derivative is $U(0)$.
The implicit-function theorem supplies the actual smooth
factor $g_1=y-z^2v(x,z)$. Division by it gives
\[
                  g_2=z^2A+x^2yzB+x^4yC,
                  \qquad A(0)C(0)\ne0.
\]
For clarity, this division retains the factors $x^2,x^4$:
subtracting the expression at $y=z^2v$ supplies one factor
$y-z^2v$; every extra substituted $y$ supplies $z^2$ and
is absorbed into the first term. The parameter Morse
centre in $z$ is divisible by $x^2y$. Substitution leaves
$x^4y$ times a unit. An actual coordinate change therefore
gives $g_2=Z^2+x^4Y$. Its finite birational normalization is
\[
                 (x,u)\longmapsto(x,-u^2,x^2u),
\]
with target coordinates $(x,Y,Z)$. It is smooth; finiteness
follows from $u^2=-Y$ and birationality from $u=Z/x^2$
away from $x=0$. The other branch $g_1$ is smooth already.
Thus the full normalization is smooth over the node.

\end{enumerate}
\end{proof}

All determinant minors, the four minimum-defect cases,
the factor substitution, and the normalization equation
are verified by \texttt{code/part4/a3\_node.py}.
The output reports PASS. Its initial failed assertion
compared an expanded polynomial with a factored expression;
that output is retained, and the corrected assertion tests
their exact zero difference.

\begin{proposition}[P4.42]\hypertarget{tag-02XG}{}\label{night:P4.42}\label{stmt:P4.42}
\StatementRecord{night:P4.42}{proved}{02XG}{\SeriesRef{night:P4.33}, \SeriesRef{night:P4.34}, \SeriesRef{night:P4.40}, \SeriesRef{night:P4.41}, \SeriesRef{night:P4R.10}, \SeriesRef{night:P4R.13}, \SeriesRef{night:P4R.14}, \SeriesRef{night:P4R.6}}

In the requested global support classes, a positive-genus
ruled-resolution support whose critical curve has generic
type $A_3$ everywhere cannot have two two-branch points as
its point-identification cycle. The two points may be
nonordinary. This includes additional
unibranch singular points of the critical curve.
\end{proposition}
\begin{proof}

\begin{enumerate}[label={$\langle1\rangle$\arabic*.},leftmargin=*]
\item After P4R.13, the remaining positive ruled case has two
rational normalized conductor components and one graph
cycle. In the stated instance the cycle is provided by
two two-branch points. We first extend the local lower
bound of P4.41 to these points.

\item At a planar curve point, write its ideal as $(z,h)$.
Its symbolic square equals its ordinary square, so the
rank-one Hessian belongs to $\operatorname{Sym}^2 E_0$,
where $E_0$ is the normalized conormal image. On a branch
put $E_0=\langle t^a e_1,t^b e_2\rangle$ and write the
primitive Hessian direction as $\ell=u e_1+v e_2$.
One of $u,v$ is a unit. Membership of
$t^d\ell^2$ in $\operatorname{Sym}^2E_0$ and projection
onto the quotient by $\ell$ give
\[
 \frac d2\ge\max\{a-\operatorname{ord}u,
                         b-\operatorname{ord}v\},\qquad
 b_{\rm proj}=\min\{a+\operatorname{ord}v,
                         b+\operatorname{ord}u\}.
\]
The sum of the maximum and minimum on the right is
$a+b$. Thus $d/2+b_{\rm proj}\ge a+b$.
The plane different formula proved in P4.34 is
$\sum(a+b)=2\delta+\deg R$. The whole local defect
is consequently at least $3\delta+q$.
A nonordinary planar two-branch point has $\delta\ge2$
and contributes at least six.

\item For a nonplanar point, the ambient Hessian vanishes: a
nonzero rank-one Hessian would give $f=X^2+g(Y,Z)$ by
the parameter Morse lemma and would put the entire
critical curve in the smooth plane $X=0$.
On a branch of multiplicity $m$, all conormal generators
have order at least $m$, and every Hessian entry has
order at least $m$. Thus $b_{\rm proj}\ge m$, $d\ge m$,
and, since $r=m-1$,
\[
                  d/2+b_{\rm proj}-r\ge1+m/2\ge1.
\]
When $\delta\ge3$, a nonplanar two-branch point
therefore contributes at least five.
The only remaining possibility has $\delta=2$.
The exact excess-one classification in P4.40 identifies
it with a cusp and transverse line,
\[
                  P=(xz,yz,h),\qquad h=y^2-x^3.
\]
Here $P^{(2)}/P^2$ is generated by $zh$ and is
annihilated by the maximal ideal. Write $f=Azh+P^2$.
Along the cusp, the normal quadratic determinant is
$4KG-A^2$, with $K(0)=0$. Its rank is one generically,
so $A(0)=0$ and in fact $f\in P^2$.
The conormal lattices on the two normalized branches
are respectively $\operatorname{diag}(t^2,t^3)$ and
$\operatorname{diag}(t,t)$. Their total colength is seven,
while $\deg R=1$ and $\delta=2$.
Applying the same symmetric-square inequality gives
defect at least $2+7-1=8$.
Every two-branch point thus contributes at least $9/2$,
and equality is possible only at the ordinary node of
P4.41, whose normalization is smooth.

\item Apply the rank-one Hessian calculation
of P4R.10 to the two normalized projective lines. On each
line the rank-two conormal kernel has degree $2+\deg R$.
The rank-one quotient has Euler characteristic
$3+\deg R-\deg L$. Combining both lines with the
normalization sequence of $C$ gives the exact formula
\[
 \chi(\mathcal O_{C_{\rm sch}})
   =8-\sum_p\left[
       \delta_p+\sum_{i\text{ over }p}
          \left(\frac{d_i}{2}+b_i-r_i\right)+q_p
                 \right].                            \tag{4.42a}
\]
Every extra unibranch singular point has a strictly positive
bracket by the local conormal inequalities of P4R.10.
At a smooth point the bracket is nonnegative. The two
two-branch points contribute at least nine together by the
local calculation just given.

\item If the horizontal section is not contracted, P4.33 gives
$\chi(\mathcal O_{C_{\rm sch}})=g-1\ge0$, contradicting
(4.42a), whose right side is at most $-1$.
If the section is contracted, that Euler characteristic is
$-1$. Equality must then hold at both two-branch points and every
other bracket must vanish. There are no additional
unibranch singular points. Both cycle points are ordinary
nodes and the normalization is smooth over them by P4.41.
At every smooth point of the
critical curve, P4R.14 gives either the product tacnode
or the displayed $A_1$ normalization degeneration. The
latter has Hessian zero order one and contributes $1/2$,
so it is excluded by the equality in (4.42a). Thus the
normalization is smooth along all of the conductor, and
it is smooth away from the conductor by P4R.6. A contracted
positive-genus horizontal section in its minimal resolution
is impossible. This proves the assertion.

\end{enumerate}
\end{proof}


\subsection{The uniform cuspidal sector}

\begin{proposition}[P4.43]\hypertarget{tag-02XH}{}\label{night:P4.43}\label{stmt:P4.43}
\StatementRecord{night:P4.43}{proved}{02XH}{\SeriesRef{night:P4R.13}, \SeriesRef{night:P4R.7}}

In the requested full-fibration classes, a support with
positive-genus ruled resolution cannot have positive
Euler number and generic transverse type $A_2$ along
every component of its critical curve.
\end{proposition}
\begin{proof}

\begin{enumerate}[label={$\langle1\rangle$\arabic*.},leftmargin=*]
\item We use the exact normalization map
$\nu:N\longrightarrow S$ and its minimal resolution
$\rho:T\longrightarrow N$. A germ of $S$ cannot have
two distinct analytic hypersurface branches. Indeed,
their intersection has dimension one in the smooth
ambient threefold. A general point of an irreducible
component of that intersection has at least two local
hypersurface branches, so its transverse curve germ is
reducible. The transverse cusp $A_2$ is irreducible.
This contradicts the stated uniform transverse type.
The points of the finite normalization over a germ are
its analytic branches. Consequently $\nu$ is bijective.
It is proper, so its continuous map from a compact space
to the Hausdorff space $S$ is a homeomorphism.

\item Let $B$ be the base curve of the minimal ruled surface,
of genus $g\ge1$, and let $H$ be the horizontal section
of coefficient two in the total pole divisor. Point
blowups preserve the fundamental group, and the ruling
gives $\pi_1(T)=\pi_1(B)$. By the fibre-tree calculation
in P4R.7, every vertical exceptional configuration is a
tree of rational curves. Collapsing any such connected
tree preserves the fundamental group: apply van Kampen
to its tubular neighbourhood and the complement, or
equivalently to the cofibration quotient by that
simply connected tree. These quotients are the actual
topological contractions of the resolution map.

\item If $H$ is not contracted by $\rho$, all resolution
exceptional configurations are of this vertical kind.
It follows that $\pi_1(N)=\pi_1(B)$. If $H$ is
contracted, its connected exceptional configuration
contains $H$. The inclusion $H\hookrightarrow T$
induces an isomorphism on fundamental groups, since
the ruling restricts to the identity section and all
point blowups preserve that isomorphism. The quotient
which collapses this configuration therefore kills
the whole group. The other contracted configurations
are rational trees, giving $\pi_1(N)=1$.

\item On the other hand, the proved positive ruled reduction,
including P4R.13, gives a conductor curve $C$ with two
rational normalized components and one graph cycle.
The conductor cofibration pushout in P4R.7 gives
$\pi_1(S)=\pi_1(C)=\mathbb Z$. This contradicts the
homeomorphism $N\simeq S$: neither the trivial group
nor the genus-$g$ surface group is $\mathbb Z$, as their
abelianizations have ranks zero and $2g$, respectively.

\end{enumerate}
\end{proof}


\subsection{Further uniform transverse sectors}

\begin{proposition}[P4.44]\hypertarget{tag-02XI}{}\label{night:P4.44}\label{stmt:P4.44}
\StatementRecord{night:P4.44}{proved}{02XI}{\SeriesRef{night:P4.33}, \SeriesRef{night:P4.42}, \SeriesRef{night:P4R.14}, \SeriesRef{night:P4R.6}}

In the remaining positive ruled case with generic transverse
type $A_3$ on every conductor component, the horizontal
section of the pole divisor is contracted by the resolution
map. The point-identification cycle occurs at a single
three-branch point of the critical curve.
\end{proposition}
\begin{proof}

\begin{enumerate}[label={$\langle1\rangle$\arabic*.},leftmargin=*]
\item P4.42 excludes a cycle at two two-branch points. The
remaining graph has two rational normalized components
and one cycle, so its cycle occurs at one three-branch
point $p$. Suppose that the horizontal section $H$ is
not contracted. It maps finitely to one normalized
conductor component $C_H^\nu=\mathbb P^1$.
The generic tacnode has two conductor preimages.
The degree of $H\to\mathbb P^1$ is greater than one,
since $H$ has genus $g\ge1$, and hence is exactly two.
No other, vertical conductor component can map onto
$C_H$: the horizontal map already supplies both generic
preimages.

\item At a point where the entire critical curve is smooth,
P4R.14 gives either the product tacnode or the germ
$X(Y^2+XZ)$. Each has two distinct normalization points
over the critical-curve point. Each point lies on the
strict conductor, by P4R.6. Thus the two preimages must
both lie on $H$, and the degree-two map is unramified
there. All ramification is over singular points of $C$.

\item By (4.42a) and P4.33 the sum of the local defects is
$9-g\le8$. A nonplanar three-branch point has
$\delta\ge2$, and the branchwise estimate in P4.42
adds at least one per branch; its defect is at least
five. A planar three-branch point has $\delta\ge3$
and contributes at least $3\delta\ge9$, which is
already impossible. Every additional unibranch singular
point contributes at least two: in the planar case
the bound is $3\delta$, and in the nonplanar case it is
at least $\delta+1$. There is consequently at most one
additional unibranch singular point.

\item The normalization $C_H^\nu$ has at most two points
above $p$, since the other global conductor component
supplies at least one of its three branches. Including
the possible additional unibranch point leaves at most
three possible branch values for $H\to\mathbb P^1$.
Every ramification point of a degree-two map has index
two, and its branch value has only that ramified
preimage. Riemann--Hurwitz requires $2g+2\ge4$ distinct
branch values. This contradiction proves that $H$ is
contracted.

\end{enumerate}
\end{proof}

\begin{proposition}[P4.45]\hypertarget{tag-02XJ}{}\label{night:P4.45}\label{stmt:P4.45}
\StatementRecord{night:P4.45}{proved}{02XJ}{\SeriesRef{night:P4.28}}

In an actual torus degeneration, a point of a planar
singular critical curve cannot have generic transverse
type $D_4$ along all its local branches.
\end{proposition}
\begin{proof}

\begin{enumerate}[label={$\langle1\rangle$\arabic*.},leftmargin=*]
\item Write the full reduced critical-curve ideal as
$P=(u,h(x,y))$, where $h$ is reduced and singular.
Generic multiplicity three gives $f\in P^{(3)}$.
Powers of this complete-intersection prime-by-components
ideal are unmixed, so $P^{(3)}=P^3$. Thus there are
holomorphic coefficients such that
\[
 f=A u^3+B u^2h+Cuh^2+Dh^3.
\]
Introduce an independent variable $H$ and the actual
holomorphic function
\[
 F(u,x,y,H)=A(u,x,y)u^3+B(u,x,y)u^2H
               +C(u,x,y)uH^2+D(u,x,y)H^3.
\]
Then $f=F(u,x,y,h)$. Arbitrarily small constant changes
of the four coefficients preserve the $D_4$ transverse
germs on a fixed boundary and make the binary cubic
at the origin have distinct roots and a simple critical
point as a polynomial in $u/H$. Its critical value is
nonzero: a zero critical value would be a repeated root.
We also choose its leading coefficient nonzero.

\item After substituting $u=Hv$, the quotient $F_u/H^2$
is holomorphic. The implicit-function theorem gives
a holomorphic critical root $v=v(x,y,H)$ and
\[
          F(Hv(x,y,H),x,y,H)=H^3\beta(x,y,H),
                       \qquad \beta(0)\ne0.
\]
Choose the local cube root of the unit $\beta$.
The map $H\mapsto H\beta(x,y,H)^{1/3}$ has nonzero
$H$ derivative at the origin and has a holomorphic
inverse with parameters $x,y$. Put
\[
             \widehat h(x,y)=h(x,y)
                           \beta(x,y,h(x,y))^{1/3}.
\]
It defines the same singular reduced plane curve as $h$.
Choose a small morsification $\widehat h_s$ with a
Morse critical point of nonzero critical value. This
exists by the positive local degree of the gradient
of the singular reduced plane function; a small constant
shift avoids zero critical values. Use the displayed
inverse to solve exactly
\[
         H_s\beta(x,y,H_s)^{1/3}=\widehat h_s,
             \qquad f_s=F(u,x,y,H_s(x,y)).
\]
On its $u$-critical branch the critical-value function
is exactly $\widehat h_s^3$. At the selected plane
Morse point, $H_s\ne0$ and $F_{uu}\ne0$. The Schur
complement of the full Hessian is
$3\widehat h_s^2\operatorname{Hess}(\widehat h_s)$.
Consequently $f_s$ has a nondegenerate three-variable
critical point with nonzero critical value.

\item All deformations above can be chosen successively small
on the fixed boundary. There the moving curve $H_s=0$
has the same ordinary-triple transverse germ: its binary
cubic has three distinct linear factors. The boundary
transport used in P4.28 therefore identifies the local
nearby fibres through nested representatives. The Morse
vanishing sphere embeds in a nearby global torus fibre
and has self-intersection $-2$, contradicting its
asphericity. This proves the assertion.

\end{enumerate}
\end{proof}

\begin{proposition}[P4.46]\hypertarget{tag-02XK}{}\label{night:P4.46}\label{stmt:P4.46}
\StatementRecord{night:P4.46}{proved}{02XK}{\SeriesRef{night:P4.31}, \SeriesRef{night:P4.33}, \SeriesRef{night:P4.34}, \SeriesRef{night:P4.39}, \SeriesRef{night:P4.40}, \SeriesRef{night:P4.45}}

In the requested full-fibration classes, a support with
positive-genus ruled resolution cannot have positive
Euler number and generic transverse type $D_4$ along
every component of its critical curve.
\end{proposition}
\begin{proof}

\begin{enumerate}[label={$\langle1\rangle$\arabic*.},leftmargin=*]
\item If the horizontal section is contracted, P4.39 already
gives a contradiction. Suppose it is not contracted.
The remaining critical curve has two rational normalized
components and one graph cycle. For the actual
conductor scheme, ordinary-triple transverse type gives
the ideal $P^{(2)}$, as in P4.31. Let $Q=P/P^{(2)}$,
and let $\tau_p$ be the full colength of its differentiated
image in the direct image of the normalized rank-two
kernel $K$. This includes both the branchwise lattice
colength and the further descent colength. Riemann--Roch
on the two normalized projective lines gives
\[
 \chi(\mathcal O_{C_{\rm sch}})
       =10-\sum_p(\delta_p+\tau_p-r_p).
\]
P4.33 identifies the left side with $g-1$, so the total
defect is $11-g\le10$.

\item We prove that a three-branch point contributes at least
eleven. A planar point contributes $5\delta\ge15$ by
the exact different calculation in P4.34. Consider a
nonplanar point. If its three branches are smooth, the
curve has multiplicity three and embedding dimension
three. A generic linear parameter has an Artinian
section of length three and embedding dimension two,
with ideal $(v,w)^2$. The defining curve ideal therefore
has exactly three minimal generators. To justify the
lifting, the parameter is a nonzerodivisor on the
Cohen--Macaulay curve ring, so the Tor term vanishes in
the ideal sequence modulo that parameter; Nakayama
then lifts precisely the three generators.
Each branchwise conormal lattice has rank two and
colength at least two, since every defining gradient
has positive order. Their sum has colength at least
six. Modulo the three branch maximal ideals its
normalization has six generator values. The image of
$Q$, generated by three elements over the local curve
ring, supplies a vector space of dimension at most
three. The additional descent colength is therefore
at least three. Since $\delta\ge2$ and $R=0$, the
full defect is at least $2+6+3=11$.

\item If one of the three branches is singular, let $m\ge4$
be total multiplicity. The branchwise gradient bound
gives $\tau\ge2m$, while $\deg R=m-3$.
The general inequality $\delta\ge m-1$ leaves only
$m=4,\delta=3$ as a possible value below the required
bound. That value has excess $\delta-(3-1)=1$.
The exact classification in P4.40 shows that its only
singular-branch model is a cusp and two independent
lines, of embedding dimension four. It does not embed
in the ambient smooth threefold. Thus $\delta\ge4$
in the present case, and the full defect is at least
$\delta+2m-(m-3)=\delta+m+3\ge11$.
A cycle at one three-branch point is impossible.

\item The other possible cycle has two two-branch points.
Each contributes at least five by P4.34; equality is
possible only for a planar ordinary node. A nonplanar
two-branch point contributes at least seven, and a
nonordinary planar point at least ten. The total bound
of ten therefore forces two ordinary planar nodes
and $g=1$. Each such node is forbidden by P4.45.
This proves the assertion.

\end{enumerate}
\end{proof}

The local Hessian identities, rank-one lattice identity,
and exact integer defect bounds are checked in
\texttt{code/part4/uniform\_sectors.py}; the saved output
reports PASS.


\subsection{The full quadratic conductor at three axes}

\begin{proposition}[P4.47]\hypertarget{tag-02XL}{}\label{night:P4.47}\label{stmt:P4.47}
\StatementRecord{night:P4.47}{proved}{02XL}{\SeriesRef{night:P4.41}}

Consider an actual torus degeneration at a point where the
entire reduced critical curve is the three coordinate axes,
with generic transverse type $A_3$ on each axis. Suppose
the local Hessian--conductor defect is nine, the three
Hessian orders are two, and the quartic equation is the
square of $a xy+b xz+c yz$ with $abc\ne0$.
Then the normalization has only rational singularities
over that point. In particular it cannot contain the
contraction of a positive-genus horizontal section.
\end{proposition}
\begin{proof}

\begin{enumerate}[label={$\langle1\rangle$\arabic*.},leftmargin=*]
\item Put $P=(xy,xz,yz)$ and $h=xyz$. The symbolic square is
$P^{(2)}=P^2+(h)$, and the actual conductor image ideal
$I$ contains $P^{(2)}$. Also $f\in P^{(2)}$.
Write $f=Ah+P^2$. On the $x$-axis, a nonzero constant
term of $A$ would give a leading rank-two transverse
Hessian, since the quadratic term is then a nonzero
multiple of $xyz$. Generic type $A_3$ has rank one.
Thus $A(0)=0$. Since $\mathfrak m h\subset P^2$,
we have $f\in P^2$.

\item The normalized conormal lattice is $tK$ on each axis,
so its projected colength is one. The curve has
$\delta=2$ and $R=0$. The stated defect therefore
forces the remaining normalization descent colength
to be one. The normalized quotient values form a
two-dimensional subspace of $\mathbb C^3$ with one
relation involving all three coordinates. Multiplying
a source generator that is a unit on one branch by
that branch's axis coordinate generates its positive
orders alone. Hence all positive-order branch functions
occur in the image, and the leading kernel vector
lifts to an actual element $g\in I$ with quadratic
part $a xy+b xz+c yz$.

\item We prove the exact ideal equality $I=(g,h)$.
On each coordinate plane, $g$ is the product of its
two coordinate functions times a unit. Thus
$I_0=(g,h)$ has radical $P$. Furthermore $P^2\subset I_0$:
modulo $h$, multiplication of $g$ by $xy$ gives a unit
times $(xy)^2$, and the other two squares follow in
the same way; all mixed products are multiples of $h$.
At the generic point of every axis, $I_0$ is the
tacnode conductor ideal, generated by a normal coordinate
and the square of the null coordinate. It agrees
generically with $I$. The ideal $I_0$ is a
codimension-two complete intersection, so
$\mathcal O/I_0$ is Cohen--Macaulay and $P/I_0$ is
torsionfree on the reduced curve. The map
$P/I_0\to P/I$ has finite-length kernel, which must
vanish. This proves $I=I_0$.

\item The actual conductor ring condition, with the hypotheses
verified in P4.41, gives $f\in I^2$ by the two-by-two
adjugate calculation of de Jong--van Straten,
Theorem~1.12. Write $f=A g^2+Bgh+Ch^2$.
Its quartic part makes $A(0)\ne0$. Completing the square
and absorbing $\sqrt A$ into the first factor gives
the exact equation
\[
                         f=G^2+h^2U.                 \tag{4.47a}
\]
The quadratic part of $G$ is nondegenerate, and
$(G,h)=(g,h)$. If $U(0)\ne0$, the two factors
$G\pm h\sqrt{-U}$ have nondegenerate quadratic
Hessians. Each is an isolated $A_1$ surface singularity,
so the full normalization has rational singularities.

\item We exclude $U(0)=0$ by a deformation that keeps the
three axes fixed. Add an arbitrarily small generic
linear form to $U$, and then replace $U$ by
$U-\epsilon$. On a fixed boundary the transverse
$A_3$ coefficient is a unit, so these changes retain
the boundary transverse type. A diagonal linear change
puts the quadratic part of $G$ in the form
$G_2=xy+xz+yz$, changing $h$ only by a nonzero constant.
Let $L$ be the resulting linear part of $U$.
Under $(x,y,z)=\epsilon(X,Y,Z)$, equation (4.47a) becomes
\[
 \frac{f_\epsilon}{\epsilon^4}
   =G_\epsilon(X,Y,Z)^2+\epsilon^3(XYZ)^2
                          \bigl(L(X,Y,Z)-1+O(\epsilon)\bigr),
 \qquad G_\epsilon=G_2+O(\epsilon).
\]
Near a nonzero point of the smooth cone $G_2=0$,
use $G_\epsilon$ as the normal coordinate. The
parameter Morse centre is $O(\epsilon^3)$.
Substitution leaves the tangential function
$\epsilon^3(\Phi+O(\epsilon))$, where
\[
                 \Phi=(XYZ)^2(L-1)|_{G_2=0}.
\]

\item We verify a nonempty open set of choices of $L$ for
which $\Phi$ has a nondegenerate critical point of
nonzero value. Parametrize the projective conic by
\[
              [X:Y:Z]=[t(t+1):t+1:-t].
\]
For $L=X+Y+Z$, set
$H(t)=-t^2(t+1)^2$ and $\ell(t)=t^2+t+1$.
In radial coordinate $r$,
$\Phi=r^6H(t)^2(r\ell(t)-1)$.
Radial criticality gives $r=6/(7\ell(t))$, and
\[
 \frac{d}{dt}\left(\frac{H}{\ell^3}\right)
  =\frac{t(t-1)(t+1)(t+2)(2t+1)}{(t^2+t+1)^4}.
\]
The three critical points $t=1,-2,-1/2$ are simple
and avoid zeros and poles. Their radial coordinates
are $2/7,2/7,8/7$. At each point,
\[
       \Phi=-\frac{1024}{823543},\qquad
       \det\operatorname{Hess}_{r,t}\Phi
                     =\frac{4718592}{1977326743}.
\]
These numbers are nonzero. Simplicity and avoidance
of zeros and poles are nonvanishing algebraic conditions
on the coefficients of $L$. The explicit example proves
that they hold on a nonempty Zariski open set, so an
arbitrarily small generic linear perturbation of any
initial linear part is sufficient.

\item After eliminating the normal Morse coordinate, divide
the two tangential critical equations by $\epsilon^3$.
They are analytic in $\epsilon$ and have invertible
Jacobian at the displayed points. The implicit-function
theorem lifts them to nondegenerate three-variable
critical points of $f_\epsilon$. Their critical values
have leading term $\epsilon^7$ times the displayed
nonzero value. They tend to the original point.
Since the deformation retains boundary $A_3$, the
nested boundary transport embeds a Morse vanishing
sphere of square $-2$ into a nearby global torus fibre.
This contradicts asphericity. Therefore $U(0)=0$
cannot occur, proving the assertion.

\end{enumerate}
\end{proof}

The conic derivative, all three exact critical values
and Hessian determinants, and the related conductor
matrices are executed in
\texttt{code/part4/a3\_three\_axes.py}.
The output reports PASS. The working record explicitly
rejects the earlier proposed change $G\mapsto G-\epsilon$:
with $h$ fixed that change replaces the boundary tacnode
by two ordinary double points and does not supply the
required boundary transport.


\subsection{The three-branch tacnode defect budget}

\begin{proposition}[P4.48]\hypertarget{tag-02XM}{}\label{night:P4.48}\label{stmt:P4.48}
\StatementRecord{night:P4.48}{proved}{02XM}{\SeriesRef{night:P4.42}, \SeriesRef{night:P4.47}}

At a planar three-branch critical-curve point with uniform
generic transverse type $A_3$, the local conductor defect
is at least ten. At an ordinary nonplanar three-branch
point, the defect is at least nine. Equality in the
second assertion forces Hessian orders $(2,2,2)$,
projected colengths $(1,1,1)$, descent colength one,
and a quartic square $(a xy+b xz+c yz)^2$ with at
least two nonzero coefficients. The case $abc\ne0$
has rational normalization by P4.47.
\end{proposition}
\begin{proof}

\begin{enumerate}[label={$\langle1\rangle$\arabic*.},leftmargin=*]
\item For a planar critical curve the ideal is $P=(u,h)$.
The rank-one nilradical $J=P/I$ of the conductor image
is generated by at most two elements over $\mathcal O_C$.
Its normalized image has at least one generator value
on each of the three branches. Modulo the three branch
maximal ideals, the source supplies a vector space of
dimension at most two, while the normalized image
supplies dimension three. Hence its descent colength
$q$ is at least one. The planar symmetric-square and
different calculation in P4.42 gives local defect
$\ge3\delta+q$. Since a planar three-branch point has
$\delta\ge3$, its defect is at least ten.

\item For an ordinary nonplanar point use actual coordinates
with $P=(xy,xz,yz)$. The argument in P4.47 gives
$f\in P^2$. On each axis, its normalized conormal
lattice is exactly $tK$, so the projected colength
is exactly one. Every Hessian order is at least two.
Write the quartic form in the basis $(xy,xz,yz)$ as
\[
 A(xy)^2+C(xz)^2+E(yz)^2
      +2B(xy)(xz)+2D(xy)(yz)+2F(xz)(yz).
\]
The rank-one Hessians give
$AC=B^2$, $AE=D^2$, $CE=F^2$.
If a diagonal coefficient vanishes, its adjacent
off-diagonal coefficients vanish and the remaining
two-by-two block has rank at most one. If all three
diagonal coefficients are nonzero, the vanishing null
cubic on the $x$-axis gives $F=BD/A$. Higher-degree
terms cannot alter the coefficient linear in $x$
of that normal cubic. Thus the whole quartic is a
square $(a xy+b xz+c yz)^2$.

\item If the quartic is zero, all three Hessian orders are
at least three, and the local defect is at least
$2+3+9/2=19/2$. If the square has just one nonzero
coefficient, say $a xy$, the $z$-axis Hessian order
is at least three. In the source-generator basis
$(xy,xz,yz)$, the three leading quotient-value rows
have the forms
\[
                      (0,1,0),\qquad(0,0,1),
                      \qquad(0,*,*).
\]
Their rank is at most two, so $q\ge1$ and again
the defect is at least $19/2$.

\item If at least two coefficients are nonzero, all three
Hessian orders are two. The leading quotient-value
matrix has rank two. For example, up to rescaling
its rows it is
\[
 \begin{pmatrix}-b&a&0\\-c&0&a\\0&-c&b\end{pmatrix}.
\]
It has zero determinant and rank two whenever at
least two of $a,b,c$ are nonzero. Therefore $q\ge1$,
and the local defect is at least $2+3+3+1=9$.
Equality forces $q=1$ and the displayed Hessian and
projection orders. P4.47 applies when all three
quadratic coefficients are nonzero.

\end{enumerate}
\end{proof}


\subsection{The remaining ordinary-axis conductor}

\begin{proposition}[P4.49]\hypertarget{tag-02XN}{}\label{night:P4.49}\label{stmt:P4.49}
\StatementRecord{night:P4.49}{proved}{02XN}{\SeriesRef{night:P4.41}, \SeriesRef{night:P4.47}, \SeriesRef{night:P4.48}}

At an ordinary three-axis critical-curve point with uniform
generic transverse type $A_3$ and local defect nine, the
normalization has rational singularities. This includes the
quartic square with exactly one missing coefficient left
by P4.47--P4.48.
\end{proposition}
\begin{proof}

\begin{enumerate}[label={$\langle1\rangle$\arabic*.},leftmargin=*]
\item Only the one-missing-coefficient case remains. Rescale
its leading conductor form to $x(y+z)$.
As in P4.47, the descent colength one lifts the leading
kernel vector to $g\in I$ with this quadratic part.
Write $g=x(\alpha y+\beta z)+yz\gamma$, where
$\alpha,\beta$ are units and $\gamma(0)=0$.
Replacing $y,z$ by $\alpha y,\beta z$ is an analytic
coordinate change preserving the axes. Put $k=y+z$,
$h=yz$. Then $g=xk-hc$, with $c(0)=0$.
The ideal $I_0=(g,xh,h^2)$ lies in the actual conductor
$I$ and has Hilbert--Burch matrix
\[
                 \begin{pmatrix}h&0\\-k&h\\c&-x\end{pmatrix}.
\]
It has radical $(xy,xz,yz)$, is Cohen--Macaulay, and
agrees with $I$ generically on each axis. The
torsionfree finite-kernel comparison used in P4.47
therefore gives the exact equality $I=I_0$.

\item There are exactly two required ideal forms. Since
$\mathcal O/(x,k,h)=\mathbb C\{y\}/(y^2)$, analytic
division writes $c=\lambda y+xA+kB+hC$ for a scalar
$\lambda$. Replace $g$ by $g+xhA+h^2C$, and then
set $x'=x-hB$. These generator and coordinate changes
give
\[
                    I=(x'k-\lambda yh,x'h,h^2).
\]
If $\lambda=0$ this is $(xk,xh,h^2)$. If
$\lambda\ne0$, rescaling $x'=\lambda X$ gives
$(Xk-yh,Xh,h^2)$. Every change is invertible near
the origin. We treat the two forms separately.

\paragraph{The ideal $(xk,xh,h^2)$.}
Appending $(a,b,c)$ to its Hilbert--Burch matrix gives
determinant $axk+bxh+ch^2$. The adjugate ring condition
of de Jong--van Straten, Theorem~1.12, applies to the
actual conductor as in P4.41. The colon ideals
\[
 I:h=(x,h),\qquad I:x=(k,h),\qquad I:k=(x,h^2)
\]
give exactly
$f=A x^2k^2+B x^2kh+Cxh^2+Dh^4$, with $A(0)\ne0$.
Complete the square in $k$, absorb $\sqrt A$ into $x$,
and replace $y$ by $y(1+\gamma z)$ to make
$k+\gamma h$ the new $y+z$. These are analytic
coordinate maps; $h$ changes only by a unit.
Absorbing that unit in the coefficients gives
\[
                       f=x^2k^2+xh^2C+h^4D.          \tag{4.49a}
\]
If $C(0)\ne0$, substitute $x=h^2v$ and divide by
$h^4$. The finite $v$-root is simple, giving a smooth
factor $x-h^2v(y,z)$. Its quotient is $xk^2+h^2U$
with $U$ a unit. The coordinate $X=x/U$ turns it
into $Xk^2+h^2$. Its finite birational normalization is
\[
                 u=h/k,\qquad X=-u^2,
                 \qquad u(y+z)=yz.
\]
The last equation is an $A_1$ quadric, so the
normalization has rational singularities.

\item Suppose $C(0)=0$. Perturb $D$ by a small constant
and the linear part of $C$ generically, retaining
$C(0)=0$, and then replace $C$ by $C-\epsilon$.
On a fixed boundary the transverse tacnode coefficients
are units, so these changes retain boundary $A_3$.
Under $x=\epsilon^3X$, $y=\epsilon Y$,
$z=\epsilon Z$, the leading scaled equation is
\[
 \frac f{\epsilon^8}=F+O(\epsilon),\qquad
 F=X^2K^2+XH^2(L-1)+D_0H^4,
 \quad K=Y+Z,\ H=YZ.
\]
For $L=K$, $D_0=1$, put $W=Y-Z$.
At $W=0$ the $X$-critical coordinate is
$X=-K^2(K-1)/32$, and the remaining derivative is
$K^5(12K^2+7K-3)/512$.
At either root of $12K^2+7K-3$, the full critical
value and Hessian determinant are respectively
\[
 \frac{K^6(K+1)(3K-1)}{1024},\qquad
 -\frac{3K^{10}(K+1)(3K-1)(28K^2+14K-5)}{32768}.
\]
They are nonzero: their numerator resultants with the
critical polynomial are $8748$ and $4923269424$.
Nondegeneracy, nonzero value and avoidance of the axes
hold on a nonempty algebraic open set of coefficients.
An arbitrarily small generic perturbation is therefore
available. The implicit-function theorem lifts these
Morse points through the higher weighted terms, with
nonzero critical values of order $\epsilon^8$.
The boundary comparison embeds their square-$-2$
vanishing spheres in a nearby torus fibre, a
contradiction. Thus $C(0)=0$ is excluded.

\paragraph{The ideal $(g,xh,h^2)$, where $g=xk-yh$.}
Append $(a,b,d)$ to
\[
                         \begin{pmatrix}h&0\\-k&h\\y&-x\end{pmatrix}.
\]
Its determinant is $ag+bxh+dh^2$. The new adjugate
entries give the conditions
\[
       ax,\ ah,\ bx+dh,\ kd+by,\ dh-ay,\ ak+bh\in I.
\]
The exact colon identities are
\[
 I:h=(x,h),\qquad I:x=(xk,h,z^2),\qquad I:x^2=(k,h).
\]
The first two conditions give $a=hA$ modulo $I$;
then $dh-ay$ gives $d=Ay+xD+hE$.
Initially $bx+dh$ permits
$b=-Ak+hB+z^2B_2$ modulo $I$.
The condition $ak+bh$ forces $z^2B_2\in(x,h)$,
so $B_2\in(x,y)$. Its $x$-part lies in $I$, since
$xz^2=h^2-xh$ modulo $I$, and its $y$-part is
absorbed in $hB$. Thus $b=-Ak+hB$ modulo $I$.
The remaining condition is $y(D+B)+kE\in(x,h)$.

\item Write $T=D+B$. Modulo $(x,h,k)$, the equality
$yT=0$ in $\mathbb C\{y\}/(y^2)$ forces $T(0)=0$.
Write $T=xW+yU+kV$. Since $k$ is a nonzerodivisor
modulo $(x,h)$, the general solution is
$E=-yU-yV+xE_1+hE_2$. Substitution in the determinant
gives $f\in(p)+I^2$, where
\[
                            p=yh^2(x-h).
\]
The exact identities
$xh^3=ygh^2-yp$ and $h^4=ygh^2-kp$ then put the
equation in the form
\[
 f=A g^2+B gxh+Cgh^2+D(xh)^2+\kappa p,
                       \qquad A(0)\ne0.              \tag{4.49b}
\]

\item We first treat $\kappa(0)\ne0$.
Use the appended coefficients
$a=Ag+Bxh+Ch^2$, $b=Dxh+\kappa yh$,
$d=-\kappa yh$. The actual conductor identifies
$\operatorname{Hom}(I,\mathcal O_S)$ with the
normalization, generated by $1,n_1,n_2$, where the
adjugate fractions are
\[
          n_1=-a/h,\qquad n_2=-(ak+bh)/h^2.
\]
Since $A$ is a unit, the same actual normalization
algebra is generated by $1,v,w$, where
$v=(n_1+Bx+Ch)/A=-g/h$ and $w=n_2$.
Direct multiplication gives the exact relations
\begin{align*}
 Av^2-Bxv-Chv+Dx^2+\kappa xy-\kappa yh&=0,\\
 xw&=Ayv-Bxy-(\kappa+C)yh,\\
 vw&=\kappa h-Dxy,\\
 w^2+[(\kappa+B)y+Ck]w-A\kappa k+ADy^2
             -D(\kappa+C)yh-B\kappa h&=0.
\end{align*}
At the origin they give $v=w=0$, so there is one
normalization point. The last equation eliminates $k$
analytically, with $k$ in the maximal-ideal square.
Thus $h=y(k-y)$ has quadratic initial term $-y^2$.
Write $a=A_0$, $b=B_0$, $d=D_0$, $u=\kappa_0$.
The three quadratic equations become
\[
 av^2-bxv+dx^2+uxy=0,\qquad
 xw-ayv+bxy=0,\qquad vw+uy^2+dxy=0.
\]
Set $V=av-(b/2)x$, $\delta=ad-b^2/4$,
$Y=\delta x+auy$, $W_0=w+(b/2)y$, and
$W=auW_0+\delta V$. This invertible linear change
turns them into
\[
                       V^2+xY=0,\qquad
                       VW+Y^2=0,\qquad xW-YV=0.
\]
They define the irreducible rational normal cubic cone.
The full tangent-cone ideal contains this prime ideal
and has the same dimension two, so it equals it.
Blowing up the normalization point is smooth along its
exceptional rational normal cubic: the two independent
initial equations have independent derivatives along
that smooth projective curve in every blowup chart.
The exceptional curve is $\mathbb P^1$ with normal
bundle $\mathcal O(-3)$. The exact sequences of its
thickenings have successive terms $\mathcal O(3n)$,
all with $H^1=0$. Formal functions gives
$R^1\rho_*\mathcal O=0$. The singularity is rational.

\item Finally suppose $\kappa(0)=0$. Perturb $D$ and the
linear part of $\kappa$ generically, and replace
$\kappa$ by $\kappa-\epsilon$. The fixed axes retain
boundary $A_3$: equation (4.49b) stays in the same
ring-condition module, and its boundary transverse
quartic coefficients remain units.
Scale $x=\epsilon^2X$, $y=\epsilon Y$, $z=\epsilon Z$.
Put $K=Y+Z$, $H=YZ$, $G=XK-YH$.
After division by $\epsilon^6$, the first terms are
\[
 A_0G^2+\epsilon(B_0GXH+C_0GH^2)
   +\epsilon^2\bigl[D_0(XH)^2+p(L-1)\bigr].
\]
Near a nonzero smooth point of $G=0$, eliminate the
normal Morse coordinate. On that cone the identities
$XH^3=-Yp$, $H^4=-Kp$ absorb the square-completion
corrections into the coefficients, leaving the leading
tangential function
\[
                    \Phi=D'(XH)^2+p(L'-1)|_{G=0}.
\]
Choose $D'=1$, $L'=-Z$ to exhibit a nonempty open
set of suitable coefficients. The parametrization
\[
 K=r,\quad Y=tr,\quad Z=(1-t)r,\quad
 X=t^2(1-t)r^2
\]
gives $\Phi=t^4(1-t)^4r^7[1+r(t^2-t+1)]$.
At $t=1/2$, $r=-7/6$, its value and Hessian are
\[
 -\frac{823543}{573308928},\qquad
 \operatorname{diag}\left(\frac{117649}{1990656},
                          \frac{15647317}{214990848}\right).
\]
The Hessian determinant is
$1840891197733/427972821516288$, nonzero.
The angular critical numerator factors as
$(2t-1)(3t^2-3t-4)$, so this point is simple and
avoids the axes and $K=0$. The generic coefficient
condition is algebraic and nonempty. After normal
elimination, divide the tangential critical equations
by $\epsilon^2$ and apply the implicit-function theorem.
The resulting full Morse points have nonzero critical
values of order $\epsilon^8$ and tend to the origin.
Boundary transport gives the forbidden square-$-2$
vanishing sphere in a nearby torus fibre. Thus
$\kappa(0)=0$ is excluded, completing both ideal forms.

\end{enumerate}
\end{proof}

All adjugates, colon ideals, normalization relations,
linear coordinate identities, and both polar Hessian
calculations are executed in
\texttt{code/part4/a3\_missing\_coefficient.py}.
The saved output reports PASS.

\begin{proposition}[P4.50]\hypertarget{tag-02XP}{}\label{night:P4.50}\label{stmt:P4.50}
\StatementRecord{night:P4.50}{proved}{02XP}{\SeriesRef{night:P4.33}, \SeriesRef{night:P4.42}, \SeriesRef{night:P4.44}, \SeriesRef{night:P4.47}, \SeriesRef{night:P4.48}, \SeriesRef{night:P4.49}, \SeriesRef{night:P4R.14}, \SeriesRef{night:P4R.16}, \SeriesRef{night:P4R.17}, \SeriesRef{night:P4R.6}}

In the requested full-fibration classes, no support
with positive-genus ruled resolution has positive
Euler number and generic transverse type $A_3$
on every critical-curve component.
\end{proposition}
\begin{proof}

\begin{enumerate}[label={$\langle1\rangle$\arabic*.},leftmargin=*]
\item P4.42 and P4.44 leave only a contracted horizontal
positive-genus section and a cycle at one three-branch
point of a conductor with two rational normalized
components. The total defect is nine by P4.33 and
(4.42a). P4.48 excludes planar triples; P4R.16
excludes the nonplanar excess-one tacnode-plus-line
case; the exact excess-two classification and conormal
calculation in P4R.17 exclude
every delta-four nonplanar triple. The branchwise bound
$\delta+3+m/2$ excludes all larger delta invariants,
where $m\ge3$ is total multiplicity.
The only remaining triple is the ordinary three-axis
point. P4.48 gives local defect at least nine, so it
uses the entire budget, and P4.47 and P4.49 give
rational normalization there.

\item Every additional singular critical-curve point has
strictly positive defect and therefore is absent.
At a smooth point of the entire critical curve,
P4R.14 gives the product tacnode or its $A_1$
normalization degeneration; both have rational
normalization. Away from the conductor, P4R.6
excludes exceptional configurations. The normalization
therefore has rational singularities everywhere.
Its minimal resolution cannot contract the horizontal
positive-genus section. This contradiction proves
the assertion.

\end{enumerate}
\end{proof}


\subsection{Local intersection forms and the quartic tangent obstruction}

\begin{proposition}[P4.51]\hypertarget{tag-02XQ}{}\label{night:P4.51}\label{stmt:P4.51}
\StatementRecord{night:P4.51}{proved}{02XQ}{\SeriesRef{night:P4R.4}}

Let $S$ be an irreducible reduced support in the torus
fibration of Part~4.2 with $b_1(S)>0$. Every local
Milnor fibre of a reduced local equation of $S$ has
identically zero rational intersection form. This applies
in particular to the remaining positive ruled support,
whose fundamental group is $\mathbb Z$.
\end{proposition}
\begin{proof}

\begin{enumerate}[label={$\langle1\rangle$\arabic*.},leftmargin=*]
\item Let $U$ be a sufficiently small full tube and $F_t$ a
nearby compact smooth torus fibre. P4R.4 identifies the
homomorphism
\[
 \pi_1(F_t)=\mathbb Z^4\longrightarrow
 \pi_1(U)=\pi_1(S)
\]
and gives a finite cyclic cokernel, of order equal to
the fibre multiplicity. Its map on rational first
homology is therefore onto. Write
\[
 W=\ker\bigl(H_1(F_t;\mathbb Q)\longrightarrow
                  H_1(S;\mathbb Q)\bigr).
\]
Then $\dim W=4-b_1(S)\le3$.

\item Choose a coordinate ball about the point, with its
closed ball contained in $U$. The local full fibre
equation is $u g^m$, with $g$ reduced and $u$ a
holomorphic unit. We take a holomorphic $m$-th root of
$u$ on the ball and absorb it into $g$. The actual
nearby fibre in the ball is consequently the disjoint
union of the local Milnor fibres
$g=\zeta t^{1/m}$. Each such fibre $M$ is an oriented
codimension-zero submanifold with boundary of $F_t$.
Its inclusion into $U$ factors through the contractible
ball. Thus every loop of $M$ maps to zero in
$H_1(U;\mathbb Q)$, and the image subgroup
\[
 H=\operatorname{im}\bigl(\pi_1(M)\longrightarrow
                                 \mathbb Z^4\bigr)
\]
satisfies $H\otimes\mathbb Q\subset W$.

\item The covering criterion lifts $M\to F_t$ to the
connected cover $F_H\to F_t$ corresponding to $H$.
This cover is a $K(H,1)$, with $H$ free abelian. The
map on second rational homology is therefore exactly
\[
 \bigwedge^2(H\otimes\mathbb Q)\longrightarrow
              \bigwedge^2 H_1(F_t;\mathbb Q).
\]
Every class from $H_2(M;\mathbb Q)$ consequently maps
into $\bigwedge^2W$. The torus intersection pairing is
exterior multiplication followed by the orientation
isomorphism $\bigwedge^4H_1(F_t;\mathbb Q)=\mathbb Q$.
It vanishes on $\bigwedge^2W$ because $\dim W\le3$.

\item This ambient pairing is the local intersection
pairing on $M$. We represent finitely many cycles
inside its interior and make them transverse there;
the orientation-preserving inclusion into $F_t$
retains all intersection points and signs. Thus the
local pairing is zero. The symbolic exterior-square
matrix and its zero Gram matrix are recorded in
\texttt{code/part4/local\_intersection.py} and its
executed output.

\end{enumerate}
\end{proof}

\begin{proposition}[P4.52]\hypertarget{tag-02XR}{}\label{night:P4.52}\label{stmt:P4.52}
\StatementRecord{night:P4.52}{proved}{02XR}{\SeriesRef{night:P4.20}, \SeriesRef{night:P4.39}, \SeriesRef{night:P4.51}, \SeriesRef{night:P4R.15}, \SeriesRef{night:P4R.6}}

In the remaining positive ruled support of Part~4.2,
the image of a contracted horizontal pole section has
ambient multiplicity exactly three. This restriction
imposes no generic-transverse-type hypothesis on
any component of the critical curve.
\end{proposition}
\begin{proof}

\begin{enumerate}[label={$\langle1\rangle$\arabic*.},leftmargin=*]
\item We first give the comparison with a leading
homogeneous polynomial. Write the reduced equation
as $g=g_d+g_{d+1}+\cdots$, and suppose that $g_d$ is
reduced. Every finite collection of cycles in
$M_d=\{g_d=1\}$, together with its intersections,
transports into a local Milnor fibre of $g$.
Indeed, choose a compact neighbourhood of the cycle
representatives. Euler's identity shows that
$dg_d\ne0$ on $M_d$. The functions
\[
 G_\epsilon(x)=\epsilon^{-d}g(\epsilon x)
\]
converge with their derivatives to $g_d$ on that
neighbourhood. We solve
\[
 dG_\epsilon(V_\epsilon)
                     =-\partial_\epsilon G_\epsilon
\]
using the Hermitian gradient divided by its squared
norm. A cutoff supported in a slightly larger
compact neighbourhood, equal to one along the
trajectories under consideration, gives a flow for
small positive real $\epsilon$. It transports those
cycles to $G_\epsilon=1$. Holomorphic scaling by
$\epsilon$ then embeds them in $g=\epsilon^d$ inside
the fixed original ball, preserving orientations
and intersections. This compares smooth nonzero
level sets on compact neighbourhoods; no assertion
about equisingularity of the zero set is needed.
By P4.51 the pairing on $M_d$ is consequently zero.

\item For a reduced homogeneous polynomial its global
fibre and a sufficiently large compact Milnor fibre
have the same homology and scalar monodromy. One
can see this directly from the squared radius on
the affine fibre. This is a proper semialgebraic
function with finitely many critical values; above
the largest, its gradient supplies a product end.
A large invariant ball therefore gives a compact
core. Rescaling identifies that core with a small
local Milnor fibre. The Milnor-fibre CW dimension
bound gives zero homology above degree two, and
reducedness gives connectedness.

\item We need one exact variation identity. For such a
homogeneous polynomial, let $\mu_i$ be the
transverse plane-curve Milnor numbers at the
singular points of its reduced projective curve,
and put $M=\sum_i\mu_i$. Write $F$ for its compact
Milnor fibre and $\partial_2F$ for the boundary
part over the singular-link circles. Let
\[
 \begin{split}
 j&:H_2(F;\mathbb Q)\longrightarrow
                         H_2(F,\partial F;\mathbb Q),\\
 j_2&:H_2(F;\mathbb Q)\longrightarrow
                         H_2(F,\partial_2F;\mathbb Q).
 \end{split}
\]
Siersma's chain-level variation identity
\cite[\S3.6, pp.~451--452; \S7.1, p.~462]
{Siersma1991Variation} is
\[
                 j_2(T-I)=\operatorname{VAR}^{II}j.
                                                            \tag{4.52a}
\]
Here $T$ is the Milnor monodromy on $H_2(F;\mathbb Q)$.
The identity does not require the variation map to
be an isomorphism in dimension two. The dimension-two
Wang sequence is explicitly treated in that source,
\S6.3. Poincar\'e--Lefschetz duality identifies $j$
with the intersection form. Its vanishing makes
$\operatorname{im}(T-I)$ lie in $\ker j_2$.
The exact sequence of $(F,\partial_2F)$ identifies
that kernel with the image of $H_2(\partial_2F)$.
Each boundary piece is the mapping torus of a
transverse Milnor fibre. Its Wang sequence gives
\[
 H_2(\partial_2F;\mathbb Q)
       =\bigoplus_i\ker(A_i-I)\subset
                        \bigoplus_iH_1(F_i;\mathbb Q),
\]
where $A_i$ is the vertical monodromy. Thus zero
intersection form implies
\[
                 \operatorname{rank}(T-I)\le M.             \tag{4.52b}
\]
An empty $\partial_2F$ gives the same assertion with
$M=0$.

\item Suppose now that the reduced projective quartic
$D=\{g_4=0\}$ has a normalization component of
positive genus. By degree, either $D$ is an
irreducible quartic with normalization genus
$g\ge1$, or it is a smooth elliptic cubic and a
line. In the first case
$\sum\delta_i=3-g\le2$, so
$M=\sum(2\delta_i-r_i+1)\le4$.
In the second case B\'ezout gives total intersection
length three and at least one intersection point.
Hence $\sum\delta_i=3$ and
$\sum(r_i-1)\ge1$, giving $M\le5$.

\item The map $\{g_4=1\}\to\mathbb P^2\setminus D$ is
an unramified fourfold cover. The normalization
sequence and plane-curve genus formula give
\[
              \chi(F)=4\chi(\mathbb P^2\setminus D)
                      =4(7-M).
\]
The three nonidentity scalar monodromies act freely
on the compact invariant core, so their Lefschetz
numbers are zero. It follows by the character
formula for the cyclic group of order four that
\[
 H_0(F;\mathbb C)-H_1(F;\mathbb C)+H_2(F;\mathbb C)
       =(7-M)\,\mathbb C[\mathbb Z/4]
\]
as virtual representations. Every nontrivial
character therefore has multiplicity at least
$7-M$ in $H_2$. The finite-order monodromy is
semisimple, and hence
\[
             \operatorname{rank}(T-I)
                   \ge3(7-M)\ge6>M.                       \tag{4.52c}
\]
This contradicts (4.52b).

\item We apply this obstruction through the actual pole
comparison. Blow up the ambient origin, normalize
the strict transform, and take its minimal
resolution $T_1$. As in P4.39, this resolution maps
to the original normalization and dominates its
minimal resolution by point blowups. The strict
transform of the positive-genus horizontal section
$H$ retains pole coefficient two. If $E$ is the
ambient exceptional divisor, put
$n=\operatorname{ord}_H(E)>0$, and let $a\ge0$
be the coefficient of $H$ in the total conductor
divisor relative to the strict-transform dualizing
sheaf. Adjunction gives exactly
\[
                          2=(d-2)n+a.                     \tag{4.52d}
\]
Effectiveness follows from the exceptional
intersection matrix argument of P4R.6. An
exceptional positive-genus curve over the new
normalization has strictly positive coefficient,
by its adjunction equation. Thus $d\ge5$ is
impossible. If $d=4$, then $n=1,a=0$; the curve
$H$ is neither exceptional over the new
normalization nor over its conductor. The
normalization is an isomorphism at its generic
point. Its image is a multiplicity-one component
of the quartic exceptional curve, birational to
$H$. It has positive genus. Such a component has
degree three or four, leaving at most one linear
factor; the entire quartic polynomial is reduced.
The preceding contradiction excludes this case.

\item For clarity, an isolated critical point cannot
hide a contracted $H$. A generic small linear
perturbation of an isolated hypersurface critical
point has a positive finite number of Morse
critical points, the local degree of its gradient.
A small constant shift makes a chosen critical
value nonzero. On a fixed sufficiently small ball
the original zero-level boundary is smooth and
transverse; compactness retains the boundary
submersion over a uniform punctured value disc.
The finite-critical-value and path-transport
argument of P4.20 then carries the Morse
vanishing sphere, of square $-2$, into the
original local Milnor fibre. P4.51 excludes it.
This argument does not assume smoothness of the
normalization.

\item At a nonisolated point with $d=2$, the Hessian is
nonzero. P4R.15 gives the exact product equation
$X^2+Y^m$. For even $m$ its two analytic branches
are smooth. For odd $m$ its normalization is the
explicit smooth parametrization
$(Y,X,Z)=(u^2,\mathrm i u^m,Z)$.
Neither case has an exceptional curve in its
minimal normalization resolution. A point with
$d=1$ is smooth as well. Therefore the image of
$H$ has multiplicity exactly three.

\item The cyclic character calculation and all integer
bounds are independently recomputed in
\texttt{code/part4/quartic\_tangent.py}; no numerical
approximation enters the argument.

\end{enumerate}
\end{proof}


\subsection{The cubic double line and its weighted leading fibre}

\begin{proposition}[P4.53]\hypertarget{tag-02XS}{}\label{night:P4.53}\label{stmt:P4.53}
\StatementRecord{night:P4.53}{proved}{02XS}{\SeriesRef{night:P4.51}, \SeriesRef{night:P4.52}}

In the remaining positive ruled support, blow up the
multiplicity-three image of the contracted horizontal
section $H$, normalize the strict transform $S_1$, and
take its minimal resolution $T_1$. The strict transform
of $H$ is exceptional over a point of $S_1$. The entire
germ of $S_1$ there is normal and has an isolated simple
elliptic singularity. In particular $H$ cannot dominate
a curve in the ambient exceptional plane. No generic
transverse type is imposed on the original critical curve.
\end{proposition}
\begin{proof}

\begin{enumerate}[label={$\langle1\rangle$\arabic*.},leftmargin=*]
\item Use the comparison maps of P4.52, and let $D_1$ be the
effective total conductor relative to $\omega_{S_1}$.
Write $n=\operatorname{ord}_H(E)>0$ and
$a=\operatorname{coeff}_H(D_1)$. The pole identity is
$2=n+a$, so $(n,a)=(2,0)$ or $(1,1)$.

\item Suppose $H$ dominates a reduced component $L$ of the
projective cubic tangent curve. At its generic point
a transverse Weierstrass presentation has degree
$m$, the multiplicity of that component in the cubic.
Factoring over the completed valuation field of the
exceptional coordinate $r$ gives the exact degree
identity
\[
 \sum_{H'}\operatorname{ord}_{H'}(r)
                 [\mathbb C(H'):\mathbb C(L)]=m.             \tag{4.53a}
\]
The factors group the roots by valuation branches;
the two factors in each summand are the ramification
index and residue-field degree. Residue extensions
are separable in characteristic zero.

\item A positive-genus component of a cubic is a smooth
elliptic cubic of multiplicity one. The strict
transform is smooth along it, with $n=1,a=0$,
contrary to $2=n+a$. Every other reduced component
is rational. The finite map from $H$ then has degree
$h\ge2$. Equation (4.53a) excludes $n=2$, since
$m\le3$. Thus $n=a=1$.

\item At a generic point of $L$, the $h$ points on $H$
give distinct branches of the transverse plane germ,
each with conductor exponent one. The branch formula
\[
             c_i=2\delta_i+\sum_{j\ne i}I_{ij}
\]
forces $h=2$ and exactly two smooth branches meeting
once, with no additional branch. Each has $r$-order
one, so their total Weierstrass degree is two. Hence
$m=2$, and the tangent polynomial is $y^2\ell$,
where $y=0$ is the double line and $\ell$ a different
linear form.

\item Choose complementary coordinates $x,z$. The
restriction $\ell(x,z)$ is nonzero; any $y$ term in
$\ell$ has higher weight for the weights used below.
In a generic blowup chart along the double line,
write $v$ for the transverse coordinate. The
strict-transform equation begins
\[
                 \ell v^2+r(f_4|_L)+\cdots.
\]
A nonzero $f_4|_L$ would give ramification order two
and residue degree one. Therefore $f_4$ is divisible
by $y$. The two branches meet once, so after
$v=ru$ their distinct leading slopes solve
\[
                         \ell u^2+G_3u+G_5=0,
\]
where $G_3,G_5$ are forms of degrees three and five
in $x,z$. The discriminant
$\Delta=G_3^2-4\ell G_5$ is nonzero and nonsquare
in $\mathbb C(\mathbb P^1)$. The actual residue-field
extension of $H\to L$ is
$\mathbb C(\mathbb P^1)(\sqrt\Delta)$.

\item Consequently the original equation has weighted
leading polynomial
\[
          P_5=\ell(x,z)y^2+G_3(x,z)y+G_5(x,z),               \tag{4.53b}
\]
of degree five for weights $(1,2,1)$ on $(x,y,z)$,
with no smaller-weight terms. This identifies the
normalized projective component birationally with
$H$ through its function field.

\item Let $D=\{P_5=0\}\subset\mathbb P(1,2,1)$.
The quadratic is irreducible over $\mathbb C(x/z)$,
since the degree-two residue cover $H$ is connected.
The only possible further factor is a common
vertical factor of $\ell,G_3,G_5$. Since $\ell$
is linear, either $P_5$ is irreducible or
$P_5=\ell Q_4$, with $Q_4$ monic quadratic in $y$.
Both alternatives are reduced.

\item In the irreducible case, let $k_i$ be the
multiplicities of the distinct zeros of the
homogeneous degree-six discriminant. The normalized
double cover has genus
$(\#\{i:k_i\text{ odd}\}-2)/2$. Positive genus
therefore requires at least four odd multiplicities,
and gives
\[
              \sum_i(k_i-1)=6-\#\{i\}\le2.                  \tag{4.53c}
\]
Away from $\ell=0$ and the weighted vertex,
completing the quadratic identifies the projective
curve germ with $w^2-s^{k_i}$, up to a unit and
holomorphic coordinates. Its transverse Milnor
number is $k_i-1$. At $\ell=0$, if $G_3\ne0$ the
finite root is simple and the other point is at
the weighted vertex. If $G_3=0$, absence of a
common factor gives $G_5\ne0$, and $\Delta$ has
a simple zero. There is no additional singularity.
In the chart $y=1$ the vertex is the quotient of
$\mathbb C^2$ by $(x,z)\mapsto(-x,-z)$; the
lifted equation $\ell+G_3+G_5$ has nonzero
differential. The original polynomial also has
nonzero derivative in the $\ell$ direction on
the $y$-axis, so that axis is not critical.
The total transverse Milnor-number sum is thus
$M\le2$.

\item In the reducible case write
$G_3=\ell A_2$, $G_5=\ell A_4$. Then
$Q_4=y^2+A_2y+A_4$ has discriminant of degree four.
A connected positive-genus double cover requires
four distinct simple roots. Its normalization is
elliptic and $Q_4$ is smooth on the weighted plane;
it avoids the weighted vertex. The rational
component $\ell=0$ meets it in total length two,
read directly from the restricted monic quadratic.
Either there are two transverse intersections,
with total Milnor number two, or one tangency of
order two, with Milnor number three. There are no
other critical rays. Hence in all cases
\[
                              M\le3.                       \tag{4.53d}
\]

\item The weighted plane has Euler number three: its
strata $x\ne0$, $x=0,z\ne0$, and $x=z=0$ are
$\mathbb C^2$, $\mathbb C$, and a point.
In the irreducible case $e(D)\le e(H)\le0$.
In the reducible case
$e(D)=0+2-\#(Q_4\cap\{\ell=0\})\le1$.
Therefore
\[
               \chi\bigl(\mathbb P(1,2,1)\setminus D\bigr)
                         \ge2.                            \tag{4.53e}
\]

\item The weighted scalar action of a primitive fifth
root, with weights $(1,2,1)$, acts freely on
$F=\{P_5=1\}$. Its quotient is
$\mathbb P(1,2,1)\setminus D$: the only possible
nontrivial stabilizer of the weighted
$\mathbb C^*$ action is at the weighted vertex,
which belongs to $D$. This is an unramified
fivefold cover. The compact-core and character
calculation of P4.52 gives
\[
 \operatorname{rank}(T-I)
 \ge4\chi\bigl(\mathbb P(1,2,1)\setminus D\bigr)
 \ge8.
\]
The variation identity (4.52a) applies to every
one-dimensional hypersurface critical locus. Its
boundary mapping-torus estimate would give
$\operatorname{rank}(T-I)\le M\le3$ if the
intersection form were zero.

\item It is zero here by the actual anisotropic scaling
comparison. Set
\[
 D_\epsilon(x,y,z)=(\epsilon x,\epsilon^2y,\epsilon z).
\]
Then $\epsilon^{-5}f(D_\epsilon(x,y,z))$ converges
with derivatives to $P_5$. The compact-cycle flow
of P4.52 transports its pairings into the actual
local Milnor fibre, whose pairing vanishes by
P4.51. This contradiction excludes every dominating
$H$.

\item Thus $H$ is exceptional over the new normalization.
Its coefficient $a$ is positive, so $n=a=1$.
Adjunction gives
\[
                    0=2g(H)-2+(D_1-H)\cdot H.
\]
It follows that $H$ is elliptic and meets no other
positive pole component. A zero-coefficient
exceptional curve cannot meet a positive one:
its own intersection equation has nonnegative
canonical term on a minimal resolution and
nonnegative remaining terms, with a strictly
positive term at such a meeting. The connected
exceptional fibre is therefore $H$ alone and is
disjoint from the strict conductor. The
normalization is an isomorphism at its image;
another analytic hypersurface branch would
intersect this branch in a curve and create
conductor meeting $H$. Hence the entire
strict-transform germ is normal, with the single
smooth elliptic exceptional curve $H$. This is
precisely an isolated simple elliptic singularity.

\item The discriminant multiplicity partitions and the
exact rank bounds are recomputed in
\texttt{code/part4/cubic\_weighted.py} and its
executed output.

\end{enumerate}
\end{proof}


\subsection{Exclusion of the contracted horizontal section}

\begin{proposition}[P4.54]\hypertarget{tag-02XT}{}\label{night:P4.54}\label{stmt:P4.54}
\StatementRecord{night:P4.54}{proved}{02XT}{\SeriesRef{night:P4.51}, \SeriesRef{night:P4.52}, \SeriesRef{night:P4.53}, \SeriesRef{night:P4R.18}}

In the remaining positive ruled support of Part~4.2,
the horizontal pole section is not contracted by the
normalization resolution. It consequently maps with
degree two or three onto a rational normalized conductor
component whose generic transverse germ is $D_4$.
No transverse-type restriction is imposed on the other
conductor component.
\end{proposition}
\begin{proof}

\begin{enumerate}[label={$\langle1\rangle$\arabic*.},leftmargin=*]
\item By P4.53, a contracted section would be the sole smooth
elliptic exceptional curve $H$ of an isolated normal
hypersurface singularity on the first ambient strict
transform. Its self-intersection is $-d<0$, and the
actual exceptional coordinate $r$ has order one along
$H$. Write $\pi:\widetilde V\to(V,0)$ for this local
resolution. We establish the following three coordinate
weights and leading degrees, retaining the actual $r$:
\[
 (1,1,1;3),\qquad(1,1,2;4),\qquad(1,2,3;6).
                                                        \tag{4.54w.1}
\]

\paragraph{1. Analytic formal functions and the exact associated graded ring}

\item Put \(L=\mathcal O_H(-H)\), of degree \(d\), and
\[
 I_m=(\pi_*\mathcal O_{\widetilde V}(-mH))_0
     =\{a\in\mathcal O_{V,0}:\operatorname{ord}_H(a)\ge m\}.
\]
For \(m\ge1\), the restrictions of
\(\mathcal O(-mH)\) to successive infinitesimal neighbourhoods
of \(H\) have graded factors \(L^{m+j}\), \(j\ge0\).
Each has zero first cohomology, by Serre duality on the elliptic
curve and its strictly positive degree. Their extensions also
have zero first cohomology.

\item The analytic formal-functions theorem is Grauert1960,
Section6, Hauptsatz IIa, printed p.59; Hauptsatz I on the same
page supplies coherence \cite[\S6, Hauptsätze I and IIa]{Grauert60}. The maximal-ideal pullback and the
ideal of \(H\) have the same radical. On a finite cover of
the compact exceptional curve, Noetherianity gives one uniform
power of the latter contained in the former. Their powers are
therefore cofinal. Applying the theorem with
\(\mathcal O(-mH)\) gives a zero completion for
\(R^1\pi_*\mathcal O(-mH)\). This coherent module is
supported at the origin, so it is zero. Consequently
\[
 I_m/I_{m+1}\simeq H^0(H,L^m)\qquad(m\ge1),             \tag{4.54w.2}
\]
through the actual restriction maps in
\[
 0\longrightarrow\mathcal O(-(m+1)H)
 \longrightarrow\mathcal O(-mH)
 \longrightarrow L^m\longrightarrow0.
\]
These maps respect multiplication. Since normality gives
\(\pi_*\mathcal O_{\widetilde V}=\mathcal O_V\), and a
function vanishes at the contracted point exactly when its
restriction to \(H\) vanishes, \(I_1=\mathfrak m\).
The degree-zero piece is \(\mathbb C\). Thus
\[
 \operatorname{gr}_{I}\mathcal O_{V,0}
   =\bigoplus_{m\ge0}H^0(H,L^m)                         \tag{4.54w.3}
\]
as graded algebras, not only as vector spaces.

\item The same formal-functions theorem in degree zero identifies
\[
 \widehat{\mathcal O}_{V,0}
      \simeq\varprojlim_m H^0(mH,\mathcal O_{mH})
      \simeq\varprojlim_m\mathcal O_{V,0}/I_m.          \tag{4.54w.4}
\]
The second equality uses the just-proved first-direct-image
vanishing. This supplies the exact completeness and separation
needed in the leading-equation argument below; no unproved equivalence between a
divisorial and an ordinary degree filtration is being used.

\item Because \(\mathfrak m^2\subset I_2\), (4.54w.2) gives a
surjection
\(\mathfrak m/\mathfrak m^2\to H^0(H,L)\).
The source has dimension three and the target dimension \(d\).
Hence \(d\le3\).

\paragraph{2. The three section rings and their actual elliptic curves}

\item For \(d=3\), Riemann--Roch shows that \(L\) separates
points and tangent vectors: subtracting any length-two
subscheme leaves a degree-one line bundle with one section.
The complete linear series therefore embeds \(H\) as a smooth
plane cubic. The restriction sequence on \(\mathbb P^2\),
using \(H^1(\mathbb P^2,\mathcal O(k))=0\), gives its whole
section ring as
\[
 \mathbb C[x,y,z]/(F_3),\qquad \deg(x,y,z)=(1,1,1),
\]
where \(F_3=0\) is that smooth cubic.

\item For \(d=2\), the complete pencil is base-point free and gives
a degree-two map \(p:H\to\mathbb P^1\), with
\(L=p^*\mathcal O(1)\). Trace splits
\(p_*\mathcal O_H=\mathcal O\oplus B^{-1}\).
Its Euler characteristic is zero, so \(\deg B=2\).
The algebra structure is therefore a double-cover equation
\(z^2=F_4(x,y)\), where the binary quartic has four distinct
zeros because the double cover is smooth. Projection formula
gives the entire section ring with generator weights
\((1,1,2)\) and that single degree-four relation.

\item For \(d=1\), write \(L=\mathcal O_H(p)\).
Riemann--Roch gives a degree-one section \(x\), one new
degree-two section \(y\), and one new degree-three section
\(z\). Dividing by powers of \(x\) gives functions of pole
orders two and three at \(p\). Their Weierstrass relation
has weighted degree six. Equivalently, after graded changes
only for describing this section ring, it is
\[
 z^2=y^3+A x^4y+B x^6,
 \qquad 4A^3+27B^2\ne0,
 \qquad\deg(x,y,z)=(1,2,3).
\]
The pole-order count shows that the monomials with at most one
factor \(z\), reduced by this relation, give every space
\(H^0(H,L^m)\). Thus this is the whole section ring. Its
weighted projective curve is the actual smooth \(H\).
The coordinate choice below does not require completing the
square or applying any of these optional graded simplifications.

\paragraph{3. Retaining the given coordinate using only permitted changes}

\item Start with analytic ambient coordinates \((r,u,v)\), retaining
the given \(r\). Their images form a basis of
\(\mathfrak m/\mathfrak m^2\). The map to
\(I_1/I_2=H^0(L)\) has the rank just computed.

\item When \(d=3\), it is an isomorphism, so any linear complement
to \(r\) supplies the other two valuation-one coordinates.

\item When \(d=2\), choose a linear coordinate \(s\) whose
leading section complements that of \(r\), and choose the
remaining independent linear coordinate \(v_0\) in the
kernel of the degree-one restriction. Thus
\(\operatorname{ord}_H(v_0)\ge2\). Its degree-two leading
section must be independent of \(r^2,rs,s^2\). Otherwise
all degree-two initial sections of analytic functions would
lie in that three-dimensional span: a Taylor expansion shows
that every term of ordinary order at least three has valuation
at least three, and the linear terms of valuation two are
accounted for by \(v_0\). This contradicts
\(h^0(L^2)=4\). Therefore \(v_0\) has valuation exactly
two and is the required indecomposable degree-two generator.
The coordinates \((r,s,v_0)\) were obtained by an invertible
linear change.

\item When \(d=1\), first choose two independent linear coordinates
\(u_0,v_0\) in the kernel of the degree-one restriction.
At least one linear combination, denoted \(s\), has a
degree-two initial section independent of \(r^2\), since
\(h^0(L^2)=2\). Choose the remaining independent linear
coordinate \(v_1\), and subtract a scalar multiple of \(s\)
so that its degree-two initial term is \(\alpha r^2\).
Set
\[
                         v_2=v_1-\alpha r^2.
\]
This is a triangular analytic coordinate change, and
\(\operatorname{ord}_H(v_2)\ge3\). Its degree-three
initial section must be independent of \(r^3,rs\).
Indeed, if it were absent or in that span, any Taylor series
of valuation three would have degree-three initial section
in that same span: the possible terms are precisely the
linear \(v_2\) term and the weighted-degree-three monomials
\(r^3,rs\); its valuation-one and valuation-two terms
must first vanish. This contradicts \(h^0(L^3)=3\).
Thus \(v_2\) has valuation exactly three and supplies the
indecomposable degree-three generator. The three ambient
coordinates remain analytically invertible, because only
an invertible linear change and the displayed quadratic
translation were performed.

\item This argument proves directly that the chosen analytic
coordinates have the prescribed orders and generate the
graded ring. It avoids a circular appeal to filtered
Nakayama or to a coordinate normalization that might move
the fixed \(r\).

\paragraph{4. The actual weighted leading equation}

\item Write the defining analytic equation in the coordinates just
constructed, with the weights in (4.54w.1). Its weighted initial
form is the single defining relation of the graded section
ring, up to a nonzero scalar. Here is an explicit proof of that
last assertion.

\item Let \(F_D\) be the graded relation from the section-ring calculation above, expressed
in the actual chosen initial generators. The analytic function
\(F_D(r,s,v)\) has valuation at least \(D+1\).
At each successive valuation, (4.54w.3) writes its initial
section as a weighted homogeneous polynomial in the three
generators. Subtract that polynomial. This produces a formal
series \(G=F_D+F_{D+1}+\cdots\) whose image in every
\(\mathcal O_{V,0}/I_m\) is zero. Since all generator
weights are at most three, the correction polynomials tend
to zero in ordinary power-series degree; their sum is a
well-defined formal power series. By (4.54w.4), its image in
\(\widehat{\mathcal O}_{V,0}\) is zero.

\item The completed hypersurface ideal is principal, generated by
the actual analytic equation \(f\). Thus \(G=h f\) for
a formal series \(h\). Taking weighted initial forms gives
\(F_D=\operatorname{in}(h)\operatorname{in}(f)\).
The polynomial \(F_D\) is irreducible because its weighted
projective curve is the smooth connected elliptic curve
\(H\). Also \(f\) is nonconstant and vanishes at the
origin. Hence \(\operatorname{in}(h)\) is a nonzero
constant, and \(\operatorname{in}(f)\) is a scalar multiple
of \(F_D\). This proves (4.54w.1) for the actual defining
equation, with all higher terms retained.

\paragraph{5. The exact descent to the original blowdown coordinates}

\item In the chart with original coordinates
\(x=r u\), \(y=r v\), every change above has the form
\[
 s=a u+b v+c r,\qquad
 v'=a'u+b'v+c'r-\alpha r^2,
 \qquad ab'-a'b\ne0.
\]
It descends to the actual holomorphic coordinate map
\[
 \begin{aligned}
 x'&=r s=a x+b y+c r^2,\\
 y'&=r v'=a'x+b'y+c'r^2-\alpha r^3,\\
 r'&=r.
 \end{aligned}                                        \tag{4.54w.5}
\]
Its Jacobian determinant at the origin is \(ab'-a'b\),
so the inverse-function theorem supplies an analytic inverse.
This is the required morphism preserving the original
parameter and blowdown structure.

\paragraph{Leading polynomials in the original coordinates.}
Apply this calculation to the simple elliptic point supplied
by P4.53. The original equation is $r^3f_1$, and its original
coordinates are $(x,y,z)=(rs,rv,r)$. The three pairs of
original weights and degrees are
\[
 (2,2,1;6),\qquad (2,3,1;7),\qquad (3,4,1;9).
\]
For a strict-transform monomial $r^as^bv^c$, multiplication
by $r^3$ gives $x^by^cz^{3+a-b-c}$. Its original weighted
degree is three plus its strict-transform weighted degree.
Analyticity of the original equation makes every
negative-power coefficient zero. Thus the actual weighted
leading term is a polynomial of the indicated degree.
On $z\ne0$, the ratios
\[
     \left(x/z^{1+w_s},\ y/z^{1+w_v}\right)
\]
identify its projective component with the strict-transform
elliptic relation on $r\ne0$. This is a birational map
of the actual residue-field curves; their smooth
normalizations are the same elliptic curve $H$.

\item Anisotropic scaling with these weights and the compact-cycle
flow in P4.52 transports every pairing of the leading
Milnor fibre into the actual local torus fibre. By P4.51,
all these pairings must vanish. We now contradict this
in each of the three cases.

\paragraph{Degree six.}
The leading polynomial is exactly
\[
                    P_6(x,y,z)=G_3(x,y,z^2),
\]
where $G_3(X,Y,Z)$ defines a smooth projective cubic.
The map $(x,y,z)\mapsto(x,y,z^2)$ gives a finite double
cover from $F=\{P_6=1\}$ to $\{G_3=1\}$, branched over
\[
                       B=\{G_3(x,y,0)=1\}.
\]
The homogeneous cubic fibre has Euler number nine: it
is a threefold unramified cover of the complement of
an elliptic cubic in $\mathbb P^2$, whose Euler number
is three. Let $r\in\{1,2,3\}$ be the number of distinct
zeros of the nonzero binary cubic $G_3(x,y,0)$ on
$\mathbb P^1$. The curve $B$ is a threefold cover of
that punctured projective line. Hence
\[
          \chi(B)=3(2-r),\qquad \chi(F)=18-\chi(B)=12+3r.
\]

\item Weighted monodromy has order six. Its only
nonidentity element with a fixed set is its cube,
which negates $z$ and fixes $x,y$; its fixed set is
$B$. The other four nonidentity elements have no
fixed point on $F$. We choose an invariant
semialgebraic compact core large enough for $F$
and $B$. The finite-order Lefschetz character of
$H_0-H_1+H_2$ is therefore $\chi(F)$ at the identity,
$\chi(B)$ at the cube, and zero elsewhere. Its
trivial-character multiplicity is
\[
                       \frac{\chi(F)+\chi(B)}6=3.
\]
Since $H_0$ has only the trivial character, the
nontrivial-character dimension in $H_2$ is at least
$\chi(F)-3=9+3r\ge12$.

\item Every nonzero critical orbit of $P_6$ lies on $z=0$
at a tangency of the cubic with $Z=0$. If the
intersection orders are $e_i$, then $\sum e_i=3$.
At a tangency, smoothness of $G_3$ makes its
$Z$-derivative a unit, and the exact transverse
germ is $u^{e_i}+z^2$, after holomorphic coordinates
and a unit. Consequently
\[
                       M=\sum(e_i-1)=3-r\le2.
\]
This also shows reducedness of $P_6$: a repeated
factor would give a two-dimensional critical set,
whereas the derivative calculation gives only
these critical orbits. The variation inequality
(4.52b) contradicts the lower bound twelve for
the nontrivial monodromy rank.

\paragraph{Weighted curve cohomology.}
We record the two curve facts needed for the last
cases, with their quotient comparisons. For
$\mathbb P(a,b,c)=\operatorname{Proj}\mathbb C[u,v,w]$
and a reduced equation of weighted degree $e$,
the exact graded sequence sheafifies to
\[
 0\longrightarrow\mathcal O_{\mathbb P}(-e)
 \longrightarrow\mathcal O_{\mathbb P}
 \longrightarrow\mathcal O_D\longrightarrow0.
\]
The three standard affine charts give its \v Cech
complex. The middle cohomology of every
$\mathcal O_{\mathbb P}(k)$ is zero. The top
cohomology of $\mathcal O_{\mathbb P}(-e)$ has
basis the Laurent monomials $u^{-i}v^{-j}w^{-k}$,
$i,j,k\ge1$, of weighted degree $-e$. Thus
\[
 h^1(\mathcal O_D)=
 \#\{\hbox{nonnegative weighted monomials of degree }
                         e-(a+b+c)\},
 \qquad h^0(\mathcal O_D)=1.                         \tag{4.54a}
\]
In particular the arithmetic genus is zero below
degree $a+b+c$, and is one in degrees $a+b+c$
and $a+b+c+1$ when there is exactly one variable
of weight one.

\item At a cyclic quotient vertex of prime order two
or three, with no invariant linear ambient
coordinate, a smooth coarse curve has a smooth
plane pullback. To prove this, its normalized
branches form one orbit. A single singular
branch of multiplicity $m\ge2$ has a faithful
normalization action of order $p$; every invariant
ambient function has order at least $2m>p$,
whereas a regular quotient requires a function
of normalization order $p$. For $p$ transitively
permuted branches the stabilizer is trivial,
and every invariant function has order at least
two, whereas a regular quotient requires order
one. Both alternatives are impossible. The
remaining single branch is smooth, and its
reduced equation has the character of a nonzero
linear ambient coordinate. These arguments
use the actual analytic invariant rings:
averaging makes the invariants functor exact,
so every quotient-curve function lifts to an
invariant ambient function.

\paragraph{Degree seven.}
Consider weights $(1,2,3)$ and degree seven.
Formula (4.54a) gives arithmetic genus zero for
curves of degree at most five and arithmetic
genus one in degrees six and seven. Thus a
positive-genus component has degree six or
seven, genus one, and multiplicity one. There
is at most one further degree-one factor;
the polynomial is reduced.

\item In degree seven the whole curve $D$ is
irreducible with arithmetic and normalization
genus one. Its total delta invariant is zero,
so it is a smooth coarse elliptic curve.
At the quotient vertices of orders two and
three, the preceding argument makes its plane
pullback smooth. Elsewhere the ambient plane
is smooth. Thus the affine cone has no
nonzero critical orbit and $M=0$.

\item In degree six write $D=E\cup\{z=0\}$, with
$z$ the weight-one coordinate. The curve $E$
is smooth coarse elliptic. It cannot contain
either quotient vertex: its degree-six
equation has trivial character there, while
both local-coordinate characters are
nontrivial, so a smooth lift would require
an impossible invariant linear equation.
The genus formula for $D$ gives total delta
one. Its two smooth components therefore
meet once transversely in the smooth part
of the ambient weighted plane. There is
one ordinary-double critical orbit and
$M=1$.

\item The weighted plane has Euler number three,
from its affine-cell stratification. In the
two cases $e(D)=0$ or one, so its complement
has Euler number three or two. Both quotient
vertices belong to $D$, since degree seven
is divisible by neither two nor three.
The weighted $\mathbb C^*$ action is free
off $D$, so $\{P_7=1\}$ is an unramified
sevenfold cover of that complement. The
free order-seven monodromy gives nontrivial
$H_2$ rank at least eighteen or twelve,
respectively. This contradicts (4.52b),
whose upper bound is zero or one.

\paragraph{Degree nine.}
For weights $(1,3,4)$, (4.54a) gives
arithmetic genus zero below degree eight
and one in degrees eight and nine.
A positive-genus component must have one
of these two degrees, and multiplicity one.

\item An elliptic degree-eight component is
impossible. It necessarily contains the
weight-three vertex, since eight is not
divisible by three. The local plane
characters are $1,1$ modulo three, whereas
its equation has character two. The
prime-order quotient argument would give
a smooth lift, whose equation must have
character one. This is a contradiction.

\item Thus the entire degree-nine curve $D$ is
irreducible and smooth coarse elliptic.
It avoids the weight-three vertex: its
equation there has character zero,
incompatible with a smooth plane lift.
At the weight-four vertex the plane
characters are $1,3$ and its equation
has character one. We check the lift
without assuming that four is prime.
The transitive orbit of normalized
branches has size one, two, or four.

\item For four branches the stabilizer is
trivial and every invariant maximal-ideal
function has order at least two on a
branch, contradicting regularity of the
quotient. For two branches the stabilizer
has order two. A singular branch gives
invariant orders at least four, exceeding
two. Two smooth branches would have a
quadratic initial equation, but the
quadratic characters of the ambient
action are only zero and two, never one.
For one branch the normalization action
is faithful of order four. Both ambient
coordinate orders are odd, since their
characters are one and three. A singular
branch therefore has multiplicity at
least three, and all invariant functions
have order at least six, exceeding four.
These contradictions leave a single
smooth branch.

\item The degree-nine cone is therefore smooth
away from its vertex. It is singular
at the vertex, because no variable has
weight nine. Its hypersurface critical
point is isolated, so the isolated-point
Morse transport of P4.52 gives a spherical
class of square $-2$ in its Milnor fibre.
This again contradicts the zero pairing
transported from the actual torus.

\item All three possibilities for the sole
simple elliptic exceptional curve have
now been excluded. Together with P4.53
this proves that the horizontal pole
section cannot be contracted. P4R.18
then gives the asserted generic $D_4$
image and degree two or three. The
other conductor component retains its
unrestricted transverse type.

\item The exact weighted cohomology counts,
order-six characters, quotient quadratic
characters and an explicit smooth cubic
are recomputed in
\path{code/part4/}\allowbreak\path{simple_elliptic_weighted.py}
and its executed outputs.

\end{enumerate}
\end{proof}


\subsection{The conductor graph and the horizontal spectral cover}

\begin{proposition}[P4.55]\hypertarget{tag-02XU}{}\label{night:P4.55}\label{stmt:P4.55}
\StatementRecord{night:P4.55}{proved}{02XU}{\SeriesRef{night:P4.31}, \SeriesRef{night:P4.34}, \SeriesRef{night:P4.54}, \SeriesRef{night:P4R.13}, \SeriesRef{night:P4R.18}, \SeriesRef{night:P4R.19}}

In the remaining positive ruled case, the image $A$ of the
horizontal pole section is a smooth rational curve. The map
$H\to A$ has degree exactly two. The third generic
normalization preimage along $A$ is supplied by one rational
vertical conductor component, mapping isomorphically to $A$.
The other image component retains its unrestricted transverse
type. A separated self-identification point on either image
component is impossible.
\end{proposition}
\begin{proof}

\begin{enumerate}[label={$\langle1\rangle$\arabic*.},leftmargin=*]
\item We write $T\to N\to S$ for the minimal normalization
resolution. By P4.54 and P4R.18, the pole section $H$ is
noncontracted, has genus $g\ge1$, and maps with degree two
or three to $A^\nu=\mathbb P^1$. The full generic transverse
germ along $A$ is $D_4$. The reduced image conductor is
$C=A\cup B$, with two rational normalized components and
one graph cycle.

\item The reduced normalization conductor $D$ has a tree as its
point-identification graph. Its positive-genus component
$H$ is smoothly embedded: the ruling descends to $N$, and
restriction of that ruling is an inverse to $H\to N$.
In particular no two points of $H$ are identified by the
normalization resolution. Every other normalized conductor
component is rational. These are the actual conductor
curves; vertical components and their identifications are
retained throughout the argument.

\item At any point of $C$, points of $N$ over it correspond to
local irreducible hypersurface branches of $S$. Two such
branches intersect along a curve, by the height bound for
two equations in a regular local ring of dimension three.
Their nonempty supports among the local analytic branches
of $C$ therefore pairwise intersect.

\paragraph{Separated self-identifications.}
Suppose $p\in B\setminus A$ has two analytic branches.
Let $D_B$ be the union of the conductor components mapping
to $B$, with $c_B$ normalized components and $s_B\ge1$
connected components. Its graph is a forest. Every component
maps finitely and surjectively to $B^\nu$, so the numbers
$r_+,r_-$ of its normalized points over the two preimages
of $p$ satisfy $r_+,r_-\ge c_B$.

\item There are $n$ normalization points of $S$ over $p$.
Their nonempty supports in $\{+,-\}$ pairwise intersect.
They have a common element: a singleton forces its element
into every support, and without a singleton all supports
are the entire two-element set. After interchanging signs,
$r_+\ge n$. The graph contribution at $p$ is consequently
\[
                  r_++r_--n\ge c_B.
\]
The whole forest has contribution $c_B-s_B<c_B$, a
contradiction. The same proof applies to a self-identification
of $A$ away from $B$, using all components mapping to $A$.
Thus the image cycle consists either of two distinct
intersections of $A$ and $B$, or of one three-branch point
belonging to both components.

\paragraph{The actual primary conductor and smoothness of $A$.}
Let $C_{A,\mathrm{sch}}$ be the primary component of the
actual image conductor supported on $A$. The primary
restriction maps of P4R.13, used in P4R.18, give
\[
                    g-1\le\chi(\mathcal O_{C_{A,\mathrm{sch}}}).
                                                            \tag{4.55a}
\]
The generic ordinary-triple conductor is the transverse
maximal ideal squared. Purity of the actual conductor
quotient identifies this primary ideal with $P_A^{(2)}$,
as in P4.31. The exact normalized conormal calculation of
P4.34 therefore gives
\[
 \chi(\mathcal O_{C_{A,\mathrm{sch}}})
       =5-\sum_p\bigl(\delta(A,p)+\tau_p-r_p\bigr).
                                                            \tag{4.55b}
\]
A planar point contributes exactly $5\delta(A,p)$, and a
nonplanar singular point with at most two branches contributes
at least seven. There cannot be three branches of $A$ at one
point: the identity
$\sum_p(r_p(C)-1)=2$ would then leave no identification for
the necessary connection with $B$; placing that connection
at the same point gives at least four branches and also
exceeds the identity.

\item Since $g\ge1$, (4.55a)--(4.55b) bound the total defect by
five. If $A$ is singular, it is planar everywhere, has total
delta invariant one, and $g=1$. It is then a locally complete
intersection curve in $Y$. Its conormal
$Q=P_A/P_A^2$ is locally free of rank two. Let
$j:\mathbb P^1\to A$ be its normalization and put $E_0=j^*Q$.
Adjunction gives
\[
             \omega_A=(\omega_Y|_A)\otimes(\det Q)^{-1}.
\]
The degree of $j^*\omega_Y$ is zero because $b_2(Y)=0$.
Also $p_a(A)=1$, so $\deg j^*\omega_A=0$. The degree of a
line bundle is preserved by finite birational normalization:
the normalization sequences for that line bundle and for
$\mathcal O_A$ have the same finite-length correction.
Consequently
\[
                              \deg E_0=0.                 \tag{4.55c}
\]

\item The actual reduced equation has generic order three along
$A$. Powers of a complete-intersection curve ideal are
unmixed, so $P_A^{(3)}=P_A^3$. Its third normal term is an
actual global section
\[
 q_3\in H^0\!\left(A,
          \operatorname{Sym}^3 Q\otimes\mathcal O_Y(S)|_A\right).
\]
The line twist has degree zero and becomes trivial on
$\mathbb P^1$. On the quotient projective bundle
$\mathbb P(E_0)$, with tautological class $\xi$ and fibre
class $F$, the pulled-back cubic has zero divisor of class
$3\xi$. Remove its common zero fibres, of total degree
$b\ge0$. The residual spectral curve $\Gamma$ has class
$3\xi-bF$. Its horizontal components are reduced, because
the generic $D_4$ cubic has three distinct roots.

\item We specify the spectral morphisms. At a generic point of
$H$, the first normal derivative of $T\to S\subset Y$
gives the direction of its smooth transverse branch and
hence a point of $\mathbb P(E_0)$. On the smooth compact
curve $H$, its two local coordinates extend after cancelling
their common power. This gives a holomorphic map over the
actual map $H\to\mathbb P^1$. Distinct generic preimages
on $H$ give distinct cubic roots. Hence its map onto the
spectral component $\Gamma_H$ is finite and birational,
and $H$ is that component's normalization. In degree three,
$\Gamma=\Gamma_H$. In degree two, precisely one remaining
rational vertical conductor component supplies the third
root, with degree one over $A$. Its direction curve
$\Gamma_R$ is a section and is isomorphic to $\mathbb P^1$.
Thus $\Gamma=\Gamma_H\cup\Gamma_R$. These exhaust the
three generic roots.

\item For a rank-two bundle of degree $e$ on $\mathbb P^1$,
\[
 \xi^2=e,\qquad \xi F=1,\qquad F^2=0,\qquad
 K_{\mathbb P(E_0)}=-2\xi+(e-2)F.
\]
Adjunction gives the exact arithmetic genus
\[
                    p_a(3\xi-bF)=3e-2-2b.               \tag{4.55d}
\]
With (4.55c) this is $-2-2b$. In degree three it is at
least $g\ge1$. In degree two the exact union sequence gives
\[
 p_a(\Gamma)=p_a(\Gamma_H)+p_a(\Gamma_R)
           +\operatorname{length}(\Gamma_H\cap\Gamma_R)-1
       \ge g-1\ge0.
\]
Both alternatives are impossible. Therefore $A$ is smooth.

\paragraph{The three-branch image point in degree three.}
Suppose now that $H\to A$ has degree three. Then $H$ is
the only conductor component mapping onto $A$. Since $A$
is smooth, a single three-branch image point consists of
one $A$ branch and two $B$ branches, denoted $+$ and $-$.

\item Let $c_B$ be the number of normalized conductor components
mapping to $B$. The entire conductor has $1+c_B$
normalized components and a tree graph, so its total
point-identification contribution is $c_B$. At the proposed
point write $r_A\ge1$, $r_+,r_-\ge c_B$ for the numbers
of conductor branches, and $n$ for the number of normalization
points. Each normalization point has at most one $A$
conductor branch, because $H$ is embedded.

\item A singleton support $\{A\}$ forces every point to carry
$A$. Then $n=r_A$, and the graph contribution is
$r_++r_-\ge2c_B>c_B$. A singleton $\{+\}$ forces
$r_+\ge n$, giving contribution at least
$r_A+r_-\ge1+c_B$; the same applies to $\{-\}$.
Thus every support has at least two elements. Consequently
$r_A+r_++r_-\ge2n$, and its graph contribution is at least
\[
       \frac{r_A+r_++r_-}{2}
                         \ge c_B+\frac{r_A}{2}>c_B.
\]
This is again impossible. Only two distinct $A$--$B$
junctions remain in degree three.

\paragraph{Both degree-three junctions are common cubic zeros.}
Let these junctions be $p,q$. If a local normalization
branch over $p$ has conductor only over $A$, pairwise
hypersurface intersections force every normalization point
over $p$ to carry $A$. Thus $n=r_A$, and the graph
contribution at $p$ is $r_B\ge c_B$, using the entire
budget. At $q$ some normalization point must carry both
$A$ and $B$: otherwise an $A$-only point and a $B$-only
point have disjoint supports. That point contributes
positively to the graph, contradicting the budget. There
is therefore no $A$-only local branch at either junction.

\item We retain all factors in the following local calculation.
Write the smooth curve as $A=\{x=y=0\}$. If its nonzero
cubic normal form has a simple root, use the normal blowup
chart $x=yu$. The implicit-function theorem on the strict
equation gives that root as $u=\alpha(y,z)$. The actual
function has the analytic factor
\[
                         f_1=x-y\alpha(y,z).
\]
This branch is smooth. Exact division writes $f=f_1f_2$,
and restriction of $f_2$ to this branch is $y^2$ times a
unit: the other cubic directions differ from the simple
root at the point. Its intersections with all other
branches are therefore supported only on $A$. It is the
forbidden $A$-only normalization branch.

\item Every nonzero binary cubic has a simple root unless it is
a cube. P4R.19 excludes a nonzero cube by its exact finite
weighted maps and compact intersection class. Hence the
whole cubic tensor vanishes at both $p$ and $q$.

\paragraph{The degree-three genus contradiction.}
For the now smooth curve $A=\mathbb P^1$, adjunction gives
$\deg N_{A/Y}=-2$, so $E=N_{A/Y}^*$ has degree two.
The actual cubic has at least two common zeros. Its residual
spectral divisor is $3\xi-bF$ with $b\ge2$, and is
irreducible with normalization $H$. Formula (4.55d) gives
\[
                    p_a(\Gamma)=4-2b\le0,
\]
contradicting $p_a(\Gamma)\ge g\ge1$. Degree three is
excluded. The degree-two alternative and its unique rational
third sheet are precisely the possibilities already identified
by the spectral morphism above.

\item The exact intersection polynomial and the graph inequalities
are recomputed in
\path{code/part4/d4_graph_spectral.py}; its saved output
reports PASS. The inequalities in this proof establish the
unrestricted graph counts.

\end{enumerate}
\end{proof}


\subsection{The other image curve with at most two generic branches}

\begin{proposition}[P4.56]\hypertarget{tag-02XV}{}\label{night:P4.56}\label{stmt:P4.56}
\StatementRecord{night:P4.56}{proved}{02XV}{\SeriesRef{night:P4.55}, \SeriesRef{night:P4R.7}, \SeriesRef{night:P4R.9}, \SeriesRef{night:P4R.19}}

In the remaining case of P4.55, the other conductor image
component $B$ has at least three generic normalization
branches. This conclusion imposes no preliminary bound on
its transverse multiplicity, conductor exponent, or tangency
order.
\end{proposition}
\begin{proof}

\begin{enumerate}[label={$\langle1\rangle$\arabic*.},leftmargin=*]
\item We retain the smooth curve $A$, its degree-two horizontal
cover $H$, and its degree-one vertical cover $R$. All other
normalized conductor components map to $B$. Let $r_B$ be
the generic normalization branch number along $B$, and
$c_B$ the number of global conductor components over it.
Their finite degrees sum to $r_B$. We exclude $r_B=1,2$.
The normalized conductor has $2+c_B$ components and a tree
graph, with total identification contribution $1+c_B$.
Distinct points on one normalized component cannot be
identified, since that would create a graph loop.

\paragraph{A length-two intersection calculation.}
Let $F,G$ be two local hypersurface branches whose
scheme-theoretic intersection is supported on the smooth
curve $A$ and has generic length $m\in\{1,2\}$ along it.
The ring
\[
                 Q=\mathbb C\{x,y,z\}/(F,G)
\]
is a one-dimensional complete intersection. Its only
associated prime is the ideal of $A$. Choose $z$ as a
coordinate on $A$. Then $Q$ is finite and torsion-free
over $\mathbb C\{z\}$, hence free of rank $m$. The
reduction map to $\mathbb C\{z\}$ splits through the
chosen coordinate.

\item For $m=1$, the ring is $\mathbb C\{z\}$, and the linear
parts of $F,G$ have rank two. Both branches are smooth.
For $m=2$, the nilradical is a free rank-one module over
$\mathbb C\{z\}$. Its square vanishes generically and
therefore vanishes, by torsion-freeness. A generator
$\epsilon$ gives the exact analytic algebra isomorphism
\[
                    Q\simeq\mathbb C\{z,\epsilon\}/(\epsilon^2).
                                                            \tag{4.56a}
\]
Thus $Q$ has embedding dimension two, so one of $dF,dG$
is nonzero. If $G$ has generic normal order two along $A$,
it belongs to the ordinary square of its smooth ideal and
$dG$ vanishes on all of $A$. Then $dF\ne0$ and $F$ is
smooth. We use this calculation on the original intersection
scheme, retaining its multiplicity.

\paragraph{Two distinct image junctions.}
At a two-branch image point, let $r_A,r'_B$ count the
normalized conductor points over its two branches, and
let $n$ count normalization points of $S$. Pairwise
intersection of the hypersurface branches gives a common
element in their supports in $\{A,B\}$. Hence
$n\le\max(r_A,r'_B)$, and the graph contribution is at
least $\min(r_A,r'_B)$. Here $r_A=2$ or three, from the
one or two $H$ points and the unique $R$ point.

\item If $c_B=2$, both degrees are one, and $r'_B=2$ at either
junction. The two contributions total at least four,
exceeding the budget three. If $c_B=1$ and its degree is
two, an unramified junction costs at least two and the
other at least one, exceeding the budget two. Thus one
component $V$ remains, with exactly one normalized point
over each junction.

\item If $B$ belongs to every branch support at a junction,
there is only one normalization point there and its
contribution is $r_A\ge2$. The other junction makes this
impossible. Therefore $A$ belongs to every support. The
budget forces contribution exactly one at each junction:
every normalization point has one $A$ branch, and the
sole $V$ point is paired with one of them. Connectedness
forces the graph path $H$--$V$--$R$, with its edges at
the two different junctions.

\item At the $R$--$V$ junction, let $F$ be the corresponding
hypersurface branch. The other branches come from $H$
and carry only $A$. If its two preimages are distinct,
intersection with either such branch has support $A$ and
generic length one, because the full generic $A$ germ
is an ordinary triple point. The preceding calculation
makes $F$ smooth. If the two $H$ preimages coalesce,
their one hypersurface branch $G$ has generic normal
order two along $A$, and the intersection with $F$ has
length two. Formula (4.56a) again makes $F$ smooth.

\item This is impossible. The branch $F$ carries all the nearby
generic $B$ preimages: either the single branch of a
singular unibranch transverse germ, or both branches
when $V\to B^\nu$ has degree two and ramifies at this
point. In either case $F$ is intrinsically nonnormal
along the nearby $B$ germ. It cannot be smooth in a
neighbourhood of the junction.

\paragraph{A three-branch image point with one $B$ component.}
The image branches are $A,B_+,B_-$, since $A$ is smooth.
If $c_B=1$, the normalized points on $V$ over the two
$B$ signs are distinct and cannot be identified in $N$.
No normalization point can therefore carry both signs.
Pairwise hypersurface intersections force every point
carrying either sign to carry $A$ as well.

\item The graph budget two now forces exactly one point over
each sign, each paired with exactly one $A$ branch, and
no other identification. The pairs are $H$--$V$ and
$R$--$V$. An extra $H$ preimage is $A$-only. At the
$R$--$V$ point, the other $H$ branch or branches intersect
its hypersurface branch only along $A$, because their
$B$ supports have the opposite sign or are empty. The
same length-one or length-two calculation makes this
branch smooth, contradicting its intrinsic nonnormality
along the $B$ germ whose generic preimages it contains.

\paragraph{A three-branch point with two $B$ components.}
Write the two degree-one components as $V,W$. If $h=1$
or two is the number of $H$ points over the $A$ value,
the total normalized branch count is
\[
                         T=(h+1)+2+2=h+5.
\]
The graph budget three gives $n\ge h+2$ normalization
points. No singleton image support is possible: a
singleton $\{A\}$ forces $n\le h+1$, and a singleton
$\{B_+\}$ or $\{B_-\}$ forces $n\le2$. Hence every
point has at least two image labels and $n\le T/2$.
This excludes $h=2$. For $h=1$ there are precisely
three points, each with exactly two conductor branches
of different image labels. Up to signs and interchange
of $V,W$, their supports are
\[
 \begin{array}{c|c}
 \text{hypersurface branch}&\text{conductor branches}\\ \hline
 F_R&R,\ B_-\text{ on }V\\
 F_H&H,\ B_+\text{ on }W\\
 F_0&B_+\text{ on }V,\ B_-\text{ on }W.
 \end{array}
                                                            \tag{4.56b}
\]
The tree is the path $H$--$W$--$V$--$R$, and $H\to A$
is ramified with local degree two at its point.

\item The intersection of $F_R$ and $F_H$ is supported on $A$
and has generic length two. The equation of $F_H$ has
normal order two along $A$, so (4.56a) makes $F_R$
smooth. Choose its equation as $X$, a coordinate $Y$
on that surface cutting $A$, and a coordinate $z$ on
$A$. Preparation and depression in $Y$, preserving the
ideal $(X,Y)$, give the actual equation
\[
                       F_H=\text{unit}\,[Y^2+X^2k(X,z)].
                                                            \tag{4.56c}
\]
The second $X$ follows from normal multiplicity two
along every nearby point of $A$.

\item The monic finite partial normalization is
\[
      N'=\{u^2+k(X,z)=0\},\qquad Y=Xu.
                                                            \tag{4.56d}
\]
It is an isomorphism away from $A$. The branch $F_H$
is generically smooth along $B_+$: the two generic
$B$ sheets belong separately to $F_H$ and $F_0$.
Thus $N'$ has at most an isolated singularity here.
Its hypersurface Cohen--Macaulay property and regularity
in codimension one make it normal, so it is the actual
normalization.

\item The reduced divisor $X=0$ is supported on $H$ and has
equation $u^2+k(0,z)=0$. The actual $H$ is smoothly
embedded and has degree two with ramification over $A$.
An even zero of $k(0,z)$ gives two branches; an odd
zero of order at least three gives a singular curve.
Therefore its zero is simple, $k_z(0,0)\ne0$, and
$N'$ is smooth. The minimal normalization resolution
is the identity in this neighbourhood.

\item Let $m\ge1$ be the generic intersection multiplicity of
the two branches along $B$. It is the same on both
local signs, because they belong to the same irreducible
image component. Let $G$ define $F_0$, and set
\[
                      b=\operatorname{ord}_{p,A}(G|_A)>0.
\]
On the smooth branch $F_R$, its divisor is $mB_-$.
Restricting to $A$ gives
\[
                 b=m I_{F_R}(A,B_-)\ge m.                \tag{4.56e}
\]
On the smooth normalization (4.56d), the divisor of $G$
is $mW$. The actual pole support is normal crossing,
so $H$ and $W$ meet transversely. Its restriction to
$H$ consequently has order $m$. But it is the pullback
of the same ambient restriction $G|_A$ under the local
degree-two map. Hence
\[
                              2b=m,
\]
contradicting (4.56e). Every case with $r_B\le2$ is
excluded.

\item The exact incidence partitions and the normalization
substitution are checked in
\texttt{code/part4/d4\_other\_two\_branches.py}.
Its corrected run reports PASS; the initial structural
symbolic-equality assertion failure is preserved separately.

\end{enumerate}
\end{proof}


\subsection{Exclusion of the two-junction spectral configuration}

\begin{proposition}[P4.57]\hypertarget{tag-02XW}{}\label{night:P4.57}\label{stmt:P4.57}
\StatementRecord{night:P4.57}{proved}{02XW}{\SeriesRef{night:P4.55}, \SeriesRef{night:P4.56}, \SeriesRef{night:P4R.19}}

The remaining degree-two configuration of P4.55--P4.56
cannot have its image-conductor graph cycle realized at
two distinct intersections of its two components.
\end{proposition}
\begin{proof}

\begin{enumerate}[label={$\langle1\rangle$\arabic*.},leftmargin=*]
\item The component $B$ has at least three generic normalization
branches, hence transverse multiplicity at least three.
Along the smooth $A$, choose coordinates $A=\{x=y=0\}$,
with $z$ on $A$. The actual equation belongs to $(x,y)^3$;
write $q_3(x,y;z)$ for its normal cubic.

\item A local branch of $B$ has first nonzero normal displacement
$(a,b)t^q$, with $(a,b)\ne0$. All ambient Hessian entries
of the actual equation vanish identically along $B$, by
its generic multiplicity. The first normal-order term of
its normal Hessian is
\[
       t^q\operatorname{Hess}_{x,y}
                    (q_3(\,\cdot\,,\,\cdot\,;0))(a,b).
\]
Terms of higher normal order and terms from varying $z$
have larger order in $t$. This matrix is therefore zero.
For a homogeneous binary cubic $Q$, a nonzero vector
annihilating this Hessian makes the directional derivative
of $Q$ the zero quadratic. In coordinates with that vector
as first axis, $Q$ is a cube of the other variable, or zero.
P4R.19 excludes the nonzero cube. Thus the cubic vanishes
at every $A$--$B$ junction, including tangent junctions.

\paragraph{The global spectral equality.}
Let $p,q$ be the two distinct junctions. The common cubic
zero divisor has degree $b\ge2$. On $\mathbb P(N_{A/Y}^*)$
its residual spectral curve is $\Gamma_H\cup\Gamma_R$,
of class $3\xi-bF$. The bundle has degree two, $\Gamma_R$
is a section, and $H$ normalizes $\Gamma_H$. Adjunction
and the exact union sequence give
\[
       g(H)+\delta(\Gamma_H)+\Gamma_H\cdot\Gamma_R=5-2b.
\]
More precisely, the displayed formula is an equality of
integers with $\delta(\Gamma_H)$ the total normalization
length. Since $g(H)\ge1$, it forces
\[
 b=2,\qquad g(H)=1,\qquad \delta(\Gamma_H)=0,
                  \qquad\Gamma_H\cdot\Gamma_R=0.          \tag{4.57a}
\]
The two common cubic zeros are simple, the double spectral
cover is smooth, and it is disjoint from the third section.

\paragraph{The graph support at a junction.}
Let $c_B$ count the global conductor components over $B$.
The graph budget is $c_B+1$. At either junction the
contribution is at least $\min(r_A,r'_B)$, where
$r_A\ge2$ and $r'_B\ge c_B$ count normalized points
over its two image branches.

\item If $c_B=2$, the two contributions total at least four,
exceeding the budget three. If $c_B=1$, a junction with
at least two $B$ preimages costs at least two, and the
other at least one, exceeding the budget two. Thus both
have one $B$ preimage. The argument of P4.56 gives the
path $H$--$V$--$R$, with one edge at each junction.
At the $R$--$V$ junction the local branch carries all
nearby generic $B$ sheets of $V$, regardless of its
finite degree. Its intersection with the $H$ branch or
branches is a length-one or length-two thickening of
smooth $A$. P4.56 makes it smooth, contradicting its
intrinsic nonnormality from the at least three $B$ sheets.

\item Therefore $c_B\ge3$. If $A$ belongs to every branch
support at a junction, that point costs at least
$r'_B\ge c_B$, and the other costs at least two.
This exceeds $c_B+1$. The two-element support argument
consequently forces $B$ to belong to every local
hypersurface branch at each junction. Branches carrying
only $B$ have been retained in this conclusion.

\paragraph{Choosing distinct normalization points.}
The ruling descends to $r_N:N\to H$. Its restriction
to $H$ is the identity, and $R$ maps to a fixed point
$b_0\in H$. If a point of $H$ and a point of $R$ have
the same image in $N$, that $H$ point is $b_0$.
This includes identifications through a contracted
exceptional chain. All such coincidences therefore lie
over the single value $a_0\in A$ which is the image of
$b_0$. Choose a junction different from $a_0$, and call
it $p$. Its $H$ preimages and its $R$ preimage are
distinct normalization points. Different $H$ points
are already distinct because $H$ is embedded.

\item The actual local equation factors into the corresponding
irreducible hypersurface factors and any factors carrying
only $B$. For an $A$-containing factor, let $b_i\ge0$
be the common $z$-order of its first nonzero normal term.
For a factor not containing $A$, let $c_j\ge1$ be the
order of its nonzero restriction to $A$ at $p$.
Gauss's lemma over $\mathbb C\{z\}$ gives
\[
                         \sum_i b_i+\sum_j c_j=1.         \tag{4.57b}
\]
Indeed primitive homogeneous polynomials have nonzero
product modulo $z$, so their common coefficient orders
add exactly. The right side is the simple cubic zero
of (4.57a).

\item If $H$ is unramified here, there are three
$A$-containing factors of normal order one. At least
two have nonzero linear normal terms, by (4.57b).
Their directions differ by (4.57a). The implicit-function
theorem makes their intersection exactly $A$, whereas
both must contain $B$. This is a contradiction.
Thus $H$ is ramified at $p$.

\item There are now two $A$-containing factors, denoted by
$F_H$ and $F_R$, of normal orders two and one. Their
residual spectral directions differ. Equality (4.57b)
has precisely the following three possibilities.

\paragraph{A factor carrying only $B$ uses the zero.}
Here $b_H=b_R=0$ and one such factor restricts to $A$
with order one. The branch $F_R$ is smooth. In coordinates
$X=F_R$, $A=\{X=Y=0\}$, restriction of $F_H$ to $X=0$
is $Y^2$ times a unit, since its residual direction
differs from the $R$ direction. Thus these two branches
intersect only along $A$. They both contain $B$, which
is impossible.

\paragraph{The quadratic factor uses the zero.}
Here $b_H=1$, $b_R=0$, and there is no other factor.
The branch $F_R$ is again smooth, and in its coordinates
exact division gives
\[
                         F_H|_{X=0}=Y^2h(Y,z).
\]
The value $h(0,z)$ has a simple zero: the residual
quadratic is nonzero on the $R$ direction, so its
common scalar factor survives with order one on this
restriction. Thus $h_z(0,0)\ne0$. Its zero curve is the
extra image branch $B$, with intersection multiplicity
one between $F_H$ and $F_R$.

\item The smooth $F_R$ branch supplies one generic $B$
normalization branch. Each branch supplied by $F_H$
has positive intersection multiplicity with it. There
can therefore be at most one such branch. The full
$B$ germ would have at most two branches, contradicting
P4.56.

\paragraph{The linear factor uses the zero.}
Finally $b_R=1$, $b_H=0$, with no other factor.
The nonzero quadratic form of $F_H$ has a double root.
Preparation and depression preserving $A$ give
\[
                   F_H=\text{unit}\,[Y^2+X^2k(X,z)].
\]
Its spectral curve is $u^2+k(0,z)=0$. By (4.57a) this
is smooth and its double cover is ramified at $p$.
Hence $k_z(0,0)\ne0$. The exact finite map
\[
 \{u^2+k(X,z)=0\}\longrightarrow\{F_H=0\},
                         \qquad Y=Xu
\]
has smooth source, is birational, and is an isomorphism
off $A$. It is the normalization. Thus $F_H$ supplies
one generic $B$ branch.

\item For $F_R$, the simple common zero of its linear normal
term makes its quadratic Taylor part
\[
                       zL(X,Y)+Q_2(X,Y),\qquad L\ne0.
\]
There is no $z^2$ term because $F_R$ vanishes on $A$.
Its Hessian has rank at least two. Rank three gives an
isolated critical point. In rank two the holomorphic
Morse lemma with a parameter gives the exact equation
$U^2+V^2+h(W)$. A nonisolated critical curve forces
$h'\equiv0$, hence $h=0$ because the factor vanishes
at the point. The resulting factorization
$(U+iV)(U-iV)$ contradicts local irreducibility.
Therefore $F_R$ has at most isolated singularities and
also supplies one generic $B$ branch. Again the full
branch number is two, contradicting P4.56.

\item All allocations in (4.57b) have been excluded, proving
the assertion. The exact genus budget and nonzero
Hessian minor are recomputed in
\texttt{code/part4/d4\_two\_junction.py}; its saved
output reports PASS.

\end{enumerate}
\end{proof}


\subsection{The completed irreducible classification and sign theorem}

\begin{proposition}[P4.58]\hypertarget{tag-02XX}{}\label{night:P4.58}\label{stmt:P4.58}
\StatementRecord{night:P4.58}{proved}{02XX}{\SeriesRef{night:P4.10}, \SeriesRef{night:P4.11}, \SeriesRef{night:P4.26}, \SeriesRef{night:P4.32}, \SeriesRef{night:P4.35}, \SeriesRef{night:P4.36}, \SeriesRef{night:P4.37}, \SeriesRef{night:P4.54}, \SeriesRef{night:P4.55}, \SeriesRef{night:P4.56}, \SeriesRef{night:P4.57}, \SeriesRef{night:P4.6}, \SeriesRef{night:P4.7}, \SeriesRef{night:P4.8}, \SeriesRef{night:P4R.13}, \SeriesRef{night:P4R.18}, \SeriesRef{night:P4R.20}}

Let $f:Y\to B$ be the fibration of Part~4.2, with
$b_1(Y)=b_2(Y)=0$, connected fibres and general fibre a
complex two-torus. Suppose every reduced fibre support
is irreducible or has normal crossings. Then
\[
 B=\mathbb P^1,\qquad c_3(Y)\in\{0,2\},
                       \qquad b_3(Y)=2-c_3(Y).
\]
Every reduced support has nonnegative Euler number.
An irreducible singular support whose normalization has
a rational resolution has Euler number two. Every other
irreducible support has Euler number zero. Such a support
is either smooth with a finite unramified torus cover,
or singular with smooth minimal elliptic ruled
normalization, or singular with smooth minimal Hopf
normalization. In particular the class-VII normalization
alternative is realized only by a smooth minimal Hopf
surface.

For a normal-crossing support the Euler number equals
its number of triple points, and is even. Its zero case
is exactly P4.8: in the reducible case the normalization
is a disjoint union of smooth minimal Hopf or non-K\"ahler
elliptic surfaces, with smooth elliptic conductor curves
and finite unramified branch maps. The equality
$c_3(Y)=0$ holds exactly when every support has one of
the corresponding zero-Euler forms.
\end{proposition}
\begin{proof}

\begin{enumerate}[label={$\langle1\rangle$\arabic*.},leftmargin=*]
\item P4.26 proves the sign for every irreducible support,
including singular normalizations, and P4.6--P4.8 give
the full normal-crossing calculation. P4.35 proves the
exact value two for the rationally resolved singular
case. P4.36--P4.37 show that every class-VII case has
smooth minimal Hopf normalization and Euler number
zero. P4.32 gives the complete irreducible zero-Euler
classification and its converses.

\item We finish the positive-genus ruled alternative without
a transverse-type assumption on its vertical conductor
branches. P4R.13 and the preceding reductions show that
a positive-Euler ruled support would have a doubled
horizontal pole section and an image conductor with two
rational normalized components and one graph cycle.
P4.54 excludes contraction of that section, using the
actual weighted normalization and blowdown maps.
P4R.18 reduces its generic image to $D_4$ with degree
two or three. P4.55 makes that image smooth and excludes
degree three, leaving the degree-two section and one
rational vertical third sheet. P4.56 forces at least
three generic branches on the other image component.

\item The graph alternatives are exhaustive by P4.55. P4.57
excludes two distinct junctions. P4R.20 excludes a sole
three-branch image point, with the exact conductor-path
and endpoint normalization calculation. Thus no
positive-Euler ruled support remains. This proves the
asserted value zero for every nonrational irreducible
support, and P4.32 supplies its actual normalization.

\item Finally Euler additivity gives $c_3(Y)$ as the sum of
these nonnegative support Euler numbers. The base,
Poincar\'e-duality calculation $e(Y)=2-b_3(Y)$, and
parity are proved in P4.26. They yield the displayed
global values. A sum of nonnegative integers vanishes
exactly when every summand vanishes, giving the equality
classification through P4.32 and P4.8.

\item The hypothesis checks are the explicit constructions
P4.10 and P4.11. Removing $b_2=0$ permits the proved
Euler value $24$. Removing the torus-fibre hypothesis
permits the Hopf-fibre example, which fails the stated
normalization equality classification. The numerical
normal-crossing sign itself follows from the stronger
calculation in P4.7 and does not require that hypothesis.

\end{enumerate}
\end{proof}


\subsection{The specialization group and the unibranch case}

\begin{proposition}[Fundamental group of an irreducible fibre]\hypertarget{tag-02Y2}{}\label{night:P4R.1}\label{stmt:P4R.1}
\StatementRecord{night:P4R.1}{proved}{02Y2}{the definitions and constructions stated in the proof}

For a proper holomorphic degeneration of complex two-tori with
smooth total threefold and central fibre $mS$, where $S$ is
irreducible and reduced, the group $\pi_1(S)$ has an abelian
normal subgroup with cyclic quotient of order dividing $m$.
\end{proposition}
\begin{proof}
Write the degeneration as $f:X\to\Delta$.
The compatible triangulation of {\L}ojasiewicz
\cite[Theorem~2, p.~464]{Lojasiewicz1964} applies to
the smooth real analytic total space and its closed analytic
subset $S$. Its hypotheses hold because the total space is
countable and finite dimensional, and the collection consisting
of $S$ is locally finite and semianalytic. Let $\tau:|K|\to X$ be the resulting homeomorphism and
$L\subset K$ the subcomplex with $|L|=\tau^{-1}(S)$.
After barycentric subdivision, $L$ is a full subcomplex.
For barycentric coordinates $\lambda_v$, put
$q=\sum_{v\in L}\lambda_v$. On the open star $q>0$, the map
\[
 r_K\Bigl(\sum_v\lambda_v v\Bigr)
   =\sum_{v\in L}\frac{\lambda_v}{q}v
\]
erases the other coordinates and divides the remaining ones
by their sum. For $0\le s\le1$, the interpolation
$H_s(x)=(1-s)x+s r_K(x)$ stays in the same star, since
$q(H_s(x))=(1-s)q(x)+s>0$, and fixes $L$.
Transport by $\tau$ gives a strong deformation
retraction $r:V\to S$. Compactness of $S$ and successive
subdivision make $V$ lie inside a relatively compact
neighbourhood of $S$.

Properness gives a small disc $\Delta_\epsilon$ with
$U=f^{-1}(\Delta_\epsilon)\subset V$: on the compact
complement of $V$ in a larger closed tube, $|f|$ has a
positive minimum. Thus the inclusion $i:S\to U$ and the
continuous restriction $r:U\to S$ satisfy $r\circ i=\id$.
The neighbourhood-retraction statement immediately after
Definition~1 of Fern\'andez de Bobadilla--Koll\'ar
\cite[p.~85]{BobadillaKollar2012} applies to the same proper
map $f:X\to\Delta$.

Write $U^*=U\setminus S$. The proper smooth fibration
$U^*\to\Delta_\epsilon^*$ has fibre $T^4$. Its homotopy
sequence and $\pi_2(\Delta_\epsilon^*)=0$ give
\[
 \pi_1(U^*)=\Lambda\rtimes\langle t\rangle,
 \qquad \Lambda=\mathbb Z^4.
\]
Choose $t$ mapping to $1$ in the infinite cyclic quotient.
Every element is uniquely $\lambda t^n$, with $\lambda\in\Lambda$;
conjugation by $t$ gives the displayed semidirect product.
The inclusion $U^*\to U$ induces a surjection on fundamental
groups. A loop can be perturbed off the real-codimension-two
smooth stratum of $S$ and off its higher-codimension singular
strata. The same general-position argument for a null-homotopy
disc avoids the singular strata and meets $S_{\rm reg}$ in
finitely many transverse points. Deleting little discs about
these intersections expresses the kernel as the normal closure
of the corresponding meridians. The smooth locus of the
irreducible $S$ is connected. Along a path in that locus, a
trivialization of its complex normal line gives an annulus of
positive meridians in $U^*$, proving their conjugacy. The complex
normal coordinate fixes their positive orientation.

At a smooth point of $S$, write $f=vz^m$ with $v$ a
holomorphic unit. On a small coordinate neighbourhood,
$z'=z\exp((\log v)/m)$ is a holomorphic normal coordinate
and $f=(z')^m$. For small $a>0$, the positive circle
$z'=a\exp(2\pi i\theta)$, $0\le\theta\le1$, maps to
$a^m\exp(2\pi im\theta)$, so its image under
$\pi_1(U^*)\to\mathbb Z$ is $m$. Let $A$ be the image of $\Lambda$
in $\pi_1(U)$. It is abelian and normal, and
$\pi_1(U)/A$ is cyclic with order dividing $m$.
The homomorphism $r_*:\pi_1(U)\to\pi_1(S)$ is surjective
because $r_*i_*=\id$. Its image $r_*(A)$ is abelian and
normal, and the induced quotient remains cyclic of order
dividing $m$. This proves the assertion for $\pi_1(S)$.

\end{proof}



\begin{proposition}[Euler number for a bijective normalization]\hypertarget{tag-02YD}{}\label{night:P4R.2}\label{stmt:P4R.2}
\StatementRecord{night:P4R.2}{proved}{02YD}{\SeriesRef{night:P4.12}, \SeriesRef{night:P4.8}, \SeriesRef{night:P4R.1}}

In the standing torus-fibration setting of this section, every irreducible singular support
whose normalization is smooth and whose normalization morphism
is bijective has nonnegative topological Euler number.
Its Euler number is zero precisely when that normalization is
a minimal ruled surface over an elliptic curve or a Hopf surface.
\end{proposition}
\begin{proof}

\begin{enumerate}[label=\arabic*.,leftmargin=*]
\item Let $\nu:T\to S$ denote this finite normalization morphism.
A continuous bijection from the compact $T$ to the Hausdorff
$S$ is a homeomorphism. It therefore induces an isomorphism
of fundamental groups and equality of Euler numbers.
\SeriesRef{night:P4R.1} makes $\pi_1(T)$ virtually abelian.

\item The effective-conductor and canonical-torsion argument of
\SeriesRef{night:P4.12} applies to this same $T$ and proves
$\kappa(T)=-\infty$. Write $q:T\to T_0$ for the contraction
to a minimal model, with $k$ point blowups. The classification
gives rational, ruled, and class-VII possibilities for $T_0$.
In the K\"ahler case, $p_g(T_0)=0$ permits an integral
K\"ahler class; the embedding given by
\cite[Main Theorem, pp.~314--315]{Kodaira1954} identifies $T_0$
with a projective surface, to which
\cite[Theorem~VI.17 and Corollary~VI.18, pp.~111--112]{Beauville1978}
applies. In the non-K\"ahler case use
\cite[Theorem~5.1, p.~27]{OVV2020}. For its elliptic alternative, Theorem~4.2, Remark~4.4 and
Proposition~4.5 there (pp.~17--22) give elliptic reduced fibres,
the cyclic translation model at multiple fibres, and a
holomorphic elliptic action. For the fibration
$\pi:T_0\to C_0$, its nowhere-zero vertical vector field gives
$Q=T_{T_0}/\OO$, and the differential induces
$Q\simeq\pi^*T_{C_0}(-R_\pi)$, where $R_\pi$ is the
multiple-fibre ramification divisor. Thus
$K_{T_0}=\pi^*K_{C_0}(R_\pi)$ and $K_{T_0}^2=0$;
the elliptic fibres also give $e(T_0)=0$. Together with
$p_g=0$, Noether's formula gives $h^1(\OO)=1$.
Theorem~2.5 there (pp.~7--9) then gives $b_1=1$.
This alternative also satisfies the class-VII conditions
$b_1=1$ and $\kappa=-\infty$.
Blowing up a smooth point preserves the fundamental group:
van Kampen compares a four-ball with the disk bundle
of $\OO_{\mathbb P^1}(-1)$, both attached along the same
simply connected boundary $S^3$. For a ruled minimal model
$p:T_0\to C$ over a genus-$g$ curve, the homotopy sequence
of its $\mathbb P^1$-bundle gives the isomorphism $p_*$
on fundamental groups. The resulting isomorphism from
$\pi_1(S)$ to $\pi_1(C)$ is $p_*q_*(\nu_*)^{-1}$.
For $g\ge2$ this group is not virtually abelian. Indeed a
finite-index subgroup corresponds to an unramified cover
of genus $1+d(g-1)>1$, and a genus-$h$ surface group maps
onto the nonabelian free group on $h$ generators by sending
the $b$ generators to the identity in its usual presentation.
Consequently $g\le1$.

\item A point blowup replaces a point by $\mathbb P^1$ and
increases Euler characteristic by one. The cell counts of
$\mathbb P^2$ and a $\mathbb P^1$-bundle over a genus-$g$
curve give $3$ and $4-4g$, respectively. For class VII,
Poincar\'e duality and $b_1(T_0)=1$ give $e(T_0)=b_2(T_0)$.
Thus
\[
 e(T)=\begin{cases}
 3+k,&T_0=\mathbb P^2,\\
 4-4g+k,&T_0\text{ ruled over a genus-}g\text{ curve},\\
 b_2(T_0)+k,&T_0\text{ of class VII}.
 \end{cases}
\]
Since $g\le1$, these formulas give
$e(S)=e(T)\ge0$. Equality in the ruled case is exactly
$g=1$ and $k=0$. Equality in class VII gives a minimal surface with
$b_2=0$. The conductor is nonempty because $S$ is singular
and nonnormal; \SeriesRef{cdp:st:A.5.3} gives its effective
conductor divisor on $T$. By \cite[Theorem~1, p.~1]{BBK2019},
there is a finite \'etale cover $a:T'\to T$ whose domain is
Hopf or Inoue. The inverse image of the conductor contains
a curve, excluding the Inoue alternative by
\cite[Propositions~2(i) and~3(i), pp.~274--275 and~277,
and the paragraph after (21), p.~279]{Inoue1974}.
The composition of $a$ with the Hopf universal cover
$\C^2\setminus\{0\}\to T'$ is a universal cover of $T$;
hence $T$ itself is Hopf. The converse follows from the same Euler
formulas and the homeomorphism $\nu$.

\end{enumerate}
\end{proof}

\subsection{The conductor pushout for ruled normalizations of higher genus}

\begin{proposition}[Conductor topology over a higher-genus ruling]\hypertarget{tag-02YF}{}\label{night:P4R.3}\label{stmt:P4R.3}
\StatementRecord{night:P4R.3}{proved}{02YF}{\SeriesRef{cdp:st:A.5.3}, \SeriesRef{night:P4.12}, \SeriesRef{night:P4.13}, \SeriesRef{night:P4R.1}}

In the standing torus-fibration setting of this section, let an irreducible singular fibre
support $S$ have smooth normalization $\nu:T\to S$ whose
minimal model is ruled over a curve of genus $g\ge2$.
Write $D$ for the conductor divisor on $T$ and
$C=\operatorname{Sing}(S)_{\rm red}$.
The inclusion $D_{\rm red}\to T$ induces an isomorphism on
fundamental groups, and the normalization pushout induces
$\pi_1(S)\simeq\pi_1(C)$. Moreover,
\[
 e(C)\ge0,\qquad e(S)>0.
\]
The assertion includes conductor components on different
ruling fibres that map onto the same curve in $C$.
\end{proposition}
\begin{proof}

\begin{enumerate}[label={$\langle1\rangle$\arabic*.},leftmargin=*]
\item Proposition~\ref{cdp:st:A.5.3}, with the identity as the
minimal resolution of the smooth surface $T$, gives
$\omega_T(D)\simeq\nu^*\omega_S$ with $D$ effective and with its
conductor multiplicities. Proposition~\ref{night:P4.12} makes this
line torsion, so $K_T+D\equiv0$.
Let $\pi:T\to T_0$ be the sequence of point blowdowns to the
minimal ruled surface $r_0:T_0\to B$, and put $D_0=\pi_*D$.
By \cite[III.10, p.~40]{Beauville1978}, write $T_0=\mathbf P_B(V)$
for a rank-two bundle $V$ on $B$. In the basis of the tautological
class $h=c_1(\OO_{\mathbf P_B(V)}(1))$ and the fibre class $F$,
$h^2=\deg V$, $h\cdot F=1$, $F^2=0$ and
$K_{T_0}=-2h+(\deg V+2g-2)F$
\cite[III.18, pp.~46--47]{Beauville1978}.
Thus $K_{T_0}+D_0\equiv0$ and $D_0\cdot F=2$.

\item There are three possibilities for the horizontal part of
$D_0$: a reduced bisection, two distinct reduced sections,
or twice one section. The first two are excluded directly.
For an irreducible horizontal component $H$ of coefficient
one, adjunction and the effectiveness of the other
components give
\[
 2p_a(H)-2=(K_{T_0}+H)\cdot H
       =-(D_0-H)\cdot H\le0.
\]
On the normalization of $H$, Riemann--Hurwitz for its
finite map of degree $d\ge1$ to $B$ gives
$2g(\widetilde H)-2\ge d(2g-2)>0$.
Since $p_a(H)\ge g(\widetilde H)$, these inequalities
are incompatible. Thus
\[
 D_0=2H+\sum_{j=1}^{v}a_jF_j,\qquad a_j\ge1,
\]
where $H$ is a section and the $F_j$ are distinct ruling
fibres. A horizontal curve of degree one is a section:
its normalization maps isomorphically to the smooth base,
and the resulting map to the projective bundle is a
section with closed image. In particular $H$ is smooth.

\item At one step $\beta:T'\to T''$ of the blowup sequence,
let $\mu$ be the multiplicity of the conductor $D''$ at the
centre and $E$ the exceptional curve.
Write $D'=D''_{\rm strict}+aE$. The canonical blowup formula
\cite[II.3(iv), pp.~17--18]{Beauville1978} gives
\[
 K_{T'}+D'
   =\beta^*(K_{T''}+D'')+(1+a-\mu)E.
\]
Intersection with $E$ and numerical triviality give
$a=\mu-1$. Effectiveness therefore places every centre
on the conductor. At a smooth conductor point of
coefficient one, the exceptional curve is absent from
the new conductor. At a smooth point of coefficient at
least two, it is a rational leaf attached at one point.
At an intersection of two conductor components it
subdivides their edge. Starting with $H$ and the
transverse fibres $F_j$, these operations preserve
normal crossings and make $D_{\rm red}$ a copy of $H$
with rational trees attached. No cycle is introduced.

\item The projection of $H$ to $B$ induces an isomorphism on
fundamental groups, and every attached rational tree
is simply connected. A point blowup replaces a four-ball by the disc bundle
of $\OO_{\mathbb P^1}(-1)$, with common boundary $S^3$.
Van Kampen therefore shows that it leaves $\pi_1(T)$ unchanged. The preceding
description also proves, at every blowup, that the
inclusion of the reduced conductor induces the same
isomorphism to the fundamental group of the surface.
Hence $\pi_1(D_{\rm red})\to\pi_1(T)$ is an
isomorphism.

\item The normalization morphism $\nu$ is finite, is an
isomorphism away from the conductor, and has smooth
source. Its nonisomorphism locus equals $C$: a normal
point has a neighbourhood identified with its
normalization and is consequently smooth. Conversely
every smooth point is normal. Thus the natural map
\[
 T\underset{D_{\rm red}}{\cup}C\longrightarrow S
\]
is a continuous bijection. Both source and target are
compact Hausdorff spaces, so it is a homeomorphism.
A triangulation compatible with the analytic pair $(T,D_{\rm red})$
and the finite curve map $D_{\rm red}\to C$ makes
$D_{\rm red}\subset T$ a subcomplex and hence a
cofibration. Replacing the map $D_{\rm red}\to C$ by
its mapping cylinder computes the same homotopy
pushout: the homotopy extension property for this
cofibration contracts the inserted cylinder without
altering the homotopy type. Van Kampen on that double
mapping cylinder therefore gives
\[
 \pi_1(S)=\pi_1(T)
       *_{\pi_1(D_{\rm red})}\pi_1(C)
       \simeq\pi_1(C).
\]
All three spaces in the amalgam are connected. This
argument includes every identification made by the
finite conductor map, including identifications
between horizontal and vertical branches.

\item Proposition~\ref{night:P4R.1} makes this group virtually abelian. Let $c$ be
the number of irreducible components of $C$, let
$g_i$ be the genera of their smooth normalizations,
and let $b$ be the first Betti number of the incidence
graph of the finite point identifications. Replacing
each identification by a small tree joining its
preimages gives a homotopy model consisting of these
closed oriented surfaces joined along a connected
graph. Choosing a spanning tree gives
\[
 \pi_1(C)\simeq
   \Bigl(*_{i=1}^{c}\pi_1(C_i^{\rm norm})\Bigr)*F_b,
 \qquad
 e(C)=c+1-2\sum_i g_i-b.
\]
The formula also covers unibranch singular points:
their normalization is a local topological
homeomorphism and they contribute no graph edge.

\item A genus-$h$ surface factor with $h\ge2$ cannot be virtually
abelian: a connected degree-$d$ cover has genus $1+d(h-1)>1$,
and its fundamental group is nonabelian. Projection onto this
surface factor is a quotient map, while every quotient of a
virtually abelian group is virtually abelian. Two nontrivial surface factors,
two free generators, or one elliptic factor together
with a free generator contain a nonabelian free
subgroup. Such a subgroup cannot lie in a virtually
abelian group: its intersection with a finite-index
abelian subgroup would have finite index and would
be abelian, whereas a finite-index subgroup of a
nonabelian free group has rank at least two. The
possibilities are therefore all $g_i=0$ with
$b\le1$, or exactly one $g_i=1$ with $b=0$.
The displayed Euler formula gives $e(C)\ge0$.
Additivity of compactly supported Euler numbers and
Proposition~\ref{night:P4.13} give
\[
 e(S)=e(S_{\rm reg})+e(C)\ge0.
\]

\item Finally equality is impossible here. It would force
$e(C)=0$, so the same formula gives $c=1$, $g_1=1$
and $b=0$. Every irreducible component of
$D_{\rm red}$ maps finitely and nonconstantly onto
this curve. A rational vertical component would
therefore give a nonconstant map from $\mathbb P^1$
to its genus-one normalization, contrary to
Riemann--Hurwitz. There are no vertical conductor
components. The initial conductor is consequently
$2H$, and there can have been no blowup: the first
centre would lie on $2H$ and would create a vertical
conductor component of coefficient one. Thus
$T=T_0$ and $D_{\rm red}=H$. It follows that
\[
 e(S_{\rm reg})=e(T\setminus H)
       =(4-4g)-(2-2g)=2-2g<0,
\]
contradicting Proposition~\ref{night:P4.13}. This proves strict positivity.

\end{enumerate}
\end{proof}

P4R.3 proves strict positivity for its stated subclass.
The later classification P4.58 excludes every positive-Euler
ruled case, including this subclass, by the complete local
and global conductor arguments P4.54--P4.57 and P4R.18--P4R.20.
The earlier live status is retained in
\path{records/part4_root_superseded_status.md}.

\begin{proposition}[Meridian quotient and rational incidence trees]\hypertarget{tag-02YG}{}\label{night:P4R.4}\label{stmt:P4R.4}
\StatementRecord{night:P4R.4}{proved}{02YG}{\SeriesRef{night:P4.13}, \SeriesRef{night:P4R.1}, \SeriesRef{night:P4R.3}}

In the setting of \SeriesRef{night:P4R.1}, choose a positive small-loop
monodromy $A$ on $\Lambda=H_1(T^4;\mathbb Z)$.
Let $t$ lift the positive base generator, with the convention
$tvt^{-1}=A(v)$ for $v\in\Lambda$.
For a suitable $\lambda\in\Lambda$, a positive
meridian of the reduced support is represented by
$\mu=\lambda t^m$ in $\Lambda\rtimes_A\langle t\rangle$.
The natural small-tube comparison gives
\[
 \pi_1(S)\simeq
 (\Lambda\rtimes_A\langle t\rangle)/\langle\!\langle\mu\rangle\!\rangle,
 \qquad
 \pi_1(S)/\operatorname{im}\Lambda\simeq\mathbb Z/m,
 \qquad
 b_1(S)=\dim_{\mathbb Q}\ker(A-I).
\]
For the support in \SeriesRef{night:P4R.3}, a singular-locus incidence
graph that is a tree of rational normalization
components gives $e(S)\ge3$.
\end{proposition}
\begin{proof}

\begin{enumerate}[label=\arabic*.,leftmargin=*]
\item Use the nested neighbourhoods in the proof of
\SeriesRef{night:P4R.1}. Choose an open-star
neighbourhood $V_0$ retracting onto $S$, a small tube
$U_0\subset V_0$, a second open-star neighbourhood
$V_1$ with $S\subset V_1\subset U_0$, and a smaller
tube $U_1\subset V_1$. Properness supplies both tube
inclusions. The punctured-disc fibrations give an
isomorphism
\[
 \pi_1(U_1\setminus S)\longrightarrow
 \pi_1(U_0\setminus S):
\]
the fibres can be identified by radial parallel
transport and both base fundamental groups are
generated by the same positive loop. As in the proof of
\SeriesRef{night:P4R.1},
the kernels of the maps to $\pi_1(U_i)$ are the
normal closures of the positive smooth-support
meridian. These meridians correspond under the
displayed isomorphism. Hence
$\pi_1(U_1)\to\pi_1(U_0)$ is an isomorphism.

\item This isomorphism factors through $\pi_1(V_1)$,
which is identified with $\pi_1(S)$ by its
retraction. Thus $\pi_1(S)\to\pi_1(U_0)$ is
surjective. The restriction to $U_0$ of the
retraction of $V_0$ proves injectivity. The inclusion
$S\hookrightarrow U_0$ and the restricted retraction are
inverse on fundamental groups. Together with the meridian
kernel calculation, this gives the asserted quotient
presentation of $\pi_1(S)$.

\item Projection to the cyclic factor sends $\mu$ to
$m$. After dividing by the image of $\Lambda$,
the remaining presentation is therefore exactly
$\langle t\mid t^m=1\rangle$. In the rational
abelianization the presentation becomes
\[
 \frac{(\Lambda/(A-I)\Lambda)\otimes\mathbb Q
       \ \oplus\ \mathbb Qt}
      {\mathbb Q(\overline\lambda+mt)}.
\]
The single displayed relation has nonzero
$t$ coefficient and reduces the dimension by
one. Explicitly, with
$V=(\Lambda/(A-I)\Lambda)\otimes\mathbb Q$, the map
$(v,c)\mapsto v-c\overline\lambda/m$ induces an isomorphism
of this quotient with $V$, whose inverse sends $v$ to
$[(v,0)]$. This proves the asserted formula for $b_1$.

\item In the rational-tree case, \SeriesRef{night:P4R.3}
gives $\pi_1(S)\simeq\pi_1(C)$ for $C=\operatorname{Sing}(S)_{\rm red}$.
Write $\nu_C:\coprod_{j=1}^c\mathbb P^1_j\to C$ for its
normalization. Replace each identification of a finite fibre
of $\nu_C$ by a star joining its preimages. Collapsing these
contractible stars gives $C$ up to homotopy; collapsing the
sphere components induces an isomorphism on fundamental
groups with the incidence graph. Van Kampen therefore gives
$\pi_1(C)=1$ and $\pi_1(S)=1$.
The cyclic quotient then forces $m=1$, and the
$b_1$ formula gives $\det(I-A)\ne0$.
\SeriesRef{night:P4.13} gives
$e(S_{\rm reg})=\det(I-A)\ge0$; this nonzero integer
is at least one.
Let $P=\{p\in C:\#\nu_C^{-1}(p)>1\}$ and put
$r_p=\#\nu_C^{-1}(p)$ for $p\in P$. The tree has
$c+\#P$ vertices and $\sum_{p\in P}r_p$ edges, so
$\sum_p(r_p-1)=c-1$. Thus
\[
 e(C)=2c-\sum_p(r_p-1)=c+1\ge2.
\]
Unibranch singular points have one preimage and contribute
zero to the subtraction. Finally
$e(S)=e(S_{\rm reg})+e(C)\ge3$.

\end{enumerate}
\end{proof}

\subsection{Hopf normalizations and the conductor graph}

\begin{proposition}[P4R.5]\hypertarget{tag-02YH}{}\label{night:P4R.5}\label{stmt:P4R.5}
\StatementRecord{night:P4R.5}{proved}{02YH}{\SeriesRef{night:P4.12}, \SeriesRef{night:P4R.1}, \SeriesRef{night:P4R.3}}

In the standing torus-fibration setting of this section, an irreducible singular
support whose normalization is smooth and has a Hopf
surface as its minimal model has nonnegative Euler
number. More precisely, for its reduced conductor
image $C$ one has $e(C)\ge0$. If the normalization is
obtained by $k$ point blowups and $v$ of the exceptional
curves belong to the conductor, then
\[
 e(S_{\rm reg})=k-v,\qquad
 e(S)=k-v+e(C)\ge0.
\]
\end{proposition}
\begin{proof}

\begin{enumerate}[label={$\langle1\rangle$\arabic*.},leftmargin=*]
\item Let $T_0$ be the minimal Hopf surface. Every
irreducible curve $E$ on $T_0$ is smooth elliptic.
Indeed $b_2(T_0)=0$ gives $E^2=K_{T_0}\cdot E=0$,
so adjunction gives $p_a(E)=1$. A rational
normalization would define a nonconstant map
$\mathbb P^1\to T_0$. It lifts to the universal
cover $\mathbb C^2\setminus\{0\}$, where its two
coordinate functions are constant. This is
impossible. Thus the normalization genus is one
and the sum of the local delta invariants is zero.
Distinct curves are disjoint, since their
intersection number is zero and every local
intersection of distinct complex curves is
positive.

\item Push the effective conductor $D$ to $D_0$ on
$T_0$. The canonical comparison of P4.12 makes
$K_{T_0}+D_0$ torsion. The divisor $D_0$ is
nonempty: otherwise a positive power of $K_{T_0}$
would be trivial, contradicting the Kodaira
dimension $-\infty$ of a Hopf surface. Its
reduced support consists of $r\ge1$ disjoint
smooth elliptic curves. The blowup coefficient
$\mu-1$ calculation of P4R.3 applies without a
ruling. At a smooth conductor point it adds a
rational leaf precisely when $\mu\ge2$; at a
node it subdivides the corresponding edge.
Consequently $D_{\rm red}$ has $r$ connected
components, each a smooth elliptic root with a
rational tree attached. With $v$ rational
components it has exactly $v$ double points,
and hence $e(D_{\rm red})=v$. Since
$e(T)=k$, normalization away from the conductor
gives $e(S_{\rm reg})=k-v$. Each rational
component was created by one of the $k$ blowups,
so $k-v\ge0$.

\item Write $s$ for the number of connected components
of $C$, $c$ for its number of irreducible
components, and $g$ for the number of its
normalizations of genus one. Every normalized
component has genus zero or one: it receives a
finite nonconstant map from a normalized
conductor component. Each elliptic component
must receive an elliptic root, and distinct
image components require distinct roots.
Thus $c\ge g$ and $r\ge g$. Let $b$ be the sum
of the cycle ranks of the point-identification
graphs of the $s$ connected components. The
normalization formula gives
\[
 e(C)=c+s-2g-b.
\]

\item We retain the disconnected conductor in the
fundamental-group comparison. Its pushout has
one vertex for $T$ and one for each connected
component of $C$, with $r$ edges supplied by
the connected components of $D_{\rm red}$.
This graph is connected, because the finite
map $D_{\rm red}\to C$ is surjective.
The van Kampen presentation has $r-s$ stable
letters, in addition to the vertex groups.
The group of each conductor component is the
group of its elliptic root. Its map into $C$
factors, on that root, through the smooth
normalization of one image component; the
attached rational tree adds no fundamental
group generator. Kill the vertex group
$\pi_1(T)$ and the normal subgroups generated
by all normalized component groups of $C$.
Every conductor-group relation then disappears.
The resulting quotient is exactly
\[
 \pi_1(S)\twoheadrightarrow F_{b+r-s}.
\]
This is also the group of the homotopy-pushout
graph obtained by making the same quotients in
the normalization graph model of $C$. P4R.1
therefore gives $b+r-s\le1$. Substitution yields
\[
 e(C)\ge c+r-2g-1\ge-1.
\]

\item We exclude the only negative possibility by
the actual local hypersurface geometry.
The value $e(C)<0$ in these inequalities
forces $c=r=g$ and $b+r-s=1$. Every image
component then has elliptic normalization.
No rational conductor component can map
finitely onto one of them, by
Riemann--Hurwitz. Thus $D_{\rm red}$ consists
of disjoint smooth elliptic curves. This
conclusion permits blowups at coefficient-one
points; their exceptional curves are absent
from the conductor and do not affect the
following argument.

\item At a point $p\in C$, suppose there were two
distinct reduced curve germs $C_1,C_2$ of
$C$. Choose normalization preimages over $p$
whose conductor germs map onto $C_1,C_2$.
They must be distinct preimages, since the
reduced conductor has just one smooth germ
at each of its points. They correspond to
distinct irreducible analytic surface
branches $S_1,S_2$ of $S$ at $p$.
The support $S$ is a hypersurface in the
smooth total threefold. Its distinct
surface branches have a one-dimensional
intersection germ through $p$: two local
nonassociate irreducible equations in the
regular three-dimensional local ring have
height two. That intersection lies in the
conductor image of both branches. The
unique conductor germ on each chosen
normalization branch therefore has the
same image germ, giving $C_1=C_2$, a
contradiction.

\item Hence $C$ is locally unibranch and its
irreducible components are disjoint. Its
point-identification graphs have $b=0$
and $s=c=g$. The Euler formula gives
$e(C)=0$, contrary to the negative
possibility. We have proved $e(C)\ge0$,
and the asserted formulas finish the proof.

\end{enumerate}
\end{proof}

\subsection{The integral pole divisor on a normalization resolution}

\begin{proposition}[Conductor and exceptional pole divisor]\hypertarget{tag-02YI}{}\label{night:P4R.6}\label{stmt:P4R.6}
\StatementRecord{night:P4R.6}{proved}{02YI}{\SeriesRef{night:P4.12}}

Let $f:Y\to B$ be a proper surjective holomorphic map with connected
fibres from a smooth compact complex threefold with $b_1(Y)=b_2(Y)=0$
to a smooth compact curve, and suppose that its smooth general fibre
is a complex two-torus. Let $mS$ be a fibre with irreducible, reduced,
singular support $S$, let $\nu:N\to S$ be its normalization, and let
$\pi:T\to N$ be its minimal resolution. Put $\rho=\nu\circ\pi$.
There is an effective integral divisor $D$ on $T$ such that
\[
 \omega_T(D)\simeq\rho^*\omega_S,
 \qquad \operatorname{Supp}D
       =\rho^{-1}(\operatorname{Sing}S)_{\rm red}.
\]
In particular $K_T+D$ is torsion, by \cref{night:P4.12}.
The strict-transform conductor has its usual positive integral
coefficients. Every connected exceptional configuration over a
singular point of $N$ has strictly positive coefficients in $D$.
Every such configuration meets the strict-transform conductor.
\end{proposition}
\begin{proof}
\begin{enumerate}[label=\arabic*.,leftmargin=*]
\item The reduced support is a Cartier hypersurface in the smooth
threefold, so $\omega_S$ is invertible. A local generator $\eta$
pulls back from $S_{\rm reg}$ to a meromorphic two-form on $T$.
Its divisorial orders are integers. Changing $\eta$ by a holomorphic
unit changes its pullback by a unit, so the negatives of these
orders define an integral divisor $D$ and the displayed
line-bundle isomorphism.

Write $\mathfrak c$ for the conductor, viewed as an ideal on $N$.
For a finite birational normalization $A\subset B$, let $Q(A)$
be the total quotient ring obtained by inverting the non-zero-divisors
of $A$. Then $B\otimes_A Q(A)=Q(A)$. Every $A$-linear map
$\varphi:B\to A$ extends to a $Q(A)$-linear endomorphism of this
common total quotient ring, hence to multiplication by $\varphi(1)$.
Thus evaluation at $1$ identifies $\operatorname{Hom}_A(B,A)$ with
$\{a\in A:aB\subset A\}$. Finite duality therefore gives
\[
 \nu_*\omega_N
   =\mathcal Hom_{\mathcal O_S}(\nu_*\mathcal O_N,\omega_S),
 \qquad
 \omega_N\simeq\mathfrak c\otimes\nu^*\omega_S.
\]
At a regular generic point of a conductor curve, write
$\mathfrak c=(z^c)$, with $c>0$. The pullback of $\eta$ is
$z^{-c}$ times a regular canonical generator. Thus the
strict-transform coefficients of $D$ are exactly the positive
integral conductor exponents. Outside the conductor and the
exceptional curves, $\rho$ identifies the forms on the smooth
loci, and the order is zero.

\item Let $E_1,\ldots,E_r$ be the exceptional curves over a singular
point $p\in N$, and put $M_{ij}=E_i\cdot E_j$.
The fibre $\pi^{-1}(p)$ is connected: a decomposition into two
disjoint compact parts would, by properness, give disjoint
neighborhoods containing $\pi^{-1}(U)$ for a sufficiently small
neighborhood $U$ of $p$. The function with values $0$ and $1$
on those neighborhoods descends to a bounded holomorphic function
on $U\setminus\{p\}$. Normality extends it across $p$, whereas
approach to the two exceptional parts gives different limiting
values. This is a contradiction.

Choose a nonzero local holomorphic function on $N$ vanishing at
$p$. Its pullback divisor is $H+\sum_i n_iE_i$, with $n_i>0$
and $H$ effective without exceptional components. The strict
transform $H$ meets the exceptional fibre by properness. Since
the divisor is principal near that fibre,
\[
 Mn=-h,\qquad h_i=H\cdot E_i\ge0,\qquad h\ne0.
\]
The off-diagonal entries of $M$ are nonnegative. With
$v_i=n_iw_i$, direct expansion gives
\[
 v^tMv
 =-\sum_{i<j}M_{ij}n_in_j(w_i-w_j)^2
   -\sum_i n_ih_iw_i^2.
\]
This is the weighted intersection identity of
\cite[p.~6]{Mumford1961}: the local principal divisor above
replaces the source's hyperplane difference, and
$v=\operatorname{diag}(n_i)w$ is the change of variables.
Connectedness and $h\ne0$ make the right-hand side strictly
negative for $v\ne0$.

For $Md=-b$, $b_i\ge0$, set $u_i=d_i/n_i$. The same divisor
equation gives the row identity
\[
 b_i=h_i u_i+\sum_{j\ne i}M_{ij}n_j(u_i-u_j).
\]
A negative minimum of the $u_i$ forces every neighboring value
to equal it and forces $h_i=0$. Propagation through the connected
graph reaches a vertex with $h_i>0$, a contradiction. Thus
$d_i\ge0$. When $b\ne0$, a zero $d_i$ forces all neighboring
$d_j$ to be zero, and propagation contradicts $Md=-b$.
Consequently $d_i>0$ in that case; equivalently, $-M^{-1}$ has
strictly positive entries on the connected block.

\item Write $D=D_{\rm str}+\sum_i d_iE_i$, where
$D_{\rm str}$ is the strict-transform conductor. Adjunction and
minimality of $\pi$ give
\[
 K_T\cdot E_i=2p_a(E_i)-2-E_i^2\ge0.
\]
Indeed $E_i^2<0$ by the preceding identity, and an integral curve
with $p_a(E_i)=0$ is smooth and rational. Minimality excludes an
exceptional curve of this kind with self-intersection $-1$.
Also $D_{\rm str}\cdot E_i\ge0$. Since $\rho(E_i)$ is a point,
the restriction of $\rho^*\omega_S$ to $E_i$ has degree zero.
The canonical divisor identity therefore gives
\[
 Md=-b,\qquad b_i=(K_T+D_{\rm str})\cdot E_i\ge0.
\]
It follows that $d_i\ge0$. A zero coefficient forces $b=0$
on the connected block. Each of its curves is then a smooth
rational $(-2)$-curve, and the block is disjoint from
$D_{\rm str}$.

\item We now prove that every exceptional block meets
$D_{\rm str}$. The identity in the first item makes
$\mathfrak c$ a rank-one reflexive ideal on the normal surface
$N$. Such an ideal is the intersection of its codimension-one
localizations in the fraction field. Hence
$\operatorname{Supp}(\mathcal O_N/\mathfrak c)$ has no isolated
point: a conductor ideal equal to $\mathcal O_N$ at all nearby
codimension-one points is equal to $\mathcal O_N$ itself.
A conductor curve through $p$ has strict transform meeting
$\pi^{-1}(p)$, by properness. Therefore a block disjoint from
$D_{\rm str}$ lies over a point where $\nu$ is an isomorphism.
Its image in $S$ would be a normal isolated hypersurface
singularity.

We exclude such a point directly. In a coordinate ball write
$f=u g^m$, where $g=0$ is the reduced equation and $u$ is a
holomorphic unit. Choose a logarithm of $u$ and put
$h=\exp((\log u)/m)g$. Then $f=h^m$ exactly. The critical point
of $h$ is isolated: a positive-dimensional critical component
through the point would have constant value zero and would lie
in the singular locus of the normal surface $\{h=0\}$.
The gradient of $h$ has positive local degree. A sufficiently
small regular value $a$ of that gradient makes
$h_a=h-\sum_j a_jz_j$ have finitely many, and at least one,
nondegenerate critical points. The degree and positivity
argument is given in \cite[Appendix~B, pp.~111--114]{Milnor1968}.
For this analytic germ the gradient is finite over a small
target neighborhood; at each regular preimage its real
Jacobian is the positive number $|\det_{\mathbb C}d(\nabla h)|^2$.

The fibre identifications can be made inside a fixed ball.
Indeed the critical locus of squared radius on
$\{h=0\}\setminus\{0\}$ is semianalytic. In a compatible locally
finite analytic triangulation from
\cite[Theorem~2, p.~464]{Lojasiewicz1964}, an accumulation at $0$
would give a connected positive-dimensional simplex with $0$
in its closure. Squared radius would be constant on it and
equal to zero, a contradiction. A small ball therefore has
boundary transverse to $\{h=0\}$. Transversality on this compact
boundary persists for small fibre values and small perturbations.
Choose a nonzero value $\delta$ and then $a$ small enough that
all critical values of $h-sa\cdot z$, $0\le s\le1$, have
absolute value less than $|\delta|/2$. Boundary-tangent
horizontal vector fields identify the fibres over $\delta$
through this proper family.

At a critical point of $h_a$, holomorphic Morse coordinates give
$h_a=c+z_1^2+z_2^2+z_3^2$. In the fibre over $c+\epsilon$,
the real sphere $\sum x_j^2=\epsilon>0$ has normal bundle
$J(TS^2)$. The tangent-then-normal order
$(e_1,e_2,Je_1,Je_2)$ has the opposite orientation to the
complex order $(e_1,Je_1,e_2,Je_2)$, so its self-intersection
is $-\chi(S^2)=-2$. Transport along a path avoiding the finite
critical values, followed by the preceding parameter-family
identification, embeds this sphere in $h^{-1}(\delta)$.
Since $f=h^m$, this local fibre embeds in the smooth global
fibre $f^{-1}(\delta^m)$, a complex two-torus. The sphere lifts
to its contractible universal cover $\mathbb C^2$, so its
homology class and self-intersection in the torus are zero.
Its normal bundle is unchanged by the inclusion of the local
fibre, contradicting the value $-2$.

Thus every exceptional block meets $D_{\rm str}$, so its vector
$b$ in the third item is nonzero. All its coefficients $d_i$
are strictly positive. Every normal point of $S$ is smooth by
the same local argument. The singular locus is therefore the
conductor locus, and its inverse image under $\rho$ consists
exactly of the strict conductor and these exceptional fibres.
This proves effectiveness and the support equality.
Finally, \cref{night:P4.12} gives $\omega_S^{\otimes m}\simeq
\mathcal O_S$, whence
\[
 (\omega_T(D))^{\otimes m}
   \simeq\rho^*(\omega_S^{\otimes m})\simeq\mathcal O_T.
\]
\end{enumerate}
\end{proof}

\begin{proposition}[P4R.7]\hypertarget{tag-02YJ}{}\label{night:P4R.7}\label{stmt:P4R.7}
\StatementRecord{night:P4R.7}{proved}{02YJ}{\SeriesRef{night:P4.13}, \SeriesRef{night:P4R.1}, \SeriesRef{night:P4R.3}, \SeriesRef{night:P4R.6}}

In Part~4.2, every irreducible singular
support whose normalization has a
resolution birationally ruled over a
curve of genus at least two has
nonnegative Euler number. The
normalization may be singular.
\end{proposition}
\begin{proof}

\begin{enumerate}[label={$\langle1\rangle$\arabic*.},leftmargin=*]
\item Use the minimal resolution and effective
integral divisor of P4R.6. Push this
divisor to the minimal ruled model
$T_0$. Its intersection with a ruling
fibre is two. The adjunction and
Riemann--Hurwitz calculation in P4R.3
therefore gives exactly
\[
 D_0=2H+\sum_j a_jF_j,
 \qquad a_j\in\mathbb Z_{>0},
\]
with $H$ a smooth section. The reverse
blowup coefficient is the integral
number $\mu-1$. At a conductor crossing
one has $\mu\ge2$, so its exceptional
curve remains in the support and
subdivides the edge. At a smooth
conductor point the operation either
adds a rational leaf or leaves the
reduced curve unchanged. Thus
$A=\operatorname{Supp}D$ is connected,
is $H$ with rational trees attached,
and $\pi_1(A)\to\pi_1(T)$ is an
isomorphism.

\item Put $C=\operatorname{Sing}(S)_{\rm red}$.
The map $\rho$ is proper and is an
isomorphism outside $A$, and P4R.6
identifies $A$ with its full inverse
image of $C$. The continuous map
$T\cup_A C\to S$ is therefore a
compact-to-Hausdorff bijection and a
homeomorphism. The same cofibration
and van Kampen argument as P4R.3 gives
$\pi_1(S)\simeq\pi_1(C)$. This retains
the exceptional components that map
to points, as well as the strict
conductor components that map to
curves. P4R.1 and the normalization
graph calculation in P4R.3 give
$e(C)\ge0$. Finally P4.13 gives
$e(S_{\rm reg})\ge0$, so
$e(S)=e(S_{\rm reg})+e(C)\ge0$.

\end{enumerate}
\end{proof}

\begin{proposition}[P4R.8]\hypertarget{tag-02YK}{}\label{night:P4R.8}\label{stmt:P4R.8}
\StatementRecord{night:P4R.8}{proved}{02YK}{\SeriesRef{night:P4R.5}, \SeriesRef{night:P4R.6}, \SeriesRef{night:P4R.7}}

In Part~4.2, every irreducible singular support
whose normalization has a resolution with Hopf
minimal model has nonnegative Euler number.
The normalization may be singular. For the
minimal normalization resolution, let $k$ be
the number of blowups from its Hopf minimal
model and $v$ the number of rational components
of the pole divisor. Then
\[
 e(S_{\rm reg})=k-v\ge0,\qquad
 e(\operatorname{Sing}S)\ge0.
\]
\end{proposition}
\begin{proof}

\begin{enumerate}[label={$\langle1\rangle$\arabic*.},leftmargin=*]
\item Use $\rho:T\to N\to S$ and the effective
integral divisor $D$ of P4R.6. Put
$A=\operatorname{Supp}D$ and
$C=\operatorname{Sing}(S)_{\rm red}$.
The Hopf-curve and blowup calculations of
P4R.5 use only integrality, effectiveness,
and $K_T+D\equiv0$. They therefore apply
to this divisor: $A$ has $r\ge1$ connected
components, each a smooth elliptic root
with rational trees attached. Each of its
$v$ rational components came from one of
the $k$ blowups; $e(A)=v$, $e(T)=k$, and
$\rho$ is an isomorphism off $A$. This
proves the displayed formula for the
regular locus.

\item Every one-dimensional component of $C$
receives a nonconstant proper map from a
component of $A$, so its normalization
has genus zero or one. There are no
isolated points of $C$: an isolated
singular point of the hypersurface is
normal by the hypersurface $S_2$ property
and regularity in codimension one, and
P4R.6's vanishing-sphere argument excludes
it. Let $s,c,g,b$ respectively count the
connected components of $C$, its
irreducible components, its genus-one
normalizations, and the cycle ranks of
its point-identification graphs.
Then $e(C)=c+s-2g-b$, $r\ge g$, and
$c\ge g$.

\item The proper-map pushout $T\cup_A C=S$
and its cofibration justification are
exactly those of P4R.7. A root which
maps to a point gives the zero
homomorphism on its fundamental group;
a noncontracted root maps through the
normalization of one component of $C$.
Every rational tree is simply connected.
Consequently killing $\pi_1(T)$ and
the normalized component groups of $C$
in the disconnected van Kampen
presentation leaves the free quotient
$F_{b+r-s}$, as in P4R.5. Virtual
abelianness gives $b+r-s\le1$, hence
\[
 e(C)\ge c+r-2g-1\ge-1.
\]

\item A negative value forces $c=r=g$.
Every root is then noncontracted, maps
onto an elliptic normalization, and
supplies one of the $c$ image components.
Every rational component of $A$ is
contracted: a nonconstant map from its
normalization $\mathbb P^1$ to an
elliptic curve contradicts
Riemann--Hurwitz. Thus every exceptional
block is a rational subtree of one
rooted component of $A$. It meets that
root in exactly one point. At least one
intersection is forced by the final
assertion of P4R.6, since the roots are
the only strict conductor components;
two intersections would make a cycle
in the rooted tree. The exceptional
fibre is connected. One can see this
directly by Stein factorization of
$T\to N$: its finite birational factor
is an isomorphism because $N$ is normal.

\item The conductor on $N$ consequently has
one local analytic branch at each of
its points. At a singular point of $N$,
each local branch has a strict transform
meeting the exceptional fibre, whereas
there is exactly one such intersection
with the smooth root. Away from those
fibres the assertion follows from the
smooth disjoint roots themselves.

\item The local hypersurface argument of
P4R.5 now applies with $N$ in place of
its smooth resolution. Two different
curve germs of $C$ would require two
different normalization points, because
each point has just one conductor germ.
Their corresponding hypersurface
branches intersect in a curve germ.
That germ belongs to both conductor
images and forces the two chosen germs
to coincide. Thus $C$ is locally
unibranch with disjoint irreducible
components. This gives $b=0$, $s=c=g$,
and $e(C)=0$, a contradiction.
Therefore $e(C)\ge0$, proving the
assertion.

\end{enumerate}
\end{proof}

\begin{proposition}[P4R.9]\hypertarget{tag-02YL}{}\label{night:P4R.9}\label{stmt:P4R.9}
\StatementRecord{night:P4R.9}{proved}{02YL}{\SeriesRef{night:P4.13}, \SeriesRef{night:P4.22}, \SeriesRef{night:P4R.6}, \SeriesRef{night:P4R.7}}

In Part~4.2, every irreducible singular
support whose normalization has a
resolution birationally ruled over an
elliptic curve has nonnegative Euler
number. When the horizontal part of
the pole divisor on the minimal ruled
model is reduced, the normalization
is smooth and minimal and the support
has Euler number zero.
\end{proposition}
\begin{proof}

\begin{enumerate}[label={$\langle1\rangle$\arabic*.},leftmargin=*]
\item Push the integral divisor of P4R.6 to
the minimal ruled surface $T_0\to B$.
Its horizontal degree is two. A reduced
horizontal component $H$ satisfies
\[
 0\le 2g(\widetilde H)-2
 \le 2p_a(H)-2
 =-(D_0-H)H\le0.
\]
The first inequality is
Riemann--Hurwitz over the elliptic
base. All inequalities are equalities:
every such component is smooth elliptic,
unramified over $B$, and disjoint from
the other components of $D_0$. Every
vertical fibre meets the horizontal
part, so no vertical component occurs.
The reduced-horizontal possibilities
are therefore one smooth elliptic
bisection or two disjoint elliptic
sections, each with coefficient one.

\item The reverse blowup formula $\mu-1$
then keeps $A=\operatorname{Supp}D$
a disjoint union of their smooth
strict transforms; each new exceptional
curve is absent from $A$. Every
exceptional curve of the normalization
resolution must belong to $A$, by
P4R.6. An exceptional component would
therefore be one of these disjoint
elliptic curves and would be disjoint
from the strict conductor, contrary
to P4R.6. Thus the normalization is
smooth. P4.22, with its hypotheses now
verified, proves minimality and
$e(S)=0$.

\item In the remaining case
$D_0=2H+\sum a_jF_j$ with $H$ a
smooth section. The same reverse
blowups make $A$ the elliptic root $H$
with rational trees attached. Its
inclusion in $T$ induces an
isomorphism on fundamental groups,
by the ruling and invariance under
point blowups. The resolution pushout
gives $\pi_1(S)=\pi_1(C)$ as in
P4R.7. The virtually abelian curve
group calculation gives $e(C)\ge0$,
and P4.13 gives $e(S_{\rm reg})\ge0$.
Their sum proves $e(S)\ge0$ also
in this remaining case.

\end{enumerate}
\end{proof}
\subsection{The Hessian line and the tacnodal equality case}

\begin{proposition}[P4R.10]\hypertarget{tag-02Y3}{}\label{night:P4R.10}\label{stmt:P4R.10}
\StatementRecord{night:P4R.10}{proved}{02Y3}{\SeriesRef{night:P4.29}, \SeriesRef{night:P4.30}}

Let $S$ be an irreducible reduced fibre support in Part~4.2.
Suppose that its reduced singular curve $C$ is irreducible and
unibranch at every point, that its normalization $j:E\to C\hookrightarrow Y$
is elliptic, and that its generic transverse germ is a tacnode
$x^2-y^4=0$. Let $C_{\mathrm{sch}}$ be the image conductor scheme.
Then
\[
 \chi(\mathcal O_{C_{\mathrm{sch}}})\leq0.
\]
Equality makes $C$ smooth and makes the transverse Hessian of
the reduced defining equation have rank one at every point of $C$.
\end{proposition}
\begin{proof}

\begin{enumerate}[label={$\langle1\rangle$\arabic*.},leftmargin=*]
\item Let $P\subset\mathcal O_Y$ be the reduced ideal of $C$, and let
$I\subset\mathcal O_Y$ be the inverse image of the conductor ideal
of $S$. The conductor depth argument of P4.29 gives a
Cohen--Macaulay quotient $\mathcal O_Y/I$ of pure dimension one.
At the generic tacnode the normalization has two smooth branches;
its conductor is $(x,y^2)$ in the displayed coordinates.
Consequently $I$ contains $P^{(2)}$. Indeed a section of $P^{(2)}$
has zero image in $\mathcal O_Y/I$ at the generic point, and any
nonzero image would generate a finite-length submodule, contrary
to depth one. We obtain
\[
 0\longrightarrow J:=P/I\longrightarrow
 \mathcal O_{C_{\mathrm{sch}}}\longrightarrow\mathcal O_C
 \longrightarrow0.
\]
Here $J^2=0$, because $P^2\subset P^{(2)}\subset I$.
The module $J$ is therefore an $\mathcal O_C$-module. It is
torsionfree of rank one: a finite-support torsion submodule would
again be a finite-length submodule of $\mathcal O_Y/I$.

\item On $E$, put
\[
 K=\ker\bigl(j^*\Omega_Y\longrightarrow\Omega_E\bigr),
 \qquad\operatorname{im}(dj)=\Omega_E(-R).
\]
These are a rank-two bundle and a line bundle respectively.
The coefficient of $R$ at a point with local parametrization
$x_i=x_i(t)$ is $r=\min_i\operatorname{ord}_t x_i'(t)$.
We have $\deg K=\deg R$, since $\deg j^*\Omega_Y=0$
and $E$ is elliptic. Differentiation identifies
$G=P/P^{(2)}$ with a subsheaf of $j_*K$, as in P4.30.
Write $E_0\subset K$ for the image of its pullback to $E$.
This is a rank-two locally free subsheaf with finite-length
quotient.

\item The Hessian of a local reduced equation $f$ of $S$ annihilates
the tangent derivative of $C$, since all first derivatives of
$f$ vanish on $C$. It thus defines a quadratic tensor in
$\operatorname{Sym}^2K$. Under a change $f\mapsto u f$ its
restriction to $C$ multiplies by $u$, so globally it is a section
\[
 q\in H^0\bigl(E,\operatorname{Sym}^2K\otimes j^*\mathcal O_Y(S)\bigr).
\]
It has generic rank one. Let $L\subset K$ be the saturation of
its generic image line. On the smooth curve $E$ this is a line
subbundle. The tensor lies in $L^{\otimes2}\otimes
j^*\mathcal O_Y(S)$ everywhere, since that is a subbundle and
the assertion holds on a dense open set. Write its zero orders
as $d_p\geq0$. The normal line of $S$ has degree zero on $E$:
locally over the base disc the divisor $mS$ is principal.
It follows that
\[
 2\deg L=\sum_p d_p.
\]

\item Generically $I/P^{(2)}$ is exactly the line of the linear
conductor equation $x$, which is the Hessian line. Projection
to $K/L$ therefore gives an injection
\[
 J\longrightarrow j_*(K/L).
\]
To check the kernel at the exceptional points, a section whose
projection vanishes generically has torsion image in $J$; the
torsionfreeness just proved makes that image zero. Let $b_p$ be
the colength of the projected lattice $E_0\to K/L$ at $p$.
Let $q_p$ be the additional nonnegative length between $J$ and
the direct image of this normalized image lattice. Riemann--Roch
on $E$ now gives the exact formula
\[
 \chi(J)=\deg R-\frac12\sum_p d_p-\sum_p b_p-\sum_p q_p.
\]
We prove locally that $d_p/2+b_p\geq r_p$.

\item At a smooth point of $C$, $r_p=0$, so the claim follows from
nonnegativity. At a singular point put
$m=\min_i\operatorname{ord}_t x_i(t)$, so $r_p=m-1$.
First suppose the embedding dimension is three. Every defining
equation of $C$ has no linear term. Its gradient along the
parametrization has all entries of order at least $m$.
Every two-by-two minor of the conormal image consequently has
order at least $2m$. Since $K$ is saturated in the ambient free
rank-three module, a generator of $\bigwedge^2K$ has a unit
Pl\"ucker coordinate. Thus
\[
 \tau_p:=\operatorname{length}(K/E_0)_p\geq2m.
\]
Let $a_p=\operatorname{length}(L/(E_0\cap L))_p$.
The exact lattice sequence gives $\tau_p=a_p+b_p$.
Choose a primitive local generator $\ell$ of $L$ and write
$q=h\ell^2$ after trivializing the normal line, so
$\operatorname{ord}_t h=d_p$.
Every column of the Hessian is the differential of one of
the first derivatives of $f$. Those first derivatives belong
to $P$, so the columns lie in $E_0$. At least one ambient
coordinate of the primitive vector $\ell$ is a unit; its
corresponding column shows that $h\ell\in E_0$.
Therefore $d_p\geq a_p$, and
\[
 \frac{d_p}{2}+b_p\geq\frac{a_p}{2}+b_p
 \geq\frac{\tau_p}{2}\geq m>m-1=r_p.
\]

\item Now suppose the embedding dimension is two. Choose coordinates
with $P=(z,H(x,y))$, where $H$ is an irreducible reduced plane
equation with no linear term. Here $P^{(2)}=P^2$.
For an explicit verification, $z,H$ is a regular sequence, and
$P/P^2$ is free over the domain $\mathcal O_Y/P$.
If $s\notin P$ and $sv\in P^2$, reduction modulo $P$ gives
$v\in P$; reduction in this free conormal module then gives
$v\in P^2$. Thus $P^2$ is $P$-primary and equals its symbolic
square. The equation $f$ belongs to this square, so
\[
 f=A z^2+B zH+C H^2.
\]
Let $u$ be a primitive generator of the plane kernel of the
derivative row. We have $dH=w u$, with
$a=\operatorname{ord}_t w\geq m$, because both plane entries
of $dH$ have order at least $m$. In the basis $(dz,u)$ of $K$
we consequently have
\[
 E_0=\langle dz,w u\rangle,
 \qquad q=2A\,dz^2+2Bw\,dz\,u+2Cw^2u^2.
\]
Write the primitive generator of $L$ as
$\ell=\ell_z dz+\ell_u u$, and again $q=h\ell^2$ with
$\operatorname{ord}_t h=d_p$.
When $\ell_z$ is a unit, the projected conormal lattice has
colength $b_p=\min(a,\operatorname{ord}_t\ell_u)$.
The coefficient of $u^2$ in the displayed tensor gives
\[
 d_p+2\operatorname{ord}_t\ell_u\geq2a,
 \qquad d_p/2+b_p\geq a\geq m.
\]
The same inequality holds when $\ell_u=0$, since then $b_p=a$.
When $\ell_z$ is not a unit, $\ell_u$ is a unit. Projection of
$dz$ generates $K/L$, so $b_p=0$; the $u^2$ coefficient gives
$d_p\geq2a$, proving the same bound. These cases cover all
singular curve points and establish the claimed local inequality.

\item We conclude that $\chi(J)\leq0$. Since normalization gives
$\chi(\mathcal O_C)=-\sum\delta(C,p)$, the conductor exact
sequence yields
\[
 \chi(\mathcal O_{C_{\mathrm{sch}}})
 =-\sum_p\delta(C,p)+\chi(J)\leq0.
\]
Equality forces every delta invariant to vanish, so $C=E$ is
smooth. Then $R=0$, $E_0=K$, and all normalization descent
defects vanish. The exact formula reduces to
$\chi(J)=-\tfrac12\sum d_p$, so equality also forces
$d_p=0$ everywhere. The Hessian has rank one at every point.
The program \texttt{code/part4/tacnode\_hessian.py} verifies the
two lattice inequalities for the stated bounds, the exact
rank-one Hessian on the singular parametrization
$(t^2,t^3,0)$, and the simple Hessian zero of the permitted
local model $zx^2+xy^2$. Its JSON and text outputs retain
the bounds; the preceding inequalities prove the unrestricted
statement.

\end{enumerate}
\end{proof}

\begin{proposition}[P4R.11]\hypertarget{tag-02Y4}{}\label{night:P4R.11}\label{stmt:P4R.11}
\StatementRecord{night:P4R.11}{proved}{02Y4}{\SeriesRef{night:P4.12}, \SeriesRef{night:P4.22}, \SeriesRef{night:P4.29}, \SeriesRef{night:P4R.10}, \SeriesRef{night:P4R.6}, \SeriesRef{night:P4R.9}}

Let $S$ be an irreducible singular support in Part~4.2 whose
normalization has ruled resolution, and suppose $e(S)=0$.
In the double-horizontal-divisor case with generic transverse
tacnode, the normalization is a smooth minimal elliptic ruled
surface. The conductor image is a smooth elliptic curve, its
reduced preimage is an unramified covering of degree two,
and every local germ is a product tacnode.
\end{proposition}
\begin{proof}

\begin{enumerate}[label={$\langle1\rangle$\arabic*.},leftmargin=*]
\item Use $T\to N\to S$, the horizontal section $H$, and
$C=\operatorname{Sing}(S)_{\mathrm{red}}$ as in P4R.9.
The fundamental-group and Euler argument of P4.29 gives an
irreducible curve $C$ with elliptic normalization, unibranch
at every point. All rational curves in the resolved conductor
are contracted, since none maps nonconstantly to an elliptic
curve. Thus $H$ is the only strict conductor component.
Its finite map to $C$ has degree two, and its genus $g$ is at
least one by Riemann--Hurwitz.

\item Every exceptional curve of $T\to N$ is vertical for the ruling
$T\to H$. The normalized-graph construction of P4.29 descends
this ruling to $N\to H$. Its Leray calculation proves
$R^1\rho_{N*}\mathcal O_T=0$, and hence
$\chi(\mathcal O_N)=1-g$. Finite duality identifies the
conductor on $N$ with $\omega_N\otimes\nu^*\omega_S^{-1}$.
The latter line twist is torsion by P4.12. Serre duality,
the rational-singularity comparison and Riemann--Roch on $T$
therefore give $\chi(I)=1-g$ and
$\chi(\mathcal O_{D_{\mathrm{sch}}})=0$.
The exact conductor square, with $\chi(\mathcal O_S)=0$,
then gives
\[
 \chi(\mathcal O_{C_{\mathrm{sch}}})=g-1.
\]
P4R.10 makes this at most zero. Thus $g=1$, $C$ is smooth,
and the reduced defining equation has rank-one Hessian along
all of $C$. Riemann--Hurwitz makes $H\to C$ unramified.

\item We verify the local analytic conclusion. Over any point of
$C$, its two preimages on $H$ cannot be identified in $N$:
each exceptional block is vertical for the descended ruling
and hence meets only one point of the section $H$.
There are therefore two local irreducible branches of $S$.
The holomorphic Morse lemma with parameters writes its reduced
equation as $X^2+g(y,z)$. Its two analytic factors are monic
linear polynomials in $X$, by Weierstrass factorization, and
their coefficients sum to zero. Each branch is a graph and
is smooth. Their intersection has support exactly the smooth
curve $C$. Consequently, in coordinates with $C=\{X=y=0\}$,
the factor coefficient is a unit times a power of $y$.
The generic tacnode fixes that power to two. Extracting a square
root of the unit in the $y$ coordinate gives
$X^2-y^4=0$ times the $z$ disc. This also verifies the actual
normalization and conductor morphisms locally.
The normalization was already smooth away from the conductor
by P4R.6. It is therefore smooth everywhere, and P4.22 makes
this elliptic ruled surface minimal.

\end{enumerate}
\end{proof}

\subsection{The components outside the pole divisor in a ruled fibre}

\begin{proposition}[P4R.12]\hypertarget{tag-02Y5}{}\label{night:P4R.12}\label{stmt:P4R.12}
\StatementRecord{night:P4R.12}{proved}{02Y5}{\SeriesRef{night:P4R.3}, \SeriesRef{night:P4R.6}, \SeriesRef{night:P4R.9}}

Let $T\to N\to S$ be the minimal normalization resolution of an
irreducible support in Part~4.2 with ruled minimal model of positive
base genus. Let $r:T\to B_0$ be its ruling, and let $D$ be the
effective integral pole divisor of P4R.6. Every component $E$ of a
reducible ruling fibre which is absent from $\operatorname{Supp}D$
is a smooth rational curve of self-intersection $-1$. It meets
exactly one pole component, transversely once, and that component
has pole coefficient one. In the double-horizontal-section case,
$E$ is disjoint from the horizontal section. Distinct such
components are disjoint; each is a leaf of the reduced ruling fibre.
These assertions retain the multiplicities of the ruling fibre
and do not require the multiplicity of $E$ to be one.
\end{proposition}
\begin{proof}

\begin{enumerate}[label={$\langle1\rangle$\arabic*.},leftmargin=*]
\item A reduced horizontal pole divisor occurs here only over an
elliptic base, by P4R.3, and P4R.9 then makes the normalization
smooth and minimal. In that case there is no reducible ruling
fibre. We may therefore work in the double-horizontal-section
case throughout the proof.
A ruling fibre on $T$ is obtained from a smooth $\mathbb P^1$
by the actual sequence of point blowups. Its reduced support is
a normal-crossing tree of smooth rational curves. Each component
of a reducible fibre has negative self-intersection. One can
verify the latter assertion without a contraction criterion:
if $F_b=\sum m_iE_i$ is the full fibre, then
$0=F_bE_i=m_iE_i^2+\sum_{j\ne i}m_jE_jE_i$;
connectedness and reducibility make the last sum positive.

\item For $E$ outside the pole support, effectiveness and adjunction give
\[
 0\le DE=-K_TE=2+E^2.
\]
Thus $E^2$ is $-1$ or $-2$. In the second case $DE=0$,
so $E$ is disjoint from the full pole support. The map to $S$
then identifies a neighbourhood of this compact curve with a
neighbourhood in $S_{\rm reg}$, by P4R.6. The local cyclic
base change which removes the multiplicity of the support is
unramified over this regular locus. Its restriction to $E$ is
a disjoint union of copies of $\mathbb P^1$, because
$\pi_1(\mathbb P^1)=0$. A chosen copy has the same normal
degree $-2$. A smooth local trivialization over a sufficiently
small parameter disc transports it into a nearby four-torus,
still with self-intersection $-2$. Every map from $S^2$ to that
torus is null-homotopic, since its universal cover is
$\mathbb R^4$; its homology class and self-intersection must
therefore be zero. This contradiction leaves $E^2=-1$ and
$DE=1$.

\item All pole coefficients and all intersection multiplicities are
positive integers at an intersection. Hence $DE=1$ says exactly
that there is one intersection, transverse and on a coefficient-one
component. A horizontal pole section of coefficient two cannot
be that component.

\item It remains to exclude an intersection of two components $E,E'$
outside the pole divisor. Both have square $-1$, and both meet
the pole support. In the full ruling fibre, let their positive
multiplicities be $m,m'$. The two fibre-intersection equations
would give
\[
 m=m'+a,\qquad m'=m+b,
\]
where $a,b$ are positive sums of the multiplicities of their
other vertical neighbours; their respective pole neighbours make both
sums positive. Adding these equations gives $a+b=0$, a
contradiction. Thus no two outside components meet. Each has
exactly the one pole neighbour already found and is a leaf.

\end{enumerate}
\end{proof}

\subsection{Restriction to an elliptic conductor component}

\begin{proposition}[P4R.13]\hypertarget{tag-02Y6}{}\label{night:P4R.13}\label{stmt:P4R.13}
\StatementRecord{night:P4R.13}{proved}{02Y6}{\SeriesRef{night:P4.30}, \SeriesRef{night:P4.31}, \SeriesRef{night:P4.33}, \SeriesRef{night:P4R.10}, \SeriesRef{night:P4R.6}, \SeriesRef{night:P4R.7}, \SeriesRef{night:P4R.9}}

In Part~4.2, let an irreducible singular support $S$ have a
normalization with ruled resolution. If its reduced image
conductor has a component with elliptic normalization, then
the normalization of $S$ is a smooth minimal elliptic ruled
surface, the image conductor is one smooth elliptic curve,
and $e(S)=0$. Every local germ along this curve is the product
of its generic transverse germ with a disc. In particular no
rational conductor component is attached to that curve.
\end{proposition}
\begin{proof}

\begin{enumerate}[label={$\langle1\rangle$\arabic*.},leftmargin=*]
\item Let $\rho_N:T\to N$ be the minimal normalization resolution,
let $\nu:N\to S$ be the finite normalization, and put
$\rho=\nu\rho_N$. The reduced-horizontal case is already
proved by P4R.9. We therefore use the actual double-horizontal
pole divisor $D$ of P4R.6. Its support is a smooth section $H$
with rational trees attached. The pushout and curve-group
calculation of P4R.7--P4R.9 show that the reduced image
conductor $C$ is connected, that its point-identification
graph has no cycles, and that it has exactly one normalized
component of positive genus, of genus one. Denote the
corresponding reduced component by $C_E$ and its normalization
by $j:E\to C_E\hookrightarrow Y$. In particular every point
of $C_E$ is unibranch. Every other normalized component of
$C$ is rational.

\item A rational component of $D$ cannot map nonconstantly to $E$.
Thus $H$ maps nonconstantly to $E$. Write $g=g(H)\ge1$.
At a general point of $C_E$, every normalization branch has
conductor exponent two, because the coefficient of $H$ in
$D$ is two. The intrinsic conductor exponent of a singular
plane branch is at least two, and each additional branch
increases it by its positive intersection multiplicity.
The possibilities are therefore a single ordinary cusp,
two smooth branches with intersection multiplicity two, or
three smooth branches with pairwise intersection multiplicity
one. These are $A_2,A_3,D_4$, respectively, and $H\to E$
has degree $1,2,3$ in the three cases.

\item We first construct a map which detects global functions on
the entire conductor, including its rational components.
Since $H$ is not contracted, P4.33 proves that $N$ has
rational singularities. Put $L_N=\nu^*\omega_S$.
Finite duality identifies its conductor ideal with
\[
 I_N=\omega_N\otimes L_N^{-1}.
\]
The canonical trace is an isomorphism
$\rho_{N*}\omega_T=\omega_N$ for rational analytic
singularities. We use Greb, Remark~2.1, p.~415
\cite{Greb2011}. The hypotheses hold here: $N$ is normal
and Cohen--Macaulay, $\rho_N$ is a proper resolution,
P4.33 gives $R^1\rho_{N*}\mathcal O_T=0$, and its
one-dimensional fibres give vanishing of the higher
structure-sheaf direct images. The integral pole identity
of P4R.6 and projection formula consequently give
\[
 \rho_{N*}\mathcal O_T(-D)
 =\rho_{N*}(\omega_T\otimes\rho_N^*L_N^{-1})
 =I_N.
\]
These are equalities of ideal subsheaves of $\mathcal O_N$:
the pole identity uses the pullback of the same residue
section which defines finite duality, and the trace agrees
with it on the birational isomorphism locus. Equality there
determines the maps between these torsionfree sheaves.
The divisor ideal sequence therefore gives an injection
\[
 \mathcal O_N/I_N\longrightarrow\rho_{N*}\mathcal O_D.
\]
The finite conductor square gives an injection from the
image conductor scheme. Composing them yields the actual
morphism of coherent $\mathcal O_S$-algebras
\[
 \mathcal O_{C_{\rm sch}}\longrightarrow
 \nu_*(\mathcal O_N/I_N)\longrightarrow\rho_*\mathcal O_D,
 \qquad\hbox{both arrows injective}.                         \tag{P4R.13.1}
\]

\item We show that a global function on $D$ is determined by
its value at the generic doubled horizontal section.
At every point blowup $\beta:T'\to T''$ in the actual
map to the minimal ruled surface, the pole divisors obey
$D'=\beta^*D''-E_\beta$. The elementary direct images are
\[
 \beta_*\mathcal O(E_\beta)=\mathcal O,
 \quad R^1\beta_*\mathcal O(E_\beta)=0,
 \quad\beta_*\mathcal O=\mathcal O,
 \quad R^1\beta_*\mathcal O=0.
\]
For the first equality, a meromorphic function allowed
a pole only on $E_\beta$ descends off the centre and
extends across that point by Hartogs. For the vanishing,
the two blowup charts, or their exceptional-curve
filtration, have successive terms
$\mathcal O_{\mathbb P^1}(j)$ with $j\ge-1$;
their first cohomology vanishes. Equivalently one applies
$0\to\mathcal O\to\mathcal O(E_\beta)
\to\mathcal O_{E_\beta}(-1)\to0$ to the usual
blowup structure-sheaf calculation. Projection formula,
the divisor ideal sequence and the absence of second
direct images then give
\[
 \beta_*\mathcal O_{D'}=\mathcal O_{D''},
 \qquad R^1\beta_*\mathcal O_{D'}=0.                         \tag{P4R.13.2}
\]
These particular maps identify $H^0(D,\mathcal O_D)$
with $H^0(D_0,\mathcal O_{D_0})$, where
\[
 D_0=2H_0+V_0,
 \qquad V_0=\sum_j a_jF_j
\]
on the minimal ruled surface. The kernel of restriction
from $D_0$ to $2H_0$ is $\mathcal O_{V_0}(-2H_0)$.
Its successive layers on each thick fibre $a_jF_j$ are
$\mathcal O_{\mathbb P^1}(-2)$, since $F_j^2=0$ and
$H_0F_j=1$. They have no global sections. Restriction
to $2H_0$ is therefore injective on global functions.
This double section is Cohen--Macaulay and has no
zero-dimensional associated point, so a section which
vanishes generically on it vanishes everywhere.

\item Let $C_{E,\rm sch}$ be the primary component of
$C_{\rm sch}$ supported on $C_E$: locally its ideal
is obtained by localizing the conductor ideal at the
generic point of $C_E$ and contracting. Primary
decomposition of the pure one-dimensional conductor
quotient makes these ideals coherent and compatible
on overlaps. Restriction is a surjection
\[
 \mathcal O_{C_{\rm sch}}\longrightarrow
 \mathcal O_{C_{E,\rm sch}}.
\]
Its map on $H^0$ is injective. Indeed, a global element
of its kernel maps by (P4R.13.1) to a function on $D$
which vanishes generically on $2H$; (P4R.13.2) and the
preceding paragraph make it zero. Its map on $H^1$
is surjective, because the kernel is supported in
dimension at most one and has no second cohomology.
We have proved
\[
 \chi(\mathcal O_{C_{\rm sch}})
 \le\chi(\mathcal O_{C_{E,\rm sch}}).                       \tag{P4R.13.3}
\]

\item We evaluate the right side using the actual primary
ideal, without removing the rational components from
the left side. In type $A_2$ the primary conductor
is reduced: a pure one-dimensional quotient which
is generically reduced has no nilpotent submodule
of finite support. Thus its Euler characteristic
is $-\sum_p\delta(C_E,p)\le0$. In type $D_4$ the
primary ideal is the symbolic square of the ideal
of $C_E$, by exactly the generic-conductor and
depth-one argument of P4.31. P4.30 gives the same
nonpositive bound, with equality only for smooth
$C_E$.

\item In type $A_3$, the primary ideal contains that
symbolic square and its nilpotent module is
torsionfree of rank one on $C_E$. The proof of
P4R.10 applies to this primary quotient: its
differentiation map, saturated Hessian line and
projection to the quotient conormal line depend
only on this ideal and on the actual equation of
$S$. The presence of other conductor components
changes none of those maps. More explicitly its
Euler characteristic is
\[
 -\sum_p\delta(C_E,p)+\deg R
 -\frac12\sum_p d_p-\sum_p b_p-\sum_p q_p\le0,
\]
with the same ramification divisor $R$, Hessian
zero orders $d_p$, projected lattice colengths
$b_p$ and descent defects $q_p$ as in P4R.10.
The local inequalities there prove this bound
at every singular point of $C_E$. Equality
makes $C_E$ smooth and its Hessian nowhere zero.

\item On the other hand P4.33 gives
$\chi(\mathcal O_{C_{\rm sch}})=g-1$. Thus
\[
 0\le g-1
 =\chi(\mathcal O_{C_{\rm sch}})
 \le\chi(\mathcal O_{C_{E,\rm sch}})\le0.                    \tag{P4R.13.4}
\]
All these quantities vanish. We obtain $g=1$,
$C_E=E$ smooth, and an unramified map $H\to E$.
In type $A_3$ its transverse Hessian has rank
one everywhere on $E$.

\item It remains to exclude attachments to this smooth
elliptic component. Put
\[
 K=\mathcal I_E/\mathcal I_E^2,
 \qquad N_E=\mathcal O_Y(S)|_E.
\]
Both $\deg K$ and $\deg N_E$ are zero. The first
follows from the conormal exact sequence,
$c_1(TY)E=0$, and $\deg\Omega_E=0$; the second
follows because a positive multiple of $S$ is
principal in a tube of this fibre.

\item In type $A_2$ let $L\subset K$ be the saturated
quadratic-image line and put $Q=K/L$. The nonzero
quadratic tensor is a section of
$L^2\otimes N_E$, so $2\deg L\ge0$.
The cubic jet on the null direction is an
intrinsic section of $Q^3\otimes N_E$.
Here is the coordinate check, including points
where the quadratic tensor vanishes. A change
of a second-order lift changes the cubic
coefficient by the polarization of the
quadratic tensor. That polarization vanishes
on the null direction. A change of tangential
lift contributes a tangential derivative of
the rank-one quadratic tensor; its evaluation
twice on the null direction is again zero.
Changing a reduced defining equation multiplies
the tensor by its transition unit. Thus the
stated line bundle and its holomorphic section
are defined on all of $E$. The section is
nonzero generically by type $A_2$, so
$-3\deg L\ge0$. Both degrees are zero,
and both sections are nowhere zero. The
parametrized Morse lemma and a cube root
of the residual unit now give the product
equation $X^2+Y^3=0$ at every point of $E$.

\item In type $A_3$ (P4R.13.4) already makes the
quadratic tensor nowhere zero, whence
$\deg L=\deg Q=0$. Its cubic null jet is
identically zero. Removing the quadratic
direction gives an intrinsic residual fourth
jet in $Q^4\otimes N_E$. In local coordinates
its relevant terms and coefficient are
\[
 a x^2+bxy^2+cy^4,
 \qquad q_4=c-\frac{b^2}{4a},\qquad a\ne0.                 \tag{P4R.13.5}
\]
Substitution of the critical coordinate
$x=-by^2/(2a)+O(y^3)$ proves the formula.
The change $x\mapsto x+vy^2$ sends $b$ to
$b+2av$ and $c$ to $c+bv+av^2$, leaving
$q_4$ unchanged. Changes of normal frames
and of equation give the claimed
$Q^4\otimes N_E$ transition. Tangential
lift changes do not affect a function
whose residual order is four. This section
is nonzero generically and has degree zero;
it never vanishes. The residual function
is $y^4$ times a unit, giving a product
tacnode at every point of $E$.

\item In type $D_4$, the equation belongs to
$\mathcal I_E^3$. This ideal is primary
because $E$ is smooth, so generic order
three proves membership everywhere.
The cubic normal jet belongs to
$\operatorname{Sym}^3K\otimes N_E$.
Its discriminant is a section of
$(\det K)^6\otimes N_E^4$, of degree
zero. It is nonzero generically, hence
nowhere zero. The three distinct normal
tangent directions give three smooth
holomorphic branches locally along $E$.
After straightening the first two to
$x=0,y=0$, the third is $ax+by=0$
with both coefficients units. Setting
$X=ax,Y=-by$ gives the product equation
$XY(X-Y)=0$.

\item In all three cases the entire singular
locus near every point of $E$ is $E$
itself. A rational conductor tree could
not meet it. Connectedness of $C$ then
forces $C=E$. The displayed local
product forms make $N$ smooth along
its conductor, and P4R.6 already makes
it smooth elsewhere. Hence $T=N$.
A finite normalization cannot contract
a curve, and a rational curve cannot
map nonconstantly to $E$. Thus $D=2H$
has no vertical component. The first
blowup of $2H_0$ in a nonempty actual
blowup sequence would create a vertical
pole component of coefficient one;
later blowups never remove its strict
transform from the pole support.
Consequently no blowup occurs and $N$
is minimal elliptic ruled. Finally
\[
 e(S)=e(N)-e(H)+e(E)=0.
\]
The exact coefficient changes, discriminant
weights and degree identities are executed
in \texttt{code/part4/elliptic\_primary\_restriction.py};
its JSON and text outputs are retained.

\end{enumerate}
\end{proof}

\subsection{The smooth critical curve of transverse type \texorpdfstring{$A_3$}{A3}}

\begin{proposition}[P4R.14]\hypertarget{tag-02Y7}{}\label{night:P4R.14}\label{stmt:P4R.14}
\StatementRecord{night:P4R.14}{proved}{02Y7}{\SeriesRef{night:P4.16}, \SeriesRef{night:P4.28}}

Let $S$ be an actual reduced support in Part~4.2. At a point
where the entire reduced critical locus of its local defining
function is one smooth curve, with generic transverse type
$A_3$, holomorphic coordinates put that function in one of
the forms
\[
 X^2+Y^4,\qquad X(Y^2+XZ).                              \tag{P4R.14.1}
\]
The normalization germs are respectively two smooth branches,
or a smooth plane and an $A_1$ quadric. No positive-genus
exceptional curve occurs over such a point.
\end{proposition}
\begin{proof}

\begin{enumerate}[label={$\langle1\rangle$\arabic*.},leftmargin=*]
\item Choose a coordinate $z$ on the critical curve and normal
coordinates $x,y$. Its quadratic transverse Hessian has generic
rank one. Diagonalization of the symmetric matrix over
$\mathbb C\{z\}$, as in P4.28, gives the quadratic term
$h(z)x^2$. The cubic coefficient on the null direction
vanishes identically, by the generic transverse type. With
weights two for $x$ and one for $y$, the weight-four terms are
\[
 h(z)x^2+a(z)xy^2+b(z)y^4,\qquad
 \Delta(z)=4h(z)b(z)-a(z)^2.                             \tag{P4R.14.2}
\]
Every remaining term has larger weight and $\Delta$ is
nonzero generically.

\item We first prove that $\Delta$ is a unit. A zero persists,
counted with multiplicity, under small independent constant
perturbations of $h,a,b$, by the argument principle on a
small parameter disc. Generic such perturbations make the
zeros simple and disjoint from $h=0$. Indeed $h=\Delta=0$
also requires $a=0$, which the independent parameters avoid;
where $h\ne0$, the derivative of $\Delta$ in the $b$
parameter is $4h\ne0$. These perturbations preserve the
critical axis and its transverse $A_3$ germ near the fixed
boundary. At a simple zero with $h\ne0$, the parametrized
Morse lemma gives
\[
 X^2+y^4U(y,z),\qquad U(0,0)=0,\quad U_z(0,0)\ne0.
\]
The actual coordinate $Z=U(y,z)$ gives $X^2+y^4Z$.
P4.16 supplies a compact spherical class with nonzero square
in its smoothing. The boundary comparison of P4.28 applies
here: near the boundary critical curve, the parametrized
Morse lemma and a fourth root of a unit give $X^2+Y^4$;
the tangential boundary coordinate makes the nonzero levels
submersive. Proper boundary transport therefore carries the
sphere to the actual nearby torus. Its universal cover is
$\mathbb R^4$, so a sphere has zero homology class and zero
square, a contradiction. Thus $\Delta$ is a unit.

\item For $h(0)\ne0$, Morse splitting and a fourth root of the
residual unit give the first form in (P4R.14.1). Suppose
$h(0)=0$. Then $a(0)\ne0$. The holomorphic function
\[
 \frac{f(y^2v,y,z)}{y^4}
\]
restricts at $y=0$ to $h(z)v^2+a(z)v+b(z)$. Its finite
root $-b(0)/a(0)$ at the origin is simple. The holomorphic
implicit-function theorem gives a root $v=v(y,z)$, and
division by the monic equation of its graph gives
\[
 X=x-y^2v(y,z),\qquad f=Xg.
\]
On $X=0$, the function $g$ is $y^2$ times a unit. Its
parameter Morse centre in $y$ is therefore zero there and
its critical value is zero. The parametrized Morse lemma,
with parameters $X,z$, gives the exact form
\[
 f=X\bigl(Y^2+XH(X,z)\bigr).                            \tag{P4R.14.3}
\]
The remaining critical value is divisible by $X$ by its
vanishing on $X=0$. The function $H(0,z)$ is not identically
zero, since the generic transverse Hessian has rank one;
at the point under consideration $H(0,0)=0$.

\item Put $G=2H+XH_X$. Off $X=0$, the critical equations of
(P4R.14.3) are $Y=0,H_z=0,G=0$. We prove $H_z(0,0)\ne0$.
Otherwise $H_z$ is a nonzero plane germ not divisible by
$X$, because $H(0,z)$ has finite order at least two.
No component of $H_z=0$ is contained in $G=0$: it would
give an additional critical-curve branch off $X=0$, whereas
the entire critical locus is the stated smooth axis. Hence
$(H_z,G)$ is an isolated complete intersection of positive
intersection multiplicity.

\item Replace $H$ by $H-\eta z-\varepsilon$. Its plane critical
equations are
\[
 H_z-\eta=0,\qquad G-2\eta z-2\varepsilon=0.              \tag{P4R.14.4}
\]
Their solutions are the fibres of the holomorphic map
\[
 \Psi(X,z)=\left(H_z,\frac{G-2zH_z}{2}\right).
\]
Its zero fibre has the positive local degree just proved.
Choose a small compact representative with no boundary zero.
That degree is unchanged on a sufficiently small open
parameter ball. The universal equations in (P4R.14.4) have
parameter Jacobian determinant two, so regular parameter
values give a nonempty finite set of transverse solutions.
We also choose such a value outside the images under $\Psi$
of $X=0$ and $H_X=0$. These exceptional images have real
four-dimensional measure zero: each source curve has
countably many smooth real two-dimensional strata, and the
smooth image of each has measure zero. The germ $H_X$ is
not identically zero. Indeed in that case $H$ would depend
only on $z$ and both $H_z$ and $G=2H$ would vanish on the
whole line $z=0$, contrary to the isolated intersection.
The two exceptional images cannot exhaust the open ball
on which the positive degree is retained. This argument
uses no analyticity assertion for the image of a nonproper
map.

\item At each chosen solution, the critical value is
\[
 f_{\eta,\varepsilon}=-\frac{X^3H_X}{2}\ne0.              \tag{P4R.14.5}
\]
For the indicated order of the two equations, put
\[
 J=\det\frac{\partial(G-2\eta z-2\varepsilon,\ H_z-\eta)}
                  {\partial(X,z)}.
\]
The $Y$-Hessian entry is $2X$, and the remaining Hessian
block at a critical point is
\[
 \begin{pmatrix}
 X(3H_X+XH_{XX})&X^2H_{Xz}\\
 X^2H_{Xz}&X^2H_{zz}
 \end{pmatrix}.
\]
Its determinant is $X^3J$, and the full Hessian determinant
is $2X^4J\ne0$. Thus these are isolated holomorphic Morse
points with nonzero critical values. Near the fixed boundary,
$H(0,z)$ remains nonzero under the small perturbation, so the
critical axis retains its transverse $A_3$ type. The same
proper boundary comparison transports a Morse vanishing
sphere, of square $-2$, into an actual nearby torus. This is
impossible. We have proved $H_z(0,0)\ne0$.

\item The holomorphic coordinate $Z=H(X,z)$, keeping $X,Y$, now
gives the second form of (P4R.14.1). Its normalization has
the plane $X=0$ and the normal quadric $Y^2+XZ=0$. The
quadric's minimal resolution has one rational $(-2)$-curve.
The first form factors as $(X+iY^2)(X-iY^2)$ and its two
normalization branches are smooth. This proves the stated
normalization conclusions.

\item The distinction between a pole coefficient and this
torus-smoothing conclusion can be checked by an exact
retained local test. Take
\[
 f=x\bigl[y^2+xz(x-z)(x-2z)\bigr].
\]
It has only the smooth $A_3$ critical axis $x=y=0$. The
second normalization branch is the double plane branched
over four distinct lines. Blowing up their common point
and normalizing gives a double cover of the exceptional
$\mathbb P^1$ branched in four points, hence an elliptic
curve of square $-2$. Its discrepancy is $-1$ and the
order of $x$ on it is one, giving pole coefficient two
for the displayed total reduced hypersurface. It is excluded
here by the polar deformation, not by the coefficient alone.
For $H=z(x-z)(x-2z)$, the deformation $H\mapsto H-1$ has
the exact Groebner basis
\[
 3x-5z^4-14z,\qquad 25z^6+50z^3-2.
\]
It gives six simple solutions; adjoining $x$, $H_x$ or
$J$ to the critical ideal gives the unit ideal in each
case. For $H\mapsto H-\varepsilon$ all six scale by
$\varepsilon^{1/3}$ and converge to the origin. The
identities and these exact rational calculations are
executed in \texttt{code/part4/smooth\_a3\_polar.py},
with JSON and text output retained.

\end{enumerate}
\end{proof}

\subsection{A nonzero Hessian and the entire critical curve}

\begin{proposition}[P4R.15]\hypertarget{tag-02Y8}{}\label{night:P4R.15}\label{stmt:P4R.15}
\StatementRecord{night:P4R.15}{proved}{02Y8}{\SeriesRef{night:P4.18}, \SeriesRef{night:P4.28}}

At every critical-curve point of an actual reduced support in
Part~4.2 with nonzero ambient Hessian, holomorphic coordinates
give the exact product germ
\[
                       X^2+Y^m,\qquad m\ge2.
\]
In particular its entire reduced critical curve is smooth.
At a singular point of the reduced critical curve, the ambient
Hessian vanishes. This assertion allows different generic
transverse types on the different branches before the exclusion.
\end{proposition}
\begin{proof}

\begin{enumerate}[label={$\langle1\rangle$\arabic*.},leftmargin=*]
\item A nonzero complex symmetric Hessian has a vector on which its
quadratic form is nonzero. The parametrized holomorphic Morse
lemma in that direction gives the exact equation
\[
                        f=X^2+g(Y,Z).
\]
The plane function $g$ is nonzero because the support is
reduced. Factor it in the convergent power-series ring as
\[
             g=u\prod_{i=1}^{r}h_i^{m_i},
             \qquad u(0)\ne0,\quad m_i\ge1,
\]
with distinct reduced irreducible factors $h_i$ vanishing at
the point. At least one exponent is at least two: a reduced
plane function has an isolated critical point, while the
critical locus here contains a curve.

\item We exclude $r\ge2$ by an actual deformation of these factors.
For two distinct factors $h_i,h_j$, the holomorphic map
\[
 (h_i,h_j):(\mathbb C^2,0)\longrightarrow(\mathbb C^2,0)
\]
has isolated zero fibre and positive finite local degree.
On a sufficiently small compact representative that degree
is retained for small target values. Regular such values
give a nonempty transverse intersection of
$h_i-s_i=0$ and $h_j-s_j=0$. Independently choose the
other constants $s_k$ so that at least one of these
intersection points avoids every other zero factor. This
is possible because the chosen intersection set is finite
and each forbidden $s_k$ belongs to a finite set of values.
All parameters can be arbitrarily small.

\item At the selected intersection, the two transverse factors
are holomorphic coordinates and every other factor is a
unit. Taking an actual local $m_i$th root of their unit
product and multiplying the first coordinate by it gives
the germ
\[
                        X^2+Y^{m_i}Z^{m_j}.
\]
P4.18 excludes every such germ by its compact spherical
class of nonzero square. The comparison here is to the
original torus degeneration, as follows. Choose a small
Milnor ball before moving the factors. Near its boundary,
the distinct reduced factors are smooth and disjoint.
The moved factor with exponent $m_i\ge2$ retains exactly
the transverse product type $A_{m_i-1}$ there; a factor
of exponent one remains a regular point of the surface.
Parametrized Morse coordinates and the appropriate root
of the remaining unit identify these boundary germs.
Solving the tangential boundary coordinate makes the
nonzero levels submersive. The proper boundary transport
used in P4.28 therefore identifies the needed nearby
levels and transports the spherical class into the actual
nearby torus. Its square cannot be nonzero there. Hence
there is only one factor.

\item Absorb the unit by writing $h=u^{1/m}h_1$, so that
$g=h^m$. The reduced plane function $h$ is irreducible.
For a singular $h$, its gradient has an isolated zero,
and the finite positive local degree of
$(h_Y,h_Z)$ gives nonempty finite regular fibres under
small constant changes of that gradient. Equivalently,
subtracting a sufficiently small generic linear form
from $h$ gives isolated Morse critical points. Add a
small constant to make their critical values nonzero.
Call the resulting function $h_s$, and choose one such
critical point with value $v\ne0$ and nonsingular Hessian
$Q$. For $f_s=X^2+h_s^m$, the critical value and Hessian
are exactly
\[
 f_s=v^m\ne0,\qquad
 \operatorname{Hess}(f_s)
   =\begin{pmatrix}2&0\\0&m v^{m-1}Q\end{pmatrix}.
\]
Thus the full Hessian determinant is
$2m^2v^{2m-2}\det Q\ne0$. The moving curve $h_s=0$
retains its transverse $A_{m-1}$ type on the fixed
boundary. The same proper comparison transports the
Morse vanishing sphere, of square $-2$, to a nearby
torus, a contradiction. Consequently $h$ is smooth,
and using it as the coordinate $Y$ gives the asserted
product equation.

\item The Hessian identity is executed symbolically for the
arbitrary positive integer $m$ in
\texttt{code/part4/hessian\_product.py}. Its retained
output also checks the moving-node example
$X^2+(YZ-\varepsilon)^m$ at a nonzero critical value.

\end{enumerate}
\end{proof}


\subsection{The tacnode and transverse-line critical curve}

\begin{proposition}[P4R.16]\hypertarget{tag-02Y9}{}\label{night:P4R.16}\label{stmt:P4R.16}
\StatementRecord{night:P4R.16}{proved}{02Y9}{\SeriesRef{night:P4.33}, \SeriesRef{night:P4.40}, \SeriesRef{night:P4.42}, \SeriesRef{night:P4.44}}

In the remaining positive ruled case with uniform transverse
type $A_3$, the three-branch point of the critical curve cannot
be the union of a plane tacnode and an independent transverse
line. Its local conductor defect is at least eleven, whereas
the total defect in the contracted-section case is nine.
\end{proposition}
\begin{proof}

\begin{enumerate}[label={$\langle1\rangle$\arabic*.},leftmargin=*]
\item The analytic ring calculation in P4.40 supplies actual ambient
coordinates for this curve. Its reduced ideal is
\[
 P=(xz,yz,h),\qquad h=y^2-x^4.
\]
Its three parametrizations are $(t,t^2,0)$, $(t,-t^2,0)$,
and $(0,0,t)$. The two tangent branches have intersection
multiplicity two, and their union meets the independent line
with intersection length one. The normalization sequence
therefore gives $\delta=3$. Every branch is immersive, so
its ramification divisor has coefficient zero.

\item We first determine the symbolic square and the actual
equation. The plane-tacnode ideal $(z,h)$ is a reduced
complete intersection. Its square is unmixed and agrees
with the square of each prime at its generic point; hence
its symbolic square is its ordinary square. Intersecting
with the square of the line ideal gives
\[
 P^{(2)}=(z,h)^2\cap(x,y)^2=P^2+(zh).
\]
Indeed an element of $(z,h)^2$ is $Az^2+Bzh+Ch^2$.
Its last two terms belong to $(x,y)^2$. The first belongs
to that ideal precisely when $A$ does, because $z$ is a
nonzerodivisor modulo $(x,y)^2$. The resulting three
generators $x^2z^2,xyz^2,y^2z^2$ are already in $P^2$.
Furthermore
\[
 xzh=(xz)h,\qquad yzh=(yz)h,\qquad
 z^2h=(yz)^2-x^2(xz)^2.
\]
The class of $zh$ modulo $P^2$ is therefore annihilated
by the ambient maximal ideal. It is nonzero, since it has
a nonzero cubic term and $P^2$ has order at least four.
Thus $P^{(2)}/P^2$ is exactly one copy of the residue field.

\item Let $f$ be the reduced defining function of the support.
Its generic transverse multiplicity is two, so
$f\in P^{(2)}$ and
\[
                       f=\kappa zh+f_0,\quad f_0\in P^2,
                       \quad\kappa\in\mathbb C.
\]
At $(t,t^2,0)$ with $t\ne0$, use $z,h$ as transverse
coordinates. The quadratic transverse part of $f$ has the
form
\[
 t^2a(t)z^2+(\kappa+t b(t))zh+c(t)h^2,
\]
where $a,b,c$ are holomorphic at $t=0$. The transverse
Hessian has rank one, so its determinant vanishes for
all sufficiently small nonzero $t$. Its constant term at
zero is $-\kappa^2$ (up to the harmless common convention
for the quadratic-form matrix). Hence $\kappa=0$, and
we have proved the exact inclusion $f\in P^2$.

\item Let $K_i$ be the saturated rank-two kernel of the derivative
of the $i$th parametrization, and let $E_i\subset K_i$ be
the span of the differentiated generators of $P$. On the
two tangent branches choose the bases
\[
                dz,\quad dy\mp2t\,dx.
\]
The generator matrices have columns
\[
 \begin{pmatrix}t&\pm t^2&0\\0&0&\pm2t^2\end{pmatrix}.
\]
Their elementary-divisor orders are $(1,2)$. On the
independent line use $dx,dy$; the matrix is
\[
                   \begin{pmatrix}t&0&0\\0&t&0\end{pmatrix},
\]
with orders $(1,1)$.

\item For completeness, the rank-one quadratic estimate of
P4.42 applies to these exact lattices as follows. Write
$E=t^a\mathbb C\{t\}e_1\oplus
t^b\mathbb C\{t\}e_2$, and write the saturated Hessian
line as $\ell=u e_1+v e_2$, with $(u,v)$ primitive.
Since $f\in P^2$, the tensor $t^d\ell^2$ belongs to
$\operatorname{Sym}^2E$. Therefore
\[
 \frac d2\ge\max(a-\operatorname{ord}u,
                         b-\operatorname{ord}v).
\]
Projection of $E$ to $K/\ell$ has colength
$b_{\rm proj}=\min(a+\operatorname{ord}v,
b+\operatorname{ord}u)$. Adding these two expressions
gives $d/2+b_{\rm proj}\ge a+b$; the identity follows
by subtracting the same two candidate quantities, with
the zero-component cases read directly. The local defect
formula of P4.42 consequently gives
\[
 \delta+\sum_i(d_i/2+b_i-r_i)+q
       \ge3+(1+2)+(1+2)+(1+1)=11.
\]
Here $q\ge0$ is the further descent colength, retained
rather than discarded from the definition. P4.44 contracts
the horizontal section, and P4.33 and (4.42a) give total
defect nine. The displayed lower bound contradicts that
total. This excludes the stated analytic curve germ.

\end{enumerate}
\end{proof}

The symbolic-square identities, all three normalized
conormal matrices, and the rank-one determinant obstruction
are checked exactly in
\texttt{code/part4/a3\_tacnode\_line.py}.


\subsection{The nonordinary three-branch tacnode sector}

\begin{proposition}[P4R.17]\hypertarget{tag-02YA}{}\label{night:P4R.17}\label{stmt:P4R.17}
\StatementRecord{night:P4R.17}{proved}{02YA}{\SeriesRef{night:P4.40}, \SeriesRef{night:P4.42}, \SeriesRef{night:P4.48}, \SeriesRef{night:P4R.10}, \SeriesRef{night:P4R.15}, \SeriesRef{night:P4R.16}}

In the remaining positive ruled case with uniform generic transverse
type $A_3$, every nonordinary three-branch critical-curve point has
local conductor defect strictly greater than nine. In particular,
the delta-four nonplanar cases have the explicit lower bounds
$13$, $21/2$, $11$, and $23/2$ below. The only three-branch germ
compatible with the total defect budget nine is the ordinary
nonplanar union of three coordinate axes.
\end{proposition}
\begin{proof}
\begin{enumerate}[label={$\langle1\rangle$\arabic*.},leftmargin=*]
\item 
Use the notation of P4R.10 and P4.42: \[
 D(C,f)=\delta(C)+
        \sum_i\left(\frac{d_i}{2}+b_i-r_i\right)+q,\qquad q\ge0.
\] Here \(d_i\) is the order of the saturated rank-one Hessian, \(b_i\) is the projected conormal colength, and \(r_i=m_i-1\) for a branch of multiplicity \(m_i\).

\par\medskip\noindent\textbf{1. Bounds without an ordinary-square assumption}\par\label{bounds-without-an-ordinary-square-assumption}

Let \(E_i\subset K_i\) be the normalized conormal lattice, with elementary divisors \((t^{e_1},t^{e_2})\), \(e_1\le e_2\). For the saturated Hessian line \(L_i\), put \[
 a_i=\operatorname{length}\bigl(L_i/(L_i\cap E_i)\bigr).
\] The exact lattice sequence gives \(a_i+b_i=e_1+e_2\). Every first derivative of \(f\) belongs to the radical ideal \(P\) of the critical curve. Thus every Hessian column belongs to \(E_i\). Write the Hessian as \(ut^{d_i}\ell_i^2\), with \(u\) a unit and \(\ell_i\) primitive. One coordinate of \(\ell_i\) is a unit; its column gives \(t^{d_i}\ell_i\in E_i\), so \(d_i\ge a_i\). Also \(b_i\ge e_1\), since \(E_i\subset t^{e_1}K_i\). Consequently \[
                  \frac{d_i}{2}+b_i
                    \ge e_1+\frac{e_2}{2}.           \tag{A}
\] This remains valid when a symbolic-square term occurs in \(f\).

Since the curve is nonplanar, \(P\) has no linear equation. Every defining gradient therefore has order at least \(m_i\) on branch \(i\), and \(b_i\ge m_i\). P4R.15 gives zero ambient Hessian, so \(f\) has order at least three and \(d_i\ge m_i\). For total multiplicity \(m=\sum m_i\), \[
                  D(C,f)\ge\delta(C)+3+\frac m2.       \tag{B}
\] Thus \(m\ge5\) gives \(D\ge19/2>9\). Only multiplicities three and four need inspection.

\par\medskip\noindent\textbf{2. Three smooth branches: the complete excess-two classification}\par\hypertarget{tag-0005}{}\label{three-smooth-branches-the-complete-excess-two-classification}

Put \[
 B=\bigoplus_{i=1}^3\mathbb C\{t_i\},\quad
 F^j=\bigoplus_{i=1}^3t_i^j\mathbb C\{t_i\},\quad
 A_0=\mathbb C_{\rm diag}+F^1.
\] The actual ring \(A\subset A_0\) has \(\operatorname{length}(A_0/A)=2\), since \(A_0\) has delta invariant two. Each branch projection of \(A\) is onto.

Let \(W_1\subset F^1/F^2=\mathbb C^3\) be its first-jet image. Each coordinate projection is nonzero. If \(W_1=\mathbb C^3\), Nakayama first gives \(\mathfrak m_AA_0=F^1\) and then, on the finite \(A\)-module \(A_0/A\), gives \(A=A_0\). This contradicts colength two. Thus \(\dim W_1=1\) or \(2\).

\par\medskip\noindent\textbf{2a. One-dimensional first jets}\par\label{a.-one-dimensional-first-jets}

Scale the branch parameters so \(W_1=\mathbb C(1,1,1)\). The two missing first jets use the entire colength, hence \(F^2\subset A\). The exact algebra is \[
 A=\mathbb C_{\rm diag}+\mathbb C(t_1,t_2,t_3)+F^2.
\] It is generated by \[
 x=(t_1,t_2,t_3),\quad y=(t_1^2,0,0),\quad z=(0,t_2^2,0).
\] Indeed \(y,z,x^2-y-z\) give all three quadratic branch ideals, and multiplication by \(x\) gives every higher order. The embedded model is exactly \[
       P=\bigl(y(y-x^2),yz,z(z-x^2)\bigr).             \tag{C}
\] The branches are \((t,t^2,0),(t,0,t^2),(t,0,0)\). Its differentiated generators give \(E_i=t_i^2K_i\) on all three branches. For example, on \((t,t^2,0)\) use the primitive basis \(dy-2t\,dx,dz\). Formula (A) gives \[
                         D\ge4+3\cdot3=13.            \tag{D}
\]

\par\medskip\noindent\textbf{2b. Two-dimensional first jets}\par\hypertarget{tag-003V}{}\label{b.-two-dimensional-first-jets}

The three branch evaluations are nonzero linear forms on \(W_1\). If no two are proportional, their squares are independent: choose the first two as \(s,t\); the third is \(as+bt\), \(ab\ne0\), whose square supplies the nonzero \(st\) coefficient. Thus products give the whole second-jet space. Multiplication by one first jet nonzero on all branches then gives every later jet. This would leave only one missing positive-jet dimension, contrary to colength two.

Two evaluations are therefore proportional. Scale and label the parameters to obtain \(W_1=\{(a,a,b)\}\). Its products give the same two-dimensional second-jet space. The one remaining missing dimension occurs there: if the second-jet image were full, multiplication by a nonvanishing first jet would make every later one full. Hence \(F^3\subset A\), and the pure second-jet image is exactly \(\{(c,c,d)\}\).

A lift of \(W_1\) has quadratic difference \(\lambda a+\mu b\). The analytic changes \(t_i\mapsto t_i+c_it_i^2\), on the first two branches, remove \(\lambda a\). Thus the exact ring has common constants, with first and second coefficients satisfying \[
                     a_1=a_2,\qquad c_1-c_2=\mu a_3, \tag{E}
\] and arbitrary independent terms of order at least three.

For \(\mu\ne0\), the two elements \[
 x=(t_1,t_2,0),\qquad y=(\mu t_1^2,0,t_3)
\] generate this whole ring. The exact identities \[
 \begin{aligned}
 xy/\mu&=(t_1^3,0,0),\\
 x^3-xy/\mu&=(0,t_2^3,0),\\
 y^3-\mu(xy)^2&=(0,0,t_3^3),\\
 y^2-\mu x^2y&=(0,0,t_3^2)
 \end{aligned}
\] give the independent cubic ideals and every permitted quadratic jet. Multiplication by \(x\) or \(y\) gives all the higher branch series. This ring is planar, contrary to the present hypothesis.

Thus \(\mu=0\). The ring is generated by \((t_1,t_2,0),(0,0,t_3),(t_1^3,0,0)\). Naming the independent transverse coordinate \(z\) gives the exact embedded model \[
                       P=(xz,yz,y(y-x^3)).            \tag{F}
\] Its branches are \((t,0,0),(t,t^3,0),(0,0,t)\), and its conormal elementary divisors are \[
                         (1,3),\quad(1,3),\quad(1,1).
\] On the second branch use \(dy-3t^2dx,dz\); the gradients give \(t^3(dy-3t^2dx)\) and \(t\,dz\). Without restricting a symbolic-square coefficient of \(f\), (A) gives \[
                  D\ge4+\frac52+\frac52+\frac32
                       =\frac{21}{2}>9.               \tag{G}
\]

These are actual analytic ring isomorphisms. In each nonplanar case the three displayed generators form a basis of the curve cotangent space. Lift them to the ambient local ring. Their linear Jacobian is invertible, so the inverse-function theorem gives an ambient coordinate isomorphism to the stated model.

\par\medskip\noindent\textbf{3. Multiplicity four: one cusp and two smooth branches}\par\hypertarget{tag-02QH}{}\label{multiplicity-four-one-cusp-and-two-smooth-branches}

The singular branch has multiplicity two and delta at least one. Joining three branches costs at least two further delta units. If the branch delta is two, all remaining gluing is minimal, and the local ring is \[
 \mathbb C_{\rm diag}
 +\mathfrak m_{\rm singular}
 +t_2\mathbb C\{t_2\}+t_3\mathbb C\{t_3\}.
\] Its embedding dimension is at least \(2+1+1=4\), impossible in a smooth threefold. A larger branch delta is impossible. Thus this branch is the ordinary cusp \(\mathbb C\{t^2,t^3\}\).

Set \[
 A_0=\mathbb C_{\rm diag}
 +\mathfrak m_{\mathbb C\{t^2,t^3\}}
 +t_2\mathbb C\{t_2\}+t_3\mathbb C\{t_3\}.
\] Now \(A_0\) has delta three, and \(A\subset A_0\) has colength one. The one-dimensional \(A\)-module \(A_0/A\) is killed by \(\mathfrak m_A\), so \(\mathfrak m_AA_0\subset A\). Its elements have zero constant term, giving \(\mathfrak m_AA_0=\mathfrak m_A\). Hence \(\mathfrak m_A\) is an ideal of \(A_0\). The local algebra \(A_0/\mathfrak m_A\) has length two, so its maximal ideal squares to zero. Therefore \[
                         \mathfrak m_0^2\subset A.     \tag{H}
\]

The four cotangent coordinates of \(A_0\), denoted \(a,b,c,d\), are represented by the cusp functions \(t^2,t^3\) and the two smooth parameters. The three-dimensional subspace \(W=\mathfrak m_A/\mathfrak m_0^2\) is a hyperplane \[
                 \lambda a+\mu b+\nu c+\xi d=0.       \tag{I}
\] Its projection to \(a,b\) is onto, since the actual branch image is the full cusp ring, and the two smooth projections are nonzero. Ambient embedding dimension three gives \(\mathfrak m_A^2=\mathfrak m_0^2\). Thus products of \(W\) must generate \(\mathfrak m_0^2/\mathfrak m_0^3\), whose coordinates are \[
                              a^2,\quad ab,\quad c^2,\quad d^2.
                                                               \tag{J}
\]

\par\medskip\noindent\textbf{3a. A nonzero coefficient of the cusp cubic}\par\hypertarget{tag-0004}{}\label{a.-a-nonzero-coefficient-of-the-cusp-cubic}

For \(\mu\ne0\), solve (I) for \(b\). Surjectivity of the cusp projection forces \((\nu,\xi)\ne(0,0)\). The three minimal ambient coordinate functions have parametrizations \[
 \begin{array}{c|c}
 \text{cusp}&(t^2-\lambda't^3,-\nu't^3,-\xi't^3)\\
 \text{second branch}&(0,t,0)\\
 \text{third branch}&(0,0,t).
 \end{array}
\] Their tangent directions are independent.

By P4R.15, \(f\in\mathfrak m^3\). Let \(G_3\) be its homogeneous cubic part. Since \(df\) vanishes on all three branches, \(dG_3\) vanishes on each independent tangent axis. Every cubic monomial other than \(xyz\) has a derivative nonzero on an axis. Hence \(G_3=cxyz\). On the \(x\)-axis its normal Hessian is \[
                         \begin{pmatrix}0&c\\c&0\end{pmatrix}.
\] A nonzero \(c\) gives a nonzero leading minor on the actual branch, contradicting generic rank one. Thus \(c=0\), and \(f\in\mathfrak m^4\).

The Hessian entries have order at least \(2m_i\) on branch \(i\), and \(b_i\ge m_i\). Its defect contribution is at least \(m_i+1\). Therefore \[
                            D\ge4+(4+3)=11>9.         \tag{K}
\]

\par\medskip\noindent\textbf{3b. A zero coefficient of the cusp cubic}\par\hypertarget{tag-003U}{}\label{b.-a-zero-coefficient-of-the-cusp-cubic}

For \(\mu=0\), surjectivity in (J) forces \(\lambda,\nu,\xi\) all to be nonzero. Otherwise two of \(a,c,d\) are proportional on \(W\), so their squares are dependent, and (J) has rank at most three. That creates an extra minimal generator, contrary to embedding dimension three. Scale the branch parameters so that \(a+c+d=0\).

The pure hyperplane generators belong to \(A\) exactly by (H). They give the parametrizations \[
                       (t^2,t^2,t^3),\quad(t,0,0),\quad(0,t,0).
\] Their defining ideal is \[
                 P=(z(x-y),xy(x-y),z^2-x^2y).        \tag{L}
\] A Hilbert--Burch matrix is \[
 \begin{pmatrix}
 xy&z\\
 -z&-x\\
 0&-(x-y)
 \end{pmatrix}.
\] Its maximal minors give (L), with alternating signs. The ideal has codimension two and is Cohen--Macaulay. Its zero set is exactly the cusp and the two axes, and it is generically reduced on each branch. A nonzero nilpotent quotient supported at the origin would violate depth one, so it is the actual reduced intersection ideal.

On the cusp use the primitive conormal basis \[
                     k_1=dx-dy,\qquad k_2=3t\,dx-2dz.
\] The differentiated generators are \(t^3k_1,t^4k_1,t^3(tk_1-k_2)\). Thus its elementary divisors are \((3,3)\). On each smooth axis they are \((1,2)\). Here \(r=1\) on the cusp and zero on the axes. By (A), \[
                  D\ge4+\left(\frac92-1\right)+2+2
                       =\frac{23}{2}>9.               \tag{M}
\]

This exhausts multiplicity four. Together with Sections 1--2 it proves the strict local bound for every nonplanar three-branch curve with delta four. Only actual coordinate isomorphisms, conormal maps, and the rank-one Hessian were used; no boundary deformation has been introduced.

\par\medskip\noindent\textbf{4. Every nonordinary triple exceeds the budget}\par\hypertarget{tag-01OO}{}\label{every-nonordinary-triple-exceeds-the-budget}

For a planar three-branch point, P4.48 gives defect at least ten, including the ordinary plane triple point. At a nonplanar point with delta at least five, (B) and \(m\ge3\) give \[
                  D\ge5+3+\frac32=\frac{19}{2}>9.
\] Delta four was exhausted above.

At delta three the excess \(\delta-(r-1)\) is one. The actual analytic ring calculation in P4.40 gives precisely an ordinary plane triple, a plane tacnode with an independent transverse line, or a cusp with two independent lines. The first is excluded by P4.48. The second has defect at least eleven by P4R.16. The third has embedding dimension four and cannot occur in the ambient smooth threefold.

Finally a three-branch point with delta two has zero excess. Its ring equals \(\mathbb C_{\rm diag}+F^1\), generated by the independent branch parameters. It is exactly the three coordinate axes, with ideal \((xy,xz,yz)\). Thus every nonordinary three-branch point is excluded from the total defect budget nine. The only remaining local possibility is the ordinary nonplanar triple treated by P4.48 and the separate full-quartic conductor calculation.


\end{enumerate}
\end{proof}

The analytic model identities, primitive conormal lattices,
Hilbert--Burch minors, cubic obstruction, and jet generators are
checked exactly in \path{code/part4/a3_delta4.py}.


\subsection{The horizontal image in the unrestricted ruled remainder}

\begin{proposition}[P4R.18]\hypertarget{tag-02YB}{}\label{night:P4R.18}\label{stmt:P4R.18}
\StatementRecord{night:P4R.18}{proved}{02YB}{\SeriesRef{night:P4.33}, \SeriesRef{night:P4.42}, \SeriesRef{night:P4R.10}, \SeriesRef{night:P4R.13}, \SeriesRef{night:P4R.14}, \SeriesRef{night:P4R.15}, \SeriesRef{night:P4R.6}}

In the remaining positive-Euler ruled case of Part~4.2, a
noncontracted horizontal pole section maps with degree two or
three to a rational normalized conductor component whose generic
transverse germ is $D_4$. This assertion imposes no transverse
type or conductor-exponent restriction on the other image component.
\end{proposition}
\begin{proof}

\begin{enumerate}[label={$\langle1\rangle$\arabic*.},leftmargin=*]
\item Write $H$ for the horizontal section, $g=g(H)\ge1$, and
$C_H$ for its image in the reduced conductor $C$. The preceding
ruled reduction gives two rational normalized components of $C$
and one point-identification cycle. The nonconstant map
\[
                         H\longrightarrow C_H^\nu=\mathbb P^1
\]
has degree at least two: a degree-one finite map between smooth
curves is an isomorphism, whereas their genera differ. Each
generic point of $H$ has conductor exponent two, by the actual
pole divisor of P4R.6. Thus the generic transverse plane curve
along $C_H$ has at least two normalization branches of exponent two.

\item For its branches $B_i$ the exact conductor formula is
\[
             c_i=2\delta(B_i)+\sum_{j\ne i}I(B_i,B_j).
\]
A singular branch with $c_i=2$ permits no other branch,
so the two specified branches are smooth. With two total
branches their intersection multiplicity is two, giving $A_3$.
With three total branches, the first two exponents force all
three pairwise intersection multiplicities to equal one;
the third branch is then smooth as well, giving $D_4$.
Four branches would give exponent at least three on every
smooth branch. These possibilities are exhaustive. In type
$A_3$, $H$ has degree exactly two and supplies both generic
normalization preimages. In type $D_4$ its degree is two or
three. We exclude the former transverse type next.

\item Assume that the transverse germ along $C_H$ is $A_3$.
Let $C_{H,\rm sch}$ be the primary component of the actual
image conductor scheme supported on $C_H$. The primary
restriction construction of P4R.13 applies before its
elliptic-genus calculation, without change. More explicitly,
P4.33 gives rational singularities of the normalization when
$H$ is not contracted, so finite duality and the canonical
trace give the injections
\[
 \mathcal O_{C_{\rm sch}}\longrightarrow
 \nu_*(\mathcal O_N/I_N)\longrightarrow\rho_*\mathcal O_D.
\]
The exact blowup relation $D'=\beta^*D''-E$ identifies
global functions on $D$ with those on $2H_0+V_0$ of the
minimal ruled model. The kernel of restriction to $2H_0$
has successive vertical-fibre layers
$\mathcal O_{\mathbb P^1}(-2)$, with no global sections.
Hence a global function is determined at the generic doubled
horizontal section. It follows that restriction from
$C_{\rm sch}$ to $C_{H,\rm sch}$ is injective on $H^0$
and surjective on $H^1$: the first assertion follows from
those injections and generic determination, and the second
from the dimension-one kernel of the restriction sheaf map.
Therefore
\[
                 g-1=\chi(\mathcal O_{C_{\rm sch}})
                    \le\chi(\mathcal O_{C_{H,\rm sch}}).
                                                            \tag{P4R.18.1}
\]
Every arrow here is the primary restriction of the same
conductor algebra; no other conductor component is deleted
from the left side.

\item We evaluate the right side on $C_H^\nu=\mathbb P^1$.
Use the actual differentiated conormal, the saturated Hessian
line and the rank-one nilradical of P4R.10. The kernel $K$
of the normalized derivative has degree $2+\deg R$, since
$c_1(Y)\cdot C_H=0$. The Hessian line has degree
$\tfrac12\sum d_p$. Riemann--Roch and the normalization
sequence give exactly
\[
 \chi(\mathcal O_{C_{H,\rm sch}})
 =4-\sum_p\left[\delta(C_H,p)
      +\sum_{i\mapsto p}(d_i/2+b_i-r_i)+q_p\right].
                                                            \tag{P4R.18.2}
\]
All terms are those of this primary curve and its primary
conductor. They remain defined at its intersections with
the other component, because the pure conductor quotient
extends its generic containment of the symbolic square.

\item Every singular point of $C_H$ contributes at least three
to the sum in (P4R.18.2). At a planar point the ideal is
$(u,h)$, its symbolic square equals its ordinary square,
and the actual equation belongs to that square. The
rank-one symmetric-square estimate of P4.42 and the exact
plane-curve different calculation give a contribution at
least $3\delta(C_H,p)+q_p\ge3$.
At a nonplanar point every defining equation has order at
least two. On a branch of multiplicity $m_i$ its projected
conormal order is at least $m_i$. The point is singular
on the entire critical curve, so P4R.15 makes the ambient
Hessian zero and gives $d_i\ge m_i$. Since $r_i=m_i-1$,
the contribution is at least
\[
                      \delta(C_H,p)+\sum_i(1+m_i/2),
\]
which is at least three at a singular curve point. The
same arguments give nonnegative contributions everywhere;
at a smooth primary point its conormal is the full kernel
and its only possible contribution is the nonnegative
Hessian zero order divided by two.
Combining (P4R.18.1)--(P4R.18.2) gives total primary
defect at most $5-g\le4$. Thus $C_H$ has at most one
singular point, and two different singular points are
impossible even in the presence of arbitrary transverse
types on the other global component.

\item We count the possible branch values of the degree-two map
$H\to\mathbb P^1$ with the identifications retained.
The normalization graph formula is
\[
                         \sum_{p\in C}(r_p-1)=2.
\]
Connectedness leaves either two two-branch points or one
three-branch point. In the two-point case, two intersections
of the global components give two points on $C_H^\nu$.
There can be at most one additional unibranch singularity
of $C_H$, so these supply at most three possible branch
values. The other connected possibility is a self-identification
on one component and one intersection of the two components.
A self-identification on $C_H$ supplies two normalization
points and already costs at least three in the primary
defect. It permits no additional singular point, so with
the intersection there are again at most three possible
branch values. A self-identification on the other component
supplies no additional point on $C_H^\nu$.

\item At a three-branch point, the two global components must
supply two and one branches. Thus $C_H^\nu$ has at most
two points over it. Two points already make $C_H$ singular
there and exclude any further singular point by the defect
bound. One point permits at most one further unibranch
singularity. This case also has at most three possible
branch values.

\item There is no ramification elsewhere. At a point where the
entire critical curve is smooth, P4R.14 gives precisely
the product tacnode or $X(Y^2+XZ)$. Both have two distinct
normalization points over that point of the conductor.
Their two generic preimages lie on $H$, so properness
places both special preimages on $H$ as well. The
degree-two map is therefore unramified there. We have
accounted for every possible ramification value.
Riemann--Hurwitz for a degree-two map from a genus-$g$
curve requires $2g+2\ge4$ distinct branch values; each
ramified fibre consists of its single point of index two.
This contradicts the count of at most three. Thus the
generic germ along $C_H$ must be $D_4$, with degree two
or three as proved at the start.

\end{enumerate}
\end{proof}

The exact connected identification patterns and their primary
defect/branch-value counts are enumerated in
\texttt{code/part4/horizontal\_primary.py}.


\subsection{A triple spectral root with nonzero cubic tensor}

\begin{proposition}[P4R.19]\hypertarget{tag-02YC}{}\label{night:P4R.19}\label{stmt:P4R.19}
\StatementRecord{night:P4R.19}{proved}{02YC}{\SeriesRef{night:P4.51}, \SeriesRef{night:P4.52}}

Let $f$ be the actual reduced local equation of the remaining
positive ruled support, and let $C_0$ be a smooth local branch
of its critical curve with generic transverse type $D_4$.
At a point where the cubic normal tensor along $C_0$ is
nonzero, that tensor is not the cube of a linear form.
Other critical-curve branches at the point are retained,
without a restriction on their generic transverse types.
\end{proposition}
\begin{proof}

\begin{enumerate}[label={$\langle1\rangle$\arabic*.},leftmargin=*]
\item Choose coordinates $(x,y,z)$ with $C_0=(x=y=0)$.
Generic normal multiplicity three gives
$f\in(x,y)^{(3)}$. Powers of this smooth
complete-intersection ideal are primary, so
$f\in(x,y)^3$ at the specified point.

\item Suppose its nonzero cubic normal tensor is a cube.
A linear normal-coordinate change makes it $x^3$, up
to a nonzero scalar. Weierstrass preparation in $x$
gives the exact equality
\[
 f=U(x,y,z)\bigl[x^3+a_2(y,z)x^2+a_1(y,z)x+a_0(y,z)\bigr],
\]
where $U$ is a holomorphic unit. The bracket belongs to
$(x,y)^3$, so $a_2,a_1,a_0$ are divisible by $y,y^2,y^3$,
respectively. Depressing the monic cubic is a holomorphic
coordinate change preserving $(x,y)$ and yields
\[
 f=U(x,y,z)\bigl[x^3+y^2B(y,z)x+y^3C(y,z)\bigr],
 \qquad B(0,0)=C(0,0)=0.                         \tag{R19.1}
\]
The unit remains in this equality and in the scaling
comparison below.

\item Set
\[
 k=\operatorname{ord}_z B(0,z),\qquad
 l=\operatorname{ord}_z C(0,z),\qquad
 \rho=\min(k/2,l/3),
\]
where the order of an identically zero series is infinity.
At least one order is finite: otherwise all generic normal
cubics would be $x^3$, contrary to generic $D_4$.
Thus $0<\rho<\infty$. Choose $L>3\rho$ and positive
rational weights
\[
       \operatorname{wt}(x)=L+\rho,\qquad
       \operatorname{wt}(y)=L,\qquad
       \operatorname{wt}(z)=1.
\]
Multiplication by a common denominator gives positive
integer weights. The weighted degree of the first face
is $3L+3\rho$. Terms of $B$ containing $y$ have weight
gap at least $L-2\rho>0$, and those of $C$ containing
$y$ have gap at least $L-3\rho>0$. Consequently the
leading polynomial, apart from the nonzero factor $U(0)$,
is exactly one of
\[
\begin{array}{ll}
 P=x(x^2+a y^2z^k),\quad a\ne0,
       & k/2<l/3,\\[2pt]
 P=x^3+b y^3z^l,\quad b\ne0,
       & l/3<k/2,\\[2pt]
 P=Q(x,yz^r),\quad
 Q(X,V)=X^3+aXV^2+bV^3,\quad a,b\ne0,
       & k=2r, l=3r, r\ge1.
\end{array}                                                   \tag{R19.2}
\]
In the last case $r$ is an integer. The depressed cubic
$Q$ is not a cube: a cube with $X^3$ coefficient one and
$X^2V$ coefficient zero would have $a=b=0$. A double
factor of $Q$ is allowed throughout the argument.

\par\medskip\noindent\textbf{The compact class used in the comparison.}
Let $B$ be a connected smooth complex curve and let
$h:B\to\mathbb C$ have two distinct simple zeros $p,q$.
Join them by a smooth embedded arc $\gamma$ whose
interior contains no zero. On the smooth surface
\[
             Z_h=\{(b,u,v)\in B\times\mathbb C^2:
                                      uv=h(b)\},
\]
consider the locus over $\gamma$ with $|u|=|v|$.
It is a circle over each interior point and collapses
to a point at each end. Near an endpoint choose the
arc so that $h$ is a fixed nonzero complex multiple
of a positive real parameter $s$. Then
$u=\sqrt{s}\,e^{\mathrm i\phi}$ and
$v=c\sqrt{s}\,e^{-\mathrm i\phi}$ are the smooth
polar coordinates of a real disc in the local $uv$
chart. The locus is thus a smoothly embedded sphere.

It is Lagrangian for a product K\"ahler form.
Writing $u=R e^{\mathrm i\phi}$,
$v=R e^{\mathrm i(\psi-\phi)}$, the sum of the two
fibre area forms restricts to
$d(R^2/2)\wedge d\psi=0$, since $R,\psi$ depend
only on the arc parameter. The base area form also
restricts to zero. Multiplication by $\mathrm i$
identifies tangent and normal bundles, with the
opposite normal orientation relative to the complex
ambient orientation. Its self-intersection is
therefore $-\chi(S^2)=-2$.
For the compactly supported Poincar\'e dual $\alpha$,
\[
                         \int_{Z_h}\alpha^2=-2.             \tag{R19.3}
\]

If $\pi:W\to Z_h$ is a finite proper holomorphic map
of degree $d>0$ between smooth complex surfaces, then
$\pi^*\alpha$ is compactly supported and
\[
              \int_W(\pi^*\alpha)^2
                       =d\int_{Z_h}\alpha^2=-2d.            \tag{R19.4}
\]
The equality follows from the proper oriented degree;
off the analytic branch locus the map has $d$ positive
sheets, and the branch locus has measure zero for
integration of compactly supported top forms.
Poincar\'e duality gives a compact homology two-class
with that square. Its representative need not itself
be a sphere.

\par\medskip\noindent\textbf{The exact finite maps.}
In the first case of (R19.2), $P=1$ forces $x\ne0$
and gives
\[
           y^2z^k=h(x),\qquad
           h(x)=\frac{1-x^3}{ax},\qquad x\in\mathbb C^*.
\]
The map
\[
 (x,y,z)\longmapsto(x,u=y^2,v=z^k)
\]
is finite proper of degree $2k$ onto $uv=h(x)$.
The target is smooth: at $u=v=0$ the zeros of $h$
are simple, and elsewhere a derivative in $u$ or
$v$ is nonzero. The source is smooth by the weighted
Euler identity on $P=1$. Join two of the three cube
roots of one by an arc in $\mathbb C^*$ avoiding the
third. Equations (R19.3)--(R19.4) give a compact
two-class of square $-4k$ on the leading fibre.

In the second case the exact finite proper map is
\[
 \begin{split}
 \{x^3+b y^3z^l=1\}&\longrightarrow
                 \{uv=(1-x^3)/b\},\\
 (x,y,z)&\longmapsto(x,u=y^3,v=z^l).
 \end{split}
\]
It has degree $3l$, and the base is $\mathbb C$.
The same three simple zeros supply a pulled-back
compact class of square $-6l$.

For the balanced case let
\[
                 B_Q=\{(X,V)\in\mathbb C^2:Q(X,V)=1\}.
\]
It is smooth by the ordinary Euler identity and is
connected even when $Q$ has a double factor. Its
projection to $\mathbb P^1$ minus the distinct zero
directions of $Q$ is an unramified cyclic cover of
degree three. Around a zero direction of multiplicity
one or two its monodromy generates $\mathbb Z/3$.
Such a direction exists because $Q$ is not a cube,
so the cover is connected.

The function $V$ on $B_Q$ has precisely the three
zeros $(\zeta,0)$ with $\zeta^3=1$, and they are
simple because $Q_X(\zeta,0)=3\zeta^2\ne0$.
The surface
\[
             Z_Q=\{(b,u,w):b\in B_Q,\ uw=V(b)\}
\]
is smooth. The exact map
\[
 \begin{split}
       \{Q(x,yz^r)=1\}&\longrightarrow Z_Q,\\
       (x,y,z)&\longmapsto
                    \bigl(b=(x,yz^r),u=y,w=z^r\bigr)
 \end{split}                                                   \tag{R19.5}
\]
is finite proper of degree $r$: $z$ satisfies the
monic equation $z^r=w$ and the other coordinates are
retained. An arc in $B_Q$ joining two $V$-zeros and
avoiding the third supplies a compact class of square
$-2r$ on the source. In particular this argument
retains $Q=(X-\lambda V)^2(X+2\lambda V)$,
$\lambda\ne0$, without a perturbation of its
repeated factor.

\par\medskip\noindent\textbf{Transport to the actual local fibre.}
Let $(w_x,w_y,w_z)$ be the positive integer weights
and let $d$ be the resulting degree. Put
\[
 D_\epsilon(x,y,z)=
       (\epsilon^{w_x}x,\epsilon^{w_y}y,\epsilon^{w_z}z),
 \qquad
 G_\epsilon=\epsilon^{-d}f\circ D_\epsilon.
\]
These actual functions converge with derivatives on
compact sets to $U(0)P$. The weighted Euler identity
makes $dP$ nonzero on $P=1$, including the repeated-factor
case on its zero set. Choose compact representatives
of the two-class above and of the transverse
perturbation computing its square. On a compact
neighbourhood of these chains, the gradient-flow
comparison of P4.52 solves
\[
       dG_\epsilon(V_\epsilon)
                         =-\partial_\epsilon G_\epsilon.
\]
The denominator is bounded away from zero on a
slightly larger neighbourhood; a cutoff gives a
flow for sufficiently small positive $\epsilon$.
It transports the chains and their intersections
to $G_\epsilon=U(0)$. Scaling by $D_\epsilon$ then
embeds them in the original fixed Milnor ball on
\[
                         f=U(0)\epsilon^d.
\]
The oriented intersection number is preserved and
equals $-4k$, $-6l$, or $-2r$. Each is nonzero,
contradicting P4.51. This proves the assertion.

\end{enumerate}
\end{proof}

\par\medskip\noindent\textbf{The two-direction exception is retained.}
The assertion cannot be enlarged to every spectral
collision. The reduced equation
\[
                       f=y(x^2+y^2z^2)
\]
has generic $D_4$ along $x=y=0$ and generic $A_1$
along $x=z=0$. Only two spectral directions collide
at their intersection. Its non-plane normalization
sheets $x=\pm\mathrm i yz$ have pole divisor
$2(y=0)+(z=0)$, with simple normal crossings.
Its affine fibre $f=1$ is exactly $(\mathbb C^*)^2$:
\[
 \begin{split}
       a&=x+\mathrm i yz,\qquad b=x-\mathrm i yz,\\
       x&=(a+b)/2,\qquad y=(ab)^{-1},\qquad
       z=(a-b)ab/(2\mathrm i).
 \end{split}
\]
The displayed maps are inverse and substitution gives
$f=1$. Its intersection form is zero: the product
of two unit circles generates its second homology,
and a radial translation gives a disjoint homologous
copy. More generally the fibre of $y(x^2+y^2z^n)$
becomes $u^2+z^n=t^3$, $t\ne0$, through
$t=1/y$, $u=x/y$. Positive radial dilation gives
a product real direction, along which any compact
cycle can be displaced from another compact cycle.
Its intersection form is also zero. Thus the
other generic transverse type remains unrestricted.

The exact weighted gaps, all three finite-map
equations, their degree-square identities, and the
inverse maps of the two-direction exception are
executed in \path{code/part4/d4_spectral_triple.py};
its \path{.json} and \path{.out} outputs report PASS.


\subsection{The three-branch image point with a double horizontal cover}

\begin{proposition}[P4R.20]\hypertarget{tag-02YE}{}\label{night:P4R.20}\label{stmt:P4R.20}
\StatementRecord{night:P4R.20}{proved}{02YE}{\SeriesRef{night:P4.55}}

In the remaining degree-two horizontal sector after P4.55,
suppose that the generic transverse curve along the other
image component $B$ has at least three normalization branches.
The unique image-graph cycle cannot be concentrated at one
point with three local critical-curve branches, namely the
smooth branch $A$ and two local branches of $B$.
\end{proposition}
\begin{proof}

\begin{enumerate}[label={$\langle1\rangle$\arabic*.},leftmargin=*]
\item Write these two local branches as $B_+,B_-$ and the
point as $p$. Let $W_1,\ldots,W_c$ be the normalized
conductor components mapping onto $B$, with degrees
$d_i>0$ over $B^\nu$. Put
\[
                         r_B=\sum_i d_i\ge3.
\]
This is the number of generic normalization branches
of the full transverse plane curve along $B$.
The normalization conductor has components
$H,R,W_1,\ldots,W_c$ and a tree graph. Its total
point-identification contribution is $c+1$.
The horizontal section $H$ maps with degree two to
$A$, while the rational vertical curve $R$ maps
with degree one. Distinct points of $H$ are not
identified in the normalization $N$. Distinct
points of $R$ are not identified either: a
contracted vertical tree effecting such an
identification would create a cycle in the ruling
fibre. All conductor branches below are counted
on their normalization.

\item We first record the exact spectral bound used in
the proof. Let $E=N_{A/Y}^*$, so $\deg E=2$ by
adjunction and $c_1(Y)\cdot A=0$. On $\mathbb P(E)$
use the quotient convention, tautological class
$\xi$, and fibre class $F$. Thus
\[
          \xi^2=2,\qquad \xi F=1,\qquad F^2=0,
          \qquad K_{\mathbb P(E)}=-2\xi.
\]
The actual normal derivative maps of P4.55 identify
$H$ with the normalization of the degree-two
spectral component $\Gamma_H$, and identify the
third sheet with a section $\Gamma_R\cong A$.
Remove the common zero divisor of the cubic tensor,
of degree $b\ge0$. For an integer $t$ write
\[
       \Gamma_R=\xi-tF,\qquad
       \Gamma_H=2\xi+(t-b)F.
\]
If $\delta_H$ is the normalization colength of
$\Gamma_H$ and $n=\Gamma_H\cdot\Gamma_R\ge0$,
adjunction and the intersection product give
\[
 \begin{split}
       g(H)+\delta_H&=p_a(\Gamma_H)=1+t-b,\\
       n&=4-t-b,\\
       g(H)+\delta_H+n&=5-2b.
 \end{split}                                                   \tag{R20.1}
\]
These are the actual residual components of the
cubic, with their common zero divisor retained.
Since $g(H)\ge1$, they imply $b\le2$.

\par\medskip\noindent\textbf{The exact graph when $c\ge2$.}
Let $r_A$ be the number of normalized conductor
points over the branch $A$ at $p$. It is two if
$H\to A$ ramifies there, and three otherwise.
Let $r_+,r_-$ count the points over $B_+,B_-$.
Every $W_i$ is surjective to $B^\nu$, so
$r_+,r_-\ge c$. Let $n_p$ be the number of
normalization points of $S$ over $p$.
Their nonempty supports in $\{A,+,-\}$ pairwise
intersect: distinct local hypersurface components
intersect in a curve.

\item A singleton support $\{A\}$ forces every point to
carry $A$. Then $n_p\le r_A$ and the graph
contribution is at least $r_++r_-\ge2c>c+1$.
A singleton $\{+\}$ forces every point to carry
$+$ and gives contribution at least
$r_A+r_-\ge c+2$; the same applies to $\{-\}$.
Thus every support has at least two elements, and
\[
 \text{contribution at }p
       \ge\frac{r_A+r_++r_-}{2}\ge c+1.
\]
The whole global budget is $c+1$. Equality is
forced throughout:
\[
                 r_A=2,\qquad r_+=r_-=c,\qquad n_p=c+1.
                                                               \tag{R20.2}
\]
Every point over $p$ has exactly two normalized
conductor branches, with different image labels,
and all graph identifications occur at $p$.
The horizontal cover ramifies over $p$.
Each $W_i$ has precisely one point over each of
$B_+,B_-$ and has valence two; $H$ and $R$ have
valence one. The connected tree is therefore
the path
\[
                    H-W_1-W_2-\cdots-W_c-R.             \tag{R20.3}
\]
Choose the label $+$ at the $H$ endpoint. The next
point on $W_1$ has label $-$. Every internal point
has support $\{+,-\}$, so it joins that point to
the $+$ point of $W_2$. Induction gives the $+$
point on the left of every $W_i$ and its $-$
point on the right. The $R$ endpoint carries $-$.

\item The endpoint hypersurface factors, denoted
$f_H,f_R$, thus have supports $\{A,+\}$ and
$\{A,-\}$. The $c-1$ internal factors $f_j$
have supports $\{+,-\}$. These are all the local
irreducible hypersurface factors.

\par\medskip\noindent\textbf{The actual length-two ribbon.}
The set-theoretic intersection of $f_H$ and $f_R$
is $A$ alone. Along generic $A$, the first factor
supplies two ordinary-triple branches and the
second supplies the third, so their intersection
multiplicity is two. Their one-dimensional
complete-intersection ring is Cohen--Macaulay,
without a zero-dimensional embedded component.
A coordinate $z$ on the smooth $A$ is a
nonzerodivisor. The ring is finite free of rank
two over $\mathbb C\{z\}$, and its quotient by
$z$ has length two.

\item If both factors had no linear term, their
restrictions modulo $z$ would belong to
$(x,y)^2$. The Artinian quotient would contain
the independent classes $1,x,y$, contrary to
length two. At least one factor is smooth.
It cannot be $f_H$: a smooth surface is its own
normalization and supplies only one generic
preimage of $A$, whereas this factor supplies
the ramified double cover. Thus $f_R$ is smooth.
Take the actual ambient coordinate $u=f_R$.
On $u=0$, the divisor of $f_H$ is exactly twice
the smooth $A$, so in coordinates
$A=(u=v=0)$ it has restriction $v^2$ times a
unit. Weierstrass preparation in $v$, followed
by depression and using $f_H\in(u,v)^2$, gives
the exact expression
\[
                    f_H=\text{unit}\,(v^2+u^2g(u,z)).     \tag{R20.4}
\]
All changes preserve $A$. The full cubic normal
tensor is therefore, up to a nonvanishing scalar,
\[
        \left(\prod_j f_j(0,0,z)\right)
                     u\bigl(v^2+u^2g(0,z)\bigr).          \tag{R20.5}
\]
The coefficient of $uv^2$ in the last factor is
one, so the common zero order is exactly the
sum of the orders of $f_j|_A$.

\par\medskip\noindent\textbf{The pole on the smooth rational branch.}
On the smooth branch $f_R=0$, its conductor is
the restriction of the product of the other
factors. The factor $f_H$ contributes $2A$.
Each $f_j$ cuts out only $B_-$ there. This
normalization branch is already smooth, and
the pole support on its minimal resolution is
SNC. Hence $A$ and its $B_-$ curve are smooth
and transverse. Choose the latter's equation
$z=0$ on this surface, with $z$ still a
coordinate on $A$, and extend it to the ambient
threefold. Then
\[
             f_j|_{f_R=0}=\text{unit}\,z^{e_j},
                  \qquad e_j\ge1.
\]
Their sum $e=\sum e_j$ is the generic conductor
exponent of the global component $W_c$ on
this branch. By (R20.5) it is also the local
common cubic-zero order. The branch on
$f_R=0$ is smooth; the plane-curve conductor
formula therefore gives
\[
          e=\sum_{\text{other branches}}I
                         \ge r_B-1\ge2.
\]
Combining $b\ge e$ with (R20.1) forces
\[
     b=e=2,\quad r_B=3,\quad
     g(H)=1,\quad \delta_H=n=0.                         \tag{R20.6}
\]
Thus $\Gamma_H$ is smooth and disjoint from
$\Gamma_R$, and there is no further common
zero of the cubic.

\par\medskip\noindent\textbf{The exact normalization at the ramified endpoint.}
The cover $H\to A$ ramifies at $p$ by (R20.2).
Equation (R20.4) presents its spectral cover as
$w^2+g(0,z)=0$. After translating its double
root, smoothness from (R20.6) gives
$g_z(0,0)\ne0$. Hence
\[
 \begin{split}
        N_H&=\{w^2+g(u,z)=0\},\\
        (u,w,z)&\longmapsto(u,v=uw,z)
 \end{split}                                                   \tag{R20.7}
\]
is the actual normalization near that point.
It is smooth, finite, and birational onto the
branch (R20.4); off $u=0$ its inverse is
$w=v/u$. Coordinates $(u,w)$ solve
$z=\psi(u,w)$, with $z|_H$ of order two in
$w$. The minimal normalization resolution
has no exceptional curve over this smooth
point, and $H$ is $u=0$.

\item The other conductor branch $W_1$ meets $H$
transversely by SNC. Thus $u$ is a coordinate
on $W_1$ and has order one on its image
$B_+$. By (R20.2), $W_1$ has only one point
over the normalized $B_+$ point, so its
global degree is $d_1=1$. The same smooth
SNC argument on $f_R=0$ gives $d_c=1$.
Since $\sum d_i=r_B=3$, the case $c=2$
would have sum two, while $c>3$ is impossible.
Therefore
\[
                 c=3,\qquad d_1=d_2=d_3=1.              \tag{R20.8}
\]
There are exactly two internal factors, and
$e_1=e_2=1$.

\par\medskip\noindent\textbf{The conductor mismatch.}
Each $f_j$ intersects $f_H$ only along $B_+$.
Its pullback to the actual smooth normalization
(R20.7) consequently has divisor supported
only on $W_1$. Its restriction to $H$ vanishes
to order two: the order on $A$ is one and
$H\to A$ has ramification index two.
Transversality of $H$ and $W_1$ makes the
divisor exactly $2W_1$. Thus the generic
intersection multiplicity of $f_H$ with
each internal factor along $B_+$ is two.

\item The path (R20.3) identifies the same global
component $W_3$ on the smooth $R$ branch
over $B_-$ and on an internal factor over
$B_+$. At $B_-$ its conductor exponent is
$e=2$. At $B_+$ it has intersection
multiplicity two with $f_H$ and positive
intersection multiplicity with the other
internal factor. Its conductor exponent
is at least three. These cannot be two
values of the coefficient on one global
irreducible conductor component.

\item The full marking is also explicit. Over
$B_+$ the factors $(f_H,f_1,f_2)$ carry
$(W_1,W_2,W_3)$; over $B_-$ the factors
$(f_1,f_2,f_R)$ carry the same ordered
global components. If $p_+,p_-\ge1$ are
the intersections of the internal factors,
the conductor vectors are
\[
         (4,2+p_+,2+p_+),\qquad
         (1+p_-,1+p_-,2).
\]
There is no missing intrinsic branch-delta
term: the exponent two of the smooth $R$
branch forces its two intersections to
equal one, which makes both other generic
$B$ branches smooth as well. The third
entries already give the contradiction.

\par\medskip\noindent\textbf{The case $c=1$.}
Now the conductor graph has $H,R,W$ and
budget two. The points of $W$ over $B_+$
and $B_-$ cannot be identified in $N$,
since that creates a loop in the conductor
tree. Points carrying the two labels are
distinct; their supports intersect, so
both carry $A$. Each contributes at least
one, exhausting the graph budget. Each
has exactly one $A$ and one $B$ branch,
and no other identification occurs.

\item If $H$ ramifies over $p$, these two
$A$ branches are $H$ and $R$. There are
only the factors with supports
$\{A,+\}$ and $\{A,-\}$. The length-two
ribbon argument makes $f_R$ smooth.
No other factor contains $B_-$, and this
smooth branch has no singularity there.
Thus $B_-$ cannot be a critical curve.

\item If $H$ is unramified, the three normalized
$A$ preimages are distinct. Two occur
at the $B$ attachments and the third
is an $A$-only normalization point.
There are three factors, each supplying
one generic $D_4$ branch. Every pair
meets only along $A$, with multiplicity
one. Its complete-intersection curve
is the reduced smooth $A$: it is pure,
generically reduced, and has no other
component. Modulo $z$, length one forces
independent normal linear terms.
All three factors are smooth near $p$
and their only pairwise intersection is
$A$. Neither $B_+$ nor $B_-$ can be a
critical branch. This final contradiction
completes the proof.

\end{enumerate}
\end{proof}

The exact spectral intersections, length-two
ribbon basis, normalization equation, scalar
cubic coefficient, endpoint-degree list and
marked conductor vectors are executed in
\path{code/part4/degree_two_triple_point.py}.
Its \path{.json} and \path{.out} certificates
report PASS.

\section{Global CDP replacement and normalization resonance}\hypertarget{tag-0348}{}\label{sec:II:12}


\begin{proposition}[A geometric replacement on an explicit fibre class]\hypertarget{tag-01FK}{}\label{cdp:st:C4a}
\StatementRecord{cdp:st:C4a}{proved}{01FK}{\SeriesRef{cdp:st:A.4.5}, \SeriesRef{cdp:st:A.1.2}, \SeriesRef{cdp:st:A.4.6}, \SeriesRef{cdp:st:C2a}}

Let $f:Y\to\PP^1$ be a connected-fibre holomorphic map from a compact connected complex threefold with $b_1(Y)=b_2(Y)=0$, and smooth general fibre a complex two-torus. Suppose every full fibre has one of the following forms: a reduced irreducible normal-crossing surface; or a smooth torus family modulo a finite free cyclic action over a cyclic base cover, so that its reduction is smooth and its multiplicity is the covering order. Let
\[
 0\ne\gamma\in H^0(\PP^1,R^1f_*\Z)
\]
be a generator, and suppose that, for every singular reduced full fibre $S$, its pullback to $H^1(\widetilde S,\Q)$ under normalization is nonzero. Then a dense open subset of every component of $\Pic(Y)$ satisfies
\[
 H^q(Y,T_Y\otimes L)=0\quad(0\le q\le3),\qquad
 \chi(Y,T_Y\otimes L)=0,\qquad c_3(Y)=0.                \tag{C.15}\hypertarget{tag-008T}{}\label{cdp:eq:C.15}
\]
The generator exists under the stated rational Betti-number assumptions. Its normalization condition is a condition on an explicit integral cohomology class, with no tangent-cohomology vanishing assumed.
\end{proposition}


\begin{proof}

\begin{enumerate}[label={$\langle1\rangle$\arabic*.},leftmargin=*]
\item The integral Leray sequence gives an injection
\[
H^0(\PP^1,R^1f_*\Z)\longrightarrow H^2(\PP^1,\Z)=\Z
\]
whose cokernel injects into the finite group $H^2(Y,\Z)$. Thus its source is infinite cyclic. The degree-one specialization argument of Part D identifies its stalk at every fibre with a nonzero integral class: a class in the preimage of a small disc that vanishes on a nearby fibre comes from the meridian of the punctured disc; a normal disc to any irreducible component makes a nonzero multiple of that meridian nullhomotopic, forcing the class to vanish. The real-to-holomorphic cohomology map on a compact connected complex space is injective on real classes, as proved there by the pluriharmonic maximum principle. Coherent Leray therefore supplies $\beta\in H^1(Y,\OO_Y)$ whose restriction to every full fibre is the image of $\gamma$.

\item Fix a real irrational number $r$ and put $M_r=\exp(r\beta)$, with the exponential normalized by $2\pi i$. On each fibre this line is represented by the unitary character $\exp(2\pi ir\gamma)$. It is nontrivial, and on a singular fibre its normalization pullback is nontrivial by the hypothesis on $\gamma$. To see this directly, triviality of an $n$th tensor power would place $nr\gamma$ in the integral lattice after applying $H^1(\R)\to H^1(\OO)$; real injectivity makes this an equality in $H^1(\R)$, which is impossible for an irrational multiple of a nonzero integral vector.

\item The same construction works in every Picard component. Each component contains a finite-order line bundle $P$: for a representative $P_0$, a positive power $P_0^m$ lies in the complex vector group $\Pic^0(Y)$, where division by $m$ is possible. Dividing $P_0$ by the resulting line makes its $m$th power trivial. Replacing $M_r$ by $M_r\otimes P$ retains nontrivial unitary restriction on each fibre and on each specified normalization. Indeed an $m$th power eliminates $P$ and reduces a purported trivialization to the preceding irrationality contradiction.

\item Set $L=K_Y\otimes(M_r\otimes P)^{-1}$. On every full fibre $F$ the Serre-dual bundle for degree two is
\[
 (T_Y\otimes L)|_F^*\otimes\omega_F
       =\Omega_Y^1|_F\otimes(M_r\otimes P)|_F,
\]
because $F$ is a Cartier fibre and $\omega_F=K_Y|_F$. At a smooth torus fibre both the conormal line and the intrinsic cotangent bundle are trivial; nontrivial unitary degree-zero lines on the torus have zero cohomology. The ambient space of sections is therefore zero. At a reduced irreducible normal-crossing fibre, Statement \ref{cdp:st:A.4.5} proves the torsion-free cotangent vanishing for every nontrivial unitary line. The normal-crossing kernel in \ref{cdp:st:A.1.2} is the normalization pushforward, and its sections vanish because the pulled-back unitary line is nontrivial. A holomorphic section of a unitary line on a compact connected complex manifold has constant norm by the maximum principle and is either zero or a trivialization. Thus the ambient group vanishes there as well.

\item On a cyclic-cover neighbourhood of a smooth multiple fibre, the covering map to $Y$ is finite etale. The pulled-back line restricts nontrivially to the covering central torus: finite etale pullback is injective on rational degree-one cohomology by its trace, and $r$ is irrational. The scheme-theoretic inverse image of the full fibre is the corresponding thickening of this torus. Each layer in its ideal filtration is the same ambient cotangent restriction tensored with this nontrivial line, since the covering base coordinate trivializes the conormal. Every layer has zero cohomology. The exact filtration sequences therefore give zero sections on the full thickening. Taking finite-group invariants, an exact functor, proves the required ambient vanishing on the original full fibre.

\item Top-degree coherent base change now gives $R^2f_*(T_Y\otimes L)=0$. For sufficiently large integer $n$, Serre vanishing on the base gives
\[
H^1\bigl(\PP^1,R^1f_*(T_Y\otimes L)\otimes\OO(n)\bigr)=0.
\]
Projection formula and the two-column Leray sequence imply
$H^2(Y,T_Y\otimes L\otimes f^*\OO(n))=0$.
This is an actual twist in the same Picard component up to the fixed finite-order class of $f^*\OO(n)$; choosing $n$ divisible by the finite order of its Chern class keeps the component unchanged. Semicontinuity in the explicit Picard family of C.1 makes the degree-two zero locus dense open in that component. Intersect it with the dense open endpoint-vanishing loci from C.1. On this intersection $\chi(T_Y\otimes L)=-h^1\le0$.

\item Statement \ref{cdp:st:A.4.6} gives Euler number zero or two for each reduced irreducible normal-crossing fibre; the positive case has rational normalization. Our normalization hypothesis excludes that case, since rational smooth surfaces have $H^1(-,\Q)=0$. Smooth torus fibres and the smooth finite etale quotients in the second fibre class have Euler number zero. Therefore \SeriesRef{cdp:st:C2a} gives $c_3(Y)=0$. The index formula in C.1 gives $\chi=0$, forcing $h^1=0$ on the same dense open set. This proves all of \eqref{cdp:eq:C.15}.

\end{enumerate}
\end{proof}






\begin{proposition}[An unavoidable failure on a rational homology six-sphere]\hypertarget{tag-01FL}{}\label{cdp:st:C4b}
\StatementRecord{cdp:st:C4b}{proved}{01FL}{\SeriesRef{cdp:st:C4a}, \SeriesRef{cdp:st:C2a}, \SeriesRef{cdp:st:A.4.6}}

Let $f:Y\to\PP^1$ be a connected-fibre holomorphic map from a compact complex threefold with $b_1=b_2=b_3=0$. For every $L\in\Pic(Y)$ there is a scheme-theoretic fibre $F$ with
\[
 H^2(F,(T_Y\otimes L)|_F)\ne0,\qquad
 H^0(F,\Omega_Y^1|_F\otimes(K_Y\otimes L^{-1})|_F)\ne0. \tag{C.16}\hypertarget{tag-008U}{}\label{cdp:eq:C.16}
\]
For torus general fibre the reduced fibre Euler numbers sum to two. Within the fibre class of \ref{cdp:st:C4a}, exactly one fibre has rational normalization; it is an irreducible reduced normal-crossing fibre of Euler number two. For that normalization,
\[
H^2(\widetilde S,K_{\widetilde S}\otimes
             \eta^*(K_Y^{-1}\otimes L))=\C
\quad\hbox{for every }L\in\Pic(Y).                    \tag{C.17}\hypertarget{tag-008V}{}\label{cdp:eq:C.17}
\]
\end{proposition}


\begin{proof}

\begin{enumerate}[label={$\langle1\rangle$\arabic*.},leftmargin=*]
\item The divisor of any fibre is nonzero and effective. The index argument C.1 therefore gives $H^2(Y,T_Y\otimes M)\ne0$ for every line bundle $M$. Suppose $R^2f_*(T_Y\otimes L)=0$. Twisting by a sufficiently large positive pullback from the base kills $H^1(R^1f_*)$, so Leray would give a line bundle $M=L\otimes f^*\OO(n)$ with $H^2(Y,T_Y\otimes M)=0$, a contradiction. Thus $R^2f_*(T_Y\otimes L)$ is nonzero. Its stalk has nonzero residue at some point, and top-degree base change gives the first nonvanishing in \eqref{cdp:eq:C.16}. Cartier-fibre adjunction and Serre duality give the second, with the exact line bundle displayed there.

\item For torus general fibre \SeriesRef{cdp:st:C2a} gives $\sum e(F_{\rm red})=e(Y)=2$. In the class of \ref{cdp:st:C4a}, \ref{cdp:st:A.4.6} makes every singular reduced contribution either zero or two, with two exactly for the rational hexagonal normalization; smooth finite torus quotients have zero Euler number. There is consequently exactly one positive fibre. Its smooth rational normalization has $\Pic^0=0$ and torsion-free integral second cohomology. Every ambient line bundle has torsion first Chern class, whose pullback is therefore zero. In particular both $\eta^*K_Y$ and $\eta^*L$ are trivial. Serre duality on the connected rational normalization gives \eqref{cdp:eq:C.17}.

\end{enumerate}
\end{proof}






\begin{proposition}[The finite singular-fibre test and a fixed unavoidable fibre]\hypertarget{tag-01FM}{}\label{cdp:st:C4c}
\StatementRecord{cdp:st:C4c}{proved}{01FM}{\SeriesRef{cdp:st:A.1.1}, \SeriesRef{cdp:st:A.3.1}, \SeriesRef{cdp:st:A.3.5}, \SeriesRef{cdp:st:D.1.11}}

Let $f:Y\to\PP^1$ be a proper connected-fibre map from a compact complex threefold with $b_1(Y)=b_2(Y)=0$ and smooth general fibre a complex two-torus. Let $\Sigma$ be its finite set of critical values, and fix a connected component $P$ of $\Pic(Y)$. For $c\in\Sigma$ retain the entire Cartier fibre $F_c$, including its multiplicities, and put
\[
 Z_c=\{L\in P:H^0(F_c,\Omega_Y^1|_{F_c}\otimes
                  (K_Y\otimes L^{-1})|_{F_c})\ne0\}.
                                                               \tag{C.18}\hypertarget{tag-008W}{}\label{cdp:eq:C.18}
\]
Then the following are equivalent: every $Z_c$ is a proper subset of $P$; and $H^2(Y,T_Y\otimes L)=0$ on a dense open subset of $P$. In that case the generic twist in $P$ satisfies
\[
 H^q(Y,T_Y\otimes L)=0\quad(q\ne1),\qquad
 \chi(Y,T_Y\otimes L)\le0,\qquad c_3(Y)\le0.          \tag{C.19}\hypertarget{tag-008X}{}\label{cdp:eq:C.19}
\]
For a reduced fibre, membership in the complement of $Z_c$ is exactly the two explicitly typed kernel-and-connecting-map conditions of \ref{cdp:st:A.1.1}, with kernel $N^*\otimes\mathcal J^{-1}$ computed in A.2. For $F_c=mS$, the equivalent complete test is the finite lifting filtration of \ref{cdp:st:A.3.1}. These tests use only sheaves and maps on the finitely many singular full fibres.

When $b_3(Y)=0$, at least one fixed critical fibre $F_c$ satisfies $Z_c=P$. Its support is irreducible, by \ref{cdp:st:A.3.5}. Under the integral assumption $H^2(Y,\Z)=0$, the Picard group is connected, so this is a single fixed fibre for which the ambient group in \eqref{cdp:eq:C.18} is nonzero for every line bundle on $Y$.
\end{proposition}


\begin{proof}

\begin{enumerate}[label={$\langle1\rangle$\arabic*.},leftmargin=*]
\item The universal Picard family constructed in C.1 restricts to a holomorphic family on the proper space $F_c\times P$. Its underlying sheaf is flat over $P$, because it is a line-bundle twist of the fixed locally free restriction $\Omega_Y^1|_{F_c}$. Proper semicontinuity therefore makes $Z_c$ a closed complex analytic subset of $P$.

\item Assume that each $Z_c$ is proper. The complement of their finite union is open dense and has full Lebesgue measure in $P\simeq\C^q$. Statement \ref{cdp:st:D.1.11} supplies a dense set of full measure of line bundles non-torsion on every full fibre. Choose $L$ in the intersection. On every smooth torus fibre, $K_Y|_F$ and the conormal are trivial. The intrinsic cotangent bundle is trivial as well, so the conormal exact sequence expresses the ambient cotangent restriction as an extension of trivial bundles. The line $L|_F$ has zero integral Chern class, since $c_1(L)$ is torsion and $H^2(F,\Z)$ is torsion-free; it is nontrivial by the chosen non-torsion condition. Its cohomology in every degree vanishes, as follows by the unitary Fourier resolution of a topologically trivial line on a torus. The ambient section group is consequently zero at every smooth fibre. At each critical fibre it is zero by the definition of the chosen complement.

\item Serre duality on each Gorenstein Cartier fibre, with $\omega_{F_c}=K_Y|_{F_c}$, and top-degree base change give $R^2f_*(T_Y\otimes L)=0$. Twisting by $f^*\OO(n)$ preserves every fibre restriction. For sufficiently large $n$, it kills $H^1(\PP^1,R^1f_*(T_Y\otimes L)\otimes\OO(n))$. We choose $n$ divisible by the finite order of $c_1(f^*\OO(1))$, so this twist stays in $P$. The two-column Leray sequence gives an actual line bundle in $P$ with zero $H^2$. Proper semicontinuity makes this zero locus dense open. Intersecting it with the dense open endpoint-vanishing loci of C.1 proves \eqref{cdp:eq:C.19}, using $\chi=c_3/2$.

\item Conversely, suppose $H^2$ vanishes on a dense open subset of $P$, and choose $L$ in that subset and in the full-measure set of \ref{cdp:st:D.1.11}. The preceding smooth-fibre argument makes $R^2f_*(T_Y\otimes L)$ a torsion sheaf. Leray gives a surjection
\[
 H^2(Y,T_Y\otimes L)\longrightarrow
 H^0(\PP^1,R^2f_*(T_Y\otimes L)).
\]
Thus that torsion sheaf is zero. Top-degree base change and fibre duality imply $L\notin Z_c$ for every $c$. In particular none of the $Z_c$ equals $P$, proving the equivalence.

\item When $b_3=0$, C.1 gives $H^2(Y,T_Y\otimes L)\ne0$ for every $L$. The equivalence therefore forces some $Z_c$ to be all of $P$. There are only finitely many $c$, and $P$ is an irreducible vector-space torsor; a finite union of proper closed analytic subsets cannot exhaust it. The integral hypothesis makes $P=\Pic(Y)$ by the exponential sequence. This proves the assertion about a fixed fibre. Finally the identifications in A.1--A.3 prove the claimed explicit forms of the fibre tests, without replacing a thickened fibre by its support.

\end{enumerate}
\end{proof}



\ref{cdp:st:C4a} is an explicit geometric instance of \ref{cdp:st:C4c}. The two vanishing groups on the support in \ref{cdp:st:B1}--\ref{cdp:st:B2} do not give its fibre test: their ambient space is one-dimensional. The compact quotient \ref{cdp:st:A.3.2} proves that the full multiple fibre cannot be replaced by one unshifted reduced-support test. \ref{cdp:st:C4c} supplies the fixed full-fibre obstruction in arbitrary torus fibrations. \ref{cdp:st:C4e} proves that its normalization target is nonzero for every twist, with arbitrary singularities and multiplicities.




\begin{proposition}[A finite transgression test for every CDP general-fibre type]\hypertarget{tag-01FN}{}\label{cdp:st:C4d}
\StatementRecord{cdp:st:C4d}{proved}{01FN}{\SeriesRef{cdp:st:D.1.9}, \SeriesRef{cdp:st:A.4.9}, \SeriesRef{cdp:st:A.4.7}, \SeriesRef{cdp:st:A.3.4}, \SeriesRef{cdp:st:A.5.2}}

Let $f:Y\to\PP^1$ satisfy the full hypotheses of CDP20 Theorem~2.2: $Y$ is a compact connected complex threefold, $H^1(Y,\Z)=H^2(Y,\Z)=0$, and $f$ is its holomorphic algebraic reduction with connected fibres. Let $\beta\in H^1(Y,\OO_Y)$ lift the integral transgression generator of \ref{cdp:st:D.1.9}, and put $M_z=\exp(2\pi i z\beta)$, with the same exponential convention as in that statement. For each critical value $c$, define
\[
 B_c=\{z\in\C:
 H^0(F_c,\Omega_Y^1|_{F_c}\otimes M_z|_{F_c})\ne0\},
 \qquad F_c=f^{-1}(c).
 \tag{C.20}\hypertarget{tag-008Z}{}\label{cdp:eq:C.20}
\]
When each $B_c$ is a proper subset of $\C$, a dense open subset of $\Pic(Y)$ satisfies the vanishing and the two inequalities in (C.19). The premise is a finite collection of tests on the actual full fibres, given by A.1--A.3; it is satisfied in particular when every reduced support is smooth.

When $c_3(Y)>0$, there is a fixed critical value $c$ with $B_c=\C$. Its reduced support is nonnormal. Thus a nonnormal reduced fibre is necessary for positive Euler characteristic under the full CDP hypotheses, with arbitrary multiplicities and arbitrary singularities allowed at the other fibres.
\end{proposition}


\begin{proof}

\begin{enumerate}[label={$\langle1\rangle$\arabic*.},leftmargin=*]
\item The transgression lift exists because the coherent Leray sequence over $\PP^1$ has no degree-one or degree-two base term for $\OO_{\PP^1}$. The integral degree-one Leray sequence identifies the integral invariant group with $\Z$ under the two stipulated integral vanishings. \ref{cdp:st:D.1.9} proves that its restriction to every full fibre, and to every nearby smooth fibre, is nonzero. On a neighbourhood $U$ of any fibre, $M_r$ for real $r$ has the unitary local system
\[
 a\longmapsto\exp(2\pi i r\langle\gamma_U,a\rangle)
\]
as its flat realization. This equality follows from the common integral lift $\gamma_U$ in the exponential sequence; no global flat realization on $Y$ is used.

\item The holomorphic family of exponentiated Cech cocycles representing $z\beta$ restricts to a flat family of coherent bundles on $F_c\times\C$. Proper semicontinuity makes $B_c$ a closed analytic subset of the complex line. A proper such subset is discrete and countable. There are finitely many critical values. We can therefore choose a real irrational $r$ outside the union of all $B_c$ when all are proper.

\item On a smooth fibre $F$, the character just displayed is the irrational exponential of a nonzero integral class. The explicit calculations for tori, bielliptic surfaces, Kodaira surfaces, Hopf surfaces and Inoue surfaces in \ref{cdp:st:A.4.9} give
\[
 H^0(F,\mathbb A_r)=H^1(F,\mathbb A_r)=0.
\]
The unitary maximum principle gives $H^0(F,M_r)=0$, and the injection \ref{cdp:st:A.4.7} gives $H^0(F,\Omega_F^1\otimes M_r)=0$. The conormal sequence of a smooth full fibre, whose conormal line is trivial, then gives
$H^0(F,\Omega_Y^1|_F\otimes M_r)=0$.
At every critical full fibre this group vanishes by our choice of $r$. Set $L=K_Y\otimes M_r^{-1}$. Cartier-fibre duality and top-degree base change give $R^2f_*(T_Y\otimes L)=0$. A sufficiently positive base twist kills $H^1(R^1f_*(T_Y\otimes L))$, so
\[
 H^2(Y,T_Y\otimes L\otimes f^*\OO(n))=0
\]
for an actual integer $n$. The exponential sequence identifies $\Pic(Y)$ with a connected complex vector space. Semicontinuity therefore makes degree-two vanishing dense open in this whole group. Intersecting with the endpoint-vanishing loci of C.1 gives (C.19).

\item For a full fibre with smooth reduction, \ref{cdp:st:A.3.4} proves the ambient vanishing for every real irrational $r$, including all the nilpotent layers. Its set $B_c$ is consequently proper. This proves the asserted geometric sufficient case.

\item Finally $c_3(Y)=2-b_3(Y)$, and the index formula is integral. Thus $c_3(Y)>0$ implies $b_3(Y)=0$. C.1 then gives $H^2(Y,T_Y\otimes L)\ne0$ for every line bundle $L$. The preceding proof shows that the finitely many $B_c$ cannot all be proper; hence one fixed $B_c$ is the entire complex line. Its reduced support cannot be smooth by \ref{cdp:st:A.3.4}. A connected normal reduced support would be irreducible. Since it is a Cartier divisor in $Y$, its dualizing line has rational first Chern class zero and its structure-sheaf Euler characteristic is zero by divisor Riemann--Roch, exactly as calculated in \ref{cdp:st:A.5.2}. That statement proves it smooth. We have therefore obtained a nonnormal support, as claimed.

\end{enumerate}
\end{proof}



The fixed fibre in \ref{cdp:st:C4d} carries an ambient section along the entire complex transgression line. For torus general fibre, \ref{cdp:st:C4c} strengthens this to all twists in the Picard group. For a general fibre with $h^1(\OO)=1$ and $h^2(\OO)=0$, \ref{cdp:st:A.5.13} identifies that line with the entire Picard group. The following proof treats both cases through their actual conductor schemes.




\begin{proposition}[The unrestricted normalization target under the integral CDP hypotheses]\hypertarget{tag-01FO}{}\label{cdp:st:C4e}
\StatementRecord{cdp:st:C4e}{proved}{01FO}{\SeriesRef{cdp:st:C.3.1}, \SeriesRef{cdp:st:C.3.7}, \SeriesRef{cdp:st:C.3.9}, \SeriesRef{cdp:st:C.3.10}, \SeriesRef{cdp:st:C4c}, \SeriesRef{cdp:st:C4d}, \SeriesRef{cdp:st:A.5.13}, \SeriesRef{cdp:st:A.3.5}, \SeriesRef{cdp:st:A.3.4}, \SeriesRef{cdp:st:A.5.2}, \SeriesRef{cdp:st:A.3.1}, \SeriesRef{cdp:st:A.4.9}, \SeriesRef{cdp:st:A.1.1}, \SeriesRef{cdp:st:A.5.4}, \SeriesRef{cdp:st:A.5.3}, \SeriesRef{cdp:eq:C.22}, \SeriesRef{cdp:st:A.5.12}, \SeriesRef{cdp:st:C2d}}

Let $f:Y\to\PP^1$ satisfy the integral CDP hypotheses and $c_3(Y)>0$. Then a fixed irreducible nonnormal reduced support $S$ has normalization $\eta:\widetilde S\to S$ satisfying
\[
 H^2\bigl(\widetilde S,\omega_{\widetilde S}\otimes
           \eta^*(K_Y^{-1}\otimes L)\bigr)\ne0
 \qquad\hbox{for every }L\in\Pic(Y).                  \tag{C.21}\hypertarget{tag-0090}{}\label{cdp:eq:C.21}
\]
The assertion concerns the reduced support; no multiplicity-one conclusion is included.
\end{proposition}


\begin{proof}

\begin{enumerate}[label={$\langle1\rangle$\arabic*.},leftmargin=*]
\item \ref{cdp:st:C.3.1} and the exclusions \ref{cdp:st:C.3.7}, \ref{cdp:st:C.3.9} and \ref{cdp:st:C.3.10} show
that the general fibre is a two-torus, a Hopf surface, or
an Inoue surface. In the last two cases it has
$h^1(\OO)=1$ and $h^2(\OO)=0$.
For torus general fibre, \ref{cdp:st:C4c} gives a fixed full fibre with a nonzero ambient cotangent section for every ambient twist. In the second case, \ref{cdp:st:C4d} gives persistence along $\C\beta$, and \ref{cdp:st:A.5.13} identifies this with $\Pic(Y)$. From this point the proof uses precisely that persistence property and does not use $c_3(Y)>0$ again. \ref{cdp:st:A.3.5} makes its support irreducible, so it is $mS$ for a single positive integer $m$. \ref{cdp:st:A.3.4} excludes a smooth support, and \ref{cdp:st:A.5.2} then excludes a normal support. Its filtration in \ref{cdp:st:A.3.1} has finitely many graded factors
\[
 \Omega_Y^1|_S\otimes M|_S\otimes\OO_Y(-jS)|_S,
 \qquad 0\le j<m.
\]
Nonvanishing on the filtered sheaf requires nonvanishing on at least one factor. Each factor's nonvanishing locus is closed analytic in the connected vector group $\Pic(Y)$. A finite union of proper such loci cannot cover that group. Hence one factor has nonzero sections for every $M$. Translating by $\OO_Y(-jS)$ shows
$H^0(S,\Omega_Y^1|_S\otimes M|_S)\ne0$ for every $M$.
\ref{cdp:st:A.4.9} kills the torsion-free quotient on a dense open set of twists. The exact sequence \ref{cdp:st:A.1.1} and proper semicontinuity then show that the inverse-Jacobian kernel itself has sections for every twist.

\item We apply the precise persistent-kernel reduction \ref{cdp:st:A.5.4}. It gives a fixed component of the normalization on whose resolution every twist of one fixed rank-one sheaf has a section. The Gauduchon degree argument there proves that the pullback of every ambient line has degree zero for every Gauduchon metric. Its conductor comparison excludes nonnegative Kodaira dimension, and its class VII argument makes every ambient line trivial on that normalization component. Rational components have the same triviality because their Picard group has no connected part. Either case gives \eqref{cdp:eq:C.21} immediately by Serre duality and descent from resolution. It remains to treat an irrational ruled resolution $p:Z\to B$.

\item Let $h:Z\to S\subset Y$. \ref{cdp:st:A.5.3} supplies an effective conductor divisor $D_Z$ with
\[
 K_Z+D_Z=h^*(K_Y+S),\qquad D_Z\cdot F=2              \tag{C.22}\hypertarget{tag-0091}{}\label{cdp:eq:C.22}
\]
for a general ruling fibre $F$. Its horizontal part is therefore either reduced of total ruling degree two, or $2C$ for a section $C$. We record the comparison through each ruling blowdown $\sigma:Z\to Z'$ used below. Put $D'=\sigma_*D_Z$, let $a$ be the coefficient of its exceptional curve $E$ in $D_Z$, and put $d=\operatorname{mult}_{\sigma(E)}D'$. Then
\[
 K_Z+D_Z=\sigma^*(K_{Z'}+D')+(1+a-d)E.
 \tag{C.22b}\hypertarget{tag-0093}{}\label{cdp:eq:C.22b}
\]
Intersection with $E$, using $c_1(K_Z+D_Z)=0$ and $E^2=-1$, gives $a=d-1$. Thus the log-canonical divisor is the exact pullback of $K_{Z'}+D'$, and the latter again has zero rational Chern class. No factorization of $h$ through the blowdown is asserted or needed.

\item We first treat the reduced horizontal part $H$. If a component is exceptional for the normalization resolution, every ambient line restricts trivially to it. Its normalization maps finitely onto $B$, and the norm composed with pullback on $\Pic^0(B)$ is multiplication by the covering degree. Hence this restriction has finite kernel; the connected ambient Picard image is trivial, giving \eqref{cdp:eq:C.21}. We may therefore assume that the horizontal components map to conductor curves. Adjunction on each of them gives
\[
 2p_a(H_i)-2+(D_Z-H_i)\cdot H_i=0.
\]
The second term is nonnegative. Since $H_i$ dominates a positive-genus base, its arithmetic genus is at least one. Thus the base has genus one, each $H_i$ is smooth elliptic, and the horizontal curves are disjoint from each other and from the vertical part. Contract ruling $(-1)$-curves to a minimal bundle. A contracted curve outside $D_Z$ has $D_Z\cdot E=1$, so cannot identify two horizontal points or create a horizontal node. A contracted curve in the vertical part is disjoint from $H$. Downstairs the horizontal curves remain smooth and disjoint, and any nonzero vertical divisor on a minimal ruled bundle meets them. Consequently the downstairs vertical part is zero. Reversing the contractions, a blowup away from the horizontal conductor would give coefficient $-1$ on its exceptional curve in \eqref{cdp:eq:C.22}, which persists under further strict transform and contradicts effectiveness. A blowup at a smooth point of a coefficient-one horizontal curve gives exceptional coefficient zero. Hence $D_Z=H$ has no vertical component. Every codimension-one singularity of $S$ is then an ordinary node, and \CompendiumRef{II}{cdp:st:A.2.3} identifies the inverse-Jacobian kernel with the normalization kernel globally, including its codimension-two points. A section of a topologically trivial line on the Kahler resolution has an effective zero divisor of zero Kahler volume, so is nowhere zero. Persistence for all twists therefore makes all ambient pullbacks trivial and proves \eqref{cdp:eq:C.21}.

\item We now have $D_Z=2C+V$, where $C$ is a section and $V$ is vertical. If $C$ is exceptional for $Z\to\widetilde S$, every ambient line is trivial on $C$. Restriction $\Pic^0(Z)\to\Pic^0(C)$ is an isomorphism, because $C$ is a ruling section. This again makes all ambient pullbacks trivial. Otherwise let $D\subset S$ be its image and let $D^\nu$ normalize that curve. When $D^\nu$ is rational, every ambient line pulls back trivially to $D^\nu$, hence to $C$, giving the same conclusion.

\item Suppose instead that $g(D^\nu)\ge1$. A vertical rational curve cannot map finitely onto $D^\nu$, so all the normalization branches over the general point of $D$ belong to $C$. Their conductor coefficient is two. We derive the possible transverse reduced plane curves rather than assuming a list of singularities. Write its equation as $f=\prod_i f_i$ in the smooth-plane local ring and let $R_i$ and $B_i$ be a branch ring and its normalization. The elements of $R=\OO/(f)$ supported only on branch $i$ form
\[
 \left(\prod_{j\ne i}f_j\right)\bigg|_{f_i}\,R_i.
\]
This equality follows from
$\bigcap_{j\ne i}(f_j)=(\prod_{j\ne i}f_j)$ in that unique
factorization domain. Multiplying by arbitrary elements of
$B_i$ therefore shows that the conductor exponent on branch $i$
is its intrinsic exponent plus $\sum_{j\ne i}I(f_i,f_j)$.
A singular branch has intrinsic exponent at least two:
exponent one would imply $\C+t\C\{t\}=\C\{t\}\subset R_i$.
Exponent two forces $R_i=\C+t^2\C\{t\}=\C\{t^2,t^3\}$,
since an element of order one would again make the ring smooth.
Any additional branch has intersection multiplicity at least
two with a singular branch and is therefore excluded. Thus a
singular branch gives exactly the cusp. When all branches are
smooth, their pairwise positive intersection multiplicities sum
to two on each branch. This gives two branches of contact two
(the tacnode), or three pairwise transverse branches (the
ordinary triple point), and exhausts the possibilities.
The only remaining alternative is a nonconstant pullback
$\Pic(Y)\to\Pic(Z)$; a constant pullback is the zero
homomorphism and already gives \eqref{cdp:eq:C.21}. In that alternative,
the coherent conductor argument \ref{cdp:st:A.5.12} gives $p_a(D)\le1$
for torus general fibre, and \ref{cdp:st:A.5.13} gives the same bound
in the second case.
Since $g(D^\nu)\ge1$, it follows that $D$ itself is smooth
elliptic. This conclusion permits nonreduced conductor image
schemes and imposes no assumption on the original singular
curve. We can therefore apply \ref{cdp:st:C2d} directly. It makes the
transverse type constant, gives the etale conductor cover
$C\to D$, and proves $C^2=0$. In particular $C$ and the
ruling base are elliptic.

\item Adjunction in \eqref{cdp:eq:C.22} gives $V\cdot C=0$. No ruling $(-1)$-curve meets $C$: for one outside $V$, its intersection with $K_Z+2C+V$ would be at least $-1+2=1$; for one in $V$, the equality $V\cdot C=0$ forbids such an intersection. Thus every contraction to the minimal ruled bundle avoids $C$ and preserves its self-intersection zero. On that bundle $K+2C$ is numerically trivial, so the effective vertical part is numerically trivial and hence zero. A first reverse blowup, necessarily away from $C$, would again produce negative conductor coefficient. There are therefore no reverse blowups: $Z$ is minimal ruled and $D_Z=2C$. There is also no exceptional curve for the normalization resolution: a negative curve $E\ne C$ would satisfy $2p_a(E)-2=E^2-2C\cdot E<0$, hence would be rational and contained in a ruling fibre; a minimal ruled bundle has only self-intersection-zero fibres. Thus the normalization itself is smooth here.

\item By the explicit inverse-Jacobian computations in \ref{cdp:st:C2d}, its pullback allows at most a simple pole along $C$. There are no other divisorial poles. Hartogs extension for the line bundle on the smooth surface gives an injection of the kernel into $h^*\OO_Y(-S)\otimes\OO_Z(C)$. The exact sequence of the section $C$ and $R^1p_*\OO_Z=0$ give
\[
 0\longrightarrow\OO_B\longrightarrow p_*\OO_Z(C)
   \longrightarrow N_{C/Z}\longrightarrow0,
 \qquad \deg N_{C/Z}=0.                              \tag{C.22a}\hypertarget{tag-0092}{}\label{cdp:eq:C.22a}
\]
For a degree-zero line $P$ on the elliptic curve, both end terms after tensoring by $P$ have zero sections unless $P=\OO_B$ or $P=N_{C/Z}^{-1}$. Thus $H^0(Z,\OO_Z(C)\otimes p^*P)$ is supported at at most two points of $\Pic^0(B)$. The ambient pullback image is a connected complex subgroup of this Picard group. Persistence for every ambient twist forces that subgroup to be a point, and it contains the identity. Therefore every ambient line pulls back trivially to $Z$ in this last case as well.

\item The equality $\pi_*\OO_Z=\OO_{\widetilde S}$ for the normal resolution descends triviality of a pulled-back line: projection formula identifies its space of sections with that of its trivial pullback, and the same statement for its inverse gives mutually inverse trivializations. Consequently one normalization component has trivial $\eta^*(K_Y^{-1}\otimes L)$ for every $L$. Serre duality on that compact normal Cohen--Macaulay surface identifies its top dualizing cohomology with $H^0(\OO)^\vee=\C$. It is a direct summand of the target in \eqref{cdp:eq:C.21}, proving the claim.

\end{enumerate}
\end{proof}



The local example \ref{cdp:st:A.2.5} has actual vertical inverse-Jacobian
poles, and \ref{cdp:st:A.2.6} has a critical curve that is not a local
complete intersection. Their local rings remain part of the
audit. The global conductor-duality arguments \ref{cdp:st:A.5.12}--\ref{cdp:st:A.5.13} prove
that a positive-genus conductor image in the persistent ruled
case is smooth, and the exact contractions above then remove
the vertical conductor. Thus neither local phenomenon leaves
an unproved step in the stated theorem. The assertion
is about the normalization of a reduced support and does not
claim that its full fibre has multiplicity one.




\begin{proposition}[An unrestricted geometric replacement for CDP20 Theorem 2.2]\hypertarget{tag-01FP}{}\label{cdp:st:C4f}
\StatementRecord{cdp:st:C4f}{proved}{01FP}{\SeriesRef{cdp:st:C4e}, \SeriesRef{cdp:st:C4c}, \SeriesRef{cdp:st:A.3.4}, \SeriesRef{cdp:st:A.3.5}, \SeriesRef{cdp:st:A.5.2}, \SeriesRef{cdp:st:A.2.2}, \SeriesRef{cdp:eq:C.VC.4}, \SeriesRef{cdp:st:A.5.13}, \SeriesRef{cdp:st:C4d}, \SeriesRef{cdp:st:D.1.11}, \SeriesRef{cdp:st:A.4.9}}

Let $f:Y\to\PP^1$ satisfy the integral CDP hypotheses.
For each critical full fibre with
irreducible nonnormal reduced support $S$, let
$\eta_S:T_S\to S$ be its normalization. Then a general
line bundle $L\in\Pic(Y)$ satisfies
\[
 H^q(Y,T_Y\otimes L)=0\quad(q\ne1)
 \tag{C.VC.1}\hypertarget{tag-015G}{}\label{cdp:eq:C.VC.1}
\]
exactly when, for every such $S$, the holomorphic pullback
homomorphism
\[
 \Pic(Y)\longrightarrow\Pic(T_S)
 \tag{C.VC.2}\hypertarget{tag-015H}{}\label{cdp:eq:C.VC.2}
\]
is nonzero. Equivalently, its differential
$H^1(Y,\OO_Y)\to H^1(T_S,\OO_{T_S})$ is nonzero.
In this case $\chi(Y,T_Y\otimes L)\le0$ and $c_3(Y)\le0$.
If one map \eqref{cdp:eq:C.VC.2} is zero, then
$H^2(Y,T_Y\otimes L)\ne0$ for every $L\in\Pic(Y)$.
All fibre multiplicities and all reduced-support
singularities are permitted in this criterion.
\end{proposition}


\begin{proof}

\begin{enumerate}[label={$\langle1\rangle$\arabic*.},leftmargin=*]
\item The classification and exclusions cited in \ref{cdp:st:C4e} leave
exactly the two cohomological cases used in that proof:
torus general fibre, or $h^1(\OO)=1$, $h^2(\OO)=0$.
The Picard group of $Y$ is the connected exponential image
of $H^1(Y,\OO_Y)$. A holomorphic group homomorphism out of
this group is zero precisely when its differential is zero.
For each fixed critical full fibre, the ambient-section
nonvanishing locus is closed analytic in this group, as
in \ref{cdp:st:C4c}. \ref{cdp:st:A.3.4} makes it proper when the support is smooth,
and \ref{cdp:st:A.3.5} makes it proper when the support is reducible.
A connected normal support is irreducible and is smooth
by \ref{cdp:st:A.5.2}, so these cases leave just the supports in the
statement.

\item Suppose the nonvanishing locus for $F=mS$ is the whole
Picard group. The proof of \ref{cdp:st:C4e}, beginning with its
fixed-fibre persistence assertion, proves that \eqref{cdp:eq:C.VC.2}
is zero. Its steps use the actual component filtration,
the inverse-Jacobian kernel, the Gauduchon bound, the
effective conductor comparison, the coherent conductor
genus bound, and the explicit constant-transverse-type
calculations. Each is independent of the Euler number
once persistence has been given.

\item Conversely, suppose \eqref{cdp:eq:C.VC.2} is zero. The last nonzero
ideal layer of the full fibre is an actual subsheaf
\[
 (\Omega_Y^1\otimes M\otimes\OO_Y(-(m-1)S))|_S
 \hookrightarrow
          (\Omega_Y^1\otimes M)|_{mS}.
 \tag{C.VC.3}\hypertarget{tag-015I}{}\label{cdp:eq:C.VC.3}
\]
The canonical normalization submodule in \ref{cdp:st:A.2.2} gives
an injection into the left side from
\[
 \eta_{S*}\bigl(\eta_S^*((M\otimes\OO_Y(-mS))|_S)\bigr).
 \tag{C.VC.4}\hypertarget{tag-015J}{}\label{cdp:eq:C.VC.4}
\]
Here the first conormal factor $\OO_Y(-S)|_S$ combines
with the displayed $(m-1)$st ideal layer; this comparison
retains every nilpotent layer needed for the injection.
Since $mS$ is a full fibre, $\OO_Y(-mS)|_S$ is a trivial
line, with its one-dimensional scalar factor determined
by a local base parameter. The pullback of $M$ to $T_S$
is trivial by \eqref{cdp:eq:C.VC.2}. Hence \eqref{cdp:eq:C.VC.4} has a nonzero global
section for every $M$, and its injections prove that the
full-fibre nonvanishing locus is the entire Picard group.
This also proves the converse directly for arbitrary $m$.

\item We have proved an exact geometric characterization of
every improper critical-fibre locus. For torus general
fibre, applying \ref{cdp:st:C4c} gives the equivalence with generic
degree-two vanishing. In the second case \ref{cdp:st:A.5.13} identifies
$\Pic(Y)$ with the transgression line, so properness of
all these loci is exactly the premise of \ref{cdp:st:C4d} and gives
generic degree-two vanishing there as well.
The endpoint vanishings of C.1 then give \eqref{cdp:eq:C.VC.1} and the
index inequalities.
For the last assertion in the torus case, generic non-torsion on every
smooth torus fibre makes $R^2f_*(T_Y\otimes L)$ a
torsion sheaf for a dense set of twists, by \ref{cdp:st:D.1.11}.
The fixed full-fibre obstruction makes it nonzero.
Its nonzero global sections are a quotient of
$H^2(Y,T_Y\otimes L)$ by the Leray sequence.
Thus degree-two cohomology is nonzero on a dense set.
Its proper semicontinuity locus is closed analytic in
the connected Picard family, so it is nonzero for every
twist. In the second case use instead the actual twists
$L_r=K_Y\otimes M_r^{-1}$ for real irrational $r$.
\ref{cdp:st:A.4.9} and the smooth-fibre conormal sequence make their
$R^2f_*(T_Y\otimes L_r)$ torsion. The fixed full-fibre
obstruction again makes that torsion sheaf nonzero,
so Leray gives $H^2(Y,T_Y\otimes L_r)\ne0$.
These twists have accumulation points in the complex
line $\Pic(Y)=\C\beta$. A closed complex analytic subset
of that line containing them is the whole line.
This proves all-twist nonvanishing in the second case
and completes both implications of the criterion.

\end{enumerate}
\end{proof}






\begin{proposition}[The necessary normalization defect in the three named test families]\hypertarget{tag-01FQ}{}\label{cdp:st:C4g}
\StatementRecord{cdp:st:C4g}{proved}{01FQ}{\SeriesRef{cdp:st:E13}, \SeriesRef{cdp:st:E14}, \SeriesRef{cdp:st:A.3.4}, \SeriesRef{cdp:st:A.4.1}, \SeriesRef{cdp:st:C.1.3}, \SeriesRef{cdp:st:C5b}}

For every admissible member of the original $(3,4)$ family,
every member of the two-parameter family in \ref{cdp:st:E13}, and every
member of the Liu--Rollenske family in \ref{cdp:st:E14}, there is a reduced
full hexagonal fibre $W$ with normalization
$\eta:\mathrm dP_6\to W$ such that
\[
 \eta^*L=\OO_{\mathrm dP_6}\quad(L\in\Pic(Y)),\qquad
 H^2\bigl(\mathrm dP_6,K_{\mathrm dP_6}\otimes
                    \eta^*(K_Y^{-1}\otimes L)\bigr)=\C.
 \tag{C.TEST.1}\hypertarget{tag-012O}{}\label{cdp:eq:C.TEST.1}
\]
In these three families the singular-fibre Euler contributions are
\[
\begin{array}{c|c|c}
\text{family}&\text{reduced singular supports}&
                 \text{their Euler numbers}\\ \hline
(3,4)&W&2\\
\text{two-parameter \ref{cdp:st:E13}}&W,\ W_1,\ W_2&2,0,0\\
\text{Liu--Rollenske \ref{cdp:st:E14}}&W,\ C_{\rm cusp}&2,0 .
\end{array}
\tag{C.TEST.2}\hypertarget{tag-012P}{}\label{cdp:eq:C.TEST.2}
\]
The other critical full fibres are smooth multiple fibres,
with multiplicities $(3,4)$, $4$, and $4$, respectively.
Every one of those full multiple fibres passes the generic
ambient vanishing test \ref{cdp:st:A.3.4}, including all its nilpotent
layers. Consequently positive $c_3$ in these examples does
not require a failure of that test on a nonreduced fibre.
The reduced full fibre $W$ already supplies the fixed
normalization obstruction, and $H^2(Y,T_Y\otimes L)\ne0$
for every ambient line bundle in each family.
\end{proposition}


\begin{proof}

\begin{enumerate}[label={$\langle1\rangle$\arabic*.},leftmargin=*]
\item The original local filling is the hexagonal toric quotient
in \ref{cdp:st:A.4.1}. The two-parameter construction in \ref{cdp:st:E13} places this
same filling at the simple nodal value where its specified
linearization vanishes; its other two simple nodal values
have rank-one fillings. The coordinate and line-bundle
comparisons of \ref{cdp:st:E14} place the hexagonal filling at its
separate simple node and the cuspidal quotient at the
remaining critical value. Those proofs construct the actual
local spaces and their identifications over the punctured
base discs. Thus the normalization of $W$ in each case is
the stated smooth rational surface.

\item For every member in question, $H^2(Y,\Z)$ is a finite
group. For the original affine family this follows from E19:
the universal coefficient sequence gives
\[
 H^2(Y,\Z)\simeq\operatorname{Ext}^1(H_1(Y,\Z),\Z)
 \simeq\Z/d\Z,\qquad d=|p|,
\]
because $H_2(Y,\Z)$ is torsion. In the two other families
the integral homology is that of the six-sphere. In
particular every ambient first Chern class is torsion.
The group $H^2(\mathrm dP_6,\Z)$ is free, so its pullback
to $\mathrm dP_6$ is zero. The exponential sequence there,
together with $H^1(\mathrm dP_6,\OO)=0$, shows that the
pulled-back line bundle is trivial. This applies to both
$L$ and $K_Y^{-1}\otimes L$, and Serre duality on
$\mathrm dP_6$ proves \eqref{cdp:eq:C.TEST.1}.

\item All these threefolds have $b_1=b_2=b_3=0$: use \CompendiumRef{II}{cdp:st:E19} for
every original affine member, and the integral homology
of the six-sphere for the other two families. Any full
fibre is a nonzero effective Cartier divisor. Statement
\ref{cdp:st:C.1.3} therefore gives
\[
 h^2(Y,T_Y\otimes L)\geq 1
 \qquad\text{for every }L\in\Pic(Y).
\]
This argument applies equally when $d>1$ and the integral
cohomology has nonzero torsion. It requires no surjectivity
of the global normalization cohomology map. Indeed \ref{cdp:st:C5b}
computes the actual exceptions to that surjectivity in
the original family.

\item The toric strata of $W$ consist of one $(\C^*)^2$ stratum,
three $\C^*$ strata and two triple points. Their compactly
supported Euler numbers give $e(W)=2$.
Each rank-one normalization is a $\PP^1$-bundle over an
elliptic curve, with its zero and infinity sections glued.
Both the normalization and these conductor curves have
Euler number zero, so normalization additivity gives
$e(W_i)=0$. At the cusp in \ref{cdp:st:E14}, normalization of the
cuspidal rational curve makes its degree-zero line bundle
trivial. The lifted contracting action on this normalized
line complement is a product of a fixed automorphism of
$\PP^1$ and multiplication by a constant of modulus
strictly less than one on $\C^*$. Its quotient is a
$\PP^1$-bundle over the elliptic curve given by that
multiplicative period. The cusp normalization is
bijective; a proper continuous bijection between these
compact Hausdorff spaces is a homeomorphism. Hence
$e(C_{\rm cusp})=0$. This proves \eqref{cdp:eq:C.TEST.2}.

\item Finally the explicit good-reduction covers in \ref{cdp:st:E13}--\ref{cdp:st:E14}
and the original root charts have the stated orders.
Their reduced central fibres are smooth finite unramified
torus quotients, and their Euler numbers are zero.
\ref{cdp:st:A.3.4} applies to their actual full fibres, proving the
asserted generic ambient vanishing. Euler localization
therefore gives $c_3(Y)=2$ in all three rows, while the
nonreduced fibres satisfy the complete multiple-fibre
test. This is the required test of the proposed additional
nonreduced-fibre necessity.

\end{enumerate}
\end{proof}




\begin{proposition}[The normalization quotient and its cohomology map]\hypertarget{tag-01FR}{}\label{cdp:st:C5a}
\StatementRecord{cdp:st:C5a}{proved}{01FR}{\citet[Lemma~6.11, p.~81]{DGKXv1}}

Let $Y$ be a smooth compact complex threefold, $S\subset Y$ a reduced normal-crossing divisor, and $\nu:\widetilde S\to Y$ the normalization followed by inclusion. This is a finite holomorphic immersion on each local branch. Put $\mathcal T=T_Y(-\log S)$. There is a canonical exact sequence of coherent $\OO_Y$-modules
\[
0\longrightarrow\mathcal T\longrightarrow T_Y
 \longrightarrow\nu_*N_\nu\longrightarrow0,
\qquad N_\nu=K_{\widetilde S}\otimes\nu^*K_Y^{-1}.      \tag{C.3}\hypertarget{tag-0094}{}\label{cdp:eq:C.3}
\]
For every holomorphic line bundle $L$ on $Y$, the image of the natural cohomology map is exactly
\[
\begin{split}
\im\bigl(H^2(Y,T_Y\otimes L)&\longrightarrow
 H^2(\widetilde S,N_\nu\otimes\nu^*L)\bigr)\\
 &=\ker\bigl(H^2(\widetilde S,N_\nu\otimes\nu^*L)
 \longrightarrow H^3(Y,\mathcal T\otimes L)\bigr).
\end{split}                                             \tag{C.4}\hypertarget{tag-0095}{}\label{cdp:eq:C.4}
\]
The target identifies by Serre duality with
$H^0(\widetilde S,\nu^*(K_Y\otimes L^{-1}))^\vee$.
\end{proposition}


The branch-normal quotient specializes to \citet[Lemma~6.11, p.~81]{DGKXv1} in the original family.

\begin{proof}
In coordinates with $S=(z_1\cdots z_r=0)$, $1\le r\le3$, the map sends a vector field $\sum a_i\partial_{z_i}$ to the normal coefficient $a_i|_{z_i=0}$ on each branch $i\le r$. The separate coefficients lift independently to holomorphic functions, making the map surjective. Its kernel consists precisely of fields with $z_i\mid a_i$ for $i\le r$, which is $T_Y(-\log S)$. Coordinate changes act on each coefficient through the normal-line transition and permute the branches; hence the map and its exactness descend canonically to the normalization. The immersion tangent sequence on $\widetilde S$ gives $\det\nu^*T_Y=\det T_{\widetilde S}\otimes N_\nu$, proving the adjunction expression. Tensoring by the locally free line bundle $L$ preserves exactness. Finite pushforward has no higher direct images, so the long exact sequence in coherent cohomology gives \eqref{cdp:eq:C.4}. The smooth compact surface $\widetilde S$ has the usual Serre pairing, which gives the stated dual target. In particular the vanishing of $H^3(Y,\mathcal T\otimes L)$ makes this cohomology map surjective.

\end{proof}






In this subsection $X=X_c$ is the period construction of \cite{AC26}, compared with the quotient construction of \cite{Engel2026}, with admissible parameter $c$, standard clockwise affine vectors
$v_1=(1,2,-4,0)^t$, $v_2=(-1,-3,3,0)^t$, and untwisted cusp collar. Thus $p=12l_0-4l_1-3l_2=-1$. We write $\mathfrak B=\PP(3,4)$, $\rho:\mathfrak B\to B$ for its coarse map, and $\bar f:X\to\mathfrak B$ for the root factorization. Let $\mathcal T=T_X(-\log W)$ and let $\nu:\mathrm dP_6\to X$ normalize $W$. The root generator is fixed by
\[
 M_n=\bar f^*\OO_{\mathfrak B}(n)
     =\OO_X(n(S_1-S_2)),\quad n\in\Z.
\]
The exponential isomorphism $H^1(X,\OO_X)\to\Pic(X)$ identifies $\Pic(X)$ with $\C$. We scale its coordinate so that $M_z=\exp(z\gamma^{0,1})$ and $M_1=\OO_X(S_1-S_2)$. The sign of this equality is fixed by the clockwise transgression $12\gamma\mapsto-p=1$. This is checked by evaluating the two finite lifts on $v_1,v_2$: their integral periods are $4l_1,3l_2$, while the cusp contributes $-12l_0$, whose sum is $-p$. The exponential sequence therefore gives $\exp(12\gamma^{0,1})=f^*\OO_B(1)$. Its injectivity and the root relation give the stated $M_1$ equality.

\begin{proposition}[Integral twists and the normalization-map exception]\hypertarget{tag-01FS}{}\label{cdp:st:C5b}
\StatementRecord{cdp:st:C5b}{proved}{01FS}{\SeriesRef{cdp:eq:C.8}, \SeriesRef{cdp:st:C5a}, \SeriesRef{cdp:eq:C.9}}

For $n\in\Z$ define
\[
 r_1(n)\equiv n\pmod3,\quad 0\le r_1(n)<3,
 \qquad r_2(n)\equiv3n\pmod4,\quad0\le r_2(n)<4,
\]
\begin{gather*}
 d(n)=\frac{n-4r_1(n)-3r_2(n)}{12},\qquad
 \epsilon_n=\begin{cases}0&3\mid n,\\1&3\nmid n.\end{cases}
\\
 u(n)=\max\{d(n)+1,0\},\qquad
 v(n)=\max\{-d(n)-1,0\}.
\end{gather*}
The coarse direct images of $T_X\otimes M_n$ are
\begin{align*}
 f_*(T_X\otimes M_n)&=\OO(d(n))\oplus\OO(d(n-1)),\\
 R^1f_*(T_X\otimes M_n)&=\OO(d(n))\oplus\OO(d(n-1))\oplus\OO(d(n+2)),\\
 R^2f_*(T_X\otimes M_n)&=\OO(d(n+6))\oplus\C_{p_1}^{\epsilon_n}.\tag{C.5}\hypertarget{tag-0096}{}\label{cdp:eq:C.5}
\end{align*}
Their global cohomology dimensions are
\begin{align*}
 h^0&=u(n)+u(n-1),\\
 h^1&=v(n)+v(n-1)+u(n)+u(n-1)+u(n+2),\\
 h^2&=v(n)+v(n-1)+v(n+2)+u(n+6)+\epsilon_n,\\
 h^3&=v(n+6).                                             \tag{C.6}\hypertarget{tag-0097}{}\label{cdp:eq:C.6}
\end{align*}
The natural map $H^2(X,T_X\otimes M_n)\to H^2(\mathrm dP_6,K_{\mathrm dP_6})=\C$ has rank one exactly for
\[
 n\in\{-6,-3,-2\}\cup\Z_{\ge0};                         \tag{C.7}\hypertarget{tag-0098}{}\label{cdp:eq:C.7}
\]
it is zero for all other integers. In particular $M_{-1}$ has tangent cohomology $(0,0,1,0)$ and a zero map to the normalization target.
\end{proposition}



\begin{proof}

\begin{enumerate}[label={$\langle1\rangle$\arabic*.},leftmargin=*]
\item We check the root-curve and logarithmic tangent inputs of \SeriesRef{fam:src:A.1.1} and \SeriesRef{fam:st:A.5.2}. The root charts at $p_1,p_2$ are the finite etale torus-family covers used in the construction; finite-group invariants over $\C$ are exact. At the cusp the local equations are $xy=q$ and $xyz=q$, and the vertical logarithmic fields are generated by $x\partial_x-y\partial_y$ and the remaining relative direction. Their affine transition derivatives depend only on the base. The invariant vertical field gives a trivial subline of the vertical bundle $E$. The quotient has twelfth power $\OO(-p_0)$: its automorphy is that of the reciprocal elliptic discriminant, whose only divisor is a simple pole at the cusp. As $\Pic(\mathfrak B)=\Z$, the quotient is $\OO(-1)$; $H^1(\mathfrak B,\OO(1))=0$ splits the extension. Thus $E=\OO\oplus\OO(-1)$.

\item The normalization complex of $W$ has global scalar complex
\[
 \C\xrightarrow{0}\C^3
 \xrightarrow{(a,b,c)\mapsto(a-b+c,a-b+c)}\C^2.
\]
The rational normalization and the three rational double curves have no positive-degree structure-sheaf cohomology. The displayed complex therefore gives $(1,2,1)$ for the fibre cohomology and its exterior cup product; the two oriented triangular faces give the nondegenerate degree-one pairing. The cyclic-cover fibres have the same dimensions. Proper flat base change on these charts gives locally free $G=R^1\bar f_*\OO_X$, with $R^2\bar f_*\OO_X=\det G$. Relative duality and the determinant of the logarithmic sequence give $\det G=\OO(-1)$. The integral class $\gamma$ maps to a nowhere-zero section of $G$: on a marked smooth torus the two identity columns of the period matrix prevent it from belonging to the span of the two holomorphic period rows, and at the cusp it is a nonzero cycle in the displayed complex. The same splitting argument yields
\[
E=G=\OO\oplus\OO(-1),\quad R^2\bar f_*\OO_X=\OO(-1),
\quad0\to\bar f^*E\to\mathcal T\to\bar f^*\OO(-5)\to0. \tag{C.8}\hypertarget{tag-0099}{}\label{cdp:eq:C.8}
\]
These are the actual root-chart and cusp comparisons needed from the earlier returns; no coarse multiple fibre was replaced by its reduction in this argument.

\item Choose root sections $x$ of weight $3$, vanishing at $p_2$, and $y$ of weight $4$, vanishing at $p_1$. The Kodaira--Spencer morphism $\OO(-5)\to E\otimes G$ has, by its weights, the form
\[
 \kappa=\begin{pmatrix}0&a y\\b y&c_0x\end{pmatrix}.
\]
Its first coefficient vanishes because $H^0(\OO(5))=0$. The logarithmic period derivative at the cusp is invertible, giving $ab\ne0$. At $p_1$, the elliptic coordinate $\tau-\rho$ has a simple zero on the root chart, so its period derivative is nonzero, giving $c_0\ne0$. Rescaling the four line frames makes the three nonzero coefficients one. Cup product with this extension class gives
\[
\delta_0=(0,y,y,x)^t,\qquad
\delta_1=\begin{pmatrix}-y&0\\-x&y\end{pmatrix}:
 \OO(-5)\oplus\OO(-6)\to\OO(-1)\oplus\OO(-2).
\]
The first is a subbundle inclusion since $x,y$ have no common zero, and its cokernel is $\OO\oplus\OO(-1)\oplus\OO(2)$. The second is injective with determinant $-y^2$; its cokernel is
$Q=\OO(-1)|_{2p_1}$, a length-two module on the order-three chart. The direct-image sequence of \eqref{cdp:eq:C.8} yields
\[
 R^0\bar f_*\mathcal T=E,\quad
 R^1\bar f_*\mathcal T=\OO\oplus\OO(-1)\oplus\OO(2),\quad
 R^2\bar f_*\mathcal T=\OO(-6)\oplus Q.
\]
The last splitting follows from $H^1(Q(6))=0$; it is not designated canonical.

\item Statement \ref{cdp:st:C5a} gives the quotient $\nu_*K_{\mathrm dP_6}$. The local cusp volume form trivializes $\nu^*K_X$, and every $\nu^*M_z$ is trivial because $H^2(X,\Z)=0$ and $\Pic^0(\mathrm dP_6)=0$. Hence this quotient has cohomology only in degree two, where it is $\C$. At $p_0$, Serre duality and the conormal computation in Part A give $h^2(W,T_X|_W)=1$. Top-degree base change identifies this dimension with the residue dimension of $R^2\bar f_*T_X$. That module has generic rank one, so Nakayama's lemma and the structure theorem over $\C\{q\}$ make it free of rank one at $p_0$. Consequently its line summand is the elementary extension of $\OO(-6)$ by $\C_{p_0}$, namely $\OO(6)$, with no new cusp torsion. We have
\[
 R^0\bar f_*T_X=E,\quad
 R^1\bar f_*T_X=\OO\oplus\OO(-1)\oplus\OO(2),\quad
 R^2\bar f_*T_X=\OO(6)\oplus Q.                       \tag{C.9}\hypertarget{tag-009A}{}\label{cdp:eq:C.9}
\]

\item For a root line, removing its two stabilizer characters gives
$\rho_*\OO(n)=\OO_B(d(n))$. Averaging the finite stabilizers is exact, so there are no higher direct images of $\rho$. Its sections are equivalently monomials $x^ay^b$ with $3a+4b=n$, and Serre duality gives $h^1\OO(n)=h^0\OO(-n-7)$. At $p_1$, the invariant exponent of $Q(n)$ is $n-1$ modulo three. Its two possible coefficients have orders zero and one, so its coarse length is $\epsilon_n$ and it is killed by the coarse parameter. Apply projection formula to \eqref{cdp:eq:C.9} and then $\rho_*$. This proves \eqref{cdp:eq:C.5}, including the support and length of every torsion summand. Leray over $B$ has columns only zero and one, proving \eqref{cdp:eq:C.6}.

\item The normalization map on $R^2$ is the quotient in
\[
0\longrightarrow\OO(n-6)\longrightarrow\OO(n+6)
 \longrightarrow\C_{p_0}\longrightarrow0.
\]
Its global-section map is nonzero precisely when a section of $\OO(n+6)$ is nonzero at $p_0$. Both $x,y$ are nonzero at that point, so this occurs precisely when $n+6$ lies in the additive semigroup generated by $3,4$. That semigroup is $\{0,3,4,6,7,8,\ldots\}$, giving \eqref{cdp:eq:C.7}. The $H^1(R^1)$ part maps to zero in the cohomology of a sheaf supported at one point, by the natural Leray filtration. At $n=-1$, \eqref{cdp:eq:C.5} has only one nonzero cohomological contribution, the reduced $\C_{p_1}$ summand in degree two, while $H^3(X,\mathcal T\otimes M_{-1})=H^1(\OO(-7))=\C$. The boundary from the normalization target is an isomorphism. This proves the stated counterexample to a surjection for every twist.

\end{enumerate}
\end{proof}






\begin{proposition}[Every nonintegral twist and its resonance lengths]\hypertarget{tag-01FT}{}\label{cdp:st:C5c}
\StatementRecord{cdp:st:C5c}{proved}{01FT}{\SeriesRef{cdp:eq:C.13}}

Retain the fixed marked functions $\tau,\mu,\beta_c$ and put $\rho=\tau(z_1)$. For $z\notin\Z$, let $\mathcal R_z$ be the set of regular base points at which integers $(n,p,q,r)$ satisfy
\[
 q\tau+r\mu=p,\qquad z=n-6q\mu-r\beta_c.               \tag{C.10}\hypertarget{tag-008O}{}\label{cdp:eq:C.10}
\]
Let
\[
 \Lambda_1=\Z+(3\beta_c(z_1)+4-2\rho)\Z,\quad
 \Lambda_2=\Z+(2\beta_c(z_2)+3-3i)\Z.
\]
Set $\ell_1(z)=1+\epsilon_n$ for $z=n+k(3\beta_c(z_1)+4-2\rho)$ with $k\ne0$, and zero off $\Lambda_1$. Set $\ell_2(z)=1$ for $z=n+k(2\beta_c(z_2)+3-3i)$ with $k\ne0$ and $n\equiv1,2\pmod4$, and zero otherwise. Then $\mathcal R_z$ is finite. With $N_z=|\mathcal R_z|$ and $J_z=2N_z+\ell_1(z)+\ell_2(z)$, we have
\[
 (h^0,h^1,h^2,h^3)(X,T_X\otimes M_z)=(0,J_z,1+J_z,0),  \tag{C.11}\hypertarget{tag-008P}{}\label{cdp:eq:C.11}
\]
and, writing $\mathfrak m_t$ for the maximal ideal of $\OO_{B,t}$,
\begin{align*}
R^1f_*(T_X\otimes M_z)&=
 \bigoplus_{t\in\mathcal R_z}\C_t^2
 \oplus\C_{p_1}^{\ell_1(z)}\oplus\C_{p_2}^{\ell_2(z)},\\
R^2f_*(T_X\otimes M_z)&=
 \C_{p_0}\oplus\bigoplus_{t\in\mathcal R_z}\OO_{B,t}/\mathfrak m_t^2
 \oplus\C_{p_1}^{\ell_1(z)}\oplus\C_{p_2}^{\ell_2(z)}.\tag{C.12}\hypertarget{tag-008Q}{}\label{cdp:eq:C.12}
\end{align*}
The natural map to $H^2(\mathrm dP_6,K_{\mathrm dP_6})$ is surjective for all these nonintegral twists.
\end{proposition}



\begin{proof}

\begin{enumerate}[label={$\langle1\rangle$\arabic*.},leftmargin=*]
\item On a marked smooth torus, the exponential sequence presents $\Pic^0$ as constant complex lattice cocycles modulo the integral cocycles and the two holomorphic period rows. The equation for $z\gamma$ to belong to this sum is
\[
 z\gamma=(n,p,q,r)+a(6\mu,\tau,1,0)+b(\beta_c,\mu,0,1).
\]
The last two coordinates give $a=-q,b=-r$, and the first two give \eqref{cdp:eq:C.10}. If $r=0$, positivity of $\operatorname{Im}\tau$ forces $p=q=0$ and $z=n$. Otherwise
\[
 \operatorname{Im}z=-r\left(\operatorname{Im}\beta_c-
           6(\operatorname{Im}\mu)^2/\operatorname{Im}\tau\right)\ne0.
\]
In particular real nonintegral twists have no regular resonance. The fixed values $\mu(z_1)=(2-\rho)/3$ and $\mu(z_2)=(1-i)/2$ give the two displayed lattices by solving the first equation. Their nonreal generators make each representation $n+k\lambda_j$ unique.

\item We give a coherent local complex to calculate every length. On a smooth simply connected torus-family disc use normalized period coordinates $\zeta=Z(s)x+y$, $x,y\in\R^2/\Z^2$, and translation-invariant forms $G_j=(dx_j)^{0,1}$. The invariant horizontal lift is $h=\partial_s+\sum_i(Z'(s)x)_i\partial_{\zeta_i}$. Since $dx_j$ is closed, $\bar\partial G_j=0$ and
$\bar\partial h=\sum Z'_{ij}e_iG_j$. After subtracting the integer cocycle of a resonance, the line multiplier has invariant Dolbeault form $a_1G_1+a_2G_2$, up to the common nonzero factor $2\pi i$. Thus the invariant relative complex has terms
\[
 C^j=(V\oplus\OO h)\otimes\Lambda^jG,\qquad j=0,1,2,
\]
\[
 d(v\eta)=v(a\wedge\eta),\qquad
 d(h\eta)=h(a\wedge\eta)+\kappa\wedge\eta,             \tag{C.13}\hypertarget{tag-008R}{}\label{cdp:eq:C.13}
\]
where
\[
 a=(z-n+6q\mu+r\beta_c,-p+q\tau+r\mu),\qquad
 \kappa=\begin{pmatrix}6\mu'&\tau'\\\beta_c'&\mu'\end{pmatrix}.
\]
This complex computes the derived direct image. To verify the comparison, decompose relative smooth forms into Fourier characters under real torus translations. The zero character gives \eqref{cdp:eq:C.13}. For every other character the relative Dolbeault differential is exterior multiplication by a nonzero covector; contraction by its normalized dual is a homotopy. On a sufficiently small closed subdisc those covectors are bounded away from zero at bounded frequency and grow linearly at large frequency. The homotopies preserve smooth Fourier series and all their derivatives. The vertical-horizontal filtration of the tangent extension, followed by the filtration in the base antiholomorphic differential, gives a finite filtration with acyclic nonzero-character quotients. Hence their full complexes are acyclic. The zero-character formulas have holomorphic coefficients, so the remaining base Dolbeault resolution is precisely the resolution of \eqref{cdp:eq:C.13}. A finite cyclic stabilizer normalizes translations, so this construction and its averaging projection are equivariant on each root chart.

\item At a regular resonance, $a'=\kappa^t(q,r)^t$ and $(q,r)\ne(0,0)$. The period determinant calculation in Part E shows that $\kappa$ is invertible at every regular point. Thus $a$ has a simple zero; choosing its direction as the second cocycle frame and changing the vertical frame gives
\[
 a=(0,s),\qquad\kappa=\begin{pmatrix}0&1\\1&0\end{pmatrix}.
\]
In degree one the two generators $e_1G_2,e_2G_2$ are killed by $s$. In degree two the relations eliminate the two vertical generators and leave $hG_1G_2$ killed by $s^2$. This gives the regular modules in \eqref{cdp:eq:C.12}.

\item At the order-three point put $t=s^3$, $\nu=\mu-(2-\tau)/3$ and $b=\beta_c+4\mu+(2/3)\tau-b(0)$. The order-three characters give $\nu=\nu_2s^2+O(s^3)$ and $b=b_3s^3+O(s^4)$. We check the joint coefficient that controls the module. With elliptic scale $f(s)=Fs^2+O(s^3)$, the Weierstrass equation gives
\[
 \nu_2^2=9F^2/512,\quad \tau_1=i\pi F^2/2,\quad
 6\mu'^2-\tau'\beta_c'=(3F^2/4)s^2+O(s^3).
\]
The first equality follows from the exact tripling coordinate $x(3P)(0)=512/81$ and $x=f^2\wp$: $3\nu$ approaches zero, so $x(3P)(0)=F^2/(9\nu_2^2)$. The determinant equality follows from the identity in Part E and $f^4g_2=12t^3(t-1)$, $f^6g_3=-8t^4(t-1)^2$. Comparing coefficients in $6\nu'^2-\tau'b'$ gives $b_3=7i/(32\pi)$. For a nonintegral resonance the Picard section is $3k(\nu,b)$ in the corresponding frames. Eliminating the unit $\tau'$ leaves transverse coefficient
\[
 3k(b-6\nu\nu'/\tau')
   =\frac{123ik}{64\pi}s^3+O(s^4),\qquad k\ne0.        \tag{C.14}\hypertarget{tag-008S}{}\label{cdp:eq:C.14}
\]
It is nonzero. Equivariant holomorphic frame operations therefore reduce \eqref{cdp:eq:C.13} to
\[
 a=(0,s^2),\qquad\kappa=\begin{pmatrix}0&s\\s&1\end{pmatrix}.
\]
Its degree-one and degree-two modules are both $\C\{s\}/(s^2)\oplus\C\{s\}/(s^3)$. Explicit degree-one cycles are $u=e_1G_2$, $v=e_2G_1+shG_2$; the third cycle is removed by the boundary $su+e_2G_2+sv$. The remaining relations are $s^2u=s^3v=0$. In degree two elimination of $e_2G_1G_2$ leaves the generators $e_1G_1G_2,hG_1G_2$ with annihilators $s^2,s^3$.

\item We retain the actual affine characters. The integral cocycle at the order-three resonance is $(n,-2k,-k,-3k)$; evaluating on $v_1$ gives $n$. At the order-four resonance it is $(n,-k,-k,-2k)$; evaluation on $v_2$ gives $-n\equiv3n\pmod4$. The invariant exponents at the order-three point are $n+2,n+2$ for the two degree-one generators and $n+2,n$ in degree two. Here an invariant exponent is the required Taylor exponent modulo the stabilizer order. Thus the length-two summand retains precisely the monomials $s^j$ with $0\le j<2$ and $j\equiv n+2\pmod3$: its invariant lengths for $n\equiv0,1,2$ are $0,1,1$. A length-three summand retains one coefficient in every residue class. Both coarse lengths are therefore $1+\epsilon_n$, and multiplication by $t=s^3$ kills them. At the order-four point the regular normal form applies. The two length-one degree-one generators have invariant exponents $3n+1,3n+2$, and the length-two degree-two generator exponent $3n+2$. Their invariant length is one exactly for $n\equiv1,2\pmod4$, and the coarse parameter $s^4$ kills them. This proves all noncusp summands.

\item At the cusp, the restriction of $M_z$ has normalization pullback trivial and dual-complex character $(a,1)$, $a=e^{2\pi iz}$. Its normalization complex is
\[
 \C\xrightarrow{(a-1,0)^t}\C^2
 \xrightarrow{(0,a-1)}\C.
\]
For $z\notin\Z$ it is acyclic. Semicontinuity and base change make $R\bar f_*M_z$ vanish on a cusp neighbourhood. The logarithmic sequence (C.8) then gives $R\bar f_*(\mathcal T\otimes M_z)=0$ there. Sequence (C.3) contributes exactly the reduced skyscraper $\C_{p_0}$ in degree two for the ordinary tangent bundle. Away from resonances a nontrivial topologically trivial line on a torus has zero cohomology: its unitary Fourier presentation has no zero eigenmode; the tangent extension by trivial bundles retains that acyclicity. We verify that this analytic support is proper. Fix a simply connected disc in the regular base. For each fixed integer tuple $(n,p,q,r)$, a resonance has $r\ne0$ and $a^\prime=\kappa^t(q,r)^t\ne0$, so its zeros are isolated. Each such zero set is countable on this disc, and the union over the countably many tuples is countable. The disc has a point outside this union. At that fibre the preceding line-bundle acyclicity and tangent extension give zero cohomology in every degree. Upper semicontinuity gives a neighbourhood with the same vanishing. Each coherent higher direct image therefore has proper analytic support on the compact base curve, hence finite support. Every regular resonance lies in that support by its computed nonzero local modules, proving that $\mathcal R_z$ is finite. Degree zero is zero because a direct image of a bundle is torsion-free and has generic rank zero. These arguments prove \eqref{cdp:eq:C.12}, and Leray proves \eqref{cdp:eq:C.11}. The logarithmic degree-two direct image is a torsion sheaf for these twists, so its first cohomology vanishes. Equation (C.4) proves the final surjectivity assertion.

\end{enumerate}
\end{proof}



Statements \ref{cdp:st:C5b}--\ref{cdp:st:C5c} reproduce, with the displayed comparisons and independent exact module computations, the groups and stalk lengths claimed in S6\_01 D.1--D.4. The extra conclusion (C.7) corrects the assertion that the map is nonzero for every twist: the target is always one-dimensional, but its global map is sometimes zero. All formulas give $\chi(T_X\otimes M_z)=1$. The generic tuple is $(0,0,1,0)$; it directly contradicts the degree-two vanishing in CDP20 Theorem 2.2(a), while its Euler characteristic and $c_3(X)=2$ contradict parts (b) and (c). A holomorphic universal Picard family is defined by $\bar\partial+z\alpha$ for any global closed Dolbeault representative $\alpha$ of the fixed generator. Proper coherent base change makes every cohomology threshold analytic. Real nonintegral parameters attain the generic tuple, so every higher threshold is a proper locally discrete analytic subset of the Picard line. This establishes the asserted jump-locus interpretation as well as the individual dimensions.



\section{Constant-monodromy period comparison and compact examples}\hypertarget{tag-0349}{}\label{sec:II:13}


\begin{proposition}[Components under trivial torus monodromy]\hypertarget{tag-01G8}{}\label{cdp:st:D.2.1}
\StatementRecord{cdp:st:D.2.1}{proved}{01G8}{\SeriesRef{cdp:st:D.1.2}, \SeriesRef{cdp:st:C.3.4}}

A smooth compact complex threefold with $b_1=b_2=0$ and a holomorphic
connected-fibre map to a curve, with general fibre a two-torus and
trivial rational monodromy, has reducible fibres. More precisely,
for every nonempty set of $s$ removed values containing all critical
values, their reduced fibres have $r=s+3$ irreducible components in
total. Every component of each reducible fibre has no compact
K\"ahler resolution. Consequently no such fibration can have all
its fibre components admitting compact K\"ahler resolutions.
\end{proposition}


\begin{proof}
\begin{enumerate}[label={$\langle1\rangle$\arabic*.},leftmargin=*]
\item 
Pullback on first rational cohomology from the base is injective because
the fibres are connected. Thus the base is $\PP^1$. Trivial monodromy on
the four-dimensional $H_1(F;\Q)$ gives $t'=4$, and the count in
Proposition~\ref{cdp:st:D.1.2} gives $r=s+3$. If all removed fibres were
irreducible one would have $r=s$, which is impossible. The assertion
about each component is Proposition~\ref{cdp:st:C.3.4}, applied to this
threefold. Such a reducible fibre prevents all fibre components from
admitting compact K\"ahler resolutions.

\end{enumerate}
\end{proof}






\begin{proposition}[Logarithmic period obstruction]\hypertarget{tag-01GA}{}\label{cdp:st:D.2.3}
\StatementRecord{cdp:st:D.2.3}{proved}{01GA}{\SeriesRef{cdp:eq:D.18}, \SeriesRef{cdp:eq:D.19}; \cite[\S2.1, p.~234; Theorem~2.18, pp.~237--238; \S2.19 and Proposition~2.20, pp.~238--239; Theorem~2.21, p.~239]{Steenbrink1976}; \cite[Proposition~1.13, pp.~233--234, and proof of Proposition~2.20, pp.~238--239]{Steenbrink1976}}


Let $h:U\to\Delta$ be a proper holomorphic map from a smooth
complex threefold, with connected fibres, smooth two-torus fibres
over $\Delta^*$, and central fibre $F=\sum_i m_iD_i$ with
normal-crossings reduced support $D=\bigcup_iD_i$.
If every normalization component of $D$ has geometric genus zero,
the monodromy on $H^2(U_t;\C)$ is nontrivial. In particular the
monodromy on $H^1(U_t;\Q)$ is nontrivial. The conclusion holds for every choice of positive integers $m_i$.
\end{proposition}


\begin{proof}

\begin{enumerate}[leftmargin=*]
\item Put
\[
\mathscr K^\bullet=
 \Omega_U^\bullet(\log D)\big/
 \bigl(h^*\Omega_\Delta^1(\log0)\wedge
       \Omega_U^{\bullet-1}(\log D)\bigr),\qquad
\mathscr H^2=\mathbb R^2h_*\mathscr K^\bullet.
\]
In \cite[\S2.1, p.~234; Theorem~2.18, pp.~237--238; \S2.19 and Proposition~2.20, pp.~238--239; Theorem~2.21, p.~239]{Steenbrink1976}, the assumptions for these assertions
are precisely a proper holomorphic map from a connected complex
manifold to a disc, smooth off zero, with normal-crossings support.
They give a locally free bundle $\mathscr H^2$,
its Gauss--Manin connection
\[
\nabla:\mathscr H^2\longrightarrow
\Omega_\Delta^1(\log0)\otimes\mathscr H^2,
\]
and the comparison
$\mathscr H^2|_{\Delta^*}\simeq
R^2h_*\C\otimes\OO_{\Delta^*}$.
The residue eigenvalues belong to $[0,1)$ and the extended
monodromy at zero is
$\exp(-2\pi i\operatorname{Res}_0\nabla)$.
Its canonical base-change map identifies the fibre at zero with
$\mathbb H^2(D,\mathscr K^\bullet/t\mathscr K^\bullet)$.
The local calculation on the coefficient ideal uses the monomial
Koszul complex of \cite[Proposition~1.13, pp.~233--234, and proof of Proposition~2.20, pp.~238--239]{Steenbrink1976}.

\item Suppose that the monodromy on $H^2$ were the identity. Its
residue has all eigenvalues zero, and its exponential is the
identity. A nilpotent matrix whose exponential is the identity
is zero. Hence the residue vanishes. In a holomorphic frame
$\nabla=d+A(t)\,dt/t$ with $A(0)=0$, so this is a holomorphic
connection across zero. It has a holomorphic horizontal frame
on the disc. Consequently the coordinates of every holomorphic
section of $\mathscr H^2$ in a flat frame, that is its periods,
are holomorphic and bounded near zero.

\item The second exterior power of the relative logarithmic cotangent bundle is
\[
\mathscr K^2\simeq
 \omega_U(D)\otimes h^*\omega_\Delta(0)^{-1}
 \simeq\omega_{U/\Delta}(D-F).
\]
The direct image $h_*\mathscr K^2$ is a torsion-free
coherent sheaf of rank one, because a smooth torus fibre has
one-dimensional space of holomorphic two-forms. Shrink the disc
and choose a generator $\Omega$ of this line bundle.
Thus $\Omega$ determines a holomorphic section $[\Omega]$ of
$\mathscr H^2$ by the inclusion of the last term of the
relative de Rham complex. Its smooth-fibre comparison is the
ordinary class $[\Omega_t]\in H^2(U_t;\C)$.

\item First suppose that $\Omega$ restricts nontrivially to a normalized
component $T$ of $D$. There is a canonical isomorphism
\[
 \mathscr K^1|_T\simeq\Omega_T^1(\log D_T),\qquad
 \mathscr K^2|_T\simeq\omega_T(D_T),
\]
where $D_T$ is the reduced conductor. To specify the map, on the
branch $z_i=0$ subtract
$\operatorname{res}_i(\eta)(dt/t)/m_i$ from an absolute logarithmic
one-form $\eta$ and restrict its tangential part. The normal
residue of $dt/t$ is $m_i$, so this kills exactly the subbundle
divided out in $\mathscr K^1$ and is independent of the lift.
The exterior square gives the second isomorphism. If every
residue of $\Omega|_T$ vanished, its local logarithmic expression
would be holomorphic on $T$. This contradicts $p_g(T)=0$.
Thus there is a smooth point of a double curve, away from triple
points, where its residue is nonzero.

\item In a small chart about that point, write $h=x^ay^b=t$ and use $z$
along the double curve. The relative two-form is
\[
 \Omega_t=g(x,y,z)\,\frac{dx}{x}\wedge dz,
 \qquad g(0,0,z_0)\ne0.
\]
Choose a fixed small $z$-disc and $\epsilon>0$ such that
$|g(x,y,z)|\ge c>0$ whenever $|x|,|y|<\epsilon$ there.
The map to the annulus
$|t|^{1/a}\epsilon^{-b/a}<|x|<\epsilon$
has $b$ sheets. Integration on its inverse image and the fixed
$z$-disc gives
\[
 \int_{U_t}\Omega_t\wedge\overline{\Omega_t}
 \ \ge\ c'\log\frac{\epsilon^{a+b}}{|t|}
 \ \longrightarrow\ +\infty.
\]
Here the integrand is a nonnegative real volume form on the
complex surface; its local expression is a positive constant
times $|g|^2|x|^{-2}\,dA(x)dA(z)$.

\item On the other hand, $\Omega_t$ is closed, and this integral is
the fixed cup-product pairing of $[\Omega_t]$ with its complex
conjugate. All its periods are bounded by the holomorphic
connection argument, so the pairing is bounded. This is a
contradiction.

\item Suppose that $\Omega$ vanishes on every
branch of $D$. Put
\[
 \mathscr J=\OO_U(-D),\qquad
 \mathscr B^\bullet=
 \mathscr J\mathscr K^\bullet/t\mathscr K^\bullet.
\]
Both $\mathscr J\Omega_U^\bullet(\log D)$ and
$t\Omega_U^\bullet(\log D)$ are subcomplexes for the absolute
exterior derivative: locally the logarithmic derivatives of
$z_1\cdots z_r$ and $t=\prod z_i^{m_i}$ are holomorphic
linear combinations of the $d\log z_i$. Hence the locally split
absolute/relative logarithmic exact sequence remains exact
when its coefficients are restricted to $\mathscr J/(t)$:
\begin{equation}
0\longrightarrow\mathscr B^\bullet[-1]
 \xrightarrow{\ (dt/t)\wedge\ }
 (\mathscr J/(t))\Omega_U^\bullet(\log D)
 \longrightarrow\mathscr B^\bullet\longrightarrow0.
\hypertarget{tag-0195}{}\label{cdp:eq:D.18}
\end{equation}
Signs in the shifted differential are chosen so that the first
arrow is a morphism of complexes. Its connecting morphism
$\rho_{\mathscr B}:\mathscr B^\bullet\to\mathscr B^\bullet$
in the derived category defines degree-preserving endomorphisms
on both cohomology sheaves and hypercohomology. The inclusion
$\mathscr J/(t)\hookrightarrow\OO_U/(t)$ gives a commutative
diagram from \eqref{cdp:eq:D.18} to the analogous exact sequence for
$\mathscr K^\bullet/t\mathscr K^\bullet$. The induced
endomorphism for the latter is precisely the Gauss--Manin
residue. Thus the map on hypercohomology intertwines these
endomorphisms.

\item At a point of $r$ branches,
write
\[
 t=\prod_{i=1}^r z_i^{m_i},\qquad
 e=\gcd(m_1,\ldots,m_r),\qquad
 y=\prod_{i=1}^r z_i^{m_i/e},\qquad t=y^e.
\]
The transverse holomorphic Poincar\'e lemma reduces the relative
complex to the Koszul complex with commuting operators
\[
 D_i=z_i\partial_{z_i}
           -\frac{m_i}{m_1}z_1\partial_{z_1},
 \qquad 2\le i\le r,
\]
on the convergent-power-series module
$(z_1\cdots z_r)/(z_1^{m_1}\cdots z_r^{m_r})$.
The monomials diagonalize these operators. A monomial with some
nonzero eigenvalue has a contracting Koszul homotopy, obtained
by division by that eigenvalue. These divisions preserve
convergence: every nonzero eigenvalue has absolute value at
least $1/m_1$. The surviving monomials have exponent vector
$a(m_1/e,\ldots,m_r/e)$. Membership in $\mathscr J$ requires
$a\ge1$, and quotienting by $t$ requires $a<e$. Consequently
each cohomology stalk is spanned by
\[
 y^a\,d\log z_{i_1}\wedge\cdots\wedge d\log z_{i_q},
 \qquad 1\le a<e,
\]
subject to the relative relation $\sum m_i d\log z_i=0$.
Since $dy^a=(a/e)y^a\,dt/t$, the connecting endomorphism
acts on these representatives by $a/e$. In particular zero is
not an eigenvalue; when $e=1$ the complex is acyclic.

\item The hypercohomology spectral sequence
\[
 H^p\bigl(D,\mathcal H^q(\mathscr B^\bullet)\bigr)
 \ \Longrightarrow\ \mathbb H^{p+q}(D,\mathscr B^\bullet)
\]
is equivariant for this derived endomorphism. There are only
finitely many branch multiplicities, so all eigenvalues on its
terms and abutment lie in a finite subset of $(0,1)$. For
example the product of the corresponding linear factors
annihilates every cohomology sheaf, and a power annihilates the
abutment. Hence $\rho_{\mathscr B}$ is invertible on
$\mathbb H^2(D,\mathscr B^\bullet)$. The residue on
$\mathscr H^2|_0$ is zero. Equivariance therefore proves that
\begin{equation}
 \mathbb H^2(D,\mathscr B^\bullet)\longrightarrow
 \mathbb H^2(D,\mathscr K^\bullet/t\mathscr K^\bullet)
       =\mathscr H^2|_0
 \quad\hbox{is the zero map}.
\hypertarget{tag-0196}{}\label{cdp:eq:D.19}
\end{equation}

\item The central class of $\Omega$ comes from the left side of
\eqref{cdp:eq:D.19}, because $\Omega\in\mathscr J\mathscr K^2$.
Thus $[\Omega]$ vanishes at zero in the locally free bundle
$\mathscr H^2$, and $[\Omega]/t$ is a holomorphic section of
that bundle. Its periods are bounded. On a smooth fibre this
section is the de Rham class of $\Omega_t/t$.

\item Primitivity of the generator $\Omega$ implies that its order
$b_i$ along at least one $D_i$ is less than $m_i$; otherwise
$\Omega/t$ would be a holomorphic section of $\mathscr K^2$
over $U$, contradicting its choice as a generator of
$h_*\mathscr K^2$. Choose a smooth point of that component
away from all other branches and all further zeros. There
are coordinates with $t=z^{m_i}$ in which
\[
 \Omega=z^{b_i}a(z,\xi,\eta)\,d\xi\wedge d\eta,
 \qquad a(0,0,0)\ne0.
\]
On one sheet above a fixed small $(\xi,\eta)$-polydisc, the
positive integral of $\Omega_t/t$ against its conjugate has
the lower bound
\[
 \int_{U_t}\frac{\Omega_t}{t}\wedge
                    \overline{\frac{\Omega_t}{t}}
 \ \ge c''|t|^{2(b_i/m_i-1)}\longrightarrow+\infty.
\]
Its expression as the fixed cup-product pairing of bounded
periods again gives a contradiction.

\item The natural
cup-product isomorphism
$\bigwedge^2H^1(U_t;\Q)\simeq H^2(U_t;\Q)$ for a two-torus
is monodromy-equivariant. Identity on first cohomology would
therefore imply identity on second cohomology.

\end{enumerate}
\end{proof}






\begin{proposition}[Euler number in the normal-crossings class]\hypertarget{tag-01GB}{}\label{cdp:st:D.2.4}
\StatementRecord{cdp:st:D.2.4}{proved}{01GB}{\SeriesRef{cdp:st:C.2.NC.5}, \SeriesRef{cdp:st:C.2.NC.3}, \SeriesRef{cdp:st:D.2.3}, \SeriesRef{cdp:st:D.2.7}}


For the threefolds and maps in \ref{cdp:st:D.2.1} whose reduced fibre supports
all have normal crossings, the Euler number is
\[
c_3(Y)=0,\qquad b_3(Y)=2.
\]
All multiplicities, including unequal multiplicities on a
reducible support, are allowed.
\end{proposition}


\begin{proof}

\begin{enumerate}[leftmargin=*]
\item Statement \ref{cdp:st:C.2.NC.5} gives $c_3(Y)=2r$, where $r$ counts the
irreducible fibre supports with rational normalization. A fibre
of that kind is necessarily $mS$ for one positive integer $m$.
Choose a small disc around its value and perform the normalized
base change $t=u^m$. Statement \ref{cdp:st:C.2.NC.3} identifies the resulting
map to the original neighbourhood as finite \emph{\'etale} of
degree $m$, and its central fibre $S'$ as reduced with normal
crossings. The total space is therefore smooth. Its nearby
fibres are the original connected tori. Stein factorization
over the normal disc shows that its central fibre is connected
as well: the finite Stein factor has generic degree one and is
the disc itself.

\item Normalization commutes with this \emph{\'etale} base change.
Each normalized component of $S'$ is an unramified finite cover
of the smooth rational normalization $T$ of $S$. Since $T$ is
simply connected, these normalized components are copies of
$T$, and all have geometric genus zero. The local monodromy
after base change is the $m$th power of the original monodromy,
so remains the identity on $H^1$. This contradicts \ref{cdp:st:D.2.3}.
Therefore $r=0$, proving $c_3(Y)=0$. Poincar\'e duality and
$b_1(Y)=b_2(Y)=0$ give $c_3(Y)=2-b_3(Y)$, hence $b_3(Y)=2$.
The example of \ref{cdp:st:D.2.7} realizes this normal-crossings class.

\end{enumerate}
\end{proof}






\begin{proposition}[Proved unrestricted sign under trivial monodromy]\hypertarget{tag-01GC}{}\label{cdp:st:D.2.5}
\StatementRecord{cdp:st:D.2.5}{proved}{01GC}{\SeriesRef{cdp:st:C4c}, \SeriesRef{cdp:st:A.3.5}, \SeriesRef{cdp:eq:D.S.2}, \SeriesRef{cdp:st:D.1.9}, \SeriesRef{cdp:st:A.4.7}, \SeriesRef{cdp:st:A.5.2}, \SeriesRef{cdp:eq:D.S.4}, \SeriesRef{cdp:st:A.5.6}, \SeriesRef{cdp:st:A.5.3}, \SeriesRef{cdp:eq:D.S.5}, \SeriesRef{cdp:st:D.2.3}; \cite[Theorem~1.6]{BierstoneMilman1997}}


Let $f:Y\to\PP^1$ be a proper surjective holomorphic map with
connected fibres from a compact connected complex threefold with
$b_1(Y)=b_2(Y)=0$. Suppose its smooth general fibre is a complex
two-torus and its monodromy representation on $H_1(F,\Q)$ is
trivial. Then
\[
 c_3(Y)\le0,\qquad b_3(Y)\ge2.
 \tag{D.S.1}\hypertarget{tag-019O}{}\label{cdp:eq:D.S.1}
\]
The fibres may have arbitrary singularities and arbitrary
multiplicities. In particular this proves the sign assertion of
D.2 for a holomorphic algebraic reduction; no existence assertion
is required for the conclusion.
\end{proposition}


\begin{proof}

\begin{enumerate}[label={$\langle1\rangle$\arabic*.},leftmargin=*]
\item Assume $c_3(Y)>0$. Poincar\'e duality gives
$c_3(Y)=2-b_3(Y)$, and the alternating nondegenerate intersection
pairing on $H^3(Y,\Q)$ makes $b_3(Y)$ even. Hence $b_3(Y)=0$.
Apply \ref{cdp:st:C4c} to the Picard component
$P=K_Y\otimes\Pic^0(Y)$. It gives a fixed critical full fibre
$F$ such that
\[
 H^0\bigl(F,\Omega_Y^1|_F\otimes M|_F\bigr)\ne0
 \qquad(M\in\Pic^0(Y)).
 \tag{D.S.2}\hypertarget{tag-019P}{}\label{cdp:eq:D.S.2}
\]
Indeed the Serre-dual parameter $K_Y\otimes L^{-1}$ runs through
$\Pic^0(Y)$ when $L$ runs through $P$.

\item Statement \ref{cdp:st:A.3.5} excludes reducible support in \eqref{cdp:eq:D.S.2}, including
all unequal-multiplicity possibilities. Thus $F=mS$ for one
irreducible reduced Cartier divisor $S$ and one positive integer
$m$. We specify why a smooth $S$ is excluded under the present
\emph{rational} Betti-number assumptions. \ref{cdp:st:D.1.9} supplies
$\beta\in H^1(Y,\OO_Y)$ whose restriction near every full fibre
comes from a nonzero integral transgression class. The line
$M_r=\exp(2\pi i r\beta)$, for real irrational $r$, has that
unitary local realization. The normalized root change $t=s^m$
near a smooth $S$ is the finite \emph{\'etale} map of A.3.4:
in a chart $t=uz^m$ it has equations $w^m=u$, $s=wz$.
Its central reduction $S'$ is smooth, and its nearby fibres are
the original tori. Proper smooth local triviality transports
the vanishing of degree-zero and degree-one cohomology of the
pulled-back unitary local system from a nearby torus to $S'$.
The injectivity \ref{cdp:st:A.4.7} and the conormal sequence then give zero
ambient cotangent sections on $S'$. Each layer of $mS'$ has the
same twist, since $s$ trivializes its conormal. The finite
filtration and exact invariants functor for the cyclic cover
give $H^0(mS,\Omega_Y^1|_{mS}\otimes M_r)=0$, contradicting
\eqref{cdp:eq:D.S.2}. This proof uses only the rational version of \ref{cdp:st:D.1.9}.
Therefore $S$ is singular. \ref{cdp:st:A.5.2}, applied using
$\omega_S=(K_Y\otimes\OO_Y(S))|_S$ and
$\chi(\OO_S)=0$, makes $S$ nonnormal.

\item The ideal filtration of $mS$ has the finitely many locally free
graded factors on $S$
\[
 E_j=(\Omega_Y^1\otimes\OO_Y(-jS))|_S,
 \qquad 0\le j<m.
 \tag{D.S.3}\hypertarget{tag-019Q}{}\label{cdp:eq:D.S.3}
\]
If all the groups $H^0(S,E_j\otimes M|_S)$ were zero, the exact
filtration sequences would make \eqref{cdp:eq:D.S.2} zero. Thus their
nonvanishing loci cover $\Pic^0(Y)$. Each locus is closed
complex analytic by the actual Picard family and proper
semicontinuity. The group $\Pic^0(Y)$ is a connected complex
vector space, since $H^1(Y,\Z)=0$. A finite union of proper
closed analytic subsets cannot cover it. Some fixed $j$
therefore satisfies
\[
 H^0(S,E_j\otimes M|_S)\ne0
 \qquad(M\in\Pic^0(Y)).
 \tag{D.S.4}\hypertarget{tag-019R}{}\label{cdp:eq:D.S.4}
\]

\item Let $\eta:T\to S$ be the finite normalization and
$\pi:Z\to T$ its minimal resolution, and put $h=\eta\pi$,
viewed also as a holomorphic map to $Y$.
Pullback of sections in \eqref{cdp:eq:D.S.4} is injective: $S$ is reduced
and irreducible, and a section of the locally free sheaf $E_j$
vanishing on its dense smooth locus is zero. Thus the fixed
vector bundle $h^*(\Omega_Y^1\otimes\OO_Y(-jS))$ has a
section after every twist by $h^*M$, $M\in\Pic^0(Y)$.
For any Gauduchon metric $\omega$ on $Z$, the uniform
negative-degree bound \ref{cdp:st:A.5.6}, applied to all positive and
negative powers of $M$, gives
\[
 \deg_\omega h^*M=0\qquad(M\in\Pic^0(Y)).
\]
Every line bundle $G$ on $Y$ has a positive power in
$\Pic^0(Y)$, because $H^2(Y,\Z)$ is finite. Additivity of
degree therefore gives $\deg_\omega h^*G=0$ for every such
$G$. No classification of the surface $Z$ is used in this step.

\item The conductor comparison \ref{cdp:st:A.5.3} gives an effective divisor
$D_Z$ with
\[
 \omega_Z(D_Z)\simeq h^*(K_Y\otimes\OO_Y(S)).
 \tag{D.S.5}\hypertarget{tag-019S}{}\label{cdp:eq:D.S.5}
\]
Moreover $D_Z\ne0$. Indeed a reduced hypersurface is
Cohen--Macaulay and hence satisfies Serre's condition $S_2$;
its nonnormality must therefore occur in codimension one.
The divisorial conductor is nonzero, and its strict transform
is part of $D_Z$.
Taking Gauduchon degree in \eqref{cdp:eq:D.S.5} yields
\[
 \deg_\omega K_Z=-\deg_\omega\OO_Z(D_Z)<0.
\]
A nonzero holomorphic section of $K_Z$ would have an effective
zero divisor, whose degree is nonnegative. This contradiction
proves $p_g(Z)=h^0(Z,K_Z)=0$.

\item Choose a small closed fibre neighbourhood $U=f^{-1}(\Delta)$.
Embedded resolution of the divisor $F$ gives a proper
holomorphic modification
\[
 q:\widehat U\longrightarrow U
\]
which is a finite composite of blowups with smooth centres in
the successive central divisors, is an isomorphism over
$\Delta^*$, has smooth total space, and makes the reduced
support of $(f\circ q)^{-1}(0)$ a simple normal-crossings
divisor. We use the analytic embedded-resolution theorem in
\cite[Theorem~1.6]{BierstoneMilman1997}; properness of the
central fibre makes its locally finite construction finite
after shrinking $\Delta$. The strict transform of $S$ is a
resolution of its normalization and has geometric genus zero
by birational invariance of holomorphic two-forms on smooth
compact surfaces. A new exceptional component is $\PP^2$ over
a point or a $\PP^1$-bundle over a smooth compact curve.
Both have geometric genus zero: restriction of a holomorphic
two-form to a ruling chart has negative fibre degree, and
$H^0(\PP^2,K_{\PP^2})=0$. Further strict transforms preserve
this property. Consequently every component of the final
central reduced divisor has geometric genus zero.

\item The map $f\circ q$ remains proper with connected fibres:
a blowup has connected fibres and its structure-sheaf direct
image is the original structure sheaf. Its smooth fibres over
$\Delta^*$ are the original two-tori under the isomorphism
$q|_{\Delta^*}$, so their $H^1$ monodromy remains the identity.
The multiplicities of the resolved central fibre are arbitrary
positive integers and are all permitted in \ref{cdp:st:D.2.3}. That
proposition now gives nontrivial monodromy on $H^2$ and hence
on $H^1$, a contradiction. This proves \eqref{cdp:eq:D.S.1}.

\end{enumerate}
\end{proof}



The sign result uses the analytic logarithmic comparison of
Steenbrink. The explicit compact construction in \ref{cdp:st:D.2.7} proves
that the stated class is nonempty and attains $c_3=0$.
Together these results settle both the existence and the sign
questions in D.2.




\begin{proposition}[Extension of the marked torus period]\hypertarget{tag-01GD}{}\label{cdp:st:D.2.6}
\StatementRecord{cdp:st:D.2.6}{proved}{01GD}{\SeriesRef{cdp:st:D.2.3}, \SeriesRef{cdp:eq:D.20}}


Let $h:U\to\Delta$ be a proper holomorphic map from a smooth
complex threefold, smooth with connected two-torus fibres over
$\Delta^*$. Suppose that its integral $H^1$ monodromy is the
identity. The marked period map
\[
 \Delta^*\longrightarrow
 \mathcal T_2
 :=\{W\in\operatorname{Gr}(2,H^1(U_t;\C)):
                   W\oplus\overline W=H^1(U_t;\C)\}
\]
extends holomorphically across zero with value in $\mathcal T_2$.
In particular the relative Albanese group family over $\Delta^*$
extends to a smooth proper group family of complex two-tori over
$\Delta$. 
\end{proposition}


\begin{proof}

\begin{enumerate}[leftmargin=*]
\item The punctured family is connected by local triviality and connectedness
of its fibres, so it lies in one open and closed component $U^0$ of $U$.
The inclusion $U^0\hookrightarrow U$ is an isomorphism over $\Delta^*$.
Apply analytic embedded resolution to the compact reduced central divisor
in $U^0$ \cite[\S0.2(2), preprint p.~3; Theorem~1.6 and Remarks~1.7,
preprint pp.~5--6]{BierstoneMilman1997}. This gives a proper composite
of blowups $\pi:\widetilde U\to U^0$ with smooth centres over that
divisor, with $\widetilde U$ smooth and the reduced central support
of $\widetilde h:=h|_{U^0}\circ\pi$ normal crossing. For $t\ne0$,
$\pi_t:\widetilde U_t\to U_t$ is a biholomorphism;
its pullback identifies holomorphic forms, the integral marking and
cup products. Transporting through these pullbacks, use the relative
logarithmic complex in the proof of Proposition~\ref{cdp:st:D.2.3}
on this model, with its full central multiplicities.
Identity on $H^1$ gives identity on $H^2=\bigwedge^2H^1$.
Consequently its logarithmic Gauss--Manin bundle has zero
residue and a holomorphic flat frame. Let $\Omega$ be a
generator of $\widetilde h_*\mathscr K^2$, as in Proposition~\ref{cdp:st:D.2.3}.
For $t\ne0$, let $\Omega_t$ be the form on $U_t$ determined by
$\pi_t^*\Omega_t=\Omega|_{\widetilde U_t}$.

\item The coefficient-ideal argument in
\eqref{cdp:eq:D.18}--\eqref{cdp:eq:D.19} shows that $\Omega$
cannot vanish on the whole reduced central support: otherwise
its divided cohomology class would have bounded periods while
its positive $L^2$ volume diverges. Thus $\Omega|_T$ is nonzero
on some normalized component $T$. The annular volume estimate
in the same proof makes every conductor residue of every
component restriction zero: a nonzero residue would give a
divergent positive volume while its cup-product expression
stays bounded. The nonzero restriction $\Omega|_T$ is therefore
a holomorphic two-form on $T$.

\item Choose a point of $T$ outside the conductor and the zeros of
this two-form. In a small chart the map is $t=z^m$ and
\[
 \Omega=a(z,\xi,\eta)\,d\xi\wedge d\eta,
 \qquad a(0,0,0)\ne0.
\]
Integrating on one sheet over a fixed small
$(\xi,\eta)$-polydisc yields a constant $c>0$ such that
\begin{equation}
 \int_{U_t}\Omega_t\wedge\overline{\Omega_t}\ge c
 \qquad(0<|t|\ll1).
\hypertarget{tag-0197}{}\label{cdp:eq:D.20}
\end{equation}
The cohomology class $v(t)=[\Omega_t]$, expressed in the flat
frame, is holomorphic across zero. Its cup pairing with its conjugate equals the
integral in \eqref{cdp:eq:D.20}, hence
\[
 v(0)\wedge\overline{v(0)}>0
 \quad\hbox{in }H^4(U_t;\R)\simeq\R,
\]
using the complex orientation. In particular $v(0)\ne0$.

\item Write $V=H^1(U_t;\C)$ with its fixed integral marking.
For $t\ne0$, $v(t)$ is a decomposable element of
$\bigwedge^2V$, spanning $\bigwedge^2H^{1,0}(U_t)$.
The decomposable cone is the closed algebraic set defined by
the Pl\"ucker equations, so $v(0)$ is decomposable. For a
nonzero decomposable $v$, the exact inverse to the Pl\"ucker
map is
\[
 W(v)=\{\alpha\in V:\alpha\wedge v=0\}.
\]
This kernel has dimension two and varies holomorphically near
$v(0)$. The nonzero pairing
$v(0)\wedge\overline{v(0)}$ says exactly that
$W(v(0))\cap\overline{W(v(0))}=0$. The extended period
therefore lies in $\mathcal T_2$.

\item Let $\mathcal W\subset V\otimes\OO_\Delta$ be the extended
rank-two subbundle $W(v(t))$, and let
$\Lambda=H_1(U_t;\Z)$ with its fixed marking. Evaluation of
cohomology on integral cycles gives a holomorphic embedding
of the constant lattice $\Lambda$ in $\mathcal W^*$.
A real cycle annihilating $\mathcal W_t$ also annihilates
$\overline{\mathcal W_t}$, hence all of $V$; it is therefore zero.
Equality of real dimensions makes the real-linear evaluation
map an isomorphism, so the image is a full lattice in every fibre,
uniformly on compact subsets of the disc. Translation by this
lattice is free and properly discontinuous. For a fixed integral
basis, the closed parallelotopes generated by its images cover
each quotient fibre; over a compact subset of the base their
union is compact. Consequently
\[
 \mathcal A=\mathcal W^*/\Lambda\longrightarrow\Delta
\]
is a smooth proper family of complex Lie groups with
two-torus fibres. On $\Delta^*$, integration of holomorphic
one-forms identifies each group fibre with the Albanese torus
of $U_t$, and these identifications vary holomorphically.


\end{enumerate}
\end{proof}



\paragraph{Scope of \ref{cdp:st:D.2.6}.}
For a fibration over $\PP^1$ with globally trivial marked
$H^1$ monodromy, the local constructions glue to a smooth
proper torus group family over $\PP^1$. The period map may
be nonconstant, as the explicit family in \ref{cdp:st:D.1.15} demonstrates.
The local filling in \ref{cdp:st:D.2.7} proves that the marked Albanese
map need not extend to the given model. Its inclusion of
a nearby fibre kills the nonzero cycle $w$ in (D.25).
An extension $U\to A\times\Delta$ of the marked fibre
isomorphisms would induce a map on $H_1$ whose composite
with that inclusion is the identity on $H_1(A;\Z)$,
which is impossible because it would send $w$ both to
zero and to $w$. This exact topological obstruction
explains why the extending group family does not imply
nonexistence in D.2.




\begin{proposition}[Proved existence]\hypertarget{tag-01GE}{}\label{cdp:st:D.2.7}
\StatementRecord{cdp:st:D.2.7}{proved}{01GE}{\SeriesRef{cdp:st:D.2.5}, \SeriesRef{cdp:eq:D.24}, \SeriesRef{cdp:eq:D.25}, \SeriesRef{cdp:eq:D.26}, \SeriesRef{cdp:eq:D.23}, \SeriesRef{cdp:eq:D.28}, \SeriesRef{cdp:eq:D.30}, \SeriesRef{cdp:eq:D.32}, \SeriesRef{cdp:eq:D.31}}


There exists a smooth compact connected complex threefold $Y$
with a holomorphic map $f:Y\to\PP^1$ with connected fibres such
that
\[
 H_1(Y;\Z)=H_2(Y;\Z)=0,\qquad
 a(Y)=1,\qquad c_3(Y)=0,\qquad b_3(Y)=2.
\tag{D.21}\hypertarget{tag-0198}{}\label{cdp:eq:D.21}
\]
Every smooth fibre is the fixed abelian surface
$A=E\times E$, where $E=\C/(\Z+\mathrm i\Z)$, and the marked
integral monodromy on $H_1(A;\Z)$ is the identity. There are
exactly three singular fibres. Each is reduced with simple
normal crossings, and consists of a primary Kodaira surface
and a primary Hopf surface meeting along one smooth elliptic
curve. Consequently the trivial-monodromy class in D.2 is
nonempty, and equality in the unrestricted sign theorem
\ref{cdp:st:D.2.5} is attained.
\end{proposition}


\begin{proof}

\begin{enumerate}[label={$\langle1\rangle$\arabic*.},leftmargin=*]
\item We give the local holomorphic filling first, including its
boundary homology maps. Put $q=e^{-2\pi}$. The map
\[
 \C/(\Z+\mathrm i\Z)\longrightarrow\C^*/q^{\Z},
 \qquad z\longmapsto e^{2\pi\mathrm i z},
\tag{D.22}\hypertarget{tag-0199}{}\label{cdp:eq:D.22}
\]
is a holomorphic group isomorphism: its first lattice generator
acts trivially and its second multiplies the exponential by
$q$. Denote these oriented integral generators by $(w,u)$,
respectively.

\item Let $C$ and $B$ be copies of $E$, with a specified product
identification $A=C\times B$. Blow up $B\times\Delta$ at
$(0,0)$ and write
\[
 b:V=\operatorname{Bl}_{(0,0)}(B\times\Delta)\longrightarrow\Delta,
 \qquad R\subset V\quad\hbox{for the exceptional curve},
 \qquad L=\OO_V(R).
\]
The notation $L^\times$ means the total space with its zero
section removed. Define the complex threefold and its map by
\[
 U=L^\times/q^{\Z}\xrightarrow{\ p\ }V\xrightarrow{\ b\ }\Delta,
\tag{D.23}\hypertarget{tag-019A}{}\label{cdp:eq:D.23}
\]
where $q^{\Z}$ acts by fibrewise scalar multiplication. This
action is free and properly discontinuous. A local
trivialization of $L$ identifies $p$ with the product with
$\C^*/q^{\Z}=C$, so $p$ is a smooth proper principal elliptic
bundle and $U$ is smooth. Both $b$ and $p$ are proper, hence
$U\to\Delta$ is proper.

\item The canonical section of $\OO_V(R)$ is nonzero off $R$.
It gives an exact holomorphic product identification
\[
 U|_{\Delta^*}\simeq C\times B\times\Delta^*
                         =A\times\Delta^*.
\tag{D.24}\hypertarget{tag-019B}{}\label{cdp:eq:D.24}
\]
The central curve of $V$ is the union of the strict transform
$B_0\simeq B$ and $R\simeq\PP^1$, meeting at one point.
Their normal-crossings equation is $xy=t$. The inverse image
of $B_0$ is the principal elliptic bundle
$\OO_B(0)^\times/q^{\Z}$ of Chern degree one. Its canonical
bundle is trivial, since the vertical tangent is trivial and
the base is elliptic. The primitive Chern transgression
kills exactly one of the two fibre one-cycles, so its first
Betti number is three. This identifies it as a primary
Kodaira surface, directly by the usual definition as a
nontrivial principal elliptic bundle over an elliptic curve.
The inverse image of $R$ is
\[
 \OO_{\PP^1}(-1)^\times/q^{\Z}
       \simeq(\C^2\setminus\{0\})/\langle z\mapsto qz\rangle,
\]
a primary Hopf surface. The two surfaces meet in the single
elliptic fibre over $B_0\cap R$. Since $p$ is a submersion,
their union has local equation $xy=t$ in the smooth
threefold $U$, and both multiplicities are one.

\item Here are the integral homology maps in the marking \eqref{cdp:eq:D.24}.
Let $\Lambda=H_1(A;\Z)$ and let $[C]=w\wedge u$ in
$H_2(A;\Z)=\bigwedge^2\Lambda$. Then inclusion of a nearby
fibre induces the exact quotients
\[
 H_1(U;\Z)=\Lambda/\Z w,\qquad
 H_2(U;\Z)=\bigwedge^2\Lambda/\Z[C].
\tag{D.25}\hypertarget{tag-019C}{}\label{cdp:eq:D.25}
\]
To verify them, $V$ retracts onto $B_0\cup R$, which has
homotopy type $B\vee S^2$. The first Chern class of $L$
pairs to zero on a nearby $B_t$ and to $-1$ on $R$.
The transgression of the principal $C$-bundle is consequently
the primitive vector $-w$ on this sphere. In degree one it
kills exactly $w$. In degree two the integral Serre sequence
has surviving terms
\[
 E_\infty^{2,0}\simeq\Z,
 \qquad E_\infty^{1,1}\simeq\Z^4,
 \qquad E_\infty^{0,2}=0.
\]
Indeed $H^2(V;\Z)$ modulo the primitive Chern class is
$\Z$, the degree-three cohomology of $V$ is zero, and the
transgression on the exterior square of $H^1(C;\Z)$ is
injective and primitive. The analogous homology sequence,
or its integral dual with these free terms, shows that the
nearby fibre maps isomorphically onto the first two graded
terms. Its fibre class $[C]$ is zero. This proves precisely
the surjectivity and the kernel in \eqref{cdp:eq:D.25}, without a torsion
ambiguity.

\item Let $\mu$ be the positive base meridian in the boundary
$A\times S^1$. In this same trivialization its inclusion
maps satisfy
\[
 \mu\longmapsto0\quad\hbox{in }H_1(U;\Z),
 \qquad
 x\wedge\mu\longmapsto0\quad\hbox{in }H_2(U;\Z)
 \quad(x\in\Lambda).
\tag{D.26}\hypertarget{tag-019D}{}\label{cdp:eq:D.26}
\]
For a cycle in $C$, choose a fixed point of $B$ different
from the blowup centre and take its product with the base
disc. For a generator of $H_1(B;\Z)$, represent it by a loop
avoiding the blowup centre and take its product with the
disc. On these products the canonical section of $L$ is
nonzero, so these are actual bounding chains in the product
trivialization \eqref{cdp:eq:D.24}. The same fixed-point disc bounds
$\mu$. This proves \eqref{cdp:eq:D.26} with its specified marking.

\item We now specify the three local fillings. Write
\[
 e_1=(1,0),\quad e_2=(\mathrm i,0),\quad
 e_3=(0,1),\quad e_4=(0,\mathrm i)
\]
for the ordered basis of $\Lambda$. The chosen elliptic
embeddings $\phi_i:E\to A$, their first lattice cycles
$w_i=\phi_i(1)$, and second cycles $u_i=\phi_i(\mathrm i)$ are
\[
\begin{array}{c|c|c|c}
 i&\phi_i(z)&w_i&u_i\\ \hline
 1&(0,z)&e_3&e_4\\
 2&(\mathrm i z,\mathrm i z)&e_2+e_4&-e_1-e_3\\
 3&(z,\mathrm i z)&e_1+e_4&e_2-e_3.
\end{array}
\tag{D.27}\hypertarget{tag-019E}{}\label{cdp:eq:D.27}
\]
In row one take the first factor as complementary $B$; in
rows two and three take the second factor. The resulting
complex $2\times2$ change-of-coordinate matrices have
determinants $-1,\mathrm i,1$, respectively, units in
$\Z[\mathrm i]$. They therefore give exact product
isomorphisms $C_i\times B_i\simeq A$, including their
integral lattices. Apply \eqref{cdp:eq:D.23} in these three markings,
obtaining $U_1,U_2,U_3$.

\item For the initial smooth $A$-bundle over $\PP^1$, use the
coordinate $t$ near zero and glue the trivial bundle there
to the trivial bundle over $\PP^1\setminus\{0\}$ by
\[
 a_{\rm out}=a_{\rm in}+(\ell(t),0),
 \qquad \ell(t)=\frac{\log t}{2\pi\mathrm i}\pmod{\Z+\mathrm i\Z}.
\tag{D.28}\hypertarget{tag-019F}{}\label{cdp:eq:D.28}
\]
This is a single-valued holomorphic translation: changing
the logarithm adds an integer. It defines a smooth proper
holomorphic principal $A$-bundle $Y_0\to\PP^1$.
Its positive clutching loop has class $v_0=e_1$.
Take disjoint discs at $0,1,\infty$. On the latter two
discs the exterior trivialization in \eqref{cdp:eq:D.28} is already a
product trivialization. Remove these three fibres of $Y_0$
and glue in $U_1,U_2,U_3$ using exactly \eqref{cdp:eq:D.24} and the
original local trivializations of $Y_0$. All gluing maps
are holomorphic isomorphisms over punctured discs. The
resulting $f:Y\to\PP^1$ is a smooth complex threefold with
the claimed three central fibres. It is Hausdorff and
proper: each disc neighbourhood is the proper model $U_i$,
and away from their centres the gluing is an isomorphism
to $Y_0$. A finite such cover of the compact base proves
compactness. Connectedness of all its fibres follows
from the product fibres and the described central unions.
All transition maps over the smooth base are translations
of the marked $A$; hence their induced integral $H_1$
monodromy is the identity.

\item We compute the global integral homology with these exact
clutchings. Let $M$ be the part above the sphere with the
interiors of the three small discs removed. It is the
product with $A$. Positive meridians $\mu_i$ about the
three deleted discs satisfy
$\mu_1+\mu_2+\mu_3=0$ in $H_1$ of that base. In the
specified local coordinates the boundary inclusion maps
to $M$ are
\[
 x\longmapsto x\quad(x\in\Lambda),\qquad
 \mu_i^{\rm loc}\longmapsto\mu_i+c_i,
 \qquad(c_1,c_2,c_3)=(e_1,0,0).
\tag{D.29}\hypertarget{tag-019G}{}\label{cdp:eq:D.29}
\]
This sign follows directly from \eqref{cdp:eq:D.28}: a positive loop
of $t$ increases $\ell(t)$ by $+1$. On mixed
two-cycles the same map is
\[
 x\wedge\mu_i^{\rm loc}
       \longmapsto x\wedge\mu_i+x\wedge c_i.
\tag{D.30}\hypertarget{tag-019H}{}\label{cdp:eq:D.30}
\]
This formula can also be obtained by the induced map on
$H_2$ of the product boundary torus; a translation loop
of class $c_i$ replaces its last degree-one generator
by $\mu_i+c_i$.

\item Use the signed Mayer--Vietoris map
$(j_{M*},-j_{U_i*})$ for each boundary. Its degree-one
cokernel is
\[
 H_1(Y;\Z)=
 \Lambda/\langle e_1,w_1,w_2,w_3\rangle=0,
\tag{D.31}\hypertarget{tag-019I}{}\label{cdp:eq:D.31}
\]
because the matrix with these four columns is
\[
 \begin{pmatrix}
 1&0&0&1\\0&0&1&0\\0&1&0&0\\0&0&1&1
 \end{pmatrix},\qquad\det=-1.
\]
There is no subsequent contribution from degree zero:
the three boundary components map injectively into the
four components consisting of $M,U_1,U_2,U_3$.
Moreover the degree-one source has rank $3(4+1)=15$
and the target has rank $(4+2)+3\cdot3=15$.
The preceding cokernel computation makes this homomorphism
surjective over $\Z$, hence an isomorphism. In particular
its kernel contributes nothing to $H_2(Y;\Z)$.

\item By \eqref{cdp:eq:D.25}, the pure fibre two-cycles map onto each
$H_2(U_i;\Z)$ with kernel $\Z[C_i]$. Eliminate those
local targets in the degree-two Mayer--Vietoris cokernel.
By \eqref{cdp:eq:D.26}, all local mixed cycles map to zero in $U_i$.
Their exterior images \eqref{cdp:eq:D.30} eliminate the eight
generators in $\Lambda\otimes H_1$ of the exterior base,
leaving exactly the relations
$x\wedge(c_1+c_2+c_3)=x\wedge e_1$. Thus the precise
remaining quotient is
\[
 H_2(Y;\Z)=
 \frac{\bigwedge^2\Lambda}
 {e_1\wedge\Lambda+\Z[C_1]+\Z[C_2]+\Z[C_3]}.
\tag{D.32}\hypertarget{tag-019J}{}\label{cdp:eq:D.32}
\]
Modulo $e_1\wedge\Lambda$ its three new relations are
\[
 [C_1]=e_3\wedge e_4,\qquad
 [C_2]=-e_2\wedge e_3+e_3\wedge e_4,\qquad
 [C_3]=-e_2\wedge e_4+e_3\wedge e_4.
\]
They form an integral basis of the rank-three quotient.
Consequently \eqref{cdp:eq:D.32} is zero, including its torsion.
The attached program constructs both full signed
Mayer--Vietoris matrices: their sizes are $15\times15$
and $29\times30$, the first has determinant $+1$, and
all nonzero Smith entries of both are $1$. It also checks
the six-dimensional relation matrix in \eqref{cdp:eq:D.32}.

\item Each central component is an elliptic bundle over a
compact curve and therefore has Euler number zero; their
intersection is elliptic and also has Euler number zero.
The three central fibres have Euler number zero. The
smooth torus fibre has Euler number zero, so additivity
over the base gives $e(Y)=0$. Gauss--Bonnet and Poincar\'e
duality now give $c_3(Y)=0$ and $b_3(Y)=2$.

\item Finally every effective divisor on $Y$ is vertical. If an
irreducible divisor dominated the base, its restriction
to a general K\"ahler torus fibre would be a nonzero
effective divisor. Its rational Chern class is zero
because $H^2(Y;\Q)=0$; integration against a K\"ahler
form gives the contradiction of a strictly positive
number equal to zero. A meromorphic function therefore
has only vertical zero and pole divisors, and its
restriction to a general compact fibre is a holomorphic
nowhere-zero function, hence constant. It consequently
descends to a meromorphic function of the base. To verify
extension at a special value with parameter $t$, choose
$N$ so that multiplication by $f^*t^N$ removes the pole
order on every component of the full fibre. The product
has no polar divisor and is holomorphic on the smooth
total space, by the local normality criterion for a
meromorphic function. Properness, connected fibres and
normality of the base give $f_*\OO_Y=\OO_{\PP^1}$ by
Stein factorization. The product therefore descends to a
holomorphic function, proving the required meromorphic
extension after division by $t^N$. Thus
$\mathcal M(Y)=f^*\C(t)$ and $a(Y)=1$.
The universal coefficient theorem also gives
$H^1(Y;\Z)=H^2(Y;\Z)=0$ from \eqref{cdp:eq:D.31}--\eqref{cdp:eq:D.32}.
This proves every assertion in \eqref{cdp:eq:D.21}.

\end{enumerate}
\end{proof}



\paragraph{Literature comparison.}
The construction uses explicit principal elliptic bundles,
one ordinary point blowup in each local quotient surface,
and the specified integral clutchings. Searches for
``complex threefold, Kodaira, Hopf, fibre, monodromy'',
``algebraic dimension, $S^3\times S^3$, torus'', and
``complex torus, blow-up, principal elliptic bundle,
degeneration'' did not identify this marked example in
the sources inspected. 
The complete construction and integral calculation above
are independent of a literature classification of such
fillings. They also show why extending the group family
of \ref{cdp:st:D.2.6} does not imply extension of its marked Albanese
map on the original model: in \eqref{cdp:eq:D.25} one integral
one-cycle of the smooth torus is killed although its
linear monodromy is the identity.




\begin{proposition}[Coherent and integral comparisons for \ref{cdp:st:D.2.7}]\hypertarget{tag-01GF}{}\label{cdp:st:D.2.8}
\StatementRecord{cdp:st:D.2.8}{proved}{01GF}{\SeriesRef{cdp:st:D.2.7}, \SeriesRef{cdp:st:D.1.11}, \SeriesRef{cdp:st:D.1.12}, \SeriesRef{cdp:eq:D.36}, \SeriesRef{cdp:st:D.1.7}}


For the threefold and markings of \ref{cdp:st:D.2.7},
\begin{equation}
 R^0f_*\OO_Y=\OO_{\PP^1},\qquad
 R^1f_*\OO_Y=\OO_{\PP^1}^{\oplus2},\qquad
 R^2f_*\OO_Y=\OO_{\PP^1},
 \qquad h^q(Y,\OO_Y)=(1,2,1,0).
\hypertarget{tag-019K}{}\label{cdp:eq:D.33}
\end{equation}
The coherent restriction map to every full fibre is surjective.
The integral comparison is
\begin{equation}
 H^0(\PP^1,R^1f_*\Z)
   =\Z(e_1^*+e_2^*-e_4^*)
   \ \subsetneq\ H^1(A;\Z)^{\pi_1}=\Z^4.
\hypertarget{tag-019L}{}\label{cdp:eq:D.34}
\end{equation}
The map in \eqref{cdp:eq:D.34} is restriction to the smooth base. The exceptional torsion locus
for restrictions of $\Pic(Y)=H^1(Y,\OO_Y)\simeq\C^2$ to
at least one fibre is exactly
\begin{equation}
 \iota\bigl(H^1(A;\Q)\bigr)\subset H^1(A,\OO_A),
 \qquad
 \iota(a,b,c,d)=\tfrac12(a+\mathrm i b,c+\mathrm i d).
\hypertarget{tag-019M}{}\label{cdp:eq:D.35}
\end{equation}
It is countable and dense. Thus the all-fibre non-torsion locus
of \ref{cdp:st:D.1.11} need not be open, even under integral
$H^1(Y;\Z)=H^2(Y;\Z)=0$.
\end{proposition}


\begin{proof}

\begin{enumerate}[leftmargin=*]
\item Write $\mathscr E=R^1f_*\OO_Y$. Proposition~\ref{cdp:st:D.1.12} gives its
local freeness, full-fibre base change, and rank two for this
threefold. A local filling $U_i$ retracts onto its central
fibre $S_i$: the retraction of the blown-up base onto
$B_0\cup R$ lifts through its principal compact torus
bundle, by the homotopy lifting property. Therefore \eqref{cdp:eq:D.25}
identifies
\begin{equation}
 H^1(S_i;\Z)=H^1(U_i;\Z)
          =\{\lambda\in\Lambda^*:\lambda(w_i)=0\}.
\hypertarget{tag-019N}{}\label{cdp:eq:D.36}
\end{equation}
The coefficient map into $H^1(U_i,\OO_{U_i})$
corresponds by coherent Leray on the Stein disc to
holomorphic sections of $\mathscr E$. On the punctured
disc, these are constant vectors in $H^1(A,\OO_A)$,
because \eqref{cdp:eq:D.24} has constant marked periods. The explicit
coefficient map is $\iota$ in \eqref{cdp:eq:D.35}: the $(0,1)$ parts
of $dx,dy$ on $E$ are $d\bar z/2$ and
$\mathrm i\,d\bar z/2$.

\item Each annihilator in \eqref{cdp:eq:D.36} has real dimension three, so
its complex span in the two-dimensional period space is
all of that space. The constant complex relation among
three chosen integral generators holds over the
punctured disc and therefore holds as an identity of
holomorphic sections over the entire disc. At zero,
their three real vectors remain independent: the map
$H^1(S_i;\R)\to H^1(S_i,\OO_{S_i})$ is real-linearly
injective by \ref{cdp:st:D.1.7}, and \ref{cdp:st:D.1.12} identifies this target
with the fibre of $\mathscr E$. Three real independent
vectors cannot lie in one complex line. Their central
complex span thus still has dimension two. Two of the
constant period vectors forming a basis consequently
extend as a frame of $\mathscr E$ across zero.

\item These local frames agree with the same constant period
frame on the smooth base. They therefore glue to
$\mathscr E\simeq\OO^{\oplus2}$ on $\PP^1$.
The cup-product isomorphism in \ref{cdp:st:D.1.12} gives
$R^2f_*\OO_Y\simeq\det\mathscr E\simeq\OO$.
Connected fibres give $f_*\OO_Y=\OO$. Leray on $\PP^1$
proves all four dimensions in \eqref{cdp:eq:D.33}, and evaluation
of the trivial rank-two bundle, followed by the
base-change isomorphism, proves coherent surjectivity.

\item The global integral sections are the intersection of
the three annihilators \eqref{cdp:eq:D.36} inside the constant
smooth-fibre lattice. Their defining equations are
$\lambda_3=0$, $\lambda_2+\lambda_4=0$,
$\lambda_1+\lambda_4=0$. The primitive solution is
$(1,1,0,-1)$, giving \eqref{cdp:eq:D.34}. The three singular fibres have two components each. Thus
$r=6$, $s=3$, and $t'=4$ give $6=3-1+4$.

\item Finally $\Pic(Y)=H^1(Y,\OO_Y)$ by the integral vanishing in \cref{cdp:st:D.2.7}
and the universal coefficient theorem. Under its coherent restriction
identification with $H^1(A,\OO_A)$, a line bundle is
torsion on a smooth fibre exactly when some positive
integral multiple of its class belongs to
$\iota(\Lambda^*)$, by the exponential sequence.
This is exactly \eqref{cdp:eq:D.35}. At the singular fibre $S_i$
the same exponential sequence and \eqref{cdp:eq:D.36} replace
$\Lambda^*$ by its sublattice $\operatorname{ann}(w_i)$.
That torsion locus is contained in \eqref{cdp:eq:D.35}, so the union
over all fibres is exactly the displayed set. The
real-linear map $\iota:\R^4\to\C^2$ is an isomorphism;
the image of $\Q^4$ is therefore countable and dense.
Its complement is dense and has full Lebesgue measure,
but has no nonempty open subset. 

\end{enumerate}
\end{proof}





\begin{proposition}[Euler equality under trivial torus monodromy]\hypertarget{tag-01G9}{}\label{cdp:st:D.2.2}
\StatementRecord{cdp:st:D.2.2}{proved}{01G9}{\SeriesRef{cdp:st:C.2.G.1}, \SeriesRef{cdp:st:D.2.5}, \SeriesRef{cdp:st:D.2.7}, \SeriesRef{cdp:st:D.2.3}}

The class of smooth compact connected complex threefolds with $b_1=b_2=0$,
a proper surjective holomorphic map with connected fibres onto a smooth compact curve,
smooth two-torus general fibre, and trivial rational monodromy is nonempty. Every member satisfies
\[
                   c_3(Y)=0,\qquad b_3(Y)=2.
\]
Every individual reduced fibre support has Euler number zero and
the geometric description of \SeriesRef{cdp:st:C.2.G.1}. Arbitrary singularities and
full-fibre multiplicities are allowed. There is an example with
integral $H_1(Y,\mathbb Z)=H_2(Y,\mathbb Z)=0$ and
algebraic dimension one.
\end{proposition}


\begin{proof}
The unrestricted lower bound of \SeriesRef{cdp:st:C.2.G.1} and the
trivial-monodromy upper bound of \ref{cdp:st:D.2.5} give $c_3(Y)=0$.
Poincar\'e duality gives $b_3(Y)=2-c_3(Y)=2$.
\SeriesRef{cdp:st:C.2.G.1} expresses $c_3$ as the sum of the nonnegative
Euler numbers of the reduced fibre supports, so each summand
vanishes and satisfies that theorem's exact equality criterion.
Statement \ref{cdp:st:D.2.7} constructs the asserted example with integral
homology vanishing and algebraic dimension one. The monodromy obstruction of \ref{cdp:st:D.2.3} gives the normal-crossings equality in \ref{cdp:st:D.2.4}.

\end{proof}






\section{Fibration stability and explicit deformation calculations}\hypertarget{tag-034A}{}\label{sec:II:14}


\begin{proposition}[The infinitesimal map and the rank-two boundary]\hypertarget{tag-01GU}{}\label{cdp:st:E5}
\StatementRecord{cdp:st:E5}{proved}{01GU}{the definitions and constructions stated in the proof}

Let $g:Y\to C$ be a proper surjective holomorphic map with connected fibres
from a smooth connected compact complex threefold to a nonsingular compact
connected complex curve. Its smooth general fibre is a two-dimensional complex
torus. There is a canonical exact sequence of vector spaces
\begin{equation}\hypertarget{tag-019Y}{}\label{cdp:eq:E-Leray-map}
0\to H^1(C,T_C)\xrightarrow{g^*}H^1(Y,g^*T_C)
 \to H^0(C,T_C\otimes R^1g_*\mathcal O_Y)\to0.
\end{equation}
On the smooth locus $U\subset C$, the relative tangent sequence
induces a sheaf morphism
\[
 \delta:T_C\otimes R^1g_*\mathcal O_Y|_U
                  \longrightarrow R^2g_*T_{Y/C}|_U.
\]
At a smooth torus fibre $F$, write its Kodaira--Spencer tensor as
$\kappa=u_1\otimes\eta_1+u_2\otimes\eta_2$, where
$\eta_1,\eta_2$ is a basis of $H^1(F,\mathcal O_F)$ and
$u_i\in\operatorname{Lie}(F)$. Then
\[
 \delta(a\eta_1+b\eta_2)
   =(bu_1-au_2)\otimes(\eta_1\wedge\eta_2).
\]
Its rank is the rank of $\kappa$. When this rank is two at one
regular value, the image of $H^1(Y,T_Y)\to H^1(Y,g^*T_C)$ lies
in $g^*H^1(C,T_C)$. For $C=\mathbb P^1$ this image is zero.
\end{proposition}


\begin{proof}

\begin{enumerate}[label={$\langle1\rangle$\arabic*.},leftmargin=*]
\item The projection formula gives
$R^qg_*g^*T_C=T_C\otimes R^qg_*\mathcal O_Y$, and the coherent
Leray spectral sequence has only columns zero and one. Its
degree-one edge sequence is exactly \eqref{cdp:eq:E-Leray-map}.
On $U$ the sequence
$0\to T_{Y/C}\to T_Y\to g^*T_C\to0$ is exact. Its extension
class on $F$ is the Kodaira--Spencer tensor. The connecting map is
cup product with that class; applying alternating multiplication to
the two displayed basis vectors gives the formula and its rank.
Exactness of the higher-direct-image sequence says that the image
of a global tangent class restricts to the kernel of $\delta$.
Rank two holds on an open dense subset after it holds at one point.
By \cite[Theorem~2.2, pp.~4--5]{Viaclovsky2026Combs}, the source $T_C\otimes R^1g_*\mathcal O_Y$ is locally free,
so a section zero on that subset is zero everywhere. This proves
the image assertion and the final statement.

\item At a critical point the sheaf map $dg$ need not be onto. Its correct
image is $\mathcal I_{\rm crit}\otimes g^*T_C$, giving
\[
0\to T_{Y/C}\to T_Y\to
 \mathcal I_{\rm crit}\otimes g^*T_C\to0.
\]
For $t=z^m$, the image ideal is $(z^{m-1})$; for $t=xy$ it is
$(x,y)$; for $t=xyz$ it is $(xy,xz,yz)$. On a root cover of a
smooth multiple fibre the ordinary tangent sequence for $z$ applies.
At a normal-crossing fibre the logarithmic map sends
$\sum a_i z_i\partial_{z_i}$ to
$(\sum a_i)t\partial_t$ and is onto the logarithmic base tangent
line. These formulas type the comparison at the special fibres.
The proof above only uses the genuine smooth-locus boundary and
the globally locally free source, so it makes no assumption about
local freeness of $R^2g_*T_{Y/C}$ at a special value.

\end{enumerate}
\end{proof}






\begin{proposition}[Stability of a rank-two torus fibration]\hypertarget{tag-01GV}{}\label{cdp:st:E6}
\StatementRecord{cdp:st:E6}{proved}{01GV}{\SeriesRef{cdp:st:E5}, \SeriesRef{cdp:eq:E-surj-embedding}, \SeriesRef{cdp:eq:E-Leray-map}, \SeriesRef{cdp:st:E1}; \cite[\S4, equations (4.3)--(4.5), p.~649; \S6, Theorem~6.1 and proof, pp.~660--661]{Horikawa1974}; \cite[Section~9, Theorems~1--2, pp.~83--85]{Douady1966}; \cite[Section~10.1, Proposition~1, p.~86]{Douady1966}; \cite[Section~10.2, Theorem~1, p.~87]{Douady1966}; \cite[Theorem 11]{Buchdahl1999}}

Let $g:Y\to\mathbb P^1$ be as in \ref{cdp:st:E5} and suppose its
Kodaira--Spencer tensor has rank two at one regular value. For every
proper deformation $\mathcal Y\to(S,0)$ of $Y$ over a complex
analytic germ, including a nonreduced germ, there is, after shrinking
$S$, a holomorphic map
\[
 G:\mathcal Y\longrightarrow\mathbb P^1\times S
\]
over $S$ whose central restriction is $g$. Every $G_s$ is surjective
with connected fibres and smooth general fibre a complex two-torus.
With $b_2(Y)=0$, every sufficiently small deformation has algebraic
dimension one.
\end{proposition}


\begin{proof}

\begin{enumerate}[label={$\langle1\rangle$\arabic*.},leftmargin=*]
\item We prove the lifting assertion at every Artin order; this is more
than the tangent-space assertion of \ref{cdp:st:E5}. The obstruction calculations
below use the Cech construction of holomorphic-map deformations
in \cite[\S4, equations (4.3)--(4.5), p.~649; \S6, Theorem~6.1 and proof, pp.~660--661]{Horikawa1974};
we spell out the portions needed here. The injectivity hypothesis
on degree-two tangent cohomology in its Theorem~6.1 is replaced
here by the explicitly proved unobstructedness of the torus source.

\item First fix a rank-two regular fibre and its embedding $i:F\hookrightarrow
Y$. The normal line is canonically $N=\mathcal O_F\otimes T_c\mathbb
P^1$. The tangent sequence and \ref{cdp:st:E5} imply that
\begin{equation}\hypertarget{tag-01A0}{}\label{cdp:eq:E-surj-embedding}
 H^1(F,T_F)\longrightarrow H^1(F,i^*T_Y)
 \quad\hbox{is surjective}.
\end{equation}
Let $A'\to A$ be a small extension of local Artin complex algebras,
with kernel $I$ annihilated by the maximal ideal. Suppose an embedding
$i_A:F_A\hookrightarrow Y_A$ deforming $i$ and a deformation
$Y_{A'}$ lifting $Y_A$ have been specified. The deformation functor
of a complex torus is smooth: with a real lattice marking its small
complex structures are graphs of arbitrary small complex linear maps
between the two complementary Hodge subspaces; constant Beltrami
tensors have zero bracket. This gives a smooth four-dimensional
miniversal germ for $F$. In particular $F_A$ has a lift $F_{A'}$.

\item Choose local lifts of $i_A$ on sufficiently small coordinate charts.
They exist because the target deformation is smooth over $A'$.
Their differences on overlaps form a one-cocycle in
$i^*T_Y\otimes I$. Changing the lifted transition functions of
$F_{A'}$ by a cocycle $v_{ij}\in T_F\otimes I$ changes this
obstruction by $di(v_{ij})$. Such changes are precisely the action
of $H^1(F,T_F)\otimes I$ on the set of lifts of $F_A$; their
cocycle law has no quadratic term because $I\mathfrak m_{A'}=0$.
Surjectivity \eqref{cdp:eq:E-surj-embedding} therefore makes the map
obstruction zero. A zero-coboundary correction glues the local maps.
The resulting proper map is a closed immersion: the map on local
rings is surjective modulo the nilpotent ideal, hence surjective
by Nakayama, and its underlying central map is the embedding $i$.
Thus the given embedded torus lifts in every ambient small extension.

\item Now suppose $g_A:Y_A\to\mathbb P^1_A$ has been lifted and that
$Y_{A'}$ is prescribed. Choose local lifts $t_i$ of the target
coordinate of $g_A$; on overlaps their discrepancies define the
holomorphic-map obstruction
\[
 \alpha\in H^1(Y,g^*T_{\mathbb P^1})\otimes I.
\]
Coordinate changes on the target give exactly the transition factors
of $g^*T_{\mathbb P^1}$, and the triple-overlap relation gives its
Cech cocycle law. A coboundary makes the local maps agree, so
$\alpha=0$ is sufficient as well as necessary.

\item For a rank-two regular value $c$, the divisor
$F_A=g_A^{-1}(c)$ is a smooth deformation of the central torus.
The preceding embedding argument provides a Cartier divisor
$F_{A'}\subset Y_{A'}$ lifting it. This forces the restriction
of $\alpha$ to $F$ to vanish, as can be checked without a comparison
of obstruction theories: choose local equations $t_i-\eta_i$ for
$F_{A'}$, where $\eta_i$ lies in $I\otimes\mathcal O_Y$.
Two such equations differ by a unit. Modulo the central equation
$t=0$, that unit difference disappears, and their overlap equality
reads
\[
 (t_i-t_j)|_F=\eta_i|_F-\eta_j|_F.
\]
This is exactly the assertion that the restricted obstruction
cocycle is a coboundary in $N\otimes I$.

\item By \eqref{cdp:eq:E-Leray-map} and $H^1(\mathbb P^1,T_{\mathbb P^1})=0$,
$\alpha$ is a section of the locally free sheaf
$T_{\mathbb P^1}\otimes R^1g_*\mathcal O_Y$, tensored with $I$.
Its evaluation at every rank-two regular fibre is zero. That set
is open dense; hence the section and $\alpha$ are zero. The map
therefore lifts across every small extension.

\item We pass from the formal assertion to an analytic map and
specify the relative parameter space, including its nilpotents.
Let $H$ be the germ at the graph of $g$ in the absolute Douady
space of $\mathcal Y\times\mathbb P^1$, and let
$\pi:\mathcal R\to H$ be its universal compact flat family.
Douady's existence and finiteness theorems
\cite[Section~9, Theorems~1--2, pp.~83--85]{Douady1966}
apply to this analytic ambient space; its local rings and those
of the parameter space may have nilpotents.

\item After shrinking $H$, one has
\[
 \pi_*\mathcal O_{\mathcal R}=\mathcal O_H,
\]
compatibly with base change. To check this without assuming a
reduced parameter space, represent $R\pi_*\mathcal O_{\mathcal R}$
near the central point by a finite complex of finite free analytic
modules. On the central fibre the unit spans
$H^0(Y,\mathcal O_Y)=\mathbb C$. Split off this unit in degree
zero. The differential is injective modulo the maximal ideal on
the complementary degree-zero summand; an invertible minor
cancels that summand with its image in degree one. The resulting
complex begins with $\mathcal O_H$ in degree zero and zero
differential from it. This proves the displayed equality and its
base-change assertion on the full analytic germ.

\item Embed a representative of $(S,0)$ in a polydisc. The coordinate
functions of the map $\mathcal R\to\mathcal Y\to S$ are now
pullbacks of unique holomorphic functions on $H$. They satisfy
the defining equations of $S$, because their pullbacks do and the
unit map just proved is injective. They therefore define a
holomorphic morphism $H\to S$, and the universal subspace lies
scheme-theoretically in
\[
 (\mathcal Y\times_S H)\times\mathbb P^1.
\]
The projection of this subspace to $\mathcal Y\times_S H$ is a
proper morphism between two proper flat families over $H$ and
is an isomorphism at the central point. By the isomorphism-locus
criterion \cite[Section~10.1, Proposition~1, p.~86]{Douady1966},
it is an isomorphism after restricting to an open neighbourhood
of that point. The proof of that criterion uses the finite
direct image, flatness and Nakayama's lemma, and applies to
nonreduced $H$. On this open locus the universal subspace is
the graph of an actual map to $\mathbb P^1$.

\item The universal property of the absolute Douady space consequently
identifies this open locus with the germ $\mathcal H\to S$ of
relative maps from $\mathcal Y/S$ to
$\mathbb P^1\times S/S$. This is the relative version of
Douady's graph construction
\cite[Section~10.2, Theorem~1, p.~87]{Douady1966}; the argument
above supplies the variation of the compact source explicitly.
The just-proved lifting property for every
small Artin extension says that this morphism is formally smooth
at $[g]$. The local analytic smoothness criterion identifies a
formally smooth morphism of finite-dimensional analytic germs with
a projection $S\times(\mathbb C^r,0)\to S$ after shrinking. One
can see the criterion on presentations of their convergent local
algebras: lifting across all square-zero quotients makes the
relative conormal obstruction zero, and the analytic implicit
function theorem eliminates the corresponding equations. The
zero section of this product supplies the required holomorphic
family of graphs. This argument applies to nonreduced $S$.

\item At one regular value, the family $G$ remains a proper submersion
near the central compact fibre. Its nearby fibres are connected
deformations of a torus. Each map $G_s$ is nonconstant, and hence
surjective. Its Stein factorization has finite degree one because
it has a connected regular fibre in this neighbourhood. Thus
every fibre is connected. Over the connected smooth locus of
$G_s$, Ehresmann's theorem identifies the differentiable fibres
with $T^4$. Every compact complex surface with this differentiable
type is a torus: its even first Betti number makes it K\"ahler
by \cite[Theorem 11]{Buchdahl1999};
its Albanese has dimension two; the exterior algebra on its four
integral degree-one classes shows the Albanese map has degree one;
and a nontrivial exceptional curve of this degree-one map would
include a $(-1)$-curve, impossible for the even intersection form
of $T^4$. Hence the general fibre is a complex two-torus.
Finally $b_2$ is constant in the differentiable deformation family,
and \ref{cdp:st:E1} proves $a(Y_s)=1$.

\end{enumerate}
\end{proof}




\begin{proposition}[Deformations of the constructed torus fibrations]\hypertarget{tag-01GW}{}\label{cdp:st:E7}
\StatementRecord{cdp:st:E7}{proved}{01GW}{\SeriesRef{cdp:st:E6}, \SeriesRef{cdp:eq:E-log-period}, \SeriesRef{cdp:st:E1}; \cite[Proposition~6.1]{Viaclovsky2026Combs}; \cite[Proposition~3.1]{Viaclovsky2026Two}; \cite[Proposition~2.5]{LiuRollenske2026}; \cite[Proposition~4.1, case~(i)]{Viaclovsky2026Combs}}

Every member of the following torus-fibred families,
including the period family $X_c$, its members with $|p|>1$,
Viaclovsky's $X_{3k,4l}$, his two-parameter family and his combs,
and the Liu--Rollenske family, satisfies the rank hypothesis of \ref{cdp:st:E6}.
Consequently every sufficiently small deformation retains a
holomorphic torus fibration over $\mathbb P^1$. Every such member
with $b_2=0$, and every sufficiently small deformation of it, has
algebraic dimension one. In particular the fibration-retention
question in \cite[Section~10(1)]{Viaclovsky2026Combs} has an
affirmative answer.
\end{proposition}


\begin{proof}

\begin{enumerate}[label=\arabic*.,leftmargin=*]
\item At their hexagonal fibre, a marking on the punctured disc puts
the Hodge periods in the form
\begin{equation}\hypertarget{tag-019Z}{}\label{cdp:eq:E-log-period}
  (I_2,Z(q)),\qquad
 Z(q)=\frac{\log q}{2\pi i}B+H(q),
 \qquad \det B\ne0,
\end{equation}
where $H$ is holomorphic at zero. Therefore
\[
 qZ'(q)=\frac{B}{2\pi i}+qH'(q)
\]
is invertible for all sufficiently small nonzero $q$. The
Kodaira--Spencer tensor of a marked torus with periods $(I,Z)$
is represented by $Z'$ after the invertible identifications of
the Lie algebra and $H^1(\mathcal O)$ with the chosen period
coordinates. Thus it has rank two at those regular fibres.

\item The local period models occur in
\cite[Proposition~6.1]{Viaclovsky2026Combs},
\cite[Proposition~3.1]{Viaclovsky2026Two}, and
\cite[Proposition~2.5]{LiuRollenske2026}.
The AC cusp has logarithmic matrix
$B_0=\left(\begin{smallmatrix}0&1\\-1&0\end{smallmatrix}\right)$.
The affine twists and finite logarithmic transformations change
the finite collars and constant gluing data, while the punctured
cusp matrix still has this same invertible logarithmic coefficient.
Engel's two deck transformations at the $I_1$-fibre have independent
valuation vectors $(1,0)$ and $(0,1)$, again giving an invertible
$B$. In the two-parameter family the designated fibre over $p_3$
is the full-lattice rank-two filling; its displayed two valuation
vectors are independent. Liu--Rollenske's compactification over
$I_1$ uses the same two independent valuation directions. Under the covering
$(z_1,z_2)\mapsto(e^{2\pi iz_1},e^{2\pi iz_2})$, lifts of the
angular loops give the two constant period columns, and logarithms
of the two deck multipliers, divided by $2\pi i$, give the columns of $Z(q)$.

\item For the combs, \cite[Proposition~4.1, case~(i)]{Viaclovsky2026Combs}
uses the flat cohomology basis $\alpha_i,\beta_i$ and the
single-valued holomorphic frame
$e_i=\alpha_i+\ell(q)\beta_i$, $f_i=\beta_i$, with
$\ell(q)=\log q/(2\pi i)$. The extended Hodge subbundle has
the rows $(I,H(q))$ in the frame $(e_i,f_i)$. The fixed value of the
smoothed period section at $q=0$ keeps it transverse to the
$f_i$-plane. Its rows therefore remain $(I,H_s(q))$ with
$H_s$ holomorphic. In the flat basis $(\alpha_i,\beta_i)$ the rows are
$(I,\ell(q)I+H_s(q))$, which is \eqref{cdp:eq:E-log-period} with
$B=I$.

\item All the sources construct the indicated members as compact
connected-fibre holomorphic maps with smooth general two-torus
fibre. \ref{cdp:st:E6} applies, and \ref{cdp:st:E1} gives the algebraic-dimension assertion.
On the differentiable six-sphere these neighbourhoods of the
listed structures constitute an open set on which $a=1$.

\end{enumerate}
\end{proof}






\begin{proposition}[The theta determinant and its exact rank locus for $X_c$]\hypertarget{tag-01GX}{}\label{cdp:st:E8}
\StatementRecord{cdp:st:E8}{proved}{01GX}{the definitions and constructions stated in the proof}

Use the period convention
$\Omega=\left(\begin{smallmatrix}\tau&\mu\\\mu&\beta/6\end{smallmatrix}\right)$,
Jacobi's odd theta function with periods $1,\tau$, and
$\beta/6=(i/\pi)\log(\vartheta_1(\mu,\tau)/\eta(\tau))+\text{constant}$.
Writing $h=\partial_\mu\log\vartheta_1$, one has the identity
\begin{equation}\hypertarget{tag-01A1}{}\label{cdp:eq:E-theta-general}
 \det\frac{d\Omega}{d\tau}
 =-\left(\mu'-\frac{i}{2\pi}h\right)^2
       -\frac{\wp(\mu,\tau)}{4\pi^2}.
\end{equation}
For the actual AC period branch $\wp(\mu,\tau)=0$, this becomes
\[
 \det\frac{d\Omega}{d\tau}
       =-\frac{g_2(\tau)^2}{36\pi^2g_3(\tau)}.
\]
On the smooth root-family base, the Kodaira--Spencer tensor has rank
one at the order-three point and rank two at every other finite
point, including the order-four point. Its determinant has order
two at the order-three point. The logarithmic tensor at the cusp
has rank two.
\end{proposition}


\begin{proof}

\begin{enumerate}[label={$\langle1\rangle$\arabic*.},leftmargin=*]
\item The symmetric matrix has a concrete comparison with the raw AC
periods. Adjoin $(\mu,\beta/6)$ to the lattice generated by
$(1,0),(0,1),(\tau,\mu),(6\mu,\beta)$. This enlarges the
lattice by index six; identity on $\mathbb C^2$ induces a finite
holomorphic isogeny of degree six to the torus with period matrix
$(I,\Omega)$. Its differential and pullback on complex first
cohomology are isomorphisms, so the Kodaira--Spencer ranks agree.
This is the comparison used in the determinant calculation.

\item The convention is
$\vartheta_j(z,\tau)=\theta_j(\pi z\mid\tau)$ and
$q=e^{\pi i\tau}$. The normally convergent theta series and
\cite[eqs.~20.2.1 and 20.13.1]{DLMFTheta2026} give
$\partial_\tau\vartheta_1=(4\pi i)^{-1}\partial_\mu^2\vartheta_1$.
Put $Q=q^2$ and
$E_2=1-24\sum_{n\ge1}nQ^n/(1-Q^n)$.
The normally convergent product
\cite[eq.~23.17.8]{DLMFTheta2026} gives
$\eta'/\eta=\pi i E_2/12$ by termwise logarithmic differentiation.
Likewise \cite[eq.~20.5.1]{DLMFTheta2026} gives
\[
 h(\mu)=\frac1\mu-\frac{\pi^2E_2}{3}\mu+O(\mu^3).
\]
Here the coefficient follows from
$\sum_{n\ge1}Q^n/(1-Q^n)^2
=\sum_{n\ge1}nQ^n/(1-Q^n)$, obtained by comparing the absolutely
convergent double series. The quasi-period equation
\cite[eq.~20.2.6]{DLMFTheta2026} makes $\partial_\mu h$ elliptic.
Consequently $-\partial_\mu h-\pi^2E_2/3$ and $\wp$
have identical poles, principal parts and zero Laurent constant.
Their difference is a holomorphic elliptic function with constant
zero, so they agree. Substituting these three
identities into
$\det\Omega'=(\beta/6)'-\mu'^2$ gives
\eqref{cdp:eq:E-theta-general}. Thus the square identity in Lead~9
uses the equation $\wp(\mu)=0$; a general section requires the
displayed additional term unless its other period terms cancel it.

\item For the zero branch differentiate $\wp(\mu,\tau)=0$. Put
$c=\pi^2E_2/3$ so $h'=-\wp-c$. Differentiating the heat equation
once in $\mu$, using $\wp''=6\wp^2-g_2/2$, and using
$E_2'=(\pi i/6)(E_2^2-E_4)$, yields at $\wp=0$
\[
 \partial_\tau\wp=\frac{h\wp'}{2\pi i}
                       +\frac{i g_2}{6\pi}.
\]
Consequently
\[
 \mu'-\frac{i h}{2\pi}=-\frac{i g_2}{6\pi\wp'},
 \qquad (\wp')^2=-g_3,
\]
which proves the second formula wherever $g_3\ne0$. It is an
identity of meromorphic expressions and therefore extends across
the branch values in a local uniformizing parameter.

\item The equation $j(\tau)=1728t$ shows that $\tau-\rho$ has order one
at the order-three root parameter, while $\tau-i$ has order two
at the order-four root parameter. The modular form $g_2$ has a
simple zero at $\rho$ and $g_3$ a simple zero at $i$, with no
other zeros modulo the modular group. Multiplication by
$(d\tau/ds)^2$ therefore gives an order-two zero at the first
point, and cancels the order-two zero of $g_3(\tau(s))$ at the
second. The tensor is nonzero at the first point because
$d\tau/ds\ne0$ there, hence has rank exactly one. Away from the
two root points the coordinate change is unramified and the
formula has nonzero determinant. At the cusp
$\tau=(\log q)/(2\pi i)+O(q)$, $\mu$ has a finite limit, and
$\beta=-\tau+O(1)$; hence the logarithmic derivative has limiting
matrix $\operatorname{diag}(1/(2\pi i),-1/(12\pi i))$.
Its determinant is nonzero.

\end{enumerate}
\end{proof}






\begin{proposition}[Deformations over a fixed genus-two curve]\hypertarget{tag-01GH}{}\label{cdp:st:E10}
\StatementRecord{cdp:st:E10}{proved}{01GH}{\SeriesRef{cdp:eq:E-Leray-map}, \SeriesRef{cdp:st:E5}}

The assertion that every small deformation of a fibration extends
over its original fixed positive-genus target curve $C$ is false even for a torus
fibration with Kodaira--Spencer rank two at a regular value.
\end{proposition}


\begin{proof}

\begin{enumerate}[label=\arabic*.,leftmargin=*]
\item Let $S\to\mathbb P^1_t$ be the rational elliptic surface obtained
by blowing up the nine base points of the Hesse pencil
$x^3+y^3+z^3-3txyz=0$; see \citet[\S2, pp.~238--241]{ArtebaniDolgachev2009}. Its critical values are
$\{1,\omega,\omega^2,\infty\}$, where $\omega^3=1$, $\omega\ne1$.
Pull another copy back by $\phi(t)=1/(t-4)$; its critical values
are $\{4,5,4+\omega,4+\omega^2\}$, disjoint from the first set.
Their fibre product $Z\to\mathbb P^1$ is a smooth projective
threefold, since at every point at least one elliptic-surface
map is a submersion. Its general fibre is a product of two
elliptic curves. Both elliptic $j$-functions are nonconstant,
so their two Kodaira--Spencer derivatives are simultaneously
nonzero off a finite set. The tensor of their product is then
diagonal of rank two.

\item Let $C_s$ be the smooth genus-two double cover
\[
 y^2=(t-10)(t-11)(t-12)(t-13)(t-14)(t-15-s)
\]
of $\mathbb P^1$, for $s$ in a small disc. All its branch values
avoid the critical values of $Z$. Thus
$Y_s=Z\times_{\mathbb P^1}C_s$ is a smooth projective deformation
of a compact threefold, and $g_s:Y_s\to C_s$ has connected torus
general fibre and rank two at a regular value. The Kodaira--Spencer
class $\theta_C\in H^1(C_0,T_{C_0})$ is nonzero: any infinitesimal
isomorphism of the hyperelliptic covers descends through their
unique hyperelliptic involutions to an infinitesimal projective
transformation fixing five distinct branch values; it is zero,
whereas the sixth branch value moves with derivative one.

\item Put
$L=\pi_C^*\mathcal O_{\mathbb P^1}(1)$ on $C_0$; the
double-cover canonical formula gives $K_{C_0}=L$, of degree two.
For either rational elliptic surface, relative duality gives
$R^0\pi_*\mathcal O=\mathcal O$ and
$R^1\pi_*\mathcal O=\mathcal O(-1)$: its relative canonical
line is $\pi^*\mathcal O(1)$ and every fibre has
$h^\bullet(\mathcal O)=(1,1)$. Relative K\"unneth and the flat
base change to $C_0$ consequently give
\[
 R^qg_*\mathcal O_Y=
 \begin{cases}
  \mathcal O_{C_0}&q=0,\\
  L^{-1}\oplus L^{-1}&q=1,\\
  L^{-2}&q=2,\\
  0&q>2.
 \end{cases}
\]
The vertical tangent sheaf is $g^*(L^{-1}\oplus L^{-1})$:
the invariant vector on each elliptic factor generates its
relative tangent sheaf, also at a node, where the kernel of
$d(xy)$ is generated by $x\partial_x-y\partial_y$.
Riemann--Roch on the genus-two curve gives
$h^0(L^{-a})=0$ and $h^1(L^{-a})=2a+1$ for $a>0$.
Projection and Leray give
\[
\begin{array}{c|rrrr}
 \mathcal F&h^0&h^1&h^2&h^3\\\hline
 \mathcal O_Y&1&2&6&5\\
 g^*T_{C_0}&0&3&10&7\\
 T_{Y/C_0}&0&6&20&14
\end{array}
\]
All higher groups vanish. In particular
$H^0(C_0,T_{C_0}\otimes R^1g_*\mathcal O_Y)=0$ and the
injection $H^1(C_0,T_{C_0})\to H^1(Y,g^*T_{C_0})$ is an
isomorphism between three-dimensional spaces.

\item Compatibility of the \v{C}ech transition equations for the family
$g_s$ gives
\[
 dg(\theta_Y)=g^*\theta_C
       \quad\text{in }H^1(Y_0,g^*T_{C_0}).
\]
The right side is nonzero because the first map of
\eqref{cdp:eq:E-Leray-map} is injective. Hence the first-order
obstruction to extending $g_0$ with fixed target $C_0$ is nonzero.
The obstruction belongs to
the base-deformation summand of \ref{cdp:st:E5}.

\end{enumerate}
\end{proof}






\begin{proposition}[All coprime $(3k,4l)$ multiplicities]\hypertarget{tag-01GI}{}\label{cdp:st:E11}
\StatementRecord{cdp:st:E11}{proved}{01GI}{\SeriesRef{cdp:st:E7}; \cite[Theorem~1.1, p.~1]{Viaclovsky2026Log}}

Fix admissible logarithmic-transformation data with multiplicities
$r=3k$, $s=4l$, where $k,l>0$ and $\gcd(r,s)=1$, in the family
constructed from the original elliptic-surface data. For every member
which is diffeomorphic to $S^6$, as in
\cite[Theorem~1.1, p.~1]{Viaclovsky2026Log},
\[
 h^\bullet(X,T_X)=(1,1,1,0),
\]
its Kuranishi germ is a smooth reduced line, and varying the nonzero
linearization scale with those data fixed is locally miniversal.
For the original $(3,4)$ affine-twist family, the same statements
hold whenever $p\ne0$, including $|p|>1$.
\end{proposition}


\begin{proof}

\begin{enumerate}[label=\arabic*.,leftmargin=*]
\item Put $\mathfrak B_{r,s}=\mathbb P(r,s)$. The root morphism
$\sigma:\mathfrak B_{r,s}\to\mathfrak B_{3,4}$ over the identical
coarse $\mathbb P^1$ sends the two root sections to their $k$th and
$l$th powers. In weighted homogeneous coordinates it is
$[x:y]\mapsto[x^l:y^k]$. In line notation
$\sigma^*\mathcal O_{\mathfrak B_{3,4}}(1)
=\mathcal O_{\mathfrak B_{r,s}}(kl)$.
On each fixed smooth root cover, the derivative of an affine
deck map is its linear part. Its translation acts as the identity
on invariant Lie fields and on fibre $\mathcal O$-cohomology.
The induced equivariant bundle comparisons are therefore unchanged. The same
cusp comparison used in \SeriesRef{fam:st:A.5.2} therefore gives the actual bundles
\[
 E=G=\mathcal O\oplus\mathcal O(-kl),\qquad
 R^2\bar f_*\mathcal O_X=\mathcal O(-kl),
\]
and
\[
0\to\bar f^*E\to\mathcal T:=T_X(-\log W)
 \to\bar f^*L\to0,
\quad L=\mathcal O(r+s-rs),
\quad K_X=\bar f^*\mathcal O(kl-r-s).
\]
For example the normal-cover coordinate at the first finite point
is $u=v^k$, and differentiation changes the ordinary root tangent
line from the old one to $T_{\mathfrak B_{r,s}}$; the vertical
bundle and its cohomology are pulled back through $u=v^k$.
At the cusp $W$ remains the same
normal-crossing divisor, so its logarithmic base correction is
$-p_0=-rs$ and its normal quotient is $\eta_*K_{\mathrm dP_6}$.

\item Locally $\sigma_*\mathcal O=\mathcal O$ and its positive direct
images vanish: the additional cyclic group has invariants
$\mathbb C\{v\}^{\mu_k}=\mathbb C\{u\}$, and averaging is
exact. Consequently the cohomology of the pulled-back lines
$\mathcal O(-kl)$ and $\mathcal O(-2kl)$ is the cohomology of
the original $\mathcal O(-1)$ and $\mathcal O(-2)$, hence zero.
Projection and Leray give
$h^\bullet(X,\bar f^*E)=(1,1,0,0)$.
For the horizontal line, its coarse direct image is
$\mathcal O_{\mathbb P^1}(-1)$: a vector field on a root disc
must have its coefficient divisible by the root parameter,
so $\rho_*T_{\mathfrak B_{r,s}}=T_{\mathbb P^1}(-p_1-p_2)$,
and the extra $-p_0$ gives the asserted degree. Moreover
$L$ and $L\otimes\mathcal O(-kl)$ have negative orbifold degree,
so have no global sections. Thus
$H^0(X,\bar f^*L)=H^1(X,\bar f^*L)=0$.
The two exact tangent sequences now give
$h^0(X,T_X)=h^1(X,T_X)=1$.

\item On the smooth torus locus,
the boundary on a relative holomorphic one-form is the transpose
of the Kodaira--Spencer tensor. The latter is generically invertible
by \ref{cdp:st:E7}. Hence every ambient holomorphic one-form on an inverse-image
base disc is horizontal. It is the pullback of a base one-form.
At a finite root point the root map is a submersion, so that
one-form extends holomorphically in the root coordinate. At a
generic smooth point of $W$ its local equation is $q=z_1$, which
excludes a pole in its base coefficient. Therefore the canonical
pullback comparison gives
\[
 \bar f_*\Omega_X^1=\Omega_{\mathfrak B_{r,s}}^1.
\]
By Serre duality and the projection formula,
\[
 h^3(X,T_X)=h^0\bigl(\mathfrak B_{r,s},
                         \mathcal O(kl-2r-2s)\bigr)=0.
\]
For the last equality, a weighted monomial of that degree would
give nonnegative integers $a,b$ with
$3k(a+2)+4l(b+2)=kl$. Reduction modulo $k$ makes $b+2$ a positive
multiple of $k$, since $\gcd(k,4l)=1$. The second term alone
is then at least $4kl$, a contradiction. On a six-sphere
$\chi(T_X)=1$, giving $h^2=1$. For the original affine-twist
family, $r=3$, $s=4$, and the horizontal Leray terms are the
acyclic lines $\mathcal O(-5)$ and $\mathcal O(-6)$ on
$\mathbb P(3,4)$. Therefore $h^\bullet(\mathcal T)=(1,1,0,0)$,
and the quotient $\eta_*K_{\mathrm dP_6}$ gives
$h^\bullet(T_X)=(1,1,1,0)$ directly for every admissible affine
twist.

\item Use the period quotient and reference linearization of
\SeriesRef{fam:st:B.1.1} and \SeriesRef{fam:st:B.1.4}. They give
$u=e^{-2\pi ic}$ and
$[\zeta_1,\zeta_2]\mapsto[\zeta_1,e^{2\pi i\zeta_2}]$.
Thus $e=\partial_{\zeta_2}$ maps to $2\pi i$ times the line-bundle
Euler field $\Xi$. The integral marking of \SeriesRef{push:src:E.1.1}
gives $\gamma=-\psi$, where $\psi$ counts the forward generator
$\Phi_u=u\phi$ covering translation by $-6P$.
For $\kappa(\Phi)=(\partial\Phi)\Phi^{-1}$, differentiation gives
\[
 \kappa(\partial_u)=u^{-1}\Xi\otimes[\psi]
       =-(2\pi iu)^{-1}e\otimes[\gamma],\qquad
 \kappa(\partial_c)=e\otimes[\gamma].
\]
Here the fibre classes extend by the coherent Leray comparison of
\SeriesRef{fam:st:A.5.2}; they are nonzero on every regular torus
fibre. A first-order trivialization of
the compact member would turn the given fibration into a
first-order deformation of the fixed map $f$ with fixed source.
Such deformations are $H^0(X,f^*T_{\mathbb P^1})
=H^0(\mathbb P^1,T_{\mathbb P^1})$, by connected fibres and
the projection formula. Thus they are precisely infinitesimal
projective transformations of the base. The two multiple critical
values and the hexagonal critical value remain the same three
constant sections in this scale family; the local equations
$t=v^r$, $t-1=v^s$, $q=xyz$ show this also over the dual numbers.
The base vector field is therefore zero. The nonzero torus class
cannot be trivialized, proving the nonzero global scale derivative.
Write the Kuranishi local ring as $\mathbb C\{z\}/I$. The scale
classifying map has nonzero derivative, so before quotienting its
substitution into $\mathbb C\{u-u_0\}$ is a local analytic
isomorphism. Every element of $I$ maps to zero; hence $I=0$. The
classifying map is locally an isomorphism, proving smoothness,
reducedness and miniversality.

\end{enumerate}
\end{proof}






\begin{proposition}[A length-one torsion direct image]\hypertarget{tag-01GJ}{}\label{cdp:st:E12}
\StatementRecord{cdp:st:E12}{proved}{01GJ}{the definitions and constructions stated in the proof}

Let $A$ be an elliptic curve with origin $0$, let
$P=\mathcal O_{A\times A}(\Delta-A\times\{0\}-\{0\}\times A)$,
and let $Y=A\times A\times\mathbb P^1$, with $g:Y\to A$
projection to the second factor. For
$V=P\boxtimes\mathcal O_{\mathbb P^1}(-2)$ one has
\[
 g_*V=R^1g_*V=0,\qquad R^2g_*V=\mathbb C_0,
 \qquad h^\bullet(Y,V)=(0,0,1,0).
\]
The sheaf in degree two has length one. For $t\ne0$ the fibre
cohomology is zero in every degree; for $t=0$ it is $(0,1,1)$.
The torsion in $R^2g_*V$ contradicts the local-freeness assertion of
\citet[Lemma~1.5, pp.~79--80]{CDP98}.
\end{proposition}


\begin{proof}

\begin{enumerate}[label=\arabic*.,leftmargin=*]
\item Each nontrivial degree-zero line on an elliptic curve has no
sections: a nonzero section would have an effective divisor of
degree zero and hence trivialize it. Riemann--Roch and Serre
duality then give zero first cohomology. The trivial line has
cohomology $(1,1)$. Put $p:A\times A\to A$ and $M=P\otimes\operatorname{pr}_1^*
\mathcal O_A(0)$. Its fibres have degree one, so
$R^1p_*M=0$ and $p_*M$ is a line bundle. The exact sequence for
the section $\{0\}\times A$ gives the two-term complex
\[
 Rp_*P=\bigl[p_*M\longrightarrow\mathcal O_A\bigr].
\]
The arrow evaluates the unique section of the moving degree-one
line at the origin. Its zero is exactly $t=0$ and is simple:
in a local coordinate the moving zero divisor is the diagonal,
with equation $x-t$. Thus the arrow is the inclusion
$\mathcal O_A(-0)\to\mathcal O_A$, and
$Rp_*P=\mathbb C_0[-1]$. Tensoring with
$R\Gamma(\mathbb P^1,\mathcal O(-2))=\mathbb C[-1]$
gives $Rg_*V=\mathbb C_0[-2]$. Fibre K\"unneth gives
the stated fibre cohomology.

\item For any proper flat family over a nonsingular curve and a locally
free sheaf $V$, multiplication by a local parameter gives the
base-change exact sequence
\[
0\to (R^qg_*V)_c/t(R^qg_*V)_c
 \to H^q(Y_c,V|_{Y_c})
 \to (R^{q+1}g_*V)_c[t]\to0.
\]
It follows directly from the long direct-image sequence of
$0\to V\xrightarrow{t}V\to V|_{Y_c}\to0$. Here $[t]$ denotes
the kernel of multiplication by $t$. Local constancy of all
fibre cohomology dimensions implies local freeness and base
change by Grauert's theorem; generic vanishing alone permits
the torsion explicitly present above.

\end{enumerate}
\end{proof}



In the supplied CDP98 arXiv source, Lemma~1.5 is used in
Proposition~2.5 to pass from a general-fibre vanishing to
$R^2f_*=0$, in the proof of Theorem~3.1 for the asserted
flatness of the two-parameter direct image, in Lemma~5.4 for
$R^1f_*\mathcal O$, in Theorem~5.6 to propagate sections from a
single fibre, and in Proposition~5.7 to remove a torsion
$R^2f_*\mathcal O$. Its first use is replaced in \CompendiumRef{II}{cdp:st:E2} by vanishing
on every full fibre; the special structure-sheaf assertion in
degree one is independently established in \cite[Theorem~2.2, pp.~4--5]{Viaclovsky2026Combs}. The other uses
require their own hypotheses, rather than the refuted lemma.




\begin{proposition}[The Kuranishi germ of the two-parameter family]\hypertarget{tag-01GK}{}\label{cdp:st:E13}
\StatementRecord{cdp:st:E13}{proved}{01GK}{\SeriesRef{cdp:eq:E13-connection}, \SeriesRef{cdp:eq:E13-residue}, \SeriesRef{cdp:st:E1}, \SeriesRef{cdp:st:E6}; \cite[Proposition~6.4, pp.~19--20]{Viaclovsky2026Two}}

For every member of the family of
\cite[Proposition~6.4, pp.~19--20]{Viaclovsky2026Two}, on the parameter
bidiscs with $\mu$ sufficiently close to zero and nonzero sufficiently
small linearization scale $\lambda$, one has
\[
 (h^0,h^1,h^2,h^3)(X_{\mu,\lambda},T_X)=(1,2,2,0).
\]
Its Kuranishi germ is a smooth reduced surface, and the displayed
$(\mu,\lambda)$-family is locally miniversal. In particular its
sufficiently small deformations all have algebraic dimension one.
\end{proposition}


\begin{proof}

\begin{enumerate}[label=\arabic*.,leftmargin=*]
\item The construction uses the equation and section
\[
 S_\mu:\quad y^2=x^3+tx+\mu t+1,\qquad
 P_\mu=(-\mu,b),\qquad b^2=1-\mu^3,
\]
where the branch of $b$ is close to one. Write $p_1,p_2,p_3$ for its
three distinct nodal values, with the zero of the linearization at
$p_3$. The order-four good-reduction cover at infinity gives the
root curve $\mathfrak B=\mathbb P(1,4)$, with coarse projection
$\rho$ and $\rho^*\mathcal O_{\mathbb P^1}(1)=\mathcal O(4)$.
Let $\bar f:X\to\mathfrak B$ denote the root factorization and
$D=W_1+W_2+W_3$ the reduced divisor of Mumford fibres.

\item The vertical logarithmic Lie bundle is
\[
 E=\mathcal O\oplus\mathcal O(-1).
\]
The first summand is generated by multiplication in the line-bundle
coordinate. For a concrete splitting of the second summand put
$\omega=dx/(2y)$ and use the relative differential
\begin{equation}\hypertarget{tag-01A5}{}\label{cdp:eq:E13-connection}
 \alpha=\frac{y+b}{x+\mu}\,\frac{dx}{2y}.
\end{equation}
Its residues are $1$ at $P_\mu$ and $-1$ at the zero section; there
is no pole at $-P_\mu$, because its numerator vanishes there.
Consequently, in the rational section of
$\mathcal O(P_\mu-O)$ having divisor $P_\mu-O$, the connection
$\nabla(hs)=(dh+\alpha h)s$ is holomorphic along each smooth
elliptic fibre. Here $s$ is that rational section and a point of
the line bundle is written $v s$; the horizontal lift of a vector
$V$ is therefore $\widetilde V=V-\alpha(V)v\partial_v$.
It is a multiplicative connection: the difference between its
pullback under multiplication and the sum of its two pullbacks is a
holomorphic one-form on the product of two elliptic curves, and its
restrictions to both coordinate axes vanish. Every holomorphic
one-form on that product is the sum of forms pulled back from those
axes, so this difference is zero. Its horizontal lift therefore
commutes with the lifted translation defining $X$.

\item At infinity put $t=r^{-4}$, $x=r^{-2}x'$, $y=r^{-3}y'$. Formula
\eqref{cdp:eq:E13-connection}, in the relative cotangent sheaf, becomes
\[
 \alpha=\frac{y'+br^3}{x'+\mu r^2}\frac{dx'}{2y'},
 \qquad y'^2=x'^3+x'+\mu r^2+r^6.
\]
This is the same holomorphic connection at the good-reduction
fibre; the limiting section $(0,0)$ is a nonzero two-torsion point,
so it remains disjoint from the zero section. Moreover
$\omega=r\,dx'/(2y')$, which gives the summand $\mathcal O(-1)$
for its dual Lie vector. At a nodal fibre the connection is
logarithmic at the two branches and its lift is a section of the
vertical logarithmic tangent bundle. These connection and transition maps give the splitting on the
entire root curve.

\item The inverse pairing associated with the nondegenerate alternating
class $\xi=u\wedge w+2\psi\wedge\delta$ in the construction induces a
holomorphic isomorphism
\[
 G:=R^1\bar f_*\mathcal O_X\simeq E.
\]
On a marked smooth fibre the periods
have cross terms in ratio $1:-2$, so the holomorphic one-form
subspace $F^1\subset H^1(F,\mathbb C)$ is Lagrangian for this pairing.
Pairing with $F^1$ therefore gives the isomorphism
$H^1(F,\mathbb C)/F^1\to(F^1)^*=\operatorname{Lie}(F)$.
At a Mumford point, subtract the logarithmic monodromy term from the
periods. In the resulting holomorphic frame the alternating pairing
is still constant and nondegenerate, and the two holomorphic
one-forms are the logarithmic coordinate differentials. Their dual
frame is the logarithmic Lie frame. The quotient frame gives
$R^1\bar f_*\mathcal O_X$: at a rank-one point its first class is
the class of the loop in the normalization diagram and its second
is the class on the elliptic normalization; at the rank-two point
both classes are loops in the normalization diagram. These are
bases, by the normalization exact sequence. Thus the pairing extends
as an isomorphism at these points. The finite good-reduction cover
gives the same conclusion at infinity, and descent uses invariance
of $\xi$. Cup product and the normalization sequences
give $R^2\bar f_*\mathcal O_X=\det G=\mathcal O(-1)$.

\item The logarithmic tangent and canonical sequences are therefore
\[
\begin{split}
 0&\longrightarrow\bar f^*E\longrightarrow T_X(-\log D)
 \longrightarrow\bar f^*\mathcal O(-7)\longrightarrow0,\\
 K_X&=\bar f^*\mathcal O(-4).
\end{split}
\]
Indeed $K_{\mathfrak B}=\mathcal O(-5)$ and
$T_{\mathfrak B}(-p_1-p_2-p_3)=\mathcal O(5-12)$; taking the
determinant of the displayed logarithmic sequence gives the
canonical formula after removing $D$ of degree $12$.
For root line bundles, cohomology is computed by weighted monomials
and Serre duality with $\mathcal O(-5)$. The projection formula and
Leray give
\[
\begin{split}
 h^\bullet(X,\bar f^*E)&=(1,1,0,0),\\
 h^\bullet(X,\bar f^*\mathcal O(-7))&=(0,1,2,1),\\
 h^\bullet(X,T_X(-\log D))&=(1,2,2,1).
\end{split}
\]
The base has dimension one, and the groups of the first displayed
row in degrees two and three vanish.

\item Let $Q=\bigoplus_{i=1}^3\eta_{i*}K_{\widetilde W_i}$, the normal
quotient of $T_X(-\log D)\to T_X$. Here
$\widetilde W_1,\widetilde W_2$ are ruled over elliptic curves and
$\widetilde W_3=dP_6$, so
\[
 h^\bullet(Q)=(0,2,3).
\]
The boundary $H^1(Q)\to H^2(T_X(-\log D))$ is a map between
two-dimensional spaces.
Under Serre duality its dual is the residue homomorphism
\begin{equation}\hypertarget{tag-01A6}{}\label{cdp:eq:E13-residue}
 H^1\bigl(X,\Omega_X^1(\log D)\otimes K_X\bigr)
 \longrightarrow
 H^1(\widetilde W_1,\mathcal O)\oplus
 H^1(\widetilde W_2,\mathcal O).
\end{equation}
The dual logarithmic sequence, tensored with $K_X$, has horizontal
term $\bar f^*\mathcal O(3)$ and vertical quotient
$\bar f^*(\mathcal O(-4)\oplus\mathcal O(-3))$.
The latter has zero cohomology in degrees zero and one: the coarse
direct images of $\mathcal O(-4)$ and $\mathcal O(-3)$ are both
$\mathcal O_{\mathbb P^1}(-1)$, and every line in their tensor
products with $G$ has negative weighted degree, so has no section.
Hence the
source of \eqref{cdp:eq:E13-residue} is canonically
\[
 H^1(X,\bar f^*\mathcal O(3))
 =H^0\bigl(\mathfrak B,\mathcal O(3)\otimes G\bigr)
 =H^0(\mathcal O(3))\oplus H^0(\mathcal O(2))=\mathbb C^2.
\]
On the affine $t$-line the generators of these last two spaces are
constant in the frame just constructed. The residue in each
summand is a nonzero scalar times evaluation followed by
$G_{p_i}\to H^1(\widetilde W_i,\mathcal O)$.

\item Write $e_i$ for the double root in the nodal cubic at $p_i$.
The line killed by this last projection is the graph-loop class.
Under $G\simeq E$ it is the toric Lie line: in the local periods
$(q\alpha,\beta),(\gamma,\lambda\delta)$ it is the direction
of the degenerating first multiplicative coordinate, whereas the
second coordinate is the elliptic normalization. If $V$ is dual to
$\omega$, its horizontal lift $\widetilde V$ has vertical
line-coordinate component $-\alpha(V)$. At the nodal boundary put
\[
 a_i=\frac{b}{e_i+\mu}.
\]
Thus the toric line is generated by
$\widetilde V+a_i v\partial_v$, and the projection to the elliptic
Lie line sends $(v\partial_v,\widetilde V)$ to $(1,-a_i)$, up to
a common nonzero factor at that point. Consequently the matrix of
\eqref{cdp:eq:E13-residue}, after nonzero row rescalings, is
\[
 \begin{pmatrix}1&-a_1\\1&-a_2\end{pmatrix},\qquad
 a_1-a_2=\frac{b(e_2-e_1)}{(e_1+\mu)(e_2+\mu)}\ne0.
\]
The inequality is exact: $b\ne0$; distinct nodal values have
distinct double roots, since $p_i=-3e_i^2$; and $e_i=-\mu$ would
place $P_\mu$ at the node and force $b=0$. The boundary therefore
has rank two.

\item Near the rank-two Mumford fibre put $q=t-p_3$. In the two
multiplicative coordinates its nonconstant periods are
$(qA_1(q),B_1(q))$ and $(A_2(q),\lambda qB_2(q))$, where all four
functions are holomorphic units; see
\cite[equation~(3.1) and Proposition~3.1, p.~6]{Viaclovsky2026Two}.
Normalize the two circle periods to the identity matrix. The remaining
period matrix is $(2\pi i)^{-1}\log q\,I_2+H(q)$ with $H$ holomorphic,
and its derivative is $(2\pi iq)^{-1}I_2+H'(q)$, of rank two for
nonzero $q$ sufficiently close to zero. In these period frames this
derivative is the cotangent Kodaira--Spencer connecting map.
The vertical part of a holomorphic one-form is therefore zero on
this open set and hence throughout the regular locus.
At a reduced singular fibre, a smooth point of one component away
from the conductor is a submersion point, so the coefficient of
the resulting base one-form extends holomorphically. At infinity
the same argument applies on the smooth root cover, and invariance
under the free affine deck action gives descent. Thus
$\bar f_*\Omega_X^1=\Omega_{\mathfrak B}^1$.
It follows that
$h^3(T_X)=h^0(\Omega_X^1\otimes K_X)
=h^0(\mathcal O(-9))=0$. The long exact sequence with quotient
$Q$ now gives $h^0(T_X)=1$, $h^1(T_X)=2$, and $h^2(T_X)=2$.

\item Finally the constructed bidisc has injective Kodaira--Spencer map.
At $\mu=0$, writing $a_i=p_i(0)$, differentiation of
$4t^3+27(\mu t+1)^2=0$ gives $p_i'(0)=-9/(2a_i)$.
An infinitesimal automorphism fixing the multiple value at infinity
has the form $At+B$; the three equalities
$-9/(2a_i)=Aa_i+B$ would make a nonzero quadratic polynomial vanish
at all three $a_i$. The cross-ratio variation is therefore nonzero,
and remains so on a smaller disc. For fixed $\mu$, a normalized
double curve is $\mathbb C^*/\langle c(\mu)\lambda\rangle$.
Its elliptic modulus has nonzero $\lambda$-derivative when
$|c(\mu)\lambda|$ is sufficiently small, as follows from
$j(q)=q^{-1}+744+O(q)$. A trivial first-order deformation of $X$
would trivialize its intrinsic anticanonical pencil, its labelled
exceptional-value configuration, and this double curve; it can
therefore kill neither of the two independent tangent directions.
Thus the bidisc maps to the two-dimensional tangent space of the
Kuranishi germ with invertible differential. Present its local
ring as $\mathbb C\{z_1,z_2\}/I$; the induced substitution into
$\mathbb C\{\mu,\lambda-\lambda_0\}$ is a local analytic
isomorphism before quotienting, by the inverse function theorem.
Every element of $I$ vanishes under it, so $I=0$. The classifying
map is a local analytic isomorphism. \ref{cdp:st:E1} and \ref{cdp:st:E6} give the algebraic-dimension assertion.

\end{enumerate}
\end{proof}






\begin{proposition}[The Liu--Rollenske family and its transverse deformation]\hypertarget{tag-01GL}{}\label{cdp:st:E14}
\StatementRecord{cdp:st:E14}{proved}{01GL}{\SeriesRef{cdp:eq:E14-elliptic-family}, \SeriesRef{cdp:st:E13}, \SeriesRef{cdp:eq:E14-canonical}, \SeriesRef{cdp:eq:E14-jet-map}, \SeriesRef{cdp:eq:E14-log-coho}, \SeriesRef{cdp:st:E6}, \SeriesRef{cdp:st:E1}; \cite[Proposition~2.1]{LiuRollenske2026}; \cite[Theorem~5.1]{LiuRollenske2026}}

Every member $Y(u)$ of the Liu--Rollenske family, with its scale
sufficiently small and nonzero as in its construction, satisfies
\[
 h^\bullet(Y(u),T_{Y(u)})=(1,2,2,0).
\]
Its Kuranishi germ is a smooth reduced surface. One parameter is
the given scale; a second parameter replaces its cuspidal fibre by
two rank-one nodal fibres. All sufficiently small deformations
retain a torus fibration and have algebraic dimension one.
\end{proposition}


\begin{proof}

\begin{enumerate}[label=\arabic*.,leftmargin=*]
\item In the Weierstrass equation of \cite[Proposition~2.1]{LiuRollenske2026}, make the explicit
coordinate comparison
\[
 t_{\rm new}=\frac{27}{4}(t_{\rm old}^{-1}-1),\quad
 x_{\rm new}=x_{\rm old}/t_{\rm old}^2,\quad
 y_{\rm new}=y_{\rm old}/t_{\rm old}^3.
\]
Put $k=27/4$. Near $t_{\rm old}=0$, the new infinity coordinates are
$x_{\rm old}/[k(1-t_{\rm old})]^2$ and
$y_{\rm old}/[k(1-t_{\rm old})]^3$, and its base coordinate is
$t_{\rm old}/[k(1-t_{\rm old})]$. Near $t_{\rm old}=\infty$, put
$v=t_{\rm old}^{-1}$; the new fibre coordinates are
$v^2x_{\rm old},v^3y_{\rm old}$ and the new base coordinate is $k(v-1)$.
These charts extend the displayed maps to an isomorphism of the
Weierstrass models and their minimal resolutions, preserving $O$ and $P$.
They identify $\mathcal O(P-O)$ and its lifted translation; using the
transported linearization preserves the scale $u$.
The new equation is $y^2=x^3+tx+t$, with section $P=(-1,i)$; the multiple
value is now infinity, the cusp is at zero and the remaining nodal
value is $-27/4$. Consider the explicit deformation
\begin{equation}\hypertarget{tag-01A8}{}\label{cdp:eq:E14-elliptic-family}
 y^2=x^3+tx+t+\epsilon,\qquad
 P_\epsilon=(-1,b),\quad b^2=\epsilon-1,
 \quad b(0)=i.
\end{equation}
Its finite discriminant is $4t^3+27(t+\epsilon)^2$, whose discriminant
as a cubic in $t$ is $314928\epsilon^3(1-\epsilon)$.
The nodal value near $-27/4$ remains simple and holomorphically
labelled. The other two nodes coalesce to the type-II cusp at
$\epsilon=0$. Infinity keeps type $III^*$; on its order-four
good-reduction cover the equation is
\[
 y'^2=x'^3+x'+r^2+\epsilon r^6,
 \qquad P'_{\epsilon}=(-r^2,br^3).
\]

\item On a disc containing the cusp
and its two nearby nodal values, the total elliptic surface in
\eqref{cdp:eq:E14-elliptic-family} is smooth over $\epsilon$:
$\partial(y^2-x^3-tx-t-\epsilon)/\partial t=-(x+1)$ is a unit
near the cusp. Translation by $-2P_\epsilon$ is an automorphism
there. Put $M_\epsilon=\mathcal O(P_\epsilon-O)$ and let $c(\epsilon)$
be the simple nodal value continuing $-27/4$. The identity
$\mathcal Hom(t_{2P}^*M,M)\simeq\mathcal O(F)$ in
\cite[Proposition~2.1]{LiuRollenske2026} persists in the smooth family
of elliptic surfaces: $H^1(\mathcal O)=H^2(\mathcal O)=0$ centrally
and nearby, so the exponential sequence identifies their Picard
groups with the fixed integral $H^2$ lattice. The two line bundles
have the same Chern class throughout the family. Their ratio is
trivial on every surface; Grauert's theorem makes this trivialization
holomorphic in $\epsilon$. The section of $\mathcal O(F_{c(\epsilon)})$
with divisor $F_{c(\epsilon)}$ therefore gives a holomorphic lifted
translation, agreeing with the original one at $\epsilon=0$.
For a uniformly small nonzero scale its contraction on the
complement of the zero section is uniform on the compact elliptic
fibres. The quotient is consequently free, properly discontinuous,
proper over the disc, and smooth over the parameter. This quotient
itself fills the cusp and the nearby rank-one nodal fibres. At the
separate node the ordinary rank-two Mumford charts vary
holomorphically. At infinity the fixed order-four torsion translation
and good-reduction comparison vary holomorphically in the displayed
equation. Gluing along their marked punctured collars therefore
constructs a proper holomorphic family over $(\epsilon,u)$ whose
central scale slice is $Y(u)$. Every member has the central
differentiable type by Ehresmann. The topology theorem
\cite[Theorem~5.1]{LiuRollenske2026} thus
identifies these differentiable manifolds with $S^6$.

\item Let $B$ be the reduced multiple fibre, so $f^*(\infty)=4B$.
The line $\mathcal O_Y(B)$, its canonical section and the isomorphism
$\mathcal O_Y(4B)\simeq f^*\mathcal O_{\mathbb P^1}(\infty)$ define
the root factorization $\bar f:Y\to\mathfrak B=\mathbb P(1,4)$.
For the coarse projection $\rho:\mathfrak B\to\mathbb P^1$ one has
$f=\rho\bar f$ and
$\rho^*\mathcal O_{\mathbb P^1}(1)=\mathcal O_{\mathfrak B}(4)$.
On the central threefold the
canonical form on the contracting quotient is the elliptic-surface
two-form wedged with $dv/v$. At the cusp it extends without a zero
or pole: from $t=(y^2-x^3)/(1+x)$ one gets
\[
 dt\wedge\frac{dx}{2y}=-\frac{dx\wedge dy}{1+x}.
\]
In a hexagonal chart $xyz=t-c(\epsilon)$, the form
$d\log x\wedge d\log y\wedge dt$ is $dx\wedge dy\wedge dz$.
At infinity the substitution
$t=r^{-4}$, $x=r^{-2}x'$, $y=r^{-3}y'$ gives the factor $-4r^{-4}$
in the surface form. The good three-form is invariant under the
free order-four affine action and descends to a nonvanishing frame.
These are the local canonical maps used in \ref{cdp:st:E13}.
Since the elliptic surface has
$K_S=\pi^*\mathcal O_{\mathbb P^1}(-1)$, these comparisons give
\begin{equation}\hypertarget{tag-01A7}{}\label{cdp:eq:E14-canonical}
 K_Y=f^*\mathcal O_{\mathbb P^1}(-1)
     =\bar f^*\mathcal O_{\mathfrak B}(-4).
\end{equation}
The calculation is simultaneous in $\epsilon$. Adjunction on every
full fibre gives trivial dualizing sheaf. \cite[Theorem~2.2, pp.~4--5]{Viaclovsky2026Combs} gives $h^0=1$;
duality gives $h^2=1$; constancy of its Euler characteristic in the
proper flat family gives $h^1=2$. Thus, also on the two-dimensional
base $(t,\epsilon)$, Grauert's theorem gives locally free $R^qf_*\mathcal O$
with base change, for $q=0,1,2$.

\item The multiplicative connection
\[
 \alpha=\frac{y+b}{x+1}\frac{dx}{2y},\qquad
 \nabla(hs)=(dh+\alpha h)s
\]
has precisely the residues of the rational section $s$ with divisor
$P_\epsilon-O$. For a vector $V$ dual to $dx/(2y)$, the horizontal lift is
$\widetilde V=V-\alpha(V)v\partial_v$. The multiplicative
connection commutes with the lifted translation, by the normalized
group-law comparison in \ref{cdp:st:E13}. On the good-reduction
chart its relative form is
\[
 \alpha=\frac{y'+br^3}{x'+r^2}\frac{dx'}{2y'},\qquad
 \frac{dx}{2y}=r\,\frac{dx'}{2y'}.
\]
Together with the scalar vector $v\partial_v$, this transition gives
the logarithmic Lie bundle $E=\mathcal O\oplus\mathcal O(-1)$. At the cusp the vector dual to $dx/(2y)$ is the derivation
$2y\partial_x+3x^2\partial_y$ on the cuspidal fibre, and its
horizontal lift is well defined because
$(y+b)/(x+1)$ is regular near the singular locus.

\item Put $G=R^1\bar f_*\mathcal O$. The inverse pairing associated with the invariant alternating
class gives $G\simeq E$ over the two-dimensional $(t,\epsilon)$
base away from its single cuspidal point: at a smooth fibre it is the pairing
$H^1(F,\mathbb C)/F^1\to(F^1)^*$; at a rank-one nodal fibre its quotient
frame consists of the graph-loop class and the elliptic normalization
class; at the hexagonal fibre both classes are graph-loop classes;
and the good-reduction cover gives the frame at infinity.
These local pairing and normalization maps are the ones in
\ref{cdp:st:E13}, applied to the displayed connection and root transition. Both $G$ and $E$ are locally free across that point, by the
preceding cohomological-flatness and connection calculations.
The matrix entries of the isomorphism and its inverse extend
across that codimension-two set by Hartogs' theorem. Their products
are the identity on a dense open set, hence everywhere. Thus
\[
 G=\mathcal O\oplus\mathcal O(-1),\qquad
 R^2\bar f_*\mathcal O=\mathcal O(-1)
\]
on the central root curve as well; the second equality also follows
directly from relative duality and \eqref{cdp:eq:E14-canonical}.

\item Let $D=C+W$ be the sum of the cuspidal and hexagonal fibres.
For the local tangent calculation choose coordinates in which
$t=y^2-x^3$. They are obtained from the earlier formula by the
invertible change $\widetilde x=x/(1+x)^{1/3}$,
$\widetilde y=y/(1+x)^{1/2}$, with the branches equal to one at
zero, and then renaming the two new coordinates only for this local
logarithmic-module calculation. The logarithmic
tangent module is freely generated by
\[
 2y\partial_x+3x^2\partial_y,\quad
 2x\partial_x+3y\partial_y,\quad\partial_z,
\]
where $z$ is a coordinate along the elliptic singular curve. The
second generator maps to $6t\partial_t$;
the first and third generate the relative logarithmic kernel.
Consequently the global exact sequence is
\[
 0\longrightarrow\bar f^*E\longrightarrow T_Y(-\log D)
 \longrightarrow\bar f^*\mathcal O(-3)\longrightarrow0,
\]
because $T_{\mathfrak B}(-C-W)=\mathcal O(5-8)$.
The vertical cohomology is $(1,1,0,0)$. Every line in the
horizontal Leray terms is $\mathcal O(-3)$ or $\mathcal O(-4)$,
which has no cohomology on $\mathbb P(1,4)$. Hence
\begin{equation}\hypertarget{tag-01AA}{}\label{cdp:eq:E14-log-coho}
 h^\bullet(T_Y(-\log D))=(1,1,0,0).
\end{equation}

\item Return to the Weierstrass coordinates. The preceding coordinate
change preserves the ideal $(x^2,y)$ and induces the identity on
$\mathbb C[x]/(x^2)$, since $\widetilde x\equiv x\bmod x^2$.
The quotient at $W$ is $\eta_*K_{dP_6}$ with cohomology $(0,0,1)$.
At $C$, the map $v\mapsto v(t)|_C$ identifies the quotient with
the Jacobian ideal
\[
 J=(x^2,y)\subset\mathcal O_C.
\]
Indeed its kernel is precisely the logarithmic tangent module and
$N_{C/Y}=\mathcal O_C$. Put $Z=V(J)\subset C$ and consider
the restriction map
\begin{equation}\hypertarget{tag-01A9}{}\label{cdp:eq:E14-jet-map}
 H^1(C,\mathcal O_C)\longrightarrow H^1(Z,\mathcal O_Z).
\end{equation}

\item In the deformation \eqref{cdp:eq:E14-elliptic-family} the critical
base near the cusp is the smooth curve
\[
 \Gamma:\quad t=-3x^2,\qquad
 \epsilon=x^2(3+2x).
\]
Its map to the $\epsilon$-disc is finite flat of degree two; its
central fibre is $\operatorname{Spec}\mathbb C[x]/(x^2)$.
The lifted translation fixes its $x$-coordinate. This follows on
the punctured curve because translation fixes the unique node of
each nodal cubic, and then follows at zero by holomorphic
continuation. On the line coordinate it is
$(x,v)\mapsto(x,c(x)v)$ with $c$ a holomorphic nowhere-zero
contraction. Thus its quotient critical locus is exactly the
elliptic family
\[
 \mathcal Z=(\Gamma\times\mathbb C^*)/
       \langle(x,v)\mapsto(x,c(x)v)\rangle.
\]
In particular $H^0(Z,\mathcal O_Z)$ and
$H^1(Z,\mathcal O_Z)$ are both free rank-one modules over
$A=\mathbb C[x]/(x^2)$, and $H^2(Z,\mathcal O_Z)=0$.

\item Restriction in this proper family is a holomorphic bundle map
$G|_\Gamma\to R^1(\mathcal Z/\Gamma)_*\mathcal O_{\mathcal Z}$.
At a node with coordinate $x$, the toric Lie line is generated by
$\widetilde V+\frac{b}{1+x}v\partial_v$: the horizontal lift has
line-coordinate component $-\alpha(V)=-b/(1+x)$ at that node.
The quotient to the elliptic Lie line, and hence the normalization
restriction under the pairing above, therefore sends
$(v\partial_v,\widetilde V)$ to
$(1,-b/(1+x))$, up to a common nonzero factor. This is the local
map in \ref{cdp:st:E13}. Thus on $\Gamma\setminus\{0\}$ restriction,
after a nonzero scalar change of target frame, has the row
\[
 \left(1,-\frac{b}{1+x}\right).
\]
The normalization of $C$ is a $\mathbb P^1$-bundle over
the elliptic singular curve $E_0$, and its normalization quotient
is $\mathcal O_{E_0}$: the first jet at the cusp transforms with
derivative one under translation on the additive smooth locus.
The normalization exact sequence, together with
$h^1(\mathcal O_C)=2$, $h^2(\mathcal O_C)=1$ and
$h^\bullet(\mathcal O_{\widetilde C})=(1,1,0)$, implies
that $H^1(\mathcal O_C)\to H^1(\mathcal O_{\widetilde C})$
is onto. Restriction from the ruled normalization to its singular
section identifies the latter group with $H^1(\mathcal O_{E_0})$.
Thus restriction to $E_0$ is onto. The common scalar in the row
above cannot have a zero at zero, since both entries would then
vanish and this restriction would be zero. It is a unit.

\item The equality of these holomorphic maps on the punctured curve
extends over all of $\Gamma$, and in particular over its
nonreduced central fibre. Hence \eqref{cdp:eq:E14-jet-map} is,
in the specified bases, the complex-linear map
\[
 \mathbb C^2\longrightarrow A,\qquad
 (a,d)\longmapsto a+d(-i+ix).
\]
Its matrix in $1,x$ has determinant $i$, so it is an isomorphism.

\item The exact sequence $0\to J\to\mathcal O_C\to\mathcal O_Z\to0$
now gives
\[
 h^\bullet(C,J)=(0,1,1).
\]
Together with \eqref{cdp:eq:E14-log-coho} and the hexagonal quotient,
the tangent exact sequence gives
$h^\bullet(T_Y)=(1,2,2,0)$.

\item In the above local trivialization
of the total space the map to the base is
\[
 f_{\epsilon}=\frac{y^2-x^3-\epsilon}{1+x}.
\]
The derivative $\partial_\epsilon f$ has image $-1+x$ in
$\mathcal O_C/J=\mathbb C[x]/(x^2)$ along the singular curve.
If the ambient first-order deformation were trivial, every
first-order variation of its map would be a global section of
$f^*T_{\mathbb P^1}$, hence the pullback of a projective vector
field, because $f_*\mathcal O_Y=\mathcal O_{\mathbb P^1}$.
Its local image in this Jacobian quotient is constant. Changing
the local ambient trivialization adds $df(v)$ and hence adds
zero in that quotient. The nonzero coefficient of $x$ therefore
proves that the $\epsilon$-direction is nontrivial and cannot be
cancelled by the scale direction, which does not change this local
map. With $\epsilon=0$, a normalized singular elliptic curve is
$\mathbb C^*/\langle c_0u\rangle$; its modulus has nonzero
$u$-derivative for sufficiently small $u$. A trivial deformation
of the total space preserves its anticanonical pencil from
\eqref{cdp:eq:E14-canonical}, its three distinguished special values,
and that elliptic singular curve. Hence the scale direction is
nontrivial too. The two directions span $H^1(T_Y)$.
Present the Kuranishi local ring as
$\mathbb C\{z_1,z_2\}/I$. The classifying map of the constructed
bidisc, followed by the embedding in $H^1(T_Y)$, has invertible
differential and is a local analytic isomorphism by the inverse
function theorem. Every element of $I$ pulls back to zero, so $I=0$.
This is the ideal argument of \ref{cdp:st:E13}, and it proves that
the displayed family is locally miniversal over the full reduced germ.
The rank-two hypothesis in \ref{cdp:st:E6} is supplied for this
Liu--Rollenske construction by \ref{cdp:st:E7}. Thus all sufficiently
small deformations retain their torus fibration.
Their second Betti number remains zero, and \ref{cdp:st:E1}
identifies their meromorphic function field with the pullback of
$\mathcal M(\mathbb P^1)$, giving algebraic dimension one.

\end{enumerate}
\end{proof}






\begin{proposition}[Tangent cohomology and Kuranishi germs of positive combs]\hypertarget{tag-01GM}{}\label{cdp:st:E15}
\StatementRecord{cdp:st:E15}{proved}{01GM}{\SeriesRef{cdp:eq:E15-products}, \SeriesRef{cdp:eq:E15-positive}, \SeriesRef{cdp:st:E5}, \SeriesRef{cdp:st:E13}, \SeriesRef{cdp:st:E1}, \SeriesRef{cdp:st:E6}; \cite{Viaclovsky2026Combs}}

For each positive-comb manifold $X_m$ constructed in
\cite{Viaclovsky2026Combs}, in particular for every $m\geq12$ in
its stated construction, one has
\[
 h^\bullet(X_m,T_{X_m})=(2m+1,\,8m+1,\,6m+1,\,0).
\]
Its Kuranishi germ is smooth and reduced of dimension $8m+1$.
A locally miniversal family is supplied by all sufficiently nearby
admissible period sections on the root curve, with the affine
attachment classes fixed. Its members have algebraic dimension one
and torus fibrations with the same three critical values,
multiplicities $3,4$, and hexagonal fibre type.
\end{proposition}


\begin{proof}

\begin{enumerate}[label={$\langle1\rangle$\arabic*.},leftmargin=*]
\item Let $\mathfrak B=\mathbb P(3,4)$, with coarse projection
$\rho:\mathfrak B\to B$. Write $p_1,p_2$ for the finite-isotropy
points and $p_0$ for the hexagonal value. On the root curve let
$\mathcal H$ be the extended cohomology bundle, $F$ the Hodge
subbundle, $G$ its quotient, and $V$ its dual:
\[
 0\longrightarrow F\longrightarrow\mathcal H
 \longrightarrow G\longrightarrow0,\qquad V=F^*.
\]
The finite-cover comparison identifies
$G=R^1\bar f_*\mathcal O_{X_m}$ by invariant pushforward of the
ordinary torus Hodge sequence. At the hexagon the two deck
characters of the infinite toric quilt map to the two generators
of $H^1(\mathcal O_W)$ in its normalization complex, establishing
the same comparison there. The bundle $\bar f^*V$ is the vertical
logarithmic tangent bundle: its translations are dual to the
holomorphic one-forms on smooth fibres, and at the hexagon the
dual logarithmic coordinate fields extend in every toric chart.
Cup product identifies $R^2\bar f_*\mathcal O$ with
$D:=\bigwedge^2G$. On finite covers this is the ordinary torus
exterior algebra; at the hexagon it is the exterior algebra of
the two cycles in the normalization complex.

\item We specify finite characters by the orders of their invariant
generators. At $p_1$ the nontrivial Hodge line has invariant
generator $w$ times an equivariant frame, while the nontrivial
line of each of $V$ and $G$ has generator $w^2$ times its frame.
At $p_2$ the corresponding orders are $3$ for $F$ and $1$ for
$V,G$. Each bundle also has one trivial-character line. These
characters persist because comb smoothing fixes the limiting
periods. The tautological quotient on each tooth is
$\mathcal O(1)^{\oplus2}$, so each tooth adds $2$ to the coarse
degree of $G$. Therefore
\[
 \deg\rho_*G=2m-1,\qquad D=\mathcal O_{\mathfrak B}(24m-1).
\]
The determinant of $\mathcal H$ has degree zero on the root
curve: its determinant local system is trivial, and its cusp
residue has trace zero. Thus
\[
 \det F=\mathcal O(1-24m),\quad
 \det V=\mathcal O(24m-1),\quad
 \deg\rho_*V=2m-1.
\]
Indeed the finite fractional correction for $V$ and $G$ is
$2/3+1/4=11/12$; subtracting it from $(24m-1)/12$ gives
$2m-1$. For $F$ the correction $1/3+3/4=13/12$ gives
$\deg\rho_*F=-2m-1$.

\item The coarse bundles $V_c=\rho_*V$ and $G_c=\rho_*G$ have
nonnegative splitting. To prove this, evaluate the unchanged
invariant real homology class $\delta$ on $F$. Its resulting
holomorphic section of $V$ never vanishes on a smooth fibre:
a real class annihilating a Hodge plane also annihilates its
conjugate. The same argument holds on the finite good-reduction
covers. At the hexagon, $\delta$ is a nonzero combination of
angular cycles and the logarithmic forms have identity angular
period matrix. Its evaluation is nonzero there too, and gives
a nowhere-zero section of $V_c$. The invariant real cohomology
class $\gamma$ likewise gives a nowhere-zero section of $G_c$;
at the cusp it belongs to $\ker N$, whose projection to the
Hodge quotient is an isomorphism in the transverse graph chart.
At each finite point these nonzero invariant central values give
nonzero values of the coarse sections. A split rank-two bundle
of positive degree on $\mathbb P^1$ with a negative summand
cannot have a nowhere-zero section: every section lies in its
other summand, which has positive degree and hence has zeros.
Since $2m-1>0$, both splittings are nonnegative. In particular
\begin{equation}\hypertarget{tag-01AB}{}\label{cdp:eq:E15-positive}
 h^0(V_c)=h^0(G_c)=2m+1,\qquad H^1(V_c)=H^1(G_c)=0.
\end{equation}

\item Set $A=\rho_*(V\otimes G)$ and $L=\mathcal O_{\mathfrak B}(-5)$.
Multiplication of invariant sections gives exact comparisons
\begin{equation}\hypertarget{tag-01AC}{}\label{cdp:eq:E15-products}
\begin{split}
 0&\longrightarrow V_c\otimes G_c\longrightarrow A
       \longrightarrow\mathbb C_{p_1}\longrightarrow0,\\
 0&\longrightarrow V_c(2m-1)\longrightarrow\rho_*(V\otimes D)
       \longrightarrow\mathbb C_{p_1}\longrightarrow0,\\
 0&\longrightarrow G_c(-1)\longrightarrow\rho_*(L\otimes G)
       \longrightarrow\mathbb C_{p_1}\longrightarrow0.
\end{split}
\end{equation}
Here is the local verification. At $p_1$, multiplying two
nontrivial invariant generators of $V,G$ gives $w^4$, whereas
their tensor character has generator $w$; multiplication is
therefore one factor $t=w^3$. At $p_2$ the product is $w^2$
and the tensor generator also has order two, giving no defect.
Products involving a trivial character have no defect. The
line $D$ has the nontrivial character of $G$ and coarse degree
$2m-1$, proving the second comparison. The invariant-generator
orders of $L$ are $1,1$, and $\rho_*L=\mathcal O_B(-1)$.
Its product with the nontrivial generator of $G$ has the same
single defect at $p_1$, proving the third comparison.

\item All summands on the left in \eqref{cdp:eq:E15-products} have zero
$H^1$, by \eqref{cdp:eq:E15-positive}; for the last sequence their
degrees are at least $-1$. Consequently
\[
\begin{array}{c|c|c}
\text{bundle on }B&h^0&h^1\\ \hline
 A&8m+1&0\\
 \rho_*(V\otimes D)&6m&0\\
 \rho_*(L\otimes G)&2m&0
\end{array}
\]
by Riemann--Roch: their degrees are $8m-3$, $6m-2$, $2m-2$
and their ranks $4,2,2$. Also
$\rho_*(L\otimes D)=\mathcal O_B(2m-1)$.
Finite-group invariant pushforward is exact, so these calculations
give the cohomology on the root curve.

\item The logarithmic tangent sequence is
\[
 0\longrightarrow\bar f^*V\longrightarrow T_{X_m}(-\log W)
 \longrightarrow\bar f^*L\longrightarrow0.
\]
Projection and Leray yield
\[
 h^\bullet(\bar f^*V)=(2m+1,8m+1,6m,0),\qquad
 h^\bullet(\bar f^*L)=(0,2m,2m,0).
\]
Its only possible connecting map between nonzero groups is
$H^0(L\otimes G)\to H^0(V\otimes D)$. This is the map of
sections of the fibrewise cup-product boundary in \ref{cdp:st:E5}. The cusp
has rank two, so its sheaf map is generically injective. A map
from a locally free sheaf on a smooth curve with zero generic
kernel is injective; hence this map on sections is injective.
It follows that
\[
 h^\bullet(T_{X_m}(-\log W))=(2m+1,8m+1,6m,0).
\]
The branch-normal quotient is $\eta_*K_{dP_6}$, since the toric
three-form trivializes the canonical bundle near the fibre.
Its cohomology is $(0,0,1)$, and its exact sequence proves the
asserted tangent-cohomology tuple.

\item We construct a family with this full tangent space. All nearby
root-period sections have normal bundle $V\otimes G$, tangent
space $H^0(A)$, and obstruction space $H^1(A)=0$. This section
germ can be realized using an ordinary coarse Grassmannian bundle
and finite analytic conditions. Split the equivariant bundle as
$F\oplus G$ on a finite-cover chart and write a nearby plane as
a graph matrix. Dividing the invariant-generator factors shows
that a holomorphic lift requires one entry to vanish at $p_1$
and two at $p_2$. At $p_1$ the entry is from the nontrivial Hodge
line to the trivial quotient line. At $p_2$ it is that entry
and the entry between the two nontrivial lines. These are linear
conditions on the values of the local graph matrix. Their normal
sheaf is $A$. The Cech gluing obstruction across a small Artin
extension lies in $H^1(A)=0$; each local graph chart and its
stated conditions is smooth. The finite-dimensional Douady
section germ with these conditions is therefore smooth of
dimension $8m+1$. This description does not require a global
manifold cover of the root curve.

\item These sections produce compact threefolds even though their
limiting periods are allowed to vary. At a finite point the
equivariant lift is a small deformation on the compact
good-reduction cover. Admissibility is open there. The affine
action has the same real formula $x\mapsto Tx+v/n$; it is
holomorphic for the new equivariant periods and stays free
because its real action is unchanged. At the hexagon the
transverse graph chart is open, and its periods are
\[
 Z(q)=\frac{\log q}{2\pi i}I+H(q)
\]
with $H$ holomorphic and uniformly bounded on a smaller parameter
polydisc. The translations with unit factors
$\exp(2\pi iH_{ij})$ preserve the same toric fan. On logarithmic
radii their action is
$r\mapsto r+(I+O(1/\log|q|))\lambda$; it is uniformly proper
and free for small $q$. A fixed finite collection of toric unit
polydiscs gives compact representatives. Thus the filling is
proper with local equation $q=x_1x_2x_3$, and its total space
is smooth over the parameter. Varying $H(0)$ varies gluing units
and preserves these estimates and local equations. On the
compact complement of the critical discs admissibility is open.
The annular gluings extend holomorphically with their fixed
integral classes: the lattice--Lie--torus exact sequence and
Cartan vanishing on an annular product extend the original
translations. This constructs a proper holomorphic family with
all $8m+1$ period directions and the specified central member.

\item Its Kodaira--Spencer map is injective. A first-order trivialization
would make its map variation a section of
$H^0(X_m,f^*T_B)=H^0(B,T_B)$. The three critical values remain
constant by the equations $w^3,w^4,x_1x_2x_3$, so that base
vector field vanishes at three points and is zero. The
trivialization would then be fibre preserving. On a marked
smooth torus the period derivative is the corresponding
constant Dolbeault tensor in
$\operatorname{Lie}(F)\otimes H^1(F,\mathcal O_F)$.
A nonzero constant tensor cannot be a Dolbeault coboundary,
since derivatives of periodic coefficients integrate to zero.
Thus a period section giving a trivial total-space deformation
vanishes on every regular fibre, and hence identically. Its
Kodaira--Spencer map is an isomorphism onto the already computed
$H^1(T_{X_m})$. The inverse-function argument of \ref{cdp:st:E13} then makes
the Kuranishi ideal zero, proving smoothness, reducedness and
miniversality. The topology theorem of the comb construction
identifies these manifolds with $S^6$; \ref{cdp:st:E1} and \ref{cdp:st:E6} give algebraic
dimension one throughout their small deformation germs.

\end{enumerate}
\end{proof}



\paragraph{Comparison with the constructed parameter space.}
The germ in \cite[Lemma labelled \texttt{paramspace} and
Proposition labelled \texttt{effective}]{Viaclovsky2026Combs}
fixes the values of the coarse period section at all three
critical points. Its normal sheaf is $N(-p_0-p_1-p_2)$, where
$N=(\rho_*F)^*\otimes G_c$. The actual inclusion of that germ
into our root-period germ has the differential induced by
\[
0\longrightarrow N(-p_0-p_1-p_2)\longrightarrow A
\longrightarrow
 \mathbb C_{p_0}^{4}\oplus\mathbb C_{p_1}^{3}
                     \oplus\mathbb C_{p_2}^{2}
\longrightarrow0.
\]
The local graph conditions above prove this exact sequence:
the old sheaf sets all four entries to zero at each point,
whereas $A$ imposes zero, one, and two conditions, respectively.
For the published combs its left term has zero $H^1$, as proved
in the cited smoothing construction. Its tangent dimension is
$8m-8$, and the quotient contributes exactly nine directions.
This explains the relation of the published lower bound to the
dimension $8m+1$ proved here by an explicit morphism of parameter
germs. The source asks for the full Kuranishi germ in
Section~10(1). Searches on 15 November 2026 for its title together
with ``Kuranishi'', ``smooth'', and ``$8m+1$'', and for the
two-parameter and Liu--Rollenske titles with ``tangent cohomology'',
located the supplied construction papers but no matching complete
computations \ref{cdp:st:E13}--\ref{cdp:st:E15}. This records the literature comparison.




\begin{proposition}[Rank divisors from the root Kodaira--Spencer morphism]\hypertarget{tag-01GN}{}\label{cdp:st:E16}
\StatementRecord{cdp:st:E16}{proved}{01GN}{\SeriesRef{cdp:eq:E.16.3}, \SeriesRef{cdp:st:E8}, \SeriesRef{cdp:st:E13}, \SeriesRef{cdp:st:E14}, \SeriesRef{cdp:st:E15}, \SeriesRef{cdp:st:E7}}

We use the logarithmic Kodaira--Spencer morphism at the specified degenerate fibres and the ordinary morphism on the finite smooth torus covers. For the original affine-twist family, including every $|p|>1$, its only rank defect is rank one at the order-three point. The logarithmic cusp has rank two. For the $(3k,4l)$ logarithmic transforms, on their full root covers the ranks at the two multiple values are respectively
\[
 \begin{cases}1&k=1,\\0&k>1,\end{cases}
 \qquad
 \begin{cases}2&l=1,\\0&l>1.\end{cases}                 \tag{E.16.1}\hypertarget{tag-01A2}{}\label{cdp:eq:E.16.1}
\]
Every other smooth point and the logarithmic cusp have rank two.

For the two-parameter family the logarithmic rank is one at each of its two rank-one nodal values, zero at the order-four root point, and two everywhere else, including the hexagonal fibre. Its determinant has simple coarse zeros at the two nodal values and order four in the root coordinate at the multiple value. For the Liu--Rollenske central family the logarithmic rank is two at both the cuspidal and hexagonal fibres and throughout the remaining smooth locus; at the order-four root point it is zero, and its determinant has order four there.

For a specified positive comb section $\sigma$ of degree parameter $m$, let $x,y$ be root coordinates of weights $3,4$, with $y=0$ at the order-three point. There is a nonzero homogeneous polynomial $P_\sigma(X,Y)$ of degree $4m$, uniquely determined up to an overall nonzero scalar by the marked period section, such that
\[
 \det\kappa_\sigma=y^2P_\sigma(x^4,y^3).             \tag{E.16.2}\hypertarget{tag-01A3}{}\label{cdp:eq:E.16.2}
\]
The rank-defect set is exactly its zero divisor. At any such point the rank is zero exactly when all four entries of the Kodaira--Spencer matrix vanish; otherwise it is one. Formula \eqref{cdp:eq:E.16.3} below computes those entries from $\sigma$. The polynomial is nonzero at the hexagonal cusp, so that cusp always has rank two.
\end{proposition}


\begin{proof}

\begin{enumerate}[label=\arabic*.,leftmargin=*]
\item Adding a linearization constant or changing only the admissible affine lifts does not change the derivative of the marked period matrix. \ref{cdp:st:E8} therefore proves the first assertion. On a root cover of a higher logarithmic transform the original local good-reduction coordinate is $s=u^k$ at the first multiple value and $s=u^l$ at the second. The lifted period map is its holomorphic pullback. The chain rule gives $k u^{k-1}\kappa(s)$ or $l u^{l-1}\kappa(s)$. \ref{cdp:st:E8} gives rank one at the original order-three centre and rank two at the original order-four centre. The covering derivative is nonzero off its centre, proving \eqref{cdp:eq:E.16.1} and the absence of further defects. The cusp coordinate is unchanged.

\item For the two-parameter family \ref{cdp:st:E13} proves, with its extension maps,
\[
 E=G=\OO\oplus\OO(-1),\qquad
 \kappa:\OO(-7)\longrightarrow E\otimes G
 \quad\text{on }\PP(1,4).
\]
Thus the four entries have weights $7,6,6,5$. Write $x$ for the weight-one coordinate vanishing at the root point and $y$ for the weight-four coordinate. Every monomial of these respective weights is divisible by $x^3,x^2,x^2,x$. Consequently the entire matrix vanishes at the root point, and its determinant is divisible by $x^4$. The determinant is a section of $\OO(12)$, whose coarse degree is three. At the three Mumford values, the local period matrix has logarithmic
part $(\log q)/(2\pi i)\operatorname{diag}(1,\epsilon)$, where
$\epsilon=0$ at the two rank-one values and $\epsilon=1$ at the
hexagonal value; see \cite[equation~(3.1), p.~6]{Viaclovsky2026Two}.
Its value under $q\partial_q$ is
$(2\pi i)^{-1}\operatorname{diag}(1,\epsilon)$ at the centre.
Thus at each of the two rank-one values the logarithmic matrix has rank one and its determinant vanishes at least once. These two coarse zeros, together with the root zero of order at least four, account for all of the degree of $\OO(12)$. The determinant is nonzero because the logarithmic rank at the hexagonal cusp is two. All the stated orders are therefore exact, with no remaining zeros.

\item For the Liu--Rollenske central member \ref{cdp:st:E14} proves the analogous map from $\OO(-3)$ on $\PP(1,4)$. Its entries have weights $3,2,2,1$; these spaces consist respectively of multiples of $x^3,x^2,x^2,x$. The determinant is a section of $\OO(4)$ and is a constant multiple of $x^4$. The constant is nonzero by the hexagonal local monodromy. Hence the only rank defect is the root point, where all four entries vanish.

\item For a positive comb, \ref{cdp:st:E15} gives
\[
 \det V=\det G=\OO(24m-1),\qquad
 \kappa_\sigma:\OO(-5)\longrightarrow V\otimes G.
\]
The original and comb period planes agree at each critical value.
A local bundle automorphism $\Psi(t)=I+O(t)$ carries one plane to
the other. On the finite covers $t=w^n$ its root lift is given by
\[
 D_3(w)=\operatorname{diag}(1,1,w,w^2),\qquad
 D_4(w)=\operatorname{diag}(1,1,w,w^3),\qquad
 \Psi^\sharp(w)=D_n(w)\Psi(w^n)D_n(w)^{-1}.
\]
The poles in $D_n^{-1}$ have order at most $n-1$, so
$\Psi^\sharp(w)=I+O(w)$ is holomorphic and equivariant, with the
same central characters. These are the maps of
\cite[Proposition~6.2, pp.~22--23]{Viaclovsky2026Combs}.
The tautological sequence on the root curve has trivial determinant
for its middle cohomology bundle: the marked flat volume extends
through the finite covers and through the determinant-one logarithmic
frame at the cusp. Taking determinants therefore identifies
$\det V$ with $\det G$. Each comb tooth contributes coarse degree
two to $G$, by \cite[equation~(5.5), p.~16]{Viaclovsky2026Combs}.
The unchanged finite characters make the determinant change the
pullback of $\mathcal O_{\mathbb P^1}(2m)$, which is
$\mathcal O_{\mathfrak B}(24m)$ on $\mathfrak B=\mathbb P(3,4)$.
The determinant is a section of $\det V\otimes\det G\otimes\OO(10)=\OO(48m+8)$. On $\PP(3,4)$ every weighted monomial $x^ay^b$ of this degree satisfies $b\equiv2\pmod3$ and $a\equiv0\pmod4$. Factoring $y^2$ therefore gives exactly the polynomial in \eqref{cdp:eq:E.16.2}, of coarse degree $4m$.

\item On a marked local root disc let the two columns of a holomorphic $4\times2$ matrix $U$ frame the Hodge subbundle of the marked cohomology extension. Let $\nabla$ be its marked Gauss--Manin connection, and let $D$ be the local generator of $T_{\mathfrak B}(-p_0)$: $\partial_s$ away from the cusp and $q\partial_q$ at the cusp. Then
\[
 \kappa_\sigma=(\nabla_DU)\bmod\operatorname{im}U,
 \qquad
 \det\kappa_\sigma\ \text{is represented by}\
 \det(U_1,U_2,\nabla_DU_1,\nabla_DU_2).              \tag{E.16.3}\hypertarget{tag-01A4}{}\label{cdp:eq:E.16.3}
\]
The first expression uses the quotient frame of $G$ and the dual frame
of $V=F^*$; it specifies every matrix entry and therefore also the
rank-zero test. On a Grassmann chart with invertible upper block $U_0$
of $U$, the frame change $U\mapsto UU_0^{-1}$ gives
$\binom{I}{Z}$. Write $A_{ij}$ for the four blocks of the connection
matrix contracted with $D$. Projection to the quotient by
$(u,v)\mapsto v-Zu$ gives the matrix
\[
 DZ+A_{21}+A_{22}Z-ZA_{11}-ZA_{12}Z.
\]
Its determinant is the four-column determinant of the frame
$UU_0^{-1}$ in \eqref{cdp:eq:E.16.3}. For a Hodge-frame change $U'=UC$, one has
$\nabla_DU'=(\nabla_DU)C+U(DC)$; wedging the four columns
multiplies the coefficient by $(\det C)^2$. Pair the quotient
determinant frame with $U_1\wedge U_2$ to the marked ambient
volume. It changes by $(\det C)^{-1}$, as does the dual Hodge
determinant frame.
For a tangent-frame change $D'=aD$, the coefficient is multiplied
by $a^2$, while the dual differential frame is multiplied by
$a^{-1}$ and its square by $a^{-2}$. The coefficient and frame
therefore give the same section of
$\det V\otimes\det G\otimes T_{\mathfrak B}(-p_0)^{-2}$.
These transition maps determine $P_\sigma$ from the marked period data. In the cusp comparison of \ref{cdp:st:E7} and \ref{cdp:st:E15},
the marked frame satisfies
$\nabla e_i=f_i\,dq/(2\pi iq)$ and $\nabla f_i=0$,
and the Hodge plane is the graph $(I,H_\sigma(q))$;
see \cite[Proposition~4.1, p.~12]{Viaclovsky2026Combs}.
Projection along this graph sends the logarithmic derivative to
$(2\pi i)^{-1}I+qH_\sigma'(q)$. At the cusp its determinant is
$(2\pi i)^{-2}\ne0$.

\end{enumerate}
\end{proof}






\begin{proposition}[Divisor twisting and embedded associated points]\hypertarget{tag-01GO}{}\label{cdp:st:E17}
\StatementRecord{cdp:st:E17}{proved}{01GO}{the definitions and constructions stated in the proof}

For an effective Cartier divisor $D$ on a compact complex
space $Z$ and a locally free sheaf $\mathcal V$, the spaces
$H^0(Z,\mathcal V(-kD))$, viewed inside $H^0(Z,\mathcal V)$,
stabilize for $k\gg0$ at
\[
 \{s\in H^0(Z,\mathcal V):\operatorname{supp}(s)\cap|D|
       =\varnothing\}.
\]
In particular they eventually vanish when $Z$ has no embedded
associated components and $D$ meets every irreducible component.
For the graph of a nonconstant meromorphic map from a smooth connected threefold
$Y$ to $\mathbb P^1$, every full fibre is isomorphic to a Cartier divisor in $Y$
and is Cohen--Macaulay with no embedded associated components.
\end{proposition}


\begin{proof}

\begin{enumerate}[label=\arabic*.,leftmargin=*]
\item Let $C=\mathbb P^1$, let $p,q$ be distinct points, and form the
square-zero complex space
\[
 |Z|=C,\qquad\mathcal O_Z=\mathcal O_C\oplus\mathbb C_p,
 \qquad(a,u)(b,v)=(ab,av+bu).
\]
Near $p$ its analytic local ring is
$\mathbb C\{z,e\}/(ze,e^2)$, and away from $p$ it is the
ordinary smooth curve. These charts glue through the quotient
by $e$, so this specifies a compact connected complex space.
It has one irreducible component, of positive dimension, and an
embedded associated point at $p$. The divisor $D=q$ is effective
Cartier: its local equation is a coordinate at $q$ and a unit
near $p$. For every integer $k$ there is the exact, split
description as an $\mathcal O_C$-module
\[
 \mathcal O_Z(-kD)=\mathcal O_C(-kq)\oplus\mathbb C_p.
\]
It gives every cohomology group:
\[
 h^0=\max(1-k,0)+1,\qquad h^1=\max(k-1,0),\qquad
 h^j=0\quad(j\geq2).
\]
In particular the nonzero section of the square-zero summand
survives for every $k\geq1$, contrary to
\citet[Proposition~1.1, p.~78]{CDP98}.

\item The descending spaces in the statement are
finite dimensional, hence stabilize. A section in their
intersection belongs to $\bigcap_k I_D^k\mathcal V_z=0$ at
every $z\in|D|$, by Krull intersection in the analytic local
ring. Its coherent support is therefore disjoint from $D$.
Conversely a section with support disjoint from $D$ vanishes
near $D$ and is a section of $\mathcal V(-kD)$ for every $k$.
When there are no
embedded associated components, every nonzero submodule of a
locally free sheaf has support containing an irreducible
component of $Z$. Such a support cannot avoid $D$ under the
stated intersection hypothesis.

\item Locally write the meromorphic map as $[a:b]$, where $a,b$ have
no common prime factor. Its graph $Z$ has equation
$av-bu=0$ in $Y\times\mathbb P^1$. At $c=[u_c:v_c]$, the maps
\[
 \{av_c-bu_c=0\}\longrightarrow Z_c,\qquad y\longmapsto(y,c),
\]
and the restriction of the projection to $Y$ are inverse holomorphic
maps of analytic spaces: both local rings are
$\mathcal O_Y/(av_c-bu_c)$. The local equation is a nonzero
element of the regular local ring of $Y$, hence a nonzerodivisor.
The quotient is Cohen--Macaulay and Gorenstein and has no embedded
associated components. This identifies the full analytic fibre,
including its nilpotents and multiplicities.

\end{enumerate}
\end{proof}



\paragraph{Audit of the printed non-holomorphic proof.}
We now give its entire chain of reductions and identify its
dependencies. CDP20, Section~3, first uses constancy of the index
under Picard twists and Riemann--Roch to deduce its parts (b)
and (c) from generic vanishing. Serre duality reduces generic
vanishing to degrees zero and two. Degree zero invokes CDP98,
Corollary~1.3; that corollary uses only Proposition~1.1 on the
smooth manifold, a multiple of an effective divisor lying in
$\operatorname{Pic}^0$, and semicontinuity. \ref{cdp:st:E17} proves the
version of Proposition~1.1 needed there, and \CompendiumRef{II}{cdp:st:E2} spells out the
exponential-family semicontinuity argument. Thus this dependency
is valid and involves none of the cotangent or monodromy steps.

For degree two the printed proof resolves the meromorphic pencil
by $\sigma:\widehat Y\to Y$, with a connected map
$h:\widehat Y\to\mathbb P^1$, and writes
\[
 h^*\mathcal O(1)=\sigma^*P\otimes\mathcal O(-E),
\]
where $E$ is the exceptional divisor of the pulled-back base
ideal. The map on second cohomology from $Y$ to $\widehat Y$
is an isomorphism, since a composition of smooth blowups has
$R\sigma_*\mathcal O=\mathcal O_Y$. The divisor sequence
reduces vanishing on $\widehat Y$ to vanishing after $-tE$
and to $H^2(tE,\sigma^*W)=0$. The printed proof asserts
the exceptional higher-image vanishing needed for this step.
For this particular pencil divisor it has the following direct
justification. Pull back the evaluation sequence on $\mathbb P^1$
and tensor by $\sigma^*P^{-t}$ to obtain, for $t\geq1$,
\[
\begin{split}
0\longrightarrow&\bigl(\mathcal O(E)\otimes
              \sigma^*P^{-(t+1)}\bigr)^{\oplus t}
 \longrightarrow(\sigma^*P^{-t})^{\oplus(t+1)}\\
 &\longrightarrow\mathcal O(-tE)\longrightarrow0.
\end{split}
\]
The evaluation sequence remains exact under pullback because it
is a sequence of vector bundles with locally free quotient.
The middle term has zero $R^2\sigma_*$, by the projection
formula and $R\sigma_*\mathcal O=\mathcal O_Y$.
Every coherent sheaf has zero $R^3\sigma_*$ because the
fibres have dimension at most two. Hence
$R^2\sigma_*\mathcal O(-tE)=0$; the divisor sequence now gives
\[
 R^1\sigma_*\mathcal O_{tE}=
 R^2\sigma_*\mathcal O_{tE}=0.
\]
The support of $\sigma_*\mathcal O_{tE}$ has dimension at most
one, so Leray and the projection formula give the required
$H^2(tE,\sigma^*W)=0$. No identification of this pushforward
with a structure sheaf is needed.

This special verification matters. The analogous assertion for
an arbitrary exceptional divisor is false. Blow up a point
$p\in\mathbb P^3$ and then a smooth plane cubic in the exceptional
$\mathbb P^2$. For the second exceptional divisor $S$, projection
to that cubic is a $\mathbb P^1$-bundle; thus
$h^\bullet(S,\mathcal O_S)=(1,1,0)$ and the composite contraction
to $\mathbb P^3$ has
$R^1\sigma_*\mathcal O_S=\mathbb C_p$.
This example, also given in \cite{LRS2019}, concerns arbitrary
exceptional divisors and does not refute the just-proved assertion
for the divisor of a resolved pencil. The graph proof of \CompendiumRef{II}{cdp:st:E2}
avoids the distinction by computing its own higher direct images
from its hypersurface equation.

The printed argument then replaces its twist $M$ by
$M\otimes P^{t+k}$ and applies Leray to $h$. It must prove
\[
 R^2h_*\sigma^*(W\otimes M\otimes P^k)=0
\]
and then the vanishing of $H^1$ of the corresponding $R^1$
twisted by $\mathcal O(t)$; the latter is ordinary Serre
vanishing on $\mathbb P^1$. For the former it uses top-degree
base change and fibrewise Serre duality. It invokes CDP98,
Proposition~1.1, first on components meeting $E$, and then
successively on components at increasing graph distance from
them. Coefficients in an auxiliary effective divisor are chosen
in reverse order of this distance, and the power of $E$ is
chosen last. \CompendiumRef{II}{cdp:st:E2} gives this argument with the actual pullback
maps, full multiplicity filtrations, and a finite-support
argument establishing uniformity. These fibres are Cartier in
a smooth threefold and have no embedded components, so \ref{cdp:st:E17}'s
proved version of the twisting lemma applies.

These are the only earlier CDP results invoked in Section~3:
CDP98, Corollary~1.3 and Proposition~1.1. Its opening reference
to CDP98, Theorem~2.1 identifies the proof being rewritten,
rather than using that theorem as a premise. CDP98, Lemma~1.5
is not invoked; its erroneous replacement of fibrewise vanishing
by generic vanishing is precisely what the full component
argument removes. No result from CDP20, Sections~4--7 is used.
In particular the cotangent-filtration and monodromy failures
treated in Parts A and D do not enter \CompendiumRef{II}{cdp:st:E2}. The supplied LRS
graph proof has the same dependency on the no-embedded-component
twisting lemma, and supplies its own higher-image calculation.
Its displayed adjunction formula has a sign misprint:
for a Cartier fibre $F\subset Y$ it is
$\omega_F=(K_Y+F)|_F$, not $(K_Y-F)|_F$.
\CompendiumRef{II}{cdp:st:E2} uses the actual dualizing line and includes nonreduced fibres,
so that misprint has no role in the proof here.




\begin{proposition}[Counterexample: The square theta formula for arbitrary sections]\hypertarget{tag-01GP}{}\label{cdp:st:E18}
\StatementRecord{cdp:st:E18}{proved}{01GP}{\SeriesRef{cdp:eq:E-theta-general}, \SeriesRef{cdp:st:E5}}

The square-only formula in Lead~9 is false as a local period
identity for arbitrary elliptic sections and linearizations.
The exact formula for the theta expression specified there is
\eqref{cdp:eq:E-theta-general}. The assertion for the actual original
period family remains true, by its equation $\wp(\mu)=0$.
\end{proposition}


\begin{proof}
\begin{enumerate}[label={$\langle1\rangle$\arabic*.},leftmargin=*]
\item 
\noindent\textbf{Step 1.} \noindent\textbf{Step 1.} We use $\vartheta_j(z,\tau)=\theta_j(\pi z\mid\tau)$ in the
notation of \cite[eqs.~20.2.1--2, 20.2.6 and 20.2.12]{DLMFTheta2026}.
Thus the nome is $q=e^{\pi i\tau}$, and the Weierstrass periods
are $2\omega_1=1$, $2\omega_3=\tau$. The products below are
\cite[eqs.~20.4.3--5 and 23.17.8]{DLMFTheta2026}, with this
specified change of theta argument.

\noindent\textbf{Step 2.} Let $\Delta$ be a sufficiently small disc about $i$ in the upper
half-plane and choose the holomorphic logarithm of the nonvanishing
function $\vartheta_1(1/2,\tau)/\eta(\tau)$ which is real at $i$.
Put
\[
 \mu(\tau)=\frac12,\qquad
 v(\tau)=\frac{i}{\pi}
       \log\frac{\vartheta_1(1/2,\tau)}{\eta(\tau)}+i,
 \qquad \beta(\tau)=6v(\tau).
\]
At $\tau=i$, the product expansions give
\[
 \frac{\vartheta_1(1/2,i)}{\eta(i)}
 =2e^{-\pi/6}\prod_{n\geq1}(1+e^{-2\pi n})^2>1.
\]
Thus $\operatorname{Im}v(i)>1$, and after shrinking $\Delta$
one has $\operatorname{Im}v>0$. Define
\[
 \Lambda_\tau=\mathbb Z(1,0)+\mathbb Z(0,1)
 +\mathbb Z(\tau,1/2)+\mathbb Z(3,\beta(\tau)),
 \qquad T_\tau=\mathbb C^2/\Lambda_\tau.
\]
The real determinant of these four periods is nonzero, since
$\operatorname{Im}\tau>0$ and $\operatorname{Im}\beta>0$.
They form a locally uniformly full lattice, so the quotients form
a proper holomorphic family of compact complex two-tori over
$\Delta$. Each fibre has
$h^\bullet(\mathcal O)=(1,2,1)$ and
$h^\bullet(T)=(2,4,2)$, with zero higher groups, by its
translation-invariant Dolbeault complex.

\noindent\textbf{Step 3.} \noindent\textbf{Step 2.} This family also has the required elliptic line-bundle
description. Quotient first by $(1,0),(0,1),(\tau,1/2)$.
Writing $z$ for the first coordinate and $a=\exp(2\pi iz_2)$
for the second gives the complement of zero in the flat line
bundle on $\mathbb C/(\mathbb Z+\tau\mathbb Z)$ with multipliers
$1$ and $-1$. For $P=1/2$, the actual meromorphic section
\[
 r(z)=\frac{\vartheta_1(z-P,\tau)}{\vartheta_1(z,\tau)}
\]
has divisor $P-O$ and satisfies $r(z+1)=r(z)$ and
$r(z+\tau)=-r(z)$. These equations identify the displayed
quotient line bundle with $\mathcal O(P-O)$. The last lattice
generator acts by
$(z,a)\mapsto(z+3,\exp(2\pi i\beta)a)$.
Its elliptic translation is $-6P=0$ in the elliptic curve;
its scalar is a holomorphic contracting linearization. This
constructs the family from an elliptic family, a section and a
linearization explicitly.

\noindent\textbf{Step 4.} \noindent\textbf{Step 3.} There is a specified relation to the symmetric period matrix
$\Omega=\left(\begin{smallmatrix}\tau&\mu\\\mu&v\end{smallmatrix}\right)$.
Adjoining $(\mu,v)$ enlarges $\Lambda_\tau$ to the lattice
generated by $(1,0),(0,1),(\tau,\mu),(\mu,v)$; its index is six,
because $6(\mu,v)=(3,\beta)$ and no smaller multiple belongs
to the original real lattice. Identity on $\mathbb C^2$ induces
a finite holomorphic isogeny of degree six between these families.
Its differential is the identity on the Lie spaces and its
pullback on complex first cohomology is an isomorphism. The
Kodaira--Spencer tensors therefore have the same rank under
these explicitly given comparisons.

\noindent\textbf{Step 5.} \noindent\textbf{Step 4.} The identity $\vartheta_1(1/2+z,\tau)=\vartheta_2(z,\tau)$
and evenness of $\vartheta_2$ give
$\partial_\mu\log\vartheta_1(1/2,\tau)=0$. Since $\mu'=0$,
the proposed square is zero. The actual determinant, by the
proved heat-equation formula, is
\cite[eqs.~20.7.5, 20.13.1, 23.3.9 and 23.6.2]{DLMFTheta2026}
\[
 \det\Omega'(i)=-\frac{\wp(1/2;i)}{4\pi^2}
 =-\frac{\vartheta_3(0,i)^4+\vartheta_4(0,i)^4}{12}\ne0.
\]
The last inequality follows from the theta products at the real
nome $e^{-\pi}$, whose factors give strictly positive theta
constants. Thus even the zero sets in the asserted formula
differ in this fully specified family. Extra holomorphic terms
in a different linearization contribute their derivatives to
the determinant; they must be retained, just as the $\wp$ term
must be retained when the section is not a zero of $\wp$.

\noindent\textbf{Step 6.} \noindent\textbf{Step 5.} For completeness, the direct images and total-space cohomology
in this local example are also determined. Write
$g:\mathscr T\to\Delta$ for the displayed proper family.
Translation-invariant forms, base change, and triviality of
holomorphic vector bundles on a disc give
\[
 R^qg_*\mathcal O_{\mathscr T}
   =\mathcal O_\Delta^{\binom2q},\qquad
 R^qg_*T_{\mathscr T/\Delta}
   =\mathcal O_\Delta^{2\binom2q}\quad(0\leq q\leq2),
\]
with all higher images zero. Shrink once more so that the
nonzero determinant stays nonzero. In the direct-image sequence
of $0\to T_{\mathscr T/\Delta}\to T_{\mathscr T}
\to g^*T_\Delta\to0$, the first boundary sends $1$ to the
nowhere-zero tensor $\kappa$, and the next boundary is the
rank-two isomorphism of \ref{cdp:st:E5}. Thus
\[
 R^qg_*T_{\mathscr T}=
 \begin{cases}
  \mathcal O_\Delta^2&q=0,\\
  \mathcal O_\Delta^3&q=1,\\
  \mathcal O_\Delta&q=2,\\
  0&q>2.
 \end{cases}
\]
\noindent\textbf{Step 6.} Since $\Delta$ is Stein, every total-space group
$H^q(\mathscr T,\mathcal F)$ for these three sheaves is exactly
the space of sections on $\Delta$ of the corresponding displayed
direct image. This specifies every cohomological degree without
assigning finite dimensions to the holomorphic-function spaces
on a disc.

\end{enumerate}
\end{proof}





\section{Connected deformation chains and two independent periods}\hypertarget{tag-034G}{}\label{sec:II:push-stability}
\begin{corollary}[arbitrary connected deformation chains]\hypertarget{tag-033C}{}\label{push:src:E.2.4}
\StatementRecord{push:src:E.2.4}{proved}{033C}{\SeriesRef{push:src:E.2.5}, \SeriesRef{push:src:E.2.24}, \SeriesRef{push:src:E.2.21}, \SeriesRef{push:src:E.2.22}, \SeriesRef{push:src:E.2.23}}

The ordered coprime pair $(3k,4l)$ is constant in every
proper smooth holomorphic family over a connected complex
analytic base which contains a constructed sphere member.
Two constructed members are deformation equivalent through
such families, allowing arbitrary intermediate members and
finite chains, if and only if their pairs agree.
Moreover different pairs cannot occur in one continuous
family in the $C^\infty$ space of integrable complex
structures on $S^6$.
\end{corollary}


\begin{proof}
Let $i:B_{\mathrm{red}}\hookrightarrow B$ be the reduction and form the
pullback $\mathcal X\times_B B_{\mathrm{red}}\to B_{\mathrm{red}}$.
The nilpotent closed immersion
$\mathcal X\times_B B_{\mathrm{red}}\hookrightarrow\mathcal X$
preserves the underlying topological spaces and the complex fibres;
properness and smoothness of the pulled-back family follow by base change. Theorem~\ref{push:src:E.2.23} therefore makes
the ordered multiplicity pair constant throughout each connected family.
At each join in a finite chain, the common member has the pair of both
adjacent families, so the pair is constant along the chain. The converse
for two constructed members with the same pair is also part of that theorem.

By Corollary~\ref{push:src:E.2.24}, the locus in the $C^\infty$ space of
integrable complex structures corresponding to a fixed pair is open and
closed. Its inverse image under a continuous family is open and closed
in the connected parameter space; if it contains one parameter, it
contains every parameter.
\end{proof}



\subsubsection*{Smooth specializations of the canonical pencil}

\begin{proposition}[a fractional canonical pencil at a smooth limit]\hypertarget{tag-033D}{}\label{push:src:E.2.5}
\StatementRecord{push:src:E.2.5}{proved}{033D}{\citet[Theorem~2.1]{CDP2020D}; \citet[Theorem~3.1 and Corollary~3.2]{LRS2019}}

Let $\pi:\mathcal X\to\Delta$ be a proper holomorphic submersion
whose fibres are diffeomorphic to $S^6$.  Suppose that for every
$t\ne0$ the fibre is one of the members of \SeriesRef{push:src:E.2.1}, with the same
coprime multiplicities $(3k,4l)$.  Set
\[
 A=12kl,\qquad D=kl-3k-4l\ne0.
\]
There is a unique holomorphic line bundle $\mathcal L$ on
$\mathcal X$, up to isomorphism, satisfying
\[
 \mathcal L^{D}=K_{\mathcal X/\Delta}^{A}.
\]
Its restriction to $X_t$ for $t\ne0$ is the coarse pencil
$f_t^*\mathcal O_{\mathbb P^1}(1)$.  The coherent sheaf
$\pi_*\mathcal L$ is locally free of rank two, and its central
fibre injects into $H^0(X_0,\mathcal L_0)$.

The central manifold has algebraic dimension one and a
holomorphic algebraic reduction $g_0:X_0\to\mathbb P^1$ with
connected fibres.  If $s_0,s_1$ are the two central sections
obtained from a local frame of $\pi_*\mathcal L$, there are a
nonconstant rational map $R:\mathbb P^1\to\mathbb P^1$ of
degree $e\ge1$ and an effective Cartier divisor $E_0$ such that
\[
 [s_0:s_1]=R\circ g_0,\qquad
 \mathcal L_0\simeq g_0^*\mathcal O(e)\otimes\mathcal O(E_0).
\]
Here $[s_0:s_1]$ denotes the holomorphic extension of the ratio;
the displayed divisor is its common fixed divisor.  The
statement does not assert that $E_0=0$.
\end{proposition}


\begin{proof}

\begin{enumerate}[label={$\langle1\rangle$\arabic*.},leftmargin=*]
\item By Ehresmann's theorem the smooth bundle over the contractible
disc is differentiably trivial.  Consequently
$H^1(\mathcal X,\mathbb Z)=H^2(\mathcal X,\mathbb Z)=0$.
The exponential sequence identifies
\[
 \operatorname{Pic}(\mathcal X)=H^1(\mathcal X,\mathcal O_{\mathcal X})
\]
as groups.  This is a complex vector space, so multiplication by
the nonzero integer $D$ is an automorphism.  This proves the
existence and uniqueness of the \emph{relative} root, including
when $D<0$.  The same argument on each fibre makes its Picard
group uniquely divisible.  \SeriesRef{push:src:E.2.1} gives
$K_{X_t}^{A}=(f_t^*\mathcal O(1))^{D}$, so uniqueness identifies
$\mathcal L_t$ with this actual pencil, not just with a numerical
class.

\item Proper direct-image coherence gives $\mathcal V=\pi_*\mathcal L$.
It is torsion-free: a local section killed by a nonzero function
of $t$ vanishes on the inverse image of an open set, hence
vanishes as a section of the invertible sheaf on the smooth
total space.  A coherent torsion-free sheaf on a smooth curve
is locally free.  Its rank is two because the noncentral
algebraic reductions have connected fibres.  Applying $\pi_*$
to
\[
 0\longrightarrow\mathcal L\xrightarrow{\,t\,}\mathcal L
 \longrightarrow i_*\mathcal L_0\longrightarrow0
\]
shows directly that $\mathcal V/t\mathcal V$ injects into
$H^0(X_0,\mathcal L_0)$.  Thus the two central sections are
linearly independent even if central cohomology jumps; no
surjectivity of base change has been assumed.

\item We use precisely the non-holomorphic-meromorphic-map theorem
of \citet[Theorem~2.1]{CDP2020D}, also written out in
\citet[Theorem~3.1 and Corollary~3.2]{LRS2019}.
For a compact complex threefold with $b_2=0$, a nonconstant
meromorphic map to a curve which fails to be holomorphic
would force $c_3=0$.  Since $c_3(X_0)[X_0]=2$, every
meromorphic function on $X_0$ extends as a holomorphic map
to $\mathbb P^1$.  This argument uses neither the different
holomorphic-fibration assertion of CDP, Theorem~2.2, nor its
Corollary~2.3.  Those assertions cannot be applied here.

\item In particular the ratio of $s_0,s_1$ gives a nonconstant
holomorphic map $h_0:X_0\to\mathbb P^1$.  Also $a(X_0)\le1$:
otherwise two algebraically independent meromorphic functions
would give a holomorphic map $X_0\to\mathbb P^1\times\mathbb P^1$
with two-dimensional image.  At a point of differential rank
two, after choosing finite affine target coordinates $u,v$
centred at its image, the global meromorphic function $u/v$
would have the local indeterminacy of the quotient of two
independent coordinates.  It cannot extend holomorphically
there.  This contradicts the preceding theorem.  Hence
$a(X_0)=1$.

\item Take the Stein factorization $h_0=R\circ g_0$.  Its intermediate
space is a smooth compact curve.  Connectedness of the fibres
injects its $H^1(\ ,\mathbb Z)$ into $H^1(X_0,\mathbb Z)=0$
by the degree-one Leray sequence, so it is $\mathbb P^1$.
Every meromorphic function on $X_0$ factors rationally through
$g_0$: the image of its holomorphic map together with $g_0$
is an analytic curve in $\mathbb P^1\times\mathbb P^1$;
normalizing this curve gives a finite factor of $g_0$, whose
degree is one because $g_0$ has connected fibres.  Thus $g_0$
is the algebraic reduction, and $R$ has some degree $e\ge1$.

\item Finally, on a target chart where the first homogeneous
coordinate of $h_0$ is nonzero, $s_1=h_0s_0$ in a local
trivialization, and the common ideal is generated by $s_0$.
The analogous assertion holds on the other chart.  These local
generators define a nonzero section of
$\mathcal L_0\otimes h_0^*\mathcal O(-1)$ and an effective
Cartier zero divisor $E_0$.  This gives the stated factorization
of $\mathcal L_0$ with no residual codimension-two base locus.

\end{enumerate}
\end{proof}



\begin{proposition}[positive canonical slope preserves the fibration]\hypertarget{tag-033E}{}\label{push:src:E.2.6}
\StatementRecord{push:src:E.2.6}{proved}{033E}{\SeriesRef{push:src:E.2.5}, \SeriesRef{push:src:E.2.22}, \SeriesRef{push:src:E.2.24}; \citet[Proposition~3.1.2]{AbashevaDeev2024AB}}

Under the hypotheses of \ref{push:src:E.2.5}, suppose $D>0$.  Then $E_0=0$
and $e=1$.  After shrinking the disc and choosing a frame of
$\pi_*\mathcal L$, its evaluation defines a proper flat
holomorphic map
\[
 g:\mathcal X\longrightarrow\mathbb P^1\times\Delta
\]
which is the intrinsic algebraic reduction on every fibre.
Its smooth general central fibre is a complex two-torus.

In particular, two families of the type in \ref{push:src:E.2.5} with positive
canonical slope cannot have the same smooth central manifold
unless their multiplicity pairs agree.  A positive-slope and
a negative-slope family cannot have a common smooth central
manifold either.  This proposition concerns a common smooth
specialization.  The additional propagation and closure
arguments for arbitrary connected chains are \ref{push:src:E.2.22}--\ref{push:src:E.2.24}.
\end{proposition}


\begin{proof}

\begin{enumerate}[label={$\langle1\rangle$\arabic*.},leftmargin=*]
\item Choose a very general smooth fibre $F$ of $g_0$, avoiding the
images of all vertical components of $E_0$ and the jumping
sets for the pluricanonical direct images below.  Its normal
bundle in $X_0$ is trivial.  Adjunction gives
$K_F=K_{X_0}|_F$, and $c_1(K_F)=0$ integrally because
$H^2(X_0,\mathbb Z)=0$.  In particular $F$ is minimal: a
$(-1)$-curve would have canonical intersection $-1$.
Restriction of the root identity in \ref{push:src:E.2.5} gives
\[
 K_F^{A}\simeq\mathcal O_F(D E_0|_F).
\]
The right hand side is effective since $D>0$, so
$\kappa(F)\ge0$.

\item We next show $\kappa(F)\le0$, without assuming that the fibres
of an algebraic reduction have algebraic dimension zero.
If $\kappa(F)>0$, some $r>0$ has $h^0(F,K_F^r)\ge2$.
Our very-general choice ensures that this is the generic rank
of the coherent sheaf $g_{0*}K_{X_0}^r$.  A sufficiently
positive twist of this sheaf on $\mathbb P^1$ is globally
generated.  Two of its global sections can therefore be
chosen to be linearly independent on a general fibre.  By
projection formula they are sections of
$K_{X_0}^r\otimes g_0^*\mathcal O(N)$; their ratio is a
meromorphic function on $X_0$ which is nonconstant on $F$.
This contradicts the algebraic-reduction assertion of \ref{push:src:E.2.5}.
The choice of a very general fibre is legitimate: for each
$r$ the exceptional set for the generic cohomology rank is
a proper analytic subset of the base curve, and there are
only countably many $r$.

\item Thus $\kappa(F)=0$.  The classification of minimal compact
complex surfaces makes $K_F$ torsion; the relevant consequence
is also explicitly recorded in the proof of
\citet[Proposition~3.1.2]{AbashevaDeev2024AB}.
An effective divisor whose line bundle is torsion is zero:
raising its canonical section to a torsion-killing power
gives a holomorphic function on a compact connected manifold,
and that function is a nonzero constant.  Therefore
$E_0|_F=0$.  Every horizontal component of $E_0$ would meet
such a general fibre in a nonempty effective divisor.
Consequently $E_0$ has only vertical components, supported
over finitely many points of $\mathbb P^1$.

\item We now use connectedness of actual divisor limits to remove
both this vertical fixed part and the possible degree $e$.
Choose a linear combination $s=as_0+bs_1$ for which its
moving central divisor is the union of the $e$ distinct
smooth fibres over $R^{-1}([-b:a])$, all disjoint from the
support of $E_0$.  Also choose it to be a regular fibre on
one specified noncentral member.  Both requirements omit
only finitely many choices in the projective line of
combinations.  Extend it by the corresponding combination
of the frame of $\pi_*\mathcal L$.

\item Its total zero divisor $\mathcal D\subset\mathcal X$ is a
relative effective Cartier divisor: the section is not
identically zero on any fibre, by independence of the two
sections, including at zero.  In local coordinates, $t$ is
a non-zero-divisor modulo its defining equation.  Hence
$\mathcal D\to\Delta$ is proper and flat.  Its general
noncentral fibre is a smooth connected torus.  There are no
vertical irreducible components of $\mathcal D$, since such
a component would have to be the whole irreducible
threefold $X_t$.  It follows that $\mathcal D$ is reduced
and irreducible: its unique generic component has multiplicity
one, and a Cartier divisor in a smooth manifold has no
embedded components.  Normalize $\mathcal D$ and take the
Stein factorization over $\Delta$.  The finite factor has
degree one over the punctured generic locus; normality of
the disc makes it the disc itself.  Its central fibre, and
therefore the image $\mathcal D_0$, is connected.

\item But the support of
\[
 \mathcal D_0=E_0+g_0^*R^*([-b:a])
\]
is a disjoint union of the $e$ chosen smooth connected fibres
and, if nonempty, the support of $E_0$ over other base points.
Connectedness forces $E_0=0$ and $e=1$.

\item The central two sections now generate $\mathcal L_0$ at
every point.  The noncentral two sections already generate
their pencils, so the evaluation
$\pi^*\pi_*\mathcal L\to\mathcal L$ is surjective after
shrinking.  It defines the asserted map $g$.  It is proper
because $\pi$ is proper.  On each $X_t$ the nonconstant
map to the smooth curve $\mathbb P^1$ is flat: its local
structure sheaf is torsion-free over the local discrete
valuation ring of the curve.  The total source is flat
over $t$, so the local criterion for flatness applied to
the parameter $t-t_0$ gives flatness over $(z,t)$.
Equivalently, the local fibre equations form a regular
sequence in the smooth source.  At a regular
point of the central base curve, properness makes $g$ a
proper submersion on a neighbourhood of the corresponding
compact fibre.  Ehresmann's theorem identifies that central
surface differentiably with nearby complex two-tori.
The classification of compact complex surfaces, or its
special case for the differentiable four-torus, then makes
the central surface a complex two-torus as well.

\item Suppose now two positive-slope families have the same
central manifold $X_0$.  Its algebraic reduction is intrinsic
up to $\operatorname{PGL}_2$, so the preceding result makes
both central pencil lines the same non-torsion line
$L=g_0^*\mathcal O(1)$.  With indices $i=1,2$ the identities
$K_{X_0}^{A_i}=L^{D_i}$ give
$L^{D_1A_2-D_2A_1}=\mathcal O$.  The line $L$ is non-torsion
because a positive power has a nonzero section with zeros.
Thus $D_1/A_1=D_2/A_2$, and the exact recovery of the ordered
pair from this slope in \SeriesRef{push:src:E.2.3} gives the assertion.

\item Finally, a common smooth limit from opposite signs would
have a nonzero section of a positive power of $K_{X_0}$
and of a negative power, by upper semicontinuity (or by
the same injective direct-image argument).  Raising them
to a common power, their nonzero product is a constant;
both are nowhere zero.  Hence $K_{X_0}$ is torsion.
The exponential sequence and $H^1(X_0,\mathbb Z)=
H^2(X_0,\mathbb Z)=0$ make $\operatorname{Pic}(X_0)$
torsion-free, so $K_{X_0}$ would be trivial.  Its nonzero
holomorphic three-form $\Omega$ is closed and exact because
$H^3(X_0,\mathbb C)=0$.  Stokes' theorem would give
$\int_{X_0}\Omega\wedge\overline\Omega=0$, whereas its
appropriate constant multiple is a positive volume form.
This contradiction excludes the common limit.

\end{enumerate}
\end{proof}



\paragraph{Scope after \ref{push:src:E.2.6}}

\ref{push:src:E.2.5} proves that the central meromorphic reduction is
holomorphic, and \ref{push:src:E.2.6} removes its fixed divisor and degree
loss for $D>0$.  The positive-slope proof alone does not
exclude the effective antipluricanonical divisor allowed
by the restriction identity when $D<0$.  \ref{push:src:E.2.8}--\ref{push:src:E.2.10}
remove that possibility by the horizontal-divisor vanishing
argument.  \ref{push:src:E.2.13} and \ref{push:src:E.2.16}--\ref{push:src:E.2.21} then recover the finite
charts, cusp, and actual smooth completion.  The separate
generic-disc propagation and closure arguments in
\ref{push:src:E.2.22}--\ref{push:src:E.2.24} prove \ref{push:src:E.2.4} for arbitrary analytic bases,
deformation chains, and continuous complex structures.
The present proposition retains its stated one-sided,
positive-slope scope.

\Uses{the hypotheses, notation and definitions stated here; \cref{push:src:E.2.6}; \cref{push:src:E.2.5}; \cref{push:src:E.2.8}; \cref{push:src:E.2.10}; \cref{push:src:E.2.13}; \cref{push:src:E.2.16}; \cref{push:src:E.2.21}; \cref{push:src:E.2.22}; \cref{push:src:E.2.24}; \cref{push:src:E.2.4}}


\begin{proposition}[an extended fibration cannot reach the zero period]\hypertarget{tag-033F}{}\label{push:src:E.2.7}
\StatementRecord{push:src:E.2.7}{proved}{033F}{\SeriesRef{push:src:E.2.5}, \SeriesRef{push:src:E.2.6}, \SeriesRef{push:src:E.2.20}, \SeriesRef{push:src:E.2.21}, \SeriesRef{push:src:E.2.22}, \SeriesRef{push:src:E.2.24}}

Under the hypotheses of \ref{push:src:E.2.5}, suppose in addition that the
intrinsic algebraic reductions extend to a holomorphic map
\[
 g:\mathcal X\longrightarrow\mathbb P^1\times\Delta
\]
whose restriction to $X_0$ is nonconstant.  This hypothesis
holds in \ref{push:src:E.2.6}.  Normalize the punctured family to fixed
discrete logarithmic data with multiplicities $(3k,4l)$, and
write its isomorphism parameter as
\[
 z=\exp\!\left(-\frac{\pi i c}{kl}\right).
\]
The resulting holomorphic function $z:\Delta^*\to\mathbb C^*$
extends across zero and satisfies $z(0)\ne0$.  Its central
value lies strictly inside the admissible period disc.

This statement concerns the period function and the extended
relative fibration.  Identification of $X_0$ with the member
at $z(0)$ is proved separately in \ref{push:src:E.2.20}--\ref{push:src:E.2.21}; propagation
to arbitrary deformation bases and continuous families is
proved in \ref{push:src:E.2.22}--\ref{push:src:E.2.24}.
\end{proposition}


\begin{proof}

\begin{enumerate}[label={$\langle1\rangle$\arabic*.},leftmargin=*]
\item Local miniversality in \SeriesRef{push:src:E.2.1} supplies holomorphic local
classifying functions for the family on $\Delta^*$.
\SeriesRef{push:src:E.1.2}--\SeriesRef{push:src:E.1.3} normalize the discrete data by holomorphic
family comparisons.  For fixed data the comparison shifts
are precisely $c\mapsto c+2kl\mathbb Z$, so the displayed
exponential is a single-valued holomorphic function on
$\Delta^*$.  An automorphism with the same marked parameter
has identity linear action on the homology of an ordinary
torus: in the necessity argument of \SeriesRef{push:src:E.1.2} its parameters
are $\alpha=1$ and $h=0$.  Consequently the homology local
system of the punctured family is also the one obtained
from these local classifying maps; possible translation
automorphisms of the comparisons add no homology monodromy.

\item The admissible parameter domain is a punctured disc
$0<|z|<R_{k,l}$, with finite positive radius.  The removable
singularity theorem therefore extends $z$ holomorphically
to $\Delta$.  It remains to exclude a zero at the centre.

\item We first control the moving three critical values without
assuming that they remain distinct at zero.  Since
$\pi:\mathcal X\to\Delta$ is a submersion, the critical
locus of $g$ is the analytic locus where its differential
has rank less than two.  Properness makes its image
\(\mathcal C\subset\mathbb P^1\times\Delta\) a closed
analytic set.  Its intersection with the central target
curve is the finite critical-value set of the nonconstant
map $g_0:X_0\to\mathbb P^1$.  In particular
$\mathcal C$ has no component equal to that whole curve.
After shrinking $\Delta$, its reduced one-dimensional
part is finite over $\Delta$ and over $\Delta^*$ consists
of the three critical-value sections.

\item Those sections are individually marked by the linear
peripheral monodromy orders $3$, $4$, and infinity.
The normalization of the closure of each section is
therefore finite of degree one over the disc.  It is
isomorphic to the disc, so the three values
$c_1(s),c_2(s),c_0(s)$ extend as holomorphic maps to
$\mathbb P^1$, with collisions allowed at $s=0$.
The unique projective transformation $A_s$ sending this
ordered triple to $(0,1,\infty)$ has meromorphic
coefficients at zero, by its explicit cross-ratio formula.
This is the critical-value argument of \SeriesRef{push:src:B.4.4} in a setting
where no vertical discriminant component is present.

\item Choose a point $v\in\mathbb P^1$ outside all central
critical values.  Properness and closedness of
$\mathcal C$ allow us to shrink the disc so that
$(v,s)\notin\mathcal C$ for every $s$.  Then
\[
 \mathcal T=g^{-1}(\{v\}\times\Delta)\longrightarrow\Delta
\]
is a proper smooth holomorphic family of surfaces.
Its noncentral fibres are ordinary complex two-tori, and
Ehresmann's theorem makes its integral first-homology
monodromy around $s=0$ the identity.

\item In the original marked base, this same family is obtained
by evaluating the period construction at
\[
 t(s)=A_s(v)\in\mathbb P^1\setminus\{0,1,\infty\}.
\]
The function $t(s)$ is meromorphic at zero.  Its loop on a
sufficiently small positive $s$-circle is therefore either
null-homotopic in the marked base or conjugate to a power
of one of its three peripheral loops.  Indeed, if its
limit is ordinary, its image lies in an ordinary small
disc; if its limit is a marked point, a local coordinate
has a zero or pole of finite order and its image lies
in the corresponding punctured disc.  This observation
also covers a constant $t(s)$.

\item Let $\Lambda$ have ordered basis
$(\widehat\gamma,\widehat u,\widehat w,\widehat\delta)$,
put $E=\widehat\delta\otimes\gamma$, and write
$S_a=I+aE$.  The two finite peripheral monodromies are
\[
 A_1=\begin{pmatrix}
 1&0&0&0\\6&0&1&0\\-6&-1&-1&0\\-2&1&0&1
 \end{pmatrix},\qquad
 A_2=\begin{pmatrix}
 1&0&0&0\\0&0&-1&0\\-6&1&0&0\\3&0&1&1
 \end{pmatrix}.
\]
They have orders three and four, respectively, even for
the larger multiplicities: the extra finite deck
transformations are translations and do not affect
$H_1$.  In the cusp marking the third monodromy is
\[
 A_0=I+N,\qquad
 N\widehat u=\widehat w,\quad
 N\widehat\gamma=-\widehat\delta,\quad
 N\widehat w=N\widehat\delta=0.
\]
All peripheral matrices preserve $\widehat\delta$ and
$\gamma$, so commute with every $S_a$.  The change to
the cusp marking preserves the same rank-one operator
$E$; conjugating a peripheral loop consequently does
not alter the cancellation test below.

\item Suppose that $z(s)$ has a zero of order $d>0$ at zero.
Continuation of a logarithm around the positive loop
changes $c$ by $-h$, where
\[
 h=2kl\,d>0.
\]
Since $\Pi_{c+h}=\Pi_cS_h$, a comparison of the identical
complex universal covers from parameter $c-h$ back to
$c$ acts on the marked lattice by $S_{-h}$.
The affine translation part of a closing comparison
acts trivially on $H_1$.  Thus the return map of the
family $\mathcal T$ has homology action, up to a common
conjugation,
\[
 S_{-h},\qquad S_{-h}A_1^n,\qquad
 S_{-h}A_2^n,\qquad\hbox{or}\qquad S_{-h}A_0^n
\]
for an integer $n$.  Reversing the convention for a
peripheral generator only replaces $n$ by $-n$.
The parameter factor and the base factor commute, so
the order in this display has no effect.

\item None of these matrices is the identity.  In the two
finite cases, taking the third or fourth power would
give $S_{-3h}=I$ or $S_{-4h}=I$, impossible because
$E\ne0$ and $h>0$.  In the cusp case, $N^2=EN=NE=0$
and hence
\[
 S_{-h}A_0^n=I+nN-hE.
\]
Its coefficient of $\widehat w$ on $\widehat u$ is
$n$, so identity would force $n=0$, and then $h=0$.
The ordinary case is $S_{-h}\ne I$.  This contradicts
the identity monodromy of the smooth family over
the full disc.  Therefore $z(0)\ne0$.

\item Finally, the extended holomorphic function satisfies
$|z(0)|\le R_{k,l}$.  Equality would make it a constant
of maximal modulus, contradicting $|z(s)|<R_{k,l}$
for $s\ne0$.  Thus $0<|z(0)|<R_{k,l}$, as claimed.
The matrix identities and the factor $h=2kld$ are
checked exactly in \path{code/e_push_E_zero_end.py}.

\end{enumerate}
\end{proof}



\paragraph{Scope and subsequent completion}

For either canonical slope, \ref{push:src:E.2.10} supplies the
fibration-persistence hypothesis of this proposition,
so a one-sided smooth specialization has an interior
period value.  Equality of that value with an existing
model does not by itself prove uniqueness of its smooth
completion.  \ref{push:src:E.2.11}--\ref{push:src:E.2.12} exclude critical values
appearing only at the central parameter; \ref{push:src:E.2.13}--\ref{push:src:E.2.20}
recover the actual finite and cusp charts, and \ref{push:src:E.2.21}
proves the full one-sided identification.  The distinct
problem of propagation through arbitrary deformation
branches is settled by \ref{push:src:E.2.22}--\ref{push:src:E.2.24}.

\Uses{the hypotheses, notation and definitions stated here; \cref{push:src:E.2.10}; \cref{push:src:E.2.11}; \cref{push:src:E.2.12}; \cref{push:src:E.2.13}; \cref{push:src:E.2.20}; \cref{push:src:E.2.21}; \cref{push:src:E.2.22}; \cref{push:src:E.2.24}}


\subsubsection*{A restriction on the negative-slope fixed divisor}

\begin{proposition}[horizontal divisors force non-K\"ahler fibre normalizations]\hypertarget{tag-033G}{}\label{push:src:E.2.8}
\StatementRecord{push:src:E.2.8}{proved}{033G}{\SeriesRef{push:src:E.2.5}, \SeriesRef{push:src:E.2.6}; \citet[Lemma~4.2]{CDP2020D}; \citet[Lemma~3.3 and Corollary~3.4]{Fujino2021E}}

Let $X$ be a compact complex threefold diffeomorphic to $S^6$, and
let $f:X\to\mathbb P^1$ be a surjective holomorphic map with
connected fibres.  Suppose that a smooth general fibre $F$ has
$b_1(F,\mathbb Q)=1$.  Then every reduced fibre of $f$ is
irreducible, and the monodromy on $H^1(F,\mathbb Q)$ is trivial.

If $H\subset X$ is a horizontal prime divisor, the normalization
of every reduced fibre lies outside Fujiki's class $\mathcal C$.
Equivalently, no such normalization admits a K\"ahler resolution.
In particular no fibre has Moishezon normalization.

Without the hypothesis $b_1(F)=1$, every irreducible component
of a \emph{reducible} fibre of any such connected map from a
complex $S^6$ also has normalization outside $\mathcal C$.

Consequently, in \ref{push:src:E.2.5} with $D<0$, a horizontal component of
$E_0$ is possible only if every reduced fibre of the central
algebraic reduction $g_0$ is irreducible and has normalization
outside $\mathcal C$.  This necessary condition does not exclude
non-K\"ahler normalizations and does not prove persistence of the
relative pencil in the negative-slope case.
\end{proposition}


\begin{proof}

\begin{enumerate}[label={$\langle1\rangle$\arabic*.},leftmargin=*]
\item Choose a nonempty finite set $\Sigma\subset\mathbb P^1$ containing
all critical values; a regular value can be added if necessary.
Write $s=|\Sigma|$, put $A=f^{-1}(\Sigma)_{\mathrm{red}}$,
and let $r$ be the number of irreducible components of $A$.
The compact analytic set $A$ is triangulable.  Each of its
irreducible two-dimensional components has its complex
fundamental class, and their intersections have real dimension
at most two.  Hence
\[
 H^4(A,\mathbb Q)\simeq\mathbb Q^r.
\]
Alexander duality in the oriented six-sphere gives
\[
 b_1(X\setminus A,\mathbb Q)=r.
\]
On the other hand, $f:X\setminus A\to U:=\mathbb P^1\setminus\Sigma$
is a proper smooth fibre bundle.  The base $U$ has the homotopy
type of a graph with $s-1$ loops.  The degree-one Leray--Serre
sequence therefore gives an exact sequence
\[
 0\longrightarrow H^1(U,\mathbb Q)
 \longrightarrow H^1(X\setminus A,\mathbb Q)
 \longrightarrow H^1(F,\mathbb Q)^{\pi_1(U)}
 \longrightarrow0.
\]
In particular, the correct formula retains the invariant
subspace:
\[
 r=s-1+\dim H^1(F,\mathbb Q)^{\pi_1(U)}\le s.
\]
Every nonempty fibre over $\Sigma$ has at least one irreducible
component, so $r\ge s$.  Equality follows throughout.  Thus
every fibre over $\Sigma$ has precisely one reduced irreducible
component and the invariant subspace is all of the
one-dimensional $H^1(F,\mathbb Q)$.  Fibres outside $\Sigma$ are
smooth and connected, hence irreducible.  This proves both
assertions without assuming monodromy triviality at the outset.
It is the rank-one repair of the component-count argument;
the unrestricted formula in \citet[Lemma~4.2]{CDP2020D} is not
being used.

\item Let $S=f^{-1}(t)_{\mathrm{red}}$.  A horizontal prime $H$ maps
surjectively to the base by properness, so $H\cap S\ne\varnothing$.
Since $S$ is irreducible, $S\subset H$ would imply $S=H$,
contrary to horizontality.  The canonical section of
$\mathcal O_X(H)$ consequently restricts to a nonzero section
on the reduced irreducible space $S$, with a nonempty zero
set.  Its zero divisor is a nonzero effective Cartier divisor
on $S$.  Indeed the section cannot vanish identically on
any local irreducible branch: its pullback would then vanish
on an open subset of the irreducible normalization, hence
everywhere by analytic continuation.  In the reduced local
ring it consequently avoids every minimal prime and is a
non-zero-divisor, including at self-intersections.

\item Let $\nu:S^\nu\to S$ be the normalization.  If $S^\nu$ were
in $\mathcal C$, it would have a compact K\"ahler resolution
$\mu:\widehat S\to S^\nu$; this precise surface statement is
\citet[Lemma~3.3 and Corollary~3.4]{Fujino2021E}.
The pulled-back line bundle
\[
 L=(\nu\mu)^*(\mathcal O_X(H)|_S)
\]
has $c_1(L)=0$ integrally, since $H^2(X,\mathbb Z)=0$.
Its pulled-back section is nonzero and has zeros: normalization
is finite surjective, and modification is proper surjective.
Thus it defines a nonzero effective divisor $B$ on
$\widehat S$.  For a K\"ahler form $\omega$,
\[
 0=\langle c_1(L)\smile[\omega],[\widehat S]\rangle
   =\int_B\omega>0,
\]
a contradiction.  The argument applies to the reduced support
even when the original fibre is nonreduced.  Moishezon spaces
are in $\mathcal C$, so the stated consequence follows.

\item For the assertion about a reducible fibre, let $S_i$ be one
of its reduced irreducible components.  The fibre is connected,
so another component $S_j$ meets $S_i$; otherwise $S_i$ would
be an open-and-closed component of its support.  Distinctness
ensures $S_i\not\subset S_j$.  Apply exactly the preceding
canonical-section and positivity argument to
$\mathcal O_X(S_j)|_{S_i}$.  Its first Chern class is zero
because $H^2(X,\mathbb Z)=0$, whereas its effective Cartier
zero divisor is nonempty.  This proves the additional assertion
without any hypothesis on the general fibre's first Betti number.

\item For the application, choose a very general smooth fibre $F$
of $g_0$ meeting a horizontal component of $E_0$ and avoiding
its vertical components.  As in \ref{push:src:E.2.6}, adjunction gives
$c_1(K_F)=0$ and $F$ is minimal.  The identity of \ref{push:src:E.2.5} now
restricts to
\[
 K_F^{-A}=\mathcal O_F((-D)E_0|_F),
 \qquad -D>0,
\]
with a nonzero effective divisor on the right.  If a positive
power $K_F^m$ had a nonzero section, multiplying a suitable
power of that section by the $m$th power of the displayed
anticanonical section would give a nonzero holomorphic
function with zeros on $F$.  Compactness forbids this.
Thus $\kappa(F)=-\infty$.  Moreover $F$ is not K\"ahler,
since its nonzero effective divisor $(-D)E_0|_F$ has zero
real Chern class.  The Enriques--Kodaira classification then
places $F$ in class VII and gives $b_1(F)=1$; see
\citet[Chapter~VI, Theorem~1.1 and Table~10, pp.~243--244]
{BHPV2004E}.  The preceding argument applies with $f=g_0$
and any horizontal prime component $H$ of $E_0$.

\end{enumerate}
\end{proof}



The proof supplies a necessary restriction on an unknown
smooth specialization.  It neither classifies the singular
non-K\"ahler normalizations nor replaces the unresolved
non-normal-fibre step in the holomorphic-fibration argument
of \citet[\S7]{CDP2020D}.  That step is not used here:
\ref{push:src:E.2.9}--\ref{push:src:E.2.10} prove the needed horizontal-divisor vanishing,
and \ref{push:src:E.2.21}--\ref{push:src:E.2.24} complete the actual specialization and
deformation-chain arguments.

\subsubsection*{Removing the negative-slope horizontal fixed part}

\begin{proposition}[no horizontal divisor when all fibres are irreducible]\hypertarget{tag-033H}{}\label{push:src:E.2.9}
\StatementRecord{push:src:E.2.9}{proved}{033H}{\SeriesRef{push:src:E.2.8}; \citet[\S3, Theorem~1, p.~81, and the complete-intersection
calculation on p.~91]{RR70push}; \cite[Chapter~III, Theorem~4.12(ii)]{BS76push}; \cite[Chapter~III, Theorem~4.1]{BS76push}}

Let $X$ be a compact complex threefold diffeomorphic to $S^6$,
and let $f:X\to\mathbb P^1$ be a surjective holomorphic map
with connected fibres.  If every reduced fibre is irreducible,
then $X$ has no horizontal prime divisor for $f$.
\end{proposition}


\begin{proof}

\begin{enumerate}[label={$\langle1\rangle$\arabic*.},leftmargin=*]
\item Suppose that $H$ is such a divisor.  It is Cartier because
$X$ is smooth.  Since $H$ is horizontal and every reduced
fibre $S$ is irreducible, $H$ meets $S$ and does not contain
it.  As in \ref{push:src:E.2.8}, restriction of its canonical section is a
non-zero-divisor on every local branch of $S$.

\item We first record the elementary vanishing used below.  On a
compact reduced irreducible complex space $S$, let $D$ be a
nonzero effective Cartier divisor and $V$ a vector bundle.
Multiplication by the powers of the canonical section
identifies
\[
 H^0(S,V(-nD))\subset H^0(S,V),\qquad n\ge0,
\]
with a decreasing sequence of finite-dimensional vector
spaces.  A section in their intersection vanishes to every
order along $D$.  Krull intersection in a local ring at a
point of $D$ makes its germ zero.  After pullback to the
irreducible normalization, analytic continuation makes the
section zero everywhere.  The decreasing sequence therefore
stabilizes at zero.  In particular
\[
 H^0(S,V(-nD))=0\quad\hbox{for all sufficiently large }n.
\]
The same conclusion holds for a vector bundle on a thickened
fibre $mS$.  Indeed $S$ is Cartier in $X$, and the finite
filtration by powers of its ideal has successive quotients
\[
 V|_S\otimes\mathcal O_S(-jS),\qquad 0\le j<m.
\]
Apply the preceding reduced-space statement to these finitely
many vector bundles, with the same nonzero divisor $H|_S$,
and then use their exact sequences.  This also proves that
the multiplication maps used here remain injective on $mS$:
its only associated component is $S$ and $H$ contains none
of its local branches.

\item Let $B$ be any holomorphic vector bundle on $X$.  The map
$f$ is flat, since its local rings are torsion-free over the
discrete valuation rings of the smooth base.  Every analytic
fibre $F_t$ is a Cartier, Cohen--Macaulay, Gorenstein surface,
possibly nonreduced.  Its dualizing line is
\[
 \omega_{F_t}=K_X|_{F_t},
\]
where the constant one-dimensional factor from the target
has been trivialized.  Serre duality on this compact
Gorenstein surface gives
\[
 H^2(F_t,B(nH)|_{F_t})^*
 \simeq H^0\bigl(F_t,(B^*\otimes K_X)(-nH)|_{F_t}\bigr).
\]
The compact-space duality being used includes singular
Cohen--Macaulay spaces; it is the absolute duality theorem
of \citet[\S3, Theorem~1, p.~81, and the complete-intersection
calculation on p.~91]{RR70push}, already specified in \SeriesRef{push:src:E.3.1},
not a smooth-fibre replacement.
The preceding filtration argument makes the right hand
side zero for all sufficiently large $n$, for each $t$.

\item For fixed $n$, the nonvanishing locus
\[
 Z_n=\{t\in\mathbb P^1:
        h^2(F_t,B(nH)|_{F_t})>0\}
\]
is closed by semicontinuity.  Multiplication on the dual
section spaces gives $Z_{n+1}\subset Z_n$.  Their intersection
is empty by fibrewise vanishing.  Compactness of $\mathbb P^1$
therefore supplies a single $N$ with $Z_N=\varnothing$.
Constant-zero cohomology and analytic cohomology-and-base-change
\cite[Chapter~III, Theorem~4.12(ii)]{BS76push} give
\[
 R^2f_*(B(NH))=0.
\]
Higher direct images vanish in degrees greater than two by
the fibre dimension.  The coherent sheaf $R^1f_*(B(NH))$
has vanishing first cohomology after a sufficiently positive
twist on $\mathbb P^1$.  The projection formula and Leray
sequence consequently give, for some integer $a\gg0$,
\[
 H^2(X,B\otimes M_+)=H^3(X,B\otimes M_+)=0,
 \qquad
 M_+=\mathcal O_X(NH)\otimes f^*\mathcal O(a).
\]
All coherence, semicontinuity and constant-zero base-change
uses are the proper analytic theorems audited in \SeriesRef{push:src:E.3.1};
the argument does not infer base change from coherence alone.

\item Independently, the reduced-space vanishing argument applied
to the connected smooth manifold $X$ and its nonzero divisor
$H$ gives
\[
 H^0(X,B\otimes M_-)=0,
 \qquad M_-=\mathcal O_X(-bH),\quad b\gg0.
\]
We must arrange these different vanishings for one line
bundle.  The exponential sequence gives
\[
 \operatorname{Pic}(X)\simeq H^1(X,\mathcal O_X)
\]
because $H^1(X,\mathbb Z)=H^2(X,\mathbb Z)=0$.
Let $\ell_-,\ell_+$ be the corresponding additive classes.
The affine line
\[
 \ell(z)=(1-z)\ell_-+z\ell_+,\qquad z\in\mathbb C,
\]
is realized by a holomorphic family of line bundles on
$X\times\mathbb C$: choose additive holomorphic Cech
cocycles representing the two classes and exponentiate their
linear combination.  Equivalently, use the integrable
Dolbeault operators obtained by a holomorphic linear
combination of fixed closed $(0,1)$-form representatives.
Thus no representability or universal-Picard-space assertion
is needed here.

\item For the proper projection $X\times\mathbb C\to\mathbb C$,
the cohomology jumping loci
\[
 T_j=\{z:h^j(X,B\otimes M_z)>0\},\qquad j=0,2,3,
\]
are closed analytic subsets of $\mathbb C$ by the proper
analytic finite-complex theorem
\cite[Chapter~III, Theorem~4.1]{BS76push}: the condition is
expressed by the rank minors in a local finite free complex.
In particular no closedness claim from the corrected
Grauert Section~7, Satz~2 is used.  The set $T_0$ omits $z=0$,
while $T_2$ and $T_3$ omit $z=1$.  Each is therefore proper,
and their finite union cannot fill the irreducible complex
line.  Choose $z$ outside that union.  Then
\[
 \chi(X,B\otimes M_z)=-h^1(X,B\otimes M_z)\le0.
\]
Every line bundle on $X$ is topologically trivial, so
Hirzebruch--Riemann--Roch makes this Euler characteristic
equal to $\chi(X,B)$.

\item Finally take $B=T_X$.  Since $H^2(X,\mathbb Z)=H^4(X,\mathbb Z)=0$,
the only contributing term in its Riemann--Roch formula is
\[
 \chi(X,T_X)=\frac12\langle c_3(X),[X]\rangle
            =\frac12\chi_{\mathrm{top}}(S^6)=1.
\]
This contradicts the preceding inequality.  Hence no
horizontal prime $H$ exists.

\end{enumerate}
\end{proof}



\begin{proposition}[persistence for either canonical slope]\hypertarget{tag-032V}{}\label{push:src:E.2.10}
\StatementRecord{push:src:E.2.10}{proved}{032V}{\SeriesRef{push:src:E.2.5}, \SeriesRef{push:src:E.2.7}, \SeriesRef{push:src:E.2.4}, \SeriesRef{push:src:E.2.6}, \SeriesRef{push:src:E.2.8}, \SeriesRef{push:src:E.2.9}}

Under the hypotheses of \ref{push:src:E.2.5}, for every nonzero
$D=kl-3k-4l$ one has $E_0=0$ and $e=1$.
After shrinking the disc, the evaluation of
$\pi_*\mathcal L$ gives a proper flat holomorphic map
\[
 g:\mathcal X\longrightarrow\mathbb P^1\times\Delta
\]
whose restriction to every fibre is its intrinsic algebraic
reduction.  A smooth general central fibre of $g$ is a
complex two-torus.  Its normalized punctured period has
an interior limit as in \ref{push:src:E.2.7}.

Two families satisfying \ref{push:src:E.2.5} cannot have the same smooth
central manifold unless their ordered multiplicity pairs
agree, regardless of the signs of their canonical slopes.
Uniqueness of a smooth completion at an interior period and
the arbitrary deformation-chain statement \ref{push:src:E.2.4} are not
asserted.
\end{proposition}


\begin{proof}

\begin{enumerate}[label={$\langle1\rangle$\arabic*.},leftmargin=*]
\item The case $D>0$ is \ref{push:src:E.2.6}.  If $D<0$ and $E_0$ had a
horizontal prime component, the final part of \ref{push:src:E.2.8} would
give $b_1(F)=1$ for a smooth general fibre of $g_0$ and
would make every reduced fibre of $g_0$ irreducible.
\ref{push:src:E.2.9} would then forbid precisely that horizontal component.
Thus for either sign $E_0$ is vertical.

\item The connected-divisor-limit argument in the middle of the
proof of \ref{push:src:E.2.6} only uses verticality, not the sign of $D$.
Its generic choice of pencil section has central zero
support equal to the disjoint union of the vertical fixed
part and $e$ distinct regular fibres.  The total divisor
is reduced and irreducible with connected central support
after normalization and Stein factorization.  Consequently
$E_0=0$ and $e=1$.  Surjectivity of evaluation, flatness and
the torus assertion follow from the remainder of that same
proof.  The hypothesis of \ref{push:src:E.2.7} is now satisfied for either
sign, and gives the interior period limit.

\item Suppose two such one-sided families have the same central
manifold $X_0$.  Its algebraic reduction is intrinsic, so
both pencil lines are the same line
$L=g_0^*\mathcal O(1)$ up to isomorphism.  With indices
$i=1,2$, the two root identities give
\[
 K_{X_0}^{A_i}=L^{D_i},\qquad
 L^{D_1A_2-D_2A_1}=\mathcal O.
\]
The line $L$ is not torsion, because a positive power has
a nonzero section with zeros.  Hence $D_1/A_1=D_2/A_2$.
\SeriesRef{push:src:E.2.3} recovers the ordered coprime multiplicity pair from
this rational number.  This proves the common-limit
assertion without identifying the limit with a member of
the given period construction.

\end{enumerate}
\end{proof}



The vanishing argument is an application of standard proper
cohomology, duality and Riemann--Roch machinery; no novelty
claim is intended.  It proves the needed absence of a
horizontal fixed divisor under the rank-one hypothesis.
It does not establish the different claim that every
holomorphic Hopf-surface fibration on a complex six-sphere
is impossible.

\subsubsection*{An extension criterion at an isolated base point}

\begin{proposition}[purity with a smooth torus Jacobian]\hypertarget{tag-032W}{}\label{push:src:E.2.11}
\StatementRecord{push:src:E.2.11}{proved}{032W}{the definitions and constructions stated in the proof}

Let $B\subset\C^2$ be a sufficiently small bidisc about $0$, put
$U=B\setminus\{0\}$, and let $g:X\to B$ be a proper flat
holomorphic map from a connected smooth complex $(n+2)$-fold,
with connected fibres of dimension $n$.  Let $G\to B$ be a
smooth proper holomorphic family of complex $n$-tori with a
holomorphic zero section and its group law.  Suppose that
$X_U\to U$ is a holomorphic $G_U$-torsor.  Then $g$ is a
submersion on all of $X$, and the given action extends uniquely
to make $X\to B$ a holomorphic $G$-torsor.

No polarization, K\"ahler metric on $X$, or section of the
given torsor on $U$ is assumed.  The smooth extension $G$ is
an explicit hypothesis.
\end{proposition}


\begin{proof}

\begin{enumerate}[label={$\langle1\rangle$\arabic*.},leftmargin=*]
\item Shrink $B$ so that the holomorphic Lie bundle of $G$ has a
frame $e_1,\ldots,e_n$.  Its period local system has a basis
of holomorphic sections
\[
 \lambda_1(b),\ldots,\lambda_{2n}(b)\in\C^n,
 \qquad
 G_b=\C^n/\sum_{a=1}^{2n}\Z\lambda_a(b).
\]
Indeed the lattice is locally constant over the simply
connected bidisc, and the local inverses of the relative
exponential map make its sections holomorphic.  The periods
are a real basis at every $b$.

\item The infinitesimal torsor action gives holomorphic vector
fields $V_1,\ldots,V_n$ on $X_U$.  Flatness makes
$Z=g^{-1}(0)$ have complex codimension two in $X$.
Hartogs extension in local coordinates therefore extends
each coefficient of each $V_j$ uniquely across $Z$; uniqueness
glues these local extensions to holomorphic vector fields
on $X$.  This use of Hartogs applies to functions on the
complement of an analytic set of codimension at least two
in a complex manifold, hence also coefficientwise to any
locally free sheaf.  It does not assert extension across a
divisor.  The identities
\[
 dg(V_j)=0,\qquad [V_i,V_j]=0
\]
hold on the dense open subset $X_U$ and consequently on $X$.

\item Every $V_j$, and every linear combination with coefficients
pulled back from $B$, is complete.  Its integral curve stays
in a compact fibre; a finite endpoint of a maximal solution
would have an accumulation point in that fibre, where the
holomorphic ordinary differential equation extends the
solution.  The flows thus define a holomorphic action of
$\C^n$ on $X$ over $B$.

\item For $a=(a_1,\ldots,a_{2n})\in\R^{2n}$ take the time-one
flow of the vertical holomorphic vector field whose Lie
coordinate is
\[
 \sum_{j=1}^{2n}a_j\lambda_j(g(x)).
\]
These maps depend smoothly on $a$ and add under composition:
all their coefficients are constant along every vertical
flow.  For $a\in\Z^{2n}$ the map is the identity on $X_U$,
hence everywhere by holomorphic uniqueness.  We obtain an
action of the fixed compact real torus
$T=\R^{2n}/\Z^{2n}$ by holomorphic automorphisms of $X$,
over $B$.  On the fibre over $b$, it is the holomorphic
action of $G_b$ obtained by factoring the $\C^n$ action by
its period lattice.  The action of $T$ is free on $X_U$.
The variable period identities also make the holomorphic
Lie-bundle action descend to a holomorphic action
$G\times_BX\to X$.  Its restriction therefore acts on
each scheme-theoretic fibre, including its nilpotent
structure.

\item We show that its stabilizers have dimension zero everywhere.
Let $x\in X$, $b=g(x)$, and let $H$ be the identity component
of the stabilizer.  Viewed in $G_b$, it is a compact complex
subgroup: it is the identity component of the inverse image
of $x$ under the holomorphic orbit map $G_b\to X_b$.
Put
\[
 K_x=\ker(dg_x)=T_x(X_b),
\]
where the right-hand side is the Zariski tangent space of
the scheme-theoretic fibre.  Differentiating the holomorphic
$G_b$ action on this fibre gives a holomorphic representation
$H\to\operatorname{GL}(K_x)$.  It is trivial, because every
holomorphic function on the compact connected complex group
$H$ is constant.  The induced representation on
$T_xX/K_x$ is also trivial: this quotient injects into
$T_bB$ by $dg_x$, and every group element acts over $B$.
Averaging a Hermitian inner product over the compact group
$H$ gives an invariant complement to $K_x$.  The full
derivative representation on $T_xX$ is therefore trivial.

\item Here the required local linearization has a direct proof.
Choose holomorphic coordinates $z$ centred at $x$ and a
sufficiently small $H$-invariant neighbourhood on which all
$z(hy)$ are defined.  Such a neighbourhood exists by
compactness of $H$ and continuity of the action.  With
normalized Haar measure set
\[
 \Psi(y)=\int_H z(hy)\,dh.
\]
This map is holomorphic, invariant under $H$, and has
derivative the identity at $x$.  It is a holomorphic
coordinate system after shrinking.  Its invariance then
forces every $h\in H$ to be the identity near $x$, and hence
on connected $X$.  This contradicts freeness on $X_U$ unless
$H$ is trivial.  Thus every stabilizer is finite, and the
$n$ vector fields $V_j$ are linearly independent everywhere.

\item Their commuting flows give local holomorphic coordinates
$(z,s)\in\C^n\times\C^2$ in which they span the $z$
directions: choose a transverse complex surface and apply
the holomorphic inverse function theorem to the flow map.
In these coordinates $g$ has the form
\[
 g(z,s)=h(s).
\]
Equidimensionality implies that the germ $h^{-1}(b)$ is
zero-dimensional.  Consequently the germ
$h:(\C^2,0)\to(B,b)$ is finite; this is the local finiteness
criterion obtained by Weierstrass division for an isolated
fibre.  If $dg_x$ failed to have rank two, the determinant
of $dh$ would vanish at $s=0$.  It is not identically zero,
since the map is a submersion away from $Z$.  Its zero set
would therefore contain a complex curve through $0$.
Finiteness sends that curve to a curve of critical values.
This contradicts smoothness over $B\setminus\{0\}$.
Hence $g$ is a submersion everywhere.

\item The $G_b$ orbits now have the full dimension of every smooth
connected fibre and are compact.  Each fibre is a single
orbit.  Its finite stabilizer acts trivially on its tangent
space, because $G_b$ is abelian; it also acts trivially on
the base tangent quotient.  Averaging as above, now over
that finite group, again makes its full derivative trivial
and then makes the action locally trivial.  Freeness on
$X_U$ forces the stabilizer to be trivial.  Thus the extended
action is free and transitive on every fibre.  Local
holomorphic sections of the submersion identify $X$ with
$G$ over small base opens, which proves the torsor assertion.
The extension is unique because it agrees with the original
one on the dense open subset $X_U$.

\end{enumerate}
\end{proof}



\paragraph{Scope and literature check}

This extension criterion has precisely the displayed hypotheses. The proof does not establish purity for every proper flat map.  In particular it does
not apply to B.4.15: that construction has vertical
threefolds over the bidisc origin and is not flat over the
bidisc.  Searches on 15 November 2026 for purity and
discriminants of complex-torus fibrations, Hartogs extension
of torus actions, and isolated critical values did not
locate the exact statement under these hypotheses without
an added K\"ahler assumption.  The complete argument is given above.

\Uses{the hypotheses, notation and definitions stated here; the argument and qualifications immediately following}


\begin{proposition}[no additional central critical values]\hypertarget{tag-032X}{}\label{push:src:E.2.12}
\StatementRecord{push:src:E.2.12}{proved}{032X}{\SeriesRef{push:src:E.2.10}, \SeriesRef{push:src:E.2.7}, \SeriesRef{push:src:E.2.11}}

In the smooth-specialization setting of \ref{push:src:E.2.10}, let
\[
 g:\mathcal X\longrightarrow\mathbb P^1\times\Delta
\]
be the extended intrinsic fibration, and let
$c_1(s),c_2(s),c_0(s)$ be the three labelled critical-value
sections of \ref{push:src:E.2.7}.  The critical values of $g_0$ are exactly
the distinct points among $c_1(0),c_2(0),c_0(0)$.
In particular, no additional critical value occurs only
on the central target curve.

If the normalizing projective transformations $A_s$ of
\ref{push:src:E.2.7} have a degenerate limit, the central regular fibres
are isomorphic complex two-tori and $g_0$ has at most two
critical values.  If the constant image of the limiting
rank-one matrix is $0$ or $1$, this torus is the
corresponding finite-root abelian surface.  If that image
is an ordinary source point, it is the ordinary torus at
that point and need not be an abelian surface.  The cusp
cannot be the constant image.
\end{proposition}


\begin{proof}

\begin{enumerate}[label={$\langle1\rangle$\arabic*.},leftmargin=*]
\item \ref{push:src:E.2.7} gives $0<|z(0)|<R_{k,l}$.  A local logarithm
therefore extends the normalized parameter $c(s)$
holomorphically to the disc.  The critical image is closed
and contains the three punctured sections, so it contains
their three limiting points.  Suppose that a different
central critical value $b_0=(v_0,0)$ exists.  Take a small
bidisc $B$ about $b_0$ which avoids all three section
closures and all other central critical values.  Put
$U=B\setminus\{b_0\}$.  The restricted family over $U$
is a proper smooth family of complex two-tori.  We verify
that its Jacobian extends to a smooth holomorphic torus
group family over $B$, so that \ref{push:src:E.2.11} applies.

\item The normalized source coordinate is $t=A_s(v)$.
Represent $A_s$ by a matrix with holomorphic entries,
after clearing poles and removing a common power of
$s$, so its central matrix is nonzero.  If that matrix
is invertible, $t$ extends holomorphically near $b_0$ and
its value there is ordinary: a marked value would put
$b_0$ on the corresponding section closure.  The
ordinary period model with parameter $c(s)$ then gives
the desired smooth Jacobian over $B$.

\item Otherwise the central matrix has rank one.  Denote its
constant image by $t_*$ and its exceptional target point
by $v_*$.  For each marked source value different from
$t_*$, its preimage tends to $v_*$.  At least two of the
three marked values are different from $t_*$, so $v_*$
belongs to at least two labelled section closures.
Consequently $v_0\ne v_*$.  After shrinking $B$, the
function $t(s,v)$ tends uniformly to $t_*$ there.
For an ordinary $t_*$ the ordinary period model again
extends directly.

\item Suppose instead that $t_*=0$ or $1$, and use a local
coordinate $\xi$ at that marked point.  The divisor of
$\xi(t(s,v))$ on $B$ is supported on $s=0$, since $B$
avoids the corresponding critical-value graph.
Consequently
\[
 \xi(t(s,v))=s^n u(s,v),\qquad n>0,
\]
where $u$ is a nowhere-zero holomorphic function.
This also explains why no zero of the leading
coefficient may be overlooked: such a zero would put
the point on the labelled graph closure.
Choose a nearby $v_1$ for which $(v_1,0)$ is regular.
The actual family over the full $s$-disc at $v_1$ is
smooth, so its integral degree-one monodromy is the
identity.  The parameter $c(s)$ has no winding.  By
the marked comparison of \ref{push:src:E.2.7}, this monodromy is
$A_1^n$ or $A_2^n$, respectively.  Their orders are
three and four, so $e\mid n$, where $e=3$ or $4$.
A holomorphic $e$th root of $u$ exists on the bidisc,
and
\[
 r(s,v)=s^{n/e}u(s,v)^{1/e}
\]
is an actual holomorphic root coordinate on $B$.
Pull back the smooth finite-root torus period model
by $(c(s),r(s,v))$.  This produces the required
smooth torus group family $G\to B$.  The larger
multiplicities $3k,4l$ do not change this Jacobian
argument: their additional deck transformations
are translations, and the linear Jacobian model
is the original order-three or order-four root model.

\item For completeness $t_*=\infty$ is impossible.  In the
cusp coordinate the same divisor argument gives
$q=s^n u$ with $n>0$.  The smooth family at $v_1$
would have monodromy $A_0^n$.  But
$A_0=I+N$, $N^2=0$, and
$N\widehat u=\widehat w\ne0$, so $A_0^n\ne I$.

\item It remains to identify $G|_U$ with the actual Jacobian,
rather than merely comparing individual tori.
The punctured bidisc $U$ is simply connected.  Its
integral first-homology local system therefore has
a global marking.  Over $s\ne0$ that marking agrees
with the reference period marking, after one constant
integral change of basis; the transition comparisons
of \SeriesRef{push:src:E.1.3} have identity linear homology action at
fixed marked parameter.  Both holomorphic period
subbundles extend on $U$, and they agree on its dense
open subset $s\ne0$.  They agree everywhere on $U$.
Equivalently, the two quotients of their holomorphic
Lie bundles by the same marked period lattice are
isomorphic torus group families on $U$.  This is a
Jacobian identification and makes no assertion that
the actual affine torsor has a section on $B$.

\item The total space is smooth and $g$ is proper and flat.
The only possible nonsmooth fibre over $B$ lies above
$b_0$, and its codimension in the total space is two.
All hypotheses of \ref{push:src:E.2.11} now hold.  That lemma makes
$g$ smooth also above $b_0$, a contradiction.  The
central critical-value set is therefore precisely
the set of labelled limits.

\item For the final assertions, in the rank-one case every
ordinary central target point different from $v_*$
has the same limiting marked torus: it is obtained
from $(c(0),t_*)$, or from the finite-root value
$(c(0),r=0)$.  Equality of the extended period
subbundles proves this assertion without assuming
separatedness for threefold completions.
All marked values except possibly $t_*$ have their
preimages tending to $v_*$.  Hence there are at most
two central critical values; if $t_*$ is ordinary,
there is exactly one.
At $t_*=0$ or $1$, the finite linear automorphism of
the root torus has complex eigenvalues $1$ and a
primitive third or fourth root of unity.  The
rational idempotents for these two distinct
cyclotomic factors split that torus, up to isogeny,
as a product of two elliptic curves.  It is therefore
an abelian surface.  This argument makes no
projectivity assertion for an ordinary limiting
torus, and supplies no polarization on the whole
central threefold.

\end{enumerate}
\end{proof}



\paragraph{Scope of \ref{push:src:E.2.11}--\ref{push:src:E.2.12}}

Purity here is obtained from the explicit good
Jacobian near a hypothetical extra critical value.
Smooth total space and flatness alone were not used
as a general purity theorem.  These results do not
give uniqueness of an interior smooth completion,
do not exclude the remaining isotrivial central
fibrations, and do not prove persistence of the
canonical pencil on an arbitrary new deformation
branch.

\Uses{the hypotheses, notation and definitions stated here; \cref{push:src:E.2.11}; \cref{push:src:E.2.12}}


\begin{proposition}[persistence of an isolated finite multiple fibre]\hypertarget{tag-032Y}{}\label{push:src:E.2.13}
\StatementRecord{push:src:E.2.13}{proved}{032Y}{\SeriesRef{push:src:E.2.10}, \SeriesRef{push:src:E.2.12}, \SeriesRef{push:src:E.2.11}}

In the setting of \ref{push:src:E.2.10}, suppose the central value of
one finite labelled critical section is distinct from
the limits of the other two labelled sections.  Write
its multiplicity as $m$, with $m=3k$ or $4l$, and its
linear monodromy order as $e=3$ or $4$, respectively.
Near that section, normalize the base change taking
an $m$th root of its defining equation.

The resulting total space is smooth, and its map to
the root bidisc is a proper smooth family of two-tori.
The deck group acts freely on the total space.
Consequently the actual central reduced finite fibre
is smooth, has exactly multiplicity $m$, and belongs
to the same smooth family of finite free affine
quotients as the punctured reduced finite fibres.

This applies to both finite fibres when the three
labelled critical values have distinct limits.  In
the two-critical-value collapsed-target case, it
applies to the isolated finite labelled fibre.
\end{proposition}


\begin{proof}

\begin{enumerate}[label={$\langle1\rangle$\arabic*.},leftmargin=*]
\item Choose target coordinates $(v,s)$ in a bidisc about
the section, and let $\Gamma$ be its graph.  Take a
local defining function $\xi$ for $\Gamma$.
On the smooth fourfold the divisor identity is
\[
 g^*\Gamma=m\mathcal D,
\]
where $\mathcal D$ is the prime closure of the
punctured reduced multiple-fibre family.  Indeed
this equality holds away from $s=0$.  The only
possible extra divisorial component would be
$X_0$, but $g_0$ is nonconstant, so $X_0$ is not
mapped into the point $\Gamma(0)$.  The closure
$\mathcal D$ is Cartier because the total space
is smooth.  Thus locally, even if its central
restriction is initially singular or reducible,
\[
 g^*\xi=a^m u
\]
for a local equation $a$ of $\mathcal D$ and a
holomorphic unit $u$.

\item Let $B'\to B$ be the smooth root bidisc
$r^m=\xi$, and let $Y$ be the normalization
of $\mathcal X_B\times_B B'$.  A local $m$th
root of $u$ factors its defining equation as
\[
 r^m-a^m u=
 \prod_{\zeta^m=1}(r-\zeta a u^{1/m}).
\]
Its normalized branches are separate copies of
the original smooth local fourfold.  Therefore
$Y\to\mathcal X_B$ is a finite etale cover,
and $Y$ is smooth.  The deck group permutes these
branches freely, including over $a=0$.
This is an assertion about the normalized cover;
the unnormalized fibre product generally has
intersecting branches.

\item The induced map $g':Y\to B'$ is proper and
equidimensional of relative dimension two.
Since both source and base are smooth, the
two pulled-back base parameters form a regular
sequence in each local ring, giving flatness.
Over $s\ne0$ this is the actual smooth root
torus family.  Its connectedness makes $Y$
connected and its generic fibre connected.
The Stein factor of $g'$ is finite birational
over the normal bidisc and hence is the
bidisc itself, so all fibres are connected.
Equivalently this follows from normality of
$Y$ and the finite birational form of its
Stein factorization.

\item By \ref{push:src:E.2.12} the original central map has no
critical value away from the labelled limits.
Our choice of the section therefore makes
$g'$ smooth away from the origin of $B'$.
It remains to supply a smooth Jacobian over
the full root bidisc.

\item If the normalizing projective transformations
$A_s$ have an invertible limit, the source
coordinate at the finite marked value has
the form
\[
 \xi_{\rm src}=\xi\,u_1(v,s)
\]
with $u_1$ a unit.  Thus its order-$e$ root
is $r^{m/e}u_1^{1/e}$, and the exact finite-root
period model with parameter $c(s)$ extends
smoothly to the full root bidisc.

\item If $A_s$ has rank-one limit, the isolated
section in the hypothesis must be the
preimage of its constant image $t_*$:
all other marked preimages tend to the
exceptional point.  Thus $t_*$ is precisely
the finite marked value under consideration.
Near the isolated section the source
coordinate has divisor equal to that section
plus a positive multiple of $s=0$, so
\[
 \xi_{\rm src}=s^n\xi\,u_1(v,s),\qquad n>0,
\]
with a nowhere-zero unit $u_1$.  Choose a
nearby central regular target point with
$\xi\ne0$.  Its actual smooth torus family
has identity parameter monodromy.  As in
\ref{push:src:E.2.12}, $c(s)$ has no winding and hence
$e\mid n$.  The required holomorphic root
coordinate on $B'$ is now
\[
 s^{n/e}r^{m/e}u_1^{1/e}.
\]
It again gives a smooth reference Jacobian
over $B'$.  On the punctured bidisc its
marking and holomorphic period subbundle
agree with the actual Jacobian by the same
simply-connected argument as \ref{push:src:E.2.12}.

\ref{push:src:E.2.11} makes $g'$ a smooth torsor over the
entire root bidisc.  In particular $r$ is
a submersion coordinate on the central
threefold upstairs.  The unramified deck
cover then shows that the original fibre
has exactly multiplicity $m$.  Restricting
to $r=0$ gives a proper smooth family of
two-tori in the parameter $s$, with free
affine deck action.  Its quotient is a
proper smooth family of the reduced finite
fibres, proving all assertions.

\end{enumerate}
\end{proof}



\paragraph{Scope of the finite-chart argument}

This argument identifies the smooth root charts
at the isolated finite critical values.  The
excluded target-collision cases are treated in
\ref{push:src:E.2.14}--\ref{push:src:E.2.15}.  The actual cusp is recovered
separately in \ref{push:src:E.2.16}--\ref{push:src:E.2.20}, and \ref{push:src:E.2.21} glues
these comparisons into the full one-sided
completion theorem.

\Uses{the hypotheses, notation and definitions stated here; \cref{push:src:E.2.14}; \cref{push:src:E.2.15}; \cref{push:src:E.2.16}; \cref{push:src:E.2.20}; \cref{push:src:E.2.21}}


\begin{proposition}[extension of the circle field and exclusion of finite-root collapse]\hypertarget{tag-032Z}{}\label{push:src:E.2.14}
\StatementRecord{push:src:E.2.14}{proved}{032Z}{\SeriesRef{push:src:E.2.10}, \SeriesRef{push:src:E.2.12}, \SeriesRef{push:src:E.2.11}}

In the setting of \ref{push:src:E.2.10}, the normalized vector field
$\delta$ of the punctured family extends to a holomorphic
relative vector field on $\mathcal X$ after shrinking
the parameter disc.  Its time-one flow is the identity,
and its real circle action on the central fibre is
equivariantly diffeomorphic to the prescribed circle
action on the punctured fibres.  In particular its
central fixed locus is a smooth rational curve with
normal bundle $\mathcal O(-1)\oplus\mathcal O(-1)$
and the same two opposite circle weights.

The rank-one limit of the normalizing projective
transformations $A_s$ cannot have constant image
$0$ or $1$.  More generally, a collapsed constant
image is impossible whenever the limiting torus's
complex $\delta$ line has two real-linearly
independent periods.  Thus the finite-root
isotrivial possibilities in \ref{push:src:E.2.12} are excluded.
\end{proposition}


\begin{proof}

\begin{enumerate}[label={$\langle1\rangle$\arabic*.},leftmargin=*]
\item Let $\pi:\mathcal X\to\Delta$ be the original
proper holomorphic submersion.  The coherent sheaf
\[
 \mathcal H=\pi_*T_{\mathcal X/\Delta}
\]
is torsion-free: multiplication by a nonzero base
function is injective on the vector bundle upstairs
and hence on its sections.  Its generic rank is
one by \SeriesRef{push:src:E.2.1}.  On the smooth disc it is therefore
locally free of rank one.  Choose a generator $V$
after shrinking the disc.  On $\Delta^*$ it has
the form
\[
 V|_{X_s}=b(s)\delta_s,
\]
with $b$ holomorphic and nowhere zero there.
The normalized fields $\delta_s$ agree on overlaps
of local classifying maps: at fixed discrete data
the comparison maps of \SeriesRef{push:src:E.1.3} have fibre-linear
part the identity and fix the marked period
$\widehat\delta$.  Their translations also
preserve the invariant vector field.
Moreover $V$ is vertical for $g$, because this
holds on the dense punctured family.

\item Choose a regular central target value and a small
bidisc around it on which $g$ is smooth.  The
ordinary or finite-root comparison used in
\ref{push:src:E.2.12} supplies its smooth marked Jacobian, with
the marked period $\widehat\delta$ extending as
a nonzero holomorphic Lie section.  The resulting
invariant field $\delta_B$ on the actual torsor
is nonvanishing there.  Comparing one nonzero
coefficient of $\delta_B$ with $V$ shows that
$b(s)$ extends holomorphically to $s=0$.
Its independence of the target coordinate follows
from that equality on $s\ne0$ and holomorphic
uniqueness.

\item If $b(0)=0$, the restriction $V|_{X_0}$ vanishes
on a nonempty open subset of connected $X_0$,
and therefore vanishes everywhere.  In local
coordinates for the submersion $\pi$, every
coefficient of $V$ is then divisible by $s$.
The quotients glue uniquely to a global
holomorphic relative vector field $V/s$.
This contradicts the fact that $V$ generates
the free module $\mathcal H$.  Hence $b(0)\ne0$,
and $V/b(s)$ gives the required global extension
of $\delta$.

\item The extended field is complete on each compact
fibre of $\pi$.  Its time-one flow is the
identity on the punctured family and hence on
the whole family by holomorphic uniqueness.
Its real flow consequently defines an action
of $\mathbb R/\mathbb Z$ on $\mathcal X$
over the parameter disc.  Average a smooth
metric and an Ehresmann horizontal distribution
over this compact group.  Parallel transport
then gives equivariant diffeomorphisms of
the fibres.  The central action is therefore
the same effective semifree action, with a
single fixed two-sphere, as on the prescribed
fibres.

\item The fixed set of a holomorphic compact-group
action is a complex submanifold, as follows
from the averaging coordinates used in
\ref{push:src:E.2.11}.  Here it is a proper smooth family
of curves of genus zero, and hence each fibre
is $\mathbb P^1$.  Its two holomorphic normal
weight lines have the same degrees as in
the punctured family.  Those degrees are
$-1,-1$ by C.3.11, or directly by the two
fixed toric edge normals.  The characters
distinguish the lines, so their central
restrictions are exactly
$\mathcal O(-1)\oplus\mathcal O(-1)$.

\item Suppose now that $A_s$ has rank-one limit
with constant image a finite marked value.
\ref{push:src:E.2.12} identifies every central regular
torus, with its marked $\delta$ direction,
with the corresponding finite-root torus
at parameter $c(0)$.  The invariant lattice
of $A_j$ has rank two and spans its
one-dimensional complex fixed subspace.
It contains $\widehat\delta$; its second
primitive period is independent over
$\mathbb R$.  The complex fixed subspace
is precisely the $\delta$ line, since
the other complex eigenvalue is a
primitive third or fourth root of unity.
Consequently, in the normalization in
which the $\delta$ period is $1$, there
is a constant second period
\[
 \kappa\in\mathbb C\setminus\mathbb R.
\]
This is a period of the invariant field
on every central regular fibre.  The
time-$\kappa$ flow is therefore the
identity on their dense union, and
hence on all of $X_0$.

\item The flow of the nonzero field $\delta_0$
thus factors through a nontrivial
holomorphic action of the compact
elliptic curve
$E=\mathbb C/(\mathbb Z+\kappa\mathbb Z)$.
Such an action has no fixed point on
a connected complex manifold unless
it is trivial.  Indeed its derivative
at a fixed point is a holomorphic
representation of the compact complex
group into a general linear group,
so is constant.  Averaging local
coordinates over the compact group
then makes the action trivial near
that point, and holomorphic uniqueness
makes it trivial everywhere.

\item The elliptic action here is nontrivial
because $\delta_0\ne0$.  It therefore
has no fixed point.  This contradicts
the fixed rational curve of the
extended real circle action: a zero
of the generating field is fixed
by its entire complex flow.  Equivalently,
the resulting nowhere-zero real vector
field would give $\chi(X_0)=0$,
whereas $X_0$ is diffeomorphic to $S^6$.
The same argument applies to an ordinary
constant-image torus if its $\delta$
line has a second independent period.

\end{enumerate}
\end{proof}



\paragraph{What the elliptic-action argument does not assume}

The argument uses the global one-dimensional
space of vector fields on the known punctured
members and the actual marked period comparison
on a central regular bidisc.  It does not
assume that all translations of an ordinary
central torus extend over the singular fibre.
For an ordinary torus the $\delta$ line may
have only the one specified real period;
such a case is not excluded by the
elliptic-action argument alone.

\Uses{the hypotheses, notation and definitions stated here; the argument and qualifications immediately following}


\begin{proposition}[an isotrivial volume obstruction]\hypertarget{tag-0330}{}\label{push:src:E.2.15}
\StatementRecord{push:src:E.2.15}{proved}{0330}{\SeriesRef{push:src:E.2.12}, \SeriesRef{push:src:E.2.10}}

Let $f:X\to\PP^1$ be a proper holomorphic map with connected
fibres from a compact connected complex threefold.  Suppose
that for some $p\in\PP^1$ the restriction over
$U=\PP^1\setminus\{p\}$ is a smooth family of mutually
isomorphic complex two-tori.  If
\[
 K_X^{\otimes A}\simeq f^*\mathcal O_{\PP^1}(D),
 \qquad A\in\Z_{>0},\quad D\in\Z,
 \tag{E.2.15.1}
\]
then necessarily $D<-A$.  In particular the ordinary-source
collapsed-target case of \ref{push:src:E.2.12} cannot occur for the
six-sphere families under consideration.
\end{proposition}


\begin{proof}

\begin{enumerate}[label={$\langle1\rangle$\arabic*.},leftmargin=*]
\item Choose an affine coordinate $t$ on $U\simeq\C$ and put
$q=t^{-1}$ near $p$.  Since $U$ is simply connected, its
integral first-homology local system has a global marking.
The holomorphic period subbundle is constant: locally the
marked periods belong to the countable integral-change-of-basis
orbit of one fixed torus, and a continuous map from a connected
base into such a countable subset of a complex Grassmannian
is constant.  For the last assertion, any two distinct image
points can be separated by a real continuous coordinate;
its connected image would contain an interval and would be
uncountable.  Thus the Jacobian group family is the product
with one fixed torus $T$.

\item Choose a nonzero translation-invariant holomorphic two-form
$\eta_T$ on $T$.  Local holomorphic origins identify our
family with its Jacobian, and two choices differ by a
translation.  Hence $\eta_T$ defines a global nonvanishing
relative canonical form $\eta$ on $X_U$.  Its fibre volume
is a fixed positive number:
\[
 \int_{X_t}\frac14\eta\wedge\overline\eta=V>0.
 \tag{E.2.15.2}
\]
This construction uses neither a global torsor section nor
a polarization on $X$.  The canonical determinant
isomorphism for a submersion gives the nonvanishing
absolute three-form $\eta\wedge dt$ on $X_U$; this notation
does not require choosing an absolute lift of $\eta$.

\item Use (E.2.15.1) to pull back a meromorphic section of
$\mathcal O(D)$ whose divisor is $D[p]$.  We obtain a
meromorphic pluricanonical form $\sigma$ with
\[
 \operatorname{div}(\sigma)=D f^*(p).
\]
On $X_U$ it is holomorphic and nowhere zero.  Its quotient
by $(\eta\wedge dt)^{\otimes A}$ is a holomorphic function
constant on each compact connected fibre.  It therefore
descends to an entire nowhere-zero function $\phi$ on $\C$:
\[
 \sigma=\phi(t)(\eta\wedge dt)^{\otimes A}.
 \tag{E.2.15.3}
\]

\item On a small disc about $p$, the meromorphic section
$q^{-D}\sigma$ is in fact a holomorphic nowhere-zero section
of $K_X^{\otimes A}$ over the whole inverse image.  This
follows directly from (E.2.15.1), or by cancellation of the
actual divisors, with their multiplicities.  Denote it by
$\Omega_A$.  On the punctured inverse image, since
$dt=-q^{-2}dq$, it is
\[
 \Omega_A=(-1)^A P(q)(\eta\wedge dq)^{\otimes A},
 \qquad P(q)=q^{-D-2A}\phi(1/q).
 \tag{E.2.15.4}
\]
The function $P$ is single-valued and holomorphic on the
punctured disc.  No extension of a comparison with another
threefold is asserted here.

\item A nowhere-zero holomorphic $A$-canonical form determines
a smooth positive density: if it is
$h(z)(dz_1\wedge dz_2\wedge dz_3)^{\otimes A}$, take
$|h|^{2/A}$ times the ordinary complex-coordinate volume
density.  This is independent of the coordinates.  The
density associated with $\Omega_A$ has finite integral
over the inverse image of a smaller closed disc, because
that inverse image is compact.  Fibre integration of
(E.2.15.4), using (E.2.15.2), consequently gives
\[
 \int_{0<|q|<\epsilon}|P(q)|^{2/A}\,d\mathrm{Area}(q)
 <\infty.
 \tag{E.2.15.5}
\]
Only a fixed positive normalization constant is suppressed.
This step concerns a smooth positive density on the actual
total space and requires no K\"ahler metric.

\item We spell out the isolated-singularity consequence of
(E.2.15.5), including when $2/A<1$.  The function
$|P|^{2/A}$ is subharmonic.  If $|q|=r$, the mean-value
inequality on the disc of radius $r/2$ centred at $q$ gives
\[
 |P(q)|^{2/A}
 \leq \frac{4}{\pi r^2}
       \int_{r/2<|z|<3r/2}|P(z)|^{2/A}\,d\mathrm{Area}(z).
\]
The last integral tends to zero as $r\to0$, by
integrability.  Hence $r^A|P(q)|\to0$, uniformly on
$|q|=r$.  The function $q^A P(q)$ therefore has a removable
singularity and vanishes at zero.  For some holomorphic
$H$ on the disc,
\[
 P(q)=q^{1-A}H(q).
 \tag{E.2.15.6}
\]
Equivalently, the pole order of $P$ is at most $A-1$.

\item Combining (E.2.15.4) and (E.2.15.6) yields
\[
 \phi(1/q)=q^{D+A+1}H(q).
 \tag{E.2.15.7}
\]
If $D\geq-A$, the right-hand side extends to zero at
$q=0$.  Thus $\phi$ extends holomorphically to $\PP^1$
and vanishes at $p$.  A holomorphic function on $\PP^1$
is constant, so $\phi=0$, contradicting its construction.
This proves $D<-A$.

\item In \ref{push:src:E.2.10} the actual canonical-root equality has
\[
 A=12kl,\qquad D=kl-3k-4l,
 \qquad
 \frac DA=\frac1{12}-\frac1{4l}-\frac1{3k}
       \geq-\frac12>-1.
\]
If the collapsed constant source point is ordinary,
\ref{push:src:E.2.12} says that the central map has precisely one
critical value and is isotrivial on its complement.
All hypotheses above apply to the actual central
threefold, and give the required contradiction.

\end{enumerate}
\end{proof}



\paragraph{Scope of the volume argument}

The proof uses the equality of line bundles in
(E.2.15.1), not just the topological vanishing of
$c_1(K_X)$.  It also uses constancy of the torus periods,
which makes the fibre volume in (E.2.15.2) constant.
It does not assert a general good-reduction or
comparison theorem for non-K\"ahler torus degenerations.
Together with the separate elliptic-action exclusion
in \ref{push:src:E.2.14}, it removes the two types of collapsed
target permitted by \ref{push:src:E.2.12}.  Identification of a
noncollapsed interior completion, and the behaviour
on arbitrary new deformation branches, remain
distinct questions.

\Uses{the hypotheses, notation and definitions stated here; \cref{push:src:E.2.14}; \cref{push:src:E.2.12}}


\begin{proposition}[the central cusp is reduced and irreducible]\hypertarget{tag-0331}{}\label{push:src:E.2.16}
\StatementRecord{push:src:E.2.16}{proved}{0331}{\SeriesRef{push:src:E.2.10}, \SeriesRef{push:src:E.2.17}, \SeriesRef{push:src:E.2.20}, \SeriesRef{push:src:E.2.12}, \SeriesRef{push:src:E.2.14}, \SeriesRef{push:src:E.2.15}, \SeriesRef{push:src:E.2.13}, \SeriesRef{push:src:E.2.7}; \cite[Corollary 3.7, arXiv version~1]{HondaViaclovsky2025D}}

In the one-sided smooth specialization of \ref{push:src:E.2.10},
the three labelled critical values remain distinct.
The two finite fibres have their original exact
multiplicities $m=3k$ and $n=4l$.  Every reduced
central critical fibre is irreducible, and the
cusp fibre has multiplicity one.  Its underlying
surface $S$ satisfies
\[
 \chi_{\rm top}(S)=2,\qquad b_1(S;\mathbb Q)=2.
\]
This supplies the reduced irreducible Cartier central
surface used in the analytic cusp identification of
\ref{push:src:E.2.17}--\ref{push:src:E.2.20}.
\end{proposition}


\begin{proof}

\begin{enumerate}[label={$\langle1\rangle$\arabic*.},leftmargin=*]
\item \ref{push:src:E.2.12}, \ref{push:src:E.2.14} and \ref{push:src:E.2.15} exclude every degenerate
limit of the normalizing projective transformations.
They therefore extend as invertible projective
transformations, and the three labelled critical
values remain distinct.  \ref{push:src:E.2.13} proves the stated
finite-fibre multiplicities.

\item Normalize the three sections to $0,1,\infty$.
Over $B^\circ=\mathbb P^1\setminus\{0,1,\infty\}$,
the relative map is a proper submersion over
$B^\circ\times\Delta$.  Hence its integral
first-homology monodromies on the central fibre
are the same matrices $A_1,A_2,A_0$ as in
\ref{push:src:E.2.7}.  The common invariant space of the dual
matrices $A_1^t,A_2^t$ has dimension one:
in the displayed marking it is generated by
the covector $\gamma$.  This is obtained by
solving their simultaneous fixed equations;
the exact calculation is recorded in
\path{code/e_push_interior_cusp.py}.

\item Let $r$ be the total number of irreducible
components of the three reduced critical
fibres, and let $X^\circ$ be their complement
in $X_0$.  Alexander duality on $S^6$ gives
\[
 b_1(X^\circ;\mathbb Q)=r.
\]
The degree-one Leray sequence over the
punctured projective line, whose second
cohomology vanishes, gives instead
\[
 b_1(X^\circ;\mathbb Q)
 =b_1(B^\circ;\mathbb Q)
   +\dim H^1(T;\mathbb Q)^{\rm mon}
 =2+1=3.
\]
There are three nonempty critical fibres,
so each has exactly one reduced component.

\item Write the cusp fibre as $dS$ with $d\ge1$.
We prove $d=1$ by a direct meridian argument.
The smooth fibration over $B^\circ$ has
connected fibres, hence
$\pi_1(X^\circ)\to\pi_1(B^\circ)$ is
surjective.  The inclusion
$X^\circ\hookrightarrow X_0$ is surjective
on fundamental groups, and its kernel is
normally generated by meridians about the
three reduced divisor components.  To see
the latter assertion, make a null-homotopy
disc transverse to their smooth loci and
disjoint from their singular loci, which
have real codimension at least four.
Removing small discs about the finitely
many intersections expresses its boundary
as a product of conjugate meridians.

\item A meridian of a multiplicity-$a$ component
maps to the $a$th power of the corresponding
base meridian, since locally the base
function is a unit times the $a$th power
of a reduced local equation.  Consequently
there is a surjection
\[
 \pi_1(X_0)\longrightarrow
 \left\langle a,b,c\ \middle|\
 a^m=b^n=c^d=abc=1\right\rangle.
 \tag{E.2.16.1}
\]
This is the usual orbifold fundamental
group, but the surjection has just been
proved without assuming a stack
factorization.  Its source is trivial.

\item For clarity the group on the right is
nontrivial for every $d\ge2$ here.
Except for $(m,n,d)=(3,4,2)$, one has
\[
 \frac1m+\frac1n+\frac1d<1.
\]
Indeed for $d\ge3$ the sum is at most
$11/12$; for $d=2$, increasing either
$m=3k$ or $n=4l$ above $3,4$ makes it
at most $23/24$.  Put
$\alpha=\pi/m$, $\beta=\pi/n$,
$\theta=\pi/d$, and choose $\ell>0$ with
\[
 \cosh\ell=
 \frac{\cos\theta+\cos\alpha\cos\beta}
      {\sin\alpha\sin\beta}>1.
\]
Let
\[
 R_\alpha=\begin{pmatrix}
 \cos\alpha&\sin\alpha\\
 -\sin\alpha&\cos\alpha
 \end{pmatrix},\qquad
 D_\ell=\operatorname{diag}(e^{\ell/2},e^{-\ell/2}).
\]
In $\operatorname{PSL}_2(\mathbb R)$ the
elements represented by $R_\alpha$ and
$D_\ell R_\beta D_\ell^{-1}$ have orders
$m,n$, and their product has trace
$-2\cos\theta$ and projective order $d$.
Taking its inverse for $c$ gives a
nontrivial representation of (E.2.16.1).
For the exceptional triple, the
permutations $a=(123)$ and $b=(1324)$
have orders three and four, and
$ab=(24)$ has order two, giving a
nontrivial representation in $S_4$.
Thus $d\ge2$ contradicts (E.2.16.1),
and the cusp is reduced.

\item The ordinary torus fibres have Euler
characteristic zero.  The finite reduced
fibres are smooth finite free quotients
of two-tori by \ref{push:src:E.2.13}, so also have
Euler characteristic zero.  Additivity
over the three singular values therefore
gives $\chi_{\rm top}(S)=\chi(X_0)=2$.
Finally, the valid irreducible-fibre
degree-one invariant-cycle statement
\cite[Corollary 3.7, arXiv version~1]{HondaViaclovsky2025D}
gives
\[
 b_1(S;\mathbb Q)=
 \dim H^1(T;\mathbb Q)^{A_0}=2.
\]
Its irreducibility hypothesis has just
been proved.  This use is separate from
the false unrestricted equality
converse audited in \SeriesRef{push:src:B.4.17}.

\end{enumerate}
\end{proof}



\begin{proposition}[persistence of the normal-crossing cusp]\hypertarget{tag-0332}{}\label{push:src:E.2.17}
\StatementRecord{push:src:E.2.17}{proved}{0332}{\SeriesRef{push:src:E.2.10}, \SeriesRef{push:src:E.2.14}, \SeriesRef{push:src:E.2.16}; \citet[Theorem 2 and its corollary, pp.~274--275]{Lazzeri1973Root}}

In the smooth fixed-pair specialization of \ref{push:src:E.2.10},
the cusp divisor is a relative normal-crossing divisor
over the parameter disc.  More precisely, after
shrinking that disc, at every point of the total cusp
divisor there are holomorphic coordinates
$(z_1,z_2,z_3,s)$, with $s$ the original parameter,
in which the cusp coordinate of the intrinsic map is
one of
\[
 q=z_1,\qquad q=z_1z_2,\qquad q=z_1z_2z_3.
 \tag{E.2.17.1}
\]
There are exactly two relative triple-point sections.
Both lie on the fixed rational curve of \ref{push:src:E.2.14}.
In particular the normalization of the total cusp
divisor is smooth and proper over the parameter
disc, and its central fibre is the normalization
of the central cusp.
\end{proposition}


\begin{proof}

\begin{enumerate}[label={$\langle1\rangle$\arabic*.},leftmargin=*]
\item By \ref{push:src:E.2.14}--\ref{push:src:E.2.16} the three labelled critical
values remain distinct, the cusp has multiplicity
one, and its reduced central fibre is irreducible.
We normalize its target section to $q=0$.
On a sufficiently small target disc, zero is the
only critical value on every parameter fibre.
The extended field $\delta$ is zero precisely on
the relative fixed rational curve $\mathcal C$,
and the two normal characters there are $+1,-1$.
For $s\ne0$ the two triple points of the prescribed
cusp belong to $\mathcal C_s$; outside that curve
the cusp is smooth or has an ordinary double crossing.

\item We first record exactly the local singularity result
we use.  If $h_s(a,b)$ is a holomorphic family and
$h_0=0$ is a reduced plane-curve germ with an
isolated critical point, choose a sufficiently
small representative of that germ.  The sum of
the multiplicities of all nearby critical points
of $h_s$, regardless of their critical values,
is $\mu(h_0)$.  One can see this directly from
the analytic critical scheme
\[
 \mathbb C\{a,b,s\}/
       ((h_s)_a,(h_s)_b).
\]
After taking a sufficiently small representative,
it is finite over the parameter disc.  Its two
equations are a regular sequence, and its
central fibre is zero-dimensional.  It is thus
Cohen--Macaulay of dimension one, with $s$ a
nonzerodivisor, and is finite flat over the
disc.  Its degree is the central Jacobian-algebra
dimension, namely $\mu(h_0)$.

\item The non-splitting theorem of
\citet[Theorem 2 and its corollary, pp.~274--275]{Lazzeri1973Root}
says that if a deformation of an isolated
hypersurface singularity has more than one
singular point in a single deformed fibre,
their summed Milnor numbers are strictly
smaller than the original Milnor number.
The strict numerical inequality is the
conclusion of the cited corollary's proof.
Consequently, if all the critical points
just counted have critical value zero,
there can be only one such point.  This
application uses a flat deformation of
plane hypersurfaces: the map
$(a,b,s)\mapsto(h_s(a,b),s)$ is flat near
the central germ because its two base
parameters form a regular sequence in
the smooth source.  No assertion of
separatedness of compact threefolds enters
this local argument.

\item Let $p$ be a central cusp point outside
$\mathcal C_0$.  A relative holomorphic
flow box for the nonzero field $\delta$
gives coordinates $(a,b,v,s)$ in which
$\delta=\partial_v$.  Invariance of the
intrinsic map gives
\[
 q=h_s(a,b).
\]
The central equation is reduced, by
\ref{push:src:E.2.16}.  A reduced plane-curve germ is
either smooth or has an isolated critical
point.  In the smooth case the implicit
function theorem already gives the first
form in (E.2.17.1).

\item In the singular case every critical point
of $h_s$ in a small representative lifts
to a critical point of the actual map
on $X_s$.  All its critical values are
therefore zero.  For $s\ne0$ each such
point is a node, since the flow box is
disjoint from the fixed curve and hence
from the two triple points.  Conservation
and non-splitting give one node and
$\mu(h_0)=1$.  Thus the central Hessian
is nondegenerate.  The implicit function
theorem gives its unique critical-point
section, whose critical value is zero.
The holomorphic Morse lemma with
parameters then gives $q=ab$ throughout
this neighbourhood, with $s$ unchanged.

\item It remains to treat the fixed curve.
Holomorphic coordinates obtained by
averaging over the extended compact
circle may be chosen in the form
$(z,x,y,s)$, with $\mathcal C$ defined
by $x=y=0$ and with circle weights
$0,+1,-1,0$.  The original parameter
$s$ is retained by this averaging.
The holomorphic invariant function $q$
vanishes on $\mathcal C$, so its
convergent power series has the form
\[
 q=F_s(T,z)=T g_s(T,z),\qquad T=xy.
 \tag{E.2.17.2}
\]
The coefficient
$Q_s(z)=g_s(0,z)$ is intrinsically a
section of the dual of the product of
the two normal weight lines along
$\mathcal C_s$.  By \ref{push:src:E.2.14} this is
$\mathcal O_{\mathbb P^1}(2)$.

\item The central section $Q_0$ is not
identically zero.  Otherwise the
central equation in every such chart
would be divisible by $x^2y^2$,
contradicting reducedness of the
central cusp.  Its zero divisor
therefore has degree two.  Suppose
that both zeros coincide at a point
$z=0$.  The plane equation
$F_0(T,z)=Tg_0(T,z)$ is reduced:
a repeated factor of $g_0$ would
remain repeated after $T=xy$, and
$T$ does not divide $g_0$ because
$Q_0$ is not identically zero.
Thus $F_0$ has an isolated plane
critical point there.

\item For $s\ne0$ the two zeros of $Q_s$
are the two distinct triple points
of the prescribed cusp.  Under the
supposed collision they both lie
in the chosen small plane
representative.  They are distinct
critical points of $F_s$ with value
zero.  Every other critical point
of $F_s$ also has value zero:
this is automatic at $T=0$, and
at $T\ne0$ it follows by lifting
to $xy=T$, where the quotient
map $(x,y,z)\mapsto(T,z)$ is a
submersion.  All sufficiently small
$T$ have such a lift in the original
coordinate chart.  A nonzero
critical value would contradict
the actual critical-value set of
the intrinsic map.

\item Conservation of all critical
multiplicities and Lazzeri's
non-splitting theorem now give a
contradiction to the two distinct
critical points.  Hence $Q_0$ has
two simple zeros.  They extend
as two disjoint holomorphic
sections over the parameter disc.
Where $Q_0\ne0$, the factor
$g_s(xy,z)$ is a unit and can
be absorbed into $y$, giving
$q=xy$.  At a simple zero,
$\partial g_0/\partial z\ne0$;
replace $z$ by $g_s(xy,z)$ to
obtain $q=xyz$, again keeping
$s$ unchanged.  This proves
all local forms in (E.2.17.1).

\item Properness makes the central cusp
compact, so finitely many such
neighbourhoods cover it and a
single shrinking of the parameter
disc suffices.  The normalization
of a relative normal-crossing
divisor separates its local
branches.  In the coordinates
above every normalized branch
is smooth over $s$, and these
descriptions agree on overlaps.
Normalization is finite for
complex analytic spaces; its
composition with the proper
cusp map is therefore proper.
Its fibres are precisely the
normalizations of the cusp
fibres in these local charts.
This proves the final assertions.

\end{enumerate}
\end{proof}



\paragraph{Literature and remaining identification}

Non-splitting is a known theorem of
Lazzeri, not a new singularity-theoretic
claim.  We opened the original paper
and use its strict numerical
corollary.  The application combines
that theorem with the actual extended
circle action, reducedness, and the
absence of extra critical values
proved above.  Searches for
non-splitting, Milnor-number
conservation and degenerations of
torus-fibred complex six-spheres
did not locate this application.  The
proposition supplies relative
normal-crossing charts and a
simultaneous normalization.  It
does not by itself identify the entire
analytic cusp neighbourhood.  The
prescribed marked Mumford completion
is recovered by \ref{push:src:E.2.18}--\ref{push:src:E.2.20}.

\Uses{the hypotheses, notation and definitions stated here; \cref{push:src:E.2.18}; \cref{push:src:E.2.20}}


\begin{proposition}[the normalized cusp and its marked hexagon]\hypertarget{tag-0333}{}\label{push:src:E.2.18}
\StatementRecord{push:src:E.2.18}{proved}{0333}{\SeriesRef{push:src:E.2.17}; \citet[\S4, no.~8(a), pp.~364--365]{Grauert1962E}; \cite[Chapter~VI, Theorem~1.1 and Table~10]{BHPV2004E}}

Let $\nu:\mathcal S\to\mathcal W$ be the
simultaneous normalization of the relative cusp
in \ref{push:src:E.2.17}, and let $\mathcal D\subset\mathcal S$
be its reduced conductor.  After shrinking the
parameter disc, each fibre pair is the standard
toric pair
\[
 (\mathcal S_s,\mathcal D_s)
 \simeq(\mathrm{dP}_6,\hbox{its toric hexagon}).
\]
The six conductor components may be labelled over
the disc so that their incidence, self-intersections
$-1$, opposite-pair involution, and endpoint
identifications are those of the punctured family.
The relative logarithmic tangent bundle has a
global holomorphic frame:
\[
 T_{\mathcal S/\Delta}(-\log\mathcal D)
 \simeq\mathcal O_{\mathcal S}^{\oplus2}.
\]
This assertion identifies the normalized pair
and supplies a relative frame; it does not by
itself identify the analytic smoothing neighbourhood.
\end{proposition}


\begin{proof}

\begin{enumerate}[label={$\langle1\rangle$\arabic*.},leftmargin=*]
\item In the relative charts of \ref{push:src:E.2.17}, the normalized
branches are smooth over the original parameter.
Their conductor branches and their intersections
are also smooth over that parameter.  The
normalization of the conductor is consequently
a proper smooth family of curves, and the
normalized intersection strata are finite
etale over the disc.  Their connected components
and incidence are constant.  Since the disc
is simply connected, label the six curve
families $\mathcal D_1,\ldots,\mathcal D_6$
and the six vertex sections by continuation
from a punctured fibre.  Each curve fibre
is a smooth curve of genus zero.  Its two
distinct vertex sections remain distinct;
the same relative local charts exclude a
new self-intersection or an additional
incidence at the centre.  Thus the central
conductor is a hexagon of six embedded
rational curves.  The degree of each
relative normal line is constant in its
proper curve family, so all six central
self-intersections equal $-1$.

\item Write $S=\mathcal S_0$ and $D_i=\mathcal D_{i,0}$.
Ehresmann transport gives $b_1(S)=0$ and
$b_2(S)=4$.  Contract the three disjoint
curves $D_1,D_3,D_5$.  The analytic
Castelnuovo contraction gives a smooth
compact complex surface $Y$; the needed
analytic statement, including smoothness
of the resulting point, is
\citet[\S4, no.~8(a), pp.~364--365]{Grauert1962E}.
Only these three contractions on the
central surface are used; no relative
contraction theorem is assumed.

\item The images $C_2,C_4,C_6$ of the remaining
curves are smooth rational curves.  Each
meets two contracted curves transversely
once, so has self-intersection
$-1+1+1=1$.  Each pair meets once, at
the image of their intervening contracted
curve; these three intersection points
are distinct and the intersections are
transverse.  Thus they form a triangle.
Also $b_1(Y)=0$ and $b_2(Y)=1$.

\item Any one of these curves, say $C=C_2$,
has nonnegative intersection with every
effective divisor: distinct irreducible
curves have nonnegative local intersection
multiplicities and $C^2=1$.  Adjunction
gives $K_Y\cdot C=-3$.  A nonzero section
of $K_Y^{\otimes r}$ for $r>0$ would
give an effective divisor of intersection
$-3r$ with $C$, impossible.  Hence
$\kappa(Y)=-\infty$.  The compact-surface
classification
\cite[Chapter~VI, Theorem~1.1 and Table~10]{BHPV2004E}
together with $b_1(Y)=0$ makes $Y$ rational:
class~VII has first Betti number one,
and a ruled surface has the first Betti
number of its base curve.  A rational
surface with $b_2=1$ is $\mathbb P^2$.
The curves $C_2,C_4,C_6$ are lines,
since their self-intersections are one.
Their three pairwise intersection points
are distinct, so they form a projective
triangle.  Blowing up its three vertices
recovers exactly the toric degree-six
del Pezzo surface and its hexagonal
boundary.  This proves the assertion
about the central pair; the punctured
pairs were already the prescribed ones.

\item The conductor involution is holomorphic
over the disc.  Its permutation of the
six labelled curve families, and its
permutation of the vertex sections,
are locally constant.  They are therefore
the same opposite-pair and endpoint
identifications as on a punctured fibre.
The remaining scalar gluing parameters
are not claimed to be determined by
this discrete observation.

\item Finally put
$\mathcal E=T_{\mathcal S/\Delta}(-\log\mathcal D)$.
For every fibre pair just identified,
the two toric infinitesimal actions
trivialize its logarithmic tangent
bundle.  Thus
\[
 \mathcal E|_{\mathcal S_s}\simeq
 \mathcal O_{\mathcal S_s}^{\oplus2},
 \qquad h^0(\mathcal S_s,\mathcal E_s)=2.
\]
Constant-cohomology base change, in the
form used in \SeriesRef{push:src:E.3.1}, makes the direct
image a rank-two holomorphic vector
bundle on the disc and identifies its
fibres with these section spaces.
Choose a frame after shrinking the
disc.  The evaluation map from its
pullback to $\mathcal E$ is an
isomorphism on every fibre, hence
is an isomorphism everywhere by
Nakayama's lemma.  This gives the
claimed relative frame without
choosing a potentially singular
family of projective coordinate
changes on the punctured disc.

\end{enumerate}
\end{proof}



\begin{proposition}[extension of the marked cusp torus action]\hypertarget{tag-0334}{}\label{push:src:E.2.19}
\StatementRecord{push:src:E.2.19}{proved}{0334}{\SeriesRef{push:src:E.2.17}, \SeriesRef{push:src:E.2.18}, \SeriesRef{push:src:E.2.16}}

Let
\[
 f:\mathcal N\longrightarrow
 B=\Delta_q\times\Delta_s
\]
be the proper cusp neighbourhood of the smooth
specialization in \ref{push:src:E.2.17}, and put
$\mathcal W=f^{-1}(q=0)$.  After shrinking the
bidisc, the rank-two vertical logarithmic tangent
bundle is trivial and is generated by two commuting
complete holomorphic vector fields.  Their action
factors through a holomorphic action of
$(\C^*)^2$ over $B$, with the integral marking
given by the periods $\mathbf w,\boldsymbol\delta$.
On the normalization of every $\mathcal W_s$
this is its usual toric action.

The marked rank-two lattice is primitive in the
period lattice of each ordinary torus fibre and
equals the full invariant lattice of the cusp
monodromy.  At the two relative triple-point
sections, the three nonzero tangent characters
are constant in $s$ and equal to those of the
prescribed Mumford cusp.  In the toric chart
lattice coordinates, transported from that
punctured model by their fixed integral change
of basis from the period marking, they are
respectively
\[
 \{(1,0),(0,1),(-1,-1)\},\qquad
 \{(-1,0),(0,-1),(1,1)\}.
 \tag{E.2.19.1}
\]
The assertion constructs the action on the actual
neighbourhood.  It does not assert that an
arbitrary equivariant quotient is smooth or
that the entire completion has already been
identified with a prescribed toric quotient.
\end{proposition}


\begin{proof}
\begin{enumerate}[label={$\langle1\rangle$\arabic*.},leftmargin=*]
\item 
\noindent\textbf{Step 1. The logarithmic bundle and its restriction.}
The coordinates in \ref{push:src:E.2.17} make $\mathcal W$ a
relative normal-crossing divisor over $s$.
Define
\[
 \mathcal L=
 \ker\left[
 d\log q:
 T_{\mathcal N/\Delta_s}(-\log\mathcal W)
       \longrightarrow\mathcal O_{\mathcal N}
 \right].
 \tag{E.2.19.2}
\]
The map is surjective: if
$q=z_1\cdots z_k$, it sends each
$z_i\partial_{z_i}$, $i\leq k$, to one.
Consequently $\mathcal L$ is locally free of
rank two.  Its sections are actual holomorphic
vector fields, tangent to the branches, and
annihilate both $q$ and $s$.

Let $\nu:\mathcal S\to\mathcal W$ and
$\mathcal D\subset\mathcal S$ be the
simultaneous normalization and conductor.
Restriction to a branch gives a canonical
isomorphism
\[
 \nu^*(\mathcal L|_{\mathcal W})
 \simeq T_{\mathcal S/\Delta_s}(-\log\mathcal D).
 \tag{E.2.19.3}
\]
This can be checked at every stratum.  For
$q=xy$, a basis is
$x\partial_x-y\partial_y,\partial_z$;
on the branch $x=0$ it becomes
$-y\partial_y,\partial_z$, a logarithmic
basis there.  For $q=xyz$, use
$x\partial_x-z\partial_z$ and
$y\partial_y-z\partial_z$.  Their
restrictions to $x=0$ have coefficient
matrix of determinant one in the basis
$y\partial_y,z\partial_z$.
The one-branch case is immediate.
The same computation keeps $s$ fixed.

By \ref{push:src:E.2.18} the right-hand side of
(E.2.19.3) has a global holomorphic
frame.  On the two preimages of each
conductor curve, the identifications
in (E.2.19.3) describe descent by an
invertible two-by-two matrix.  Its
entries are holomorphic on a proper
$\mathbb P^1$ family over $\Delta_s$,
so are constant along each curve
fibre and are holomorphic functions
of $s$.  Denote these matrices by
$G_j(s)$.

For $s\ne0$ both invariant logarithmic
fields in the prescribed Mumford model
descend.  Indeed, its period maps in
multiplicative coordinates multiply
the two coordinates by units depending
only on the base.  Their derivative
on the two vertical logarithmic
fields is the identity.  The toric
small-resolution charts preserve
these same fields.  Thus the entire
two-dimensional space of sections
on the normalized pair descends for
$s\ne0$.  In any chosen frame this
means $G_j(s)=I$, and the identity
theorem gives $G_j(0)=I$ as well.

There is no additional descent
condition hidden at a triple point.
For a normal-crossing local ring,
functions on the normalization
descend precisely when their
restrictions to every pair of
branches agree.  This is the usual
equalizer for
$\C\{x,y,z\}/(xyz)$, and likewise
for one or two branches.  Tensoring
with the locally free bundle
$\mathcal L|_{\mathcal W}$ proves
the same assertion for its sections.
Consequently the chosen normalized
frame descends, and
\[
 \mathcal L|_{\mathcal W}
       \simeq\mathcal O_{\mathcal W}^{\oplus2}.
 \tag{E.2.19.4}
\]

\noindent\textbf{Step 2. A frame on the whole neighbourhood.}
The map $f$ is proper and flat; this
also follows locally from the equations
$(q,s)=(z_1\cdots z_k,s)$.
The bundle $\mathcal L$ is flat over
$B$.  On each ordinary fibre it is
the tangent bundle of a complex
two-torus.  On each $q=0$ fibre it
is trivial by (E.2.19.4).  All these
fibres are compact and connected,
so in every case
\[
 h^0(\mathcal N_{q,s},
       \mathcal L|_{\mathcal N_{q,s}})=2,
\]
and those sections generate the
bundle pointwise.  The proper
constant-cohomology base-change
theorem, in the form established
in \SeriesRef{push:src:E.3.1}, makes $f_*\mathcal L$
a rank-two vector bundle on $B$.
Its evaluation map
\[
 f^*f_*\mathcal L\longrightarrow\mathcal L
 \tag{E.2.19.5}
\]
is an isomorphism on every fibre,
and hence everywhere by Nakayama's
lemma.  Shrink the bidisc and
choose a holomorphic frame $V_1,V_2$.

These fields commute.  On every
ordinary compact torus, every
holomorphic vector field is a
translation field, so their
bracket vanishes there.  The
ordinary locus is dense, and
the bracket is holomorphic.
Each field is complete because
it is vertical and every fibre
of the proper map is compact:
the holomorphic local flow can
be continued for every complex
time on that compact invariant
fibre.  Holomorphic dependence
on initial points and parameters
gives a holomorphic action
\[
 \operatorname{Act}:\C^2\times
       \mathcal N\longrightarrow\mathcal N
\]
over $B$.

\noindent\textbf{Step 3. Holomorphic extension of the two toric periods.}
Choose a smooth point $p_0$ of the
central cusp away from its conductor.
The map $f$ is a submersion there,
so there is a holomorphic section
$p(q,s)$ through $p_0$.
On the central normalized toric
pair the two fields are a basis
of the toric logarithmic Lie
algebra.  Its open orbit is
$(\C^*)^2$.  Let
$v_1^0,v_2^0\in\C^2$ be a
basis of its full period lattice
in the coordinates $V_1,V_2$.
They are complex-linearly
independent.

Near each $v_i^0$, write
$\operatorname{Act}(v,p(q,s))=p(q,s)$
using two vertical coordinates
in a submersion chart around
$p_0$.  The derivative in $v$
is invertible at $(v_i^0,0,0)$,
since the infinitesimal action
is an isomorphism on that open
orbit.  The holomorphic implicit
function theorem gives a unique
holomorphic vector $v_i(q,s)$
near $v_i^0$ solving this equation.
For $q\ne0$, fixing $p(q,s)$
means that this translation is
the identity on the entire
torus fibre.  The equality of
holomorphic maps then extends
over $q=0$.  Thus $v_i(q,s)$
is a period on every fibre.
Their complex-linear
independence persists after
shrinking $B$.

Use these two columns as a
change of the frame.  Because
their coefficients are pulled
back from $B$, the new fields
still commute and are complete.
They have period one, and the
action consequently factors
through
$\C^2/\Z^2\simeq(\C^*)^2$,
using the exponential
$z\mapsto e^{2\pi i z}$
in each coordinate.

\noindent\textbf{Step 4. Primitivity and the actual marking.}
On $B^\circ=\Delta_q^*\times\Delta_s$
the period lattices of the
translation action are the
integral first-homology local
system of the torus family.
The continuously varying
periods $v_1,v_2$ represent
flat integral classes: in a
local differentiable
trivialization their flow
loops vary continuously, and
integral homology is discrete.
They generate a rank-two
invariant sublattice $L$.

Suppose $L$ were not primitive
at some ordinary fibre.  There
would be a rational pair
$(a_1,a_2)\notin\Z^2$ such
that $a_1v_1+a_2v_2$ is an
integral period there.  This
is an equality of flat
integral classes after
multiplication by a common
denominator, so it persists
throughout the connected
ordinary base.  In particular
the corresponding fixed
finite-order element of
$(\C^*)^2$ acts identically
there.  By the identity
theorem it acts identically
on the central open orbit.
This contradicts the choice
of $v_1^0,v_2^0$ as the full
central toric period lattice.
Thus $L$ is primitive.

The marked cusp monodromy,
constant in $s$ by \ref{push:src:E.2.16},
has the form
\[
 \widehat\gamma\longmapsto
       \widehat\gamma-\widehat\delta,
 \qquad
 \widehat u\longmapsto\widehat u+\widehat w,
 \qquad
 \widehat w\longmapsto\widehat w,
 \qquad
 \widehat\delta\longmapsto\widehat\delta.
 \tag{E.2.19.6}
\]
Its full integral invariant
lattice is
$\Z\widehat w\oplus\Z\widehat\delta$.
Since $L$ is a primitive
rank-two sublattice of this
rank-two lattice, equality
follows.  An integral
unimodular change of the two
periods therefore gives
exactly the marking
$\mathbf w,\boldsymbol\delta$.
This also proves that no
extra finite kernel has been
silently introduced.

\noindent\textbf{Step 5. Characters at the triple points.}
The two triple-point sections
are fixed by the torus action:
on the normalization they
are toric vertices.  Restrict
the derivative representation
to the compact real torus
$(S^1)^2$ and average a
Hermitian metric along either
section.  Its characters
are integral pairs and
vary continuously in $s$;
hence their multiset is
locally constant.  Equivalently,
Fourier averaging gives the
individual character
subbundles.  They coincide
at $s=0$ with the known
characters for $s\ne0$.

For clarity, the two unit
triangles in the prescribed
cocharacter lattice have
height-one rays
\[
 \{(0,0,1),(1,0,1),(0,1,1)\},
 \qquad
 \{(1,0,1),(0,1,1),(1,1,1)\}.
\]
Inverting their unimodular
ray matrices gives the two
character triples in
(E.2.19.1); the sum in
each triple is zero, as
required by invariance of
$q=z_1z_2z_3$.
The calculation is recorded
in \path{code/d_push_cusp_torus_action.py}.
All action and marking
assertions follow.

\end{enumerate}
\end{proof}



\begin{proposition}[identification of the actual analytic cusp]\hypertarget{tag-0336}{}\label{push:src:E.2.20}
\StatementRecord{push:src:E.2.20}{proved}{0336}{\SeriesRef{push:src:E.2.10}, \SeriesRef{push:src:E.2.7}, \SeriesRef{push:src:E.2.16}, \SeriesRef{push:src:E.2.18}, \SeriesRef{push:src:E.2.19}, \SeriesRef{push:src:E.2.12}; \cite[p.~30]{AC26Root}}

In the smooth specialization of \ref{push:src:E.2.10}, normalize
the three distinct critical-value sections to
$0,1,\infty$ and write $q=1/t$ near the cusp.
Let $c(s)$ be the holomorphic interior period
lift supplied by \ref{push:src:E.2.7}.  After shrinking the
$(q,s)$-bidisc, its inverse image in the actual
family is biholomorphic, over that bidisc, to
the prescribed marked Mumford cusp family at
parameter $c(s)$.  The isomorphism induces
the given identification of the regular
torus Jacobians and their integral markings.
Its affine origin may be fixed by choosing
a holomorphic section through a smooth
point of the reduced cusp.
\end{proposition}


\begin{proof}
\begin{enumerate}[label={$\langle1\rangle$\arabic*.},leftmargin=*]
\item 
Write $f:\mathcal N\to\Delta_q\times\Delta_s$
for the proper actual cusp neighbourhood and
$\mathcal W=f^{-1}(q=0)$.  \ref{push:src:E.2.16}--\ref{push:src:E.2.18}
give its reduced relative normal-crossing
divisor, its normalized toric hexagon, and
the unchanged opposite-side identifications.
\ref{push:src:E.2.19} gives an actual holomorphic action
of $H=(\C^*)^2$ over the bidisc.  Its
period marking is the primitive vanishing
lattice
\[
 \Lambda_{\rm tor}
   =\Z\widehat w\oplus\Z\widehat\delta.
 \tag{E.2.20.1}
\]
We reconstruct the neighbourhood from
this action, keeping track of the
extension of its characters.

\noindent\textbf{Step 1. The topology and the exact covering.}
The relative coordinates
$q=x,xy,xyz$ give smooth real lifts
of the two real parameter directions
of $s$ which annihilate $dq$ and
are tangent to every cusp stratum.
Away from the cusp such lifts exist
because $f$ is a submersion.
A smooth partition of unity preserves
both these linear requirements.
On the inverse image of a smaller
closed $q$-disc the lifts integrate
along a sufficiently short parameter
path: $q$ remains constant, and
properness prevents escape.
Thus their flows identify the cusp
neighbourhoods, all strata, and
the regular-fibre inclusion maps
by $q$-preserving diffeomorphisms.

For the known Mumford cusp we
check the covering calculation
directly.  Each maximal toric
chart above $|q|<\epsilon$ is
\[
 \{(x,y,z)\in\C^3:|xyz|<\epsilon\},
\]
which is star-shaped and simply
connected.  Every intersection
contains the same connected dense
open torus $(\C^*)^2\times\Delta_q^*$.
Each connected component of an
intersection is open and must meet
that dense set, so the intersection
is connected.  The same holds for
the intersection with a finite
union of preceding charts.
Van Kampen proves simple
connectedness of each finite
union; a loop lies in such a
union by compactness.  The full
toric cover is therefore simply
connected.  The known free,
properly discontinuous period
action has deck group
\[
 \overline\Lambda=\Lambda/\Lambda_{\rm tor}
                 \simeq\Z^2.
\]
On a regular fibre its covering
is $\C^2/\Lambda_{\rm tor}\to
\C^2/\Lambda$.  The regular-fibre
map to the cusp body's fundamental
group is exactly the projection
to $\overline\Lambda$, and the
$q$-meridian dies through the
section $(x,y,z)=(q,1,1)$.
These statements also agree
with the original manuscript's
Proposition~4.7 and Corollary~4.8,
whose original PDF was checked
\cite[p.~30]{AC26Root}; the
preceding argument verifies
the step we use.  The
$q$-preserving transport transfers
this exact description to the
actual neighbourhood.

Take the universal covering
\[
 \widetilde{\mathcal N}\longrightarrow\mathcal N.
\]
Its deck group is
$\overline\Lambda$, and its
central normalized components
form the planar unfolding of
the hexagon.  The $H$ action
lifts.  Indeed the fundamental
loops of $H$ at a central open
orbit lie in one normalized
del Pezzo component, which is
simply connected, and hence die
in $\pi_1(\mathcal N)$.
The identity lift is unique
and satisfies the action law.
Choose a holomorphic section
$p(q,s)$ through a smooth
central cusp point and lift
it to the cover.  On $q\ne0$
its orbit map is an isomorphism
\[
 H\times(\Delta_q^*\times\Delta_s)
       \xrightarrow{\ \sim\ }
 \widetilde{\mathcal N}|_{q\ne0}.
 \tag{E.2.20.2}
\]
The stabilizer has been removed
exactly by (E.2.20.1), using
the just-verified regular-fibre
inclusion map.  In particular
the two characters of $H$
are globally defined on this
open covering space.

Connectedness of $H$ makes
its lifted action commute
with every deck transformation:
conjugation acts continuously
on the discrete deck group
and is the identity at the
identity of $H$.  On
(E.2.20.2), the two deck
generators therefore act by
\[
 (h,q,s)\longmapsto(P_i(q,s)h,q,s),
 \qquad i=1,2.
 \tag{E.2.20.3}
\]
The functions $P_i$ are the
actual marked period multipliers.
Changing the affine origin
does not change them.
The marked period comparison
of \ref{push:src:E.2.7} and \ref{push:src:E.2.12} identifies
them with the prescribed
$q$-monomials times holomorphic
units at $c(s)$.

\noindent\textbf{Step 2. Character extension across adjacent components.}
For a central component $D_v$,
choose a holomorphic base
section $a_v(q,s)$ through
its dense $H$ orbit.  Its
orbit map is an open
holomorphic embedding
\[
 H\times(\Delta_q\times\Delta_s)
       \longrightarrow\widetilde{\mathcal N}.
\]
It is injective on the
central open orbit and on
each regular fibre, and
its differential is invertible
along those loci.  Its
image contains the whole
regular covering.  Two
such origins differ for
$q\ne0$ by a holomorphic
translation $h_{vw}(q,s)$
in $H$.  We must prove
that this function has no
essential singularity at
$q=0$.

If both origins lie in
the same central component,
their ratio extends as a
holomorphic $H$-valued
unit.  The central open
orbit is free and transitive,
so this follows from the
holomorphic inverse-function
theorem for the orbit map,
with both parameters retained.

For two adjacent components,
use a crossing chart with
$q=ab$ and an orbit coordinate
$z\ne0$.  The stabilizer
of a general conductor
point is a primitive
connected one-dimensional
subtorus of $H$.  Its
two normal characters are
$+1,-1$: on each normalized
toric component these are
the primitive boundary
characters, and invariance
of $q$ makes them opposite.
Compact-circle averaging
gives these exact local
linear coordinates.  A
complementary torus direction
provides the orbit coordinate;
an invariant unit in $q$
can be absorbed while
preserving the weights.

Fix small nonzero constants
$a_0,b_0,z_0$ in this chart
and take the two special
origins
\[
 (a,b,z)=(q/b_0,b_0,z_0),\qquad
 (a,b,z)=(a_0,q/a_0,z_0).
 \tag{E.2.20.4}
\]
The stabilizer cocharacter
evaluated at $a_0b_0/q$
sends the first to the
second.  This equality
is justified by the local
action itself: a path
in that subgroup can
increase the modulus of
$a$ and decrease that
of $b$, with $ab=q$
fixed, staying inside
the crossing chart;
the phase can be varied
without changing either
modulus.  Thus the
possibly large group
parameter causes no
unjustified use of a
local coordinate outside
its domain.

Comparing (E.2.20.4) with
arbitrary chosen origins
adds only holomorphic
$H$-valued units, by the
same-component argument.
The adjacent comparison
therefore has the form
\[
 h_{vw}(q,s)=q^{n_{vw}}u_{vw}(q,s),
 \qquad n_{vw}\in\Z^2,
 \tag{E.2.20.5}
\]
where $u_{vw}$ is holomorphic
with values in $H$ on the
whole bidisc.  Every central
component is reached by
a finite adjacent path.
Multiplying these expressions
proves the same assertion
for its comparison with
the initial origin.
Different paths agree
on the regular cover,
and hence define identical
meromorphic characters.
The exponents are locally
constant in $s$, since
(E.2.20.5) holds with a
nonvanishing unit on the
full bidisc.

\noindent\textbf{Step 3. The fan and the analytic coordinate maps.}
At either triple-point
section the whole $H$
fixes the point.
Compact-torus averaging
gives three eigen-coordinates
with the constant integer
characters of \ref{push:src:E.2.19}.
Their only primitive
positive invariant relation
is their product.
The reduced normal-crossing
equation therefore writes
$q$ as that product times
a holomorphic invariant
unit.  The unit is a
function of the product
and $s$; absorb it into
one eigen-coordinate to
obtain exactly $q=xyz$.

The orders of the two
global characters along
$D_v$, together with
$\operatorname{ord}_{D_v}q=1$,
give a ray $(v,1)$.
At a triple point the
three incident rays
form a unimodular cone.
Indeed all adjacent
orders and tangent
characters are constant
in $s$ and coincide
with the known ones
for $s\ne0$.
The entire resulting
fan is consequently
the prescribed $A_2$
triangulation, with
its fixed labels.
No additional refinement
is being introduced.

In a triple chart a
global character is
a Laurent monomial in
$x,y,z$ times a function
of $q,s$.  Evaluate it
on a section keeping
two nonzero coordinates
fixed and solving the
third from $q$.
The origin comparisons
above show that the
remaining function is
a power of $q$ times
a holomorphic unit.
Thus the dual monomial
coordinates of the
unimodular fan cone
pull back to $x,y,z$
times holomorphic units
in $q,s$.  They extend
holomorphically and
give a local
biholomorphism through
the triple point.
The crossing and
open-orbit charts give
the same conclusion
on the other strata.
The maps agree on
the dense regular
cover and therefore
glue.

We finally check the
common neighbourhood
and the deck action.
We do not take an
infimum of infinitely
many chart radii.
Choose finitely many
charts covering one
compact normalized
central del Pezzo
component and finitely
many overlaps for
the two deck generators.
Shrink the bidisc once
for this finite data.
On the regular cover
the character map
intertwines the actual
deck action with the
reference multipliers
$P_i(q,s)$ in
(E.2.20.3).  Both maps
extend in the charts
just constructed, so
their equality extends
to the central strata.
The reference multipliers
preserve the prescribed
invariant fan.  Deck
translates of the finite
chart collection cover
the entire central
unfolding and use the
same base bidisc.

We have obtained local
biholomorphisms at all
central strata, bijections
on their labelled torus
orbits, and the original
isomorphism on the whole
regular cover.  The
inverse coordinate maps
are the same unimodular
monomial charts with
the inverse units.
Quotienting by the
identical free properly
discontinuous deck
actions gives the
claimed biholomorphism
of actual analytic
cusp neighbourhoods.

\end{enumerate}
\end{proof}



\paragraph{Scope and source comparison}

The toric covering and its
period action are the known
Mumford construction.  The
point here is to recover
that construction from
the actual smooth limit,
using the action proved
in \ref{push:src:E.2.19}.  In particular
the absence of an essential
singularity in the character
comparison is proved in
(E.2.20.4)--(E.2.20.5);
it is not an assumed
meromorphic comparison.
The local identification
does not alone identify
the global threefold or
settle arbitrary deformation
chains.  \ref{push:src:E.2.21} supplies the
finite affine and global collar
comparison; \ref{push:src:E.2.22}--\ref{push:src:E.2.24}
provide the additional propagation
and closure arguments for general
bases and continuous structures.

\Uses{the hypotheses, notation and definitions stated here; \cref{push:src:E.2.19}; \cref{push:src:E.2.21}; \cref{push:src:E.2.22}; \cref{push:src:E.2.24}}


\begin{proposition}[global identification of a smooth completion]\hypertarget{tag-0337}{}\label{push:src:E.2.21}
\StatementRecord{push:src:E.2.21}{proved}{0337}{\SeriesRef{push:src:E.2.5}, \SeriesRef{push:src:E.2.7}, \SeriesRef{push:src:E.2.10}, \SeriesRef{push:src:E.2.12}, \SeriesRef{push:src:E.2.14}, \SeriesRef{push:src:E.2.15}, \SeriesRef{push:src:E.2.13}, \SeriesRef{push:src:E.2.20}}

Let $\pi:\mathcal X\to\Delta$ satisfy the hypotheses of
\ref{push:src:E.2.5}, with fixed coprime multiplicities $(3k,4l)$ on
every noncentral fibre.  Then $X_0$ is a member of that
same constructed family.  Its isomorphism parameter is
the interior value $z(0)$ of \ref{push:src:E.2.7}.

After choosing a holomorphic lift $c(s)$ of that parameter
and shrinking the disc, the actual family is isomorphic
to the prescribed family at $c(s)$.  The ambiguity in
this comparison is the connected translation automorphism
group already determined in E.1.
\end{proposition}


\begin{proof}
\begin{enumerate}[label={$\langle1\rangle$\arabic*.},leftmargin=*]
\item 
\ref{push:src:E.2.10} extends the intrinsic pencil to a proper flat map
\[
 g:\mathcal X\longrightarrow\mathbb P^1\times\Delta.
\]
\ref{push:src:E.2.7} supplies its interior nonzero period $z(s)$ and hence
a holomorphic logarithmic lift $c(s)$.  \ref{push:src:E.2.12} excludes
additional critical values, and \ref{push:src:E.2.14}--\ref{push:src:E.2.15} exclude
target collapse.  Normalize the three distinct labelled
sections to $0,1,\infty$ by a holomorphic projective
transformation over the disc.

On the ordinary locus the actual marked Jacobian has
period matrix $\Pi_{c(s)}(t)$.  More intrinsically, its
holomorphic period kernel in the complexified marked
lattice equals that of this matrix.  This equality holds
for $s\ne0$ by the given family identification and extends
over $s=0$ by the identity theorem.  The local marking
choices agree: an automorphism of a fixed marked member
has identity action on ordinary $H_1$, by \SeriesRef{push:src:E.1.2}.  The
actual smooth family at any fixed ordinary target point
also has identity homology return around $s=0$.

\ref{push:src:E.2.13} gives smooth actual normalized root covers at the
two finite values, of degrees $3k$ and $4l$, with their
free cyclic deck actions.  Their marked Jacobians are
the corresponding reference root Jacobians.  \ref{push:src:E.2.20}
identifies the actual cusp neighbourhood over the full
$(q,s)$-bidisc with the prescribed marked Mumford
neighbourhood, including its period multipliers.
It remains to identify the finite affine actions and
the global translations on the overlaps.

\noindent\textbf{Step 1. The finite affine classes.}
Choose local holomorphic origins on a finite root torus
family.  Every deck map is affine on the compact torus
fibres.  Its linear part is the fixed marked finite
monodromy.  Subtracting this part gives a holomorphic
translation section.  Lift it to the Lie bundle on
simply connected refinements of the root charts.
The finite-order relation says that the sum of its
successive linear translates is an integral period.
This lattice vector is locally constant in $s$:
the lifted translation and period vectors vary
continuously, whereas the integral coefficients are
discrete.  Changing the lift changes it by precisely
the integral coboundary allowed in E.1.

Consequently the finite affine cocycle has the same
integral class at $s=0$ as at $s\ne0$.  After matching
these classes, the remaining difference is a cocycle
in the complex vector sheaf.  The explicit finite
averaging in E.1 removes it, since division by the
finite group order is allowed in that sheaf.
This constructs the finite-root comparisons with
holomorphic dependence on $s$.  Equivalently one
may use a finite chart cover and its deck translates;
the same averaging solves the equivariant comparison.
The reduced finite fibres are the resulting free
quotients; they have not been assumed to be tori.

\noindent\textbf{Step 2. The integral collar classes.}
Use the two finite-root neighbourhoods, the cusp
neighbourhood, ordinary simply connected charts,
and fixed annular collars to cover the target.
Choose local holomorphic origins on the ordinary
torus charts.  A transition between two such
descriptions is fibrewise affine.  The common
integral marking makes its linear part the
prescribed one, leaving a translation section.

Lift these sections to the Lie bundles on simply
connected refinements of the collars.  Differences
on deck translates and sums on triple overlaps
are integral periods.  Every resulting integral
vector is constant in $s$.  This comparison takes
place on fixed collars away from the singular
centres; it does not evaluate a logarithm at a
singular fibre.  The finite and cusp chart
identifications specify exactly which integral
changes of origins are allowed.

For $s\ne0$ these classes are those of the
prescribed logarithmic-transform data, modulo
the integral changes of origins in E.1.
They are therefore the same at $s=0$.
One can normalize them once on the connected
parameter disc.  This includes the cusp
winding class, since \ref{push:src:E.2.20} identifies the
actual marked period action there.  Thus
there is no additional central integral
translation or fan datum.

\noindent\textbf{Step 3. Solving the holomorphic remainder.}
After this discrete normalization, the remaining
comparison is the additive cocycle in the
extension sheaf $\mathcal V$ of E.1.  Its
finite extension conditions are exactly
those of the already identified root
charts.  Its cusp extension condition
is the actual toric collar condition
proved in \ref{push:src:E.2.20}.  The calculation in
E.1 therefore applies to this actual
cover and gives
\[
 0\longrightarrow\mathcal O_{\mathbb P^1}
       \boldsymbol\delta
 \longrightarrow\mathcal V
 \longrightarrow\mathcal O_{\mathbb P^1}(-\infty)
 \longrightarrow0,
 \qquad H^1(\mathbb P^1,\mathcal V)=0.
\]
The explicit finite averaging and Laurent-series
splitting in E.1 kill the additive cocycle.
The resulting translations glue the local
comparisons to a biholomorphism of $X_0$
with the prescribed member at $c(0)$.

The same solution works over a smaller parameter
disc.  On the fixed collars, Cauchy integral
formulas for the Laurent coefficients preserve
holomorphic dependence on $s$, as does finite
averaging.  Thus the solved translations and
their glued comparison are holomorphic in
$s$.  Their ambiguity is the connected
scalar translation group of E.1.  This
proves both the central and family assertions.

\end{enumerate}
\end{proof}



\paragraph{Scope of \ref{push:src:E.2.21}}

The known-pair locus is open by \SeriesRef{push:src:E.2.2} and
is now stable under full punctured-disc
smooth specializations.  These two facts
alone do not prove that an arbitrary open
known locus in a reduced analytic base
is topologically closed.  A holomorphic
disc through its boundary need not have
its entire punctured image in that locus.
The needed generic propagation argument
is proved in \ref{push:src:E.2.22}.  \ref{push:src:E.2.23}--\ref{push:src:E.2.24} then
establish the arbitrary-base and continuous-family
conclusions asserted in \ref{push:src:E.2.4}.

\Uses{the hypotheses, notation and definitions stated here; \cref{push:src:E.2.21}; \cref{push:src:E.2.22}; \cref{push:src:E.2.23}; \cref{push:src:E.2.24}; \cref{push:src:E.2.4}}


\begin{proposition}[propagation from a known open subset of a disc]\hypertarget{tag-0338}{}\label{push:src:E.2.22}
\StatementRecord{push:src:E.2.22}{proved}{0338}{\SeriesRef{push:src:E.2.21}, \SeriesRef{push:src:E.2.5}, \SeriesRef{push:src:E.2.18}, \SeriesRef{push:src:E.2.19}, \SeriesRef{push:src:E.2.20}; \cite[Chapter~II, Theorem~8.8]{Demailly2012E}; \cite[Chapter~VI, Theorem~1.1 and Table~10]{BHPV2004E}; \cite[Chapter~II, Definition~7.9 and Theorem~7.12]{Demailly2012E}}

Let $\pi:\mathcal X\to\Delta$ be a proper holomorphic
submersion with fibres diffeomorphic to $S^6$.  Suppose
there is a nonempty open subset $U\subset\Delta$ such
that every $X_s$, $s\in U$, is a member of the fixed
coprime multiplicity family $(3k,4l)$ of \SeriesRef{push:src:E.2.1}.
Then every fibre over $\Delta$ is a member of that
same family.
\end{proposition}


\begin{proof}
\begin{enumerate}[label={$\langle1\rangle$\arabic*.},leftmargin=*]
\item 
Local miniversality in \SeriesRef{push:src:E.2.1} supplies actual
holomorphic comparisons with the prescribed family
locally over $U$, with locally constant integral
markings.  Thus the known open set has the relative
root and normal-crossing charts used below; only
pointwise isomorphisms are not being substituted
for these relative comparisons.

The proof first propagates the actual geometric data
outside finitely many successive discrete analytic
exceptional sets.  It then fills these sets in reverse
order by \ref{push:src:E.2.21}.  Later exceptional sets are only
claimed to be discrete on the domain then under
consideration; they may accumulate at points removed
at an earlier stage.

\noindent\textbf{Step 1. The intrinsic line and the relative pencil.}
Put $A=12kl$ and $D=kl-3k-4l\ne0$.
The argument of \ref{push:src:E.2.5} constructs the unique line
bundle $\mathcal L$ on the whole total space with
\[
 \mathcal L^D=K_{\mathcal X/\Delta}^{A}.
\]
Indeed Ehresmann transport over the contractible disc
gives $H^1(\mathcal X,\Z)=H^2(\mathcal X,\Z)=0$,
so the exponential sequence identifies its Picard
group with a complex vector space.  This construction
needs no hypothesis on the fibres outside $U$.

Proper coherence makes $\mathcal V=\pi_*\mathcal L$
a torsion-free, hence locally free, sheaf on the
disc.  Its rank is two because it has that rank
on $U$.  At every parameter its fibre injects
into $H^0(X_s,\mathcal L_s)$ by the multiplication
exact sequence used in \ref{push:src:E.2.5}.  Choose a frame
of $\mathcal V$ on the disc.  The common zero
locus of its two evaluated sections is closed
analytic in $\mathcal X$.  Its image is closed
analytic by properness and Remmert's theorem
\cite[Chapter~II, Theorem~8.8]{Demailly2012E}.
It misses $U$, so is a discrete analytic subset.
Remove it temporarily.

On the remaining connected parameter curve $S$
the sections define a proper holomorphic map
\[
 g:\mathcal X_S\longrightarrow\mathbb P^1\times S,
 \qquad \mathcal L=g^*\mathcal O(1).
 \tag{E.2.22.1}
\]
Each $g_s$ is nonconstant because its two sections
are independent.  Its fibres have pure dimension
two; the smooth total space and smooth target
then give flatness.  The fibres are connected:
the finite Stein factor is generically of degree
one, as it is so above $U$, and a finite
birational map to the normal target is an
isomorphism.  Its source is irreducible, so
there is no additional vertical component
in the Stein factor.  We retain
\[
 K_{\mathcal X_S/S}^{A}=g^*\mathcal O(D).
 \tag{E.2.22.2}
\]

\noindent\textbf{Step 2. Critical values and ordinary tori.}
The relative critical locus of $g$ is analytic
and proper over the target.  Its image has
finite intersection with each target fibre:
the nonconstant proper map $g_s$ to a curve
has finitely many critical values.  Thus the
reduced critical image is finite over $S$.
Its one-dimensional part has degree three,
as can be read over $U$; every such component
dominates $S$.  Remove the images of its
zero-dimensional components and its branch
and singular points.  These are further
proper analytic discrete subsets, absent
over the known open set.

A finite etale parameter cover now labels
the three critical sections.  Pointwise
classification after this cover descends;
images of later discrete exceptional sets
remain discrete because this cover is finite.
One may work on each connected component,
each of which meets the inverse image of $U$.
Send the labelled sections to $0,1,\infty$.
There are exactly these three critical values
on the current parameter domain.

The ordinary family over
$B^\circ\times S$, where
$B^\circ=\mathbb P^1\setminus\{0,1,\infty\}$,
is smooth and proper.  Its compact surface
fibres have the topology of the known two-tori
by Ehresmann's theorem and connectedness.
Adjunction gives integral $c_1(K_F)=0$,
since $F$ is a smooth fibre in a complex
six-sphere and its normal line is trivial.
These surfaces are complex tori.  To spell
out the classification step, $b_1(F)=4$
is even, so $F$ is K\"ahler.  The zero
canonical class excludes exceptional curves
and ruled surfaces.  On a K\"ahler surface,
a nonzero effective divisor cannot have
zero real Chern class; hence a nonzero
pluricanonical section has no zeros and
cannot give positive Kodaira dimension.
The compact-surface classification leaves
the Kodaira-dimension-zero list, in which
$b_1=4$ identifies a two-torus
\cite[Chapter~VI, Theorem~1.1 and Table~10]{BHPV2004E}.

\noindent\textbf{Step 3. The finite root charts.}
For any labelled critical section $\Gamma_j$,
every component of the Cartier divisor $g^*\Gamma_j$
dominates $S$.  A vertical divisor component would
be an entire fibre $X_s$, making $g_s$ constant.
Every dominating component meets $U$, where the
reduced fibre is irreducible.  Thus there is
exactly one such component, and its multiplicity
is the known one.  In particular
\[
 g^*\Gamma_1=(3k)\mathcal D_1,
 \qquad g^*\Gamma_2=(4l)\mathcal D_2,
 \qquad g^*\Gamma_0=\mathcal W,
 \tag{E.2.22.3}
\]
with reduced total hypersurfaces on the right.
For clarity, two distinct dominating components
could not both restrict to the known irreducible
reduced fibre.  Each is a Cartier prime in
the smooth total space, is torsion-free and
hence flat over the smooth parameter curve,
and has pure two-dimensional fibres.  Thus
its fibre over any parameter in $U$ contains
the whole known irreducible reduced surface.
At a general smooth point of
that fibre the known local equation has only
one reduced total branch.  Properness forces
every dominating component to meet that open
part of the known fibre, a contradiction.

For a finite multiplicity $m$, write locally
$q=a^m u$, with $a$ a reduced local equation
of $\mathcal D_j$ and $u$ a unit.  The normalization
of the cyclic root base change $r^m=q$ is
locally the disjoint union of the graphs
\[
 r=\zeta a u^{1/m},\qquad \zeta^m=1.
\]
It is therefore smooth and etale over the
actual total space, with free cyclic deck
action.  On $r\ne0$ its map to the root
disc times $S$ is already smooth.  Its
critical locus is supported over $r=0$,
is proper over $S$, and is empty over $U$.
Delete its discrete analytic image for
each of the two roots.  The resulting
proper smooth root families have connected
fibres by their Stein factors and are
families of two-tori by the same topology,
adjunction and surface argument above.

\noindent\textbf{Step 4. A finite analytic test for relative normal crossings.}
The total cusp $\mathcal W$ in (E.2.22.3)
is a reduced embedded hypersurface in
the smooth fourfold.  Take its finite
analytic normalization
\[
 \nu:\widetilde{\mathcal W}\longrightarrow
       \mathcal W\longrightarrow\mathcal X_S.
\]
Existence and finiteness are the analytic
normalization theorem
\cite[Chapter~II, Definition~7.9 and Theorem~7.12]{Demailly2012E}.
The normalization is proper over $S$.
First remove the proper images of its
singular locus, and then of the critical
locus of its map to $S$.  Each is analytic
and misses $U$.  On the resulting domain
the normalization is smooth over $S$.
Next remove the proper image of the
rank-drop locus of
\[
 T_{\widetilde{\mathcal W}/S}
       \longrightarrow\nu^*T_{\mathcal X_S/S}.
\]
This locus too is analytic and absent over
$U$.  The finite map $\nu$ is now an immersion.

For $r=2,3,4$, let $R_r$ be the locus
of ordered, pairwise distinct preimages
in the $r$-fold fibre product of
$\widetilde{\mathcal W}$ over $\mathcal X_S$.
Its partial diagonals are both open and
closed: a finite immersion is locally
an embedding, so its diagonal is open,
and separatedness makes it closed.
Consequently $R_r$ is closed analytic
in the proper fibre product and is
proper over $S$.

At a point of $R_r$ the $r$ local sheets
have relative conormal lines in the same
rank-three bundle $T^*_{\mathcal X_S/S}$.
On $R_2$ and $R_3$, delete the proper
images of the rank-drop loci of these
two or three lines.  Delete also the
image of $R_4$.  These are closed
analytic images and miss $U$, hence
are discrete.  There are only finitely
many such deletions.

At a remaining point the total cusp
has at most three local sheets, each
embedded smoothly and submersive over
$S$, with independent vertical normal
covectors.  Since the \emph{total}
$\mathcal W$ is a reduced embedded
hypersurface, its local ideal is exactly
the product of the distinct local
equations of those sheets.  No finite
subring or seminormality assumption
is being used.  Independence of the
vertical normals makes these equations,
the parameter on $S$, and the remaining
coordinates a holomorphic coordinate
system.  Thus
\[
 q=x,\qquad q=xy,\qquad\hbox{or}\qquad q=xyz
 \tag{E.2.22.4}
\]
after absorbing a unit.  These are
actual relative normal-crossing charts;
in particular the remaining cusp fibres
are reduced.  This proof did not assume
that normalization commutes with base
change before obtaining the charts.

\noindent\textbf{Step 5. The normalized toric pair and its action.}
The smooth proper normalization and
its conductor strata have the same
topology, incidence and normal degrees
as over $U$.  A further finite etale
cover can label these finitely many
strata.  The argument of \ref{push:src:E.2.18} applies
to every fibre: the six conductor
curves are a hexagon of $(-1)$-curves,
three alternate contractions give a
triangle in $\mathbb P^2$, and the
normalized pair is $(\mathrm{dP}_6,D)$.
Its relative logarithmic tangent bundle
is the evaluation pullback of a rank-two
bundle on the parameter curve.

In this evaluation description the
conductor descent maps are holomorphic
endomorphisms of that rank-two bundle.
They are constant along the compact
conductor curves.  They equal the identity
over $U$, where both toric logarithmic
fields descend.  The holomorphic identity
theorem on the connected parameter curve
makes them the identity everywhere.
This is the descent proof of \ref{push:src:E.2.19};
it requires a nonempty known open set,
not specifically a punctured disc.
Grauert evaluation, commutation on the
ordinary tori, completeness and the
implicit extension of the two primitive
toric periods then give the actual
marked cusp torus action locally on
the current parameter curve.  The
integral monodromy and triple-point
characters are constant by transport
and are the prescribed ones.

\noindent\textbf{Step 6. Continuation of the full period kernel.}
Fix an ordinary target point $t_*$.
The integral local system over
$B^\circ\times S$ has constant marked
peripheral representation, up to a
choice of marking: its matrices are
transported from those over $U$.
Parameter monodromy commutes with the
base representation.  It therefore
preserves the cusp invariant lattice
\[
 \ker(A_0-I)=\Z\widehat w\oplus\Z\widehat\delta.
\]
The wedge of the two corresponding
holomorphic period vectors at $t_*$
has a well-defined zero locus on $S$:
changing its integral basis multiplies
the wedge by a nonzero determinant.
It is nonzero over $U$, so that zero
locus is discrete analytic.  Remove
it as one final exceptional layer.

On a connected universal cover of
the remaining parameter curve, normalize
the Lie frame at $t_*$ by these two
columns.  The $\widehat\gamma$ column
then defines a holomorphic function
$c(s)$ by subtracting the fixed
reference value of $\beta^{\rm part}(t_*)$.
All other column identities at $t_*$
hold over a lift of $U$ and hence
everywhere by the identity theorem.

Now work on the universal cover of
the whole ordinary base times this
parameter cover.  The actual holomorphic
period kernel is a rank-two subbundle
of the complexified marked rank-four
lattice.  Compare it with
\[
 \ker\Pi_{c(s)}(t).
 \tag{E.2.22.5}
\]
The latter matrix always has complex
rank two, since its $\widehat w$ and
$\widehat\delta$ columns are the
standard basis.  These two holomorphic
Grassmannian maps agree over the
product with a lift of $U$.  They
therefore agree on the entire
connected ordinary cover.
This comparison does not assume
that the two actual columns were
independent at every $t$; equality
of the kernels proves it afterwards.
In a local actual Lie frame the
identity says that the full period
matrix is $W(t,s)\Pi_{c(s)}(t)$,
where $W$ is its two marked columns.
If $W$ had rank at most one, every
period column would have that rank,
contradicting the real rank four
of the compact torus lattice.

The actual ordinary and finite-root
fibres are compact tori.  The exact
period determinant is nonzero there.
Its sign is the negative sign on
$U$ and stays negative by connectedness
and nonvanishing.  Equality of the
period kernels also extends to the
smooth finite root charts; their
nonzero determinant preserves this
sign at the root point where the
maximum in \SeriesRef{fam:src:C.2.3} is attained.  Thus
$\operatorname{Im}c(s)+M<0$ by
the precise threshold of \SeriesRef{fam:src:C.2.3}.
This places $c(s)$ in the
admissible interior domain, including
strictness at its boundary.  The
cusp multipliers are the corresponding
marked reference multipliers.  The
proof of \ref{push:src:E.2.20} now reconstructs
the actual cusp at each parameter:
relative normal-crossing transport
gives its known cover and incidence,
the action gives the same characters,
and adjacent-origin comparison gives
the monomial times unit extensions.
The proof works along any path from
$U$ and locally in the current
parameter; it does not need an
already known punctured neighbourhood
at the point being classified.

The finite affine lattice classes
and integral collar classes are
locally constant on this connected
domain, hence remain those over
$U$.  The additive comparison in
E.1 and \ref{push:src:E.2.21} identifies every
remaining actual threefold with
the reference member at $c(s)$.
The pointwise conclusion descends
from every finite parameter cover
used to label the critical or
conductor strata.

\noindent\textbf{Step 7. Filling the exceptional layers.}
We have made only finitely many
successive deletions.  At each stage
the deleted subset was closed
analytic and discrete on the
parameter domain then present.
Finite covers do not change that
property after taking their proper
images.  On the final domain all
fibres are known members of the
fixed pair.  At a point of the
last deleted set, choose a small
disc containing no other point of
that set and avoiding all earlier
deleted sets.  Its punctured disc
consists of known members.  \ref{push:src:E.2.21}
identifies the central fibre as
well.  Fill the entire last layer
in this way, then the preceding
one, and continue in reverse order.
This is valid even when a later
layer accumulates at an earlier
deleted point.  Finally every
fibre over the original disc is
known.  No analytic curve selection
for an arbitrary open locus has
been assumed.

\end{enumerate}
\end{proof}



\begin{theorem}[separation under arbitrary deformation chains]\hypertarget{tag-0339}{}\label{push:src:E.2.23}
\StatementRecord{push:src:E.2.23}{proved}{0339}{\SeriesRef{push:src:E.2.4}, \SeriesRef{push:src:E.2.22}; \citet[Chapter~II, Theorem~4.19, pp.~95--96]{Demailly2012E}; \cite[Chapter~II, Theorems~7.6 and~7.12]{Demailly2012E}}

In a proper smooth holomorphic family of compact complex
manifolds over a connected reduced complex analytic base,
if one fibre is a sphere member of \SeriesRef{push:src:E.2.1} with coprime
multiplicities $(3k,4l)$, then every fibre is a constructed
member with that same pair.

Consequently two of the constructed sphere members are
deformation equivalent, allowing chains of proper smooth
families over connected complex analytic bases, if and
only if their ordered multiplicity pairs agree.  In
particular the distinct coprime pairs give distinct
deformation classes, and \ref{push:src:E.2.4} is resolved by this theorem.
\end{theorem}


\begin{proof}

\begin{enumerate}[label={$\langle1\rangle$\arabic*.},leftmargin=*]
\item Fix the pair and let $U$ be the locus of its constructed
members in the base $B$.  It is open by the actual local
miniversality of \SeriesRef{push:src:E.2.1} and \SeriesRef{push:src:E.2.2}.  We prove that it is
closed without assuming a punctured analytic disc can
be selected inside an arbitrary open subset.

\item First consider any proper smooth family over a connected
Riemann surface.  If its known-pair locus is nonempty,
Ehresmann's theorem makes all its fibres diffeomorphic
to $S^6$, and that locus is open.  At any point of its closure choose a
coordinate disc meeting that locus.  The intersection
contains a nonempty open subset of the disc, so \ref{push:src:E.2.22}
makes the entire disc known.  Thus the known-pair locus
is closed and equals the connected Riemann surface.
Only this disc propagation, not a special choice of
punctured disc, is used below.

\item Let $b\in\overline U\subset B$.  If $b\in U$ there
is nothing to prove.  Otherwise choose a small local
representative and an irreducible branch $T$ through
$b$ which meets $U$ arbitrarily near $b$.  There are
only finitely many such local branches.  Its dimension
is positive.  The local parametrization theorem gives
a finite proper projection
\[
 \rho:T\longrightarrow\Delta^d,
 \qquad \rho^{-1}(0)=\{b\},
 \tag{E.2.23.1}
\]
after shrinking the representative and translating
$b$ to the origin.  This is the finite projection
of \citet[Chapter~II, Theorem~4.19, pp.~95--96]{Demailly2012E}.
The isolated central fibre is part of choosing
the generic projection for the germ; the
Weierstrass representative is then taken small
enough to contain no other point over zero.

\item Choose $a\in U\cap T$ in this representative,
with $\rho(a)\ne0$.  Let $L$ be the complex
line through $0$ and $\rho(a)$, restricted
to the polydisc.  An irreducible component
$C$ of the reduced analytic set
$\rho^{-1}(L)$ which contains $a$ has
dimension one.  Indeed $d-1$ linear
equations give local dimension at least
one, while finiteness of $\rho$ gives
dimension at most one.  Its finite
proper image in $L$ is one-dimensional,
so is all of $L$.  Therefore $C$ contains
the unique point $b$ over zero.

\item Normalize $C$.  Its normalization
$\widetilde C$ is a connected Riemann
surface, finite over $C$; analytic
normalization exists and a normal
curve is smooth by
\cite[Chapter~II, Theorems~7.6 and~7.12]{Demailly2012E}.
It has points above both $a$ and $b$.
Pull back the given proper smooth
family to $\widetilde C$.  Its
known-pair locus contains the
inverse image of the open set
$U\cap C$ near $a$.  The preceding
Riemann-surface argument makes
every fibre on $\widetilde C$
known, in particular the fibre
above $b$.  Thus $b\in U$.

\item We have proved that $U$ is open
and closed.  Since $B$ is
connected and $U$ is nonempty,
$U=B$.  A nonreduced base
does not add fibre isomorphism
classes: restriction to its
reduction preserves both the
fibres and smoothness of the
family.  Thus the same conclusion
also applies to such a base
when it is admitted in the
definition of deformation
equivalence.

\item Applying this result to each
link in a finite deformation
chain preserves the pair.
Conversely, E.1 identifies
all admissible sphere data
with a fixed pair with the
connected period family,
using its holomorphic
comparisons and its period
lattice $2kl\Z$.
The admissible half-plane
and its punctured-disc
quotient are connected.
Thus all members with
that pair are deformation
equivalent.  This proves
both directions.

\end{enumerate}
\end{proof}



\begin{corollary}[separation for continuous complex structures]\hypertarget{tag-033A}{}\label{push:src:E.2.24}
\StatementRecord{push:src:E.2.24}{proved}{033A}{\SeriesRef{push:src:E.2.23}; \citet[Section~3.2, Theorem~3.4 and Remark~3.6,
pp.~898--899]{Meersseman2019E}; \cite[Section~2.1, pp.~883--884]{Meersseman2019E}}

Let $\mathcal I(S^6)$ be the space of smooth integrable
almost-complex structures on the fixed smooth sphere,
with its $C^\infty$ topology.  For a fixed coprime pair
$(3k,4l)$, let $\mathcal I_{k,l}$ consist of the structures
whose complex manifolds are biholomorphic to constructed
members with that pair.  Then $\mathcal I_{k,l}$ is open
and closed in $\mathcal I(S^6)$.

Consequently a continuous family of such structures
over any connected topological parameter space cannot
contain constructed members with two different pairs.
The same separation holds in the quotient spaces by
$\operatorname{Diff}_0(S^6)$ or by all orientation-preserving
diffeomorphisms, with their quotient topologies.
This statement does not identify the number of marked
connected components belonging to a single pair.
\end{corollary}


\begin{proof}

\begin{enumerate}[label={$\langle1\rangle$\arabic*.},leftmargin=*]
\item We use the classical Kuranishi local-completeness
theorem at an arbitrary $J_0\in\mathcal I(S^6)$,
including an initially unconstructed structure.
The complex orientation is locally constant, so the
oriented formulation applies on each such neighbourhood.
The precise $C^\infty$ version used here is
\citet[Section~3.2, Theorem~3.4 and Remark~3.6,
pp.~898--899]{Meersseman2019E}.  It supplies a connected
open neighbourhood $\mathcal U$ of $J_0$, a
finite-dimensional Kuranishi analytic space $K$
containing $J_0$, and a continuous retraction
$\Xi:\mathcal U\to K$ such that every $J\in\mathcal U$
is obtained from $\Xi(J)$ by a diffeomorphism isotopic
to the identity.  In that theorem the local product
map is first constructed in Sobolev completions;
elliptic regularity gives precisely this restriction
to smooth operators.  The topology on these operators
is the $C^\infty$ topology of
\cite[Section~2.1, pp.~883--884]{Meersseman2019E}.

\item The Kuranishi space carries its proper smooth
holomorphic deformation family.  We can take $K$
connected: in the stated local theorem it is the
continuous retract of the connected $\mathcal U$.
If necessary replace $K$ by its reduction, which
changes neither its points nor the complex manifolds
in its fibres.  Suppose $\mathcal U$ contains some
$J\in\mathcal I_{k,l}$.  Its Kuranishi representative
$\Xi(J)$ has the same complex manifold up to
biholomorphism.  \ref{push:src:E.2.23} applied to the connected
Kuranishi family therefore makes every point of
$K$ a member of the same pair.  Since all points
of $\mathcal U$ are biholomorphic to their
Kuranishi representatives, it follows that
$\mathcal U\subset\mathcal I_{k,l}$.

\item Thus at every $J_0$ the neighbourhood $\mathcal U$
is either contained in $\mathcal I_{k,l}$ or
disjoint from it.  This proves both openness and
closedness.  The inverse image of this set under
a continuous map from a connected parameter space
is then empty or the whole space, proving the
continuous-family assertion.  Finally
$\mathcal I_{k,l}$ is saturated under the indicated
diffeomorphism groups, because its definition uses
only the biholomorphism class.  Its image in either
quotient is open and closed by the definition of
the quotient topology.  No assertion that the
saturated set itself is connected is needed.

\end{enumerate}
\end{proof}





\subsection{E.5. The two-parameter family and its actual singular fibres}
\begin{proposition}[the input configuration]\hypertarget{tag-033N}{}\label{push:src:E.5.1}
\StatementRecord{push:src:E.5.1}{proved}{033N}{\SeriesRef{push:src:E.5.2}; \cite{ViaTwo26push}; \cite[Lemma 2.1]{ViaTwo26push}}

The proposed configuration with three rank-one reduced singular fibres
does not describe the construction of \cite{ViaTwo26push}. The elliptic surface
$y^2=x^3+tx+1$ has fibre $III^*$ at infinity.  The constructed
threefold has two rank-one fibres at $p_1,p_2$, a rank-two
hexagonal fibre at $p_3$, and the multiple bielliptic fibre
$4B_*$ at infinity.
\end{proposition}


\begin{proof}
\begin{enumerate}[label={$\langle1\rangle$\arabic*.},leftmargin=*]
\item 
\noindent\textbf{Step 1.} Put $s=t^{-1}$, $X=s^2x$, and $Y=s^3y$.  The minimal equation is
$Y^2=X^3+s^3X+s^6$, with coefficient and discriminant orders
$(3,6,9)$; this is the $III^*$ equation, as in
\cite[Lemma 2.1]{ViaTwo26push}.  The discriminant at finite
points is $-16(4t^3+27)$ with three simple zeros.  In the
threefold's multiplicative periods the orders at these points
are respectively $(1,0),(1,0),(1,1)$:
\[
 h(z,v)=(q\alpha(q)z,\beta(q)v),\qquad
 k(z,v)=(\gamma(q)z,\lambda q^\epsilon\delta(q)v),
\]
where all four coefficients are units, $\epsilon=0$ at $p_1,p_2$,
and $\epsilon=1$ at $p_3$.  Thus the first two torus ranks
drop by one and the last by two.  The explicit matrices of
\ref{push:src:E.5.2} independently have nilpotent ranks $1,1,2$; their exact
calculation rules out the assertion that all three are rank one.
These are the types stated in Proposition 3.1 and Theorem 1.1
of the cited source, not a correction to that source.

\end{enumerate}
\end{proof}



\begin{proposition}[monodromy, centralizer, and integral lattices]\hypertarget{tag-033O}{}\label{push:src:E.5.2}
\StatementRecord{push:src:E.5.2}{proved}{033O}{\cite{ViaTwo26push}}

Use the source's left-action marking
$H^1(F,\mathbb Z)=\mathbb Z\langle\psi,u,w,\delta\rangle$
and the ordered meridians $*,1,2,3$ with product one.
The cohomological matrices are
\[
\begin{aligned}
 T_*&=\begin{pmatrix}1&0&-2&-1\\0&0&-1&0\\0&1&0&-1\\0&0&0&1\end{pmatrix},&
 T_1&=\begin{pmatrix}1&0&2&2\\0&1&1&1\\0&0&1&0\\0&0&0&1\end{pmatrix},\\
 T_2&=\begin{pmatrix}1&0&0&0\\0&1&0&0\\0&-1&1&0\\0&0&0&1\end{pmatrix},&
 T_3&=\begin{pmatrix}1&0&0&-1\\0&1&1&0\\0&0&1&0\\0&0&0&1\end{pmatrix}.
\end{aligned}
\]
On ordinary dual homology put $H_i=(T_i^{-1})^t$,
$\ell=(1,-1,-1,0)^t$ and $\widehat\delta=(0,0,0,1)^t$.
Then
\[
 \operatorname{Cent}_{M_4(\mathbb Q)}(H_*,H_1,H_2,H_3)
 =\mathbb Q I\oplus\mathbb Q E_{41},
\]
and the cohomological centralizer is
$\mathbb Q I\oplus\mathbb Q E_{14}$.
The fixed and norm lattices at the multiple fibre are
\[
 \begin{aligned}
 \Lambda^{H_*}&=\mathbb Z\ell\oplus\mathbb Z\widehat\delta,\\
 N_4\Lambda&=4\mathbb Z\ell\oplus2\mathbb Z\widehat\delta,\\
 \Lambda^{H_*}/N_4\Lambda&\simeq\mathbb Z/4\oplus\mathbb Z/2.
 \end{aligned}
\]
The primitive vanishing lattices at the three reduced fibres are
\[
 K_1=\mathbb Z(\widehat w+\widehat\delta),\qquad
 K_2=\mathbb Z\widehat u,\qquad
 K_3=\mathbb Z\widehat w\oplus\mathbb Z\widehat\delta.
\]
Their invariant lattices are
\[
 \begin{split}
 \Lambda^{H_1}&=\mathbb Z(\widehat\psi-2\widehat u)
       \oplus\mathbb Z\widehat w\oplus\mathbb Z\widehat\delta,\\
 \Lambda^{H_2}&=\mathbb Z\widehat\psi\oplus
       \mathbb Z\widehat u\oplus\mathbb Z\widehat\delta,\\
 \Lambda^{H_3}&=K_3.
 \end{split}
\]
\end{proposition}


\begin{proof}
The matrices are those of \citet[Proposition~4.2]{ViaTwo26push}, whose
derivation uses the elliptic monodromies, the winding of $P$,
and the primitive rank-two cusp.  In particular it determines
the winding integer by the saturation index $d^2=1$.  Direct multiplication gives
$T_*T_1T_2T_3=I$, $T_*^4=I$ and
$(T_i-I)^2=0$ for $i=1,2,3$.
The dual homology matrices are
\[
\begin{aligned}
 H_*&=\begin{pmatrix}1&0&0&0\\-2&0&-1&0\\0&1&0&0\\1&1&0&1\end{pmatrix},&
 H_1&=\begin{pmatrix}1&0&0&0\\0&1&0&0\\-2&-1&1&0\\-2&-1&0&1\end{pmatrix},\\
 H_2&=\begin{pmatrix}1&0&0&0\\0&1&1&0\\0&0&1&0\\0&0&0&1\end{pmatrix},&
 H_3&=\begin{pmatrix}1&0&0&0\\0&1&0&0\\0&-1&1&0\\1&0&0&1\end{pmatrix}.
\end{aligned}
\]
Solving the sixty-four rational equations $CH_i=H_iC$
gives $C=aI+bE_{41}$ with free $a,b$. Transposition gives the
linear bijection between the two centralizers: the equation
$CH_i=H_iC$ is equivalent to $C^tT_i=T_iC^t$, and
$E_{41}^t=E_{14}$.  Solving $(H_*-I)v=0$ gives
$v=(a,-a,-a,d)^t$ integrally, and
\[
 N_4=I+H_*+H_*^2+H_*^3
 =\begin{pmatrix}4&0&0&0\\-4&0&0&0\\-4&0&0&0\\0&2&-2&4\end{pmatrix}.
\]
In the basis $(\ell,\widehat\delta)$ its columns are
$(4,0),(0,2),(0,-2),(0,4)$, giving exactly the displayed
norm group.  Each $K_i$ is the integral image of
$H_i-I$; its Smith form has only unit nonzero entries.
The kernels of these three matrices give the remaining fixed
lattices.

\end{proof}



We now specify the parameter domain, because the source proves
a family on a sufficiently small bidisc, not on the entire
$\mu$-plane.  Set
\[
 S_\mu:\ y^2=x^3+tx+\mu t+1,\qquad
 P_\mu=(-\mu,\sqrt{1-\mu^3}),
\]
with the square root equal to one at zero.  Label the three roots
$p_i(\mu)$ of $4t^3+27(\mu t+1)^2$ holomorphically near zero;
$p_3$ is the zero of the chosen linearization.
Choose a holomorphic generator of its one-dimensional space of
sections vanishing at $p_3$ and scale it by $\lambda$.
For the added homology generator take the inverse of the
contracting lift.  On the universal cover of the four-punctured
base its period columns then have the form
\[
 \Pi_{\mu,c}=
 \begin{pmatrix}2z_\mu&\tau_\mu&1&0\\
 \beta_\mu+c&z_\mu&0&1\end{pmatrix},
 \qquad \lambda=e^{-2\pi i c},\qquad
 \partial_c\Pi=e\otimes\psi.
\]
Here $z_\mu$ is a lifted elliptic logarithm in the chosen marking;
changing the generator of the linearization changes
$\beta_\mu$ and the origin of $c$ by a holomorphic function of
$\mu$, and changes neither the integer-difference criterion
below nor the original $\lambda$-labelled bidisc.
This choice of inverse lift fixes the sign of $c$.
At $p_3$ its second period has jump $-1$, exactly canceled by
$H_3\widehat\psi=\widehat\psi+\widehat\delta$.
The covariance identities for all four matrices, including
$\tau_*=-1/\tau$, $z_*=(z-1)/\tau$,
$\beta_* =\beta-1-2(z-1)^2/\tau$, are checked symbolically
in the same program.

\begin{theorem}[classification on the specified parameter domain]\hypertarget{tag-033P}{}\label{push:src:E.5.3}
\StatementRecord{push:src:E.5.3}{proved}{033P}{\SeriesRef{push:src:E.5.1}, \SeriesRef{push:src:E.5.2}; \cite{ViaTwo26push}}

There are $\epsilon>0$ and a real constant $C$ such that the
source construction, with its fixed logarithmic vector $\ell$,
defines a proper holomorphic family
$X_{\mu,c}$ over
$\{ |\mu|<\epsilon\}\times\{\operatorname{Im}c<C\}$.
It has the singular fibres of \ref{push:src:E.5.1} and the markings above.
Two of its members are biholomorphic if and only if
\[
 \mu'=\mu,\qquad c'-c\in2\mathbb Z.
\]
Every such biholomorphism covers the identity of the elliptic
base.  The quotient parameter in the second variable is
$v=e^{-\pi i c}$, with $\lambda=v^2$.
On a sufficiently small bidisc centred at any
$(0,\lambda_0)$ in Theorem 1.1 of \cite{ViaTwo26push}, choose
one logarithm of $\lambda$.  Within that actual bidisc,
\[
 X_{\mu,\lambda}\simeq X_{\mu',\lambda'}
 \quad\Longleftrightarrow\quad
 (\mu,\lambda)=(\mu',\lambda').
\]
\end{theorem}


\begin{proof}

\begin{enumerate}[label={$\langle1\rangle$\arabic*.},leftmargin=*]
\item The equations of Proposition 6.4 in \cite{ViaTwo26push}
give the simultaneous elliptic family, line bundle, linearization,
and good-reduction comparison.  On a smaller closed $\mu$-disc
all unit factors in its four local fillings are uniformly bounded
above and below.  Its contraction estimates therefore have a
common nonzero bound $|\lambda|<\eta$.  Pulling back by
$\lambda=e^{-2\pi i c}$ gives the indicated half-plane.
The integral markings and the four-torsion section $\ell/4$
continue uniquely on this simply connected parameter domain.
On every relatively compact subset the source's finite root
quotients and the fixed rank-one and hexagonal fans give
holomorphic charts with uniform bounds.  Their comparisons are
the specified section-preserving ones, so they agree on
intersections and define a family.  Properness is local on
the parameter domain and follows from the compact four-piece
construction there; smoothness follows from the same free
finite action and unimodular charts.  We make no assertion
about continuation to arbitrary $\mu$.

\item The intrinsic fibration is the anticanonical pencil.  To verify
the needed input from Proposition 6.1 of the source, on the
smooth core the line-coordinate form $dv/v$ wedged with a
base two-form identifies $K_X$ with $f^*\mathcal O(-1)$.
The forms $dq\wedge d\log z\wedge d\log v$ extend without
zeros or poles in both rank-one charts $q=ab$ and rank-two
charts $q=abc$.  At the finite root at infinity the source's
coordinate change gives
\[
 ds\wedge\frac{dX}{Y}\wedge\frac{dv}{v}
 =4\,dr\wedge\frac{dx'}{y'}\wedge\frac{dv'}{v'},
 \qquad s=r^4.
\]
It is invariant under the actual affine deck map and descends
through its free quotient.  Thus the same canonical formula
holds everywhere.  Connected fibres and projection formula
give $H^0(-K_X)=f^*H^0(\mathcal O(1))$, without base points.
A biholomorphism therefore induces an affine automorphism of
the elliptic base, fixing its distinguished multiple point
at infinity and its distinguished hexagonal point $p_3$.
It may initially interchange the two rank-one points.

\item Put $Q(\mu)=(p_1-p_3)/(p_2-p_3)$.
The source computes $p_i'(0)=-9/(2p_i(0))$.
If $Q'(0)$ vanished, the three velocities would have the form
$Ap_i(0)+B$; the polynomial $Az^2+Bz+9/2$ would vanish
at three distinct numbers, a contradiction.
Also $Q(0)^2-Q(0)+1=0$, so $Q(0)\ne Q(0)^{-1}$.
Shrink $\epsilon$ so that $Q$ is injective and its image
is disjoint from the image of $Q^{-1}$.  Preservation or
interchange of the rank-one values gives respectively
$Q(\mu')=Q(\mu)$ or $Q(\mu')=Q(\mu)^{-1}$.
The second possibility is excluded, and the first gives
$\mu'=\mu$.  The base automorphism now fixes three distinct
points individually and is the identity.

\item On regular fibres a biholomorphism is affine, and its integral
linear map commutes with all four $H_i$.  \ref{push:src:E.5.2} gives
$aI+bE_{41}$.  Comparing the constant third and fourth
period columns makes its complex linear part $aI$;
unimodularity gives $a=\alpha\in\{1,-1\}$.
The first column then gives $h=c'-c\in\mathbb Z$ and
lattice map $\alpha(I-hE_{41})$.
Normalize after the actual fourth-root base change.  Its
affine generator is the one with fourth power $\ell$, not
an arbitrary generator of its cyclic subgroup.  Summing the
translation defect around this generator gives
\[
 (I+hE_{41})\ell-\alpha\ell
 =(1-\alpha)\ell+h\widehat\delta\in N_4\Lambda.
\]
The norm lattice in \ref{push:src:E.5.2} forces $\alpha=1$ and $h$ even.
This proves necessity, including the multiple fibre's full
non-reduced structure.

\item We prove sufficiency with all four collar conditions.
Let $h$ be even and put
\[
 b_*=-\tfrac h4(\widehat u+\widehat w),\qquad
 L_*=-\tfrac h2\widehat u.
\]
Direct multiplication gives
\[
 (I-H_*)b_*=\tfrac h4\widehat\delta+L_*.
\]
Consequently $\zeta\mapsto\zeta+\Pi_{\mu,c}b_*$ on the
complete root torus family conjugates the two affine actions,
modulo the integral arrow $L_*$.  It and its inverse extend
over $r=0$ and descend.  On the punctured root collar its
middle translation includes
$h(\log r)/(2\pi i)e$; this is the derivative-free collar
term for the change of the invariant translation
$\Pi\ell/4$, not a logarithm used on the central chart.

\item At a rank-one point the vanishing coordinate is the first
toric coordinate in its chain cover, while the other
direction is the universal cover of the elliptic factor.
An extending translation has integral winding only in $K_j$.
After removing the corresponding power of $q$, its toric
coefficient is a holomorphic unit and its elliptic translation
extends holomorphically.  These are the full analytic collar
conditions.  Necessity follows by lifting a transverse disc
through a smooth central point to this free cover: its residual
defect in $\Lambda/K_j$ is zero.  The coordinate ratios on
the cover then give the asserted unit and holomorphic
remainder.  Conversely an integral power of $q$ translates
the infinite chain by an integer number of components;
the holomorphic remainders give torus and elliptic translations.
These maps and their inverses extend at both branches of every
node.  At the hexagonal point the same description uses two
toric windings in $K_3$, and integer translation of the fan
extends over its three double curves and both triple points.
This criterion includes regularity of the remainder; membership
of the defect in $K_j$ alone would not exclude an essential
singularity of a preassigned punctured translation.

\item Choose the three integral defects
\[
 L_1=0,\qquad L_2=\tfrac h2\widehat u,\qquad
 L_3=-\tfrac h2\widehat w.
\]
They belong to $K_1,K_2,K_3$, respectively, and are realized
by the preceding toric translations with initially constant
unit remainder.  In primitive cocharacter coordinates these
are $1$, $q_2^{h/2}$, and $(q_3^{-h/2},1)$, where the
monodromy defect of $q^a$ is defined as $a$ times the chosen
vanishing generator.  The required sphere relation is
\[
 L_*+H_*L_1+H_*H_1L_2+H_*H_1H_2L_3=0.
\]
It holds as an exact integer identity, checked in the program.
Thus the four local choices have compatible additive defects
and define a torsor of allowed middle translations.

\item We verify the vanishing of its analytic obstruction, including
the difference sheaf.  The root curve has just the order-four
point at infinity and is $\mathfrak C=\mathbb P(1,4)$,
with $K_{\mathfrak C}=\mathcal O(-5)$ and coarse pullback
$\mathcal O_B(1)=\mathcal O_{\mathfrak C}(4)$.
Vertical logarithmic fields on a rank-one chart are generated
by $a\partial_a-b\partial_b$ and the elliptic vector;
on a hexagonal chart they are the two height-zero toric fields.
The root chart has the ordinary vertical tangent frame.
Their linear transition functions depend only on the base;
they define a rank-two bundle $\mathcal E$ on $\mathfrak C$.
The fixed line-coordinate field supplies a subline
$\mathcal O\subset\mathcal E$ everywhere in these logarithmic
frames.  The logarithmic determinant sequence and cancellation
of the three reduced fibre divisors give
$K_X=\bar f^*(K_{\mathfrak C}\otimes(\det\mathcal E)^{-1})$.
The verified canonical formula and
$\bar f_*\mathcal O_X=\mathcal O_{\mathfrak C}$ therefore
give $\det\mathcal E=\mathcal O(-1)$.
The extension splits since
$H^1(\mathfrak C,\mathcal O(1))=0$, by the same two-chart
Laurent calculation with weights $1,4$.
The coarse difference sheaf is $V_B=\rho_*\mathcal E$:
on the root chart this is precisely finite equivariance,
and on each reduced collar it is precisely the holomorphic
remainder just described.  Finite averaging is exact, so
\[
 H^1(B,V_B)=H^1(\mathfrak C,\mathcal E)
 =H^1(\mathcal O\oplus\mathcal O(-1))=0.
\]
The torsor's Cech class vanishes.  Correcting the four local
translations by its holomorphic zero-cochain gives a global
middle translation and exact equality with all four local
extensions.  Their inverses have the opposite allowed
translations and glue as well.  This constructs the required
biholomorphism over $B$.

\item Finally choose a single logarithm of $\lambda$ on the source's
small bidisc.  Its $c$-values can be placed in a strip of
width less than two.  The criterion just proved then says
$c'=c$ and $\mu'=\mu$, which is equivalent to equality of
the original two parameters.  Conversely equality gives the
identity, with the fixed root, rank-one and hexagonal charts.

\end{enumerate}
\end{proof}



The matrix data and the independent-variation argument are in
\cite[Proposition 4.2 and Proposition 6.4]{ViaTwo26push}; the
source states dimension at least two and does not state the
integer classification on the logarithmic lift.  We derived
the latter here from the marked norm and collar calculations.
For E.1--E.5 our searches on 15 November 2026 included the
exact arXiv identifiers together with ``classification'',
``Kuranishi'', ``deformation'', and ``parameters'', and the
references and version-specific records of the attached
papers.  No further paper resolving the general deformation
equivalence in \ref{push:src:E.2.4} was located.  This is a record of the
search performed, a record with the stated search scope;
the query responses and opened-source metadata are retained
in \texttt{evidence/E\_search*.json}.


\section{Comb invariants and the precise vector-field correction}\hypertarget{tag-034D}{}\label{sec:II:combs}
\begin{proposition}[Canonical bundle, exact tangent cohomology and Hodge growth on the combs]\hypertarget{tag-01SV}{}\label{inv:E2.2}
\StatementRecord{inv:E2.2}{proved}{01SV}{\SeriesRef{cdp:st:E15}}

For the same $X_m$, the root lift $\bar f_m:X_m\to\mathcal B$ has
vertical bundle $E_m$ and $G_m=R^1\bar f_{m*}\OO_{X_m}$ satisfying
\[
 0\longrightarrow\OO\longrightarrow E_m
       \longrightarrow\OO(24m-1)\longrightarrow0,
 \qquad
 0\longrightarrow\OO\longrightarrow G_m
       \longrightarrow\OO(24m-1)\longrightarrow0.      \tag{E.2}\hypertarget{tag-01TR}{}\label{invII:eq:E.2}
\]
The canonical bundle and all positive-degree holomorphic forms obey
\[
 K_{X_m}=\bar f_m^*\OO(-24m-6),\qquad
 h^{1,0}=h^{2,0}=h^{3,0}=0.                            \tag{E.3}\hypertarget{tag-01TS}{}\label{invII:eq:E.3}
\]
The tangent cohomology is that of Proposition~\ref{cdp:st:E15}.
The following proof gives a distinct calculation of those dimensions.
More precisely, writing $\mathcal L=\OO_{\mathcal B}(-5)$,
\[
 \begin{split}
 f_{m*}T_{X_m}&=\rho_*E_m,\\
 R^1f_{m*}T_{X_m}&=
  \operatorname{coker}\{\rho_*\mathcal L\xrightarrow{\delta_0}
                                      \rho_*(E_m\otimes G_m)\},\\
 R^2f_{m*}T_{X_m}&\simeq\OO_B(2m)\oplus Q_m,
 \qquad\operatorname{length}Q_m=4m,
 \end{split}                                        \tag{E.4a}\hypertarget{tag-01TT}{}\label{invII:eq:E.4a}
\]
where $Q_m=\operatorname{coker}\{\rho_*(\mathcal L\otimes G_m)
\xrightarrow{\delta_1}\rho_*(E_m\otimes\det G_m)\}$ has no support
at the cusp. The splitting in the last line need not be canonical.
There is a rank-three vector bundle $K_m$ on $B$ for which
\[
 \boxed{h^{1,1}(X_m)=4m+1+h^0(K_m),\qquad
 h^{1,2}(X_m)=4m+h^0(K_m).}                            \tag{E.5}\hypertarget{tag-01TU}{}\label{invII:eq:E.5}
\]
In particular both remaining middle Hodge numbers grow at least
linearly with $m$.
\end{proposition}


\begin{proof}

\begin{enumerate}[label={$\langle1\rangle$\arabic*.},leftmargin=*]
\item The construction retains the integral local system and the limiting
period planes. The invariant real homology class $\widehat\delta$
therefore gives a vertical translation field on every regular torus.
It is nowhere zero there, since the real period map is an isomorphism.
At a finite root point the same reasoning applies on the smooth
cover, and at the cusp its nonzero angular class gives a nonzero
logarithmic vertical vector. The local vertical transitions, which
depend only on the base exactly as in \SeriesRef{inv:D1.2}, define $E_m$ and give its
trivial subbundle. The invariant real cohomology class $\gamma$ gives
a nowhere-zero section of $G_m$, by the same smooth-fibre and cusp
comparison in \SeriesRef{inv:D1.2}.

\item The determinant of $G_m$ has the same two stabilizer characters as
that of $G_0$: invariant exponents $2$ at $p_1$ and $1$ at $p_2$.
Its coarse determinant has degree $2m-1$ by \CompendiumRef{II}{inv:E2.1}. The explicit descent
formula in \SeriesRef{inv:D1.1} therefore identifies
$\det G_m=\OO(12(2m-1)+8+3)=\OO(24m-1)$.
The rank-four flat bundle on the root base has trivial determinant:
the integral monodromies have determinant one, the finite-root charts
have no residual ramification, and the logarithmic cusp monodromy
has nilpotent residue of trace zero. Its flat volume section thus
extends without a zero or pole. The exact Hodge-plane sequence
$0\to E_m^*\to\mathcal H\to G_m\to0$ gives
$\det E_m=\det G_m$. Taking quotients by the two nonvanishing
sections proves \eqref{invII:eq:E.2}. The determinant of the log tangent sequence
$0\to\bar f_m^*E_m\to T_{X_m}(-\log W)
\to\bar f_m^*\OO(-5)\to0$ then gives \eqref{invII:eq:E.3}'s canonical formula.

\item The logarithmic Kodaira--Spencer map is invertible near the cusp.
This does not follow merely from fixing a limiting plane:
Viaclovsky's Proposition 6.1 gives the actual local period expression
$(\log q)I/(2\pi i)+H_m(q)$ with $H_m$ holomorphic. Its logarithmic
derivative has invertible leading matrix $I/(2\pi i)$, or
$B_0/(2\pi i)$ after the retained marking change. Thus the map is
generically invertible. The direct image of the logarithmic
cotangent sequence consequently has kernel exactly $\OO(5)$.
It has no section by \SeriesRef{inv:D1.1}, so $h^{1,0}=0$. To treat two-forms, take
the second exterior power of that sequence. Its possible global
sections lie in
$H^0(E_m^*\otimes\OO(5))$ or
$H^0((\det E_m)^{-1})$.
The latter is zero, having weight $1-24m<0$. A morphism
$E_m\to\OO(5)$ vanishes on the subline $\OO$, since
$H^0(\OO(5))=0$, and therefore factors through $\OO(24m-1)$.
The resulting section would have weight $6-24m<0$, so it too
vanishes. Absolute two-forms inject into logarithmic ones, which
proves $h^{2,0}=0$. The negative canonical weight proves $h^{3,0}=0$.

\item The first connecting map from the horizontal line to
$E_m\otimes G_m$ is nonzero near the cusp, so it is injective as a
sheaf map. The same normal quotient (D.10) has no degree-zero direct
image, giving $f_{m*}T_{X_m}=\rho_*E_m$.
Exact invariants applied to \eqref{invII:eq:E.2} give, for each of $V=E_m,G_m$,
\[
 0\longrightarrow\OO_B\longrightarrow\rho_*V
 \longrightarrow\OO_B(2m-1)\longrightarrow0.
\]
These bundles have vanishing $H^1$ and $h^0=2m+1$.
They also split into two nonnegative-degree lines: if a splitting
had one negative summand, the section of its saturated $\OO_B$
subline would have to lie in the positive summand. Every nonzero
section of that summand has a zero, contradicting saturation.

\item The natural map
\[
 (\rho_*E_m)\otimes(\rho_*G_m)
              \longrightarrow\rho_*(E_m\otimes G_m)
\]
is injective, with length-one cokernel at $p_1$ and no other
cokernel. Indeed both bundles have invariant exponents $(0,2)$ at
$p_1$. The product of their exponent-two generators is $s^4$,
whereas the product representation has invariant generator $s$;
the comparison on that one summand is multiplication by $t=s^3$.
All other exponents add without a carry. At $p_2$ the exponents
are $(0,1)$, and no sum reaches four. At the cusp there is no
stabilizer. The source is a sum of nonnegative-degree line bundles,
so its $H^1$ is zero. The length-one upper modification consequently
has $H^1=0$, rank four and degree
$2(2m-1)+2(2m-1)+1=8m-3$; it has $8m+1$ global sections.

\item In the direct-image sequence of the log tangent extension, the
second cup map
$\delta_1:\mathcal L\otimes G_m\to E_m\otimes\det G_m$
is generically invertible, by the invertible cusp Kodaira--Spencer
matrix. It is therefore injective as a sheaf map. The first cup map
$\delta_0$ is injective for the same reason. Consequently
$R^1\bar f_{m*}T(-\log W)=\operatorname{coker}\delta_0$.
The normal quotient $\nu_*K_P$ in (D.10) has zero direct images in
degrees zero and one, so the same formula holds for $R^1T$.
Since $\rho_*\mathcal L=\OO_B(-1)$ is acyclic, its cokernel has
$h^0=8m+1$ and $h^1=0$. Leray and $H^1(\rho_*E_m)=0$ now give
the exact value $h^1(T)=8m+1$.

\item For completeness the degree-two calculation gives Proposition~\ref{cdp:st:E15} directly,
with the full torsion correction retained. The coarse degrees of
$\mathcal L\otimes G_m$ and $E_m\otimes\det G_m$ are $2m-2$
and $6m-2$. Their invariant exponents at $p_1$ are respectively
$(1,0)$ and $(2,1)$, and at $p_2$ both are $(1,2)$.
Their root determinant weights are $24m-11$ and $72m-3$;
subtracting the character contributions and dividing by twelve
gives exactly these coarse degrees. Thus $Q_m$ has length $4m$.
It has no cusp contribution, where the cup map is invertible.
On the root base there is an exact sequence
\[
 0\longrightarrow\operatorname{coker}\delta_1
 \longrightarrow R^2\bar f_{m*}T(-\log W)
 \longrightarrow\mathcal L\otimes\det G_m
                 =\OO(24m-6)\longrightarrow0.
\]
It splits, since the extension group is the first cohomology of a
zero-dimensional sheaf. The cusp normal quotient then raises the
degree of the free line by $p_0$, exactly as in \SeriesRef{fam:st:A.5.2}. The local
fibre, its normalization and the logarithmic residue are unchanged;
in particular $h^2(T|_W)=1$ still excludes a split torsion summand at
the cusp. Hence the new free line is $\OO(24m+6)$, which descends to
$\OO_B(2m)$, proving \eqref{invII:eq:E.4a}. Leray gives
$h^2(T)=(2m+1)+4m=6m+1$ and $h^3(T)=0$.
As an independent check, these values satisfy the full tangent
index $\chi(T)=1$. The source family of Viaclovsky's Proposition 8.3
also gives
\[
 8m-8\le\dim\operatorname{Def}(X_m)\le8m+1;
\]
the lower bound concerns the actual classifying-family dimension,
while the upper bound is the tangent dimension now computed.

\item We give the bundle specifying the remaining Hodge number. Set
\[
 K_m=\ker\left\{\rho_*(E_m^*\otimes G_m)
  \xrightarrow{\ \kappa_1\ }
       \rho_*(\OO(5)\otimes\det G_m)\right\}.         \tag{E.6}\hypertarget{tag-01TV}{}\label{invII:eq:E.6}
\]
The map is cup product with the actual logarithmic
Kodaira--Spencer class, in the category of coherent sheaves on $B$.
It is generically onto, so the kernel is a rank-three vector bundle.
The preceding connecting map
\[
 \rho_*E_m^*\xrightarrow{\ \kappa_0\ }
                 \rho_*(\OO(5)\otimes G_m)
\]
is injective. The source and target have degrees $-2m-1$ and
$2m-1$. To check these degrees, use their fixed local character
exponents $(0,1)$ versus $(2,1)$ at the order-three point and
$(0,3)$ versus $(3,0)$ at the order-four point, together with the
root determinant degrees. The invariant determinant comparison in
\SeriesRef{inv:D1.1} gives precisely the displayed integers. Its cokernel $A_m$
therefore has length $4m$, irrespective of how its zeros are
distributed among regular and special fibres.

\item Writing $\mathcal A_m=\Omega^1_{X_m}(\log W)$, the direct-image
sequence gives
$0\to A_m\to R^1f_{m*}\mathcal A_m\to K_m\to0$.
The residue sequence
$0\to\Omega^1_{X_m}\to\mathcal A_m\to\nu_*\OO_P\to0$
has a surjective degree-zero residue map at $p_0$: the logarithmic
horizontal form $dq/q$ has residue one on every branch. Its kernel
is the one elementary modification
$f_{m*}\Omega^1=\OO_B(-2)$ of
$f_{m*}\mathcal A_m=\rho_*\OO(5)=\OO_B(-1)$.
The rationality of $P$ makes the higher direct images of the
residue quotient zero. Thus
$R^1f_{m*}\Omega^1=R^1f_{m*}\mathcal A_m$.
Since $H^1(A_m)=0$, its global sections have dimension
$4m+h^0(K_m)$, and Leray adds
$h^1(\OO_B(-2))=1$ in total degree one. This proves the formula
for $h^{1,1}$. Finally the sphere Chern numbers give
$\chi(\Omega^1)=-1$, and the vanishing just proved, together with
Serre duality, gives $h^{1,1}-h^{1,2}=1$. This proves \eqref{invII:eq:E.5}.

\end{enumerate}
\end{proof}



\Needspace{12\baselineskip}
\begin{proposition}[The kernel on an open family of twelve-tooth smoothings]\hypertarget{tag-03LZ}{}\label{inv:E2.5}
\StatementRecord{inv:E2.5}{proved}{03LZ}{\SeriesRef{inv:E2.2}, \SeriesRef{cdp:st:E15}, \SeriesRef{invII:eq:E.6}; \cite{Viaclovsky2026Combs}}
Use the original marked section of the comb construction on
$\mathcal B=\mathbb P(3,4)$, with coarse base $B=\mathbb P^1$, the
order-three and order-four points $p_1,p_2$, and the hexagonal cusp $p_0$.
Its bundles are $E_0=G_0=\mathcal O\oplus\mathcal O(-1)$ and
$\mathcal L=\mathcal O(-5)$. There is a nonempty open family of
admissible positive-comb manifolds $X_{12}$ obtained by embedded
smoothings of twelve deformed twistor conics attached at distinct
ordinary points, fixing the three marked period planes, for which the
actual kernel in \eqref{invII:eq:E.6} satisfies
\[
 K_{12}\simeq\mathcal O_B(-8)\oplus\mathcal O_B(-9)^{\oplus2},
 \qquad \operatorname{coker}\kappa_1=0.
\]
On this family,
\[
 h^0(K_{12})=0,\qquad h^1(K_{12})=23,\qquad
 h^{1,1}(X_{12})=49,\qquad h^{1,2}(X_{12})=48.
\]
\end{proposition}

\begin{proof}
\begin{enumerate}[label={$\langle1\rangle$\arabic*.},leftmargin=*]
\item 
\textbf{The actual seed kernel.}
Put $F_0=E_0^*$ and $D_0=\det G_0$. Write $x,y$ for the homogeneous
coordinates of weights $3,4$, respectively. Thus $y$ vanishes at the
order-three point. In row-$G$ and column-$F$ convention the actual
logarithmic period derivative has the form
\[
 T_0=\begin{pmatrix}0&ay\\by&cx\end{pmatrix},\qquad ab\ne0.
\]
Indeed its entries have weights $5,4,4,3$; the corresponding section
spaces are $0,\mathbb Cy,\mathbb Cy,\mathbb Cx$. Its invertible
logarithmic residue at the cusp forces $a,b$ to be nonzero.
For a tensor matrix $\eta$ in row-$F$ and column-$G$ convention define
\[
 \langle T,\eta\rangle
   =\sum_{i=1}^2(T_{2i}\eta_{i1}-T_{1i}\eta_{i2})
   =\operatorname{tr}(T\eta J),\qquad
 J=\begin{pmatrix}0&1\\-1&0\end{pmatrix}.
\]
This is the actual determinant-valued cup pairing. The seed row is
$(by,0,cx,-ay)$, with source weights $0,-1,1,0$ and target weight $4$.
Exact finite invariants give
\[
 \rho_*\mathcal O(4)=\mathcal O_B,
 \qquad \rho_*\mathcal O(1)=\rho_*\mathcal O(-1)=\mathcal O_B(-1).
\]
For the last two identities the invariant sheaves are lines with zero
$H^0$ and $H^1$, as follows from weighted homogeneous sections and
Serre duality with $K_{\mathcal B}=\mathcal O(-7)$.
The invariant section $y$ is the global coarse frame of the target.
Let $\sigma_x:\rho_*\mathcal O(1)\to\rho_*\mathcal O(4)$ be the
actual multiplication-by-$x$ map. The coarse row is
$(b,0,c\sigma_x,-a)$, which is onto. Its kernel isomorphism is
\[
 (u,v,w)\longmapsto
 (au/b-c\sigma_xv/b,\ w,\ v,\ u),
 \qquad \eta\longmapsto(\eta_{22},\eta_{21},\eta_{12}).
\]
Consequently $K_0=\mathcal O_B\oplus\mathcal O_B(-1)^2$.
This computation includes the possible $cx$ column. The determinant
$-ab y^2$ also proves invertibility of the actual handle derivative at
every ordinary attachment point.

\item \textbf{The root family and the flat relative kernel.}
Choose twelve distinct ordinary points. The conics chosen below lie
in the fixed-marking embedded smoothing construction of
\cite[Proposition~5.3]{Viaclovsky2026Combs}.
Here is the actual family lift at the finite points. In equivariant
frames $f_i,g_j$, the invariant generators are $w^{a_i^F}f_i$ and
$w^{a_j^G}g_j$. A coarse graph matrix $C(b,s)$ lifts by
\[
 Z_{ji}(w,s)=w^{a_j^G-a_i^F}C_{ji}(w^n,s).
\]
This follows by dividing its invariant graph vector by
$w^{a_i^F}$; multiplying back gives the inverse graph map.
At $p_1$ the order is $3$ and the exponent matrix is
$\left(\begin{smallmatrix}0&-1\\2&1\end{smallmatrix}\right)$.
At $p_2$ the order is $4$ and it is
$\left(\begin{smallmatrix}0&-3\\1&-2\end{smallmatrix}\right)$.
Since the entire period plane stays fixed along each marked section,
the required entries have the joint holomorphic divisions
$C_{12}=bD_{12}$, and also $C_{22}=bD_{22}$ at $p_2$.
Their lifts are respectively
\[
 w^2D_{12}(w^3,s),\qquad
 wD_{12}(w^4,s),\qquad w^2D_{22}(w^4,s).
\]
These are precisely the finite-character conditions of
\ref{cdp:st:E15}; all other exponents are nonnegative.
They construct the actual family Hodge and quotient bundles, with
their unchanged finite characters. At the cusp the chart in that
proposition is $(\log q)I/(2\pi i)+H(q,s)$ with $H$ jointly
holomorphic and bounded on a smaller parameter polydisc. Its residue
is unchanged. The same finite-cover actions, toric unit translations
and annular attachments give the compact family as in
\ref{cdp:st:E15}. In particular its coarse cusp cup row stays onto.

Choose a smoothing direction nonzero in every node-smoothing space.
Its regular total surface has coordinates $xy=s$ at each node,
with handle $y=0$ and tooth $x=0$. In an actual flat cohomology
frame write the period plane as the graph of $Z(x,y)$.
The relative dualizing generator $dx/x=-dy/y$ gives the derivative
\[
 \Theta=xZ_x-yZ_y.
\]
The relative cup takes values in
$D\otimes\omega_{\mathcal C/\Delta}(p_0)$, with the finite roots
along the marked sections understood. Put $H=Z_x(0,0)$ and
$B=Z_y(0,0)$. Our conic choices make the two functionals
$\langle H,-\rangle$ and $\langle B,-\rangle$ independent.
The cup-row entries then generate $(x,y)$ by their independent
linear terms and Nakayama's lemma. Invertible local operations
reduce the row to $(x,y,0,0)$, with free kernel generated by
$(-y,x,0,0)$ and the last two coordinate vectors.
Thus the relative kernel is locally free at every node.
Away from the nodes the central coarse row is onto; on the teeth
this follows from the conic calculation in Step 4. Its unit
coefficients persist after shrinking the family. Exact finite
invariants commute with parameter restriction in the character
frames just constructed. Finally $(x,y)$ is the unit ideal on a
noncentral fibre. We obtain a locally free coarse relative kernel
whose smooth restriction is the actual $K_{12}$ and whose smooth
coarse cokernel is zero.

\item \textbf{The complete node condition.}
Let $\eta_h(x),\eta_t(y)$ be the two branch sections, with common
ordinary value $\eta_0$. Their value equations are
\[
 \langle H,\eta_0\rangle=\langle B,\eta_0\rangle=0.
\]
Write $Z=Z_h(x)+Z_t(y)-Z(0,0)+xyR(x,y)$. The coefficient of $xy$
in $xZ_x-yZ_y$ is zero, including the arbitrary mixed graph term.
The mixed coefficient in the kernel equation is therefore
\[
 \langle H,\eta_t'(0)\rangle
       -\langle B,\eta_h'(0)\rangle=0.                 \tag{E.12.1}\hypertarget{tag-03M0}{}\label{invII:eq:E12-node}
\]
This is also sufficient to all analytic orders. Over
$R_0=\mathbb C\{x,y\}$ let $\varphi:R_0^4\to R_0$ be the actual
row, with image $(x,y)$. A lift of the branch pair has
$\varphi(\widetilde\eta)=xyh$, and $h(0)$ is exactly the left side
of \eqref{invII:eq:E12-node}. Adding $xy\zeta$ changes $h$ by
$\varphi(\zeta)$, whose values range over $(x,y)$. The displayed
vanishing therefore gives an actual analytic correction with
$\varphi(\zeta)=-h$. The corrected lift is a section of the
total kernel with the prescribed restrictions. This proves that
the one mixed-jet equation is the entire remaining condition.
The kernel of the central row exceeds the central fibre of the
flat kernel by just one node copy of $\mathbb C$.
Changing a determinant or dualizing generator adds multiples of
the two zero value equations. Tensor-frame changes transport both
derivatives with their actual frame derivatives retained.

\item \textbf{Conics and their actual two-jet map.}
In the flat decomposition $F_b\oplus G_b$ the through-point conics
under consideration have homogeneous inclusion and quotient
\[
 S_{[u:v]}=\binom{uI+vA}{vB},\qquad
 Q_{[u:v]}=(-vB,\ uI+vBAB^{-1}),\qquad \det B\ne0.
\]
Their product is zero. At $v\ne0$ the indicated bottom and left
blocks have full rank, and at $v=0$ the other blocks have full rank.
Thus they give the exact sequence with bundles
$F|_{C_i}=\mathcal O(-1)^2$ and $G|_{C_i}=\mathcal O(1)^2$.
The bottom-block chart recovers $u/v$ from the plane; moreover
$QS'=B$ near $y=v/u=0$. They are embedded immersed conics.
Set $P=I+yA$ and $C=BPB^{-1}$. Then
\[
 Z_t=yBP^{-1},\qquad S=(I;Z_t)P,
 \qquad Q=C(-Z_t,I),\qquad C Z_t'P=B.
\]
The tensor transition from the polynomial conic frames to graph
frames is $\eta_{\rm graph}=P\eta_{\rm conic}C^{-t}$.
The determinant factor is fixed by the exact identity
\[
 \langle CTP,P^{-1}\eta C^t\rangle
       =\det(C)\langle T,\eta\rangle,
\]
which follows from $C^tJC=(\det C)J$.
The ordinary tooth cup is the constant row $\langle B,-\rangle$
on $F\otimes G=\mathcal O^4$, so its kernel is $\mathcal O^3$.
Its global sections are constant in these frames. The derivative
forced by their value at the node is consequently
\[
 \eta_t'(0)=A\eta_0-\eta_0(BAB^{-1})^t.
\]
Substituting this exact frame derivative in \eqref{invII:eq:E12-node}
gives the two conditions on an ordinary handle-kernel section:
\[
 \begin{split}
 \langle B,\eta_h(0)\rangle&=0,\\
 \langle B,\eta_h'(0)\rangle
  &=\langle H,A\eta_h(0)-\eta_h(0)(BAB^{-1})^t\rangle.
 \end{split}                                        \tag{E.12.2}\hypertarget{tag-03M1}{}\label{invII:eq:E12-jets}
\]
The eight parameters $(A,B)$, modulo
$(A,B)\mapsto(rA+tI,rB)$ with $r\ne0$, are the actual
six-dimensional through-point conic chart: a polynomial basis
of its tautological $\mathcal O(-1)^2$ has constant blocks $I,0$
and linear blocks $A,B$. The actual source map is $y=r\widetilde y/(1+t\widetilde y)$.
Multiplying its tautological columns by $1+t\widetilde y$ gives
exactly the displayed matrix action. The two-jet kernel in
\eqref{invII:eq:E12-jets} is unchanged by these reparameterizations.
Lemma~5.1 of \cite{Viaclovsky2026Combs} gives the same local deformation space and
a nonempty complex-open set of admissible compact conics.

\item \textbf{Dominance onto the cyclic two-jet quotients.}
The perfect pairing $\langle-,-\rangle$ identifies
$B\mapsto\langle B,-\rangle|_{\ker\langle H,-\rangle}$ with the
quotient of a four-dimensional dual space by the line spanned by
$H$. It gives every first covector, subject to a dense open restriction.
For the second coefficient the following frame map is explicit.
Under $F'=FU$, $G'=GV$ one has
\[
 \begin{gathered}
 T'=V^{-1}TU,\qquad \eta'=U^{-1}\eta V^{-t},\\
 \langle T',\eta'\rangle=(\det V)^{-1}\langle T,\eta\rangle,
 \qquad A'=U^{-1}AU,\quad B'=V^{-1}BU.
 \end{gathered}
\]
Choose an eigenvector matrix $R$ for $H^{-1}B$ and put $U=R,V=HR$.
The actual derivative becomes $I$, the actual conic direction becomes
$\operatorname{diag}(\beta_1,\beta_2)$, and the tensor becomes
$R^{-1}\eta H^{-t}R^{-t}$. Work on the nonempty open set
\[
 \beta_1\beta_2(\beta_1-\beta_2)(\beta_1+\beta_2)\ne0.
\]
The two value equations then give
$\eta_0=\operatorname{diag}(u,v)$. For
$A=\left(\begin{smallmatrix}p&q\\r&s\end{smallmatrix}\right)$,
the second coefficient in these actual frames is
\[
 \langle I,A\eta_0-\eta_0(BAB^{-1})^t\rangle
   =r(1+\beta_2/\beta_1)u-q(1+\beta_1/\beta_2)v.
\]
Its Jacobian in $r,q$ has determinant
$-(\beta_1+\beta_2)^2/(\beta_1\beta_2)$, so the two coefficients
are arbitrary. The scalar part of $A$ contributes zero, as required
by the source reparameterization.
In any fixed local handle frame the resulting quotient is
\[
 \eta\longmapsto\ell(\eta(0))+
 t\bigl(\ell(\eta'(0))+\psi(\eta(0))\bigr)\pmod{t^2},
 \qquad \ell\ne0,
\]
where $\psi$ is arbitrary modulo $\ell$. The actual derivative of
the chosen handle frame translates $\psi$. This proves dominance
from actual conics to the four-dimensional cyclic length-two
quotient space, with its complete local frame map.

\item \textbf{An attained balanced open set.}
Let $K'$ be the subsheaf of $K_0$ defined by the twelve quotients
\eqref{invII:eq:E12-jets}. Their parameter space is irreducible.
There is an exact witness for its balanced locus: assign four of
the distinct nodes to each summand of
$K_0=\mathcal O\oplus\mathcal O(-1)^2$, and take the coordinate
quotient modulo the square of the local parameter at each assigned
point. Multiplication by each product of four squared linear
factors gives its actual kernel inclusion. The resulting bundle is
$\mathcal O(-8)\oplus\mathcal O(-9)^2$, so
\[
 H^0(K'(7))=0,\qquad H^1(K'(8))=0.                    \tag{E.12.3}\hypertarget{tag-03M2}{}\label{invII:eq:E12-open}
\]
These conditions define a nonempty Zariski-open subset by
semicontinuity. Step 5 pulls it back to a nonempty open subset of
actual conic data; no particular coordinate witness is asserted to
be an admissible conic.

The original handle normal bundle has degrees $(-1,0,0,1)$.
With the three fixed marks, positive modifications distributed as
$(4,3,3,2)$ give a globally generated modified normal bundle with
zero $H^1$. The complete normal-sheaf proof in
\cite[Proposition~5.3 and \S9.2]{Viaclovsky2026Combs}
shows that this is a second nonempty Zariski-open condition in the
actual tangent directions. The conic tangent map is dominant, so
both conditions pull back to dense open subsets of the same
irreducible conic chart. Their intersection meets the nonempty
complex-open set of admissible compact conics.
For these choices the cited fixed-marking deformation theorem
gives an embedded smoothing with a direction nonzero at every
node. The curve has genus zero and degree one over $B$, so each
smooth curve is an actual section. It fixes all three marked
period planes and has the root and compact filling of Step 2.

\item \textbf{Collapse, cohomology and the smooth bundle.}
Let $K_C$ denote the central fibre of the locally free relative
kernel. Its handle restriction maps into $K_0$ by the first
elementary modification at each node. On a tooth its image in the
ordinary $\mathcal O^3$ kernel is cut out by the additional nonzero
value functional $\langle H,-\rangle$, and hence
\[
 K_C|_{C_i}=\mathcal O(-1)\oplus\mathcal O^2.
\]
This follows by choosing that functional as the first coordinate:
the first coordinate vanishes at the node and the other two are
constant. Its $H^1$ vanishes, and its two global sections evaluate
injectively into the three-dimensional relative-kernel node fibre.
Their image retains exactly the derivative constraint of Step 3.

Collapse the teeth by the actual projection $\pi:C\to B$.
The normalization sequence for $K_C$, the zero tooth $H^1$ and the
surjective local evaluation of the handle onto each whole node
fibre give $R^1\pi_*K_C=0$. A handle section extends uniquely over
each tooth precisely by the two equations in
\eqref{invII:eq:E12-jets}, including the sufficient mixed condition.
Thus restriction gives the exact sheaf identification
$\pi_*K_C=K'$. For the actual pullback line from the base this gives
\[
 H^j(C,K_C\otimes\pi^*\mathcal O_B(k))
       =H^j(B,K'(k)),\qquad j=0,1.
\]
The projection of the embedded family supplies this single line
throughout the smoothing: it is trivial on the central teeth and
is the required $\mathcal O_B(k)$ on each smooth section.
Flatness and semicontinuity therefore transport
\eqref{invII:eq:E12-open} to all sufficiently small noncentral
members of the chosen smoothing.
The central Euler characteristic is
\[
 (1-12)+12\cdot2-12\cdot3=-23,
\]
so its rank-three degree is $-26$. Equivalently, the smooth
surjective cup row has source degree $-2$ and target degree $24$.

For completeness, these two vanishings determine the bundle by
explicit maps. Riemann--Roch gives $h^0(K_{12}(8))=1$. Its section
is nowhere zero, since a zero would give a section of $K_{12}(7)$.
The quotient of $K_{12}(8)$ by this copy of $\mathcal O$ has rank
two, degree $-2$ and zero $H^0,H^1$. Evaluation of its twist by
one is onto at every point; Riemann--Roch gives two source
dimensions, so that evaluation is an isomorphism. The quotient
is $\mathcal O(-1)^2$, and the extension by $\mathcal O$ splits
because $H^1(\mathcal O(1))^2=0$. Twisting back proves the stated
kernel isomorphism. The two vanishings and surjectivity are open
in the actual nearby fixed-marking section family, giving the
nonempty open family in the statement. Formula
\eqref{invII:eq:E.5} supplies its two Hodge numbers.

\end{enumerate}
\end{proof}

\paragraph{Relation to the original construction.}
The conic geometry, fixed-marking embedded smoothing and sufficient
twelve-tooth bound are those of
\cite[Lemma~5.1, Proposition~5.3 and \S9.2]{Viaclovsky2026Combs}.
The proof above computes the actual flat cup kernel and its two-jet
modifications, then transports their cohomology through that
smoothing. It determines the stated open twelve-tooth family in the
remaining-Hodge-number question of \cite[Section~10(2)]{Viaclovsky2026Combs}.
The general and special-member splitting problem remains in
\ref{inv:E2.3}.


\begin{proposition}[Counterexample: The universal two-vector-field bound]\hypertarget{tag-01SX}{}\label{inv:E2.4}
\StatementRecord{inv:E2.4}{proved}{01SX}{\SeriesRef{inv:E2.2}; \citet[Corollary~1.3, preprint p.~2]{HKP2000}; \citet[Remarks 1.2 and 9.26]{AC2026}}

The assertion that every complex structure on $S^6$ has
$h^0(T)\le2$, printed in
\citet[Corollary~1.3, preprint p.~2]{HKP2000}
and repeated in \citet[Remarks 1.2 and 9.26]{AC2026}, is false
for the comb members.
\end{proposition}


\begin{proof}
Proposition~\ref{inv:E2.2} gives $h^0(T_{X_m})=2m+1\ge25$ for $m\ge12$.
In the proof of \citet[Corollary~1.3, preprint p.~2]{HKP2000},
a third field is expressed as a meromorphic linear combination of two fields;
the coefficients are then declared constant using the
Campana--Demailly--Peternell assertion that every meromorphic function
on a complex six-sphere is constant. For the combs, the surjective algebraic reduction
$f_m$ induces an injection $f_m^*:\C(t)\hookrightarrow\mathcal M(X_m)$,
so its pullback of $t$ is nonconstant.
\end{proof}





\section{Uniform non-torsion on every full fibre}\hypertarget{tag-034H}{}\label{sec:II:uniform-nontorsion}
\begin{theorem}[Uniform non-torsion]
\hypertarget{tag-03DQ}{}\label{wb:thm:uniform-nontorsion-audit}
\hypertarget{tag-01TQ}{}\label{inv:wb:uniform-nontorsion-provenance}
\StatementRecord{wb:thm:uniform-nontorsion-audit}{proved}{03DQ}{\SeriesRef{inv:wb:eq:9.19}, \SeriesRef{inv:wb:eq:9.23}, \SeriesRef{inv:wb:eq:9.22}, \SeriesRef{inv:wb:eq:9.24}, \SeriesRef{inv:wb:eq:9.25}}

There is a line bundle \(L\in\operatorname{Pic}(X)\) such that
\(L|_{X_b}\) is non-torsion for every scheme-theoretic fibre
\(X_b=f^*(b)\), and \(L|_{S_j}\) is non-torsion for \(j=1,2\).
More precisely, identify \(\operatorname{Pic}(X)\) with \(\mathbb C\),
\(z\mapsto L_z\).  Every \(z\) outside the union of \(\mathbb R\) with a
countable set has these properties.
\end{theorem}



\begin{proof}

\begin{enumerate}[label=\arabic*.,leftmargin=*]
\item The pullback \(f^*\mathcal O_B(1)\) is nontrivial: by the projection
formula and \(f_*\mathcal O_X=\mathcal O_B\),
\[
 h^0(X,f^*\mathcal O_B(1))=h^0(B,\mathcal O_B(1))=2,
\]
whereas \(h^0(X,\mathcal O_X)=1\). Write \(M_u=\exp(ug)\) for the
additive Picard isomorphism of \SeriesRef{inv:C3.1}. There is a unique
\(u_0\ne0\) with \(M_{u_0}=f^*\mathcal O_B(1)\). The complex-linear
isomorphism
\[
 \phi:\mathbb C\longrightarrow\mathbb C,\qquad
 \phi(z)=u_0z/12,\qquad \phi^{-1}(u)=12u/u_0
\]
defines \(L_z=M_{\phi(z)}\). On a fibre \(F_q\), put
\(\iota_q(z)=(u_0z/12)g|_{F_q}\); its exponential is \(L_z|_{F_q}\).
Multiplication by the nonzero scalar \(u_0/12\) preserves the image
line \(\mathbb C\gamma^{0,1}\) of the restriction map. Thus
\[
 L_{12}=f^*\mathcal O_B(1).
\tag{WBAPII.9.19}\hypertarget{tag-01TJ}{}\label{inv:wb:eq:9.19}
\]

\item Fix a smooth fibre \(F_q\).  Put
\[
 \Lambda_q=\iota_q^{-1}\!\left(
 \operatorname{im}[H^1(F_q;\mathbb Z)\to
 H^1(F_q,\mathcal O_{F_q})]\right)\subset\mathbb C.
\tag{WBAPII.9.20}\hypertarget{tag-01TK}{}\label{inv:wb:eq:9.20}
\]
The integral image is a lattice in the underlying real vector space of
\(H^1(F_q,\mathcal O_{F_q})\), so \(\Lambda_q\) is discrete.  It contains
\(12\), because \eqref{inv:wb:eq:9.19} is trivial on every fibre.  The exponential
sequence on \(F_q\) shows
\[
 L_z|_{F_q}\text{ is torsion}
 \quad\Longleftrightarrow\quad z\in\mathbb Q\Lambda_q.
\tag{WBAPII.9.21}\hypertarget{tag-01TL}{}\label{inv:wb:eq:9.21}
\]
If \(\operatorname{rk}\Lambda_q=1\), its nonzero element \(12\) forces
\(\Lambda_q\subset\mathbb R\); all its bad parameters are real.

\item Suppose \(\operatorname{rk}\Lambda_q=2\).  By
\SeriesRef{inv:C3.1}, the image of \(\iota_q\) is the line
generated by the image of \(\gamma\).  Thus there is an integral class
\[
 \lambda=a\gamma+bu+cw+d\delta\notin\mathbb Q\gamma
\tag{WBAPII.9.22}\hypertarget{tag-01TM}{}\label{inv:wb:eq:9.22}
\]
whose \((0,1)\)-image lies on that line.  On the covering family, the four
marked lattice columns are
\[
 \pi_{\widehat\gamma}=(6\mu,\beta_c),\quad
 \pi_{\widehat u}=(\tau,\mu),\quad
 \pi_{\widehat w}=(1,0),\quad
 \pi_{\widehat\delta}=(0,1).
\tag{WBAPII.9.23}\hypertarget{tag-01TN}{}\label{inv:wb:eq:9.23}
\]
Every real cohomology class has a unique expression
\[
 2\operatorname{Re}(\bar\xi\,d\zeta_1+
                       \bar\upsilon\,d\zeta_2).
\]
Evaluating it on \eqref{inv:wb:eq:9.23} gives
\[
\begin{aligned}
2\operatorname{Re}\bar\xi&=c,&
2\operatorname{Re}\bar\upsilon&=d,\\
2\operatorname{Re}(\bar\xi\tau+\bar\upsilon\mu)&=b,&
2\operatorname{Re}(6\bar\xi\mu+\bar\upsilon\beta_c)&=a.
\end{aligned}
\tag{WBAPII.9.24}\hypertarget{tag-01TO}{}\label{inv:wb:eq:9.24}
\]
For \(\gamma\), the first two right sides vanish.  The third equation then
shows that its \((0,1)\)-line is represented, up to a nonzero scalar, by
\((\operatorname{Im}\mu,-\operatorname{Im}\tau)\); the fourth equation
only fixes the scalar.  Write
\(\bar\xi=c/2+ir\), \(\bar\upsilon=d/2+is\).  The condition that the class
in \eqref{inv:wb:eq:9.22} have \((0,1)\)-image on the \(\gamma\)-line is
\[
 \bar\xi\operatorname{Im}\tau+
 \bar\upsilon\operatorname{Im}\mu=0.
\]
Its real part, together with the third equation of \eqref{inv:wb:eq:9.24}, is equivalent
to the single holomorphic equation
\[
 c\tau+d\mu-b=0.
\tag{WBAPII.9.25}\hypertarget{tag-01TP}{}\label{inv:wb:eq:9.25}
\]
Neither \(a\) nor \(\beta_c\) occurs, as adding a multiple of \(\gamma\)
must not change the condition.

\item The functions \(1,\tau,\mu\) are linearly independent.  Applying the cusp
transformation \((\tau,\mu)\mapsto(\tau-1,\mu)\) to a constant linear
relation first kills the coefficient of \(\tau\).  If the remaining
relation made \(\mu\) constant, the order-four law
\(\mu\circ g_2=1+\mu/\tau\), with nonconstant \(\tau\), would force an
impossible identity.  Hence all coefficients vanish.

\item For every integral \(\lambda\notin\mathbb Q\gamma\), equation \eqref{inv:wb:eq:9.25} is
therefore a nonzero holomorphic function on the covering Riemann surface.
Its zero set is discrete and countable.  There are countably many
\(\lambda\), so their union, and its image in \(B^\circ\), are countable.
This image contains every \(q\) for which
\(\operatorname{rk}\Lambda_q=2\).  Each such fibre contributes the
countable bad set \(\mathbb Q\Lambda_q\); all rank-one fibres together
contribute only parameters in \(\mathbb R\).  Thus the smooth fibres
exclude only \(\mathbb R\) together with a countable set.

\item There remain the three scheme-theoretic special fibres.  By
\SeriesRef{inv:C3.1}, restriction
\[
 \mathbb C\longrightarrow H^1(X_b,\mathcal O_{X_b})
\]
is complex-linear and injective for each of them.  The exponential sequence
and countability of \(H^1(X_b;\mathbb Z)\) show that the parameters for
which \(L_z|_{X_b}\) is torsion form a countable set.

\item Finally consider the reductions \(S_j\).  The restriction
\(\varphi_j:\mathbb C\to H^1(S_j,\mathcal O_{S_j})\) is complex-linear.
It is nonzero.  Write \(\mathcal O_X(S_j)=L_{z_j}\).  Since
\[
 \mathcal O_X(m_jS_j)=f^*\mathcal O_B(1)=L_{12},
\]
the additive parameter gives \(z_j=12/m_j\), while
the root model of \SeriesRef{fam:src:A.1.2} identifies
\(\mathcal N_j:=\mathcal O_X(S_j)|_{S_j}\) with the normal line of
the reduced fibre. Its connected degree-\(m_j\) cover
\(\widehat S_j\to S_j\) has free deck generator \(a_j\). In the
root coordinate \(s_j\), the generator acts by
\(s_j\mapsto\zeta_j s_j\), with \(\zeta_j\) primitive of order
\(m_j\). The derivative therefore identifies the line with
\[
 \mathcal N_j\simeq
 (\widehat S_j\times\mathbb C)/
       ((x,v)\sim(a_jx,\zeta_jv)).
\]
A trivializing section of \(\mathcal N_j^{\otimes r}\) pulls back
to a nowhere-zero holomorphic function \(h\) on the compact connected
torus \(\widehat S_j\), with \(h(a_jx)=\zeta_j^r h(x)\).
The function is constant, so \(m_j\mid r\). Conversely the constant
section descends for \(r=m_j\); hence the order is exactly \(m_j\).
In particular,
\[
 L_{12/m_j}|_{S_j}=\mathcal N_j\ne\mathcal O_{S_j}.
\]  If \(\varphi_j\) were zero, the exponential sequence and
the vanishing of the restricted first Chern class would make every
\(L_z|_{S_j}\) trivial, a contradiction.  Thus \(\varphi_j\) is
injective, and the same exponential-sequence argument makes its torsion
parameters countable.  Removing these two additional countable sets proves
the theorem.

\end{enumerate}
\end{proof}






\Uses{the hypotheses, notation and definitions stated here; the argument and qualifications immediately following}



% Complete working-edition supplement. All source files remain unchanged.
% Reader and writer: GPT-6 Astra (OpenAI, ChatGPT Pro mode), effort Pro,
% 17 November 2026 UTC. Full source and label mapping accompanies this file.

\section{Analytic duality and the exact Picard restriction maps}
\begin{lemma}[Identification of the torsion subsheaf]
\hypertarget{tag-03IF}{}\label{wbII:lem:torsion-identification}
\StatementRecord{wbII:lem:torsion-identification}{proved}{03IF}{the definitions and constructions stated in the proof}

If $S$ is a reduced effective Cartier divisor in a smooth complex manifold,
then
\[
\ker\bigl(\Omega^1_S\longrightarrow j_*\Omega^1_{S_{\rm reg}}\bigr)
  =\operatorname{Tor}(\Omega^1_S).
\]
\end{lemma}



\begin{proof}
Every irreducible component of a reduced hypersurface is generically smooth,
so $S_{\rm reg}$ is dense and its complement is a proper analytic subset.
If a local section of $\Omega^1_S$ vanishes on $S_{\rm reg}$, its support avoids
all minimal primes of the reduced local ring.  Prime avoidance therefore gives
an element of its annihilator outside every minimal prime; that element is a
nonzerodivisor, so the section is torsion.  Conversely,
$\Omega^1_S|_{S_{\rm reg}}$ is locally free, hence has no torsion, and every
torsion section restricts to zero there.

\end{proof}



\subsubsection*{Local proof of the conormal filtration}


Let $S$ be a reduced effective Cartier divisor in a smooth complex
threefold $X$.  There are natural short exact sequences
\[
0\longrightarrow N^*_{S/X}\longrightarrow\mathscr K_S
\longrightarrow\mathscr T_S\longrightarrow0,
\]
\[
0\longrightarrow\mathscr K_S\longrightarrow\Omega_X^1|_S
\longrightarrow\widetilde\Omega_S^1\longrightarrow0.
\]



\begin{proof}
\begin{enumerate}[label={$\langle1\rangle$\arabic*.},leftmargin=*]
\item 
The conormal sequence is right exact.  It is also left exact here.  Locally
write $S=(g=0)$ with $g$ reduced.  If $h\,dg|_S=0$, then $h$ vanishes at the
generic smooth point of every irreducible component, because $dg$ is nonzero
there.  The ring $\mathcal O_S$ is reduced, so the intersection of its
minimal primes is zero; hence $h=0$.  Thus
\[
0\to N^*_{S/X}\to\Omega_X^1|_S\to\Omega_S^1\to0
\]
is exact.  By definition, $\mathscr K_S$ is the inverse image of
$\mathscr T_S$ under the last surjection.  Restriction to that inverse image
is onto $\mathscr T_S$ and has kernel $N^*_{S/X}$, proving the first
sequence.  Composing with the quotient by $\mathscr T_S$ proves the second.

\end{enumerate}
\end{proof}

\begin{proof}[The minimal-prime formulation of conormal injectivity]
\begin{enumerate}[label={$\langle1\rangle$\arabic*.},leftmargin=*]
\item 
Right exactness is the conormal sequence for a closed immersion.  We check
the left end.  Work at a point of $S$ and let $R$ be the regular local ring
of $X$, with reduced local equation $g$ for $S$.  The map at the left is
\[
 R/(g)\longrightarrow (\Omega^1_R/g\Omega^1_R),
 \qquad \bar h\longmapsto h\,dg.
 \tag{W.145}\hypertarget{tag-03DW}{}\label{wbII:eq:CN:2.2}
\]
Let $\mathfrak p$ be a minimal prime of $(g)$.  Since $R$ is regular and
has characteristic zero, the height-one local ring $R_{\mathfrak p}$ is a
discrete valuation ring.  Reducedness of $(g)$ gives
$\operatorname{ord}_{\mathfrak p}(g)=1$, so $g$ is a unit times a
uniformizer there; in characteristic zero the differential of that
uniformizer is nonzero modulo $\mathfrak p$.  Hence $dg$ is nonzero in
$\Omega^1_{R_{\mathfrak p}}\otimes k(\mathfrak p)$, so \eqref{wbII:eq:CN:2.2} sends $\bar h$
to zero only if $h\in\mathfrak p$.  This holds for every minimal prime.
Because $R/(g)$ is reduced, the intersection of its minimal primes is zero;
therefore $\bar h=0$.  Thus the conormal sequence in \cref{cdp:st:A.1.1} is exact.

\end{enumerate}
\end{proof}

\subsubsection*{Cohomology proof of the exact section criterion}


Let $A$ be a line bundle on $S$.  If
\[
H^0(S,\widetilde\Omega_S^1\otimes A)=0,
\]
then
\[
H^0(S,\Omega_X^1|_S\otimes A)
=H^0(S,\mathscr K_S\otimes A).
\]
If also $H^0(S,N^*_{S/X}\otimes A)=0$, then this group is canonically
\[
\ker\!\left[
H^0(S,\mathscr T_S\otimes A)
\xrightarrow{\delta_A}H^1(S,N^*_{S/X}\otimes A)
\right].
\]
Hence the two vanishings printed in the p.~692 reduction imply the desired
ambient vanishing if and only if $\delta_A$ is injective.



\begin{proof}
\begin{enumerate}[label={$\langle1\rangle$\arabic*.},leftmargin=*]
\item 
Tensoring the second sequence of
\cref{cdp:st:A.1.1} by the locally free sheaf $A$
preserves exactness; its global-section sequence gives the first equality.
Tensor the first sequence and take global sections.  If the conormal
$H^0$ vanishes, exactness identifies the middle group with the displayed
kernel.  No condition on the torsion-free quotient controls that kernel.

\end{enumerate}
\end{proof}

\subsection{Local algebra of a reduced normal-crossings fibre}

Let $g:\mathcal X\to\Delta$ be a holomorphic map from a smooth complex
threefold to a disc, and put $S=g^{-1}(0)$, assumed nonempty.  Assume $S$ is reduced and has
normal crossings: at every point there are coordinates $(x,y,z)$ in which
\[
g=ux,\qquad g=uxy,\qquad\text{or}\qquad g=uxyz
\tag{W.3}\hypertarget{tag-03HS}{}\label{wbII:eq:W:5.3}
\]
for a unit $u$.  Let $\eta:\widetilde S\to S$ be the normalization,
$D=S_{\mathrm{sing,red}}$, and $\widetilde D=\eta^{-1}(D)_{\mathrm{red}}$.

\subsubsection*{Coordinate proof for normal-crossing fibres}

\ In the setting above, $N^*_{S/\mathcal X}\simeq\mathcal O_S$,
and its generator maps to $dg|_S$.  Multiplication by $dg|_S$ is injective,
so
\[
0\longrightarrow\mathcal O_S
\xrightarrow{\,h\mapsto h\,dg\,}
\Omega^1_{\mathcal X}|_S\longrightarrow\Omega^1_S\longrightarrow0
\tag{W.4}\hypertarget{tag-03HX}{}\label{wbII:eq:W:5.4}
\]
is exact.



\begin{proof}
\begin{enumerate}[label={$\langle1\rangle$\arabic*.},leftmargin=*]
\item 
The fibre ideal is the principal ideal $(g)$ and is reduced, so
$(g)/(g^2)\simeq\mathcal O_S$.  Multiplication by a unit does not affect
injectivity, and the term $f\,du$ vanishes after restriction to $f=0$.
It therefore suffices to use $f=x,xy,xyz$.

For $f=xy$, if a lift $\widetilde h$ satisfies
$\widetilde h(y\,dx+x\,dy)=0$ modulo $(xy)$, then
$\widetilde h y,\widetilde h x\in(xy)$.  Hence
$\widetilde h\in(x)\cap(y)=(xy)$ in the regular local ring, so $h=0$ on
$S$.  For $f=xyz$, the three coefficient conditions say
$\widetilde h yz,\widetilde h xz,\widetilde h xy\in(xyz)$; hence
$\widetilde h\in(x)\cap(y)\cap(z)=(xyz)$.  The case $f=x$ is immediate.
Thus the conormal map is injective in every local model, proving \eqref{wbII:eq:W:5.4}.

\end{enumerate}
\end{proof}

\subsubsection*{Coordinate proof of the normalization--conductor square}


\ The sequence
\[
0\longrightarrow\mathcal O_S
\xrightarrow{\alpha}
\eta_*\mathcal O_{\widetilde S}\oplus\mathcal O_D
\xrightarrow{\beta}
\eta_*\mathcal O_{\widetilde D}
\longrightarrow0,
\tag{W.5}\hypertarget{tag-03HY}{}\label{wbII:eq:W:5.5}
\]
\[
\alpha(h)=(\eta^*h,h|_D),\qquad
\beta(\widetilde h,k)=
\widetilde h|_{\widetilde D}-\eta^*k,
\]
is exact.



\begin{proof}
\begin{enumerate}[label={$\langle1\rangle$\arabic*.},leftmargin=*]
\item 
The assertion is local.  Away from $D$, $\eta$ is an isomorphism.
At a double point $S=\{xy=0\}$, a member of $\ker\beta$ is
\[
\bigl((h_1(y,z),h_2(x,z)),k(z)\bigr),
\qquad
h_1(0,z)=h_2(0,z)=k(z).
\]
The function
\[
h=h_1(y,z)+h_2(x,z)-h_1(0,z)
\]
on $\{xy=0\}$ restricts to $h_1,h_2$ on the two branches and to $k$ on
their intersection.  This proves middle exactness in the double model.

At a triple point $S=\{xyz=0\}$, write the branch functions as
$h_1(y,z)$, $h_2(x,z)$, $h_3(x,y)$.  Membership in $\ker\beta$ is exactly
\[
h_1(0,z)=h_2(0,z),\quad
h_1(y,0)=h_3(0,y),\quad
h_2(x,0)=h_3(x,0).
\]
Then
\[
\begin{split}
h={}&h_1(y,z)+h_2(x,z)+h_3(x,y)\\
&-h_1(0,z)-h_1(y,0)-h_2(x,0)+h_1(0,0)
\end{split}
\]
restricts to the prescribed function on every branch.  Injectivity of
$\alpha$ follows because $S$ is reduced and $\eta$ is an isomorphism over a
dense open set.  Surjectivity of $\beta$ follows locally by lifting a
function on the branch conductor to the smooth branch and taking $k=0$.
This proves exactness in every local model.

\end{enumerate}
\end{proof}

\begin{theorem}[Conductor-supported ambient differential]
\hypertarget{tag-03IK}{}\label{wbII:prop:conductor}
\StatementRecord{wbII:prop:conductor}{proved}{03IK}{\SeriesRef{wbII:eq:W:5.3}, \SeriesRef{wbII:eq:W:5.5}, \SeriesRef{wbII:eq:W:5.4}}
\label{wbII:prop:conductor-differential}
In the setting \eqref{wbII:eq:W:5.3}, let $A$ be a line bundle on $S$ with
$\eta^*A\simeq\mathcal O_{\widetilde S}$, and choose a nowhere-vanishing
section $e$ of $\eta^*A$.  There is a unique nonzero
\[
s\in H^0(S,\Omega^1_{\mathcal X}|_S\otimes A)
\tag{W.6}\hypertarget{tag-03HZ}{}\label{wbII:eq:W:5.6}
\]
such that
\[
\eta^*s=\eta^*(dg|_S)\otimes e,\qquad s|_D=0.
\tag{W.7}\hypertarget{tag-03I0}{}\label{wbII:eq:W:5.7}
\]
Its image $\rho(s)$ in $\Omega^1_S\otimes A$ is supported on $D$ and
therefore torsion; its image in
$\widetilde\Omega^1_S\otimes A$ is zero.  Moreover:
\begin{enumerate}[label=(\alph*)]
\item $\rho(s)=0$ if and only if $e=\eta^*\sigma$ for a nowhere-vanishing
$\sigma\in H^0(S,A)$;
\item if $S$ is compact and connected and $A\not\simeq\mathcal O_S$, then
$H^0(S,A)=0$ and $s$ does not come from the conormal term in \eqref{wbII:eq:W:5.4}.
\end{enumerate}
\end{theorem}



\begin{proof}

\begin{enumerate}[label=\arabic*.,leftmargin=*]
\item In the double model, $d(xy)=y\,dx+x\,dy$ has coefficients in the reduced
ideal $(x,y)$ of $D$.  In the triple model,
\[
d(xyz)=yz\,dx+xz\,dy+xy\,dz
\]
has coefficients in $(xy,xz,yz)$, the reduced ideal of the union of the
three double curves.  Unit factors in \eqref{wbII:eq:W:5.3} do not change this after
restriction to $S$.  Hence $\eta^*(dg|_S)$ restricts to zero on
$\widetilde D$.  Tensor \eqref{wbII:eq:W:5.5} by the locally free sheaf
$E=\Omega^1_{\mathcal X}|_S\otimes A$.  The compatible pair
\[
\bigl(\eta^*(dg|_S)\otimes e,0\bigr)
\]
therefore descends uniquely to $s$.  On the smooth locus of $S$, $dg|_S$
is a nonzero conormal vector and $e$ is nowhere zero, so $s\ne0$.

\item On $S\setminus D$, the normalization is an isomorphism and
$s=dg|_S\otimes e$.  By \eqref{wbII:eq:W:5.4} its intrinsic image vanishes there.  Thus
$\rho(s)$ is supported on $D$ and lies in the torsion subsheaf; its image
modulo torsion is zero.  Here, for $j:S_{\rm reg}\hookrightarrow S$,
$\widetilde\Omega^1_S=\Omega^1_S/\operatorname{Tor}(\Omega^1_S)$
is identified with the image in $j_*\Omega^1_{S_{\rm reg}}$ by
$[\omega]\mapsto\omega|_{S_{\rm reg}}$; its inverse sends a local
restriction to the class of any lift.  The equality of the restriction
kernel and the torsion subsheaf is
\cref{wbII:lem:torsion-identification}.

\item If $\rho(s)=0$, exactness of \eqref{wbII:eq:W:5.4} after tensoring by $A$ gives
$s=dg|_S\otimes\sigma$ for a section $\sigma$ of $A$.  Pulling back and
comparing with \eqref{wbII:eq:W:5.7} on the dense open set where $dg\ne0$ gives
$\eta^*\sigma=e$ there, hence everywhere by holomorphic continuation.
Since $e$ never vanishes and $\eta$ is surjective, $\sigma$ never vanishes
and trivializes $A$.  The converse is immediate.

\item Finally suppose $S$ is compact and connected and $\sigma\in H^0(S,A)$.  On each compact
connected component of $\widetilde S$, write
$\eta^*\sigma=c_i e$; the holomorphic function $c_i$ is constant.  If one
$c_i$ is zero and another is nonzero, connectedness of the finite union of
irreducible components of $S$ provides two such components meeting at a
point.  The value of $\sigma$ at that point would be both zero and nonzero,
a contradiction.  Hence either every $c_i=0$, so $\sigma=0$, or every
$c_i\ne0$, so $\sigma$ is nowhere zero and $A$ is trivial.  When $A$ is
nontrivial only the first case occurs.  Since
$N^*_{S/\mathcal X}\simeq\mathcal O_S$, the conormal global-section space is
$H^0(S,A)=0$, while $s\ne0$.

\end{enumerate}
\end{proof}



\begin{corollary}[Nonzero cotangent connecting kernel]
\hypertarget{tag-03DU}{}\label{wbII:cor:normal-crossing-kernel}
\StatementRecord{wbII:cor:normal-crossing-kernel}{proved}{03DU}{\SeriesRef{wbII:prop:conductor}, \SeriesRef{cdp:st:A.1.1}}

If, in addition, $H^0(S,N^*_{S/X}\otimes A)=0$ and
$H^0(S,\widetilde\Omega^1_S\otimes A)=0$, then the section in
\cref{wbII:prop:conductor} proves that $\delta_A$ in
the section-kernel isomorphism of \cref{cdp:st:A.1.1} is not injective.
\end{corollary}



\begin{proof}
The section is nonzero and maps to zero in the quotient
$\widetilde\Omega^1_S=\Omega^1_S/\mathscr T_S$, so it
lies in the left-hand side of the section-kernel isomorphism of \cref{cdp:st:A.1.1}.
Under that isomorphism it is sent to its intrinsic image $\rho(s)$.
The hypothesis $H^0(S,N^*_{S/X}\otimes A)=0$ makes this map injective,
and the ambient lift $s$ gives $\delta_A(\rho(s))=0$.  \cref{cdp:st:A.1.1}
identifies that group with $\ker\delta_A$, which is therefore nonzero.

\end{proof}



\subsection{A proper counterexample to the printed inference}

\subsubsection*{Conductor proof of the nodal endpoint counterexample}

\ The p.~692 implication described after \CompendiumRef{II}{wbII:eq:W:5.2} is false as a formal
sheaf-theoretic statement, even for a proper family with smooth total space
and an irreducible reduced normal-crossings central fibre.



\begin{proof}
\begin{enumerate}[label={$\langle1\rangle$\arabic*.},leftmargin=*]
\item 
Let $\Delta$ be a sufficiently small disc with coordinate $t$ and put
\[
\mathcal C=
\{([x:y:z],t)\in\mathbb P^2\times\Delta:
y^2z=x^2(x+z)+tz^3\}.
\tag{W.8}\hypertarget{tag-03I1}{}\label{wbII:eq:W:5.8}
\]
This is proper over $\Delta$.  Its total space is smooth.  On $z\ne0$, the
partial derivative of the defining equation with respect to $t$ is
$-z^3\ne0$.  On $z=0$, the equation forces $x=0$; then $y\ne0$ and the
$z$-derivative is $y^2\ne0$.  The central fibre $C_0$ is the irreducible
nodal cubic $y^2z=x^2(x+z)$, with node $[0:0:1]$.

Set
\[
\mathcal X=\mathcal C\times\mathbb P^1,\qquad
S=C_0\times\mathbb P^1,
\]
and let $g:\mathcal X\to\Delta$ be projection.  Then $\mathcal X$ is a
smooth threefold, $g$ is proper, and $S$ is an irreducible reduced
normal-crossings surface.  Its normalization is
$\widetilde S\simeq\mathbb P^1\times\mathbb P^1$.

Let $p,q$ be the two inverse images of the node in the normalization of
$C_0$.  Fix $\lambda\in\mathbb C^*$ which is not a root of unity (for
example, $\lambda=2$).  Glue the fibres of
the trivial line bundle on the normalization at $p$ and $q$ by the scalar
$\lambda$.  This defines a line bundle
$A_\lambda\in\operatorname{Pic}^0(C_0)$ whose pullback is trivial.  It has
infinite order: its $n$th tensor power has gluing scalar $\lambda^n$, which
is never one.  It is not trivial: a global section pulls back to a constant
$c$ satisfying
$c=\lambda c$, so
\[
H^0(C_0,A_\lambda)=0.
\tag{W.9}\hypertarget{tag-03I2}{}\label{wbII:eq:W:5.9}
\]
Put $A=\operatorname{pr}_1^*A_\lambda$ on $S$.  Then
$\eta^*A\simeq\mathcal O_{\widetilde S}$ and $H^0(S,A)=0$.
Because $S$ is the full fibre,
$N^*_{S/\mathcal X}\simeq\mathcal O_S$, so the conormal-twist global
sections vanish.

The natural map
$\widetilde\Omega^1_S\to\eta_*\Omega^1_{\widetilde S}$ is injective:
its kernel is supported where normalization is not an isomorphism, while
the source is torsion-free by definition.  Projection formula and
$\eta^*A\simeq\mathcal O_{\widetilde S}$ therefore give
\[
H^0(S,\widetilde\Omega^1_S\otimes A)
\hookrightarrow
H^0(\widetilde S,\Omega^1_{\widetilde S})=0,
\tag{W.10}\hypertarget{tag-03EA}{}\label{wbII:eq:W:5.10}
\]
the last equality following from
$\widetilde S=\mathbb P^1\times\mathbb P^1$.
Nevertheless \cref{wbII:prop:conductor} produces
\[
0\ne s\in H^0(S,\Omega^1_{\mathcal X}|_S\otimes A).
\]
Thus both vanishings invoked in the printed reduction hold while the desired
ambient vanishing fails.  This calculation chooses $A$ directly on $S$ and disproves the
claimed \emph{formal implication between the three sheaf-cohomology groups}
for a non-torsion line bundle. The explicit extending choice of $A$ is
constructed in the next subsection, which also treats the total-space
line-bundle condition on p.~692.  The family \eqref{wbII:eq:W:5.8} also does not satisfy all hypotheses of CDP
Theorem~2.2, so this is not a counterexample to that theorem.

\end{enumerate}
\end{proof}

\subsection{The explicit extending choice in the nodal family}
Continue with
\[
\mathcal C=\{y^2z=x^2(x+z)+tz^3\}\subset\PP^2\times\Delta,
\qquad \mathcal X=\mathcal C\times\PP^1,
\qquad S=C_0\times\PP^1.
\]
Write $n=[0:0:1]$, let $D=\{n\}\times\PP^1$, and let
$\nu:\PP^1\to C_0$ denote the normalization described next.
\subsubsection*{Explicit normalization and crossing coordinates}

The projection $g$ is proper, $\mathcal X$ is a smooth complex threefold,
  and $S$ is a reduced irreducible normal-crossing (self-crossing) surface whose normalization
is $\PP^1\times\PP^1$.  Its singular locus is
\(D=\{[0:0:1]\}\times\PP^1\), and locally along $D$ the pair $(\mathcal X,S)$
has the model $t=uv$; away from $D$ the fibre is smooth and has the usual
local model $(\C^3,\{u=0\})$.


\begin{proof}
\begin{enumerate}[label={$\langle1\rangle$\arabic*.},leftmargin=*]
\item 
Properness follows because $\mathcal C$ is closed in the projective
projection $\PP^2\times\Delta\to\Delta$.  Write the defining equation as
\[
 F=y^2z-x^2(x+z)-tz^3.
\]
On the chart $z\ne0$, $\partial F/\partial t=-z^3\ne0$, so $\mathcal C$
is smooth there.  On $z=0$, the equation forces $x=0$ and $y\ne0$; at such
points $\partial F/\partial z=y^2\ne0$.  Hence $\mathcal C$, and therefore
$\mathcal X$, is smooth.  The only finite discriminant values of the fibres
are $t=0$ and $t=-4/27$; after shrinking $\Delta$ to $|t|<4/27$, every
noncentral fibre is smooth.

The central curve is
\(C_0:\ y^2z=x^2(x+z)\).  Its only singular point is
\(n=[0:0:1]\): on $z=1$ the equations $F=F_x=F_y=0$ give $(x,y)=(0,0)$,
and the point at infinity $[0:1:0]$ is smooth because $F_z=1$ there
(Euler's identity also gives the remaining projective singularity equation).
The parametrization
\[
 \nu:\PP^1\longrightarrow C_0,\qquad
 [u:v]\longmapsto
 [(u^2-v^2)v:\,u(u^2-v^2):\,v^3]
 \tag{W.160}\hypertarget{tag-03DX}{}\label{wbII:eq:CN:4.3}
\]
maps the two points $p=[1:1]$, $q=[-1:1]$ to $n$.  On the chart $v\ne0$, set
$r=u/v$; after scaling, its image is
$[r^2-1:r(r^2-1):1]$.  Conversely, for an affine point of $C_0$ with
$x\ne0$, the equation gives $r=y/x$ and $x=r^2-1$, so it is hit by
$[r:1]$.  The only affine point with $x=0$ is the node $(0,0)$, hit by
$[1:1]$ and $[-1:1]$, and the point at infinity $[0:1:0]$ is hit by
$[1:0]$.  Thus $\nu(\PP^1)=C_0$; as a nonconstant morphism to a projective
curve it is finite and generically one-to-one, hence the normalization.  In
particular $C_0$ is irreducible and rational.  The product normalization
is $\eta=\nu\times\id_{\PP^1}:\PP^1\times\PP^1\to S$.

At $n$ the affine equation is
\(t=y^2-x^2-x^3\).  On a sufficiently small neighbourhood choose the
holomorphic square root $h(x)=\sqrt{1+x}$ and set
\(u=y-xh(x)\), \(v=y+xh(x)\).  Then $uv=t$, and the map
$(x,y)\mapsto(u,v)$ has invertible $2\times2$ Jacobian at $(0,0)$, so $u,v$
are local coordinates on the surface $\mathcal C$.  If $s$ is a local
coordinate on $\PP^1$, then $(u,v,s)$ are coordinates on $\mathcal X$ and
$t=uv$.  Hence the pair has the asserted normal-crossing (self-crossing)
local model, with singular locus $D$ as stated.

\end{enumerate}
\end{proof}

\begin{lemma}[A line bundle with non-torsion gluing]
\hypertarget{tag-03IB}{}\label{wbII:lem:gluing-bundle}
\StatementRecord{wbII:lem:gluing-bundle}{proved}{03IB}{the definitions and constructions stated in the proof}

Fix $\lambda\in\C^*$ that is not a root of unity.  Gluing the trivial line
bundle on $\PP^1$ at $p$ and $q$ by the scalar $\lambda$ defines a line bundle
\(A_\lambda\in\Pic(C_0)\) with
\[
 \nu^*A_\lambda\simeq\cO_{\PP^1},\qquad
 H^0(C_0,A_\lambda)=0,\qquad
 A_\lambda^{\otimes m}\not\simeq\cO_{C_0}\quad(m\ge1).
 \tag{W.161}\hypertarget{tag-03DY}{}\label{wbII:eq:CN:4.4}
\]
\end{lemma}



\begin{proof}
For a nodal curve, a line bundle with trivial pullback to the normalization
is obtained by identifying the two fibres over the branches of the node.
Use the identification whose scalar is $\lambda$.  A global section pulls
back to a constant $c$ on the projective line and must satisfy
$c=\lambda c$; since $\lambda\ne1$, $c=0$.  The $m$th tensor power in \eqref{wbII:eq:CN:4.4} has
gluing scalar $\lambda^m$, which is never $1$ because $\lambda$ is not a
root of unity.

\end{proof}



\begin{lemma}[An extending gluing bundle]
\hypertarget{tag-03I9}{}\label{wbII:lem:extending-gluing}
\StatementRecord{wbII:lem:extending-gluing}{proved}{03I9}{\SeriesRef{cdp:st:B1}, \SeriesRef{wbII:lem:gluing-bundle}}

After shrinking $\Delta$, the sections
\[
 O(t)=[0:1:0],\qquad
 P(t)=[3:\sqrt{36+t}:1]
\tag{W.162}\hypertarget{tag-03E0}{}\label{wbII:eq:CN:4.5a}
\]
of $\mathcal C\to\Delta$ are disjoint and holomorphic.  If
$B=\cO_{\mathcal C}(P-O)$, then the line bundle
$\widetilde A=\operatorname{pr}_1^*B$ on $\mathcal X=\mathcal C\times\PP^1$
restricts to $\operatorname{pr}_1^*A_{\lambda}$ on $S$ for
$\lambda=1/3$ (or its inverse, according to the gluing convention).  The bundle in
\cref{cdp:st:B1} with this value of $\lambda$ therefore extends to
$\mathcal X$.
\end{lemma}



\begin{proof}
The equation of $\mathcal C$ is satisfied by both sections; the square root
is holomorphic for $|t|$ small, and the two sections are disjoint because
$O$ has $z=0$ whereas $P$ has $z=1$.  On the normalization of $C_0$, write
$r=u/v$.  Then $O(0)$ corresponds to $r=\infty$, $P(0)$ to $r=2$, and the
two branches over the node correspond to $r=1$ and $r=-1$.  Hence
\[
 \nu^*(B|_{C_0})\simeq\cO_{\PP^1}([2]-[\infty]).
\]
The rational function $f(r)=r-2$ has precisely this divisor and trivializes
the pullback.  Its values at the two branches are $f(1)=-1$ and
$f(-1)=-3$, so the induced gluing multiplier is $f(1)/f(-1)=1/3$,
up to inversion.  This is not a root of unity, so its gluing bundle is non-torsion by
Lemma~\ref{wbII:lem:gluing-bundle}.

\end{proof}



\subsubsection*{The local torsion calculation in the nodal example}

Let $A=\operatorname{pr}_1^*A_\lambda$ on
$S=C_0\times\PP^1$, with $\lambda$ as in Lemma~\ref{wbII:lem:gluing-bundle}.
Then
\[
 H^0(S,N^*_{S/\mathcal X}\otimes A)=0,
 \qquad
 H^0(S,\widetilde\Omega^1_S\otimes A)=0,
 \tag{W.163}\hypertarget{tag-03DZ}{}\label{wbII:eq:CN:4.5}
\]
but
\[
 H^0(S,\Omega^1_{\mathcal X}|_S\otimes A)\ne0.
 \tag{W.164}\hypertarget{tag-03E1}{}\label{wbII:eq:CN:4.6}
\]
In particular, the two endpoint vanishings in \eqref{wbII:eq:CN:4.5} do not imply the
vanishing of \eqref{wbII:eq:CN:4.6}, even for a proper morphism over a disc with smooth total
space and an irreducible reduced normal-crossing (self-crossing) central
fibre.


\begin{proof}
\begin{enumerate}[label={$\langle1\rangle$\arabic*.},leftmargin=*]
\item 
Because $S$ is the principal fibre $(t=0)$, its conormal line is trivial:
\(N^*_{S/\mathcal X}\simeq\cO_S\).  The projection
\(\operatorname{pr}_1:S\to C_0\) has
\(\operatorname{pr}_{1*}\cO_S\simeq\cO_{C_0}\), so \eqref{wbII:eq:CN:4.4} gives
\[
 H^0(S,A)\simeq H^0(C_0,A_\lambda)=0.
 \tag{W.165}\hypertarget{tag-03E2}{}\label{wbII:eq:CN:4.7}
\]
This proves the first vanishing in \eqref{wbII:eq:CN:4.5}.

We next check the second.  At a normal-crossing (self-crossing) point with local equation
\(uv=0\), the pullback map
\(\Omega^1_S\to\eta_*\Omega^1_{\widetilde S}\) has kernel exactly
\(\mathscr T_S\).  Indeed, in
\(R=\C\{u,v,w\}/(uv)\) the relation is
\(v\,du+u\,dv=0\).  If
\(a\,du+b\,dv+c\,dw\) restricts to zero on both branches, then
\(a\in(v)\), \(b\in(u)\), and
\(c\in(u)\cap(v)=(uv)=0\).  Modulo the relation the class is therefore a
multiple of $\tau=v\,du=-u\,dv$.  This class is nonzero: if
\(v\,du=h(v\,du+u\,dv)\), then coefficient comparison gives
  \((1-h)v=0\) and $hu=0$.  Since
\(\operatorname{Ann}_R(v)=(u)\) and \(\operatorname{Ann}_R(u)=(v)\),
\(1-h\in(u)\) and \(h\in(v)\), which would imply
\(1=(1-h)+h\in(u)+(v)\), impossible because $(u)+(v)$ is proper.
Moreover $(u+v)\tau=0$.  Since $u+v$ is a non-zero-divisor in
$R$ (the associated primes are $(u)$ and $(v)$), $\tau$ is genuine
torsion.
Conversely, each summand of the normalization module is a module over one of
the domains $R/(u)$ and $R/(v)$.  A non-zero-divisor of the reduced ring $R$
avoids both minimal primes and therefore acts injectively on both summands;
the target is torsion-free.  Thus the kernel of the normalization map is
exactly the torsion subsheaf, and there is an injection
\[
 \widetilde\Omega^1_S\hookrightarrow\eta_*\Omega^1_{\PP^1\times\PP^1}.
 \tag{W.166}\hypertarget{tag-03E3}{}\label{wbII:eq:CN:4.8}
\]
Since $\eta^*A\simeq\cO_{\PP^1\times\PP^1}$, the projection formula and \eqref{wbII:eq:CN:4.8}
give
\[
 H^0(S,\widetilde\Omega^1_S\otimes A)
 \hookrightarrow
 H^0(\PP^1\times\PP^1,\Omega^1_{\PP^1\times\PP^1})=0.
 \tag{W.167}\hypertarget{tag-03E4}{}\label{wbII:eq:CN:4.9}
\]
The last equality follows from
\(\Omega^1_{\PP^1\times\PP^1}\simeq
\operatorname{pr}_1^*\Omega^1_{\PP^1}\oplus
\operatorname{pr}_2^*\Omega^1_{\PP^1}\) and
\(H^0(\PP^1,\Omega^1_{\PP^1})=0\).

Finally, along the conductor $D$ the pullback of $dg|_S$ is zero.  Locally
$g=h\,uv$ with $h$ a unit, so after restriction to $S$ one has
$i^*(dg)=h(v\,du+u\,dv)$; both coefficients vanish on the preimage of $D$.
Proposition~\ref{wbII:prop:conductor}
therefore gives a nonzero section in \eqref{wbII:eq:CN:4.6}.  This proves the theorem.

\end{enumerate}
\end{proof}

\paragraph{what the theorem does and does not refute}

\hypertarget{tag-03IS}{}\label{wbII:rem:scope}
The bundle $A$ is non-torsion on $S$: restriction to
\(C_0\times\{r\}\) would otherwise make $A_\lambda$ torsion.  Thus the
example meets the non-torsion condition attached to the p.~692 reduction.
In fact Lemma~\ref{wbII:lem:extending-gluing} gives an extending choice of $A$ on
the displayed smooth total space.  If $\widetilde L:=\omega_{\mathcal X}\otimes
\widetilde A^{-1}$, then
$(\widetilde L^\vee\otimes\omega_{\mathcal X})|_S\simeq A$, so the example
also has the Serre-dual twist appearing in the CDP argument.  This family
nevertheless does not have
the global topological and algebraic-dimension hypotheses of CDP20
Theorem~2.2.  The conclusion is exactly the formal one stated in
Theorem~\ref{cdp:st:B1}: the two displayed endpoint vanishings do not
control the torsion-supported lifting obstruction $\delta_A$.

\Uses{the hypotheses, notation and definitions stated here; \cref{wbII:lem:extending-gluing}; \cref{cdp:st:B1}}


\subsection{Cartier-divisor duality without a normality hypothesis}

The next two results are explicit re-derivations of classical analytic
Cartier duality.  The trace duality and complete-intersection identification
are supplied by \citet[Theorem~2, pp.~82--84, and Proposition~2, p.~91]{RR70push}.
The proofs retain the full scheme structure of the fibre and state
separately where compactness is needed.

\begin{theorem}[Duality for a compact Cartier divisor]
\hypertarget{tag-03IY}{}\label{wbII:thm:cartier-duality}
\StatementRecord{wbII:thm:cartier-duality}{proved}{03IY}{\cite[Proposition~2 and pp.~90--91]{RR70push}; \cite[Theorem~2, pp.~82--83]{RR70push}}

Let $X$ be a compact smooth complex threefold,
$i:W\hookrightarrow X$ an effective Cartier divisor, and $E$ a holomorphic
vector bundle.  Put $F=E|_W$ and
\[
\omega_W=(\omega_X\otimes\mathcal O_X(W))|_W.
\]
For $q=0,1,2$ there are natural perfect pairings
\[
H^q(W,F)^\vee
\simeq H^{2-q}(W,F^\vee\otimes\omega_W).
\tag{W.11}\hypertarget{tag-03EB}{}\label{wbII:eq:W:5.11}
\]
Moreover $H^q(W,F)=0$ for $q>2$.
\end{theorem}



\begin{proof}

\begin{enumerate}[label=\arabic*.,leftmargin=*]
\item Let \(i:W\hookrightarrow X\) be the given Cartier inclusion.  The
length-one locally free resolution
\[
 0\longrightarrow\mathcal O_X(-W)
 \xrightarrow{\;\cdot g\;}\mathcal O_X
 \longrightarrow i_*\mathcal O_W\longrightarrow0
\]
is locally defined by a nonzerodivisor \(g\); changing \(g\) by a unit
gives the usual transition functions for \(\mathcal O_X(W)\).
Applying sheaf Hom into \(\omega_X\), the degree-zero map is
injective and its cokernel is
\[
 \mathcal E xt^1_X(i_*\mathcal O_W,\omega_X)
 =i_*\bigl((\omega_X\otimes\mathcal O_X(W))|_W\bigr)
 =i_*\omega_W.
\]
The other sheaf Ext groups vanish.  Equivalently, \(W\) is a
two-dimensional Cohen--Macaulay Gorenstein complex space and its
analytic dualizing complex is
\[
 K_W^\bullet\simeq\omega_W[2].
\]
This is the hypersurface calculation in
\cite[Proposition~2 and pp.~90--91]{RR70push}.

\item For the locally free sheaf \(F=E|_W\), the analytic duality trace of
\cite[Theorem~2, pp.~82--83]{RR70push} supplies the
functorial pairing, and identifies the dual of \(H_c^q(W,F)\) with
\[
 \operatorname{Ext}^{-q}_W(F,K_W^\bullet)
 =\operatorname{Ext}^{2-q}_W(F,\omega_W)
 =H^{2-q}(W,F^\vee\otimes\omega_W).
\]
The last equality follows because the positive sheaf Ext groups
of a locally free first argument vanish.  Since \(W\) is compact,
\(H_c^q(W,F)=H^q(W,F)\); coherent cohomology is finite-dimensional,
so this pairing is perfect in the ordinary algebraic sense.
Under this trace pairing, the isomorphism
\eqref{wbII:eq:W:5.11} is contravariant in \(F\) by precomposition in the Ext
argument.

\item For \(q>2\), the displayed Ext degree \(2-q\) is negative.  The
derived Hom of the locally free sheaf \(F\), concentrated in degree
zero, into the sheaf \(\omega_W\), concentrated in degree zero, has
no negative cohomology.  Hence the duality theorem gives
\(H^q(W,F)=0\).  This also follows directly from coherent
cohomological dimension on the compact surface \(W\).

\end{enumerate}
\end{proof}



\begin{proposition}[absolute duality on a compact Cartier fibre]
\hypertarget{tag-03IJ}{}\label{wbII:prop:cartier-duality-absolute}
\StatementRecord{wbII:prop:cartier-duality-absolute}{proved}{03IJ}{\cite[Proposition~2 and the complete-intersection calculation,
pp.~90--91]{RR70push}; \cite[Theorem~2, pp.~82--83]{RR70push}; \cite{Grauert1960D}}

Let $X$ be a smooth complex threefold,
and let $i:W\hookrightarrow X$ be a compact effective Cartier divisor that
is the full scheme-theoretic fibre of a holomorphic map $f:X\to B$ to a
smooth curve.  For every vector bundle $E$ on $X$ there is a functorial
perfect pairing
\[
 H^2(W,E|_W)\times
 H^0(W,E^\vee|_W\otimes\omega_X|_W)\longrightarrow\mathbb C.
\]
Equivalently,
\[
 H^2(W,E|_W)^\vee\simeq
 H^0(W,E^\vee|_W\otimes\omega_X|_W).
\tag{W.12}\hypertarget{tag-03EC}{}\label{wbII:eq:W:5.15a}
\]
\end{proposition}



\begin{proof}

\begin{enumerate}[label=\arabic*.,leftmargin=*]
\item An effective Cartier divisor in a smooth threefold is locally a
two-dimensional hypersurface quotient of a regular local ring.  Hence $W$
is Cohen--Macaulay and Gorenstein.  If $K_W^\bullet$ denotes its analytic
dualizing complex, the complete-intersection calculation of Ramis--Ruget
gives
\[
 K_W^\bullet\simeq\omega_W[2]
\]
with $\omega_W$ invertible
\cite[Proposition~2 and the complete-intersection calculation,
pp.~90--91]{RR70push}.  Cartier adjunction gives
\[
 \omega_W\simeq
 (\omega_X\otimes\mathcal O_X(W))|_W.
\]
Because $W$ is the full fibre over $b\in B$, one has
$\mathcal O_X(W)\simeq f^*\mathcal O_B(b)$ locally along $W$; its
restriction to $W$ is the pullback of a one-dimensional vector space.
Choosing a nonzero vector identifies
$\mathcal O_X(W)|_W\simeq\mathcal O_W$ and
$\omega_W\simeq\omega_X|_W$.  A different choice rescales the final trace
pairing by a nonzero scalar.

\item Ramis--Ruget duality pairs $H_c^p(W,\mathcal G)$ with
$\operatorname{Ext}^{-p}_W(\mathcal G,K_W^\bullet)$
\cite[Theorem~2, pp.~82--83]{RR70push}.  Since $W$ is compact,
$H_c^p=H^p$.  Taking $p=2$ and $\mathcal G=E|_W$, the shift and local
freeness give
\[
 \operatorname{Ext}^{-2}_W(\mathcal G,K_W^\bullet)
 \simeq\operatorname{Hom}_W(\mathcal G,\omega_W)
 =H^0(W,E^\vee|_W\otimes\omega_W).
\]
Grauert finiteness makes both coherent cohomology spaces finite-dimensional
\cite{Grauert1960D}; the strong analytic duality is therefore the ordinary
algebraic perfect pairing \eqref{wbII:eq:W:5.15a}.  Its variance is contravariant in $E|_W$
by precomposition on the Hom term.

\end{enumerate}
\end{proof}



\subsection{Flatness and top-degree analytic base change}

\begin{lemma}[Flatness and Cartier fibres over a curve]
\hypertarget{tag-03I8}{}\label{wbII:lem:curve-flatness}
\StatementRecord{wbII:lem:curve-flatness}{proved}{03I8}{\cite[Tag~0AUW]{StacksProject}}

Let $f:X\to B$ be a nonconstant holomorphic map from a connected
complex manifold to a Riemann surface.  Then $f$ is flat, every holomorphic
vector bundle on $X$ is $f$-flat, and every full scheme-theoretic fibre is an
effective Cartier divisor of pure codimension one.  If $X$ is compact, then
$f$ is proper.
\end{lemma}



\begin{proof}

\begin{enumerate}[label=\arabic*.,leftmargin=*]
\item Fix $x\in X$, $c=f(x)$, and put
$R=\mathcal O_{B,c}$, $S=\mathcal O_{X,x}$.  The ring $R$ is a discrete
valuation ring and $S$ is a regular local domain.  The homomorphism
$R\to S$ is injective.  Indeed, if a nonzero germ $h\in R$ mapped to zero,
$h\circ f$ would vanish on a neighbourhood of $x$.  The zero set of a
nonzero one-variable holomorphic germ is isolated, so $f$ would be constant
on that neighbourhood and therefore, by the identity theorem on connected
$X$, constant everywhere.

\item Thus every nonzero element of $R$ acts on the domain $S$ as a nonzero
element and hence as a nonzerodivisor.  Therefore $S$ is torsion-free as an
$R$-module.  A module over a discrete valuation ring is flat if and only if
it is torsion-free \cite[Tag~0AUW]{StacksProject}; hence $S$ is $R$-flat.
This holds at every point, so $f$ is flat.  A vector bundle is locally a
finite free $\mathcal O_X$-module, and a finite direct sum of $R$-flat
modules is flat; hence it is $f$-flat.  If $t$ is a local parameter at
$c$, then $f^*t$ is a nonzero element of the domain $S$ and thus a
nonzerodivisor.  The full scheme-theoretic fibre is therefore locally cut
out by one nonzerodivisor, so it is an effective Cartier divisor of pure
codimension one.  Finally, when $X$ is compact, the underlying continuous
map to the Hausdorff surface $B$ is proper.

\end{enumerate}
\end{proof}



\begin{proposition}[Analytic base change in degree two]
\hypertarget{tag-03II}{}\label{wbII:prop:analytic-base-change}
\StatementRecord{wbII:prop:analytic-base-change}{proved}{03II}{\cite[pp.~120--121]{BS76push}; \cite[Lemma~1.5]{CDP98}}

Let $f:X\to B$ be proper between complex spaces, let
$\mathcal F$ be coherent and $f$-flat, and fix $c\in B$.  If
\[
H^3(X_c,\mathcal F|_{X_c})=0,
\tag{W.13}\hypertarget{tag-03ED}{}\label{wbII:eq:W:5.16}
\]
then the canonical fibre map
\[
\beta_c^2:
(R^2f_*\mathcal F)_c\otimes_{\mathcal O_{B,c}}\mathbb C(c)
\longrightarrow H^2(X_c,\mathcal F|_{X_c})
\tag{W.14}\hypertarget{tag-03EE}{}\label{wbII:eq:W:5.17}
\]
is an isomorphism.
\end{proposition}



\begin{proof}
Apply B\u{a}nic\u{a}--St\u{a}n\u{a}\c{s}il\u{a},
Chapter~III, \S3, Corollary~3.5, with $q=3$
\cite[pp.~120--121]{BS76push}.  Its hypotheses are precisely
properness, coherence, and flatness.  Condition \eqref{wbII:eq:W:5.16} first forces the
degree-three direct image to vanish near $c$.  The corollary then applies the
equivalent base-change conditions of its preceding Theorem~3.4 in degree
$q-1=2$ and states that the canonical map \eqref{wbII:eq:W:5.17} is bijective.  The top-degree application also appears in
\cite[proof of Lemma~1.5, part~(c), p.~80]{CDP98}.

\end{proof}



\begin{corollary}[Top-degree base change for surface fibres]
\hypertarget{tag-03DV}{}\label{wbII:cor:surface-top-base-change}
\StatementRecord{wbII:cor:surface-top-base-change}{proved}{03DV}{\SeriesRef{wbII:eq:W:5.17-U}, \SeriesRef{wbII:prop:analytic-base-change}; \cite[Chapter~III, \S4, Theorem~4.1, p.~124]{BS76push}}

Let $f:X\to B$ be a proper flat holomorphic map to a Riemann
surface whose full scheme-theoretic fibres have pure complex dimension two,
and let $E$ be a holomorphic vector bundle on $X$.  Then
\[
 R^qf_*E=0\quad(q>2)
 \tag{W.15}\hypertarget{tag-03FU}{}\label{wbII:eq:W:5.17-V}
\]
and, for every $b\in B$, including a point with nonreduced fibre, the
canonical map
\[
 (R^2f_*E)_b\otimes_{\mathcal O_{B,b}}\mathbb C(b)
 \xrightarrow{\ \sim\ }H^2(X_b,E|_{X_b})
 \tag{W.16}\hypertarget{tag-03FS}{}\label{wbII:eq:W:5.17-T}
\]
is an isomorphism.
\end{corollary}



\begin{proof}
On a sufficiently small neighbourhood $U$ of $b$, the finite-complex theorem
used in analytic cohomology and base change represents $Rf_*E|_U$ by a
bounded complex $K^\bullet$ of finite free $\mathcal O_U$-modules
\cite[Chapter~III, \S4, Theorem~4.1, p.~124]{BS76push}.  Flatness of
$E$ over $B$ identifies its derived fibre with the ordinary restriction to
the full scheme-theoretic fibre.  Over the discrete valuation ring
$A=\mathcal O_{B,b}$, the fibre map of
\cref{wbII:prop:analytic-base-change} occurs in the universal-coefficient sequence
\[
0\longrightarrow H^q(K^\bullet)_b\otimes_A\mathbb C(b)
 \longrightarrow H^q(K^\bullet\otimes_A\mathbb C(b))
 \longrightarrow
 \operatorname{Tor}_1^A(H^{q+1}(K^\bullet)_b,\mathbb C(b))
 \longrightarrow0.
\tag{W.17}\hypertarget{tag-03FT}{}\label{wbII:eq:W:5.17-U}
\]
If some $R^qf_*E$ with $q>2$ were nonzero, choose the largest such $q$ and a
point of its support.  Nakayama's lemma makes the left term of \eqref{wbII:eq:W:5.17-U}
nonzero there, while maximality kills the Tor term.  This would give nonzero
coherent cohomology in degree $q>2$ on a compact two-dimensional complex
space, a contradiction.  Thus \eqref{wbII:eq:W:5.17-V} holds.  For $q=2$, the only
obstruction in \eqref{wbII:eq:W:5.17-U} is
$\operatorname{Tor}_1^A((R^3f_*E)_b,\mathbb C(b))$, which is zero by
\eqref{wbII:eq:W:5.17-V}, proving \eqref{wbII:eq:W:5.17-T}.

\end{proof}



\subsection{Marked data used in the restriction arguments}
Here $f:X\to B=\mathbb P^1$ is the standard family of Paper I,
$B^\circ=B\setminus\{p_0,p_1,p_2\}$, and $F$ is a smooth fibre.
The cusp is $W=f^{-1}(p_0)$ and $i:W\hookrightarrow X$ is its inclusion.
Its normalization is $\eta:\widetilde W\simeq\mathrm{dP}_6\to W$.
Use the ordered dual lattices
\[
 V=H^1(F;\mathbb Z)=\mathbb Z\langle\gamma,u,w,\delta\rangle,
 \qquad
 \Lambda=H_1(F;\mathbb Z)
 =\mathbb Z\langle\widehat\gamma,\widehat u,\widehat w,\widehat\delta\rangle.
\]
Vectors are columns. The cohomology monodromies and their dual homology
actions are
\[
T_1=\begin{pmatrix}1&0&-6&2\\0&-1&1&1\\0&-1&0&1\\0&0&0&1\end{pmatrix},
\quad
T_2=\begin{pmatrix}1&6&0&-3\\0&0&-1&1\\0&1&0&0\\0&0&0&1\end{pmatrix},
\quad A_j=(T_j^{-1})^{\mathsf t}.
\]
Put $T_0=(T_1T_2)^{-1}$ and let $G$ denote the generated monodromy group.
The finite-filling orders, translations and invariant forms are
\[
\begin{gathered}
(m_1,m_2)=(3,4),\qquad
v_1=(1,2,-4,0)^{\mathsf t},\quad v_2=(-1,-3,3,0)^{\mathsf t},\\
\eta_1=2u+w+3\delta,\quad \eta_2=u+w+2\delta,
\qquad q=u\wedge w+6\gamma\wedge\delta.
\end{gathered}
\]
The actual period functions of Paper I define
\[
\Pi(z)=\begin{pmatrix}6\mu(z)&\tau(z)&1&0\\
\beta(z)&\mu(z)&0&1\end{pmatrix},
\quad
\operatorname{Im}\beta-
\frac{6(\operatorname{Im}\mu)^2}{\operatorname{Im}\tau}<0,
\]
so $\Pi(z):\Lambda_{\mathbb R}\to\mathbb C^2$ is a real isomorphism.
The marking is $\Xi(z,x+\Lambda)=[z,\Pi(z)x]$.
The finite overlap on the $s_j$-cover is translation by
$(\log s_j)/(2\pi i)\,\Pi(s_j)v_j$.
At the cusp use the unchanged zero-section gluing, with no additional
vertical translation. These are the data constructed and proved in
\CompendiumRef{I}{wb:prop:period-input-audit}, \CompendiumRef{I}{wb:thm:candidate-cusp-geometry}, \CompendiumRef{I}{wb:prop:candidate-finite-fillings}. The arguments
below do not replace these data by a topologically equivalent model.

\subsection{A topology-independent \texorpdfstring{$\operatorname{Pic}^0$}{Pic0} route and three all-Pic routes}

The vanishing of the normalized restriction of every global line bundle
for this threefold is already established in
\citet[Theorem~10.5, proof Step~1, equation~(10.5), p.~88]{AC2026}.
The complete arguments in
\cref{wbII:prop:leray-edge-restriction}, \cref{wbII:prop:direct-boundary-picard}, \cref{wbII:thm:analytic-all-picard}
are retained as three distinct proofs of that known conclusion.  The
first uses the integral Leray edge, the second uses the boundary norm
identities and local specialization, and the third uses holomorphic
factors and toric support functions.  The preliminary
\(\operatorname{Pic}^0\) argument isolates the weaker hypotheses already
sufficient for topologically trivial twists.

Let $i:W\hookrightarrow X$ denote the cusp-fibre inclusion and
$\eta:\widetilde W\simeq\mathrm{dP}_6\to W$ its normalization.

\begin{proposition}[Pullback of topologically trivial line bundles]
\hypertarget{tag-03IO}{}\label{wbII:prop:pic0-pullback}
\StatementRecord{wbII:prop:pic0-pullback}{proved}{03IO}{the definitions and constructions stated in the proof}

Assume that $\omega_X$ is trivial on a neighbourhood of $W$.
For every $L\in\operatorname{Pic}^0(X)$, put
\[
 A_L=(L^\vee\otimes\omega_X)|_W.
\]
Then
\[
 \eta^*A_L\simeq\mathcal O_{\mathrm{dP}_6}.
\tag{W.18}\hypertarget{tag-03FV}{}\label{wbII:eq:W:5.17a}
\]
\end{proposition}



\begin{proof}
Every line bundle in $\operatorname{Pic}^0(X)$ has $c_1(L)=0$.  Naturality
therefore gives $c_1(\eta^*i^*L)=0$.  The degree-six del Pezzo surface is the
blow-up of $\mathbb P^2$ in three noncollinear points, so
$H^1(\mathrm{dP}_6,\mathcal O)=0$.  The exponential sequence consequently
makes
\[
 c_1:\operatorname{Pic}(\mathrm{dP}_6)\longrightarrow
 H^2(\mathrm{dP}_6;\mathbb Z)
\]
injective.  Hence $\eta^*i^*L$ is trivial.  The assumed neighbourhood
trivialization of $\omega_X$ also pulls back trivially, which proves
\eqref{wbII:eq:W:5.17a}.

\end{proof}



\begin{proposition}[Restriction of line bundles to the fibre normalization]
\hypertarget{tag-03IG}{}\label{wbII:prop:all-pic-restriction}
\StatementRecord{wbII:prop:all-pic-restriction}{proved}{03IG}{\SeriesRef{wbII:prop:conductor}, \SeriesRef{wbII:prop:pic0-pullback}}

Define
\[
 r=\eta^*i^*:\operatorname{Pic}(X)\longrightarrow
 \operatorname{Pic}(\mathrm{dP}_6).
\]
Under the same local canonical trivialization,
\[
 \eta^*A_L\simeq r(L)^\vee
 \qquad(L\in\operatorname{Pic}(X)).
\tag{W.19}\hypertarget{tag-03GL}{}\label{wbII:eq:W:5.17b}
\]
Thus the trivial-pullback hypothesis used by
\cref{wbII:prop:conductor} holds simultaneously for every line bundle exactly
when $r=0$.  It always holds on $\operatorname{Pic}^0(X)$.
\end{proposition}



\begin{proof}
Pullback commutes with tensor product and duals, while the local
trivialization gives $\eta^*i^*\omega_X\simeq\mathcal O$.  This is \eqref{wbII:eq:W:5.17b},
and its uniform equivalence with $r=0$ is immediate.  The final assertion is
\cref{wbII:prop:pic0-pullback}.

\end{proof}



\begin{proposition}[Leray-edge vanishing of normalized cusp restrictions]
\hypertarget{tag-03IM}{}\label{wbII:prop:leray-edge-restriction}
\StatementRecord{wbII:prop:leray-edge-restriction}{proved}{03IM}{\SeriesRef{wbII:eq:W:5.17d}, \SeriesRef{wbII:eq:W:5.17e}, \SeriesRef{wbII:prop:pic0-pullback}; \cite[Tags~012O, 012U and~012W]{StacksProject}; \cite[Tag~09V6]{StacksProject}; \cite[Tags~011N, 011O and~012U]{StacksProject}}

For the constructed map $f:X\to B$ and every
$L\in\operatorname{Pic}(X)$, the degree-two Leray edge section satisfies
\[
 e(c_1(L))=0\in H^0(B,R^2f_*\mathbb Z),
 \qquad
 \eta^*i^*L\simeq\mathcal O_{\mathrm{dP}_6}.
\tag{W.20}\hypertarget{tag-03GM}{}\label{wbII:eq:W:5.17c}
\]
\end{proposition}



\begin{proof}

\begin{enumerate}[label=\arabic*.,leftmargin=*]
\item For every $b\in B$,
proper base change and functoriality of the filtered-complex spectral
sequence give a commutative square
\[
\begin{CD}
 H^2(X;\mathbb Z) @>{e}>> H^0(B,R^2f_*\mathbb Z)\\
 @V{i_b^*}VV @VV{\operatorname{ev}_b}V\\
 H^2(X_b;\mathbb Z) @= H^2(X_b;\mathbb Z).
\end{CD}
\tag{W.21}\hypertarget{tag-03GN}{}\label{wbII:eq:W:5.17d}
\]
Indeed, pull the Leray filtered complex back along $\{b\}\to B$; the
spectral sequence over a point has only column zero, so its edge is the
identity, while the right vertical map is the proper-base-change
identification.  This uses the functoriality and convergence statements in
\cite[Tags~012O, 012U and~012W]{StacksProject} and topological proper base
change \cite[Tag~09V6]{StacksProject}.  The connecting homomorphism in the
exponential sequence is natural under $i_b^*$, so \eqref{wbII:eq:W:5.17d} sends
$c_1(L)$ to $c_1(L|_{X_b})$.

\item Put $s_L=e(c_1(L))$.  The image of an edge morphism consists of permanent
cycles, hence
\[
 s_L\in E_\infty^{0,2}\subseteq\ker
 \bigl(d_2^{0,2}:E_2^{0,2}\to E_2^{2,1}\bigr);
\tag{W.22}\hypertarget{tag-03GO}{}\label{wbII:eq:W:5.17e}
\]
see \cite[Tags~011N, 011O and~012U]{StacksProject}.  For the standard family, \cite[Propositions~7.26--7.27, pp.~61--63]{AC2026}
give $E_2^{0,2}=\mathbb Z(2q)$ and
$E_2^{2,1}=\mathbb Z\,\omega_{\rm AC}[\delta]$, with
$d_2^{0,2}(2q)=\varepsilon p\,\omega_{\rm AC}[\delta]$,
where $\varepsilon\in\{1,-1\}$ is the common differential sign and
$p=12\ell_0-4\ell_1-3\ell_2=-1$.
Define the generator $\omega:=\varepsilon p\,\omega_{\rm AC}$; equivalently,
$\omega[\delta]$ is the actual element $d_2^{0,2}(2q)$.  It is primitive
because $\varepsilon p\in\{1,-1\}$.  With this generator,
\[
 E_2^{0,2}=\mathbb Z(2q),\qquad
 E_2^{2,1}=\mathbb Z\,\omega[\delta],\qquad
 d_2^{0,2}(2q)=\omega[\delta].
\]
Thus $d_2^{0,2}$ is an isomorphism and \eqref{wbII:eq:W:5.17e} forces $s_L=0$.
Taking $b=p_0$ in \eqref{wbII:eq:W:5.17d} gives $c_1(i^*L)=0$, and hence
$c_1(\eta^*i^*L)=0$.  Finally, the injectivity of
$c_1:\operatorname{Pic}(\mathrm{dP}_6)\to H^2(\mathrm{dP}_6;\mathbb Z)$
proved in \cref{wbII:prop:pic0-pullback} makes $\eta^*i^*L$ trivial.

\end{enumerate}
\end{proof}



\begin{lemma}[The pair-of-pants defect class]
\hypertarget{tag-03IC}{}\label{wbII:lem:picard-crossed-defect}
\StatementRecord{wbII:lem:picard-crossed-defect}{proved}{03IC}{\SeriesRef{wbII:eq:W:5.17j}; \cite[Example~3.5, p.~4]{SchnellLecture3}}

Choose the compact pair-of-pants core $P\subset B^\circ$ and
the zero-section whiskers used in \CompendiumRef{I}{inv:B3.1}, and put
$J_P=f^{-1}(P)$.  Orient its boundary meridians clockwise.  Write
$\rho_0$ for the projection of the zero-section cusp meridian
$\widetilde a_0$.  The standard family uses the tautological cusp
gluing, whose affine regluing class is $\mu=0$.  Thus its flat-marked cusp lift
has no vertical translation and $\ell_0=\gamma(\mu)=0$.  The base loops and their chosen lifts therefore satisfy
\[
 \rho_1\rho_2\rho_0=1,
 \qquad
 \widehat\rho_1\widehat\rho_2\widetilde a_0=1.
\tag{W.23}\hypertarget{tag-03GQ}{}\label{wbII:eq:W:5.17g}
\]
For $L\in\operatorname{Pic}(X)$, write
$E=c_1(L|_F)\in\bigwedge^2V$.  Then $E\in(\bigwedge^2V)^G$, and there is a
crossed homomorphism
\[
 z:\pi_1(P)\longrightarrow V_{\mathbb R},\qquad
 z(gh)=z(g)+T_gz(h),
\tag{W.24}\hypertarget{tag-03GR}{}\label{wbII:eq:W:5.17h}
\]
whose class records
$c_1(L|_{J_P})-[\Omega_E]$.  Here $A_g$ is the linear part of the
left deck action $(b,x)\mapsto(gb,A_gx)$, and
$T_g=(A_g^{-1})^{\mathsf t}$ is its action on coefficient covectors.
A clockwise loop lifts from $b$ to $gb$; the endpoint return map $g^{-1}$
acts on homology by $A_g^{-1}$ and on cohomology by $A_g^{\mathsf t}=T_g^{-1}$.
The crossed-homomorphism convention above uses the left deck action $T_g$.  Here $\Omega_E$ is the closed, translation-
invariant vertical two-form induced by $E$ in the flat marking.  If
$d_i=z(\rho_i)$ and a bar denotes the image in
$V_G\otimes\mathbb R$, then
\[
 \bar d_1+\bar d_2+\bar d_0=0.
\tag{W.25}\hypertarget{tag-03GS}{}\label{wbII:eq:W:5.17i}
\]
The three barred classes, and the number
$\langle d_0,\widehat\delta\rangle$, do not depend on the chosen crossed-
homomorphism representative.  More precisely, use the base-first convention
for the edge injection in \eqref{wbII:eq:W:5.17j}.  For any oriented loop $\rho\subset P$
and any $A_\rho$-invariant $\nu\in\Lambda$, its flat-marked sweep satisfies
\[
 \bigl\langle z(\rho),\nu\bigr\rangle
 =-\bigl\langle c_1(L|_{J_P})-[\Omega_E],
 a_*([S^1_\nu]\times[\widetilde\rho])\bigr\rangle.
\tag{W.26}\hypertarget{tag-03GU}{}\label{wbII:eq:W:5.17j'}
\]
Here $a$ is the fibrewise translation action and $\widetilde\rho$ is the
flat-marked lift of $\rho$.
\end{lemma}



\begin{proof}

\begin{enumerate}[label=\arabic*.,leftmargin=*]
\item The flat marking identifies the pullback of $J_P$ to the universal cover of
$P$ with
$\widetilde P\times(\Lambda_{\mathbb R}/\Lambda)$.  Since $E$ is the
restriction of a class on the total space, it is fixed by every deck coefficient action
$T_g$.  The constant vertical form representing $E$ is therefore invariant
under the deck action and descends to the asserted form $\Omega_E$ on
$J_P$.

\item The pair of pants retracts onto a graph.  Hence the Serre spectral sequence
of the marked real torus bundle has only columns zero and one and gives the
short exact edge sequence
\[
0\longrightarrow H^1(P;V_{\mathbb R})
 \longrightarrow H^2(J_P;\mathbb R)
 \longrightarrow H^0(P;\textstyle\bigwedge^2V_{\mathbb R})
 \longrightarrow0.
\tag{W.27}\hypertarget{tag-03GT}{}\label{wbII:eq:W:5.17j}
\]
Both $c_1(L|_{J_P})$ and $[\Omega_E]$ map to $E$, so their difference is
the image of a unique class in $H^1(P;V_{\mathbb R})$.  The standard group-
cohomology model for a graph identifies this group with crossed
homomorphisms modulo principal crossed homomorphisms
\cite[Example~3.5, p.~4]{SchnellLecture3}; this gives \eqref{wbII:eq:W:5.17h}.

\item For the slant convention in \eqref{wbII:eq:W:5.17j'}, give an oriented
edge of a graph spine of $P$ a coordinate $s$ increasing with its based
loop and represent the corresponding $V_{\mathbb R}$-valued cellular
one-cochain by a base-first form $ds\wedge\beta_s$.  On the sweep of an
$A_\rho$-invariant fibre circle $S^1_\nu$, Fubini's formula gives
\[
\int_{S^1_s\times S^1_\nu}ds\wedge\beta_s
 =\int_{S^1_s}\langle\beta_s,\nu\rangle\,ds
 =\langle z(\rho),\nu\rangle.
\]
The sweep in \eqref{wbII:eq:W:5.17j'} is written in the fibre-first order
$[S^1_\nu]\times[\widetilde\rho]$, which differs from this base-first
cross-product orientation by $(-1)^{1\cdot1}=-1$; this is the displayed
minus sign.
Changing the cellular representative by a principal cochain changes both
sides by the same boundary term, which vanishes because $A_\rho\nu=\nu$.

\item Applying $z$ to \eqref{wbII:eq:W:5.17g} gives
\[
 d_1+T_1d_2+T_1T_2d_0=0.
\]
All $T_g$ act trivially on the global coinvariants, proving \eqref{wbII:eq:W:5.17i}.  A
change of representative has the form
$z(g)\mapsto z(g)+(T_g-I)a$.  Its image in $V_G\otimes\mathbb R$ is zero.
Moreover the displayed matrices in
\CompendiumRef{I}{def:common-conventions} fix $\widehat\delta$ under every dual
action $A_g$.  Therefore
\[
 \langle(T_g-I)a,\widehat\delta\rangle
 =\langle a,(A_g^{-1}-I)\widehat\delta\rangle=0.
\]

\end{enumerate}
\end{proof}



\begin{lemma}[The finite-cover norm identity]
\hypertarget{tag-03IE}{}\label{wbII:lem:picard-finite-norm}
\StatementRecord{wbII:lem:picard-finite-norm}{proved}{03IE}{\SeriesRef{wbII:lem:picard-crossed-defect}, \SeriesRef{wbII:eq:W:5.17g}, \SeriesRef{wbII:eq:W:5.17m'}, \SeriesRef{wbII:eq:W:5.17l}, \SeriesRef{wbII:eq:W:5.17l''}, \SeriesRef{wbII:eq:W:5.17l'''}}

Retain the notation of
\cref{wbII:lem:picard-crossed-defect} and the normalized root covers in \CompendiumRef{I}{fam:src:A.1.2}.  For
$j=1,2$, let
\[
 N_{T_j}=I+T_j+\cdots+T_j^{m_j-1}.
\]
With the clockwise conventions of \eqref{wbII:eq:W:5.17g}, the pullback to the
$m_j$-fold $s_j$-cover gives
\[
 N_{T_j}d_j=\iota_{v_j}E
 \quad\text{in }V_{\mathbb R},
\qquad
 m_j\bar d_j=[\iota_{v_j}E]
 \quad\text{in }V_G\otimes\mathbb R.
\tag{W.28}\hypertarget{tag-03GV}{}\label{wbII:eq:W:5.17k}
\]
\end{lemma}



\begin{proof}

\begin{enumerate}[label=\arabic*.,leftmargin=*]
\item Restriction of the crossed homomorphism to the subgroup generated by
$\rho_j^{m_j}$ has value
\[
 z(\rho_j^{m_j})
 =\sum_{r=0}^{m_j-1}T_j^rz(\rho_j)=N_{T_j}d_j.
\tag{W.29}\hypertarget{tag-03GW}{}\label{wbII:eq:W:5.17l}
\]
This is the cohomological pullback of the annular defect to the
$s_j$-cover, expressed by the norm $N_{T_j}$.  A principal change
$d_j\mapsto d_j+(T_j-I)a$ does not alter this value, since
$N_{T_j}(T_j-I)a=(T_j^{m_j}-I)a=0$.

\item Let $Q_j:J'_j\to N_j$ be the finite quotient map.  The global flat marking
$\Xi$ over the full disc $\Delta_j$ makes the constant vertical form
$\Omega_E$ well defined on $J'_j$.  Radial retraction onto the central
torus makes
\[
 H^2(J'_j;\mathbb R)\longrightarrow H^2(A_0^{(j)};\mathbb R)
\]
an isomorphism.  On the central fibre the quotient is exactly $\pi_j$, and
naturality gives
\[
 (Q_j|_{A_0^{(j)}})^*c_1(L|_{S_j})
 =\pi_j^*c_1(L|_{S_j})=E.
\tag{W.30}\hypertarget{tag-03GX}{}\label{wbII:eq:W:5.17l'}
\]
The second equality follows from the normalized root cover in
\CompendiumRef{I}{fam:src:A.1.2} and the radial restriction above; on a nearby fibre the gluing map $h_j$
is a translation and therefore acts trivially on $H^2$ of that fibre.
Consequently the two classes
$c_1(Q_j^*(L|_{N_j}))$ and $[\Omega_E]$ have the same central restriction,
and the retraction isomorphism proves
\[
 c_1(Q_j^*(L|_{N_j}))-[\Omega_E]=0
 \quad\text{in }H^2(J'_j;\mathbb R).
\tag{W.31}\hypertarget{tag-03GY}{}\label{wbII:eq:W:5.17l''}
\]

\item On $\Delta_j^*$, the connector in
\CompendiumRef{I}{pub:I:eq:PD8} satisfies
$h_j\widetilde g_j^{\log}h_j^{-1}=\widetilde g_j^0$.  This identity shows that $h_j$
descends from the logarithmic quotient $N_j^*$ to the linear/global
quotient $J^*$.  Thus
\[
 Q_j^*(L|_{N_j})|_{\Delta_j^*}
 \simeq h_j^*(L|_{J^*}).
\tag{W.32}\hypertarget{tag-03GZ}{}\label{wbII:eq:W:5.17l'''}
\]
The logarithmic overlap is the holomorphic translation
\[
 h_j[s,\zeta]=[s,\zeta+\alpha(s)\Pi(s)v_j],
 \qquad \alpha(s)=\frac{\log s_j}{2\pi i}.
\]
To compute its effect on a real cohomology class, retain the exact
real marking
\[
 \Xi(s,x+\Lambda)=[s,\Pi(s)x],\qquad
 J_s=\Pi(s)^{-1}\circ i\circ\Pi(s)
 \quad\text{on }\Lambda_{\mathbb R}.
\]
The period matrix \(\Pi(s)\) is a real isomorphism by the displayed period inequality;
thus \(J_s\) is a well-defined smooth family of real complex
structures.  On a branch of the logarithm the real
translation vector is
\[
 a_j(s)=\operatorname{Re}\alpha(s)\,v_j+
        \operatorname{Im}\alpha(s)\,J_s v_j.
 \tag{W.33}\hypertarget{tag-03H0}{}\label{wbII:eq:W:5.17m}
\]
Indeed, applying \(\Pi(s)\) gives precisely
\(\alpha(s)\Pi(s)v_j\).  A different logarithm branch changes
\(a_j\) by an integral multiple of \(v_j\), so its torus translation
is branch-independent.  The real and imaginary parts in \eqref{wbII:eq:W:5.17m} therefore give
the translation in the real marking.

\item Write \(E=\frac12\sum_{a,b}E_{ab}\,dx^a\wedge dx^b\) in the
ordered real coordinates dual to the marked lattice; the matrix
\((E_{ab})\) is alternating and is constant over the covering
disc.  Direct substitution in the real map
\(h_j(s,x)=(s,x+a_j(s))\) gives
\[
 \begin{split}
 h_j^*\Omega_E-\Omega_E
  ={}&\sum_{a,b}E_{ab}\,d a_j^a\wedge dx^b\\
    &+\frac12\sum_{a,b}E_{ab}\,d a_j^a\wedge d a_j^b.
 \end{split}
 \tag{W.34}\hypertarget{tag-03H1}{}\label{wbII:eq:W:5.17m'}
\]
The last term has both arguments in the base directions.
In particular, it restricts to zero when the base is an oriented
circle.

\item Take one full clockwise circle in the \(s_j\)-cover,
\[
 s_j(\theta)=r e^{-2\pi i\theta},\qquad 0\le\theta\le1.
\]
Choose the logarithm continuously along that circle, so
\(\alpha(1)=\alpha(0)-1\).  The endpoints represent the same
point of the covering disc.  The marking is defined on the full
disc, and therefore \(J_{s_j(1)}=J_{s_j(0)}\).  Consequently
\[
 a_j(1)-a_j(0)=-v_j.
 \tag{W.35}\hypertarget{tag-03H2}{}\label{wbII:eq:W:5.17m''}
\]
For any \(\lambda\in\Lambda\), restrict \eqref{wbII:eq:W:5.17m'} to the sweep of
the marked circle \(x(u)=u\lambda\), using the base-first
orientation \(d\theta\wedge du\).  Its integral is exactly
\[
 \int_0^1 E\!\left(\frac{d a_j}{d\theta},\lambda\right)d\theta
 =E(a_j(1)-a_j(0),\lambda)
 =-E(v_j,\lambda).
\]
The horizontal term contributes zero.  Since marked lattice
circles span \(H_1(F;\mathbb R)\), these integrals determine the
mixed component in
\[
 H^1(S^1;\mathbb R)\otimes H^1(F;\mathbb R).
\]
Thus the mixed component of \(h_j^*\Omega_E-\Omega_E\) is
\(-\iota_{v_j}E\), with \((\iota_{v_j}E)(\lambda)=E(v_j,\lambda)\).

\item The pullback of the defect
\(c_1(L|_{J^*})-[\Omega_E]\) has zero fibre component and mixed
component \(N_{T_j}d_j\) by \eqref{wbII:eq:W:5.17l}.  Translation by \(h_j\) does
not change this mixed component: on the covering circle such
a fibre-zero class is represented by
\(d\theta\wedge\beta\), with \(\beta\) a constant fibre one-form,
and substituting \(dx\mapsto dx+d a_j\) adds only a base two-form,
which vanishes on the circle.  Equations \eqref{wbII:eq:W:5.17l''} and
\eqref{wbII:eq:W:5.17l'''} state that the total pulled-back class extends over
the covering disc and has zero mixed component.  It follows that
\[
 N_{T_j}d_j-\iota_{v_j}E=0.
\]
In the global coinvariants every \(T_j\) acts as the identity,
so \(m_j\bar d_j=[\iota_{v_j}E]\).

\end{enumerate}
\end{proof}



\begin{lemma}[The cusp boundary defect and Chern-class specialization]
\hypertarget{tag-03ID}{}\label{wbII:lem:picard-cusp-defect}
\StatementRecord{wbII:lem:picard-cusp-defect}{proved}{03ID}{\SeriesRef{wbII:lem:picard-crossed-defect}, \SeriesRef{wbII:eq:W:5.17g}, \SeriesRef{wbII:eq:W:5.17j'}, \SeriesRef{wbII:eq:W:5.17p}}

In the notation of
\cref{wbII:lem:picard-crossed-defect}, one has
\[
 \bar d_0=0\quad\text{in }V_G\otimes\mathbb R.
\tag{W.36}\hypertarget{tag-03H3}{}\label{wbII:eq:W:5.17n}
\]
Moreover, for every line bundle $L$ on the cusp neighbourhood $N_0$,
nearby-cycle specialization is compatible with first Chern classes:
\[
 \operatorname{sp}^2\bigl(c_1(L|_W)\bigr)=c_1(L|_F).
\tag{W.37}\hypertarget{tag-03H4}{}\label{wbII:eq:W:5.17o}
\]
\end{lemma}



\begin{proof}

\begin{enumerate}[label=\arabic*.,leftmargin=*]
\item Let $\widetilde a_0$ be the flat-marked clockwise zero-section meridian in
the lifted relation of \eqref{wbII:eq:W:5.17g}, and let
\[
 \xi_{\widehat\delta}
 =a_*([S^1_{\widehat\delta}]\times[\widetilde a_0])
\]
be its sweep under fibrewise translation by $\widehat\delta$.  The boundary
$\rho_0$ of $P$ was chosen to be its projection.  The zero-section
lift $\widetilde a_0$ has constant fibre coordinate by the overlap
calculation in \CompendiumRef{I}{inv:B3.1}.  Hence the
vertical form $\Omega_E$ vanishes on the swept torus, and the slant identity
\eqref{wbII:eq:W:5.17j'} gives
\[
 \langle d_0,\widehat\delta\rangle
 =-\langle c_1(L|_{J_P})-[\Omega_E],
          \xi_{\widehat\delta}\rangle
 =-\langle c_1(L|_{J_P}),\xi_{\widehat\delta}\rangle.
\tag{W.38}\hypertarget{tag-03H5}{}\label{wbII:eq:W:5.17p}
\]
Combine the exponential overlap coordinates $(x_1,x_2,t)$ given by $x_k=e^{2\pi i\zeta_k}$ and $t=e^{2\pi is}$ with
the unimodular cusp chart of \CompendiumRef{I}{mapsI:src:A.1.2}
\[
 (u,x_1,x_2)=\left(\frac{t}{x_1x_2},x_1,x_2\right).
\]
The solid torus
\[
 \mathbb S_{\widehat\delta}
 =\{|u|=\rho,\ x_1=1,\ |x_2|\le1\}
\]
has boundary parametrization
\[
 (\theta,\phi)\longmapsto
 \bigl(\rho e^{i(\theta-\phi)},1,e^{i\phi}\bigr)
\tag{W.39}\hypertarget{tag-03H7}{}\label{wbII:eq:W:5.17p0}
\]
and hence $t=\rho e^{i\theta}$; its $\phi$-circle is translation by
$\widehat\delta$.  Direct evaluation in the proved cusp marking therefore
identifies $\partial\mathbb S_{\widehat\delta}$ with
$\xi_{\widehat\delta}$.  Reversing $\theta$ supplies the fixed clockwise
convention and only changes the common orientation sign.
Since $L$ extends across $N_0$, its Chern class evaluates to zero on that
boundary.  Thus the left side of \eqref{wbII:eq:W:5.17p} is zero.  By
\cite[Proposition~7.26, pp.~61--62]{AC2026}, $V_G=\mathbb Z[\delta]$, and evaluation on the
globally invariant dual vector $\widehat\delta$ identifies
$V_G\otimes\mathbb R$ with $\mathbb R$.

\item For \eqref{wbII:eq:W:5.17o}, let $\ell=c_1(L|_{N_0})$, let
$\iota_0:W\hookrightarrow N_0$, and let
$\iota_t:F_t\hookrightarrow N_0$ be a sufficiently near smooth cusp fibre.
The neighbourhood realization of the nearby-cycle unit
$\mathbb Z_W\to\psi_{f_0}\mathbb Z$ in \cite[Theorem~B.1, pp.~101--105, especially (N1)--(N2), p.~101]{AC2026} gives the commutative square
\[
\begin{CD}
 H^2(N_0;\mathbb Z) @>{\iota_0^*}>> H^2(W;\mathbb Z)\\
 @V{\iota_t^*}VV @VV{\operatorname{sp}^2}V\\
 H^2(F_t;\mathbb Z) @= H^2(F_t;\mathbb Z).
\end{CD}
\tag{W.40}\hypertarget{tag-03H6}{}\label{wbII:eq:W:5.17p'}
\]
Applying this square to $\ell$ gives
$\operatorname{sp}^2(c_1(L|_W))=c_1(L|_{F_t})$.  Flat parallel transport
in the marked smooth family identifies the latter with $c_1(L|_F)=E$.
Equivalently, $c_1$ is the connecting morphism in the exponential sequence,
and restriction to the shrinking neighbourhoods and to $F_t$ commutes with
that morphism before passage to the direct limit.

\end{enumerate}
\end{proof}



\begin{proposition}[Direct boundary-Picard vanishing]
\hypertarget{tag-03IL}{}\label{wbII:prop:direct-boundary-picard}
\StatementRecord{wbII:prop:direct-boundary-picard}{proved}{03IL}{\SeriesRef{wbII:lem:picard-finite-norm}, \SeriesRef{wbII:eq:W:5.17i}, \SeriesRef{wbII:eq:W:5.17n}, \SeriesRef{wbII:eq:W:5.17o}, \SeriesRef{wbII:prop:pic0-pullback}}

For the standard threefold and every analytic line bundle
$L\in\operatorname{Pic}(X)$,
\[
 c_1(L|_F)=0,
 \qquad
 \eta^*i^*L\simeq\mathcal O_{\mathrm{dP}_6}.
\tag{W.41}\hypertarget{tag-03H8}{}\label{wbII:eq:W:5.17q}
\]
\end{proposition}



\begin{proof}

\begin{enumerate}[label=\arabic*.,leftmargin=*]
\item The common invariant calculation in \cite[Proposition~7.26, pp.~61--62]{AC2026} first gives
\[
 E=c_1(L|_F)=kq.
\]
Because $L$ extends across the order-four filling, its fibre class belongs
to the specialization lattice
\[
 \pi_2^*H^2(S_2;\mathbb Z)
 =\{a\gamma\eta_2+bq:a\equiv b\pmod2\}
\]
of \cite[Lemma~A.4, pp.~99--100]{AC2026}, through the normalized cover of \CompendiumRef{I}{fam:src:A.1.2}.  Taking
coefficient independence in the displayed invariant basis forces
$a=0$ and $b=k$; hence $k$ is even.  Write $E=2nq$.

\item Contraction with the two affine translation vectors is an exact integral
calculation:
\[
 \iota_{v_1}(2q)=4\eta_1,
 \qquad
 \iota_{v_2}(2q)=-6\eta_2.
\tag{W.42}\hypertarget{tag-03H9}{}\label{wbII:eq:W:5.17r}
\]
Indeed $v_1=(1,2,-4,0)^t$, $v_2=(-1,-3,3,0)^t$ in the basis dual to
$(\gamma,u,w,\delta)$, while
$q=uw+6\gamma\delta$.  In the global coinvariants,
\[
 [\eta_1]=3[\delta],\qquad [\eta_2]=2[\delta].
\]
Applying \cref{wbII:lem:picard-finite-norm} with $m_1=3$, $m_2=4$ therefore
gives
\[
 \bar d_1=4n[\delta],\qquad
 \bar d_2=-3n[\delta].
\tag{W.43}\hypertarget{tag-03HA}{}\label{wbII:eq:W:5.17s}
\]
Combine the balance relation \eqref{wbII:eq:W:5.17i} with the cusp vanishing \eqref{wbII:eq:W:5.17n}:
\[
 0=\bar d_1+\bar d_2+\bar d_0=n[\delta].
\]
The coinvariant group is torsion-free, so $n=0$ and hence $E=0$.

\item Finally \eqref{wbII:eq:W:5.17o} and the injectivity of the local integral specialization in
\cite[Theorem~B.1, pp.~101--105]{AC2026} imply
$c_1(L|_W)=0$.  Pulling back to the normalization gives
$c_1(\eta^*i^*L)=0$.  The injectivity of
$c_1:\operatorname{Pic}(\mathrm{dP}_6)\to
H^2(\mathrm{dP}_6;\mathbb Z)$ proved in
\cref{wbII:prop:pic0-pullback} yields the second assertion of \eqref{wbII:eq:W:5.17q}.

\end{enumerate}
\end{proof}



\subsection{A local analytic obstruction to a cusp-only Picard argument}
\hypertarget{tag-03IU}{}\label{wbII:subsec:local-cusp-picard-obstruction}

The relative Picard space of the cusp family cannot by itself force the
normalised central restriction to vanish.  The obstruction is visible
directly in the toric model of \CompendiumRef{I}{wb:thm:candidate-cusp-geometry}.  We give
the complete factor-of-automorphy calculation, including the descent
cocycle and the sign of the generic Chern class.

Retain the cusp notation of
\CompendiumRef{I}{wb:prop:period-input-audit}; \CompendiumRef{I}{wb:thm:candidate-cusp-geometry}:
\[
 B_0=\begin{pmatrix}0&1\\-1&0\end{pmatrix},
 \qquad
 Z(s)=sB_0+C(t),
 \qquad
 C(t)=\begin{pmatrix}6\mu(t)&h(t)\\ b(t)-h(t)&\mu(t)\end{pmatrix},
 \qquad t=e^{2\pi i s},
\tag{W.61}\hypertarget{tag-03HB}{}\label{wbII:eq:W:5.17t}
\]
\[
 c_{\lambda}(t)=\exp(2\pi iC(t)\lambda),
 \qquad
 \exp(2\pi iZ(s)\lambda)=t^{B_0\lambda}c_\lambda(t)
 \quad(\lambda\in\bar\Lambda\simeq\mathbb Z^2),
\tag{W.62}\hypertarget{tag-03HC}{}\label{wbII:eq:W:5.17u}
\]
and
\[
 \Phi_\lambda(x,t)=(t^{B_0\lambda}x,t),
 \qquad
 \Psi_\lambda(p)=(c_\lambda(t(p)),1)\cdot\Phi_\lambda(p),
 \qquad
 \Psi_\lambda\Psi_\mu=\Psi_{\lambda+\mu}.
\tag{W.63}\hypertarget{tag-03HD}{}\label{wbII:eq:W:5.17v}
\]
Write \(E_v\), \(v\in\mathbb Z^2\), for the toric prime divisors of
\(Y_\epsilon\), so that \(E_0\simeq\mathrm{dP}_6\), and let
\[
 \pi:Y_\epsilon\longrightarrow N_0=Y_\epsilon/\bar\Lambda,
 \qquad
 \nu=\pi|_{E_0}:E_0\longrightarrow W
\]
be the quotient and the normalization of the central fibre.
Put
\[
 N'=\mathbb Z^2\oplus\mathbb Z,
 \qquad
 M'=\operatorname{Hom}(N',\mathbb Z).
\]
For \(m=(\ell,k)\in M'=\mathbb Z^2\oplus\mathbb Z\), fix the character
convention
\[
 \chi^m(x,t)=x^\ell t^k,
 \qquad
 \operatorname{div}(\chi^m)=
 \sum_{v\in\mathbb Z^2}\langle m,(v,1)\rangle E_v.
\tag{W.64}\hypertarget{tag-03HE}{}\label{wbII:eq:W:5.17w}
\]

\begin{proposition}[An explicit nonzero local cusp restriction]
\hypertarget{tag-03IN}{}\label{wbII:prop:local-cusp-picard-obstruction}
\StatementRecord{wbII:prop:local-cusp-picard-obstruction}{proved}{03IN}{\SeriesRef{wbII:eq:W:5.17t}, \SeriesRef{wbII:eq:W:5.17ad}, \SeriesRef{wbII:eq:W:5.17u}, \SeriesRef{wbII:eq:W:5.17v}, \SeriesRef{wbII:eq:W:5.17ai}, \SeriesRef{wbII:eq:W:5.17ag}, \SeriesRef{wbII:eq:W:5.17aj}, \SeriesRef{wbII:eq:W:5.17ak}, \SeriesRef{wbII:eq:W:5.17al}, \SeriesRef{wbII:eq:W:5.17am}, \SeriesRef{wbII:eq:W:5.17an}, \SeriesRef{wbII:eq:W:5.17ao}, \SeriesRef{wbII:eq:W:5.17ae}, \SeriesRef{wbII:eq:W:5.17ax}}

Let
\[
 H=\begin{pmatrix}1&0\\0&-6\end{pmatrix},
 \qquad
 P=B_0^{\mathsf t}H=\begin{pmatrix}0&6\\1&0\end{pmatrix},
 \qquad
 Q(v)=v^{\mathsf t}Hv=v_1^2-6v_2^2.
\tag{W.65}\hypertarget{tag-03HF}{}\label{wbII:eq:W:5.17x}
\]
Then
\[
 \mathcal D_Q=\sum_{v\in\mathbb Z^2}Q(v)E_v
\tag{W.66}\hypertarget{tag-03HG}{}\label{wbII:eq:W:5.17y}
\]
is a locally finite Cartier divisor on \(Y_\epsilon\).
The line bundle \(\mathcal O_{Y_\epsilon}(\mathcal D_Q)\) has an explicit
\((\bar\Lambda,\Psi)\)-equivariant structure and descends to a holomorphic
line bundle \(\mathcal L_Q\) on \(N_0\).  If
\[
 (a_1,a_2,a_3,a_4,a_5,a_6)
 =(e_1,e_2,e_2-e_1,-e_1,-e_2,e_1-e_2),
 \qquad C_i=E_0\cap E_{a_i},
\]
then
\[
 \nu^*(\mathcal L_Q|_W)
 \simeq
 \mathcal O_{E_0}\!\left(\sum_{i=1}^6Q(a_i)C_i\right),
\qquad
 \left(\deg_{C_i}\nu^*(\mathcal L_Q|_W)\right)_{i=1}^6
 =(-12,2,0,-12,2,0).
\tag{W.67}\hypertarget{tag-03HH}{}\label{wbII:eq:W:5.17z}
\]
In particular \(\nu^*(\mathcal L_Q|_W)\) is nontrivial.

On every nonzero fibre
\[
 F_t=\mathbb C^2/(Z(s)\mathbb Z^2+\mathbb Z^2),
\]
the first Chern class of \(\mathcal L_Q|_{F_t}\), in the ordered lattice
basis
\((\widehat\gamma,\widehat u,\widehat w,\widehat\delta)\) indexing the
columns of \([Z\mid I_2]\), is exactly
\[
 c_1(\mathcal L_Q|_{F_t})=-2q,
 \qquad
 q=u\wedge w+6\gamma\wedge\delta.
\tag{W.68}\hypertarget{tag-03FW}{}\label{wbII:eq:W:5.17aa}
\]
\end{proposition}



\begin{proof}
\begin{enumerate}[label=\arabic*.,leftmargin=*]
\item \emph{Local finiteness and Cartier status.}
The fan of \(Y\) consists of the cones over the cells of the locally finite
\(A_2\)-triangulation at height one.  If
\(U_\sigma\simeq\mathbb C^3\) is the affine chart belonging to a triangle
\(\sigma=\{v_0,v_1,v_2\}\), then \(U_\sigma\) meets precisely
\(E_{v_0},E_{v_1},E_{v_2}\).  Indeed, by the orbit--cone
correspondence, \(U_\sigma\cap E_v\ne\varnothing\) exactly when
\(\mathbb R_{\geq0}(v,1)\) is a face of the cone over \(\sigma\).
Thus \(\{E_v\}_{v\in\mathbb Z^2}\) is locally finite and
\(\mathcal D_Q\) is a well-defined locally finite Weil divisor.

Every triangle of the fan is unimodular, as in
\CompendiumRef{I}{mapsI:src:A.1.2}.  If \(z_0,z_1,z_2\) are the
characters dual to
\((v_0,1),(v_1,1),(v_2,1)\), then
\[
 \mathcal D_Q|_{U_\sigma}
 =\operatorname{div}\!
   \left(z_0^{Q(v_0)}z_1^{Q(v_1)}z_2^{Q(v_2)}\right).
\tag{W.69}\hypertarget{tag-03FX}{}\label{wbII:eq:W:5.17ab}
\]
Negative exponents merely make the displayed monomial meromorphic, as
permitted for a local equation of a Cartier divisor.  Hence
\(\mathcal D_Q\) is Cartier on every chart and therefore on
\(Y_\epsilon\).

\item \emph{The signed divisor pullback identity.}
For \(\lambda\in\bar\Lambda=\mathbb Z^2\), set
\[
 b_\lambda=B_0\lambda,
 \qquad
 \ell_\lambda=2Hb_\lambda,
 \qquad
 k_\lambda=Q(b_\lambda),
 \qquad
 m_\lambda=(\ell_\lambda,k_\lambda)\in M'.
\tag{W.70}\hypertarget{tag-03FY}{}\label{wbII:eq:W:5.17ac}
\]
The fan automorphism sends \(E_v\) to \(E_{v+b_\lambda}\).
Consequently \(\Phi_\lambda^*E_w=E_{w-b_\lambda}\), and
\begin{align*}
 \Phi_\lambda^*\mathcal D_Q-\mathcal D_Q
 &=\sum_{v\in\mathbb Z^2}
   \bigl(Q(v+b_\lambda)-Q(v)\bigr)E_v\\
 &=\sum_{v\in\mathbb Z^2}
   \bigl(2v^{\mathsf t}Hb_\lambda+
         b_\lambda^{\mathsf t}Hb_\lambda\bigr)E_v
 =\operatorname{div}(\chi^{m_\lambda}).
\end{align*}
The \(t\)-dependent horizontal torus translation
\(p\mapsto(c_\lambda(t(p)),1)\cdot p\) preserves every
torus-invariant prime divisor \(E_v\).  Hence
\[
 \boxed{\ 
 \Psi_\lambda^*\mathcal D_Q
 =\mathcal D_Q+\operatorname{div}(\chi^{m_\lambda})\ }.
\tag{W.71}\hypertarget{tag-03FZ}{}\label{wbII:eq:W:5.17ad}
\]

\item \emph{The complete \(\Psi\)-cocycle.}
The matrices in \eqref{wbII:eq:W:5.17t} and \eqref{wbII:eq:W:5.17x} give
\[
 PC(t)=
 \begin{pmatrix}
 6(b(t)-h(t))&6\mu(t)\\
 6\mu(t)&h(t)
 \end{pmatrix},
\tag{W.72}\hypertarget{tag-03G0}{}\label{wbII:eq:W:5.17ae}
\]
which is symmetric.  Thus
\[
 A(t):=2PC(t)
\tag{W.73}\hypertarget{tag-03G1}{}\label{wbII:eq:W:5.17af}
\]
is a symmetric holomorphic matrix on the full cusp disc.  Define the
holomorphic units \(h_\lambda\) and meromorphic functions \(f_\lambda\) by
\[
 h_\lambda(t)=
 \exp\!\left(\pi i\lambda^{\mathsf t}A(t)\lambda\right),
 \qquad
 f_\lambda=h_\lambda(t)\chi^{m_\lambda}.
\tag{W.74}\hypertarget{tag-03G2}{}\label{wbII:eq:W:5.17ag}
\]
Equation \eqref{wbII:eq:W:5.17ad} gives
\[
 \operatorname{div}(f_\lambda)
 =\Psi_\lambda^*\mathcal D_Q-\mathcal D_Q.
\tag{W.75}\hypertarget{tag-03G3}{}\label{wbII:eq:W:5.17ah}
\]

The cocharacter map underlying \(\Phi_\mu\) is
\[
 \varphi_\mu(y,y_3)=(y+y_3b_\mu,y_3),
\]
so its pullback on characters is
\[
 \varphi_\mu^*(\ell,k)
 =(\ell,k+\ell^{\mathsf t}b_\mu).
\tag{W.76}\hypertarget{tag-03G4}{}\label{wbII:eq:W:5.17ai}
\]
The quadratic identity for \(Q\) gives
\[
 m_{\lambda+\mu}=m_\mu+\varphi_\mu^*m_\lambda.
\tag{W.77}\hypertarget{tag-03G5}{}\label{wbII:eq:W:5.17aj}
\]
Indeed, the first coordinates add, while the last coordinate on the
right is
\[
 Q(b_\mu)+Q(b_\lambda)+2b_\lambda^{\mathsf t}Hb_\mu
 =Q(b_\lambda+b_\mu).
\]
Using \eqref{wbII:eq:W:5.17u}, \eqref{wbII:eq:W:5.17v}, and \eqref{wbII:eq:W:5.17ai}, one obtains
\begin{align*}
 \Psi_\mu^*\chi^{m_\lambda}
 &=\chi^{m_\lambda}(c_\mu(t),1)
   \chi^{\varphi_\mu^*m_\lambda}\\
 &=\exp\!\left(2\pi i\ell_\lambda^{\mathsf t}C(t)\mu\right)
   \chi^{\varphi_\mu^*m_\lambda}\\
 &=\exp\!\left(2\pi i\lambda^{\mathsf t}A(t)\mu\right)
   \chi^{\varphi_\mu^*m_\lambda}.
\tag{W.78}\hypertarget{tag-03G6}{}\label{wbII:eq:W:5.17ak}
\end{align*}
The last equality is exact because
\[
 \ell_\lambda^{\mathsf t}C\mu
 =(2HB_0\lambda)^{\mathsf t}C\mu
 =2\lambda^{\mathsf t}B_0^{\mathsf t}HC\mu
 =2\lambda^{\mathsf t}PC\mu
 =\lambda^{\mathsf t}A\mu.
\]
Since \(A=A^{\mathsf t}\), \eqref{wbII:eq:W:5.17ag}, \eqref{wbII:eq:W:5.17aj}, and \eqref{wbII:eq:W:5.17ak} imply
\begin{align*}
 f_\mu\,\Psi_\mu^*f_\lambda
 &=h_\mu h_\lambda
   \exp(2\pi i\lambda^{\mathsf t}A\mu)
   \chi^{m_\mu+\varphi_\mu^*m_\lambda}\\
 &=\exp\!\left(
   \pi i(\lambda+\mu)^{\mathsf t}A(\lambda+\mu)\right)
   \chi^{m_{\lambda+\mu}}
 =f_{\lambda+\mu}.
\tag{W.79}\hypertarget{tag-03G7}{}\label{wbII:eq:W:5.17al}
\end{align*}

Use the divisor-sheaf convention
\[
 \mathcal O(D)(U)=
 \{g\in\mathcal M_Y(U):
   \operatorname{div}(g)+D|_U\geq0\}.
\]
If \(D'=D+\operatorname{div}(f)\), multiplication by \(f\) maps
\(\mathcal O(D')\) isomorphically to \(\mathcal O(D)\).  Thus
\eqref{wbII:eq:W:5.17ad} and \eqref{wbII:eq:W:5.17ag} define
\[
 \theta_\lambda:
 \Psi_\lambda^*\mathcal O(\mathcal D_Q)
 \xrightarrow{\ \times f_\lambda\ }
 \mathcal O(\mathcal D_Q).
\tag{W.80}\hypertarget{tag-03G8}{}\label{wbII:eq:W:5.17am}
\]
Because
\((\Psi_\lambda\Psi_\mu)^*=\Psi_\mu^*\Psi_\lambda^*\),
equation \eqref{wbII:eq:W:5.17al} is exactly
\[
 \theta_{\lambda+\mu}
 =\theta_\mu\circ\Psi_\mu^*\theta_\lambda.
\tag{W.81}\hypertarget{tag-03G9}{}\label{wbII:eq:W:5.17an}
\]

\item \emph{Analytic descent.}
The cusp action of \CompendiumRef{I}{pub:I:eq:PD11}, proved in
\CompendiumRef{I}{fam:src:A.3.1}, is
free and properly discontinuous.  Every point therefore has an open
neighbourhood \(U\) whose distinct translates are disjoint and on which
\(\pi|_U\) is a biholomorphism onto \(\pi(U)\).  Transport
\(\mathcal O(\mathcal D_Q)|_U\) to \(\pi(U)\), and on overlaps arising
from translates glue by \eqref{wbII:eq:W:5.17am}.  Equation \eqref{wbII:eq:W:5.17an} is precisely the
triple-overlap condition.  The resulting holomorphic line bundle
\(\mathcal L_Q\) on \(N_0\) satisfies
\[
 \pi^*\mathcal L_Q
 \simeq\mathcal O_{Y_\epsilon}(\mathcal D_Q).
\tag{W.82}\hypertarget{tag-03GA}{}\label{wbII:eq:W:5.17ao}
\]

\item \emph{The normalized central restriction.}
Let \(j_0:E_0\hookrightarrow Y_\epsilon\).  Since
\(\nu=\pi\circ j_0\), \eqref{wbII:eq:W:5.17ao} gives
\[
 \nu^*(\mathcal L_Q|_W)
 =j_0^*\pi^*\mathcal L_Q
 \simeq\mathcal O_{E_0}(\mathcal D_Q|_{E_0}).
\tag{W.83}\hypertarget{tag-03GB}{}\label{wbII:eq:W:5.17ap}
\]
The only \(E_v\), \(v\ne0\), meeting \(E_0\) are those indexed by the
six \(a_i\) in the statement, and the intersections are the six toric
\((-1)\)-curves \(C_i\).  Since \(Q(0)=0\), there is no
\(E_0|_{E_0}\) term, so
\[
 \mathcal D_Q|_{E_0}
 =\sum_{i=1}^6Q(a_i)C_i,
 \qquad
 (Q(a_1),\ldots,Q(a_6))=(1,-6,-5,1,-6,-5).
\tag{W.84}\hypertarget{tag-03GC}{}\label{wbII:eq:W:5.17aq}
\]
The hexagonal intersection table is
\[
 C_i^2=-1,
 \qquad
 C_i\cdot C_{i\pm1}=1,
 \qquad
 C_i\cdot C_j=0
 \quad(j\notin\{i,i\pm1\}),
\tag{W.85}\hypertarget{tag-03GD}{}\label{wbII:eq:W:5.17ar}
\]
with indices modulo six.  Therefore
\[
 \deg_{C_i}\mathcal O_{E_0}
 \!\left(\sum_jQ(a_j)C_j\right)
 =-Q(a_i)+Q(a_{i-1})+Q(a_{i+1}),
\]
and direct substitution gives
\[
 \begin{array}{c|rrrrrr}
 i&1&2&3&4&5&6\\ \hline
 -Q(a_i)+Q(a_{i-1})+Q(a_{i+1})
 &-12&2&0&-12&2&0.
 \end{array}
\tag{W.86}\hypertarget{tag-03GE}{}\label{wbII:eq:W:5.17as}
\]
The degree on \(C_1=C_{e_1}\) is \(-12\), proving nontriviality.

\item \emph{The generic-fibre factor and its exact Chern class.}
Fix \(t\ne0\), choose \(s\) with \(t=e^{2\pi i s}\), and write
\(x=e^{2\pi i\zeta}\), \(\zeta\in\mathbb C^2\).
Equation \eqref{wbII:eq:W:5.17u} identifies the deck transformation indexed by
\(\lambda\in\mathbb Z^2\) with
\[
 \zeta\longmapsto\zeta+Z(s)\lambda
\]
modulo the horizontal translations
\(\zeta\mapsto\zeta+n\), \(n\in\mathbb Z^2\).  The matrix identities
\[
 HB_0=P^{\mathsf t},
 \qquad
 B_0^{\mathsf t}HB_0=PB_0
\tag{W.87}\hypertarget{tag-03GF}{}\label{wbII:eq:W:5.17at}
\]
give
\[
 \chi^{m_\lambda}(x,t)
 =x^{2P^{\mathsf t}\lambda}
  t^{\lambda^{\mathsf t}PB_0\lambda},
\qquad
 h_\lambda(t)
 =\exp(2\pi i\lambda^{\mathsf t}PC(t)\lambda).
\]
Since \(Z=sB_0+C\), it follows that
\[
 f_\lambda(\zeta)
 =\exp\!\left(
 2\pi i\lambda^{\mathsf t}PZ\lambda+
 4\pi i\lambda^{\mathsf t}P\zeta\right).
\tag{W.88}\hypertarget{tag-03GG}{}\label{wbII:eq:W:5.17au}
\]
The matrix \(PZ=sPB_0+PC\) is symmetric:
\[
 PB_0=\operatorname{diag}(-6,1),
\]
and \(PC\) is symmetric by \eqref{wbII:eq:W:5.17ae}.

Under the descent convention \eqref{wbII:eq:W:5.17am}, a pulled-back local section
written as \(\varphi(\zeta)\) in the canonical frame satisfies
\[
 \varphi(\zeta+Z\lambda)
 =f_\lambda(\zeta)^{-1}\varphi(\zeta).
\]
Thus the factor of automorphy for the lattice translation
\(\zeta\mapsto\zeta+Z\lambda+n\) is
\[
 a_{(\lambda,n)}(\zeta)
 =\exp(2\pi i g_{(\lambda,n)}(\zeta)),
\qquad
 g_{(\lambda,n)}(\zeta)
 =-\lambda^{\mathsf t}PZ\lambda
  -2\lambda^{\mathsf t}P\zeta.
\tag{W.89}\hypertarget{tag-03GH}{}\label{wbII:eq:W:5.17av}
\]
This is independent of \(n\) as a factor of automorphy: replacing
\(\zeta\) by \(\zeta+n\) changes \(g\) by the integer
\(-2\lambda^{\mathsf t}Pn\).

For \(l=(\lambda,n)\), \(l'=(\mu,m)\), the exponential connecting
cocycle is
\begin{align*}
 c(l,l')
 &=g_l(\zeta+Z\mu+m)+g_{l'}(\zeta)-g_{l+l'}(\zeta)\\
 &=-2\lambda^{\mathsf t}Pm\in\mathbb Z,
\tag{W.90}\hypertarget{tag-03GI}{}\label{wbII:eq:W:5.17aw}
\end{align*}
where the symmetry of \(PZ\) cancels the two mixed \(Z\)-terms.
The corresponding alternating integral form is
\[
 E(l,l')=c(l,l')-c(l',l)
 =-2\lambda^{\mathsf t}Pm+2\mu^{\mathsf t}Pn.
\tag{W.91}\hypertarget{tag-03GJ}{}\label{wbII:eq:W:5.17ax}
\]
This is \(c_1(\mathcal L_Q|_{F_t})\).  In the ordered column basis of
\([Z\mid I_2]\), the cross-block of the class
\(q=u\wedge w+6\gamma\wedge\delta\) from
\((\widehat\gamma,\widehat u)\) to
\((\widehat w,\widehat\delta)\) is
\[
 \begin{pmatrix}
 q(\widehat\gamma,\widehat w)&
 q(\widehat\gamma,\widehat\delta)\\
 q(\widehat u,\widehat w)&
 q(\widehat u,\widehat\delta)
 \end{pmatrix}
 =\begin{pmatrix}0&6\\1&0\end{pmatrix}=P.
\tag{W.92}\hypertarget{tag-03GK}{}\label{wbII:eq:W:5.17ay}
\]
Hence \eqref{wbII:eq:W:5.17ax} is \(-2q\).

\end{enumerate}
\end{proof}



\paragraph{the exact obstruction to the purely analytic route}

\hypertarget{tag-03IQ}{}\label{wbII:rem:local-cusp-picard-obstruction}
The proposition proves that the local homomorphism
\[
 \operatorname{Pic}(N_0)\longrightarrow\operatorname{Pic}(E_0),
 \qquad
 \mathcal L\longmapsto\nu^*(\mathcal L|_W),
\]
is not zero.  Thus the toric cusp geometry, conductor descent on \(W\),
or local relative-Picard representability cannot alone prove the desired
global vanishing
\[
 \nu^*i^*:\operatorname{Pic}(X)\longrightarrow
 \operatorname{Pic}(E_0).
\]
There is no contradiction with
\cref{wbII:prop:leray-edge-restriction}, \cref{wbII:prop:direct-boundary-picard}: the local
bundle above has nonzero generic class \(-2q\), whereas the two global
proofs use the order-three and order-four fillings and their common
boundary relation to exclude every nonzero multiple of \(q\).

More generally, Bingener proves that if a holomorphic map is proper,
flat, and cohomologically flat in dimension zero, then its relative
Picard functor is represented by a \emph{not necessarily separated}
complex space
\cite[Introduction, p.~318; Satz~6.2, pp.~337--338]{Bingener1980}.
The same introduction explains that the exponential-sequence quotient
works directly under topological local triviality and cohomological
flatness in dimension one, but that passage to the quotient fails in
general.  The cusp map of \CompendiumRef{I}{wb:thm:candidate-cusp-geometry} is not
topologically locally trivial across \(0\).  Bingener's theorem therefore
supplies representability but neither separatedness nor the
component-specialization assertion needed here.

Consequently a purely analytic completion must add global compatibility
with both finite fillings and must control the boundary support function;
neither relative-Picard representability nor conductor descent supplies
those statements.  The later calculation in
\cref{wbII:prop:analytic-ah-boundary}, \cref{wbII:lem:analytic-support-specialization}, \cref{wbII:thm:analytic-all-picard} now supplies them: it first forces the generic
Appell--Humbert form to vanish using the two finite fillings, and then
proves that its mixed block is the polarization of the three normalized
central degrees.  Thus this proposition remains the sharp cusp-local
obstruction, while the global purely analytic route is complete.

\Uses{the hypotheses, notation and definitions stated here; \cref{wbII:prop:leray-edge-restriction}, \cref{wbII:prop:direct-boundary-picard}; \cite[Introduction, p.~318; Satz~6.2, pp.~337--338]{Bingener1980}; \cref{wbII:prop:analytic-ah-boundary}, \cref{wbII:lem:analytic-support-specialization}, \cref{wbII:thm:analytic-all-picard}}



\subsection{An analytic completion of the all-Picard certificate}
\hypertarget{tag-03IT}{}\label{wbII:subsec:analytic-all-picard}

The local example in \cref{wbII:prop:local-cusp-picard-obstruction} shows that
the cusp neighbourhood alone cannot annihilate normalized central
restrictions.  We now close the stronger analytic route by retaining the
order-three and order-four fillings.  The proof has two independent parts.
First, logarithms of the actual \(\mathcal O^*\)-factors give the boundary
cocycle and its two finite-filling norm identities.  Second, a
standard-family support-function calculation reads every possible
central restriction directly from the mixed Appell--Humbert block on a
nearby torus.  No global Leray differential, nearby-cycle specialization,
or Chern-class evaluation on a cusp bounding chain occurs.

\begin{proposition}[Analytic Appell--Humbert boundary cocycle]
\hypertarget{tag-03IH}{}\label{wbII:prop:analytic-ah-boundary}
\StatementRecord{wbII:prop:analytic-ah-boundary}{proved}{03IH}{\SeriesRef{wbII:eq:W:5.17-A8}, \SeriesRef{wbII:eq:W:5.17av}, \SeriesRef{wbII:eq:W:5.17ax}, \SeriesRef{wbII:eq:W:5.17-A12}, \SeriesRef{wbII:eq:W:5.17-A11}, \SeriesRef{wbII:eq:W:5.17-A13}, \SeriesRef{wbII:eq:W:5.17-A14a}, \SeriesRef{wbII:eq:W:5.17-A14c}, \SeriesRef{wbII:eq:W:5.17-A14}, \SeriesRef{wbII:eq:W:5.17-A14d}, \SeriesRef{wbII:eq:W:5.17-A14e}, \SeriesRef{wbII:eq:W:5.17-A15}, \SeriesRef{wbII:eq:W:5.17-A17}, \SeriesRef{wbII:eq:W:5.17-A19a}, \SeriesRef{wbII:eq:W:5.17-A16}, \SeriesRef{wbII:eq:W:5.17-A9}, \SeriesRef{wbII:eq:W:5.17-A20}; \cite[equations~(3.3)--(3.10), pp.~3--7]{SchnellLecture3}; \cite[\S2, fundamental theorem~B, p.~80]{Cartan1957}}

Let \(P\subset B^\circ\) be the compact pair-of-pants core,
orient its boundary meridians clockwise, and retain the zero-section lifts
of \CompendiumRef{I}{inv:B3.1}.  Thus
\[
 \rho_1\rho_2\rho_0=1,
 \qquad
 \widehat\rho_1\widehat\rho_2\widetilde a_0=1,
\tag{W.93}\hypertarget{tag-03EF}{}\label{wbII:eq:W:5.17-A1}
\]
and the second relation has no terminal vertical translation.  For
\(L\in\operatorname{Pic}(X)\), let
\[
 E=c_1(L|_F)\in\operatorname{Alt}^2(\Lambda,\mathbb Z)
   =\bigwedge^2V
\]
be its fibrewise Appell--Humbert form.  Put
\(V_{\mathbb R}=\operatorname{Hom}(\Lambda,\mathbb R)\), and use
\[
 A_g=(T_g^{-1})^{\mathsf t},
 \qquad
 \langle T_g\alpha,\lambda\rangle
 =\langle\alpha,A_g^{-1}\lambda\rangle.
\tag{W.94}\hypertarget{tag-03EY}{}\label{wbII:eq:W:5.17-A2}
\]
Here $A_g$ is the linear part of the left deck action and $T_g=A_g^{-\mathsf t}$
is its coefficient action.  A clockwise loop ends at the deck translate $g$;
its endpoint return $g^{-1}$ acts on homology by $A_g^{-1}$ and on
cohomology by $A_g^{\mathsf t}=T_g^{-1}$.  The cocycle below uses $T_g$.
There is a crossed homomorphism, unique up to a principal crossed
homomorphism,
\[
 d:\pi_1(P)\longrightarrow V_{\mathbb R},
 \qquad
 d(gh)=d(g)+T_gd(h),
 \qquad
 d(g)\sim d(g)+(T_g-I)a.
\tag{W.95}\hypertarget{tag-03F9}{}\label{wbII:eq:W:5.17-A3}
\]
If \(d_i=d(\rho_i)\), then
\[
 d_1+T_1d_2+T_1T_2d_0=0,
 \qquad
 \bar d_1+\bar d_2+\bar d_0=0
 \quad\text{in }V_G\otimes\mathbb R.
\tag{W.96}\hypertarget{tag-03FM}{}\label{wbII:eq:W:5.17-A4}
\]
For the finite filling of order \(m_j\), \(j=1,2\),
\[
 N_{T_j}d_j=\iota_{v_j}E,
 \qquad
 m_j\bar d_j=[\iota_{v_j}E],
 \qquad
 N_{T_j}=I+T_j+\cdots+T_j^{m_j-1},
\tag{W.97}\hypertarget{tag-03FN}{}\label{wbII:eq:W:5.17-A5}
\]
where \((\iota_vE)(\lambda)=E(v,\lambda)\).  At the cusp,
\[
 \langle d_0,\widehat\delta\rangle=0,
 \qquad
 \bar d_0=0\quad\text{in }V_G\otimes\mathbb R.
\tag{W.98}\hypertarget{tag-03FO}{}\label{wbII:eq:W:5.17-A6}
\]
\end{proposition}



\begin{proof}

\begin{enumerate}[label=\arabic*.,leftmargin=*]
\item Choose an open Riemann-surface neighbourhood \(P^{\rm an}\Subset B^\circ\)
which deformation retracts onto \(P\), and pull the period family to its
universal cover.  Before quotienting by the fibre lattice, the resulting
space is
\[
 \widetilde{\mathcal J}
 =\widetilde P^{\rm an}\times\mathbb C^2.
\]
It is Stein and contractible, so the pullback of \(L\) is holomorphically
trivial.  In a nowhere-zero frame, write the factor belonging to a deck
transformation \(\alpha\) as
\[
 a_\alpha\in\mathcal O^*(\widetilde{\mathcal J}),
 \qquad
 a_{\alpha\beta}(x)=a_\alpha(\beta x)a_\beta(x).
\tag{W.99}\hypertarget{tag-03FP}{}\label{wbII:eq:W:5.17-A7}
\]
The space is simply connected, hence \(a_\alpha=e^{2\pi i h_\alpha}\) for
a holomorphic logarithm.  Therefore
\[
 c(\alpha,\beta)
 =h_\alpha(\beta x)+h_\beta(x)-h_{\alpha\beta}(x)\in\mathbb Z
\tag{W.100}\hypertarget{tag-03FQ}{}\label{wbII:eq:W:5.17-A8}
\]
is independent of \(x\), and associativity makes \(c\) an integral group
two-cocycle.

\item Let \(\tau_\lambda\) denote fibre translation by
\(\Pi(z)\lambda\), \(\lambda\in\Lambda\), and let \(\widehat g\) be the
zero-section lift of a base deck transformation.  Period equivariance gives
\[
 \widehat g\,\tau_\lambda\,\widehat g^{-1}
 =\tau_{A_g\lambda}.
\tag{W.101}\hypertarget{tag-03FR}{}\label{wbII:eq:W:5.17-A9}
\]
On fibre translations, the commutator of \eqref{wbII:eq:W:5.17-A8} is
\[
 E(\lambda,\mu)
 =c(\tau_\lambda,\tau_\mu)-c(\tau_\mu,\tau_\lambda).
\tag{W.102}\hypertarget{tag-03EG}{}\label{wbII:eq:W:5.17-A10}
\]
The four oriented sides of the \((\lambda,\mu)\)-parallelogram identify
this integer with the exponential connecting class of \(L|_F\); this is
the Appell--Humbert convention used in
\eqref{wbII:eq:W:5.17av}--\eqref{wbII:eq:W:5.17ax}, and agrees with the standard factor-of-automorphy
calculation
\cite[equations~(3.3)--(3.10), pp.~3--7]{SchnellLecture3}.
Conjugating the commutator by \(\widehat g\) yields
\[
 E(A_g\lambda,A_g\mu)=E(\lambda,\mu).
\tag{W.103}\hypertarget{tag-03EH}{}\label{wbII:eq:W:5.17-A11}
\]

\item Push the logarithmic central extension defined by \(c\) out along
\(\mathbb Z\hookrightarrow\mathbb R\), and denote its central element of
weight \(r\) by \(\mathbf c(r)\).  Choose lifts \(U_\lambda\) of the fibre
translations satisfying the real Appell--Humbert normalization
\[
 U_\lambda U_\mu
 =\mathbf c\!\left(\frac12E(\lambda,\mu)\right)U_{\lambda+\mu}.
\tag{W.104}\hypertarget{tag-03EI}{}\label{wbII:eq:W:5.17-A12}
\]
To see that this normalization exists, start with arbitrary lifts.  The
difference between their real cocycle and \(E/2\) has zero commutator, so
it defines an abelian extension of the free abelian group \(\Lambda\) by
\(\mathbb R\).  Choose lifts of an integral basis and extend
multiplicatively; this splits that abelian extension and gives \eqref{wbII:eq:W:5.17-A12}.
If \(U'_\lambda\) is a second normalized system, then
\(U'_\lambda U_\lambda^{-1}\) is central, say
\(\mathbf c(\langle a,\lambda\rangle)\).  Comparing the two copies of
\eqref{wbII:eq:W:5.17-A12} proves additivity in \(\lambda\), so
\(a\in V_{\mathbb R}\).  Thus every two normalized systems differ by exactly
the character change used below.

\item For \(g\in\pi_1(P)\), choose a lift \(G_g\) of \(\widehat g\) and define
\(d(g)\in V_{\mathbb R}\) by
\[
 G_gU_{A_g^{-1}\lambda}G_g^{-1}
 =\mathbf c\bigl(\langle d(g),\lambda\rangle\bigr)U_\lambda.
\tag{W.105}\hypertarget{tag-03EJ}{}\label{wbII:eq:W:5.17-A13}
\]
The exponent is additive in \(\lambda\): conjugate \eqref{wbII:eq:W:5.17-A12} and use
\eqref{wbII:eq:W:5.17-A11}.  The lifts of \(gh\) and \(G_gG_h\) differ centrally, so
\[
\begin{aligned}
 \langle d(gh),\lambda\rangle
 &=\langle d(g),\lambda\rangle
   +\langle d(h),A_g^{-1}\lambda\rangle\\
 &=\langle d(g)+T_gd(h),\lambda\rangle.
\end{aligned}
\]
This proves the crossed-homomorphism law.  Replacing
\(U_\lambda\) by
\(\mathbf c(\langle a,\lambda\rangle)U_\lambda\) changes \eqref{wbII:eq:W:5.17-A13} by
\((T_g-I)a\).  Replacing \(G_g\) by a central multiple does not change
\eqref{wbII:eq:W:5.17-A13}, so this proves the stated independence and uniqueness.  Applying
the crossed law to
\eqref{wbII:eq:W:5.17-A1} proves the first identity in \eqref{wbII:eq:W:5.17-A4}; passage to coinvariants
proves the second.

\item For the order-\(m_j\) filling, put
\[
 \widetilde J_j=\Delta_j\times\mathbb C^2,\qquad
 r_j:\widetilde J_j\longrightarrow J'_j
\]
for the lattice quotient, whose deck translations are \(\tau_\lambda\).
On the vector cover define
\[
 Q_{j,\log}:=\widetilde g_j^{\,v_j/m_j},\qquad
 Q_{j,\mathrm{lin}}:=\widetilde g_j^{\,0}.
\]
They satisfy
\[
\begin{aligned}
 Q_{j,\log}\tau_\lambda Q_{j,\log}^{-1}
  &=\tau_{A_j\lambda},&
 (Q_{j,\log})^{m_j}&=\tau_{v_j},\\
 Q_{j,\mathrm{lin}}\tau_\lambda Q_{j,\mathrm{lin}}^{-1}
  &=\tau_{A_j\lambda},&
 (Q_{j,\mathrm{lin}})^{m_j}&=1.
\end{aligned}
\tag{W.106}\hypertarget{tag-03EK}{}\label{wbII:eq:W:5.17-A14}
\]
The positive sign is fixed by
\(s_j\circ g_j=e^{-2\pi i/m_j}s_j\) and \CompendiumRef{I}{pub:I:eq:PD7}: after exactly
\(m_j\) clockwise iterations, the logarithmic vector-cover lift translates
by \(+v_j\).

\item Let \(q_{j,\log},q_{j,\mathrm{lin}}\) be the descended automorphisms of \(J'_j\),
and put \(J_j^{\prime *}=J'_j|_{\Delta_j^*}\).  Both descended automorphisms have
order dividing \(m_j\), because \(\tau_{v_j}\) is trivial after the lattice
quotient; each covers \(g_j\), whose order is exactly \(m_j\), so each
descended order is exactly \(m_j\).
The branch-independent torus translation
\[
 h_j:J_j^{\prime *}\longrightarrow J_j^{\prime *}
\]
constructed in \CompendiumRef{I}{pub:I:eq:PD8} satisfies
\[
 h_jq_{j,\log}h_j^{-1}=q_{j,\mathrm{lin}}.
\]
Define
\[
\begin{aligned}
 p_{j,\log}:J_j^{\prime *}&\longrightarrow
  N_j^*:=J_j^{\prime *}/\langle q_{j,\log}\rangle,\\
 p_{j,\mathrm{lin}}:J_j^{\prime *}&\longrightarrow
 J_j^*:=J_j^{\prime *}/\langle q_{j,\mathrm{lin}}\rangle .
\end{aligned}
\]
The conjugacy induces a biholomorphism
\[
 \bar h_j:N_j^*\longrightarrow J_j^*,\qquad
 p_{j,\mathrm{lin}}\circ h_j=\bar h_j\circ p_{j,\log}.
\]
The gluing of \(X\) identifies
\(L|_{N_j^*}\simeq\bar h_j^*(L|_{J_j^*})\); hence the
line-bundle square gives
\[
 p_{j,\log}^*(L|_{N_j^*})
 \simeq h_j^*p_{j,\mathrm{lin}}^*(L|_{J_j^*}).
\tag{W.107}\hypertarget{tag-03EL}{}\label{wbII:eq:W:5.17-A14a}
\]

\item On the filled vector cover, let
\[
 P_{j,\log}:J'_j\longrightarrow N_j,\qquad
 \pi_j^{\mathrm{vec}}:=P_{j,\log}\circ r_j:
 \widetilde J_j\longrightarrow N_j .
\]
Here \(P_{j,\log}\) is the full quotient and
\(P_{j,\log}|_{J_j^{\prime *}}=p_{j,\log}\).
The deck group is
\[
 \Gamma_j^{\log}
 =\langle\tau_\lambda,Q_{j,\log}:\lambda\in\Lambda\rangle.
\]
The pullback \((\pi_j^{\mathrm{vec}})^*(L|_{N_j})\) is trivial on the
contractible Stein manifold \(\widetilde J_j\); its factors define an
integral central extension, which we push out along
\(\mathbb Z\hookrightarrow\mathbb R\).  Over \(\Delta_j^*\), the
line-bundle square \eqref{wbII:eq:W:5.17-A14a} identifies each fibre with the global factor
extension by a translation, so its commutator form is \(E\).  That form is
integer-valued and locally constant on the connected disc \(\Delta_j\);
therefore it is \(E\) on the filled cover as well.  Choose vertical lifts
\(U_\lambda^{(j)}\) normalized by
\[
 U_\lambda^{(j)}U_\mu^{(j)}
 =\mathbf c\!\left(\frac12E(\lambda,\mu)\right)
  U_{\lambda+\mu}^{(j)}.
\]
For \(\alpha\in\Gamma_j^{\log}\), define \(A_\alpha,T_\alpha\) by
\[
 \alpha\tau_\lambda\alpha^{-1}=\tau_{A_\alpha\lambda},\qquad
 T_\alpha=(A_\alpha^{-1})^{\mathsf t}.
\]
For a lift \(\widetilde G_\alpha\), define
\(\widetilde d_j(\alpha)\in V_{\mathbb R}\) through
\[
 \widetilde G_\alpha U_{A_\alpha^{-1}\lambda}^{(j)}
 \widetilde G_\alpha^{-1}
 =\mathbf c\!\left(\langle\widetilde d_j(\alpha),\lambda\rangle\right)
  U_\lambda^{(j)}.
\tag{W.108}\hypertarget{tag-03EM}{}\label{wbII:eq:W:5.17-A14c}
\]
Indeed, before naming \(\widetilde d_j(\alpha)\), the quotient of the two
sides of \eqref{wbII:eq:W:5.17-A14c} is central because both cover the same translation.
Conjugating the normalized product and using
\[
 E(A_\alpha\lambda,A_\alpha\mu)=E(\lambda,\mu)
\]
shows that its exponent \(r_\alpha(\lambda)\) satisfies
\[
 r_\alpha(\lambda+\mu)=r_\alpha(\lambda)+r_\alpha(\mu).
\]
Thus it is the pairing with a unique element of \(V_{\mathbb R}\), as
asserted.  Since \(A_{\alpha\beta}=A_\alpha A_\beta\), choose the lift of
\(\alpha\beta\) to be \(\widetilde G_\alpha\widetilde G_\beta\); any other
choice differs centrally.  Then
\[
\begin{aligned}
 &\widetilde G_\alpha\widetilde G_\beta
 U_{A_\beta^{-1}A_\alpha^{-1}\lambda}^{(j)}
 \widetilde G_\beta^{-1}\widetilde G_\alpha^{-1}\\
 &\quad=
 \mathbf c\!\left(
  \langle\widetilde d_j(\beta),A_\alpha^{-1}\lambda\rangle
  +\langle\widetilde d_j(\alpha),\lambda\rangle\right)
 U_\lambda^{(j)}.
\end{aligned}
\]
The defining dual action
\(\langle T_\alpha v,\lambda\rangle
=\langle v,A_\alpha^{-1}\lambda\rangle\) therefore proves the crossed law
\[
 \widetilde d_j(\alpha\beta)
 =\widetilde d_j(\alpha)+T_\alpha\widetilde d_j(\beta).
\tag{W.109}\hypertarget{tag-03EN}{}\label{wbII:eq:W:5.17-A14d}
\]

\item Put
\[
 \widetilde J_j^*:=\Delta_j^*\times\mathbb C^2,
 \qquad
 r_j^*:=r_j|_{\widetilde J_j^*}:
 \widetilde J_j^*\longrightarrow J_j^{\prime *}
\]
and define
\[
\begin{aligned}
 \Phi_{j,\log}
  &:=p_{j,\log}\circ r_j^*,\\
 \Phi_{j,\mathrm{glob}}
  &:=p_{j,\mathrm{lin}}\circ h_j\circ r_j^*
   =\bar h_j\circ\Phi_{j,\log}.
\end{aligned}
\]
Both maps are invariant under \(\Gamma_j^{\log}\).  This is immediate for
every \(\tau_\lambda\), while
\[
\begin{aligned}
 \Phi_{j,\log}Q_{j,\log}
 &=p_{j,\log}q_{j,\log}r_j^*
  =\Phi_{j,\log},\\
 \Phi_{j,\mathrm{glob}}Q_{j,\log}
 &=p_{j,\mathrm{lin}}h_jq_{j,\log}r_j^*
  =p_{j,\mathrm{lin}}q_{j,\mathrm{lin}}h_jr_j^*
  =\Phi_{j,\mathrm{glob}}.
\end{aligned}
\]
Thus \(\widetilde J_j^*\) is the same \(\Gamma_j^{\log}\)-cover of the two
biholomorphic quotients.  Pulling the gluing isomorphism to this cover gives
\[
 \Theta_j:
 \Phi_{j,\log}^*(L|_{N_j^*})
 \xrightarrow{\ \sim\ }
 \Phi_{j,\mathrm{glob}}^*(L|_{J_j^*}).
\]
For the comparison of logarithms, pass one further time to the universal
cover
\[
 \widehat J_j^*=\widetilde{\Delta_j^*}\times\mathbb C^2
 \longrightarrow \widetilde J_j^*.
\]
The deck group over the punctured quotient is
\[
 \widehat\Gamma_j=
 \langle\tau_\lambda,q_j\mid
 q_j\tau_\lambda q_j^{-1}=\tau_{A_j\lambda}\rangle.
\]
The map from this universal cover to the filled vector cover induces the
explicit homomorphism
\[
 \sigma_j:\widehat\Gamma_j\longrightarrow\Gamma_j^{\log},
 \qquad \tau_\lambda\longmapsto\tau_\lambda,
 \quad q_j\longmapsto Q_{j,\log}.
\]
Its kernel is generated by
\(q_j^{m_j}\tau_{-v_j}\): this element is one deck turn of
\(\widetilde{\Delta_j^*}\to\Delta_j^*\), while its image on the
filled vector cover is the identity by \eqref{wbII:eq:W:5.17-A14}.
Indeed every word is \(\tau_\lambda q_j^r\); its filled image can be
identity only when \(r=m_jk\), and then only when
\(\lambda=-kv_j\).  Thus the kernel assertion follows directly.

\item Let \(\mathscr E_j\) be the real central extension of
\(\Gamma_j^{\log}\) defined by the filled factors.  Pull it back along
\(\sigma_j\), explicitly as the fibre-product group
\[
 \sigma_j^*\mathscr E_j=
 \{(e,g):e\in\mathscr E_j,\ g\in\widehat\Gamma_j,
           \ p(e)=\sigma_j(g)\}.
\]
The global factor extension restricted to the based boundary group is
an extension of the same \(\widehat\Gamma_j\), with the generator
\(q_j\) identified with \(\widehat\rho_j\) through the fixed whisker.
The endpoint calculation below proves this precise identification.
Pull \(\Theta_j\) to \(\widehat J_j^*\).  This space is simply
connected, so the holomorphic unit comparing the two frames has a
holomorphic logarithm.  Substitution into \eqref{wbII:eq:W:5.17-A8}
changes the two integral cocycles by a coboundary.  Consequently it gives
an isomorphism from \(\sigma_j^*\mathscr E_j\) to the restricted global
extension, sending every \(\mathbf c(r)\) to the same
\(\mathbf c(r)\).  In the comparisons below,
\(U_\lambda^{(j)}\) means its paired lift
\((U_\lambda^{(j)},\tau_\lambda)\), and
\(\widetilde G_{Q_{j,\log}}\) means
\((\widetilde G_{Q_{j,\log}},q_j)\), when the isomorphism
\(\Theta_{j,*}\) is applied.  Their conjugation exponents remain the
ones in the filled extension.  
\item Let \(\gamma_j^J:[0,1]\to J_j^*\) be the zero-section image of the
unwhiskered clockwise boundary meridian, and let \(\omega_j\) be the fixed
zero-section whisker from the global basepoint to \(\gamma_j^J(0)\).  Thus
the based total-space loop retained in \eqref{wbII:eq:W:5.17-A1} is
\[
 \widehat\rho_j=\omega_j*\gamma_j^J*\omega_j^{-1}.
\]
Let $s_j$ be the root point selected by the lift of the fixed zero-section
whisker, and choose a logarithm branch $\alpha=\log s_j/(2\pi i)$.  For the vector-cover connector
$H_j(s,\zeta)=(s,\zeta+\alpha(s)\Pi(s)v_j)$, set
\[
 \widetilde x_j=(s_j,-\alpha(s_j)\Pi(s_j)v_j),\qquad
 x_j=r_j^*(\widetilde x_j),\qquad x_{j,\mathrm{lin}}=h_j(x_j).
\]
Thus $H_j(\widetilde x_j)=(s_j,0)$ and
$p_{j,\mathrm{lin}}(x_{j,\mathrm{lin}})=\gamma_j^J(0)$.  Let
\(\widetilde\gamma_{j,\mathrm{lin}}\) be the
\(p_{j,\mathrm{lin}}\)-lift of \(\gamma_j^J\) beginning at
\(x_{j,\mathrm{lin}}\).  For the chosen clockwise generator,
\[
 \widetilde\gamma_{j,\mathrm{lin}}(1)
 =q_{j,\mathrm{lin}}(x_{j,\mathrm{lin}}).
\]
Set
\[
 \gamma_j^N:=\bar h_j^{-1}\circ\gamma_j^J,
 \qquad
 \widetilde\gamma_{j,\log}
 :=h_j^{-1}\circ\widetilde\gamma_{j,\mathrm{lin}}.
\]
Then the quotient square gives the point-valued identities
\[
\begin{aligned}
 p_{j,\log}\circ\widetilde\gamma_{j,\log}
  &=\gamma_j^N,\\
 \widetilde\gamma_{j,\log}(0)&=x_j,\\
 \widetilde\gamma_{j,\log}(1)&=q_{j,\log}(x_j).
\end{aligned}
\]
Let $\widetilde\gamma_{j,\mathrm{vec}}$ be the $r_j^*$-lift of
$\widetilde\gamma_{j,\log}$ beginning at the selected $\widetilde x_j$.
For the clockwise lift, $\alpha_1=\alpha_0-1/m_j$.  Writing
$\zeta_0=-\alpha_0\Pi(s_j)v_j$, period equivariance and $A_jv_j=v_j$ give
\[
 \zeta_1=-(\alpha_0-1/m_j)\Pi(g_js_j)v_j
 =R_{g_j}(s_j)\zeta_0+\Pi(g_js_j)v_j/m_j.
\]
By \CompendiumRef{I}{pub:I:eq:PD7}, this gives
\[
 \widetilde\gamma_{j,\mathrm{vec}}(1)
 =Q_{j,\log}(\widetilde x_j),
\]
for this selected vector lift.
Fibre loops remain \(\tau_\lambda\), since \(h_j\) is a fibre translation.

\item Lift $\omega_j$ along the zero section on the linear cover and use
$H_j^{-1}$ at its local endpoint; this gives the selected
$\widetilde x_j$ above.  The factor of $L$ along this fixed lift identifies the
target factor over the local boundary segment with the globally based
factor.  Let \(G_{\rho_j}^{(j)}\) be the target lift of the local boundary
deck element whose whiskered image is the chosen global lift
\(G_{\rho_j}\).  Thus it covers \(\widehat\rho_j\) after whiskering and,
by \eqref{wbII:eq:W:5.17-A13}, satisfies
\[
 G_{\rho_j}^{(j)}U_{A_j^{-1}\lambda}(G_{\rho_j}^{(j)})^{-1}
 =\mathbf c(\langle d_j,\lambda\rangle)U_\lambda,
 \qquad d_j=d(\rho_j).
\]
Consequently, for some
\(a_j\in V_{\mathbb R}\) and \(b_j\in\mathbb R\),
\[
 \Theta_{j,*}(U_\lambda^{(j)})
 =\mathbf c(\langle a_j,\lambda\rangle)U_\lambda,
 \qquad
 \Theta_{j,*}(\widetilde G_{Q_{j,\log}})
 =\mathbf c(b_j)G_{\rho_j}^{(j)}.
\]
Here \(U_\lambda\) is the normalized global system in \eqref{wbII:eq:W:5.17-A12}, and
\(\mathbf c(b_j)\) disappears under conjugation.
Therefore
\[
\begin{aligned}
 \langle\widetilde d_j(Q_{j,\log}),\lambda\rangle
 +\langle a_j,\lambda\rangle
 &=\langle d_j,\lambda\rangle
   +\langle a_j,A_j^{-1}\lambda\rangle\\
 &=\langle d_j+T_ja_j,\lambda\rangle,
\end{aligned}
\]
so
\[
 \widetilde d_j(Q_{j,\log})=d_j+(T_j-I)a_j.
\tag{W.110}\hypertarget{tag-03EO}{}\label{wbII:eq:W:5.17-A14e}
\]

\item For a translation choose
\(\widetilde G_{\tau_v}=U_v^{(j)}\).  The filled normalization above gives
\[
 U_v^{(j)}U_\lambda^{(j)}(U_v^{(j)})^{-1}
 =\mathbf c(E(v,\lambda))U_\lambda^{(j)},
\]
and hence
\[
 \widetilde d_j(\tau_v)=\iota_vE,\qquad
 (\iota_vE)(\lambda)=E(v,\lambda).
\tag{W.111}\hypertarget{tag-03EP}{}\label{wbII:eq:W:5.17-A14f}
\]
Here
\[
 A_{Q_{j,\log}}=A_j,\qquad T_{Q_{j,\log}}=T_j.
\]
Since \(N_{T_j}(T_j-I)=T_j^{m_j}-I=0\), equations
\eqref{wbII:eq:W:5.17-A14}, \eqref{wbII:eq:W:5.17-A14d}, and \eqref{wbII:eq:W:5.17-A14e} give
\[
\begin{aligned}
 N_{T_j}d_j
 &=N_{T_j}\widetilde d_j(Q_{j,\log})\\
 &=\widetilde d_j((Q_{j,\log})^{m_j})\\
 &=\widetilde d_j(\tau_{v_j})\\
 &=\iota_{v_j}E.
\end{aligned}
\tag{W.112}\hypertarget{tag-03EQ}{}\label{wbII:eq:W:5.17-A15}
\]
Changing a line-bundle frame gives an isomorphic factor extension; changing
\(\widetilde G_\alpha\) by a central factor does not change conjugation; and
every second normalized vertical system differs by a character, changing
\(\widetilde d_j(\alpha)\) by \((T_\alpha-I)a\).  Finally,
\(\ell\mapsto\ell+2\pi in\) changes a vector-cover lift of \(h_j\) by
\[
 H_j\longmapsto H'_j=\tau_{nv_j}H_j.
\]
At the factor-extension level this means
\[
 \Theta'_{j,*}
 =\operatorname{Ad}(U_{nv_j})\circ\Theta_{j,*}.
\]
Consequently the comparison character changes by
\(n\iota_{v_j}E\), because
\[
 U_{nv_j}
 \bigl(\mathbf c(\langle a_j,\lambda\rangle)U_\lambda\bigr)
 U_{nv_j}^{-1}
 =
 \mathbf c(\langle a_j+n\iota_{v_j}E,\lambda\rangle)U_\lambda.
\]
Since \(A_jv_j=v_j\) and \(E\) is \(A_j\)-invariant,
\[
 \langle T_j(\iota_{v_j}E),\lambda\rangle
 =E(v_j,A_j^{-1}\lambda)
 =E(A_jv_j,\lambda)
 =E(v_j,\lambda),
\]
so this branch change leaves both \eqref{wbII:eq:W:5.17-A14e} and \eqref{wbII:eq:W:5.17-A15} unchanged.
The translation \(\tau_{nv_j}\) also commutes with
\(Q_{j,\mathrm{lin}}\) on the vector cover because \(A_jv_j=v_j\); on the
lattice quotient it is already the identity.  Hence it changes only the
central lift choice there.
The descended \(h_j,\bar h_j\), and the quotient square are branch
independent.  Thus \eqref{wbII:eq:W:5.17-A15} is choice-independent and proves \eqref{wbII:eq:W:5.17-A5},
including its sign.

\item Since
\(A_0\widehat\delta=\widehat\delta\), equation \eqref{wbII:eq:W:5.17-A13} says that
\[
 [G_0,U_{\widehat\delta}]
 =\mathbf c\bigl(\langle d_0,\widehat\delta\rangle\bigr).
\tag{W.113}\hypertarget{tag-03ER}{}\label{wbII:eq:W:5.17-A16}
\]
This logarithmic commutator is the clutching winding with the base loop
first.  Thus, if \(b\) and \(f\) denote the oriented base and fibre loops,
respectively,
\[
 \operatorname{wind}_{b\wedge f}(L)
  =\langle d_0,\widehat\delta\rangle,
 \qquad
 \operatorname{wind}_{\xi_{\widehat\delta}}(L)
 =-\langle d_0,\widehat\delta\rangle.
\tag{W.114}\hypertarget{tag-03ES}{}\label{wbII:eq:W:5.17-A17}
\]
The second equality uses the fibre-first convention
\(\xi_{\widehat\delta}=[S^1_{\widehat\delta}]\times[\rho_0]\), hence the
single sign change.
Use the unimodular cusp chart of \CompendiumRef{I}{mapsI:src:A.1.2}
\[
 (u,x_1,x_2)=\left(\frac{t}{x_1x_2},x_1,x_2\right),
 \qquad t=ux_1x_2.
\tag{W.115}\hypertarget{tag-03ET}{}\label{wbII:eq:W:5.17-A18}
\]
For \(0<\rho\ll1\), the boundary parametrization
\[
 (\theta,\phi)\longmapsto
 \bigl(\rho e^{i(\theta-\phi)},1,e^{i\phi}\bigr),
 \qquad t=\rho e^{i\theta},
\tag{W.116}\hypertarget{tag-03EU}{}\label{wbII:eq:W:5.17-A19}
\]
has base coordinate \(\theta\) and fibre translation
\(\widehat\delta\) in the \(\phi\)-direction.  The latter fixed-base loop
does \emph{not} itself bound the \(x_2\)-disc.  To identify the correct
meridian, put \(\alpha=\theta-\phi\), let \(a_u\) be the positively
oriented \(u\)-longitude, and let \(m_{x_2}\) be the positively oriented
\(x_2\)-meridian at fixed \(u\).  In the basis
$ (a_u,m_{x_2}) $, the counterclockwise base loop $b_+$ and the fibre loop are
\[
 b_+=a_u,
 \qquad
 f=m_{x_2}-a_u,
 \qquad
 b_+\wedge f=a_u\wedge m_{x_2}.
\tag{W.117}\hypertarget{tag-03EV}{}\label{wbII:eq:W:5.17-A19a}
\]
The determinant of this change of basis is $+1$.  The clockwise loop in
\eqref{wbII:eq:W:5.17-A17} is $b=-b_+$, so
$b\wedge f=-a_u\wedge m_{x_2}$.  The winding computed below is zero
in either orientation.  The loop $m_{x_2}$ bounds the
analytic disc
\[
 z\longmapsto(\rho e^{i\alpha_0},1,z),
 \qquad |z|\le1.
\tag{W.118}\hypertarget{tag-03EW}{}\label{wbII:eq:W:5.17-A19b}
\]

\item Choose a thin annulus \(A_u\) about \(|u|=\rho\), a disc \(D_{x_1}\)
about \(1\), and a disc \(D_{x_2}\) containing \(\{|x_2|\le1\}\), small
enough that
\[
 A_u\times D_{x_1}\times D_{x_2}
 \subset\{|ux_1x_2|<\epsilon\}.
\tag{W.119}\hypertarget{tag-03EX}{}\label{wbII:eq:W:5.17-A19c}
\]
Write \(U=A_u\times D_{x_1}\times D_{x_2}\) for this neighbourhood.
It is Stein and retracts onto \(S^1\), so \(H^2(U,\mathbb Z)=0\).
The exponential sequence is
\[
0\longrightarrow\mathbb Z\longrightarrow\mathcal O_U
 \xrightarrow{\,h\mapsto e^{2\pi i h}\,}\mathcal O_U^*
 \longrightarrow1.
\]
Its exactness follows from the existence of a holomorphic logarithm locally
for every nonvanishing holomorphic function, and from the integral kernel
of the exponential. Cartan's theorem~B gives
\(H^1(U,\mathcal O_U)=H^2(U,\mathcal O_U)=0\)
\cite[\S2, fundamental theorem~B, p.~80]{Cartan1957}.
The long cohomology sequence therefore identifies
\(\operatorname{Pic}(U)=H^1(U,\mathcal O_U^*)\) with
\(H^2(U,\mathbb Z)=0\). Pull \(L|_{N_0}\) back to
the cusp cover and choose one holomorphic frame on this entire
neighbourhood.  Compute the clutching cocycle with this extending frame on
both sides of a cut in the \(u\)-longitude.  Its transition unit on
\(D_{x_2}\) is then identically \(F=1\), so by \eqref{wbII:eq:W:5.17-A19a} its
base--fibre commutator winding is
\[
 \frac1{2\pi i}\int_{|x_2|=1}\frac{dF}{F}=0.
\tag{W.120}\hypertarget{tag-03EZ}{}\label{wbII:eq:W:5.17-A20}
\]
This computes the same alternating logarithmic cocycle as \eqref{wbII:eq:W:5.17-A16}.
Indeed, changing from the global Appell--Humbert frame to the extending
chart frame changes \(c\) by a group coboundary.  Put
\[
 g=\widehat g_0,
 \qquad
 h=\tau_{\widehat\delta}
\]
in the underlying deck group.  Equation \eqref{wbII:eq:W:5.17-A9} and
\(A_0\widehat\delta=\widehat\delta\) give \(gh=hg\).  If
\(c'=c+\delta\beta\) is the cocycle after a frame change, then
\[
 \bigl[c'(g,h)-c'(h,g)\bigr]
 -\bigl[c(g,h)-c(h,g)\bigr]
 =\beta(hg)-\beta(gh)=0.
\]
Thus the alternating cocycle is frame-invariant without assuming that its
central-extension lifts commute.  Equations \eqref{wbII:eq:W:5.17-A17}, \eqref{wbII:eq:W:5.17-A19a}, and
\eqref{wbII:eq:W:5.17-A20} therefore give
\(\langle d_0,\widehat\delta\rangle=0\).
Finally, the matrices \(A_1,A_2\) fix
\(\widehat\delta\), so the evaluation is unchanged by every principal
change in \eqref{wbII:eq:W:5.17-A3}.  Since
\(V_G=\mathbb Z[\delta]\) and
\(\delta(\widehat\delta)=1\), its vanishing is equivalent to
\(\bar d_0=0\).

\end{enumerate}
\end{proof}



\begin{lemma}[The toric support function detects central restriction]
\hypertarget{tag-03I6}{}\label{wbII:lem:analytic-support-specialization}
\StatementRecord{wbII:lem:analytic-support-specialization}{proved}{03I6}{\SeriesRef{wbII:prop:local-cusp-picard-obstruction}, \SeriesRef{wbII:eq:W:5.17ab}, \SeriesRef{wbII:eq:W:5.17-A24}, \SeriesRef{wbII:eq:W:5.17-A30}, \SeriesRef{wbII:eq:W:5.17-A29}, \SeriesRef{wbII:eq:W:5.17-A31}, \SeriesRef{wbII:eq:W:5.17-A32}, \SeriesRef{wbII:eq:W:5.17-A34b}, \SeriesRef{wbII:eq:analytic-residual-descent-factor}, \SeriesRef{wbII:eq:W:5.17am}, \SeriesRef{wbII:eq:W:5.17-A34c}, \SeriesRef{wbII:eq:W:5.17-A26}; \cite[2010 text, Propositions~4.1.1--4.1.2, p.~171; Theorems~4.2.8,
p.~181 and~4.2.12, p.~184]{CoxLittleSchenck2011}; \cite[2010 text, Example~1.2.21, p.~31]{CoxLittleSchenck2011}}

Let \(M\in\operatorname{Pic}(N_0)\), let
\(\pi:Y_\epsilon\to N_0\) be the cusp covering, and let
\(\nu:E_0\to W\) be the normalization.  If the Appell--Humbert form of
\(M|_{F_t}\), \(t\ne0\), is \(E_M\), then the three independent degrees
of \(\nu^*(M|_W)\) determine an integral symmetric form
\(B:\mathbb Z^2\times\mathbb Z^2\to\mathbb Z\) satisfying
\[
 E_M\bigl((\lambda,0),(0,m)\bigr)
 =-B(B_0\lambda,m)
 \qquad(\lambda,m\in\mathbb Z^2).
\tag{W.121}\hypertarget{tag-03F0}{}\label{wbII:eq:W:5.17-A21}
\]
Consequently
\[
 E_M=0
 \quad\Longrightarrow\quad
 \nu^*(M|_W)\simeq\mathcal O_{E_0}.
\tag{W.122}\hypertarget{tag-03F1}{}\label{wbII:eq:W:5.17-A22}
\]
\end{lemma}



\begin{proof}

\begin{enumerate}[label=\arabic*.,leftmargin=*]
\item Put \(\mathcal M=\pi^*M\).  The identity
\(\pi\circ\Psi_\lambda=\pi\) gives \(\mathcal M\) its \(\Psi\)-equivariant descent structure.  In the notation of
\cref{wbII:prop:local-cusp-picard-obstruction}, order the six rays in the star
of the origin as
\[
 (r_1,r_2,r_3,r_4,r_5,r_6)
 =(e_1,e_2,e_2-e_1,-e_1,-e_2,e_1-e_2),
 \qquad C_i=E_0\cap E_{r_i}.
\tag{W.123}\hypertarget{tag-03F2}{}\label{wbII:eq:W:5.17-A23}
\]
The quotient identifies opposite conductor branches.  Therefore, for
\(A=\mathcal M|_{E_0}=\nu^*(M|_W)\),
\[
 d_i=\deg_{C_i}A,
 \qquad
 (\deg_{C_1}A,\ldots,\deg_{C_6}A)
 =(d_1,d_2,d_3,d_1,d_2,d_3).
\tag{W.124}\hypertarget{tag-03F3}{}\label{wbII:eq:W:5.17-A24}
\]

\item Define the integer-valued quadratic function
\[
 a(x,y)
 =(d_2+d_3)\binom{x}{2}+d_3xy+(d_1+d_3)\binom{y}{2}.
\tag{W.125}\hypertarget{tag-03F4}{}\label{wbII:eq:W:5.17-A25}
\]
Its polarization is
\[
 B(v,w)=a(v+w)-a(v)-a(w),
 \qquad
 [B]=
 \begin{pmatrix}
  d_2+d_3&d_3\\ d_3&d_1+d_3
 \end{pmatrix}.
\tag{W.126}\hypertarget{tag-03F5}{}\label{wbII:eq:W:5.17-A26}
\]
The locally finite divisor
\[
 D_a=\sum_{v\in\mathbb Z^2}a(v)E_v
\tag{W.127}\hypertarget{tag-03F6}{}\label{wbII:eq:W:5.17-A27}
\]
is Cartier: on every maximal chart, the three height-one ray generators
form an integral basis, so the three assigned coefficients are the values
of one integral character.  In the standard toric convention the associated
support function has ray values \(-a(v)\); the divisor coefficients here are
the values \(a(v)\), consistent with
\(\operatorname{div}(\chi^m)=\sum_\rho\langle m,u_\rho\rangle D_\rho\)
\cite[2010 text, Propositions~4.1.1--4.1.2, p.~171; Theorems~4.2.8,
p.~181 and~4.2.12, p.~184]{CoxLittleSchenck2011}.
This is the same chartwise Cartier argument as
in \eqref{wbII:eq:W:5.17ab}, and it applies to an arbitrary integer-valued \(a\).

\item Direct substitution gives
\[
 (a(r_1),\ldots,a(r_6))
 =(0,0,d_2,d_2+d_3,d_1+d_3,d_1).
\tag{W.128}\hypertarget{tag-03F7}{}\label{wbII:eq:W:5.17-A28}
\]
Using \(C_i^2=-1\) and \(C_iC_{i\pm1}=1\), the degree vector of
\(\mathcal O(D_a)|_{E_0}\) is exactly the vector in \eqref{wbII:eq:W:5.17-A24}.  More
generally, restrict the principal identity
\(\operatorname{div}(t)=\sum_wE_w\) to \(E_v\).  Since \(E_v\) meets
precisely \(E_{v+r_i}\), it gives
\[
 \mathcal O(D_a)|_{E_v}
 \simeq
 \mathcal O_{E_v}\!\left(
   \sum_{i=1}^6(a(v+r_i)-a(v))C_{v,i}
 \right).
\tag{W.129}\hypertarget{tag-03F8}{}\label{wbII:eq:W:5.17-A29}
\]
The coefficient difference is
\[
 a(v+r_i)-a(v)=a(r_i)+B(v,r_i).
\tag{W.130}\hypertarget{tag-03FA}{}\label{wbII:eq:W:5.17-A30}
\]
The second summand in \eqref{wbII:eq:W:5.17-A30} is the divisor of the toric character
with exponent \(B(v,-)\).  Thus every restriction in \eqref{wbII:eq:W:5.17-A29} has the
same six degrees as \(\mathcal O(D_a)|_{E_0}\).

\item Because \(B_0\bar\Lambda=\mathbb Z^2\), the deck action is transitive on
the components \(E_v\).  The descent structure on \(\mathcal M\)
therefore transports \eqref{wbII:eq:W:5.17-A24}, with its ordered six boundary directions,
to every \(E_v\).  It follows that
\[
 \mathcal R:=\mathcal M\otimes\mathcal O(-D_a)
 \quad\text{satisfies}\quad
 \mathcal R|_{E_v}\simeq\mathcal O_{E_v}
 \quad(v\in\mathbb Z^2).
\tag{W.131}\hypertarget{tag-03FB}{}\label{wbII:eq:W:5.17-A31}
\]
Indeed, the six toric curves detect \(\operatorname{Pic}(\mathrm{dP}_6)\):
in the blow-up basis \(H,E_1,E_2,E_3\), vanishing of the degrees on the
three exceptional curves kills the three \(E_i\)-coefficients, and
vanishing on any strict transform \(H-E_i-E_j\) kills the remaining
\(H\)-coefficient.  Since
\(H^1(\mathrm{dP}_6,\mathcal O)=0\), equality of those degrees implies
isomorphism, not merely numerical equivalence.

\item For \(b_\lambda=B_0\lambda\), the torus-translation part of
\(\Psi_\lambda\) preserves each toric prime divisor, while its fan part
translates the index by \(b_\lambda\).  Hence
\[
 \Psi_\lambda^*D_a-D_a
 =\sum_v(a(v+b_\lambda)-a(v))E_v
 =\operatorname{div}(\chi^{m_\lambda}),
\tag{W.132}\hypertarget{tag-03FC}{}\label{wbII:eq:W:5.17-A32}
\]
where
\[
 m_\lambda=(\ell_\lambda,k_\lambda),
 \qquad
 \ell_\lambda=B(b_\lambda,-),
 \qquad
 k_\lambda=a(b_\lambda).
\tag{W.133}\hypertarget{tag-03FD}{}\label{wbII:eq:W:5.17-A33}
\]
Thus multiplication by \(\chi^{m_\lambda}\) gives an individual
isomorphism
\(\Psi_\lambda^*\mathcal O(D_a)\simeq\mathcal O(D_a)\).
These individual isomorphisms separate the horizontal winding of the
descent factor of $\mathcal M$ in the tensor-product map below.

\item For a ray \(\mathbb R_{\ge0}(v,1)\), its toric open set in
\(Y_\epsilon\) is
\[
 U_v\simeq\Delta_\epsilon\times(\mathbb C^*)^2,
 \qquad E_v^\circ=\{0\}\times(\mathbb C^*)^2.
\tag{W.134}\hypertarget{tag-03FE}{}\label{wbII:eq:W:5.17-A34}
\]
Indeed \((e_1,0),(e_2,0),(v,1)\) is an integral basis and \(t\) is the
character dual to the last vector; compare
\cite[2010 text, Example~1.2.21, p.~31]{CoxLittleSchenck2011}.
The set \(U_v\) is Stein and retracts
onto \(E_v^\circ\).  By \eqref{wbII:eq:W:5.17-A31},
\(c_1(\mathcal R|_{U_v})=0\); the exponential sequence and Cartan's
theorem~B therefore trivialize \(\mathcal R|_{U_v}\).  Choose a global
nowhere-zero frame \(s_v\) of \(\mathcal R|_{E_v}\).  A frame on \(U_v\) can be
multiplied by a unit pulled back from \((\mathbb C^*)^2\) so that its
restriction to \(E_v^\circ\) agrees with \(s_v|_{E_v^\circ}\); denote the
resulting frame by \(e_v\).

\item For the comparison of these frames on the generic torus, let \(v,w\) be
adjacent vertices and choose a triangle \(\sigma=\{v,w,u\}\) containing them,
write \(z_v,z_w,z_u\) for the coordinates dual to its three height-one ray
generators, and put
\[
 S_\sigma:=U_\sigma\cap Y_\epsilon
 =\{(z_v,z_w,z_u)\in\mathbb C^3:
       |z_vz_wz_u|<\epsilon\}.
\tag{W.135}\hypertarget{tag-03FF}{}\label{wbII:eq:W:5.17-A34a}
\]
Because \(\rho_v<\sigma\), localization of the maximal affine chart gives
\begin{equation}
 U_{\rho_v}=D(z_wz_u)\subset U_\sigma,\qquad
 U_{\rho_w}=D(z_vz_u)\subset U_\sigma,\qquad
 U_{\rho_v}\cap U_{\rho_w}=D(z_vz_wz_u)=U_0.
\hypertarget{tag-03I5}{}\label{wbII:eq:analytic-toric-open-directions}
\end{equation}
Thus \(U_\sigma\) is not contained in \(U_{\rho_v}\cap U_{\rho_w}\), and
the coordinate discs in \(S_\sigma\) cannot be applied directly to the
ratio \(e_w/e_v\).

\item The space $S_\sigma$ is the closed graph
\[
 \{(z_v,z_w,z_u,t)\in\mathbb C^3\times\Delta_\epsilon:
   t=z_vz_wz_u\}
\]
in the Stein space \(\mathbb C^3\times\Delta_\epsilon\), and hence is
Stein.  Radial scaling of the three \(z\)-coordinates contracts it to the
origin.  The exponential sequence and Cartan's theorem~B therefore give
\(\operatorname{Pic}(S_\sigma)=0\).  Choose a frame \(e_\sigma\) of
\(\mathcal R|_{S_\sigma}\).

\item Let \(s_v\) be the global frame of
\(\mathcal R|_{E_v}\simeq\mathcal O_{E_v}\) used to normalize \(e_v\).
On \(U_{\rho_v}\cap Y_\epsilon\) put
\[
 A_{\sigma,v}=e_v/e_\sigma.
\]
At \(t=0\), this is the restriction to \((\mathbb C^*)^2\) of the
holomorphic unit
\[
 G_{\sigma,v}
 =\frac{s_v}{e_\sigma|_{E_v\cap S_\sigma}}
 \in\mathcal O^*(E_v\cap S_\sigma),
 \qquad E_v\cap S_\sigma=\{z_v=0\}\simeq\mathbb C^2.
\]
Consequently its winding on every horizontal cocharacter is zero at
\(t=0\).  Under
\(U_{\rho_v}\cap Y_\epsilon\simeq
\Delta_\epsilon\times(\mathbb C^*)^2\), that winding is an integer locally
constant in \(t\), and is therefore zero on every nearby nonzero fibre.
The same argument applies to \(A_{\sigma,w}\).  On their common
domain \(U_0\cap Y_\epsilon\),
\[
 e_w/e_v=A_{\sigma,w}/A_{\sigma,v},
\]
so this ratio has zero horizontal winding.  The one-skeleton of the
\(A_2\)-triangulation is connected, and multiplication along an edge path
gives
\[
 \operatorname{wind}_m(e_w/e_v)=0
 \qquad(v,w,m\in\mathbb Z^2)
\tag{W.136}\hypertarget{tag-03FG}{}\label{wbII:eq:W:5.17-A34b}
\]
on every nonzero fibre.

\item Let
\(\theta^{\mathcal M}_\lambda:\Psi_\lambda^*\mathcal M\to\mathcal M\)
be the descent isomorphism, let
\(\mu^+_\lambda:\Psi_\lambda^*\mathcal O(D_a)\to\mathcal O(D_a)\)
be multiplication by \(\chi^{m_\lambda}\) from \eqref{wbII:eq:W:5.17-A32}, and define the
inverse-dual isomorphism
\[
 \mu^-_\lambda:=((\mu^+_\lambda)^{-1})^\vee:
 \Psi_\lambda^*\mathcal O(-D_a)\longrightarrow\mathcal O(-D_a),
 \qquad s\longmapsto\chi^{-m_\lambda}s.
\]
The ordinary dual of \(\mu^+_\lambda\) has the opposite arrow; the
inverse-dual above is the map required by the pullback tensor product.  Then
\[
 \theta^{\mathcal R}_\lambda
 :=\theta^{\mathcal M}_\lambda\otimes\mu^-_\lambda:
 \Psi_\lambda^*\mathcal R\longrightarrow\mathcal R.
\]
Since \(\Psi_\lambda(U_v)=U_{v+b_\lambda}\), define
\(r_{\lambda,v}\in\mathcal O^*(U_v)\) by the identity
\begin{equation}
 \theta^{\mathcal R}_\lambda
   (\Psi_\lambda^*e_{v+b_\lambda})
 =r_{\lambda,v}e_v.
\hypertarget{tag-03I4}{}\label{wbII:eq:analytic-residual-descent-factor}
\end{equation}
At \(t=0\) it extends to an isomorphism between two trivial line bundles
on the compact surface \(E_v\).  Its factor in the chosen global frames is
therefore a nonzero constant.  The winding of
\(r_{\lambda,v}(t,-)\) around every horizontal one-parameter subgroup is
an integer locally constant in \(t\); at \(t=0\) it is zero, hence it is
zero for every nearby \(t\ne0\).

\item To express this isomorphism in one frame on the generic torus, put
\(q_{\lambda,v}=e_{v+b_\lambda}/e_v\).  Equation \eqref{wbII:eq:W:5.17-A34b} gives zero
horizontal winding for \(q_{\lambda,v}\); because
\(\Psi_\lambda|_{Y_t}\) is multiplication by a fixed torus point, it
induces the identity on horizontal \(H_1\), and
\(\Psi_\lambda^*q_{\lambda,v}\) also has zero winding.  Relative to the
single frame \(e_v\), the residual multiplier is therefore
\[
 h_{\lambda,v}
 =\frac{r_{\lambda,v}}{\Psi_\lambda^*q_{\lambda,v}},
 \qquad
 \operatorname{wind}_m(h_{\lambda,v})=0.
\tag{W.137}\hypertarget{tag-03FH}{}\label{wbII:eq:W:5.17-A34c}
\]

\item On \(U_0\) let \(\mathbf1_a\) be the canonical toric frame of
\(\mathcal O(D_a)\), and put \(F_v=e_v\otimes\mathbf1_a\).  Equations
\cref{wbII:eq:analytic-residual-descent-factor} and \eqref{wbII:eq:W:5.17-A32} give
\[
 \theta^{\mathcal M}_\lambda(\Psi_\lambda^*F_v)
 =h_{\lambda,v}\chi^{m_\lambda}F_v.
\]
On a nonzero fibre, \(\chi^{m_\lambda}=x^{\ell_\lambda}t^{k_\lambda}\).
With the section-descent convention \eqref{wbII:eq:W:5.17am}, the factor of automorphy is
therefore
\[
 \alpha_\lambda=(h_{\lambda,v}\chi^{m_\lambda})^{-1}.
\]
The residual factor has zero horizontal winding by \eqref{wbII:eq:W:5.17-A34c}, and hence
\begin{equation}
 \operatorname{wind}_m(\alpha_\lambda)
 =-\ell_\lambda(m)=-B(B_0\lambda,m).
\hypertarget{tag-03I3}{}\label{wbII:eq:analytic-factor-horizontal-winding}
\end{equation}

\item This winding is exactly the mixed Chern-class entry.  Pull \(F_v\) back to
the universal cover
\[
 p_t:\mathbb C^2\longrightarrow Y_t,
 \qquad \zeta\longmapsto(e^{2\pi i\zeta},t).
\]
The period indexed by \((\lambda,n)\in\mathbb Z^2\oplus\mathbb Z^2\)
acts by \(\zeta\mapsto\zeta+Z(s)\lambda+n\).  The horizontal factor is one,
because \(F_v\) is already defined on \(Y_t\).  Choose a logarithm
\(g_\lambda\in\mathcal O(\mathbb C^2)\) with
\[
 p_t^*\alpha_\lambda=e^{2\pi i g_\lambda}.
\]
For every horizontal period \(n\),
\[
 g_\lambda(\zeta+n)-g_\lambda(\zeta)
 =\operatorname{wind}_n(\alpha_\lambda).
\]
In the exponential connecting cocycle
\[
 c(l,l')=g_l(\zeta+l')+g_{l'}(\zeta)-g_{l+l'}(\zeta),
\]
take the horizontal logarithm to be zero and the logarithm for the combined
period \((\lambda,n)\) to be \(g_\lambda\).  Then
\[
 c((\lambda,0),(0,n))
 =\operatorname{wind}_n(\alpha_\lambda),
 \qquad
 c((0,n),(\lambda,0))=0.
\]
Therefore the alternating exponential connecting integer between
the vertical deck
translation indexed by \(\lambda\) and the horizontal loop indexed by
\(m\) is
\[
 E_M\bigl((\lambda,0),(0,m)\bigr)
 =-\ell_\lambda(m)
 =-B(B_0\lambda,m),
\]
which proves \eqref{wbII:eq:W:5.17-A21}.  This sign also recovers the explicit check in
\cref{wbII:prop:local-cusp-picard-obstruction}: for
\(a(v)=v^{\mathsf t}Hv\), one has \(B=2H\), so \eqref{wbII:eq:W:5.17-A21} is the
cross-block \(-2B_0^{\mathsf t}H=-2P\), namely \(-2q\).

\item If \(E_M=0\), equation \eqref{wbII:eq:W:5.17-A21} and the unimodularity of \(B_0\) force
\(B=0\).  The matrix in \eqref{wbII:eq:W:5.17-A26} then gives
\(d_3=0\), \(d_2=0\), and \(d_1=0\).  Thus all six degrees of \(A\)
vanish, and the Picard detection argument following \eqref{wbII:eq:W:5.17-A31} gives
\(A\simeq\mathcal O_{E_0}\).

\end{enumerate}
\end{proof}



\begin{theorem}[Analytic vanishing of the global restriction]
\hypertarget{tag-03IW}{}\label{wbII:thm:analytic-all-picard}
\StatementRecord{wbII:thm:analytic-all-picard}{proved}{03IW}{\SeriesRef{wbII:eq:W:5.17-A5}, \SeriesRef{wbII:eq:W:5.17-A4}, \SeriesRef{wbII:eq:W:5.17-A6}, \SeriesRef{wbII:lem:analytic-support-specialization}}

For the standard threefold,
\[
 r=\eta^*i^*:\operatorname{Pic}(X)\longrightarrow
 \operatorname{Pic}(\mathrm{dP}_6)
\]
is the zero homomorphism.
\end{theorem}



\begin{proof}
Let \(L\in\operatorname{Pic}(X)\).  By
\cite[Proposition~7.26, pp.~61--62]{AC2026}, its invariant fibre class is
\[
 E=c_1(L|_F)=kq,
 \qquad k\in\mathbb Z,
 \qquad q=uw+6\gamma\delta.
\tag{W.138}\hypertarget{tag-03FI}{}\label{wbII:eq:W:5.17-A35}
\]
The contractions are
\[
 \iota_{v_1}q=2\eta_1,
 \qquad
 \iota_{v_2}q=-3\eta_2,
 \qquad
 [\eta_1]=3[\delta],
 \qquad
 [\eta_2]=2[\delta].
\tag{W.139}\hypertarget{tag-03FJ}{}\label{wbII:eq:W:5.17-A36}
\]
Apply \eqref{wbII:eq:W:5.17-A5} with \((m_1,m_2)=(3,4)\).  In the real global
coinvariants,
\[
 \bar d_1=2k[\delta],
 \qquad
 \bar d_2=-\frac32k[\delta].
\tag{W.140}\hypertarget{tag-03FK}{}\label{wbII:eq:W:5.17-A37}
\]
Equations \eqref{wbII:eq:W:5.17-A4} and \eqref{wbII:eq:W:5.17-A6} therefore give
\[
 0=\bar d_1+\bar d_2+\bar d_0
   =\frac{k}{2}[\delta].
\tag{W.141}\hypertarget{tag-03FL}{}\label{wbII:eq:W:5.17-A38}
\]
Since \(V_G\otimes\mathbb R=\mathbb R[\delta]\), one has \(k=0\), so
\(E=0\).  Apply \cref{wbII:lem:analytic-support-specialization} to
\(M=L|_{N_0}\).  It gives
\[
 \eta^*i^*L=\nu^*(L|_W)\simeq\mathcal O_{E_0}.
\]

\end{proof}




The nonvanishing conclusion of the next corollary is contained in
\citet[Theorem~10.5(a), pp.~88--89]{AC2026}, which treats every global line
bundle.  This proof records the smaller set of inputs sufficient for
\(\operatorname{Pic}^0\)-twists.

\begin{corollary}[Standard-family nonvanishing for \texorpdfstring{$\operatorname{Pic}^0$}{Pic0}]
\hypertarget{tag-03DS}{}\label{wbII:cor:candidate-pic0-route}
\StatementRecord{wbII:cor:candidate-pic0-route}{proved}{03DS}{\SeriesRef{wbII:prop:pic0-pullback}, \SeriesRef{wbII:prop:conductor}, \SeriesRef{wbII:prop:cartier-duality-absolute}, \SeriesRef{wbII:lem:curve-flatness}, \SeriesRef{wbII:prop:analytic-base-change}}

For the standard threefold and every
$L\in\operatorname{Pic}^0(X)$,
\[
 H^0(W,\Omega_X^1|_W\otimes A_L)\ne0,
 \qquad
 H^2(W,(T_X\otimes L)|_W)\ne0,
\]
and
\[
 R^2f_*(T_X\otimes L)\ne0.
\tag{W.44}\hypertarget{tag-03GP}{}\label{wbII:eq:W:5.17f}
\]
\end{corollary}



\begin{proof}
The full cusp fibre $W=f^{-1}(p_0)$ is reduced by
\CompendiumRef{I}{def:common-conventions}.
At $k=l=1$, the cusp charts in the proof of
\CompendiumRef{I}{push:src:E.3.1} are $q=z_0z_1z_2$,
with the two-factor and smooth charts on the other strata,
and the normalization is $\eta:\mathrm{dP}_6\to W$.
The canonical-bundle isomorphism
$\omega_X\simeq\mathcal O_X(-2S_2)$ in
\CompendiumRef{I}{fam:src:A.1.1} gives a trivialization on
$X\setminus S_2$, which contains $W$: the two supports
are fibres over the distinct points $p_0$ and $p_2$. By
\cref{wbII:prop:pic0-pullback}, $\eta^*A_L$ is trivial, so
\cref{wbII:prop:conductor} gives the first nonvanishing.  Apply
\cref{wbII:prop:cartier-duality-absolute} to $E=T_X\otimes L$; its dual space is
exactly the preceding nonzero section space, giving the fibrewise $H^2$
nonvanishing.  The fibre has complex dimension two, so its coherent
$H^3$ vanishes.  The global construction makes $X$ compact and $f$
nonconstant; \cref{wbII:lem:curve-flatness} therefore gives properness and
$f$-flatness of $T_X\otimes L$.  The base-change isomorphism of
\cref{wbII:prop:analytic-base-change} identifies the fibre of the $R^2$
stalk with the nonzero $H^2$ group.  Hence the stalk, and therefore the
direct-image sheaf, is nonzero.

\end{proof}



\paragraph{reconciliation of the three restriction routes}

The Pic$^0$ argument proves \eqref{wbII:eq:W:5.17f} without global homology,
Mayer--Vietoris, or Leray.  The all-Pic statement
\cref{wbII:prop:leray-edge-restriction} avoids citing the packaged theorem
$H^2(X;\mathbb Z)=0$, but it reuses the same verified global integral
Leray and specialization machinery from which that theorem is deduced.
By contrast, \cref{wbII:prop:direct-boundary-picard} supplies a second all-Pic
proof without any global Leray differential: it recovers the primitive
$4-3=1$ obstruction directly from the finite logarithmic gluing and the
cusp boundary.  It still uses the local integral cusp-specialization
theorem.  The third proof in \cref{wbII:thm:analytic-all-picard} replaces both
that specialization and the cusp Chern-class bounding-chain evaluation by
logarithmic $\mathcal O^*$ factors, the argument principle, and the exact
$\Psi$-equivariant support function.  Thus the purely analytic
strengthening is now complete; the three proofs have different
dependency sets.

\Uses{the hypotheses, notation and definitions stated here; \cref{wbII:eq:W:5.17f}; \cref{wbII:prop:leray-edge-restriction}; \cref{wbII:prop:direct-boundary-picard}; \cref{wbII:thm:analytic-all-picard}}


\subsection{Abstract conductor theorem for the standard fibre}

The conductor mechanism below is the argument of
\citet[Lemma~10.3 and Theorem~10.5(a), pp.~87--89]{AC2026}, with its sufficient
hypotheses separated explicitly.  In particular the constructed-family
conclusion of \cref{wbII:cor:candidate-conductor-actual} is the known
nonvanishing assertion of that theorem.  The complete proof is retained
so that its normalization, duality and base-change inputs can be checked
individually.

\begin{theorem}[Conductor nonvanishing for the standard family]
\hypertarget{tag-03IX}{}\label{wbII:thm:candidate-conductor}
\StatementRecord{wbII:thm:candidate-conductor}{proved}{03IX}{\SeriesRef{wbII:eq:W:5.3}, \SeriesRef{wbII:prop:conductor}, \SeriesRef{wbII:thm:cartier-duality}, \SeriesRef{wbII:eq:W:5.20}, \SeriesRef{wbII:lem:curve-flatness}, \SeriesRef{wbII:prop:analytic-base-change}; \cite[Theorem~9.2.3, 2010 text, p.~410]{CoxLittleSchenck2011}; \cite[\S12, Theorem~1, p.~19]{Serre1956GAGA}}

Assume:
\begin{enumerate}[label=(H\arabic*)]
\item $X$ is a connected compact smooth complex threefold and
$f:X\to B$ is a nonconstant holomorphic map to a smooth compact curve;
\item $W=f^{-1}(p_0)$ is the full scheme-theoretic fibre, is reduced, and
has the normal-crossings local models \eqref{wbII:eq:W:5.3};
\item the normalization $\eta:\widetilde W\to W$ is the degree-six del
Pezzo surface $\mathrm{dP}_6$;
\item $H^2(X;\mathbb Z)=0$.
\end{enumerate}
Then, for every $L\in\operatorname{Pic}(X)$,
\[
H^2(W,(T_X\otimes L)|_W)\ne0
\quad\text{and}\quad
R^2f_*(T_X\otimes L)\ne0.
\tag{W.45}\hypertarget{tag-03HI}{}\label{wbII:eq:W:5.18}
\]
\end{theorem}



\begin{proof}

\begin{enumerate}[label=\arabic*.,leftmargin=*]
\item Fix $L$ and put
\[
\mathcal F_L=T_X\otimes L,\qquad
A_L=(L^\vee\otimes\omega_X)|_W.
\]
Hypothesis (H4) gives $c_1(L^\vee\otimes\omega_X)=0$.  Pullback therefore
gives $c_1(\eta^*A_L)=0$. The surface $\mathrm{dP}_6$ is the blow-up of
$\mathbb P^2$ at three noncollinear points. A projective change of coordinates
moves these points to the three coordinate points, so this is the smooth
complete toric surface obtained by subdividing the three two-dimensional
cones of the fan of $\mathbb P^2$. Its fan has support $\mathbb R^2$,
and the zero Cartier divisor is nef. Demazure vanishing therefore gives
$H^1(\mathrm{dP}_6,\mathcal O)=0$ in algebraic cohomology
\cite[Theorem~9.2.3, 2010 text, p.~410]{CoxLittleSchenck2011}.
Since this surface is projective and $\mathcal O$ is coherent,
Serre's comparison theorem identifies that group with analytic cohomology
\cite[\S12, Theorem~1, p.~19]{Serre1956GAGA}.
In the exponential sequence, the kernel
of
\[
c_1:\operatorname{Pic}(\mathrm{dP}_6)\longrightarrow
H^2(\mathrm{dP}_6;\mathbb Z)
\]
is the image of $H^1(\mathrm{dP}_6,\mathcal O)$ and is therefore zero.
Consequently
\[
\eta^*A_L\simeq\mathcal O_{\widetilde W}.
\tag{W.46}\hypertarget{tag-03HJ}{}\label{wbII:eq:W:5.19}
\]
\Cref{wbII:prop:conductor} gives
\[
H^0(W,\Omega^1_X|_W\otimes A_L)\ne0.
\tag{W.47}\hypertarget{tag-03HK}{}\label{wbII:eq:W:5.20}
\]

\item Because $W=f^*(p_0)$, one has
$\mathcal O_X(W)=f^*\mathcal O_B(p_0)$, whose restriction to the fibre is
trivial.  Adjunction therefore gives $\omega_W\simeq\omega_X|_W$.
Apply \cref{wbII:thm:cartier-duality} to $E=\mathcal F_L$:
\[
\begin{split}
H^2(W,\mathcal F_L|_W)^\vee
&\simeq H^0(W,(\mathcal F_L|_W)^\vee\otimes\omega_W)\\
&=H^0(W,\Omega^1_X|_W\otimes A_L).
\end{split}
\]
By \eqref{wbII:eq:W:5.20}, the dual space is nonzero, giving the first nonvanishing in \eqref{wbII:eq:W:5.18}.
The same duality theorem gives
$H^3(W,\mathcal F_L|_W)=0$.

\item Since $X$ is compact, $f$ is proper.  By
\cref{wbII:lem:curve-flatness}, $f$ is flat and the vector bundle
$\mathcal F_L$ is $f$-flat.  Thus every hypothesis of
\cref{wbII:prop:analytic-base-change} is satisfied at $p_0$, and
\[
(R^2f_*\mathcal F_L)_{p_0}\otimes\mathbb C(p_0)
\simeq H^2(W,\mathcal F_L|_W)\ne0.
\]
The stalk, and hence the coherent direct-image sheaf in
\eqref{wbII:eq:W:5.18}, is nonzero.

\end{enumerate}
\end{proof}



\begin{corollary}[Application to the standard threefold]
\hypertarget{tag-03DR}{}\label{wbII:cor:candidate-conductor-actual}
\StatementRecord{wbII:cor:candidate-conductor-actual}{proved}{03DR}{\SeriesRef{wbII:prop:direct-boundary-picard}, \SeriesRef{wbII:cor:candidate-pic0-route}, \SeriesRef{wbII:prop:leray-edge-restriction}}

For the global construction of
\CompendiumRef{I}{fam:src:A.3.1} and every
$L\in\operatorname{Pic}(X)$,
\[
H^2(W,(T_X\otimes L)|_W)\ne0,
\qquad R^2f_*(T_X\otimes L)\ne0.
\]
\end{corollary}



\begin{proof}
The cusp charts in \CompendiumRef{I}{push:src:E.3.1} give the compact
full Cartier fibre and its $\mathrm{dP}_6$ normalization.  The isomorphism
$\omega_X\simeq\mathcal O_X(-2S_2)$ in
\CompendiumRef{I}{fam:src:A.1.1} trivializes $\omega_X$ near $W$, since
$W$ and $S_2$ lie over distinct base points.
By the local canonical trivialization and
\cref{wbII:prop:direct-boundary-picard}, $\eta^*A_L$ is trivial for every $L$.
The conductor proposition then gives
$H^0(W,\Omega_X^1|_W\otimes A_L)\ne0$; absolute Cartier-fibre duality turns
this into the displayed $H^2$ nonvanishing, and the uniform top-degree
base-change isomorphism turns it into a nonzero $R^2$ stalk.  This chain is proved in \cref{wbII:cor:candidate-pic0-route}, with the all-Pic
restriction supplied by \cref{wbII:prop:direct-boundary-picard}.  The restriction also follows from
\cref{wbII:prop:leray-edge-restriction}.

\end{proof}



\subsection{The corrected theorem that survives both defects}

The original hypotheses do not force the CDP conclusion: the constructed
threefold above is a counterexample.  There is nevertheless a sharp useful
replacement.  It isolates the exact extra global vanishing needed by the
intact part of the argument, and the following fibrewise criterion replaces
the invalid p.~692 reduction.

\begin{lemma}[Generic endpoints in an analytic parameter space]
\hypertarget{tag-03IA}{}\label{wbII:lem:generic-endpoints}
\StatementRecord{wbII:lem:generic-endpoints}{proved}{03IA}{\cite[Chapter~III, \S4, Theorem~4.1, p.~124]{BS76push}}

Let $Y$ be a connected compact complex manifold of dimension $n$, let $D$ be
a nonzero effective Cartier divisor with $\cO_Y(D)\in\Pic^0(Y)$, and let
$B$ be a holomorphic vector bundle.  Let $\mathcal P\subset\Pic^0(Y)$ be a
connected irreducible reduced complex analytic parameter space carrying a
Poincar\'e line bundle on $Y\times\mathcal P$.  Assume that dualization
preserves $\mathcal P$, and suppose it
  contains points $P_0=\cO_Y(-mD)$ and $Q_0=\cO_Y(-\ell D)$ for some
  $m,\ell\ge1$ with
  $H^0(Y,B\otimes P_0)=0$ and
  $H^0(Y,B^\vee\otimes\omega_Y\otimes Q_0)=0$.  Then there is a dense open subset
$U\subset\mathcal P$ such that, for every $P'\in U$,
\[
 H^0(Y,B\otimes P')=H^n(Y,B\otimes P')=0.
 \tag{W.169}\hypertarget{tag-03E5}{}\label{wbII:eq:CN:5.2}
\]
\end{lemma}



\begin{proof}
  The chosen point $P_0$ satisfies
  $H^0(Y,B\otimes P_0)=0$.  The Poincar\'e bundle gives a proper flat family over
$\mathcal P$.  Locally, its cohomology is computed by a finite free complex
\cite[Chapter~III, \S4, Theorem~4.1, p.~124]{BS76push}.
The rank conditions on its differentials make
$\{P':h^0(Y,B\otimes P')\ge1\}$ a determinantal locus; these local loci glue
to a closed analytic subset of $\mathcal P$.  It is proper because it misses $P_0$, and irreducibility of
$\mathcal P$ makes its complement dense open.  The other chosen point satisfies
$H^0(Y,B^\vee\otimes\omega_Y\otimes Q_0)=0$; dualization stability puts
$Q_0^\vee$ in $\mathcal P$.  Serre duality then gives
$H^n(Y,B\otimes Q_0^\vee)=0$.  The corresponding top-degree jump locus is
therefore also proper with dense open complement.  Intersecting the two
dense open complements proves the claim.

\end{proof}



The following endpoint argument is a re-derivation of
\citet[Proposition~1.1 and Corollaries~1.2--1.3]{CDP98}, with Serre duality
for the top degree.  Its formulation makes the effective-divisor
\(\operatorname{Pic}^0\) hypothesis explicit in arbitrary dimension.

\begin{lemma}[Generic endpoints on \texorpdfstring{$\operatorname{Pic}^0$}{Pic0}]\hypertarget{tag-03I7}{}\label{wbII:lem:cdp-generic-endpoints}
\StatementRecord{wbII:lem:cdp-generic-endpoints}{proved}{03I7}{\cite[Proposition~1.1, p.~78]{CDP98}; \cite[Chapter~III, \S4, Theorem~4.1, p.~124]{BS76push}}

Let \(Y\) be a connected compact complex manifold, let \(D\) be
a nonzero effective Cartier divisor, and suppose
\(\mathcal O_Y(D)\in\operatorname{Pic}^0(Y)\).  For every holomorphic vector
bundle \(B\), a general \(P\in\operatorname{Pic}^0(Y)\) satisfies
\[
 H^0(Y,B\otimes P)=0,\qquad
 H^{\dim Y}(Y,B\otimes P)=0.
\tag{W.50}\hypertarget{tag-03HL}{}\label{wbII:eq:W:5.23}
\]
\end{lemma}



\begin{proof}
By \cite[Proposition~1.1, p.~78]{CDP98}, applied to the connected smooth
space \(Y\), the divisor \(D\), and the locally free sheaf \(B\), one has
\[
 V_k=H^0(Y,B\otimes\mathcal O_Y(-kD))
 \subseteq H^0(Y,B),
\]
with \(V_k=0\) for all sufficiently large \(k\).
Since
\(\mathcal O_Y(-kD)\in\operatorname{Pic}^0(Y)\), local finite free complexes for the proper family over a
local universal parameter space of topologically trivial line bundles
\cite[Chapter~III, \S4, Theorem~4.1, p.~124]{BS76push} show that
\[
 \{P:h^0(Y,B\otimes P)>0\}
\]
is a proper analytic subset of the connected complex Lie group
\(\operatorname{Pic}^0(Y)\), so its complement is dense open.
Serre duality identifies the dual
of the second group in \eqref{wbII:eq:W:5.23} with
\[
 H^0(Y,B^\vee\otimes\omega_Y\otimes P^\vee).
\]
Apply the first part to \(B^\vee\otimes\omega_Y\) and use that inversion is
an automorphism of \(\operatorname{Pic}^0(Y)\).

\end{proof}



\begin{theorem}[Global vanishing from one degree-two vanishing]
\hypertarget{tag-03J0}{}\label{wbII:thm:corrected-cdp-global}
\StatementRecord{wbII:thm:corrected-cdp-global}{proved}{03J0}{\SeriesRef{wbII:lem:cdp-generic-endpoints}, \SeriesRef{wbII:eq:W:5.28}}
\label{wbII:thm:corrected-global}
Let \(X\) be a smooth compact connected complex threefold such that
\[
 H^1(X;\mathbb Z)=H^2(X;\mathbb Z)=0,\qquad a(X)=1,
\]
and suppose its algebraic reduction is a holomorphic map
\(f:X\to C\).  If there exists \(L_0\in\operatorname{Pic}(X)\) with
\[
 H^2(X,T_X\otimes L_0)=0,
\tag{W.51}\hypertarget{tag-03HM}{}\label{wbII:eq:W:5.24}
\]
then, for \(M\) in a dense open subset of \(\operatorname{Pic}(X)\),
\[
 H^i(X,T_X\otimes M)=0\quad(i\ne1),
\tag{W.52}\hypertarget{tag-03HN}{}\label{wbII:eq:W:5.25}
\]
\[
 \chi(X,T_X\otimes M)\le0,\qquad
 \left\langle c_3(T_X),[X]\right\rangle\le0.
\tag{W.53}\hypertarget{tag-03HO}{}\label{wbII:eq:W:5.26}
\]
\end{theorem}



\begin{proof}

\begin{enumerate}[label=\arabic*.,leftmargin=*]
\item The exponential sequence and the two integral vanishings give
\[
 \operatorname{Pic}(X)=\operatorname{Pic}^0(X)
 \simeq H^1(X,\mathcal O_X).
\tag{W.54}\hypertarget{tag-03HP}{}\label{wbII:eq:W:5.27}
\]
Thus this Picard group is a connected irreducible complex vector group.
Choose additive \v{C}ech cocycles $a^{(r)}_{ij}$ representing a basis of
$H^1(X,\mathcal O_X)$. For coordinates $\lambda_r$ on that vector space,
\[
 g_{ij}(\lambda)=\exp\!\left(\sum_r\lambda_r a^{(r)}_{ij}\right)
\]
satisfy the multiplicative cocycle condition and define a holomorphic
line bundle on $X\times H^1(X,\mathcal O_X)$. Its restriction at every
parameter represents the corresponding line bundle, and the projection to the parameter space is proper.  Any fibre divisor \(D\) of
\(f\) is nonzero and effective, and its line bundle lies in
\(\operatorname{Pic}^0(X)\) by \(H^2(X;\mathbb Z)=0\).  Applying
\cref{wbII:lem:cdp-generic-endpoints} to \(B=T_X\) shows that the degree-zero and
degree-three groups in \eqref{wbII:eq:W:5.25} vanish away from two proper analytic subsets
of \(\operatorname{Pic}(X)\).

\item The cohomology-jump locus
\[
 \{M\in\operatorname{Pic}(X):h^2(X,T_X\otimes M)\ge1\}
\]
is locally determinantal in this family
\cite[Chapter~III, \S4, Theorem~4.1, p.~124]{BS76push}.  It is proper because \eqref{wbII:eq:W:5.24} gives a
point outside it.  Its complement meets the two endpoint
vanishing loci in a dense open set, where \eqref{wbII:eq:W:5.25} holds and
\[
 \chi(X,T_X\otimes M)=-h^1(X,T_X\otimes M)\le0.
\tag{W.55}\hypertarget{tag-03HQ}{}\label{wbII:eq:W:5.28}
\]

\item Every \(M\in\operatorname{Pic}(X)\) has rationally trivial first Chern
class.  The hypotheses also give \(c_1(T_X)=0\) in rational cohomology and,
by Poincar\'e duality, \(H^4(X;\mathbb Q)=0\), so
\(c_2(T_X)=0\) rationally.  The degree-six term of
\(\operatorname{ch}(T_X\otimes M)\operatorname{td}(X)\) is consequently
\(\tfrac12c_3(T_X)\).  Hirzebruch--Riemann--Roch and \eqref{wbII:eq:W:5.28} give the inequalities
in \eqref{wbII:eq:W:5.26}:
\[
 \chi(X,T_X\otimes M)
 =\frac12\left\langle c_3(T_X),[X]\right\rangle\le0.
\]

\end{enumerate}
\end{proof}



\begin{corollary}[Relative vanishing criterion]\hypertarget{tag-03DT}{}\label{wbII:cor:corrected-cdp-relative}
\StatementRecord{wbII:cor:corrected-cdp-relative}{proved}{03DT}{\SeriesRef{wbII:thm:corrected-cdp-global}, \SeriesRef{wbII:eq:W:5.24}; \cite[Proposition~2.4, p.~680]{CDP20}}
\label{wbII:cor:relative}
In the setting of
\cref{wbII:thm:corrected-cdp-global}, it is enough that one line bundle
\(L\in\operatorname{Pic}(X)\) satisfy
\[
 R^2f_*(T_X\otimes L)=0.
\tag{W.56}\hypertarget{tag-03HR}{}\label{wbII:eq:W:5.29}
\]
The relative criterion is \cite[Proposition~2.4, p.~680]{CDP20}.
\end{corollary}



\begin{proof}
Let \(A\) be a positive-degree line bundle on the compact Riemann surface
\(C\), and put \(E=T_X\otimes L\).  Projection formula gives
\[
 R^2f_*(E\otimes f^*A^m)=R^2f_*E\otimes A^m=0.
\]
For \(m\gg0\), Serre vanishing on \(C\) gives
\[
 H^1(C,R^1f_*E\otimes A^m)=0.
\]
The Leray spectral sequence has only columns zero and one, so its
total-degree-two terms yield
\[
 H^2(X,E\otimes f^*A^m)=0.
\]
Thus \(L_0=L\otimes f^*A^m\) satisfies \eqref{wbII:eq:W:5.24}, and
\cref{wbII:thm:corrected-cdp-global} applies.

\end{proof}



\begin{theorem}[Reduced-fibre vanishing criterion]
\hypertarget{tag-03IZ}{}\label{wbII:thm:corrected-cdp-fibrewise}
\StatementRecord{wbII:thm:corrected-cdp-fibrewise}{proved}{03IZ}{\SeriesRef{wbII:thm:corrected-cdp-global}, \SeriesRef{wbII:eq:W:5.29}, \SeriesRef{wbII:eq:W:5.25}, \SeriesRef{wbII:eq:W:5.26}, \SeriesRef{cdp:st:A.1.1}, \SeriesRef{wbII:cor:surface-top-base-change}, \SeriesRef{wbII:cor:corrected-cdp-relative}}

In the setting of
\cref{wbII:thm:corrected-cdp-global}, assume every full scheme-theoretic fibre
\(W_c=f^{-1}(c)\) is reduced, and fix \(L\in\operatorname{Pic}(X)\).  Put
\[
 A_c=(\omega_X\otimes L^\vee)|_{W_c},\qquad
 \mathscr T_c=\ker\!\left(
 \Omega^1_{W_c}\longrightarrow
 j_*\Omega^1_{(W_c)_{\rm reg}}\right).
\]
For each \(c\), let
\[
 \delta_c:
 H^0(W_c,\mathscr T_c\otimes A_c)
 \longrightarrow
 H^1(W_c,N^*_{W_c/X}\otimes A_c)
\tag{W.57}\hypertarget{tag-03HT}{}\label{wbII:eq:W:5.30}
\]
be the connecting map of the cotangent filtration in \cref{cdp:st:A.1.1}.  If, for every
\(c\),
\[
 H^0(W_c,\widetilde\Omega^1_{W_c}\otimes A_c)=0,\qquad
 H^0(W_c,N^*_{W_c/X}\otimes A_c)=0,\qquad
 \delta_c\ \text{is injective},
\tag{W.58}\hypertarget{tag-03HU}{}\label{wbII:eq:W:5.31}
\]
then \eqref{wbII:eq:W:5.29}, \eqref{wbII:eq:W:5.25}, and \eqref{wbII:eq:W:5.26} hold.
\end{theorem}



\begin{proof}
Apply \cref{cdp:st:A.1.1} to the whole reduced
Cartier fibre \(W_c\).  Conditions \eqref{wbII:eq:W:5.31}
give
\[
 H^0(W_c,\Omega_X^1|_{W_c}\otimes A_c)
 =\ker\delta_c=0.
\tag{W.59}\hypertarget{tag-03HV}{}\label{wbII:eq:W:5.32}
\]
The map \(f\) is flat: a local parameter at \(c\) pulls back to a
non-zero-divisor in every regular local ring of the connected manifold \(X\),
and torsion-freeness over the base DVR is flatness.  Moreover, the normal
line of the full fibre is the pullback of the one-dimensional vector space
\(T_cC\), hence is trivial.  Adjunction therefore identifies
\(\omega_{W_c}\simeq\omega_X|_{W_c}\).  With \(E=T_X\otimes L\), absolute
Cartier-fibre duality gives
\[
 H^2(W_c,E|_{W_c})^\vee
 \simeq H^0(W_c,\Omega_X^1|_{W_c}\otimes A_c)=0.
\]
The uniform top-degree base-change isomorphism of
\cref{wbII:cor:surface-top-base-change} therefore gives
\[
 (R^2f_*E)_c\otimes_{\mathcal O_{C,c}}\mathbb C(c)=0
\]
for every \(c\).  Nakayama's lemma kills every stalk of the coherent sheaf
\(R^2f_*E\), proving \eqref{wbII:eq:W:5.29}.  The remaining conclusions follow from
\cref{wbII:cor:corrected-cdp-relative}.

\end{proof}



\paragraph{scope and sharpness of the repair}

Under the two endpoint vanishings in \eqref{wbII:eq:W:5.31}, injectivity of \(\delta_c\) is
equivalent to the desired ambient vanishing \eqref{wbII:eq:W:5.32}; it is the exact lifting condition.  For a reducible reduced fibre the
map must be computed on the whole Cartier fibre.  A nonreduced fibre requires
additional nilpotent-thickening layers and is not covered by
\cref{wbII:thm:corrected-cdp-fibrewise}.  The global criterion
\cref{wbII:thm:corrected-cdp-global} and the relative criterion
\cref{wbII:cor:corrected-cdp-relative} do allow nonreduced fibres.

The restandard family violates both added conditions rather than being
silently excluded by terminology.  Since \(c_3(X)=2\), the contrapositive of
\cref{wbII:thm:corrected-cdp-global} proves
\[
 H^2(X,T_X\otimes L_0)\ne0
 \quad\text{for every }L_0\in\operatorname{Pic}(X).
\]
Independently, \cref{wbII:cor:candidate-conductor-actual} proves
\[
 R^2f_*(T_X\otimes L)\ne0
 \quad\text{for every }L\in\operatorname{Pic}(X),
\]
and its cusp conductor supplies a nonzero element of the relevant
\(\ker\delta_c\).  The first conclusion rules out the global hypothesis; the
second independently rules out the stronger relative sufficient condition.
Thus the corrected theorem is compatible with, and
precisely separates, the verified counterexample to CDP20
Theorem~2.2 as stated.

\Uses{the hypotheses, notation and definitions stated here; \cref{wbII:eq:W:5.31}; \cref{wbII:eq:W:5.32}; \cref{wbII:thm:corrected-cdp-fibrewise}; \cref{wbII:thm:corrected-cdp-global}; \cref{wbII:cor:corrected-cdp-relative}; \cref{wbII:cor:candidate-conductor-actual}; \cref{wbII:eq:W:5.33}}


\paragraph{The exact monodromy repair.}
The split-extension argument can reach only the invariant integral lattice
\(H^1(F;\mathbb Z)^{\pi_1(C^\circ)}\).  Its final complex-linear inference is
valid precisely when
\[
 \operatorname{Span}_{\mathbb C}\operatorname{im}\!\left(
 H^1(F;\mathbb Z)^{\pi_1(C^\circ)}
 \longrightarrow H^1(F,\mathcal O_F)\right)
 =H^1(F,\mathcal O_F).
\tag{W.60}\hypertarget{tag-03HW}{}\label{wbII:eq:W:5.33}
\]
Trivial monodromy on \(H_1(F;\mathbb Z)\) is a stronger sufficient
hypothesis.  Neither \eqref{wbII:eq:W:5.33} nor trivial monodromy follows from the hypotheses
of CDP20 Theorem~2.2.  The corrected theorem above bypasses this rank route
entirely; it does not relabel the false direct-sum formula as proved.

\begin{corollary}[Positive Euler number and cohomology nonvanishing]
\hypertarget{tag-03IV}{}\label{wbII:test:cdp-comparison}
\StatementRecord{wbII:test:cdp-comparison}{proved}{03IV}{\SeriesRef{wbII:thm:corrected-cdp-global}, \SeriesRef{wbII:cor:corrected-cdp-relative}, \SeriesRef{wbII:cor:candidate-conductor-actual}; \citet[Section~10.3, Theorem~10.5 and Corollary~10.6, pp.~88--89]{AC2026}}

For the standard compact threefold $X$, every $L\in\operatorname{Pic}(X)$
satisfies
\[
 H^2(X,T_X\otimes L)\ne0,
 \qquad R^2f_*(T_X\otimes L)\ne0.
\]
The second conclusion also follows directly from the cusp conductor.
\end{corollary}


\begin{proof}
The standard compact smooth threefold constructed in Paper I has
$H^1(X;\mathbb Z)=H^2(X;\mathbb Z)=0$, holomorphic algebraic reduction
to $\mathbb P^1$, and $\langle c_3(T_X),[X]\rangle=2$.
If the displayed degree-two cohomology vanished for one line bundle,
\cref{wbII:thm:corrected-cdp-global} would give $c_3(X)\le0$, a contradiction.
If the second direct image vanished for one twist,
\cref{wbII:cor:corrected-cdp-relative} would give such a global twist after
tensoring by the pullback of a sufficiently positive base line bundle,
again a contradiction. Independently,
\cref{wbII:cor:candidate-conductor-actual} constructs a nonzero cusp-fibre
cohomology class and identifies it with the fibre of the second direct
image by top-degree base change. The resulting counterexample to the CDP conclusion is due to \citet[Section~10.3, Theorem~10.5 and Corollary~10.6, pp.~88--89]{AC2026}.

\end{proof}


\subsection{Ordinary relative one-forms on the actual local models}
The following computation concerns the relative cotangent module on the
smooth total space, before restriction to a fibre. All base maps and local
rings are specified in the statement.
\begin{proposition}[Crossing and multiple-fibre modules]
\hypertarget{tag-03IP}{}\label{wbII:prop:relative-local}
\StatementRecord{wbII:prop:relative-local}{proved}{03IP}{the definitions and constructions stated in the proof}

The symbols $R_2$ and $R_3$ denote copies of the same regular local ring,
with different maps from the one-dimensional base:
let $R_2=\C\{x,y,z\}$ with base coordinate $t=xy$, and let
$R_3=\C\{x,y,z\}$ with base coordinate $t=xyz$.  Then
\[
\begin{aligned}
 M_2&:=\Omega^1_{R_2/\C\{t\}}
 =\frac{R_2\,dx\oplus R_2\,dy\oplus R_2\,dz}
        {R_2\,(y\,dx+x\,dy)},\\[2pt]
 M_3&:=\Omega^1_{R_3/\C\{t\}}
 =\frac{R_3\,dx\oplus R_3\,dy\oplus R_3\,dz}
        {R_3\,(yz\,dx+xz\,dy+xy\,dz)}.
\end{aligned}
\tag{W.176}\hypertarget{tag-03E6}{}\label{wbII:eq:CN:6.1}
\]
Both modules are torsion-free, but neither is locally free at the origin.
They embed in free modules by
\[
\begin{aligned}
 M_2&\hookrightarrow R_2e\oplus R_2dz,
 &dx&\longmapsto xe,&dy&\longmapsto-ye,&dz&\longmapsto dz,\\
 M_3&\hookrightarrow R_3e_1\oplus R_3e_2,
 &dx&\longmapsto xe_1,&dy&\longmapsto ye_2,&dz&\longmapsto-z(e_1+e_2),
\end{aligned}
\tag{W.177}\hypertarget{tag-03E7}{}\label{wbII:eq:CN:6.2}
\]
with images
\[
 R_2dz\oplus (x,y)e,
 \qquad
 R_3xe_1+R_3ye_2+R_3z(e_1+e_2).
\tag{W.178}\hypertarget{tag-03E8}{}\label{wbII:eq:CN:6.3}
\]
If $R_m=\C\{s,y_1,y_2\}$ and the base coordinate is $t=s^m$ with
$m\ge2$, then
\[
\Omega^1_{R_m/\C\{t\}}
\simeq R_m\,dy_1\oplus R_m\,dy_2
       \oplus R_m/(s^{m-1})\,ds,
\tag{W.179}\hypertarget{tag-03E9}{}\label{wbII:eq:CN:6.4}
\]
and the final summand is exactly its torsion submodule.
\end{proposition}



\begin{proof}

\begin{enumerate}[label=\arabic*.,leftmargin=*]
\item Differentiating $t=xy$ and $t=xyz$ gives the presentations in \eqref{wbII:eq:CN:6.1}.  For
$M_2$, the first map in \eqref{wbII:eq:CN:6.2} sends $(a,b,c)$ to $(ax-by,c)$.  If this is
zero, then $c=0$ and $ax=by$.  Since $x$ and $y$ are relatively prime in the
UFD $R_2$, there is an $h$ with $b=xh$ and $a=yh$; hence its kernel is exactly
the displayed relation.  For $M_3$, the second map sends $(a,b,c)$ to
$(ax-cz,by-cz)$.  If it is zero, relative primeness first gives $x\mid c$,
then $y\mid(c/x)$; writing $c=xyh$ gives
$(a,b,c)=h(yz,xz,xy)$, again exactly the defining relation.  Thus both maps
are injective, proving torsion-freeness and \eqref{wbII:eq:CN:6.3}.

\item At the maximal ideal $(x,y,z)$ each defining relation has all coefficients in
that maximal ideal.  Consequently the quotient of either module by the
maximal ideal has dimension $3$ over $\C$, while its generic rank is $2$;
a locally free module of rank $2$ would need only $2$ generators.  Finally, for $t=s^m$ one has
$dt=m s^{m-1}ds=0$ relatively.  The integer $m$ is a unit in $\C$, so the
decomposition \eqref{wbII:eq:CN:6.4} follows.  Its last summand is annihilated by
$s^{m-1}$, while the other two summands are free; after inverting $s$ it
vanishes.  It is therefore exactly the torsion submodule.

\end{enumerate}
\end{proof}



\paragraph{scope of the local correction}

\hypertarget{tag-03IR}{}\label{wbII:rem:relative-local-scope}
The modules in \eqref{wbII:eq:CN:6.1} are modules on the smooth total-space chart, and their
torsion-freeness is asserted before restricting to a fibre.  They must be
distinguished both from the intrinsic sheaf $\Omega^1_S$ of the singular
fibre and from the restricted relative module
$M_i\otimes_{R_i}R_i/(t)$ ($i=2,3$): tensoring with the fibre ring can create
conductor-supported torsion, as the $\tau$ calculation in the proof of
Theorem~\ref{cdp:st:B1} shows.  If a local quotient is finite
surjective \'{e}tale over a chart in \eqref{wbII:eq:CN:6.4}, relative differentials commute with
this base change and faithful flatness then reflects torsion of the pulled-back
module; a ramified quotient
requires a separate invariant-module calculation.  Thus the bare $t=xyz$
chart has a torsion-free but non-locally-free total-space relative module for
$p=1$; this fact alone does not decide torsion after fibre restriction or
quotient.  The multiple-fibre charts do have the total-space torsion in
\eqref{wbII:eq:CN:6.4}.  No claim about torsion in $\bigwedge^2\Omega^1$ or
$\bigwedge^3\Omega^1$, and no global direct-image or Hodge conclusion, is
made here.

\Uses{the hypotheses, notation and definitions stated here; \cref{wbII:eq:CN:6.1}; \cref{cdp:st:B1}; \cref{wbII:eq:CN:6.4}}





\section{The Liu--Rollenske parameter covering}

\begin{theorem}[The retained covering-parameter result]\hypertarget{tag-01RO}{}\label{ii:LR:covering}
\StatementRecord{ii:LR:covering}{proved}{01RO}{\SeriesRef{cdp:st:E1}, \SeriesRef{night:P3F.41}, \SeriesRef{night:P3F.47}, \SeriesRef{iv:correction:P3K.35}, \SeriesRef{night:P3K.74}; \cite[Lemma~2.7, Corollary~2.8, Proposition~2.12, Proposition~3.6 and Definition~3.7]{LiuRollenske2026}}
Retain the exact scalar $u$ in the Liu--Rollenske action $B_u=uB_0$,
the good fourth root at the $III^*$ source fibre, the specified
logarithmic vector $a=e_4/4$, its fixed line character, and
the retained periodic cusp fan with its source and line attachment maps.
The resulting construction is a holomorphic family on the universal
logarithmic cover of a sufficiently small punctured $u$-disc.
The threefold reached after one positive $u$-turn is not
biholomorphic to the initial threefold. The actual family descends
to the degree-two parameter cover
\[
                         u=w^2.\tag{LR.1}
\]
This degree is minimal. The conclusion holds with every retained
constant root-line character.
\end{theorem}
\begin{proof}
\par


\begin{enumerate}[label={$\langle1\rangle$\arabic*.},leftmargin=*]
\item \textbf{The actual central torus and its scalar}

\par\nobreak

The source's equation (2.3) is the lift of translation by
\(-2P\), multiplied by \(u\). On the minimal good root the marked
section specializes to \((1+i)/2\) on
\(E=\mathbb C/(\mathbb Z+i\mathbb Z)\). Its line is the nonzero
order-two line with character
\[
                  m+in\longmapsto(-1)^{m+n}.          \tag{LR.4}
\]
These are the explicit statements and coordinate maps in
\cite[Lemma~2.7 and Corollary~2.8, pp.~6--7 of the attached-source edition]{LiuRollenske2026}. The same marked root and corrected line are constructed by the source and Cartier maps of Paper~IV, \SeriesRef{night:P3F.41} and \SeriesRef{night:P3F.47}.

Since \(2P\) specializes to zero, the restriction of \(B_0\) to
this line complement is multiplication by a constant \(c\ne0\).
The constant is independent of \(u\): every source, line, and root
comparison was fixed before multiplying the action by \(u\).
Thus its actual scalar multiplier is
\[
              \alpha=cu,\qquad
              \tau=\frac{\log(cu)}{2\pi i},
              \qquad\operatorname{Im}\tau>0.         \tag{LR.5}
\]
The line can be presented on the universal cover of \(E\) as
\((z,v)\sim(z+m+in,(-1)^{m+n}v)\). Writing
\(v=e^{2\pi i\eta}\), and then quotienting by \(B_u\), gives
exactly the central abelian surface
\[
 A_\tau=\mathbb C^2/\Lambda_\tau,
 \quad
 \Lambda_\tau=
 \left\langle
  (2,0),(1+i,0),(1,1/2),(0,\tau)
 \right\rangle_{\mathbb Z}.                           \tag{LR.6}
\]
The displayed lattice map is literal: its generators are the two
periods in the kernel of (LR.4), the remaining character period,
and the actual scalar quotient period. In particular
\[
 \delta=(0,1)=2(1,1/2)-(2,0),\qquad e_4=(0,\tau).
\]

The root deck acts on the source by \(z\mapsto-iz\). In this flat
line presentation an actual lift is
\((z,\eta)\mapsto(-iz,\eta+\kappa)\), with
\(4\kappa\in\mathbb Z\). Indeed (LR.4) is invariant under
\(-i\); the ratio of any two lifts is a nowhere-zero function on
the compact elliptic curve, hence a scalar; its fourth power is one.
The flat lift displayed in the proof of \cite[Proposition~3.9, p.~10 of the attached-source edition]{LiuRollenske2026} has \(\kappa=1/4\). Keeping \(\kappa\) records this actual lift and every retained constant finite line character. The logarithmic vector
is exactly \(a=(0,\tau/4)\), by \cite[Definition~3.7, p.~10 of the attached-source edition]{LiuRollenske2026}. Therefore the reduced multiple fibre is
\[
 B_{\tau,\kappa}=A_\tau/\langle R\rangle,
 \qquad R(z,\eta)=(-iz,\eta+\tau/4+\kappa).           \tag{LR.7}
\]
This action has order four and is free: a nontrivial power has a
nonintegral multiple of \(\operatorname{Im}\tau\) in its second
coordinate, whereas the imaginary second-coordinate periods of
\(\Lambda_\tau\) are integral multiples of
\(\operatorname{Im}\tau\). Its fourth power is the lattice period
\((0,\tau+4\kappa)\).

For reference, the linear deck in the ordered lattice basis of
(LR.6) is
\[
 D=\begin{pmatrix}
 1&1&0&0\\-2&-1&-1&0\\0&0&1&0\\0&0&0&1
 \end{pmatrix}.
\]
Its invariant lattice is exactly
\(\mathbb Z\delta\oplus\mathbb Z e_4\), agreeing with the
\cite[Proposition~3.6, pp.~9--10 of the attached-source edition]{LiuRollenske2026}.
A positive \(u\)-turn has the exact comparison
\[
 \tau\longmapsto\tau+1,
 \quad \Lambda_{\tau+1}=\Lambda_\tau,
 \quad a\longmapsto a+\delta/4.                      \tag{LR.8}
\]
The unchanged abelian surface in (LR.8) does not identify the affine
quotients: the torsion translation in (LR.7) has changed.

\par

\item \textbf{An intrinsic invariant of the reduced fibre}

\par\nobreak

The actual holomorphic projection
\[
 \pi_{\tau,\kappa}:B_{\tau,\kappa}\longrightarrow
 E_{\tau,\kappa}:=
 \frac{\mathbb C}{\frac12\mathbb Z+
                       (\frac\tau4+\kappa)\mathbb Z},
 \qquad [(z,\eta)]\longmapsto[\eta]                  \tag{LR.9}
\]
is well-defined by every generator in (LR.6) and (LR.7).
It has connected fibres. More explicitly, the image of the full
universal-cover deck group under its second-coordinate translation
is precisely the displayed rank-two lattice; after choosing a lift
of a target point, its zero-second-coordinate subgroup has elliptic
kernel generated by \((2,0),(1+i,0)\). This is the actual fibre of
(LR.9), so it is connected.

We also verify that (LR.9) is the Albanese map, without a choice of
marking. Holomorphic one-forms on \(A_\tau\) are constant linear
forms: the coefficients of their pullback to \(\mathbb C^2\) are
periodic entire functions and hence constant by compactness. The
\(R\)-invariant subspace is exactly \(\mathbb C\,d\eta\), since
\(R^*dz=-i\,dz\) and \(R^*d\eta=d\eta\). The periods of
\(d\eta\) on the full fundamental group of \(B_{\tau,\kappa}\)
are exactly \(\frac12\mathbb Z+(\frac\tau4+\kappa)\mathbb Z\):
this follows directly from the same affine deck generators. The
integration definition of the Albanese torus therefore gives
precisely (LR.9), up to a target translation. In particular its
elliptic-curve isomorphism class is intrinsic to \(B_{\tau,\kappa}\).

Normalize its first period to one by \(\eta\mapsto2\eta\). Its
Tate parameter is
\[
 Q=\exp\bigl(2\pi i(\tau/2+2\kappa)\bigr)
    =e^{4\pi i\kappa}\sqrt{cu}.                     \tag{LR.10}
\]
The factor \(e^{4\pi i\kappa}\) is a fixed sign. Under (LR.8)
this parameter changes to \(-Q\).

For completeness, the classical Weierstrass lattice invariant has
\(j(Q)=Q^{-1}+744+O(Q)\). The leading term alone suffices and can
be obtained directly from the convergent lattice sums: summing the
geometric series in the cotangent partial fraction gives
\[
 E_4(Q)=1+240\sum_{n\ge1}\sigma_3(n)Q^n,
 \quad E_6(Q)=1-504\sum_{n\ge1}\sigma_5(n)Q^n,
 \quad
 j(Q)=\frac{1728E_4(Q)^3}{E_4(Q)^3-E_6(Q)^2}.
\]
Their denominator is \(1728Q+O(Q^2)\), so \(Qj(Q)\) extends
holomorphically with value one at zero. Consequently
\[
 Q\bigl(j(Q)-j(-Q)\bigr)
\]
extends holomorphically with value two. It is nonzero on a
sufficiently small disc. Thus
\[
            B_{\tau,\kappa}\not\simeq
            B_{\tau+1,\kappa}                       \tag{LR.11}
\]
for every sufficiently small nonzero \(u\), even without preserving
any of their markings. A biholomorphism would induce an isomorphism
of their Albanese elliptic curves, whereas (LR.10) has different
Weierstrass invariant after the turn.

\par

\item \textbf{Every threefold endpoint must preserve this fibre}

\par\nobreak

We give the intrinsic argument independently of any canonical-bundle
formula. The actual source/cap Mayer--Vietoris calculation in Paper~IV, \SeriesRef{night:P3F.41} and \SeriesRef{night:P3F.47},
proves \(H^2(Y;\mathbb R)=0\) for these sphere members. Let
\(f:Y\to\mathbb P^1\) be their constructed connected torus fibration.
Every effective divisor on \(Y\) is vertical. Otherwise a horizontal
irreducible component meets a general smooth torus fibre \(F\) in
a nonzero effective divisor \(C\). Its real first Chern class is
zero, because it is the restriction of the corresponding class on
\(Y\). Every compact complex torus admits a translation-invariant
Kähler form \(\omega\), and the divisor identity gives
\(\int_F c_1(\mathcal O(C))\wedge\omega=\int_C\omega>0\), a
contradiction.

Let \(g\) be a meromorphic function on \(Y\). Its poles are
therefore vertical. On a general fibre it has no divisorial poles;
normality and Hartogs extension remove the remaining codimension-two
indeterminacy. It is a holomorphic function on a compact connected
torus and is constant. This gives a meromorphic function on the
regular part of the base. Near a bad base value \(p\), multiply
\(g\) by a sufficiently large integer power of \(f^*(t-p)\) to
cancel every polar order on the finitely many fibre components.
The product is holomorphic by normality. It is constant on compact
connected fibres and hence descends holomorphically: its descended
function is bounded near \(p\) by properness, and the removable
singularity theorem applies. Dividing back proves
\[
                   \mathbb C(Y)=f^*\mathbb C(t).      \tag{LR.12}
\]

A biholomorphism between two such threefolds induces an automorphism
of this meromorphic function field, hence a Möbius map of the base.
Indeed a rational function with a rational compositional inverse
has degree one. Thus it fits into an actual commutative square with
their maps \(f\). Multiplicities of fibres are preserved, because a
Möbius map has nonzero local derivative. The unique multiple fibre
has multiplicity four, so its reduced support must map biholomorphically
to the unique reduced multiple fibre in the target. This contradicts
(LR.11) after a single positive \(u\)-turn. This argument excludes
all possible holomorphic endpoint maps, including source/base
symmetries and source inversion; no assumed translation normal form
for an endpoint is needed.

\par


\item \textbf{The actual map for the double turn.}\par\nobreak
We display the retained construction, including its scalar norm.
Write $G=P$ for the marked source section and
$\mathcal M=\mathcal O(G-O)$ for the line used in the source.
On the generic elliptic fibre the rational function
\[
 f=\frac1{x-x(G)},\qquad
 \operatorname{div}(f)=O+O-G-(-G)
                       =t_G^*(G-O)-(G-O)
\]
gives a line isomorphism covering translation by $+G$.
Let $U$ be its meromorphic bundle lift, with the prescribed
vertical Cartier corrections retained. Since $B_0^{-1}$
covers translation by $+2G$, their exact ratio is a rational
base function
\[
              \rho(t)=B_0^{-1}/U^2.\tag{LR.13}
\]
Indeed the ratio is a nonvanishing holomorphic function on the
generic compact elliptic fibre, hence constant there; the
explicit divisor maps make that constant rational in $t$.
At the selected $I_1$ value, write $e$ for the vertical order
of $U$. The actual linearization has a simple zero there, so
\[
                  \operatorname{ord}\rho=-1-2e.\tag{LR.14}
\]
At an uncapped retained fibre the corresponding order is $-2e$.
The integral component-intersection kernel is generated by the
full reduced fibre; thus removing that integral whole-fibre
power makes $U$ invertible there.

On the good fourth root let $z$ be the source-root coordinate,
let $D$ be its deck, and let $L$ be the extending lift of $+G$.
The source and line correspondence of Paper~IV constructs this
lift on the actual corrected root line. The point $G(0)$ has
order two. The fibre weights of $\mathcal O(G-O)$ at $G(0)$
and $O(0)$ are the deck derivative and its inverse. Their
ratio is $-1$, since the derivative has order four. A retained
constant character multiplies both weights by the same number
and leaves the ratio unchanged. Hence
\[
                 DLD^{-1}=-L.\tag{LR.15}
\]
Initially the commutator may be $-c(z)$ with $c(0)=1$ and
$\prod_{j=0}^3c(i^jz)=1$. Put $g=\log c$ and
$\mathcal D h(z)=h(-iz)$. The explicit function
$h=-\tfrac14\sum_{j=0}^3j\mathcal D^jg$ satisfies
$(\mathcal D-I)h=-g$. Rescaling the target of $L$ by $e^h$
gives (LR.15). The unit $L^2/B_0^{-1}$ is then deck invariant;
its invariant analytic square root further normalizes
\[
                              L^2=B_0^{-1}.\tag{LR.16}
\]
This rescaling preserves (LR.15).

Write $U=z^eh(z)L$ on the root line, with $h$ a nonzero unit.
Because $U$ is defined over the original base, its deck
conjugation gives $e\equiv2\pmod4$ and $h(iz)=h(z)$.
In particular
\[
                    \operatorname{ord}_{t=0}\rho=-e/2
                         \equiv1\pmod2.\tag{LR.17}
\]
Let $h_0(t)$ be the actual positive one-gamma closing multiplier
of the degree-four parameter family $u=r^4$. It has order
$-1$ at this cap, order $+1$ at the selected $I_1$, and no
other orders except the retained ordinary framing orders,
which are divisible by four. Equations (LR.14) and (LR.17)
show that $h_0\rho$ has even order at every point of
$\mathbb P^1$. Factoring its zeros and poles and choosing the
specified cusp value therefore gives a rational function
$g_0$ with
\[
                             g_0^2=h_0\rho.\tag{LR.18}
\]
The actual ordinary-chart endpoint covering $r\mapsto-r$ is
\[
                    \widetilde D_r=-r^{-2}g_0(t)U.\tag{LR.19}
\]
The minus sign is the scalar norm correction and is part of
this map.

On the good-root cap the apparent $z^{-2}$ factor in
(LR.19) is cancelled by the actual ordinary-to-cap gluing
change under $r\mapsto-r$. The remaining map is an invariant
unit times $L$, so it extends invertibly through the cap.
Writing its action as $A_r=rA_0$, equation (LR.15) gives
\[
               LA_rL^{-1}=-A_r=A_{-r}.\tag{LR.20}
\]
It also commutes with $B_u$, since $(-r)^4=r^4$.
At the retained order-one datum the deck is the identity,
and the ordinary comparison extends directly. These are
maps of the actual quotients, with their finite characters.

The cusp comparison includes both source and line periods.
Put $v=\log r/(2\pi i)$ and write the actual bounded period
expressions as
\[
 H=s+U_0(q),\qquad
 \gamma=(\Gamma,C-s-4v),\qquad
 \tau(v)=\frac{H\delta+\gamma}{2}.
\]
The endpoint lift is $\tau(v)+\delta/2$.
Its two steps satisfy the literal identity
\[
 [\tau(v)+\delta/2]+[\tau(v+1/2)+\delta/2]
                       =H\delta+\gamma.\tag{LR.21}
\]
Every cusp Cox character is multiplied by its bounded
holomorphic unit; this extends through every retained cone
and face of the chosen periodic fan. On the ordinary chart
the identical norm is
\[
 \widetilde D_{-r}\widetilde D_r
    =r^{-4}g_0^2U^2
    =h_0 B_u^{-1}.\tag{LR.22}
\]
Thus the second iterate is exactly the prescribed one-gamma
gluing, including the nonconstant multiplier $h_0$.

First quotient the universal $v$-family by this one-gamma
map to obtain the proper $r$-family. Equation (LR.21) then
makes (LR.19) an order-two action on that family, covering
$r\mapsto-r$. It is free because $r\ne0$. Its quotient
is a proper holomorphic family on the punctured $w$-disc,
where $w=r^2$ and $u=w^2$. The cap and cusp extensions above
prove smoothness of its total space. Steps 2--3 exclude
every biholomorphic endpoint for one positive $u$-turn, so
the degree two is minimal.

\end{enumerate}
\end{proof}

\begin{refuted}[The historical all-fan clause of P3K.36]\hypertarget{tag-02UH}{}\label{night:P3K.36}\label{stmt:P3K.36}
\StatementRecord{night:P3K.36}{false}{02UH}{\SeriesRef{night:P3K.73}, \SeriesRef{night:P3K.74}}
\par\noindent\textbf{Refined receiving item:} \SeriesRef{iv:correction:P3K.36}. The counterexample below uses the displayed marked fan and its actual source and cap maps.\par

The earlier clause assigned the value
$3(S-1)(S-3)/8\pmod {28}$ to every fan of the corrected
$u=w^2$ family, where $S=p^2-pq+q^2$ was the norm of the
physical delta cocharacter. The actual fan $T_2$ has $S=1$,
so that expression gives zero. The retained degree-four
page has $(H,\chi)=(2,1)$ and its degree-two quotient has
$(J,\widehat H,\widehat\chi)=(2,1,1/2)$.
The actual fan maps in Paper~IV, \SeriesRef{night:P3K.73} and
\SeriesRef{night:P3K.74}, give
\[
 (\widehat\chi-\widehat H/2)-J\,2(2-1)/2
                           =-2=26\pmod {28}.
\]
This is the stated counterexample to the all-fan clause.
The original fixed reference fan retains its separately
proved zero value; the single $u$-turn has no endpoint
identification by \ref{ii:LR:covering}.
\end{refuted}

\clearpage
\part{The sign theorem for every degeneration}
\section{The degeneration and its normalization maps}
In the global sign statements of this part, $f:Y\to B$ is a proper holomorphic map with connected fibres from a smooth compact complex threefold. The sign theorem assumes $b_1(Y)=b_2(Y)=0$ and a smooth general fibre that is a complex two-torus. Local comparison constructions state their own hypotheses; they are used with precisely the displayed maps. A full fibre is its Cartier scheme. Its reduced support is denoted $S$, its normalization by $\nu:T\to S$, and a resolution by $\pi:Z\to T$. Euler characteristic means topological Euler characteristic unless $\chi(\mathcal O)$ is written.

\section{Unrestricted Euler topology and component covers}\hypertarget{tag-0342}{}\label{sec:II:06}


\begin{proposition}[Every torus-degeneration fibre has even Euler number]\hypertarget{tag-01CM}{}\label{cdp:st:C.2.UN.1}
\StatementRecord{cdp:st:C.2.UN.1}{proved}{01CM}{\SeriesRef{cdp:eq:C.UN.3}, \SeriesRef{cdp:eq:C.UN.4}}

Let $f:Y\to B$ be a proper holomorphic map from a smooth complex
threefold to a smooth complex curve, with connected fibres and smooth
general fibre a complex two-torus. For every $b\in B$,
\[
                    e((Y_b)_{\rm red})\in2\Z.       \tag{C.UN.1}\hypertarget{tag-0152}{}\label{cdp:eq:C.UN.1}
\]
The fibres may have arbitrary multiplicities and singularities.
\end{proposition}


\begin{proof}

\begin{enumerate}[label=\arabic*.,leftmargin=*]
\item Take a sufficiently small closed coordinate disc $\Delta$ about $b$,
containing no other critical value. Put $U=f^{-1}(\Delta)$ and
$M=\partial U$. Then $U$ is a compact oriented smooth real six-manifold
with boundary, and $M$ is the mapping torus of the orientation-preserving
monodromy of a real four-torus. The analytic central support $S$ is a
deformation retract of $U$. One precise construction uses a controlled
gradient for the real analytic function $|f|^2$ on a Whitney
stratification: it decreases $|f|$, is tangent to the strata over zero,
and has trajectories with a limit on the zero set by the local
analytic gradient inequality. Compactness of the zero set gives a
common small disc, and the resulting continuous limit homotopy fixes
the zero set. Equivalently the triangulated proper local tube is a
regular neighbourhood of the zero fibre. Nilpotents do not change this
retraction. Thus $e(S)=\chi(U)$.

\item We give the parity calculation entirely over $\Q$. Write
$s_i=\operatorname{rank}(H_i(U)\to H_i(U,M))$.
Exactness of the pair sequence, with Poincar\'e--Lefschetz duality,
gives
\[
 b_0(M)+b_1(M)+b_2(M)
       =\sum_{i=0}^{5}b_i(U)-s_0-s_3-2s_1-2s_2.
                                                        \tag{C.UN.2}\hypertarget{tag-0153}{}\label{cdp:eq:C.UN.2}
\]
Indeed, the summand $b_i(M)$ is the sum of the ranks of
$H_i(M)\to H_i(U)$ and $H_{i+1}(U,M)\to H_i(M)$;
substitution of $b_{i+1}(U,M)=b_{5-i}(U)$ gives the formula.
Here $s_0=0$ and $b_6(U)=0$. The map defining $s_3$ is represented,
after duality, by the middle-dimensional intersection pairing of $U$.
That pairing is skew-symmetric because the middle dimension is three,
so $s_3$ is even. Consequently
\[
 \chi(U)\equiv b_0(M)+b_1(M)+b_2(M)\pmod2.            \tag{C.UN.3}\hypertarget{tag-0154}{}\label{cdp:eq:C.UN.3}
\]

\item Let $T$ act on $H^1(F,\Q)$ for the four-torus fibre $F$, and set
$a_i=\dim\ker(\Lambda^iT-I)$. Cutting the mapping torus along one
fibre gives the Wang exact sequence; hence
\[
 b_0(M)=1,\quad b_1(M)=1+a_1,\quad b_2(M)=a_1+a_2.
                                                        \tag{C.UN.4}\hypertarget{tag-0155}{}\label{cdp:eq:C.UN.4}
\]
Orientation preservation says $\det T=1$. The transformation
$G=\Lambda^2T$ preserves the nondegenerate symmetric wedge pairing
on $\Lambda^2\Q^4$ and has determinant $(\det T)^3=1$.
The rank $\operatorname{rank}(G-I)$ is even. Complexify and split off the generalized
$(-1)$-eigenspace. Its dimension is even: all eigenvalues outside
$\{1,-1\}$ occur in reciprocal pairs with equal algebraic
multiplicity, and $\det G=1$. This generalized eigenspace is
nondegenerate and orthogonal to its invariant complement. On it
$G-I$ is invertible, of even rank. On the complement $G+I$ is
invertible, and
\[
                         (G-I)(G+I)^{-1}
\]
is skew-adjoint for the symmetric wedge pairing. Multiplication by
that nondegenerate pairing represents it as a skew-symmetric matrix,
whose rank is even. Its rank equals that of $G-I$ on the complement.
Since $\dim\Lambda^2\Q^4=6$, $a_2$ is
even. Equations \eqref{cdp:eq:C.UN.3}--\eqref{cdp:eq:C.UN.4} prove \eqref{cdp:eq:C.UN.1}.

\end{enumerate}
\end{proof}






\begin{proposition}[All even Euler numbers and an unbounded zero-Euler family]\hypertarget{tag-01CX}{}\label{cdp:st:C.2.UN.2}
\StatementRecord{cdp:st:C.2.UN.2}{proved}{01CX}{\SeriesRef{cdp:st:C2b}, \SeriesRef{cdp:st:C.2.UN.1}}

The set of Euler numbers of reduced fibres among compact smooth
threefolds fibred in two-tori is exactly $2\Z$. For each $k\ge1$
there is a smooth projective such threefold with one singular reduced
normal-crossing fibre of Euler number zero consisting of a blown-up
abelian surface, $k$ ruled surfaces over a genus-two curve, and $k$
projective planes. Their incidence graph is a chain of the ruled
components attached to the abelian component, together with $k$
projective-plane leaves attached to that component.
\end{proposition}


\begin{proof}

\begin{enumerate}[label=\arabic*.,leftmargin=*]
\item Use the explicit embedded genus-two curve $C_2\subset A$ of \ref{cdp:st:C2b}.
In $A\times\PP^1$ let
$R=C_2\times\PP^1$ be the smooth horizontal surface. First blow up $C_2\times\{0\}$, as in \ref{cdp:st:C2b}.
The strict transform of $R$ is still $R$, because its centre is a
Cartier divisor in $R$. It meets the exceptional ruled surface in
the section corresponding to the base-normal direction; this
section is disjoint from the intersection with the strict abelian
component. The map to the base is a submersion near this section.

\item Blow up that section and repeat, always using the intersection of
the strict transform of $R$ with the newest exceptional component.
The next centre is a smooth genus-two curve in the smooth locus of
the reduced central fibre, disjoint from its double curves. Its
normal sequence splits: the normal line to the curve inside $R$
maps isomorphically under $df$ to the tangent line of the base,
and supplies a splitting of the normal sequence of the central
surface. The two line summands give two disjoint sections on the
new exceptional projective bundle. One meets the previous
component; the other is the next centre. This proves by induction
that all $k$ curve blowups are defined and that the reduced central
support has the stated chain. In local coordinates the comparison
is always $t=u\tau$ as in \ref{cdp:st:C2b}, so every new component has
multiplicity one and no triple point is introduced.

\item Finally blow up $m\ge0$ distinct points of the strict abelian
component away from the genus-two double curve. They introduce
$m$ projective-plane leaves, each meeting the blown-up abelian
component in one projective line. All maps are proper holomorphic
blowdowns followed by the original projection. Every nonzero
fibre is biholomorphic to $A$.

\item Each curve blowup increases the Euler number of the central
support by $e(C_2)=-2$: the exceptional projective bundle has Euler
number $2e(C_2)$ and its intersection with the old component has
Euler number $e(C_2)$. Each point blowup increases it by two,
because the old surface gains one, the exceptional plane adds
three, and their intersection subtracts two. Thus
\[
                         e(S)=2(m-k).                 \tag{C.UN.5}\hypertarget{tag-0156}{}\label{cdp:eq:C.UN.5}
\]
Given any integer $n$, take $k\ge\max(1,-n)$ and $m=k+n$.
This realizes $2n$; \ref{cdp:st:C.2.UN.1} supplies the opposite inclusion.
Taking $m=k$ gives the asserted zero-Euler family.

\item The integral cohomology of the total space is torsion-free. The
excision and projective-bundle comparison of \ref{cdp:st:C2b} gives
\[
 (b_0,\ldots,b_6)
      =(1,4,7+k+m,8+4k,7+k+m,4,1).                  \tag{C.UN.6}\hypertarget{tag-0157}{}\label{cdp:eq:C.UN.6}
\]
Its Euler characteristic is independently $2(m-k)$. The structure
sheaf direct images of each blowup are $\OO$ in degree zero and
zero in positive degrees, so
$(h^0(\OO),h^1(\OO),h^2(\OO),h^3(\OO))=(1,2,1,0)$.


\end{enumerate}
\end{proof}






\begin{proposition}[The same punctured torus family has non-bimeromorphic zero-Euler fillings]\hypertarget{tag-01D8}{}\label{cdp:st:C.2.UN.3}
\StatementRecord{cdp:st:C.2.UN.3}{proved}{01D8}{\SeriesRef{cdp:st:D.2.7}, \SeriesRef{cdp:eq:C.UC.2}, \SeriesRef{cdp:eq:C.UC.1}, \SeriesRef{cdp:eq:C.UC.4}, \SeriesRef{cdp:eq:C.UC.3}, \SeriesRef{cdp:eq:C.UC.6}, \SeriesRef{cdp:eq:C.UC.8}, \SeriesRef{cdp:eq:C.UC.9}, \SeriesRef{cdp:eq:C.UC.10}, \SeriesRef{cdp:eq:C.UC.7}; \cite[Section~1.2.4 and Lemma~1.3.1]{AKMW2002}}

Let $A=E\times E$, $E=\C/(\Z+\mathrm i\Z)$.
There are two proper smooth-total-space fillings of
$A\times\Delta^*$ over a disc whose reduced central
supports both have Euler number zero, but for which
no bimeromorphic map over the disc exists. One is
$A\times\Delta$. The other is the local filling
\eqref{cdp:eq:D.23}, whose central support consists of a primary
Kodaira surface and a primary Hopf surface meeting
along one elliptic curve. The nonexistence persists
after every finite base change and resolution of the
normalized base change, with the given punctured
family fixed.

These fillings occur in compact smooth threefolds:
the principal $A$-bundle $Y_0$ and the threefold $Y$
constructed in \SeriesRef{cdp:st:D.2.7} are isomorphic away from the
three specified fibres. That isomorphism does not
extend bimeromorphically, even after a finite base
change of their base curve. In particular smooth
torus fillings, together with modifications by
blowups and blowdowns over the central fibre, do
not generate all zero-Euler fillings of this
punctured family.

More strongly, the same punctured product has
infinitely many pairwise non-bimeromorphic
zero-Euler fillings. They are obtained by replacing
$\OO_V(R)$ in \eqref{cdp:eq:C.UC.2} by $\OO_V(nR)$,
$n\ge1$. This distinction also persists after
arbitrary finite base change and resolution.
\end{proposition}


\begin{proof}

\begin{enumerate}[label=\arabic*.,leftmargin=*]
\item We first record the precise birational comparison
used in the proof. Let $\phi:M\dashrightarrow N$
be a bimeromorphic map between smooth complex
threefolds, and let $P\subset M$ be a compact
irreducible divisor whose smooth resolution is
not covered by rational curves. Then $\phi$
cannot contract $P$ to a subset of dimension
at most one. Resolve the inverse meromorphic
map by smooth-centre blowups:
\[
 \widehat N\xrightarrow{\ a\ }N,
 \qquad b:\widehat N\longrightarrow M,
 \qquad b=\phi^{-1}a.
\tag{C.UC.1}\hypertarget{tag-014C}{}\label{cdp:eq:C.UC.1}
\]
The elimination-of-indeterminacy construction
has this form in the analytic category as
well as the algebraic one; see
\cite[Section~1.2.4 and Lemma~1.3.1]{AKMW2002}.
For maps of proper disc models, the same
construction is made on a smaller disc; its
centres above the compact central fibre are
compact. For the compact threefolds in the
statement the cited construction applies
directly.

\item Since $b$ is a proper modification of the
smooth target $M$, the strict transform
$\widehat P$ maps birationally onto $P$.
The analytic justification is as follows. The locus in
$M$ over which $b$ has positive-dimensional
fibres is analytic by properness. It cannot
contain a divisor: its preimage would then
have dimension at least three while being
a proper analytic subset of the irreducible
threefold $\widehat N$. Thus $b$ is finite
over a dense open subset of $P$. A finite
modification of the normal manifold $M$
is an isomorphism, so its degree there is
one. If $\phi$ contracted
$P$, then $a$ would contract $\widehat P$.
It would be an exceptional divisor of the
composition of smooth-centre blowups $a$.
Every such divisor is birational to the
exceptional divisor created at one of those
steps. In dimension three that divisor is
$\PP^2$ for a point centre, or a
$\PP^1$-bundle over a smooth curve for a
curve centre. A divisor centre gives an
isomorphism and creates no exceptional
divisor. Hence $\widehat P$, and therefore
$P$, would be covered by rational curves,
a contradiction. Applying the same argument
to the inverse shows that the surviving
divisors correspond birationally.

\item Put
$q=e^{-2\pi}$ and
\[
 V=\operatorname{Bl}_{(0,0)}(E\times\Delta),
 \qquad R=\operatorname{Exc}(V\to E\times\Delta),
 \qquad U=\OO_V(R)^\times/q^{\Z}.
\tag{C.UC.2}\hypertarget{tag-014E}{}\label{cdp:eq:C.UC.2}
\]
The quotient map $U\to V$ is a principal
elliptic bundle with fibre $\C^*/q^{\Z}=E$.
It is smooth and proper, since each bundle
chart is the product with this compact
elliptic curve. Composition with $V\to\Delta$
is therefore a proper map from a smooth
threefold. The canonical section of
$\OO_V(R)$ trivializes it on the punctured
disc, giving the specified isomorphism
\[
                         U|_{\Delta^*}=A\times\Delta^*.
\tag{C.UC.3}\hypertarget{tag-014F}{}\label{cdp:eq:C.UC.3}
\]
This is exactly the construction in \SeriesRef{cdp:st:D.2.7},
including its holomorphic gluing to the
compact threefold.

\item The strict transform $E_0$ of $E\times\{0\}$
and the exceptional $R$ meet transversely at
one point. The restrictions of $\OO_V(R)$
to them have degree $+1$ and $-1$,
respectively. Their inverse images in $U$
are
\[
 P=\OO_E(0)^\times/q^{\Z},\qquad
 Q=\OO_{\PP^1}(-1)^\times/q^{\Z}
       =(\C^2\setminus\{0\})/\langle q\rangle.
\tag{C.UC.4}\hypertarget{tag-014G}{}\label{cdp:eq:C.UC.4}
\]
They meet in a single elliptic fibre $J$,
and the local equation of the map is
$xy=t$ along $J$. The first surface is a
primary Kodaira surface with $b_1(P)=3$:
in its principal elliptic bundle the primitive
Chern class kills one of the two elliptic
fibre cycles. The second is the indicated
primary Hopf surface, with $b_1(Q)=1$.
The identifications and the integral
transgression are calculated in
\eqref{cdp:eq:D.23}--\eqref{cdp:eq:D.25}.

\item Neither surface contains a rational curve.
A holomorphic map $\PP^1\to P$ projects
constantly to its elliptic base and then
constantly to its elliptic fibre. A map
$\PP^1\to Q$ lifts to $\C^2\setminus\{0\}$,
since $\PP^1$ is simply connected; its two
holomorphic coordinate functions are constant.
Thus their resolutions are not covered by
rational curves. A bimeromorphic comparison
from $U$ to $A\times\Delta$ over $\Delta$
would, by \eqref{cdp:eq:C.UC.1}, have to carry each of
$P,Q$ birationally onto a central divisor.
The product has only the one central
divisor $A$, and neither $P$ nor $Q$ is
birational to it. Indeed the first Betti
number of a smooth compact surface is
unchanged by point blowups: a ball is
replaced by the simply connected exceptional
neighbourhood, and van Kampen gives the
same fundamental group. Resolution of a
surface bimeromorphism therefore preserves
$b_1$, whereas $b_1(A)=4$ differs from
both values in \eqref{cdp:eq:C.UC.4}. This proves that
there is no such bimeromorphic comparison.

\item All these supports have Euler number zero.
For $P$ and $Q$, their principal elliptic
bundle projections are locally trivial
proper bundles, so their Euler numbers are
the product of the Euler numbers of base
and fibre, hence zero. The intersection
$J$ also has Euler number zero. Additivity
gives
\[
                    e(P\cup Q)=e(P)+e(Q)-e(J)=0=e(A).
\tag{C.UC.5}\hypertarget{tag-014H}{}\label{cdp:eq:C.UC.5}
\]

\item Consider a finite base change.
Locally such a change has equation
$t=u(s)s^m$, $u(0)\ne0$. Taking a local
$m$th root of the unit replaces its
parameter by one with $t=s^m$.
At a general point of either $P$ or $Q$
the original map is a submersion with
coordinate $x=t$. Its base change is the
smooth chart
\[
                          x=s^m,
\tag{C.UC.6}\hypertarget{tag-014I}{}\label{cdp:eq:C.UC.6}
\]
with the two surface coordinates unchanged.
The component over $s=0$ thus maps with
degree one to the original component.
At their intersection the normalized base
change has equation $xy=s^m$, with one
additional smooth elliptic coordinate.
This is normal, with transverse rational
double point $A_{m-1}$ when $m>1$.
A resolution adds ruled exceptional
surfaces above $J$; the strict transforms
of $P$ and $Q$ remain birational to the
surfaces in \eqref{cdp:eq:C.UC.4}. No choice of further
resolution changes those birational
classes. The target base change remains
the smooth product $A\times\Delta_s$.
The argument using \eqref{cdp:eq:C.UC.1} and their
first Betti numbers again excludes every
bimeromorphic comparison over $\Delta_s$.

\item Finally use the compact gluing of \SeriesRef{cdp:st:D.2.7}.
The original principal $A$-bundle $Y_0$
has a smooth abelian central fibre at
each of the three chosen base points.
The replacement $Y$ has the local model
\eqref{cdp:eq:C.UC.2} at each point and is identified
with $Y_0$ elsewhere by \eqref{cdp:eq:C.UC.3}.
An extension of this isomorphism would
respect the base, since the two base
maps agree on its dense open graph.
Its restriction near any of the three
fibres would be the already excluded
comparison. The same argument applies
above every preimage of these points
under a finite cover of the compact base,
using \eqref{cdp:eq:C.UC.6}. This proves the compact
and base-change assertions.

\item For the infinite family put
\[
 U_n=\OO_V(nR)^\times/q^{\Z},\qquad
 P_n=\OO_E(n[0])^\times/q^{\Z},\qquad
 Q_n=\OO_{\PP^1}(-n)^\times/q^{\Z}.
\tag{C.UC.7}\hypertarget{tag-014J}{}\label{cdp:eq:C.UC.7}
\]
The canonical section again trivializes the
bundle off the central fibre. Thus the punctured
isomorphism is the same exact product \eqref{cdp:eq:C.UC.3}.
The central support is $P_n\cup_J Q_n$;
all three terms have Euler number zero by their
principal elliptic bundle projections. Each
$P_n$ contains no rational curve by its
projection to $E$ and then to its elliptic
fibre. For $Q_n$, the usual quotient
$\OO_{\PP^1}(-n)^\times
=(\C^2\setminus\{0\})/\boldsymbol\mu_n$
is explicit: a nonzero vector $v$ represents
$v^{\otimes n}$ in the $n$th tensor power of
the tautological line. Multiplication by $q$
on that tensor power lifts to multiplication
by the positive real number $q^{1/n}$ on
$\C^2\setminus\{0\}$. Hence
\[
 Q_n=(\C^2\setminus\{0\})/
       \langle\boldsymbol\mu_n,q^{1/n}\rangle,
 \qquad
 \pi_1(Q_n)=\Z\oplus\Z/n\Z.
\tag{C.UC.8}\hypertarget{tag-014K}{}\label{cdp:eq:C.UC.8}
\]
The two commuting generators are independent,
because a nonzero power of $q^{1/n}$ has
absolute value different from one. The action
is free and properly discontinuous, and its
covering space retracts to the simply connected
$S^3$. This proves the fundamental-group
calculation and also proves that every map
$\PP^1\to Q_n$ is constant by lifting it to
the cover.

\item For $P_n$, the corresponding calculation
also determines the extension class. A Hermitian
metric on $\OO_E(n[0])$ identifies its nonzero
total space smoothly with its unit-circle bundle
times $\R_{>0}$. As $q$ is positive real, its
quotient is the unit-circle bundle times
$\R_{>0}/q^{\Z}\simeq S^1$. The circle bundle
has Chern number $n$. Lifting the two generators
of $\pi_1(E)$, the attaching map of the base
two-cell has winding number $n$ in the circle
fibre. With a consistent orientation this gives
the exact presentation
\[
 \pi_1(P_n)=
 \langle a,b,c,d\mid
 [a,b]=c^n,\ [a,c]=[b,c]=1,
 [a,d]=[b,d]=[c,d]=1\rangle.
\tag{C.UC.9}\hypertarget{tag-014L}{}\label{cdp:eq:C.UC.9}
\]
There are no further relations: the homotopy
sequence of the circle bundle has injective
fibre group because $\pi_2(E)=0$, and the last
factor is the displayed direct $S^1$ factor.
Abelianizing \eqref{cdp:eq:C.UC.8}--\eqref{cdp:eq:C.UC.9} yields
\[
 H_1(P_n,\Z)=\Z^3\oplus\Z/n\Z,
 \qquad H_1(Q_n,\Z)=\Z\oplus\Z/n\Z.
\tag{C.UC.10}\hypertarget{tag-014D}{}\label{cdp:eq:C.UC.10}
\]
Integral first homology, including its torsion,
is invariant under a bimeromorphism of smooth
compact surfaces, by the same point-blowup
comparison already used above. By \eqref{cdp:eq:C.UC.1},
a bimeromorphism $U_n\dashrightarrow U_{n'}$
over the disc must carry the two non-uniruled
central divisors birationally to the two such
divisors on the target. The ranks in \eqref{cdp:eq:C.UC.10}
force $P_n$ to correspond to $P_{n'}$ and
$Q_n$ to $Q_{n'}$. The torsion orders then
force $n=n'$. After any finite base change,
\eqref{cdp:eq:C.UC.6} still leaves these components
birationally unchanged, and the resolution of
$xy=s^m$ adds only ruled surfaces over $J$.
Further resolutions add only uniruled
exceptional divisors. They cannot replace a
non-uniruled component under \eqref{cdp:eq:C.UC.1}, so
the same torsion argument proves persistence.
Finally all $U_n$ have compact realizations:
replace $V$ by
$\operatorname{Bl}_{(0,0)}(E\times\PP^1)$
and use the same line bundle $\OO(nR)$.
Its elliptic quotient is a compact smooth
threefold, and the projection to $\PP^1$
restricts to \eqref{cdp:eq:C.UC.7} near zero.

\end{enumerate}
\end{proof}



The comparison obstruction specifies the limit
of a particular proposed classification route.
In \cite[Sections~5.1--5.3]{KKN2008}, a
nondegenerate logarithmic torus pairing is
given first, and a suitable complete fan
produces a proper model of that logarithmic
complex torus. The theorem there supplies
those models; an extension of an arbitrary
punctured-family isomorphism to a
bimeromorphic map of its chosen fillings
is an additional requirement. Statement
\ref{cdp:st:C.2.UN.3} refutes that additional requirement
for the displayed fillings, even though
the punctured family is a product.




\begin{proposition}[Every even Euler number occurs above each fixed torus filling]\hypertarget{tag-01DJ}{}\label{cdp:st:C.2.UN.4}
\StatementRecord{cdp:st:C.2.UN.4}{proved}{01DJ}{\SeriesRef{cdp:eq:C.VR.4}, \SeriesRef{cdp:eq:C.VR.6}, \SeriesRef{cdp:eq:C.VR.8}, \SeriesRef{cdp:eq:C.VR.9}, \SeriesRef{cdp:st:C.2.UN.1}}

Let $f:Y\to B$ be a proper holomorphic map from a smooth complex
threefold to a smooth complex curve, with connected fibres and smooth
general fibre a complex two-torus. Fix $b\in B$ and put
$S=(Y_b)_{\rm red}$. For every integer $k$ there is a finite composition
of blowups with smooth compact centres in successive fibres over $b$,
\[
 \pi:\widetilde Y\longrightarrow Y,\qquad
 \widetilde f=f\pi,\qquad
 e((\widetilde Y_b)_{\rm red})=2k.                 \tag{C.VR.1}\hypertarget{tag-015W}{}\label{cdp:eq:C.VR.1}
\]
The map $\pi$ is an isomorphism over $B\setminus\{b\}$.
In particular every such filling is dominated by a smooth filling
whose central reduced support has Euler number zero. More precisely,
every further smooth model obtained by vertical blowups admits a
zero-Euler refinement by additional vertical blowups. Thus zero-Euler
models are cofinal for this domination relation, in the explicit sense
that every object is dominated by one of them.
Above every given filling there are such zero-Euler refinements with
arbitrarily many irreducible components of their central support.

An elementary operation lowering the Euler number by two consists of
one point blowup and one smooth genus-three curve blowup. On integral
cohomology its additional summands are $\Z^2$ in degrees two and four
and $\Z^6$ in degree three. In particular, whenever the Betti numbers
are finite,
\[
 (\Delta b_2,\Delta b_3,\Delta b_4)=(2,6,2),
 \qquad \Delta b_j=0\quad(j\notin\{2,3,4\}).       \tag{C.VR.2}\hypertarget{tag-015X}{}\label{cdp:eq:C.VR.2}
\]
The positive full-fibre multiplicities are unchanged on
their strict transforms.
\end{proposition}


\begin{proof}

\begin{enumerate}[label=\arabic*.,leftmargin=*]
\item Let $Z\subset S$ be a smooth compact
centre of complex codimension $r$ in the smooth ambient threefold,
and let $\pi_Z:Y'\to Y$ be its blowup. Write $S'=(f\pi_Z)^{-1}(b)_{\rm red}$.
The equality of underlying sets is $S'=\pi_Z^{-1}(S)$, and the
restrictions of the blowup give
\[
 S'\setminus\pi_Z^{-1}(Z)\simeq S\setminus Z,
 \qquad \pi_Z^{-1}(Z)=\PP(N_{Z/Y}).                \tag{C.VR.3}\hypertarget{tag-015Y}{}\label{cdp:eq:C.VR.3}
\]
Compact complex analytic spaces admit finite triangulations, compatible
with these closed analytic subsets. Euler additivity, applied with
compact supports on the complements, and the projective-bundle
formula therefore give
\[
 \begin{split}
 e(S')-e(S)
 &=e(\PP(N_{Z/Y}))-e(Z)\\
 &=(r-1)e(Z).
 \end{split}                                      \tag{C.VR.4}\hypertarget{tag-015Z}{}\label{cdp:eq:C.VR.4}
\]
The projective-bundle Euler equality also follows from its locally
trivial fibre $\PP^{r-1}$, whose Euler number is $r$.
Equation \eqref{cdp:eq:C.VR.4} applies to the full inverse image, including
intersections of the centre with other components and the singular support.

\item Choose a smooth point $p$ of one component of $S$, away from every
other component. Such points form a nonempty open subset of each
irreducible component. If the full fibre has multiplicity $m>0$ on
that component, local holomorphic coordinates may be chosen with
\[
                    t\circ f=x^m.                \tag{C.VR.5}\hypertarget{tag-0160}{}\label{cdp:eq:C.VR.5}
\]
Indeed the local equation initially has the form $x^m a$, where $a$
is a holomorphic unit; a local holomorphic $m$th root of $a$ replaces
$x$ by $x a^{1/m}$. Blow up $p$. The exceptional surface
$E\simeq\PP^2$ meets the strict component in the line
$L=\{X=0\}\subset\PP(T_pY)$, with homogeneous coordinates
$[X:Y:Z]$ induced by $(x,y,z)$. It meets no other central component.
In the chart $y=r$, $x=ru$, $z=rv$, the map is
\[
                         t=(ru)^m.                \tag{C.VR.6}\hypertarget{tag-0161}{}\label{cdp:eq:C.VR.6}
\]
Consequently both $E=(r=0)$ and the strict component $(u=0)$
have multiplicity $m$. Formula \eqref{cdp:eq:C.VR.4}, with $r=3$ and centre a
point, shows that this step increases the central Euler number by two.

\item Inside $E$ take the Fermat quartic
\[
 C=\{X^4+Y^4+Z^4=0\}\subset\PP^2.                \tag{C.VR.7}\hypertarget{tag-0162}{}\label{cdp:eq:C.VR.7}
\]
Its three partial derivatives cannot vanish simultaneously at a
projective point, so it is smooth. It is connected: if the quartic
were a disjoint union of distinct plane components, two such
components would meet by their positive intersection number; their
intersection would be a singular point. Adjunction gives
$K_C=\OO_C(1)$ and $2g(C)-2=4$, hence $g(C)=3$ and $e(C)=-4$.
On $L$ its equation is $Y^4+Z^4=0$. It has four distinct roots,
none at $Z=0$, and the two gradients of the quartic and $X$ are
independent there. Thus $C$ meets $L$ transversely in exactly four
points.

\item Now blow up the smooth curve $C$ in the threefold obtained after the
first blowup. Its exceptional divisor is the projective line bundle
$\PP(N_{C/Y_1})\to C$. Equation \eqref{cdp:eq:C.VR.4}, with $r=2$, gives Euler
change $e(C)=-4$. The two successive changes sum to $-2$.
The ambient blowup is smooth at the four points of $C\cap L$. Near any
one of them, preserve the local equations $r=0$ of $E$ and $u=0$ of
the strict old component. Transversality gives a third coordinate $v$
such that $C=(r=v=0)$. Absorbing a remaining unit in $r$ preserves
these equations and gives $t=(ru)^m$ as in \eqref{cdp:eq:C.VR.6}. The two charts of
the curve blowup are
\[
 \begin{array}{c|c|c}
 \text{chart}&\text{blowdown substitution}&\text{base map}\\ \hline
 \text{$r$ chart}&v=rw&t=(ru)^m\\
 \text{$v$ chart}&r=vw&t=(uvw)^m .
 \end{array}                                      \tag{C.VR.8}\hypertarget{tag-0163}{}\label{cdp:eq:C.VR.8}
\]
In the second chart the new exceptional divisor, the strict transform
of $E$, and the strict old component have equations $v=0$, $w=0$,
and $u=0$, respectively. They meet with normal crossings and all
have multiplicity $m$. Away from $C\cap L$ the comparison is the
same two-coordinate blowup of $t=r^m$. The strict transform of $E$
is isomorphic to $E$, since $C$ is Cartier in $E$; the strict old
component is its point blowup at the four transverse intersection
points. These identifications follow directly from the charts in
\eqref{cdp:eq:C.VR.8}. Outside the specified neighbourhoods the maps are
isomorphisms, so any other singularities of the original support
are unchanged.

\item For a smooth centre of codimension $r$, let $j:E_Z\hookrightarrow Y'$
and $p:E_Z\to Z$ denote the exceptional inclusion and projection,
and put $\xi=c_1(\OO_{E_Z}(1))$. The blowup homomorphism is
\[
 \begin{split}
 H^q(Y,\Z)\oplus\bigoplus_{i=1}^{r-1}H^{q-2i}(Z,\Z)
   &\longrightarrow H^q(Y',\Z),\\
 (\alpha,\beta_1,\ldots,\beta_{r-1})
   &\longmapsto \pi_Z^*\alpha+
       \sum_{i=1}^{r-1}j_*(\xi^{i-1}p^*\beta_i).
 \end{split}                                      \tag{C.VR.9}\hypertarget{tag-0164}{}\label{cdp:eq:C.VR.9}
\]
It is an isomorphism. To check it integrally, excision identifies the
relative cohomology of $(Y',E_Z)$ with that of $(Y,Z)$. The
projective-bundle theorem makes $H^q(E_Z,\Z)$ the direct sum of
$\xi^i p^*H^{q-2i}(Z,\Z)$ for $0\le i<r$, and the self-intersection
formula gives $j^*j_*\eta=-\xi\eta$. Thus the extra terms of
\eqref{cdp:eq:C.VR.9} map isomorphically, with sign minus, onto the summands with
$i>0$. Comparing the relative long exact sequences proves
surjectivity; their remaining $i=0$ term is exactly the pullback from
$Y$. The proper degree-one trace makes $\pi_Z^*$ injective, and
restriction to $E_Z$ then proves independence of all the displayed
extra summands. This proves \eqref{cdp:eq:C.VR.9}. The same argument applies on a
compact fibre tube with boundary, since the centre is in its interior
and the blowup is unchanged on the boundary.

\item The point blowup adds $\Z$ in each of degrees two and four. The
genus-three curve blowup adds $H^{q-2}(C,\Z)$, whose ranks in these
three successive degrees are $1,6,1$. This proves \eqref{cdp:eq:C.VR.2} and shows
that pre-existing integral torsion is unchanged.

\item Both steps are proper holomorphic maps from smooth threefolds. A
blowup has connected fibres and pushes its structure sheaf to the
structure sheaf of the original smooth space. Stein factorization
therefore proves that their composition with $f$ again has connected
fibres. The centres lie over $b$, so the resulting map is the same
proper torus family over $B\setminus\{b\}$.

\item Proposition~\ref{cdp:st:C.2.UN.1} gives $e(S)=2n$ for an integer $n$. To obtain $2k$
with $k\ge n$, use $k-n$ point blowups. To obtain $2k$ with $k<n$,
use $n-k$ of the two-step operations just proved. At each stage a
smooth point of a single component exists, and the construction
can be performed there; only finitely many centres are needed.
Applying this same procedure to any given vertical blowup model
proves cofinality with the displayed domination maps.
After reaching Euler number zero, compose one further point blowup
with one of the two-step decreasing operations, at smooth
points. Their net Euler change is zero, while they add three
irreducible central components. Repetition gives arbitrarily many
components above the original filling. Each repetition adds
$(3,6,3)$ to $(b_2,b_3,b_4)$ when those numbers are finite; the
maps are the successive blowdowns.

\end{enumerate}
\end{proof}






\begin{proposition}[The same non-uniruled components occur with every even Euler number]\hypertarget{tag-01CR}{}\label{cdp:st:C.2.UN.14}
\StatementRecord{cdp:st:C.2.UN.14}{proved}{01CR}{\SeriesRef{cdp:st:C.2.UN.4}, \SeriesRef{cdp:eq:C.EC.2}}

Fix a proper holomorphic family $p^\times:\mathcal T^\times\to\Delta^\times$
of connected complex two-tori. A marked smooth filling is a proper
holomorphic map $f:Y\to\Delta$ from a smooth complex threefold with
connected fibres, together with a biholomorphism
$f^{-1}(\Delta^\times)\simeq\mathcal T^\times$ over $\Delta^\times$.
For each irreducible component $S_i$ of $(Y_0)_{\rm red}$, let
$Z_i\to S_i^\nu\to S_i$ be a smooth compact resolution of its
normalization. Write $\mathcal E(f)$ for the finite multiset of
bimeromorphism classes $[Z_i]$ for which $Z_i$ is not uniruled.
Here a smooth compact surface is uniruled when it admits a dominant
meromorphic map $C\times\PP^1\dashrightarrow Z$ for a smooth compact
curve $C$. Repetitions in the multiset record distinct components.

For every integer $k$, the following equality holds between sets of
finite multisets:
\[
 \begin{split}
 \{\mathcal E(f): f\text{ is a marked smooth filling of }p^\times\}
 ={}&\{\mathcal E(f): f\text{ is such a filling},\\[-2pt]
 &\hspace{48mm}e((Y_0)_{\rm red})=2k\}.
 \end{split}                                                     \tag{C.EC.1}\hypertarget{tag-00FO}{}\label{cdp:eq:C.EC.1}
\]
For each filling in the left-hand set, the representative on the
right is obtained by the explicit proper holomorphic vertical
blowup map of \ref{cdp:st:C.2.UN.4}. That map induces a bijection between the
original non-uniruled components and the resulting non-uniruled
components, and a bimeromorphism of their normalization resolutions.
Every original full-fibre multiplicity is preserved by this bijection.

The same assertion holds for compact smooth total spaces over a
fixed compact base curve, with the marked family fixed away from a
finite set $Q$. One can prescribe an arbitrary even Euler number at
each point of $Q$ simultaneously, preserving the multiset and
multiplicity at each point. In particular, the collection of
non-uniruled surface classes which occur in arbitrary compact
torus-fibred fillings is exactly the collection which occur in
compact fillings all of whose reduced fibres have Euler number zero.
\end{proposition}


\begin{proof}

\begin{enumerate}[label={$\langle1\rangle$\arabic*.},leftmargin=*]
\item We specify the component comparison for one of the smooth blowups
$\pi:Y'\to Y$ used in \ref{cdp:st:C.2.UN.4}. Its centre $C$ is a point or a
smooth compact curve in the central support. The inverse image of
that support consists of the strict transforms $S_i'$ of its
original components together with the exceptional divisor
\[
 E=\PP(N_{C/Y})\longrightarrow C.                    \tag{C.EC.2}\hypertarget{tag-00FP}{}\label{cdp:eq:C.EC.2}
\]
Each restriction $S_i'\to S_i$ is a proper holomorphic map and is an
isomorphism over the complement of $C$. Since $C$ has smaller
dimension than $S_i$, this complement is dense. Thus this actual
map is bimeromorphic. Normalizing both surfaces and resolving them
gives a bimeromorphism $Z_i'\dashrightarrow Z_i$: explicitly, take
the closure of the graph of the isomorphism on the dense normal
open sets, normalize it, and resolve the resulting compact analytic
surface. Its two proper holomorphic projections resolve the stated
meromorphic map. The pullback of the defining equation of the full
fibre has its original vanishing order at the generic point of
$S_i'$, because $\pi$ is an isomorphism there. Hence its
multiplicity is unchanged.

\item Uniruledness is preserved by these bimeromorphisms. Indeed, composing
a dominant meromorphic map $C_0\times\PP^1\dashrightarrow Z_i$
with either bimeromorphic inverse gives such a map for $Z_i'$; its
dominance follows on the dense open set where that inverse is an
isomorphism. Resolving the graph makes the composition a
well-defined meromorphic map. The same argument in the reverse
direction proves the equivalence.

\item The newly introduced divisor in \eqref{cdp:eq:C.EC.2} is uniruled. For a point
centre it is $\PP^2$, which is rational: the map
$\PP^1\times\PP^1\dashrightarrow\PP^2$ given on the affine chart
by $(u,v)\mapsto[1:u:v]$ is bimeromorphic. For a curve centre it is
a projective line bundle over $C$. Its rank-two bundle is
meromorphically trivial on $C$: tensoring by a sufficiently positive
line bundle supplies two global sections independent at the
generic point, and their ratio gives the meromorphic
trivialization. Projectivizing it gives a bimeromorphism
$C\times\PP^1\dashrightarrow\PP(N_{C/Y})$.
Its subsequent strict transforms remain uniruled by the same
composition argument. Therefore a finite sequence of these blowups
introduces no new non-uniruled class and loses none of the original
classes. It preserves their repetitions and multiplicities as well.

\item Starting from any marked filling, \ref{cdp:st:C.2.UN.4} constructs a finite
sequence of exactly these maps whose final reduced Euler number is
$2k$. Every centre lies over zero, so the blowdown is a
biholomorphism over $\Delta^\times$ and transports the given marking
by composition. We have just proved that its multiset is the
original $\mathcal E(f)$. This gives inclusion of the left side of
\eqref{cdp:eq:C.EC.1} in the right. Inclusion in the other direction follows
because each filling with Euler number $2k$ is among the same
marked fillings defining the left side. This also proves the
assertion when there are no fillings: both sets are empty.

\item For the compact version, apply the same construction successively
at the finitely many points of $Q$. Every centre is a compact
analytic submanifold of the current compact threefold. Its global
blowup is compact and smooth, and the composite with the fixed base
map is proper holomorphic with connected fibres, by the
structure-sheaf and Stein-factor comparison in \ref{cdp:st:C.2.UN.4}. Its
restriction off that point is an isomorphism. Thus a construction
at one point preserves every fibre and every marking already
treated at another point. There are only finitely many critical
values of a compact proper map to a curve which is smooth
generically: the image of its compact analytic critical locus is
a proper analytic subset of the curve. Taking $Q$ to contain them
all and taking $k=0$ at each of them proves the final assertion;
the remaining smooth torus fibres already have Euler number zero.

\end{enumerate}
\end{proof}



For completeness, a surface dominated meromorphically by
$C\times\PP^1$ has no nonzero pluricanonical section. The pullback
of such a section is nonzero, because a dominant meromorphic map
between smooth surfaces is generically of full rank. It is
holomorphic away from the finitely many indeterminacy points and
extends over them by Hartogs' theorem. Its restriction to any
$\{c\}\times\PP^1$ has degree $-2m$ in pluricanonical degree $m>0$
and must vanish. This contradicts its being nonzero. Thus the
equality applies in particular to a specified general-type
normalization class. It makes the occurrence problem for such a
component part of the unrestricted zero-Euler problem through the
explicit maps just proved. It does not supply a comparison between
two independently specified marked fillings. The construction
uses \ref{cdp:st:C.2.UN.4} only and invokes no assertion that such a comparison
extends across the central fibre.




\begin{proposition}[No component of a torus degeneration has general-type normalization]\hypertarget{tag-01CS}{}\label{cdp:st:C.2.UN.15}
\StatementRecord{cdp:st:C.2.UN.15}{proved}{01CS}{\SeriesRef{cdp:eq:C.SC.3}, \SeriesRef{cdp:eq:C.SC.6}; \cite[Theorem~3.3, with its proof]{Guillot2024}; \cite{Guillot2019}}

Let $f:U\to\Delta$ be a proper holomorphic map from
a smooth complex threefold to a disc, with connected
fibres and smooth fibres complex two-tori over
$\Delta^*$. For every irreducible component $S_i$
of the reduced central support, its smooth
normalization resolution $Z_i$ satisfies
\[
                            \kappa(Z_i)\le1.
 \tag{C.SC.1}\hypertarget{tag-010L}{}\label{cdp:eq:C.SC.1}
\]
No K\"ahler, projectivity, Betti-number or Euler-number
assumption is imposed on $U$ or its central fibre.
\end{proposition}


\begin{proof}

\begin{enumerate}[label={$\langle1\rangle$\arabic*.},leftmargin=*]
\item We first make the chosen component reduced by an
actual finite base change. If its full-fibre
multiplicity is $m_i$, put $t=s^{m_i}$ and normalize
the fibre product. At a generic smooth point of
$S_i$, away from the other components, the original
equation is $t=u z^{m_i}$ with $u$ a unit. The
normalized base change has the explicit local maps
\[
                  w^{m_i}=u,\qquad z=s/w.
 \tag{C.SC.2}\hypertarget{tag-010M}{}\label{cdp:eq:C.SC.2}
\]
It is therefore smooth at those points, and every
central component above $S_i$ has multiplicity one
and maps generically finitely onto $S_i$. Choose one
such component and resolve the normalized total
space and central support, isomorphically over
$\Delta^*$ and generically along that component.
Write $f':U'\to\Delta_s$ for the resulting proper
map and $E$ for the smooth strict component. The
space $U'$ is smooth, $E$ still has multiplicity
one, and the smooth nearby fibres are the original
connected tori. Connectedness of all full fibres
follows also from Stein factorization over the
normal disc, since its generic degree is one.

\item If $Z_i$ were of general type, so would be $E$.
Here is the precise comparison. The map $E\to S_i$
factors through normalization and gives a
generically finite meromorphic map $E\dashrightarrow
Z_i$. Resolve its graph to obtain a birational
map $a:\widehat E\to E$ and a generically finite
holomorphic map $b:\widehat E\to Z_i$ between
smooth compact surfaces. The actual Jacobian map
gives
\[
                 K_{\widehat E}=b^*K_{Z_i}+R,
                 \qquad R\ge0.                 \tag{C.SC.3}\hypertarget{tag-010N}{}\label{cdp:eq:C.SC.3}
\]
Pullback of pluricanonical sections is injective.
The generically finite pullback of a big canonical
line is big; equivalently a generically finite
pluricanonical map of $Z_i$ remains generically
finite after composition with $b$. Equation
\eqref{cdp:eq:C.SC.3} therefore makes $K_{\widehat E}$ big.
Birational invariance gives $\kappa(E)=2$.
In particular $E$ is projective, as is every
smooth compact complex surface of general type.

\item We construct a nonzero meromorphic vector field
on $E$ with single-valued solutions. Define the
coherent relative tangent sheaf by the actual
kernel
\[
 \mathcal T_{U'/\Delta_s}
 =\ker\bigl[df':T_{U'}\longrightarrow f'^*T_{\Delta_s}\bigr].
 \tag{C.SC.4}\hypertarget{tag-010O}{}\label{cdp:eq:C.SC.4}
\]
Its proper direct image is torsion-free on the
disc: a section annihilated by a power of $s$
vanishes on the dense open punctured tube and
hence vanishes identically. Its rank is two,
because the global holomorphic vector fields
of a two-torus are its two translation fields.
It is thus a locally free rank-two sheaf after
shrinking the disc. Choose a nonzero section
$v$, viewed by adjunction as an actual
holomorphic vertical vector field on $U'$.
At the generic point of $E$ its order of
vanishing is a finite integer $k\ge0$.
Set
\[
                              w=s^{-k}v.
 \tag{C.SC.5}\hypertarget{tag-010P}{}\label{cdp:eq:C.SC.5}
\]
This is a global meromorphic vector field on
$U'$, with poles confined to the central
divisor. At a generic point of $E$, the
multiplicity-one assertion makes $s$ a local
transverse coordinate. Since $v(s)=0$, write
$v=a(x,y,s)\partial_x+b(x,y,s)\partial_y$.
The integer $k$ is the minimum of the two
coefficient orders along $s=0$. Thus $w$
is holomorphic and nonzero generically on
$E$ and has a well-defined nonzero tangent
meromorphic restriction $w_E$ to $E$.

\item We use the precise closed-univalence theorem
\cite[Theorem~3.3, with its proof]{Guillot2024}:
for a holomorphic vector field on a complex
manifold of any dimension, the locus where
its maximal solutions are single-valued is
closed. Such a field is called semicomplete
when that locus is the whole manifold.
The closure assertion is a path-continuation
statement. A failure at a point is witnessed
by a nonclosed path in a solution leaf whose
time-form integral is zero. Lifting the
time parametrization of this compact path
through finitely many flow charts produces
the same nonclosed zero-time path through
all sufficiently nearby initial points.
This proves openness of the failure locus
and records exactly the comparison used here.

\item Let $P$ be the polar divisor of $w$ and put
$W=U'\setminus P$. On the dense open subset
$U'^*\subset W$, $w$ restricts to a
holomorphic vector field on each compact
torus, multiplied by the constant $s^{-k}$
on that fibre. Each such field is complete:
its complex-time solution stays in the
compact fibre, so local existence extends
it for all times. Consequently $w$ is
semicomplete on all of $W$ by the stated
closedness theorem.

\item The hypersurface $E\cap W$ is invariant.
Indeed $w(s)=0$, and at any point where
$w$ is holomorphic, factoring the local
equation of $s$ into its irreducible
components shows that $w$ preserves the
ideal of each component. A solution
starting in $E\cap W$ therefore stays
there, and its single-valuedness is the
one inherited from $w$. The curve removed
from $E$ by the ambient polar divisor may
contain removable poles of $w_E$. Apply
the same closedness theorem once more on
$E$ minus the actual polar divisor of
$w_E$. Its dense open subset $E\cap W$
is already semicomplete. Hence $w_E$ is
a nonzero semicomplete meromorphic vector
field on the entire compact surface $E$,
in the precise sense of single-valuedness
on its holomorphic locus.

\item We finish by proving that a projective
surface of general type admits no such
field. Suppose first that one maximal
solution of $w_E$ has Zariski-dense image.
The Main Theorem of \cite{Guillot2019},
including the clarification immediately
following it, gives a birational model
which is an abelian surface or admits a
rational or elliptic fibration. The
holomorphic alternative is covered there
by Proposition~22; the dense leaf excludes
a first integral in that alternative.
All these models have Kodaira dimension
at most one, contradicting $\kappa(E)=2$.
Only single-valued maximality, already
proved, is needed in this application;
the solution need not be entire.

\item Here is the other case, with its actual
fibration map. If every solution leaf is
contained in a proper analytic curve,
the closures of the regular leaves are
irreducible algebraic invariant curves
on the projective surface $E$. They
cover its open regular foliation locus.
Use the countably many Hilbert polynomials
and flattening strata of their universal
families to parametrize reduced irreducible
curves. The locus of such integral curves
is constructible. After clearing the
polar divisor, invariance is the vanishing
of the induced normal-component map of
the vector field on the universal curve;
on these strata it is a constructible
condition. Thus we have countably many
irreducible locally closed parameter
strata of integral invariant curves.
Their incidence images in $E$ are
constructible. If none were dense, their
Zariski closures would be countably many
proper closed analytic subsets covering
the regular foliation locus, contrary
to the Baire theorem. Choose a stratum
whose incidence dominates $E$. Its
generic curves are reduced and
irreducible. At a regular point of the
foliation exactly one such irreducible
curve can pass: two would contain the
same local solution germ and would
coincide. Thus the parameter has dimension
one and the incidence map to $E$ is
generically one-to-one. Compactify and
normalize this parameter curve and
close the incidence graph. Inverting
its birational projection to $E$ and
projecting to the parameter gives an
actual meromorphic first integral.
Resolve its indeterminacies and take its
Stein factor to obtain a morphism
\[
                         q:E'\longrightarrow C
 \tag{C.SC.6}\hypertarget{tag-010Q}{}\label{cdp:eq:C.SC.6}
\]
with connected generic fibre a smooth
compact curve. The lifted meromorphic
field remains semicomplete: it agrees
with the old field on the birational
open set, and the same closedness theorem
extends single-valuedness over its full
holomorphic locus. It is a nonzero
tangent meromorphic field on a generic
fibre $G$ of \eqref{cdp:eq:C.SC.6}.

\item Such a field on a compact curve has no
pole. A pole of order $r\ge1$ would
have the local form $z^{-r}a(z)\partial_z$
with $a(0)\ne0$. Its time form is
$z^r a(z)^{-1}dz$, whose primitive has
a zero of order $r+1$. The inverse
solution therefore has $r+1$ branches
around that time value, all lying in
the punctured coordinate disc. This
contradicts semicompleteness. The field
on $G$ is consequently holomorphic
and nonzero, so
$\deg T_G=2-2g(G)\ge0$ and $g(G)\le1$.

\item This also contradicts general type
directly through pluricanonical maps.
Adjunction on the smooth generic fibre
gives $K_{E'}|_G=K_G$, of degree at
most zero. Every linear system
$|nK_{E'}|$ therefore restricts to a
constant map on $G$: a negative-degree
line has no sections, and a degree-zero
line on a connected compact curve has
at most one independent section. A
generically finite pluricanonical map
of the general-type surface $E'$ cannot
contract its generic fibres in this
way. This excludes the all-algebraic-
leaves case as well.

\item Both cases contradict the general-type
consequence of \eqref{cdp:eq:C.SC.3}. The chosen
$Z_i$ cannot have Kodaira dimension
two. Since $S_i$ was arbitrary,
\eqref{cdp:eq:C.SC.1} follows for every component.

\end{enumerate}
\end{proof}






\begin{proposition}[Two commuting semicomplete fields on the actual component cover]\hypertarget{tag-01CT}{}\label{cdp:st:C.2.UN.16}
\StatementRecord{cdp:st:C.2.UN.16}{proved}{01CT}{\SeriesRef{cdp:st:C.2.UN.15}, \SeriesRef{cdp:eq:C.CF.3}, \SeriesRef{cdp:eq:C.CF.5}, \SeriesRef{cdp:eq:C.CF.7}, \SeriesRef{cdp:eq:C.CF.8}; \cite[Theorem~3.3]{Guillot2024}; \cite[Theorem~9.2]{Guillot2024}; \cite[Theorem~0.3]{DOT2001}; \cite[pp.~256 and~260]{DOT2000}; \cite[Propositions~2(ii), 3(ii) and~5]{Inoue1974}}

Let $f:U\to\Delta$ be a proper holomorphic map from a smooth
complex threefold, with connected fibres and with smooth complex
two-tori over $\Delta^*$. Fix an irreducible central component $S_i$
of multiplicity $m_i$. Make the actual base change $t=s^{m_i}$,
normalize, and resolve as in \ref{cdp:st:C.2.UN.15}. Every chosen smooth strict
component $E$ above $S_i$ has multiplicity one and carries a
two-dimensional complex vector space
\[
 \mathfrak a_E\subset
 H^0(E,T_E\otimes_{\mathcal O_E}\mathcal M_E)
 \tag{C.CF.1}\hypertarget{tag-00BN}{}\label{cdp:eq:C.CF.1}
\]
of commuting meromorphic vector fields. Every member of
$\mathfrak a_E$ is semicomplete on its holomorphic locus.
The map $E\to S_i$ factors through normalization and is
proper and generically finite.

There are two exhaustive geometric alternatives for this particular
space of fields. Its generic rank is two, in which case any complex
basis is a commuting meromorphic frame; or its generic rank is one,
in which case an actual resolution of a meromorphic first integral
and its Stein factor give
\[
 \widehat E\xrightarrow{\ q\ }B,
 \qquad \widehat E\xrightarrow{\rho}E,
 \tag{C.CF.2}\hypertarget{tag-00BO}{}\label{cdp:eq:C.CF.2}
\]
where $\rho$ is proper birational, $B$ is a smooth compact curve,
and the connected general fibre of $q$ is rational or elliptic.

When $a(S_i^\nu)=0$, the unique minimal model of $E$ is a
two-torus or a Hopf surface. This last conclusion concerns the
specified generically finite covering component; the maps giving
it are part of the assertion.
\end{proposition}



\begin{proof}

\begin{enumerate}[label={$\langle1\rangle$\arabic*.},leftmargin=*]
\item Write $f':U'\to\Delta_s$ for the resolved base change.
The local equations
$w^{m_i}=u$, $z=s/w$ in (C.SC.2), starting from
$t=u z^{m_i}$, prove both multiplicity one and generic
finiteness. Properness and the universal property of normalization
give the asserted factorization $E\to S_i^\nu\to S_i$.
These are the actual maps used below.

\item Put $\mathcal O=\mathbb C\{s\}$ and
$K=\operatorname{Frac}(\mathcal O)$. Let
\[
 M=(f'_*\mathcal T_{U'/\Delta_s})_0,
 \qquad
 L=\{v\in M\otimes_{\mathcal O}K:
                     \operatorname{ord}_E(v)\ge0\}.
 \tag{C.CF.3}\hypertarget{tag-00BP}{}\label{cdp:eq:C.CF.3}
\]
The order in this formula is the minimum of the coefficient
orders in a local tangent frame at the generic point of $E$.
It is independent of that frame. The direct image is coherent
by properness, is torsion-free because its sections are actual
vector fields on $U'$, and has rank two by base change on the
smooth torus locus. Consequently $M$ is a free rank-two
$\mathcal O$-module. Its sections are identified with actual
holomorphic vertical fields by the definition of direct image.
Any two of them commute: their restrictions to a smooth torus
are translation fields, so their bracket vanishes on the dense
punctured tube and hence everywhere.

\item Here is a finite bound for the saturation in \eqref{cdp:eq:C.CF.3}.
Choose an $\mathcal O$-basis $v_1,v_2$ of $M$.
At a generic point of $E$, multiplicity one makes $s$ a
coordinate, so the two vertical tangent coordinates $x,y$
give
\[
 (v_1\ v_2)=(\partial_x\ \partial_y)A,
 \qquad d=\operatorname{ord}_E(\det A)<\infty.
 \tag{C.CF.4}\hypertarget{tag-00BQ}{}\label{cdp:eq:C.CF.4}
\]
All entries of $A$ are holomorphic there. The determinant is
nonzero because the evaluation of the two translation fields
is an isomorphism on each nearby torus. In particular $d\ge0$.
For $v=a_1(s)v_1+a_2(s)v_2\in L$, the column $Aa$ has
nonnegative order along $E$. The identity
$a=(\det A)^{-1}\operatorname{adj}(A)(Aa)$ shows that
each $a_j$ has order at least $-d$. The order of a function
of $s$ along $E$ equals its $s$-adic order, again by
multiplicity one. We obtain actual inclusions
\[
                        M\subset L\subset s^{-d}M.
 \tag{C.CF.5}\hypertarget{tag-00BR}{}\label{cdp:eq:C.CF.5}
\]
The set $L$ is an $\mathcal O$-submodule. Since $\mathcal O$
is noetherian, \eqref{cdp:eq:C.CF.5} makes it finite; since it is torsion-free
of rank two over the discrete valuation ring $\mathcal O$,
it is free of rank two.

\item Every member of $L$ is a global meromorphic vertical field
on a sufficiently small representative of $U'$, regular at
the generic point of $E$. Its restriction is therefore a
global meromorphic section of $T_E$. Tangency follows from
$v(s)=0$ at the generic smooth point and then as an identity
of meromorphic sections. Restriction defines the exact map
\[
 \operatorname{res}_E:L\longrightarrow
 H^0(E,T_E\otimes\mathcal M_E),
 \qquad
 \ker(\operatorname{res}_E)=sL.
 \tag{C.CF.6}\hypertarget{tag-00BS}{}\label{cdp:eq:C.CF.6}
\]
Indeed, zero restriction means order at least one along $E$;
division by $s$ then belongs to $L$. Conversely, every member
of $sL$ restricts to zero. Thus
\[
                  L/sL\simeq\mathbb C^2
                    \hookrightarrow
                  H^0(E,T_E\otimes\mathcal M_E).
 \tag{C.CF.7}\hypertarget{tag-00BT}{}\label{cdp:eq:C.CF.7}
\]
Define $\mathfrak a_E$ to be its image. Since vertical fields
annihilate functions of $s$, brackets of members of
$M\otimes K$ vanish. Hence the fields in \eqref{cdp:eq:C.CF.7} commute.

\item For completeness, semicompleteness applies to the entire
space, not only to a selected nonzero element. Given
$V\in\mathfrak a_E$, choose its representative $v\in L$.
Its poles are confined to the central fibre. On every
punctured fibre it is a holomorphic translation field
multiplied by constants, and is complete. Remove the
ambient polar divisor. The closedness of the
uniformizability locus for a holomorphic vector field
\cite[Theorem~3.3]{Guillot2024} makes $v$ semicomplete
there, since the punctured tube is dense. The component
$E$ is invariant: locally $v(s)=0$ and factorization of
the local equation of $s$ show that each irreducible
component ideal is invariant wherever $v$ is holomorphic.
Solutions in $E$ inherit single-valuedness. Some removed
curves may be removable poles for $V$; the same closedness
theorem on $E$ minus its actual polar divisor extends
semicompleteness across those curves. This proves the
assertion for every constant linear combination in
$\mathfrak a_E$.

\item Take a complex basis $V_1,V_2$ of $\mathfrak a_E$.
When $V_1\wedge V_2$ is nonzero, it is the announced
meromorphic frame. Otherwise there is a globally defined
meromorphic function $h$ with
\[
                  V_2=hV_1,
                  \qquad V_1(h)=0.
 \tag{C.CF.8}\hypertarget{tag-00BU}{}\label{cdp:eq:C.CF.8}
\]
The local coefficient ratios agree on their overlaps,
which proves that $h$ is global. It is nonconstant by
the complex-linear independence in \eqref{cdp:eq:C.CF.7}. The second
identity follows from the vanishing bracket.
Resolve the graph of $h:E\dashrightarrow\mathbb P^1$
and take the Stein factor of its projection to
$\mathbb P^1$. This gives exactly \eqref{cdp:eq:C.CF.2}, with a
finite map $B\to\mathbb P^1$. The lifted field is
tangent to $q$ and nonzero on a general fibre.
Its semicompleteness passes through the birational
open isomorphism and then over its full holomorphic
locus by the same closedness theorem.

\item A nonzero semicomplete meromorphic field on a smooth
compact curve has no pole. At a pole of order $r\ge1$
its expression is $z^{-r}a(z)\partial_z$, $a(0)\ne0$.
The time form $z^r a(z)^{-1}dz$ has a primitive
vanishing to order $r+1$; its inverse solution has
distinct analytic continuations around that time
value. They all lie in the punctured coordinate
disc, contradicting single-valuedness. The restricted
field is therefore holomorphic, so
$\deg T_G=2-2g(G)\ge0$ for a general fibre $G$.
This proves the rational-or-elliptic assertion.

\item Finally suppose $a(S_i^\nu)=0$. The map
$E\to S_i^\nu$ has a finite Stein factor
$E\to T'\to S_i^\nu$, whose first map is
bimeromorphic. For a meromorphic function on $T'$, the
elementary symmetric functions of its values over
the finite map are meromorphic on $S_i^\nu$.
They give a monic algebraic equation over
$\mathcal M(S_i^\nu)$. Thus the meromorphic function
field extension is algebraic and
\[
                            a(E)=a(S_i^\nu)=0.
 \tag{C.CF.9}\hypertarget{tag-00BV}{}\label{cdp:eq:C.CF.9}
\]
In particular \eqref{cdp:eq:C.CF.8} cannot occur.

\item Apply the compact-surface classification of semicomplete
meromorphic fields \cite[Theorem~9.2]{Guillot2024}
separately to $V_1$ and $V_2$. Its fibration alternative
would give a nonconstant meromorphic function on $E$;
therefore each field is holomorphic on some smooth
bimeromorphic model. An algebraic-dimension-zero
surface is not ruled. Its smooth minimal model $E_0$
is unique, and every smooth bimeromorphic model maps
to it by successive point blowdowns. Holomorphic
fields descend under these actual maps: for a
blowdown $\pi$ the equality
$\pi_*\mathcal O=\mathcal O$ sends the derivation
$g\mapsto V(\pi^*g)$ to a holomorphic derivation on
the base. On the common open isomorphism it is the
original field. The two resulting holomorphic fields
on the same $E_0$ remain complex-linearly independent
and commute.

\item The complete list of minimal compact surfaces
admitting a nonzero holomorphic field is
\cite[Theorem~0.3]{DOT2001}. In algebraic dimension
zero, its possibilities are two-tori, Hopf surfaces,
Inoue surfaces, and minimal class-VII surfaces with
positive second Betti number. The last class has
at most one independent holomorphic field
\cite[pp.~256 and~260]{DOT2000}; Inoue surfaces
also have at most one
\cite[Propositions~2(ii), 3(ii) and~5]{Inoue1974}.
The two independent fields on $E_0$ exclude both
of those alternatives. Hence $E_0$ is a two-torus
or a Hopf surface, as claimed.

\end{enumerate}
\end{proof}






\begin{proposition}[Canonical support forces the numerically anticanonical class VII case]\hypertarget{tag-01CU}{}\label{cdp:st:C.2.UN.17}
\StatementRecord{cdp:st:C.2.UN.17}{proved}{01CU}{\SeriesRef{cdp:st:A.5.3}, \SeriesRef{cdp:eq:C.CA.4}, \SeriesRef{cdp:eq:C.CA.3}, \SeriesRef{cdp:eq:C.CA.5}, \SeriesRef{cdp:eq:C.CA.7}, \SeriesRef{cdp:eq:C.CA.9}, \SeriesRef{cdp:eq:C.CA.6}; \cite[Remark~2.7]{DlousskyTeleman2012}; \cite[\S1, pp.~4--5]{Dloussky2006}; \cite[Theorem~0.2]{Dloussky2006}}

Let $f:U\to\Delta$ be a proper holomorphic map from a smooth
complex threefold, with connected fibres and smooth two-torus
fibres over $\Delta^*$. Write its full central fibre as
$F=\sum_i m_iS_i$. For each component let
$h_i:Z_i\to S_i$ be its normalization followed by minimal
resolution, and let $D_i\ge0$ be the intrinsic resolution
conductor of \SeriesRef{cdp:st:A.5.3}. There are nonnegative integers $a_j$ with
$K_U\simeq\OO_U(\sum_j a_jS_j)$. Put
$E_{ij}=h_i^*(S_j)$ for $j\ne i$, using the effective
Cartier pullback. The exact line-bundle identity is
\[
 K_{Z_i}^{\otimes m_i}\otimes\OO_{Z_i}(m_iD_i)
 \simeq
 \OO_{Z_i}\!\left(
  \sum_{j\ne i}\bigl[m_i a_j-m_j(a_i+1)\bigr]E_{ij}
 \right).                                                \tag{C.CA.1}\hypertarget{tag-00BC}{}\label{cdp:eq:C.CA.1}
\]
In particular a positive tensor power of the canonical line on the minimal
model of $Z_i$ has a meromorphic section whose divisor is
supported on the images of these interface and conductor
curves.

When this minimal model $M$ is of class VII, it has the
following restrictions. If $b_2(M)=0$, it is a Hopf surface.
If $b_2(M)>0$, there are an integer $n>0$, a nonzero
effective integral divisor $A$ and a complex-flat line
bundle $L$ such that
\[
                 K_M^{\otimes n}\otimes\OO_M(A)
                         \simeq L.                       \tag{C.CA.2}\hypertarget{tag-00BF}{}\label{cdp:eq:C.CA.2}
\]
Consequently $M$ contains a global spherical shell by the
proved numerically anticanonical theorem. In particular,
the unrestricted normalization problem does not require
the classification of class VII surfaces without a
numerically anticanonical divisor. These are restrictions
on the components of the actual map $f$; no occurrence of
every surface with \eqref{cdp:eq:C.CA.2} is asserted.
\end{proposition}


\begin{proof}

\begin{enumerate}[label={$\langle1\rangle$\arabic*.},leftmargin=*]
\item After shrinking the disc, $f_*K_U$ is a locally free
rank-one sheaf. Coherence follows from properness, and
torsion-freeness follows because a holomorphic canonical
section killed by a power of the base coordinate vanishes
on a dense open subset. Its rank is one since every nearby
torus has a one-dimensional space of canonical forms.
Choose a frame and evaluate it on $U$. Over $\Delta^*$,
Grauert base change identifies its restriction with a
nonzero torus canonical form, which is nowhere zero.
The evaluation is therefore a nonzero holomorphic section
of $K_U$ with divisor $\sum_j a_jS_j$, where $a_j\ge0$.
This constructs the claimed line-bundle isomorphism.

\item Let $N_i=h_i^*\OO_U(S_i)$. The canonical sections defining
$S_j$ do not vanish identically on $S_i$ for $j\ne i$.
Their pullbacks define the actual effective divisors
$E_{ij}$, including any exceptional orders. Pullback of
the principal full fibre gives
\[
 N_i^{\otimes m_i}\otimes
 \OO_{Z_i}\!\left(\sum_{j\ne i}m_jE_{ij}\right)
                         \simeq\OO_{Z_i}.               \tag{C.CA.3}\hypertarget{tag-00BG}{}\label{cdp:eq:C.CA.3}
\]
Adjunction for the individual reduced hypersurface $S_i$,
followed by the conductor comparison of \SeriesRef{cdp:st:A.5.3}, gives
\[
 \begin{split}
 K_{Z_i}\otimes\OO_{Z_i}(D_i)
   &\simeq h_i^*\bigl((K_U\otimes\OO_U(S_i))|_{S_i}\bigr)\\
   &\simeq N_i^{\otimes(a_i+1)}\otimes
      \OO_{Z_i}\!\left(\sum_{j\ne i}a_jE_{ij}\right).
 \end{split}                                             \tag{C.CA.4}\hypertarget{tag-00BH}{}\label{cdp:eq:C.CA.4}
\]
Raise \eqref{cdp:eq:C.CA.4} to the $m_i$-th power and substitute
\eqref{cdp:eq:C.CA.3}. This proves \eqref{cdp:eq:C.CA.1}, with its signs and every
full-fibre multiplicity retained.

\item Let $\tau:Z_i\to M$ be the sequence of point blowdowns
to a minimal model. Identity \eqref{cdp:eq:C.CA.1} provides a
meromorphic section of $K_{Z_i}^{\otimes m_i}$ with the
displayed signed divisor, less $m_iD_i$. On the open set
where $\tau$ is an isomorphism this is a meromorphic
$m_i$-canonical section on $M$. In local coordinates at
each omitted point, the coefficient of that section is a
meromorphic function on a punctured smooth surface and
extends meromorphically. Thus the section extends across
those points. Its divisor is precisely the pushforward:
exceptional curves disappear, and every other curve
keeps its order. Equivalently, at an individual blowup
the exact canonical formula
$K_{\widehat M}=\tau^*K_M+E$ removes the corresponding
$m_iE$ from the pulled-back pluricanonical divisor.
We have therefore constructed a signed integral divisor
$B$ on $M$ and an actual isomorphism
\[
                         K_M^{\otimes m_i}\simeq\OO_M(B).
                                                               \tag{C.CA.5}\hypertarget{tag-00BI}{}\label{cdp:eq:C.CA.5}
\]
Every component of $B$ is the image of an interface or
intrinsic conductor curve from \eqref{cdp:eq:C.CA.1}.

\item Suppose now that $M$ is of class VII. If it contained no
curve, \eqref{cdp:eq:C.CA.5} would have $B=0$ and would trivialize a
positive power of $K_M$. This contradicts
$\kappa(M)=-\infty$. When $b_2(M)=0$, choose a nonempty
curve. A positive multiple of it is homologically trivial,
since its integral class is torsion. The classification
of such effective divisors in
\cite[Remark~2.7]{DlousskyTeleman2012} gives a Hopf
surface in this zero-Betti case. This proves the first
alternative without an assertion about unknown class VII
surfaces.

\item For $b_2(M)>0$, let $C_1,\ldots,C_r$ be the distinct
curves in the support of $B$, and write
\[
 c_1(K_M)_{\mathbb R}=\sum_{\ell=1}^r c_\ell[C_\ell],
 \qquad c_\ell\in\Q.                                  \tag{C.CA.6}\hypertarget{tag-00BJ}{}\label{cdp:eq:C.CA.6}
\]
The intersection form on $H^2(M,\mathbb R)$ is negative
definite and $K_M^2=-b_2(M)$. Moreover
$K_M\cdot C_\ell\ge0$ for every curve on the minimal
surface; these standard class VII facts are stated with
their adjunction argument in \cite[\S1, pp.~4--5]{Dloussky2006}.
For completeness, adjunction gives
$K_M\cdot C=2p_a(C)-2-C^2$. A curve of arithmetic genus
zero is a smooth rational curve; minimality excludes
self-intersection $-1$, and self-intersection zero would
give zero real class and hence $K_M\cdot C=0$, contradicting
adjunction. Every other curve has arithmetic genus at
least one and $C^2\le0$, so also has nonnegative canonical
intersection.

\item Let $Q=(C_\ell\cdot C_k)$ and
$v_\ell=K_M\cdot C_\ell$. Then
\[
                         Qc=v,\qquad v\ge0.             \tag{C.CA.7}\hypertarget{tag-00BK}{}\label{cdp:eq:C.CA.7}
\]
The matrix $Q$ is negative semidefinite, and all its
off-diagonal entries are nonnegative. Separate it into
blocks according to the connected components of the
intersection graph. For a negative-definite block, write
$c=c_+-c_-$ with disjoint nonnegative coordinate supports.
If $c_+\ne0$, then
\[
 c_+^tQc=c_+^tQc_+-c_+^tQc_-<0,
 \qquad c_+^tv\ge0,                                    \tag{C.CA.8}\hypertarget{tag-00BL}{}\label{cdp:eq:C.CA.8}
\]
contradicting \eqref{cdp:eq:C.CA.7}. Thus every coefficient in this
block is nonpositive.

\item We also treat every singular block, rather than assuming
linear independence of the curve classes. Choose a nonzero
kernel vector $u$. Replacing it by its coordinatewise
absolute value cannot decrease its quadratic form, because
the off-diagonal entries are nonnegative. Both quadratic
forms must be zero by negative semidefiniteness. Hence
$|u|$ is again in the kernel. A zero coordinate of $|u|$
would, in the corresponding kernel equation, force every
neighboring coordinate to be zero. Connectedness would
then force the whole vector to vanish. Thus the block has
a strictly positive kernel vector. Pairing \eqref{cdp:eq:C.CA.7} with
this vector gives $v=0$ on the block, and then $Qc=0$.
The real divisor class $\beta=\sum c_\ell[C_\ell]$ of
that block satisfies $\beta^2=0$. The ambient negative
definite form forces $\beta=0$.

\item Discard the singular blocks and negate the coefficients
on the negative-definite blocks. This gives an effective
rational divisor $A_{\Q}$ on the actual curves of $M$
such that
\[
                   c_1(K_M+A_{\Q})_{\mathbb R}=0.        \tag{C.CA.9}\hypertarget{tag-00BM}{}\label{cdp:eq:C.CA.9}
\]
It is nonzero, since otherwise $c_1(K_M)_{\mathbb R}=0$
would contradict $K_M^2=-b_2(M)<0$.

\item Choose a positive multiple clearing the denominators in
$A_{\Q}$ and then another positive multiple killing the
torsion integral Chern class in \eqref{cdp:eq:C.CA.9}. We obtain an
integer $n>0$, a nonzero effective integral divisor $A$,
and a topologically trivial holomorphic line
\[
                       L=K_M^{\otimes n}\otimes\OO_M(A).
                                                               \tag{C.CA.10}\hypertarget{tag-00BD}{}\label{cdp:eq:C.CA.10}
\]
We specify why this line is complex-flat. On a class VII
surface $p_g=0$, $c_2(M)=b_2(M)$ and $K_M^2=-b_2(M)$,
so Noether's formula gives $\chi(\OO_M)=0$ and
$h^1(\OO_M)=1$. Also $\dim_{\C}H^1(M,\C)=1$.
The natural map $H^1(M,\R)\to H^1(M,\OO_M)$ is
injective: if a real closed one-form $\alpha$ has
$\alpha^{0,1}=\bar\partial(p+\mathrm i q)$, with $p,q$
real, then
$\alpha=dp+\mathrm i(\bar\partial q-\partial q)$ and
closedness gives $\partial\bar\partial q=0$. The
pluriharmonic maximum principle on the compact connected
surface makes $q$ constant, so $\alpha$ is exact.
Consequently the complex-linear map
$H^1(M,\C)\to H^1(M,\OO_M)$ is an isomorphism.
The exponential sequences for constant and holomorphic
coefficients therefore give the surjection
\[
             H^1(M,\C^*)\longrightarrow\Pic^0(M).       \tag{C.CA.11}\hypertarget{tag-00BE}{}\label{cdp:eq:C.CA.11}
\]
This is the actual flat-character map, taking constant
transition functions to the same holomorphic transition
functions. It supplies a complex-flat realization of $L$.

\item The canonical section of $\OO_M(A)$ now becomes a
nonzero section of $K_M^{-n}\otimes L$. These are exactly
the hypotheses of \cite[Theorem~0.2]{Dloussky2006},
which proves existence of a global spherical shell.
This proves \eqref{cdp:eq:C.CA.2} and the stated normalization
restriction. The argument also shows directly why a
positive-Betti class VII surface whose every curve class
is zero in real cohomology cannot be a minimal
normalization model here: \eqref{cdp:eq:C.CA.6} would contradict
$K_M^2=-b_2(M)$.

\end{enumerate}
\end{proof}



\ref{cdp:st:C.2.UN.19} subsequently excludes the positive-$b_2$ alternative
altogether by the actual cyclic base-change quotient and its
descent to a Hopf surface. The canonical-support and
numerically anticanonical comparison above is retained as an
independent necessary reduction. It does not supply an example
of a positive-$b_2$ class VII component.




\begin{proposition}[Every class VII normalization component has Hopf minimal model]\hypertarget{tag-01CW}{}\label{cdp:st:C.2.UN.19}
\StatementRecord{cdp:st:C.2.UN.19}{proved}{01CW}{\SeriesRef{cdp:eq:C.HQ.2}, \SeriesRef{cdp:st:C.2.UN.17}, \SeriesRef{cdp:st:C.2.UN.16}, \SeriesRef{cdp:eq:C.HQ.3}, \SeriesRef{cdp:eq:C.HQ.4}, \SeriesRef{cdp:eq:C.HQ.6}, \SeriesRef{cdp:eq:C.HQ.7}, \SeriesRef{cdp:eq:C.HQ.8}, \SeriesRef{cdp:eq:C.HQ.10}; \cite[\S1, universal and canonical desingularization]{BierstoneMilman1997}; \cite[p.~256]{DOT2000}; \cite[\S1.1]{Nakamura2008VII}; \cite[Theorem~2.1]{Khaled2024Hopf}; \cite[Remark~2.7]{DlousskyTeleman2012}}

For every proper holomorphic degeneration of smooth complex
two-tori with smooth three-dimensional total space, each
central component whose normalization resolution is of class
VII has a Hopf surface as its minimal model. This holds for
arbitrary multiplicities, arbitrary central singularities,
and arbitrary Euler number. In particular no Inoue surface
and no minimal class VII surface with positive second Betti
number is a normalization minimal model of such a component.

The comparison proving this assertion is the cyclic cover
obtained from $t=s^{m_i}$ for that component. Its stabilizer
acts on the minimal model of the covering component, and
the quotient is bimeromorphic to the original normalization.
\end{proposition}



\begin{proof}

\begin{enumerate}[label={$\langle1\rangle$\arabic*.},leftmargin=*]
\item We first identify the quotient in the actual base change.
Normalize $U\times_{\Delta_t}\Delta_s$, where
$t=s^{m_i}$. Its cyclic group $\mu_{m_i}$ acts over $U$.
Choose a resolution of the total space and central divisor
equivariant for this finite group. Canonical analytic
resolution has this property: the invariants defining its
centres are preserved by every local isomorphism, so the
action lifts through every blowup
\cite[\S1, universal and canonical desingularization]{BierstoneMilman1997}.
The algorithm is an isomorphism over the punctured tube
and at the generic smooth points used in (C.SC.2).

\item Let $E$ be a chosen smooth strict component above $S_i$,
and let $G\subset\mu_{m_i}$ be its stabilizer. At a
generic point of $S_i$ the equations are exactly
\[
                 w^{m_i}=u,\qquad z=s/w.
 \tag{C.HQ.1}\hypertarget{tag-00MQ}{}\label{cdp:eq:C.HQ.1}
\]
Consequently the reduced inverse image over a dense
smooth open subset of $S_i$ is a $\mu_{m_i}$-torsor.
The cyclic group permutes its connected components
transitively. The stabilizer of a selected component
acts freely and transitively on each fibre of that
component. Equivalently, the degree of $E\to S_i$
at the generic point is $|G|$, and quotienting by
$G$ gives degree one. The proper map to normalization
therefore descends to a proper bimeromorphic map
\[
                         E/G\longrightarrow S_i^\nu.
 \tag{C.HQ.2}\hypertarget{tag-00MS}{}\label{cdp:eq:C.HQ.2}
\]
This argument remains valid when the root cover is
disconnected over $S_i$; the stabilizer, rather than
the whole base-change group, is the group in \eqref{cdp:eq:C.HQ.2}.

\item Write $Z_i$ for a smooth normalization resolution and
$M_i$ for its minimal model. The case $b_2(M_i)=0$
already has Hopf minimal model by \ref{cdp:st:C.2.UN.17}. Every
minimal class VII surface with $b_2>0$ has algebraic
dimension zero \cite[p.~256]{DOT2000}. It therefore
suffices to treat $a(M_i)=0$, by an argument which
also covers the zero-Betti case. The field comparison
in (C.CF.9) gives
$a(E)=0$, and \ref{cdp:st:C.2.UN.16} says that the unique smooth
minimal model $E_0$ is a torus or a Hopf surface.

\item Every element of $G$ induces a birational automorphism
of $E_0$. On a nonruled minimal compact surface this
is holomorphic, by uniqueness of the minimal model.
Thus the contraction $\rho:E\to E_0$ is equivariant:
the identity $\rho g=g_0\rho$ holds on the common
open isomorphism and hence everywhere. It descends
to a proper bimeromorphic map
\[
                    E/G\longrightarrow E_0/G.
 \tag{C.HQ.3}\hypertarget{tag-00MT}{}\label{cdp:eq:C.HQ.3}
\]
The actions on $E_0$ are faithful after removing
their kernel, which does not change the quotient.
Resolving the graph of \eqref{cdp:eq:C.HQ.2}--\eqref{cdp:eq:C.HQ.3} gives
actual proper bimeromorphic comparisons between
$Z_i$ and a resolution of $E_0/G$.

\item We exclude the torus alternative by its first
cohomology. For a finite holomorphic quotient
$q:A\to A/G$, proper base change gives
$R^jq_*\mathbb Q=0$ for $j>0$, and averaging gives
$(q_*\mathbb Q)^G=\mathbb Q$. Taking invariants is
exact over $\mathbb Q$, so the actual pullback map
identifies
\[
                     H^1(A/G,\mathbb Q)
                          \simeq H^1(A,\mathbb Q)^G.
 \tag{C.HQ.4}\hypertarget{tag-00MU}{}\label{cdp:eq:C.HQ.4}
\]
For a complex torus the holomorphic action preserves
$H^{1,0}\oplus H^{0,1}$ and conjugation interchanges
the two invariant summands. Hence the dimension in
\eqref{cdp:eq:C.HQ.4} is even. Each singularity of a cyclic surface
quotient has a resolution with a connected chain
of rational curves: after linearization, the
quotient cone is the first quadrant in
$\mathbb Z^2+\frac1n(a,b)\mathbb Z$, and a regular
fan subdivision gives that resolution. Removing
pseudoreflections first only replaces the local
coordinates by their invariant powers. Thus every
exceptional fibre has $H^1=0$ and is connected.
Proper base change and the degree-one part of
Leray show that resolving $A/G$ leaves $H^1$
unchanged. Birational maps between smooth compact
surfaces preserve $H^1$, as follows from their
factorization into point blowups. The comparisons
\eqref{cdp:eq:C.HQ.2}--\eqref{cdp:eq:C.HQ.3} would therefore make $b_1(Z_i)$
even, contradicting $b_1(Z_i)=1$.

\item It remains to prove that the cyclic quotient of
the Hopf surface $H=E_0$ is again Hopf. We give
the fixed-point and quotient argument explicitly.
Let
\[
            H=(\mathbb C^2\setminus\{0\})/\Gamma.
 \tag{C.HQ.5}\hypertarget{tag-00MV}{}\label{cdp:eq:C.HQ.5}
\]
The deck group contains an infinite cyclic
subgroup of finite index generated by a
holomorphic contraction. Equivalently $H$ has
a finite unramified primary Hopf cover; the
normal forms and the finite-cover description
are recalled in \cite[\S1.1]{Nakamura2008VII}
and \cite[Theorem~2.1]{Khaled2024Hopf}.
Every deck transformation extends across the
origin by Hartogs, as does its inverse. These
extensions fix the origin and are biholomorphic.
The same is true for any lift $\widetilde g$
of a holomorphic automorphism $g$ of $H$.

\item Define the real character
\[
 \delta:\Gamma\longrightarrow\mathbb R,
 \qquad \delta(\gamma)=\log|\det(d\gamma)_0|.
 \tag{C.HQ.6}\hypertarget{tag-00MW}{}\label{cdp:eq:C.HQ.6}
\]
It is nonzero, since its value on a contraction
is negative. For any automorphism lift,
conjugacy of derivatives gives
\[
 \delta(\widetilde g\gamma\widetilde g^{-1})
                         =\delta(\gamma).
 \tag{C.HQ.7}\hypertarget{tag-00MX}{}\label{cdp:eq:C.HQ.7}
\]
A primary Hopf cover has rational Betti numbers
$(1,1,0,1,1)$: this follows directly from the
covering spectral sequence for
$\mathbb C^2\setminus\{0\}\simeq S^3$ and the
infinite cyclic deck group, whose action on
$H^3(S^3,\mathbb Q)$ preserves orientation.
Transfer from the finite primary cover gives
$b_1(H)\le1$ and $b_2(H)=0$, while \eqref{cdp:eq:C.HQ.6}
gives $b_1(H)\ge1$. Hence $b_1(H)=1$.
Equation \eqref{cdp:eq:C.HQ.7} says that every holomorphic
automorphism acts identically on $H^1(H,\mathbb R)$.
Poincar\'e duality and its complex orientation
then give the identity action on $H^3$ as well.
Its Lefschetz number is therefore
\[
                 L(g)=1-1+0-1+1=0.
 \tag{C.HQ.8}\hypertarget{tag-00MY}{}\label{cdp:eq:C.HQ.8}
\]

\item For a nonidentity finite-order automorphism,
the fixed locus is a disjoint union of smooth
compact curves and isolated points, by finite
holomorphic linearization. Each fixed curve $C$
is smooth. Since $H^2(H,\mathbb Q)=0$, one has
$C^2=K_H\cdot C=0$. The actual adjunction
isomorphism $K_C=(K_H\otimes\mathcal O_H(C))|_C$
therefore gives $2g(C)-2=0$. Thus all the
curve components of the fixed locus have Euler
number zero. For a finite-order
map the topological Lefschetz number equals
the Euler number of its fixed locus. Indeed,
use an equivariant triangulation and barycentric
subdivision; each invariant simplex is then
fixed pointwise, so the alternating trace
on simplicial cochains counts exactly the
fixed simplices. Equation \eqref{cdp:eq:C.HQ.8} consequently
shows that there are no isolated fixed points.

\item At any point of $H$, the stabilizer in the
effective cyclic group $G$ is cyclic. Its
generator has a fixed curve through that
point and, in linearizing coordinates, acts
by
\[
                (x,y)\longmapsto(x,\zeta y),
 \qquad
 \mathbb C\{x,y\}^{\langle\zeta\rangle}
                    =\mathbb C\{x,y^{|G_x|}\}.
 \tag{C.HQ.9}\hypertarget{tag-00MZ}{}\label{cdp:eq:C.HQ.9}
\]
Here $G_x$ denotes the stabilizer, and
$\zeta$ has its full order; a finite-order
germ with identity derivative is the identity.
This proves that the quotient $Q=H/G$ is
smooth, with the displayed local quotient
map even at its branch curves.

\item The invariant-cohomology comparison in
\eqref{cdp:eq:C.HQ.4}, now in every degree and combined
with \eqref{cdp:eq:C.HQ.7}, gives
\[
                         b_1(Q)=1,
                         \qquad b_2(Q)=0.
 \tag{C.HQ.10}\hypertarget{tag-00MR}{}\label{cdp:eq:C.HQ.10}
\]
The finite map $H\to Q$ injects every space
of pluricanonical sections by the actual
Jacobian pullback. Since a Hopf surface
has $\kappa=-\infty$, this proves
$\kappa(Q)=-\infty$. Also $Q$ is minimal:
an exceptional curve of self-intersection
$-1$ would have a nonzero real $H^2$ class,
contrary to \eqref{cdp:eq:C.HQ.10}.

\item There is a curve on $Q$. Explicitly, the
primary Hopf normal form
$(z_1,z_2)\mapsto(\alpha z_1+a z_2^r,\beta z_2)$
contains the elliptic curve given by $z_2=0$;
its finite image in $H$, and then in $Q$,
is a nonempty curve. A positive multiple
of that curve is homologically trivial by
\eqref{cdp:eq:C.HQ.10}. The divisor classification
\cite[Remark~2.7]{DlousskyTeleman2012} applied
to the minimal class VII surface $Q$ now
gives a Hopf surface: its Enoki alternative
has positive second Betti number and is
excluded by \eqref{cdp:eq:C.HQ.10}. We have proved the
cyclic-quotient assertion by actual quotient
maps, including their branch points.

\item Finally \eqref{cdp:eq:C.HQ.2}--\eqref{cdp:eq:C.HQ.3} make $Z_i$
bimeromorphic to this smooth Hopf surface
$Q$. Uniqueness of the class VII minimal
model proves that $M_i\simeq Q$. This
proves the assertion for every component
in the original, unaltered degeneration.

\end{enumerate}
\end{proof}






\begin{proposition}[The rank-one component fibration descends to the original support]\hypertarget{tag-01CY}{}\label{cdp:st:C.2.UN.20}
\StatementRecord{cdp:st:C.2.UN.20}{proved}{01CY}{\SeriesRef{cdp:st:C.2.UN.16}, \SeriesRef{cdp:st:C.2.UN.19}, \SeriesRef{cdp:eq:C.FD.6}}

In the setting and with the actual component cover of \ref{cdp:st:C.2.UN.16},
the generic-rank-one branch gives a fibration on a resolution of
the original normalization. More precisely, for the stabilizer
$G$ of the chosen covering component $E$, there are proper
birational maps and quotient maps in a commutative diagram
\[
\begin{array}{ccc}
 \widehat E&\xrightarrow{\ q\ }&B\\
 \downarrow&&\downarrow\\
 \widehat E/G&\longrightarrow&B/G,
\end{array}                                                    \tag{C.FD.1}\hypertarget{tag-00K2}{}\label{cdp:eq:C.FD.1}
\]
where $B$ is a smooth compact curve, $q$ has connected general
fibre, and $\widehat E/G$ is bimeromorphic to $S_i^\nu$.
Resolving this last quotient and taking its Stein factor gives
an actual proper fibration
\[
        Z'_i\longrightarrow B/G
        \quad\hbox{with connected rational or elliptic general fibre}.
                                                               \tag{C.FD.2}\hypertarget{tag-00K3}{}\label{cdp:eq:C.FD.2}
\]
The base in \eqref{cdp:eq:C.FD.2} is already the Stein base; in particular
there is no additional finite cover of $B/G$ implicit in this
description. For every original component with nonnegative
Kodaira dimension, its general fibre in \eqref{cdp:eq:C.FD.2} is elliptic.
The same conclusion holds for every original component that
is not uniruled.
\end{proposition}


\begin{proof}

\begin{enumerate}[label={$\langle1\rangle$\arabic*.},leftmargin=*]
\item Use a resolution equivariant for the finite cyclic group of the
normalized base change, as in \ref{cdp:st:C.2.UN.19}. Its stabilizer $G$
preserves $E$. Recall the actual modules
\[
 M=(f'_*\mathcal T_{U'/\Delta_s})_0,
 \qquad
 L=\{v\in M\otimes_{\mathbb C\{s\}}\mathbb C(\{s\}):
                  \operatorname{ord}_E(v)\ge0\}.
                                                               \tag{C.FD.3}\hypertarget{tag-00K4}{}\label{cdp:eq:C.FD.3}
\]
An element $g\in G$ covers $s\mapsto\zeta_gs$. Pushforward of
vector fields is therefore a semilinear automorphism of $M$
and of its meromorphic localization. Since $g(E)=E$, it
preserves the valuation defining $L$. It also preserves $sL$,
because $s$ changes by a nonzero constant. Reduction gives
a complex-linear representation on $L/sL$, and the restriction
map of (C.CF.6) is equivariant:
\[
 g_*\bigl(v|_E\bigr)=(g_*v)|_E.
                                                               \tag{C.FD.4}\hypertarget{tag-00K5}{}\label{cdp:eq:C.FD.4}
\]
Consequently the two-dimensional space $\mathfrak a_E$ and
its rank-one meromorphic tangent distribution are preserved
by $G$.

\item For a basis $V_1,V_2$ write $V_2=hV_1$. The function $h$ is
nonconstant, and both fields annihilate it. If the matrix of
$g_*$ in this basis is
\[
       g_*V_1=aV_1+bV_2,\qquad
       g_*V_2=cV_1+dV_2,
\]
then the actual coefficient comparison gives
\[
              h\circ g^{-1}=\frac{c+dh}{a+bh}.
                                                               \tag{C.FD.5}\hypertarget{tag-00K6}{}\label{cdp:eq:C.FD.5}
\]
This identity holds where the displayed denominator is nonzero
and hence as an identity of meromorphic functions. It supplies
the finite projective action on the target of
$h:E\dashrightarrow\mathbb P^1$. Resolve its graph
equivariantly, and then take the Stein factor. Uniqueness of
the Stein factor lifts the action on $\mathbb P^1$ to an action
on the smooth compact curve $B$, and gives an equivariant map
$q:\widehat E\to B$. This is the same connected-leaf
fibration as in \ref{cdp:st:C.2.UN.16}. Its general fibre is rational or
elliptic, by the semicomplete one-dimensional field calculation
there. Taking the finite quotient produces \eqref{cdp:eq:C.FD.1}.

\item We spell out the general quotient fibre. Let
$K\subset G$ be the kernel of the action on $B$. Outside
the finitely many fixed points of the nonidentity elements
of $G/K$, a point $b\in B$ has stabilizer exactly $K$.
The inverse image of its orbit in $B/G$ is the disjoint union
of the fibres $q^{-1}(gb)$, and $G$ acts transitively on these
fibres. Hence the quotient fibre is exactly
\[
                       q^{-1}(b)/K.
                                                               \tag{C.FD.6}\hypertarget{tag-00K7}{}\label{cdp:eq:C.FD.6}
\]
It is connected. A finite quotient of a smooth compact curve
is smooth, as follows locally from the invariant coordinate
$z\mapsto z^r$ at a cyclic stabilizer. Riemann--Hurwitz
shows that the genus of \eqref{cdp:eq:C.FD.6} is at most the genus of
$q^{-1}(b)$, and is therefore zero or one.

\item For clarity, the generic fibre remains unchanged when the
surface quotient is resolved. A stabilizer of a point in a
general fibre is contained in $K$. It fixes the base coordinate
and, after finite holomorphic linearization, acts by
$(x,y)\mapsto(x,\xi y)$. Its invariant ring is
$\mathbb C\{x,y^r\}$, so the surface quotient is smooth
along a general fibre. A resolution can be an isomorphism
there. The resulting map to $B/G$ has connected general
fibre by \eqref{cdp:eq:C.FD.6}. Its Stein factor is a finite map of normal
compact curves of degree one, hence an isomorphism. This
proves the assertion about the base in \eqref{cdp:eq:C.FD.2}.

\item The quotient comparisons (C.HQ.2)--(C.HQ.3), or directly
the generic $G$-torsor in (C.HQ.1), identify
$\widehat E/G$ bimeromorphically with $S_i^\nu$.
Resolving their graph gives the proper birational comparison
of \eqref{cdp:eq:C.FD.2} with a smooth normalization resolution of the
original component. Thus the fibration has been constructed
on the original birational component, through the specified
finite quotient, rather than inferred from images of a
covering family of curves.

\item Finally a rational general fibre $F$ of \eqref{cdp:eq:C.FD.2} has trivial
normal line in $Z'_i$, and adjunction gives
$K_{Z'_i}|_F=K_F=\mathcal O_{\mathbb P^1}(-2)$.
Every positive pluricanonical section restricts to zero on
every general fibre and therefore vanishes identically.
This makes $\kappa(Z'_i)=-\infty$, so the nonnegative
Kodaira-dimension cases have elliptic general fibre.
More generally the rational fibration itself is a family
of rational curves covering the surface; it is precisely
an actual ruling on a proper birational model and is
excluded when the original component is not uniruled.

\end{enumerate}
\end{proof}






\begin{proposition}[A component cover with a rank-two translation frame has coherent Euler characteristic zero]\hypertarget{tag-01D0}{}\label{cdp:st:C.2.UN.22}
\StatementRecord{cdp:st:C.2.UN.22}{proved}{01D0}{\SeriesRef{cdp:st:C.2.UN.16}; \cite[Theorem~3.3]{Guillot2024}; \cite[Theorem~6.2, pp.~564--566; Section~8(iv)(11), p.~585]{Kodaira1963II}; \citet[Section~2.3, pp.~5--6]{UlmerUrzua2021}}

Let $X$ be a smooth compact complex surface with
$\kappa(X)\ge0$. Suppose that it carries a two-dimensional
complex vector space $\mathfrak a$ of commuting meromorphic
vector fields, of generic rank two, every member of which
is semicomplete on its holomorphic locus. Then
\[
                         \chi(\mathcal O_X)=0.
                                                               \tag{C.RF.5}\hypertarget{tag-00YY}{}\label{cdp:eq:C.RF.5}
\]
More precisely, $X$ has a two-torus birational model,
or it has an actual relatively minimal elliptic fibration
on a proper birational model, over a smooth compact curve
of genus at most one, with identity integral $H_1$ monodromy
and with every fibre a positive multiple of a smooth elliptic
curve. Its smooth fibres are torsors under one fixed elliptic
curve. All actual multiple fibres and their multiplicities
are retained in this conclusion.

This applies in particular to every covering component
$E$ of \ref{cdp:st:C.2.UN.16} with $\kappa(E)\ge0$ in its
generic-rank-two branch. Such a covering component cannot
have a K3 or Enriques minimal model. The assertion concerns
the actual specified root-cover component; quotients of that
component are compared separately.
\end{proposition}



\begin{proof}[The elliptic-first-integral step in \ref{cdp:st:C.2.UN.22}]

\begin{enumerate}[label={$\langle1\rangle$\arabic*.},leftmargin=*]
\item Let $X$ be a smooth compact complex surface and let $V,Y$ be
commuting meromorphic vector fields with $V\wedge Y\ne0$.
Assume that both fields are semicomplete on their holomorphic
loci and that an actual resolution of a meromorphic first
integral of $V$, followed by its Stein factorization, gives
a map
\[
 q:\widehat X\longrightarrow B
 \tag{C.RF.1}\hypertarget{tag-00YU}{}\label{cdp:eq:C.RF.1}
\]
with connected elliptic general fibre.  The proof below does
not require $X$ to be K\"ahler.  We use the meromorphic lifts
of the two fields on $\widehat X$; their brackets still vanish
on the common open isomorphism and hence everywhere, and
their semicompleteness extends over their holomorphic loci
by \cite[Theorem~3.3]{Guillot2024}.

\item First, $V$ has no horizontal polar curve.  On a general
smooth fibre $F$ of $q$, its restriction is a nonzero
semicomplete meromorphic field on the compact elliptic
curve.  The time-form argument in \ref{cdp:st:C.2.UN.16} excludes poles
on $F$.  Thus this restriction is holomorphic and nowhere
zero, and its flow acts by all translations of $F$.
The commutation relation implies
$(\Phi_V^a)_*Y=Y$ as an identity of meromorphic fields.
Consequently the intersection of the polar divisor of $Y$
with $F$ is invariant under all translations of $F$.
It is a finite set, so it is empty.  This proves that
$Y$ also has no horizontal polar curve.

\item The field $Y$ projects to a nonzero meromorphic field $Z$
on $B$.  In fact, for each local meromorphic function $g$
on $B$,
\[
 V\bigl(Y(q^*g)\bigr)
   =Y\bigl(V(q^*g)\bigr)=0.
\]
The meromorphic function $Y(q^*g)$ is therefore constant
on the connected general fibres; proper Stein descent
makes it the pullback of a meromorphic function on $B$.
These functions define $Z$.  Were $Z$ zero, both fields
would be tangent to the one-dimensional fibres, contrary
to $V\wedge Y\ne0$.

\item The projected field has no pole.  To verify this without
an unproved descent assertion about semicompleteness,
suppose that $Z$ has a pole of order $r\ge1$ at $b$.
The primitive of its time form has a zero of order $r+1$
there.  Its inverse solution consequently has distinct
continuations after a loop about the corresponding time
value.  Choose that loop and its lifted base paths in a
small punctured disc about $b$ avoiding all other bad
values.  The inverse image of the compact union of these
paths is compact.  Both $q$ and $Y$ are holomorphic there,
and $Y$ has no pole there, since all its polar curves are
vertical and the disc was chosen accordingly.  The usual
continuation theorem for holomorphic ordinary differential
equations therefore lifts the whole time loop to a
solution of $Y$.  Its endpoint has a different image in
$B$.  This contradicts semicompleteness of $Y$.
Hence $Z$ is a nonzero holomorphic field on the compact
curve $B$, and $g(B)\le1$.

\item Remove finitely many values so that $q$ is a smooth
elliptic fibration, $V,Y$ are holomorphic, $V$ is nowhere
zero vertically, and $Z$ is nonzero.  On this open base
$B^\circ$, let $\alpha_b$ be the time form of $V|_{F_b}$.
The flow of $Y$, defined for a uniform small time along
each compact fibre, carries whole fibres biholomorphically
to whole fibres and carries $V$ to itself.  It therefore
preserves their time forms.  Integration gives the exact
horizontal map from the integral homology local system
\[
 \mathcal H_1(q)\big|_{B^\circ}
   \longrightarrow\underline{\mathbb C},
 \qquad
 \gamma\longmapsto\int_\gamma\alpha_b.
 \tag{C.RF.2}\hypertarget{tag-00YV}{}\label{cdp:eq:C.RF.2}
\]
The map is injective on each stalk: the periods of a
nonzero holomorphic one-form on an elliptic curve are a
rank-two lattice in $\mathbb C$.  The preceding flow
comparison proves that its values are locally constant.
Analytic continuation of any integral cycle around a
loop preserves its complex period; injectivity then
proves that every monodromy operator on $H_1(F_b,\mathbb Z)$
is the identity.  The image lattice is fixed, say
$\Lambda\subset\mathbb C$, and the fibres over
$B^\circ$ are torsors under the fixed elliptic curve
$\mathbb C/\Lambda$.

\item Contract the vertical exceptional curves of the first
kind to obtain the actual relatively minimal elliptic
map
\[
 \widehat X\xrightarrow{\ \mu\ }X_0
       \xrightarrow{\ q_0\ }B.
 \tag{C.RF.3}\hypertarget{tag-00YW}{}\label{cdp:eq:C.RF.3}
\]
These contractions preserve the preceding integral
monodromy representation.  Kodaira's classification
\cite[Theorem~6.2, pp.~564--566; Section~8(iv)(11), p.~585]{Kodaira1963II}, together with the explicit positive-loop matrices in \citet[Section~2.3, pp.~5--6]{UlmerUrzua2021},
now shows that every fibre has smooth elliptic support.
Here is the comparison for multiplicities, to specify
the use of the table.  The only multiple types in that
classification are $mI_b$.  Locally write such a full
fibre as $t=u h^m$, where $h=0$ is its reduced Cartier
support and $u$ is a unit.  In the normalization of the
actual base change $t=s^m$, the charts
$w^m=u$, $s=wh$ are \emph{\'etale} over the original
smooth surface.  The covering family has a reduced
central fibre and its linear monodromy is the $m$th
power of the original one, hence the identity.  It
remains relatively minimal: every central component
downstairs has canonical degree zero by the displayed
Kodaira types; the \emph{\'etale} canonical pullback
gives the same zero degree on every central prime
curve upstairs, and adjunction excludes an exceptional
rational curve of self-intersection $-1$ upstairs.
Among the reduced relatively minimal types in the
cited table, identity monodromy is the smooth type
$I_0$.  Hence the original support is smooth as well.
The argument applies to every connected covering
component and makes no connectedness assumption on
the root cover.

\item In particular, for pairwise distinct base values the
fibre supports $C_b$ are disjoint smooth elliptic
curves, and $C_b^2=0$, since some positive multiple
of $C_b$ is a full fibre.  The Euler characteristic
is zero by additivity over the smooth base stratum
and the finitely many special fibres.  To compute
the canonical square, choose a nonzero meromorphic
one-form $\eta$ on $B$.  The meromorphic field $V$
descends through $\mu$, and at the generic point
there is a unique meromorphic two-form $\Omega$
satisfying
\[
 \iota_V\Omega=q_0^*\eta.
 \tag{C.RF.4}\hypertarget{tag-00YX}{}\label{cdp:eq:C.RF.4}
\]
Both sides are interpreted as actual meromorphic
differentials.  Since $V$ is tangent to $q_0$, the
right-hand side annihilates $V$, which proves
existence and uniqueness of $\Omega$ by a local
tangent-frame calculation.  Its divisor is
vertical: on a general compact fibre $V$ is
holomorphic and nowhere zero, and $\eta$ is
holomorphic and nonzero at its base value.
Thus a canonical divisor is a finite integral
combination of the disjoint square-zero curves
$C_b$, and $K_{X_0}^2=0$.  The surface
Riemann--Roch formula gives
$\chi(\mathcal O_{X_0})=(K_{X_0}^2+e(X_0))/12=0$.
Proper birational invariance of the coherent
Euler characteristic under the displayed point
blowups and blowdowns yields
$\chi(\mathcal O_X)=0$.

\item No vanishing of $Z$ at a multiple value is asserted.
Indeed, a logarithmic elliptic translation can absorb
the cyclic root monodromy.  Multiple smooth fibres
are included in the conclusion.

\end{enumerate}
\end{proof}

\begin{proof}[Completion of \ref{cdp:st:C.2.UN.22}]

\begin{enumerate}[label={$\langle1\rangle$\arabic*.},leftmargin=*]
\item We first treat the nonprojective cases by actual
meromorphic maps. If $a(X)=1$, choose a nonconstant
$h\in\mathcal M(X)$, resolve its graph, and take the
Stein factor $q:\widehat X\to B$. Every
$g\in\mathcal M(X)$ is algebraic over $\mathbb C(h)$.
On a connected general fibre of $q$ it therefore takes
finitely many values and is constant. Proper meromorphic
descent consequently gives the exact field identity
\[
                       \mathcal M(X)=q^*\mathcal M(B).
                                                               \tag{C.RF.9}\hypertarget{tag-00Z2}{}\label{cdp:eq:C.RF.9}
\]
For example, this descent can be seen by mapping the
proper graph of $g$ to $B\times\mathbb P^1$; its image
is the graph of the descended meromorphic function.
Each field in $\mathfrak a$ differentiates this field
of functions and therefore induces a meromorphic vector
field on $B$. The induced fields commute. They cannot
both vanish, since that would make both original fields
tangent to the one-dimensional fibres and violate
generic rank two. On a curve any two commuting fields
with one nonzero are constant multiples, by \eqref{cdp:eq:C.RF.8}.
A nonzero constant linear combination in $\mathfrak a$
is thus vertical. On the compact general fibre its
semicompleteness excludes poles, so the genus is at most
one; genus zero again contradicts $\kappa(X)\ge0$ by
adjunction. This gives an elliptic fibration and reduces
the case $a(X)=1$ to \eqref{cdp:eq:C.RF.1}--\eqref{cdp:eq:C.RF.4}.

\item If $a(X)=0$, apply the final argument of \ref{cdp:st:C.2.UN.16}.
That argument uses only the given two-dimensional
space of semicomplete fields: the fibration alternative
of \cite[Theorem~9.2]{Guillot2024} is impossible in
algebraic dimension zero, and the fields on their
respective holomorphic birational models descend to
two independent holomorphic fields on the same unique
nonruled minimal model. The compact-surface field
classification and dimension bounds used there give
a two-torus or a Hopf surface. Kodaira dimension
excludes the latter. This gives the two-torus alternative
and \eqref{cdp:eq:C.RF.5}.

\item It remains to treat $a(X)=2$. A compact complex surface
of algebraic dimension two has a smooth projective
birational model. Passing through the common resolution
preserves the bracket, generic rank, and coherent Euler
characteristic; semicompleteness on the full holomorphic
locus follows from the same closedness comparison as in
\ref{cdp:st:C.2.UN.16}. We may therefore work on that projective model.
Choose a basis $V_1,V_2$ of $\mathfrak a$. The nonzero
meromorphic field $V_1$ has either a leaf with Zariski-dense
image or only algebraic leaves. In the latter case the
constructible integral-curve incidence argument in
\ref{cdp:st:C.2.UN.15} gives an actual resolution of its first integral
and a Stein fibration. The semicomplete time-form calculation
on its compact general fibre gives genus at most one.
A rational general fibre would force every pluricanonical
section to vanish, by adjunction as in \ref{cdp:st:C.2.UN.20}, contrary
to $\kappa(X)\ge0$. Thus its general fibre is elliptic,
and the preceding proof applies to $V_1,V_2$.

\item It remains to treat a Zariski-dense leaf of $V_1$.
Semicompleteness supplies the univalent maximal solution
required by \cite[Main Theorem and the clarification immediately
following it, p.~2]{Guillot2019}. On an actual smooth
birational model, either the surface is Abelian, when
\eqref{cdp:eq:C.RF.5} follows at once, or there is a fibration
\[
                         q:X'\longrightarrow B
                                                               \tag{C.RF.6}\hypertarget{tag-00YZ}{}\label{cdp:eq:C.RF.6}
\]
over a rational or elliptic curve, with rational or elliptic
general fibre, which $V_1$ preserves and whose fibres
contain every polar curve of $V_1$. Again the rational-fibre
case contradicts $\kappa(X')\ge0$. Thus \eqref{cdp:eq:C.RF.6} is an
elliptic fibration. The projection $Z_1=q_*V_1$ is nonzero:
otherwise every leaf of $V_1$ would lie in a compact fibre,
contrary to the selected dense leaf. Both lifted fields
remain semicomplete on their holomorphic loci by the same
birational-open comparison and closedness argument used
in \ref{cdp:st:C.2.UN.16}.

\item We prove that the polar divisor of $V_2$ is also vertical.
At a generic point where $V_1$ is holomorphic and nonzero,
the identity $[V_1,V_2]=0$ says that its local flow preserves
$V_2$ as a meromorphic tensor and therefore preserves its
polar divisor. Equivalently, in a flow coordinate
$V_1=\partial_z$, the coefficients of $V_2$ are independent
of $z$. Thus a horizontal polar component would be a
compact invariant multisection of the foliation of $V_1$.
We now rule out such a multisection with the actual finite
maps needed for this conclusion.

\item Shrink to a nonempty open base $B^\circ$ on which the
fibration is smooth, $V_1$ is holomorphic on each whole
fibre, and $Z_1$ is nonzero. Its flow is defined uniformly
for small time on any one compact fibre, and maps that
whole fibre biholomorphically to a neighboring fibre.
Consequently the elliptic curves are isotrivial over
$B^\circ$. Suppose that $C\subset X'$ were a compact
horizontal invariant curve. Its normalization $\widetilde C$
maps finitely to $B$. In the normalized fibre product with
$\widetilde C$, the diagonal curve supplies an actual
section; it is horizontal for the lifted field. Resolve
this fibre product, taking a resolution which is an
isomorphism above the smooth open locus.

\item Fix an elliptic curve $A$ isomorphic to a smooth fibre.
Over a smaller open subset of $\widetilde C$, the
origin-preserving isomorphisms from $A$ to the fibres
form a finite \'etale torsor with structure group
$\operatorname{Aut}(A,0)$. This is a concrete finite
cover: in a Weierstrass equation with constant $j$, the
change of coordinates is obtained by adjoining the
scalars $\lambda$ with
$g_2=\lambda^4g_{2,A}$ and
$g_3=\lambda^6g_{3,A}$, omitting the equation with
zero constant coefficient when $j=0$ or $1728$.
Its ambiguity is exactly the finite origin-preserving
automorphism group. Normalize the compact curve
$\widetilde C$ in a connected component of this finite
cover. On the resulting compact curve $B'$ there is
an actual product isomorphism over a nonempty open set,
and resolving its graph gives proper maps
\[
 \begin{array}{ccc}
 \widetilde X'&\longrightarrow&X'\\
 \downarrow&&\downarrow q\\
 B'&\longrightarrow&B,
 \end{array}
 \qquad
 \widetilde X'\dashrightarrow A\times B'
 \quad\hbox{birational}.
                                                               \tag{C.RF.7}\hypertarget{tag-00Z0}{}\label{cdp:eq:C.RF.7}
\]
Here all curves and surfaces are projective, so the
isomorphism over the algebraic open set defines the
displayed rational map; its graph closure is proper.
The section is carried to the zero section.

\item On the product open set the lifted field has the form
$a(b)\partial_b+c(b)\partial_u$, where $\partial_u$
is a fixed translation field on $A$. Indeed, its vertical
coefficient is holomorphic on each whole compact elliptic
fibre, so is constant on that fibre. Tangency to the zero
section gives $c=0$. Every local leaf therefore lies in
$\{u\}\times B'$. The compact closure of this section
in the proper graph model maps to an algebraic invariant
curve in $X'$. Through a general point this contains
the entire leaf by uniqueness of analytic continuation.
This contradicts the dense leaf of $V_1$. This argument
also explains the invariant-multisection step in
\cite[\S6.3]{Guillot2019}; no extension of an arbitrary
transcendental horizontal trivialization is being used.

\item It follows that $V_2$ has no horizontal polar curve.
Hence it is holomorphic on every fibre over a suitable
nonempty open subset of $B^\circ$. Its normal component
there is a holomorphic function on a compact elliptic
curve and is constant. It consequently projects to a
meromorphic field $Z_2$ on $B$; the coefficients descend
through the proper connected-fibre map, or can be
recovered by the trace along any finite multisection.
The bracket identity gives $[Z_1,Z_2]=0$. For nonzero
$Z_1$ on a connected curve this is the precise equation
\[
              Z_1(Z_2/Z_1)=0,
              \qquad Z_2=cZ_1\quad(c\in\mathbb C).
                                                               \tag{C.RF.8}\hypertarget{tag-00Z1}{}\label{cdp:eq:C.RF.8}
\]
Thus $W=V_2-cV_1$ is vertical. It is nonzero and
$W\wedge V_1\ne0$ by the generic-rank-two hypothesis.
It is semicomplete because it is a constant member of
$\mathfrak a$, and it commutes with $V_1$. Apply the
preceding elliptic-first-integral proof to $W,V_1$
and \eqref{cdp:eq:C.RF.6}. It gives the stated relatively minimal
model, fixed period lattice, identity integral monodromy,
and \eqref{cdp:eq:C.RF.5}.

\item Finally $\chi(\mathcal O)=2$ on a K3 surface and
$\chi(\mathcal O)=1$ on an Enriques surface. These
values and coherent Euler characteristic are preserved
under the actual point blowups and blowdowns comparing
smooth birational surface models. They exclude those
two minimal models for the rank-two covering component.

\end{enumerate}
\end{proof}






\begin{proposition}[A cyclic torus quotient with nonnegative Kodaira dimension is a torus, bielliptic, or K3]\hypertarget{tag-01D1}{}\label{cdp:st:C.2.UN.23}
\StatementRecord{cdp:st:C.2.UN.23}{proved}{01D1}{\SeriesRef{cdp:st:C.2.UN.16}, \SeriesRef{cdp:eq:C.CQ.3}, \SeriesRef{cdp:eq:C.CQ.4}, \SeriesRef{cdp:eq:C.CQ.6}, \SeriesRef{cdp:eq:C.CQ.9}, \SeriesRef{cdp:eq:C.CQ.11}, \SeriesRef{cdp:st:C.2.UN.20}}

Let $A$ be a compact complex two-torus, let $G$ be a finite
cyclic group of holomorphic automorphisms acting effectively
on $A$, and let $Z\to A/G$ be a smooth resolution.
Every positive plurigenus of $Z$ is at most one. There are
the following exhaustive alternatives for a generator
$g(z)=Lz+b$:
\[
\begin{array}{c|c}
\text{linear part and action}&\text{minimal resolution type}\\ \hline
L=I&\text{a two-torus}\\
\dim_{\C}\ker(L-I)=1,\ G\text{ acts freely}
   &\text{a bielliptic surface}\\
\ker(L-I)=0,\ \det L=1&\text{a K3 surface}\\
\text{all other cases}&P_m(Z)=0\text{ for every }m>0 .
\end{array}                                                    \tag{C.CQ.1}\hypertarget{tag-00CJ}{}\label{cdp:eq:C.CQ.1}
\]
In particular an Enriques surface is not bimeromorphic to
a cyclic quotient of a complex two-torus.

For the actual component cover in \ref{cdp:st:C.2.UN.16}, suppose its
minimal model is a two-torus. The quotient by its cyclic
stabilizer, followed by the specified proper birational
comparison to the original normalization, satisfies
\eqref{cdp:eq:C.CQ.1}. Thus an original component of nonnegative
Kodaira dimension in this branch has minimal model a
torus, a bielliptic surface, or a K3 surface.
\end{proposition}


\begin{proof}

\begin{enumerate}[label={$\langle1\rangle$\arabic*.},leftmargin=*]
\item The derivative of a holomorphic torus automorphism is
constant in the translation trivialization: its entries
are holomorphic functions on a compact torus. Lifting to
the universal cover consequently writes the automorphism
as $z\mapsto Lz+b$, with $L\Lambda=\Lambda$ for the
lattice $\Lambda$. For a finite-order automorphism, $L$
is diagonalizable with roots of unity as eigenvalues.

\item We first specify the canonical comparison used in every
case. Resolve the rational lift from $A$ to $Z$:
\[
\begin{array}{ccc}
A'&\xrightarrow{\ v\ }&Z\\
\downarrow{\scriptstyle u}&&\downarrow\\
A&\longrightarrow&A/G .
\end{array}                                                    \tag{C.CQ.2}\hypertarget{tag-00CN}{}\label{cdp:eq:C.CQ.2}
\]
Here $u$ is a sequence of point blowups and $v$ is
generically finite. Pullback of a pluricanonical form
through $v$ is holomorphic. On the open set where $u$
is an isomorphism it gives a pluricanonical form on
$A$; its coefficients extend across the finitely many
missing points by Hartogs. Hence there is an injective
map
\[
 H^0(Z,mK_Z)\longrightarrow H^0(A,mK_A)
                  =\C\,\sigma^{\otimes m},
 \qquad \sigma=dz_1\wedge dz_2 .
                                                               \tag{C.CQ.3}\hypertarget{tag-00CO}{}\label{cdp:eq:C.CQ.3}
\]
Injectivity follows on the dense locus of the dominant
map. The image is $G$-invariant, so a nonzero image
also requires $(\det L)^m=1$.

\item If the quotient map $A\to A/G$ is ramified along a
curve, its generic local form is
$(x,y)\mapsto(x^e,y)$ with $e>1$. The Jacobian
pullback of a holomorphic $m$-canonical form vanishes
to order at least $m(e-1)$ along that curve. This is
incompatible with the nonzero constant form in
\eqref{cdp:eq:C.CQ.3}. We have proved
\[
 \text{a branch curve in }A\longrightarrow A/G
       \quad\Longrightarrow\quad P_m(Z)=0\quad(m>0).
                                                               \tag{C.CQ.4}\hypertarget{tag-00CP}{}\label{cdp:eq:C.CQ.4}
\]
This conclusion uses the actual ramification map,
including when the quotient is smooth along the
branch image.

\item When $L=I$, the group consists of translations. The
action is free and its quotient is the torus obtained
by enlarging $\Lambda$ by the finitely many translation
vectors. This is the first row of \eqref{cdp:eq:C.CQ.1}.

\item Suppose that $L$ has exactly one eigenvalue equal to
one. Every power of $L$ has that eigenvalue. A
nonidentity element of $G$ with a fixed point therefore
has a fixed curve, unless its derivative is the
identity; in the latter case it is a translation and
has no fixed point. Thus a nonfree action gives
\eqref{cdp:eq:C.CQ.4}. A free action gives a smooth finite
\'etale quotient. Its canonical line has finite order,
since the finite-order character $\det L$ is the
action on $\sigma$; it has no nonzero canonical
form. Averaging over the finite cover gives
\[
 h^1(Z,\OO_Z)=1,\qquad h^0(Z,K_Z)=0,\qquad
                  b_1(Z)=2 .
                                                               \tag{C.CQ.5}\hypertarget{tag-00CQ}{}\label{cdp:eq:C.CQ.5}
\]
There is also an explicit product-elliptic comparison.
Let
$A_+=\ker(L-I)^0$ and $A_-=\operatorname{im}(L-I)$.
They are elliptic subtori, their tangent spaces are
complementary, and the addition map
\[
                  A_+\times A_-\longrightarrow A
                                                               \tag{C.CQ.6}\hypertarget{tag-00CR}{}\label{cdp:eq:C.CQ.6}
\]
is a finite \'etale isogeny. The finite group of lifts
of $G$ to this product acts freely, by translations
on the first factor and affine automorphisms on the
second; its quotient is exactly $A/G$. To see the
lift explicitly, lift the translation vector $b$
through \eqref{cdp:eq:C.CQ.6} and use the linear action
$(1,L|_{A_-})$. Its power equal to the order of $g$
is a deck translation of \eqref{cdp:eq:C.CQ.6}, so the group
generated with the finite deck group is finite.
This proves the stated finite free product quotient,
the usual bielliptic presentation, without inferring
it only from the numerical invariants. In particular
this quotient is projective. This proves the second
row of \eqref{cdp:eq:C.CQ.1}.

\item It remains to treat $\ker(L-I)=0$. The endomorphism
$I-L$ of $A$ is an isogeny, because its derivative
is invertible. Solve $(I-L)p=b$ and translate the
origin to $p$. The generator becomes the linear
map $L$ and has a fixed point. Its order equals
the order of $L$.

\item Assume first that there is no branch curve. At the
fixed origin, no nontrivial power of $L$ can have
exactly one eigenvalue one; such a power would fix
a divisor. If the order is $n$, each eigenvalue
of $L$ is therefore a primitive $n$th root.
After changing the cyclic generator, its local
representation is
\[
             (z_1,z_2)\longmapsto
                     (\zeta z_1,\zeta^r z_2),
       \quad 1\le r<n,\quad \gcd(r,n)=1 .
                                                               \tag{C.CQ.7}\hypertarget{tag-00CS}{}\label{cdp:eq:C.CQ.7}
\]
We give the precise local pluricanonical test. The
quotient is the affine toric surface for the first
quadrant in
\[
 N=\Z^2+\Z(1/n,r/n),\qquad
 M=\{(a,b)\in\Z^2:a+rb\equiv0\pmod n\}.
                                                               \tag{C.CQ.8}\hypertarget{tag-00CT}{}\label{cdp:eq:C.CQ.8}
\]
These are the cocharacter and character lattices:
the invariant monomials are exactly
$z_1^az_2^b$ with $(a,b)\in M$. If
$(\det L)^m=1$, the form in \eqref{cdp:eq:C.CQ.3} descends on
the smooth locus as
\[
 (z_1z_2)^m
       \left(\frac{dz_1}{z_1}\wedge
             \frac{dz_2}{z_2}\right)^{\otimes m}.
\]
Its order on the divisor of a primitive toric ray
$v=(v_1,v_2)\in N$ is
\[
                       m(v_1+v_2-1).
                                                               \tag{C.CQ.9}\hypertarget{tag-00CU}{}\label{cdp:eq:C.CQ.9}
\]
Indeed a unimodular toric chart expresses the
monomial's order as its character pairing with $v$,
while the logarithmic volume form has one simple
pole along that divisor. The same formula holds
on a refinement, by its explicit monomial map.

\item The vector $(1/n,r/n)$ is primitive in $N$ and
can be included in a smooth fan subdivision of
the quadrant. To obtain such a subdivision, first
insert this ray and then insert primitive lattice
rays until every adjacent cone has determinant
one in a lattice basis; the Euclidean algorithm
terminates this two-dimensional procedure.
When $r<n-1$, \eqref{cdp:eq:C.CQ.9} is negative on this ray.
The descended form cannot extend to that smooth
resolution. Smooth resolutions have the same
pluricanonical sections under their actual common
point-blowup resolution, by the Hartogs comparison
already used in \eqref{cdp:eq:C.CQ.3}. Therefore
\[
 \ker(L-I)=0,\quad \det L\ne1
      \quad\Longrightarrow\quad P_m(Z)=0\quad(m>0),
                                                               \tag{C.CQ.10}\hypertarget{tag-00CK}{}\label{cdp:eq:C.CQ.10}
\]
also including the branch-curve case by \eqref{cdp:eq:C.CQ.4}.

\item Finally let $\det L=1$. Every stabilizer acts with
weights $(1,-1)$ after choosing its generator, so
there are no branch curves. For a stabilizer of
order $n$, the rays
\[
 (1,0),\ ((n-1)/n,1/n),\ldots,
                 (1/n,(n-1)/n),\ (0,1)
                                                               \tag{C.CQ.11}\hypertarget{tag-00CL}{}\label{cdp:eq:C.CQ.11}
\]
give a smooth subdivision. Consecutive determinants
are $1/n$, the covolume of $N$, and all these rays
have height $v_1+v_2=1$. Formula \eqref{cdp:eq:C.CQ.9} shows
that the invariant two-form extends without a
zero or pole across every exceptional divisor.
Thus the simultaneous resolution $Z_0$ has
$K_{Z_0}\simeq\OO_{Z_0}$. Its exceptional curves
are the chains of rational curves given by
\eqref{cdp:eq:C.CQ.11}.

\item For clarity, its simple connectedness follows
through the quotient map as well. Remove the
discrete set of points in $\C^2$ fixed by
nontrivial elements of
$\Gamma=\Lambda\rtimes\langle L\rangle$.
The complement is simply connected: a compact
loop bounds a disc in $\R^4$, and a perturbation
relative to its boundary avoids the finitely
many removed points met in a compact region.
It is the universal cover of the smooth
quotient locus. Filling in a quotient point
kills the normal closure of its finite
stabilizer in $\Gamma$, by van Kampen.
For every $\lambda\in\Lambda$, the element
$(\lambda,L)$ has a fixed point, since $I-L$
is invertible. Thus all $(\lambda,L)$ are
killed, including $(0,L)$; their products
kill every translation $(\lambda,I)$ as well.
The resulting fundamental group is trivial.
Replacing a quotient point by the rational
chain in \eqref{cdp:eq:C.CQ.11} imposes the same relation:
its plumbing neighbourhood is simply connected
and has the same punctured boundary as the
quotient ball. Hence $\pi_1(Z_0)=1$.
Together with $K_{Z_0}=\OO$, this identifies
$Z_0$ as a K3 surface. It is minimal, since
adjunction excludes a rational $(-1)$-curve
on a surface with trivial canonical line.
This proves the third row.

\item An Enriques surface has a nonzero even
pluricanonical form, but has no nonzero
canonical form and has $b_1=0$. None of
the three nonnegative cases in \eqref{cdp:eq:C.CQ.1}
has those properties: the torus and K3
have a canonical form, and the bielliptic
case has $b_1=2$. All the other cases
have no positive pluricanonical form.
This proves the exclusion.

\item We finish with the component comparison.
The stabilizer of a component in the
cyclic base change of \ref{cdp:st:C.2.UN.16} is cyclic,
and its quotient is properly bimeromorphic
to the original normalization, as in
(C.HQ.1)--(C.HQ.3). If the minimal model
of the covering component $E$ is a torus
$A$, its actual minimal contraction
$E\to A$ is equivariant. One way to
verify the descent is to factor all the
exceptional rational curves: their maps
to the torus are constant, and the
equality $\rho_*\OO_E=\OO_A$ descends
the coordinate functions and hence
every conjugate action. Thus the induced
action on $A$ is holomorphic, and
\[
 \begin{array}{ccc}
 E&\longrightarrow&A\\
 \downarrow&&\downarrow\\
 E/G&\longrightarrow&A/G
 \end{array}
 \qquad
             E/G\dashrightarrow S_i^\nu
                                                               \tag{C.CQ.12}\hypertarget{tag-00CM}{}\label{cdp:eq:C.CQ.12}
\]
has a proper birational horizontal quotient
map and the specified proper birational
comparison on the right. Resolving their
graph identifies the pluricanonical
spaces and minimal birational type.
Applying \eqref{cdp:eq:C.CQ.1} gives exactly the last
assertion. This statement concerns this
actual torus-cover branch; the separate
rank-one fibration branch is handled by
the different comparison of \ref{cdp:st:C.2.UN.20}.

\end{enumerate}
\end{proof}



For comparison with the primary literature, the projective
quotient classification is already given by
\cite[Theorem~2.1 and Lemma~3.4]{Yoshihara1995}.
Its Enriques cases have noncyclic holonomy: the Klein
four group or a dihedral group of order eight
\cite[Example~2.3 and Lemma~3.6]{Yoshihara1995}.
The proof above includes nonprojective tori and supplies
the actual cyclic component-cover comparison needed here.




\begin{proposition}[Every compact elliptic surface satisfies the intrinsic rank-one field condition]\hypertarget{tag-01D4}{}\label{cdp:st:C.2.UN.26}
\StatementRecord{cdp:st:C.2.UN.26}{proved}{01D4}{\SeriesRef{cdp:st:C.2.UN.16}, \SeriesRef{cdp:eq:C.FL.3}; \cite[Theorem~3.3]{Guillot2024}}

Let $q:X\to B$ be a proper holomorphic map from a
smooth compact connected complex surface to a smooth
compact curve, with connected elliptic general fibre.
There are commuting meromorphic vector fields $V_1,V_2$
which are linearly independent over $\mathbb C$, have
generic rank one, and whose every constant linear
combination is semicomplete on its full holomorphic
locus. Their ratio has an actual graph resolution and
Stein factor birationally equal to $q$.

If a finite cyclic group $G$ acts holomorphically on
$X$ and preserves $q$, the two-dimensional space can
be chosen $G$-stable. Thus the rank-one field-space
condition of \ref{cdp:st:C.2.UN.16}, even with this equivariance,
is satisfied by every specified compact elliptic
surface and cyclic fibration action. Any additional
restriction on its occurrence as a degeneration
component must use further data of that degeneration.
\end{proposition}



\begin{proof}

\begin{enumerate}[label={$\langle1\rangle$\arabic*.},leftmargin=*]
\item The relative canonical line is the actual invertible
sheaf
\[
 \omega_{X/B}=K_X\otimes q^*K_B^{-1},
 \qquad \mathcal E=q_*\omega_{X/B}.
                                                               \tag{C.FL.1}\hypertarget{tag-00KR}{}\label{cdp:eq:C.FL.1}
\]
Proper coherence makes $\mathcal E$ coherent. It is
torsion-free: a section of a line bundle annihilated
by a nonzero base function vanishes over the dense
complement of its zero fibre and hence vanishes
identically. On a smooth elliptic fibre the canonical
space has dimension one, so $\mathcal E$ has rank one.
It is therefore a line bundle on $B$.

\item Choose a nonzero meromorphic section $\sigma$ of
$\mathcal E$. This is an actual global choice:
Riemann--Roch on $B$ gives a nonzero holomorphic
section of $\mathcal E(D)$ for an effective divisor
$D$ of sufficiently large degree, which is the
required meromorphic section of $\mathcal E$.
Evaluation gives a nonzero meromorphic section
\[
 a:\mathcal O_X\dashrightarrow\omega_{X/B}
       \quad\hbox{from}\quad q^*\mathcal E\longrightarrow
                                      \omega_{X/B}.
                                                               \tag{C.FL.2}\hypertarget{tag-00KS}{}\label{cdp:eq:C.FL.2}
\]
Here the dashed notation records poles of the section,
not a replacement of the displayed sheaf evaluation.
On the smooth part of $q$, the canonical differential
identification $\omega_{X/B}=\Omega^1_{X/B}$ makes
$a$ an actual relative one-form.

\item Its dual defines a global meromorphic vertical field
$V$. The coordinate comparison at critical fibres
can be specified as follows. On a coordinate disc
of $B$ with coordinate $z$, set
$\Omega_z=a\otimes q^*(dz)$, a meromorphic section
of $K_X$. Define $V$ by
\[
                   \iota_V\Omega_z=-q^*(dz).
                                                               \tag{C.FL.3}\hypertarget{tag-00KT}{}\label{cdp:eq:C.FL.3}
\]
Contraction by a nonzero meromorphic two-form is an
isomorphism between the meromorphic tangent and
cotangent bundles, so this has a unique meromorphic
solution. Under $z'=\phi(z)$ both sides are multiplied
by $q^*\phi'(z)$; the solution is unchanged and
therefore glues. The alternating identity in
\eqref{cdp:eq:C.FL.3} gives $dq(V)=0$. On a regular fibre its
relative time form is $a$, so $a(V)=1$ with the
usual convention $\Omega_z=q^*(dz)\wedge a$.

\item Remove finitely many base values so that $q$ is
smooth and the evaluated section is holomorphic
and nonzero on every remaining fibre. Such a
removal is possible by generic base change for
$q_*\omega_{X/B}$ and the finite zero and pole
divisors of $\sigma$. A nonzero canonical form
on a compact elliptic curve is nowhere zero.
Thus $V$ restricts on each of these fibres to
a nonzero holomorphic translation field and is
complete there.

\item Choose any nonconstant meromorphic function
$h$ on $B$, and put
\[
              V_1=V,\qquad V_2=(q^*h)V.
                                                               \tag{C.FL.4}\hypertarget{tag-00KU}{}\label{cdp:eq:C.FL.4}
\]
Their bracket vanishes because $V(q^*h)=0$.
They are independent over $\mathbb C$ since $h$
is nonconstant, and have generic rank one by
the formula. For $c,d\in\mathbb C$, the field
$W=(c+dq^*h)V$ is complete on the dense union
of the preceding compact smooth fibres away
from the poles of $h$. Its coefficient is
constant on each such fibre; zero coefficients
give constant complete solutions. Apply
\cite[Theorem~3.3]{Guillot2024} to the complex
manifold obtained by deleting the actual polar
divisor of $W$. The univalence locus is closed
there and contains this dense union. It is
therefore the whole holomorphic locus, proving
the asserted semicompleteness for every $W$.

\item The ratio is $q^*h$, so the resulting map is
the actual composite $X\xrightarrow q B
\xrightarrow h\mathbb P^1$. The second arrow
is finite. Its Stein factor is $q$ itself:
$q_*\mathcal O_X=\mathcal O_B$, and $B$ is
the normalization in the field of functions
constant on the connected general fibres.
This proves the precise first-integral comparison.

\item For the equivariant assertion, the finite orbit
of $V$ spans a nonzero finite-dimensional
representation of $G$ by meromorphic vertical
fields. A cyclic complex representation has
an eigenvector $V'\ne0$. On a general smooth
fibre each orbit field is holomorphic and
vertical, and therefore a translation field.
The same is true of $V'$; after removing
finitely many further values it is nonzero
on every remaining fibre. Choose a nonconstant
meromorphic function $h$ on the quotient
curve $B/G$, and pull it back to $B$.
Then the space
\[
                  \mathbb C V'+\mathbb C(q^*h)V'
                                                               \tag{C.FL.5}\hypertarget{tag-00KV}{}\label{cdp:eq:C.FL.5}
\]
is $G$-stable, with the same eigencharacter
on its two displayed generators. Its members
are semicomplete by the dense compact-fibre
and closedness argument just proved, and its
ratio has the same connected-fibre Stein
factor $q$. This proves the full statement.

\end{enumerate}
\end{proof}






\begin{proposition}[The sum of the normalization geometric genera in one torus fibre is at most one]\hypertarget{tag-01D6}{}\label{cdp:st:C.2.UN.28}
\StatementRecord{cdp:st:C.2.UN.28}{proved}{01D6}{\SeriesRef{cdp:eq:C.PG.2}, \SeriesRef{cdp:eq:C.PG.6}, \SeriesRef{cdp:eq:C.PG.5}}

Let $f:U\to\Delta$ be a proper holomorphic map from a
smooth complex threefold, with connected two-torus
fibres over $\Delta^*$. Write its full central fibre
as $\sum_i m_iS_i$, and let $Z_i\to S_i^\nu$ be
smooth resolutions of the normalizations. Then
\[
                         \sum_i p_g(Z_i)\le1.
                                                               \tag{C.PG.1}\hypertarget{tag-00UZ}{}\label{cdp:eq:C.PG.1}
\]
The assertion has no K\"ahler, Betti-number, reducedness,
or normal-crossing hypothesis. In particular at most
one normalization component in a given fibre has a
nonzero holomorphic two-form, and that component has
geometric genus one.

The same bound holds for any collection of multiplicity-one
components in a smooth resolved finite base change of
this family. The proof compares their compact homology
cores with one actual nearby torus, preserving their
intersection forms.
\end{proposition}



\begin{proof}

\begin{enumerate}[label={$\langle1\rangle$\arabic*.},leftmargin=*]
\item Choose $N$ divisible by every $m_i$, make the actual
base change $t=s^N$, normalize, and resolve the total
space and its central support, isomorphically over
$s\ne0$ and at a generic point of every strict
component. This produces a proper map from a smooth
threefold
\[
                         g:\widetilde U\to\Delta_s
                                                               \tag{C.PG.2}\hypertarget{tag-00V0}{}\label{cdp:eq:C.PG.2}
\]
whose nearby fibre is exactly $f^{-1}(s^N)$.
Every strict component above an $S_i$ has multiplicity
one. To see the map explicitly at a generic point,
write $t=u x^{m_i}$ with $u$ a unit. The normalized
root charts are
\[
       w^{m_i}=u,\qquad x=s^{N/m_i}/w.
                                                               \tag{C.PG.3}\hypertarget{tag-00V1}{}\label{cdp:eq:C.PG.3}
\]
The unit-root equation is nonsingular at $s=0$;
there $s$ is a normal parameter. Subsequent
resolution is an isomorphism at this generic
point and preserves the multiplicity one.

\item Select one smooth strict component $E_i$ above each
$S_i$. Its map to $S_i^\nu$ is generically finite.
Resolve the actual rational comparison to $Z_i$:
\[
              E_i\xleftarrow{\rho_i}W_i
                            \xrightarrow{v_i}Z_i.
                                                               \tag{C.PG.4}\hypertarget{tag-00V2}{}\label{cdp:eq:C.PG.4}
\]
The first map is a sequence of point blowups,
and the second is generically finite. Pullback
of holomorphic two-forms through $v_i$ is
injective, since its Jacobian is nonzero on
a dense open set. Point-blowup invariance of
holomorphic two-forms follows by restriction
and Hartogs extension across the blown-down
points. Hence
\[
                         p_g(E_i)\ge p_g(Z_i).
                                                               \tag{C.PG.5}\hypertarget{tag-00V3}{}\label{cdp:eq:C.PG.5}
\]
It suffices to prove the bound for the selected
components $E_i$ of \eqref{cdp:eq:C.PG.2}.

\item Let $D_i$ be the reduced intersection of $E_i$
with all the other central components of the
resolved support. It is a compact normal-crossing
curve, or is empty. The restriction of $g$ is
a submersion near $E_i\setminus D_i$: in the
normal-crossing local equation only this
multiplicity-one component remains, so $g$
is a unit times its local defining coordinate.

\item We first specify the relevant homology map on
any such compact smooth surface $E=E_i$ and
boundary curve $D=D_i$. Poincar\'e--Alexander
duality, applied to the triangulable compact
curve $D$ in the oriented real four-manifold
$E$, identifies the following portion of the
relative homology sequence:
\[
 H_2(E\setminus D,\mathbb R)\longrightarrow
 H_2(E,\mathbb R)\xrightarrow{\ I_D\ }
 H_2(E,E\setminus D;\mathbb R)
       \simeq H^2(D,\mathbb R).
                                                               \tag{C.PG.6}\hypertarget{tag-00V4}{}\label{cdp:eq:C.PG.6}
\]
There is one basis vector in $H^2(D,\mathbb R)$
for each irreducible curve of $D$. In these
coordinates $I_D$ is precisely the vector of
intersection numbers with those curves: under
Poincar\'e duality the map is restriction of
$H^2(E)$ to $D$. Thus its kernel is exactly
the image from the open complement in
\eqref{cdp:eq:C.PG.6}. This statement includes all
intersections between boundary curves; no
disjointness of those curves is assumed.

\item Every holomorphic two-form $\alpha$ on $E$
is closed, because both its $(3,0)$ derivative
and its antiholomorphic derivative vanish.
The real classes of these forms give a
positive subspace
\[
 P_E=\{[\operatorname{Re}\alpha]:
                       \alpha\in H^0(E,K_E)\}
       \subset H^2(E,\mathbb R),
 \qquad \dim_{\mathbb R}P_E=2p_g(E).
                                                               \tag{C.PG.7}\hypertarget{tag-00V5}{}\label{cdp:eq:C.PG.7}
\]
Indeed, with the complex orientation,
\[
 \int_E(\operatorname{Re}\alpha)^2
       =\tfrac12\int_E\alpha\wedge\overline\alpha>0
                      \quad (\alpha\ne0).
                                                               \tag{C.PG.8}\hypertarget{tag-00V6}{}\label{cdp:eq:C.PG.8}
\]
This also proves injectivity of the real-part
map on the underlying real vector space of
holomorphic forms. Every such class has zero
integral on every boundary curve, since the
pullback of a two-form to the normalization
of a curve is zero. Its Poincar\'e-dual
homology class therefore lies in
$\ker I_D$ and lifts to
$H_2(E\setminus D,\mathbb R)$ by \eqref{cdp:eq:C.PG.6}.
No Hodge decomposition or K\"ahler form has
been used.

\item Choose a real basis of each $P_{E_i}$,
choose lifts of its Poincar\'e-dual classes,
and represent those lifts by finite compact
smooth singular cycles in $E_i\setminus D_i$.
Include their transverse perturbations needed
to compute all intersection numbers in a
compact set $K_i$ of that complement. These
are homology-cycle representatives; they
need not be embedded surfaces.

\item There is an actual orientation-preserving
smooth embedding of a neighbourhood of
$K_i$ into a nearby fibre of \eqref{cdp:eq:C.PG.2}.
Here is the local construction. On a
neighbourhood of $K_i$, $dg$ is surjective.
Choose smooth horizontal lifts of the two
coordinate fields of the real base disc;
a partition of unity preserves their
specified images under $dg$. Multiply by
a cutoff outside a slightly larger compact
neighbourhood and flow successively for
the real and imaginary parts of $s$.
For sufficiently small $s$, the trajectories
stay where the specified horizontal equations
hold. Their composition maps a neighbourhood
of $K_i$ in $E_i$ diffeomorphically to an
open subset of $g^{-1}(s)$. The image is
in that exact fibre by the two equations
for $dg$, and its orientation agrees with
the complex orientation by continuity
from $s=0$.

\item The finitely many compact sets $K_i$ lie
in different central components away
from all intersections. Choose their
ambient neighbourhoods disjoint, and
then choose the same sufficiently small
$s\ne0$ for all the flows. Their images
in the single torus $F=g^{-1}(s)$ remain
disjoint. Transverse intersection numbers
of the chosen cycles are preserved by
the local embeddings, and cycles from
different components have zero intersection.
Consequently their images give mutually
orthogonal positive subspaces of
$H_2(F,\mathbb R)$, of dimensions
$2p_g(E_i)$. Injectivity on each positive
subspace, and on their direct sum,
follows from preservation of its positive
intersection form.

\item The intersection form of an oriented
real four-torus is $3H$, where $H$ is
the hyperbolic plane. This follows
directly from the exterior-algebra
identification of its integral
cohomology: the three complementary
degree-two pairs give three hyperbolic
blocks. Its positive index is therefore
three. We have proved
\[
                    2\sum_i p_g(E_i)\le3.
                                                               \tag{C.PG.9}\hypertarget{tag-00V7}{}\label{cdp:eq:C.PG.9}
\]
The sum is integral, so it is at most
one. Together with \eqref{cdp:eq:C.PG.5} this is
\eqref{cdp:eq:C.PG.1}. The compact-core argument
used only multiplicity one, and hence
also proves the stated assertion for
any collection of such components in
a smooth finite-base-change model.

\end{enumerate}
\end{proof}






\begin{proposition}[Rank-one elliptic component covers have only smooth multiple fibres]\hypertarget{tag-01D7}{}\label{cdp:st:C.2.UN.29}
\StatementRecord{cdp:st:C.2.UN.29}{proved}{01D7}{\SeriesRef{cdp:st:C.2.UN.16}, \SeriesRef{cdp:st:C.2.UN.22}, \SeriesRef{cdp:eq:C.RG.6}, \SeriesRef{cdp:eq:C.RG.8}, \SeriesRef{cdp:eq:C.RG.5}, \SeriesRef{cdp:eq:C.RG.12}, \SeriesRef{cdp:st:C.2.UN.28}, \SeriesRef{cdp:st:C.2.UN.20}, \SeriesRef{cdp:st:C.2.UN.15}, \SeriesRef{cdp:st:C.2.UN.19}; \cite[Theorem~6.2; Section~8(iv)(11), p.~585]{Kodaira1963II}; \citet[Section~2.3, pp.~5--6]{UlmerUrzua2021}}

Let $E$ be an actual multiplicity-one covering component
constructed in \ref{cdp:st:C.2.UN.16}, with its saturated two-dimensional
space of limiting translation fields. In the generic-rank-one
branch, take the actual graph resolution and Stein factor
of their ratio. When its connected general fibre is elliptic,
the resulting relatively minimal elliptic surface
\[
                         q:E_0\longrightarrow B
                                                               \tag{C.RG.1}\hypertarget{tag-00Z3}{}\label{cdp:eq:C.RG.1}
\]
has identity integral $H_1$ monodromy. Every fibre is a
positive multiple of a smooth elliptic curve, and all smooth
fibres are torsors under one fixed elliptic curve. Moreover
\[
 R^1q_*\OO_{E_0}\simeq\OO_B,\qquad
 p_g(E)=g(B)\le1,\qquad \chi(\OO_E)=0.
                                                               \tag{C.RG.2}\hypertarget{tag-00Z8}{}\label{cdp:eq:C.RG.2}
\]
The maps to $E_0$ are the specified point blowups resolving
the first integral followed by vertical exceptional-curve
contractions. No extension of that first integral to an
ambient neighbourhood is assumed.

Together with \ref{cdp:st:C.2.UN.22}, this proves that every such root-cover
component of nonnegative Kodaira dimension has coherent
Euler characteristic zero. In particular its minimal model
is neither K3 nor Enriques. Every original component of
Kodaira dimension one has an isotrivial elliptic Iitaka
fibration over a curve of genus at most one; its possible
singular fibres are retained by the actual cyclic quotient
comparison, and are not asserted to be smooth.
\end{proposition}


\begin{proof}

\begin{enumerate}[label={$\langle1\rangle$\arabic*.},leftmargin=*]
\item We work on the actual root-change and resolution model
$f:U\to\Delta_s$ in \ref{cdp:st:C.2.UN.16}. Its smooth general fibre
is still a complex two-torus. The chosen component $E$
has multiplicity one. Its two ambient meromorphic fields
$V_1,V_2$ are vertical, commute, are holomorphic and
independent on each general torus, and have no pole along
the generic point of $E$. Their restrictions are independent
over $\C$ and have rank one over the meromorphic function
field of $E$.

\item Resolve the meromorphic ratio on $E$ by point blowups
and take its connected-fibre Stein factor $q:E\to B$.
Perform the same point blowups in the smooth ambient
threefold. This gives the actual ambient extension of
these modifications, retains multiplicity one on the
strict transform of $E$, and lifts both fields
meromorphically. Their generic restrictions and bracket
are unchanged through the common open isomorphism.
We use the same letter $E$ for this modified component.
There is a nonconstant meromorphic function $H$ on $B$
such that
\[
                    V_2|_E=(q^*H)V_1|_E.
                                                               \tag{C.RG.3}\hypertarget{tag-00Z9}{}\label{cdp:eq:C.RG.3}
\]
This is precisely the function and Stein map given by
the original ratio, rather than a chosen ambient extension
of it.

\item Remove finitely many values of $B$ so that $q$ has smooth
connected elliptic fibres and $V_1|_E$ is a nonzero
holomorphic translation field on each of them. The latter
assertion follows from the compact-curve time-form argument
in C.2.UN.16: a semicomplete meromorphic field on an
elliptic curve has no pole. After excluding the finitely
many fibres on which its restriction is identically zero,
it is nowhere zero on every remaining fibre. Remove also
the poles and critical points of $H$.

\item Let $D$ be the union on $E$ of its intersections with
the other central divisors and with the actual ambient
polar divisors of the two fields. Each remaining general
elliptic fibre meets $D$ in finitely many points, after
discarding any finitely many fibres contained in $D$.
The fields are holomorphic on $U^*$; hence all their
ambient polar divisors lie over zero. In particular
outside this explicitly removed set the ambient fields
and the submersion $f$ are holomorphic. At a suitable
point $p$ of a remaining fibre $F_{b_0}$ there are
holomorphic coordinates $(u,z,s)$ with
\[
 f=s,\qquad V_1=\partial_u,\qquad
 V_2=h(z,s)\partial_u+A(z,s)\partial_z,
 \qquad h(z,0)=H(z).
                                                               \tag{C.RG.4}\hypertarget{tag-00ZA}{}\label{cdp:eq:C.RG.4}
\]
Here $z$ is a local coordinate on $B$ along $E$.
Indeed the holomorphic flow-box theorem gives $V_1=\partial_u$,
verticality preserves $s$, and commutation makes both
coefficients of $V_2$ independent of $u$. The local coordinate
$z$ is only used in this one flow box. It need not extend
around the compact elliptic fibre.

\item Since the ambient fields are independent on the nearby
tori, $A$ is not identically zero. Since their restrictions
to $E$ have rank one, $A(z,0)=0$. Choosing $p$ away
from the zeros of its leading coefficient gives
\[
 A(z,s)=s^d\bigl(a(z)+O(s)\bigr),\qquad d\ge1,
 \qquad a(z)\ne0,\qquad H'(z)\ne0
                                                               \tag{C.RG.5}\hypertarget{tag-00ZB}{}\label{cdp:eq:C.RG.5}
\]
on a smaller disc. The order $d$ is finite; it is the
order of the ambient wedge along this component in the
chosen nonvanishing frame.

\par\medskip\noindent\textbf{Actual elliptic return maps.}
Let $\lambda_1(z),\lambda_2(z)$ be the periods of the
time form of $V_1|_{F_z}$, for an oriented integral
homology basis continued over a small disc about $b_0$.
These are holomorphic functions; a local holomorphic
section and the ordinary integration construction of
an elliptic curve supply the marked periods. Put
$c_i=\lambda_i(b_0)$.

\item The time covering $\C\to F_{b_0}$ of the translation
field pulls the finite bad set $D\cap F_{b_0}$ back to
a discrete subset of $\C$. Choose a piecewise smooth
time path from $0$ to $c_i$ avoiding this discrete set.
Its image is a compact path on which both ambient fields
are holomorphic and $V_1$ is nonzero. A finite chain of
holomorphic flow charts therefore continues the solution
along this entire path for every initial point in a
sufficiently small ambient neighbourhood of $p$.
Its endpoint is again near $p$, since the central path
closes. This constructs a genuine holomorphic return
germ $\Phi_i$; it does not deform a compact elliptic
orbit or assert that such an orbit exists on a nearby
torus.

\item Every local continuation of the flow of $V_1$ preserves
$f,V_1,V_2$, by verticality and the bracket identity.
Thus the return germ in the single flow box has the form
\[
 \Phi_i(u,z,s)=\bigl(u+\delta_i(z,s),\phi_i(z,s),s\bigr).
                                                               \tag{C.RG.6}\hypertarget{tag-00ZC}{}\label{cdp:eq:C.RG.6}
\]
On $E$ the first integral is constant along its elliptic
orbits, and comparison with the continued period gives
\[
 \phi_i(z,0)=z,\qquad
 \delta_i(z,0)=c_i-\lambda_i(z).
                                                               \tag{C.RG.7}\hypertarget{tag-00ZD}{}\label{cdp:eq:C.RG.7}
\]
To see the second equality directly, follow the fixed
time path of total time $c_i$ on $F_z$, then return
to the local section by the small flow segment of
time $-\delta_i(z,0)$. The resulting closed path is
the continued integral cycle, so its time integral is
$c_i-\delta_i(z,0)=\lambda_i(z)$. This also specifies
which integral cycle and period enter the return map.

\item Pushing $V_2$ forward by \eqref{cdp:eq:C.RG.6} and equating it to
$V_2$ at the endpoint gives the two exact identities
\[
 \begin{split}
 A(\phi_i(z,s),s)&=A(z,s)\,\partial_z\phi_i(z,s),\\
 h(\phi_i(z,s),s)&=h(z,s)+A(z,s)\,\partial_z\delta_i(z,s).
 \end{split}                                                   \tag{C.RG.8}\hypertarget{tag-00ZE}{}\label{cdp:eq:C.RG.8}
\]
These are identities of holomorphic germs on the actual
ambient threefold.

\par\medskip\noindent\textbf{The leading normal coefficients.}
The second identity of \eqref{cdp:eq:C.RG.8}, together with $H'$ being
a unit and \eqref{cdp:eq:C.RG.5}, first gives
$\phi_i-z=O(s^d)$. For instance, factor
$h(\phi_i,s)-h(z,s)=(\phi_i-z)h_1(z,s)$;
the factor $h_1$ is a unit because $h_1(z,0)=H'(z)$.
Write
\[
                \phi_i=z+s^d\psi_i(z)+O(s^{d+1}).
                                                               \tag{C.RG.9}\hypertarget{tag-00ZF}{}\label{cdp:eq:C.RG.9}
\]
In the first identity of \eqref{cdp:eq:C.RG.8}, all coefficients
of $A$ below order $2d$ cancel. At order $2d$ the
remaining equality is
\[
                         a'\psi_i=a\psi_i'.
                                                               \tag{C.RG.10}\hypertarget{tag-00Z4}{}\label{cdp:eq:C.RG.10}
\]
Terms $s^{d+k}a_k(z)$ with $k\ge1$ first contribute
to the difference at order $2d+k$, so they do not
enter this calculation. Since $a$ is a unit,
$(\psi_i/a)'=0$. Hence $\psi_i=k_i a$ for a constant
$k_i\in\C$. At order $d$, the second identity gives
\[
 H'\psi_i=a\,\partial_z\delta_i(z,0)
          =-a\lambda_i',\qquad
             \lambda_i'= -k_i H'.
                                                               \tag{C.RG.11}\hypertarget{tag-00Z5}{}\label{cdp:eq:C.RG.11}
\]
Consequently there are constants $d_i\in\C$ with
\[
                         \lambda_i=d_i-k_iH
                    \quad(i=1,2)
                                                               \tag{C.RG.12}\hypertarget{tag-00Z6}{}\label{cdp:eq:C.RG.12}
\]
on the small disc. This uses only the two return germs
and one flow box. In particular it does not require
an ambient extension of $q$ or of $H$ around a fibre.
The formal coefficient extraction is independently
certified in \path{code/A_return_germs.py}; its executed
output records the higher-coefficient cancellations.

\par\medskip\noindent\textbf{Global periods and the singular fibres.}
Let $B^\circ$ be a connected complement of finitely
many values containing the disc used above, on which
the elliptic fibration, time form, and $H$ are regular.
Continue the chosen homology basis and periods on its
universal cover. Both sides of \eqref{cdp:eq:C.RG.12} are holomorphic
there, so the identity theorem gives the same formula
everywhere, with the same constants. Since $H$ is an
actual single-valued meromorphic function on $B$, each
continued period is single-valued on $B^\circ$.
The integration map from $H_1(F_b,\Z)$ to $\C$ is
injective for the time form of a nonzero elliptic
translation field. Thus returning both periods
unchanged forces every integral monodromy matrix to
be the identity.

\item The oriented ratio $\tau=\lambda_2/\lambda_1$ is
therefore a single-valued holomorphic map
$B^\circ\to\mathbb H$. The denominator never vanishes,
because a nonzero integral elliptic cycle has nonzero
period. The Cayley transform $(\tau-\mathrm i)/(\tau+\mathrm i)$
is bounded, extends holomorphically across all finitely
many omitted points, and is constant on the compact
curve $B$. Hence $\tau$ is constant. Integrating a
locally normalized time form identifies the smooth
fibres with torsors under one fixed elliptic curve.

\item Contract the vertical exceptional curves to obtain
\eqref{cdp:eq:C.RG.1}. The integral monodromy is unchanged. The
exact Kodaira-table and multiple-root comparison in
(C.RF.3) applies: for a full fibre $mC$, write
$t=u h^m$, make the actual change $t=s^m$, and
normalize by $w^m=u$, $s=wh$. This covering surface
is \'etale over the original one. Its reduced
relatively minimal fibre has identity monodromy,
so is the smooth type $I_0$. It follows that every
original reduced support $C$ is smooth elliptic.
This is precisely the comparison established from
\cite[Theorem~6.2; Section~8(iv)(11), p.~585]{Kodaira1963II} and \citet[Section~2.3, pp.~5--6]{UlmerUrzua2021} in
\ref{cdp:st:C.2.UN.22}; it includes all multiple fibres.

\par\medskip\noindent\textbf{The direct image and the genus bound.}
Over a small smooth base disc, a local holomorphic
section and integration of the marked differential
give an actual product with the fixed elliptic curve.
Since the integral monodromy is the identity, the
origin-preserving linear part of every overlap map
is the identity. The overlap maps are translations,
which act trivially on $H^1(\OO)$.
At a multiple value use the actual normalized root
family just described. It is smooth with constant
elliptic modulus. The cyclic action on its central
elliptic fibre is free, since the map to $E_0$ is
\'etale, and is therefore a translation. Its linear
part on the neighbouring marked elliptic fibres
is the same discrete automorphism and hence is
the identity. On $H^1(\OO)$ it acts trivially.
Taking invariants in the local root disc consequently
gives $R^1q_*\OO=\OO$ at the multiple value as well:
the coefficient functions invariant under
$s\mapsto\zeta_m s$ are exactly the holomorphic
functions of $t=s^m$. These local generators agree
with the preceding translation-invariant generators
over the punctured discs. This proves the first
assertion of \eqref{cdp:eq:C.RG.2} with its actual descent map.

\item Connected fibres give $q_*\OO_{E_0}=\OO_B$.
The coherent Leray sequence on the curve now gives
\[
 H^2(E_0,\OO_{E_0})\simeq H^1(B,\OO_B),\qquad
 h^1(E_0,\OO_{E_0})=g(B)+1.
                                                               \tag{C.RG.13}\hypertarget{tag-00Z7}{}\label{cdp:eq:C.RG.13}
\]
Serre duality and birational invariance give
$p_g(E)=g(B)$ and $\chi(\OO_E)=0$.
The independent compact-core inequality
\ref{cdp:st:C.2.UN.28} gives $p_g(E)\le1$, and therefore
$g(B)\le1$.

\item Every root-cover component of nonnegative Kodaira
dimension has either generic rank two, covered by
\ref{cdp:st:C.2.UN.22}, or rank one. In the latter case a rational
general first-integral fibre would make the surface
uniruled and give Kodaira dimension $-\infty$ by
the actual adjunction argument in \ref{cdp:st:C.2.UN.20}.
Its general fibre is therefore elliptic and the
present proof applies. This proves the stated
coherent-Euler and K3/Enriques exclusions on the
covering components.

\item Finally let an original component have Kodaira
dimension one. Its actual generically finite
component cover has Kodaira dimension at least
one by pluricanonical pullback, and at most one
by \ref{cdp:st:C.2.UN.15}. Its elliptic fibration is therefore
its connected-fibre Iitaka fibration. The cyclic
stabilizer acts on this fibration, since it acts
on the pluricanonical field, and the proper
quotient and graph comparisons in \ref{cdp:st:C.2.UN.19}--20
give the descended fibration on a smooth
birational model of the original component.
Its base is the quotient of $B$, so has genus
at most one. On a general fibre, the subgroup
fixing a general base point must act by
translations: an origin-preserving nontrivial
linear part would give a rational general quotient
fibre. Its normal line is trivial, so adjunction
restricts every positive pluricanonical line to
$\OO_{\mathbb P^1}(-2m)$; all pluricanonical sections
would then vanish on the dense family of these
fibres, contrary to Kodaira dimension one. The finite translation
subgroup of the fixed elliptic curve is locally
constant, because it consists of points of
bounded finite order. Thus the general quotient
fibre is one fixed elliptic curve as well.
This proves isotriviality through the actual
finite quotient, while leaving its exceptional
and singular fibres in place.

\end{enumerate}
\end{proof}






\begin{proposition}[Nonnegative-Kodaira-dimension components have tightly specified cyclic covers]\hypertarget{tag-01DA}{}\label{cdp:st:C.2.UN.31}
\StatementRecord{cdp:st:C.2.UN.31}{proved}{01DA}{\SeriesRef{cdp:st:C.2.UN.24}, \SeriesRef{cdp:st:C.2.UN.16}, \SeriesRef{cdp:st:C.2.UN.15}, \SeriesRef{cdp:st:C.2.UN.22}, \SeriesRef{cdp:st:C.2.UN.20}, \SeriesRef{cdp:eq:C.CS.2}, \SeriesRef{cdp:st:C.2.UN.29}, \SeriesRef{cdp:st:C.2.UN.28}, \SeriesRef{cdp:eq:C.CS.4}, \SeriesRef{cdp:st:C.3.1}, \SeriesRef{cdp:eq:C.CS.5}, \SeriesRef{cdp:st:C.2.UN.23}, \SeriesRef{cdp:st:C.2.UN.32}}

Let $f:U\to\Delta$ be an unrestricted proper degeneration
of connected complex two-tori from a smooth threefold.
Let $Z$ be the smooth minimal model of the normalization
of one central component, with $\kappa(Z)\ge0$.
There is an actual generically finite component cover
$E\to S^\nu$, obtained by the normalized root base change
and resolution, with cyclic stabilizer $G$, such that
$E/G$ is properly birational to $S^\nu$. Its smooth
minimal model $E_0$ satisfies one of these descriptions:

\begin{itemize}
\item $E_0$ is a complex two-torus;
\item $E_0$ has an isotrivial elliptic fibration over a
curve $B$ with $g(B)\le1$, identity integral $H_1$
monodromy, and only positive multiples of smooth
elliptic fibres. In this case
$R^1q_*\mathcal O_{E_0}=\mathcal O_B$,
$p_g(E_0)=g(B)$, and $\chi(\mathcal O_{E_0})=0$.
\end{itemize}
In both cases $\chi(\mathcal O_{E_0})=0$, and the
action of $G$ descends holomorphically to $E_0$.

For the original minimal component the consequences are
\[
 p_g(Z)\le1,
 \qquad \chi(\mathcal O_Z)\in\{0,1,2\}.
                                                               \tag{C.CS.1}\hypertarget{tag-00DH}{}\label{cdp:eq:C.CS.1}
\]
Every component with $\kappa(Z)=1$ has an actual
isotrivial elliptic Iitaka fibration on $Z$, over a
curve of genus at most one. An Enriques component
has an isotrivial elliptic fibration over $\mathbb P^1$.
A K3 component is either the cyclic torus quotient
of \ref{cdp:st:C.2.UN.24}, with the exact orders and configurations
listed there, or has an isotrivial elliptic fibration
over $\mathbb P^1$ through the non-torus-cover branch.
The singular and multiple fibres on these original
quotient surfaces are retained by the quotient maps;
the smooth-multiple assertion concerns the specified
cover $E_0$.
\end{proposition}



\begin{proof}

\begin{enumerate}[label={$\langle1\rangle$\arabic*.},leftmargin=*]
\item Take the actual component, cyclic stabilizer, and
proper birational quotient comparison in \ref{cdp:st:C.2.UN.16}
and (C.CQ.12). Pluricanonical pullback through a
resolution of the generically finite graph gives
\[
             \kappa(E)\ge\kappa(Z)\ge0.
                                                               \tag{C.CS.2}\hypertarget{tag-00DI}{}\label{cdp:eq:C.CS.2}
\]
The compact-component theorem \ref{cdp:st:C.2.UN.15} applies
also to this finite base-change family, so
$\kappa(E)\le1$. For the saturated limiting field
space, the generic-rank-two branch is covered
by \ref{cdp:st:C.2.UN.22}. In rank one the first integral
has rational or elliptic general fibre. Rational
fibres would force every positive pluricanonical
section to vanish, by the exact adjunction
argument of \ref{cdp:st:C.2.UN.20}, contradicting \eqref{cdp:eq:C.CS.2}.
Thus \ref{cdp:st:C.2.UN.29} applies. These two statements
give precisely the asserted torus or smooth-multiple
elliptic alternatives and coherent Euler characteristic
zero.

\item In the elliptic alternative the actual root-disc
invariant calculation of (C.RG.2)--(C.RG.13) gives
$R^1q_*\mathcal O=\mathcal O_B$ and
$p_g(E)=g(B)$. That calculation requires only
the fixed elliptic modulus, identity integral
monodromy and the smooth-multiple local charts,
so it applies to the elliptic alternative of
\ref{cdp:st:C.2.UN.22} as well. The independent compact-core
bound \ref{cdp:st:C.2.UN.28} gives $p_g(E)\le1$, and hence
$g(B)\le1$ in both branches.

\item All proper birational maps between the displayed
smooth surface models factor through actual
point blowups and contractions. On a nonruled
surface the minimal model is unique. Equivalently,
the exceptional curves of a minimal contraction
form a negative-definite configuration; its
$(-1)$-curves are disjoint, so their finite
$G$-orbits can be contracted equivariantly,
one stage at a time. This gives the holomorphic
action on $E_0$ and the exact comparison
\[
 \begin{array}{ccc}
 E&\longrightarrow&E_0\\
 \downarrow&&\downarrow\\
 E/G&\longrightarrow&E_0/G
 \end{array}
 \qquad E_0/G\dashrightarrow Z
       \quad\hbox{properly birational}.
                                                               \tag{C.CS.3}\hypertarget{tag-00DJ}{}\label{cdp:eq:C.CS.3}
\]
It is this action and quotient which are used
below.

\item Suppose first that $\kappa(Z)=1$. Then
\eqref{cdp:eq:C.CS.2} and \ref{cdp:st:C.2.UN.15} give $\kappa(E_0)=1$,
so the torus alternative is impossible. The
elliptic fibration is the connected-fibre
Iitaka fibration. To specify the invariance,
the group acts on every space of pluricanonical
sections. An equivariant graph resolution of
a sufficiently large pluricanonical map and
its Stein factor therefore give a $G$-equivariant
map to $B$. The same Stein base is the one
of the preceding elliptic fibration: pluricanonical
ratios restrict constantly to its elliptic
fibres, and Kodaira dimension one makes the
induced map of bases finite. Taking the
connected-fibre Stein factor removes that
finite map.

\item The actual quotient fibration is
\[
                   E'/G\longrightarrow B/G,
                                                               \tag{C.CS.4}\hypertarget{tag-00DK}{}\label{cdp:eq:C.CS.4}
\]
after this equivariant graph resolution.
For a general base point, let
$K=\ker(G\to\operatorname{Aut}B)$.
Its fibre is exactly $F/K$, as in \ref{cdp:st:C.2.UN.20},
and is connected. Every $F$ is isomorphic
to the same elliptic curve $A$.
The subgroup $K$ acts by translations:
a finite automorphism of an elliptic curve
with nontrivial linear part has a fixed
point, and its quotient is rational by
Riemann--Hurwitz. A rational general fibre
in \eqref{cdp:eq:C.CS.4} would annihilate every positive
pluricanonical section, contrary to
$\kappa(Z)=1$. The translation points
belong to the finite subgroup $A[|K|]$,
so they are locally constant under the
marked product identifications. Therefore
the quotient general fibre is the one
fixed elliptic curve $A/K$.
The quotient base has genus at most one
by Riemann--Hurwitz and $g(B)\le1$.

\item This fibration is the Iitaka fibration
of $Z$. Its map is holomorphic on the
minimal surface. One can also check the
descent on an actual common resolution
$Z\xleftarrow{\rho}X\xrightarrow q B/G$:
for a smooth general elliptic fibre $F$,
\[
 0=K_X\cdot F
   =\rho^*K_Z\cdot F+\sum_j a_jE_j\cdot F,
                   \qquad a_j>0.
                                                               \tag{C.CS.5}\hypertarget{tag-00DL}{}\label{cdp:eq:C.CS.5}
\]
On a minimal surface of Kodaira dimension
one the canonical class is nef; this is
the elliptic canonical-bundle formula
in the surface classification used in
\SeriesRef{cdp:st:C.3.1}. Every term in \eqref{cdp:eq:C.CS.5} is
nonnegative. Hence every exceptional
curve is vertical. Connected exceptional
fibres and $\rho_*\mathcal O_X=\mathcal O_Z$
then descend the actual proper graph to
a holomorphic map on $Z$. This proves
the claimed isotrivial Iitaka fibration.

\item Now let $Z$ be Enriques. Algebraic
dimension is preserved by the actual
generically finite component map, by
the symmetric-function field argument
in (C.K3.2)--(C.K3.3). Thus $E_0$ is
projective. If it has Kodaira dimension
one, the preceding intrinsic Iitaka
quotient argument again gives an
isotrivial elliptic fibration on a
birational model of $Z$. If it has
Kodaira dimension zero, the surface
classification and $\chi(\mathcal O_{E_0})=0$
make it Abelian or bielliptic.
The Abelian case is excluded by the
exact cyclic quotient theorem \ref{cdp:st:C.2.UN.23}.
For the bielliptic case, use the actual
orbit fibration of its connected
automorphism group. In its finite free
product presentation $(A_1\times A_2)/H$,
$H$ acts by translations on $A_1$ and
has nontrivial linear action on $A_2$.
The invariant tangent fields are exactly
the one-dimensional translation space
of $A_1$. Thus the image of the compact
elliptic translation group $A_1$ is
the whole connected automorphism group,
and its orbit map is the explicit
elliptic fibration to $A_2/H\simeq\mathbb P^1$.
Its fibres are smooth elliptic multiples;
its integral monodromy is the identity,
since all actions on the fibre direction
are translations. Every automorphism
normalizes the connected automorphism
group, so preserves this orbit map and
acts on its connected-fibre base. It
is therefore $G$-equivariant. The quotient
\eqref{cdp:eq:C.CS.4} is available again. Its
general fibre must be elliptic, since
a rational general fibre would force
the nonzero even pluricanonical form
of an Enriques surface to vanish.

\item In either case, \eqref{cdp:eq:C.CS.5} descends the
elliptic map to the minimal Enriques
surface: now $K_Z$ is torsion, so its
intersection term is zero. The Leray
injection of $H^1$ of the base into
$H^1(Z,\mathcal O_Z)=0$ makes the
base $\mathbb P^1$. This proves the
Enriques assertion through the actual
quotient fibration.

\item For a K3 target, the nonzero canonical
form gives $p_g(E_0)\ge1$. If $E_0$
is a torus, \ref{cdp:st:C.2.UN.23}--24 give exactly
the cyclic symplectic torus quotient
already stated. Otherwise it has an
elliptic fibration of the specified
fixed-modulus type. If $\kappa(E_0)=1$,
use the equivariant Iitaka fibration.
If $a(E_0)=1$, use the intrinsic
field identity $\mathcal M(E_0)=q^*\mathcal M(B)$
and the equivariant Stein factor in
(C.K3.4). The remaining projective
Kodaira-dimension-zero case would be
bielliptic, but has $p_g=0$ and is
excluded. Algebraic dimension zero
has only the torus alternative by
\ref{cdp:st:C.2.UN.16}. Thus these are all cases.
Taking the quotient gives an isotrivial
elliptic fibration, since its general
fibre cannot be rational on a K3
birational model. Formula \eqref{cdp:eq:C.CS.5}
with $K_Z=0$ descends it to $Z$,
and $H^1(Z,\mathcal O_Z)=0$ again
makes its base $\mathbb P^1$.

\item Finally \ref{cdp:st:C.2.UN.28} gives $p_g(Z)\le1$.
The compact surface classification
gives $\chi(\mathcal O_Z)\ge0$ when
$\kappa(Z)\ge0$: in Kodaira dimension
zero its values are the explicit
torus, bielliptic, Kodaira, K3 and
Enriques values; in Kodaira dimension
one the relatively minimal elliptic
surface has $K_Z^2=0$ and nonnegative
Euler number by Kodaira's fibre list,
so surface Riemann--Roch gives the
inequality. On the other hand
\[
        \chi(\mathcal O_Z)=1-h^1(Z,\mathcal O_Z)+p_g(Z)
                              \le2.
                                                               \tag{C.CS.6}\hypertarget{tag-00DM}{}\label{cdp:eq:C.CS.6}
\]
It is an integer, giving \eqref{cdp:eq:C.CS.1}.
The later exact quotient theorem
\ref{cdp:st:C.2.UN.32} further identifies every K3
component, in all algebraic dimensions,
with the cyclic torus quotient family.
The preceding non-torus-cover comparison
is retained as its explicit intermediate
fibration map.

\item All these comparisons retain the
finite quotient's singular fibres;
they do not identify them with the
smooth fibres of its root cover.

\end{enumerate}
\end{proof}






\begin{proposition}[A properly elliptic normalization component has one of eleven positive-coherent-Euler fibre configurations]\hypertarget{tag-01DC}{}\label{cdp:st:C.2.UN.33}
\StatementRecord{cdp:st:C.2.UN.33}{proved}{01DC}{\SeriesRef{cdp:st:C.2.UN.16}, \SeriesRef{cdp:st:C.2.UN.15}, \SeriesRef{cdp:st:C.2.UN.22}, \SeriesRef{cdp:st:C.2.UN.29}, \SeriesRef{cdp:st:C.2.UN.19}, \SeriesRef{cdp:eq:C.EQ.3}, \SeriesRef{cdp:st:C.2.UN.32}, \SeriesRef{cdp:eq:C.EQ.4}, \SeriesRef{cdp:eq:C.EQ.5}, \SeriesRef{cdp:eq:C.EQ.7}, \SeriesRef{cdp:eq:C.EQ.8}, \SeriesRef{cdp:eq:C.EQ.10}, \SeriesRef{cdp:eq:C.EQ.11}, \SeriesRef{cdp:st:C.2.UN.28}; \citet[Section~2.3, pp.~5--6]{UlmerUrzua2021}; \citet[Theorem~6.2 and Section~8(iv)]{Kodaira1963II}}

Let $M$ be the smooth minimal model of a normalization
component of a fibre of an unrestricted two-torus degeneration,
and suppose $\kappa(M)=1$. Its connected elliptic Iitaka map
$p:M\to B_0$ is isotrivial, and $g(B_0)\le1$.
Its coherent Euler characteristic belongs to $\{0,1,2\}$.
When it is zero, every fibre of $p$ is a positive multiple
of a smooth elliptic curve. When it is positive,
$B_0\simeq\PP^1$, $q(M)=0$, and the complete configuration
of the fibres of $p$ with singular reduced support is one
of the following rows:
\[
\begin{array}{c|l|c}
\chi(\OO_M)&\text{singular reduced fibres of }p&e(M)\\ \hline
1&\mathrm{II}+\mathrm{II}^*&12\\
1&\mathrm{III}+\mathrm{III}^*&12\\
1&\mathrm{IV}+\mathrm{IV}^*&12\\
1&2\mathrm I_0^*&12\\
1&3\mathrm{IV}&12\\
1&2\mathrm{III}+\mathrm I_0^*&12\\
1&\mathrm{II}+\mathrm{IV}+\mathrm I_0^*&12\\ \hline
2&4\mathrm I_0^*&24\\
2&3\mathrm{IV}^*&24\\
2&2\mathrm{III}^*+\mathrm I_0^*&24\\
2&\mathrm{II}^*+\mathrm{IV}^*+\mathrm I_0^*&24 .
\end{array}                                                   \tag{C.EQ.1}\hypertarget{tag-00IC}{}\label{cdp:eq:C.EQ.1}
\]
In each row $p_g(M)=\chi(\OO_M)-1$. All additional
smooth multiple fibres are retained as part of the specified
surface; the table lists its nonsmooth reduced fibres.
This is a necessary classification of the displayed
invariant for actual components. The statement does not
assert occurrence of every abstract surface with these data.
The map $p$ is the elliptic map on the component surface;
the original threefold map has the degeneration disc as target.
\end{proposition}



\begin{proof}

\begin{enumerate}[label={$\langle1\rangle$\arabic*.},leftmargin=*]
\item Take the actual cyclic component cover of \ref{cdp:st:C.2.UN.16}.
Pluricanonical pullback gives Kodaira dimension at least
one, and \ref{cdp:st:C.2.UN.15} excludes dimension two. Its minimal model
$E$ therefore has Kodaira dimension one. The results
\ref{cdp:st:C.2.UN.22} and \ref{cdp:st:C.2.UN.29} give its connected elliptic Iitaka map
\[
 q:E\longrightarrow B,\qquad g(B)\le1,                 \tag{C.EQ.2}\hypertarget{tag-00IF}{}\label{cdp:eq:C.EQ.2}
\]
with a fixed elliptic fibre, identity integral monodromy,
and only smooth multiple fibres. The cyclic stabilizer
$G$ preserves $q$, because it acts on the pluricanonical
field. The normalization of the original component is
birational to the actual finite quotient $E/G$, through
the equivariant graph and minimal-model comparisons of
\ref{cdp:st:C.2.UN.19}--20. Thus its relatively minimal elliptic model is
the model of that quotient over $B/G$.

\item Let $H$ be the kernel of $G\to\operatorname{Aut}(B)$.
On a general elliptic fibre it acts by translations.
Indeed a nontrivial linear elliptic automorphism has a
fixed point and its quotient is rational by
Riemann--Hurwitz. A rational general fibre on the
quotient would make every positive pluricanonical
section vanish by adjunction, contrary to $\kappa(M)=1$.
The finite translations form a fixed subgroup of the
constant elliptic curve, since points of bounded order
are discrete. We must retain the action at multiple
fibres when taking the quotient by $H$.

\par\medskip\noindent\textbf{The quotient by fibre translations.}
At a full fibre $mF$ of $q$, the actual normalized root
chart is
\[
 (F'\times\Delta_z)/\langle\delta\rangle,
 \qquad \delta(v,z)=(v+\tau,\zeta_m z),
 \qquad \operatorname{ord}(\tau)=m,
 \qquad t=z^m.                                        \tag{C.EQ.3}\hypertarget{tag-00IG}{}\label{cdp:eq:C.EQ.3}
\]
The root family is a smooth fixed-elliptic family.
The cyclic deck action is a translation on every fibre,
and its free action on the central fibre gives exactly
the indicated order of $\tau$. This is the normalized
root comparison proved in \ref{cdp:st:C.2.UN.22} and \ref{cdp:st:C.2.UN.29}, with the
constant-translation gauge written in (C.KX.4).

\item Lift a base-fixing element of $H$ to this chart.
Its multiplier on $z$ is a power of $\zeta_m$;
composition with that inverse power of $\delta$
makes it fix $z$. The appropriate finite power
is a deck transformation fixing $z$, and hence
the identity. Its fibre action is therefore a
translation by a finite-order constant point.
The lifts so chosen generate a finite translation
subgroup $T\subset F'$, commuting with $\delta$.
First quotient by $T$. Write $\bar F=F'/T$ and
$m'=\operatorname{ord}(\bar\tau)$, so $m'\mid m$.
The remaining action is
$(\bar v,z)\mapsto(\bar v+\bar\tau,\zeta_m z)$.
Its $m'$-th power is a pure rotation of the disc.
Quotienting that subgroup replaces $z$ by
$w=z^{m/m'}$; the remaining free action is
\[
 (\bar v,w)\longmapsto(\bar v+\bar\tau,\zeta_{m'}w),
 \qquad t=w^{m'}.                                    \tag{C.EQ.4}\hypertarget{tag-00IH}{}\label{cdp:eq:C.EQ.4}
\]
This proves directly that $E_1=E/H$ is smooth near
the fibre and has the full smooth multiple fibre
$m'(\bar F/\langle\bar\tau\rangle)$. It is smooth
elsewhere by the same local translation quotient.
Thus $E_1\to B$ still has a fixed elliptic fibre,
identity integral monodromy, and only smooth multiples.
The remaining cyclic group $Q=G/H$ acts faithfully
on $B$.

\par\medskip\noindent\textbf{The local action at a base branch value.}
Let a stabilizer of order $s$ act at $b\in B$, and
let $\beta$ be its linear action on the fixed elliptic
curve. The linear action is constant along the regular
base, because it belongs to a discrete finite automorphism
group. Suppose first that $\beta\ne1$.
We prove that the fibre of $E_1$ at $b$ has multiplicity
one; this prevents a hidden change of the singular-fibre
type in the quotient.

\item Use \eqref{cdp:eq:C.EQ.3} for a putative multiplicity $m$.
The generator lifts to the normalized root chart,
with normal multiplier $\rho$ satisfying
$\rho^m=\zeta_s$, and affine fibre derivative $\beta$.
It commutes with the deck action: its conjugation
preserves the multiplier $\zeta_m$ on the root
coordinate, and hence preserves the corresponding
deck generator. Restriction to the central elliptic
curve gives $\beta\tau=\tau$. The finite-order
endomorphism $\beta\ne1$ satisfies
$1+\beta+\cdots+\beta^{s-1}=0$. Applying it to
$\tau$ yields $s\tau=0$, so $m\mid s$.

\item The induced affine map on the reduced elliptic
fibre has a fixed point, because $\beta-1$ is a
surjective elliptic isogeny. Choose a lift of that
point, and compose the lifted generator with a
deck transformation to make it fix the chosen
lift. Its $s$-th power is a deck transformation
fixing a point of the central elliptic curve;
the deck action there is free, so that power is
the identity. Therefore $\rho^s=1$.
Together with $\rho^m$ having order $s$, this
gives $\gcd(m,s)=1$. Since $m\mid s$, we have
\[
                 m=1\qquad(\beta\ne1).               \tag{C.EQ.5}\hypertarget{tag-00II}{}\label{cdp:eq:C.EQ.5}
\]
This is the same explicit local comparison used
in \ref{cdp:st:C.2.UN.32}, and the argument did not use a canonical
form or a K3 hypothesis.

\item When $\beta=1$, every element with a fixed point
on a reduced elliptic fibre fixes that whole fibre:
an elliptic translation with a fixed point is the
identity. Its stabilizer acts only on the normal
coordinate. A finite holomorphic action with
identity derivative is locally the identity,
by the averaging linearization, so every nontrivial
stabilizer is a normal reflection. The quotient
surface is smooth and the reduced fibre remains
smooth. Equivalently, the translation and root
quotients \eqref{cdp:eq:C.EQ.3}--\eqref{cdp:eq:C.EQ.4} display that fibre
as a smooth multiple. Hence these values have
zero topological Euler contribution.

\item For $\beta\ne1$, \eqref{cdp:eq:C.EQ.5} gives a smooth product
near the source fibre. An equivariant local section
through a fixed point removes the translation part
of the finite action. To construct it, lift its
translation cocycle to small holomorphic complex-valued
functions near the fixed point and average over the
finite group. The cocycle sum is zero because the
action has finite order and the lift vanishes at
the fixed point. In additive notation, for an operator
$P$ of order $s$ and a cocycle $c$ with
$\sum_{k=0}^{s-1}P^kc=0$, the function
$s^{-1}\sum_{k=0}^{s-1}kP^kc$ solves
$(P-1)b=c$. This supplies the claimed holomorphic
section and gives a product action.
Quotient the subgroup acting trivially on the
elliptic factor. It consists of pure normal
reflections. The remaining action has order
$d\in\{2,3,4,6\}$ and the exact form
\[
             (v,z)\longmapsto(\zeta_d^r v,\zeta_d z),
             \qquad r\in\{1,d-1\}.                   \tag{C.EQ.6}\hypertarget{tag-00IJ}{}\label{cdp:eq:C.EQ.6}
\]
Translations of an origin have already been removed.

\item Here is an explicit comparison with the Kodaira
types, including the direction of the character.
Write the fixed elliptic curve as
$y^2=x^3+Ax+B$, so multiplication by $\zeta_d^r$
acts on $x,y$ by $\zeta_d^{-2r},\zeta_d^{-3r}$.
The invariant meromorphic functions
\[
          T=z^d,\qquad X=z^{2r}x,\qquad Y=z^{3r}y
                                                               \tag{C.EQ.7}\hypertarget{tag-00IK}{}\label{cdp:eq:C.EQ.7}
\]
generate the quotient function field. They give
the following minimal Weierstrass equations:
\[
\begin{array}{c|c|c|c}
d&r&\text{equation}&\text{Kodaira type}\\ \hline
2&1&Y^2=X^3+AT^2X+BT^3&\mathrm I_0^*\\
3&1&Y^2=X^3+BT^2&\mathrm{IV}\\
3&2&Y^2=X^3+BT^4&\mathrm{IV}^*\\
4&1&Y^2=X^3+ATX&\mathrm{III}\\
4&3&Y^2=X^3+AT^3X&\mathrm{III}^*\\
6&1&Y^2=X^3+BT&\mathrm{II}\\
6&5&Y^2=X^3+BT^5&\mathrm{II}^* .
\end{array}                                                   \tag{C.EQ.8}\hypertarget{tag-00IL}{}\label{cdp:eq:C.EQ.8}
\]
For $d=3,6$ the curve has $A=0$ and $B\ne0$;
for $d=4$ it has $B=0$ and $A\ne0$.
The equations are minimal, since their indicated
orders are below the Weierstrass rescaling threshold.
The primitive monodromy matrices and their characters distinguish the
unstarred and starred types; see \citet[Section~2.3, pp.~5--6]{UlmerUrzua2021}
and \citet[Theorem~6.2 and Section~8(iv)]{Kodaira1963II}.
For the displayed minimal equations the discriminant orders are computed
directly as $12r/d$, and the corresponding Kodaira configurations have
the same Euler numbers.
The map \eqref{cdp:eq:C.EQ.7} is birational over the base
after taking the finite quotient. Resolving its
graph and contracting vertical exceptional curves
gives the same relatively minimal elliptic model;
therefore these are the fibre types of the actual
quotient. This verifies the reduction to that table,
including the two different types of each primitive
order. Their Euler numbers are $12r/d$.

\par\medskip\noindent\textbf{All faithful cyclic base actions.}
On $B=\PP^1$, a nontrivial cyclic action is conjugate
to $z\mapsto\zeta_nz$. Its only branch values are
zero and infinity, where positive local generators
are inverse. A nontrivial fibre character of order
$d$ consequently gives the two opposite exponents
in \eqref{cdp:eq:C.EQ.8}. The resulting pairs are
\[
 (\mathrm{II},\mathrm{II}^*),\quad
 (\mathrm{III},\mathrm{III}^*),\quad
 (\mathrm{IV},\mathrm{IV}^*),\quad
 (\mathrm I_0^*,\mathrm I_0^*).                         \tag{C.EQ.9}\hypertarget{tag-00IM}{}\label{cdp:eq:C.EQ.9}
\]
Every pair has Euler sum twelve. A trivial fibre
character gives only smooth multiples.

\item On an elliptic $B$, a finite cyclic group either
acts by translations, in which case it is free,
or has a fixed point. In the latter case translate
that point to the origin. The action is linear of
order $n=2,3,4,6$, as follows from its action on
the rank-two integral period lattice. The branch
indices of $B\to B/Q\simeq\PP^1$ are
\[
\begin{array}{c|c}
n&\text{branch indices}\\ \hline
2&(2,2,2,2)\\
3&(3,3,3)\\
4&(4,4,2)\\
6&(6,3,2).
\end{array}                                                   \tag{C.EQ.10}\hypertarget{tag-00ID}{}\label{cdp:eq:C.EQ.10}
\]
For completeness, the number fixed by a nontrivial
power is $|1-\zeta_n^k|^2$, the determinant of
that power minus one on the lattice. Removing
the points with larger stabilizer and dividing
by the corresponding orbit sizes gives exactly
\eqref{cdp:eq:C.EQ.10}. Each row also satisfies the degree-$n$
Riemann--Hurwitz equation.

\item Write the fibre character as $\zeta_n^a$,
$0\le a<n$. At a branch index $e$, its exponent
is $a$ modulo $e$. Exponent zero gives a smooth
multiple. Otherwise put $c=\gcd(a,e)$; quotienting
the normal reflection kernel gives the entry of
\eqref{cdp:eq:C.EQ.8} with $d=e/c$ and $r=(a\bmod e)/c$.
Enumerating these actual characters gives
\[
\begin{array}{c|c|l}
n&a&\text{nonsmooth fibre configuration}\\ \hline
2&1&4\mathrm I_0^*\\
3&1&3\mathrm{IV}\\
3&2&3\mathrm{IV}^*\\
4&1&2\mathrm{III}+\mathrm I_0^*\\
4&2&2\mathrm I_0^*\\
4&3&2\mathrm{III}^*+\mathrm I_0^*\\
6&1&\mathrm{II}+\mathrm{IV}+\mathrm I_0^*\\
6&2,4&\mathrm{IV}+\mathrm{IV}^*\\
6&3&2\mathrm I_0^*\\
6&5&\mathrm{II}^*+\mathrm{IV}^*+\mathrm I_0^* .
\end{array}                                                   \tag{C.EQ.11}\hypertarget{tag-00IE}{}\label{cdp:eq:C.EQ.11}
\]
The zero character gives no nonsmooth fibre.
These finite counts, the local reflection reductions,
and the Euler sums are independently executed in
\path{code/root_cyclic_elliptic_types.py} and its
saved output.

\item We finish by identifying the coherent invariants
without adding any unproved occurrence assertion.
The singular values just listed are all the possible
ones: a free base orbit is locally unchanged, and
the source had only smooth multiple fibres.
A free cyclic action on an elliptic base also leaves
only smooth multiples. Thus positive Euler occurs
only over $B_0=\PP^1$ and with nontrivial fibre
linear character. The exact smooth-multiple comparison
of \ref{cdp:st:C.2.UN.29} applies to $E_1$ and gives $p_g(E_1)=g(B)$.
The generically finite quotient comparison, resolved
by its actual graph, injects holomorphic two-forms
on $M$ into the invariant forms on $E_1$, as in
(C.PG.4)--(C.PG.5).
For $B=\PP^1$ this immediately gives $p_g(M)=0$.
For elliptic $B$, the unique two-form on $E_1$ is
the wedge of a base differential and the fibre
differential with fixed normalized periods. At a
smooth multiple fibre the same form extends with
the zero of order $m-1$ supplied by the root
coordinate; hence it is a global holomorphic form.
Its character is $\zeta_n^{a+1}$. Consequently
$p_g(M)=0$ whenever $a+1\not\equiv0\pmod n$.
Every Euler-twelve row of \eqref{cdp:eq:C.EQ.11} has precisely
that nontrivial character. This argument uses
pullback through the actual graph and Hartogs
extension, and does not infer extension of a
form from numerical invariance alone.

\item On a relatively minimal elliptic surface the canonical
divisor is vertical. An actual meromorphic differential
on the base and a meromorphic relative differential
give such a divisor; restriction to a general compact
elliptic fibre has no zero or pole. The components
of the displayed reducible singular fibres are $(-2)$
rational curves. A type $\mathrm{II}$ fibre has arithmetic
genus one and square zero, and a smooth reduced multiple
fibre is an elliptic curve of square zero. Adjunction gives
$K_M\cdot C=0$ for each of them, and hence $K_M^2=0$.
Analytic Riemann--Roch gives
$12\chi(\OO_M)=e(M)$. Euler localization for the
elliptic map $p$ and \eqref{cdp:eq:C.EQ.8}--\eqref{cdp:eq:C.EQ.11} give
exactly the rows of \eqref{cdp:eq:C.EQ.1}, or zero. For Euler
twelve we proved $p_g=0$, so $\chi=1$ gives $q=0$.
For Euler twenty-four, \ref{cdp:st:C.2.UN.28} gives $p_g\le1$, while
$\chi=2=1-q+p_g$ and $q\ge0$ force $p_g=1,q=0$.
In the zero case
the local list has no nonsmooth reduced fibre,
since every such entry has positive Euler number.
This proves all assertions while preserving every
smooth multiple fibre and every actual quotient map.

\end{enumerate}
\end{proof}






\begin{proposition}[Compact component cores give a multiplicity-weighted invariant lattice in one nearby torus]\hypertarget{tag-01DF}{}\label{cdp:st:C.2.UN.36}
\StatementRecord{cdp:st:C.2.UN.36}{proved}{01DF}{\SeriesRef{cdp:eq:C.CL.8}, \SeriesRef{cdp:eq:C.CL.9}, \SeriesRef{cdp:eq:C.CL.7}, \SeriesRef{cdp:st:C.2.UN.33}, \SeriesRef{cdp:st:C.2.UN.11}, \SeriesRef{cdp:eq:C.CL.10}; \cite[Theorem~3.44, pp.~254--255; \S3.G, p.~321]{Hatcher2002}; \cite[Chapter~1, Proposition~3.5, p.~17]{HuybrechtsK3Draft}; \cite[Chapter~1, \S1, pp.~5--6; \S2, p.~8; \S3, p.~9]{Peters1990}}

Let $f:U\to\Delta$ be a proper holomorphic map from
a smooth complex threefold, with connected two-torus
fibres over $\Delta^*$. Write its full central fibre
as $\sum_i m_iS_i$, and put $S=\bigcup_iS_i$ with
its reduced structure. For each component choose a
resolution of its normalization
$h_i:Z_i\to S_i\subset S$ which is an isomorphism
over $S_i\setminus\operatorname{Sing}S$. Set
\[
 \begin{gathered}
 D_i=h_i^{-1}(\operatorname{Sing}S)_{\rm red},\qquad
 R_i=\operatorname{span}_{\mathbb R}
       \{[C]:C\text{ is a curve component of }D_i\},
 \\
 K_i=\{v\in H_2(Z_i,\Z)/\operatorname{tors}:v\cdot R_i=0\}.
 \end{gathered}
                                                               \tag{C.CL.1}\hypertarget{tag-00BW}{}\label{cdp:eq:C.CL.1}
\]
The set $D_i$ may also have isolated points; they
contribute no curve class to $R_i$. The pairing on
$K_i$ is the restriction of the ordinary oriented
surface intersection form and may have a radical.

Choose a common base-change order $N$ divisible by
all $m_i$ and by the semisimple monodromy order.
Let $U_1$ be the monodromy on the integral first
homology of the actual nearby torus after this
base change. There is a homomorphism, obtained by
choosing compact representatives and their transfers,
\[
 \Psi:\bigoplus_iK_i\longrightarrow
       H_2(F,\Z)^{\Lambda^2U_1},\qquad
 \Psi(v)\cdot\Psi(w)=\sum_i m_i(v_i\cdot w_i).
                                                               \tag{C.CL.2}\hypertarget{tag-00C1}{}\label{cdp:eq:C.CL.2}
\]
Thus every multiplicity appears in the preserved
bilinear form. The kernel is contained in the
radical of that form. If its source is nondegenerate
of rank six, its image has finite integral index
in $H_2(F,\Z)$ and
\[
 \left|\det\left(\bigoplus_i m_iQ_{K_i}\right)\right|
                       \text{ is an integer square}.
                                                               \tag{C.CL.3}\hypertarget{tag-00C2}{}\label{cdp:eq:C.CL.3}
\]

Let $(p_i,q_i,z_i)$ be the positive, negative and
radical dimensions of the real form on $K_i$.
Let $(p_U,q_U,z_U)$ be those of the intersection
form restricted to
$\ker(\Lambda^2U_1-I)\subset H_2(F,\mathbb R)$.
Then
\[
             \sum_i p_i\le p_U\le3,\qquad
             \sum_i q_i\le q_U\le3.                 \tag{C.CL.4}\hypertarget{tag-00C3}{}\label{cdp:eq:C.CL.4}
\]
In particular, if $r_i=\dim R_i$, then
\[
 \sum_i\max\{b_2^-(Z_i)-r_i,0\}\le q_U\le3,
 \qquad
                 2\sum_i p_g(Z_i)\le p_U.           \tag{C.CL.5}\hypertarget{tag-00C4}{}\label{cdp:eq:C.CL.5}
\]
A K3 component whose chosen resolution is obtained
from its minimal model by $n$ point blowups must
therefore have at least $16+n$ independent boundary
curve classes. For an Enriques component or a
properly elliptic component of coherent Euler
characteristic one, the corresponding lower bound
is $6+n$.

The existing monodromy theorem gives
$(U_1-I)^2=0$. In an oriented real basis of
$H_1(F,\mathbb R)$, all its possible Jordan forms
give the following exact invariant forms:
\[
\begin{array}{c|c|c|c}
\operatorname{rank}(U_1-I)&\text{nonzero values of }U_1-I
 &\dim\ker(\Lambda^2U_1-I)&(p_U,q_U,z_U)\\ \hline
0&\text{none}&6&(3,3,0)\\
1&e_2\mapsto e_1&4&(1,1,2)\\
2&e_2\mapsto e_1,\ e_4\mapsto e_3&4&(1,2,1)\\
2&e_2\mapsto e_1,\ e_4\mapsto-e_3&4&(2,1,1).
\end{array}                                                   \tag{C.CL.6}\hypertarget{tag-00C5}{}\label{cdp:eq:C.CL.6}
\]
The two last rows record the two oriented real
forms of that Jordan type. In particular rank-one
unipotent monodromy forces $p_g(Z_i)=0$ for every
normalization component. All these assertions
apply to arbitrary actual central fibres, with
no normal-crossing, reducedness or K\"ahler
hypothesis on the original fibre.
\end{proposition}



\begin{proof}

\begin{enumerate}[label={$\langle1\rangle$\arabic*.},leftmargin=*]
\item The open set
$Z_i^\circ=Z_i\setminus D_i$ maps isomorphically
to $S_i\setminus\operatorname{Sing}S$. It is
therefore precisely the part of this component
where the reduced support is one smooth embedded
surface. The compact analytic set $D_i$ consists
of curves and possibly isolated points. Its
degree-two integral cohomology is free, with
one generator for each curve component. Indeed
normalization of a compact reduced curve changes
only the identifications at finitely many points,
which do not change its degree-two cohomology;
isolated points have none.

\item Poincar\'e--Alexander duality
\cite[Theorem~3.44, pp.~254--255]{Hatcher2002}
and the relative homology sequence give the exact map
\[
 H_2(Z_i^\circ,\Z)\longrightarrow H_2(Z_i,\Z)
 \longrightarrow H_2(Z_i,Z_i^\circ;\Z)
       \simeq H^2(D_i,\Z).                          \tag{C.CL.7}\hypertarget{tag-00C6}{}\label{cdp:eq:C.CL.7}
\]
The second arrow, under these identifications,
is the vector of intersections with the curve
components of $D_i$. This follows by applying
Poincar\'e duality to the restriction map to
that compact analytic set. Torsion maps to
zero since the target is free. Consequently
every element of $K_i$ has an integral lift
represented by a compact finite cycle in
$Z_i^\circ$. Choose such lifts for an integral
basis of each $K_i$. Their mutual intersections
can be computed by compact transverse smooth
chain representatives in that same open set.
Include all the required perturbations in
one compact subset for each component.

\item Make the actual change $t=s^N$, normalize,
and resolve the total space and central support,
isomorphically over the smooth open sets just
specified and over $s\ne0$. Let
$g:\widetilde U\to\Delta_s$ be the resulting
proper map from a smooth threefold. Over a
point of $S_i\setminus\operatorname{Sing}S$
the original local equation is
$t=u x^{m_i}$, with $u$ a unit. Its normalized
root charts are
\[
       w^{m_i}=u,\qquad x=s^{N/m_i}/w.               \tag{C.CL.8}\hypertarget{tag-00C7}{}\label{cdp:eq:C.CL.8}
\]
At $s=0$ the unit-root equation has exactly
$m_i$ local solutions, each nonsingular.
The full collection of strict components
above $S_i$ therefore restricts to a finite
\'etale covering
\[
        \pi_i:\widetilde Z_i^\circ
                       \longrightarrow Z_i^\circ
                  \quad\text{of total degree }m_i.  \tag{C.CL.9}\hypertarget{tag-00C8}{}\label{cdp:eq:C.CL.9}
\]
It may be disconnected. Every one of these
strict components has multiplicity one;
the root normal coordinate is $s$.
No modification used in the resolution
changes \eqref{cdp:eq:C.CL.8}--\eqref{cdp:eq:C.CL.9}.

\item Transfer an integral compact cycle by taking
the sum of all its lifts under \eqref{cdp:eq:C.CL.9},
subdividing its finite singular chain into
evenly covered pieces. This is the ordinary
integral covering transfer
\cite[\S3.G, p.~321]{Hatcher2002}
and is a chain map. For two transverse compact cycles,
every signed intersection point downstairs
has exactly $m_i$ lifts with the same local
orientation. Hence their transferred
intersection is exactly $m_i$ times the
original intersection. This proves the
multiplicity factor directly, including
disconnected root covers.

\item The transferred cycles lie in the smooth
single-component locus of the central
support of $\widetilde U$. The map $g$ is
a submersion on a neighbourhood of their
compact union. Choose smooth horizontal
lifts of the two real coordinate fields
on the base disc. A partition of unity
preserves their specified images under
$dg$. After a cutoff outside a slightly
larger neighbourhood, their successive
real and imaginary flows give, for every
sufficiently small $s$, an actual smooth
embedding of a neighbourhood of the compact
cycles into $g^{-1}(s)$. The image lies in
that exact fibre, and the embedding preserves
the complex orientation by continuity.
Choose disjoint ambient neighbourhoods for
the compact subsets in distinct strict
components, and the same small disc for
all flows. Their images then remain
disjoint in one nearby torus $F$.

\item Transport the selected transferred basis
cycles and extend by integer linearity.
This defines the homomorphism $\Psi$.
The transfer intersection calculation and
the disjoint neighbourhoods give the
bilinear equality of \eqref{cdp:eq:C.CL.2}. The choice
of lifts is part of this construction;
no canonical splitting of \eqref{cdp:eq:C.CL.7} is
needed.

\item The image is invariant under the actual
nearby monodromy. In fact the two-flow
construction defines a smooth family of
each finite cycle over the \emph{whole}
small $s$-disc. Its restriction to a
circle in that disc returns to the
identical parametrized cycle. In any
local differentiable trivialization of
the proper smooth family over $s\ne0$,
this continuous family has a locally
constant integral homology class.
Thus the class returns unchanged under
Gauss--Manin transport. The actual
identification
$H_2(F,\Z)\simeq\Lambda^2H_1(F,\Z)$
intertwines transport, so its image is
in the group written in \eqref{cdp:eq:C.CL.2}.
Only these explicitly transported
cycles enter the assertion.

\item If $\Psi(v)=0$, the preserved pairing
with every $\Psi(w)$ shows that $v$
is in the radical of the source form.
Choose a real complement of that radical
on which the restricted form is
nondegenerate. Its image is injective
and has the same positive and negative
indices. Maximal positive and negative
subspaces of the target restricted form
therefore give \eqref{cdp:eq:C.CL.4}. If the integral
source is nondegenerate of rank six,
$\Psi$ is an injection of full-rank
free lattices into $H_2(F,\Z)$.
The latter has the unimodular form
$3H$, of determinant $-1$. If the
image has index $c$, the change-of-basis
determinant gives source determinant
$-c^2$. This proves \eqref{cdp:eq:C.CL.3}.

\item Intersect a maximal negative definite
subspace of $H_2(Z_i,\mathbb R)$ with
the codimension-$r_i$ subspace $K_i\otimes\mathbb R$.
Its dimension is at least
$b_2^-(Z_i)-r_i$, so its negative index
is at least the corresponding nonnegative
maximum in \eqref{cdp:eq:C.CL.5}. For the second
inequality there, the real and imaginary
parts of holomorphic two-forms give the
positive $2p_g(Z_i)$-dimensional subspace
proved in (C.PG.7)--(C.PG.8). Their
integrals on every boundary curve
vanish by type; their dual homology
classes therefore belong to
$K_i\otimes\mathbb R$ by \eqref{cdp:eq:C.CL.7}.
This proves \eqref{cdp:eq:C.CL.5}.

\item A K3 minimal surface has intersection
signature $(3,19)$
\cite[Chapter~1, Proposition~3.5, p.~17]{HuybrechtsK3Draft};
$n$ point blowups add $n$ orthogonal classes of square $-1$.
The point-blowup splitting, and the Enriques invariants
and signature formula giving signature $(1,9)$, are recorded in
\cite[Chapter~1, \S1, pp.~5--6; \S2, p.~8; \S3, p.~9]{Peters1990}. A properly
elliptic component with $\chi(\OO)=1$
has $q=p_g=0$ and $e=12$ by \ref{cdp:st:C.2.UN.33},
and signature $(1,9)$ as well:
$b_1=0$, $b_2=10$, and the signature
formula with $K^2=0$ gives signature
$-8$. Applying the negative-index
bound to $19+n$ or $9+n$ gives
the asserted boundary ranks.

\item It remains to compute the exact
monodromy bound. \ref{cdp:st:C.2.UN.11} proves
that $N_0=U_1-I$ has square zero
and rank at most two. In rank one,
choose a vector and its nonzero
image, then complete the basis in
the kernel. Changing the sign of
one unused basis vector makes the
basis positively oriented, giving
the second row of \eqref{cdp:eq:C.CL.6}.
In rank two, choose two image
vectors and their preimages. If
their resulting ordered basis is
negative, reverse one preimage.
This gives the two oriented forms
\[
          N_0e_2=e_1,\qquad N_0e_4=\epsilon e_3,
                     \qquad\epsilon\in\{1,-1\}.     \tag{C.CL.10}\hypertarget{tag-00BX}{}\label{cdp:eq:C.CL.10}
\]
Write $e_{ij}=e_i\wedge e_j$.
In the oriented wedge basis
$(e_{12},e_{13},e_{14},e_{23},e_{24},e_{34})$
the only nonzero complementary
pairings are
\[
 e_{12}\cdot e_{34}=1,\qquad
 e_{13}\cdot e_{24}=-1,\qquad
 e_{14}\cdot e_{23}=1.                              \tag{C.CL.11}\hypertarget{tag-00BY}{}\label{cdp:eq:C.CL.11}
\]
For rank one the invariant space
has basis $(e_{12},e_{13},e_{14},e_{34})$;
its form is one hyperbolic plane
and a two-dimensional radical.
For \eqref{cdp:eq:C.CL.10} an invariant basis is
\[
       e_{12},\ e_{34},\ e_{13},\ e_{14}-\epsilon e_{23}.
                                                               \tag{C.CL.12}\hypertarget{tag-00BZ}{}\label{cdp:eq:C.CL.12}
\]
Its Gram matrix is
\[
 \begin{pmatrix}
 0&1&0&0\\
 1&0&0&0\\
 0&0&0&0\\
 0&0&0&-2\epsilon
 \end{pmatrix}.                                               \tag{C.CL.13}\hypertarget{tag-00C0}{}\label{cdp:eq:C.CL.13}
\]
These are exactly the rows of
\eqref{cdp:eq:C.CL.6}; the identity row is
the whole form $3H$.
The four exterior-square kernels and their restricted forms
are computed over the rationals by
\path{code/current_audits/audit_sign/code/independent_sign_checks.py};
the bases and Gram matrices are recorded in
\path{code/current_audits/audit_sign/independent_sign_checks.json}.
For the rank-two rows the recorded basis is
$(e_{12},e_{13},-\epsilon e_{14}+e_{23},e_{34})$.
Reordering its columns as $(1,4,2,3)$ and multiplying the last
column by $-\epsilon$ gives exactly the basis in
\eqref{cdp:eq:C.CL.12}, and congruence gives
\eqref{cdp:eq:C.CL.13}. The four inertias are respectively
$(3,3,0)$, $(1,1,2)$, $(1,2,1)$ and $(2,1,1)$. Finally
the rank-one row has $p_U=1$,
so \eqref{cdp:eq:C.CL.5} forces every
$p_g(Z_i)$ to vanish, as claimed.

\end{enumerate}
\end{proof}






\begin{proposition}[Every curve of a smooth minimal algebraic-dimension-zero Hopf normalization lies in the fibre boundary]\hypertarget{tag-01DH}{}\label{cdp:st:C.2.UN.38}
\StatementRecord{cdp:st:C.2.UN.38}{proved}{01DH}{\SeriesRef{cdp:eq:C.HB.2}, \SeriesRef{cdp:eq:C.HB.3}, \SeriesRef{cdp:eq:C.HB.5}, \SeriesRef{cdp:eq:C.HB.4}; \cite[Theorem~2.1]{Khaled2024Hopf}; \cite[\S1.1]{Nakamura2008VII}}

Let $f:U\to\Delta$ be a proper degeneration of
connected complex two-tori from a smooth threefold,
and write its full central fibre as $\sum_i m_iS_i$.
On the open subset of $S_i$ where the reduced
support is smooth, its canonical line has an
actual trivial $m_i$-th power in a neighborhood
of every compact subset. In particular, a
smooth compact elliptic curve $C$ in that open
subset satisfies
\[
                         N_{C/S_i}^{\otimes m_i}
                                      \simeq\OO_C.    \tag{C.HB.1}\hypertarget{tag-00LF}{}\label{cdp:eq:C.HB.1}
\]

Suppose the normalization of a central component
is itself a smooth minimal Hopf surface of
algebraic dimension zero. Then every elliptic
curve on that normalization is contained in
the inverse image of the singular locus of
the reduced support. For a diagonal primary
Hopf surface this includes both axes. For a
non-diagonal primary Hopf surface it includes
its unique axis; in this case the conclusion
uses the full canonical neighborhood and
the explicit higher-jet calculation below.
The same conclusion descends through every
actual finite primary Hopf covering.
\end{proposition}



\begin{proof}

\begin{enumerate}[label={$\langle1\rangle$\arabic*.},leftmargin=*]
\item The proper direct image $f_*K_U$ is a
torsion-free coherent sheaf of generic rank
one on the disc and hence is a line bundle.
Choose a generator after shrinking the disc.
Its evaluation is a holomorphic section
$\omega$ of $K_U$. On each nonzero fibre,
proper base change identifies this evaluation
with a nonzero canonical form on that torus,
up to the fixed base differential. It has
no zero there. Consequently its exact divisor
is vertical:
\[
                   \operatorname{div}(\omega)
                               =\sum_i k_iS_i,
                        \qquad k_i\ge0.              \tag{C.HB.2}\hypertarget{tag-00LG}{}\label{cdp:eq:C.HB.2}
\]
Let $W\subset U$ be a neighborhood of a
compact subset in the indicated regular
part of $S=S_i$, small enough to avoid all
other components and the singularities of
$S$. On $W$ the actual fibre equation has
divisor $mS$, where $m=m_i$. Thus it gives
the trivialization of $\OO_W(mS)$, and
\eqref{cdp:eq:C.HB.2} gives
\[
 \begin{gathered}
 K_U|_W\simeq\OO_W(k_iS),\\
 K_S|_{S\cap W}\simeq
                  \OO_W((k_i+1)S)|_{S\cap W},\\
 (K_S|_{S\cap W})^{\otimes m}\simeq\OO_{S\cap W}.
 \end{gathered}
                                                               \tag{C.HB.3}\hypertarget{tag-00LH}{}\label{cdp:eq:C.HB.3}
\]
These are line-bundle isomorphisms on the
neighborhood, obtained from the actual
canonical and fibre sections. Adjunction
for an elliptic $C\subset S\cap W$ gives
$N_{C/S}=K_C\otimes K_S^{-1}|_C$.
Since $K_C=\OO_C$, \eqref{cdp:eq:C.HB.1} follows.

\item Consider first the diagonal primary
Hopf surface
\[
 H_{\alpha,\beta}=(\C^2\setminus\{0\})/
       \langle(z,w)\mapsto(\alpha z,\beta w)\rangle,
                    \qquad0<|\alpha|,|\beta|<1.
                                                               \tag{C.HB.4}\hypertarget{tag-00LI}{}\label{cdp:eq:C.HB.4}
\]
The axis $C_\alpha$ is $\C^*/\alpha^{\Z}$,
and its actual normal line is the quotient
of $\C^*\times\C$ by
$(z,v)\mapsto(\alpha z,\beta v)$.
If a positive tensor power $n$ of this
line is trivial, a nowhere-zero holomorphic
function $a$ on $\C^*$ satisfies
$a(\alpha z)=\beta^n a(z)$.
In its Laurent expansion each nonzero
coefficient of $z^j$ requires
$\alpha^j=\beta^n$. Since $|\alpha|<1$,
there is at most one such exponent, and
there must be one. Its exponent is positive
because $|\beta|<1$ and $n>0$.
Conversely that equality gives the
explicit trivialization $a(z)=z^j$.
It also gives the nonconstant invariant
meromorphic function $w^n/z^j$ on the
Hopf surface. Algebraic dimension zero
therefore makes the normal of each axis
non-torsion. Formula \eqref{cdp:eq:C.HB.1} excludes
an entire such axis from the regular
part of the reduced support.

\item The non-diagonal primary normal form
requires a further argument. The primary
normal forms and finite unramified cover
used here are the explicit ones recalled
in \cite[Theorem~2.1]{Khaled2024Hopf}
and \cite[\S1.1]{Nakamura2008VII}.
Write the non-diagonal contraction
as
\[
 H=(\C^2\setminus\{0\})/\langle F\rangle,
 \qquad F(z,w)=(\beta^r z+w^r,\beta w),
                  \quad r\ge1,\quad0<|\beta|<1,
                                                               \tag{C.HB.5}\hypertarget{tag-00LJ}{}\label{cdp:eq:C.HB.5}
\]
and let $C$ be the image of $w=0$.
The invariant meromorphic two-form
$dz\wedge dw/w^{r+1}$ gives the actual
canonical divisor
\[
                         K_H\simeq\OO_H(-(r+1)C).
                                                               \tag{C.HB.6}\hypertarget{tag-00LK}{}\label{cdp:eq:C.HB.6}
\]
Its normal line has torsion, so \eqref{cdp:eq:C.HB.1}
alone gives no contradiction. Instead,
\eqref{cdp:eq:C.HB.3} on a neighborhood of $C$ would
make $\OO(m(r+1)C)$ trivial there.
Its canonical divisor section would
then be a holomorphic function with
positive divisor $m(r+1)C$ on that
neighborhood. We show directly that
no such function exists.

\item Lift it to the invariant neighborhood
of $\C^*\times\{0\}$ in the covering
space and expand in the normal variable:
\[
 h(z,w)=\sum_{l\ge j}h_l(z)w^l,
       \qquad j>0,\qquad h_j\ne0,
       \qquad h_l\in\OO(\C^*),
       \qquad h\circ F=h.                           \tag{C.HB.7}\hypertarget{tag-00LL}{}\label{cdp:eq:C.HB.7}
\]
The radius of convergence may depend
on $z$; the Taylor coefficients are
nevertheless holomorphic on all of
$\C^*$ because the lifted neighborhood
contains that entire axis. The leading
coefficient satisfies
$\beta^j h_j(\beta^r z)=h_j(z)$.
Its Laurent expansion forces
\[
                      j=rk,\qquad
                      h_j(z)=c z^{-k},
                      \qquad k\ge1,\quad c\ne0.
                                                               \tag{C.HB.8}\hypertarget{tag-00LM}{}\label{cdp:eq:C.HB.8}
\]
At order $j+r$, the same invariant
function equation is exactly
\[
 \beta^{j+r}h_{j+r}(\beta^r z)-h_{j+r}(z)
                    =k c\beta^{-r}z^{-k-1}.          \tag{C.HB.9}\hypertarget{tag-00LN}{}\label{cdp:eq:C.HB.9}
\]
Indeed the only extra contribution
at that order is the first derivative
of $h_j(\beta^r z+w^r)\beta^j$;
all other Taylor terms start at
higher order. In the Laurent
expansion of the left side, the
coefficient multiplying $z^{-k-1}$
is zero, since
$\beta^{j+r+r(-k-1)}-1=0$.
The right side has the nonzero
coefficient $k c\beta^{-r}$.
This contradiction proves the
neighborhood obstruction for
the non-diagonal axis as well.

\item Finally let $\nu:H\to S$ be the
given smooth minimal Hopf
normalization. Its conductor
and the inverse images of
the other central components
have curve support. Outside
this support $\nu$ is the
actual isomorphism onto the
regular locus of the reduced
fibre. For clarity, there
is no extra isolated
normalization defect: a
hypersurface is Cohen--Macaulay,
so its nonnormal locus is
the conductor curve locus;
where it is normal, the
normalization is an isomorphism,
and smoothness of $H$ then
gives smoothness there.
An intersection with another
central component is a
Cartier curve divisor and
its finite pullback also
has curve support.

\item We also verify the curve assertion
used in applying the obstruction.
The inverse image of a compact
curve in the covering punctured
plane extends analytically across
the origin by the removable-set
theorem for analytic sets, since
its dimension is one and the
added set has dimension zero.
Its germ has a reduced equation
$a$ satisfying $a\circ F=v a$
for a holomorphic unit $v$.
After shrinking to a contracting
ball, the product solving the
unit cocycle converges: if
$u=v/v(0)$, then
$b=\prod_{j\ge0}u\circ F^j$
is a holomorphic unit and
$b\circ F=b/u$.
Thus $h=ba$ satisfies
$h\circ F=v(0)h$.
For the diagonal contraction,
two distinct nonzero monomials
in $h$ would give
$\alpha^i\beta^j=1$ for a
nonzero integer pair; because
both moduli are less than one,
this is a positive-power
relation between $\alpha$
and $\beta$ and gives a
nonconstant meromorphic
invariant. In algebraic
dimension zero there is
therefore just one monomial,
so the curve consists of
axes. For \eqref{cdp:eq:C.HB.5}, decompose
the convergent power series
into weighted degrees with
weights $r,1$ on $z,w$.
The contraction preserves
each finite-dimensional
weighted piece and acts
there with sole eigenvalue
$\beta^d$ in degree $d$.
Their distinct moduli force
an eigenfunction into one
degree. Write it as
$w^dP(z/w^r)$; after taking
out $\beta^d$, invariance
requires the polynomial
$P$ to be periodic under
translation by $\beta^{-r}$.
It is constant. Thus only
the axis $w=0$ occurs.
Analytic continuation along
the contraction orbits
gives this assertion on
the entire punctured plane.
Hence the only curves
are the two disjoint axes
in \eqref{cdp:eq:C.HB.4}, or the single
axis in \eqref{cdp:eq:C.HB.5}.
Therefore an axis not
contained in that boundary
is disjoint from it. The
preceding neighborhood
contradictions apply and
force every axis into the
boundary. For a secondary
Hopf surface, pull back
the neighborhood and its
canonical trivialization
along its finite unramified
primary covering. The
covering has the same
algebraic dimension:
meromorphic functions
descend through their
elementary symmetric
polynomials. A curve
outside the boundary
would lift to an axis
outside the pulled-back
boundary, contradicting
the already proved
primary calculation.
This proves the stated
descent and the theorem.

\end{enumerate}
\end{proof}






\begin{proposition}[Curve genera, triple points and the finite monodromy constant decide zero Euler number]\hypertarget{tag-01DI}{}\label{cdp:st:C.2.UN.39}
\StatementRecord{cdp:st:C.2.UN.39}{proved}{01DI}{\SeriesRef{cdp:st:C.2.UN.11}, \SeriesRef{cdp:eq:C.ZC.5}, \SeriesRef{cdp:eq:C.ZC.6}, \SeriesRef{cdp:eq:C.ZC.8}, \SeriesRef{cdp:eq:C.ZC.7}, \SeriesRef{cdp:eq:C.ZC.9}, \SeriesRef{cdp:st:C.2.UN.4}; \cite[\S2, pp.~237--243; \S3, pp.~243--245]{ACampo1975}}

Let $f:U\to\Delta$ be an actual proper holomorphic
two-torus degeneration with smooth total threefold,
and let $S=f^{-1}(0)_{\rm red}$. Choose an embedded
resolution $\pi:\widetilde U\to U$ as a sequence
of blowups with smooth centres contained in the
successive central supports. Write its full
central divisor as $\sum_jN_jE_j$ with simple
normal crossings. Let $\mathcal D$ be the set
of connected double curves of $\bigcup_jE_j$,
and let $t_3$ be its number of triple points.
Count every point centre at every stage of
$\pi$, obtaining an integer $p_\pi$, and list
every connected smooth curve centre as $C_\alpha$.
Define the integer
\[
 I_\pi=
     \sum_{D\in\mathcal D}\bigl(g(D)-1\bigr)+t_3
       +p_\pi+\sum_\alpha\bigl(1-g(C_\alpha)\bigr).
                                                               \tag{C.ZC.1}\hypertarget{tag-018W}{}\label{cdp:eq:C.ZC.1}
\]
It is independent of the chosen embedded
resolution, including additional vertical
blowups used to obtain another such resolution.

Let $T$ be the actual monodromy on $H^1(F,\Z)$
of a nearby torus, and let $M$ be the order of
its semisimple part. Put
\[
             j(T)=\frac1M\sum_{k=1}^{M}\det(I-T^k).
                                                               \tag{C.ZC.2}\hypertarget{tag-018X}{}\label{cdp:eq:C.ZC.2}
\]
Then
\[
                e(S)=j(T)-2I_\pi,\qquad
                e(S)=0\quad\Longleftrightarrow\quad
                         I_\pi=\frac{j(T)}2 .        \tag{C.ZC.3}\hypertarget{tag-018Y}{}\label{cdp:eq:C.ZC.3}
\]
The constant $j(T)$ belongs to $\{0,4,6,8\}$.
For all nineteen monodromy polynomials its
value and the required zero index are
\[
 \begin{array}{c|r|r|r}
\det(xI-T)&M&j(T)&I_\pi\text{ for }e(S)=0\\ \hline
\Phi_{1}^{4}&1&0&0\\
\Phi_{1}^{2} \Phi_{2}^{2}&2&0&0\\
\Phi_{1}^{2} \Phi_{3}&3&0&0\\
\Phi_{1}^{2} \Phi_{4}&4&0&0\\
\Phi_{1}^{2} \Phi_{6}&6&0&0\\
\Phi_{2}^{4}&2&8&4\\
\Phi_{3}^{2}&3&6&3\\
\Phi_{2}^{2} \Phi_{4}&4&4&2\\
\Phi_{4}^{2}&4&6&3\\
\Phi_{5}&5&4&2\\
\Phi_{2}^{2} \Phi_{3}&6&4&2\\
\Phi_{2}^{2} \Phi_{6}&6&4&2\\
\Phi_{3} \Phi_{6}&6&4&2\\
\Phi_{6}^{2}&6&6&3\\
\Phi_{8}&8&4&2\\
\Phi_{10}&10&4&2\\
\Phi_{3} \Phi_{4}&12&4&2\\
\Phi_{4} \Phi_{6}&12&4&2\\
\Phi_{12}&12&4&2\\
\end{array}
                                                               \tag{C.ZC.4}\hypertarget{tag-018Z}{}\label{cdp:eq:C.ZC.4}
\]
Thus this is a necessary and sufficient
geometric test for the zero Euler number of
every actual fibre, using only the indicated
curve genera, triple points, resolution centres
and finite monodromy constant. Arbitrary
component and centre multiplicities are
included through the full divisor and the
actual blowup maps. The statement classifies
this numerical property of actual fibres;
it makes no occurrence assertion for abstract
genus or monodromy data and does not enumerate
analytic isomorphism classes of supports.
\end{proposition}



\begin{proof}

\begin{enumerate}[label={$\langle1\rangle$\arabic*.},leftmargin=*]
\item An embedded resolution of the specified kind
exists by the analytic resolution theorem used
in \ref{cdp:st:C.2.UN.11}, with all centres over the central
support and with an isomorphism over
$\Delta^*$. Smooth divisor centres need
not be included, since their blowups are
isomorphisms. The remaining connected centres
in a threefold are compact smooth curves or
points. A disconnected centre can be counted
by its disjoint connected components; their
blowups commute. In particular every genus
and count in \eqref{cdp:eq:C.ZC.1} is attached to an actual
compact centre or curve in this resolution.

\item We first compute the resolved support without
any reduction of its multiplicities. Put
\[
             E_j^\circ=E_j\setminus\bigcup_{i\ne j}E_i.
\]
The proper analytic A'Campo comparison
\cite[\S2, pp.~237--243; \S3, pp.~243--245]{ACampo1975}, with
its actual nearby-cycle maps as established
in (C.MS.6), says for every positive $k$ that
\[
         \det(I-T^k)=\sum_{N_j\mid k}N_j e(E_j^\circ).
                                                               \tag{C.ZC.5}\hypertarget{tag-0190}{}\label{cdp:eq:C.ZC.5}
\]
Here the full integer $N_j$ is used for each
component. Choose $L$ divisible by $M$ and
by every $N_j$. The determinant on the left
depends only on the eigenvalues of $T$,
so is periodic with period $M$ even when
its unipotent part is nontrivial. Average
\eqref{cdp:eq:C.ZC.5} over $1\le k\le L$. Exactly $L/N_j$
of these integers are divisible by $N_j$.
It follows that
\[
 \begin{split}
 j(T)
 &=\frac1L\sum_{k=1}^{L}\det(I-T^k)\\
 &=\sum_j\frac{N_j}{L}\frac{L}{N_j}e(E_j^\circ)
   =\sum_j e(E_j^\circ).
 \end{split}                                                   \tag{C.ZC.6}\hypertarget{tag-0191}{}\label{cdp:eq:C.ZC.6}
\]
This calculation includes components whose
multiplicity does not divide $M$; choosing
$L$ retains them before the averaging.

\item Let $\widetilde S=\bigcup_jE_j$. Its
remaining strata are the double curves
with their triple points removed, and
the triple points themselves. Every
triple point lies on exactly three
double curves. Since each connected
smooth compact double curve $D$ has
Euler number $2-2g(D)$, additivity gives
\[
 \begin{split}
 e(\widetilde S)
 &=\sum_j e(E_j^\circ)
       +\sum_{D\in\mathcal D}(2-2g(D))-3t_3+t_3\\
 &=j(T)-2\left(
           \sum_{D\in\mathcal D}(g(D)-1)+t_3\right).
 \end{split}                                                   \tag{C.ZC.7}\hypertarget{tag-0192}{}\label{cdp:eq:C.ZC.7}
\]
The surface Euler numbers have disappeared:
the open surface contribution is the finite
monodromy average computed in \eqref{cdp:eq:C.ZC.6}.

\item We now compare the original support and
the resolved one through each actual
blowup. Let $\sigma:U'\to U_0$ blow up
a connected smooth centre $Z$ contained
in a current reduced central support
$S_0$. Its reduced inverse image is
the next central support $S'$. Away
from $Z$ the map is an isomorphism;
over every point of $Z$ its fibre is
$\mathbb P^{c-1}$, where $c$ is the
complex codimension of $Z$ in $U_0$.
The entire projective fibre belongs
to $S'$, because the centre is contained
in $S_0$. Therefore the actual proper
map, stratified by $Z$ and its complement,
gives
\[
                    e(S')-e(S_0)=(c-1)e(Z).
                                                               \tag{C.ZC.8}\hypertarget{tag-0193}{}\label{cdp:eq:C.ZC.8}
\]
For a point this is $2$; for a smooth
connected curve it is $2-2g(Z)$.
This computation is valid for every
multiplicity of the full fibre along
the centre. More explicitly, if the
generic order of its local equation
along the centre is $\mu$, its pullback
has the exceptional factor of order
$\mu$. That factor changes the full
divisor coefficient but not the
projective fibre in its reduced
inverse image. No value of $\mu$
was assumed in \eqref{cdp:eq:C.ZC.8}.

\item Iterating the actual proper maps in
$\pi$ gives
\[
 e(\widetilde S)=e(S)+2p_\pi
                        +2\sum_\alpha(1-g(C_\alpha)).
                                                               \tag{C.ZC.9}\hypertarget{tag-0194}{}\label{cdp:eq:C.ZC.9}
\]
Combining \eqref{cdp:eq:C.ZC.7} and \eqref{cdp:eq:C.ZC.9}
proves \eqref{cdp:eq:C.ZC.3} with exactly the
integer in \eqref{cdp:eq:C.ZC.1}. Every such
resolution is an isomorphism over
the punctured disc and therefore
has the same actual monodromy.
The formula consequently identifies
$I_\pi$ with $(j(T)-e(S))/2$ for
each choice. This proves independence
by the actual blowup comparisons;
no comparison between unrelated
fillings is required.

\item The nineteen possible characteristic
polynomials and their one-period
determinants were computed in \ref{cdp:st:C.2.UN.11}.
Applying the finite average \eqref{cdp:eq:C.ZC.2}
gives the entire table \eqref{cdp:eq:C.ZC.4}.
Equivalently, $j=0$ when the polynomial
has a factor $\Phi_1$; $j=8$ for
$\Phi_2^4$; $j=6$ for
$\Phi_3^2,\Phi_4^2,\Phi_6^2$;
and $j=4$ for the ten other rows.
The exact integer-matrix calculation in
\path{code/current_audits/audit_sign/code/UN24_UN39_exact_certificate_08.py}
computes every one-period determinant, the M\"obius coefficients,
all nineteen averages and their zero indices. Its complete output is
\path{code/current_audits/audit_sign/UN24_UN39_exact_certificate_08.json}.
For each characteristic polynomial the program uses a direct sum
of cyclotomic companion matrices. Over $\C$, triangularizing the actual
$T$ gives $\det(I-T^k)=\prod_\lambda(1-\lambda^k)$, with every eigenvalue
multiplicity retained; this identity identifies the computed determinant
with the required one even in the presence of a unipotent part.
The average agrees with the sum of the A'Campo coefficients in every row.
All arithmetic is integral or rational.

\item For a concrete check of the geometric
terms, start from a smooth torus support
with identity monodromy. A point blowup
creates one rational double curve and
no triple point. The new support has
index $-1$ and Euler number $2$.
Follow it by the genus-three quartic
blowup of \ref{cdp:st:C.2.UN.4}. There are then five
rational double curves, one double
curve of genus three, and four triple
points. Their index is
$-5+(3-1)+4=1$, giving Euler number
$-2$. Viewed instead as a resolution
of the original smooth support, these
two centres contribute
$1+(1-3)=-1$ to the centre term,
so \eqref{cdp:eq:C.ZC.1} returns index zero for
that original support. These exact
charts verify both the sign of the
centre correction and its compatibility
with the universal balancing operations.

\end{enumerate}
\end{proof}






\begin{proposition}[The fundamental group and the component-incidence constraint of every torus fibre]\hypertarget{tag-01DO}{}\label{cdp:st:C.2.UN.44}
\StatementRecord{cdp:st:C.2.UN.44}{proved}{01DO}{\SeriesRef{cdp:st:C.2.UN.39}, \SeriesRef{cdp:eq:C.FG.9}, \SeriesRef{cdp:eq:C.FG.10}, \SeriesRef{cdp:eq:C.FG.14}, \SeriesRef{cdp:eq:C.FG.15}, \SeriesRef{cdp:eq:C.FG.12}, \SeriesRef{cdp:eq:C.FG.17}, \SeriesRef{cdp:eq:C.FG.18}; \cite[Theorem (LI), p.~775]{Kurdyka1998}; \cite[\S3, Theorem~2]{Lojasiewicz1964}}

Let $f:U\to\Delta$ be a proper holomorphic degeneration of connected
complex two-tori, with smooth total threefold. Write
\[
 f^{-1}(0)=\sum_{i=1}^r m_iS_i,\qquad
 S=\bigcup_iS_i,\qquad d=\gcd(m_1,\ldots,m_r).
                                                        \tag{C.FG.1}\hypertarget{tag-00K8}{}\label{cdp:eq:C.FG.1}
\]
All fundamental groups below are based; changing the paths to the
chosen base point changes the indicated presentations by their actual
change-of-base-point isomorphisms. Choose a lift $t$ of the positive
generator of $\pi_1(\Delta^*)$, put $\Lambda=\pi_1(F)=\Z^4$, and
write $T$ for conjugation by $t$ on $\Lambda$. There are vectors
$v_i\in\Lambda$, obtained from positively oriented meridians of the
actual components, for which
\[
 \Gamma:=\pi_1(S)
  \simeq
  \bigl(\Lambda\rtimes_T\langle t\rangle\bigr)
       \big/\langle\!\langle x^{v_i}t^{m_i}:1\le i\le r
                    \rangle\!\rangle .                \tag{C.FG.2}\hypertarget{tag-00KJ}{}\label{cdp:eq:C.FG.2}
\]
Here $x^v$ denotes the element $v$ of the additive lattice written
multiplicatively. Its image $A$ in $\Gamma$ is abelian, and there is
an exact sequence
\[
       1\longrightarrow A\longrightarrow\Gamma
          \longrightarrow\Z/d\Z\longrightarrow1 .    \tag{C.FG.3}\hypertarget{tag-00KK}{}\label{cdp:eq:C.FG.3}
\]
In particular the index is exactly $d$ and $A$ is generated by at most
four elements. There is an exact integral description as follows.
Set $n_i=m_i/d$, choose integers $c_i$ with $\sum_i c_in_i=1$, and
put $v_0=\sum_i c_iv_i$. Then
\[
 A\simeq\Lambda/R_d,\qquad
 R_d=(T^d-I)\Lambda+
       \left\langle T^jv_i-n_iv_0:
                         1\le i\le r,\ 0\le j<d\right\rangle_{\Z}.
                                                        \tag{C.FG.4}\hypertarget{tag-00KL}{}\label{cdp:eq:C.FG.4}
\]
For $d=1$, the whole group is abelian and this specializes to
\[
 \pi_1(S)\simeq
   \Lambda\big/\left((T-I)\Lambda+
                   \langle v_i-m_iv_0\rangle_{\Z}\right).
                                                        \tag{C.FG.5}\hypertarget{tag-00KM}{}\label{cdp:eq:C.FG.5}
\]
The normalized base change $s^d=f$ is an actual connected degree-$d$
\emph{\'etale} cover of the total space near $S$. Its reduced central
support $S_d$ is connected, its restriction $S_d\to S$ is the
connected cyclic covering belonging to the subgroup $A$, and
$\pi_1(S_d)=A$ under that covering injection.

For every embedded resolution of the central divisor by smooth
vertical blowups, let $\mathcal D$ be the dual complex of its simple
normal-crossing support. There is a surjection
\[
                  \pi_1(S)\longrightarrow\pi_1(\mathcal D).
                                                        \tag{C.FG.6}\hypertarget{tag-00KN}{}\label{cdp:eq:C.FG.6}
\]
Its target has an abelian normal subgroup of index dividing $d$,
generated by at most four elements and of rank at most two.
Moreover the cup-product map
$\bigwedge^*H^1(\mathcal D,\Q)\to H^*(\mathcal D,\Q)$ is injective.
When the resolved support has
no triple points, its connected incidence graph is a tree or has
exactly one independent cycle. In particular two components of an
actual simple normal-crossing fibre with no triple points can meet
in at most two connected curves.

For an actual simple normal-crossing support with no triple points,
write $e_1$ for the number of its connected double curves. The zero
Euler criterion of \ref{cdp:st:C.2.UN.39} takes the following joint form:
\[
 \begin{gathered}
 \#\{
      \text{independent incidence cycles}\}\le1,\\
 e(S)=0\ \Longleftrightarrow\
       \sum_{D\ {\rm double}}g(D)=e_1+\frac{j(T)}2 .
 \end{gathered}
                                                        \tag{C.FG.7}\hypertarget{tag-00KO}{}\label{cdp:eq:C.FG.7}
\]
Here $j(T)$ is unchanged by replacing the action on $H^1(F)$ by
the corresponding action on $H_1(F)$ or its inverse. If all $m_i$
are equal, $j(T)=0$. If $1$ is not an eigenvalue of $T$, a
no-triple incidence graph must be a tree.
These are constraints and exact group formulas
for actual degenerations; they include arbitrary multiplicities
and do not assert the existence of a smoothing of abstract gluing data.
\end{proposition}



\begin{proof}\leavevmode\par



 
\begin{enumerate}[label={$\langle1\rangle$\arabic*.},leftmargin=0pt,labelsep=.5em]
\item \par\medskip\noindent\textbf{The actual small-tube retraction.}
Shrink the disc so that zero is its only critical value, and take
a compact closed tube $U_\epsilon=f^{-1}(|z|\le\epsilon)$ inside
the original open total space. Put $\varphi=|f|^2$ and choose a
smooth Riemannian metric near this tube. The local analytic gradient
inequality, in the form of
\cite[Theorem (LI), p.~775]{Kurdyka1998}, gives near each point of
$S$ an exponent smaller than one. Compactness, a finite coordinate
cover, and uniform equivalence of the coordinate Euclidean metrics
with the chosen metric give constants $c>0$ and $0<\theta<1$ with
\[
                 \|\nabla\varphi\|\ge c\varphi^\theta
                   \quad\hbox{on }U_\epsilon\setminus S.
                                                        \tag{C.FG.8}\hypertarget{tag-00KP}{}\label{cdp:eq:C.FG.8}
\]
Indeed take the maximum of the finitely many exponents, then shrink
to $0\le\varphi<1$; properness permits a whole smaller tube to be
contained in the resulting neighbourhood of $S$. There is no
critical point of $\varphi$ off $S$, since $df$ is surjective there.


\begin{enumerate}[label={$\langle2\rangle$\arabic*.},leftmargin=1.3em,labelsep=.5em]
\item  On that complement use the vector field
$V=-\nabla\varphi/\|\nabla\varphi\|^2$. Along its flow the value
of $\varphi$ decreases at speed one. A trajectory starting at level
$a$ therefore exists up to every time smaller than $a$: each
positive closed level band is compact and the vector field is
smooth there. Its remaining length below level $b$ is at most
\[
          \int_0^b\frac{du}{c u^\theta}
                 =\frac{b^{1-\theta}}{c(1-\theta)}.    \tag{C.FG.9}\hypertarget{tag-00KQ}{}\label{cdp:eq:C.FG.9}
\]
It consequently has a limit in $S$. Let $H(x,q)$ be the flow for
time $q\varphi(x)$ when $q<1$, and its limiting point when $q=1$;
put $H(x,q)=x$ on $S$. The estimate \eqref{cdp:eq:C.FG.9} proves continuity
at points of $S$ and uniform convergence as $q\to1$. For $q<1$
away from $S$, ordinary smooth dependence of flows proves
continuity. Thus $H$ is an actual strong deformation retraction
of $U_\epsilon$ onto $S$. This also proves the needed retraction
for every later resolved or covered tube. No invariant-cycle
surjectivity is involved.


\end{enumerate}

\item \par\medskip\noindent\textbf{The exact normal subgroup of meridians.}
Write $U^*=U_\epsilon\setminus S$. Proper horizontal transport
over the punctured disc gives a smooth torus bundle. Its homotopy
sequence, and $\pi_2(\Delta^*)=0$, identify
\[
       \pi_1(U^*)=\Lambda\rtimes_T\langle t\rangle,
       \qquad \pi_1(U^*)\longrightarrow\pi_1(\Delta^*)=\Z
       \text{ sends }x^v t^a\text{ to }a .             \tag{C.FG.10}\hypertarget{tag-00K9}{}\label{cdp:eq:C.FG.10}
\]
The splitting exists because an extension onto the infinite cyclic
group splits after choosing a lift of its generator.


\begin{enumerate}[label={$\langle2\rangle$\arabic*.},leftmargin=1.3em,labelsep=.5em]
\item  General position makes every loop in the smooth six-manifold
$U_\epsilon$ avoid the divisor $S$. A null-homotopy of a loop from
$U^*$ can be made transverse to the smooth part of $S$ and disjoint
from its singular set: the latter has real dimension at most two,
so a transverse real two-disc misses it. Removing little discs
around the finitely many transverse intersections expresses the
original loop as a product of conjugates of positive or negative
meridians. Conversely each such meridian bounds its normal disc
in $U_\epsilon$.


\item  One meridian for each global irreducible component suffices. The
smooth locus of an irreducible component, with its intersections
with other components removed, is path connected. This follows
also by taking its connected normalization resolution and moving
paths off the proper complex analytic curves and points. A path
there moves positive normal circles through positive normal circles,
giving conjugate based meridians in $U^*$. Near a chosen smooth
point of $S_i$ the actual equation is $f=u z^{m_i}$ with $u$ a
unit. The meridian has winding $m_i$, hence is exactly
$x^{v_i}t^{m_i}$ in \eqref{cdp:eq:C.FG.10}, for a uniquely determined $v_i$
after the indicated choices. The retraction and this kernel
calculation prove \eqref{cdp:eq:C.FG.2}.


\item  The image $A$ of $\Lambda$ is normal and abelian. Quotienting
\eqref{cdp:eq:C.FG.2} by this image leaves precisely the presentation
$\langle t\mid t^{m_i}=1\text{ for all }i\rangle$.
This is $\Z/d\Z$, which proves the exact index in \eqref{cdp:eq:C.FG.3}.
For later comparison, abelianizing the same actual presentation
also gives
\[
 H_1(S,\Z)\simeq
 \frac{\bigl(\Lambda/(T-I)\Lambda\bigr)\oplus\Z}
      {\langle ([v_i],m_i):1\le i\le r\rangle_{\Z}}.
                                                        \tag{C.FG.11}\hypertarget{tag-00KA}{}\label{cdp:eq:C.FG.11}
\]


\end{enumerate}

\item \par\medskip\noindent\textbf{The connected normalized root cover.}
Work momentarily in the open complex neighbourhood of the closed
tube. Since all divisor multiplicities of $f$ are divisible by
$d$, unique factorization in the regular local analytic ring
gives $f=u h^d$ near every point, including singular points of the
support. The normalization of $s^d=u h^d$ is the explicit map
\[
       \{v^d=u\}\longrightarrow\{s^d=u h^d\},
                         \qquad s=vh .                \tag{C.FG.12}\hypertarget{tag-00KB}{}\label{cdp:eq:C.FG.12}
\]
Its source is finite and \'etale over the smooth original chart:
$v$ is a unit and $d v^{d-1}$ never vanishes. It is normal and is
isomorphic to the base change over $h\ne0$, hence is exactly that
base change's normalization. These normalization maps glue.


\begin{enumerate}[label={$\langle2\rangle$\arabic*.},leftmargin=1.3em,labelsep=.5em]
\item  Over $U^*$ this is the cyclic covering given by the winding number
modulo $d$ in \eqref{cdp:eq:C.FG.10}. That homomorphism is onto, so the cover
is connected. Its complement in the normalization is a divisor;
the full normalized cover is therefore connected as well.
The tube retraction $H$ lifts uniquely to this actual covering.
Its lift starts at the identity, fixes the inverse image of $S$
pointwise, and at time one has image in that inverse image. It is
therefore a strong deformation retraction onto $S_d$, proving
connectedness of $S_d$ and the asserted covering injection.
The covering subgroup in $\Gamma$ is the kernel of the winding
modulo $d$ map, which is exactly $A$ by \eqref{cdp:eq:C.FG.3}. In particular
this construction does not presume that the cover of an individual
$S_i$ is connected or that it has degree one.


\item  There is also a useful integral calculation using the punctured
part of this same cover. Its group is
$\Lambda\rtimes_{T^d}\langle w\rangle$, where $w=t^d$.
The kernel of its map to $\pi_1(S_d)$ is normally generated by
all the lifted meridians
\[
                x^{T^jv_i}w^{n_i},
                  \qquad 1\le i\le r,\quad0\le j<d .  \tag{C.FG.13}\hypertarget{tag-00KC}{}\label{cdp:eq:C.FG.13}
\]
To see completeness of this list, use the coset representatives
$1,t,\ldots,t^{d-1}$ for the winding cover in \eqref{cdp:eq:C.FG.10}.
Conjugating the original normal generators by these representatives,
and then by elements of the covering subgroup, gives exactly its
kernel. Thus no choice of the connected components over $S_i$
has omitted a meridian.


\item  The relations with $j=0$ imply $w^{n_i}=x^{-v_i}$.
Multiplying their $c_i$th powers gives $w=x^{-v_0}$. Consequently
the covering group is abelian, $(T^d-I)\Lambda$ is killed, and
the remaining relations are precisely
$T^jv_i-n_iv_0=0$. Conversely the map $x^v\mapsto[v]$,
$w\mapsto-[v_0]$ from this presentation to $\Lambda/R_d$
respects every semidirect and meridian relation. It is inverse to
the map induced by $\Lambda$. This proves \eqref{cdp:eq:C.FG.4}. For $d=1$
the same argument proves \eqref{cdp:eq:C.FG.5}; in particular it first proves
$t\in A$ and only then imposes $T=I$ on that quotient.
The Smith-normal-form comparisons in
\path{code/A_meridian_groups.py} check the two integral
presentations and their torsion in explicit test matrices; the
presentation is proved here for every actual monodromy.


\end{enumerate}

\item \par\medskip\noindent\textbf{Resolution and the actual dual-complex map.}
A smooth vertical point or curve blowup induces an isomorphism on
the fundamental group of the total tube. For a curve centre, cut
out a small normal disc bundle. Its boundary is an $S^3$-bundle
over the curve, and the exceptional neighbourhood retracts to the
$\PP^1$-bundle over that same curve. The boundary inclusion
induces the same isomorphism on fundamental groups as the old
disc bundle, since both fibre fundamental groups are zero. For
a point centre the boundary is $S^5$ and the exceptional fibre
is $\PP^2$, again both simply connected. Van Kampen with the
unchanged complement proves the assertion for the actual blowdown
map. Iteration, followed by the two tube retractions, gives
\[
              \pi_1(S^{\rm res})\xrightarrow{\ \simeq\ }\pi_1(S).
                                                        \tag{C.FG.14}\hypertarget{tag-00KD}{}\label{cdp:eq:C.FG.14}
\]
The old multiplicities are retained on strict transforms. Every
new multiplicity is the order of the old fibre equation along its
centre and is divisible by $d$, so their gcd remains $d$.


\begin{enumerate}[label={$\langle2\rangle$\arabic*.},leftmargin=1.3em,labelsep=.5em]
\item  For an SNC support $E=\bigcup E_i$, form the space $P$ from the
disjoint components $E_i$, one $D\times[0,1]$ for each connected
double curve, and one two-simplex for each triple point. The faces
are glued by the actual intersection inclusions. There are actual
continuous maps
\[
                   E\ \xleftarrow{\ p\ }\ P
                       \xrightarrow{\ q\ }\mathcal D,             
                                                        \tag{C.FG.15}\hypertarget{tag-00KE}{}\label{cdp:eq:C.FG.15}
\]
where $p$ forgets the simplex coordinate and $q$ collapses every
connected stratum to its associated cell. Compatible triangulation
of the finite analytic family of components and intersections makes
all the inclusions closed subcomplex inclusions; see
\cite[\S3, Theorem~2]{Lojasiewicz1964}. The map $p$ is a homotopy
equivalence. Explicitly, for two subcomplexes the double mapping
cylinder projects to their union by the homotopy-pushout comparison
for a cofibration. Inductively add one component, applying that
comparison also to its intersection with the previous union.
The intersection's mapping cylinders and the triple-point
triangles are precisely the pieces already included in $P$.
This proves the comparison for the finite union and supplies a
homotopy inverse of $p$.


\item  Van Kampen on these mapping cylinders gives a particularly precise
description of $q_*$. Start with the vertex fundamental groups and
edge paths, impose the relations induced by the two inclusions of
each $\pi_1(D)$, and then the triangle boundary relations.
Killing the vertex groups leaves exactly the presentation of
$\pi_1(\mathcal D)$. Thus
\[
 \pi_1(\mathcal D)\simeq
 \pi_1(E)\big/
       \langle\!\langle\operatorname{im}\pi_1(E_i):i
                              \rangle\!\rangle .       \tag{C.FG.16}\hypertarget{tag-00KF}{}\label{cdp:eq:C.FG.16}
\]
The map $q\circ p^{-1}$, where $p^{-1}$ is this homotopy inverse,
and \eqref{cdp:eq:C.FG.14} give \eqref{cdp:eq:C.FG.6}. In particular it is an actual
topological quotient map on fundamental groups, including when
two components have several connected double curves.


\item  The image of $A$ is an abelian normal subgroup whose quotient is
a quotient of $\Z/d\Z$. If there are no triple points,
$\mathcal D$ is a finite connected graph. Let its free-group rank
be $b$ and let the index of that abelian subgroup be $a$.
The corresponding connected graph covering has $a$ times as many
vertices and edges, hence free fundamental group of rank
$1+a(b-1)$. An abelian free group has rank at most one. Thus
$b\le1$, proving the tree-or-one-cycle assertion and its two-vertex
consequence.


\end{enumerate}

\item \par\medskip\noindent\textbf{The rank-two bound, with the actual finite-cover comparison.}
For any of these tubes, inclusion of a nearby fibre followed by
the constructed retraction is a continuous map
$\sigma:F\to S$. Restriction to its fundamental group has image
$A$ of finite index. Therefore
\[
 \sigma^*:H^1(S,\Q)\longrightarrow H^1(F,\Q)
           \quad\hbox{is injective}.                  \tag{C.FG.17}\hypertarget{tag-00KG}{}\label{cdp:eq:C.FG.17}
\]
Indeed a homomorphism $\Gamma\to\Q$ vanishing on $A$ factors through
the finite group $\Gamma/A$ and is zero. This proves the stated
direction of the cohomology map. The map $\sigma^*$ in all degrees
is a ring homomorphism because it comes from that actual continuous
map. Its degree-one image is $T$-invariant, as is also immediate
from the semidirect presentation.


\begin{enumerate}[label={$\langle2\rangle$\arabic*.},leftmargin=1.3em,labelsep=.5em]
\item  Apply this to the resolved tube $W$ with central support
$E=S^{\rm res}$. The map $q\circ p^{-1}:E\to\mathcal D$ of
\eqref{cdp:eq:C.FG.15} induces an injection on $H^1$ because its fundamental
group map is onto. The composite
\[
 H^*(\mathcal D,\Q)\longrightarrow H^*(E,\Q)
               \xrightarrow{\ \sigma^*\ }H^*(F,\Q)
                  =\bigwedge^*H^1(F,\Q)               \tag{C.FG.18}\hypertarget{tag-00KH}{}\label{cdp:eq:C.FG.18}
\]
is again a ring homomorphism and is injective in degree one.
The exterior algebra on any linear subspace of $H^1(F,\Q)$
injects into this exterior algebra. Consequently cup product
on $\mathcal D$ is injective on
$\bigwedge^*H^1(\mathcal D,\Q)$. Since $\mathcal D$ is a
two-dimensional cell complex, its third cohomology vanishes,
and three independent degree-one classes would give a
contradiction. Thus $b_1(\mathcal D)\le2$.


\item  This also bounds the rank of its finite-index abelian subgroup,
rather than only the rank of its abelianization. Put
$\Pi=\pi_1(\mathcal D)$ and let $B$ be the image of $A$ in $\Pi$.
The index $a=[\Pi:B]$ divides $d$. Take the actual connected
degree-$a$ covering $\mathcal D'\to\mathcal D$ belonging to $B$,
and pull it back along the continuous map
$W\to E\to\mathcal D$ just constructed. The pullback is connected
because $\pi_1(W)\to\Pi$ is onto. Its covering subgroup is the
inverse image of $B$ in $\Gamma$. The quotient $\Pi/B$ is cyclic
of order $a$, generated by the image of $t$, so this inverse
image is precisely the kernel of the winding number modulo $a$.
The uniqueness of a based covering with a given subgroup
identifies this pullback, over $W$, with the actual normalized
root cover
\[
                  s^a=f\circ\pi:W_a\longrightarrow W. \tag{C.FG.19}\hypertarget{tag-00KI}{}\label{cdp:eq:C.FG.19}
\]
It is \'{e}tale by \eqref{cdp:eq:C.FG.12}, since $a$ divides every resolved
component multiplicity. Its lifted tube retraction has connected
central support $E_a$. The lifted map
$E_a\to\mathcal D'$ is onto on fundamental groups: its image is
exactly $B$, by the inverse-image description of the covering
subgroup. Regarded over the $s$-disc, $W_a$ has connected nearby
fibre equal to the actual torus $F_{s^a}$. Hence \eqref{cdp:eq:C.FG.17}--\eqref{cdp:eq:C.FG.18}
apply to this map and give $b_1(\mathcal D')\le2$. Since
$\pi_1(\mathcal D')=B$ is finitely generated abelian, this is
exactly $\operatorname{rank}B\le2$, as claimed.


\item  If a no-triple incidence graph had one cycle, its nonzero
$H^1$ would inject by \eqref{cdp:eq:C.FG.18} into the $T$-invariant part of
$H^1(F,\Q)$. Thus $1$ must be an eigenvalue of $T$ in that case.


\item  Finally, for the original support already SNC with no triple
points, \ref{cdp:st:C.2.UN.39} has no resolution-centre terms and gives
$e(S)=j(T)-2\sum_D(g(D)-1)$, proving \eqref{cdp:eq:C.FG.7}.
If all multiplicities are $m$, the full trace formula is
$\det(I-T^k)=m\sum_i e(S_i^\circ)$ when $m\mid k$ and zero
otherwise. Taking $k$ also divisible by the semisimple order makes
$T^k$ unipotent, so this determinant is zero. Hence the common
sum is zero and all traces vanish; in particular $j(T)=0$.



\end{enumerate}

\end{enumerate}
\end{proof}






\begin{proposition}[Positive geometric genus controls the whole support group]\hypertarget{tag-01DR}{}\label{cdp:st:C.2.UN.47}
\StatementRecord{cdp:st:C.2.UN.47}{proved}{01DR}{\SeriesRef{cdp:st:C.2.UN.33}, \SeriesRef{cdp:st:C.2.UN.44}, \SeriesRef{cdp:eq:C.PF.5}, \SeriesRef{cdp:st:C.2.UN.28}, \SeriesRef{cdp:eq:C.PF.4}, \SeriesRef{cdp:eq:C.PF.7}, \SeriesRef{cdp:eq:C.PF.8}, \SeriesRef{cdp:eq:C.PF.6}, \SeriesRef{cdp:eq:C.PF.12}, \SeriesRef{cdp:eq:C.PF.10}}

Let $f:U\to\Delta$ be a proper holomorphic two-torus degeneration
with smooth connected total threefold, let $S=f^{-1}(0)_{\rm red}$,
and put $\Gamma=\pi_1(S)$. For an irreducible component $S_i$ let
$h_i:Z_i\to S_i\hookrightarrow U$ be a resolution of its normalization.
The following are simultaneous constraints on this actual fibre.

Every $Z_i$ with $p_g(Z_i)>0$ has $p_g(Z_i)=1$, and its actual
map induces a surjection
\[
                         \pi_1(Z_i)\longrightarrow\Gamma .
                                                        \tag{C.PF.1}\hypertarget{tag-00UN}{}\label{cdp:eq:C.PF.1}
\]
Consequently every SNC dual complex of an embedded resolution of
the fibre is simply connected whenever such a component occurs.
A component with K3 minimal model forces $\Gamma=1$.
More generally,
\[
                    \chi(\OO_{Z_i})=2
                         \quad\Longrightarrow\quad\Gamma=1 .
                                                        \tag{C.PF.2}\hypertarget{tag-00UR}{}\label{cdp:eq:C.PF.2}
\]
This includes the Kodaira-dimension-one components of holomorphic
Euler characteristic two in \ref{cdp:st:C.2.UN.33}.

For a component with $\chi(\OO_{Z_i})=1$, its fundamental-group
image in $\Gamma$ is zero, with just the following possible
group configuration:
\[
 \Gamma\simeq\Z/2\Z,\qquad
 \operatorname{im}\bigl(\pi_1(Z_i)\to\Gamma\bigr)=\Gamma .
                                                        \tag{C.PF.3}\hypertarget{tag-00US}{}\label{cdp:eq:C.PF.3}
\]
In that configuration the actual connected double cover of $Z_i$
induced from the tube has
$\chi(\OO)=2$, $p_g=1$ and $q=0$. At most one component with
$\chi(\OO)=1$ can have nontrivial fundamental-group image.
The assertion \eqref{cdp:eq:C.PF.3} is a restriction on groups of actual fibres,
and is not an assertion that every surface or gluing with these
invariants occurs.
\end{proposition}



\begin{proof}

\begin{enumerate}[label={$\langle1\rangle$\arabic*.},leftmargin=*]
\item \par\medskip\noindent\textbf{Every finite quotient gives another actual torus degeneration.}
Use the strong deformation retraction of \ref{cdp:st:C.2.UN.44} on a sufficiently
small tube; its annular complement also retracts radially by proper
horizontal transport. Thus an open tube and its closed smaller
tube both have group $\Gamma$. Given a finite quotient
$\rho:\Gamma\twoheadrightarrow Q$ of order $N$, its connected
regular topological covering carries the unique complex structure
obtained by lifting the holomorphic coordinate charts. This gives
an actual finite \'{e}tale map
\[
                     p:U_Q\longrightarrow U,\qquad \deg p=N.
                                                        \tag{C.PF.4}\hypertarget{tag-00UT}{}\label{cdp:eq:C.PF.4}
\]
It is proper: a finite covering of a locally compact Hausdorff
space is proper, as follows by covering a compact set with
finitely many relatively compact evenly covered neighbourhoods.
Consequently $f\circ p$ is proper. Its total space is smooth and
connected.

\item The punctured covering $U_Q^*$ is connected, since
$\pi_1(U^*)\to\Gamma\to Q$ is onto by \ref{cdp:st:C.2.UN.44}. Let $A$ be the image
of the nearby fibre group in $\Gamma$ and let $B=\rho(A)$.
The connected components over a fixed nearby fibre are indexed
by $Q/B$. The base generator permutes them transitively, and
$Q/B$ is cyclic, because $\Gamma/A$ is cyclic. If
$a=[Q:B]$, the connected base covering in the Stein factor of
$f\circ p$ is therefore the punctured degree-$a$ cover of the disc.
Its normal completion has one point above zero and is the disc
with actual coordinate map $t=s^a$. One can obtain this
identification by the classification of coverings of the
punctured disc and then extend the two inverse maps across zero:
each is bounded over a smaller disc, and the completed normal
one-dimensional spaces have the removable-singularity property.
Equivalently it is the normalization of this finite disc cover.
The proper Stein factor is this completion, because it is finite,
normal and agrees over the dense punctured disc. We have thus
constructed a proper map with connected fibres
\[
                        f_Q:U_Q\longrightarrow\Delta_s .
                                                        \tag{C.PF.5}\hypertarget{tag-00UU}{}\label{cdp:eq:C.PF.5}
\]
Its reduced central support is exactly $(p^{-1}S)_{\rm red}$,
because the completed Stein disc has just one point over zero.
Every smooth general fibre of $f_Q$ is a connected finite \'{e}tale
cover of $F$. Writing $F=\C^2/\Lambda$, its actual lifted universal
cover identifies it with $\C^2/\Lambda'$ for a finite-index
sublattice $\Lambda'\subset\Lambda$. Hence \eqref{cdp:eq:C.PF.5} is another
two-torus degeneration, to which \ref{cdp:st:C.2.UN.28} applies. This comparison
does not require the original cover to be a pure base cover.

\par\medskip\noindent\textbf{The component pullbacks and their precise degree.}
Fix $Z=Z_i$ and take its fibre product with \eqref{cdp:eq:C.PF.4}. It is a
smooth compact, possibly disconnected surface
\[
               Z_Q=Z\times_U U_Q\longrightarrow Z
                       \quad\hbox{finite \'{e}tale of degree }N .
                                                        \tag{C.PF.6}\hypertarget{tag-00UV}{}\label{cdp:eq:C.PF.6}
\]
Let $H_Q=\rho(h_{i*}\pi_1(Z))$, let $e=|H_Q|$, and put
$k=[Q:H_Q]=N/e$. Path lifting gives exactly $k$ connected
components $Z_1',\ldots,Z_k'$, each a connected cover of $Z$
of degree $e$.

\item These are resolutions of distinct actual normalization components
of the central support of $U_Q$. To verify this comparison, remove
the singular set of the original support and the exceptional
curves of $Z\to S_i$. Over the remaining dense open set the map
$Z\to U$ is its embedding as the smooth single-component locus.
Its base change is therefore an isomorphism onto the corresponding
open subset of the central support of $U_Q$. Removing those proper
analytic curves and points from a connected smooth $Z_j'$ leaves
it connected. Its proper image is consequently an irreducible
central component, and the displayed isomorphism over the dense
open set proves that $Z_j'$ maps birationally onto its normalization.
Different $Z_j'$ cannot give the same component, since their dense
open isomorphisms have disjoint images. Thus the compact-core
bound of \ref{cdp:st:C.2.UN.28} gives the exact useful inequality
\[
                         \sum_{j=1}^k p_g(Z_j')\le1 .
                                                        \tag{C.PF.7}\hypertarget{tag-00UW}{}\label{cdp:eq:C.PF.7}
\]
Other central components can only add nonnegative terms to that
bound.

\item Pullback of a holomorphic two-form under a finite \'{e}tale
cover is injective, by restriction to any open set on which
that cover is biholomorphic. If $p_g(Z)>0$, \ref{cdp:st:C.2.UN.28} on the
original degeneration already gives $p_g(Z)=1$. Every
$Z_j'$ then contributes at least one in \eqref{cdp:eq:C.PF.7}, so $k=1$.
We have proved that
\[
       \rho(h_{i*}\pi_1(Z))=Q
                 \quad\hbox{for every finite quotient }\rho
                         \hbox{ of }\Gamma .           \tag{C.PF.8}\hypertarget{tag-00UX}{}\label{cdp:eq:C.PF.8}
\]

\par\medskip\noindent\textbf{The group-theoretic passage retains the actual subgroup.}
\ref{cdp:st:C.2.UN.44} supplies a normal finitely generated abelian subgroup
$A\subset\Gamma$ of finite index. Every subgroup $H$ of $\Gamma$
is separable by finite quotients in the sense needed here.
If $HA\ne\Gamma$, its image already fails to be all of
the finite quotient $\Gamma/A$. Otherwise, for a proper
$H$, choose $b\in A\setminus H$. The finitely generated
abelian group $A/(H\cap A)$ has a finite quotient detecting
the class of $b$; equivalently some positive integer $n$
satisfies
\[
                         b\notin(H\cap A)+nA .
                                                        \tag{C.PF.9}\hypertarget{tag-00UY}{}\label{cdp:eq:C.PF.9}
\]
The subgroup $nA$ is invariant under every automorphism of $A$,
hence is normal in $\Gamma$, and $\Gamma/nA$ is finite.
Moreover $HN\cap A=(H\cap A)+nA$ for $N=nA$.
Thus the image of $H$ in this finite quotient misses $b$.
This proves directly that \eqref{cdp:eq:C.PF.8} implies
$h_{i*}\pi_1(Z)=\Gamma$, establishing \eqref{cdp:eq:C.PF.1}.

\item The same elementary argument gives residual finiteness of
$\Gamma$. An element outside $A$ is detected in $\Gamma/A$.
A nonzero element of $A$ is detected in $A/nA$ for some $n$,
and therefore in the finite quotient $\Gamma/nA$.
No residual-finiteness assertion for a general degeneration
group is being imported.

\par\medskip\noindent\textbf{The dual complex and the K3 consequence.}
Choose any embedded resolution of the central divisor.
The strict transform of $S_i$ is a smooth resolution of its
normalization, with the same geometric genus as $Z$.
For completeness, the graph of the birational correspondence
between these two smooth resolutions has a common resolution
by point blowups. Pullback and Hartogs extension of holomorphic
two-form coefficients across a point show that each point blowup
preserves $p_g$. The same blowup maps induce isomorphisms on
$\pi_1$, by the $S^3$ and $\PP^1$ neighbourhood comparison.
The two maps from the common resolution to $U$ agree on a
dense open set and hence everywhere. Thus the image of the
strict-transform vertex group in $\Gamma$ is exactly the
image in \eqref{cdp:eq:C.PF.1}.

\item By (C.FG.14), the resolved support has group $\Gamma$.
By (C.FG.16), passage to its dual complex kills the image of
every component group. In particular it kills this surjective
vertex group, so the dual-complex fundamental group is trivial.
If $Z$ has K3 minimal model, its fundamental group is zero:
the K3 is simply connected and the point blowups do not change
its fundamental group. Equation \eqref{cdp:eq:C.PF.1} then gives $\Gamma=1$.

\par\medskip\noindent\textbf{Holomorphic Euler characteristic gives stronger restrictions.}
The degree in \eqref{cdp:eq:C.PF.6} also gives
\[
       \sum_{j=1}^k\chi(\OO_{Z_j'})=N\chi(\OO_Z).
                                                        \tag{C.PF.10}\hypertarget{tag-00UO}{}\label{cdp:eq:C.PF.10}
\]
This follows, for example, by applying surface Riemann--Roch to
the actual \'{e}tale tangent-bundle pullback:
$12\chi(\OO)=\int(c_1^2+c_2)$, and integration of a pulled-back
top-degree class multiplies by the covering degree. Combining
this equality with
$p_g(Z_j')=\chi(\OO_{Z_j'})-1+q(Z_j')$ gives
\[
       1\ \ge\ \sum_j p_g(Z_j')
           \ \ge\ N\chi(\OO_Z)-k,\qquad k\le N .
                                                        \tag{C.PF.11}\hypertarget{tag-00UP}{}\label{cdp:eq:C.PF.11}
\]
For $\chi(\OO_Z)=2$, the right side is at least $N$.
Every finite quotient of $\Gamma$ therefore has order one.
Residual finiteness proves \eqref{cdp:eq:C.PF.2}.

\item For $\chi(\OO_Z)=1$, a finite quotient in which $H_Q$ is
nontrivial has $e\ge2$ and satisfies
\[
                    1\ge N-\frac Ne\ge\frac N2 .
                                                        \tag{C.PF.12}\hypertarget{tag-00UQ}{}\label{cdp:eq:C.PF.12}
\]
It follows that $N=e=2$, so the only such finite quotient
is a quotient of order two. If the image
$H=h_{i*}\pi_1(Z)$ in $\Gamma$ is nontrivial, residual
finiteness provides one such quotient
$\rho:\Gamma\to\Z/2\Z$ detecting a nonidentity element of $H$.
If $\ker\rho$ contained a nonidentity element, another finite
quotient would detect that element. The image of the product
of these two quotient maps would then be a finite quotient
of order greater than two still detecting the original
element of $H$, contradicting \eqref{cdp:eq:C.PF.12}. Thus $\ker\rho=1$,
which proves precisely \eqref{cdp:eq:C.PF.3}.

\item In this case \eqref{cdp:eq:C.PF.6} is a connected double cover of $Z$,
and \eqref{cdp:eq:C.PF.10} gives $\chi=2$. Equation \eqref{cdp:eq:C.PF.7} and
$p_g=\chi-1+q$ give $p_g=1$ and $q=0$. If two different
components with $\chi=1$ both had nontrivial image, the
same actual double cover of the tube would contain their
two distinct normalization resolutions, each of geometric
genus one. This contradicts \ref{cdp:st:C.2.UN.28}. It proves the final
joint assertion.

\end{enumerate}
\end{proof}






\section{Irreducible and reducible torus-fibre Euler criteria}\hypertarget{tag-0343}{}\label{sec:II:07}


\begin{proposition}[Every normal irreducible support is a smooth torus quotient]\hypertarget{tag-01DU}{}\label{cdp:st:C.2.UN.5}
\StatementRecord{cdp:st:C.2.UN.5}{proved}{01DU}{\SeriesRef{cdp:st:A.5.17}, \SeriesRef{cdp:eq:C.UNN.2}; \cite[Theorem~11]{Buchdahl1999}; \cite[Theorem~4.6]{Catanese2002}; \cite[Definition~3.1 and Theorems~3.2--3.3]{Catanese2002}}

Let $f:Y\to B$ be a proper holomorphic map from a smooth
complex threefold to a smooth curve, with connected fibres
and smooth general fibre a complex two-torus. Let $mS$ be
a full fibre whose reduced support is irreducible and normal.
Then $S$ is smooth and $e(S)=0$. More precisely, the
normalized base change $t=s^m$ gives a connected two-torus
$A$ and a finite unramified cyclic map
\[
                         A\longrightarrow S
\]
of degree $m$. Thus $S$ is a two-torus or a bielliptic
surface. In the bielliptic case the order $n$ of $K_S$ is
$2,3,4$, or $6$, and $n\mid m$; in particular a reduced
normal fibre is a torus.

Every two-torus occurs with every positive multiplicity.
Every bielliptic surface of canonical order $n$ occurs
with every multiplicity divisible by $n$. These
realizations can be made in smooth compact threefolds
over $\PP^1$.
\end{proposition}


\begin{proof}

\begin{enumerate}[label={$\langle1\rangle$\arabic*.},leftmargin=*]
\item A normal surface hypersurface has only isolated
singularities. At a singular point of $S$, write the
actual map as $t=u h^m$, with $u$ a holomorphic unit
and $h$ a reduced local equation of $S$. A local
$m$th root of $u$ gives $t=(u^{1/m}h)^m$. The
isolated-hypersurface assertion of \SeriesRef{cdp:st:A.5.17} supplies an
embedded vanishing two-sphere of self-intersection
$-2$ in a small nonzero Milnor fibre of $u^{1/m}h$.
This fibre is an open part of the actual nearby
torus fibre of $f$. A sphere in a torus is nullhomotopic,
and hence has zero homology class and zero
self-intersection. The contradiction proves that $S$
is smooth. This argument uses no Betti-number or
K\"ahler hypothesis on the ambient threefold.

\item Shrink the base to a disc with coordinate $t$. Near
any point of $S$ write $t=u z^m$ with $z$ a smooth
equation of $S$. The normalization $U'$ of the base
change $t=s^m$ is given locally by the explicit maps
\[
        w^m=u,\qquad s=wz,
        \qquad p:U'\longrightarrow U.
 \tag{C.UNN.1}\hypertarget{tag-015B}{}\label{cdp:eq:C.UNN.1}
\]
Since $u$ never vanishes, $p$ is finite and unramified
of degree $m$; these local normalizations glue to the
normalization of the actual fibre product. The map
$h:U'\to\Delta_s$ has $ds=w\,dz$ along its central
support, so that support is smooth and reduced.
For each $s\ne0$, the map $p$ identifies $h^{-1}(s)$
with $f^{-1}(s^m)$, a connected torus. The Stein
factor of $h$ has generic degree one over the normal
disc and is therefore the disc. In particular $U'$
and its central fibre $A$ are connected. After
shrinking the disc, $h$ is a proper submersion.
The deck action is $w\mapsto\zeta_m w$, so it is
free also on $A$, and its quotient there is exactly
$S$.

\item We verify that $A$ is a torus through its actual
Albanese map. Ehresmann's theorem identifies its
integral cohomology ring with that of a real
four-torus. Thus $b_1(A)=4$, and
\cite[Theorem~11]{Buchdahl1999} makes $A$ K\"ahler.
The period-integral map
$a:A\to\operatorname{Alb}(A)$ is consequently
defined by its two independent closed holomorphic
one-forms and induces an isomorphism on integral
$H^1$. Both surfaces have the exterior algebra on
this lattice as their integral cohomology ring.
The pullback of a top-degree generator is therefore
a generator, so $a$ has degree one; holomorphicity
fixes the positive sign. Its pullback on $H^2$ is
also an isomorphism. A curve contracted by $a$
would pair to zero with every pullback class and
hence with a K\"ahler class on $A$, contradicting
positive area. Thus $a$ is finite of degree one
and is an isomorphism. This proves the torus
identification, without imposing a metric on $U'$.
It also agrees with the deformation theorem
\cite[Theorem~4.6]{Catanese2002}.

\item Write a generator of the deck group on
$A=\mathbb C^2/\Lambda$ as $g(z)=Lz+b$.
Every holomorphic torus automorphism has this affine
form: the derivative of a lift is periodic and hence
constant on the compact quotient. The matrix $L$
preserves $\Lambda$ and has finite order. If $1$
were not an eigenvalue of $L$, the equation
$(I-L)z=b$ in $\mathbb C^2$ would give a fixed
point of $g$, contrary to the free deck action.
If $L=I$, the deck group acts by translations and
$S$ is the torus quotient by the resulting finite
subgroup.

\item Otherwise the eigenvalues on $\mathbb C^2$ are
$1$ and a nontrivial root of unity $\zeta$.
The corresponding integral real representation
has characteristic polynomial
$(x-1)^2(x-\zeta)(x-\bar\zeta)$, so the order
$n$ of $\zeta$ is $2,3,4$, or $6$. The images of
$I-L$ and of $I+L+\cdots+L^{n-1}$ are elliptic
subtori with complementary tangent lines. Their
addition map is an isogeny from a product of
elliptic curves to $A$. The kernel of the deck
group's linear representation consists of
translations. Quotienting by it first gives a
torus $A_1$ with a free cyclic action of order
$n$ and quotient $S$. On the canonical form
this action has character $\zeta$, so $K_S$
has exact order $n$, while the invariant
holomorphic one-form space has dimension one.
Equivalently, lifting to the product isogeny
and removing the pure translation kernels gives
the bielliptic presentation
$(E\times F)/G$, with $G$ acting faithfully
by translations on $E$ and with $F/G\simeq\PP^1$.
These are precisely the bielliptic surfaces
of \cite[Definition~3.1 and Theorems~3.2--3.3]{Catanese2002}.
The linear order divides the deck order $m$.
In both cases the unramified cover gives
$m\,e(S)=e(A)=0$.

\item We finish with realizations and their multiplicities.
Given any free cyclic torus quotient
$A\to S=A/\langle g\rangle$ of degree $m$, set
\[
 Y=(A\times\PP^1)/\langle(a,z)\mapsto
                         (g(a),\zeta_m z)\rangle,
 \qquad f([a,z])=z^m.
 \tag{C.UNN.2}\hypertarget{tag-015C}{}\label{cdp:eq:C.UNN.2}
\]
The action is free, hence $Y$ is smooth and compact.
For $t\ne0,\infty$ its $m$ slices are permuted
transitively, and the quotient fibre is exactly
$A$. The full fibres at $0$ and $\infty$ are
$mS$. Thus these are actual compact realizations
of the displayed support and multiplicity.
For a prescribed torus $S$, an index-$m$ sublattice
with cyclic quotient gives the required cyclic
isogeny $A\to S$ for every $m\ge1$.

\item For a prescribed bielliptic surface, use its
product presentation $(E\times F)/G$. The linear
character $G\to\mu_n$ is surjective, and its
kernel acts by translations on both factors.
Consequently its canonical cyclic cover is a
torus. Moreover the character factors through
$\pi_1(E/G)\simeq\mathbb Z^2$, because $G$
acts on $E$ by translations. A surjective
character $\mathbb Z^2\to\mathbb Z/n$ lifts
to a primitive integral homomorphism
$\mathbb Z^2\to\mathbb Z$: choose integer
lifts of its two coefficients, and change one
by a multiple of $n$ to make their gcd one.
The possible prime divisors of the other
coefficient are avoided by the Chinese
remainder theorem; primes dividing $n$ already
cannot divide both coefficients. Pull this
homomorphism back along the connected elliptic
fibration $S\to E/G$ and reduce it modulo
$m=kn$. Its kernel defines a connected cyclic
unramified cover of degree $m$ that factors
through the canonical torus cover. A connected
finite unramified cover of a torus is the
quotient of $\mathbb C^2$ by a finite-index
sublattice, so this new cover is again a
torus. Applying \eqref{cdp:eq:C.UNN.2} proves every claimed
bielliptic multiplicity.

\end{enumerate}
\end{proof}






The statements \ref{cdp:st:C.2.CR.1}--\ref{cdp:st:C.2.CR.6} below retain the
proved local and global reductions. The later
\ref{cdp:st:C.2.UN.7} proves that the contracted irrational
configuration has $e(S)\ge2$ and $\pi_1(S)=1$.
Together with the CDP transgression, or the zero
Euler conclusion \ref{cdp:st:C.2.MH.3}, this excludes the
configuration entirely under the CDP hypotheses.
The unrestricted genus assertion is separately
proved in \ref{cdp:st:C.2.UN.10}.

\begin{proposition}[Topology and coefficients in the contracted-section case]\hypertarget{tag-01BQ}{}\label{cdp:st:C.2.CR.1}
\StatementRecord{cdp:st:C.2.CR.1}{proved}{01BQ}{\SeriesRef{cdp:st:A.5.3}, \SeriesRef{cdp:st:C.2.UN.7}, \SeriesRef{cdp:eq:C.CR.4}, \SeriesRef{cdp:st:A.5.17}, \SeriesRef{cdp:eq:C.CR.6}, \SeriesRef{cdp:eq:C.CR.7}, \SeriesRef{cdp:st:D.1.9}, \SeriesRef{cdp:eq:C.CR.8}}

Let $Y$ be a smooth compact complex threefold with
$b_1(Y)=b_2(Y)=0$, and let $f:Y\to\PP^1$ be a proper
holomorphic surjection with connected fibres and smooth
general fibre a complex two-torus. Let $mS$
be a full fibre with irreducible nonnormal reduced support, and
write $h=\eta\pi:Z\to T\to S$ for its normalization and minimal
normalization resolution. Suppose that $p:Z\to B$ is an
irrational ruling, that the effective divisor in \SeriesRef{cdp:st:A.5.3} is
$D=2C+V$, where $C$ is a section and $V$ is vertical, and
that $\pi$ contracts $C$. Put $g=g(B)$, let $\Gamma$ be the
conductor-image scheme, and let $n$ be the number of irreducible
components of $\Gamma_{\rm red}$.

Then the following are exact consequences of these hypotheses:
\[
\begin{gathered}
 H^1(T;\mathbb R)=H^1(T,\OO_T)=H^2(T,\OO_T)=0,
 \qquad \operatorname{length}R^1\pi_*\OO_Z=g,\\
 \pi_1(S)\simeq\pi_1(\Gamma_{\rm red})\simeq\mathbb Z,
 \qquad e(\Gamma_{\rm red})=n,\\
 \chi(\OO_\Gamma)=-1,\qquad
 \chi(\OO_{\Gamma_{\rm red}})\in\{-1,0\},\qquad
 \sum_{x\in\Gamma_{\rm red}}\bigl(\delta_x-r_x+1\bigr)\le1.
\end{gathered}
\tag{C.CR.1}\hypertarget{tag-00CV}{}\label{cdp:eq:C.CR.1}
\]
Here $r_x$ is the number of normalization branches at $x$.
The full inverse image $D_{\rm red}$ is a smooth genus-$g$
section with trees of smooth rational curves attached; all
its intersections are ordinary double points. For adjacent
components of coefficients $a,b$ one has
\[
                         \gcd(a-1,b-1)=1.          \tag{C.CR.2}\hypertarget{tag-00D6}{}\label{cdp:eq:C.CR.2}
\]
Choose a specified sequence of ruling blowdowns to a bundle
$Z_0\to B$. Write its pushed-forward divisor as
$2C_0+\sum_{j=1}^{v}a_jF_j$, with distinct fibres $F_j$
and $a_j>0$, and let $N_0$ count the reverse blowups whose
centre has multiplicity one in the current divisor. Then
\[
                     e(S)=2-2g-v+N_0+n.           \tag{C.CR.3}\hypertarget{tag-00DA}{}\label{cdp:eq:C.CR.3}
\]
These necessary equations are retained in the Euler
calculation. \ref{cdp:st:C.2.UN.7} later excludes their simultaneous
geometric occurrence under the stated CDP hypotheses.
\end{proposition}


\begin{proof}

\begin{enumerate}[label={$\langle1\rangle$\arabic*.},leftmargin=*]
\item The finite-duality calculation (A.RT.5), whose proof does
not assume rationality, gives $H^2(T,\OO_T)=0$. The
degree-one Leray sequence is consequently
\[
 0\longrightarrow H^1(T,\OO_T)\longrightarrow H^1(Z,\OO_Z)
 \longrightarrow H^0(T,R^1\pi_*\OO_Z)\longrightarrow0.
\tag{C.CR.4}\hypertarget{tag-00DB}{}\label{cdp:eq:C.CR.4}
\]
The middle vector space has dimension $g$. At the point
to which $C$ is contracted, formal functions and the
successive infinitesimal exceptional fibres give a
surjection from the completed stalk of $R^1\pi_*\OO_Z$
onto the $H^1$ of the scheme-theoretic exceptional fibre.
Indeed the restriction maps on $H^1$ of these thickenings
are surjective, since their kernels are supported in
dimension at most one. The completed stalk equals the
stalk, because this higher direct image has finite support.
Restriction from the exceptional fibre to $C$ is likewise
surjective on $H^1$. Thus its length is at least
$h^1(C,\OO_C)=g$, whereas \eqref{cdp:eq:C.CR.4} bounds the total length
by $g$. Equality proves the coherent assertions in the
first line of \eqref{cdp:eq:C.CR.1}.

\item Connectedness of resolution fibres makes the constant-sheaf
Leray map $H^1(T;\mathbb R)\to H^1(Z;\mathbb R)$ injective.
Restriction from $Z$ to its section $C$ is an isomorphism
on $H^1$, because a ruled bundle has simply connected
fibres and point blowups do not change $H^1$. The image of
$H^1(T;\mathbb R)$ restricts to zero on the contracted
curve $C$. This proves its vanishing.

\item We describe the actual divisor comparisons. At every ruling
blowdown $\sigma:W\to W'$ with exceptional curve $E$, let
$a$ be the coefficient of $E$ and let $d$ be the multiplicity
of the pushed-forward boundary at the centre. The identity
\[
 K_W+D_W=\sigma^*(K_{W'}+D_{W'})+(1+a-d)E
\tag{C.CR.5}\hypertarget{tag-00DC}{}\label{cdp:eq:C.CR.5}
\]
and numerical triviality give $a=d-1$. The minimal model
boundary is therefore $2C_0+\sum a_jF_j$. Its reduced
support is the section with disjoint rational fibre leaves.
Every reverse blowup preserves the asserted simple-normal-
crossings tree structure: a positive exceptional coefficient
adds a leaf at a smooth boundary point or subdivides an edge
at a boundary node; a zero exceptional coefficient leaves
the reduced boundary isomorphic to its predecessor. A centre
off the current boundary would create coefficient $-1$,
which persists on its strict transform and contradicts the
effectiveness of the original divisor. In particular
$D_{\rm red}$ is connected, and its fundamental group maps
isomorphically to that of $Z$ by inclusion, with both
identified with $\pi_1(B)$ through the section.

\item This also proves \eqref{cdp:eq:C.CR.2} without discarding coefficients.
Initially every edge has coefficients $(2,a_j)$ and hence
coprime shifted coefficients. A node blowup inserts the
coefficient $a+b-1$, so its shifted coefficient is
$(a-1)+(b-1)$; the two new edge gcds equal the old gcd.
A positive-coefficient blowup at a smooth point of a
coefficient-$a$ boundary creates coefficient $a-1$ and
an edge with shifted coefficients $a-1,a-2$, whose gcd
is one. These are all the possible operations.

\item On a minimal normalization resolution, a coefficient-zero
exceptional curve $E$ has $K_Z\cdot E\ge0$. Since
$(K_Z+D)\cdot E=0$, it follows that $D\cdot E=0$.
Such a curve cannot meet $D_{\rm red}$. Connectedness of
exceptional fibres and of $D_{\rm red}$ now gives the same
support equality as in A.5.16:
\[
 D_{\rm red}=h^{-1}(\Gamma_{\rm red}).
\tag{C.CR.6}\hypertarget{tag-00DD}{}\label{cdp:eq:C.CR.6}
\]
Any exceptional fibre outside this set would correspond to
an isolated normal hypersurface singularity of $S$, because
the normal conductor ideal is reflexive and has no isolated
support. \SeriesRef{cdp:st:A.5.17} excludes every such point in the actual
torus smoothing. Thus $h$ is an isomorphism off the two
closed sets in \eqref{cdp:eq:C.CR.6}.

\item It follows that $S$ is the topological pushout of
$D_{\rm red}\hookrightarrow Z$ and
$D_{\rm red}\to\Gamma_{\rm red}$. This is the actual
pushout: its map to $S$ is a continuous bijection from a
compact space to a Hausdorff space. Triangulation of the
pair makes the first map a cofibration. Van Kampen applied
to its double mapping cylinder therefore gives
\[
 \pi_1(S)=\pi_1(Z)*_{\pi_1(D_{\rm red})}
                         \pi_1(\Gamma_{\rm red})
          \simeq\pi_1(\Gamma_{\rm red}).
\tag{C.CR.7}\hypertarget{tag-00DE}{}\label{cdp:eq:C.CR.7}
\]
Indeed the first amalgamating map is an isomorphism, and
the second map is trivial, since the section $C$ generates
$\pi_1(D_{\rm red})$ and $h(C)$ is a point. Every component
of $\Gamma_{\rm red}$ is dominated by a vertical rational
curve, so its normalization is $\PP^1$. Its fundamental
group is the free group of rank $b_1(\Gamma_{\rm red})$.

\item For the exact local comparison with the nearby torus, take
a triangulated regular neighbourhood retracting onto $S$
and put a sufficiently small full fibre tube inside it.
Perturbing loops off the divisor gives a surjection from
the fundamental group of its punctured torus bundle onto
that of $S$. The former is $\mathbb Z^4\rtimes\mathbb Z$,
which is metabelian. A free group of rank at least two
is not metabelian, since it surjects onto $S_4$.
Consequently the rank in \eqref{cdp:eq:C.CR.7} is zero or one. The
nonzero integral class on every support supplied by \SeriesRef{cdp:st:D.1.9}
excludes zero. This proves the fundamental-group assertion.

\item Normalization of the rational image components gives the
exact identities
\[
 b_1(\Gamma_{\rm red})
   =\sum_x(r_x-1)-n+1=1,\qquad
 e(\Gamma_{\rm red})=2n-\sum_x(r_x-1)=n.
\tag{C.CR.8}\hypertarget{tag-00DF}{}\label{cdp:eq:C.CR.8}
\]
Put $\epsilon=h^0(S,\omega_S)\in\{0,1\}$.
The vanishings already proved, the normalization exact
sequence and conductor duality give
\[
 h^0(\eta_*\OO_T/\OO_S)=1+\epsilon,\qquad
 h^1(\eta_*\OO_T/\OO_S)=\epsilon,
 \qquad\chi(\OO_\Gamma)=-1.
\tag{C.CR.9}\hypertarget{tag-00DG}{}\label{cdp:eq:C.CR.9}
\]
The nilradical sequence and curve duality give
$\chi(\OO_{\Gamma_{\rm red}})\ge-1$ exactly as in
C.2.VII.2: for $\epsilon=0$ the relevant twisted $H^1$
has dimension at most one, and for $\epsilon=1$ the twist
is trivial and connectedness supplies $h^0=1$.
On a connected curve with rational normalizations,
$\chi(\OO_{\Gamma_{\rm red}})=n-\sum_x\delta_x$.
Together with \eqref{cdp:eq:C.CR.8} and $\delta_x\ge r_x-1$, this
proves the remaining assertions of \eqref{cdp:eq:C.CR.1}.

\item Finally $e(Z_0)=4-4g$ and
$e((D_0)_{\rm red})=2-2g+v$. A reverse blowup of
multiplicity at least two increases the Euler numbers
of both the surface and the reduced boundary by one,
leaving their difference unchanged. A multiplicity-one
centre leaves the boundary isomorphic to its predecessor
and increases the complement Euler number by one.
Therefore
\[
 e(Z\setminus D_{\rm red})=2-2g-v+N_0.
\]
Equation \eqref{cdp:eq:C.CR.6}, the excluded off-conductor exceptionals,
and additivity of Euler number now give \eqref{cdp:eq:C.CR.3}.

\end{enumerate}
\end{proof}






\begin{proposition}[The ruling genus is at most two]\hypertarget{tag-01BR}{}\label{cdp:st:C.2.CR.2}
\StatementRecord{cdp:st:C.2.CR.2}{proved}{01BR}{\SeriesRef{cdp:st:C.2.CR.1}, \SeriesRef{cdp:eq:C.CR.10}, \SeriesRef{cdp:st:A.5.3}, \SeriesRef{cdp:eq:C.CR.12}; \cite[Sections~2.2.2--2.2.5]{vanStraten1987}}

In the situation of \ref{cdp:st:C.2.CR.1}, the ruling genus is $g=1$ or
$g=2$. Let $p=\pi(C)$ and $q=\eta(p)$. When $g=2$, the
normalization has just the one point $p$ over $q$, and the
full conductor-image scheme $\Gamma$ is smooth and reduced
at $q$. Its component through $q$ has generic transversal
either an ordinary node or an ordinary cusp.
\end{proposition}


\begin{proof}

\begin{enumerate}[label={$\langle1\rangle$\arabic*.},leftmargin=*]
\item Work over a small neighbourhood of $q$ and fix a generator
of the invertible dualizing sheaf $\omega_S$. Finite
duality then identifies the conductor on the normal branch
$(T,p)$ with its canonical module:
\[
 I_p=\omega_{T,p}\subset\OO_{T,p}.
\]
The conductor ideal on $S$ is the direct sum of these
normalization-branch ideals; thus $I_p$ itself embeds
as an ideal of $R=\OO_{S,q}$ by putting zero on the other
normalization branches.

\item Put $J=(\pi_*\omega_Z)_p$. Proper duality and
Grauert--Riemenschneider vanishing give
\[
 \operatorname{length}(I_p/J)
   =\operatorname{length}(R^1\pi_*\OO_Z)_p=g.
 \tag{C.CR.10}\hypertarget{tag-00CW}{}\label{cdp:eq:C.CR.10}
\]
These are the analytic duality and vanishing maps described
in \cite[Sections~2.2.2--2.2.5]{vanStraten1987}.
For precision, the defect of $\OO_T\to R\pi_*\OO_Z$ is
the finite-length module $R^1\pi_*\OO_Z$ in degree one;
duality turns it into the quotient
$\omega_T/\pi_*\omega_Z$, since the positive direct
images of $\omega_Z$ vanish. Finite-length duality
preserves length.

\item The residue sequence for the smooth exceptional section is
\[
 0\longrightarrow\omega_Z\longrightarrow\omega_Z(C)
   \longrightarrow\omega_C\longrightarrow0 .
 \tag{C.CR.11}\hypertarget{tag-00CX}{}\label{cdp:eq:C.CR.11}
\]
The same vanishing theorem makes its pushforward quotient
at $p$ equal to $H^0(C,\omega_C)$, of dimension $g$.
The sheaf $\pi_*\omega_Z(C)$ embeds in $\omega_T$:
away from the exceptional set its sections are holomorphic
two-forms, and the canonical module of a normal surface
extends these across its isolated singular points.
Comparing lengths with \eqref{cdp:eq:C.CR.10} proves
\[
 I_p=(\pi_*\omega_Z(C))_p,\qquad
 I_p/J\simeq H^0(C,\omega_C).
 \tag{C.CR.12}\hypertarget{tag-00CY}{}\label{cdp:eq:C.CR.12}
\]
In the fixed dualizing trivialization the exact comparison
of \SeriesRef{cdp:st:A.5.3} reads
$\omega_Z=\OO_Z(-2C-V)$, so that
$\omega_Z(C)=\OO_Z(-C-V)$. The quotient in \eqref{cdp:eq:C.CR.12}
is consequently the leading coefficient of order one
along $C$, with the vertical vanishing conditions retained.

\item Let $\mathfrak m$ be the maximal ideal of $R$. Every
element of $\mathfrak m$ pulls back to a function vanishing
on $C$, because $h(C)=q$. Its square therefore vanishes
to order at least two on $C$. Formula \eqref{cdp:eq:C.CR.12}, including
all vertical conditions in $V$, implies
\[
                 I_p\cap\mathfrak m^2\subset J.
 \tag{C.CR.13}\hypertarget{tag-00CZ}{}\label{cdp:eq:C.CR.13}
\]
Indeed an element of $I_p$ is already a section of
$\OO_Z(-C-V)$; an additional zero along $C$ puts it
in $\OO_Z(-2C-V)$.

\item It follows that
\[
 g=\dim(I_p/J)
 \le \dim\bigl(I_p/(I_p\cap\mathfrak m^2)\bigr)
 \le \dim(\mathfrak m/\mathfrak m^2)\le3.
 \tag{C.CR.14}\hypertarget{tag-00D0}{}\label{cdp:eq:C.CR.14}
\]
If the middle dimension were three, the ideal $I_p$
would generate the whole cotangent space. Nakayama's
lemma applied to $\mathfrak m/I_p$ would give
$I_p=\mathfrak m$. The full conductor ideal, which
contains $I_p$, would then be $\mathfrak m$, making its
image zero-dimensional. This contradicts the pure
one-dimensional Cohen--Macaulay conductor image.
Thus the middle dimension is at most two, proving
$g\le2$.

\item When $g=2$, choose $x,y\in I_p$ whose linear parts
are independent, and complete them to coordinates
$(x,y,z)$ on the smooth ambient threefold. The full
conductor ideal contains $(x,y)$ in $R$. The defining
equation of $S$ belongs to $(x,y)$, since the conductor
has dimension one. The quotient by $(x,y)$ is therefore
$\mathbb C\{z\}$. A strictly larger proper conductor
ideal would give a zero-dimensional quotient of this
ring. Purity hence forces the full conductor ideal
to be exactly $(x,y)$, proving smoothness and
reducedness of $\Gamma$ at $q$.

\item There is no other normalization point over $q$.
The two elements $x,y$ were chosen in the single
summand $I_p$, so both vanish identically on any other
local normalization branch. Such a two-dimensional
hypersurface branch cannot be contained in the
one-dimensional smooth germ $\{x=y=0\}$.

\item Finally conductor duality gives
$Q\simeq\omega_\Gamma\otimes\omega_S^{-1}|_\Gamma$.
Near the smooth reduced curve $\Gamma$ this is a
line bundle, so the transverse normalization quotient
has length one. A reduced plane curve of delta
invariant one is a node or a cusp. One can see this
directly: multiple branches require exactly two
smooth branches with intersection multiplicity one;
a single branch has normalization
$\mathbb C\{u\}$ and a colength-one subring, necessarily
$\mathbb C\{u^2,u^3\}$.

\end{enumerate}
\end{proof}






\begin{proposition}[The genus-two node alternative is excluded]\hypertarget{tag-01BS}{}\label{cdp:st:C.2.CR.3}
\StatementRecord{cdp:st:C.2.CR.3}{proved}{01BS}{\SeriesRef{cdp:st:C.2.CR.2}}

In the genus-two situation of \ref{cdp:st:C.2.CR.2}, the component
of $\Gamma$ through $q=h(C)$ cannot have generic
transversal an ordinary node. Consequently the
remaining genus-two alternative has an ordinary cusp
at its generic point, and exactly one strict conductor
curve of coefficient two and degree one lies above
that component.
\end{proposition}


\begin{proof}

\begin{enumerate}[label={$\langle1\rangle$\arabic*.},leftmargin=*]
\item We first record an additional consequence of the
specified boundary blowups in C.2.CR.1: every
coefficient-one component of $D_{\rm red}$ is a leaf.
This holds on the initial section-and-fibres star.
A smooth boundary blowup creates coefficient one
only when the old coefficient is two, and creates
it as a leaf. A node blowup could create coefficient
one only from an edge $(1,1)$, which (C.CR.2)
excludes. Subsequent blowups preserve the valence
one of an existing coefficient-one vertex: at its
unique boundary node they merely subdivide that edge,
and at a smooth point its new exceptional coefficient
is zero.

\item Let $E_0$ be the connected exceptional fibre containing
$C$. In the boundary tree, the first nonexceptional
vertex on every path from $C$ meets $E_0$, so its image
passes through $p$. Its further image is therefore
a component of $\Gamma$ through $q$. By \ref{cdp:st:C.2.CR.2}
there is only one such component, denoted $G$.

\item Suppose its generic transversal were a node.
Every strict conductor curve above $G$ then has
coefficient one. Thus each first nonexceptional
vertex just found is a leaf. Every strict conductor
curve is consequently such a leaf, attached at a
single point to $E_0$. No additional exceptional
tail can meet it, since that would raise its
boundary valence. The reduced conductor on $T$
is therefore a union of rational curves whose only
possible multibranch point is $p$; every component
passes through $p$, is smooth away from $p$, and
has only one branch at $p$.

\item All their images pass through $q$, where $\Gamma$
is smooth, so the conductor image has just the one
global component $G$. It also has no multibranch
point elsewhere. Indeed, at $x\ne q$ each point
of the normalization over $x$ receives just one
local branch of $G$. If $G$ had two distinct local
branches there, their receiving normalization
points would be distinct. The corresponding
irreducible hypersurface branches of $S$ must
intersect in a curve, because their distinct
local equations generate a height-two ideal in
the regular local ring of the threefold. That
intersection curve is a conductor-image branch
received by both normalization points, contradicting
their different singleton branch sets.

\item The normalization of $G$ is rational and all its
germs are unibranch. The normalization map is a
continuous bijection from a compact space to a
Hausdorff space and hence a homeomorphism.
Thus $\pi_1(G)=1$, contradicting (C.CR.1).
The generic cusp alternative is the only one left.
Its transverse normalization is unibranch with
conductor exponent two, so it contributes exactly
one strict coefficient-two curve, of degree one
over $G$.

\end{enumerate}
\end{proof}






\begin{proposition}[A smooth critical curve with generic cusp is a constant cusp]\hypertarget{tag-01BT}{}\label{cdp:st:C.2.CR.4}
\StatementRecord{cdp:st:C.2.CR.4}{proved}{01BT}{\SeriesRef{cdp:st:A.5.17}, \SeriesRef{cdp:eq:C.CR.17}, \SeriesRef{cdp:st:A.5.15}, \SeriesRef{cdp:eq:C.CR.18}}

Let a reduced hypersurface germ $\{f=0\}$ in a smooth
complex threefold occur in a fibre whose nearby smooth
fibres are complex two-tori. Suppose that the full reduced
critical locus of $f$ near the point is a smooth curve
and that the generic transverse singularity is an ordinary
cusp. Then there are local holomorphic coordinates in which
\[
                            f=X^2+Y^3.             \tag{C.CR.15}\hypertarget{tag-00D1}{}\label{cdp:eq:C.CR.15}
\]
An analytic unit and a common full-fibre multiplicity may
be retained in the original map: they do not change this
assertion about the reduced hypersurface.
\end{proposition}


\begin{proof}

\begin{enumerate}[label={$\langle1\rangle$\arabic*.},leftmargin=*]
\item Choose coordinates $(x,y,z)$ with critical curve
$\{x=y=0\}$. The rank-one quadratic form along this curve
can be written
\[
 a(z)x^2.
\]
To obtain these coordinates without a meromorphic change,
divide the matrix of its coefficients by its common
vanishing order at $z=0$. A primitive symmetric rank-one
matrix has a unit entry, factors as a unit times the
square of a primitive linear form, and that linear form
extends to an invertible holomorphic change of the normal
coordinates. The intrinsic cubic coefficient on its
kernel is then $b(z)y^3$. The generic cusp hypothesis says
that both $a$ and $b$ are nonzero functions; shrink the
disc so that their only possible zero is at $z=0$.

\item If $a(0)\ne0$, the Morse lemma with parameters gives
\[
 f=X^2+y^3u(y,z),\qquad u(0,z)=b(z)
\]
up to multiplication of $b$ by a holomorphic unit.
If $u(0,0)=0$, the reduced support of $y^3u(y,z)$ is
singular: $u$ has a component through the origin
different from $y=0$, because $u(0,z)$ is not identically
zero. The arbitrary-suspension clause of \SeriesRef{cdp:st:A.5.17} then
puts a sphere of nonzero self-intersection into a
nearby actual torus fibre, which is impossible.
Thus $u$ is a unit, and replacing $y$ by a holomorphic
cube root of $u$ times $y$ gives \eqref{cdp:eq:C.CR.15}.

\item It remains to exclude $a(0)=0$. Consider the actual
holomorphic perturbation
\[
                  f_{\epsilon,\delta}
                      =f+\epsilon x^2+\delta y^3.
 \tag{C.CR.16}\hypertarget{tag-00D2}{}\label{cdp:eq:C.CR.16}
\]
Choose small parameters so that $a(z)+\epsilon$ has a
simple zero $z_\epsilon$ and
$b(z_\epsilon)+\delta\ne0$. Such choices exist:
regular values of the nonconstant function $a$ are
dense near zero, and finitely many values of $\delta$
are then excluded.

\item We specify the germ produced at
$(0,0,z_\epsilon)$. Weierstrass preparation in $y$,
followed by a Tschirnhaus substitution, writes it as
a holomorphic unit times
\[
             Y^3+x^2\bigl(\beta(x,z)Y+\gamma(x,z)\bigr).
 \tag{C.CR.17}\hypertarget{tag-00D3}{}\label{cdp:eq:C.CR.17}
\]
Here the absence of a linear transverse jet gives
the divisibility of the constant term by $x^2$,
and the rank-one quadratic form gives the divisibility
of the coefficient of $Y$ by $x^2$. More explicitly,
before the Tschirnhaus substitution the coefficient
of $y^2$ is divisible by $x$, while those of $y$
and of $1$ are divisible by $x^2$; these divisibilities
are preserved by that substitution.
The function $\gamma(0,z)$ is a unit multiple of
$a(z)+\epsilon$. Consequently
\[
             Z=\beta(x,z)Y+\gamma(x,z)
\]
is an invertible change of the curve coordinate at
the chosen simple zero. Equation \eqref{cdp:eq:C.CR.17} is thus
a unit times $Y^3+x^2Z$. This remaining unit can
also be removed: this polynomial is homogeneous
of degree three, so scaling all three coordinates
by a holomorphic cube root of the unit gives the
exact equation
\[
                            Y^3+x^2Z.
 \tag{C.CR.18}\hypertarget{tag-00D4}{}\label{cdp:eq:C.CR.18}
\]
\SeriesRef{cdp:st:A.5.15} supplies a compact spherical cycle of square $-4$
in an arbitrarily small nonzero level of this germ.
Indeed its Hirzebruch formula is
$-2(b/\gcd(b,c))(c/\gcd(b,c))$, here with
$(b,c)=(2,1)$. The assertion is about the homology
class of the explicit continuous map from $S^2$;
no embedded representative is needed.

\item For completeness we prove the boundary comparison
needed to put this spherical cycle in the original fibre.
Choose a sufficiently small closed ball whose
boundary meets the smooth critical curve transversely
and is transverse to the smooth stratum of $f=0$.
On the compact link of the critical curve, both
$a$ and $b$ are units. For all sufficiently small
$(\epsilon,\delta)$ the parametrized Morse and
cusp-coordinate changes therefore give holomorphic
product-cusp charts there, retaining $z$ as curve
coordinate. In each such chart the vector
$\partial/\partial z$ is tangent to every level
of $f_{\epsilon,\delta}$. Along the critical link,
an appropriate real radial vector in this direction
is transverse to the fixed sphere. This remains
true on a neighbourhood of that link, uniformly
for the small perturbations. Finitely many charts
suffice. On the rest of the boundary ordinary
compactness and transversality at parameter zero
give the same conclusion for small target values.
Thus all sufficiently small nonzero levels have
uniformly transverse boundary. The full-critical-
locus hypothesis ensures that no other boundary
singular curve has been omitted from this argument.

\item For fixed small parameters, the nonzero critical
values inside the ball are finite. The local
spherical cycle from \eqref{cdp:eq:C.CR.18} can therefore be transported
through regular target values to a fixed small
value $t_1$, avoiding those critical values.
Choose the parameters small enough that every
critical value throughout the segment joining
them to zero has modulus less than $|t_1|/2$.
This is possible by compactness: every limit of
critical points is a critical point of $f$, and
its value is zero; the just-proved product charts
control the boundary as well. The family of
$t_1$-levels along this parameter segment is a
proper submersion with transverse boundary.
Ehresmann transport identifies it with the
original $f=t_1$ level and preserves the square
of the transported spherical class.

\item The original level lies in a smooth nearby torus
fibre. A two-torus is aspherical, so every sphere
in it is null-homotopic and has zero homology
class and self-intersection. This contradicts
the nonzero square obtained above. Hence $a(0)\ne0$,
and the first part of the proof gives \eqref{cdp:eq:C.CR.15}.
If the original map is a unit times $f^m$, absorb
the unit by its local $m$th root before this
comparison; the resulting sphere lies in the
actual level with parameter $t_1^m$.

\end{enumerate}
\end{proof}






\begin{proposition}[A contracted horizontal section is elliptic]\hypertarget{tag-01BU}{}\label{cdp:st:C.2.CR.5}
\StatementRecord{cdp:st:C.2.CR.5}{proved}{01BU}{\SeriesRef{cdp:st:C.2.CR.1}, \SeriesRef{cdp:st:C.2.CR.2}, \SeriesRef{cdp:st:C.2.CR.3}, \SeriesRef{cdp:st:C.2.CR.4}}

In the situation of \ref{cdp:st:C.2.CR.1} one has
\[
                              g(B)=1.              \tag{C.CR.19}\hypertarget{tag-00D5}{}\label{cdp:eq:C.CR.19}
\]
In particular the remaining exact Euler expression is
\[
                            e(S)=n-v+N_0.          \tag{C.CR.20}\hypertarget{tag-00D7}{}\label{cdp:eq:C.CR.20}
\]
\end{proposition}


\begin{proof}
\begin{enumerate}[label={$\langle1\rangle$\arabic*.},leftmargin=*]
\item 
\noindent\textbf{Step 1.} By \ref{cdp:st:C.2.CR.2} the only other possibility is $g=2$.
The conductor image is then smooth and reduced at
$q=h(C)$, and its generic transversal is a node or
a cusp. \ref{cdp:st:C.2.CR.3} excludes the node. For the cusp,
the reduced critical locus near $q$ is precisely
that smooth conductor curve: every positive-dimensional
singular locus of a reduced hypersurface surface
belongs to its nonnormal locus, and there are no
other isolated points after the neighbourhood is
shrunk. \ref{cdp:st:C.2.CR.4} makes the germ a product cusp.
Its normalization is smooth, contrary to the
exceptional genus-two curve contracted to the
normalization point $p$. Thus $g=2$ is excluded.
Substitution in (C.CR.3) gives \eqref{cdp:eq:C.CR.20}.

\end{enumerate}
\end{proof}






\begin{proposition}[The elliptic contraction has a planar conductor and only double transverse branches at its image]\hypertarget{tag-01BV}{}\label{cdp:st:C.2.CR.6}
\StatementRecord{cdp:st:C.2.CR.6}{proved}{01BV}{\SeriesRef{cdp:st:C.2.CR.1}, \SeriesRef{cdp:eq:C.CR.22}}

In the remaining case $g=1$ of \ref{cdp:st:C.2.CR.1}, let
$p=\pi(C)$ and $q=\eta(p)$. The conductor-image
scheme is a local complete intersection at $q$.
More precisely, ambient coordinates can be chosen
so that its ideal is
\[
                         (x,g(y,z)).               \tag{C.CR.21}\hypertarget{tag-00D8}{}\label{cdp:eq:C.CR.21}
\]
Every global conductor component passing through
$q$ has generic transverse type $A_k$ for some
$k\ge1$. For type $A_{2s-1}$, its two transverse
normalization branches have conductor coefficient
$s$ each; their global covering degrees sum to two.
For type $A_{2s}$ there is one coefficient-$2s$
normalization branch, of covering degree one.
\end{proposition}


\begin{proof}

\begin{enumerate}[label={$\langle1\rangle$\arabic*.},leftmargin=*]
\item The residue comparison (C.CR.12) is valid also
when $g=1$. By (C.CR.13) its nonzero one-dimensional
quotient requires an element of the conductor
summand $I_p$ with nonzero ambient linear part.
Call this element $x$ and extend it to coordinates
$(x,y,z)$ on the smooth threefold.

\item The full conductor ideal contains $x$ and has a
pure one-dimensional Cohen--Macaulay quotient.
After quotienting the ambient local ring by $x$,
its remaining ideal is an unmixed height-one
ideal in the regular two-dimensional ring
$\mathbb C\{y,z\}$. Such an ideal is principal:
it is the intersection of its height-one primary
localizations, and each is a power of a principal
prime in this factorial ring. Equivalently its
quotient has depth one, so the ideal is a rank-one
reflexive module on this smooth surface and
is invertible. This proves \eqref{cdp:eq:C.CR.21}, retaining
all multiplicities in $g$.

\item We give the precise adjunction comparison behind
the transverse multiplicity assertion. For a
reduced plane curve germ $D\subset W$, let
$\nu:\widetilde D\to D$ be its normalization,
and let $\rho:\widetilde W\to W$ be an embedded
resolution with smooth disjoint strict transform
$\widehat D$. Finite curve duality identifies
\[
 \mathfrak c_D\,\omega_D=\nu_*\omega_{\widetilde D}.
\]
The residue sequence on $\widetilde W$ and
Grauert--Riemenschneider vanishing identify the
right side with the residue of
$\rho_*\omega_{\widetilde W}(\widehat D)$.
In a local generator $du\wedge dv/f_D$ of the
ambient meromorphic residue description, a
conductor element $a$ therefore lifts to a form
\[
                     a\,du\wedge dv/f_D
\]
whose pullback has no exceptional pole. On the
first plane blowup, with exceptional curve $E$,
the canonical divisor has coefficient one and
the total transform of $D$ has coefficient
$m=\operatorname{mult}_0 D$. The coefficient
of this pulled-back form along $E$ is
$\operatorname{ord}_0(a)+1-m$. Hence
\[
 \mathfrak c_D\subset\mathfrak m_D^{\,m-1}
 \quad\hbox{in the sense that every ambient
 representative has order at least }m-1.
 \tag{C.CR.22}\hypertarget{tag-00D9}{}\label{cdp:eq:C.CR.22}
\]
Adding a multiple of $f_D$ does not change this
order bound, since $f_D$ has order $m$.

\item At a general smooth point of any conductor
component through $q$, the ambient coordinate
$x$ remains in its conductor ideal and has
a nonzero transverse linear part: its zero
set is a smooth ambient plane containing the
conductor curve. Applying \eqref{cdp:eq:C.CR.22} to a
transverse plane slice gives $m\le2$.
The slice is singular, so $m=2$.
Weierstrass preparation, completion of the
square, and absorption of a unit then put
a reduced multiplicity-two plane curve in
the form $u^2-v^{k+1}=0$.

\item For $k=2s-1$, the two branches are
$u=\pm v^s$; their intersection multiplicity
is $s$, and the conductor on each normalized
branch is $(v^s)$. For $k=2s$, the normalization
is $v=t^2$, $u=t^{2s+1}$, whose semigroup
has conductor $2s$, giving $(t^{2s})$.
These calculations specify both the coefficients
and the numbers of transverse branches; finite
normalization over the global image component
expresses those numbers as the asserted sums
of covering degrees.

\end{enumerate}
\end{proof}






\begin{proposition}[A smooth full critical curve with generic node is a product node]\hypertarget{tag-01EZ}{}\label{cdp:st:C.2.UN.9}
\StatementRecord{cdp:st:C.2.UN.9}{proved}{01EZ}{\SeriesRef{cdp:eq:C.UNNODE.2}, \SeriesRef{cdp:st:A.5.15}, \SeriesRef{cdp:eq:C.UNNODE.3}}

Let a reduced hypersurface germ $\{f=0\}$ in a smooth
complex threefold occur in a fibre of a proper torus-surface
degeneration. Suppose that the full reduced critical locus
of $f$ is a smooth curve and its generic transverse
singularity is an ordinary node. There are local
holomorphic coordinates in which
\[
                            f=XY.                  \tag{C.UNNODE.1}\hypertarget{tag-015D}{}\label{cdp:eq:C.UNNODE.1}
\]
The conclusion about the reduced hypersurface also holds
when the actual map has a holomorphic unit and a common
positive multiplicity.
\end{proposition}


\begin{proof}

\begin{enumerate}[label={$\langle1\rangle$\arabic*.},leftmargin=*]
\item Choose coordinates $(x,y,z)$ with critical curve
$\Gamma=\{x=y=0\}$. Since $f$ vanishes on $\Gamma$ and
its first derivatives vanish there, its transverse
quadratic term is
\[
                  (x,y)Q(z)(x,y)^{\mathsf t},
\]
where $Q$ is a holomorphic symmetric $2\times2$ matrix.
Its determinant is not identically zero. Shrink the
curve disc so that $\det Q(z)\ne0$ for $z\ne0$.
If $Q(0)$ is invertible, the holomorphic Morse lemma
with parameter $z$ gives $f=X^2+Y^2$, and a linear
change gives \eqref{cdp:eq:C.UNNODE.1}.

\item Suppose instead that $Q(0)$ is singular. We construct
arbitrarily small constant symmetric matrices $E$
for which $Q(z)+E$ has a rank-one value at a simple
zero of its determinant near the origin. Choose
$z_0\to0$ with $Q'(z_0)\ne0$; such points exist
because the nonconstant determinant makes $Q$
nonconstant. Choose a rank-one symmetric matrix
$R$ close to $Q(0)$ with
\[
                  \operatorname{tr}
                      (\operatorname{adj}(R)Q'(z_0))\ne0.
\]
This choice is possible also when $Q(0)=0$: for a
rank-one symmetric matrix with kernel spanned by
$v$, its adjugate is a nonzero scalar multiple of
$vv^{\mathsf t}$, and the polynomial
$v^{\mathsf t}Q'(z_0)v$ is not identically zero.
Thus the forbidden kernels form a proper finite
set in the projective line, and the rank-one
matrix can be chosen as close to $Q(0)$ as desired.
Set $E=R-Q(z_0)$. Then $E\to0$, and
\[
 \det(Q(z_0)+E)=0,\qquad
 \left.\frac{d}{dz}\det(Q(z)+E)\right|_{z=z_0}\ne0.
 \tag{C.UNNODE.2}\hypertarget{tag-015E}{}\label{cdp:eq:C.UNNODE.2}
\]

\item Consider the actual perturbation
$f_E=f+(x,y)E(x,y)^{\mathsf t}$ at $(0,0,z_0)$.
The rank-one Morse lemma, retaining $(y,z)$ as
parameters after a linear normal-coordinate change,
writes
\[
                         f_E=X^2+y^2h(y,z).
\]
The factor $y^2$ follows from $f_E(0,0,z)=0$ and
the vanishing of its normal first derivative along
the same curve. Equation \eqref{cdp:eq:C.UNNODE.2} says exactly
that $h(0,z_0)=0$ and $\partial h/\partial z(0,z_0)
\ne0$, up to a nowhere-zero coefficient from the
nondegenerate quadratic direction. Hence the actual
coordinate substitution $Z=h(y,z)$ gives
\[
                            f_E=X^2+y^2Z.          \tag{C.UNNODE.3}\hypertarget{tag-015F}{}\label{cdp:eq:C.UNNODE.3}
\]
The explicit construction in \SeriesRef{cdp:st:A.5.15} gives a compact
spherical class in a small nonzero level of
\eqref{cdp:eq:C.UNNODE.3}. Its square is
$-2(2/\gcd(2,1))(1/\gcd(2,1))=-4$.
The construction supplies a continuous map from
$S^2$ representing this class; embeddedness is
not required.

\item Here is the precise transport into an original
fibre. Choose a small closed ball whose boundary
is transverse to $\Gamma$ and to the smooth
stratum of $f=0$. Along the compact boundary link
of $\Gamma$, the matrix $Q$ is invertible.
For all sufficiently small $E$, the parameter
Morse lemma therefore gives product-node charts
there, retaining the curve coordinate. In these
charts $\partial/\partial z$ is tangent to every
level of $f_E$. A real radial curve-direction
vector is transverse to the fixed boundary at
the critical link and remains transverse on its
neighbourhood, uniformly for small $E$. On the
remaining compact boundary, ordinary transversality
gives the same uniform conclusion for small target
values. The full-critical-locus hypothesis ensures
that these charts cover every boundary singular
point requiring this argument.

\item For a fixed small $E$, transport the spherical
class from \eqref{cdp:eq:C.UNNODE.3} through regular target values
to a fixed small value $t_1$, avoiding the finitely
many nonzero critical values in the ball. Choose
$E$ small enough that every critical value along
the segment from $E$ to $0$ has modulus less than
$|t_1|/2$. This follows by compactness: a limit of
critical points is a critical point of $f$ with
value zero; the just-proved boundary charts rule
out loss of transversality at the boundary.
The $t_1$-levels along that segment form a proper
submersion with boundary. Ehresmann transport
therefore maps the actual spherical class to a
class of square $-4$ in the original $f=t_1$ level.
Its inclusion into the nearby torus preserves
the square, since both the compact representative
and its intersection push-off lie in that open
level. But this inclusion also maps its representing
$S^2$ to an aspherical torus, where it is
nullhomotopic and has zero homology class. This
contradiction proves $\det Q(0)\ne0$.

\item If the original map is $u f^m$, replace the reduced
equation by $u^{1/m}f$ using a local holomorphic
root of the unit. Its critical-set germ is still
$\Gamma$: every critical component through the
point has constant value zero, and the singular
locus of the reduced zero hypersurface is unchanged.
The argument then places the
same class in the actual nearby level with base
parameter $t_1^m$. This proves the last assertion.

\end{enumerate}
\end{proof}






\begin{proposition}[A contracted irrational horizontal section is elliptic without ambient Betti assumptions]\hypertarget{tag-01CN}{}\label{cdp:st:C.2.UN.10}
\StatementRecord{cdp:st:C.2.UN.10}{proved}{01CN}{\SeriesRef{cdp:st:C.2.CR.1}, \SeriesRef{cdp:st:C.2.CR.2}, \SeriesRef{cdp:st:C.2.UN.9}, \SeriesRef{cdp:st:C.2.CR.4}}

Let $f:Y\to B$ be a proper holomorphic map from a
smooth complex threefold to a smooth curve, with
connected fibres and smooth general fibre a complex
two-torus. Let $mS$ be an irreducible nonnormal full
fibre. Write $h=\eta\pi:Z\to T\to S$ for the
normalization and its minimal resolution. Suppose
that $Z$ is ruled over a curve of positive genus,
that its effective adjunction divisor is
$D=2C+V$, with $C$ a section and $V$ vertical,
and that $\pi$ contracts $C$. Then the ruling
base has genus one.
\end{proposition}


\begin{proof}

\begin{enumerate}[label={$\langle1\rangle$\arabic*.},leftmargin=*]
\item We first supply the comparisons that previously
came from the global Betti assumptions. On a
small full fibre tube $U$, the rank-one torsion-free
sheaf $f_*K_U$ is a line bundle. A generator has
no horizontal zeros, since its restriction to a
smooth torus is a nonzero canonical form. Its
zero divisor is consequently $kS$ for some
$k\ge0$, after shrinking away other special
fibres. Thus the evaluation map gives the actual
isomorphism
\[
 K_U\simeq\OO_U(kS),\qquad
 \omega_S\simeq\OO_U((k+1)S)|_S.                \tag{C.UNCR.1}\hypertarget{tag-0158}{}\label{cdp:eq:C.UNCR.1}
\]
The divisor $mS$ is principal, so both these line
bundles are torsion. The map to the smooth curve
is flat, and the holomorphic Euler characteristic
of its full fibres is constant. The filtration
of $\OO_{mS}$ has quotients
$\OO_S(-jS)$ for $0\le j<m$; each is a torsion
line-bundle twist of $\OO_S$. Analytic
Riemann--Roch gives the same Euler characteristic
for these twists, since their rational Chern
characters are one. It follows that
\[
            0=\chi(\OO_{mS})=m\,\chi(\OO_S).
 \tag{C.UNCR.2}\hypertarget{tag-0159}{}\label{cdp:eq:C.UNCR.2}
\]
The exact normalization and resolution duality
maps therefore give $K_Z+D=h^*\omega_S$, with
the same effective divisor and torsion right-hand
side as in \ref{cdp:st:C.2.CR.1}--\ref{cdp:st:C.2.CR.2}.

\item Here are the parts of that argument needed here,
with no transgression-class assumption. Finite
duality embeds $\omega_T$ into the torsion line
$\eta^*\omega_S$, with its sections vanishing on
the nonempty conductor. A nonzero section of a
torsion line on a connected compact normal space
is nowhere zero: taking a trivializing power
makes it a nonzero constant. Consequently
$H^2(T,\OO_T)=0$. The coherent Leray sequence
then bounds $\operatorname{length}R^1\pi_*\OO_Z$
by $h^1(Z,\OO_Z)=g$, whereas formal functions
and restriction to the contracted section $C$
give length at least $h^1(C,\OO_C)=g$ at its
image point $p$. Thus equality holds at $p$ and
there are no further contributions.

\item At $q=\eta(p)$, fix a local generator of $\omega_S$.
Finite duality identifies the branch conductor
ideal $I_p$ with $\omega_{T,p}$, embedded as an
actual ideal summand in $\OO_{S,q}$. Proper
duality and the residue sequence of $C$ give
exactly
\[
 I_p=(\pi_*\omega_Z(C))_p,\qquad
 I_p/(\pi_*\omega_Z)_p\simeq H^0(C,\omega_C).
 \tag{C.UNCR.3}\hypertarget{tag-015A}{}\label{cdp:eq:C.UNCR.3}
\]
The equality of lengths follows from the preceding
formal-functions computation. In the fixed local
dualizing trivialization these sheaves are
$\OO_Z(-C-V)$ and $\OO_Z(-2C-V)$.
The maximal ideal of $\OO_{S,q}$ vanishes on $C$;
its square supplies an additional order along
$C$ without losing any prescribed order along
$V$. Hence
$I_p\cap\mathfrak m_q^2\subset(\pi_*\omega_Z)_p$.
The embedding dimension three and Nakayama
argument of \ref{cdp:st:C.2.CR.2} now give $g\le2$. When
$g=2$ they give the actual conductor ideal
$(x,y)$, a unique normalization point $p$ above
$q$, and generic transverse delta invariant
one. Thus the full conductor-image curve is
smooth and reduced at $q$, and the generic
transversal is a node or a cusp.

\item Every positive-dimensional critical component
of the reduced surface is part of that conductor
curve; after shrinking, its full reduced critical
locus is exactly the smooth curve just obtained.
The node case is a product node by \ref{cdp:st:C.2.UN.9}.
Its normalization has two smooth points above
$q$, contradicting the unique point $p$.
The cusp case is a product cusp by \ref{cdp:st:C.2.CR.4}.
Its explicit normalization
$(u,z)\mapsto(X,Y,z)=(u^3,-u^2,z)$, after
adjusting the harmless sign in $X^2+Y^3$,
is smooth. This contradicts a minimal
normalization resolution contracting a
positive-genus curve to $p$. Both genus-two
possibilities are excluded, and $g=1$ follows.

\end{enumerate}
\end{proof}






\begin{proposition}[The reduced horizontal case is a minimal elliptic ruled double-curve fibre]\hypertarget{tag-01C3}{}\label{cdp:st:C.2.IRR.1}
\StatementRecord{cdp:st:C.2.IRR.1}{proved}{01C3}{\SeriesRef{cdp:st:A.5.3}, \SeriesRef{cdp:st:A.5.17}, \SeriesRef{cdp:st:C4e}, \SeriesRef{cdp:st:A.5.15}}

Let $f:Y\to C_0$ be a proper holomorphic map from a smooth compact
complex threefold with $H^2(Y,\mathbb Q)=0$, with connected fibres
and smooth general fibre a complex two-torus. Let $mS$ be a full
fibre with irreducible nonnormal reduced support. Let
$h=\eta\pi:Z\to T\to S$ be the minimal resolution of its
normalization, and suppose that $p:Z\to B$ is an irrational
ruling. Write the effective divisor of \SeriesRef{cdp:st:A.5.3} as
$D_Z=H+V$, where $V$ is vertical and the horizontal part $H$
is reduced, of total degree two over $B$.

Then $B$ is elliptic, $T=Z$ is a smooth minimal ruled surface
over $B$, and $D_Z=H$ is a disjoint union of smooth elliptic
curves. The reduced conductor image is a disjoint union of
smooth elliptic curves. Its inverse image $H$ maps onto it
by a finite etale covering of total degree two over each
component. The support $S$ has only ordinary double-curve
normal crossings and
\[
                         e(S)=0.                 \tag{C.IR.1}\hypertarget{tag-00OH}{}\label{cdp:eq:C.IR.1}
\]
No assumption on the initial normalization singularities or
on the ambient Picard pullback is made.
\end{proposition}


\begin{proof}

\begin{enumerate}[label={$\langle1\rangle$\arabic*.},leftmargin=*]
\item The canonical comparison of
\SeriesRef{cdp:st:A.5.3} gives
\[
 K_Z+D_Z=h^*\omega_S,\qquad c_1(K_Z+D_Z)=0
 \quad\hbox{in }H^2(Z,\mathbb Q).
 \tag{C.IR.2}\hypertarget{tag-00OI}{}\label{cdp:eq:C.IR.2}
\]
If $g(B)\ge2$, that proposition makes the horizontal part
twice a section.  Since the given horizontal part is reduced,
$g(B)=1$.  Contract the vertical ruling blowups to a minimal
ruled surface $\rho:Z\to Z_0$, and write
$D_0=\rho_*D_Z=H_0+V_0$.  Choose a section $C_0$ with
$C_0^2=-n$ and a ruling fibre $F$.  Numerically,
\[
 K_{Z_0}=-2C_0-nF,\qquad H_0=2C_0+bF,\qquad V_0=vF,
 \qquad b+v=n.
\]
Here $v\ge0$.  Adjunction gives $p_a(H_0)=1-v$.
An irreducible $H_0$ has normalization of genus at least one
by Riemann--Hurwitz.  If $H_0$ has two components, they are
sections, and its arithmetic genus is one plus their
intersection number.  Hence $v=0$.  In the irreducible case
$H_0$ is smooth elliptic; in the other case its two smooth
elliptic sections are disjoint.

\item We justify the numerical comparison
under every blowdown and retain its exact divisor map.
For a point blowup $\sigma:W\to W'$ with exceptional curve
$E$, let $a$ be its coefficient in $D_W$, let
$D'=\sigma_*D_W$, and let $d$ be the multiplicity of $D'$
at the centre.  Then
\[
 K_W+D_W=\sigma^*(K_{W'}+D')+(1+a-d)E.
 \tag{C.IR.4}\hypertarget{tag-00OJ}{}\label{cdp:eq:C.IR.4}
\]
Intersection with $E$ gives $a=d-1$, so the numerically
trivial canonical-boundary class descends.  This proves
the numerical assertion used in Step~1 without requiring
$h$ to factor through a ruling blowdown.

In the reverse sequence, effectivity forces every centre
to lie on the current divisor: a centre off it would create
coefficient $-1$, which persists on its strict transform.
Starting with the disjoint smooth coefficient-one curves
$H_0$, each permitted centre therefore has $d=1$.  Its
exceptional coefficient is zero, and its new boundary is
the strict transform of $H_0$.  Induction gives
\[
                         D_Z=H,\qquad V=0.
 \tag{C.IR.5}\hypertarget{tag-00OK}{}\label{cdp:eq:C.IR.5}
\]
In particular $H$ is a disjoint union of smooth elliptic
curves.  Coefficient-zero ruling exceptional curves have
not yet been discarded.

\item Let $E$ be a curve exceptional for
the minimal normalization resolution $\pi$ and not contained
in $H$.  Its negative self-intersection and adjunction give
\[
 2p_a(E)-2=E^2-H\cdot E<0.
\]
Thus $p_a(E)=0$, and $E$ is smooth rational.  Minimality
excludes $E^2=-1$.  Since $H\cdot E=2+E^2\ge0$, it follows
that $E^2=-2$ and $H\cdot E=0$.  A positive-coefficient
exceptional curve is one of the mutually disjoint elliptic
components of $H$.  Consequently every connected exceptional
fibre is disjoint from the strict conductor curves.

The normalization conductor ideal is the rank-one reflexive
sheaf $\omega_T\otimes\eta^*\omega_S^{-1}$ embedded in
$\mathcal O_T$ by finite duality.  Any point on its divisorial
support lies on a conductor curve, whose strict transform
must meet the resolution fibre over that point.  The fibres
just described therefore lie outside that support.  There
the conductor ideal equals $\mathcal O_T$, by reflexivity,
and $\eta$ is an isomorphism.  A nontrivial such resolution
fibre would give an isolated normal hypersurface singularity
of $S$.  Proposition~\SeriesRef{cdp:st:A.5.17} excludes it in the
actual torus smoothing, with its unit and full multiplicity
retained.  Hence $\pi$ is the identity and $T=Z$ is smooth.

\item The conductor ideal of $T$ now
defines the reduced smooth divisor $H$.  Its contraction to
$\mathcal O_S$ is radical, since the quotient embeds into
$\mathcal O_H$.  Let $G$ be this reduced conductor image.
At a generic transverse plane curve, conductor exponent one
on a normalization branch forces that branch to be smooth
and to meet exactly one other smooth branch transversely:
the branchwise conductor formula is twice the intrinsic
delta invariant plus the sum of its intersection
multiplicities with the other branches.  This is the
formula used in \SeriesRef{cdp:st:C4e}.  Thus $H\to G$ has total
generic degree two over each component.

The sheet involution extends to the smooth compact curves
$H$ over the normalization of $G$; a degree-two map of
smooth curves has its holomorphic deck involution, including
the exchange of two degree-one components.  The involution
preserves the finite fibre over every point of $G$, since
its composites to $G$ agree on a dense open subset and
therefore everywhere.

\item We prove the local ring comparison
at every point of $G$.  Let $A$ be the reduced hypersurface
ring, let $B$ be its finite normalization ring, and let $I$
be the conductor.  The ring $B$ is a product of regular
two-dimensional local rings.  Define
\[
 A'=\{b\in B:b\bmod I\in(B/I)^\iota\},
\]
where $\iota$ is the actual involution of the preceding
step.  There is an inclusion $A\subset A'\subset B$.
The two first rings agree in codimension one by the
ordinary-node calculation, so $A'/A$ has finite length.
The module $A'$ is finite and torsion-free over $A$.
Local cohomology of $0\to A\to A'\to A'/A\to0$ injects
$A'/A$ into $H^1_{\mathfrak m}(A)$.  The latter vanishes
because a hypersurface surface ring has depth two.
It follows that $A'=A$, and therefore
$A/I=(B/I)^\iota$.

The image curve $G$ is consequently the quotient of a
smooth curve by its involution and is smooth.  The exact
unramified local pullback ring is
\[
 \{(a(x,z),b(y,z)):a(0,z)=b(0,z)\}
       \simeq\mathbb C\{x,y,z\}/(xy).
\]
At a ramification point it is
\[
 \{a(u,t):a(0,t)\in\mathbb C\{t^2\}\}
 =\mathbb C\{u,ut,t^2\}
 \simeq\mathbb C\{x,y,z\}/(x^2-y^2z),
 \tag{C.IR.6}\hypertarget{tag-00OL}{}\label{cdp:eq:C.IR.6}
\]
with $(x,y,z)=(ut,u,t^2)$.  Proposition~\SeriesRef{cdp:st:A.5.15}
excludes the latter local map, including a unit times the
$m$th power.  Therefore all conductor maps are unramified,
their image components are disjoint smooth elliptic curves,
and $S$ has ordinary double-curve normal crossings.

\item Apply the normal-crossings
comparison \CompendiumRef{II}{cdp:st:C.2.NC.1} to the actual full fibre
$mS$.  It gives
\[
 0=e(Y_t)=m\,e(T\setminus H)=m\,e(T),
 \tag{C.IR.9}\hypertarget{tag-00OM}{}\label{cdp:eq:C.IR.9}
\]
because $e(H)=0$.  A minimal ruled surface over the elliptic
curve $B$ has Euler number zero, while every point blowup
adds one.  Thus $\rho$ has no blowups and $T=Z$ is minimal.
Finally the actual normalization is an isomorphism off
$H$ and $G$, so $e(S)=e(T)-e(H)+e(G)=0$.

\end{enumerate}
\end{proof}



The elliptic ruled geometry in \ref{cdp:st:C.2.IRR.1} has a direct literature
comparison. Kamenova's Theorem 3.1 in \cite{KamenovaTori} treats
semistable Kahler degenerations with trivial total-space canonical
bundle; its elliptic ruled cases have elliptic double curves arranged
along a path or a cycle. Our finite normalization $T\to S$ specifies
the self-identification of the double curves in the irreducible case.
The proof above establishes this geometry from the actual conductor
and the full-fibre smoothing, retaining the multiplicity and proving
the necessary local pinch exclusion. We have not supplied a Kahler
form or a trivialization of the total-space canonical bundle needed
to apply that theorem directly to our tube. 




\begin{proposition}[Proved anticanonical support comparison]\hypertarget{tag-01F0}{}\label{cdp:st:C.2.VII.1}
\StatementRecord{cdp:st:C.2.VII.1}{proved}{01F0}{\SeriesRef{cdp:st:C.3.8}; \cite[Theorem~0.2, Proposition~2.29 and Theorem~2.38]{Dloussky2006}; \cite[\S1, pp.~284--285, and \S3, p.~296]{DOT2003}; \cite[Corollary~2.25 and Proposition~2.27]{Dloussky2006}; \cite[Introduction, p.~283]{DOT2003}; \cite[Lemma~2.22 and the following paragraph; Proposition~2.27]{Dloussky2006}}


Let $Z$ be a smooth compact surface of class VII, not
necessarily minimal, and let $D>0$ be an effective integral
divisor such that $K_Z+D$ is torsion in $\Pic(Z)$.
Write $R=D_{\mathrm{red}}$, and let $r(R)$ be its number
of rational irreducible components. Then
\[
 e(R)=r(R),\qquad r(Z)\le b_2(Z),
\tag{C.VII.1}\hypertarget{tag-015K}{}\label{cdp:eq:C.VII.1}
\]
where $r(Z)$ counts all rational curves of $Z$.
Every connected component of $R$ has arithmetic genus one.
In particular a connected component containing an elliptic
curve consists of that single elliptic curve with trees of
rational curves attached; it cannot contain a second
elliptic curve or a path joining two distinct points of the
first one.
\end{proposition}


\begin{proof}

\begin{enumerate}[label={$\langle1\rangle$\arabic*.},leftmargin=*]
\item We specify the established surface results used in this
comparison. For a minimal class VII surface with positive
second Betti number, existence of a numerically
anticanonical divisor implies existence of a global
spherical shell by
\cite[Theorem~0.2, Proposition~2.29 and Theorem~2.38]{Dloussky2006}.
The reduced support is one rational cycle with attached
trees, one or two rational cycles, or, in the parabolic
Inoue case, a disjoint rational cycle and a smooth elliptic
curve. The intermediate graph and its actual curve
configuration are described in
\cite[\S1, pp.~284--285, and \S3, p.~296]{DOT2003}.
The fact that no connected piece of the anticanonical
support is an isolated tree is
\cite[Corollary~2.25 and Proposition~2.27]{Dloussky2006};
Lemma~2.22 there shows that a curve meeting that support
belongs to it. These statements apply to our divisor with
index one, since its twisting line is torsion and therefore
flat. They give, on the minimal model, arithmetic genus
one for each connected component and Euler number equal
to the number of its rational components. The rational
curve count is at most $b_2$, as stated and used in
\cite[Introduction, p.~283]{DOT2003}.
These are consequences of the proved theorem about
anticanonical divisors; no classification assertion about
class VII surfaces without such a divisor is used.

\item There is a coefficient issue in the parabolic Inoue case
which the torsion hypothesis resolves explicitly. Write
$E$ for its elliptic curve and $C$ for its reduced rational
cycle. The exact canonical relation is
$K_{Z_0}+E+C=0$; the curves $E$ and $C$ are disjoint,
$E^2=-b_2(Z_0)$, and the intersection matrix of the
components of $C$ has kernel generated by their sum
\cite[Lemma~2.22 and the following paragraph; Proposition~2.27]{Dloussky2006}.
Consequently an effective integral divisor with
$c_1(K_{Z_0}+D_0)_{\mathbb R}=0$ has the form
$D_0=E+kC$, $k\ge0$: intersection with $E$ fixes its
coefficient at one, and the cycle intersection matrix
makes all cycle coefficients equal. Our stronger
torsion relation then gives
$K_{Z_0}+D_0=(k-1)C$ torsion. Choose $n>0$ trivializing $\OO((k-1)C)^{\otimes n}$.
A trivialization compared with the canonical meromorphic
section gives a meromorphic function with divisor
$n(k-1)C$. For $k>1$ it has no poles and a nonzero
effective zero divisor, impossible on the compact
connected surface. For $k=0$ its reciprocal has that
property. Hence $k=1$. In particular both connected
pieces of the stated reduced support really occur.
This deduction uses the torsion relation, not merely
the vanishing of its real first Chern class.

\item For a minimal model with $b_2=0$, the nonzero effective
divisor excludes the curve-free Inoue surfaces. The
Hopf alternative has only smooth elliptic curves, which
are pairwise disjoint. This follows as well from the
explicit contraction description used in \SeriesRef{cdp:st:C.3.8} and its
finite unramified quotient: the preimages of two distinct
curves are disjoint elliptic curves, and the quotient
is locally biholomorphic. All assertions are then
immediate on that model.

\item We verify the transfer through the actual blowdowns.
Let $\sigma:Z\to Z_0$ be a sequence of point blowups
from a minimal model. A torsion line is represented by
a finite unitary character of the fundamental group,
and a point blowup induces an isomorphism on fundamental
groups. Thus the torsion line $K_Z+D$ descends to $Z_0$.
Pushing forward a meromorphic canonical section with pole
divisor $D$, or its torsion-trivializing power, gives an
effective divisor $D_0$ with $K_{Z_0}+D_0$ torsion.
It is nonzero, since a torsion canonical bundle would
give Kodaira dimension zero. At each individual blowup
the canonical divisor formula and the pushforward section
give the exact divisor equality
\[
 D_{\mathrm{new}}=\sigma^*D_{\mathrm{old}}-E.
\tag{C.VII.2}\hypertarget{tag-015O}{}\label{cdp:eq:C.VII.2}
\]
The coefficient of $E$ is
$\operatorname{mult}_pD_{\mathrm{old}}-1$.
Effectivity forces the centre to lie on $D_{\mathrm{old}}$.
If this coefficient is positive, the reduced support
is its entire reduced inverse image: a smooth point
acquires a rational leaf, or an ordinary intersection
is subdivided by the exceptional curve. These operations
preserve arithmetic genus and increase both the Euler
number and the rational-component count by one.
If the coefficient is zero, the centre is a smooth point
of a multiplicity-one component and the new support is
its strict transform, isomorphic to the old support.
Thus all the stated support properties transfer in the
required direction. Finally each blowup increases $b_2$
by one and adds precisely its exceptional rational
curve; every other rational curve is the strict transform
of one on the preceding surface. The curve-count
inequality in \eqref{cdp:eq:C.VII.1} follows.

\end{enumerate}
\end{proof}






\begin{proposition}[Proved global conductor Euler bound]\hypertarget{tag-01F1}{}\label{cdp:st:C.2.VII.2}
\StatementRecord{cdp:st:C.2.VII.2}{proved}{01F1}{\SeriesRef{cdp:st:A.5.12}, \SeriesRef{cdp:st:A.5.3}, \SeriesRef{cdp:st:C.2.VII.1}, \SeriesRef{cdp:eq:C.VII.5}, \SeriesRef{cdp:eq:C.VII.6}}


Let $S$ be a compact irreducible reduced hypersurface in
a smooth complex threefold, with $\omega_S$ torsion and
$\chi(S,\OO_S)=0$. Let $\eta:T\to S$ be normalization,
let $\Gamma$ be its conductor image scheme, and put
$C=\Gamma_{\mathrm{red}}$. Then
\[
 \chi(C,\OO_C)\ge-1.
\tag{C.VII.3}\hypertarget{tag-015P}{}\label{cdp:eq:C.VII.3}
\]
Suppose in addition that the minimal resolution
$\pi:Z\to T$ is a class VII surface. Then
\[
 e(C)\ge0.
\tag{C.VII.4}\hypertarget{tag-015Q}{}\label{cdp:eq:C.VII.4}
\]
The conductor scheme may be nonreduced and its generic
transverse exponent is unrestricted.
\end{proposition}


\begin{proof}

\begin{enumerate}[label={$\langle1\rangle$\arabic*.},leftmargin=*]
\item If the conductor is empty, both conclusions are immediate.
We therefore assume that it is nonempty.
Put $L=\omega_S$ and
$\varepsilon=h^0(S,L)\in\{0,1\}$.
A nonzero section of a torsion line on a compact
connected reduced space trivializes it, by taking a
torsion-trivializing power. Serre duality and
$\chi(\OO_S)=0$ therefore give
$h^1(S,\OO_S)=1+\varepsilon$, with
$\varepsilon=1$ exactly when $L=\OO_S$.
The normalization quotient
$\mathcal Q=\eta_*\OO_T/\OO_S$ has
$H^0(\mathcal Q)\hookrightarrow H^1(S,\OO_S)$,
because $S$ and $T$ have the same constant holomorphic
functions. The actual conductor duality of \SeriesRef{cdp:st:A.5.12} is
\[
 \mathcal Q=\omega_\Gamma\otimes L^{-1}|_\Gamma .
\]
Serre duality on the Cohen--Macaulay conductor curve
gives $h^1(\Gamma,L|_\Gamma)\le1+\varepsilon$.
The restriction onto $C$ is a quotient of coherent
sheaves on a one-dimensional space. Its kernel has
no second cohomology, so
$h^1(C,L|_C)\le1+\varepsilon$.
Since $L|_C$ is torsion, it has degree zero on every
irreducible component and
$\chi(C,L|_C)=\chi(C,\OO_C)$.
If $\varepsilon=0$, Riemann--Roch gives
$\chi(\OO_C)=h^0(C,L|_C)-h^1(C,L|_C)\ge-1$.
If $\varepsilon=1$, then $L=\OO_S$ and
$h^0(C,\OO_C)\ge1$, while $h^1(C,\OO_C)\le2$.
This proves \eqref{cdp:eq:C.VII.3}, including all nilpotents of
the original conductor scheme in the comparison.

\item Now assume that $Z$ is class VII. \SeriesRef{cdp:st:A.5.3} gives a
nonzero effective integral divisor $D$ with
$K_Z+D=(\eta\pi)^*\omega_S$.
Each curve in the reduced strict conductor on $Z$
is rational or elliptic by \ref{cdp:st:C.2.VII.1}. Its finite image
in $C$ has normalization of genus at most one, by
Riemann--Hurwitz. Let $n$ be the number of irreducible
components of $C$, let $g_i\in\{0,1\}$ be the genera
of their normalizations, and let $r_x$ and $\delta_x$
be the branch number and delta invariant at a singular
point of $C$. The exact normalization sequence and
the topological normalization map give
\[
 \chi(\OO_C)=\sum_i(1-g_i)-\sum_x\delta_x,\qquad
 e(C)=\sum_i(2-2g_i)-\sum_x(r_x-1).
\tag{C.VII.5}\hypertarget{tag-015R}{}\label{cdp:eq:C.VII.5}
\]
For every reduced curve germ,
$\delta_x\ge r_x-1$: its normalization permits
independent constants on the branches, whereas a
function on the original germ has one common constant.
If $n_0$ is the number of rational-normalization
components, \eqref{cdp:eq:C.VII.3}--\eqref{cdp:eq:C.VII.5} imply
$e(C)\ge n_0-1$.

\item The only possible negative value is consequently
$e(C)=-1$, with every normalized component elliptic.
Then \eqref{cdp:eq:C.VII.3} says $\sum\delta_x\le1$.
Negativity requires two branches, so there is exactly
one singular point, it has two branches and
$\delta=1$, and the local curve ring is
\[
 \{(a(u),b(v))\in\C\{u\}\oplus\C\{v\}:a(0)=b(0)\}.
\tag{C.VII.6}\hypertarget{tag-015S}{}\label{cdp:eq:C.VII.6}
\]
Indeed that subring already has colength one in the
normalization, and the curve ring is contained in it
with the same colength. Thus the point is an ordinary
node. Its two branches come either from distinct
elliptic normalizations or from two distinct points
of one elliptic normalization.

\item We show that this node is incompatible with the
specified hypersurface and resolution maps. At least
one normalization point $t\in T$ over the node receives
preimages of both local branches of $C$.
Otherwise partition the irreducible local branches
of the hypersurface germ $S$ according to the one
conductor branch their normalization receives. Both
parts are nonempty. Two hypersurface branches from
opposite parts intersect in a curve: in the regular
three-dimensional local ambient ring, the ideal of
their two distinct equations has height two.
That intersection curve is contained in the conductor
image and therefore is one of its two branches.
Both surface branches must receive it, contrary to
the partition. With only one local surface branch
the assertion is immediate.

\item Every $\pi$-exceptional curve not contained in
$D_{\mathrm{red}}$ is disjoint from it. To see this,
write $E$ for such a curve. Minimality of the
resolution and adjunction give $K_Z\cdot E\ge0$,
while $D\cdot E\ge0$, and
$(K_Z+D)\cdot E=0$ forces equality. In particular
$D\cdot E=0$. A resolution fibre over a point of
the conductor must meet its strict transform, by
properness. Since its fibre is connected, this
disjointness proves
$\pi^{-1}(t)\subset D_{\mathrm{red}}$ for the
normalization point just found.

\item The strict curves over the two elliptic-normalized
branches are elliptic: a rational curve cannot map
nonconstantly onto an elliptic curve. They meet the
connected fibre $\pi^{-1}(t)$. If the curves are
distinct, one connected component of $D_{\mathrm{red}}$
therefore contains two elliptic curves. If they are
the same elliptic curve, their meeting points with
the fibre are distinct. This follows from the exact
factorization of a map from a normal curve through
the normalization of its image: one point cannot
map to both points over the node in \eqref{cdp:eq:C.VII.6}.
The connected resolution fibre consequently gives
a path in $D_{\mathrm{red}}$ joining these distinct
points. Both possibilities contradict C.2.VII.1:
the first has two positive-genus components, and
the second adds a graph cycle to an elliptic
component, giving arithmetic genus at least two.
This excludes the only negative possibility and
proves \eqref{cdp:eq:C.VII.4}.

\end{enumerate}
\end{proof}






\begin{proposition}[Proved sign for class VII normalization]\hypertarget{tag-01F2}{}\label{cdp:st:C.2.VII.3}
\StatementRecord{cdp:st:C.2.VII.3}{proved}{01F2}{\SeriesRef{cdp:st:A.5.12}, \SeriesRef{cdp:st:A.5.3}, \SeriesRef{cdp:st:C.2.VII.1}, \SeriesRef{cdp:st:C.2.VII.2}}


Let $f:Y\to B$ be a proper connected-fibre holomorphic
map from a smooth compact complex threefold with
$b_2(Y)=0$ to a smooth curve, whose general fibre is a
complex two-torus. Let $mS$ be a full fibre with
irreducible reduced nonnormal support, and suppose
that the minimal resolution $Z$ of its normalization
is class VII. Then
\[
 e(S)\ge0.
\tag{C.VII.7}\hypertarget{tag-015T}{}\label{cdp:eq:C.VII.7}
\]
No restriction on the singularity type or conductor
multiplicity of $S$ is made.
\end{proposition}


\begin{proof}

\begin{enumerate}[label={$\langle1\rangle$\arabic*.},leftmargin=*]
\item The canonical evaluation and adjunction comparison
(A.CD.1) give $\omega_S^{\otimes m}=\OO_S$.
The divisor sequence and Hirzebruch--Riemann--Roch
give $\chi(\OO_S)=0$, because $c_1(\OO_Y(S))=0$
rationally. These arguments require only $b_2(Y)=0$;
the integral condition and Picard-variation assumption
in the later part of \SeriesRef{cdp:st:A.5.12} are not needed here.
Let $\eta:T\to S$, $\pi:Z\to T$, and
$D>0$ be the maps and divisor of \SeriesRef{cdp:st:A.5.3}. Put
$C=\Gamma_{\mathrm{red}}$ and
$A=(\eta^{-1}C)_{\mathrm{red}}\subset T$.
Thus $K_Z+D$ is torsion, and \ref{cdp:st:C.2.VII.1}--\ref{cdp:st:C.2.VII.2}
apply.

\item We compute $e(T\setminus A)$ through the actual
resolution. As in \ref{cdp:st:C.2.VII.2}, an exceptional curve
$E$ with coefficient zero in $D$ has
$K_Z\cdot E=D\cdot E=0$.
Adjunction and negative self-intersection then give
$p_a(E)=0$ and $E^2=-2$.
These curves are smooth rational and disjoint
from $D_{\mathrm{red}}$.
Their negative-definite intersection matrix makes
each connected configuration a tree: a repeated
edge or a cycle would give a nonnegative square
for a nonzero vector in the matrix. The intersections
are transverse, since their local intersection
multiplicities are one; three curves cannot meet
at a point, since that would give a triangle.
Connectedness of resolution fibres shows that a
fibre either consists entirely of these curves or
lies entirely in $D_{\mathrm{red}}$.
The zero-coefficient fibres lie outside $A$, since
a fibre over $A$ must meet the strict conductor.

\item Let $a$ be the total number of zero-coefficient
exceptional curves, and let $b$ be the number of
exceptional fibres lying in $D_{\mathrm{red}}$
whose images in $T$ lie outside $A$. On
$Z\setminus D_{\mathrm{red}}$ the former fibres
are still present and the latter are absent.
A tree with $a_j$ rational curves has Euler
number $a_j+1$, so replacing it by its image
point subtracts $a_j$. Adding back the image
point of an absent positive-coefficient fibre
adds one. Away from these finitely many fibres,
$\pi$ is the required isomorphism of complements.
Additivity of compactly supported Euler
characteristic therefore gives the exact equality
\[
 \begin{aligned}
 e(T\setminus A)
   &=e(Z)-e(D_{\mathrm{red}})-a+b\\
   &=b_2(Z)-r(D_{\mathrm{red}})-a+b
   \ \ge0 .
 \end{aligned}
\tag{C.VII.8}\hypertarget{tag-015U}{}\label{cdp:eq:C.VII.8}
\]
The inequality follows because the rational curves
counted by $r(D_{\mathrm{red}})$ and by $a$ are
distinct, and their total is at most $b_2(Z)$ by
\ref{cdp:st:C.2.VII.1}.
Finally $\eta$ is a biholomorphism
$T\setminus A\to S\setminus C$. Hence
\[
 e(S)=e(T\setminus A)+e(C)\ge0
\]
by \ref{cdp:st:C.2.VII.2}. This proves the assertion for the
original reduced support and all full multiplicities.

\end{enumerate}
\end{proof}






\begin{proposition}[Proved equality comparison]\hypertarget{tag-01F3}{}\label{cdp:st:C.2.VII.4}
\StatementRecord{cdp:st:C.2.VII.4}{proved}{01F3}{\SeriesRef{cdp:st:C.2.VII.3}, \SeriesRef{cdp:st:C.2.VII.2}, \SeriesRef{cdp:st:A.5.12}, \SeriesRef{cdp:st:C.2.VII.1}, \SeriesRef{cdp:st:C.2.INT.3}}


Retain the irreducible class VII situation of
\ref{cdp:st:C.2.VII.3}. Let $C=\Gamma_{\rm red}$, let $n_0$
be the number of its rational-normalized components,
and put $v(C)=\sum_x(s_x-1)$, where $s_x$ is the
number of branches at a singular point of $C$.
On the minimal resolution $Z$ let $D$ be the
resolution conductor, let $r$ count the rational
curves in $D_{\rm red}$, and let $a$ count the
zero-coefficient exceptional curves lying over
points outside the conductor. Let $b$ count the
positive-coefficient exceptional fibres over
such points, as in \ref{cdp:st:C.2.VII.3}. Then $e(S)=0$
is equivalent to the following explicit
normalization and curve data:
\[
 b=0,\qquad b_2(Z)=r+a,
 \qquad
 \begin{cases}
 n_0=0, v(C)=0,\quad\text{or}\\
 n_0=1, v(C)=2.
 \end{cases}
 \tag{C.VII.9}\hypertarget{tag-015V}{}\label{cdp:eq:C.VII.9}
\]
In the first curve alternative all components
have elliptic normalization and the sum of
their delta invariants is at most one. In the
second the total delta invariant is exactly
two, and every singular point satisfies
$\delta_x=s_x-1$.

If the actual homomorphism
$\Pic^0(Y)\to\Pic^0(Z)$ is nonconstant,
there is the sharper geometric equivalence:
$e(S)=0$ exactly when its normalization is a
smooth minimal Hopf surface and the reduced
conductor image and inverse image are disjoint
unions of smooth elliptic curves, with finite
unramified maps between them.
\end{proposition}


\begin{proof}

\begin{enumerate}[label={$\langle1\rangle$\arabic*.},leftmargin=*]
\item The exact comparison already proved in
\ref{cdp:st:C.2.VII.3} is
\[
 e(S)=b_2(Z)-r-a+b+e(C),
 \qquad b_2(Z)-r-a\ge0,\quad b\ge0,\quad e(C)\ge0.
 \tag{C.VII.10}\hypertarget{tag-015L}{}\label{cdp:eq:C.VII.10}
\]
Every normalized component of $C$ has genus
zero or one. Its topological normalization
formula is $e(C)=2n_0-v(C)$. The stronger
bound in \ref{cdp:st:C.2.VII.2} is $e(C)\ge n_0-1$.
Thus $e(C)=0$ permits precisely $n_0=0,1$,
with $v(C)=0,2$, respectively. This proves
both directions of \eqref{cdp:eq:C.VII.9}, using actual
branch numbers rather than an assumption
about a local singularity type.
The same proof gives
$\chi(\OO_C)=n_0-\sum\delta_x\ge-1$.
For $n_0=0$ this bounds the total delta
invariant by one. For $n_0=1$, the value
$v(C)=2$ and $\delta_x\ge s_x-1$ force
$\sum\delta_x=2$ and equality at every point.
These are exact consequences for fibres which
satisfy the criterion; no separate claim
that all such abstract curve patterns admit
a torus smoothing is used.

\item Now suppose that the Picard pullback is
nonconstant. The argument of \SeriesRef{cdp:st:A.5.12} applies
to $\Pic^0(Y)$ even with only rational
$H^2(Y,\Q)=0$. Namely its nonzero differential
is the actual composite
\[
 H^1(Y,\OO_Y)\longrightarrow H^1(S,\OO_S)
 \longrightarrow H^1(T,\OO_T)
 \longrightarrow H^1(Z,\OO_Z).
 \tag{C.VII.11}\hypertarget{tag-015M}{}\label{cdp:eq:C.VII.11}
\]
The middle arrow consequently has positive
rank. Put $\varepsilon=h^0(S,\omega_S)$;
the torsion canonical relation makes
$\varepsilon\in\{0,1\}$ and
$h^1(S,\OO_S)=1+\varepsilon$.
The normalization sequence then gives
$h^0(\mathcal Q)\le\varepsilon$.
Conductor duality and restriction onto $C$
therefore give
\[
 h^1(C,\omega_S|_C)\le\varepsilon.
 \tag{C.VII.12}\hypertarget{tag-015N}{}\label{cdp:eq:C.VII.12}
\]
If $\varepsilon=0$, degree-zero
Riemann--Roch gives $\chi(\OO_C)\ge0$.
If $\varepsilon=1$, the torsion line
$\omega_S$ is trivial, and
$h^0(C,\OO_C)\ge1$ gives the same inequality.
The exact formula (C.INT.26) now implies
that $e(C)=0$ holds precisely when $C$ is
a disjoint union of smooth elliptic curves.

\item Assume $e(S)=0$. Every rational curve in
$D_{\rm red}$ must be resolution-exceptional,
because its image lies in $C$ and cannot
dominate an elliptic curve. A minimal class
VII model with positive second Betti number
and a torsion anticanonical divisor is one
of the proved configurations used in
\ref{cdp:st:C.2.VII.1}. An intermediate or
Inoue--Hirzebruch model has no elliptic
curves, and therefore could not dominate
the nonempty elliptic conductor. A
parabolic Inoue model has a rational cycle
of square zero and exactly $b_2$ rational
curves. Its point blowups therefore still
have exactly $b_2(Z)$ rational curves.
The equality $b_2(Z)=r+a$ makes every one
of them resolution-exceptional, including
any zero-coefficient blowup curves. Thus
the full total transform of the minimal
rational cycle, whose square remains zero,
would be exceptional. This contradicts
negative definiteness of a resolution's
exceptional intersection matrix.
Consequently the minimal model is Hopf.
All rational curves on $Z$ are now the
exceptional curves of its point blowups,
and their number is $b_2(Z)$.
The equality $b_2(Z)=r+a$ makes all of
them resolution-exceptional. As in the
last paragraph of \ref{cdp:st:C.2.INT.3}, a nonempty
blowup sequence would have a last
$(-1)$ curve, contradicting minimality
of the resolution. Thus the normalization
is already a smooth minimal Hopf surface.
Its reduced conductor curves are elliptic;
their finite maps to the smooth elliptic
curves of $C$ are unramified by
Riemann--Hurwitz.
Conversely those normalization and
conductor data have Euler number zero
by the exact normalization formula.
This proves the sharper equivalence.

\end{enumerate}
\end{proof}






\begin{proposition}[Euler signs for two of the rational boundary classes]\hypertarget{tag-01CD}{}\label{cdp:st:C.2.RAT.1}
\StatementRecord{cdp:st:C.2.RAT.1}{proved}{01CD}{\SeriesRef{cdp:st:A.5.16}, \SeriesRef{cdp:eq:C.RA.7}, \SeriesRef{cdp:eq:C.RA.8}, \SeriesRef{cdp:eq:C.RA.9}, \SeriesRef{cdp:eq:C.RA.10}, \SeriesRef{cdp:st:A.5.5}, \SeriesRef{cdp:st:A.5.15}, \SeriesRef{cdp:st:A.5.17}, \SeriesRef{cdp:eq:C.RA.14}, \SeriesRef{cdp:eq:C.RA.15}, \SeriesRef{cdp:eq:C.RA.6}, \SeriesRef{cdp:eq:C.RA.13}, \SeriesRef{cdp:eq:C.RA.2}}

Let $f:Y\to B$ be a proper holomorphic map from a smooth compact
complex threefold with $b_2(Y)=0$, with connected fibres and smooth
general fibre a complex two-torus. Let $mS$ be a full fibre with
irreducible nonnormal reduced support. Write
$h=\eta\pi:Z\to T\to S$ for the minimal resolution of its
normalization and suppose that $Z$ is rational. Use the effective
divisor $K_Z+D=0$ of \SeriesRef{cdp:st:A.5.16}.

If a contraction to $\mathbf P^2$ sends $D$ to three concurrent
reduced lines, then $e(S)\ge1$. More precisely,
$e(S)=-1+N_0+e(\Gamma_{\rm red})\ge1+N_0$, with $N_0$ as in
(A.RT.12).

Consider also a specified morphism $\rho:Z\to\mathbb F_n$,
$n\ge1$, which is a sequence of point blowups in reverse and has
\[
 \rho_*D=2C_0+\sum_{j=1}^r a_jF_j,\qquad
 C_0^2=-n,\qquad
 a_j\ge1,\quad \sum_j a_j=n+2,                  \tag{C.RA.1}\hypertarget{tag-00WQ}{}\label{cdp:eq:C.RA.1}
\]
where the $F_j$ are distinct ruling fibres. Let $C$ be the
strict transform of $C_0$. Suppose $h(C)$ is a curve and the
generic transverse singularity there is a cusp $A_2$.
Then $r\le3$ and $e(S)\ge1$. If $h(C)$ is smooth, then
$r\le2$. These assertions allow singular normalization,
arbitrary vertical conductor multiplicities and a nonreduced
conductor-image scheme.
\end{proposition}


\begin{proof}

\begin{enumerate}[label={$\langle1\rangle$\arabic*.},leftmargin=*]
\item Put $G=\Gamma_{\rm red}$. \SeriesRef{cdp:st:A.5.16} has already proved
\[
 e(S)=e(Z\setminus D_{\rm red})+e(G).             \tag{C.RA.2}\hypertarget{tag-00X1}{}\label{cdp:eq:C.RA.2}
\]
In either boundary class under consideration, $D_{\rm red}$
is simply connected, all its components are rational, and
(A.RT.18)--(A.RT.20) give
\[
 \pi_1(S)=\pi_1(G)\in\{1,\mathbb Z\}.
 \tag{C.RA.3}\hypertarget{tag-00XC}{}\label{cdp:eq:C.RA.3}
\]
Write $n_G$ for the number of irreducible components of $G$
and $\beta=b_1(G)\le1$. The normalization maps from their
rational normalizations are homeomorphisms except for the
identification of finitely many points. Consequently
\[
                   e(G)=n_G+1-\beta\ge n_G\ge1.
 \tag{C.RA.4}\hypertarget{tag-00XH}{}\label{cdp:eq:C.RA.4}
\]
The exact two-branch incidence argument (A.RT.21) strengthens
this to $e(G)\ge2$ in the present boundary classes.
For three concurrent lines, (A.RT.12)--(A.RT.13) therefore
prove the first assertion.

\item We prove the cusp assertion. Put $H=h(C)_{\rm red}$.
The transverse $A_2$ curve has a single normalized branch.
Thus $C\to H$ is the normalization map of degree one, and
no other strict conductor component maps onto $H$.
In particular $H^\nu=C=\mathbf P^1$. The restriction
$H^1(G,\OO_G)\to H^1(H,\OO_H)$ is surjective, since its
kernel has one-dimensional support and no second cohomology.
\SeriesRef{cdp:st:A.5.16} therefore gives
\[
              \delta:=p_a(H)\le2.              \tag{C.RA.5}\hypertarget{tag-00XI}{}\label{cdp:eq:C.RA.5}
\]
Since $H^\nu$ is rational, $\delta$ is the sum of its local
delta invariants. Write $R$ for the union of the other irreducible components of
$G$. The exact constant-sheaf sequence
\[
 0\longrightarrow\mathbb Q_G\longrightarrow
 \mathbb Q_H\oplus\mathbb Q_R\longrightarrow
 \mathbb Q_{H\cap R}\longrightarrow0
\]
gives a surjection
$H^1(G;\mathbb Q)\to H^1(H;\mathbb Q)\oplus H^1(R;\mathbb Q)$,
since $H\cap R$ is finite. Hence $b_1(H)\le b_1(G)\le1$.
If $s$ denotes the number of points of $C$ above singular
points of $H$, normalization of the reduced curve gives
\[
 s=\sum_{p\in\operatorname{Sing}H}r_p
   =b_1(H)+\#\operatorname{Sing}H\le1+\delta
 \quad(\delta>0),\qquad s=0\quad(\delta=0).       \tag{C.RA.6}\hypertarget{tag-00XJ}{}\label{cdp:eq:C.RA.6}
\]

\item We need the precise low-delta curve analysis, including its
two possible non-lci cases. A reduced curve germ of delta
invariant at most two has at most three branches.
Three branches require delta two, and the ring is exactly
the subring of its normalization consisting of triples
with equal constant value. This follows by equality of
colengths, both equal to two. Its topological branch
identification alone contributes two loops to a globally
irreducible curve, so $b_1(H)\le1$ excludes it here.

\item For two branches, the exact sequence for their union gives
$\delta=\delta_1+\delta_2+\ell(\OO_{H_1\cap H_2})$.
Two smooth branches therefore give a node, or, at delta
two, the fibre product matching their constant and linear
terms. The latter has the plane tacnode ring
$\mathbb C\{x,y\}/(y^2-x^4)$.
The only remaining two-branch possibility is a cusp and
a transverse smooth branch, whose ring is
$\mathbb C\{t^2,t^3\}\times_{\mathbb C}\mathbb C\{u\}$.
In embedding dimension three its ideal can be written
\[
                  I=(xz,yz,y^2-x^3).            \tag{C.RA.7}\hypertarget{tag-00XK}{}\label{cdp:eq:C.RA.7}
\]
For one branch, the delta invariant equals the number of
gaps of its valuation semigroup. With at most two gaps,
the singular possibilities are
$\langle2,3\rangle$, $\langle2,5\rangle$, and
$\langle3,4,5\rangle$. The first two are plane cusps.
For the last, the conductor is $t^3\mathbb C\{t\}$,
and colength two makes the ring
$\mathbb C+t^3\mathbb C\{t\}$ exactly. Its ideal is
\[
 I=(g_1,g_2,g_3)
   =(y^2-xz,\ yz-x^3,\ z^2-x^2y),\qquad
 (x,y,z)=(t^3,t^4,t^5).                         \tag{C.RA.8}\hypertarget{tag-00XL}{}\label{cdp:eq:C.RA.8}
\]
For completeness, the asserted plane form at semigroup
$\langle2,5\rangle$ follows by making an order-two
generator equal to $t^2$ by a change of parameter.
Its conductor is then $t^4\mathbb C\{t\}$, and the ring
is $\mathbb C+\mathbb Ct^2+t^4\mathbb C\{t\}$.
These conductor descriptions also follow in the completed
local rings by eliminating successive higher terms; finite
normalization and faithful flatness of completion give the
same conductor in the convergent rings. Thus there are no
unlisted cases. A germ of embedding dimension two is lci
in the smooth threefold: choose one regular transverse
ambient coordinate in its ideal and then its single plane
equation. We now exclude \eqref{cdp:eq:C.RA.7}--\eqref{cdp:eq:C.RA.8} for our cusp
image.

\item For \eqref{cdp:eq:C.RA.8}, put
\[
 F=x^5+xy^3+z^3-3x^2yz .
\]
The symbolic square is $I^{(2)}=I^2+(F)$, with
\[
 xF=g_2^2-g_1g_3,\quad
 yF=g_2g_3+xg_1^2,\quad
 zF=g_3^2+xg_1g_2.                              \tag{C.RA.9}\hypertarget{tag-00XM}{}\label{cdp:eq:C.RA.9}
\]
The saturation computation $I^2:x^\infty=I^2+(F)$,
including the complete exact Groebner certificate, is
given in \path{code/A_low_delta_cusp.py} and its output.
It computes the symbolic square because the curve is lci
away from the origin and $x$ vanishes on it only at the
origin. The same equality in the analytic local ring
follows by flatness of the convergent extension and the
displayed finite saturation.

\item Let $N$ be the dual of the torsion-free pullback to $C$ of
the ambient conormal module of $H$. At \eqref{cdp:eq:C.RA.8} the
normalized conormal basis and its dual lifts are
\[
\begin{gathered}
 a=dg_1,\qquad b=dg_2+t\,dg_1,\qquad
 dg_3=-2t^2a+tb,\\
 A=-t^{-3}\partial_z,\qquad
 B=(t^{-5}\partial_y+2t^{-4}\partial_z)/3 .
\end{gathered}                                  \tag{C.RA.10}\hypertarget{tag-00WR}{}\label{cdp:eq:C.RA.10}
\]
Their wedge is $a\wedge b=t^6\,\iota_{\gamma'}(dx\wedge
dy\wedge dz)$, where $\gamma=(t^3,t^4,t^5)$.
Thus the determinant formula is
$\det N=K_C\otimes h_C^*K_Y^{-1}(6p)$.
Since every ambient line has degree zero on every compact
curve when $b_2(Y)=0$, its degree here is four.

\item Write the local equation of $S$ as $cF+f_2$, $f_2\in I^2$.
Its normalized quadratic form is the sum of a regular
quadratic form and $c$ times
\[
 q_F=\frac3t a^2-\frac3{t^2}ab+\frac1{t^3}b^2,
 \qquad \det(q_F)=\frac3{4t^4}.                  \tag{C.RA.11}\hypertarget{tag-00WS}{}\label{cdp:eq:C.RA.11}
\]
If $c(0)\ne0$, the determinant has an uncancellable
pole of order four: a cross term with the regular form
has pole at most three. The generic cusp requires rank
one, so $c(0)=0$. Equations \eqref{cdp:eq:C.RA.9} then give
$cF\in I^2$, and the quadratic form is regular.

\item The intrinsic cubic on its radical has pole order at most
ten at $p$. Here is a direct bound for arbitrary coefficients,
not only for the polynomial $F$. In the lifts \eqref{cdp:eq:C.RA.10},
the differentials $dg_i$ are regular, the second derivatives
$D^2g_i(Au+Bv)^2$ have poles at most ten, and a derivative
of a holomorphic coefficient has pole at most five.
The formula for $D^3(b_{ij}g_i g_j)$ therefore has pole
at most ten. Terms containing an undifferentiated $g_i$
vanish on the curve. The exact derivatives are in the cited computation.
The radical cubic is independent of a tangent component
of the lift: differentiating the identically zero quadratic
form on its radical along the smooth part of the curve
gives this invariance, and equality extends meromorphically.

\item The same argument is shorter for \eqref{cdp:eq:C.RA.7}. Put
$g_1=xz$, $g_2=yz$, $g_3=y^2-x^3$. Then
$I^{(2)}=I^2+(z g_3)$ and
\[
 x(zg_3)=g_1g_3,\quad y(zg_3)=g_2g_3,\quad
 z(zg_3)=g_2^2-xg_1^2.
 \tag{C.RA.12}\hypertarget{tag-00WT}{}\label{cdp:eq:C.RA.12}
\]
On the cusp branch $(t^2,t^3,0)$ use $a=dg_1$, $b=dg_3$.
The quadratic form of $z g_3$ is $t^{-2}ab$, whose
determinant has pole four and cannot be cancelled by a
regular quadratic form. Thus its coefficient again
vanishes at the origin and the surface equation belongs
to $I^2$. The determinant corrections are four on the
cusp branch and two on the smooth branch $(0,0,t)$.
The total determinant degree is again four.
Dual lifts on the cusp branch are
$t^{-2}\partial_z$ and $\frac12t^{-3}\partial_y$;
on the smooth branch they are
$t^{-1}\partial_x,t^{-1}\partial_y$.
The same product derivative formula bounds the cubic
poles by six and two respectively, of total degree eight.

\item In either non-lci case let $R\subset N$ be the saturated
radical of the regular rank-one quadratic form and let
$Q=N/R$. Its nonzero square coefficient is a section of
$Q^{-2}\otimes M$, where $M=\OO_Y(S)|_C$ has degree
zero. Hence $\deg Q\le0$ and $\deg R=4-\deg Q\ge4$.
The nonzero intrinsic cubic of a generic $A_2$ is a
section of $R^{-3}\otimes M$ with poles of total degree
at most ten, or eight. It would give respectively
$3\deg R\le10$ or $3\deg R\le8$, both impossible.
This excludes the two non-lci cases, without assuming
that the conductor-image scheme is reduced.

\item We may consequently use the lci conormal bundle along $H$.
Let $\Xi$ be its normalization conductor divisor, of degree
$2\delta$. The determinant formula and the
conductor-cleared intrinsic cubic from \SeriesRef{cdp:st:A.5.5} give
\[
\begin{gathered}
 \det N=K_C(\Xi)\otimes h_C^*K_Y^{-1},\\
 q\in H^0(C,Q^{-2}\otimes M),\qquad
 k\in H^0(C,R^{-3}\otimes M(2\Xi)),\\
 0\ne q^3k^2\in
 H^0\bigl(C,(\det N)^{-6}\otimes M^5(4\Xi)\bigr),\\
 \deg\bigl((\det N)^{-6}\otimes M^5(4\Xi)\bigr)
                         =12-4\delta .
\end{gathered}                                  \tag{C.RA.13}\hypertarget{tag-00WU}{}\label{cdp:eq:C.RA.13}
\]
The section is nonzero because the generic transverse
singularity is a cusp.

\item We give the local comparison that strengthens the zero
count at a smooth point of $H$. A simple zero of $q$ at
which $k$ is nonzero is impossible in a torus smoothing.
Choose coordinates with $H=\{x=y=0\}$, radical direction
$\partial_x$, and coordinate $z$ along $H$.
The nonzero cubic permits Weierstrass preparation in $x$:
\[
 f_S=U\bigl(x^3+a_2(y,z)x^2+a_1(y,z)x+a_0(y,z)\bigr).
 \tag{C.RA.14}\hypertarget{tag-00WV}{}\label{cdp:eq:C.RA.14}
\]
The critical-curve equations and the rank-one Hessian
along it give $a_2\in(y)$ and $a_1,a_0\in(y^2)$.
After depressing the cube by $x'=x+a_2/3$, this is
\[
 f_S=U'\bigl(x'^3+y^2(B_1(y,z)x'+B_0(y,z))\bigr).
 \tag{C.RA.15}\hypertarget{tag-00WW}{}\label{cdp:eq:C.RA.15}
\]
The simple zero of the Hessian coefficient says
$\partial_zB_0(0,0)\ne0$. Thus
$z'=B_1(y,z)x'+B_0(y,z)$ is an actual holomorphic
coordinate with $x',y$, and the reduced equation is a
unit times $x'^3+y^2z'$. The full fibre equation is a
unit times its $m$th power. \SeriesRef{cdp:st:A.5.15} excludes precisely
this Hirzebruch germ, including that unit and multiplicity.

\item Every point of $C$ above the smooth part of $H$ where a
vertical conductor component meets it has $q=0$.
Indeed a nonzero Hessian would, by \SeriesRef{cdp:st:A.5.17}, give a constant
transverse cusp product with smooth normalization and pure
conductor $2C$; its minimal normalization resolution could
have no such vertical component. At a zero of $q$ the
cubic is holomorphic. If it also vanishes, the section
$q^3k^2$ has zero order at least five. If it does not
vanish, \eqref{cdp:eq:C.RA.14}--\eqref{cdp:eq:C.RA.15} force the zero of $q$ to
have order at least two, and the order is at least six.

\item Each of the $r$ original ruling fibres in \eqref{cdp:eq:C.RA.1} has
a positive-boundary total-transform component meeting $C$,
at a different point of $C$. This remains true at each
blowup: a blowup at its intersection with the coefficient-two
section inserts an exceptional curve of positive coefficient.
Using \eqref{cdp:eq:C.RA.6} and the degree in \eqref{cdp:eq:C.RA.13} now gives
\[
 \begin{array}{c|c|c|c}
 \delta& s\text{ at most}&
 \text{bad points above }H_{\rm reg}\text{ at most}&r\text{ at most}\\ \hline
 0&0&2&2\\
 1&2&1&3\\
 2&3&0&3 .
 \end{array}                                    \tag{C.RA.16}\hypertarget{tag-00WX}{}\label{cdp:eq:C.RA.16}
\]
This proves the asserted bound on $r$.

\item Finally (A.RT.12) and (A.RT.16) give
$e(Z\setminus D_{\rm red})=2-r+N_0$.
The forest argument of (A.RT.21) applies to this entire
star class, whether its minimal complement is negative
or not, and gives $e(G)\ge2$.
Together with $r\le3$, \eqref{cdp:eq:C.RA.2} now gives directly
\[
                 e(S)\ge4-r+N_0\ge1+N_0.
 \tag{C.RA.17}\hypertarget{tag-00WY}{}\label{cdp:eq:C.RA.17}
\]

\end{enumerate}
\end{proof}



The two classes in \ref{cdp:st:C.2.RAT.1} are proved without a restriction
on nilpotent conductor structure. The other horizontal
coefficient-two types and contraction of that horizontal
section are not included in this statement.




\begin{proposition}[A rational double cover of a smooth tacnode image]\hypertarget{tag-01CE}{}\label{cdp:st:C.2.RAT.2}
\StatementRecord{cdp:st:C.2.RAT.2}{proved}{01CE}{\SeriesRef{cdp:st:C.2.RAT.1}, \SeriesRef{cdp:st:A.5.17}, \SeriesRef{cdp:eq:C.RA.19}}

Retain the torus-fibre, rational-normalization and specified
ruled-boundary morphism of \ref{cdp:st:C.2.RAT.1}. Suppose that $C$ is not
contracted, that $H=h(C)_{\rm red}$ is a smooth rational curve,
that the generic transversal along $H$ is an ordinary tacnode,
and that $C\to H$ has degree two. Then $r\le2$ and
\[
                         e(S)\ge2+N_0.
\tag{C.RA.18}\hypertarget{tag-00WZ}{}\label{cdp:eq:C.RA.18}
\]
The assertion allows arbitrary vertical conductor and singular
normalization away from the generic points of $C$.
\end{proposition}


\begin{proof}

\begin{enumerate}[label={$\langle1\rangle$\arabic*.},leftmargin=*]
\item The two branches of a generic tacnode are both supplied by
$C\to H$; no other strict conductor component can dominate $H$.
Apply the exact normal-jet construction (A.T.5)--(A.T.7) on
$H\simeq\PP^1$. Here the curve conductor divisor $\Xi$ is zero.
The quotient $Q$ of its rank-one Hessian has degree $q_0\le0$.
The quadratic coefficient $a$ and the quartic invariant
$\Delta=ae-b^2/4$ are nonzero sections with
\[
 A:=\deg\operatorname{div}(a)=-2q_0,\qquad
 \deg\operatorname{div}(\Delta)=8+2q_0=8-A.
\tag{C.RA.19}\hypertarget{tag-00X0}{}\label{cdp:eq:C.RA.19}
\]
In particular $A$ is a nonnegative even integer. These are
sections of the actual lines $M\otimes Q^{-2}$ and
$(M\otimes Q^{-1}\otimes R^{-2})^2$, with
$\deg M=0$ and $\deg R=-2-q_0$.

\item The normalized conductor cover is obtained by adjoining
$\sqrt{-\Delta}$, by the quadratic equation for the ratio of
the first normal coordinate to the square of the radical
coordinate. Thus its branch divisor consists exactly of the
odd-order zeros of $\Delta$, as proved in (A.T.7). Since
$C\simeq\PP^1$ and the map has degree two, Riemann--Hurwitz
gives two distinct branch points on $H$. At each of them
$a$ vanishes: a nonzero Hessian would give the constant
transverse $A_3$ product by \SeriesRef{cdp:st:A.5.17} and hence an unramified
conductor cover. The coordinate comparison
(A.T.10)--(A.T.12) excludes a simple zero of $\Delta$ at
such a point. Its odd zero order is consequently at least
three. Equations \eqref{cdp:eq:C.RA.19} now give
\[
                        A\ge2,\qquad 8-A\ge6.
\tag{C.RA.20}\hypertarget{tag-00X2}{}\label{cdp:eq:C.RA.20}
\]
Therefore $A=2$. The two branch points are exactly the
zeros of $a$, each simple, and the zeros of $\Delta$ have
order three there, with no further zeros.

\item Every original fibre $F_j$ of (C.RA.1) gives a positive
vertical boundary component meeting $C$. The positivity
is preserved under blowup at its intersection with the
coefficient-two horizontal section, and these $r$
intersection points on $C$ are distinct. At any such point
the Hessian coefficient at its image must vanish. Otherwise
\SeriesRef{cdp:st:A.5.17} supplies the constant $A_3$ product, whose two
normalization branches are smooth and have no vertical
conductor; the minimal resolution $\pi$ is an isomorphism
near them. Each of the two zeros of $a$ is a branch point
of $C\to H$, and consequently has just one point above it
on $C$. It follows that $r\le2$.

\item Finally the exact complement calculation and (A.RT.21) give
\[
 e(S)=2-r+N_0+e(\Gamma_{\rm red})
                  \ge4-r+N_0\ge2+N_0,
\]
which proves \eqref{cdp:eq:C.RA.18}.

\end{enumerate}
\end{proof}






\begin{proposition}[Every uncontracted degree-one horizontal image has nonnegative Euler number]\hypertarget{tag-01CF}{}\label{cdp:st:C.2.RAT.3}
\StatementRecord{cdp:st:C.2.RAT.3}{proved}{01CF}{\SeriesRef{cdp:st:C.2.RAT.1}, \SeriesRef{cdp:eq:C.RA.22}, \SeriesRef{cdp:eq:C.RA.23}, \SeriesRef{cdp:eq:C.RA.24}}

Retain the torus-fibre, rational-normalization and ruled-boundary
morphism of \ref{cdp:st:C.2.RAT.1}. Suppose that $C$ is not contracted and
that its actual map to the normalization of
$H=h(C)_{\rm red}$ has degree one. Let $r_{\rm dom}$ be the
number of original ruling fibres $F_j$ whose positive vertical
boundary contains a strict conductor component mapping onto
$H$. Then
\[
       0\le r_{\rm dom}\le2,\qquad
       e(S)\ge2-r_{\rm dom}\ge0.
\tag{C.RA.21}\hypertarget{tag-00X3}{}\label{cdp:eq:C.RA.21}
\]
There is no assumption on the generic transverse type, on the
local complete intersection property of $H$, or on the
nilpotent structure of the conductor image.
\end{proposition}


\begin{proof}

\begin{enumerate}[label={$\langle1\rangle$\arabic*.},leftmargin=*]
\item The conductor coefficient on $C$ is two, because $C$ is
the strict transform of the coefficient-two section in
(C.RA.1) and is not $\pi$-exceptional. At a general point
of $H$, its degree-one map supplies one branch of the reduced
transverse plane curve. If that branch has intrinsic
conductor exponent $c_i^{\rm int}$, its exponent in the
whole plane curve is
\[
             2=c_i^{\rm int}+\sum_{j\ne i}I(C_i,C_j).
\tag{C.RA.22}\hypertarget{tag-00X4}{}\label{cdp:eq:C.RA.22}
\]
Here $I(C_i,C_j)$ is the positive intersection multiplicity
of two different plane branches. To verify the displayed
comparison, write the reduced equation as $u\prod f_j$.
On the normalization of the $i$th branch the residue
description of its dualizing generator gains the factor
$\prod_{j\ne i}f_j|_{C_i}$; its valuation is the sum in
\eqref{cdp:eq:C.RA.22}. The remaining factor is exactly the conductor
of the individual branch. This gives the actual inclusion
of its conductor ideal into its normalization, with the
stated exponent. Thus there are at most two other generic
branches. Each vertical strict component dominating $H$
supplies at least one such branch, counted with its covering
degree. This proves $r_{\rm dom}\le2$.

\item For an original fibre $F_j$, let $B_j$ be the positive
vertical boundary in its total transform. Its reduced
support is connected and meets $C$ at a single point.
These assertions follow at each point blowup: a centre
at the intersection of positive components inserts a
positive exceptional component; a centre at a smooth
coefficient-one point creates only a coefficient-zero
leaf; and a centre at a smooth point of coefficient
at least two creates a positive leaf. These are all the
possibilities in the exact divisor formula
$D_{\rm new}=\sigma^*D_{\rm old}-E$.
The $r$ points of $C$ so obtained are distinct.

\item If every component of $B_j$ is $\pi$-exceptional, there
has been at least one multiplicity-one blowup centre
above $F_j$. Indeed, in the absence of such a centre,
every component of the entire total transform of $F_j$
has positive conductor coefficient. That complete
ruling fibre, with its positive multiplicities, has
self-intersection zero. It cannot be supported on the
negative-definite exceptional matrix of a normal surface
resolution. Distinct original fibres require distinct
blowup events. Hence at most $N_0$ of the $B_j$ can be
entirely $\pi$-exceptional.

\item Discard those fibres and the $r_{\rm dom}$ fibres counted
above. There remain at least
\[
                        k\ge r-N_0-r_{\rm dom}
\tag{C.RA.23}\hypertarget{tag-00X5}{}\label{cdp:eq:C.RA.23}
\]
fibres with a surviving strict conductor component and
none dominating $H$. Follow the connected positive
boundary from its point on $C$ to the first component
which is not $\pi$-exceptional. Its image is a curve
different from $H$ meeting $H$ at that point's image.
Every intervening connected chain is contracted by
$\pi$, so this is an actual intersection in $S$, not
an identification imposed on an auxiliary graph.

\item Use the normalization-incidence graph of
$G=\Gamma_{\rm red}$: its vertices are the rational
normalization components and the singular points with
more than one normalization branch, and its edges are
the individual branch preimages. The preceding $k$
intersections give $k$ distinct edges adjacent to the
vertex belonging to $H$. They remain distinct even
when two of their images are the same singular point:
$C\to H^\nu$ is an isomorphism, and the points on $C$
were distinct. Remove that vertex and all its adjacent
edges. Each of the $c$ remaining connected components
which receives one of these edges contains a
normalization vertex of a curve other than $H$.
In particular $c\le n_G-1$, where $n_G$ is the number
of irreducible components of $G$.

\item Contract a spanning tree in each of those $c$ graph
components. The $k$ retained edges then give a connected
multistar with $c+1$ vertices and $k$ edges. Its first
Betti number is $k-c$, so the first Betti number
$\beta$ of the full incidence graph satisfies
\[
             k-c\le\beta\le1,
          \qquad n_G\ge k+1-\beta\ge k.
\tag{C.RA.24}\hypertarget{tag-00X6}{}\label{cdp:eq:C.RA.24}
\]
The upper bound on $\beta$ is the actual local torus-group
quotient (A.RT.18)--(A.RT.20). Adding genus-zero
normalization components changes no graph first Betti
number, and the normalization quotient gives
$e(G)=n_G+1-\beta\ge n_G$. Combining \eqref{cdp:eq:C.RA.23}--\eqref{cdp:eq:C.RA.24}
with the exact surface Euler formula yields
\[
\begin{split}
 e(S)&=2-r+N_0+e(G)\\
     &\ge2-r+N_0+n_G
      \ge2-r+N_0+(r-N_0-r_{\rm dom})
       =2-r_{\rm dom}.
\end{split}
\tag{C.RA.25}\hypertarget{tag-00X7}{}\label{cdp:eq:C.RA.25}
\]
This proves the assertion.

\end{enumerate}
\end{proof}



The degree-one result accounts for every generic transverse
type allowed by conductor exponent two on that branch. It
does not replace the other branches by a presumed ordinary
triple: those branches may have larger conductor exponents.
The cases in which the horizontal section is contracted, or
maps with higher degree to a singular conductor image, require
the further arguments identified separately below.




\begin{proposition}[Every rational degree-two tacnode cover has positive Euler number]\hypertarget{tag-01CG}{}\label{cdp:st:C.2.RAT.4}
\StatementRecord{cdp:st:C.2.RAT.4}{proved}{01CG}{\SeriesRef{cdp:st:C.2.RAT.1}, \SeriesRef{cdp:st:A.5.18}, \SeriesRef{cdp:eq:C.RA.26}, \SeriesRef{cdp:st:A.5.17}, \SeriesRef{cdp:eq:C.RA.28}, \SeriesRef{cdp:eq:C.RA.30}, \SeriesRef{cdp:eq:C.RA.31}, \SeriesRef{cdp:st:C.2.RAT.3}}

Retain the torus-fibre, rational-normalization and specified
ruled-boundary morphism of \ref{cdp:st:C.2.RAT.1}. Suppose $C$ is not
contracted, $H=h(C)_{\rm red}$ has generic transverse type
$A_3$, and the actual map $C\to H^\nu$ has degree two.
Then $e(S)\ge1$. No smoothness or local complete intersection
assumption on the image curve is required.
\end{proposition}


\begin{proof}

\begin{enumerate}[label={$\langle1\rangle$\arabic*.},leftmargin=*]
\item Write $\Sigma=H^\nu\simeq\PP^1$, $G=\Gamma_{\rm red}$,
$\delta=p_a(H)\le2$, and
$s=\#\nu^{-1}(\operatorname{Sing}H)\le3$ as in
(C.RA.5)--(C.RA.6). There are exactly two generic
normalization branches above $H$, both supplied by the
degree-two map from $C$. Thus no vertical strict conductor
component dominates $H$.

\item We first justify using the lci jet construction without
assuming that property. The low-delta classification in the
proof of \ref{cdp:st:C.2.RAT.1} applies to this same reduced image curve.
Its three-independent-axis germ is excluded by
$b_1(H)\le b_1(G)\le1$. The remaining non-lci possibilities
are the space cusp with semigroup $\langle3,4,5\rangle$
and a cusp with a transverse smooth branch. The
genus-zero double-cover calculation at the end of \SeriesRef{cdp:st:A.5.18}
excludes both: the space-cusp inequalities force
$q_0=-1$ and two simple zeros of the cleared quartic,
one over $H_{\rm reg}$, contradicting (A.T.12); for the
cusp-plus-line ring its cleared quartic has nonpositive
degree and cannot have two branch points. Therefore $H$
is lci at all its points.

\item Let $\Xi$ be the normalization conductor divisor of $H$,
and use the actual jet lines and sections in
(A.T.5)--(A.T.7). Put
\[
 A=\deg\operatorname{div}(a)=-2\deg Q,\qquad
 \sigma\in H^0(\Sigma,L^2(4\Xi)),\qquad
 \deg\operatorname{div}(\sigma)=8-A.
\tag{C.RA.26}\hypertarget{tag-00X8}{}\label{cdp:eq:C.RA.26}
\]
Here $A$ is a nonnegative even integer and $\sigma$ is
the conductor-cleared version of $\Delta=ae-b^2/4$.
The degree-two cover $C\to\Sigma$ has two branch points,
exactly the odd-order zeros of $\sigma$. Let $B_s$ be
the number of these branch points lying over singular
points of $H$; there are $2-B_s$ over its smooth locus.
Every smooth branch point has zero order at least three
in $\sigma$, by the actual coordinate exclusion
(A.T.10)--(A.T.12).

\item The section $a$ must vanish at every point of $\Sigma$
over a singular point of $H$. To see this, the Hessian
matrix has entries in $\OO_H$. If its normalized scalar
$a$ is nonzero on one branch above $p$, its primitive
rank-one factor has a unit component there. A diagonal
entry consequently has a nonzero value at $p$. That
same entry has the same nonzero value on every branch,
so $a$ is nonzero at all normalization points above $p$.
In that event (A.T.13) removes the conductor denominator:
$\Delta$ is holomorphic in $L^2$ there. The orders of
$\sigma$ above $p$ are at least those of $4\Xi$,
which are even. Its total zero divisor must additionally
contain the two parity contributions of the two branch
points, whether or not they lie over $p$. Thus
\[
 \deg\operatorname{div}(\sigma)
       \ge4\deg(\Xi|_{\nu^{-1}(p)})+2
       =8\delta_p+2\ge10,
\tag{C.RA.27}\hypertarget{tag-00X9}{}\label{cdp:eq:C.RA.27}
\]
contrary to \eqref{cdp:eq:C.RA.26}. This proves the claimed vanishing.

\item It follows in particular that
\[
                              A\ge2s-2.
\tag{C.RA.28}\hypertarget{tag-00XA}{}\label{cdp:eq:C.RA.28}
\]
For $s\le1$ this is automatic; for $s=2$ the two
zeros give $A\ge2$. If $s=3$, the bounds
$\delta\le2$ and $b_1(H)\le1$ force exactly two
singular points, each of delta invariant one: one node
and one cusp. The two node preimages contribute at
least one each. At the cusp the ring is
$\mathbb C\{t^2,t^3\}$, so any vanishing Hessian entry
has order at least two. Dividing such a diagonal entry
by the square of a unit component of the primitive
normalized factor shows that $a$ has order at least
two there. Thus $A\ge1+1+2=4=2s-2$.

\item We now check exactly which smooth image points can
receive two vertical boundary attachments on their two
conductor preimages. At a smooth point with $a\ne0$,
\SeriesRef{cdp:st:A.5.17} gives the constant $A_3$ product and no vertical
attachment. Suppose instead that $a=0$ and $\Delta$
is a unit. In the actual coordinates (A.T.10), the
coefficient $b$ is then a unit. Substitute $x=y^2\lambda$
in that equation and divide by $y^4$. The result is
holomorphic in $(\lambda,y,z)$ and has, at the point
in question, the simple root
$\lambda=-e(0)/b(0)$. The implicit-function theorem
gives a holomorphic $\lambda(y,z)$ and an actual
factorization
\[
 f_0=F_1F_2,\qquad
 F_1=x-y^2\lambda(y,z),\qquad
 F_2|_{F_1=0}=y^2\,u(y,z),\quad u(0)\ne0.
\tag{C.RA.29}\hypertarget{tag-00XB}{}\label{cdp:eq:C.RA.29}
\]
The last equality is the restriction of
$\partial f_0/\partial x$ to $F_1=0$; its quotient
by $y^2$ has constant value $b(0)$.
The branch $F_1=0$ is smooth, and its conductor in the
whole reduced hypersurface is exactly the divisor
$2\{y=0\}$, by the intersection factor in
(C.RA.22). There is therefore no vertical conductor
near this normalization point. Its minimal resolution
is an isomorphism there as well. Since $\Delta$ is a
unit, this is one of the two distinct conductor
preimages over the given point of $\Sigma$.
Consequently a smooth nonbranch value with two vertical
attachments has a positive even zero order of $\sigma$,
at least two.

\item Let $d_{\rm reg}$ count distinct smooth values with
two such attachments. The two branch points and these
distinct nonbranch zeros give
\[
             3(2-B_s)+B_s+2d_{\rm reg}\le8-A.
\tag{C.RA.30}\hypertarget{tag-00XD}{}\label{cdp:eq:C.RA.30}
\]
Among the $s$ points above the singular locus, the
$B_s$ branch points each have only one preimage on
$C$. Thus at most $s-B_s$ of these values can have
two attachments. If $d$ counts all values on
$\Sigma$ with two attachments, \eqref{cdp:eq:C.RA.28}--\eqref{cdp:eq:C.RA.30}
give
\[
 d\le s-B_s+d_{\rm reg}
       \le s+1-\frac A2\le2.
\tag{C.RA.31}\hypertarget{tag-00XE}{}\label{cdp:eq:C.RA.31}
\]

\item In fact $d\le1$. If $d=2$, \eqref{cdp:eq:C.RA.31} requires
$A\le2s-2$. For $s=0$ this is impossible, and for
$s=1$ it contradicts the vanishing of $a$ over the
singular point. For $s=2$, it forces $A=2$ and
the two singular preimages exhaust the zero divisor
of $a$. For $s=3$, it forces $A=4$, exhausted by
the node and cusp orders $1+1+2$ established above.
In both remaining cases $a$ is nowhere zero over
$H_{\rm reg}$. The constant $A_3$ product there
has neither a cover branch point nor a vertical
attachment. Thus both cover branch points lie
over $\operatorname{Sing}H$, and only the other
$s-2\le1$ singular preimages can have two
attachments. This again contradicts $d=2$.
Consequently
\[
                            d\le1.
\tag{C.RA.31a}\hypertarget{tag-00XF}{}\label{cdp:eq:C.RA.31a}
\]

\item It remains to connect this local count to the exact
surface Euler number. As in \ref{cdp:st:C.2.RAT.3}, at least
$K\ge r-N_0$ original positive vertical boundaries
have a component which is not $\pi$-exceptional.
None maps onto $H$, so they give $K$ distinct points
of $C$ anchored to image curves other than $H$.
After projection to $\Sigma$ their number of distinct
normalization-incidence edges is $k=K-d'$ where
$d'\le d\le1$: a degree-two fibre has at most two
points, and the loss of one occurs precisely at a
value with two selected attachments. Apply the graph
construction of (C.RA.24) to these $k$ distinct
edges. It gives $n_G\ge k$ and $e(G)\ge n_G$.
The actual Euler decomposition therefore yields
\[
\begin{split}
 e(S)&=2-r+N_0+e(G)
       \ge2-K+e(G)\\
     &\ge2-K+k=2-d'\ge1.
\end{split}
\tag{C.RA.32}\hypertarget{tag-00XG}{}\label{cdp:eq:C.RA.32}
\]
Every contraction, covering fibre and graph edge in
this count comes from the specified maps
$Z\to T\to S$ and $C\to\Sigma\to H$.

\end{enumerate}
\end{proof}






\begin{proposition}[A locally unibranch tangent-root component gives positive Euler number]\hypertarget{tag-01CH}{}\label{cdp:st:C.2.RAT.5}
\StatementRecord{cdp:st:C.2.RAT.5}{proved}{01CH}{\SeriesRef{cdp:st:C.2.RAT.1}, \SeriesRef{cdp:eq:C.RD.2}, \SeriesRef{cdp:eq:C.RD.4}, \SeriesRef{cdp:st:C.2.RAT.3}}

Retain the torus-fibre, rational-normalization and specified
ruled-boundary morphism of \ref{cdp:st:C.2.RAT.1}. Suppose $C$ is not
contracted and maps with degree $d=2$ or $3$ onto a smooth
rational curve $H$ with generic ordinary-triple transversal.
Let $q_3\in H^0(H,\operatorname{Sym}^3N^*\otimes M)$ be
its cubic normal tensor, where $N=N_{H/Y}$ and
$M=\OO_Y(S)|_H$. Suppose $q_3$ is nonzero as a tensor at
every point of $H$. Its tangent-root divisor
$R_3\subset\PP(N)$ is then finite of degree three over $H$.
Let $R_d$ be the component corresponding to $C$ over the
generic point. If $R_d$ is locally unibranch at every point,
then
\[
                           e(S)\ge d-1>0.
\tag{C.RD.1}\hypertarget{tag-00XZ}{}\label{cdp:eq:C.RD.1}
\]
In particular this applies when $d=2$ and the only singularities
of $R_2$ are an $A_4$ cusp or two $A_2$ cusps.
\end{proposition}


\begin{proof}

\begin{enumerate}[label={$\langle1\rangle$\arabic*.},leftmargin=*]
\item The nonzero cubic has three distinct roots over the generic
point and is never identically zero on a projective fibre.
It therefore defines a reduced finite divisor $R_3$.
The function-field algebra of its roots agrees with that of
the normalized conductor curves dominating $H$: a branch
of an ordinary triple is distinguished by its tangent line.
Taking the integral closure of $\OO_H$ in that algebra gives
the exact finite-cover identification
\[
                            C\simeq R_d^\nu
                            \quad\text{over }H.
\tag{C.RD.2}\hypertarget{tag-00Y9}{}\label{cdp:eq:C.RD.2}
\]
Thus two distinct points of $C$ over one point $p\in H$
give distinct root points of $R_d$, because normalization
of a locally unibranch reduced curve is bijective on its
underlying points.

\item We check what a simple root implies for the actual
hypersurface, including its conductor. In coordinates
$H=(x,y)=0$, with parameter $z$ on $H$, the local equation
$f_0$ of $S$ belongs to $(x,y)^3$. The lower normal jets
vanish on the reduced smooth curve, and the graded pieces
of its regular-sequence ideal are locally free, which
justifies this membership. Choose the normal coordinates
so a specified simple tangent root is finite in $x/y$.
Then
\[
 f_0(y\lambda,y,z)/y^3
       =q_3(\lambda,1,z)+y\,A(\lambda,y,z).
\tag{C.RD.3}\hypertarget{tag-00YK}{}\label{cdp:eq:C.RD.3}
\]
The simple root gives a holomorphic solution
$\lambda=\lambda(y,z)$ by the implicit-function theorem.
Division by the monic linear factor yields an actual
factorization
\[
 f_0=(x-y\lambda(y,z))\,G(x,y,z),\qquad
 G|_{x=y\lambda(y,z)}=y^2 u(y,z),\quad u(0)\ne0.
\tag{C.RD.4}\hypertarget{tag-00YT}{}\label{cdp:eq:C.RD.4}
\]
The unit is the derivative of the cubic with respect to
its root coordinate, at the chosen simple root. Hence
this normalization branch is smooth and its conductor
in $S$ is exactly twice the curve lying over $H$.
There is no vertical conductor component there, and
the minimal resolution $\pi$ is an isomorphism nearby.

\item Two distinct $C$-points over $p$ cannot both receive a
vertical boundary attachment. Their root points are
distinct by \eqref{cdp:eq:C.RD.2}. If both could receive attachments,
\eqref{cdp:eq:C.RD.4} would make both roots multiple, requiring total
cubic degree at least four. Thus each point of $H$
receives at most one attachment from $C$.

\item There are at most $3-d$ generic normalization branches
left to be supplied by vertical strict conductor
components. In the notation of \ref{cdp:st:C.2.RAT.3} this gives
$r_{\rm dom}\le3-d$. At least
$K\ge r-N_0-r_{\rm dom}$ original positive vertical
boundaries therefore survive with an image different
from $H$. Their $K$ points of $C$ give $K$ distinct
normalization-incidence edges at $H$, by the preceding
paragraph. Equation (C.RA.24) gives $n_G\ge K$, and
the exact Euler decomposition gives
\[
 e(S)\ge2-r+N_0+K
       \ge2-r_{\rm dom}\ge d-1.
\]
This proves the assertion. The final examples are locally
unibranch by their displayed cusp equations.

\end{enumerate}
\end{proof}






\begin{proposition}[Every primitive rational double ordinary-triple cover has nonnegative Euler number]\hypertarget{tag-01CJ}{}\label{cdp:st:C.2.RAT.7}
\StatementRecord{cdp:st:C.2.RAT.7}{proved}{01CJ}{\SeriesRef{cdp:st:C.2.RAT.1}, \SeriesRef{cdp:st:A.5.19}, \SeriesRef{cdp:eq:C.RD.14}, \SeriesRef{cdp:eq:C.RD.15}}

Retain the torus-fibre, rational-normalization and specified
ruled-boundary morphism of \ref{cdp:st:C.2.RAT.1}. Suppose $C$ is not
contracted, its image $H$ has generic ordinary-triple
transversal, and $C\to H^\nu$ has degree two. Suppose
the cubic normal tensor along $H$ has no common zero of
its coefficients. Then $e(S)\ge0$. This holds for every
splitting type of the normal bundle $N_{H/Y}$.
\end{proposition}


\begin{proof}

\begin{enumerate}[label={$\langle1\rangle$\arabic*.},leftmargin=*]
\item Statement \SeriesRef{cdp:st:A.5.19} makes $H$ smooth: over a singular
rational ordinary-triple image every connected normalized
conductor cover has degree one. Thus $H\simeq\PP^1$.
Set $N=N_{H/Y}$, $W=\PP(N)$, and
$\xi=c_1(\OO_W(1))$, where $\PP(N)$ parameterizes lines.
The ambient Chern-class hypothesis gives
$\deg N=-2$ and $\OO_Y(S)|_H\simeq\OO_H$.
The cubic is an actual section of $\OO_W(3)$.
The assumed absence of a common coefficient zero makes
its root divisor finite of degree three over $H$:
\[
 R=R_1+R_2\sim3\xi,\qquad
 \xi^2=2,\qquad K_W=-2\xi,\qquad p_a(R)=4.
\tag{C.RD.11}\hypertarget{tag-00Y0}{}\label{cdp:eq:C.RD.11}
\]
Here $R_1$ is a section and hence is smooth rational;
$R_2$ has normalization $C\simeq\PP^1$, by the exact
integral-closure comparison (C.RD.2). Both are reduced
and they are different. The normalization sequence gives
\[
                  \sum_p\delta_p(R)=5.
\tag{C.RD.12}\hypertarget{tag-00Y1}{}\label{cdp:eq:C.RD.12}
\]
This equality also holds if the two components are
disjoint: the arithmetic-genus convention is
$p_a(R)=1-\chi(\OO_R)$ and the normalization has Euler
characteristic two.

\item Form the cyclic triple cover
$p_0:X_0\to W$ with algebra
$\OO_W\oplus\OO_W(-1)\oplus\OO_W(-2)$ and multiplication
given by the cubic. It is normal because its branch
divisor $R$ is reduced, and is connected because a
meridian around a smooth branch point permutes its
three sheets cyclically. Resolve its isolated
singularities and its reduced boundary to obtain
$X$, with $D=X\setminus U$, where
$U=\{v\in N:q_3(v)=1\}$ is the smooth normal-cone
level surface. The isomorphism over $W\setminus R$
is the exact radial-vector/dual-tautological-coordinate
comparison (C.D4.6).

\item The projective-bundle cohomology calculation gives
\[
 H^1(W,\OO_W)=H^1(W,\OO_W(-1))
             =H^1(W,\OO_W(-2))=0.
\tag{C.RD.13}\hypertarget{tag-00Y2}{}\label{cdp:eq:C.RD.13}
\]
Indeed both relative cohomologies of $\OO_W(-1)$
vanish, while $p_*\OO_W(-2)=0$ and
$R^1p_*\OO_W(-2)=\det N=\OO_{\PP^1}(-2)$.
Thus $H^1(X_0,\OO_{X_0})=0$ for every normal-bundle
splitting. Resolution Leray and \CompendiumRef{II}{cdp:st:C.2.RAT.6} give
$q(X)\le\sum p_g(X_0,x)\le1$, so $b_1(X)\le2$.

\item We spell out the boundary calculation with its number
of connected components retained. Write
\[
 \nu=\sum_p(r_p(R)-1),\qquad
 c=b_0(R),\qquad
 \lambda=\sum_x b_1(\operatorname{link}(X_0,x))\le2.
\]
Normalization of $R$ gives $e(R)=4-\nu$.
Since $U\to W\setminus R$ is etale of degree three,
$e(U)=3\nu$. The resolved boundary has the same
number $c$ of connected components. Replacing each
singular point of $R$ by the connected exceptional
divisor of its cubic suspension gives, by the
Mayer--Vietoris sequence with its individual branch
attachment points,
\[
                 b_1(D)=\nu-2+c+\lambda.
\tag{C.RD.14}\hypertarget{tag-00Y3}{}\label{cdp:eq:C.RD.14}
\]
The exceptional first Betti number is the link's
first Betti number: the exceptional intersection
matrix in its plumbed neighbourhood is invertible
over $\Q$. This is precisely the local comparison
proved with (C.D4.7). Additional boundary-resolution
blowups add rational trees and leave \eqref{cdp:eq:C.RD.14}
unchanged.

\item Let $d_D$ be the number of irreducible components
of $D$, and let $V\subset H_2(X,\Q)$ be the span
of their classes. The exceptional classes are
independent. Pushing a relation to $W$ and pulling
back a relation from $R_1,R_2$ shows that
\[
 \varepsilon:=d_D-\dim V
      =2-\dim\operatorname{span}\{[R_1],[R_2]\}
      \in\{0,1\}.
\]
This proves the relation count rather than
assuming the two branch classes are independent.
The space $V$ contains the pullback of $R$,
whose square is $3R^2=54>0$. The Hodge index
theorem therefore makes its intersection pairing
nondegenerate. The divisor Gysin sequence gives
$b_1(U)=b_1(X)+\varepsilon\le3$; its degree-two
part identifies the image of $H_2(U)$ in
$H_2(X)$ with $V^\perp$. Consequently the rank
of the compact intersection pairing on $U$ is
$b_2(X)-d_D+\varepsilon$.
Now use $e(X)=2-2b_1(X)+b_2(X)$ and
$e(D)=c-b_1(D)+d_D$. Equations \eqref{cdp:eq:C.RD.14} and
$e(U)=3\nu$ yield the exact identity
\[
\begin{split}
 \operatorname{rank}I_U
   &=e(U)-2+2b_1(X)+c-b_1(D)+\varepsilon\\
   &=2\nu+2b_1(X)-\lambda+\varepsilon
     \ge2\nu-2.
\end{split}
\tag{C.RD.15}\hypertarget{tag-00Y4}{}\label{cdp:eq:C.RD.15}
\]
In particular this formula includes disjoint root
components; no connected-boundary assumption has
been introduced.

\item If $\nu\ge2$, \eqref{cdp:eq:C.RD.15} gives a nonzero compact
intersection pairing on a surface with $b_1(U)\le3$.
The compact-core transport (C.D4.2) embeds that
pairing into an actual nearby torus. Its induced
first-homology image has dimension at most three,
so its second-homology image has zero wedge pairing,
as established immediately after (C.D4.2). This
contradiction proves $\nu\le1$.

\item Finally suppose a value of $H$ receives vertical
boundary attachments at both its $C$-preimages.
If their tangent-root points in $R_2$ were distinct,
at least one would be a simple root of the cubic,
and (C.RD.3)--(C.RD.4) would give pure conductor
there, a contradiction. They therefore map to a
single multibranch point of $R_2$. Distinct values
give distinct such points, each contributing at
least one to $\nu$. Hence there is at most one
value with two attachments. With
$K\ge r-N_0-r_{\rm dom}$ surviving original
vertical boundaries, at least $K-1$ distinct
normalization-incidence edges remain at $H$.
There is only one remaining generic triple branch,
so $r_{\rm dom}\le1$. The graph comparison
(C.RA.24) and exact Euler decomposition now give
\[
 e(S)\ge2-r+N_0+(K-1)
                         \ge1-r_{\rm dom}\ge0.
\tag{C.RD.16}\hypertarget{tag-00Y5}{}\label{cdp:eq:C.RD.16}
\]
This completes the proof.

\end{enumerate}
\end{proof}






\begin{proposition}[A single common cubic zero in the rational double-cover case]\hypertarget{tag-01CK}{}\label{cdp:st:C.2.RAT.8}
\StatementRecord{cdp:st:C.2.RAT.8}{proved}{01CK}{\SeriesRef{cdp:st:C.2.RAT.1}, \SeriesRef{cdp:st:C.2.RAT.7}, \SeriesRef{cdp:eq:C.RD.17}, \SeriesRef{cdp:eq:C.RD.18}, \SeriesRef{cdp:eq:C.RD.20}, \SeriesRef{cdp:eq:C.RD.21}, \SeriesRef{cdp:eq:C.RD.22}; \cite[Theorem 2.1.1 and Corollary 2.1.3]{MerleTeissier1977}}

Retain the torus-fibre, rational-normalization and specified
ruled-boundary morphism of \ref{cdp:st:C.2.RAT.1}. Suppose $C$ is not
contracted, its image has generic ordinary-triple transversal,
and $C\to H^\nu$ has degree two. If the common zero divisor
of the cubic normal coefficients has degree one, then
$e(S)\ge0$.
\end{proposition}


\begin{proof}

\begin{enumerate}[label={$\langle1\rangle$\arabic*.},leftmargin=*]
\item As in \ref{cdp:st:C.2.RAT.7}, $H$ is smooth rational. On $W=\PP(N)$
write the complete root divisor as
\[
 B=R_1+R_2+F\sim3\xi,
 \qquad B_h=R_1+R_2\sim3\xi-F.
\tag{C.RD.17}\hypertarget{tag-00Y6}{}\label{cdp:eq:C.RD.17}
\]
Here $F$ is the ruling fibre over the unique simple
common zero. The horizontal divisor $B_h$ is finite
over $H$, $R_1$ is a section, and the normalization of
$R_2$ is the rational degree-two cover $C$.
Since $K_W=-2\xi$, $\xi^2=2$, and $\xi F=1$,
\[
 p_a(B_h)=2,
 \qquad \sum_p\delta_p(B_h)=3,
 \qquad I(F,B_h)=3.
\tag{C.RD.18}\hypertarget{tag-00Y7}{}\label{cdp:eq:C.RD.18}
\]
The full divisor $B$ has total delta six. We will
use its particular decomposition \eqref{cdp:eq:C.RD.17}, not a
statement about arbitrary delta-six plane germs.

\paragraph{The local cubic-suspension genus along $F$.}
At a point of $F$, put $i=I_p(F,B_h)$. If $i=1$,
the germ of $B$ is an ordinary node. If $i=2$,
the horizontal germ is either smooth with contact
two with $F$, two smooth branches transverse to
$F$, or a multiplicity-two cusp transverse to $F$.
The bound $\delta_p(B_h)\le3$ makes the resulting
plane types respectively
\[
 A_3,\qquad D_{2k+2}\ (1\le k\le3),\qquad
                 D_{2k+3}\ (1\le k\le3).
\tag{C.RD.19}\hypertarget{tag-00Y8}{}\label{cdp:eq:C.RD.19}
\]
For two smooth branches $k$ is their contact;
for the cusp $k$ is its delta invariant.
The exact analytic identifications are the graph
straightening and square-completion comparisons
in \CompendiumRef{II}{cdp:st:C.2.RAT.6}. Their cubic suspensions have
geometric genus zero or one by (C.RD.10).

\item If $i=3$, the point contains $R_1$ and all the
intersection of $R_2$ with $F$. The section $R_1$
is transverse to $F$. A smooth branch of $R_2$
with intersection two is tangent to $F$ and
transverse to $R_1$, giving $D_6$. If $R_2$ has
two local branches, both are smooth and transverse
to $F$. The horizontal delta bound three forces
all three horizontal branches to be pairwise
transverse. The full germ is therefore four
distinct transverse smooth branches.
The last possibility is one singular branch
of $R_2$. Its multiplicity is two, since its
intersection with $F$ is two. Its intersection
with $R_1$ is at least two; hence \eqref{cdp:eq:C.RD.18}
forces it to be an ordinary cusp, transverse
to both smooth branches $F$ and $R_1$.
These alternatives exhaust $i=3$.

\item For four transverse branches, the degree-four
initial plane polynomial has four distinct
linear factors. Its cubic suspension is Newton
nondegenerate with weights $(1/4,1/4,1/3)$;
the only positive lattice point on or below
its compact Newton face is $(1,1,1)$. Its
geometric genus is therefore one by
\cite[Theorem 2.1.1 and Corollary 2.1.3]{MerleTeissier1977}.

\item We give the precise same calculation for the cusp
and two transverse branches, retaining higher
terms of the actual equation. Straighten $F,R_1$
to coordinate axes and make a linear change so
the cusp tangent is $u=0$. Up to units, the
equation is then
\[
 r(u,v)=v(u-v)\bigl(u^2+c v^3+
           u\,O((u,v)^2)+O((u,v)^4)\bigr),
                  \qquad c\ne0.
\tag{C.RD.20}\hypertarget{tag-00YA}{}\label{cdp:eq:C.RD.20}
\]
The nonzero radical cubic is the ordinary-cusp
condition. The compact Newton facets of
$z^3-r(u,v)$ have normalized inward vectors
\[
             (1/4,1/4,1/3),\qquad
             (3/10,1/5,1/3).
\tag{C.RD.21}\hypertarget{tag-00YB}{}\label{cdp:eq:C.RD.21}
\]
Their plane edges have vertex sets
$\{(3,1),(2,2)\}$ and
$\{(2,2),(0,5)\}$. The higher terms in
\eqref{cdp:eq:C.RD.20} are strictly above these faces.
The first edge polynomial is a nonzero multiple
of $u^2v(u-v)$; the second is a nonzero multiple
of $v^2(u^2+c v^3)$. Both have smooth zero sets
in $(\C^*)^2$. Each two-dimensional compact
face additionally contains the term $z^3$,
whose $z$-derivative is nonzero on $(\C^*)^3$.
This proves the required Newton nondegeneracy.
Every noncompact face inequality is strictly
satisfied by positive integral points; the only
additional one beyond coordinate positivity is
$v+z/3\ge1$.
For either vector in \eqref{cdp:eq:C.RD.21}, the point
$(1,1,1)$ has weight $5/6$, and increasing any
one coordinate makes its weight greater than
one. Thus it is the only positive lattice
point outside the interior of the Newton
polyhedron. The same adjunction theorem gives
$p_g(z^3-r)=1$.

\item This proves that every suspension singularity
over $F$ has geometric genus at most one.
A nonzero value requires $i\ge2$ and
$\delta_p(B_h)\ge1$: the cases with horizontal
delta zero were the ordinary node and $A_3$.
Since the total intersection with $F$ is three,
there is at most one such nonzero value over
$F$. Away from $F$, \CompendiumRef{II}{cdp:st:C.2.RAT.6} shows that a
nonzero geometric genus costs horizontal delta
at least three. It cannot coexist with a
nonzero value over $F$, by \eqref{cdp:eq:C.RD.18}.
Consequently the cyclic cover $X_0$ branched
over $B$ satisfies
\[
                        \sum_x p_g(X_0,x)\le1.
\tag{C.RD.22}\hypertarget{tag-00YC}{}\label{cdp:eq:C.RD.22}
\]

\item We also obtain the required link bound without
assuming a delta-six monodromy table. For a
normal surface singularity, let $E$ be the
reduced exceptional divisor of a good resolution.
The restriction maps from infinitesimal
neighbourhoods of $E$ onto $E$ are surjective
on $H^1$, since the kernels have one-dimensional
support and no $H^2$. Formal functions therefore
gives $h^1(E,\OO_E)\le p_g$.
For an SNC exceptional curve,
$b_1(E)=2\sum g_i+\beta\le
2(\sum g_i+\beta)=2h^1(E,\OO_E)$.
The invertible exceptional intersection matrix
identifies this $b_1(E)$ with the link's first
Betti number, as in the proof of \CompendiumRef{II}{cdp:st:C.2.RAT.6}.
Equation \eqref{cdp:eq:C.RD.22} yields
\[
              \lambda:=\sum_x b_1(\operatorname{link}(X_0,x))
                            \le2.
\tag{C.RD.23}\hypertarget{tag-00YD}{}\label{cdp:eq:C.RD.23}
\]

\paragraph{The compact pairing and the attachment count.}
The cyclic algebra and projective-bundle cohomology
are unchanged from (C.RD.13), so its projective
resolution $X$ has $b_1(X)\le2$.
The divisor $B$ is connected, since $F$ meets
both horizontal components. Its three classes
span a two-dimensional space: the section
$R_1$ and the fibre $F$ are independent.
Thus the boundary-relation number is exactly
$\varepsilon=1$, and $b_1(U)=b_1(X)+1\le3$.
Put $\nu=\sum_p(r_p(B)-1)$.
Normalization of its three rational components
and the resolved-boundary comparison give
\[
 e(U)=3(\nu-2),\qquad
 b_1(D)=\nu-2+\lambda,
 \qquad
 \operatorname{rank}I_U
       =2\nu-4+2b_1(X)-\lambda\ge2\nu-6.
\tag{C.RD.24}\hypertarget{tag-00YE}{}\label{cdp:eq:C.RD.24}
\]
The rank identity is the same Gysin, Hodge-index
and Euler calculation as (C.RD.15), with all
three boundary components and their one relation
included. Compact-core transport into the
nearby torus therefore forces $\nu\le3$.

\item Let $p\in H$ be the common zero value and
$a=\#(C\to H)^{-1}(p)\in\{1,2\}$.
At every point of $F\cap B_h$, the full divisor
has the branch of $F$ together with the
individual horizontal normalization branches.
It follows exactly that
\[
        \nu_F:=\sum_{x\in F}(r_x(B)-1)=1+a.
\tag{C.RD.25}\hypertarget{tag-00YF}{}\label{cdp:eq:C.RD.25}
\]
The contribution one is the section $R_1$;
the other $a$ branches are the normalization
points of $R_2$ above $p$. This identity retains
coinciding tangent roots and ramification.

\item The loss of distinct horizontal anchor values
at $p$ is at most $a-1=\nu_F-2$. Away from
$p$, the cubic is primitive. The argument at
the end of \ref{cdp:st:C.2.RAT.7} makes each duplicate
attachment there a multibranch point of $R_2$,
costing at least one in
$\nu_{\rm out}=\nu-\nu_F$. Hence the total
loss $\ell$ of distinct normalization-incidence
edges satisfies
\[
                    \ell\le\nu_F-2+\nu_{\rm out}
                              =\nu-2\le1.
\tag{C.RD.26}\hypertarget{tag-00YG}{}\label{cdp:eq:C.RD.26}
\]
Finally $r_{\rm dom}\le1$ and
$K\ge r-N_0-r_{\rm dom}$. Projecting these
$K$ surviving attachments to $H$ retains at
least $K-\ell$ distinct edges. Equation
(C.RA.24) and the Euler decomposition give
\[
                 e(S)\ge2-r_{\rm dom}-\ell\ge0.
\tag{C.RD.27}\hypertarget{tag-00YH}{}\label{cdp:eq:C.RD.27}
\]
The two Newton facets, their exact lower
lattice-point counts and the attachment
budgets are recorded in
\path{code/A_cubic_scalar_one.py} and its output.

\end{enumerate}
\end{proof}






\begin{proposition}[Every rational double ordinary-triple cover has nonnegative Euler number]\hypertarget{tag-01CL}{}\label{cdp:st:C.2.RAT.9}
\StatementRecord{cdp:st:C.2.RAT.9}{proved}{01CL}{\SeriesRef{cdp:st:C.2.RAT.1}, \SeriesRef{cdp:st:A.5.19}, \SeriesRef{cdp:st:C.2.RAT.7}, \SeriesRef{cdp:st:C.2.RAT.8}, \SeriesRef{cdp:eq:C.RD.28}, \SeriesRef{cdp:eq:C.RD.34}}

Retain the torus-fibre, rational-normalization and specified
ruled-boundary morphism of \ref{cdp:st:C.2.RAT.1}. Suppose $C$ is not
contracted, its image has generic ordinary-triple transversal,
and $C\to H^\nu$ has degree two. Then $e(S)\ge0$,
including all common zeros and their multiplicities in
the cubic normal tensor.
\end{proposition}


\begin{proof}

\begin{enumerate}[label={$\langle1\rangle$\arabic*.},leftmargin=*]
\item Again \SeriesRef{cdp:st:A.5.19} makes $H$ smooth rational. Let $A$ be the
common zero divisor of the cubic coefficients and put
$v=\deg A$. The cases $v=0,1$ are \ref{cdp:st:C.2.RAT.7}--8.
We first verify that $v\le2$. Write
$N=\OO(a)\oplus\OO(-2-a)$ with $a\ge-1$.
For $a\ge0$ and $v>0$, the coefficient of the cube
of the first normal coordinate has degree $-3a-v$
and vanishes. Thus the other coordinate is a factor.
The remaining quadratic is irreducible over the
function field, because it describes the connected
degree-two cover. Its first-coordinate-square
coefficient must be nonzero and has degree
$2-a-v$. This gives $a+v\le2$.
For $a=-1$, on $\PP^1\times\PP^1$ the section
$R_1$ and quadratic component have bidegrees
$(1,d)$ and $(2,3-v-d)$, with $d\ge0$ and the
first entry denoting degree over $H$. An integral
quadratic component requires $3-v-d\ge1$:
a divisor of bidegree $(2,0)$ is a union of
constant sections over $\C$. Thus $v\le2$ here
as well.

\item It remains to consider $v=2$. With the notation
of \ref{cdp:st:C.2.RAT.7}, the horizontal divisor is
\[
 B_h=R_1+R_2\sim3\xi-2F,
       \qquad p_a(B_h)=0,
       \qquad \sum_p\delta_p(B_h)=1.
\tag{C.RD.28}\hypertarget{tag-00YI}{}\label{cdp:eq:C.RD.28}
\]
Both normalization components are rational.
Consequently its only possible singularity is
one ordinary node or one ordinary cusp. An
intersection between $R_1$ and $R_2$, if present,
is the node. Otherwise a node or cusp belongs
to $R_2$. Every other point of the horizontal
divisor is smooth.

\paragraph{Two distinct common zero values.}
First let $A=p+q$ with $p\ne q$. The full branch
divisor is reduced:
$B=R_1+R_2+F_p+F_q\sim3\xi$.
Use the ordinary cyclic algebra of (C.RD.13).
The local analysis in the proof of \ref{cdp:st:C.2.RAT.8}
applies at each of the two fibres. A nonzero
local geometric genus there requires positive
horizontal delta. The total horizontal delta
in \eqref{cdp:eq:C.RD.28} is one, and the two fibres are
disjoint. Away from these fibres the local
plane delta is at most one and its cubic
suspension is rational. Therefore
\[
 P:=\sum_x p_g(X_0,x)\le1,\qquad
 \lambda:=\sum_x b_1(\operatorname{link}(X_0,x))\le2.
\tag{C.RD.29}\hypertarget{tag-00YJ}{}\label{cdp:eq:C.RD.29}
\]

\item In this case the irregularity is zero, not
merely at most one. The finite-cover dualizing
formula gives
$\omega_{X_0}=p_0^*(\omega_W\otimes\OO_W(2))
\simeq\OO_{X_0}$, and (C.RD.13) gives
$H^1(X_0,\OO)=0$. Thus $\chi(\OO_{X_0})=2$.
For a resolution $\rho:X\to X_0$, local duality
and Grauert--Riemenschneider vanishing identify
$\rho_*\omega_X\subset\omega_{X_0}$ with an
ideal of colength $p_g(X_0,x)$ at each singular
point. Explicitly this is the dual of the
finite-length $R^1\rho_*\OO_X$ in the surface
duality sequence. If $P>0$, the global generator
of $\omega_{X_0}$ does not belong to that ideal
at a point with positive colength. No nonzero
global canonical form on $X$ can then exist,
because its pushforward would be a nonzero
constant multiple of that generator.
Resolution Leray now gives
\[
 \chi(\OO_X)=2-P,
 \qquad q(X)=P-1\quad(P>0),
 \qquad q(X)=0\quad(P=0).
\tag{C.RD.30}\hypertarget{tag-00YL}{}\label{cdp:eq:C.RD.30}
\]
Since $P\le1$, this proves $b_1(X)=0$.

\item The four branch classes span the two-dimensional
space generated by a section and a fibre, so
$\varepsilon=2$ and $b_1(U)=2$.
The branch divisor and resolved boundary are
connected. For
$\nu=\sum_x(r_x(B)-1)$, normalization and the
exact pairing calculation give
\[
 e(U)=3(\nu-4),\qquad b_1(D)=\nu-3+\lambda,
 \qquad \operatorname{rank}I_U=2\nu-8-\lambda.
\tag{C.RD.31}\hypertarget{tag-00YM}{}\label{cdp:eq:C.RD.31}
\]
The compact torus comparison forces this rank
to be zero; otherwise its nondegenerate compact
pairing would embed into a rank-at-most-two
first-homology image in the torus. Hence
$\nu=4+\lambda/2\le5$.

\item Let $a_p,a_q\in\{1,2\}$ be the numbers of
$C$-preimages of these two common values.
The exact branch count (C.RD.25) gives
$\nu_{F_p}+\nu_{F_q}=2+a_p+a_q$.
The loss of distinct attachment values on
these fibres is at most $a_p+a_q-2$.
Away from them each loss costs a multibranch
point of $R_2$. Consequently
\[
                    \ell\le\nu-4\le1.
\tag{C.RD.32}\hypertarget{tag-00YN}{}\label{cdp:eq:C.RD.32}
\]
The graph count gives
$e(S)\ge2-r_{\rm dom}-\ell\ge0$.

\paragraph{One double common zero value.}
Now let $A=2p$. Normalize the cyclic cover
with branch multiplicities $R_1+R_2+2F_p$.
Its exact character algebra is
\[
 (p_0)_*\OO_{X_0}
       =\OO_W\oplus L^{-1}\oplus L^{-2}(F_p),
                  \qquad L=\OO_W(1).
\tag{C.RD.33}\hypertarget{tag-00YO}{}\label{cdp:eq:C.RD.33}
\]
The additional integral generator at $F_p$
is $w^2/s$ in $w^3=s^2u$. Its divisorwise
character line extends reflexively over the
smooth base, which proves the sheaf identity
at the finitely many intersections as well;
this is also (C.DD.16). Relative duality gives
$R^1p_*L^{-2}(F_p)=\OO_H(-1)$.
Thus $H^1(X_0,\OO_{X_0})=0$.

\item We check that every singularity of this
normalized cover is rational, including the
horizontal node and cusp allowed by \eqref{cdp:eq:C.RD.28}.
For this purpose write a cyclic branch divisor
as $D_1+2D_2$, with root line $L$.
At a point blowup $\sigma$ let
$m_i=\operatorname{mult}D_i$,
$n=m_1+2m_2=3k+j$, $0\le j<3$, and let $E$
be the new exceptional curve. Keeping the
same function field gives root line
$L'=\sigma^*L(-kE)$ and new exceptional
branch weight $j$. Its two character sheaves
are exactly
\[
\begin{split}
 (L')^{-1}&=\sigma^*(L^{-1})(kE),\\
 (L')^{-2}(D'_2)
   &=\sigma^*(L^{-2}(D_2))
       \bigl((2k-m_2+\mathbf1_{j=2})E\bigr).
\end{split}
\tag{C.RD.34}\hypertarget{tag-00YP}{}\label{cdp:eq:C.RD.34}
\]
Here $D'_2=\sigma^*D_2-m_2E+\mathbf1_{j=2}E$.
These identities follow by replacing the
cyclic root by its quotient by the $k$th
power of a local exceptional equation.

\item At a smooth horizontal branch the contact
with $F_p$ is one or two. The weight-$(2,1)$
contact calculation (C.DD.18)--(C.DD.20)
separates the branches with character
increments $(1,1)$ at each blowup.
If the horizontal node lies on $F_p$, there
are three smooth local branches of weights
$(2,1,1)$. They are transverse except possibly
for contact two between $F_p$ and the
quadratic branch in the case of an intersection
$R_1\cap R_2$. A blowup has $n=4,m_2=1$,
so both character increments in \eqref{cdp:eq:C.RD.34}
are one. In the transverse case it separates
the three branches. In the contact-two case
one remaining ordinary triple has the same
weights $(2,1,1)$; the same blowup separates it.

\item If the horizontal cusp lies on $F_p$, its
local form is a weight-one cusp transverse
to the weight-two fibre:
\[
                        w^3=s^2(x^2-s^3).
\tag{C.RD.35}\hypertarget{tag-00YQ}{}\label{cdp:eq:C.RD.35}
\]
The first blowup has $n=4,m_2=1$ and
increments $(1,1)$. The strict cusp becomes
smooth with contact two with the new
weight-one exceptional curve. Blow up that
weight-$(1,1)$ tangency; here $n=2,m_2=0$
and the increments are $(0,1)$.
The remaining ordinary triple has weights
$(1,1,2)$, so a final blowup has increments
$(1,1)$. The resulting branch has normal
crossings. At a crossing of weights $a,b$
in $\{1,2\}$, normalization of
$w^3=x^a y^b$ is the quotient of $\C^2$
under the subgroup
\[
 \{(\zeta^i,\zeta^j):ai+bj=0\pmod3\},
 \qquad x=u^3,\quad y=v^3,\quad w=u^av^b.
\]
These are the rational cyclic quotient
singularities $1/3(1,1)$ or $1/3(1,2)$.
Their standard resolutions have respectively
one $(-3)$-curve or a two-curve $(-2)$ chain,
and $R^1\OO=0$. The horizontal node or cusp
away from $F_p$ likewise has a rational
cubic suspension by (C.RD.10).

\item For completeness, none of the preceding
base blowups adds a geometric-genus defect.
Every exceptional character increment in
\eqref{cdp:eq:C.RD.34} was zero or one. A point blowup
has $\sigma_*\OO(E)=\OO$ and
$R^1\sigma_*\OO(E)=0$, from
$\OO(E)|_E=\OO_{\PP^1}(-1)$ and the
structure-sheaf direct images; the same
holds for increment zero. Projection formula
therefore makes the first direct image of
each normalized character algebra vanish.
The normalized cyclic covers over successive
base blowups map properly and birationally
to the preceding cover through their same
function field. Finite direct image and
Leray identify their $R^1\OO$ with the
just computed eigensheaf direct images;
faithfulness of finite pushforward gives
$R^1\OO=0$. Resolving the final rational
cyclic quotients preserves this property.
This proves the claimed local rationality.

\item We consequently have $b_1(X)=0$ and
$\lambda=0$. The reduced branch support
$B=R_1+R_2+F_p$ has three rational
normalization components, is connected,
and has class rank two. Thus
\[
 b_1(U)=1,\qquad e(U)=3(\nu-2),\qquad
 b_1(D)=\nu-2,\qquad
          \operatorname{rank}I_U=2\nu-4.
\tag{C.RD.36}\hypertarget{tag-00YR}{}\label{cdp:eq:C.RD.36}
\]
Again compact transport forces rank zero,
hence $\nu=2$.
Equation (C.RD.25) then gives exactly one
$C$-preimage over the common value and no
multibranch point of $R_2$ away from it.
There is no loss of distinct attachment
values. The graph inequality yields
\[
                    e(S)\ge2-r_{\rm dom}\ge1.
\tag{C.RD.37}\hypertarget{tag-00YS}{}\label{cdp:eq:C.RD.37}
\]
This proves the degree-two-common-divisor
case and, together with \ref{cdp:st:C.2.RAT.7}--8 and
the degree bound at the start, every case
in the statement.
The additional weighted node/cusp charts,
their exact character increments and both
attachment budgets are checked in
\path{code/A_cubic_scalar_two.py} and its output.

\end{enumerate}
\end{proof}






\begin{proposition}[Global exclusion of connected positive-genus triple covers over a rational image]\hypertarget{tag-01BW}{}\label{cdp:st:C.2.D4.1}
\StatementRecord{cdp:st:C.2.D4.1}{proved}{01BW}{\SeriesRef{cdp:eq:C.D4.2}, \SeriesRef{cdp:eq:C.D4.7}, \SeriesRef{cdp:eq:C.D4.6}, \SeriesRef{cdp:eq:C.D4.5}, \SeriesRef{cdp:eq:C.D4.9}, \SeriesRef{cdp:eq:C.D4.12}, \SeriesRef{cdp:eq:C.D4.4}; \cite[Theorem 2.1.1 and Corollary 2.1.3]{MerleTeissier1977}; \cite[pp. 165--167]{Steenbrink1989Spectrum}; \cite[\S1 and Example 2.4(a)]{Nemethi1991Join}}

Let $Y$ be a smooth compact complex threefold with
$b_2(Y;\Q)=0$, and let $f:Y\to B$ be a holomorphic map to a
smooth curve whose smooth nearby fibres are complex two-tori.
Fix a fibre $mS$, with $S$ reduced and irreducible, and restrict
the base to a small coordinate disc about it. Let
$H\subset S$ be a smooth rational curve along which the generic
transversal of $S$ is an ordinary triple point. Suppose that the
three generic normalization branches form one connected degree-three
cover $C\to H$, where $C$ denotes its smooth complete model.
Then $C$ is rational, and the cubic normal jet of $S$ vanishes
identically on at least one normal fibre over $H$.
In particular, a positive-genus connected degree-three cover is
impossible, and a primitive cubic normal jet is impossible for
any connected degree-three cover.
\end{proposition}


\begin{proof}

\begin{enumerate}[label={$\langle1\rangle$\arabic*.},leftmargin=*]
\item We construct all comparison maps, including the map into the
actual nearby fibre. Put $N=N_{H/Y}$ and
$M=\OO_Y(S)|_H$. The ambient Chern-class hypothesis gives
$M\simeq\OO_H$ and $\deg N=-2$. The order-three normal jet is
therefore an actual homogeneous function
\[
 P\in H^0(H,\operatorname{Sym}^3N^*),\qquad
 U=\{v\in N:P(v)=1\}.
\tag{C.D4.1}\hypertarget{tag-00E1}{}\label{cdp:eq:C.D4.1}
\]
The generic cubic has three distinct roots. The fibrewise Euler
identity $dP(v)(v)=3P(v)$ proves that $U$ is smooth everywhere,
including above a base point where the cubic degenerates. Above
a point where the entire cubic is zero there is no point of $U$.

\paragraph{Comparison with the actual fibre.}
Let $t$ be the base coordinate, so $f$ is a holomorphic function
on a neighbourhood of $H$. Locally $f=u s^m$, where $s=0$
defines $S$ and $u$ is a unit. The degree-$3m$ normal jet of
$f$ is $P^m$, after rescaling the chosen trivialization of $M$
by a constant. Indeed the ratio of the two jets is an invertible
holomorphic function on $H\simeq\PP^1$, hence constant.
Take a smooth tubular map $\tau$ from a neighbourhood of the
zero section in $N$ to $Y$, whose derivative on normal fibres
is the specified complex-linear identification. For positive
real $s\to0$, Taylor's formula gives
\[
 s^{-3m}f\bigl(\tau(sv)\bigr)\longrightarrow P(v)^m
                 \quad\hbox{in }C^\infty\hbox{ on compact subsets of }N.
\tag{C.D4.2}\hypertarget{tag-00EC}{}\label{cdp:eq:C.D4.2}
\]
The order of vanishing is at least $3m$ along every point of
$H$: its ideal is a regular codimension-two ideal, and its
ordinary powers are primary along the smooth curve, so the
order verified at its generic point extends over $H$.

\item Choose any compact codimension-zero smooth submanifold
$K\subset U$. On a neighbourhood of $K$, the differential of
$P^m$ has a smooth right inverse as a real map to $\C$.
The implicit-function theorem, uniformly on $K$, or the
corresponding right-inverse vector field, transports this
neighbourhood to the level
$s^{-3m}f(\tau(sv))=1$. Composing with $v\mapsto\tau(sv)$
gives an orientation-preserving embedding into the actual
fibre $f^{-1}(s^{3m})$. In particular, intersections of compact
two-cycles, computed with parallel perturbations inside $K$,
are preserved. This construction uses the actual function
$f$ and its leading normal jet; it makes no assertion that
arbitrary changes of local equations preserve a Milnor fibre.

\item The algebraic surface $U$ admits a compact smooth core whose
inclusion is a homotopy equivalence. For example, resolve the
boundary in a projective compactification, triangulate the
resulting pair, take the complement of a sufficiently small
regular neighbourhood of the boundary and round its corners.
The radial collars of this regular neighbourhood give the
retraction. We choose this core for $K$. Thus the preceding
embedding induces one map carrying all of $H_1(U)$ and
$H_2(U)$, together with their compact intersection pairing.

\item Here is the precise torus obstruction we will use. For a map
$K\to T^4$, let $W$ be the rational image of $H_1(K)$ in
$H_1(T^4)$. The map lifts to the covering corresponding to
its fundamental-group image, a space homotopy equivalent to
$T^r$, where $r=\dim W$. Its image on $H_2$ consequently lies
in $\Lambda^2W\subset\Lambda^2\Q^4$. For $r\le3$ the wedge
pairing on this image is zero, because $\Lambda^4W=0$.
Consequently a compact core with $b_1(U)\le3$ and nonzero
intersection pairing cannot embed by \eqref{cdp:eq:C.D4.2} into a nearby
torus. This is a comparison of the actual homology maps.

\paragraph{The normal bundle and the root divisor.}
Write
\[
                  N=\OO_H(a)\oplus\OO_H(-2-a),\qquad a\ge-1.
\tag{C.D4.3}\hypertarget{tag-00EE}{}\label{cdp:eq:C.D4.3}
\]
If $a>0$, the coefficient of the cube of the first fibre
coordinate vanishes, since it has degree $-3a$. The other
coordinate is then a global linear factor of the cubic.
This contradicts connectedness of the degree-three root
cover. Thus $a=-1$ or $a=0$.

\item For $a=0$, irreducibility of the generic cubic forces the
coefficient of the first coordinate cubed to be a nonzero
constant. Along the trivial summand of $N$, the equation
$P=1$ therefore has three constant sections. Each is a
compact $\PP^1$ in $U$ with normal bundle $\OO_{\PP^1}(-2)$:
the derivative in the trivial summand is nonzero, so the
implicit-function theorem identifies its normal direction
with the other summand. Its homology square is $-2$.
Transporting its tubular neighbourhood by \eqref{cdp:eq:C.D4.2} gives a
sphere of square $-2$ in a torus, whereas the torus has
$\pi_2=0$. This is impossible.

\item We are left with $N=\OO(-1)^{\oplus2}$ and
$E=\PP(N)\simeq\PP^1\times\PP^1$. Let $A$ be the common
zero divisor of the cubic coefficients on the base, of degree
$v$. Removing that scalar factor gives an integral horizontal
root divisor $R$ of bidegree $(3,3-v)$, where the first entry
is its degree over $H$. Its normalization is the given
cover $C$: at a general point the maps identify a normalization
branch with its distinct tangent line, and both complete
normal covers are the integral closure in this same function
field. Thus, writing $g=g(C)$,
\[
                       0\le g\le p_a(R)=4-2v.
\tag{C.D4.4}\hypertarget{tag-00EF}{}\label{cdp:eq:C.D4.4}
\]
For a primitive cubic $v=0$. For a positive-genus cover with
$v>0$, the bound forces $v=1$. We treat these two cases.
Their total root divisors are respectively
\[
 B=R_{(3,3)},\qquad
 B=R_{(3,2)}+F_{(0,1)}.
\tag{C.D4.5}\hypertarget{tag-00EG}{}\label{cdp:eq:C.D4.5}
\]
Both are reduced divisors of class $(3,3)$; the second has
one simple vertical component. They are connected. Let
$c=1$ or $2$ be their number of components, and set
$\nu=\sum_p(r_p(B)-1)$, where $r_p(B)$ is the number of
analytic branches at a singular point. Their total delta
invariants are $4-g$ and $5-g$, respectively, so in the
cases under consideration they are at most four.

\paragraph{The cyclic compactification and its local invariants.}
Put $L=\OO_E(1,1)$ and form the finite cyclic triple cover
\[
 X_0=\operatorname{Spec}_E
       (\OO_E\oplus L^{-1}\oplus L^{-2}),\qquad w^3=b,
\tag{C.D4.6}\hypertarget{tag-00EH}{}\label{cdp:eq:C.D4.6}
\]
where $b$ is the section of $L^3$ defining $B$. The algebra
multiplication is given by $b$. Because $B$ is reduced,
$X_0$ is normal: its hypersurface local rings are Cohen--Macaulay
and regular in codimension one. The cover is connected,
as a meridian around any smooth point of $B$ acts by a
three-cycle. Above $E\setminus B$ it is exactly $U$.
Indeed a line $[v]\in E\setminus B$ has three radial vectors
with $P(v)=1$; equivalently, the dual tautological coordinate
has cube $b$. This gives the claimed isomorphism, including
the transition functions.

\item Let $\rho:X\to X_0$ resolve the surface singularities and
then the reduced boundary $D=X\setminus U$. Let $Q_p$ denote
the reduced exceptional divisor above a singular point of
$X_0$. We need the following finite computation:
\[
  \sum_p p_g(X_0,p)\le1,\qquad
  \sum_p b_1(Q_p;\Q)\le2.
\tag{C.D4.7}\hypertarget{tag-00EI}{}\label{cdp:eq:C.D4.7}
\]
Here are its equations and derivation. A reduced plane germ
with delta invariant at most four has one of the forms
\[
\begin{array}{c|c|c}
 \hbox{type}&r(x,y)&\hbox{range}\\ \hline
 A_n&x^2+y^{n+1}&1\le n\le8\\
 D_n&x^2y+y^{n-1}&4\le n\le7\\
 E_6&x^3+y^4&\\
 E_7&x^3+xy^3&\\
 E_8&x^3+y^5&
\end{array}
\]
For completeness, multiplicity at least four gives
$\delta\ge\binom42=6$. At multiplicity two, Weierstrass
preparation and completing the square give $A_n$, with
$\delta=\lfloor(n+1)/2\rfloor$. At multiplicity three,
three different tangent lines give $D_4$. A double tangent
and a transverse line give that smooth line times an
$A_k$ branch or pair, adding intersection multiplicity two;
the delta bound gives $D_5,D_6,D_7$. With a single tangent
line, an irreducible branch has initial Puiseux pair $(3,4)$
or $(3,5)$, giving $E_6,E_8$; a smooth branch and a cusp
have total delta at least $1+3=4$, and equality gives $E_7$.
Three smooth branches with one tangent line already give
delta at least six. These reductions give the displayed
analytic forms; the higher terms in the last three forms
are eliminated by their positive weighted Jacobian terms.

\item The singularity of $X_0$ is $w^3=r(x,y)$. Normalize the
weights of $r$ to degree one and give $w$ weight $1/3$.
The adjunction calculation of
\cite[Theorem 2.1.1 and Corollary 2.1.3]{MerleTeissier1977}
counts $p_g$ by the positive lattice points under the
weight-one plane. Equivalently, it counts monomials
$x^iy^jw^k$ with
$(i+1)w_x+(j+1)w_y+(k+1)/3\le1$.
The normalized monodromy weights are those of the Jacobian
monomials, shifted by the sum of the variable weights
\cite[pp. 165--167]{Steenbrink1989Spectrum}.
The following table records the result; all entries are
also computed with rational arithmetic in
\path{code/root_cubic_normal_cone.py}.
\[
\begin{array}{c|c|c|c}
 \hbox{plane type}&\delta&p_g(w^3-r)&b_1(Q_p;\Q)\\ \hline
 A_1,A_2&1&0&0\\
 A_3,A_4&2&0&0\\
 A_5&3&1&2\\
 A_6&3&1&0\\
 A_7,A_8&4&1&0\\
 D_4&3&1&2\\
 D_5&3&1&0\\
 D_6,D_7&4&1&0\\
 E_6&3&1&0\\
 E_7,E_8&4&1&0
\end{array}
\]
To verify the last column directly, use the Jacobian bases
$\{y^j:0\le j<n\}$ for $A_n$;
$\{1,y,\ldots,y^{n-2},x\}$ for $D_n$;
$\{x^iy^j:0\le i<2,0\le j<3\}$ for $E_6$;
$\{1,y,\ldots,y^4,x,xy\}$ for $E_7$; and
$\{x^iy^j:0\le i<2,0\le j<4\}$ for $E_8$.
Count weights $(i+1)w_x+(j+1)w_y$ congruent to $1/3$ or
$2/3$ modulo one. The tensor-product monodromy for adding
$w^3$ identifies this count with the eigenvalue-one
multiplicity of the surface monodromy
\cite[\S1 and Example 2.4(a)]{Nemethi1991Join}.
The variation sequence of its simply connected Milnor fibre
identifies that multiplicity with the rational first Betti
number of its link. The negative-definite exceptional
intersection matrix then identifies $H_1$ of the link with
$H_1(Q_p)$: in the homology sequence of its plumbed
neighbourhood, the map on second homology is that invertible
matrix. This proves the column without replacing the actual
exceptional divisor by an unspecified graph.
Every nonzero entry in either last column costs delta at
least three. Since the total delta is at most four,
\eqref{cdp:eq:C.D4.7} follows.

\paragraph{Homology and the compact intersection pairing.}
The finite algebra in \eqref{cdp:eq:C.D4.6} and the K\"unneth formula give
$H^1(X_0,\OO_{X_0})=0$: the three summands on $E$ are
$\OO$, $\OO(-1,-1)$ and $\OO(-2,-2)$, all with $H^1=0$.
The resolution Leray sequence injects $H^1(X,\OO_X)$ into
$H^0(X_0,R^1\rho_*\OO_X)$, of dimension
$\sum p_g(X_0,p)$. Since $X$ is projective, \eqref{cdp:eq:C.D4.7} gives
\[
                         b_1(X;\Q)=2q(X)\le2.
\tag{C.D4.8}\hypertarget{tag-00EJ}{}\label{cdp:eq:C.D4.8}
\]
The irreducible components of $D$ have independent homology
classes. The exceptional classes have a negative-definite
intersection matrix. In the first case of \eqref{cdp:eq:C.D4.5} the
remaining class pushes forward to the nonzero class of $R$;
in the second, the two remaining classes push forward to
$(3,2)$ and $(0,1)$, which are independent on $E$.
Moreover their span contains the pullback of the ample
class $B=(3,3)$, whose square is positive. The Hodge index
theorem implies that the intersection form on this span is
nondegenerate: its part orthogonal to the positive pullback
class is negative definite.

\item Write $n_D$ for the number of boundary components. The exact
homology sequence of the pair $(X,U)$ identifies the map
$H_2(X;\Q)\to H_2(X,U;\Q)=\Q^{n_D}$ with intersection
against these components. Its surjectivity gives
$H_1(U;\Q)\simeq H_1(X;\Q)$. Its kernel is the orthogonal
complement of their span, and is the image of $H_2(U)$.
The pairing on that image is nondegenerate, of dimension
$b_2(X)-n_D$. As $D$ is connected,
$e(D)=n_D+1-b_1(D)$, and consequently
\[
\begin{split}
 b_1(U)&=b_1(X)\le2,\\
 \operatorname{rank}\bigl(\hbox{intersection pairing on }H_2(U)\bigr)
     &=b_2(X)-n_D\\
     &=e(U)-1+2b_1(X)-b_1(D).
\end{split}
\tag{C.D4.9}\hypertarget{tag-00EK}{}\label{cdp:eq:C.D4.9}
\]
Here the equality of ranks follows also for a degenerate
pairing on $H_2(U)$: its kernel includes the kernel of the
map to $H_2(X)$, and the induced form on its image is the
nondegenerate form just identified.

\item There is a useful exact version allowing relations among
the strict boundary classes. If there are $\varepsilon$
independent relations, while their span still contains a
positive-square class, the same pair sequence and the same
Hodge-index argument give
\[
\begin{split}
 b_1(U)&=b_1(X)+\varepsilon,\\
 \operatorname{rank}\bigl(\hbox{intersection pairing on }H_2(U)\bigr)
 &=e(U)-1+2b_1(X)-b_1(D)+\varepsilon.
\end{split}
\tag{C.D4.10}\hypertarget{tag-00E2}{}\label{cdp:eq:C.D4.10}
\]
These identities concern the actual classes in the indicated
cyclic cover. They will also be used for the degree-two
root divisor below.

\item To compute $b_1(D)$, start with the normalizations of the
$c$ components of $B$, of total genus $g$, and attach the
connected exceptional divisor over each singular point.
There are $r_p(B)$ attaching branches there. The incidence
graph contributes $\nu-c+1$ cycles. Subsequent point
blowups resolving the boundary subdivide edges or add trees
and preserve this first Betti number. Thus \eqref{cdp:eq:C.D4.7} gives
\[
                         b_1(D)\le2g+\nu-c+3.
\tag{C.D4.11}\hypertarget{tag-00E3}{}\label{cdp:eq:C.D4.11}
\]
The unramified cover $U\to E\setminus B$ has degree three,
so additivity and multiplicativity of Euler characteristic
give the exact value
\[
                         e(U)=3(4-2c+2g+\nu).
\tag{C.D4.12}\hypertarget{tag-00E4}{}\label{cdp:eq:C.D4.12}
\]
For $c=1$, \eqref{cdp:eq:C.D4.9}--\eqref{cdp:eq:C.D4.12} give intersection rank at
least $3+4g+2\nu>0$. For $c=2$ they give rank at least
$4g+2\nu-2$. In this second case $g\ge1$, and $R$ meets
$F$ with intersection number three, so $\nu\ge1$ and the
rank is again positive. Both cases have $b_1(U)\le2$.
They contradict the compact-core comparison with $T^4$
proved immediately after \eqref{cdp:eq:C.D4.2}.

\item This excludes $v=0$ for every connected triple cover and
$v=1$ for every positive-genus one. By \eqref{cdp:eq:C.D4.4} these are
all possibilities for a positive-genus cover, proving the
statement. The genus-zero cases with $v=1$ or $2$ are not
asserted to be excluded by this argument.

\end{enumerate}
\end{proof}






\begin{proposition}[A rational ordinary-triple image has no positive-genus double branch cover]\hypertarget{tag-01BX}{}\label{cdp:st:C.2.D4.2}
\StatementRecord{cdp:st:C.2.D4.2}{proved}{01BX}{\SeriesRef{cdp:st:C.2.D4.1}, \SeriesRef{cdp:st:A.5.19}, \SeriesRef{cdp:eq:C.DD.3}, \SeriesRef{cdp:eq:C.DD.6}, \SeriesRef{cdp:eq:C.DD.2}, \SeriesRef{cdp:eq:C.DD.9}, \SeriesRef{cdp:eq:C.DD.10}, \SeriesRef{cdp:eq:C.DD.8}, \SeriesRef{cdp:eq:C.DD.4}, \SeriesRef{cdp:eq:C.DD.16}, \SeriesRef{cdp:eq:C.DD.7}, \SeriesRef{cdp:eq:C.DD.20}, \SeriesRef{cdp:eq:C.DD.5}}

Let $f:Y\to C_0$ have the torus-fibre and ambient hypotheses
of \ref{cdp:st:C.2.D4.1}. Let $H\subset S$ be a smooth rational curve
with generic ordinary-triple transversal. Suppose two of
the three generic normalization branches form a connected
degree-two cover $C\to H$, and the remaining branch forms
a degree-one cover. Then the smooth complete curve $C$
is rational. The assertion includes all zeros, with their
multiplicities, of the entire cubic normal tensor.

Consequently, in the irrational ruled setup of \SeriesRef{cdp:st:A.5.19},
a noncontracted coefficient-two horizontal conductor
with generic ordinary-triple transversal has smooth
elliptic image. Its normalization is smooth minimal
elliptic ruled, the horizontal cover is unramified of
degree three, and $e(S)=0$.
\end{proposition}


\begin{proof}

\begin{enumerate}[label={$\langle1\rangle$\arabic*.},leftmargin=*]
\item Suppose $g=g(C)\ge1$. Use the actual normal tensor and
normal-cone comparison (C.D4.1)--(C.D4.2), and write
\label{layout:II:2}
\[
\begin{gathered}
N=N_{H/Y}=\OO(a)\oplus\OO(-2-a),\quad a\ge-1,
 \\
\qquad P\in H^0(H,\operatorname{Sym}^3N^*),
 \qquad U=\{P=1\}\subset N.
\end{gathered}
\tag{C.DD.1}\hypertarget{tag-00EL}{}\label{cdp:eq:C.DD.1}
\]
The comparison embeds a compact core of $U$, including
its compact intersection pairing, into an actual nearby
torus. As proved there, $b_1(U)\le3$ and a nonzero
intersection pairing contradict this embedding.

\item Put $E=\PP(N)$, using lines in $N$, and let
$L=\OO_E(1)$, $\xi=c_1(L)$ and $F$ be the fibre class.
These conventions give
\[
 \xi^2=-\deg N=2,\qquad \xi F=1,\quad F^2=0,
 \qquad K_E=-2\xi.
\tag{C.DD.2}\hypertarget{tag-00EW}{}\label{cdp:eq:C.DD.2}
\]
Let $A$ be the common zero divisor of the coefficients
of $P$, and put $v=\deg A$. The divisor of the induced
section of $L^3$ is
\[
                    B_{\rm tot}=R_1+R_2+p^*A.
\tag{C.DD.3}\hypertarget{tag-00EZ}{}\label{cdp:eq:C.DD.3}
\]
Here $R_1$ is a section and $R_2$ is an integral
degree-two curve over $H$. The section is smooth:
its rational tangent-root map extends over the smooth
complete base into the proper bundle $E$.
The normalization of $R_2$ is precisely $C$. Over the
generic point a normalization branch is identified
with its distinct tangent root; taking the integral
closure in this same quadratic function field identifies
the two complete normal covers, as in (C.RD.2).

\paragraph{All possible common divisors.}
For $a\ge0$ and $v>0$, the pure cube of the coordinate
in $\OO(a)$ has coefficient of degree $-3a-v$, and
therefore vanishes. The other coordinate is a factor.
The quadratic factor is irreducible over the function
field. Its coefficient of the first coordinate squared
must be nonzero; that coefficient has degree
\[
                              2-a-v.
\tag{C.DD.4}\hypertarget{tag-00F0}{}\label{cdp:eq:C.DD.4}
\]
Thus $a+v\le2$. In particular, $v\le2$, and $v=2$
forces $a=0$.

\item For $a=-1$, identify $E$ with the product of the base
and the projective direction line. Write $d\ge0$ for
the degree of the map defining the section $R_1$.
The bidegrees, with degree over $H$ first, are
\[
 R_1:(1,d),\qquad R_2:(2,3-v-d),\qquad
 p_a(R_2)=2-v-d.
\tag{C.DD.5}\hypertarget{tag-00F1}{}\label{cdp:eq:C.DD.5}
\]
The inequality $g\le p_a(R_2)$ gives $v+d\le1$.
Consequently the only common-divisor possibilities
for our positive-genus cover are $v=0$, $v=1$, and
$v=2$ with $N=\OO\oplus\OO(-2)$.

\paragraph{The exact cyclic-cover homology formulas.}
We record the comparison with relations and disconnected
boundary retained. First suppose that \eqref{cdp:eq:C.DD.3} is reduced,
and let $B=B_{\rm tot}$. The cyclic cover
\[
 \pi:X_0=\operatorname{Spec}_E
       (\OO_E\oplus L^{-1}\oplus L^{-2})\longrightarrow E
\tag{C.DD.6}\hypertarget{tag-00F2}{}\label{cdp:eq:C.DD.6}
\]
has multiplication given by $P$. It is normal, since
the branch divisor is reduced. Its complement over
$B$ is exactly $U$: if $w\in\ell^*$ is a nonzero
cube root of $P|_\ell$, its corresponding vector in
the line $\ell\subset N$ is the unique vector $z$
with $w(z)=1$. This description specifies inverse
maps and their bundle transitions.

\item All three summands in \eqref{cdp:eq:C.DD.6} have vanishing $H^1$.
Indeed the relative fibre cohomology of $L^{-1}$
vanishes, while relative duality gives
\[
 p_*L^{-2}=0,\qquad R^1p_*L^{-2}=\det N=\OO(-2).
\tag{C.DD.7}\hypertarget{tag-00F3}{}\label{cdp:eq:C.DD.7}
\]
Thus $H^1(E,L^{-2})=H^0(H,\OO(-2))=0$;
$H^1(E,\OO_E)=0$ as well. For a resolution
$\rho:X\to X_0$, including resolution of the reduced
boundary $D=X\setminus U$, Leray therefore gives
\[
 b_1(X)=2q(X)\le2\sum_x p_g(X_0,x).
\tag{C.DD.8}\hypertarget{tag-00F4}{}\label{cdp:eq:C.DD.8}
\]
The resolution is projective, so the equality here
is Hodge decomposition on the actual surface $X$.

\item Write $c$ for the number of components of $B$,
$b_0=b_0(B)$, and
$\nu=\sum_x(r_x(B)-1)$. Its normalized components
have total genus $g$. Let $r$ be the rank of their
classes in $H_2(E,\Q)$ and put $\varepsilon=c-r$.
The boundary classes on $X$ have exactly
$\varepsilon$ relations. To see the equality, modulo
the independent exceptional classes, the strict
component over $B_i$ is $\pi^*[B_i]/3$.
Indeed its ramification index is three, and the
remaining terms in the divisor pullback are exceptional.
The pullback of $H_2(E,\Q)$ is injective and its
intersection pairing is multiplied by three.
This proves both directions of the asserted relation
count. The span of the boundary classes is
nondegenerate: it contains the positive-square
pullback of $B_{\rm tot}=3\xi$, and the Hodge index
theorem makes the part orthogonal to that class
negative definite.

\item The pair sequence, exactly as in (C.D4.9)--(C.D4.10),
now gives the version with $b_0$ retained:
\[
\begin{split}
 b_1(U)&=b_1(X)+\varepsilon,\\
 \operatorname{rank}\langle\ ,\ \rangle_U
 &=e(U)+b_0-2+2b_1(X)-b_1(D)+\varepsilon.
\end{split}
\tag{C.DD.9}\hypertarget{tag-00F5}{}\label{cdp:eq:C.DD.9}
\]
The second identity follows from
$e(D)=b_0+n_D-b_1(D)$ and the nondegenerate
orthogonal complement of the boundary span. Here
$\langle\ ,\ \rangle_U$ is the compact intersection
pairing, so it is precisely the form preserved by
the normal-cone embedding.

\item If $\lambda$ is the sum of the first Betti numbers
of the reduced exceptional divisors over the original
surface singularities, normalization of $B$ and the
attaching-branch graph give
\[
 b_1(D)\le2g+\nu-c+b_0+\lambda,
 \qquad e(U)=3(4-2c+2g+\nu).
\tag{C.DD.10}\hypertarget{tag-00EM}{}\label{cdp:eq:C.DD.10}
\]
Each exceptional divisor is connected. Its attachments
number $r_x(B)$ at a singular point, giving graph
cycle rank $\nu-c+b_0$. Further boundary blowups
subdivide edges or add trees. The Euler identity
uses the actual unramified degree-three cover
$U\to E\setminus B$ and
$e(B)=2c-2g-\nu$.

\paragraph{No common zero.}
For $v=0$, the divisor $B=R_1+R_2$ is reduced,
has $c=2$, and has arithmetic genus four by
\eqref{cdp:eq:C.DD.2}. Its total delta invariant is therefore
\[
                         \delta(B)=5-g\le4.
\tag{C.DD.11}\hypertarget{tag-00EN}{}\label{cdp:eq:C.DD.11}
\]
The local calculation (C.D4.7), or \CompendiumRef{II}{cdp:st:C.2.RAT.6},
gives $\sum p_g\le1$ and $\lambda\le2$.
The rank $r$ is one or two, so
$\varepsilon=2-r\le1$. Thus $b_1(U)\le3$.
Substitution in \eqref{cdp:eq:C.DD.9}--\eqref{cdp:eq:C.DD.10} gives
\[
                  \operatorname{rank}\langle\ ,\ \rangle_U
                        \ge4g+2\nu-r>0.
\tag{C.DD.12}\hypertarget{tag-00EO}{}\label{cdp:eq:C.DD.12}
\]
This calculation includes a disconnected $B$ and
dependent section classes. It contradicts the
normal-cone comparison.

\paragraph{One simple common zero.}
For $v=1$, one has $B=R_1+R_2+F_0$, still reduced.
It is connected because the fibre meets both
horizontal components. Here $c=3$, $b_0=1$,
$\nu\ge2$, and
$\delta(B)=6-g\le5$.
Statement \CompendiumRef{II}{cdp:st:C.2.RAT.6} supplies
$\sum p_g\le1$ and $\lambda\le2$.
The section and the fibre have independent classes,
so $r=2$ and $\varepsilon=1$.
Equations \eqref{cdp:eq:C.DD.8}--\eqref{cdp:eq:C.DD.10} yield
\[
 b_1(U)\le3,\qquad
 \operatorname{rank}\langle\ ,\ \rangle_U
                         \ge4g+2\nu-6\ge2.
\tag{C.DD.13}\hypertarget{tag-00EP}{}\label{cdp:eq:C.DD.13}
\]
This is the same contradiction.

\paragraph{The exact form when the common divisor has degree two.}
Now $N=\OO\oplus\OO(-2)$ by \eqref{cdp:eq:C.DD.4}.
Let $x$ and $y$ be its projections, valued in
$\OO$ and $\OO(-2)$ respectively. The coefficient
of $x^2$ in the irreducible quadratic is a nonzero
constant. Completing its square is a global bundle
automorphism: the mixed coefficient is a section
of $\OO(2)$ and therefore defines an actual map
$\OO(-2)\to\OO$. After a nonzero constant rescaling,
the tensor has the form
\[
 P=a\,y\,(x^2-b y^2),\qquad
 a\in H^0(H,\OO(2)),\quad b\in H^0(H,\OO(4)).
\tag{C.DD.14}\hypertarget{tag-00EQ}{}\label{cdp:eq:C.DD.14}
\]
Its horizontal root divisors are
$R_1=\{y=0\}$ and $R_2=\{x^2-b y^2=0\}$.
They are disjoint. The quadratic cover has branch
degree at most four. Riemann--Hurwitz and $g\ge1$
force $g=1$, branch degree four, and four distinct
simple zeros of $b$. Thus $R_2$ is already a smooth
elliptic curve. A ruling fibre meets it either at
two transverse points or at a single point of
contact two. This verifies the local models used
in both remaining cases.

\item If $p^*A=F_1+F_2$ with distinct fibres, the reduced
branch divisor has four components. Its only plane
singularities are ordinary nodes and tangencies of
order two, of types $A_1$ and $A_3$.
Their cubic suspensions have geometric genus zero
and exceptional first Betti number zero, by
(C.RD.10). Consequently $b_1(X)=0$.
The branch divisor is connected, its class rank
is two, and $4\le\nu\le6$.
Thus $\varepsilon=2$, $b_1(D)=\nu-1$, and
\[
 b_1(U)=2,\qquad
 \operatorname{rank}\langle\ ,\ \rangle_U
                       =2\nu-4\ge4.
\tag{C.DD.15}\hypertarget{tag-00ER}{}\label{cdp:eq:C.DD.15}
\]
Again the torus comparison is impossible.

\paragraph{A double common zero: normalization of the cyclic cover.}
It remains to take $p^*A=2F_0$. The branch divisor
with multiplicities is $R_1+R_2+2F_0$; its reduced
support $B$ has three components. We use the
normalization of the cyclic cover. Its exact
eigensheaf algebra has underlying module
\[
       \pi_*\OO_{X_0}
       =\OO_E\oplus L^{-1}\oplus L^{-2}(F_0).
\tag{C.DD.16}\hypertarget{tag-00ES}{}\label{cdp:eq:C.DD.16}
\]
At the generic point of $F_0$, write the equation
as $w^3=s^2u$ with $u$ a unit. The extra integral
generator is $w^2/s$. This gives the last summand
in \eqref{cdp:eq:C.DD.16}. At the other branch components the
usual generators $1,w,w^2$ apply. These comparisons
extend across the finitely many crossings: the
normal finite surface is Cohen--Macaulay, so its
finite direct image over the smooth surface $E$
is locally free, and each character summand is
the rank-one reflexive extension of its divisorwise
description. This proves \eqref{cdp:eq:C.DD.16} as a sheaf
identity, not only over the nonsingular locus.

\item Relative duality as in \eqref{cdp:eq:C.DD.7} now gives
\[
 R^1p_*\bigl(L^{-2}(F_0)\bigr)
                       =\det N\otimes\OO(1)=\OO(-1).
\tag{C.DD.17}\hypertarget{tag-00ET}{}\label{cdp:eq:C.DD.17}
\]
Thus every summand in \eqref{cdp:eq:C.DD.16} has vanishing
$H^1$, and $H^1(X_0,\OO_{X_0})=0$.
We next prove that all singularities of this
normalized cover are rational.

\item The only intersections are between the weight-two
branch $F_0$ and a weight-one smooth branch. Their
contact is one or two. The following explicit
resolution works for every contact $k\ge1$.
In coordinates the branch equations are
\[
                F_0=(s=0),\qquad R=(s-u^k=0),
                \qquad w^3=s^2(s-u^k).
\tag{C.DD.18}\hypertarget{tag-00EU}{}\label{cdp:eq:C.DD.18}
\]
Blow up their common point until the strict
branches separate, requiring $k$ blowups. At
each step their total branch multiplicity is
$2+1=3$. The exceptional divisor is therefore
unbranched after normalization, and the new
root line and weight-two divisor are exactly
\[
 L'=\sigma^*L(-E_0),\qquad
 F'=\sigma^*F-E_0.
\tag{C.DD.19}\hypertarget{tag-00EV}{}\label{cdp:eq:C.DD.19}
\]
For example $s=u s_1$, $w=u w_1$ changes the
equation to
$w_1^3=s_1^2(s_1-u^{k-1})$; the extra normalized
generator changes by
$w^2/s=u(w_1^2/s_1)$.
Both nontrivial eigensheaves consequently satisfy
the exact identities
\[
 (L')^{-1}=\sigma^*(L^{-1})(E_0),\qquad
 (L')^{-2}(F')=\sigma^*(L^{-2}(F))(E_0).
\tag{C.DD.20}\hypertarget{tag-00EX}{}\label{cdp:eq:C.DD.20}
\]
For a point blowup of a smooth surface,
$\sigma_*\OO(E_0)=\OO$ and
$R^1\sigma_*\OO(E_0)=0$.
The first follows by extension of functions across
the centre; the second follows from
$0\to\OO\to\OO(E_0)\to\OO_{\PP^1}(-1)\to0$
and the structure-sheaf direct images of the
point blowup. Projection formula applied to
\eqref{cdp:eq:C.DD.20} therefore makes $R^1$ of the finite
cyclic algebra vanish at every step.

\item For precision, the normalized covers over the
successive blown-up bases have proper birational
maps to $X_0$, induced by their common function
field. Finite direct image and Leray identify the
pushforward of the $R^1$ of each such map with
the just computed $R^1$ of its eigensheaf algebra.
The finite pushforward is faithful, so that $R^1$
is zero. At the end the two branch curves are
smooth and disjoint, and the normalized cover
is smooth: a single branch of multiplicity one
or two is parameterized by $s=t^3$, with $w=t$
or $w=t^2$, up to a unit root. This is a resolution
with $R^1\OO=0$ and proves the claimed rationality.
Its connected exceptional divisors have rational
components and tree incidence graphs: their
$H^1(\OO)$ is a quotient of that $R^1$ and is zero.
Hence their first rational Betti numbers are zero.

\item In our global double-fibre case it follows that
$b_1(X)=0$ and $\lambda=0$.
The comparison $U\to E\setminus B$ is still an
unramified degree-three cover, since normalization
does not change it off the branch divisor.
Each reduced branch component still has one
strict transform with ramification index three,
so the relation and intersection calculations
\eqref{cdp:eq:C.DD.9}--\eqref{cdp:eq:C.DD.10} apply to this reduced support.
It is connected, $c=3$, its class rank is two,
and $\nu=2$ or $3$. Thus
\[
 b_1(U)=1,\qquad b_1(D)=\nu,\qquad e(U)=3\nu,
 \qquad\operatorname{rank}\langle\ ,\ \rangle_U
                                      =2\nu>0.
\tag{C.DD.21}\hypertarget{tag-00EY}{}\label{cdp:eq:C.DD.21}
\]
This is the final contradiction. Every common
divisor allowed by \eqref{cdp:eq:C.DD.4}--\eqref{cdp:eq:C.DD.5} has now
been treated, proving that $C$ cannot have
positive genus.

\item For the stated irrational ruled consequence,
\SeriesRef{cdp:st:A.5.19} leaves a smooth rational or smooth elliptic
image. A rational image would have horizontal
degree two or three. The preceding proof excludes
the first, and \ref{cdp:st:C.2.D4.1} excludes the second.
The elliptic conclusion, its exact minimality
comparison, and its Euler number are those already
proved in \SeriesRef{cdp:st:A.5.19}.

\end{enumerate}
\end{proof}



The splitting bounds, all intersection-rank lower
bounds, the exact scalar-zero coordinate example,
and the normalized cyclic-cover blowup substitutions
are retained in \path{code/E_cubic_double.py} and
its saved output. The local delta-five table is
independently computed in \path{code/A_cubic_delta5.py}.




\begin{proposition}[Rational connected triple covers have positive fibre Euler number]\hypertarget{tag-01BY}{}\label{cdp:st:C.2.D4.3}
\StatementRecord{cdp:st:C.2.D4.3}{proved}{01BY}{\SeriesRef{cdp:st:C.2.RAT.1}, \SeriesRef{cdp:st:C.2.D4.1}, \SeriesRef{cdp:st:C.2.RAT.3}, \SeriesRef{cdp:st:C.2.RAT.5}, \SeriesRef{cdp:eq:C.D4.13}, \SeriesRef{cdp:eq:C.D4.19}}

Retain the torus-fibre, rational-normalization and specified
Hirzebruch-boundary morphism of \ref{cdp:st:C.2.RAT.1}. Suppose the horizontal
section $C$ is not contracted, its image $H$ is smooth rational,
the generic transversal along $H$ is an ordinary triple point,
and $C\to H$ has degree three. Then $e(S)\ge1$.
\end{proposition}


\begin{proof}

\begin{enumerate}[label={$\langle1\rangle$\arabic*.},leftmargin=*]
\item Use the normal-cone construction of \ref{cdp:st:C.2.D4.1}. That result forces
$N=\OO(-1)^{\oplus2}$ and a common cubic-coefficient divisor
of degree $v=1$ or $2$. The horizontal root curve is integral
of bidegree $(3,3-v)$, with rational normalization $C$.
No vertical strict conductor component dominates $H$, since
all three generic branches are supplied by $C$.

\item We first recall exactly how a finite attachment bound gives
the Euler inequality. In the notation of \ref{cdp:st:C.2.RAT.1} let $r$
be the number of original positive ruling fibres and let
$N_0$ count coefficient-zero blowup events. At least
$K\ge r-N_0$ of their positive boundary trees contain a
strict conductor curve. The negative-definite argument of
\ref{cdp:st:C.2.RAT.3} proves this bound without an assumption about
normalization singularities. They yield $K$ distinct
attachment points on $C$. Their images on $H$ give $k$
distinct normalization-incidence edges. Put $\ell=K-k$.
The incidence graph has first Betti number at most one,
so the exact graph count (C.RA.24) gives
$e(\Gamma_{\rm red})\ge k$. Consequently
\[
             e(S)=2-r+N_0+e(\Gamma_{\rm red})\ge2-\ell.
\tag{C.D4.13}\hypertarget{tag-00E5}{}\label{cdp:eq:C.D4.13}
\]

\item At a point where the cubic is primitive, a simple tangent
root gives a smooth normalization branch with conductor
exactly twice the horizontal curve. Indeed the ambient
blowup chart divides the reduced equation by the third
power of its normal coordinate; the nonzero root derivative
solves the remaining coordinate holomorphically. The
remaining quadratic factor has nonzero value at the chosen
simple root. Its restriction therefore gives the conductor
factor $y^2$ times a unit, even when its own two roots coincide. This is the explicit
factorization and conductor comparison of \ref{cdp:st:C.2.RAT.5}.
Such a point cannot carry a vertical boundary attachment.
If the horizontal root curve is locally unibranch, a
multiple root has just one point on its normalization.
A cubic has at most one multiple root unless all its roots
are simple. Thus off the common divisor each image value
receives at most one attachment when that root curve is
locally unibranch.

\paragraph{One simple common zero.}
Here the total reduced branch divisor is
$B=R_{(3,2)}+F_{(0,1)}$, with two rational normalizations
and total delta five. The complete computation \CompendiumRef{II}{cdp:st:C.2.RAT.6}
gives $P:=\sum p_g\le1$ and
$J:=\sum b_1(Q_p)\le2$. The classes of $R,F$ are
independent. As in (C.D4.8)--(C.D4.12), writing
$\nu=\sum(r_p(B)-1)$, we obtain
\[
 b_1(U)\le2,\qquad e(U)=3\nu,\qquad
 \operatorname{rank}(\hbox{intersection pairing})
                         =2\nu+2b_1(X)-J.
\tag{C.D4.14}\hypertarget{tag-00E6}{}\label{cdp:eq:C.D4.14}
\]
The torus comparison excludes $\nu\ge2$. Since $R\cdot F=3$,
we have $\nu\ge1$, so $\nu=1$. All branches of $R$ are
then locally unibranch: a multibranch point would contribute
to $\nu$ in addition to the intersection with $F$, or would
raise its contribution at that intersection to at least two.
There is exactly one point of $R\cap F$, and it has just
one preimage on $C$. Thus the common value also receives
at most one attachment. The preceding simple-root argument
handles every other value. Hence $\ell=0$ and \eqref{cdp:eq:C.D4.13}
gives $e(S)\ge2$.

\paragraph{Two distinct common zeros.}
Now $B=R_{(3,1)}+F_1+F_2$. The curve $R$ is smooth rational,
and each $R\cap F_i$ has total intersection number three.
A contact of order $k=1,2,3$ is the plane type $A_{2k-1}$.
Let $P$ be the number of order-three contacts. The local
table gives precisely
\[
              \sum p_g=P,\qquad \sum b_1(Q_p)=2P,
                         \qquad 0\le P\le2.
\tag{C.D4.15}\hypertarget{tag-00E7}{}\label{cdp:eq:C.D4.15}
\]
The cyclic cover $X_0$ has $\omega_{X_0}\simeq\OO_{X_0}$
and $\chi(\OO_{X_0})=2$, directly from its cyclic algebra
and adjunction. Proper duality identifies the quotient
$\omega_{X_0}/\rho_*\omega_X$ with a finite-length dual
module of length $P$. If $P>0$, a generator of the trivial
canonical line does not belong to the proper adjunction
ideal at at least one point. Thus no nonzero global
canonical section extends to $X$, and $p_g(X)=0$.
Leray gives $\chi(\OO_X)=2-P$. If $P=0$, the canonical
section extends and the same calculation gives $q(X)=0$.
Together these identities yield
\[
                         q(X)=\max(P-1,0).
\tag{C.D4.16}\hypertarget{tag-00E8}{}\label{cdp:eq:C.D4.16}
\]
There is exactly one relation among the three strict
boundary classes: $F_1,F_2$ have the same class, while
$R$ and a fibre are independent. The exceptional matrix
is negative definite, so no further relation occurs.
Thus (C.D4.10) gives $b_1(U)=2q(X)+1\le3$. With
$\nu=\#(R\cap(F_1\cup F_2))$, the boundary topology
and Euler number give
\[
 \operatorname{rank}(\hbox{intersection pairing})
          =2\nu-4-2P+4\max(P-1,0).
\tag{C.D4.17}\hypertarget{tag-00E9}{}\label{cdp:eq:C.D4.17}
\]
We check every possibility. For $P=0$ each fibre has at
least two distinct intersections, so $\nu\ge4$ and the
rank is positive, which is excluded. For $P=1$, one fibre
has a single order-three contact. The other has either
three simple intersections or one contact of order two
and one simple intersection. The first has $\nu=4$ and
positive rank; the only remaining value is $\nu=3$.
For $P=2$, each fibre has one point of intersection and
$\nu=2$. These are all the cases.

\item Since $R$ is smooth, the simple-root argument gives at
most one attachment at any noncommon value. A common
value has as many points on $C$ as the corresponding
intersection with $F_i$. In the remaining $P=1$ case
these numbers are one and two, so $\ell\le1$. In the
$P=2$ case they are both one, so $\ell=0$.
Equation \eqref{cdp:eq:C.D4.13} gives $e(S)\ge1$ and $e(S)\ge2$,
respectively.

\paragraph{A double common zero.}
Finally the cubic divisor is $R_{(3,1)}+2F$. The normalized
cyclic cover has the exact eigensheaf algebra
\[
 \pi_*\OO_{X_0}
       =\OO_E\oplus L^{-1}\oplus L^{-2}(F),
       \qquad L=\OO_E(1,1).
\tag{C.D4.18}\hypertarget{tag-00EA}{}\label{cdp:eq:C.D4.18}
\]
One obtains the last summand by dividing the square of the
cyclic coordinate by a local equation of the weight-two
branch; this integral element satisfies its monic cubic
relation. These summands describe the normalization in
codimension one and hence everywhere by normality. Their
$H^1$ groups vanish, since their bidegrees are $(0,0)$,
$(-1,-1)$ and $(-2,-1)$.

\item We verify that the normalization singularities are rational.
At a contact of order $k\le3$ the local cover is the
normalization of
\[
                            w^3=s^2(s-u^k).
\tag{C.D4.19}\hypertarget{tag-00EB}{}\label{cdp:eq:C.D4.19}
\]
Blow up the common point of the two smooth branch curves
repeatedly, $k$ times, until their strict transforms are
disjoint. At each step $\sigma$, the branch multiplicity
is three, so the new cyclic line and weight-two branch
are $L'=\sigma^*L(-E)$ and $F'=\sigma^*F-E$; the exceptional
curve has branch weight zero. Both nontrivial normalized
eigensheaves on the blown-up base are therefore the
pullbacks of their predecessors tensored with $\OO(E)$.
For a point blowup,
$\sigma_*\OO(E)=\OO$ and $R^1\sigma_*\OO(E)=0$, as follows
from $0\to\OO\to\OO(E)\to\OO_E(-1)\to0$.
Finite-cover Leray then proves that the induced birational
map of the normalized covers has $R^1\OO=0$ at each step.
The final cover is smooth: each remaining branch is smooth,
the branches are disjoint, and a single branch of either
weight one or two has smooth normalization. Composing
these maps proves rationality of every singularity in
\eqref{cdp:eq:C.D4.19}. In particular $q(X)=0$, and their exceptional
configurations have $b_1=0$ (rational singularities have
rational-tree exceptional divisors, also seen here from
the indicated successive cyclic covers of $E$).

\item The reduced boundary has two components with independent
classes, and contains the pullback of the ample class
$(3,2)$. If $\nu=\#(R\cap F)$, the same exact homology
calculation now gives
\[
        b_1(U)=0,\qquad e(U)=3\nu,\qquad
        b_1(D)=\nu-1,\qquad
        \operatorname{rank}(\hbox{intersection pairing})=2\nu>0.
\tag{C.D4.20}\hypertarget{tag-00ED}{}\label{cdp:eq:C.D4.20}
\]
This contradicts (C.D4.2), and excludes the double common
zero. The three cases exhaust $v=1,2$, proving $e(S)\ge1$.

\end{enumerate}
\end{proof}






\begin{proposition}[The irreducible local comparison and its rational and ruled consequences]\hypertarget{tag-01EG}{}\label{cdp:st:C.2.UN.7}
\StatementRecord{cdp:st:C.2.UN.7}{proved}{01EG}{\SeriesRef{cdp:st:D.1.9}, \SeriesRef{cdp:eq:C.UI.2}, \SeriesRef{cdp:eq:C.UI.3}, \SeriesRef{cdp:st:A.5.3}, \SeriesRef{cdp:eq:C.UI.4}, \SeriesRef{cdp:eq:C.UI.5}, \SeriesRef{cdp:eq:C.UI.6}, \SeriesRef{cdp:eq:C.UI.7}, \SeriesRef{cdp:st:C.2.MH.2}, \SeriesRef{cdp:st:A.5.11}, \SeriesRef{cdp:st:C.2.MH.1}, \SeriesRef{cdp:st:C.2.UN.5}, \SeriesRef{cdp:st:A.5.16}, \SeriesRef{cdp:st:A.5.17}, \SeriesRef{cdp:eq:C.UI.12}, \SeriesRef{cdp:eq:C.UI.11}, \SeriesRef{cdp:st:C.2.UN.1}, \SeriesRef{cdp:st:C.2.IRR.1}, \SeriesRef{cdp:st:A.5.14}, \SeriesRef{cdp:st:A.5.18}, \SeriesRef{cdp:st:A.5.20}, \SeriesRef{cdp:st:A.5.5}, \SeriesRef{cdp:st:A.5.19}, \SeriesRef{cdp:st:C.2.D4.1}, \SeriesRef{cdp:st:C.2.D4.2}, \SeriesRef{cdp:eq:C.CR.5}, \SeriesRef{cdp:eq:C.CR.6}, \SeriesRef{cdp:st:C.2.CR.1}, \SeriesRef{cdp:st:A.5.10}, \SeriesRef{cdp:eq:C.UI.15}; \cite[\S2 and \S3]{ACampo1975}}

Let $f:Y\to B$ be a proper holomorphic map from a smooth
complex threefold to a smooth curve, with connected fibres
and smooth general fibre a complex two-torus. Let $mS$ be a full
fibre with irreducible reduced support, and let $F$ be a nearby
smooth fibre. Write $T_1$ for the local monodromy on
$H^1(F,\mathbb Q)$. Then
\[
 m\,e(S_{\rm reg})=\det(1-T_1^m)\ge0.
 \tag{C.UI.1}\hypertarget{tag-014M}{}\label{cdp:eq:C.UI.1}
\]
For its normalization and minimal normalization resolution
$h=\eta\pi:Z\to T\to S$, the following statements hold.

\begin{enumerate}
\item If $Z$ is rational, then $e(S)\ge2$. More precisely,
      whenever $H^1(S,\mathbb Q)\ne0$, one has $e(S)=2$.
\item Suppose $Z$ is irrational ruled. A reduced horizontal
      conductor gives a smooth minimal elliptic ruled
      normalization, with disjoint smooth elliptic conductor
      curves and their finite etale double-cover maps to the
      conductor image. The support has ordinary double-curve
      normal crossings and Euler number zero.
\item Suppose the horizontal conductor is twice a noncontracted
      section. The normalization is again smooth minimal
      elliptic ruled. Its conductor is twice that section;
      its image is smooth elliptic, and the map from the
      section is etale of degree $1$, $2$, or $3$.
      The respective constant transverse types are the
      ordinary cusp, tacnode, and ordinary triple point.
      Every such support has Euler number zero.
\item Suppose the horizontal conductor is twice a section
      contracted by $\pi$. Then $e(S)\ge2$ and
      $\pi_1(S)=1$. In particular this case cannot occur
      under the full CDP hypotheses, since \SeriesRef{cdp:st:D.1.9} supplies
      a nonzero integral degree-one class on every fibre
      support there.
\end{enumerate}
The alternatives in the last three items exhaust the
irrational ruled case. No vanishing of $b_1(Y)$ or $b_2(Y)$
is used in the first four conclusions.
\end{proposition}


\begin{proof}\leavevmode\par



 
\begin{enumerate}[label={$\langle1\rangle$\arabic*.},leftmargin=0pt,labelsep=.5em]
\item \medskip\noindent\textbf{The actual canonical and Chern-class comparisons.}\quad
Choose a coordinate $t$ on a small disc $\Delta$ about the
chosen value and put $U=f^{-1}(\Delta)$. The torsion-free
rank-one sheaf $f_*\omega_{U/\Delta}$ is a line bundle on
$\Delta$. Its rank is one by base change on the punctured
disc, where a two-torus has a one-dimensional space of
holomorphic top forms. A generator evaluates to a
holomorphic section of $\omega_{U/\Delta}$ with no zeros
on any smooth nearby fibre: on that fibre it is a nonzero
holomorphic two-form. Its divisor is therefore $kS$ for
some integer $k\ge0$. The coordinate $dt$ identifies the
relative canonical line with $K_U$. The actual fibre
equation and adjunction give
\[
 K_U\simeq\OO_U(kS),\qquad
 \OO_U(mS)\simeq\OO_U,\qquad
 \omega_S^{\otimes m}\simeq\OO_S.
 \tag{C.UI.2}\hypertarget{tag-014U}{}\label{cdp:eq:C.UI.2}
\]
This restricts a section of the relative canonical line;
it does not pull a three-form back as a differential form
to a surface.


\begin{enumerate}[label={$\langle2\rangle$\arabic*.},leftmargin=1.3em,labelsep=.5em]
\item  Here is the Euler calculation on this actual tube.
The map is flat: its local rings are torsion-free
over the regular one-dimensional base ring.
Proper flatness makes the holomorphic Euler
characteristic of its full fibres constant, so
$\chi(\OO_{mS})=\chi(\OO_F)=0$.
The Cartier filtration of $\OO_{mS}$ has successive
quotients $\OO_S(-jS)$, $0\le j<m$; by \eqref{cdp:eq:C.UI.2}
these are torsion line-bundle twists.
Each such twist has the same Euler characteristic
as $\OO_S$. One can check this without a compact
ambient manifold: use the normalization sequence
and the resolution Leray sequence to compute the
Euler characteristic on the compact smooth surface
$Z$, the conductor curve, and the finite-support
sheaf $R^1\pi_*\OO_Z$. Smooth-surface
Riemann--Roch leaves the first term unchanged
under a torsion twist; curve Riemann--Roch,
applied to the generic-rank filtration of the
conductor quotient, leaves the second unchanged;
the finite lengths in the third are unchanged.
These exact sequences prove the assertion on $S$.
Consequently, writing $c_i=c_i(TU)$, one has
\[
 \begin{split}
 0&=\chi(\OO_{mS})=m\chi(\OO_S),\\
 c_1(h^*TU)&=0\quad\hbox{in }H^2(Z,\mathbb Q),
 \qquad c_2(h^*TU)=0\quad\hbox{in }H^4(Z,\mathbb Q).
 \end{split}
 \tag{C.UI.3}\hypertarget{tag-014V}{}\label{cdp:eq:C.UI.3}
\]
The first Chern-class assertion follows from \eqref{cdp:eq:C.UI.2}.
For the second we specify the local cycle comparison.
The inverse image of a radial segment $[0,t_0]$
is a compact triangulable real five-dimensional
space. Over the open segment it has the orientation
of the complex fibres followed by the radial
orientation. At a smooth point of $S$, the exact
local equation $t=u z^m$ gives $m$ radial sheets
with the same boundary orientation on $S$.
The remaining central singular locus has real
dimension at most two and contributes no
four-dimensional boundary. The resulting oriented
fundamental chain therefore has boundary
$[F]-m[S]$. This proves $[F]=m[S]$ in
$H_4(U,\mathbb Q)$, with its multiplicity retained.
The proper map $h$ has degree one onto $S$, hence
\[
 \int_Z h^*c_2=\langle c_2,[S]\rangle
       =\frac1m\int_F c_2(TU|_F)=0.
\]
The last equality follows from
$0\to T_F\to TU|_F\to\OO_F\to0$ and the
triviality of $T_F$. Top rational cohomology of
the connected closed oriented four-manifold $Z$
is one-dimensional, proving the second class
equality in \eqref{cdp:eq:C.UI.3}.


\item  Finite duality and the minimal-resolution calculation
\SeriesRef{cdp:st:A.5.3} now give the actual effective integral divisor $D$
with
\[
                    K_Z(D)\simeq h^*\omega_S.
 \tag{C.UI.4}\hypertarget{tag-014W}{}\label{cdp:eq:C.UI.4}
\]
Thus that line is torsion. Equations \eqref{cdp:eq:C.UI.2}--\eqref{cdp:eq:C.UI.4}
also give degree zero for both $K_Y$ and $\OO_Y(S)$
on every compact curve in this fibre. These are precisely
the ambient numerical inputs used below.


\end{enumerate}

\item \par\medskip\noindent\textbf{The proper analytic Lefschetz formula.}\quad
Take an embedded resolution
$r:\widetilde U\to U$ of the pair, isomorphic over
$U\setminus\operatorname{Sing}S$, with a simple
normal-crossings central divisor
$\sum_i a_iE_i$. Its strict transform $E_0$ has
multiplicity $a_0=m$. Every exceptional multiplicity
satisfies
\[
                    a_i\ge2m\quad(i\ne0).
 \tag{C.UI.5}\hypertarget{tag-014X}{}\label{cdp:eq:C.UI.5}
\]
Indeed a first centre outside the existing exceptional
locus lies in the singular locus of the reduced strict
transform, where its order is at least two. A centre
inside the existing exceptional locus already receives
the coefficient of an exceptional component, at least
$2m$. The coefficient of a new exceptional divisor is
the order of the current pullback of $f$ along its centre,
so the other contributions are nonnegative. Induction
proves \eqref{cdp:eq:C.UI.5}. The resolution can be chosen with all
centres above $\operatorname{Sing}S$. Its strict open
stratum is consequently
$E_0^\circ\simeq S_{\rm reg}$.


\begin{enumerate}[label={$\langle2\rangle$\arabic*.},leftmargin=1.3em,labelsep=.5em]
\item  We specify the proper version of the trace formula used
here. A'Campo's original Theorem~1 is stated for a local
hypersurface germ. His \S2 begins with a proper analytic
map from a smooth analytic space to a disc, smooth off
the central simple-normal-crossings divisor, with no
Kahler assumption. The global boundary fibration and
its cohomology spectral sequence are constructed there;
the stalk-trace computation is on pp.~243--244
\cite[\S2 and \S3]{ACampo1975}. In the present notation
that construction gives
\[
 L(T^j):=\sum_q(-1)^q\operatorname{Tr}
       (T^j|H^q(F,\mathbb Q))
       =\sum_{a_i\mid j}a_i e(E_i^\circ).
 \tag{C.UI.6}\hypertarget{tag-014Y}{}\label{cdp:eq:C.UI.6}
\]
For clarity, the local and global maps producing this
identity are as follows. The real oriented blowup of
the divisor, followed by the angular map to $S^1$,
has a boundary fibre homeomorphic to $F$, as in
\CompendiumRef{II}{cdp:st:C.2.NC.1}. Above an $r$-branch stratum its fibre is
\[
 \{(u_1,\ldots,u_r)\in(S^1)^r:
                 u_1^{a_1}\cdots u_r^{a_r}=e^{i\theta}\}.
 \tag{C.UI.7}\hypertarget{tag-014Z}{}\label{cdp:eq:C.UI.7}
\]
An integral change of angular basis identifies this with
$\gcd(a_1,\ldots,a_r)$ copies of $(S^1)^{r-1}$.
Monodromy permutes these copies and acts by translations
on them. For $r\ge2$ the alternating trace is zero,
whether the permutation fixes a copy or not. For $r=1$
it is $a_1$ when $a_1\mid j$ and zero otherwise.
The constructible higher-direct-image sheaves of this
boundary projection carry those stalk actions. Their
cohomology spectral sequence computes $H^*(F)$ with
its actual monodromy. Triangulate the compact central
divisor compatibly with the strata and these local
systems. On every open cell the alternating trace of
compactly supported cochains is its signed cell
contribution times the stalk trace. Additivity and
the spectral sequence prove \eqref{cdp:eq:C.UI.6}. The proper
modification is an isomorphism over $\Delta^*$, so
these are the nearby fibre and monodromy of the
original map.


\item  This description also proves quasi-unipotence without
an extra hypothesis. If $N$ is a common multiple of
the finitely many $a_i$, then $T^N$ is the identity
on every stalk cohomology group in \eqref{cdp:eq:C.UI.7}, since
the finite permutation disappears and translations
act trivially on cohomology. It is the identity on
the corresponding constructible cohomology sheaves,
hence on the second page of their cohomology spectral
sequence. The finite filtration on $H^q(F)$ has
identity action on its associated graded. Thus
$T^N$ is unipotent on $H^q(F)$.


\item  Put $j=m$ in \eqref{cdp:eq:C.UI.6}. By \eqref{cdp:eq:C.UI.5} only $E_0$
contributes. The cohomology ring of the real
four-torus is the exterior algebra of its first
cohomology, and a monodromy diffeomorphism acts by
the exterior powers of $T_1$. Therefore
\[
 m e(S_{\rm reg})
   =\sum_{q=0}^4(-1)^q\operatorname{Tr}(\Lambda^qT_1^m)
   =\det(1-T_1^m).
 \tag{C.UI.8}\hypertarget{tag-0150}{}\label{cdp:eq:C.UI.8}
\]
Every eigenvalue of $T_1^m$ is a root of unity by
the preceding quasi-unipotence. A real eigenvalue
is $1$ or $-1$ and contributes $0$ or $2$; each
nonreal conjugate pair contributes $|1-\lambda|^2$.
Their product is nonnegative. This proves \eqref{cdp:eq:C.UI.1}.


\end{enumerate}

\item \par\medskip\noindent\textbf{The local degree-one class and the finite-character comparison.}\quad
We record the exact local replacement for the global
transgression used in \ref{cdp:st:C.2.MH.2}. Every nonzero class
$\gamma_S\in H^1(S,\mathbb Z)$ extends to a sufficiently
small tube and restricts nontrivially over $\mathbb Q$
to $F$. To see the extension, triangulate the analytic
pair $(Y,S)$ as in \SeriesRef{cdp:st:A.5.11} and take a regular
neighbourhood $N$ with a retraction $N\to S$.
Properness puts $f^{-1}(\Delta)$ inside $N$ after
shrinking $\Delta$. Pullback under that retraction
gives the specified class $\gamma_U$ on the tube,
whose restriction to $S$ is $\gamma_S$.


\begin{enumerate}[label={$\langle2\rangle$\arabic*.},leftmargin=1.3em,labelsep=.5em]
\item  For injectivity upon restriction to $F$, the inclusion
$U^*=U\setminus S\hookrightarrow U$ surjects on
fundamental groups: loops can be perturbed off the
real codimension-two divisor, avoiding its higher
codimension singular strata. If a rational
degree-one class on $U$ restricts to zero on $F$,
the exact homotopy sequence of the punctured-disc
bundle makes its restriction to $U^*$ a scalar
multiple $a$ of the base-circle class. A small
meridian at a smooth point of $S$ is null-homologous
in $U$ and maps to $m$ times that base circle.
It follows that $ma=0$, hence $a=0$. Surjectivity
on fundamental groups now makes the original
class zero. In particular $\gamma_U|_F\ne0$.


\item  Choose a prime $\ell$ for which both restrictions
to $S$ and $F$ stay nonzero modulo $\ell$.
Their values on one loop in each space exclude
only finitely many primes. The reduction of
$\gamma_U$ defines the actual cyclic etale cover
of $U$, with rational augmentation local system
$V$ of rank $d=\ell-1$. It is self-dual,
$H^0(S,V)=0$, and every nontrivial complex character
in it is nontrivial on $F$. The Koszul complex
of $\mathbb Z^4$ is then contracted by the inverse
of a nonzero character difference. Hence
\[
 H^*(F,V)=H^*(U^*,V)=0.
 \tag{C.UI.9}\hypertarget{tag-0151}{}\label{cdp:eq:C.UI.9}
\]
The second equality is the bundle Leray comparison.


\item  When $Z$ is rational or ruled over a curve of genus
$g$, all strata of an embedded resolution of the
central fibre are projective: the strict transform
is a blowup of $Z$, and every other surface arises
by blowups from a projective bundle over a compact
curve or from $\mathbb P^2$. \ref{cdp:st:C.2.MH.1} thus applies
to this actual local system. Its support comparison
makes $H^3(S,V)$ pure of weight three. The exact
normalization-conductor descent sequence has a
kernel which is a quotient of the degree-two
cohomology of a compact curve, hence pure of weight
two; strictness makes that kernel zero. The
negative-definite exceptional intersection matrices
then give the actual isomorphisms
\[
 H^3(S,V)\simeq H^3(Z,V),\qquad
 H^4(S,V)\simeq H^4(Z,V).
 \tag{C.UI.10}\hypertarget{tag-014N}{}\label{cdp:eq:C.UI.10}
\]
These are precisely the maps (C.MH.10)--(C.MH.12);
none of their proofs uses ambient cohomology.
The acyclic punctured tube and local support duality
give, in the direct limit of shrinking tubes,
$H^j(S,V)\simeq H_{6-j}(S,V)$. Thus
\[
 d e(S)=2h^4(Z,V)-h^3(Z,V)
 =\begin{cases}
 2d,&Z\text{ rational},\\
 (2-2g)d,&Z\text{ ruled over genus }g.
 \end{cases}
 \tag{C.UI.11}\hypertarget{tag-014O}{}\label{cdp:eq:C.UI.11}
\]
For the second line, the local system pulls back
from the ruling base. If its character is trivial
there, the two dimensions on the left are $d$
and $2gd$; otherwise they are $0$ and $(2g-2)d$.
This proves the precise local character version
of \ref{cdp:st:C.2.MH.2} and the additional assertion in item~1.


\end{enumerate}

\item \par\medskip\noindent\textbf{Strict positivity for a rational normalization.}\quad
If $S$ is normal, \ref{cdp:st:C.2.UN.5} already makes it a
smooth torus or bielliptic surface, so its
resolution is neither rational nor irrational
ruled. Thus the remaining surface cases under
consideration are nonnormal.
The proof of \SeriesRef{cdp:st:A.5.16} uses the ambient assumption
only for $\omega_S$ torsion and $\chi(\OO_S)=0$,
now supplied by \eqref{cdp:eq:C.UI.2}--\eqref{cdp:eq:C.UI.3}. Consequently
its exact support and genus conclusions apply:
\[
 \begin{gathered}
 K_Z+D\simeq\OO_Z,\qquad D_{\rm red}\text{ connected},
 \qquad p_a(D_{\rm red})\le1,\\
 \Gamma_{\rm red}\text{ connected},\qquad
 p_a(\Gamma_{\rm red})\le2,\qquad
 e(\Gamma_{\rm red})\ge0,\\
 D_{\rm red}=h^{-1}(\Gamma_{\rm red}),\qquad
 S_{\rm reg}\simeq Z\setminus D_{\rm red}.
 \end{gathered}
 \tag{C.UI.12}\hypertarget{tag-014P}{}\label{cdp:eq:C.UI.12}
\]
The final isomorphism retains all exceptional curves:
any exceptional fibre outside the conductor would
give an isolated normal hypersurface singularity,
excluded by the actual Milnor-sphere comparison
in \SeriesRef{cdp:st:A.5.17}. Thus
$e(S)=e(S_{\rm reg})+e(\Gamma_{\rm red})\ge0$.


\begin{enumerate}[label={$\langle2\rangle$\arabic*.},leftmargin=1.3em,labelsep=.5em]
\item  Suppose it were zero. Both summands would be zero.
If $D_{\rm red}$ has an elliptic-normalized component,
the bound on its arithmetic genus makes that component
smooth and unique. Follow a sequence of rational
surface blowdowns to $\mathbb P^2$ or a Hirzebruch
surface. The divisor comparison at each step is
$D_{\rm new}=\sigma^*D_{\rm old}-E$; a centre of
multiplicity one increases the Euler number of the
complement by one, and a centre of multiplicity at
least two preserves it. A centre outside the
boundary is impossible by effectiveness. The
elliptic component survives these blowdowns.
On $\mathbb P^2$ the resulting anticanonical cubic
is a smooth elliptic cubic. On a Hirzebruch
surface the elliptic component is a smooth
bisection: a section is rational, and for a
coefficient-one bisection $H$ with vertical part
$V$ adjunction gives
$2p_a(H)-2+H\cdot V=0$, forcing $V=0$.
The minimal complement Euler numbers are therefore
$3$ and $4$, respectively. Their blowup comparison
contradicts $e(S_{\rm reg})=0$.


\item  All normalized components of $D_{\rm red}$ are
therefore rational. Its first Betti number is at
most one, since the normalization sequence gives
$b_1(D_{\rm red})\le p_a(D_{\rm red})\le1$.
Every component of $\Gamma_{\rm red}$ is also
rational-normalized, being the finite image of a
strict conductor curve. Let their number be
$n\ge1$. For a connected compact curve with
rational normalizations the exact normalization
topology gives
$e(\Gamma_{\rm red})=1+n-b_1(\Gamma_{\rm red})$.
The supposed equality therefore gives
$b_1(\Gamma_{\rm red})=n+1\ge2$.


\item  Here the required comparison with $S$ is an
actual pushout, not a presumed identification
of cohomology groups. By \eqref{cdp:eq:C.UI.12}, the natural
map from
$Z\cup_{D_{\rm red}}\Gamma_{\rm red}$ to $S$
is a continuous bijection from a compact space
to a Hausdorff space, hence a homeomorphism.
Triangulation makes $D_{\rm red}\hookrightarrow Z$
a cofibration, so its double mapping cylinder
gives the cohomology sequence of that pushout.
Connectedness and $H^1(Z,\mathbb Q)=0$ give
\[
 H^1(S,\mathbb Q)
  \simeq\ker\bigl[H^1(\Gamma_{\rm red},\mathbb Q)
                       \to H^1(D_{\rm red},\mathbb Q)\bigr].
 \tag{C.UI.13}\hypertarget{tag-014Q}{}\label{cdp:eq:C.UI.13}
\]
Its dimension is at least $n>0$. The proved
local finite-character calculation \eqref{cdp:eq:C.UI.11}
now gives $e(S)=2$, contradicting zero.
Thus $e(S)>0$, and \ref{cdp:st:C.2.UN.1} makes it an even
integer. This proves item~1.


\end{enumerate}

\item \par\medskip\noindent\textbf{All noncontracted irrational ruled alternatives.}\quad
On a general ruling fibre $R$, \eqref{cdp:eq:C.UI.4} gives
$D\cdot R=2$. Thus its horizontal part is reduced
of total degree two, or is $2C$ for a smooth
section $C$. The proof of \ref{cdp:st:C.2.IRR.1} uses the
ambient condition only for the zero degrees
in (C.IR.2), (C.IR.7)--(C.IR.8) and for the
torsion-canonical, zero-$\chi$ genus bound \SeriesRef{cdp:st:A.5.14}.
All are furnished by \eqref{cdp:eq:C.UI.2}--\eqref{cdp:eq:C.UI.4}.
Its exact branch-reception, pinch-point and
blowdown comparisons therefore give item~2,
including its assertion that no horizontal
component is contracted.


\begin{enumerate}[label={$\langle2\rangle$\arabic*.},leftmargin=1.3em,labelsep=.5em]
\item  Consider $D=2C+V$ with $C$ not contracted.
Every resolution-exceptional curve is vertical.
Indeed a positive-coefficient horizontal one
would have to be $C$. A coefficient-zero
exceptional curve $E$ satisfies
$K_Z\cdot E=D\cdot E=0$ by minimality and
adjunction. A horizontal exceptional curve over
a positive-genus base has $p_a(E)\ge1$ and
$E^2<0$, so $K_Z\cdot E>0$, a contradiction.
These are exactly the exceptional-curve
hypotheses of \SeriesRef{cdp:st:A.5.18}--\SeriesRef{cdp:st:A.5.20}.


\item  Let $H$ be the reduced image of $C$ and
$d=\deg(C\to H^\nu)$. The branchwise conductor
formula is its intrinsic plane-branch conductor
plus the sum of its intersections with the
other branches. Its value on each horizontal
branch is two. If $d=1$, $H^\nu$ has the
positive genus of $C$, and no vertical rational
curve can dominate it. The sole branch with
conductor two is the ordinary cusp. If $d=2$,
the two horizontal branches must be smooth:
a singular branch would contribute at least two
intrinsically and at least one from the other
branch. Their allowed intersections give either
contact two between those branches, or simple
contact between them and with one additional
smooth branch. These are exactly a tacnode and
an ordinary triple point. For $d=3$ all three
branches must be smooth and pairwise transverse,
and there is no further branch. For $d\ge4$
the three other intersection contributions
already exceed two. This proves exhaustion
of the transverse cases used in the following
comparisons.


\item  The local numerical hypotheses of \SeriesRef{cdp:st:A.5.14} are
exactly $\chi(\OO_S)=0$ and torsion $\omega_S$.
Its genus bound and delta-one complete-intersection
comparison consequently remain valid. The
conductor-cleared jets \SeriesRef{cdp:st:A.5.5} use degree zero for
$K_Y$ and $\OO_Y(S)$ on the image curve.
The cusp proof \SeriesRef{cdp:st:A.5.20} and the complete tacnode
proof \SeriesRef{cdp:st:A.5.18} therefore apply. The latter's
rational-image fundamental-group restriction
uses only the local quotient of
$\mathbb Z^4\rtimes\mathbb Z$, not an ambient
degree-one class. For the triple case, \SeriesRef{cdp:st:A.5.19}
excludes singular rational images with a
positive-genus covering component. \cref{cdp:st:C.2.D4.1}, \cref{cdp:st:C.2.D4.2}
exclude the connected degree-three and degree-two
covers of a smooth rational image. Their normal
cone is the actual cubic tensor in
$\operatorname{Sym}^3N_{H/Y}^*$; on
$H\simeq\mathbb P^1$ the needed identities
$\OO_Y(S)|_H=\OO_H$ and $\deg N_{H/Y}=-2$
follow from \eqref{cdp:eq:C.UI.2}. Their compact-core
transport is into the original nearby torus,
so no global Picard condition enters. The
independent ruled Chern comparison \SeriesRef{cdp:st:A.5.10} also
extends, since its additional input
$c_2(h^*TY)=0$ is \eqref{cdp:eq:C.UI.3}.


\item  Consequently the image in each surviving case
is smooth elliptic and the transverse analytic
type is constant. No vertical rational curve
dominates that image, so all its generic
branches belong to $C$; the actual local
normalizations give etale degrees $1$, $2$,
or $3$, respectively. The normal line of $C$
has degree zero. The divisor comparison
(A.T.22), together with this square-zero
section and the absence of further conductor
along it, contracts no point of $C$ and makes
the minimal ruled boundary exactly $2C$.
A first reverse blowup away from $C$ would
create coefficient $-1$; a blowup on $C$
would create a conductor component meeting
it, contrary to the constant local type.
Thus the normalization is smooth minimal
elliptic ruled, exactly as proved in
\SeriesRef{cdp:st:A.5.18}--\SeriesRef{cdp:st:A.5.20}. The normalization Euler
comparison $e(S)=e(T)-e(C)+e(H)$ gives zero.
This proves item~3.


\end{enumerate}

\item \par\medskip\noindent\textbf{A contracted section and its fundamental group.}\quad
Finally let $D=2C+V$ with $C$ contracted and
let $g\ge1$ be the ruling genus. The exact
boundary blowdown relation \eqref{cdp:eq:C.CR.5}, now
supplied by \eqref{cdp:eq:C.UI.4}, shows that $D_{\rm red}$
is the smooth section $C$ with trees of smooth
rational curves attached. All its crossings
are transverse nodes. Its inclusion in $Z$
induces an isomorphism on fundamental groups,
both identified with $\pi_1(C)$.
The proof of \eqref{cdp:eq:C.CR.6} uses only this
connectedness, minimality and the isolated
singularity exclusion. It gives
\[
 D_{\rm red}=h^{-1}(\Gamma_{\rm red}),\qquad
 S_{\rm reg}\simeq Z\setminus D_{\rm red}.
 \tag{C.UI.14}\hypertarget{tag-014R}{}\label{cdp:eq:C.UI.14}
\]
Every component of the connected curve
$\Gamma_{\rm red}$ is dominated by a vertical
rational curve and hence has normalization
$\mathbb P^1$. The same actual topological
pushout as above and van Kampen give
\[
 \pi_1(S)\simeq
 \pi_1(Z)*_{\pi_1(D_{\rm red})}\pi_1(\Gamma_{\rm red})
 \simeq\pi_1(\Gamma_{\rm red}).
 \tag{C.UI.15}\hypertarget{tag-014S}{}\label{cdp:eq:C.UI.15}
\]
The second amalgamating map is trivial because
$C$ generates $\pi_1(D_{\rm red})$ and is
contracted. This group is free of rank
$\beta=b_1(\Gamma_{\rm red})$.


\begin{enumerate}[label={$\langle2\rangle$\arabic*.},leftmargin=1.3em,labelsep=.5em]
\item  For its comparison with the nearby torus,
take the same regular neighbourhood and
small tube used above. A loop in $S$ lies
in the tube, and the neighbourhood retraction
is the identity on $S$. Thus the tube group
surjects onto $\pi_1(S)$. Perturbing loops
off $S$ makes it a quotient of
$\pi_1(U^*)=\mathbb Z^4\rtimes\mathbb Z$.
The latter is metabelian, whereas a free
group of rank at least two surjects onto
the nonmetabelian group $S_4$. Hence
$\beta\le1$. If $n\ge1$ is the number
of conductor-image components, its exact
normalization topology gives
\[
 e(\Gamma_{\rm red})=n+1-\beta\ge n\ge1,
 \qquad e(S)=e(S_{\rm reg})+e(\Gamma_{\rm red})>0.
 \tag{C.UI.16}\hypertarget{tag-014T}{}\label{cdp:eq:C.UI.16}
\]
Here \eqref{cdp:eq:C.UI.1} proves the last inequality.
Parity makes $e(S)\ge2$. If $\beta=1$,
\eqref{cdp:eq:C.UI.15} supplies a nonzero integral
degree-one class on $S$; \eqref{cdp:eq:C.UI.11} then
gives $e(S)=2-2g\le0$, a contradiction.
Thus $\beta=0$, proving $\pi_1(S)=1$.
The transgression of \SeriesRef{cdp:st:D.1.9} excludes this
case under the full CDP hypotheses. This
proves item~4 and completes the proof.



\end{enumerate}

\end{enumerate}
\end{proof}



The assertion $e(S)=2$ in item~1 is deliberately tied
to the proved nonzero local-character comparison.
\ref{cdp:st:C.2.UN.8} constructs an irreducible rational-normalized
support with $\pi_1(S)=1$ and $e(S)=6$, so that
local character cannot be inserted without its
existence argument. The classification of all
irreducible zero-Euler supports is assembled in
\ref{cdp:st:C.2.UN.12} using this result, \ref{cdp:st:C.2.UN.5}, and the
class VII conductor comparisons. Reducible
unrestricted supports are not covered by this
irreducible statement.




\begin{proposition}[The unrestricted irreducible zero-Euler criterion]\hypertarget{tag-01CP}{}\label{cdp:st:C.2.UN.12}
\StatementRecord{cdp:st:C.2.UN.12}{proved}{01CP}{\SeriesRef{cdp:st:C.2.UN.5}, \SeriesRef{cdp:st:C.2.UN.7}, \SeriesRef{cdp:st:C.2.UN.19}, \SeriesRef{cdp:st:A.5.20}, \SeriesRef{cdp:st:A.5.18}, \SeriesRef{cdp:st:C.2.D4.1}, \SeriesRef{cdp:st:C.2.D4.2}, \SeriesRef{cdp:st:C.2.IRR.1}, \SeriesRef{cdp:st:C.2.VII.1}, \SeriesRef{cdp:st:C.2.VII.2}, \SeriesRef{cdp:st:A.5.17}, \SeriesRef{cdp:eq:C.IZ.6}, \SeriesRef{cdp:eq:C.IZ.5}, \SeriesRef{cdp:eq:C.IZ.8}, \SeriesRef{cdp:eq:C.IZ.7}}

Let $f:Y\to B$ be a proper holomorphic map from a smooth
complex threefold to a smooth complex curve, with connected
fibres and smooth general fibre a complex two-torus. Fix an
irreducible full fibre $mS$, and write
$h=\eta\pi:Z\to T\to S$ for normalization followed by its
minimal resolution. Every such support satisfies $e(S)\ge0$.
Its Euler number is zero precisely in the following geometric
cases, stated for these actual maps.

A normal support is a smooth torus or a smooth bielliptic
surface, with multiplicities as in \ref{cdp:st:C.2.UN.5}.

A nonnormal support with irrational ruled normalization has
smooth minimal normalization ruled over an elliptic curve.
Its conductor image and reduced inverse image are disjoint
unions of smooth elliptic curves. The transversal types and
their finite conductor maps are exactly those in the table:
\[
\begin{array}{c|c|c}
 \text{normalization conductor}&\text{transversal type}
       &\text{degree over the image}\\ \hline
 H\text{ reduced, ruling degree }2&xy=0&2\\
 2C,\ C\text{ a section}&x^2-y^3=0&1\\
 2C,\ C\text{ a section}&x^2-y^4=0&2\\
 2C,\ C\text{ a section}&xy(x-y)=0&3
\end{array}                                                   \tag{C.IZ.1}\hypertarget{tag-00ON}{}\label{cdp:eq:C.IZ.1}
\]
Every displayed map of elliptic curves is unramified; the
degree-one map is an isomorphism. The transversal type is
constant along each conductor component. In the first row,
degree two means the total degree over each image component;
its inverse image may have two connected components. There are no
contracted horizontal sections in the zero-Euler case.

For a nonnormal support whose normalization resolution
has a Hopf surface as its minimal model, put $R=D_{\rm red}$,
where $D$ is the effective
resolution conductor with $K_Z+D=h^*\omega_S$. Let $r$ be
the number of rational components of $R$. Put
$C=\Gamma_{\rm red}$ for the conductor image, let $n_0$ be
the number of its rational-normalized components, and set
$v(C)=\sum_x(s_x-1)$, using the actual number $s_x$ of
normalization branches at each singular point. The exact
zero criterion is
\[
 b_2(Z)=r,\qquad
 \begin{cases}
 n_0=0,\quad v(C)=0,\quad\sum_x\delta_x\le1,\quad\text{or}\\
 n_0=1,\quad v(C)=2,\quad\sum_x\delta_x=2.
 \end{cases}                                                   \tag{C.IZ.2}\hypertarget{tag-00OO}{}\label{cdp:eq:C.IZ.2}
\]
Every other normalized conductor component is elliptic.
In the second line every singular point satisfies
$\delta_x=s_x-1$. No normalization singularity lies outside
the conductor in either line. An irreducible support with
rational normalization resolution always has $e(S)\ge2$.

This is an equivalence for fibres of the specified map.
The displayed curve data retain their normalization and
resolution morphisms; no assertion of smoothability of
arbitrary abstract surface gluings is included.
\end{proposition}


\begin{proof}

\begin{enumerate}[label={$\langle1\rangle$\arabic*.},leftmargin=*]
\item We first check the classification's scope without an ambient
Chern-class assumption. On a sufficiently small fibre tube
$U=f^{-1}(\Delta)$, the evaluation of a generator of the
rank-one torsion-free sheaf $f_*K_U$ gives
$K_U\simeq\OO_U(kS)$, as proved with its divisor comparison
in \ref{cdp:st:C.2.UN.7}. Since $mS$ is principal, $K_U$ and
$\omega_S=\OO_U((k+1)S)|_S$ are torsion. The actual
$\OO_{mS}$ filtration and flatness give
$\chi(\OO_S)=0$. For nonnormal $S$, finite duality and
resolution adjunction give the nonempty effective divisor
$D$ and the stated torsion line $K_Z+D$.

\item Every positive pluricanonical section on $Z$ would, after
taking a power trivializing that torsion line, give a
nonzero section of $\OO_Z(-aD)$ for some $a>0$.
Its inclusion in $\OO_Z$ would be a nonzero holomorphic
function on the compact connected surface vanishing along
$D$, a contradiction. Thus $\kappa(Z)=-\infty$.
The surface classification, together with the actual
minimal-model blowdown, makes $Z$ rational, irrational
ruled, or class VII. In the last case \ref{cdp:st:C.2.UN.19} proves,
through the specified cyclic cover and quotient maps,
that its minimal model is Hopf. Thus the formulation
above includes every class VII possibility for the
actual degeneration. Normal supports have already been
classified with their finite root covers in \ref{cdp:st:C.2.UN.5}.

\ref{cdp:st:C.2.UN.7} proves strict positivity for rational $Z$.
It also proves strict positivity when an irrational
horizontal section is contracted, and extends every
noncontracted irrational ruled comparison to this tube.
The precise replacements for the earlier ambient $b_2=0$
uses are the torsion canonical and normal lines, the
vanishing of $\chi(\OO_S)$, and
$\int_Zh^*c_2(T_U)=0$, all proved there through restriction
to the nearby torus.

\item For clarity, the horizontal possibilities exhaust the
ruling degree two of $D$. A reduced horizontal part is
covered by C.2.IRR.1: its adjunction comparison makes its
curves smooth elliptic and disjoint, its normalization
smooth minimal ruled, and its conductor map an unramified
double cover with product-node germs. A nonreduced
horizontal part is $2C$ for a section $C$. At a generic
conductor point, the plane-curve conductor exponent of a
normalization branch equals its intrinsic conductor plus
its intersections with the other branches. Exponent two
permits one cusp branch, two smooth branches meeting with
multiplicity two, or three pairwise transverse smooth
branches. These are the last three rows of \eqref{cdp:eq:C.IZ.1}.
An additional branch would give exponent at least three.
The exact cusp, tacnode and cubic-cover comparisons in
\SeriesRef{cdp:st:A.5.20}, \SeriesRef{cdp:st:A.5.18} and \cref{cdp:st:C.2.D4.1}, \cref{cdp:st:C.2.D4.2}, extended in \ref{cdp:st:C.2.UN.7},
give the asserted elliptic image, constant type and
degrees one, two and three. A vertical rational curve
cannot dominate that elliptic image, so every generic
branch is accounted for by the specified horizontal map.
All these comparisons remove vertical conductor and
normalization exceptional curves. Their Euler formula is
\[
 e(S)=e(T)-e((\eta^{-1}C)_{\rm red})+e(C)=0,             \tag{C.IZ.3}\hypertarget{tag-00OP}{}\label{cdp:eq:C.IZ.3}
\]
since each term vanishes. This proves both directions of
the ruled part of the criterion.

\item It remains to check and extend the class VII part. The
proof of \ref{cdp:st:C.2.VII.1} uses only the actual torsion relation
$K_Z+D$; it gives
\[
 e(R)=r,\qquad r\le b_2(Z).                              \tag{C.IZ.4}\hypertarget{tag-00OQ}{}\label{cdp:eq:C.IZ.4}
\]
The proof of \ref{cdp:st:C.2.VII.2} uses only this relation,
$\omega_S$ torsion and $\chi(\OO_S)=0$. They have just
been established on $U$, so its finite-duality and
curve-normalization maps give
\[
 \chi(\OO_C)\ge-1,\qquad e(C)\ge0.                      \tag{C.IZ.5}\hypertarget{tag-00OR}{}\label{cdp:eq:C.IZ.5}
\]
Its normalizations are rational or elliptic, because they
are finite images of the rational or elliptic curves in
$R$. These statements have no remaining ambient Betti
assumption.

\item We sharpen the exceptional terms in the earlier comparison.
At a point of $T$ outside the conductor, $\eta$ is an
isomorphism. A singular point there would consequently be
an isolated normal hypersurface singularity of $S$ in
the smooth tube. The isolated Milnor-fibre argument
\SeriesRef{cdp:st:A.5.17}, retaining the local unit and multiplicity of
the full fibre, excludes that point. Thus $T$ is smooth
outside the conductor. Every resolution-exceptional curve
of coefficient zero in $D$ is disjoint from $R$: minimal
resolution adjunction gives $K_Z\cdot E\ge0$, while
$D\cdot E\ge0$ and $(K_Z+D)\cdot E=0$ force both
numbers to vanish. Each exceptional fibre over the
conductor meets its strict transform and is connected.
Hence it cannot contain such a disjoint curve, and its
whole inverse image lies in $R$. There are no exceptional
fibres outside the conductor, by the preceding local
argument. We have therefore identified the actual
complement isomorphism
\[
 Z\setminus R\xrightarrow{\ \simeq\ }
 T\setminus(\eta^{-1}C)_{\rm red}
 \xrightarrow{\ \simeq\ }S\setminus C.                \tag{C.IZ.6}\hypertarget{tag-00OS}{}\label{cdp:eq:C.IZ.6}
\]
Its first map is the resolution restricted to this open
set, and its second is the finite normalization restricted
where its conductor ideal is the unit ideal.

\item A class VII surface and its point blowups have
$b_1=b_3=1$, so $e(Z)=b_2(Z)$. Euler additivity in
\eqref{cdp:eq:C.IZ.6} now yields the exact formula
\[
 e(S)=b_2(Z)-r+e(C),\qquad
 b_2(Z)-r\ge0,\quad e(C)\ge0.                           \tag{C.IZ.7}\hypertarget{tag-00OT}{}\label{cdp:eq:C.IZ.7}
\]
The curve normalization sequence identifies the remaining
term, rather than replacing its singular points by a
model:
\[
 e(C)=2n_0-v(C),\qquad
 \chi(\OO_C)=n_0-\sum_x\delta_x,\qquad
 \delta_x\ge s_x-1.                                    \tag{C.IZ.8}\hypertarget{tag-00OU}{}\label{cdp:eq:C.IZ.8}
\]
Equations \eqref{cdp:eq:C.IZ.5} and \eqref{cdp:eq:C.IZ.8} imply $e(C)\ge n_0-1$.
Consequently $e(C)=0$ allows exactly $n_0=0,v=0$ or
$n_0=1,v=2$. In the first case the coherent bound gives
$\sum\delta_x\le1$. In the second, that bound and
$\delta_x\ge s_x-1$ give $\sum\delta_x=2$, with equality
at every singular point. Conversely either of these
two curve patterns makes $e(C)=0$ in the exact
topological formula. Equation \eqref{cdp:eq:C.IZ.7} proves precisely
\eqref{cdp:eq:C.IZ.2}, and proves the sign for every class VII
support. All three nonnormal surface classes have now
been included, which proves the assertion.

\end{enumerate}
\end{proof}



The class VII surface input is the proved anticanonical
configuration theorem checked in \ref{cdp:st:C.2.VII.1}, with its exact
references to \cite{Dloussky2006,DOT2003}. \ref{cdp:st:C.2.UN.19} further
identifies the minimal model in this branch as Hopf;
the existing conductor comparison applies through its
actual point-blowup map. The result here
transfers that input to arbitrary irreducible torus fibres
through the canonical evaluation, isolated-singularity
exclusion and the complement map \eqref{cdp:eq:C.IZ.6}. The unreduced
fibre multiplicity is retained throughout. The remaining
unrestricted classification problem in C.2 concerns
reducible supports.




\begin{proposition}[Proved elliptic-image restriction]\hypertarget{tag-01C0}{}\label{cdp:st:C.2.INT.1}
\StatementRecord{cdp:st:C.2.INT.1}{proved}{01C0}{\SeriesRef{cdp:eq:C.INT.3}, \SeriesRef{cdp:st:D.1.7}, \SeriesRef{cdp:eq:C.INT.1}, \SeriesRef{cdp:eq:C.INT.2}; \cite[Lemma~2.2 and \S2]{BrinzanescuMoraru2005}; \cite[Remark~2.7]{DlousskyTeleman2012}}


Let $Y$ be a smooth compact complex threefold with
$H^2(Y,\Q)=0$, let $S_i$ be an irreducible component of
a connected reducible divisor $S$, and let
$h:Z\to S_i\hookrightarrow Y$ be the minimal resolution
of its normalization. Suppose that the minimal model
$\sigma:Z\to Z_0$ is a Hopf surface or a non-K\"ahler
elliptic surface. Every strict transform $E\subset Z$ of
an elliptic curve on $Z_0$ has a one-dimensional image
under $h$, and the normalization of that image is elliptic.
The actual restriction homomorphism
$\Pic^0(Y)\to\Pic^0(E)$ is nonconstant for each such $E$.
In particular such a component cannot meet a component
whose minimal model is an Enoki surface.
\end{proposition}


\begin{proof}

\begin{enumerate}[label={$\langle1\rangle$\arabic*.},leftmargin=*]
\item Choose an adjacent component $S_j$. The actual pullback
of its canonical divisor section gives a nonzero effective
Cartier divisor
\[
 A=h^*S_j>0,\qquad \OO_Z(A)=h^*\OO_Y(S_j).
 \tag{C.INT.1}\hypertarget{tag-00NQ}{}\label{cdp:eq:C.INT.1}
\]
Its first Chern class is torsion. For some $q>0$,
$\OO_Y(qS_j)\in\Pic^0(Y)$, so the restriction homomorphism
$\rho:\Pic^0(Y)\to\Pic^0(Z)$ has in its image the line
$\OO_Z(qA)$. That line has strictly positive degree for
every Gauduchon metric on $Z$.

\item We spell out the blowdown comparison. Put $A_0=\sigma_*A$.
The difference $A-\sigma^*A_0$ is exceptional and has
zero real Chern class: both $A$ and its pushforward have
zero real class. The negative definite exceptional
intersection matrix makes every coefficient of this
difference zero. Thus
\[
 A=\sigma^*A_0,\qquad A_0>0.
 \tag{C.INT.2}\hypertarget{tag-00O1}{}\label{cdp:eq:C.INT.2}
\]
Nonzeroness follows by the same negative definiteness if
$A_0$ were zero. Pullback is an isomorphism
$\Pic^0(Z_0)\simeq\Pic^0(Z)$; the integral blowup
formula and $R\sigma_*\OO_Z=\OO_{Z_0}$ identify both
the lattice and the coefficient space defining it.
We therefore regard $\rho$ as a map into $\Pic^0(Z_0)$.
Its image contains $\OO_{Z_0}(qA_0)$.

\item We first prove the stronger assertion that the composite
restriction to $E$ cannot be zero. Suppose that it is zero.
This situation would in particular arise if $h(E)$ were
a point or had rational normalization.
In the latter case the map from the normal curve $E$
factors as $E\to\PP^1\to Y$. In both cases every line
from $\Pic^0(Y)$ restricts trivially to $E$. Here the
pullback to $\PP^1$ is in $\Pic^0(\PP^1)=0$; this is
the specified map giving the vanishing. The restriction
of $\sigma$ to $E$ is an isomorphism onto its smooth
elliptic image $E_0$.

\item For a Hopf $Z_0$ the homomorphism
\[
 H^1(Z_0,\R)\longrightarrow H^1(E_0,\R)
 \tag{C.INT.3}\hypertarget{tag-00OA}{}\label{cdp:eq:C.INT.3}
\]
is nonzero. On a primary Hopf surface this follows from
its universal cover $\C^2\setminus\{0\}$: a finite image
of $\pi_1(E_0)$ in $\pi_1(Z_0)=\Z$ would lift a finite
unramified cover of $E_0$ to that cover. Its two coordinate
functions would be constant, contradicting the nonconstant
inclusion of the curve. For a secondary Hopf surface,
pull back along a finite unramified primary cover. A
component above $E_0$ is an elliptic curve. Transfer makes
$H^1(Z_0,\R)\to H^1(Z'_0,\R)$ an isomorphism between
one-dimensional spaces, and the preceding primary
argument proves \eqref{cdp:eq:C.INT.3} after this pullback, hence
before it. The real-to-holomorphic coefficient injection
of \SeriesRef{cdp:st:D.1.7}, applied on $E_0$, shows that
$H^1(Z_0,\OO)\to H^1(E_0,\OO)$ is nonzero. Since
$\Pic^0(Z_0)=\C^*$, the identity component of the
kernel of $\Pic^0(Z_0)\to\Pic^0(E_0)$ is zero.
Connectedness of $\Pic^0(Y)$ would therefore make
$\rho$ zero, contradicting \eqref{cdp:eq:C.INT.1}--\eqref{cdp:eq:C.INT.2}.

\item For a non-K\"ahler elliptic $p:Z_0\to B$, every curve
on its minimal model is a reduced elliptic fibre.
The actual covering construction of
\cite[Lemma~2.2 and \S2]{BrinzanescuMoraru2005}
provides a finite cyclic map $\psi:Z'_0\to Z_0$ and a
principal elliptic bundle $p':Z'_0\to B'$ such that a
reduced fibre $E'$ above $E_0$ maps to $E_0$ by a finite
isogeny. To include multiple fibres, take a base cover
whose branching orders are divisible by their
multiplicities and normalize the fibre product. In the
local cyclic presentation of a smooth multiple fibre,
this removes the transverse multiplicity and leaves a
translation quotient on the reduced fibre, giving exactly
the stated isogeny. Extra branch points over ordinary
fibres make the total map branched there, but its map on
the reduced elliptic fibre is an isomorphism.

\item On this principal bundle the source's explicit
presentation is
\[
 Z'_0=\Theta^*/\langle\tau\rangle,
 \quad |\tau|>1,\quad \deg\Theta>0,
 \qquad
 \Pic^0(Z'_0)=p'^*\Pic^0(B')\times\C^*.
 \tag{C.INT.4}\hypertarget{tag-00OB}{}\label{cdp:eq:C.INT.4}
\]
In the last factor a scalar $\alpha$ is the automorphy
factor of a line bundle. Its restriction to
$E'=\C^*/\langle\tau\rangle$ is that same automorphy
line; its kernel is the discrete subgroup
$\langle\tau\rangle$. Consequently the identity
component of the restriction kernel is precisely the
compact group $p'^*\Pic^0(B')$.
Our assumed trivial restriction on $E$ makes
$\psi^*\rho(\Pic^0(Y))$ lie in that identity component.
Gauduchon degree is a continuous additive homomorphism
to $\R$, and is zero on any compact subgroup. On the
other hand this image contains
$\OO_{Z'_0}(q\psi^*A_0)$, the line of a nonzero effective
divisor, which has strictly positive Gauduchon degree.
This contradiction proves the first assertion in both
cases, and proves the stated nonconstant restriction.
Riemann--Hurwitz gives genus at most one for the
normalization of the nonconstant image of $E$, so it is
elliptic.

\item For the last assertion, an effective divisor of zero
real Chern class on a minimal class VII surface is,
by \cite[Remark~2.7]{DlousskyTeleman2012}, an elliptic
divisor on a Hopf surface or a positive multiple of the
rational cycle on an Enoki surface. If the source states
integral homological triviality, first multiply the
divisor by an integer killing its torsion integral
class and use Poincar\'e duality; the resulting support
and coefficient conclusion descends to the original
divisor. Formula \eqref{cdp:eq:C.INT.2} also applies to an Enoki
minimal model. Thus the adjacent divisor on its
resolution is the full pullback of a rational cycle;
all its curve images have rational normalization.
On a Hopf or non-K\"ahler elliptic neighbouring
resolution, the adjacent divisor is the full pullback
of a nonzero divisor consisting of elliptic curves.
At least one of their strict transforms must map
into the same intersection. It cannot map to a point
or a rational-normalized curve by the result just
proved. This is the required contradiction.

\end{enumerate}
\end{proof}






\begin{proposition}[Exclusion of Enoki components]\hypertarget{tag-01C1}{}\label{cdp:st:C.2.INT.2}
\StatementRecord{cdp:st:C.2.INT.2}{proved}{01C1}{\SeriesRef{cdp:st:C.3.4}, \SeriesRef{cdp:st:C.2.INT.1}, \SeriesRef{cdp:st:A.5.3}, \SeriesRef{cdp:st:A.4.4}, \SeriesRef{cdp:eq:C.INT.5}, \SeriesRef{cdp:eq:C.INT.7}, \SeriesRef{cdp:eq:C.INT.6}; \cite[Lemma~2.22 and the following paragraph]{Dloussky2006}}


Let $f:Y\to B$ be a proper connected-fibre map from a
smooth compact complex threefold with $H^2(Y,\Q)=0$
and general fibre a complex two-torus. No component
of a reducible fibre has normalization resolution
with Enoki minimal model.
\end{proposition}


\begin{proof}

\begin{enumerate}[label=\arabic*.,leftmargin=*]
\item Write the full fibre as $F=\sum_i m_iS_i$. Its
normalization resolutions $Z_i$ are non-K\"ahler by
\SeriesRef{cdp:st:C.3.4}. A neighbouring component gives on each $Z_i$
a nonzero effective divisor $E_{ij}=h_i^*S_j$ of
zero real Chern class. On a class VII minimal model
the precise classification in \ref{cdp:st:C.2.INT.1} makes that
model Hopf or Enoki. The other non-K\"ahler minimal
models are elliptic surfaces, including both types
of Kodaira surfaces. \ref{cdp:st:C.2.INT.1} rules out an interface
between either such type and an Enoki component.
Thus connectedness of the fibre shows that, if an
Enoki component occurs, every component has Enoki
minimal model.

\item Let $D_i^{\rm self}\ge0$ be the effective resolution
conductor for $S_i$ from \SeriesRef{cdp:st:A.5.3}, and set
$D_i=D_i^{\rm self}+\sum_{j\ne i}E_{ij}$. This is
also the resolution conductor for the total reduced
support. Adjunction specifies the exact comparisons
\[
 K_{Z_i}+D_i=h_i^*(K_Y+\sum_j S_j),\qquad
 K_{Z_i}+D_i^{\rm self}=h_i^*(K_Y+S_i).
 \tag{C.INT.5}\hypertarget{tag-00OC}{}\label{cdp:eq:C.INT.5}
\]
The effective numerically anticanonical divisor on
an Enoki minimal model makes it parabolic Inoue;
this is the Enoki case in
\cite[Lemma~2.22 and the following paragraph]{Dloussky2006}.
Write $E_i^0$ for its elliptic curve and $C_i$ for
its reduced rational cycle. They are disjoint and
\[
 K_{Z_i^0}+E_i^0+C_i=0.
 \tag{C.INT.6}\hypertarget{tag-00OD}{}\label{cdp:eq:C.INT.6}
\]
The pushforward of $D_i$ has the form
$E_i^0+k_i' C_i$, with $k_i'\ge0$, by intersection
with $E_i^0$ and the kernel of the cycle intersection
matrix. Each nonzero adjacent divisor has the exact
form $E_{ij}=n_{ij}\sigma_i^*C_i$, $n_{ij}\ge1$.
In particular $k_i'\ge\sum_{j\ne i}n_{ij}\ge1$.

\item The canonical evaluation used in \SeriesRef{cdp:st:A.4.4} gives, on a
neighbourhood of this fibre,
$K_Y=\OO_Y(\sum_i k_iS_i)$ for integers $k_i\ge0$;
the line from the base is trivial on the chosen disc.
Put $r_i=(k_i+1)/m_i$. The principal full-fibre
relation yields the following equality in
$\Pic(Z_i)\otimes\Q$:
\[
 K_{Z_i}+D_i
   =\sum_{j\ne i}m_j(r_j-r_i)E_{ij}.
 \tag{C.INT.7}\hypertarget{tag-00OE}{}\label{cdp:eq:C.INT.7}
\]
Indeed subtract $r_i h_i^*F=0$ from the first
formula \eqref{cdp:eq:C.INT.5}. Choose $i$ for which $r_i$ is
maximal. The negative of the right side of
\eqref{cdp:eq:C.INT.7} is an effective rational divisor $A_i$.
Thus $K_{Z_i}+D_i+A_i$ is torsion after clearing
denominators. Push to its minimal model and use
\eqref{cdp:eq:C.INT.6}. Taking Gauduchon degree gives
\[
 k_i'+\sum_{j\ne i}m_j(r_i-r_j)n_{ij}=1,
 \tag{C.INT.8}\hypertarget{tag-00OF}{}\label{cdp:eq:C.INT.8}
\]
since $\deg\OO(C_i)>0$. All terms after $k_i'$
are nonnegative, and $k_i'\ge\sum n_{ij}\ge1$.
Equality forces $k_i'=1$, exactly one neighbour,
$n_{ij}=1$, and $r_j=r_i$. That neighbour is again
a maximum. Propagating along the connected
intersection graph proves that there are exactly
two components $S_1,S_2$, all ratios are equal,
and on both minimal models
\[
 (D_i^{\rm self})_0=E_i^0,
 \qquad (E_{ij})_0=C_i.
 \tag{C.INT.9}\hypertarget{tag-00OG}{}\label{cdp:eq:C.INT.9}
\]

\item Every topologically
torsion line descends through a point blowup, by the
integral Picard blowup formula. Applying the
canonical divisor formula to the second equality
\eqref{cdp:eq:C.INT.5} therefore gives at every blowup
\[
 (D_i^{\rm self})_{\rm new}
   =\sigma^*(D_i^{\rm self})_{\rm old}-E_{\rm exc}.
 \tag{C.INT.10}\hypertarget{tag-00NR}{}\label{cdp:eq:C.INT.10}
\]
Starting with the smooth multiplicity-one elliptic
curve $E_i^0$, effectivity permits centres only on
its current strict transform, and the new
exceptional coefficient is zero. Hence
$D_i^{\rm self}$ is that elliptic strict transform,
and all these centres are disjoint from $C_i$.
The external cycle $\sigma_i^*C_i=C_i$ is therefore
unchanged, and disjoint from the self conductor.

\item It follows that $S_i$ is normal near the external
intersection. Any
resolution-exceptional curve outside $D_i$ is
disjoint from it: adjunction and minimality give
$K_{Z_i}\cdot E\ge0$, while $D_i\cdot E\ge0$,
so $(K_{Z_i}+D_i)\cdot E=0$ forces $D_i\cdot E=0$.
Connected resolution fibres cannot join the
elliptic self conductor to the external rational
cycle. The conductor on the normal surface is a
reflexive ideal, by the duality calculation in
\SeriesRef{cdp:st:A.5.3}, so it has no isolated additional support
on that cycle. Thus normalization is an isomorphism
near its image. Resolution over this image is
crepant by the second formula \eqref{cdp:eq:C.INT.5}, since
$D_i^{\rm self}$ is disjoint from it. Its
exceptional curves are consequently $(-2)$
curves in $C_i$. A whole cycle cannot be contracted,
because its intersection matrix is not negative
definite. The connected fibres which are
contracted are proper chains in that cycle.

\item The intersection scheme $D=S_1\cap S_2$ is a
complete-intersection curve in $Y$. Its pullback
has reduced divisor $C_i$, so it is generically
reduced; its Cohen--Macaulay property then makes
it reduced everywhere. It is connected, since it
is the image of $C_i$. The equality of all $r_i$
means $K_Y+S_1+S_2=rF$ in $\Pic(Y)\otimes\Q$
on the fibre neighbourhood. Adjunction gives
$\omega_D$ torsion, hence $p_a(D)=1$.
Topologically the map $C_i\to D$ contracts just
the proper chains described above and is an
isomorphism elsewhere. It preserves the graph
circle and all its remaining components are
rational. Thus its normalized components are
$\PP^1$, its graph has exactly one cycle, and
$\sum\delta_x=\sum(r_x-1)$. Each contracted
chain has only two ends; every singular point
therefore has two branches and $\delta=1$.
The colength-one ring comparison \eqref{cdp:eq:C.VII.6}
proves that these are ordinary nodes. At least
one node occurs.

\item At such a node the complete-intersection curve
has embedding dimension two. At least one of
the two surface equations has a nonzero linear
part; take it to be $z$. On $z=0$ the other
equation has an ordinary node. The holomorphic
Morse lemma with parameter $z$ therefore writes
it as $xy+h(z)$. Normality of that surface near
$D$ excludes $h\equiv0$, and a change of the
$z$ coordinate writes $h(z)=z^n$, $n\ge1$.
Consequently the original fibration has local
equation
\[
 f=u(x,y,z)\,z^a(xy+z^n)^b,
 \qquad a,b,n\ge1,
 \tag{C.INT.11}\hypertarget{tag-00NS}{}\label{cdp:eq:C.INT.11}
\]
where $u$ is a unit and $a,b$ are the two full
fibre multiplicities. The unit can be removed
by a biholomorphic change of variables.
Put $d=2(a+nb)$ and choose a local logarithm of
$u$. For $f_0=z^a(xy+z^n)^b$, weighted homogeneity
gives
\[
 \Phi(x,y,z)=(x u^{n/d},y u^{n/d},z u^{2/d}),
 \qquad f_0\circ\Phi=u f_0.
 \tag{C.INT.12}\hypertarget{tag-00NT}{}\label{cdp:eq:C.INT.12}
\]
The derivative of $\Phi$ at the origin is
invertible, proving the required equivalence
of the local fibres.

\item For $t\ne0$ small, take a simply connected
sector in the punctured $z$-plane and one
branch of $(t z^{-a})^{1/b}$. An open part
of the smooth fibre is
\[
 xy=g(z):=(t z^{-a})^{1/b}-z^n.
 \tag{C.INT.13}\hypertarget{tag-00NU}{}\label{cdp:eq:C.INT.13}
\]
Consecutive zeros on this logarithmic branch
have angular separation $2\pi b/(a+nb)<2\pi$;
two such zeros fit in one sector. They are
simple, since at either one
$g'(z)=-(a/b+n)z^{n-1}\ne0$.
Join them by a smooth embedded path with no
other zero and take the circle $|x|=|y|$
over its interior. Explicitly, for a continuous
choice of $\arg g$ along that interior, it is
\[
 x=|g|^{1/2}e^{i(\arg g/2+\theta)},\qquad
 y=|g|^{1/2}e^{i(\arg g/2-\theta)}.
 \tag{C.INT.14}\hypertarget{tag-00NV}{}\label{cdp:eq:C.INT.14}
\]
The circle collapses to a smooth disc at each
simple zero, giving an embedded sphere.
The two $x,y$ terms of the standard K\"ahler
form cancel on this sphere, and the $z$ term
vanishes on the one-dimensional path. The
sphere is Lagrangian; its normal bundle is
$JTS^2$ with the opposite Euler orientation,
so its self-intersection is $-2$.
The whole sphere lies in the local chart:
$|z|=O(|t|^{1/(a+nb)})$ and
$|x|,|y|=O(|t|^{n/(2(a+nb))})$.
Its image in the compact smooth torus fibre
would be null-homotopic, because a torus has
contractible universal cover. Its integral
homology class would therefore be zero and
its self-intersection zero, contradicting
$-2$. This contradiction excludes the Enoki
alternative.

\end{enumerate}
\end{proof}






\begin{proposition}[Proved sign and equality for every reducible fibre]\hypertarget{tag-01C2}{}\label{cdp:st:C.2.INT.3}
\StatementRecord{cdp:st:C.2.INT.3}{proved}{01C2}{\SeriesRef{cdp:st:C.2.INT.2}, \SeriesRef{cdp:st:C.2.INT.1}, \SeriesRef{cdp:st:A.5.3}, \SeriesRef{cdp:eq:C.INT.18}, \SeriesRef{cdp:eq:C.INT.21}, \SeriesRef{cdp:eq:C.INT.20}, \SeriesRef{cdp:eq:C.INT.22}, \SeriesRef{cdp:eq:C.INT.26}, \SeriesRef{cdp:eq:C.INT.27}}


Let $f:Y\to B$ be a proper connected-fibre map from a smooth
compact complex threefold with $H^2(Y;\mathbb Q)=0$ and general
fibre a complex two-torus, as in \ref{cdp:st:C.2.INT.2}.
Let $S$ be a connected reducible fibre support.
Then
\[
 e(S)\ge0.
 \tag{C.INT.15}\hypertarget{tag-00NW}{}\label{cdp:eq:C.INT.15}
\]
Equality holds precisely when the normalization of $S$
is a disjoint union of smooth minimal Hopf or
non-K\"ahler elliptic surfaces, and the reduced
conductor image and its reduced inverse image are
disjoint unions of smooth elliptic curves. The maps
between these elliptic curves are finite unramified
covers. This criterion allows arbitrary conductor
scheme structure and arbitrary unequal multiplicities
of the full fibre.
\end{proposition}


\begin{proof}

\begin{enumerate}[label={$\langle1\rangle$\arabic*.},leftmargin=*]
\item Let $\nu:T=\coprod T_i\to S$ be normalization,
$\pi_i:Z_i\to T_i$ its minimal resolutions, and
$\sigma_i:Z_i\to Z_i^0$ the minimal-model maps.
By \ref{cdp:st:C.2.INT.1}--\ref{cdp:st:C.2.INT.2} the $Z_i^0$ are Hopf or
non-K\"ahler elliptic surfaces. They have no rational
curves and have both Euler characteristics
$e(Z_i^0)=\chi(\OO_{Z_i^0})=0$.
Every rational curve of $Z_i$ is exceptional for
$\sigma_i$; their total number is exactly the number
$N_i$ of point blowups, and $e(Z_i)=N_i$.
The elliptic strict transforms cannot be contracted
by $\pi_i$, by \ref{cdp:st:C.2.INT.1}. Thus every
$\pi_i$-exceptional curve is $\sigma_i$-exceptional.

\item The connected fibres of $\pi_i$ consequently map to
points under $\sigma_i$. Normality and
$(\pi_i)_*\OO_{Z_i}=\OO_{T_i}$ descend the latter
map to a proper birational morphism
$\tau_i:T_i\to Z_i^0$, with
\[
 \sigma_i=\tau_i\pi_i.
 \tag{C.INT.16}\hypertarget{tag-00NX}{}\label{cdp:eq:C.INT.16}
\]
This is an actual factorization, not a replacement
of the normalization by its minimal model.
Since point blowups have $R\sigma_{i*}\OO_{Z_i}
=\OO_{Z_i^0}$, the low-degree Leray sequence gives
\[
 0\longrightarrow R^1\tau_{i*}\OO_{T_i}
 \longrightarrow 0
 \longrightarrow \tau_{i*}R^1\pi_{i*}\OO_{Z_i}
 \longrightarrow R^2\tau_{i*}\OO_{T_i}=0.
 \tag{C.INT.17}\hypertarget{tag-00NY}{}\label{cdp:eq:C.INT.17}
\]
The last vanishing follows from fibre dimension at
most one. The middle sheaf $R^1\pi_{i*}\OO_{Z_i}$
is supported at finitely many points, and its
pushforward under $\tau_i$ detects it. Hence it is
zero. This proves that $T_i$ has rational
singularities, and in particular
$\chi(\OO_{T_i})=\chi(\OO_{Z_i})=0$.

\item Let $\Gamma$ be the full conductor image scheme,
$A_i\subset T_i$ its full inverse-image conductor,
and let $D_i\ge0$ be the resolution conductor from
\SeriesRef{cdp:st:A.5.3}. Its line relation
$K_{Z_i}+D_i=h_i^*(K_Y+S)$ has torsion topological
first Chern class. That line descends through
$\sigma_i$. Consequently the exact canonical
divisor formula gives
\[
 D_i=\sigma_i^*D_i^0-K_{Z_i/Z_i^0},
 \qquad D_i^0=(\sigma_i)_*D_i.
 \tag{C.INT.18}\hypertarget{tag-00NZ}{}\label{cdp:eq:C.INT.18}
\]
The divisor $D_i^0$ is nonzero, since $D_i$ contains
a nonzero adjacent component divisor whose
minimal-model pushforward is nonzero by (C.INT.2).
It is a sum of positive multiples of disjoint
smooth elliptic curves. At each point blowup,
\eqref{cdp:eq:C.INT.18} has coefficient
$\operatorname{mult}_pD_{\rm old}-1$ on the new
exceptional curve. Effectivity makes the centre
lie on the old support. Starting from the disjoint
elliptic curves, the reduced support consequently
acquires only rational leaves and subdivisions of
edges. Each of its connected components is one
elliptic curve with rational trees attached.
Writing $r_i$ for its number of rational curves,
we obtain
\[
 e((D_i)_{\rm red})=r_i.
 \tag{C.INT.19}\hypertarget{tag-00O0}{}\label{cdp:eq:C.INT.19}
\]

\item We next prove a simultaneous cohomological
comparison on the entire, possibly nonreduced,
conductor. Relative duality applied to
$R\pi_{i*}\OO_{Z_i}=\OO_{T_i}$ gives
$R\pi_{i*}\omega_{Z_i}=\omega_{T_i}$.
Finite conductor duality identifies
$\omega_{T_i}\otimes\nu^*\omega_S^{-1}$ with
the conductor ideal $\mathcal I_{A_i/T_i}$.
The divisor sequence on $Z_i$, using
$\OO_{Z_i}(-D_i)=\omega_{Z_i}\otimes h_i^*\omega_S^{-1}$,
therefore gives
\[
 R\pi_{i*}\OO_{D_i}=\OO_{A_i}.
 \tag{C.INT.20}\hypertarget{tag-00O2}{}\label{cdp:eq:C.INT.20}
\]
The identical calculation for $\sigma_i$, using
\eqref{cdp:eq:C.INT.18}, gives
\[
 R\sigma_{i*}\OO_{D_i}=\OO_{D_i^0},
 \qquad R\tau_{i*}\OO_{A_i}=\OO_{D_i^0}.
 \tag{C.INT.21}\hypertarget{tag-00O3}{}\label{cdp:eq:C.INT.21}
\]
For completeness the canonical pushforward for a
point blowup is elementary: its relative canonical
line is $\OO(E)$, and
$0\to\OO\to\OO(E)\to\OO_E(-1)\to0$
has zero cohomology in its last term. It gives
$R\sigma_*\omega=\omega$; composition proves the
version used in \eqref{cdp:eq:C.INT.21}.

\item For $M\in\Pic^0(Y)$ its pullback to $Z_i$ descends
to a line $M_i^0\in\Pic^0(Z_i^0)$, by the Picard
blowup isomorphism. Projection formula and
\eqref{cdp:eq:C.INT.20}--\eqref{cdp:eq:C.INT.21} identify the actual spaces
\[
 H^0(A_i,\nu^*M|_{A_i})
   =H^0(D_i^0,M_i^0|_{D_i^0}).
 \tag{C.INT.22}\hypertarget{tag-00O4}{}\label{cdp:eq:C.INT.22}
\]
For a component $nE$ of $D_i^0$, the filtration by
its powers of the defining ideal has graded lines
$M_i^0|_E\otimes\OO_E(-jE)$, $0\le j<n$.
They have degree zero, since $E^2=0$ on the
minimal model. \ref{cdp:st:C.2.INT.1} proves that the map from
$\Pic^0(Y)$ to $\Pic^0(E)$ is nonconstant. We can
therefore choose one $M$ avoiding, for each of
the finitely many such curves and integers $j$,
the prescribed value $\OO_E(jE)$.
Each excluded locus is a proper closed analytic
subset of the connected complex Lie group
$\Pic^0(Y)$; their finite union does not fill it.
All the graded lines are nontrivial of degree
zero on an elliptic curve and have no sections.
The filtration and \eqref{cdp:eq:C.INT.22} prove
\[
 H^0(A_i,\nu^*M)=0\quad\hbox{for every }i.
 \tag{C.INT.23}\hypertarget{tag-00O5}{}\label{cdp:eq:C.INT.23}
\]

\item The conductor inclusion is
$\OO_\Gamma\hookrightarrow\nu_*\OO_A$, where
$A=\coprod A_i$. It now gives $H^0(\Gamma,M)=0$.
The normalization sequence and the preceding
rational-singularity computation give
$\chi(\mathcal Q)=0$ for
$\mathcal Q=\nu_*\OO_T/\OO_S$: here
$\chi(\OO_S)=0$ is the Cartier-divisor
Riemann--Roch calculation under $H^2(Y,\Q)=0$.
Conductor duality is
$\mathcal Q=\omega_\Gamma\otimes\omega_S^{-1}|_\Gamma$.
Twisting on this compact Cohen--Macaulay curve by
an ambient topologically torsion line preserves
Euler characteristic; its degree on every
component is zero. Serre duality thus gives
\[
 \chi(\OO_\Gamma)=-\chi(\mathcal Q)=0.
 \tag{C.INT.24}\hypertarget{tag-00O6}{}\label{cdp:eq:C.INT.24}
\]
The same holds after twisting by $M$, so
$H^0(\Gamma,M)=0$ implies $H^1(\Gamma,M)=0$.
Restriction to $C=\Gamma_{\rm red}$ surjects
on $H^1$, since its kernel is supported in
dimension at most one and has no $H^2$.
Consequently
\[
 \chi(\OO_C)=\chi(M|_C)=h^0(C,M|_C)\ge0.
 \tag{C.INT.25}\hypertarget{tag-00O7}{}\label{cdp:eq:C.INT.25}
\]
This is where the full conductor scheme is used;
no passage to its reduction is assumed harmless.

\item Each normalized component of $C$ has genus zero
or one, since it is dominated by a curve of
$(D_i)_{\rm red}$. Let $n_0$ be the number of
rational-normalized components, and let
$\delta_x$ and $s_x$ be the delta invariant and
branch number of a singular point of $C$.
The exact curve normalization formulas give
\[
 e(C)=\chi(\OO_C)+n_0+
                \sum_x\bigl(\delta_x-(s_x-1)\bigr)\ge0.
 \tag{C.INT.26}\hypertarget{tag-00O8}{}\label{cdp:eq:C.INT.26}
\]
Every summand is nonnegative, because
$\delta_x\ge s_x-1$. Equality holds precisely
when $C$ is a disjoint union of smooth elliptic
curves. Indeed equality first forces $n_0=0$;
then $\chi(\OO_C)=-\sum\delta_x\ge0$
forces every $\delta_x=0$. The converse is immediate.

\item We finish the topological comparison on each
normalization component. A $\pi_i$-exceptional
curve outside $(D_i)_{\rm red}$ is a $(-2)$ curve
disjoint from that support, by the adjunction
argument following (C.INT.10). Each connected
configuration of such curves is a tree: a cycle
or a repeated intersection would contradict
negative definiteness of its matrix with
all diagonal entries $-2$. Such fibres lie
off $A_i$, because a fibre over its conductor
meets the strict conductor and is connected.
Let $a_i$ count their irreducible curves, and
let $b_i$ count exceptional fibres contained
in $(D_i)_{\rm red}$ whose image lies outside
$A_i$. Removing the conductor and contracting
these specified fibres gives the exact equality
\[
 e(T_i\setminus(A_i)_{\rm red})
       =N_i-r_i-a_i+b_i\ge0.
 \tag{C.INT.27}\hypertarget{tag-00O9}{}\label{cdp:eq:C.INT.27}
\]
A tree of $a$ rational curves has Euler number
$a+1$, so its contraction removes $a$; a
positive-conductor fibre removed upstairs but
represented by a point downstairs adds one.
These are all changes, since normalization is
an isomorphism off the conductor. Nonnegativity
follows because the $r_i$ support curves and
the $a_i$ disjoint exceptional curves are
distinct rational curves among the $N_i$
rational curves of $Z_i$.
The normalization is an isomorphism
$T\setminus A\simeq S\setminus C$. Additivity
with \eqref{cdp:eq:C.INT.26}--\eqref{cdp:eq:C.INT.27} proves \eqref{cdp:eq:C.INT.15}.

\item If $e(S)=0$, every nonnegative summand in these
formulas is zero. In particular $C$ is a
union of smooth elliptic curves, $b_i=0$,
and $r_i+a_i=N_i$. Every rational curve in
$(D_i)_{\rm red}$ must then be $\pi_i$-exceptional:
a rational curve cannot map nonconstantly to
one of these elliptic curves. The $a_i$ curves
are exceptional by definition. Thus all $N_i$
rational curves of $Z_i$ are $\pi_i$-exceptional.
If $N_i>0$, the last point blowup in
$\sigma_i$ supplies a $(-1)$ curve on $Z_i$.
Its being $\pi_i$-exceptional contradicts the
minimality of that resolution. Hence $N_i=0$.
There are now no rational curves to contract,
and \ref{cdp:st:C.2.INT.1} excludes contracting elliptic
curves. Therefore $Z_i=T_i=Z_i^0$ is smooth
and minimal. Its reduced conductor is a
disjoint union of elliptic curves. A finite
map between smooth elliptic curves is
unramified by Riemann--Hurwitz, proving the
stated map comparison.
Conversely these normalization and conductor
conditions give $e(T_i)=e((A_i)_{\rm red})=e(C)=0$;
the normalization formula gives $e(S)=0$.
The underlying spaces, and all comparisons
just used, are independent of the positive
full-fibre multiplicities. This completes the
proof including the equality criterion.

\end{enumerate}
\end{proof}






\begin{proposition}[The analytic support comparison used here]\hypertarget{tag-01C4}{}\label{cdp:st:C.2.MH.1}
\StatementRecord{cdp:st:C.2.MH.1}{proved}{01C4}{\SeriesRef{cdp:eq:C.MH.2}, \SeriesRef{cdp:eq:C.MH.3}, \SeriesRef{cdp:eq:C.MH.5}; \cite[Section~3, Theorem~1 and
Corollary~2]{Steenbrink1999}; \cite[Section~4, equation~(3)]{Steenbrink1999}}

Let $U$ be a smooth complex threefold, let $S\subset U$ be a
compact divisor, and suppose that an embedded resolution
$r:U'\to U$ is a sequence of blowups with smooth compact
centres in $S$, with $D=(r^{-1}S)_{\rm red}$ a simple
normal crossings divisor whose components are projective.
Let $V$ be a finite-monodromy rational local system on $U$.
Suppose that the cohomology of the punctured sufficiently
small fibre tubes around $S$, with coefficients in $V$,
vanishes in every degree. Then the mixed Hodge structure
on $H^k(S,V)$ obtained from a projective hyperresolution
is pure of weight $k$.

In this assertion the support-to-restriction morphism is
the actual morphism
\[
 R\Gamma_D(U',r^*V)\longrightarrow
 \underset{D\subset W}{\mathop{\rm colim}}\,
 R\Gamma(W,r^*V)=R\Gamma(D,r^*V).
 \tag{C.MH.1}\hypertarget{tag-00QA}{}\label{cdp:eq:C.MH.1}
\]
Neither a K\"ahler metric on $U$ nor a projective
compactification of $U$ is needed for this comparison.
\end{proposition}


\begin{proof}

\begin{enumerate}[label={$\langle1\rangle$\arabic*.},leftmargin=*]
\item First take constant rational coefficients. Write $D_I$
for the intersection of the components indexed by a
nonempty set $I$, and $D^{[p]}=\coprod_{|I|=p+1}D_I$.
All these spaces are smooth and projective. The ordinary
\v Cech de Rham complex on $D^{[\bullet]}$, with its
\v Cech weight filtration and its differential-form
filtration, is therefore a cohomological mixed Hodge
complex $K(D)$ for $\mathbb Q_D$. Its cohomology in
degree $k$ has weights at most $k$.

\item Here is the precise analytic construction of the other
complex and of the arrow. For each stratum $C=D_I$, put
\[
 K(C\log D)_{\mathbb C}
   =\Omega_{U'}^\bullet(\log D)\otimes\OO_C.
 \tag{C.MH.2}\hypertarget{tag-00QJ}{}\label{cdp:eq:C.MH.2}
\]
The tensor product in \eqref{cdp:eq:C.MH.2} is a complex: in local
normal crossings coordinates, the ideal of $C$ is
preserved by the logarithmic differential. Endow it with
the quotient residue filtration and the filtration by
form degree. The quotient map from the restriction of
$\Omega_{U'}^\bullet(\log D)$ is a filtered
quasi-isomorphism to \eqref{cdp:eq:C.MH.2}. Its rational realization is
$i_C^*Rj_*\mathbb Q_{U'\setminus D}$, where $j$ is the
open inclusion. To see the comparison, apply the
iterated Poincar\'e residue: its residue-degree-$a$
quotient is a direct sum of the de Rham complexes of
$D_J\cap C$, shifted by $[-a]$ and with Tate twist
$(-a)$. The ideal and its differentials restrict these
forms to precisely $D_J\cap C$. This is the local
coordinate proof of \cite[Section~3, Theorem~1 and
Corollary~2]{Steenbrink1999}; it uses holomorphic
coordinates and the projectivity, or compact K\"ahler
property, of the displayed strata. It does not use a
compactification of their ambient neighbourhood.

\item Totalizing \eqref{cdp:eq:C.MH.2} over $D^{[\bullet]}$ gives the
cohomological mixed Hodge complex $K(D\log D)$, with
rational realization $i_D^*Rj_*\mathbb Q_{U'\setminus D}$.
The natural inclusion $K(D)\to K(D\log D)$ is a
morphism of mixed Hodge complexes. Its mixed cone,
shifted by $[-1]$, computes cohomology with support in
$D$; the cone projection is exactly \eqref{cdp:eq:C.MH.1}. This is
the excision comparison of
\cite[Section~4, equation~(3)]{Steenbrink1999}, applied
on the compact divisor. All complexes just used are
defined on that divisor, and all purity checks are on
its projective strata.

\item For completeness, the weight bound on the source can
be checked without invoking an invariant-cycle theorem.
The closed-support Mayer--Vietoris complex, together
with the holomorphic Thom isomorphisms, gives
\[
 E_1^{-p,q}=
 \bigoplus_{|I|=p+1}
 H^{q-2p-2}(D_I,\mathbb Q)(-p-1)
 \ \Longrightarrow\ H_D^{q-p}(U',\mathbb Q),
 \qquad p\ge0.
 \tag{C.MH.3}\hypertarget{tag-00QK}{}\label{cdp:eq:C.MH.3}
\]
The first differential is the alternating sum of Gysin
maps for $D_I\hookrightarrow D_{I\setminus\{i\}}$.
The comparison with the preceding mixed cone can be
verified by induction on the number of components.
For one smooth divisor the residue exact sequence
identifies
\[
 \operatorname{Cone}\bigl[\Omega_{U'}^\bullet
       \longrightarrow\Omega_{U'}^\bullet(\log D_i)\bigr][-1]
       \simeq\Omega_{D_i}^\bullet[-2](-1),
\]
where $(-1)$ records both the shift of the Hodge
filtration by one and the rational $2\pi i$ residue.
For a union $A\cup B$, use the actual support triangle
$R\Gamma_{A\cap B}\to R\Gamma_A\oplus R\Gamma_B
\to R\Gamma_{A\cup B}$. Iterated residues commute
with its inclusion maps, with the alternating signs
fixed by the ordering of the omitted coordinates.
Induction identifies its totalization with the same
mixed cone. These comparisons agree on overlaps,
since they are the residue maps for the holomorphic
divisor line bundles. Thus \eqref{cdp:eq:C.MH.3} compares the actual support
complex, including its arrow to $K(D)$, rather than
only its unfiltered vector spaces. Its displayed term
has pure weight $q$. In total degree $k=q-p$ this is
weight $k+p\ge k$. It follows that $H_D^k(U',\mathbb Q)$
has weights at least $k$.

\item Finite coefficients are handled by taking a finite
\'{e}tale cover on which $V$ is constant and then a
rational direct summand of a finite direct sum of its
regular representation.
The inverse images of the $D_I$ are still projective:
a finite holomorphic \'{e}tale cover of a
projective manifold is projective, since the pullback
of an ample bundle is positive. The deck action acts
on the displayed mixed Hodge complexes. Rational
idempotents consequently give exactly the same
comparison and weight bounds with coefficients in $V$.

\item By the assumed punctured-tube acyclicity, the support
long exact sequence makes \eqref{cdp:eq:C.MH.1} an isomorphism on
cohomology. Properness of $r$ identifies the punctured
tubes upstairs and downstairs, since its centres lie
in $S$. The map is a morphism of mixed Hodge structures,
whose source has weights at least $k$ and whose target
has weights at most $k$. Strictness therefore proves
that $H^k(D,r^*V)$ is pure of weight $k$.

\item We finally make the comparison with the original
support explicit. The proper degree-one map between
the smooth ambient manifolds has the trace map, defined
by oriented Poincar\'e duality and proper pushforward,
and the projection formula gives
\[
 r_*r^*=\operatorname{id}:H^k(U,V)\longrightarrow H^k(U,V).
 \tag{C.MH.4}\hypertarget{tag-00QL}{}\label{cdp:eq:C.MH.4}
\]
The same identity holds on every smaller ambient
neighbourhood. Passing to the direct limit of these
identities gives an injection
\[
                  r^*:H^k(S,V)\hookrightarrow H^k(D,r^*V).
 \tag{C.MH.5}\hypertarget{tag-00QM}{}\label{cdp:eq:C.MH.5}
\]
The identifications in this limit are the usual
continuity identifications for a compact triangulable
subset. They can equally be made using cofinal
triangulated regular neighbourhoods retracting onto
$S$ and $D$. Properness ensures that inverse images of
neighbourhoods are cofinal; thus the map in \eqref{cdp:eq:C.MH.5}
is the pullback by $D\to S$.

\item This last map is a mixed Hodge morphism as well. One
may construct the projective hyperresolution of $S$
from its projective normalization resolution and
resolutions of the conductor and exceptional curves.
To compare it with $D^{[\bullet]}$, resolve their
fibre products over $S$. Each such fibre product is a
closed analytic subspace of a product of projective
spaces already occurring in the two resolutions;
Chow's theorem makes it algebraic and projective,
and its resolution is projective. The resulting
common hyperresolution realizes \eqref{cdp:eq:C.MH.5} by pullbacks
between projective strata. Cohomological descent then
identifies this with the stated mixed Hodge structures.
Injectivity into the pure target proves the assertion.

\end{enumerate}
\end{proof}






\begin{proposition}[Euler number from a ruled normalization]\hypertarget{tag-01C5}{}\label{cdp:st:C.2.MH.2}
\StatementRecord{cdp:st:C.2.MH.2}{proved}{01C5}{\SeriesRef{cdp:st:D.1.9}, \SeriesRef{cdp:st:D.1.8}, \SeriesRef{cdp:st:C.2.MH.1}, \SeriesRef{cdp:eq:C.MH.9}, \SeriesRef{cdp:eq:C.MH.8}, \SeriesRef{cdp:eq:C.MH.13}, \SeriesRef{cdp:eq:C.MH.11}, \SeriesRef{cdp:eq:C.MH.14}}

Let $f:Y\to\PP^1$ satisfy the full CDP hypotheses
$b_1(Y)=b_2(Y)=0$ with smooth torus general fibre.
Let $mS$ be a full fibre with irreducible reduced support,
and let $Z\to T\to S$ be a resolution of its
normalization. If $Z$ is rational, then
\[
                              e(S)=2.
 \tag{C.MH.6}\hypertarget{tag-00QN}{}\label{cdp:eq:C.MH.6}
\]
If $Z$ is ruled over a curve of genus $g\ge1$, then
\[
                              e(S)=2-2g.
 \tag{C.MH.7}\hypertarget{tag-00QO}{}\label{cdp:eq:C.MH.7}
\]
The calculation retains all multiplicities in the full
fibre and places no normal crossings assumption on $S$.
\end{proposition}


\begin{proof}

\begin{enumerate}[label={$\langle1\rangle$\arabic*.},leftmargin=*]
\item The integral class of \SeriesRef{cdp:st:D.1.9} restricts nontrivially to
$S$ and, by \SeriesRef{cdp:st:D.1.8}, nontrivially to the nearby torus.
More precisely, restrict the section
$\gamma\in H^0(\PP^1,R^1f_*\mathbb Z)$ to a small
disc $\Delta$ about this fibre. The integral Leray edge
isomorphism
$H^1(f^{-1}\Delta,\mathbb Z)
 \simeq H^0(\Delta,R^1f_*\mathbb Z)$ supplies the
actual integral class on the tube: the two adjacent
base groups $H^1(\Delta,\mathbb Z)$ and
$H^2(\Delta,\mathbb Z)$ vanish.
Choose a prime $\ell$ for which its restriction to that
torus is nonzero modulo $\ell$. On a sufficiently small
fibre tube $U$, its reduction defines a connected cyclic
\'{e}tale cover with group $\mathbb Z/\ell$. Let $V$ be
the rational augmentation local system, that is, the
kernel of the trace from the regular representation
to the constant representation. Its rank is
$d=\ell-1$, it is self-dual, and $H^0(S,V)=0$.
Over $\mathbb C$ it is the sum of the nontrivial
characters of this cyclic group. Since $\ell$ is prime,
each of these characters restricts nontrivially to
the nearby torus $F$. The group-cochain complex for
$\mathbb Z^4$ with such a character is acyclic: one of
the four commuting operators $\rho(e_i)-1$ is
invertible, and its inverse gives a contraction of the
Koszul complex. Thus
\[
 H^j(F,V)=0\quad\hbox{for every }j,
 \qquad H^j(U\setminus S,V)=0\quad\hbox{for every }j.
 \tag{C.MH.8}\hypertarget{tag-00QP}{}\label{cdp:eq:C.MH.8}
\]
The second assertion is the Leray sequence of the
punctured-disc torus bundle. This argument uses the
smooth fibres of the original map, so it also applies
when the full central fibre is $mS$.

\item Choose an embedded resolution of $(U,S)$. Its strict
transform is a blowup of $Z$, hence projective. Every
other exceptional surface is obtained by blowups from
a projective bundle over a compact smooth curve, or
from $\PP^2$. It is therefore projective. All centres
are compact and lie in the central fibre. \ref{cdp:st:C.2.MH.1}
applies and gives
\[
                         H^3(S,V)\text{ pure of weight }3.
 \tag{C.MH.9}\hypertarget{tag-00QQ}{}\label{cdp:eq:C.MH.9}
\]

\item Let $A\subset T$ and $\Gamma\subset S$ be the reduced
conductor inverse image and image. The finite
normalization square is an actual topological pushout.
Its cohomological descent sequence, with the pulled
back local system, gives an exact sequence of mixed
Hodge structures
\[
\begin{gathered}
0\longrightarrow C_V\longrightarrow H^3(S,V)
  \longrightarrow H^3(T,V)\longrightarrow0,\\
C_V=\operatorname{coker}\bigl[H^2(T,V)\oplus H^2(\Gamma,V)
                         \longrightarrow H^2(A,V)\bigr].
\end{gathered}
 \tag{C.MH.10}\hypertarget{tag-00QB}{}\label{cdp:eq:C.MH.10}
\]
This also covers the normal case, with the conductor
terms zero. The left term is pure of weight two:
$H^2(A,V)$ is the top cohomology of a compact curve
with a finite local system, and hence pure of weight
two. Strictness and \eqref{cdp:eq:C.MH.9} make that term zero.
Consequently
\[
                           H^3(S,V)\simeq H^3(T,V).
 \tag{C.MH.11}\hypertarget{tag-00QC}{}\label{cdp:eq:C.MH.11}
\]

\item For the resolution $\pi:Z\to T$, the restriction
$H^2(Z,V)\to H^2(E,V)$ to its exceptional curve is
surjective. Indeed $V$ is constant on each exceptional
fibre, since it is pulled back from $T$; the Gysin
classes of its irreducible exceptional curves restrict
by their negative-definite intersection matrix. That
matrix is invertible over $\mathbb Q$, separately
for each exceptional fibre and each coefficient
vector. The resolution descent sequence therefore
gives
\[
 H^3(T,V)\simeq H^3(Z,V),\qquad
 H^4(S,V)\simeq H^4(T,V)\simeq H^4(Z,V).
 \tag{C.MH.12}\hypertarget{tag-00QD}{}\label{cdp:eq:C.MH.12}
\]
The last two identifications also follow directly
because all the omitted loci have real dimension two.

\item There is a second, purely topological consequence of
\eqref{cdp:eq:C.MH.8}. The support sequence and oriented local
Poincar\'e duality give
\[
 H^k(S,V)\simeq H_S^k(U,V)
             \simeq H_{6-k}(S,V).
 \tag{C.MH.13}\hypertarget{tag-00QE}{}\label{cdp:eq:C.MH.13}
\]
Again these are the natural direct-limit maps on
sufficiently small tubes. Since $S$ has real
dimension four, \eqref{cdp:eq:C.MH.13} gives $H^1(S,V)=0$ and,
using the self-duality of $V$,
$h^2(S,V)=h^4(S,V)$. Therefore \eqref{cdp:eq:C.MH.11}--\eqref{cdp:eq:C.MH.13}
give the exact numerical identity
\[
 d\,e(S)=\chi(S,V)=2h^4(Z,V)-h^3(Z,V).
 \tag{C.MH.14}\hypertarget{tag-00QF}{}\label{cdp:eq:C.MH.14}
\]

\item If $Z$ is rational, the pullback of $V$ is constant,
$h^4(Z,V)=d$, and $h^3(Z,V)=0$; this proves \eqref{cdp:eq:C.MH.6}.
If $Z$ is ruled over $B$, blowups leave its fundamental
group and its third cohomology unchanged, and the
pulled-back local system comes from $B$. The
projective-bundle formula gives
$H^3(Z,V)=H^1(B,V)(-1)$. If the cyclic character is
trivial on $B$, its dimensions are
$h^4(Z,V)=d$ and $h^3(Z,V)=2gd$. Otherwise every
nontrivial character of the prime-order group remains
nontrivial on $B$, so $h^4(Z,V)=0$ and
$h^3(Z,V)=(2g-2)d$. Both substitutions in \eqref{cdp:eq:C.MH.14}
give \eqref{cdp:eq:C.MH.7}.

\end{enumerate}
\end{proof}






\begin{proposition}[The contracted irrational section contributes zero]\hypertarget{tag-01C6}{}\label{cdp:st:C.2.MH.3}
\StatementRecord{cdp:st:C.2.MH.3}{proved}{01C6}{\SeriesRef{cdp:st:C.2.CR.1}, \SeriesRef{cdp:st:C.2.CR.5}, \SeriesRef{cdp:st:C.2.MH.2}, \SeriesRef{cdp:st:C.2.UN.7}}

In the contracted-horizontal-section situation of
\ref{cdp:st:C.2.CR.1}, the normalization resolution is ruled over
an elliptic curve and
\[
                    e(S)=0,\qquad n=v-N_0.
 \tag{C.MH.15}\hypertarget{tag-00QG}{}\label{cdp:eq:C.MH.15}
\]
\end{proposition}


\begin{proof}
\begin{enumerate}[label={$\langle1\rangle$\arabic*.},leftmargin=*]
\item 
\noindent\textbf{Step 1.} \ref{cdp:st:C.2.CR.5} proves that the ruling genus is one.
\ref{cdp:st:C.2.MH.2} then gives $e(S)=0$ for the actual fibre.
The exact equality (C.CR.3) gives $n=v-N_0$.
The later unconditional calculation \ref{cdp:st:C.2.UN.7}
gives $e(S)\ge2$ and $\pi_1(S)=1$ for this
contracted irrational configuration. Thus its
comparison with the necessary equations here
excludes the case entirely under the CDP
hypotheses; this reduction supplies no example
of an actual CDP fibre.

\end{enumerate}
\end{proof}






\begin{proposition}[Every irreducible class VII normalization contributes zero]\hypertarget{tag-01C7}{}\label{cdp:st:C.2.MH.4}
\StatementRecord{cdp:st:C.2.MH.4}{proved}{01C7}{\SeriesRef{cdp:st:C.2.VII.4}, \SeriesRef{cdp:st:C.2.MH.2}, \SeriesRef{cdp:eq:C.MH.17}, \SeriesRef{cdp:st:C.2.UN.1}, \SeriesRef{cdp:st:C.2.VII.3}}

Under the full CDP torus-fibration hypotheses, an
irreducible fibre support whose normalization resolution
is a surface of class VII has
\[
                              e(S)=0.
 \tag{C.MH.16}\hypertarget{tag-00QH}{}\label{cdp:eq:C.MH.16}
\]
Consequently every such actual fibre satisfies the
equality data of \ref{cdp:st:C.2.VII.4}.
\end{proposition}


\begin{proof}

\begin{enumerate}[label={$\langle1\rangle$\arabic*.},leftmargin=*]
\item Use the prime-order rational augmentation local system
$V$ of \ref{cdp:st:C.2.MH.2}, of rank $d=\ell-1$. Its construction,
the torus acyclicity (C.MH.8), and the support duality
(C.MH.13) do not require a projective normalization.
They give $h^0(S,V)=h^1(S,V)=0$ and
$h^2(S,V)=h^4(S,V)$. Conductor descent still gives the
surjection $H^3(S,V)\to H^3(T,V)$ and the isomorphism
$H^4(S,V)\simeq H^4(T,V)$. The negative-definite
exceptional Gysin matrix in the proof of (C.MH.12)
is valid for the resolution of any normal complex
surface. It gives
\[
 \begin{gathered}
 h^3(S,V)\ge h^3(Z,V),\qquad h^4(S,V)=h^4(Z,V),\\
 d\,e(S)\le 2h^4(Z,V)-h^3(Z,V).
 \end{gathered}
 \tag{C.MH.17}\hypertarget{tag-00QI}{}\label{cdp:eq:C.MH.17}
\]
No mixed Hodge assertion is used in these comparisons.

\item If the cyclic character restricts nontrivially to $Z$,
all its nontrivial powers remain nontrivial because
$\ell$ is prime. Poincar\'e duality on the compact
oriented surface gives $h^4(Z,V)=0$, and \eqref{cdp:eq:C.MH.17}
gives $e(S)\le0$. If its restriction is trivial,
$h^4(Z,V)=d$ and $h^3(Z,V)=d$: a class VII surface
has $b_1=b_3=1$, and point blowups preserve $b_3$.
Thus $e(S)\le1$ in this second case. \ref{cdp:st:C.2.UN.1} proves
that the Euler number of every actual torus-degeneration
fibre is even, so the second case also gives
$e(S)\le0$. The nonnegative lower bound of \ref{cdp:st:C.2.VII.3}
now gives \eqref{cdp:eq:C.MH.16} in both cases. \ref{cdp:st:C.2.VII.4} supplies
its exact normalization-and-conductor equality data.

\end{enumerate}
\end{proof}






\section{Normal-crossing topology and canonical-flat component comparisons}\hypertarget{tag-034B}{}\label{sec:II:all:08}


\begin{proposition}[Euler localization for torus fibrations]\hypertarget{tag-01FG}{}\label{cdp:st:C2a}
\StatementRecord{cdp:st:C2a}{proved}{01FG}{the definitions and constructions stated in the proof}

Let $g:Y\to C$ be a proper surjective holomorphic map from a connected compact complex manifold to a compact Riemann surface, with connected general fibre $F$. Let $\Sigma\subset C$ contain every critical value. Then $\Sigma$ can be chosen finite, and
\[
 e(Y)=e(C\setminus\Sigma)e(F)+\sum_{c\in\Sigma}e((Y_c)_{\rm red}). \tag{C.1}\hypertarget{tag-008N}{}\label{cdp:eq:C.1}
\]
For a complex threefold with general fibre a complex two-torus this becomes
$\int_Yc_3(T_Y)=e(Y)=\sum_{c\in\Sigma}e((Y_c)_{\rm red})$.
\end{proposition}


\begin{proof}
\begin{enumerate}[label={$\langle1\rangle$\arabic*.},leftmargin=*]
\item 
\noindent\textbf{Step 1.} The critical locus is a closed analytic subset of $Y$. Properness and the proper mapping theorem make its image analytic in $C$; generic smoothness over characteristic zero makes this image proper. Thus it is finite. Over $C\setminus\Sigma$, the underlying differentiable map is a proper submersion. Ehresmann's theorem makes it a smooth locally trivial fibre bundle. Euler characteristic with compact supports is additive for the finite stratification by this open part and the fibres; triangulating the base and trivializing over its cells shows that the open part contributes $e_c(C\setminus\Sigma)e(F)$. The real dimensions of a complex curve and its punctured strata are even, so $e_c(C\setminus\Sigma)=e(C\setminus\Sigma)$. Nilpotents change no underlying topological space. This proves \eqref{cdp:eq:C.1}. A two-torus has the cohomology of the real four-torus and Euler number $(1-1)^4=0$. Finally the top Chern class of the complex tangent bundle is the Euler class for the complex orientation, giving the last identity.

\end{enumerate}
\end{proof}






\begin{proposition}[Counterexample: Nonnegativity of every fibre in unrestricted torus fibrations]\hypertarget{tag-01FH}{}\label{cdp:st:C2b}
\StatementRecord{cdp:st:C2b}{proved}{01FH}{the definitions and constructions stated in the proof}

There is a smooth projective threefold $Y$ with a proper connected-fibre map $Y\to\PP^1$, smooth general fibre an abelian surface, and a reduced normal-crossing fibre of Euler number $-2$. There is also such a threefold with a singular reduced fibre of Euler number zero whose three irreducible components are a blown-up abelian surface, a ruled surface over a genus-two curve, and $\PP^2$.
\end{proposition}


\begin{proof}

\begin{enumerate}[label={$\langle1\rangle$\arabic*.},leftmargin=*]
\item Let $C_2$ be the smooth projective completion of $w^2=z^6-1$. Its six simple branch points and the unramified points at infinity give genus two by Riemann--Hurwitz. Let $A=\operatorname{Jac}(C_2)$ and embed $C_2$ in $A$ by $p\mapsto\OO(p-p_*)$. This is injective because a principal divisor $p-q$ for distinct points would define a degree-one map $C_2\to\PP^1$. It is an immersion because the canonical system of a genus-two curve has no base point; the explicit differentials $dz/w,z\,dz/w$ verify that property, including at infinity. A proper injective immersion is an embedding. Adjunction on the torus gives $N_{C_2/A}=K_{C_2}$.

\item Blow up $A\times\PP^1$ along $C_2\times\{0\}$, and denote the resulting smooth projective threefold by $Y$. Compose its blowdown with projection to $\PP^1$. The exceptional divisor is $E=\PP(K_{C_2}\oplus\OO_{C_2})$. The strict transform of $A\times\{0\}$ is $A$ itself, since blowing up a Cartier divisor in a smooth surface is an isomorphism. These two components meet along a section isomorphic to $C_2$. In local coordinates $(u,v,t)$ with centre $(u,t)=(0,0)$, a blowup chart has $t=u\tau$, so the fibre is the reduced normal-crossing divisor $u\tau=0$. The other chart and the complement give the claimed global map and connectedness of its fibres. Euler additivity gives
\[
 e((Y_0)_{\rm red})=e(A)+e(E)-e(C_2)=0+2(-2)-(-2)=-2.
\]
Every other fibre is $A$.

\item Now blow up one point of the strict-transform component $A$ away from $C_2$. The new central fibre has components $\operatorname{Bl}_pA$, $E$, and $\PP^2$. Its pairwise intersections are $C_2$ and $\PP^1$, with no triple intersection. Its Euler number is
\[
1-4+3-(-2)-2=0.
\]
Both modifications have smooth total space and reduced normal-crossing central fibre. For completeness the integral cohomology is torsion-free. The Betti numbers of $A\times\PP^1$ are $(1,4,7,8,7,4,1)$; the blowup formula adds $H^{k-2}(C_2,\Z)$ in degree $k$, giving $(1,4,8,12,8,4,1)$. The point blowup adds one generator in degrees two and four, giving $(1,4,9,12,9,4,1)$. These formulas follow from excision for the exceptional projective bundle and its hyperplane class, whose restriction identifies the new relative summands. Their Euler numbers are respectively $-2$ and zero, independently checking the fibre computation. The exact arithmetic is in \texttt{code/C\_euler.py}.

\end{enumerate}
\end{proof}






\begin{proposition}[Euler numbers in semistable toric rank-two models]\hypertarget{tag-01FI}{}\label{cdp:st:C2c}
\StatementRecord{cdp:st:C2c}{proved}{01FI}{\cite[Section~6, especially (6.3)--(6.7)]{Mumford1972}}

Consider a rank-two Mumford model with smooth three-dimensional total space obtained from a unimodular periodic triangulation of $\R^2$, divided by a rank-two translation lattice acting freely on the total model. Let $V,E,F$ be the numbers of vertices, edges and triangles in the quotient triangulation. Its reduced central fibre has Euler number
\[
 e(S)=F=2V\in2\Z_{>0}.                                \tag{C.2}\hypertarget{tag-008Y}{}\label{cdp:eq:C.2}
\]
Every positive even number occurs in this class. A rank-one model assembled from $\PP^1$-bundles over elliptic curves with their boundary elliptic sections identified has Euler number zero.
\end{proposition}


\begin{proof}

\begin{enumerate}[label={$\langle1\rangle$\arabic*.},leftmargin=*]
\item The toric orbit strata of a rank-two central fibre have dimensions two, one and zero. Its two-dimensional strata are copies of $(\C^*)^2$, its one-dimensional strata copies of $\C^*$, and its zero-dimensional strata correspond to triangles of the quotient triangulation. The first two have compactly supported Euler number zero. Hence the Euler number is $F$. The quotient complex is a real two-torus, so $V-E+F=0$, and every edge has two incident triangles, giving $3F=2E$, with incidences counted even when a face has identified edges. Thus $F=2V$.

\item For existence, use the lattice $N=\mathbb Z^2\oplus\mathbb Z$
and the standard triangulation of $\mathbb R^2$ whose triangles are
\[
 [u,u+e_1,u+e_1+e_2],\qquad
 [u,u+e_2,u+e_1+e_2],\qquad u\in\mathbb Z^2.
\]
These are the Delaunay triangles for the positive integral quadratic
form $q(u_1,u_2)=u_1^2-u_1u_2+u_2^2$.
Take the cones over these triangles at height one and all their
faces. They form a fan. Every three-dimensional cone is generated
by a basis of $N$, because its triangle is unimodular. Gluing its
affine toric charts therefore gives a smooth complex threefold
$\mathcal X$ with a holomorphic map to the affine line. On each
maximal chart there are coordinates $x_1,x_2,x_3$ in which the map
is exactly $t=x_1x_2x_3$. The central support is consequently
reduced and has normal crossings. Over $t\ne0$ the fibre is the
common algebraic torus $(\mathbb C^*)^2$.

\item Let $\Lambda_N=\mathbb Z(Ne_1)\oplus\mathbb Z e_2$.
For $\ell\in\Lambda_N$, the lattice shear
\[
 (u,k)\longmapsto(u+k\ell,k)
\]
preserves the fan and induces an automorphism over the base.
On the punctured torus chart it is
\[
 (z_1,z_2,t)\longmapsto
 (t^{\ell_1}z_1,t^{\ell_2}z_2,t).
\]
Restrict to $|t|<\epsilon<1$. This action is free: away from the
central fibre, a fixed point would imply $t^{\ell_i}=1$ for both
coordinates and hence $\ell=0$; on the central fibre, a fixed
point would require its bounded indexing face of the triangulation
to be invariant under translation by $\ell$, which again forces
$\ell=0$.

\item Here is also the required analytic properness check. For $t\ne0$,
put $v_i=\log|z_i|/\log|t|$. In the chart of a triangle with
vertices $u_1,u_2,u_3$, write
\[
 v=\sum_{j=1}^3\lambda_j u_j,\qquad
 \lambda_j=\frac{\log|x_j|}{\log|t|},\qquad
 \sum_j\lambda_j=1.
\]
On a compact subset of finitely many charts the $|x_j|$ are
bounded. As $t$ tends to zero, this gives a uniform lower bound
for every $\lambda_j$, and their sum then gives a uniform upper
bound. Thus $v$ remains bounded there. Away from zero the same
conclusion follows from compactness in the torus charts.
Translation by $\ell$ sends $v$ to $v+\ell$. Consequently only
finitely many elements can carry a fixed compact set into itself,
or make two fixed compact sets intersect. On the central fibre
this is also immediate from the finitely many indexing faces.
The action is properly discontinuous, so its quotient is a smooth
complex threefold.

\item The quotient map to the disc is proper. Indeed, translate any
punctured-fibre point so that $v$ belongs to a fixed compact
fundamental parallelogram for $\Lambda_N$. Only finitely many
triangles meet that parallelogram. Choose one containing $v$;
then all its barycentric coordinates are nonnegative, so
$|x_j|=|t|^{\lambda_j}\le1$. These finitely many closed unit
polydiscs provide compact representatives for every sequence
approaching the central fibre. The central fibre itself has
finitely many toric components modulo $\Lambda_N$, and each is
compact: the star of a vertex is a complete two-dimensional fan.
These facts prove compactness over every closed smaller disc.
For $t\ne0$ the quotient fibre is explicitly
\[
 \mathbb C^*/t^{N\mathbb Z}\ \times\
 \mathbb C^*/t^{\mathbb Z},
\]
a product of elliptic curves. This supplies the actual convergent
analytic degeneration, including its smooth total space and
properness, without inferring analytic convergence from a formal
model. The quadratic form supplies the usual polarization data:
its difference under every lattice translation is an integral
affine function. Thus the construction is the analytic realization
of the polyhedral model described in
\cite[Section~6, especially (6.3)--(6.7)]{Mumford1972}.
That source formulates its construction over a complete discrete
valuation ring; the preceding chart and compactness argument is
the complex-disc comparison used here.

\item The quotient triangulation has $V=N$ and $F=2N$, as required.
For rank one every normalization component is a $\PP^1$-bundle over an elliptic curve, with Euler number zero, and the identified boundary sections also have Euler number zero. Additivity gives the result.

\end{enumerate}
\end{proof}




\begin{proposition}[All rank-four monodromy stratum signatures]\hypertarget{tag-01CO}{}\label{cdp:st:C.2.UN.11}
\StatementRecord{cdp:st:C.2.UN.11}{proved}{01CO}{\SeriesRef{cdp:eq:C.MS.6}, \SeriesRef{cdp:st:C.2.UN.1}, \SeriesRef{cdp:st:C.2.UN.4}; \cite[\S2, pp.~237--243, and \S3]{ACampo1975}; \cite[Theorem~1, p.~235,
and proof on pp.~243--244]{ACampo1975}; \cite[Theorem~1.6]{BierstoneMilman1997}}

Let $f:U\to\Delta$ be a proper holomorphic map from a
smooth complex threefold, with connected fibres and
smooth fibres complex two-tori away from zero. Take
an embedded resolution, unchanged away from the
central fibre, and write its full central divisor as
$\sum_jN_jE_j$, with simple normal crossings. Put
\[
 E_j^\circ=E_j\setminus\bigcup_{i\ne j}E_i,
 \qquad a_n=\sum_{N_j=n}e(E_j^\circ).
 \tag{C.MS.1}\hypertarget{tag-00RH}{}\label{cdp:eq:C.MS.1}
\]
Let $T$ be the actual monodromy on $H^1(F,\Z)$,
and let $M$ be the order of its semisimple part.
Then $T\in\operatorname{SL}_4(\Z)$ is
quasi-unipotent, and the following exact maps
determine the entire sequence $(a_n)$:
\[
 \begin{aligned}
 L_k&=\det(I-T^k)
       =\sum_{n\mid k}n a_n,\\
 a_n&=\frac1n\sum_{d\mid n}\mu(n/d)\det(I-T^d).
 \end{aligned}                                  \tag{C.MS.2}\hypertarget{tag-00RJ}{}\label{cdp:eq:C.MS.2}
\]
Here $\mu$ is the integer M\"obius function.
The complete table follows; $\Phi_d$ denotes
the $d$th cyclotomic polynomial and every
unlisted $a_n$ is zero.

\begin{center}
\small
\begin{tabular}{ccl}
$P_T$ & $M$ & Nonzero $a_n$\\\hline
$\Phi_{1}^{4}$ & 1 & $\text{all zero}$ \\
$\Phi_{1}^{2}\Phi_{2}^{2}$ & 2 & $\text{all zero}$ \\
$\Phi_{1}^{2}\Phi_{3}$ & 3 & $\text{all zero}$ \\
$\Phi_{1}^{2}\Phi_{4}$ & 4 & $\text{all zero}$ \\
$\Phi_{1}^{2}\Phi_{6}$ & 6 & $\text{all zero}$ \\
$\Phi_{2}^{4}$ & 2 & $a_{1}=16,\quad a_{2}=-8$ \\
$\Phi_{3}^{2}$ & 3 & $a_{1}=9,\quad a_{3}=-3$ \\
$\Phi_{2}^{2}\Phi_{4}$ & 4 & $a_{1}=8,\quad a_{2}=-4$ \\
$\Phi_{4}^{2}$ & 4 & $a_{1}=4,\quad a_{2}=6,\quad a_{4}=-4$ \\
$\Phi_{5}$ & 5 & $a_{1}=5,\quad a_{5}=-1$ \\
$\Phi_{2}^{2}\Phi_{3}$ & 6 & $a_{1}=12,\quad a_{2}=-6,\quad a_{3}=-4,\quad a_{6}=2$ \\
$\Phi_{2}^{2}\Phi_{6}$ & 6 & $a_{1}=4,\quad a_{2}=-2,\quad a_{3}=4,\quad a_{6}=-2$ \\
$\Phi_{3}\Phi_{6}$ & 6 & $a_{1}=3,\quad a_{2}=3,\quad a_{3}=-1,\quad a_{6}=-1$ \\
$\Phi_{6}^{2}$ & 6 & $a_{1}=1,\quad a_{2}=4,\quad a_{3}=5,\quad a_{6}=-4$ \\
$\Phi_{8}$ & 8 & $a_{1}=2,\quad a_{2}=1,\quad a_{4}=3,\quad a_{8}=-2$ \\
$\Phi_{10}$ & 10 & $a_{1}=1,\quad a_{2}=2,\quad a_{5}=3,\quad a_{10}=-2$ \\
$\Phi_{3}\Phi_{4}$ & 12 & $a_{1}=6,\quad a_{2}=3,\quad a_{3}=-2,\quad a_{4}=-3,\quad a_{6}=-1,\quad a_{12}=1$ \\
$\Phi_{4}\Phi_{6}$ & 12 & $a_{1}=2,\quad a_{2}=5,\quad a_{3}=2,\quad a_{4}=-3,\quad a_{6}=-3,\quad a_{12}=1$ \\
$\Phi_{12}$ & 12 & $a_{1}=1,\quad a_{3}=1,\quad a_{4}=2,\quad a_{6}=2,\quad a_{12}=-2$ \\
\end{tabular}
\end{center}


In particular $a_n=0$ unless $n\mid M$, where
\[
                 M\in\{1,2,3,4,5,6,8,10,12\}.
 \tag{C.MS.3}\hypertarget{tag-00RK}{}\label{cdp:eq:C.MS.3}
\]
On $H^1(F,\Q)$ one also has
\[
                  (T^M-I)^2=0,
                  \qquad\operatorname{rank}(T^M-I)\le2.
 \tag{C.MS.10}\hypertarget{tag-00RI}{}\label{cdp:eq:C.MS.10}
\]
All $a_n$ vanish exactly when $1$ is an eigenvalue
of $T$. The table is the complete range of \eqref{cdp:eq:C.MS.2}
on orientation-preserving quasi-unipotent integral
rank-four matrices. Every actual filling has one
of these signatures. Every row is realized by a
map from a smooth projective threefold to $\PP^1$
whose general fibre is an abelian surface and
whose every fibre has Euler number zero. It
can be chosen to have exactly two singular
fibres, both with reducible simple-normal-crossings
support.

There are two useful consequences for the original
fibre before resolution. Write it as $\sum_i m_iS_i$,
put $m_0=\min_i m_i$, and let $S_i^\circ$ be the
smooth locus of $S_i$ with all its intersections
with the other components removed. Then
\[
 m_0\sum_{m_i=m_0}e(S_i^\circ)
                    =\det(I-T^{m_0})\ge0.        \tag{C.MS.4}\hypertarget{tag-00RL}{}\label{cdp:eq:C.MS.4}
\]
If every component has multiplicity greater than
one, the entire signature is zero. For an
irreducible full fibre $mS$ one has
\[
 \begin{gathered}
 m>1\quad\Longrightarrow\quad e(S_{\rm reg})=0,\\
 e(S_{\rm reg})\in
       \{0,1,2,3,4,5,6,8,9,12,16\},\\
 e(S_{\rm reg})>0\quad\Longrightarrow\quad m=1.
 \end{gathered}                                  \tag{C.MS.5}\hypertarget{tag-00RM}{}\label{cdp:eq:C.MS.5}
\]
These constraints apply in particular to every
reducible or irreducible zero-Euler filling.
\end{proposition}


\begin{proof}

\begin{enumerate}[label={$\langle1\rangle$\arabic*.},leftmargin=*]
\item We specify first the comparison with the analytic
family. The resolution is an isomorphism over
$\Delta^*$, so it identifies the actual nearby
torus and its monodromy before and after resolution.
The proper global construction in
\cite[\S2, pp.~237--243, and \S3]{ACampo1975}
gives the nearby-cycle cohomology sheaves on the
resolved central support and their spectral sequence
to $H^*(F,\Q)$. In a chart
$z_1^{N_1}\cdots z_r^{N_r}=t$, a local nearby
fibre has $\gcd(N_1,\ldots,N_r)$ components,
each with the homotopy type of $(S^1)^{r-1}$.
Monodromy permutes these components cyclically;
after their permutation is killed its action
is a translation on the torus and induces the
identity on cohomology. A common multiple of all
the $N_j$ therefore acts identically on every
stalk of every nearby-cycle cohomology sheaf.
It acts identically on the corresponding
cohomology groups of those sheaves and hence
on the associated graded of the finite
spectral-sequence filtration of $H^*(F,\Q)$.
Its difference from the identity is nilpotent
on that cohomology. This proves quasi-unipotence
using the actual analytic normal-crossings
comparison; no ambient K\"ahler assumption enters.
In degree one this filtration has only the
two possible terms $H^1(E,\mathcal H^0\psi)$
and $H^0(E,\mathcal H^1\psi)$, where
$E=\bigcup_jE_j$ and $\mathcal H^q\psi$
denotes the nearby-cycle cohomology sheaf
already described. Thus for a
common multiplicity multiple $N$ one has
$(T^N-I)^2=0$. The semisimple order $M$
divides $N$, and $U=T^M$ is unipotent.
Writing
$U^{N/M}-I=(U-I)(I+U+\cdots+U^{N/M-1})$,
the second factor is invertible over $\Q$
and commutes with $U$. Hence $(U-I)^2=0$.
Its image is contained in its kernel on
the four-dimensional space, proving
the rank bound in \eqref{cdp:eq:C.MS.10}.

\item The same local model computes the alternating
stalk trace of the $k$th monodromy power. At a
point of $E_j^\circ$ it is $N_j$ if $N_j\mid k$
and zero otherwise. At a point lying on two or
more components it is zero, since the torus
factor has positive dimension and Euler number
zero. Taking the alternating trace through the
constructible spectral sequence, or its finite
cell filtration, gives
\[
 \sum_q(-1)^q\operatorname{Tr}
             (T^k;H^q(F,\Q))
       =\sum_{N_j\mid k}N_j e(E_j^\circ).
 \tag{C.MS.6}\hypertarget{tag-00RN}{}\label{cdp:eq:C.MS.6}
\]
This is the global form of the explicit local
calculation in \cite[Theorem~1, p.~235,
and proof on pp.~243--244]{ACampo1975}.
The cup-product map
$\Lambda^qH^1(F,\Z)\to H^q(F,\Z)$ is an
isomorphism for a torus and commutes with the
actual monodromy. Thus the left side of
\eqref{cdp:eq:C.MS.6} is exactly $\det(I-T^k)$.
The complex orientation of the smooth fibre
is preserved by transport. The map on its
top cohomology is consequently the identity;
under the fourth exterior power this says
$\det T=1$. Integer M\"obius inversion now
gives both equations of \eqref{cdp:eq:C.MS.2}.

\item We compute their complete range. The characteristic
polynomial of a quasi-unipotent integral rank-four
matrix is a product of cyclotomic polynomials
whose degrees sum to four. The possible individual
orders are
$1,2,3,4,5,6,8,10,12$. Since $\det T=1$,
the multiplicity of the eigenvalue $-1$ is even.
There are exactly nineteen resulting products:
the five containing $\Phi_1$ and the fourteen
remaining products displayed in the table.
This enumeration follows directly from
$\deg\Phi_1=\deg\Phi_2=1$,
$\deg\Phi_3=\deg\Phi_4=\deg\Phi_6=2$,
and degree four for $\Phi_5,\Phi_8,\Phi_{10},
\Phi_{12}$. The sum of the degree-one
multiplicities is even, so the $\Phi_1$
multiplicity is even as well.

\item The determinants $\det(I-T^k)$ depend only on
these eigenvalues, including their multiplicities;
Jordan blocks make no change to them. Each listed
product is represented by the block diagonal
matrix of its integral cyclotomic companion
matrices, which has determinant one. Its traces
therefore compute the exact full range on
matrices. Applying \eqref{cdp:eq:C.MS.2} gives the displayed
integers.

\item There is also a direct verification that no
unlisted index can be concealed by cancellation.
The actual exterior-power identity gives
\[
 \begin{aligned}
 \exp\!\left(\sum_{k\ge1}\frac{L_k}{k}t^k\right)
 &=\prod_{q=0}^4
       \det\bigl(I-t\Lambda^qT\bigr)^{(-1)^{q+1}}\\
 &=\prod_{n\ge1}(1-t^n)^{-a_n}.
 \end{aligned}                                  \tag{C.MS.7}\hypertarget{tag-00RO}{}\label{cdp:eq:C.MS.7}
\]
Every eigenvalue of every exterior power has
order dividing $M$. Hence the first product
factors into cyclotomic polynomials $\Phi_d(t)$
with $d\mid M$. The identity
$\Phi_d(t)=\prod_{n\mid d}(1-t^n)^{\mu(d/n)}$
for $d>1$, with the harmless sign at $d=1$,
expresses it using only indices dividing $M$.
Finite products of the factors $1-t^n$ have
unique integer exponent vectors: the coefficient
at the smallest differing degree in their
formal logarithms detects any difference.
Thus $a_n=0$ for every $n\nmid M$, including
all possible cancellations. If $T$ has an
eigenvalue one, every $L_k$ is zero and so is
the signature. Conversely, without that
eigenvalue, the real eigenvalues are $-1$
and the others occur in complex-conjugate
root-of-unity pairs. Their factors in
$\det(I-T)$ are respectively $2$ and
$|1-\lambda|^2$, all positive. Thus $a_1=L_1>0$.
The same factorization proves $L_k\ge0$ for
every $k$.

\item For \eqref{cdp:eq:C.MS.4}, choose the resolution with centres
over the singular locus of the original reduced
support. Its first exceptional multiplicities
are at least $2m_0$: at such a centre either
one support component has order at least two
or at least two components occur. All later
exceptional multiplicities are sums of the
orders of the components containing their
centres. If the generic point of a centre is
contained in an existing exceptional component,
that component already contributes at least
$2m_0$. Otherwise the map is locally an
isomorphism at its generic point and the
original reduced support still has order at
least two there. Thus this case has the same
lower bound, including successive resolutions
over disjoint original singular points.
The open stratum of a strict
transform is exactly the specified
$S_i^\circ$. In \eqref{cdp:eq:C.MS.6} for $k=m_0$, only
strict transforms of multiplicity $m_0$
contribute; this proves \eqref{cdp:eq:C.MS.4}.

\item If $m_0>1$, the resolved divisor has no
multiplicity-one component. Hence $a_1=0$,
so $T$ has eigenvalue one and the entire
signature vanishes. For an irreducible
fibre $mS$, \eqref{cdp:eq:C.MS.4} is
$m e(S_{\rm reg})=L_m$. When $m>1$ it
is zero by the preceding argument. When
$m=1$, its possible values are the
$a_1$ entries of the table, exactly the
set in \eqref{cdp:eq:C.MS.5}.

\item We now realize every row by an actual compact
family. The fifteen products of two degree-two
blocks are obtained from products of elliptic
curves. The five blocks
\[
                 \Phi_1^2,\quad\Phi_2^2,\quad
                 \Phi_3,\quad\Phi_4,\quad\Phi_6
\]
come respectively from the identity, negation,
and multiplication by a primitive third, fourth
or sixth root of unity on an elliptic curve
with the corresponding triangular or square
lattice. Taking the product automorphism gives
exactly the indicated product on integral
first cohomology.

\item For $\Phi_5$, let $A$ be the Jacobian of
the smooth genus-two curve
\[
                           C:\ y^2=x^5-1.
\]
Its double cover of $\PP^1$ has six simple
branch points, including infinity, so the
genus assertion follows from Riemann--Hurwitz.
The automorphism $(x,y)\mapsto(\zeta_5x,y)$
fixes infinity and acts on the basis
$dx/y,x\,dx/y$ of holomorphic differentials
with eigenvalues $\zeta_5,\zeta_5^2$.
The Abel--Jacobi map based at infinity
intertwines this automorphism with its induced
automorphism of $A$ and induces the actual
isomorphism on integral $H^1$. Including the
conjugate eigenvalues gives $\Phi_5$.
Composing the Jacobian automorphism with
negation gives order ten and characteristic
polynomial $\Phi_{10}$. Here negation is also
the map induced by the curve's hyperelliptic
involution: the divisor of $x-x(p)$ gives
$[p-\infty]+[\iota(p)-\infty]=0$ on the
Jacobian, verifying that comparison explicitly.

\item For the remaining two rows take
$A=E_i\times E_i$, where
$E_i=\C/(\Z+i\Z)$. The Gaussian-integer
matrices
\[
 G_8=\begin{pmatrix}0&i\\1&0\end{pmatrix},
 \qquad
 G_{12}=\begin{pmatrix}0&1\\1&i\end{pmatrix}
 \tag{C.MS.8}\hypertarget{tag-00RP}{}\label{cdp:eq:C.MS.8}
\]
define actual automorphisms of this torus.
One has $G_8^2=iI$, and its holomorphic
eigenvalues are $\zeta_8,\zeta_8^5$.
The characteristic polynomial of $G_{12}$
over $\C$ is $x^2-ix-1$, with roots
$(i+\sqrt3)/2=\zeta_{12}$ and
$(i-\sqrt3)/2=\zeta_{12}^5$.
Together with the conjugate eigenvalues,
these give the integral characteristic
polynomials $\Phi_8$ and $\Phi_{12}$ and
the exact orders eight and twelve.
Both matrices preserve the displayed lattice;
their determinants are Gaussian units.
This accounts for all nineteen rows with
finite-order automorphisms $g$ of projective
two-tori $A$.

\item Let $M$ be the order of the selected $g$.
Take the projective finite quotient
\[
 X=(A\times\PP^1)/
       \langle(a,z)\mapsto(g(a),\zeta_Mz)\rangle,
 \qquad f([a,z])=z^M.                           \tag{C.MS.9}\hypertarget{tag-00RQ}{}\label{cdp:eq:C.MS.9}
\]
The action is free over $z\ne0,\infty$.
Over a nonzero finite base value the $M$
slices are permuted transitively, and their
quotient is the original connected torus $A$.
The actual positive meridian lifts from $z$
to $\zeta_Mz$, so under this identification
the monodromy is $g^{-1}$. Inverting all
the eigenvalues preserves each cyclotomic
polynomial and hence the displayed signature.
Resolve the projective quotient and then its
two central supports by projective blowups
over $0$ and $\infty$, using
\cite[Theorem~1.6]{BierstoneMilman1997}.
The resulting threefold is smooth and projective,
with SNC supports, and the punctured family is
unchanged. Its Stein factor has generic
degree one over the normal curve $\PP^1$,
so every fibre is connected. When $M=1$
the construction is simply the product,
which already has zero-Euler fibres.

\item For the remaining cases, each of the two
possibly singular supports has even Euler
number by \ref{cdp:st:C.2.UN.1}. Apply the exact vertical
blowups of \ref{cdp:st:C.2.UN.4} to make each number
zero. A point blowup changes it by $+2$;
a point blowup followed by the blowup of
a smooth plane quartic in its exceptional
$\PP^2$ changes it by $-2$. These
operations are projective and are performed
only over the selected value, so they
preserve the actual punctured family and
the signature \eqref{cdp:eq:C.MS.2}.

\item They preserve SNC supports as well. Choose
the point on a single component away from
the others. Its exceptional plane meets
the old strict transform along a line,
and choose the quartic transverse to that
line. Near any of their four intersection
points write the map as $t=(ru)^m$,
with exceptional component $r=0$, old
component $u=0$, and quartic centre
$(r,v)=(0,0)$. Its blowup charts
$v=rw$ and $r=vw$ give respectively
$t=(ru)^m$ and $t=(vwu)^m$.
Thus every resulting component is still
a coordinate hyperplane; away from these
four points the same assertion uses only
two coordinates. These explicit charts
verify the claimed SNC zero-Euler
realization of every row.
After the Euler number is zero, one additional
$+2$ operation and one $-2$ operation over the
same value leave it zero and introduce
exceptional components. Performing this at
both $0$ and $\infty$ ensures precisely two
reducible singular fibres, including the
identity-monodromy row. All other fibres
remain the original smooth tori.

\item The exact computation is saved in
\path{code/D_monodromy_signatures.py}, its
output, and
\path{records/D_monodromy_signatures.csv}.
It verifies every determinant and exterior-power
zeta identity over the integers or $\Q(t)$,
and the M\"obius inversions use exact integers.
The matrix-to-signature map itself forgets
information: for example $I_4$ and
$I_4+E_{12}$ have the same characteristic
polynomial and zero signature, while their
minimal polynomials differ. The table therefore
records precisely the stated invariant and
does not supply an analytic isomorphism
classification of fillings.

\end{enumerate}
\end{proof}






\begin{proposition}[K3 components are elliptic in positive algebraic dimension; algebraic dimension zero has an exact occurrence criterion]\hypertarget{tag-01D2}{}\label{cdp:st:C.2.UN.24}
\StatementRecord{cdp:st:C.2.UN.24}{proved}{01D2}{\SeriesRef{cdp:st:C.2.UN.16}, \SeriesRef{cdp:eq:C.K3.2}, \SeriesRef{cdp:st:C.2.UN.20}, \SeriesRef{cdp:st:C.2.UN.22}, \SeriesRef{cdp:eq:C.K3.3}, \SeriesRef{cdp:st:C.3.1}, \SeriesRef{cdp:st:C.2.UN.23}, \SeriesRef{cdp:eq:C.K3.7}, \SeriesRef{cdp:eq:C.K3.4}, \SeriesRef{cdp:eq:C.K3.5}, \SeriesRef{cdp:eq:C.K3.6}, \SeriesRef{cdp:st:C.2.UN.1}, \SeriesRef{cdp:st:C.2.UN.4}; \cite[Chapter~1, Theorem~3.1, p.~10]{Peters1990}; \cite[Theorem~II, Corollary~II]{Lagrange1770}; \cite[\S3, Corollary~3, with the reflection
argument of \S6]{PiatetskiShapiroShafarevich1971}; \cite[\S3.1, p.~7659; \S4.1, Theorem~27, p.~7663]{Roulleau2019GeneralizedKummer}}

Let $f:U\to\Delta$ be a proper holomorphic map from a
smooth complex threefold, with connected two-torus fibres
over $\Delta^*$. Let $S_i$ be an irreducible component of
its central support, and suppose that a resolution of
$S_i^\nu$ has a K3 minimal model $M$.
If $a(M)>0$, there is an actual elliptic fibration
$M\to\mathbb P^1$.

Among K3 surfaces of algebraic dimension zero, the minimal
birational classes which occur in this way are exactly the
minimal resolutions of cyclic quotients of complex two-tori
by nontrivial symplectic automorphisms without eigenvalue
one on the tangent space. After translating the origin, the generator
is linear of order $n=2,3,4$, or $6$. The quotient singularities
and exceptional configurations are exactly
\[
\begin{array}{c|c|c|c}
n&\text{quotient singularities}&
 \text{number of exceptional curves}&e(A/\langle g\rangle)\\ \hline
2&16A_1&16&8\\
3&9A_2&18&6\\
4&4A_3+6A_1&18&6\\
6&A_5+4A_2+5A_1&18&6.
\end{array}                                                    \tag{C.K3.1}\hypertarget{tag-00OY}{}\label{cdp:eq:C.K3.1}
\]
Here the cyclic quotient in the criterion is required to
have algebraic dimension zero; the table applies to the
cyclic K3 quotient models in every algebraic dimension.

Every cyclic quotient K3 minimal class just described is
realized by a compact smooth threefold over $\mathbb P^1$
with connected torus general fibre and with both special
reduced supports of Euler characteristic zero. The central
support has a component whose normalization resolution has
exactly that K3 minimal model. The total space and construction
are projective when the given torus is projective.
\end{proposition}



\begin{proof}

\begin{enumerate}[label={$\langle1\rangle$\arabic*.},leftmargin=*]
\item Use the actual normalized root base change and equivariant
resolution of \ref{cdp:st:C.2.UN.16}. For the chosen component, let $E$
be a smooth strict covering component and $G$ its cyclic
stabilizer. The comparison is
\[
 E\longrightarrow S_i^\nu,
 \qquad E/G\dashrightarrow S_i^\nu\dashrightarrow M,
                                                               \tag{C.K3.2}\hypertarget{tag-00P0}{}\label{cdp:eq:C.K3.2}
\]
where the first map is generically finite, the quotient
comparison is proper birational, and the remaining arrows
are resolved by their proper graphs. In particular, a
nonzero holomorphic two-form on $M$ pulls back, by its
actual Jacobian map, to a nonzero holomorphic two-form on
a resolution of the graph in \eqref{cdp:eq:C.K3.2}. Point-blowup
invariance then gives
\[
                  p_g(E)\ge1,\qquad \kappa(E)\ge0.
                                                               \tag{C.K3.3}\hypertarget{tag-00P1}{}\label{cdp:eq:C.K3.3}
\]
Also $a(E)=a(M)$. To check this equality through the
specified map, a meromorphic function on $E$ has elementary
symmetric functions along the finite generic sheets.
They descend meromorphically to $M$, so that the original
function satisfies a monic polynomial over $\mathcal M(M)$.
Thus the two meromorphic function fields have the same
transcendence degree.

\item The two-dimensional space of fields on $E$ in \ref{cdp:st:C.2.UN.16}
has generic rank one or two. In rank one, \ref{cdp:st:C.2.UN.20} gives
an actual rational or elliptic fibration on a birational
model of $M$. The rational alternative would make all
positive pluricanonical sections vanish, by restricting
them to its smooth general fibre. The pulled-back K3
form excludes it. Hence the quotient fibration is elliptic.

\item In rank two, \ref{cdp:st:C.2.UN.22} gives $\chi(\mathcal O_E)=0$ and
its precise torus or elliptic-fibration alternatives.
First suppose $a(M)=a(E)=1$. In the elliptic alternative,
the connected-fibre map $q:E'\to B$ has
\[
                     \mathcal M(E')=q^*\mathcal M(B).
                                                               \tag{C.K3.4}\hypertarget{tag-00P2}{}\label{cdp:eq:C.K3.4}
\]
Indeed every meromorphic function is algebraic over any
one nonconstant function on $B$ and hence constant on
the connected general fibres; the proper graph descends
it, as in (C.RF.9). The equality makes this field
intrinsic and therefore $G$-invariant. Resolve the graph
of $q$ equivariantly and take its Stein factor. The group
acts on that factor, and the actual finite quotient is a
map to $B/G$ whose general fibre is $F/K$, with
$K=\ker(G\to\operatorname{Aut}B)$. This is the same
quotient and Stein comparison proved in \ref{cdp:st:C.2.UN.20}.
Its general fibre is connected and rational or elliptic;
the K3 form again excludes the rational case. In the
torus alternative the minimal model of $E$ is a torus
and the cyclic quotient argument below applies.

\item Now suppose $a(M)=a(E)=2$. We work on the smooth
projective minimal model $E_0$, to which the finite
action descends. In rank two, \eqref{cdp:eq:C.K3.3} and
$\chi(\mathcal O_{E_0})=0$, together with the compact
surface classification already used in \SeriesRef{cdp:st:C.3.1}
\cite[Chapter~1, Theorem~3.1, p.~10]{Peters1990}, leave the following
possibilities. The precise models in \ref{cdp:st:C.2.UN.22}
already give $\kappa(E_0)\le1$: on an elliptic general fibre
$F$ of its stated model, adjunction gives
$K|_F=K_F=\OO_F$, so every pluricanonical map is constant
on $F$ and has image of dimension at most one; a torus
has trivial canonical line. If $\kappa(E_0)=0$, it is an Abelian
surface: K3 and Enriques surfaces have coherent Euler
characteristic $2$ and $1$, and bielliptic surfaces have
$p_g=0$. If $\kappa(E_0)=1$, its elliptic Iitaka
fibration is intrinsic. More explicitly, $G$ acts on
each space of pluricanonical sections, so a sufficiently
large pluricanonical map, its equivariant graph resolution,
and its Stein factor are all $G$-equivariant. The
connected general fibres are elliptic. Taking the
quotient gives the actual fibration with general fibre
$F/K$ just described. It is elliptic because its total
quotient has the K3 minimal model. Thus the only
remaining positive-dimensional case is a cyclic quotient
of a two-torus.

\item For use in all these branches, an elliptic fibration on
a smooth birational model of $M$ descends to $M$ itself.
Take an actual common resolution
\[
                  M\xleftarrow{\rho}X\xrightarrow{q}B
                                                               \tag{C.K3.5}\hypertarget{tag-00P3}{}\label{cdp:eq:C.K3.5}
\]
with connected general elliptic fibre $F$. Since $K_M=0$,
the point-blowup formula gives
$K_X=\sum_j a_jE_j$ with $a_j>0$ on every
$\rho$-exceptional prime curve. Adjunction gives
$K_X\cdot F=0$. Each $E_j\cdot F\ge0$, so all these
intersections are zero and every exceptional curve is
vertical for $q$. Every connected fibre of $\rho$ is
therefore sent to a point. The proper graph, together
with $\rho_*\mathcal O_X=\mathcal O_M$, descends $q$
to a holomorphic map $M\to B$ with connected fibres.
The Leray injection
$H^1(B,\mathcal O_B)\hookrightarrow H^1(M,\mathcal O_M)=0$
gives $B\simeq\mathbb P^1$.

\item We next compute the cyclic quotient branch. By
\ref{cdp:st:C.2.UN.23}, a torus minimal model $A$ of $E$ has an
actual holomorphic cyclic action whose quotient is
birational to $M$. The K3 alternative there says that
its generator $g(z)=Lz+b$ has no eigenvalue one and
$\det L=1$. Since $I-L:A\to A$ is an isogeny, the
equation $(I-L)p=b$ has a solution. Translation by
$p$ conjugates $g$ to the linear automorphism $L$.
Write its eigenvalues as $\lambda,\lambda^{-1}$,
with $\lambda$ a primitive $n$th root. The integral
action on $H_1(A,\mathbb Z)$ has characteristic polynomial
\[
 (x^2-sx+1)^2
   =x^4-2sx^3+(s^2+2)x^2-2sx+1,
 \qquad s=\lambda+\lambda^{-1}\in[-2,2].
                                                               \tag{C.K3.6}\hypertarget{tag-00P4}{}\label{cdp:eq:C.K3.6}
\]
Its integral coefficients imply $2s\in\mathbb Z$
and $s^2\in\mathbb Z$. Writing $s=k/2$ shows that
$k$ is even, hence $s\in\mathbb Z$. The value $s=2$
is excluded by the absence of eigenvalue one.
Thus $s=-2,-1,0,1$, giving precisely
$n=2,3,4,6$.

\item For every nonidentity power, its fixed points are the
kernel of an isogeny and their exact number is
\[
 \#\operatorname{Fix}(g^k)
 =\det_{\mathbb R}(I-L^k)
 =|\det_{\mathbb C}(I-L^k)|^2
 =\bigl(2-2\cos(2\pi k/n)\bigr)^2.
                                                               \tag{C.K3.7}\hypertarget{tag-00P5}{}\label{cdp:eq:C.K3.7}
\]
At a point with stabilizer of order $d$, the local
action is $(u,v)\mapsto(\zeta_du,\zeta_d^{-1}v)$.
Its invariant coordinates $x=u^d,y=v^d,z=uv$ give
exactly $xy=z^d$. The height-one toric resolution
in (C.CQ.11) gives its chain of $d-1$ rational
$(-2)$-curves. Formula \eqref{cdp:eq:C.K3.7} now gives the table
by subtracting points with larger stabilizers and
dividing by the remaining orbit sizes. For $n=4$
there are $4$ full-stabilizer points and
$(16-4)/2=6$ order-two orbits. For $n=6$ the
counts are $1$, $(9-1)/2=4$, and
$(16-1)/3=5$ for stabilizer orders $6,3,2$.
The order-two and order-three cases have $16$ and
$9$ fixed points respectively. The finite quotient
Euler formula gives the last column; adding the
exceptional ranks gives $24$ in every row.
These cyclic configurations agree with the classical generalized-Kummer
list in \cite[\S3.1, p.~7659; \S4.1, Theorem~27, p.~7663]{Roulleau2019GeneralizedKummer},
which records the preceding cyclic-quotient work.
The exact integer matrices, every power's fixed-point count,
exact-stabilizer subtraction, quotient Euler average and resolution
Euler correction are computed in
\path{code/current_audits/audit_sign/code/UN24_UN39_exact_certificate_08.py};
the full four-row output and the formal identities
\eqref{cdp:eq:C.K3.8}--\eqref{cdp:eq:C.K3.9} are recorded in
\path{code/current_audits/audit_sign/UN24_UN39_exact_certificate_08.json}.

\item If this quotient K3 $M$ is projective, the table gives
an explicit elliptic pencil criterion. Choose one
exceptional $(-2)$-curve from each of four different
singular points, and call them $C_1,\ldots,C_4$.
They are pairwise disjoint. For an ample integral
class $H$ on $M$, put
\[
 H'=H+\sum_{i=1}^4\frac{H\cdot C_i}{2}C_i,
 \quad H'\cdot C_i=0,
 \quad H'^2=H^2+\frac12\sum_{i=1}^4(H\cdot C_i)^2>0.
                                                               \tag{C.K3.8}\hypertarget{tag-00P6}{}\label{cdp:eq:C.K3.8}
\]
Lagrange's four-square theorem
\cite[Theorem~II, Corollary~II]{Lagrange1770}
expresses the positive rational number $H'^2/2$ as
$\sum_i x_i^2$ with $x_i\in\mathbb Q$: for a
positive fraction $a/b$, apply the integral theorem
to $ab$ and divide the four summands by $b$.
Consequently
\[
 D=H'-\sum_{i=1}^4x_iC_i,
 \qquad D^2=0,
 \qquad D\cdot H'=H'^2>0.
                                                               \tag{C.K3.9}\hypertarget{tag-00P7}{}\label{cdp:eq:C.K3.9}
\]
Clear denominators to obtain a nonzero isotropic
integral divisor class. The original elliptic-pencil
criterion \cite[\S3, Corollary~3, with the reflection
argument of \S6]{PiatetskiShapiroShafarevich1971}
then produces an actual base-point-free elliptic
pencil on $M$. This proves the remaining projective
branch. If the quotient has algebraic dimension
one, then $a(A)=1$ as well. Choose a nonconstant
meromorphic function on $A$, and let $D$ be a
general member of its pencil after removing fixed
components. Its self-intersection is nonnegative,
by intersecting two distinct members. In fact
$D^2=0$. To verify this, average the positive
integration current $[D]$ under translations of
the compact torus. The result is the constant
semipositive $(1,1)$-form representing $c_1(D)$.
If $D^2>0$, this form is positive definite, so
the integral class is a polarization and the
torus is projective, contrary to $a(A)=1$.
The equality $D^2=0$ shows that two distinct
members have no common point. The pencil thus
defines a holomorphic map, and its Stein factor
$q:A\to B$ has connected smooth general fibre
with trivial normal line. Adjunction with
$K_A=0$ gives genus one. The field identity
\eqref{cdp:eq:C.K3.4} makes this fibration $G$-equivariant;
its quotient and \eqref{cdp:eq:C.K3.5} give the actual
elliptic pencil on $M$. This completes the
assertion for $a(M)>0$.

\item If $a(M)=0$, equality of algebraic dimensions in
\eqref{cdp:eq:C.K3.2} gives $a(E)=0$. The final assertion of
\ref{cdp:st:C.2.UN.16} identifies its minimal model as a torus
or a Hopf surface. The nonzero canonical form
in \eqref{cdp:eq:C.K3.3} excludes Hopf. Thus \eqref{cdp:eq:C.K3.2} and
\ref{cdp:st:C.2.UN.23} give exactly the cyclic symplectic torus
quotient of \eqref{cdp:eq:C.K3.6}--\eqref{cdp:eq:C.K3.7}. This proves necessity
of the stated algebraic-dimension-zero criterion.

\item Finally we prove the occurrence assertion with its
specified compact comparison. Start with any torus
$A$ and one of the actual cyclic symplectic actions
above, of order $n$. Form the finite analytic quotient
\[
 W=(A\times\mathbb P^1_s)/
       \langle(a,s)\mapsto(g(a),\zeta_ns)\rangle,
 \qquad W\longrightarrow\mathbb P^1_t,
 \qquad t=s^n.
                                                               \tag{C.K3.10}\hypertarget{tag-00OZ}{}\label{cdp:eq:C.K3.10}
\]
Over $t\ne0,\infty$, the $n$ sheets of the base
change are permuted transitively, and their quotient
fibre is exactly $A$. The reduced fibres over zero
and infinity are exactly $A/\langle g\rangle$.
Choose an analytic resolution of the total quotient
and its two central supports, isomorphic over the
punctured base and at a generic point of each of
these two original components. The strict transform
of the latter maps properly birationally to the
specified torus quotient. Its normalization resolution
therefore has the K3 minimal model already proved
in \ref{cdp:st:C.2.UN.23}. The resulting smooth compact total
space has connected fibres: its Stein factor over
$\mathbb P^1$ has generic degree one and hence is
the identity finite map to the normal base.

\item Each of its two reduced special supports has even
Euler characteristic by \ref{cdp:st:C.2.UN.1}. Apply the actual
vertical point and quartic-curve blowup operations
of \ref{cdp:st:C.2.UN.4} separately over zero and infinity.
They change the reduced Euler number by $+2$ or
$-2$ until it is zero, preserve the punctured torus
family, and leave the original strict component
birational to the same torus quotient. Resolution
and these blowups are projective morphisms when
$A$ is projective. This supplies all asserted
compact realizations, including every class of
algebraic dimension zero in the criterion.

\item The algebraic-dimension-zero occurrence criterion
is thus an equivalence. In positive algebraic
dimension the theorem gives the elliptic-fibration
restriction and realizes its cyclic quotient
subfamily by the explicit compact construction.

\end{enumerate}
\end{proof}






\begin{proposition}[K3 components in arbitrary torus fillings are exactly the cyclic generalized Kummer surfaces]\hypertarget{tag-01DB}{}\label{cdp:st:C.2.UN.32}
\StatementRecord{cdp:st:C.2.UN.32}{proved}{01DB}{\SeriesRef{cdp:st:C.2.UN.24}, \SeriesRef{cdp:st:C.2.UN.16}, \SeriesRef{cdp:st:C.2.UN.19}, \SeriesRef{cdp:st:C.2.UN.28}, \SeriesRef{cdp:st:C.2.UN.23}, \SeriesRef{cdp:st:C.2.UN.22}, \SeriesRef{cdp:st:C.2.UN.29}, \SeriesRef{cdp:st:C.2.UN.20}, \SeriesRef{cdp:eq:C.KX.4}, \SeriesRef{cdp:eq:C.KX.8}, \SeriesRef{cdp:eq:C.KX.10}, \SeriesRef{cdp:eq:C.KX.7}, \SeriesRef{cdp:eq:C.KX.13}, \SeriesRef{cdp:eq:C.KX.1}}

The K3 minimal birational classes occurring as
normalization components of proper two-torus
degenerations are exactly the minimal resolutions
of cyclic quotients of complex two-tori by
nontrivial symplectic automorphisms.
The generator has no eigenvalue one on the tangent space
and, after an actual translation of the origin,
is linear of order $2,3,4$, or $6$.
The exceptional configurations are those of
(C.K3.1). This criterion holds in every
algebraic dimension.

Every such class occurs in a smooth compact
filling with both special reduced supports
of Euler number zero, by the explicit quotient
and balancing construction of \ref{cdp:st:C.2.UN.24}.
In particular the algebraic-dimension-zero
equivalence of that statement extends to
all K3 components.
\end{proposition}


\begin{proof}

\begin{enumerate}[label={$\langle1\rangle$\arabic*.},leftmargin=*]
\item Let $M$ be the K3 minimal model of an original
component. Take the actual multiplicity-one
component cover $E$ and its cyclic stabilizer
$G$ from \ref{cdp:st:C.2.UN.16}. The proper graph comparisons
of \ref{cdp:st:C.2.UN.19} give
\[
                         E/G\dashrightarrow M
                                                               \tag{C.KX.1}\hypertarget{tag-00PW}{}\label{cdp:eq:C.KX.1}
\]
as a proper birational comparison after
resolving the graph. Pullback of the K3
two-form is injective, whereas the
compact-core bound \ref{cdp:st:C.2.UN.28} gives $p_g(E)\le1$.
Thus
\[
                       p_g(E)=1,\qquad\kappa(E)\ge0.
                                                               \tag{C.KX.2}\hypertarget{tag-00Q2}{}\label{cdp:eq:C.KX.2}
\]
Pass to its unique smooth minimal model
$E_0$; the finite action descends through
the actual point blowdowns.
If $E_0$ is a torus, \ref{cdp:st:C.2.UN.23} already gives
the required cyclic K3 quotient.

\item In the other cases, \ref{cdp:st:C.2.UN.22} and \ref{cdp:st:C.2.UN.29} give
an actual relatively minimal elliptic map
\[
                   q:E_0\longrightarrow B
                                                               \tag{C.KX.3}\hypertarget{tag-00Q3}{}\label{cdp:eq:C.KX.3}
\]
with identity integral monodromy and
only smooth multiple fibres. Its
relative Jacobian is one fixed elliptic
curve $F$. The local invariant
direct-image calculation (C.RG.13)
applies and gives
$p_g(E_0)=g(B)$, so $B$ is elliptic.
The map can be taken $G$-equivariant.
In the rank-one branch this is \ref{cdp:st:C.2.UN.20}.
In algebraic dimension one its
connected-fibre function field is
intrinsic, by (C.RG.13) and
(C.K3.4). In algebraic dimension
two a non-torus minimal model
with $p_g=1$ and $\chi(\OO)=0$
has Kodaira dimension one, so
its elliptic map is its intrinsic
Iitaka fibration, as in \ref{cdp:st:C.2.UN.24}.
Algebraic dimension zero was
already the torus branch of \ref{cdp:st:C.2.UN.16}.
These are the same equivariant
graph and Stein comparisons
used there.

\item We first describe every multiple
fibre with its actual cyclic
uniformization. For a fibre $mC$,
the normalized root change
$t=u^m$ is \'etale over the
original smooth surface and
has connected smooth elliptic
central fibre. The identity
monodromy and constant period
ratio make this root family
a torsor under the fixed $F$
over a disc. Its torsor is
trivial: the sheaf sequence
$0\to\Lambda_F\to\OO\to F\to0$
on a disc has $H^1(\OO)=0$
and $H^2(\Lambda_F)=0$.
Thus it is an actual product.
The generator of its root
deck group has translation
linear part and can be
written
$(u,z)\mapsto(\zeta_m u,z+\tau+a(u))$,
where $\tau=a_{\rm total}(0)$
and $a(0)=0$.
Lift $a$ holomorphically to
$\C$. The finite-order
identity gives
$\sum_{j=0}^{m-1}a(\zeta_m^ju)=0$.
Replacing $z$ by $z+b(u)$,
with
$b(u)=-m^{-1}\sum_{j=0}^{m-1}j\,a(\zeta_m^ju)$,
removes $a(u)$.
The local model is therefore
\[
 (F\times\Delta_u)/\langle\delta\rangle,
 \qquad
 \delta(u,z)=(\zeta_m u,z+\tau),\quad
 \operatorname{ord}(\tau)=m,\quad t=u^m .
                                                               \tag{C.KX.4}\hypertarget{tag-00Q4}{}\label{cdp:eq:C.KX.4}
\]
The exact order of $\tau$
follows from freeness of
the \'etale root covering
on its central fibre.

\item Let $H$ be the kernel of the
action of $G$ on $B$.
Every element of $H$ acts
on a general elliptic fibre
by a translation. Indeed
a nonidentity linear part
would make the general
finite quotient fibre
rational. The actual
quotient and Stein map
would then give a rational
fibration on a birational
model of $M$, contradicting
its nonzero canonical form
by restriction to a general
rational fibre.
The translation subgroup
on the fixed $F$ is
constant: its points have
bounded finite order and
belong to a discrete
finite torsion set.

\item The quotient
\[
              E_1=E_0/H,\qquad Q=G/H
                                                               \tag{C.KX.5}\hypertarget{tag-00Q5}{}\label{cdp:eq:C.KX.5}
\]
is smooth. To check this
also at multiple fibres,
use \eqref{cdp:eq:C.KX.4}. A fibrewise
finite automorphism has
linear part one on the
reduced elliptic fibre,
so its restriction is
a translation. If it
fixes one point, it
fixes that curve pointwise.
Finite holomorphic
linearization then gives
$(x,z)\mapsto(\zeta_e x,z)$;
its quotient coordinates
are $(x^e,z)$. Every
point stabilizer is
cyclic, so these charts
prove smoothness of the
full quotient. They also
show exactly how the
fibre multiplicity can
decrease: $t=x^m$ becomes
$t=(x^e)^{m/e}$, with
$e\mid m$. The reduced
support remains a smooth
elliptic curve. The
general fibre is the
fixed isogenous elliptic
curve obtained by the
constant translation
subgroup. Its integral
linear monodromy remains
the identity.

\item Thus $q_1:E_1\to B$ is
again a smooth-multiple
elliptic fibration with
constant Jacobian.
The same direct-image
calculation gives
\[
 p_g(E_1)=1,\quad q(E_1)=2,\quad
                 \chi(\OO_{E_1})=0 .
                                                               \tag{C.KX.6}\hypertarget{tag-00Q6}{}\label{cdp:eq:C.KX.6}
\]
The actual quotient
$E_1/Q=E_0/G$ remains
birational to $M$.
The group $Q$ acts
faithfully on the base.
Its quotient fibration
has connected elliptic
general fibre. The
K3 descent argument
(C.K3.5) makes its
base $B/Q\simeq\PP^1$.
Consequently $Q$ is
nontrivial and its
action on $B$ is not
a group of translations.
After translating the
origin on $B$, its
generator is a linear
elliptic automorphism.
The derivative character
$\lambda:Q\to\C^*$ is
faithful.

\item We record the exact
canonical character.
Choose a nonzero
one-form $\omega_B$
on the elliptic base
and a time form $dz$
for the fixed Jacobian.
Their wedge glues
over the translation
trivializations. In
\eqref{cdp:eq:C.KX.4} it is
\[
 \Omega=\omega_B\wedge dz
        =m\,c(u^m)u^{m-1}\,du\wedge dz ,
                  \qquad c(0)\ne0 .
                                                               \tag{C.KX.7}\hypertarget{tag-00Q7}{}\label{cdp:eq:C.KX.7}
\]
It is invariant under
$\delta$ and extends
holomorphically. Its
divisor is exactly
$\sum_b(m_b-1)C_b$.
The group acts on
the fixed elliptic
Jacobian with a
constant linear
character $\beta$;
translation terms
add only base
differentials to $dz$.
Hence
$g^*\Omega=\lambda(g)\beta(g)\Omega$.
The nonzero K3 form
pulls back through
the actual graph to
the unique canonical
form on $E_1$.
It is invariant, so
\[
                   \lambda(g)\beta(g)=1
                       \quad(g\in Q).                \tag{C.KX.8}\hypertarget{tag-00Q8}{}\label{cdp:eq:C.KX.8}
\]

\item Here is the local
multiple-fibre
obstruction, with
all covering maps
retained. Let a
base point have
stabilizer generated
by $g$ of order
$s>1$, and suppose
its fibre has
multiplicity $m$.
In \eqref{cdp:eq:C.KX.4} choose
a lift of $g$:
\[
 \widetilde g(u,z)
       =(\rho u,\beta z+a(u)),\qquad
       \rho^m=\lambda .
                                                               \tag{C.KX.9}\hypertarget{tag-00Q9}{}\label{cdp:eq:C.KX.9}
\]
Such a lift is the
functorial lift of
the automorphism
to the normalized
root change, after
choosing an $m$th
root of $\lambda$.
Conjugation preserves
the root deck group.
Its action on $u$
shows that
$\widetilde g\delta
 \widetilde g^{-1}=\delta$.
Restricting this
identity to $u=0$
gives
\[
                         \beta\tau=\tau .
                                                               \tag{C.KX.10}\hypertarget{tag-00PX}{}\label{cdp:eq:C.KX.10}
\]
By \eqref{cdp:eq:C.KX.8}, $\beta$
is nonidentity and
has order $s$.
The sum of its
endomorphisms
$1+\beta+\cdots+\beta^{s-1}$
is zero. Applying
it to $\tau$ in
\eqref{cdp:eq:C.KX.10} gives
$s\tau=0$.
Since $\tau$ has
exact order $m$,
we obtain $m\mid s$.

\item On the reduced
elliptic fibre
$F/\langle\tau\rangle$,
the automorphism
$g$ has nonidentity
linear part $\beta$.
It has a fixed
point, because
$1-\beta$ is an
isogeny of that
elliptic curve.
Choose a lift of
this point and
multiply the
chosen $\widetilde g$
by a root deck
transformation so
that it fixes
that lift.
Then $\widetilde g^s$
is a root deck
transformation
fixing that point.
The deck action
is free, so
$\widetilde g^s=1$.
In particular
$\rho^s=1$, while
$\rho^m=\lambda$
has exact order
$s$. This implies
$\gcd(m,s)=1$.
Together with
$m\mid s$, it
proves
\[
          \text{a branched base value has }m=1 .
                                                               \tag{C.KX.11}\hypertarget{tag-00PY}{}\label{cdp:eq:C.KX.11}
\]

\item At a base point
with trivial
stabilizer, the
quotient is locally
an isomorphism
on a neighbourhood
of one fibre in
the orbit. Its
smooth elliptic
support therefore
retains multiplicity
$m$ and the zero
order $m-1$ of
\eqref{cdp:eq:C.KX.7}.
On a common smooth
resolution of
$E_1/Q$ and $M$,
the strict transform
of this curve is
elliptic and cannot
be contracted by
a birational map
of smooth surfaces:
all curves
exceptional for
successive point
blowdowns are
rational.
At its generic
point the pullback
of the nowhere-zero
K3 form must
therefore have
zero order zero.
Its equality with
the descended form
\eqref{cdp:eq:C.KX.7} forces
$m-1=0$.
This treats every
free base orbit.
We have proved
that $q_1$ has
no multiple fibre
at all.

\item We finish by an
explicit Chern-class
comparison, without
assuming that a
finite image or
preimage is
K\"ahler.
The smooth map
$q_1:E_1\to B$
is now a principal
bundle under its
fixed elliptic
curve $F_1$.
Its Chern vector is
\[
        c\in H^2(B,\Lambda_{F_1})
                  \simeq\Lambda_{F_1}.
                                                               \tag{C.KX.12}\hypertarget{tag-00PZ}{}\label{cdp:eq:C.KX.12}
\]
A bundle
automorphism
covering $g_B$
and having
fibre linear
part $\beta$
gives
$g_B^*c=\beta_*c$.
For the precise
map, lift the
translation
transition
functions on
a base cover
to $\C$.
Their triple
differences are
the lattice
cocycle for $c$.
Conjugating the
trivializations
multiplies that
cocycle by $\beta$,
up to a lattice
coboundary; this
is the displayed
equality.
The holomorphic
automorphism of
the elliptic
base has degree
one, so it acts
as the identity
on $H^2(B,\Z)$.
Thus $c=\beta c$.
For a generator,
$\beta\ne1$, and
$\beta-I$ is
invertible on
$\Lambda_{F_1}
\otimes\R$.
The lattice is
torsion-free, so
$c=0$.

\item The sheaf sequence
\[
 0\longrightarrow\Lambda_{F_1}
 \longrightarrow\OO_B
 \longrightarrow\mathcal F_1
 \longrightarrow0
                                                               \tag{C.KX.13}\hypertarget{tag-00Q0}{}\label{cdp:eq:C.KX.13}
\]
now lifts the
principal-bundle
class to
$H^1(B,\OO_B)$.
Write $B=\C/\Gamma$.
Every class in
that one-dimensional
complex space
is represented
by a constant
$(0,1)$-form
$a\,d\overline z$.
On the universal
cover its
primitive is
$a\overline z$,
and its changes
under $\gamma\in\Gamma$
are the constants
$a\overline\gamma$.
Consequently its
additive torsor,
and its image
under \eqref{cdp:eq:C.KX.13},
have the exact
quotient presentation
\[
 E_1\simeq
 \C^2\big/
 \left\langle
 (z,w)\mapsto
       (z+\gamma,w+a\overline\gamma),\
 (z,w)\mapsto(z,w+\lambda)
 \right\rangle_{
       \gamma\in\Gamma,\,
       \lambda\in\Lambda_{F_1}} .
                                                               \tag{C.KX.14}\hypertarget{tag-00Q1}{}\label{cdp:eq:C.KX.14}
\]
Equality of the
Cech classes gives
this actual
isomorphism:
the difference
of two transition
systems is a
holomorphic
zero-cochain,
and translating
the fibre
coordinate by
that cochain
glues the local
maps.
The displayed
group is a
discrete rank-four
lattice of
translations,
with compact
fundamental
parallelepiped.
Thus $E_1$
is a complex
two-torus.

\item Finally $Q$ is
cyclic and
$E_1/Q=E_0/G$
has the actual
birational
comparison
\eqref{cdp:eq:C.KX.1} to $M$.
\ref{cdp:st:C.2.UN.23} applies
to this torus,
and gives
exactly its
symplectic K3
quotient case.
The order and
exceptional
configuration
calculation
(C.K3.6)--(C.K3.7)
gives $2,3,4,6$
and (C.K3.1).
Conversely every
one of those
actual torus
quotients is
realized by
(C.K3.10), its
proper resolution,
and the two
explicit balancing
operations of
\ref{cdp:st:C.2.UN.24}. This
proves both
directions in
every algebraic
dimension.

\end{enumerate}
\end{proof}






\begin{proposition}[Terminal ruled pairs, their actual contractions and the full normal extension]\hypertarget{tag-01EA}{}\label{cdp:st:C.2.UN.64}
\StatementRecord{cdp:st:C.2.UN.64}{proved}{01EA}{\SeriesRef{cdp:eq:C.EL.6}, \SeriesRef{cdp:eq:C.EL.7}, \SeriesRef{cdp:eq:C.EL.9}, \SeriesRef{cdp:eq:C.EL.11}; \cite[p.~484]{Nakano1970}; \cite[\S2, Theorem~2' and Proposition~9']{FujikiNakano1971}}

Let $f:U\to\Delta$ be a proper degeneration of connected complex
two-tori from a smooth threefold. Assume that the full central fibre
$S$ is reduced and SNC. Let $T$ be one component, meeting the rest
of $S$ in exactly one smooth connected curve $C$, contained in a
single other component $P$. Let $k_T,k_P$ be the orders of the
canonical evaluation along these components. Suppose that either
$q:T\to B$ is a specified holomorphic $\PP^1$-bundle over a smooth
compact curve, or $T\simeq\PP^2$.

For a ruled $T$, precisely one of the following two numerical and
marked-pair descriptions holds. Either $C$ is a section,
\[
 K_T=-2C,\qquad k_T-k_P=1,\qquad
 N_{C/T}\simeq K_B^{-1},\qquad N_{C/P}\simeq K_B,
                                                        \tag{C.EL.1}\hypertarget{tag-00FY}{}\label{cdp:eq:C.EL.1}
\]
where the section identifies $C$ with $B$; or $B\simeq\PP^1$,
$C$ is a smooth anticanonical elliptic bisection and
\[
 K_T=-C,\qquad k_T=k_P,\qquad
 T\simeq\mathbb F_e\quad\hbox{with }e\in\{0,1,2\}.
                                                        \tag{C.EL.2}\hypertarget{tag-00G1}{}\label{cdp:eq:C.EL.2}
\]
For $T\simeq\PP^2$, either $C$ is a line and $k_T-k_P=2$,
or $C$ is a smooth cubic and $k_T=k_P$.
These assertions classify the stated terminal pairs in actual
degenerations; they make no occurrence assertion for the
anticanonical alternatives.

In the section case \eqref{cdp:eq:C.EL.1} there is an actual proper holomorphic
map $\pi:U\to U'$ onto a smooth threefold, equal to an isomorphism
off $T$, which contracts the specified ruling to a smooth curve
$B\subset U'$. It identifies $U$ with the ordinary blowup of $U'$
along $B$. The map $f$ descends to a proper holomorphic torus
degeneration $f':U'\to\Delta$, smooth along $B$; the restriction
$P\to P'$ is an isomorphism and identifies $C$ with $B$.
The full new fibre $S'$ is again reduced and SNC. Its exact normal
bundle and Euler comparisons are
\[
\begin{gathered}
 0\longrightarrow K_B\longrightarrow V:=N_{B/U'}
        \xrightarrow{\,df'\,}\OO_B\longrightarrow0,\qquad
 (T,C)\simeq\bigl(\PP_B(V),\PP_B(K_B)\bigr),\\
 N_{T/U}\simeq\OO_{\PP_B(V)}(-1)\simeq\OO_T(-C),\qquad
 e(S)=e(S')+2-2g(B).
\end{gathered}                                                   \tag{C.EL.3}\hypertarget{tag-00G2}{}\label{cdp:eq:C.EL.3}
\]
Here $\PP_B(V)$ parametrizes lines. The extension class belongs
to the actual one-dimensional space $H^1(B,K_B)$ and is retained,
with its endpoint identifications, by the contraction record.

In the plane-line case the corresponding map contracts $T$ to a
smooth point of a smooth reduced central component, identifies
$U$ with its point blowup, and gives
\[
       P=\operatorname{Bl}_{p}P',\qquad e(S)=e(S')+2.
                                                        \tag{C.EL.4}\hypertarget{tag-00G3}{}\label{cdp:eq:C.EL.4}
\]
These contractions also exist on a compact total threefold over a
compact curve. Iteration on eligible leaves terminates, retaining
the punctured torus family, with exact Euler comparison
\[
 e(S)=e(S_{\rm rem})+2a+\sum_{j=1}^{b}\bigl(2-2g(B_j)\bigr),
                                                        \tag{C.EL.5}\hypertarget{tag-00G4}{}\label{cdp:eq:C.EL.5}
\]
where $a$ is the number of recorded point contractions and the
$B_j$ are the recorded curve centres. Reversing those actual
ordinary blowups restores the original marked support and all
its neighbourhood maps.

For every specified elliptic curve $B$, both the split and the
nonzero-extension pairs in \eqref{cdp:eq:C.EL.3} occur in a compact torus
degeneration over $\PP^1$ with exactly one singular fibre. That
fibre has two components: a two-torus and the prescribed ruled
pair. Its Euler number is zero, and every other fibre is a smooth
two-torus.
\end{proposition}



\begin{proof}

\begin{enumerate}[label={$\langle1\rangle$\arabic*.},leftmargin=*]
\item \par\medskip\noindent\textbf{The actual canonical and normal restrictions.}
A local generator of the rank-one direct image of the canonical
bundle evaluates to a holomorphic three-form on $U$, using the
disc coordinate to trivialize $K_\Delta$. Its divisor is vertical:
on every smooth torus fibre the corresponding nonzero canonical
form has no zeros. Thus, after shrinking the disc, its divisor
is $\sum_v k_vT_v$. Adding a common integer to the $k_v$ does not
change the following comparisons. Since $S$ is the full reduced
principal fibre, restriction to $T$ gives actual line isomorphisms
\[
 \OO_U(T)|_T\simeq\OO_T(-C),\qquad
 K_T\simeq\OO_T\bigl((k_P-k_T-1)C\bigr).
                                                        \tag{C.EL.6}\hypertarget{tag-00G5}{}\label{cdp:eq:C.EL.6}
\]
The second map is adjunction applied to the evaluation section
after dividing by a local equation of $T$ to its order $k_T$.
It also specifies a meromorphic canonical section with exactly
the indicated divisor. At $C$ the reduced equation is $xy=t$;
its product trivialization gives
$N_{C/T}\otimes N_{C/P}\simeq\OO_C$.

\item For a ruling fibre $L\simeq\PP^1$, put $d=C\cdot L$ and
$a=k_P-k_T-1$. Intersecting \eqref{cdp:eq:C.EL.6} with $L$ gives
$-2=ad$. Thus $d>0$, $a<0$, and the integer pairs are exactly
$(d,a)=(1,-2)$ and $(2,-1)$. The map $C\to B$ is finite:
it is a nonconstant proper map of smooth compact curves. In the
degree-one case it is an isomorphism and hence a section.
Adjunction along this section gives
$K_C=\OO_C(-C)$, proving both normal identifications in
\eqref{cdp:eq:C.EL.1}, including their actual dual-normal map.

\item In the degree-two case $K_T+C$ is trivial, so adjunction gives
$K_C\simeq\OO_C$ and $g(C)=1$. Riemann--Hurwitz restricts
$g(B)$ to zero or one. We exclude the latter by the actual
cohomology maps. The projection satisfies
$q_*\OO_T=\OO_B$ and $R^1q_*\OO_T=0$, by the cohomology of
$\PP^1$ in a local bundle chart. The restriction
\[
 H^1(T,\OO_T)\longrightarrow H^1(C,\OO_C)
                                                        \tag{C.EL.7}\hypertarget{tag-00G6}{}\label{cdp:eq:C.EL.7}
\]
is, under this Leray isomorphism, pullback by $C\to B$.
For $g(B)=1$ that degree-two map is etale; its trace composed
with pullback is multiplication by two. Hence \eqref{cdp:eq:C.EL.7} is
nonzero. On the other hand the residue exact sequence
\[
 0\longrightarrow K_T\longrightarrow K_T(C)
       \longrightarrow K_C\longrightarrow0
                                                        \tag{C.EL.8}\hypertarget{tag-00G7}{}\label{cdp:eq:C.EL.8}
\]
has $H^0(K_T)=0$, $H^0(K_T(C))=\C$ and $H^0(K_C)=\C$.
Its residue map is consequently an isomorphism and its
connecting map is zero. Serre duality identifies that connecting
map with the dual of \eqref{cdp:eq:C.EL.7}, by the residue pairing. This
contradiction excludes $g(B)=1$. Thus $B=\PP^1$ and
$T=\mathbb F_e$. The negative section $C_0$ has
$(-K_T)\cdot C_0=2-e$. A distinct irreducible effective
anticanonical curve has nonnegative intersection with $C_0$,
so $e\le2$. This proves \eqref{cdp:eq:C.EL.2}.

\item On the plane, write $C\sim dH$. Equation \eqref{cdp:eq:C.EL.6} instead
gives $ad=-3$. Thus $(d,a)$ is $(1,-3)$ or $(3,-1)$,
giving exactly the stated line and cubic alternatives.

\par\medskip\noindent\textbf{The corrected analytic contraction theorem and its application.}
We use the Main Theorem of \cite[p.~484]{Nakano1970}, in its
corrected and strengthened form in
\cite[\S2, Theorem~2' and Proposition~9']{FujikiNakano1971}.
For a smooth divisor which is a holomorphic projective-space
bundle, the condition that its ambient normal line restricts to
$\OO(-1)$ on every projective fibre produces a smooth analytic
blowdown with that bundle as its exceptional divisor. The
supplement repairs the cohomology argument and proves that this
normal condition suffices. These are the versions read and used
here; their local blowdown maps glue by the actual isomorphisms
off the exceptional divisor.

\item In the section case, \eqref{cdp:eq:C.EL.6} restricts to $\OO_L(-1)$ on
every ruling fibre. The theorem therefore supplies the stated
proper holomorphic blowdown $\pi:U\to U'$ along the given
bundle projection. It is the ordinary blowup of a smooth curve
$B$ in a smooth threefold. The equality
$\pi_*\OO_U=\OO_{U'}$ gives a unique holomorphic $f'$ with
$f=f'\circ\pi$. This equality can also be checked in the two
ordinary blowup charts by extending holomorphic functions
across the smooth codimension-two centre. Properness follows
from $f'^{-1}(K)=\pi(f^{-1}(K))$ for every compact
$K\subset\Delta$.

\item The coefficient of $T$ in $f^{-1}(0)$ is one, so the first
normal term of $f'$ along $B$ is a section of $N_{B/U'}^*$.
Its zero set on $\PP_B(N_{B/U'})$ is exactly $C$.
It is nonzero at every point of $B$: a zero at $b$ would
make that zero set contain the entire ruling fibre over $b$,
whereas $C$ is a section. Its normal differential is therefore
surjective along $B$. Thus $f'$ is a submersion there and
$P'=f'^{-1}(0)$ is smooth near the centre. The restriction of
the blowup to $P'$ is the blowup of its Cartier curve $B$,
which is an isomorphism. This proves $P\simeq P'$ with
the specified curve identification. Every other component
misses the exceptional divisor, so no other central component
meets this centre. The new full fibre remains reduced and SNC.

\item The normal exact sequence of $B\subset P'\subset U'$ is
precisely the first row of \eqref{cdp:eq:C.EL.3}: its quotient is trivialized
by $df'$, and its subline is $N_{C/P}\simeq K_B$ from
\eqref{cdp:eq:C.EL.1}. The exceptional divisor of the ordinary blowup is
the lines projectivization of this very normal bundle. Its
intersection with the strict central component is
$\PP_B(K_B)$. The quotient $V\to\OO_B$ defines this section
as the zero divisor of the corresponding section of
$\OO_{\PP_B(V)}(1)$, which gives the second exact line
identification in \eqref{cdp:eq:C.EL.3}. Serre duality gives
$H^1(B,K_B)\simeq H^0(B,\OO_B)^*\simeq\C$.
In local normal frames the transition matrices have the form
\[
        \begin{pmatrix}k_{ij}&a_{ij}\\0&1\end{pmatrix}.
                                                        \tag{C.EL.9}\hypertarget{tag-00G8}{}\label{cdp:eq:C.EL.9}
\]
Changing a lift of the quotient changes $a_{ij}$ by its
actual Cech coboundary; rescaling the subline multiplies the
extension class by that scalar. Thus there are the zero and
nonzero orbits when endpoint frames are allowed to change.
The contraction retains the matrices \eqref{cdp:eq:C.EL.9}, rather than
discarding this extension class.

\item In the plane-line case the normal line is $\OO_{\PP^2}(-1)$.
The same corrected theorem, with centre a point, produces an
ordinary point blowdown onto a smooth threefold. The exceptional
coefficient one says that the descended $f'$ has a nonzero
linear term at that point. It is consequently a submersion
there. Its central surface pulls back to its point blowup,
with exceptional curve the specified line $C$, proving
\eqref{cdp:eq:C.EL.4}. The normal-dual map gives $C^2_P=-1$, consistently
with this actual surface blowdown.

\par\medskip\noindent\textbf{Euler and neighbourhood comparisons.}
For a curve contraction, the map is an isomorphism off $T$,
with fibre $\PP^1$ over each point of $B$. Additivity gives
$e(S)-e(S')=e(T)-e(B)=e(B)=2-2g(B)$.
For a point contraction it gives $e(T)-1=3-1=2$.
The same maps on a compact total space give the identical
Euler changes there. The descended family is unchanged over
the punctured base, and its fibres there are the original
marked two-tori. Each contraction removes one central
component; every iteration therefore terminates. Summing
these equalities gives \eqref{cdp:eq:C.EL.5}. The record consists of the
actual smooth centre, its normal extension, the local ordinary
blowup charts and their unchanged exterior identifications.
Reversing the maps restores the initial threefold near its
support, so this comparison retains all neighbourhood data.

\par\medskip\noindent\textbf{Both elliptic normal extensions occur compactly.}
Fix the specified elliptic $B$, choose another elliptic curve
$E$, and set $A=B\times E$. Use the compact smooth twistor
torus family (C.PT.13) with period order $R=1$ or $R=2$.
Its central torus is exactly $A$. Fix the curve $B\times\{0\}$
in that central member. In the actual primitive period cover
near this curve, (C.PR.13) writes the deck transformation as
\[
 (p,w,t)\longmapsto
       (p+\delta(t),\gamma(t)w,t),\qquad
 \delta(t)=t^R d(t),\quad d(0)\ne0,
                                                        \tag{C.EL.10}\hypertarget{tag-00FZ}{}\label{cdp:eq:C.EL.10}
\]
where $B=\C^*/\gamma(0)^{\Z}$ and $0<|\gamma(0)|<1$.
The curve is $p=t=0$. Differentiating these actual deck
coordinates gives its normal transition over this elliptic
cover:
\[
 V:\quad
 \begin{pmatrix}p\\t\end{pmatrix}\longmapsto
 \begin{pmatrix}1&\delta'(0)\\0&1\end{pmatrix}
 \begin{pmatrix}p\\t\end{pmatrix}.
                                                        \tag{C.EL.11}\hypertarget{tag-00G0}{}\label{cdp:eq:C.EL.11}
\]
The term $\gamma'(0)w$ is tangent to the central elliptic
curve and hence vanishes in this normal quotient. The other
primitive period has identity normal transition.

\item For $R=2$, \eqref{cdp:eq:C.EL.11} is the identity, so the normal extension
is split. For $R=1$, its upper entry is nonzero. A splitting
would supply a holomorphic $b$ on $\C^*$ satisfying
$b(\gamma(0)w)-b(w)=\delta'(0)$, with one possible opposite
sign from changing the frame convention. The Laurent
expansion on $\C^*$ has zero constant coefficient on the
left and a nonzero one on the right. No such splitting
exists. Thus its extension class is nonzero in $H^1(B,\OO_B)$.
To prescribe the nonzero class while keeping the endpoint
frames and the quotient $df'$ fixed, use instead the actual
period family $V_t(w)=w+a tJ\overline w$, with $a\in\C^*$.
Its central marked torus is unchanged, and differentiation
multiplies the cocycle in \eqref{cdp:eq:C.EL.11} by $a$. Since
$H^1(B,K_B)$ is one-dimensional and the cocycle just proved
is nonzero, this realizes the prescribed class exactly.
The real period determinant is $(1+|at|^2)^2$. At infinity,
dividing the period frame by $at$ gives
$J\overline w+(s/a)w$, with $s=1/t$; this is the same
holomorphic compact completion used in (C.PT.13).

\item Blow up this actual central curve in the compact smooth
twistor family. The exceptional marked pair is exactly
$\bigl(\PP_B(V),\PP_B(\OO_B)\bigr)$, by the ordinary
normal-bundle description just proved. The strict central
torus is unchanged because its centre is Cartier there.
Ordinary coordinates for the ideal $(p,t)$ have $t=xy$,
so the full new fibre has its two components with
multiplicity one. Its Euler number is zero by \eqref{cdp:eq:C.EL.3},
and all other fibres are the unchanged smooth two-tori.

\item The exact degree alternatives, line-frame comparisons,
ordinary blowup Jacobians, normal-extension matrices and
Euler identities are computed in
\path{code/root_exceptional_leaves.py}.

\end{enumerate}
\end{proof}






\begin{proposition}[Trivial local canonical class, triple points and the complete zero case]\hypertarget{tag-01EB}{}\label{cdp:st:C.2.UN.65}
\StatementRecord{cdp:st:C.2.UN.65}{proved}{01EB}{\SeriesRef{cdp:eq:C.CZ.3}, \SeriesRef{cdp:eq:C.CZ.5}, \SeriesRef{cdp:eq:C.CZ.7}, \SeriesRef{cdp:eq:C.CZ.4}, \SeriesRef{cdp:st:C.2.UN.5}, \SeriesRef{cdp:st:C.3.1}, \SeriesRef{cdp:eq:C.CZ.8}, \SeriesRef{cdp:eq:C.CZ.9}, \SeriesRef{cdp:st:C.2.UN.19}, \SeriesRef{cdp:st:C.2.UN.43}, \SeriesRef{cdp:st:C.2.UN.59}, \SeriesRef{cdp:st:C.2.UN.57}, \SeriesRef{cdp:eq:C.CZ.12}, \SeriesRef{cdp:st:C.2.UN.46}; \cite[Lemma 2.1]{Kulikov1977}; \cite[Lemma 2.2]{Kulikov1977}}

Let $f:X\to\Delta$ be a proper degeneration of connected complex
two-tori, with smooth three-dimensional total space and full
reduced simple-normal-crossing central fibre
$S=\sum_iV_i$. Suppose that the orders of the actual canonical
evaluation along all $V_i$ are equal. Equivalently, the canonical
line of a neighbourhood of $S$ is trivial. Let $t$ be the number
of triple points and $r$ the number of rational double curves.
Then
\[
                         e(S)=t,\qquad 2r=3t.
\tag{C.CZ.1}\hypertarget{tag-00DN}{}\label{cdp:eq:C.CZ.1}
\]
In particular $e(S)$ is a nonnegative even integer.

The complete zero-Euler list in this class is as follows.
The fibre is either a smooth two-torus, or is a cycle of smooth
unblown components, each a marked diagonal primary or secondary
Hopf surface, or a marked surface
\[
                   \PP(\OO_E\oplus L),
          \qquad E\text{ elliptic},\quad L\in\operatorname{Pic}^0(E).
\tag{C.CZ.2}\hypertarget{tag-00DT}{}\label{cdp:eq:C.CZ.2}
\]
The boundary curves are exactly the two Hopf axes or the two
ruled sections. Their actual normal lines are dual across
every edge, and the actual integral complement-lattice product
is $M=I$.

Conversely every specified marked cycle in this list with
the stated dual-normal data and $M=I$ has a compact realization
over $\PP^1$, retaining all components and elliptic gluing
maps, with smooth total space and trivial canonical line
near the specified fibre. Its other singular fibre is the
conjugate marked cycle and all reduced fibre supports have
Euler number zero. The construction chooses normal-product
constants. A smooth two-torus is realized by its product
with $\PP^1$.
\end{proposition}



\begin{proof}

\begin{enumerate}[label={$\langle1\rangle$\arabic*.},leftmargin=*]
\item \par\medskip\noindent\textbf{The volume form and the actual double-curve line maps.}
Shrink the disc to remove every other critical value.
The rank-one canonical direct image has an evaluation
generator $\Omega$ whose divisor is supported on $S$.
If all its orders are $k$, the quotient
$\Omega/f^k$ is a nowhere-zero holomorphic three-form.
This is the actual trivialization of $K_X$ used below.
Conversely a nowhere-zero three-form generates this
direct image: every other section is its product with
a holomorphic function constant on the compact connected
fibres, hence pulled back from the disc. Its evaluation
orders are all zero. This proves the equivalence in the
statement by the displayed multiplication and division maps.

\item Put $D_i=\sum_{j\ne i}(V_i\cap V_j)$, with each connected
double curve included once. Adjunction and the principal
divisor $S=(f)$ give the actual isomorphism
\[
                           K_{V_i}\simeq\OO_{V_i}(-D_i).
\tag{C.CZ.3}\hypertarget{tag-00DU}{}\label{cdp:eq:C.CZ.3}
\]
For a connected double curve $C\subset V_i\cap V_j$,
let $T_C$ be its reduced divisor of triple points.
Adjunction on $V_i$ gives
\[
             K_C=(K_{V_i}+C)|_C\simeq\OO_C(-T_C),
             \qquad 2g(C)-2=-\#T_C.
\tag{C.CZ.4}\hypertarget{tag-00DV}{}\label{cdp:eq:C.CZ.4}
\]
Every double curve is therefore elliptic and has no
triple points, or is rational and has exactly two.
Each triple point lies on exactly three double curves.
Counting these actual incidences proves $2r=3t$.

\item Restriction of the trivial line $\OO_X(S)$ to $C$
also gives the exact triple-point comparison
\[
 N_{C/V_i}\otimes N_{C/V_j}\otimes\OO_C(T_C)
                                  \simeq\OO_C.
\tag{C.CZ.5}\hypertarget{tag-00DW}{}\label{cdp:eq:C.CZ.5}
\]
Indeed $\OO_X(V_i)|_C=N_{C/V_j}$,
$\OO_X(V_j)|_C=N_{C/V_i}$, and every remaining
component contributes its actual transverse triple
point. In degrees this is
$C_{V_i}^2+C_{V_j}^2=-\#T_C$.
This is also the direct line-bundle proof of
\cite[Lemma 2.1]{Kulikov1977} in the present notation.

\par\medskip\noindent\textbf{The exact Euler identity.}
By \eqref{cdp:eq:C.CZ.3}, $K_{V_i}^2=D_i^2$.
The self-intersections of the double curves sum,
by \eqref{cdp:eq:C.CZ.5}, to $-3t$. The cross terms in the
boundary squares contribute $2$ on each of the
three surfaces through a triple point. Consequently
\[
                         \sum_iK_{V_i}^2=3t.
\tag{C.CZ.6}\hypertarget{tag-00DX}{}\label{cdp:eq:C.CZ.6}
\]
The alternating restriction sequence on the actual
SNC strata is
\[
 0\longrightarrow\OO_S\longrightarrow\bigoplus_i\OO_{V_i}
  \longrightarrow\bigoplus_C\OO_C
  \longrightarrow\bigoplus_p\C_p\longrightarrow0.
\tag{C.CZ.7}\hypertarget{tag-00DY}{}\label{cdp:eq:C.CZ.7}
\]
Proper flatness gives $\chi(\OO_S)=\chi(\OO_{F_t})=0$.
Only the rational double curves contribute to
$\sum_C\chi(\OO_C)$, so \eqref{cdp:eq:C.CZ.7} gives
\[
                         \sum_i\chi(\OO_{V_i})=r-t.
\tag{C.CZ.8}\hypertarget{tag-00DZ}{}\label{cdp:eq:C.CZ.8}
\]
The sequence and equality also agree with
\cite[Lemma 2.2]{Kulikov1977}.

\item Hirzebruch--Riemann--Roch and complex Gauss--Bonnet
on each compact complex surface give
$e(V_i)=12\chi(\OO_{V_i})-K_{V_i}^2$.
Topological inclusion-exclusion on the same strata
then yields
\[
 \begin{aligned}
 e(S)&=\sum_i e(V_i)-\sum_C e(C)+t\\
     &=12(r-t)-3t-2r+t\\
     &=10r-14t=t,
 \end{aligned}
\tag{C.CZ.9}\hypertarget{tag-00E0}{}\label{cdp:eq:C.CZ.9}
\]
where the last equality uses $2r=3t$.
These formulas require no K\"ahler assumption.
The count $2r=3t$ also makes $t$ even.

\par\medskip\noindent\textbf{The zero-case surfaces are minimal.}
Assume $e(S)=0$. Equations \eqref{cdp:eq:C.CZ.1} and \eqref{cdp:eq:C.CZ.4}
give no triple points and only elliptic double
curves. If there is just one component, the
central fibre is smooth and reduced. The
actual smooth-family comparison in \ref{cdp:st:C.2.UN.5} makes
it a complex two-torus. We therefore treat
the connected reducible case, where every
$D_i$ is nonempty.

\item Equation \eqref{cdp:eq:C.CZ.3} gives
$H^0(V_i,mK_{V_i})=H^0(V_i,\OO(-mD_i))=0$
for every $m>0$: its image in $H^0(V_i,\OO)=\C$
would be a constant vanishing on $D_i$.
Thus every component has Kodaira dimension
$-\infty$. The compact-surface classification
in \SeriesRef{cdp:st:C.3.1} makes its minimal model ruled or
of class VII.

\item For a ruled minimal model over a curve $B$
of genus at least two, every elliptic curve
on $V_i$ maps constantly to $B$, by
Riemann--Hurwitz. A general ruling fibre
$F\simeq\PP^1$ avoids their vertical images
and every blowup centre. It would satisfy
$D_i\cdot F=0$, whereas \eqref{cdp:eq:C.CZ.3} and
adjunction give $D_i\cdot F=-K_{V_i}\cdot F=2$.
This contradiction leaves ruled genus zero
or one, and class VII.

\item Their coherent Euler characteristics are,
respectively, one, zero and zero. Equation
\eqref{cdp:eq:C.CZ.8} is now $\sum_i\chi(\OO_{V_i})=0$,
so no rational component occurs. A ruled
component over an elliptic curve has Euler
number equal to its number of point blowups.
A class VII component has $b_1=b_3=1$ and
$e(V_i)=b_2(V_i)\ge0$. Since all double
curves are elliptic, \eqref{cdp:eq:C.CZ.9} is also
$0=e(S)=\sum_i e(V_i)$. Every term is
therefore zero. The ruled components have
no blowups. Every class VII component is
minimal with $b_2=0$, and its actual
minimal-model restriction in \ref{cdp:st:C.2.UN.19} makes
it a Hopf surface. This proves the claimed
absence of blowups and identifies both
remaining kinds of surface.

\par\medskip\noindent\textbf{The ruled boundary consists of two actual sections.}
Let $\pi:V\to E$ be one of the minimal ruled
surfaces over an elliptic curve. Each component
of $D$ is elliptic and maps finitely to $E$:
it cannot be contained in a rational ruling
fibre. The map is unramified by
Riemann--Hurwitz. Since $D\cdot F=2$, the
only possibilities are two degree-one
sections or one connected elliptic
bisection $C$.

\item The latter is excluded by an exact residue
map. Suppose $K_V+C\simeq\OO_V$.
The sequence
\[
 0\longrightarrow K_V\longrightarrow K_V(C)
      \xrightarrow{\operatorname{res}}K_C\longrightarrow0
\tag{C.CZ.10}\hypertarget{tag-00DO}{}\label{cdp:eq:C.CZ.10}
\]
has a nonzero residue of the nowhere-zero
section of $K_V(C)$: in a local equation
$x=0$ it is the residue of a unit times
$dx\wedge dz/x$. Thus
$H^0(K_V(C))\to H^0(K_C)$ is onto and
its connecting map to $H^1(K_V)$ is zero.
Under Serre duality that connecting map
is dual to the actual restriction
$H^1(V,\OO_V)\to H^1(C,\OO_C)$.
But $\pi^*:H^1(E,\OO_E)\to H^1(V,\OO_V)$
is an isomorphism, and its composite with
restriction is pullback by the degree-two
unramified map $C\to E$. The trace of
that pullback is multiplication by two,
so it is nonzero. This contradicts the
zero connecting map and rules out the
bisection.

\item The two disjoint sections give two
quotient lines of the rank-two ruling
bundle. Their combined quotient map is
an isomorphism on every fibre, because
the sections are distinct. It is therefore
an actual bundle isomorphism with the
sum of those two lines. Tensoring by
one inverse line gives
$V\simeq\PP(\OO_E\oplus L)$ with the
two specified sections. At this point
$L$ may have any degree; that degree
is determined below by the actual
nearby fundamental group.

\par\medskip\noindent\textbf{The Hopf boundary is the individually marked diagonal pair.}
Pull $K_H+D\simeq\OO_H$ back to an actual
central primary Hopf cover, as in \ref{cdp:st:C.2.UN.43}.
For a nonlinear contraction the only
elliptic curve is its axis $C$, with
$K=-(a+1)C$, $a\ge1$. The effective
axis line has infinite order. Thus a
reduced nonempty pullback boundary
cannot satisfy $K+D=0$ there.

\item For a diagonal contraction, the lifted
anticanonical section is
$g(z,w)\partial_z\wedge\partial_w$,
where Hartogs extends $g$ across zero
and
\[
                        g(\alpha z,\beta w)
                                      =\alpha\beta g(z,w).
\tag{C.CZ.11}\hypertarget{tag-00DP}{}\label{cdp:eq:C.CZ.11}
\]
In the nonresonant case its only
possible monomial is $zw$.
In the resonant case a further central
power has multipliers $(\gamma^p,\gamma^q)$,
with coprime positive $p,q$.
Its monomials have weight $p+q$.
For $p,q>1$ this again gives only
$zw$. For $p>1,q=1$ the equation is
$g=w(az+bw^p)$, and reducedness
requires $a\ne0$. The triangular
coordinate change $z\mapsto z+(b/a)w^p$
commutes with the contraction and
gives the two axes. The symmetric
case is identical. When $p=q=1$,
the reduced quadratic is a product
of two distinct linear forms, and
an invertible linear change commuting
with the scalar contraction gives
the axes. Thus the pulled boundary
is exactly two axes in every case.

\item Their unique logarithmic two-form,
up to a constant, is
$dz/z\wedge dw/w$. A deck lift
preserving both axes individually
is diagonal by the contraction-unit
argument of \ref{cdp:st:C.2.UN.59} and fixes this
form. A lift interchanging the axes
is a diagonal map followed by the
interchange and negates the form.
The nonzero section of $K_H(D)$
pulls back to a deck-invariant
section, so no interchange occurs.
Consequently the two images are
distinct elliptic curves, the full
deck group is diagonal, and the
original marked pair is precisely
one of the diagonal primary or
secondary Hopf pairs of \ref{cdp:st:C.2.UN.57}.

\par\medskip\noindent\textbf{The integral obstruction to nonzero ruled degree.}
Every component now has exactly two
boundary curves. The connected
incidence graph is a cycle. At a
ruled vertex its complement is the
actual nonzero line $L^\times$ over
$E$. For $d=\deg L$ its group is
\[
 \pi_1(L^\times)=
 \langle a,b,c\mid [a,c]=[b,c]=1,
                               \ [a,b]=c^d\rangle .
\tag{C.CZ.12}\hypertarget{tag-00DQ}{}\label{cdp:eq:C.CZ.12}
\]
This is the actual circle-bundle
extension: trivialize over an
oriented elliptic parallelogram;
the boundary commutator is the
degree-$d$ clutching loop of its
fibre. The fibre generator $c$
has infinite order, since the
homotopy sequence injects
$\pi_1(\C^*)$ when $\pi_2(E)=0$.

\item For a given smoothing take the
component cores and the actual
$xy=t$ collars. A Hermitian radial
retraction identifies a ruled
core with an annulus bundle in
$L^\times$; either boundary
inclusion induces the isomorphism
onto \eqref{cdp:eq:C.CZ.12}. For a diagonal
Hopf core, the same statement is
its actual rank-three logarithmic
lattice comparison. The dual-normal
maps identify the intervening
circle bundles. The canonical
divisor sections and a small smooth
unit root give these product-preserving
tubular maps for the arbitrary
given smoothing, exactly as in
\ref{cdp:st:C.2.UN.46}. Flatness of $L$ is not used
for this topological comparison.

\item The resulting cycle of group
isomorphisms is a graph of groups
whose vertex groups inject into
the total group. Explicitly cut
one edge, compose the remaining
isomorphisms, and obtain a single
vertex group with one stable
letter implementing the closing
automorphism. Its normal form
contains that vertex group as
a subgroup. Since the actual
nearby group is $\Z^4$, every
injected vertex group is abelian.
Equation \eqref{cdp:eq:C.CZ.12} and the
infinite order of $c$ force
$d=0$ at every ruled vertex.
This proves \eqref{cdp:eq:C.CZ.2} by the
actual collar and clutching maps.

\item The remaining component groups
are all their actual rank-three
abelian lattices. The same cycle
calculation gives
$\pi_1(F_t)=\Lambda\rtimes_M\Z$.
Its injected lattice and the
abelian group $\Z^4$ force
$M=I$. We have obtained exactly
the marked data of \ref{cdp:st:C.2.UN.57}, with
no omitted unblown component
or boundary alternative.

\par\medskip\noindent\textbf{Realization with a nowhere-zero local volume.}
For every such marked cycle,
\ref{cdp:st:C.2.UN.57} gives the actual proper
plumbing, the compact two-ended
period family, and its exact
punctured-family gluing. It
retains all marked components
and elliptic maps. We check
the required local canonical
condition on this construction.

\item On its common period charts
put
\[
                       \Omega=dt\wedge d\eta_1\wedge d\eta_2.
\tag{C.CZ.13}\hypertarget{tag-00DR}{}\label{cdp:eq:C.CZ.13}
\]
Here $\eta$ is the common
vector coordinate obtained
from each component chart
by its actual moving frame
$P_i(t)$. Thus every overlap
has identity linear part in
this common frame.
Every period and affine edge
translation depends holomorphically
on $t$; its differential changes
the vector differentials only by
multiples of $dt$. Thus this
three-form descends through the
actual lattice and collar maps.
It is nonzero on component cores.
At a double curve the relative
two-form is a holomorphic unit
times $dx/x\wedge dz$, where
$z$ is an elliptic coordinate
and $xy=t$. Hence
\[
             dt\wedge(dx/x\wedge dz)
                               =dy\wedge dx\wedge dz.
\tag{C.CZ.14}\hypertarget{tag-00DS}{}\label{cdp:eq:C.CZ.14}
\]
The right side is holomorphic
and nowhere zero. The moving
Hopf and ruled frames have
nonzero determinant at the
central value, so the displayed
unit never vanishes after
shrinking the disc. The form
therefore extends across all
of $S$ as a nowhere-zero
three-form. Its evaluation
orders are equal to zero.
The compact realization in
the statement now follows
from \ref{cdp:st:C.2.UN.57}, with the same
verification at its conjugate
fibre. The smooth torus case
is its literal product family.

\end{enumerate}
\end{proof}



\par\medskip\noindent\textbf{Exact relation to the primary modification literature.}
Kulikov's definition of a good degeneration in
\cite[original Russian ed., pp.~1008--1009]{Kulikov1977} is the smooth-total,
reduced SNC condition used here. Its Lemmas 2.1--2.2
give the same principal-divisor and coherent-stratum
maps as \eqref{cdp:eq:C.CZ.5}--\eqref{cdp:eq:C.CZ.8}. The general canonical
modification statement is
\cite[Theorem 1$'$]{Kulikov1981}: a good degeneration
of surfaces with trivial canonical class admits a
good modification with trivial total canonical
class precisely under its stated minimum-order
condition on class VII components. The present
equal-order hypothesis gives the actual identity
comparison to such a target through $\Omega/f^k$.
The zero-case classification above is proved on
that given neighbourhood and supplies the explicit
marked-cycle maps and compact occurrence. The
1981 appendix also records the correction to the
1977 Enriques modification assertion; the cited
line and Euler comparisons are the displayed
1977 lemmas.

Kamenova's author-hosted primary text
\cite[Theorem 3.1]{KamenovaTori} assumes a
K\"ahler total space and trivial $K_X$.
Its necessary list contains a smooth torus,
elliptic-ruled paths and cycles, and a rational
case with a triangulated torus as dual complex.
The residue and finite-trace maps in
\eqref{cdp:eq:C.CZ.10} exclude a terminal elliptic bisection,
so the path alternative contributes no
additional zero-Euler support in the stated
smooth-component setting. This sharpens that
necessary list. In the K\"ahler case Hopf
components are also excluded by restriction
of the K\"ahler form: their vanishing real
$H^2$ would make its positive square integrate
to zero by Stokes. The present proof keeps
the actual component and collar maps without
a K\"ahler hypothesis and supplies compact
occurrence of every marked zero-Euler cycle
specified in the statement.




\begin{proposition}[Actual leaf contractions, elliptic cores and the intrinsic Euler formula]\hypertarget{tag-01EC}{}\label{cdp:st:C.2.UN.66}
\StatementRecord{cdp:st:C.2.UN.66}{proved}{01EC}{\SeriesRef{cdp:st:C.2.UN.64}, \SeriesRef{cdp:st:C.2.UN.65}, \SeriesRef{cdp:st:C.2.UN.61}, \SeriesRef{cdp:st:C.2.UN.62}, \SeriesRef{cdp:eq:C.MR.3}, \SeriesRef{cdp:st:C.2.UN.5}, \SeriesRef{cdp:eq:C.MR.5}, \SeriesRef{cdp:st:C.2.UN.19}, \SeriesRef{cdp:st:C.2.UN.58}, \SeriesRef{cdp:st:C.2.UN.50}, \SeriesRef{cdp:st:C.2.UN.29}, \SeriesRef{cdp:st:C.2.UN.22}, \SeriesRef{cdp:st:C.2.UN.28}, \SeriesRef{cdp:eq:C.MR.10}, \SeriesRef{cdp:eq:C.MR.11}, \SeriesRef{cdp:st:C.2.UN.38}, \SeriesRef{cdp:st:C.2.UN.31}, \SeriesRef{cdp:eq:C.MR.12}, \SeriesRef{cdp:st:C.2.UN.33}, \SeriesRef{cdp:eq:C.MR.8}, \SeriesRef{cdp:eq:C.MR.4}, \SeriesRef{cdp:eq:C.MR.9}, \SeriesRef{cdp:st:C.2.UN.47}, \SeriesRef{cdp:st:C.2.UN.54}, \SeriesRef{cdp:st:C.2.UN.44}, \SeriesRef{cdp:eq:C.MR.15}, \SeriesRef{cdp:eq:C.MR.16}}

Let $f:U\to\Delta$ be a proper degeneration of connected complex
two-tori from a smooth threefold. Suppose that its full central
fibre $S=\sum_vT_v$ is reduced and SNC, has no triple points,
and that every component is either a smooth minimal compact
complex surface or a holomorphic $\PP^1$-bundle over a smooth
compact curve. The latter convention includes $\mathbb F_1$.
All components, intersection curves and neighbourhood maps refer
to this actual degeneration.

There is a finite sequence of the ordinary holomorphic
contractions of \ref{cdp:st:C.2.UN.64}, each along a terminal component, which
retains the punctured torus family and ends in one of the
following marked supports:
\begin{enumerate}
\item a smooth two-torus;
\item a marked diagonal Hopf and degree-zero elliptic ruled
cycle of \ref{cdp:st:C.2.UN.65}, with its actual dual normal maps and
complement-lattice product $M=I$;
\item a rooted elliptic tree passing the complete test of
\ref{cdp:st:C.2.UN.61}, including its actual zero-Chern inverse torsor;
\item a unicyclic genus-zero elliptic quasi-bundle support
passing the complete test of \ref{cdp:st:C.2.UN.62}.
\end{enumerate}
Every original surface and neighbourhood is recovered by
reversing the recorded ordinary blowups. In particular the
normal extension at every curve centre is retained.

No rational component occurs in $S$. Let $\mathcal R_{\ge2}$
be the set of its ruled components whose ruling bases $B_v$
have genus at least two. Then
\[
 e(S)=-2\sum_{v\in\mathcal R_{\ge2}}\bigl(g(B_v)-1\bigr)
                                                        \le0.
\tag{C.MR.1}\hypertarget{tag-00R1}{}\label{cdp:eq:C.MR.1}
\]
Consequently $e(S)=0$ exactly when $\mathcal R_{\ge2}$ is
empty. In that case all the contractions in the preceding
reduction have elliptic curve centres. Their exceptional
marked pairs are precisely
\[
 \left(\PP_B(V),\PP_B(\OO_B)\right),
 \qquad
 0\longrightarrow\OO_B\longrightarrow V
             \longrightarrow\OO_B\longrightarrow0,
 \qquad g(B)=1,
\tag{C.MR.2}\hypertarget{tag-00R9}{}\label{cdp:eq:C.MR.2}
\]
with their actual split or nonzero extension and endpoint
identifications.

Conversely, start with an actual torus degeneration having
one of the displayed cores and perform any finite sequence of
ordinary blowups along smooth compact elliptic curves in the
smooth locus of its full reduced central fibre, avoiding its
double curves at each stage. The result is again a full
reduced SNC torus fibre without triple points, its components
are in the stated minimal-or-bundle class, and its Euler
number is zero. These are actual blowup maps on the chosen
neighbourhood. The reduction does not identify an arbitrary
other realization of the same core support with that
neighbourhood or discard the recorded normal extensions.

There is also a complete extension with no initial minimality
or no-triple assumption. For any full reduced SNC torus fibre
whose connected double curves are all elliptic, $e(S)\ge0$.
Its Euler number is zero exactly when it has no triple
points and every component is minimal and nonrational.
Every such zero-Euler fibre is therefore covered by the
preceding reduction and exact reconstruction.
\end{proposition}



\begin{proof}\leavevmode\par



 
\begin{enumerate}[label={$\langle1\rangle$\arabic*.},leftmargin=0pt,labelsep=.5em]
\item \par\medskip\noindent\textbf{Canonical orders and actual contractions.}
For the connected reducible case let $k_v$ be the orders of
a canonical evaluation generator. Its divisor is vertical,
and adjunction with the full principal fibre gives a
meromorphic canonical section on each component with exact
divisor
\[
 K_{T_v}=
       \sum_{e=(v,w)}(k_w-k_v-1)C_e.
\tag{C.MR.3}\hypertarget{tag-00RA}{}\label{cdp:eq:C.MR.3}
\]
The notation denotes this section and its line isomorphism.
The curves in the sum are pairwise disjoint because there
are no triple points. Along each edge the actual equation
$xy=t$ gives
\[
                 N_{C_e/T_v}\otimes N_{C_e/T_w}
                                        \simeq\OO_{C_e}.
\tag{C.MR.4}\hypertarget{tag-00RB}{}\label{cdp:eq:C.MR.4}
\]


\begin{enumerate}[label={$\langle2\rangle$\arabic*.},leftmargin=1.3em,labelsep=.5em]
\item  Apply \ref{cdp:st:C.2.UN.64} whenever a terminal ruled component meets its
parent in a section, or a terminal plane meets its parent
in a line. Each is the actual ordinary blowdown of the
specified smooth curve or point, and the new family is
proper, smooth-total, reduced and SNC. Its other components
and marked curves are unchanged except for the parent
comparison proved there.


\item  This procedure preserves the surface convention in the
statement. At a curve contraction the parent surface is
unchanged. At a plane-line contraction the parent contains
the boundary $(-1)$-curve. A minimal surface has none.
Among the allowed projective-line bundles, a rational
curve on a bundle over a positive-genus base is a ruling
fibre and has square zero. On $\mathbb F_n$ the only
$(-1)$-curve occurs on $\mathbb F_1$ and is its negative
section: an irreducible curve different from the negative
section has class $aC_0+bf$ with $b\ge na$ when $a>0$,
and hence nonnegative square; the fibres have square zero.
Thus that parent is $\mathbb F_1$, and its specified
surface blowdown is $\PP^2$. It remains in the allowed
class. Every other boundary was disjoint from the
contracted curve, so no new intersection or triple point
is introduced.


\item  Each operation removes one component, so it terminates.
The canonical order differences on surviving components
are unchanged: the ordinary curve and point blowup
Jacobians have exceptional orders one and two,
respectively, as in (C.EL.1) and (C.EL.4). The remaining
family therefore still has \eqref{cdp:eq:C.MR.3}, and no eligible
\ref{cdp:st:C.2.UN.64} leaf. We call its support $S_0$.
For one smooth reduced component, \ref{cdp:st:C.2.UN.5} already identifies
the fibre as a two-torus. Assume below that $S_0$ is
reducible.


\end{enumerate}

\item \par\medskip\noindent\textbf{A nonconstant maximum is a Hopf leaf.}
At a component $T$ of maximum canonical order, put
$a_e=k_T-k_w+1\ge1$. Equation \eqref{cdp:eq:C.MR.3} reads
\[
                         -K_T=\sum_e a_eC_e.
\tag{C.MR.5}\hypertarget{tag-00RC}{}\label{cdp:eq:C.MR.5}
\]
Its boundary is nonempty. If a positive pluricanonical
power had a nonzero section, its product with the
corresponding power of the nonzero section in \eqref{cdp:eq:C.MR.5}
would be a nonzero holomorphic function with a nonempty
zero divisor on the compact connected surface. This is
impossible. Thus $\kappa(T)=-\infty$. The surface
classification and the actual class-VII restriction \ref{cdp:st:C.2.UN.19}
leave a projective-line bundle, $\PP^2$, or a Hopf
surface.


\begin{enumerate}[label={$\langle2\rangle$\arabic*.},leftmargin=1.3em,labelsep=.5em]
\item  For a ruled $T\to B$, intersect \eqref{cdp:eq:C.MR.5} with a general
ruling fibre. No boundary curve is vertical: a vertical
curve would be a ruling fibre and would meet every
horizontal boundary, whereas some horizontal boundary
is required by the following equality. Hence
\[
              2=\sum_e a_ed_e,\qquad
              d_e=C_e\cdot(\text{ruling fibre})>0.
\tag{C.MR.6}\hypertarget{tag-00RD}{}\label{cdp:eq:C.MR.6}
\]
The integer partitions are exactly one section of
coefficient two, one bisection of coefficient one,
or two sections of coefficient one each. The first
is an eligible contraction. The terminal-pair
classification in \ref{cdp:st:C.2.UN.64} makes the second a rational
ruled surface with its smooth anticanonical elliptic
bisection; its neighbour has the same canonical order.
In the third case adjunction along either disjoint
section gives
\[
       2g(B)-2=(K_T+C_e)\cdot C_e=0.
\]
Thus $B$ is elliptic, and both neighbours again have
the same order as $T$. On $\PP^2$ disjoint effective
boundary curves can have only one component, by
Bezout. The two possibilities are the eligible
plane-line contraction or the anticanonical cubic
with an equal-order neighbour.


\item  For a Hopf surface of algebraic dimension one use
its actual elliptic map over $\PP^1$, the connected
normal-root chart and the invariant relative
differential of \ref{cdp:st:C.2.UN.58}. All its curves are fibre
supports. The canonical order formula gives
\[
 d_{T,e}:=\frac{k_w-k_T}{m_e}\in\Z_{\le0},
 \qquad
              \sum_{e\ni T}d_{T,e}
                                  =\operatorname{val}(T)-2.
\tag{C.MR.7}\hypertarget{tag-00RE}{}\label{cdp:eq:C.MR.7}
\]
An unmarked multiple fibre is excluded by the same
root formula $m\operatorname{ord}\theta+m-1=0$.
Thus the valence is one or two. Valence one gives
$k_T-k_w=m_e>0$; valence two gives two equal-order
neighbours. A nonlinear Hopf surface has its single
elliptic axis, so is a leaf; its actual canonical
axis form gives the positive drop equal to its
resonance, including the secondary quotient in \ref{cdp:st:C.2.UN.50}.


\item  For completeness, a diagonal algebraic-dimension-zero
Hopf surface cannot give another maximum case. Pull
\eqref{cdp:eq:C.MR.5} to its primary contraction cover. Its only
curves are the two axes. If both occur, the
anticanonical eigenfunction equation gives
$\alpha^{a_1-1}\beta^{a_2-1}=1$.
Since $|\alpha|,|\beta|<1$ and $a_i\ge1$, both
exponents are zero. If only one occurs, the same
equation gives a positive integral resonance
between $\alpha$ and $\beta$, contradicting
algebraic dimension zero. Thus this case has
exactly the two equal-order boundaries.


\item  It follows that every noneligible maximum component
which is not a Hopf leaf has all of its neighbours
at that same maximum. An equal-order edge cannot
lead to one of the strict-drop leaves just listed.
Following these edges therefore gives an entire
connected component of the whole incidence graph.
Connectedness of $S_0$ forces all orders to be equal.
In that event \ref{cdp:st:C.2.UN.65} applies; since there are no triple
points, its complete zero-case classification is
exactly the displayed cycle or smooth torus.
If the orders are not all equal, a maximum Hopf
leaf exists. Its boundary is elliptic and its normal
line has finite exact order.


\end{enumerate}

\item \par\medskip\noindent\textbf{The actual field space aligns a vertical elliptic fibration.}
We record the comparison needed to propagate from that
leaf. On any component $T$ of the current reduced smooth
SNC family, the actual multiplicity-one construction
(C.CF.3)--(C.CF.7) gives two independent commuting
meromorphic fields, all of whose constant combinations
are semicomplete. Their polar divisors are confined to
the boundary of $T$. Indeed the original module consists
of holomorphic fields on $U$; multiplying its basis by
meromorphic functions of $t$ introduces poles only in
the central fibre. Cancelling the generic order on
$T$ makes the restriction holomorphic at every point
of $T$ away from its intersections with the other
components. This is a restriction on the original
component, with no different covering surface substituted.


\begin{enumerate}[label={$\langle2\rangle$\arabic*.},leftmargin=1.3em,labelsep=.5em]
\item  Let $q:T\to B$ be an actual connected-fibre elliptic
map for which every boundary is vertical. Over the
complement of finitely many values the fields are
holomorphic on the entire compact fibres. Their
coefficients in a base direction are holomorphic
functions on each such fibre and hence constant.
Proper Stein descent gives commuting meromorphic
fields on $B$.


\item  The projected fields have no poles, by the exact
continuation argument in (C.RF.1). At a pole of
order $r$, the primitive of the one-dimensional
time form vanishes to order $r+1$. Continue one
inverse branch around its time value, using a
small loop whose lifted base paths avoid all
other bad values. The preimage under $q$ of
the compact union of these paths is compact
and contained in the holomorphic locus.
The upstairs solution therefore continues
along that whole loop. Its endpoint has a
different base value, contradicting
semicompleteness. This proves holomorphy of
the projected fields without an assumed
descent theorem for semicompleteness.


\item  If a projection is nonzero, any other commuting
projection is a constant multiple of it:
the derivative of their meromorphic ratio
along the nonzero field is zero. Subtracting
that constant multiple upstairs gives a
nonzero vertical commuting field. The
actual period-preservation map (C.RF.2)
for this transverse/vertical pair shows
that the integral elliptic periods are
constant and monodromy is the identity.
If both projections are zero, the two
fields are translation fields on each
regular elliptic fibre. Their nonconstant
meromorphic ratio is constant there,
so its connected Stein fibres are
exactly those of $q$. Applying \ref{cdp:st:C.2.UN.29}
to this actual first integral gives
the same conclusion.


\item  Thus the fibration just constructed has
one fixed Jacobian and identity integral
monodromy. Its multiple-root comparison
in \ref{cdp:st:C.2.UN.22}/29 gives only smooth multiple
supports, and its invariant cohomology
frame gives
\[
             R^1q_*\OO_T=\OO_B,\qquad
                         p_g(T)=g(B)\le1.
\tag{C.MR.8}\hypertarget{tag-00RF}{}\label{cdp:eq:C.MR.8}
\]
The last inequality is the actual simultaneous
compact-core bound \ref{cdp:st:C.2.UN.28}.


\end{enumerate}

\item \par\medskip\noindent\textbf{A torsion-normal elliptic boundary on a ruled component.}
Consider a component $T$ meeting a known elliptic
boundary $C$ whose normal line is torsion. It has
$C^2=0$. A minimal rational surface contains no
such elliptic curve. On $\PP^2$ a smooth elliptic
curve has degree three and square nine. On a
Hirzebruch surface write an irreducible elliptic
curve as $aC_0+bf$, with $a>0$. Its arithmetic
genus equation is
\[
           (a-1)(2b-na-2)=2.
\tag{C.MR.9}\hypertarget{tag-00RG}{}\label{cdp:eq:C.MR.9}
\]
It is distinct from $C_0$ and has $b\ge na$.
The only possibilities have $a=2$, $b=n+2$,
$n\le2$, with square eight, or $a=3,n=1,b=3$,
with square nine. These include $\mathbb F_1$.
None has square zero. On a ruled surface over
a base of genus at least two an elliptic curve
cannot be horizontal by Riemann--Hurwitz, and
cannot be contained in a rational ruling fibre.
Those components are excluded as well.


\begin{enumerate}[label={$\langle2\rangle$\arabic*.},leftmargin=1.3em,labelsep=.5em]
\item  Suppose that $T\to E$ is ruled over an elliptic
curve. Its finite map $C\to E$ is unramified.
Pull the bundle back along this actual map.
The diagonal copy of $C$ is a section whose
normal line is the original torsion line $L$.
Its relative degree-one line $\OO(C)$ has
rank-two direct image, and the actual
fibrewise evaluation identifies this
pulled-back surface with the corresponding
projective bundle. After twisting that
rank-two bundle by its section subline,
it has the exact form
\[
     0\longrightarrow\OO_C\longrightarrow V
                         \longrightarrow L\longrightarrow0.
\tag{C.MR.10}\hypertarget{tag-00R2}{}\label{cdp:eq:C.MR.10}
\]
If $L\ne\OO_C$, its extension group is
$H^1(C,L^{-1})=H^0(C,L)^*=0$, so it splits.
If $L=\OO_C$, the complement of this section
is the additive torsor of the actual extension
class in $H^1(C,\OO_C)$.


\item  There is a second boundary curve on $T$.
Otherwise \ref{cdp:st:C.2.UN.64} would make this terminal pair
a section leaf: its connected anticanonical
bisection over an elliptic ruling base is
excluded by the residue/finite-trace map.
It would already have been contracted.
Pull a second boundary back and normalize
one component. It is a compact curve mapping
finitely onto $C$ and lying in the complement
of the section. It supplies a section of
the pulled-back additive torsor. Finite
trace composed with pullback on $H^1(\OO)$
is multiplication by its positive degree,
so the original additive class is zero.
Thus \eqref{cdp:eq:C.MR.10} splits in this case also.


\item  Kill the torsion $L$ by a further finite
etale cover. The resulting actual surface
is $E'\times\PP^1$, with its finite Galois
covering map to $T$. The base cover is
an isogeny after selecting origins; its
deck group acts by translations on $E'$.
A deck map on the product is a translation
on $E'$ and a holomorphic map
$E'\to\operatorname{PGL}_2(\C)$ on the
projective fibre. The latter is constant:
$\operatorname{PGL}_2$ is affine, so its
coordinate functions are constant on the
compact elliptic curve. These constants
give an actual finite action on $\PP^1$.
Taking its quotient constructs
\[
        T=(E'\times\PP^1)/G
                         \longrightarrow\PP^1/G.
\tag{C.MR.11}\hypertarget{tag-00R3}{}\label{cdp:eq:C.MR.11}
\]
Its regular connected fibre is $E'/K$,
where $K$ is the kernel of the projective
action. All fibre actions are translations.
At a projective stabilizer the base action
on $E'$ is free, so the local quotient
is exactly the smooth multiple chart
of (C.QT.2). This proves fixed Jacobian
and identity integral monodromy directly.


\item  Every other boundary pulls back to a
compact curve disjoint from the fixed
section of $E'\times\PP^1$. Its second
coordinate is a holomorphic function
with values in the affine complement
of that section, and is constant.
It is therefore a constant section.
Consequently every original boundary
is a fibre support of \eqref{cdp:eq:C.MR.11}, with
its actual normal-root cover retained.


\end{enumerate}

\item \par\medskip\noindent\textbf{A torsion-normal elliptic boundary on a nonruled component.}
A Hopf component with a torsion-normal
elliptic curve is either a nonlinear
leaf, or is an elliptic Hopf component
whose curves are fibres of its fixed
Jacobian map over $\PP^1$. The
algebraic-dimension-zero diagonal case
has nontorsion axis normals and cannot
occur here, by the exact characters
in \ref{cdp:st:C.2.UN.38} and \ref{cdp:st:C.2.UN.58}.


\begin{enumerate}[label={$\langle2\rangle$\arabic*.},leftmargin=1.3em,labelsep=.5em]
\item  For a minimal component of nonnegative
Kodaira dimension, the common-root
comparison in \ref{cdp:st:C.2.UN.31} has degree one in
this full reduced SNC family. Thus
the component itself is a two-torus
or has the actual smooth-multiple,
fixed-Jacobian elliptic map of that
theorem. There is no further surface
quotient in this use of \ref{cdp:st:C.2.UN.31}.
On a torus the elliptic curve $C$
is a translate of an elliptic subtorus.
Its actual quotient makes $C$ a fibre,
and every disjoint boundary is vertical.


\item  Let $q_0:T\to B_0$ be the elliptic
alternative. For $N$ divisible by
all its fibre multiplicities, its
translation-root differential gives
the actual canonical line formula
\[
 K_T^{\otimes N}
  =q_0^*\left(K_{B_0}^{\otimes N}
       \otimes\OO_{B_0}
          \left(\sum_iN(1-1/m_i)b_i\right)\right).
\tag{C.MR.12}\hypertarget{tag-00R4}{}\label{cdp:eq:C.MR.12}
\]
Put $\delta=2g(B_0)-2+\sum_i(1-1/m_i)$.
In Kodaira dimension one $\delta>0$,
as follows from the exact base
pluricanonical pushdown. Since
$C$ is elliptic and $C^2=0$,
adjunction gives $K_T\cdot C=0$.
Equation \eqref{cdp:eq:C.MR.12} then makes $C$
vertical. Every disjoint boundary
is also vertical, since a horizontal
curve meets this entire fibre support.


\item  In Kodaira dimension zero, $\delta=0$.
For an elliptic $B_0$ there are no
multiple fibres, so $q_0$ is a
principal elliptic torsor. A horizontal
$C$ maps finitely and unramifiedly
onto $B_0$ and gives a section after
that actual base pullback. The integral
Chern vector therefore pulls back to
zero. Pullback is multiplication by
the positive covering degree on the
free lattice group $H^2(B_0,\Lambda_F)$,
so this vector is zero. The additive
torsor comparison (C.PT.7)--(C.PT.8)
identifies this very surface with
a complex two-torus. This is already
the preceding torus case. If $C$ is
vertical there is nothing to change.


\item  For $B_0=\PP^1$, $\delta=0$ gives
exactly the four multiplicity lists
\[
       (2,2,2,2),\quad(3,3,3),\quad
                       (2,4,4),\quad(2,3,6).
\tag{C.MR.13}\hypertarget{tag-00R5}{}\label{cdp:eq:C.MR.13}
\]
There is an actual elliptic covering
$E\to\PP^1$ with precisely these
ramification indices. For four
specified double points it is the
double cover branched at those
points. For the three-point lists
it is the elliptic quotient by
the cyclic group of order three,
four or six, with the three
branch values moved to the
specified points by a projective
coordinate. These are the actual
quotient maps and stabilizer lists
in \ref{cdp:st:C.2.UN.33}.


\item  Normalize the pullback of $T$ by
this map. In each chart (C.QT.2)
it is the finite etale surface
map $F\times\Delta_s\to
(F\times\Delta_s)/\mu_{m_i}$;
the free translation removes
the multiple fibre. Thus the
global pullback $T'\to T$ is
an actual finite etale Galois
cover and $T'\to E$ is a
principal fixed-elliptic torsor.
If $C$ was horizontal, its
pullback components are elliptic
and map finitely to $E$.
The same trace/Chern argument
and (C.PT.8) make $T'$ an
actual two-torus.


\item  The preimage of $C$ in $T'$ is
a finite disjoint union of
parallel elliptic translates.
The deck group preserves their
elliptic subtorus and therefore
acts on its elliptic quotient.
The actual descended map
\[
 T=T'/G\longrightarrow
                  (T'/F')/G
\tag{C.MR.14}\hypertarget{tag-00R6}{}\label{cdp:eq:C.MR.14}
\]
has connected elliptic fibres.
Indeed the kernel on the base
acts by translations on $F'$:
a nontranslation finite elliptic
automorphism would have a fixed
point on a fibre, contradicting
the etale surface action.
The same argument at a base
stabilizer gives precisely
smooth multiple fibres.
The curve $C$ is one such
fibre support, and every
disjoint boundary is vertical.
The field-space comparison
leading to \eqref{cdp:eq:C.MR.8} now applies
to this very descended map,
and removes any possible
nontrivial linear monodromy.
This completes the nonruled
alignment by actual finite
cover and quotient morphisms.


\end{enumerate}

\item \par\medskip\noindent\textbf{Propagation and the exhaustive remaining cores.}
Start with the maximum Hopf
leaf when the canonical orders
are not constant. Its boundary
has finite normal order, and
\eqref{cdp:eq:C.MR.4} gives the same property
on its neighbour. The preceding
comparisons put that neighbour
in the fixed-Jacobian elliptic
class, with all its boundaries
vertical, unless it is itself
a nonlinear leaf. Every vertical
boundary again has finite exact
normal order by (C.QT.2).
Repeat across the connected
incidence graph. No rational
or higher-genus ruled component
can occur in this propagation,
by \eqref{cdp:eq:C.MR.9} and
Riemann--Hurwitz. A one-boundary
elliptic ruled surface would be
an eligible \ref{cdp:st:C.2.UN.64} section leaf
and is absent from $S_0$.
Consequently the entire $S_0$
consists of the displayed
elliptic quasi-bundles and
nonlinear Hopf leaves.


\begin{enumerate}[label={$\langle2\rangle$\arabic*.},leftmargin=1.3em,labelsep=.5em]
\item  Equation \eqref{cdp:eq:C.MR.8} and the
simultaneous bound \ref{cdp:st:C.2.UN.28} give
at most one genus-one base.
If there is one, its component
has $p_g=1$. The actual
component-group and dual-complex
comparison \ref{cdp:st:C.2.UN.47} makes the
incidence graph simply connected.
It is therefore a tree. The
entire support is in the exact
scope of the extended \ref{cdp:st:C.2.UN.61}
and satisfies all its necessary
tests, including its integral
inverse-torsor obstruction.
If there is no genus-one base,
all vertices have coherent
dimensions $(1,1,0)$ and the
nonzero boundary restrictions
of (C.QT.3) or \ref{cdp:st:C.2.UN.54}.
The actual graph quotient of
\ref{cdp:st:C.2.UN.44} has at most one cycle.
Leaf elimination would give
$H^2(\OO_{S_0})=0$ on a tree,
contradicting upper
semicontinuity from a torus.
Thus it is unicyclic and the
whole support is in the exact
scope of \ref{cdp:st:C.2.UN.62}, with all of
its necessary tests satisfied.
This proves that the
contraction procedure reaches
one of the stated cores.
No support-only inverse
logarithm has been used to
assert a new smoothing:
the intermediate degenerations
are the actual \ref{cdp:st:C.2.UN.64} blowdowns,
and \ref{cdp:st:C.2.UN.61}/62 are applied to
their actual final family.


\end{enumerate}

\item \par\medskip\noindent\textbf{Rational components cannot persist or disappear.}
We now use the proved endpoint
description to exclude rational
components in the original
support. A nonrational
component in the allowed class
has no rational boundary.
For a Hopf surface a rational
curve would lift to
$\C^2\setminus\{0\}$ and its
coordinate functions would be
constant. On a ruled surface
over a positive-genus base a
rational curve is a ruling
fibre. If it were a boundary,
disjointness would leave no
horizontal boundary. A general
ruling fibre would then have
intersection zero with the
right side of \eqref{cdp:eq:C.MR.3}, while
its canonical degree is $-2$.
This is impossible.


\begin{enumerate}[label={$\langle2\rangle$\arabic*.},leftmargin=1.3em,labelsep=.5em]
\item  For a minimal component with
$\kappa\ge0$, use its degree-one
\ref{cdp:st:C.2.UN.31} description. A torus has
no rational curve by its
universal covering. In the
elliptic alternative \eqref{cdp:eq:C.MR.12}
has $\delta\ge0$. A rational
curve cannot be vertical in
a smooth elliptic support.
If its map to the base had
degree $d>0$, each local
ramification order over $b_i$
would be divisible by $m_i$,
from the actual equation
$t=u x^{m_i}$. Riemann--Hurwitz
would therefore give
\[
 -2\ \ge\
 d\left(2g(B_0)-2+
              \sum_i(1-1/m_i)\right)
                              =d\delta\ge0,
\tag{C.MR.15}\hypertarget{tag-00R7}{}\label{cdp:eq:C.MR.15}
\]
a contradiction. This proves
the asserted absence of
rational boundaries at every
stage of the reduction.


\item  Whenever a rational component
is removed by \ref{cdp:st:C.2.UN.64}, its
section or line boundary is
rational. Its parent must
therefore be rational too.
The operation cannot remove
the last rational component.
Every final core, however,
has no rational component:
this is explicit in \ref{cdp:st:C.2.UN.65},
and in \ref{cdp:st:C.2.UN.61}/62 every
component has coherent
Euler zero. Thus no rational
component existed initially.
In particular no plane-line
contraction occurred.


\end{enumerate}

\item \par\medskip\noindent\textbf{The intrinsic Euler identity and the complete zero description.}
All contractions are now
curve contractions, with
base genus at least one.
All original ruled components
of base genus at least two
must be contracted, since
none occurs in the final
cores. A curve contraction
does not modify its parent
surface, and there are no
point contractions. Hence
these are exactly the
original higher-genus ruled
components, each counted once
with its original ruling base.
Every other contracted
component has elliptic base.
The exact Euler comparison
(C.EL.5), and the zero Euler
number of each final core,
give precisely \eqref{cdp:eq:C.MR.1}.
This also proves the stated
zero criterion independently
of any choice of contraction
order.


\begin{enumerate}[label={$\langle2\rangle$\arabic*.},leftmargin=1.3em,labelsep=.5em]
\item  In the zero case every
centre is elliptic.
Equation (C.EL.3) then
gives \eqref{cdp:eq:C.MR.2}, including
its full normal extension.
The recorded ordinary
blowups and unchanged
exterior identifications
recover the original
neighbourhood exactly.


\item  For the converse construction,
the displayed cores are in
the stated surface class.
For the quasi-bundle vertices
their actual smooth-multiple
fibration gives $e=\chi(\OO)=0$.
In nonnegative Kodaira dimension
the argument \eqref{cdp:eq:C.MR.15} excludes
exceptional rational curves.
In Kodaira dimension $-\infty$,
the surface classification,
$\chi(\OO)=0$ and $e=0$
give an unblown elliptic ruled
surface or a Hopf surface.
The root and the canonical-flat
cores are already minimal by
\ref{cdp:st:C.2.UN.61} and \ref{cdp:st:C.2.UN.65}.


\item  Now
an elliptic centre in the
smooth central locus lies
in exactly one component.
The full fibre is reduced,
so the map is a submersion
near it. Blowing up its
two-generator ideal gives
the ordinary local equation
$t=xy$, with two full
multiplicity-one components.
The strict parent is the
blowup of a Cartier curve
and hence is unchanged;
the new component is the
projective normal bundle
over the elliptic centre.
No other boundary meets the
centre, so there are no
new triple points.
Properness and the punctured
torus family are preserved
by this actual blowup.
Its Euler change is
$2-2g(B)=0$. Iterating
proves the converse in
the statement and retains
every actual neighbourhood
map and extension.


\item  The finite maximum partitions,
the ruled elliptic-curve
intersection alternatives,
the Euclidean multiplicity
lists and the Euler
contraction identities are
checked in
\path{code/A_minimal_reduction.py}
and its saved output.


\end{enumerate}

\item \par\medskip\noindent\textbf{All elliptic double curves force minimality in the zero case.}
Finally drop the component-minimality and no-triple
assumptions, retaining a full reduced SNC torus fibre
whose connected double curves are all elliptic.
The single-component case is the smooth two-torus
case of \ref{cdp:st:C.2.UN.5}. In the reducible case every component
has a nonempty boundary, and \eqref{cdp:eq:C.MR.3} remains the
actual adjunction formula even at triple points.


\begin{enumerate}[label={$\langle2\rangle$\arabic*.},leftmargin=1.3em,labelsep=.5em]
\item  No component can be birationally ruled over a
curve of genus at least two. Its map through
the actual contraction to that minimal ruled
surface and then to the base sends each
elliptic boundary to a point, by
Riemann--Hurwitz. A general ruling fibre
avoids those finitely many base values and
all point-blowup centres. Its strict
transform is a rational curve of
canonical degree $-2$ and has zero
intersection with every term of
\eqref{cdp:eq:C.MR.3}, a contradiction.


\item  Write $Z_v$ for the actual smooth minimal
model of $T_v$, and $n_v$ for the number
of point blowups in that contraction.
If its Kodaira dimension is nonnegative,
the degree-one comparison of \ref{cdp:st:C.2.UN.31} gives
a torus or the smooth-multiple elliptic
minimal model, so $e(Z_v)=0$. If it has
Kodaira dimension $-\infty$, the preceding
paragraph leaves rational, elliptic ruled
or class-VII minimal models. Their Euler
numbers are respectively strictly positive,
zero and zero; the last assertion uses
the actual Hopf restriction \ref{cdp:st:C.2.UN.19}.
Thus every $e(Z_v)$ is nonnegative,
and it is positive exactly for a
rational component. The actual
point-blowup comparisons give
$e(T_v)=e(Z_v)+n_v$.


\item  Let $t$ denote the number of triple
points. Topological inclusion-exclusion
on the actual SNC strata gives
\[
 \begin{aligned}
 e(S)&=\sum_v e(T_v)-\sum_C e(C)+t\\
     &=\sum_v e(Z_v)+\sum_v n_v+t
                                      \ge0.
 \end{aligned}
\tag{C.MR.16}\hypertarget{tag-00R8}{}\label{cdp:eq:C.MR.16}
\]
The triple term is retained until this
equality is used. If $e(S)=0$, all
nonnegative terms vanish: $t=0$,
every $n_v=0$, and no component is
rational. These are exactly the
initial hypotheses of the proved
reduction, with no higher-genus ruled
component. Conversely those same
conditions make every term in
\eqref{cdp:eq:C.MR.16} zero. This proves the
extension and its asserted exact
zero description.



\end{enumerate}

\end{enumerate}
\end{proof}






\begin{proposition}[A torsion-normal elliptic boundary cannot lie on a rational component]\hypertarget{tag-01ED}{}\label{cdp:st:C.2.UN.67}
\StatementRecord{cdp:st:C.2.UN.67}{proved}{01ED}{\SeriesRef{cdp:eq:C.RH.1}, \SeriesRef{cdp:eq:C.RH.4}, \SeriesRef{cdp:st:C.2.UN.16}, \SeriesRef{cdp:eq:C.RH.8}, \SeriesRef{cdp:eq:C.RH.7}, \SeriesRef{cdp:st:C.2.UN.37}, \SeriesRef{cdp:st:C.2.UN.50}, \SeriesRef{cdp:st:C.2.UN.53}; \cite[Theorem~6.2; Section~8(iv)(11), p.~585]{Kodaira1963II}; \citet[Section~2.3, pp.~5--6]{UlmerUrzua2021}}

Let \(f:U\to\Delta\) be a proper degeneration of connected
complex two-tori, with smooth total threefold and full
reduced SNC central fibre. Let \(P\) be a rational
component, with no minimality assumption. It cannot
have a smooth elliptic boundary curve \(J\) whose normal
line \(N_{J/P}\) is torsion and which is disjoint from
every other boundary curve on \(P\).

In particular a rational component, including a rational
elliptic surface and any point blowups of a rational
surface, cannot be adjacent to an elliptic Hopf component
or a nonlinear Hopf component. No Hopf leaf can have a
rational parent. The assertion permits rational double
curves and triple points elsewhere in the central fibre.
It imposes neither a zero-Euler hypothesis nor a
minimality hypothesis on the other components.

The obstruction is obtained on the actual component.
Its torsion normal produces the connected elliptic
pencil \(|mJ|\), where \(m=\operatorname{ord}N_{J/P}\).
The actual limiting translation fields force this
pencil to have identity integral monodromy. After its
vertical exceptional curves are contracted, it has
only smooth multiple fibres and coherent Euler
characteristic zero, contradicting
\(\chi(\OO_P)=1\).
\end{proposition}



\begin{proof}

\begin{enumerate}[label={$\langle1\rangle$\arabic*.},leftmargin=*]
\item \par\medskip\noindent\textbf{The exact pencil made by the torsion normal.}
Suppose that such a component and curve exist, and put
\(N=N_{J/P}\), with exact order \(m\ge1\).
The degree of \(N\) is zero, hence \(J^2=0\).
For every \(k\ge1\) use the actual divisor restriction
sequence
\[
 0\longrightarrow\OO_P((k-1)J)
 \longrightarrow\OO_P(kJ)
 \longrightarrow N^{\otimes k}\longrightarrow0.
 \tag{C.RH.1}\hypertarget{tag-00ZG}{}\label{cdp:eq:C.RH.1}
\]
For \(1\le k<m\), the last line is a nontrivial
degree-zero line on the elliptic curve. It has
both \(H^0\) and \(H^1\) zero: a nonzero section
would be nowhere zero, and Serre duality gives
the \(H^1\) assertion. Since \(P\) is rational,
\(H^1(P,\OO_P)=0\) and \(H^0(P,\OO_P)=\C\).
Induction in the actual long exact sequences gives
\[
 h^0(P,\OO_P(kJ))=1,\qquad
 H^1(P,\OO_P(kJ))=0\quad(0\le k<m).
 \tag{C.RH.2}\hypertarget{tag-00ZI}{}\label{cdp:eq:C.RH.2}
\]
At \(k=m\), the restriction in \eqref{cdp:eq:C.RH.1} is
onto \(H^0(J,\OO_J)=\C\). Thus
\[
 h^0(P,\OO_P(mJ))=2,\qquad
 H^0(P,\OO_P(mJ))\longrightarrow H^0(J,\OO_J)
                         \text{ is onto}.
 \tag{C.RH.3}\hypertarget{tag-00ZJ}{}\label{cdp:eq:C.RH.3}
\]
One section has divisor \(mJ\); a second restricts
to a nonzero constant on \(J\). Their common
zero locus is empty: on \(J\) the second
does not vanish, and off \(J\) the first
does not vanish. They define the actual
holomorphic map
\[
                   q_0:P\longrightarrow\PP^1,\qquad
                             q_0^{-1}(\infty)=mJ .
 \tag{C.RH.4}\hypertarget{tag-00ZK}{}\label{cdp:eq:C.RH.4}
\]

\item We check its connectedness through the Stein map,
rather than assuming that a pencil has connected
fibres. Write \(P\xrightarrow q B\xrightarrow\nu\PP^1\)
for its Stein factor, and let \(d=\deg\nu\).
The injection \(H^1(B,\OO_B)\to H^1(P,\OO_P)\)
gives \(B\simeq\PP^1\).
Since the support of \eqref{cdp:eq:C.RH.4} is connected,
\(\nu^{-1}(\infty)\) has just one point.
It is totally ramified of order \(d\), and
the actual fibre of \(q\) there is
\((m/d)J\). Its principal local equation
trivializes \(\OO_P((m/d)J)|_J\), so
\[
                   N^{\otimes(m/d)}=\OO_J.
 \tag{C.RH.5}\hypertarget{tag-00ZL}{}\label{cdp:eq:C.RH.5}
\]
Exact order \(m\) forces \(d=1\).
Thus \(q=q_0\) has connected fibres.
Adjunction gives \(K_P\cdot J=0\);
a general member \(F\sim mJ\) is smooth
by generic smoothness and satisfies
\[
                 2g(F)-2=F\cdot(K_P+F)=0.
 \tag{C.RH.6}\hypertarget{tag-00ZM}{}\label{cdp:eq:C.RH.6}
\]
It is a smooth connected elliptic curve.
This constructs the required elliptic
map on the original nonminimal surface.

\item Every other boundary component \(C\) on
\(P\) is disjoint from \(J\) by hypothesis.
Consequently \(C\cdot q^*(\infty)=mC\cdot J=0\).
If \(q|_C\) were nonconstant, that intersection
would be its positive degree. Hence every
boundary is vertical for this very map \(q\).

\par\medskip\noindent\textbf{The actual translation fields align this pencil.}
The component has full multiplicity one.
The saturated relative translation module
of \ref{cdp:st:C.2.UN.16} supplies two commuting meromorphic
fields \(V_1,V_2\) on \(P\), independent over
\(\C\), all of whose constant combinations
are semicomplete. Their polar divisors
are contained in the boundary of \(P\).
Indeed they are obtained from holomorphic
ambient relative fields by meromorphic
base coefficients; cancelling the generic
order on \(P\) leaves a holomorphic
restriction at every point of \(P\)
away from the other central components.
This is the original surface and its
actual fields, not a replacement cover.

\item Remove the finitely many base values
supporting a boundary, a singular fibre
or a vertical zero of the relevant field.
On each remaining compact elliptic fibre
the fields are holomorphic. Their base
coefficients are holomorphic functions
on that fibre and hence are constant.
The connected-fibre map therefore
descends them to commuting meromorphic
fields on \(B=\PP^1\).

\item These projected fields have no poles.
For clarity, at a pole of order \(a>0\),
the primitive of the one-dimensional
time form has a zero of order \(a+1\).
Continue one inverse branch around
its corresponding time value. Choose
the small loop so that its lifted
base paths avoid every other bad
value. Their compact union has a
compact inverse image under the proper
map \(q\), contained in the holomorphic
locus of the upstairs field. The
upstairs solution therefore continues
along the entire loop. Its final
base value is a different inverse
branch, contradicting the
single-valued maximal solution of
a semicomplete field. This is the
actual continuation map used in
(C.RF.1) and (C.MR.8); it does
not require a minimal surface.

\item If one projection is nonzero,
the other is a constant multiple,
because the derivative of their
meromorphic ratio along a nonzero
commuting curve field is zero.
Subtract that constant multiple.
The resulting nonzero vertical
field is a translation field on
each compact regular elliptic fibre.
Its commuting transverse field
identifies neighbouring fibres and
preserves its time form. Its two
marked integral periods are
therefore locally constant and,
by analytic continuation on the
connected smooth base, globally
single-valued. Integration of a
nonzero elliptic time form is
injective on \(H_1(F,\Z)\);
the integral monodromy is the
identity.

\item If both projections vanish,
the fields are vertical translation
fields on each regular fibre.
Their meromorphic ratio is a
nonconstant function pulled
back from \(B\), since the two
fields are independent over
\(\C\). Its connected Stein
fibres are exactly those of
the map \(q\) just constructed.
The actual ambient return germs
(C.RG.4)--(C.RG.12) apply to
this same first integral:
\(\lambda_i=d_i-k_iH\) for its
two continued integral periods.
The right-hand side is
single-valued, and injectivity
of elliptic period integration
again gives identity integral
monodromy. No elliptic pencil
on an ambient neighbourhood
has been assumed.

\par\medskip\noindent\textbf{Relative minimization and the exact singular-fibre comparison.}
Contract only the vertical
\((-1)\)-curves of \(q\).
Each contraction is the
actual smooth surface
blowdown of a rational
exceptional curve, and \(q\)
descends because that curve
is one fibre of the blowdown
and maps to a point of \(B\).
There are finitely many such
steps, since each lowers
\(b_2\). The resulting map
\[
               q_{\min}:P_{\min}\longrightarrow\PP^1
 \tag{C.RH.7}\hypertarget{tag-00ZN}{}\label{cdp:eq:C.RH.7}
\]
is relatively minimal. It
has the same marked smooth
fibres and the same identity
monodromy. Moreover
\(\chi(\OO_{P_{\min}})=\chi(\OO_P)=1\),
by the direct-image equality
for every ordinary point
blowup.

\item We now use the relatively
minimal singular-fibre
classification, precisely in
its stated scope:
\cite[Theorem~6.2; Section~8(iv)(11), p.~585]{Kodaira1963II} and \citet[Section~2.3, pp.~5--6]{UlmerUrzua2021}.
Its nonmultiple non-smooth
types have monodromies
conjugate to nontrivial
unipotent matrices, their
negatives, or the nonidentity
finite-order matrices of
orders \(2,3,4,6\).
Only the smooth type \(I_0\)
has identity monodromy.

\item For a full multiple fibre
write its local function as
\(t=u h^a\), with \(u\) a
unit and the divisor of
\(h\) primitive. Normalize
the actual base change
\(t=s^a\) by
\[
                         w^a=u,\qquad s=wh .
 \tag{C.RH.8}\hypertarget{tag-00ZO}{}\label{cdp:eq:C.RH.8}
\]
The resulting surface map
is \'{e}tale over the
original surface, because
\(w\) is a root of a unit.
Its monodromy is the
\(a\)-th power of the
original one and is
therefore the identity.
It is still relatively
minimal by the following
canonical intersection
comparison. For every
irreducible component \(C\)
of a fibre \(F\) below,
\(F\cdot C=0\) gives
\(C^2\le0\). Adjunction
gives
\(K_{P_{\min}}\cdot C
=2p_a(C)-2-C^2\ge0\),
except possibly for a
smooth rational curve of
square zero or minus one.
The latter is excluded
by relative minimality.
In the square-zero case,
all intersections with
the other fibre components
vanish. Connectedness
would give \(F=bC\),
contradicting
\(K_{P_{\min}}\cdot F=0\).
Thus every component has
nonnegative canonical
intersection. Their
positive weighted sum
is \(K_{P_{\min}}\cdot F=0\),
so each intersection
is zero. For the actual
\'{e}tale covering map
\(\pi\) in \eqref{cdp:eq:C.RH.8},
\(K_{\rm up}=\pi^*K_{\rm down}\).
Any vertical \((-1)\)-curve
\(E\) upstairs would then
satisfy
\(K_{\rm up}\cdot E
=\deg(E\to C)K_{\rm down}\cdot C=0\),
contrary to adjunction,
which gives minus one.
This uses the cycle
pushforward and does
not assume that the
image curve is embedded
or nonsingular.
Its primitive Kodaira
fibre must consequently
be \(I_0\). Thus every
fibre of \(q_{\min}\)
is a positive multiple
of a smooth elliptic
curve.

\item The oriented period ratio
is now a single-valued
holomorphic map from a
punctured \(\PP^1\) to
the upper half-plane.
Its bounded Cayley
transform extends across
the missing points and
is constant on \(\PP^1\).
All regular fibres have
one fixed Jacobian.
Their transition maps
are translations, since
the integral monodromy
is the identity. At a
multiple value use
\eqref{cdp:eq:C.RH.8}. The free
central elliptic inertia
is a translation, and
its action on \(H^1(\OO)\)
is trivial. Taking the
actual invariant coefficient
functions gives
\[
 (q_{\min})_*\OO_{P_{\min}}=\OO_{\PP^1},
 \qquad R^1(q_{\min})_*\OO_{P_{\min}}=\OO_{\PP^1}.
 \tag{C.RH.9}\hypertarget{tag-00ZP}{}\label{cdp:eq:C.RH.9}
\]
The coherent Leray
sequence now yields
\(\chi(\OO_{P_{\min}})=1-1=0\),
contradicting the preserved
value one in \eqref{cdp:eq:C.RH.7}.
This proves the obstruction
without applying a
minimal-only assertion
to the original \(P\).

\par\medskip\noindent\textbf{The Hopf interface.}
On an elliptic Hopf
component every curve
is a fibre support and
has normal of finite
exact order, by the
actual normalized-root
map in \ref{cdp:st:C.2.UN.37}. A nonlinear
Hopf has its one axis,
whose normal has exact
resonance order, including
the secondary comparison
of \ref{cdp:st:C.2.UN.50}. Distinct Hopf
curves are disjoint.
There can therefore be
no triple point on an
interface with a Hopf
component: such a point
would make two different
boundary curves meet
inside that Hopf surface.
The actual \(xy=t\)
normal-product map gives
\[
                     N_{J/P}\otimes N_{J/H}=\OO_J.
 \tag{C.RH.10}\hypertarget{tag-00ZH}{}\label{cdp:eq:C.RH.10}
\]
It makes \(N_{J/P}\)
torsion, and the same
absence of triple points
makes \(J\) disjoint
from every other
boundary curve on \(P\).
The obstruction just
proved applies.

\item Every terminal Hopf
pair has one of these
torsion-normal axes,
by the exact
one-boundary canonical
and character comparison
in \ref{cdp:st:C.2.UN.53}. The two
nontorsion axes of
an algebraic-dimension-zero
diagonal Hopf cannot
form such a leaf.
This proves all the
asserted parent and
interface exclusions.

\end{enumerate}
\end{proof}






\begin{proposition}[Rational toric component pairs and the exact dual torus]\hypertarget{tag-01EE}{}\label{cdp:st:C.2.UN.68}
\StatementRecord{cdp:st:C.2.UN.68}{proved}{01EE}{\SeriesRef{cdp:st:C.2.UN.65}, \SeriesRef{cdp:eq:C.TC.3}, \SeriesRef{cdp:st:C.2.UN.19}, \SeriesRef{cdp:eq:C.TC.2}, \SeriesRef{cdp:eq:C.TC.5}, \SeriesRef{cdp:eq:C.TC.6}, \SeriesRef{cdp:eq:C.TC.9}, \SeriesRef{cdp:eq:C.TC.8}, \SeriesRef{cdp:eq:C.TC.11}, \SeriesRef{cdp:eq:C.TC.12}; \cite[Definition 1.2 and Proposition 1.3]{GHK2015}}

Let $f:X\to\Delta$ have the full reduced simple-normal-crossing
central fibre $S=\sum_{i=1}^{v}V_i$ of \ref{cdp:st:C.2.UN.65}, and let its
canonical evaluation orders be equal. Suppose that $e(S)>0$.
Then every $V_i$ is rational and its actual boundary
$D_i=\sum_{j\ne i}V_i\cap V_j$ is the full toric boundary
for an actual algebraic two-torus action on $V_i$.
The dual complex of this given fibre is a closed real
two-torus. If $r$ counts its connected double curves and
$t$ its triple points, then
\[
                  r=3v,\qquad t=2v,\qquad e(S)=2v.
\tag{C.TC.1}\hypertarget{tag-012C}{}\label{cdp:eq:C.TC.1}
\]
In particular $v\ge3$ and $e(S)\ge6$ in the embedded
simple-normal-crossing setting of the statement.

Together with \ref{cdp:st:C.2.UN.65} this identifies every component
pair and the dual topological space of a canonical-flat
torus degeneration. The assertion concerns the actual
given degeneration and its actual component pairs.
No smoothing of arbitrary proposed toric gluing data
is inferred from these necessary comparisons.
\end{proposition}



\begin{proof}

\begin{enumerate}[label={$\langle1\rangle$\arabic*.},leftmargin=*]
\item \par\medskip\noindent\textbf{A rational boundary cycle forces a rational component.}
Use the actual three-form $\Omega/f^k$ of \ref{cdp:st:C.2.UN.65}.
It gives
\[
               K_{V_i}+D_i=0,\qquad
               K_C=\OO_C(-T_C)
\tag{C.TC.2}\hypertarget{tag-012G}{}\label{cdp:eq:C.TC.2}
\]
by its component residue and the principal-fibre line maps.
Thus a boundary component is either elliptic and disjoint
from the remaining boundary, or rational and meets it
at exactly two transverse points. Consequently the
rational boundary components form disjoint nodal cycles.
The positive Euler identity $e(S)=t$ of \ref{cdp:st:C.2.UN.65} supplies
at least one such cycle.

\item We check how an actual point blowdown on a component
acts on this divisor. Let $E$ be a smooth rational
$(-1)$-curve and $\pi:V\to V'$ its ordinary blowdown.
Adjunction gives
\[
                         D\cdot E=1.
\tag{C.TC.3}\hypertarget{tag-012H}{}\label{cdp:eq:C.TC.3}
\]
If $E$ is not in $D$, it meets exactly one smooth
boundary point transversely. It cannot meet a node,
which would contribute at least two to \eqref{cdp:eq:C.TC.3}.
The image $D'=\pi_*D$ is again reduced nodal and
$K_{V'}+D'=0$. Locally this is the ordinary blowup
at a smooth point of $D'$ with $D$ its strict transform.

\item If $E$ is in $D$, then
$(D-E)\cdot E=2$. Its two intersection points are
distinct and transverse. The blowdown identifies
their two transverse tangent directions and creates
one ordinary node on $D'$. The two branches may
belong to the same irreducible curve. This is the
actual inverse of the blowup of a boundary node;
again $K_{V'}+D'=0$ by the ordinary canonical
Jacobian formula. In particular a rational cycle
persists. When it has become an irreducible nodal
rational curve, it cannot be the smooth exceptional
curve of a further ordinary point blowdown.

\item These two maps preserve the separation between
the rational cycles and the elliptic boundary
components. Contracting the point blowups to
a minimal model therefore retains a rational
cycle whenever one was present. Each component
has $\kappa=-\infty$, because $-K=D$ is
nonempty effective. A nonrational ruled minimal
model has rational curves only in its ruling
fibres: a map from $\PP^1$ to the positive-genus
base is constant. Those fibres are mutually
disjoint smooth rational curves and cannot
contain a nodal rational cycle. The other
nonrational possibility is the Hopf minimal
model of the actual class-VII component,
by \ref{cdp:st:C.2.UN.19}. A rational curve on a Hopf surface
would lift to $\C^2\setminus\{0\}$, where
its coordinate functions are constant. It
has no rational cycle either. Hence a
component whose boundary contains a rational
cycle is rational.

\item On a rational component the exact sequence
\[
 0\longrightarrow K_V\longrightarrow\OO_V
                 \longrightarrow\OO_D\longrightarrow0
\tag{C.TC.4}\hypertarget{tag-012I}{}\label{cdp:eq:C.TC.4}
\]
has $H^0(K_V)=H^1(K_V)=0$, the second vanishing
being Serre dual to $H^1(\OO_V)=0$.
It follows that $H^0(\OO_D)=\C$, so $D$ is
connected. If $D$ contains a rational curve,
it is therefore exactly one rational nodal
cycle. A neighbouring component contains
that same rational double curve and hence
is rational by the preceding argument.
Propagation through the connected fibre
shows that every $V_i$ is rational and
every $D_i$ is one rational cycle.

\par\medskip\noindent\textbf{The actual dual complex and its cohomology.}
Let $\mathcal B$ be the dual complex with
one vertex per $V_i$, one edge per connected
double curve and one triangle per triple
point. Multiple edges and triangles are
retained. Every edge belongs to two
triangles by \eqref{cdp:eq:C.TC.2}, and the link of
each vertex is exactly the boundary cycle
on that component. Thus the actual
topological realization of $\mathcal B$
is a connected closed real surface.

\item All positive-degree structure-sheaf
cohomology of the rational components
and rational double curves vanishes.
The SNC restriction sequence therefore
identifies, by its alternating restriction
maps, the full complex computing
$H^*(S,\OO_S)$ with
\[
           \C^{v}\longrightarrow\C^{r}
                                  \longrightarrow\C^{t},
\tag{C.TC.5}\hypertarget{tag-012J}{}\label{cdp:eq:C.TC.5}
\]
the cellular cochain complex of
$\mathcal B$. The principal-fibre
adjunction and the actual volume form
give $\omega_S\simeq\OO_S$.
Serre duality on the compact Gorenstein
surface $S$ consequently gives
\[
             H^2(S,\OO_S)\simeq
                         H^0(S,\OO_S)^*\simeq\C.
\tag{C.TC.6}\hypertarget{tag-012K}{}\label{cdp:eq:C.TC.6}
\]
Here $H^0(S,\OO_S)=\C$ follows from
compactness, connectedness and reducedness.
Equations \eqref{cdp:eq:C.TC.5}--\eqref{cdp:eq:C.TC.6} make
$\mathcal B$ orientable. Proper flatness
gives $\chi(\OO_S)=0$, hence
$\chi(\mathcal B)=0$. A connected closed
oriented surface with this Euler
characteristic is a real two-torus.

\item The same cochain counts give
$v-r+t=0$. Together with the exact
incidence count $2r=3t$ of \ref{cdp:st:C.2.UN.65},
this proves \eqref{cdp:eq:C.TC.1}. A triple point
in a simple-NC divisor lies on three
distinct embedded components, so
$t>0$ implies $v\ge3$.

\par\medskip\noindent\textbf{An elementary charge test with its actual equality morphisms.}
For a smooth rational surface with an
anticanonical nodal rational cycle
$D=D_1+\cdots+D_n$, allowing an
irreducible nodal rational curve when
$n=1$, put
\[
                        Q(V,D)=e(V)-n=e(V\setminus D).
\tag{C.TC.7}\hypertarget{tag-012L}{}\label{cdp:eq:C.TC.7}
\]
Indeed a rational nodal cycle has
Euler characteristic $n$, including
the one-component case. We prove
\[
 Q(V,D)\ge0,\qquad Q(V,D)=0
       \quad\Longleftrightarrow\quad
       (V,D)\text{ is an actual toric pair}.
\tag{C.TC.8}\hypertarget{tag-012M}{}\label{cdp:eq:C.TC.8}
\]

\item Follow any sequence of ordinary
$(-1)$-curve contractions to a
minimal rational surface. The two
maps already proved after \eqref{cdp:eq:C.TC.3}
keep an anticanonical nodal cycle.
At a contraction outside the boundary,
$e(V)$ drops by one and $n$ is
unchanged. At a boundary contraction,
both drop by one. Thus
\[
 Q(V,D)=Q(V_{\min},D_{\min})+a,
\tag{C.TC.9}\hypertarget{tag-012N}{}\label{cdp:eq:C.TC.9}
\]
where $a$ is the number of actual
contracted exceptional curves outside
the current boundary.

\item On $\PP^2$ the boundary has total
degree three, so $n\le3=e(\PP^2)$.
Equality means three distinct lines;
their nodal cycle condition makes
them nonconcurrent. The invertible
matrix of their linear equations is
an actual projective transformation
to the coordinate triangle.

\item On a minimal rational ruled surface
$q:\mathbb F_b\to\PP^1$, the total
horizontal degree of $D$ is two.
There are at most two horizontal
components. Each vertical prime is
a complete ruling fibre and meets
the horizontal part in two points
counted with multiplicity. The cycle
has only two incidences at each
horizontal component, so there are
at most two vertical components.
Hence $n\le4=e(\mathbb F_b)$.
Equality forces two degree-one
sections and two ruling fibres,
with the sections disjoint.
Their two quotient line maps split
the ruling bundle. A base coordinate
puts the two marked fibres at zero
and infinity. The resulting actual
pair is
\[
 \left(\PP_{\PP^1}(\OO\oplus\OO(b)),
        \text{two sections plus the fibres over }0,\infty\right).
\tag{C.TC.10}\hypertarget{tag-012D}{}\label{cdp:eq:C.TC.10}
\]
In a line coordinate the transitions
are $w_\infty=z^{-b}w_0$.
The maps $(z,w_0)\mapsto(\lambda z,\mu w_0)$
and $(z^{-1},w_\infty)\mapsto
(\lambda^{-1}z^{-1},\lambda^{-b}\mu w_\infty)$
give its actual two-torus action
and precisely the four specified
boundary components.

\item The two minimal computations prove
$Q(V_{\min},D_{\min})\ge0$.
If the initial charge is zero,
\eqref{cdp:eq:C.TC.9} says that $a=0$ and
the minimal pair is one of the
two explicit toric pairs just
identified. Reversing the recorded
contractions therefore blows up
only boundary nodes. Each such
node is fixed by the torus action;
the invariant ideal $(x,y)$ gives
the actual lifted action on its
blowup charts $y=xu$ and $x=yv$.
The new full boundary is the
toric boundary. Induction gives
an actual torus action on the
original marked pair $(V,D)$.
Conversely a smooth complete
toric surface has one fixed
point per boundary node; its
orbit decomposition gives
$e(V)=n$, proving the reverse
implication in \eqref{cdp:eq:C.TC.8}.

\item This direct argument agrees with
the actual toric-model comparison
of \cite[Definition 1.2 and Proposition 1.3]{GHK2015},
whose hypotheses are exactly a
smooth rational projective surface
and an anticanonical nodal rational
cycle. The displayed contractions
prove the equality case without
discarding the original pair.

\par\medskip\noindent\textbf{The global charge is zero.}
Write $n_i$ for the number of
curves in $D_i$. Every double
curve belongs to two component
boundaries, so
\[
                         \sum_i n_i=2r=3t.
\tag{C.TC.11}\hypertarget{tag-012E}{}\label{cdp:eq:C.TC.11}
\]
The exact HRR and boundary-square
calculation in \ref{cdp:st:C.2.UN.65} gives
\[
            \sum_i e(V_i)
             =12(r-t)-3t=3t.
\tag{C.TC.12}\hypertarget{tag-012F}{}\label{cdp:eq:C.TC.12}
\]
Subtracting \eqref{cdp:eq:C.TC.11} from
\eqref{cdp:eq:C.TC.12} yields
$\sum_i Q(V_i,D_i)=0$.
Every summand is nonnegative by
\eqref{cdp:eq:C.TC.8}, so each is zero.
The actual equality morphisms
above make every specified
$(V_i,D_i)$ a toric pair.
This proves the asserted
classification on the given
degeneration.

\end{enumerate}
\end{proof}






\section{Compact occurrence of prescribed component classes}\hypertarget{tag-0345}{}\label{sec:II:09}


\begin{proposition}[Zero-Euler product fillings with essential Kodaira-dimension-one components]\hypertarget{tag-01E5}{}\label{cdp:st:C.2.UN.6}
\StatementRecord{cdp:st:C.2.UN.6}{proved}{01E5}{\SeriesRef{cdp:eq:C.ER.4}, \SeriesRef{cdp:eq:C.ER.5}, \SeriesRef{cdp:eq:C.ER.7}, \SeriesRef{cdp:eq:C.ER.8}, \SeriesRef{cdp:eq:C.ER.14}, \SeriesRef{cdp:eq:C.ER.9}, \SeriesRef{cdp:eq:C.ER.3}}

Put $q=e^{-2\pi}$ and $E=\C^*/q^{\Z}\simeq\C/(\Z+\mathrm i\Z)$.
For every integer $m\ge2$, the fixed punctured product
$(E\times E)\times\Delta^*$ has a proper filling
$f_m:U_m\to\Delta$ with smooth total space and full central fibre
\[
                    (U_m)_0=P_m+mQ_m.             \tag{C.ER.1}\hypertarget{tag-00IN}{}\label{cdp:eq:C.ER.1}
\]
Here $Q_m$ is a primary Hopf surface, and $P_m$ is a smooth compact
non-K\"ahler elliptic surface over $E$ with precisely one multiple
fibre $mJ$. The surfaces meet transversely along the smooth elliptic
curve $J$. One has
\[
 \begin{gathered}
 e((U_m)_0{}_{\rm red})=0,\qquad H_1(P_m,\Z)=\Z^3,\\
 K_{P_m}\simeq\OO_{P_m}((m-1)J),\qquad \kappa(P_m)=1,\\
 h^0(P_m,K_{P_m}^{\otimes r})=
 \begin{cases}
 1,&d_r=0,\\
 d_r,&d_r>0,
 \end{cases}
 \qquad d_r=\left\lfloor\frac{r(m-1)}m\right\rfloor,
 \quad r\ge1.
 \end{gathered}                                   \tag{C.ER.2}\hypertarget{tag-00IX}{}\label{cdp:eq:C.ER.2}
\]
Neither $P_m$ nor $Q_m$ can be contracted by a bimeromorphic map
between smooth threefold fillings. The $U_m$ are pairwise
non-bimeromorphic over the disc. After every finite base change and
resolution of its normalization, a central component remains
birational to $P_m$. Thus a properly elliptic component of Kodaira
dimension one is forced to survive in these examples, including
after finite base change. All these fillings have compact smooth
threefold realizations over $\PP^1$.
\end{proposition}


\begin{proof}

\begin{enumerate}[label={$\langle1\rangle$\arabic*.},leftmargin=*]
\item We give the construction as an actual hypersurface and a free
quotient. Let
\[
 V=\operatorname{Bl}_{(0,0)}(E\times\Delta),\quad
 R=\operatorname{Exc}(V\to E\times\Delta),\quad
 L=\OO_V(R),\quad \sigma\in H^0(V,L)
                                                        \tag{C.ER.3}\hypertarget{tag-00IY}{}\label{cdp:eq:C.ER.3}
\]
with $\sigma$ the canonical section. Define
\[
 W_m=\{(v,\ell,a):v\in V,\ \ell\in L_v^\times,\ a\in\C,
                         \ \sigma(v)=a^m\ell\}.
                                                        \tag{C.ER.4}\hypertarget{tag-00IZ}{}\label{cdp:eq:C.ER.4}
\]
The generator of a $\Z$-action is
\[
             (v,\ell,a)\longmapsto(v,q^m\ell,q^{-1}a),
 \qquad U_m=W_m/\Z.                               \tag{C.ER.5}\hypertarget{tag-00J0}{}\label{cdp:eq:C.ER.5}
\]
The action preserves \eqref{cdp:eq:C.ER.4}. Near $R$, choose a holomorphic frame
of $L$ in which $\sigma=xe$ and write $\ell=be$, $b\ne0$.
Then \eqref{cdp:eq:C.ER.4} is $x=a^m b$ and eliminates the coordinate $x$;
thus $W_m$ is smooth. It is also smooth off $R$, where $a\ne0$
and $\ell=\sigma/a^m$. A Hermitian norm on $L$ is multiplied by
$q^m$ under \eqref{cdp:eq:C.ER.5}. It is bounded above and bounded away from zero
on each compact subset of $W_m$, proving proper discontinuity.
No nonzero power fixes a point, since $\ell\ne0$. Consequently
$U_m$ is a smooth complex threefold.

\item We verify properness and the local maps explicitly. With $y$ a
second base coordinate, the map
\[
 (z,w,y)\longmapsto
 (x=z^m,y;\ b=w^m,\ a=z/w),\qquad w\in\C^*,
                                                        \tag{C.ER.6}\hypertarget{tag-00J1}{}\label{cdp:eq:C.ER.6}
\]
identifies the local $W_m$ with the quotient by
$(z,w,y)\mapsto(\zeta z,\zeta w,y)$, $\zeta^m=1$.
It is onto because an $m$th root of $b$ determines $w$ and then
$z=aw$; its fibres are exactly the displayed orbits. The action is
free. The generator \eqref{cdp:eq:C.ER.5} lifts to $(z,w,y)\mapsto(z,qw,y)$.
Hence the local projection $\rho_m:U_m\to V$ is exactly
\[
 (\Delta_z\times E\times\Delta_y)/\boldsymbol\mu_m
        \longrightarrow\Delta_x\times\Delta_y,
 \quad [z,[w],y]\longmapsto(z^m,y),
 \quad \zeta:(z,[w],y)\longmapsto(\zeta z,[\zeta w],y).
                                                        \tag{C.ER.7}\hypertarget{tag-00J2}{}\label{cdp:eq:C.ER.7}
\]
The last finite action is free: $[\zeta w]=[w]$ would imply
$\zeta=q^j$ for an integer $j$, forcing $j=0$ and $\zeta=1$.
The map before this finite quotient is the product of a finite
disc map with a compact elliptic curve. It is proper, so \eqref{cdp:eq:C.ER.7}
is proper. These charts prove that $\rho_m$ is proper everywhere.
Its fibres are connected elliptic curves, including the reduced
fibre over $R$. Thus composition with $V\to\Delta$ is proper
with connected fibres.

\item The comparison with the punctured product is particularly simple.
For $v\notin R$ and $c\in\C^*$ set
\[
       (v,[c])\longmapsto[(v,\sigma(v)c^m,c^{-1})].
                                                        \tag{C.ER.8}\hypertarget{tag-00J3}{}\label{cdp:eq:C.ER.8}
\]
Changing $c$ to $qc$ changes the representative by \eqref{cdp:eq:C.ER.5}.
Conversely $[c]=[a^{-1}]$ is its inverse. Thus \eqref{cdp:eq:C.ER.8} is an
actual isomorphism $(V\setminus R)\times E\simeq
\rho_m^{-1}(V\setminus R)$, for every $m$, and over $\Delta^*$
it gives the specified fixed product $(E\times E)\times\Delta^*$.

\item Let $E_0\simeq E$ be the strict transform of $E\times\{0\}$.
It meets $R\simeq\PP^1$ at one point, and
$L|_{E_0}=\OO_E([0])$, $L|_R=\OO_{\PP^1}(-1)$.
The inverse image of $R$, with its reduction, is $a=0$ in
\eqref{cdp:eq:C.ER.4}, so it is
\[
 Q_m=L|_R^\times/(q^m)^{\Z}
     \simeq(\C^2\setminus\{0\})/\langle q^m\operatorname{id}\rangle.
                                                        \tag{C.ER.9}\hypertarget{tag-00J4}{}\label{cdp:eq:C.ER.9}
\]
The first isomorphism sends the nonzero tautological vector to
the same vector in $\C^2\setminus\{0\}$. It identifies $Q_m$
as a primary Hopf surface. The reduced inverse image $P_m$ of
$E_0$ is smooth by \eqref{cdp:eq:C.ER.7}, since $E_0$ meets $R$ transversely.
Its map $p:P_m\to E_0$ has smooth elliptic fibres off $[0]$ and
one full fibre $mJ$ over $[0]$. The reduced fibre is
\[
                    J=\C^*/(q^m)^{\Z}.
                                                        \tag{C.ER.10}\hypertarget{tag-00IO}{}\label{cdp:eq:C.ER.10}
\]
For example the isomorphism $E/\boldsymbol\mu_m\to J$ in
\eqref{cdp:eq:C.ER.7} is $[w]\mapsto[w^m]$, which also specifies the finite
covering between the two elliptic curves.

\item Near $E_0\cap R$, choose coordinates with $R=(x=0)$,
$E_0=(y=0)$ and base map $t=xy$. In the smooth chart before
the free quotient the equation is $t=z^m y$; in \eqref{cdp:eq:C.ER.4} it is
\[
                         t=a^m b y,\qquad b\ne0.
                                                        \tag{C.ER.11}\hypertarget{tag-00IP}{}\label{cdp:eq:C.ER.11}
\]
It follows that the full central fibre is exactly $P_m+mQ_m$,
and its reduction has the two transverse components stated in
\eqref{cdp:eq:C.ER.1}. This calculation also retains the ramification at $J$.

\item We compute the canonical line without a classification assertion.
Choose a nowhere-zero holomorphic one-form $\omega$ on $E_0$.
On a local cyclic chart for $P_m$ as in \eqref{cdp:eq:C.ER.7}, the form
\[
                         p^*\omega\wedge\frac{dw}{w}
                                                        \tag{C.ER.12}\hypertarget{tag-00IQ}{}\label{cdp:eq:C.ER.12}
\]
is invariant under both $w\mapsto qw$ and the finite action.
Changing a local frame changes $dw/w$ by a form from the
one-dimensional base, which has zero wedge with $p^*\omega$.
It therefore glues to a global holomorphic two-form on $P_m$.
If $\omega=h(x)\,dx$ near $[0]$, its pullback is
\[
 m h(z^m)z^{m-1}\,dz\wedge\frac{dw}{w}.
                                                        \tag{C.ER.13}\hypertarget{tag-00IR}{}\label{cdp:eq:C.ER.13}
\]
It has precisely the divisor $(m-1)J$, proving the canonical
isomorphism in \eqref{cdp:eq:C.ER.2}.

\item For every integer $k\ge0$ there is a precise pushforward equality
\[
        p_*\OO_{P_m}(kJ)
          =\OO_E\bigl(\lfloor k/m\rfloor[0]\bigr).
                                                        \tag{C.ER.14}\hypertarget{tag-00IS}{}\label{cdp:eq:C.ER.14}
\]
A section on the left is a meromorphic function whose only allowed
poles are along $J$. Off $[0]$ it is holomorphic on every compact
elliptic fibre, so it is constant on that fibre and descends to a
holomorphic function on the base. On the cyclic chart at $[0]$,
its Laurent coefficients in $z$ are holomorphic functions on the
compact elliptic curve and hence constant. Finite-group invariance
forces the powers of $z$ to be multiples of $m$. The pole bound
$k$ is consequently exactly a pole bound $\lfloor k/m\rfloor$
in $x=z^m$. This proves both inclusions in \eqref{cdp:eq:C.ER.14}.
Applying it to $k=r(m-1)$ and applying Riemann--Roch on the elliptic
curve proves the displayed plurigenera. Their growth is linear
with positive slope $(m-1)/m$, so $\kappa(P_m)=1$.

\item Here is also an integral topological check that this surface is
non-K\"ahler. Choose a Hermitian metric on $L|_{E_0}$ and put
\[
 M_m=\{(v,\ell,a)\in W_m|_{E_0}:\|\ell\|=1\}.
                                                        \tag{C.ER.15}\hypertarget{tag-00IT}{}\label{cdp:eq:C.ER.15}
\]
The maps $\ell=e^{mu}\ell_0$, $a=e^{-u}a_0$ identify
$W_m|_{E_0}=M_m\times\R$ smoothly. The generator \eqref{cdp:eq:C.ER.5}
translates $u$ by $\log q$, so
$P_m\simeq M_m\times S^1$ smoothly. Over the complement of a
small disc about $[0]$, the canonical section trivializes $M_m$
as the product with a circle, whose coordinate is the phase of $a$.
On the disc choose the metric to be constant in the frame used
above. Then $M_m$ has the solid-torus coordinates
$a=\rho e^{\mathrm i\theta}$, $b=e^{\mathrm i\varphi}$, and
\[
              x=a^m b,\qquad \arg x=m\theta+\varphi.
                                                        \tag{C.ER.16}\hypertarget{tag-00IU}{}\label{cdp:eq:C.ER.16}
\]
Let $\delta$ denote the boundary of the punctured elliptic base
and let $h$ denote its product circle. The meridian of this solid
torus maps to $m\delta+h$; its longitude maps to $\delta$.
The class $\delta$ is zero in the first homology of the once-punctured
elliptic curve. Mayer--Vietoris therefore kills $h$ and the solid-torus
longitude, leaving exactly the two free base classes. All spaces
and their boundary torus are connected, so the degree-zero part
adds no further class. Hence
$H_1(M_m,\Z)=\Z^2$ and $H_1(P_m,\Z)=\Z^3$.
The odd first Betti number excludes a K\"ahler metric, by Hodge
decomposition on a compact K\"ahler manifold.

\item The Euler calculation retains the multiple fibre. Over
$E_0\setminus\{[0]\}$, $P_m$ is locally a product with an elliptic
curve; its remaining reduced fibre $J$ is also elliptic.
Stratified Euler additivity gives $e(P_m)=0$.
The Hopf surface in \eqref{cdp:eq:C.ER.9} is a principal elliptic bundle over
$\PP^1$, so $e(Q_m)=0$. Finally $e(J)=0$, and consequently
$e(P_m\cup_JQ_m)=0$.

\item Every map $\PP^1\to P_m$ is constant: it first maps constantly
to the elliptic base $E_0$, and then constantly to the reduced
elliptic fibre containing its image. Every map $\PP^1\to Q_m$
lifts to $\C^2\setminus\{0\}$ and is constant there by its
coordinate functions. Thus neither component is covered by
rational curves. The inverse-map resolution in (C.UC.1) applies
and proves that neither can be contracted by a bimeromorphic
comparison of smooth fillings. The Kodaira dimension of $Q_m$
is $-\infty$: a lifted pluricanonical form extends across the
origin by Hartogs, and its invariance under the contraction
$q^m\operatorname{id}$ forces all its power-series coefficients
to vanish. It cannot correspond birationally to $P_{m'}$.
Therefore any bimeromorphic comparison $U_m\dashrightarrow U_{m'}$
would identify $P_m$ birationally with $P_{m'}$.
Plurigenera are unchanged by a point blowup of a smooth compact
surface: a pluricanonical form descends off the centre and its
coefficients extend across that point by Hartogs. Resolution of
a surface bimeromorphism proves their birational invariance.
The different slopes $(m-1)/m$ in \eqref{cdp:eq:C.ER.2} force $m=m'$.

\item For the persistence assertion, take a finite base change and choose
a local parameter in which it is $t=s^d$; a unit is absorbed by
its local $d$th root. At every point of the nonempty open subset
$P_m\setminus J$, the map $f_m$ is a submersion. Its base change
is the smooth chart $t=s^d$ with the two coordinates along $P_m$
unchanged. Thus a unique central component maps with degree one
onto this open subset. Its closure, followed by normalization and
any resolution, is birational to $P_m$. The preceding
non-uniruled-divisor comparison prevents its contraction under
any further bimeromorphic modification between smooth fillings.
This proves the asserted survival without imposing a description
on the other components of the base change.

\item Finally replace $\Delta$ throughout \eqref{cdp:eq:C.ER.3} by $\PP^1$ and use
the same exceptional divisor, line bundle, and canonical section.
The local proofs of smoothness and properness still apply. Since
the base $V$ is now compact, the quotient is a compact smooth
threefold, with the same central fibre and with every other fibre
the product $E\times E$. This gives the claimed compact realizations.
The exact quotient-chart, canonical, homology, and plurigenus
calculations are retained in \path{code/E_elliptic_roots.py} and
its output.

\end{enumerate}
\end{proof}



\paragraph{Literature comparison and an explicit source audit.}
The occurrence of a Kodaira-dimension-one component in a degeneration
of abelian surfaces already appears in
\cite[Example~3.4(1), p.~99]{Fujimoto1988Log}.
The preceding statement supplies its own quotient construction,
plurigenus distinctions, and finite-base-change comparison.

There is a precise reason not to use the divisibility assertion in
\cite[Theorem~2.8 and Corollary~2.10, pp.~96--97]{Fujimoto1988Log}
as an obstruction to \eqref{cdp:eq:C.ER.4}. For every $m\ge2$ consider the
diagonal Hopf threefold and the displayed map
\[
 \begin{split}
 H_m&=(\C^3\setminus\{0\})/
       \langle(a,b,c)\longmapsto(qa,q^m b,q^m c)\rangle,\\
 \Phi_m:H_m&\longrightarrow\PP^2,
       \qquad[a,b,c]\longmapsto[a^m:b:c].
 \end{split}                                             \tag{C.ER.17}\hypertarget{tag-00IV}{}\label{cdp:eq:C.ER.17}
\]
The action is free and properly discontinuous. A compact annulus
between a small sphere and its image under the contraction meets every
orbit, so $H_m$ is compact and smooth. The three projective coordinates
in \eqref{cdp:eq:C.ER.17} are multiplied by the same factor $q^m$, proving descent.
The map is holomorphic, surjective, and proper. Its fibre over
$[1:Y:Z]$ is parametrized by
$a\mapsto(a,Ya^m,Za^m)$ and is $\C^*/q^{\Z}$; its fibre over
$[0:Y:Z]$ is parametrized by
$s\mapsto(0,sY,sZ)$ and is $\C^*/(q^m)^{\Z}$.
In particular every fibre is connected and elliptic. These are all
the hypotheses in the definition of an elliptic threefold in
\cite[pp.~91--92]{Fujimoto1988Log}.
Near a point of the line $B=(X=0)$, use a chart with $b\ne0$ or
$c\ne0$. In the first chart $X/Y=a^m/b$, where $b$ is a unit.
It follows that
\[
             \Phi_m^*B=m\{a=0\},\qquad B^2=1,
             \qquad\pi_1(\PP^2\setminus B)=0.             \tag{C.ER.18}\hypertarget{tag-00IW}{}\label{cdp:eq:C.ER.18}
\]
Thus \eqref{cdp:eq:C.ER.17} is an explicit counterexample to both the asserted
divisibility and the asserted surjection to $\Z/m$ in that theorem
as printed. This conclusion follows from the verified map, rather
than from an inference that an incomplete proof makes a statement
false.

The problematic covering step can also be located exactly. For an
effective divisor $D$ with canonical section $s_D$, the construction
in that proof takes the equation $\xi^m=s_D^m$ in the total space of
$\OO(D)$. Its normalization is the disjoint union of the $m$ global
graphs $\xi=\zeta s_D$, $\zeta^m=1$, each isomorphic to the original
space. It therefore supplies no connected branched covering of the
base. Finally, \eqref{cdp:eq:C.ER.17} is itself a special case of \eqref{cdp:eq:C.ER.4}: take
$V=\PP^2$, $L=\OO(1)$ and $\sigma=X$. A nonzero functional
$\ell$ on a projective line determines its unique vector $v$ with
$\ell(v)=1$; the equation $\sigma=a^m\ell$ becomes
$v=(a^m,b,c)$. Under this identification \eqref{cdp:eq:C.ER.5} is the inverse
of the contraction in \eqref{cdp:eq:C.ER.17}, giving an actual isomorphism of
the two quotient constructions.




\paragraph{Credit and coordinate comparison.}
The cuspidal Weierstrass self-product and its projective small resolutions
are the construction of Namikawa described in
\cite[Proposition~4, pp.~10--11]{Rossi2017}; the ambient Hodge numbers
$h^{1,1}=21$ and $h^{2,1}=3$ are also recorded in
\cite[Remark~5, p.~20]{Rossi2017}.
The coordinate exchange $(x_{\rm Rossi},y_{\rm Rossi})=(y,x)$
and the coefficient $B(t)=t^6-1$ identify the underlying Weierstrass
surface and its fibre self-product with the equations below.
We retain the explicit two blowup ideals, the full support, and its
integral collapse and normalization maps in the proof.

\begin{proposition}[A simply connected irreducible rational-normalized fibre]\hypertarget{tag-01ER}{}\label{cdp:st:C.2.UN.8}
\StatementRecord{cdp:st:C.2.UN.8}{proved}{01ER}{\SeriesRef{cdp:st:C.2.MH.2}, \SeriesRef{cdp:eq:C.CP.5}, \SeriesRef{cdp:eq:C.CP.7}, \SeriesRef{cdp:eq:C.CP.6}, \SeriesRef{cdp:eq:C.CP.4}, \SeriesRef{cdp:eq:C.CP.10}; \cite[Proposition~4, pp.~10--11, and Remark~5, p.~20]{Rossi2017}}

There is a smooth projective threefold $Y$ with $K_Y\simeq\OO_Y$
and a connected-fibre map $f:Y\to\PP^1$ whose general fibre is a
product of elliptic curves. It has six singular fibres. Each is
reduced and irreducible, has rational normalization resolution, and
satisfies
\[
 \begin{gathered}
 \pi_1(S)=1,\qquad e(S)=6,\\
 (H^i(S,\Z))_{i=0}^4=(\Z,0,\Z^4,0,\Z),\qquad
 (h^i(S,\OO_S))_{i=0}^2=(1,2,1).
 \end{gathered}
 \tag{C.CP.1}\hypertarget{tag-00C9}{}\label{cdp:eq:C.CP.1}
\]
The ambient cohomology is
\[
 \begin{gathered}
 (b_i(Y))_{i=0}^6=(1,0,21,8,21,0,1),\qquad e(Y)=36,\\
 (h^i(Y,\OO_Y))_{i=0}^3=(1,0,0,1),\qquad
 h^{1,1}(Y)=21,\quad h^{2,1}(Y)=3,
 \end{gathered}
 \tag{C.CP.2}\hypertarget{tag-00CB}{}\label{cdp:eq:C.CP.2}
\]
where the Betti numbers use rational coefficients. In particular
the nonzero local character used in \ref{cdp:st:C.2.MH.2} need not exist after
the global CDP hypotheses are removed, and its Euler conclusion
$e(S)=2$ does not extend to every irreducible rational-normalized
torus degeneration.
\end{proposition}


\begin{proof}

\begin{enumerate}[label={$\langle1\rangle$\arabic*.},leftmargin=*]
\item Let $B=\PP^1$, let $L=\OO_B(1)$, and take the Weierstrass
surface $g:R\to B$ with zero section and equation
\[
                  y^2z=x^3+p(t)z^3,
             \qquad p(t)=t^6-1.                 \tag{C.CP.3}\hypertarget{tag-00CC}{}\label{cdp:eq:C.CP.3}
\]
Here $x,y$ have the usual transition weights $L^2,L^3$;
thus $p$ is a section of $L^6$. The projective Weierstrass
bundle supplies the compactification. At infinity the transformed
coefficient is $1-t'^6$, so that fibre is smooth. The six
zeros of $p$ are simple. The only singular fibres are the six
cuspidal cubics, and the surface is smooth: in the affine
Weierstrass chart a possible singular point has $x=y=0$, while
the derivative $p'(t)$ there is nonzero. The zero-section chart
is smooth directly from its equation. Adjunction in the
Weierstrass bundle gives
\[
                         K_R=g^*\OO_B(-1).
 \tag{C.CP.4}\hypertarget{tag-00CD}{}\label{cdp:eq:C.CP.4}
\]
This formula can also be read from the local canonical form
$dt\wedge dx/y$ and the transition weights just specified.

\item Form the projective fibre product $X=R\times_B R$. Its only
singularities are the six points where both factors are cusp
points. At one of them use the base coordinate $\tau=p(t)$,
and write
\[
 y_1^2-x_1^3=\tau=y_2^2-x_2^3.
\]
Fix a primitive cube root $\rho$. With
\[
 \begin{gathered}
 u=y_1-y_2,\quad v=y_1+y_2,\quad
 a=x_1-x_2,\quad b=x_1-\rho x_2,\\
 c=x_1-\rho^2x_2=\alpha a+\beta b,
 \qquad \alpha=-\rho,\quad\beta=-\rho^2,
 \end{gathered}
\]
the exact threefold equation is
\[
                         uv=abc.                \tag{C.CP.5}\hypertarget{tag-00CE}{}\label{cdp:eq:C.CP.5}
\]
Both $\alpha$ and $\beta$ are nonzero. We resolve \eqref{cdp:eq:C.CP.5}
by two specified blowups. First blow up the ideal $(u,a)$.
In its $u$-chart, $a=ur$ and
\[
                  v=rb(\alpha ur+\beta b),
 \tag{C.CP.6}\hypertarget{tag-00CF}{}\label{cdp:eq:C.CP.6}
\]
so $(u,r,b)$ are smooth coordinates. In its $a$-chart put
$u=as$. The equation is $sv=b(\alpha a+\beta b)$; its
only singular point is the origin, an ordinary threefold
node. The strict transform of $(u,b)$ has ideal $(s,b)$
there. Blowing up that ideal gives the two smooth charts
\[
 \begin{array}{c|c}
 \text{coordinates}&\text{blowdown formulas}\\ \hline
 (a,s,h)&b=sh,\quad u=as,\quad
                    v=h(\alpha a+\beta sh),\\
 (b,j,v)&s=bj,\quad a=(jv-\beta b)/\alpha,\quad u=abj.
 \end{array}                                    \tag{C.CP.7}\hypertarget{tag-00CG}{}\label{cdp:eq:C.CP.7}
\]
The inverse image of the original singular point is a chain
of two projective lines. In the first chart of \eqref{cdp:eq:C.CP.7}
the old and new curves have equations $(a,h)=(0,0)$
and $(a,s)=(0,0)$; they meet once transversely. Charts
\eqref{cdp:eq:C.CP.6}--\eqref{cdp:eq:C.CP.7} complete each to a projective line.
No exceptional divisor is introduced. Outside the origin
the respective ideals are Cartier, so these maps are
isomorphisms there.

\item The blowups globalize without an analytic gluing choice.
The ideal $(u,a)$ is the ideal of the Weil diagonal
$R\hookrightarrow X$. The ideal $(u,b)$ is the Weil
graph of the global automorphism
$(x,y,t)\mapsto(\rho x,y,t)$ of $R$; after the first
blowup take its strict transform. Each of these divisors
is Cartier off the six specified singular points.
Their projective blowups therefore give a smooth
projective threefold $Y$ and a proper small map
$\pi:Y\to X$, with precisely the local charts above.
Set $f=g\pi$, using the common projection of the fibre
product. It agrees with the original product of elliptic
curves away from the six values. Its fibres are connected,
because the product fibres and every fibre of the two
blowdowns are connected.

\item The original central support at each of the six values
is the product $S_0=C\times C$ of a cuspidal cubic with
itself. It is reduced and irreducible. Since $\pi$ has
no exceptional divisors, the full fibre upstairs is its
single strict transform, still reduced. Its normalization
resolution is rational: the function field is unchanged
from that of the finite birational parametrization
$\PP^1\times\PP^1\to S_0$. This assertion is birational;
we do not discard the exceptional curves of the actual
normalization resolution.

\item Topologically a cuspidal cubic is $S^2$. Its finite
normalization is bijective, so compactness makes it a
homeomorphism. Hence $S_0$ is homeomorphic to
$S^2\times S^2$. The restriction $S\to S_0$ is an
isomorphism off the cusp-pair point and replaces that
point by the chain $E$ of two projective lines. The
quotient $S/E\to S_0$ is a continuous bijection from a
compact space to a Hausdorff space, hence a homeomorphism.
Choose a finite triangulation of $(S,E)$. Its cofibration
sequence identifies the reduced homology of $S/E$ with
the relative homology of $(S,E)$. Since
$H_2(E,\Z)=\Z^2$ and $E$ is simply connected, the
long exact sequence gives
\[
 \begin{gathered}
 H_4(S,\Z)=\Z,\qquad H_3(S,\Z)=H_1(S,\Z)=0,\\
 0\longrightarrow\Z^2\longrightarrow H_2(S,\Z)
                 \longrightarrow\Z^2\longrightarrow0.
 \end{gathered}                                  \tag{C.CP.8}\hypertarget{tag-00CH}{}\label{cdp:eq:C.CP.8}
\]
The last sequence splits as a sequence of abelian groups.
Van Kampen for the same cofibration shows
$\pi_1(S/E)=\pi_1(S)/\langle\!\langle\pi_1(E)\rangle\!\rangle$.
Thus $\pi_1(S)=1$. Universal coefficients now give the
integral cohomology in \eqref{cdp:eq:C.CP.1}, and its Euler number
is $1+4+1=6$.

\item The fibre product is a local complete intersection in
$R\times R$, cut out by the inverse image of the base
diagonal. Adjunction and \eqref{cdp:eq:C.CP.4} give
\[
 \omega_X=p_1^*K_R\otimes p_2^*K_R\otimes f^*T_B
                         \simeq\OO_X.
 \tag{C.CP.9}\hypertarget{tag-00CI}{}\label{cdp:eq:C.CP.9}
\]
Its generator pulls back to a nowhere-vanishing canonical
form on $Y$ away from the small exceptional curves.
Its coefficients and the coefficients of its inverse
extend across that codimension-two set by Hartogs'
theorem in a local canonical frame. Hence $K_Y=\OO_Y$.
Adjunction makes $\omega_S=\OO_S$. Proper flatness over
the smooth base preserves the holomorphic Euler
characteristic, so $\chi(\OO_S)=\chi(\OO_{E_t^2})=0$.
Connectedness and Serre duality then give
$(h^0,h^1,h^2)(\OO_S)=(1,2,1)$.

\item We include the ambient groups, with their comparison
maps. A type-II local elliptic monodromy has characteristic
polynomial $x^2-x+1$, so $T-I$ is an integral isomorphism.
This also follows directly from the cyclic sixth-root
trivialization of $y^2=x^3+\tau$. The smooth part of
the fibration has fundamental group generated by its
fibre lattice and base meridians. Its inclusion into
$Y$ is surjective on fundamental groups, by perturbing
loops off the central divisor. The pair of zero sections
extends over every special value and kills each base
meridian in $Y$. Conjugation by one such meridian
therefore kills the image of $T-I$ on each of the two
elliptic lattices. Since those two maps are unimodular,
the whole fibre lattice is killed. Thus $\pi_1(Y)=1$.
Projectivity and $K_Y=\OO_Y$ give
$(h^0,h^1,h^2,h^3)(\OO_Y)=(1,0,0,1)$ by Hodge theory
and Serre duality. The identical argument for $R$ gives
$b_1(R)=0$. Its six cuspidal fibres give $e(R)=12$,
hence $b_2(R)=10$. Moreover $K_R=g^*\OO(-1)$ makes
every positive plurigenus zero. The surface classification
therefore makes $R$ rational, so its Picard group has
rank ten.

\item Let $K=\C(t)$ and let $E_K$ be the generic elliptic
curve of $R$. Restriction of divisors gives a surjection
$\Pic(R)\to\Pic(E_K)$: close a generic divisor in the
smooth surface. Its kernel is generated, up to linear
equivalence, by the full fibre, because every fibre is
irreducible. The zero section splits the degree map on
$\Pic(E_K)$, whose degree-zero part is $E_K(K)$.
Consequently $E_K(K)$ has rank eight.

\item Apply the same divisor restriction to $Y$ and its
generic fibre $A_K=E_K\times E_K$. It is surjective
and its vertical kernel has rank one, since all six
singular fibres of $Y$ are irreducible. The degree-zero
Picard group of $A_K$ is its dual abelian variety's
$K$-points, which have rank sixteen under the product
principal polarization. Its N\'eron--Severi group has
rank four. To verify this last number explicitly, use
the two factors, the diagonal, and the graph of the
automorphism $\rho$. Their intersection matrix is
\[
 \begin{pmatrix}
 0&1&1&1\\1&0&1&1\\1&1&0&3\\1&1&3&0
 \end{pmatrix}.                                \tag{C.CP.10}\hypertarget{tag-00CA}{}\label{cdp:eq:C.CP.10}
\]
Its determinant is $-3$. The entry three is the degree of
$1-\rho$, whose kernel on the triangular elliptic
lattice has order three. Here is the exact comparison
for the upper bound on the generic N\'eron--Severi rank.
Over $K'=K(u)$ with $u^6=p(t)$, the algebraic
coordinate substitution $x=u^2X$, $y=u^3Y$ identifies
$E_{K'}$ with the constant elliptic curve
$E_0:\ Y^2=X^3+1$. Base change to an algebraic closure
identifies the geometric N\'eron--Severi group with that
of $(E_0\times E_0)_{\overline K}$. For the constant
complex abelian surface, the first Chern class embeds
this group into the integral $(1,1)$ classes and
$h^{1,1}(E_0^2)=4$. Algebraically closed extension
preserves the N\'eron--Severi group: the Picard scheme
and its identity component commute with field extension,
and each of its connected components over $\C$ remains
geometrically connected. The four classes
in \eqref{cdp:eq:C.CP.10} are already defined over $K$, so its rank
over $K$ is exactly four. Thus
$\operatorname{rank}\Pic(Y)=1+16+4=21$.
Since $h^1(\OO_Y)=h^2(\OO_Y)=0$, the exponential
sequence identifies this rank with $b_2(Y)$.

\item Euler localization now gives $e(Y)=6\cdot6=36$.
Poincar\'e duality and $b_1=0$ yield $b_3=8$;
the remaining Betti numbers and Hodge numbers follow
as displayed in \eqref{cdp:eq:C.CP.2}. These rational Betti assertions
make no assertion about integral torsion in degrees
not already computed.

\item Finally a sufficiently small compact fibre tube retracts
onto $S$. Its fundamental group is therefore trivial.
There is no nontrivial finite local system there, and
no nonzero integral degree-one class to restrict to
the nearby torus. The example consequently disproves
the proposed unrestricted replacement for the character
construction itself, with its actual comparison map.
The exact coordinate identities, intersection matrix,
and all displayed integer ranks are checked in
\path{code/root_cusp_product.py} and its saved output.

\end{enumerate}
\end{proof}



The compact construction in this proof is the Namikawa
cuspidal self-product. Rossi~\cite[\S3, Proposition~4]{Rossi2017}
gives the same two graph blowups and exceptional chains;
his Remark~5 gives $h^{1,1}=21$ and $h^{2,1}=3$.
The coordinate,
integral fibre-cohomology and Picard calculations above
derive the precise local-character obstruction used here.





\begin{proposition}[A projective zero-Euler filling with an essential Kummer component]\hypertarget{tag-01CQ}{}\label{cdp:st:C.2.UN.13}
\StatementRecord{cdp:st:C.2.UN.13}{proved}{01CQ}{\SeriesRef{cdp:eq:C.KM.4}, \SeriesRef{cdp:eq:C.KM.5}, \SeriesRef{cdp:st:C.2.UN.4}, \SeriesRef{cdp:eq:C.KM.S2}, \SeriesRef{cdp:st:C.2.UN.3}; \cite[Chapter~1, Example~1.3(iii)]{HuybrechtsK3Draft}}

Fix an abelian surface $A$. There is a smooth projective threefold
$W$ with a proper map $f_W:W\to\PP^1$ whose fibres over
$\C^*$ are isomorphic to $A$. It has precisely two special fibres,
over $0$ and $\infty$. Each has full divisor
\[
               2P+\sum_{i=1}^{16}E_i,\qquad
 P=\operatorname{Km}(A),\quad E_i\simeq\PP^2.       \tag{C.KM.1}\hypertarget{tag-00PI}{}\label{cdp:eq:C.KM.1}
\]
Here $P$ is the smooth Kummer K3 surface, the $E_i$ are disjoint,
and $P\cap E_i$ is a projective line on $E_i$ and a curve of
self-intersection $-2$ on $P$. The reduced divisor is simple
normal crossing and has Euler number $40$.

There is a composition of smooth vertical blowups $Y\to W$,
disjoint from both K3 components, for which both special reduced
fibres have Euler number zero. Each has $57$ irreducible
components, including its original unchanged Kummer surface.
The maps to the base and the nearby tori are retained. Their
rational Betti numbers and coherent dimensions are
\[
 \begin{array}{c|c|c|c}
 & (b_i)_{i=0}^6 & e & (h^i(\OO))_{i=0}^3\\ \hline
 W&(1,0,39,0,39,0,1)&80&(1,0,1,0)\\
 Y&(1,0,119,240,119,0,1)&0&(1,0,1,0).
 \end{array}                                             \tag{C.KM.2}\hypertarget{tag-00PL}{}\label{cdp:eq:C.KM.2}
\]
For either special value, let $S_W$ and $S_Y$ denote its
reduced supports in these two models. Their cohomology is
\[
 \begin{gathered}
 (H^i(S_W,\Q))_{i=0}^4=(\Q,0,\Q^{22},0,\Q^{17}),\\
 (H^i(S_Y,\Q))_{i=0}^4=(\Q,0,\Q^{62},\Q^{120},\Q^{57}),\\
 (h^i(S_W,\OO_{S_W}))_{i=0}^2
   =(h^i(S_Y,\OO_{S_Y}))_{i=0}^2=(1,0,1).
 \end{gathered}                                          \tag{C.KM.S1}\hypertarget{tag-00PT}{}\label{cdp:eq:C.KM.S1}
\]
Neither K3 component can be contracted by a bimeromorphic map
between smooth fillings. For comparison, after the quadratic
base change the normalization of the original quotient is
exactly $A\times\PP^1$, and the corresponding base changes
of $W$ and $Y$ are birational modifications of that product.
Thus the asserted persistence concerns bimeromorphic comparison
over the original base; its relation to base change is given
by this specified map.
\end{proposition}


\begin{proof}

\begin{enumerate}[label={$\langle1\rangle$\arabic*.},leftmargin=*]
\item On $A\times\PP^1_s$ take the involution
\[
       \iota(a,s)=(-a,-s),\qquad
 X=(A\times\PP^1_s)/\langle\iota\rangle,
       \qquad t=s^2.                                  \tag{C.KM.3}\hypertarget{tag-00PM}{}\label{cdp:eq:C.KM.3}
\]
The quotient is projective and normal. Its fixed points are
exactly $A[2]\times\{0,\infty\}$, so there are $32$ of them.
Away from them the action is free. In local coordinates centred
at any fixed point, including the coordinate $s^{-1}$ at
infinity, the action is scalar multiplication by $-1$ on
$\C^3$. Thus the only singularities of $X$ are those $32$
quotient points.

\item Blow up these points in $X$, and call the resulting map
$\pi:W\to X$. We verify its smoothness and its exact relation
to the quotient before using it. Write the coordinates of a
covering chart as $(a,b,s)$. The invariant local ring is
generated by their six quadratic monomials. In the blowup
chart corresponding to $a^2$, put
\[
             a=r,\qquad b=ru,\qquad s=rv,
             \qquad R=r^2.                            \tag{C.KM.4}\hypertarget{tag-00PN}{}\label{cdp:eq:C.KM.4}
\]
All invariant generators pull back as
\[
 (a^2,ab,as,b^2,bs,s^2)
           =(R,Ru,Rv,Ru^2,Ruv,Rv^2).                  \tag{C.KM.5}\hypertarget{tag-00PO}{}\label{cdp:eq:C.KM.5}
\]
Consequently $(R,u,v)$ are smooth coordinates on the blowup.
The analogous charts for $b^2$ and $s^2$ cover it: the
exceptional projectivized invariant cone is the quadratic
Veronese image of $\PP^2$, on which at least one of the
three squared coordinates is nonzero. Equivalently, blow up
the origin of $\C^3$ first. The lifted involution fixes its
exceptional $\PP^2$ and sends its normal coordinate $r$ to
$-r$; its quotient has precisely \eqref{cdp:eq:C.KM.4}--\eqref{cdp:eq:C.KM.5}.
The transitions for $R=r^2$ identify its normal line with
$\OO_{\PP^2}(-2)$. This proves that $W$ is smooth and
projective, with exceptional divisors $E_i\simeq\PP^2$.

\item The base map in \eqref{cdp:eq:C.KM.4} is exactly
\[
                             t=R v^2.                 \tag{C.KM.6}\hypertarget{tag-00PP}{}\label{cdp:eq:C.KM.6}
\]
Here $E_i=(R=0)$ and the strict central component is
$P=(v=0)$. In the chart where $s$ is the exceptional
coordinate the map is $t=R$ and only $E_i$ occurs.
These charts prove the multiplicities and the transverse
intersection in \eqref{cdp:eq:C.KM.1}. They also show that
$N_{E_i/W}|_{P\cap E_i}=\OO_{\PP^1}(-2)$, which is
the normal line of that intersection in $P$.
The same calculation applies over infinity. Off those two
values, choose either square root $s$ of $t$; the map
$a\mapsto[a,s]$ identifies the fibre with $A$, and switching
the square root composes this identification with $a\mapsto-a$.
In particular the map has connected fibres and its monodromy
on $H^1(A,\Z)$ is $-I_4$.

\item The restriction $P\to A/\{\pm1\}$ resolves exactly the
sixteen quotient surface nodes. It is the classical Kummer
construction; see \cite[Chapter~1, Example~1.3(iii)]{HuybrechtsK3Draft}.
Here the two required properties can also be checked in the
same charts. A translation-invariant two-form on $A$ is
invariant under $-1$. In a surface chart $a=r$, $b=ru$,
its pullback is
\[
                da\wedge db=r\,dr\wedge du
                            =\tfrac12\,dR\wedge du.   \tag{C.KM.7}\hypertarget{tag-00PQ}{}\label{cdp:eq:C.KM.7}
\]
It therefore descends to a nowhere-zero form on $P$.
The finite covering $\operatorname{Bl}_{A[2]}A\to P$
injects $H^1(P,\OO_P)$ into the invariant part of
$H^1(A,\OO_A)$, by the averaging trace. This invariant
part is zero because $-1$ acts as $-1$ on $H^1(A,\OO_A)$.
Thus $K_P=\OO_P$ and $h^1(P,\OO_P)=0$, verifying that
$P$ is a K3 surface.

\item The finite quotient map identifies rational cohomology with
invariants. Explicitly, its higher direct images of the
constant rational sheaf vanish by proper base change and
the finiteness of its fibres; averaging identifies the
constant sheaf downstairs with the invariant direct image.
Applied to $A/\{\pm1\}$, this gives Betti numbers
$(1,0,6,0,1)$ and Euler number eight. Resolving each of its
sixteen points replaces that point by $\PP^1$, so
$e(P)=8+16=24$. Hence the actual inclusion--exclusion
comparison in \eqref{cdp:eq:C.KM.1} gives
\[
          e\left(P\cup\bigcup_{i=1}^{16}E_i\right)
                    =24+16\cdot3-16\cdot2=40.          \tag{C.KM.8}\hypertarget{tag-00PR}{}\label{cdp:eq:C.KM.8}
\]
Mayer--Vietoris gives the support groups as well. The restriction
map from the direct sum of $H^2(P,\Q)$ and the sixteen
$H^2(E_i,\Q)$ onto the sixteen $H^2(P\cap E_i,\Q)$ is
surjective: the hyperplane class on each $E_i$ restricts to
the degree-one class of its intersection line. Its kernel
has dimension $22+16-16=22$, and the same exact sequence
gives $H^1(S_W,\Q)=H^3(S_W,\Q)=0$ and
$H^4(S_W,\Q)=\Q^{17}$. The coherent sequence
\[
 0\longrightarrow\OO_{S_W}\longrightarrow
 \OO_P\oplus\bigoplus_i\OO_{E_i}
 \longrightarrow\bigoplus_i\OO_{P\cap E_i}
 \longrightarrow0
\]
has a surjective difference map on $H^0$. The lines and planes
have no positive coherent cohomology, so it gives
$(h^0,h^1,h^2)(\OO_{S_W})=(1,0,1)$.

\item We give the ambient cohomology maps as well. For the finite
quotient $A\times\PP^1\to X$, the same invariant-sheaf
argument gives
\[
 H^k(X,\Q)\simeq H^k(A\times\PP^1,\Q)^{\iota},
 \qquad (b_k(X))_{k=0}^6=(1,0,7,0,7,0,1).              \tag{C.KM.9}\hypertarget{tag-00PS}{}\label{cdp:eq:C.KM.9}
\]
Indeed $-1$ acts by $(-1)^j$ on $H^j(A)$, while the
base involution acts trivially on both nonzero cohomology
groups of $\PP^1$.

\item Let $j_i:E_i\hookrightarrow W$ and let
$h_i=c_1(\OO_{E_i}(1))$. Collapsing the $32$ exceptional
divisors separately gives exactly $X$. The relative
cohomology sequence of this quotient identifies the
kernel of restriction to the $E_i$ with the pullback
of $H^k(X,\Q)$. The self-intersection formula gives
\[
 j_i^*j_{i*}(1)=-2h_i,
       \qquad j_i^*j_{i*}(h_i)=-2h_i^2.                 \tag{C.KM.10}\hypertarget{tag-00PJ}{}\label{cdp:eq:C.KM.10}
\]
Thus restriction is surjective in degrees two and four,
with explicit rational sections $-j_{i*}(1)/2$ and
$-j_{i*}(h_i)/2$. In these two degrees pullback is
injective because the preceding odd cohomology of the
exceptional divisors vanishes. The same relative
sequence then leaves the other positive-degree groups
unchanged. Each exceptional point consequently adds one
class in degrees two and four and none in any other
degree. This proves the row for $W$ in \eqref{cdp:eq:C.KM.2}.

\item For coherent cohomology the finite quotient and averaging
give
$H^k(X,\OO_X)=H^k(A\times\PP^1,\OO)^{\iota}$,
of dimensions $(1,0,1,0)$. The local resolution just
used has $R^q\pi_*\OO_W=0$ for $q>0$. To verify this
without suppressing the quotient singularity, filter the
infinitesimal neighbourhoods of $E_i$ by its ideal.
Their successive quotients are $\OO_{\PP^2}(2r)$,
$r\ge0$, whose positive-degree coherent cohomology
vanishes. Formal functions therefore makes the completed
stalks of those higher direct images zero, and coherence
makes the stalks themselves zero. Normality gives
$\pi_*\OO_W=\OO_X$. Leray proves the coherent row for
$W$.

\item We now perform exactly twenty of the negative-two
operations of \ref{cdp:st:C.2.UN.4} over each of $0$ and $\infty$.
Every operation starts at a fresh smooth point on the
strict transform of one fixed $E_i$, away from its
intersection with $P$ and from all other components.
The first blowup creates a new $\PP^2$, and the second
blows up a smooth plane quartic of genus three inside
it, transverse to its old intersection line. These
centres and their entire modifications are disjoint
from $P$. The exact charts of \ref{cdp:st:C.2.UN.4} retain simple
normal crossings and show that both new components
have multiplicity one. Each operation changes the
reduced-fibre Euler number by $-2$ and adds two
components. Thus both final supports have Euler number
$40-20\cdot2=0$ and $17+20\cdot2=57$ components,
with their unchanged multiplicity-two K3 surfaces.

\item We retain the support comparison under each of these blowups,
including the four boundary intersections of the quartic.
Let $\sigma:\widehat W\to W'$ be one blowup, with smooth
centre $C$ of ambient codimension $c$, let $S$ be the current
reduced support, and put $S'=\sigma^{-1}(S)_{\rm red}$ and
$E=\PP(N_{C/W'})$. Collapsing $E$ in $S'$ and $C$ in $S$
gives the actual homeomorphism $S'/E\simeq S/C$.
Hence pullback identifies the relative cohomology of the
pairs $(S,C)$ and $(S',E)$. With $p_E:E\to C$ and
$\xi=c_1(\OO_E(1))$, the projective-bundle decomposition is
\[
 H^k(E,\Q)=\bigoplus_{i=0}^{c-1}
                \xi^i p_E^*H^{k-2i}(C,\Q).
\]
For $i\ge1$, the restriction to $S'$ of the ambient Gysin
class $j_*(\xi^{i-1}p_E^*\alpha)$ restricts further to $E$
as $-\xi^i p_E^*\alpha$, because $N_{E/\widehat W}=\OO_E(-1)$.
These classes split every positive projective-bundle summand.
The two relative exact sequences now identify the remaining
summand with the pullback of $H^k(S,\Q)$ and give
\[
 H^k(S',\Q)\simeq H^k(S,\Q)\oplus
                \bigoplus_{i=1}^{c-1}H^{k-2i}(C,\Q).
                                                        \tag{C.KM.S2}\hypertarget{tag-00PU}{}\label{cdp:eq:C.KM.S2}
\]
This argument uses the full inverse image and does not require
the centre to avoid the other support components at every point.
A point step therefore adds one class in degrees two and four.
The quartic step adds ranks one, six and one in degrees two,
three and four. After twenty blocks, \eqref{cdp:eq:C.KM.S2} gives exactly
the second support row of \eqref{cdp:eq:C.KM.S1}.

\item Each chosen centre is generically contained in just one
reduced component, of multiplicity one in the reduced support.
Consequently the equality of Cartier divisors
$\sigma^*S=S'$ holds, also at the four crossing points as
the explicit equations show. Apply derived pushforward to
$0\to\OO_{\widehat W}(-S')\to\OO_{\widehat W}\to
\OO_{S'}\to0$. The projection formula and
$R\sigma_*\OO_{\widehat W}=\OO_{W'}$ give
\[
                         R\sigma_*\OO_{S'}=\OO_S.       \tag{C.KM.S3}\hypertarget{tag-00PV}{}\label{cdp:eq:C.KM.S3}
\]
Thus all support coherent cohomology groups are unchanged,
proving the remaining assertion of \eqref{cdp:eq:C.KM.S1}.

\item The pullback and exceptional Gysin maps in \ref{cdp:st:C.2.UN.4}
give cohomology increments $(\Delta b_2,\Delta b_3,
\Delta b_4)=(2,6,2)$ per operation. Forty operations
give the second Betti row in \eqref{cdp:eq:C.KM.2}. Smooth-centre
blowups preserve coherent cohomology, by their direct
image $\OO$ and vanishing higher direct images, so
the second coherent row is unchanged. This also checks
$e(Y)=0$ directly from its displayed Betti numbers.

\item Finally, $P$ is not uniruled. A dominating ruled surface
would pull its nonzero holomorphic two-form back to a
nonzero two-form, whereas a ruled surface has no such
form. The inverse-map resolution in \ref{cdp:st:C.2.UN.3} therefore
prevents its contraction by any bimeromorphic comparison
between smooth threefold fillings. The finite-base-change
comparison is different and explicit: the map
\[
 A\times\PP^1_s\longrightarrow
 X\times_{\PP^1_t}\PP^1_s,
         \qquad(a,s)\longmapsto([a,s],s)               \tag{C.KM.11}\hypertarget{tag-00PK}{}\label{cdp:eq:C.KM.11}
\]
is finite and generically one-to-one. Its source is
normal, so it is the normalization of the reduced
base change. The normalized base changes of $W$ and
$Y$ factor properly and birationally through this
normalization. This specifies exactly how the Kummer
component is related to a torus after that base change.
The construction is based on the classical Kummer quotient.
The explicit invariant-character charts and transitions, descended
two-form coefficients, invariant and Gysin matrices, ambient and support
Betti vectors, quartic-block Euler changes and quadratic-normalization
inverse are computed in
\path{code/current_audits/audit_sign/code/UN13_current_exact_certificate_08.py}.
The complete output is
\path{code/current_audits/audit_sign/UN13_current_exact_certificate_08.json}.
These calculations use the finite-sheaf, relative-pair and
derived-pushforward maps proved above.

\end{enumerate}
\end{proof}






\begin{proposition}[Opposite logarithmic data produce a projective essential component]\hypertarget{tag-01CV}{}\label{cdp:st:C.2.UN.18}
\StatementRecord{cdp:st:C.2.UN.18}{proved}{01CV}{\SeriesRef{cdp:eq:C.PE.5}, \SeriesRef{cdp:eq:C.PE.6}, \SeriesRef{cdp:eq:C.PE.7}, \SeriesRef{cdp:eq:C.PE.12}, \SeriesRef{cdp:eq:C.PE.8}, \SeriesRef{cdp:eq:C.PE.9}, \SeriesRef{cdp:eq:C.PE.16}, \SeriesRef{cdp:eq:C.PE.18}, \SeriesRef{cdp:st:C.2.UN.3}, \SeriesRef{cdp:st:C.2.UN.6}; \cite[Example~3.4(1)]{Fujimoto1988Log}}

Put $q=e^{-2\pi}$ and $E=\C^*/q^{\Z}$. Choose distinct
points $p_+,p_-\in E$ whose difference is torsion.
For every $m\ge2$ the fixed punctured family
$(E\times E)\times\Delta^*$ has a smooth proper filling
with full central fibre
\[
                         P_m+mQ_++mQ_-.                 \tag{C.PE.1}\hypertarget{tag-00U2}{}\label{cdp:eq:C.PE.1}
\]
Here $Q_+$ and $Q_-$ are disjoint primary Hopf surfaces,
$P_m$ is a smooth projective surface of Kodaira dimension
one, and $J_\pm=P_m\cap Q_\pm$ are disjoint smooth
elliptic curves. The support is a simple normal crossing
surface of Euler number zero. The map $P_m\to E$ is an
elliptic fibration with exactly the two multiple fibres
$mJ_+$ and $mJ_-$, and
\[
 \begin{gathered}
 K_{P_m}\simeq\OO_{P_m}((m-1)(J_++J_-)),\qquad
 H_1(P_m,\Z)\simeq\Z^4,\\
 h^0(P_m,K_{P_m}^{\otimes r})=
 \begin{cases}1,&d_r=0,\\d_r,&d_r>0,\end{cases}
 \qquad d_r=2\left\lfloor\frac{r(m-1)}m\right\rfloor.
 \end{gathered}                                         \tag{C.PE.2}\hypertarget{tag-00UD}{}\label{cdp:eq:C.PE.2}
\]
The surface $P_m$ admits a finite \'etale cover by a
product of two projective curves, one elliptic and the
other of genus at least two. All three components contain
no rational curve and survive birationally in every
bimeromorphic comparison between smooth fillings. The
multiplicity-one component $P_m$ also survives birationally
after every finite base change and resolution. The
fillings for different $m$ are not bimeromorphic over the
disc.

The construction extends to a smooth compact threefold
$Y_m\to\PP^1$ with just this special fibre. All integral
cohomology groups below are free. Writing $S_m$ for its
reduced special support, their ranks and coherent
cohomology dimensions are
\[
 \begin{array}{c|c|c}
 &\operatorname{rank}H^k(-,\Z)&h^k(\OO)\\ \hline
 P_m&(1,4,6,4,1)&(1,2,1)\\
 Q_\pm&(1,1,0,1,1)&(1,1,0)\\
 S_m&(1,3,6,7,3)&(1,2,1)\\
 Y_m&(1,3,7,10,7,3,1)&(1,2,1,0).
 \end{array}                                           \tag{C.PE.3}\hypertarget{tag-00UG}{}\label{cdp:eq:C.PE.3}
\]
\end{proposition}


\begin{proof}

\begin{enumerate}[label={$\langle1\rangle$\arabic*.},leftmargin=*]
\item Let $V$ be the blowup of $E\times\Delta$ at
$(p_+,0)$ and $(p_-,0)$, and let $R_+,R_-$ be its
disjoint exceptional curves. Put
\[
 D=R_+-R_-,\qquad L=\OO_V(D),\qquad
 \operatorname{div}(\sigma)=D,                         \tag{C.PE.4}\hypertarget{tag-00UH}{}\label{cdp:eq:C.PE.4}
\]
where $\sigma$ is the canonical meromorphic section of
$L$. We define a smooth covering space $W_m$ using two
actual charts. On $V\setminus R_-$ use
\[
 \sigma=a^m\ell,\qquad \ell\in L^\times,\quad a\in\C;
                                                               \tag{C.PE.5}\hypertarget{tag-00UI}{}\label{cdp:eq:C.PE.5}
\]
on $V\setminus R_+$ use
\[
 \sigma^{-1}=(a')^m\ell',\qquad
 \ell'\in(L^{-1})^\times,\quad a'\in\C.                 \tag{C.PE.6}\hypertarget{tag-00UJ}{}\label{cdp:eq:C.PE.6}
\]
On the common open set both $a,a'$ are nonzero, and the
gluing is $a'=a^{-1}$, $\ell'=\ell^{-1}$. Thus $a$
is globally a $\PP^1$-valued coordinate. The generator
of the free $\Z$-action is
\[
 (\ell,a)\longmapsto(q^m\ell,q^{-1}a),\qquad
 (\ell',a')\longmapsto(q^{-m}\ell',qa').                 \tag{C.PE.7}\hypertarget{tag-00UK}{}\label{cdp:eq:C.PE.7}
\]
Near $R_+$, a frame of $L$ makes \eqref{cdp:eq:C.PE.5} the equation
$x=a^m b$, $b\ne0$. Near $R_-$, a frame of $L^{-1}$
makes \eqref{cdp:eq:C.PE.6} the equation $x=(a')^m b'$, $b'\ne0$.
Each eliminates $x$ and proves smoothness. A Hermitian
norm of the globally defined nonzero vector $\ell$ is
multiplied by $q^m$ in \eqref{cdp:eq:C.PE.7}; it proves proper
discontinuity and the absence of stabilizers. The
quotient $U_m=W_m/\Z$ is therefore a smooth threefold.

\item The local proper maps are explicit. In the positive
chart write $x=z^m$, $b=w^m$, $a=z/w$ and quotient
by $(z,w)\mapsto(\zeta z,\zeta w)$. The generator
\eqref{cdp:eq:C.PE.7} lifts to $w\mapsto qw$. In the negative
chart use $x=z^m$, $b'=(w')^m$, $a'=z/w'$ and the
same diagonal finite action; the generator lifts to
$w'\mapsto q^{-1}w'$. In each case the local map to
$V$ is a free finite quotient of
\[
 \Delta_z\times E\times\Delta_y
   \longrightarrow\Delta_x\times\Delta_y,\qquad
 (z,[w],y)\longmapsto(z^m,y).                           \tag{C.PE.8}\hypertarget{tag-00UL}{}\label{cdp:eq:C.PE.8}
\]
Multiplication by a nontrivial $m$th root of unity is a
nontrivial translation of $E$, so the finite action is
free. The map before this quotient is proper; hence
the resulting $\rho:U_m\to V$ is proper. Away from
$R_+\cup R_-$ the actual product identification is
\[
 (v,[c])\longmapsto[(v,\ell=\sigma(v)c^m,a=c^{-1})].     \tag{C.PE.9}\hypertarget{tag-00UM}{}\label{cdp:eq:C.PE.9}
\]
It gives exactly the stated fixed punctured product.

\item Write $B\simeq E$ for the strict central curve of $V$.
The base function has divisor $B+R_++R_-$. At the two
intersections it has equation $t=xy$, so its pullbacks
are respectively
\[
                   t=a^mby,\qquad t=(a')^m b'y.         \tag{C.PE.10}\hypertarget{tag-00U3}{}\label{cdp:eq:C.PE.10}
\]
These are transverse products with multiplicities
$m,1$. They prove \eqref{cdp:eq:C.PE.1}, smoothness of $P_m$ and
the stated intersections. Over $R_+$ the component is
$L|_{R_+}^{\times}/q^{m\Z}$; over $R_-$ it is
$(L^{-1}|_{R_-})^{\times}/q^{-m\Z}$. Both line bundles
are $\OO_{\PP^1}(-1)$. Sending a nonzero tautological
vector to the same vector in $\C^2\setminus0$ identifies
both components with
$(\C^2\setminus0)/q^{m\Z}$, a primary Hopf surface.
The two reduced multiple fibres $J_\pm$ are each
$\C^*/q^{m\Z}$.

\item We construct the finite projective cover of $P_m$.
The restriction of $L$ to $B$ is
$L_B=\OO_E(p_+-p_-)$, a torsion line. Choose a torsion
line $N$ with $N^{\otimes m}\simeq L_B$; multiplication
by $m$ is onto on the torsion subgroup of $\Pic^0(E)$.
Normalize the cyclic cover described in the projective
bundle $\PP(\OO_E\oplus N)$ by
\[
                             y^m=\sigma|_B.              \tag{C.PE.11}\hypertarget{tag-00U4}{}\label{cdp:eq:C.PE.11}
\]
Here $y$ is the tautological meromorphic $N$-coordinate.
At $p_+$ this is $y^m=z$, and at $p_-$ its reciprocal
is $(y^{-1})^m=z$, after local frame choices. The
normalization $C$ is a connected smooth projective
curve: the order-one zero of $\sigma$ prevents the
cyclic equation from splitting, and the local maps
are the stated fully ramified ones. Riemann--Hurwitz
gives $2g(C)-2=2(m-1)$, so $g(C)=m$.

\item Pull back a finite \'etale cover of $E$ trivializing
$N$, and choose a connected component of its fibre
product with $C$. This gives a finite \'etale cover
$C'\to C$ and a map $\tau:C'\to E$ such that
$\tau^*N$ has a nowhere-zero frame $n$. Write
$y=y_0n$. There is a holomorphic map
\[
 \begin{split}
 \Phi:C'\times E&\longrightarrow P_m,\\
 (z,[u])&\longmapsto
       [\ell=u^m n(z)^m,\ a=y_0(z)/u].
 \end{split}                                             \tag{C.PE.12}\hypertarget{tag-00U5}{}\label{cdp:eq:C.PE.12}
\]
At a pole of $y_0$ this formula is interpreted in
\eqref{cdp:eq:C.PE.6}, so $a'=u/y_0$ and $\ell'=u^{-m}n^{-m}$.
It is well defined under $u\mapsto qu$ by \eqref{cdp:eq:C.PE.7}.
It is onto and has finite fibres: away from the two
branch values its restriction to each elliptic fibre
is an isomorphism; at a branch value it is the finite
map $[u]\mapsto[u^m]$ from $E$ to
$\C^*/q^{m\Z}$. Thus it is finite and proper.

\item In fact $\Phi$ is \'etale everywhere. Local nonvanishing
frame factors are absorbed into the coordinate $u$.
At a zero branch choose the root parameter $\zeta$;
in the actual surface coordinates
$(a,b)$ its map is
\[
             (\zeta,u)\longmapsto(\zeta/u,u^m),
             \qquad \det d\Phi=m u^{m-2}\ne0.           \tag{C.PE.13}\hypertarget{tag-00U6}{}\label{cdp:eq:C.PE.13}
\]
At a pole use $y_0=\zeta^{-1}$ and the map
$(\zeta,u)\mapsto(a'=\zeta u,b'=u^{-m})$, whose
Jacobian is $-m u^{-m}\ne0$. Away from those branches
the assertion follows from the \'etale base map and
the fibre isomorphism. This proves the claimed finite
\'etale product cover. Its source is projective, and
a finite holomorphic image of a projective surface is
Moishezon; a smooth compact Moishezon surface is
projective. The comparison can also be made by taking
the finite \'etale Galois closure of \eqref{cdp:eq:C.PE.12}, a
projective surface, and its finite algebraic quotient.

\item Let $\omega$ be a nowhere-zero one-form on $E$.
Off the multiple fibres, the form
$p^*\omega\wedge d\log\ell$ is well defined: a change
of frame adds a base one-form to $d\log\ell$, whose
wedge with $p^*\omega$ is zero. Near $J_+$, write
$\omega=h(x)dx$ and $x=a^m b$; then this two-form is
\[
                m h(a^m b)a^{m-1}\,da\wedge db.         \tag{C.PE.14}\hypertarget{tag-00U7}{}\label{cdp:eq:C.PE.14}
\]
Near $J_-$ the dual formula has a minus sign and the
same vanishing order $m-1$. It has no further zero.
This proves the canonical line in \eqref{cdp:eq:C.PE.2}. Constancy
on compact elliptic fibres and the invariant local
root charts give the exact direct image
\[
 p_*\OO_{P_m}(k_+J_++k_-J_-)
 =\OO_E\!\left(
  \left\lfloor\frac{k_+}m\right\rfloor p_+
 +\left\lfloor\frac{k_-}m\right\rfloor p_-\right).         \tag{C.PE.15}\hypertarget{tag-00U8}{}\label{cdp:eq:C.PE.15}
\]
A permitted Laurent term in a root parameter descends
precisely when its order is divisible by $m$.
Substitute $k_+=k_-=r(m-1)$ and use Riemann--Roch
on the elliptic curve. This proves the plurigenera
and $\kappa(P_m)=1$.

\item We give the integral comparisons in \eqref{cdp:eq:C.PE.3}. The
real action
$(\ell,a)\mapsto(\lambda^m\ell,\lambda^{-1}a)$,
$\lambda>0$, preserves the equations and translates
$\log\|\ell\|$. Since $q$ is positive real, its
quotient gives actual differentiable products
\[
                   P_m\simeq M_3\times S^1,\qquad
                   Y_m\simeq M_5\times S^1,             \tag{C.PE.16}\hypertarget{tag-00U9}{}\label{cdp:eq:C.PE.16}
\]
where the $M$'s are the level sets $\|\ell\|=1$ in
the corresponding covering spaces. For $Y_m$ use
the construction over
$V=\operatorname{Bl}_{(p_+,0),(p_-,0)}(E\times\PP^1)$.
The local proper maps \eqref{cdp:eq:C.PE.8} make this quotient
compact and smooth. Off $0$ the product comparison
\eqref{cdp:eq:C.PE.9} applies, including at infinity.

\item For $M_3$, remove small discs around $p_+,p_-$.
Over their complement the angular coordinate of
$a$ trivializes its circle fibration. Let $h$ be
the positive angular $a$-loop and let $\delta_\pm$
be positive base boundary loops with $a$ fixed.
In the positive solid torus the phases satisfy
$\delta_+=m\arg(a)+\arg(b)$; in the negative one
they satisfy $\delta_-=m\arg(a')+\arg(b')$ and
$\arg(a)=-\arg(a')$. The two meridians are
\[
                         m\delta_++h,\qquad
                         m\delta_- -h.                 \tag{C.PE.17}\hypertarget{tag-00UA}{}\label{cdp:eq:C.PE.17}
\]
In the abelianization of the punctured elliptic base,
$\delta_++\delta_-=0$. These relations give
$h=-m\delta_+$ and $\delta_-=-\delta_+$, leaving
$H_1(M_3,\Z)=\Z^3$, generated by the two
elliptic-base loops and $\delta_+$. There is no
torsion. Poincar\'e duality on this closed oriented
three-manifold and \eqref{cdp:eq:C.PE.16} give all the asserted
integral groups of $P_m$.

\item For the compact five-manifold put
$Q=(E\times\PP^1)\setminus(B^4_+\sqcup B^4_-)$,
with two disjoint open balls. The exterior is
$Q\times S^1$, while each filling is
$S^3\times D^2$. Near an exceptional curve regard
a nonzero tautological vector as
$w\in\C^2\setminus0$. On the unit norm slice
$w\in S^3$, and the blowdown is
$u=a^m w$ or $u=(a')^m w$. The boundary maps are
\[
 \begin{split}
 (w,e^{\mathrm i\theta})&\longmapsto
       (e^{\mathrm i m\theta}w,e^{\mathrm i\theta}),\\
 (w,e^{\mathrm i\theta})&\longmapsto
       (e^{\mathrm i m\theta}w,e^{-\mathrm i\theta}).
 \end{split}                                             \tag{C.PE.18}\hypertarget{tag-00UB}{}\label{cdp:eq:C.PE.18}
\]
The pair sequence for the two removed balls gives
free homology of $Q$ of ranks $(1,2,2,3,0)$; its
two boundary sphere classes are opposite
primitive classes $\delta,-\delta$ in $H_3(Q,\Z)$.
In Mayer--Vietoris for \eqref{cdp:eq:C.PE.18}, the map on
$H_1(S^3\times S^1)^{\oplus2}$ has primitive
rank-one image, the circle generator of the
exterior. The map on $H_3$ is injective with
primitive image, since each of the two filling
$S^3$ summands receives its own identity map.
The map on $H_4$ has primitive rank-one image
generated by $\delta\times[S^1]$. Every kernel
and cokernel used in the sequence is therefore
free, and it gives
\[
       \operatorname{rank}H_k(M_5,\Z)=(1,2,5,5,2,1).
                                                               \tag{C.PE.19}\hypertarget{tag-00UC}{}\label{cdp:eq:C.PE.19}
\]
K\"unneth and \eqref{cdp:eq:C.PE.16} give the $Y_m$ row of
\eqref{cdp:eq:C.PE.3}, including integral torsion-freeness.

\item Each Hopf surface is differentiably $S^3\times S^1$.
For the support use the actual union
$S_m=P_m\cup Q_+\cup Q_-$ with intersections
$J_+\sqcup J_-$. In the preceding $M_3$ presentation,
the core angular loop of $J_+$ is $\delta_+$ and
that of $J_-$ is $-\delta_+$. Their radial loops
are the $S^1$ factor of \eqref{cdp:eq:C.PE.16}. Inclusion in
$Q_\pm$ kills the angular loop and takes the
radial loop to its Hopf generator. Consequently
the integral map
\[
 \begin{split}
 H_1(J_+\sqcup J_-,\Z)\longrightarrow
 H_1(P_m,\Z)\oplus H_1(Q_+,\Z)\oplus H_1(Q_-,\Z)
 \end{split}                                             \tag{C.PE.20}\hypertarget{tag-00UE}{}\label{cdp:eq:C.PE.20}
\]
has primitive rank three: its columns are, up to
orientation signs, $\delta_+$, $z-q_+$,
$-\delta_+$, $z-q_-$. On $H_2$, the two elliptic
fundamental classes map to the same primitive
fibre class of $P_m$ and to zero on the Hopf
surfaces. Primitivity follows directly from
$[J_\pm]=\pm\delta_+\times[S^1]$ with the appropriate
product orientations. More explicitly, the dual chart at
$J_-$ reverses both its angular loop and its radial loop:
its complex orientation is
$(-\delta_+)\times(-z)=\delta_+\times z$.
Using the same global radial generator $z$ at both
ends instead reverses the displayed product orientation
at $J_-$. The restriction map on
cohomology thus has rank one in degree two and
rank three in degree one. Mayer--Vietoris gives
the free support ranks $(1,3,6,7,3)$ and Euler zero.

\item The coherent groups also retain exact maps.
On every chart \eqref{cdp:eq:C.PE.8}, finite direct image and
the exact finite-group invariant functor give
\[
 \rho_*\OO_{Y_m}=\OO_V,\qquad
 R^1\rho_*\OO_{Y_m}=\OO_V,\qquad
 R^j\rho_*\OO_{Y_m}=0\quad(j\ge2).                       \tag{C.PE.21}\hypertarget{tag-00UF}{}\label{cdp:eq:C.PE.21}
\]
Translations act trivially on $H^1(E,\OO_E)$.
The trivialization of $R^1$ is global: the
constant cocycle assigning $1$ to the generator
\eqref{cdp:eq:C.PE.7} gives a class restricting to the
nonzero logarithmic cocycle on each elliptic
fibre, including the multiple fibres in the
finite-quotient chart. It is a nowhere-zero
frame. Since $h^k(V,\OO_V)=(1,1,0)$, Leray gives
$(1,2,1,0)$ for $Y_m$. The same local calculation
on $P_m\to E$ gives $(1,2,1)$ there, consistent
with \eqref{cdp:eq:C.PE.2}.

\item For a primary Hopf surface, a holomorphic
two-form lifts and extends to $\C^2$; invariance
under its strict scalar contraction forces
every power-series coefficient to vanish.
Hence $p_g=0$. Its topological $e=0$ and
$H^2(-,\Z)=0$ give $\chi(\OO)=0$ by Noether's
formula, so its coherent dimensions are
$(1,1,0)$. Restriction
$H^1(Q_\pm,\OO)\to H^1(J_\pm,\OO)$ is nonzero:
the deck cocycle $1$ restricts to the same
cocycle on $\C^*/q^{m\Z}$. It is not a
coboundary, since the constant Laurent
coefficient of $g(q^m u)-g(u)$ is zero for
every holomorphic $g$ on $\C^*$.
The coherent normalization sequence for
$S_m$ therefore has a surjective difference
map in degrees zero and one. Its remaining
terms give $h^k(S_m,\OO)=(1,2,1)$.

\item Finally, a rational curve in $P_m$ would lift
through the finite \'etale map \eqref{cdp:eq:C.PE.12}, since
$\PP^1$ is simply connected. Its two projections
to the positive-genus curves would be constant.
The Hopf components have no rational curve
either: a map from $\PP^1$ lifts to
$\C^2\setminus0$, and its coordinate functions
are constant. The inverse-map resolution
comparison in \ref{cdp:st:C.2.UN.3} prevents contraction
of any of these surfaces in a bimeromorphic
map between smooth threefold fillings. A
finite base change $t=s^d$ is generically,
along $P_m$, the graph $x=s^d$ for the local
transverse coordinate $x=t$. Its normalized
strict component maps with degree one to
$P_m$; resolution preserves that birational
component. The distinct plurigenus slopes
$2(1-1/m)$ distinguish the unique properly
elliptic components, and hence the fillings
for different $m$.

\item This is an explicit signed logarithmic-transform
example. The prior Kodaira-dimension-one
occurrences discussed after \ref{cdp:st:C.2.UN.6}, including
\cite[Example~3.4(1)]{Fujimoto1988Log}, are retained
as the literature comparison.

\end{enumerate}
\end{proof}






\begin{proposition}[A free half-period quotient realizes secondary Kodaira normalization]\hypertarget{tag-01CZ}{}\label{cdp:st:C.2.UN.21}
\StatementRecord{cdp:st:C.2.UN.21}{proved}{01CZ}{\SeriesRef{cdp:st:C.2.UN.3}, \SeriesRef{cdp:eq:C.SKD.4}, \SeriesRef{cdp:eq:C.SKD.5}, \SeriesRef{cdp:eq:C.SKD.7}, \SeriesRef{cdp:eq:C.SKD.9}, \SeriesRef{cdp:eq:C.SKD.6}, \SeriesRef{cdp:st:C.2.UN.18}, \SeriesRef{cdp:eq:C.SKD.12}}

Let $q=e^{-2\pi}$ and $E=\C^*/q^{\Z}$, and use the
origin of its group structure. There is a smooth compact
threefold $Y\to\PP^1$ whose smooth general fibre is
$E\times E$ and whose two special full fibres are
\[
                        2P+2Q\quad\hbox{and}\quad 2B.
                                                            \tag{C.SKD.1}\hypertarget{tag-011N}{}\label{cdp:eq:C.SKD.1}
\]
The surface $P$ is a secondary Kodaira surface, $Q$ is
a primary Hopf surface, and they meet transversely along
a smooth elliptic curve $J$. The surface $B$ is smooth
bielliptic. Both reduced special supports have Euler
number zero. In particular $P$ occurs as an essential
component in a zero-Euler torus filling.

The canonical bundle of $P$ has exact order two. For
$S=P\cup Q$, the rational cohomology ranks and coherent
dimensions are
\[
 \begin{array}{c|c|c}
 &\dim H^k(-,\Q)&h^k(\OO)\\ \hline
 P&(1,1,0,1,1)&(1,1,0)\\
 Q&(1,1,0,1,1)&(1,1,0)\\
 J&(1,2,1)&(1,1)\\
 S&(1,1,1,3,2)&(1,1,0)\\
 B&(1,2,2,2,1)&(1,1,0)\\
 Y&(1,1,2,4,2,1,1)&(1,1,0,0).
 \end{array}                                             \tag{C.SKD.2}\hypertarget{tag-011U}{}\label{cdp:eq:C.SKD.2}
\]
The full schemes in \eqref{cdp:eq:C.SKD.1}, in contrast to their
reductions, each have coherent dimensions $(1,2,1)$.
The primary Kodaira double cover of $P$ is recovered
by the actual quadratic base change at the first
special value.
\end{proposition}


\begin{proof}

\begin{enumerate}[label={$\langle1\rangle$\arabic*.},leftmargin=*]
\item Start with the compact version of \ref{cdp:st:C.2.UN.3} at $n=1$:
\[
 V=\operatorname{Bl}_{(0,0)}(E\times\PP^1_t),\qquad
 L=\OO_V(R),\qquad
 \widetilde Y=L^\times/q^{\Z},
                                                            \tag{C.SKD.3}\hypertarget{tag-011V}{}\label{cdp:eq:C.SKD.3}
\]
where $R$ is the exceptional curve. Let $\sigma$ be
the canonical section of $L$. The projection
$\widetilde Y\to V\to\PP^1_t$ is proper and smooth
off its reduced central divisor
$\widetilde P+\widetilde Q$. The first component is
the primary Kodaira surface
$\OO_E([0])^\times/q^{\Z}$, the second is the Hopf
surface $\OO_{\PP^1}(-1)^\times/q^{\Z}$, and their
intersection is an elliptic curve. These are the
actual line-bundle quotient maps of \ref{cdp:st:C.2.UN.3}.

\item The involution $(p,t)\mapsto(-p,-t)$ fixes the blowup
centre and lifts to an involution $\iota$ of $V$.
It has the canonical order-two lift $\widehat\iota$
to $L$ preserving $\sigma$. Define an automorphism
of $\widetilde Y$ by
\[
             g[\ell]=[q^{1/2}\widehat\iota(\ell)].
                                                            \tag{C.SKD.4}\hypertarget{tag-011W}{}\label{cdp:eq:C.SKD.4}
\]
Its square is multiplication by $q$, so $g^2=1$
on $\widetilde Y$. It acts freely. Indeed at a
fixed base point the order-two linearization on
the line $L_v$ has eigenvalue $+1$ or $-1$.
Neither $q^{1/2}$ nor $-q^{1/2}$ belongs to
$q^{\Z}$. Thus \eqref{cdp:eq:C.SKD.4} has no fixed point on
any elliptic fibre, including those over $R$.
The quotient
\[
              \pi:\widetilde Y\longrightarrow Y
                        =\widetilde Y/\langle g\rangle
                                                            \tag{C.SKD.5}\hypertarget{tag-011X}{}\label{cdp:eq:C.SKD.5}
\]
is finite \'etale of degree two, and $Y$ is smooth
and compact. The function $T=t^2$ defines its map
to $\PP^1_T$. For $T\ne0,\infty$, the two values
of $t$ are interchanged by $g$, so each quotient
fibre is exactly one of the original product
tori.

\item The strict central base curve and $R$ are
invariant. The local equation upstairs at their
intersection is $t=xy$, hence $T=x^2y^2$.
Since \eqref{cdp:eq:C.SKD.5} is \'etale, this proves the first
full fibre and the normal-crossing assertion
in \eqref{cdp:eq:C.SKD.1}. At infinity the original smooth
fibre is $E\times E$, and the action is base
inversion on the first elliptic factor and
translation by the class of $q^{1/2}$ on the
second. It is free. Its quotient is a
bielliptic surface $B$, and the transverse
coordinate is squared, giving the full fibre
$2B$.

\item We identify the central quotients more precisely.
The canonical linearization on
$L|_R=\OO_{\PP^1}(-1)$ acts by $-1$ on its
tautological vector: the derivative at the
blown-up point is the scalar $-I$ on $\C^2$.
Thus the quotient of $\widetilde Q$ is
\[
                 Q=(\C^2\setminus0)/
                         (-q^{1/2})^{\Z},              \tag{C.SKD.6}\hypertarget{tag-011Y}{}\label{cdp:eq:C.SKD.6}
\]
again a primary Hopf surface. On $\widetilde P$
the involution covers inversion on $E$.
The nowhere-zero canonical form of the primary
Kodaira surface is the wedge of the base
one-form and the fibre logarithmic differential.
The first changes sign; the second does not.
Consequently $g$ acts by $-1$ on that form.
Its square descends to a nowhere-zero section
of $K_P^{\otimes2}$, while no canonical form
descends. Thus $K_P$ has exact order two.
The finite \'etale cover $\widetilde P\to P$
and the computed $b_1(P)=1$ below identify $P$
as a secondary Kodaira surface.

\item We give the topology through actual mapping
tori. Choose an $\iota$-invariant Hermitian
metric on $L$, by averaging. Let $M_5$ be its
unit circle bundle. The positive-real norm
coordinate gives
$\widetilde Y\simeq M_5\times S^1$. In that
description \eqref{cdp:eq:C.SKD.4} is
$(x,u)\mapsto(\iota_Mx,u+\tfrac12)$, where
the period of $u$ is one. Hence the quotient
is the mapping torus
\[
          Y\simeq T(\iota_M)
            =(M_5\times\R)/
               ((x,u+1)\sim(\iota_Mx,u)).               \tag{C.SKD.7}\hypertarget{tag-011Z}{}\label{cdp:eq:C.SKD.7}
\]
For example the map $[x,u]\mapsto[x,2u]$
identifies the two quotient descriptions.

\item The integral circle-bundle Gysin sequence
over $V$ has Euler class $[R]$. The base has
free cohomology of ranks $(1,2,3,2,1)$.
Cup product with $[R]$ has ranks $1,0,1$
in degrees zero, one and two; the nonzero
images are primitive because $R^2=-1$.
Thus $M_5$ has free homology of ranks
\[
                         (1,2,4,4,2,1).                \tag{C.SKD.8}\hypertarget{tag-0120}{}\label{cdp:eq:C.SKD.8}
\]
Its involution acts as $-I$ on $H_1$ and,
by oriented Poincar\'e duality, on $H_4$.
On $H_2$ it has an integral diagonal
decomposition with two $+1$ and two $-1$
summands. To specify that decomposition,
lift $E\times\{\infty\}$ and
$\{p_*\}\times\PP^1$, where $p_*\ne0$ is
a two-torsion point, using the normalized
canonical section of $L$. These are
invariant cycles, avoid $R$, and generate
the two base summands. The preimages of
the two elliptic-base loops at infinity
are circle tori; they generate the two
remaining summands and are sent to their
negatives. The homological Gysin sequence
identifies these four cycles with an
integral basis. Duality gives the same
two-positive, two-negative decomposition
on $H_3$. The actions on $H_0,H_5$ are
the identity. The Wang sequence of
\eqref{cdp:eq:C.SKD.7} therefore gives the rational
$Y$ row of \eqref{cdp:eq:C.SKD.2}.

\item This also determines its integral groups,
not only their ranks. Writing $T_2=(\Z/2)^2$,
the integral cohomology of $Y$ is
\[
 \begin{array}{c|ccccccc}
 k&0&1&2&3&4&5&6\\ \hline
 H^k(Y,\Z)&
 \Z&\Z&\Z^2\oplus T_2&\Z^4\oplus T_2&
 \Z^2\oplus T_2&\Z\oplus T_2&\Z .
 \end{array}                                             \tag{C.SKD.9}\hypertarget{tag-0121}{}\label{cdp:eq:C.SKD.9}
\]
Indeed the Wang maps are zero or multiplication
by $-2$ on the stated integral summands.
Their kernel terms are free, so the short
exact sequences split as abelian groups;
the universal coefficient sequence gives
\eqref{cdp:eq:C.SKD.9}.

\item On $\widetilde P$, the unit circle bundle
$M_3\to E$ has Euler number one and free
homology ranks $(1,2,2,1)$. Its involution
acts by $-I$ in degrees one and two and
by $+1$ in degrees zero and three.
The quotient $P$ is its mapping torus.
The same Wang maps give
\[
 \begin{array}{c|ccccc}
 k&0&1&2&3&4\\ \hline
 H^k(P,\Z)&
 \Z&\Z&T_2&\Z\oplus T_2&\Z .
 \end{array}                                             \tag{C.SKD.10}\hypertarget{tag-011O}{}\label{cdp:eq:C.SKD.10}
\]
In particular its rational ranks are those
of \eqref{cdp:eq:C.SKD.2}. The Hopf ranks follow from
\eqref{cdp:eq:C.SKD.6}, which is differentiably
$S^3\times S^1$.

\item For the support use the exact union
$S=P\cup Q$ with intersection $J$. The
radial deck loop of $J$ maps to the
mapping-torus generator on $P$ and the
Hopf generator on $Q$. Its angular loop
maps to zero in rational first homology:
upstairs it is the circle of an
Euler-number-one bundle, hence is already
zero in the abelianization. Therefore
the map
\[
           H^1(P,\Q)\oplus H^1(Q,\Q)
                         \longrightarrow H^1(J,\Q)
                                                            \tag{C.SKD.11}\hypertarget{tag-011P}{}\label{cdp:eq:C.SKD.11}
\]
has rank one. Both surfaces have zero
rational $H^2$, so the next restriction
map has rank zero. Mayer--Vietoris gives
the support ranks $(1,1,1,3,2)$.
There is also an exact integral comparison at the only
potential torsion restriction. In the unit circle bundle
$M_3\to E$, the core fibre over zero bounds an
$\iota_M$-invariant oriented surface. To construct it,
take the normalized canonical section over the complement
of a small invariant disc at zero. Choose a local
coordinate and an equivariant frame with
$\sigma=z e$; the lift is
$(z,u)\mapsto(-z,-u)$ in the unit-circle coordinates.
The section on the disc boundary is
$(z,u)=(\varepsilon e^{i\theta},e^{i\theta})$.
Attach the annulus
\[
 (r,\theta)\longmapsto(re^{i\theta},e^{i\theta}),
                 \qquad 0\le r\le\varepsilon .
\]
Its inner boundary is the core circle. The involution
is $\theta\mapsto\theta+\pi$ on this annulus and
preserves its orientation; it also preserves the
orientation of the section over the punctured elliptic
curve. The joined surface, or its invariant rounded
chain, therefore has an oriented mapping torus whose
boundary is $J$. Thus
$H_2(J,\Z)\to H_2(P,\Z)$ is zero integrally,
including its possible torsion target. The angular
first-homology class is already zero in $M_3$;
the other generator maps to a primitive mapping-torus
class on $P$ and to a primitive Hopf class on $Q$.
The integral Mayer--Vietoris sequence now gives
\[
 \begin{array}{c|ccccc}
 k&0&1&2&3&4\\ \hline
 H_k(S,\Z)&
 \Z&\Z\oplus T_2&\Z\oplus T_2&\Z^3&\Z^2\\
 H^k(S,\Z)&
 \Z&\Z&\Z\oplus T_2&\Z^3\oplus T_2&\Z^2 .
 \end{array}                                             \tag{C.SKD.14}\hypertarget{tag-011S}{}\label{cdp:eq:C.SKD.14}
\]
The free kernel terms split as abelian groups.
Any torsion summand in the radial image is absorbed
by eliminating the primitive Hopf generator, so
no choice of radial lift changes these groups.
The invariant cohomology of the free
quotient of $E\times E$ at infinity
gives the row for $B$. In both cases
the displayed Euler number is zero.
For completeness, $B$ is also the product of an
angular circle with the mapping torus of inversion
on the first elliptic curve. The Wang maps on
that curve are $0,-2I_2,0$. Its integral
cohomology is therefore
\[
 H^k(B,\Z)=
 \bigl(\Z,\ \Z^2,\ \Z^2\oplus T_2,\
                  \Z^2\oplus T_2,\ \Z\bigr)_k .
                                                            \tag{C.SKD.15}\hypertarget{tag-011T}{}\label{cdp:eq:C.SKD.15}
\]
The integral groups of $J$ and $Q$ are free with
the ranks already displayed.

\item For coherent cohomology, the elliptic
line-bundle quotient $\widetilde Y\to V$
has direct images $\OO_V$ in degrees
zero and one and zero above one. The
degree-one frame is the constant
logarithmic deck cocycle, as in
(C.PE.21). Its involution action is
$+1$, while the action on
$H^1(V,\OO_V)=H^1(E,\OO_E)$ is $-1$.
The finite \'etale invariant-sheaf
comparison for \eqref{cdp:eq:C.SKD.5}, applied to
Leray's filtration, gives
$h^k(Y,\OO_Y)=(1,1,0,0)$. The identical
calculation for $\widetilde P\to E$
gives $h^k(P,\OO_P)=(1,1,0)$, consistent
with the canonical character.
The Hopf and elliptic rows were
computed from their actual deck cocycles
in \ref{cdp:st:C.2.UN.18}. In particular
$H^1(Q,\OO_Q)\to H^1(J,\OO_J)$ is
nonzero and hence onto. The coherent
normalization sequence gives
$h^k(S,\OO_S)=(1,1,0)$. The quotient
at infinity has the same base-negative,
fibre-positive action, giving
$h^k(B,\OO_B)=(1,1,0)$.

\item We finally verify the full scheme fibres.
Upstairs let $\widetilde S$ be the
reduced central fibre of the original
map to $t$. Its coherent normalization
sequence gives $(1,2,1)$, the same
dimensions as a nearby product torus.
Grauert's theorem therefore makes all
$R^qf_*\OO_{\widetilde Y}$ locally free
near $t=0$, of ranks $1,2,1$, with
base change. The actual inverse image
of the full fibre $2S$ under the
\'etale map \eqref{cdp:eq:C.SKD.5} is $2\widetilde S$,
defined by $t^2=0$. Thus
\[
 H^q(2S,\OO_{2S})
  \simeq
 \left((R^qf_*\OO_{\widetilde Y})/
                   t^2(R^qf_*\OO_{\widetilde Y})\right)^g .
                                                            \tag{C.SKD.12}\hypertarget{tag-011Q}{}\label{cdp:eq:C.SKD.12}
\]
Choose a basis of eigenvectors on the
central fibre of each direct-image
bundle. Lift each vector and average
with $g$ with its specified sign; the
resulting holomorphic sections remain
a frame near zero and satisfy
$g(e)=\pm e$. Since $g(t)=-t$, the
invariants in each summand
$\C[t]/(t^2)e$ have dimension one:
they are generated by $e$ for sign
$+1$ and by $te$ for sign $-1$.
Equation \eqref{cdp:eq:C.SKD.12} gives precisely
$(1,2,1)$. The same argument at
infinity proves the assertion for
$2B$. This comparison retains both
the nilpotent direction and its sign.

\item No rational curve lies in $P$:
its map from $\PP^1$ would lift to
the primary Kodaira cover, whose
elliptic base and elliptic fibres
force that map to be constant.
The same lifting argument to
$\C^2\setminus0$ treats $Q$.
The inverse-map resolution comparison
of \ref{cdp:st:C.2.UN.3} therefore proves their
essentiality. At $T=0$, the finite
map
\[
 \widetilde Y\longrightarrow
             Y\times_{\PP^1_T}\PP^1_t,\qquad
 y\longmapsto(\pi(y),f(y)),\quad T=t^2,
                                                            \tag{C.SKD.13}\hypertarget{tag-011R}{}\label{cdp:eq:C.SKD.13}
\]
is generically one-to-one and has
normal source. It identifies
$\widetilde Y$ with the normalization
of the reduced fibre product.
Consequently the primary Kodaira
component is recovered by this
actual quadratic base change.
The construction supplies an
explicit secondary Kodaira component;
no birational persistence of its
secondary type under that base change
is asserted.

\end{enumerate}
\end{proof}






\begin{proposition}[An essential Enriques surface occurs in a compact zero-Euler torus filling]\hypertarget{tag-01D3}{}\label{cdp:st:C.2.UN.25}
\StatementRecord{cdp:st:C.2.UN.25}{proved}{01D3}{\SeriesRef{cdp:st:C.2.UN.14}, \SeriesRef{cdp:st:C.2.UN.18}, \SeriesRef{cdp:eq:C.EN.7}, \SeriesRef{cdp:eq:C.EN.8}, \SeriesRef{cdp:st:C.2.UN.23}, \SeriesRef{cdp:eq:C.EN.15}, \SeriesRef{cdp:st:C.2.UN.4}, \SeriesRef{cdp:eq:C.EN.S2}, \SeriesRef{cdp:eq:C.EN.S1}, \SeriesRef{cdp:eq:C.EN.S4}; \cite[Theorem~2.1 and Example~2.3]{Yoshihara1995}}

There is a smooth compact complex threefold over
$\PP^1$ with connected two-torus general fibre and
a simple-normal-crossing special support of Euler
number zero containing an unchanged smooth Enriques
surface as an irreducible component.
Before six explicit vertical modifications, its full
central fibre is
\[
 2M+2A+4Q_1+4Q_2+E_1+E_2+E_3+E_4,                 \tag{C.EN.1}\hypertarget{tag-00G9}{}\label{cdp:eq:C.EN.1}
\]
where $M$ is Enriques, $A\simeq\PP^1\times\PP^1$,
$Q_1,Q_2$ are primary Hopf surfaces, and every
$E_i$ is the Hirzebruch surface $\mathbb F_2$.
That reduced support has Euler number $12$.
All six modifications can be made disjoint from $M$.
The Enriques component is therefore essential in the
sense of \ref{cdp:st:C.2.UN.14}.
The reduced supports before and after balancing have
respectively
\[
 \begin{array}{c|c|c}
 & (b_0,b_1,b_2,b_3,b_4)&(h^0(\OO),h^1(\OO),h^2(\OO))\\ \hline
 S&(1,1,8,4,8)&(1,1,0)\\
 S_{\rm bal}&(1,1,20,40,20)&(1,1,0).
 \end{array}                                             \tag{C.EN.18}\hypertarget{tag-00GI}{}\label{cdp:eq:C.EN.18}
\]
\end{proposition}


\begin{proof}

\begin{enumerate}[label={$\langle1\rangle$\arabic*.},leftmargin=*]
\item We first give a compact base with specified local
coordinates. Take the Weierstrass elliptic surface
over $\PP^1_T$
\[
 y^2=x^3+(T^4-3)x+(T^6+T+2).                       \tag{C.EN.2}\hypertarget{tag-00GJ}{}\label{cdp:eq:C.EN.2}
\]
The coefficients are sections of $\OO(4)$ and
$\OO(6)$. Its discriminant, up to the usual nonzero
constant factor, is
\[
 \begin{split}
 \Delta(T)={}&31T^{12}-36T^8+54T^7+108T^6\\
             &+108T^4+27T^2+108T .
 \end{split}                                       \tag{C.EN.3}\hypertarget{tag-00GK}{}\label{cdp:eq:C.EN.3}
\]
The exact Euclidean algorithm gives
$\gcd(\Delta,\Delta')=1$; also
$\Delta'(0)=108$ and its leading coefficient is
$31$. These checks are saved in
\texttt{code/E\_enriques\_filling.py}.
Thus every singular fibre is nodal of type $I_1$,
the fibre at zero is nodal, and infinity is smooth.
The total Weierstrass surface is smooth.

\item Here is an explicit chart at the node of the
zero fibre. Put $\xi=x-1$ and
$U=1+T^3(1+\xi)+T^5$. On the surface,
\[
 X=\frac{y-\xi\sqrt{3+\xi}}{U},\qquad
 Y=y+\xi\sqrt{3+\xi}
 \quad\Longrightarrow\quad XY=T .                  \tag{C.EN.4}\hypertarget{tag-00GL}{}\label{cdp:eq:C.EN.4}
\]
The coordinate Jacobian at $\xi=y=T=0$ is
nonzero, so these are actual local coordinates.
Make the quadratic base change $T=t^2$ and
resolve $XY=t^2$. One resolution chart is
\[
       (X,Y,t)=(u,uv^2,uv),\qquad
       h(u,v)=(u,-v),                              \tag{C.EN.5}\hypertarget{tag-00GM}{}\label{cdp:eq:C.EN.5}
\]
where $h$ is the deck involution. The other chart
interchanges $X,Y$ and replaces $v$ by $v^{-1}$.
The central fibre of the smooth base surface $V_0$
is $R_0+R_1$ of type $I_2$. Its strict nodal
normalization $R_0\simeq\PP^1$ is fixed pointwise
by $h$. The exceptional curve $R_1\simeq\PP^1$
has coordinate $z=v$, with $h(z)=-z$, and
meets $R_0$ exactly at $z=0,\infty$.

\item Blow up the four points $z=1,2,-1,-2$ of $R_1$.
Write their exceptional curves as
$D_1,D_2,D_{-1},D_{-2}$, and denote the resulting
compact surface by $V$. The involution lifts
and exchanges $D_j,D_{-j}$. Set
\[
 D=D_1+D_2-D_{-1}-D_{-2},\qquad L=\OO_V(D).
                                                               \tag{C.EN.6}\hypertarget{tag-00GN}{}\label{cdp:eq:C.EN.6}
\]
The canonical meromorphic section $\sigma$ of
$L$ supplies the section-preserving identification
$h^*L=L^{-1}$. Its square is the identity.

\item Apply the explicit signed root construction of
\ref{cdp:st:C.2.UN.18} with order $m=2$ at these four disjoint
curves. Thus on the complement of the negative
curves use $\sigma=a^2\ell$, with $\ell\in L^\times$;
on the complement of the positive curves use
$\sigma^{-1}=(a')^2\ell'$, where
$a'=a^{-1}$ and $\ell'=\ell^{-1}$. Divide by
\[
       (\ell,a)\longmapsto(q^2\ell,q^{-1}a),
                  \qquad 0<q<1 .                  \tag{C.EN.7}\hypertarget{tag-00GO}{}\label{cdp:eq:C.EN.7}
\]
Every local root chart is the already verified
free finite quotient of an elliptic product in
(C.PE.8), so the resulting compact threefold
$\widetilde Y$ is smooth and proper over $V$.
Away from the four exceptional curves the
actual product coordinate is
$\ell=\sigma c^2,\ a=c^{-1}$, with
$[c]\in E=\C^*/q^{\Z}$.

\item Dualization gives an involution of
$\widetilde Y$: it takes the global meromorphic
coordinates $(\ell,a)$ to $(\ell^{-1},a^{-1})$
over $h$. Equivalently, a positive chart maps
to the target negative chart with unchanged
numerical coordinates $(\ell',a')=(\ell,a)$.
It conjugates \eqref{cdp:eq:C.EN.7} to its inverse, so
descends to the quotient and still squares
to one. In the product coordinate it is
\[
                       (v,[c])\longmapsto(hv,[c^{-1}]).
                                                               \tag{C.EN.8}\hypertarget{tag-00GP}{}\label{cdp:eq:C.EN.8}
\]
It covers $t\mapsto-t$.

\item We identify the component over the strict
transform of $R_1$. The restriction of $L$
there is trivial; in a compatible frame,
\[
 \sigma|_{R_1}=\phi(z)
   =\frac{(z-1)(z-2)}{(z+1)(z+2)},\qquad
                    \phi(-z)=\phi(z)^{-1}.          \tag{C.EN.9}\hypertarget{tag-00GQ}{}\label{cdp:eq:C.EN.9}
\]
Let $C$ be the smooth elliptic curve
\[
 C:\quad Y^2=(z^2-1)(z^2-4),\qquad
 y=\frac{Y}{(z+1)(z+2)},\quad y^2=\phi(z).
                                                               \tag{C.EN.10}\hypertarget{tag-00GA}{}\label{cdp:eq:C.EN.10}
\]
Its four simple branch values are $1,2,-1,-2$.
The exact finite product-cover map from \ref{cdp:st:C.2.UN.18} is
\[
 C\times E\longrightarrow B,\qquad
 (z,y,[u])\longmapsto[\ell=u^2,\ a=y/u].
                                                               \tag{C.EN.11}\hypertarget{tag-00GB}{}\label{cdp:eq:C.EN.11}
\]
It is \'etale, including at the four root values,
by the two Jacobians in (C.PE.13). Its deck
involution is
\[
 \eta:(z,y,[u])\longmapsto(z,-y,[-u]).
                                                               \tag{C.EN.12}\hypertarget{tag-00GC}{}\label{cdp:eq:C.EN.12}
\]
The translation $[u]\mapsto[-u]$ is a nonzero
two-torsion translation, so $\eta$ is free.
Thus $B=(C\times E)/\langle\eta\rangle$ is
bielliptic. The involution \eqref{cdp:eq:C.EN.8} lifts to
\[
 k:(z,y,[u])\longmapsto(-z,y^{-1},[u^{-1}])
       \quad\text{or}\quad
       (z,Y,[u])\longmapsto(-z,Y,[u^{-1}]).
                                                               \tag{C.EN.13}\hypertarget{tag-00GD}{}\label{cdp:eq:C.EN.13}
\]
These formulas commute with $\eta$.

\item On $C$, the involution $j:(z,Y)\mapsto(z,-Y)$
and $k_C:(z,Y)\mapsto(-z,Y)$ each act by $-1$
on $dz/Y$. The latter has four fixed points,
the two points over each of $0,\infty$.
Their product $jk_C$ has no fixed point:
over zero its equation would require $Y=0$,
and at infinity it exchanges the two points.
It is consequently a nonzero two-torsion
translation of the elliptic curve $C$.
It follows that $\eta$ and $\eta k$ are
free on $C\times E$, whereas $k$ has exactly
$4\cdot4=16$ fixed points. They occur in
eight free $\eta$-orbits. At each, the
derivative of $k$ is $(-1,-1)$.

\item The quotient $B/\langle h\rangle$ therefore
has exactly eight $A_1$ singularities and
no branch curve. Its minimal resolution
$M$ is Enriques. One exact verification
uses the intermediate quotient
$(C\times E)/\langle k\rangle$. Its
minimal resolution is K3 by \ref{cdp:st:C.2.UN.23}.
The involution $\eta$ extends over its
sixteen exceptional curves, permuting
them in pairs. It is free away from
them because both $\eta$ and $\eta k$
are free upstairs, and free on them
because they are paired. Its quotient
is precisely $M$. Equivalently, the
canonical form has character $-1$
under $\eta$ and $+1$ under $k$;
its square descends through the
explicit crepant $A_1$ resolutions,
while $p_g(M)=q(M)=0$. Thus $K_M$
has exact order two and $e(M)=12$.
This is an actual noncyclic product
quotient, consistent with the two
Enriques holonomy cases of
\cite[Theorem~2.1 and Example~2.3]{Yoshihara1995}.

\item It remains to realize this surface resolution
inside the threefold, not just on an abstract
component. Near $R_0$ the four root curves
are absent, so \eqref{cdp:eq:C.EN.8} is a product.
Its fixed locus over zero consists of four
disjoint copies of $R_0$, indexed by $E[2]$.
Near either node of $R_0+R_1$, choose a
local coordinate $w$ at the corresponding
two-torsion point of $E$, so inversion is
$w\mapsto-w$. The ambient action is
\[
 (u,v,w)\longmapsto(u,-v,-w),\qquad
                         T=u^2v^2 .
                                                               \tag{C.EN.14}\hypertarget{tag-00GE}{}\label{cdp:eq:C.EN.14}
\]
Its quotient has invariant coordinates
$A_0=v^2,B_0=vw,C_0=w^2$, with
$A_0C_0=B_0^2$ and $T=u^2A_0$.
The two smooth resolution charts are
\[
 \begin{array}{lll}
 A_0=r,&B_0=rs,&C_0=rs^2,\quad T=u^2r,\\
 A_0=r's'^2,&B_0=r's',&C_0=r',\quad T=u^2r's'^2 .
 \end{array}                                             \tag{C.EN.15}\hypertarget{tag-00GF}{}\label{cdp:eq:C.EN.15}
\]
They glue by $s'=s^{-1}$ and $r'=rs^2$.
Globally they resolve the four fixed
curves. Restricting to $u=0$ is exactly
the minimal $A_1$ resolution of the
eight surface singularities just found.
In particular the strict central component
is the smooth Enriques surface $M$.

\item The component over $R_0$ becomes
$A=R_0\times(E/\{\pm1\})\simeq\PP^1\times\PP^1$.
Each exceptional surface of \eqref{cdp:eq:C.EN.15}
is the projective normal-direction bundle
of a fixed curve. That normal bundle is
$\OO_{R_0}(-2)\oplus\OO_{R_0}$: the
product trivialization is valid on an
entire neighbourhood of $R_0$, and
$R_0^2=-2$. Hence these surfaces are
$E_i\simeq\mathbb F_2$.
The four root Hopf components are
exchanged in two pairs and give
$Q_1,Q_2$. Squaring the base parameter,
and the exact monomials in \eqref{cdp:eq:C.EN.15},
give every coefficient in \eqref{cdp:eq:C.EN.1}.
Those same monomials show that its
reduced support is simple normal
crossing.

\item The full incidence calculation is
\[
 \begin{array}{c|c}
 \text{intersection}&\text{reduced space}\\ \hline
 M\cap A&\PP^1\sqcup\PP^1\\
 M\cap Q_j&\text{one elliptic curve}\quad(j=1,2)\\
 A\cap E_i&\PP^1\quad(i=1,\ldots,4)\\
 M\cap E_i&\PP^1\sqcup\PP^1\quad(i=1,\ldots,4)\\
 M\cap A\cap E_i&\text{two points}\quad(i=1,\ldots,4).
 \end{array}                                             \tag{C.EN.16}\hypertarget{tag-00GG}{}\label{cdp:eq:C.EN.16}
\]
There are no other nonempty intersections.
Thus inclusion--exclusion gives, with
every nonzero term displayed,
\[
 e(S)=12+4+4\cdot4-(4+4\cdot2+4\cdot4)+8=12 .
                                                               \tag{C.EN.17}\hypertarget{tag-00GH}{}\label{cdp:eq:C.EN.17}
\]

\item For the compact construction, the base involution
also fixes the smooth elliptic fibre at infinity.
The ambient fixed loci there are four elliptic
curves, and the same transverse $A_1$ charts
resolve them. Their exceptional surfaces and
all intersections are elliptic ruled and
have Euler number zero. At the remaining
nodal base values the elliptic-product
coordinate gives a nodal curve times $E$,
also of Euler number zero. The quotient
and the resolutions are proper and compact.
Over a general $T$, the two $t$-sheets
are exchanged by $h$, so their quotient
is exactly one product of two elliptic
curves. The Stein factor has generic
degree one over the normal curve
$\PP^1$, proving connectedness.

\item Finally perform six of the point-then-quartic
operations of \ref{cdp:st:C.2.UN.4}, all over $T=0$ and
away from $M$. They can begin at fresh
smooth points of an $E_i$ outside its
intersections. Their exceptional plane
quartics and subsequent centres remain
disjoint from $M$. Each changes the
reduced Euler number by $-2$, preserves
the smooth total space and the normal
crossing support, and leaves the
punctured family unchanged. The resulting
support has Euler number zero, with
the same smooth Enriques component.
Its nonzero even pluricanonical form
precludes a covering family of rational
curves, so it is non-uniruled.
The exact inverse-map resolution argument
of \ref{cdp:st:C.2.UN.14} proves essentiality. This completes
the compact occurrence construction.
\par\medskip\noindent\textbf{The reduced support and its exact cohomology.}
Let $S$ denote the reduced support before the six balancing
operations. Label the Enriques component $M$, the component over
$R_0$ by $A\simeq\PP^1\times\PP^1$, the two Hopf components by
$Q_a,Q_b$, and the four exceptional ruled surfaces by
$E_1,\ldots,E_4\simeq\mathbb F_2$.
The actual quotient and resolution charts give
\[
\begin{array}{c|c|c}
\text{components}&\text{double curve}&\text{number}\\ \hline
M,A&D_0,D_\infty\simeq\PP^1&2\\
M,Q_j&J_j\text{ elliptic},\ j=a,b&2\\
A,E_i&Z_i\simeq\PP^1&4\\
M,E_i&U_{i0},U_{i\infty}\simeq\PP^1&8.
\end{array}                                                   \tag{C.EN.S1}\hypertarget{tag-00H7}{}\label{cdp:eq:C.EN.S1}
\]
There are precisely eight triple points
$D_\epsilon\cap Z_i\cap U_{i\epsilon}$, one for
$1\le i\le4$ and $\epsilon=0,\infty$, and no higher
intersection. On $A$, the curves $D_0,D_\infty$ are two
fibres of one projection and the $Z_i$ are sections of that
projection, all with the same divisor class. On $E_i$, $Z_i$
is the negative section and $U_{i0},U_{i\infty}$ are two
ruling fibres. These assertions retain the actual intersections:
the local quotient equation is $AC=B^2$, with base function
$T=u^2A$, and its $C$-chart after resolution has
$A=C\delta^2$, $B=C\delta$, hence $T=u^2C\delta^2$.
The three reduced components are $u=0$, $C=0$, and
$\delta=0$. Thus this support has simple normal crossings.

\item The elliptic map $M\to\PP^1$, induced by $z\mapsto z^2$,
has at $0,\infty$ the actual fibre divisors
\[
 \begin{gathered}
 F_0=2D_0+\sum_{i=1}^4U_{i0},\qquad
 F_\infty=2D_\infty+\sum_{i=1}^4U_{i\infty},\\
 [F_0]=[F_\infty]=2[J_a]=2[J_b]
       \quad\text{in }H_2(M,\Q).
 \end{gathered}                                          \tag{C.EN.S2}\hypertarget{tag-00H8}{}\label{cdp:eq:C.EN.S2}
\]
Indeed near a fixed base value the elliptic involution has
four fixed points; the displayed resolution gives one central
rational curve of multiplicity two and four exceptional arms
of multiplicity one. At the two root values the chart
$x=a^2b$ gives the two smooth double fibres $2J_a,2J_b$.
All these are fibres of the specified proper elliptic map,
which proves the equality of their rational fundamental classes.

\item The ten rational curves in \eqref{cdp:eq:C.EN.S2} have precisely one rational
homology relation, namely $F_0-F_\infty=0$. To check this,
their intersection matrix is the direct sum of two affine
$D_4$ matrices: each central curve and its four arms has
self-intersection $-2$, each arm meets the central curve
once, and all remaining intersections are zero. Each matrix
has rank four and kernel generated by $(2,1,1,1,1)$.
Any homology relation lies in this two-dimensional kernel.
The common fibre class is nonzero, since its intersection
with an ample divisor of the projective surface $M$ is
positive. Thus the only kernel combination which is zero in
homology is the difference of the two fibres. In particular
these ten curves span a nine-dimensional subspace of
$H_2(M,\Q)$.

\item Here is the complete constant-sheaf comparison. If $S^{[p]}$
denotes the disjoint union of the $(p+1)$-fold intersections
listed in \eqref{cdp:eq:C.EN.S1}, the alternating restriction maps give
an exact complex
\[
 0\longrightarrow\Q_S\longrightarrow\Q_{S^{[0]}}
 \longrightarrow\Q_{S^{[1]}}
 \longrightarrow\Q_{S^{[2]}}\longrightarrow0.           \tag{C.EN.S3}\hypertarget{tag-00H9}{}\label{cdp:eq:C.EN.S3}
\]
At a point this is the augmented cochain complex of the
simplex formed by the branches through that point, proving
its exactness. Its cohomology spectral sequence has
$E_1^{p,q}=H^q(S^{[p]},\Q)$.
In degree zero the dimensions are $8,16,8$. The first
map has rank seven, because the incidence graph is connected.
The second has rank eight: the boundary of each triple
point contains its own edge $U_{i\epsilon}$, which occurs
in no other triangle. Consequently its degree-zero cohomology
has dimensions $(1,1,0)$ in columns $0,1,2$.

\item Only the two Hopf components contribute to $H^1(S^{[0]},\Q)$,
of total dimension two; only $J_a,J_b$ contribute to
$H^1(S^{[1]},\Q)$, of total dimension four.
Each Hopf restriction has rank one. In the actual presentation
$(\C^2\setminus0)/q^{2\Z}$, inclusion of its elliptic
attaching curve preserves the radial deck generator and
kills its angular generator. The two images lie in separate
elliptic summands. This map therefore has rank two.

\item The degree-two restriction map has dimensions $20\to16$
and rank fourteen. We verify its rank through the dual map
on fundamental curve classes. Give coefficients
$x_0,x_\infty$ to $D_0,D_\infty$, coefficients $y_a,y_b$
to $J_a,J_b$, coefficients $z_i$ to $Z_i$, and coefficients
$u_{i0},u_{i\infty}$ to $U_{i0},U_{i\infty}$, orienting
edges from the first component to the second in \eqref{cdp:eq:C.EN.S1}.
The $E_i$ equations force
$z_i=0$ and $u_{i\infty}=-u_{i0}$; the $A$ equation
forces $x_\infty=-x_0$. Equation \eqref{cdp:eq:C.EN.S2} and the unique
rational-curve relation then force
\[
 y_a+y_b=0,\qquad u_{i0}=x_0/2\quad(1\le i\le4).
                                                               \tag{C.EN.S4}\hypertarget{tag-00HA}{}\label{cdp:eq:C.EN.S4}
\]
For completeness, the common half-fibre class is independent
of the four-dimensional span of the differences
$D_0-D_\infty,U_{i0}-U_{i\infty}$: in the curve vector
space modulo the sole relation $F_0-F_\infty$, a multiple
of $F_0$ cannot be a difference vector unless that multiple
is zero. Thus \eqref{cdp:eq:C.EN.S4} lists all kernel elements. Its two
parameters are the relation $J_a-J_b$ and the relation
$2D_0-2D_\infty+\sum_i(U_{i0}-U_{i\infty})$.
This proves rank fourteen.

\item Finally the components contribute dimension two in degree
three, from the Hopf components, and dimension eight in
degree four. Intersections have no cohomology in those
degrees. The sole possible higher differential with nonzero
target would land in $E_2^{2,0}$, which was proved zero;
all groups $E_2^{2,q}$ for $q>0$ vanish because triple
intersections are points. The preceding actual maps give
\[
 \bigl(b_i(S)\bigr)_{i=0}^4=(1,1,8,4,8),\qquad e(S)=12.
                                                               \tag{C.EN.S5}\hypertarget{tag-00HB}{}\label{cdp:eq:C.EN.S5}
\]

\item The identical alternating complex with structure sheaves
is exact in the same normal crossing charts. All components
have $H^2(\OO)=0$; the only nonzero component $H^1(\OO)$
groups are those of the two Hopf surfaces, one-dimensional
each. Their restriction to $H^1(J_j,\OO_{J_j})$ is an
isomorphism. One can see this through the actual Hopf
elliptic fibration over $\PP^1$: its transition maps on the
fixed elliptic curve are translations, so its $R^1\OO$
is the trivial line; the Leray edge map is an isomorphism
and evaluation at any fibre is an isomorphism. The degree-zero
incidence calculation is unchanged. Hence
\[
 \bigl(h^i(S,\OO_S)\bigr)_{i=0}^2=(1,1,0).
                                                               \tag{C.EN.S6}\hypertarget{tag-00HC}{}\label{cdp:eq:C.EN.S6}
\]

\item The six negative balancing blocks may be placed away from
$M$, keeping the Enriques surface itself untouched. For
each point or curve blowup, the support is the full inverse
image. The relative-pair comparison and the restrictions
of the ambient exceptional Gysin classes give exactly
(C.KM.S2); each negative block adds $(2,6,2)$ in degrees
two, three, and four. For coherent cohomology the chosen
centre has generic reduced-support multiplicity one, and
the displayed blowup charts give $\sigma^*S=S'$ as
Cartier divisors, including the four crossing points of
the quartic step. Applying the projection formula to the
Cartier sequence, with $R\sigma_*\OO=\OO$, gives
$R\sigma_*\OO_{S'}=\OO_S$ as in (C.KM.S3).
Thus the balanced support $S_{\rm bal}$ has
\[
 \bigl(b_i(S_{\rm bal})\bigr)_{i=0}^4=(1,1,20,40,20),
 \qquad
 \bigl(h^i(S_{\rm bal},\OO_{S_{\rm bal}})\bigr)_{i=0}^2
       =(1,1,0),\qquad e(S_{\rm bal})=0.                 \tag{C.EN.S7}\hypertarget{tag-00HD}{}\label{cdp:eq:C.EN.S7}
\]
The complete incidence matrices, curve-class map, and
their exact ranks are executed in
\path{code/A_enriques_support.py} and retained in its
\texttt{.out} file.

\end{enumerate}
\end{proof}






\begin{proposition}[A fixed punctured torus family has infinitely many zero-Euler fillings with essential projective elliptic components of coherent Euler characteristic one]\hypertarget{tag-01D5}{}\label{cdp:st:C.2.UN.27}
\StatementRecord{cdp:st:C.2.UN.27}{proved}{01D5}{\SeriesRef{cdp:st:C.2.UN.25}, \SeriesRef{cdp:eq:C.EN.R7}, \SeriesRef{cdp:eq:C.EN.R10}, \SeriesRef{cdp:eq:C.EN.R11}, \SeriesRef{cdp:eq:C.EN.R8}, \SeriesRef{cdp:st:C.2.UN.4}, \SeriesRef{cdp:eq:C.EN.G1}, \SeriesRef{cdp:st:C.2.UN.14}}

For every integer $r\ge2$, the compact construction
of \ref{cdp:st:C.2.UN.25} has a version whose distinguished smooth
component $M_r$ is a minimal projective elliptic
surface with exactly two fibres of type $I_0^*$
and $r$ double smooth elliptic fibres $2J_1,\ldots,2J_r$.
Writing $F$ for a full fibre, its actual canonical
line and plurigenera are
\[
 \begin{gathered}
 K_{M_r}\simeq\OO_{M_r}\left(\sum_{j=1}^rJ_j-F\right),
 \qquad
 2K_{M_r}\simeq\OO_{M_r}((r-2)F),
 \\
 P_m(M_r)=\max\{0,\,1-m+r\lfloor m/2\rfloor\}.
 \end{gathered}
                                                               \tag{C.EN.R1}\hypertarget{tag-00GV}{}\label{cdp:eq:C.EN.R1}
\]
It has
\[
 (b_0,b_1,b_2,b_3,b_4)=(1,0,10,0,1),\qquad
 (h^0(\OO),h^1(\OO),h^2(\OO))=(1,0,0).
                                                               \tag{C.EN.R2}\hypertarget{tag-00GZ}{}\label{cdp:eq:C.EN.R2}
\]
The case $r=2$ is the Enriques component of \ref{cdp:st:C.2.UN.25}.
Every $r\ge3$ gives Kodaira dimension one and
coherent Euler characteristic one.

Before balancing, the full central fibre is
\[
 2M_r+2A+4\sum_{j=1}^r Q_j+\sum_{i=1}^4E_i,
                                                               \tag{C.EN.R3}\hypertarget{tag-00H0}{}\label{cdp:eq:C.EN.R3}
\]
with $A\simeq\PP^1\times\PP^1$, primary Hopf
$Q_j$, and $E_i\simeq\mathbb F_2$.
Its reduced support has Euler number $12$.
Six of the explicit negative-two operations,
all disjoint from $M_r$, make that number zero.
The reduced supports before and after these
operations have respectively
\[
 \begin{array}{c|c|c}
 &(b_0,b_1,b_2,b_3,b_4)&(h^0(\OO),h^1(\OO),h^2(\OO))\\ \hline
 S_r&(1,1,r+6,2r,r+6)&(1,1,0)\\
 S_{r,\mathrm{bal}}&(1,1,r+18,2r+36,r+18)&(1,1,0).
 \end{array}                                             \tag{C.EN.R4}\hypertarget{tag-00H1}{}\label{cdp:eq:C.EN.R4}
\]
All these fillings have the same specified punctured
torus family. Their distinguished essential components
are pairwise non-bimeromorphic, with canonical growth
slope $r/2-1$.
\end{proposition}


\begin{proof}

\begin{enumerate}[label={$\langle1\rangle$\arabic*.},leftmargin=*]
\item Use exactly the compact Weierstrass base, quadratic
base change, resolution, and involution of
(C.EN.2)--(C.EN.5). On its exceptional $R_1$,
blow up the $2r$ points
$z=1,\ldots,r,-1,\ldots,-r$.
Replace (C.EN.6) by
\[
 D_r=\sum_{j=1}^r(D_j-D_{-j}),\qquad L_r=\OO(D_r).
                                                               \tag{C.EN.R5}\hypertarget{tag-00H2}{}\label{cdp:eq:C.EN.R5}
\]
The canonical section and the involution give
$h^*L_r=L_r^{-1}$ with its actual order-two
section-preserving dualization. The two signed
root charts still have order two, and their
free $q$-quotient and properness are exactly
(C.EN.7)--(C.EN.8). The formula
$\ell=\sigma c^2,\ a=c^{-1}$ supplies an
actual product identification away from the
central divisor, and the involution there
is always $(v,[c])\mapsto(hv,[c^{-1}])$.
It is independent of $r$. This proves the
assertion about the fixed punctured family.

\item On the strict transform of $R_1$, use
\[
 \phi_r(z)=\prod_{j=1}^r\frac{z-j}{z+j},\qquad
 C_r:\quad Y^2=\prod_{j=1}^r(z^2-j^2),\qquad
 y=\frac{Y}{\prod_{j=1}^r(z+j)}.
                                                               \tag{C.EN.R6}\hypertarget{tag-00H3}{}\label{cdp:eq:C.EN.R6}
\]
The $2r$ branch points are distinct, so $C_r$
is smooth and has genus $r-1$, by the
degree-two Riemann--Hurwitz formula.
The actual finite \'etale product cover
of the corresponding component is again
$\ell=u^2,\ a=y/u$ from $C_r\times E$,
with $E=\C^*/q^{\Z}$. Its deck involution
and the lifted dualization are
\[
 \begin{split}
 \eta(z,Y,[u])&=(z,-Y,[-u]),\\
 k(z,Y,[u])&=(-z,(-1)^rY,[u^{-1}]).
 \end{split}                                             \tag{C.EN.R7}\hypertarget{tag-00H4}{}\label{cdp:eq:C.EN.R7}
\]
They commute and generate a Klein four group.
The first action is free. The second formula
is forced by the actual identity
$\phi_r(-z)=\phi_r(z)^{-1}$, so it includes
the sign for odd $r$.

\item We check both parities at the fixed base values.
At infinity, $Y/z^r=\pm1$, and $k_C$
fixes the two points while $jk_C$ exchanges
them. At zero, the two values of $Y$ are
nonzero. When $r$ is even, $k_C$ fixes
them and $jk_C$ exchanges them; when $r$
is odd, the roles are reversed. Therefore
\[
 \begin{array}{c|cc}
 &\#\operatorname{Fix}(k_C)&
       \#\operatorname{Fix}(jk_C)\\ \hline
 r\text{ even}&4&0\\
 r\text{ odd}&2&2 .
 \end{array}                                             \tag{C.EN.R8}\hypertarget{tag-00H5}{}\label{cdp:eq:C.EN.R8}
\]
Each nontrivial fixed-point action on the
smooth curve has derivative $-1$.
On $E$, both $[u]\mapsto[u^{-1}]$
and $[u]\mapsto[-u^{-1}]$ have four
fixed points and derivative $-1$.
Consequently the two elements $k,\eta k$
have a total of sixteen fixed points,
all with local representation $(-1,-1)$.
Their sets are disjoint because $\eta$
is free. Dividing by the remaining free
involution pairs them into eight quotient
points. Thus
\[
 X_r=(C_r\times E)/\langle\eta,k\rangle
                                                               \tag{C.EN.R9}\hypertarget{tag-00H6}{}\label{cdp:eq:C.EN.R9}
\]
has precisely eight $A_1$ singularities,
and no branch curve. Let $M_r$ be its
minimal $A_1$ resolution. The same
ambient charts (C.EN.14)--(C.EN.15)
realize that exact resolution as a
central component. For odd $r$ the
fixed lifts over zero are the
$\eta k$ lifts; downstairs the local
involution, the four fixed curves,
and all monomials of those charts
are unchanged.

\item The map $w=z^2$ gives the elliptic map
$p:M_r\to\PP^1_w$. Above a general $w$,
the four points of $C_r$ are acted on
transitively by $\langle\eta,k\rangle$,
so its fibre is exactly $E$. At $w=j^2$,
the local double-root chart and the free
translation in $\eta$ give a smooth
double fibre $2J_j$. At zero and infinity,
the involution has four fixed points on
the elliptic fibre; its $A_1$ resolutions
give the full divisor
$2D+\sum_{i=1}^4 U_i$ of type $I_0^*$,
as in (C.EN.S2). These are all special
values. In particular the morphism $p$
has connected fibres.

\item We compute its canonical line using an
actual invariant meromorphic form,
instead of deducing it from numerical
invariants. On $C_r\times E$ use
\[
                       \Omega=dz\wedge\frac{du}{u}.
                                                               \tag{C.EN.R10}\hypertarget{tag-00GW}{}\label{cdp:eq:C.EN.R10}
\]
It is fixed by both generators in
\eqref{cdp:eq:C.EN.R7}. On $C_r$, $dz$ has a simple
zero at each of the $2r$ branch points
and a double pole at each of the two
points over infinity, and no other
zero or pole. The quotient is free
in codimension one. The root charts
therefore give one zero along every
$J_j$. At infinity, the local invariant
coordinate is $w^{-1}$, whose divisor
pulls back with order two at both
points of $C_r$.

\item The resolution comparison at the isolated
quotient points retains these orders.
In local coordinates $(s,v)$ where the
involution is $(-s,-v)$, its crepant
volume form satisfies
$ds\wedge dv=\tfrac12\,dr\wedge d\delta$
in the chart $s^2=r,\ sv=r\delta$.
Thus the double pole $s^{-2}$ of
\eqref{cdp:eq:C.EN.R10} becomes a simple pole
on the exceptional curve, and in the
other chart has orders one and two
on the exceptional and strict central
curves. These are exactly the
multiplicities of the full $I_0^*$
fibre. At zero the form is regular
and nowhere zero, and its crepant
extension is likewise nowhere zero.
Consequently the descended form has
the exact divisor
\[
                \operatorname{div}(\Omega)
                         =\sum_{j=1}^r J_j-F_\infty.
                                                               \tag{C.EN.R11}\hypertarget{tag-00GX}{}\label{cdp:eq:C.EN.R11}
\]
This proves the first formula of
\eqref{cdp:eq:C.EN.R1}. Since $2J_j$ is linearly
equivalent to a full fibre, it also
proves its second formula.

\item Every meromorphic section of $mK_{M_r}$
is a meromorphic function times
$\Omega^{\otimes m}$. A holomorphic
such section restricts to a constant
multiple on a general elliptic fibre,
because that fibre's canonical line
is trivial and $\Omega$ has no zero
or pole there. Its ratio is therefore
a function from the connected-fibre
base. The exact allowed orders at
the double fibres in \eqref{cdp:eq:C.EN.R11}
give the sheaf equality
\[
 p_*\OO_{M_r}(mK_{M_r})
  \simeq
 \OO_{\PP^1}\left(
      \sum_{j=1}^r\lfloor m/2\rfloor[j^2]
                       -m[\infty]\right).
                                                               \tag{C.EN.R12}\hypertarget{tag-00GY}{}\label{cdp:eq:C.EN.R12}
\]
The ordinary full fibre at infinity,
including all its components, is
Cartier and contributes exactly the
displayed $-m$. Taking sections gives
the plurigenus formula in \eqref{cdp:eq:C.EN.R1}.
For $r=2$, it gives $P_m=1$ for even
$m$ and zero for odd $m$, recovering
\ref{cdp:st:C.2.UN.25}. For $r\ge3$, its linear growth
has positive slope $r/2-1$, so
$\kappa(M_r)=1$. The line $2K_{M_r}$
is a nonnegative multiple of a fibre
and is nef. Adjunction therefore
excludes a $(-1)$-curve, proving that
$M_r$ is minimal for all $r\ge2$.

\item The component is projective because
it is the resolution of the finite
quotient of the projective surface
$C_r\times E$. Its coherent groups
can also be read through that exact
quotient. On $H^1(C_r,\OO)$ the
hyperelliptic involution $j$ is
$-I$; this follows directly from
the basis $z^a dz/Y$, $0\le a\le r-2$,
of canonical forms and duality.
The action $\eta$ therefore removes
this summand, whereas $k$ removes
the remaining elliptic summand.
Thus $q(M_r)=0$. The $A_1$ resolution
does not change coherent cohomology:
the explicit exceptional $(-2)$-curve
has positive successive conormal
powers, whose $H^1$ vanish, so
formal functions gives $R^1\pi_*\OO=0$.
The formula at $m=1$ gives $p_g=0$.
The finite quotient Euler formula is
$(0+0+16)/4=4$, for both parity
patterns in \eqref{cdp:eq:C.EN.R8}. Resolving
eight $A_1$ points adds eight,
giving $e(M_r)=12$. Hodge symmetry
and Poincar\'e duality on this
projective surface now yield
\eqref{cdp:eq:C.EN.R2}.

\item The ambient fixed curves over zero
are still the four copies of $R_0$,
and their normal bundles are still
$\OO(-2)\oplus\OO$. The $2r$
root Hopf components are paired
by the involution into $r$ components.
The exact quotient and resolution
monomials give \eqref{cdp:eq:C.EN.R3}, with no
parity change in its coefficients.
The incidence data of \ref{cdp:st:C.2.UN.25} are
unchanged except for the $r$
elliptic attaching curves $M_r\cap Q_j$.
Those curves all have Euler number
zero. Hence the same calculation
gives $e(S_r)=12$, and six \ref{cdp:st:C.2.UN.4}
operations away from $M_r$ give
zero while preserving $M_r$ itself.
The normal crossing and compactness
arguments are the actual charts
of \ref{cdp:st:C.2.UN.25}; no new analytic gluing
is required.

\par\medskip\noindent\textbf{Cohomology of all the paired-root supports.}\par\noindent
The support calculation (C.EN.S1)--(C.EN.S7) extends through
the same maps, with $r$ Hopf components instead of two.
We give the dimensions and the changed maps explicitly.
There are $r+6$ components: the projective elliptic surface
$M$, $A\simeq\PP^1\times\PP^1$, four surfaces
$E_i\simeq\mathbb F_2$, and $Q_1,\ldots,Q_r$.
The double curves are the fourteen rational curves
$D_0,D_\infty,Z_i,U_{i0},U_{i\infty}$ already listed
in (C.EN.S1), and the $r$ elliptic curves
$J_j=M\cap Q_j$. The same eight triple points remain.
On $M$, the actual elliptic fibration gives
\[
 [2D_0+\textstyle\sum_iU_{i0}]
 =[2D_\infty+\textstyle\sum_iU_{i\infty}]
 =2[J_j]\quad(1\le j\le r).
                                                               \tag{C.EN.G1}\hypertarget{tag-00GR}{}\label{cdp:eq:C.EN.G1}
\]
Its two $I_0^*$ configurations still span dimension nine
in $H_2(M,\Q)$, by the same affine-$D_4$ intersection
calculation and nonzero fibre class. Here $b_2(M)=10$,
$b_1(M)=b_3(M)=0$, and $h^1(M,\OO_M)=h^2(M,\OO_M)=0$
by the invariant and resolution computations above.

\item For the exact constant-sheaf complex (C.EN.S3), the
degree-zero row now has dimensions
$r+6,r+14,8$. Its two ranks are $r+5$ and eight:
each extra Hopf vertex and attaching edge is a leaf,
and each of the eight triangles retains its private
edge $U_{i\epsilon}$. Thus its cohomology is still
$(1,1,0)$. In degree one the map has dimensions
$r\to2r$ and rank $r$, since each Hopf radial
generator restricts nontrivially to its own elliptic
attaching curve. In particular this row contributes
dimension $r$ to $H^2$ of the support.

\item The degree-two restriction map has dimensions
$20\to r+14$ and rank fourteen. The $E_i$ and $A$
equations in (C.EN.S4) are unchanged. The $M$
equation, using \eqref{cdp:eq:C.EN.G1}, becomes
\[
 \sum_{j=1}^r y_j=0,\qquad
 u_{i0}=x_0/2,\quad u_{i\infty}=-x_0/2,\quad
 x_\infty=-x_0,\quad z_i=0.
                                                               \tag{C.EN.G2}\hypertarget{tag-00GS}{}\label{cdp:eq:C.EN.G2}
\]
These are all solutions by the already proved
independence of the common half-fibre and the
rational-curve difference subspace. The kernel
therefore has dimension $(r-1)+1=r$, proving
rank fourteen. The component degree-three and
degree-four groups have dimensions $r$ and $r+6$,
respectively. As before, every possible higher
differential has zero target. Consequently
\[
 \bigl(b_i(S_r)\bigr)_{i=0}^4=(1,1,r+6,2r,r+6),
 \qquad e(S_r)=12.
                                                               \tag{C.EN.G3}\hypertarget{tag-00GT}{}\label{cdp:eq:C.EN.G3}
\]
For structure sheaves the $r$ Hopf-to-elliptic
restrictions are isomorphisms, and every component
has zero $H^2(\OO)$. The same exact complex gives
$(h^i(S_r,\OO_{S_r}))_{i=0}^2=(1,1,0)$.
Applying six negative blocks away from $M$, with
the support pullback/Gysin maps and derived Cartier
comparison used in (C.EN.S7), gives
\[
 \begin{gathered}
 \bigl(b_i(S_{r,\mathrm{bal}})\bigr)_{i=0}^4
      =(1,1,r+18,2r+36,r+18),\qquad e(S_{r,\mathrm{bal}})=0,\\
 \bigl(h^i(S_{r,\mathrm{bal}},\OO_{S_{r,\mathrm{bal}}})\bigr)_{i=0}^2
      =(1,1,0).
 \end{gathered}                                                \tag{C.EN.G4}\hypertarget{tag-00GU}{}\label{cdp:eq:C.EN.G4}
\]
The parameterized incidence and curve-class matrices
are included in \path{code/A_enriques_support.py}; its
executed output checks both the original case and
the first two properly elliptic cases.

\item Finally a covering family of rational
curves would kill every positive
pluricanonical section by its actual
dominant ruled parametrization.
Thus $M_r$ is non-uniruled, and
\ref{cdp:st:C.2.UN.14} gives its essentiality under
proper birational comparison of
smooth fillings. The canonical
growth slopes in \eqref{cdp:eq:C.EN.R1} are
distinct, so these components
are pairwise non-bimeromorphic.
For $r\ge3$ each is the unique
component of Kodaira dimension
one; the other original components
are rational or Hopf, and the
additional blowup components are
ruled. The same inverse-map
comparison therefore also
distinguishes the corresponding
fillings over the fixed punctured
family.

\end{enumerate}
\end{proof}






\begin{proposition}[The fixed Kummer punctured family has zero-Euler fillings with arbitrarily large canonical growth on a projective elliptic component]\hypertarget{tag-01D9}{}\label{cdp:st:C.2.UN.30}
\StatementRecord{cdp:st:C.2.UN.30}{proved}{01D9}{\SeriesRef{cdp:st:C.2.UN.18}, \SeriesRef{cdp:st:C.2.UN.23}, \SeriesRef{cdp:st:C.2.UN.13}, \SeriesRef{cdp:st:C.2.UN.25}, \SeriesRef{cdp:eq:C.ET.11}, \SeriesRef{cdp:st:C.2.UN.14}}

Fix the square elliptic curve
$E_0=\C/(\Z+i\Z)$ and $E=\C^*/q^{\Z}$,
where $q=e^{-2\pi}$. For every $r\ge1$ there is
a smooth compact threefold over $\PP^1$,
with the same punctured product-torus family
as the diagonal involution quotient
\[
 (E_0\times E\times\PP^1_t)/
      \langle(p,[c],t)\mapsto(-p,[c^{-1}],-t)\rangle,
                                                               \tag{C.ET.1}\hypertarget{tag-00J5}{}\label{cdp:eq:C.ET.1}
\]
whose two special reduced supports have Euler
number zero. The support over zero contains an
unchanged smooth minimal projective properly
elliptic surface $M_r$ with four $I_0^*$ fibres
and $r$ double smooth elliptic fibres $2J_j$.
For its connected elliptic map $p_r:M_r\to\PP^1$
and a full fibre $F$,
\[
 \begin{aligned}
 K_{M_r}&\simeq\OO_{M_r}\left(\sum_{j=1}^rJ_j\right),
 \qquad 2K_{M_r}\simeq\OO_{M_r}(rF),\\
 (p_r)_*\OO(mK_{M_r})
   &\simeq\OO_{\PP^1}\left(
                \sum_{j=1}^r\lfloor m/2\rfloor[b_j]\right),
 \\
 P_m(M_r)&=1+r\lfloor m/2\rfloor .
 \end{aligned}                                             \tag{C.ET.2}\hypertarget{tag-00J9}{}\label{cdp:eq:C.ET.2}
\]
Here $b_j$ is the value of $2J_j$. In particular
$\chi(\OO_{M_r})=2$, $p_g(M_r)=1$, $q(M_r)=0$,
and
\[
            (b_0,b_1,b_2,b_3,b_4)(M_r)
                         =(1,0,22,0,1).              \tag{C.ET.3}\hypertarget{tag-00JA}{}\label{cdp:eq:C.ET.3}
\]
These essential components, and the associated
fillings over the fixed punctured family, are
pairwise non-bimeromorphic.

Before balancing, the full fibre over zero is
\[
                  2M_r+4\sum_{j=1}^rQ_j
                           +\sum_{\nu=1}^{16}P_\nu,
                                                               \tag{C.ET.4}\hypertarget{tag-00JB}{}\label{cdp:eq:C.ET.4}
\]
where the $Q_j$ are primary Hopf surfaces and
$P_\nu\simeq\PP^2$. It has Euler number $40$.
Its reduced support and its balanced version
have respectively
\[
 \begin{array}{c|c|c}
 &(b_0,b_1,b_2,b_3,b_4)&
                 (h^0(\OO),h^1(\OO),h^2(\OO))\\ \hline
 S_r&(1,0,21+r,2r-1,17+r)&(1,0,1)\\
 S_{r,\mathrm{bal}}&(1,0,61+r,2r+119,57+r)&(1,0,1).
 \end{array}                                             \tag{C.ET.5}\hypertarget{tag-00JC}{}\label{cdp:eq:C.ET.5}
\]
\end{proposition}


\begin{proof}

\begin{enumerate}[label={$\langle1\rangle$\arabic*.},leftmargin=*]
\item Choose an odd prime $N>2r$ and the torsion
points $p_j=j/N$ of $E_0$, for $1\le j\le r$.
Their pairs $\{p_j,-p_j\}$ are disjoint and
avoid the two-torsion points. Blow up
$E_0\times\PP^1_t$ at $(p_j,0),(-p_j,0)$,
and call the exceptional curves $R_j,R_{-j}$.
The involution $(p,t)\mapsto(-p,-t)$ lifts
to the resulting surface $V$, exchanging
the paired curves. Put
\[
       D=\sum_{j=1}^r(R_j-R_{-j}),\qquad L=\OO_V(D).
                                                               \tag{C.ET.6}\hypertarget{tag-00JD}{}\label{cdp:eq:C.ET.6}
\]
The canonical meromorphic section $\sigma$
gives $h^*L=L^{-1}$ with its order-two
section-preserving dualization.

\item Perform the signed order-two root construction
of \ref{cdp:st:C.2.UN.18} along these disjoint curves:
$\sigma=a^2\ell$ on positive charts and
$\sigma^{-1}=(a')^2\ell'$ on negative charts,
with $a'=a^{-1},\ell'=\ell^{-1}$.
Divide by $(\ell,a)\mapsto(q^2\ell,q^{-1}a)$.
The explicit finite free charts (C.PE.8)
prove that the resulting compact threefold
$\widetilde Y$ is smooth and proper over $V$.
Dualization over $h$ takes the global
meromorphic pair $(\ell,a)$ to
$(\ell^{-1},a^{-1})$ and conjugates the
$q$-action to its inverse. It therefore
defines an actual involution of
$\widetilde Y$.

\item Away from the root curves the formula
$\ell=\sigma c^2,\ a=c^{-1}$ identifies the
family with $V\times E$. Under that actual
identification the involution is
$(v,[c])\mapsto(hv,[c^{-1}])$.
All modifications lie over $t=0$.
Consequently the quotient family over
$T=t^2$, on the complement of zero and
infinity, is exactly the punctured family
in \eqref{cdp:eq:C.ET.1}, independently of $r$.
The quotient and every resolution are proper over
$\PP^1$; the Stein factor has generic degree one
and is therefore the identity finite map to the
normal base. All fibres are connected.

\item Let $B_r$ be the component above the
strict central elliptic curve $E_0$.
It is an elliptic surface over $E_0$
with $2r$ smooth double fibres, at
$p_j,-p_j$. It is projective; we
give the actual finite product cover
which proves this. The restricted line
\[
 L|_{E_0}\simeq
       \OO_{E_0}\left(\sum_{j=1}^r[p_j]-[-p_j]\right)
                                                               \tag{C.ET.7}\hypertarget{tag-00JE}{}\label{cdp:eq:C.ET.7}
\]
is torsion because all $p_j$ are torsion.
Choose a torsion square root $\mathcal N$
of this line, and form the normalization
$C$ of $y^2=\sigma|_{E_0}$ in the
corresponding line-bundle compactification.
It is a connected double cover of $E_0$
branched simply at the $2r$ specified
points, so $g(C)=r+1$.
Choose a finite \'etale cover $C'\to C$
trivializing the pulled-back torsion
line $\mathcal N$, and a nowhere-zero
frame $n$ there. Write $y=y_0n$.
Then
\[
 C'\times E\longrightarrow B_r,\qquad
 (z,[u])\longmapsto
             [\,\ell=u^2n^2,\ a=y_0/u\,]             \tag{C.ET.8}\hypertarget{tag-00JF}{}\label{cdp:eq:C.ET.8}
\]
is the exact finite \'etale map of \ref{cdp:st:C.2.UN.18}.
At a positive root its coordinate Jacobian
is $2$; in the negative chart it is
$-2u^{-2}$. Thus it is \'etale also
over the multiple fibres. The finite
product-cover comparison proves
projectivity of $B_r$, and hence of
the finite quotient and its surface
resolution used below.

\item The involution on $B_r$ covers inversion
on $E_0$. At each of the four two-torsion
base points, the root curves are absent,
and the actual product formula gives
inversion on the elliptic fibre.
There are exactly sixteen fixed points
on $B_r$, each with derivative $(-1,-1)$.
There is no fixed curve. Its quotient
has sixteen $A_1$ singularities; let
$M_r$ be their minimal resolution.
The induced base is
$E_0/\{\pm1\}\simeq\PP^1$.
The four fixed base values give $I_0^*$
fibres by the same explicit resolution
as (C.EN.S2). The paired root values
give the $r$ smooth double fibres
$2J_j$. Away from those values the
quotient is a smooth elliptic fibration,
so this is its complete fibre list.

\item Let $\omega_0$ be a nonzero holomorphic
one-form on $E_0$. On $B_r$ the
meromorphic expression
\[
                        \Omega=\omega_0\wedge d\log\ell
                                                               \tag{C.ET.9}\hypertarget{tag-00JG}{}\label{cdp:eq:C.ET.9}
\]
is global. A change of frame of $L$
adds a base logarithmic differential
to $d\log\ell$; its wedge with
$\omega_0$ is zero. Multiplication
by the constant $q^2$ also leaves
the expression unchanged. In a
positive root chart $z=a^2b$ it is
a unit times $2a\,da\wedge db$;
the negative dual chart has the
opposite sign and the same simple
zero. Thus its divisor on $B_r$
is precisely the sum of the
$2r$ reduced multiple fibres.
The involution changes both
$\omega_0$ and $d\log\ell$ by a
minus sign, up to the harmless
base differential already described.
Hence $\Omega$ is invariant.

\item At its sixteen fixed points the form
is regular and nowhere zero. The
height-one $A_1$ charts of \ref{cdp:st:C.2.UN.23}
extend it without a zero or pole.
Paired root fibres are exchanged
freely and become $J_j$. The
descended form consequently has
the exact divisor
\[
                      \operatorname{div}(\Omega)
                                      =\sum_{j=1}^rJ_j .
                                                               \tag{C.ET.10}\hypertarget{tag-00J6}{}\label{cdp:eq:C.ET.10}
\]
This proves the first two line-bundle
identities in \eqref{cdp:eq:C.ET.2}.
The ratio of any holomorphic
$m$-canonical section to
$\Omega^{\otimes m}$ is constant
on a connected general elliptic
fibre. It descends meromorphically
to the base. Its allowed pole
order at $b_j$ is exactly
$\lfloor m/2\rfloor$, because
pullback of a base parameter
has divisor $2J_j$; there are
no other allowed poles.
This proves the actual direct-image
equality in \eqref{cdp:eq:C.ET.2}, and hence
all its plurigenera.
Their growth proves $\kappa(M_r)=1$.
The line $2K_{M_r}=rF$ is nef,
so adjunction excludes a rational
$(-1)$-curve; $M_r$ is minimal.

\item The surface $B_r$ has Euler number
zero: every fibre over its elliptic
base has elliptic reduced support,
and the map is locally a smooth
elliptic bundle off finitely many
values. The finite involution
quotient therefore has Euler
number $(0+16)/2=8$. Resolving
sixteen $A_1$ points gives
$e(M_r)=24$. Its canonical
divisor is a sum of disjoint
square-zero fibres, so
$K_{M_r}^2=0$. Surface
Riemann--Roch gives
$\chi(\OO_{M_r})=24/12=2$.
The first plurigenus in \eqref{cdp:eq:C.ET.2}
is one; hence $q(M_r)=0$.
Projectivity, Hodge symmetry and
Poincar\'e duality give \eqref{cdp:eq:C.ET.3}.

\item We now specify the ambient resolution.
The only fixed base points of
$h$ lie at $E_0[2]\times\{0,\infty\}$,
and none was blown up. Each has
four fixed elliptic-fibre points.
The thirty-two ambient fixed
points therefore have local
linear action $(-1,-1,-1)$.
The actual invariant quadratic
blowup charts of \ref{cdp:st:C.2.UN.13} resolve
each quotient point smoothly.
For example, write upstairs
$(x,t,w)=(s,su,sv)$ and put
$R=s^2$. The resolved quotient
chart has coordinates $(R,u,v)$
and
\[
                              T=Ru^2 .                \tag{C.ET.11}\hypertarget{tag-00J7}{}\label{cdp:eq:C.ET.11}
\]
The exceptional surface is $\PP^2$.
Its intersection with the strict
central component is a line;
restriction is exactly the
minimal $A_1$ surface resolution.
Thus the strict component is
the same smooth $M_r$ proved
above, and the monomial gives
its multiplicity two and the
exceptional multiplicity one.
The paired root Hopf surfaces
give $r$ components of
multiplicity four, exactly as
in \ref{cdp:st:C.2.UN.25}. This proves \eqref{cdp:eq:C.ET.4}
with all full-fibre coefficients.

\item The reduced support is a star:
each $P_\nu$ meets $M_r$ along
one $\PP^1$, each $Q_j$ meets
$M_r$ along the elliptic curve
$J_j$, and there are no other
intersections or triple points.
It is simple normal crossing
by \eqref{cdp:eq:C.ET.11} and the root charts.
Therefore
\[
                 e(S_r)=24+16\cdot3-16\cdot2=40.
                                                               \tag{C.ET.12}\hypertarget{tag-00J8}{}\label{cdp:eq:C.ET.12}
\]
At infinity the same ambient
resolution gives the original
product-Kummer support
$2K+16\PP^2$ of \ref{cdp:st:C.2.UN.13}, also
of Euler number forty.
Twenty negative-two blocks
on each special support,
all beginning at fresh
smooth points of its plane
components away from the
K3 or $M_r$, make both
Euler numbers zero and
leave those distinguished
surfaces unchanged.

\item For completeness the support
cohomology in \eqref{cdp:eq:C.ET.5} follows
through exact restriction maps.
The constant-sheaf union
sequence has component term
$M_r\sqcup\bigsqcup Q_j
\sqcup\bigsqcup P_\nu$
and curve term
$\bigsqcup J_j\sqcup
\bigsqcup\PP^1_\nu$.
Its degree-zero incidence
graph is a tree. In degree
one the $r$ Hopf generators
map injectively, each onto
the radial one-dimensional
subspace of its elliptic
attaching curve. The cokernel
has dimension $r$.
The degree-two restriction
has source dimension $38$
and target dimension $16+r$.
Its dual curve-class map
has kernel of dimension
$r-1$: the plane components
force each of their sixteen
edge coefficients to vanish;
the remaining elliptic
classes all equal the
nonzero class $\tfrac12F$
on $M_r$ (a full fibre meets
an ample divisor positively), leaving exactly
the relations whose
coefficients sum to zero.
The map thus has rank
seventeen. Its kernel
and cokernel have dimensions
$21$ and $r-1$.
The components contribute
$r$ in degree three and
$17+r$ in degree four.
The resulting long exact
sequence gives the first
Betti row of \eqref{cdp:eq:C.ET.5}.

\item The analogous structure-sheaf
sequence has isomorphisms
$H^1(Q_j,\OO)\to H^1(J_j,\OO)$,
proved by the actual Hopf
elliptic Leray map in \ref{cdp:st:C.2.UN.25}.
The sole component
$H^2(\OO)$ is the one
on $M_r$. Hence the
coherent dimensions are
$(1,0,1)$.
Each of the twenty
balancing blocks adds
$(2,6,2)$ to support
cohomology in degrees
two, three and four,
by the restricted
exceptional Gysin maps
of (C.KM.S2), and
preserves coherent
cohomology by the
actual Cartier pullback
in (C.KM.S3). This
proves the second row.

\item The canonical growth slopes
$r/2$ are distinct and
positive. Therefore these
components are non-uniruled,
essential by \ref{cdp:st:C.2.UN.14}, and
pairwise non-bimeromorphic.
They are the unique components
of Kodaira dimension one
in their respective supports;
the other components are
rational, Hopf, K3, or
the ruled blowup components.
The inverse-map resolution
comparison consequently
also distinguishes the
fillings over their same
specified punctured family.

\end{enumerate}
\end{proof}






\begin{proposition}[Every fixed-modulus elliptic quasi-bundle over a rational or elliptic curve occurs unchanged]\hypertarget{tag-01DD}{}\label{cdp:st:C.2.UN.34}
\StatementRecord{cdp:st:C.2.UN.34}{proved}{01DD}{\SeriesRef{cdp:st:C.2.UN.22}, \SeriesRef{cdp:st:C.2.UN.29}, \SeriesRef{cdp:eq:C.SF.4}, \SeriesRef{cdp:eq:C.SF.3}, \SeriesRef{cdp:eq:C.SF.5}, \SeriesRef{cdp:eq:C.SF.7}, \SeriesRef{cdp:eq:C.SF.8}, \SeriesRef{cdp:eq:C.SF.1}, \SeriesRef{cdp:eq:C.SF.9}, \SeriesRef{cdp:eq:C.SF.10}, \SeriesRef{cdp:eq:C.SF.6}, \SeriesRef{cdp:eq:C.SF.12}, \SeriesRef{cdp:eq:C.SF.2}, \SeriesRef{cdp:eq:C.SF.11}, \SeriesRef{cdp:eq:C.SF.13}}

Let $q:S\to B$ be a proper holomorphic map from a
smooth compact connected complex surface to a smooth
compact curve with $g(B)\le1$. Assume that every
fibre is a positive multiple of a smooth elliptic
curve, that its smooth fibres have one fixed elliptic
modulus, and that its integral $H_1$ monodromy is
the identity. Then there are a smooth compact complex
threefold $Y$, a proper holomorphic map
$f:Y\to\mathbb P^1$, and a specified isomorphism
from $S$ to a component of its reduced zero fibre,
with the following properties:

\begin{itemize}
\item the component $S$ is unchanged and has
multiplicity one in the full fibre;
\item over $\mathbb P^1\setminus\{0\}$ the map
is an actual product with a fixed product of two
elliptic curves, and hence has connected torus
fibres and identity monodromy;
\item the reduced zero fibre has simple normal
crossings, all its components are smooth elliptic
quasi-bundles, all its double curves are smooth
elliptic, and it has no triple points;
\item every reduced fibre, including the zero
fibre, has Euler characteristic zero.
\end{itemize}
This realizes every elliptic covering model in
\ref{cdp:st:C.2.UN.22} and \ref{cdp:st:C.2.UN.29}. The prescribed surface
is retained through its full holomorphic torsor
class, including its degree-zero parameters and
its individual multiple-fibre translations.
\end{proposition}



\begin{proof}

\begin{enumerate}[label={$\langle1\rangle$\arabic*.},leftmargin=*]
\item Fix the elliptic curve $A=\mathbb C/\Lambda$ of
the smooth fibres, and choose an ordered lattice
basis $\lambda_1,\lambda_2$. Let the multiple
values be $p_i\in B$, with full fibre
$m_iC_i$ and $m_i>1$. Introduce the actual root
curve
\[
                 \mathcal B=B[(p_i,m_i)_i].
                                                               \tag{C.SF.1}\hypertarget{tag-010R}{}\label{cdp:eq:C.SF.1}
\]
Its local chart at $p_i$ is
$[\Delta_w/\mu_{m_i}]\to\Delta_z$,
$z=w^{m_i}$, with the indicated standard action
on $w$. We will use these charts to specify
all coarse-space maps in the construction.
The line and section $\mathcal O_S(C_i),s_{C_i}$,
whose $m_i$th power is the pullback of
$\mathcal O_B(p_i),s_{p_i}$, give the actual
lift of $q$ to this root curve.

\item First the given surface is the coarse space
of a principal $A$-torsor $\mathcal T$ over
$\mathcal B$. Over the smooth base, choose
local holomorphic sections and integrate a
marked differential. Constant modulus and
identity integral monodromy make every
overlap map a translation of $A$, giving
the actual principal torsor. At a multiple
value, normalize the actual base change
$z=w^{m_i}$ of $q$. The local equations
$z=u h^{m_i}$ become $v^{m_i}=u$, $w=vh$;
the normalized cover is \emph{\'etale} over
$S$ and is a smooth elliptic family. Its
fibres are connected: the proper Stein
factor over the root disc has generic
degree one and is the identity over
that normal disc. Its
deck action is free on the central elliptic
curve and hence is a translation of exact
order $m_i$, denoted $\tau_i\in A[m_i]$.
The translation can be made constant
on the whole root chart. If its product
expression is $a\mapsto a+b(w)$, lift
$b(w)-\tau_i$ to a holomorphic function
$\widetilde b(w)$ with value zero at
the origin. The order-$m_i$ relation
gives $\sum_{k=0}^{m_i-1}\widetilde b(\zeta^kw)=0$.
The coordinate change $a'=a+c(w)$, where
$c(w)=-m_i^{-1}\sum_{k=0}^{m_i-1}
 k\widetilde b(\zeta^kw)$, changes this
translation to exactly $\tau_i$.
The marked product description on the root
disc therefore has the local form
\[
 (A\times\Delta_w)/
 \langle(a,w)\mapsto(a+\tau_i,\zeta_{m_i}w)\rangle.
                                                               \tag{C.SF.2}\hypertarget{tag-010W}{}\label{cdp:eq:C.SF.2}
\]
This is precisely the original multiple-fibre
chart, including its torsor gluing on the
punctured disc, by the normalized-root and
translation comparison in \ref{cdp:st:C.2.UN.29}.
Thus the assertion identifies the given
$S$, rather than only another surface with
the same multiplicities.

\item We next lift the entire torsor to two
line bundles. The holomorphic group map
\[
 \begin{split}
 \Phi:(\mathbb C^*)^2&\longrightarrow A,\\
 (u,v)&\longmapsto
 \left[\frac{\lambda_1\log u+\lambda_2\log v}{2\pi\mathrm i}
                                                    \right]
 \end{split}
                                                               \tag{C.SF.3}\hypertarget{tag-010X}{}\label{cdp:eq:C.SF.3}
\]
is independent of both logarithm choices.
It is surjective, since its derivative
is surjective and its image is therefore
an open subgroup of the connected group
$A$. Its kernel is the closed copy of
$\mathbb C$ given by
\[
 w\longmapsto
 \bigl(\exp(2\pi\mathrm i\lambda_2w),
                  \exp(-2\pi\mathrm i\lambda_1w)\bigr).
                                                               \tag{C.SF.4}\hypertarget{tag-010Y}{}\label{cdp:eq:C.SF.4}
\]
To verify exactness, take logarithms of
a point in the kernel. Their linear
combination lies in $\Lambda$; subtract
the corresponding integral logarithm
periods to make that combination zero.
The resulting pair lies in the line
parametrized by \eqref{cdp:eq:C.SF.4}. Its exponential
map is injective, since an integral
relation between $\lambda_1$ and
$\lambda_2$ is zero. The image is the
kernel of \eqref{cdp:eq:C.SF.3}, hence is closed;
local logarithms give its holomorphic
inverse. This proves the exact group
sequence and the corresponding exact
sequence of analytic sheaves on
$\mathcal B$,
\[
 0\longrightarrow\mathcal O_{\mathcal B}
   \longrightarrow(\mathcal O_{\mathcal B}^*)^2
   \xrightarrow{\Phi}\mathcal A_{\mathcal B}
   \longrightarrow0,
                                                               \tag{C.SF.5}\hypertarget{tag-010Z}{}\label{cdp:eq:C.SF.5}
\]
where $\mathcal A$ is the sheaf of
holomorphic maps to the fixed group $A$.

\item The coarse root-curve map
$\pi:\mathcal B\to B$ has
$\pi_*\mathcal O=\mathcal O_B$ and
no higher coherent direct images.
Locally this is exactly the invariants
functor of a finite group in characteristic
zero; averaging makes it exact. Hence
$H^2(\mathcal B,\mathcal O)=H^2(B,\mathcal O_B)=0$.
The long exact sequence of the abelian
sheaves in \eqref{cdp:eq:C.SF.5} now gives a surjection
\[
 \operatorname{Pic}(\mathcal B)^2
                 \longrightarrow H^1(\mathcal B,\mathcal A).
                                                               \tag{C.SF.6}\hypertarget{tag-0110}{}\label{cdp:eq:C.SF.6}
\]
Choose a pair of line bundles
$\mathcal L_1,\mathcal L_2$ whose
image is the \emph{given} torsor
$\mathcal T$. This equality is an
equality of its analytic $H^1$ class,
so fixes its full holomorphic gluing.

\item Let $\mathcal R_i$ be the tautological
root line with
$\mathcal R_i^{m_i}=\pi^*\mathcal O_B(p_i)$.
Each line bundle on $\mathcal B$ has
an expression
\[
 \mathcal L_a\simeq\pi^*M_a\otimes
                  \bigotimes_i\mathcal R_i^{e_{ai}},
                  \qquad a=1,2.
                                                               \tag{C.SF.7}\hypertarget{tag-0111}{}\label{cdp:eq:C.SF.7}
\]
Indeed choose $e_{ai}$ to be its
inertia character modulo $m_i$.
After tensoring by the inverse root
characters it has trivial inertia;
its invariant local frames descend
to the ordinary line bundle $M_a$
on the coarse curve. This also
proves \eqref{cdp:eq:C.SF.7} by actual local
frame identifications.

\item Represent each $M_a$ by an ordinary
signed divisor
\[
                        M_a\simeq\mathcal O_B\left(
                                    \sum_p d_{ap}p\right)
                                                               \tag{C.SF.8}\hypertarget{tag-0112}{}\label{cdp:eq:C.SF.8}
\]
whose support avoids the finitely
many marked points $p_i$. It may
also avoid any one further chosen
point. For an exact choice, take
an effective divisor away from
that finite set, twist $M_a$ by
a sufficiently large multiple,
and use a section which is nonzero
at all points of the finite set;
curve Riemann--Roch and global
generation provide it. Subtracting
the chosen twist gives \eqref{cdp:eq:C.SF.8}.
This represents degree-zero classes
as well as classes of nonzero
degree.

\item We construct the auxiliary compact
surface. If $B$ is elliptic, take
$V_0=B\times\mathbb P^1$ with its
projection and selected component
$B=B\times\{0\}$. If $B=\mathbb P^1$,
start with any product
$A_0\times\mathbb P^1$ for an
elliptic curve $A_0$, blow up one
point over zero, and identify its
exceptional curve with the prescribed
$B$. Its single intersection with
the strict central elliptic curve
may be chosen away from all $p_i$;
choose \eqref{cdp:eq:C.SF.8} away from it as well.
In both cases the selected component
has multiplicity one and the
auxiliary central divisor has
simple normal crossings.

\item Blow up the distinct points of
the selected $B$ consisting of
all $p_i$ and all supports in
\eqref{cdp:eq:C.SF.8}. Call the resulting
surface $V$, its map to
$\mathbb P^1$ $h$, and the
exceptional curves $T_i,T_p$.
They are pairwise disjoint and
vertical, each meets the strict
transform of $B$ transversely
at its specified point, and
\[
 \mathcal O_V(T_i)|_B=\mathcal O_B(p_i),
 \qquad\mathcal O_V(T_p)|_B=\mathcal O_B(p).
                                                               \tag{C.SF.9}\hypertarget{tag-0113}{}\label{cdp:eq:C.SF.9}
\]
These are the actual restriction
maps of the exceptional divisors.
Every component of $h^{-1}(0)$
still has multiplicity one; its
reduced support is a tree of
smooth curves with no triple
points. The strict $B$ remains
isomorphic to the given curve.

\item Root $V$ along the disjoint
$T_i$ with their specified
orders, obtaining
$\mathcal V=V[(T_i,m_i)_i]$.
The inverse image of the strict
$B$ is exactly \eqref{cdp:eq:C.SF.1}, since
the intersections are transverse.
Write $\widetilde{\mathcal R}_i$
for the corresponding root lines
on $\mathcal V$. The vertical
extensions
\[
 \widetilde{\mathcal L}_a
   =\mathcal O_V\left(\sum_p d_{ap}T_p\right)
       \otimes\bigotimes_i
                       \widetilde{\mathcal R}_i^{e_{ai}}
                          \qquad(a=1,2)
                                                               \tag{C.SF.10}\hypertarget{tag-010S}{}\label{cdp:eq:C.SF.10}
\]
restrict, through \eqref{cdp:eq:C.SF.9}, to
exactly $\mathcal L_a$. Here
ordinary line bundles are pulled
back to $\mathcal V$ in the
displayed formula.

\item Push the pair in \eqref{cdp:eq:C.SF.10}
out by the actual group map
\eqref{cdp:eq:C.SF.3}. Explicitly, if their
nonvanishing local transition
functions are $g^1_{\alpha\beta}$
and $g^2_{\alpha\beta}$, glue
the products with $A$ by
\[
 (x,a)_\beta=(x,a+
     \Phi(g^1_{\alpha\beta}(x),
          g^2_{\alpha\beta}(x)))_\alpha.
                                                               \tag{C.SF.11}\hypertarget{tag-010T}{}\label{cdp:eq:C.SF.11}
\]
The triple-overlap equality follows
from the line-bundle cocycles and
the homomorphism property of $\Phi$.
Local logarithm choices change the
translation only by an element
of $\Lambda$. This produces a
principal $A$-torsor over
$\mathcal V$ whose restriction
is isomorphic to $\mathcal T$
by \eqref{cdp:eq:C.SF.6}--\eqref{cdp:eq:C.SF.10}.

\item Its coarse space $Y$ is smooth.
In a root chart transverse to
$T_i$ it is precisely
\[
 (A\times\Delta_w\times\Delta_x)/
 \langle(a,w,x)\mapsto
             (a+\tau_i,\zeta_{m_i}w,x)\rangle,
 \qquad z=w^{m_i}.
                                                               \tag{C.SF.12}\hypertarget{tag-010U}{}\label{cdp:eq:C.SF.12}
\]
The translation is the same
$\tau_i$ as in the original
torsor, because the inertia
characters in \eqref{cdp:eq:C.SF.7} are
preserved. Its exact order
is $m_i$, so the finite action
in \eqref{cdp:eq:C.SF.12} is free, including
on $w=0$. Thus the quotient
is a complex manifold. There
are no simultaneous root
intersections. Elsewhere it
is an ordinary elliptic bundle.
These charts also prove
properness over $V$: locally
the map is a finite quotient
of a product with the compact
curve $A$. Since $V$ is
compact, $Y$ is compact.

\item Compose with $h$ to obtain
$f:Y\to\mathbb P^1$.
The component over the strict
$B$ is exactly the coarse
torsor $S$ in \eqref{cdp:eq:C.SF.2}, with
the specified isomorphism
coming from its identical
$H^1$ class. At a generic
point no root divisor is
present and the auxiliary
component has multiplicity
one, so $S$ also has
multiplicity one in the
full fibre of $f$.

\item Every ordinary vertical line
in \eqref{cdp:eq:C.SF.10} has its canonical
meromorphic divisor section,
nonvanishing off the central
fibre. Every root line has
its tautological section,
also nonvanishing there.
Thus both line bundles are
canonically trivial on the
entire complement of zero,
and \eqref{cdp:eq:C.SF.11} gives the
actual product identification
\[
 f^{-1}(\mathbb P^1\setminus\{0\})
   \simeq A\times A_0\times
                         (\mathbb P^1\setminus\{0\}),
                                                               \tag{C.SF.13}\hypertarget{tag-010V}{}\label{cdp:eq:C.SF.13}
\]
where $A_0=B$ in the elliptic
base case. The fibres of
$Y\to V$ are connected
elliptic curves, including
the finite translation
quotients over root divisors.
Since $h$ has connected
fibres, so does $f$.

\item Finally the auxiliary central
support is an SNC curve
with no triple points.
Its inverse image in the
free quotient charts
\eqref{cdp:eq:C.SF.12} has coordinate
normal crossings. For
example a crossing whose
root coordinate is $z$
has local full equation
$w^{m_i}x=0$. Every
component of the reduced
support is a smooth
elliptic quasi-bundle over
one auxiliary curve; its
special fibres still have
smooth elliptic reduction.
Its Euler characteristic
is therefore zero by
additivity over the smooth
base stratum and the
finitely many multiple
values. Every double
curve is one smooth
elliptic translation
quotient, and there are
no triple points.
Inclusion--exclusion gives
$e(f^{-1}(0)_{\mathrm{red}})=0$.
All other fibres in
\eqref{cdp:eq:C.SF.13} are smooth tori,
so have the same zero
Euler characteristic.
No balancing modification
of the given $S$ is needed.

\end{enumerate}
\end{proof}






\begin{proposition}[Every prescribed cyclic smooth-multiple elliptic quotient occurs in a zero-Euler torus filling]\hypertarget{tag-01DE}{}\label{cdp:st:C.2.UN.35}
\StatementRecord{cdp:st:C.2.UN.35}{proved}{01DE}{\SeriesRef{cdp:st:C.2.UN.23}, \SeriesRef{cdp:st:C.2.UN.32}, \SeriesRef{cdp:st:C.2.UN.29}, \SeriesRef{cdp:eq:C.EO.1}, \SeriesRef{cdp:eq:C.EO.3}, \SeriesRef{cdp:eq:C.EO.5}, \SeriesRef{cdp:eq:C.EO.H3}, \SeriesRef{cdp:eq:C.EO.H1}, \SeriesRef{cdp:eq:C.EO.H4}, \SeriesRef{cdp:eq:C.EO.H2}, \SeriesRef{cdp:eq:C.EO.H7}, \SeriesRef{cdp:eq:C.EO.H5}, \SeriesRef{cdp:eq:C.EO.H6}, \SeriesRef{cdp:eq:C.EO.7}, \SeriesRef{cdp:eq:C.EO.9}, \SeriesRef{cdp:eq:C.EO.10}, \SeriesRef{cdp:eq:C.EO.4}, \SeriesRef{cdp:eq:C.EO.11}, \SeriesRef{cdp:eq:C.EO.14}, \SeriesRef{cdp:eq:C.EO.6}, \SeriesRef{cdp:eq:C.EO.12}, \SeriesRef{cdp:eq:C.EO.13}, \SeriesRef{cdp:st:C.2.UN.1}, \SeriesRef{cdp:st:C.2.UN.4}, \SeriesRef{cdp:st:C.2.UN.31}, \SeriesRef{cdp:st:C.2.UN.22}, \SeriesRef{cdp:st:C.2.UN.20}, \SeriesRef{cdp:st:C.2.UN.24}; \cite[Theorem~1.9]{Steinberg1965}}

Let $q:E\to B$ be a connected compact smooth elliptic
surface with a fixed elliptic Jacobian $F$, identity
integral $H_1$ monodromy, and only smooth multiple
fibres. Suppose $g(B)\le1$. Let $Q$ be a finite
cyclic group acting holomorphically on $E$, preserving
$q$, and acting faithfully on $B$. The trivial group
is allowed. The smooth minimal model of $E/Q$ occurs
as the minimal normalization model of a component of
a fibre in a smooth compact threefold over $\PP^1$
whose general fibre is a complex two-torus. It also
occurs in such a fibre with reduced-support Euler
number zero. The construction retains the given
surface and its given action before taking the
quotient: the distinguished component of the
constructed covering threefold is $Q$-equivariantly
isomorphic to $E$ over $B$.

Consequently the minimal normalization models of
nonnegative Kodaira dimension occurring in arbitrary
two-torus degenerations are exactly the following
two explicitly specified classes:
\begin{itemize}
\item the minimal resolutions of finite cyclic
quotients of complex two-tori that have nonnegative
Kodaira dimension;
\item the minimal models of $E/Q$ in the preceding
construction that have nonnegative Kodaira dimension.
\end{itemize}
The same two classes occur among zero-Euler fillings.
The torus-quotient class is described completely in
\ref{cdp:st:C.2.UN.23} and \ref{cdp:st:C.2.UN.32}. In the second class, all multiplicities,
translation gluing, and the cyclic action are part
of the specified input; none is discarded by the
occurrence construction.
\end{proposition}



\begin{proof}\leavevmode\par

\begin{enumerate}[label={$\langle1\rangle$\arabic*.},leftmargin=0pt,labelsep=.5em]
\item  We give the extension on actual finite root charts,
and then prove the exhaustive assertion.


\item \par\medskip\noindent\textbf{The actual orbicurve torsor.}
Write the multiple fibres as $m_iC_i$ over $p_i\in B$,
with $m_i>1$. The normalized root chart proved in
\ref{cdp:st:C.2.UN.29} and (C.KX.3)--(C.KX.5) is
\[
 (F\times\Delta_u)/\langle\delta_i\rangle,
 \qquad
 \delta_i(z,u)=(z+\tau_i,\zeta_{m_i}u),
 \qquad \operatorname{ord}(\tau_i)=m_i,
 \qquad t=u^{m_i}.
                                                               \tag{C.EO.1}\hypertarget{tag-00HE}{}\label{cdp:eq:C.EO.1}
\]
The free translation on $F$ is essential here.
Let $\mathcal B=B[(p_i,m_i)]$ denote the root
orbicurve, whose chart at $p_i$ is
$[\Delta_u/\mu_{m_i}]$ with $t=u^{m_i}$.
Changes of local origins of the fixed elliptic
curve are translations: the linear part is the
identity because the integral monodromy is the
identity. Thus these charts and the ordinary
product charts exhibit the given $E$ as a
principal $F$-torsor over $\mathcal B$.
Its total space is the original smooth surface,
because all the root inertia actions in \eqref{cdp:eq:C.EO.1}
are free. The action of $Q$ induces the functorial
action on the root orbicurve, and its linear action
on $F$ is a constant finite character
\[
                 \beta:Q\longrightarrow
                            \operatorname{Aut}(F,0).
                                                               \tag{C.EO.2}\hypertarget{tag-00HM}{}\label{cdp:eq:C.EO.2}
\]
The constancy follows from the discreteness of
$\operatorname{Aut}(F,0)$ and connectedness of
the regular base. The actual action on the
torsor includes its translation part.


\item \par\medskip\noindent\textbf{A lift that retains the whole action.}
Put $F=\C/\Lambda$ and let
$\ell:\Lambda\otimes_{\Z}\C\to\C$
be the complex-linear extension of the inclusion
of the period lattice. There is an exact sequence
of complex Lie groups
\[
 0\longrightarrow K\longrightarrow
 T=(\Lambda\otimes_{\Z}\C)/\Lambda
 \xrightarrow{\,\Phi\,}F\longrightarrow0,
 \qquad K=\ker\ell\simeq(\C,+),
 \qquad T\simeq(\C^*)^2.                              \tag{C.EO.3}\hypertarget{tag-00HN}{}\label{cdp:eq:C.EO.3}
\]
For a basis $(\lambda_1,\lambda_2)$ of $\Lambda$,
the last map is explicitly
\[
 \Phi(z_1,z_2)=
 \frac{\lambda_1\log z_1+\lambda_2\log z_2}{2\pi i}
                    \pmod\Lambda .                 \tag{C.EO.4}\hypertarget{tag-00HO}{}\label{cdp:eq:C.EO.4}
\]
Changing either logarithm changes the numerator
by the corresponding period. The map is locally
split and surjective. Its kernel is exactly
the image of the complex line $\ker\ell$:
a vector mapping to a period can first be
changed by that integral lattice vector.
The intersection $\ker\ell\cap\Lambda$ is zero,
and the kernel is closed because $F$ is Hausdorff.
This proves the exactness, including the actual
additive kernel, without a choice of a cyclic
subgroup of the elliptic periods.


\begin{enumerate}[label={$\langle2\rangle$\arabic*.},leftmargin=1.3em,labelsep=.5em]
\item  Every $\beta(g)$ acts on $\Lambda$ by an integral
matrix and hence on $T$ by the corresponding
monomial automorphism. It preserves $K$; on
that line its character is $\beta(g)^{-1}$,
because its determinant on $\Lambda$ is one.
Sequence \eqref{cdp:eq:C.EO.3} is therefore $Q$-equivariant.
On the quotient orbicurve $\mathcal C=[\mathcal B/Q]$
it gives an exact sequence of sheaves of
holomorphic sections
\[
 0\longrightarrow\mathcal K\longrightarrow
 \mathcal T\longrightarrow\mathcal F\longrightarrow0.
                                                               \tag{C.EO.5}\hypertarget{tag-00HP}{}\label{cdp:eq:C.EO.5}
\]
Here $\mathcal K$ is a holomorphic line bundle
with its additive group law; the other sheaves
retain the indicated finite monodromy. The given
equivariant elliptic torsor is an element of
$H^1(\mathcal C,\mathcal F)$.


\item  One has $H^2(\mathcal C,\mathcal K)=0$.
Indeed finite-group invariants on complex vector
spaces are exact by averaging. On each finite
root chart the higher invariant direct images
therefore vanish, and the invariant direct image
of a coherent sheaf is coherent on the coarse
curve. Coherent cohomology in degree two on a
compact complex curve is zero. The same averaging
for $Q$ proves the asserted vanishing on
$\mathcal C$. The long exact sequence of
\eqref{cdp:eq:C.EO.5} consequently supplies an actual
$\mathcal T$-torsor $P$ and an isomorphism
\[
                  P\mathbin{\times^{T}}F
                           \simeq E                 \tag{C.EO.6}\hypertarget{tag-00HQ}{}\label{cdp:eq:C.EO.6}
\]
in the equivariant orbicurve category. In
particular the descent isomorphisms of $P$
satisfy the full cocycle identity; they are
not merely isomorphisms of underlying line
bundle classes.


\end{enumerate}

\item \par\medskip\noindent\textbf{An invariant meromorphic frame.}
The faithful action on $B$ makes
$L=\C(B)$ a cyclic extension of
$k=\C(B/Q)$. At the generic point the root
inertia is absent. The monomial descent of
$(\C^*)^2_L$ defines an algebraic $k$-torus
$T_k$, and the generic fibre of $P$ is an
algebraic torsor under it. To see algebraicity
directly, pull back to $\mathcal B$ and choose
meromorphic sections of its two associated
orbifold line bundles. Such sections exist:
after separating the finitely many root weights,
this is the ordinary meromorphic-section theorem
for line bundles on a compact curve. The
equivariant descent isomorphisms are then
monomials multiplied by meromorphic functions
on $B$, which belong to $L$. Their existing
cocycle identities give precisely this algebraic
generic torsor.


\begin{enumerate}[label={$\langle2\rangle$\arabic*.},leftmargin=1.3em,labelsep=.5em]
\item  It has a $k$-rational point. This is the
rank-two torus case of Steinberg's theorem
\cite[Theorem~1.9]{Steinberg1965}; the independent
curve-field and cyclic-lattice proof below
establishes the required vanishing for exactly
the representations used here.

\end{enumerate}

\item \par\medskip\noindent\textbf{An independent calculation of the generic torus torsor.}
Here is a direct verification of the required vanishing
for the actual cyclic rank-two torus. Put
$K=\mathbb C(B/Q)$. The splitting action on its
character lattice has order $d\in\{1,2,3,4,6\}$;
if the base group is larger, pass first to the fixed
field of the kernel of this lattice action. This gives
the actual cyclic splitting extension $L/K$ of degree $d$.
We prove
\[
                         H^1(K,T)=0.                 \tag{C.EO.H1}\hypertarget{tag-00HU}{}\label{cdp:eq:C.EO.H1}
\]
The cohomology here is the Galois cohomology of the
algebraic torus obtained by descent of the two-dimensional
split multiplicative group, and classifies its actual
torsors over $K$.


\begin{enumerate}[label={$\langle2\rangle$\arabic*.},leftmargin=1.3em,labelsep=.5em]
\item  First every function field $K'$ of a smooth compact
complex curve has the following elementary property:
a homogeneous polynomial of degree $a$ in $b>a$
variables has a nonzero zero over $K'$. To verify it,
let $g$ be the genus of the curve, choose a point $P$,
and choose an effective divisor $D$ bounding the poles
of all coefficients. Seek each variable in
$H^0(\OO(rP))$. For large $r$ there are
$b(r+1-g)$ scalar unknowns, whereas the polynomial
value belongs to $H^0(\OO(arP+D))$, of dimension
$ar+\deg D+1-g$. Its vanishing imposes at most that
many homogeneous equations of positive degree.
For large $r$ the number of unknowns is strictly
larger. The height bound for an ideal generated by
that many elements gives an affine homogeneous zero
set of dimension at least one. Over $\mathbb C$ it
contains a point other than the origin. This yields
the asserted nonzero zero over $K'$.


\item  Consequently, for every finite field extension
$L'/K'$ the norm map is surjective:
\[
               N_{L'/K'}:(L')^*\longrightarrow(K')^*.
                                                               \tag{C.EO.H2}\hypertarget{tag-00HV}{}\label{cdp:eq:C.EO.H2}
\]
Indeed, for $c\in(K')^*$, choose a $K'$-basis of
$L'$ of size $a=[L':K']$ and apply the preceding
property to
$N_{L'/K'}(x_1e_1+\cdots+x_ae_a)-cz^a$.
This has degree $a$ in $a+1$ variables.
A nonzero zero cannot have $z=0$, because the norm
of a field element vanishes only for the zero element.
Dividing by $z$ gives an element of norm $c$.
Every intermediate field used below is again the
function field of a smooth compact complex curve,
so the same proof applies to it.


\item  The elliptic lattice action has, after an integral
choice of basis, the following matrices:
\[
\begin{array}{c|c}
d&M_d\\ \hline
1&I_2\\
2&-I_2\\
3&\begin{pmatrix}0&-1\\1&-1\end{pmatrix}\\
4&\begin{pmatrix}0&-1\\1&0\end{pmatrix}\\
6&\begin{pmatrix}0&-1\\1&1\end{pmatrix}.
\end{array}                                                   \tag{C.EO.H3}\hypertarget{tag-00HW}{}\label{cdp:eq:C.EO.H3}
\]
For completeness, a rank-two lattice preserved by a
primitive automorphism of order $3,4$, or $6$ is a
rank-one module over the corresponding ring
$\mathbb Z[\zeta_d]$. After tensoring with $\mathbb Q$
and choosing a nonzero lattice vector, it becomes
a fractional ideal of that ring. The Gaussian and
Eisenstein rings are Euclidean for their usual norm,
so this ideal is principal. The same argument applies
to the dual character lattice. It gives the matrices
in \eqref{cdp:eq:C.EO.H3}; replacing a generator by its inverse
gives an integrally conjugate matrix. The order-two
case is scalar multiplication by $-1$.


\item  If $d=1$, the torus is split and Hilbert's Theorem 90
gives \eqref{cdp:eq:C.EO.H1}. For $d=2$, it is the product of two
quadratic norm-one tori, since each character lattice
is the sign representation. For $d=3$, it is the
cubic norm-one torus. To specify this identification,
the exact sequence
\[
 1\longrightarrow R^1_{L/K}\mathbb G_m
 \longrightarrow R_{L/K}\mathbb G_m
 \xrightarrow{\ N\ }\mathbb G_m\longrightarrow1
                                                               \tag{C.EO.H4}\hypertarget{tag-00HX}{}\label{cdp:eq:C.EO.H4}
\]
has character quotient
$\mathbb Z[C_3]/\mathbb Z(1+g+g^2)$.
In the basis $[1],[g]$, multiplication by $g$ is
exactly $M_3$. In degree two the same quotient is
the sign lattice. Galois descent of characters
therefore gives the stated isomorphisms of tori.
The long exact sequence of \eqref{cdp:eq:C.EO.H4} and Hilbert 90
identify its $H^1$ with $K^*/N_{L/K}L^*$, which is
zero by \eqref{cdp:eq:C.EO.H2}.


\item  For $d=4$, let $K_2$ be the degree-two intermediate
field. The torus is
\[
        T\simeq R_{K_2/K}
                  \bigl(R^1_{L/K_2}\mathbb G_m\bigr).
                                                               \tag{C.EO.H5}\hypertarget{tag-00HY}{}\label{cdp:eq:C.EO.H5}
\]
This is an exact character-lattice identification:
induce the sign module of $\langle g^2\rangle$
to $\langle g\rangle$; with basis $e,ge$, the
relations $g(e)=ge$ and $g(ge)=-e$ give $M_4$.
Weil restriction and the corresponding Shapiro
map identify its $H^1$ over $K$ with the $H^1$
of the norm-one torus over $K_2$. That group is
$K_2^*/N_{L/K_2}L^*$ and vanishes by \eqref{cdp:eq:C.EO.H2}.


\item  For $d=6$, let $K_2,K_3$ be the degree-two and
degree-three intermediate fields. On restriction
to $K_2$, a generator acts by $M_6^2$; on restriction
to $K_3$ it acts by $M_6^3=-I_2$. The integral matrix
\[
 P=\begin{pmatrix}1&-1\\0&1\end{pmatrix},\qquad
                  P^{-1}M_6^2P=M_3                  \tag{C.EO.H6}\hypertarget{tag-00HZ}{}\label{cdp:eq:C.EO.H6}
\]
identifies the first restricted torus with the
cubic norm-one torus, and the second with a product
of quadratic norm-one tori. Thus both restricted
$H^1$ groups vanish.


\item  We spell out the transfer used to return to $K$.
For a finite extension $K'/K$, put
$\mathcal G=\operatorname{Gal}(K^{\rm sep}/K)$,
$\mathcal H=\operatorname{Gal}(K^{\rm sep}/K')$,
and $A=T(K^{\rm sep})$, written additively for this
calculation. Restriction is restriction of continuous
one-cocycles to $\mathcal H$. Choose left coset
representatives $R$ for $\mathcal G/\mathcal H$;
write $\sigma r=r_\sigma h_{\sigma,r}$.
For a one-cocycle $c$ on $\mathcal H$ the transfer
is represented by
\[
       (\operatorname{cor}c)(\sigma)
            =\sum_{r\in R}r_\sigma\,
                         c(h_{\sigma,r}).             \tag{C.EO.H7}\hypertarget{tag-00I0}{}\label{cdp:eq:C.EO.H7}
\]
The cocycle identity verifies that this is a cocycle.
For $c=\operatorname{res}b$, the identity
$b(\sigma r)=b(\sigma)+\sigma b(r)
=b(r_\sigma)+r_\sigma b(h_{\sigma,r})$
shows that \eqref{cdp:eq:C.EO.H7} differs from
$[K':K]b(\sigma)$ by the coboundary of
$\sum_{r\in R}b(r)$. Thus
$\operatorname{cor}\circ\operatorname{res}$ is
multiplication by $[K':K]$ on $H^1(K,T)$.
The restrictions to $K_2$ and $K_3$ just proved
to vanish therefore kill every class by both
two and three. Their coprimality gives \eqref{cdp:eq:C.EO.H1}.


\item  This proves the required generic torsor vanishing
in every possible elliptic character, with the
actual restriction and norm maps specified.
The integral identities \eqref{cdp:eq:C.EO.H3}, \eqref{cdp:eq:C.EO.H5},
and \eqref{cdp:eq:C.EO.H6}, together with the compatible
elliptic quotient and additive-kernel characters,
are independently executed in
\path{code/A_equivariant_torus.py} and its saved output.


\item  The curve-field proof in that calculation also
checks the precise field hypothesis in the
printed statement of Steinberg's theorem.
It applies to every finite extension of $k$,
since that extension is another complex curve
function field. A central division algebra
of degree $d>1$ over any such field has a
reduced norm of degree $d$ in $d^2$ variables.
The proved $C_1$ property would give a nonzero
element of reduced norm zero, impossible in
a division algebra. Thus all the indicated
division algebras are commutative, which is
the cohomological-dimension-at-most-one
condition used in \cite[Theorem~1.9]{Steinberg1965}.
Pull back that rational point to $L$.
It is a meromorphic section $s$ of $P$ on
$\mathcal B$ satisfying the exact equality
\[
                         g(s(b))=s(gb)
                              \qquad(g\in Q).        \tag{C.EO.7}\hypertarget{tag-00HR}{}\label{cdp:eq:C.EO.7}
\]
Thus the translation cocycle has been retained
in a genuine invariant meromorphic frame.


\item  Let $M=\operatorname{Hom}(T,\C^*)$ be the
character lattice. For $\chi\in M$, write
$L_\chi$ for the associated line bundle and
$s_\chi$ for its meromorphic section induced
by $s$. Its divisor defines an additive map
\[
 D:M\longrightarrow\operatorname{Div}(\mathcal B),
       \qquad D(\chi)=\operatorname{div}(s_\chi).
                                                               \tag{C.EO.8}\hypertarget{tag-00HS}{}\label{cdp:eq:C.EO.8}
\]
This map is $Q$-equivariant with its actual
action on characters and divisors, by \eqref{cdp:eq:C.EO.7}.
Choose two basis characters. Their supports,
together with the root points, lie in one
finite $Q$-invariant subset $\Sigma\subset B$.
In terms of coarse points, every coefficient
of $D(\chi)$ is integral away from the root
points and belongs to $m_i^{-1}\Z$ at $p_i$.
We have the canonical multiplicative
identifications
\[
          (L_\chi,s_\chi)
                   \simeq(\OO_{\mathcal B}(D(\chi)),1),
                                                               \tag{C.EO.9}\hypertarget{tag-00HT}{}\label{cdp:eq:C.EO.9}
\]
where $1$ denotes the canonical meromorphic
divisor section. They commute with $Q$ and
with addition of characters. There is no
remaining scalar linearization or group
two-cocycle in \eqref{cdp:eq:C.EO.9}.


\end{enumerate}

\item \par\medskip\noindent\textbf{An equivariant vertical extension.}
We construct a smooth projective surface
$V_0\to\PP^1_t$ containing $B$ in its
central fibre and extending its given
$Q$-action. Let $n=|Q|$.
For elliptic $B$, choose any faithful
character $\xi:Q\to\C^*$ when $n>1$
(and $\xi=1$ for the trivial group),
and take $V_0=B\times\PP^1_t$
with action $(b,t)\mapsto(gb,\xi(g)t)$,
and use $B\times\{0\}$.
For $B\simeq\PP^1$, first diagonalize
its given cyclic action as
$z\mapsto\lambda(g)z$ and define
$\xi=\lambda^{-1}$. Faithfulness makes
this a faithful disc character too.
For the trivial group use both
characters equal to one. Choose any
elliptic curve $A$ and point $o\in A$.
Blow up $(o,0)$ in $A\times\PP^1_t$,
where the action is
$(a,t)\mapsto(a,\xi(g)t)$.
The exceptional curve is identified with
the given $B$: in a local coordinate $x$
at $o$, its affine coordinate $z=x/t$
has exactly the required action. The
other chart supplies its point at infinity.


\begin{enumerate}[label={$\langle2\rangle$\arabic*.},leftmargin=1.3em,labelsep=.5em]
\item  Blow up every point of $\Sigma$ on the
distinguished strict curve $B$, in complete
$Q$-orbits. Call the resulting disjoint
exceptional curves $T_p$, $p\in\Sigma$,
and call the new surface $V$. The strict
curve is still isomorphic to $B$, and
\[
 \OO_V(T_p)|_B=\OO_B(p),\qquad
                     g(T_p)=T_{gp}.                \tag{C.EO.10}\hypertarget{tag-00HF}{}\label{cdp:eq:C.EO.10}
\]
These identities include the canonical
sections and their equivariant structures.
There is no restriction at a fixed point
or at the node of the original central
fibre. At a node with coordinates $xy=0$
and $B=(x=0)$, its blowup has chart
$x=uv,y=v$. Here the strict $B$ is
$u=0$, $T_p$ is $v=0$, and the
restriction of its canonical section
is the simple parameter $v$ on $B$.
This proves \eqref{cdp:eq:C.EO.10} also in that case.
The other central branch meets $T_p$
at the other point of this exceptional
curve; it does not meet the strict $B$.


\item  Root $V$ along each $T_{p_i}$ to order
$m_i$, obtaining the orbifold surface
$\mathcal V$. Its local chart at a
point of $T_{p_i}$ is
\[
 [\Delta_u\times\Delta_v/\mu_{m_i}],
           \qquad x=u^{m_i},\qquad
           \zeta\cdot(u,v)=(\zeta u,v).
                                                               \tag{C.EO.11}\hypertarget{tag-00HG}{}\label{cdp:eq:C.EO.11}
\]
The $T_{p_i}$ are disjoint, so there
is only this one inertia group in
each such chart. Its restriction to
$B$ is precisely the previous root
orbicurve $\mathcal B$, including the
functorial $Q$-action. The transverse
intersection in \eqref{cdp:eq:C.EO.10} is what
identifies its root line and root
section with those on $\mathcal B$.


\item  Replace each coarse point $p$ in
$D(\chi)$ by $T_p$, with the same
coefficient. Denote the resulting
orbifold Cartier divisor by
$\widetilde D(\chi)$. All its
coefficients are legitimate integral
root-divisor coefficients on
$\mathcal V$. The map
$\chi\mapsto\widetilde D(\chi)$ is
additive and $Q$-equivariant because
the replacement $p\mapsto T_p$ is.
The character lines
$\OO_{\mathcal V}(\widetilde D(\chi))$
therefore define a principal $T$-torsor
$\widetilde P$ with its exact monomial
$Q$-action. Define the principal
elliptic bundle
\[
              \widetilde Y
                =\widetilde P\mathbin{\times^T}F
                     \longrightarrow\mathcal V .   \tag{C.EO.12}\hypertarget{tag-00HH}{}\label{cdp:eq:C.EO.12}
\]
Concretely, on every line-bundle
trivialization its transition map is
translation on $F$ by the value under
\eqref{cdp:eq:C.EO.4} of the two multiplicative
transition functions. Equations
\eqref{cdp:eq:C.EO.9}--\eqref{cdp:eq:C.EO.10} give an actual
$Q$-equivariant isomorphism
\[
                   \widetilde Y|_{\mathcal B}
                                   \simeq E,         \tag{C.EO.13}\hypertarget{tag-00HI}{}\label{cdp:eq:C.EO.13}
\]
over the given base, not merely
an equality of topological Chern
classes. Locally the meromorphic
frame identification extends
holomorphically because the two
divisors agree: their quotient is
a nowhere-zero holomorphic unit.
The same assertion holds in the
root coordinates. Since both
frames in \eqref{cdp:eq:C.EO.9} are exactly
equivariant, these local isomorphisms
also retain the given $Q$-action.
More explicitly, a lift of a base
map to $x=u^{m_i}$ involves a
choice of an $m_i$-th root of
its nonzero local multiplier.
Two choices differ by the
root inertia in \eqref{cdp:eq:C.EO.11}.
The character lines simultaneously
apply the corresponding elliptic
translation in \eqref{cdp:eq:C.EO.14}.
Hence the two lifts induce the
same map on the total coarse
quotient. Compositions of lifts
agree there for the same reason.
This is the finite-chart form
of the canonical section-preserving
action, and proves that no local
root choice changes its group law.


\end{enumerate}

\item \par\medskip\noindent\textbf{Smoothness, properness and the generic torus.}
Let $(a_1,a_2)\pmod {m_i}$ be the
two weights of $P$ at $p_i$.
In \eqref{cdp:eq:C.EO.11} the total elliptic
bundle has the finite chart
\[
 (\Delta_u\times\Delta_v\times F)/\mu_{m_i},
 \quad
 \zeta_{m_i}(u,v,z)
     =(\zeta_{m_i}u,v,z+\tau_i),\quad
 \tau_i=\frac{a_1\lambda_1+a_2\lambda_2}{m_i}
                                           \pmod\Lambda .
                                                               \tag{C.EO.14}\hypertarget{tag-00HJ}{}\label{cdp:eq:C.EO.14}
\]
The order of $\tau_i$ is exactly
$m_i$ by \eqref{cdp:eq:C.EO.1} and \eqref{cdp:eq:C.EO.6}.
Every nonidentity inertia element
therefore acts freely on the total
chart, including $u=0$ and every
central crossing $v=0$. Thus the
coarse space in \eqref{cdp:eq:C.EO.12} is a
smooth compact complex threefold;
we continue to call it
$\widetilde Y$. Properness follows
on these finite charts from
properness of the compact factor
$F$ and glues over the compact
surface $V$. The principal-bundle
transition functions in \eqref{cdp:eq:C.EO.12}
give the gluing maps explicitly.


\begin{enumerate}[label={$\langle2\rangle$\arabic*.},leftmargin=1.3em,labelsep=.5em]
\item  All divisors used to construct
$\widetilde P$ are vertical over
$t=0$. Their canonical meromorphic
sections are therefore nowhere
zero on the complement of that
fibre. There \eqref{cdp:eq:C.EO.12} is the
actual product
\[
 \widetilde Y|_{t\ne0}
       \simeq A\times F\times(\PP^1\setminus\{0\}),
                                                               \tag{C.EO.15}\hypertarget{tag-00HK}{}\label{cdp:eq:C.EO.15}
\]
where $A=B$ in the elliptic-base
case and $A$ is the chosen auxiliary
elliptic curve in the rational-base
case. The action is the prescribed
action on $A$ (the identity in the
rational-base construction),
the linear action $\beta$ on $F$,
and multiplication by $\xi$ on
$t$. This formula follows from
the invariant divisor frames,
so no additional translation
cocycle survives off zero.


\item  Take the finite quotient by $Q$
and resolve its singularities.
Finite holomorphic actions can
be linearized locally, and these
cyclic quotients have compatible
proper resolutions; a resolution
of the compact quotient, chosen
over its singular locus, is
enough. Its map to
$\PP^1_T$, $T=t^n$, has general
fibre exactly $A\times F$:
the $n$ covering fibres are
permuted simply transitively
for $T\ne0,\infty$. The quotient
is smooth there, so the
resolution leaves those fibres
unchanged. The Stein factor is
a finite map of degree one
to the normal curve $\PP^1$,
and is the identity. Thus
the resulting proper map has
connected fibres.


\item  The distinguished component
before the quotient is exactly
the given $E$ by \eqref{cdp:eq:C.EO.13}.
Its image in the finite quotient
is $E/Q$: taking invariants
is exact in characteristic zero,
so the restriction of the
quotient structure sheaf to
this image is
$(\OO_E)^Q$. It is normal
because $E$ is normal. The
strict component on the ambient
resolution maps properly and
birationally to $E/Q$.
Resolving that component and
contracting its surface
$(-1)$-curves gives exactly
the prescribed minimal model.
Finally \ref{cdp:st:C.2.UN.1} gives even Euler
number for the reduced support.
The actual vertical point and
quartic blowups of \ref{cdp:st:C.2.UN.4} change
that number by $+2$ and $-2$
respectively. Finitely many
such operations make it zero
and preserve the minimal
normalization model through
their explicit proper
birational maps. This proves
the occurrence assertion.


\end{enumerate}

\item \par\medskip\noindent\textbf{Exhaustion of the nonnegative classes.}
Let $Z$ be a minimal normalization
model of nonnegative Kodaira
dimension in an arbitrary
two-torus degeneration. \ref{cdp:st:C.2.UN.31}
provides its actual cyclic
component cover and equivariant
minimal contraction. A torus
cover gives exactly the first
class in the statement. In
the remaining case \ref{cdp:st:C.2.UN.22} and
\ref{cdp:st:C.2.UN.29} give a fixed-Jacobian,
identity-monodromy fibration
with only smooth multiples
over a rational or elliptic
curve. We specify why an
appropriate such fibration
is preserved by the cyclic
action.


\begin{enumerate}[label={$\langle2\rangle$\arabic*.},leftmargin=1.3em,labelsep=.5em]
\item  For a cover of Kodaira
dimension one it is the
intrinsic connected Iitaka
fibration. In algebraic
dimension one its base is
the intrinsic meromorphic
function field, and the
equivariant Stein factor
of \ref{cdp:st:C.2.UN.20} applies. The
non-torus projective case
of Kodaira dimension zero
is bielliptic. Use its
actual finite free product
description
$(F_1\times F_2)/\Gamma$,
with $\Gamma$ acting by
translations on $F_1$
and with a nontrivial
linear action on $F_2$.
The invariant holomorphic
vector fields on the
product are exactly the
constant fields in the
$F_1$ direction. These
integrate to the connected
automorphism group of
the quotient, an elliptic
quotient of $F_1$.
Its orbit map is
\[
 (F_1\times F_2)/\Gamma
                   \longrightarrow F_2/\Gamma\simeq\PP^1.
                                                               \tag{C.EO.16}\hypertarget{tag-00HL}{}\label{cdp:eq:C.EO.16}
\]
Every automorphism normalizes
the connected automorphism
group and hence preserves
this orbit fibration. Its
local stabilizers act by
free translations on
$F_1$; its root charts
are \eqref{cdp:eq:C.EO.1}. Therefore
it has the required
identity monodromy and
only smooth multiples.
This is the preserved
fibration used in the
bielliptic case.


\item  Let $H$ be the kernel of
the cyclic action on the
chosen base. Its action
on a general elliptic
fibre is by translations:
a nontrivial linear
elliptic automorphism
has rational quotient,
which would make the
quotient surface ruled
and annihilate all
positive pluricanonical
sections, contrary to
$\kappa(Z)\ge0$.
The exact translation
and reflection charts
(C.EQ.3)--(C.EQ.4)
show that quotienting
the cover by $H$ is
again a smooth elliptic
surface with fixed
Jacobian, identity
monodromy, and only
smooth multiples.
The residual cyclic
group acts faithfully
on the base. Its
quotient is the same
actual surface quotient
as before, through
the finite quotient
maps. It is therefore
one of the inputs
already constructed
above.


\item  Conversely, the
construction proves
occurrence for every
member of the second
class. For the first,
take the compact
diagonal quotient
$(A\times\PP^1_t)/G$
using a faithful
character on $t$,
resolve it, and use
the same \ref{cdp:st:C.2.UN.4} balancing;
the exact component
comparison is \ref{cdp:st:C.2.UN.24}.
All occurrences may
therefore be chosen
with zero-Euler
support, proving
both asserted
equalities of classes.



\end{enumerate}

\end{enumerate}
\end{proof}






\begin{proposition}[Occurrence of every Hopf surface of algebraic dimension one]\hypertarget{tag-01DG}{}\label{cdp:st:C.2.UN.37}
\StatementRecord{cdp:st:C.2.UN.37}{proved}{01DG}{\SeriesRef{cdp:st:C.2.UN.34}, \SeriesRef{cdp:eq:C.EH.2}, \SeriesRef{cdp:eq:C.EH.4}, \SeriesRef{cdp:eq:C.EH.6}, \SeriesRef{cdp:eq:C.EH.7}}

Let $H$ be any primary or secondary Hopf surface with $a(H)=1$.
There is a proper holomorphic map from a smooth compact complex
threefold to $\mathbb P^1$, with connected fibres and smooth
two-tori away from one value, whose fibre over that value contains
$H$ itself as a smooth component of multiplicity one. Its reduced
special fibre has simple normal crossings and Euler number zero.
Over the complement of the special value the family is the
actual product of two fixed elliptic curves with the base.

The map used to apply \ref{cdp:st:C.2.UN.34} is the actual connected elliptic
reduction
\[
                         q:H\longrightarrow\mathbb P^1.
\tag{C.EH.1}\hypertarget{tag-00FQ}{}\label{cdp:eq:C.EH.1}
\]
Every fibre of $q$ is a multiple of a smooth elliptic curve,
the modulus is constant, and the integral $H^1$ monodromy is
the identity. The assertion includes secondary Hopf surfaces
without imposing a cyclic condition on their finite primary cover.
\end{proposition}



\begin{proof}

\begin{enumerate}[label={$\langle1\rangle$\arabic*.},leftmargin=*]
\item We first construct \eqref{cdp:eq:C.EH.1} and check its full fibres directly
on the given surface. For any Hopf surface one has
\[
                    b_1(H)=1,\qquad H^2(H,\mathbb Q)=0.
\tag{C.EH.2}\hypertarget{tag-00FR}{}\label{cdp:eq:C.EH.2}
\]
These statements follow through the actual finite unramified
primary cover: the primary cover has the rational cohomology
of $S^1\times S^3$, transfer gives $b_1\leq1$ and $b_2=0$,
and the nonzero deck character
$\gamma\mapsto\log|\det(d\gamma)_0|$ gives $b_1\geq1$.
This is the comparison proved in (C.HQ.5)--(C.HQ.7).
Thus it applies to secondary surfaces as well.

\item Choose a nonconstant meromorphic function $h$ on $H$.
Its effective zero and pole divisors have no common
irreducible component and have intersection number zero
by \eqref{cdp:eq:C.EH.2}. Local intersection multiplicities of distinct
effective curves on a smooth surface are positive.
Consequently those two divisors are disjoint. In a local
coprime expression $h=f/g$, an indeterminacy point would
belong to both $(f=0)$ and $(g=0)$, so no such point exists.
The function therefore defines an actual holomorphic map
$h:H\to\mathbb P^1$. Its Stein factorization is
\[
               H\xrightarrow{q}B\xrightarrow{\nu}\mathbb P^1,
\tag{C.EH.3}\hypertarget{tag-00FS}{}\label{cdp:eq:C.EH.3}
\]
where $B$ is a smooth compact connected curve, $q$ has
connected fibres, and $\nu$ is finite. The degree-one
Leray edge map injects $H^1(B,\mathbb Q)$ into
$H^1(H,\mathbb Q)$, since $q_*\mathbb Q=\mathbb Q$.
Thus $2g(B)\leq1$, and $B\simeq\mathbb P^1$.

\item Every irreducible compact curve $C\subset H$ is smooth
elliptic. Indeed, adjunction and \eqref{cdp:eq:C.EH.2} give
\[
                 2p_a(C)-2=(K_H+C)\cdot C=0.
\tag{C.EH.4}\hypertarget{tag-00FT}{}\label{cdp:eq:C.EH.4}
\]
Its normalization cannot be rational: a holomorphic map
$\mathbb P^1\to H$ lifts through the universal cover
$\mathbb C^2\setminus\{0\}\to H$, and both coordinate
functions of the lift are constant. If $g$ is the genus
of the normalization and $\delta$ is its normalization
defect, \eqref{cdp:eq:C.EH.4} says $g+\delta=1$. Hence $g=1$ and
$\delta=0$. Furthermore two distinct irreducible curves
are disjoint, again by their zero intersection number
and local positivity. Since $q$ has connected fibres,
each full fibre is therefore
\[
                              q^{-1}(b)=m_b C_b,
\tag{C.EH.5}\hypertarget{tag-00FU}{}\label{cdp:eq:C.EH.5}
\]
with $C_b$ a smooth connected elliptic curve. The scheme
fibre is the displayed Cartier divisor: a map from the
smooth surface to a curve has no embedded fibre component.

\item We check the monodromy by an actual finite map, including
at a multiple fibre. Take a small disc with coordinate $t$
at $b$, write $m=m_b$, and normalize the base change
$t=s^m$. If $x=0$ defines $C_b$ locally, then
$t=u x^m$ for a holomorphic unit $u$. The normalized
cover is given locally by
\[
                     w^m=u,\qquad s=wx.
\tag{C.EH.6}\hypertarget{tag-00FV}{}\label{cdp:eq:C.EH.6}
\]
It is a finite unramified cover of the original surface
tube: locally it is the disjoint union of the $m$ choices
of a holomorphic root of $u$. Denote it by $V\to U$.
The deck group $\mu_m$ acts freely, and $V/\mu_m=U$.
The punctured base-changed family has connected fibres
and connected total space. The normalization has no
component supported over zero. Its Stein factor over
$\Delta_s$ is finite and generically degree one and
therefore is the disc itself. In particular the central
fibre $C'$ is connected. It is smooth by \eqref{cdp:eq:C.EH.6}, and
its multiplicity as the fibre of $V\to\Delta_s$ is one.
Thus this map is a submersion on a neighbourhood of $C'$.
The actual central comparison is the finite unramified
quotient
\[
                         C'\longrightarrow C'/\mu_m=C_b.
\tag{C.EH.7}\hypertarget{tag-00FW}{}\label{cdp:eq:C.EH.7}
\]

\item A fixed-point-free finite-order automorphism of an elliptic
curve is a translation. To see this, choose an origin and
write it as $z\mapsto\beta z+a$. If $\beta\ne1$, the
nonzero elliptic endomorphism $1-\beta$ is an isogeny,
so $(1-\beta)z=a$ has a solution. Hence every nonidentity
element of the deck group in \eqref{cdp:eq:C.EH.7} is a translation.
Choose a $\mu_m$-invariant smooth horizontal connection
for the submersion $V\to\Delta_s$, by finite averaging.
Radial parallel transport then gives an equivariant
smooth trivialization. A positive meridian in the
$t$-disc acts on $H^1$ by the corresponding deck
translation or its inverse, according to the transport
convention; both induce the identity. We have proved
\[
                  T_b=I\quad\hbox{on }H^1(q^{-1}(t),\mathbb Z).
\tag{C.EH.8}\hypertarget{tag-00FX}{}\label{cdp:eq:C.EH.8}
\]
The finitely many meridians generate the fundamental
group of the punctured base $\mathbb P^1$, so the entire
integral monodromy representation is trivial.

\item The holomorphic elliptic $j$-function extends across
each omitted value because the actual cover
$V\to\Delta_s$ is a smooth elliptic family there;
its $j$-function is deck invariant and descends to
the $t$-disc. Thus $j$ is holomorphic on $\mathbb P^1$
with values in $\mathbb C$ and is constant. The map
\eqref{cdp:eq:C.EH.1} has exactly the fixed-modulus, identity-monodromy,
smooth-multiple-fibre structure required in \ref{cdp:st:C.2.UN.34}.
Apply that theorem to the given surface $H$ and this
specific map $q$. Its root-curve torsor is the one
defined by \eqref{cdp:eq:C.EH.6}--\eqref{cdp:eq:C.EH.7}, and its complete analytic
torsor class is retained by the two-line-bundle lift.
Consequently the distinguished component in that
construction is isomorphic to $H$, with multiplicity
one, rather than merely having its topological or
birational invariants. The proper compact map, exact
punctured product and Euler-zero simple-normal-crossing
support are those constructed in (C.SF.8)--(C.SF.13).

\end{enumerate}
\end{proof}






\begin{proposition}[Every diagonal primary Hopf surface occurs unchanged in a compact torus filling]\hypertarget{tag-01DL}{}\label{cdp:st:C.2.UN.41}
\StatementRecord{cdp:st:C.2.UN.41}{proved}{01DL}{\SeriesRef{cdp:eq:C.DH.5}, \SeriesRef{cdp:eq:C.DH.6}, \SeriesRef{cdp:eq:C.DH.7}, \SeriesRef{cdp:eq:C.DH.13}, \SeriesRef{cdp:eq:C.DH.14b}, \SeriesRef{cdp:eq:C.DH.12}, \SeriesRef{cdp:st:C.2.UN.40}; \cite[\S7.6, pp.~33--34]{SchuttShioda2010}; \cite[Theorem~2.1, p.~2]{Papikian2004Tate}}

For every pair $\alpha,\beta\in\C$ with
$0<|\alpha|,|\beta|<1$, the diagonal Hopf surface
\[
 H_{\alpha,\beta}=(\C^2\setminus\{0\})/
       \langle(z,w)\mapsto(\alpha z,\beta w)\rangle
                                                               \tag{C.DH.1}\hypertarget{tag-00F6}{}\label{cdp:eq:C.DH.1}
\]
occurs as an unchanged smooth component of a
reduced normal-crossing fibre of a proper map
from a smooth compact complex threefold to
$\PP^1$, with general fibre a two-torus.
The distinguished fibre consists of two copies
of \eqref{cdp:eq:C.DH.1}, joined along both axis elliptic
curves by inversion. Every reduced fibre
support in this compact construction has
Euler number zero. This includes all
multiplicatively independent pairs, whose
Hopf surfaces have algebraic dimension zero.

The distinguished support $S$ has torsion-free
integral cohomology of ranks
\[
                  (1,2,3,4,2),\qquad
                  h^\bullet(S,\OO_S)=(1,2,1).        \tag{C.DH.2}\hypertarget{tag-00FG}{}\label{cdp:eq:C.DH.2}
\]
Each normalization component has integral
cohomology of ranks $(1,1,0,1,1)$ and
$h^\bullet(\OO)=(1,1,0)$.
Locally the nearby torus has the exact
multiplicative period presentation
\[
 A_t=(\C^*)^2/
        \langle(\alpha,\beta),(t,t^{-1})\rangle
                    \qquad(0<|t|\ll1).             \tag{C.DH.3}\hypertarget{tag-00FH}{}\label{cdp:eq:C.DH.3}
\]
In the compact construction the constants
$\alpha,\beta$ in this formula vary
holomorphically with unchanged limiting
values, and $t$ is replaced by a holomorphic
parameter with a simple zero. The explicit
comparison below retains the torus fibres.
\end{proposition}



\begin{proof}\leavevmode\par

\begin{enumerate}[label={$\langle1\rangle$\arabic*.},leftmargin=0pt,labelsep=.5em]
\item  We first construct the proper local filling
and its actual torus comparison, then glue
it into a compact suspension.


\item \par\medskip\noindent\textbf{Exact neighborhoods and the two plumbing maps.}
The two axes of $H=H_{\alpha,\beta}$ are
$C_\alpha=\C^*/\alpha^{\Z}$ and
$C_\beta=\C^*/\beta^{\Z}$. The open
neighborhood $z\ne0$ of the first is
the entire normal line bundle
\[
 L_\alpha=(\C^*\times\C)/
         ((z,x)\sim(\alpha z,\beta x));
                                                               \tag{C.DH.4}\hypertarget{tag-00FI}{}\label{cdp:eq:C.DH.4}
\]
the second has the analogous presentation
with $\alpha,\beta$ interchanged.
Inversion on $C_\alpha$ pulls this
normal line to its inverse. Explicitly
the rank-two plumbing bundle over that
elliptic curve is
\[
 \mathcal N_\alpha=(\C^*\times\C^2)/
       ((z,x,y)\sim(\alpha z,\beta x,\beta^{-1}y)),
                         \qquad t=xy.               \tag{C.DH.5}\hypertarget{tag-00FJ}{}\label{cdp:eq:C.DH.5}
\]
Its branch $y=0$ maps to the first
Hopf copy by $(z,x)$, and its branch
$x=0$ maps to the second by
$(Z,W)=(z^{-1},y)$.
For the other axis use
\[
 \mathcal N_\beta=(\C^*\times\C^2)/
       ((w,x,y)\sim(\beta w,\alpha x,\alpha^{-1}y)),
                         \qquad t=xy,                \tag{C.DH.6}\hypertarget{tag-00FK}{}\label{cdp:eq:C.DH.6}
\]
whose branches map to the first
copy by $(z,w)=(x,w)$ and to the
second by $(Z,W)=(y,w^{-1})$.
These formulas prove the two dual-normal
identifications, including their
base and fibre transition functions.


\begin{enumerate}[label={$\langle2\rangle$\arabic*.},leftmargin=1.3em,labelsep=.5em]
\item  Put $A=-\log|\alpha|>0$ and
$B=-\log|\beta|>0$. On the common
open torus of a Hopf copy define
\[
 h=B\log|z|-A\log|w|,\qquad
 \nu_\alpha=|w||z|^{-B/A}=e^{-h/A},\qquad
 \nu_\beta=|z||w|^{-A/B}=e^{h/B}.        \tag{C.DH.7}\hypertarget{tag-00FL}{}\label{cdp:eq:C.DH.7}
\]
All three expressions are invariant
under the Hopf deck generator.
The two $\nu$ are the norms of
flat Hermitian metrics on the
normal lines. Neighborhoods
$\nu_\alpha<R$ and $\nu_\beta<R$
are disjoint for $R<1$.
Fix $0<r<R<1$ and $0<\delta<r^2$.
For each copy take the product
of the open compact core
$U_i=\{\nu_\alpha>r,\nu_\beta>r\}$
with $\Delta_\delta$.
In \eqref{cdp:eq:C.DH.5} take
\[
 \|x\|=|x||z|^{-B/A}<R,
 \qquad \|y\|=|y||z|^{B/A}<R,
 \qquad |xy|<\delta;                                \tag{C.DH.8}\hypertarget{tag-00FM}{}\label{cdp:eq:C.DH.8}
\]
take the analogous open set in
\eqref{cdp:eq:C.DH.6}. Glue it to the first
core where $\|x\|>r$, using the
first branch coordinate and
$t=xy$, and to the second core
where $\|y\|>r$, using the second
branch coordinate. The product
of the two norms is exactly
$|t|$. Thus the two core
overlaps in any one neck are
disjoint, since $|t|<r^2$.
The coordinate maps are
biholomorphisms on the stated
overlaps; their inverses insert
the coordinate $t/x$ or $t/y$.


\item  This finite gluing is Hausdorff.
Indeed at the boundary $\|x\|=r$
of a core overlap the core
point is absent, and at its
other boundary $\|x\|=R$ the
neck point is absent. Thus
the graph of each identification
is closed relative to the
two pieces. The same holds
for $y$; distinct necks have
disjoint overlaps in a given
core, and the preceding
$r^2$ inequality excludes
an additional identification
through the opposite core.
This also proves separation
for the resulting finite
equivalence relation. The
quotient is a smooth complex
threefold $\mathcal Y$ with
the holomorphic map $t$.
At the two double curves it
has the exact coordinates
$xy=t$; at zero its fibre
is the two unchanged Hopf
surfaces with the asserted
two inversion identifications.


\item  The map is proper after the
indicated restriction to the
disc. Choose
$r<r'<R'<R$.
Over a closed subdisc,
the two compact sets
$\{\nu_\alpha\ge r',\nu_\beta\ge r'\}$
in the Hopf copies, together
with the two closed neck sets
$\|x\|,\|y\|\le R'$, cover
the entire inverse image.
For a core point missing its
compact set, the small normal
coordinate is less than $r'$
and the other neck coordinate
has norm at most $\delta/r<r$.
For a neck point with one
norm greater than $R'$,
the corresponding core point
belongs to its compact set.
All these closed sets are
compact: their bases are
compact elliptic curves and
their vector-bundle coordinates
are bounded. They cover the
closed inverse image, proving
properness. This argument
also explicitly verifies
compactness of the fibres.


\end{enumerate}

\item \par\medskip\noindent\textbf{The exact nearby torus.}
For $t\ne0$, the punctured
identifications are
\[
 \phi_\alpha(z,w)=(z^{-1},t/w),\qquad
 \phi_\beta(z,w)=(t/z,w^{-1}).                       \tag{C.DH.9}\hypertarget{tag-00FN}{}\label{cdp:eq:C.DH.9}
\]
Each conjugates the Hopf deck
generator to its inverse.
Their difference is the
additional period
$(z,w)\mapsto(tz,t^{-1}w)$.
The four additive period
vectors of \eqref{cdp:eq:C.DH.3} are
\[
 (2\pi i,0),\ (0,2\pi i),\
 (\log\alpha,\log\beta),\
                         (\log t,-\log t).          \tag{C.DH.10}\hypertarget{tag-00F7}{}\label{cdp:eq:C.DH.10}
\]
They are real-linearly
independent: after quotienting
the two imaginary directions,
the real determinant is
\[
       -\log|t|\bigl(\log|\alpha|+\log|\beta|\bigr)
                                      \ne0.          \tag{C.DH.11}\hypertarget{tag-00F8}{}\label{cdp:eq:C.DH.11}
\]
Their integral span is
therefore a discrete
cocompact lattice in $\C^2$.
This proves directly that
$A_t$ is a compact complex
two-torus. A turn of $t$
adds $(2\pi i,-2\pi i)$
to the last vector, which
is already in the lattice,
so these quotients form
the stated family over
the punctured disc.


\begin{enumerate}[label={$\langle2\rangle$\arabic*.},leftmargin=1.3em,labelsep=.5em]
\item  There is an actual
holomorphic comparison
$\Psi_t:\mathcal Y_t\to A_t$.
On the first core use
$[z,w]$, and on the
second use $[Z^{-1},t/W]$.
On the first neck use
$[z,x]$; on the second
use $[x,w]$. On its
second-core overlap
the latter gives instead
$[x/t,tw]$, differing by
the inverse of the period
$(t,t^{-1})$. Thus the
maps agree exactly on
all the overlaps. They
are local biholomorphisms.
The proper source makes
$\Psi_t$ a finite
unramified covering
onto the connected
torus $A_t$.


\item  Its degree is one.
For completeness the
height in \eqref{cdp:eq:C.DH.7}
gives a direct check.
Put $L=-\log|t|$ and
$c_r=-\log r$, $c_R=-\log R$.
The first core has
height interval
$(-Bc_r,Ac_r)$.
Under the chosen comparison
the second core has interval
$(AL-Ac_r,AL+Bc_r)$.
The first neck fills
$(Ac_R,AL-Ac_R)$,
and the other fills,
modulo the period
$(A+B)L$,
$(AL+Bc_R,(A+B)L-Bc_R)$.
These intervals cover
one height circle;
their consecutive
overlaps are exactly
the gluing annuli.
Since $L>2c_r$,
height zero has one
preimage in the first
core and none in the
other three pieces,
apart from the already
identified overlaps.
On that core the only
remaining identification
is the original Hopf
deck generator. Thus
the covering degree
is one, proving the
actual biholomorphism
in \eqref{cdp:eq:C.DH.3}.


\item  The entire construction
also works for holomorphic
$a(t),b(t)$ with limiting
values $\alpha,\beta$
and for a holomorphic
smoothing parameter $Q(t)$
with a simple zero.
Use the smooth family
of Hopf quotients with
deck multipliers $a(t),b(t)$
and put $xy=Q(t)$ in
the two necks. Their
quotient and gluing maps
are the same holomorphic
formulas. The norms and
height vary continuously,
and the preceding bounds
hold uniformly on a
smaller disc. Since
$Q'(0)\ne0$, the total
space is still smooth.
The comparison is now
\[
 \mathcal Y_t\simeq
 (\C^*)^2/
       \langle(a(t),b(t)),(Q(t),Q(t)^{-1})\rangle.
                                                               \tag{C.DH.12}\hypertarget{tag-00F9}{}\label{cdp:eq:C.DH.12}
\]


\end{enumerate}

\item \par\medskip\noindent\textbf{A compact suspension with these exact punctured fibres.}
We construct a compact elliptic
surface $J\to\PP^1$ with a
section $P$ whose value on a
nodal fibre is the specified
$\alpha$. Start with the
nodal plane cubic
$y^2z=x^2(x+z)$, with origin
$O=[0:1:0]$. On its
normalization put
$v=y/x$; its smooth group
coordinate is
$\zeta=(v-1)/(v+1)$,
with $O$ corresponding to
$\zeta=1$. Choose the
smooth point $P_0$ with
$\zeta=\alpha$.
Take a second cubic
through $O,P_0$, otherwise
general. Its intersection
with the nodal cubic
consists of nine distinct
smooth points, with
transverse intersections,
and misses the node.
These requirements are
nonempty open conditions
on its restriction to
the cubic and its first
jets; the restriction
map of cubic polynomials
is surjective, so a
general such choice
satisfies them.
Blow up the nine
base points of the
pencil. The resulting
smooth rational surface
$J$ has a connected
elliptic pencil, zero
section from $O$, and
section $P$ from $P_0$.
Its central member is
the original nodal
cubic, a fibre of type
$I_1$. At the node the
second cubic is nonzero,
so the pencil parameter
smooths the node to
first order. The full
fibre class is
$F=3H-\sum E_i=-K_J$.
Consequently there is
no vertical rational
$(-1)$-curve, by
adjunction, and this
is the smooth relatively
minimal elliptic model.


\begin{enumerate}[label={$\langle2\rangle$\arabic*.},leftmargin=1.3em,labelsep=.5em]
\item  Translation by $P$ is an
automorphism $T_P$ of
this entire surface.
Here is the required
extension argument.
The group law first
defines a birational
map over the base,
with inverse translation
by $-P$. Resolve its
graph minimally,
$J\xleftarrow{p}W\xrightarrow{q}J$.
Every exceptional curve
is vertical. If an
exceptional $(-1)$-curve
$D$ of $p$ is not
exceptional for $q$,
then $q^*K_J\cdot D=0$
because $K_J=-F$.
The discrepancy formula
$K_W=q^*K_J+\sum a_iD_i$
has $a_i>0$ and
$D_i\cdot D\ge0$,
contradicting $K_W\cdot D=-1$.
Thus $D$ is exceptional
for both maps and can
be contracted in their
common graph, contrary
to its minimality.
There are no exceptional
curves, so the translation
and its inverse are
holomorphic automorphisms.
This agrees with the explicit
section-translation extension
in \cite[\S7.6, pp.~33--34]{SchuttShioda2010}.


\item  Put $q_0=\alpha\beta$.
The quotient
\[
 Y^{\rm std}=(\C^*_u\times J)/
        \langle(u,y)\mapsto(q_0u,T_P(y))\rangle
                                                               \tag{C.DH.13}\hypertarget{tag-00FA}{}\label{cdp:eq:C.DH.13}
\]
is a smooth compact
threefold. Projection to
$E_{q_0}=\C^*/q_0^{\Z}$
is a locally trivial
holomorphic bundle with
compact fibre $J$:
local choices of a
logarithm give its
trivializations, and
their transition is
the indicated integral
power of $T_P$.
The deck action is
free and properly
discontinuous because
its first coordinate
is multiplication by
$q_0$, with $|q_0|<1$.
It preserves the
elliptic-pencil parameter,
so \eqref{cdp:eq:C.DH.13} maps
properly to $\PP^1$.


\item  On a small punctured
disc at its $I_1$ value
there is the actual
analytic Tate presentation
\[
 J_t\simeq\C^*/Q(t)^{\Z},\qquad
 Q(t)=t\,u(t),\quad u(0)\ne0,
 \qquad P(t)=[a(t)],\quad a(0)=\alpha.
                                                               \tag{C.DH.14}\hypertarget{tag-00FB}{}\label{cdp:eq:C.DH.14}
\]
The presentation can be
obtained by normalizing
the nonvanishing cycle
period to one. Its
other period has the
form $(2\pi i)^{-1}\log t$
plus a holomorphic
function; exponentiation
gives $Q(t)$ with a
simple zero. The
standard nodal Tate
chart identifies its
smooth central group
with $\C^*$, so the
section through $P_0$
has the holomorphic
nonzero lift $a(t)$
with that limiting
value. This is an
isomorphism of the
elliptic groups and
their smooth central
charts, not only an
identity of their
$j$-invariants.
For clarity, the central
chart and its sign can
also be checked directly.
The complex uniformization in \cite[Theorem~2.1 and equation~(2.1), p.~2]{Papikian2004Tate} gives
\[
 \begin{gathered}
 y^2+xy=x^3+a_4(q)x+a_6(q),\\
 a_4(q)=-5q+O(q^2),\quad
 a_6(q)=-q+O(q^2),\quad
 \Delta(q)=q+O(q^2).
 \end{gathered}                                 \tag{C.DH.14a}\hypertarget{tag-00FC}{}\label{cdp:eq:C.DH.14a}
\]
Here $a_4=-5\sum_{n\ge1}n^3q^n/(1-q^n)$ and
$a_6=-\sum_{n\ge1}(5n^3+7n^5)q^n/(12(1-q^n))$.
The multiplicative
uniformization has the
central values
\[
 x=\frac{u}{(1-u)^2},\qquad
 y=\frac{u^2}{(1-u)^3},\qquad
 X=4x,\quad Y=8y+4x,\qquad
 \frac{Y/X-1}{Y/X+1}=u.                          \tag{C.DH.14b}\hypertarget{tag-00FD}{}\label{cdp:eq:C.DH.14b}
\]
Its equation becomes
$Y^2=X^2(X+1)$, precisely
the chosen nodal cubic.
The series converge
normally on compact
subsets of the smooth
central chart. Invert
the holomorphic series
$1/j(q)=q+O(q^2)$ at
the given simple zero
of $1/j(t)$ to obtain
$Q(t)$. After completing
the square and cube,
the two short Weierstrass
coefficients are units
at zero. Their equal
$j$-invariants identify
them by a holomorphic
unit fourth root; choose
that root at zero to
match the sixth-weight
coefficient. The identity
of the squared ratios
then matches that
coefficient throughout
the disc. This supplies
the actual Weierstrass
coordinate isomorphism.
Choose its central value
as in \eqref{cdp:eq:C.DH.14b}.
The inverse function
theorem at the smooth
point $u=\alpha$ gives
the holomorphic lift
$a(t)$ with exactly
$a(0)=\alpha$.


\item  Let $b(t)=q_0/a(t)$.
It tends to $\beta$,
and both multipliers
have modulus less than
one on a sufficiently
small disc. The actual
coordinate change
\[
                        z=y,\qquad w=u/y
                                                               \tag{C.DH.15}\hypertarget{tag-00FE}{}\label{cdp:eq:C.DH.15}
\]
on the multiplicative
cover of \eqref{cdp:eq:C.DH.13}
turns its two deck
generators into
$(a(t),b(t))$ and
$(Q(t),Q(t)^{-1})$.
Thus its punctured
family is exactly
\eqref{cdp:eq:C.DH.12}. Replace
the local filling
of \eqref{cdp:eq:C.DH.13} by
the proper plumbing
just constructed,
using that explicit
isomorphism on the
punctured disc.
This is gluing over
the open cover
$\PP^1\setminus\{0\}$
and a disc at zero.
The total spaces
are identified by
biholomorphisms over
their common base,
so the result is
Hausdorff and smooth;
properness is local
on the base and
holds on each member
of this cover.
The compact base
therefore gives a
compact total space.
The distinguished
central components
are exactly the
prescribed Hopf
surfaces. General
fibres are the
actual tori of
\eqref{cdp:eq:C.DH.12}, or,
elsewhere, the same
translation suspensions
of smooth elliptic
fibres of $J$.
The latter are tori
by their additive
four-period lattice,
or by the direct
constant-translation
bundle presentation
in \eqref{cdp:eq:C.DH.13}.


\item  The support over zero
has Euler number
$2e(H)-2e(C)=0$.
At every other value
the support is a
locally trivial bundle
over $E_{q_0}$ with
fibre the corresponding
reduced elliptic-pencil
support; its monodromy
is $T_P$. Its Euler
number is the product
with $e(E_{q_0})=0$.
Thus every reduced
support has Euler
number zero.


\end{enumerate}

\item \par\medskip\noindent\textbf{The support cohomology and its maps.}
For the two-component
support there is the
actual normalization
Mayer--Vietoris sequence
with middle map the
difference of the two
restrictions to
$C_\alpha\sqcup C_\beta$.
In degree zero its
matrix is
$\left(\begin{smallmatrix}1&-1\\1&-1\end{smallmatrix}\right)$,
of rank one and
primitive image.
Choose on each axis
the circle of its
Hopf deck period and
its angular circle.
The integral $H^1$
generator of each
Hopf surface restricts
as $(1,0)$ on each
axis. The inversion
gluing reverses these
two generators on
the second side.
Consequently the
degree-one restriction
matrix is
\[
                     \begin{pmatrix}
                       1&1\\0&0\\1&1\\0&0
                     \end{pmatrix},                \tag{C.DH.16}\hypertarget{tag-00FF}{}\label{cdp:eq:C.DH.16}
\]
again of rank one
with primitive image.
The Hopf components
have $H^2=0$, and
the two elliptic
curves have $H^2=\Z^2$
and no higher
cohomology. Exactness
therefore gives
\eqref{cdp:eq:C.DH.2}, including
torsion-freeness.
For coherent cohomology,
$H^1(H,\OO_H)=\C$
is represented by
the constant deck
cocycle one. Its
restriction to each
axis is the nonzero
constant cocycle on
that elliptic curve;
inversion changes
its sign. The
degree-one coherent
restriction matrix
has both rows $(1,1)$.
Together with
$h^\bullet(\OO_H)=(1,1,0)$
and the degree-zero
matrix already given,
the coherent
normalization sequence
gives $(1,2,1)$.
All maps in these
calculations are
the stated axis
restrictions and
inversion pullbacks.
The curve calculation in \ref{cdp:st:C.2.UN.40} leaves only the two axes when
the multipliers are multiplicatively independent. A nonconstant meromorphic
function would have infinitely many distinct nonempty generic level divisors,
which is impossible when the only compact curves are those two axes.
Consequently the Hopf component in this case has algebraic dimension zero.



\end{enumerate}
\end{proof}






\begin{proposition}[Every nonlinear primary Hopf surface occurs unchanged over a rational base]\hypertarget{tag-01DM}{}\label{cdp:st:C.2.UN.42}
\StatementRecord{cdp:st:C.2.UN.42}{proved}{01DM}{\SeriesRef{cdp:eq:C.NH.6}, \SeriesRef{cdp:eq:C.NH.5}, \SeriesRef{cdp:eq:C.NH.4}, \SeriesRef{cdp:eq:C.NH.10}, \SeriesRef{cdp:eq:C.NH.12}, \SeriesRef{cdp:eq:C.NH.11}, \SeriesRef{cdp:eq:C.NH.8}, \SeriesRef{cdp:st:C.2.UN.40}, \SeriesRef{cdp:eq:C.NH.14}, \SeriesRef{cdp:eq:C.NH.13}}

For every integer $r\ge1$ and every $\beta$ with $0<|\beta|<1$,
the primary Hopf surface
\[
 H_r(\beta)=(\C^2\setminus\{0\})/
   \langle(z,w)\mapsto(\beta^r z+w^r,\beta w)\rangle
                                                        \tag{C.NH.1}\hypertarget{tag-00S9}{}\label{cdp:eq:C.NH.1}
\]
occurs unchanged as a smooth multiplicity-one component of a
reduced SNC fibre of a proper map $f:Y\to\PP^1$, where $Y$ is
a smooth compact complex threefold and the general fibre is a
complex two-torus. Every reduced fibre support in this
construction has Euler number zero.

The distinguished full fibre is $S=H_r(\beta)+P_r$, with one
smooth elliptic intersection $J$. The second component is a
smooth elliptic surface over an elliptic curve, with exactly
one multiple fibre $rJ$ when $r>1$, and
\[
 \begin{gathered}
 K_{P_r}=\OO_{P_r}((r-1)J),\qquad
 b(P_r)=(1,3,4,3,1),\qquad h^\bullet(\OO_{P_r})=(1,2,1),\\
 H^\bullet(S,\Z)\text{ is torsion-free of ranks }(1,3,5,5,2),
 \qquad h^\bullet(S,\OO_S)=(1,2,1).
 \end{gathered}                                         \tag{C.NH.2}\hypertarget{tag-00SJ}{}\label{cdp:eq:C.NH.2}
\]
The Hopf component has integral ranks $(1,1,0,1,1)$ and
coherent dimensions $(1,1,0)$. The construction works for
every prescribed $\beta$, without a smallness condition
beyond $|\beta|<1$.
\end{proposition}



\begin{proof}

\begin{enumerate}[label={$\langle1\rangle$\arabic*.},leftmargin=*]
\item \par\medskip\noindent\textbf{A compact base with the required translation jet.}
Let $B$ be the elliptic cubic $v^2=u^3-u+1$, with origin the
point at infinity. Its cubic has distinct nonzero roots. Let
$C$ be the smooth projective completion of
\[
 v^2=t^{6r}-t^{2r}+1,
 \qquad \psi:C\longrightarrow B,\quad
          (t,v)\longmapsto(t^{2r},v).
                                                        \tag{C.NH.3}\hypertarget{tag-00SK}{}\label{cdp:eq:C.NH.3}
\]
The polynomial has distinct roots: a repeated root would be
nonzero and would give a common root of $u^3-u+1$ and its
derivative. At infinity the two charts
$s=t^{-1}$, $v s^{3r}=\pm\sqrt{1-s^{4r}+s^{6r}}$ give two
distinct smooth points $c_+,c_-$. Thus $C$ is a smooth compact
curve of genus $3r-1$. The involution $h_C(t,v)=(t,-v)$ has
quotient $\PP^1_t$, interchanges $c_+,c_-$, and satisfies
$\psi h_C=-\psi$. Its fixed points are the $6r$ finite roots
of the displayed polynomial.

\item The elliptic uniformizer $u/v$ has order $r$ at both
$c_\pm$. Choose a local additive coordinate $p$ on $B$ at
zero. At each of the two points a holomorphic coordinate
$s$ on $C$ can therefore be chosen so that
\[
 2\psi(s)=s^r\quad\hbox{in the additive coordinate on }B.
                                                        \tag{C.NH.4}\hypertarget{tag-00SL}{}\label{cdp:eq:C.NH.4}
\]
Taking the $r$-th root of its nonzero holomorphic unit proves
this coordinate equality. On the compact surface $B\times C$
consider the commuting automorphisms
\[
 A(p,c)=(p+2\psi(c),c),\qquad
 A_{1/2}(p,c)=(p+\psi(c),c),\qquad
 h(p,c)=(-p,h_Cc).
                                                        \tag{C.NH.5}\hypertarget{tag-00SM}{}\label{cdp:eq:C.NH.5}
\]
Here $A_{1/2}^2=A$, $h^2=1$, and the equality
$\psi h_C=-\psi$ proves that $h$ commutes with both
translations. All three preserve the pair of points
$(0,c_+),(0,c_-)$.

\par\medskip\noindent\textbf{The exact weighted blowup and its smooth nonzero line.}
Perform the weighted blowup of weights $(r,1)$ at both
points, retaining its smooth orbifold charts. Denote the
resulting compact orbifold surface by $\mathcal V$, its
exceptional divisors by $R_+,R_-$, and put
$L=\OO_{\mathcal V}(R_++R_-)$. Locally its weighted blowup
has the explicit presentation
\[
 \left[((\C^2_{z,w}\setminus\{0\})\times\C_x)/\C^*\right],
 \quad \lambda(z,w,x)=(\lambda^r z,\lambda w,\lambda^{-1}x),
 \quad p=x^r z,\quad s=xw.
                                                        \tag{C.NH.6}\hypertarget{tag-00SN}{}\label{cdp:eq:C.NH.6}
\]
The exceptional divisor is $x=0$. The line $L$ has fibre
coordinate $\ell$ of weight $-1$ and canonical section
$\ell=x$. The quotient chart $w\ne0$ is the smooth chart
$p=s^r(z/w^r)$; the chart $z\ne0$ has the residual finite
action
\[
 (w,x,\ell)\longmapsto
      (\zeta w,\zeta^{-1}x,\zeta^{-1}\ell),\qquad \zeta^r=1.
                                                        \tag{C.NH.7}\hypertarget{tag-00SO}{}\label{cdp:eq:C.NH.7}
\]
It is free when $\ell\ne0$. More explicitly, on $L^\times$
use $\lambda=\ell$ in \eqref{cdp:eq:C.NH.6} to set $\ell=1$. This gives
the actual smooth coordinates
$(z,w,x)\in(\C^2\setminus\{0\})\times\C$, after replacing
the coordinates by $(\ell^rz,\ell w,\ell^{-1}x)$.
Thus the coarse nonzero line space $L^\times$ is a smooth
complex threefold, including the orbifold points.

\item All three maps in \eqref{cdp:eq:C.NH.5} lift canonically. In the
coordinates \eqref{cdp:eq:C.NH.4}, $A$ is $p\mapsto p+s^r$ and its lift
in \eqref{cdp:eq:C.NH.6} is $(z,w,x)\mapsto(z+w^r,w,x)$.
The weighted ideals are therefore preserved. The half
translation replaces the coefficient by $1/2$; $h$
interchanges the two centres and preserves their weighted
orders. The line lifts preserve its canonical divisor
section, so their group laws hold exactly, including
$\widetilde A_{1/2}^{\,2}=\widetilde A$ and
$\widetilde h^{\,2}=1$.

\item Multiply the line fibre by $\beta$ and compose with
$\widetilde A$. On the smooth coordinates just constructed
the resulting automorphism is
\[
 g(z,w,x)=
       (\beta^r(z+w^r),\beta w,\beta^{-1}x),
 \qquad s=xw.
                                                        \tag{C.NH.8}\hypertarget{tag-00SP}{}\label{cdp:eq:C.NH.8}
\]
It preserves the map to $C$. On $x=0$, the change
$Z=\beta^{-r}z$ turns its first coordinate into
$\beta^r Z+w^r$, proving the exact identification with
\eqref{cdp:eq:C.NH.1}.

\par\medskip\noindent\textbf{Uniform properness for every contraction.}
We verify the global quotient, rather than inferring its
compactness from a local contraction. Choose a Hermitian
orbifold metric on $L$. Write $\sigma$ for its canonical
section. Near either centre, its norm off the exceptional
divisor is comparable, with constants independent of the
point, to
\[
       \rho(p,s)=\bigl(d_B(p,0)^2+|s|^{2r}\bigr)^{1/(2r)}.
                                                        \tag{C.NH.9}\hypertarget{tag-00SQ}{}\label{cdp:eq:C.NH.9}
\]
Here $d_B$ is a fixed flat elliptic distance. To see the
comparison, normalize the weighted coordinates so that
$|z|^2+|w|^{2r}=1$. Then
$(|p|^2+|s|^{2r})^{1/(2r)}=|x|$, which is the norm of the
canonical section up to a smooth positive bounded factor
on the compact weighted exceptional curve. Away from the
two centres all these norms are bounded above and below.
More precisely this comparison holds on $B$ times a small
$c_\pm$-disc: outside the chosen $p$-coordinate neighborhood,
both $\rho$ and $\|\sigma\|$ are bounded above and away from
zero. It therefore also covers an iterate $A^n$ that moves
$p$ out of that local coordinate chart.

\item For every integer $n$, the triangle inequality on the flat
elliptic curve gives
\[
 d_B(p+2n\psi(s),0)
       \le d_B(p,0)+C_0|n||s|^r.
\]
The same inequality applied to the inverse translation
gives the reverse comparison. Since the canonical lift
sends $\sigma(p,s)$ to $\sigma(p+2n\psi(s),s)$, there is
a global constant $C_1\ge1$ such that
\[
 C_1^{-1}(1+|n|)^{-1/r}\|\ell\|
 \le \|\widetilde A^{\,n}\ell\|
 \le C_1(1+|n|)^{1/r}\|\ell\|.
                                                        \tag{C.NH.10}\hypertarget{tag-00SA}{}\label{cdp:eq:C.NH.10}
\]
The comparisons first hold off the exceptional divisors.
The ratio of line norms extends continuously and positively
on the entire orbifold line, so they hold there as well.
On the compact part of $C$ away from both centres the
canonical section is nowhere zero and supplies the same
bounds with a fixed constant. This proves the uniform
global assertion.

\item Consequently $\|g^n\ell\|$ tends uniformly to zero on the
unit circle bundle as $n\to+\infty$, and to infinity as
$n\to-\infty$: multiply \eqref{cdp:eq:C.NH.10} by $|\beta|^n$.
The action is properly discontinuous because any two
compact subsets have line norms bounded above and away
from zero. It is free, since a fixed point of a nonzero
power would also be fixed by arbitrarily large positive
powers, contradicting the norm limit.

\item Choose $N$ such that $g^N$ decreases every norm by a factor
strictly less than one. Its minimum norm ratio is a positive
number $c_N$. Every $g^{N\Z}$-orbit meets the compact annulus
$c_N\le\|\ell\|\le1$: take the first iterate whose norm is
at most one. This annulus is compact over the compact
orbifold base. It follows that
\[
                 \widetilde Y=L^\times/\langle g\rangle
                                                        \tag{C.NH.11}\hypertarget{tag-00SB}{}\label{cdp:eq:C.NH.11}
\]
is a smooth compact threefold, with a proper map to $C$.
This proof applies to every $0<|\beta|<1$.

\item Off the two exceptional fibres, the canonical section
trivializes $L$ and gives the actual quotient
\[
 (B\times\C^*)/
   \langle(p,\ell)\mapsto(p+2\psi(c),\beta\ell)\rangle.
                                                        \tag{C.NH.12}\hypertarget{tag-00SC}{}\label{cdp:eq:C.NH.12}
\]
If $B=\C/\Lambda_B$, its additive covering lattice is
generated by $(\lambda,0)$ for $\lambda\in\Lambda_B$,
$(0,2\pi i)$, and $(a_c,\log\beta)$, where $a_c$ is any
lift of $2\psi(c)$. These four real generators are
independent because $\log|\beta|\ne0$. Changing the lift
of $a_c$ changes a generator by $\Lambda_B$. This proves
that \eqref{cdp:eq:C.NH.12} is precisely a compact two-torus.

\item At either distinguished value the exact function is
$s=xw$. Its full fibre is reduced SNC, with components
$H:x=0$ and $P:w=0$, meeting along
\[
          J=(\C^*_z/\beta^{r\Z}),\qquad x=w=0.
                                                        \tag{C.NH.13}\hypertarget{tag-00SD}{}\label{cdp:eq:C.NH.13}
\]
Both components and their intersection are compact since
they are closed in \eqref{cdp:eq:C.NH.11}. On $P$ the map to the
unchanged elliptic base $B$ has the exact local form
\[
  p=x^r z,\qquad (z,x)\sim(\beta^r z,\beta^{-1}x).
                                                        \tag{C.NH.14}\hypertarget{tag-00SE}{}\label{cdp:eq:C.NH.14}
\]
Away from $p=0$ its canonical frame is the product
$B\setminus\{0\}\times E_\beta$; on the overlap the
frame coordinate is $\ell=x^{-1}$ and $z=p\ell^r$.
These equalities identify the surface with the one-fibre
logarithmic model in (C.ER.10)--(C.ER.16), with $m=r$.
They also give the same model for $r=1$ directly.
In particular every fibre support of $P\to B$ is a
smooth elliptic curve, and the special full fibre is
$rJ$. Stratified Euler additivity gives $e(P)=0$.
The Hopf surface and $J$ also have Euler number zero,
so the entire distinguished support has Euler number zero.
The local holomorphic canonical evaluation form is
$x^r\,dz\wedge dw\wedge dx$. The Jacobian of \eqref{cdp:eq:C.NH.8}
is $\beta^r$, and the multiplier of $x^r$ is $\beta^{-r}$,
so this form is invariant. Its holomorphic vanishing orders
along $(H,P)$ are exactly $(r,0)$. This realizes the
strict canonical-order difference in \ref{cdp:st:C.2.UN.40} without
replacing a holomorphic generator by a meromorphic
full-fibre shift.

\par\medskip\noindent\textbf{The free quotient giving the rational base.}
Fix either square root $\beta^{1/2}$. The commuting
canonical lifts define
\[
           G=\beta^{1/2}\widetilde A_{1/2}\widetilde h,
                    \qquad G^2=g.
                                                        \tag{C.NH.15}\hypertarget{tag-00SF}{}\label{cdp:eq:C.NH.15}
\]
Thus $G$ induces an involution on $\widetilde Y$. A fixed
point there would have $h_Cc=c$. Such a point is not one
of the two centres. Moreover $\psi(c)$ is a two-torsion
point of $B$, so $2\psi(c)=0$. In the canonical line frame,
\[
 G(p,c,\ell)=(-p-\psi(c),c,\beta^{1/2}\ell),
 \qquad g^n(p,c,\ell)=(p,c,\beta^n\ell).
                                                        \tag{C.NH.16}\hypertarget{tag-00SG}{}\label{cdp:eq:C.NH.16}
\]
The equality $G(p,c,\ell)=g^n(p,c,\ell)$ is impossible by
the absolute values of the line coordinates: it would
give $|\beta|^{1/2}=|\beta|^n$. The involution is therefore
free, and
\[
                 Y=\widetilde Y/\langle G\rangle
                  \longrightarrow C/\langle h_C\rangle=\PP^1
                                                        \tag{C.NH.17}\hypertarget{tag-00SH}{}\label{cdp:eq:C.NH.17}
\]
is a smooth compact threefold with the required map.
The two Hopf components over $c_+,c_-$ are exchanged.
Restriction of the quotient to either one is an
isomorphism onto its image, so the prescribed Hopf
surface is unchanged. The entire distinguished fibre is
isomorphic to either $H+P$ above, and remains reduced.
At a fixed value of $h_C$ the full fibre is twice a free
quotient of a two-torus; its reduced support has Euler
number zero. At every other ordinary value it is a
two-torus. All fibres are connected, as is seen directly
from these quotient descriptions and from $H\cap P=J$.
This proves all the stated occurrence and Euler assertions.

\par\medskip\noindent\textbf{The component and support groups.}
The exact local frame in \eqref{cdp:eq:C.NH.14} gives a global canonical
form $dp\wedge d\log\ell$ off $J$. In that chart it is
$-x^{r-1}dz\wedge dx$, and it is invariant under the
displayed deck map. Hence its divisor is $(r-1)J$.
The direct-image equality proved from this same chart is
$p_*\OO_P(kJ)=\OO_B(\lfloor k/r\rfloor[0])$.
It gives $h^0(K_P)=1$. Since $J^2=0$ and $e(P)=0$,
surface Riemann--Roch gives $\chi(\OO_P)=0$ and
$h^\bullet(\OO_P)=(1,2,1)$.

\item For the integral groups use the exact Seifert filling in
(C.ER.16). Its meridian is $r\delta+h$ and its longitude
is $\delta$, where $\delta$ is the boundary of the
once-punctured elliptic base and $h$ the product circle.
The base relation kills $\delta$ in first homology, and
the meridian then kills $h$. Thus the closed oriented
three-manifold $M_r$ has $H_1(M_r,\Z)=\Z^2$; all its
homology is free, by integral Poincare duality and the
universal coefficient sequence. The radial line coordinate
gives $P\simeq M_r\times S^1$ smoothly. For a nonreal
$\beta$ its angular rotation is isotopic to the identity
through the same circle action, so the mapping torus has
this same product identification. This proves the ranks
of $P$ in \eqref{cdp:eq:C.NH.2}.
That circle isotopy preserves the Seifert core. Thus it
identifies the pair $(P,J)$ with
$(M_r\times S^1,\mathrm{core}\times S^1)$, not merely
the ambient cohomology groups.

\item The angular loop of $J$ is the solid-torus longitude and
is zero in $H_1(M_r,\Z)$; its radial loop maps primitively
to the last $S^1$. Therefore $[J]=0$ in $H_2(P,\Z)$,
including integrally. In the Hopf component its angular
loop is zero and its radial loop is the primitive deck
generator. The integral normalization sequence of
$S=P\cup_JH$ consequently has the restriction maps
\[
 H^0(P\sqcup H,\Z)\to H^0(J,\Z):\ (a,b)\mapsto a-b,
 \qquad
 H^1(P\sqcup H,\Z)\to H^1(J,\Z):
 \begin{pmatrix}0&0&1&-1\\0&0&0&0\end{pmatrix},
                                                        \tag{C.NH.18}\hypertarget{tag-00SI}{}\label{cdp:eq:C.NH.18}
\]
up to simultaneous orientation choices for the radial
loops, and its degree-two map is zero. Both displayed
images are primitive. The exact sequence gives precisely
the free groups of ranks $(1,3,5,5,2)$.

\item Finally the constant deck cocycle represents the nonzero
class in $H^1(H,\OO_H)$. It is not a coboundary: a
holomorphic solution of $a\circ F-a=1$ on $\C^2\setminus0$
extends across zero and gives $0=1$ there. Its restriction
to $J$ is still nonzero, since the constant Laurent
coefficient of $a(\beta^r z)-a(z)$ is zero. Thus the
coherent degree-one restriction onto $H^1(J,\OO_J)$
is surjective. Together with the degree-zero surjection
and the component groups already computed, the coherent
normalization sequence gives $h^\bullet(S,\OO_S)=(1,2,1)$.
All support maps are the actual inclusions in \eqref{cdp:eq:C.NH.13}--\eqref{cdp:eq:C.NH.14}.

\end{enumerate}
\end{proof}






\begin{proposition}[Every secondary Hopf surface of algebraic dimension zero occurs unchanged]\hypertarget{tag-01DN}{}\label{cdp:st:C.2.UN.43}
\StatementRecord{cdp:st:C.2.UN.43}{proved}{01DN}{\SeriesRef{cdp:st:C.2.UN.40}, \SeriesRef{cdp:eq:C.SH.3}, \SeriesRef{cdp:st:C.2.UN.41}, \SeriesRef{cdp:eq:C.SH.9}, \SeriesRef{cdp:eq:C.SH.10}, \SeriesRef{cdp:st:C.2.UN.1}, \SeriesRef{cdp:st:C.2.UN.4}, \SeriesRef{cdp:st:C.2.UN.42}; \cite[Theorem~2.1]{Khaled2024Hopf}; \cite[\S3, pp.~230--232]{Kato1975Hopf}}

Every secondary Hopf surface $H$ of algebraic dimension zero
has an actual finite quotient presentation of one of the forms
\[
\begin{array}{ll}
 H_{\alpha,\beta}/\langle(z,w)\mapsto
          (\zeta_n^a z,\zeta_n^b w)\rangle,
 &\gcd(a,n)=\gcd(b,n)=1,\\[2pt]
 H_r(\beta)/\langle(z,w)\mapsto
          (\zeta_n^r z,\zeta_n w)\rangle,
 &\gcd(r,n)=1.
\end{array}
\tag{C.SH.1}\hypertarget{tag-0114}{}\label{cdp:eq:C.SH.1}
\]
The first primary contraction is $(\alpha z,\beta w)$ with no
multiplicative resonance, and the second is
$(\beta^r z+w^r,\beta w)$, where $0<|\alpha|,|\beta|<1$.
This is an exhaustive presentation, without a uniqueness
assertion about the displayed parameters.

Every prescribed such surface occurs itself as a smooth
multiplicity-one component of a reduced SNC fibre of a proper
map from a smooth compact complex threefold to a smooth
compact curve, with connected fibres and general fibre a
complex two-torus. All reduced fibre supports can have Euler
number zero. In the diagonal construction the base is the
finite torsion-cover curve $B'$ constructed below; no claim
that this curve is rational is needed.
\end{proposition}



\begin{proof}

\begin{enumerate}[label={$\langle1\rangle$\arabic*.},leftmargin=*]
\item \par\medskip\noindent\textbf{The actual deck group.}
Write $W=\mathbb C^2\setminus\{0\}$ and $H=W/\Gamma$.
The deck group has an infinite cyclic contraction subgroup of
finite index. Its normal core, the intersection of its finitely
many conjugates, is a normal infinite cyclic subgroup
$\langle F\rangle$ of finite index. Conjugation acts on it by
$1$ or $-1$. The nonzero conjugation-invariant character
$\gamma\mapsto\log|\det(d\gamma)_0|$ excludes the second
possibility. Thus $F$ is central and the actual primary cover is
\[
                 P=W/\langle F\rangle\longrightarrow H.
\tag{C.SH.2}\hypertarget{tag-0119}{}\label{cdp:eq:C.SH.2}
\]
Taking elementary symmetric functions on its finite generic
sheets gives $a(P)=a(H)=0$.

\item The primary normal forms and their full centralizers are
\cite[Theorem~2.1]{Khaled2024Hopf}. In the nonresonant diagonal
case a commuting power series has the form $(Az,Bw)$:
the eigenmonomial argument of \ref{cdp:st:C.2.UN.40} leaves only $z$ and
$w$ in its respective coordinate functions. In the
non-diagonal case the full centralizer is
\[
       (z,w)\longmapsto(A^r z+B w^r,Aw),\qquad A\ne0.
\tag{C.SH.3}\hypertarget{tag-011A}{}\label{cdp:eq:C.SH.3}
\]
Indeed write a commuting map as $(P,Q)$. Its second
coordinate satisfies $Q\circ F=\beta Q$. The weighted
eigenfunction calculation of \ref{cdp:st:C.2.UN.40} forces $Q=Aw$.
The first coordinate then satisfies
$P\circ F=\beta^rP+A^r w^r$. On a weighted-degree-$d$
piece with $d\ne r$, the operator $F^*-\beta^r$ is
invertible, since its only eigenvalue is
$\beta^d-\beta^r\ne0$. Thus $P=Cz+Bw^r$.
Comparing its $w^r$ coefficients gives $C=A^r$;
invertibility requires $A\ne0$. This proves the full
formula \eqref{cdp:eq:C.SH.3} directly.
The parameters $(A,B/A^r)$ in this formula have product
$(A_1,c_1)(A_2,c_2)=(A_1A_2,c_1+c_2)$. Both centralizers
are abelian. Since every element of $\Gamma$ commutes
with $F$, the actual deck group is abelian.

\item A finitely generated abelian group containing $\mathbb Z$
with finite index has the form $\mathbb Z\oplus T$, with $T$
its finite torsion subgroup. Choose a generator $f$ of the
infinite factor with the contracting sign. Some positive power
is $F^k$ times an element of $T$, $k>0$. Thus both diagonal
eigenvalue moduli of $f$ are less than one. In \eqref{cdp:eq:C.SH.3} it has
$|A|<1$ and nonzero additive parameter, since the parameter
of $F$ is nonzero. Rescaling $z$ makes its resonant coefficient
one. We obtain the actual holomorphic quotient maps
\[
       H=(W/\langle f\rangle)/T,\qquad
       \Gamma=\langle f\rangle\times T.
\tag{C.SH.4}\hypertarget{tag-011B}{}\label{cdp:eq:C.SH.4}
\]
For diagonal $T\subset(\mathbb C^*)^2$, each coordinate
character is injective: its kernel would fix a nonzero point
of the corresponding axis. Hence $T$ is cyclic and both
weights are units modulo its order. In \eqref{cdp:eq:C.SH.3}, a finite-order
element has additive parameter zero and acts as
$(A^r z,Aw)$. Its finite subgroup is cyclic, and freeness
on the first axis forces $\gcd(r,n)=1$. This proves \eqref{cdp:eq:C.SH.1}
and the finite-primary-cover comparison, including the
exclusion of a nonabelian group in algebraic dimension zero.
The abelian calculation also appears in
\cite[\S3, pp.~230--232]{Kato1975Hopf}; no topological
classification is used here.

\item Conversely the indicated finite maps are unramified.
An equation $\tau^k(z,w)=f^\ell(z,w)$ at a nonzero coordinate
forces $\ell=0$ by the modulus of the contraction multiplier;
the coprimality conditions then force $k=0$ modulo $n$.
In the non-diagonal case use $w$ when it is nonzero and
$z$ on its complement. Thus these are the exact free
finite actions to be extended.

\par\medskip\noindent\textbf{A prescribed torsion section at the nodal fibre.}
For the diagonal case start with the compact elliptic surface
$J\to\mathbb P^1$, zero section, nodal fibre at zero and
section $P$ of \ref{cdp:st:C.2.UN.41}. Its multiplicative value at that
fibre is the prescribed $\alpha$. Explicitly, on
$y^2=x^2(x+1)$ the point with group coordinate $\alpha$ is
\[
 \left(\frac{4\alpha}{(1-\alpha)^2},
       \frac{4\alpha(1+\alpha)}{(1-\alpha)^3}\right).
\tag{C.SH.5}\hypertarget{tag-011C}{}\label{cdp:eq:C.SH.5}
\]
Blowing up the nine distinct smooth base points of a general
cubic pencil through this point and the origin gives that
surface and its section, with an unchanged $I_1$ fibre.

\item Let $B^\circ$ be the smooth elliptic-fibre locus.
The exact-order-$n$ part of $J[n]|_{B^\circ}$ is a finite
etale algebraic cover. Select a component by the following
actual local section. The smooth group locus of the nodal
family has central group $\mathbb C^*$; multiplication by
$n$ has nonzero derivative at $\zeta_n$. The holomorphic
implicit-function theorem gives a unique local $n$-torsion
section specializing to $\zeta_n$, necessarily of exact order.
Take the connected cover containing its punctured restriction,
and normalize $\mathbb P^1$ in its function field. This gives
a smooth connected projective curve and finite map
\[
                             B'\longrightarrow\mathbb P^1.
\tag{C.SH.6}\hypertarget{tag-011D}{}\label{cdp:eq:C.SH.6}
\]
At the selected point $b_0$ this map is unramified: the local
sheet is the graph of the just constructed section and has
degree one over the disc. The actual $I_1$ fibre, its simple
smoothing parameter and $P(b_0)=\alpha$ remain unchanged.

\item Resolve the compact algebraic base-changed elliptic surface
and contract vertical $(-1)$-curves, away from this unchanged
tube, to obtain a smooth relatively minimal surface $J'\to B'$.
The pulled-back section and the tautological torsion point
extend as sections $P',R$: rational maps from the smooth
projective base curve to the proper regular surface extend by
the valuative criterion. They satisfy
\[
                 P'(b_0)=\alpha,\qquad R(b_0)=\zeta_n,
                 \qquad \operatorname{ord}(R)=n.
\tag{C.SH.7}\hypertarget{tag-011E}{}\label{cdp:eq:C.SH.7}
\]
Their translations extend to automorphisms of $J'$.
The minimal-graph argument of \ref{cdp:st:C.2.UN.41} applies because
$K_{J'}\cdot C\ge0$ for every vertical curve, by adjunction
and relative minimality, or by Kodaira's fibre components.
A $(-1)$-curve exceptional for one graph projection but not
the other would have nonnegative intersection with its
canonical pullback and all other effective discrepancy
curves, contradicting canonical intersection $-1$.
Both translations and inverses are therefore holomorphic;
their commutation on the smooth fibres proves commutation
everywhere.

\par\medskip\noindent\textbf{The diagonal action and the actual compact quotient.}
Put $q_0=\alpha\beta$ and form the compact suspension
\[
 Y^{\rm std}=(\mathbb C^*_u\times J')/
       \langle(u,y)\mapsto(q_0u,y+P')\rangle.
\tag{C.SH.8}\hypertarget{tag-011F}{}\label{cdp:eq:C.SH.8}
\]
Its proper compact bundle comparison is the one in
\ref{cdp:st:C.2.UN.41}. The commuting order-$n$ automorphism
\[
               (u,y)\longmapsto
                  (\zeta_n^{a+b}u,y+aR)
\tag{C.SH.9}\hypertarget{tag-011G}{}\label{cdp:eq:C.SH.9}
\]
descends. On the selected Tate tube $R=[\zeta_n]$ exactly,
by uniqueness of the torsion section in the smooth group
locus. In the multiplicative coordinate change $z=y$, $w=u/y$,
\eqref{cdp:eq:C.SH.9} is precisely
\[
                    (z,w)\longmapsto
                    (\zeta_n^a z,\zeta_n^b w).
\tag{C.SH.10}\hypertarget{tag-0115}{}\label{cdp:eq:C.SH.10}
\]
Replace the tube by \ref{cdp:st:C.2.UN.41}'s two-Hopf plumbing using
its exact punctured-family isomorphism. The action extends
on the first Hopf copy by \eqref{cdp:eq:C.SH.10}, on the second by the
inverse weights, and on the first neck by
\[
 (z,x,y)\longmapsto
     (\zeta_n^a z,\zeta_n^b x,\zeta_n^{-b}y),
                         \qquad xy=t.
\tag{C.SH.11}\hypertarget{tag-0116}{}\label{cdp:eq:C.SH.11}
\]
Interchange $a,b$ and the axes for the other neck.
Substitution in $(z,w)\mapsto(z^{-1},t/w)$ and
$(z,w)\mapsto(t/z,w^{-1})$ proves equivariance of both
gluing maps. The action is free on the whole central
support: on a double curve the equation
$\zeta_n^{ak}z=\alpha^\ell z$, or its counterpart, forces
$\ell=0$ by modulus and then $k=0$ by coprimality.
Off the double curves use the already verified Hopf
freeness. On a nearby torus it is a finite translation,
also free. Hence the quotient is smooth on this entire
tube and its two central components are the prescribed
secondary Hopf surface itself. The inverse weights
generate the same subgroup.

\item Take the finite quotient of the entire compact threefold
and resolve it over the other finitely many bad values,
leaving this smooth tube unchanged. General fibres are
quotients of tori by translations, hence tori. Connectedness
passes through the finite quotient and the connected fibres
of the birational resolution. The distinguished reduced
SNC fibre has two Hopf components and two elliptic double
curves, and therefore Euler number zero. Every other
reduced fibre has even Euler number by \ref{cdp:st:C.2.UN.1}.
Apply \ref{cdp:st:C.2.UN.4} only over those other values to make their
Euler numbers zero; none of these operations changes
the specified Hopf surface or its distinguished tube.
This proves compact diagonal secondary occurrence over
the actual curve $B'$.

\par\medskip\noindent\textbf{The nonlinear scalar action.}
Use the intermediate compact family $\widetilde Y\to C$ in
\eqref{cdp:eq:C.NH.11} for the prescribed $H_r(\beta)$, before the
additional involution in \eqref{cdp:eq:C.NH.15}. Its smooth nonzero
tautological line space has coordinates
\[
 (z,w,x)\in(\mathbb C^2\setminus\{0\})\times\mathbb C,
                   \qquad p-p_0=x^r z,\quad s=xw.
\tag{C.SH.12}\hypertarget{tag-0117}{}\label{cdp:eq:C.SH.12}
\]
Multiplication of that line by $\zeta_n$ is the actual map
\[
 (z,w,x)\longmapsto
       (\zeta_n^r z,\zeta_n w,\zeta_n^{-1}x).
\tag{C.SH.13}\hypertarget{tag-0118}{}\label{cdp:eq:C.SH.13}
\]
It commutes with the lifted base translations and the
contracting deck map. It is free when $\gcd(r,n)=1$:
a nonzero $w$ or $x$ forces the corresponding root to be
one; if both vanish then $z\ne0$ and its $r$-th power
must be one. It stays free after the infinite cyclic
quotient. Away from the exceptional component the
canonical line trivialization turns equality with a
deck power into $\zeta_n^k=\beta^\ell$, so $\ell=0$.
On the Hopf component use its nonzero $w$ coordinate,
or its $z$ coordinate if $w=0$, for the same modulus
argument. Coprimality then finishes the proof.

\item The resulting quotient is a smooth compact threefold.
Its desired component is exactly the second quotient
in \eqref{cdp:eq:C.SH.1}, with no change to that surface. On every
general torus fibre the action is a line-coordinate
translation, so the quotient is a torus. The distinguished
full fibre remains reduced SNC, and its Euler number
zero is preserved by the free finite quotient. Any
remaining values can be balanced by \ref{cdp:st:C.2.UN.4} away
from that fibre.

\item The base for this construction is the specified compact
curve in \ref{cdp:st:C.2.UN.42}. No rational-base globalization is
needed for this occurrence assertion.

\end{enumerate}
\end{proof}






\begin{proposition}[Every compact curve is the ruling base of an actual normalization component]\hypertarget{tag-01DP}{}\label{cdp:st:C.2.UN.45}
\StatementRecord{cdp:st:C.2.UN.45}{proved}{01DP}{\SeriesRef{cdp:st:C.2.UN.4}, \SeriesRef{cdp:eq:C.RU.7}, \SeriesRef{cdp:eq:C.RU.8}, \SeriesRef{cdp:eq:C.RU.5}}

For every smooth compact connected complex curve $C$, there is a smooth
compact complex threefold $Y$ and a proper map $f:Y\to\PP^1$ with connected
fibres, whose restriction over $\PP^1\setminus\{0\}$ is the product with a
fixed two-torus, and whose reduced zero fibre has Euler number zero and
contains a smooth multiplicity-one component
\[
                         R=\PP(N_{C/U})\longrightarrow C.
\tag{C.RU.1}\hypertarget{tag-010C}{}\label{cdp:eq:C.RU.1}
\]
Here $C\hookrightarrow U$ is the explicit vertical embedding constructed
below. The component $R$ is retained unchanged by the Euler-balancing
operations. Its birational class is that of $C\times\PP^1$.
Consequently every ruled surface birational class occurs in an actual
compact zero-Euler torus fibre. A rational component is also obtained
unchanged as the exceptional plane of a point blowup followed by balancing
away from that plane.

Writing $g=g(C)$, the integral cohomology of $R$ is torsion-free and has ranks
\[
              (b_0,b_1,b_2,b_3,b_4)(R)=(1,2g,2,2g,1),
              \qquad h^\bullet(R,\OO_R)=(1,g,0).
\tag{C.RU.2}\hypertarget{tag-010D}{}\label{cdp:eq:C.RU.2}
\]
The construction proves occurrence of every ruled birational class; it
does not prescribe an arbitrary rank-two bundle as the normal bundle in
\eqref{cdp:eq:C.RU.1}.
\end{proposition}



\begin{proof}

\begin{enumerate}[label={$\langle1\rangle$\arabic*.},leftmargin=*]
\item Choose a nonconstant meromorphic function $x$ on $C$, supplied by
Riemann--Roch. It gives a finite map $C\to\PP^1$, so the meromorphic
function field of $C$ is a finite separable extension of $\C(x)$.
Choose a primitive element $y$. The meromorphic map
\[
                    C\longrightarrow\PP^2,\qquad [1:x:y]
\tag{C.RU.3}\hypertarget{tag-010E}{}\label{cdp:eq:C.RU.3}
\]
extends at every point: multiply its three local meromorphic coordinates
by a power of a uniformizer so that their minimum vanishing order is zero.
At least one resulting holomorphic coordinate is nonzero. The image is an
integral projective plane curve $D$, and the map is finite and birational,
because $\C(x,y)=\mathcal M(C)$. Thus $C$ is the actual normalization of
$D$. For $C=\PP^1$ one may choose $x=s,y=s^2$, giving a smooth conic.
In every case choose projective coordinates so that $D$ is not the fixed
line used below.

\item Start with $A\times\PP^1$, where $A$ is any fixed smooth complex two-torus,
and blow up a point of $A\times\{0\}$. Its exceptional divisor is an
actual $\PP^2$. Place $D$ in this plane. Successively blow up singular
points of its strict transform, performing each point blowup in the
ambient smooth threefold. The strict transform of the plane is precisely
its surface point blowup: in local coordinates this is the strict
transform of the coordinate plane in the three ordinary blowup charts.
The full map to $\PP^1$ is unchanged over its punctured base.

\item Here the point-resolution process is finite without an unproved
termination assertion. For an integral curve $D_0$ on a smooth surface,
normalization gives
\[
  p_a(D_0)-g(C)=\operatorname{length}
                (\nu_*\OO_C/\OO_{D_0})\ge0.
\tag{C.RU.4}\hypertarget{tag-010F}{}\label{cdp:eq:C.RU.4}
\]
At a singular point the multiplicity $m$ is at least two. Under a surface
point blowup, $D_1=\pi^*D_0-mE$ and
$K_{V_1}=\pi^*K_{V_0}+E$. Adjunction gives
\[
                        p_a(D_1)=p_a(D_0)-\frac{m(m-1)}2.
\tag{C.RU.5}\hypertarget{tag-010G}{}\label{cdp:eq:C.RU.5}
\]
The normalization is unchanged through this proper birational map.
Therefore its total nonnegative normalization length decreases by at
least one at every singular-point blowup. After finitely many such
blowups the strict transform is smooth and is the specified $C$ itself.
This gives a genuine closed embedding $C\hookrightarrow U$ into the
resulting smooth compact threefold.

\item The strict transform of the original exceptional plane has full-fibre
multiplicity one, since that was its coefficient after the first point
blowup and later strict transforms retain the coefficient. The curve
$C$ is not any boundary curve of this plane: it dominates the original
$D$, which is different from the line of intersection with the original
torus component, whereas every subsequently exceptional boundary curve
maps to a point of that plane. Consequently a generic point of $C$ lies
only in this multiplicity-one component. Blow up the smooth curve $C$ in
$U$. The exceptional divisor is exactly the projectivized normal bundle
in \eqref{cdp:eq:C.RU.1}, with coefficient one in the full fibre. This follows also
from the local charts of the blowup of $(x,y)$: its exceptional fibre at
each point of $C$ is the projective space of normal directions.

\item Choose two meromorphic sections of $N_{C/U}$ which form a basis at the
generic point. Such sections exist by twisting by a sufficiently positive
line bundle on the curve and applying Riemann--Roch. They give an actual
bundle trivialization away from a finite set, and hence a birational map
\[
                         \PP(N_{C/U})\dashrightarrow C\times\PP^1.
\tag{C.RU.6}\hypertarget{tag-010H}{}\label{cdp:eq:C.RU.6}
\]
The same meromorphic-frame map applies to every rank-two bundle on $C$.
It therefore identifies this birational class with that of every ruled
surface over this specified curve. In genus zero it also identifies the
rational class, for example through the affine-coordinate birational map
$\PP^2\dashrightarrow\PP^1\times\PP^1$.

\item We retain the exact Euler changes, including all singularities along the
centres. The reduced inverse image of a point centre replaces that point
by $\PP^2$, adding two to the support Euler number. The reduced inverse
image of the smooth curve $C$ replaces each point by $\PP^1$, adding
$e(C)=2-2g$. These are the proper-map comparisons of (C.ZC.8), which do
not impose a normal-crossing or reduced full-fibre hypothesis. If $p$ is
the total number of point blowups, including the initial one, the support
after blowing up $C$ has
\[
                              e=2p+2-2g.
\tag{C.RU.7}\hypertarget{tag-010I}{}\label{cdp:eq:C.RU.7}
\]
It is even. There is a smooth open part of the original torus component
disjoint from $R$, since only finitely many proper curves and points have
been affected there. Perform the positive point blocks or negative
point--quartic blocks of \ref{cdp:st:C.2.UN.4} in such an open part and in the new
exceptional components they create. Each changes Euler number by exactly
$2$ or $-2$. Choosing the indicated finite number makes \eqref{cdp:eq:C.RU.7} zero.
All these centres are disjoint from $R$, so the displayed smooth component,
its normal projection to $C$, and its coefficient one are unchanged.
All other fibres remain the original smooth tori. Every blowup has
connected fibres, so the composite map retains connected fibres as well.
Compactness and smoothness follow at each step from the actual point or
smooth-curve blowup of the compact smooth threefold.

\item For the cohomology comparison use the lines convention in \eqref{cdp:eq:C.RU.1} and
put $\xi=c_1(\OO_{\PP(N)}(1))$. Pullback and cup product give the integral
projective-bundle map
\[
 H^k(C,\Z)\oplus H^{k-2}(C,\Z)
       \longrightarrow H^k(R,\Z),\qquad
 (a,b)\longmapsto\pi^*a+\xi\smile\pi^*b.
\tag{C.RU.8}\hypertarget{tag-010J}{}\label{cdp:eq:C.RU.8}
\]
It is an isomorphism: on every fibre the classes $1,\xi$ are the standard
integral basis of $H^*(\PP^1,\Z)$, so the fibrewise basis glues and the
Leray--Hirsch comparison has these exact maps. The relation for this lines
convention is $\xi^2+\pi^*c_1(N)\xi=0$. Since the curve cohomology is
torsion-free, \eqref{cdp:eq:C.RU.8} proves the integral assertion. The fibrewise
calculation $H^0(\PP^1,\OO)=\C$, $H^1(\PP^1,\OO)=0$ gives
$\pi_*\OO_R=\OO_C$, $R^1\pi_*\OO_R=0$, with their restriction and
evaluation maps. Leray therefore gives the coherent tuple in \eqref{cdp:eq:C.RU.2}.

\item There is an explicit all-genus computation supporting the construction.
For $g\ge1$ take the plane curve
\[
                 y^2z^{2g-1}=x^{2g+1}-z^{2g+1}.
\tag{C.RU.9}\hypertarget{tag-010K}{}\label{cdp:eq:C.RU.9}
\]
It is smooth on $z\ne0$. At its single point at infinity, use
$u=x/y,v=z/y$; the equation is $v^{2g-1}(1+v^2)=u^{2g+1}$.
The holomorphic change $v'=v(1+v^2)^{1/(2g-1)}$ identifies its germ with
the plane branch $v'^{2g-1}=u^{2g+1}$. The two exponents are coprime.
Successive point blowups subtract the smaller exponent from the larger,
and the strict curve is smooth when the smaller exponent becomes one.
Summing \eqref{cdp:eq:C.RU.5} through these actual charts gives total normalization
defect $2g(g-1)$, so its normalization has genus $g$.
The executed code \path{code/root_ruled_occurrence.py} checks all exponent
charts, genus drops, Euler changes, balancing counts and the integral
projective-bundle ranks for $g=1,\ldots,50$. The displayed subtraction
and adjunction formulas prove the same identities for every $g$.

\end{enumerate}
\end{proof}






\section{Marked Hopf, ruled and elliptic supports with all multiplicities}\hypertarget{tag-0346}{}\label{sec:II:10}


\begin{proposition}[Hopf-axis adjacency with arbitrary multiplicities and point blowups]\hypertarget{tag-01DK}{}\label{cdp:st:C.2.UN.40}
\StatementRecord{cdp:st:C.2.UN.40}{proved}{01DK}{\SeriesRef{cdp:eq:C.HA.3}, \SeriesRef{cdp:eq:C.HA.4}; \cite[Theorem~2.1]{Khaled2024Hopf}}

Let $f:U\to\Delta$ be a proper degeneration of smooth complex
two-tori with smooth total space, and suppose that its reduced
central support has simple normal crossings. Write its full
fibre as $\sum_i m_iE_i$. Choose a generator of the rank-one
sheaf $f_*K_U$ and write its divisor as $\sum_i k_iE_i$.
Set $r_i=(k_i+1)/m_i$.

Suppose that the minimal model $H_i$ of a component $E_i$
is a Hopf surface of algebraic dimension zero. The contraction
$\rho_i:E_i\to H_i$ may consist of arbitrarily many point
blowdowns. Choose an actual finite unramified primary Hopf
cover $p:P\to H_i$.

If $P$ has diagonal contraction $(z,w)\mapsto(\alpha z,\beta w)$,
the images of both coordinate axes are boundary curves on
$H_i$: their strict transforms in $E_i$ meet another central
component. For each such strict transform, let $E_j$ be its
unique adjacent component at the generic point. Then
\[
                 m_i(k_j+1)-m_j(k_i+1)=0,
                 \qquad r_j=r_i.
\tag{C.HA.1}\hypertarget{tag-00L9}{}\label{cdp:eq:C.HA.1}
\]
If $P$ has non-diagonal contraction
$(z,w)\mapsto(\beta^r z+w^r,\beta w)$, $r\geq1$, the image
of its unique elliptic curve is a boundary curve. Its unique
generic neighbour satisfies
\[
       m_i(k_j+1)-m_j(k_i+1)=-m_i r,
       \qquad m_j(r_j-r_i)=-r<0.
\tag{C.HA.2}\hypertarget{tag-00LA}{}\label{cdp:eq:C.HA.2}
\]
These are integral divisor-class consequences on the actual
components. In particular a non-diagonal algebraic-dimension-zero
Hopf component cannot meet another algebraic-dimension-zero
Hopf component along this elliptic curve. Diagonal such
components can meet along their axes only at equal ratios.
\end{proposition}



\begin{proof}

\begin{enumerate}[label={$\langle1\rangle$\arabic*.},leftmargin=*]
\item The direct image $f_*K_U$ is coherent, torsion-free and of
rank one on the disc, because a smooth torus fibre has one
holomorphic two-form. Shrink the disc and choose a generator.
Its evaluation is an actual holomorphic three-form on $U$.
On every punctured fibre it is a nonzero translation volume
form, so its zero divisor is vertical and gives
$K_U\simeq\mathcal O_U(\sum_i k_iE_i)$, with $k_i\geq0$.
For $j\ne i$ put $D_{ij}=E_j|_{E_i}$ and
$D_i=\sum_{j\ne i}D_{ij}$. Each is the reduced intersection
divisor specified by the simple-normal-crossing coordinates.
The full fibre is principal. Adjunction and its restriction
to $E_i$ therefore give the following integral identity:
\[
 m_i(K_{E_i}+D_i)
 \sim\sum_{j\ne i}
       \bigl[m_i(k_j+1)-m_j(k_i+1)\bigr]D_{ij}.
\tag{C.HA.3}\hypertarget{tag-00LB}{}\label{cdp:eq:C.HA.3}
\]
Indeed, first use
$K_{E_i}=((k_i+1)E_i+\sum_{j\ne i}k_jE_j)|_{E_i}$,
then replace $m_i E_i|_{E_i}$ by
$-\sum_{j\ne i}m_jD_{ij}$, and finally add $m_iD_i$.
Thus \eqref{cdp:eq:C.HA.3} has no rational divisor or divisibility convention.

\item We specify its pushforward through the point blowdowns.
For a proper birational map of smooth surfaces the common
meromorphic function field sends the divisor of a function
to the divisor of the same function. Sending every prime
curve dominating a curve to that curve and every exceptional
prime curve to zero consequently induces a homomorphism
of divisor-class groups. It sends $K_{E_i}$ to $K_{H_i}$:
pull back a meromorphic two-form on $H_i$ and observe that
its additional Jacobian divisor is exceptional. This is a
divisor-class homomorphism; no equality with the coherent
direct image of an arbitrary line bundle is asserted.
Applying it to \eqref{cdp:eq:C.HA.3} gives
\[
 m_i(K_{H_i}+B_i)
 \sim\sum_{j\ne i}
       \bigl[m_i(k_j+1)-m_j(k_i+1)\bigr]\rho_{i*}D_{ij},
 \qquad B_i=\rho_{i*}D_i.
\tag{C.HA.4}\hypertarget{tag-00LC}{}\label{cdp:eq:C.HA.4}
\]
The divisor $B_i$ is reduced. A point-blowup morphism is an
isomorphism at the generic point of every curve on its target,
so there is exactly one strict transform dominating that curve.
At its generic point the SNC equation has exactly two branches,
$x_i x_j=0$; an additional branch can meet it only at isolated
triple points. Hence this strict transform occurs in exactly
one $D_{ij}$, with coefficient one.

\item The pullback of \eqref{cdp:eq:C.HA.4} through the specified unramified
cover $p:P\to H_i$ is an identity of integral divisor classes,
and $p^*K_{H_i}=K_P$. We record precisely the curve and line
classes on $P$. The primary normal forms are those in
\cite[Theorem~2.1]{Khaled2024Hopf}. The finite meromorphic
field comparison proves $a(P)=a(H_i)=0$, by taking elementary
symmetric functions on the finite generic sheets, as in
(C.CF.9).

\item Here is also a direct check of the curves needed in that
comparison. The inverse image of a compact curve in
$\mathbb C^2\setminus\{0\}$ has analytic closure across
the origin by the Remmert--Stein extension theorem. Its
divisor on $\mathbb C^2$ is principal, say $(F=0)$.
For the contraction $\gamma$ one has
$F\circ\gamma=uF$ with $u$ a holomorphic unit. Write
$u=u(0)e^h$, $h(0)=0$. The locally uniformly convergent sum
$a=\sum_{n\geq0}h\circ\gamma^n$ satisfies
$a-a\circ\gamma=h$. Replacing $F$ by $e^aF$ makes it an
eigenfunction of $\gamma^*$ with constant eigenvalue.

\item For a diagonal contraction, distinct monomials can have
the same eigenvalue only when $\alpha^a\beta^b=1$ for
some nonzero integer pair. Because $|\alpha|,|\beta|<1$,
such a pair has opposite signs and gives a nonconstant
invariant meromorphic monomial, contrary to $a(P)=0$.
Thus each eigenfunction is one monomial, and the only
curves are the two axes $C_1,C_2$. In the non-diagonal
case use weighted degree $\deg z=r$, $\deg w=1$.
On the finite-dimensional space of weighted degree $d$
the contraction has eigenvalue $\beta^d$ times the
translation operator on polynomials in $z/w^r$.
Its eigenvectors are precisely the constants in that
variable. The eigenvalues $\beta^d$ for different $d$
are distinct, so every eigenfunction is a power of $w$.
There is exactly the curve $C=(w=0)$.

\item For completeness, denote by $L_\lambda$ the line bundle
whose lifted fibre coordinate is multiplied by $\lambda$
under the contraction. The map
$\lambda\mapsto L_\lambda$ is injective. A trivialization
would be a nowhere-zero holomorphic eigenfunction on
$\mathbb C^2\setminus\{0\}$. Hartogs extends it and its
inverse across zero, and evaluation there forces
$\lambda=1$. Coordinate functions and the Jacobian
therefore give the exact classes
\[
\begin{array}{ll}
 \mathcal O(C_1)=L_\alpha,
 \quad\mathcal O(C_2)=L_\beta,
 &K_P=\mathcal O(-C_1-C_2),\\[2pt]
 \mathcal O(C)=L_\beta,
 &K_P=\mathcal O(-(r+1)C)
 \quad\hbox{in the non-diagonal case}.
\end{array}
\tag{C.HA.5}\hypertarget{tag-00LD}{}\label{cdp:eq:C.HA.5}
\]
The first two classes are integrally independent: a relation,
or even a torsion relation, would give the forbidden
positive-power resonance after raising to a power.
The class of $C$ in the second line has infinite order
because $0<|\beta|<1$.

Every curve on $H_i$ has nonempty reduced inverse image
on $P$, since $p$ is unramified and surjective. Thus the
pullback of $B_i$ is a reduced subset of the indicated
axis curves. In the diagonal case, compare coefficients
of their independent classes in the pullback of \eqref{cdp:eq:C.HA.4}.
If an axis were absent, its coefficient on the left
would be $-m_i$ and on the right zero. Each axis is
therefore present. Its coefficient on the left is then
zero, and the coefficient on the right is exactly that
of its unique generic neighbour. This proves \eqref{cdp:eq:C.HA.1},
including when one boundary curve has more than one
component in its finite inverse image.

In the non-diagonal case, absence of $C$ would instead
give the impossible nonzero class $-m_i(r+1)C\sim0$.
It is present, with coefficient one. Comparing its
infinite-order class gives
\[
       -m_i r=m_i(k_j+1)-m_j(k_i+1),
\tag{C.HA.6}\hypertarget{tag-00LE}{}\label{cdp:eq:C.HA.6}
\]
which is \eqref{cdp:eq:C.HA.2}. This proves both assertions on the
original component after all its point blowups, with
the actual finite-primary-cover comparison included.

Finally, the shared strict axis is a smooth elliptic
curve. It cannot be contracted by a point-blowdown
morphism on the neighbouring smooth component, since
all exceptional curves of such a morphism are rational.
If that neighbour also has an algebraic-dimension-zero
Hopf minimal model, this same curve is one of its
minimal axes. Its diagonal rule would force equal
ratios, contradicting \eqref{cdp:eq:C.HA.2}; its non-diagonal rule
would force the opposite strict inequality. Neither
is possible. The diagonal--diagonal case merely repeats
the equality \eqref{cdp:eq:C.HA.1}. The theorem gives these exact
adjacency restrictions and does not infer a global
nonexistence assertion from a local decrease.

\end{enumerate}
\end{proof}






\begin{proposition}[The diagonal Hopf cycle criterion and the complete two-component case]\hypertarget{tag-01DQ}{}\label{cdp:st:C.2.UN.46}
\StatementRecord{cdp:st:C.2.UN.46}{proved}{01DQ}{\SeriesRef{cdp:st:C.2.UN.41}, \SeriesRef{cdp:st:C.2.UN.43}, \SeriesRef{cdp:st:C.2.UN.40}, \SeriesRef{cdp:eq:C.HC.3}, \SeriesRef{cdp:eq:C.HC.4}, \SeriesRef{cdp:eq:C.HC.7}, \SeriesRef{cdp:eq:C.HC.5}, \SeriesRef{cdp:eq:C.HC.8}, \SeriesRef{cdp:eq:C.HC.10}, \SeriesRef{cdp:eq:C.HC.9}}

Consider a reduced SNC fibre of an actual smooth-total
degeneration of two-tori, all of whose irreducible components
are unblown diagonal Hopf surfaces of algebraic dimension zero.
Primary and secondary Hopf surfaces are both allowed.
Its components form a cycle: each of the two axis elliptic
curves of each component is a double curve, and there are
no triple points.

For an oriented cycle of specified diagonal Hopf surfaces
$H_1,\ldots,H_\ell$, specify the actual elliptic gluing
isomorphisms and dual-normal isomorphisms along its edges.
Let
\[
 G_i=H_i\setminus(C_i^-\cup C_i^+)
       =\mathbb C^2/\Lambda_i,\qquad \Lambda_i\simeq\mathbb Z^3,
\tag{C.HC.1}\hypertarget{tag-00LO}{}\label{cdp:eq:C.HC.1}
\]
where $\Lambda_i$ is the inverse image of its diagonal
Hopf deck group under coordinatewise exponentiation.
The edge data determine isomorphisms of these integral
lattices and a cycle product $M:\Lambda_1\to\Lambda_1$.
For the finite proper local plumbing constructed below,
the exact criterion is
\[
 \text{$Y_t$ is a two-torus}\quad(0<|t|\ll1)
       \quad\Longleftrightarrow\quad M=I.
\tag{C.HC.2}\hypertarget{tag-00LR}{}\label{cdp:eq:C.HC.2}
\]
Every actual torus degeneration in the stated SNC branch
has $M=I$. The definition of $M$, its independence from
the choices of logarithms, and the actual nearby quotient
are given in the proof.

For two components, the criterion is more explicit.
The two Hopf surfaces are isomorphic. After an actual
diagonal automorphism of one of them, both gluing maps
are inversion on the corresponding axis elliptic curves.
Thus the reduced fibre is precisely two copies of the
same diagonal Hopf surface joined by the two inversion
maps. Every such reduced fibre occurs in a compact
zero-Euler torus filling, by \ref{cdp:st:C.2.UN.41} and \ref{cdp:st:C.2.UN.43}.
This completely determines the reduced supports in
this two-component branch. The constants in the
dual-normal isomorphisms remain explicit local
smoothing parameters in the quotient below.

The same local lattice criterion applies to the
specified operation of identifying the two distinct
axes of one Hopf surface. Its single transition
cannot have $M=I$, so that particular self-gluing
admits no smooth-total torus smoothing.
\end{proposition}



\begin{proof}

\begin{enumerate}[label={$\langle1\rangle$\arabic*.},leftmargin=*]
\item \par\medskip\noindent\textbf{The cycle and its actual edge data.}
\ref{cdp:st:C.2.UN.40} forces both axis curves into the boundary of
every component of the given SNC fibre. An
algebraic-dimension-zero diagonal primary Hopf surface
has exactly these two curves, and its free diagonal
finite quotients have their two distinct images, by
\ref{cdp:st:C.2.UN.43}. They are disjoint. Consequently there is
no triple point, and at the generic point of either
axis the SNC equation has exactly one other branch.
Every vertex of the connected incidence graph has
degree two, so the graph is a cycle, with a pair
of edges allowed between two vertices.

\item Orient this cycle. On $H_i$ choose diagonal coordinates
so that the incoming axis $C_i^-$ has coordinate $z_i$
and the outgoing axis $C_i^+$ has coordinate $w_i$.
Thus the outgoing normal coordinate is $z_i$, and
the incoming normal coordinate is $w_i$.
For the edge from $i$ to $i+1$ write its data as
\[
 \phi_i:C_i^+\xrightarrow{\sim}C_{i+1}^-,
 \qquad
 \theta_i:N_{C_i^+/H_i}\otimes
       \phi_i^*N_{C_{i+1}^-/H_{i+1}}
           \xrightarrow{\sim}\mathcal O_{C_i^+}.
\tag{C.HC.3}\hypertarget{tag-00LS}{}\label{cdp:eq:C.HC.3}
\]
These are holomorphic isomorphisms. In an actual
reduced fibre the second comes from the defining
function: locally $f=xy$ times a unit, and its
transverse bilinear term is precisely \eqref{cdp:eq:C.HC.3}.
There is no extra divisor in this isomorphism,
because there are no triple points.

\item Each axis neighbourhood is an entire holomorphic
normal line bundle. For a primary diagonal Hopf
surface the outgoing example is
\[
 (\mathbb C^*_{w}\times\mathbb C_z)/
       ((w,z)\sim(\beta w,\alpha z)).
\tag{C.HC.4}\hypertarget{tag-00LT}{}\label{cdp:eq:C.HC.4}
\]
For a secondary surface take also the specified
finite diagonal quotient. Its action on the base
elliptic curve is free, since both finite characters
are faithful; the quotient of \eqref{cdp:eq:C.HC.4} is therefore
an actual line bundle on the quotient elliptic curve.
Removing its zero section gives exactly $G_i$.
Its universal cover is $\mathbb C^2$ in logarithmic
coordinates. The lattice $\Lambda_i$ includes
$2\pi i\mathbb Z^2$ and one logarithm of a contracting
deck transformation, together with the finite
diagonal roots. It is a discrete rank-three lattice,
and it spans $\mathbb C^2$ over $\mathbb C$.

\par\medskip\noindent\textbf{The affine edge maps and the integral product.}
For $t\ne0$, the normal-product equation and
\eqref{cdp:eq:C.HC.3} define a global biholomorphism
$\Phi_i(t):G_i\to G_{i+1}$: send the base by $\phi_i$
and send a nonzero normal vector to its dual vector
whose product is $t$. Lift it to the logarithmic
universal covers. An elliptic isomorphism lifts to
an affine map, so its base coordinate is
$c_i\log w_i+d_i$, with $c_i\ne0$.
The normal coordinate has the form
$-\log z_i+h_i(\log w_i)+\log t$.
The nonzero holomorphic fibre multiplier has a
holomorphic logarithm $h_i$ on the simply connected
base cover. For every elliptic period $\lambda$,
$h_i(u+\lambda)-h_i(u)$ is constant, since both
normal bundles in \eqref{cdp:eq:C.HC.4} have constant factors of
automorphy. Hence $h_i'$ descends to a holomorphic
function on a compact elliptic curve and is constant.
We obtain the exact lift
\[
 \binom{\log z_{i+1}}{\log w_{i+1}}
 =
 L_i\binom{\log z_i}{\log w_i}
 +\binom{d_i}{e_i+\log t},
 \qquad
 L_i=\begin{pmatrix}0&c_i\\-1&k_i\end{pmatrix}.
\tag{C.HC.5}\hypertarget{tag-00LU}{}\label{cdp:eq:C.HC.5}
\]
This derivation uses the actual entire line bundles
and their isomorphism; it does not replace an
analytic neighbourhood by a formal normal bundle.

\item Since $\Phi_i(t)$ is a global biholomorphism,
$L_i\Lambda_i=\Lambda_{i+1}$. Choosing integral bases
gives actual matrices $M_i\in\operatorname{GL}_3(\mathbb Z)$.
Put
\[
 M=M_\ell\cdots M_1,\qquad L=L_\ell\cdots L_1.
\tag{C.HC.6}\hypertarget{tag-00LV}{}\label{cdp:eq:C.HC.6}
\]
Changing the initial basis conjugates $M$.
Changing any logarithm changes only its translation
by a lattice vector. Multiplying $\theta_i$ by a
nonzero constant changes $e_i$, not $L_i$.
Thus the condition $M=I$ is well defined.
It is equivalent to $L=I$, because $\Lambda_1$
spans the complex vector space.

\par\medskip\noindent\textbf{Actual finite plumbing and properness.}
Use \eqref{cdp:eq:C.HC.4} to choose disjoint tubular neighbourhoods
of the two axes in each Hopf component. They have
flat Hermitian normal metrics. The dual-normal
isomorphism in \eqref{cdp:eq:C.HC.3} identifies the two flat
metrics up to a positive constant: the ratio of
compatible flat metrics is the exponential of a
harmonic function on the compact elliptic curve.
Adjust that constant so that the product of the
two normal norms is $|t|$ in each neck.

\item For each edge take the open part of
$N_{C_i^+/H_i}\oplus N_{C_i^+/H_i}^{-1}$
with bounded fibre coordinates $x,y$, and the
actual holomorphic function $xy=t$.
Its two branches are identified with the two
specified axis neighbourhoods using $\phi_i,\theta_i$.
Glue these finitely many necks to the products
of the compact component cores with a disc.
Explicitly choose $0<r_0<R_0$ and $|t|<r_0^2$,
use $\|x\|,\|y\|<R_0$, and attach where one of
these norms is greater than $r_0$.
The inverse map inserts $t/x$ or $t/y$.
The two ends of a neck cannot overlap, because
$\|x\|\|y\|=|t|<r_0^2$; the two axis collars in
a component are disjoint as well.

\item These are actual biholomorphic gluing maps.
Their graphs are closed relative to the pieces:
at one boundary of an attaching annulus its core
point is absent, and at the other boundary its
neck point is absent. Each equivalence class
meets at most one neck and one core. The finite
gluing is therefore Hausdorff. The intermediate
closed cores and closed smaller neck disc bundles
used in \ref{cdp:st:C.2.UN.41} cover the inverse image of every
closed subdisc; each is compact over a compact
elliptic curve or compact Hopf core. This proves
properness. The result is a smooth complex
threefold over a disc, with exactly the prescribed
reduced normal-crossing central surface.

\par\medskip\noindent\textbf{The exact fundamental group, including actual given smoothings.}
The open group $G_i=\mathbb C^2/\Lambda_i$ is
smoothly $T^3\times\mathbb R$.
For example, a linear height annihilating
$\operatorname{span}_{\mathbb R}\Lambda_i$ has
compact three-torus levels; the usual axis norms
are exponentials of this height. A component
core and each annular neck consequently retract
onto a three-torus. Every collar inclusion
induces the indicated isomorphism on its
rank-three fundamental group. Van Kampen,
first along a spanning tree of the cycle and
then at its last edge, gives the actual group
\[
                       \pi_1(Y_t)=\Lambda_1\rtimes_M\mathbb Z.
\tag{C.HC.7}\hypertarget{tag-00LW}{}\label{cdp:eq:C.HC.7}
\]
The copy of $\Lambda_1$ is injected, as in the
displayed semidirect product presentation.

\item For an arbitrary actual smoothing of the same
reduced SNC data, this calculation still uses
actual smooth maps. Near a double curve $C$,
the divisor line bundles of its two branches
have canonical sections $s_1,s_2$, and $f$
trivializes their tensor product. Choose a
smooth retraction to $C$ and smooth line-bundle
identifications with the pullbacks of their
restrictions, choosing the second identification
so that the product trivialization is preserved.
The map consisting of the retraction and
$(s_1,s_2)$ has invertible derivative along $C$.
After shrinking the neighbourhood it is a
smooth tubular diffeomorphism to the two
normal line bundles, and its actual function
is $xy$. These are the normal bundles and
the product isomorphism \eqref{cdp:eq:C.HC.3}.
On a compact core away from these neighbourhoods
$f$ is a submersion; horizontal transport
identifies its nearby cores. The attaching
maps are isotopic to the normal maps just
specified, by radial contraction of the
tubular identifications. Thus their maps on
$\pi_1$ are exactly $M_i$, proving \eqref{cdp:eq:C.HC.7}
for the given smoothing as well. No holomorphic
linearization of its full neighbourhood is
asserted or needed.

\item If the given $Y_t$ is a two-torus, its fundamental
group is abelian. The injected copy of $\Lambda_1$
in \eqref{cdp:eq:C.HC.7} then forces $M=I$.
This proves the necessary half of \eqref{cdp:eq:C.HC.2}.

\par\medskip\noindent\textbf{Closed one-forms give the sufficient torus comparison.}
Suppose $M=I$, equivalently $L=I$.
The vector of closed holomorphic one-forms
$(d\log z_i,d\log w_i)$ on $G_i$ transforms
by the constant matrix $L_i$ on each overlap.
Choose constant invertible matrices $B_i$
successively so that $B_{i+1}L_i=B_i$.
The identity of the cycle product makes this
choice consistent on its last edge.
These forms therefore glue on each actual
nearby compact fibre, including all its necks,
to two everywhere independent closed holomorphic
one-forms.

\item Their dual holomorphic vector fields commute,
because the two dual one-forms are closed.
On a compact complex manifold holomorphic
vector fields have entire complex-time flows:
a finite coordinate cover supplies a uniform
local existence interval, which can be iterated.
The two commuting flows give a transitive
holomorphic $\mathbb C^2$-action; every orbit
is open, and the fibre is connected.
Its stabilizer is discrete and cocompact,
since the infinitesimal action is an
isomorphism and the fibre is compact.
The orbit map is therefore the actual
biholomorphism
\[
                  Y_t\simeq\mathbb C^2/\Lambda(t)
\tag{C.HC.8}\hypertarget{tag-00LX}{}\label{cdp:eq:C.HC.8}
\]
for a rank-four real lattice.

\item One may read this lattice without discarding
the translation data. Compose the affine
maps \eqref{cdp:eq:C.HC.5} on a chosen logarithm branch:
\[
                    \Phi_\ell(t)\cdots\Phi_1(t)(v)
                          =v+b_0+v_0\log t.
\tag{C.HC.9}\hypertarget{tag-00LY}{}\label{cdp:eq:C.HC.9}
\]
In the normalization $B_1=I$, the period
integrals on the core generators are
$\Lambda_1$ and the cycle generator has
period $b_0+v_0\log t$. Equation \eqref{cdp:eq:C.HC.7}
lists all generators, so the exact lattice is
\[
                 \Lambda(t)=
                    \Lambda_1+\mathbb Z(b_0+v_0\log t).
\tag{C.HC.10}\hypertarget{tag-00LP}{}\label{cdp:eq:C.HC.10}
\]
Its discreteness and full real rank follow
from the preceding transitive-action proof,
rather than being assumed. A turn of $\log t$
changes the last vector by $2\pi i v_0\in\Lambda_1$,
as follows directly from the globally defined
edge maps. Hence \eqref{cdp:eq:C.HC.8}--\eqref{cdp:eq:C.HC.10} descend
to the actual punctured-disc family. This proves
the sufficient half and records every smoothing
translation parameter.

\par\medskip\noindent\textbf{One and two components.}
A single transition in \eqref{cdp:eq:C.HC.5} cannot be the
identity, because its upper-left entry is zero.
Thus the specified self-identification of the
two distinct axes cannot give a torus smoothing.

\item For two components multiply the two matrices:
\[
 L_2L_1=
 \begin{pmatrix}
   -c_2&c_2k_1\\
   -k_2&-c_1+k_2k_1
 \end{pmatrix}.
\tag{C.HC.11}\hypertarget{tag-00LQ}{}\label{cdp:eq:C.HC.11}
\]
It is the identity precisely when
$c_1=c_2=-1$ and $k_1=k_2=0$.
Consequently $L_1=L_2=-S$, where $S$ exchanges
the two coordinates, and
$\Lambda_2=-S\Lambda_1$.
Exponentiation gives
$\Gamma_2=(S\Gamma_1)^{-1}=S\Gamma_1$,
since these are groups. The actual coordinate
exchange $(z,w)\mapsto(w,z)$ therefore descends
to a biholomorphism $H_1\to H_2$, including
their finite diagonal subgroups.

\item After this coordinate exchange, the two
elliptic identifications have the form
$z\mapsto a/z$ and $w\mapsto b/w$.
The actual diagonal automorphism
$(Z,W)\mapsto(a^{-1}Z,b^{-1}W)$ of the second
Hopf copy commutes with its complete deck group
and makes both identifications inversion.
It follows that the reduced central surface
is isomorphic to exactly the support in
\ref{cdp:st:C.2.UN.41}, or its free finite quotient in
\ref{cdp:st:C.2.UN.43}. Those constructions give the asserted
compact realizations and retain this entire
reduced surface. The constants in $\theta_i$
are recorded by \eqref{cdp:eq:C.HC.9}--\eqref{cdp:eq:C.HC.10}; the
compact realization assertion here identifies
the reduced support, while the explicit finite
plumbing realizes each specified local choice.

\end{enumerate}
\end{proof}






\begin{proposition}[The full compatible cycle is realized over a compact rational base]\hypertarget{tag-01DS}{}\label{cdp:st:C.2.UN.48}
\StatementRecord{cdp:st:C.2.UN.48}{proved}{01DS}{\SeriesRef{cdp:st:C.2.UN.46}, \SeriesRef{cdp:eq:C.HG.3}, \SeriesRef{cdp:eq:C.HG.5}, \SeriesRef{cdp:st:C.2.UN.40}, \SeriesRef{cdp:st:C.2.UN.43}, \SeriesRef{cdp:eq:C.HG.7}, \SeriesRef{cdp:st:C.2.UN.41}, \SeriesRef{cdp:eq:C.HG.6}, \SeriesRef{cdp:eq:C.HG.12}, \SeriesRef{cdp:eq:C.HG.10}, \SeriesRef{cdp:eq:C.HG.13}, \SeriesRef{cdp:eq:C.HG.14}, \SeriesRef{cdp:eq:C.HG.11}}

Take the specified cycle of primary or secondary diagonal
Hopf surfaces of algebraic dimension zero in \ref{cdp:st:C.2.UN.46},
with its specified elliptic gluing isomorphisms.
Suppose its dual-normal bundles are isomorphic to the
trivial bundle, and compute the integral edge product
from those actual isomorphisms as in that statement.
Every cycle whose product is the identity occurs,
with precisely its given components and gluing maps,
as a reduced SNC fibre of a holomorphic map
\[
                 f:Y\longrightarrow\PP^1
\]
from a smooth compact complex threefold with connected
two-torus general fibre. Every reduced fibre of this
map has Euler characteristic zero.

In particular the necessary criterion of \ref{cdp:st:C.2.UN.46}
is sufficient for compact occurrence of the whole
reduced support, for arbitrary cycle length and for
secondary components as well. The construction makes
specified choices of the normal-product constants;
it preserves the entire reduced surface, while
recording those choices and the actual punctured-family
comparison below.

For example there is a three-component compact
realization with primary Hopf multipliers
\[
 \begin{gathered}
 (e^{-1},e^{-\sqrt2}),\qquad
 (e^{-\sqrt2},e^{-1-\sqrt2}),\qquad
 (e^{-1-\sqrt2},e^{-1}),\\
 L_1=\begin{pmatrix}0&-1\\-1&-1\end{pmatrix},\qquad
 L_2=\begin{pmatrix}0&1\\-1&1\end{pmatrix},\qquad
 L_3=\begin{pmatrix}0&-1\\-1&1\end{pmatrix}.
 \end{gathered}
 \tag{C.HG.1}\hypertarget{tag-00LZ}{}\label{cdp:eq:C.HG.1}
\]
Here $L_3L_2L_1=I$ and the actual monodromy period
on the common open group is $(t^{-2},t^2)$.
Thus the longer-cycle occurrence assertion includes
an explicit case beyond the two-component inversion
construction.
\end{proposition}



\begin{proof}\leavevmode\par



 
\begin{enumerate}[label={$\langle1\rangle$\arabic*.},leftmargin=0pt,labelsep=.5em]
\item \par\medskip\noindent\textbf{The common lattice and the sign of the normal periods.}
Retain the exact logarithmic coordinates of \ref{cdp:st:C.2.UN.46}.
The product identity permits constant matrices $P_i$
identifying every open group with one common group
\[
                  G=\C^2/\Lambda,\qquad \Lambda\simeq\Z^3,
                  \qquad P_{i+1}L_i=P_i.
\tag{C.HG.2}\hypertarget{tag-00M8}{}\label{cdp:eq:C.HG.2}
\]
Take indices cyclically. Write $n_i=P_i(2\pi i,0)$
for the outgoing normal meridian. Since
$L_i(1,0)=(0,-1)$, the incoming meridian on the
$i$th component is $-n_{i-1}$, and
\[
 P_i=\left(\frac{n_i}{2\pi i},
                   -\frac{n_{i-1}}{2\pi i}\right),
 \qquad
 2\pi i v_0=-\sum_i n_i=:u\in\Lambda.
\tag{C.HG.3}\hypertarget{tag-00M9}{}\label{cdp:eq:C.HG.3}
\]
The last equality follows by adding the actual
edge translations $P_{i+1}(d_i,e_i+\log t)$.
It retains their signs, rather than choosing a
meridian sign afterwards.


\begin{enumerate}[label={$\langle2\rangle$\arabic*.},leftmargin=1.3em,labelsep=.5em]
\item  Put $W=\operatorname{span}_{\R}\Lambda$, and choose
a nonzero real linear functional $\ell:\C^2\to\R$
with kernel $W$. On the $i$th Hopf surface choose
its actual contracting generator with multipliers
$\alpha_i,\beta_i$, including when a finite diagonal
subgroup is also present. Its common logarithm
\[
 c_i=P_i(\log\alpha_i,\log\beta_i)
       =\frac{\log\alpha_i}{2\pi i}n_i
          -\frac{\log\beta_i}{2\pi i}n_{i-1}
\]
belongs to $\Lambda$. Applying $\ell$ and using
$\ell(n_i)=\ell(n_{i-1})=0$ gives the exact real
identity
\[
 \log|\alpha_i|\,\ell(i n_i)
       =\log|\beta_i|\,\ell(i n_{i-1}).
\tag{C.HG.4}\hypertarget{tag-00MA}{}\label{cdp:eq:C.HG.4}
\]
Both logarithms are strictly negative.
Neither displayed value of $\ell$ is zero: if
both vanished, $\ell$ would vanish on the complex
span of $n_i,n_{i-1}$, which is $\C^2$ by
\eqref{cdp:eq:C.HG.3}. The identity would make both zero if
one were zero. It follows successively around
the cycle that all $\ell(i n_i)$ have the same
nonzero sign. Consequently
\[
                       \ell(iu)=-\sum_i\ell(i n_i)\ne0.
\tag{C.HG.5}\hypertarget{tag-00MB}{}\label{cdp:eq:C.HG.5}
\]
In particular $u\ne0$, and $\C u\cap W=\R u$.


\item  Write $u=d p$, where $p$ is a primitive element
of the free lattice $\Lambda$ and $d>0$.
The real quotient map
$W/\R p\to\C^2/\C p$ is an isomorphism, by
\eqref{cdp:eq:C.HG.5}. Therefore the image of $\Lambda$ in
$\C^2/\C p$ is a discrete rank-two lattice.
Its quotient is an actual elliptic curve.
Extend $p$ to an integral basis $p,\lambda,c$
of $\Lambda$. Choose a linear quotient
coordinate $X$ with $X(\lambda)=2\pi i$.
The two projected periods are real-linearly
independent, so $\Re X(c)\ne0$; reverse $c$
if necessary to make it negative. Choose a
second complex linear coordinate $Z$ with
$Z(p)=2\pi i$ and $Z(\lambda)=0$. These
conditions determine it because $p,\lambda$
are complex-linearly independent. In these
actual linear coordinates the common lattice
is
\[
 \Lambda=\left\langle
     (0,2\pi i),\ (2\pi i,0),\
     (\log q_0,\log a)\right\rangle_{\Z},
 \qquad 0<|q_0|<1,\qquad a\in\C^*,
 \qquad u=d(0,2\pi i).
\tag{C.HG.6}\hypertarget{tag-00MC}{}\label{cdp:eq:C.HG.6}
\]
The logarithms in this formula are the actual
chosen coordinates of $c$.


\item  Here $a\ne1$. Indeed, if $a=1$, subtract an
integral multiple of the first period from $c$
so that its second coordinate is zero. The map
\[
 \C/\langle2\pi i,\log q_0\rangle_{\Z}
          \longrightarrow G,\qquad [X]\longmapsto[(X,0)]
\tag{C.HG.7}\hypertarget{tag-00MD}{}\label{cdp:eq:C.HG.7}
\]
is well defined and injective: equality in $G$
forces the coefficient of $(0,2\pi i)$ to vanish,
leaving exactly the two periods in its source.
It is an immersion and, by compactness, an embedded
elliptic curve. But $G$ is the complement of the
two axes in any of the prescribed Hopf surfaces.
The curve calculation of \ref{cdp:st:C.2.UN.40} and \ref{cdp:st:C.2.UN.43}
says that these are all its compact curves.
The image of \eqref{cdp:eq:C.HG.7} contradicts that calculation.
This proves the precise specialization needed
for the pencil construction below.


\end{enumerate}

\item \par\medskip\noindent\textbf{A compact elliptic suspension with the required limiting periods.}
The cubic construction in \ref{cdp:st:C.2.UN.41} works for
every specified $a\in\C^*\setminus\{1\}$:
its choice of the smooth nodal point uses no
inequality on $|a|$. Explicitly use
\[
 \begin{gathered}
 y^2z=x^2(x+z),\qquad O=[0:1:0],\\
 P_0=\left(\frac{4a}{(1-a)^2},
                  \frac{4a(1+a)}{(1-a)^3}\right)
 \end{gathered}
\tag{C.HG.8}\hypertarget{tag-00ME}{}\label{cdp:eq:C.HG.8}
\]
in the affine chart $z=1$, whose smooth group
coordinate is $a$. The point is distinct from
$O$ and from the node. Choose a second cubic
through $O,P_0$, transverse at nine distinct
smooth points of the nodal cubic and missing
its node. The jet-surjectivity and nonempty
open conditions in \ref{cdp:st:C.2.UN.41} give this choice.
Blowing up the nine actual base points gives
a compact rational elliptic surface with
sections $O,P$ and the selected $I_1$ fibre.


\begin{enumerate}[label={$\langle2\rangle$\arabic*.},leftmargin=1.3em,labelsep=.5em]
\item  Make the global base change $\PP^1_s\to\PP^1$
which is $s\mapsto s^d$ at that value.
Normalize, resolve and take the relatively
minimal elliptic model $J\to\PP^1_s$.
Near the chosen node the equation before
resolution is $xy=s^d$. Its resolution has
charts $(u_j,v_j)$, $0\le j<d$, with exact maps
\[
 x=u_j^{j+1}v_j^j,\qquad
 y=u_j^{d-j-1}v_j^{d-j},\qquad s=u_jv_j.
\]
Adjacent charts glue by
$u_{j+1}=v_j^{-1}$, $v_{j+1}=u_jv_j^2$.
These maps are isomorphisms off the node and
insert the chain of $d-1$ projective lines
between its two branches; each displayed
central equation is $u_jv_j=0$.
Together with the old component
this gives an $I_d$ fibre. The sections avoid
these charts and retain their given values.
At all other values these operations are
proper vertical modifications of a compact
elliptic surface. The section translation by
$P$ extends to an automorphism of this
relatively minimal model, by the exact graph
and relative-canonical-nef argument of
\ref{cdp:st:C.2.UN.41} and \ref{cdp:st:C.2.UN.43}.


\item  The analytic Tate comparison at this fibre
gives, on a sufficiently small punctured disc,
\[
       J_s=\C^*/Q(s)^{\Z},\qquad
       \operatorname{ord}_0Q=d,\qquad P_s=[a(s)],
       \qquad a(0)=a.
\tag{C.HG.9}\hypertarget{tag-00MF}{}\label{cdp:eq:C.HG.9}
\]
The lift $a(s)$ is a holomorphic unit, since
the pulled-back section meets the smooth
identity component. Choose the holomorphic
$d$th root of the unit $Q(s)/s^d$ and use the
actual coordinate
$t=s\bigl(Q(s)/s^d\bigr)^{1/d}$.
It has nonzero derivative at zero and makes
$Q=t^d$ exactly. Denote the section lift in
this coordinate by $a(t)$, and fix the
holomorphic logarithm with value $\log a$.


\item  Now take the actual suspension
\[
 Y^{\rm std}=
  (\C^*_b\times J)/
      \langle(b,y)\longmapsto(q_0 b,y+P)\rangle
          \longrightarrow\PP^1_s .
\tag{C.HG.10}\hypertarget{tag-00M0}{}\label{cdp:eq:C.HG.10}
\]
The action is free and properly discontinuous
because its first coordinate has modulus
multiplier $|q_0|<1$. A closed annulus in
$\C^*$ times compact $J$ meets every orbit,
so the quotient is compact and smooth.
Its general connected fibres are two-tori.
Near the distinguished value their exact
logarithmic periods are
\[
 \Lambda(t)+\Z(0,d\log t),\qquad
 \Lambda(t)=
 \left\langle(0,2\pi i),(2\pi i,0),
                  (\log q_0,\log a(t))\right\rangle_{\Z}.
\tag{C.HG.11}\hypertarget{tag-00M1}{}\label{cdp:eq:C.HG.11}
\]
This is the displayed suspension quotient,
not just equality of characteristic polynomials.


\end{enumerate}

\item \par\medskip\noindent\textbf{Transporting the prescribed Hopf surfaces and the whole cycle.}
View each original $n_i$ and contracting
$c_i$ as its actual integral combination
of the three basis periods in \eqref{cdp:eq:C.HG.6}.
Use those same integer coefficients in
$\Lambda(t)$ to define $n_i(t),c_i(t)$.
All integral relations are retained exactly.
In particular $\sum_i n_i(t)=-d(0,2\pi i)$.
Set
\[
 \begin{gathered}
 P_i(t)=\left(\frac{n_i(t)}{2\pi i},
                    -\frac{n_{i-1}(t)}{2\pi i}\right),\\
 \Lambda_i(t)=P_i(t)^{-1}\Lambda(t),
 \quad L_i(t)=P_{i+1}(t)^{-1}P_i(t).
 \end{gathered}
\tag{C.HG.12}\hypertarget{tag-00M2}{}\label{cdp:eq:C.HG.12}
\]
These matrices are holomorphic and invertible
after shrinking the disc. The first column
of $L_i(t)$ is $(0,-1)$, so it has exactly the
form
$\left(\begin{smallmatrix}0&c_i^*(t)\\-1&k_i^*(t)\end{smallmatrix}\right)$
required for the normal-product map.
Its cycle product is the identity.


\begin{enumerate}[label={$\langle2\rangle$\arabic*.},leftmargin=1.3em,labelsep=.5em]
\item  The group $\Lambda_i(t)/(2\pi i\Z^2)$ is
the same abstract finitely generated group
as at zero, by the fixed integral relations.
Its selected free generator is the actual
diagonal contraction with logarithms
$P_i(t)^{-1}c_i(t)$. Both moduli remain
strictly between zero and one on a smaller
disc. Every finite torsion generator has
a relation $n v=A(2\pi i,0)+B(0,2\pi i)$
with the same integers as at zero. Therefore
its two roots of unity and their coprimality
conditions remain unchanged. These are
actual holomorphic families of the prescribed
primary or secondary Hopf quotients. The
contractions are uniformly contracting over
a closed smaller disc, and their fixed
finite groups remain free. Thus the
quotients of $(\C^2\setminus\{0\})\times\Delta$
are smooth proper families, with precisely
the original surfaces at zero.


\item  Let $d_i$ be the constants specifying the
original elliptic gluing maps in the
coordinates of \ref{cdp:st:C.2.UN.46}. Keep these constants.
Their normal multipliers have freely
adjustable nonzero constants, whose logarithms
we denote by $e_i(t)$. The full affine cycle
translation in common coordinates is
\[
 \sum_iP_{i+1}(t)
        \binom{d_i}{e_i(t)+\log t}
       =b_0(t)+(0,d\log t).
\tag{C.HG.13}\hypertarget{tag-00M3}{}\label{cdp:eq:C.HG.13}
\]
Two adjacent coefficient vectors of the
unknowns $e_i(t)$ are
$-n_i(t)/(2\pi i)$ and $-n_{i+1}(t)/(2\pi i)$.
They are complex-linearly independent by
\eqref{cdp:eq:C.HG.12}. Fix all the other $e_i(t)$, and
solve the resulting invertible two-by-two
linear system holomorphically to obtain
\[
                              b_0(t)=0.
\tag{C.HG.14}\hypertarget{tag-00M4}{}\label{cdp:eq:C.HG.14}
\]
Only the normal-product constants change.
The central surfaces and every given
elliptic gluing map remain exactly the
prescribed ones.


\item  Use the actual varying normal line bundles
and these affine edge maps to perform the
finite plumbing of \ref{cdp:st:C.2.UN.46}, with $xy=t$
in each neck. The flat metrics, disjoint
axis neighbourhoods and compact cores may
be chosen uniformly on a smaller parameter
disc; the strict inequalities used in that
proof persist by continuity. The result
$\mathcal Y\to\Delta_t$ is a smooth proper
threefold whose reduced central fibre is
the given cycle, with every component
of multiplicity one.


\item  There is an actual punctured-family
biholomorphism to \eqref{cdp:eq:C.HG.10}. To see the map
without suppressing affine data, choose one
core as origin in the common logarithmic
coordinates. On consecutive cores use
$P_i(t)$ and subtract the accumulated edge
translations from \eqref{cdp:eq:C.HG.13}. Use the same
formulas on the necks, where $xy=t$.
They agree on successive overlaps. On the
last overlap their difference is the
period $(0,d\log t)$ by \eqref{cdp:eq:C.HG.14};
changing logarithm branches changes them
by $\Lambda(t)$. Hence they define a
holomorphic map to the quotient with
periods \eqref{cdp:eq:C.HG.11}. Fibrewise this is the
closed-coframe comparison of \ref{cdp:st:C.2.UN.46},
which is a biholomorphism: its periods
are all the generators of the actual
fundamental group. Its differential is
invertible in the fibre directions and
it is the identity on the base, so it
is a local biholomorphism of total
spaces as well. The fibrewise inverses
therefore give the actual holomorphic
inverse over the punctured disc.


\end{enumerate}

\item \par\medskip\noindent\textbf{Compact replacement and Euler numbers.}
Replace the distinguished tube in
$Y^{\rm std}$ by $\mathcal Y$, using this
biholomorphism on a punctured annulus.
More precisely remove a closed smaller
disc from the base of the first family
and glue the second over a larger disc.
Their annular overlap is the displayed
isomorphism. Properness of both families
and two nested radii give a Hausdorff
gluing: at either boundary of the annulus
only one of the two pieces contains the
limiting base point. Closed smaller
pieces in these two proper families
cover the resulting space, proving
compactness. Smoothness and the
holomorphic map to the same $\PP^1_s$
are preserved. Its fibres are connected
and its general fibre is the specified
suspension two-torus.


\begin{enumerate}[label={$\langle2\rangle$\arabic*.},leftmargin=1.3em,labelsep=.5em]
\item  The new central support has Euler number
zero, since every Hopf component and every
double elliptic curve has Euler number
zero and there are no triple points.
For every other base value the support
is unchanged from \eqref{cdp:eq:C.HG.10}. Its map to
the actual elliptic curve $\C^*/q_0^{\Z}$
is a locally trivial bundle with fibre
the reduced curve $J_s^{\rm red}$ and
transition the automorphism induced by
$P$. A finite triangulation, or the
Euler characteristic of its finite
local-system cellular complex, gives
Euler number
$e(\C^*/q_0^{\Z})e(J_s^{\rm red})=0$.
Thus every reduced fibre has Euler
number zero, as asserted.


\item  Finally \eqref{cdp:eq:C.HG.1} is verified by taking,
in the common angular lattice,
\[
 n_1=2\pi i(1,0),\quad
 n_2=2\pi i(1,-1),\quad
 n_3=2\pi i(0,-1),\quad
 c=(-1,-\sqrt2).
\tag{C.HG.15}\hypertarget{tag-00M5}{}\label{cdp:eq:C.HG.15}
\]
The three chosen contracting logarithms
are $P_1^{-1}c$, $-P_2^{-1}c$ and
$-P_3^{-1}c$, exactly the pairs in
\eqref{cdp:eq:C.HG.1}. Their ratios are irrational,
so the three Hopf surfaces have
algebraic dimension zero. The edge
maps in multiplicative coordinates,
with constant factors one, are
\[
 \begin{gathered}
 (z,w)\mapsto(w^{-1},t z^{-1}w^{-1}),\\
 (z,w)\mapsto(w,t z^{-1}w),\\
 (z,w)\mapsto(w^{-1},t z^{-1}w).
 \end{gathered}
\tag{C.HG.16}\hypertarget{tag-00M6}{}\label{cdp:eq:C.HG.16}
\]
Their composite is $(z,w)\mapsto(t^{-2}z,t^2w)$.
The exact integral matrices, negative
contraction logarithms, primitive
period divisor $d=2$ and the two-constant
translation solution are checked in
\path{code/D_hopf_cycle_matrices.py}
and its saved output.


\item  The finite-secondary divisibility is also
explicit. For a repeated two-component
cycle of length $2k$ with finite weights
$(a,b)$ modulo $n$, the common angular
lattice is
$\Z^2+\Z(a,b)/n$, and the period vector is
$k(-1,1)$ in units of $2\pi i$.
Its intersection with $\R(-1,1)$ is
$\Z(-1,1)/\gcd(n,a+b)$: the sum of the
coordinates requires $n\mid j(a+b)$ for
a representative $j(a,b)/n$, and the
coprimality of $a$ to $n$ then gives
exactly those fractions on this line.
Consequently the precise positive divisor is
\[
                    d=k\gcd(n,a+b).
\tag{C.HG.17}\hypertarget{tag-00M7}{}\label{cdp:eq:C.HG.17}
\]
The sign is fixed by the actual outgoing
vectors $(1,0)$ and $(0,-1)$.
For example weights $(1,4)$ modulo $5$
give $d=5$ for two components.
In the three-component example, the
finite weights $(1,2)$ modulo $5$ become
$(1,2),(3,2),(2,4)$ on its three actual
Hopf quotients. They are all free, the
common period still has divisor $2$,
and the same construction gives this
compact secondary three-cycle.
The saved computation checks 12,180
repeated-cycle cases and 1,376 free
secondary three-cycles by exact
Hermite normal forms.



\end{enumerate}

\end{enumerate}
\end{proof}






\begin{proposition}[Exact nonlinear Hopf star criterion and compact realization]\hypertarget{tag-01DT}{}\label{cdp:st:C.2.UN.49}
\StatementRecord{cdp:st:C.2.UN.49}{proved}{01DT}{\SeriesRef{cdp:eq:C.NS.1}, \SeriesRef{cdp:eq:C.NS.3}, \SeriesRef{cdp:eq:C.NS.2}, \SeriesRef{cdp:eq:C.NS.5}, \SeriesRef{cdp:eq:C.NS.4}, \SeriesRef{cdp:eq:C.NS.6}, \SeriesRef{cdp:eq:C.NS.9}, \SeriesRef{cdp:eq:C.NS.10}, \SeriesRef{cdp:eq:C.NS.8}, \SeriesRef{cdp:eq:C.NS.11}, \SeriesRef{cdp:eq:C.NS.15}, \SeriesRef{cdp:eq:C.NS.18}, \SeriesRef{cdp:eq:C.NS.20}, \SeriesRef{cdp:eq:C.NS.19}}

Let $S$ be a specified compact SNC surface with components
$P,H_1,\ldots,H_n$, $n\ge1$, and no triple points. Suppose
$q:P\to B$ is a connected fixed-Jacobian elliptic quasi-bundle:
its only nonsmooth scheme fibres are multiples of smooth
elliptic curves, and its integral $H^1$ monodromy is the
identity. Each $H_i=H_{r_i}(\beta_i)$ is a nonlinear primary
Hopf surface as in (C.NH.1). It meets $P$ along its unique
axis, identified by a specified holomorphic isomorphism
with a distinct reduced fibre $J_i$ of $q$. There are no
other intersections.

The following assertions are equivalent.

\begin{enumerate}
\item $S$ is the reduced full central fibre of a proper
one-parameter degeneration of two-tori with smooth total space.
\item $S$ is the reduced full special fibre of a smooth compact
threefold over $\PP^1$, all of whose other fibres are smooth
two-tori.
\item The normal lines on every double curve are dual;
$B$ is elliptic; all $r_i$ equal one integer $r$; the full
marked fibres are $rJ_i$ and there are no other multiple
fibres; and the actual inverse logarithmic transform
$T_0\to B$ defined in \eqref{cdp:eq:C.NS.1}--\eqref{cdp:eq:C.NS.3} has zero integral
Chern vector in $H^2(B,\Lambda_F)$.
\end{enumerate}

Here $F$ is the fixed elliptic Jacobian and $\Lambda_F$ its
integral period lattice. The inverse transform retains
the entire analytic torsor of $P$ and the given axis
identifications. It is unique up to the explicit product
translations described below. Thus the criterion concerns
the specified surface $S$, including its component
moduli and gluing maps. Every fibre support in assertion
(2) has Euler number zero. No extension of the original
smoothing's elliptic quotient is assumed.
\end{proposition}



\begin{proof}\leavevmode\par



 
\begin{enumerate}[label={$\langle1\rangle$\arabic*.},leftmargin=0pt,labelsep=.5em]
\item \par\medskip\noindent\textbf{The actual inverse transform.}
We first define the operation whenever the marked
multiplicity is $r_i$ and the two normal lines are dual.
Write the Hopf axis as $\C^*/\beta_i^{r_i\Z}$, with
coordinate $z$. Its normal has multiplier $\beta_i$.
The matching neighborhood in $P$ has the precise model
\[
 p=x^{r_i}z,\qquad
 (z,x)\sim(\beta_i^{r_i}z,\beta_i^{-1}x).
                                                        \tag{C.NS.1}\hypertarget{tag-00TA}{}\label{cdp:eq:C.NS.1}
\]
Indeed, normalize the $r_i$-th root of $q$ at this fibre.
The result is a smooth principal $F$-bundle over a disc,
hence is trivial there. The free root action is an
order-$r_i$ translation on $F$, together with the
primitive normal character. A finite averaging gauge
makes that translation independent of the root
parameter. The dual-normal isomorphism identifies this
connected cyclic cover with the one coming from the
Hopf axis. In multiplicative coordinates the cover is
$z=w^{r_i}$, $x=u/w$, with deck $w\mapsto\beta_i w$;
its quotient is exactly \eqref{cdp:eq:C.NS.1}, with $p=u^{r_i}$.
This also proves that the normal of a smooth
multiple fibre has exact order its multiplicity.


\begin{enumerate}[label={$\langle2\rangle$\arabic*.},leftmargin=1.3em,labelsep=.5em]
\item  On the punctured base disc the map
\[
       (z,x)\longmapsto(p,w)=(x^{r_i}z,x^{-1}),
       \qquad z=pw^{r_i},
       \qquad [w]\in F_i=\C^*/\beta_i^{\Z}
                                                        \tag{C.NS.2}\hypertarget{tag-00TL}{}\label{cdp:eq:C.NS.2}
\]
is an actual product trivialization. Replace \eqref{cdp:eq:C.NS.1}
by $F_i\times\Delta_p$ through \eqref{cdp:eq:C.NS.2}, at every
marked point. The original torsor on the complement
and these exact transition maps define
\[
                             q_0:T_0\longrightarrow B.
                                                        \tag{C.NS.3}\hypertarget{tag-00TN}{}\label{cdp:eq:C.NS.3}
\]
When there are no unmarked multiple fibres, this is a
smooth principal $F$-bundle. All identifications of
$F_i$ with $F$ are the ones furnished by the given
normal covers and the original torsor.


\item  The construction is independent of the permitted local
choices. Two dual-normal isomorphisms differ by a
nonzero holomorphic function on the compact elliptic
axis and hence by a constant. This changes $w$ by a
constant translation of $F_i$, which extends over the
product disc. An automorphism of a root-trivialized
principal $F$-bundle that extends across its central
curve is a holomorphic translation there, followed by
a fixed linear automorphism. If it preserves the
specified normal cover, that translation is invariant
under the root action and descends to a holomorphic
translation on the ordinary disc. It again extends
over the product filling. The specified axis map
fixes the remaining linear identification. Thus
\eqref{cdp:eq:C.NS.3} retains the full holomorphic torsor class;
it is not defined merely by a rational Chern number.


\end{enumerate}

\item \par\medskip\noindent\textbf{The integral boundary matrix and the smoothing topology.}
We compute all meridians before applying a smoothing
argument. In \eqref{cdp:eq:C.NS.1}, let $a$ be the positive angular
loop of $z$, let $b$ be its deck loop, and let
$\mu_P$ be the positive angular loop of $x$. On the
punctured product \eqref{cdp:eq:C.NS.2}, let $\delta$ be the
positive base loop with $w$ fixed, $e$ the angular
loop of $w$, and $h$ its $\beta_i$-deck loop. Then
\[
 \begin{pmatrix}a&b&\mu_P\end{pmatrix}_{(\delta,e,h)}
       =
 \begin{pmatrix}
  1&0&r_i\\
  0&0&-1\\
  0&1&0
 \end{pmatrix};
 \qquad
 a=\delta,\quad b=h,\quad \mu_P=r_i\delta-e.
                                                        \tag{C.NS.4}\hypertarget{tag-00TO}{}\label{cdp:eq:C.NS.4}
\]
The matrix is integral of determinant one. For example,
holding $p$ fixed while $w$ winds once gives
$x=w^{-1}$ and $z=pw^{r_i}$, proving
$e=r_i a-\mu_P$. The formulas retain the normal
character, not just its order.


\begin{enumerate}[label={$\langle2\rangle$\arabic*.},leftmargin=1.3em,labelsep=.5em]
\item  The nonlinear Hopf pair is smoothly equivalent, for
this boundary calculation, to the diagonal pair with
contraction $(\beta_i^{r_i}z,\beta_i w)$. One exact
comparison on the complement is
\[
 U=\frac{z}{w^{r_i}},\qquad
 V=U-\frac{\beta_i^{-r_i}}{\log|\beta_i|}\log|w|.
                                                        \tag{C.NS.5}\hypertarget{tag-00TP}{}\label{cdp:eq:C.NS.5}
\]
The deck map sends $U$ to $U+\beta_i^{-r_i}$ and
$w$ to $\beta_iw$, so $V$ is invariant. This gives
$H_i\setminus J_i\simeq\C_V\times F_i$ smoothly.
The map
$z\mapsto z-\beta_i^{-r_i}w^{r_i}\log|w|/\log|\beta_i|$
extends continuously as the identity on the axis and
preserves orientation; it is smooth on every boundary
torus. Equivalently, vary the resonant coefficient from
one to zero in the contraction and use a uniform
compact fundamental annulus. This is a smooth proper
family of pairs, so the same integral boundary
identification follows by relative Ehresmann.


\item  Let $\mu_H$ be the Hopf normal meridian. The
smoothing equation is $xw=t$, so
$\mu_H=-\mu_P$. The regular fibre cycles on the
Hopf side are
\[
                   e_H=r_i a+\mu_H=e,\qquad h_H=b=h.
                                                        \tag{C.NS.6}\hypertarget{tag-00TQ}{}\label{cdp:eq:C.NS.6}
\]
The complementary solid torus $F_i\times D^2$ in
\eqref{cdp:eq:C.NS.5} kills $a=\delta$. Consequently replacing
the tubular neighborhood in $P$ by the Hopf
complement is precisely the inverse logarithmic
transform \eqref{cdp:eq:C.NS.2}, with its indicated integral
framing.


\item  For completeness, this calculation applies to every
actual SNC smoothing of the specified support.
Remove small tubular neighborhoods of the disjoint
double curves. On each remaining compact component
core the original map is a submersion through zero;
Ehresmann with boundary identifies that core with
its nearby copy. On the collars the local equations
are exactly $xw=t$ after absorbing a holomorphic
unit. The product of the two normal lines is
trivialized by this same function. Two such
trivializations along a compact double curve differ
by a constant in $\C^*$, whose rotation is isotopic
to the identity. Thus the nearby four-manifold is
the indicated tubular plumbing, and \eqref{cdp:eq:C.NS.4}--\eqref{cdp:eq:C.NS.6}
identify it with the underlying four-manifold of
$T_0$. Under this identification, the oriented
generator of $H_2(H_i\setminus J_i,\Z)$ is the
ordinary fibre class of $q_0$, for every $i$.
The identifications between the normal covers are
holomorphic elliptic isomorphisms, hence have degree
one and preserve that orientation. In particular all
these transported core cycles are the same integral
two-cycle in the nearby fibre, even when their
primitive angular directions in $\Lambda_F$ differ.


\item  Here is an explicit global version of the Chern test.
Trivialize the principal $F$-bundle on the punctured
base and compare that trivialization with \eqref{cdp:eq:C.NS.2}
around the $i$-th disc. Let $\nu_i\in\Lambda_F$ be
the winding of its translation transition, with
positive disc-boundary orientation. Then
\[
 c(T_0)=\sum_i\nu_i,\qquad
 \mu_{P,i}=r_i\delta_{\rm ext}
                    +(r_i\nu_i-e_i).
                                                        \tag{C.NS.7}\hypertarget{tag-00TR}{}\label{cdp:eq:C.NS.7}
\]
The first equality is the integral obstruction to
extending these boundary trivializations over the
discs. It can alternatively be used as the sign
convention for the Chern vector. The second equality
is \eqref{cdp:eq:C.NS.4} after $\delta_i=\delta_{\rm ext}+\nu_i$.
Thus the rational Seifert vector in this convention
is $\sum_i(\nu_i-e_i/r_i)$, and the inverse transform
adds exactly $\sum_i e_i/r_i$. The holomorphic
transition functions themselves, rather than only
their windings, remain in \eqref{cdp:eq:C.NS.3}.


\end{enumerate}

\item \par\medskip\noindent\textbf{Canonical orders, the elliptic base, and equal resonance.}
Suppose first that $S$ occurs as in assertion (1).
The equation $xw=t$ proves duality of the normal
lines. Their exact torsion orders then give the
marked multiplicities $r_i$.


\begin{enumerate}[label={$\langle2\rangle$\arabic*.},leftmargin=1.3em,labelsep=.5em]
\item  Choose a generator of the rank-one direct image of
the relative canonical line over the disc. Its
evaluation has a vertical divisor
$k_PP+\sum_i k_iH_i$. It is nowhere zero on every
nonzero torus fibre. Adjunction on the Hopf component
and $P|_{H_i}=J_i$, $H_i|_{H_i}=-J_i$ give
\[
 (k_P-k_i)J_i=K_{H_i}+J_i=-r_iJ_i,
       \qquad k_i-k_P=r_i.
                                                        \tag{C.NS.8}\hypertarget{tag-00TS}{}\label{cdp:eq:C.NS.8}
\]
The last equality is an equality of integers:
$\OO_{H_i}(J_i)$ has infinite order, since a
positive effective divisor on a compact complex
surface cannot have a trivial positive tensor power.
Adjunction on $P$ now gives the actual line equality
\[
                         K_P=\OO_P\Bigl(\sum_i(r_i-1)J_i\Bigr).
                                                        \tag{C.NS.9}\hypertarget{tag-00TT}{}\label{cdp:eq:C.NS.9}
\]


\item  For the given quasi-bundle the root charts give
$q_*K_P=K_B$ and
$K_P=q^*K_B+\sum_{\text{all multiple fibres}}(m-1)J$.
Indeed, the local form is the base differential
wedged with the invariant elliptic differential;
on the root chart its exact zero is $(m-1)J$,
and taking the fibrewise constant Laurent
coefficients gives the displayed pushforward.
Equation \eqref{cdp:eq:C.NS.9} implies $h^0(K_P)\ge1$, so
$g(B)\ge1$. Subtracting the marked terms from
the two canonical formulas gives
\[
 q^*K_B+\sum_{\text{unmarked multiple fibres}}(m-1)J=0
                    \quad\hbox{as a line bundle}.       \tag{C.NS.10}\hypertarget{tag-00TB}{}\label{cdp:eq:C.NS.10}
\]
For $g(B)\ge2$ a nonzero canonical section on $B$
has a nonzero effective divisor, and its pullback
together with the other terms would be a nonzero
effective divisor linearly equivalent to zero on
the compact surface. This is impossible. Therefore
$B$ is elliptic, and the same argument in \eqref{cdp:eq:C.NS.10}
removes every unmarked multiple fibre.


\item  Let $\omega_t$ be the evaluated relative two-form
divided by $t^{k_P}$. On a compact core of $H_i\setminus J_i$
its leading term is $t^{r_i}\eta_i$, where
\[
 \eta_i=c_i\,\frac{dz\wedge dw}{w^{r_i+1}},
                         \qquad c_i\ne0.                \tag{C.NS.11}\hypertarget{tag-00TC}{}\label{cdp:eq:C.NS.11}
\]
To verify this exact pole order, write the total
canonical evaluation in a double chart as
$x^{k_i}w^{k_P}$ times a holomorphic unit and a
coordinate three-form, then divide by $dt$ with
$t=xw$. Equation \eqref{cdp:eq:C.NS.8} gives \eqref{cdp:eq:C.NS.11}.
The resulting coefficient on the entire Hopf
component is a nowhere-zero section of
$K_{H_i}((r_i+1)J_i)=\OO_{H_i}$, so it is the
nonzero constant $c_i$.


\item  On the logarithmic cover of $H_i\setminus J_i$,
put $U=z/w^{r_i}$ and $W=\log w$. Its two deck
periods are
$(0,2\pi i)$ and $(\beta_i^{-r_i},\log\beta_i)$.
Thus the integral of $\eta_i=c_i\,dU\wedge dW$
on its oriented core torus is
$\pm2\pi i c_i\beta_i^{-r_i}\ne0$.
Transport that compact torus through the
submersion near the core. By \eqref{cdp:eq:C.NS.6} all such
cycles are the same class $\gamma$ in the
nearby fibre. Therefore the single holomorphic
period satisfies, for every $i$,
\[
        \int_\gamma\omega_t
             =t^{r_i}\bigl(\pm2\pi i c_i\beta_i^{-r_i}
                              +O(t)\bigr).              \tag{C.NS.12}\hypertarget{tag-00TD}{}\label{cdp:eq:C.NS.12}
\]
Its vanishing order is simultaneously every
$r_i$. Hence all resonances equal $r$.
This uses the transported two-cycle; it makes
no assumption about parallel primitive angular
directions or an elliptic quotient of the
nearby complex torus.


\item  By the plumbing comparison the underlying
four-manifold of $T_0$ is a real four-torus.
For a principal elliptic bundle over the
elliptic curve, the homotopy sequence injects
$\Lambda_F$ into its fundamental group, and
the exact presentation is
\[
 \langle a,b,e,h:\ e,h\text{ central},\
             [a,b]=e^{c_1}h^{c_2}\rangle,
                         \quad(c_1,c_2)=c(T_0).          \tag{C.NS.13}\hypertarget{tag-00TE}{}\label{cdp:eq:C.NS.13}
\]
Its fundamental group is abelian only if
$c(T_0)=0$. This proves every necessity in (3).


\end{enumerate}

\item \par\medskip\noindent\textbf{The inverse transform is then an actual torus.}
Assume (3). The sheaf sequence
$0\to\Lambda_F\to\OO_B\to F\to0$ lifts the
principal torsor class of $T_0$ to $H^1(B,\OO_B)$
because its connecting class is precisely
$c(T_0)=0$. On $B=\C/\Gamma_B$ represent this
class by a constant $(0,1)$-form. Integrating it
on the universal cover gives the actual
translation lattice
\[
 (p,z)\mapsto(p+\gamma,z+a\bar\gamma),\quad
 (p,z)\mapsto(p,z+\lambda),
       \qquad \gamma\in\Gamma_B,\quad\lambda\in\Lambda_F.
                                                        \tag{C.NS.14}\hypertarget{tag-00TF}{}\label{cdp:eq:C.NS.14}
\]
Changes of the chosen Cech cocycle are fibre
translations and give an isomorphism with the
original torsor, exactly as in (C.KX.14).
The lattice is discrete of real rank four.
Consequently $T_0$ is a complex two-torus,
with its specified principal elliptic quotient
and all the marked fibres retained.


\item \par\medskip\noindent\textbf{A compact twistor family with the exact local period presentations.}
Write $T_0=\C^2/\Lambda$ in coordinates $(p,z)$
so that the common elliptic fibre direction is
$p=0$. On the affine $t$-chart of $\PP^1$ put
\[
 \ell_t(p,z)=(p+t\bar z,\ z-t\bar p).
                                                        \tag{C.NS.15}\hypertarget{tag-00TG}{}\label{cdp:eq:C.NS.15}
\]
This is real-linear and invertible for every
$t$, since its inverse has denominator
$1+|t|^2$. The values $\ell_t(v)$ for
$v\in\Lambda$ are holomorphic sections of
$\OO_{\PP^1}(1)^{\oplus2}$: on the other chart
they are $(s p+\bar z,s z-\bar p)$, $s=t^{-1}$.
Translate the vector bundle by these sections
and take the quotient. Each fibre lattice has
real rank four. The inverse real-linear maps
give local uniform fundamental parallelepipeds
over the compact base, proving a free proper
action and a smooth compact quotient
$\mathcal T\to\PP^1$. Its fibres are two-tori
and its fibre at zero is the actual $T_0$.
Pull this family back by $t=s^r$; the pullback
remains smooth and compact because the
original map is a submersion.


\begin{enumerate}[label={$\langle2\rangle$\arabic*.},leftmargin=1.3em,labelsep=.5em]
\item  For the $i$-th Hopf identification choose the
oriented fibre-lattice basis
$\lambda_i=(0,a_i)$, $\mu_i=(0,b_i)$ so that
$\beta_i=\exp(2\pi i b_i/a_i)$. This is the
basis supplied by \eqref{cdp:eq:C.NS.2}, not a replacement
of the given Hopf modulus. On the vector
space with coordinates $(P,Z)$ after
\eqref{cdp:eq:C.NS.15}, define
\[
 p_i=P-s^r\frac{\bar a_i}{a_i}Z,\qquad
 z_i=\frac{2\pi i}{a_i}Z.
                                                        \tag{C.NS.16}\hypertarget{tag-00TH}{}\label{cdp:eq:C.NS.16}
\]
The $\lambda_i$-period is $(0,2\pi i)$ and
the $\mu_i$-period is
\[
   \bigl(\delta_i(s),\log\beta_i\bigr),\qquad
   \delta_i(s)=s^r
        \left(\bar b_i-\frac{\bar a_i b_i}{a_i}\right).
                                                        \tag{C.NS.17}\hypertarget{tag-00TI}{}\label{cdp:eq:C.NS.17}
\]
The coefficient is nonzero because $a_i,b_i$
are an oriented real basis of a complex line.


\item  Complete this pair to an integral basis of
$\Lambda$ by two lifts of the base lattice.
Their $p_i$-periods define a smooth local
family of elliptic curves $B_{i,s}$, and
their exponentiated $z_i$-periods define a
holomorphic family of flat degree-zero line
bundles $M_i$ on $B_{i,s}$. Thus, on a
sufficiently small disc, the actual torus
family has the exact presentation
\[
 \mathcal T_s
   =M_{i,s}^{\times}/\langle g_i\rangle,\qquad
 g_i(p,\ell)=(p+\delta_i(s),\beta_i\ell).
                                                        \tag{C.NS.18}\hypertarget{tag-00TJ}{}\label{cdp:eq:C.NS.18}
\]
This equality follows by quotienting the
same four displayed lattice generators,
first by the three excluding $\mu_i$.
It retains the flat part of the torsor,
even when that part is not torsion.


\end{enumerate}

\item \par\medskip\noindent\textbf{The local weighted insertion with the varying flat line.}
At the marked point of $B_{i,0}$ perform
the weight-$(r,1)$ blowup of the relative
elliptic surface $B_i\to\Delta_s$, and replace
$M_i$ by its pullback tensored with
$\OO(R_i)$. The nonzero line has the same
smooth coarse chart as (C.NH.6)--(C.NH.7).
If $\delta_i(s)=d_i s^r$, the lifted deck
map in a flat local frame is exactly
\[
 (z,w,x)\mapsto
   \bigl(\beta_i^r(z+d_iw^r),
              \beta_iw,\beta_i^{-1}x\bigr),
           \qquad p-p_i=x^rz,\quad s=xw.
                                                        \tag{C.NS.19}\hypertarget{tag-00TK}{}\label{cdp:eq:C.NS.19}
\]
The change $Z=z/(\beta_i^r d_i)$ gives
the specified $H_r(\beta_i)$ on $x=0$.
On $w=0$ the chart is precisely \eqref{cdp:eq:C.NS.1};
hence the other component is obtained
from $T_0$ by the inverse of \eqref{cdp:eq:C.NS.2}.


\begin{enumerate}[label={$\langle2\rangle$\arabic*.},leftmargin=1.3em,labelsep=.5em]
\item  We check properness with the varying
flat line. On a small closed base disc,
choose smoothly varying unitary flat
metrics on $M_{i,s}$. If its period
characters are $e^{Z_j(s)}$ and base
periods $P_j(s)$, solve the two real
equations
$\operatorname{Re}(Z_j-c(s)P_j)=0$.
Their determinant is nonzero, so $c(s)$
is smooth and bounded. In the resulting
metric, the norm multiplier of the flat
lift in \eqref{cdp:eq:C.NS.18} is
$|\beta_i|\exp(-\operatorname{Re}(c(s)\delta_i(s)))$.
It is uniformly less than one after
shrinking the disc. The additional
exceptional line has the two-sided
polynomial distortion (C.NH.9)--(C.NH.10).
Exponential contraction still dominates
that polynomial uniformly on the closed
disc. The compact-annulus proof of
(C.NH.11) therefore gives a smooth
proper filling $\mathcal Y_i\to\Delta_s$.
Its infinite cyclic action is free,
and every prescribed $\beta_i$ is
retained.


\item  Off the exceptional divisor the canonical
section of $\OO(R_i)$ identifies the new
line with $M_i$, equivariantly for the
two deck maps. Consequently there is
the actual isomorphism
\[
       \mathcal Y_i\setminus H_i
            \simeq \mathcal T|_{\Delta_s}\setminus F_i,
                                                        \tag{C.NS.20}\hypertarget{tag-00TM}{}\label{cdp:eq:C.NS.20}
\]
where $F_i$ is the marked central
elliptic fibre in $T_0$. On the remaining
central component this is exactly the
punctured product map \eqref{cdp:eq:C.NS.2}.


\end{enumerate}

\item \par\medskip\noindent\textbf{Simultaneous replacement and compactness.}
The central curves $F_i$ are disjoint.
We justify that \eqref{cdp:eq:C.NS.20} allows
independent replacements near them,
although its generic tori need not
have an elliptic quotient.


\begin{enumerate}[label={$\langle2\rangle$\arabic*.},leftmargin=1.3em,labelsep=.5em]
\item  If a sequence approaches $H_i$ in
$\mathcal Y_i$, use a compact Hopf
fundamental domain in \eqref{cdp:eq:C.NS.19}.
Then $z,w$ stay bounded, do not
simultaneously approach zero, and
$x\to0$. Its image in the original
line presentation has
$p-p_i=x^rz$, $s=xw$, and flat
line coordinate proportional to $x^{-1}$.
To put that coordinate in a fixed
compact norm annulus requires
$O(1+|\log|x||)$ deck iterates.
Their total base displacement is
$O(|s|^r(1+|\log|x||))$, which
tends to zero. Hence every cluster
point of the images lies in $F_i$.
Conversely properness of $\mathcal Y_i$
shows that images approaching a point
outside $F_i$ have only their unique
inverse-image limit outside $H_i$.


\item  Choose disjoint relatively compact
neighborhoods $W_i$ of $F_i$ in
$\mathcal T$ over a small base disc.
The set consisting of $H_i$ and the
inverse image of $W_i\setminus F_i$
under \eqref{cdp:eq:C.NS.20} is open, by the
cluster-set calculation. Glue these
open replacements to
$\mathcal T\setminus\bigcup_i F_i$.
Their identification graphs are
relatively closed by that same
calculation, and the replacements
have disjoint central neighborhoods.
Thus the finite gluing is Hausdorff
and smooth. Nested compact
neighborhoods inside the $W_i$,
properness of each $\mathcal Y_i$,
and the compact original complement
cover every closed inverse image
of a sufficiently small base disc.
The resulting map is proper.


\item  The central main component is exactly
$P$: outside the marked fibres its
torsor is unchanged, and at every
marked fibre \eqref{cdp:eq:C.NS.19} reverses
the precise map \eqref{cdp:eq:C.NS.2}.
The Hopf components are the prescribed
$H_i$. Any constant axis translation
in their identification is realized
by the Hopf automorphism
$(z,w)\mapsto(a^rz,aw)$; the linear
part has already been retained in
the chosen fibre-lattice basis.
Thus all specified elliptic gluing
maps, not just their degrees, are
the original ones. The full central
equations are $xw=s$, so the fibre
is reduced SNC.


\item  Finally carry out these replacements
inside the compact pulled-back
twistor family over $\PP^1$ and
leave its complement unchanged.
The preceding proper local gluing
and the original compact complement
give a smooth compact total space.
Every fibre away from zero remains
the original two-torus. The central
support has Euler number zero since
each elliptic quasi-bundle, Hopf
component, and double elliptic curve
has Euler number zero. We have
proved (3)$\Rightarrow$(2), including
the full compact-total-space and
torsor comparisons. Assertion
(2)$\Rightarrow$(1) is restriction
to a disc, completing the equivalence.



\end{enumerate}

\end{enumerate}
\end{proof}






\begin{proposition}[Complete primary and secondary nonlinear Hopf star criterion]\hypertarget{tag-01DV}{}\label{cdp:st:C.2.UN.50}
\StatementRecord{cdp:st:C.2.UN.50}{proved}{01DV}{\SeriesRef{cdp:st:C.2.UN.49}, \SeriesRef{cdp:st:C.2.UN.43}, \SeriesRef{cdp:eq:C.SN.4}, \SeriesRef{cdp:eq:C.SN.5}, \SeriesRef{cdp:eq:C.SN.6}, \SeriesRef{cdp:eq:C.SN.9}}

In \ref{cdp:st:C.2.UN.49} allow each prescribed leaf to be
\[
 H_i=H_{r_i}(\beta_i)/\mu_{n_i},\qquad
 \zeta\cdot(z,w)=(\zeta^{r_i}z,\zeta w),\qquad
                         \gcd(r_i,n_i)=1,               \tag{C.SN.1}\hypertarget{tag-0122}{}\label{cdp:eq:C.SN.1}
\]
including $n_i=1$. These are all nonlinear primary and
secondary Hopf surfaces of algebraic dimension zero,
by \ref{cdp:st:C.2.UN.43}. The same equivalence between an actual
smooth-total torus smoothing and compact realization
over $\PP^1$ holds. Its necessary-and-sufficient tests
are dual normal lines, an elliptic central base with
no unmarked multiple fibres, one common resonance
$r_i=r$, and zero integral Chern vector of the
inverse logarithmic transform defined below.

The marked full fibres in $P$ still have multiplicity
$r_i$, while the finite integer $n_i$ enters the
exact inverse-transform framing. In a global
punctured-base trivialization its meridian is
\[
       \mu_{P,i}=r_i\delta_{\rm ext}
                         +(r_i\nu_i-n_i e_i),\qquad
                         c(T_0)=\sum_i\nu_i.             \tag{C.SN.2}\hypertarget{tag-0124}{}\label{cdp:eq:C.SN.2}
\]
Thus its rational Seifert vector is
$\sum_i(\nu_i-n_i e_i/r_i)$ in this convention.
The full integer $n_i$ is retained; its residue
modulo $r_i$ alone does not determine the criterion.
All original component moduli and elliptic gluing
maps are preserved by the compact realization.
\end{proposition}



\begin{proof}

\begin{enumerate}[label={$\langle1\rangle$\arabic*.},leftmargin=*]
\item We give every change to the exact maps of \ref{cdp:st:C.2.UN.49}.
Fix one leaf and suppress its index. Choose integers
$k,l$ with
\[
                         rk-nl=1,\qquad \gamma=\beta^n.
                                                        \tag{C.SN.3}\hypertarget{tag-0125}{}\label{cdp:eq:C.SN.3}
\]
On the primary cover of the paired $P$ neighborhood
use $p=x^rz$, the infinite deck
$(z,x)\mapsto(\beta^rz,\beta^{-1}x)$, and the finite
action $(z,x)\mapsto(\zeta^rz,\zeta^{-1}x)$.
The following are genuine invariant coarse coordinates:
\[
 u=z^n,\qquad X=xz^k,\qquad
 p=X^ru^{-l},\qquad
 (u,X)\sim(\gamma^ru,\gamma^lX).
                                                        \tag{C.SN.4}\hypertarget{tag-0126}{}\label{cdp:eq:C.SN.4}
\]
Here $u\ne0$. Since $rk-1=nl$, both coordinates
are finite-group invariant, and the finite quotient
map has no ramification. Conversely, a root of $u$
recovers $z$ and $x=Xz^{-k}$; all choices are exactly
the original finite action. Thus \eqref{cdp:eq:C.SN.4} is the
actual coarse neighborhood, not a formal coordinate
substitution.

\item The normal of $J$ in $P$ has multiplier $\gamma^l$
over $\C^*/\gamma^{r\Z}$. Its order is exactly $r$,
since $\gcd(l,r)=1$. On the normal root cover write
$u=y^r$, $X=s y^l$; then $p=s^r$.
The generator $(y,s)\mapsto(\zeta y,\zeta^{-l}s)$
preserves these coordinates. Normalize its normal
character to $\zeta$: the fibre translation is then
$y\mapsto\zeta^n y$, because $-ln\equiv1\pmod r$.
This proves the exact torsion-character comparison,
including its dependence on $n$.
The dual normal on the secondary Hopf axis is
the inverse character, so its normal-root
identification with the given $P$ gives precisely
\eqref{cdp:eq:C.SN.4}.

\item The punctured product map used for the inverse
transform is
\[
 (u,X)\longmapsto(p,W)
       =(X^ru^{-l},\,X^{-n}u^k),\qquad
 u=p^nW^r,\quad X=p^kW^l,\qquad
                       F=\C^*/\gamma^{\Z}.               \tag{C.SN.5}\hypertarget{tag-0127}{}\label{cdp:eq:C.SN.5}
\]
In the ordered variables $(u,X)$ and $(p,W)$, the exponent
matrix is $\left(\begin{smallmatrix}-l&r\\k&-n\end{smallmatrix}\right)$
and has determinant $nl-rk=-1$, hence is unimodular.
Consequently these are inverse holomorphic maps
on the punctured chart. On the primary cover
they are $W=x^{-n}$ and $u=z^n=p^nW^r$.
Replacing the chart by $F\times\Delta_p$ through
\eqref{cdp:eq:C.SN.5} defines the actual principal torsor $T_0$.
It retains the original analytic transitions
exactly as in (C.NS.3).

\item Let $a$ be the primitive positive angular loop of
$u$ with $X$ fixed, $b$ the infinite deck loop,
and $\mu_P$ the normal loop of $X$. In the
punctured product basis $(\delta,e,h)$,
\eqref{cdp:eq:C.SN.5} gives the full integral matrix
\[
 \begin{pmatrix}a&b&\mu_P\end{pmatrix}_{(\delta,e,h)}
   =
 \begin{pmatrix}
   -l&0&r\\
    k&0&-n\\
    0&1&0
 \end{pmatrix},
 \qquad
 \delta=n a+k\mu_P,\quad e=r a+l\mu_P,\quad h=b.
                                                        \tag{C.SN.6}\hypertarget{tag-0128}{}\label{cdp:eq:C.SN.6}
\]
Its determinant is one. In particular the normal
meridian is $r\delta-ne$ and is primitive even
when $n>1$.

\item On the Hopf side the invariant normal coordinate
is $Y=wz^{-k}$. Its normal-line multiplier is $\gamma^{-l}$,
and the diagonal reference pair has reciprocal
base coordinate $Y^ru^l=w^r/z$ and fibre coordinate
$w^n=Y^nu^k$. The complementary solid torus
therefore has meridian $n a-k\mu_H$ and fibre
cycles $(r a-l\mu_H,b)$. The actual smoothing
identifies $\mu_H=-\mu_P$. Formula \eqref{cdp:eq:C.SN.6}
then identifies its meridian with $\delta$
and its fibre cycles with $(e,h)$. This proves
the inverse-transform topology and the common
integral fibre-cycle identification for every
secondary leaf. Substituting
$\delta_i=\delta_{\rm ext}+\nu_i$ proves \eqref{cdp:eq:C.SN.2}.
The nonlinear-to-diagonal comparison (C.NS.5)
is finite-group invariant, so its boundary
calculation applies to this actual quotient.

\item The canonical form
$dz\wedge dw/w^{r+1}$ is invariant under both
the infinite contraction and $\mu_n$. It gives
$K_H=-(r+1)J$ on the quotient. Thus the
canonical-order calculation, the effectivity
argument forcing elliptic $B$ and no unmarked
multiple fibres, and the Chern-zero necessity
are unchanged. The only change to the core
period is explicit. In $U=z/w^r$, $V=\log w$,
the two primitive periods are
\[
             (0,2\pi i/n),\qquad
             (\beta^{-r},\log\beta).
                                                        \tag{C.SN.7}\hypertarget{tag-0129}{}\label{cdp:eq:C.SN.7}
\]
The leading two-form has core integral
$\pm2\pi i c\beta^{-r}/n$, still nonzero.
All cores are the same integral fibre class
by \eqref{cdp:eq:C.SN.6}, so their one global period
again forces all resonances $r_i$ to coincide.
The finite orders $n_i$ need not coincide.

\item For compact sufficiency perform the exact
inverse transformations and use their
Chern-zero torsor $T_0$, which is an actual
torus by (C.NS.14). In the twistor family
(C.NS.15), choose the $i$-th fibre basis
so that its multiplicative modulus is
$\gamma_i=\beta_i^{n_i}$. After $t=s^r$,
the same four-period calculation gives
\[
 \mathcal T_s=M_{i,s}^{\times}/
   \langle(p,\ell)\mapsto
        (p+d_i s^r,\gamma_i\ell)\rangle,\qquad d_i\ne0.
                                                        \tag{C.SN.8}\hypertarget{tag-012A}{}\label{cdp:eq:C.SN.8}
\]
At the marked point perform the weight-$(r,1)$
blowup, but now tensor the flat line by
$\OO(n_iR_i)$. Its local line coordinate
has weight $-n_i$. Setting it equal to one
leaves the finite action
\[
 (z,w,x)\mapsto
      (\zeta^{r}z,\zeta w,\zeta^{-1}x),
                    \qquad \zeta^{n_i}=1.                \tag{C.SN.9}\hypertarget{tag-012B}{}\label{cdp:eq:C.SN.9}
\]
It is free: a nonzero $w$ or $x$ forces
$\zeta=1$, and if both vanish then $z\ne0$
and $\gcd(r,n_i)=1$ forces the same conclusion.
This proves directly that the coarse nonzero
line is smooth. Equivalently the weighted
orbifold inertia has order $r$ and acts
on the line with the faithful weight $n_i$.

\item Choose the given $n_i$-th root $\beta_i$
of $\gamma_i$. The infinite line deck
lifts to the exact map
\[
 (z,w,x)\mapsto
   \bigl(\beta_i^r(z+d_iw^r),\beta_iw,\beta_i^{-1}x\bigr)
                                                        \tag{C.SN.10}\hypertarget{tag-0123}{}\label{cdp:eq:C.SN.10}
\]
on \eqref{cdp:eq:C.SN.9}. It commutes with the finite
action. Its exceptional component is
exactly $H_r(\beta_i)/\mu_{n_i}$ after
$Z=z/(\beta_i^r d_i)$.
Off that component the flat-line
coordinate is $W=x^{-n_i}$; consequently
the remaining central chart and its
outer product are exactly
\eqref{cdp:eq:C.SN.4}--\eqref{cdp:eq:C.SN.5}.

\item The polynomial norm distortion is now
$C(1+|j|)^{n_i/r}$ for the $j$-th
translation lift. The flat-line
contraction multiplier is
$|\gamma_i|\exp(-\operatorname{Re}(c_i(s)d_i s^r))$,
uniformly less than one on a small disc.
It still dominates that polynomial.
The same compact norm annulus proves
properness, and comparison with arbitrarily
large powers proves freeness of the
infinite cyclic quotient. Thus there
is again an actual isomorphism of the
complement of the new Hopf component
with the original twistor tube minus
its marked central elliptic curve.

\item In the cluster calculation the original
flat coordinate is $x^{-n_i}$, so only
$O(1+|\log|x||)$ deck iterates are
needed. Their base displacement remains
$O(|s|^r(1+|\log|x||))$. All cluster
points therefore lie on the prescribed
central fibre curve. The simultaneous
disjoint replacements and compact gluing
of \ref{cdp:st:C.2.UN.49} apply with these exact
coarse maps. The main component is the
original $P$, and the finite diagonal
actions commute with the axis-translation
automorphisms used there. Every original
gluing map is retained. The full
central equations are still $xw=s$;
the finite action preserves this
function and is free, so the full
fibre remains reduced SNC.
This proves the equivalence and
compact realization asserted here.

\end{enumerate}
\end{proof}






\begin{proposition}[The mixed-cycle criterion and compact occurrence with Hopf components]\hypertarget{tag-01DW}{}\label{cdp:st:C.2.UN.51}
\StatementRecord{cdp:st:C.2.UN.51}{proved}{01DW}{\SeriesRef{cdp:eq:C.MX.3}, \SeriesRef{cdp:st:C.2.UN.46}, \SeriesRef{cdp:eq:C.MX.4}, \SeriesRef{cdp:eq:C.MX.5}, \SeriesRef{cdp:st:C.2.UN.48}, \SeriesRef{cdp:eq:C.MX.8}, \SeriesRef{cdp:eq:C.MX.9}}

Specify an oriented cycle of surfaces, each of which is either
a diagonal primary or secondary Hopf surface of algebraic dimension
zero, marked by its two axis curves, or
\[
                  R_i=\PP(\OO_{E_i}\oplus L_i),
                  \qquad E_i\ \hbox{elliptic},\quad
                         L_i\in\operatorname{Pic}^0(E_i),
                                                        \tag{C.MX.1}\hypertarget{tag-00RR}{}\label{cdp:eq:C.MX.1}
\]
marked by its zero and infinity sections. Specify the elliptic
isomorphism joining the outgoing curve to the next incoming curve,
and a trivialization of their tensor product of normal lines.
The associated proper normal-crossing plumbing has two-torus
general fibre exactly when the cycle product $M$ of its actual
rank-three integral complement lattices is the identity.
Every actual reduced smooth-total two-torus smoothing of the
specified support satisfies this criterion.

For every such cycle containing a Hopf component, $M=I$ is also
equivalent to occurrence of the entire specified reduced support
in a smooth compact threefold over $\PP^1$, with two-torus
general fibre and Euler number zero for every reduced fibre
support. This retains all component surfaces and every elliptic
gluing map. The construction chooses the normal-product constants.
A compatible cycle containing Hopf components contains at least
two of them.

The local criterion also includes all-ruled cycles and the
identification of the two marked sections of one ruled surface.
In that case the exact punctured period lattice has the form
\[
                 \Lambda+\Z\bigl(b_0+v_0\log t\bigr),
           \qquad \Lambda\simeq\Z^3\subset\C^2,          \tag{C.MX.2}\hypertarget{tag-00RT}{}\label{cdp:eq:C.MX.2}
\]
where the projection of $b_0$ to the elliptic quotient of the
common complement is the actual total translation of the base
gluing maps. It is retained in this formula.
\end{proposition}



\begin{proof}

\begin{enumerate}[label={$\langle1\rangle$\arabic*.},leftmargin=*]
\item \par\medskip\noindent\textbf{The complements and their actual line-bundle maps.}
For a ruled component, removing the two sections gives the
nonzero line bundle $L_i^\times$. Choose an origin on $E_i$
and a rigidification of $L_i$ there. Since $\deg L_i=0$, it
has constant factors of automorphy on the universal cover
of $E_i$. Taking a logarithm of its fibre coordinate gives
the exact presentation
\[
 L_i^\times=\C^2/\Lambda_i,\qquad
 \Lambda_i=\left\langle
       (\gamma_1,\log\chi_i(\gamma_1)),
       (\gamma_2,\log\chi_i(\gamma_2)),(0,2\pi i)
                         \right\rangle_{\Z} .             \tag{C.MX.3}\hypertarget{tag-00RU}{}\label{cdp:eq:C.MX.3}
\]
Here $\gamma_1,\gamma_2$ are an oriented integral period
basis of $E_i$, and $\chi_i$ is the chosen constant
factor character. The lattice is discrete of real rank
three and spans $\C^2$ over $\C$. Each marked-section
neighbourhood is the entire normal line bundle obtained
by removing the other section. Inverting the line
coordinate identifies the two descriptions on their
common complement. Their positive normal meridians are
opposite primitive elements of $\Lambda_i$.

\item For Hopf components use the actual entire normal
line-bundle neighbourhoods and logarithmic lattices
of (C.HC.4)--(C.HC.5), including their finite diagonal
deck groups. Every edge is the same explicit operation:
apply its elliptic base isomorphism, and send a nonzero
normal vector to its dual vector whose product is $t$.
It is a global biholomorphism of the two complements.
On universal covers it is affine. Indeed its base lift
is affine, and the derivative of its logarithmic normal
multiplier is doubly periodic and holomorphic, hence
constant. Thus it has a constant complex linear part
which maps $\Lambda_i$ isomorphically to $\Lambda_{i+1}$.
The product defines the actual integral automorphism
$M$ up to conjugacy. Changes of logarithms change only
translations by lattice vectors.

\item Every complement is smoothly $T^3\times\R$.
For \eqref{cdp:eq:C.MX.3}, a unitary flat metric on $L_i$ makes
the logarithm of its norm a linear height, with compact
three-torus level sets. The same description for a
Hopf complement was proved in \ref{cdp:st:C.2.UN.46}. Hence all component
cores and annular necks have the same rank-three
fundamental group, with collar maps exactly the
indicated integral maps. The finite normal-product
plumbing of \ref{cdp:st:C.2.UN.46} applies without a changed local map:
choose disjoint section neighbourhoods, use $xy=t$
in each dual-line neck, and glue where one normal
norm exceeds a fixed inner radius. The two ends
are disjoint for sufficiently small $|t|$.
Closed intermediate cores and necks give properness;
the opposite open ends give the same Hausdorff
comparison as in that proof.

\item Van Kampen consequently gives
\[
                        \pi_1(Y_t)=\Lambda_1\rtimes_M\Z .
                                                        \tag{C.MX.4}\hypertarget{tag-00RV}{}\label{cdp:eq:C.MX.4}
\]
The rank-three lattice injects in this presentation.
For an actual given smoothing, the divisor sections
and the product-preserving smooth tubular map used
in \ref{cdp:st:C.2.UN.46} identify its collars with these very normal
products. Radial contraction of the tubular maps
gives the same maps on $\pi_1$. Thus \eqref{cdp:eq:C.MX.4} holds
also for that smoothing. A two-torus has abelian
fundamental group, forcing $M=I$.

\item Conversely, when $M=I$, the complex linear cycle
product is the identity because $\Lambda_1$
spans $\C^2$. Constant changes of the two logarithmic
one-forms on successive pieces therefore glue them
to a closed holomorphic coframe on each compact
nearby fibre. The commuting dual fields have
complete complex flows and give the actual
biholomorphism with $\C^2$ divided by its period
lattice, exactly as in (C.HC.8). The core periods
and the one cycle period are all the fundamental
group generators, giving \eqref{cdp:eq:C.MX.2}. This proves
the complete local criterion for every stated
combination of components.

\par\medskip\noindent\textbf{Common normal periods and the Hopf restriction.}
Assume $M=I$ and identify all complements with
$G=\C^2/\Lambda$ using their actual linear maps.
Let $n_i\in\Lambda$ be the positive outgoing normal
meridian. The incoming normal meridian of component
$i$ is $-n_{i-1}$. On a ruled component the two
meridians are opposite, so
\[
                      n_i=n_{i-1}
                  \quad\hbox{at every ruled vertex}.   \tag{C.MX.5}\hypertarget{tag-00RW}{}\label{cdp:eq:C.MX.5}
\]
At a Hopf vertex $n_i,n_{i-1}$ are complex-linearly
independent: its two logarithmic coordinate
meridians are a complex basis. A cycle with
exactly one Hopf vertex would have equality of
these two meridians by following \eqref{cdp:eq:C.MX.5} through
all the other vertices. This proves its impossibility.

\item Put $W=\operatorname{span}_{\R}\Lambda$ and choose
a nonzero real linear functional $\ell$ with
kernel $W$. At a Hopf vertex, its contracting
logarithm gives
\[
 \log|\alpha_i|\,\ell(i n_i)
        =\log|\beta_i|\,\ell(i n_{i-1}).
                                                        \tag{C.MX.6}\hypertarget{tag-00RX}{}\label{cdp:eq:C.MX.6}
\]
Both coefficients are negative, and neither value
can be zero because the two meridians span the
complex plane. At a ruled vertex their values
are equal by \eqref{cdp:eq:C.MX.5}. Thus, as soon as there is
a Hopf vertex, all $\ell(i n_i)$ have one nonzero
sign. The actual logarithmic cycle translation has
\[
                  2\pi i v_0=u=-\sum_i n_i,\qquad
                              \ell(iu)\ne0 .            \tag{C.MX.7}\hypertarget{tag-00RY}{}\label{cdp:eq:C.MX.7}
\]
Write $u=d p$ with $p$ primitive in $\Lambda$ and
$d>0$. The projection along $\C p$ takes $\Lambda$
to a genuine elliptic lattice, since
$\C p\cap W=\R p$. Exactly as in (C.HG.6), actual
linear coordinates give
\[
 \Lambda=\left\langle
 (0,2\pi i),(2\pi i,0),(\log q_0,\log a)
                     \right\rangle_{\Z},\qquad
              |q_0|<1,\quad u=d(0,2\pi i).             \tag{C.MX.8}\hypertarget{tag-00RZ}{}\label{cdp:eq:C.MX.8}
\]
Here $a\ne1$: otherwise the explicit embedded
elliptic curve (C.HG.7) would lie in the common
complement of an algebraic-dimension-zero Hopf
surface, contradicting its two-curve calculation.

\par\medskip\noindent\textbf{Transporting both sorts of component.}
Use the compact elliptic surface and suspension
constructed in \ref{cdp:st:C.2.UN.48} with the exact periods
\[
 \Lambda(t)+\Z(0,d\log t),\qquad
 \Lambda(t)=\left\langle
 (0,2\pi i),(2\pi i,0),(\log q_0,\log a(t))
                         \right\rangle_{\Z},             \tag{C.MX.9}\hypertarget{tag-00S0}{}\label{cdp:eq:C.MX.9}
\]
where $a(0)=a$ and the actual Tate parameter is
$Q=t^d$. All the prescribed normal meridians and
contracting logarithms have fixed integral
coordinates in \eqref{cdp:eq:C.MX.8}. Retain those coordinates
in \eqref{cdp:eq:C.MX.9}.

\item At a Hopf vertex use the moving coordinate basis
$n_i(t)/(2\pi i),-n_{i-1}(t)/(2\pi i)$.
It is invertible near zero, retains both
contracting multipliers and all finite diagonal
characters, and gives the actual Hopf family
of \ref{cdp:st:C.2.UN.48}.

\item At a ruled vertex those two columns are
dependent, so a different actual frame is
required. Its primitive normal period $n_i$
extends to a lattice basis $n_i,\lambda_i,c_i$.
Projection along its complex line gives the
elliptic base lattice of this particular
ruled component. Choose a base coordinate
with the projected $\lambda_i$-period $2\pi i$,
and a logarithmic fibre coordinate having
normal period $2\pi i$ and zero $\lambda_i$-period.
Transport $n_i,\lambda_i,c_i$ by their fixed
integer coordinates. The first two remain
complex-linearly independent. They give
holomorphic coordinates in which the complement
lattice is
\[
        \left\langle
          (0,2\pi i),(2\pi i,0),
                      (\gamma_i(t),b_i(t))
                                  \right\rangle_{\Z} .   \tag{C.MX.10}\hypertarget{tag-00RS}{}\label{cdp:eq:C.MX.10}
\]
Its projected periods remain real-linearly
independent on a smaller disc. They define
a proper elliptic family $E_i(t)$, and the
factor $\exp b_i(t)$ defines an actual
degree-zero line bundle $L_i(t)$ on it.
The resulting
$\PP(\OO_{E_i(t)}\oplus L_i(t))$ is a smooth
proper family of the prescribed ruled
components, with their marked sections.
At zero its line-bundle presentation is
the original one in \eqref{cdp:eq:C.MX.3}, so it is
an isomorphism of marked surfaces.

\item Every moving edge linear map sends the
outgoing normal vector to minus the next
incoming normal vector. Its induced base
map is therefore a holomorphic elliptic
isomorphism, and its nonzero normal map
is the reciprocal linear-fibre map with
the prescribed constant factors. Thus it
extends to exactly the required dual
line-bundle data. Keep the constants of
all original elliptic gluing maps.
The coefficients of the adjustable
normal-product logarithms in the total
common translation are
$-n_i(t)/(2\pi i)$. At a Hopf vertex
the two adjacent such vectors are
complex-linearly independent. Fix the
other constants and solve the resulting
two-by-two holomorphic system to make
the constant cycle translation zero.
The total period is then exactly
$(0,d\log t)$ in \eqref{cdp:eq:C.MX.9}, while every
central component and elliptic gluing
map is unchanged.

\item The uniformly chosen flat metrics, cores
and normal-product necks now give a
smooth proper local filling. On its
punctured disc the piecewise affine
maps into the common period quotient
agree modulo precisely \eqref{cdp:eq:C.MX.9};
their changes of logarithm are integral
normal periods. Their fibrewise inverses
are the closed-coframe inverses already
proved above, so they give an actual
biholomorphism of the punctured
families with the compact suspension.
The nested-annulus replacement of
\ref{cdp:st:C.2.UN.48} produces a smooth compact total
space over $\PP^1$ with exactly the
specified mixed central support.
Each of its components and double
curves has Euler number zero. All
other reduced supports are the
elliptic-base bundles in that
suspension and also have Euler
number zero. This proves the
compact assertion.

\par\medskip\noindent\textbf{The retained all-ruled translation.}
When every vertex is ruled, \eqref{cdp:eq:C.MX.5}
makes all normal periods equal to
one primitive $n$. The common
elliptic quotient is
$\C^2/(\C n+\Lambda)$. The coefficient
vectors of all normal-product constants
lie in the single line $\C n$,
whereas the projection of the full
affine cycle constant $b_0$ is the
composition of the actual base
translations. Consequently varying
normal-product constants changes
only its vertical part; the period
formula \eqref{cdp:eq:C.MX.2} retains the projected
translation. The local torus proof
above made no assumption that it
vanishes. In particular the
single ruled self-identification
is included by its actual
normal-line inversion, rather than
the Hopf coordinate interchange.
The exact examples and moving-frame comparisons in
\path{code/A_mixed_cycles.py}
check 169 mixed cycles with two Hopf vertices and
arbitrary numbers of ruled insertions in the
displayed range. Their primitive cycle divisor is
$\gcd(a,b)$, and their real period determinant is
$\log|t|\bigl(a\log|\beta|+b\log|\alpha|\bigr)>0$.
The same calculation retains the all-ruled
horizontal translation exactly.

\end{enumerate}
\end{proof}






\begin{proposition}[The weighted root criterion, exact inertia and compact occurrence]\hypertarget{tag-01DX}{}\label{cdp:st:C.2.UN.52}
\StatementRecord{cdp:st:C.2.UN.52}{proved}{01DX}{\SeriesRef{cdp:eq:C.WH.1}, \SeriesRef{cdp:eq:C.WH.2}, \SeriesRef{cdp:eq:C.WH.3}, \SeriesRef{cdp:eq:C.WH.4}, \SeriesRef{cdp:st:C.2.UN.40}, \SeriesRef{cdp:st:C.2.UN.46}, \SeriesRef{cdp:st:C.2.UN.51}, \SeriesRef{cdp:st:C.2.UN.43}, \SeriesRef{cdp:eq:C.WH.9}, \SeriesRef{cdp:eq:C.WH.8}, \SeriesRef{cdp:st:C.2.UN.48}, \SeriesRef{cdp:eq:C.WH.12}, \SeriesRef{cdp:eq:C.WH.11}, \SeriesRef{cdp:eq:C.WH.15}}

Let $S$ be a simple normal-crossing oriented cycle
with at least two embedded components, which are specified unblown
diagonal primary or secondary Hopf surfaces $H_i$
of algebraic dimension zero, with their two axes
joined by specified elliptic isomorphisms between distinct components.
Assign positive integers $m_i$ to its components.
The compatible weighted normal data defined in
\eqref{cdp:eq:C.WH.1}--\eqref{cdp:eq:C.WH.2} give an explicit smooth proper
local plumbing with full fibre $\sum_i m_iH_i$.

There is the following complete test for its
general fibre to be a connected two-torus.
Put $N=\operatorname{lcm}_i(m_i)$, form the finite
cyclic component covers \eqref{cdp:eq:C.WH.3}, and join their
axes by the normalized root necks \eqref{cdp:eq:C.WH.4}.
Their graph must be connected, and the product
$M$ of their actual rank-three integral complement
lattice maps must be the identity. Every actual
smooth-total torus degeneration with the given
weighted support supplies these data and satisfies
this test. Conversely every weighted support
passing this test occurs, with its specified
reduced surface and component multiplicities,
over a compact rational base, with a smooth
compact total threefold and Euler number zero
for every reduced fibre support. The compact
construction chooses normal-product constants;
it does not identify different full analytic
neighbourhoods without a comparison.

The finite-root comparison also determines the
inertia. There is a common number $r$ of connected
strict covers of every $H_i$, with
$r\mid\gcd_i m_i$. Write
\[
 g_i=\gcd(m_i,m_{i+1}),\qquad
 h_i=\frac{N}{\operatorname{lcm}(m_i,m_{i+1})}.
\]
Each strict cover of $H_i$ has degree $m_i/r$;
each connected cover of its outgoing axis has
degree $g_i/r$ and is \'{e}tale. If $K$ is the
stabilizer of a strict component and $I_i$
its generic inertia subgroup, then
\[
 \begin{gathered}
 |K|=N/r,\qquad |I_i|=N/m_i,\\
 |I_iI_{i+1}|=N/g_i,\qquad
 |I_i\cap I_{i+1}|=h_i.
 \end{gathered}
\tag{C.WH.0}\hypertarget{tag-016X}{}\label{cdp:eq:C.WH.0}
\]
The root resolution inserts exactly $h_i-1$
ruled elliptic components along each lifted edge.

There are compact examples with multiplicities
$(1,n)$ for every $n\ge1$, described explicitly
below. There is also a compact three-component
example with multiplicities $(1,2,2)$ in which
all three components are primary Hopf surfaces.
\end{proposition}



\begin{proof}\leavevmode\par



 
\begin{enumerate}[label={$\langle1\rangle$\arabic*.},leftmargin=0pt,labelsep=.5em]
\item \par\medskip\noindent\textbf{The actual weighted data and smooth plumbing.}
Let $C_i^-$ and $C_i^+$ be the incoming and
outgoing axes on $H_i$, and let
$\phi_i:C_i^+\to C_{i+1}^-$ be the specified map.
The data consist of a holomorphic line
$\mathcal L_i$ and a section
\[
 \sigma_i\in H^0(H_i,\mathcal L_i^{-m_i}),
 \qquad
 \operatorname{div}(\sigma_i)=
       m_{i-1}C_i^-+m_{i+1}C_i^+,
\tag{C.WH.1}\hypertarget{tag-016Y}{}\label{cdp:eq:C.WH.1}
\]
together with the actual edge identifications
\[
 \begin{aligned}
 \mathcal L_i|_{C_i^+}
   &\simeq \phi_i^*N_{C_{i+1}^-/H_{i+1}},\\
 \phi_i^*(\mathcal L_{i+1}|_{C_{i+1}^-})
   &\simeq N_{C_i^+/H_i}.
 \end{aligned}
\tag{C.WH.2}\hypertarget{tag-0177}{}\label{cdp:eq:C.WH.2}
\]
The leading coefficients of $\sigma_i$ and
$\sigma_{i+1}$ at this axis must agree under
these identifications. Equivalently they
give the same trivialization of
$N_{C_i^+/H_i}^{m_{i+1}}\otimes
 \phi_i^*N_{C_{i+1}^-/H_{i+1}}^{m_i}$.
This is an equality of specified holomorphic
line maps, including their nonzero constants.


\begin{enumerate}[label={$\langle2\rangle$\arabic*.},leftmargin=1.3em,labelsep=.5em]
\item  These are precisely the maps furnished by
an actual defining function $f$:
take $\mathcal L_i=\OO_Y(H_i)|_{H_i}$,
divide $f$ by the $m_i$th power of the
canonical divisor section, and restrict
to $H_i$. The divisor of the resulting
section is \eqref{cdp:eq:C.WH.1}. On a double curve
the two canonical sections give \eqref{cdp:eq:C.WH.2}
and the single equation $f=x^{m_i}y^{m_{i+1}}$
after absorbing a unit. \ref{cdp:st:C.2.UN.40} forces
both axes into the boundary in every
actual case; their disjointness gives
the stated cycle and no triple points.


\item  The line $\mathcal L_i$ is flat. Its
$m_i$th tensor power is the inverse of
the indicated axis divisor line, which
has the constant factors of automorphy
calculated in \ref{cdp:st:C.2.UN.40}. Transfer that flat
connection through \eqref{cdp:eq:C.WH.1}. In a local
frame $e$, divide the connection form
on $e^{m_i}$ by $m_i$; the frame-change
formula divides in the same way, giving
an actual flat holomorphic connection
on $\mathcal L_i$.
On either entire axis line-bundle chart
of the Hopf surface it is therefore the
pullback of its axis restriction:
the projection is a homotopy equivalence,
and the two flat holomorphic lines have
the same monodromy character. This gives
an actual holomorphic line isomorphism,
not a formal neighbourhood assertion.


\item  On a component core use the total line
$\mathcal L_i$ and the function
$\sigma_i(v^{m_i})$. On the edge use the
two normal lines in \eqref{cdp:eq:C.WH.2}, with function
$x^{m_{i+1}}y^{m_i}$ in the indicated
weighted trivialization. The preceding
entire line-bundle maps glue these
functions exactly. Choose bounded normal
norms and disjoint axis collars as in
\ref{cdp:st:C.2.UN.46}. If both attaching norms exceed
the inner radius $\epsilon$, their
weighted product has modulus at least
$\epsilon^{m_i+m_{i+1}}$; take the base
disc smaller than all these bounds.
Thus the two ends of each neck are
disjoint. The same closed-graph and
compact-core comparison proves a
Hausdorff smooth proper threefold.
Its central divisor has exactly the
assigned multiplicities. Away from
zero its function is a submersion.


\end{enumerate}

\item \par\medskip\noindent\textbf{The exact component covers and normalized necks.}
For each component form the finite cyclic
cover
\[
 \mathcal E_i=
 \operatorname{Norm}
 \{\xi\in\mathcal L_i^{-1}:
                         \xi^{m_i}=\sigma_i\}.
\tag{C.WH.3}\hypertarget{tag-0178}{}\label{cdp:eq:C.WH.3}
\]
It has degree $m_i$, possibly with more
than one connected component. It is
the actual strict cover in the
normalization of $s^N=f$:
on the component core the formula
$\xi=s^{N/m_i}/v$ gives \eqref{cdp:eq:C.WH.3}.
The root generator acts on it by
$\xi\mapsto\zeta_N^{N/m_i}\xi$.


\begin{enumerate}[label={$\langle2\rangle$\arabic*.},leftmargin=1.3em,labelsep=.5em]
\item  For an edge put $a=m_i$, $b=m_{i+1}$,
$g=\gcd(a,b)$, $A=a/g$, $B=b/g$ and
$h=N/\operatorname{lcm}(a,b)$.
The normalization of $x^a y^b=s^N$
has $g$ local branches. After choosing
one of their constant roots of unity,
each is exactly
\[
      uv=s^h,\qquad x=u^B,\quad y=v^A.
\tag{C.WH.4}\hypertarget{tag-0179}{}\label{cdp:eq:C.WH.4}
\]
Indeed $(x^Ay^B)^g=s^{gABh}$ first
gives its $g$ branches. In one branch
the displayed finite map is
birational, because $\gcd(A,B)=1$
makes the choices of the roots
$u^B=x$, $v^A=y$, $uv=s^h$ unique
off $xy=0$. Its source is the normal
$A_{h-1}$ hypersurface, so it is the
normalization. Units in the original
equation are absorbed in local
holomorphic roots before this map.
The maps glue by uniqueness of
normalization.


\item  Every connected component of
$\mathcal E_i$ is a smooth diagonal
Hopf surface of algebraic dimension
zero. Here is the actual quotient
comparison. Its open part is a
connected finite \'{e}tale cover of
$G_i=\C^2/\Lambda_i$, hence is
$\C^2/\Lambda_i'$ for a finite-index
sublattice $\Lambda_i'\subset\Lambda_i$.
Let $q_+,q_-$ be the smallest positive
multiples of the two angular meridians
belonging to $\Lambda_i'$.
The local normalization
$\xi^{m_i}=z^{m_{i+1}}$ shows
$q_+=m_i/\gcd(m_i,m_{i+1})$;
the incoming formula is identical.
Divide the two logarithmic coordinates
by $q_+,q_-$. The resulting lattice
contains the two usual angular periods,
and exponentiating it gives a diagonal
deck group containing a contraction.
A finite element fixing a nonzero axis
point would give a smaller positive
angular multiple in $\Lambda_i'$,
contrary to the definition of $q_\pm$.
Thus its quotient of $\C^2\setminus0$
is a smooth primary or secondary
diagonal Hopf surface.


\item  The coordinate-power map
$(z,w)\mapsto(z^{q_+},w^{q_-})$
is equivariant for these actual deck
groups. It is finite on the compact
quotients, agrees with the open
lattice covering, and has precisely
the local branch maps above. Uniqueness
of the finite normalization therefore
identifies it with the given component
of \eqref{cdp:eq:C.WH.3}, including its two axes.
Finite pullback preserves algebraic
dimension, so it is still of algebraic
dimension zero. Each preimage axis is
one connected elliptic curve.


\item  The local branches \eqref{cdp:eq:C.WH.4} join these
covers by their actual root maps.
Every vertex has degree two. Resolving
$uv=s^h$ inserts the explicit $h-1$
projective-line bundles over the
corresponding elliptic axis cover,
with the resolution charts of
(C.HG.9)'s proof. They have two
disjoint sections and no triple
points. Thus the connected components
of the root support graph are cycles;
the compressed graph retains only
the Hopf vertices and labels each
edge by its positive integer $h$.
On the resolved family every central
component has multiplicity one.
Its Stein factor therefore has degree
one over each connected root disc.
It follows that the original general
fibre is connected exactly when this
root graph is connected.


\end{enumerate}

\item \par\medskip\noindent\textbf{The integral criterion and the exact root inertia.}
On a compressed edge the two normal
lines are dual and the product equation
is $uv=s^h$. Its logarithmic map is
the one in (C.HC.5), with $\log s$
replaced by $h\log s$. It identifies
the actual rank-three lattices of
the two open Hopf covers. Let $M$
be their cycle product. The resolved
bridges give the same linear product:
their two normal meridians are
opposite, and their successive open
maps compose to this exact compressed
edge map. \ref{cdp:st:C.2.UN.46} and \ref{cdp:st:C.2.UN.51}, or their
same core-and-neck Van Kampen
calculation, give
\[
                  \pi_1(F_s)=\Lambda\rtimes_M\Z.
\tag{C.WH.5}\hypertarget{tag-017A}{}\label{cdp:eq:C.WH.5}
\]
Consequently a two-torus forces $M=I$.
Conversely product identity glues
the closed logarithmic coframes.
Their complete commuting flows give
an actual complex two-torus, as in
(C.HC.8). The root fibre over a
nonzero $s$ maps isomorphically to
the original fibre over $s^N$.
This proves the stated local criterion.


\begin{enumerate}[label={$\langle2\rangle$\arabic*.},leftmargin=1.3em,labelsep=.5em]
\item  For an actual smoothing the same
calculation uses its actual normal
data. On cores, the canonical divisor
section and a smooth tubular
identification make its function
$v^{m_i}\sigma_i$: the ratio to this
expression is one on the zero
section and has a unique small smooth
$m_i$th root. On double curves use
both divisor sections and their
weighted product trivialization.
The resulting tubular maps give the
same covering and collar maps on
fundamental groups. Equivalently
resolve its explicit root charts
\eqref{cdp:eq:C.WH.4} and use the actual reduced
SNC comparison of \ref{cdp:st:C.2.UN.51}. No holomorphic
linearization of an arbitrary given
neighbourhood is inferred.


\item  Assume the root graph is connected.
Orient the original cycle and use
the actual factorization of each
strict root cover through its
original normalization $H_i$.
Its two axes map to the two distinct
marked axis branches on that $H_i$.
The root deck group lies over the
identity on the original space,
so it preserves these labels and
the lifted cyclic orientation.
It acts on the compressed graph
by rotations. This excludes a
reflection by the actual covering
map; an unrelated automorphism of
a Hopf cycle can reverse its
orientation.


\item  The cyclic root group acts transitively
on the strict covers of each original
component. In a connected oriented
cycle the stabilizer of any vertex
is the kernel of the rotation action
and is the same subgroup $K$ for all
vertices. If $r=N/|K|$, there are
exactly $r$ strict covers of every
$H_i$. At a generic point its normal
ramification index is $N/m_i$.
Counting the $N$ sheets gives
\[
               \deg(E_i\to H_i)=m_i/r,
                    \qquad r\mid\gcd_i m_i.
\tag{C.WH.6}\hypertarget{tag-017B}{}\label{cdp:eq:C.WH.6}
\]
At a generic point of the original
double curve \eqref{cdp:eq:C.WH.4} has $g_i$
points above it. There are $r$
connected lifted edges, so their
axis covers have degree $g_i/r$.
They are \'{e}tale by that same
local chart, including at every
point of the compact elliptic axis.


\item  The pointwise axis stabilizer $J_i$
has order $N/g_i$. It contains the
two component inertia subgroups,
whose orders are $N/m_i$ and
$N/m_{i+1}$. In a cyclic group their
product and intersection have orders
the least common multiple and
greatest common divisor of these
numbers. These are $N/g_i$ and $h_i$,
respectively. This proves \eqref{cdp:eq:C.WH.0}
and the entire degree comparison.


\end{enumerate}

\item \par\medskip\noindent\textbf{The finite action on the common open group is by translations.}
Assume now $M=I$. Identify all open
Hopf covers with $G=\C^2/\Lambda$
using the actual linear edge maps.
Let $n_j$ be the outgoing normal
meridian of its $j$th vertex, and
let $h_j$ be the compressed edge
length. The common period is
\[
                u=-\sum_j h_j n_j\in\Lambda.
\tag{C.WH.7}\hypertarget{tag-017C}{}\label{cdp:eq:C.WH.7}
\]
At every vertex the contracting
generator proves (C.HG.4). The
numbers $\ell(i n_j)$ therefore
have one nonzero sign. All $h_j$
are positive, so $\ell(iu)\ne0$.
Writing $u=\delta p$ with $p$
primitive gives the actual elliptic
quotient $G\to E$ along $\C p$.
It is the nonzero bundle of a
degree-zero line $L$ on $E$.


\begin{enumerate}[label={$\langle2\rangle$\arabic*.},leftmargin=1.3em,labelsep=.5em]
\item  The line $L$ is not torsion. If
$L^k=\OO_E$, its pullback by the
multiplication-$k$ elliptic isogeny
is trivial. The actual finite map
$E'\times\C^*\to L^\times$ would
then send $E'\times\{1\}$ to a
compact elliptic curve in $G$.
But $G$ is the axis complement of
an algebraic-dimension-zero Hopf
surface, and its proved curve list
excludes such a curve.


\item  The deck generator is affine on
each common open group. Indeed
isomorphisms of these diagonal
Hopf covers preserving the two
labelled axes lift to diagonal
linear maps on their universal
$\C^2\setminus0$ covers, by the
nonresonant calculation of \ref{cdp:st:C.2.UN.43}.
Their logarithmic maps are affine.
Compatibility with the actual
normal edge maps makes their
common linear part a single map
$R$ preserving $\Lambda$. It has
finite order. The already proved
rotation property gives
$R n_j=n_{\sigma j}$ and
$h_{\sigma j}=h_j$, hence $Ru=u$.
In particular $Rp=p$.


\item  Thus $R$ acts as the identity on
the kernel $\C^*$ of $G\to E$
and induces a finite origin-preserving
elliptic automorphism $\beta$ with
$\beta^*L\simeq L$. If $\beta\ne1$,
the map $\beta^*-1$ on
$\operatorname{Pic}^0(E)$ is an
isogeny, so its finite kernel
would make $L$ torsion. If
$\beta=1$, $R$ is a shear with
both eigenvalues one; a finite
order shear is the identity.
We have therefore proved
\[
                              R=I.
\tag{C.WH.8}\hypertarget{tag-017D}{}\label{cdp:eq:C.WH.8}
\]


\end{enumerate}

\item \par\medskip\noindent\textbf{Integral descent of the period and holomorphic finite transport.}
Write the compressed edge maps in
the common coordinates as
\[
 v\longmapsto v+b_j+l_j\log s,
 \qquad l_j=-h_jn_j/(2\pi i),
\]
and the root generator as
$(j,v,s)\mapsto(\sigma j,v+\tau_j,\zeta_Ns)$.
Its exact compatibility and order give
lattice elements $\lambda_j,\mu_j$ with
\[
 \begin{aligned}
 \tau_{j+1}-\tau_j+b_j-b_{\sigma j}
                   -l_j\,2\pi i/N&=\lambda_j\in\Lambda,\\
 \sum_{a=0}^{N-1}\tau_{\sigma^a j}&=\mu_j\in\Lambda.
 \end{aligned}
\tag{C.WH.9}\hypertarget{tag-017E}{}\label{cdp:eq:C.WH.9}
\]
These identities are obtained by composing
the actual affine maps, with
$\log\zeta_N=2\pi i/N$.
Summing the first identity around the
cycle gives
\[
                   u/N\in\Lambda,\qquad
                 \delta=Nq\quad\hbox{for an integer }q>0.
\tag{C.WH.10}\hypertarget{tag-016Z}{}\label{cdp:eq:C.WH.10}
\]
This is an integral assertion in the
actual rank-three lattice.


\begin{enumerate}[label={$\langle2\rangle$\arabic*.},leftmargin=1.3em,labelsep=.5em]
\item  The constants $b_j$ can first be
adjusted to have sum zero, retaining
the entire central surface and its
deck action. Adding the same multiple
of $n_j$ to $b_j$ on its whole
$\sigma$-orbit changes only a
normal-product constant and does
not change either equation in
\eqref{cdp:eq:C.WH.9}. At a Hopf vertex the
adjacent normal vectors are
complex-linearly independent.
They belong to different rotation
orbits, because \eqref{cdp:eq:C.WH.8} makes
the normal vector constant on
each orbit. These two adjustable
orbit constants span $\C^2$,
so they set $\sum_j b_j=0$.


\item  Normalize the common lattice as in
(C.HG.6). Then $u=\delta(0,2\pi i)$,
$|q_0|<1$ and the section value
$a$ is different from one by the
same explicit embedded-elliptic
argument (C.HG.7).
Construct the compact $I_q$
elliptic suspension of \ref{cdp:st:C.2.UN.48},
where $q=\delta/N$. Its periods are
\[
 \begin{gathered}
 \Lambda(t)+\Z(0,q\log t),\\
 \Lambda(t)=
 \langle(0,2\pi i),(2\pi i,0),
                       (\log q_0,\log a(t))\rangle_{\Z} .
 \end{gathered}
\tag{C.WH.11}\hypertarget{tag-0170}{}\label{cdp:eq:C.WH.11}
\]
Put $t=s^N$. Transport every normal
period, contracting vector and
lattice element in \eqref{cdp:eq:C.WH.9} using
its same integral coefficients in
$\Lambda(t)$. This gives the actual
Hopf families of \ref{cdp:st:C.2.UN.48}, with all
finite characters retained.


\item  There are holomorphic $b_j(t),\tau_j(t)$
through their central values satisfying
\eqref{cdp:eq:C.WH.9} and $\sum b_j(t)=0$.
Here is the precise linear-algebra
reason. In common coordinates the
unknowns occur in a matrix with
constant integer coefficients.
Its right side is a rational linear
combination of the three transported
lattice vectors, because
$l_j\,2\pi i/N=-h_jn_j/N$.
Every rational left-kernel relation
on that matrix vanishes on the
right side at zero. The three
original lattice vectors are
$\Q$-linearly independent, so each
of its rational coefficients is
zero separately. The relation
therefore remains zero after
transport. A fixed linear right
inverse on the matrix image gives
a holomorphic solution, and adding
the original kernel vector makes
its value at zero the prescribed
one. The extra equation $\sum b_j=0$
is included in this matrix.


\item  Use these maps in the finite
compressed plumbing, retaining the
literal hypersurfaces $uv=s^{h_j}$
in its necks. The holomorphic
finite action exists on this entire
normal threefold: \eqref{cdp:eq:C.WH.9} is
exactly its overlap equation and
its $N$th power is a deck period.
On component cores it is the
diagonal translation isomorphism
between the transported Hopf
quotients. On necks it is the
induced line map preserving the
product up to $\zeta_N^{h_j}$.
Thus it extends across their
zero sections and respects the
actual equation. All central
components, axis maps and finite
characters are the original ones.


\end{enumerate}

\item \par\medskip\noindent\textbf{The quotient is smooth with the original multiplicities.}
This assertion is checked on the
actual local invariant rings.
At an axis put $a=gA$, $b=gB$,
$\gcd(A,B)=1$, $h=N/(gAB)$.
The pointwise axis stabilizer
$J$ has order $hAB$. In the chart
$uv=s^h$, its subgroup of order
$h$ fixes $u,v$ and rotates $s$.
The invariant ring of this
subgroup is $\C\{u,v\}$.
Its remaining quotient acts on
$u,v$ as the independent groups
of orders $B,A$. Therefore
\[
 \bigl(\C\{u,v,s\}/(uv-s^h)\bigr)^J
           =\C\{u^B,v^A\}=\C\{x,y\},
 \qquad t=s^N=x^a y^b .
\tag{C.WH.12}\hypertarget{tag-0171}{}\label{cdp:eq:C.WH.12}
\]
For an explicit generator choose
integers $c,d$ with $cA+dB=1$.
It has characters
$(\zeta_B^c,\zeta_A^d,\zeta_{hAB})$
on $(u,v,s)$, so their product
indeed matches the character of
$s^h$. These characters are
unchanged in the transported
finite family: they are finite
roots of unity and have the
specified central values.
The stabilizer action along the
axis is trivial; the residual
axis covering group acts freely.
Thus \eqref{cdp:eq:C.WH.12}, followed by that
\'{e}tale quotient, is a smooth
threefold chart. At a general
component point the inertia is
the normal reflection of order
$N/m_i$, so the quotient is
smooth there as well. Away from
the central support the action
moves $s$, hence is free.
This proves smoothness everywhere
near the fibre and retains
precisely the multiplicities
$m_i$. Properness follows from
the finite quotient of the
proper compressed plumbing.


\begin{enumerate}[label={$\langle2\rangle$\arabic*.},leftmargin=1.3em,labelsep=.5em]
\item  The central quotient is the
given reduced surface. The
transport left its cover surfaces,
axis maps and all finite action
maps unchanged at zero; these
are the actual normalization
and quotient maps constructed
from \eqref{cdp:eq:C.WH.1}--\eqref{cdp:eq:C.WH.4}.
The adjusted orbit constants
change only the normal-product
maps, not this quotient surface.


\item  The punctured quotient has the
periods \eqref{cdp:eq:C.WH.11}, since
$\log s=(\log t)/N$ and the
change of branch adds $u/N\in\Lambda$.
The logarithmic coframe descends
because of \eqref{cdp:eq:C.WH.8}.
For clarity, a possible translation
torsor does not obstruct the
actual comparison. The quotient
is a holomorphic torsor under
the marked family of tori with
these periods: local origins
differ by translations, since
the coframe is fixed. On a
punctured disc the sheaf sequence
\[
 0\longrightarrow\mathscr L
   \longrightarrow\OO^{\oplus2}
   \longrightarrow\mathscr A
   \longrightarrow0
\tag{C.WH.13}\hypertarget{tag-0172}{}\label{cdp:eq:C.WH.13}
\]
has $H^1(\OO^{\oplus2})=0$ by
the holomorphic Cousin theorem
on the Stein punctured disc,
and $H^2(\mathscr L)=0$ because
it retracts to a circle.
Thus $H^1(\mathscr A)=0$.
Lifting the translation cocycle
through this sequence produces
an actual holomorphic section
and an isomorphism with the
period family. This identifies
the punctured quotient with
the compact suspension in
\eqref{cdp:eq:C.WH.11}, as a holomorphic
family, including all affine
translation data.


\item  The proper nested-annulus
replacement of \ref{cdp:st:C.2.UN.48} now
globalizes this smooth quotient
over the same $\PP^1$.
Every central Hopf component
and double elliptic curve
has Euler number zero.
Every other reduced support
is unchanged from the elliptic
suspension and also has Euler
number zero. This completes
the compact assertion and
the full criterion.


\end{enumerate}

\item \par\medskip\noindent\textbf{Explicit unequal multiplicities.}
Fix nonresonant contractions
$\alpha,\beta$ and put
\[
 H=(\C^2\setminus0)/
  \langle(\alpha z,\beta w),
              (\zeta_n z,\zeta_n^{-1}w)\rangle .
\tag{C.WH.14}\hypertarget{tag-0173}{}\label{cdp:eq:C.WH.14}
\]
Take the two-copy inversion
plumbing of this exact secondary
Hopf surface. Let $g(s)=\zeta_n s$,
let $g$ be the identity on the
first copy, and on the second
put $g(Z,W)=(Z,\zeta_n W)$.
The first neck map
$(z,w)\mapsto(z^{-1},s/w)$
intertwines these actions.
The other neck gives instead
$(Z,W)\mapsto(\zeta_n Z,W)$;
the two differ by the actual
anti-diagonal deck element.
Thus they define a genuine
order-$n$ action.


\begin{enumerate}[label={$\langle2\rangle$\arabic*.},leftmargin=1.3em,labelsep=.5em]
\item  At every double curve it is
$(x,y)\mapsto(x,\zeta_n y)$,
in the actual line-bundle chart.
Its invariant coordinates are
$x,Y=y^n$ and the full quotient
equation is
\[
                              t=s^n=x^nY .
\tag{C.WH.15}\hypertarget{tag-0174}{}\label{cdp:eq:C.WH.15}
\]
The quotient is smooth and its
two multiplicities are $(1,n)$.
The first component is exactly
$H$. On the second, the finite
anti-diagonal group and this
extra action generate the two
independent scalar groups.
The coordinate-power map
$(Z,W)\mapsto(Z^n,W^n)$
therefore identifies its
quotient with the primary
Hopf surface
$H(\alpha^n,\beta^n)$.
The axis identifications are
the actual maps $z\mapsto z^{-n}$
from the corresponding axis
quotients; they are isomorphisms,
since their kernel is exactly
the given finite angular group.


\item  This example also has an
independent compact comparison.
Its common coordinates are
$(b,y)=(zw,w^n)$, with periods
$q_0=\alpha\beta$ and
$a=\beta^n$. The normal cycle
has primitive divisor $n$.
Start with the compact $I_1$
suspension with section $a(t)$,
pull back by $t=s^n$, and use
$\beta(s)=a(s^n)^{1/n}$ with
the given value at zero,
$\alpha(s)=q_0/\beta(s)$.
The two-copy replacement and
the displayed finite action
are equivariant. On the
punctured common quotient this
action is exactly the base
deck action. Its smooth
quotient consequently replaces
a tube in the original compact
suspension, giving \eqref{cdp:eq:C.WH.15}
over a compact rational base.


\item  For the primary triangle
(C.HG.1), let the base action
be $s\mapsto-s$ and the three
component actions be
\[
            (z,w),\qquad(z,-w),\qquad(-z,w).
\tag{C.WH.16}\hypertarget{tag-0175}{}\label{cdp:eq:C.WH.16}
\]
The three actual maps
(C.HG.16) intertwine them.
The action on the middle
double curve is a free
elliptic translation; on
the other two curves it
is a normal reflection.
The smooth quotient has
multiplicities $(1,2,2)$
and primary multipliers
\[
 (e^{-1},e^{-\sqrt2}),\quad
 (e^{-\sqrt2},e^{-2-2\sqrt2}),\quad
 (e^{-2-2\sqrt2},e^{-1}).
\tag{C.WH.17}\hypertarget{tag-0176}{}\label{cdp:eq:C.WH.17}
\]
Its axis maps are inversion,
identity and inversion.
The root primitive divisor
is two, so the same quadratic
base change of a compact
$I_1$ suspension gives its
compact replacement.


\item  The exact normalization,
inertia and invariant-ring
exponents, 300 resolution
charts, sixty anti-diagonal
secondary indices, and the
three equivariance identities
are checked in
\path{code/D_weighted_hopf_cycles.py}
and its saved output.



\end{enumerate}

\end{enumerate}
\end{proof}






\begin{proposition}[Complete reduced Hopf-star support criterion]\hypertarget{tag-01DY}{}\label{cdp:st:C.2.UN.53}
\StatementRecord{cdp:st:C.2.UN.53}{proved}{01DY}{\SeriesRef{cdp:st:C.2.UN.38}, \SeriesRef{cdp:st:C.2.UN.43}, \SeriesRef{cdp:eq:C.AS.6}, \SeriesRef{cdp:eq:C.AS.7}, \SeriesRef{cdp:eq:C.AS.9}, \SeriesRef{cdp:st:C.2.UN.28}, \SeriesRef{cdp:eq:C.AS.13}, \SeriesRef{cdp:st:C.2.UN.31}, \SeriesRef{cdp:st:C.2.UN.36}, \SeriesRef{cdp:st:C.2.UN.49}, \SeriesRef{cdp:eq:C.AS.12}, \SeriesRef{cdp:eq:C.AS.15}, \SeriesRef{cdp:eq:C.AS.16}, \SeriesRef{cdp:eq:C.AS.18}, \SeriesRef{cdp:eq:C.AS.17}, \SeriesRef{cdp:st:C.2.UN.50}}

Let $S$ be a specified compact connected SNC surface with smooth
components $P,H_1,\ldots,H_n$, $n\ge1$. Each $H_i$ is a primary
or secondary Hopf surface. Its intersection with $P$ is one
connected smooth curve $J_i$, the $J_i$ are disjoint on $P$,
and there are no other intersections. We impose no initial
classification or elliptic-fibration hypothesis on $P$.

The following two occurrence assertions are equivalent: $S$
is the full reduced central fibre of a proper two-torus
degeneration with smooth total space; and $S$ is such a fibre
of a smooth compact threefold over $\PP^1$, with every other
fibre a smooth two-torus. They hold precisely for the following
explicit data.

Every leaf pair is isomorphic to one of
\[
 \begin{gathered}
 (H^{\rm diag}_{r_i}(\beta_i)/\mu_{n_i},J_i),\qquad
 (H^{\rm nl}_{r_i}(\beta_i)/\mu_{n_i},J_i),\\
 r_i\ge1,\qquad 0<|\beta_i|<1,\qquad \gcd(r_i,n_i)=1 ,
 \end{gathered}                                             \tag{C.AS.1}\hypertarget{tag-00AG}{}\label{cdp:eq:C.AS.1}
\]
where the two contractions on $\C^2\setminus\{0\}$ are,
respectively,
\[
 (z,w)\longmapsto(\beta_i^{r_i}z,\beta_iw),\qquad
 (z,w)\longmapsto(\beta_i^{r_i}z+w^{r_i},\beta_iw),
 \quad
 \zeta(z,w)=(\zeta^{r_i}z,\zeta w),                         \tag{C.AS.2}\hypertarget{tag-00AR}{}\label{cdp:eq:C.AS.2}
\]
and $J_i$ is the image of $w=0$. The normal lines on every
double curve are dual. The central component has a
fixed-Jacobian, identity-integral-monodromy elliptic
quasi-bundle
\[
                         q:P\longrightarrow B             \tag{C.AS.3}\hypertarget{tag-00AS}{}\label{cdp:eq:C.AS.3}
\]
over an elliptic curve. Its marked full fibres are $r_iJ_i$;
there are no other multiple fibres and every fibre support
is smooth. The actual inverse logarithmic transform $T_0$
defined by (C.SN.3)--(C.SN.6), retaining all the specified
axis identifications and the entire analytic torsor,
satisfies
\[
                         c(T_0)=0\in H^2(B,\Lambda_F).     \tag{C.AS.4}\hypertarget{tag-00AT}{}\label{cdp:eq:C.AS.4}
\]
Finally there is a positive integer $R$ such that
\[
 \begin{cases}
 r_i=R,&H_i\text{ is of the nonlinear form in \eqref{cdp:eq:C.AS.2}},\\
 r_i<R,&H_i\text{ is of the diagonal form in \eqref{cdp:eq:C.AS.2}}.
 \end{cases}                                             \tag{C.AS.5}\hypertarget{tag-00AU}{}\label{cdp:eq:C.AS.5}
\]
When every leaf is diagonal, the last test is satisfied by
any $R>\max_i r_i$. When a nonlinear leaf occurs, it
determines $R$, and every other nonlinear leaf has that
same resonance.

The conclusion retains the full isomorphism class of $S$.
In particular it retains arbitrary flat torsor parameters,
finite secondary orders, independent primitive angular
directions, and the given elliptic gluing maps. Every support
in the compact realization has Euler characteristic zero.
\end{proposition}



\begin{proof}\leavevmode\par



 
\begin{enumerate}[label={$\langle1\rangle$\arabic*.},leftmargin=0pt,labelsep=.5em]
\item \par\medskip\noindent\textbf{The canonical identity on an arbitrary Hopf leaf.}
We first use an actual reduced smoothing. Let the divisor
of a local canonical evaluation generator be
$k_PP+\sum_i k_iH_i$. Its existence and the fact that it
has only vertical zeroes are the exact direct-image
calculation in \ref{cdp:st:C.2.UN.38}. Adjunction on one leaf gives
\[
                  K_H=\OO_H(-aJ),\qquad
                  a=k_H-k_P+1\in\Z.                      \tag{C.AS.6}\hypertarget{tag-00AV}{}\label{cdp:eq:C.AS.6}
\]
We classify this equality on the actual pair $(H,J)$,
including the full secondary deck group.


\begin{enumerate}[label={$\langle2\rangle$\arabic*.},leftmargin=1.3em,labelsep=.5em]
\item  Take a normal primary cyclic subgroup of the Hopf deck
group $\Gamma$. Its generator $F$ is central: conjugation
can send it only to $F$ or $F^{-1}$, while the nonzero
logarithm of the modulus of its derivative determinant
excludes the latter. This is the exact normal-core
argument of \ref{cdp:st:C.2.UN.43}. We therefore have a finite
unramified primary cover $H'\to H$, and the full inverse
image $D$ of $J$ is a disjoint union of smooth elliptic
curves. The finite quotient acts transitively on these
curves, since $J$ is connected. Pullback gives
\[
                           K_{H'}=-aD.                  \tag{C.AS.7}\hypertarget{tag-00AW}{}\label{cdp:eq:C.AS.7}
\]


\item  For a nonlinear primary contraction the curve computation
of \ref{cdp:st:C.2.UN.38} gives its unique axis, and the full
centralizer computation of \ref{cdp:st:C.2.UN.43} gives exactly the
nonlinear quotient in \eqref{cdp:eq:C.AS.1}--\eqref{cdp:eq:C.AS.2}.
The invariant form $dz\wedge dw/w^{r+1}$ has divisor
$-(r+1)J$, on the quotient as well. Since $\OO_H(J)$ has
infinite order, \eqref{cdp:eq:C.AS.6} forces $a=r+1$.


\item  Consider a diagonal primary contraction. Its invariant
curve computation can be read directly from an equation
on $\C^2$. Closure across the origin and removal of the
nonconstant unit in its transformation give an eigenfunction
of $F$. Comparison of its power-series monomials gives
the axes in the absence of a multiplicative resonance.
In that case the complete centralizer is diagonal, so
it cannot interchange the axes. Thus $D$ is one axis.
If, for example, it is $w=0$, the exact automorphy
characters in \eqref{cdp:eq:C.AS.7} give
$\alpha\beta=\beta^a$, which is the resonance
$\alpha=\beta^{a-1}$. The contraction moduli make
$a-1>0$. This contradicts the assumed absence of a
resonance and accounts for that case through an actual
character equality.


\item  In the resonant case pass to a further central cyclic
subgroup to remove the finite root-of-unity discrepancy.
Its generator has the form
\[
                 F(z,w)=(\gamma^p z,\gamma^q w),
       \qquad p,q\ge1,\quad\gcd(p,q)=1,\quad |\gamma|<1 .
                                                               \tag{C.AS.8}\hypertarget{tag-00AX}{}\label{cdp:eq:C.AS.8}
\]
Here the axes have defining weights $p,q$, and every
other irreducible curve has equation $z^q=cw^p$ of
weight $pq$. This follows from the same eigenfunction
calculation: after an axis monomial is removed, a
weighted homogeneous polynomial factors into these
binomials over $\C$. Its irreducible factor is unique
because $p,q$ are coprime. The curves are smooth away
from the origin, as required for the elliptic curves
on $H'$.


\item  The canonical weight is $p+q$. If $D$ is an orbit of
$m$ non-axis curves, \eqref{cdp:eq:C.AS.7} therefore gives
\[
                              ampq=p+q.                 \tag{C.AS.9}\hypertarget{tag-00AY}{}\label{cdp:eq:C.AS.9}
\]
The same calculation with a single axis replaces $mpq$
by $p$ or $q$. These are genuine character equalities:
a purported extra holomorphic unit and its inverse
extend across the origin, and their nonzero value there
forces the identical character. In particular $a>0$.


\item  For $p,q>1$, coprimality gives $pq>p+q$, so \eqref{cdp:eq:C.AS.9}
is impossible. A single axis would require $p\mid q$
or $q\mid p$, also impossible. In this case the
centralizer preserves each axis, so an orbit containing
both axes cannot occur. For $p=r>1,q=1$, the curves
$z=cw^r$, including $z=0$, all have weight $r$.
The triangular centralizer can permute these curves;
an orbit of $m$ of them has total weight $mr$ and
would require $amr=r+1$, hence $r\mid1$.
The only possibility is the weight-one axis $w=0$,
and then $a=r+1$.


\item  We retain the full deck group in this last calculation.
Commuting with $(\gamma^rz,\gamma w)$ gives exactly
\[
              (z,w)\longmapsto(Az+Cw^r,Bw).             \tag{C.AS.10}\hypertarget{tag-00AH}{}\label{cdp:eq:C.AS.10}
\]
Its induced action on $z/w^r$ is affine. The central
cyclic subgroup acts trivially there; hence the image
is a finite affine group. Averaging the orbit of zero
gives an actual common fixed point. The corresponding
coordinate change $z\mapsto z-cw^r$ makes every element
diagonal, while retaining the marked axis. Equation
\eqref{cdp:eq:C.AS.6} then says, for every deck element,
$AB=B^{r+1}$, so $A=B^r$. Thus $\Gamma$ is the
weighted scalar group with an infinite contracting
generator $\beta$ and finite roots of unity
$\mu_n$. Freeness on the marked axis gives
$\gcd(r,n)=1$. This is exactly the diagonal pair in
\eqref{cdp:eq:C.AS.1}--\eqref{cdp:eq:C.AS.2}.


\item  It remains to check the scalar case $p=q=1$,
where the full secondary group need not initially
be diagonal. Every component of $D$ lifts to a
line, and \eqref{cdp:eq:C.AS.7} gives $am=2$. If $m=1$,
the finite projective group fixes that line.
On its affine complement, averaging gives a
second fixed point, so the linear group becomes
diagonal in coordinates retaining the first line.
The character equality $AB=B^2$ forces $A=B$
for every element. We obtain the case $r=1$
of \eqref{cdp:eq:C.AS.2}, with scalar finite action.
If $m=2$, then $a=1$ and an element of the full
deck group interchanges the two lines. Choose
them as the two coordinate axes. The unique
log-canonical form up to a constant is
\[
                             \frac{dz}{z}\wedge\frac{dw}{w}.
                                                               \tag{C.AS.11}\hypertarget{tag-00AI}{}\label{cdp:eq:C.AS.11}
\]
Every such interchange multiplies this form
by $-1$. A form with exactly the same divisor
differs by a nowhere-zero holomorphic function
on the compact primary Hopf surface, hence by
a constant. It therefore cannot descend.
This contradicts the actual equality
$K_H+J=0$ in \eqref{cdp:eq:C.AS.6}. Thus this last orbit
does not satisfy \eqref{cdp:eq:C.AS.6}. We have proved,
without an a priori restriction on the secondary
group, precisely the leaf-pair list \eqref{cdp:eq:C.AS.1},
and in every surviving case
\[
                              k_i-k_P=r_i\ge1.          \tag{C.AS.12}\hypertarget{tag-00AJ}{}\label{cdp:eq:C.AS.12}
\]


\end{enumerate}

\item \par\medskip\noindent\textbf{The central component is forced into the stated class.}
Adjunction on $P$ now gives the exact effective
canonical divisor
\[
                         K_P=\sum_i(r_i-1)J_i .         \tag{C.AS.13}\hypertarget{tag-00AK}{}\label{cdp:eq:C.AS.13}
\]
Thus $\kappa(P)\ge0$ and $p_g(P)\ge1$.
\ref{cdp:st:C.2.UN.28} gives $p_g(P)=1$. The surface $P$
is minimal. Indeed, a $(-1)$-curve is rational,
so is distinct from every elliptic $J_i$.
Its intersections with all of them are
nonnegative, contradicting $K_P\cdot E=-1$
and \eqref{cdp:eq:C.AS.13}.


\begin{enumerate}[label={$\langle2\rangle$\arabic*.},leftmargin=1.3em,labelsep=.5em]
\item  We now specify the covering comparison needed
to apply \ref{cdp:st:C.2.UN.31} to this original component.
Its full-fibre multiplicity is one. In the
normalized root base change $t=s^N$, a generic
chart is $x=s^N$ after a holomorphic unit is
absorbed. There is exactly one strict covering
component and its map to $P$ has degree one.
Equivalently, the total generic degree of the
strict component covers is the original
multiplicity, as calculated in \ref{cdp:st:C.2.UN.36}.
After resolution it remains birational to $P$.
Both minimal models are therefore the same
smooth surface. The actual root-cover conclusion
of \ref{cdp:st:C.2.UN.31} thus applies to $P$ itself:
it is a two-torus, or a fixed-Jacobian,
identity-monodromy elliptic quasi-bundle over
a curve of genus at most one, with only smooth
multiple fibres. In the latter description
$p_g(P)=g(B)$, so the base is elliptic.
No assertion about an arbitrary quotient
component has been substituted for this
degree-one comparison.


\item  We also verify that all the given curves
$J_i$ are fibres of one such fibration.
On a torus, an embedded elliptic curve is a
translate of an elliptic subgroup: lift its
map to universal covers and differentiate;
the derivative is periodic and holomorphic,
hence constant. Two elliptic subgroups with
distinct tangent lines give a surjective
isogeny from their product to the two-torus.
Every two translates then meet. The disjoint
$J_i$ consequently have one common tangent
line and are fibres of the actual elliptic
quotient.


\item  In the quasi-bundle description over an
elliptic base the canonical divisor is
$\sum(m-1)F_{\rm red}$, by the exact root
chart calculation in \ref{cdp:st:C.2.UN.49}.
If a multiple fibre occurs, a horizontal
$J_i$ would have positive intersection with
this canonical divisor. But its torsion
normal line gives $J_i^2=0$, and adjunction
gives $K_P\cdot J_i=0$. Thus every $J_i$
is vertical. If no multiple fibre occurs,
we have a principal elliptic bundle.
A nonzero integral Chern vector excludes
a horizontal elliptic curve: its finite
unramified map to the base would pull the
bundle back to one with a section, while
its Chern vector would be multiplied by
the nonzero covering degree. If that Chern
vector is zero, the exact Cech lattice
calculation (C.NS.14) identifies $P$ with
a torus, reducing to the preceding argument.
This proves \eqref{cdp:eq:C.AS.3} with every marked
curve retained.


\item  The normal of the leaf axis has exact
torsion order $r_i$, including its finite
secondary character by (C.SN.3)--(C.SN.4).
Normal duality in the reduced smoothing
therefore gives full marked fibre $r_iJ_i$
on $P$. Subtracting the marked terms of
its canonical root-chart formula from
\eqref{cdp:eq:C.AS.13} leaves an effective sum of
unmarked multiple fibres linearly
equivalent to zero. Such a nonzero
effective divisor on a compact connected
complex surface cannot be principal:
its defining meromorphic function would
have no poles and hence be constant.
Consequently there are no unmarked
multiple fibres.


\end{enumerate}

\item \par\medskip\noindent\textbf{The exact topology and the one period shared by all leaves.}
The primary and secondary inverse-transform
charts (C.SN.3)--(C.SN.6) apply to both
contractions in \eqref{cdp:eq:C.AS.2}: they use their
identical axis normal character. For a
nonlinear leaf the explicit deformation
of its resonant coefficient to zero gives
the relative smooth-pair comparison
proved in \ref{cdp:st:C.2.UN.49}. For a diagonal leaf
the product complement is already exact.
Thus every Hopf complement kills the
same displayed complementary meridian,
and each oriented core two-torus becomes
the ordinary fibre class of the actual
inverse transform $T_0\to B$.
The nearby smoothing is its underlying
four-manifold. Its fundamental group is
the principal-bundle presentation
(C.NS.13), with the unaltered integer
Chern vector. A two-torus has abelian
fundamental group; the injected fibre
lattice in that presentation therefore
forces \eqref{cdp:eq:C.AS.4}.


\begin{enumerate}[label={$\langle2\rangle$\arabic*.},leftmargin=1.3em,labelsep=.5em]
\item  Let $\gamma$ be this one transported
integral core class. Divide the relative
canonical generator by $t^{k_P}$.
On the $i$th Hopf core its leading
term is $t^{r_i}\eta_i$, with
\[
                   \eta_i=c_i\,\frac{dz\wedge dw}{w^{r_i+1}},
                          \qquad c_i\ne0.               \tag{C.AS.14}\hypertarget{tag-00AL}{}\label{cdp:eq:C.AS.14}
\]
This is the actual coefficient calculation
of (C.NS.11), together with \eqref{cdp:eq:C.AS.12}.
Put $U=z/w^{r_i}$ and $V=\log w$.
The primitive angular period of the
secondary quotient is $(0,2\pi i/n_i)$.
For a nonlinear leaf the other period is
$(\beta_i^{-r_i},\log\beta_i)$, so
\[
       \int_{\gamma}\omega_t
          =t^{r_i}\left(\frac{\pm2\pi i c_i\beta_i^{-r_i}}{n_i}
                         +O(t)\right).                 \tag{C.AS.15}\hypertarget{tag-00AM}{}\label{cdp:eq:C.AS.15}
\]
Its coefficient is nonzero. For a diagonal
leaf the other period is
$(0,\log\beta_i)$: the core lies at
constant $U$, and $c_i\,dU\wedge dV$
has zero integral there. Consequently
\[
                         \int_\gamma\omega_t=O(t^{r_i+1})
                  \quad\hbox{at a diagonal leaf}.        \tag{C.AS.16}\hypertarget{tag-00AN}{}\label{cdp:eq:C.AS.16}
\]
These are periods of the same actual
class in the same nearby fibre. They
are holomorphic functions: integration
of a closed relative holomorphic form
over the transported compact cycle is
holomorphic by differentiation and
Stokes' theorem. The displayed leading
terms also give their holomorphic
extension at zero.


\item  If a nonlinear leaf occurs, \eqref{cdp:eq:C.AS.15}
gives one finite vanishing order $R$.
Every nonlinear $r_i$ equals it, and
\eqref{cdp:eq:C.AS.16} gives $r_i<R$ for every
diagonal leaf. If all leaves are
diagonal, no finite vanishing order
is asserted for the given smoothing;
the explicit numerical test \eqref{cdp:eq:C.AS.5}
is satisfied by choosing
$R=1+\max_i r_i$. This proves its
necessity without imposing an unproved
nonzero period on that family.


\end{enumerate}

\item \par\medskip\noindent\textbf{One compact construction realizes all the specified data.}
We now use data \eqref{cdp:eq:C.AS.1}--\eqref{cdp:eq:C.AS.5}.
The zero-Chern inverse transform is
the actual torus $T_0$ with its marked
fibres, by (C.NS.14). Take its explicit
compact $\OO_{\PP^1}(1)^{\oplus2}$
twistor family (C.NS.15), and pull it
back by $t=s^R$. For the $i$th marked
curve the exact four-period coordinate
change (C.NS.16)--(C.NS.18), with
$\gamma_i=\beta_i^{n_i}$, gives
\[
 \mathcal T_s=M_{i,s}^{\times}/
   \langle(p,\ell)\mapsto(p+d_i s^R,\gamma_i\ell)\rangle,
                    \qquad d_i\ne0 .                  \tag{C.AS.17}\hypertarget{tag-00AO}{}\label{cdp:eq:C.AS.17}
\]
The non-torsion flat line and every
original primitive fibre direction
are retained by these actual period
maps.


\begin{enumerate}[label={$\langle2\rangle$\arabic*.},leftmargin=1.3em,labelsep=.5em]
\item  At this marked point use weight
$(r_i,1)$ and tensor the flat line
by $\OO(n_iE_i)$, where $E_i$ is
the weighted exceptional divisor.
The same free coarse quotient as
(C.SN.9) has coordinates
$p-p_i=x^{r_i}z$, $s=xw$.
The exact infinite deck map is now
\[
 \begin{split}
 (z,w,x)\longmapsto\bigl(&
   \beta_i^{r_i}(z+d_i x^{R-r_i}w^R),\\
   &\beta_iw,\ \beta_i^{-1}x\bigr).
 \end{split}                                             \tag{C.AS.18}\hypertarget{tag-00AP}{}\label{cdp:eq:C.AS.18}
\]
The exponent $R-r_i$ is nonnegative.
The added monomial has weighted degree
$r_i$, so this map commutes with the
finite action of weights
$(r_i,1,-1)$ modulo $n_i$. It preserves
$s=xw$ and changes the original base
coordinate by exactly $d_i s^R$.
Thus \eqref{cdp:eq:C.AS.18} is the actual lift
of \eqref{cdp:eq:C.AS.17}.


\item  On $x=0$, equality $R=r_i$ gives the
specified nonlinear contraction after
$Z=z/(\beta_i^{r_i}d_i)$. Strict
inequality $R>r_i$ gives the specified
diagonal contraction. On $w=0$ the
map is always
$(z,x)\mapsto(\beta_i^{r_i}z,\beta_i^{-1}x)$,
so the main component and its outer
trivialization are exactly
(C.SN.4)--(C.SN.5), independently of
$R$. Every given central component
and every marked normal cover is
therefore the one prescribed.


\item  We record the analytic bounds to
verify the simultaneous compact
replacement in this varying-weight
case. On a small closed disc,
$|s|^R\le C|s|^{r_i}$ for every $i$.
The weighted norm comparison of
(C.NH.9)--(C.NH.10) consequently
gives at most
$C(1+|j|)^{n_i/r_i}$ distortion
for the $j$th translation lift.
The unitary flat metric used in
\ref{cdp:st:C.2.UN.50} has scalar multiplier
$|\gamma_i|\exp(-\operatorname{Re}(c_i(s)d_i s^R))$,
uniformly less than one after
shrinking the disc. Its exponential
contraction dominates that polynomial.
The same compact norm annulus proves
properness and freeness of the
infinite quotient. Coprimality in
\eqref{cdp:eq:C.AS.1} makes its finite coarse
quotient smooth.


\item  Outside the new Hopf component the
canonical exceptional section gives
the actual isomorphism with the
original twistor family minus its
marked central elliptic curve.
In a compact Hopf fundamental domain,
$z,w$ remain bounded and $x\to0$.
The old flat coordinate is proportional
to $x^{-n_i}$. Putting it in a compact
norm annulus uses
$O(1+|\log|x||)$ iterates, with total
base displacement
\[
                   O\bigl(|s|^R(1+|\log|x||)\bigr)
                          \longrightarrow0,
                         \qquad s=xw.                   \tag{C.AS.19}\hypertarget{tag-00AQ}{}\label{cdp:eq:C.AS.19}
\]
All cluster points consequently lie
on the specified original central
elliptic curve. The disjoint-curve
gluing proof of (C.NS.20) therefore
applies to all these insertions
simultaneously. It proves a smooth
Hausdorff proper local family, and
the unchanged compact twistor
complement makes the total space
compact over $\PP^1$.


\item  The central main component is
precisely $P$, because all its
analytic torsor transitions have
been reversed by the exact inverse
charts. The diagonal automorphisms
$(z,w)\mapsto(a^{r_i}z,aw)$ commute
with either contraction and with
its finite action; they retain
every specified constant in the
axis gluing. The linear parts were
already retained by the chosen
normal-cover lattice bases.
Thus the full central support is
the specified $S$. Its equations
are $xw=s$ in the free coarse
charts, so every component has
multiplicity one. All other fibres
are the original smooth two-tori.
The Euler number of the central
support is zero by additivity:
each quasi-bundle, Hopf leaf and
double elliptic curve has Euler
number zero.


\item  The exact symbolic checks in
\path{code/root_mixed_hopf_stars.py}
retain the weighted base map, its
Jacobian, the finite-action weights
and the common period exponent.
The finite integer search there
checks the degree exclusions in
\eqref{cdp:eq:C.AS.9}; the preceding integer
argument proves them for all weights.



\end{enumerate}

\end{enumerate}
\end{proof}






\begin{proposition}[Every reduced all-Hopf support has one cycle, with an exact algebraic-dimension-zero classification]\hypertarget{tag-01DZ}{}\label{cdp:st:C.2.UN.54}
\StatementRecord{cdp:st:C.2.UN.54}{proved}{01DZ}{\SeriesRef{cdp:st:C.2.UN.48}, \SeriesRef{cdp:st:C.2.UN.37}, \SeriesRef{cdp:st:C.2.UN.38}, \SeriesRef{cdp:st:C.2.UN.53}, \SeriesRef{cdp:st:C.2.UN.43}, \SeriesRef{cdp:st:C.2.UN.44}, \SeriesRef{cdp:eq:C.HR.6}, \SeriesRef{cdp:eq:C.HR.7}, \SeriesRef{cdp:st:C.2.UN.46}}

Let $f:U\to\Delta$ be a proper degeneration of connected complex
two-tori with smooth total space, whose full central fibre
$S=\bigcup_{v\in V}H_v$ is reduced and simple normal crossing.
Suppose that every $H_v$ is a primary or secondary Hopf surface.
Then every connected double curve is elliptic, there are no triple
points, and the connected incidence graph has exactly one independent
cycle.

For an oriented edge $e$ from $v$ to $w$, let $C_e$ be its actual
double curve and let
\[
 r_{v,e}:H^1(H_v,\OO_{H_v})\longrightarrow H^1(C_e,\OO_{C_e})
                                                        \tag{C.HR.1}\hypertarget{tag-00N0}{}\label{cdp:eq:C.HR.1}
\]
be restriction through the specified inclusion (and through the
specified gluing map at the other endpoint). Every such map is a
nonzero map between one-dimensional complex vector spaces. The
weighted incidence map
\[
 d_1:\bigoplus_vH^1(H_v,\OO_{H_v})
       \longrightarrow\bigoplus_eH^1(C_e,\OO_{C_e}),\qquad
 (d_1a)_e=r_{v,e}a_v-r_{w,e}a_w
                                                        \tag{C.HR.2}\hypertarget{tag-00N1}{}\label{cdp:eq:C.HR.2}
\]
has rank $|V|-1$. In any choices of nonzero bases, the product of
$r_{v,e}/r_{w,e}$ around the oriented cycle is one. This product
retains the actual restriction maps; it is unchanged by changing
the chosen bases. The full coherent dimensions are
\[
                    h^q(S,\OO_S)=(1,2,1),\qquad q=0,1,2.
                                                        \tag{C.HR.3}\hypertarget{tag-00N2}{}\label{cdp:eq:C.HR.3}
\]

The one-dimensional space $H^0(S,\omega_S)$ has a more precise
description. Its nonzero section is represented by logarithmic
two-forms on the components. Their residues are nonzero on exactly
the cycle edges and zero on every edge of an attached tree. The
forms themselves are nonzero on exactly the cycle components.
For a canonical evaluation generator with vertical orders $k_v$,
the minimum of the $k_v$ is therefore attained exactly at the
cycle components.

In particular, if every $H_v$ has algebraic dimension zero, the
graph has no attached trees and every component is a diagonal Hopf
surface. There are at least two components. This gives the following
complete criterion for this entire reduced branch, without initially
assuming either diagonality or a cycle. A specified compact reduced
SNC surface all of whose components are algebraic-dimension-zero
Hopf surfaces occurs as the full fibre of a smooth-total torus
degeneration precisely when its components are diagonal, its graph
is a cycle, its actual axis normal lines are dual, and the integral
edge product $M$ of (C.HC.3)--(C.HC.6) is the identity. Every such
specified support has the compact realization of \ref{cdp:st:C.2.UN.48}, with
all support Euler characteristics zero.
\end{proposition}



\begin{proof}

\begin{enumerate}[label={$\langle1\rangle$\arabic*.},leftmargin=*]
\item \par\medskip\noindent\textbf{The curves and the actual restriction maps.}
A Hopf surface has $H^2(H,\Q)=0$. For an irreducible compact curve
$C\subset H$, adjunction gives $p_a(C)=1$. A map from $\PP^1$ to
$H$ lifts to its universal cover $\C^2\setminus\{0\}$, and its two
coordinate functions are constant. The normalization of $C$ hence
cannot be rational. The equality
$p_a(C)=g(C^\nu)+\delta(C)$ now gives $g(C^\nu)=1$ and
$\delta(C)=0$. Thus $C$ is smooth elliptic. Two distinct such
curves have intersection number zero, since $H^2(H,\Q)=0$;
positivity of every local intersection forces them to be disjoint.
These are the same explicit curve comparisons used in \ref{cdp:st:C.2.UN.37}.
It follows at once that the present support has only elliptic
double curves and no triple point.

\item We prove nonvanishing in \eqref{cdp:eq:C.HR.1}, including secondary surfaces.
First let $H'=(\C^2\setminus\{0\})/\langle F\rangle$ be a primary
Hopf surface. The constant additive deck cocycle $1$ represents a
nonzero class in $H^1(H',\OO_{H'})$. Indeed, a trivialization would
be a holomorphic function $b$ on $\C^2\setminus\{0\}$ satisfying
$b(Fz)-b(z)=1$. Hartogs extension to the origin and evaluation there
give a contradiction. Since $h^1(H',\OO_{H'})=1$, this class spans
the group.

\item For every elliptic curve $C'\subset H'$, its lift can have finitely
many branches permuted by $F$. Choose a positive power $F^d$
preserving one branch and removing any root-of-unity discrepancy
in the diagonal multipliers. The actual curve list used in
\ref{cdp:st:C.2.UN.38} and \ref{cdp:st:C.2.UN.53} gives its normalized parameter $\C^*$,
with $F^d$ acting by $u\mapsto\lambda u$, where
$0<|\lambda|<1$. For an axis this is its displayed diagonal
multiplier; for a resonant binomial curve $z^q=cw^p$ it follows
from $(z,w)=(a u^p,b u^q)$. The scalar and nonlinear-axis cases
are included. The actual map
$\C^*/\langle F^d\rangle\to C'$ induced by this inclusion is a
finite unramified elliptic covering. Pulling the constant deck
cocycle $1$ to the subgroup generated by $F^d$ gives the cocycle
$d$. Thus, if the original restriction were zero, its pullback
would give a holomorphic function on $\C^*$ satisfying
\[
                         b(\lambda u)-b(u)=d.          \tag{C.HR.4}\hypertarget{tag-00N3}{}\label{cdp:eq:C.HR.4}
\]
In its Laurent expansion on $\C^*$ the constant coefficient on
the left is zero, whereas $d>0$. This proves nonvanishing without
assuming that $F$ fixes every lifted branch individually.

\item For a secondary $H$, take the actual finite unramified primary
cover $\pi:H'\to H$ of \ref{cdp:st:C.2.UN.43}. A component $C'$ of the inverse
image of $C$ maps by a finite unramified elliptic covering to $C$.
The trace identities
$\operatorname{Tr}\circ\pi^*=(\deg\pi)\operatorname{id}$, on $H$
and on this curve covering, show that both pullbacks on $H^1(\OO)$
are injective. The surface pullback is an isomorphism of
one-dimensional spaces. Naturality of the restriction square,
together with the primary calculation, proves that the original
$H^1(H,\OO_H)\to H^1(C,\OO_C)$ is nonzero. Thus every coefficient
in \eqref{cdp:eq:C.HR.2} is nonzero, with the actual gluing maps included.

\par\medskip\noindent\textbf{Normalization, the graph, and its exact cycle product.}
The no-triple SNC equation gives the exact normalization sequence
\[
 0\longrightarrow\OO_S
 \longrightarrow\bigoplus_v\OO_{H_v}
 \xrightarrow{\ d\ }\bigoplus_e\OO_{C_e}
 \longrightarrow0,                                    \tag{C.HR.5}\hypertarget{tag-00N4}{}\label{cdp:eq:C.HR.5}
\]
where $d$ is the difference of the actual two restrictions. On
$H^0$ it is the ordinary signed incidence matrix $d_0$. The
connectedness of $S$ gives $\operatorname{rank}d_0=|V|-1$.
Since $H^2(H_v,\OO_{H_v})=0$, its long exact sequence gives
\[
 \begin{gathered}
 0\longrightarrow\operatorname{coker}d_0
 \longrightarrow H^1(S,\OO_S)
 \longrightarrow\ker d_1\longrightarrow0,\\
             H^2(S,\OO_S)\simeq\operatorname{coker}d_1.
 \end{gathered}                                        \tag{C.HR.6}\hypertarget{tag-00N5}{}\label{cdp:eq:C.HR.6}
\]
Here $\dim\operatorname{coker}d_0$ is the first Betti number of
the incidence graph.

\item Choose a spanning tree and a nonzero value at one vertex.
Its edge equations $r_{v,e}a_v=r_{w,e}a_w$ successively determine
nonzero values at all other vertices. Thus the tree rows of $d_1$
are independent and have a one-dimensional kernel. In particular
if the whole graph were a tree, $d_1$ would be surjective and
$H^2(S,\OO_S)=0$.

\item The map $f$ is flat: multiplication by its nonzero local parameter
is injective in each regular local ring of the smooth total space,
and torsion-free modules over the base discrete valuation ring
are flat. Upper semicontinuity for this proper flat family gives
$h^2(S,\OO_S)\ge h^2(F,\OO_F)=1$. The central fibre is the
reduced $S$ by hypothesis, so no thickened-fibre cohomology has
been substituted here. The tree case is impossible.

\ref{cdp:st:C.2.UN.44} gives a surjection from $\pi_1(S)$ to the incidence
graph group, and $\pi_1(S)$ is a quotient of $\Z^4$ because the
full fibre is reduced. The free graph group is therefore abelian
and has rank at most one. Together with the preceding paragraph
this proves that the graph is unicyclic and $|E|=|V|$.

\item There is now just one row left after the spanning-tree equations.
Writing its endpoints using the successive nonzero values shows
that it vanishes precisely when the product of the actual ratios
$r_{v,e}/r_{w,e}$ around the cycle is one. If it did not vanish,
$d_1$ would have rank $|V|$ and again $H^2(S,\OO_S)=0$.
The product is consequently one, the rank is $|V|-1$, and
\eqref{cdp:eq:C.HR.6} gives exactly \eqref{cdp:eq:C.HR.3}. A vertex-basis change cancels
between the incoming and outgoing ratios; an edge-basis change
cancels inside its ratio. This proves the asserted independence
without discarding the analytic restriction coefficients.

\par\medskip\noindent\textbf{The dualizing form and the canonical-order minimum.}
Let $D_v$ be the disjoint union of the double curves on $H_v$.
The residue exact sequence is
\[
 0\longrightarrow K_{H_v}
 \longrightarrow K_{H_v}(D_v)
 \xrightarrow{\operatorname{Res}}\bigoplus_{e\ni v}K_{C_e}
 \longrightarrow0.                                    \tag{C.HR.7}\hypertarget{tag-00N6}{}\label{cdp:eq:C.HR.7}
\]
Since $H^0(K_{H_v})=0$, a logarithmic two-form is determined by
its residues. Serre duality identifies the connecting map from
these residues to $H^1(K_{H_v})$ with the dual of the actual
restriction maps \eqref{cdp:eq:C.HR.1}. The signs may be fixed explicitly:
in a double chart $xy=0$ the two residues of a dualizing form
sum to zero. Orienting the edge takes its residue on the tail
and the negative on the head. Thus the global residue vector of
a section of $\omega_S$ belongs exactly to $\ker d_1^\vee$.
Conversely a vector in that kernel lifts uniquely by \eqref{cdp:eq:C.HR.7} on
each component, and the opposite-residue rule glues these forms
to a section of $\omega_S$. These are inverse maps.

\item At a leaf of an attached tree, the corresponding row of
$d_1^\vee$ contains only its one nonzero coefficient. The
residue on that edge is zero. Delete that leaf and repeat.
Every tree-edge residue is zero. On the remaining cycle each
vertex relates its two adjacent residues by nonzero coefficients;
its one-dimensional kernel has a nonzero residue on every cycle
edge. On an off-cycle component all residues vanish, so the
form is an ordinary holomorphic two-form and is zero. On a cycle
component its two nonzero residues make it nonzero. Vanishing
residue along any other boundary curve removes the logarithmic
pole there, directly from the local expression
$a(y,z)\,dy\wedge dz/y$.

\item Finally take the canonical evaluation section of \ref{cdp:st:C.2.UN.38} and
write its vertical divisor as $\sum k_vH_v$. Put $k_0=\min k_v$.
Because $S$ is the full reduced divisor of $t$, division by
$t^{k_0}$ is an actual holomorphic section of the relative
canonical line. Its restriction to $S$ is a nonzero section of
$\omega_S$. At an edge joining $v,w$, the explicit adjunction
calculation, with $t=xy$, is
\[
 \frac{x^{k_v-k_0}y^{k_w-k_0}\,dx\wedge dy\wedge dz}{xy}
 \quad\longmapsto\quad
 y^{k_w-k_0-1}\,dy\wedge dz\quad\text{on }x=0
                                                        \tag{C.HR.8}\hypertarget{tag-00N7}{}\label{cdp:eq:C.HR.8}
\]
when $k_v=k_0$, times a holomorphic unit; when $k_v>k_0$ the
restriction on that component is zero. The local logarithmic
expression is the usual restriction of the relative dualizing
line, or equivalently the residue of the displayed meromorphic
three-form. It follows that its nonzero components are exactly
those with $k_v=k_0$. The preceding residue calculation identifies
these with the cycle components, as claimed.

\par\medskip\noindent\textbf{The full algebraic-dimension-zero branch.}
The actual primary and secondary curve computations in
\ref{cdp:st:C.2.UN.38} and \ref{cdp:st:C.2.UN.43} give two possibilities in algebraic
dimension zero: a nonlinear Hopf surface has just one elliptic
curve, and a diagonal Hopf surface has just its two distinct
axis elliptic curves. The diagonal finite action cannot exchange
them: such an exchange would force equal positive powers of the
two contracting multipliers, giving a nonconstant meromorphic
invariant and algebraic dimension one. Thus every vertex of the
present graph has degree at most two.

\item A connected unicyclic finite graph with all degrees at most two
is a cycle and has every degree equal to two. Indeed its equality
$|E|=|V|$ gives average degree two, so no smaller degree is possible.
Consequently there is no nonlinear component and no attached
tree. There is no loop at a single vertex, since the components
of the specified SNC surface are smooth embedded components:
such a loop would be a self-identification, not a double curve
between distinct components. Hence there are at least two
components.

\item For this now established diagonal cycle, the actual normal-product
trivialization of $t=xy$ gives the dual-normal maps
(C.HC.3). Removing the zero sections gives the actual open groups
$G_v=\C^2/\Lambda_v$ and the affine logarithmic edge maps
(C.HC.5). Their integral lattice product is precisely $M$, and
\ref{cdp:st:C.2.UN.46} proves that an actual torus smoothing forces $M=I$.
Conversely the specified axis maps and any compatible dual-normal
maps with $M=I$ give the proper plumbing of \ref{cdp:st:C.2.UN.46}; the actual
varying-period comparison in \ref{cdp:st:C.2.UN.48} extends it to a compact
smooth threefold over a compact curve. That theorem retains every
component and every specified gluing constant, so the central
support is isomorphic to the specified $S$, rather than merely
having the same graph. Its other supports have Euler zero.

\item The exact matrix calculation in
\path{code/E_hopf_graph_reduction.py} checks the spanning-tree
elimination, the cycle product, both coherent dimensions, and the
support of the dual residue vector on several branched graphs.
The argument above proves each comparison for every finite graph.

\end{enumerate}
\end{proof}






\begin{proposition}[The complete all-ruled cycle criterion with its full elliptic translation]\hypertarget{tag-01E0}{}\label{cdp:st:C.2.UN.55}
\StatementRecord{cdp:st:C.2.UN.55}{proved}{01E0}{\SeriesRef{cdp:st:C.2.UN.51}, \SeriesRef{cdp:eq:C.TL.3}, \SeriesRef{cdp:eq:C.TL.1}, \SeriesRef{cdp:eq:C.TL.2}, \SeriesRef{cdp:eq:C.TL.5}, \SeriesRef{cdp:eq:C.TL.4}, \SeriesRef{cdp:eq:C.TL.7}, \SeriesRef{cdp:eq:C.TL.8}}

Specify an oriented cycle of $N\geq1$ marked ruled elliptic surfaces
$R_i=\PP(\OO_{E_i}\oplus L_i)$, with
$L_i\in\operatorname{Pic}^0(E_i)$, elliptic identifications of
successive marked sections, and duality of their actual normal
lines. For $N=1$ identify the two distinct marked sections of
the same normalization. The following are equivalent:
the integral complement-lattice product $M$ of \ref{cdp:st:C.2.UN.51} is the
identity; the specified support has an actual smooth-total
two-torus smoothing; and it occurs as a full reduced fibre of
a smooth compact complex threefold over $\PP^1$ with two-torus
general fibre.

The compact realization retains every marked surface and every
elliptic gluing map, including an arbitrary total elliptic
translation around the cycle. It chooses the normal-product
constants. Its only singular fibres are the specified cycle and
another reduced cycle of $N$ ruled elliptic components. Every
reduced fibre support has Euler number zero.
\end{proposition}



\begin{proof}

\begin{enumerate}[label={$\langle1\rangle$\arabic*.},leftmargin=*]
\item The necessary condition and the actual local equivalence are
\ref{cdp:st:C.2.UN.51}. We prove the compact assertion without removing the
horizontal affine translation.

\par\medskip\noindent\textbf{An actual common lattice and its retained translation.}
Use the actual logarithmic linear identifications of \ref{cdp:st:C.2.UN.51} to put
all complements in $G=\C^2/\Lambda$, where $\Lambda$ has real
rank three. All positive outgoing normal meridians are the same
primitive $n\in\Lambda$. Put $p=-n$. The normal-line quotient of
each ruled surface is an elliptic curve. Consequently
$\C p\cap W=\R p$, where $W=\operatorname{span}_{\R}\Lambda$.
The complex line $W\cap iW$ and $p$ are a complex basis of
$\C^2$. The complex linear map sending this line to the first
coordinate line and $p$ to $(0,1)$ puts $W$ exactly in
$\C\times\R$.

\item Extend the primitive $p$ to an integral basis of $\Lambda$.
After ordering the other two vectors, its exact form is
\[
 \begin{gathered}
 w_1=(a_1,c_1),\qquad w_2=(a_2,c_2),\qquad p=(0,1),\\
 c_1,c_2\in\R,\qquad
 C=\operatorname{Im}(\overline{a_1}a_2)>0.
 \end{gathered}
 \tag{C.TL.1}\hypertarget{tag-013Q}{}\label{cdp:eq:C.TL.1}
\]
This changes coordinates, not the lattice or its integral
subgroups. In common coordinates every actual edge map is a
translation with its normal logarithmic term. Choose logarithmic
lifts of its constants $b_i$ and put $b=\sum_i b_i$.
Its first coordinate, modulo
$\Z a_1+\Z a_2$, is exactly the total elliptic gluing translation.
A normal-product constant changes $b_i$ only in the line
$\C p$. Adjust one such constant so that
\[
                         b=(\sigma,i).
 \tag{C.TL.2}\hypertarget{tag-013R}{}\label{cdp:eq:C.TL.2}
\]
The prescribed elliptic maps, hence the horizontal class of
$\sigma$, are unchanged. No restriction on that class has been
imposed.

\par\medskip\noindent\textbf{A holomorphic period family on the entire punctured sphere.}
For $z\in\C$ and a fixed vector $w=(w^{(1)},w^{(2)})\in\C^2$
define
\[
 V_z(w)=
 \bigl(w^{(1)}+z\overline{w^{(2)}},
       w^{(2)}-z\overline{w^{(1)}}\bigr).
 \tag{C.TL.3}\hypertarget{tag-013S}{}\label{cdp:eq:C.TL.3}
\]
For each fixed $w$ this is a holomorphic vector-valued function
of $z$. For each fixed $z$ it is a real linear isomorphism,
with real determinant $(1+|z|^2)^2$. On a branch of $\log z$,
put
\[
 \begin{gathered}
 P_1(z)=V_z(w_1),\quad P_2(z)=V_z(w_2),\quad
 P_3(z)=V_z(p),\\
 P_4(z)=V_z(b)+\frac{N}{2\pi i}\log z\,V_z(p).
 \end{gathered}
 \tag{C.TL.4}\hypertarget{tag-013T}{}\label{cdp:eq:C.TL.4}
\]
Changing the logarithm by $2\pi i$ replaces $P_4$ by
$P_4+NP_3$. Thus the subgroup generated by these four vectors
is single-valued. We check its real rank everywhere, including
where the rank-three group alone has no elliptic normal quotient.

\item Write $r=|z|$ and $\tau=N\log z/(2\pi i)$. The actual inverse
of \eqref{cdp:eq:C.TL.3} gives
\[
 V_z^{-1}\bigl(iV_z(p)\bigr)
       =\frac{i(2z,1-r^2)}{1+r^2}.
 \tag{C.TL.5}\hypertarget{tag-013U}{}\label{cdp:eq:C.TL.5}
\]
The first three vectors in \eqref{cdp:eq:C.TL.1} span $\C\times\R$, so
their real determinant with any fourth vector $v$ is
$C\operatorname{Im}v^{(2)}$. Equation \eqref{cdp:eq:C.TL.2} therefore gives
$\det_{\R}(w_1,w_2,p,b)=C$. The real multiple
$\operatorname{Re}\tau\,V_z(p)$ has no effect on a determinant.
Using \eqref{cdp:eq:C.TL.5} for the remaining multiple yields the exact value
\[
 \begin{aligned}
 D(z)&=\det_{\R}(P_1(z),P_2(z),P_3(z),P_4(z))\\
 &=C(1+r^2)
       \left(1+r^2+\frac{N}{2\pi}(r^2-1)\log r\right)>0
                         \qquad (0<r<\infty).
 \end{aligned}
 \tag{C.TL.6}\hypertarget{tag-013V}{}\label{cdp:eq:C.TL.6}
\]
Indeed $(r^2-1)\log r\geq0$, while $1+r^2>0$.
In particular the retained horizontal $\sigma$ and the real
characters $c_1,c_2$ require no inequality.

\item The four holomorphic period sections on a simply connected
base chart act on $\C^2$ by translations. Their real determinant
is nonzero; on each compact base subset their real matrices and
inverses are uniformly bounded. Thus their integral translation
action is properly discontinuous, and a continuously varying
closed real parallelepiped gives properness of the quotient.
The integral change $P_4\mapsto P_4+NP_3$ identifies these
quotients on overlaps. This constructs a smooth proper family
$\mathcal T\to\C^*$ of connected two-tori, with these actual
periods.

\par\medskip\noindent\textbf{The prescribed filling at zero.}
For $z$ sufficiently close to zero, the line
$\C V_z(p)$ is transverse to the real rank-three lattice in
the same sense as at zero. Indeed its transverse determinant
is $C(1+r^2)(1-r^2)$, which is nonzero for $r<1$.
Use the separate ruled moving frame of (C.MX.10):
the normal vector $V_z(n)$ and a transverse integral lattice
lift define an actual elliptic family and a holomorphic
degree-zero line on it. Its projective completion gives the
marked ruled family through each prescribed $R_i$.
The original integral identifications and all finite changes
of logarithm are retained.

\item Choose the moving edge translation to be $V_z(b_i)$ in the
common coordinates. It induces the actual translation on the
elliptic quotient and the reciprocal normal-line map; these
extend over the marked sections by their line-bundle formulas.
The normal logarithmic coefficient of each edge is
$-V_z(n)/(2\pi i)=V_z(p)/(2\pi i)$. Consequently the entire
cycle has precisely the fourth period in \eqref{cdp:eq:C.TL.4}, not only
its leading asymptotic term. The central value of each edge
is its prescribed elliptic gluing map.

\item Apply the proper normal-product plumbing of \ref{cdp:st:C.2.UN.51} to these
moving line charts, with actual equation $xy=z$ in each neck.
The moving component cores are smooth over the disc; the
normal-product equation is smooth at $x=y=z=0$.
Uniform flat norms and finitely many bounded cores and necks
give the same Hausdorff proper local filling as in that proof.
Its full central divisor is reduced, and its support is the
specified marked cycle.

\item On a punctured base chart, send each core's logarithmic
coordinate to the common vector space, adding the successive
edge translations. Across each neck these maps agree modulo
the periods in \eqref{cdp:eq:C.TL.4}. Across a change of branch of $\log z$
their changes are integral multiples of $P_3$. Going around
the cycle adds exactly $P_4$. They therefore give a holomorphic
map to $\mathcal T$. Its fibre maps are the closed-coframe
isomorphisms of \ref{cdp:st:C.2.UN.51}: all four fundamental-group generators
are exactly the displayed periods. Their local inverses
depend holomorphically on the base, so this is an actual
biholomorphism of punctured families.

\par\medskip\noindent\textbf{The second proper filling at infinity.}
Put $s=z^{-1}$ and change the two-dimensional vector-bundle
frame by $\eta_\infty=\eta/z$. In this frame \eqref{cdp:eq:C.TL.3} becomes
\[
 V^\infty_s(w)
       =\bigl(\overline{w^{(2)}},-\overline{w^{(1)}}\bigr)
                          +s w.
 \tag{C.TL.7}\hypertarget{tag-013W}{}\label{cdp:eq:C.TL.7}
\]
This is an actual holomorphic frame change on $\C^*$, and
$V^\infty_0$ is an invertible antilinear map. The periods
in the new frame are exactly
\[
 \begin{gathered}
 V^\infty_s(w_1),\quad V^\infty_s(w_2),\quad V^\infty_s(p),\\
 V^\infty_s(b)
          -\frac{N}{2\pi i}\log s\,V^\infty_s(p).
 \end{gathered}
 \tag{C.TL.8}\hypertarget{tag-013X}{}\label{cdp:eq:C.TL.8}
\]
The ambiguity of the logarithm again adds an integral period.

\item Near $s=0$ the rank-three lattice and the complex line through
$V^\infty_s(p)$ give a genuine elliptic quotient; the
transversality at zero follows by applying the invertible
antilinear map in \eqref{cdp:eq:C.TL.7} to the original quotient.
Construct $N$ marked ruled components from this actual
nonzero line bundle. Give each positive outgoing meridian
the vector $n_\infty(s)=V^\infty_s(p)$, and each incoming
meridian its negative. Choose the sum of their full affine
edge constants to be $V^\infty_s(b)$, for example by putting
this constant on one edge and zero on the others.
Every such constant induces an actual elliptic translation
and its compatible normal-line map.

\item The normal-period sum of this opposite-oriented cycle is
$-N V^\infty_s(p)$. Its normal-product plumbing with equation
$xy=s$ therefore has precisely the periods \eqref{cdp:eq:C.TL.8}.
The same finite-core construction is a smooth proper local
filling with reduced central divisor. Its exact logarithmic
maps identify the punctured family with the frame change of
$\mathcal T$ just displayed. This explicitly constructs the
second filling; no contraction of a global auxiliary line
bundle is involved.

\par\medskip\noindent\textbf{Compact gluing and the Euler comparison.}
Choose small disjoint discs about $0$ and $\infty$ and a middle
annulus covering the rest of $\PP^1$. Glue their two proper
fillings to $\mathcal T$ by the actual punctured-family
biholomorphisms above. Over every point of the base all pieces
are identified with one Hausdorff proper local model. Hence
the glued space is Hausdorff and the map is proper locally
on the base. A finite closed refinement of the three base
charts makes its inverse image a finite union of compact
sets, proving properness over the compact $\PP^1$.
The total space is smooth by its displayed local models.
All other fibres are the tori of \eqref{cdp:eq:C.TL.4}.

\item Each ruled elliptic component has Euler number
$e(\PP^1)e(E)=0$, and each double curve has Euler number zero.
There are no triple points, so both singular supports have
Euler number zero. This proves every asserted compact and
marked-data comparison, including the one-component
self-identification.

\item The symbolic determinant, frame-change, inverse-map and
monodromy identities are certified in
\path{code/A_twistor_log_cycles.py} and its saved output.

\end{enumerate}
\end{proof}






\begin{proposition}[The complete weighted mixed-cycle criterion]\hypertarget{tag-01E1}{}\label{cdp:st:C.2.UN.56}
\StatementRecord{cdp:st:C.2.UN.56}{proved}{01E1}{\SeriesRef{cdp:st:C.2.UN.52}, \SeriesRef{cdp:eq:C.WM.4}, \SeriesRef{cdp:eq:C.WM.5}, \SeriesRef{cdp:eq:C.WM.6}, \SeriesRef{cdp:st:C.2.UN.51}}

The criterion and compact realization of \ref{cdp:st:C.2.UN.52}
remain valid when some vertices are marked ruled surfaces
\[
        R_i=\PP(\OO_{E_i}\oplus L_i),\qquad
        E_i\ \hbox{elliptic},\qquad L_i\in\operatorname{Pic}^0(E_i),
\tag{C.WM.1}\hypertarget{tag-017F}{}\label{cdp:eq:C.WM.1}
\]
with their zero and infinity sections, provided at least
one vertex is a diagonal primary or secondary Hopf surface
of algebraic dimension zero. The support is a simple-NC
cycle of embedded components, with edges between distinct
components. All marked surfaces, elliptic gluing maps and
positive component multiplicities are retained.

More precisely choose the actual compatible weighted normal
lines, sections and boundary maps (C.WH.1)--(C.WH.2) on
these marked surfaces. Form their normalized common root
covers and the compressed necks $uv=s^{h_i}$ as in \ref{cdp:st:C.2.UN.52}.
The resulting root graph must be connected and its actual
rank-three integral lattice product must be the identity.
These conditions are necessary and sufficient for connected
two-torus fibres in the explicit smooth weighted plumbing.
Every weighted support passing them occurs over $\PP^1$
with smooth compact total space and Euler number zero for
every reduced fibre support. Every actual torus degeneration
in this specified branch gives these data and passes the test.

There are explicit additional necessary relations at a ruled
vertex of multiplicity $m$, with neighbouring multiplicities
$a,b$:
\[
 \begin{gathered}
      m\mid a+b,\qquad
      \deg(\mathcal L|_{\PP^1})=-\frac{a+b}{m},\\
      \gcd(m,a)=\gcd(m,b),\qquad
      q=\frac{m}{\gcd(m,a)}
         =\frac{m}{\gcd(m,b)}.
 \end{gathered}
\tag{C.WM.2}\hypertarget{tag-017G}{}\label{cdp:eq:C.WM.2}
\]
Every connected strict root cover of this ruled component
is exactly a marked surface
$\PP(\OO_{E'}\oplus L')$, with an \'{e}tale elliptic map
$\pi:E'\to E$ and an actual line isomorphism
\[
                   (L')^q\simeq\pi^*L.
\tag{C.WM.3}\hypertarget{tag-017H}{}\label{cdp:eq:C.WM.3}
\]
Its finite map has fibre degree $q$, ramification order
$q$ at each marked section, and base degree
$\gcd(m,a)/r$, where $r$ is the common number of strict
covers in \ref{cdp:st:C.2.UN.52}. In particular its total degree is $m/r$.
All the component, axis and inertia formulas (C.WH.0)
hold unchanged.

A compatible cycle in this branch has at least two
original Hopf vertices. Every ruled line $L_i$ in it
has infinite order in $\operatorname{Pic}^0(E_i)$.
\end{proposition}



\begin{proof}

\begin{enumerate}[label={$\langle1\rangle$\arabic*.},leftmargin=*]
\item \par\medskip\noindent\textbf{The weighted normal line on an actual ruled component.}
Let $C^-,C^+$ be its two marked sections and write
\[
  \mathcal L^{m}
       \simeq \OO_R(-aC^- -bC^+),\qquad
       \operatorname{div}(\sigma)=aC^-+bC^+.
\tag{C.WM.4}\hypertarget{tag-017I}{}\label{cdp:eq:C.WM.4}
\]
Restriction to any ruling fibre gives
$m\deg(\mathcal L|_{\PP^1})=-(a+b)$,
proving the first two assertions of \eqref{cdp:eq:C.WM.2}.
The congruence $b\equiv-a\pmod m$ gives
the equality of the two gcd's and the displayed
value of $q$.

\item The local line-bundle comparison needed for the weighted
plumbing is holomorphic, including when $\mathcal L$
has nonzero degree on ruling fibres.
Each section neighbourhood obtained by removing the
opposite section is the entire total space of its
normal degree-zero line on $E$.
On this chart, the right side of \eqref{cdp:eq:C.WM.4} is the
pullback of a degree-zero line on $E$:
the removed divisor is trivial there, and the divisor
of the zero section is the pullback of its normal line.
Thus $\mathcal L^m$ is flat on this chart.
Dividing its transferred flat connection by $m$
gives an actual flat connection on $\mathcal L$ there.
The chart retracts to its zero section, so its
monodromy character agrees with the pullback of
the restricted character. This gives the actual
holomorphic isomorphism with the pulled-back
axis restriction. The weighted component cores
and normal-product necks consequently glue exactly
as in \ref{cdp:st:C.2.UN.52}. No assertion that $\mathcal L$ is
globally flat on the ruled surface is used.

\item For an alternative global description of this
normal line, take $\OO_R(C^+)$ as a relative
degree-one line and put
$d=\deg(\mathcal L|_{\PP^1})$.
The line $\mathcal L\otimes\OO_R(-dC^+)$
has trivial restriction to every ruling fibre.
Its direct image is a line on $E$, and evaluation
is an isomorphism on every fibre and hence everywhere,
by proper base change and the local Nakayama lemma.
Therefore
$\mathcal L=\OO_R(dC^+)\otimes\pi_R^*A$.
Equation \eqref{cdp:eq:C.WM.4} and the fact that
$C^--C^+$ is pulled back from $\operatorname{Pic}^0(E)$
give $\deg A=0$. This describes the same local
flat pullbacks by an explicit bundle map.

\par\medskip\noindent\textbf{The complete ruled root-cover comparison.}
Remove the two sections. Its actual open group is
$G=L^\times=\C^2/\Lambda$, with the rank-three
period lattice (C.MX.3). Its positive normal
meridian $n$ is primitive in $\Lambda$.
Let $\Lambda'\subset\Lambda$ be the finite-index
sublattice of one connected cyclic root component
over this open set, and let $q$ be the smallest
positive integer for which $qn\in\Lambda'$.
Projection along the complex line $\C n$
gives the actual elliptic base covering
\[
 E'=\C^2/(\C n+\Lambda')
      \longrightarrow
 E=\C^2/(\C n+\Lambda).
\tag{C.WM.5}\hypertarget{tag-017J}{}\label{cdp:eq:C.WM.5}
\]
This is finite and \'{e}tale: the two projected
rank-two lattices are discrete, and one has
finite index in the other. The source open
group maps to $E'$ with kernel
$\C n/\Z qn\simeq\C^*$.
Taking logarithms of its fibre coordinate
gives constant factors of automorphy, hence
an actual degree-zero line $L'$ on $E'$.

\item The homomorphism of open groups induced by
$\Lambda'\subset\Lambda$ acts on this
$\C^*$ kernel by the power $z\mapsto z^q$.
Equivalently it gives exactly \eqref{cdp:eq:C.WM.3}.
For example choose lifts of the two base
periods to the source logarithmic lattice.
If their normal coordinates are $b_j'$,
the image coordinates are $qb_j'$ modulo
integral angular periods; exponentiation
is the line isomorphism in \eqref{cdp:eq:C.WM.3}.
Thus the open finite map extends to the
actual projective-bundle map
\[
 \PP(\OO_{E'}\oplus L')
          \longrightarrow \PP(\OO_E\oplus L),
\tag{C.WM.6}\hypertarget{tag-017K}{}\label{cdp:eq:C.WM.6}
\]
using the given base map and the homogeneous
power map on each ruling fibre.
This map is finite, has ramification order
$q$ at both marked sections, and agrees
with the original cover on the dense open
group. Its source is smooth and normal.
Uniqueness of finite normalization therefore
identifies it with the actual strict root
component, including both sections.

\item This also proves directly that the two
ramification orders are equal. Applying
the local normalization
$\xi^m=z^a$ and $\xi^m=z^b$ to the two
sections identifies their common value
with \eqref{cdp:eq:C.WM.2}. The lattice index splits as
\[
 [\Lambda:\Lambda']
     =q\,[\Lambda:\Lambda'+\Z n].
\tag{C.WM.7}\hypertarget{tag-017L}{}\label{cdp:eq:C.WM.7}
\]
The second factor is the degree of \eqref{cdp:eq:C.WM.5}.
Since \ref{cdp:st:C.2.UN.52}'s component count gives total
degree $m/r$, it equals $\gcd(m,a)/r$.
The two section covers have that same
degree, proving all the asserted maps
and degrees without replacing the root
surface by an unspecified ruled model.

\par\medskip\noindent\textbf{The connected root cycle and its finite action.}
Each strict root component is now either
one of the actual diagonal Hopf covers of
\ref{cdp:st:C.2.UN.52} or the marked ruled surface \eqref{cdp:eq:C.WM.6}.
Every one has two connected marked elliptic
curves. The local root charts are still
exactly (C.WH.4), and their resolutions
insert the same two-section ruled bridges.
Thus the root graph is a union of cycles.
Connectedness, reduced-root Stein factorization
and the actual core-and-neck group calculation
give the same criterion
\[
                         M=I
\tag{C.WM.8}\hypertarget{tag-017M}{}\label{cdp:eq:C.WM.8}
\]
for a connected two-torus general fibre.
This is also the mixed-cycle criterion
of \ref{cdp:st:C.2.UN.51}, applied to the resolved root,
with the bridges compressed by their
exact $uv=s^h$ maps.

\item The root deck group preserves the actual
map to each original normalization and
its two labelled marked branches. It
therefore preserves the lifted cyclic
orientation and acts by rotations.
Its common vertex stabilizer, connected
cover counts, generic inertia and axis
stabilizers are exactly those computed
in (C.WH.6) and (C.WH.0). The local
degrees at a ruled vertex just proved
are consistent with, and refine, those
same actual sheet counts.

\item When $M=I$, identify all open complements
with one $G=\C^2/\Lambda$.
On a strict ruled component the two
normal meridians are opposite, so
the common outgoing periods satisfy
$n_j=n_{j-1}$ there. At every Hopf
component they are complex-linearly
independent and satisfy the same
contracting-height identity (C.HG.4).
The positive compressed edge lengths
therefore give
\[
                  u=-\sum_j h_j n_j,\qquad \ell(iu)\ne0.
\tag{C.WM.9}\hypertarget{tag-017N}{}\label{cdp:eq:C.WM.9}
\]
The common group has a Hopf open chart,
so it contains no compact elliptic curve.
Its quotient along the primitive part
of $u$ is the same nontorsion degree-zero
elliptic line used in \ref{cdp:st:C.2.UN.52}.

\item The finite root action is affine also
on the ruled open charts. A marked
projective-bundle isomorphism has an
affine elliptic base map and a
nonzero fibre multiplier. The logarithmic
derivative of that multiplier is
holomorphic and doubly periodic,
hence constant, exactly as in \ref{cdp:st:C.2.UN.51}.
The common linear part $R$ fixes $u$
because the actual action is a
rotation. The nontorsion line argument
of (C.WH.8) consequently proves
$R=I$ without any change.
Summing the exact affine equations
(C.WH.9) again gives
$u/N\in\Lambda$.

\item Every ruled line $L_i$ has infinite
order. Otherwise its finite pullback
and $q$th root \eqref{cdp:eq:C.WM.3} would make
$L'$ torsion as well. Its nonzero
bundle would contain the image of
an elliptic curve under a finite
trivializing isogeny. This would be
a compact elliptic curve in the
common Hopf open group, a contradiction.

\item There must be at least two original
Hopf vertices. If there were only one,
start at a lifted Hopf vertex and
follow the cycle to the next lifted
vertex above that same original Hopf.
Every intermediate strict vertex is
ruled, so its incoming and outgoing
normal periods agree; compressing
or resolving a neck has the same
property. Hence the common period
immediately preceding the next Hopf
is the outgoing normal period of
the first Hopf, in the notation of
(C.MX.5). A root deck rotation takes
the first marked Hopf to the second.
Its linear part is the identity,
so their respective outgoing normal
periods agree as well. At that
second Hopf the two consecutive
normal periods would therefore be
equal, contradicting their complex
linear independence. This proves
the final component-count assertion.

\par\medskip\noindent\textbf{Compact realization preserves every marked surface.}
The finite transport of \ref{cdp:st:C.2.UN.52} applies
to this mixed root with one precise
additional frame choice. At a Hopf
vertex transport its two independent
normal periods. At a ruled vertex,
transport instead its primitive
normal period, a transverse base
lattice lift and a third lattice
basis vector, as in (C.MX.10).
This gives an actual family
$\PP(\OO_{E'(t)}\oplus L'(t))$
of the given marked ruled surface,
rather than attempting to use its
dependent pair of normal meridians.

\item These frames can be chosen identically
on each root-deck orbit: the common
linear part is the identity and
the corresponding normal periods
and lattice subgroups agree.
Translations of the common open
group extend to the marked ruled
compactification. Indeed its
projection to $E'$ is a base
translation, and the group law on
the rigidified degree-zero line
gives the fibre isomorphism;
this extends to its zero and
infinity sections. Hence the
transported affine deck maps
extend to these ruled families
as actual holomorphic maps.

\item All equations for the edge
translations and deck translations
are still the constant-integer
system (C.WH.9).
The two adjacent independent
directions at a Hopf vertex belong
to distinct rotation orbits, so
the two orbitwise normal constants
still set the total translation
to zero without changing any
central elliptic gluing map.
The rational-left-kernel argument
and fixed right inverse of \ref{cdp:st:C.2.UN.52}
give the full holomorphic finite
action on the moving compressed
root plumbing.

\item On every neck retain $uv=s^{h_i}$.
Its actual pointwise stabilizer
has the same invariant ring
$\C\{u^B,v^A\}$ as (C.WH.12),
and its residual axis action is
\'{e}tale. Component inertia is
again a normal reflection.
Thus the quotient total space is
smooth and has exactly the
original full multiplicities.
At zero the components and all
their marked maps and finite
quotients are the original ones,
by \eqref{cdp:eq:C.WM.6} and the Hopf comparison.
The punctured quotient is the
marked torus torsor with period
index $\operatorname{div}(u)/N$;
the exact sequence (C.WH.13)
gives its actual holomorphic
isomorphism to the compact
elliptic suspension.

\item Performing the same proper
annular replacement therefore
gives the stated compact family
over $\PP^1$. Ruled elliptic
components, Hopf components and
their double elliptic curves all
have Euler number zero. Every
other reduced support is unchanged
from the compact suspension and
also has Euler number zero.
This proves both sufficiency and
the asserted exact retention.
Necessity for every actual
smoothing uses its actual weighted
normal data and the tubular
comparison of \ref{cdp:st:C.2.UN.52}, now with
the explicit ruled maps above.

\end{enumerate}
\end{proof}






\begin{proposition}[The exact marked-cycle criterion without an algebraic-dimension restriction]\hypertarget{tag-01E2}{}\label{cdp:st:C.2.UN.57}
\StatementRecord{cdp:st:C.2.UN.57}{proved}{01E2}{\SeriesRef{cdp:st:C.2.UN.46}, \SeriesRef{cdp:st:C.2.UN.51}, \SeriesRef{cdp:eq:C.AD.1}, \SeriesRef{cdp:eq:C.AD.2}, \SeriesRef{cdp:eq:C.AD.3}, \SeriesRef{cdp:st:C.2.UN.55}, \SeriesRef{cdp:eq:C.AD.6}, \SeriesRef{cdp:eq:C.AD.10}, \SeriesRef{cdp:eq:C.AD.8}, \SeriesRef{cdp:eq:C.AD.9}}

Specify a connected oriented cycle of marked surfaces, each either a
diagonal primary or secondary Hopf surface with its two distinct axes,
or $\PP(\OO_E\oplus L)$ with $L\in\operatorname{Pic}^0(E)$ and its
zero and infinity sections. Specify the elliptic gluing maps and
duality of the actual axis or section normal lines. Components are
smooth embedded components, except that a single ruled normalization
with its two sections identified is also allowed.

The actual rank-three complement-lattice product $M$ is the identity
precisely for those specified supports which admit a smooth-total
two-torus smoothing. Every such support occurs as a full reduced fibre
of a smooth compact complex threefold over $\PP^1$. The construction
retains each component and every specified elliptic gluing map. Its
only other singular fibre is the complex conjugate marked support,
and every reduced fibre support has Euler number zero. No restriction
on the algebraic dimension of the diagonal Hopf components is imposed.
The construction chooses normal-product constants.
\end{proposition}



\begin{proof}

\begin{enumerate}[label={$\langle1\rangle$\arabic*.},leftmargin=*]
\item \par\medskip\noindent\textbf{The actual complement maps and the local criterion.}
For a diagonal Hopf component the exponential preimage of its full
diagonal deck group is an actual rank-three lattice $\Lambda_i$ in
$\C^2$. Its complement of both axes is $\C^2/\Lambda_i$. This
description retains the finite secondary group, including its two
angular characters. Both small normal meridians are primitive: a
smaller angular lattice vector would give a nonidentity finite deck
element fixing a nonzero point of the other axis, contrary to the
freeness of this Hopf quotient. For a ruled component the complement
is exactly $L^\times$; its degree-zero line-bundle factors give the
rank-three lattice and primitive normal meridian in (C.MX.3).

\item An elliptic gluing map lifts to an affine map of universal covers.
A dual-normal line isomorphism lifts to a nowhere-zero function;
the derivative of its logarithm is periodic under the two elliptic
periods, hence constant. Its logarithm is therefore affine.
Consequently the actual normal-product neck has exactly the affine
logarithmic maps of (C.HC.5) and (C.MX.2), including their constants.
These statements use the displayed axes and sections; they use no
absence of further curves on a component.

\item The component cores and the $xy=t$ necks give the actual fundamental
group
\[
                         \pi_1(F_t)=\Lambda\rtimes_M\Z.
                                                        \tag{C.AD.1}\hypertarget{tag-009B}{}\label{cdp:eq:C.AD.1}
\]
For an arbitrary specified smoothing, this is the same comparison
as in \ref{cdp:st:C.2.UN.46} and \ref{cdp:st:C.2.UN.51}: its canonical divisor sections and a smooth
tubular retraction identify the component cores with their actual
normal circle bundles, and its two divisor sections identify each
double collar with $xy=t$. The changes of these identifications
give the specified elliptic and normal-line maps on fundamental
groups. Thus \eqref{cdp:eq:C.AD.1} also holds for the actual smoothing, without
asserting that its holomorphic neighbourhood is linear. The
abelian group of a two-torus forces $M=I$.

\item Conversely the actual flat normal lines and affine necks give the
proper normal-product plumbing of those proofs. Their argument
requires only the entire line-bundle axis charts, disjoint marked
curves and the displayed constant factors, all present here.
When $M=I$, the common logarithmic coframe glues on every nearby
compact fibre. Its two dual holomorphic fields commute and are
complete. Their $\C^2$ action has discrete cocompact stabilizer
and identifies that fibre with a complex two-torus. This proves
the full local equivalence in the present generality.

\par\medskip\noindent\textbf{The positive total normal period.}
Assume $M=I$ and identify the complement groups with
$G=\C^2/\Lambda$ by the actual linear edge maps. Write $n_i$
for the outgoing normal meridian at vertex $i$; the incoming
one is $-n_{i-1}$. At a Hopf vertex these two vectors are
complex-linearly independent, and its actual logarithmic frame is
\[
       P_i=\left(\frac{n_i}{2\pi i},
                         -\frac{n_{i-1}}{2\pi i}\right).
                                                        \tag{C.AD.2}\hypertarget{tag-009E}{}\label{cdp:eq:C.AD.2}
\]
At a ruled vertex $n_i=n_{i-1}$, and its separate frame is the
one in (C.MX.10).

\item Put $W=\operatorname{span}_{\R}\Lambda$ and choose a nonzero
real linear functional $\ell$ with kernel $W$. A contracting
Hopf period with logarithms $\log\alpha_i,\log\beta_i$ lies
in $W$. Applying $\ell$ to \eqref{cdp:eq:C.AD.2} times that period gives
\[
 \log|\alpha_i|\,\ell(i n_i)
       =\log|\beta_i|\,\ell(i n_{i-1}).                 \tag{C.AD.3}\hypertarget{tag-009F}{}\label{cdp:eq:C.AD.3}
\]
Both logarithms are negative. Neither transverse value can
vanish: their simultaneous vanishing would put both complex
normal lines, hence all $\C^2$, inside the real three-plane
$W$. Thus they have the same nonzero sign. At a ruled vertex
the two values coincide. In an all-ruled cycle their common
value is nonzero because the actual elliptic quotient satisfies
$\C n\cap W=\R n$. Hence all values around every present
cycle have the same nonzero sign.

\item The exact fourth period is
\[
 b+\frac{\log t}{2\pi i}u,\qquad
             u=-\sum_i n_i\in\Lambda,
                                                        \tag{C.AD.4}\hypertarget{tag-009G}{}\label{cdp:eq:C.AD.4}
\]
where $b$ is the sum of the actual affine edge constants.
Equation \eqref{cdp:eq:C.AD.3} gives $\ell(iu)\ne0$. Write $u=d p$ with
$p$ primitive and $d>0$. The actual complex linear coordinate
change of (C.TL.1) therefore gives
\[
 \begin{gathered}
 W=\C\times\R,\qquad p=(0,1),\\
 \Lambda=\langle (a_1,c_1),(a_2,c_2),p\rangle_{\Z},
 \quad c_1,c_2\in\R,\quad
 C=\operatorname{Im}(\overline{a_1}a_2)>0.
 \end{gathered}                                        \tag{C.AD.5}\hypertarget{tag-009H}{}\label{cdp:eq:C.AD.5}
\]
Every $n_i$ has nonzero second coordinate, by the transverse
calculation. Changing one normal-product constant changes
its affine edge constant by an arbitrary multiple of $n_i$.
Choose it to make $b=(\sigma,i)$. This changes no elliptic
gluing map. In the all-ruled case it leaves the full horizontal
translation $\sigma$ unchanged, exactly as in \ref{cdp:st:C.2.UN.55}.

\par\medskip\noindent\textbf{The entire punctured sphere of actual tori.}
Use $V_z$ from (C.TL.3), and set
\[
 \begin{gathered}
 \Lambda(z)=V_z\Lambda,\qquad
 B(z)=V_z(b)+\frac{d}{2\pi i}\log z\,V_z(p),\\
                   T_z=\C^2/(\Lambda(z)+\Z B(z)).
 \end{gathered}                                        \tag{C.AD.6}\hypertarget{tag-009I}{}\label{cdp:eq:C.AD.6}
\]
These are holomorphic period sections for every $z\in\C^*$.
The change of a logarithm adds the integral period $dV_z(p)$.
The exact real determinant proved in (C.TL.5)--(C.TL.6), with
$N$ replaced by the positive integer $d$, is
\[
 C(1+|z|^2)\left(1+|z|^2+
          \frac{d}{2\pi}(|z|^2-1)\log|z|\right)>0.
                                                        \tag{C.AD.7}\hypertarget{tag-009J}{}\label{cdp:eq:C.AD.7}
\]
Thus \eqref{cdp:eq:C.AD.6} gives an actual smooth proper family of two-tori
on all $\C^*$, by the locally uniformly bounded lattice
matrices and parallelepipeds in \ref{cdp:st:C.2.UN.55}. No elliptic quotient
of this entire middle family is required.

\par\medskip\noindent\textbf{The exact moving components at zero.}
Transport every lattice vector by $V_z$, retaining its integral
coefficients. At a Hopf vertex put
\[
 P_i(z)=\left(\frac{V_z(n_i)}{2\pi i},
                       -\frac{V_z(n_{i-1})}{2\pi i}\right),
 \qquad \Lambda_i(z)=P_i(z)^{-1}\Lambda(z).
                                                        \tag{C.AD.8}\hypertarget{tag-009K}{}\label{cdp:eq:C.AD.8}
\]
The frame is invertible near zero. Its two distinguished
angular periods are literally $(2\pi i,0)$ and $(0,2\pi i)$;
the finite angular enlargement is retained. The finite quotient characters
are unchanged, since all their integer relations were
transported. An original contracting generator becomes a
holomorphic pair of logarithms, still with negative real
parts on a smaller disc. Exponentiation therefore gives an
actual smooth family of the marked primary or secondary
diagonal Hopf quotients. Its central member is the original
marked surface, including every finite character.

\item At a ruled vertex use its primitive normal vector and a
transverse integral lattice lift, as in (C.MX.10). Transporting
both gives the actual moving elliptic base, degree-zero line
and projective completion. This is valid on a smaller disc
because its two base periods remain real-linearly independent.
Its two marked sections are retained.

\item If $b_i$ is the original common affine edge constant after
the one permitted normal-product change, choose its moving
value to be $V_z(b_i)$. The common lattice and the paired
normal meridians give the required linear edge maps. Their
affine constants induce actual elliptic translations and
normal-line isomorphisms, extending over the marked curves.
The sum of the edge translations is $V_z(b)$; the sum of
their logarithmic coefficients is $dV_z(p)/(2\pi i)$.
Thus the proper moving plumbing with equation $xy=z$ has
exactly the periods \eqref{cdp:eq:C.AD.6}. Its logarithmic chart maps
identify its entire punctured family with that family:
each collar change is the stated affine map, every branch
change is an integral period, and the complete circuit is
exactly $B(z)$. The inverse is the closed-coframe inverse
already used in the local criterion. These maps and their
inverses vary holomorphically with $z$.

\par\medskip\noindent\textbf{The conjugate Hopf and ruled filling at infinity.}
Put $s=z^{-1}$ and divide the common vector frame by $z$.
Write
\[
 W_s(w)=J\overline w+s w,\qquad
              J=\begin{pmatrix}0&1\\-1&0\end{pmatrix}.
                                                        \tag{C.AD.9}\hypertarget{tag-009L}{}\label{cdp:eq:C.AD.9}
\]
The rank-three periods become $W_s\Lambda$, and the fourth
period is
\[
                 W_s(b)-\frac{d}{2\pi i}\log s\,W_s(p).
                                                        \tag{C.AD.10}\hypertarget{tag-009C}{}\label{cdp:eq:C.AD.10}
\]
Choose the outgoing normal at vertex $i$ to be
$n_i^\infty(s)=-W_s(n_i)$. Its incoming normal is
$-n_{i-1}^\infty(s)$ and its total normal period is
$-dW_s(p)$, exactly as required by \eqref{cdp:eq:C.AD.10}.

\item For a Hopf vertex the corresponding frame satisfies the
following exact identities at zero:
\[
 \begin{gathered}
 P_i^\infty(0)=J\overline{P_i},\\
 (P_i^\infty(0))^{-1}W_0\Lambda
                         =\overline{\Lambda_i}.
 \end{gathered}                                        \tag{C.AD.11}\hypertarget{tag-009D}{}\label{cdp:eq:C.AD.11}
\]
The signs follow from $\overline{2\pi i}=-2\pi i$.
Consequently the actual contracting logarithms are the
complex conjugates of the original ones, whose real parts
are still negative. The finite angular characters are
their conjugates and remain free. For small $s$, the same
holomorphic lattice construction as \eqref{cdp:eq:C.AD.8} gives the
smooth marked Hopf family through this conjugate surface.

\item For a ruled vertex the separate frame transforms by the
same $J$ and conjugation, with its angular normal changed
by the displayed minus sign. The resulting central ruled
surface is exactly the conjugate marked projective bundle;
the actual line-bundle lattice proves this identification.
Choose the affine constant on edge $i$ to be $W_s(b_i)$.
At zero its component-coordinate expression is the complex
conjugate of the original expression, so it induces the
conjugate elliptic gluing map. Its logarithmic coefficient
is $W_s(n_i)/(2\pi i)$, and their sum is exactly the one
in \eqref{cdp:eq:C.AD.10}.

\item The same proper normal-product plumbing, now with $xy=s$,
therefore has central support equal to the specified
complex conjugate support. Its actual punctured-family
logarithmic map identifies it with \eqref{cdp:eq:C.AD.10}. This proves
both the component and the entire-family comparisons at
infinity; complex conjugation is used only in the fixed
central coefficients of holomorphic period sections.

\item Finally glue these two proper fillings to \eqref{cdp:eq:C.AD.6} on
their punctured discs. The exact frame change \eqref{cdp:eq:C.AD.9}
and the two logarithmic biholomorphisms give the transition
maps. The finite closed-base refinement argument of \ref{cdp:st:C.2.UN.55}
proves the resulting total space Hausdorff and proper
over $\PP^1$, hence compact. Its displayed local models
are smooth. All middle fibres are smooth two-tori.
Every marked component and double elliptic curve has
Euler number zero, so both reduced singular supports
have Euler number zero as well.

\item The symbolic real determinant and the two conjugated
Hopf-frame identities are checked in
\path{code/root_all_diagonal_cycles.py} and its output.

\end{enumerate}
\end{proof}






\begin{proposition}[The complete canonical and period reduction for all-Hopf trees]\hypertarget{tag-01E3}{}\label{cdp:st:C.2.UN.58}
\StatementRecord{cdp:st:C.2.UN.58}{proved}{01E3}{\SeriesRef{cdp:st:C.2.UN.46}, \SeriesRef{cdp:st:C.2.UN.48}, \SeriesRef{cdp:st:C.2.UN.54}, \SeriesRef{cdp:st:C.2.UN.38}, \SeriesRef{cdp:st:C.2.UN.37}, \SeriesRef{cdp:eq:C.HT.4}, \SeriesRef{cdp:st:C.2.UN.50}, \SeriesRef{cdp:eq:C.HT.5}, \SeriesRef{cdp:eq:C.HT.7}, \SeriesRef{cdp:eq:C.HT.6}, \SeriesRef{cdp:eq:C.HT.9}, \SeriesRef{cdp:st:C.2.UN.49}, \SeriesRef{cdp:eq:C.HT.10}, \SeriesRef{cdp:eq:C.HT.11}}

Let $S$ be a full reduced SNC fibre of a smooth-total degeneration
of connected complex two-tori, with every irreducible component a
primary or secondary Hopf surface. If one component is a diagonal
Hopf surface of algebraic dimension zero, the whole support is
one of the diagonal cycles classified by \ref{cdp:st:C.2.UN.46} and \ref{cdp:st:C.2.UN.48}.

In every remaining case the incidence graph is the unicyclic graph
of \ref{cdp:st:C.2.UN.54}. Every cycle vertex and every nonterminal tree vertex
has algebraic dimension one. Its actual elliptic reduction is
\[
                    q_v:H_v\longrightarrow\PP^1,        \tag{C.HT.1}\hypertarget{tag-00N8}{}\label{cdp:eq:C.HT.1}
\]
with fixed Jacobian, identity integral monodromy and smooth multiple
supports. The other possible vertices are nonlinear
algebraic-dimension-zero Hopf leaves. Every multiple fibre of
\eqref{cdp:eq:C.HT.1} is a double curve of $S$.

For each edge $e$ the normal line has a finite exact order $m_e$.
It is the full fibre multiplicity on both elliptic Hopf endpoints,
and is the resonance at a nonlinear endpoint. The connected
normal-root covers of the two endpoints are identified by the
actual dual-normal map. Their Jacobians therefore identify along
every edge, and the product of their linear identifications around
the cycle is the identity. The actual elliptic cycle product can
still be a translation, whose value is retained by these maps.

Let $k_v$ be the canonical evaluation orders, and let $k_0$ be their
common minimum on the cycle. Root each attached tree at that cycle.
The heights are determined by the weighted graph:
\[
 h_v:=k_v-k_0=
 \begin{cases}
 0,&v\text{ is on the cycle},\\
 \displaystyle\sum_{e\text{ on the path from the cycle to }v}m_e,
       &v\text{ is in an attached tree}.
 \end{cases}                                           \tag{C.HT.2}\hypertarget{tag-00NC}{}\label{cdp:eq:C.HT.2}
\]
More precisely, after its invariant relative elliptic one-form is
factored out, the canonical form on a cycle vertex is a nonzero
constant multiple of $dz/z$, where the two cycle points are
$0,\infty$. On a tree vertex it is a nonzero constant multiple of
$dz$, where its parent point is $\infty$.

All regular elliptic-fibre two-tori on the algebraic-dimension-one
vertices, together with the real core two-tori on the nonlinear
leaves, represent one transported integral class $\gamma$ in the
nearby fibre. If a nonlinear leaf occurs, the single period
\[
                     P(t)=\int_{\gamma_t}\omega_t       \tag{C.HT.3}\hypertarget{tag-00ND}{}\label{cdp:eq:C.HT.3}
\]
of the relative canonical evaluation divided by $t^{k_0}$ has a
finite positive vanishing order $R$. Every nonlinear leaf has
height $R$, and every algebraic-dimension-one vertex has height
strictly smaller than $R$. These are assertions about the actual
common period. When there is no nonlinear leaf its period is
allowed to vanish identically.
\end{proposition}



\begin{proof}

\begin{enumerate}[label={$\langle1\rangle$\arabic*.},leftmargin=*]
\item \par\medskip\noindent\textbf{Separation by the actual normal lines.}
The axes of an algebraic-dimension-zero diagonal primary Hopf
surface have nontorsion normal lines. For example, torsion of the
normal of $w=0$ for $(z,w)\mapsto(\alpha z,\beta w)$ would give
$\beta^n=\alpha^a$ for some integers $n>0,a$. The corresponding
monomial quotient is a nonconstant meromorphic invariant, contrary
to algebraic dimension zero. This is the exact automorphy
calculation of \ref{cdp:st:C.2.UN.38}. Pullback to the finite primary cover
proves nontorsion on a secondary quotient as well.

\item Every curve on an algebraic-dimension-one Hopf surface is a fibre
support of \eqref{cdp:eq:C.HT.1}: a horizontal curve would meet a regular elliptic
fibre, contradicting $H^2(H,\Q)=0$ and positivity of intersection.
Its normal has finite order, as checked explicitly below. The
nonlinear Hopf axis also has finite order equal to its resonance,
by (C.SN.3)--(C.SN.4). Normal duality at an actual reduced SNC edge
preserves torsion and nontorsion. Thus a diagonal
algebraic-dimension-zero vertex can meet only another such vertex.
Connectedness then puts every component in that class, and
\ref{cdp:st:C.2.UN.54} gives precisely the already classified diagonal cycle.

\item In the remaining case every algebraic-dimension-zero vertex is
nonlinear and has just its one elliptic curve. It is a leaf.
All other vertices have \eqref{cdp:eq:C.HT.1} by \ref{cdp:st:C.2.UN.37}. In particular every
cycle or nonterminal tree vertex has that fibration.

\par\medskip\noindent\textbf{The root charts and the invariant relative form.}
At a full fibre $mJ$ of \eqref{cdp:eq:C.HT.1}, its actual normalized $m$th root
cover is a connected smooth elliptic family, finite unramified over
the original surface tube. The fixed modulus and identity linear
monodromy give a local principal bundle with fixed elliptic fibre
$F$. It is holomorphically trivial on a disc: the exponential
sequence $0\to\Lambda_F\to\OO\to F\to0$, together with
$H^1(\Delta,\OO)=H^2(\Delta,\Lambda_F)=0$, kills its translation
torsor. Thus the original tube is
\[
       (F\times\Delta_s)/\mu_m,\qquad
       (z,s)\longmapsto(z+\tau,\zeta_m s),\qquad
       q=s^m,\quad \operatorname{ord}(\tau)=m.          \tag{C.HT.4}\hypertarget{tag-00NE}{}\label{cdp:eq:C.HT.4}
\]
Here the translation can be made constant by an actual gauge.
If its initially holomorphic expression is $\tau+A(s)$, choose a
small complex lift with $\sum_{j=0}^{m-1}A(\zeta_m^j s)=0$ and put
$b(s)=-m^{-1}\sum_{j=0}^{m-1}jA(\zeta_m^j s)$.
Then $b(\zeta_ms)-b(s)=-A(s)$, so $z\mapsto z+b(s)$ gives \eqref{cdp:eq:C.HT.4}.
The exact order follows from the free deck action on the connected
central elliptic curve.

\item The normal of $J=F/\langle\tau\rangle$ is the associated line for
the primitive character $\zeta_m$. It has exact order $m$:
if a smaller tensor power were trivial, its pullback trivialization
would be a nowhere-zero holomorphic function on the compact $F$
with a nontrivial character, which is impossible. The connected
normal-root cover is exactly $F\to J$. Duality of the two normal
lines at an edge identifies their exact orders and their connected
root covers. Hence the number $m_e$ in the statement is well defined,
including at a nonlinear leaf via the exact secondary charts of
\ref{cdp:st:C.2.UN.50}.

\item In every root chart write $dz$ for a fixed nonzero invariant
elliptic one-form. Transition maps between principal elliptic
charts are translations; their differentials differ by a multiple
of the base differential. Consequently
\[
                         q^*\theta\wedge dz             \tag{C.HT.5}\hypertarget{tag-00NF}{}\label{cdp:eq:C.HT.5}
\]
is a globally defined meromorphic two-form whenever $\theta$ is a
meromorphic differential on $\PP^1$. The finite translations in
\eqref{cdp:eq:C.HT.4} preserve it. This constructs the needed form rather than
postulating a global absolute elliptic one-form.

\item Every meromorphic two-form whose divisor is vertical has precisely
the form \eqref{cdp:eq:C.HT.5}. Its ratio to one such reference form is a
meromorphic function with vertical divisor. On a regular compact
elliptic fibre away from that divisor the ratio is a nonzero
holomorphic function and is constant. It therefore descends on
the regular base. In \eqref{cdp:eq:C.HT.4} its descent is a meromorphic function
of $s$ invariant under $s\mapsto\zeta_ms$, and hence is a
meromorphic function of $s^m$. This proves descent across every
multiple fibre as well. In this identification the exact order is
\[
 \operatorname{ord}_J(q^*\theta\wedge dz)
               =m\,\operatorname{ord}_{q(J)}\theta+m-1 .
                                                        \tag{C.HT.6}\hypertarget{tag-00NG}{}\label{cdp:eq:C.HT.6}
\]

\par\medskip\noindent\textbf{The canonical equation and the leaf induction.}
The actual canonical evaluation gives on every component a
nonzero meromorphic two-form $\eta_v$ with divisor
\[
 \operatorname{div}(\eta_v)
       =\sum_{e=(v,w)}(k_w-k_v-1)J_e .                 \tag{C.HT.7}\hypertarget{tag-00NH}{}\label{cdp:eq:C.HT.7}
\]
This is adjunction with the specified vertical orders. It can also
be read in the double chart $t=xy$ by dividing the evaluation by
$t^{k_v}$ and taking the residue after division by $t$. Away from
the boundary there is no zero or pole, since the evaluation has
only its central component divisors. Equation \eqref{cdp:eq:C.HT.7} is a divisor
of the actual section, not just numerical equivalence.

\item Apply \eqref{cdp:eq:C.HT.6} at an algebraic-dimension-one vertex. An unmarked
multiple fibre would have order zero in \eqref{cdp:eq:C.HT.7}, requiring
$m-1\equiv0\pmod m$ and hence $m=1$. There are therefore no
unmarked multiple fibres. At a marked fibre \eqref{cdp:eq:C.HT.6} says
\[
 d_{v,e}:=\frac{k_w-k_v}{m_e}
       =\operatorname{ord}_{q_v(J_e)}\theta_v+1\in\Z.
                                                        \tag{C.HT.8}\hypertarget{tag-00NI}{}\label{cdp:eq:C.HT.8}
\]
Every other point of $\PP^1$ has order zero. Taking the degree of
the divisor of this rational differential proves
\[
                    \sum_{e\ni v}d_{v,e}
                           =\operatorname{val}(v)-2.   \tag{C.HT.9}\hypertarget{tag-00NJ}{}\label{cdp:eq:C.HT.9}
\]
Opposite orientations give opposite $d$ because the same exact
normal order $m_e$ occurs at the two endpoints.

\ref{cdp:st:C.2.UN.54} identifies the cycle vertices exactly as the minimum
canonical-order vertices. Their $k_v$ are one common $k_0$,
so both cycle-edge differences are zero. A terminal
algebraic-dimension-one vertex has just its parent edge, and
\eqref{cdp:eq:C.HT.9} gives $k_{\rm leaf}-k_{\rm parent}=m_e$.
For a nonlinear leaf the invariant form
$dz\wedge dw/w^{r+1}$ and the infinite order of its effective
axis line give exactly the same equality, with $m_e=r$,
as in (C.AS.12).

\item Now work inward along any finite attached tree. Suppose all the
$c$ child edges of a vertex have outgoing $d=1$. At that vertex
the one parent term is the negative of the parent's outgoing
$d_{\rm parent}$, and \eqref{cdp:eq:C.HT.9} gives
$-d_{\rm parent}+c=c-1$. Thus $d_{\rm parent}=1$.
Every nonterminal tree vertex has algebraic dimension one, so this
induction reaches the cycle. It proves \eqref{cdp:eq:C.HT.2} on every edge.

\item At a cycle vertex the two cycle terms have $d=0$, giving simple
poles of $\theta_v$ there; every child term has $d=1$ and order
zero. There are no other zeros or poles. A coordinate taking the
cycle points to $0,\infty$ therefore gives $\theta_v=c\,dz/z$.
At a tree vertex the parent term has $d=-1$ and order $-2$;
all child terms have order zero. A coordinate with its parent
point at infinity gives $\theta_v=c\,dz$. In each case $c\ne0$.
This proves the more precise differential assertion.

\par\medskip\noindent\textbf{The actual Jacobian comparison around the cycle.}
The same root charts show $R^1q_{v*}\OO_{H_v}=\OO_{\PP^1}$.
On $F\times\Delta_s$ the cohomology frame is the fixed
$H^1(F,\OO_F)$; finite deck translations act trivially, so its
invariant functions are exactly functions of $s^m$. On overlaps
the identity linear monodromy and translations give the same
frame. Leray on $\PP^1$ therefore identifies
$H^1(H_v,\OO_{H_v})$ with this one-dimensional constant fibre
cohomology.

\item Restriction to $J_e$ followed by pullback along the actual finite
cover $F\to J_e$ agrees with evaluation on the regular elliptic
fibre. This follows from the commuting pullback diagram of the
surface tube and \eqref{cdp:eq:C.HT.4}, where both maps are restriction of the
same constant cohomology frame. Thus the scalar cycle product
in (C.HR.2) is exactly the action on $H^1(F,\OO_F)$ of the product
of the lifted fibre isomorphisms. It is one by \ref{cdp:st:C.2.UN.54}.
An elliptic automorphism whose action on $H^1(\OO)$ is the identity
has linear part one: after choosing an origin, its lift is
$z\mapsto az+b$, and the nonzero invariant one-form detects $a$
(or its conjugate on $H^1(\OO)$). Hence that product is a
translation. Different choices of a lift through a finite
normal-root cover alter it only by deck translations. The
Jacobian product is therefore the identity with all maps retained.

\par\medskip\noindent\textbf{One integral core class and its nonzero nonlinear periods.}
We describe the cycle comparison before integrating. Near an
edge the two surface tubes are the disk bundles of the actual
torsion normal lines in \eqref{cdp:eq:C.HT.4}. Their dual-normal isomorphism
identifies their connected root covers. On that cover the
normal-product collar is
\[
                  xy=t,\qquad q_v=x^{m_e},\quad
                  q_w=y^{m_e}.                         \tag{C.HT.10}\hypertarget{tag-00N9}{}\label{cdp:eq:C.HT.10}
\]
Multiplying one base coordinate by a nonzero constant makes these
equalities exact: the two trivializations of the $m_e$th normal
power differ by a holomorphic unit on a compact elliptic curve,
and hence by a constant. The regular fibre $x^{m_e}=c$ on the
first side is sent to $y^{m_e}=t^{m_e}/c$ on the other. Both are
the same connected elliptic cover of $J_e$, with its complex
orientation. Varying the nonzero normal level supplies the
oriented isotopy in the neck.

\item For an arbitrary given smoothing, the product-preserving smooth
tubular comparison used in \ref{cdp:st:C.2.UN.49} carries these normal collars
to the actual $xy=t$ collars. Compact cores are compared by
Ehresmann transport. Consequently \eqref{cdp:eq:C.HT.10} is an equality of
the transported integral two-cycle classes in the actual nearby
fibre, without a holomorphic linearization of its neighbourhood.
Regular fibres on one component are isotopic through its regular
base, so every algebraic-dimension-one vertex gives the same
class $\gamma$. At a nonlinear leaf the exact unimodular
boundary maps (C.SN.6)--(C.SN.8), together with the displayed
smooth-pair deformation, identify its real core with this same
oriented class. They retain the full secondary order.

\item Write $\omega_t$ for the evaluation divided by $t^{k_0}$.
On a compact regular fibre of $H_v$ its leading coefficient is
$t^{h_v}\eta_v$. The restriction of the meromorphic two-form
$\eta_v$ to this holomorphic curve is zero. Therefore
\[
                     P(t)=O(t^{h_v+1})
       \quad\text{at every algebraic-dimension-one vertex}.
                                                        \tag{C.HT.11}\hypertarget{tag-00NA}{}\label{cdp:eq:C.HT.11}
\]
These are actual periods of a closed relative holomorphic form
over the same compact class. They are holomorphic by
differentiation and Stokes' theorem; the local compact-core
transport gives the displayed extension and leading coefficient.

\item At a nonlinear leaf
$H_r(\beta)/\mu_n$ its unique-axis form is
$\eta=c\,dz\wedge dw/w^{r+1}$, $c\ne0$. In
$U=z/w^r$, $V=\log w$, the two primitive real core periods are
$(0,2\pi i/n)$ and $(\beta^{-r},\log\beta)$. Hence
\[
 P(t)=t^{h_{\rm leaf}}
       \left(\frac{\pm2\pi i c\beta^{-r}}{n}+O(t)\right),
                    \qquad c\ne0.                     \tag{C.HT.12}\hypertarget{tag-00NB}{}\label{cdp:eq:C.HT.12}
\]
The coefficient is nonzero. One such leaf fixes the finite
vanishing order $R=h_{\rm leaf}>0$ of the single period.
Every other nonlinear leaf has that same height, while
\eqref{cdp:eq:C.HT.11} gives $h_v<R$ at every elliptic Hopf vertex.
No nonvanishing conclusion is imposed when there is no nonlinear
leaf.

\item The exact root congruence, tree equations, height sums, collar
monomials and common-period inequalities are checked in
\path{code/E_hopf_tree_heights.py} and its saved output. The
argument proves these comparisons for every actual support in
the stated branch. It does not assert that these necessary
equations alone smooth arbitrary abstract attached-tree data.

\end{enumerate}
\end{proof}






\begin{proposition}[All multiplicities and all diagonal Hopf and ruled vertices]\hypertarget{tag-01E4}{}\label{cdp:st:C.2.UN.59}
\StatementRecord{cdp:st:C.2.UN.59}{proved}{01E4}{\SeriesRef{cdp:st:C.2.UN.40}, \SeriesRef{cdp:st:C.2.UN.43}, \SeriesRef{cdp:st:C.2.UN.52}, \SeriesRef{cdp:st:C.2.UN.56}, \SeriesRef{cdp:st:C.2.UN.51}, \SeriesRef{cdp:st:C.2.UN.55}, \SeriesRef{cdp:eq:C.EW.3}, \SeriesRef{cdp:eq:C.EW.6}, \SeriesRef{cdp:eq:C.EW.9}, \SeriesRef{cdp:eq:C.EW.10}, \SeriesRef{cdp:eq:C.EW.13}, \SeriesRef{cdp:eq:C.EW.15}, \SeriesRef{cdp:eq:C.EW.14}, \SeriesRef{cdp:st:C.2.UN.57}, \SeriesRef{cdp:eq:C.EW.12}, \SeriesRef{cdp:eq:C.EW.1}, \SeriesRef{cdp:eq:C.EW.19}, \SeriesRef{cdp:eq:C.EW.20}, \SeriesRef{cdp:eq:C.EW.17}, \SeriesRef{cdp:eq:C.EW.18}}

Specify a simple-NC cycle of smooth embedded marked surfaces, with
edges between distinct components. Each vertex is a diagonal primary
or secondary Hopf surface with its two axes, or
$\PP(\OO_E\oplus L)$ with $E$ elliptic and
$L\in\operatorname{Pic}^0(E)$, marked by its two sections.
Specify their elliptic gluing maps, positive multiplicities $m_i$,
and the actual compatible weighted normal lines, sections and
boundary maps of (C.WH.1)--(C.WH.2).

Put $N=\operatorname{lcm}_i m_i$. Form the actual normalized
$N$th-root component covers and the compressed necks $uv=s^{h_i}$.
The root graph is connected and its rank-three integral
complement-lattice product is the identity precisely when the
specified weighted normal plumbing has connected two-torus fibres.
Every actual smooth-total torus degeneration with the specified
marked support yields these data and passes this test.

Every weighted support passing the test occurs over $\PP^1$ in a
smooth compact complex threefold. The construction retains all
marked components, elliptic gluing maps and multiplicities. Its
only other singular fibre is the complex conjugate weighted
support. All reduced fibre supports have Euler number zero.
Normal-product constants are chosen in this compact realization.
There is no algebraic-dimension restriction on the Hopf vertices,
and an all-ruled cycle is included.

The complete component, axis and inertia comparisons in
(C.WH.0), and the ruled-cover identities (C.WM.2)--(C.WM.7),
hold in this scope. In common root-complement coordinates the
finite deck linear part fixes the primitive total normal period.
Its other eigenvalue is a root of unity of order $1,2,3,4$ or $6$.
The compact construction retains this linear part and every
affine deck constant.
\end{proposition}



\begin{proof}\leavevmode\par



 
\begin{enumerate}[label={$\langle1\rangle$\arabic*.},leftmargin=0pt,labelsep=.5em]
\item \par\medskip\noindent\textbf{Actual root maps and the local criterion.}
The weighted line on a diagonal Hopf vertex has a flat tensor
power, since both axis divisors have the explicit constant
automorphy factors of \ref{cdp:st:C.2.UN.40} and \ref{cdp:st:C.2.UN.43}. Dividing the transferred
connection by its tensor exponent gives the actual flat line
used in \ref{cdp:st:C.2.UN.52}. Its root cover on the axis complement is the
finite-index logarithmic lattice cover constructed there.
Coordinate powers extend it across both axes, and the actual
contracting generator and free finite angular characters identify
its finite normalization with a smooth diagonal Hopf surface.
This comparison does not require nonresonant multipliers.


\begin{enumerate}[label={$\langle2\rangle$\arabic*.},leftmargin=1.3em,labelsep=.5em]
\item  On a ruled vertex the proof (C.WM.4)--(C.WM.7) applies without
any Hopf assumption. It gives the actual finite normalization
$\PP(\OO_{E'}\oplus L')\to\PP(\OO_E\oplus L)$, with
an \'{e}tale base map $\pi$, fibre power $q$, and
$(L')^q\simeq\pi^*L$. In particular
$m\mid a+b$ and $q=m/\gcd(m,a)=m/\gcd(m,b)$.
The weighted normal line is pulled back on each entire section
chart by its transferred flat connection; its nonzero global
degree on a ruling fibre is retained as in \ref{cdp:st:C.2.UN.56}.


\item  At an original double curve put
$a=gA$, $b=gB$, $\gcd(A,B)=1$ and
$h=N/(gAB)$. The finite normalization is still exactly
\[
             uv=s^h,\qquad x=u^B,\quad y=v^A,
             \qquad t=s^N=x^a y^b.
\tag{C.EW.1}\hypertarget{tag-00JH}{}\label{cdp:eq:C.EW.1}
\]
The actual root component and axis covers, their connected
degrees, and all stabilizers therefore have precisely the
maps and orders proved in \ref{cdp:st:C.2.UN.52} and \ref{cdp:st:C.2.UN.56}. Each strict component
has two connected marked elliptic curves, so the compressed
root graph is a union of cycles. Resolving a neck inserts the
specified degree-zero ruled bridges and changes no complement
lattice product.


\item  Connectedness of this graph is equivalent to connectedness of
the root family, by the proper Stein factor and its reduced
central fibre. On a connected graph the actual core-and-neck
calculation gives $\pi_1(F_s)=\Lambda\rtimes_M\Z$.
Thus a torus forces $M=I$. Conversely $M=I$ glues the closed
logarithmic coframe; its complete commuting fields identify
the nearby compact fibre with $\C^2$ modulo its rank-four
period lattice. For an arbitrary actual smoothing the
canonical divisor sections and their small smooth unit roots
give the same weighted tubular and collar maps, exactly as
in \ref{cdp:st:C.2.UN.52}. This proves the asserted local criterion for the
actual maps, as well as for the explicit plumbing.


\end{enumerate}

\item \par\medskip\noindent\textbf{The exact finite linear and affine data.}
Assume that the test passes and identify all open root
complements with $G=\C^2/\Lambda$ by the actual linear
edge maps. Let $n_j$ be the outgoing normal meridian at
vertex $j$, and let $h_j$ be its outgoing compressed
neck exponent. The incoming meridian is $-n_{j-1}$.
The root generator lies over the identity on each original
marked normalization and preserves its two marked branches.
It consequently induces an orientation-preserving rotation
$\sigma$ of the connected root cycle. The common number
$r$ of strict covers and every inertia order are therefore
exactly those of (C.WH.0).


\begin{enumerate}[label={$\langle2\rangle$\arabic*.},leftmargin=1.3em,labelsep=.5em]
\item  The generator is affine on the common complement. For clarity,
an isomorphism of diagonal Hopf surfaces preserving both axes
lifts to a biholomorphism of $\C^2\setminus0$, extends with
its inverse across zero, and has coordinates $zU(z,w)$ and
$wV(z,w)$, with $U,V$ units. Conjugating a contracting deck
generator and evaluating at zero makes the corresponding
multiplier of each unit equal to one. Iteration of the
contraction then makes each unit constant. Thus this lift
is diagonal even for resonant multipliers. For a ruled
vertex the actual base isomorphism and marked line-bundle
isomorphism have affine logarithmic lifts by the constant
factors of automorphy, as in \ref{cdp:st:C.2.UN.51}. Compatibility with the
actual neck maps makes the linear part a single map $R$.
Write the generator and the edge maps as
\[
 \begin{aligned}
 g(j,v,s)&=(\sigma j,Rv+\tau_j,\zeta_Ns),\\
 e_j(v,s)&=v+b_j+l_j\log s,\\
 l_j&=-\frac{h_j n_j}{2\pi i}.
 \end{aligned}
\tag{C.EW.2}\hypertarget{tag-00JS}{}\label{cdp:eq:C.EW.2}
\]
Here $R\Lambda=\Lambda$, $R^N=I$,
$Rn_j=n_{\sigma j}$ and $h_{\sigma j}=h_j$.
Actual composition of these maps, with
$\log\zeta_N=2\pi i/N$, gives lattice elements
$\lambda_j,\mu_j\in\Lambda$ satisfying
\[
 \begin{aligned}
 \tau_{j+1}+Rb_j-\tau_j-b_{\sigma j}
             +\frac{h_j n_{\sigma j}}N&=\lambda_j,\\
 \sum_{a=0}^{N-1}R^{N-1-a}\tau_{\sigma^a j}&=\mu_j.
 \end{aligned}
\tag{C.EW.3}\hypertarget{tag-00JV}{}\label{cdp:eq:C.EW.3}
\]
These are equalities in the actual logarithmic vector space.


\item  Put
\[
             b=\sum_jb_j,\qquad
             u=-\sum_jh_jn_j=d p\in\Lambda,
\tag{C.EW.4}\hypertarget{tag-00JW}{}\label{cdp:eq:C.EW.4}
\]
where $p$ is primitive and $d>0$. The contraction-height
calculation (C.AD.3) applies at every diagonal Hopf vertex;
at a ruled vertex the consecutive outgoing normals agree.
Thus for $W=\operatorname{span}_{\R}\Lambda$ and
$\ker\ell=W$, all numbers $\ell(i n_j)$ have one
nonzero sign. In the all-ruled case the same assertion
is the actual elliptic quotient calculation of \ref{cdp:st:C.2.UN.55}.
Consequently $\ell(iu)\ne0$. Rotation and the displayed
normal maps give $Ru=u$, hence $Rp=p$.


\item  The complex line $W\cap iW$ is $R$-invariant and is
complementary to $\C p$. Choose the complex coordinates
given by these two lines, with $p=(0,1)$. Then
\[
 \begin{gathered}
 W=\C\times\R,\qquad R=\operatorname{diag}(\rho,1),\\
 \Lambda=\langle w_1=(a_1,c_1),\ w_2=(a_2,c_2),\ p\rangle_{\Z},\\
 c_1,c_2\in\R,\qquad
 C=\operatorname{Im}(\overline{a_1}a_2)>0.
 \end{gathered}
\tag{C.EW.5}\hypertarget{tag-00JX}{}\label{cdp:eq:C.EW.5}
\]
The projection of $\Lambda$ to the first coordinate is
the actual elliptic lattice. Since $R$ preserves it,
$\rho$ has order $1,2,3,4$ or $6$: its integral real
two-dimensional matrix has determinant one, trace
$2\operatorname{Re}\rho\in\Z$, and finite order.
Also $\rho^N=1$. The stabilizer of a root vertex acts
as the deck translations of its finite open-group cover,
so in fact $R^r=I$.


\item  Summing \eqref{cdp:eq:C.EW.3} retains the full affine descent relation
\[
                         (R-I)b-u/N\in\Lambda.
\tag{C.EW.6}\hypertarget{tag-00JY}{}\label{cdp:eq:C.EW.6}
\]
When $R=I$ this recovers $u/N\in\Lambda$ from \ref{cdp:st:C.2.UN.52}.
Equation \eqref{cdp:eq:C.EW.6}, with its possibly nontrivial $R$, is
the relation used below.


\item  Adding $c n_j$ to $b_j$ on a whole $\sigma$-orbit,
with the same $c\in\C$, changes only the normal-product
constants. It preserves both equations \eqref{cdp:eq:C.EW.3}, since
$Rn_j=n_{\sigma j}$. The second coordinates of all
$n_j$ are real and have the same nonzero sign, by the
height calculation. The sum on any such orbit therefore
has nonzero second coordinate. Choose $c$ on one orbit
so that
\[
                              b=(\sigma_0,i).
\tag{C.EW.7}\hypertarget{tag-00JZ}{}\label{cdp:eq:C.EW.7}
\]
Every original elliptic gluing map and every finite
action map on the central support is retained.
The orbit-normal sum is also $R$-fixed; when
$\rho\ne1$ its first coordinate is zero. Thus the
entire common horizontal affine constant is retained
in that case as well.


\end{enumerate}

\item \par\medskip\noindent\textbf{A finite-equivariant period family on the punctured sphere.}
Choose $k\in\{1,\ldots,N\}$ with
$\rho=\zeta_N^k$, using $k=N$ for $\rho=1$.
Let
\[
 J=\begin{pmatrix}0&1\\-1&0\end{pmatrix},\qquad
 V_s(w)=w+s^kJ\overline w.
\tag{C.EW.8}\hypertarget{tag-00K0}{}\label{cdp:eq:C.EW.8}
\]
For each fixed vector $w$ this is holomorphic in $s$.
The exact relation $RJ\overline R^{-1}=\rho J$
gives
\[
                           R V_s=V_{\zeta_Ns}R.
\tag{C.EW.9}\hypertarget{tag-00K1}{}\label{cdp:eq:C.EW.9}
\]
Transport all integral vectors by $V_s$, and define
\[
 \begin{gathered}
 \Lambda_s=V_s\Lambda,\qquad
 B_s=V_s b+\frac{d\log s}{2\pi i}V_s p,\\
                         T_s=\C^2/(\Lambda_s+\Z B_s).
 \end{gathered}
\tag{C.EW.10}\hypertarget{tag-00JI}{}\label{cdp:eq:C.EW.10}
\]
A logarithm change adds $dV_s p$, an actual integral
period. The determinant calculation of \ref{cdp:st:C.2.UN.55} with
$z=s^k$ and logarithmic coefficient $d\log s/(2\pi i)$
gives, for $r_s=|s|>0$,
\[
 C(1+r_s^{2k})
 \left(1+r_s^{2k}
       +\frac d{2\pi}(r_s^{2k}-1)\log r_s\right)>0.
\tag{C.EW.11}\hypertarget{tag-00JJ}{}\label{cdp:eq:C.EW.11}
\]
Thus these are actual compact two-tori for every
$s\in\C^*$, and form a smooth proper holomorphic
family there. Properness follows from the uniformly
bounded period matrices and inverses on compact
base charts, exactly as in \ref{cdp:st:C.2.UN.55}.


\begin{enumerate}[label={$\langle2\rangle$\arabic*.},leftmargin=1.3em,labelsep=.5em]
\item  The transported finite action and edge data are explicitly
\[
 \begin{aligned}
 b_j(s)&=V_s b_j,& n_j(s)&=V_s n_j,\\
 l_j(s)&=-h_jn_j(s)/(2\pi i),&
 \tau_j(s)&=V_{\zeta_Ns}\tau_j.
 \end{aligned}
\tag{C.EW.12}\hypertarget{tag-00JK}{}\label{cdp:eq:C.EW.12}
\]
Notice that the translation is expressed in the
\emph{target} frame. Applying \eqref{cdp:eq:C.EW.9} to the first
equation of \eqref{cdp:eq:C.EW.3} gives exactly
\[
 \begin{aligned}
 &\tau_{j+1}(s)+Rb_j(s)-\tau_j(s)-b_{\sigma j}(\zeta_Ns)\\
 &\qquad-l_{\sigma j}(\zeta_Ns)\,2\pi i/N
       =V_{\zeta_Ns}\lambda_j.
 \end{aligned}
\tag{C.EW.13}\hypertarget{tag-00JL}{}\label{cdp:eq:C.EW.13}
\]
The logarithm correction transports because
$h_jn_{\sigma j}/N$ is a real rational multiple
of an integral vector. For the order equation,
\[
 R^{N-1-a}V_{\zeta_N^{a+1}s}
                         =V_sR^{N-1-a},
\]
so the $N$th iterate translates by $V_s\mu_j$.
This proves the full finite action, including all
overlap and order equalities. Likewise
\[
                    R B_s-B_{\zeta_Ns}
                   =V_{\zeta_Ns}\bigl((R-I)b-u/N\bigr)
                   \in\Lambda_{\zeta_Ns}.
\tag{C.EW.14}\hypertarget{tag-00JM}{}\label{cdp:eq:C.EW.14}
\]


\item  Here is an explicit global comparison with \eqref{cdp:eq:C.EW.10}.
On a logarithm chart put
$\theta_j(s)=\sum_{a<j}(b_a(s)+l_a(s)\log s)$.
The common coordinate of the $j$th component chart is
$v_j-\theta_j(s)$. The generator on the period quotient
therefore has the affine translation
\[
 R\theta_j(s)+\tau_j(s)-\theta_{\sigma j}(\zeta_Ns).
\tag{C.EW.15}\hypertarget{tag-00JN}{}\label{cdp:eq:C.EW.15}
\]
Changing $j$ changes this by an integral period, by
\eqref{cdp:eq:C.EW.13}; passing the last edge adds exactly $B_s$.
Changing the logarithm changes each partial sum by
an integral vector. Thus \eqref{cdp:eq:C.EW.15} defines a single
holomorphic finite action on the entire punctured
period family. Its order is the order already proved
in \eqref{cdp:eq:C.EW.13}, and its linear and period maps are
\eqref{cdp:eq:C.EW.9} and \eqref{cdp:eq:C.EW.14}.


\end{enumerate}

\item \par\medskip\noindent\textbf{The actual finite filling at zero.}
Near zero construct the moving strict components from
their transported lattices. At a Hopf vertex use
\[
 P_j(s)=\left(\frac{V_s n_j}{2\pi i},
                        -\frac{V_s n_{j-1}}{2\pi i}\right).
\tag{C.EW.16}\hypertarget{tag-00JO}{}\label{cdp:eq:C.EW.16}
\]
It is invertible near zero, retains the distinguished
angular periods and their finite angular enlargement,
and transports the original contracting logarithms
holomorphically. Their real parts remain negative.
It therefore defines the actual marked Hopf family.
At a ruled vertex use its primitive normal and an
independent projected elliptic lattice lift, as in
\ref{cdp:st:C.2.UN.56} and \ref{cdp:st:C.2.UN.57}. This gives the actual moving
degree-zero line and its projective completion.


\begin{enumerate}[label={$\langle2\rangle$\arabic*.},leftmargin=1.3em,labelsep=.5em]
\item  Use the edge constants \eqref{cdp:eq:C.EW.12} and keep the literal
compressed hypersurfaces $uv=s^{h_j}$. They give the
proper moving plumbing, whose exact periods are
\eqref{cdp:eq:C.EW.10}. The component logarithmic maps and their
closed-coframe inverses give a holomorphic
biholomorphism of its entire punctured family with
that period family. Formula \eqref{cdp:eq:C.EW.15} makes this
biholomorphism equivariant.


\item  The finite maps extend across all marked curves.
At a Hopf vertex
$RP_j(s)=P_{\sigma j}(\zeta_Ns)$, so the logarithmic
action in these component coordinates is a translation;
its exponential is the actual diagonal map on the
marked Hopf family. At a ruled vertex the same
normal-period map and its elliptic base map give the
actual line-bundle map and homogeneous projective
extension from \ref{cdp:st:C.2.UN.56}. The overlap equalities
\eqref{cdp:eq:C.EW.13} extend them across the compressed necks.


\item  The smoothness of the finite quotient is the same
local ring calculation as (C.WH.12), which uses
the actual inertia characters rather than $R=I$.
At \eqref{cdp:eq:C.EW.1} the pointwise axis stabilizer has order
$hAB$. Its order-$h$ subgroup fixes $u,v$ and rotates
$s$; the remaining independent characters have
orders $B,A$. These finite characters retain their
specified central values throughout the moving
families. We give the actual character coordinates.
The stabilizer preserves each adjacent component,
so it sends each of $u,v$ to itself times a unit.
Reynolds-average these functions to their specified
central eigencharacters. The results are $u'=uU$
and $v'=vV$, with $U,V$ units. Average the coordinate
along the pointwise fixed axis to an invariant
coordinate; its derivative there is unchanged.
The character identity
$\chi_u\chi_v=\chi_s^h$ and the equation
$u'v'=s^hUV$ show that $A=UV$ is invariant.
Replace $u'$ by $u'/A$. The product is then exactly
$s^h$, both eigencharacters are retained, and the
actual base coordinate $s$ has not changed. The
coordinate change has invertible differential at
the axis and is an analytic local isomorphism.
In these exact coordinates the ring is
\[
 \bigl(\C\{u,v,s\}/(uv-s^h)\bigr)^{J_{\rm axis}}
                         =\C\{u^B,v^A\}
                         =\C\{x,y\},
 \qquad t=x^a y^b.
\tag{C.EW.17}\hypertarget{tag-00JP}{}\label{cdp:eq:C.EW.17}
\]
The residual axis cover is \'{e}tale. At a general
component point the inertia is the normal reflection
of order $N/m_i$. Thus the quotient is smooth at
every point and has exactly the original full
multiplicities. At zero the covers, central maps and
finite action are the original ones; the quotient
support is therefore the specified marked surface.


\end{enumerate}

\item \par\medskip\noindent\textbf{The conjugate finite filling at infinity.}
Put $w=s^{-1}$ and divide the vector frame by $s^k$.
For a fixed vector $c$ write
\[
 W_w(c)=J\overline c+w^k c.
\tag{C.EW.18}\hypertarget{tag-00JQ}{}\label{cdp:eq:C.EW.18}
\]
The periods in this frame are $W_w\Lambda$ and
\[
                    W_wb-\frac{d\log w}{2\pi i}W_wp.
\tag{C.EW.19}\hypertarget{tag-00JR}{}\label{cdp:eq:C.EW.19}
\]
Choose outgoing normals $-W_wn_j$, the same neck
exponents $h_j$, and affine edge constants $W_wb_j$.
Their logarithmic coefficients are
$h_jW_wn_j/(2\pi i)$, so their complete circuit is
exactly \eqref{cdp:eq:C.EW.19}.


\begin{enumerate}[label={$\langle2\rangle$\arabic*.},leftmargin=1.3em,labelsep=.5em]
\item  The actual transformed deck generator is
\[
 \begin{aligned}
 (j,v,w)&\longmapsto
     (\sigma j,R_\infty v+\tau_{j,\infty}(w),\zeta_N^{-1}w),\\
 R_\infty&=\rho^{-1}R=J\overline R J^{-1},\\
 \tau_{j,\infty}(w)
     &=J\overline{\tau_j}+\rho^{-1}w^k\tau_j
       =W_{\zeta_N^{-1}w}\tau_j.
 \end{aligned}
\tag{C.EW.20}\hypertarget{tag-00JT}{}\label{cdp:eq:C.EW.20}
\]
Indeed these formulas follow by multiplying the image
vector $Rv+\tau_j(s)$ by $(\zeta_Ns)^{-k}$.
They retain all finite equations, since they are the
actual frame change of the action already constructed.


\item  At $w=0$ the Hopf frame is exactly
$J\overline{P_j(0)}$, as in (C.AD.11).
The angular normals have the displayed minus sign
because conjugation changes $2\pi i$ to $-2\pi i$.
Thus the contracting logarithms and every finite
angular character are their actual complex
conjugates. The separate ruled frame gives the
conjugate line bundle and marked sections by the
same logarithmic calculation. Each affine edge
constant, in its component coordinates, is the
conjugate of the original one. Equation \eqref{cdp:eq:C.EW.20}
is precisely the conjugate central action under
the anti-linear comparison $v\mapsto J\overline v$,
with the conjugate base rotation.


\item  The resulting proper compressed filling at infinity
is therefore the complex conjugate root plumbing
with its actual finite action. Its punctured-family
map has exactly \eqref{cdp:eq:C.EW.19} and \eqref{cdp:eq:C.EW.20}.
Its quotient rings are the conjugates of
\eqref{cdp:eq:C.EW.17}, so the quotient is smooth and its full
central fibre has the conjugate support with the
same multiplicities.


\item  Finally glue the two finite-equivariant proper
fillings to the entire period family \eqref{cdp:eq:C.EW.10},
using their explicit logarithmic biholomorphisms
and the frame change \eqref{cdp:eq:C.EW.18}. The closed-base
refinement argument of \ref{cdp:st:C.2.UN.55} proves this normal
root total space Hausdorff and proper over
$\PP^1_s$. Take the actual finite quotient by
the constructed action. The base quotient is
$\PP^1_t$ with $t=s^N$. Its total space is smooth
by \eqref{cdp:eq:C.EW.17} at the two special values and by
freeness of the base action elsewhere. It is
compact and proper. A general quotient fibre is
identified with any one of its $N$ root fibres,
so is the specified connected two-torus.
Every component and every double curve of either
special support has Euler number zero. Additivity
therefore proves the final Euler assertion.


\end{enumerate}

\item \par\medskip\noindent\textbf{Explicit nontrivial finite linear actions.}
For $N=2,3,4,6$, choose an elliptic curve $E$ with
an origin-preserving automorphism $\rho$ of exact
order $N$. Start with a root cycle of $2N$ copies
of $E\times\PP^1$, glued by the identity on their
elliptic sections. In the coordinates above take
\[
 \begin{gathered}
 \Lambda=\Lambda_E\oplus\Z p,\quad
 n_j=-p,\quad h_j=1,\quad
 b_j=ip/(2N),\quad \tau_j=j p/N,\\
 \sigma(j)=j+2\pmod{2N},\qquad
 R=\operatorname{diag}(\rho,1).
 \end{gathered}
\tag{C.EW.21}\hypertarget{tag-00JU}{}\label{cdp:eq:C.EW.21}
\]
The first equation \eqref{cdp:eq:C.EW.3} is zero except at the
cyclic endpoint, where it is $-2p$.
The $N$th translation sum is $(N-1)p$ on even
vertices and $Np$ on odd vertices. These are
actual integral periods, so this is the specified
finite action. It is free on the central graph.
The quotient has two components $E\times\PP^1$,
both with full multiplicity $N$. One elliptic
gluing is the identity; the other is $\rho^{-1}$,
as follows by identifying vertex $2$ with vertex
$0$ using the inverse generator. The root cycle
has identity product, while this latter gluing
retains the nontrivial finite linear part.
The compact construction above realizes these
exact weighted supports and their conjugates.


\begin{enumerate}[label={$\langle2\rangle$\arabic*.},leftmargin=1.3em,labelsep=.5em]
\item  The equivariance, target-frame affine equations,
infinity identities, positive determinant and
all four finite-action examples are checked
exactly in \path{code/D_all_weighted_cycles.py}
and its saved output.



\end{enumerate}

\end{enumerate}
\end{proof}






\begin{proposition}[Exact pruning and compact reinsertion classify every reduced all-Hopf support]\hypertarget{tag-01E6}{}\label{cdp:st:C.2.UN.60}
\StatementRecord{cdp:st:C.2.UN.60}{proved}{01E6}{\SeriesRef{cdp:st:C.2.UN.46}, \SeriesRef{cdp:st:C.2.UN.50}, \SeriesRef{cdp:eq:C.PR.2}, \SeriesRef{cdp:eq:C.PR.4}, \SeriesRef{cdp:st:C.2.UN.43}, \SeriesRef{cdp:eq:C.PR.6}, \SeriesRef{cdp:eq:C.PR.3}, \SeriesRef{cdp:st:C.2.UN.49}, \SeriesRef{cdp:st:C.2.UN.58}, \SeriesRef{cdp:eq:C.PR.5}, \SeriesRef{cdp:st:C.2.UN.57}, \SeriesRef{cdp:st:C.2.UN.54}, \SeriesRef{cdp:eq:C.PR.9}, \SeriesRef{cdp:st:C.2.UN.48}, \SeriesRef{cdp:eq:C.PR.10}, \SeriesRef{cdp:eq:C.PR.12}, \SeriesRef{cdp:eq:C.PR.13}, \SeriesRef{cdp:eq:C.PR.14}, \SeriesRef{cdp:eq:C.PR.16}, \SeriesRef{cdp:st:C.2.UN.42}}

Let $S$ be a specified compact connected reduced SNC surface whose
smooth components are primary or secondary Hopf surfaces. Its
elliptic gluing maps are part of the specified surface. The following
finite geometric test is necessary and sufficient for $S$ to occur
as the full central fibre of a proper degeneration of connected
complex two-tori with smooth total space.

\begin{enumerate}
\item The actual normal lines on every double curve are dual, and
the incidence graph has exactly one independent cycle.
\item If a diagonal algebraic-dimension-zero component occurs, the
whole support is its diagonal cycle and its integral lattice product
$M$ is the identity, as in \ref{cdp:st:C.2.UN.46}.
\item In the remaining case every algebraic-dimension-zero vertex
is a nonlinear leaf. Every other vertex has its actual elliptic
reduction to $\PP^1$, and every multiple fibre of that reduction
is one of its marked double curves.
\item Let $m_e$ be the exact order of either normal line at an edge,
and let $h_v$ be zero on the cycle and the sum of the $m_e$ along
the path from the cycle to $v$. There is an integer $R>0$ such that
every nonlinear leaf has $h_v=R$ and every other vertex has $h_v<R$.
If there is no nonlinear leaf, any $R>\max h_v$ passes this test.
\item Perform the exact leaf-pruning operations specified below,
including the full integer secondary framing and all torsor maps.
The resulting marked Hopf/ruled core cycle has integral
complement-lattice product $M=I$.
\end{enumerate}

In the second case the tests referring to elliptic tree vertices are
unneeded. In the remaining case the pruning operations always exist
under the preceding geometric requirements. Each terminal step is
one of two explicit operations: the inverse logarithmic chart of
\ref{cdp:st:C.2.UN.50} for a primary or secondary Hopf pair; or removal of the
product pair $(F\times\PP^1,F\times\{\infty\})$. The latter leaves
the parent surface unchanged. After its children have been removed,
every vertex has exactly one of these two forms. The pruned core
and its test are independent, up to the stated actual marked
isomorphisms, of the order of pruning disjoint branches.

Every support passing the test has a realization over $\PP^1$
with smooth compact total threefold. It retains all the original
components, their analytic torsors and every specified elliptic
gluing map. Its other singular support can be taken to be the
complex conjugate of its pruned core cycle. All reduced fibre
supports have Euler characteristic zero.
\end{proposition}



\begin{proof}\leavevmode\par



 
\begin{enumerate}[label={$\langle1\rangle$\arabic*.},leftmargin=0pt,labelsep=.5em]
\item \par\medskip\noindent\textbf{The exact one-point torsor and its marked compact surface.}
We first prove the classification used at a terminal pruning step.
Let $q:T\to\PP^1$ be a fixed-Jacobian, identity-monodromy elliptic
quasi-bundle, with just one possible multiple fibre $mJ$, marked
over infinity. Its actual normal-root charts make it a principal
$F$-torsor on
\[
                 \mathcal B=\PP^1[(\infty,m)].          \tag{C.PR.1}\hypertarget{tag-00VH}{}\label{cdp:eq:C.PR.1}
\]
Choose an oriented period basis $\lambda_1,\lambda_2$ of
$F=\C/\Lambda_F$. The exact sheaf sequence used here is
\[
 \begin{gathered}
 0\longrightarrow\OO_{\mathcal B}
 \xrightarrow{\ a\mapsto
 (e^{2\pi i\lambda_2a},e^{-2\pi i\lambda_1a})\ }
 (\OO_{\mathcal B}^{*})^2
 \xrightarrow{\ \Phi\ }F_{\mathcal B}\longrightarrow0,\\
 \Phi(u,v)=
 \frac{\lambda_1\log u+\lambda_2\log v}{2\pi i}
                          \pmod{\Lambda_F}.            \tag{C.PR.2}\hypertarget{tag-00VR}{}\label{cdp:eq:C.PR.2}
 \end{gathered}
\]
Changing a logarithm changes the numerator by a lattice period.
Conversely a pair in the kernel can be lifted to logarithms,
adjusted by integral periods, and then lies on the displayed
complex line. The intersection of that line with the integral
two-period lattice is zero because
$\lambda_2/\lambda_1$ is not real. Thus the additive kernel and
every map in \eqref{cdp:eq:C.PR.2} are exact.


\begin{enumerate}[label={$\langle2\rangle$\arabic*.},leftmargin=1.3em,labelsep=.5em]
\item  Finite stabilizer invariants are exact in characteristic zero, so
$H^j(\mathcal B,\OO)=H^j(\PP^1,\OO)=0$ for $j=1,2$.
Moreover
\[
 \operatorname{Pic}(\mathcal B)=\Z[\mathcal R],
                 \qquad \mathcal R^m=\OO_{\PP^1}(\infty).
                                                        \tag{C.PR.3}\hypertarget{tag-00VS}{}\label{cdp:eq:C.PR.3}
\]
Indeed every line has a unique stabilizer weight modulo $m$;
subtract that power of the canonical root line and it descends
to $\PP^1$. The displayed relation then gives the asserted
infinite cyclic group, with its canonical sections retained.
The long exact sequence of \eqref{cdp:eq:C.PR.2} therefore identifies the
actual $F$-torsor class with a unique vector in
$\operatorname{Pic}(\mathcal B)^2=\Z^2$.


\item  If this vector is zero, the torsor is the product
$F\times\mathcal B$. The inertia must act freely on its total
space because the actual connected root cover of $mJ$ is
unramified over the original surface tube. A nontrivial inertia
would act trivially on its product fibre. Thus $m=1$ and
$T=F\times\PP^1$, with $J=F\times\{\infty\}$.


\item  For a nonzero vector, choose its primitive direction and an
oriented period basis so that it is $(-n,0)$, $n>0$. This is
an integral change of the actual fibre basis, not removal of
its direction. The torsor is the pushout of
$(\mathcal R^{-n},\OO)$ by \eqref{cdp:eq:C.PR.2}. Gauge the second nonzero
line coordinate to one. The remaining ambiguity is multiplication
of the first coordinate by
$\gamma^{\Z}$, where
$\gamma=\exp(2\pi i\lambda_2/\lambda_1)$ and $|\gamma|<1$.
Consequently the actual total surface is
\[
                         (\mathcal R^{-n})^\times/\gamma^{\Z}.
                                                        \tag{C.PR.4}\hypertarget{tag-00VT}{}\label{cdp:eq:C.PR.4}
\]
Write $\mathcal B=\PP(m,1)$ with coordinates $(Z,W)$ of
weights $(m,1)$ and coarse coordinate $U=Z/W^m$.
The nonzero total line of $\mathcal R^{-n}$ is
\[
 [((\C^2\setminus0)\times\C^*_\ell)/
              \lambda\cdot(Z,W,\ell)
                =(\lambda^mZ,\lambda W,\lambda^{-n}\ell)].
                                                        \tag{C.PR.5}\hypertarget{tag-00VU}{}\label{cdp:eq:C.PR.5}
\]
Setting $\ell=1$ leaves the finite group $\mu_n$ with weights
$(m,1)$. A scalar multiplication by $\gamma$ lifts, after choosing
$\beta^n=\gamma$, to $(Z,W)\mapsto(\beta^mZ,\beta W)$.
Thus \eqref{cdp:eq:C.PR.4} is exactly
$H_m^{\rm diag}(\beta)/\mu_n$, with the marked axis $W=0$.
Freeness of the original inertia is exactly
$\gcd(m,n)=1$. Changing the chosen root of $\gamma$ changes
the infinite generator by the already included finite group.
Every primitive fibre basis, the full integer $n$, and the marked
axis identification have been retained.


\item  This also explains why reduction of $n$ modulo $m$ is impermissible:
the entire root-line degree in \eqref{cdp:eq:C.PR.4}, rather than only its
stabilizer character, is the torsor class.


\end{enumerate}

\item \par\medskip\noindent\textbf{The actual inverse operation at a Hopf leaf.}
For a nonlinear leaf use its full normal form from \ref{cdp:st:C.2.UN.43};
for an elliptic terminal surface use \eqref{cdp:eq:C.PR.4}. Both have the
same marked normal-line data. Put $r=m$, choose $k,l$ with
$mk-nl=1$, and write $\gamma=\beta^n$. On the parent the
actual coarse root chart and its inverse product chart are
\[
 \begin{gathered}
 (u,X)\sim(\gamma^m u,\gamma^lX),\qquad p=X^m u^{-l},\\
 W=X^{-n}u^k,\qquad
                u=p^nW^m,\quad X=p^kW^l.
 \end{gathered}                                        \tag{C.PR.6}\hypertarget{tag-00VV}{}\label{cdp:eq:C.PR.6}
\]
These are precisely (C.SN.3)--(C.SN.5), with the given dual-normal
map and elliptic gluing used to choose the same connected root
cover. Replace this marked multiple-fibre tube by
$(p,W)\in\Delta\times(\C^*/\gamma^{\Z})$ through the displayed
isomorphism on $p\ne0$. This is the actual inverse logarithmic
transform. Its base point becomes an ordinary elliptic fibre.
Every other chart and transition of the parent is unchanged.
In particular the parent remains the same fixed-Jacobian
identity-monodromy quasi-bundle off this one point, with its full
analytic torsor retained by the displayed transition map.


\begin{enumerate}[label={$\langle2\rangle$\arabic*.},leftmargin=1.3em,labelsep=.5em]
\item  The choices do not discard gluing data. A change of a root
trivialization or a normal-product constant changes \eqref{cdp:eq:C.PR.6} by
a translation gauge extending over the inserted product disc.
The primitive nonzero direction in \eqref{cdp:eq:C.PR.3} fixes the angular
direction; the other period can change by an integral multiple,
which changes the generator by its same finite deck group.
Translations of the marked axis are retained in the transition.
This is the exact choice comparison already proved for the
full inverse torsor in \ref{cdp:st:C.2.UN.49} and \ref{cdp:st:C.2.UN.50}.


\item  Remove terminal vertices inward from every attached tree.
All original nonterminal vertices have the elliptic reduction
of \ref{cdp:st:C.2.UN.58}. Removing a child removes exactly its marked
multiple fibre by \eqref{cdp:eq:C.PR.6}, or leaves an ordinary fibre unchanged
in the product-cap case. Hence, after all children of a vertex
are removed, its only possible multiple fibre is its parent
mark. The one-point classification above applies to that
actual modified surface. This proves that every terminal step
has one of the two asserted forms.


\item  Operations in distinct discs of one base keep the same transition
atlas on their complement and hence commute. Operations in disjoint
subtrees do likewise. Every allowable pruning order differs by
interchanging such operations; a parent is processed only after
its children. This proves independence through actual marked
torsor isomorphisms, not merely through the same degree vector.
Keep the entire maps \eqref{cdp:eq:C.PR.6}, the primitive periods, the integer
$n$ and all remaining marked points as the record of each step.


\end{enumerate}

\item \par\medskip\noindent\textbf{The exact two-point core models.}
At the end each cycle vertex is still a connected fixed-$F$
quasi-bundle over $\PP^1$, with only its two marked possible
multiple fibres, of orders $a,b$. Put $g=\gcd(a,b)$,
$A=a/g$, $B=b/g$. There is the actual representable finite
etale orbicurve map
\[
 \PP^1[(0,A),(\infty,B)]
       \longrightarrow \PP^1[(0,a),(\infty,b)],
                         \qquad t=z^g .                \tag{C.PR.7}\hypertarget{tag-00VW}{}\label{cdp:eq:C.PR.7}
\]
At zero its root uniformizers have map $v=u$ and inertia
$\mu_A\subset\mu_a$; the infinity chart has the analogous
inclusion. Thus it is etale as a map of these root curves.
The source, since $A,B$ are coprime, is the weighted projective
line $\PP(B,A)$, with coarse coordinate
$z=X^A/Y^B$. Its Picard group is generated by $\OO(1)$:
the two root divisors have classes $B$ and $A$, and Bezout
expresses their generator. Its higher coherent groups vanish.


\begin{enumerate}[label={$\langle2\rangle$\arabic*.},leftmargin=1.3em,labelsep=.5em]
\item  Pull back the actual $F$-torsor. The same sequence \eqref{cdp:eq:C.PR.2}
therefore describes it by an integral vector. A nonzero vector
$(-n,0)$ gives the explicit nonzero line $\OO(-n)$, and the
gauge \eqref{cdp:eq:C.PR.5}, now of weights $(B,A)$, gives the actual diagonal
Hopf cover with contraction $(\beta^BX,\beta^AY)$ and finite
action $\mu_n(B,A)$. The inertia is free precisely when
$\gcd(n,A)=\gcd(n,B)=1$. Its cyclic descent preserves both
marked axes. Lifts to a central primary cover commute with a
contraction and preserve each axis, so they have the form
$(Xf(X,Y),Yh(X,Y))$; contraction equivariance forces the
nonvanishing units $f,h$ to be constant. Thus the full quotient
is the actual marked diagonal primary or secondary Hopf surface.


\item  For the zero vector, free inertia forces $A=B=1$ and the cover
is $F\times\PP^1$. Its cyclic descent is
$(y,z)\mapsto(y+\tau,\zeta_g z)$: a holomorphic map
$\PP^1\to F$ is constant, and the principal-$F$ action excludes
a linear part. Freeness at $0,\infty$ makes $\tau$ have exact
order $g$. The quotient is exactly
\[
        \PP(\OO_E\oplus L),\quad
        E=F/\langle\tau\rangle,\quad
        L=(F\times\C)/((y,v)\sim(y+\tau,\zeta_gv)).
                                                        \tag{C.PR.8}\hypertarget{tag-00VX}{}\label{cdp:eq:C.PR.8}
\]
This exhibits its two marked sections, the exact torsion order
of $L$, the full fibre map $q=z^g$ on the cover, and both
elliptic identifications. It also gives the actual finite
projective-bundle map from $F\times\PP^1$.
Thus the pruned core is precisely one of the marked diagonal
Hopf/ruled cycles of \ref{cdp:st:C.2.UN.57}, with all its data retained.


\end{enumerate}

\item \par\medskip\noindent\textbf{Necessity of the core lattice test through actual group maps.}
For an actual torus degeneration, the graph, normal and height
requirements are \ref{cdp:st:C.2.UN.54} and \ref{cdp:st:C.2.UN.58}. It remains to compare
the nearby group to that of the pruned core; no contraction of
the given holomorphic total space is required.


\begin{enumerate}[label={$\langle2\rangle$\arabic*.},leftmargin=1.3em,labelsep=.5em]
\item  The complement of the marked axis in a terminal diagonal Hopf
pair is exactly $\C\times F$. For a nonlinear pair the explicit
proper smooth-pair deformation of \ref{cdp:st:C.2.UN.49} identifies its compact
core and boundary with the diagonal pair. The product ruled cap
has the same complement $F\times\C$. In the primitive boundary
basis of \ref{cdp:st:C.2.UN.50} the actual homomorphism of this compact
piece is
\[
        \Z^3_{\delta,e,h}\longrightarrow\Z^2_{e,h},
                  \delta\longmapsto0,\quad
                  e\longmapsto e,\quad h\longmapsto h. \tag{C.PR.9}\hypertarget{tag-00VY}{}\label{cdp:eq:C.PR.9}
\]
For the Hopf case the exact matrix
$[(-l,k,0),(0,0,1),(m,-n,0)]$ has determinant one;
the normal reversal in $xy=t$ turns the killed Hopf meridian
into precisely $\delta$. The inverse product filling
$F\times D^2$ has the identical homomorphism \eqref{cdp:eq:C.PR.9},
with the same two retained generators. For a product cap this
is the literal product-disc boundary map.


\item  Van Kampen applied to the unchanged exterior and this boundary
gives two pushout diagrams with identical maps. Their resulting
groups are canonically isomorphic by the identity on the exterior
and the two displayed fibre generators. Repeat the comparison,
retaining every remaining collar. The actual tubular comparison
of \ref{cdp:st:C.2.UN.49} carries these maps to an arbitrary original smoothing.
Consequently the nearby group of the core plumbing is the
original $\Z^4$. The exact core calculation of \ref{cdp:st:C.2.UN.57} gives
$\Lambda\rtimes_M\Z$, with $\Lambda$ injected. Its abelianness
forces $M=I$. This proves necessity using the actual maps needed
for the test.


\end{enumerate}

\item \par\medskip\noindent\textbf{The compact core family with a prescribed remaining period order.}
For sufficiency first handle the all-algebraic-dimension-zero
diagonal branch by \ref{cdp:st:C.2.UN.48}. Otherwise prune as above and assume
the core test passes. Its components have retained elliptic
quasi-bundles. In the common complement
$G=\C^2/\Lambda$ of \ref{cdp:st:C.2.UN.57} their regular fibres are translates
of the compact elliptic subgroup
\[
       F=(W\cap iW)/(\Lambda\cap W\cap iW),
                     \qquad W=\operatorname{span}_{\R}\Lambda.
                                                        \tag{C.PR.10}\hypertarget{tag-00VI}{}\label{cdp:eq:C.PR.10}
\]
Their lattice is primitive of rank two: their actual quotient
in the open two-boundary component is $\C^*$.
The root-cover maps already recorded identify every later
inserted regular fibre with this same elliptic group.


\begin{enumerate}[label={$\langle2\rangle$\arabic*.},leftmargin=1.3em,labelsep=.5em]
\item  In \ref{cdp:st:C.2.UN.57} replace $V_z(w)=w+zJ\overline w$ by
$V_{t^R}(w)=w+t^RJ\overline w$, retaining the term
$d\log t$ and the normal neck equation $xy=t$.
The exact real determinant becomes
\[
 C(1+|t|^{2R})
 \left(1+|t|^{2R}
       +\frac{d}{2\pi}(|t|^{2R}-1)\log|t|\right)>0.
                                                        \tag{C.PR.11}\hypertarget{tag-00VJ}{}\label{cdp:eq:C.PR.11}
\]
The product $(|t|^{2R}-1)\log|t|$ is nonnegative.
Thus the same actual period parallelepipeds and both endpoint
plumbings give a smooth compact torus family on $\PP^1$ with
the prescribed core at zero and its conjugate at infinity.
At infinity divide the frame by $t^R$; the displayed frame is
$J\overline w+s^Rw$, and the logarithmic term still has exactly
the sign and coefficient for the reduced neck $xy=s$.


\item  Let $\lambda,\mu=c\lambda$ be an oriented primitive period basis
of the compact group in \eqref{cdp:eq:C.PR.10}, with $\operatorname{Im}c>0$.
The ordinary determinant in the actual vector frame is
\[
 \det_{\C}(V_{t^R}\lambda,V_{t^R}\mu)
       =2i\operatorname{Im}(c)\,\|\lambda\|^2 t^R\ne0
                      \quad(t\ne0).                    \tag{C.PR.12}\hypertarget{tag-00VK}{}\label{cdp:eq:C.PR.12}
\]
Here $\det(\lambda,J\overline\lambda)=-\|\lambda\|^2$.
This is the exact common core period, including its integral
class, rather than a chosen formal order.


\end{enumerate}

\item \par\medskip\noindent\textbf{The recursive local model and the exact order drops.}
The actual logarithmic charts near a regular compact fibre give
a neighbourhood in the following rank-two period quotient:
\label{layout:II:3}
\[
\begin{gathered}
\mathcal G=
 (\C_p\times\C_w^*\times\Delta_t)/
 \langle(p,w,t)\mapsto(p+\delta(t),\gamma(t)w,t)\rangle,
 \quad
 \\
\delta(t)=t^\ell d(t),\quad d(0)\ne0,\quad|\gamma(0)|<1 .
\end{gathered}
\tag{C.PR.13}\hypertarget{tag-00VL}{}\label{cdp:eq:C.PR.13}
\]
Indeed choose the primitive first period as $(0,2\pi i)$ in
holomorphic linear coordinates; the second is then
$(\delta(t),\log\gamma(t))$. Its first coordinate is a unit
times their determinant. The charts of the core plumbing select
a bounded transverse disc in this cover and identify it with
the actual neighbourhood of the chosen central elliptic fibre.
Initially \eqref{cdp:eq:C.PR.12} gives $\ell=R$.


\begin{enumerate}[label={$\langle2\rangle$\arabic*.},leftmargin=1.3em,labelsep=.5em]
\item  Every prescribed new primitive fibre basis is allowed. If
$\lambda'=a\lambda+b\mu$, $\mu'=c\lambda+d\mu$ with
$ad-bc=1$, the coordinate
$p'=p-b\delta V/(2\pi ia+b\log\gamma)$ and the rescaled
logarithmic coordinate
$V'=2\pi iV/(2\pi ia+b\log\gamma)$ make $\lambda'=(0,2\pi i)$.
The second displacement is
$2\pi i\,\delta/(2\pi ia+b\log\gamma)$.
The denominator is a holomorphic unit near zero.
Thus its order is unchanged and the full new fibre basis is
retained by an actual linear period map.


\item  For a recorded Hopf step of weight $m$ and secondary integer $n$,
perform the weight-$(m,1)$ blowup of the base coordinates $(p,t)$
and tensor its nonzero line by $\OO(nE)$. In the free coarse
chart of \ref{cdp:st:C.2.UN.50} write $p=x^mZ$, $t=xW$.
A chosen holomorphic root $\beta(t)^n=\gamma(t)$ gives the exact
infinite map
\[
 (Z,W,x)\longmapsto
 \left(\beta(t)^m
    \bigl(Z+d(t)x^{\ell-m}W^\ell\bigr),
       \beta(t)W,\beta(t)^{-1}x\right).
                                                        \tag{C.PR.14}\hypertarget{tag-00VM}{}\label{cdp:eq:C.PR.14}
\]
The finite action has weights $(m,1,-1)$ modulo $n$ and is free,
because $\gcd(m,n)=1$. The added monomial has weight $m$,
so these maps commute. The inequality $\ell\ge m$ follows
from the height test. At $\ell=m$ the exceptional component
is exactly the recorded nonlinear Hopf surface after a constant
rescaling of $Z$. At $\ell>m$ it is exactly the recorded
diagonal pair. On the new child's complement of its parent curve
put $U=Z/W^m$. The actual deck map becomes
\[
        (U,W^n)\longmapsto
       \left(U+\frac{\delta(t)}{t^m},\,\gamma(t)W^n\right).
                                                        \tag{C.PR.15}\hypertarget{tag-00VN}{}\label{cdp:eq:C.PR.15}
\]
It has remaining displacement order $\ell-m$.
At any other marked point use $U-U_0$ as its base coordinate.
This retains the original marked point and every later child
transition. No extension of an arbitrary original smoothing's
elliptic quotient is being asserted.


\item  For a recorded product cap, $m=1$ and $\ell\ge2$, because this
vertex was an elliptic vertex of height strictly smaller than $R$.
Blow up the actual ideal $(p,t)$ along its compact elliptic
curve in \eqref{cdp:eq:C.PR.13}. The deck map preserves this ideal.
On its exceptional normal directions it is the identity,
since $\delta'(0)=0$; hence the exceptional is precisely
$F\times\PP^1$ with its parent section. The chart $p=tU$
has the exact map
\[
                    (U,w)\longmapsto
          (U+\delta(t)/t,\gamma(t)w).                   \tag{C.PR.16}\hypertarget{tag-00VO}{}\label{cdp:eq:C.PR.16}
\]
Thus this step has exactly the same order drop one. The strict
parent surface is unchanged, since its centre is a Cartier
curve on that surface. The two other blowup coordinates give
the reduced normal-crossing equation $t=xy$.


\item  Inductively, the model at a vertex of height $h$ has remaining
order $R-h$. All vertices with descendants are elliptic and
have positive remaining order. Only a nonlinear terminal step
has order zero after insertion. The height test is therefore
precisely what makes every recorded step in \eqref{cdp:eq:C.PR.14}--\eqref{cdp:eq:C.PR.16}
holomorphic with the required central type.


\end{enumerate}

\item \par\medskip\noindent\textbf{Local properness and simultaneous compact replacement.}
We give the needed bounds also for the affine, noncompact base
of \eqref{cdp:eq:C.PR.13}. On its weight-$(m,1)$ blowup use the actual
canonical-section norm
\[
                  \rho(p,t)=
                  (|p|^2+|t|^{2m})^{1/(2m)}.           \tag{C.PR.17}\hypertarget{tag-00VP}{}\label{cdp:eq:C.PR.17}
\]
In the chart $p=x^mZ$, $t=xW$ this is
$|x|(|Z|^2+|W|^{2m})^{1/(2m)}$, the nonvanishing smooth
metric factor on the weighted exceptional charts. Translation
by $j\delta(t)$ gives two-sided norm distortion at most
$C(1+|j|)^{1/m}$, because $|\delta(t)|\le C|t|^m$.
For $\OO(nE)$ the exponent is $n/m$. On a sufficiently small
closed disc the scalar $|\gamma(t)|$ is uniformly smaller
than one. Its exponential contraction dominates this polynomial.
The infinite action is free and properly discontinuous, and
a high fixed iterate gives a norm-annulus representative for
every orbit, exactly as in \ref{cdp:st:C.2.UN.42}.


\begin{enumerate}[label={$\langle2\rangle$\arabic*.},leftmargin=1.3em,labelsep=.5em]
\item  The old line coordinate is the new line divided by the $n$th
power of the canonical exceptional section. If an old image
stays in a compact old annular chart, putting both line norms
in annuli changes the base by at most a constant times
$|\delta(t)|(1+|\log\rho|)$. For $\rho<1$ use $\rho\ge|t|$;
this quantity is bounded and tends to zero as $t\to0$.
For $\rho\ge1$, a new base representative consequently obeys
\[
 |p_{\rm new}|\le B+C|\delta(t)|\log(2+|p_{\rm new}|).
                                                        \tag{C.PR.18}\hypertarget{tag-00VQ}{}\label{cdp:eq:C.PR.18}
\]
It is uniformly bounded: a linear function eventually dominates
the logarithm. Thus the new annulus over the required base set
is compact, including the compact exceptional Hopf component.
This proves the local compactness comparison without assuming
that the affine base itself is compact.


\item  Off the new Hopf component the canonical exceptional section
gives the actual isomorphism with the old model minus its
marked central elliptic curve. In a compact Hopf fundamental
domain, $x\to0$ gives $p=x^mZ$, $t=xW$ and old line coordinate
proportional to $x^{-n}$. Normalizing that coordinate takes
$O(1+|\log|x||)$ iterates. Every old base cluster point is zero,
because
$|x|^m+|t|^\ell(1+|\log|x||)\to0$.
The normal-annulus coordinate has a compact limit on the old
elliptic curve. Hence every old cluster point lies on that
curve and on no other point of the old model.


\item  Choose disjoint small neighbourhoods of the finitely many
curves at each insertion stage. The new open piece is the
inserted compact component together with the inverse image
of the punctured old neighbourhood. The preceding cluster
comparison proves it is open and that all gluing graphs are
relatively closed. The resulting complex space is Hausdorff.
The bounded-annulus comparison over smaller closed
neighbourhoods and the unchanged compact exterior prove
properness. Product-cap steps are ordinary proper blowups.
Finite iteration therefore gives a smooth proper family over
a disc, glued by actual biholomorphisms to the unchanged
compact core family away from zero.


\end{enumerate}

\item \par\medskip\noindent\textbf{Restoration of the prescribed support.}
On the strict parent of every Hopf insertion the chart is
exactly \eqref{cdp:eq:C.PR.6}, in the opposite direction: the old coordinate
is $W=x^{-n}$ and the full secondary integer is unchanged.
It restores the complete recorded elliptic torsor, not just
its rational Chern class. All untouched charts and all
other marked fibres remain their recorded charts. On a cap
step the actual product exceptional pair and its parent
section are the recorded ones. Fibre-basis maps retain
the linear parts of every elliptic gluing. The constant
translations are retained in the recorded transition maps;
equivalently the automorphisms
$(Z,W)\mapsto(a^mZ,aW)$ of either Hopf contraction realize
every required marked-axis translation. They commute with
the finite action. Product caps have the corresponding
elliptic translations and projective automorphisms fixing
the parent section.


\begin{enumerate}[label={$\langle2\rangle$\arabic*.},leftmargin=1.3em,labelsep=.5em]
\item  Reverse the complete pruning record. At each step these
actual maps recover exactly the surface and marked gluing
removed at that step. Induction therefore restores the
specified $S$ in its full analytic isomorphism class.
The normal equations are $t=xW$ in the free weighted
coarse charts and $t=xy$ for an ordinary blowup, so every
component has full multiplicity one. The other fibres
are unchanged from the compact period family, and the
filling at infinity is its conjugate pruned core.
All these reduced supports have Euler zero: their
components are Hopf or elliptic ruled surfaces, their
double curves are elliptic, and there are no triple points.
This proves sufficiency and the compact assertion.


\item  The exact comparisons are reproduced in
\path{code/E_hopf_tree_pruning.py} and its saved output:
primitive integral fibre-basis changes, both-point root Picard
lattices, the complete weighted deck and finite-action maps,
the nested remaining-order identities, and the product-cap
blowup. These checks retain the full integers in the pruning
record; the proof establishes them for arbitrary finite trees.



\end{enumerate}

\end{enumerate}
\end{proof}






\begin{proposition}[The exact rooted elliptic-tree criterion and compact occurrence]\hypertarget{tag-01E7}{}\label{cdp:st:C.2.UN.61}
\StatementRecord{cdp:st:C.2.UN.61}{proved}{01E7}{\SeriesRef{cdp:st:C.2.UN.60}, \SeriesRef{cdp:st:C.2.UN.5}, \SeriesRef{cdp:st:C.2.UN.54}, \SeriesRef{cdp:eq:C.PT.3}, \SeriesRef{cdp:st:C.2.UN.28}, \SeriesRef{cdp:eq:C.PT.5}, \SeriesRef{cdp:eq:C.PT.6}, \SeriesRef{cdp:st:C.2.UN.31}, \SeriesRef{cdp:eq:C.PT.7}, \SeriesRef{cdp:eq:C.PT.8}, \SeriesRef{cdp:st:C.2.UN.58}, \SeriesRef{cdp:eq:C.PT.9}, \SeriesRef{cdp:st:C.2.UN.50}, \SeriesRef{cdp:eq:C.PT.11}}

Let $S$ be a specified compact connected reduced SNC surface whose
incidence graph is a nontrivial tree. Distinguish one smooth component
$P$, with no initial restriction on its geometric genus or surface
type. Every other component is either a primary or secondary Hopf
surface, or is equipped with a specified fixed-Jacobian elliptic
quasi-bundle over $\PP^1$, with identity integral monodromy,
connected fibres and only smooth multiple fibre supports.
In the latter case every boundary curve is a fibre support of
the specified map. Retain the actual elliptic gluing maps
and the complete analytic torsors. The following finite
geometric test is necessary and sufficient for $S$ to occur
as a full fibre of a smooth-total degeneration of connected complex
two-tori.

\begin{enumerate}
\item The actual normal lines at every double curve are dual.
\item The root $P$ has a fixed-Jacobian elliptic quasi-bundle
$q_P:P\to B$ over a smooth elliptic curve, with identity integral
monodromy and only smooth multiple supports. Its marked boundary
curves are fibre supports, and every multiple fibre is marked.
Every other component is either a nonlinear algebraic-dimension-zero
Hopf leaf or an elliptic quasi-bundle over $\PP^1$ of the stated
type, whose multiple fibres are all marked double curves.
\item For each edge $e$, let $m_e$ be the exact finite order of
its normal line. It is the full elliptic multiplicity at each
elliptic endpoint and the resonance at a nonlinear endpoint.
Set $h_P=0$ and
\[
                  h_v=\sum_{e\in[P,v]}m_e.
                  \tag{C.PT.1}\hypertarget{tag-00VZ}{}\label{cdp:eq:C.PT.1}
\]
There is an integer $R>0$ such that every nonlinear leaf has
$h_v=R$ and every elliptic nonroot vertex has $h_v<R$.
In the absence of nonlinear leaves any $R>\max_vh_v$ passes.
\item Perform the exact leaf-pruning operations of \ref{cdp:st:C.2.UN.60},
retaining the complete inverse logarithmic charts, primitive fibre
directions, secondary integers and torsor transitions. The final
surface is a principal elliptic torsor $T_0\to B$. Its integral
Chern vector in $H^2(B,\Lambda_F)$ is zero.
\end{enumerate}

The normal-root maps identify the one fixed Jacobian $F$ along
the tree. The last condition refers to the actual analytic
inverse torsor; it does not replace it by its rational Chern
data. The pruning result is independent, up to the actual
marked isomorphisms of \ref{cdp:st:C.2.UN.60}, of the order of disjoint operations.

Every support passing the test has a realization over $\PP^1$
with smooth compact total threefold. Its only singular fibre
is the given reduced $S$; every other fibre is a smooth
two-torus. All reduced fibre supports have Euler number zero.
Every component and every specified elliptic gluing map is
retained, including the complete torsor of every original
properly elliptic nonroot component.

In every actual occurrence the root is minimal,
$h^q(P,\OO_P)=h^q(S,\OO_S)=(1,2,1)$, and its actual canonical
divisor is
\[
                 K_P=\sum_{e\ni P}(m_e-1)J_e.
                 \tag{C.PT.2}\hypertarget{tag-00W5}{}\label{cdp:eq:C.PT.2}
\]
It is the unique minimum-order component for canonical evaluation;
the differences of the other orders are exactly \eqref{cdp:eq:C.PT.1}.
The zero-edge case, excluded from the nontrivial tree here,
is the smooth torus fibre already treated in \ref{cdp:st:C.2.UN.5}.
\end{proposition}



\begin{proof}\leavevmode\par



 
\begin{enumerate}[label={$\langle1\rangle$\arabic*.},leftmargin=0pt,labelsep=.5em]
\item \par\medskip\noindent\textbf{The root cohomology is forced by the actual incidence maps.}
Every double curve has a nonroot endpoint. At a Hopf
endpoint it is smooth elliptic by the curve argument of
\ref{cdp:st:C.2.UN.54}; at a supplied quasi-bundle endpoint it is a smooth
elliptic fibre support by hypothesis. Distinct boundary
curves on either such endpoint are disjoint, so there
are no triple points. In particular the curves $J_e$
on $P$ are also pairwise disjoint. Write $T_v$ for a
nonroot component. The exact normalization sequence is
\[
 0\longrightarrow\OO_S
 \longrightarrow\OO_P\oplus\bigoplus_{v\ne P}\OO_{T_v}
 \longrightarrow\bigoplus_e\OO_{J_e}\longrightarrow0.
 \tag{C.PT.3}\hypertarget{tag-00W6}{}\label{cdp:eq:C.PT.3}
\]
For a Hopf component, $h^q(\OO_{T_v})=(1,1,0)$ and
every boundary restriction on $H^1$ is nonzero by \ref{cdp:st:C.2.UN.54}.
For a supplied quasi-bundle, (C.QT.2)--(C.QT.3) give
the same dimensions and the actual commuting square:
restriction to $J_e$ followed by its connected normal-root
pullback is evaluation of the fixed $H^1(F,\OO_F)$
frame. Thus its boundary restriction is also nonzero.
These are the actual restriction and gluing maps.
Root the tree at $P$ and order the nonroot columns from the
leaves inward. The submatrix of the $H^1$ incidence map
on those columns is square, with one column and row per
nonroot vertex. Elimination at a leaf has its nonzero
restriction as pivot. Repeating gives an invertible
matrix. The full $H^1$ incidence map is therefore onto,
regardless of the root restriction maps.


\begin{enumerate}[label={$\langle2\rangle$\arabic*.},leftmargin=1.3em,labelsep=.5em]
\item  The $H^0$ incidence map of a tree is also onto. Since
$H^2(\OO_{T_v})=0$ on every nonroot component, the actual long exact
sequence of \eqref{cdp:eq:C.PT.3} gives
\[
 H^2(S,\OO_S)\simeq H^2(P,\OO_P),\qquad
                 h^1(S,\OO_S)=h^1(P,\OO_P).
 \tag{C.PT.4}\hypertarget{tag-00W7}{}\label{cdp:eq:C.PT.4}
\]
The smooth total map is flat over the disc. Upper
semicontinuity gives $h^2(S,\OO_S)\ge1$, while the
compact-core bound \ref{cdp:st:C.2.UN.28} gives $p_g(P)\le1$. Thus
$p_g(P)=1$. Coherent Euler characteristic in this proper
flat family is zero. Additivity in \eqref{cdp:eq:C.PT.3}, since every
nonroot component and every elliptic curve has coherent Euler zero,
then gives $\chi(\OO_P)=0$ and $q(P)=2$. This proves
the asserted coherent dimensions without an initial
geometric-genus hypothesis on $P$.


\end{enumerate}

\item \par\medskip\noindent\textbf{The actual dualizing form and the unique minimum.}
Use the residue sequence (C.HR.7), or the identical
sequence of the actual boundary divisor, at every
nonroot vertex. At a terminal vertex its one nonzero restriction
coefficient forces the sole residue of a global
dualizing form to be zero. Since this component has no
holomorphic two-form, the form on that component is zero.
Remove that leaf and repeat. All nonroot forms and every
edge residue are zero. The remaining form on $P$ is
an ordinary holomorphic two-form. Conversely such a form,
with zero on all nonroot components, glues by the zero-residue
rule. These are inverse maps, giving the actual isomorphism
\[
                  H^0(S,\omega_S)\simeq H^0(P,K_P).
 \tag{C.PT.5}\hypertarget{tag-00W8}{}\label{cdp:eq:C.PT.5}
\]


\begin{enumerate}[label={$\langle2\rangle$\arabic*.},leftmargin=1.3em,labelsep=.5em]
\item  Take a canonical evaluation generator, with vertical orders
$k_v$, and divide it by $t^{\min k_v}$. Its nonzero restriction
to the reduced fibre spans \eqref{cdp:eq:C.PT.5}. The local adjunction map
(C.HR.8) identifies its nonzero components precisely with the
minimum-order components. Hence the minimum is attained
only at $P$. Writing $k_P=k_0$, the same actual residue
calculation on $P$ gives a holomorphic two-form with divisor
\[
          K_P=\sum_{e=(P,w)}(k_w-k_0-1)J_e,\qquad
                          k_w>k_0.
 \tag{C.PT.6}\hypertarget{tag-00W9}{}\label{cdp:eq:C.PT.6}
\]
All coefficients are nonnegative integers; there is no
additional zero away from the boundary, because the
evaluation has only its vertical component divisors.
If $E$ were a $(-1)$-curve, adjunction would give
$K_P\cdot E=-1$. The effective divisor in \eqref{cdp:eq:C.PT.6}
is supported on elliptic curves, none equal to $E$,
and has nonnegative intersection with $E$. This
contradiction proves that $P$ is minimal.


\end{enumerate}

\item \par\medskip\noindent\textbf{The actual elliptic fibration on the root.}
Its nonzero canonical form gives $\kappa(P)\ge0$.
Every component has full multiplicity one, and the
given total space and central support are already
smooth and SNC. Thus the actual common-root map in
\ref{cdp:st:C.2.UN.31} has degree one and its component over $P$ is
$P$ itself. That theorem makes $P$ a two-torus or
gives on it a fixed-modulus, identity-monodromy
elliptic quasi-bundle with only smooth multiples.
Its base has genus $p_g(P)=1$.


\begin{enumerate}[label={$\langle2\rangle$\arabic*.},leftmargin=1.3em,labelsep=.5em]
\item  We record an exact principal-torsor comparison used twice
below. For a principal $F=\C/\Lambda_F$-torsor over a
smooth elliptic curve $B$, the sequence
\[
 0\longrightarrow\Lambda_F
 \longrightarrow\OO_B\longrightarrow F_B
 \longrightarrow0
 \tag{C.PT.7}\hypertarget{tag-00WA}{}\label{cdp:eq:C.PT.7}
\]
maps its actual analytic class to its integral Chern
vector. If that vector is zero, the class comes from
$H^1(B,\OO_B)$. On $B=\C/\Lambda_B$, a representative
$a\,d\overline z$ has additive transition functions
$a\overline\gamma$ for $\gamma\in\Lambda_B$.
Its actual total torsor is therefore
\[
 \C^2/
 \left\langle
  (z,w)\mapsto(z+\gamma,w+a\overline\gamma),\
  (z,w)\mapsto(z,w+\lambda):
   \gamma\in\Lambda_B,\ \lambda\in\Lambda_F
 \right\rangle .
 \tag{C.PT.8}\hypertarget{tag-00WB}{}\label{cdp:eq:C.PT.8}
\]
The four real periods are independent by projection
to $\Lambda_B$ and then to $\Lambda_F$; they give
a discrete cocompact lattice. This is a complex
two-torus, with the original torsor class represented
by the same additive transitions. A different lift
through \eqref{cdp:eq:C.PT.7} gives its isomorphic original
$F$-torsor. No section of the original torsor has
been assumed.


\item  We now identify every marked elliptic curve on $P$
with a fibre of one such elliptic fibration.
On a torus an elliptic curve is a translate of an
elliptic subtorus: lifting its map to universal
covers makes its derivative a periodic holomorphic
function, hence constant. Two different elliptic
directions give an isogeny from their product to
the two-torus, so their translates intersect.
The pairwise disjoint boundary curves therefore
have the same direction and are fibres of the
actual quotient by that elliptic subtorus.


\item  In the other alternative let $q:P\to B$ be the
fibration supplied by \ref{cdp:st:C.2.UN.31}. If it has a multiple
fibre, any nonvertical marked elliptic curve would
map finitely to the elliptic $B$. Riemann--Hurwitz
makes that map unramified. But at a multiple fibre
the actual local function is $s^m$, $m>1$,
so the restriction to that curve has positive
ramification, a contradiction.
If there is no multiple fibre, $q$ is a principal
$F$-torsor. A nonvertical elliptic curve would give
a section after its finite base pullback. Pullback
on $H^2(B,\Lambda_F)$ multiplies its Chern vector
by the positive covering degree. The vector group
is torsion-free, so the original vector would be
zero. Formula \eqref{cdp:eq:C.PT.8} then makes $P$ a torus,
already handled above. Otherwise all marked
curves are vertical. Since every reduced fibre
is one smooth elliptic curve, each is exactly
a marked fibre support. This proves the required
actual fibration $q_P$ in all cases.


\end{enumerate}

\item \par\medskip\noindent\textbf{Finite normal roots and the exact height equations.}
The normal-root chart (C.HT.4) gives each root
boundary normal a finite exact order equal to its
full elliptic multiplicity. Normal duality carries
that order to the adjacent component.
An algebraic-dimension-zero diagonal Hopf has only
nontorsion axis normals, so it cannot be adjacent
to $P$. At an elliptic Hopf all curve normals have
finite order, and a nonlinear Hopf has only its
single torsion-normal axis. A supplied quasi-bundle
has finite boundary normal orders by (C.QT.2).
Following any path from $P$ therefore excludes
its first algebraic-dimension-zero diagonal
vertex, and hence all such vertices. A nonlinear
vertex is terminal because it has only one curve.
Every remaining nonroot vertex is an elliptic
quasi-bundle of the required genus-zero type.


\begin{enumerate}[label={$\langle2\rangle$\arabic*.},leftmargin=1.3em,labelsep=.5em]
\item  For each elliptic component, the actual root
charts give the global meromorphic form
$q^*\theta\wedge dz$ and its exact order
$m\operatorname{ord}\theta+m-1$, as in
(C.HT.5)--(C.HT.6) and (C.QT.6). These formulas
use translation charts, not the Kodaira dimension
of the component. The canonical residue form
on that component has divisor only on its
marked curves. At an unmarked multiple fibre
its order would be zero, which is incompatible
with $m\operatorname{ord}\theta+m-1=0$ for
$m>1$. This proves the absence of all unmarked
multiples on $P$ and on every elliptic nonroot
vertex.


\item  The dual-normal maps identify the connected
normal-root elliptic covers. Beginning with
the fixed Jacobian of $q_P$ and following the
unique path in the tree gives its actual
identification with every later Jacobian and
with each nonlinear axis root cover. In
particular the exact normal order $m_e$ is
the same at both endpoints.


\item  At every elliptic nonroot vertex the residue
calculation gives the integer equation
\[
 \sum_{e=(v,w)}
       \frac{k_w-k_v}{m_e}
                        =\operatorname{val}(v)-2.
 \tag{C.PT.9}\hypertarget{tag-00WC}{}\label{cdp:eq:C.PT.9}
\]
A terminal elliptic vertex gives the outward
increment $m_e$. A terminal nonlinear vertex
gives the same increment from its exact
axis form $dz\wedge dw/w^{m_e+1}$ and its
infinite-order effective axis line.
The finite leaf induction of \ref{cdp:st:C.2.UN.58} applies
without an equation at the root: once every
child term is one, \eqref{cdp:eq:C.PT.9} makes its parent
term minus one. Therefore every outward
increment is $m_e$, proving \eqref{cdp:eq:C.PT.1}.
Substituting the root increments into
\eqref{cdp:eq:C.PT.6} proves \eqref{cdp:eq:C.PT.2}.


\item  The actual root collars $xy=t$ and the full
secondary boundary matrices of \ref{cdp:st:C.2.UN.50} identify
every regular elliptic fibre and every
nonlinear real core with one transported
integral two-cycle $\gamma_t$ in the nearby
torus. They are exactly the maps
(C.HT.10)--(C.HT.12); the genus of the root
base does not change those collar maps.
For the canonical evaluation divided by
$t^{k_0}$, its leading form on an elliptic
vertex restricts to zero on the elliptic
curve. Its period is therefore
$O(t^{h_v+1})$, including at $P$.
At a nonlinear leaf it is
\[
 t^{h_v}
       \left(\frac{\pm2\pi i c\beta^{-m_e}}{n}
                         +O(t)\right),\qquad c\ne0.
 \tag{C.PT.10}\hypertarget{tag-00W0}{}\label{cdp:eq:C.PT.10}
\]
These are the same holomorphic period, by the
actual integral collar identifications.
One nonlinear leaf fixes its finite positive
order $R$; all nonlinear heights equal $R$,
and every elliptic vertex has smaller height.
This proves the stated height test.


\end{enumerate}

\item \par\medskip\noindent\textbf{The exact inverse torsor and its integral obstruction.}
Perform the operations (C.PR.6) from the leaves
inward. After its children have been removed,
an elliptic nonroot vertex retains its actual
fixed-Jacobian map over $\PP^1$: each operation
replaced exactly the specified multiple-fibre
tube by the ordinary product tube and left
all other transitions unchanged. Its only
possible multiple fibre is now the parent
mark. It is therefore the actual one-mark
torsor of (C.PR.1)--(C.PR.5), hence is a
recorded diagonal Hopf pair or the product
ruled cap. The root-line classification retains
its entire torsor, even when the original
unpruned surface had Kodaira dimension one.
Nonlinear leaves use the same
inverse product map, with their full
secondary integer and primitive fibre
direction. The root always remains its
fixed-Jacobian quasi-bundle over the original
elliptic $B$. Each operation replaces exactly
one marked multiple-fibre tube by its
displayed ordinary product tube; a product
cap leaves that ordinary tube unchanged.
There are no unmarked multiples. After all
nonroot vertices have been removed, the final
root is therefore a principal $F$-torsor
$T_0\to B$, with its full analytic transition
class retained.


\begin{enumerate}[label={$\langle2\rangle$\arabic*.},leftmargin=1.3em,labelsep=.5em]
\item  The exact boundary homomorphism of every
removed compact piece is (C.PR.9): it kills
the one meridian $\delta$ and retains the
two generators $e,h$. The inverse product
filling has those identical maps. Applying
the actual Van Kampen pushouts successively,
as in \ref{cdp:st:C.2.UN.60}, gives
\[
                         \pi_1(T_0)\simeq\pi_1(F_t)
                                             \simeq\Z^4.
 \tag{C.PT.11}\hypertarget{tag-00W1}{}\label{cdp:eq:C.PT.11}
\]
Only these smooth collar maps are used for
an arbitrary original degeneration; no
holomorphic contraction of its neighbourhood
is asserted.


\item  For a principal elliptic bundle over an
elliptic curve the homotopy sequence has
an injected fibre lattice, since
$\pi_2(B)=0$. The group is the central
extension
\[
 1\longrightarrow\Lambda_F
 \longrightarrow\pi_1(T_0)
 \longrightarrow\Lambda_B\longrightarrow1.
 \tag{C.PT.12}\hypertarget{tag-00W2}{}\label{cdp:eq:C.PT.12}
\]
The commutator of lifts of an oriented
base basis is its integral Chern vector.
This follows directly by trivializing over
the usual base parallelogram: its boundary
commutator is the clutching loop of the
fibre torus. Equation \eqref{cdp:eq:C.PT.11} makes that
commutator zero in the injected lattice.
Thus the Chern vector is zero. Formula
\eqref{cdp:eq:C.PT.8} identifies this very inverse
torsor, including its analytic class,
with a complex two-torus. This proves
necessity of the last test. Independence
of disjoint pruning orders is the actual
chart comparison of \ref{cdp:st:C.2.UN.60}.


\end{enumerate}

\item \par\medskip\noindent\textbf{Compact reconstruction of the complete prescribed tree.}
Assume all the tests pass and retain the
complete pruning record. Formula \eqref{cdp:eq:C.PT.8}
gives the actual initial torus
$T_0=\C^2/\Gamma$. Its elliptic fibre
lattice is a primitive rank-two subgroup
of $\Gamma$, because the quotient base
lattice is free. Choose $R$ as in the
height test and put
\[
 V_t(w)=w+t^R J\overline w,\qquad
              J=\begin{pmatrix}0&1\\-1&0\end{pmatrix}.
 \tag{C.PT.13}\hypertarget{tag-00W3}{}\label{cdp:eq:C.PT.13}
\]
For any real basis of $\Gamma$ its image
has determinant $(1+|t|^{2R})^2$ times
the original nonzero determinant.
These holomorphic period sections
therefore give a proper smooth torus
family on the finite chart. At infinity
the frame divided by $t^R$ gives
$J\overline w+s^Rw$, $s=t^{-1}$.
It is again real-linearly invertible
at every point, including $s=0$.
The two charts consequently give an
actual compact smooth torus family
over $\PP^1$, with central member
exactly $T_0$.


\begin{enumerate}[label={$\langle2\rangle$\arabic*.},leftmargin=1.3em,labelsep=.5em]
\item  For an oriented primitive fibre basis
$\lambda,\mu=c\lambda$,
$\operatorname{Im}c>0$, its exact period
in this frame is
\[
 \det_{\C}(V_t\lambda,V_t\mu)
       =2i\operatorname{Im}(c)\|\lambda\|^2t^R.
 \tag{C.PT.14}\hypertarget{tag-00W4}{}\label{cdp:eq:C.PT.14}
\]
The coefficient is nonzero. Near each
recorded central elliptic fibre use
the actual rank-two period cover and
the primitive-basis change (C.PR.13).
This gives
$(p,w)\mapsto(p+\delta(t),\gamma(t)w)$
with $\delta(t)=t^R d(t)$, $d(0)\ne0$,
and the original marked Jacobian and
primitive fibre direction.


\item  Reverse the pruning record. At a
recorded Hopf step of weight $m$ use
the free weighted chart (C.PR.14),
including its full secondary integer.
Its child's affine pencil has the
exact displacement $\delta(t)/t^m$.
At a product-cap step use the ordinary
blowup (C.PR.16), with displacement
$\delta(t)/t$. Thus the remaining
order at every vertex is $R-h_v$.
The height test makes each recorded
elliptic child linear and every
recorded nonlinear terminal child
nonlinear, with its prescribed
resonance and contraction. The
recorded primitive basis changes
retain all linear gluing maps,
and the retained transition charts
restore all elliptic translations
and every analytic torsor. In particular an
original properly elliptic nonroot surface
is restored by reversing its successive
root-line and ordinary-tube changes. No
replacement by a surface with only the same
degrees or normal characters occurs.


\item  All local compactness comparisons
are precisely (C.PR.17)--(C.PR.18).
Exponential contraction dominates
the polynomial exceptional norm
distortion, and the old base
displacement is
$O(|\delta(t)|\,|\log|x||)$.
Its cluster set lies in the old
marked elliptic curve. These exact
bounds give disjoint Hausdorff
proper replacements at each stage;
ordinary caps are proper blowups.
There are finitely many stages.
Each is the identity on the
punctured family and outside its
chosen small central neighbourhood.
They therefore preserve the
compact family over the rest of
$\PP^1$, including infinity.


\item  Induction through the recorded
inverse charts restores precisely
the specified $P$, every nonroot
component and each prescribed
elliptic gluing map. Their normal
equations are the free coarse
$t=xW$ charts or the ordinary
$t=xy$ charts, so the full fibre
is reduced and the total space
is smooth. Every other fibre is
the unchanged smooth torus.
The root quasi-bundle and every supplied
genus-zero quasi-bundle have Euler number
zero by additivity over their smooth
elliptic fibre supports. Every remaining
Hopf component and every double elliptic
curve also has Euler number zero. Additivity proves
the stated zero Euler number
for $S$ and for all other supports.


\item  The rooted incidence matrices,
the integral Chern commutator,
the exact torsor period matrix
and the height equations are
checked in
\path{code/A_rooted_hopf_trees.py};
the primitive period and nested
order identities are also
certified in
\path{code/A_tree_periods.py}.



\end{enumerate}

\end{enumerate}
\end{proof}






\begin{proposition}[Unicyclic quasi-bundle supports with all torsors retained]\hypertarget{tag-01E8}{}\label{cdp:st:C.2.UN.62}
\StatementRecord{cdp:st:C.2.UN.62}{proved}{01E8}{\SeriesRef{cdp:st:C.2.UN.60}, \SeriesRef{cdp:eq:C.QT.2}, \SeriesRef{cdp:st:C.2.UN.54}, \SeriesRef{cdp:st:C.2.UN.50}, \SeriesRef{cdp:st:C.2.UN.44}, \SeriesRef{cdp:eq:C.QT.3}, \SeriesRef{cdp:eq:C.QT.7}, \SeriesRef{cdp:eq:C.QT.6}, \SeriesRef{cdp:eq:C.QT.8}, \SeriesRef{cdp:eq:C.QT.5}, \SeriesRef{cdp:eq:C.QT.9}, \SeriesRef{cdp:st:C.2.UN.57}, \SeriesRef{cdp:eq:C.QT.11}}

Let $S$ be a specified compact connected SNC surface with smooth
embedded components. Each component is either a nonlinear primary
or secondary Hopf surface of algebraic dimension zero, or is
equipped with an actual map
\[
                         q_v:T_v\longrightarrow\PP^1
\tag{C.QT.1}\hypertarget{tag-00WD}{}\label{cdp:eq:C.QT.1}
\]
with connected fibres, every fibre a multiple of a smooth elliptic
curve, fixed elliptic Jacobian, and identity monodromy on the
integral first cohomology of the regular fibre. Every boundary
curve on $T_v$ is a fibre support of this specified map. The
elliptic gluing maps, the surfaces and their complete analytic
torsors are part of the data.

The following geometric test is necessary and sufficient for
$S$ to occur as a full reduced fibre of a smooth-total
degeneration of connected complex two-tori.
\begin{enumerate}
\item The actual boundary normal lines are dual, and the incidence
graph has exactly one independent cycle. Nonlinear Hopf vertices
are leaves.
\item Every multiple fibre of every map \eqref{cdp:eq:C.QT.1} is marked by
a boundary curve. Let $m_e$ be the exact order of its normal
line, also the normal order at the other endpoint. At a
nonlinear endpoint this is its resonance.
\item Set $h_v=0$ on the cycle and set $h_v$ equal to the sum
of the $m_e$ along its path from the cycle. There is an integer
$R>0$ for which every nonlinear leaf has height $R$ and every
elliptic quasi-bundle vertex has height strictly smaller than
$R$. With no nonlinear leaves one may choose any
$R>\max_v h_v$.
\item Apply the exact inverse-logarithmic and product-cap
pruning maps of \ref{cdp:st:C.2.UN.60}, retaining their full integer
framings and all torsor identifications. The resulting
marked diagonal Hopf or degree-zero ruled core cycle has
integral complement-lattice product $M=I$.
\end{enumerate}
All pruning operations in this test exist under its preceding
requirements. Their marked result is independent of the order
of disjoint branch operations. Every support passing the test
has a realization over $\PP^1$ with smooth compact total space,
retaining every original component, analytic torsor and
specified elliptic gluing map. The only other singular support
is the complex conjugate of its pruned core. Every reduced
fibre support has Euler number zero.

In particular the statement includes elliptic quasi-bundles
of Kodaira dimension one and coherent Euler characteristic
zero, together with the genus-zero quasi-bundle components
of Kodaira dimension zero and the Hopf and ruled components
in the same explicitly specified class.
\end{proposition}



\begin{proof}\leavevmode\par



 
\begin{enumerate}[label={$\langle1\rangle$\arabic*.},leftmargin=0pt,labelsep=.5em]
\item \par\medskip\noindent\textbf{The actual root chart and its cohomology maps.}
Let $mJ$ be a fibre of \eqref{cdp:eq:C.QT.1}. The normalized local root
$t=s^m$ is finite unramified over the surface tube, by
the exact chart $s=wx$, $w^m=u$ for $t=ux^m$.
Its generic fibre is connected. Its finite Stein factor
over the root disc has degree one, so the central
elliptic cover is connected as well. The fixed Jacobian
and identity linear monodromy make the root family a
translation torsor under the fixed elliptic curve $F$.
The sequence $0\to\Lambda_F\to\OO\to F\to0$ and
$H^1(\Delta,\OO)=H^2(\Delta,\Lambda_F)=0$ trivialize
this actual torsor. The finite translation gauge in
(C.HT.4) therefore gives the surface tube exactly as
\[
 (F\times\Delta_s)/\mu_m,
 \qquad(z,s)\longmapsto(z+\tau,\zeta_ms),
 \qquad\operatorname{ord}(\tau)=m.
\tag{C.QT.2}\hypertarget{tag-00WI}{}\label{cdp:eq:C.QT.2}
\]
The last order is exact because the central action is
free. Its associated normal character is primitive.
Thus $N_{J/T_v}$ has exact order $m$, and its connected
normal-root cover is precisely $F\to J$.


\begin{enumerate}[label={$\langle2\rangle$\arabic*.},leftmargin=1.3em,labelsep=.5em]
\item  The fixed cohomology frame $H^1(F,\OO_F)$ on this
root chart is unchanged by the finite translations.
Taking invariant functions gives the ordinary
coordinate $s^m$. The regular frames glue by the
specified identity monodromy and by translations.
Consequently the actual direct images and their
cohomology are
\[
 q_{v*}\OO_{T_v}=\OO_{\PP^1},\qquad
 R^1q_{v*}\OO_{T_v}=\OO_{\PP^1},\qquad
 (h^0,h^1,h^2)(T_v,\OO)=(1,1,0).
\tag{C.QT.3}\hypertarget{tag-00WJ}{}\label{cdp:eq:C.QT.3}
\]
At a marked fibre $J$, restriction to $J$ followed
by the finite pullback $F\to J$ is evaluation on this
same cohomology frame. The resulting map
$H^1(T_v,\OO)\to H^1(F,\OO_F)$ is an isomorphism.
In particular the original restriction to $J$ is
nonzero. The commuting square is the actual root
chart \eqref{cdp:eq:C.QT.2}; no identification of different
cohomology groups by dimensions alone is used.


\item  A nonlinear Hopf leaf has the same values $(1,1,0)$
and the nonzero restriction map proved in \ref{cdp:st:C.2.UN.54}.
Duality of boundary normal lines identifies their
exact orders and their connected normal-root covers.
It thus identifies the endpoint Jacobians by actual
finite-cover maps, including the secondary character
of a nonlinear leaf from \ref{cdp:st:C.2.UN.50}.


\end{enumerate}

\item \par\medskip\noindent\textbf{Graph, residues and the canonical minimum.}
Distinct fibre supports on each $T_v$ are disjoint,
and a nonlinear Hopf has only its one elliptic curve.
Hence the specified SNC surface has no triple points.
For an actual degeneration its full fibre is reduced,
so the meridian quotient of \ref{cdp:st:C.2.UN.44} makes $\pi_1(S)$
abelian. The actual map onto the incidence graph
then makes that graph's free fundamental group
abelian. It has at most one independent cycle.


\begin{enumerate}[label={$\langle2\rangle$\arabic*.},leftmargin=1.3em,labelsep=.5em]
\item  The normalization sequence is exactly
\[
 0\longrightarrow\OO_S\longrightarrow
   \bigoplus_v\OO_{T_v}\longrightarrow
   \bigoplus_e\OO_{J_e}\longrightarrow0,
\tag{C.QT.4}\hypertarget{tag-00WK}{}\label{cdp:eq:C.QT.4}
\]
with nonlinear Hopf terms included in the vertex
sum. The induced $H^1$ matrix is a weighted
incidence matrix whose every endpoint coefficient
is nonzero, by \eqref{cdp:eq:C.QT.3}. A tree gives a surjective
matrix by leaf elimination, and therefore
$H^2(S,\OO_S)=0$. Flatness over the smooth curve
and upper semicontinuity from a torus instead
give $h^2(S,\OO_S)\ge1$. Thus the graph is
unicyclic. Its spanning-tree minor has rank
$\#V-1$; the remaining row must vanish on that
minor's kernel. This is exactly product-one
holonomy for the actual endpoint restriction
maps, and
\[
                    (h^0,h^1,h^2)(S,\OO_S)=(1,2,1).
\tag{C.QT.5}\hypertarget{tag-00WL}{}\label{cdp:eq:C.QT.5}
\]


\item  Serre duality identifies the transpose incidence
map with the component residue obstruction maps.
Its kernel has zero residues on every attached
tree, by leaf elimination, and nonzero residues
on every cycle edge. Since every vertex has
$H^0(K_{T_v})=0$, a component section with all
residues zero is zero. The unique global section
of $\omega_S$ is consequently nonzero exactly
on the cycle components. Divide the actual
canonical evaluation by its minimum full-fibre
power. Its nonzero reduction is this section.
Its minimum vertical order $k_0$ is therefore
attained exactly on the cycle, as in the explicit
evaluation comparison of \ref{cdp:st:C.2.UN.54}.


\end{enumerate}

\item \par\medskip\noindent\textbf{The exact differential and height calculation.}
On a vertex \eqref{cdp:eq:C.QT.1}, invariant fibre differentials
give actual meromorphic two-forms
$q_v^*\theta\wedge dz$. Translation charts
change $dz$ only by a multiple of the base
differential, which disappears in this wedge.
The forms therefore glue. Every meromorphic
two-form with vertical divisor has this form:
its ratio to one reference form is constant
on the regular compact fibres and descends;
\eqref{cdp:eq:C.QT.2} makes its descended meromorphic
function invariant under $s\mapsto\zeta_ms$,
so it is meromorphic in $s^m$ across a multiple
fibre. The local order is exactly
\[
 \operatorname{ord}_{J}(q_v^*\theta\wedge dz)
             =m\operatorname{ord}_{q_v(J)}\theta+m-1.
\tag{C.QT.6}\hypertarget{tag-00WM}{}\label{cdp:eq:C.QT.6}
\]


\begin{enumerate}[label={$\langle2\rangle$\arabic*.},leftmargin=1.3em,labelsep=.5em]
\item  Restriction of the canonical evaluation divided
by $t^{k_v}$ gives the actual meromorphic
component form with divisor
\[
           \sum_{e=(v,w)}(k_w-k_v-1)J_e.
\tag{C.QT.7}\hypertarget{tag-00WN}{}\label{cdp:eq:C.QT.7}
\]
An unmarked multiple fibre would have order zero
in \eqref{cdp:eq:C.QT.7}, and \eqref{cdp:eq:C.QT.6} would require
$m\operatorname{ord}\theta+m-1=0$, hence $m=1$.
At a marked fibre these same equations give
\[
 d_{v,e}:=\frac{k_w-k_v}{m_e}
       =\operatorname{ord}_{q_v(J_e)}\theta+1\in\Z,
 \qquad\sum_{e\ni v}d_{v,e}=\operatorname{val}(v)-2.
\tag{C.QT.8}\hypertarget{tag-00WO}{}\label{cdp:eq:C.QT.8}
\]
The last equality is the degree $-2$ of an
actual rational differential on $\PP^1$.
Opposite orientations negate $d$ because their
normal-root covers have the same exact order.


\item  At a terminal elliptic vertex \eqref{cdp:eq:C.QT.8} gives
height difference $m_e$ from its parent. A
nonlinear terminal vertex gives precisely the
same difference by its invariant axis form
$dz\wedge dw/w^{m_e+1}$ and the exact secondary
normal comparison in \ref{cdp:st:C.2.UN.50}. Induct inward along
each tree. If the $c$ child terms are one,
\eqref{cdp:eq:C.QT.8} gives $-d_{\rm parent}+c=c-1$,
so the parent term is also one. The cycle
terms are zero by the canonical minimum
already identified. Hence $k_v-k_0$ is exactly
the stated path sum $h_v$.


\item  This proof also gives the actual base differentials:
on a cycle vertex it is $c\,dz/z$, with the two
cycle points $0,\infty$; on a tree vertex it is
$c\,dz$, with its parent point at infinity.
The other marked points have order zero.
No Kodaira-dimension assumption was used in
this calculation.


\end{enumerate}

\item \par\medskip\noindent\textbf{The common integral period.}
Dual normal-root charts give on an edge
$xy=t$, $q_v=x^{m_e}$ and $q_w=y^{m_e}$,
with the same connected elliptic cover and
the actual normal-product trivialization.
The level $x^{m_e}=c$ corresponds to
$y^{m_e}=t^{m_e}/c$. Varying its nonzero
normal level gives the oriented isotopy
between the two regular-fibre two-tori.
The product-preserving tubular comparison
transfers this equality to the actual
nearby torus, as in (C.HT.10).
At a nonlinear leaf the unimodular boundary
map of \ref{cdp:st:C.2.UN.50} identifies its real core with
this same integral class. All vertices
therefore give one class $\gamma$.


\begin{enumerate}[label={$\langle2\rangle$\arabic*.},leftmargin=1.3em,labelsep=.5em]
\item  The finite-root restriction square in
\eqref{cdp:eq:C.QT.3} also shows that the coherent
cycle holonomy in \eqref{cdp:eq:C.QT.5} is the linear
action of the actual product of elliptic
root-cover isomorphisms. Its value one
forces that product to be a translation.
This retains its possible translation
and all Jacobian identifications.


\item  The period of the evaluation divided by
$t^{k_0}$ is holomorphic. On a regular
holomorphic fibre curve of $T_v$ its
leading two-form restricts to zero, so
\[
                  P(t)=O(t^{h_v+1}).
\tag{C.QT.9}\hypertarget{tag-00WP}{}\label{cdp:eq:C.QT.9}
\]
At a nonlinear leaf of secondary order
$n$ and contraction parameter $\beta$,
the exact primitive core computation is
\[
 P(t)=t^{h_v}
       \left(\frac{\pm2\pi i c\beta^{-m_e}}n+O(t)\right),
                         \qquad c\ne0.
\tag{C.QT.10}\hypertarget{tag-00WE}{}\label{cdp:eq:C.QT.10}
\]
Thus every nonlinear leaf has the one
finite order $R$ of this actual common
period, and \eqref{cdp:eq:C.QT.9} makes every elliptic
vertex's height smaller. Without such a
leaf this condition places no further
restriction: an integer above the finite
set of heights exists. This proves all
the height requirements for actual fibres.


\end{enumerate}

\item \par\medskip\noindent\textbf{Pruning with the full torsors.}
The pruning proof of \ref{cdp:st:C.2.UN.60} applies to the
maps \eqref{cdp:eq:C.QT.1} themselves. We specify the
comparisons that retain the present
broader surface class. Delete a terminal
nonlinear Hopf by the inverse chart
(C.PR.6), with its full primitive period
basis and integer $n$. After the children
of an elliptic vertex have been removed,
its actual modified map still has fixed
Jacobian and identity monodromy over
$\PP^1$; each inverse operation replaced
one multiple fibre by the displayed
ordinary product tube and changed no
other transition. Its only possible
multiple fibre is now its parent mark.


\begin{enumerate}[label={$\langle2\rangle$\arabic*.},leftmargin=1.3em,labelsep=.5em]
\item  It is therefore the actual $F$-torsor
on $\PP^1[(\infty,m)]$ classified by
(C.PR.1)--(C.PR.5). The exact logarithm
sequence and $H^1(\OO)=H^2(\OO)=0$
identify the entire torsor with its
two root-line classes. A nonzero
vector gives the literal weighted
Hopf quotient, including $n$ and
the primitive fibre direction;
a zero vector forces $m=1$ and
the product ruled cap. This conclusion
concerns the actual modified surface,
even when its original unmodified
surface had Kodaira dimension one.
Remove that pair by its corresponding
map and retain the entire inverse
record.


\item  At the end the two-mark root-stack
cover (C.PR.7) and its actual line
pushout identify each core vertex
with a marked diagonal Hopf surface
or the torsion ruled bundle (C.PR.8).
These are exact maps on the original
transition atlases. Disjoint branch
operations commute on those atlases,
so their marked result has the
independence stated in \ref{cdp:st:C.2.UN.60}.


\item  For an actual degeneration, every
terminal complement and its inverse
product cap have the identical
boundary homomorphism (C.PR.9),
retaining both elliptic generators
and killing exactly the specified
primitive meridian. The van Kampen
pushouts are identical on the
unchanged exterior. Their iteration
identifies the core plumbing group
with the original $\Z^4$. The
injected-lattice semidirect product
for that core consequently has
$M=I$. This proves necessity of
the last test by actual group maps.


\end{enumerate}

\item \par\medskip\noindent\textbf{Compact restoration of every prescribed component.}
Conversely suppose the test passes.
Use the compact marked core of \ref{cdp:st:C.2.UN.57}
with the power $V_{t^R}$ and the
unchanged logarithmic coefficient
of (C.PR.11). Its necks remain
$xy=t$ and its opposite filling
is the conjugate pruned core.
The compact elliptic subgroup of
its common open group is the
actual regular-fibre group, and
(C.PR.12) makes its transverse
period determinant a nonzero
constant times $t^R$.


\begin{enumerate}[label={$\langle2\rangle$\arabic*.},leftmargin=1.3em,labelsep=.5em]
\item  Reverse the entire pruning record.
At a vertex of height $h$, the
local quotient (C.PR.13) has
displacement of order $R-h$.
The recorded primitive basis is
implemented by the displayed
unimodular linear period map
following (C.PR.13). At a Hopf
step, the weight-$(m,1)$ insertion
with $\OO(nE)$ has exactly the
free finite chart (C.PR.14).
The new child's coordinate
$U=Z/W^m$ gives the actual
displacement $\delta(t)/t^m$,
as in (C.PR.15). At a product
cap the ordinary blowup has
the exact displacement
$\delta(t)/t$, as in (C.PR.16).
The height test makes all these
maps holomorphic and gives
precisely the prescribed diagonal
or nonlinear central type.


\item  All operations use the entire
recorded inverse transitions.
In particular an original
properly elliptic vertex is
restored by reversing the
successive changes of its
root lines and ordinary tubes;
its complete analytic torsor,
all multiple fibres and all
marked points reappear, rather
than only its degree or normal
characters. The strict-parent
map is exactly (C.PR.6) in the
opposite direction. The same
argument applies to every
other original quasi-bundle.
Fibre translations and the
recorded linear maps retain
each specified elliptic gluing.


\item  The norm and affine-base bounds
(C.PR.17)--(C.PR.18) prove the
actual local properness and
the old-cluster comparison for
each replacement. Their constants
depend on the local displacement
order and contraction, not on the
Kodaira dimension of the surface
being restored. The finitely many
disjoint replacements therefore
give the same Hausdorff compact
gluing as \ref{cdp:st:C.2.UN.60}. The weighted
free coordinates have $t=xW$;
ordinary cap charts have $t=xy$.
Every final component consequently
has full multiplicity one.


\item  Each component \eqref{cdp:eq:C.QT.1} has Euler
number zero by additivity over
its elliptic fibre supports;
each nonlinear Hopf also has
Euler number zero. Every double
curve is elliptic and there are
no triple points. Both special
supports thus have Euler zero,
as do the general tori. This
proves the complete criterion
and the compact assertion.


\end{enumerate}

\item \par\medskip\noindent\textbf{An explicit properly elliptic component.}
Fix $F=\C^*/\beta^{\Z}$, $0<|\beta|<1$.
Start with the two-component
ruled core $F\times\PP^1$
glued along its two marked
sections by the identity.
Its lattice product is the
identity. On one component
choose four distinct ordinary
fibre points away from those
sections. Perform four recorded
weight-three, secondary-order-one
insertions, with the common
period order $R=3$.
They add four nonlinear primary
Hopf leaves $H_3(\beta)$.


\begin{enumerate}[label={$\langle2\rangle$\arabic*.},leftmargin=1.3em,labelsep=.5em]
\item  The modified central component
has an explicit independent
description. Put
\[
 \mathcal B=\PP^1[(p_1,3),\ldots,(p_4,3)],\qquad
 L=\OO_{\mathcal B}(D_1+\cdots+D_4),
 \qquad 3D_i=p_i.
\tag{C.QT.11}\hypertarget{tag-00WF}{}\label{cdp:eq:C.QT.11}
\]
It is the smooth coarse surface
of $L^\times/\beta^{\Z}$.
Each inertia acts by a primitive
fibre translation, so its coarse
surface is smooth and its four
fibres are exactly triple smooth
elliptic curves. In a root chart
$p=s^3$ its line coordinate is
$y=sW$, where $W$ is the old
product coordinate. The inverse
chart (C.PR.6) with $m=3,n=1$
has precisely $W=y/s$.
Thus this given line pushout
is the component produced by
the four insertions, including
its full torsor.
Explicitly the positive divisor
line has frames $e_{\rm in}=s^{-1}e_{\rm out}$.
The equality of the actual vector
$W e_{\rm out}=y e_{\rm in}$
is precisely $y=sW$; its canonical
section has local coefficient $s$.
This fixes the sign of each $D_i$
in \eqref{cdp:eq:C.QT.11}.


\item  For any map of the class
\eqref{cdp:eq:C.QT.1}, the actual invariant
fibre differential and the
root chart give the plurigenus
formula
\[
 h^0(T_v,\omega_{T_v}^{\otimes a})
 =\max\left\{0,\ 1-2a+
       \sum_i\left\lfloor
                    a(1-1/m_i)\right\rfloor\right\}.
\tag{C.QT.12}\hypertarget{tag-00WG}{}\label{cdp:eq:C.QT.12}
\]
Indeed a pluricanonical section
is a base meromorphic
$a$-differential times the
$a$th invariant fibre factor;
holomorphy at a root chart
permits precisely the displayed
pole order. Descent and its
inverse are those actual
invariant-coordinate maps.
For \eqref{cdp:eq:C.QT.11} this gives
\[
 h^0(T,\omega_T^{\otimes a})
    =\max\{0,1-2a+4\lfloor2a/3\rfloor\},
 \qquad \omega_T^{\otimes3}\simeq q^*\OO_{\PP^1}(2).
\tag{C.QT.13}\hypertarget{tag-00WH}{}\label{cdp:eq:C.QT.13}
\]
In particular $p_g(T)=0$,
$q(T)=1$, $\chi(\OO_T)=0$
and $h^0(T,\omega_T^{\otimes3b})=2b+1$.
Thus this is an actual properly
elliptic component. The complete
compact reduced fibre consists
of it, the unchanged ruled core
component, and the four Hopf
leaves, with every gluing given
by the preceding maps.
The root-line maps, the incidence
matrix, the root-coordinate pole
bounds and these exact plurigenera
are checked in
\path{code/D_quasibundle_trees.py}
and its saved output.



\end{enumerate}

\end{enumerate}
\end{proof}






\begin{proposition}[The weighted-root criterion for elliptic and Hopf trees]\hypertarget{tag-01E9}{}\label{cdp:st:C.2.UN.63}
\StatementRecord{cdp:st:C.2.UN.63}{proved}{01E9}{\SeriesRef{cdp:st:C.2.UN.59}, \SeriesRef{cdp:eq:C.WT.7}, \SeriesRef{cdp:st:C.2.UN.37}, \SeriesRef{cdp:eq:C.WT.10}, \SeriesRef{cdp:st:C.2.UN.40}, \SeriesRef{cdp:st:C.2.UN.28}, \SeriesRef{cdp:eq:C.WT.9}, \SeriesRef{cdp:st:C.2.UN.61}, \SeriesRef{cdp:st:C.2.UN.62}, \SeriesRef{cdp:st:C.2.UN.47}, \SeriesRef{cdp:st:C.2.UN.58}, \SeriesRef{cdp:st:C.2.UN.50}, \SeriesRef{cdp:eq:C.WT.14}, \SeriesRef{cdp:eq:C.WT.8}, \SeriesRef{cdp:eq:C.WT.15}, \SeriesRef{cdp:st:C.2.UN.51}, \SeriesRef{cdp:st:C.2.UN.55}, \SeriesRef{cdp:eq:C.WT.18}, \SeriesRef{cdp:eq:C.WT.19}, \SeriesRef{cdp:eq:C.WT.23}, \SeriesRef{cdp:st:C.2.UN.57}, \SeriesRef{cdp:eq:C.WT.30}, \SeriesRef{cdp:eq:C.WT.32}, \SeriesRef{cdp:st:C.2.UN.60}, \SeriesRef{cdp:eq:C.WT.34}, \SeriesRef{cdp:eq:C.WT.31}, \SeriesRef{cdp:eq:C.WT.28}, \SeriesRef{cdp:eq:C.WT.36}, \SeriesRef{cdp:eq:C.WT.33}, \SeriesRef{cdp:eq:C.WT.38}, \SeriesRef{cdp:eq:C.WT.42}, \SeriesRef{cdp:eq:C.WT.25}}

Specify a connected SNC surface with embedded smooth components,
edges between distinct components, and positive component
multiplicities \(m_i\). Each component is a primary or secondary
Hopf surface, or is supplied with its actual fixed-Jacobian,
identity-integral-monodromy elliptic quasi-bundle over a smooth
compact curve, with connected fibres and only smooth multiple
supports. At a supplied quasi-bundle every marked double curve
is a fibre support. Retain every component, analytic torsor and
elliptic gluing map.

Specify the actual weighted line data
\[
 \begin{gathered}
 \sigma_i\in H^0(S_i,\mathcal L_i^{-m_i}),\qquad
 \operatorname{div}\sigma_i=\sum_{j\sim i}m_jJ_{ij},\\
 \mathcal L_i|_{J_{ij}}\simeq
       \phi_{ij}^{*}N_{J_{ji}/S_j},\qquad
 \phi_{ij}^{*}\mathcal L_j|_{J_{ji}}
       \simeq N_{J_{ij}/S_i},
 \end{gathered}
 \tag{C.WT.1}\hypertarget{tag-017O}{}\label{cdp:eq:C.WT.1}
\]
with the equal leading-coefficient line maps at each edge.
Form the finite normalizations
\[
 \mathcal E_i=\operatorname{Norm}
       \{\xi\in\mathcal L_i^{-1}:\xi^{m_i}=\sigma_i\},
 \qquad N=\operatorname{lcm}_i m_i,
 \tag{C.WT.2}\hypertarget{tag-017Z}{}\label{cdp:eq:C.WT.2}
\]
their actual elliptic boundary maps and their cyclic
\(\mu_N\)-action. At an edge of multiplicities \(a=gA,b=gB\),
where \(\gcd(A,B)=1\), retain the literal normalized neck
\[
       uv=s^h,\qquad x=u^B,\quad y=v^A,\qquad
       h=\frac{N}{gAB},\qquad t=s^N=x^a y^b.
 \tag{C.WT.3}\hypertarget{tag-018B}{}\label{cdp:eq:C.WT.3}
\]
These are compressed necks; no unspecified contraction of a
reconstructed ruled chain is part of the test.

The following test is necessary and sufficient for the given
weighted support to occur as a fibre of a smooth-total
degeneration of connected complex two-tori.

\begin{enumerate}
\item The root graph is connected. The branch containing an
algebraic-dimension-zero diagonal Hopf component is the complete
two-axis cycle and passes \ref{cdp:st:C.2.UN.59}. In every remaining branch,
each strict elliptic root cover has base genus at most one, and
at most one such cover has genus one.
\item If all elliptic root bases have genus zero, the compressed
root graph is unicyclic. If there is a genus-one root component,
the graph is a tree rooted there. All other root components are
the supplied genus-zero elliptic quasi-bundles or nonlinear
Hopf leaves. There are no unmarked multiple elliptic fibres.
The actual normal lines at each root edge are dual.
\item Let \(m_e\) be their common exact torsion order, and let
\(h_e\) be the exponent in \eqref{cdp:eq:C.WT.3}. Give the cycle, or the
genus-one root, height zero, and set
\[
       H_v=\sum_{e\text{ on the rooted path to }v}h_e m_e.
 \tag{C.WT.4}\hypertarget{tag-018M}{}\label{cdp:eq:C.WT.4}
\]
All nonlinear leaves have one common height \(R>0\);
every elliptic vertex has height strictly less than \(R\).
When there is no nonlinear leaf, \(R\) may be chosen larger
than all the heights.
\item Perform the exact inverse operations below in whole
cyclic-group orbits, retaining the complete torsors, the full
secondary integers, primitive fibre bases and all marked maps.
The final cycle has actual rank-three complement product
\(M=I\), or the final principal elliptic torsor over its
elliptic base has integral Chern vector zero.
\item Let \(D\) be the finite linear part on this actual pruned
core. For a cycle, it is the linear part on its common
two-dimensional logarithmic vector space; for a torus core,
it is the linear part of its actual affine torus action.
For the root generator with base action \(s\mapsto\zeta_Ns\),
the height satisfies
\[
                         \zeta_N^R=\det_{\C}D.
 \tag{C.WT.5}\hypertarget{tag-018R}{}\label{cdp:eq:C.WT.5}
\]
In the absence of nonlinear leaves one chooses any sufficiently
large positive \(R\) in this residue class.
\end{enumerate}

Every step uses finite covers, actual line and torsor maps,
integral lattices, and the displayed inverse charts. In
particular \eqref{cdp:eq:C.WT.5} is an equality of the actual canonical
character. A cycle action may reverse orientation through
vertices; such reflections are included.

Every support passing the test has a realization over
\(\PP^1\) with smooth compact total threefold. The construction
retains its original components, gluing maps and all
multiplicities. It can have exactly the prescribed weighted
support at zero and its complex conjugate at infinity, with
all other fibres smooth connected two-tori. Every reduced
fibre support has Euler number zero. Normal-product constants
are chosen in the realization; no equivalence of different
full analytic neighbourhoods is asserted.
\end{proposition}



\begin{proof}\leavevmode\par



 
\begin{enumerate}[label={$\langle1\rangle$\arabic*.},leftmargin=0pt,labelsep=.5em]
\item \par\medskip\noindent\textbf{The actual component covers and their local inertia.}
An actual map supplies \eqref{cdp:eq:C.WT.1} by dividing its defining
function by the \(m_i\)-th power of the canonical divisor
section and restricting to \(S_i\). On its strict root
component the exact formula is
\(\xi=s^{N/m_i}/v_i\), where \(v_i\) is that normal
divisor coordinate. Thus the root generator is
\[
                  \xi\longmapsto
                         \zeta_N^{N/m_i}\xi.
 \tag{C.WT.6}\hypertarget{tag-018S}{}\label{cdp:eq:C.WT.6}
\]
At a double curve, absorbing a unit in one normal
coordinate gives \(x^a y^b=s^N\). Its \(g\) local
branches have precisely the finite birational
normal maps \eqref{cdp:eq:C.WT.3}, as proved in (C.WH.4).
The equal leading-coefficient maps in \eqref{cdp:eq:C.WT.1}
identify their actual curve covers and normal
characters. This constructs the compressed root
data without replacing a nonlinear Hopf
neighbourhood by its normal line.


\begin{enumerate}[label={$\langle2\rangle$\arabic*.},leftmargin=1.3em,labelsep=.5em]
\item  Every strict component cover is smooth. Locally along
a marked curve its equation is
\(\xi^m=z^b\) times a unit. Write
\(d=\gcd(m,b)\), \(q=m/d\), \(B=b/d\).
On each branch its normalization is
\[
                         z=w^q,\qquad \xi=cw^B,
 \tag{C.WT.7}\hypertarget{tag-018T}{}\label{cdp:eq:C.WT.7}
\]
with \(c\) a fixed root of unity. This is a smooth
coordinate chart with the unchanged curve
coordinate. Away from these curves the map is
finite \'{e}tale. The induced finite maps of
reduced boundary curves are therefore \'{e}tale.
Every stabilizer preserving such a connected
elliptic curve acts on it by deck translations.
The same is true on its connected normal-root
elliptic cover.


\item  For a supplied elliptic component, remove the finitely
many marked base values. A connected component of
its finite surface cover has a fixed finite-index
fibre subgroup \(\Lambda'\subset\Lambda_F\).
Take the actual connected-fibre Stein base cover.
On it the fibre is the fixed elliptic curve
\(F'=\C/\Lambda'\), with the actual isogeny
\(\pi:F'\to F\). A lifted transition \(a\)
satisfies
\[
                            \pi a=T_b\pi.
 \tag{C.WT.8}\hypertarget{tag-018U}{}\label{cdp:eq:C.WT.8}
\]
Its derivative is therefore the identity, so it
is a translation. This proves identity integral
monodromy on the connected-fibre base cover,
including loops which permuted the original
components of a fibre cover.


\item  Near a deleted value lift the actual chart
\((F\times\Delta)/\langle(\tau,\zeta_m)\rangle\)
and apply \eqref{cdp:eq:C.WT.7} to its normal power. Quotient
first by the subgroup trivial on \(F'\); it
acts by normal pseudoreflections. The remaining
inertia acts freely by translations on \(F'\).
The result is again a smooth multiple-fibre
chart
\((F'\times\Delta)/\langle(\tau',\zeta_{m'})\rangle\).
The finite translation gauge of (C.HT.4) makes
\(\tau'\) constant. Thus the compact strict cover
has an actual fixed-Jacobian quasi-bundle over
the smooth compact Stein base \(B'\), with only
smooth multiple supports. Its invariant
cohomology frame gives
\[
                  R^1q_*\OO=\OO_{B'},\qquad
                         p_g=g(B').
 \tag{C.WT.9}\hypertarget{tag-018V}{}\label{cdp:eq:C.WT.9}
\]
This argument also applies to every
algebraic-dimension-one Hopf component, using
its actual fibration from \ref{cdp:st:C.2.UN.37}.


\item  For a nonlinear Hopf component there is a
different explicit map. On its primary cover
write the contraction
\[
                    (z,w)\longmapsto
                     (\beta^r z+c w^r,\beta w),
                 \qquad c\ne0 .
 \tag{C.WT.10}\hypertarget{tag-017P}{}\label{cdp:eq:C.WT.10}
\]
Its only curve is \(w=0\). The power of the
weighted normal line is the constant-automorphy
axis line; dividing its transferred connection
gives a flat connection on the line itself.
On \(\C^2\setminus0\) the root section is a
constant times \(w^b\). Indeed its quotient
by \(w^b\) extends as a holomorphic unit at
zero; its automorphy factor is one by
evaluation there, and contraction makes that
unit constant. Formula \eqref{cdp:eq:C.WT.7} is consequently
the actual map \((z,v)\mapsto(z,v^q)\).
A positive power of \eqref{cdp:eq:C.WT.10} preserving one
branch lifts to
\[
                 (z,v)\longmapsto
                    (\alpha^{qr}z+c'v^{qr},\alpha v),
                     \qquad c'\ne0 .
 \tag{C.WT.11}\hypertarget{tag-017Q}{}\label{cdp:eq:C.WT.11}
\]
Its branch stabilizer is free on
\(\C^2\setminus0\). A fixed point would project
to a fixed point of the original free deck
group; an element over its identity fixing
both \(v^q\) and \(v^B\) fixes \(v\), since
\(\gcd(q,B)=1\). Hence the finite normalization
is exactly a primary or secondary nonlinear
Hopf surface, with the displayed map and
resonance \(qr\).


\item  The nontorsion axis normals of a diagonal
algebraic-dimension-zero Hopf cannot satisfy
the weighted normal equality with a torsion
normal. Its axis covers remain nonresonant:
a multiplicative relation after taking finite
coordinate roots would give a multiplicative
relation for the original multipliers.
\ref{cdp:st:C.2.UN.40} forces both its axes into the boundary.
Thus that connected branch is the complete
cycle already covered by \ref{cdp:st:C.2.UN.59}. We henceforth
treat the torsion-normal branch.


\item  Resolve the \(A_{h-1}\) necks only to apply
the compact-core inequality \ref{cdp:st:C.2.UN.28}. It gives
\(\sum p_g\le1\) for all strict components.
By \eqref{cdp:eq:C.WT.9}, their base genera are at most one
and at most one is one. No assertion that
a many-marked Hopf root cover is still Hopf
has been made.


\end{enumerate}

\item \par\medskip\noindent\textbf{Compressed canonical heights and the actual period character.}
The central fibre of \(uv=s^h\) is reduced,
and its abstract two-component crossing is
SNC. Resolving the neck inserts the actual
elliptic ruled chain, with the same normal
root maps. Either description gives the
same component and edge cohomology maps.
The exact restriction square (C.QT.2)--(C.QT.3)
and the normalization and residue arguments
of \ref{cdp:st:C.2.UN.61}--\ref{cdp:st:C.2.UN.62} therefore give the following.
With no genus-one vertex the root graph is
unicyclic, and the minimum canonical order
is attained exactly on its cycle. With one
genus-one vertex the graph is a tree and
that vertex is the unique minimum. One can
also apply \ref{cdp:st:C.2.UN.47} to the latter assertion.
The connectedness here follows from the
proper Stein factor of the actual root
family, whose nearby fibres are connected.


\begin{enumerate}[label={$\langle2\rangle$\arabic*.},leftmargin=1.3em,labelsep=.5em]
\item  Here is the precise change caused by a
compressed neck. Use a curve coordinate
\(z\) and the relative canonical frame
\(du/u\wedge dz=-dv/v\wedge dz\).
If the leading form on the \(u\)-side has
curve order \(d\), then
\[
 s^k u^{d+1}\frac{du}{u}\wedge dz
      =-s^{k+h(d+1)}v^{-d-1}
                         \frac{dv}{v}\wedge dz .
 \tag{C.WT.12}\hypertarget{tag-017R}{}\label{cdp:eq:C.WT.12}
\]
Thus the canonical-order difference is
\(h(d+1)\), rather than \(d+1\).
At an elliptic multiple fibre of order \(m\),
the exact wedge order is
\(d=m\operatorname{ord}\theta+m-1\).
Consequently the edge integer is
\[
       \frac{k_w-k_v}{h_e m_e}
                        =\operatorname{ord}_{p_e}\theta+1.
 \tag{C.WT.13}\hypertarget{tag-017S}{}\label{cdp:eq:C.WT.13}
\]
At an unmarked multiple fibre the required
order is zero; the formula excludes every
unmarked \(m>1\). The degree of the actual
base differential gives the same
valence-minus-two equation as \ref{cdp:st:C.2.UN.58}.
The finite leaf induction now has outward
increment \(h_e m_e\). This proves exactly
\eqref{cdp:eq:C.WT.4}, including product caps with
normal order one. Nonlinear terminal
forms give the same increment, by \eqref{cdp:eq:C.WT.10}.


\item  The root collar sends the connected cover
\(u^{m_e}=c\) to
\(v^{m_e}=s^{h_e m_e}/c\). These are the
same oriented integral two-cycle. Together
with the full secondary boundary matrix
of \ref{cdp:st:C.2.UN.50}, they identify every regular
elliptic fibre and every nonlinear core
with one integral class \(\gamma_s\).
The canonical period has order greater
than \(H_v\) on an elliptic vertex, and
has order exactly \(H_v\), with the
nonzero coefficient of (C.PT.10), at a
nonlinear leaf. This proves the height
test even on an arbitrary actual
smoothing.


\item  Choose a finite-equivariant relative
canonical generator: a nonconstant unit
in its finite character is removed by
the finite averaging of its logarithm.
Divide by the minimum power of \(s\).
The inverse charts used below carry this
form to the nonzero canonical frame on
the pruned core. Its finite character
there is exactly \(\chi=\det_{\C}D\).
The deck maps preserve the oriented
regular-fibre class, so
\[
 P(\zeta_Ns)=\chi P(s),\qquad
 P(s)=\int_{\gamma_s}\Omega_s .
 \tag{C.WT.14}\hypertarget{tag-017T}{}\label{cdp:eq:C.WT.14}
\]
If a nonlinear leaf is present, its
nonzero leading coefficient proves
\(\zeta_N^R=\chi\). Without such a leaf
one may choose a larger positive exponent
in this character class. This proves
\eqref{cdp:eq:C.WT.5} with the actual minimum-form
character, without inserting an
unproved parity assumption.


\end{enumerate}

\item \par\medskip\noindent\textbf{Equivariant inverse pruning and the exact core obstruction.}
Use the one-mark root-torsor classification
(C.PR.1)--(C.PR.8), or its identical application
to the supplied quasi-bundles in \ref{cdp:st:C.2.UN.61}--\ref{cdp:st:C.2.UN.62}.
After all children of a vertex have been
removed, its only possible multiple fibre
is its parent mark. Its actual torsor is
therefore a recorded diagonal Hopf pair,
with its full integer \(n\), or the product
ruled cap. At a Hopf step choose
\(mk-nl=1\). The actual inverse maps are
\[
 \begin{gathered}
 (u,X)\sim(\gamma^m u,\gamma^lX),\qquad
 p=X^m u^{-l},\qquad W=X^{-n}u^k,\\
                       u=p^nW^m,\qquad X=p^kW^l.
 \end{gathered}
 \tag{C.WT.15}\hypertarget{tag-017U}{}\label{cdp:eq:C.WT.15}
\]
They identify the punctured parent tube
with the specified product tube. The
parameter \(p\) is its actual base
coordinate, since its pullback is the
displayed multiple-fibre function.
The tensor and torsor maps therefore
extend across the new ordinary fibre,
with all other base charts unchanged.
The identity
\[
             d\log p\wedge d\log W
                         =d\log X\wedge d\log u
 \tag{C.WT.16}\hypertarget{tag-017V}{}\label{cdp:eq:C.WT.16}
\]
also proves the canonical-frame comparison
used in \eqref{cdp:eq:C.WT.14}: the old order \(m-1\)
becomes the nonzero frame \(dp\wedge
d\log W\). The normal-product neck
exponent \(h\) changes the order in the
family parameter, but does not change
this central inverse map.


\begin{enumerate}[label={$\langle2\rangle$\arabic*.},leftmargin=1.3em,labelsep=.5em]
\item  Carry out these operations simultaneously
on whole finite-group orbits. A stabilizer
of a connected boundary curve acts there
by a deck translation, as proved after
\eqref{cdp:eq:C.WT.7}. On the connected normal-root
elliptic cover its linear part is also
one, by the isogeny derivative argument
of \eqref{cdp:eq:C.WT.8}. It consequently fixes the
entire primitive fibre lattice. Choose
one primitive basis and the full integer
\(n\) on each orbit and transport it.
The maps \eqref{cdp:eq:C.WT.15}, their canonical
divisor sections and their inverses
then give the finite maps on the
replacement tubes. Their group law
holds on the punctured tubes and hence
on their holomorphic extensions.
This retains the stabilizer characters;
no independent choice of roots of
transition multipliers is made.
Disjoint replacements commute by their
actual atlases, so a group-invariant
postorder gives the required exact
equivariant pruning record.


\item  The root action cannot invert an edge:
its two endpoints lie over two distinct
original components, each fixed by the
root generator. It can, however, reflect
a cycle through vertices, interchanging
two different preimage curves over one
original boundary. Tree parent
orientations are preserved.


\item  For necessity no analytic contraction
of the given family is used.
At a fixed \(s\ne0\) the neck
\(uv=s^h\) has the same boundary
homomorphisms as \(uv=s\).
The full secondary matrix in \ref{cdp:st:C.2.UN.50}
shows that removal of the compact
leaf piece kills precisely the
primitive meridian \(\delta\) and
retains the two fibre generators.
The inverse product filling has
those identical maps. The actual
Van Kampen pushouts of (C.PR.9)
therefore preserve the nearby
fundamental group throughout.
For a cycle the final injected
semidirect product is
\(\Lambda\rtimes_M\Z\); its equality
with \(\Z^4\) forces \(M=I\).
For a genus-one core the injected
fibre lattice in (C.PT.12) has
commutator its integral Chern
vector, so that vector is zero.
The explicit Cech lattice
(C.PT.8) identifies this very
principal torsor with a complex
two-torus. This proves the fourth
test while retaining the full
analytic class.


\item  In the cycle case the exact two-mark
classification gives marked diagonal
Hopf or degree-zero ruled components.
Their regular fibres are the common
compact elliptic subgroup of their
open group \(\C^2/\Lambda\).
The finite maps on that group are
affine. For a diagonal Hopf pair,
an axis-preserving lift has form
\((zU,wV)\); contraction makes
the units \(U,V\) constant, exactly
as in \ref{cdp:st:C.2.UN.59}. An axis-interchanging
lift is the composition of such
a map with the coordinate
exchange. Marked ruled maps have
the actual affine logarithmic
lifts of \ref{cdp:st:C.2.UN.51}. Hence the finite
linear part \(D\) is a single
map on the common vector space.
For a torus core its derivative
is constant by periodicity,
so its actual finite action is
affine as well.


\end{enumerate}

\item \par\medskip\noindent\textbf{Signed cycle equations, including reflections.}
Orient the pruned cycle and let \(n_j\)
be the outgoing normal meridian,
\(h_j\) its retained neck exponent.
Let \(\epsilon=1\) or \(-1\) be the
orientation character of the root
generator, and \(\sigma\) its
permutation of vertices. The image
edge is \(k=\sigma(j)\) for
\(\epsilon=1\), and
\(k=\sigma(j+1)\) for \(\epsilon=-1\).
In actual logarithmic coordinates
write
\[
 \begin{gathered}
 g(j,v,s)=(\sigma j,Dv+\tau_j,\zeta_Ns),\\
 e_j(v,s)=v+b_j+l_j\log s,\qquad
       l_j=-\frac{h_jn_j}{2\pi i},\\
 Dn_j=\epsilon n_k,\qquad
 Dl_j=\epsilon l_k.
 \end{gathered}
 \tag{C.WT.17}\hypertarget{tag-017W}{}\label{cdp:eq:C.WT.17}
\]
For a reverse edge the comparison
is with \(e_k^{-1}\), not \(e_k\).
Taking \(\log\zeta_N=2\pi i/N\),
the actual compositions give
\[
 \begin{aligned}
 \tau_{j+1}+Db_j-\tau_j
       -\epsilon b_k-\epsilon l_k\,2\pi i/N
                                      &\in\Lambda,\\
 \sum_{a=0}^{N-1}D^{N-1-a}
                       \tau_{\sigma^a j}&\in\Lambda.
 \end{aligned}
 \tag{C.WT.18}\hypertarget{tag-017X}{}\label{cdp:eq:C.WT.18}
\]
For example when \(\epsilon=-1\),
\(g_{j+1}e_j=e_k^{-1}g_j\) gives
the first line directly. These
equalities retain every affine
constant of the actual central
action.


\begin{enumerate}[label={$\langle2\rangle$\arabic*.},leftmargin=1.3em,labelsep=.5em]
\item  Put \(b=\sum b_j\) and
\(u=-\sum h_jn_j=dp\), with \(p\)
primitive and \(d>0\). The exact
contraction-height calculation
(C.AD.3) gives
\(\ell(in_j)\) of one nonzero
sign, where
\(\ker\ell=\operatorname{span}_{\R}\Lambda\).
For ruled vertices this is the
actual elliptic quotient
calculation of \ref{cdp:st:C.2.UN.55}.
Thus \(\ell(iu)\ne0\) and
\(Du=\epsilon u\).
Choose the common elliptic
line and \(\C p\) as coordinates.
Then
\[
 \begin{gathered}
 W=\operatorname{span}_{\R}\Lambda=\C\times\R,\quad
 p=(0,1),\quad D=\operatorname{diag}(\rho,\epsilon),\\
 \Lambda=\langle(a_1,c_1),(a_2,c_2),p\rangle_{\Z},
 \qquad c_1,c_2\in\R,\qquad
        \operatorname{Im}(\overline{a_1}a_2)>0,\\
              (D-\epsilon I)b-\epsilon u/N\in\Lambda.
 \end{gathered}
 \tag{C.WT.19}\hypertarget{tag-017Y}{}\label{cdp:eq:C.WT.19}
\]
The last line is the sum of
\eqref{cdp:eq:C.WT.18}. The finite elliptic
linear part \(\rho\) has order
\(1,2,3,4\) or \(6\), by its
integral trace, as in \ref{cdp:st:C.2.UN.59}.
In particular
\(\det D=\rho\epsilon\), including
the sign at a reflected core.


\item  Adding \(c n_j\) to \(b_j\) on
a whole edge orbit, with the
same \(c\in\C\), preserves
\eqref{cdp:eq:C.WT.18}, since
\(Dn_j=\epsilon n_k\).
The second coordinates of all
the \(n_j\) have the same
nonzero sign. Such an orbit
therefore allows the second
coordinate of \(b\) to be
set to \(i\). This changes
only a normal-product constant:
the translation along the
normal meridian restricts
to zero on the double curve.
It preserves the specified
elliptic gluing maps and the
entire finite action.


\end{enumerate}

\item \par\medskip\noindent\textbf{The finite-equivariant compact core.}
For a cycle use the coordinates
in \eqref{cdp:eq:C.WT.19}; for a torus choose
a finite-invariant Hermitian
metric and an orthonormal
complex basis. In either case
\(D\) is unitary and
\[
 D J\overline D^{-1}=(\det D)J,
       \qquad J=\begin{pmatrix}0&1\\-1&0\end{pmatrix}.
 \tag{C.WT.20}\hypertarget{tag-0180}{}\label{cdp:eq:C.WT.20}
\]
This follows by multiplying
the two-by-two matrices, or
from the alternating form
on a two-dimensional space.
With \(R\) as in \eqref{cdp:eq:C.WT.5} set
\[
 V_s(w)=w+s^R J\overline w,\qquad
             D V_s=V_{\zeta_Ns}D .
 \tag{C.WT.21}\hypertarget{tag-0181}{}\label{cdp:eq:C.WT.21}
\]
For a torus core the four
real period vectors \(V_s\Gamma\)
remain independent because
\(V_s\) has real determinant
\((1+|s|^{2R})^2\).
Transport an affine constant
to the target frame:
\(\tau(s)=V_{\zeta_Ns}\tau\).
All finite lattice and order
equalities are then preserved
exactly, so these are an actual
finite action on the torus
family.


\begin{enumerate}[label={$\langle2\rangle$\arabic*.},leftmargin=1.3em,labelsep=.5em]
\item  For a cycle transport
\[
 b_j(s)=V_s b_j,\qquad
 \tau_j(s)=V_{\zeta_Ns}\tau_j,\qquad
 B_s=V_s b+\frac{d\log s}{2\pi i}V_s p .
 \tag{C.WT.22}\hypertarget{tag-0182}{}\label{cdp:eq:C.WT.22}
\]
The correction
\(l_j\,2\pi i/N=-h_jn_j/N\)
in \eqref{cdp:eq:C.WT.18} is a real-rational
period vector. Applying \(V_s\)
therefore carries the signed
equations to the identical
equations in the target frame.
Summing gives
\[
                 D B_s-\epsilon B_{\zeta_Ns}
                               \in V_{\zeta_Ns}\Lambda .
 \tag{C.WT.23}\hypertarget{tag-0183}{}\label{cdp:eq:C.WT.23}
\]
The finite action consequently
preserves the actual rank-four
nearby period lattice. The
global map itself can be
written on the actual middle
charts. Choose a cut and put
\[
 \begin{gathered}
 \theta_j(s)=\sum_{a<j}
      \bigl(b_a(s)+l_a(s)\log s\bigr),\qquad
 \eta=v_j-\theta_j(s),\\
 \eta\longmapsto
 D\eta+D\theta_j(s)+\tau_j(s)
                         -\theta_{\sigma j}(\zeta_Ns),
 \qquad l_a(s)=-h_aV_s(n_a)/(2\pi i).
 \end{gathered}
 \tag{C.WT.23a}\hypertarget{tag-0184}{}\label{cdp:eq:C.WT.23a}
\]
Changing an adjacent chart
changes this expression by
the actual lattice element
in \eqref{cdp:eq:C.WT.18}. Changing the
cut adds \(B_s\) on the source
and \(\epsilon B_{\zeta_Ns}\)
on the target, so \eqref{cdp:eq:C.WT.23}
proves that the expression
is unchanged as a map of
the middle torus. A logarithm
change adds integral normal
periods. The second equation
of \eqref{cdp:eq:C.WT.18} proves its
\(N\)-th-order relation.
Thus this is an actual
global finite holomorphic
action, including reflections,
not only an invariant lattice.
The
normal maps are the literal
\(uv=s^{h_j}\) maps; reversing
an edge exchanges their two
coordinates. The moving
Hopf and ruled frame maps
of \ref{cdp:st:C.2.UN.57}--\ref{cdp:st:C.2.UN.59} give these
actual charts also for a
reflection through a vertex.


\item  The core is compact over the
whole parameter sphere. In
the cycle case its determinant
is a positive constant times
\[
 (1+r^{2R})
 \left[1+r^{2R}
    +\frac d{2\pi}(r^{2R}-1)\log r\right],
                \qquad r=|s|>0 .
 \tag{C.WT.24}\hypertarget{tag-0185}{}\label{cdp:eq:C.WT.24}
\]
Each factor is positive.
This is the exact computation
of \ref{cdp:st:C.2.UN.57} with \(V_s\) replaced
by \(V_{s^R}\), keeping
\(d\log s\) unchanged.
It gives compact connected
two-tori on every punctured
fibre. Finite core and neck
collars give the proper
filling at zero. No torus
rank is inferred merely
from the fundamental group.


\item  At infinity put \(z=s^{-1}\)
and divide the vector frame
by \(s^R\). The periods become
\[
              W_z(w)=J\overline w+z^Rw,\qquad
              D_{\infty}=(\det D)^{-1}D
                       =J\overline D J^{-1}.
 \tag{C.WT.25}\hypertarget{tag-0186}{}\label{cdp:eq:C.WT.25}
\]
The base character is
\(\zeta_N^{-1}\). An affine
constant becomes
\(J\overline\tau+(\det D)^{-1}z^R\tau
=W_{\zeta_N^{-1}z}(\tau)\).
The normal meridians are
\(-J\overline{n_j}\); changing
\(\log s\) to \(-\log z\)
gives the conjugate core
with the same positive neck
exponents. Thus the infinity
maps, including every affine
constant and reflection, are
the actual complex conjugates
of the central core maps.
This provides the second
proper compact end.


\end{enumerate}

\item \par\medskip\noindent\textbf{The precise finite jet gauge at a surgery orbit.}
The common regular elliptic fibre has a
primitive period basis \(\lambda,\mu=c\lambda\),
\(\operatorname{Im}c>0\). In the core frame
its exact determinant is
\[
 \det_{\C}(V_s\lambda,V_s\mu)
              =2i\operatorname{Im}(c)\|\lambda\|^2s^R .
 \tag{C.WT.26}\hypertarget{tag-0187}{}\label{cdp:eq:C.WT.26}
\]
Consequently the actual rank-two period
cover near a marked fibre has the
semi-torus presentation
\[
 (p,w,s)\longmapsto(p+\delta(s),\gamma(s)w,s),
 \qquad \delta(s)=s^{\ell} d(s),\quad d(0)\ne0,
                     \quad|\gamma(0)|<1 .
 \tag{C.WT.27}\hypertarget{tag-0188}{}\label{cdp:eq:C.WT.27}
\]
Initially \(\ell=R\). The exact primitive
basis change (C.PR.13) multiplies \(\delta\)
by a holomorphic unit and retains its
order. Every later chart below has this
same presentation.


\begin{enumerate}[label={$\langle2\rangle$\arabic*.},leftmargin=1.3em,labelsep=.5em]
\item  Consider the stabilizer of one surgery
curve, of order \(d_0\), with base
character \(\xi\). This is the stabilizer
order, not the order of the whole
graph group when that curve moves.
As proved after \eqref{cdp:eq:C.WT.15}, its linear
action on the actual fibre lattice
is the identity. On the universal
logarithmic cover its finite lift
therefore has an upper-triangular
affine form
\[
 (p,V,s)\longmapsto
    (a(s)p+b(s),\,V+c(s)p+e(s),\,\xi s),
                     \qquad V=\log w .
 \tag{C.WT.28}\hypertarget{tag-0189}{}\label{cdp:eq:C.WT.28}
\]
The lift may have \(d_0\)-th power
equal to the \(k\)-th infinite deck
translation, with an integral
angular deck translation as well.
Both integers are retained from the
actual root action. The \(p\)-part
of that equality is
\[
 \prod_{j=0}^{d_0-1}a(\xi^j s)=1,\qquad
                 g^{d_0}(p)=p+k\delta(s).
 \tag{C.WT.29}\hypertarget{tag-018A}{}\label{cdp:eq:C.WT.29}
\]
The norm-one unit cocycle for \(a\)
is made constant \(a_0=a(0)\) by
a holomorphic unit gauge, equal
to one at zero. Explicitly take
a logarithm of \(a/a_0\), whose
finite orbit sum is zero, and
use the same finite weighted
sum as in (C.HT.4). Its
coboundary is that logarithm.
Commutation with the infinite
period then gives
\[
 \delta(\xi s)=a_0\delta(s),\qquad
 \sum_{j=0}^{d_0-1}a_0^{d_0-1-j}b(\xi^j s)
                                      =k\delta(s).
 \tag{C.WT.30}\hypertarget{tag-018C}{}\label{cdp:eq:C.WT.30}
\]


\item  For any coefficient \(b_j s^j\)
of order \(j<\ell\), the sum
in \eqref{cdp:eq:C.WT.30} is zero unless
\(\xi^j=a_0\); in that resonant
case it is
\(d_0 a_0^{d_0-1}b_j s^j\),
so \(b_j=0\). Every nonresonant
coefficient is removed by
\(p\mapsto p+c_j s^j\), since
its change is
\((\xi^j-a_0)c_j s^j\).
Thus, whenever \(\ell\ge hm\),
one makes
\[
                               b(s)=O(s^{hm})
 \tag{C.WT.31}\hypertarget{tag-018D}{}\label{cdp:eq:C.WT.31}
\]
by an actual finite polynomial
jet gauge. This preserves the
central marked point and every
resonant coefficient. For an
intermediate diagonal Hopf
component a constant nontrivial
affine action is first centered
at its finite fixed point;
this is its actual triangular
Hopf automorphism and is
recorded on all remaining
marks.


\item  In particular the nonlinear
translation is never discarded.
On a nonlinear Hopf universal
cover put \(U=z/w^m,V=\log w\).
Its infinite generator is
\((U,V)\mapsto(U+c,V+\log\beta)\),
with \(c\ne0\). A parent-axis-
preserving lift commuting with
a positive central contraction
has \(U\mapsto U+b\):
an affine coefficient \(a\)
must satisfy \(ac=c\).
If its \(d_0\)-th power is the
\(k\)-th infinite generator
times a finite angular element,
that element fixes \(U\), so
\[
                                d_0b=kc.
 \tag{C.WT.32}\hypertarget{tag-018E}{}\label{cdp:eq:C.WT.32}
\]
This is exactly the term which
survives \eqref{cdp:eq:C.WT.30} at a terminal
insertion. Indeed if \(\ell=hm\),
then \(a_0=\xi^{\ell}\) and the
leading coefficient there is
\(b_{\ell}=k a_0d(0)/d_0\).
After dividing by \((\xi s)^{hm}\)
it is \(k d(0)/d_0\), as required
by \eqref{cdp:eq:C.WT.32}. For a diagonal
component the corresponding
equation is
\(b(1+a+\cdots+a^{d_0-1})=0\);
it forces \(b=0\) when \(a=1\)
and permits exactly the finite
fixed-point gauge otherwise.
This proves the full central
triangular comparison.


\end{enumerate}

\item \par\medskip\noindent\textbf{The literal compressed weighted insertion.}
For a recorded Hopf step of
normal order \(m\), full secondary
integer \(n\), and neck exponent
\(h\), use the actual normalized
pullback, by \(t=s^h\), of the
weight-\((m,1)\) base modification
in \ref{cdp:st:C.2.UN.60}. Its root chart is
\[
       p=x^m Z,\qquad xW=s^h,\qquad (Z,W)\ne(0,0).
 \tag{C.WT.33}\hypertarget{tag-018F}{}\label{cdp:eq:C.WT.33}
\]
Retain the pullback of the old
tautological line \(\OO(nE_{\rm old})\).
Its canonical local section is
\(x^n\); since \(x=s^h/W\) at a
generic point of the new component,
\(\operatorname{div}(x)=hE_{\rm new}\).
Thus the pulled line has coefficient
\(nh\), not \(n\), on this generic
divisor. Replacing it by
\(\OO(nE_{\rm new})\) would change
the actual root action and is not
performed here.


\begin{enumerate}[label={$\langle2\rangle$\arabic*.},leftmargin=1.3em,labelsep=.5em]
\item  Choose the recorded \(\beta(s)\)
with \(\beta(s)^n=\gamma(s)\).
The nonzero line has the free
finite chart
\[
       (Z,W,x)\longmapsto
                   (\eta^mZ,\eta W,\eta^{-1}x),
                           \qquad\eta\in\mu_n.
 \tag{C.WT.34}\hypertarget{tag-018G}{}\label{cdp:eq:C.WT.34}
\]
It is free because
\(\gcd(m,n)=1\), including
where one of \(Z,W\) vanishes.
The infinite deck map is exactly
\[
 \begin{split}
 (Z,W,x,s)\longmapsto\bigl(
   &\beta(s)^m[Z+d(s)s^{\ell-hm}W^m],\\
   &\beta(s)W,\ \beta(s)^{-1}x,\ s\bigr).
 \end{split}
 \tag{C.WT.35}\hypertarget{tag-018H}{}\label{cdp:eq:C.WT.35}
\]
It preserves \(xW=s^h\) and
changes \(p=x^mZ\) by precisely
\(d(s)s^{\ell}=\delta(s)\).
Every term has the same finite
character as \(Z\) in
\eqref{cdp:eq:C.WT.34}. The central child
is nonlinear exactly when
\(\ell=hm\), and diagonal
when \(\ell>hm\). Its recorded
nonzero nonlinear coefficient
is matched by the actual
rescaling of \(Z\); the finite
translation in \eqref{cdp:eq:C.WT.32}
scales with it.


\item  Away from its parent curve put
\(U=Z/W^m\) and use the finite
coarse fibre coordinate \(W^n\).
The exact new deck is
\[
            (U,W^n)\longmapsto
            (U+\delta(s)/s^{hm},\,\gamma(s)W^n).
 \tag{C.WT.36}\hypertarget{tag-018I}{}\label{cdp:eq:C.WT.36}
\]
This is the required remaining
order \(\ell-hm\), at every
other marked child fibre.
There is no change of the
given analytic torsor outside
the one recorded tube.


\item  The finite action also lifts
exactly. The gauge \eqref{cdp:eq:C.WT.31}
makes its base map preserve
the weighted ideal
\((p,s^{hm})\), and hence its
normalized weighted modification.
The canonical divisor section
gives the action on the pulled
tautological line. In a chart,
an \(n\)-th root of a nonzero
multiplier may be chosen;
different choices are precisely
the \(\mu_n\)-action \eqref{cdp:eq:C.WT.34}.
Thus the coarse maps are
independent of the choice.
Their group law holds on the
unchanged open set, and then
on the normal chart. This
proves the actual finite action,
including \(g^{d_0}=G^k\)
on a period lift, without
a new linearization cocycle.


\item  More explicitly the new
transverse coordinate is
\(U=p/s^{hm}\). Its finite
map is
\[
             U\longmapsto
       \frac{a_0}{\xi^{hm}}U
                 +\frac{b(s)}{\xi^{hm}s^{hm}} .
 \tag{C.WT.37}\hypertarget{tag-018J}{}\label{cdp:eq:C.WT.37}
\]
On logarithmic fibre coordinates,
the substitution adds \(nh\log s\);
the new off-diagonal coefficient
is \(c(s)s^{hm}\). Both are
holomorphic in the new chart.
Thus subsequent stabilizers
again have the exact affine
form \eqref{cdp:eq:C.WT.28}, with the
remaining displacement in
\eqref{cdp:eq:C.WT.36}. This proves the
induction for the complete
triangular action, not just
its action on homology.


\item  A product ruled cap uses the
actual normalized blowup of
\((p,s^h)\). Its two charts are
\[
                  p=s^hU,\qquad\text{and}\qquad xW=s^h .
 \tag{C.WT.38}\hypertarget{tag-018K}{}\label{cdp:eq:C.WT.38}
\]
The height test has \(\ell>h\);
the central projective normal
action is therefore the identity
for the infinite deck map.
The cap is the specified
product \(F\times\PP^1\), and
its child displacement is
\(\delta(s)/s^h\).
The same finite jet gauge,
now through order \(h-1\),
gives its actual equivariant
blowup; its central affine
map is the recorded finite
projectivized normal map.


\item  Each step subtracts exactly
\(h_em_e\) from the period
order, with \(m_e=1\) for
a product cap. The height
test consequently supplies
the required holomorphic
inequalities at every step
and stops with zero remaining
order only at a nonlinear
terminal leaf. Reverse the
complete group-invariant
pruning record. Equations
\eqref{cdp:eq:C.WT.15}, \eqref{cdp:eq:C.WT.33}--\eqref{cdp:eq:C.WT.38}
restore every original root
component, all its child marks,
and its full torsor and
linear and translation
gluing maps.


\end{enumerate}

\item \par\medskip\noindent\textbf{Properness of the compressed replacements.}
The affine-base estimate must
retain \(h\). Put
\[
 \rho(p,s)=
       (|p|^2+|s|^{2hm})^{1/(2m)}
       =|x|(|Z|^2+|W|^{2m})^{1/(2m)} .
 \tag{C.WT.39}\hypertarget{tag-018L}{}\label{cdp:eq:C.WT.39}
\]
It satisfies \(\rho\ge|s|^h\).
Since \(\ell\ge hm\), the
translation by \(j\delta(s)\)
distorts this norm by at most
a uniform polynomial in
\(1+|j|\), in both directions.
The pulled line has the same
local fibre \(x^{-n}\) as
before base change. Its norm
distortion is a polynomial
of exponent at most \(n/m\).
Uniform exponential contraction
by \(\gamma(s)^j\) dominates
that polynomial, so a compact
norm annulus meets every orbit
and the action is proper
and free. The finite-group
orbit of a chosen norm gives
the same estimates uniformly.


\begin{enumerate}[label={$\langle2\rangle$\arabic*.},leftmargin=1.3em,labelsep=.5em]
\item  For comparison with the old
annular chart, the required
normalizing deck exponent is
\(O(1+|\log\rho|)\). Its base
displacement is bounded by
\[
              C|\delta(s)|(1+|\log\rho|).
 \tag{C.WT.40}\hypertarget{tag-018N}{}\label{cdp:eq:C.WT.40}
\]
If \(\rho<1\), use
\(\rho\ge|s|^h\). If \(\rho\ge1\),
an old compact base bound
gives
\(|p_{\rm new}|\le B+
C|\delta(s)|\log(2+|p_{\rm new}|)\),
so \(p_{\rm new}\) is uniformly
bounded. The required new
annular domain is consequently
compact over every old compact
domain, even though the
universal affine base is
not compact.


\item  Finally, on a compact Hopf
fundamental domain, \(x\to0\)
implies \(|s|^h=|xW|=O(|x|)\).
The old base coordinate after
normalization is
\[
          O(|x|^m+|s|^\ell|\log|x||)\longrightarrow0.
 \tag{C.WT.41}\hypertarget{tag-018O}{}\label{cdp:eq:C.WT.41}
\]
Every cluster point is therefore
on the single removed compact
elliptic curve. The closed
overlap graphs and nested
compact annuli used in \ref{cdp:st:C.2.UN.60}
now give the same Hausdorff
proper replacement, retaining
the literal \(A_{h-1}\) neck.
Product caps are proper
normalized blowups. Distinct
central marks have disjoint
finite-group-invariant tubes;
their finite simultaneous
replacements therefore glue.
There are only finitely many
stages.


\end{enumerate}

\item \par\medskip\noindent\textbf{The smooth weighted quotient and the compact second end.}
The constructed root family
has only its recorded neck
singularities. At a generic
point of a strict component
over \(S_i\), the inertia has
order \(N/m_i\), fixes the
component and acts on its
normal parameter by that
primitive character. Its
quotient is smooth and the
full fibre multiplicity is
exactly \(m_i\).


\begin{enumerate}[label={$\langle2\rangle$\arabic*.},leftmargin=1.3em,labelsep=.5em]
\item  At an original double curve
of multiplicities \(a=gA,b=gB\),
the stabilizer of a point of
one lifted curve has order
\(hAB=N/g\). Its subgroup
\(\mu_h\) fixes \(u,v\) and
acts primitively on \(s\).
Quotienting \(uv=s^h\) first
by it gives the smooth
\((u,v)\)-chart. The residual
cyclic group of order \(AB\)
has the two independent
characters of orders \(B,A\).
Since \(\gcd(A,B)=1\), its
invariant ring is exactly
\[
 \bigl(\C\{z,u,v,s\}/(uv-s^h)\bigr)^{\,I}
       =\C\{z,u^B,v^A\},\qquad
                  t=(u^B)^a(v^A)^b .
 \tag{C.WT.42}\hypertarget{tag-018P}{}\label{cdp:eq:C.WT.42}
\]
Here \(I\) is that inertia
group. The curve coordinate
\(z\) is invariant, because
the actual curve map is
\'{e}tale. Finite character
averaging gives coordinates
of the form \(uU,vV,z+\cdots\)
with their original characters.
Their product is \(s^h\)
times an invariant unit;
dividing one coordinate by
that unit keeps both the
characters and the literal
equation \(uv=s^h\).
This proves \eqref{cdp:eq:C.WT.42} for the
actual finite action, not
only its tangent representation.


\item  There is no edge inversion
to add an unrecorded quotient
of this chart. At interior
points the strict component
maps are \'{e}tale after the
generic normal inertia, so
there are no other quotient
singularities. The resulting
total space near zero is
smooth and its weighted
central support is the
original one.


\item  For completeness the weighted
normal line is retained as
well. On the root complement
the original normal coordinate
is \(v_i=s^{N/m_i}/\xi\).
The character eigensheaves
of the finite component map
extend by \eqref{cdp:eq:C.WT.7}; explicitly
their \(k\)-th pieces are
\[
 \mathcal L_i^{\,k}\otimes
  \OO_{S_i}\!\left(
       \sum_j\left\lfloor\frac{k m_j}{m_i}\right\rfloor J_{ij}
             \right),\qquad 0\le k<m_i .
 \tag{C.WT.43}\hypertarget{tag-018Q}{}\label{cdp:eq:C.WT.43}
\]
This follows by dividing
\(\xi^k\) by its integral
normal zero order in that
chart. The retained full
section and its orders
recover \(\mathcal L_i\)
and the specified edge line
maps; no reduction of an
integer modulo \(m_i\)
has occurred. The chosen
normal-product constants
are the only freedom stated
in the theorem.


\item  Carry out the conjugate
pruning record at infinity
as well. Formula \eqref{cdp:eq:C.WT.25}
has base character
\(\zeta_N^{-1}\) and canonical
character \((\det D)^{-1}\);
its primitive elliptic period
again has exact order \(R\).
Thus the conjugate height
and finite-jet comparisons
hold with precisely the
conjugate characters.
All charts, including the
compressed necks and
\eqref{cdp:eq:C.WT.42}, are the actual
conjugates of those at zero.
The quotient is therefore
smooth at that end too and
has the conjugate prescribed
weighted support. The two
sets of replacements have
disjoint parameter
neighbourhoods.


\item  The compact core and the
proper replacements give
a compact normal root
threefold. Its finite
quotient is smooth by the
charts just proved. Away
from zero and infinity
the base action is free;
one fibre over each orbit
maps biholomorphically to
the quotient fibre. They
are consequently connected
complex two-tori. Each
original component and
each double elliptic curve
has Euler number zero,
so both reduced central
supports have Euler zero.
This completes necessity,
sufficiency and the compact
comparison, including all
finite triangular actions.



\end{enumerate}

\end{enumerate}
\end{proof}






\section{Actual contractions, canonical reduction and finite generation}\hypertarget{tag-0347}{}\label{sec:II:11}


\begin{proposition}[Maximum components without a minimality assumption]\hypertarget{tag-01EF}{}\label{cdp:st:C.2.UN.69}
\StatementRecord{cdp:st:C.2.UN.69}{proved}{01EF}{\SeriesRef{cdp:st:C.2.UN.64}, \SeriesRef{cdp:st:C.2.UN.50}, \SeriesRef{cdp:st:C.2.UN.58}, \SeriesRef{cdp:st:C.2.UN.19}, \SeriesRef{cdp:st:C.2.UN.15}, \SeriesRef{cdp:st:C.2.UN.31}, \SeriesRef{cdp:eq:C.NM.4}, \SeriesRef{cdp:eq:C.NM.7}, \SeriesRef{cdp:eq:C.NM.6}, \SeriesRef{cdp:st:C.2.UN.66}, \SeriesRef{cdp:eq:C.NM.9}, \SeriesRef{cdp:eq:C.NM.3}, \SeriesRef{cdp:eq:C.NM.8}, \SeriesRef{cdp:eq:C.NM.10}}

Let $f:U\to\Delta$ be a proper degeneration of connected complex
two-tori from a smooth threefold. Let its full central fibre
$S=\sum_vT_v$ be reduced and simple normal crossing, without triple
points. No minimality assumption is made on any $T_v$.
Write $k_v$ for the actual canonical evaluation orders.

At a local maximum $T_v$ of these orders, precisely one of the
following descriptions applies:
\begin{enumerate}
\item Every adjacent component has order $k_v$.
\item The component is an unblown terminal ruled surface meeting
its parent in a section, with order drop one, or an unblown
terminal plane meeting its parent in a line, with order drop two.
These are the actual contractions of \ref{cdp:st:C.2.UN.64}.
\item The component is an unblown terminal Hopf surface, with
the exact axis normal order and positive canonical drop of
\ref{cdp:st:C.2.UN.50} and \ref{cdp:st:C.2.UN.58}.
\end{enumerate}
In the second alternative the ruled base has positive genus.
Thus, when the orders are not all equal, a global maximum is
one of the two kinds of strict-drop leaves.

Moreover every rational double curve can be removed by a finite
sequence of the actual plane-line contractions of \ref{cdp:st:C.2.UN.64}.
At each step the other double curves and their neighbourhoods
are unchanged. If $r_0$ is the number of connected rational
double curves of the original support, the resulting actual
torus degeneration has no rational double curve and satisfies
\[
             e(S)=e(S_{\rm nr})+2r_0 .
\tag{C.NM.1}\hypertarget{tag-00SR}{}\label{cdp:eq:C.NM.1}
\]
There are exactly $r_0$ point contractions. Reversing their
recorded ordinary blowups recovers the original support,
including its original smooth threefold neighbourhood.
In particular the rational double edges are removed as terminal
edges of the actual incidence graph; none belongs to a cycle.
\end{proposition}



\begin{proof}

\begin{enumerate}[label={$\langle1\rangle$\arabic*.},leftmargin=*]
\item \par\medskip\noindent\textbf{The exact restrictions at an edge.}
Canonical evaluation, adjunction and the actual equation $xy=t$
give the same line maps as (C.MR.3)--(C.MR.4), without a
minimality assumption:
\[
 \begin{split}
 K_{T_v}&=\operatorname{div}(\omega_v)
       =\sum_{e=(v,w)}(k_w-k_v-1)C_e,\\
 N_{C_e/T_v}\otimes N_{C_e/T_w}&\simeq\OO_{C_e},\\
 K_{C_e}&\simeq N_{C_e/T_v}^{\,k_w-k_v}.
 \end{split}
\tag{C.NM.2}\hypertarget{tag-00ST}{}\label{cdp:eq:C.NM.2}
\]
The last equality uses that the other boundaries are disjoint
from $C_e$. In particular, for $g_e=g(C_e)$,
\[
                  (k_w-k_v)C_e^2=2g_e-2.
\tag{C.NM.3}\hypertarget{tag-00SU}{}\label{cdp:eq:C.NM.3}
\]
These formulas retain the actual meromorphic section and
dual-normal trivialization.

\par\medskip\noindent\textbf{A positive-square rational curve on a component.}
We shall use the following implication for any component, not
only a maximum:
\[
 \text{a smooth rational curve of positive square on }T
                         \quad\Longrightarrow\quad T
                         \text{ is rational}.
\tag{C.NM.4}\hypertarget{tag-00SV}{}\label{cdp:eq:C.NM.4}
\]
Here is the comparison with its actual minimal-model map.
For a nonrational ruled minimal model, every rational curve
maps to a point of the positive-genus ruling base. A surviving
such curve is a ruling fibre, of square zero; its strict
transform has nonpositive square. Every exceptional curve
also has negative square. On the Hopf minimal model supplied
by \ref{cdp:st:C.2.UN.19} there are no rational curves, since a map from
$\PP^1$ lifts to $\C^2\setminus\{0\}$ and has constant
coordinate functions. Hence a rational curve on a blowup of
that model is exceptional and has negative square.

\item The remaining actual minimal models have Kodaira dimension
zero or one by \ref{cdp:st:C.2.UN.15}. In this multiplicity-one setting \ref{cdp:st:C.2.UN.31}
gives a torus or the smooth-multiple elliptic model. Its
canonical line is numerically nonnegative on every curve:
on a torus it is trivial; in the elliptic case its positive
power is the pullback of the actual base line of degree
\[
                   2g(B)-2+\sum_b(1-1/m_b)\ge0.
\tag{C.NM.5}\hypertarget{tag-00SW}{}\label{cdp:eq:C.NM.5}
\]
The inequality follows from the pluricanonical pushdown and
the nonnegative Kodaira dimension. For its point-blowup map
$\pi:T\to T_{\min}$, one has
$K_T=\pi^*K_{T_{\min}}+\sum b_EE$, with $b_E>0$.
A nonexceptional rational curve $C$ would therefore have
$K_T\cdot C\ge0$. Positive square instead gives
$K_T\cdot C=-2-C^2<0$ by adjunction. An exceptional curve
has negative square. This proves \eqref{cdp:eq:C.NM.4} through the
actual minimal-model morphism.

\item If $C\subset T$ is rational with square $s=1$ or $2$,
blow up $s$ distinct smooth points of $C$. The strict
transform $C'$ has square zero. Since $T$ is rational,
$H^1(\OO)=0$ on this blown-up surface. The restriction
sequence for $\OO(C')$, with
$\OO_{C'}(C')\simeq\OO_{\PP^1}$, gives a second section
nonzero everywhere on $C'$. Together with the canonical
section cutting out $C'$ it gives an actual basepoint-free
pencil
\[
                         q:T'\longrightarrow\PP^1.
\tag{C.NM.6}\hypertarget{tag-00SX}{}\label{cdp:eq:C.NM.6}
\]
There are no basepoints outside $C'$ because the canonical
section is nonzero there. The fibre $C'$ is reduced and
connected. In the Stein factorization a finite map of
degree greater than one would either give more than one
point above this value or a multiplicity greater than one
on $C'$. Thus $q$ has connected fibres. Adjunction gives
genus zero on a smooth fibre. The ordinary relative
minimal-model contractions identify this rational fibration
with a projective-line bundle followed by point blowups
in fibres; in particular every irreducible vertical curve
is rational. Every original curve disjoint from $C$
is untouched by the selected blowups and has intersection
zero with $q^*\OO_{\PP^1}(1)=\OO(C')$, so it is vertical.
This supplies the precise pencil used below.

\par\medskip\noindent\textbf{A maximum with a rational boundary is already minimal.}
At a local maximum put $a_e=k_v-k_w+1\ge1$. Then
\[
             -K_T=\sum_e a_eC_e,\qquad
             (1-a_e)C_e^2=2g(C_e)-2.
\tag{C.NM.7}\hypertarget{tag-00SY}{}\label{cdp:eq:C.NM.7}
\]
The boundary is nonempty in the reducible case. A nonzero
pluricanonical section multiplied by the corresponding
power of this anticanonical section would be a nonzero
holomorphic function having zeros on a compact connected
surface. Consequently $\kappa(T)=-\infty$.

\item For a rational boundary, \eqref{cdp:eq:C.NM.7} gives exactly
\[
                     (a_e,C_e^2)=(3,1)\ \text{or}\ (2,2).
\tag{C.NM.8}\hypertarget{tag-00SZ}{}\label{cdp:eq:C.NM.8}
\]
Equation \eqref{cdp:eq:C.NM.4} first makes $T$ rational; no rational
pencil is used before this conclusion. The actual pencil
\eqref{cdp:eq:C.NM.6} then makes every other disjoint boundary rational.
Every such boundary would again have positive square by
\eqref{cdp:eq:C.NM.7}, whereas the Hodge index theorem permits at most
one positive-square class among mutually disjoint curves.
Thus there is only one boundary $C$, and $K_T=-aC$.
Any smooth rational $(-1)$-curve $E$ would satisfy
$a(C\cdot E)=1$, which is impossible for $a=2$ or $3$.
The surface is therefore minimal. The terminal-pair
classification and actual maps in \ref{cdp:st:C.2.UN.64} give its ruled
section or plane-line contraction.

\par\medskip\noindent\textbf{All other blowdowns occur only at coefficient-one elliptic boundaries.}
Suppose now that no boundary is rational. A $(-1)$-curve
$E$ is not in the boundary, and adjunction gives
\[
                     1=(-K_T)\cdot E
                              =\sum_e a_e(C_e\cdot E).
\tag{C.NM.9}\hypertarget{tag-00T0}{}\label{cdp:eq:C.NM.9}
\]
Every summand is a nonnegative integer. Exactly one is
one: $E$ meets one boundary transversely once, that
boundary has coefficient one, and $E$ misses all others.
The coefficient-one boundary is elliptic by \eqref{cdp:eq:C.NM.7}.
Its ordinary surface blowdown keeps this boundary smooth
and keeps all boundary curves disjoint. Pushforward of
the actual meromorphic canonical section gives the same
canonical divisor with the same coefficients on the
new surface. This is a statement about the surface
morphism; no ambient contraction of $E$ is invoked.
Repeat to the minimal model.

\item The minimal possibilities with \eqref{cdp:eq:C.NM.7} are exactly those
computed in (C.MR.6)--(C.MR.7) and \ref{cdp:st:C.2.UN.64}. On a ruled
surface, intersection with a ruling fibre gives the three
partitions of two: one coefficient-two section, one
coefficient-one bisection, or two coefficient-one sections.
The connected elliptic bisection over an elliptic base is
excluded by the actual residue/trace map (C.EL.7)--(C.EL.8).
On the plane the possibilities are a coefficient-three
line or a coefficient-one cubic. On an elliptic Hopf
surface the invariant relative differential gives
$\sum (k_w-k_T)/m_e=\operatorname{val}(T)-2$ with every
summand nonpositive; thus it is either a strict-drop
one-boundary leaf or a two-boundary equal-height component.
The nonlinear Hopf axis formula gives the same strict
leaf alternative, including its exact secondary framing.
For an algebraic-dimension-zero diagonal Hopf surface,
the two-axis character equation and $|\alpha|,|\beta|<1$
make both coefficients one; a single axis would impose
an integral resonance. These are actual marked pairs
by \ref{cdp:st:C.2.UN.50}, \ref{cdp:st:C.2.UN.58} and the proof of \ref{cdp:st:C.2.UN.66}.

\item In particular whenever a coefficient is greater than
one, there is just one boundary, with that coefficient.
No coefficient-one boundary then exists. Since the
surface blowdowns above preserve every boundary and
its coefficient, \eqref{cdp:eq:C.NM.9} shows that no such blowdown
could have taken place. The original strict-drop maximum
was already minimal. In all remaining cases every
coefficient is one and every neighbour has the same
height. A component at the global maximum with this
property cannot be adjacent to a strict-drop maximum
leaf. Its equal-height neighbours have the same property,
so connectedness propagates it to the whole support.
This proves the maximum assertion, except for the
rational ruled-section exclusion proved next.

\par\medskip\noindent\textbf{A rational section centre cannot lie in the smooth torus-fibre locus.}
In a ruled-section contraction of \ref{cdp:st:C.2.UN.64} with base
$B\simeq\PP^1$, the actual descended family is a
submersion near the compact centre $B$. Its containing
central component $P'$ is smooth there and
\[
                    N_{B/P'}\simeq K_B\simeq\OO_{\PP^1}(-2).
\tag{C.NM.10}\hypertarget{tag-00SS}{}\label{cdp:eq:C.NM.10}
\]
Choose a relatively compact tubular neighbourhood of
$B$ inside that smooth central surface and a slightly
larger compact neighbourhood in the submersion locus.
A smooth metric gives horizontal lifts of the two
coordinate vector fields on the disc. For sufficiently
small time their flows exist on the whole smaller
neighbourhood and embed it in an actual nearby fibre.
The induced diffeomorphism preserves the oriented
normal bundle and its Euler number. Thus the nearby
two-torus would contain an embedded sphere of
self-intersection minus two. Its universal cover is
the contractible space $\C^2$, so this sphere map is
nullhomotopic, its homology class is zero, and its
self-intersection is zero. This contradiction excludes
the rational ruled-section alternative.

\par\medskip\noindent\textbf{Following rational edges gives actual point contractions.}
Take a rational double curve and orient it toward its
higher-height endpoint. Equation \eqref{cdp:eq:C.NM.3} makes the
height difference nonzero; at that higher endpoint
its square is positive, equal to one or two. By
\eqref{cdp:eq:C.NM.4} the endpoint surface is rational. The pencil
\eqref{cdp:eq:C.NM.6} makes every other boundary on it rational.
Hodge index leaves at most one positive-square
boundary, the incoming one. Every other boundary
therefore has negative square by \eqref{cdp:eq:C.NM.3}, and its
neighbour has strictly higher height. Follow one
of those edges when it exists.

\item The heights strictly increase and there are finitely
many components. The path ends at a rational
component having exactly its one incoming boundary.
It has $K=-aC$, with $(a,C^2)$ as in \eqref{cdp:eq:C.NM.8}.
The same integer-intersection argument excludes every
$(-1)$-curve, so it is an actual \ref{cdp:st:C.2.UN.64} leaf. Its ruled
section possibility has just been excluded by
\eqref{cdp:eq:C.NM.10}. It is therefore a plane-line leaf, and
\ref{cdp:st:C.2.UN.64} contracts it to a smooth point of the smooth
central locus of its parent.

\item Every other boundary misses that line, so the actual
blowdown is an isomorphism on its neighbourhood.
It removes exactly one rational double curve, creates
no new double curve or triple point, retains the
punctured torus family, and decreases $e(S)$ by two.
Apply the argument to the new actual degeneration
until no rational double curve remains. Exactly
$r_0$ steps have been made; \eqref{cdp:eq:C.NM.1} and the recorded
ordinary-blowup reconstruction follow.

\end{enumerate}
\end{proof}






\begin{proposition}[Torsion-normal parent alignment with all vertical blowups retained]\hypertarget{tag-01EH}{}\label{cdp:st:C.2.UN.70}
\StatementRecord{cdp:st:C.2.UN.70}{proved}{01EH}{\SeriesRef{cdp:st:C.2.UN.67}, \SeriesRef{cdp:st:C.2.UN.15}, \SeriesRef{cdp:st:C.2.UN.31}, \SeriesRef{cdp:eq:C.NP.3}, \SeriesRef{cdp:st:C.2.UN.19}, \SeriesRef{cdp:st:C.2.UN.37}, \SeriesRef{cdp:eq:C.NP.7}, \SeriesRef{cdp:eq:C.NP.8}, \SeriesRef{cdp:eq:C.NP.9}}

Let \(P\) be a component of a full reduced SNC fibre of
a smooth-total degeneration of connected complex two-tori.
Suppose it has a smooth elliptic boundary curve \(J\)
disjoint from its other boundary curves, with torsion
normal line. Suppose its coefficient in the actual
canonical divisor on \(P\) is nonnegative:
\[
 K_P=\operatorname{div}\omega_P
       =a_JJ+\sum_C a_C C,\qquad a_J\ge0 ,
 \tag{C.NP.1}\hypertarget{tag-00T1}{}\label{cdp:eq:C.NP.1}
\]
where the sum is over the other boundary curves and
the equality includes the actual meromorphic canonical
section induced by the ambient evaluation.

Then \(P\) is not rational. There is an actual proper
birational morphism
\[
 \pi:P\longrightarrow Z,\qquad q_Z:Z\longrightarrow B,
                    \qquad q=q_Z\pi ,
 \tag{C.NP.2}\hypertarget{tag-00T2}{}\label{cdp:eq:C.NP.2}
\]
with the following properties. The morphism \(\pi\)
is a finite sequence of point blowdowns and is an
isomorphism on a neighbourhood of \(J\). The map
\(q_Z\) is a relatively minimal fixed-Jacobian elliptic
quasi-bundle with identity integral monodromy, only
smooth multiple supports, and base genus at most one.
The image of \(J\) is one of these fibre supports,
of full multiplicity exactly
\(\operatorname{ord}N_{J/P}\).
Every original boundary curve on \(P\) is vertical
for \(q\). Every exceptional curve of \(\pi\) is
vertical; the complete remaining nonminimal data
are the recorded point blowups of fibres of \(q_Z\),
with their actual positions and normal maps.
In particular every rational boundary curve on
\(P\) is contracted by \(\pi\).

This applies, in particular, to the parent of a
Hopf leaf: its canonical coefficient is
\(a_J=m-1\ge0\), with \(m\) the exact normal
order of the Hopf axis. No contraction of any
of these surface exceptional curves in the
ambient threefold is asserted.
\end{proposition}



\begin{proof}

\begin{enumerate}[label={$\langle1\rangle$\arabic*.},leftmargin=*]
\item Rational \(P\) is excluded by \ref{cdp:st:C.2.UN.67}. A component
of general type is excluded by \ref{cdp:st:C.2.UN.15}. Its remaining
minimal model is therefore ruled, Hopf, or has
Kodaira dimension zero or one, by the actual
component comparisons already established.
We construct the maps separately and then
apply the actual fields on \(P\).

\par\medskip\noindent\textbf{Nonnegative Kodaira dimension: every blowdown misses \(J\).}
Let \(\pi:P\to Z\) be its ordinary minimal-model
map and let \(J_0=\pi(J)\). The surface \(Z\)
is the actual minimal model supplied by
\ref{cdp:st:C.2.UN.31} in this multiplicity-one case. It is
a torus, or has the smooth-multiple elliptic
map with the canonical line formula
(C.MR.12). In the latter case that formula
has base degree
\[
       \delta=2g(B_0)-2+\sum_i(1-1/m_i)\ge0 .
 \tag{C.NP.3}\hypertarget{tag-00T3}{}\label{cdp:eq:C.NP.3}
\]
The inequality follows from its exact
pluricanonical pushdown and
\(\kappa(Z)\ge0\). Thus \(K_Z\) has
nonnegative intersection with every
curve: it is zero on a vertical
curve, and its intersection with
a horizontal curve is its positive
base degree times \(\delta\).
For a torus \(K_Z\) is trivial.
This proves the needed intersection
property by the explicit base-line
map, without a new nefness assumption.

\item For the actual sequence of point
blowups write
\(K_P=\pi^*K_Z+\sum_E b_EE\),
where every exceptional coefficient
\(b_E\) is a positive integer.
Adjunction and \(J^2=0\) give
\[
       0=K_P\cdot J
          =K_Z\cdot J_0+\sum_E b_E(E\cdot J).
 \tag{C.NP.4}\hypertarget{tag-00T4}{}\label{cdp:eq:C.NP.4}
\]
Every term is nonnegative; \(J\)
is not an exceptional rational
curve. Thus each \(E\cdot J=0\).
The exceptional locus is disjoint
from \(J\), so \(\pi\) is an
actual isomorphism of its
neighbourhood. In particular
\(J_0\) is smooth, elliptic,
of square zero, with the
same torsion normal.

\item The exact nonruled alignment
(C.MR.12)--(C.MR.14) applies
to this marked pair \((Z,J_0)\).
In Kodaira dimension one,
\eqref{cdp:eq:C.NP.3} is positive, so
adjunction makes \(J_0\)
vertical for the given map.
On a torus it is a translate
of an elliptic subtorus and
its quotient gives the
required map. In the remaining
Kodaira-dimension-zero case,
the finite translation-root
cover in that alignment is a
principal elliptic torsor.
A horizontal lift of \(J_0\)
trivializes its integral Chern
vector after a finite base
pullback; the actual Cech
period lattice then makes
the cover a torus. Its
parallel lifted elliptic
curves define the quotient
pencil, and the finite
group preserving them
descends it to \(Z\).
These are the actual
finite-cover and quotient
maps of (C.MR.14).
They make \(J_0\) a full
reduced fibre support.
Every other noncontracted
boundary image is disjoint
from \(J_0\) and is therefore
vertical. Contracted curves
map to points and are
vertical as well.

\par\medskip\noindent\textbf{Hopf minimal model.}
Suppose the actual minimal
model \(Z\) is Hopf, as in
\ref{cdp:st:C.2.UN.19}. Every irreducible
curve on it is a smooth
elliptic curve with square
zero. The ordinary blowup
square formula gives
\[
                 0=J^2=J_0^2-\sum_\nu r_\nu^2.
 \tag{C.NP.5}\hypertarget{tag-00T5}{}\label{cdp:eq:C.NP.5}
\]
Thus no centre lies on
any successive image of
\(J\); again the map
preserves its entire
neighbourhood. Its normal
is torsion. This excludes
the nontorsion axes of an
algebraic-dimension-zero
diagonal Hopf.

\item If \(Z\) were nonlinear,
its only curve would be
its axis \(J_0\). All
other boundaries would
be contracted. Push the
actual meromorphic form
in \eqref{cdp:eq:C.NP.1} through the
common open isomorphism.
Its divisor on \(Z\)
would be \(a_JJ_0\).
But \(K_Z=-(r+1)J_0\)
with \(r\ge1\). This
would trivialize the
line of the positive
effective divisor
\((a_J+r+1)J_0\), contrary
to constancy of holomorphic
functions on the compact
surface. Thus \(Z\) is
an elliptic Hopf surface.
Its actual map from \ref{cdp:st:C.2.UN.37}
makes every surviving
boundary a smooth multiple
fibre support. Composing
with \(\pi\) gives \eqref{cdp:eq:C.NP.2},
with every exceptional
curve vertical.

\par\medskip\noindent\textbf{Ruled minimal model: the contraction preserving all lifts of \(J\).}
A ruled minimal base of
genus at least two cannot
occur. The normalization
of the image of \(J\)
would map nontrivially
from an elliptic curve
to that base, contradicting
Riemann--Hurwitz; a
vertical curve in a
rational ruling fibre
cannot have elliptic
normalization. The
rational case was
excluded already.
Thus the original ruling
extends to a rational
fibration
\[
                             p:P\longrightarrow E
 \tag{C.NP.6}\hypertarget{tag-00T6}{}\label{cdp:eq:C.NP.6}
\]
over an elliptic curve.
Its restriction \(J\to E\)
is a finite unramified
map. Pull back this
actual fibration along
that cover. After origins
are chosen, it is an
isogeny and is Galois.
On the smooth pulled-back
surface \(P'\), the
preimage of \(J\) consists
of disjoint sections
\(J_0',\ldots,J_{d-1}'\),
each of square zero and
with the original torsion
normal. Its deck group
permutes these sections
and acts freely on the
elliptic base.

\item There is a vertical
\((-1)\)-curve disjoint
from \(J_0'\) whenever
a ruling fibre is reducible.
To prove this, let
\(F=\sum_\alpha m_\alpha E_\alpha\)
be such a fibre. The
section meets just one
component, with multiplicity
one, and misses all nodes.
All \(E_\alpha\) are
smooth rational curves;
adjunction and \(K\cdot F=-2\)
give
\[
             \sum_\alpha
                    m_\alpha(2+E_\alpha^2)=2 .
 \tag{C.NP.7}\hypertarget{tag-00T7}{}\label{cdp:eq:C.NP.7}
\]
If the only possible
\((-1)\)-curve were the
component meeting the
section, its contribution
would be one and all
others would be nonpositive.
If it were not a
\((-1)\)-curve, all terms
would be nonpositive.
Both alternatives contradict
\eqref{cdp:eq:C.NP.7}. A vertical
\((-1)\)-curve disjoint
from the section therefore
exists.

\item Contract such curves
successively. This gives
a \(\PP^1\)-bundle \(Z'\)
over the same elliptic
base, preserving a
neighbourhood of \(J_0'\).
Consider any other
lifted section \(J_i'\).
Its final image is
disjoint from \(J_0'\).
The latter is a nonzero
nef square-zero curve,
so the Hodge index
theorem gives square
at most zero for that
image. The blowdown
square formula gives
square at least zero.
Equality holds; hence
every contracted curve
also missed every
\(J_i'\).

\item This proves equivariance
of the contraction, not
just existence of a
choice after base change.
In each original ruling
fibre, the unique component
meeting \(J_0'\) survives
and all the others are
contracted. Since all the
other sections were
preserved, they meet
that same surviving
component. It is
therefore the intrinsic
component meeting the
full set of lifted
sections. The exceptional
locus, its connected
fibres and the resulting
contraction are preserved
by the whole deck group.
An isomorphism on their
common open complement
extends uniquely through
the proper contraction,
so the action descends
to \(Z'\).
Its action on the base
is free, hence its
action on \(Z'\) is
free. Taking the actual
quotients gives a
birational morphism
\(\pi:P\to Z=Z'/G\),
where \(Z\) is a smooth
\(\PP^1\)-bundle over
\(E\). This morphism
is a sequence of point
blowdowns and is an
isomorphism near \(J\).
These facts also follow
by contracting the
disjoint deck orbits
of the indicated
exceptional curves.

\item At least one other
boundary survives on
\(Z\). Otherwise the
pushforward of the
actual canonical section
would give \(K_Z=a_JJ_0\).
Its intersection with
a rational ruling fibre
would be nonnegative,
whereas that canonical
intersection is minus two.
This uses the
nonnegative coefficient
in the statement.
Every surviving boundary
is disjoint from \(J_0\).

\par\medskip\noindent\textbf{The actual split bundle and descended elliptic map.}
On \(Z'\) the diagonal
section has torsion
normal \(L\). Its actual
section subline gives,
after tensoring, an
extension
\[
           0\longrightarrow\OO
             \longrightarrow V
             \longrightarrow L\longrightarrow0,
                  \qquad Z'=\PP(V).
 \tag{C.NP.8}\hypertarget{tag-00T8}{}\label{cdp:eq:C.NP.8}
\]
If \(L\ne\OO\), its
extension group
\(H^1(L^{-1})\) is zero.
If \(L=\OO\), the
complement of the section
is the additive torsor
of this extension class.
Normalize a component
of the pullback of the
second surviving boundary.
It maps finitely to the
elliptic base and lies
in that complement,
so it supplies a section
of the pulled additive
torsor. Trace composed
with pullback on
\(H^1(\OO)\) is its
positive degree.
The original extension
class is therefore zero.
Thus \eqref{cdp:eq:C.NP.8} splits
in both cases.

\item Take a further finite
etale elliptic base cover
which kills \(L\) and
is Galois over \(E\).
The actual pulled-back
surface is \(E'\times\PP^1\).
Its finite deck maps have
the form
\((b,z)\mapsto(b+\tau,A(b)z)\).
The holomorphic map
\(A:E'\to\operatorname{PGL}_2(\C)\)
is constant, since the
target is affine and
the source compact.
Consequently the
actual quotient has
the map
\[
        Z=(E'\times\PP^1)/G'
                   \longrightarrow\PP^1/G'.
 \tag{C.NP.9}\hypertarget{tag-00T9}{}\label{cdp:eq:C.NP.9}
\]
The regular connected
fibre is \(E'/K\), where
\(K\) is the kernel of
the projective action.
At a projective stabilizer
the elliptic translation
is free. Thus the chart
is precisely a smooth
multiple-fibre chart,
with identity integral
monodromy and one fixed
Jacobian. Every boundary
lift disjoint from the
fixed section has constant
projective coordinate:
on its compact
normalization that
coordinate is a
holomorphic function
with values in the
affine complement.
Thus every surviving
boundary is a fibre
support of \eqref{cdp:eq:C.NP.9}.
Every contracted boundary
maps to a point, so
all original boundaries
are vertical for the
actual composite \(q\).

\par\medskip\noindent\textbf{The original fields and the remaining vertical blowups.}
In all cases we have
constructed an actual
connected elliptic map
on \(P\), with every
boundary vertical, and
the specified isomorphism
near \(J\). Apply the
field alignment proof
of (C.MR.8), in the
explicit nonminimal
form of \ref{cdp:st:C.2.UN.67}, to this
map on \(P\). Its
saturated fields are
holomorphic on entire
general elliptic fibres.
The projected-field
continuation and the
actual return germs
give identity integral
monodromy. No field
pushforward through a
surface contraction is
needed for this step.

\item The constructed \(Z\)
is already relatively
minimal for its map.
In the nonruled case
it is the minimal
surface; Hopf surfaces
contain no rational
curves, and an elliptic
ruled minimal surface
contains no rational
\((-1)\)-curves. The
exact relatively minimal
unit-root comparison in
\ref{cdp:st:C.2.UN.67} consequently gives
only smooth multiple
fibres and
\(R^1(q_Z)_*\OO_Z=\OO_B\).
The compact-core bound
then gives
\(g(B)=p_g(Z)=p_g(P)\le1\).
Its connected normal-root
chart identifies the full
multiplicity at \(J_0\)
with the exact order of
its normal line. Since
\(\pi\) preserves that
neighbourhood, this is
the original order.

\item All exceptional curves
of \(\pi\) map to
points of \(Z\), hence
are vertical for \(q\).
Conversely a rational
boundary curve cannot
map onto a fibre support
of \(q_Z\), which is
smooth elliptic; it
must be contracted.
Factor the proper
birational surface map
into its actual point
blowdowns and retain
the inverse centres
and charts. This is
exactly the claimed
record of vertical
point blowups. It
does not contract
anything in the
ambient threefold.
For a Hopf leaf,
the exact terminal
canonical relation
gives \(a_J=m-1\),
so all hypotheses
hold. This proves
the marked parent
alignment.

\end{enumerate}
\end{proof}






\begin{proposition}[All positive canonical-flat central log spaces over a disc]\hypertarget{tag-01EI}{}\label{cdp:st:C.2.UN.71}
\StatementRecord{cdp:st:C.2.UN.71}{proved}{01EI}{\SeriesRef{cdp:st:C.2.UN.68}, \SeriesRef{cdp:eq:C.TG.4}, \SeriesRef{cdp:st:C.2.UN.65}, \SeriesRef{cdp:eq:C.TG.8}, \SeriesRef{cdp:eq:C.TG.3}, \SeriesRef{cdp:eq:C.TG.6}, \SeriesRef{cdp:eq:C.TG.10}, \SeriesRef{cdp:eq:C.TG.12}, \SeriesRef{cdp:eq:C.TG.11}, \SeriesRef{cdp:eq:C.TG.13}; \cite[Theorems 4.14, 5.2 and 5.4]{GrossSiebert2006I}; \cite[Example 3.30 and Theorem 3.22]{GrossSiebert2006I}}

The positive canonical-flat central fibres of \ref{cdp:st:C.2.UN.68},
with their actual divisorial log structure and marked
dual complex, have the following complete description.
Take a rank-two lattice $N$, a finite-index sublattice
$\Gamma\subset N$, a $\Gamma$-periodic unimodular
triangulation $\mathscr P$ of $N_{\R}$ with vertex set
$N$, and a homomorphism
\[
       c:\Gamma\longrightarrow T_N
          :=\operatorname{Hom}(N^*,\C^*).
\tag{C.TG.1}\hypertarget{tag-012Q}{}\label{cdp:eq:C.TG.1}
\]
Require that each triangle have three distinct vertices
in $N/\Gamma$; this is the embedded simple-NC condition.
The central component at $p\in N/\Gamma$ is the actual
toric surface of the star fan
$\{\R_{\ge0}(\sigma-p):p\in\sigma\in\mathscr P\}$.
The homomorphism in \eqref{cdp:eq:C.TG.1} supplies all the remaining
logarithmic gluing constants, modulo the explicit torus
changes of the local charts described below.

Every such datum gives a smooth proper analytic
degeneration of connected complex two-tori over a
sufficiently small disc, with trivial canonical line,
and with exactly the specified marked central log space.
Every actual fibre in \ref{cdp:st:C.2.UN.68} is obtained in this way.
In a basis of $N$, let the columns of the integral
matrix $Q$ be a basis of $\Gamma$, and choose logarithms
of its two values under $c$ as the columns of $C$.
The actual nearby torus is
\[
 F_t=\C^2\big/
       \left(2\pi i\Z^2+(\log t\,Q+C)\Z^2\right),
       \qquad 0<|t|<\epsilon .
\tag{C.TG.2}\hypertarget{tag-012X}{}\label{cdp:eq:C.TG.2}
\]
The central Euler number is $2[N:\Gamma]$.
This is a proper local occurrence and central-log-space
classification. It does not identify the full analytic
germ of an arbitrary given degeneration with this
chosen quotient, and it makes no compact-base
globalization assertion.
\end{proposition}



\begin{proof}

\begin{enumerate}[label={$\langle1\rangle$\arabic*.},leftmargin=*]
\item \par\medskip\noindent\textbf{The actual integral-affine charts.}
On each toric component $V_i$ of \ref{cdp:st:C.2.UN.68} write
$N_i=\pi_1(V_i\setminus D_i)$, and let $n_{ij}$
be the primitive angular meridian of its boundary
curve $C_{ij}$. A triple point $ijk$ gives the
unimodular triangle
\[
                 0,\ n_{ij},\ n_{ik}\quad\hbox{in }(N_i)_{\R}.
\tag{C.TG.3}\hypertarget{tag-012Y}{}\label{cdp:eq:C.TG.3}
\]
To pass from the chart at $j$ to that at $i$,
use the unique integral-affine map which sends
$0$ to $n_{ij}$ and whose linear part $P_{ij}$ satisfies
\[
 P_{ij}n_{ji}=-n_{ij},\qquad
 P_{ij}n_{jk}=n_{ik}-n_{ij}.
\tag{C.TG.4}\hypertarget{tag-012Z}{}\label{cdp:eq:C.TG.4}
\]
These are the actual angular maps. Indeed in a
triple chart $f=xyz$, the coordinates on $V_i$
are $(y,z)$ and those on $V_j$ are $(x,z)$.
On the nearby fibre $x=t/(yz)$, a loop varying
$x$ with $z$ fixed has angular vector $-n_{ij}$,
and a loop varying $z$ with $x$ fixed has vector
$n_{ik}-n_{ij}$. Multiplication of these coordinates
by holomorphic units leaves these integral maps
unchanged after the chart is shrunk.

\item There are exactly two triple points on $C_{ij}$.
Write their third vertices as $k,l$, retaining
their two distinct incidences if the same global
component occurs twice. The toric fan relation on
$V_i$ is
\[
 n_{ik}+n_{il}=-(C_{ij}^2)_{V_i}n_{ij}.
\tag{C.TG.5}\hypertarget{tag-0130}{}\label{cdp:eq:C.TG.5}
\]
The corresponding relation at $j$ agrees with
\eqref{cdp:eq:C.TG.4} at the other endpoint precisely because
\[
             (C_{ij}^2)_{V_i}+(C_{ij}^2)_{V_j}=-2,
\tag{C.TG.6}\hypertarget{tag-0131}{}\label{cdp:eq:C.TG.6}
\]
the actual normal-product and two-triple-point
formula of \ref{cdp:st:C.2.UN.65}. Thus there is no affine singularity
in an edge. The complete smooth toric fan at each
vertex makes its star an affine neighbourhood of
the origin. The triangle charts and \eqref{cdp:eq:C.TG.4}
therefore give an integral-affine structure on
the entire actual dual two-torus $B$.

\par\medskip\noindent\textbf{The angular lattices inject into the actual nearby torus.}
The nowhere-zero three-form of \ref{cdp:st:C.2.UN.65} is, in an
actual triple chart,
\[
       \Omega=a(x,y,z)\,dx\wedge dy\wedge dz,
          \quad a(0)\ne0,
          \qquad f=xyz.
\tag{C.TG.7}\hypertarget{tag-0132}{}\label{cdp:eq:C.TG.7}
\]
The induced relative two-form on $F_t$ is
$a(x,y,t/(xy))\,d\log x\wedge d\log y$.
Its integral over the real angular two-torus
$|x|=|y|=|z|=|t|^{1/3}$ is exactly
\[
 (2\pi i)^2\sum_{q\ge0}a_{q,q,q}t^q
                  =(2\pi i)^2(a(0)+O(t))\ne0.
\tag{C.TG.8}\hypertarget{tag-0133}{}\label{cdp:eq:C.TG.8}
\]
This follows by the convergent power-series expansion
and the two circle integrals; only the diagonal
monomials survive.

\item For a fixed component, its toric complement retracts
onto such an angular two-torus. Any finite collection
of its loops, their homotopies and the connecting
paths to the corner lies in a compact subset away
from the boundary. The submersion $f$ near this
compact subset gives actual smooth horizontal
transports into $F_t$. Consequently its angular
lattice has a well-defined homomorphism
$N_i\to\pi_1(F_t)$, up to conjugation, and the
corner homomorphism is this same map. Since
$\pi_1(F_t)=\Z^4$ is abelian, conjugation makes
no ambiguity. If this map had rank at most one,
the fundamental class of the angular torus would
map to zero in $H_2(F_t,\Z)$: the actual complex
torus is a $K(\Z^4,1)$ and its degree-two homology
map is the exterior square of its degree-one map.
This contradicts \eqref{cdp:eq:C.TG.8}. Thus every $N_i$
injects into the actual $\Z^4$.

\item The corner maps and the compact double-curve collars
identify the injected lattices by exactly \eqref{cdp:eq:C.TG.4}.
Transport around any loop in $B$ gives a linear
holonomy $A$ with $\iota_iA=\iota_i$ in $\Z^4$.
Injectivity gives $A=I$. This proves trivial linear
holonomy by maps into the given smooth fibre; no
invariant-cycle surjectivity is being assumed.

\item The affine charts now differ by translations.
Their Euclidean metrics give a complete flat metric
on compact $B$. Extending straight lines on its
universal cover proves that the developing map is
an actual affine isomorphism with $N_{\R}$.
The deck group is a full lattice of translations
$\Gamma\subset N$, because all the transition
translations in \eqref{cdp:eq:C.TG.3}--\eqref{cdp:eq:C.TG.4} are integral.
Every developed triangle is unimodular. Conversely
every lattice point of $N$ is a vertex, since a
unimodular triangle contains no other lattice point
and the triangles cover $N_{\R}$. Thus the actual
dual complex is $(N_{\R},\mathscr P)/\Gamma$,
with exactly $[N:\Gamma]$ vertices.

\par\medskip\noindent\textbf{The exact logarithmic gluing comparison.}
We spell out the primary-source comparison used to
retain the given gluing constants. The divisorial
log structure of $(X,S)$ restricted to $S$ is locally
generated by the actual factors of $f$. At a triple
point its chart is $\mathbb N^3$, and the base generator
is $(1,1,1)$. At a double curve it is the analogous
two-factor chart. The normalization is the disjoint
union of the actual toric $V_i$; its conductor is
the reduced full toric boundary, with the exact
cartesian and cocartesian SNC conductor square.
The central surface is Gorenstein and this log
map is smooth at every point, directly from
$xyz=t$, $xy=t$ and the submersion chart. Thus
the log-singular set is empty.

\item The precise gluing result is
\cite[Theorems 4.14, 5.2 and 5.4]{GrossSiebert2006I}.
Its use here does not import a polarization or an
algebraic total space. The following description
checks the analytic-to-algebraic comparison in the
proof. Fix the toric identifications of every
normalization stratum with its actual counterpart.
Every boundary identification is an automorphism
of $\PP^1$ preserving its two specified corner
points, and hence, in the specified coordinates,
is multiplication by a nonzero constant. These
are the closed gluing maps of that source.

\item At a triple point, restrict an actual log chart
to each of the three toric affine planes. A
monomial $z^p$ is multiplied by a unit $h_p$
in its analytic local ring, with
$h_{p+q}=h_ph_q$. The map
\[
       p\longmapsto h_pz^p
       \quad\longmapsto\quad
       p\longmapsto h_p(0)z^p
\tag{C.TG.9}\hypertarget{tag-0134}{}\label{cdp:eq:C.TG.9}
\]
is exactly the units-to-residues map in the proof
of Theorem 4.14. Restricting to a shared coordinate
axis and evaluating at its endpoint commute.
The resulting constant monoid maps therefore
intertwine the actual closed gluing diagrams.
Their colimits give the actual open toric charts
and the isomorphism with the given central space,
as in the commutative squares of that proof.
All these maps are defined by the same convergent
local rings in the analytic setting. The constant
residues are algebraic monomial maps; thus the
resulting finite open atlas and its actual
equivalence relation also give an algebraic
space whose analytification is the given $S$.

\item The remaining log data are the nowhere-zero
sections of the lines
$N_{C/V_i}\otimes N_{C/V_j}\otimes\OO_C(p+q)$
on the double $\PP^1$'s, where $p,q$ are their
two triple points, with the actual triple
compatibilities. These corrected normal-product
lines have degree zero by \eqref{cdp:eq:C.TG.6}, and
their sections are nonzero constants in their
trivializations. They are algebraic as well.
This is the exact log-smooth-structure torsor
in \cite[Example 3.30 and Theorem 3.22]{GrossSiebert2006I};
the constant maps just described retain its
given section, not merely the underlying surface.
Consequently the log central space satisfies
Definition 4.3 of that source over $\C$ with
empty singular set. Positivity says these
sections have no poles and is immediate.
The affine monodromy polytopes are points;
in Definition 1.60 simplicity is the case
$p=0$, whose remaining conditions are vacuous.

\item Theorems 5.2 and 5.4 now supply actual torus
changes of the open charts, preserving the
marked dual complex, which turn these data
into a Cech cocycle for the constant sheaf
$N\otimes\C^*=T_N$ on $B$.
The maps of a cohomologous change are explicitly
$s'_{\alpha\beta}=a_\alpha^{-1}
s_{\alpha\beta}a_\beta$, acting on each chart
by its torus element $a_\alpha$; their pullbacks
also preserve the normalized log section.
Thus they are actual isomorphisms of central
log spaces. On the universal cover
$N_{\R}$ a cocycle is trivialized by path
products, and its only remaining values are
the two commuting deck holonomies. Therefore
\[
 H^1(B,N\otimes\C^*)
       =\operatorname{Hom}(\Gamma,T_N),
\tag{C.TG.10}\hypertarget{tag-012R}{}\label{cdp:eq:C.TG.10}
\]
with precisely the marked gluing datum \eqref{cdp:eq:C.TG.1}.
This is the actual comparison to be used in
the construction, rather than an assertion
that arbitrary analytic neighbourhoods are toric.

\par\medskip\noindent\textbf{The smooth toric threefold and its exact deck maps.}
For every triangle $\sigma$ of $\mathscr P$, take
the cone over $\sigma\times\{1\}$ in
$N_{\R}\oplus\R$. These cones and their faces
form a fan $\Sigma$. Its three primitive
generators $(p_0,1),(p_1,1),(p_2,1)$ are an
integral basis because the triangle is
unimodular. Its toric chart is therefore
literally $\C^3$ with coordinates $x_0,x_1,x_2$,
and the character of the last dual coordinate is
\[
                  t=x_0x_1x_2,
       \qquad z^m=\prod_{a=0}^2x_a^{\langle m,p_a\rangle}
                       \quad(m\in N^*).
\tag{C.TG.11}\hypertarget{tag-012S}{}\label{cdp:eq:C.TG.11}
\]
The charts glue by their actual Laurent
monomial localizations. The fan intersection
condition makes this a Hausdorff smooth toric
analytic threefold. Although the fan has
infinitely many cones, each central orbit has
the ordinary finite star of its cone, and the
open torus is a single common open chart.
These are actual analytic neighbourhoods.

\item A translation $\lambda\in\Gamma$ acts on the
lattice $N\oplus\Z$ by the integral shear
$(n,h)\mapsto(n+h\lambda,h)$. It preserves
the fan and induces the actual monomial map
between its charts. Multiply this map by
$c_\lambda$ on the torus factor. On the
punctured open torus the result is exactly
\[
              g_\lambda(z,t)=(c_\lambda t^\lambda z,t).
\tag{C.TG.12}\hypertarget{tag-012T}{}\label{cdp:eq:C.TG.12}
\]
The $c_\lambda$ are a homomorphism, and
constant torus multiplication commutes with
the shears. Thus $g_\lambda g_\mu=g_{\lambda+\mu}$
on every chart. On the central fibre these
maps identify the component and boundary
charts by the cocycle in \eqref{cdp:eq:C.TG.10}. The
quotient central log space is consequently
the exact space specified earlier, including
its given log section.

\par\medskip\noindent\textbf{Freeness, properness and the actual compact fibres.}
Extend the homomorphism
$\lambda\mapsto\log|c_\lambda|$ real-linearly
to a map $K:N_{\R}\to N_{\R}$.
For $0<|t|<\epsilon$ put $L=\log|t|$ and
\[
                 q=\frac{\log|z|}{L}.
\tag{C.TG.13}\hypertarget{tag-012U}{}\label{cdp:eq:C.TG.13}
\]
Then \eqref{cdp:eq:C.TG.12} sends $q$ to
$q+(I+K/L)\lambda$.
Choose $\epsilon<1$ so small that $I+K/L$
and its inverse are uniformly bounded and
invertible for $0<|t|<\epsilon$.
The action on every nonzero fibre is
therefore free and properly discontinuous;
its real logarithmic quotient has a compact
fundamental parallelogram. Its angular
directions already form a compact two-torus.
Taking logarithms gives exactly \eqref{cdp:eq:C.TG.2},
whose four real periods are independent
because $LQ+\operatorname{Re}C$ is invertible.

\item At $t=0$ a nonzero translation cannot fix
any orbit: it would preserve a bounded cell
of the triangulation, which is impossible.
The central action is thus free. To verify
proper discontinuity also at the central
fibre, use a finite collection of fixed
toric charts about a proposed compact set.
For points in those charts tending to zero,
the coordinates in \eqref{cdp:eq:C.TG.11} show that $q$
stays bounded: each $\log|x_a|/L$ is bounded
below by a quantity tending to zero and
the three coefficients sum to one.
If both a point and its $g_\lambda$-translate
stay in such fixed compact sets, \eqref{cdp:eq:C.TG.13}
and the uniform inverse bound force
$\lambda$ to stay bounded. Only finitely
many lattice elements occur. The same
argument at central points uses the finite
sets of their translated cells. Hence the
quotient is Hausdorff and smooth.

\item Finally choose the representatives of $q$
in a uniformly bounded fundamental
parallelogram for $(I+K/L)\Gamma$.
Only finitely many triangles of
$\mathscr P$ meet this bounded set.
In a triangle containing $q$, its
barycentric coordinates are precisely
\[
             \ell_a=\frac{\log|x_a|}{L},
             \qquad \ell_a\ge0,\quad\sum_a\ell_a=1.
\tag{C.TG.14}\hypertarget{tag-012V}{}\label{cdp:eq:C.TG.14}
\]
Consequently $|x_a|\le1$ in that chart.
Any sequence of fibres tending to zero
has, after its actual deck translations,
a subsequence in one of these finitely
many closed unit polydiscs, and hence a
convergent subsequence. The central fibre
itself is the finite union of the complete
toric star surfaces. This proves properness
over a smaller disc. The punctured fibres
are connected by their literal quotient
presentation, and the central dual torus
is connected as well.

\item The triangle-distinctness condition in
the statement prevents a component from
meeting a translate of itself. The
quotient therefore retains smooth embedded
components and the full reduced
simple-NC equation in \eqref{cdp:eq:C.TG.11}.
The differential
\[
                  dt\wedge d\log z_1\wedge d\log z_2
\tag{C.TG.15}\hypertarget{tag-012W}{}\label{cdp:eq:C.TG.15}
\]
is invariant under \eqref{cdp:eq:C.TG.12}, because all
additional terms are multiples of $dt$.
In each unimodular chart \eqref{cdp:eq:C.TG.11} it is,
up to the basis determinant $\pm1$,
$dx_0\wedge dx_1\wedge dx_2$.
It extends nowhere vanishingly across
the central fibre. Thus the constructed
total space has trivial canonical line.
The exact central-space comparison and
\ref{cdp:st:C.2.UN.68} give its Euler number and complete
both the occurrence and exhaustion
assertions in the stated local scope.

\end{enumerate}
\end{proof}






\begin{proposition}[The prescribed toric fibre and its conjugate over the projective line]\hypertarget{tag-01EK}{}\label{cdp:st:C.2.UN.73}
\StatementRecord{cdp:st:C.2.UN.73}{proved}{01EK}{\SeriesRef{cdp:st:C.2.UN.71}, \SeriesRef{cdp:eq:C.TX.4}, \SeriesRef{cdp:eq:C.TX.8}, \SeriesRef{cdp:eq:C.TX.10}, \SeriesRef{cdp:eq:C.TX.12}}

Every marked central log space classified in \ref{cdp:st:C.2.UN.71} occurs
unchanged in a smooth compact complex threefold mapping
properly to $\PP^1$, with connected complex two-torus
fibres away from two values. Its two full singular fibres
are the prescribed reduced simple-NC log space and its
complex conjugate. Both neighbourhoods have a nowhere-zero
canonical form. If its lattice index is $d=[N:\Gamma]$,
each singular support has Euler number $2d$ and the
compact total space has Euler number $4d$.

The local base coordinates at the two values, and the
maps restoring the prescribed log sections, are part
of the construction below. The assertion preserves
the specified central log spaces; it does not assert
that an arbitrary given full degeneration germ is
the germ of this compact construction.
\end{proposition}



\begin{proof}

\begin{enumerate}[label={$\langle1\rangle$\arabic*.},leftmargin=*]
\item \par\medskip\noindent\textbf{The exact initial periods and a positive symmetric part.}
Use the actual data $(N,\Gamma,\mathscr P,c)$ of \ref{cdp:st:C.2.UN.71}.
Choose $N=\Z^2$, write $\Gamma=Q\Z^2$ with
$Q\in\operatorname{Mat}_2(\Z)$ invertible over $\R$,
and choose the logarithms $C$ of its two torus
holonomies. In the vector coordinate
$(2\pi i)^{-1}Q^{-1}\log z$, the angular lattice is
\[
        N_0=Q^{-1}\Z^2\supset\Z^2,
             \qquad b=(2\pi i)^{-1}Q^{-1}C.
\tag{C.TX.1}\hypertarget{tag-013Y}{}\label{cdp:eq:C.TX.1}
\]
The two other periods of the local model are
$b_j+(\log u)/(2\pi i)e_j$. Write $b=U+iA$
with $U,A$ real matrices. Choose $R>0$ large enough
that the symmetric part $S$ of $A+RI$ is positive
definite, and put
\[
 b'=b+iRI,
 \qquad \kappa=e^{-2\pi R}>0,
 \qquad A+RI=S+aJ,
 \qquad J=\begin{pmatrix}0&1\\-1&0\end{pmatrix}.
\tag{C.TX.2}\hypertarget{tag-0144}{}\label{cdp:eq:C.TX.2}
\]
In particular
$\exp(2\pi iQb'_j)=\kappa^{Qe_j}c_{Qe_j}$.
This changes only a recorded base scale; the
actual coordinate maps undo it at both ends.
We do not commute this imaginary shift through
the real-linear operation used next.

\par\medskip\noindent\textbf{A holomorphic family on the entire punctured line.}
Fix any integer $k\ge1$. For each constant
vector $w\in\C^2$, let
\[
 V_t(w)=w+t^kJ\overline w,
 \qquad B_t=I+t^kJ,
 \qquad 0\ne t\in\C.
\tag{C.TX.3}\hypertarget{tag-0145}{}\label{cdp:eq:C.TX.3}
\]
This is holomorphic in $t$ for every fixed $w$.
Define the period lattice by its actual generators
\[
 \Lambda_t=V_t(N_0)
   +\sum_{j=1}^2\Z
       \left(V_t(b'_j)+\frac{\log t}{2\pi i}B_te_j\right).
\tag{C.TX.4}\hypertarget{tag-0146}{}\label{cdp:eq:C.TX.4}
\]
Changing the logarithm by $2\pi i$ adds $B_te_j$,
which is an angular period because $\Z^2\subset N_0$.
Thus these lattices are globally single-valued
over $\C^*$, with the displayed integral change
of basis on overlaps of logarithm branches.

\item We prove their real rank is four everywhere.
It suffices to use $\Z^2$ in place of its finite
angular enlargement $N_0$. Subtract the real
angular combinations supplied by $U$ and by
$\operatorname{Re}(\log t/(2\pi i))$ from the
last two columns. Set
\[
 x=\operatorname{Re}(t^k),\quad
 w=\operatorname{Im}(t^k),\quad
 r=|t|^k,\quad
 y=-\frac{\log|t|}{2\pi}.
\tag{C.TX.5}\hypertarget{tag-0147}{}\label{cdp:eq:C.TX.5}
\]
With $A'=S+aJ=A+RI$, the resulting real
period matrix is exactly
\[
 \begin{pmatrix}
 I+xJ & wJ(A'-yI)\\
 wJ   & (I-xJ)A'+y(I+xJ)
 \end{pmatrix}.
\tag{C.TX.6}\hypertarget{tag-0148}{}\label{cdp:eq:C.TX.6}
\]
Its determinant, by the displayed block
multiplication, is
\[
 \begin{aligned}
 \Delta(t)={}&(1+r^2)^2\det S
       +(1-r^2)(1+r^2)y\operatorname{tr}S\\
       &+y^2(1-r^2)^2
       +\bigl((1+r^2)a+2xy\bigr)^2>0.
 \end{aligned}
\tag{C.TX.7}\hypertarget{tag-0149}{}\label{cdp:eq:C.TX.7}
\]
Indeed $(1-r^2)y\ge0$, the first term is
strictly positive, and the other three
terms are nonnegative. The identity uses
$r^2=x^2+w^2$ and holds for every argument
of $t$; it also holds where $B_t$ fails
to be complex-invertible. Thus \eqref{cdp:eq:C.TX.4}
is a discrete rank-four real lattice
for every $t\in\C^*$.

\item The actual quotient of the holomorphic
vector bundle $\C^2\times\C^*$ by these
locally holomorphic period sections is
therefore a smooth family of connected
complex two-tori. On every compact
annulus, the real period matrices and
their inverses are uniformly bounded.
Their real fundamental parallelepipeds
give compact representatives. This
proves properness of this middle family
over the annulus, and hence over $\C^*$.

\par\medskip\noindent\textbf{The exact chart and prescribed log section at zero.}
Near zero $B_t$ is complex-invertible.
In the middle vector coordinate $v$ use
the actual exponential map
\[
                   z=\exp(2\pi iQ B_t^{-1}v).
\tag{C.TX.8}\hypertarget{tag-014A}{}\label{cdp:eq:C.TX.8}
\]
It quotients exactly the angular lattice
$B_tN_0$, since $QB_t^{-1}B_tN_0=\Z^2$.
The other two period translations act by
\[
 z\longmapsto c^0_{Qe_j}(t)\,t^{Qe_j}z,
 \qquad
 c^0_{Qe_j}(t)=
       \exp\bigl(2\pi iQ B_t^{-1}V_t(b'_j)\bigr).
\tag{C.TX.9}\hypertarget{tag-014B}{}\label{cdp:eq:C.TX.9}
\]
These are holomorphic nonzero vectors
near zero; extend them multiplicatively
to a homomorphism $c^0(t):\Gamma\to T_N$.
Their exact central values are
$c^0_\lambda(0)=\kappa^\lambda c_\lambda$.

\item Use the base coordinate $u=\kappa t$.
The same action is then
\[
 z\longmapsto d^0_\lambda(u)u^\lambda z,
 \qquad
 d^0_\lambda(u)=\kappa^{-\lambda}
                    c^0_\lambda(u/\kappa),
 \qquad d^0_\lambda(0)=c_\lambda.
\tag{C.TX.10}\hypertarget{tag-013Z}{}\label{cdp:eq:C.TX.10}
\]
This is an actual equality of holomorphic
maps, not a replacement of its limiting
constant by a numerical invariant.

\item Apply the cone-chart construction of
\ref{cdp:st:C.2.UN.71} to these variable holomorphic
units $d^0_\lambda(u)$. Multiplication
by them is a holomorphic torus action
on every toric chart over the same
base. Their homomorphism identity
makes the shear action exact. The
proof of freeness and properness is
unchanged: the real-linear maps
$\lambda\mapsto\log|d^0_\lambda(u)|$
are uniformly bounded as $u\to0$,
so the perturbation of the tropical
translation lattice tends uniformly
to the identity. The same finitely
many barycentric unit polydiscs give
properness at zero. Thus \eqref{cdp:eq:C.TX.8}
identifies this filling with the
middle family on a punctured disc.
Its actual central log constants in
the coordinate $u$ are exactly $c$,
by \eqref{cdp:eq:C.TX.10}, so the specified
marked log space is retained.

\par\medskip\noindent\textbf{The exact conjugate filling at infinity.}
Put $s=t^{-1}$ and change the vector
frame by $v_\infty=s^kv$. The angular
periods become
\[
 B^\infty_sN_0,
       \qquad B^\infty_s=s^kI+J.
\tag{C.TX.11}\hypertarget{tag-0140}{}\label{cdp:eq:C.TX.11}
\]
Reverse the two nonangular period
generators. Their resulting exact
expressions are
\[
 -(s^kb'_j+J\overline{b'_j})
                +\frac{\log s}{2\pi i}B^\infty_se_j.
\tag{C.TX.12}\hypertarget{tag-0141}{}\label{cdp:eq:C.TX.12}
\]
The integer generator reversal is
permitted because it changes neither
the subgroup nor the actual quotient.
The matrix $B^\infty_s$ is invertible
near $s=0$. The exponential coordinates
$z_\infty=\exp(2\pi iQ(B^\infty_s)^{-1}v_\infty)$
therefore give the exact deck constants
\[
 c^\infty_{Qe_j}(s)=
  \exp\bigl(-2\pi iQ(B^\infty_s)^{-1}
                  (s^kb'_j+J\overline{b'_j})\bigr).
\tag{C.TX.13}\hypertarget{tag-0142}{}\label{cdp:eq:C.TX.13}
\]
At zero this is
\[
 c^\infty_\lambda(0)
       =\overline{c^0_\lambda(0)}
       =\kappa^\lambda\overline{c_\lambda}.
\tag{C.TX.14}\hypertarget{tag-0143}{}\label{cdp:eq:C.TX.14}
\]
The sign in \eqref{cdp:eq:C.TX.12} is essential:
$\exp(-2\pi i\overline{Qb'})$
is the conjugate of $\exp(2\pi iQb')$.

\item Choose the base coordinate $v=\kappa s$
at infinity. Define
$d^\infty_\lambda(v)=\kappa^{-\lambda}
c^\infty_\lambda(v/\kappa)$.
Then $d^\infty_\lambda(0)=\overline{c_\lambda}$,
and the exact action is
$z_\infty\mapsto d^\infty_\lambda(v)v^\lambda z_\infty$.
The two base charts satisfy $uv=\kappa^2$;
they are literal holomorphic coordinates
at the two points of one projective line.
Use the same integral fan $\Sigma$ and
its variable-constant cone quotient
at this end. All its monomial chart
maps are defined over $\Z$, while its
central torus gluing constants are
the complex conjugates of the given
ones. Hence its marked central log
space is exactly the conjugate of
the prescribed one, with the stated
local base coordinate.

\par\medskip\noindent\textbf{Compact gluing and canonical neighbourhoods.}
The exponential maps \eqref{cdp:eq:C.TX.8} and
their infinity counterparts are
actual isomorphisms of the
punctured torus families, since
they first divide by precisely
the angular lattice and then
by the displayed remaining
period group. Glue the two
proper cone fillings to the
middle family along these
isomorphisms. Disjoint smaller
base discs and a compact
intermediate annulus show
directly that the glued total
space is Hausdorff and proper
over $\PP^1$. It is smooth,
and its fibres are connected.
Properness over the compact
base makes the total space
compact. There are no other
singular fibres, because the
middle family is smooth.

\item In both end fillings the form
$du\wedge d\log z_1\wedge d\log z_2$
(respectively with $v,z_\infty$)
is invariant under the actual
variable-constant shear action:
the additional logarithmic
derivatives are multiples of
the base differential and
vanish in this wedge. The
unimodular cone calculation
of \ref{cdp:st:C.2.UN.71} extends it to a
nowhere-zero three-form.
Thus both end neighbourhoods
are canonical-flat with the
full reduced simple-NC
fibres specified in the
statement. Finally Euler
additivity over the base,
with smooth torus Euler
number zero, gives
$e(Y)=2d+2d=4d$.
This completes the compact
occurrence construction.

\end{enumerate}
\end{proof}






In this subsection the marked weighted data include the reduced
central analytic space, its component multiplicities, the central
log space of the normalized cyclic root, and the old cone cells
before their semistable subdivision.  A coarsening is an actual
proper toric morphism between the displayed monoid charts.  These
data retain more than the dual incidence complex.

\begin{proposition}[Exact weighted toric descent and compact occurrence]\hypertarget{tag-01EN}{}\label{cdp:st:C.2.UN.76}
\StatementRecord{cdp:st:C.2.UN.76}{proved}{01EN}{\SeriesRef{cdp:st:C.2.UN.71}, \SeriesRef{cdp:eq:C.WF.9}, \SeriesRef{cdp:eq:C.WF.11}, \SeriesRef{cdp:eq:C.WF.28}, \SeriesRef{cdp:eq:C.WF.5}, \SeriesRef{cdp:eq:C.WF.6}, \SeriesRef{cdp:eq:C.WF.4}, \SeriesRef{cdp:eq:C.WF.7}, \SeriesRef{cdp:st:C.2.UN.68}, \SeriesRef{cdp:eq:C.WF.8}, \SeriesRef{cdp:st:C.2.UN.73}, \SeriesRef{cdp:eq:C.WF.10}, \SeriesRef{cdp:eq:C.WF.14}, \SeriesRef{cdp:eq:C.WF.15}, \SeriesRef{cdp:eq:C.WF.17}, \SeriesRef{cdp:eq:C.WF.19}, \SeriesRef{cdp:eq:C.WF.22}, \SeriesRef{cdp:eq:C.WF.25}, \SeriesRef{cdp:eq:C.WF.24}, \SeriesRef{cdp:eq:C.WF.26}, \SeriesRef{cdp:eq:C.WF.18}}

Let \(f:X\to\Delta_t\) be a proper degeneration of connected complex
two-tori, with smooth total space and full fibre
\(F=\sum_i m_iV_i\) whose reduced support is simple normal crossing.
Suppose there is a triple point and that the actual orders of a
generator of \(f_*K_X\) satisfy
\[
 \operatorname{div}(\Omega)=\sum_i k_iV_i,\qquad
 \frac{k_i+1}{m_i}=\lambda\quad\hbox{for every }i.
\tag{C.WF.1}\hypertarget{tag-0165}{}\label{cdp:eq:C.WF.1}
\]
Then the following conclusions hold.

The normalized cyclic root, followed by the explicit height-one
subdivision below, is a positive canonical-flat toric degeneration
of \ref{cdp:st:C.2.UN.71}.  The original marked weighted support is exactly its
retained coarse-cone quotient by the root group.  Thus this branch
is described by the toric data \((N,\Gamma,\mathscr P,c)\) of \ref{cdp:st:C.2.UN.71},
the retained periodic coarse triangles, and the actual cyclic
monomial maps \eqref{cdp:eq:C.WF.9}--\eqref{cdp:eq:C.WF.11}, with their specified smooth
coarse invariant charts \eqref{cdp:eq:C.WF.28}.

There are precisely two possible common ratios:
\[
 \lambda=1,\quad k_i=m_i-1;
 \qquad\hbox{or}\qquad
 \lambda=\tfrac12,\quad k_i=\tfrac{m_i}{2}-1.
\tag{C.WF.2}\hypertarget{tag-016G}{}\label{cdp:eq:C.WF.2}
\]
In the second case every \(m_i\) is even.  Every positive-dimensional
open stratum of the original support has Euler number zero.
Consequently
\[
             e(F_{\mathrm{red}})
               =\#\{\hbox{original triple points}\}>0.
\tag{C.WF.3}\hypertarget{tag-016Q}{}\label{cdp:eq:C.WF.3}
\]
In particular this entire weighted branch contains no zero-Euler
support.

Every such marked weighted log support has a realization in a
smooth compact complex threefold over \(\PP^1\), with connected
two-torus fibres away from two values.  The two marked weighted
supports are the given one and its complex conjugate.  The actual
old-cone contractions and finite invariant-ring maps are retained
at both ends, so the reduced analytic support, all its gluing maps,
and every full-fibre multiplicity are unchanged.  The total Euler
number of this compact realization is twice \eqref{cdp:eq:C.WF.3}.

The exact original nonreduced fibre and its entire neighbourhood
are recovered by the original coarsening and quotient morphisms.
The compact occurrence assertion preserves the marked reduced
log support and multiplicity cycle; it does not identify an
arbitrary original higher infinitesimal thickening or full
degeneration germ with the new compact realization.
\end{proposition}



\begin{proof}\leavevmode\par



 
\begin{enumerate}[label={$\langle1\rangle$\arabic*.},leftmargin=0pt,labelsep=.5em]
\item \par\medskip\noindent\textbf{The actual normalized-root and coarse-cone maps.}
Choose \(M\) divisible by all \(m_i\), and use \(t=s^M\).
Locally a holomorphic root of the nonzero unit in \(f\) gives
actual coordinates in which
\[
 t=x_1^{m_1}\cdots x_r^{m_r},\qquad 1\le r\le3.
\tag{C.WF.4}\hypertarget{tag-016R}{}\label{cdp:eq:C.WF.4}
\]
The coordinate change is invertible because its derivative in
the changed normal direction is a unit.  Keep its choice and
all its transition units in the normalization diagram.
The normalized root has \(\gcd(m_1,\ldots,m_r)\) local components.
Each is the normal toric chart with cocharacter lattice and cone
\[
 L=\{(v,h)\in\Z^r\oplus\Z:\ \textstyle\sum_i m_iv_i=Mh\},
 \qquad \sigma=\{v_i\ge0\}.
\tag{C.WF.5}\hypertarget{tag-016S}{}\label{cdp:eq:C.WF.5}
\]
The remaining \(3-r\) coordinates are unchanged.  Its old rays
are \((M e_i/m_i,1)\), which are primitive because their height
is one.  Formula \eqref{cdp:eq:C.WF.5}, its saturated dual monoid, and the
characters \(x_i,s\) are the actual normalization maps; no
unramified branches are discarded.


\begin{enumerate}[label={$\langle2\rangle$\arabic*.},leftmargin=1.3em,labelsep=.5em]
\item  The height-one slice is an interval or a lattice triangle.
Subdivide its edges at every lattice point and triangulate
using all its lattice points.  Every resulting empty lattice
triangle is unimodular: its Euclidean lattice area is \(1/2\)
by the elementary lattice-area formula.  Equivalently its
three lifted height-one vectors are an integral basis of \(L\).
Thus the resulting charts are literally
\(\C^3\), with \(s=u_1u_2u_3\), or their two-factor and
one-factor versions.


\item  For the actual original charts, perform this refinement by
successive height-one star subdivisions.  Each is the normalized
blowup of its monomial ideal.  Changing a boundary equation by a
holomorphic unit multiplies each monomial generator by a unit
and leaves that ideal unchanged.  The normalized blowup maps
therefore glue with the original tangent-dependent transition
units retained.  We have not treated those original overlaps
as constant toric maps.  After all height-one lattice points
have been inserted the same empty-triangle argument gives the
unimodular refinement.


\item  The subdivisions can be chosen equivariantly.  A stabilizer of an old
coarse cone in the actual cyclic root fixes its old rays
individually: they lie over the distinct original SNC
components at that stratum.  It therefore acts trivially on
the height-one simplex.  Choose one subdivision on each orbit
of old triangles after the common edge subdivisions, and
translate it to the other members of the orbit.  Equivalently
perform the corresponding star operations on their whole
orbits; the old component labels and the trivial simplex
stabilizers give the same restrictions on every shared face.
The proper normalized blowup maps are compatible on overlaps.  They give
an actual proper equivariant morphism
\[
 \widetilde X\longrightarrow X_M
       :=\operatorname{Norm}(X\times_{\Delta_t}\Delta_s),
 \qquad X_M/\mu_M=X.
\tag{C.WF.6}\hypertarget{tag-016T}{}\label{cdp:eq:C.WF.6}
\]
The equality on the right is the actual finite analytic
quotient: invariant meromorphic functions are those on \(X\),
and normality of the smooth \(X\) makes its integral invariant
algebra exactly \(\OO_X\).  The retained coarsening in
\eqref{cdp:eq:C.WF.6} is essential; taking the quotient of the subdivided
space alone would usually retain additional divisors.
The semistable support is embedded SNC globally.  Its old
rays retain the distinct original component labels.  A new
ray retains its original stratum and its lattice position
in that stratum's old simplex.  Transition maps and root
inertia fix these positions on each old cone.  Rays occurring
together in a subdivided cone therefore cannot become two
branches of one global prime.  The star charts of each
prime glue as the smooth toric boundary charts of that
labelled ray.


\end{enumerate}

\item \par\medskip\noindent\textbf{The canonical orders, with their exact normalization.}
At a height-one ray \((v,1)\), direct differentiation of
\eqref{cdp:eq:C.WF.4} gives
\[
 \operatorname{ord}_{v}(\rho^*\Omega)
       =\sum_i(k_i+1)v_i-1
       =M\lambda-1.
\tag{C.WF.7}\hypertarget{tag-016U}{}\label{cdp:eq:C.WF.7}
\]
For an old ray this is \(M(k_i+1)/m_i-1\).
For every new ray it is the affine interpolation of those
old values.  This also proves the formula in a two-factor
chart after wedging the unchanged tangent differential.
The pulled-back nonzero coefficient is a unit throughout
each normalization chart.  Hence
\[
       \widetilde\Omega
          =s^{-(M\lambda-1)}\rho^*\Omega
\tag{C.WF.8}\hypertarget{tag-016V}{}\label{cdp:eq:C.WF.8}
\]
is a nowhere-zero holomorphic three-form on \(\widetilde X\).
Its exponent is an integer by \eqref{cdp:eq:C.WF.7}.
An old triple cone contains at least one unimodular triangle,
so the root fibre has a triple point.  \ref{cdp:st:C.2.UN.68} and \ref{cdp:st:C.2.UN.71} apply to
this actual reduced semistable root.  In particular its marked
central log space is the toric space of
\((N,\Gamma,\mathscr P,c)\).  The old cones give a periodic
coarsening \(\mathscr P_0\) of this same developed triangulation.


\begin{enumerate}[label={$\langle2\rangle$\arabic*.},leftmargin=1.3em,labelsep=.5em]
\item  Notice the normalization of \eqref{cdp:eq:C.WF.1}.  In a weighted fibre it
is not necessary that some \(k_i\) be zero.  Positivity of
the \(k_i+1\) gives \(\lambda>0\).  If \(\lambda>1\), every
integer \(k_i-m_i=(\lambda-1)m_i-1\) is nonnegative, so
\(\Omega/t\) is holomorphic on the smooth \(X\).  This
contradicts that \(\Omega\) generates \(f_*K_X\).
Thus \(0<\lambda\le1\).


\end{enumerate}

\item \par\medskip\noindent\textbf{The exact central finite action.}
Write \(\zeta=\exp(2\pi i/M)\).  Lift the root generator
to the developed affine plane and to the path-product
trivialization of the toric cocycle in \ref{cdp:st:C.2.UN.71}.  The resulting
maps on the open torus are
\[
 g(z,s)=(a\,s^{b_0}z^A,\zeta s),\qquad
 A\in\operatorname{GL}(N),\quad b_0\in N,\quad a\in T_N,
 \qquad p\longmapsto Ap+b_0.
\tag{C.WF.9}\hypertarget{tag-016W}{}\label{cdp:eq:C.WF.9}
\]
They preserve both \(\mathscr P\) and the marked coarsening
\(\mathscr P_0\).  Here \(z^A\) denotes the map whose
pullback on characters is \(A^{\mathsf T}\).


\begin{enumerate}[label={$\langle2\rangle$\arabic*.},leftmargin=1.3em,labelsep=.5em]
\item  We justify the central comparison, including its constants.
Each component map preserves its full toric boundary, so on
its dense torus it is a torus translation followed by the
integral map of its fan.  This follows directly by pulling
back a character: its divisor is a specified integral
combination of boundary primes, and the quotient of two
functions with that divisor is a holomorphic unit on the
compact component, hence a constant.  In the global
characters at a vertex \(p\), its dense component coordinates
are \(z^m s^{-\langle m,p\rangle}\).  Substitution in
\eqref{cdp:eq:C.WF.9} gives the component constant
\(a\,\zeta^{-(Ap+b_0)}\).
Conversely the actual double-curve map fixes the tangential
character comparison, and its marked normal-product log
section fixes the other character comparison.  In the
three-factor chart these are exactly the two substitutions
just displayed.  Thus after the \ref{cdp:st:C.2.UN.71} path-product
trivialization the quantities
\(a_p\zeta^{Ap+b_0}\) agree across every edge and equal one
constant \(a\).  This proves \eqref{cdp:eq:C.WF.9} on the entire central
log atlas.  It is a central-log comparison; no conjugacy of
the arbitrary original whole germ has been asserted.


\item  The lifted \(g\) need not have order \(M\) on this developed
cover.  Instead it normalizes the actual period maps
\(g_\gamma(z,s)=(c_\gamma s^\gamma z,s)\) and satisfies
\[
 \begin{split}
 A\Gamma&=\Gamma,\qquad A^M=I,\\
 c_\gamma^A&=\zeta^{A\gamma}c_{A\gamma},\\
 g^M&=g_{\gamma_0}\quad\hbox{for a retained }\gamma_0\in\Gamma .
 \end{split}
\tag{C.WF.10}\hypertarget{tag-0166}{}\label{cdp:eq:C.WF.10}
\]
All are equalities of monomial maps.  For clarity the last
one means
\[
 \gamma_0=\sum_{j=0}^{M-1}A^{M-1-j}b_0,\qquad
 \prod_{j=0}^{M-1}(a^{A^{M-1-j}})
       \ \zeta^{\sum_{j=0}^{M-1}jA^{M-1-j}b_0}
       =c_{\gamma_0}.
\tag{C.WF.11}\hypertarget{tag-0167}{}\label{cdp:eq:C.WF.11}
\]
The central log comparison preserves these relations.
Indeed two monomial maps with the same affine map and base
character differ by a constant torus translation; their
actions on a height-one component determine that constant.
Equality of the actual central maps therefore implies the
displayed equality on every toric chart.


\item  The relative residue of \eqref{cdp:eq:C.WF.8} transforms under the root
generator by \(\zeta^{-M\lambda}\).  On the developed toric
log atlas it is a common nonzero scalar times
\(d\log z_1\wedge d\log z_2\), and \eqref{cdp:eq:C.WF.9} transforms this
residue by \(\det A\).  The scalar is common because the
developed transition maps are translations and the double
residues agree.  Consequently
\[
                  \zeta^{-M\lambda}=\det A\in\{1,-1\}.
\tag{C.WF.12}\hypertarget{tag-0168}{}\label{cdp:eq:C.WF.12}
\]
Together with \(0<\lambda\le1\), this proves \eqref{cdp:eq:C.WF.2},
including the evenness of every \(m_i\) in the second case.


\end{enumerate}

\item \par\medskip\noindent\textbf{Euler number on the actual original strata.}
Let \(V_i^\circ\) be an original component with all double
curves removed.  The strict root components above it have
finite unramified maps to \(V_i^\circ\); locally these are
exactly the branches of \(x_i^{m_i}=s^M\) with the
nonzero units retained.  Their total degree is \(m_i\).
The semistable subdivisions change nothing there.
On every old root component this inverse image is its full
toric complement \((\C^*)^2\).  The sum of their Euler
numbers is therefore zero, and finite-cover multiplicativity
gives \(e_c(V_i^\circ)=0\).


\begin{enumerate}[label={$\langle2\rangle$\arabic*.},leftmargin=1.3em,labelsep=.5em]
\item  For a connected original double curve \(C\), remove its
original triple points and write \(C^\circ\).  The normalized
root over this open curve is locally the two-dimensional
monomial root times \(C^\circ\).  Its subdivision is the
corresponding chain times that curve.  It creates no triple
point over the interior of \(C^\circ\).  Choose a boundary
curve of an old strict root component which dominates \(C\);
boundary curves supported over an original triple point
are not used here.  Every such dominating curve is a toric
\(\PP^1\), and the
preimage of the removed original endpoints is precisely
its two toric endpoints.  Thus its open part is \(\C^*\)
and maps finitely and unramifiedly to \(C^\circ\).
This remains true component by component when the root
cover is disconnected over \(C^\circ\).  Its positive total
degree gives \(e_c(C^\circ)=0\).  Additivity on the original
SNC stratification now proves \eqref{cdp:eq:C.WF.3}.  This argument
uses the retained old cells, rather than the Euler number
of the subdivided root fibre.


\end{enumerate}

\item \par\medskip\noindent\textbf{The period equations retaining the finite action.}
Choose \(N=\Z^2\), \(\Gamma=Q\Z^2\) and logarithms \(C\)
of \(c\), as in \ref{cdp:st:C.2.UN.73}.  Initially put
\[
 \begin{gathered}
 \Lambda= \Z^2,\quad \Lambda_0=Q^{-1}\Z^2,\quad
 B=(2\pi i)^{-1}Q^{-1}C,\\
 D=Q^{-1}AQ,\quad
 d=Q^{-1}b_0,\quad
 a_0=(2\pi i)^{-1}Q^{-1}\log a .
 \end{gathered}
\tag{C.WF.13}\hypertarget{tag-0169}{}\label{cdp:eq:C.WF.13}
\]
The additive action is \(w\mapsto Dw+d\,l+a_0\),
where \(l=\log s/(2\pi i)\).  Replace these real lattices
by their actual images under one real invertible matrix
which makes the finite-order \(D\) orthogonal for an
averaged positive metric.  Replace \(B,D\) by their
conjugates and \(d,a_0\) by their images under the same
matrix.  We retain the names
\(\Lambda\subset\Lambda_0\), \(B,D,d,a_0\) for these
actual images.  In particular no integral lattice is
replaced by the abstract standard lattice in this step.
Let \(T\) be the actual real matrix taking this vector
coordinate back to \((2\pi i)^{-1}\log z\).


\begin{enumerate}[label={$\langle2\rangle$\arabic*.},leftmargin=1.3em,labelsep=.5em]
\item  For \(\ell\in\Lambda\), \eqref{cdp:eq:C.WF.10} becomes
\[
       (DB-BD-D/M)\ell\in\Lambda_0.
\tag{C.WF.14}\hypertarget{tag-016A}{}\label{cdp:eq:C.WF.14}
\]
Define
\[
 d_M=\sum_{j=0}^{M-1}D^{M-1-j}d\in\Lambda,\quad
 a_M=\sum_{j=0}^{M-1}D^{M-1-j}a_0,\quad
 e_M=\sum_{j=0}^{M-1}jD^{M-1-j}d .
\]
The order relation is precisely
\[
                   a_M+e_M/M-Bd_M\in\Lambda_0 .
\tag{C.WF.15}\hypertarget{tag-016B}{}\label{cdp:eq:C.WF.15}
\]
These are lattice-valued identities obtained by taking
the chosen actual logarithms of \eqref{cdp:eq:C.WF.10}--\eqref{cdp:eq:C.WF.11}.


\end{enumerate}

\item \par\medskip\noindent\textbf{An equivariant compact middle family.}
In the orthogonal coordinates put
\[
 J=\begin{pmatrix}0&1\\-1&0\end{pmatrix},\qquad
 DJ=(\det D)JD.
\]
Choose \(k\ge1\) with \(\zeta^k=\det D\).
Choose \(R>0\) so that the symmetric part \(S\) of
\(\operatorname{Im}B+RI\) is positive definite, and put
\[
 B'=B+iRI,\quad a'_0=a_0+iRd,\quad
 \kappa=e^{-2\pi R},\qquad
 V_s(q)=q+s^kJ\overline q,\quad H_s=I+s^kJ.
\tag{C.WF.16}\hypertarget{tag-016C}{}\label{cdp:eq:C.WF.16}
\]
The simultaneous shift of \(a_0\) is required: it is
exactly the substitution \(u=\kappa s\) in the original
affine map.  Equations \eqref{cdp:eq:C.WF.14}--\eqref{cdp:eq:C.WF.15} remain true
with \(B',a'_0\), since their added terms cancel.


\begin{enumerate}[label={$\langle2\rangle$\arabic*.},leftmargin=1.3em,labelsep=.5em]
\item  On the whole punctured line use the actual period lattice
\[
 \mathcal L_s=V_s(\Lambda_0)+
       \{\,V_s(B'\ell)+l(s)V_s(\ell):\ell\in\Lambda\,\}.
\tag{C.WF.17}\hypertarget{tag-016D}{}\label{cdp:eq:C.WF.17}
\]
Changing the logarithm adds an angular period.  The
determinant calculation of \ref{cdp:st:C.2.UN.73} applies in the chosen
real coordinates, with a fixed positive basis factor.
Writing \(\operatorname{Im}B'=S+aJ\),
\(r=|s|^k\), \(x=\operatorname{Re}(s^k)\) and
\(y=-\log|s|/(2\pi)\), its variable factor is
\[
 (1+r^2)^2\det S+(1-r^2)(1+r^2)y\operatorname{tr}S
 +y^2(1-r^2)^2+\bigl((1+r^2)a+2xy\bigr)^2>0 .
\tag{C.WF.18}\hypertarget{tag-016E}{}\label{cdp:eq:C.WF.18}
\]
Thus \eqref{cdp:eq:C.WF.17} gives a proper smooth torus family over
\(\C^*\), including the points where \(H_s\) is singular.


\item  The finite action on this family is the actual map
\[
 (v,s)\longmapsto
 \bigl(Dv+V_{\zeta s}(a'_0)
                  +l(s)V_{\zeta s}(d),\ \zeta s\bigr).
\tag{C.WF.19}\hypertarget{tag-016F}{}\label{cdp:eq:C.WF.19}
\]
The identity \(DV_s=V_{\zeta s}D\) follows from the
choice of \(k\).  Equation \eqref{cdp:eq:C.WF.14} proves that
\eqref{cdp:eq:C.WF.19} preserves the actual period subgroups.
Iterating it \(M\) times gives the period
\[
 V_s(a'_M+e_M/M)+l(s)V_s(d_M)
       =V_s(B'd_M+n)+l(s)V_s(d_M),\qquad n\in\Lambda_0,
\tag{C.WF.20}\hypertarget{tag-016H}{}\label{cdp:eq:C.WF.20}
\]
by \eqref{cdp:eq:C.WF.15}.  Hence its induced action has order
dividing \(M\).  Its base action has order \(M\), so
the action is faithful and free over \(\C^*\).
All these computations use only real coefficients
\(j/M\) inside \(V_s\); the complex number \(l(s)\)
is kept outside.  No complex-linearity of \(V_s\)
is presumed.


\end{enumerate}

\item \par\medskip\noindent\textbf{The exact zero-end filling and original action.}
Near zero use
\[
                    z=\exp(2\pi iT H_s^{-1}v).
\tag{C.WF.21}\hypertarget{tag-016I}{}\label{cdp:eq:C.WF.21}
\]
This divides by exactly the angular lattice.
The other periods give the variable toric holonomies
of \ref{cdp:st:C.2.UN.73}, with central values \(\kappa^\gamma c_\gamma\).
In this coordinate \eqref{cdp:eq:C.WF.19} is
\[
 (z,s)\longmapsto
   \bigl(a(s)s^{b_0}z^A,\zeta s\bigr),\qquad
 a(s)=\exp\bigl(2\pi iT H_{\zeta s}^{-1}
                              V_{\zeta s}(a'_0)\bigr).
\tag{C.WF.22}\hypertarget{tag-016J}{}\label{cdp:eq:C.WF.22}
\]
The coefficient is a holomorphic nonzero vector and
\(a(0)=\kappa^{b_0}a\).
Using the actual local base \(u=\kappa s\)
therefore restores exactly \(a\), \(c_\gamma\),
and the given central root action.  Both
\(\mathscr P\) and its marked old coarsening are
preserved.  Apply the variable-unit toric chart
construction of \ref{cdp:st:C.2.UN.73}, first with the subdivision
and then with its proper coarsening.  The monomial
maps \eqref{cdp:eq:C.WF.22} extend to both atlases and commute
with the actual coarsening morphism.


\item \par\medskip\noindent\textbf{The reflected and conjugate infinity end.}
Set \(q=s^{-1}\), replace the vector frame by
\(v_\infty=q^kv\), and put
\[
 H^\infty_q=q^kI+J,\qquad
 V^\infty_q(w)=q^kw+J\overline w .
\]
Reverse the two nonangular period generators as
in \ref{cdp:st:C.2.UN.73}.  The new period group has generators
\[
 H^\infty_q\Lambda_0,\qquad
 -V^\infty_q(B'\ell)+l(q)H^\infty_q\ell
                 \quad(\ell\in\Lambda).
\tag{C.WF.23}\hypertarget{tag-016K}{}\label{cdp:eq:C.WF.23}
\]
The exact finite map in this frame is
\[
 (v_\infty,q)\longmapsto
 \bigl((\det D)Dv_\infty+
       V^\infty_{\zeta^{-1}q}(a'_0)
       -l(q)V^\infty_{\zeta^{-1}q}(d),\
       \zeta^{-1}q\bigr).
\tag{C.WF.24}\hypertarget{tag-016L}{}\label{cdp:eq:C.WF.24}
\]
Since \((\det D)D H^\infty_q=
H^\infty_{\zeta^{-1}q}D\), the exponential
coordinate with frame \(H^\infty_q\) has central
affine vertex map \(p\mapsto Ap-b_0\).
Use the reflected old and subdivided fans
\(-\mathscr P_0\) and \(-\mathscr P\).
After the local base change \(v_{\mathrm{base}}=\kappa q\),
the central map is exactly
\[
 (z_\infty,v_{\mathrm{base}})
  \longmapsto
  (\overline a^{-1}v_{\mathrm{base}}^{-b_0}
                   z_\infty^A,\ \zeta^{-1}v_{\mathrm{base}}),
\tag{C.WF.25}\hypertarget{tag-016M}{}\label{cdp:eq:C.WF.25}
\]
and its period constants are \(\overline c_\gamma\).
Here \(\overline{a'_0}=\overline{a_0}-iRd\)
and \(-d\,l(q)=-d\,l(v_{\mathrm{base}})+iRd\),
so the recorded shifts cancel exactly.


\begin{enumerate}[label={$\langle2\rangle$\arabic*.},leftmargin=1.3em,labelsep=.5em]
\item  Finally apply the actual torus inversion
\(z_\infty\mapsto z_\infty^{-1}\) and relabel the
period group by \(\gamma\mapsto-\gamma\).
This sends the reflected fans to the original
ones, sends the map \eqref{cdp:eq:C.WF.25} to the conjugate
of \eqref{cdp:eq:C.WF.9}, and leaves the conjugate period
constants correct because
\(\overline c_{-\gamma}^{-1}=\overline c_\gamma\).
Thus the full coarse and subdivided central
log spaces, including the finite action, are
the complex conjugates of the prescribed data.
The two base coordinates satisfy
\(u\,v_{\mathrm{base}}=\kappa^2\).


\end{enumerate}

\item \par\medskip\noindent\textbf{Finite stabilizers and smooth coarse quotients.}
We check the local smoothness assertion rather than
inferring it from an abstract quotient.  At a point
of a retained coarse central chart, take its finite
point stabilizer after quotienting by \(\Gamma\).
The actual Kummer chart \eqref{cdp:eq:C.WF.4}--\eqref{cdp:eq:C.WF.5} shows that
its inertia fixes each normal ray and is the identity
on the tangent stratum: the latter cover is unramified.
This property is retained by the exact central log
comparison.  Its monomial lattice map is therefore
identity on the ray span and on its quotient, and,
being of finite order, is the identity on the whole
lattice.  We use a sufficiently small invariant
neighbourhood of this point.  A stabilizer of a
lower-dimensional cone which moves points on its
torus stratum is not used as a point stabilizer.
In the newly constructed chart the point stabilizer's
action on a character \(\chi^m\)
is multiplication by a holomorphic unit \(c_m(s)\).
These coefficients depend only on \(s\), by
\eqref{cdp:eq:C.WF.22} or \eqref{cdp:eq:C.WF.24}, and are multiplicative in \(m\).
Their central constants are the original stabilizer
characters.  For a generator of order \(d\), with
base character \(\xi\), let
\[
 L_m(s)=\log\!\frac{c_m(s)}{c_m(0)},\qquad
 B_m(s)=\frac1d\sum_{j=0}^{d-1}jL_m(\xi^js).
\tag{C.WF.26}\hypertarget{tag-016N}{}\label{cdp:eq:C.WF.26}
\]
The logarithm is the unique one vanishing at zero.
The finite-order identity gives
\(\sum_jL_m(\xi^js)=0\), and hence
\[
 B_m(\xi s)-B_m(s)=L_m(s).
\tag{C.WF.27}\hypertarget{tag-016O}{}\label{cdp:eq:C.WF.27}
\]
The coordinate change
\(\chi^m\mapsto e^{-B_m(s)}\chi^m\)
therefore makes every stabilizer coefficient constant.
The functions \(B_m\) are additive in \(m\), so this is
an actual automorphism of the saturated monoid algebra,
with the displayed inverse.  The base character has
constant coefficient, so its \(B_m\) is zero and the
change fixes \(s\).  Averaging has only been used for
this finite chart stabilizer, not for the developed
lift with \(g^M=g_{\gamma_0}\).


\begin{enumerate}[label={$\langle2\rangle$\arabic*.},leftmargin=1.3em,labelsep=.5em]
\item  We can now use precisely the original coarse invariant
monoid chart, retained in \eqref{cdp:eq:C.WF.5}--\eqref{cdp:eq:C.WF.6}:
\[
 \left(\OO_{X_M,\mathrm{coarse}}\right)^{\mathrm{Stab}}
        =\C\{x_1,x_2,x_3\},\qquad
       s^M=x_1^{m_1}\cdots x_r^{m_r}.
\tag{C.WF.28}\hypertarget{tag-016P}{}\label{cdp:eq:C.WF.28}
\]
The same equation holds in the new chart after the
explicit unit change \eqref{cdp:eq:C.WF.26}.  This proves smoothness
of the quotient total space and its exact component
multiplicities.  It also proves the embedded SNC
condition for its reduced fibre, including the original
component labels.  Reduction and finite invariants
commute in characteristic zero, so the central
coarsening and quotient identify the actual original
reduced analytic support and all its gluing maps.


\end{enumerate}

\item \par\medskip\noindent\textbf{Compactness and the precise retained scope.}
The determinant \eqref{cdp:eq:C.WF.18} and real fundamental
parallelepipeds give a proper middle torus family
on every compact annulus.  The bounded-unit
toric filling argument of \ref{cdp:st:C.2.UN.73} gives proper
fillings at both ends.  A toric coarsening has
the same fan support and is proper; the coarse
fillings therefore remain proper.  The exact
exponential maps above identify them with the
middle family on punctured discs.  They glue
to a Hausdorff compact normal root space over
\(\PP^1_s\).  Its finite action is the actual
\eqref{cdp:eq:C.WF.19}, extended by the displayed monomial
maps at both ends.  Its quotient is proper
over \(\PP^1_s/\mu_M=\PP^1\) and is smooth
by \eqref{cdp:eq:C.WF.28}; it is free over the punctured
base, where each quotient fibre is one of
the actual connected two-tori.


\begin{enumerate}[label={$\langle2\rangle$\arabic*.},leftmargin=1.3em,labelsep=.5em]
\item  The two marked reduced supports and their
multiplicities are the prescribed one and
its conjugate.  Euler additivity and the
zero Euler number of a smooth torus give
the stated total Euler number.  On the
original family, \eqref{cdp:eq:C.WF.6} recovers its full
nonreduced fibre by its actual invariant
analytic algebra, not just by a cycle.
The constructed compact family instead
uses the explicit variable period units
in \eqref{cdp:eq:C.WF.17}.  Their higher \(s\)-jets were
not identified with those of an arbitrary
original germ.  No such identification is
needed for the stated compact occurrence
of the marked reduced log support and
full multiplicity cycle.  This completes
all the asserted comparisons.



\end{enumerate}

\end{enumerate}
\end{proof}






\begin{proposition}[A section and ruling fibres give the exact transverse curve blowup]\hypertarget{tag-01EL}{}\label{cdp:st:C.2.UN.74}
\StatementRecord{cdp:st:C.2.UN.74}{proved}{01EL}{\SeriesRef{cdp:st:C.2.UN.64}, \SeriesRef{cdp:eq:C.RC.8}, \SeriesRef{cdp:eq:C.RC.9}, \SeriesRef{cdp:eq:C.RC.10}, \SeriesRef{cdp:eq:C.RC.12}; \cite[\S2, Theorem~2' and Proposition~9']{FujikiNakano1971}}

Let $f:U\to\Delta$ be a proper degeneration of connected
complex two-tori from a smooth threefold, with SNC reduced
central support and arbitrary positive full multiplicities.
Let $T$ be a component with a specified holomorphic
projective-line bundle $p:T\to B$ over a smooth compact
curve. Suppose its boundary in the reduced support is
exactly
\[
           C+F_1+\cdots+F_r,\qquad
           C\ \text{a section},\quad F_j=p^{-1}(b_j),
\tag{C.RC.1}\hypertarget{tag-00XN}{}\label{cdp:eq:C.RC.1}
\]
where the $b_j$ are distinct. The section is the
intersection with one component $P$, and $F_j$ is the
intersection with a component $Q_j$. The components
$Q_j$ may repeat at different $b_j$, but none is $P$.
All these are actual reduced transverse intersection
divisors; their intersections $C\cap F_j$ are the
specified triple points. Let the full multiplicities
of $T,P,Q_j$ be respectively $m_T,m_P,n_j$.

There is an ordinary smooth curve contraction along
the ruling, inverse to an ordinary blowup and equal
to an isomorphism off $T$, exactly when
$m_T=m_P$. In particular this always holds for a
full reduced fibre. Set their common value equal
to $m$. The resulting actual map is
\[
              \pi:U\longrightarrow U',\qquad
              f=f'\pi,\qquad \pi(T)=B\subset P'.
\tag{C.RC.2}\hypertarget{tag-00XR}{}\label{cdp:eq:C.RC.2}
\]
The threefold $U'$ is smooth and $f'$ is proper with
the same punctured torus family. Its reduced central
support is SNC. The map $P\to P'$ is an isomorphism.
Each $Q_j\to Q'_j$ is the ordinary point blowup at
$b_j\in B$, with all indicated points blown up when
the same component repeats. The curve $B$ lies in
the smooth surface $P'$, avoids every triple point
of the descended support, and meets its double
curves transversely, exactly at the $b_j$.

All full multiplicities are retained: the parent
has multiplicity $m$ and $Q'_j$ has multiplicity
$n_j$. At $b_j$ there are actual coordinates
$(p_0,q_0,r_0)$ with
\[
 \begin{gathered}
 P'=(p_0=0),\quad Q'_j=(q_0=0),\quad
 B=(p_0=r_0=0),\qquad f'=p_0^m q_0^{\,n_j},\\
 p_0=xy,\quad r_0=x,\quad q_0=z,\qquad
                         f=x^m y^m z^{\,n_j}
 \end{gathered}
\tag{C.RC.3}\hypertarget{tag-00XS}{}\label{cdp:eq:C.RC.3}
\]
in the blowup chart containing $C\cap F_j$.
Every other component and every neighbourhood away
from the indicated exceptional configurations are
unchanged.

Write
$L=N_{B/P'}$ and $M=N_{P'/U'}|_B$.
The actual normal extension and full-fibre
trivialization are
\[
 \begin{gathered}
 0\longrightarrow L\longrightarrow
       V=N_{B/U'}\longrightarrow M\longrightarrow0,
 \qquad (T,C)=(\PP_B(V),\PP_B(L)),\\
 N_{T/U}\simeq\OO_{\PP_B(V)}(-1),\qquad
 M^{\otimes m}\simeq
             \OO_B\left(-\sum_{j=1}^r n_jb_j\right).
 \end{gathered}
\tag{C.RC.4}\hypertarget{tag-00XT}{}\label{cdp:eq:C.RC.4}
\]
Thus the full fibre weights on the ruling-fibre
boundary enter the actual root line $M$; they are
not discarded. In particular $m$ divides
$\sum_jn_j$. For actual canonical orders $k_P,k_{Q_j}$
on the descended family, the section subline is
\[
 L\simeq K_B\otimes M^{-(k_P+1)}
                    \otimes\OO_B\left(-\sum_jk_{Q_j}b_j\right),
             \qquad k_T-k_P=1.
\tag{C.RC.5}\hypertarget{tag-00XU}{}\label{cdp:eq:C.RC.5}
\]
The entire extension in \eqref{cdp:eq:C.RC.4}, with its endpoint
identifications, is retained.

For a full reduced fibre, put $Z=\sum_jb_j$. Then
these maps specialize to
\[
 \begin{gathered}
 M\simeq\OO_B(-Z),\qquad
 L\simeq K_B\otimes
              \OO_B\left(\sum_j(k_P-k_{Q_j}+1)b_j\right),\\
 df'|_V:\ V\twoheadrightarrow
            \OO_B(-Z)\hookrightarrow\OO_B.
 \end{gathered}
\tag{C.RC.6}\hypertarget{tag-00XV}{}\label{cdp:eq:C.RC.6}
\]
If in addition $T$ is a local maximum of the
canonical order and $r>0$, the extension splits:
$V\simeq L\oplus M$, with its specified subline.

For the underlying reduced supports,
\[
              e(S)=e(S')+2-2g(B),\qquad
              t(S)=t(S')+r,
\tag{C.RC.7}\hypertarget{tag-00XW}{}\label{cdp:eq:C.RC.7}
\]
where $t$ counts triple points. The same contraction
exists on a compact total space. Conversely an
ordinary blowup of a smooth curve in a single
central component, transverse to its double curves
and avoiding the triple locus, has exactly the
maps, markings and multiplicities displayed above.
\end{proposition}



\begin{proof}

\begin{enumerate}[label={$\langle1\rangle$\arabic*.},leftmargin=*]
\item \par\medskip\noindent\textbf{The actual normal degree and the corrected contraction theorem.}
Restrict the principal full central fibre to $T$.
For a ruling fibre $\ell$ away from the marked
values, the resulting line isomorphism gives
\[
       m_T\deg(N_{T/U}|_\ell)=-m_P,\qquad
       N_{T/U}^{\,m_T}
           \simeq\OO_T\left(-m_PC-\sum_jn_jF_j\right).
\tag{C.RC.8}\hypertarget{tag-00XX}{}\label{cdp:eq:C.RC.8}
\]
If $m_T=m_P$, the restriction is precisely
$\OO_{\PP^1}(-1)$. The corrected theorem of
\cite[\S2, Theorem~2' and Proposition~9']{FujikiNakano1971},
in the form used in \ref{cdp:st:C.2.UN.64}, supplies the actual
smooth blowdown along $p:T\to B$. It identifies
$U$ with the ordinary blowup of a smooth curve
in a smooth threefold $U'$. Conversely an
ordinary such blowup has normal degree minus
one, so \eqref{cdp:eq:C.RC.8} forces $m_T=m_P$.

\item The equality $\pi_*\OO_U=\OO_{U'}$ gives the
unique holomorphic $f'$ in \eqref{cdp:eq:C.RC.2}. Its
properness follows because the inverse image
of a compact set in the disc is the image
under $\pi$ of the corresponding compact
set for $f$. The smooth fibres are unchanged.
It remains to prove the claimed structure
at the centre rather than infer it from a
numerical normal degree.

\par\medskip\noindent\textbf{The parent image is smooth and its whole surface is retained.}
The proper image $P'=\pi(P)$ is a hypersurface
in the smooth target. Its actual Cartier
pullback has the form
\[
                            \pi^*P'=P+aT.
\tag{C.RC.9}\hypertarget{tag-00XY}{}\label{cdp:eq:C.RC.9}
\]
Restriction to an ordinary ruling fibre
gives $0=1-a$, since $P\cap T=C$ is a
section and $T|_\ell=\OO_\ell(-1)$.
Thus $a=1$.

\item A local defining section of $P'$ vanishes
to first order along the generic centre.
Its first normal term on $\PP_B(N_{B/U'})$
has zero divisor exactly $P\cap T=C$,
by \eqref{cdp:eq:C.RC.9}. This term cannot vanish at
any whole fibre: $C$ is a section and
contains none. Its differential in the
normal directions is therefore nonzero
at every point of $B$. Hence $P'$ is
smooth along $B$. Since the centre is
a Cartier curve in $P'$, the actual
restriction of its blowup is an
isomorphism $P\to P'$. In particular
it identifies the section $C$ with
the original curve $B$.

\par\medskip\noindent\textbf{Each ruling-fibre boundary becomes a transverse double curve.}
Consider one of the other image components
$Q'=\pi(Q)$. It does not contain $B$:
otherwise its generic first normal
divisor would meet every ruling fibre,
whereas $Q\cap T$ is a finite union
of fibres. More explicitly, writing
$\pi^*Q'=Q+aT$ and restricting to a
generic ruling fibre gives $a=0$.
Thus $\pi^*Q'=Q$.

\item Restrict a local defining section of
$Q'$ to $B$. Its pullback to $T$
vanishes exactly on the specified
reduced fibres belonging to $Q$.
At $b_j$ that section therefore
has a simple zero on $B$. Its
differential is nonzero on the
tangent line of $B$. Consequently
$Q'$ is smooth at $b_j$ and is
transverse to $B$. The independent
normal differential of $P'$ makes
$P'$ and $Q'$ transverse as well.

\item There is only one such other
component at each $b_j$. A second
would have the same fibre in
$T\cap Q$, contradicting the
given SNC intersections and
\eqref{cdp:eq:C.RC.1}. Any further image
component passing through $b_j$
would likewise have a boundary
component on the exceptional
fibre; hence none is missing
from this count. Thus there is
no triple point of the descended
support on $B$.

\item Choose coordinates with $P'$ and
$Q'$ the first two coordinate
planes and $B$ the stated
coordinate curve. Transversality
just proved permits this choice.
The full fibre is a unit times
$p_0^m q_0^{n_j}$; a local root
of that unit absorbs it into
$p_0$ without changing these
marked loci. Blow up $(p_0,r_0)$.
The exact chart $p_0=xy,r_0=x$
gives \eqref{cdp:eq:C.RC.3}. On $Q'$ this
ideal is the maximal ideal
of $b_j$, so $Q$ is precisely
its ordinary point blowup,
with exceptional fibre $F_j$.
On $P'$ the ideal is Cartier,
giving the already established
isomorphism. These are the
claimed original neighbourhood
maps.

\par\medskip\noindent\textbf{The full normal extension and the weighted line.}
The normal sequence of
$B\subset P'\subset U'$
is the first row of \eqref{cdp:eq:C.RC.4}.
The exceptional divisor of
the ordinary blowup is the
lines projectivization of
this very $V$, with
ambient normal $\OO(-1)$.
The strict parent intersects
it in the specified subline
$\PP_B(L)$. In particular,
with these line conventions,
\[
                 \OO_T(C)\simeq
                      \OO_{\PP_B(V)}(1)\otimes p^*M.
\tag{C.RC.10}\hypertarget{tag-00XO}{}\label{cdp:eq:C.RC.10}
\]
The principal full fibre on
$U'$ is $mP'+\sum_Qm_QQ'$.
Restriction to $B$ gives
exactly the last isomorphism
of \eqref{cdp:eq:C.RC.4}, since each
$Q'$ meets $B$ simply at
its listed $b_j$. This
also follows by substituting
\eqref{cdp:eq:C.RC.10} into \eqref{cdp:eq:C.RC.8}.
Taking degrees proves the
stated divisibility. The
line $M$ and its actual
power trivialization are
retained, including any
torsion information.

\item Restrict the actual vertical
canonical evaluation section
on $U'$ to $B$. Its line
isomorphism is
\[
 K_{U'}|_B\simeq
       M^{k_P}\otimes
             \OO_B\left(\sum_jk_{Q_j}b_j\right).
\tag{C.RC.11}\hypertarget{tag-00XP}{}\label{cdp:eq:C.RC.11}
\]
Other components have a
nonvanishing defining section
on $B$ and contribute their
specified trivializations.
Adjunction for $B\subset P'$
then proves the formula for
$L$ in \eqref{cdp:eq:C.RC.5}. The
ordinary curve-blowup
Jacobian contributes $T$
once to $K_U$; the generic
centre lies only in $P'$.
Thus $k_T=k_P+1$.
Changing the canonical
generator by a power of $f'$
adds the full multiplicities
to the orders, and \eqref{cdp:eq:C.RC.4}
shows that \eqref{cdp:eq:C.RC.5} remains
the same line isomorphism.

\item In the full reduced case
$m=n_j=1$, so $M=\OO_B(-Z)$.
The derivative of
$f'=p_0q_0$ along the
normal bundle of $B$
is the quotient normal
map to $M$, followed by
its inclusion in $\OO_B$
with simple zero at $b_j$.
Away from the marked
points it is surjective.
This proves all maps in
\eqref{cdp:eq:C.RC.6}. For $m>1$ the
ordinary derivative of
$f'$ is zero along the
generic centre; the
quotient in \eqref{cdp:eq:C.RC.4}
instead is the actual
normal map of $P'$,
and has already been
specified independently
of that derivative.

\item For the splitting assertion,
the extension group is the
actual space $H^1(B,L\otimes
M^{-1})$. Equations \eqref{cdp:eq:C.RC.6}
and Serre duality identify
its dual with
\[
 H^0\left(B,\OO_B\left(
                   \sum_j(k_{Q_j}-k_P-2)b_j\right)\right).
\tag{C.RC.12}\hypertarget{tag-00XQ}{}\label{cdp:eq:C.RC.12}
\]
At a local maximum
$k_{Q_j}\le k_T=k_P+1$.
Every coefficient in
\eqref{cdp:eq:C.RC.12} is at most
minus one. When $r>0$
the line has negative
degree and no section.
The extension group is
zero, which proves the
actual splitting with
the marked subline.
When $r=0$ the extension
is precisely the one
retained in \ref{cdp:st:C.2.UN.64}.

\par\medskip\noindent\textbf{Euler, triples, the converse and compact descent.}
The actual blowup gives
$S\setminus T\simeq S'\setminus B$,
whereas $T$ is a
projective-line bundle
of Euler number $2e(B)$.
Hence $e(S)-e(S')=e(B)$,
independently of the
full multiplicities.
At each $b_j$ the chart
\eqref{cdp:eq:C.RC.3} has exactly one
triple point upstairs
and only a smooth double
curve downstairs. There
are no other changes,
which proves \eqref{cdp:eq:C.RC.7}.

\item Conversely, blowing up the
indicated transverse curve
uses exactly these charts.
The exceptional multiplicity
is the generic parent
multiplicity $m$, and
every other weight is
unchanged. The section
and ruling fibres are
exactly the intersections
in \eqref{cdp:eq:C.RC.1}, so all
the displayed maps and
markings are recovered.
As in \ref{cdp:st:C.2.UN.64}, the actual
local blowdowns glue to
the identity off $T$.
This proves the compact
version as well.

\end{enumerate}
\end{proof}






An \emph{overlap-compatible modification record} below retains the
actual analytic charts, all centres and blowup maps, and the
isomorphisms used to reglue them. In a local replacement
$X=U\cup V$, the modified chart $U'$ has a specified overlap
$W'\simeq W=U\cap V$; this isomorphism is the restriction of
the recorded composite map on $U$. The replacement is
$X'=U'\cup_W V$. Thus this datum specifies an analytic
operation, including its neighbourhoods. It is not merely a
list of component graphs or normal degrees. The records used
here are the finite local $\sigma$-records in
\cite[\S\S4.2--4.14]{Kulikov1977}, with the corrections and
extra fibre blowups in \cite[\S\S8--11]{Kulikov1981}.

\begin{proposition}[The Hopf-minimum branch is generated by canonical-flat cores and exact sigma records]\hypertarget{tag-01EM}{}\label{cdp:st:C.2.UN.75}
\StatementRecord{cdp:st:C.2.UN.75}{proved}{01EM}{\SeriesRef{cdp:st:C.2.UN.65}, \SeriesRef{cdp:st:C.2.UN.71}, \SeriesRef{cdp:st:C.2.UN.74}, \SeriesRef{cdp:st:C.2.UN.73}, \SeriesRef{cdp:eq:C.KG.6}, \SeriesRef{cdp:st:C.2.UN.19}, \SeriesRef{cdp:st:C.2.UN.64}, \SeriesRef{cdp:eq:C.KG.8}, \SeriesRef{cdp:st:C.2.UN.68}; \cite[Theorem~1']{Kulikov1981}; \cite[original Russian ed., pp.~1008--1009]{Kulikov1977}; \cite[\S4.4]{Kulikov1977}; \cite[\S\S4.10.7--4.10.8,4.14.5]{Kulikov1977}; \cite[\S8]{Kulikov1981}; \cite[\S10]{Kulikov1981}; \cite[\S4.17]{Kulikov1977}; \cite[\S2]{Kulikov1981}; \cite[\S4.10.8, original Russian ed., pp.~1032--1035]{Kulikov1977}}

Let $f:X\to\Delta$ be a proper degeneration of connected
complex two-tori, with smooth total threefold and full
reduced simple-normal-crossing central fibre $S=\sum_iV_i$.
Choose a generator $\Omega$ of $f_*K_X$ near zero and write
\[
             \operatorname{div}(\Omega)=\sum_i k_iV_i,
                  \qquad k_i\ge0,\quad \min_i k_i=0.
\tag{C.KG.1}\hypertarget{tag-00P8}{}\label{cdp:eq:C.KG.1}
\]
Let $\mathcal H$ be the set of indices whose component has
a Hopf surface as its smooth minimal model. The following
two conditions are equivalent:
\begin{enumerate}
\item $k_i=0$ for every $i\in\mathcal H$.
\item There is a finite overlap-compatible modification
record from $f$ to a smooth proper torus degeneration
$f_0:X_0\to\Delta$ whose full central fibre $S_0$ is
reduced simple NC and whose canonical line is trivial.
\end{enumerate}
The constructed comparison is an actual bimeromorphic map
over the same disc and induces an isomorphism
\[
                X\setminus S\ \simeq\ X_0\setminus S_0
                         \quad\hbox{over }\Delta^*.
\tag{C.KG.2}\hypertarget{tag-00PA}{}\label{cdp:eq:C.KG.2}
\]
Intermediate charts retain their normalization branches;
in particular a type-II move may temporarily give a
self-intersection of a global component.

The endpoint has precisely one of the following forms,
with the actual markings and line maps of the indicated
statements:
\begin{enumerate}
\item a smooth two-torus;
\item a cycle of the marked diagonal Hopf and degree-zero
elliptic ruled pairs in \ref{cdp:st:C.2.UN.65}, with dual edge normals and
the actual complement-lattice product $M=I$;
\item a positive toric central log space in \ref{cdp:st:C.2.UN.71}, given by
$(N,\Gamma,\mathscr P,c)$, with Euler number
$2[N:\Gamma]$ and the specified full toric component pairs.
\end{enumerate}
Every degeneration in the stated branch is recovered by
reversing its finite record on its actual endpoint
neighbourhood. Conversely every such reverse record with
a smooth full reduced simple-NC final fibre gives a
degeneration in this branch. These are inverse analytic
constructions, retaining every centre, overlap and normal
extension.

Here is the exact Euler bookkeeping for these records.
Orient a recorded blowup arrow positively from its target
to its blowup, and negatively in the reverse direction.
For its actual centre $L$ and ideal $I$, put
\[
 \begin{split}
  E_I&=\operatorname{Proj}_L\left(\bigoplus_{a\ge0}
                                      I^a/I^{a+1}\right),\\
  \delta(I)&=e_c(E_I)-e_c(L)
       =\int_L\bigl(e((E_I)_x)-1\bigr)\,de_c(x).
 \end{split}
\tag{C.KG.3}\hypertarget{tag-00PB}{}\label{cdp:eq:C.KG.3}
\]
Local terms are evaluated on a finite constructible
subdivision subordinate to the retained charts. Restored
overlaps cancel. If $\epsilon_\nu\in\{1,-1\}$ records
the directions in the forward reduction to $S_0$, then
\[
             e(S)=e(S_0)-\sum_\nu\epsilon_\nu\delta(I_\nu).
\tag{C.KG.4}\hypertarget{tag-00PC}{}\label{cdp:eq:C.KG.4}
\]
For an ordinary smooth centre in a smooth threefold,
$\delta=2$ for a point and $\delta=2-2g(L)$ for a compact
curve. Both type-I and type-II moves have total defect
zero. An ordinary ruled-section contraction in \ref{cdp:st:C.2.UN.74} has
Euler change $-(2-2g(B))$ and removes its specified number
of triple points. Thus the exact zero-Euler criterion is
\[
 \sum_\nu\epsilon_\nu\delta(I_\nu)=
 \begin{cases}
 0,&S_0\text{ is of type 1 or 2},\\
 2[N:\Gamma],&S_0\text{ is of type 3}.
 \end{cases}
\tag{C.KG.5}\hypertarget{tag-00PD}{}\label{cdp:eq:C.KG.5}
\]
The normal-cone formula also applies to auxiliary
singular-centre charts; their contribution is not
replaced by a smooth-curve formula.

For a given compact total space over a compact curve,
the reduction and its inverse take place on that same
compact base, by the actual isomorphism \eqref{cdp:eq:C.KG.2} on an
annulus. A reverse record on an actual compact core
realization of \ref{cdp:st:C.2.UN.65} or \ref{cdp:st:C.2.UN.73} likewise gives a compact
realization. Its total Euler change is the same as
the central-fibre change in \eqref{cdp:eq:C.KG.4}.
\end{proposition}



\begin{proof}
\begin{enumerate}[label={$\langle1\rangle$\arabic*.},leftmargin=*]
\item 
\par\medskip\noindent\textbf{The canonical generator and its normalized orders.}
The proper direct image $f_*K_X$ is a torsion-free
rank-one coherent sheaf on the disc, hence a line.
On $\Delta^*$ its restriction and base change are the
one-dimensional spaces of top forms on the tori.
A generator therefore has no zero away from $S$.
Its divisor is the effective divisor in \eqref{cdp:eq:C.KG.1}.
If every $k_i$ were positive, $\Omega/f$ would be
holomorphic on the smooth total space, since the
\noindent\textbf{Step 1. full} fibre is reduced. This would contradict
that $\Omega$ generates the direct-image line.
Consequently its minimum is zero. Another generator
is $a(t)\Omega$ for a holomorphic unit $a$, so these
orders are independent of that choice.

For later comparison, let $\phi:X\dashrightarrow X'$
be an actual bimeromorphic map over the disc between
smooth total spaces. Pullback of a holomorphic top
form is holomorphic wherever $\phi$ is holomorphic.
The indeterminacy has codimension at least two;
Hartogs extension of the coefficient in smooth
coordinate volume frames gives a holomorphic form
on all of $X$. The inverse map gives the inverse
pullback, since their compositions agree on their
common biholomorphic open set. The construction
works over every open subdisc, and gives the actual
isomorphism
\[
                f'_*K_{X'}\xrightarrow{\ \phi^*\ } f_*K_X.
\tag{C.KG.6}\hypertarget{tag-00PE}{}\label{cdp:eq:C.KG.6}
\]
In particular normalized generators agree under
pullback up to a base unit. At the generic point
of a divisor which survives the comparison,
the map is a biholomorphism of the corresponding
divisorial local rings, so its canonical order is
unchanged.

\par\medskip\noindent\textbf{Why every Hopf-birational component is retained.}
A holomorphic map from $\PP^1$ to a Hopf surface
is constant: lift it to the universal cover
$\C^2\setminus\{0\}$ and use the two constant
coordinate functions. Every rational curve on
a point blowup of that surface is therefore
exceptional over the Hopf minimal model. A finite
sequence of point blowups has only finitely many
such curves. These surfaces cannot be covered by
rational curves. They are also not projective:
their first Betti number is one, unchanged by
point blowups, whereas a projective surface has
even first Betti number.

In the recorded local $\sigma$-maps, a divisorial
exceptional component over a general point of a
curve centre is ruled by complete rational
fibres. A component lying over an isolated
point is projective, as a component of a
projectivized normal cone. Here the contradiction
with the original compact component uses the
proper composite graphs, not a comparison of
local analytic patches. Suppose a strict
transform were removed for the first time.
For a curve centre, choose a general complete
rational exceptional fibre through a generic
strict-transform point. The proper preceding
graph maps this curve meromorphically back to
the original compact component. The map is
nonconstant, since at that generic point the
strict-transform comparison is a local
biholomorphism. Varying the fibre gives
rational curves through a dense open subset
of the original component, contradicting the
finite exceptional-curve description above.
For a point centre, restrict the proper
preceding graph to the compact projective
exceptional component containing the generic
strict-transform piece. Its image back in
the original compact component is closed
and contains that open piece, and therefore
is the whole component. The comparison has
generic degree one. Resolving this surface
graph gives a bimeromorphic comparison with
a projective surface; point blowups preserve
$b_1$, again contradicting $b_1=1$.
Thus the strict transforms through the
finite chart record retain a prime divisor
of that birational type. The restored
overlap maps identify these strict transforms
by their actual generic-point maps. Hence a
component in $\mathcal H$ survives to the
endpoint.

At a canonical-flat smooth endpoint a nowhere-zero
three-form generates $f_{0*}K_{X_0}$, so every
canonical order is zero. Equation \eqref{cdp:eq:C.KG.6} then
proves $k_i=0$ for every $i\in\mathcal H$.
This proves the necessary minimum condition for
the full Hopf birational class, including
arbitrary point blowups. No identification of
different conventions for the notation
\noindent\textbf{Step 2. VII$_0$} is used here.

\par\medskip\noindent\textbf{The finite reduction and the actual allowed charts.}
\ref{cdp:st:C.2.UN.19} identifies every actual class-VII minimal
model in a torus degeneration as Hopf. Thus the
minimum condition just stated implies the
condition in \cite[Theorem~1']{Kulikov1981}.
That theorem applies to proper analytic good
degenerations with trivial fibre canonical
class; its hypotheses do not include
projectivity. The definition of good in
\cite[original Russian ed., pp.~1008--1009]{Kulikov1977} is exactly the smooth
total space and full reduced NC fibre with
smooth components used here.

We specify the comparison maps in that proof.
For a nonboundary $(-1)$-curve $A\subset V_i$
meeting one double curve transversely once
and avoiding the remaining boundary, the
principal-fibre normal sequence is
\[
 0\longrightarrow\OO_{\PP^1}(-1)
 \longrightarrow N_{A/X}
 \longrightarrow\OO_{\PP^1}(-1)
 \longrightarrow0.
\tag{C.KG.7}\hypertarget{tag-00PF}{}\label{cdp:eq:C.KG.7}
\]
It splits because $H^1(\PP^1,\OO)=0$.
Blow up $A$; the exceptional divisor is
$\PP^1\times\PP^1$ with normal
$\OO(-1,-1)$. Contract its other ruling
using the corrected contraction theorem in
\ref{cdp:st:C.2.UN.64}. This transfers the point blowup to
the neighbour. The exact adjunction equation
gives $k_i=k_j$ for these two components.
This is the actual type-I map.

For a double curve which is $(-1)$ on both
sides, the same normal bundle calculation
gives the same two-ruling construction.
The principal-fibre line has exactly two
triple points on that curve. If the opposite
components there have orders $k_3,k_4$,
adjunction gives
\[
                         k_1+k_2=k_3+k_4.
\tag{C.KG.8}\hypertarget{tag-00PG}{}\label{cdp:eq:C.KG.8}
\]
The ordinary blowup has exceptional full
multiplicity two; the other contraction
removes the old double curve and introduces
the new one between the two opposite
branches. If those branches belong to one
global component, both normalization branches
are retained, as in
\cite[\S4.4]{Kulikov1977}. One does not
discard that self-incidence or call this
intermediate fibre embedded simple NC.

The other contraction charts are the actual
ruled section-and-fibre maps of \ref{cdp:st:C.2.UN.74} and
the plane-line map of \ref{cdp:st:C.2.UN.64}. For a reducible
ruling fibre, the local contractions and
inverse-order restorations in
\cite[\S\S4.10.7--4.10.8,4.14.5]{Kulikov1977}
give the retained finite chart record.
The corrected proof in
\cite[\S8]{Kulikov1981} also blows up the
specified missing fibre components before
contracting the chosen ruled surface.
The new ruled components have canonical
order $k+1$ when the removed component has
order $k$. For a maximum-order component
no such additions are needed.

Here termination is an actual finite
induction, rather than a decrease of an
unspecified complexity. Let $R=\max k_i$.
Every removed old component has level
$k\in\{1,\ldots,R\}$. It produces only
finitely many new components, all at level
$k+1$, and produces none when $k=R$.
There are finitely many level-one removals;
therefore finitely many level-two objects
are ever produced. Continuing through the
fixed finite set of levels proves that all
removals terminate. The components left
have order zero. The corrected endpoint
argument in \cite[\S10]{Kulikov1981}, using
\cite[\S4.17]{Kulikov1977}, gives a smooth
total space and smooth reduced components.
In particular \eqref{cdp:eq:C.KG.8} prevents a minimum
component from acquiring a self-incidence
in a type-II move whose removed edge has
an order above the minimum. The endpoint
generator is therefore nowhere zero.

\par\medskip\noindent\textbf{The source's actual overlap maps glue global graphs.}
The bare two-cover definition of modification
in \cite[\S2]{Kulikov1981} is not used to
infer a global map. The construction itself
provides the following additional data.
In \cite[\S4.10.8, original Russian ed., pp.~1032--1035]{Kulikov1977}
one keeps $V$ fixed, performs the displayed
finite contractions and restorations in
$U$, and obtains $U'$ with an overlap
$W'$ isomorphic to the original $W$ by
the restriction of that same composite.
Section4.14.5 repeats this construction.
The added local fibre blowups of
\cite[\S8]{Kulikov1981} leave this overlap
unchanged.

For a finite $\sigma$-chain, its composite
graph is obtained as the closure of the
common biholomorphic graph in the proper
image of the fibre product of its arrow
graphs. Both projections are proper.
Over $W$ this graph is precisely the
graph of the restored isomorphism.
Glue it to the identity graph over $V$.
The graphs agree on their overlap, and
properness is local on each projection
target. The glued graph is consequently
proper over both $U\cup V$ and
$U'\cup_W V$. It gives actual mutually
inverse bimeromorphic maps. All its arrows
are over the same disc and all their
centres are central. Their restriction
over $\Delta^*$ is an isomorphism.
Iteration proves \eqref{cdp:eq:C.KG.2} and justifies
the generator comparison \eqref{cdp:eq:C.KG.6}.

This argument does not assert a single
global factorization into ordinary
blowups. The finite open covers and
their induced overlap maps are part
of the stated modification record.
Auxiliary singular charts in the
restoration are also retained rather
than silently replaced by smooth ones.

\par\medskip\noindent\textbf{Endpoint classification and the converse record.}
Apply \ref{cdp:st:C.2.UN.65} to the actual endpoint.
Its zero-Euler alternatives give types1
and2, including all the marked line
and integral-lattice comparisons.
Its positive case is \ref{cdp:st:C.2.UN.68}; \ref{cdp:st:C.2.UN.71}
identifies its actual central log
space with the specified toric data
by the displayed chart and residue
maps. This proves the stated complete
endpoint list.

Reverse each recorded arrow and each
overlap isomorphism. They reconstruct
the original neighbourhood literally.
Starting instead with any actual
canonical-flat endpoint and such a
reverse record, all off-central maps
remain isomorphisms; the general
fibres are the same connected tori.
At a smooth full reduced final fibre
the argument following \eqref{cdp:eq:C.KG.6} shows
that every Hopf-birational component
has order zero. This proves both
directions of the generation statement.
In this statement the endpoint is an
actual analytic family. The maps do
not replace its neighbourhood by a
different realization merely because
their central log spaces agree.

\par\medskip\noindent\textbf{The exact Euler defect, including singular centres.}
For an actual blowup with centre ideal
$I$, its exceptional inverse image is
the projectivized normal cone in
\eqref{cdp:eq:C.KG.3}. Outside $L$ it is an
isomorphism. Since $L$ is central,
the reduced central inverse image
replaces precisely $L$ by $E_I$.
Additivity of compactly supported
Euler characteristic proves \eqref{cdp:eq:C.KG.3}
for this arrow, including its
constructible-fibre integral.
For these computations refine the
retained analytic charts by relatively
compact subanalytic coordinate
neighbourhoods, with closures inside
the original charts. Compactness of
the central fibre gives a finite
subcover. Subdivide its overlaps and
centres compatibly, and use inverse
images under the retained proper
graphs; all pieces have finite
subanalytic triangulations. This
refinement also defines the local
terms when a displayed centre is an
open part of a curve. Euler numbers
of arbitrary analytic open covers
are not used.

In one replacement, evaluated on this
finite-type refinement,
$U\cup V\rightsquigarrow
U'\cup_WV$, inclusion--exclusion gives
\[
 \begin{split}
 e(S')-e(S)
   &=e_c(S'\cap U')-e_c(S\cap U)\\
   &=\sum_{\nu\text{ in this chart}}
                         \epsilon_\nu\delta(I_\nu),
 \end{split}
\tag{C.KG.9}\hypertarget{tag-00PH}{}\label{cdp:eq:C.KG.9}
\]
because the unchanged $V$ and the
isomorphic restored overlaps cancel.
Equivalently one applies this equality
to the finite subdivision just
described and sums; no boundary
piece is counted twice. The local
arrow contributions telescope, and
then the global replacements
telescope, proving \eqref{cdp:eq:C.KG.4}.

For a smooth centre of dimension $a$
in a smooth threefold the normal
cone is its rank-$(3-a)$ normal
bundle, so its exceptional fibre
is $\PP^{2-a}$. Thus
\[
                \delta(I)=(2-a)e_c(L).
\tag{C.KG.10}\hypertarget{tag-00P9}{}\label{cdp:eq:C.KG.10}
\]
For typeI and typeII the common
blowup replaces a compact $\PP^1$
by $\PP^1\times\PP^1$ from either
side, so the two contributions
$+2$ and $-2$ cancel. The ruled
and plane changes follow from
\ref{cdp:st:C.2.UN.74} and \ref{cdp:st:C.2.UN.64}. In an auxiliary
singular-centre chart \eqref{cdp:eq:C.KG.3}
remains the formula: for example,
the blowup of the three coordinate
axes in $\C^3$, with ideal
$(xy,xz,yz)$, has fibre $\PP^1$
over a general point of each axis
and fibre $\PP^2$ over their
common origin. The last defect
is two rather than one. This
illustrates why the full normal
cone was retained.
Substitution of the three endpoint
Euler numbers proves \eqref{cdp:eq:C.KG.5}.

\par\medskip\noindent\textbf{Compact scope.}
Restrict the given compact family
to a small disc at the chosen
value. On a smaller annulus use
the actual isomorphism \eqref{cdp:eq:C.KG.2}
to glue the new proper filling
to the unchanged complement.
The two proper pieces give a
Hausdorff proper analytic family;
their smooth charts coincide on
the overlap. This produces the
claimed compact replacement,
and the inverse record recovers
the original compact family.
The same construction starts
from the explicit compact cores
of \ref{cdp:st:C.2.UN.65} and \ref{cdp:st:C.2.UN.73} and any reverse
record on their actual tubes.
Euler additivity cancels the
unchanged complement and annulus,
so the total-space Euler change
is precisely the central one.

\end{enumerate}
\end{proof}






\begin{proposition}[Actual minimalization at the largest canonical order]\hypertarget{tag-01EO}{}\label{cdp:st:C.2.UN.77}
\StatementRecord{cdp:st:C.2.UN.77}{proved}{01EO}{\SeriesRef{cdp:st:C.2.UN.75}, \SeriesRef{cdp:eq:C.HM.2}, \SeriesRef{cdp:eq:C.HM.4}, \SeriesRef{cdp:eq:C.HM.3}, \SeriesRef{cdp:eq:C.HM.7}, \SeriesRef{cdp:eq:C.HM.9}, \SeriesRef{cdp:eq:C.HM.5}, \SeriesRef{cdp:eq:C.HM.8}; \cite[\S4.11.2(III)]{Kulikov1977}; \cite[\S4.5]{Kulikov1977}; \cite[\S\S4.10.7--4.10.8,4.14]{Kulikov1977}; \cite[\S8]{Kulikov1981}; \cite[Lemma2.19]{Kulikov1977}; \cite[\S4.11.2]{Kulikov1977}; \cite[\S4.13]{Kulikov1977}; \cite[Lemma4.14]{Kulikov1977}; \cite[\S4.16.2]{Kulikov1977}; \cite[\S9]{Kulikov1981}}

Let \(f:X\to\Delta\) be a proper degeneration of connected
complex two-tori with smooth total space and full reduced
simple-NC central fibre.  Normalize its canonical generator
as in \ref{cdp:st:C.2.UN.75} and put
\[
        \operatorname{div}(\Omega)=\sum k_iV_i,\qquad
        \min k_i=0,\qquad M=\max k_i .
\tag{C.HM.1}\hypertarget{tag-00MG}{}\label{cdp:eq:C.HM.1}
\]
There is a finite overlap-compatible modification record,
over the same disc and unchanged over its punctured part,
whose endpoint has smooth total space and full reduced
normal-crossing fibre, such that every component with
canonical order \(M\) and Hopf minimal model is itself a
smooth minimal Hopf surface.  These Hopf components are
embedded, their boundary curves are disjoint smooth
elliptic curves, and no triple point meets them.
Other endpoint components may retain the two normalized
branches of an NC self-intersection.

All original Hopf-birational prime divisors survive, with
their original orders.  Every whole-component removal block
removes a ruled component of positive order \(q\);
its preliminary additions, when needed, are finitely many
ruled components of order \(q+1\le M\).  There are no such
additions when \(q=M\).  Thus the retained reduced degeneration stages have
no global component of order greater than \(M\).
Auxiliary common blowups in a type-I or type-II diagram
retain their actual multiplicities and orders; the bound
just stated does not apply to those auxiliary divisors.

The record retains its actual centres, local contractions,
normalization branches and restored overlap maps.  Its
reverse recovers the original family.  For a given compact
family the same operations take place on that compact
base, by the actual annular comparison.  Their Euler
changes are exactly the normal-cone terms of \ref{cdp:st:C.2.UN.75}.

The endpoint assertion concerns the maximal-order Hopf
components.  It makes no holomorphic inverse-logarithm
assertion for an arbitrary neighbourhood of a remaining
Hopf leaf.
\end{proposition}



\begin{proof}\leavevmode\par



 
\begin{enumerate}[label={$\langle1\rangle$\arabic*.},leftmargin=0pt,labelsep=.5em]
\item \par\medskip\noindent\textbf{The surface divisor and coefficient-one curves.}
At every smooth-total NC stage, let \(H\) be the smooth
normalization of a component of order \(M\) whose minimal
model is Hopf.  The actual adjunction map, with both
preimages of each self-double curve retained, gives
\[
       -K_H=D_H=\sum_C a_C C,\qquad
       a_C=M-k_{\mathrm{other}(C)}+1\ge1 .
\tag{C.HM.2}\hypertarget{tag-00MI}{}\label{cdp:eq:C.HM.2}
\]
A self-double branch has coefficient one.  The equality
is the restriction of the common meromorphic volume
form; it retains its line and divisor, rather than only
its numerical class.


\begin{enumerate}[label={$\langle2\rangle$\arabic*.},leftmargin=1.3em,labelsep=.5em]
\item  Let \(\pi:H\to H_{\min}\) be the actual sequence of point
blowdowns to its Hopf minimal surface.  A map from
\(\PP^1\) to a Hopf surface is constant: lift it to
\(\C^2\setminus\{0\}\) and use its two constant coordinate
functions.  Thus every rational curve on \(H\) is one
of the finitely many exceptional curves of \(\pi\).
Every compact curve on \(H_{\min}\) is a smooth elliptic
curve.  Indeed rational normalization is impossible,
and \(H^2(H_{\min},\Q)=0\) and adjunction give arithmetic
genus one; the normalization genus therefore equals
one and its delta invariant is zero.  Distinct such
curves are disjoint, since their intersection number
is zero and all local intersection multiplicities are
positive.


\item  Push forward \eqref{cdp:eq:C.HM.2} through the point blowdowns.  At
each reverse point blowup \(\sigma\) one has the exact
divisor formula
\[
             D_{\mathrm{new}}=\sigma^*D_{\mathrm{old}}-E.
\tag{C.HM.3}\hypertarget{tag-00MJ}{}\label{cdp:eq:C.HM.3}
\]
Every intermediate divisor is effective: a negative
coefficient would persist on its strict transform in
the final effective divisor.  Starting from the
disjoint elliptic support on \(H_{\min}\), all these
blowups consequently have centres on the positive
support.  The positive support stays SNC.  At a smooth
point of a component of coefficient \(a\), the new
coefficient is \(a-1\); at a node of coefficients
\(a,b\), it is \(a+b-1\).  Coefficient-zero exceptional
curves are omitted from the positive support, but are
retained as curves of the surface.


\item  It follows by this actual point-map induction that
\[
 \begin{gathered}
 \text{two coefficient-one positive curves never meet;}\\
 \text{a rational coefficient-one curve is a leaf
 of the positive boundary,}\\
 \text{and an elliptic coefficient-one curve is
 disjoint from that boundary.}
 \end{gathered}
\tag{C.HM.4}\hypertarget{tag-00MK}{}\label{cdp:eq:C.HM.4}
\]
For the first assertion a smooth blowup creates
\((a,a-1)\), never \((1,1)\); a node blowup could
create a new \((1,1)\) only from an old \((1,1)\).
For the second assertion a new rational coefficient
one is first created at a smooth point of coefficient
two, and has valence one.  Subsequent node blowups
preserve that valence, and a smooth blowup at its
points has coefficient zero.  An original coefficient
one elliptic curve likewise gains no positive boundary
at a smooth blowup.


\item  These facts also rule out a rational self-double curve
on a maximal Hopf component at any NC stage in the
retained class.  Its two normalized preimages, as
specified by \cite[\S4.11.2(III)]{Kulikov1977}, have
coefficient one.  Each is rational and has exactly one
boundary intersection by \eqref{cdp:eq:C.HM.4}; hence the image has
one triple point.  Each rational curve has negative
self-intersection on the blown-up Hopf, whereas the
actual two-normal-bundle formula would give
\[
              C_1^2+C_2^2=-1,
\tag{C.HM.5}\hypertarget{tag-00ML}{}\label{cdp:eq:C.HM.5}
\]
an impossibility.  A triple point with three branches
of this same component would already give two
coefficient-one curves meeting, also impossible.
The modifications below can create new self-double
curves only from rational exceptional curves.
They therefore create no elliptic self-double curve
on a maximal Hopf starting from the given embedded
components.  At every subsequent smooth-total NC
stage these Hopf components are consequently embedded.


\end{enumerate}

\item \par\medskip\noindent\textbf{The exact last-point alternatives.}
Suppose such an \(H\) is not minimal and choose a last
point blowup in the actual map \(\pi\).  Its exceptional
curve \(E\) is a smooth \((-1)\)-curve.  Formula
\eqref{cdp:eq:C.HM.3} leaves exactly three cases.


\begin{enumerate}[label={$\langle2\rangle$\arabic*.},leftmargin=1.3em,labelsep=.5em]
\item  First, \(E\) is not a boundary curve.  It was created
at a smooth point of a coefficient-one boundary \(C\),
and meets that curve once.  The ambient normal
sequence is
\[
 0\longrightarrow\OO_{\PP^1}(-1)
 \longrightarrow N_{E/X}
 \longrightarrow\OO_{\PP^1}(-1)\longrightarrow0 .
\tag{C.HM.6}\hypertarget{tag-00MM}{}\label{cdp:eq:C.HM.6}
\]
Its splitting and the ordinary two-ruling diagram of
\ref{cdp:st:C.2.UN.75} give the actual type-I move.  The neighbour across
\(C\) has order \(M\).  It is not Hopf-birational.
If \(C\) is elliptic, its square on \(H\) is at most
\(-1\) and it contains no triple point; its square on
the neighbour is therefore positive.  If \(C\) is
rational, it was already exceptional before this last
blowup, so its current square is at most \(-2\);
\eqref{cdp:eq:C.HM.4} gives one triple point, and its square on
the neighbour is again positive.  Such squares are
impossible for the corresponding elliptic or rational
curves on a blown-up Hopf.  The move removes one
point blowup from \(H\), adding it to a non-Hopf
component.


\item  Second, \(E\) was created at a node of two positive
boundary curves of coefficients \(a,b\).  It has
coefficient \(a+b-1\) and two triple points.  The
two-normal-bundle equation makes its square on the
other component also \(-1\), so the actual type-II
move of \ref{cdp:st:C.2.UN.75} applies.  It blows down \(E\) on both
sides and blows up the two opposite normalized
branches.  If an opposite component has order \(M\),
the associated curve \(C\) on \(H\) has coefficient
one.  It is a rational leaf by \eqref{cdp:eq:C.HM.4}, and its
square after the last node blowup is at most \(-2\).
Its square on that opposite component is
\[
                 -1-C_H^2\ge1 .
\tag{C.HM.7}\hypertarget{tag-00MN}{}\label{cdp:eq:C.HM.7}
\]
Hence that opposite component is not Hopf-birational.
This proves the same assertion when both opposite
branches belong to one global component: their actual
normalization branches are used in \eqref{cdp:eq:C.HM.7}.


\item  Let
\[
       h=\sum_{\substack{k_i=M\\
                     V_i\ \mathrm{Hopf\text{-}birational}}}
                              b_2(V_i).
\tag{C.HM.8}\hypertarget{tag-00MO}{}\label{cdp:eq:C.HM.8}
\]
At the smooth NC stages just described, this is the
total number of surface point blowups on these Hopfs.
Each of the two moves above strictly decreases \(h\).
No assertion that \(h\) decreases inside the longer
ruled-removal blocks below is needed.


\item  Third, \(E\) was created at a smooth point of a
positive curve of coefficient \(a\ge2\).  It has
coefficient \(a-1\) and one triple point.  Write \(Q\)
for its neighbour and \(P\) for the third component.
The exact normal degrees and canonical orders are
\[
 \begin{split}
 E_H^2&=-1,\qquad E_Q^2=0,\\
 q:=k_Q&=M-a+2,\qquad k_P=q-1,\\
                         &1\le q\le M .
 \end{split}
\tag{C.HM.9}\hypertarget{tag-00MP}{}\label{cdp:eq:C.HM.9}
\]
The complete rational curve of square zero gives
the actual ruling on the compact normalization of
\(Q\).  For clarity, its local complete deformations
are disjoint, since their intersection number is zero;
the resulting compact curve family gives the ruling
after the point blowdowns of its singular fibres.
Equivalently this is the ruled-surface comparison in
\cite[\S4.5]{Kulikov1977}.  Its one boundary intersection
is the section belonging to \(P\).  Any other
horizontal boundary would meet a general complete
fibre in this same local family, which it does not.
All other boundary curves on \(Q\) are vertical.
This is the precise ruled section-and-fibre
configuration, including any nonminimal fibres.


\end{enumerate}

\item \par\medskip\noindent\textbf{Ruled removal and cleanup of its singular charts.}
We use the actual finite maps in
\cite[\S\S4.10.7--4.10.8,4.14]{Kulikov1977}, with
the missing-fibre correction in
\cite[\S8]{Kulikov1981}.  Their overlaps and proper
composite graphs are those verified in \ref{cdp:st:C.2.UN.75}.
They remove the ruled component in \eqref{cdp:eq:C.HM.9}.
If its boundary is disconnected, they first insert
the specified missing fibre components; these
exceptional ruled surfaces have order \(q+1\).
When \(q=M\), the canonical divisor on its
normalization has only negative boundary
coefficients, and \cite[Lemma2.19]{Kulikov1977}
makes that boundary connected.  No preliminary
addition is then made.  Thus all added orders are
at most the original fixed \(M\).


\begin{enumerate}[label={$\langle2\rangle$\arabic*.},leftmargin=1.3em,labelsep=.5em]
\item  The result of this removal need not immediately
have smooth total space.  We retain precisely the
finite sigma-atlas of
\cite[\S4.11.2]{Kulikov1977}: its singular sets are
contained in finitely many disjoint charts; in each
chart the recorded local exceptional branches
contract successively to a smooth NC chart.  We
now describe an actual cleanup, before applying
another last-point argument to a Hopf.


\item  If one of these lists is nonempty, choose a
component of largest canonical order among the
globally occurring last locally blown-up branches.
Use an auxiliary resolution \(a:\widehat X\to X'\)
as in \cite[\S4.13]{Kulikov1977} to read its
old strict component \(\widehat V\).  This
resolution is a diagnostic model for the
hypotheses of Lemma4.14; it is not inserted
as an additional degeneration stage in the
modification record.  In particular its new
exceptional divisors receive no card labels,
and no bound on their canonical orders is used.
Each branch has a complete local exceptional
\(\PP^1\)-fibre.  These complete fibres extend to
a ruling of the old smooth surface
\(\widehat V\).  In the
generic local exceptional chart there is exactly
one horizontal boundary section.  Completeness
of those fibres shows that there can be no
additional global horizontal boundary; all the
other boundary, including the diagnostic
exceptional curves over the finite exceptional
points, is vertical.  This supplies
conditions (a), (b) and (d) of
\cite[Lemma4.14]{Kulikov1977}, by their actual
normalization and section maps.


\item  If another curve on this chosen component had
been locally blown up, its exceptional branch
would have order one larger.  If that child
is not last, follow its finite local descendant
list.  The order increases at each step, so
a last descendant has strictly larger order
than the chosen component.  This contradicts
the choice of largest last-branch order.
Thus the choice gives condition
(e).  This is the exact comparison in
\cite[\S4.16.2]{Kulikov1977}.  The reordering in
its \S4.12 makes its specified branches last
without changing their restored overlaps.
The correction in \cite[\S8]{Kulikov1981}
inserts any missing vertical links in the
specified smooth open charts, leaving their
overlap fixed.  Lemma4.14 then removes this
global ruled component, and the result again
has the same kind of sigma-atlas.


\item  Every cleanup removal has positive order:
a locally exceptional component has the
actual order of its parent plus one.  It
adds only finitely many components at the
next order, and none at \(M\).  Here we
use precisely the maximal-old-order clause
of \cite[\S8]{Kulikov1981}, including its
last-locally-blown-up case: its boundary
is connected and no missing links are
inserted.  This does not assert that
every divisor on the auxiliary resolution
has order at most \(M\).  A Hopf
component cannot be one of these locally
exceptional global components: its complete
local rational fibres would give rational
curves through an open set of its compact
normalization, contrary to the finite
exceptional-curve description above.
The strict Hopf divisors survive by the
proper graph argument in \ref{cdp:st:C.2.UN.75}.


\item  This cleanup terminates without inserting an
ambient resolution into the record or estimating
orders of
arbitrary resolution divisors.  Label every
current global component by its order in
\(\{0,\ldots,M\}\).  A removal deletes a
positive label \(q\), adds finitely many
labels \(q+1\) when \(q<M\), and adds none
when \(q=M\).  There are finitely many
removals at level one, hence finitely many
objects ever created at level two, and
inductively at every level through \(M\).
This is the finite game of
\cite[\S9]{Kulikov1981}, now applied also
to the cleanup removals.


\item  When cleanup stops, no recorded locally
exceptional branch remains.  The defining
contraction lists of the sigma-atlas are
empty; their terminal smooth NC charts
are therefore the current charts.
The total space is smooth and its full
fibre is reduced NC.  Its normalization
components are smooth.  Other global
components may retain self-incidences.
At a maximal Hopf the intrinsic argument
\eqref{cdp:eq:C.HM.2}--\eqref{cdp:eq:C.HM.5} excludes rational
self-incidences.  No elliptic self-incidence
was introduced: type-I and type-II changes
use rational exceptional curves; in a
local ruled contraction the horizontal
parent has order one lower, so a
maximal-order Hopf can only meet that
contraction in vertical rational curves.
Thus the maximal Hopfs are again smooth
embedded components, and their ordinary
point maps to their Hopf minimal models
are available for the next step.


\end{enumerate}

\item \par\medskip\noindent\textbf{Termination of the whole procedure.}
Perform a type-I or type-II step whenever
the chosen last point gives the first or
second case.  In the third case perform
the finite ruled removal and cleanup
just specified.  Across the entire
procedure there are only finitely many
divisorial removal blocks: the same
fixed-level game counts all their
positive component labels together.
Between these blocks every step strictly
decreases the finite integer \eqref{cdp:eq:C.HM.8}.
More explicitly, if \(c_j\) counts the
retained global components of order \(j\),
the block-level potential is
\[
                 (c_1,c_2,\ldots,c_M,h)
                 \quad\hbox{in lexicographic order.}
\tag{C.HM.10}\hypertarget{tag-00MH}{}\label{cdp:eq:C.HM.10}
\]
A whole-component removal decreases its
first changed \(c_j\), whatever finite
next-level additions or Hopf point
changes occur inside the block.
A type-I or type-II step chosen on a
last Hopf point leaves all the \(c_j\)
unchanged and decreases \(h\).
Lexicographic order on this fixed
finite product of nonnegative integers
has no infinite strictly descending
sequence, by the same successive-coordinate
induction.  This potential is attached
to the specified finite blocks, not to
every auxiliary arrow in their graphs.
An infinite procedure would therefore
have either infinitely many divisorial
removals, contradicting that game, or
an infinite tail of strict decreases of
\(h\).  Both are impossible.


\begin{enumerate}[label={$\langle2\rangle$\arabic*.},leftmargin=1.3em,labelsep=.5em]
\item  At the endpoint a nonminimal maximal
Hopf would still have a last point
blowup and one of the three alternatives,
so every such Hopf is minimal.  Its
curves are the disjoint smooth elliptic
curves proved above, and hence no
triple point meets it.  The original
Hopf primes survive with unchanged
orders by \ref{cdp:st:C.2.UN.75}, and each additional
prime used in the record is ruled.
The same actual proper graphs and
restored overlap isomorphisms give
the asserted reversible comparison.
Their annular gluing and normal-cone
Euler terms are exactly those of \ref{cdp:st:C.2.UN.75}.
This proves the statement in its
specified scope.



\end{enumerate}

\end{enumerate}
\end{proof}






\begin{proposition}[The \((-1,-2)\) curve with canonical orders \((0,0,1)\) occurs in a compact zero-Euler torus fibre]\hypertarget{tag-01EQ}{}\label{cdp:st:C.2.UN.79}
\StatementRecord{cdp:st:C.2.UN.79}{proved}{01EQ}{\SeriesRef{cdp:st:C.2.UN.70}, \SeriesRef{cdp:st:C.2.UN.41}, \SeriesRef{cdp:eq:C.BC.4}, \SeriesRef{cdp:st:C.2.UN.74}, \SeriesRef{cdp:eq:C.BC.5}, \SeriesRef{cdp:eq:C.BC.8}, \SeriesRef{cdp:eq:C.BC.9}, \SeriesRef{cdp:eq:C.BC.13}}

There is a smooth compact projective threefold \(X\), a proper map
\(f:X\to\PP^1\) with connected general fibre a complex two-torus,
and a full reduced SNC fibre \(S=f^{-1}(0)\) with \(e(S)=0\),
containing a smooth rational curve \(F\) with the following exact
properties. The curve belongs to just one component \(P\) generically,
is not a double curve, has \(F_P^2=-1\), and meets two other
components transversely, once each, at distinct points. Its entire
analytic neighbourhood, including \(f\), is the germ along the zero
section of
\[
 \operatorname{Tot}\bigl(\OO_{\PP^1}(-1)_a
                 \oplus\OO_{\PP^1}(-2)_b\bigr),\qquad
 f=bUV,\quad [U:V]\in\PP^1.                         \tag{C.BC.1}\hypertarget{tag-00AZ}{}\label{cdp:eq:C.BC.1}
\]
Here \(P=(b=0)\), and the other two components are the fibres
\(U=0\) and \(V=0\) of the displayed vector bundle.
The actual generator of \(f_*K_X\) has orders \(0,1,0\)
on these three components in that order. Thus, after ordering the
two neighbours by canonical height, the tuple is \((0,0,1)\).
In particular
\[
                    N_{F/X}=\OO(-1)\oplus\OO(-2).            \tag{C.BC.2}\hypertarget{tag-00B4}{}\label{cdp:eq:C.BC.2}
\]

The construction retains an actual point-blowdown map from \(P\)
to \(\PP^1\times E\), where \(E\) is an elliptic curve.
It maps \(F\) isomorphically onto \(\PP^1\times\{e_0\}\).
For the elliptic map obtained by composing this blowdown with
\(\PP^1\times E\to\PP^1\), the curve \(F\) is therefore
horizontal. The specified exceptional locus of this surface map
does not contain \(F\). This records the scope relevant to \ref{cdp:st:C.2.UN.70}:
the example concerns an arbitrary nonboundary \((-1)\)-curve,
with its exact comparison to the chosen elliptic minimalization.
\end{proposition}



\begin{proof}

\begin{enumerate}[label={$\langle1\rangle$\arabic*.},leftmargin=*]
\item \par\medskip\noindent\textbf{An actual compact elliptic surface with a toric \(I_3\) chart.}
Take a compact rational elliptic surface with a simple nodal fibre.
Such a surface, with an explicit cubic-pencil construction, was
used in \ref{cdp:st:C.2.UN.41}. Alternatively the Weierstrass equation
\[
       y^2=x^3+(t^4+1)x+t^6                         \tag{C.BC.3}\hypertarget{tag-00B5}{}\label{cdp:eq:C.BC.3}
\]
gives one: its discriminant, up to a nonzero constant, is
\(31t^{12}+12t^8+12t^4+4\), which is square-free and has no
zero at infinity. Choose one nodal value, put it at \(t=0\)
by a base coordinate, and take the global cubic base change
\(t=s^3\). Normalize and minimally resolve. The resulting
compact smooth projective elliptic surface \(J\to\PP^1_s\)
has an \(I_3\) fibre
\[
                         C_0+C_1+C_2 .             \tag{C.BC.4}\hypertarget{tag-00B6}{}\label{cdp:eq:C.BC.4}
\]
The components are smooth rational curves, meeting in a triangle
at three distinct nodes, with no triple point. The base change
is unramified at the other singular values; at a smooth value
it is the ordinary pullback of a smooth family. Thus these
operations give a smooth total surface after resolving the
single local \(A_2\) equation at the chosen nodal value.

\item Here is the required local comparison, with its entire two-chart
neighbourhood retained. At the original nodal point the map is
\(xy=t\), by the holomorphic Morse coordinates. The root is
\(xy=s^3\). Two consecutive charts of its actual minimal
resolution are
\[
 \begin{aligned}
 (x,y,s)&=(b,b^2z^3,bz),\qquad (b,z)\in\C^2,\\
 (x,y,s)&=(u^2c,uc^2,uc),\qquad (u,c)\in\C^2,\\
                 u&=z^{-1},\qquad c=z^2b .
 \end{aligned}                                             \tag{C.BC.5}\hypertarget{tag-00B7}{}\label{cdp:eq:C.BC.5}
\]
The curve \(C_1\) is \(b=0\) in the first chart and
\(c=0\) in the second. Its adjacent curves are \(z=0\)
and \(u=0\). The two-chart union is the actual total space
of \(\OO(-2)\), restricted to a neighbourhood if necessary;
its map is exactly \(s=bz=cu\).
The displayed monomials satisfy \(xy=s^3\), glue on their
overlap, and are the successive point-resolution charts,
so this comparison does not identify a neighbourhood merely
from its normal degree.

\item The two-form \(dx\wedge ds/x\) is \(db\wedge dz\)
on the first chart and \(du\wedge dc\) on the second.
An original canonical generator is a nonvanishing
holomorphic multiple of \(dx\wedge dy\) at the node.
Pull it back and divide by \(s^2\); the resulting
form is a nonvanishing holomorphic multiple of the
displayed two-form and extends through the unchanged
smooth charts of the original fibre. The three
components have the same original root order two.
This gives a nowhere-zero canonical generator
near \eqref{cdp:eq:C.BC.4}. Hence the actual
canonical evaluation orders on all three components of
\(J\) near this fibre are zero.

\par\medskip\noindent\textbf{One global curve blowup produces the required germ and orders.}
Fix an elliptic curve \(E\) and a point \(e_0\in E\).
Set \(X_0=J\times E\), with map \(f_0\) induced by \(s\).
Its general fibre is the product of two elliptic curves.
The central components are \(P_i=C_i\times E\), all
of canonical order zero. Blow up the complete smooth
rational curve
\[
                      B_0=C_0\times\{e_0\}\subset P_0 .
                                                               \tag{C.BC.6}\hypertarget{tag-00B8}{}\label{cdp:eq:C.BC.6}
\]
The centre is contained generically in just \(P_0\) and
meets each of its two double curves transversely at one
point. The inverse chart of \ref{cdp:st:C.2.UN.74}, or the following direct
calculation, shows that the new fibre is full reduced SNC.
Near any such point write the original map as \(xy\),
the centre as \((x,a)\), and the blowup charts as
\(a=xA\) and \(x=aX\). Their maps are respectively
\(xy\) and \(aXy\). All displayed factors have
multiplicity one. The new exceptional component \(R_0\)
has canonical order one, because the Jacobian of this
smooth curve blowup has a simple exceptional zero.

\item Choose the labelling in \eqref{cdp:eq:C.BC.5} so that \(z=0\) is
the branch of \(C_0\). Let \(a\) be a local coordinate
of \(E\) at \(e_0\). Near the whole curve
\(C_1\times\{e_0\}\), the centre is \((z,a)\)
near its first endpoint and is absent elsewhere.
In the \(z\)-chart of its blowup put
\[
               a=zA,\qquad s=zb .                          \tag{C.BC.7}\hypertarget{tag-00B9}{}\label{cdp:eq:C.BC.7}
\]
The strict transform \(F\) is \(A=b=0\). In the other
chart of \eqref{cdp:eq:C.BC.5} no modification occurs near this curve;
write its elliptic coordinate as \(a_\infty\). The actual
transition is
\[
       u=z^{-1},\qquad A=u a_\infty,\qquad b=u^2c,
             \qquad s=zb=uc .                              \tag{C.BC.8}\hypertarget{tag-00BA}{}\label{cdp:eq:C.BC.8}
\]
These are exactly the two vector-bundle charts in \eqref{cdp:eq:C.BC.1}.
They prove \eqref{cdp:eq:C.BC.2} and \(F_P^2=-1\), with the entire
map retained. The two boundary intersections occur at
\(z=0\) and \(u=0\), whereas \(F\) is contained
generically only in \(b=0\).

\item Up to its retained nonvanishing holomorphic factor, the actual
product canonical form pulls back as
\[
 db\wedge dz\wedge da
       =z\,db\wedge dz\wedge dA .                          \tag{C.BC.9}\hypertarget{tag-00BB}{}\label{cdp:eq:C.BC.9}
\]
At the other endpoint this is a nonzero unit times
\(du\wedge dc\wedge da_\infty\). Thus it has a simple
zero exactly on \(R_0=(z=0)\), with zero orders on
\(P\) and its other old neighbour. This is the requested
canonical tuple for the actual direct-image generator.
It remains a generator: the curve blowup has the actual
canonical pushforward \(\pi_*K=K\) through its Jacobian,
or directly through Hartogs across its smooth codimension-two
centre.

\item The surface \(P\) is \(\operatorname{Bl}_p(\PP^1\times E)\),
where \(p\) is the endpoint of
\(\PP^1\times\{e_0\}\). The strict transform of that
rational curve is \(F\). The old double curve through
\(p\) has the other tangent direction in \(P_i\), so
its strict transform meets the exceptional line at a
different point from \(F\). In particular the full
strict transform of \(P_0\) is disjoint from \(F\).
The actual surface blowdown maps \(F\) isomorphically
to \(\PP^1\times\{e_0\}\); its composition with
the first projection has degree one on \(F\).

\par\medskip\noindent\textbf{Exact zero-Euler balancing away from the entire curve.}
We make the remaining centres explicit. Use
\[
 E:\ y^2=x(x-1)(x-2),\qquad e_0=(0,0),\quad e_1=(1,0).
                                                               \tag{C.BC.10}\hypertarget{tag-00B0}{}\label{cdp:eq:C.BC.10}
\]
The strict \(P_0\) is still the identical surface
\(\PP^1\times E\): restriction of the first blowup
to this component is the blowup of its Cartier divisor
\(B_0\), hence an isomorphism. Blow up the disjoint
curve \(B_1=C_0\times\{e_1\}\). Again \(P_0\)
is unchanged and all full multiplicities are one.
The second centre is disjoint from \(F\), so this
operation preserves a neighbourhood of the whole
compact curve.

\item Choose the coordinate \(z\) on \(C_0\simeq\PP^1\)
so its two old boundary points are \(0,\infty\). Put
\[
 P(z)=z^2+1,\qquad Q(z)=3z^2+z+5,\qquad
 D=\{P(z)x=Q(z)\}\subset\PP^1\times E .
                                                               \tag{C.BC.11}\hypertarget{tag-00B1}{}\label{cdp:eq:C.BC.11}
\]
This is the divisor of a section of
\(\OO_{\PP^1}(2)\boxtimes\OO_E(2\cdot\infty_E)\).
Its map to \(\PP^1_z\) is the exact smooth double
cover
\[
          Y^2=P Q(Q-P)(Q-2P),\qquad Y=yP^2 .                  \tag{C.BC.12}\hypertarget{tag-00B2}{}\label{cdp:eq:C.BC.12}
\]
The four quadratic factors have respectively discriminants
\(-4,-59,-31,-11\), are pairwise coprime, and have
nonzero leading coefficients. Their eight simple roots
give a smooth connected genus-three curve. The chart
at \(P=0\) is the ordinary elliptic infinity chart;
since these zeros are simple it gives the same smooth
double cover, not additional components. At \(z=0\)
and \(\infty\) the values \(x=5,3\) are outside
\(\{0,1,2,\infty\}\). Thus \(D\) meets each old
elliptic boundary transversely in two points and avoids
all four existing triple points, which have elliptic
coordinate \(e_0\) or \(e_1\). The simple zeros of
\(Q\) and \(Q-P\) prove transverse intersection
with \(B_0\) and \(B_1\), respectively. These facts
also verify all its boundary intersections directly.

\item Blow up this \(D\) in the threefold. The same two
charts \(xy\), \(aXy\) prove reduced SNC at its
boundary crossings; elsewhere it is a smooth-centre
blowup in one reduced component. It has canonical
order one generically, since its ambient centre lies
in the zero-order component \(P_0\). The centre is
disjoint from \(F\), so \eqref{cdp:eq:C.BC.8}--\eqref{cdp:eq:C.BC.9} remain
unchanged on an open neighbourhood of the entire
curve, including every analytic transverse jet.

\item Initially \(e((C_0+C_1+C_2)\times E)=0\).
For a smooth curve centre contained in the fibre,
the actual blowup replaces that curve by its
\(\PP^1\)-bundle exceptional divisor and is an
isomorphism on the complement. Its reduced-fibre
Euler defect is therefore exactly \(e\) of the
centre. Consequently our three operations give
\[
             e(S)=0+e(B_0)+e(B_1)+e(D)
                         =0+2+2-4=0 .                       \tag{C.BC.13}\hypertarget{tag-00B3}{}\label{cdp:eq:C.BC.13}
\]
For an independent stratum check, the six final components
have Euler numbers \(0,4,4,6,6,-8\). The double curves
have total Euler characteristic \(24\), and there are
\(12\) triple points. Thus their actual SNC
inclusion--exclusion gives \(12-24+12=0\), in
agreement with \eqref{cdp:eq:C.BC.13}.

\item All blowups are along closed smooth curves in a
compact projective threefold. Their total spaces
are smooth, compact and projective. They preserve
the map outside the single fibre; properness and
connectedness follow from their connected blowup
fibres and the original connected elliptic product.
The other fibres remain unchanged products with
\(E\). Finally, every additional point blowup
induced on \(P\) is away from \(F\), so the
same composite surface blowdown and its degree-one
map on \(F\) prove the stated horizontal
qualification on the final zero-Euler model.

\end{enumerate}
\end{proof}






\begin{proposition}[Every point-exceptional prime of an aligned parent is an actual boundary]\hypertarget{tag-01ES}{}\label{cdp:st:C.2.UN.80}
\StatementRecord{cdp:st:C.2.UN.80}{proved}{01ES}{\SeriesRef{cdp:st:C.2.UN.70}, \SeriesRef{cdp:st:C.2.UN.16}, \SeriesRef{cdp:eq:C.EP.6}, \SeriesRef{cdp:eq:C.EP.7}, \SeriesRef{cdp:eq:C.EP.3}, \SeriesRef{cdp:eq:C.EP.10}, \SeriesRef{cdp:eq:C.EP.9}, \SeriesRef{cdp:eq:C.EP.11}}

Let \(P\), \(\pi:P\to Z\), and \(q_Z:Z\to B\) be the
actual component and maps of \ref{cdp:st:C.2.UN.70}. Write \(D_P\)
for the reduced union of its intersections with the
other central components. There is a nonzero constant
linear combination of the two restricted saturated
ambient translation fields which is the meromorphic
lift of a nowhere-zero holomorphic vertical field
\[
 W_Z\in H^0(Z,T_Z),\qquad dq_Z(W_Z)=0.
 \tag{C.EP.1}\hypertarget{tag-00I1}{}\label{cdp:eq:C.EP.1}
\]
On every smooth elliptic fibre its complex-time period
lattice is the same fixed lattice. The field extends
without a zero or a pole across every smooth multiple
fibre of \(q_Z\).

Every prime exceptional for the specified point map
\(\pi\) is a polar divisor of its meromorphic lift
\(W_P\). Consequently
\[
                 \operatorname{Exc}(\pi)\subset D_P.
 \tag{C.EP.2}\hypertarget{tag-00I4}{}\label{cdp:eq:C.EP.2}
\]
In particular, a last point-exceptional curve for
\(\pi\) is an actual boundary curve. This assertion
concerns the exceptional primes of this actual map:
it makes no assertion that every rational
\((-1)\)-curve on \(P\) is exceptional for \(\pi\).
The component and its ambient neighbourhood are not
changed.
\end{proposition}



\begin{proof}

\begin{enumerate}[label={$\langle1\rangle$\arabic*.},leftmargin=*]
\item \par\medskip\noindent\textbf{The actual field space and its polar locus.}
Use the two-dimensional saturated field space of
\ref{cdp:st:C.2.UN.16} on the original multiplicity-one component
\(P\). In this full reduced situation no root cover
of \(P\) is needed. Denote a basis of its restrictions
by \(X,Y\). They are independent over \(\C\), commute,
and are meromorphic fields on \(P\).

\item Their actual ambient representatives are linear
combinations of holomorphic relative fields with
coefficients meromorphic in the disc coordinate.
They therefore have polar divisors only in the
central fibre. Saturation removes their pole along
the generic point of \(P\). On \(P\setminus D_P\)
the central fibre has just this one smooth branch,
so restriction gives
\[
       X,Y\ \hbox{holomorphic on }P\setminus D_P .
 \tag{C.EP.3}\hypertarget{tag-00I5}{}\label{cdp:eq:C.EP.3}
\]
This is the exact restriction comparison used in
(C.MR.8) and \ref{cdp:st:C.2.UN.70}. It applies to every constant
linear combination. All boundary and exceptional
curves of the specified map \(\pi\) are vertical
for \(q=q_Z\pi\).

\item Remove finitely many values from \(B\) so that
\(\pi\) is an isomorphism over their complement
\(B^\circ\), all fibres are smooth, and the fields
and their projected coefficients are regular.
On a compact elliptic fibre the latter coefficients
are holomorphic functions and hence constants.
Thus \(X,Y\) project to meromorphic vector fields
on \(B\), exactly as in the alignment proof.

\par\medskip\noindent\textbf{Rank one: the period identity gives a constant combination.}
Suppose the meromorphic rank of \(\langle X,Y\rangle\)
is one. Alignment identifies its connected elliptic
first integral with the actual map \(q\). Thus
\[
                       Y=(q^*H)X
 \tag{C.EP.4}\hypertarget{tag-00I6}{}\label{cdp:eq:C.EP.4}
\]
for a nonconstant meromorphic function \(H\) on
\(B\). The actual return-germ calculation
(C.RG.4)--(C.RG.12), performed on the original
ambient representatives, gives time periods
\(\lambda_1,\lambda_2\) of \(X\) and constants
\(d_i,k_i\in\C\) such that
\[
            \lambda_i=d_i-k_iH,\qquad
            \lambda_2/\lambda_1=\tau\in\mathbb H .
 \tag{C.EP.5}\hypertarget{tag-00I7}{}\label{cdp:eq:C.EP.5}
\]
The formula continues on \(B^\circ\) with these
same constants. This uses the actual global ratio
\(H\), the continued integral cycles, and the
identity theorem; it does not extend \(q\) to a
threefold neighbourhood.

\item Take
\[
                       W=d_1X-k_1Y=\lambda_1X.
 \tag{C.EP.6}\hypertarget{tag-00I8}{}\label{cdp:eq:C.EP.6}
\]
A nonzero integral cycle on an elliptic curve has
nonzero time period, so \(W\) is nonzero on every
fibre under consideration. Multiplying a field by
\(\lambda_1\) divides its time periods by
\(\lambda_1\); the two periods of \(W\) are exactly
\(1,\tau\). Thus \eqref{cdp:eq:C.EP.6} is a constant linear
combination with a fixed nonzero period lattice.
In an integration coordinate on each fibre it
is the fixed invariant translation field.

\par\medskip\noindent\textbf{Rank two: commuting transverse transport preserves the periods.}
Suppose instead that the meromorphic rank is two.
The two projected fields on the curve \(B\)
commute. Two nonzero commuting meromorphic fields
on a connected curve have constant ratio:
writing them locally as \(a(z)\partial_z\) and
\(b(z)\partial_z\), the bracket equation is
\(ab'-ba'=0\). The ratio therefore continues as
a constant. They cannot both project to zero,
since vertical fields on an elliptic fibre have
meromorphic rank at most one. After interchanging
the basis if necessary there is consequently a
constant \(c\) for which
\[
             W=Y-cX\ne0,\qquad q_*W=0,\qquad
                         q_*X\ne0 .
 \tag{C.EP.7}\hypertarget{tag-00I9}{}\label{cdp:eq:C.EP.7}
\]
On the compact fibres over \(B^\circ\), \(W\)
is a nonzero translation field. The local flow
of the transverse field \(X\), constructed on
an entire compact fibre for sufficiently small
time, identifies neighbouring elliptic fibres.
Since \([X,W]=0\), this identification preserves
the time form and its two integral periods.
Thus their lattice is locally constant and
hence constant on connected \(B^\circ\).
This is the actual compact-fibre flow comparison
of (C.RF.1)--(C.RF.2). A constant rescaling of
\(W\) makes one nonzero period equal to \(1\).
Once again \(W\) is a constant combination of
the original two restricted fields.

\par\medskip\noindent\textbf{Extension on the specified minimal elliptic model.}
On \(q_Z^{-1}(B^\circ)\), push \(W\) through
the open isomorphism of \(\pi\). Integration of
its fixed lattice identifies each small smooth
fibre tube with a torsor under one fixed elliptic
curve \(F\). Every overlap has translation
linear part, so its derivative preserves the
same invariant field of \(F\).

\item At a multiple value, use the actual normalized
root chart of \ref{cdp:st:C.2.UN.70} and (C.RG.13):
\[
 (F\times\Delta_z)/\langle
       (u,z)\longmapsto(u+\eta,\zeta_mz)\rangle,
       \qquad t=z^m .
 \tag{C.EP.8}\hypertarget{tag-00IA}{}\label{cdp:eq:C.EP.8}
\]
The translation \(\eta\) gives the free cyclic
action on the total space. The invariant
vertical field on \(F\times\Delta_z\) is fixed
by this action and is nowhere zero. On the
punctured root disc our field agrees with it:
its coefficient relative to the invariant
frame has the fixed nonzero time periods and
therefore is the same constant. It follows
that the agreement extends across \(z=0\)
and descends through the actual etale quotient.
These extensions agree on overlaps with the
previously fixed field. This proves \eqref{cdp:eq:C.EP.1},
including every zero or pole that might have
been hidden at an omitted base value.

\item The meromorphic lift of \(W_Z\) through \(\pi\)
agrees with the constant combination \eqref{cdp:eq:C.EP.6}
or \eqref{cdp:eq:C.EP.7} on their common open isomorphism,
and therefore agrees everywhere as a
meromorphic vector field. In particular its
polar divisors are contained in \(D_P\), by
\eqref{cdp:eq:C.EP.3}.

\par\medskip\noindent\textbf{Every exceptional valuation is polar.}
Let \(E\) be any prime exceptional for \(\pi\),
with image \(p\in Z\). The nowhere-zero field
\(W_Z\) has a holomorphic flow box \((z,w)\)
centred at \(p\) in which
\[
                   W_Z(z)=1,\qquad W_Z(w)=0.
 \tag{C.EP.9}\hypertarget{tag-00IB}{}\label{cdp:eq:C.EP.9}
\]
Both pullbacks vanish to positive order along
\(E\). Work at a generic smooth point of \(E\)
and let \(s\) be its local equation. Write
\[
            \pi^*w=s^b u,\qquad
            \pi^*z=s^a v,\qquad a,b\ge1,
 \tag{C.EP.10}\hypertarget{tag-00I2}{}\label{cdp:eq:C.EP.10}
\]
where \(u,v\) are holomorphic units after
shrinking this generic neighbourhood.

\item Suppose the lifted derivation \(D=W_P\) were
holomorphic there. The identity
\(D(\pi^*w)=0\), divided by \(s^{b-1}\) and
reduced modulo \(s\), gives
\[
                    b\,u\,D(s)=0\pmod{s}.
 \tag{C.EP.11}\hypertarget{tag-00I3}{}\label{cdp:eq:C.EP.11}
\]
Characteristic zero and the fact that \(u\)
is a unit imply \(D(s)\in(s)\). Thus \(D\)
preserves the ideal of \(E\), and applying it
to the second expression in \eqref{cdp:eq:C.EP.10} gives
\(D(\pi^*z)\in(s)\). This contradicts the
exact pulled-back identity \(D(\pi^*z)=1\).
Hence \(E\) is a polar divisor. The argument
applies directly to every exceptional prime,
including all infinitely-near point divisors;
it does not assume that a first pole survives
an unexamined modification.

\item Combining this with \eqref{cdp:eq:C.EP.3} proves \eqref{cdp:eq:C.EP.2}.
The exact period rescalings, the root-chart
invariance, and finite toric valuation
certificates for \eqref{cdp:eq:C.EP.9}--\eqref{cdp:eq:C.EP.11} are
executed in \path{code/A_translation_exceptions.py}.
The general exceptional-prime assertion is
the local derivation proof above, without a
toric hypothesis on the actual point map.

\end{enumerate}
\end{proof}






\begin{proposition}[Finite removal of all vertical point defects at the marked Hopf interfaces]\hypertarget{tag-01ET}{}\label{cdp:st:C.2.UN.81}
\StatementRecord{cdp:st:C.2.UN.81}{proved}{01ET}{\SeriesRef{cdp:st:C.2.UN.77}, \SeriesRef{cdp:st:C.2.UN.75}, \SeriesRef{cdp:st:C.2.UN.70}, \SeriesRef{cdp:st:C.2.UN.80}, \SeriesRef{cdp:eq:C.AP.5}, \SeriesRef{cdp:eq:C.AP.3}, \SeriesRef{cdp:eq:C.AP.6}, \SeriesRef{cdp:st:C.2.UN.74}, \SeriesRef{cdp:st:C.2.UN.31}, \SeriesRef{cdp:eq:C.AP.9}, \SeriesRef{cdp:eq:C.AP.10}, \SeriesRef{cdp:eq:C.AP.11}, \SeriesRef{cdp:st:C.2.UN.64}, \SeriesRef{cdp:eq:C.AP.2}, \SeriesRef{cdp:eq:C.AP.4}, \SeriesRef{cdp:eq:C.AP.14}; \cite[\S\S4.10.7--4.10.8,4.14]{Kulikov1977}; \cite[\S8]{Kulikov1981}}

Let \(f:X\to\Delta\) be a proper degeneration of connected
complex two-tori with smooth total space and full reduced
SNC central fibre. We also allow as the starting recorded model
the smooth-total NC endpoint obtained from such a fibre by
\ref{cdp:st:C.2.UN.77}, retaining its two labelled normalization branches at
each self-double curve. Normalize the actual canonical generator by
\[
 \operatorname{div}(\Omega)=\sum k_iV_i,\qquad
 \min k_i=0,\qquad M=\max k_i.
 \tag{C.AP.1}\hypertarget{tag-009M}{}\label{cdp:eq:C.AP.1}
\]
Suppose that the components of order \(M\) with Hopf minimal
model have already been made minimal by the actual record of
\ref{cdp:st:C.2.UN.77}. Mark any finite collection of their interfaces
\(J_\alpha=H_\alpha\cap P_\alpha\) for which the normal line of
\(J_\alpha\) is torsion and \(k_{P_\alpha}<M\).
Retain repeated markings on the same parent.

There is a finite overlap-compatible modification record,
unchanged over \(\Delta^*\), whose endpoint has smooth total
space and full reduced NC central fibre, such that every
designated parent is an embedded minimal fixed-Jacobian
elliptic quasi-bundle. Its full elliptic fibres are multiples
of smooth elliptic curves, its integral linear monodromy is
the identity, and every one of its boundary curves is a
smooth elliptic fibre support. In particular no triple point
meets a designated parent at the endpoint.

All maximal-order minimal Hopf components and an open
neighbourhood of their union are unchanged by the record.
In particular every marked \(J_\alpha\), its normal map and
its gluing to \(H_\alpha\) are retained. The designated parent
prime divisors survive. Other endpoint components can retain
NC self-incidences; their two normalization branches are
part of the record.

Every retained global component has canonical order between
\(0\) and the original \(M\). Whole-component removals delete
a component of positive order \(q\); their preliminary
additions, when needed, have order \(q+1\le M\), and there
are no additions at \(q=M\). The record retains its actual
centres, local maps, normalization branches and restored
overlaps, so reversing it recovers the original family.
For a given compact family the same modifications are
performed on that compact base. Their Euler changes are
the exact normal-cone terms of \ref{cdp:st:C.2.UN.75}.

No assertion is made here that every maximal Hopf component
has one of the marked interfaces. In particular an
equal-order interface with a nontorsion normal line is not
among these markings.
\end{proposition}



\begin{proof}\leavevmode\par



 
\begin{enumerate}[label={$\langle1\rangle$\arabic*.},leftmargin=0pt,labelsep=.5em]
\item \par\medskip\noindent\textbf{The exact marked parent maps.}
All curves on a minimal Hopf surface are disjoint smooth
elliptic curves. Thus each marked curve is isolated from
the other boundaries on both sides. Its actual normal-product
isomorphism makes the normal on the parent torsion as well.
On the normalization of a designated parent \(P\), of order
\(p\), the marked curve has coefficient
\(M-p-1\ge0\) in the actual canonical divisor. \ref{cdp:st:C.2.UN.70} therefore
gives a specified point map and elliptic map
\[
 P\xrightarrow{\pi_P} Z_P\xrightarrow{q_P}B_P,
 \qquad R^1(q_P)_*\OO_{Z_P}=\OO_{B_P},
 \qquad g(B_P)\le1.
 \tag{C.AP.2}\hypertarget{tag-009S}{}\label{cdp:eq:C.AP.2}
\]
The point map is an isomorphism on a neighbourhood of every
marked curve. Every boundary is vertical for \(q_P\pi_P\),
and every rational boundary is exceptional for \(\pi_P\).
The exact constant-period field of \ref{cdp:st:C.2.UN.80}, equations
(C.EP.1)--(C.EP.11), proves the additional inclusion
\[
                 \operatorname{Exc}(\pi_P)\subset D_P.
 \tag{C.AP.3}\hypertarget{tag-009T}{}\label{cdp:eq:C.AP.3}
\]
Here \(D_P\) is the full actual reduced boundary. It includes
both preimages of a self-double curve at an NC stage.


\begin{enumerate}[label={$\langle2\rangle$\arabic*.},leftmargin=1.3em,labelsep=.5em]
\item  These assertions also apply at the intermediate smooth-total
NC stages used below. Indeed the actual adjunction divisor
on the normalization counts those two preimages; the marked
elliptic neighbourhood stays embedded and unchanged. The
saturated ambient fields restrict as meromorphic fields to
that normalization and are holomorphic off its full boundary.
The surface maps in \ref{cdp:st:C.2.UN.70} and the direct exceptional-valuation
argument in \ref{cdp:st:C.2.UN.80} consequently apply to this actual smooth
normalization, without identifying its two boundary branches.


\item  The target \(Z_P\) is a minimal surface in its birational
class. Thus
\[
 h_P=b_2(P)-b_2(Z_P)\ge0
 \tag{C.AP.4}\hypertarget{tag-009U}{}\label{cdp:eq:C.AP.4}
\]
is exactly the number of point blowups in \(\pi_P\). The
minimal second Betti number is unchanged under the subsequent
birational maps: in the ruled case it is two, and in the
other cases the minimal model is unique. We count each
designated parent once in \(h=\sum_P h_P\).


\end{enumerate}

\item \par\medskip\noindent\textbf{The full fibre graph and the receiving degrees.}
The point map \(\pi_P\) is a sequence of point blowups of
smooth elliptic fibre supports. For one such fibre, retain
its entire reduced inverse image, whether or not its strict
elliptic component is an actual boundary. Its graph is an
elliptic root with rational trees. Write \(C^2=-b_C\) and let
\(n_C\) be the number of its intersections with the other
components of this entire reduced fibre. The point maps give
\[
 n_C\le b_C+1\quad\hbox{for a rational component},\qquad
 n_C\le b_C\quad\hbox{for the elliptic component}.
 \tag{C.AP.5}\hypertarget{tag-009V}{}\label{cdp:eq:C.AP.5}
\]
The elliptic root starts with \((b,n)=(0,0)\). A new rational
curve has \((b,n)=(1,1)\) at a smooth point and \((1,2)\) at
a node. Blowing up a smooth point of an old component raises
both \(b\) and \(n\) by one; a node blowup raises \(b\) by one
and leaves its valence unchanged. These explicit updates
prove \eqref{cdp:eq:C.AP.5} at every infinitely-near point.


\begin{enumerate}[label={$\langle2\rangle$\arabic*.},leftmargin=1.3em,labelsep=.5em]
\item  Choose a last point blowup of \(\pi_P\), with exceptional
curve \(E\). By \eqref{cdp:eq:C.AP.3}, \(E\) is an actual boundary curve.
It meets zero, one or two other actual boundary curves.
These are exactly its triple points in the threefold.
The complete fibre graph has at most two neighbours at this
last point, so there is no additional case.


\item  Suppose first that there are two. On the other component
\(Q\), the normal-product formula gives \(E_Q^2=-1\), because
\(E_P^2=-1\). Hence the actual type-II two-ruling diagram of
\ref{cdp:st:C.2.UN.75} applies. It blows down \(E\) on both adjacent surfaces
and blows up the two opposite normalized branches at their
specified points.


\item  Neither receiving branch belongs to a designated parent.
To prove this, let \(C\) be the opposite boundary on \(P\),
and write \(-b_0\) and \(n_0\) for its square and entire-fibre
valence just before the last node blowup. Afterwards its
square is \(-b_0-1\), and its number \(t_C\) of actual triple
points is at most \(n_0\). On the opposite receiving surface
\(R\), the exact normal map therefore gives
\[
 C_R^2=b_0+1-t_C\ge b_0+1-n_0\ge0,
 \quad\hbox{and } C_R^2\ge1\hbox{ if }C\hbox{ is elliptic}.
 \tag{C.AP.6}\hypertarget{tag-009W}{}\label{cdp:eq:C.AP.6}
\]
On a designated parent every rational boundary is exceptional
and has negative square; every elliptic boundary is the
strict transform of a fibre support and has nonpositive
square. Both alternatives contradict \eqref{cdp:eq:C.AP.6}. This argument
uses the individual normalized branch when the two receivers
are branches of one global component. Thus a type-II move
leaves the component counts unchanged and strictly decreases
\(h\); it does not create a self-double curve on a designated
parent.


\end{enumerate}

\item \par\medskip\noindent\textbf{One triple point: an actual ruled-removal block.}
Suppose that \(E\) has one triple point. Then \(E_Q^2=0\).
Its complete rational deformations give the actual ruling
on the compact normalization of \(Q\), after its vertical
point contractions. The boundary passing through this one
triple point is its unique section; every other boundary is
vertical. This is the same complete-curve and actual ruling
comparison used in (C.HM.9), not a ruling inferred from a
numerical class alone. If \(q=k_Q\) and \(R\) is the section
component, restriction of the actual canonical divisor to
a complete ruling fibre gives
\[
                         k_R=q-1,
                         1\le q\le M.
 \tag{C.AP.7}\hypertarget{tag-009X}{}\label{cdp:eq:C.AP.7}
\]
The actual finite ruled-removal and cleanup block in \ref{cdp:st:C.2.UN.77}
therefore applies. On a projective-line bundle its contraction
is precisely \ref{cdp:st:C.2.UN.74}; for nonminimal fibres it is the retained
sigma-atlas construction of
\cite[\S\S4.10.7--4.10.8,4.14]{Kulikov1977}, with the
missing-fibre correction of \cite[\S8]{Kulikov1981}.


\begin{enumerate}[label={$\langle2\rangle$\arabic*.},leftmargin=1.3em,labelsep=.5em]
\item  For clarity, the numerical and analytic properties of this
block used here are the ones already proved in \ref{cdp:st:C.2.UN.75} and
\ref{cdp:st:C.2.UN.77}. It removes a positive-order component \(q\), adds
finitely many components of order \(q+1\) when necessary,
and adds none when \(q=M\). The diagnostic resolution used
to check Lemma4.14 is not inserted into the component list.
Its actual sigma-atlas is cleaned up by choosing a largest
last local branch, following its finite descendants and
performing the same ruled removals. Empty lists give smooth
NC charts. All overlaps are restored by the restrictions of
the same proper local graph maps. The positive component
card vector strictly decreases across this whole block.
No monotonicity of \(h\) inside the block is required.


\end{enumerate}

\item \par\medskip\noindent\textbf{Zero triple points: the exact plane comparison.}
Suppose that \(E\) has no triple point. Its square on \(Q\)
is \(+1\), and every other boundary on \(Q\) is disjoint
from \(E\). Here rationality does not require \(Q\) to be
a canonical maximum. A rational curve of positive square
cannot be exceptional for a minimal-model point map. On a
minimal model of nonnegative Kodaira dimension, the explicit
canonical lines in \ref{cdp:st:C.2.UN.31} and its quotient classification are
nef. Their pullback and the positive point discrepancies
would give \(K_Q\cdot E\ge0\), contrary to adjunction
\(K_Q\cdot E=-3\). A Hopf minimal model contains no rational
curve. On an irrational ruled minimal model every rational
curve is vertical for its ruling and has nonpositive square,
and point blowups cannot increase that square. The component
classification therefore leaves exactly the rational case.
One can specify its required plane
map without choosing a birational presentation. The sequence
\[
 0\longrightarrow\OO_Q\longrightarrow\OO_Q(E)
       \longrightarrow\OO_E(1)\longrightarrow0
 \tag{C.AP.8}\hypertarget{tag-009Y}{}\label{cdp:eq:C.AP.8}
\]
and \(H^1(Q,\OO_Q)=0\) give three sections. They have no base
point on \(E\), and the section defining \(E\) is nonzero
off \(E\). Thus \(|E|\) gives an actual morphism
\[
 \rho:Q\longrightarrow\PP^2,
 \qquad \rho^*\OO_{\PP^2}(1)=\OO_Q(E),
 \qquad \rho|_E:E\xrightarrow{\sim}\ell .
 \tag{C.AP.9}\hypertarget{tag-009Z}{}\label{cdp:eq:C.AP.9}
\]
Its image is a surface since \(E^2=1\); its degree is one
by the same intersection number. Hence it is a birational
morphism, factoring into the actual point blowdowns away
from \(E\). Every other boundary is contracted, since it
has zero intersection with \(E\).


\begin{enumerate}[label={$\langle2\rangle$\arabic*.},leftmargin=1.3em,labelsep=.5em]
\item  The canonical restriction to \(E\) gives
\(k_Q=k_P+2\). More generally call a current pair \((Q,E)\)
of the form \eqref{cdp:eq:C.AP.9}, with all other boundary disjoint from
\(E\), a marked plane pair, and write \(q=k_Q\). Its actual
canonical divisor has the precise form
\[
 K_Q=-3E+\sum_C a_CC,
 \qquad \sum_C a_CC=K_Q-\rho^*K_{\PP^2},
 \qquad a_C>0.
 \tag{C.AP.10}\hypertarget{tag-009N}{}\label{cdp:eq:C.AP.10}
\]
Indeed choose on \(\PP^2\) the meromorphic two-form whose
divisor is \(-3\ell\). The ratio of the actual form on \(Q\)
to its pullback has divisor supported on the exceptional
curves. On \(\PP^2\) this same meromorphic function has no
divisorial zero or pole, so it is a nonzero constant. Thus
the equality in \eqref{cdp:eq:C.AP.10} is an equality of actual divisors,
with the positive discrepancies of the specified point map.
It shows that every exceptional prime is an actual
boundary, and that there is no additional boundary of
coefficient zero. The only coefficient on the marked curve
is \(-3\). A self-double curve away from that marking would
have coefficient \(-1\), so such a pair stays embedded at
the smooth-total NC stages under consideration.


\item  Choose a last point blowup of \(\rho\). Its exceptional
boundary curve \(L\) lies away from \(E\). At a smooth point
of a discrepancy component of coefficient \(a>0\) its new
coefficient is \(a+1\); at a node of coefficients \(a,b>0\)
it is \(a+b+1\); at a point off the discrepancy divisor it
is one. These are the actual point-map identities for
\(K_Q+3E\). They give the following exhaustive alternatives:
\[
\begin{array}{c|c|c}
 \text{triple points of }L&\text{order of its neighbour}
                       &\text{actual next operation}\\ \hline
 0&q+2&\text{another marked plane pair}\\
 1&q+a+2\ge q+3&\text{a ruled-removal block}\\
 2&q+a+b+2&\text{a type-II move}.
\end{array}
 \tag{C.AP.11}\hypertarget{tag-009O}{}\label{cdp:eq:C.AP.11}
\]
In the last row the two receiving orders are
\(q+a+1\) and \(q+b+1\), both at least \(q+2\). The type-II
move decreases \(b_2(Q)\) by one; every other increase in
surface second Betti number occurs at a strictly higher
order. Its centre and all its points on \(Q\) are away from
the marked \(E\), so \eqref{cdp:eq:C.AP.9} persists on the modified \(Q\).
In the first row the same complete linear system constructs
the next marked plane pair, at a strictly greater order.


\item  This yields a finite progress block. Until its first whole
component removal, let
\[
                    b_j=\sum_{k_V=j}b_2(V^\nu),
                    \qquad 0\le j\le M,
 \tag{C.AP.12}\hypertarget{tag-009P}{}\label{cdp:eq:C.AP.12}
\]
where the sum uses the actual smooth normalizations. Each
type-II step in \eqref{cdp:eq:C.AP.11} strictly decreases
\((b_0,\ldots,b_M)\) lexicographically. Passing to another
marked plane pair makes no modification and strictly raises
its order. There can be at most \(M\) such successive calls
before an actual step. Since lexicographic order on this
fixed finite product of nonnegative integers is well founded,
these steps cannot continue indefinitely without a component
removal. At a marked plane pair with no exceptional curve,
the surface is \(\PP^2\), its sole boundary is \(E=\ell\),
and the principal full-fibre restriction gives
\[
                 N_{Q/X}=\OO_{\PP^2}(-1).
 \tag{C.AP.13}\hypertarget{tag-009Q}{}\label{cdp:eq:C.AP.13}
\]
The actual point contraction of \ref{cdp:st:C.2.UN.64} removes it to a smooth
point. The adjacent surface is its ordinary point blowup,
and no other component passes through that point. In the
one-triple alternative the positive ruled-removal block
already supplies a component removal.


\item  Stop this progress block at its first whole-component
removal, and perform the finite sigma cleanup just described.
Its net positive card vector decreases. It is unnecessary
to claim that the first marked plane pair or its marking
survives an arbitrary later cleanup: after this completed
block the outer argument uses the actual designated parent
maps again. This distinction retains every map while avoiding
an unsupported global preservation assertion.


\end{enumerate}

\item \par\medskip\noindent\textbf{Preserved maximal Hopf neighbourhoods.}
Every type-II centre used on a designated parent is vertical
for \eqref{cdp:eq:C.AP.2}, hence disjoint from its marked elliptic curves.
An opposite curve of a type-II move cannot lie on a minimal
maximal Hopf: a rational curve is impossible there, and the
elliptic alternative has the positive square \eqref{cdp:eq:C.AP.6}.
For the plane progress moves all opposite boundary curves
are rational. They likewise cannot meet a minimal Hopf.


\begin{enumerate}[label={$\langle2\rangle$\arabic*.},leftmargin=1.3em,labelsep=.5em]
\item  A ruled component removed in one of the actual blocks has
one boundary section, whose neighbouring order is \(q-1\),
and all other boundary is rational and vertical. It cannot
meet a minimal Hopf of order \(M\): that Hopf contains no
rational curve, while its being the section parent would
require \(q=M+1\). The same statement applies to every
cleanup removal and missing-fibre insertion. Plane point
contractions also have rational boundary. Thus these actual
closed centres and local sigma charts avoid the entire
union of maximal minimal Hopfs. Their restored overlap maps
are the identity on an open neighbourhood of that union.


\item  The designated parent prime divisors cannot be removed:
each meets one of these unchanged Hopfs. Their marked
elliptic neighbourhoods therefore persist. At every smooth
NC stage the exact maps \eqref{cdp:eq:C.AP.2}--\eqref{cdp:eq:C.AP.4} can be applied
again to their current normalizations. No holomorphic
contraction of a Hopf neighbourhood has been invoked.


\end{enumerate}

\item \par\medskip\noindent\textbf{The fixed finite termination potential.}
Let \(c_j\) count the retained global components of order
\(j\), for \(1\le j\le M\). Use the block-level potential
\[
                  (c_1,c_2,\ldots,c_M,h).
 \tag{C.AP.14}\hypertarget{tag-009R}{}\label{cdp:eq:C.AP.14}
\]
All coordinates are nonnegative integers. A direct type-II
step on a last designated-parent point leaves the cards
unchanged and decreases \(h\) by \eqref{cdp:eq:C.AP.6}. Either of the
other last-point alternatives supplies a finite progress
block with a strict decrease of the first changed card.
Any finite defect increase inside that whole block is
irrelevant to this lexicographic decrease. The missing-fibre
and cleanup comparisons keep all retained orders at most
the original \(M\); their auxiliary common blowups are not
additional retained cards. Hence \eqref{cdp:eq:C.AP.14} has no infinite
strictly descending sequence.


\begin{enumerate}[label={$\langle2\rangle$\arabic*.},leftmargin=1.3em,labelsep=.5em]
\item  When this procedure stops, a positive designated-parent
defect would still have a last point curve. \ref{cdp:st:C.2.UN.80} and the
three exhaustive alternatives would supply another strict
decrease. Thus every such defect is zero. Each designated
normalization is the actual minimal elliptic model in
\eqref{cdp:eq:C.AP.2}, and every boundary is a fibre support. All new
self-incidences created by the type-I/type-II and ruled
sigma charts are rational exceptional curves. None can
remain on these parents, since their rational boundary
would have to be exceptional for a point map of defect
zero. No elliptic self-incidence was introduced: a ruled
removal has only one horizontal boundary section, and all
its other changes are at its rational fibre curves. The
designated parents are consequently embedded, with disjoint
elliptic boundaries and no triple point.


\item  Every arrow used above is one of the actual two-ruling,
curve, point, or restored sigma-atlas maps already specified.
The inverse record, including all normalization branches
and overlap restrictions, reconstructs the given family.
The proper graph gluing and annular comparison of \ref{cdp:st:C.2.UN.75}
give the same record on a compact base, and its normal-cone
Euler formula applies at every retained step. This proves
the stated marked-interface reduction in its exact scope.



\end{enumerate}

\end{enumerate}
\end{proof}






\begin{proposition}[Every positive maximal Hopf after minimalization is a torsion-normal leaf]\hypertarget{tag-01EU}{}\label{cdp:st:C.2.UN.82}
\StatementRecord{cdp:st:C.2.UN.82}{proved}{01EU}{\SeriesRef{cdp:st:C.2.UN.75}, \SeriesRef{cdp:st:C.2.UN.77}, \SeriesRef{cdp:st:C.2.UN.70}, \SeriesRef{cdp:st:C.2.UN.19}, \SeriesRef{cdp:eq:C.ML.5}, \SeriesRef{cdp:eq:C.ML.6}, \SeriesRef{cdp:eq:C.ML.8}, \SeriesRef{cdp:eq:C.ML.3}}

Let \(f:X\to\Delta\) be a proper degeneration of connected
complex two-tori from a smooth threefold, with full reduced
NC central fibre and smooth normalization components. Retain
every normalized branch of a self-double curve, as in \ref{cdp:st:C.2.UN.75}.
For the primitive canonical evaluation write
\[
       \operatorname{div}(\Omega)=\sum_i k_iV_i,
       \qquad \min_i k_i=0,\qquad M=\max_i k_i .
\tag{C.ML.1}\hypertarget{tag-00QR}{}\label{cdp:eq:C.ML.1}
\]
On the normalization \(T\) of a component of order \(M\),
suppose that one boundary branch \(C\) is smooth elliptic,
is disjoint from all its other boundary branches, and has
coefficient one in its induced anticanonical divisor.
Then every central component has order zero and there
are no triple points.

Consequently consider the actual endpoint of \ref{cdp:st:C.2.UN.77} for
an original full reduced SNC family with \(M>0\).
Every Hopf component \(H\) of this maximal order is a
terminal component. Its boundary is one smooth elliptic
curve \(C=H\cap P\), and its exact canonical and normal
data satisfy
\[
 \begin{gathered}
       -K_H=aC,\qquad a=M-k_P+1\ge2,\\
       N_{C/H}^{\otimes(a-1)}\simeq\OO_C,
       \qquad N_{C/P}\simeq N_{C/H}^{-1},\\
       k_P\le M-1,\qquad
       \operatorname{coeff}_C(K_P)=a-2\ge0 .
 \end{gathered}
\tag{C.ML.2}\hypertarget{tag-00QT}{}\label{cdp:eq:C.ML.2}
\]
The normal-line equality uses the actual local branches
on \(P\) when another part of that component is nonnormal.
Each such interface therefore supplies precisely the
marked torsion-normal, nonnegative-coefficient data of
\ref{cdp:st:C.2.UN.70} on that normalization. No neighbourhood is replaced
in this conclusion; the ambient comparison remains the
finite record of \ref{cdp:st:C.2.UN.77}.
\end{proposition}



\begin{proof}

\begin{enumerate}[label={$\langle1\rangle$\arabic*.},leftmargin=*]
\item \par\medskip\noindent\textbf{The induced divisor at the maximal order.}
The normalization adjunction map for the reduced NC
fibre, including both branches over every self-double
curve, gives the actual meromorphic canonical section
on \(T\) with
\[
      -K_T=D_T=\sum_B (M-k_{\mathrm{other}(B)}+1)B.
\tag{C.ML.3}\hypertarget{tag-00QU}{}\label{cdp:eq:C.ML.3}
\]
Every displayed coefficient is a positive integer.
For a self-double branch its coefficient is one.
This is the restriction of the common volume form and
the principal full-fibre trivialization; it retains
the line isomorphism, rather than just intersections.
The support of \(D_T\) is exactly the full marked
boundary. The curve \(C\) in the statement is one
isolated component of this effective divisor, with
coefficient one.

\item We first prove the following surface consequence of
these data: the whole divisor \(D_T\) is reduced and
is a disjoint union of smooth elliptic curves. The
minimal model of \(T\) is ruled or Hopf. Indeed the
nonzero effective anticanonical divisor makes every
positive pluricanonical section vanish: through
\(K_T^{\otimes n}\simeq\OO_T(-nD_T)\), such a
section would be a holomorphic function on the compact
connected surface vanishing on \(D_T\). Thus
\(\kappa(T)=-\infty\). The surface classification
and the actual component restriction \ref{cdp:st:C.2.UN.19} give the
stated ruled or Hopf alternatives.

\par\medskip\noindent\textbf{Rational surfaces: the entire divisor is the marked elliptic curve.}
For a rational \(T\), the exact divisor sequence is
\[
 0\longrightarrow K_T\longrightarrow\OO_T
          \longrightarrow\OO_{D_T}\longrightarrow0.
\tag{C.ML.4}\hypertarget{tag-00QV}{}\label{cdp:eq:C.ML.4}
\]
Here \(H^0(K_T)=0\) and
\(H^1(K_T)=H^1(\OO_T)^*=0\) by Serre duality.
Consequently \(H^0(\OO_{D_T})=\C\). An effective
divisor with disconnected support has independent
constant functions on its disjoint components, so
its support is connected. Since \(C\) is isolated
in that support and has coefficient one, the exact
divisor is \(D_T=C\).

\par\medskip\noindent\textbf{Keeping the marked curve through the surface point maps.}
Take an actual sequence of point blowdowns
\(\pi:T\to T_0\) to a smooth minimal ruled or Hopf
surface. Push the meromorphic canonical section
through their common function-field identification.
The resulting anticanonical divisors remain effective.
At a reverse point blowup \(\sigma\), their exact
comparison is
\[
                 D_{\mathrm{new}}=\sigma^*D_{\mathrm{old}}-E.
\tag{C.ML.5}\hypertarget{tag-00QW}{}\label{cdp:eq:C.ML.5}
\]
The curve \(C\) is never exceptional, since it is
elliptic. Its image remains a smooth isolated
coefficient-one component at every blowdown. To
check this assertion at a last exceptional curve
\(E\), first suppose its coefficient in the
effective divisor is positive. It is then disjoint
from \(C\), so the contraction is an isomorphism
near \(C\). If its coefficient is zero, \eqref{cdp:eq:C.ML.5}
says that the old divisor has total multiplicity
exactly one at the centre. When \(E\) meets \(C\),
this forces the image of \(C\) to be the only
positive branch there, of multiplicity one. It is
smooth and no other boundary image meets it.
Equivalently, a curve of multiplicity \(r\) at
the centre has total intersection \(r\) with
the exceptional divisor; here \(r=1\).
This proves the assertion by induction through
the specified point maps. Write \(C_0=\pi(C)\).

\par\medskip\noindent\textbf{Irrational ruled surfaces: two coefficient-one sections.}
Let \(p:T_0\to B\) be the minimal ruling. The
elliptic curve \(C_0\) cannot be contained in a
rational ruling fibre. Its map to \(B\) is finite,
so Riemann--Hurwitz gives \(g(B)\le1\).
The rational case has already been proved; hence
\(B\) is elliptic and that finite map is etale.
For a ruling fibre \(F\),
\[
                 D_{T_0}\cdot F=-K_{T_0}\cdot F=2.
\tag{C.ML.6}\hypertarget{tag-00QX}{}\label{cdp:eq:C.ML.6}
\]
Every positive vertical curve would meet \(C_0\),
contrary to its isolation. Thus there is no such
vertical part. If \(C_0\) has degree two over
\(B\), \eqref{cdp:eq:C.ML.6} gives \(D_{T_0}=C_0\).
This connected anticanonical bisection is impossible
by the following exact residue and trace comparison.
The nowhere-zero section of \(K_{T_0}(C_0)\)
has nonzero residue on \(C_0\), so in
\[
 0\longrightarrow K_{T_0}\longrightarrow K_{T_0}(C_0)
       \xrightarrow{\operatorname{res}}K_{C_0}\longrightarrow0
\tag{C.ML.7}\hypertarget{tag-00QY}{}\label{cdp:eq:C.ML.7}
\]
the connecting map \(H^0(K_{C_0})\to H^1(K_{T_0})\)
is zero. Its Serre-dual map is restriction
\(H^1(\OO_{T_0})\to H^1(\OO_{C_0})\).
Composed with the ruling isomorphism from
\(H^1(\OO_B)\), this is the pullback of the
degree-two etale map \(C_0\to B\). Trace composed
with that pullback is multiplication by two,
which is nonzero. This contradiction excludes
the bisection.

\item The remaining possibility in \eqref{cdp:eq:C.ML.6} is
\[
                     D_{T_0}=C_0+C'_0,
\tag{C.ML.8}\hypertarget{tag-00QZ}{}\label{cdp:eq:C.ML.8}
\]
with two disjoint degree-one sections, each of
coefficient one. A degree-one irreducible curve
over the smooth ruling base is the graph of
its section and is smooth. Thus both curves
are smooth elliptic.

\item Pull \eqref{cdp:eq:C.ML.8} through the actual reverse point
blowups. Every intermediate divisor is effective:
a negative coefficient would remain negative
on its later strict transform. A centre away
from the positive support would create coefficient
\(-1\) in \eqref{cdp:eq:C.ML.5}, so no such centre occurs.
Starting with the two disjoint coefficient-one
curves, a centre on either has total positive
multiplicity one and creates coefficient zero.
The positive strict transforms remain disjoint
and smooth. Induction shows that no positive
exceptional coefficient ever appears. Therefore
\(D_T\) consists exactly of their two smooth
elliptic strict transforms, still with coefficient
one. All these are comparisons through the actual
surface morphism \(\pi\).

\par\medskip\noindent\textbf{Hopf surfaces: one simple factor makes the anticanonical divisor reduced.}
Use the actual unramified primary cover of the
Hopf minimal surface \(T_0\). The coefficient-one
curve \(C_0\) has reduced unramified inverse
image and retains that coefficient. In a primary
normal form the pulled-back anticanonical section
is \(g(z,w)\partial_z\wedge\partial_w\).
Hartogs extends \(g\) across zero, and its
functional equation is the constant-Jacobian
equation from (C.HA.5) and (C.CZ.11).

\item For a nonlinear contraction
\((z,w)\mapsto(\beta^r z+w^r,\beta w)\), its
only possible eigenfunction is a nonzero multiple
of \(w^{r+1}\), with \(r\ge1\). Its single
curve has coefficient at least two, contradicting
the given simple factor.

\item For a diagonal contraction without a resonance,
the section is a multiple of \(zw\). In the
resonant case pass through the actual finite
unramified central-power cover, so its multipliers
are \((\gamma^p,\gamma^q)\), with coprime
positive \(p,q\). Its monomials have weight
\(p+q\). The exact equation
\(pi+qj=p+q\) gives
\(i=1+q\ell\), \(j=1-p\ell\) for an integer
\(\ell\), by coprimality. For \(p,q>1\),
nonnegative exponents force \(\ell=0\), so the
only monomial is \(zw\). For \(p>1,q=1\), the section has form
\[
                        g=w(az+bw^p).
\tag{C.ML.9}\hypertarget{tag-00R0}{}\label{cdp:eq:C.ML.9}
\]
When \(a=0\) its only factor has order \(p+1\),
so the given coefficient one forces \(a\ne0\).
Then its two distinct smooth factors each have
order one. The case \(p=1,q>1\) is the same
with the coordinates interchanged. For \(p=q=1\),
the nonzero quadratic is either the square of
one linear form or the product of two distinct
linear forms. The simple factor excludes the
first alternative. Thus the entire pulled-back
divisor is reduced, consisting of two disjoint
elliptic curves on the primary Hopf surface.
Unramified descent retains every coefficient,
so the original \(D_{T_0}\) is reduced as well.
Its curves are disjoint smooth elliptic curves.

\item The same point-map induction following \eqref{cdp:eq:C.ML.8}
now applies to \(T\): every positive coefficient
is one, positive branches never meet, and the
additional exceptional coefficients are zero.
It proves the asserted surface consequence for
every blown-up Hopf component too.

\par\medskip\noindent\textbf{Propagation through the actual central support.}
The preceding arguments apply to the smooth
normalization of every maximal component carrying
one isolated coefficient-one elliptic branch.
They show that all its boundary branches have
these same properties. By \eqref{cdp:eq:C.ML.3}, the component
across each one again has order \(M\). Since the
branch has no triple point, its counterpart on
that normalization is again an isolated smooth
elliptic branch of coefficient one. The argument
can therefore be repeated across every incident
edge. Self-incidences retain both normalization
branches and introduce no new component order.

\item The finite incidence graph of the connected
central support is connected. The resulting
set of components of order \(M\) has no boundary
edge to its complement; it is consequently the
whole support. Every normalized boundary is
elliptic and isolated, so no triple point occurs.
Primitivity of the original direct-image generator
gives \(\min_i k_i=0\), hence \(M=0\).
This proves the first assertion.

\par\medskip\noindent\textbf{The remaining maximal Hopfs are the required leaves.}
At the endpoint of \ref{cdp:st:C.2.UN.77} the maximal Hopf components
are actual embedded minimal Hopf surfaces, and
their boundary curves are disjoint smooth elliptic
curves. The comparison preserves their orders and
the primitive canonical direct-image line. For
\(M>0\), the first assertion excludes coefficient
one on any such boundary.

\item The eigenfunction classification just proved
also shows that an effective anticanonical divisor
on a minimal Hopf surface either is reduced or
is supported on one elliptic curve: the nonlinear
form is a pure power, and the only nonreduced
diagonal alternatives above are pure powers of
one smooth factor. This conclusion descends
through the finite unramified primary cover.
Thus the nonempty boundary here is exactly one
curve \(C\), with coefficient \(a\ge2\).
Its unique neighbour has order \(M-a+1\).

\item Adjunction on this actual elliptic curve gives
\[
     \OO_C\simeq K_C=(K_H+C)|_C
                          \simeq N_{C/H}^{\otimes(1-a)}.
\tag{C.ML.10}\hypertarget{tag-00QS}{}\label{cdp:eq:C.ML.10}
\]
There are no triple points on \(C\), so the
full reduced-fibre trivialization gives
\(N_{C/H}\otimes N_{C/P}\simeq\OO_C\).
Finally the exact coefficient on \(P\) is
\(M-k_P-1=a-2\). These are precisely all the
claims in \eqref{cdp:eq:C.ML.2}. The argument uses only the
retained normalized branches, actual surface
point maps, and finite primary covers. It
does not prune a Hopf neighbourhood or assert
a different smoothing of the remaining support.
The exact weight spaces, polynomial factor cases,
nonlinear translation kernels, labelled point-map
forests and maximal-leaf order identities are
independently checked in
\path{code/root_maximal_hopf_leaf.py}; its executed
output is retained beside it. The proof above supplies
the general maps and unbounded arguments.

\end{enumerate}
\end{proof}






\begin{proposition}[The negative-coefficient interfaces have exact elliptic or one-section models]\hypertarget{tag-01EV}{}\label{cdp:st:C.2.UN.83}
\StatementRecord{cdp:st:C.2.UN.83}{proved}{01EV}{\SeriesRef{cdp:st:C.2.UN.74}, \SeriesRef{cdp:st:C.2.UN.77}, \SeriesRef{cdp:st:C.2.UN.81}, \SeriesRef{cdp:st:C.2.UN.67}, \SeriesRef{cdp:st:C.2.UN.19}, \SeriesRef{cdp:st:C.2.UN.31}, \SeriesRef{cdp:st:C.2.UN.70}, \SeriesRef{cdp:st:C.2.UN.80}, \SeriesRef{cdp:eq:C.TI.5}, \SeriesRef{cdp:eq:C.TI.12}, \SeriesRef{cdp:eq:C.TI.9}, \SeriesRef{cdp:eq:C.TI.13}}

Let \(T\) be a smooth normalization component of a full
reduced NC fibre of a smooth-total degeneration of connected
complex two-tori. Retain the two normalized branches of
every self-double curve. Let \(J\subset T\) be a smooth
elliptic boundary branch, disjoint from every other boundary
branch, whose exact normal line is torsion. No sign is
imposed on its coefficient in the actual meromorphic
canonical divisor. Then \(T\) is not rational, and one
of the following explicit comparisons holds.

\begin{enumerate}
\item There is an actual sequence of point blowdowns
\[
 T\xrightarrow{\pi}Z\xrightarrow{q_Z}B
 \tag{C.TI.1}\hypertarget{tag-0135}{}\label{cdp:eq:C.TI.1}
\]
which is an isomorphism on a neighbourhood of \(J\).
The map \(q_Z\) is a relatively minimal fixed-Jacobian
elliptic quasi-bundle with identity integral monodromy
and base genus at most one. Its fibres have smooth
elliptic supports. Every original boundary and every
exceptional curve of \(\pi\) is vertical, and every
exceptional prime is an actual boundary.

\item There is an actual point map \(\pi:T\to H\),
an isomorphism near \(J\), onto a nonlinear minimal
Hopf surface. Its only curve is \(J_0=\pi(J)\).
For an integer \(a\ge2\), the actual canonical
comparison is
\[
 K_H=-aJ_0,\qquad
 K_T=-aJ+\sum_E b_EE,\qquad b_E>0.
 \tag{C.TI.2}\hypertarget{tag-013B}{}\label{cdp:eq:C.TI.2}
\]
The other boundary support is exactly the exceptional
set of \(\pi\). Every neighbouring order along such
an exceptional curve exceeds \(k_T+1\).

\item There is an actual point map \(\pi:T\to R\),
an isomorphism near \(J\), onto a projective-line
bundle over an elliptic curve. The image \(J_0\)
is a section, and
\[
 K_R=-2J_0,\qquad N_{J_0/R}\simeq\OO_{J_0},\qquad
 K_T=-2J+\sum_E b_EE,\quad b_E>0.
 \tag{C.TI.3}\hypertarget{tag-013C}{}\label{cdp:eq:C.TI.3}
\]
Every other boundary is exactly one of these
exceptional primes. After tensoring its defining
rank-two bundle by a line, \(R\) is
\(\PP(V)\) for an extension
\[
              0\longrightarrow\OO\longrightarrow V
                    \longrightarrow\OO\longrightarrow0.
 \tag{C.TI.4}\hypertarget{tag-013D}{}\label{cdp:eq:C.TI.4}
\]
The extension is retained and need not split.
If \(P\) is the component across \(J\), then
\(k_T=k_P+1>0\). This is the actual one-section
ruled-removal situation of \ref{cdp:st:C.2.UN.74} and the finite
restored ruled block of \ref{cdp:st:C.2.UN.77}/\ref{cdp:st:C.2.UN.81}.
\end{enumerate}

In the ruled birational case, alternative (3) occurs
only when \(J\) is a degree-one section of the actual
ruling and no other horizontal boundary survives the
point contractions disjoint from \(J\). A higher-degree
elliptic multisection, or a second surviving horizontal
boundary, produces the actual elliptic map in (1).
All comparisons retain the original marked normal line
and the actual boundary branches. No replacement of an
arbitrary threefold neighbourhood by a prescribed
logarithmic model is used.
\end{proposition}



\begin{proof}

\begin{enumerate}[label={$\langle1\rangle$\arabic*.},leftmargin=*]
\item \par\medskip\noindent\textbf{The exact canonical data and the rational exclusion.}
Write the actual restricted canonical section as
\[
              K_T=a_JJ+\sum_C a_CC.
 \tag{C.TI.5}\hypertarget{tag-013E}{}\label{cdp:eq:C.TI.5}
\]
The sum includes every other actual boundary branch,
even when its coefficient is zero. All intersections
of \(J\) with that support are empty. Its torsion
normal gives \(J^2=0\) and \(K_T\cdot J=0\).
\ref{cdp:st:C.2.UN.67} already excludes a rational component under
exactly these hypotheses, without a coefficient
sign assumption. Its actual restriction sequences
produce \(|mJ|\), with
\(m=\operatorname{ord}N_{J/T}\); the exact order
forces connectedness, and the limiting ambient
fields give the coherent-Euler contradiction.
The same proof works on this smooth normalization:
its fields are holomorphic off the full normalized
boundary, and its marked neighbourhood is embedded.

\item The actual component restriction \ref{cdp:st:C.2.UN.19} and \ref{cdp:st:C.2.UN.31}
leaves a Hopf, a nonrational ruled, or a
nonnegative-Kodaira-dimension minimal model.

\par\medskip\noindent\textbf{Nonnegative Kodaira dimension.}
Use the ordinary minimal-model morphism
\(\pi:T\to Z\). Its explicit canonical base-line
formula, or the trivial canonical line in the
torus case, gives \(K_Z\cdot C\ge0\) for every
curve \(C\). The exact point-map formula is
\[
 K_T=\pi^*K_Z+\sum_E b_EE,\qquad b_E>0.
 \tag{C.TI.6}\hypertarget{tag-013F}{}\label{cdp:eq:C.TI.6}
\]
The curve \(J\) is not exceptional. Adjunction
therefore gives
\[
 0=K_T\cdot J
   =K_Z\cdot\pi(J)+\sum_E b_E(E\cdot J).
 \tag{C.TI.7}\hypertarget{tag-013G}{}\label{cdp:eq:C.TI.7}
\]
Every summand is nonnegative, so the point map is
an isomorphism near \(J\). Its image is a smooth
elliptic curve with the same torsion normal.
The actual marked comparisons (C.NP.3)--(C.NP.4)
and (C.MR.12)--(C.MR.14) now apply without using
the sign of \(a_J\). In Kodaira dimension one,
the positive canonical base degree makes this
curve vertical. On a torus its quotient by the
parallel elliptic subtorus supplies the map.
In the remaining dimension-zero case, a
horizontal marked curve trivializes the Chern
vector after the specified finite translation
root cover; the actual period lattice is a
torus lattice, and the parallel marked curves
give the descended elliptic quotient.

\item Thus there is an actual elliptic map making
the image of \(J\) a full fibre support.
Every other noncontracted boundary is disjoint
from that support and hence vertical; every
contracted curve is vertical too. The ambient
field alignment and the translation-root
comparison of \ref{cdp:st:C.2.UN.70} then give \eqref{cdp:eq:C.TI.1}.
The proof of \ref{cdp:st:C.2.UN.80} uses only this actual map
and its vertical boundary, so gives its final
exceptional-prime assertion.

\par\medskip\noindent\textbf{Hopf minimal model.}
For a Hopf minimal model \(H\), every curve
is a smooth elliptic curve of square zero.
The blowup square formula
\[
                  0=J^2=J_0^2-\sum_\nu r_\nu^2
 \tag{C.TI.8}\hypertarget{tag-013H}{}\label{cdp:eq:C.TI.8}
\]
shows that no centre lies on a successive
image of \(J\). Thus \(\pi\) is again an
isomorphism near it. A diagonal Hopf axis
has torsion normal precisely in the resonant
elliptic case, through the actual finite
primary cover and its multiplier character.
In that case the actual Hopf elliptic map
makes every surviving boundary vertical;
its composition with \(\pi\), followed by
the original ambient-field comparison,
gives (1).

\item If \(H\) is nonlinear, its unique elliptic
curve is \(J_0\), so every other boundary
is contracted. Push \eqref{cdp:eq:C.TI.5} through the
actual common open isomorphism. On a
primary cover, the nonlinear normal form
has canonical divisor \(-(r+1)J_0'\) for
an integer \(r\ge1\). Comparing divisors
of the lifted actual meromorphic section
forces \(a_J=-(r+1)\): a different exponent
would give a nonzero principal divisor
supported on just one effective curve,
impossible on a compact surface. Unramified
descent preserves this coefficient. Put
\(a=r+1\). The point-map formula now gives
\eqref{cdp:eq:C.TI.2}. Since \eqref{cdp:eq:C.TI.5} has no other
noncontracted support, every exceptional
prime has its positive discrepancy as
coefficient and is an actual boundary.
Conversely every other boundary is
exceptional. Its neighbour has order
\(k_T+b_E+1\ge k_T+2\).

\par\medskip\noindent\textbf{The actual ruling and its elliptic multisection.}
It remains to treat a ruled birational
class. Its actual ruling is a morphism
\[
                       p:T\longrightarrow B_0
 \tag{C.TI.9}\hypertarget{tag-013I}{}\label{cdp:eq:C.TI.9}
\]
obtained by composing the specified point
map with a minimal ruling. The image of
the elliptic curve \(J\) is not a point:
all components of a ruling fibre are
rational. Riemann--Hurwitz makes \(B_0\)
elliptic and \(p|_J\) an etale isogeny,
of degree \(d\ge1\). The rational base
case was already excluded.

\item First let \(d\ge2\). Make the actual
finite etale base change
\[
 T'=T\times_{B_0}J\longrightarrow T.
 \tag{C.TI.10}\hypertarget{tag-0136}{}\label{cdp:eq:C.TI.10}
\]
The isogeny is Galois. Its group \(G\)
acts freely on \(T'\), and \(T'/G=T\).
The inverse image of \(J\) consists of
\(d\) disjoint sections \(J_1,\ldots,J_d\)
of \(p':T'\to J\), each with the pulled
torsion normal and square zero.

\item For two of these sections, Hodge index
gives numerical equivalence. Indeed their
intersection and both squares are zero;
an ample class determines a linear
combination orthogonal to that ample
class, still of square zero, which is
numerically zero. Intersection with a
general ruling fibre then makes its two
coefficients equal. On this point-blown-up
elliptic ruled surface \(H^2(T',\Z)\) is
torsion-free, and a numerically trivial
line lies in \(\operatorname{Pic}^0(T')\).
The actual ruling pullback
\[
        (p')^*:\operatorname{Pic}^0(J)
                       \xrightarrow{\sim}\operatorname{Pic}^0(T')
 \tag{C.TI.11}\hypertarget{tag-0137}{}\label{cdp:eq:C.TI.11}
\]
is an isomorphism: it is so for a
projective-line bundle by its direct
images, and each actual point blowup
preserves this exponential-sequence map.
Consequently
\(\OO(J_1-J_2)=(p')^*L\). Restriction to
the section \(J_1\) identifies \(L\)
with \(N_{J_1/T'}\), because \(J_2\)
is disjoint from \(J_1\).

\item Let \(m\) be the exact order of that
normal line. Then the two disjoint
divisors \(mJ_1,mJ_2\) are linearly
equivalent. Their sections have no
common zero and give an actual pencil
\[
                  q':T'\longrightarrow\PP^1.
 \tag{C.TI.12}\hypertarget{tag-0138}{}\label{cdp:eq:C.TI.12}
\]
It has connected fibres. To check this,
take its Stein factor with finite degree
\(e\). The inverse image of the value
with divisor \(mJ_1\) has just one
point on the Stein base, so that finite
map is totally ramified there of
order \(e\). The connected-fibre map
has fibre \((m/e)J_1\). Its local
principal equation gives
\(N_{J_1/T'}^{m/e}=\OO\), and exact
order \(m\) forces \(e=1\).
Adjunction, \(J_1^2=K_{T'}\cdot J_1=0\),
then makes the connected general fibre
of \eqref{cdp:eq:C.TI.12} elliptic.

\item Every \(J_i\) is vertical for \(q'\).
It is the whole reduced support of
its fibre: since \(J_i^2=0\), the
intersection of that full fibre
with \(J_i\) forces every other
component to be disjoint from it,
and connectedness excludes such
additional components.
For \(g\in G\) and a general connected
fibre \(F\), the intersection
\((gF)\cdot J_1=F\cdot g^{-1}J_1=0\)
makes \(gF\) vertical too. Thus \(G\)
acts on the actual Stein base and
\(q'\) descends through \(T'/G=T\).
The kernel of the base action acts
freely on a general elliptic fibre;
it therefore acts by translations,
and the descended general fibre
is elliptic. This gives an actual
connected elliptic map on \(T\).
The marked \(J\) is its whole reduced
fibre support. Every other original
boundary pulls back disjointly from
all the \(J_i\), hence is vertical
for \(q'\) by intersection and is
vertical after descent. The original
ambient fields can now be applied
on \(T\) itself. The rank-one return
identities or the rank-two transverse
flow give identity monodromy, and
vertical minimization gives exactly
(1), isomorphic near \(J\).

\par\medskip\noindent\textbf{A section: the contraction retaining its actual normal line.}
Now let \(d=1\), so \(J\) is a section
of \eqref{cdp:eq:C.TI.9}. Every reducible ruling
fibre admits a rational \((-1)\)-curve
disjoint from \(J\). Here is a direct
comparison with its actual point maps.
Take a last blowup in its finite
construction. If the last exceptional
does not meet \(J\), use it. Otherwise
the centre is the smooth point where
the old section meets the old fibre:
a section never passes through a
node or a multiple component.
If the old fibre is irreducible,
its strict transform is the required
disjoint \((-1)\)-curve. If it is
reducible, induction supplies an old
disjoint \((-1)\)-curve, unchanged
by this last blowup at the section.
The chosen curve has valence at most
two; the actual point blowdown
restores the corresponding smooth
or nodal fibre chart. Repeating
decreases the finite number of
components and is an isomorphism
near the section throughout.

\item Doing this over its finitely many
singular values gives
\[
 T\xrightarrow{\pi_0}R=\PP(V_0)
                    \xrightarrow{p_0}B_0,
 \qquad J_0=\pi_0(J),
 \qquad N_{J_0/R}=N_{J/T}.
 \tag{C.TI.13}\hypertarget{tag-0139}{}\label{cdp:eq:C.TI.13}
\]
All centres and their actual maps
are retained. This construction
does not push two previously
intersecting marked branches
together at the section.
Every other surviving boundary
is disjoint from \(J_0\).
A vertical ruling fibre meets
the section, so any such surviving
boundary is horizontal.

\item Suppose that one exists. Kill the
torsion normal by the actual finite
etale cover \(\widetilde B_0\to B_0\).
After a line twist, the pulled
projective bundle has section
normal \(\OO\); its complement
of that section is an additive
\(\OO_{\widetilde B_0}\)-torsor,
of class
\(\xi\in H^1(\widetilde B_0,\OO)\).
Normalize a connected component
of the pulled horizontal boundary.
It gives a finite curve map
\(a:C\to\widetilde B_0\) and an
actual section of the pulled
torsor, hence \(a^*\xi=0\).
The finite trace satisfies
\(\operatorname{tr}\circ a^*
 =(\deg a)\operatorname{id}\)
on \(H^1(\OO)\). Therefore
\(\xi=0\), including a possibly
ramified map \(a\).

\item The pulled marked bundle is thus
\(\widetilde B_0\times\PP^1\).
The finite deck maps preserve the
section at infinity. Their affine
coefficients on the other coordinate
are holomorphic functions on the
compact base, hence constants.
The finite affine group has a
common fixed point by averaging,
so its action preserves the actual
elliptic projection to \(\PP^1\).
The quotient gives an elliptic
quasi-bundle on \(R\), with \(J_0\)
a fibre support. Each other compact
horizontal boundary lifts into
\(\widetilde B_0\times\C\); its
second coordinate is constant
on its compact normalization.
It is therefore another fibre
of this elliptic map. Compose
with \(\pi_0\). All exceptional
curves and all original boundaries
are vertical. Applying the original
ambient fields and then minimizing
this elliptic map proves (1).

\item Finally suppose there is no other
surviving horizontal boundary.
Push the actual canonical section
\eqref{cdp:eq:C.TI.5} through \eqref{cdp:eq:C.TI.13}.
Its divisor is supported just on
\(J_0\). Intersecting a ruling
fibre gives its coefficient \(-2\).
Thus \(K_R=-2J_0\) as the exact
divisor of this meromorphic
section. Adjunction on the
elliptic section gives
\[
 \OO_{J_0}=K_{J_0}
          =(K_R+J_0)|_{J_0}=N_{J_0/R}^{-1}.
 \tag{C.TI.14}\hypertarget{tag-013A}{}\label{cdp:eq:C.TI.14}
\]
The section normal is actually
trivial. The defining quotient
of the rank-two bundle, after
tensoring by its kernel line,
is therefore the extension
\eqref{cdp:eq:C.TI.4}; no vanishing of its
extension class has been used.

\item The actual point-map formula is
now \eqref{cdp:eq:C.TI.3}. Every exceptional
prime has positive discrepancy,
so comparison with \eqref{cdp:eq:C.TI.5}
forces it to be a boundary.
Every other boundary is
exceptional by construction.
The coefficient \(-2\) on \(J\)
is \(k_P-k_T-1\), so
\(k_T=k_P+1>0\). The surviving
ruled model has exactly its
one horizontal boundary section;
the remaining full boundary on
\(T\) consists of vertical
exceptional configurations.
These are precisely the actual
one-section inputs to the
restored finite ruled-removal
block of \ref{cdp:st:C.2.UN.77}/\ref{cdp:st:C.2.UN.81}, whose
bundle stage is \ref{cdp:st:C.2.UN.74}.
This proves (3) with all
normal and extension data.

\item All three cases use actual point
maps, finite covers, trace maps,
and descended Stein fibrations.
The finite section-fibre graph
checks and exact normal-order
divisibility certificates are
executed in
\path{code/A_torsion_interfaces.py}.

\end{enumerate}
\end{proof}






\begin{proposition}[The highest Hopf order can be made the global maximum]\hypertarget{tag-01EW}{}\label{cdp:st:C.2.UN.84}
\StatementRecord{cdp:st:C.2.UN.84}{proved}{01EW}{\SeriesRef{cdp:st:C.2.UN.82}, \SeriesRef{cdp:st:C.2.UN.75}, \SeriesRef{cdp:st:C.2.UN.77}, \SeriesRef{cdp:eq:C.TK.4}, \SeriesRef{cdp:eq:C.TK.5}, \SeriesRef{cdp:st:C.2.UN.19}, \SeriesRef{cdp:eq:C.TK.3}, \SeriesRef{cdp:eq:C.TK.6}; \cite[\S7]{Kulikov1981}; \cite[\S\S4.8.5--4.8.12,4.16.2--4.16.6]{Kulikov1977}; \cite[Lemma2.19]{Kulikov1977}}

Let \(f:X\to\Delta\) be a proper degeneration of connected
complex two-tori with smooth total space and full reduced
SNC central fibre. Use its actual primitive canonical
generator and write
\[
 \operatorname{div}(\Omega)=\sum_i k_iV_i,\qquad
 \min_i k_i=0,\qquad R=\max_i k_i.
 \tag{C.TK.1}\hypertarget{tag-013J}{}\label{cdp:eq:C.TK.1}
\]
Let \(\mathcal H\) be the set of prime components whose
smooth minimal model is Hopf, and suppose that
\[
                  h=\max_{i\in\mathcal H}k_i>0.
 \tag{C.TK.2}\hypertarget{tag-013K}{}\label{cdp:eq:C.TK.2}
\]
There is a finite overlap-compatible modification record,
over the same disc and isomorphic on its punctured part,
whose endpoint has smooth total space and full reduced
NC fibre, such that its global maximum canonical order
is exactly \(h\). Every original Hopf-birational prime
survives with its original order, and no new such prime
appears. Every component of order \(h\) with Hopf minimal
model can, in the same finite record, be made an embedded
minimal Hopf surface. These maximal Hopfs are the
torsion-normal elliptic leaves of \ref{cdp:st:C.2.UN.82}.

At all retained global stages the canonical orders lie
between zero and the original \(R\). A whole-component
removal has positive order \(q\); its possible added
components have order \(q+1\), bounded by the current
maximum, and there are no additions when \(q\) is that
maximum. The removed \(q\) is not required to exceed
\(h\). The termination potential therefore counts every
positive order, including orders at most \(h\).

The record retains its actual centres, normalization
branches, local sigma maps and restored overlaps.
Reversing it recovers the given analytic family. On
an original compact family it gives the same actual
modification over its compact base. Its Euler changes
are the signed normal-cone terms of \ref{cdp:st:C.2.UN.75}.
\end{proposition}



\begin{proof}

\begin{enumerate}[label={$\langle1\rangle$\arabic*.},leftmargin=*]
\item \par\medskip\noindent\textbf{The precise input to the component selection.}
We use the actual selection argument, rather than
applying the conclusion of Kulikov's canonical-flat
theorem to a family that does not satisfy its
Hopf-minimum hypothesis. The relevant distinction
is explicit in \cite[\S7]{Kulikov1981}: the search
for a component whose normalization is ruled with
one boundary section uses only that the current
maximal components are ruled. Section8 removes
the earlier dual-complex hypothesis and supplies
the missing-fibre blowups. Thus a lower-order Hopf
component is permitted during this selection.

\item Here are the actual maps and the order information
from that argument which we retain. The boundary
analysis at the maximal ruled components in
\cite[\S\S4.8.5--4.8.12,4.16.2--4.16.6]{Kulikov1977}
produces either a plane with one line or a ruled
normalization with one section and vertical boundary.
The two-section cases are converted by the specified
finite type-I/type-II moves along their section
chains or cycles; a rational tree gives the complete
one-triple ruling curve. At a singular sigma chart,
the last locally blown-up branch of largest order
supplies the same one-section configuration. These
are the actual two-ruling, complete-curve and
sigma-atlas maps recorded in \ref{cdp:st:C.2.UN.75} and \ref{cdp:st:C.2.UN.77}.

\item The ruled candidate need not be maximal. If its
order is \(q\), restriction of the common canonical
section to a generic complete ruling fibre gives
the section coefficient \(-2\), hence
\[
       k_{\mathrm{section\ parent}}=q-1.
 \tag{C.TK.3}\hypertarget{tag-013L}{}\label{cdp:eq:C.TK.3}
\]
For a plane candidate, the line coefficient \(-3\)
gives the parent order \(q-2\). Consequently every
removed component has positive order. When a ruled
candidate has disconnected boundary, the corrected
procedure inserts the missing fibre components.
Their actual curve-blowup Jacobians give order
\(q+1\). If \(q\) is the current maximum, the
effective anticanonical one-section comparison
of \cite[Lemma2.19]{Kulikov1977} makes the boundary
connected and no insertion is made. In every case
\[
 \begin{split}
 &\text{one component of order }q>0\text{ is removed},\\
 &\text{only finitely many components of order }q+1
       \text{ may be added},\\
 &q+1\le R_{\mathrm{current}}\text{ whenever an addition occurs}.
 \end{split}
 \tag{C.TK.4}\hypertarget{tag-013M}{}\label{cdp:eq:C.TK.4}
\]
This is an assertion about the retained degeneration
stages. It does not assign these bounds to auxiliary
common blowups in a two-ruling diagram.

\item The local contractions are followed by the exact
restoration of the overlap, using the same composite
map. The proper-graph comparison in (C.KG.2) and
(C.KG.6) therefore gives an actual bimeromorphic
map over the disc, not merely equal central graphs.
If singular sigma charts remain, we use the
largest last locally blown-up branch and its finite
descendant cleanup from \ref{cdp:st:C.2.UN.77}. Such branches are
ruled by their actual complete exceptional fibres.
This cleanup uses no hypothesis on the other
Hopf orders and has the same positive-level updates
\eqref{cdp:eq:C.TK.4}. It ends at a smooth-total NC stage with
the labelled normalization branches retained.

\par\medskip\noindent\textbf{Why selection remains available above \(h\).}
The proper composite-graph argument in \ref{cdp:st:C.2.UN.75} proves
that every Hopf-birational prime survives every
recorded removal. A general complete rational
exceptional fibre cannot cover a point-blown-up
Hopf surface; a compact projective point exceptional
cannot map bimeromorphically onto it. The actual
canonical direct-image isomorphism (C.KG.6) preserves
its canonical order. Every added prime is ruled or
projective over the specified centre, so it is not
Hopf-birational. It follows that the set of Hopf
orders, with its original labels, is unchanged.
In particular its highest value remains \(h\).

\item Consider a current smooth-total NC stage with
maximum \(R'>h\). On the smooth normalization
\(T\) of a maximal component, actual adjunction is
\[
       -K_T=\sum_C (R'-k_{\mathrm{other}(C)}+1)C.
 \tag{C.TK.5}\hypertarget{tag-013N}{}\label{cdp:eq:C.TK.5}
\]
Both preimages of a self-double curve are counted;
their coefficients are one. The divisor is nonzero:
the connected support contains a component of Hopf
order \(h<R'\), so no maximal component is an
isolated component of the whole support. Every
coefficient in \eqref{cdp:eq:C.TK.5} is positive. A positive
pluricanonical section would be a holomorphic
function on the compact connected surface vanishing
on this effective divisor, which is impossible.
Thus \(\kappa(T)=-\infty\).

\item The actual component theorem \ref{cdp:st:C.2.UN.19} makes a
nonruled component of this Kodaira dimension
Hopf-birational. No such component has order
\(R'>h\). Hence every current maximal
normalization is ruled or rational, including
the possible plane. This is exactly the input
to the selection proved above. Controlled
singular sigma charts are first handled by
their specified last-local ruled branches;
they do not require an anticanonical formula
on a singular ambient chart.

\par\medskip\noindent\textbf{The stopped finite-level induction.}
For each \(1\le j\le R\), let \(c_j\) be the
number of retained global components of order
\(j\), and order the tuples
\[
                     (c_1,c_2,\ldots,c_R)
 \tag{C.TK.6}\hypertarget{tag-013O}{}\label{cdp:eq:C.TK.6}
\]
lexicographically, with the smaller index first.
Each whole-component block \eqref{cdp:eq:C.TK.4} strictly
decreases this tuple: it deletes a card of
order \(q\), leaves all smaller orders unchanged,
and only increases the next order. Type-I and
type-II preparatory moves preserve the card
counts. A fixed finite product of copies of
\(\mathbb N\) in this order has no infinite
strict descending sequence. The finite cleanup
blocks obey the same comparison.

\item We continue only while the maximum exceeds
\(h\), and finish the finite cleanup belonging
to the last selected block. If the maximum
still exceeded \(h\), \eqref{cdp:eq:C.TK.5} would supply
another selection. The decreasing tuple
therefore forces a finite stop with maximum
at most \(h\). The retained Hopf prime of
order \(h\) forces equality. No order-zero
prime has been removed, by \eqref{cdp:eq:C.TK.3} and its
plane counterpart, so the transported canonical
generator remains primitive with minimum zero.

\item Counting only the components above \(h\)
would not prove this induction. For example,
removing an order-one ruled candidate and
adding order-two components changes none of
those upper counts when \(h\ge2\). Such a
candidate can be the ruled neighbour produced
by the one-triple step. The full tuple
\eqref{cdp:eq:C.TK.6} is strictly decreased by the exact
update and is the potential actually used.

\par\medskip\noindent\textbf{The actual maximal Hopf components at the stopping stage.}
The original central components were embedded.
The record has introduced no elliptic self-double
curve on a surviving Hopf component. Type-I and
type-II changes create incidences on their
rational exceptional curves. A removed ruled
component has one horizontal boundary section;
its generic contraction cannot identify two
elliptic horizontal branches. Its other changes
are on its rational fibre curves, and its
restored overlaps retain the original elliptic
maps. Point contractions likewise only change
the indicated rational boundary configuration.

\item At the stopping stage a Hopf normalization
of order \(h\) has the effective divisor
\eqref{cdp:eq:C.TK.5}, now with \(R'=h\). The actual
point-map induction (C.HM.3)--(C.HM.4)
shows that a rational coefficient-one boundary
is a leaf, and that two coefficient-one
boundaries never meet. A rational self-double
curve would have two coefficient-one preimages,
each with negative self-square and exactly
one boundary intersection. Their actual normal
product would consequently give
\[
                        C_1^2+C_2^2=-1,
 \tag{C.TK.7}\hypertarget{tag-013P}{}\label{cdp:eq:C.TK.7}
\]
contradicting the two negative integer squares.
A triple point with three branches of this same
component would already give two intersecting
coefficient-one curves. Thus these maximal
Hopfs are embedded at the stopped NC stage.

\item All the prerequisites of the last-point and
finite-cleanup proof of \ref{cdp:st:C.2.UN.77} now hold on this
actual recorded stage. That proof uses its
smooth normalizations, the labelled NC branches
elsewhere and the embedded maximal Hopfs; it
does not require the other components to be
embedded. Apply those exact moves to remove
their point defects. The retained global orders
stay at most \(h\), every Hopf prime and order
is preserved, and the endpoint maximal Hopfs
are smooth minimal surfaces. \ref{cdp:st:C.2.UN.82} then supplies
their single elliptic boundaries and exact
torsion-normal line maps.

\item Finally the local proper graphs glue with the
recorded overlap isomorphisms, exactly as in
\ref{cdp:st:C.2.UN.75}. They give the punctured-family isomorphism
and the inverse record on its actual endpoint
neighbourhood. Annular gluing to the unchanged
outside performs the same operation on an
original compact family. The signed normal-cone
Euler formula (C.KG.3)--(C.KG.4) applies to
every retained arrow. No chosen smoothing of
the endpoint central space replaces that
actual neighbourhood.

\end{enumerate}
\end{proof}






\begin{proposition}[Exact root, birational action and coarsening records for the Hopf-minimum branch]\hypertarget{tag-01EP}{}\label{cdp:st:C.2.UN.78}
\StatementRecord{cdp:st:C.2.UN.78}{proved}{01EP}{\SeriesRef{cdp:st:C.2.UN.75}, \SeriesRef{cdp:st:C.2.UN.39}, \SeriesRef{cdp:eq:C.RK.9}, \SeriesRef{cdp:eq:C.RK.11}, \SeriesRef{cdp:eq:C.RK.16}, \SeriesRef{cdp:eq:C.RK.21}, \SeriesRef{cdp:eq:C.RK.22}, \SeriesRef{cdp:st:C.2.UN.71}, \SeriesRef{cdp:st:C.2.UN.73}; \cite[Theorems~1.6 and1.10]{BierstoneMilman1997}}

Let $f:X\to\Delta_t$ be a proper degeneration of connected
complex two-tori from a smooth threefold, with arbitrary
central support $S=f^{-1}(0)_{\rm red}$. Take an actual
embedded resolution $\pi:Y\to X$ by smooth central
point and curve blowups, and write
\[
 \operatorname{div}(f\pi)=\sum_i m_iD_i,
 \qquad \operatorname{div}(\pi^*\Omega)=\sum_i k_iD_i,
 \qquad r_i=\frac{k_i+1}{m_i},
 \qquad N=\operatorname{lcm}_i(m_i),
\tag{C.RK.1}\hypertarget{tag-00ZQ}{}\label{cdp:eq:C.RK.1}
\]
where $\Omega$ generates $f_*K_X$. The $D_i$ are the
smooth components of the resolved SNC support. Form
the actual normalization and a toric semistable
resolution
\[
 \begin{gathered}
 Y_N=\operatorname{Norm}(Y\times_{\Delta_t}\Delta_s),
                \qquad t=s^N,\\
 Z\xrightarrow{\ q\ }Y_N\xrightarrow{\ \rho\ }Y,
                 \qquad G=\mu_N.
 \end{gathered}
\tag{C.RK.2}\hypertarget{tag-0101}{}\label{cdp:eq:C.RK.2}
\]
These maps can be chosen $G$-equivariantly, with $Z$
smooth and full reduced simple-NC central fibre.
All components added by $q$ are ruled or rational.

The following is an explicit finite-cover condition
on this resolution record: every strict root-cover
component whose smooth minimal model is Hopf lies
over an index $i$ with
\[
                            r_i=\min_j r_j.
\tag{C.RK.3}\hypertarget{tag-0105}{}\label{cdp:eq:C.RK.3}
\]
For exactly the records satisfying \eqref{cdp:eq:C.RK.3}, the
semistable root $Z$ admits the finite compatible
canonical-flat reduction of \ref{cdp:st:C.2.UN.75}. Denote its
actual comparison by $\phi:Z\dashrightarrow X_0$.
The flat endpoint has one of the three complete
core types in \ref{cdp:st:C.2.UN.75}, and it carries the actual
finite \emph{birational} action
\[
                 \gamma_0=\phi\gamma\phi^{-1}
                 \quad (\gamma\in G).
\tag{C.RK.4}\hypertarget{tag-0106}{}\label{cdp:eq:C.RK.4}
\]
Its group laws hold as equalities of meromorphic
maps; it acts over $s\mapsto\zeta s$ and is regular
on the punctured family. Retain the full record
$\phi$, the toric coarsening $q$, and $\pi$.
Their inverse and quotient operations reconstruct
the given original degeneration exactly:
\[
 X_0\ \dashrightarrow\ Z\ \xrightarrow{q}\ Y_N
       \ \longrightarrow\ Y_N/G=Y
       \ \xrightarrow{\pi}\ X.
\tag{C.RK.5}\hypertarget{tag-0107}{}\label{cdp:eq:C.RK.5}
\]
In particular the inverse record is applied
\emph{before} the finite quotient. No quotient of
an undefined central birational action is used.
Conversely the same operations on any actual
flat-core record with these retained regularization,
root-coarsening and smooth blowdown maps produce
a torus degeneration in the stated branch.
This gives complete finite generation for this
explicit resolution-and-root branch, without
an equivariant canonical-flat reduction assertion.

The Euler correction for the root is explicit.
For a nonempty index set $I$ with nonempty stratum
\[
 D_I^\circ=
    \left(\bigcap_{i\in I}D_i\right)
           \setminus\bigcup_{j\notin I}D_j,
 \qquad d_I=\gcd_{i\in I}m_i,
 \qquad \nu_I=\frac{N^{|I|-1}d_I}{\prod_{i\in I}m_i},
\tag{C.RK.6}\hypertarget{tag-0108}{}\label{cdp:eq:C.RK.6}
\]
one has $1\le|I|\le3$ and $\nu_I$ is a positive
integer. Put
\[
 \begin{split}
 R_N&=\sum_I(d_I\nu_I-1)e_c(D_I^\circ),\\
 \Delta_\pi&=2p_\pi+
                   \sum_\alpha(2-2g(C_\alpha)),
 \end{split}
\tag{C.RK.7}\hypertarget{tag-0109}{}\label{cdp:eq:C.RK.7}
\]
where $p_\pi$ and $C_\alpha$ are the actual point
and connected curve centres of $\pi$. Let
$\Delta_K$ be the signed normal-cone defect of
the forward \ref{cdp:st:C.2.UN.75} record from $Z$ to $X_0$.
Then the exact original-support formula is
\[
                e(S)=e(S_0)-\Delta_K-R_N-\Delta_\pi.
\tag{C.RK.8}\hypertarget{tag-010A}{}\label{cdp:eq:C.RK.8}
\]
Thus zero Euler number is equivalent to
$\Delta_K+R_N+\Delta_\pi=e(S_0)$, where the
right side is zero for the smooth or marked
cycle core, and $2[N_0:\Gamma_0]$ for the
toric core. The symbols $N_0,\Gamma_0$ here
are its two-dimensional lattices, separate
from the root exponent $N$ in \eqref{cdp:eq:C.RK.1}.
Formula \eqref{cdp:eq:C.RK.8} agrees with the universal
original-support formula of \ref{cdp:st:C.2.UN.39} by the
actual root and blowup maps.

The test \eqref{cdp:eq:C.RK.3} is in fact intrinsic to
the original degeneration. Write its original
full fibre and generator orders as
$\sum_\alpha n_\alpha S_\alpha$ and
$\sum_\alpha a_\alpha S_\alpha$.
On the normalization $T_\alpha$ of $S_\alpha$
form the finite normal cover $C_\alpha$
by the following actual local equations
and transition maps:
\[
 \begin{gathered}
 f=g_j^{\,n_\alpha}u_j,\qquad
 w_j^{\,n_\alpha}=u_j|_{T_\alpha},\\
 w_k=(g_j/g_k)|_{T_\alpha}\,w_j .
 \end{gathered}
\tag{C.RK.18}\hypertarget{tag-00ZZ}{}\label{cdp:eq:C.RK.18}
\]
Here $g_j$ is a local defining equation
of the Cartier component $S_\alpha$;
the polynomial covers are normalized
after gluing. Define
\[
 \tau=\sup\left\{c\in\R:
       |f|^{-2c}|\Omega|^2\text{ is locally integrable
                              near }S\right\}.
\tag{C.RK.19}\hypertarget{tag-0100}{}\label{cdp:eq:C.RK.19}
\]
The notation $|\Omega|^2$ denotes its
positive top-degree density, so this
definition requires no metric choice.
Then \eqref{cdp:eq:C.RK.3} is equivalent to
\[
 \frac{a_\alpha+1}{n_\alpha}=\tau
 \quad\text{for every }\alpha\text{ for which a connected
 component of }C_\alpha\text{ has Hopf minimal model}.
\tag{C.RK.20}\hypertarget{tag-0102}{}\label{cdp:eq:C.RK.20}
\]
Consequently its truth is independent
of the embedded resolution, its
semistable triangulation and the
chosen root exponent divisible by
all resolved multiplicities.

For an original compact family the record can
be made on a finite compact cover of its base,
and \eqref{cdp:eq:C.RK.5} recovers the original compact
family. Reverse generation on compact actual
core records with the displayed finite action
and coarsening likewise remains compact.
The original endpoint neighbourhood, not only
its central log space, is retained throughout.
\end{proposition}



\begin{proof}

\begin{enumerate}[label={$\langle1\rangle$\arabic*.},leftmargin=*]
\item \par\medskip\noindent\textbf{Embedded resolution and the exact normalized root charts.}
Use analytic embedded resolution by central
smooth centres, as in \ref{cdp:st:C.2.UN.39} and
\cite[Theorems~1.6 and1.10]{BierstoneMilman1997}.
Divisor centres give identity blowups and are
omitted. The remaining centres in a smooth
threefold are points and compact curves.
The actual pullback of $\Omega$ is the
canonical generator on $Y$: apply the
Hartogs direct-image comparison (C.KG.6)
to the ordinary blowup maps.

\item Near a stratum with $r$ incident components,
absorb the unit in the fibre equation into
a coordinate. The exact local equation and
volume form have the shape
\[
 t=x_1^{m_1}\cdots x_r^{m_r},\qquad
 \pi^*\Omega=u\,x_1^{k_1}\cdots x_r^{k_r}
                         dx_1\wedge dx_2\wedge dx_3,
                         \quad u\text{ a unit}.
\tag{C.RK.9}\hypertarget{tag-010B}{}\label{cdp:eq:C.RK.9}
\]
The normalization of $t=s^N$ has $d_I$
local branches. Each is the affine toric
chart with cocharacter lattice
\[
 \mathcal L_I=\left\{(v,h)\in\Z^r\oplus\Z:
                         \sum_i m_iv_i=Nh\right\}
\tag{C.RK.10}\hypertarget{tag-00ZR}{}\label{cdp:eq:C.RK.10}
\]
and cone generated by the height-one rays
$((N/m_i)e_i,1)$. The omitted $3-r$
coordinates are unchanged. Local roots of
the unit, if not absorbed in \eqref{cdp:eq:C.RK.9},
give the same normalization maps and
the actual finite transition characters.
The group $G$ permutes the $d_I$ branches
and acts torically on each stabilised
branch, over $s\mapsto\zeta s$.

\item The height-one slice has dimension at most
two. Subdivide its edges at every lattice
point, and triangulate each two-dimensional
lattice triangle using all its lattice
points. An empty lattice triangle is
unimodular, by Pick's formula. Thus these
subdivisions give smooth toric charts.
They agree on their subdivided edges.
The refinement can be obtained by
successively inserting the height-one
lattice points as star subdivisions.
Each step is the normalized blowup
of its actual monomial ideal. On an
overlap, replacing a boundary equation
by a holomorphic unit times that
equation multiplies every monomial
generator by a unit and leaves this
ideal unchanged. Thus the normalized
blowup maps glue through the original
tangent-dependent transition units;
those transitions have not been
replaced by constant torus maps.
The stabilizer of an old coarse cone fixes
each old ray: the original SNC components
are individually labelled, and the group
acts over their identity map in $Y$.
Its action on this lattice slice is
therefore the identity. Choose the
triangulations orbitwise. The actual
toroidal transition maps, which preserve
these labelled rays, glue the subdivisions
to a proper $G$-equivariant map $q$.

\item Every new ray still has height one,
so $s$ has order one on every central
component of $Z$. The resulting toric
boundary is simple NC with smooth
components. For the embedded assertion
retain every old component and stratum
label as well as each subdivision ray.
Transition units fix those labels and
the ray vector; distinct rays in one
old cone are never identified. The
local components in a unimodular
chart therefore belong to distinct
global components, and their smooth
star charts glue without a global
self-incidence. Over a double curve the
added surfaces are projective-line
bundles over finite covers of that
curve, with their indicated modifications;
over a triple point they are rational
toric surfaces. None has Hopf minimal
model. The normalized fibre product
over $\Delta_s^*$ is connected because
its compact torus fibres and its base
are connected. The proper Stein factor
has degree one there and hence on the
normal disc. Thus $Z\to\Delta_s$ has
connected full central fibre as well.

\par\medskip\noindent\textbf{Canonical ratios and the explicit minimum test.}
At a primitive height-one ray $(v,1)$,
the toric Jacobian of \eqref{cdp:eq:C.RK.9} gives
\[
 \operatorname{ord}_{(v,1)}(q^*\rho^*\pi^*\Omega)
                 =\sum_i(k_i+1)v_i-1.
\tag{C.RK.11}\hypertarget{tag-00ZS}{}\label{cdp:eq:C.RK.11}
\]
For a strict ray this is
$N(k_i+1)/m_i-1$. Any height-one
point is a convex combination of the
old vertices, so \eqref{cdp:eq:C.RK.11} is the
same affine combination of their
orders. The minimum among all rays is
therefore
\[
                  a_{\min}=N\min_i r_i-1.
\tag{C.RK.12}\hypertarget{tag-00ZT}{}\label{cdp:eq:C.RK.12}
\]
It is an integer, since a minimum is
attained at an old vertex. The form
\[
            \Omega_Z=s^{-a_{\min}}
                             q^*\rho^*\pi^*\Omega
\tag{C.RK.13}\hypertarget{tag-00ZU}{}\label{cdp:eq:C.RK.13}
\]
is holomorphic and generates the
canonical direct image on the reduced
semistable family. Its component
orders have minimum zero. Since all
new components are ruled or rational,
\ref{cdp:st:C.2.UN.75}'s Hopf-minimum condition is
exactly \eqref{cdp:eq:C.RK.3} for the strict covers.
This proves both directions of the
test for the chosen resolution record.
It also retains the actual generator
character
$\gamma^*\Omega_Z=\zeta^{-a_{\min}}\Omega_Z$.

\par\medskip\noindent\textbf{The actual finite birational action and its quotient order.}
Apply \ref{cdp:st:C.2.UN.75} to $Z$. Its comparison
$\phi$ is a bimeromorphic map over
$\Delta_s$, with a specified inverse
record and an actual punctured
isomorphism. Formula \eqref{cdp:eq:C.RK.4}
therefore defines meromorphic maps
on $X_0$. Their products satisfy
the group laws on the dense
punctured set and consequently
as meromorphic maps. They have
order $N$ on the total family,
as seen on the base. Choosing
the endpoint volume by
$\phi^*\Omega_0=\Omega_Z$ also
gives exactly the character in
the preceding paragraph under
this birational action.

\item To reconstruct, first reverse
the retained $\phi$-record.
The action is then the original
regular $G$-action on $Z$.
Apply its equivariant toric
coarsening $q$ to get the actual
normal space $Y_N$. Its quotient
is literally $Y$. Indeed on a
normal local ring $A$ of $Y$,
the root ring is the integral
closure $B$ of
$A[s]/(s^N-f\pi)$. The invariants
of its total ring of fractions are
$\operatorname{Frac}(A)$.
An invariant element of $B$ is
integral over $A$ and lies in
that fraction field, hence lies
in $A$. Thus
\[
                              B^G=A.
\tag{C.RK.14}\hypertarget{tag-00ZV}{}\label{cdp:eq:C.RK.14}
\]
This argument applies to the
normalization on the connected
total space; the local split
branches are permuted by $G$
and give the same invariant
ring equality. The retained
ordinary blowdowns $\pi$ now
recover $X$. These exact maps
prove \eqref{cdp:eq:C.RK.5} and its converse.
An equivariant replacement of
$\phi$ has not been used.

\par\medskip\noindent\textbf{The root correction from actual stratum fibres.}
Over $D_I^\circ$ the reduced
central support of $Y_N$ is an
\emph{actual} finite etale cover
of degree $d_I$. The local root
equation in \eqref{cdp:eq:C.RK.9} proves the
degree; unit-root transition
maps glue the cover and may
permute its sheets. At a point
of this cover the toric fibre
of $q$ has Euler number $\nu_I$.
To see this, use its torus-orbit
stratification. All positive-
dimensional torus orbits have
Euler number zero. The zero-
dimensional orbits are exactly
the maximal unimodular cones
in the subdivision, so their
number is the normalized
volume of the height-one
simplex.

\item Here is its exact lattice
volume. Set $a_i=m_i/d_I$ and
choose an integral vector $b$
with $\sum_i a_i b_i=1$.
The difference lattice of
the slice is
$\ker(a:\Z^r\to\Z)$.
Appending $b$ to a basis of
this kernel gives a unimodular
basis of $\Z^r$. Consequently
the normalized volume is the
absolute determinant obtained
by appending $b$ to the $r-1$
vertex differences. Expansion
gives
\[
 \left|\det\left[
 \frac{N}{m_1}e_1-\frac{N}{m_r}e_r,
 \ldots,
 \frac{N}{m_{r-1}}e_{r-1}-\frac{N}{m_r}e_r,
 b\right]\right|
               =\frac{N^{r-1}d_I}{\prod_i m_i}.
\tag{C.RK.15}\hypertarget{tag-00ZW}{}\label{cdp:eq:C.RK.15}
\]
For $r=1$ both sides are one.
This proves the formula and
integrality of $\nu_I$ without
a choice of triangulation.

\item Proper stratified Euler
pushforward now gives
\[
 \begin{split}
 e(S_Z)&=\sum_I d_I\nu_I e_c(D_I^\circ),\\
 e(S_Y)&=\sum_I e_c(D_I^\circ),\\
 e(S_Z)-e(S_Y)&=R_N.
 \end{split}
\tag{C.RK.16}\hypertarget{tag-00ZX}{}\label{cdp:eq:C.RK.16}
\]
This also checks the finite
quotient directly. A deck
element $\gamma^j$ fixes the
$d_I$-sheeted central cover
over this stratum exactly
when $d_I\mid j$. Therefore
\[
 \begin{split}
 e(S_{Y_N}^{\gamma^j})
       &=\sum_{d_I\mid j}d_I e_c(D_I^\circ),\\
 \frac1N\sum_{j=0}^{N-1}
          e(S_{Y_N}^{\gamma^j})
       &=\sum_I e_c(D_I^\circ)=e(S_Y).
 \end{split}
\tag{C.RK.17}\hypertarget{tag-00ZY}{}\label{cdp:eq:C.RK.17}
\]
All these equalities hold for
the actual covers even when
they are globally connected
and their sheet labels have
nontrivial monodromy.

\item The original resolution has
$e(S_Y)-e(S)=\Delta_\pi$ by
its ordinary point and curve
maps in \ref{cdp:st:C.2.UN.39}. The forward
\ref{cdp:st:C.2.UN.75} record has
$e(S_0)-e(S_Z)=\Delta_K$.
Combine these two equalities
with \eqref{cdp:eq:C.RK.16} to obtain
\eqref{cdp:eq:C.RK.8}, including its zero
criterion. Each summand has
been compared through an
actual proper map or a
retained overlap-compatible
record.

\par\medskip\noindent\textbf{Intrinsic covers and the intrinsic minimum.}
The component $S_\alpha$ is Cartier
in the original smooth threefold.
Its local equations satisfy
$g_k=v_{kj}g_j$ for holomorphic
units, and therefore
$u_k=v_{kj}^{-n_\alpha}u_j$.
This is precisely the transition
in \eqref{cdp:eq:C.RK.18}, including its cocycle
identity. The corresponding monic
finite algebras glue before taking
their finite normalizations.
On the open smooth component
away from other components this
is the $\mu_{n_\alpha}$-torsor
of roots of the normal-line
trivialization supplied by $f$.
No global meromorphic frame for
that normal line is chosen.

\item For any root exponent divisible
by $n_\alpha$, the exact map from
the strict central root cover is
\[
                         w_j=\frac{s^{N/n_\alpha}}{g_j}.
\tag{C.RK.21}\hypertarget{tag-0103}{}\label{cdp:eq:C.RK.21}
\]
At its generic point numerator
and denominator have the same
order, and their quotient satisfies
\eqref{cdp:eq:C.RK.18}. On a smooth strict
component of $Z$ this meromorphic
quotient is integral over its
holomorphic local ring, by that
monic equation. Normality makes
it holomorphic. The map therefore
extends to a proper map to the
glued finite cover, and factors
through its normalization. It
has generic degree one onto a
connected component of $C_\alpha$.
Every such component is obtained,
since on the dense open both
covers have the same complete
set of $n_\alpha$ roots.
Thus the strict surfaces are
actual birational resolutions
of these fixed covers for every
choice of embedded resolution
and divisible root exponent.

\item Every exceptional surface of
the original embedded resolution
is projective: initially it is
$\PP^2$ or a projective-line
bundle over a compact curve,
and its subsequent transforms
are projective modifications.
The ideals on these projective
surfaces are algebraic by GAGA.
A finite normal analytic cover
of a projective surface is also
projective: its finite coherent
algebra and multiplication
algebraize by GAGA, and the
finite pullback of an ample
line is ample. Its surface
resolutions remain bimeromorphic
to projective surfaces. Such
root components cannot have
Hopf minimal model. The new
semistable components have
already been shown ruled or
rational. Hence \eqref{cdp:eq:C.RK.21}
accounts for every Hopf root
component, with no exceptional
ones omitted.

\item Finally apply the literal change
of variables under $\pi$ to the
density in \eqref{cdp:eq:C.RK.19}, outside
the exceptional sets of measure
zero. Near a resolved crossing
it is a bounded positive unit
times the product of the radial
densities
$|x_i|^{2(k_i-cm_i)}\,dA(x_i)$,
with smooth tangent density.
On finitely many relatively
compact charts about the compact
central fibre, the integral is
finite exactly when
$k_i-cm_i>-1$ for every $i$.
At equality it diverges already
at a general point of that
divisor. Thus
\[
                             \tau=\min_i\frac{k_i+1}{m_i}.
\tag{C.RK.22}\hypertarget{tag-0104}{}\label{cdp:eq:C.RK.22}
\]
This is the original density's
integrability threshold, so it
is independent of resolution.
The strict transform of an
original component retains
its original pair
$(a_\alpha,n_\alpha)$.
Together with \eqref{cdp:eq:C.RK.21} and the
exceptional-component comparison,
\eqref{cdp:eq:C.RK.22} proves exactly
\eqref{cdp:eq:C.RK.20} and the stated
choice independence.

\par\medskip\noindent\textbf{Compact scope.}
For a compact base $B$ choose
a meromorphic function with
a simple zero at the selected
point; Riemann--Roch supplies
one with its poles away from
that point. Normalize the
cyclic cover $s^N=h$ of $B$.
It is a compact curve with
the exact local parameter
map used in \eqref{cdp:eq:C.RK.2}.
Additional branch values
are retained in this base
record. Perform the normal
base change and equivariant
resolution globally, and
the \ref{cdp:st:C.2.UN.75} replacement near
the chosen central fibre.
Its actual annular map
glues to the unchanged
outside as in \ref{cdp:st:C.2.UN.75}. Conjugate
the global finite action
through that comparison;
it is a finite birational
action on this compact
family. Reversing the
record regularizes it,
and the compact proper
coarsening and quotient
in \eqref{cdp:eq:C.RK.5} recover the
original compact family.
The converse uses those
same proper maps on an
actual compact core record.
No replacement of a full
neighbourhood follows only
from the central-log-space
classification of \ref{cdp:st:C.2.UN.71}
or its occurrence theorem
\ref{cdp:st:C.2.UN.73}.

\end{enumerate}
\end{proof}






\begin{proposition}[All no-triple fibres are ordinary vertical blowups of the classified elliptic cores]\hypertarget{tag-01EJ}{}\label{cdp:st:C.2.UN.72}
\StatementRecord{cdp:st:C.2.UN.72}{proved}{01EJ}{\SeriesRef{cdp:st:C.2.UN.65}, \SeriesRef{cdp:st:C.2.UN.57}, \SeriesRef{cdp:st:C.2.UN.61}, \SeriesRef{cdp:st:C.2.UN.62}, \SeriesRef{cdp:st:C.2.UN.69}, \SeriesRef{cdp:st:C.2.UN.70}, \SeriesRef{cdp:eq:C.NT.5}, \SeriesRef{cdp:eq:C.NT.6}, \SeriesRef{cdp:st:C.2.UN.64}, \SeriesRef{cdp:st:C.2.UN.5}, \SeriesRef{cdp:st:C.2.UN.7}, \SeriesRef{cdp:st:C.2.UN.39}, \SeriesRef{cdp:st:C.2.UN.66}}

Let $f:U\to\Delta$ be a proper degeneration of connected complex
two-tori from a smooth threefold. Suppose that its full central
fibre $S=\sum_vT_v$ is reduced and simple normal crossing and
has no triple points. The components may be nonminimal.

There is a finite sequence of actual ordinary holomorphic
blowdowns over the disc, each the inverse of a point blowup
or a blowup of a smooth compact curve of positive genus in
the smooth locus of the current central fibre, which ends
in an actual degeneration with one of these marked cores:
\begin{enumerate}
\item a smooth complex two-torus;
\item a marked diagonal Hopf and degree-zero elliptic ruled
cycle with the exact normal maps and lattice product $M=I$
of \ref{cdp:st:C.2.UN.65} and \ref{cdp:st:C.2.UN.57};
\item a rooted elliptic tree satisfying the full criterion
of \ref{cdp:st:C.2.UN.61}, including its original analytic torsors,
secondary framings and zero-Chern inverse principal torsor;
\item a unicyclic support of genus-zero fixed-Jacobian
identity-monodromy elliptic quasi-bundles and nonlinear
Hopf leaves satisfying the full criterion of \ref{cdp:st:C.2.UN.62}.
\end{enumerate}
The original neighbourhood is recovered by reversing the
recorded ordinary blowups on this actual core degeneration.
Each centre avoids every double curve then present.
For every curve centre $B$ the record retains the exact
normal extension and marked exceptional pair
\[
 \begin{gathered}
 0\longrightarrow K_B\longrightarrow V=N_{B/U'}
                 \xrightarrow{\,df'\,}\OO_B\longrightarrow0,\\
 (T,C)=\bigl(\PP_B(V),\PP_B(K_B)\bigr),\qquad
 N_{T/U}\simeq\OO_{\PP_B(V)}(-1).
 \end{gathered}
\tag{C.NT.1}\hypertarget{tag-00TU}{}\label{cdp:eq:C.NT.1}
\]
Here $\PP_B(V)$ parametrizes lines. The actual endpoint
identifications in \eqref{cdp:eq:C.NT.1} are part of the record.

Conversely, take any actual torus degeneration with one of
these cores and perform any finite sequence of ordinary
point blowups and ordinary blowups of smooth compact
positive-genus curves in the smooth locus of its current
central fibre, avoiding all current double curves. The
result is a full reduced SNC torus fibre without triple
points. This gives the complete reconstruction in the
stated class. It operates on actual neighbourhoods; a
different realization of the same abstract marked core
is not substituted for the given neighbourhood.
For a compact total space over a compact curve all the
recorded contractions and blowups have the same compact
version.

If $a$ is the number of point blowups and $B_1,\ldots,B_b$
are the recorded curve centres, then
\[
                 e(S)=2a-2\sum_{j=1}^b\bigl(g(B_j)-1\bigr).
\tag{C.NT.2}\hypertarget{tag-00TV}{}\label{cdp:eq:C.NT.2}
\]
Thus the complete zero-Euler reconstruction is obtained
exactly by the balance
\[
                       a=\sum_{j=1}^b\bigl(g(B_j)-1\bigr).
\tag{C.NT.3}\hypertarget{tag-00TW}{}\label{cdp:eq:C.NT.3}
\]
Equivalently, let $r_0$ count the connected rational
double curves of the original support and let
$\mathcal E_{\ge2}$ be its connected double curves of
genus at least two. These intrinsic data satisfy
\[
 \begin{split}
 a&=r_0,\qquad
 \{B_j:g(B_j)\ge2\}
       \ \simeq\ \{C_e:e\in\mathcal E_{\ge2}\},\\
 e(S)&=2r_0-2\sum_{e\in\mathcal E_{\ge2}}
                                      \bigl(g(C_e)-1\bigr).
 \end{split}
\tag{C.NT.4}\hypertarget{tag-00TX}{}\label{cdp:eq:C.NT.4}
\]
The bijection is the one induced by the actual contractions
on the double curves. Every nonelliptic edge is a terminal
edge at the stage at which it is removed; in particular
every cycle of the original incidence graph consists of
elliptic double curves.
\end{proposition}



\begin{proof}

\begin{enumerate}[label={$\langle1\rangle$\arabic*.},leftmargin=*]
\item \par\medskip\noindent\textbf{First remove every rational double curve by actual point contractions.}
Apply \ref{cdp:st:C.2.UN.69} to the given degeneration. Its strict-height
ascent ends in actual plane-line leaves, and each ordinary
blowdown removes one rational double curve and no other
double curve. A rational ruled-section leaf was excluded
there by transporting its square-minus-two centre into
the actual nearby torus. Thus this first stage makes
exactly $r_0$ point contractions. It leaves a proper
smooth-total torus degeneration, still full reduced
SNC without triple points, and with no rational double
curve. The maps are isomorphisms on neighbourhoods of
every surviving double curve.

\par\medskip\noindent\textbf{An increasing-height path cannot change from higher genus to elliptic.}
Use the actual canonical orders of the current family.
For every remaining double curve the exact adjunction
map is (C.NM.2). In particular a curve of genus at least
two has endpoints of different heights. Orient that
edge toward its higher endpoint and call that component
$P$.

\item Suppose $P$ has a neighbour $Q$ of strictly greater
height, and write $J=P\cap Q$ for the corresponding
connected boundary. If $J$ were elliptic, (C.NM.2)
would give the actual isomorphism
\[
      N_{J/P}^{\,k_Q-k_P}\simeq K_J\simeq\OO_J,
        \qquad a_J=k_Q-k_P-1\ge0 .
\tag{C.NT.5}\hypertarget{tag-00TY}{}\label{cdp:eq:C.NT.5}
\]
Its normal is therefore torsion, with the actual
canonical coefficient required by \ref{cdp:st:C.2.UN.70}. All boundary
curves on $P$ are disjoint. That theorem supplies
the actual map
\[
        \pi:P\longrightarrow Z,\qquad
        q_Z:Z\longrightarrow B,
\tag{C.NT.6}\hypertarget{tag-00TZ}{}\label{cdp:eq:C.NT.6}
\]
which is a sequence of point blowdowns followed by a
relatively minimal elliptic quasi-bundle with only
smooth multiple supports; every original boundary
is vertical for $q_Z\pi$.

\item No smooth boundary of genus at least two can be
vertical for this map. An exceptional prime of the
point-blowdown map is rational. A nonexceptional
boundary has the same function field as its image
under this birational morphism. A vertical image is
one of the smooth elliptic fibre supports of $q_Z$,
so its smooth strict transform has genus one.
Both alternatives contradict the higher-genus
boundary already on $P$. Hence $J$ is not elliptic.
There are no rational double curves at this stage,
so every strictly higher neighbour is reached by
another higher-genus double curve.

\item This uses the actual maps \eqref{cdp:eq:C.NT.5}--\eqref{cdp:eq:C.NT.6}. It does
not contract any arbitrary surface exceptional curve
inside the threefold or assume an elliptic map on an
ambient neighbourhood.

\par\medskip\noindent\textbf{Every higher-genus edge leads to an actual section-leaf contraction.}
Continue along strictly increasing heights whenever
possible. There are finitely many components, so the
path terminates at a local maximum $T$ having a
strictly lower neighbour along a higher-genus curve.
The maximum theorem \ref{cdp:st:C.2.UN.69} applies without any
minimality premise. Its equal-height alternative is
impossible because of that lower neighbour. Its Hopf
leaf alternative has an elliptic boundary, and its
plane-line alternative has a rational boundary.
Thus $T$ is exactly an unblown projective-line
bundle whose single boundary is a section of genus
at least two. The actual \ref{cdp:st:C.2.UN.64} map contracts it
to the smooth curve identified with that section.

\item The parent surface is unchanged by this curve
contraction, since its restriction is the blowup
of a Cartier curve. Every other component misses
the exceptional divisor. Consequently the new
central fibre is still full reduced SNC without
triple points; no rational double curve has appeared,
and every surviving boundary neighbourhood is
unchanged. Canonical order differences on surviving
components are retained by the ordinary blowup
Jacobian, as in \ref{cdp:st:C.2.UN.64}.

\item Repeat this finite procedure. Each step removes
one component and precisely one higher-genus double
curve. It therefore removes all such curves and
gives the exact bijection in the first line of
\eqref{cdp:eq:C.NT.4}. Each curve centre and its normal extension
are exactly \eqref{cdp:eq:C.NT.1}. At this point every remaining
double curve is elliptic.

\par\medskip\noindent\textbf{The remaining support has zero Euler number and is already minimal.}
Denote this remaining full support by $S_{\rm ell}$.
If it is smooth, the smooth-support comparison \ref{cdp:st:C.2.UN.5}
identifies it with a two-torus. Suppose it is
reducible. The actual nearby-cycle Euler formula
for a reduced SNC fibre is
\[
        e(F)=\sum_v e(T_v^\circ),\qquad
        T_v^\circ=T_v\setminus
                          \bigcup_{w\ne v}(T_v\cap T_w).
\tag{C.NT.7}\hypertarget{tag-00U0}{}\label{cdp:eq:C.NT.7}
\]
This is the proper analytic formula of A'Campo used
in \ref{cdp:st:C.2.UN.7} and \ref{cdp:st:C.2.UN.39}. Its local comparison is especially
simple here: over a smooth component stratum there
is one nearby sheet, while a double stratum has
the annular local fibre $xy=t$, of Euler number zero.
There are no triple strata. Compact triangulations
and the actual tubular comparison give the displayed
additivity, without a Kähler assumption.

\item Writing $C_e$ for the double curves, all of which
are elliptic, gives
\[
 \begin{split}
 0=e(F)&=\sum_v e(T_v)-2\sum_e e(C_e)
                            =\sum_v e(T_v),\\
 e(S_{\rm ell})&=\sum_v e(T_v)-\sum_e e(C_e)=0.
 \end{split}
\tag{C.NT.8}\hypertarget{tag-00U1}{}\label{cdp:eq:C.NT.8}
\]
Now use the all-elliptic extension (C.MR.16), whose
proof applies to these actual possibly nonminimal
components. Its nonnegative minimal-model Euler
terms and positive point-blowup terms force every
component of $S_{\rm ell}$ to be minimal and
nonrational. Thus the actual support has reached
the precise scope of \ref{cdp:st:C.2.UN.66}, rather than merely a
support with a prescribed numerical Euler value.

\ref{cdp:st:C.2.UN.66} now makes its further ordinary contractions.
All their curve centres are elliptic: the boundary
curves present are elliptic and no rational or
higher-genus edge is introduced. The remaining
actual support is one of the four cores in the
statement. Its full torsor, lattice, secondary and
normal identifications are those of \ref{cdp:st:C.2.UN.61}, \ref{cdp:st:C.2.UN.62}
or \ref{cdp:st:C.2.UN.65}, as appropriate. This establishes the
endpoint before any reconstruction claim.

\par\medskip\noindent\textbf{Exact reversal and the converse construction.}
Every preceding ambient map is one of the ordinary
holomorphic blowdowns of \ref{cdp:st:C.2.UN.64}. Reverse the finite
record on the actual endpoint family. A point
step replaces a smooth point of one central
component by its exceptional plane. A curve
step has the exact normal sequence \eqref{cdp:eq:C.NT.1}.
These comparisons are the actual local blowup
charts and glue with the original isomorphism
off each exceptional divisor. They recover the
given $U$ over the disc, not just the birational
classes of its components.

\item For clarity, the same charts prove the converse
directly. At a curve centre choose coordinates
$(z,w,t)$ with the centre $(z,t)=0$ and $f=t$.
The blowup charts $t=zs$ and $z=tr$ have fibre
equations $zs=0$ and $t=0$. At a point centre
blow up the ideal $(x,y,t)$; the chart
$y=xu,\ t=xv$ has fibre equation $xv=0$,
and the other two charts give the same ordinary
normal-crossing comparison. The centre avoids
all pre-existing double curves, so no third
component meets either exceptional pair.
The full exceptional multiplicity is one, since
$t$ has order one along the centre.
Thus each step retains a smooth total space
and a full reduced SNC fibre without triple points.
Its punctured family is unchanged. Properness
is retained by the proper blowup map.

\item On a curve centre $B$ in the smooth central
locus the actual canonical divisor is supported
on boundaries away from $B$. Its restriction
trivializes $K_{P'}|_B$. Adjunction therefore gives
$N_{B/P'}\simeq K_B$, and the normal sequence
with the quotient $df'$ is precisely \eqref{cdp:eq:C.NT.1}.
No curve of genus zero can be such a centre:
the compact submersion transport (C.NM.10)
would again put its square-minus-two sphere
into the nearby torus. Points and smooth curves
of positive genus are exactly the centre
types used in the reconstruction.

\item When the original total space is compact, the
actual local blowdowns glue to the identity
off their exceptional divisors by \ref{cdp:st:C.2.UN.64}.
Their targets remain compact and the map to
the compact base descends. Conversely every
listed compact centre has its ordinary global
blowup. This proves the compact version without
assuming an isomorphism between unrelated
realizations of a marked core.

\par\medskip\noindent\textbf{Euler balance and the intrinsic double-curve record.}
All four cores have Euler number zero. A point
blowup replaces one point by $\PP^2$ and adds two
to the Euler number of the support. A curve
blowup replaces $B$ by its projective-line bundle
of Euler number $2e(B)$, and so adds
$e(B)=2-2g(B)$. Additivity along the actual
blowup maps gives \eqref{cdp:eq:C.NT.2}, and its zero
condition is exactly \eqref{cdp:eq:C.NT.3}.

\item Every such blowup attaches a single new double
edge: a rational line for a point, or the actual
curve $B$ for a curve centre. It leaves all
pre-existing double curves and their
neighbourhoods unchanged. The cores have only
elliptic double curves. The rational-edge
removal stage made exactly $r_0$ point steps,
and the higher-genus removal stage identifies
each of its centres with exactly the removed
original higher-genus double curve. These
actual maps prove all of \eqref{cdp:eq:C.NT.4} and its
independence of the chosen contraction order.
They also show that every nonelliptic edge
is removed as a leaf, which proves the
assertion about cycles. This completes the
classification in the full reduced
no-triple class.

\end{enumerate}
\end{proof}






\begin{proposition}[Every reduced SNC torus degeneration has a reversible finite reduction to a classified core]\hypertarget{tag-01EX}{}\label{cdp:st:C.2.UN.85}
\StatementRecord{cdp:st:C.2.UN.85}{proved}{01EX}{\SeriesRef{cdp:st:C.2.UN.65}, \SeriesRef{cdp:st:C.2.UN.68}, \SeriesRef{cdp:st:C.2.UN.71}, \SeriesRef{cdp:st:C.2.UN.61}, \SeriesRef{cdp:st:C.2.UN.62}, \SeriesRef{cdp:st:C.2.UN.75}, \SeriesRef{cdp:st:C.2.UN.84}, \SeriesRef{cdp:st:C.2.UN.57}, \SeriesRef{cdp:st:C.2.UN.73}, \SeriesRef{cdp:st:C.2.UN.82}, \SeriesRef{cdp:st:C.2.UN.77}, \SeriesRef{cdp:st:C.2.UN.81}, \SeriesRef{cdp:st:C.2.UN.83}, \SeriesRef{cdp:eq:C.GP.5}, \SeriesRef{cdp:st:C.2.UN.80}, \SeriesRef{cdp:eq:C.GP.4}, \SeriesRef{cdp:eq:C.GP.8}, \SeriesRef{cdp:st:C.2.UN.74}, \SeriesRef{cdp:st:C.2.UN.50}, \SeriesRef{cdp:st:C.2.UN.28}, \SeriesRef{cdp:st:C.2.UN.47}}

Let \(f:X\to\Delta\) be a proper degeneration of connected
complex two-tori, with smooth total space and full reduced
SNC central fibre. There is a finite overlap-compatible
modification record over the same disc, isomorphic over
\(\Delta^*\), ending in one of the following actual
smooth-total reduced degeneration germs.

\begin{enumerate}
\item A canonical-flat core. Its central fibre is a smooth
two-torus, one of the exact no-triple elliptic-ruled/diagonal
Hopf cycles of \ref{cdp:st:C.2.UN.65}, or one of the marked positive toric
log spaces of \ref{cdp:st:C.2.UN.68} and \ref{cdp:st:C.2.UN.71}. The toric alternative has
Euler number equal to its positive triple count.

\item An embedded no-triple core whose components are
minimal fixed-Jacobian elliptic quasi-bundles with identity
integral monodromy, together with any minimal nonlinear
Hopf leaves. All boundary curves are smooth elliptic fibre
supports or the specified Hopf axes. Every elliptic
quasi-bundle base has genus zero or one, and at most one
has genus one. The exact support, including its original
torsors, normal characters, translations and gluing maps,
is one of the rooted-tree or unicyclic supports classified
in \ref{cdp:st:C.2.UN.61} and \ref{cdp:st:C.2.UN.62}. Its Euler number is zero.
\end{enumerate}

If every Hopf-birational prime has minimum canonical order,
the first alternative is supplied by \ref{cdp:st:C.2.UN.75}. If some such
prime has positive order, the second alternative can be
obtained while retaining a fixed maximal-order minimal
Hopf component and its entire neighbourhood after the
preliminary reduction of \ref{cdp:st:C.2.UN.84}. No nonlinear Hopf
neighbourhood is contracted or replaced during this
reduction.

The record retains all actual centres, two-ruling maps,
point and curve contractions, local sigma charts,
normalization branches and restored overlap maps.
Its inverse reconstructs the original analytic family
from this actual endpoint germ. On an originally compact
family all operations take place on that compact base.
If \(S_0\) is the original reduced support, \(S_*\) the
endpoint support, and \(\delta_j\) the signed normal-cone
Euler change of its \(j\)-th retained step as in \ref{cdp:st:C.2.UN.75},
then
\[
                 e(S_0)=e(S_*)-\sum_j\delta_j.
 \tag{C.GP.1}\hypertarget{tag-00KY}{}\label{cdp:eq:C.GP.1}
\]

The classified core central spaces have the compact
occurrence constructions of \ref{cdp:st:C.2.UN.57}, \ref{cdp:st:C.2.UN.61}, \ref{cdp:st:C.2.UN.62} and \ref{cdp:st:C.2.UN.73}.
The inverse record here is applied to its actual endpoint
neighbourhood. Identical central log data in a different
smoothing are not asserted to identify its higher analytic
jets or to carry that particular inverse record.
\end{proposition}



\begin{proof}\leavevmode\par



 
\begin{enumerate}[label={$\langle1\rangle$\arabic*.},leftmargin=0pt,labelsep=.5em]
\item \par\medskip\noindent\textbf{The initial dichotomy and the protected component.}
Normalize the primitive canonical evaluation as
\[
       \operatorname{div}(\Omega)=\sum_i k_iV_i,
       \qquad \min_i k_i=0.
 \tag{C.GP.2}\hypertarget{tag-00L0}{}\label{cdp:eq:C.GP.2}
\]
If every Hopf-birational prime has order zero, \ref{cdp:st:C.2.UN.75}
gives its actual canonical-flat endpoint. \ref{cdp:st:C.2.UN.65} and
\ref{cdp:st:C.2.UN.68}--\ref{cdp:st:C.2.UN.71} give exactly the alternatives in (1),
including the actual marked log maps in the positive
case. This uses the proved component restriction that
every class-VII minimal model is Hopf. There is no
unexamined additional class-VII case.


\begin{enumerate}[label={$\langle2\rangle$\arabic*.},leftmargin=1.3em,labelsep=.5em]
\item  Suppose instead that a Hopf-birational prime has
positive order. Apply the finite truncated record
\ref{cdp:st:C.2.UN.84}, including its maximal Hopf minimalization. At
its smooth-total NC endpoint put
\[
                  M=\max_i k_i>0.
 \tag{C.GP.3}\hypertarget{tag-00L1}{}\label{cdp:eq:C.GP.3}
\]
The minimum is still zero and all orders are at
most this fixed \(M\). Choose one of the surviving
minimal Hopf components \(H\) of order \(M\).
\ref{cdp:st:C.2.UN.82} shows that it has exactly one elliptic
boundary curve \(J\), with torsion normal, and
that its neighbour has smaller order. The marked
neighbourhood is embedded. No elliptic self-double
curve has been introduced in the preliminary
record: its new incidences are on rational
exceptional curves, as proved in \ref{cdp:st:C.2.UN.84}.


\item  All later whole-component removal blocks will be
the positive-order blocks of \ref{cdp:st:C.2.UN.77} and \ref{cdp:st:C.2.UN.81}. Such
a ruled component of order \(q\le M\) has one
horizontal boundary section, whose neighbour has
order \(q-1\), and every other boundary is rational
and vertical. It cannot meet \(H\): a rational
curve cannot lie on a minimal Hopf, and the
section case would require \(q=M+1\). A plane
contraction also has only rational boundary.
The same argument applies to the cleanup blocks
and their missing-fibre insertions. Direct moves
below will likewise avoid \(H\). Thus its entire
neighbourhood is protected and provides the same
starting interface after every completed block.


\end{enumerate}

\item \par\medskip\noindent\textbf{The recorded models, not a numerical eligibility class.}
At a smooth-total NC stage, call a component
\emph{recognized} only when one of the actual
comparisons of \ref{cdp:st:C.2.UN.83} has been recorded for its
smooth normalization \(T\), together with its
entrance elliptic curve and its point map
\[
                         \pi_T:T\longrightarrow Z_T.
 \tag{C.GP.4}\hypertarget{tag-00L2}{}\label{cdp:eq:C.GP.4}
\]
There are three recorded types: an elliptic
quasi-bundle model; a nonlinear Hopf model with
its one axis; or a one-section elliptic ruled
model with its actual rank-two bundle extension.
The point map is an isomorphism near the entrance
curve. All boundary branches, including both
branches of a self-double, are retained.


\begin{enumerate}[label={$\langle2\rangle$\arabic*.},leftmargin=1.3em,labelsep=.5em]
\item  Let \(d_T\) be the exact number of point blowups
in this map. The following properties are part
of those actual comparisons:
\[
\begin{gathered}
 d_T\ge0,\qquad
 \operatorname{Exc}(\pi_T)\subset D_T,\\
 \text{every rational boundary is exceptional
 and has negative square},\\
 \text{every elliptic boundary has nonpositive square}.
\end{gathered}
 \tag{C.GP.5}\hypertarget{tag-00L3}{}\label{cdp:eq:C.GP.5}
\]
For the elliptic type, the boundary is vertical
for the composed elliptic map; \eqref{cdp:eq:C.GP.5} follows
from \ref{cdp:st:C.2.UN.80} and the fibre squares. For the other
two types the entrance curve has square zero,
all point centres are away from it, and every
other boundary is exactly a positive discrepancy
exceptional curve, by (C.TI.2)--(C.TI.3).


\item  Initially recognize \(H\) alone. \ref{cdp:st:C.2.UN.83} applied
to its actual isolated torsion-normal axis gives
the elliptic or nonlinear Hopf type, and its
point defect is zero. Let
\[
\begin{gathered}
 u=\#\{\text{current unrecognized global component primes}\},
 \\
 d=\sum_{T\ \mathrm{recognized}}d_T.
\end{gathered}
 \tag{C.GP.6}\hypertarget{tag-00L4}{}\label{cdp:eq:C.GP.6}
\]
Recognition changes no complex space: it records
the actual maps of \ref{cdp:st:C.2.UN.83} and their marked curves.
In particular the word does not assert that an
unrecognized receiving surface can never acquire
such a comparison at a later step.


\end{enumerate}

\item \par\medskip\noindent\textbf{A direct move on a recorded elliptic model.}
Suppose that a recognized elliptic-type component
has positive defect. Choose a last exceptional
curve of \eqref{cdp:eq:C.GP.4}. It is an actual boundary by
\ref{cdp:st:C.2.UN.80}. The complete elliptic fibre graph and the
exact receiving-degree comparison (C.AP.5)--
(C.AP.6) apply without any additional sign
assumption on the entrance coefficient.


\begin{enumerate}[label={$\langle2\rangle$\arabic*.},leftmargin=1.3em,labelsep=.5em]
\item  Here is the exact comparison with an embedded
curve in the total space at an NC stage. On a
complete fibre graph of the point map, write
\(a_C\) for the coefficient of each component
in the actual canonical divisor on the
normalization, assigning coefficient zero
also to a component absent from the actual
boundary. Adjacent shifted coefficients
\(a_C+1\) are relatively prime. There is no
initial edge on the smooth elliptic fibre;
a blowup at a smooth point creates consecutive
shifted coefficients, and a node blowup inserts
the sum of the two old shifted coefficients.
Thus the assertion follows through the actual
point map, including negative initial coefficients.
In particular no edge has coefficients
\((-1,-1)\).


\item  A self-double branch has coefficient \(-1\).
If the image of the chosen last \((-1)\)-curve
were nodal at an actual triple point with two
local branches of its source component, its
two neighbours on the normalization would be
the two self-double branches. The last-node
coefficient formula would also give
\(a_E=(-1)+(-1)+1=-1\), contradicting that
edge invariant. A triple with three branches
of the source gives adjacent \(-1\) branches
directly. With only one local source branch,
two image branches would already make a node
on the chosen curve in that smooth source
normalization, contrary to its being the
smooth last exceptional \(\PP^1\).
The actual image curve is consequently
smooth. For a self-double its normalized
degree-two cover is unramified, by the local
NC coordinates. A connected such cover by
\(\PP^1\) is impossible by Riemann--Hurwitz;
the cover therefore splits into two distinct
normalized branches. Over the resulting
simply connected image \(\PP^1\), those
labels give actual two smooth surface germs.
When the triple count is two, the normal-product
map gives the two normal degrees \(-1,-1\).
The actual type-II neighbourhood and its two
ruling maps therefore apply even if these
two surface germs belong to one global prime.
For the positive-discrepancy models below,
all rational coefficients are positive, so
the same self-double obstruction is immediate.


\item  With two triple points, the actual type-II move
blows down this curve on its two adjacent
normalized surfaces. Each opposite receiving
boundary has square at least zero, or at least
one when it is elliptic. By \eqref{cdp:eq:C.GP.5} neither
receiver was recognized before this move. Thus
the move adds no point defect to a recognized
surface.


\item  If its second source surface was recognized,
the contracted rational boundary was exceptional
for its recorded point map, again by \eqref{cdp:eq:C.GP.5}.
That map factors through the actual blowdown:
it is constant on the contracted curve, and
the proper surface contraction and its direct
image of the structure sheaf give this descent.
The resulting birational surface morphism is
the remaining sequence of point blowups. Its
defect decreases by one. The same factorization
applies to the first source. All recorded
entrance curves are disjoint from these
exceptional curves and are unchanged. Hence
the recognized set and its models persist,
\(u\) is unchanged, and \(d\) strictly decreases.
If both source germs belong to one global
recognized component, their two normalized
curves are disjoint. The same point map
factors through both ordinary contractions;
its defect drops by two while the global
component is counted only once.


\item  If a last curve has one triple point, use the
actual positive ruled-removal block of (C.AP.7).
If it has none, use the exact plane progress
block (C.AP.8)--(C.AP.13). Each gives a strict
component-card decrease after its finite cleanup.
No nonboundary last-point case is admitted: its
exclusion is the actual field statement \ref{cdp:st:C.2.UN.80},
not an inference from nonzero canonical coefficients.


\end{enumerate}

\item \par\medskip\noindent\textbf{A direct move on either positive-discrepancy model.}
Now let the recognized model be nonlinear Hopf
or one-section ruled, with positive defect.
The actual canonical divisor is
\[
 K_T=-aJ_T+\sum_E b_EE,
 \qquad b_E>0,
 \qquad a\ge2,
 \tag{C.GP.7}\hypertarget{tag-00L5}{}\label{cdp:eq:C.GP.7}
\]
with \(a=2\) for the ruled type. Every point
centre is away from the preserved \(J_T\),
and the full other boundary is exactly the
exceptional set. The graph of each exceptional
tree starts with a rational curve of square
\(-1\) and valence zero. If \(C^2=-b\) and
\(n\) is its valence in this entire exceptional
set, the actual smooth and node blowups give
\[
                              n\le b+1.
 \tag{C.GP.8}\hypertarget{tag-00L6}{}\label{cdp:eq:C.GP.8}
\]
At a smooth point both \(n\) and \(b\) increase
by one on an old component; at a node \(b\)
increases and its old valence is unchanged.
A newly inserted curve has \((b,n)=(1,0)\),
\((1,1)\), or \((1,2)\), according as its
centre is off the exceptional support, on
one component, or at a node. These are
exactly the updates proving \eqref{cdp:eq:C.GP.8}.
At the last node the opposite
square is \(b+1-t\ge0\), since the actual
triple count \(t\) is at most the old valence.
All opposite curves here are rational.
Thus no receiver is already recognized by
\eqref{cdp:eq:C.GP.5}, and the exact point-map factorization
in the preceding paragraph again strictly
decreases \(d\). The positive discrepancies
also give the precise height comparison:
if the two old coefficients are \(b_1,b_2>0\)
and the current component has order \(p\),
the receiving orders are
\[
                 p+b_1+1,\qquad p+b_2+1.
 \tag{C.GP.9}\hypertarget{tag-00L7}{}\label{cdp:eq:C.GP.9}
\]


\begin{enumerate}[label={$\langle2\rangle$\arabic*.},leftmargin=1.3em,labelsep=.5em]
\item  A last curve with no triple point has a
rational \(+1\) receiver of order \(p+2\),
and (C.AP.8)--(C.AP.13) supply the finite
plane progress block. With one triple point
its receiver is ruled with one section and
has positive order, so the same ruled block
applies. These changes concern only the
actual rational exceptional configurations,
not the nonlinear Hopf neighbourhood itself.


\item  When a one-section ruled model reaches defect
zero, it is the actual bundle \(\PP(V)\)
with sole boundary section \(J_T\), including
its possibly nonzero extension class. \ref{cdp:st:C.2.UN.74}
gives the ordinary curve contraction along
its ruling. Its order is
\(k_T=k_{\mathrm{neighbour}}+1>0\), by \ref{cdp:st:C.2.UN.83}.
This is another whole-card decrease.


\end{enumerate}

\item \par\medskip\noindent\textbf{Recognition across the next actual interface.}
We always eliminate positive recorded defects
before recognizing another component. We also
remove a defect-zero one-section ruled model
by the preceding actual contraction. If neither
action is available, every recognized component
is a minimal elliptic quasi-bundle or a minimal
nonlinear Hopf leaf.


\begin{enumerate}[label={$\langle2\rangle$\arabic*.},leftmargin=1.3em,labelsep=.5em]
\item  Such a component is embedded and has disjoint
smooth elliptic boundary curves. Indeed no
rational boundary can remain when its point
defect is zero. The original fibre was SNC,
and the retained type-II and one-section
sigma records create new self-incidences only
on rational fibre or exceptional curves.
Their generic one-section contraction cannot
identify two elliptic horizontal branches.
Thus there is no new elliptic self-double
curve on one of these minimal components.


\item  If an unrecognized component remains, the
connected incidence graph of the central
support has an edge from the recognized
region to its complement. Starting with
\(H\) and recording each new entrance edge
makes that region connected; the direct
moves have retained these entrance curves.
The edge is therefore a smooth elliptic
curve, isolated on its recognized side.
There is no triple point on it, so its
other normalization branch is isolated
as well. Its normal is torsion: on a
quasi-bundle it is the exact normal root
of a smooth multiple fibre, and on a
nonlinear Hopf it is its axis character.
The actual reduced-fibre normal-product
map gives the inverse torsion line on
the other side.


\item  Apply \ref{cdp:st:C.2.UN.83} to that actual other normalization.
It supplies one of the three point models
\eqref{cdp:eq:C.GP.4}, with the entrance neighbourhood
unchanged and every original boundary
retained. Record it as recognized. Then
\(u\) decreases by one. The newly recorded
defect may be any finite nonnegative integer;
we make no assertion that this recognition
alone decreases \(d\).


\end{enumerate}

\item \par\medskip\noindent\textbf{Whole blocks may reset the recognized region.}
A whole-component removal and its sigma
cleanup need not preserve the point models
on the previously recognized region. For
example a missing-fibre insertion can create
additional point blowups on a section parent.
We allow all such changes explicitly.
After this completed block discard every
recognition except the unchanged maximal
Hopf \(H\), and restart with its same
torsion-normal interface. This is possible
because its entire neighbourhood is protected
as proved above. There is no assertion that
the rest of the processed region was protected.


\begin{enumerate}[label={$\langle2\rangle$\arabic*.},leftmargin=1.3em,labelsep=.5em]
\item  Let \(c_j\), for \(1\le j\le M\), count
the retained global components of order \(j\).
The full fixed potential is
\[
                 (c_1,c_2,\ldots,c_M,u,d).
 \tag{C.GP.10}\hypertarget{tag-00KZ}{}\label{cdp:eq:C.GP.10}
\]
Every direct type-II move preserves the card
vector and \(u\), and strictly decreases \(d\).
Recognition preserves the cards and strictly
decreases \(u\), regardless of its finite
new defect. Every whole-component block
strictly decreases its first changed card,
regardless of the finite reset of \(u,d\).
The actual missing-fibre and cleanup rules
add only finitely many cards at the next
order and add none at \(M\). Thus all
retained orders stay at most the fixed \(M\).


\item  Lexicographic order on this fixed finite
product of nonnegative integers is well
founded. The internal plane progress and
sigma cleanup blocks are finite by \ref{cdp:st:C.2.UN.81}
and \ref{cdp:st:C.2.UN.77}. Consequently this procedure
cannot continue indefinitely.


\item  At its endpoint no recorded defect is
positive, no one-section ruled model
remains, and no component is unrecognized.
Otherwise the preceding paragraphs give
another strict decrease. All components
are therefore the minimal elliptic
quasi-bundles and nonlinear Hopf leaves
in (2), and every boundary is smooth
elliptic and isolated. In particular the
endpoint is embedded SNC and has no
triple point.


\end{enumerate}

\item \par\medskip\noindent\textbf{Identification of the endpoint and the retained analytic scope.}
Every endpoint component and double curve
has topological Euler number zero, so
\(e(S_*)=0\). The exact invariant relative
cohomology frame gives
\(p_g(Z)=g(B)\) for a quasi-bundle. The
common regular Jacobian is retained by
the actual normal-root maps: along an
edge joining two quasi-bundle vertices,
the dual normal lines have the same
exact order, and inversion identifies
their connected root covers of the
elliptic double curve. These are exactly
the two regular elliptic fibres, by
(C.QT.2)--(C.QT.3); their translations
are retained in the gluing record.
At a nonlinear Hopf leaf the same root
cover is instead its exact period cover
of \ref{cdp:st:C.2.UN.50}. Such leaves do not interrupt
propagation through the connected
quasi-bundle region, and their actual
leaf periods and finite characters are
retained by that comparison.
The
compact-core bound \ref{cdp:st:C.2.UN.28} therefore permits
at most one genus-one base. If there is
one, that normalization has \(p_g=1\).
\ref{cdp:st:C.2.UN.47} gives the actual surjection from its
fundamental group onto that of the full
support. The graph-of-spaces comparison
gives a surjection from the support
fundamental group to that of its incidence
graph, and kills the image of this
component's fundamental group.
The graph fundamental group is therefore
trivial, so the connected graph is a tree.
The exact coherent restriction and
fundamental-group comparisons of \ref{cdp:st:C.2.UN.61}
then give its rooted-tree support, with the
actual integral Chern and recursive
period conditions. If there is none,
\ref{cdp:st:C.2.UN.62} gives the unicyclic support and
its exact torsor, height and common-period
conditions. Its nonlinear leaves retain
the secondary finite characters of \ref{cdp:st:C.2.UN.50}.
These applications keep the original
surfaces and gluing maps; they do not
replace the actual endpoint smoothing.


\begin{enumerate}[label={$\langle2\rangle$\arabic*.},leftmargin=1.3em,labelsep=.5em]
\item  All prior operations are the actual
finite diagrams already used in \ref{cdp:st:C.2.UN.75},
\ref{cdp:st:C.2.UN.77}, \ref{cdp:st:C.2.UN.81} and \ref{cdp:st:C.2.UN.84}. Their proper graphs
glue by the retained overlap maps, so
their inverse recovers this given
analytic family. Equation \eqref{cdp:eq:C.GP.1}
is the telescoping normal-cone formula
on that finite record, including its
restored-overlap subdivision. Annular
gluing to an unchanged compact outside
gives the same actual modifications
when the starting total space is compact.


\item  The separate compact constructions for
the classified core support data prove
their occurrence. They do not identify
arbitrary higher transverse jets of
this actual endpoint with another core
smoothing. No inverse sigma record is
transferred through an unproved such
identification. This is the claimed
finite reduction and exact reconstruction
of every given reduced SNC torus
degeneration.



\end{enumerate}

\end{enumerate}
\end{proof}






The data in the following construction include analytic
neighbourhoods and their transition maps. This is essential
when an inverse operation uses a normal extension or higher
transverse terms. We first specify those data without
referring to a family to be reconstructed.

\paragraph{Finite admissible data.}
The following are the inputs and the checks on them.
\begin{enumerate}
\item An actual proper smooth-total core family
\(f_*:X_*\to\Delta_s\), with its labelled analytic
charts and overlaps and its punctured torus-family
comparison, is given. Its full reduced central space
is one of the cores in \ref{cdp:st:C.2.UN.85}: a smooth torus; a
canonical-flat marked Hopf/degree-zero-ruled cycle;
a positive periodic toric log space; or an embedded
rooted-tree or unicyclic elliptic quasi-bundle/Hopf
support of \ref{cdp:st:C.2.UN.61}--\ref{cdp:st:C.2.UN.62}. The core chart functions,
including their transverse terms, are retained.

\item A finite inverse overlap-compatible sigma record
on that actual core is given. Each arrow includes
the local blowup or contraction map, its centre
and ideal, the target charts, and the isomorphism
on the restored overlap. The maps and their inverse
comparisons are checked on these charts as in \ref{cdp:st:C.2.UN.75}.
The resulting space \(Z\) is smooth and proper over
\(\Delta_s\), with full reduced SNC central divisor.
The record is an isomorphism on the punctured family.

\item A finite cyclic birational action on the core,
over \(s\mapsto\zeta s\), is given, where \(\zeta\)
has order \(N\). Conjugating it by the specified
inverse record gives holomorphic maps on every
chart of \(Z\), including the central divisor.
Their compositions satisfy the cyclic group law.
Thus it is an actual regular action of
\(G=\mu_N\) on \(Z\), not only a central action
or an action on its meromorphic function field.

\item A labelled equivariant toric coarsening
\(q:Z\to Y_N\) is given. Its local semigroup
charts, units and maps are explicit normalized
monomial-ideal blowups in the reverse direction.
The finite quotient charts of \(Y_N\) have
invariant ring \(\OO_Y\), with smooth coordinates
in which the descended function is
\[
                  t=u\prod_{i=1}^r x_i^{m_i},
              \qquad r\le3,\quad u\in\OO_Y^*,
              \quad m_i\mid N.
 \tag{C.AR.1}\hypertarget{tag-00A0}{}\label{cdp:eq:C.AR.1}
\]
The normal charts upstairs are exactly the
normalization of \(s^N=t\) in these coordinates.
The labelled height-one subdivisions defining
\(q\) are unimodular, and their overlap maps
and finite action are retained. These ring and
coordinate conditions verify the smooth quotient
\(Y=Y_N/G\) and its full SNC divisor with
the displayed multiplicities.

\item A finite sequence \(\pi:Y\to X\) of
ordinary inverse point or smooth-curve blowups
is given, with their actual smooth target charts
and gluing maps. Their centres are in the central
fibre. The function \(t\) descends through each
map, as is checked by its displayed pullback.
The final total space \(X\) is smooth. Its
central support is allowed arbitrary singularities
and arbitrary full multiplicities.
\end{enumerate}
All quotient, contraction and overlap requirements here
are checks on the supplied finite maps. They do not
refer to the existence of an additional original family.
An ordinary arrow can equivalently be supplied by its
Rees-algebra charts and the corresponding proper map.
For compact data, the core family and base, the cyclic
base cover and finite action, and all maps outside these
central charts are supplied globally; the local overlaps
agree with them on an annulus. This is the compact
version of the same definition.

\begin{proposition}[Complete rooted generation and the zero-Euler criterion for arbitrary fibres]\hypertarget{tag-01EY}{}\label{cdp:st:C.2.UN.86}
\StatementRecord{cdp:st:C.2.UN.86}{proved}{01EY}{\SeriesRef{cdp:eq:C.AR.1}, \SeriesRef{cdp:st:C.2.UN.84}, \SeriesRef{cdp:st:C.2.UN.85}, \SeriesRef{cdp:st:C.2.UN.78}, \SeriesRef{cdp:eq:C.AR.10}, \SeriesRef{cdp:st:C.2.UN.75}, \SeriesRef{cdp:st:C.2.UN.57}, \SeriesRef{cdp:st:C.2.UN.61}, \SeriesRef{cdp:st:C.2.UN.63}, \SeriesRef{cdp:st:C.2.UN.71}, \SeriesRef{cdp:st:C.2.UN.73}; \cite{BierstoneMilman1997}}

Every proper degeneration of connected complex two-tori
from a smooth threefold, with arbitrary full central
fibre, is the output of the finite admissible construction
above. Conversely every such input produces a proper
degeneration of connected complex two-tori from a smooth
threefold. The inverse record recovers the supplied
family exactly. For compact threefolds over compact
curves, the same assertions hold for the compact data.

Here is an explicit formula selecting precisely the
zero-Euler outputs. Let \(S_*=f_*^{-1}(0)\), let
\(D_i\) be the quotient components in \eqref{cdp:eq:C.AR.1}, and put
\[
 \begin{split}
 D_I^\circ&=\bigcap_{i\in I}D_i
                    \setminus\bigcup_{j\notin I}D_j,\\
 d_I&=\gcd_{i\in I}m_i,\qquad
 \nu_I=\frac{N^{|I|-1}d_I}{\prod_{i\in I}m_i},\\
 R_N&=\sum_I(d_I\nu_I-1)e_c(D_I^\circ).
 \end{split}
 \tag{C.AR.2}\hypertarget{tag-00A8}{}\label{cdp:eq:C.AR.2}
\]
For the ordinary resolution record \(\pi\), define
\[
 \Delta_\pi=2p_\pi+\sum_\alpha(2-2g(C_\alpha)).
 \tag{C.AR.3}\hypertarget{tag-00A9}{}\label{cdp:eq:C.AR.3}
\]
Let \(\Delta_\sigma\) be the signed normal-cone Euler
change in the forward direction \(Z\dashrightarrow X_*\)
of the retained sigma record. Then its output satisfies
\[
 e\bigl(f^{-1}(0)_{\mathrm{red}}\bigr)
       =e(S_*)-\Delta_\sigma-R_N-\Delta_\pi.
 \tag{C.AR.4}\hypertarget{tag-00AA}{}\label{cdp:eq:C.AR.4}
\]
Every zero-Euler fibre is therefore obtained, and only
such fibres are obtained, by the admissible data with
\[
 \Delta_\sigma+R_N+\Delta_\pi=
 \begin{cases}
  2[N_0:\Gamma_0],&\text{for the positive toric core},\\
  0,&\text{for each other core}.
 \end{cases}
 \tag{C.AR.5}\hypertarget{tag-00AB}{}\label{cdp:eq:C.AR.5}
\]
Here \(N_0\) and \(\Gamma_0\) are the two-dimensional
affine lattices of the toric core; the root exponent
\(N\) in \eqref{cdp:eq:C.AR.1} is a separate integer. The construction
in \eqref{cdp:eq:C.AR.5} specifies spaces, bundles, finite maps,
centres and analytic gluings, as well as their Euler
numbers. Reversing a retained record recovers every
specified normal extension and every transverse jet.

The choice between a canonical-flat core and a positive
Hopf elliptic-network core is intrinsic to the original
degeneration. Write the original full fibre and the
primitive generator as
\[
        \operatorname{div}(f)=\sum_\alpha n_\alpha S_\alpha,
        \qquad \operatorname{div}(\Omega)
                       =\sum_\alpha a_\alpha S_\alpha.
 \tag{C.AR.6}\hypertarget{tag-00AC}{}\label{cdp:eq:C.AR.6}
\]
Let \(C_\alpha\to T_\alpha\) be the actual finite normal
root cover of the normalization of \(S_\alpha\) defined
by (C.RK.18), and let \(\tau\) be the density threshold
(C.RK.19). Set
\[
 H=\max\left\{\frac{a_\alpha+1}{n_\alpha}:
       C_\alpha\text{ has a component with Hopf minimal model}
       \right\},
 \tag{C.AR.7}\hypertarget{tag-00AD}{}\label{cdp:eq:C.AR.7}
\]
with \(H=\tau\) when this set is empty. Then \(H\ge\tau\).
For every actual divisible semistable root, its highest
Hopf canonical order, when such a component exists, is
\[
                         h=N(H-\tau).
 \tag{C.AR.8}\hypertarget{tag-00AE}{}\label{cdp:eq:C.AR.8}
\]
Thus \(H=\tau\) gives the canonical-flat reduction,
and \(H>\tau\) gives the positive-Hopf reduction of
\ref{cdp:st:C.2.UN.84}--\ref{cdp:st:C.2.UN.85}. The covers, threshold and this branch test
do not depend on an embedded resolution or on its
chosen divisible root exponent.
\end{proposition}



\begin{proof}

\begin{enumerate}[label={$\langle1\rangle$\arabic*.},leftmargin=*]
\item \par\medskip\noindent\textbf{The actual forward construction for an arbitrary fibre.}
Shrink the base to a disc about the chosen value.
The proper direct image \(f_*K_X\) is a torsion-free
rank-one sheaf, so it has a primitive generator
\(\Omega\). Its zero divisor is vertical, giving
\eqref{cdp:eq:C.AR.6}. Take an actual embedded resolution
\(\pi:Y\to X\), by smooth point and curve blowups
in the central fibre, with smooth SNC components
\(D_i\). Retain every centre and every blowup map.
The analytic embedded-resolution comparison used
in \ref{cdp:st:C.2.UN.78}, from \cite{BierstoneMilman1997}, gives
\[
 \operatorname{div}(f\pi)=\sum_i m_iD_i,
 \qquad \operatorname{div}(\pi^*\Omega)=\sum_i k_iD_i,
 \qquad N=\operatorname{lcm}_i m_i.
 \tag{C.AR.9}\hypertarget{tag-00AF}{}\label{cdp:eq:C.AR.9}
\]

\item Normalize the literal fibre product for \(t=s^N\)
and use the equivariant semistable maps of \ref{cdp:st:C.2.UN.78}:
\[
 Z\xrightarrow{q}Y_N
     =\operatorname{Norm}(Y\times_{\Delta_t}\Delta_s)
                   \longrightarrow Y.
 \tag{C.AR.10}\hypertarget{tag-00A1}{}\label{cdp:eq:C.AR.10}
\]
For clarity, the local root lattice is
\[
 \left\{(v,h)\in\Z^r\oplus\Z:
                      \sum_i m_iv_i=Nh\right\}.
 \tag{C.AR.11}\hypertarget{tag-00A2}{}\label{cdp:eq:C.AR.11}
\]
Its old height-one rays are \(((N/m_i)e_i,1)\).
In dimensions at most three, subdividing each
height-one segment or triangle at its lattice
points gives unimodular cones. These refinements
are implemented by normalized blowups of
monomial ideals, which are unchanged when a
boundary equation is multiplied by a unit.
Their labelled maps therefore glue. The root
group fixes the old rays within a chart and
permutes the charts, so orbitwise choices give
the actual equivariant map \eqref{cdp:eq:C.AR.10}. Retaining
the old component and stratum labels prevents
ray identifications within an old cone. Thus
\(Z\) is smooth and its full fibre is reduced
SNC, including globally embedded components.

\item The normalized root is connected. Its generic
fibre over \(s\ne0\) is the connected original
torus; all components of the normalized fibre
product dominate the base, so there cannot be
an additional component. Equivalently its
Stein factor has degree one over the normal
disc. All the toric maps are proper and are
isomorphisms away from zero.

\item Apply the complete reduced theorem \ref{cdp:st:C.2.UN.85} to
\(Z\), and retain its actual comparison
\(\phi:Z\dashrightarrow X_*\). This gives
one of the listed cores and all the inverse
chart maps in the input. The root action
becomes the actual birational action
\[
            \gamma_* =\phi\gamma\phi^{-1}
                            \quad(\gamma\in\mu_N).
 \tag{C.AR.12}\hypertarget{tag-00A3}{}\label{cdp:eq:C.AR.12}
\]
Its group law holds on the dense punctured
family, hence as meromorphic maps. Reversing
\(\phi\) first restores the original regular
action on \(Z\). Applying the retained
equivariant coarsening then gives \(Y_N\).
For a normal local ring \(A\) of \(Y\),
let \(B\) be the integral closure of
\(A[s]/(s^N-f\pi)\). An invariant element
of \(B\) belongs to \(\operatorname{Frac}(A)\)
and is integral over \(A\), so it belongs
to \(A\). The reverse inclusion is immediate.
This remains true when local root branches
split and are permuted by the group. Thus
\(B^{\mu_N}=A\), and the full comparison is
\[
 X_*\dashrightarrow Z\xrightarrow{q}Y_N
              \longrightarrow Y_N/\mu_N=Y
                       \xrightarrow{\pi}X.
 \tag{C.AR.13}\hypertarget{tag-00A4}{}\label{cdp:eq:C.AR.13}
\]
Every arrow and overlap in this diagram has
been retained. In particular the finite
quotient follows the inverse record; no
central quotient of an indeterminate
birational action is taken.

\par\medskip\noindent\textbf{The converse from the specified finite input.}
Begin with the data at the start of the
subsection, independently of any original
\(X\). The inverse sigma maps give the
specified smooth proper \(Z\), and the
chart checks give an actual regular finite
action. Its equivariant coarsening and
invariant-coordinate charts construct the
smooth \(Y\) and its function \eqref{cdp:eq:C.AR.1}.
The subsequent ordinary contraction maps
construct the smooth \(X\) and its descended
function. Properness follows from the
proper core and retained proper maps;
for a local overlap replacement it is the
proper-graph and annular comparison of
\ref{cdp:st:C.2.UN.75}. The ordinary target maps and finite
quotient are proper, so the final family
is proper after the common disc has been
shrunk if required.

\item For a fixed nonzero \(t\), the \(N\)
choices of \(s\) are permuted freely and
transitively by \(\mu_N\). Inclusion of
any one \(Z_s\) into their union induces
a biholomorphism onto the quotient fibre.
The coarsening and ordinary maps are
isomorphisms there. Consequently the
final general fibre is exactly that
connected complex two-torus. Since the
total space is smooth, the nonconstant
base function is a non-zero-divisor;
the map to the smooth curve is flat.
Its Stein factor has generic degree one
and its base is normal, so all its fibres
are connected. These deductions use only
the supplied maps and their checks.

\item The inverse elementary record is therefore
a constructive description of the resulting
analytic support, with its actual ambient
neighbourhood. Forward construction shows
that every original support has such a
record. The converse shows that every
admissible record produces such a support.
Isomorphic labelled atlases and commuting
record maps give biholomorphic outputs
by their explicit glued isomorphism.
No uniqueness of the reduction record is
required for this exact generation statement.

\par\medskip\noindent\textbf{Original support Euler numbers, including the finite quotient.}
The actual calculation (C.RK.15)--(C.RK.17)
applies to \eqref{cdp:eq:C.AR.10} without a canonical-order
restriction. Over \(D_I^\circ\), the central
normal root is a finite etale cover of degree
\(d_I\). The toric resolution fibre has Euler
number \(\nu_I\): its zero-dimensional torus
orbits are the unimodular top cones, whose
number is the normalized volume of the
height-one simplex. The exact Bezout
determinant in (C.RK.15) gives the value
in \eqref{cdp:eq:C.AR.2}. Thus
\[
       e(S_Z)-e(S_Y)=R_N.
 \tag{C.AR.14}\hypertarget{tag-00A5}{}\label{cdp:eq:C.AR.14}
\]
The fixed-point average for the actual
finite root group gives the same quotient
comparison, including globally nontrivial
sheet permutations. This is not a free-cover
Euler division at the central fibre.

\item The ordinary blowups give
\(e(S_Y)-e(S)=\Delta_\pi\), and the
forward sigma record gives
\(e(S_*)-e(S_Z)=\Delta_\sigma\).
For an arrow with centre ideal \(I\),
its signed term is the difference between
the Euler numbers of the actual
\(\operatorname{Proj}(\bigoplus I^a/I^{a+1})\)
and its centre, on the finite constructible
subdivision of \ref{cdp:st:C.2.UN.75}. Restored overlaps
cancel. Combining these three exact
comparisons proves \eqref{cdp:eq:C.AR.4}. The core
Euler numbers are those of \ref{cdp:st:C.2.UN.85}, proving
\eqref{cdp:eq:C.AR.5}. Since both construction
directions have already been established,
this selects the full class of zero-Euler
outputs, including singular and nonreduced
original fibres.

\par\medskip\noindent\textbf{Intrinsic root data and the two endpoint branches.}
The local finite algebras defining
\(C_\alpha\) are glued from
\(w_j^{n_\alpha}=u_j|_{T_\alpha}\), where
\(f=g_j^{n_\alpha}u_j\); their transition
is \(w_k=(g_j/g_k)w_j\). Their finite
normalization is an actual compact
normal cover. No global meromorphic
frame of a nonalgebraic normal line
is required. For every divisible
root, the strict-cover comparison is
the integral local function
\(w_j=s^{N/n_\alpha}/g_j\), as proved
in (C.RK.21). Every connected component
of this fixed cover is obtained, with
its actual proper birational resolution.

\item All exceptional surfaces of the original
embedded resolution are projective, and
their finite normal covers are projective
by the coherent-algebra GAGA comparison
in \ref{cdp:st:C.2.UN.78}. The additional semistable
components are ruled or rational.
Thus no Hopf component is omitted by
the original covers \(C_\alpha\).

\item Proper change of variables in the actual
density gives
\[
                   \tau=\min_i\frac{k_i+1}{m_i}.
 \tag{C.AR.15}\hypertarget{tag-00A6}{}\label{cdp:eq:C.AR.15}
\]
At a height-one root ray \((v,1)\),
the pulled canonical form has order
\(\sum_i(k_i+1)v_i-1\). Its minimum
is \(N\tau-1\); the old strict ray
over \(D_i\) has order
\(N(k_i+1)/m_i-1\). Dividing by
\(s^{N\tau-1}\) gives the primitive
root generator, whose strict orders
are therefore
\[
                  N\left(\frac{k_i+1}{m_i}-\tau\right).
 \tag{C.AR.16}\hypertarget{tag-00A7}{}\label{cdp:eq:C.AR.16}
\]
Original strict transforms retain
\((k_i,m_i)=(a_\alpha,n_\alpha)\).
The fixed-cover comparison and exclusion
of exceptional Hopfs now give exactly
\eqref{cdp:eq:C.AR.8}, including its integrality.
When no Hopf occurs, \ref{cdp:st:C.2.UN.75} applies
directly. Otherwise all Hopf root
orders are zero exactly when
\(H=\tau\). If \(H>\tau\), \ref{cdp:st:C.2.UN.84}
and \ref{cdp:st:C.2.UN.85} give the positive-Hopf core.
Both the threshold and the covers
were defined on the original family,
so this branch test is independent
of the resolution and exponent.

\par\medskip\noindent\textbf{Compact data and the extent of the core parametrizations.}
For an original compact base \(B\), let
\(c\) be the selected value and let
\(\Sigma\) be the finite critical-value set.
Choose a point \(p\notin\Sigma\cup\{c\}\).
For sufficiently large \(d\),
\[
 H^1\!\left(B,\OO_B(dp-2c-
                    \sum_{v\in\Sigma\setminus\{c\}}v)\right)=0.
\]
Indeed its degree is greater than \(2g(B)-2\),
so Serre duality gives the vanishing. The
actual restriction sequence consequently
surjects from \(H^0(B,\OO_B(dp))\) onto
the prescribed first jet at \(c\) and
values at the other critical points.
Choose the meromorphic function \(a\)
with jet \(a=t+O(t^2)\) in the fixed
coordinate at \(c\), and value one at
every other point of \(\Sigma\).
Its only possible pole is \(p\), and
all its other zeros and poles are regular
values of the original family.

\item Normalize the compact cyclic base cover
\(s^N=a\). The simple zero at \(c\)
makes it connected and totally ramified
there. It is etale over every other
critical value. Near its additional
branch values the original map is a
submersion, so its base change is locally
a product with the smooth normalized
cover curve and is smooth. Over the
other critical values the base change
is etale over the original smooth total
space. The normalized total space is
therefore smooth away from the selected
root fibre. Apply the preceding root
resolution and modification maps there,
retaining the actual identity comparisons
on its annulus and on the complement.
The local function \(a\) is a coordinate
at \(c\). Write \(a=t u(t)\), with
\(u(0)=1\), and choose its analytic
\(N\)-th root near \(c\). The map
\[
 (t,s)\longmapsto
 \left(t,\sigma=\frac{s}{u(t)^{1/N}}\right),
 \qquad \sigma^N=t,
\]
is an actual equivariant isomorphism
between the global-cover root chart
and the retained local root model;
its inverse is \(s=\sigma u(t)^{1/N}\).

\item The finite action is conjugated by that
actual global modification and becomes
regular again upon reversal. The quotient
and original blowdowns recover the
original compact family. Conversely the
globally supplied compact admissible
data give a compact output by the same
proper maps.

\item The central support parameters of the
flat cycles and elliptic networks are
the complete marked line, torsor,
finite-character and gluing data in
\ref{cdp:st:C.2.UN.57} and \ref{cdp:st:C.2.UN.61}--\ref{cdp:st:C.2.UN.63}. The positive
toric core is parametrized by its
periodic affine triangulation and
normalized log gluing data in \ref{cdp:st:C.2.UN.71}.
Their compact occurrence is proved
in those statements and \ref{cdp:st:C.2.UN.73}.
Those constructions supply actual
compact starting families for records
defined on their neighbourhoods.
For recovery of a different given
family, its actual endpoint atlas
and higher transverse terms remain
in the input. The proof uses its
specified overlap isomorphisms and
does not transfer a contraction
through an equality of central log
data alone. This states exactly
which parameters are classified
centrally and which analytic data
the reversible reconstruction retains.

\end{enumerate}
\end{proof}






\section{Global CDP replacement and normalization resonance}\hypertarget{tag-034C}{}\label{sec:II:all:12}


\begin{theorem}[The Euler sign for every degeneration and its equality criterion]
\hypertarget{tag-01BZ}{}\label{cdp:st:C.2.G.1}
\StatementRecord{cdp:st:C.2.G.1}{proved}{01BZ}{\SeriesRef{cdp:st:D.1.9}, \SeriesRef{cdp:st:A.3.4}, \SeriesRef{cdp:st:A.5.2}, \SeriesRef{cdp:st:A.5.3}, \SeriesRef{cdp:st:C.2.MH.1}, \SeriesRef{cdp:st:C.2.MH.2}, \SeriesRef{night:P4.26}, \SeriesRef{night:P4.32}, \SeriesRef{night:P4.36}, \SeriesRef{night:P4.37}, \SeriesRef{cdp:st:C.2.INT.1}, \SeriesRef{cdp:st:C.2.INT.2}, \SeriesRef{cdp:st:C.2.INT.3}, \SeriesRef{cdp:st:C2a}; \cite[Proposition~2.1 and Theorem~2.2, pp.~842--843]{Peters2021}}

Let $Y$ be a connected smooth compact complex threefold with
$b_1(Y)=b_2(Y)=0$, and let $f:Y\to B$ be a proper surjective
holomorphic map with connected fibres to a smooth compact curve.
Assume that the smooth general fibre is a complex two-torus.
The full fibres may have arbitrary positive component multiplicities
and their reduced supports may have arbitrary singularities. Then
\[
 B\simeq\PP^1,\qquad
 c_3(Y)=e(Y)\in\{0,2\},\qquad b_3(Y)=2-c_3(Y).
 \tag{C.G.1}\hypertarget{tag-00KW}{}\label{cdp:eq:C.G.1}
\]
Every reduced fibre support has nonnegative Euler number.
An irreducible singular support with rational normalization
resolution has Euler number two; every other irreducible support
has Euler number zero.

The equality $c_3(Y)=0$ holds exactly when every support has one
of the following descriptions. A smooth irreducible support has
a finite unramified torus cover. A singular irreducible support
has smooth minimal elliptic ruled or smooth minimal Hopf
normalization. A reducible support has normalization a disjoint
union of smooth minimal Hopf or non-K\"ahler elliptic surfaces;
its reduced conductor image and inverse image are disjoint
unions of smooth elliptic curves, and the conductor branch maps
are finite unramified covers. Each description concerns the
normalization and resolution maps of the actual fibre. In
particular the full-fibre and conductor multiplicities are
retained.
\end{theorem}


\begin{proof}
\begin{enumerate}[label={$\langle1\rangle$\arabic*.},leftmargin=*]
\item 
\noindent\textbf{Step 1. The base is rational and the local
finite character is available on every full fibre.}
The beginning of connected-fibre Leray injects
$H^1(B,\Q)$ into $H^1(Y,\Q)=0$, so $B=\PP^1$.
Proposition~\SeriesRef{cdp:st:D.1.9} supplies an integral section
of $R^1f_*\Z$ which restricts nontrivially to each reduced
support and to each nearby torus. On a small disc its
Leray lift is an integral class on the entire tube.
Reduction modulo a prime which does not divide its fibre
restriction gives the finite cyclic character used in
\ref{cdp:st:C.2.MH.2}. The construction is on the original
full fibre and requires no change of its multiplicity.

\medskip\noindent\textbf{Step 2. A normal irreducible support
is smooth and has a finite unramified torus cover.}
Proposition~\SeriesRef{cdp:st:A.5.2} excludes a singular normal
support. For a smooth support $S$ in a full fibre $mS$,
the normalized change $t=s^m$ is, in a local equation
$t=uz^m$, the map
\[
 w^m=u,\qquad s=wz.
\]
These charts give a finite unramified central cover of $S$
and a proper holomorphic submersion over the new disc,
as in \SeriesRef{cdp:st:A.3.4}. Its nearby fibre is the
original torus. Smooth proper transport gives $b_1=4$
and rational first Chern class zero on the central cover.
It is K\"ahler by the even-$b_1$ criterion, and the
surface classification identifies it as a torus; the
precise classification facts are
\cite[Proposition~2.1 and Theorem~2.2, pp.~842--843]{Peters2021}.
Thus $e(S)=0$.

\medskip\noindent\textbf{Step 3. The irreducible nonnormal
normalization falls into three explicit resolution classes.}
Write $Z\to T\to S$ for the minimal resolution of the
normalization. The canonical evaluation and
\SeriesRef{cdp:st:A.5.3} give a nonempty effective divisor $D$
with $K_Z+D$ torsion. A positive pluricanonical section,
raised to a power killing this torsion, would give a
nonzero holomorphic function vanishing along a positive
multiple of $D$. Compactness excludes that function.
Hence $\kappa(Z)=-\infty$. The classification in
\cite[Theorem~2.2, p.~843]{Peters2021} gives a rational
resolution, a resolution ruled over a curve of positive
genus, or a resolution of class VII. The actual maps
$Z\to T\to S$ remain the maps in every subsequent
comparison.

\medskip\noindent\textbf{Step 4. Every irreducible support has
the claimed exact Euler value and normalization.}
For a rational resolution,
\ref{cdp:st:C.2.MH.2} gives $e(S)=2$. For a resolution
ruled over a curve of genus $g\ge1$, the same proposition
gives $e(S)=2-2g$. The local irreducible inequality in
\SeriesRef{night:P4.26} gives $e(S)\ge0$, and therefore $g=1$
and $e(S)=0$. Its proof uses the effective conductor,
the local trace inequality, and the fundamental group
of this support; no assertion about the other fibre
supports enters that local inequality. The zero-Euler
comparison in \SeriesRef{night:P4.32} then identifies $T$ as
a smooth minimal elliptic ruled surface.

For a class-VII resolution, \SeriesRef{night:P4.37} excludes
every minimal model with positive second Betti number
by its saturated-foliation Chern identity and the
negative-definite curve form. The curve-free Inoue
case cannot carry the nonempty effective conductor.
The remaining Hopf case is treated by
\SeriesRef{night:P4.36}: its derivative and foliation
calculations force the actual normalization to be
smooth and minimal, with $e(S)=0$. These arguments
are local comparisons for the given irreducible
support; they do not invoke the irreducible-or-normal-
crossings hypothesis on all fibres in
\SeriesRef{night:P4.58}. This proves the asserted exact
irreducible classification in the present scope.

\medskip\noindent\textbf{Step 5. Arbitrary reducible supports
have nonnegative Euler number, with the stated equality data.}
For each normalization resolution, an adjacent global
component supplies a nonzero effective divisor with
rationally trivial first Chern class. Its positive
pairing with a proposed K\"ahler form excludes a
K\"ahler resolution. Proposition~\ref{cdp:st:C.2.INT.1}
gives the actual Picard restriction to each elliptic
curve and excludes a Hopf--Enoki or elliptic--Enoki
interface. Proposition~\ref{cdp:st:C.2.INT.2} excludes
the remaining all-Enoki possibility using the full
multiplicities, the canonical coefficient maximum,
and the local vanishing sphere of square $-2$.

Proposition~\ref{cdp:st:C.2.INT.3} now applies to the
entire reduced support. With $A_i$ its normalization
conductor and $C$ its reduced image, its actual
normalization map gives
\[
 e(S)=\sum_i e(T_i\setminus A_i)+e(C),
 \qquad e(T_i\setminus A_i)\ge0,\quad e(C)\ge0.
\]
The proof retains the scheme conductor in the
coherent-duality argument and all contractions in
the topological calculation. It proves both
directions of the equality characterization:
the normalizations are the stated smooth minimal
surfaces, the conductor curves are smooth elliptic,
and their maps are finite unramified covers.

\medskip\noindent\textbf{Step 6. The global sum is zero or two.}
The finite-critical-value and Euler-additivity
calculation of \ref{cdp:st:C2a}, together with
complex-oriented Chern--Gauss--Bonnet, gives
\[
 c_3(Y)=e(Y)=\sum_{b\in B}e((Y_b)_{\rm red})\ge0.
 \tag{C.G.2}\hypertarget{tag-00KX}{}\label{cdp:eq:C.G.2}
\]
The smooth torus contributes zero over the regular
base. Poincar\'e duality gives $e(Y)=2-b_3(Y)$.
The nonsingular middle cup-product pairing is
skew-symmetric over $\Q$, so $b_3(Y)$ is even.
Its nonnegativity now proves \eqref{cdp:eq:C.G.1}.

\medskip\noindent\textbf{Step 7. Vanishing is equivalent to
the listed geometric data on every fibre.}
Every term in \eqref{cdp:eq:C.G.2} is nonnegative;
the sum vanishes precisely when every term does.
Steps 2 and 4 give the irreducible forms, and
Step 5 gives the reducible forms with their exact
conductor maps. Conversely each of these forms has
zero Euler number by the same local calculations,
so their occurrence on every support makes the
global sum zero.

\end{enumerate}
\end{proof}










\clearpage
\part{Results not yet established}
\section{The remaining splitting problem}\hypertarget{tag-000B}{}\label{app:II:programs}
\begin{problem}[The remaining comb bundle splitting]\hypertarget{tag-01SW}{}\label{inv:E2.3}
\StatementRecord{inv:E2.3}{open}{01SW}{\SeriesRef{invII:eq:E.6}, \SeriesRef{cdp:st:E15}, \SeriesRef{invII:eq:E.4a}, \SeriesRef{invII:eq:E.5}, \SeriesRef{invII:eq:E.7}}

Proposition~\ref{inv:E2.5} determines the actual splitting
$K_{12}=\mathcal O_B(-8)\oplus\mathcal O_B(-9)^2$ and its zero
space of sections on a nonempty open family of twelve-tooth smoothings.
The remaining problem is to determine the splitting of the rank-three
bundle $K_m$ in \eqref{invII:eq:E.6} for general tooth counts and for
special members outside that proved open family, together with its
exact number of sections.
The Kuranishi germ and its vanishing obstruction map are established by
Proposition~\ref{cdp:st:E15}. The four tangent
cohomology dimensions themselves are fully determined by Proposition~\ref{cdp:st:E15}.
\end{problem}


\begin{proof}[The remaining finite data and the computations made]

\begin{enumerate}[label={$\langle1\rangle$\arabic*.},leftmargin=*]
\item If $e_m$ is the length of $\operatorname{coker}\kappa_1$ in \eqref{invII:eq:E.6},
then its source has rank four and degree $-2$, its target is
$\OO_B(2m)$, and
\[
 \deg K_m=-2m-2+e_m.                                  \tag{E.7}\hypertarget{tag-01TW}{}\label{invII:eq:E.7}
\]
The degree $-2$ follows from the two pairs of opposite stabilizer
characters of $E_m^*\otimes G_m$; its root determinant is trivial,
and each finite root point subtracts one from the coarse determinant.
Although \eqref{invII:eq:E.4a} determines the tangent dimensions independently of
the individual zeros of the Kodaira--Spencer determinant, the
support and elementary divisors of $Q_m$ still depend on that
period-section map. The tangent dimensions alone do not determine obstruction maps.
The complete period-section family and its tangent isomorphism in
Proposition~\ref{cdp:st:E15} establish that the Kuranishi ideal is zero.

\item We evaluated the universal exact degree and dimension expressions
for $m=12,13,20,100$ in \texttt{d\_exact.py} and
\texttt{d\_comb.py}; the first case already
gives $h^{1,1}\ge49$, $h^{1,2}\ge48$, $h^0(T)=25$,
$h^1(T)=97$ and $h^2(T)=73$. Proposition~\ref{inv:E2.5} now computes the actual kernel on a
nonempty open family of smoothed twelve-tooth period sections. Its
flat relative kernel, exact conic two-jet maps and tooth-collapse
comparison give the splitting
$(-8,-9,-9)$, the vanishing cokernel and the attained values
$h^{1,1}=49$, $h^{1,2}=48$. This proves the open-family case without
assigning transition matrices to an unspecified individual member.
The remaining calculation concerns general tooth counts and special
members outside that open family; formulas
\eqref{invII:eq:E.5}--\eqref{invII:eq:E.7} and the uniform bounds
remain the constraints for those calculations.

\end{enumerate}
\end{proof}




\paragraph{Bounds obtained from all possible coarse splittings.}
The available bundle data give the further bounds
\[
 0\le h^0(K_m)\le 2m+2,\qquad
 4m+1\le h^{1,1}(X_m)\le6m+3,\qquad
 4m\le h^{1,2}(X_m)\le6m+2.
\]
Here is a complete derivation. Put $S=2m-1$, and write
\[
 V_c=\rho_*E_m=\OO(a)\oplus\OO(S-a),\qquad
 G_c=\rho_*G_m=\OO(b)\oplus\OO(S-b).
\]
The already proved nonnegative splitting implies $0\le a,b\le S$.
Reordering each pair gives $a,b\le S/2$.
Taking invariants defines a natural injective morphism
\[
 \rho_*\mathcal Hom(E_m,G_m)\longrightarrow
 \mathcal Hom(V_c,G_c).
\]
It is an isomorphism away from the two stack points, and its source is
torsion free, which proves injectivity. In the invariant local bases,
the map from the nontrivial-character source line to the trivial-character
target line has to vanish at the stack point; the other entries have no
additional condition. Thus there is one elementary condition at each
finite point, consistently with the source degree $-2$ in
\eqref{invII:eq:E.7}.

The four contributions to $h^0(\mathcal Hom(V_c,G_c))$ are
\[
 \max\{b-a+1,0\},\quad S-a-b+1,\quad
 \max\{a+b-S+1,0\},\quad\max\{a-b+1,0\}.
\]
If $a+b<S$, the third is zero and their sum is at most
$S-a-b+1+|a-b|+2\le S+3$. If $a+b=S$, the two inequalities
$a,b\le S/2$ imply $a=b=S/2$, and the sum is four, also at most
$S+3$. Since $K_m$ is a subbundle of
$\rho_*\mathcal Hom(E_m,G_m)$, this proves $h^0(K_m)\le S+3=2m+2$.
Substitution in \eqref{invII:eq:E.5} proves the two Hodge bounds.
In particular $0\le h^0(K_{12})\le26$.
Proposition~\ref{inv:E2.5} attains the lower endpoint on its open
family, and therefore attains the lower Hodge bounds for $m=12$.
The derivation of the uniform bounds above uses the stated
splitting constraints. Their upper endpoints have not been shown
sharp by this calculation.

\paragraph{Why the numerical constraints do not determine the answer.}
Set $D=\OO^{\oplus2}\oplus\OO(-1)^{\oplus2}$ and
$L=\OO(2m)$ on $\PP^1$, with homogeneous coordinates $[X:Y]$.
The two morphisms
\[
 \phi_A=(X^{2m},Y^{2m},0,0),\qquad
 \phi_B=(0,0,X^{2m+1},Y^{2m+1}):D\longrightarrow L
\]
are surjective, because their two nonzero entries have no common
projective zero. The degree-$2m$ pencil sequence, together with the
unused summands, gives
\[
 \ker\phi_A=\OO(-2m)\oplus\OO(-1)^{\oplus2}.
\]
The degree-$(2m+1)$ pencil sequence tensored by $\OO(-1)$ gives
\[
 \ker\phi_B=\OO^{\oplus2}\oplus\OO(-2m-2).
\]
Both maps therefore have rank-four source of degree $-2$, target of
degree $2m$, and cokernel of length zero. Their rank-three kernels
both have degree $-2m-2$, but their section dimensions are zero and
two. For $m=12$ the splitting triples are $(-24,-1,-1)$ and
$(0,0,-26)$. These exact vector-bundle examples show that the known
rank, degree and cokernel length alone cannot determine the kernel.
We do not assert that either morphism is the cup map of an admissible
comb period section, or that different such geometric sections have
different Hodge numbers.

There is an exact relation between the tangent and cotangent cup maps.
On the root curve put $F=E_m^*$, $G=G_m$, $D=\det G$ and
$\mathcal L=\OO(-5)$, and set
$N=\operatorname{coker}(\delta_0:\mathcal L\to E_m\otimes G)$.
The wedge isomorphism
\[
 J:G\longrightarrow G^*\otimes D,\qquad J(g)(h)=h\wedge g,
\]
has matrix $\left(\begin{smallmatrix}0&1\\-1&0\end{smallmatrix}\right)$
in an ordered frame. With the same Cech boundary convention,
dualizing the tangent extension replaces its off-diagonal block
$\beta$ by $-\beta^t$. Consequently its cotangent cup map is
\[
 \kappa_{\mathrm{root}}
 =-(\delta_0^*\otimes I_D)(I_F\otimes J).
\]
Indeed, for $\delta_0(\ell)=u_1\otimes(ag_1+bg_2)
+u_2\otimes(cg_1+dg_2)$, the row of the right-hand side
in the ordered basis $(f_1g_1,f_1g_2,f_2g_1,f_2g_2)$ is
$(b,-a,d,-c)$, exactly the contracted dual-extension cup product.
Dualizing the cokernel sequence and then taking finite-stabilizer
invariants gives the explicit kernel identification
\[
 K_m\simeq\rho_*(N^*\otimes D),\qquad D=\OO(24m-1).
\]
Here the identification is induced by $I_F\otimes J$; possible
torsion in $N$ is killed by the dual, so no saturation of
$\operatorname{im}\delta_0$ is assumed. Finite invariants preserve
the kernel because averaging over an order-$r$ stabilizer is the
projection $r^{-1}\sum_g g$.
The tangent calculation determines the cohomology of $\rho_*N$.
To split $\rho_*(N^*\otimes D)$, the remaining task is to choose
an actual admissible period section and compute its transition
matrices and derivative. The existence and smoothing construction
supplies an open set of such sections. The exact duality above,
the estimates and the algebraic examples preserve this remaining
geometric choice and cup-map calculation.

\clearpage



\clearpage
\section*{Programs}
\addcontentsline{toc}{section}{Programs}
The following files contain the indicated computations at commit
\begin{center}\small\ttfamily\SsixRepositoryCommit\end{center}
Each link identifies the complete repository path.

{\footnotesize
\renewcommand{\arraystretch}{.9}
\setlength{\tabcolsep}{4pt}
\begin{longtable}{@{}>{\raggedright\arraybackslash}p{.57\textwidth}>{\raggedright\arraybackslash}p{.39\textwidth}@{}}
\caption{Programs and their mathematical scope.}\label{tab:II:programs}\\
\toprule Computation and statement & Source file \\
\midrule\endfirsthead
\toprule Computation and statement & Source file \\
\midrule\endhead
\bottomrule\endfoot
Tangent index; integral monodromy, including Wang's four matrices and integral coinvariant quotient; Kodaira and Hopf surfaces; 120 unit quaternions, 14,400 products and a surjective commutator map over $\Q(\sqrt5)$ represented by rational pairs (\cref{cdp:st:C.1.1,cdp:st:C.3.2,cdp:st:C.3.3,cdp:st:D.1.6}) & \href{https://raw.githubusercontent.com/KokunoYumeto/s6-complex-structure-workbench/\SsixRepositoryCommit/pro_returns/S6_03_CDP_checkpoint_13_20261116T0136Z/S6_03_CDP/code/D_exact_checks.py}{\nolinkurl{D_exact_checks.py}} \\[3pt]
Normalization, inverse Jacobian module and residue poles (\cref{cdp:st:A.2.5}) & \href{https://raw.githubusercontent.com/KokunoYumeto/s6-complex-structure-workbench/\SsixRepositoryCommit/pro_returns/S6_03_CDP_checkpoint_13_20261116T0136Z/S6_03_CDP/code/A_vertical_pole.py}{\nolinkurl{A_vertical_pole.py}} \\[3pt]
Curve equations, transverse jets and conormal valuations (\cref{cdp:st:A.2.6}) & \href{https://raw.githubusercontent.com/KokunoYumeto/s6-complex-structure-workbench/\SsixRepositoryCommit/pro_returns/S6_03_CDP_checkpoint_13_20261116T0136Z/S6_03_CDP/code/A_non_lci_cusp.py}{\nolinkurl{A_non_lci_cusp.py}} \\[3pt]
Exterior-square parity examples and Euler numbers of blowup families (\cref{cdp:st:C.2.UN.1,cdp:st:C.2.UN.2}) & \href{https://raw.githubusercontent.com/KokunoYumeto/s6-complex-structure-workbench/\SsixRepositoryCommit/pro_returns/S6_03_CDP_checkpoint_13_20261116T0136Z/S6_03_CDP/code/root_unrestricted_euler.py}{\nolinkurl{root_unrestricted_euler.py}} \\[3pt]
Euler and integral cohomology of vertical point and curve blowups (\cref{cdp:st:C.2.UN.4}) & \href{https://raw.githubusercontent.com/KokunoYumeto/s6-complex-structure-workbench/\SsixRepositoryCommit/pro_returns/S6_03_CDP_checkpoint_13_20261116T0136Z/S6_03_CDP/code/E_vertical_refinements.py}{\nolinkurl{E_vertical_refinements.py}} \\[3pt]
Exact rational jet restrictions at orders $2,3,4,6,8$ and conductor ideals of the two genus-one plane cubics (\cref{night:P4.14}) & \href{https://raw.githubusercontent.com/KokunoYumeto/s6-complex-structure-workbench/\SsixRepositoryCommit/pro_returns/S6_06_BLIND_checkpoint_15_20261117T1457Z/S6_06_BLIND/code/part4/genus_one_conductor.py}{\nolinkurl{genus_one_conductor.py}} \\[3pt]
Jet covariance, binary-cubic discriminants and inverse-Jacobian pole coefficients (\cref{cdp:st:C2d}) & \href{https://raw.githubusercontent.com/KokunoYumeto/s6-complex-structure-workbench/\SsixRepositoryCommit/pro_returns/S6_03_CDP_checkpoint_13_20261116T0136Z/S6_03_CDP/code/C_jets.py}{\nolinkurl{C_jets.py}} \\[3pt]
Integral compact intersection lattices, primitive boundary kernels, blowup charts and the alternative two-blowup compactification (\cref{night:P4.16,night:P4.17}) & \href{https://raw.githubusercontent.com/KokunoYumeto/s6-complex-structure-workbench/\SsixRepositoryCommit/pro_returns/S6_06_BLIND_checkpoint_15_20261117T1457Z/S6_06_BLIND/code/part4/milnor_intersection.py}{\nolinkurl{milnor_intersection.py}} \\[3pt]
Coprime exponent matrices, transfer pairings, normalization minors and inverse coordinates on the affine nearby fibre (\cref{night:P4.18,night:P4.19}) & \href{https://raw.githubusercontent.com/KokunoYumeto/s6-complex-structure-workbench/\SsixRepositoryCommit/pro_returns/S6_06_BLIND_checkpoint_15_20261117T1457Z/S6_06_BLIND/code/part4/binomial_milnor.py}{\nolinkurl{binomial_milnor.py}} \\[3pt]
Integral monodromy matrices, rational centralizers, fixed lattices, norm lattices and primitive vanishing cycles (\cref{push:src:E.5.2}) & \href{https://raw.githubusercontent.com/KokunoYumeto/s6-complex-structure-workbench/\SsixRepositoryCommit/pro_returns/S6_02_PUSH_final_20261115T2217Z/S6_02_PUSH/code/e_push_two_parameter.py}{\nolinkurl{e_push_two_parameter.py}} \\[3pt]
Conic coordinate maps, component multiplicities, weighted scalings and oriented sphere Euler numbers (\cref{night:P4.25}) & \href{https://raw.githubusercontent.com/KokunoYumeto/s6-complex-structure-workbench/\SsixRepositoryCommit/code/S6_SERIES_PUBLISH_II_A_02WX_conic_maps.py}{\nolinkurl{S6_SERIES_PUBLISH_II_A_02WX_conic_maps.py}} \\[3pt]
Transverse two-jets and critical-value remainders of the double and triple deformations (\cref{night:P4.25}) & \href{https://raw.githubusercontent.com/KokunoYumeto/s6-complex-structure-workbench/\SsixRepositoryCommit/code/S6_SERIES_PUBLISH_II_A_02WX_deformation_jets_postrelease.py}{\nolinkurl{S6_SERIES_PUBLISH_II_A_02WX_deformation_jets_postrelease.py}} \\[3pt]
Alternating period pairing, integral monodromy, weighted line cohomology and the branch-value Jacobian (\cref{cdp:st:E13}) & \href{https://raw.githubusercontent.com/KokunoYumeto/s6-complex-structure-workbench/\SsixRepositoryCommit/code/S6_SERIES_PUBLISH_II_01GK_prereading_certificate.py}{\nolinkurl{S6_SERIES_PUBLISH_II_01GK_prereading_certificate.py}} \\[3pt]
Restriction-map inverses, residue-frame comparison and the Laurent connection identity (\cref{cdp:st:E13}) & \href{https://raw.githubusercontent.com/KokunoYumeto/s6-complex-structure-workbench/\SsixRepositoryCommit/code/S6_SERIES_PUBLISH_II_01GK_postrelease_inverse_certificate.py}{\nolinkurl{S6_SERIES_PUBLISH_II_01GK_postrelease_inverse_certificate.py}} \\[3pt]
Barycentric retraction identities and the symbolic meridian winding map (\cref{night:P4R.1}) & \href{https://raw.githubusercontent.com/KokunoYumeto/s6-complex-structure-workbench/\SsixRepositoryCommit/code/S6_SERIES_PUBLISH_II_A_02Y2_neighborhood_certificate.py}{\nolinkurl{S6_SERIES_PUBLISH_II_A_02Y2_neighborhood_certificate.py}} \\[3pt]
Integral marking, period and root-coordinate maps, and weighted-line cohomology (\cref{cdp:st:E11}) & \href{https://raw.githubusercontent.com/KokunoYumeto/s6-complex-structure-workbench/\SsixRepositoryCommit/code/S6_SERIES_PUBLISH_II_01GI_prereading_weight_certificate.py}{\nolinkurl{S6_SERIES_PUBLISH_II_01GI_prereading_weight_certificate.py}} \\[3pt]
Hilbert-series recurrence and the degree-three Riemann--Roch expansion (\cref{cdp:st:E11}) & \href{https://raw.githubusercontent.com/KokunoYumeto/s6-complex-structure-workbench/\SsixRepositoryCommit/code/S6_SERIES_PUBLISH_II_01GI_postrelease_hilbert_certificate.py}{\nolinkurl{S6_SERIES_PUBLISH_II_01GI_postrelease_hilbert_certificate.py}} \\[3pt]
Euler identities from curve and sphere cells, point blowups and class-VII Betti ranks (\cref{night:P4R.2}) & \href{https://raw.githubusercontent.com/KokunoYumeto/s6-complex-structure-workbench/\SsixRepositoryCommit/code/S6_SERIES_PUBLISH_II_A_02YD_euler_cells.py}{\nolinkurl{S6_SERIES_PUBLISH_II_A_02YD_euler_cells.py}} \\[3pt]
Meridian relation-space isomorphism, determinant sign and incidence-tree Euler count (\cref{night:P4R.4}) & \href{https://raw.githubusercontent.com/KokunoYumeto/s6-complex-structure-workbench/\SsixRepositoryCommit/code/S6_SERIES_PUBLISH_II_A_02YG_relation_certificate.py}{\nolinkurl{S6_SERIES_PUBLISH_II_A_02YG_relation_certificate.py}} \\[3pt]
Cusp normalization jets, conductor cochains, weighted cohomology and restriction-map inverses (\cref{cdp:st:E14}) & \href{https://raw.githubusercontent.com/KokunoYumeto/s6-complex-structure-workbench/\SsixRepositoryCommit/code/S6_SERIES_PUBLISH_II_01GL_prereading_exact.py}{\nolinkurl{S6_SERIES_PUBLISH_II_01GL_prereading_exact.py}} \\[3pt]
Critical-value norm and exact cusp jet basis changes (\cref{cdp:st:E14}) & \href{https://raw.githubusercontent.com/KokunoYumeto/s6-complex-structure-workbench/\SsixRepositoryCommit/code/S6_SERIES_PUBLISH_II_01GL_postrelease_critical_norm.py}{\nolinkurl{S6_SERIES_PUBLISH_II_01GL_postrelease_critical_norm.py}} \\[3pt]
Weighted determinant polynomials, integral period pairings and the Grassmann quotient matrix (\cref{cdp:st:E16}) & \href{https://raw.githubusercontent.com/KokunoYumeto/s6-complex-structure-workbench/\SsixRepositoryCommit/code/S6_SERIES_PUBLISH_II_01GN_prereading_certificate_v2.py}{\nolinkurl{S6_SERIES_PUBLISH_II_01GN_prereading_certificate_v2.py}} \\[3pt]
Hodge and tangent frame covariance and the expansion into two-by-two minors (\cref{cdp:st:E16}) & \href{https://raw.githubusercontent.com/KokunoYumeto/s6-complex-structure-workbench/\SsixRepositoryCommit/code/S6_SERIES_PUBLISH_II_01GN_postrelease_frame_certificate.py}{\nolinkurl{S6_SERIES_PUBLISH_II_01GN_postrelease_frame_certificate.py}} \\[3pt]
Weighted intersection identities, sphere orientation and quadratic Hessian determinant (\cref{night:P4R.6}) & \href{https://raw.githubusercontent.com/KokunoYumeto/s6-complex-structure-workbench/\SsixRepositoryCommit/code/S6_SERIES_PUBLISH_II_A_02YI_interpolation_certificate.py}{\nolinkurl{S6_SERIES_PUBLISH_II_A_02YI_interpolation_certificate.py}} \\[3pt]
Quadratic-fibre coordinate inverses and local intersection orientations (\cref{night:P4R.6}) & \href{https://raw.githubusercontent.com/KokunoYumeto/s6-complex-structure-workbench/\SsixRepositoryCommit/code/S6_SERIES_PUBLISH_II_A_02YI_A1_transport_postrelease.py}{\nolinkurl{S6_SERIES_PUBLISH_II_A_02YI_A1_transport_postrelease.py}} \\[3pt]
\end{longtable}}

\begin{thebibliography}{99}
\providecommand{\natexlab}[1]{#1}
\providecommand{\doi}[1]{\url{https://doi.org/#1}}
\bibitem[{Claude} and {L. Alp{\"o}ge}(2026)]{AC2026}
{Claude} and {L. Alp{\"o}ge}.

\newblock The (3,4,infinity) modular family of 2-tori, completed at its three
  special points, is a complex structure on {$S^6$}.

\newblock 108-page manuscript, 2026.

\newblock URL \url{https://alpo.ge/s6.pdf}.

\bibitem[{Claude} and {L. Alp{\"o}ge}(2026{\natexlab{b}})]{AC26}
{Claude} and {L. Alp{\"o}ge}.

\newblock The $(3,4,\infty)$ modular family of 2-tori, completed at its three
  special points, is a complex structure on $s^6$.

\newblock Attached original manuscript, 108 pages, 2026{\natexlab{b}}.

\newblock URL \url{https://alpo.ge/s6.pdf}.
\newblock \emph{Verified}.

\bibitem[{Claude} and Alp{\"o}ge(2026)]{AC26Root}
{Claude} and L.~Alp{\"o}ge.

\newblock The $(3,4,\infty)$ modular family of $2$-tori, completed at its three
  special points, is a complex structure on $s^6$, 2026.

\newblock URL \url{https://alpo.ge/s6.pdf}.
\newblock \emph{Verified}.

\bibitem[A'Campo(1975)]{ACampo1975}
Norbert A'Campo.

\newblock La fonction z\^eta d'une monodromie.

\newblock \emph{Commentarii Mathematici Helvetici}, 50:\penalty0 233--248,
  1975.

\newblock \doi{10.1007/BF02565748}.

\newblock URL
  \url{https://www.e-periodica.ch/cntmng?bot=1&pid=com-001%3A1975%3A50%3A%3A27}.
\newblock \emph{Verified}.

\bibitem[Abramovich et~al.(2002)Abramovich, Karu, Matsuki, and
  W{\l}odarczyk]{AKMW2002}
Dan Abramovich, Kalle Karu, Kenji Matsuki, and Jaros{\l}aw W{\l}odarczyk.

\newblock Torification and factorization of birational maps.

\newblock \emph{Journal of the American Mathematical Society}, 15\penalty0
  (3):\penalty0 531--572, 2002.

\newblock URL \url{https://arxiv.org/pdf/math/9904135}.
\newblock \emph{Verified}.

\bibitem[Abasheva and D{\'e}ev(2024)]{AbashevaDeev2024AB}
Anna Abasheva and Rodion D{\'e}ev.

\newblock Complex surfaces with many algebraic structures.

\newblock \emph{International Mathematics Research Notices}, 2024\penalty0
  (9):\penalty0 7379--7400, 2024.

\newblock \doi{10.1093/imrn/rnad190}.

\newblock URL \url{https://arxiv.org/pdf/2303.10764}.
\newblock \emph{Verified}.

\bibitem[Albanese(2020)]{Albanese2020}
Michael Albanese.

\newblock The {Yamabe} invariants of {Inoue} surfaces, {Kodaira} surfaces, and
  their blowups, 2020.

\newblock URL \url{https://arxiv.org/pdf/2005.09494}.
\newblock \emph{Verified}.

\bibitem[Buonerba et~al.(2019)Buonerba, Bogomolov, and Kurnosov]{BBK2019}
Federico Buonerba, Fedor Bogomolov, and Nikon Kurnosov.
\newblock Classifying {VII}$_0$ surfaces with $b_2=0$ via group theory.
\newblock arXiv:1709.00062v3, 20 May 2019.
\newblock URL \url{https://arxiv.org/pdf/1709.00062v3}.


\bibitem[Barth et~al.(2004)Barth, Hulek, Peters, and {Van de Ven}]{BHPV2004E}
Wolf~P. Barth, Klaus Hulek, Chris A.~M. Peters, and Antonius {Van de Ven}.

\newblock \emph{Compact Complex Surfaces}, volume~4 of \emph{Ergebnisse der
  Mathematik und ihrer Grenzgebiete. 3. Folge}.

\newblock Springer, Berlin, Heidelberg, 2 edition, 2004.

\newblock \doi{10.1007/978-3-642-57739-0}.

\newblock URL \url{https://link.springer.com/book/10.1007/978-3-642-57739-0}.
\newblock \emph{Verified}.

\bibitem[Biswas et~al.(2012)Biswas, Mj, and Seshadri]{BMS2012}
Indranil Biswas, Mahan Mj, and Harish Seshadri.

\newblock 3-manifold groups, {K\"ahler} groups and complex surfaces.

\newblock \emph{Communications in Contemporary Mathematics}, 14\penalty0
  (6):\penalty0 1250038, 2012.

\newblock \doi{10.1142/S0219199712500381}.

\newblock URL \url{https://mathweb.tifr.res.in/~mahan/kahler3fld.pdf}.
\newblock Cited author preprint: 27 October 2011; preprint page and statement
locators refer to the linked author edition.

\bibitem[B{\u a}nic{\u a} and St{\u a}n{\u a}{\c s}il{\u a}(1976)]{BS76push}
Constantin B{\u a}nic{\u a} and Octavian St{\u a}n{\u a}{\c s}il{\u a}.

\newblock \emph{Algebraic Methods in the Global Theory of Complex Spaces}.

\newblock Editura Academiei and John Wiley and Sons, 1976.


\bibitem[Bierstone and Milman(1997)]{BierstoneMilman1997}
Edward Bierstone and Pierre~D. Milman.

\newblock Canonical desingularization in characteristic zero by blowing up the
  maximum strata of a local invariant.

\newblock \emph{Inventiones Mathematicae}, 128:\penalty0 207--302, 1997.

\newblock \doi{10.1007/s002220050141}.

\newblock URL \url{https://arxiv.org/pdf/alg-geom/9508005v2}.
\newblock Cited preprint edition: arXiv:alg-geom/9508005v2,
16 August 1995; preprint page locators refer to this edition.

\bibitem[Bingener(1980)]{Bingener1980}
J\"urgen Bingener.

\newblock Darstellbarkeitskriterien f\"ur analytische funktoren.

\newblock \emph{Annales scientifiques de l'\'Ecole Normale Sup\'erieure},
  13\penalty0 (3):\penalty0 317--347, 1980.

\newblock \doi{10.24033/asens.1385}.
\newblock \emph{Verified}.

\bibitem[Birkenhake and Lange(1999)]{BirkenhakeLange1999}
Christina Birkenhake and Herbert Lange.

\newblock \emph{Complex Tori}, volume 177 of \emph{Progress in Mathematics}.

\newblock Birkh\"auser, 1999.

\newblock \doi{10.1007/978-1-4612-1566-0}.
\newblock \emph{Verified}.

\bibitem[Fern\'andez~de Bobadilla and Koll\'ar(2012)]{BobadillaKollar2012}
Javier Fern\'andez~de Bobadilla and J\'anos Koll\'ar.

\newblock Homotopically trivial deformations.

\newblock \emph{Journal of Singularities}, 5:\penalty0 85--93, 2012.

\newblock \doi{10.5427/jsing.2012.5f}.

\newblock URL \url{https://www.journalofsing.org/volume5/bobadilla-kollar.pdf}.

\bibitem[Borcea(1984)]{Borcea1984}
Ciprian Borcea.

\newblock Moduli for {Kodaira} surfaces.

\newblock \emph{Compositio Mathematica}, 52\penalty0 (3):\penalty0 373--380,
  1984.

\newblock URL \url{https://www.numdam.org/article/CM_1984__52_3_373_0.pdf}.


\bibitem[Br{\^i}nz{\u a}nescu and Moraru(2005)]{BrinzanescuMoraru2005}
Vasile Br{\^i}nz{\u a}nescu and Ruxandra Moraru.

\newblock Holomorphic rank-2 vector bundles on non-{K\"ahler} elliptic
  surfaces.

\newblock \emph{Annales de l'Institut Fourier}, 55\penalty0 (5):\penalty0
  1659--1683, 2005.

\newblock \doi{10.5802/aif.2135}.

\newblock URL \url{https://arxiv.org/pdf/math/0306191}.
\newblock \emph{Verified}.

\bibitem[Brown(1982)]{Brown1982}
Kenneth~S. Brown.

\newblock \emph{Cohomology of Groups}, volume~87 of \emph{Graduate Texts in
  Mathematics}.

\newblock Springer, New York, 1982.

\newblock \doi{10.1007/978-1-4684-9327-6}.
\newblock \emph{Verified}.

\bibitem[Buchdahl(1999)]{Buchdahl1999}
Nicholas Buchdahl.

\newblock On compact {K\"ahler} surfaces.

\newblock \emph{Annales de l'Institut Fourier}, 49\penalty0 (1):\penalty0
  287--302, 1999.

\newblock \doi{10.5802/aif.1674}.

\newblock URL \url{https://www.numdam.org/article/AIF_1999__49_1_287_0.pdf}.

\bibitem[Campana et~al.(2020{\natexlab{a}})Campana, Demailly, and
  Peternell]{CDP20}
Fr{\'e}d{\'e}ric Campana, Jean-Pierre Demailly, and Thomas Peternell.

\newblock The algebraic dimension of compact complex threefolds with vanishing
  second betti number.

\newblock \emph{Compositio Mathematica}, 156\penalty0 (4):\penalty0 679--696,
  2020{\natexlab{a}}.

\newblock \doi{10.1112/S0010437X19007802}.
\newblock \emph{Verified}.

\bibitem[Campana et~al.(2020)Campana, Demailly, and Peternell]{CDP2020}
Fr{\'e}d{\'e}ric Campana, Jean-Pierre Demailly, and Thomas Peternell.

\newblock The algebraic dimension of compact complex threefolds with vanishing
  second betti number.

\newblock \emph{Compositio Mathematica}, 156\penalty0 (4):\penalty0 679--696,
  2020.

\newblock \doi{10.1112/S0010437X19007802}.
\newblock \emph{Verified}.

\bibitem[Campana et~al.(2020{\natexlab{c}})Campana, Demailly, and
  Peternell]{CDP2020D}
Fr{\'e}d{\'e}ric Campana, Jean-Pierre Demailly, and Thomas Peternell.

\newblock The algebraic dimension of compact complex threefolds with vanishing
  second betti number.

\newblock \emph{Compositio Mathematica}, 156\penalty0 (4):\penalty0 679--696,
  2020{\natexlab{c}}.

\newblock \doi{10.1112/S0010437X19007802}.

\newblock URL \url{https://arxiv.org/pdf/1904.11179}.
\newblock \emph{Verified}.

\bibitem[Campana et~al.(1998)Campana, Demailly, and Peternell]{CDP98}
Fr{\'e}d{\'e}ric Campana, Jean-Pierre Demailly, and Thomas Peternell.

\newblock The algebraic dimension of compact complex threefolds with vanishing
  second betti number.

\newblock \emph{Compositio Mathematica}, 112\penalty0 (1):\penalty0 77--91,
  1998.

\newblock \doi{10.1023/A:1000313214795}.

\bibitem[Cartan(1957)]{Cartan1957}
Henri Cartan.

\newblock Vari\'et\'es analytiques r\'eelles et vari\'et\'es analytiques
  complexes.

\newblock \emph{Bulletin de la Soci\'et\'e Math\'ematique de France},
  85:\penalty0 77--99, 1957.

\newblock \doi{10.24033/bsmf.1481}.

\newblock URL \url{https://www.numdam.org/article/BSMF_1957__85__77_0.pdf}.

\bibitem[Catanese(2002)]{Catanese2002}
Fabrizio Catanese.

\newblock Deformation types of real and complex manifolds, 2002.

\newblock URL \url{https://arxiv.org/pdf/math/0111245}.
\newblock \emph{Verified}.

\bibitem[Catanese(2004)]{Catanese2004}
Fabrizio Catanese.

\newblock Deformation in the large of some complex manifolds, {I}.

\newblock \emph{Annali di Matematica Pura ed Applicata}, 183:\penalty0
  261--289, 2004.

\newblock \doi{10.1007/s10231-003-0096-y}.

\newblock URL
  \url{https://www.mathe8.uni-bayreuth.de/de/team/prof-fabrizio-catanese-old/pdf/085.pdf}.
\newblock \emph{Verified}.

\bibitem[Cox et~al.(2011)Cox, Little, and Schenck]{CoxLittleSchenck2011}
David~A. Cox, John~B. Little, and Hal Schenck.

\newblock \emph{Toric Varieties}, volume 124 of \emph{Graduate Studies in
  Mathematics}.

\newblock American Mathematical Society, Providence, RI, 2011.

\newblock \doi{10.1090/gsm/124}.
\newblock Page locators refer to the authors' 2010 text, whose preface is
  dated September 2010.
\newblock URL \url{https://dacox.people.amherst.edu/toric.html}.

\bibitem[Du et~al.(2026)Du, Guo, Kusakabe, and Xie]{DGKXv1}
Yun-Heng Du, Bin Guo, Yuta Kusakabe, and Song-Yan Xie.

\newblock Oka geometry and deformations of torus-fibred complex six-spheres,
  2026.

\newblock URL \url{https://arxiv.org/abs/2610.00250v1}.
\newblock \emph{Verified}.

\bibitem[{National Institute of Standards and Technology}(2026)]{DLMFTheta2026}
{National Institute of Standards and Technology}.

\newblock {NIST Digital Library of Mathematical Functions}, chapters 20 and 23,
  2026.

\newblock URL \url{https://dlmf.nist.gov/}.
\newblock \emph{Verified}.

\bibitem[Dloussky et~al.(2000)Dloussky, Oeljeklaus, and Toma]{DOT2000}
Georges Dloussky, Karl Oeljeklaus, and Matei Toma.

\newblock Surfaces de la classe {VII}$_0$ admettant un champ de vecteurs.

\newblock \emph{Commentarii Mathematici Helvetici}, 75\penalty0 (2):\penalty0
  255--270, 2000.

\newblock \doi{10.1007/PL00000374}.

\newblock URL \url{https://ems.press/content/serial-article-files/42865?nt=1}.
\newblock \emph{Verified}.

\bibitem[Dloussky et~al.(2001)Dloussky, Oeljeklaus, and Toma]{DOT2001}
Georges Dloussky, Karl Oeljeklaus, and Matei Toma.

\newblock Surfaces de la classe {VII}$_0$ admettant un champ de vecteurs, {II}.

\newblock \emph{Commentarii Mathematici Helvetici}, 76\penalty0 (4):\penalty0
  640--664, 2001.

\newblock \doi{10.1007/s00014-001-8323-x}.

\newblock URL \url{https://ems.press/content/serial-article-files/42912?nt=1}.
\newblock \emph{Verified}.

\bibitem[Dloussky et~al.(2003)Dloussky, Oeljeklaus, and Toma]{DOT2003}
Georges Dloussky, Karl Oeljeklaus, and Matei Toma.

\newblock Class {$\mathrm{VII}_0$} surfaces with {$b_2$} curves.

\newblock \emph{Tohoku Mathematical Journal, Second Series}, 55\penalty0
  (2):\penalty0 283--309, 2003.

\newblock \doi{10.2748/tmj/1113246942}.

\newblock URL
  \url{https://www.jstage.jst.go.jp/article/tmj/55/2/55_2_283/_pdf/-char/en}.
\newblock \emph{Verified}.

\bibitem[de~Jong and van Straten(1990)]{DeJongVanStraten1990Normalization}
Theo de~Jong and Duco van Straten.

\newblock Deformations of the normalization of hypersurfaces.

\newblock \emph{Mathematische Annalen}, 288:\penalty0 527--547, 1990.

\newblock \doi{10.1007/BF01444547}.

\newblock URL
  \url{https://download.uni-mainz.de/mathematik/Algebraische%20Geometrie/DvS-Publikationen/Deformations%20of%20the%20normalization%20of%20hypersurfaces%20-%201990.pdf}.
\newblock \emph{Verified}.

\bibitem[Demailly(2012)]{Demailly2012E}
Jean-Pierre Demailly.

\newblock Complex analytic and differential geometry, 2012.

\newblock URL
  \url{https://www-fourier.univ-grenoble-alpes.fr/~demailly/manuscripts/agbook.pdf}.
\newblock \emph{Verified}.

\bibitem[Demailly(2002)]{DemaillyHodgeAuthor2002}
Jean-Pierre Demailly.

\newblock $l^2$ hodge theory and vanishing theorems.

\newblock In \emph{Introduction to Hodge Theory}, volume~8 of \emph{SMF/AMS
  Texts and Monographs}. American Mathematical Society and Soci\'et\'e
  Math\'ematique de France, 2002.

\newblock URL
  \url{https://www-fourier.univ-grenoble-alpes.fr/~demailly/manuscripts/hodge-ams.pdf}.

\newblock Author-hosted PDF body: Theorem 3.15, p. 15 (physical PDF p. 27), and
  Theorems 7.1--7.2, p. 32 (physical PDF p. 44);\allowbreak{} opened 19 November 2026.
  Scope: classical Green and Hodge inputs. The receiving finite-kernel
  analytic-minor argument is derived separately.
\newblock \emph{Verified}.

\bibitem[Dimca(2004)]{Dimca2004}
Alexandru Dimca.

\newblock \emph{Sheaves in Topology}.

\newblock Universitext. Springer, 2004.

\newblock \doi{10.1007/978-3-642-18868-8}.
\newblock \emph{Verified}.

\bibitem[Dloussky(2006)]{Dloussky2006}
Georges Dloussky.

\newblock On surfaces of class {$\mathrm{VII}_0^+$} with numerically
  anticanonical divisor.

\newblock \emph{American Journal of Mathematics}, 128\penalty0 (3):\penalty0
  639--670, 2006.

\newblock \doi{10.1353/ajm.2006.0021}.

\newblock URL \url{https://arxiv.org/pdf/math/0406387}.
\newblock \emph{Verified}.

\bibitem[Dloussky and Teleman(2012)]{DlousskyTeleman2012}
Georges Dloussky and Andrei Teleman.

\newblock Infinite bubbling in non-{K\"ahlerian} geometry.

\newblock \emph{Mathematische Annalen}, 353:\penalty0 1283--1314, 2012.

\newblock \doi{10.1007/s00208-011-0713-9}.

\newblock URL \url{https://arxiv.org/pdf/1012.5247}.
\newblock \emph{Verified}.

\bibitem[Douady(1966)]{Douady1966}
Adrien Douady.

\newblock Le probl\`eme des modules pour les sous-espaces analytiques compacts
  d'un espace analytique donn\'e.

\newblock \emph{Annales de l'Institut Fourier}, 16\penalty0 (1):\penalty0
  1--95, 1966.

\newblock \doi{10.5802/aif.226}.

\newblock URL \url{https://www.numdam.org/item/AIF_1966__16_1_1_0/}.
\newblock \emph{Verified}.

\bibitem[Engel(2026)]{Engel2026}
Philip Engel.

\newblock Complex structures on {$S^6$}.

\newblock arXiv:2609.38442v1, 2026.

\newblock URL \url{https://arxiv.org/abs/2609.38442}.
\newblock \emph{Verified}.

\bibitem[Fujiki and Nakano(1971)]{FujikiNakano1971}
Akira Fujiki and Shigeo Nakano.

\newblock Supplement to {On the Inverse of Monoidal Transformation}.

\newblock \emph{Publications of the Research Institute for Mathematical
  Sciences}, 7\penalty0 (3):\penalty0 637--644, 1971.

\newblock \doi{10.2977/PRIMS/1195193401}.

\newblock URL \url{https://ems.press/content/serial-article-files/41704}.
\newblock \emph{Verified}.

\bibitem[Fujimoto(1988)]{Fujimoto1988Log}
Yoshio Fujimoto.

\newblock Logarithmic transformations on elliptic fiber spaces.

\newblock \emph{Journal of Mathematics of Kyoto University}, 28\penalty0
  (1):\penalty0 91--110, 1988.

\newblock URL
  \url{https://scispace.com/pdf/logarithmic-transformations-on-elliptic-fiber-spaces-1uinrvp8jj.pdf}.
\newblock \emph{Verified}.

\bibitem[Fujino(2021)]{Fujino2021E}
Osamu Fujino.

\newblock Minimal model theory for log surfaces in {Fujiki}'s class {$\mathcal
  C$}.

\newblock \emph{Nagoya Mathematical Journal}, 244:\penalty0 256--282, 2021.

\newblock \doi{10.1017/nmj.2020.14}.

\newblock URL
  \url{https://www.cambridge.org/core/journals/nagoya-mathematical-journal/article/minimal-model-theory-for-log-surfaces-in-fujikis-class-mathcalc/86AAC54D7E1F3B189DE91969A584D16E}.
\newblock \emph{Verified}.

\bibitem[Furushima(1993)]{Furushima1993}
Mikio Furushima.

\newblock The complete classification of compactifications of {$\mathbb C^3$}
  which are projective manifolds with the second betti number one.

\newblock Technical Report MPI/93-6, Max-Planck-Institut f\"ur Mathematik,
  1993.

\newblock URL
  \url{https://archive.mpim-bonn.mpg.de/id/eprint/2170/1/preprint_1993_6.pdf}.
\newblock \emph{Verified}.

\bibitem[Gross et~al.(2015)Gross, Hacking, and Keel]{GHK2015}
Mark Gross, Paul Hacking, and Sean Keel.

\newblock Mirror symmetry for log {Calabi--Yau} surfaces {I}.

\newblock \emph{Publications Math\'ematiques de l'IH\'ES}, 122:\penalty0
  65--168, 2015.

\newblock \doi{10.1007/s10240-015-0073-1}.

\newblock URL \url{https://www.numdam.org/articles/10.1007/s10240-015-0073-1/}.
\newblock \emph{Verified}.

\bibitem[Granger and Schulze(2014)]{GS2014}
Michel Granger and Mathias Schulze.

\newblock Normal crossing properties of complex hypersurfaces via logarithmic
  residues.

\newblock \emph{Compositio Mathematica}, 150\penalty0 (9):\penalty0 1607--1622,
  2014.

\newblock \doi{10.1112/S0010437X13007860}.

\newblock URL \url{https://arxiv.org/pdf/1109.2612}.
\newblock \emph{Verified}.

\bibitem[Grauert(1960{\natexlab{a}})]{Grauert1960D}
Hans Grauert.

\newblock Ein theorem der analytischen garbentheorie und die modulräume
  komplexer strukturen.

\newblock \emph{Publications Mathématiques de l'IHÉS}, 5:\penalty0 5--64,
  1960{\natexlab{a}}.

\newblock \doi{10.1007/BF02684746}.

\newblock URL \url{https://www.numdam.org/article/PMIHES_1960__5__5_0.pdf}.

\newblock Reopened in this continuation on19 November2026: Section7 Satz1,
  printed p60 (Remmert proper-image theorem), with the reduced-space and
  proper-holomorphic hypotheses matched to the compact meromorphic graph.
  Verified in this writing on 18 November 2026 at these loci only. Section6
  HauptsatzeI,II,IIap59;\allowbreak{}Section7 Satz3p62,Satz5p64. Scope: Coherence, analytic
  formal functions, upper semicontinuity and constant-dimension base change,
  with flatness hypothesis retained Section 7, Satz 5, pp. 63–64. Scope:
  Opened full Satz 5. Also read corrigendum 1963 pp. 35–36
  (https://www.numdam.org/item/10.1007/BF02684291.pdf);\allowbreak{} it does not alter this
  statement. Author/title and edition matched to the opened item. Verification
  applies only to its recorded locator and edition;\allowbreak{} aliasing does not transfer
  preprint pagination to a published edition.
\newblock \emph{Verified}.

\bibitem[Grauert(1962)]{Grauert1962E}
Hans Grauert.

\newblock {\"U}ber modifikationen und exzeptionelle analytische mengen.

\newblock \emph{Mathematische Annalen}, 146:\penalty0 331--368, 1962.

\newblock \doi{10.1007/BF01441136}.

\newblock URL \url{https://link.springer.com/article/10.1007/BF01441136}.
\newblock \emph{Verified}.

\bibitem[Grauert(1960{\natexlab{b}})]{Grauert60}
Hans Grauert.

\newblock Ein theorem der analytischen garbentheorie und die modulräume
  komplexer strukturen.

\newblock \emph{Publications Mathématiques de l'IHÉS}, 5:\penalty0 5--64,
  1960{\natexlab{b}}.

\newblock \doi{10.1007/BF02684746}.

\newblock URL \url{https://www.numdam.org/article/PMIHES_1960__5__5_0.pdf}.
\newblock \emph{Verified}.

\bibitem[Greb(2011)]{Greb2011}
Daniel Greb.

\newblock Rational singularities and quotients by holomorphic group actions.

\newblock \emph{Annali della Scuola Normale Superiore di Pisa, Classe di
  Scienze}, 10\penalty0 (2):\penalty0 413--426, 2011.

\newblock URL \url{https://www.numdam.org/item/ASNSP_2011_5_10_2_413_0.pdf}.
\newblock \emph{Verified}.

\bibitem[Gross and Siebert(2006)]{GrossSiebert2006I}
Mark Gross and Bernd Siebert.

\newblock Mirror symmetry via logarithmic degeneration data {I}.

\newblock \emph{Journal of Differential Geometry}, 72\penalty0 (2):\penalty0
  169--338, 2006.

\newblock URL \url{https://arxiv.org/pdf/math/0309070}.
\newblock \emph{Verified}.

\bibitem[Guillot(2019)]{Guillot2019}
Adolfo Guillot.

\newblock Meromorphic vector fields with single-valued solutions on complex
  surfaces.

\newblock \emph{Advances in Mathematics}, 354:\penalty0 106742, 2019.

\newblock \doi{10.1016/j.aim.2019.106742}.

\newblock URL \url{https://arxiv.org/pdf/1603.02288}.
\newblock \emph{Verified}.

\bibitem[Guillot(2024)]{Guillot2024}
Adolfo Guillot.

\newblock On the singularities of complete holomorphic vector fields in
  dimension two.

\newblock In \emph{Handbook of Geometry and Topology of Singularities VI:
  Foliations}. Springer, 2024.

\newblock \doi{10.1007/978-3-031-54172-8_1}.

\newblock URL
  \url{https://www.matem.unam.mx/~guillot/guillot-singularities.pdf}.
\newblock \emph{Verified}.

\bibitem[Huckleberry et~al.(2000)Huckleberry, Kebekus, and Peternell]{HKP2000}
Alan~T. Huckleberry, Stefan Kebekus, and Thomas Peternell.

\newblock Group actions on {$S^6$} and complex structures on {$\mathbb P^3$}.

\newblock \emph{Duke Mathematical Journal}, 102\penalty0 (1):\penalty0
  101--124, 2000.

\newblock \doi{10.1215/S0012-7094-00-10214-1}.

\newblock URL \url{https://arxiv.org/abs/math/9812076v1}.
\newblock Cited preprint edition: arXiv:math/9812076v1, 13 December 1998;
preprint page locators refer to this edition.

\bibitem[Hatcher(2002{\natexlab{a}})]{Hatcher2002}
Allen Hatcher.

\newblock \emph{Algebraic Topology}.

\newblock Cambridge University Press, 2002{\natexlab{a}}.

\newblock URL \url{https://pi.math.cornell.edu/~hatcher/AT/AT.pdf}.
\newblock \emph{Verified}.

\bibitem[Honda and Viaclovsky(2025)]{HondaViaclovsky2025D}
Nobuhiro Honda and Jeff Viaclovsky.

\newblock Some remarks on fibrations in complex geometry.

\newblock \emph{Journal of Mathematical Study}, 58\penalty0 (4):\penalty0
  552--574, 2025.

\newblock \doi{10.4208/jms.v58n4.25.11}.

\newblock URL \url{https://arxiv.org/abs/2504.14420}.
\newblock \emph{Verified}.

\bibitem[Horikawa(1974)]{Horikawa1974}
Eiji Horikawa.

\newblock On deformations of holomorphic maps. {II}.

\newblock \emph{Journal of the Mathematical Society of Japan}, 26\penalty0
  (4):\penalty0 647--667, 1974.

\newblock \doi{10.2969/jmsj/02640647}.

\newblock URL
  \url{https://www.jstage.jst.go.jp/article/jmath1948/26/4/26_4_647/_pdf/-char/en}.
\newblock \emph{Verified}.

\bibitem[Huybrechts(2016)]{HuybrechtsK3Draft}
Daniel Huybrechts.

\newblock \emph{Lectures on {K3} Surfaces}.

\newblock Cambridge University Press, 2016.

\newblock URL \url{https://www.math.uni-bonn.de/people/huybrech/K3Global.pdf}.
\newblock \emph{Verified}.

\bibitem[Inoue(1974)]{Inoue1974}
Masahisa Inoue.

\newblock On surfaces of class {$VII_0$}.

\newblock \emph{Inventiones Mathematicae}, 24\penalty0 (4):\penalty0 269--310,
  1974.

\newblock \doi{10.1007/BF01425563}.



\bibitem[Kajiwara et~al.(2008)Kajiwara, Kato, and Nakayama]{KKN2008}
Takeshi Kajiwara, Kazuya Kato, and Chikara Nakayama.

\newblock Logarithmic abelian varieties, part {I}: Complex analytic theory.

\newblock \emph{Journal of Mathematical Sciences, The University of Tokyo},
  15:\penalty0 69--193, 2008.

\newblock URL \url{https://www.ms.u-tokyo.ac.jp/journal/pdf/jms150104.pdf}.
\newblock \emph{Verified}.

\bibitem[Kamenova(n.d.)]{KamenovaTori}
Ljudmila Kamenova.

\newblock Degenerations of 2-dimensional tori, n.d.

\newblock URL
  \url{https://www.math.stonybrook.edu/~kamenova/homepage_files/degen.pdf}.
\newblock \emph{Verified}.

\bibitem[Kato(1975)]{Kato1975Hopf}
Masahide Kato.

\newblock Topology of {Hopf} surfaces.

\newblock \emph{Journal of the Mathematical Society of Japan}, 27\penalty0
  (2):\penalty0 222--238, 1975.

\newblock \doi{10.2969/jmsj/02720222}.

\newblock URL
  \url{https://www.jstage.jst.go.jp/article/jmath1948/27/2/27_2_222/_pdf/-char/en}.
\newblock \emph{Verified}.

\bibitem[Kawai(1969)]{Kawai1969}
Soichi Kawai.

\newblock On compact complex analytic manifolds of complex dimension 3, {II}.

\newblock \emph{Journal of the Mathematical Society of Japan}, 21\penalty0
  (4):\penalty0 604--616, 1969.

\newblock \doi{10.2969/jmsj/02140604}.

\newblock URL
  \url{https://www.jstage.jst.go.jp/article/jmath1948/21/4/21_4_604/_pdf/-char/en}.

\bibitem[Khaled(2024)]{Khaled2024Hopf}
Zahraa Khaled.

\newblock Real structures on primary {Hopf} surfaces.

\newblock \emph{Complex Manifolds}, 11\penalty0 (1), 2024.

\newblock \doi{10.1515/coma-2024-0007}.

\newblock URL \url{https://arxiv.org/pdf/2310.05265}.
\newblock \emph{Verified}.

\bibitem[Khaled and Teleman(2026{\natexlab{a}})]{KhaledTeleman2026}
Zahraa Khaled and Andrei Teleman.

\newblock On the classification of {Inoue} surfaces, 2026{\natexlab{a}}.

\newblock Published in \emph{Journal of Mathematical Sciences, The University
of Tokyo}, 33(3):283--357, 2026.
\newblock Cited preprint edition: arXiv:2406.15158v4, 13 July 2026;
preprint page and statement locators refer to this edition.
\newblock URL \url{https://arxiv.org/pdf/2406.15158v4}.

\bibitem[Khaled and Teleman(2026{\natexlab{b}})]{KhaledTeleman2026Journal}
Zahraa Khaled and Andrei Teleman.

\newblock On the classification of {Inoue} surfaces.

\newblock \emph{Journal of Mathematical Sciences, The University of Tokyo},
  33\penalty0 (3):\penalty0 283--357, 2026{\natexlab{b}}.

\newblock URL \url{https://www.ms.u-tokyo.ac.jp/journal/33-3-1.pdf}.


\bibitem[Kodaira(1954)]{Kodaira1954}
Kunihiko Kodaira.
\newblock On K\"ahler varieties of restricted type.
\newblock \emph{Proceedings of the National Academy of Sciences of the
United States of America}, 40\penalty0 (5):\penalty0 313--316, 1954.
\newblock \doi{10.1073/pnas.40.5.313}.


\bibitem[Kodaira(1963)]{Kodaira1963II}
Kunihiko Kodaira.

\newblock On compact analytic surfaces: {II}.

\newblock \emph{Annals of Mathematics}, 77\penalty0 (3):\penalty0 563--626,
  1963.

\newblock \doi{10.2307/1970131}.

\newblock URL \url{https://www.jstor.org/stable/1970131}.
\newblock \emph{Verified}.

\bibitem[Kulikov(1977)]{Kulikov1977}
Viktor~S. Kulikov.

\newblock Degenerations of {$K3$} surfaces and enriques surfaces.

\newblock \emph{Mathematics of the USSR-Izvestiya}, 11\penalty0 (5):\penalty0
  957--989, 1977.

\newblock \doi{10.1070/IM1977v011n05ABEH001753}.

\newblock URL
  \url{https://www.mathnet.ru/php/getFT.phtml?jrnid=im&option_lang=eng&paperid=1877&what=fullteng}.
\newblock \emph{Verified}.

\bibitem[Kulikov(1981)]{Kulikov1981}
Viktor~S. Kulikov.

\newblock On modifications of degenerations of surfaces with {$\varkappa=0$}.

\newblock \emph{Mathematics of the USSR-Izvestiya}, 17\penalty0 (2):\penalty0
  339--342, 1981.

\newblock \doi{10.1070/IM1981v017n02ABEH001361}.

\newblock URL
  \url{https://www.mathnet.ru/php/getFT.phtml?jrnid=im&option_lang=eng&paperid=1953&what=fullteng}.
\newblock \emph{Verified}.

\bibitem[Kurdyka(1998)]{Kurdyka1998}
Krzysztof Kurdyka.

\newblock On gradients of functions definable in o-minimal structures.

\newblock \emph{Annales de l'Institut Fourier}, 48\penalty0 (3):\penalty0
  769--783, 1998.

\newblock \doi{10.5802/aif.1638}.

\newblock URL \url{https://www.numdam.org/article/AIF_1998__48_3_769_0.pdf}.
\newblock \emph{Verified}.

\bibitem[Lehn et~al.(2019)Lehn, Rollenske, and Schinko]{LRS2019}
Christian Lehn, S{\"o}nke Rollenske, and Caren Schinko.

\newblock The complex geometry of a hypothetical complex structure on {$S^6$},
  2019.

\newblock URL \url{https://arxiv.org/abs/1912.09719}.
\newblock \emph{Verified}.

\bibitem[Lagrange(1868)]{Lagrange1770}
Joseph-Louis Lagrange.

\newblock D{\'e}monstration d'un th{\'e}or{\`e}me d'arithm{\'e}tique.

\newblock In \emph{{\OE}uvres de Lagrange}, volume III, pages 189--201.
  Gauthier-Villars, Paris, 1868.

\newblock URL
  \url{https://yuvalfilmus.cs.technion.ac.il/Manuscripts/Lagrange%20four%20square%20theorem.pdf}.
\newblock \emph{Verified}.

\bibitem[Lazzeri(1973)]{Lazzeri1973Root}
Fulvio Lazzeri.

\newblock A theorem on the monodromy of isolated singularities.

\newblock In \emph{Singularit{\'e}s {\`a} Carg{\`e}se}, volume 7--8 of
  \emph{Ast{\'e}risque}, pages 269--275. Soci{\'e}t{\'e} Math{\'e}matique de
  France, 1973.

\newblock URL \url{https://numdam.org/item/AST_1973__7-8__269_0/}.
\newblock \emph{Verified}.

\bibitem[Liu and Rollenske(2012)]{LiuRollenske2012SLC}
Wenfei Liu and S{\"o}nke Rollenske.
\newblock Two-dimensional semi-log-canonical hypersurfaces.
\newblock \emph{Le Matematiche}, 67(2):185--202, 2012.
\newblock doi:\href{https://doi.org/10.4418/2012.67.2.14}{10.4418/2012.67.2.14}.

\bibitem[Liu and Rollenske(2026)]{LiuRollenske2026}
Wenfei Liu and S{\"o}nke Rollenske.

\newblock Yet another family of complex structures on {$S^6$}.

\newblock arXiv:2609.33826v2, 2026.

\newblock URL \url{https://arxiv.org/abs/2609.33826v2}.

\bibitem[{\L}ojasiewicz(1964)]{Lojasiewicz1964}
Stanis{\l}aw {\L}ojasiewicz.

\newblock Triangulation of semi-analytic sets.

\newblock \emph{Annali della Scuola Normale Superiore di Pisa, Classe di
  Scienze}, ser.~3, 18\penalty0 (4):\penalty0 449--474, 1964.

\newblock URL \url{https://www.numdam.org/item/ASNSP_1964_3_18_4_449_0.pdf}.

\bibitem[Michel et~al.(2005)Michel, Pichon, and Weber]{MPW2005Hirzebruch}
Fran{\c c}oise Michel, Anne Pichon, and Claude Weber.

\newblock The boundary of the {Milnor} fiber of {Hirzebruch} surface
  singularities, 2005.

\newblock URL \url{https://www.maths.ed.ac.uk/~v1ranick/papers/mpw.pdf}.
\newblock \emph{Verified}.

\bibitem[Meersseman(2019)]{Meersseman2019E}
Laurent Meersseman.

\newblock The {Teichm\"uller} and {Riemann} moduli stacks.

\newblock \emph{Journal de l'\'Ecole polytechnique---Math\'ematiques},
  6:\penalty0 879--945, 2019.

\newblock \doi{10.5802/jep.108}.

\newblock URL \url{https://www.numdam.org/item/10.5802/jep.108.pdf}.
\newblock \emph{Verified}.

\bibitem[Merle and Teissier(1976--1977)]{MerleTeissier1977}
Michel Merle and Bernard Teissier.

\newblock Conditions d'adjonction, d'apr\`es {Du Val}.

\newblock \emph{S\'eminaire sur les singularit\'es des surfaces}, pages 1--16,
  1976--1977.

\newblock URL \url{https://www.numdam.org/item/SSS_1976-1977____A16_0.pdf}.
\newblock \emph{Verified}.

\bibitem[Milnor(1968)]{Milnor1968}
John Milnor.

\newblock \emph{Singular Points of Complex Hypersurfaces}, volume~61 of
  \emph{Annals of Mathematics Studies}.

\newblock Princeton University Press and University of Tokyo Press, 1968.

\bibitem[Mumford(1961)]{Mumford1961}
David Mumford.

\newblock The topology of normal singularities of an algebraic surface and a
  criterion for simplicity.

\newblock \emph{Publications Mathématiques de l'IHÉS}, 9:\penalty0 5--22,
  1961.

\newblock \doi{10.1007/BF02698717}.

\newblock URL \url{https://www.numdam.org/item/PMIHES_1961__9__5_0.pdf}.

\bibitem[Mumford(1972)]{Mumford1972}
David Mumford.

\newblock An analytic construction of degenerating abelian varieties over
  complete rings.

\newblock \emph{Compositio Mathematica}, 24\penalty0 (3):\penalty0 239--272,
  1972.

\newblock URL
  \url{https://www.dam.brown.edu/people/mumford/alg_geom/papers/1972b--AnalConstrDegAV-Numdam.pdf}.
\newblock \emph{Verified}.

\bibitem[Nakamura(2008)]{Nakamura2008VII}
Iku Nakamura.

\newblock Survey on {VII}$_0$ surfaces, 2008.

\newblock URL
  \url{https://www.math.sci.hokudai.ac.jp/~nakamura/70surfaces080306.pdf}.
\newblock \emph{Verified}.

\bibitem[Nakano(1970)]{Nakano1970}
Shigeo Nakano.

\newblock On the inverse of monoidal transformation.

\newblock \emph{Publications of the Research Institute for Mathematical
  Sciences}, 6\penalty0 (3):\penalty0 483--502, 1970.

\newblock \doi{10.2977/PRIMS/1195193917}.

\newblock URL \url{https://ems.press/content/serial-article-files/41677}.
\newblock \emph{Verified}.

\bibitem[N{\'e}methi(1991)]{Nemethi1991Join}
Andr{\'a}s N{\'e}methi.

\newblock Generalized local and global {Sebastiani--Thom} type theorems.

\newblock \emph{Compositio Mathematica}, 80\penalty0 (1):\penalty0 1--14, 1991.

\newblock URL \url{https://www.numdam.org/article/CM_1991__80_1_1_0.pdf}.
\newblock \emph{Verified}.

\bibitem[Ornea et~al.(2020)Ornea, Verbitsky, and Vuletescu]{OVV2020}
Liviu Ornea, Misha Verbitsky, and Victor Vuletescu.
\newblock Classification of non-{K}{\"a}hler surfaces and locally conformally
  {K}{\"a}hler geometry.
\newblock arXiv:1810.05768v3, 26 June 2020.
\newblock URL \url{https://arxiv.org/pdf/1810.05768v3}.


\bibitem[Papikian(2004)]{Papikian2004Tate}
Mihran Papikian.

\newblock Rigid analytic geometry and abelian varieties.

\newblock Author notes, 13 June 2004, 2004.

\newblock URL \url{https://math.stanford.edu/~vakil/snowbird/mihranjun21.pdf}.
\newblock \emph{Verified}.

\bibitem[Peters(1990)]{Peters1990}
Chris A.~M. Peters.

\newblock Compact complex surfaces, 1990.

\newblock URL
  \url{https://cpeters1.win.tue.nl/ConfsAndSchools/Canada/surface.pdf}.
\newblock \emph{Verified}.

\bibitem[Peters(2021)]{Peters2021}
Chris Peters.

\newblock On complex surfaces with definite intersection form.

\newblock \emph{New York Journal of Mathematics}, 27:\penalty0 840--847, 2021.

\newblock URL \url{https://nyjm.albany.edu/j/2021/27-32p.pdf}.

\bibitem[Piatetski-Shapiro and
  Shafarevich(1971)]{PiatetskiShapiroShafarevich1971}
I.~I. Piatetski-Shapiro and I.~R. Shafarevich.

\newblock A {Torelli} theorem for algebraic surfaces of type {K3}.

\newblock \emph{Mathematics of the USSR-Izvestiya}, 5\penalty0 (3):\penalty0
  547--588, 1971.

\newblock \doi{10.1070/IM1971v005n03ABEH001075}.

\newblock URL \url{https://www.mathnet.ru/eng/im2021}.
\newblock \emph{Verified}.

\bibitem[Ramis and Ruget(1970)]{RR70push}
Jean-Pierre Ramis and Gabriel Ruget.

\newblock Complexe dualisant et th{\'e}or{\`e}mes de dualit{\'e} en
  g{\'e}om{\'e}trie analytique complexe.

\newblock \emph{Publications Math{\'e}matiques de l'IH{\'E}S}, 38:\penalty0
  77--91, 1970.

\newblock \doi{10.1007/BF02684652}.

\newblock URL \url{https://www.numdam.org/item/PMIHES_1970__38__77_0.pdf}.
\newblock \emph{Verified}.

\bibitem[Reid(1994)]{Reid1994Nonnormal}
Miles Reid.

\newblock Nonnormal del {Pezzo} surfaces.

\newblock \emph{Publications of the Research Institute for Mathematical
  Sciences}, 30\penalty0 (5):\penalty0 695--727, 1994.

\newblock \doi{10.2977/prims/1195165581}.

\newblock URL \url{https://ems.press/content/serial-article-files/40537}.
\newblock \emph{Verified}.

\bibitem[Rossi(2017)]{Rossi2017}
Michele Rossi.

\newblock A small and non-simple geometric transition.

\newblock \emph{Mathematical Physics, Analysis and Geometry}, 20\penalty0
  (2):\penalty0 15, 2017.

\newblock \doi{10.1007/s11040-017-9247-z}.

\newblock URL \url{https://arxiv.org/pdf/1211.6329}.
\newblock \emph{Verified}.

\bibitem[Roulleau(2019)]{Roulleau2019GeneralizedKummer}
Xavier Roulleau.

\newblock On generalized Kummer surfaces and the orbifold Bogomolov--Miyaoka--Yau inequality.

\newblock \emph{Transactions of the American Mathematical Society}, 371(11):7651--7668, 2019.

\newblock \doi{10.1090/tran/7507}.
\newblock \emph{Verified}.

\bibitem[Schr{\"o}er and Siebert(2002)]{SS2002}
Stefan Schr{\"o}er and Bernd Siebert.

\newblock Irreducible degenerations of primary {Kodaira} surfaces.

\newblock In \emph{Complex Geometry: Collection of Papers Dedicated to Hans
  Grauert}, pages 193--222. Springer, 2002.

\newblock \doi{10.1007/978-3-642-56202-0_14}.
\newblock Cited preprint edition: arXiv:math/0001077v2, 28 October 2001;
preprint page and statement locators refer to this edition.
\newblock URL \url{https://arxiv.org/pdf/math/0001077v2}.

\bibitem[Schnell(n.d.)]{SchnellLecture3}
Christian Schnell.

\newblock Lecture 3: Cohomology of compact complex tori.

\newblock MAT615 lecture notes, Stony Brook University;\allowbreak{} lecture dated February
  4, year not specified in the document, n.d.

\newblock URL
  \url{https://www.math.stonybrook.edu/~cschnell/mat615/lectures/lecture3.pdf}.

\bibitem[Sch{\"u}tt and Shioda(2010)]{SchuttShioda2010}
Matthias Sch{\"u}tt and Tetsuji Shioda.

\newblock Elliptic surfaces.

\newblock \emph{Algebraic Geometry in East Asia---Seoul 2008}, 60:\penalty0
  51--160, 2010.

\newblock URL \url{https://arxiv.org/pdf/0907.0298}.
\newblock \emph{Verified}.

\bibitem[Serre(1956)]{Serre1956GAGA}
Jean-Pierre Serre.

\newblock G\'eom\'etrie alg\'ebrique et g\'eom\'etrie analytique.

\newblock \emph{Annales de l'Institut Fourier}, 6:\penalty0 1--42, 1956.

\newblock \doi{10.5802/aif.59}.

\bibitem[Siersma(1991)]{Siersma1991Variation}
Dirk Siersma.

\newblock Variation mappings on singularities with a 1-dimensional critical
  locus.

\newblock \emph{Topology}, 30\penalty0 (3):\penalty0 445--469, 1991.

\newblock \doi{10.1016/0040-9383(91)90025-Y}.

\newblock URL \url{https://www.maths.ed.ac.uk/~v1ranick/papers/siersma.pdf}.
\newblock \emph{Verified}.

\bibitem[Simoni\v{c}(2016)]{Simonic2016Hartogs}
Aleksander Simoni\v{c}.

\newblock Elementary approach to the {Hartogs} extension theorem, 2016.

\newblock URL \url{https://arxiv.org/abs/1608.00950}.

\newblock arXiv:1608.00950v1, 1 August 2016;\allowbreak{} Theorem 1, p. 1.
\newblock \emph{Verified}.

\bibitem[{The Stacks Project Authors}(2026)]{StacksProject}
{The Stacks Project Authors}.

\newblock The stacks project, 2026.

\newblock URL \url{https://stacks.math.columbia.edu}.

\bibitem[Steenbrink(1976)]{Steenbrink1976}
Joseph Steenbrink.

\newblock Limits of Hodge structures.

\newblock \emph{Inventiones Mathematicae}, 31:\penalty0 229--257, 1976.

\newblock \doi{10.1007/BF01403146}.


\bibitem[Steenbrink(1989)]{Steenbrink1989Spectrum}
Joseph H.~M. Steenbrink.

\newblock The spectrum of hypersurface singularities.

\newblock In \emph{Th\'eorie de Hodge---Luminy, Juin 1987}, number 179--180 in
  Ast\'erisque, pages 163--184. Soci\'et\'e math\'ematique de France, 1989.

\newblock URL
  \url{https://www.numdam.org/article/AST_1989__179-180__163_0.pdf}.
\newblock \emph{Verified}.

\bibitem[Steenbrink(1999)]{Steenbrink1999}
J.~H.~M. Steenbrink.

\newblock Cup products and mixed {Hodge} structures, 1999.

\newblock URL \url{https://arxiv.org/pdf/math/9902134}.
\newblock \emph{Verified}.

\bibitem[Steinberg(1965)]{Steinberg1965}
Robert Steinberg.

\newblock Regular elements of semi-simple algebraic groups.

\newblock \emph{Publications Math\'ematiques de l'IH\'ES}, 25:\penalty0 49--80,
  1965.

\newblock \doi{10.1007/BF02684397}.

\newblock URL \url{https://www.numdam.org/item/PMIHES_1965__25__49_0.pdf}.
\newblock \emph{Verified}.

\bibitem[Teleman(2006)]{Teleman2006}
Andrei Teleman.

\newblock The pseudo-effective cone of a non-{K}\"ahlerian surface and
  applications.

\newblock \emph{Mathematische Annalen}, 335\penalty0 (4):\penalty0 965--989,
  2006.

\newblock \doi{10.1007/s00208-006-0782-3}.

\newblock URL \url{https://arxiv.org/pdf/0704.2948}.
\newblock \emph{Verified}.

\bibitem[Teleman(2008)]{Teleman2008}
Andrei Teleman.

\newblock Gauge theoretical methods in the classification of
  non-{K}{\"a}hlerian surfaces, 2008.

\newblock URL
  \url{https://www.i2m.univ-amu.fr/perso/andrei.teleman/documents/postnikov.pdf}.
\newblock \emph{Verified}.

\bibitem[Ulmer and Urz{\'u}a(2021)]{UlmerUrzua2021}
Douglas Ulmer and Giancarlo Urz{\'u}a.

\newblock Bounding tangencies of sections on elliptic surfaces.

\newblock \emph{International Mathematics Research Notices}, 2021\penalty0
  (6):\penalty0 4768--4802, 2021.

\newblock \doi{10.1093/imrn/rnaa222}.

\newblock URL \url{https://dlulmer.github.io/research/papers/2021a.pdf}.
\newblock \emph{Verified}.

\bibitem[Viaclovsky(2026{\natexlab{a}})]{ViaTwo26push}
Jeff Viaclovsky.

\newblock A two-parameter family of complex structures on {$S^6$},
  2026{\natexlab{a}}.

\newblock URL \url{https://arxiv.org/abs/2609.33785}.
\newblock \emph{Verified}.

\bibitem[Viaclovsky(2026{\natexlab{b}})]{Viaclovsky2026Combs}
Jeff Viaclovsky.

\newblock Combs and complex structures on the six-sphere, 2026{\natexlab{b}}.

\newblock arXiv:2610.10451v1; manuscript dated 7 October 2026.

\newblock URL \url{https://arxiv.org/abs/2610.10451v1}.

\bibitem[Viaclovsky(2026{\natexlab{c}})]{Viaclovsky2026Log}
Jeff Viaclovsky.

\newblock Logarithmic transformations on the six-sphere, 2026{\natexlab{c}}.

\newblock arXiv:2610.39102v1, 29 October 2026; manuscript dated 28 October 2026.
\newblock \doi{10.48550/arXiv.2610.39102}.
\newblock URL \url{https://arxiv.org/abs/2610.39102v1}.
\newblock Cited edition: the version-one preprint; its Theorem~1.1 is on p.~1.


\bibitem[Viaclovsky(2026{\natexlab{d}})]{Viaclovsky2026Two}
Jeff Viaclovsky.

\newblock A two-parameter family of complex structures on {$S^6$},
  2026{\natexlab{d}}.

\newblock arXiv:2609.33785v1; manuscript dated 26 September 2026.

\newblock URL \url{https://arxiv.org/abs/2609.33785v1}.

\bibitem[Wang(2026)]{Wang2026}
Zichang Wang.

\newblock A complex structure on a rational homology six-sphere with
  two-torsion.

\newblock arXiv:2609.11541v2, 2026.

\newblock URL \url{https://arxiv.org/abs/2609.11541v2}.
\newblock Cited edition: version~2, 8 November 2026; page and statement locators refer to this edition.

\bibitem[Yoshihara(1995)]{Yoshihara1995}
Hisao Yoshihara.

\newblock Quotients of abelian surfaces.

\newblock \emph{Publications of the Research Institute for Mathematical
  Sciences}, 31\penalty0 (1):\penalty0 135--143, 1995.

\newblock URL
  \url{https://www.jstage.jst.go.jp/article/kyotoms1969/31/1/31_1_135/_pdf}.
\newblock \emph{Verified}.

\bibitem[van Straten(1987)]{vanStraten1987}
Duco van Straten.

\newblock \emph{Weakly Normal Surface Singularities and Their Improvements}.

\newblock PhD thesis, Rijksuniversiteit Leiden, 1987.

\newblock URL
  \url{https://download.uni-mainz.de/mathematik/Algebraische%20Geometrie/DvS-Publikationen/DvS_1987_weakly%20normal%20Surface%20Singularities%20and%20their%20improvements.pdf}.
\newblock \emph{Verified}.

\bibitem[Hassett(2009)]{Hassett2009Rational}
B.~Hassett, \emph{Rational surfaces over nonclosed fields},
in \emph{Arithmetic Geometry}, H.~Darmon, D.~Ellwood, B.~Hassett and Y.~Tschinkel (eds.),
Clay Mathematics Proceedings \textbf{8}, American Mathematical Society and Clay Mathematics Institute, 2009, pp.~155--209.
\url{https://www.claymath.org/wp-content/uploads/2022/03/cmip08c.pdf}.

\bibitem[Artebani and Dolgachev(2009)]{ArtebaniDolgachev2009}
Michela Artebani and Igor~V. Dolgachev.
\newblock The Hesse pencil of plane cubic curves.
\newblock \emph{L'Enseignement Math\'ematique}, 55(3/4):235--273, 2009.
\newblock \doi{10.4171/LEM/55-3-3}.
\newblock URL \url{https://ems.press/content/serial-article-files/44190?nt=1}.

\bibitem[Gim\'enez Conejero and Pint\'er(2024)]{GCP2024}
Roberto Gim\'enez Conejero and Gerg\H{o} Pint\'er.
\newblock Signature of the Milnor fiber of parametrized surfaces.
\newblock arXiv:2401.14951v2, 30 January 2024.
\newblock URL \url{https://arxiv.org/abs/2401.14951v2}.

\bibitem[Beauville(1978)]{Beauville1978}
Arnaud Beauville.
\newblock \emph{Surfaces alg\'ebriques complexes}.
\newblock Ast\'erisque, volume 54. Soci\'et\'e math\'ematique de France, 1978.
\newblock URL \url{https://www.numdam.org/item/AST_1978__54__R1_0/}.

\end{thebibliography}

\end{document}
